% Punjabi Shahmukhi edition of The Open Logic Text, CC BY 4.0.
% Assemble with tools/build_cumulative_reader.py.
\PassOptionsToPackage{bidi=default}{babel}
\documentclass[12pt,a4paper,oneside,openany]{memoir}
\errorcontextlines=60
\usepackage[margin=23mm]{geometry}
\usepackage{fontspec}
\newcommand*{\olpath}{upstream}
\input{\olpath/sty/open-logic.sty}
\babelprovide[import]{urdu}
\babelfont{rm}[Path=fonts/,BoldFont=NotoSerif-Bold.ttf,
  BoldFeatures={Ligatures=TeXOff},ItalicFont=NotoSerif-Italic.ttf,
  BoldItalicFont=NotoSerif-BoldItalic.ttf]{NotoSerif-Regular.ttf}
\babelfont[urdu]{rm}[Path=fonts/,Script=Arabic,
  BoldFont=NotoNastaliqUrdu-Regular.ttf,BoldFeatures={FakeBold=1.5}]{NotoNastaliqUrdu-Regular.ttf}
\linespread{1.65}
% Punjabi Shahmukhi reader labels and retained OpenLogic tokens.
\ProvidesPackage{pnb-morphology}
\newcommand{\pnboblique}[1]{\begingroup
  \definetoken{p}{element}{عنصراں}%
  #1\endgroup}
% Load after original OpenLogic token definitions.
\settexttoken{element}{عنصر}{عنصر}[عنصر][عنصر]
\definetoken{a}{element}{}
\definetoken{A}{element}{}
\settexttoken{formula}{فارمولا}{فارمولے}[فارمولا][فارمولے]
\definetoken{a}{formula}{}
\definetoken{A}{formula}{}
\settexttoken{sentence}{جملہ}{جملے}[جملہ][جملے]
\definetoken{a}{sentence}{}
\definetoken{A}{sentence}{}
\settexttoken{constant}{ثابتہ نشان}{ثابتہ نشان}[ثابتہ نشان][ثابتہ نشان]
\definetoken{a}{constant}{}
\definetoken{A}{constant}{}
\settexttoken{function}{فنکشن نشان}{فنکشن نشان}[فنکشن نشان][فنکشن نشان]
\definetoken{a}{function}{}
\definetoken{A}{function}{}
\settexttoken{predicate}{محمول نشان}{محمول نشان}[محمول نشان][محمول نشان]
\definetoken{a}{predicate}{}
\definetoken{A}{predicate}{}
\settexttoken{derivation}{اخذ}{اخذ}[اخذ][اخذ]
\definetoken{a}{derivation}{}
\definetoken{A}{derivation}{}
\settexttoken{derive}{اخذ}{اخذ}[اخذ][اخذ]
\settexttoken{free for}{بدل لئی آزاد}{بدل لئی آزاد}[بدل لئی آزاد][بدل لئی آزاد]
\settexttoken{discharge}{خارج}{خارج}[خارج][خارج]
\settexttoken{discharged}{خارج}{خارج}[خارج][خارج]
\settexttoken{undischarged}{غیر خارج}{غیر خارج}[غیر خارج][غیر خارج]
\definetoken{a}{undischarged}{}
\definetoken{A}{undischarged}{}
\settexttoken{identity}{برابری دی علامت}{برابری دیاں علامتاں}[برابری دی علامت][برابری دیاں علامتاں]
\definetoken{a}{identity}{}
\definetoken{A}{identity}{}
\settexttoken{structure}{ساخت}{ساختاں}[ساخت][ساختاں]
\definetoken{a}{structure}{}
\definetoken{A}{structure}{}
\settexttoken{variable}{متغیر}{متغیر}[متغیر][متغیر]
\definetoken{a}{variable}{}
\definetoken{A}{variable}{}
\settexttoken{domain}{دائرہ}{دائرے}[دائرہ][دائرے]
\definetoken{a}{domain}{}
\definetoken{A}{domain}{}
\settexttoken{complete}{کامل}{کامل}[کامل][کامل]
\definetoken{a}{complete}{}
\definetoken{A}{complete}{}
\settexttoken{decidable}{فیصلہ کرن جوگ}{فیصلہ کرن جوگ}[فیصلہ کرن جوگ][فیصلہ کرن جوگ]
\definetoken{a}{decidable}{}
\definetoken{A}{decidable}{}
\settexttoken{axiomatizable}
  {مسلّماتی طور متعین ہون جوگ}{مسلّماتی طور متعین ہون جوگ}
  [مسلّماتی طور متعین ہون جوگ][مسلّماتی طور متعین ہون جوگ]
\definetoken{a}{axiomatizable}{}
\definetoken{A}{axiomatizable}{}
\settexttoken{axiomatizability}
  {مسلّماتی طور متعین ہون دی اہلیت}{مسلّماتی طور متعین ہون دی اہلیت}
  [مسلّماتی طور متعین ہون دی اہلیت][مسلّماتی طور متعین ہون دی اہلیت]
\definetoken{a}{axiomatizability}{}
\definetoken{A}{axiomatizability}{}
\settexttoken{represents}{نمائندگی}{نمائندگی}[نمائندگی][نمائندگی]
\definetoken{a}{represents}{}
\definetoken{A}{represents}{}
\settexttoken{enumerable}
  {گنتی وار فہرست بنن جوگ}{گنتی وار فہرست بنن جوگ}
  [گنتی وار فہرست بنن جوگ][گنتی وار فہرست بنن جوگ]
\settexttoken{nonenumerable}
  {گنتی وار فہرست نہ بنن جوگ}{گنتی وار فہرست نہ بنن جوگ}
  [گنتی وار فہرست نہ بنن جوگ][گنتی وار فہرست نہ بنن جوگ]
\settexttoken{denumerable}
  {گنتی وار لامتناہی}{گنتی وار لامتناہی}
  [گنتی وار لامتناہی][گنتی وار لامتناہی]
\settexttoken{surjective}
  {ہر ہدف قدر لین والا}{ہر ہدف قدر لین والے}
  [ہر ہدف قدر لین والا][ہر ہدف قدر لین والے]
\definetoken{a}{surjective}{}
\definetoken{A}{surjective}{}
\settexttoken{computably enumerable}
  {قابلِ حساب گنتی وار فہرست بنن جوگ}
  {قابلِ حساب گنتی وار فہرست بنن جوگ}
  [قابلِ حساب گنتی وار فہرست بنن جوگ]
  [قابلِ حساب گنتی وار فہرست بنن جوگ]
\definetoken{a}{computably enumerable}{}
\definetoken{A}{computably enumerable}{}

% Preserve logical Unicode text behind the shaped Shahmukhi glyphs in XeTeX PDFs.
\ifdefined\XeTeXgenerateactualtext
  \XeTeXgenerateactualtext=1\relax
\fi

% Apply after OpenLogic loads the TikZ automata library.
\AtBeginDocument{%
  \ifdefined\tikzset
    \pgfkeysifdefined{/tikz/initial by arrow/.@cmd}{%
      \tikzset{initial by arrow/.append style={initial text={شروع}}}%
    }{}%
  \fi
}
\settexttoken{parameter}{پیرامیٹر}{پیرامیٹر}[پیرامیٹر][پیرامیٹر]
\settexttoken{lambda define}
  {لامبڈا نال متعین}{لامبڈا نال متعین کردا اے}
  [لامبڈا نال متعین][لامبڈا نال متعین کردا اے]
\settexttoken{lambda defined}
  {متعین}{متعین}[متعین][متعین]
\settexttoken{lambda definable}
  {لامبڈا نال متعین ہون جوگ}{لامبڈا نال متعین ہون جوگ}
  [لامبڈا نال متعین ہون جوگ][لامبڈا نال متعین ہون جوگ]
\definetoken{a}{lambda definable}{}
\definetoken{A}{lambda definable}{}
\settexttoken{propositional variable}
  {قضیاتی متغیر}{قضیاتی متغیراں}[قضیاتی متغیر][قضیاتی متغیراں]
\definetoken{a}{propositional variable}{}
\definetoken{A}{propositional variable}{}
\settexttoken{valuation}{قدر بندی}{قدر بندیاں}[قدر بندی][قدر بندیاں]
\definetoken{a}{valuation}{}
\definetoken{A}{valuation}{}
\settexttoken{truth}{صدق}{صدق}[صدق][صدق]
\settexttoken{falsity}{کذب}{کذب}[کذب][کذب]
\settexttoken{conditional}{شرطیہ}{شرطیے}[شرطیہ][شرطیے]
\settexttoken{biconditional}{دو شرطیہ}{دو شرطیے}[دو شرطیہ][دو شرطیے]
\definetoken{a}{truth}{}
\definetoken{A}{truth}{}
\definetoken{a}{falsity}{}
\definetoken{A}{falsity}{}
\definetoken{a}{conditional}{}
\definetoken{A}{conditional}{}
\definetoken{a}{biconditional}{}
\definetoken{A}{biconditional}{}
\settexttoken{derivable}
  {اخذ ہون جوگ}{اخذ ہون جوگ}[اخذ ہون جوگ][اخذ ہون جوگ]
\settexttoken{derivability}
  {اخذ ہون دی اہلیت}{اخذ ہون دی اہلیت}[اخذ ہون دی اہلیت][اخذ ہون دی اہلیت]
\definetoken{a}{derivable}{}
\definetoken{A}{derivable}{}
\definetoken{a}{derivability}{}
\definetoken{A}{derivability}{}
\settexttoken{language}{زبان}{زباناں}[زبان][زباناں]
\definetoken{a}{language}{}
\definetoken{A}{language}{}

% Reader text tokens used in the translated chapters.
\settexttoken{tableau}{تابلو}{تابلو}[تابلو][تابلو]
\definetoken{a}{tableau}{}
\definetoken{A}{tableau}{}
\settexttoken{signed formula}{علامت لایا فارمولا}{علامت لائے فارمولے}[علامت لایا فارمولا][علامت لائے فارمولے]
\definetoken{a}{signed formula}{}
\definetoken{A}{signed formula}{}
\settexttoken{operator}{منطقی عامل}{منطقی عاملاں}[منطقی عامل][منطقی عاملاں]
\definetoken{a}{operator}{}
\definetoken{A}{operator}{}
\settexttoken{main operator}{مرکزی عامل}{مرکزی عامل}[مرکزی عامل][مرکزی عامل]
\definetoken{a}{main operator}{}
\definetoken{A}{main operator}{}
\settexttoken{subformula}{ذیلی فارمولا}{ذیلی فارمولے}[ذیلی فارمولا][ذیلی فارمولے]
\definetoken{a}{subformula}{}
\definetoken{A}{subformula}{}
\settexttoken{value}{قدر}{قدراں}[قدر][قدراں]
\definetoken{a}{value}{}
\definetoken{A}{value}{}
\settexttoken{injective}{انجیکٹو}{انجیکٹو}[انجیکٹو][انجیکٹو]
\definetoken{a}{injective}{}
\definetoken{A}{injective}{}
\settexttoken{injection}{اک بہ اک نقشہ بندی}{اک بہ اک نقشہ بندیاں}[اک بہ اک نقشہ بندی][اک بہ اک نقشہ بندیاں]
\definetoken{a}{injection}{}
\definetoken{A}{injection}{}
\settexttoken{surjection}{ہر ہدف قدر لین والی نقشہ بندی}{ہر ہدف قدر لین والیاں نقشہ بندیاں}[ہر ہدف قدر لین والی نقشہ بندی][ہر ہدف قدر لین والیاں نقشہ بندیاں]
\definetoken{a}{surjection}{}
\definetoken{A}{surjection}{}
\settexttoken{bijective}{بائجیکٹو}{بائجیکٹو}[بائجیکٹو][بائجیکٹو]
\definetoken{a}{bijective}{}
\definetoken{A}{bijective}{}
\settexttoken{bijection}{دوطرفی نقشہ بندی}{دوطرفی نقشہ بندیاں}[دوطرفی نقشہ بندی][دوطرفی نقشہ بندیاں]
\definetoken{a}{bijection}{}
\definetoken{A}{bijection}{}
\settexttoken{c.e.}{قابلِ حساب گنتی وار فہرست بنن جوگ}{قابلِ حساب گنتی وار فہرست بنن جوگ}[قابلِ حساب گنتی وار فہرست بنن جوگ][قابلِ حساب گنتی وار فہرست بنن جوگ]
\definetoken{a}{c.e.}{}
\definetoken{A}{c.e.}{}
\settexttoken{nonderivability}{اخذ نہ ہو سکّن}{اخذ نہ ہو سکّن}[اخذ نہ ہو سکّن][اخذ نہ ہو سکّن]
\definetoken{a}{nonderivability}{}
\definetoken{A}{nonderivability}{}
\settexttoken{axiomatized}{مسلّمیاں نال متعین}{مسلّمیاں نال متعین}[مسلّمیاں نال متعین][مسلّمیاں نال متعین]
\definetoken{a}{axiomatized}{}
\definetoken{A}{axiomatized}{}


% The Urdu Babel locale supplies Arabic-script shaping and paragraph direction.
% Names, prose, and the edition's language metadata are Punjabi Shahmukhi.
\setlocalecaption{urdu}{contents}{فہرست}
\setlocalecaption{urdu}{chapter}{باب}
\setlocalecaption{urdu}{part}{حصہ}
\setlocalecaption{urdu}{figure}{شکل}
\setlocalecaption{urdu}{table}{جدول}
\setlocalecaption{urdu}{proof}{ثبوت}
\setlocalecaption{urdu}{history}{تاریخی نوٹ}
\setlocalecaption{urdu}{reading}{ہور پڑھائی}
\setlocalecaption{urdu}{appendix}{ضمیمہ}
\setlocalecaption{urdu}{appendices}{ضمیمے}
\setlocalecaption{urdu}{theorem}{قضیہ}
\setlocalecaption{urdu}{example}{مثال}
\setlocalecaption{urdu}{definition}{تعریف}
\setlocalecaption{urdu}{lemma}{لما}
\setlocalecaption{urdu}{proposition}{نتیجہ}
\setlocalecaption{urdu}{corollary}{لازمی نتیجہ}
\setlocalecaption{urdu}{problem}{مشق}
\setlocalecaption{urdu}{problems}{مشقاں}
\setlocalecaption{urdu}{remark}{تبصرہ}
\setlocalecaption{urdu}{axiom}{مسلّمہ}
\setlocalecaption{urdu}{note}{نوٹ}
\setlocalecaption{urdu}{case}{صورت}
\setlocalecaption{urdu}{convention}{رواج}
\setlocalecaption{urdu}{OLP}{اوپن لاجک پروجیکٹ}
\setlocalecaption{urdu}{OLT}{اوپن لاجک متن}
\setlocalecaption{urdu}{OLTlink}{https://github.com/OpenLogicProject/OpenLogic}
\setlocalecaption{urdu}{license}{کریئیٹو کامنز انتساب ۴٫۰ بین الاقوامی لائسنس}

\newcommand*{\pnbcrefname}[3]{\crefname{#1}{#2}{#3}\Crefname{#1}{#2}{#3}}
\addto\extrasurdu{%
  \pnbcrefname{part}{حصہ}{حصے}%
  \pnbcrefname{chapter}{باب}{باب}%
  \pnbcrefname{section}{حصہ}{حصے}%
  \pnbcrefname{subsection}{ذیلی حصہ}{ذیلی حصے}%
  \pnbcrefname{page}{صفحہ}{صفحے}%
  \pnbcrefname{thm}{قضیہ}{قضیے}%
  \pnbcrefname{ex}{مثال}{مثالاں}%
  \pnbcrefname{defn}{تعریف}{تعریفاں}%
  \pnbcrefname{lem}{لما}{لمے}%
  \pnbcrefname{prop}{نتیجہ}{نتیجے}%
  \pnbcrefname{cor}{لازمی نتیجہ}{لازمی نتیجے}%
  \pnbcrefname{prob}{مشق}{مشقاں}%
  \pnbcrefname{probd}{مشق}{مشقاں}%
  \pnbcrefname{rem}{تبصرہ}{تبصرے}%
  \pnbcrefname{axiom}{مسلّمہ}{مسلّمے}%
  \pnbcrefname{note}{نوٹ}{نوٹ}%
  \pnbcrefname{case}{صورت}{صورتاں}%
  \pnbcrefname{conv}{رواج}{رواج}%
  \pnbcrefname{figure}{شکل}{شکلاں}%
  \pnbcrefname{table}{جدول}{جدولاں}%
  \pnbcrefname{equation}{مساوات}{مساواتاں}%
  \def\crefpairconjunction{ تے }%
  \def\crefmiddleconjunction{، }%
  \def\creflastconjunction{ تے }%
}
\DeclareRobustCommand{\pnbLatin}[1]{\foreignlanguage{english}{#1}}
\DeclareTextFontCommand{\pnbemphasis}{\upshape\bfseries}
\let\emph\pnbemphasis
\RenewDocumentCommand{\indcase}{s t{!} m m +m}{%
  \DeclareDocumentMacro\indfrm{#3}%
  \DeclareDocumentMacro\indfrmp{#3}%
  \DeclareDocumentMacro\indcomplex{#4}%
  \IfBooleanTF{#1}{$#3$ جوہری اے۔ }{$#3 \ident #4$۔ }%
  \IfBooleanTF{#2}{مشق۔}{#5}%
}
\pdfstringdefDisableCommands{%
  \def\pnbLatin#1{#1}%
  \def\foreignlanguage#1#2{#2}%
  \def\Th#1{#1}\def\Log#1{#1}%
}

% Keep dvipdfmx content delimiters in left-to-right order inside each picture.
\makeatletter
\let\PNBpgfBeginPicture\pgfsys@beginpicture
\let\PNBpgfEndPicture\pgfsys@endpicture
\def\pgfsys@beginpicture{\beginL\PNBpgfBeginPicture}
\def\pgfsys@endpicture{\PNBpgfEndPicture\endL}
\NewDocumentEnvironment{pnbDiagramLine}{+b}
  {\par\noindent\beginL\makebox[\linewidth][c]{\ignorespaces#1\unskip}\endL\par}
  {}
\AtBeginDocument{%
  \@ifundefined{ifpnbDiagramReader}{%
    \@ifundefined{tikz@initial@text}{}{%
      \def\tikz@initial@text{\foreignlanguage{urdu}{شروع}}%
    }%
  }{}%
}

% Keep propositional references to shared material on the corresponding
% first-order label. The actual label and reference type come from cleveref.
\newcommand*{\pnbEnableSharedLabels}{%
  \let\pnbOriginalLabel\label
  \renewcommand*{\label}[1]{%
    \pnbOriginalLabel{##1}%
    \begingroup
    \edef\pnbCurrentLabel{##1}%
    \expandafter\pnbSharedLabel\pnbCurrentLabel::\@nil
    \endgroup
  }%
}
\def\pnbSharedLabel#1:#2:#3\@nil{%
  \ifstrequal{#1}{fol}{%
    \ifstrempty{#2}{}{%
      \ifinlist{#2}{\pnbSharedChapters}{%
        \expandafter\pnbWriteSharedLabel\pnbCurrentLabel\@nil
      }{}%
    }%
  }{}%
}
\def\pnbWriteSharedLabel#1:#2\@nil{\pnbOriginalLabel{pl:#2}}
\newcommand*{\pnbSharedChapters}{}
\forcsvlist{\listadd\pnbSharedChapters}{prf,seq,ntd,tab,axd,com}
\makeatother

% Localized text tokens used by the chapter bodies.
\settexttoken{proof}{ثبوت}{ثبوت}[ثبوت][ثبوت]
\definetoken{a}{proof}{}
\definetoken{A}{proof}{}
\settexttoken{height}{اُچائی}{اُچائیاں}[اُچائی][اُچائیاں]
\definetoken{a}{height}{}
\definetoken{A}{height}{}
\settexttoken{prove}{ثابت}{ثابت}[ثابت][ثابت]
\definetoken{a}{prove}{}
\definetoken{A}{prove}{}
\settexttoken{proving}{ثابت کرن}{ثابت کرن}[ثابت کرن][ثابت کرن]
\definetoken{a}{proving}{}
\definetoken{A}{proving}{}
\settexttoken{provable}{ثابت ہو سکّن جوگ}{ثابت ہو سکّن جوگ}[ثابت ہو سکّن جوگ][ثابت ہو سکّن جوگ]
\definetoken{a}{provable}{}
\definetoken{A}{provable}{}
\settexttoken{introduction}{شامل کرن دا}{شامل کرن دے}[شامل کرن دا][شامل کرن دے]
\definetoken{a}{introduction}{}
\definetoken{A}{introduction}{}
\settexttoken{elimination}{ہٹاؤن دا}{ہٹاؤن دے}[ہٹاؤن دا][ہٹاؤن دے]
\definetoken{a}{elimination}{}
\definetoken{A}{elimination}{}
\settexttoken{relational model}{تعلقاں والا نمونہ}{تعلقاں والے نمونے}[تعلقاں والا نمونہ][تعلقاں والے نمونے]
\definetoken{a}{relational model}{}
\definetoken{A}{relational model}{}
\settexttoken{depth}{ڈونگھائی}{ڈونگھائیاں}[ڈونگھائی][ڈونگھائیاں]
\definetoken{a}{depth}{}
\definetoken{A}{depth}{}
\settexttoken{sentences}{جملے}{جملے}[جملے][جملے]
\definetoken{a}{sentences}{}
\definetoken{A}{sentences}{}
\settexttoken{colorC}{لال}{لال}[لال][لال]
\definetoken{a}{colorC}{}
\definetoken{A}{colorC}{}
\settexttoken{colorD}{نیلا}{نیلا}[نیلا][نیلا]
\definetoken{a}{colorD}{}
\definetoken{A}{colorD}{}
\settexttoken{colorE}{ہرا}{ہرا}[ہرا][ہرا]
\definetoken{a}{colorE}{}
\definetoken{A}{colorE}{}
\settexttoken{hp}{hp}{hp}[hp][hp]
\definetoken{a}{hp}{}
\definetoken{A}{hp}{}

\hypersetup{hidelinks,unicode,pdflang={pnb-Arab-PK},
  pdftitle={اوپن لاجک متن — پنجابی شاہ مکھی},
  pdfauthor={Open Logic Project},
  pdfsubject={The Open Logic Text, Punjabi Shahmukhi AI translation},
  pdfkeywords={OpenAI Codex — GPT-5.6 Sol (Ultra), OpenAI Codex — GPT-6 Sol (Ultra), OpenAI Codex — GPT-6.1 Sol (Ultra)}}
\setlength{\parindent}{0pt}
\setlength{\parskip}{0.4em}
\emergencystretch=3em
\let\cleardoublepage\clearpage
\pagestyle{plain}

\begin{document}
\selectlanguage{urdu}
\pnbEnableSharedLabels
\XeTeXgenerateactualtext=0\relax
\frontmatter
\begin{titlingpage}
\centering
\vspace*{0.15\textheight}
{\Huge\bfseries اوپن لاجک متن\par}
\vspace{1em}
{\Large پنجابی شاہ مکھی ترجمہ\par}
\vspace{2em}
{\large \href{https://openlogicproject.org/people/}{اوپن لاجک پروجیکٹ}\par}
\vfill
ایہ اوپن لاجک پروجیکٹ دے \textit{اوپن لاجک متن} دا اے آئی نال کیتا
گیا پنجابی شاہ مکھی ترجمہ اے۔ ریاضی دیاں کجھ اصطلاحاں ہجے عارضی نیں۔
\par
{\small ترجمے دی مُڈھلی لکھائی — \pnbLatin{OpenAI Codex — GPT-5.6 Sol (Ultra)}۔}\par
{\small ترجمے دے اگلے حصیاں دی لکھائی — \pnbLatin{OpenAI Codex — GPT-6 Sol (Ultra)}۔}\par
{\small باقی ترجمہ، ترجمے دیاں اصلاحاں تے ایڈیشن دی صفحہ بندی — \pnbLatin{OpenAI Codex — GPT-6.1 Sol (Ultra)}۔}
\par
\href{https://creativecommons.org/licenses/by/4.0/}
  {کریئیٹو کامنز انتساب ۴٫۰ بین الاقوامی لائسنس}
\par
\href{https://github.com/OpenLogicProject/OpenLogic/tree/9620cc73f9c8e0ad003c514a5d3748f29611c4c0}
  {انگریزی ماخذ}
\quad
\href{https://github.com/KokunoYumeto/OpenLogic-translations}{ہور ترجمے}
\end{titlingpage}
% OLP-0001 content/open-logic-about.tex
\chapter*{اوپن لاجک پروجیکٹ بارے}
\addcontentsline{toc}{chapter}{اوپن لاجک پروجیکٹ بارے}


\textit{اوپن لاجک متن} رسمی ماورائے منطق \pnbLatin{(formal meta-logic)} تے رسمی طریقہ کار
\pnbLatin{(formal methods)} دی اک کھلے ماخذ والی، سانجھے کم نال تیار کیتی درسی کتاب اے۔ ایہ درمیانی درجے توں شروع
ہوندی اے (یعنی رسمی منطق دے تعارفی کورس توں پچھوں)۔ بھاویں ایہ ریاضی
دے ماہر نہ ہون والے پڑھنے والیاں لئی اے (خاص طور تے فلسفے تے کمپیوٹر
سائنس دے طالب علماں لئی)، پر ایہ منطقی سختی نال لکھی گئی اے۔

ہُن شامل کجھ موضوعاں دا بیان ہالے پورا نہیں ہویا، تے کئی حصیاں نوں ہالے
وی وڈی نظرثانی دی لوڑ اے۔ اسی اگے چل کے ہور موضوع شامل کرن لئی متن
نوں ودھاؤن دا ارادہ رکھدے آں۔ اسی متن وچ اصطلاحاں دی فرہنگ، ہور پڑھائی
دی فہرست، تاریخی نوٹ، تصویراں، بہتر وضاحتاں، نتائج دی فلسفے، کمپیوٹر
سائنس تے ریاضی لئی اہمیت دسن والے حصے، تے ہور مشقاں تے مثالاں وی
شامل کرنا چاہندے آں۔ جے تہانوں کوئی غلطی لبھے یا تہاڈے کول کوئی صلاح ہووے،
تاں \href{https://github.com/OpenLogicProject/OpenLogic/wiki/Contributing}{مہربانی کر کے پروجیکٹ دی ٹیم نوں دسو}۔

ایہ پروجیکٹ کھلے ماخذ دے جذبے نال کم کردا اے۔ متن صرف مفت دستیاب ای
نہیں، بلکہ اسی LaTeX دا ماخذ وی کریئیٹو کامنز انتساب لائسنس دے تحت
دیندے آں۔ ایہ لائسنس ہر کسے نوں ساڈا کم ڈاؤن لوڈ کرن، ورتن، بدلّن،
نویں ترتیب دین، ہور صورت وچ لے جان تے اگے ونڈن دا حق دیندا اے، بشرطیکہ
اوہ مناسب انتساب کرے۔ ہور معلومات لئی اوپن لاجک پروجیکٹ دی ویب سائٹ
\href{http://openlogicproject.org/}{openlogicproject.org} ویکھو۔
% /OLP-0001
\setpnumwidth{3em}
\setrmarg{4em}
\tableofcontents*
\mainmatter
% OLP-0002 content/content.tex

\clearpage

\begin{editorial}
ایہ فائل اوپن لاجک پروجیکٹ وچ شامل سارا مواد لوڈ کردی اے۔ ایہو جہے
ادارتی نوٹ، جے دکھائی دین، تاں ایہ دس دے نیں کہ فائل نوں ایس گل دا
خیال کیتے بغیر کمپائل کیتا گیا اے کہ مواد کِویں پیش ہونا چاہیدا اے۔
جے تسیں ایہ پڑھ سکدے او، تاں غالباً ایس پی ڈی ایف توں پڑھانا یا پڑھنا
\emph{مناسب نہیں}۔

اوپن لاجک پروجیکٹ کئی طریقے دیندا اے جنہاں نال پڑھاؤن یا اپنے آپ
پڑھن لئی زیادہ مناسب متن تیار کیتا جا سکدا اے۔ مثال لئی، پہلے توں طے
ترتیب وچ متن سارے منطقی عامل \pnbLatin{(logical operators)} بنیادی ماندا اے تے ثبوتاں وچ ہر عامل
دیاں ساریاں صورتاں نوں حل کردا اے۔ پر ایہناں وچوں کجھ صورتاں نوں مشقاں
دے طور تے چھڈنا زیادہ چنگا اے۔ اوپن لاجک پروجیکٹ اُتے ہالے کم وی جاری
اے۔ سانجھے کم تے بہتری نوں ہلا شیری دین لئی مواد اوہدوں وی شامل کیتا
جاندا اے جدوں اوہ صرف مسودہ ہووے، اوہدیاں مشقاں نہ ہون، وغیرہ۔ کِسے
کورس لئی تیار کیتی پی ڈی ایف وچ ایہ حصے شامل نہیں ہون گے۔

پڑھاؤن تے پڑھائی لئی زیادہ مناسب پی ڈی ایف لبھن لئی
\href{http://builds.openlogicproject.org/}{اوپن لاجک دی ویب سائٹ اُتے دستیاب نمونے دے کورس}
ویکھو۔ اپنا متن تیار کرن لئی تسیں
\href{https://github.com/OpenLogicProject/OpenLogic/tree/master/courses/sample}{نمونے دی ڈرائیور فائل \pnbLatin{(sample driver file)}}
توں شروع کر سکدے او، یا ہور نفیس تے اگلی سطح دیاں مثالاں لئی ایس متن
توں بنیاں درسی کتاباں دے ماخذ ویکھ سکدے او۔
\end{editorial}

% OLP-0003 content/sets-functions-relations/sets-functions-relations-complete.tex
\begingroup
\def\olfilename{sets-functions-relations-complete}
\expandafter\def\csname ol@sectioncs\endcsname{section}

\olpart{sfr}{سادہ سیٹ تھیوری \pnbLatin{(Naive Set Theory)}}

\begin{editorial}
  ایس حصے دا مواد سادہ سیٹ تھیوری دے بنیادی خیالاں دا تعارف اے۔ ٹِم بٹن
  دی اوپن سیٹ تھیوری شامل ہون نال، ایس وچ عددی نظاماں دی تعمیر تے
  لامتناہیت بارے بحث وی آ گئی اے۔ ایہ موضوع اوپن لاجک پروجیکٹ دے منطقی
  حصیاں لئی ضروری نہیں۔
\end{editorial}

% OLP-0004 content/sets-functions-relations/sets/sets.tex
\begingroup
\def\olfilename{sets}
\expandafter\def\csname ol@sectioncs\endcsname{section}

\olchapter{sfr}{set}{سیٹ}

% OLP-0005 content/sets-functions-relations/sets/basics.tex
\begingroup
\def\olfilename{basics}
\expandafter\def\csname ol@sectioncs\endcsname{section}

\olfileid{sfr}{set}{bas}
\olsection{عنصراں نال تعیّن دا اصول \pnbLatin{(Extensionality)}}

\emph{سیٹ} شےواں دا اک اکٹھ اے، جہنوں {اکو} شے منّیا جاندا اے۔ سیٹ بناؤن والیاں
شےواں نوں اوس سیٹ دے \emph{عنصر} یا \emph{رکن} (element/member) آکھیا جاندا
اے۔ جے $x$ سیٹ~$A$ دا
!!a{element} ہووے، تاں اسی $x \in A$ لکھدے آں؛ جے نہ ہووے، تاں $x \notin A$
لکھدے آں۔ جس سیٹ دے کوئی !!{element}s نہ ہون، اوہنوں \emph{خالی} سیٹ
\pnbLatin{(empty set)} آکھدے نیں
تے اوہنوں~``$\emptyset$'' نال ظاہر کردے نیں۔

\begin{explain}
ایس گل نال کوئی فرق نہیں پیندا کہ اسی سیٹ نوں کِویں \emph{دسدے} آں، اوہدے
\pnboblique{!!{element}s} نوں کِہڑی \emph{ترتیب} وچ رکھدے آں، یا اوہدے \pnboblique{!!{element}s} نوں
\emph{کِنّی واری} گِندے آں۔ اصل گل صرف ایہ اے کہ اوہدے !!{element}s کِہڑے نیں۔
اسی ایس گل نوں ہِیٹھلے اصول دی صورت وچ بیان کردے آں۔
\end{explain}

\begin{defn}[عنصراں نال تعیّن دا اصول]
  جے $A$ تے $B$ سیٹ نیں، تاں $A = B$ اُدوں تے صرف اُدوں
  \pnbLatin{(if and only if)} ہوندا اے جدوں
  $A$ دا ہر !!{element}، $B$ دا وی !!a{element} ہووے، تے ایس توں اُلٹ وی۔
\end{defn}

عنصراں نال تعیّن دا اصول کجھ علامتاں ورتن دا جواز دیندا اے۔ عام طور تے، جدوں
ساڈے کول کجھ شےواں $a_{1}$، \dots، $a_{n}$ ہون، تاں $\{a_{1}, \dots, a_{n}\}$
\emph{اوہی} سیٹ اے جہدے !!{element}s، $a_1, \ldots, a_n$ نیں۔ اسی لفظ
``\emph{اوہی}'' اُتے زور دیندے آں، کیوں جو عنصراں نال تعیّن دا اصول دسدا اے کہ
ایہو جہا سیٹ صرف \emph{اک} ہو سکدا اے۔ ایس اصول توں ایہ وی جائز اے:
  \[
    \{a, a, b\} = \{a, b\} = \{b,a\}.
  \]
ایس نال اوہ گل واضح ہوندی اے کہ سیٹاں بارے سوچدیاں سانوں ایہدی پرواہ نہیں
ہوندی کہ اوہناں دے !!{element}s دی ترتیب کی اے، یا اوہ کِنّی واری دسے گئے نیں۔

\begin{tagblock}{novice}
\begin{ex}
جدوں وی تہاڈے کول کجھ شےواں ہون، تسیں اوہناں نوں اک سیٹ وچ اکٹھیاں کر سکدے او۔
مثال لئی، رچرڈ دے بھین بھراواں دا سیٹ اک ایہو جہا سیٹ اے جہدے وچ اک بندہ اے،
تے اسی اوہنوں $S=\{\textrm{Ruth}\}$ لکھ سکدے آں۔ $4$ توں گھٹ مثبت صحیح
اعداد \pnbLatin{(positive integers)} دا سیٹ $\{1, 2, 3\}$ اے، پر اوہنوں $\{3, 2, 1\}$ یا
$\{1, 2, 1, 2, 3\}$ وی لکھیا جا سکدا اے۔ عنصراں نال تعیّن دے اصول موجب ایہ
سارے اکو سیٹ نیں۔ کیوں جو $\{1, 2, 3\}$ دا ہر !!{element}،
$\{3, 2, 1\}$ دا (تے $\{1, 2, 1, 2, 3\}$ دا وی) !!a{element} اے، تے ایس
توں اُلٹ وی۔
\end{ex}
\end{tagblock}

اکثر اسی سیٹ نوں اوہدے !!{element}s دی سانجھی خاصیت نال دسّاں گے۔ ایس لئی
اسی ایہ مختصر علامتی لکھت \pnbLatin{(set-builder notation)} ورتاں گے:
$\Setabs{x}{\phi(x)}$، جتھے $\phi(x)$
اوہ خاصیت اے جہڑی $x$ وچ ہونی چاہیدی اے تاں جو اوہ سیٹ دے \pnboblique{!!{element}s} وچوں
اک گنیا جا سکے۔

\begin{tagblock}{novice}
\begin{ex}
اپنی مثال وچ اسی $S$ نوں ایویں وی دس سکدے ساں:
\[
S = \Setabs{x}{x \text{ رچرڈ دی بھین یا اوہدا بھرا اے}}.
\]
\end{ex}
\end{tagblock}

\begin{tagblock}{math}
\begin{ex}
کِسے عدد نوں \emph{کامل} \pnbLatin{(perfect number)} اُدوں تے صرف اُدوں آکھدے نیں جدوں
اوہ اپنے اوہناں ونڈن والے اعداد \pnbLatin{(proper divisors)} دے جوڑ دے برابر ہووے جو
عدد نوں پورا ونڈدے ہون پر آپ اوہدے
برابر نہ ہون۔ مثال لئی، $6$ کامل اے، کیوں جو اوہنوں ونڈن والے اعداد، جو آپ
اوہدے برابر نہیں، $1$، $2$ تے~$3$ نیں، تے $6 = 1 + 2 + 3$ اے۔ دراصل،
$6$، $10$ توں گھٹ اکو اک مثبت صحیح عدد اے جو کامل اے۔ سو، عنصراں نال تعیّن
دا اصول ورت کے اسی آکھ سکدے آں:
\[
	\{6\} = \Setabs{x}{x\text{ کامل اے تے }0 \leq x \leq 10}
\]
سجے پاسے دی علامتی لکھت نوں اسی ایویں پڑھدے آں: ``اوہناں $x$ دا سیٹ جیہناں لئی
$x$ کامل اے تے $0 \leq x \leq 10$''۔ ایتھے برابری دی ایہ لکھت تصدیق کردی
اے کہ سیٹاں بارے سوچدیاں سانوں ایہدی پرواہ نہیں ہوندی کہ اوہناں نوں کِویں
دسیا گیا اے۔ ہور عام طور تے، عنصراں نال تعیّن دا اصول ایہ یقینی بناندا اے کہ
اوہناں $x$ دا سیٹ جنہاں لئی $\phi(x)$ ہووے، ہمیشہ اکو ہوندا اے۔
سو، عنصراں نال تعیّن دا اصول $\Setabs{x}{\phi(x)}$ نوں اوہناں $x$ دا
\emph{اوہی} سیٹ آکھن دا جواز دیندا اے جنہاں لئی~$\phi(x)$ ہووے۔
\end{ex}
\end{tagblock}

عنصراں نال تعیّن دا اصول ایہ وکھاؤن دا طریقہ دیندا اے کہ سیٹ اکو نیں: $A = B$
وکھاؤن لئی ایہ وکھاؤ کہ جدوں وی $x \in A$ ہووے، تاں $x \in B$ وی ہووے،
تے جدوں وی $y \in B$ ہووے، تاں $y \in A$ وی ہووے۔

\begin{prob}
ثابت کرو کہ ودھ توں ودھ اک خالی سیٹ ہو سکدا اے، یعنی ایہ وکھاؤ کہ جے $A$ تے
$B$ ایہو جہے سیٹ نیں جنہاں دے کوئی !!{element}s نہیں، تاں $A = B$ اے۔
\end{prob}
\endgroup
% /OLP-0005

% OLP-0006 content/sets-functions-relations/sets/subsets.tex
\begingroup
\def\olfilename{subsets}
\expandafter\def\csname ol@sectioncs\endcsname{section}

\olfileid{sfr}{set}{sub}
\olsection{ذیلی سیٹ تے ذیلی سیٹاں دے سیٹ \pnbLatin{(Subsets and Power Sets)}}

\begin{explain}
سانوں اکثر سیٹاں دا آپس وچ ٹاکرا کرن دی لوڑ پئے گی۔ ٹاکرے دی اک صاف صورت
ایہ اے: \emph{اک سیٹ وچ جیہڑی وی شے اے، اوہ دوجے وچ وی اے}۔ ایہ صورت اینی
اہم اے کہ ایس لئی کجھ نویاں علامتاں متعارف کرائیاں جان۔
\end{explain}

\begin{defn}[ذیلی سیٹ (Subset)]
جے سیٹ $A$ دا ہر !!{element}، $B$ دا وی !!a{element} ہووے، تاں اسی آکھدے آں
کہ $A$، $B$ دا \emph{ذیلی سیٹ} اے، تے $A \subseteq B$ لکھدے آں۔ جے
$A$، $B$ دا ذیلی سیٹ نہ ہووے، تاں اسی $A \not\subseteq B$ لکھدے آں۔
جے $A \subseteq B$ ہووے پر $A \neq B$ ہووے، تاں اسی $A \subsetneq B$
لکھدے آں تے آکھدے آں کہ $A$، $B$ دا \emph{حقیقی ذیلی سیٹ}
\pnbLatin{(proper subset)} اے۔
\end{defn}

\begin{ex}
ہر سیٹ اپنے آپ دا ذیلی سیٹ اے، تے $\emptyset$ ہر سیٹ دا ذیلی سیٹ اے۔ قدرتی
جفت اعداد \pnbLatin{(even natural numbers)} دا سیٹ قدرتی اعداد دے سیٹ دا ذیلی سیٹ اے۔
ایس طرح
$\{ a, b \} \subseteq \{ a, b, c \}$۔ پر $\{ a, b, e \}$،
$\{ a, b, c \}$ دا ذیلی سیٹ نہیں اے۔
\end{ex}

\begin{ex}
عدد $2$ صحیح اعداد (integers) دے سیٹ دا !!{element} اے، جد کہ جفت اعداد دا
سیٹ صحیح اعداد دے سیٹ دا ذیلی سیٹ اے۔ پر ایہ وی ہو سکدا اے کہ اک سیٹ کِسے ہور سیٹ
دا !!a{element} وی ہووے تے اوہدا ذیلی سیٹ وی، یعنی \emph{دونیں} گلاں ہون؛
مثال لئی، $\{0\} \in \{0, \{0\}\}$ تے نال ہی
$\{0\} \subseteq \{0, \{0\}\}$۔
\end{ex}

عنصراں نال تعیّن دا اصول سیٹاں دے اکو ہون دی کسوٹی دیندا اے: $A = B$ اُدوں
تے صرف اُدوں ہوندا اے جدوں $A$ دا ہر !!{element}، $B$ دا وی !!a{element}
ہووے تے ایس توں اُلٹ وی۔ ``ذیلی سیٹ'' دی تعریف $A \subseteq B$ نوں عین ایس
کسوٹی دے پہلے ادھ دے طور تے دسدی اے: $A$ دا ہر !!{element}، $B$ دا وی
!!a{element} اے۔ بے شک، جے اسی $A$ تے $B$ نوں آپس وچ وٹا دیئیے، تاں وی
تعریف لاگو ہوندی اے: یعنی $B \subseteq A$ اُدوں تے صرف اُدوں جدوں $B$ دا
ہر !!{element}، $A$ دا وی !!a{element} ہووے۔ تے ایہ عین عنصراں نال تعیّن
دے اصول دا ``ایس توں اُلٹ وی'' والا حصہ اے۔ ہور لفظاں وچ، ایس اصول توں
نکلدا اے کہ سیٹ اُدوں تے صرف اُدوں برابر نیں جدوں اوہ اک دوجے دے ذیلی سیٹ ہون۔

\begin{prop}
$A = B$ اُدوں تے صرف اُدوں ہوندا اے جدوں $A \subseteq B$ تے
$B \subseteq A$، دونیں ہون۔
\end{prop}

ہُن کجھ ہور مددگار علامتاں متعارف کراؤن دا وی چنگا موقع اے۔ ایہ تعریف
کردیاں کہ $A$، $B$ دا ذیلی سیٹ کدوں اے، اسی آکھیا سی کہ ``$A$ دا ہر
!!{element} \dots اے''، تے ``$\dots$'' دی تھاں ``$B$ دا !!a{element}''
رکھیا سی۔ پر عبارت دی ایہ \emph{بنت} اینی عام اے کہ ایس لئی کجھ رسمی
علامتاں متعارف کرانا مددگار ہووے گا۔

\begin{defn}\ollabel{forallxina}
$(\forall x \in A)\phi$، $\forall x(x \in A \lif
\phi)$ دی مختصر لکھت اے۔ ایس طرح، $(\exists x \in A)\phi$،
$\exists x(x \in A \land \phi)$ دی مختصر لکھت اے۔
\end{defn}

ایہ علامتاں ورت کے، اسی آکھ سکدے آں کہ $A \subseteq B$ اُدوں تے صرف اُدوں
ہوندا اے جدوں $(\forall x \in A)x \in B$ ہووے۔

ہُن اسی سیٹ دی اک خاص قسم ول آندے آں: کِسے دتے ہوئے سیٹ دے سارے ذیلی
سیٹاں دا سیٹ۔

\begin{defn}[ذیلی سیٹاں دا سیٹ \pnbLatin{(Power Set)}]
سیٹ~$A$ دے سارے ذیلی سیٹاں توں بنے سیٹ نوں $A$ دا \emph{ذیلی سیٹاں دا سیٹ}
\pnbLatin{(power set)} آکھدے نیں تے $\Pow{A}$ لکھدے نیں۔
  \[
    \Pow{A} = \Setabs{B}{B \subseteq A}
  \]
\end{defn}

\begin{ex}
$\{ a, b, c \}$ دے سارے ممکن ذیلی سیٹ کِہڑے نیں؟ اوہ ایہ نیں:
$\emptyset$، $\{a \}$، $\{b\}$، $\{c\}$، $\{a, b\}$، $\{a, c\}$،
$\{b, c\}$، $\{a, b, c\}$۔ ایہناں سارے ذیلی سیٹاں دا سیٹ
$\Pow{\{a,b,c\}}$ اے:
\[
\Pow{\{ a, b, c \}} = \{\emptyset, \{a \}, \{b\}, \{c\}, \{a, b\},
\{b, c\}, \{a, c\}, \{a, b, c\}\}
\]
\end{ex}

\begin{prob}
$\{a, b, c, d\}$ دے سارے ذیلی سیٹ لکھو۔
\end{prob}

\begin{prob}
وکھاؤ کہ جے $A$ دے $n$ !!{element}s ہون، تاں $\Pow{A}$ دے $2^n$
!!{element}s ہون گے۔
\end{prob}
\endgroup
% /OLP-0006

% OLP-0007 content/sets-functions-relations/sets/important-sets.tex
\begingroup
\def\olfilename{important-sets}
\expandafter\def\csname ol@sectioncs\endcsname{section}

\olfileid{sfr}{set}{imp}
\olsection{کجھ اہم سیٹ}

\begin{ex}
ساڈا واسطہ زیادہ تر اوہناں سیٹاں نال پئے گا جیہناں دے !!{element}s ریاضی
دیاں شےواں نیں۔ ایہو جہے چار سیٹ اینے اہم نیں کہ اوہناں دے خاص ناں نیں:
\begin{multline*}
    \Nat = \{0, 1, 2, 3, \ldots\} \\
    \shoveright{\text{قدرتی اعداد (natural numbers) دا سیٹ}}\\
    \shoveleft{\Int = \{\ldots, -2, -1, 0, 1, 2, \ldots\}} \\
    \shoveright{\text{صحیح اعداد (integers) دا سیٹ}}\\
    \shoveleft{\Rat = \Setabs{\nicefrac{m}{n}}{m, n \in \Int\text{ تے }n \neq 0}}\\
    \shoveright{\text{ناطق اعداد (rationals) دا سیٹ}}\\
    \shoveleft{\Real = (-\infty, \infty)}\\
    \text{حقیقی اعداد (real numbers) دا سیٹ؛ تسلسل (continuum)}
\end{multline*}
ایہ سارے \emph{لامتناہی (infinite)} سیٹ نیں، یعنی ہر اک دے !!{element}s دی گنتی
لامتناہی اے۔

ایہناں سیٹاں وچ اگے ودھدیاں، اسی اپنے کول \emph{ہور} عدد شامل کردے آں۔ ایہ
صاف ہونا چاہیدا اے کہ $\Nat \subseteq \Int \subseteq \Rat \subseteq \Real$:
آخر ہر قدرتی عدد صحیح عدد اے؛ ہر صحیح عدد ناطق عدد اے؛ تے ہر ناطق عدد
حقیقی عدد اے۔ ایس طرح ایہ وی صاف ہونا چاہیدا اے کہ
$\Nat \subsetneq \Int \subsetneq \Rat$، کیوں جو $-1$ صحیح عدد اے پر
قدرتی عدد نہیں، تے $\nicefrac{1}{2}$ ناطق عدد اے پر صحیح عدد نہیں۔ ایہ گل
اینی صاف نہیں کہ $\Rat \subsetneq \Real$، یعنی کجھ حقیقی عدد ایہو جہے
نیں جو ناطق نہیں\oliflabeldef{sfr:arith:real:realline}{، پر اسی
\olref[arith][real]{realline} وچ ایس ول مُڑ آواں گے}{}۔

کدی کدی اسی مثبت صحیح اعداد دا سیٹ $\PosInt = \{1, 2, 3, \dots\}$ تے
صرف پہلے دو قدرتی اعداد والا سیٹ $\Bin = \{0, 1\}$ وی ورتاں گے۔
\end{ex}

\begin{tagblock}{compsci}
\begin{ex}[لڑیاں (Strings)]
اک ہور دلچسپ مثال سیٹ $A^{*}$ اے، جہڑا حرفاں دے مجموعے (alphabet) $A$ اُتے
\emph{محدود لڑیاں \pnbLatin{(finite strings)}} دا سیٹ اے: $A$ دے عنصراں دا کوئی وی محدود
سلسلہ \pnbLatin{(finite sequence)}، $A$ اُتے اک لڑی اے۔ اسی \emph{خالی لڑی \pnbLatin{(empty string)}
$\Lambda$} نوں وی $A$ اُتے بنن والیاں لڑیاں وچ
شامل کردے آں، ہر حرفی مجموعے~$A$ لئی۔ مثال لئی،
\begin{multline*}
\Bin^*
=\{\Lambda,0,1,00,01,10,11,\\
000,001,010,011,100,101,110,111,0000,\ldots\}.
\end{multline*}
جے $x=x_{1}\ldots x_{n}\in A^{*}$ اک لڑی ہووے جہدے وچ $n$ ``حرف'' ہون،
جو $A$ توں لئے گئے ہون، تاں اسی آکھدے آں کہ لڑی دی \emph{لمبائی (length)}~$n$ اے تے
$\len{x}=n$ لکھدے آں۔
\end{ex}
\end{tagblock}

\begin{ex}[لامتناہی سلسلے \pnbLatin{(Infinite sequences)}]
کِسے وی سیٹ $A$ لئی اسی لامتناہی سلسلیاں دا سیٹ~$A^\omega$ وی ویکھ سکدے آں،
جیہڑا $A$ دے \pnboblique{!!{element}s} دے سلسلیاں توں بنیا اے۔ اک لامتناہی سلسلہ
$a_1a_2a_3a_4\dots$ شےواں دی اک پاسے ول لامتناہی فہرست \pnbLatin{(one-way infinite list)} اے، جہدی ہر
شے $A$ دا !!a{element} اے۔
\end{ex}
\endgroup
% /OLP-0007

% OLP-0008 content/sets-functions-relations/sets/unions-and-intersections.tex
\begingroup
\def\olfilename{unions-and-intersections}
\expandafter\def\csname ol@sectioncs\endcsname{section}

\olfileid{sfr}{set}{uni}
\olsection{اتحاد تے اشتراک \pnbLatin{(Unions and Intersections)}}

\begin{explain}
\olref[sfr][set][bas]{sec} وچ اسی خاصیت راہیں سیٹاں دی تعریف
\pnbLatin{(definition by abstraction)} متعارف کرائی
سی، یعنی $\Setabs{x}{\phi(x)}$ دی صورت والیاں تعریفاں۔ ایتھے اسی کوئی
خاصیت~$\phi$ ورتدے آں، تے ایس خاصیت وچ اوہناں سیٹاں دا ذکر ہو سکدا اے
جنہاں دی تعریف اسی پہلاں کر چکے آں۔ مثال لئی، جے $A$ تے~$B$ سیٹ نیں،
تاں سیٹ $\Setabs{x}{x \in A \lor x \in B}$ اوہناں ساریاں شےواں دا بنیا اے
جو $A$ یا~$B$ دے !!{element}s نیں؛ یعنی ایہ اوہ سیٹ اے جو $A$ تے~$B$ دے
\pnboblique{!!{element}s} نوں اکٹھے کردا اے۔ اسی ایہنوں \olref{fig:union}
وانگ شکل وچ ویکھ سکدے آں، جتھے اُبھاریا ہویا علاقہ دوہاں سیٹاں $A$ تے~$B$
دے \pnboblique{!!{element}s} نوں اکٹھے ظاہر کردا اے۔

\begin{figure}
  \olasset{assets/diagrams/union.tikz}
  \caption{دو سیٹاں دا اتحاد $A \cup B$، $A$ تے~$B$ دے
   \pnboblique{!!{element}s} نوں اکٹھے لے کے بنن والا سیٹ اے۔}
  \ollabel{fig:union}
\end{figure}

سیٹاں اُتے ایہ عمل---اوہناں نوں اکٹھا کرنا---بہت فائدےمند تے عام اے، ایس
لئی اسی ایہنوں اک رسمی ناں تے اک علامت دیندے آں۔
\end{explain}

\begin{defn}[اتحاد (Union)]
دو سیٹاں $A$ تے $B$ دا \emph{اتحاد (union)}، جہنوں $A \cup B$ لکھدے نیں،
اوہناں ساریاں شےواں دا سیٹ اے جو $A$، $B$ یا دوہاں دے !!{element}s نیں۔
\[
A \cup B = \Setabs{x}{x \in A \lor x \in B}
\]
\end{defn}

\begin{ex}
کیوں جو \pnboblique{!!{element}s} نوں وار وار گِنن دا کوئی فرق نہیں پیندا،
دو ایہو جہے سیٹاں دا اتحاد جنہاں دا !!a{element} سانجھا ہووے، اوس
!!{element} نوں صرف اک واری رکھدا اے؛ مثال لئی،
$\{ a, b, c\} \cup \{ a, 0, 1\} = \{a, b, c, 0, 1\}$۔

کِسے سیٹ تے اوہدے کِسے ذیلی سیٹ دا اتحاد بس وڈا سیٹ ای اے:
$\{a, b, c \} \cup \{a \} = \{a, b, c\}$۔

سیٹ دا خالی سیٹ نال اتحاد اوہی سیٹ اے:
$\{a, b, c \} \cup \emptyset = \{a, b, c \}$۔
\end{ex}

\begin{prob}
ثابت کرو کہ جے $A \subseteq B$ ہووے، تاں $A \cup B = B$ اے۔
\end{prob}

\begin{explain}
اسی اتحاد دے مقابل اک ``دوہرا عمل'' \pnbLatin{(dual operation)} وی ویکھ سکدے آں۔ ایہ اوہ عمل
اے جو سارے اوہناں \pnboblique{!!{element}s} دا سیٹ بناندا اے جو
$A$ دے !!{element}s وی ہون تے $B$ دے !!{element}s وی۔ ایس عمل نوں
\emph{اشتراک (intersection)} آکھدے نیں، تے ایہنوں \olref{fig:intersection} وانگ وکھایا
جا سکدا اے۔
\begin{figure}
  \olasset{assets/diagrams/intersection.tikz}
  \caption{دو سیٹاں دا اشتراک $A \cap B$ اوہناں دے سانجھے
    \pnboblique{!!{element}s} دا سیٹ اے۔}
  \ollabel{fig:intersection}
\end{figure}
\end{explain}

\begin{defn}[اشتراک (Intersection)]
دو سیٹاں $A$ تے $B$ دا \emph{اشتراک (intersection)}، جہنوں $A \cap B$ لکھدے نیں،
اوہناں ساریاں شےواں دا سیٹ اے جو $A$ تے~$B$ دوہاں دے !!{element}s نیں۔
\[
A \cap B = \Setabs{x}{x \in A \land x \in B}
\]
دو سیٹاں نوں \emph{بے سانجھ (disjoint)} آکھدے نیں جے اوہناں دا اشتراک خالی ہووے۔
ایس دا مطلب اے کہ اوہناں دے کوئی !!{element}s سانجھے نہیں۔
\end{defn}

\begin{ex}
جے دو سیٹاں دے کوئی !!{element}s سانجھے نہ ہون، تاں اوہناں دا اشتراک
خالی اے: $\{ a, b, c\} \cap \{ 0, 1\} = \emptyset$۔

جے دو سیٹاں دے کجھ !!{element}s سانجھے ہون، تاں اوہناں دا اشتراک اوہناں
ساریاں دا سیٹ اے: $\{a, b, c \} \cap \{a, b, d \} = \{a, b\}$۔

کِسے سیٹ دا اوہدے کِسے ذیلی سیٹ نال اشتراک بس چھوٹا سیٹ ای اے:
$\{a, b, c\} \cap \{a, b\} = \{a, b\}$۔

کِسے وی سیٹ دا خالی سیٹ نال اشتراک خالی اے:
$\{a, b, c \} \cap \emptyset = \emptyset$۔
\end{ex}

\begin{prob}
باقاعدہ دلیل دے کے ثابت کرو کہ جے $A \subseteq B$ ہووے، تاں
$A \cap B = A$ اے۔
\end{prob}

\begin{explain}
اسی دو توں ودھ سیٹاں دا اتحاد یا اشتراک وی بنا سکدے آں۔ عام صورت وچ ایس
لئی اک سوہنا طریقہ ایہ اے: منّ لو تسیں اوہناں سارے سیٹاں نوں اکو سیٹ
وچ اکٹھا کر لیندے او جنہاں دا اتحاد (یا اشتراک) بنانا اے۔ فیر اسی اپنے
اصل سیٹاں دے اتحاد نوں اوہناں ساریاں شےواں دے سیٹ دے طور تے متعین کر
سکدے آں جو اکٹھے کیتے سیٹ دے گھٹ توں گھٹ اک !!{element} وچ ہون، تے
اشتراک نوں اوہناں ساریاں شےواں دے سیٹ دے طور تے جو اوہدے ہر !!{element}
وچ ہون۔
\end{explain}

\begin{defn}
جے $A$ سیٹاں دا سیٹ اے، تاں $\bigcup A$، سیٹ~$A$ دے
\pnboblique{!!{element}s} دے \pnboblique{!!{element}s} دا سیٹ اے:
\begin{align*}
\bigcup A & = \Setabs{x}{x \text{، دتے ہوئے سیٹ دے کسے !!a{element} دا عنصر اے: } A},
\text{ یعنی،}\\
& = \Setabs{x}{\text{کوئی } B \in A
  \text{ ایہو جہا اے کہ } x \in B}
\end{align*}
\end{defn}

\begin{defn}
جے $A$ سیٹاں دا سیٹ اے، تاں $\bigcap A$ اوہناں شےواں دا سیٹ اے جو
$A$ دے سارے عنصراں وچ سانجھیاں نیں:
\begin{align*}
\bigcap A & = \Setabs{x}{x \text{، دتے ہوئے سیٹ دے ہر !!{element} وچ اے: } A},
\text{ یعنی،}\\
 & = \Setabs{x}{\text{ہر } B \in A, x \in B}
\end{align*}
\end{defn}

\begin{ex}
منّ لو کہ $A = \{ \{ a, b \}, \{ a, d, e \}, \{ a, d \} \}$۔
فیر $\bigcup A = \{ a, b, d, e \}$ تے $\bigcap A = \{ a \}$ اے۔
\end{ex}
\begin{prob}
	وکھاؤ کہ جے $A$ اک سیٹ اے تے $A \in B$، تاں $A \subseteq \bigcup B$ اے۔
\end{prob}

اسی سیٹاں دے سلسلے $A_1$، $A_2$، \dots لئی وی ایہو کم کر سکدے آں:
\begin{align*}
\bigcup_i A_i & = \Setabs{x}{x \text{، دتے ہوئے سیٹاں وچوں کسے اک وچ شامل اے: } A_i}\\
\bigcap_i A_i & = \Setabs{x}{x \text{، دتے ہوئے ہر سیٹ وچ شامل اے: } A_i}.
\end{align*}

جدوں ساڈے کول سیٹاں دا \emph{اشاریہ (index)} ہووے، یعنی کوئی سیٹ $I$ جس لئی
اسی $A_i$ اُتے غور کردے ہاں، ہر $i \in I$ لئی، تاں اسی ایہ مختصر
لکھتاں وی ورت سکدے آں:
\begin{align*}
	\bigcup_{i \in I} A_i & = \bigcup \Setabs{A_i }{i \in I}\\
	\bigcap_{i \in I} A_i & = \bigcap\Setabs{A_i}{i \in I}
\end{align*}

آخر وچ، اسی سارے اوہناں \pnboblique{!!{element}s} دے سیٹ بارے سوچ سکدے
آں جو $A$ وچ نیں پر $B$ وچ نہیں۔ ایہنوں \olref{difference} وانگ وکھایا
جا سکدا اے۔

\begin{figure}
  \olasset{assets/diagrams/difference.tikz}
  \caption{دو سیٹاں دا فرق $A \setminus B$ اوہناں \pnboblique{!!{element}s}
    دا سیٹ اے جو $A$ وچ نیں پر $B$ دے !!{element}s نہیں۔}
  \ollabel{difference}
\end{figure}

\begin{defn}[فرق (Difference)]
\emph{سیٹاں دا فرق \pnbLatin{(set difference)}}~$A \setminus B$ سارے اوہناں \pnboblique{!!{element}s}
دا سیٹ اے جو $A$ وچ نیں پر $B$ دے !!{element}s نہیں، یعنی
\[
A\setminus B = \Setabs{x}{x\in A \text{ تے } x \notin B}.
\]
\end{defn}

\begin{prob}
	ثابت کرو کہ جے $A \subsetneq B$ ہووے، تاں $B \setminus A \neq \emptyset$ اے۔
\end{prob}
\endgroup
% /OLP-0008

% OLP-0009 content/sets-functions-relations/sets/pairs-and-products.tex
\begingroup
\def\olfilename{pairs-and-products}
\expandafter\def\csname ol@sectioncs\endcsname{section}

\olfileid{sfr}{set}{pai}
\olsection{جوڑے، ٹپل تے کارتیسی حاصلِ ضرب \pnbLatin{(Pairs, Tuples, Cartesian Products)}}

\begin{explain}
عنصراں نال تعیّن دے اصول توں نکلدا اے کہ سیٹ دے عنصراں دی کوئی ترتیب نہیں
ہوندی۔ سو، جے اسی ترتیب نوں ظاہر کرنا چاہنے آں، تاں \emph{ترتیب وار جوڑے
\pnbLatin{(ordered pairs)}} $\tuple{x, y}$ ورتدے آں۔ بے ترتیب جوڑے \pnbLatin{(unordered pair)}
$\{x, y\}$ وچ ترتیب دا کوئی فرق نہیں پیندا: $\{x, y\} = \{y, x\}$۔ ترتیب
وار جوڑے وچ فرق پیندا اے: جے
$x \neq y$ ہووے، تاں $\tuple{x, y} \neq \tuple{y, x}$ اے۔

سیٹاں دے نظریے وچ ترتیب وار جوڑیاں بارے ساڈا تصور کی ہونا چاہیدا اے؟ بنیادی
گل ایہ اے کہ اسی ایہ خیال قائم رکھنا چاہندے آں کہ ترتیب وار جوڑے اُدوں تے
صرف اُدوں اکو ہوندے نیں جدوں اوہناں دا پہلا عنصر اکو ہووے تے دوجا عنصر وی
اکو ہووے، یعنی:
\[
  \tuple{a, b}= \tuple{c, d}\text{ اُدوں تے صرف اُدوں، جدوں }a = c \text{ تے }b=d.
\]
سیٹاں دے نظریے وچ اسی وینر--کوراتوسکی دی تعریف نال ترتیب وار جوڑے دی
تعریف کر سکدے آں۔
\end{explain}

\begin{defn}[ترتیب وار جوڑا \pnbLatin{(Ordered pair)}]\ollabel{wienerkuratowski}
	$\tuple{a, b} = \{\{a\}, \{a, b\}\}$۔
\end{defn}

\begin{prob}
	\olref[sfr][set][pai]{wienerkuratowski} نوں ورت کے ثابت کرو کہ
	$\tuple{a, b}= \tuple{c, d}$ اُدوں تے صرف اُدوں ہوندا اے جدوں
	$a = c$ وی ہووے تے $b=d$ وی۔
\end{prob}

\begin{explain}
ترتیب وار جوڑے دی تعریف طے کر کے اسی اوہنوں ہور سیٹاں دی تعریف لئی ورت
سکدے آں۔ مثال لئی، کدی کدی سانوں دو توں ودھ چیزاں دے ترتیب وار سلسلے
\pnbLatin{(ordered sequences)} وی چاہیدے ہوندے نیں، جیویں \emph{تِن-ٹپل (triples)}
$\tuple{x, y, z}$، \emph{چار-ٹپل (quadruples)} $\tuple{x, y, z, u}$، تے
ایس طرح اگے۔ تِن-ٹپلاں نوں اسی خاص ترتیب وار جوڑے
سمجھ سکدے آں، جنہاں دا پہلا عنصر آپ اک ترتیب وار جوڑا اے:
$\tuple{x, y, z}$، $\tuple{\tuple{x, y},z}$ اے۔ چار-ٹپلاں لئی وی ایہو گل اے:
$\tuple{x,y,z,u}$، $\tuple{\tuple{\tuple{x,y},z},u}$ اے، تے ایس طرح اگے۔
عام طور تے اسی \emph{ترتیب وار $n$-ٹپل \pnbLatin{(ordered n-tuples)}}
$\tuple{x_1, \dots, x_n}$ دی گل کردے آں۔

ترتیب وار جوڑیاں، یا ہور ترتیب وار $n$-ٹپلاں دے کجھ سیٹ ساڈے کم آؤن گے۔
\end{explain}

\begin{defn}[کارتیسی حاصلِ ضرب \pnbLatin{(Cartesian product)}]
دتے ہوئے سیٹاں $A$ تے $B$ دا \emph{کارتیسی حاصلِ ضرب \pnbLatin{(Cartesian product)}}
$A \times B$ ایویں متعین ہوندا اے:
\[
  A \times B = \Setabs{\tuple{x, y}}{x \in A \text{ تے } y \in B}.
\]
\end{defn}

\begin{ex}
جے $A = \{0, 1\}$ تے $B = \{1, a, b\}$ ہون، تاں اوہناں دا حاصلِ ضرب اے:
\[
A \times B = \{ \tuple{0, 1}, \tuple{0, a}, \tuple{0, b},
    \tuple{1, 1}, \tuple{1, a}, \tuple{1, b} \}.
\]
\end{ex}

\begin{ex}
جے $A$ اک سیٹ اے، تاں $A$ دا اپنے آپ نال حاصلِ ضرب، $A \times A$،
$A^2$ وی لکھیا جاندا اے۔ ایہ اوہناں \emph{سارے} جوڑیاں $\tuple{x, y}$ دا
سیٹ اے جنہاں لئی $x, y \in A$ اے۔ سارے تِن-ٹپلاں $\tuple{x, y, z}$ دا
سیٹ $A^3$ اے، تے ایس طرح اگے۔ اسی اک بازگشتی تعریف \pnbLatin{(recursive definition)}
دے سکدے آں:
\begin{align*}
  A^1 & = A\\
  A^{k+1} & = A^k \times A
\end{align*}
\end{ex}

\begin{prob}
سارے !!{element}s لکھو جو $\{1, 2, 3\}^3$ وچ نیں۔
\end{prob}

\begin{prop}\ollabel{cardnmprod}
جے $A$ دے $n$ !!{element}s تے $B$ دے $m$ !!{element}s ہون، تاں
$A \times B$ دے $n\cdot m$ عنصر ہون گے۔
\end{prop}

\begin{proof}
ہر !!{element}~$x$ لئی جو $A$ وچ اے، $m$ ایہو جہے !!{element}s نیں
جنہاں دی صورت $\tuple{x, y} \in A \times B$ اے۔ منّ لو کہ
$B_x = \Setabs{\tuple{x, y}}{y \in B}$۔ کیوں جو جدوں وی $x_1 \neq x_2$
ہووے، تاں $\tuple{x_1, y} \neq \tuple{x_2, y}$ اے، ایس لئی
$B_{x_1} \cap B_{x_2} = \emptyset$۔ پر جے $A = \{x_1, \dots, x_n\}$
ہووے، تاں $A \times B = B_{x_1} \cup \dots \cup B_{x_n}$ اے، تے ایس لئی
اوہدے $n\cdot m$ !!{element}s نیں۔

ایس گل نوں شکل وچ ویکھن لئی، \pnboblique{!!{element}s} نوں جو $A \times B$
وچ نیں، قطاراں تے کالماں والے جال (grid) وچ ترتیب نال رکھو:
\[
\begin{array}{rcccc}
  B_{x_1} = & \{\tuple{x_1, y_1} & \tuple{x_1, y_2} & \dots & \tuple{x_1, y_m}\}\\
  B_{x_2} = & \{\tuple{x_2, y_1} & \tuple{x_2, y_2} & \dots & \tuple{x_2, y_m}\}\\
  \vdots & & \vdots\\
  B_{x_n} = & \{\tuple{x_n, y_1} & \tuple{x_n, y_2} & \dots & \tuple{x_n, y_m}\}
\end{array}
\]
کیوں جو سارے $x_i$ وکھ وکھ نیں، تے سارے $y_j$ وی وکھ وکھ نیں، ایس جال
وچ کوئی دو جوڑے اکو نہیں، تے اوہناں دی گنتی $n\cdot m$ اے۔
\end{proof}

\begin{prob}
$k$ اُتے ریاضیاتی استقرا \pnbLatin{(mathematical induction)} نال وکھاؤ کہ سارے
$k \ge 1$ لئی، جے $A$ دے
$n$ !!{element}s ہون، تاں $A^k$ دے $n^k$ !!{element}s ہون گے۔
\end{prob}

\begin{ex}
جے $A$ اک سیٹ اے، تاں $A$ اُتے \emph{لفظ (word)}، $A$ دے
\pnboblique{!!{element}s} دا کوئی وی سلسلہ (sequence) اے۔ ہر سلسلے نوں
اک $n$-ٹپل سمجھیا جا سکدا اے، جہدے !!{element}s $A$ دے ہون۔ مثال لئی، جے
$A = \{a, b, c\}$ ہووے، تاں سلسلے ``$bac$'' نوں تِن-ٹپل~$\tuple{b, a, c}$
دے طور تے سمجھیا
جا سکدا اے۔ لفظ، یعنی علامتاں دے سلسلے، کمپیوٹر سائنس وچ بنیادی اہمیت
رکھدے نیں۔ رواج موجب، اسی \pnboblique{!!{element}s} نوں جو $A$ وچ نیں،
لمبائی~$1$ والے سلسلے گِندے آں، تے $\emptyset$ نوں لمبائی~$0$ والا سلسلہ۔
فیر $A$ اُتے \emph{سارے} لفظاں دا سیٹ اے:
\[
A^* = \{\emptyset\} \cup A \cup A^2 \cup A^3 \cup \dots
\]
\end{ex}
\endgroup
% /OLP-0009

% OLP-0010 content/sets-functions-relations/sets/russells-paradox.tex
\begingroup
\def\olfilename{russells-paradox}
\expandafter\def\csname ol@sectioncs\endcsname{section}

\olfileid{sfr}{set}{rus}
\olsection{رسل دا پیراڈوکس \pnbLatin{(Russell's Paradox)}}

عنصراں نال تعیّن دا اصول علامتی لکھت $\Setabs{x}{\phi(x)}$ دا جواز دیندا
اے، جہڑی اوہناں $x$ دا \emph{اوہی} سیٹ دسدی اے جنہاں لئی~$\phi(x)$
ہووے۔ پر عنصراں نال تعیّن دا اصول \emph{اصل وچ} صرف ایس خیال دا جواز
دیندا اے: \emph{جے} اک ایہو جہا سیٹ موجود ہووے جہدے رکن سارے اوہ تے صرف
اوہ ہون جیہڑے $\phi$ خاصیت پوری کردے ہون، \emph{تاں} ایہو جہا سیٹ اکو اے۔
ہور لفظاں وچ: کِسے~$\phi$ نوں طے کرن مگروں، سیٹ
$\Setabs{x}{\phi(x)}$ یکتا (unique) اے،
\emph{جے اوہ موجود ہووے}۔

پر ایہ شرط بڑی اہم اے!{} بنیادی گل ایہ اے کہ ہر خاصیت نال
\emph{خاصیت توں سیٹ بنانا (comprehension)} جائز نہیں۔ یعنی کجھ خاصیتاں
\emph{سیٹ دی تعریف نہیں} کردیاں۔ جے ساریاں کردیاں، تاں اسی کھلے تضاداں
وچ پھس جانا سی۔ ایس دی سب توں مشہور مثال رسل دا پیراڈوکس اے۔

سیٹ ہور سیٹاں دے !!{element}s ہو سکدے نیں---مثال لئی، سیٹ~$A$ دا ذیلی
سیٹاں دا سیٹ، سیٹاں توں ای بنیا ہوندا اے۔ سو ایہ پچھنا یا کھوجنا بجا اے
کہ کوئی سیٹ دوجے سیٹ دا !!a{element} اے یا نہیں۔ کی کوئی سیٹ اپنے آپ دا
رکن ہو سکدا اے؟ سیٹ دے خیال وچ کوئی گل ایس نوں روکندی نہیں لگدی۔ مثال
لئی، جے \emph{سارے} سیٹ چیزاں دا اک اکٹھ بناندے نیں، تاں کوئی ایہ سوچ
سکدا اے کہ اوہناں نوں اکو سیٹ وچ اکٹھا کیتا جا سکدا اے---سارے سیٹاں دا
سیٹ۔ تے اوہ، آپ اک سیٹ ہون پاروں، سارے سیٹاں دے سیٹ دا !!a{element}
ہووے گا۔

رسل دا پیراڈوکس اُدوں پیدا ہوندا اے جدوں اسی اپنے آپ نوں !!a{element} نہ
رکھن دی خاصیت، یعنی \emph{اپنے آپ دا رکن نہ ہون (non-self-membered)} دی
خاصیت، اُتے غور کردے آں۔ جے اسی منّ لئیے کہ اوہناں سارے سیٹاں دا سیٹ موجود اے جو اپنے آپ نوں
!!a{element} نہیں رکھدے، تاں کی ہووے گا؟ کی
\[
R = \Setabs{x}{x \notin x}
\]
موجود اے؟ گل ایہ اے کہ اسی ثابت کر سکدے آں کہ ایہ موجود نہیں۔

\begin{thm}[رسل دا پیراڈوکس \pnbLatin{(Russell's Paradox)}]\ollabel{thm:russells-paradox}
	کوئی سیٹ $R = \Setabs{x}{x \notin x}$ موجود نہیں۔
\end{thm}

\begin{proof}
جے $R = \Setabs{x}{x \notin x}$ موجود ہووے، تاں
$R \in R$ اُدوں تے صرف اُدوں جدوں $R \notin R$، تے ایہ اک تضاد
(contradiction) اے۔
\end{proof}

\begin{tagblock}{novice}
\begin{explain}
آؤ ایس ثبوت نوں ہولی ہولی ویکھیے۔ جے $R$ موجود اے، تاں ایہ پُچھنا بجا اے
کہ $R \in R$ اے یا نہیں۔ منّ لو کہ واقعی $R \in R$ اے۔ ہُن $R$ دی
تعریف اوہناں سارے سیٹاں دے سیٹ دے طور تے کیتی گئی سی جو اپنے آپ دے
!!{element}s نہیں۔ سو، جے $R \in R$، تاں $R$ آپ اوہ خاصیت نہیں رکھدا
جس نال $R$ دی تعریف کیتی گئی اے۔ پر صرف اوہ سیٹ $R$ وچ نیں جو ایہ
خاصیت رکھدے نیں، ایس لئی $R$، $R$ دا !!a{element} نہیں ہو سکدا، یعنی
$R \notin R$۔ پر $R$، $R$ دا !!a{element} ہووے وی تے نہ وی ہووے، ایہ
نہیں ہو سکدا؛ سو ساڈے کول اک تضاد اے۔

کیوں جو ایہ فرض کرنا کہ $R \in R$ تضاد ول لے جاندا اے، ایس لئی
$R \notin R$ اے۔ پر ایہ وی تضاد ول لے جاندا اے!{} کیوں جو جے $R \notin R$
ہووے، تاں $R$ آپ اوہ خاصیت رکھدا اے جس نال $R$ دی تعریف کیتی گئی اے،
تے ایس لئی $R$، $R$ دا !!a{element} ہووے گا، بالکل اوہناں ہور سیٹاں
وانگ جو اپنے آپ دے رکن نہیں۔ تے فیر، ایہ نہیں ہو سکدا کہ اوہ $R$ دا
!!a{element} نہ وی ہووے تے ہووے وی۔
\end{explain}
\end{tagblock}

\begin{digress}
اسی سیٹاں دا ایہو جہا نظریہ کِویں بنائیے جو رسل دے پیراڈوکس وچ نہ پھسے،
یعنی ایہ \emph{غیر موافق (inconsistent)} دعویٰ نہ کرے کہ
$R = \Setabs{x}{x \notin x}$ موجود اے؟ ایس لئی سانوں مسلّمے (axioms)
مقرر کرنے پَین گے، جو ایہ دسن لئی بالکل
درست شرطاں دین کہ سیٹ کدوں موجود ہوندے نیں (تے کدوں نہیں ہوندے)۔

ایس باب وچ سیٹاں دے نظریے دا جیہڑا خاکہ دتا گیا اے، اوہ ایہ کم نہیں کردا۔
ایہ \emph{سچ مچ سادہ \pnbLatin{(genuinely naive)}} سیٹ نظریہ اے۔ ایہ صرف دسدا اے کہ سیٹ عنصراں نال
تعیّن دے اصول اُتے چلدے نیں تے، جے تہاڈے کول کجھ سیٹ ہون، تاں تسیں اوہناں
دا اتحاد، اشتراک وغیرہ بنا سکدے او۔ سیٹاں دے نظریے نوں ایس توں زیادہ
باقاعدگی نال بنایا جا سکدا اے۔ \oliflabeldef{cumul:::part}{اوہ باقاعدگی
حصہ \olref[cumul][][]{part} لئی چھڈی جائے گی۔ ہُن لئی، اسی سادہ ڈھنگ نال
(naively) اگے ودھاں گے تے احتیاط نال تضاداں توں بچن دی کوشش کراں گے۔}{}
\end{digress}
\endgroup
% /OLP-0010

% OLP-0721 content/sets-functions-relations/sets/proofs-about-sets.tex
\begingroup
\def\olfilename{proofs-about-sets}
\expandafter\def\csname ol@sectioncs\endcsname{section}

\olfileid{sfr}{set}{prf}
\olsection{سیٹاں بارے ثبوت}

\begin{editorial}
ایس حصے دی تھاں \verb|content/method/proofs| نے لے لئی اے تے ایہ عام ترتیب وچ ایس باب وچوں ہٹا دِتّا گیا اے۔
\end{editorial}

\begin{explain}
سیٹ تے اوہ نشانیاں جیہڑیاں اسیں ہُن تک دَسّیاں نیں، سانوں سیٹ دَسّن تے اوہناں وچکار تعلق ظاہر کرن دے سوکھے چھوٹے روپ دیندیاں نیں۔ اکثر ایہو جیہے تعلقاں بارے دعوے ثابت کرن دی وی لوڑ ہووے گی۔ جے تسی ریاضی دے ثبوتاں توں واقف نہیں، تاں ہو سکدا اے ایہ تہاڈے لئی نویں گل ہووے۔ ایس لئی اسیں اک سادی مثال قدم بہ قدم ویکھاں گے۔ اسیں ثابت کراں گے کہ کسے وی دو سیٹاں $X$ تے $Y$ لئی ہمیشہ $X \cap
(X \cup Y) = X$۔ سیٹاں وچ ایہو جیہی برابری کِنج ثابت کردے او؟ چیتے رکھو کہ سیٹ صرف اپنے !!{element}s نال طے ہوندے نیں، یعنی سیٹ اوہدوں تے صرف اوہدوں برابر نیں جدوں اوہناں دے !!{element}s اوہی ہون۔ سو ایس صورت وچ سانوں ثابت کرنا پینا اے کہ (a) ہر !!{element}، $X \cap (X \cup
Y)$ دا، !!{element} وی اے $X$ دا، تے الٹے پاسے (b) ہر !!{element}، $X$ دا، !!{element} وی اے $X \cap (X \cup Y)$ دا۔ دوجے لفظاں وچ اسیں (a) $X \cap (X \cup Y) \subseteq X$ تے (b) $X \subseteq X \cap (X \cup Y)$ دونے وکھاؤنے آں۔

عام دعوے دا ثبوت، جیویں «ہر !!{element}~$z$، $X \cap (X
\cup Y)$ دا، !!{element} وی اے $X$ دا»، پہلاں منّن نال ہوندا اے کہ کوئی من مرضی دا $z \in X \cap (X \cup Y)$ دِتّا اے، تے ایس مفروضے توں ثابت کرن نال کہ $z \in X$۔ ہو سکدا اے تسی ایس طریقے نوں «عام شرطیہ ثبوت» دے ناں نال جاندے ہوو۔ ایس ثبوت وچ سانوں مفروضے تے نتیجے وچ ورتیاں تعریفاں وی ورتنیاں پین گئیاں، جیویں ایتھے «$\cap$» تے «$\cup$» دیاں۔ سو (a) دی دلیل ایویں ہووے گی:

\begin{quote}
(a) پہلاں اسیں وکھاؤنا چاہنے آں کہ $X \cap (X \cup Y) \subseteq X$، یعنی $\subseteq$ دی تعریف نال، جے $z \in X \cap (X \cup Y)$ تاں $z \in X$، کسے وی~$z$ لئی۔ سو منّ لو $z \in X \cap (X \cup
Y)$۔ کیونکہ $z$ دو سیٹاں دے اشتراک دا !!{element} اوہدوں تے صرف اوہدوں اے جدوں اوہ دوناں سیٹاں دا !!{element} اے، ایس لئی نتیجہ کڈھ سکدے آں کہ $z \in X$ تے نالے $z \in X \cup Y$۔ خاص کر، $z \in X$، جیہڑا اسیں وکھاؤنا چاہندے ساں۔
\end{quote}

ایس نال ثبوت دا پہلا ادھ پورا ہو گیا۔ ویکھو کہ آخری قدم اُتے اسیں ایہ گل ورتی سی کہ جے اک عطف ($z \in X$ تے $z \in X
\cup Y$) کسے مفروضے توں نکلدا اے، تاں عطف دا ہر جزو اوسے مفروضے توں نکلدا اے۔ ہو سکدا اے تسی ایس قاعدے نوں «عطف ہٹاؤنا»، یا $\land$ ہٹاؤنا، دے ناں نال جاندے ہووے۔ ہُن (b) ثابت کریے:

\begin{quote}
(b) ہُن اسیں ثابت کردے آں کہ $X \subseteq X \cap (X \cup Y)$، یعنی $\subseteq$ دی تعریف نال، جے $z \in X$ تاں $z \in X \cap
(X \cup Y)$ وی، کسے وی~$z$ لئی۔ منّ لو $z \in X$۔ ایہ وکھاؤن لئی کہ $z \in X
\cap (X \cup Y)$، سانوں «$\cap$» دی تعریف نال وکھاؤنا پینا اے کہ (i) $z \in X$ تے (ii) $z \in X \cup Y$ وی۔ ایتھے (i) بس ساڈا مفروضہ اے، سو ہور کجھ ثابت کرن نوں نہیں۔ (ii) لئی چیتے رکھو کہ $z$ سیٹاں دے اتحاد دا !!{element} اوہدوں تے صرف اوہدوں اے جدوں اوہ اوہناں وچوں گھٹّو گھٹ اک سیٹ دا !!{element} اے۔ کیونکہ $z \in X$، تے $X \cup Y$، $X$ تے $Y$ دا اتحاد اے، ایس لئی ایتھے ایہ شرط پوری اے۔ سو $z \in X \cup Y$۔ اسیں (i) $z \in X$ تے (ii) $z \in X \cup Y$ دونے وکھا دِتّے نیں۔ ایس لئی «$\cap$» دی تعریف نال $z \in X \cap (X \cup Y)$۔
\end{quote}

ایہ ذرا لمّی گل ہو گئی، پر ایہ وکھاؤندی اے کہ اسیں سیٹاں تے اوہناں دے تعلقاں بارے دلیل کِنج دیندے آں۔ عام طور اُتے اسیں اینا کھولھ کے نہیں لکھدے۔ خاص کر ہو سکدا اے اسیں ساریاں تعریفاں مُڑ نہ دسیے۔ اگلے درجے دی کتاب وچ ساڈے نتیجے دا ثبوت بہت چھوٹا ہووے گا۔ ہو سکدا اے ایہ کجھ ایویں لگے۔
\end{explain}

\begin{prop}[جذب]
سارے سیٹاں $X$، $Y$ لئی،
\[
X \cap (X \cup Y) = X
\]
\end{prop}

\begin{proof}
(a) منّ لو $z \in X \cap (X \cup Y)$۔ تاں $z \in X$، سو $X \cap (X
  \cup Y) \subseteq X$۔

(b) ہُن منّ لو $z \in X$۔ تاں $z \in X \cup Y$ وی، تے ایس لئی $z \in X \cap (X \cup Y)$ وی۔ سو $X \subseteq X \cap (X \cup
  Y)$۔
\end{proof}

\begin{prob}
تفصیل نال ثابت کرو کہ $X \cup (X \cap Y) = X$۔ فیر اک چھوٹا، گھُٹواں ثبوت دیو۔ (اشارہ۔ $X \cup (X \cap Y) \subseteq X$ والی سمت لئی تہانوں صورتاں نال ثبوت، یعنی $\lor$ ہٹاؤنا، دی لوڑ پئے گی۔)
\end{prob}
\endgroup
% /OLP-0721
\OLEndChapterHook
\endgroup
% /OLP-0004

% OLP-0011 content/sets-functions-relations/relations/relations-complete.tex
\begingroup
\def\olfilename{relations-complete}
\expandafter\def\csname ol@sectioncs\endcsname{section}

\olchapter{sfr}{rel}{تعلق \pnbLatin{(Relations)}}

% OLP-0012 content/sets-functions-relations/relations/relations-as-sets.tex
\begingroup
\def\olfilename{relations-as-sets}
\expandafter\def\csname ol@sectioncs\endcsname{section}

\olfileid{sfr}{rel}{set}
\olsection{سیٹاں دے روپ وچ تعلق \pnbLatin{(Relations as Sets)}}

\begin{explain}
\olref[sfr][set][imp]{sec} وچ اسی کجھ اہم سیٹاں دا ذکر کیتا سی:
$\Nat$, $\Int$, $\Rat$, $\Real$۔ تہانوں ضرور یاد ہووے گا کہ
ایہناں وچوں کجھ سیٹاں دے \pnboblique{!!{element}s} وچکار کجھ دلچسپ تعلق ہوندے نیں۔
مثال دے طور تے، ایہناں وچوں ہر سیٹ اُتے اک بالکل معیاری \emph{ترتیبی
تعلق \pnbLatin{(order relation)}} ہوندا اے۔ اک \emph{اکو ہون دا \pnbLatin{(is identical with)}} تعلق وی اے، جیہڑا ہر شے دا
اپنے آپ نال ہوندا اے تے ہور کسے شے نال نہیں ہوندا۔ اگے سانوں
ہور کئی دلچسپ تعلق ملن گے، تے ممکن تعلق تاں اوہناں توں وی ودھ نیں۔
پر اوہناں دا جائزہ لین توں پہلاں اسی ایس گل توں شروع کراں گے کہ
تعلقاں نوں اک خاص قسم دیاں سیٹاں دے روپ وچ ویکھیا جا سکدا اے۔
  
ایس لئی \olref[sfr][set][pai]{sec} دیاں دو گلاں یاد کرو۔ پہلی،
\emph{ترتیب وار جوڑے} دا تصور یاد کرو: جے $a$ تے $b$ دتے ہون، تاں اسی
$\tuple{a, b}$ بنا سکدے آں۔ اہم گل ایہ اے کہ ایتھے عنصراں دی ترتیب
\emph{واقعی} فرق پاؤندی اے۔ سو جے $a \neq b$، تاں $\tuple{a, b} \neq \tuple{b, a}$۔
(ایہنوں بے ترتیب جوڑیاں، یعنی $2$ عنصراں والیاں سیٹاں، نال ملا کے ویکھو، جتھے
$\{a, b\}=\{b, a\}$۔) دوجی، \emph{کارتیسی
حاصلِ ضرب} دا تصور یاد کرو: جے $A$ تے $B$ سیٹ ہون، تاں اسی $A \times B$ بنا
سکدے آں؛ ایہ اوہناں سارے جوڑیاں $\tuple{x, y}$ دا سیٹ اے جنہاں لئی $x \in A$ تے
$y \in B$ ہووے۔ خاص طور تے، $A^{2}= A \times A$، $A$ توں بنن والے سارے ترتیب وار
جوڑیاں دا سیٹ اے۔
  
ہن اسی اک سیٹ اُتے اک خاص تعلق ویکھاں گے: $<$ دا تعلق، جیہڑا قدرتی عدداں
دے سیٹ $\Nat$ اُتے اے۔ عدداں دے اوہ سارے جوڑے $\tuple{n, m}$ لو جتھے $n<m$، یعنی
\[
  R=\Setabs{\tuple{n, m}}{n, m \in \Nat \text{ تے } n<m}.
\]
$n$ دا $m$ توں چھوٹا ہونا تے جوڑے $\tuple{n, m}$ دا $R$ دا رکن ہونا،
ایہناں وچکار ڈونگھا جوڑ اے، یعنی:
\[
      n<m\text{ اُدوں تے صرف اُدوں جدوں }\tuple{n, m} \in R.
\]
درحقیقت، معلومات دا کوئی حصہ گوائے بغیر اسی سیٹ $R$ نوں ای
$<$ دا تعلق \emph{منّ} سکدے آں، جیہڑا $\Nat$ اُتے اے۔

ایسے طرح عدداں وچکار کسے وی تعلق لئی اسی $\Nat^{2}$ دا اک ذیلی سیٹ بنا
سکدے آں۔ پُٹھے پاسے، جے عدداں دے جوڑیاں دا کوئی وی سیٹ
$S \subseteq \Nat^{2}$ دتا ہووے، تاں اوہدے مطابق عدداں وچکار اک تعلق ہوندا اے:
اوہ تعلق جیہڑا $n$ دا $m$ نال اُدوں تے صرف اُدوں ہووے جدوں
$\tuple{n, m} \in S$۔ ایس بنا اُتے اگلی تعریف جائز اے:
\end{explain}
  
\begin{defn}[دو رکنی تعلق \pnbLatin{(Binary relation)}]
اک سیٹ $A$ اُتے \emph{دو رکنی تعلق \pnbLatin{(binary relation)}}، $A^{2}$ دا اک ذیلی سیٹ ہوندا اے۔ جے $R
\subseteq A^{2}$، $A$ اُتے اک دو رکنی تعلق ہووے تے $x, y \in A$ ہون، تاں
اسی کدی کدی $Rxy$ (یا $xRy$)، $\tuple{x, y} \in R$ دی تھاں لکھدے آں۔
\end{defn}
  
\begin{ex}
  \ollabel{relations}
قدرتی عدداں دے جوڑیاں دا سیٹ $\Nat^{2}$، دو جہتی میٹرکس وچ ایس طرح
لکھیا جا سکدا اے:
\[
  \begin{array}{ccccc}
  \mathbf{\tuple{ 0,0 }} & \tuple{ 0,1 } &
    \tuple{ 0,2 } & \tuple{ 0,3 } & \ldots\\
  \tuple{ 1,0 } & \mathbf{\tuple{ 1,1 }} &
    \tuple{ 1,2 } & \tuple{ 1,3 } & \ldots\\
  \tuple{ 2,0 } & \tuple{ 2,1 } &
    \mathbf{\tuple{ 2,2 }} & \tuple{ 2,3 } & \ldots\\
  \tuple{ 3,0 } & \tuple{ 3,1 } & \tuple{ 3,2 } &
    \mathbf{\tuple{ 3,3 }} & \ldots\\
  \vdots & \vdots & \vdots & \vdots & \mathbf{\ddots}
  \end{array}
\]
ایتھے اسی ترچھے قطر نوں موٹے اکھراں وچ وکھایا اے، کیوں جو $\Nat^2$ دا
اوہ ذیلی سیٹ جیہدے جوڑے ایس قطر اُتے پئے نیں، یعنی
\[
  \{\tuple{0,0 }, \tuple{ 1,1 }, \tuple{ 2,2 }, \dots\},
  \]
$\Nat$ اُتے \emph{اکو ہون دا تعلق \pnbLatin{(identity relation)}} اے۔ (اکو ہون دا تعلق بہت ورتیا جاندا اے،
سو کسے وی سیٹ لئی $\Id{A}=\Setabs{\tuple{ x,x }}{x \in
A}$ دی تعریف کر لئیے، جتھے سیٹ $A$ اے۔) قطر توں اُتے پئے سارے جوڑیاں دا
ذیلی سیٹ، یعنی
\[
  L = \{\tuple{ 0,1 },\tuple{ 0,2 },\ldots,\tuple{ 1,2 },
  \tuple{ 1,3 }, \dots, \tuple{ 2,3 }, \tuple{ 2,4 },\ldots\},
\]
\emph{توں چھوٹا ہون} دا تعلق اے، یعنی $Lnm$ اُدوں تے صرف اُدوں جدوں $n<m$۔
قطر توں تھلے پئے جوڑیاں دا ذیلی سیٹ، یعنی
\[
  G=\{ \tuple{ 1,0 },\tuple{ 2,0 },\tuple{
    2,1 }, \tuple{ 3,0 },\tuple{ 3,1 },\tuple{ 3,2 }, \dots\},
\]
\emph{توں وڈا ہون} دا تعلق اے، یعنی $Gnm$ اُدوں تے صرف اُدوں جدوں $n>m$۔
$L$ دا $I$ نال اتحاد، جیہنوں اسی $K=L\cup I$ آکھ سکدے آں، \emph{توں چھوٹا
یا برابر ہون} دا تعلق اے: $Knm$ اُدوں تے صرف اُدوں جدوں $n \le m$۔ ایسے طرح
$H=G \cup I$، \emph{توں وڈا یا برابر ہون دا تعلق} اے۔ ایہ تعلق
$L$, $G$, $K$ تے $H$ اک خاص قسم دے تعلق نیں جنہاں نوں
\emph{ترتیباں (orders)} آکھدے نیں۔ $L$ تے $G$ دی ایہ خاصیت اے کہ کسے وی عدد دا اپنے
آپ نال نہ $L$ دا تعلق اے تے نہ $G$ دا (یعنی ہر $n$ لئی، نہ $Lnn$ تے نہ
$Gnn$)۔ ایس خاصیت والے تعلقاں نوں \emph{اپنے آپ نال تعلق نہ ہون والے (irreflexive)} تعلق
آکھدے نیں؛ تے جے اوہ ترتیباں وی ہون، تاں اوہناں نوں \emph{سخت ترتیباں \pnbLatin{(strict orders)}}
آکھدے نیں۔
\end{ex}

\begin{explain}
بھاویں ترتیباں تے اکو ہون دا تعلق اہم تے قدرتی تعلق نیں، پر ایس گل اُتے
زور دینا چاہیدا اے کہ ساڈی تعریف مطابق $A^{2}$ دا \emph{کوئی وی} ذیلی سیٹ،
$A$ اُتے اک تعلق اے، بھاویں اوہ کِنّا وی غیر فطری یا گھڑیا ہویا لگے۔ خاص
طور تے، $\emptyset$ کسے وی سیٹ اُتے اک تعلق اے (\emph{خالی تعلق \pnbLatin{(empty relation)}}، جیہڑا
عنصراں دے کسے وی جوڑے وچکار نہیں ہوندا)، تے خود $A^{2}$ وی $A$ اُتے اک
تعلق اے (جیہڑا ہر جوڑے وچکار ہوندا اے)، جیہنوں \emph{ہمہ گیر تعلق \pnbLatin{(universal relation)}} آکھدے
نیں۔ پر $E=\Setabs{\tuple{n, m}}{n>5 \text{ یا } m \times n \ge 34}$ ورگی
شے وی اک تعلق گنی جاندی اے۔
\end{explain}
  
\begin{prob}
  سیٹ اُتے تعلق $\subseteq$ دے !!{element}s دی فہرست بناؤ، جتھے سیٹ
  $\Pow{\{a, b, c\}}$ اے۔
\end{prob}
\endgroup
% /OLP-0012

% OLP-0013 content/sets-functions-relations/relations/reflections.tex
\begingroup
\def\olfilename{reflections}
\expandafter\def\csname ol@sectioncs\endcsname{section}

\olfileid{sfr}{rel}{ref}
\olsection{فلسفیانہ سوچ بچار \pnbLatin{(Philosophical Reflections)}}

\olref[set]{sec} وچ اسی تعلقاں دی تعریف کجھ خاص سیٹاں دے طور تے کیتی سی۔
سانوں ذرا رک کے اک فلسفیانہ سوال پچھنا چاہیدا اے: ایہو جیہی تعریف
\emph{کر کی رہی اے}؟ ایس گل اُتے بہت شک اے کہ اسی ایہ کہنا چاہواں گے کہ
اسی مابعد الطبیعیاتی یکسانیت \pnbLatin{(metaphysical identity)} دے کجھ حقائق \emph{لبھ لَئے} نیں؛ مثال دے
طور تے، کہ $\Nat$ اُتے ترتیبی تعلق اوہ سیٹ $R=  \Setabs{\tuple{n,m}}{n, m \in \Nat\text{ تے }
n < m}$ \emph{نکلیا} جیہدی تعریف اسی \olref[set]{sec} وچ کیتی سی۔
ایہ گل کیوں مشکوک اے، ایہدیاں تِن وجہاں نیں۔

پہلی: \olref[set][pai]{wienerkuratowski} وچ اسی $\tuple{a, b} =
\{\{a\}, \{a, b\}\}$ دی تعریف کیتی سی۔ ایہدی تھاں ایہ تعریف ویکھو:
$\lVert a, b\rVert = \{\{b\}, \{a, b\}\} = \tuple{b,a}$۔ جدوں $a \neq b$ ہووے،
تاں $\tuple{a, b} \neq \lVert a,b\rVert$ ہوندا اے۔ پر اسی $\lVert a,b\rVert$
نوں وی ترتیب وار جوڑے دی اپنی تعریف منّ سکدے ساں، بجائے $\tuple{a,b}$ دے۔
دونویں تعریفاں اکو جِنّیاں چنگیاں طرح کم کردیاں۔ سو ہن ساڈے کول قدرتی عدداں
اُتے ترتیبی تعلق ``ہون'' دے دو اکو جِنّے چنگے امیدوار نیں، یعنی:
\begin{align*}
		R &= \Setabs{\tuple{n,m}}{n, m \in \Nat \text{ تے }n < m}\\
		S &= \Setabs{\lVert n,m\rVert}{n, m \in \Nat \text{ تے }n < m}.
\end{align*}
عنصراں نال تعیّن دے اصول دی بنا اُتے $R \neq S$ اے؛ سو صاف اے کہ اوہ
\emph{دونویں} $\Nat$ اُتے ترتیبی تعلق دے نال اکو نہیں ہو سکدے۔ پر ایہ
دعویٰ کرنا کہ دراصل $R$، نہ کہ $S$ (یا ایہدے اُلٹ)، ترتیبی تعلق
\emph{اے}، محض من مرضی دی گل ہووے گی، تے ایس لئی کجھ شرمندگی والی وی۔
(ایہ سیٹ نظریاتی تحویلیت \pnbLatin{(set-theoretic reductionism)} دے خلاف اوس دلیل دی اک بہت سادہ مثال اے جیہنوں
بیناسیراف (Benacerraf) نے \citeyear{Benacerraf1965} وچ مشہور کیتا سی۔ اسی ایہنوں کئی
واری مُڑ ویکھاں گے۔)

دوجی: جے اسی سمجھدے آں کہ \emph{ہر} تعلق نوں اک سیٹ دے نال اکو منّنا
چاہیدا اے، تاں خود سیٹ وچ رکن ہون دا تعلق \pnbLatin{(set-membership relation)}، $\in$، اک خاص سیٹ ہونا چاہیدا
اے۔ دراصل، اوہنوں سیٹ $\Setabs{\tuple{x,y}}{x \in y}$ ہونا پئے گا۔ پر کی
ایہ سیٹ موجود اے؟ رسل دے پیراڈوکس \pnbLatin{(Russell's Paradox)} نوں سامنے رکھدیاں، ایہ کوئی معمولی دعویٰ
نہیں کہ ایہو جیہا سیٹ موجود اے۔ درحقیقت،
\oliflabeldef{cumul:::part}{سیٹاں دا اوہ نظریہ جیہڑا اسی
\olref[cumul][][]{part} وچ تیار کراں گے، ایس سیٹ دی ہوند دا \emph{انکار}
کرے گا۔\footnote{ذرا اگے دی گل کر کے وجہ ویکھ لئیے۔ خلف توں ثبوت لئی فرض
کرو کہ $I = \Setabs{\tuple{x,y}}{x \in y}$ موجود اے۔ فیر $\bigcup \bigcup I$
ہمہ گیر سیٹ اے، جیہڑا \olref[sfr][z][sep]{thm:NoUniversalSet} دے خلاف اے۔}}{سیٹ
نظریے نوں اک مسلّمی نظریے \pnbLatin{(axiomatic theory)} دے طور تے پوری سختی نال تیار کرنا ممکن اے، تے
اوہ نظریہ واقعی ایس سیٹ دی ہوند دا انکار کرے گا۔}
سو بھاویں کجھ تعلقاں نوں سیٹ منّ کے کم چلایا جا سکدا اے، سیٹ وچ رکن ہون
دے تعلق نوں اک خاص صورت منّنا پئے گا۔

تیجی: جدوں اسی تعلقاں نوں سیٹاں دے نال ``اکو'' منّیا، تاں اسی آکھیا سی کہ
اسی $Rxy$ نوں $\tuple{x,y} \in R$ دی تھاں لکھن دی اجازت دے دیاں گے۔ ایہ
ٹھیک اے، بشرطیکہ سیٹ وچ رکن ہون دے تعلق \pnbLatin{(set-membership relation)}، ``$\in$''، نوں محمول (predicate) \emph{دے طور تے}
لیا جاوے۔ پر جے اسی سمجھدے آں کہ ``$\in$'' کسے خاص قسم دے سیٹ دا ناں اے،
تاں عبارت ``$\tuple{x,y} \in R$'' صرف تِن واحدی اصطلاحاں \pnbLatin{(singular terms)} توں بنی اے
جیہڑیاں سیٹاں نوں ناں دیندیاں نیں: ``$\tuple{x,y}$''، ``$\in$'' تے ``$R$''۔ تے
ناواں دی ایہو جیہی فہرست اک قضیہ (proposition) بیان کرن جوگی اوہدے نالوں ودھ نہیں جِنّی
ایہ بے معنی لڑی اے: ``پیالی قلم دان میز''۔ فیر اوہی گل اے: بھاویں کجھ
تعلقاں نوں سیٹ منّ کے کم چلایا جا سکدا اے، سیٹ دا رکن ہون دے تعلق نوں اک
خاص صورت منّنا پئے گا۔ (ایہ فریگے دے \emph{تصور گھوڑا} پیراڈوکس \pnbLatin{(concept horse paradox)} دی اک
سادہ صورت تے وٹگنسٹائن (Wittgenstein) دے رسل اُتے کیتے اک مشہور اعتراض نوں اکٹھا کردا اے۔)

سو ہن اسی کیہڑے نتیجے اُتے پہنچے آں؟ عدداں اُتے تعلقاں نوں سیٹ آکھن وچ
کوئی \emph{غلطی} نہیں۔ بس سانوں سمجھنا چاہیدا اے کہ ایہ گل کس بھاوَ نال
آکھی جاندی اے۔ اسی مابعد الطبیعیاتی یکسانیت دا کوئی حقیقتی دعویٰ نہیں کر
رہے۔ اسی صرف ایہ نوٹ کر رہے آں کہ کجھ سیاقاں وچ اسی (کجھ) تعلقاں نوں کجھ
خاص سیٹاں دے طور تے \emph{ورت سکدے آں} (تے ایس طرح ورت کے کم کراں گے)۔
\endgroup
% /OLP-0013

% OLP-0014 content/sets-functions-relations/relations/special-properties.tex
\begingroup
\def\olfilename{special-properties}
\expandafter\def\csname ol@sectioncs\endcsname{section}

\olfileid{sfr}{rel}{prp}
\olsection{تعلقاں دیاں خاص خاصیتاں \pnbLatin{(Special Properties of Relations)}}

\begin{intro}
تعلقاں دیاں کجھ قسماں اینیاں عام نیں کہ اوہناں نوں خاص ناں دتے گئے نیں۔
مثال دے طور تے، $\le$ تے $\subseteq$ دونویں اپنے اپنے دائرے دے عنصراں نوں
رلدے ملدے طریقیاں نال جوڑدے نیں (مثلاً $\Nat$، جدوں تعلق $\le$ ہووے،
تے $\Pow{A}$، جدوں تعلق $\subseteq$ ہووے)۔ ایہ ٹھیک ٹھیک سمجھن لئی کہ
ایہ تعلق کِویں رلدے ملدے نیں تے کِویں وکھرے نیں، اسی اوہناں دی ونڈ
کجھ خاص خاصیتاں دے حساب نال کردے آں جیہڑیاں تعلقاں وچ ہو سکدیاں نیں۔
ایہناں وچوں کجھ خاصیتاں (دے میل) خاص اہم نکلدے نیں: ترتیباں (orders) تے
ہم ارزیت دے تعلق \pnbLatin{(equivalence relations)}۔
\end{intro}

\begin{defn}[انعکاسیت (Reflexivity)]
اک تعلق $R \subseteq A^2$ اُدوں تے صرف اُدوں \emph{انعکاسی (reflexive)} ہوندا اے جدوں
ہر $x \in A$ لئی $Rxx$ ہووے۔
\end{defn}

\begin{defn}[تعدیت (Transitivity)]
اک تعلق $R \subseteq A^2$ اُدوں تے صرف اُدوں \emph{متعدی (transitive)} ہوندا اے جدوں،
جد وی $Rxy$ تے $Ryz$ ہون، تاں $Rxz$ وی ہووے۔
\end{defn}

\begin{defn}[تناظر (Symmetry)]
اک تعلق $R \subseteq A^2$ اُدوں تے صرف اُدوں \emph{متناظر (symmetric)} ہوندا اے جدوں،
جد وی $Rxy$ ہووے، تاں $Ryx$ وی ہووے۔
\end{defn}

\begin{defn}[ضد تناظر (Anti-symmetry)]
اک تعلق $R \subseteq A^2$ اُدوں تے صرف اُدوں \emph{ضد تناظری (anti-symmetric)} ہوندا اے جدوں،
جد وی $Rxy$ تے $Ryx$ دونویں ہون، تاں $x=y$ ہووے (یا ہور لفظاں وچ:
جے $x\neq y$ ہووے، تاں یا $\lnot Rxy$ یا $\lnot Ryx$ ہووے)۔
\end{defn}

\begin{explain}
متناظر تعلق وچ $Rxy$ تے $Ryx$ ہمیشہ اکٹھے درست ہوندے نیں، یا دونویں
درست نہیں ہوندے۔ ضد تناظری تعلق وچ $Rxy$ تے $Ryx$ اکٹھے صرف اوہدوں
درست ہو سکدے نیں جدوں $x = y$ ہووے۔ خیال رکھو کہ ایہ اوہناں دا
اکٹھے درست ہونا \emph{لازم} نہیں کردا: $Rxy$ تے $Ryx$ دا $x = y$ ہون اُتے
درست ہونا صرف خارج نہیں کیتا گیا۔ سو ضد تناظری تعلق انعکاسی ہو سکدا
اے، پر ایہ گل نہیں کہ ہر ضد تناظری تعلق انعکاسی ہووے۔ ایہ وی یاد رکھو
کہ ضد تناظری ہونا تے صرف متناظر نہ ہونا دو وکھریاں شرطاں نیں۔ درحقیقت،
اک تعلق اکو ویلے متناظر تے ضد تناظری دونویں ہو سکدا اے (مثلاً اکو ہون
دا تعلق ایہو جیہا اے)۔
\end{explain}

\begin{defn}[باہمی موازنہ پذیری (Connectivity)]
اک تعلق $R \subseteq A^2$ \emph{باہم قابلِ موازنہ (connected)} ہوندا اے جے ہر $x,y\in
A$ لئی، جے $x \neq y$، تاں یا $Rxy$ یا $Ryx$ ہووے۔
\end{defn}

\begin{prob}
ایہو جہے تعلقاں دیاں مثالاں دیو جیہڑے (الف) انعکاسی تے متناظر ہون پر
متعدی نہ ہون، (ب) انعکاسی تے ضد تناظری ہون، (ج) ضد تناظری تے متعدی
ہون پر انعکاسی نہ ہون، تے (د) انعکاسی، متناظر تے متعدی ہون۔ عدداں یا
سیٹاں اُتے تعلق نہ ورتو۔
\end{prob}
 
\begin{defn}[اپنے آپ نال تعلق نہ ہونا (Irreflexivity)]
اک تعلق $R \subseteq A^2$ \emph{اپنے آپ نال تعلق نہ ہون والا (irreflexive)} آکھیا جاندا
اے جے ہر $x \in A$ لئی $Rxx$ نہ ہووے۔
\end{defn}

\begin{defn}[یک طرفہ ہونا (Asymmetry)]
اک تعلق $R \subseteq A^2$ \emph{یک طرفہ (asymmetric)} آکھیا جاندا اے جے کسے وی جوڑے
$x,y\in A$ لئی $Rxy$ تے $Ryx$ دونویں اکٹھے نہ ہون۔
\end{defn}

خیال رکھو کہ جے $A \neq \emptyset$ ہووے، تاں $A$ اُتے اپنے آپ نال تعلق
نہ ہون والا کوئی وی تعلق انعکاسی نہیں ہوندا، تے $A$ اُتے ہر یک طرفہ تعلق
ضد تناظری وی ہوندا اے۔ پر ایہو جہے $R \subseteq A^2$ موجود نیں جیہڑے
انعکاسی وی نہیں تے ہر عنصر دا اپنے آپ نال بے تعلق ہونا وی لازم نہیں کردے؛ تے ایہو جہے ضد تناظری
تعلق وی نیں جیہڑے یک طرفہ نہیں۔
\endgroup
% /OLP-0014

% OLP-0015 content/sets-functions-relations/relations/equivalence-relations.tex
\begingroup
\def\olfilename{equivalence-relations}
\expandafter\def\csname ol@sectioncs\endcsname{section}

\olfileid{sfr}{rel}{eqv}

\olsection{ہم ارزیت دے تعلق \pnbLatin{(Equivalence Relations)}}

اک سیٹ اُتے اکو ہون دا تعلق \pnbLatin{(identity relation)} انعکاسی (reflexive)،
متناظر (symmetric) تے متعدی (transitive) ہوندا اے۔ جیہڑے تعلقاں $R$
وچ ایہ تِنّے خاصیتاں ہون، اوہ بہت عام نیں۔

\begin{defn}[ہم ارزیت دا تعلق \pnbLatin{(Equivalence relation)}]
اک تعلق $R \subseteq A^2$ نوں، جے اوہ انعکاسی، متناظر تے متعدی ہووے،
\emph{ہم ارزیت دا تعلق} آکھیا جاندا اے۔ !!^{element}s $x$ تے $y$، جیہڑے
$A$ وچ نیں، \emph{$R$ دے لحاظ نال ہم ارز (R-equivalent)} آکھے جاندے
نیں جے $Rxy$ ہووے۔
\end{defn}

ہم ارزیت دے تعلقاں توں \emph{ہم ارزیت دے طبقے \pnbLatin{(equivalence class)}} دا تصور
نکلدا اے۔ ہم ارزیت دا تعلق دائرۂ تعریف (domain) نوں اک ونڈ (partition)
دے وکھ وکھ حصیاں وچ ``ونڈ دیندا'' اے۔ ونڈ دے ہر حصے اندر ساریاں چیزاں
اک دوجی نال تعلق رکھدیاں نیں؛ تے وکھرے حصیاں دیاں چیزاں آپس وچ تعلق
نہیں رکھدیاں۔ کدی کدی ایہناں حصیاں بارے \emph{سدھے طور تے} گل کرنا فائدے
مند ہوندا اے۔ ایس مقصد لئی اسی اک تعریف دیندے آں:

\begin{defn}\ollabel{def:equivalenceclass}
فرض کرو کہ $R \subseteq A^2$ اک ہم ارزیت دا تعلق اے۔ ہر $x \in A$ لئی،
$x$ دا $A$ وچ \emph{ہم ارزیت دا طبقہ} اوہ سیٹ $\equivrep{x}{R}
= \Setabs{y \in A}{Rxy}$ اے۔ $A$ دا، $R$ دے تحت، \emph{خارج قسمت سیٹ
(quotient)}، $\equivclass{A}{R} = \Setabs{\equivrep{x}{R}}{x \in A}$،
ایہناں ہم ارزیت دے طبقیاں دا سیٹ اے۔
\end{defn}

اگلا نتیجہ ہم ارزیت دے طبقے دی تعریف نوں درست ٹھہراؤندا اے: ایہ ثابت
کردا اے کہ ہم ارزیت دے طبقے واقعی $A$ دی ونڈ دے حصے نیں:

\begin{prop}
جے $R \subseteq A^2$ ہم ارزیت دا تعلق ہووے، تاں $Rxy$ اُدوں تے صرف اُدوں
جدوں $\equivrep{x}{R} = \equivrep{y}{R}$۔
\end{prop}

\begin{proof}
کھبے توں سجے پاسے دی سمت لئی، فرض کرو کہ $Rxy$، تے $z \in
\equivrep{x}{R}$ لو۔ تعریف مطابق فیر $Rxz$ اے۔ کیوں جو $R$ ہم ارزیت دا
تعلق اے، ایس لئی $Ryz$ اے۔ (کھول کے دسیے: کیوں جو $Rxy$ اے تے $R$
متناظر اے، ساڈے کول $Ryx$ اے؛ تے کیوں جو $Rxz$ اے تے $R$ متعدی اے،
ساڈے کول $Ryz$ اے۔) سو $z \in \equivrep{y}{R}$۔ ایس گل نوں عام کر کے،
$\equivrep{x}{R} \subseteq \equivrep{y}{R}$۔ پر بالکل ایسے طرح
$\equivrep{y}{R} \subseteq \equivrep{x}{R}$ وی اے۔ سو عنصراں نال تعیّن دے اصول
(Extensionality) مطابق، $\equivrep{x}{R} = \equivrep{y}{R}$۔

سجے توں کھبے پاسے دی سمت لئی، فرض کرو کہ $\equivrep{x}{R} =
\equivrep{y}{R}$۔ کیوں جو $R$ انعکاسی اے، ایس لئی $Ryy$، سو $y \in
\equivrep{y}{R}$۔ ایس توں $y \in \equivrep{x}{R}$ وی اے، کیوں جو اسی
$\equivrep{x}{R} = \equivrep{y}{R}$ فرض کیتا اے۔ سو $Rxy$۔
\end{proof}

\begin{ex}
ہم ارزیت دے تعلقاں دی اک چنگی مثال باقی اُتے مبنی حساب \pnbLatin{(modular arithmetic)}
توں ملدی اے۔
کسے وی $a$، $b$ تے $n \in \PosInt$ لئی، اسی $a \equiv_n b$ اُدوں تے
صرف اُدوں آکھدے آں جدوں $a$ نوں $n$ نال ونڈن توں اوہی باقی بچے جیہڑا
$b$ نوں $n$ نال ونڈن توں بچدا اے۔ (کجھ ہور علامتاں نال: $a \equiv_n b$
اُدوں تے صرف اُدوں جدوں کسے $k \in \Int$ لئی $a - b = kn$ ہووے۔) ہن
$\equiv_n$ ہر $n$ لئی ہم ارزیت دا تعلق اے۔ تے ٹھیک $n$ وکھرے ہم ارزیت دے
طبقے تعلق $\equiv_n$ توں بن دے نیں؛ یعنی $\equivclass{\Nat}{\equiv_n}$ دے
$n$ !!{element}s نیں۔ ایہ اوہ نیں: اوہناں عدداں دا سیٹ جیہڑے $n$ نال
بغیر باقی دے ونڈے جان، یعنی $\equivrep{0}{\equiv_n}$؛ اوہناں عدداں دا سیٹ
جیہناں نوں $n$ نال ونڈن اُتے باقی $1$ بچے، یعنی $\equivrep{1}{\equiv_n}$؛
\ldots؛ تے اوہناں عدداں دا سیٹ جیہناں نوں $n$ نال ونڈن اُتے باقی $n-1$
بچے، یعنی $\equivrep{n-1}{\equiv_n}$۔
\end{ex}

\begin{prob}
ثابت کرو کہ $\equiv_n$ ہر $n \in \PosInt$ لئی ہم ارزیت دا تعلق اے، تے
$\equivclass{\Nat}{\equiv_n}$ دے ٹھیک $n$ رکن نیں۔
\end{prob}
\endgroup
% /OLP-0015

% OLP-0016 content/sets-functions-relations/relations/orders.tex
\begingroup
\def\olfilename{orders}
\expandafter\def\csname ol@sectioncs\endcsname{section}

\olfileid{sfr}{rel}{ord}
\olsection{ترتیباں \pnbLatin{(Orders)}}

\begin{explain}
ساڈے کئی موازنے کجھ چیزاں نوں، کسے خاص لحاظ نال، ہور چیزاں توں
«گھٹ»، اوہناں دے «برابر»، یا اوہناں توں «ودھ» دسن نال تعلق رکھدے نیں۔
ایہناں وچ \emph{ترتیبی تعلق \pnbLatin{(order relations)}} آؤندے نیں۔ پر ترتیبی تعلقاں دیاں وکھ وکھ
قسماں نیں۔ مثال دے طور تے، کجھ وچ ضروری اے کہ کوئی وی دو چیزاں آپس وچ
قابلِ موازنہ ہون، پر کجھ وچ نہیں۔ کجھ اکو ہون (identity) نوں شامل کردے نیں
(جیویں~$\le$) تے کجھ اوہنوں کڈھ دیندے نیں (جیویں~$<$)۔
ایتھے ایہناں دی درجہ بندی (taxonomy) ساڈے کم آوے گی۔
\end{explain}

\begin{defn}[پیش ترتیب (Preorder)]
اک تعلق نوں، جے اوہ انعکاسی (reflexive) تے متعدی (transitive) دونویں ہووے،
\emph{پیش ترتیب (preorder)} آکھیا جاندا اے۔
\end{defn}

\begin{defn}[جزوی ترتیب \pnbLatin{(Partial order)}]
اک پیش ترتیب جیہڑی ضد تناظری (anti-symmetric) وی ہووے، \emph{جزوی ترتیب
\pnbLatin{(partial order)}} آکھلاندی اے۔
\end{defn}

\begin{defn}[خطی ترتیب \pnbLatin{(Linear order)}]\ollabel{def:linearorder}
اک جزوی ترتیب جیہڑی باہم قابلِ موازنہ (connected) وی ہووے، \emph{کُل ترتیب
\pnbLatin{(total order)}} یا \emph{خطی ترتیب \pnbLatin{(linear order)}} آکھلاندی اے۔
\end{defn}

\begin{ex}
ہر خطی ترتیب جزوی ترتیب وی ہوندی اے، تے ہر جزوی ترتیب پیش ترتیب وی
ہوندی اے، پر اُلٹ گلاں درست نہیں۔ $A$ اُتے ہمہ گیر تعلق \pnbLatin{(universal relation)} پیش ترتیب اے، کیوں
جے اوہ انعکاسی تے متعدی اے۔ پر جے $A$ وچ اک توں ودھ !!{element} ہووے،
تاں ہمہ گیر تعلق ضد تناظری نہیں ہوندا، تے ایس لئی
جزوی ترتیب وی نہیں۔
\end{ex}

\begin{ex}
\emph{لمبائی وچ ودھ نہ ہون \pnbLatin{(no longer than)}} دے تعلق $\preccurlyeq$ اُتے غور
کرو، جیہڑا $\Bin^*$ اُتے اے: $x \preccurlyeq y$ اُدوں تے صرف اُدوں جدوں
$\len{x} \le \len{y}$۔
ایہ پیش ترتیب اے (انعکاسی تے متعدی)، بلکہ باہم قابلِ موازنہ وی اے، پر
جزوی ترتیب نہیں، کیوں جے ایہ ضد تناظری نہیں۔ مثال لئی، $01
\preccurlyeq 10$ تے $10 \preccurlyeq 01$، پر $01 \neq 10$۔
\end{ex}

\begin{ex}
اک اہم جزوی ترتیب، سیٹاں دے کسے سیٹ اُتے تعلق $\subseteq$ اے۔ ایہ عام
طور تے خطی ترتیب نہیں، کیوں جے جے $a \neq b$ ہووے تے اسی
$\Pow{\{a, b\}} = \{\emptyset, \{a\}, \{b\}, \{a,b\}\}$ اُتے غور کریے،
تاں ویکھدے آں کہ $\{a\} \nsubseteq \{b\}$ تے $\{a\} \neq \{b\}$ تے
$\{b\} \nsubseteq \{a\}$۔
\end{ex}

\begin{ex}
\emph{بغیر باقی پورا ونڈن \pnbLatin{(divisibility without remainder)}} دا تعلق سانوں اک ایہو جیہی جزوی ترتیب
دیندا اے جیہڑی خطی ترتیب نہیں۔ صحیح عدداں $n$ تے~$m$ لئی، اسی
$n \mid m$ ایہ مطلب دسن لئی لکھدے آں کہ $n$، عدد $m$ نوں پورا ونڈدا
اے، یعنی اُدوں تے صرف اُدوں جدوں کوئی صحیح عدد~$k$ ایہو جیہا ہووے کہ
$m = kn$۔ $\Nat$ اُتے ایہ جزوی ترتیب اے، پر خطی ترتیب نہیں: مثال لئی،
$2 \nmid 3$ تے نالے $3 \nmid 2$۔ جے $\Int$ اُتے تعلق دے طور تے ویکھیے،
تاں پورا ونڈن دا تعلق صرف پیش ترتیب اے، کیوں جے ایہ ضد تناظری نہیں:
$1 \mid -1$ تے $-1 \mid 1$، پر $1 \neq -1$۔
\end{ex}

\begin{ex}
سلسلیاں دے سیٹ~$A^*$ اُتے \emph{ودھاؤ (extension)} دا تعلق ایہ اے: $s \sqsubseteq s'$
اُدوں تے صرف اُدوں جدوں $s = \emptyseq$ (خالی سلسلہ)، $s = s'$، یا
$s = \tuple{s_1, \dots, s_n}$ تے $s' =
\tuple{s_1, \dots, s_n, s_{n+1}, \dots, s_m}$ ہووے۔ جے $s \sqsubseteq s'$
ہووے، تاں اسی ایہ وی آکھدے آں کہ $s$، سلسلے~$s'$ دا \emph{مُڈھلا ٹکڑا
\pnbLatin{(initial segment)}}
اے۔ $A^*$ اُتے ودھاؤ دا تعلق جزوی ترتیب اے پر خطی ترتیب نہیں؛ مثال لئی،
جے $a \neq b$ ہووے، تاں $ab \not\sqsubseteq ba$ تے $ba \not\sqsubseteq ab$۔
\end{ex}

\begin{defn}[سخت ترتیب \pnbLatin{(Strict order)}]
\emph{سخت ترتیب \pnbLatin{(strict order)}} اک اوہ تعلق اے جیہڑا اپنے آپ نال تعلق
نہ ہون والا (irreflexive)، یک طرفہ (asymmetric)، تے متعدی (transitive) ہوندا اے۔
\end{defn}

\begin{defn}[سخت خطی ترتیب \pnbLatin{(Strict linear order)}]\ollabel{def:strictlinearorder}
اک سخت ترتیب جیہڑی باہم قابلِ موازنہ وی ہووے، \emph{سخت کُل ترتیب
\pnbLatin{(strict total order)}} یا \emph{سخت خطی ترتیب \pnbLatin{(strict linear order)}} آکھلاندی اے۔
\end{defn}

\begin{ex}
$\le$ اوہ خطی ترتیب اے جیہڑی سخت خطی ترتیب~$<$ دے مطابق اے۔
$\subseteq$ اوہ جزوی ترتیب اے جیہڑی سخت ترتیب~$\subsetneq$ دے مطابق اے۔
\end{ex}

کوئی وی سخت ترتیب $R$ جیہڑی~$A$ اُتے ہووے، اوہدے وچ قطر (diagonal) $\Id{A}$ شامل
کر کے، یعنی سارے جوڑے~$\tuple{x, x}$ شامل کر کے، جزوی ترتیب وچ بدلی
جا سکدی اے۔ (ایہنوں $R$ دی \emph{انعکاسی تکمیل \pnbLatin{(reflexive closure)}} آکھدے نیں۔) اُلٹ پاسے،
جزوی ترتیب توں شروع کر کے، $\Id{A}$ کڈھ کے سخت ترتیب حاصل کیتی جا
سکدی اے۔ اگلے دو نتیجے ایہ گل ٹھیک ٹھیک بیان کردے نیں۔

\begin{prop}\ollabel{prop:stricttopartial}
جے $R$، سیٹ~$A$ اُتے سخت ترتیب اے، تاں $R^+ = R \cup \Id{A}$ جزوی
ترتیب اے۔ نالے، جے $R$ سخت خطی ترتیب اے، تاں $R^+$ خطی ترتیب اے۔
\end{prop}

\begin{proof}
فرض کرو کہ $R$ سخت ترتیب اے، یعنی $R \subseteq A^2$ تے $R$ اپنے آپ نال
تعلق نہ ہون والا، یک طرفہ، تے متعدی اے۔ رکھو کہ $R^+ = R \cup \Id{A}$۔
سانوں ثابت کرنا اے کہ $R^+$ انعکاسی، ضد تناظری، تے متعدی اے۔

$R^+$ صاف طور تے انعکاسی اے، کیوں جے $\tuple{x, x} \in \Id{A} \subseteq
R^+$ ہر $x \in A$ لئی درست اے۔

ایہ ثابت کرن لئی کہ $R^+$ ضد تناظری اے، برہانِ خلف (reductio) لئی فرض کرو کہ $R^+xy$
تے $R^+yx$، پر $x \neq y$۔ کیوں جے $\tuple{x,y} \in R \cup \Id{A}$، پر
$\tuple{x, y} \notin \Id{A}$، ایس لئی لازماً $\tuple{x, y} \in R$، یعنی
$Rxy$۔ ایس طرح~$Ryx$ وی۔ پر ایہ ایس مفروضے دے اُلٹ اے کہ $R$ یک طرفہ اے۔

تعدیت ثابت کرن لئی، من لو کہ $R^+xy$ تے $R^+yz$۔ جے $\tuple{x, y} \in R$
تے $\tuple{y,z} \in R$ دونویں ہون، تاں $\tuple{x, z} \in R$، کیوں جے
$R$~متعدی اے۔ نہیں تاں، یا $\tuple{x, y} \in \Id{A}$، یعنی $x = y$، یا
$\tuple{y, z} \in \Id{A}$، یعنی $y = z$۔ پہلی صورت وچ، مفروضے مطابق
$R^+yz$، تے $x = y$، سو $R^+xz$۔ دوسری صورت وچ وی ایس طرح۔ ہر صورت وچ
$R^+xz$، سو $R^+$ متعدی وی اے۔

«نالے» والی گل لئی، فرض کرو کہ $R$ باہم قابلِ موازنہ وی اے۔ سو ہر
$x \neq y$ لئی، یا $Rxy$ یا~$Ryx$، یعنی یا $\tuple{x, y} \in R$ یا
$\tuple{y, x} \in R$۔ کیوں جے $R \subseteq R^+$، ایہ گل $R^+$ لئی وی
درست رہندی اے، سو $R^+$ باہم قابلِ موازنہ وی اے۔
\end{proof}

\begin{prop}\ollabel{prop:partialtostrict}
جے $R$، سیٹ~$A$ اُتے جزوی ترتیب اے، تاں $R^- = R \setminus \Id{A}$ سخت
ترتیب اے۔ نالے، جے $R$ خطی ترتیب اے، تاں $R^-$ سخت خطی ترتیب اے۔
\end{prop}

\begin{proof}
ایہ ثبوت مشق دے طور تے چھڈّیا جاندا اے۔
\end{proof}

\begin{prob}
\olref[sfr][rel][ord]{prop:partialtostrict} دا ثبوت دیو۔
\end{prob}

اگلا سادہ نتیجہ ایہ ثابت کردا اے کہ سخت خطی ترتیباں، عنصراں نال تعیّن دے اصول
(Extensionality) ورگی اک خاصیت پوری کردیاں نیں:

\begin{prop}\ollabel{prop:extensionality-strictlinearorders}
جے $<$، سیٹ $A$ اُتے سخت خطی ترتیب اے، تاں:
\[
  (\forall a, b \in A)((\forall x \in A)(x < a \liff x < b) \lif a = b).
\]
\end{prop}

\begin{proof}
من لو کہ $(\forall x \in A)(x < a \liff x < b)$۔ جے $a < b$، تاں
$a < a$، جیہڑا ایس گل دے اُلٹ اے کہ $<$ اپنے آپ نال تعلق نہ ہون والا
اے؛ سو $a \nless b$۔ عین ایس طرح $b \nless a$۔ سو $a = b$، کیوں جے
$<$ باہم قابلِ موازنہ اے۔
\end{proof}
\endgroup
% /OLP-0016

% OLP-0017 content/sets-functions-relations/relations/graphs.tex
\begingroup
\def\olfilename{graphs}
\expandafter\def\csname ol@sectioncs\endcsname{section}

\olfileid{sfr}{rel}{grp}
\olsection{گراف \pnbLatin{(Graphs)}}

\emph{گراف} اک خاکہ اے جیہدے وچ نقطے---جنہاں نوں «گرہاں» (nodes) یا
«راساں» (vertices؛ واحد vertex نوں «راس») آکھدے نیں---کناریاں (edges)
نال جڑے ہوندے نیں۔ گراف گسستہ ریاضی \pnbLatin{(discrete mathematics)} تے کمپیوٹر
سائنس وچ تھاں تھاں ورتے جان والے اوزار نیں۔ ایہ
تعلقاں تے بناوٹاں دی نمائندگی کرن تے اوہناں نوں اکھاں اگے لیاؤن لئی
بہت کم آؤندے نیں؛ مختلف قسماں دے نیٹ ورکاں ورگیاں ٹھوس چیزاں توں لے کے
فیصلیاں دے ممکنہ نتیجیاں ورگیاں مجرد بناوٹاں تیکر۔ علمی لکھتاں وچ
گرافاں دیاں بہت وکھ وکھ قسماں ملدیاں نیں۔ اوہ، مثال دے طور تے، ایہناں
گلاں وچ وکھریاں ہوندیاں نیں کہ کناریاں دی سمت اے یا نہیں، اوہناں اُتے
لیبل نیں یا نہیں، کسے گرہ توں اوسے گرہ ول کنارا ہو سکدا اے یا نہیں،
اوہناں ای گرہاں دے وچکار اک توں ودھ کنارے ہو سکدے نیں یا نہیں، وغیرہ۔
\emph{سمت دار گرافاں} \pnbLatin{(directed graphs)} دا تعلقاں نال خاص جوڑ اے۔

\begin{defn}[سمت دار گراف \pnbLatin{(Directed graph)}]
اک \emph{سمت دار گراف} \pnbLatin{(directed graph)} $G = \tuple{V, E}$،
\emph{راساں} (vertices) دے اک سیٹ~$V$ تے \emph{کناریاں} (edges) دے اک
سیٹ~$E \subseteq V^2$ توں بندا اے۔
\end{defn}

\begin{explain}
ساڈی تعریف مطابق، گراف بس اک سیٹ تے اوس سیٹ اُتے اک تعلق اے۔ بے شک،
جدوں گرافاں دی گل کریے، تاں ایہ توقع فطری اے کہ اوہناں نوں خاکے دی
صورت وچ وکھایا جاوے: اسی دو راساں~$v_1$ تے $v_2$ نوں تیر نال اُدوں
تے صرف اُدوں جوڑ کے گراف بنا سکدے آں جدوں $\tuple{v_1, v_2} \in E$۔
اپنے آپ وچ اک تعلق تے اک گراف وچ فرق صرف ایہ اے کہ گراف راساں دا
سیٹ وی دسدا اے، یعنی گراف وچ الگ تھلگ راس \pnbLatin{(isolated vertices)} وی ہو
سکدے نیں۔ پر اہم گل
ایہ اے کہ ہر تعلق~$R$ جیہڑا کسے سیٹ~$X$ اُتے ہووے، سمت دار گراف
$\tuple{X, R}$ دے طور تے ویکھیا جا سکدا اے؛ تے اُلٹ پاسے، سمت دار
گراف~$\tuple{V, E}$ نوں تعلق $E \subseteq V^2$ دے طور تے ویکھیا جا سکدا
اے، جیہدے نال سیٹ $V$ صاف طور تے دسیا ہووے۔
\end{explain}

\begin{ex}
گراف $\tuple{V, E}$، جیہدے لئی $V = \{1, 2, 3, 4\}$ تے $E =
\{\tuple{1,1}, \allowbreak \tuple{1, 2}, \allowbreak \tuple{1, 3},
\allowbreak \tuple{2, 3}\}$، ایویں دکھائی دیندا اے:
\begin{align*}
& \begin{tikzpicture}[->,node distance=2cm]
  \node[draw,circle] (A) {$1$};
  \node[draw,circle] (B) [right of=A] {$2$};
  \node[draw,circle] (C) [below of=B] {$3$};
  \node[draw,circle] (D) [right of=B] {$4$};
  \draw (A) to [loop above]  (A);
  \draw (A) to  (B);
  \draw (A) to  (C);
  \draw (B) to  (C);
  \end{tikzpicture}
\intertext{ایہ گراف، گراف $\tuple{V', E}$ توں وکھرا اے جیہدے لئی $V' =
  \{1, 2, 3\}$؛ اوہ ایویں دکھائی دیندا اے:}
& \begin{tikzpicture}[->,node distance=2cm]
  \node[draw,circle] (A) {$1$};
  \node[draw,circle] (B) [right of=A] {$2$};
  \node[draw,circle] (C) [below of=B] {$3$};
  \draw (A) to [loop above]  (A);
  \draw (A) to  (B);
  \draw (A) to  (C);
  \draw (B) to  (C);
  \end{tikzpicture}
\end{align*}
\end{ex}

\begin{prob}
  گھٹ یا برابر ہون دے تعلق~$\le$ نوں، سیٹ $\{1, 2, 3, 4\}$ اُتے، گراف
  دے طور تے ویکھو تے اوہدا خاکہ بناؤ۔
\end{prob}
\endgroup
% /OLP-0017

% OLP-0018 content/sets-functions-relations/relations/trees.tex
\begingroup
\def\olfilename{trees}
\expandafter\def\csname ol@sectioncs\endcsname{section}

\olfileid{sfr}{rel}{tre}
\olsection{درخت \pnbLatin{(Trees)}}

جزوی ترتیب \pnbLatin{(partial order)} دی اک خاص قسم، جیہڑی منطق دے سارے حصیاں وچ
اہم کم کردی اے، \emph{درخت} (tree) اے۔ محدود درخت \pnbLatin{(finite trees)} منطق
دے ابتدائی حصیاں وچ آؤندے نیں: مثال لئی، !!{formula}s نوں اوہناں دے
نحوی درخت \pnbLatin{(syntax tree)} وچ ٹوٹن دے لحاظ نال سمجھیا جا سکدا اے، جد کہ
!!{derivation}s (derivations) کئی !!{derivation} نظاماں \pnbLatin{(derivation systems)} وچ وی محدود درختاں دی صورت اختیار کردے نیں۔
%
لامتناہی درخت \pnbLatin{(infinite trees)} قضیاتی منطق \pnbLatin{(propositional logic)} تے پہلے
درجے دی منطق \pnbLatin{(first-order logic)} دے تمامیت دے قضیاں \pnbLatin{(completeness theorems)} دے ثبوتاں وچ ای آ جاندے نیں، تے ریاضیاتی منطق \pnbLatin{(mathematical logic)} وچ تھاں تھاں ورتے جاندے نیں۔

سیٹ نظریاتی تصور \pnbLatin{(set-theoretic concept)} دے طور تے درخت، گراف نظریے
\pnbLatin{(graph theory)} دے درخت نال گہرا تعلق رکھدا اے۔ اک (محدود) درخت دی
تصویر ایہ اے:

\begin{center}
\begin{tikzpicture}[nodes={draw, circle}, -]
\node{$r$} [grow'=up]
    child { node {$a$} 
        child { node {$c$} }
        child { node {$d$} }
        child { node {$e$} }
    }
    child { node {$b$} };
\end{tikzpicture}
\end{center}

سب توں تھلّے والی گرہ (node)~$r$ جڑ (root) اے۔ $r$ توں ہٹ کے ہر گرہ
دی عین تھلّے صرف اک والد گرہ \pnbLatin{(parent node)} ہوندی اے۔ اسی گرہ~$x$ دا
گرہ~$y$ نال اوہ تعلق، جیہڑا اُدوں ہووے جدوں $y$ تیکر~$x$ توں کناریاں
(edges) دے نال نال اُتّے نوں جا کے پہنچ سکّیے، ایویں سمجھ سکدے آں کہ
$x$، گرہ~$y$ دا \emph{پُرکھ} (ancestor) اے۔

درخت وچ پُرکھ ہون دا تعلق سخت جزوی ترتیب \pnbLatin{(strict partial order)} اے۔ ایہ
گل سیٹ نظریے والی تعریف دا محرک اے۔ اوہ تعریف دسن لئی سانوں دو تصور
چاہیدے نیں۔ کسے سیٹ~$A$ وچ، جیہنوں~$\le$ نال جزوی ترتیب دتی گئی ہووے،
\emph{سب توں چھوٹا عنصر} \pnbLatin{(least element)} اوہ !!a{element} $x \in A$
اے جیہدے لئی ہر $y \in A$ اُتے~$x \le y$ ہووے۔ اک سیٹ، تعلق~$\le$ دے
لحاظ نال \emph{حُسنِ ترتیب والا} (well-ordered) ہوندا اے جے اوہدے ہر
غیر خالی ذیلی سیٹ دا سب توں چھوٹا عنصر ہووے۔

\begin{defn}[درخت (Tree)]
اک \emph{درخت} (tree) جوڑا $T = \tuple{A, \le}$ اے جیہدے وچ $A$ اک
سیٹ تے $\le$، سیٹ~$A$ اُتے جزوی ترتیب \pnbLatin{(partial order)} ہووے، تے اک
اکلوتا سب توں چھوٹا عنصر \pnbLatin{(unique least element)} $r \in A$ (جیہنوں
\emph{جڑ} (root) آکھدے نیں) ہووے، ایس طرح کہ ہر $x \in A$ لئی، سیٹ
$\Setabs{y}{y \le x}$، تعلق~$\le$ دے لحاظ نال حُسنِ ترتیب والا
(well-ordered) ہووے۔
\end{defn}

\begin{defn}[جانشین (Successors)]
من لو کہ $T = \tuple{A, \le}$ اک درخت اے۔ جے $x,y \in A$، $x < y$،
تے کوئی وی $z \in A$ ایہو جیہا نہ ہووے کہ $x < z < y$، تاں اسی آکھدے
آں کہ $y$، گرہ~$x$ دا \emph{جانشین} (successor) اے۔
\end{defn}

گرہ $x \in A$ دے جانشین اوہدے \emph{بچے} (children) وی آکھلاندے نیں۔
جے $y$، گرہ~$x$ دا جانشین اے، تاں اسی $x$ نوں گرہ~$y$ دا
\emph{پیشرو} (predecessor) یا \emph{والد گرہ} (parent) آکھدے آں۔

\begin{prop}
جے $\tuple{A,\le}$ اک درخت اے، تاں جڑ توں ہٹ کے ہر $x \in A$ دا ودھ
توں ودھ اک پیشرو ہوندا اے۔
\end{prop}

\begin{proof}
  من لو کہ $y_1 < x$ تے $y_2 < x$ تے $y_1 \neq y_2$۔ تاں $\{y_1,
  y_2\} \subseteq \Setabs{z}{z<x}$۔ کیوں جے $\Setabs{z}{z<x}$،
  تعلق~$\le$ دے لحاظ نال حُسنِ ترتیب والا اے، اوہدے ذیلی سیٹ
  $\{y_1, y_2\}$ دا سب توں چھوٹا عنصر اے، جیہڑا صاف طور تے یا $y_1$
  یا~$y_2$ ہونا چاہیدا اے۔ سو یا $y_1 \le y_2$ یا $y_2 \le y_1$۔ اسی
  منّیا سی کہ $y_1 \neq y_2$، سو اصل وچ یا $y_1 < y_2$ یا $y_2 < y_1$۔
  کیوں جے اسی منّیا سی کہ $y_1 < x$ تے $y_2 < x$، ایس لئی نالے یا
  $y_1 < y_2 < x$ یا $y_2 < y_1 < x$۔ سو $y_1$ تے $y_2$ دونویں،
  گرہ~$x$ دے پیشرو نہیں ہو سکدے۔
\end{proof}

\begin{defn}
درخت $T = \tuple{A, \le}$ \emph{لامتناہی} (infinite) آکھلاندا اے جے
$A$ لامتناہی سیٹ ہووے، تے نہیں تاں \emph{محدود} (finite)۔ جے $T$ ایہو
جیہا ہووے کہ ہر $x \in A$ دے صرف محدود گنتی وچ جانشین ہون، تاں اسی
آکھدے آں کہ $T$ \emph{ہر راس دے محدود جانشیناں والا} \pnbLatin{(finitely branching)} اے۔
\end{defn}

\begin{defn}[شاخاں (Branches)]
دتے ہوئے درخت $T = \tuple{A, \le}$ دی \emph{شاخ} (branch)، $T$ وچ اک
ہور نہ ودھ سکن والی زنجیر \pnbLatin{(maximal chain)} اے، یعنی $T$ وچ اک سیٹ
$B \subseteq A$ جیہڑا ایہو جیہا ہووے کہ ہر $x, y \in B$ لئی یا
$x \le y$ یا $y \le x$، تے ہر $z \in A \setminus B$ لئی کوئی
$u \in B$ ایہو جیہا موجود ہووے کہ نہ $z \le u$ تے نہ $u \le z$۔
ماخذ وچ ایتھے غیر متعین X دی تھاں درخت دا حامل سیٹ A ورتیا گیا اے
(PNBREL-001)۔
%
اسی $[T]$ نال $T$ دیاں ساریاں شاخاں دا سیٹ ظاہر کردے آں۔
\end{defn}

\begin{ex}
ہر راس دے محدود جانشیناں والے \pnbLatin{(finitely branching)} درخت دی اک مشہور
مثال، $0$ تے~$1$ دے محدود سلسلیاں دا \emph{لامتناہی دو شاخی درخت}
\pnbLatin{(infinite binary tree)} اے، جیہنوں کدی $\{0,1\}^*$ یا~$\Bin^*$ نال
ظاہر کردے نیں، تے ودھاؤ دے تعلق \pnbLatin{(extension relation)} $\sqsubseteq$
نال ترتیب دیندے نیں (مثلاً $101 \sqsubseteq 101101$)۔ کیوں جے کوئی وی
دو علامتی لڑی \pnbLatin{(binary string)} آخر اُتے $0$ یا $1$ لا کے ہمیشہ ودھائی
جا سکدی اے، ایس درخت وچ لامتناہی عنصر نیں: ہر عنصر~$s$ دے عین دو
جانشین، $s0$ تے~$s1$، نیں۔ اوہدی جڑ خالی سلسلہ \pnbLatin{(empty sequence)}
$\emptyseq$ اے۔
\end{ex}

\begin{ex}
تھوڑا ہور عام طور تے، قدرتی عدداں دے محدود سلسلیاں دا سیٹ~$\Nat^*$،
ودھاؤ دے تعلق~$\sqsubseteq$ نال، وی اک درخت اے۔ صاف ظاہر اے کہ ایہ ہر
راس دے محدود جانشیناں والا \pnbLatin{(finitely branching)} نہیں: ہر $s \in \Nat^*$ دے لامتناہی
جانشین~$sn$ نیں، ہر $n \in \Nat$ لئی اک۔ ہر $A \subseteq \Nat^*$
جیہڑا~$\sqsubseteq$ دے تحت بند ہووے، $\Nat^*$ دا \emph{ذیلی درخت}
(subtree) اے۔
(یعنی $A$ ایہو جیہا اے کہ جے $s \in A$ تے $s' \sqsubseteq s$، تاں
$s' \in A$ وی۔) سارے محدود درخت $\Nat^*$ دے محدود ذیلی درختاں دے طور
تے ظاہر کیتے جا سکدے نیں۔
\end{ex}

\begin{prop}[K\H{o}nig دا معاون قضیہ (König's lemma)]
جے $T = \tuple{A,\le}$، ہر راس دے محدود جانشیناں والا \pnbLatin{(finitely branching)} لامتناہی درخت اے،
تاں $T$ دی اک لامتناہی شاخ اے۔
\end{prop}

K\H{o}nig دے معاون قضیے دی اک خاص صورت، جیہڑی قابلِ حسابیت دے نظریے
وچ بہت ورتی جاندی اے تے \emph{K\H{o}nig دا کمزور معاون قضیہ} (weak
König's lemma) آکھلاندی اے، ایہ اے: $\{0,1\}^*$ دے ہر لامتناہی ذیلی
درخت دی اک لامتناہی شاخ اے۔
\endgroup
% /OLP-0018

% OLP-0019 content/sets-functions-relations/relations/operations.tex
\begingroup
\def\olfilename{operations}
\expandafter\def\csname ol@sectioncs\endcsname{section}

\olfileid{sfr}{rel}{ops}
\olsection{تعلقاں اُتے عمل \pnbLatin{(Operations on Relations)}}

تعلقاں نوں بدلنا یا رلانا اکثر کم آؤندا اے۔
\olref[sfr][rel][ord]{prop:stricttopartial} وچ اسی تعلقاں دے \emph{اتحاد}
(union) اُتے غور کیتا سی؛ ایہ بس دو تعلقاں نوں جوڑیاں (pairs) دے سیٹ
سمجھ کے اوہناں
دا اتحاد لینا اے۔ ایس طرح \olref[sfr][rel][ord]{prop:partialtostrict}
وچ اسی تعلقاں دے نسبتی فرق \pnbLatin{(relative difference)} اُتے غور کیتا سی۔ ہُن
تعلقاں اُتے کیتے جا سکن والے کجھ ہور عمل (operations) ایہ نیں۔

\begin{defn}\ollabel{relationoperations} 
من لو کہ $R$، $S$ تعلق نیں، تے $A$ کوئی وی سیٹ اے۔

$R$ دا \emph{اُلٹا تعلق} (inverse)، $R^{-1} = \Setabs{\tuple{y, x}}{\tuple{x,
    y} \in R}$ اے۔

$R$ تے $S$ دا \emph{نسبتی حاصل ضرب} \pnbLatin{(relative product)}، $(R \mid S) =
\{\tuple{x, z} : \exists y(Rxy \land Syz)\}$ اے۔

$R$ دی $A$ تک \emph{تحدید} (restriction)، $\funrestrictionto{R}{A}= R
\cap A^2$ اے۔

$R$ دا $A$ اُتے \emph{اطلاق} (application)، $\funimage{R}{A} = \{y :
(\exists x \in A)Rxy\}$ اے۔
\end{defn}

\begin{ex}
من لو کہ $S \subseteq \Int^2$، سیٹ~$\Int$ اُتے جانشین ہون دا تعلق
\pnbLatin{(successor relation)} اے،
یعنی $S = \Setabs{\tuple{x, y} \in \Int^2}{x + 1 = y}$، سو $Sxy$ اُدوں
تے صرف اُدوں جدوں $x + 1 = y$۔

$S^{-1}$، سیٹ $\Int$ اُتے پیشرو ہون دا تعلق \pnbLatin{(predecessor relation)} اے، یعنی
$\Setabs{\tuple{x,y}\in\Int^2}{x -1 =y}$۔

$S\mid S$ ایہ اے:
$ \Setabs{\tuple{x,y}\in\Int^2}{x + 2 =y}$

$\funrestrictionto{S}{\Nat}$، سیٹ~$\Nat$ اُتے جانشین ہون دا تعلق
\pnbLatin{(successor relation)} اے۔

$\funimage{S}{\{1,2,3\}}$، سیٹ $\{2, 3, 4\}$ اے۔
\end{ex}

\begin{defn}[تعدی تکمیل \pnbLatin{(Transitive closure)}]من لو کہ $R \subseteq A^2$
دو رکنی تعلق \pnbLatin{(binary relation)} اے۔

تعلق~$R$ دی \emph{تعدی تکمیل} \pnbLatin{(transitive closure)}، $R^+ = \bigcup_{0 < n \in
\Nat} R^n$ اے، جتھے اسی بازگشتی طور تے (recursively) $R^1 = R$ تے $R^{n+1} = R^n
\mid R$ تعریف کردے آں۔

$R$ دی \emph{انعکاسی تعدی تکمیل} \pnbLatin{(reflexive transitive closure)}، $R^* = R^+ \cup
\Id{A}$ اے۔
نوٹ: ایتھے R+ اوپر والی مقامی تعریف دے مطابق تعدی تکمیل اے؛ ترتیب والے
پچھلے حصے وچ ایہو علامت انعکاسی تکمیل لئی ورتی گئی سی۔
\end{defn}

\begin{ex}
جانشین ہون دا تعلق $S \subseteq \Int^2$ لو۔ $S^2xy$ اُدوں تے صرف
اُدوں جدوں $x + 2 = y$، $S^3xy$ اُدوں تے صرف اُدوں جدوں $x + 3 = y$،
وغیرہ۔ سو $S^+xy$ اُدوں تے صرف اُدوں جدوں $x + n = y$ کسے
$n \geq 1$ لئی ہووے۔ ہور لفظاں وچ، $S^+xy$ اُدوں تے صرف اُدوں جدوں
$x < y$، تے $S^*xy$ اُدوں تے صرف اُدوں جدوں $x \le y$۔
\end{ex}

\begin{prob}
وکھاؤ کہ $R$ دی تعدی تکمیل \pnbLatin{(transitive closure)} واقعی متعدی (transitive) اے۔
\end{prob}
\endgroup
% /OLP-0019

\OLEndChapterHook



\par\noindent\emph{متبادل ترتیب۔} تعلق دا مُڈھلا باب (\texttt{relations.tex})۔ ایہ ترتیب تدوین جوگ ماخذ فائل وچ موجود اے۔
\endgroup
% /OLP-0011

% OLP-0020 content/sets-functions-relations/functions/functions.tex
\begingroup
\def\olfilename{functions}
\expandafter\def\csname ol@sectioncs\endcsname{section}

\olchapter{sfr}{fun}{فنکشن \pnbLatin{(Functions)}}

% OLP-0021 content/sets-functions-relations/functions/function-basics.tex
\begingroup
\def\olfilename{function-basics}
\expandafter\def\csname ol@sectioncs\endcsname{section}

\olfileid{sfr}{fun}{bas}
\olsection{بنیادی گلاں \pnbLatin{(Basics)}}

\begin{explain}
اک \emph{فنکشن} (function) ایہو جیہی نقشہ بندی (mapping) اے جیہڑی دتے ہوئے سیٹ (set) دے ہر
!!{element} نوں کسے (ہور) دتے ہوئے سیٹ دے اک متعین !!{element} ول
گھلّدی اے۔ مثال لئی، $1$ جمع کرن دا عمل اک فنکشن (function) تعریف کردا اے: ہر
عدد~$n$ نوں اک اکلوتے عدد~$n+1$ نال جوڑیا جاندا اے۔

ہور عام طور تے، فنکشن (function) جوڑے (pairs)، تِکے (triples) وغیرہ انپٹ (inputs) دے طور تے لے سکدے نیں
تے کسے قسم دا آؤٹ پٹ (output) دیندے نیں۔ کئی فنکشن سانوں بنیادی حساب توں
واقف نیں۔ مثال لئی، جمع تے ضرب فنکشن نیں۔ اوہ دو عدد لیندے نیں تے
تیجا عدد دیندے نیں۔

ایس ریاضیاتی، مجرد معنی \pnbLatin{(abstract sense)} وچ، فنکشن (function) اک \emph{کالا ڈبّا} \pnbLatin{(black box)} اے: اہم صرف
ایہ اے کہ کیہڑا آؤٹ پٹ (output) کیہڑے انپٹ (input) نال جوڑیا ہویا اے، نہ کہ آؤٹ پٹ
کڈھن دا طریقہ۔
\end{explain}

\begin{defn}[فنکشن (Function)]
اک \emph{فنکشن} (function) $f \colon A \to B$، سیٹ (set)~$A$ دے ہر !!{element} دی
سیٹ~$B$ دے اک !!{element} ول نقشہ بندی (mapping) اے۔

اسی $A$ نوں فنکشن~$f$ دا \emph{دائرۂ تعریف} (domain) تے $B$ نوں فنکشن~$f$ دا
\emph{ہدف سیٹ} (codomain) آکھدے آں۔ سیٹ~$A$ دے !!{element}s، فنکشن~$f$ دے انپٹ
(inputs) یا \emph{دلیلاں} (arguments) آکھلاندے نیں؛ تے سیٹ~$B$ دا اوہ !!{element} جیہنوں
دلیل~$x$ نال فنکشن~$f$ جوڑدا اے، \emph{فنکشن~$f$ دی قدر} (value) آکھلاندا اے
جدوں دلیل~$x$ ہووے، تے اوہنوں~$f(x)$ لکھدے آں۔

\emph{حاصل قدراں دا سیٹ} (range) $\ran{f}$، فنکشن~$f$ دے ہدف سیٹ دا اوہ ذیلی
سیٹ (subset) اے جیہدے وچ فنکشن~$f$ دیاں اوہ قدراں (values) شامل نیں جیہڑیاں کسے دلیل
لئی حاصل ہوندیاں نیں؛ $\ran{f} = \Setabs{f(x)}{x \in A}$۔
\end{defn}

\olref{fig:function} والا خاکہ فنکشناں بارے سوچن وچ مدد دے سکدا اے۔
کھبے پاسے دا بیضہ نما دائرہ فنکشن دا \emph{دائرۂ تعریف} (domain) ظاہر کردا اے؛
سجے پاسے دا بیضہ نما دائرہ فنکشن دا \emph{ہدف سیٹ} (codomain) ظاہر کردا اے؛ تے
تیر دائرۂ تعریف دی اک \emph{دلیل} (argument) توں ہدف سیٹ دی اوہدے نال جڑی
\emph{قدر} (value) ول اشارہ کردا اے۔

\begin{figure}
  \olasset{assets/diagrams/function.tikz}
  \caption{فنکشن (function) اک سیٹ (set) دے ہر !!{element} دی ہور سیٹ دے !!a{element}
    ول نقشہ بندی (mapping) اے۔ تیر دائرۂ تعریف (domain) دی اک دلیل (argument) توں ہدف سیٹ (codomain) دی
    اوہدے نال جڑی قدر (value) ول اشارہ کردا اے۔}
  \ollabel{fig:function}
\end{figure}

\begin{ex}
ضرب قدرتی عدداں \pnbLatin{(natural numbers)} دے جوڑے (pairs) انپٹ (inputs) دے طور تے لیندی اے تے اوہناں نوں
قدرتی عدداں نال آؤٹ پٹ (outputs) دے طور تے جوڑدی اے، سو ایہ $\Nat \times \Nat$
(دائرۂ تعریف (domain)) توں $\Nat$ (ہدف سیٹ (codomain)) ول جاندی اے۔ حاصل قدراں دا سیٹ (range) وی
$\Nat$ ای نکلدا اے، کیوں جے ہر $n \in \Nat$، حاصل ضرب $n \times 1$ اے۔
\end{ex}

\begin{ex}
ضرب اک فنکشن (function) اے، کیوں جے ایہ ہر انپٹ (input)---قدرتی عدداں دے ہر جوڑے---نوں
اکو آؤٹ پٹ (output) نال جوڑدی اے: $\times \colon \Nat^2 \to \Nat$۔ ایس دے
اُلٹ، قدرتی عدد~$n$ نوں ایہو جہے حقیقی عدد \pnbLatin{(real number)}~$x$ نال جوڑنا کہ $x^2 = n$
ہووے، فنکشن نہیں بناؤندا، کیوں جے ہر مثبت صحیح عدد \pnbLatin{(positive integer)}~$n$ دے دو مربع
جذر نیں: $\sqrt{n}$ تے~$-\sqrt{n}$۔ صرف غیر منفی (اصل) مربع جذر \pnbLatin{(nonnegative principal square root)} واپس
دے کے اسی ایہنوں فنکشن بنا سکدے آں: $\sqrt{\phantom{X}} \colon \Nat \to \Real$۔
% Source correction OLFUN-002: English says "positive" although Nat includes 0.
\end{ex}

\begin{ex}
اوہ تعلق (relation) جیہڑا جماعت دے ہر طالب علم نوں اوہدے آخری گریڈ نال جوڑدا
اے، فنکشن (function) اے---اکو جماعت وچ کسے طالب علم نوں دو وکھ وکھ آخری گریڈ
نہیں مل سکدے۔ اوہ تعلق جیہڑا جماعت دے ہر طالب علم نوں اوہدے ماں پیو
نال جوڑدا اے، فنکشن نہیں: طالب علماں دے صفر، یا دو، یا دو توں ودھ
ماں پیو ہو سکدے نیں۔
\end{ex}

\begin{explain}
اسی ہر ممکن دلیل (argument) لئی فنکشن دی قدر (value) نوں کسے ٹھیک ٹھیک طریقے نال دس
کے فنکشن (function) تعریف کر سکدے آں۔ ایہ کرن دے وکھ وکھ طریقے نیں: فارمولا
دینا، قدر کڈھن دا طریقہ بیان کرنا، یا ہر دلیل لئی قدر دی فہرست (list) دینا۔
فنکشن نوں جِویں وی تعریف کریے، سانوں ایہ یقینی بناؤنا چاہیدا اے کہ
ہر دلیل لئی اک، تے صرف اک، قدر دسی گئی ہووے۔
\end{explain}


\begin{ex}
من لو کہ $f \colon \Nat \to \Nat$ ایویں تعریف کیتا گیا اے کہ
$f(x) = x+1$۔ ایہ تعریف $f$ نوں ایہو جہے فنکشن دے طور تے دسدی اے
جیہڑا قدرتی عدد \pnbLatin{(natural number)} لیندا اے تے قدرتی عدد آؤٹ پٹ (output) دے طور تے دیندا اے۔
ایہ سانوں دسدی اے کہ قدرتی عدد~$x$ دتے جان اُتے، $f$ اوہدا
جانشین (successor)~$x+1$ دے گا۔ ایس صورت وچ، ہدف سیٹ (codomain) $\Nat$، فنکشن~$f$ دا حاصل
قدراں دا سیٹ (range) نہیں، کیوں جے قدرتی عدد~$0$ کسے وی قدرتی عدد دا جانشین
نہیں۔ فنکشن~$f$ دا حاصل قدراں دا سیٹ سارے مثبت صحیح عدداں \pnbLatin{(positive integers)} دا سیٹ،
$\Int^{+}$، اے۔
\end{ex}

\begin{ex}\ollabel{examplefunext}
من لو کہ $g \colon \Nat \to \Nat$ ایویں تعریف کیتا گیا اے کہ
$g(x) = x+2-1$۔ ایہ سانوں دسدا اے کہ $g$ اک فنکشن اے جیہڑا قدرتی
عدد \pnbLatin{(natural number)} لیندا اے تے قدرتی عدد دیندا اے۔ قدرتی عدد~$x$ دتے جان اُتے، $g$
عدد~$x$ دے جانشین (successor) دے جانشین دا پیشرو (predecessor)، یعنی $x+1$، دے گا۔
% Source correction OLFUN-003: English switches from input n to x; target consistently uses x.
\end{ex}

\begin{explain}
اسی ہُنے دو فنکشناں (functions)، $f$ تے $g$، اُتے غور کیتا اے جنہاں دیاں
\emph{تعریفاں} وکھریاں نیں۔ پر ایہ \emph{اکو فنکشن} نیں۔ آخر، ہر
قدرتی عدد~$n$ لئی ساڈے کول $f(n) = n+1 = n+2-1 = g(n)$ اے۔ ہور
لفظاں وچ، $f$ تے~$g$ دیاں ساڈیاں تعریفاں وکھ وکھ مساواتاں راہیں
اکو نقشہ بندی (mapping) دسدیاں نیں۔ سو اسی ضمنی طور تے فنکشناں لئی عنصراں نال
تعیّن دے اک اصول \pnbLatin{(principle of extensionality)} اُتے بھروسا کر رہے آں:
\[
  \text{جے }\forall x\, f(x) = g(x)\text{، تاں }f = g
\]
بشرطیکہ $f$ تے~$g$ دا دائرۂ تعریف تے ہدف سیٹ اکو ہون۔
\end{explain}

\begin{ex}
اسی صورتاں دے حساب \pnbLatin{(by cases)} نال وی فنکشن (function) تعریف کر سکدے آں۔ مثال لئی، اسی
$h \colon \Nat \to \Nat$ نوں ایویں تعریف کر سکدے آں:
\[
h(x) =
\begin{cases}
  \frac{x}{2} & \text{جے $x$ جفت (even) اے} \\
  \frac{x+1}{2} & \text{جے $x$ طاق (odd) اے۔}
\end{cases}
\]
کیوں جے ہر قدرتی عدد \pnbLatin{(natural number)} یا جفت اے یا طاق، ایس فنکشن دا آؤٹ پٹ (output) ہمیشہ
قدرتی عدد ہووے گا۔ بس ایہ یاد رکھو کہ جے تسی صورتاں دے حساب نال
فنکشن تعریف کرو، تاں ہر ممکن انپٹ (input) عین اک صورت وچ آؤنا چاہیدا اے۔
کجھ موقعیاں اُتے ایس لئی ثبوت دینا پئے گا کہ صورتاں سارے امکان
گھیردیاں (exhaustive) نیں تے آپس وچ اکٹھیاں لاگو نہیں ہوندیاں (exclusive)۔
\end{ex}
\endgroup
% /OLP-0021

% OLP-0022 content/sets-functions-relations/functions/function-kinds.tex
\begingroup
\def\olfilename{function-kinds}
\expandafter\def\csname ol@sectioncs\endcsname{section}

\olfileid{sfr}{fun}{kin}
\olsection{فنکشناں دیاں قسماں \pnbLatin{(Kinds of Functions)}}

\begin{explain}
ساڈے کم آوے گا جے اسی عام طور تے ملن والیاں فنکشناں (functions) دیاں کجھ
قسماں دی درجہ بندی متعارف کرا دیئے۔

شروع وچ اسی اوہناں فنکشناں اُتے غور کر سکدے آں جنہاں دے ہدف سیٹ
(codomain) دا ہر رکن فنکشن دی اک قدر ہوندا اے۔ ایہو جہے فنکشناں نوں
!!{surjective} (ہر ہدف قدر لین والے؛ surjective) آکھیا جاندا اے، تے
\olref{fig:surjective} وچ دتے خاکے وانگ وکھایا جا سکدا اے۔

\begin{figure}
  \olasset{assets/diagrams/surjective.tikz}
  \caption{!!^a{surjective} فنکشن دے ہدف سیٹ (codomain) دا ہر !!{element}
    (element/member) اوہدی
    اک قدر ہوندا اے۔}
  \ollabel{fig:surjective}
\end{figure}
\end{explain}

\begin{defn}[!!^{surjective} فنکشن (ہر ہدف قدر لین والا فنکشن)]
اک فنکشن $f \colon A \rightarrow B$ \emph{!!{surjective}} عین اوس
ویلے اے جدوں $B$ وی~$f$ دا حاصل قدراں دا سیٹ (range) ہووے، یعنی ہر
$y \in B$ لئی گھٹ توں گھٹ اک $x \in A$ ہووے جیہدے لئی~$f(x) = y$؛
علامتاں وچ:
\[
  (\forall y \in B)(\exists x \in A)f(x) = y.
\]
اسی ایہو جہے فنکشن نوں $A$ توں $B$ ول !!a{surjection} (سرجیکشن)
آکھدے آں۔
\end{defn}

\begin{explain}
جے تسی وکھاؤنا چاہندے او کہ $f$ اک !!a{surjection} اے، تاں ایہ
وکھاؤنا پئے گا کہ $f$ دے ہدف سیٹ (codomain) دی ہر شے، $f(x)$ دی قدر
(output/value) اے، کسے انپٹ (input) $x$ لئی۔

خیال رکھو کہ ہر فنکشن اک !!a{surjection} (سرجیکشن) \emph{پیدا کردا} اے۔ آخر،
دتے ہوئے فنکشن $f \colon A \to B$ لئی $f' \colon A \to \ran{f}$ نوں
$f'(x) = f(x)$ راہیں تعریف کرو۔ کیوں جے $\ran{f}$ نوں
\emph{تعریف ای} $\Setabs{f(x) \in B}{x \in A}$ کیتا گیا اے، ایس
لئی فنکشن $f'$ دا !!a{surjection} (سرجیکشن) ہونا یقینی اے۔
\end{explain}

\begin{explain}
ہُن ہر فنکشن ہر ممکن انپٹ (input) نوں اک اکلوتے آؤٹ پٹ (output) نال جوڑدا اے۔ پر
ایہو جہے فنکشن وی نیں جیہڑے دو وکھ وکھ انپٹاں نوں کدے اکو آؤٹ پٹ
نال نہیں جوڑدے۔ ایہو جہے فنکشناں نوں !!{injective} (اک بہ اک؛ injective)
آکھیا جاندا اے،
تے \olref{fig:injective} وچ دتے خاکے وانگ وکھایا جا سکدا اے۔
\begin{figure}
  \olasset{assets/diagrams/injective.tikz}
  \caption{!!^a{injective} فنکشن (اک بہ اک فنکشن) دو وکھ وکھ دلیلاں نوں کدے اکو قدر
    نال نہیں جوڑدا۔}
  \ollabel{fig:injective}
\end{figure}
\end{explain}

\begin{defn}[!!^{injective} فنکشن (اک بہ اک فنکشن)]
اک فنکشن $f \colon A \rightarrow B$ \emph{!!{injective}} عین اوس
ویلے اے جدوں ہر $y \in B$ لئی ودھ توں ودھ اک $x \in A$ ہووے جیہدے
لئی~$f(x) = y$۔ اسی ایہو جہے فنکشن نوں $A$ توں~$B$ ول
!!a{injection} (انجیکشن) آکھدے آں۔
\end{defn}

\begin{explain}
جے تسی وکھاؤنا چاہندے او کہ $f$ اک !!a{injection} (انجیکشن) اے، تاں ایہ
وکھاؤنا پئے گا کہ کنہاں وی !!{element}s (elements) $x$ تے $y$ لئی، جیہڑے
$f$ دے دائرۂ تعریف (domain) وچ نیں، جے $f(x)=f(y)$ ہووے، تاں $x=y$ ہووے۔
\end{explain}

\begin{ex}
مستقل فنکشن \pnbLatin{(constant function)} $f\colon \Nat \to \Nat$، جیہڑا $f(x) = 1$ نال دتا
گیا اے، نہ !!{injective} (injective) اے، نہ !!{surjective} (surjective)۔

شناختی فنکشن \pnbLatin{(identity function)} $f\colon \Nat \to \Nat$، جیہڑا $f(x) = x$ نال دتا
گیا اے، !!{injective} وی اے تے !!{surjective} وی۔

جانشین فنکشن \pnbLatin{(successor function)} $f \colon \Nat \to \Nat$، جیہڑا $f(x) = x+1$ نال
دتا گیا اے، !!{injective} اے پر !!{surjective} نہیں۔

ایویں صورتاں دے حساب نال تعریف کیتا فنکشن \pnbLatin{(function defined by cases)}
$f \colon \Nat \to \Nat$:
\[
  f(x) =
  \begin{cases}
    \frac{x}{2} & \text{جے $x$ جفت اے} \\
    \frac{x+1}{2} & \text{جے $x$ طاق اے۔}
  \end{cases}
\]
!!{surjective} (surjective) اے، پر !!{injective} (injective) نہیں۔
\end{ex}

\begin{explain}
اکثر سانوں اوہناں فنکشناں اُتے غور کرنا ہوندا اے جیہڑے
!!{injective} وی ہون تے !!{surjective} وی۔ اسی ایہو جہے فنکشناں نوں
!!{bijective} (اک بہ اک تے ہر ہدف قدر لین والے؛ bijective) آکھدے آں۔ اوہ
\olref{fig:bijective} وچ وکھائے فنکشن
وانگ ہوندے نیں۔ !!^{bijection}s نوں کدی کدی
\emph{اک بہ اک مطابقتاں \pnbLatin{(one-to-one correspondences)}} وی آکھیا جاندا اے، کیوں جے اوہ ہدف سیٹ (codomain) دے
عنصراں نوں دائرۂ تعریف دے عنصراں نال اکلوتے طور تے جوڑدے نیں۔
\begin{figure}
  \olasset{assets/diagrams/bijective.tikz}
  \caption{!!^a{bijective} فنکشن \pnbLatin{(bijective function)} ہدف سیٹ دے عنصراں نوں دائرۂ تعریف
    دے عنصراں نال اکلوتے طور تے جوڑدا اے۔}
  \ollabel{fig:bijective}
\end{figure}
\end{explain}

\begin{defn}[!!^{bijection} (بائجیکشن)]
اک فنکشن $f \colon A \to B$ \emph{!!{bijective}} عین اوس ویلے اے
جدوں اوہ !!{surjective} وی ہووے تے !!{injective} وی۔ اسی ایہو جہے
فنکشن نوں $A$ توں~$B$ ول (یا $A$ تے~$B$ دے وچکار) !!a{bijection}
(بائجیکشن) آکھدے آں۔
\end{defn}
\endgroup
% /OLP-0022

% OLP-0023 content/sets-functions-relations/functions/functions-relations.tex
\begingroup
\def\olfilename{functions-relations}
\expandafter\def\csname ol@sectioncs\endcsname{section}

\olfileid{sfr}{fun}{rel}

\olsection{فنکشن تعلقاں دے طور تے \pnbLatin{(Functions as Relations)}}

\begin{explain}
اوہ فنکشن (function) جیہڑا~$A$ دے !!{element}s (elements) نوں~$B$ دے
!!{element}s نال جوڑدا اے، ظاہر اے کہ $A$ تے~$B$ دے وچکار اک تعلق
(relation) تعریف کردا اے:
اوہی تعلق جیہڑا $x$ تے $y$ دے وچکار عین اوس ویلے قائم اے جدوں
$f(x) = y$ ہووے۔ اصل وچ، جے مثال لئی ساڈی دلچسپی ریاضی دے بنیادی
اجزا گھٹاؤن وچ ہووے، تاں اسی فنکشن~$f$ نوں ایس تعلق نال، یعنی جوڑیاں
(pairs) دے اک سیٹ نال، \emph{اکو سمجھ (identify)} وی سکدے آں۔ فیر ایہ سوال پیدا ہوندا
اے: کیہڑے تعلق ایس طریقے نال فنکشن تعریف کردے نیں؟
\end{explain}

\begin{defn}[فنکشن دا گراف \pnbLatin{(Graph of a function)}] من لو کہ $f\colon A \to B$ اک فنکشن اے۔
فنکشن~$f$ دا \emph{گراف (graph)} اوہ تعلق (relation) $R_f \subseteq A \times B$ اے
جیہڑا ایویں تعریف ہوندا اے:
\[
R_f = \Setabs{\tuple{x,y}}{f(x) = y}.
\]
\end{defn}

\begin{explain}
عنصراں نال تعیّن دے اصول (extensionality) کرکے فنکشن دا گراف اکلوتے طور تے طے ہوندا
اے۔ ہور ایہ کہ، سیٹاں اُتے عنصراں نال تعیّن دا اصول \pnbLatin{(extensionality on sets)} فنکشناں لئی
ضمنی اصول نوں فوراً ثابت کردا اے: جے $f$ تے~$g$ دا دائرۂ تعریف (domain) تے
ہدف سیٹ (codomain) سانجھا ہووے، تاں ہر قدر (value) اُتے متفق ہون دی صورت وچ اوہ اکو
فنکشن نیں۔

ایسے طرح، جے کوئی تعلق ``فنکشنی'' \pnbLatin{(functional relation)} ہووے، تاں اوہ
کسے فنکشن دا گراف اے۔
\end{explain}

\begin{prop}\ollabel{prop:graph-function}
من لو کہ $R \subseteq A \times B$ اک فنکشنی تعلق \pnbLatin{(functional relation)} ایس طرح اے کہ:
\begin{enumerate}
\item جے $Rxy$ تے $Rxz$، تاں $y = z$؛ تے
\item ہر $x \in A$ لئی کوئی $y \in B$ اے جیہدے لئی $\tuple{x,
y} \in R$۔
\end{enumerate}
فیر $R$ اوس فنکشن $f\colon A \to B$ دا گراف اے جیہڑا
$f(x) = y$ عین اوس ویلے جدوں $Rxy$ راہیں تعریف ہوندا اے۔
\end{prop}

\begin{proof}
من لو کہ کوئی $y$ اے جیہدے لئی $Rxy$۔ جے کوئی ہور $z \neq y$ وی
ہوندا جیہدے لئی $Rxz$، تاں~$R$ اُتے شرط دی خلاف ورزی ہوندی۔ ایس
لئی جے کوئی $y$ اے جیہدے لئی $Rxy$، تاں ایہ $y$ اکلوتا اے، تے ایس
واسطے $f$ صحیح طور تے تعریف شدہ (well-defined) اے۔ ظاہر اے کہ $R_f = R$۔
\end{proof}

\begin{explain}
ہر فنکشن $f\colon A \to B$ دا اک گراف \pnbLatin{(graph of a function)} ہوندا اے، یعنی دائرۂ تعریف
تے ہدف سیٹ دے وچکار اک تعلق (relation؛ برابر طور تے $A \times B$ دا اک ذیلی سیٹ) جیہڑا
$f(x) = y$ راہیں تعریف ہوندا اے۔ دوجے پاسے، ہر تعلق~$R
\subseteq A \times B$ جیہدے وچ \olref{prop:graph-function} والیاں
خاصیتاں ہون، اک فنکشن~$f \colon A \to B$ دا گراف اے۔ فنکشناں تے
اوہناں دے گرافاں دے ایس نیڑے رشتے کرکے اسی فنکشن نوں بس اوہدا گراف
سمجھ سکدے آں۔ ہور لفظاں وچ، فنکشناں نوں کجھ مخصوص تعلقاں، یعنی
ٹپلاں (tuples) دے کجھ مخصوص سیٹاں نال شناخت (identify) کیتا جا سکدا اے۔
% Source correction OLFUN-004: relation between A and B / subset of A times B, not a relation on A times B.
\oliflabeldef{sfr:rel:ref:sec}{پر ایہ یاد رکھو کہ ایس ``شناخت'' دی
روح \olref[sfr][rel][ref]{sec} والی اے: ایہ فنکشناں دی مابعد الطبیعی
حقیقت بارے دعویٰ نہیں، بلکہ ایہ مشاہدہ اے کہ فنکشناں نوں کجھ سیٹاں
دے طور تے \emph{ورتنا} سوکھا اے۔ ایس سہولت دی اک وجہ ایہ اے کہ }{}
ہُن اسی فنکشناں اُتے وی اوہو جہے عمل (operations) کر سکدے آں جیہڑے اسی تعلقاں
اُتے کیتے سن (ویکھو \olref[sfr][rel][ops]{sec})۔ خاص طور تے:
\end{explain}

\begin{defn}\ollabel{defn:funimage}
من لو کہ $f \colon A \to B$ اک فنکشن (function) اے تے $C\subseteq A$ اک ذیلی سیٹ (subset) اے۔

فنکشن~$f$ دی~$C$ تک \emph{تحدید (restriction)} اوہ
فنکشن~$\funrestrictionto{f}{C}\colon C \to B$ اے جیہڑا
$(\funrestrictionto{f}{C})(x) = f(x)$ راہیں ہر $x \in C$ لئی
تعریف ہوندا اے۔ ہور لفظاں وچ، $\funrestrictionto{f}{C} =
\Setabs{\tuple{x, y} \in R_f}{x \in C}$۔

فنکشن~$f$ دا~$C$ اُتے \emph{اطلاق (application)} $\funimage{f}{C} =
\Setabs{f(x)}{x \in C}$ اے۔ اسی ایہنوں~$C$ دا~$f$ دے تحت
\emph{عکس (image)} وی آکھدے آں۔
\end{defn}

\begin{explain}
ایہناں تعریفاں توں نکلدا اے کہ حاصل قدراں دا سیٹ (range) $\ran{f} =
\funimage{f}{\dom{f}}$، ہر فنکشن~$f$ لئی۔
\oliflabeldef{sfr:rel:ops:sec}{ایہ خیال تعلقاں نال فنکشناں دی شناخت
تے \olref[sfr][rel][ops]{sec} دیاں تعریفاں نال اوہو طرح جُڑدے نیں
جِویں توقع اے۔ پر اتھے فنکشن دی تحدید صرف دائرۂ تعریف نوں ذیلی سیٹ
تک گھٹاؤندی اے تے ہدف سیٹ نوں اوہو رکھدی اے؛ ایس لئی ایہ پہلے
دتے یکساں سیٹ اُتے تعلق دی دو پاسیاں والی تحدید دا حرف بہ حرف
اوہی عمل نہیں۔ دوجے دو عمل---اُلٹے (inverses) تے نسبتی حاصل ضرب
\pnbLatin{(relative products)}---لئی کجھ ہور
تفصیل دی لوڑ اے۔ اسی اوہ \olref[inv]{sec} تے
\olref[cmp]{sec} وچ دیواں گے۔}{}
% Source correction OLFUN-005: domain-only function restriction; earlier homogeneous relation restriction is two-sided.
\end{explain}
\endgroup
% /OLP-0023

% OLP-0024 content/sets-functions-relations/functions/inverses.tex
\begingroup
\def\olfilename{inverses}
\expandafter\def\csname ol@sectioncs\endcsname{section}

\olfileid{sfr}{fun}{inv}
\olsection{فنکشناں دے اُلٹ \pnbLatin{(Inverses of Functions)}}

\begin{explain}
اسی فنکشناں نوں نقشہ بندیاں (maps) سمجھدے آں۔ سو فنکشناں بارے اک
قدرتی سوال ایہ اے کہ کی نقشہ بندی (mapping) نوں ``اُلٹا'' چلایا جا
سکدا اے۔ مثال لئی، جانشین فنکشن \pnbLatin{(successor function)}
$f(x) = x + 1$ نوں ایس معنی وچ اُلٹا چلایا جا سکدا اے کہ فنکشن
$g(y) = y - 1$، $f$ دے کیتے نوں ``واپس موڑدا'' (undoes) اے۔

پر سانوں احتیاط کرنا پئے گا۔ بھاویں~$g$ دی تعریف اک فنکشن
$\Int \to \Int$ تعریف کردی اے، پر ایہ اک \emph{فنکشن} $\Nat
\to \Nat$ تعریف نہیں کردی، کیوں جے $g(0) \notin \Nat$۔ سو سادیاں
صورتاں وچ وی ایہ صاف نہیں کہ فنکشن نوں اُلٹا چلایا جا سکدا اے یا
نہیں؛ ایہ دائرۂ تعریف (domain) تے ہدف سیٹ (codomain) اُتے منحصر
ہو سکدا اے۔

اُلٹ فنکشن \pnbLatin{(inverse function)} دا تصور ایس گل نوں ہور ٹھیک ٹھیک
بناؤندا اے۔
\end{explain}

\begin{defn}
اک فنکشن $g \colon B \to A$، فنکشن $f \colon A \to B$ دا
\emph{اُلٹ} (inverse) اے جے $f(g(y)) = y$ تے $g(f(x)) = x$،
ہر $x \in A$ تے ہر $y \in B$ لئی۔
\end{defn}

جے $f$ دا اُلٹ~$g$ ہووے، تاں اُلٹ دی علامت \pnbLatin{(inverse notation)} وچ
اسی اکثر $f^{-1}$ لکھدے آں،~$g$ دی تھاں۔

\begin{explain}
ہُن اسی طے کراں گے کہ فنکشناں دے اُلٹ (inverses) کدوں ہوندے نیں۔
$f\colon A \to B$ دے اُلٹ لئی اک فطری امیدوار $g\colon B \to A$
اے جیہڑا ایویں ``تعریف کیتا'' جاوے:
\[
g(y) = \text{``اوہ'' $x$ جیہدے لئی $f(x) = y$۔}
\]
پر ``تعریف کیتا'' (تے ``اوہ'') دے گرد ڈراؤنے واوین دسدے نیں کہ ایہ
اجے تعریف نہیں۔ گھٹ توں گھٹ، ایہ پوری عمومیت نال ہمیشہ کم نہیں
کرے گی۔ ایس تعریف راہیں فنکشن ملن لئی اک، تے صرف اک،~$x$ ہونا
چاہیدا اے جیہدے لئی $f(x) = y$---یعنی~$g$ دا آؤٹ پٹ اکلوتے طور
تے طے ہونا چاہیدا اے۔ نالے ہر $y \in B$ لئی ایہ طے ہونا چاہیدا اے۔
جے $x_1$ تے $x_2 \in A$ ایہو جہے ہون کہ $x_1 \neq x_2$ پر
$f(x_1) = f(x_2)$، تاں $g(y)$، $y = f(x_1) = f(x_2)$ لئی اکلوتے
طور تے طے نہیں ہووے گا۔ تے جے کوئی~$x$ ای نہیں جیہدے لئی
$f(x) = y$، تاں $g(y)$ بالکل طے نہیں ہوندا۔ ہور لفظاں وچ،~$g$ دی
تعریف لئی~$f$ دا !!{injective} (اک بہ اک؛ injective) وی تے
!!{surjective} (ہر ہدف قدر لین والا؛ surjective) وی ہونا ضروری اے۔

آؤ ہولی ہولی چلیے۔ اسی سوال نوں دو حصیاں وچ ونڈاں گے: دتا ہویا
فنکشن~$f\colon A \to B$ ہووے، تاں فنکشن $g\colon B \to A$ کدوں
ہوندا اے جیہدے لئی $g(f(x)) = x$؟ ایہو جہا~$g$، $f$ دے اثر نوں
``واپس موڑدا'' (undoes) اے تے~$f$ دا \emph{کھبا اُلٹ}
\pnbLatin{(left inverse)} آکھلاندا اے۔ دوجا،
فنکشن $h\colon B \to A$ کدوں ہوندا اے جیہدے لئی $f(h(y)) = y$؟
ایہو جہا~$h$، $f$ دا \emph{سجا اُلٹ} \pnbLatin{(right inverse)} آکھلاندا
اے---$f$، $h$ دے اثر نوں ``واپس موڑدا'' (undoes) اے۔
\end{explain}

\begin{prop}
جے دائرۂ تعریف (domain) خالی نہیں تے $f\colon A \to B$،
!!{injective} (اک بہ اک؛ injective) اے، تاں اک
\emph{کھبا اُلٹ} \pnbLatin{(left inverse)} $g\colon B \to A$،~$f$ دا اے جیہدے لئی
$g(f(x)) = x$، ہر $x \in A$ اُتے۔
% Source correction OLFUN-001: require nonempty A; exact condition is A nonempty or B empty.
\end{prop}

\begin{proof}
من لو کہ دائرۂ تعریف خالی نہیں تے $f\colon A \to B$،
!!{injective} (اک بہ اک؛ injective) اے۔ اک $y \in B$ اُتے غور
کرو۔ جے $y \in \ran{f}$، تاں کوئی $x \in A$ اے جیہدے لئی
$f(x) = y$۔ کیوں جے $f$، !!{injective} اے، اک بہ اک ہون کرکے
ایہو جہا اکو~$x \in A$ اے۔ فیر
اسی تعریف کر سکدے آں: $g(y) = x$، یعنی
$g(y)$ ``اوہ'' $x \in A$ اے جیہدے لئی $f(x) = y$۔ جے
$y \notin \ran{f}$، تاں اسی ایہنوں کسے وی~$a \in A$ نال جوڑ سکدے
آں۔ دائرۂ تعریف خالی نہ ہون دی شرط کرکے اسی اک $a \in A$ چُن
سکدے آں تے
$g\colon B \to A$ نوں ایویں تعریف کر سکدے آں:
\[
g(y) = \begin{cases}
    x & \text{جے $f(x) = y$}\\
    a & \text{جے $y \notin \ran{f}$۔}
\end{cases}
\]
ایہ ہر $y \in B$ لئی تعریف شدہ اے، کیوں جے ایہو جہے ہر
$y \in \ran{f}$ لئی عین اک $x \in A$ اے جیہدے لئی $f(x) = y$۔
تعریف موجب، جے $y = f(x)$، تاں $g(y) = x$، یعنی
$g(f(x)) = x$۔
\end{proof}

\begin{prob}
وکھاؤ کہ جے $f\colon A \to B$ دا کھبا اُلٹ \pnbLatin{(left inverse)}~$g$
اے، تاں~$f$، !!{injective} (اک بہ اک؛ injective) اے۔
\end{prob}

\begin{prop}
    جے $f\colon A \to B$، !!{surjective} (ہر ہدف قدر لین والا؛
    surjective) اے، تاں اک \emph{سجا اُلٹ} \pnbLatin{(right inverse)}
    $h\colon B \to A$،~$f$ دا اے جیہدے لئی
    $f(h(y)) = y$، ہر~$y \in B$ اُتے۔
\end{prop}

\begin{proof}
من لو کہ $f\colon A \to B$، !!{surjective} (ہر ہدف قدر لین والا؛
surjective) اے۔ اک $y \in B$ اُتے غور کرو۔ کیوں جے~$f$،
!!{surjective} اے، یعنی ہر ہدف قدر لیندا اے، ایس لئی کوئی $x_y \in A$
اے جیہدے لئی $f(x_y) = y$۔ فیر اسی تعریف کر سکدے آں:
$h(y) = x_y$، یعنی ہر $y \in B$ لئی اسی کوئی $x \in A$ چُندے آں
جیہدے لئی $f(x) = y$؛ کیوں جے~$f$، !!{surjective} اے، ایس لئی
چُنن واسطے ہمیشہ گھٹ توں گھٹ اک ہوندا اے۔\footnote{کیوں جے $f$،
!!{surjective} اے، ہر~$y \in B$ لئی سیٹ $\Setabs{x}{f(x) = y}$
غیر خالی اے۔~$h$ دی ساڈی تعریف لئی سانوں ایہناں وچوں ہر سیٹ توں
اک $x$ چُننا پینا اے۔ ایہ گل کہ ہمیشہ ایویں چُنیا جا سکدا اے، اصل
وچ صاف نہیں---ایہ چوناں کر سکن دی صلاحیت بس اک مسلّم دے طور تے منی
گئی اے۔ ہور لفظاں وچ، ایہ قضیہ نام نہاد چُناؤ دے مسلّم
\pnbLatin{(Axiom of Choice)} نوں مندا اے؛
ایہ مسئلہ اسی اگے \oliflabeldef{sth:choice::chap}{
\olref[sth][choice][]{chap} وچ فیر چکاں گے}{نظر انداز کر دیواں گے}۔
پر کئی خاص صورتاں وچ، مثال لئی جدوں $A = \Nat$ ہووے یا محدود ہووے،
یا جدوں $f$، !!{bijective} ہووے، چُناؤ دے مسلّم دی لوڑ نہیں۔ (خاص
صورت وچ جدوں $f$، !!{bijective} ہووے، ہر $y \in B$ لئی سیٹ
$\Setabs{x}{f(x) = y}$ دا عین اک !!{element} ہوندا اے، ایس لئی
چُنن دی کوئی گل ای نہیں رہندی۔)}
تعریف موجب، جے $x = h(y)$، تاں $f(x) = y$، یعنی کسے وی
$y \in B$ لئی $f(h(y)) = y$۔
\end{proof}

\begin{prob}
وکھاؤ کہ جے $f\colon A \to B$ دا سجا اُلٹ \pnbLatin{(right inverse)}~$h$
اے، تاں~$f$، !!{surjective} (ہر ہدف قدر لین والا؛ surjective) اے۔
\end{prob}

\begin{explain}
  پچھلے ثبوت دے خیالاں نوں اکٹھا کرکے سانوں ہُن ایہ ملدا اے کہ ہر
  !!{bijection} (اک بہ اک تے ہر ہدف قدر لین والا؛ bijection) دا اک
  اُلٹ ہوندا اے، یعنی اک اکو فنکشن جیہڑا~$f$ دا کھبا اُلٹ
  \pnbLatin{(left inverse)} وی اے تے سجا اُلٹ \pnbLatin{(right inverse)} وی۔
\end{explain}

\begin{prop}\ollabel{prop:bijection-inverse}
جے $f\colon A \to B$، !!{bijective} (اک بہ اک تے ہر ہدف قدر لین
والا؛ bijective) اے، تاں اک اُلٹ فنکشن \pnbLatin{(inverse function)}
$f^{-1}\colon B \to A$ اے جیہدے لئی ہر $x \in A$ اُتے
$f^{-1}(f(x)) = x$ تے ہر $y \in B$ اُتے $f(f^{-1}(y)) = y$۔
\end{prop}

\begin{proof}
مشق (Exercise)۔
\end{proof}

\begin{prob}
\olref[sfr][fun][inv]{prop:bijection-inverse} ثابت کرو (Prove)۔
تہانوں~$f^{-1}$ تعریف کرنا اے، وکھاؤنا اے کہ ایہ فنکشن اے، تے
وکھاؤنا اے کہ ایہ~$f$ دا اُلٹ (inverse) اے، یعنی
$f^{-1}(f(x)) = x$ تے
$f(f^{-1}(y)) = y$، ہر $x \in A$ تے $y \in B$ لئی۔
\end{prob}

\begin{explain}
اُلٹ کڈھن دا اک ذرا ہور عام طریقہ وی اے۔ اسی \olref[kin]{sec} وچ
ویکھیا سی کہ ہر فنکشن $f$ اک !!a{surjection} (ہر ہدف قدر لین والا
فنکشن؛ surjection)
$f' \colon A \to \ran{f}$ پیدا کردا اے، $f'(x) = f(x)$ رکھ کے ہر
$x \in A$ لئی۔ صاف اے
کہ جے $f$، !!{injective} اے، تاں $f'$، !!{bijective} اے، سو
\olref{prop:bijection-inverse} موجب اوہدا اک اکلوتا اُلٹ اے۔ علامت
دی اک بہت چھوٹی جیہی ڈھیل \pnbLatin{(minor abuse of notation)} نال اسی کدی
کدی $f'$ دے اُلٹ نوں بس ``~$f$ دا اُلٹ'' آکھ دیندے آں۔
\end{explain}

\begin{prop}\ollabel{prop:left-right}%
  وکھاؤ کہ جے $f\colon A \to B$ دا کھبا اُلٹ \pnbLatin{(left inverse)}~$g$
  تے سجا اُلٹ \pnbLatin{(right inverse)}~$h$ اے، تاں $h = g$۔
\end{prop}

\begin{proof}
  مشق (Exercise)۔
\end{proof}

\begin{prob}
  \olref[sfr][fun][inv]{prop:left-right} ثابت کرو (Prove)۔
\end{prob}

\begin{prop}\ollabel{prop:inverse-unique}
ہر فنکشن~$f$ دا ودھ توں ودھ اک اُلٹ \pnbLatin{(at most one inverse)} ہوندا اے۔
\end{prop}

\begin{proof}
  من لو کہ $g$ تے $h$ دونیں~$f$ دے اُلٹ (inverses) نیں۔ فیر خاص
  طور تے $g$،~$f$ دا کھبا اُلٹ \pnbLatin{(left inverse)} اے تے $h$ سجا اُلٹ
  \pnbLatin{(right inverse)}۔
  \olref{prop:left-right} موجب، $g = h$۔
\end{proof}
\endgroup
% /OLP-0024

% OLP-0025 content/sets-functions-relations/functions/composition.tex
\begingroup
\def\olfilename{composition}
\expandafter\def\csname ol@sectioncs\endcsname{section}

\olfileid{sfr}{fun}{cmp}
\olsection{فنکشناں دی ترکیب \pnbLatin{(Composition of Functions)}}

\begin{explain}
\oliflabeldef{sfr:fun:inv:sec}{اسی \olref[inv]{sec} وچ ویکھیا سی
کہ اُلٹ~$f^{-1}$، !!a{bijection} (بائجیکشن؛ bijection)~$f$ دا، آپ وی
اک فنکشن اے۔ فنکشناں اُتے اک ہور عمل فنکشن دی ترکیب \pnbLatin{(function composition)} اے: ا}{ا}سی دو فنکشناں، $f$ تے~$g$، دی ترکیب نال نواں
فنکشن تعریف کر سکدے آں: پہلے~$f$ تے اوہدے بعد~$g$ لاگو کریے۔ ایہ
اوہدوں ای ممکن اے جدوں حاصل قدراں دے سیٹ (ranges) تے دائرۂ تعریف
(domains) آپس وچ میل کھان؛ یعنی~$f$ دا حاصل قدراں دا سیٹ (range)،
~$g$ دے دائرۂ تعریف (domain) دا ذیلی سیٹ (subset) ہونا چاہیدا اے۔
\oliflabeldef{sfr:rel:ops:sec}{فنکشناں اُتے ایہ عمل، تعلقاں (relations)
اُتے \olref[rel][ops]{sec} وچ دتے نسبتی حاصل ضرب \pnbLatin{(relative product)}
دے عمل دا مماثل (analogue) اے۔}{}

اک خاکہ فنکشناں دی ترکیب (composition) دا خیال سمجھاؤن وچ مدد دے
سکدا اے۔ \olref{fig:composition} وچ اسی دو فنکشن $f \colon A \to B$
تے $g \colon B \to C$، تے اوہناں دی ترکیب~$(\comp{f}{g})$ وکھاؤندے
آں۔ فنکشن $(\comp{f}{g}) \colon A \to C$،~$A$ دے ہر !!{element}
(element/member) نوں~$C$ دے !!a{element} نال جوڑدا اے۔~$C$ دے کیہڑے
!!{element} نال~$A$ دا !!a{element} جوڑیا جاندا اے، ایہ ایویں
طے ہوندا اے: انپٹ (input) $x \in A$ دتا ہووے، تاں پہلے فنکشن~$f$
نوں~$x$ اُتے لاگو کرو؛ اوہ $f(x) = y \in B$ آؤٹ پٹ (output) دیندا
اے۔ فیر فنکشن~$g$ نوں~$y$ اُتے لاگو کرو؛ اوہ
$g(f(x)) = g(y) = z \in C$ آؤٹ پٹ دیندا اے۔
\begin{figure}
  \olasset[2\olphotowidth]{assets/diagrams/composition.tikz}
  \caption{ترکیب (composition) $g \circ f$، دو فنکشناں $f$ تے~$g$ دی۔}
  \ollabel{fig:composition}
\end{figure}
\end{explain}

\begin{defn}[ترکیب (Composition)]
من لو کہ $f\colon A \to B$ تے $g\colon B \to C$ فنکشن نیں۔ پہلے
$f$ تے فیر~$g$ لا کے بنائی \emph{ترکیب} (composition)،
$\comp{f}{g} \colon A \to C$ اے، جتھے
$(\comp{f}{g})(x) = g(f(x))$۔
\end{defn}

\begin{ex}
فنکشن $f(x) = x + 1$ تے $g(x) = 2x$ اُتے غور کرو۔ کیوں جے
$(\comp{f}{g})(x) = g(f(x))$، ہر انپٹ (input)~$x$ لئی پہلے اوہدا
جانشین (successor) لینا، تے فیر حاصل نوں دو نال ضرب دینا اے۔ ایس لئی
اوہناں دی ترکیب (composition) $(\comp{f}{g})(x) = 2(x+1)$ نال دتی
جاندی اے۔
\end{ex}

\begin{prob}
وکھاؤ کہ جے $f \colon A \to B$ تے $g \colon B \to C$ دونیں
!!{injective} (اک بہ اک؛ injective) ہون، تاں
$\comp{f}{g}\colon A \to C$ وی !!{injective} (اک بہ اک؛ injective)
اے۔
\end{prob}

\begin{prob}
وکھاؤ کہ جے $f \colon A \to B$ تے $g \colon B \to C$ دونیں
!!{surjective} (ہر ہدف قدر لین والے؛ surjective) ہون، تاں
$\comp{f}{g}\colon A \to C$ وی !!{surjective} (ہر ہدف قدر لین والا؛
surjective) اے۔
\end{prob}

\begin{prob}
من لو کہ $f \colon A \to B$ تے $g \colon B \to C$۔ وکھاؤ کہ
ترکیبی فنکشن~$\comp{f}{g}$ دا گراف (graph) $R_f \mid R_g$ اے۔
\end{prob}
\endgroup
% /OLP-0025

%\olimport{isomorphic-functions}

% OLP-0026 content/sets-functions-relations/functions/partial-functions.tex
\begingroup
\def\olfilename{partial-functions}
\expandafter\def\csname ol@sectioncs\endcsname{section}

\olfileid{sfr}{fun}{par}

\olsection{جزوی فنکشن \pnbLatin{(Partial Functions)}}

\begin{explain}
کدی کدی فنکشن دی تعریف وچ ایہ شرط ہولی کرنا فائدہ مند ہوندا اے کہ
ہر ممکن انپٹ (input) لئی فنکشن دا آؤٹ پٹ (output) تعریف شدہ ہووے۔
ایہو جہیاں نقشہ بندیاں نوں \emph{جزوی فنکشن} \pnbLatin{(partial functions)}
آکھیا جاندا اے۔
\end{explain}

\begin{defn}
اک \emph{جزوی فنکشن} \pnbLatin{(partial function)} $f \colon A \pto B$ ایہو
جیہی نقشہ بندی (mapping) اے جیہڑی~$A$ دے ہر !!{element}
(element/member) نوں~$B$ دے ودھ توں ودھ اک !!{element} نال جوڑدی
اے۔ جے $f$،~$B$ دا اک عنصر $x \in A$ نال جوڑدا اے، تاں اسی آکھدے
آں کہ $f(x)$ \emph{تعریف شدہ} (defined) اے، تے نہیں تاں
\emph{غیر تعریف شدہ} (undefined) اے۔ جے $f(x)$ تعریف شدہ ہووے، تاں
اسی $f(x) \fdefined$ لکھدے آں، نہیں تاں $f(x) \fundefined$۔ جزوی
فنکشن~$f$ دا \emph{دائرۂ تعریف} (domain)،~$A$ دا اوہ ذیلی سیٹ
(subset) اے جتھے ایہ تعریف شدہ اے، یعنی
$\dom{f} = \Setabs{x \in A}{f(x) \fdefined}$۔
\end{defn}

\begin{ex}
ہر فنکشن $f\colon A \to B$ اک جزوی فنکشن \pnbLatin{(partial function)} وی اے۔
اوہ جزوی فنکشن جیہڑے~$A$ اُتے ہر تھاں تعریف شدہ ہون---یعنی جنہاں
نوں اسی ہُنے تک بس فنکشن آکھدے آئے آں---\emph{کُل} (total) فنکشن
وی آکھلاندے نیں۔
\end{ex}

\begin{ex}
جزوی فنکشن \pnbLatin{(partial function)} $f \colon \Real \pto \Real$، جیہڑا
$f(x) = 1/x$ نال دتا اے، $x = 0$ اُتے غیر تعریف شدہ (undefined) تے
ہور ہر تھاں تعریف شدہ (defined) اے۔
\end{ex}

\begin{prob}
$f\colon A \pto B$ دتا ہووے، تاں جزوی فنکشن \pnbLatin{(partial function)}
$g\colon B \pto A$ ایویں تعریف کرو: کسے وی $y \in B$ لئی، جے
اکلوتا $x \in A$ اے جیہدے لئی $f(x) = y$، تاں $g(y) = x$؛ نہیں
تاں $g(y) \fundefined$۔ وکھاؤ کہ جے $f$ اک بہ اک (injective) اے،
تاں $g(f(x)) = x$ ہر $x \in \dom{f}$ لئی، تے $f(g(y)) = y$ ہر
$y \in \ran{f}$ لئی۔ ایتھے دائرۂ تعریف (domain) تے حاصل قدراں دا
سیٹ (range) اپنے اوپرلے علامتی معنی وچ ورتے گئے نیں۔
\end{prob}

\begin{defn}[جزوی فنکشن دا گراف \pnbLatin{(Graph of a partial function)}]
من لو کہ $f\colon A \pto B$ اک جزوی فنکشن \pnbLatin{(partial function)} اے۔
~$f$ دا \emph{گراف} (graph) اوہ تعلق (relation)
$R_f \subseteq A \times B$ اے جیہڑا ایویں تعریف ہوندا اے:
\[
R_f = \Setabs{\tuple{x,y}}{f(x) = y}.
\]
\end{defn}

\begin{prop}
من لو کہ $R \subseteq A \times B$ وچ ایہ خاصیت اے کہ جدوں وی $Rxy$
تے $Rxy'$ ہون، تاں $y = y'$۔ فیر $R$ اوس جزوی فنکشن \pnbLatin{(partial function)} $f\colon A \pto B$ دا گراف اے جیہڑا ایویں تعریف ہوندا اے:
جے کوئی $y$ اے جیہدے لئی $Rxy$، تاں $f(x) = y$، نہیں تاں
$f(x) \fundefined$۔ جے $R$ نالے \emph{سریئل} (serial) ہووے، یعنی ہر
$x \in A$ لئی کوئی $y \in B$ ہووے جیہدے لئی $Rxy$، تاں $f$ کُل
(total) اے۔
\end{prop}

\begin{proof}
من لو کہ کوئی $y$ اے جیہدے لئی $Rxy$۔ جے کوئی ہور $y' \neq y$
ہوندا جیہدے لئی $Rxy'$، تاں $R$ اُتے دتی شرط ٹٹ جاندی۔ ایس لئی، جے
کوئی $y$ اے جیہدے لئی $Rxy$، تاں اوہ $y$ اکلوتا اے، تے سو $f$
صحیح طور تے تعریف شدہ (well-defined) اے۔ صاف اے کہ $R_f = R$، تے
$f$ کُل (total) اے جے~$R$ سریئل (serial) ہووے۔
\end{proof}
\endgroup
% /OLP-0026

% OLP-0717 content/sets-functions-relations/functions/isomorphic-functions.tex
\begingroup
\def\olfilename{isomorphic-functions}
\expandafter\def\csname ol@sectioncs\endcsname{section}

\olfileid{sfr}{fun}{iso}
\olsection{ہم ساختی}

\begin{explain}
\emph{ہم ساختی} اوہ دو طرفہ میل اے جیہڑا متعلقہ سیٹاں دی بناوٹ قائم رکھدا اے۔ بناوٹ دا مطلب ایہ اے کہ سیٹاں دے !!{element}s وچ کِہڑے تعلق قائم نیں۔ ایہ دو سیٹ $X=\{1,2,3\}$ تے $Y=\{4,5,6\}$ ویکھو۔ دوناں سیٹاں دی بناوٹ اگلا ہون، گھٹ ہون تے ودھ ہون دے تعلقاں نال بندی اے۔ دوہاں سیٹاں وچ ہم ساختی اوہ !!a{bijection} اے جیہڑی اوہ بناوٹاں قائم رکھدی اے۔ سو !!a{bijective} فنکشن $f \colon X
\to Y$ ہم ساختی اے جے $i<j$ اوہدوں تے صرف اوہدوں جدوں $f(i)<f(j)$، $i>j$ اوہدوں تے صرف اوہدوں جدوں $f(i)>f(j)$، تے $j$، $i$ دا اگلا اے اوہدوں تے صرف اوہدوں جدوں $f(j)$، $f(i)$ دا اگلا اے۔
\end{explain}

\begin{defn}[ہم ساختی]
منّ لو $U$ جوڑا $\langle X, R\rangle$ اے تے $V$ جوڑا $\langle
Y, S\rangle$ اے، جتھے $X$ تے $Y$ سیٹ نیں تے $R$ تے $S$، بالترتیب $X$ تے $Y$ اُتے تعلق نیں۔ !!{bijection} $f$، $X$ توں $Y$ تک، $U$ توں $V$ تک \emph{ہم ساختی} اے اوہدوں تے صرف اوہدوں جدوں اوہ تعلقاں دی بناوٹ قائم رکھدی اے۔ یعنی کسے وی $x_{1}$ تے $x_{2}$ لئی، جیہڑے $X$ وچ نیں، $\tuple{x_1,x_2} \in R$ اوہدوں تے صرف اوہدوں جدوں $\tuple{f(x_1),f(x_2)} \in S$۔
\end{defn}

\begin{ex}
ایہ دو سیٹ $X=\{1,2,3\}$ تے $Y=\{4,5,6\}$ تے گھٹ ہون تے ودھ ہون دے تعلق ویکھو۔ فنکشن $f\colon X \to
Y$، جتھے $f(x) = 7-x$، $\tuple{X,<}$ تے $\tuple{Y,>}$ وچکار ہم ساختی اے۔
\end{ex}
\endgroup
% /OLP-0717
\OLEndChapterHook
\endgroup
% /OLP-0020

% OLP-0027 content/sets-functions-relations/size-of-sets/size-of-sets-complete.tex
\begingroup
\def\olfilename{size-of-sets-complete}
\expandafter\def\csname ol@sectioncs\endcsname{section}

\olchapter{sfr}{siz}{سیٹاں دا سائز \pnbLatin{(The Size of Sets)}}

\begin{editorial}
ایہ باب گنتی وار فہرستاں (enumerations)، قابلِ شمار ہون
(countability) تے ناقابلِ شمار ہون (uncountability) اُتے گل کردا اے۔
کئی حصے دو روپاں وچ دتے گئے نیں: اک ہور بنیادی روپ \pnbLatin{(more elementary version)}، جیہڑا گنتی وار فہرستاں نوں لسٹاں (lists)، یا $\PosInt$ توں
ہر ہدف قدر لین والے فنکشناں (surjections) دے طور تے لیندا اے؛ تے اک
ہور مجرد روپ \pnbLatin{(more abstract version)}، جیہڑا گنتی وار فہرستاں نوں
$\Nat$ نال اک بہ اک مطابقتاں (bijections) دے طور تے تعریف کردا اے۔
\end{editorial}

% OLP-0028 content/sets-functions-relations/size-of-sets/introduction.tex
\begingroup
\def\olfilename{introduction}
\expandafter\def\csname ol@sectioncs\endcsname{section}

\olfileid{sfr}{siz}{int}

\olsection{تعارف \pnbLatin{(Introduction)}}

جدوں گیورگ کانتور نے 1870واں دے دہاکے وچ سیٹ تھیوری \pnbLatin{(set theory)} نوں
ترقی دِتّی، اوہدے مقصداں وچوں اک لامتناہی (infinite) اکٹھ دے خیال نوں
قابلِ قبول بنانا سی---اوہ بالفعل لامتناہیت \pnbLatin{(actual infinity)}، جیویں
وچلے زمانے دے اہلِ علم آکھدے سن۔ ایس کم دا اک اہم حصہ وکھ وکھ سیٹاں
دے \emph{سائز} (size) نال اوہدا ورتاؤ سی۔ جے $a$،
$b$ تے $c$ سارے اک دوجے توں وکھ نیں، تاں سیٹ $\{a, b, c\}$ بدیہی
طور اُتے $\{a, b\}$ نالوں \emph{وڈا} اے۔ پر لامتناہی سیٹاں بارے کیہ؟
کی اوہ سارے اک دوجے جڈے وڈے نیں؟ پتا لگدا اے کہ نہیں۔

ایتھے پہلا اہم خیال گنتی وار فہرست (enumeration) دا اے۔ اسی ہر متناہی
(finite) سیٹ دے سارے !!{element}s، یعنی عنصراں (elements)، نوں فہرست
وچ لکھ کے اوہدی فہرست بنا سکدے آں۔ کجھ لامتناہی سیٹاں دے سارے
!!{element}s وی فہرست وچ لکھ سکدے آں، جے اسی فہرست نوں آپ لامتناہی
ہون دی اجازت دیئیے۔ ایہو جہے سیٹ !!{enumerable}، یعنی گنتی وار فہرست
بنن جوگ (enumerable)، آکھے جاندے نیں۔ کانتور دا حیران کرن والا نتیجہ،
جیہنوں اسی ایس باب دے آخر تک پوری طرح سمجھ لاواں گے، ایہ سی کہ کجھ
لامتناہی سیٹ !!{enumerable} نہیں ہوندے۔
\endgroup
% /OLP-0028

% OLP-0029 content/sets-functions-relations/size-of-sets/enumerability.tex
\begingroup
\def\olfilename{enumerability}
\expandafter\def\csname ol@sectioncs\endcsname{section}

\olfileid{sfr}{siz}{enm}

\olsection{گنتی وار فہرستاں \pnbLatin{(Enumerations)} تے \usetoken{S}{enumerable} سیٹ}

\begin{editorial}
  ایہ حصہ سیٹاں دیاں گنتی وار فہرستاں (enumerations) بارے گل کردا اے
  تے اوہناں نوں $\PosInt$ توں ہر ہدف قدر لین والے فنکشناں
  (surjections) دے طور تے تعریف کردا اے۔ گھٹ ریاضیاتی پس منظر والے
  قاریاں لئی ایہ گل ہولی ہولی سمجھائی گئی اے۔ اک متبادل، مختصر روپ
  \olref[enm-alt]{sec} وچ دتا گیا اے، جیہڑا گنتی وار فہرستاں نوں وکھرے
  ڈھنگ نال، یعنی $\Nat$ نال اک بہ اک تے ہر ہدف قدر لین والیاں مطابقتاں
  (bijections) (یا اوہدے کسے مُڈھلے ٹکڑے نال ایہو جہیاں مطابقتاں) دے
  طور تے تعریف کردا اے۔
\end{editorial}

\begin{explain}
اسی سیٹاں دے !!{element}s دی فہرست (list) بنا کے پہلے وی اوہناں دیاں مثالاں
دتیاں نیں۔ آؤ ہُن عام لفظاں وچ گل کریے کہ کسے سیٹ دے !!{element}s دی
فہرست کِویں تے کدوں بن سکدی اے، بھاویں اوہ سیٹ لامتناہی ای کیوں نہ ہووے۔
\end{explain}

\begin{defn}[گنتی وار فہرست (Enumeration)، غیر رسمی طور تے]
غیر رسمی طور تے، سیٹ~$A$ دی اک \emph{گنتی وار فہرست (enumeration)}، $A$ دے
!!{element}s دی اک فہرست (جیہڑی لامتناہی وی ہو سکدی اے) اے، جیہدے
وچ $A$ دا ہر !!{element} کسے متناہی مقام \pnbLatin{(finite position)} اُتے آوے۔
جے $A$ دی کوئی گنتی وار فہرست اے، تاں $A$ نوں
\emph{!!{enumerable}} (گنتی وار فہرست بنن جوگ؛ enumerable) آکھیا جاندا اے۔
\end{defn}

\begin{explain}
گنتی وار فہرستاں (enumerations) بارے کجھ گلاں:
\begin{enumerate}
\item اسی صرف اوہناں فہرستاں نوں گنتی وار فہرست منّنے آں جیہڑیاں
  کسے تھاں توں شروع ہوندیاں نیں، تے جیہناں وچ پہلے عنصر توں
  بغیر ہر !!{element} دے عین پہلاں اک اکلوتا !!{element} ہوندا اے۔
  دوجے لفظاں وچ، فہرست دے پہلے !!{element} تے کسے وی ہور !!{element}
  دے وچکار صرف متناہی گنتی دے عنصر ہوندے نیں۔ خاص طور تے، ایس دا
  مطلب اے کہ گنتی وار فہرست دے ہر !!{element} دا مقام (position) متناہی اے: پہلے
  !!{element} دا مقام~$1$، دوجے دا مقام~$2$، تے ایویں اگے۔
\item اکّو سیٹ~$A$ دیاں وکھ وکھ گنتی وار فہرستاں ہو سکدیاں نیں، جیہناں
  وچ !!{element}s دا آون دا کرم وکھ ہووے: $4$، $1$، $25$، $16$،~$9$
  پہلے پنج مربع عدداں دے (سیٹ) دی اوہنی ای چنگی گنتی وار فہرست اے
  جِنّی $1$، $4$، $9$، $16$،~$25$ اے۔
\item وادھو دہراؤ والیاں (redundant) فہرستاں وی گنتی وار فہرستاں ای نیں: $1$،
  $1$، $2$، $2$، $3$، $3$،~\dots{} اوسے سیٹ دی گنتی وار فہرست اے
  جیہدی $1$، $2$، $3$،~\dots{} اے۔
\item جدوں اسی کوئی خاص گنتی وار فہرست دسدے آں، تاں کرم (order) تے
  دہراؤ (repetition)
  \emph{واقعی} اہم ہوندے نیں: اسی مثبت صحیح عدداں دی گنتی وار فہرست
  $1$، $2$، $3$، $1$، \dots{} نال شروع کر سکدے آں، پر جدوں اوہناں نوں
  معیاری ڈھنگ نال $1$، $2$، $3$، $4$،~\dots لکھیا جاوے تاں نمونہ
  ودھ سوکھا دِسدا اے۔
\item گنتی وار فہرست دا اک مُڈھ ہونا لازمی اے: \dots، $3$، $2$، $1$
  مثبت صحیح عدداں دی گنتی وار فہرست نہیں، کیوں جے ایس دا کوئی پہلا
  !!{element} نہیں۔ ایہ گل غیر رسمی تعریف توں کِویں نکلدی اے، ویکھن
  لئی اپنے آپ نوں پچھو، ”عدد 76 فہرست دے کیہڑے مقام اُتے آندا اے؟“
\item اگلی فہرست مثبت صحیح عدداں دی گنتی وار فہرست نہیں:
  $1$، $3$، $5$، \dots، $2$، $4$، $6$، \dots\@ مسئلہ ایہ اے کہ
  جفت عدد متناہی مقاماں دی بجائے $\infty + 1$، $\infty + 2$،
  $\infty + 3$ مقاماں اُتے آندے نیں۔
\item خالی سیٹ گنتی وار فہرست بنن جوگ اے: خالی فہرست اوہدی گنتی وار فہرست اے!{}
\end{enumerate}
\end{explain}

\begin{prop}
  جے $A$ دی کوئی گنتی وار فہرست اے، تاں اوہدی دہراؤ توں بغیر وی اک
  گنتی وار فہرست اے۔
\end{prop}

\begin{proof}
  فرض کرو $A$ دی اک گنتی وار فہرست $x_1$، $x_2$، \dots{} اے، جیہدے
  وچ ہر $x_i$، $A$ دا !!{element} اے۔ اسی دہرائے ہوئے !!{element}s
  ہٹا کے فہرست دا دہراؤ ختم کر سکدے آں۔ مثال لئی، اسی ایس نوں ایہو
  جہی نویں گنتی وار فہرست بنا سکدے آں جیہدے وچ $x_i$ نوں اوس ویلے
  لکھیے جے اوہ $A$ دا !!a{element} ہووے جیہڑا $x_1$، \dots،
  $x_{i-1}$ وچ شامل نہیں؛ تے جے $x_i$ پہلے ای $x_1$، \dots،~$x_{i-1}$
  وچ آ چکیا ہووے، تاں اوہنوں فہرست توں ہٹا دیئیے۔
\end{proof}

پچھلی دلیل وکھاؤندی اے کہ گنتی وار فہرستاں تے !!{enumerable} سیٹاں
نوں چنگی طرح سمجھن تے اوہناں بارے ثبوت دین لئی سانوں ہور ٹھیک تعریف
دی لوڑ اے۔ اگلی تعریف اوہ کم کردی اے۔

\begin{defn}[گنتی وار فہرست (Enumeration)، رسمی طور تے]
سیٹ $A \neq \emptyset$ دی اک \emph{گنتی وار فہرست} کوئی وی
!!{surjective} (ہر ہدف قدر لین والا؛ surjective) فنکشن
$f \colon \PosInt \to A$ اے۔
\end{defn}

\begin{explain}
آؤ تسلی کر لئیے کہ رسمی تعریف تے ممکنہ طور تے لامتناہی فہرست والی
غیر رسمی تعریف اک دوجی دے برابر نیں۔ پہلے، $\PosInt$ توں سیٹ~$A$ ول
ہر !!{surjective} (ہر ہدف قدر لین والا؛ surjective) فنکشن، $A$ دی
گنتی وار فہرست بناؤندا اے۔ ایہو جہا
فنکشن اُتلی غیر رسمی تعریف والی فہرست $f(1)$، $f(2)$، $f(3)$، \dots
متعین کردا اے۔ کیوں جے $f$، !!{surjective} (ہر ہدف قدر لین والا) اے،
ایس لئی $A$ دا ہر
!!{element} لازماً $f(n)$ دی قدر اے، کسے~$n \in \PosInt$ لئی۔ ایس
لئی $A$ دا ہر !!{element} فہرست دے کسے متناہی مقام اُتے آندا اے۔
فنکشن ممکن اے !!{injective} (اک بہ اک؛ injective) نہ ہووے، ایس لئی
فہرست وچ دہراؤ ہو سکدا
اے، پر ایہ قابلِ قبول اے (جیویں اُتے دسیا گیا اے)۔

دوجے پاسے، $A$ دے سارے !!{element}s دی گنتی وار فہرست دتی ہووے، تاں
اسی !!a{surjective} (ہر ہدف قدر لین والا؛ surjective) فنکشن
$f\colon \PosInt \to A$ ایویں تعریف کر
سکدے آں: $f(n)$، فہرست دا $n$واں !!{element} ہووے؛ تے جے فہرست دا
کوئی $n$واں !!{element} نہ ہووے تاں اوہدا آخری !!{element} ہووے۔
ایہ صرف اوس صورت وچ !!a{surjective} (ہر ہدف قدر لین والا) فنکشن نہیں
بناؤندا جدوں $A$ خالی اے تے، ایس لئی، فہرست وی خالی اے۔ سو ہر غیر
خالی فہرست !!a{surjective} فنکشن $f\colon \PosInt \to A$ متعین کردی اے۔
\end{explain}

\begin{defn}
  \ollabel{defn:enumerable}
  سیٹ~$A$، !!{enumerable} (گنتی وار فہرست بنن جوگ؛ enumerable) عین
  اوس ویلے اے جدوں اوہ خالی ہووے یا اوہدی کوئی گنتی وار فہرست ہووے۔
\end{defn}

\begin{ex}
مثبت صحیح عدداں ($\PosInt$) دی گنتی وار فہرست بناؤن والا فنکشن سادہ
شناختی فنکشن $f(n) = n$ اے۔ قدرتی عدداں $\Nat$ دی گنتی وار فہرست
بناؤن والا فنکشن $g(n) = n - 1$ اے۔
\end{ex}

\begin{ex}
اگلے دتے ہوئے فنکشن $f\colon \PosInt \to \PosInt$ تے
$g \colon \PosInt \to \PosInt$:
\begin{align*}
f(n) & = 2n \text{ تے}\\
g(n) & = 2n - 1
\end{align*}
ترتیب وار جفت مثبت صحیح عدداں تے طاق مثبت صحیح عدداں دیاں گنتی وار
فہرستاں بناؤندے نیں۔ پر دوہاں وچوں کوئی وی فنکشن $\PosInt$ دی گنتی
وار فہرست نہیں، کیوں جے دوہاں وچوں کوئی وی !!{surjective}
(ہر ہدف قدر لین والا؛ surjective) نہیں۔
\end{ex}

\begin{prob}
مثبت مربع عدداں $1$، $4$، $9$، $16$، \dots دی گنتی وار فہرست تعریف کرو۔
\end{prob}

\begin{ex}
فنکشن $f(n) = (-1)^{n} \lceil \frac{(n-1)}{2}\rceil$ (جتھے
$\lceil x \rceil$، \emph{چھت (ceiling)} فنکشن اے، جیہڑا $x$ نوں اوس
سب توں چھوٹے صحیح عدد تک اُتے پورا کردا اے جیہڑا x توں گھٹ نہ ہووے)
صحیح عدداں دے سیٹ~$\Int$ دی گنتی
وار فہرست بناؤندا اے۔ دھیان دیو کہ $f$، مثبت تے منفی صحیح عدداں دے
وچکار ”ٹپدے“ ہوئے $\Int$ دیاں قدراں پیدا کردا اے:
\[
\begin{array}{c c c c c c c c}
f(1) & f(2) & f(3) & f(4) & f(5) & f(6) & f(7) & \dots \\ \\
- \lceil \tfrac{0}{2} \rceil & \lceil \tfrac{1}{2}\rceil & - \lceil \tfrac{2}{2} \rceil & \lceil \tfrac{3}{2} \rceil & - \lceil \tfrac{4}{2} \rceil  & \lceil \tfrac{5}{2}
\rceil & - \lceil \tfrac{6}{2} \rceil & \dots \\ \\
0 & 1 & -1 & 2 & -2 & 3 & -3 & \dots
\end{array}
\]
% OLSIZ-001 ماخذ دی درستی: f(7) دی قدر -3 اے؛ انگریزی ماخذ دی آخری قطار وچ ایہ قدر رہ گئی سی۔
تسیں $f$ نوں اگلی صورتاں نال تعریف کیتا ہویا وی سمجھ سکدے او:
\[
f(n) = \begin{cases}
  0 & \text{جے $n = 1$ اے}\\
  n/2 & \text{جے $n$ جفت اے}\\
  -(n-1)/2 & \text{جے $n$ طاق تے $>1$ اے}
  \end{cases}
\]
\end{ex}

\begin{prob}
  وکھاؤ کہ جے $A$ تے $B$، !!{enumerable} نیں، تاں $A \cup B$ وی
  گنتی وار فہرست بنن جوگ اے۔ ایہ کرن لئی، فرض کرو کہ !!{surjective} فنکشن
  $f\colon \PosInt \to A$ تے $g\colon \PosInt \to B$ نیں، تے اک
  !!a{surjective} فنکشن~$h\colon \PosInt \to A \cup B$ تعریف کر کے
  ثابت کرو کہ اوہ !!{surjective} اے۔ اوہ صورتاں وی ویکھو جدوں $A$ یا
  ~$B = \emptyset$ ہووے۔
\end{prob}

\begin{prob}
  وکھاؤ کہ جے $B \subseteq A$ تے $A$، !!{enumerable} اے، تاں~$B$ وی
  گنتی وار فہرست بنن جوگ اے۔ ایہ کرن لئی، فرض کرو کہ !!a{surjective} فنکشن
  $f\colon \PosInt \to A$ اے۔ اک !!a{surjective} فنکشن
  ~$g\colon \PosInt \to B$ تعریف کر کے ثابت کرو کہ اوہ
  !!{surjective} اے۔ جے $B = \emptyset$ ہووے تاں کی ہوندا اے؟
\end{prob}

\begin{prob}
  $n$ اُتے استقرا نال وکھاؤ کہ جے $A_1$، $A_2$، \dots، $A_n$ سارے
  !!{enumerable} نیں، تاں $A_1 \cup \dots \cup A_n$ وی گنتی وار فہرست بنن جوگ
  اے۔ تسیں ایہ حقیقت منّ سکدے او کہ جے دو سیٹ $A$ تے~$B$،
  !!{enumerable} نیں، تاں~$A \cup B$ وی گنتی وار فہرست بنن جوگ اے۔
\end{prob}

سیٹ دے !!{element}s دی فہرست بناؤندیاں پہلے $1$ویں !!{element} توں
گنتی شروع کرنا شاید ودھ قدرتی لگے، پر ریاضی دان چیزاں گنن لئی قدرتی
عدد~$\Nat$ ورتنا پسند کردے نیں۔ اوہ فہرست دے $0$ویں، $1$ویں، $2$ویں
وغیرہ !!{element}s دی گل کردے نیں۔ ایس دے مطابق، اسی گنتی وار فہرست
نوں $\Nat$ توں~$A$ ول !!a{surjective} فنکشن دے طور تے تعریف کر سکدے
آں۔ بے شک، دوویں تعریفاں اک دوجی دے برابر نیں۔

\begin{prop}\ollabel{prop:enum-shift}
  !!a{surjection} $f\colon \PosInt \to A$ عین اوس ویلے موجود اے جدوں
  !!a{surjection} $g\colon \Nat \to A$ موجود اے۔
\end{prop}

\begin{proof}
  !!a{surjection} $f\colon \PosInt \to A$ دتا ہووے، تاں اسی
  $g(n) = f(n+1)$ تعریف کر سکدے آں، سارے $n \in \Nat$ لئی۔ ایہ ویکھنا
  سوکھا اے کہ $g\colon \Nat \to A$، !!{surjective} اے۔ اُلٹ پاسے،
  !!a{surjection} $g\colon \Nat \to A$ دتا ہووے تاں
  $f(n) = g(n-1)$ تعریف کرو۔
\end{proof}

ایس توں سانوں اگلا نتیجہ ملدا اے:

\begin{cor}\ollabel{cor:enum-nat}
سیٹ $A$، !!{enumerable} عین اوس ویلے اے جدوں اوہ خالی ہووے یا کوئی
!!a{surjective} فنکشن $f\colon \Nat \to A$ موجود ہووے۔
\end{cor}

اسی اُتے گل کیتی سی کہ سیٹ~$A$ دے !!{element}s دی فہرست نوں دہراؤ
توں بغیر فہرست وچ بدلیا جا سکدا اے۔ ایہ گل گنتی وار فہرستاں لئی وی
سچ اے، پر اینوں پورے رسمی ڈھنگ نال بیان تے ثابت کرنا ذرا اوکھا اے۔ ہر
فنکشن $f\colon \PosInt \to A$ دا ہر $n \in \PosInt$ لئی تعریف شدہ
ہونا لازمی اے۔ جے $A$ وچ صرف متناہی !!{element}s ہون، تاں ظاہر اے کہ
اسی $\PosInt$ دے لامتناہی !!{element}s اُتے تعریف شدہ ایہو جہا فنکشن
نہیں بنا سکدے جیہڑا $A$ دے سارے !!{element}s نوں قدر دے طور تے لے پر
اکّو قدر کدے دو واری نہ لے۔ ایس صورت وچ، یعنی جدوں دہراؤ توں بغیر
فہرست متناہی ہووے، سانوں $f$ لئی ہور دائرۂ تعریف چُننا پئے گا، جیہدے
وچ صرف متناہی !!{element}s ہون۔ دہراؤ نہ ہون دا مطلب اے کہ $f$،
!!{injective} ہووے۔ کیوں جے اوہ !!{surjective} وی اے، ایس لئی اسی کسے
متناہی سیٹ $\{1, \dots, n\}$ یا $\PosInt$ تے~$A$ دے وچکار
!!a{bijection} لبھ رہے آں۔

\begin{prop}\ollabel{prop:enum-bij}
جے $f\colon \PosInt \to A$، !!{surjective} اے (یعنی $A$ دی گنتی وار
فہرست اے)، تاں !!a{bijection} $g\colon Z \to A$ موجود اے جتھے $Z$،
یا~$\PosInt$ اے یا $\{1, \dots, n\}$ اے، کسے~$n \in \PosInt$ لئی۔
\end{prop}

\begin{proof}
  اسی فنکشن $g$ نوں بازگشتی (recursively) طور تے تعریف کردے آں:
  $g(1) = f(1)$
  رکھو۔ جے $g(i)$ پہلے تعریف ہو چکیا اے، تاں $g(i+1)$ نوں $f(1)$،
  $f(2)$، \dots{} دی اوہ پہلی قدر رکھو جیہڑی پہلے $g(1)$، \dots،
  $g(i)$ وچ شامل نہ ہووے، جے ایہو جہی قدر موجود ہووے۔ جے $A$ دے
  ٹھیک $n$ !!{element}s ہون، تاں $g(1)$، \dots، $g(n)$ سارے تعریف
  شدہ نیں، تے ایس طرح اسی فنکشن $g\colon \{1, \dots, n\} \to A$
  تعریف کر دتا اے۔ جے $A$ دے لامتناہی !!{element}s ہون، تاں ہر $i$
  لئی $A$ دا کوئی !!a{element}، گنتی وار فہرست $f(1)$، $f(2)$،
  \dots وچ ضرور ہووے گا، جیہڑا پہلے $g(1)$، \dots، $g(i)$ وچ شامل
  نہ ہووے۔ ایس صورت وچ اسی فنکشن $g\colon \PosInt \to A$ تعریف کر
  دتا اے۔

  فنکشن $g$، !!{surjective} اے، کیوں جے $A$ دا ہر عنصر، $f(1)$،
  $f(2)$، \dots{} وچ شامل اے (کیوں جے $f$، !!{surjective} اے) تے
  ایس لئی آخرکار $g(i)$ دی قدر بنے گا، کسے~$i$ لئی۔ ایہ !!{injective}
  وی اے، کیوں جے جے کوئی $j < i$ ایہو جہا ہوندا کہ $g(j) = g(i)$،
  تاں $g(i)$ پہلے ای $g(1)$، \dots، $g(i-1)$ وچ شامل ہوندا، جیہڑا
  $g$ دی ساڈی تعریف دے اُلٹ اے۔
\end{proof}

\begin{cor}\ollabel{cor:enum-nat-bij}
سیٹ $A$، !!{enumerable} عین اوس ویلے اے جدوں اوہ خالی ہووے یا
!!a{bijection} $f\colon N \to A$ موجود ہووے، جتھے یا $N = \Nat$ اے
یا $N = \{0, \dots, n\}$ اے، کسے $n \in \Nat$ لئی۔
\end{cor}

\begin{proof}
$A$، !!{enumerable} عین اوس ویلے اے جدوں $A$ خالی ہووے یا
!!a{surjective} $f\colon \PosInt \to A$ موجود ہووے۔
\olref{prop:enum-bij} مطابق، پچھلی گل عین اوس ویلے اے جدوں اک
!!a{bijective} فنکشن~$f\colon Z \to A$ موجود ہووے، جتھے
$Z = \PosInt$ یا $Z = \{1, \dots, n\}$ اے، کسے $n \in \PosInt$
لئی۔ \olref{prop:enum-shift} دے ثبوت والی اوسے دلیل نال، اگوں ایہ عین
اوس ویلے اے جدوں اک !!a{bijection} $g\colon N \to A$ موجود ہووے،
جتھے یا $N = \Nat$ اے یا $N = \{0, \dots, n-1\}$ اے۔
\end{proof}

\begin{prob}
  \olref[sfr][siz][enm]{defn:enumerable} مطابق، سیٹ $A$ گنتی وار فہرست
  بنن جوگ عین اوس ویلے اے جدوں $A = \emptyset$ ہووے یا !!a{surjective}
  $f\colon \PosInt \to A$ موجود ہووے۔ ”!!{enumerable} سیٹ“ نوں ٹھیک
  ٹھیک ایویں تعریف کرنا وی ممکن اے: سیٹ عین اوس ویلے گنتی وار فہرست
  بنن جوگ اے جدوں !!a{injective} فنکشن $g\colon A \to \PosInt$ موجود
  ہووے۔ وکھاؤ کہ دوویں تعریفاں اک دوجی دے برابر نیں، یعنی وکھاؤ کہ
  !!a{injective} فنکشن $g\colon A \to \PosInt$ عین اوس ویلے موجود اے
  جدوں یا $A = \emptyset$ ہووے یا !!a{surjective}
  $f\colon \PosInt \to A$ موجود ہووے۔
\end{prob}
\endgroup
% /OLP-0029

% OLP-0030 content/sets-functions-relations/size-of-sets/zig-zag.tex
\begingroup
\def\olfilename{zig-zag}
\expandafter\def\csname ol@sectioncs\endcsname{section}

\olfileid{sfr}{siz}{zigzag}

\olsection{کانتور دا زگ زیگ طریقہ \pnbLatin{(Cantor's Zig-Zag Method)}}

\begin{explain}
اسی کجھ ”سوکھیاں“ گنتی وار فہرستاں (enumerations) پہلے ای ویکھ چکے
آں۔ ہُن اسی کجھ ذرا اوکھا ویکھاں گے۔ قدرتی عدداں دیاں ترتیب وار
جوڑیاں \pnbLatin{(ordered pairs)} دے سیٹ اُتے غور
کرو\oliflabeldef{sfr:set:pai:sec}{، جیہنوں اسی \olref[set][pai]{sec}
وچ ایویں تعریف کیتا سی:}{، جیہڑا ایویں تعریف اے:}
\[
\Nat \times \Nat = \Setabs{\tuple{n,m}}{n,m \in \Nat}
\]
اسی ایہناں ترتیب وار جوڑیاں نوں اک \emph{جدول} (array) وچ ایویں سجا سکدے آں:
\[
\begin{array}{ c | c | c | c | c | c}
& \mathbf 0 & \mathbf 1 & \mathbf 2 & \mathbf 3 & \dots \\
\hline
\mathbf 0 & \tuple{0,0} & \tuple{0,1} & \tuple{0,2} & \tuple{0,3} & \dots \\
\hline
\mathbf 1 & \tuple{1,0} & \tuple{1,1} & \tuple{1,2} & \tuple{1,3} & \dots \\
\hline
\mathbf 2 & \tuple{2,0} & \tuple{2,1} & \tuple{2,2} & \tuple{2,3} & \dots \\
\hline
\mathbf 3 & \tuple{3,0} & \tuple{3,1} & \tuple{3,2} & \tuple{3,3} & \dots \\
\hline
\vdots & \vdots & \vdots & \vdots & \vdots & \ddots\\
\end{array}
\]
صاف اے کہ $\Nat \times \Nat$ دی ہر ترتیب وار جوڑی جدول وچ ٹھیک اک
واری آوے گی۔ خاص طور تے، $\tuple{n,m}$، $n$ویں قطار (row) تے $m$ویں
کالم (column) وچ آوے گی۔ پر اسی ایہو جہے جدول دے عنصراں (elements)
نوں ”اک جہتی“ (one-dimensional) فہرست (list) وچ کِویں سجاواں؟
ہیٹھلے جدول دا نمونہ ایہ کرن دا اک طریقہ وکھاؤندا اے (بھاویں ہور کئی
طریقے وی موجود نیں):
\[
\begin{array}{ c | c | c | c | c | c | c}
& \mathbf 0 & \mathbf 1 & \mathbf 2 & \mathbf 3 & \mathbf 4 &\dots \\
\hline
\mathbf 0 & 0  & 1& 3 & 6& 10 &\ldots \\
\hline
\mathbf 1 &2 & 4& 7 & 11 & \dots &\ldots \\
\hline
\mathbf 2 & 5 & 8 & 12 & \ldots & \dots&\ldots \\
\hline
\mathbf 3 & 9 & 13 & \ldots & \ldots & \dots & \ldots \\
\hline
\mathbf 4 & 14 & \ldots & \ldots & \ldots & \dots & \ldots \\
\hline
\vdots & \vdots & \vdots & \vdots & \vdots&\ldots & \ddots\\
\end{array}
\]\noindent
ایس نمونے نوں \emph{کانتور دا زگ زیگ طریقہ} \pnbLatin{(Cantor's zig-zag method)}
آکھیا جاندا اے۔ ایہ $\Nat \times \Nat$ دی گنتی وار فہرست
(enumeration) ایویں بناؤندا اے:
\[
\tuple{0,0}, \tuple{0,1}, \tuple{1,0}, \tuple{0,2}, \tuple{1,1},
\tuple{2,0}, \tuple{0,3}, \tuple{1,2}, \tuple{2,1}, \tuple{3,0}, \dots
\]
تے ایس توں اگلی گل قائم ہوندی اے:
\end{explain}

\begin{prop}\ollabel{natsquaredenumerable}
$\Nat \times \Nat$ دی گنتی وار فہرست بن سکدی اے؛ یعنی ایہہ
!!{enumerable} اے۔
\end{prop}

\begin{proof}
فنکشن $f \colon \Nat \to \Nat\times\Nat$، ہر $k \in \Nat$ نوں اوہ
ترتیب وار جوڑی $\tuple{n,m} \in \Nat \times \Nat$ تے بھیجدا اے جیہدے
لئی $k$، کانتور دے زگ زیگ جدول دی $n$ویں قطار (row) تے $m$ویں کالم
(column) دے ملاپ اُتے لکھی قدر اے۔
\end{proof}

\begin{explain}
ایس تکنیک نوں بڑے سوہنے ڈھنگ نال عام وی کیتا جا سکدا اے۔ مثال لئی،
اسی قدرتی عدداں دیاں ترتیب وار تِکّیاں \pnbLatin{(ordered triples)} دے سیٹ دی
گنتی وار فہرست بناؤن لئی اینوں ورت سکدے آں، یعنی:
\[
\Nat \times \Nat \times \Nat = \Setabs{\tuple{n,m,k}}{n,m,k \in \Nat}
\]
اسی $\Nat \times \Nat \times \Nat$ نوں $\Nat \times \Nat$ تے
$\Nat$ دا کارتیسی حاصل ضرب \pnbLatin{(Cartesian product)} سمجھدے آں، یعنی
\[
\Nat^3 = (\Nat \times \Nat) \times \Nat =
\Setabs{\tuple{\tuple{n,m},k}}{n, m, k
  \in \Nat }
\]
تے ایس طرح اسی $\Nat^3$ دی گنتی وار فہرست اک جدول (array) نال بنا
سکدے آں: اک محور (axis) اُتے $\Nat$ دی گنتی وار فہرست دے لیبل لا
دیو، تے دوجے محور اُتے $\Nat^2$ دی گنتی وار فہرست دے:
\[
\begin{array}{ c | c | c | c | c | c}
& \mathbf 0 & \mathbf 1 & \mathbf 2 & \mathbf 3 & \dots \\
\hline
\mathbf{\tuple{0,0}} & \tuple{0,0,0} & \tuple{0,0,1} & \tuple{0,0,2} & \tuple{0,0,3} & \dots \\
\hline
\mathbf{\tuple{0,1}} & \tuple{0,1,0} & \tuple{0,1,1} & \tuple{0,1,2} & \tuple{0,1,3} & \dots \\
\hline
\mathbf{\tuple{1,0}} & \tuple{1,0,0} & \tuple{1,0,1} & \tuple{1,0,2} & \tuple{1,0,3} & \dots \\
\hline
\mathbf{\tuple{0,2}} & \tuple{0,2,0} & \tuple{0,2,1} & \tuple{0,2,2} & \tuple{0,2,3} & \dots\\
\hline
\vdots & \vdots & \vdots & \vdots & \vdots & \ddots \\
\end{array}
\]
سو کانتور دے زگ زیگ طریقے ورگا طریقہ ورت کے اسی اوسے طرح $\Nat^3$
دی گنتی وار فہرست (enumeration) حاصل کر سکدے آں۔ تے اسی ایویں اگے
ودھدے ہوئے، $\Nat^n$ دیاں گنتی وار فہرستاں حاصل کر سکدے آں، ہر قدرتی
عدد $n$ لئی۔ سو ساڈے کول اگلا نتیجہ اے:
\end{explain}

\begin{prop}
$\Nat^n$ دی، ہر $n \in \Nat$ لئی، گنتی وار فہرست بن سکدی اے؛ یعنی
ایہہ !!{enumerable} اے۔
\end{prop}

\begin{prob}\label{sfr:siz:zigzag:prob:posint-n}
وکھاؤ کہ $(\PosInt)^n$ دی، ہر $n \in \Nat$ لئی، گنتی وار فہرست بن
سکدی اے؛ یعنی ایہہ !!{enumerable} اے۔
\end{prob}

\begin{prob}\label{sfr:siz:zigzag:prob:posint-star}
  وکھاؤ کہ $(\PosInt)^*$ دی گنتی وار فہرست بن سکدی اے؛ یعنی ایہہ
  !!{enumerable} اے۔ تسیں \cref{sfr:siz:zigzag:prob:posint-n}
  منّ سکدے او۔
\end{prob}
\endgroup
% /OLP-0030

% OLP-0031 content/sets-functions-relations/size-of-sets/pairing.tex
\begingroup
\def\olfilename{pairing}
\expandafter\def\csname ol@sectioncs\endcsname{section}

\olfileid{sfr}{siz}{pai}
\olsection{جوڑی فنکشن تے کوڈ \pnbLatin{(Pairing Functions and Codes)}}

\begin{explain}
کانتور دا زگ زیگ طریقہ \pnbLatin{(Cantor's zig-zag method)}، $\Nat^n$ دے گنتی وار
فہرست بنن جوگ ہون (enumerability) نوں اکھیں ویکھن جوگ بنا دیندا اے۔
پر آؤ اپنے $\Nat^2$ نوں وکھاؤن والے جدول (array) اُتے دھیان دیئیے۔
جدول وچ زگ زیگ لکیر دے نال نال چل کے تے مقام گن کے اسی جانچ سکدے آں
کہ $\tuple{1,2}$ نوں عدد~$7$ نال جوڑیا گیا اے۔ پر جے اسی ایہ گل ہور
سدھے ڈھنگ نال حساب کر سکدے، تاں چنگا ہوندا۔ یعنی ساڈے کول زگ زیگ
گنتی وار فہرست دا \emph{اُلٹ} (inverse)،
$g\colon \Nat^2 \to \Nat$، ہونا چنگا ہوندا، ایہو جہا کہ
\[
g(\tuple{0,0}) = 0, \;
g(\tuple{0,1}) = 1, \;
g(\tuple{1,0}) = 2, \; \dots,
g(\tuple{1,2}) = 7, \; \dots
\]
ایس نال اسی ٹھیک ٹھیک حساب کر سکدے آں کہ ترتیب وار جوڑی
$\tuple{n, m}$ ساڈی گنتی وار فہرست (enumeration) وچ کتھے آوے گی۔

اصل وچ اسی دو گلاں اُتے غور کر کے $g$ نوں سدھے ڈھنگ نال تعریف کر سکدے
آں۔ پہلی گل: جے $n$ویں قطار (row) تے $m$ویں کالم (column) وچ قدر~$v$
اے، تاں $(n+1)$ویں قطار تے $(m-1)$ویں کالم وچ قدر $v + 1$ اے۔ دوجی
گل: ساڈی گنتی وار فہرست دی پہلی قطار تکونی عدداں \pnbLatin{(triangular numbers)}
اُتے مشتمل اے، جیہڑی
$0$، $1$، $3$، $6$ وغیرہ توں شروع ہوندی اے۔ $k$واں تکونی عدد،
$< k$ قدرتی عدداں دا جوڑ اے، جیہڑا $k(k+1)/2$ دے طور تے حساب کیتا جا
سکدا اے۔ ایہ دوویں گلاں کٹھیاں کر کے اگلے فنکشن اُتے غور کرو:
\[
  g(n,m) = \frac{(n+m+1)(n+m)}{2} + n
\]
اسی اکثر صرف $g(n, m)$ لکھدے آں، $g(\tuple{n, m})$ دی بجائے، کیوں جے
ایہ اکھاں نوں سوکھا لگدا اے۔ ایہ تہانوں آکھدا اے کہ پہلاں
$(n+m)^\text{th}$ تکونی عدد کڈھو، تے فیر اوہدے وچ $n$ جوڑ دیو۔ ایہ
جدول نوں ٹھیک اوسے ڈھنگ نال بھردا اے جیویں اسی چاہندے آں۔ سو خاص طور
تے جوڑی $\tuple{1, 2}$ نوں $\frac{4 \times 3}{2} +
1 = 7$ ول بھیجیا جاندا اے۔

ایہ فنکشن $g$، جوڑیاں دے سیٹ دی اک گنتی وار فہرست دا \emph{اُلٹ}
(inverse) اے۔ ایہو جہے فنکشناں نوں \emph{جوڑی فنکشن} \pnbLatin{(pairing functions)} آکھیا جاندا اے۔
\end{explain}

\begin{defn}[جوڑی فنکشن \pnbLatin{(Pairing function)}]
  فنکشن $f\colon A \times B \to \Nat$ اک حسابی
  \emph{جوڑی فنکشن} \pnbLatin{(pairing function)} اے جے $f$ اک بہ اک؛ injective
  ہووے۔ اسی ایہ وی آکھدے آں کہ $f$، $A \times B$ نوں \emph{کوڈ کردا}
  (encodes) اے، تے $f(x,y)$، $\tuple{x,y}$ دا \emph{کوڈ} (code) اے۔
\end{defn}

\begin{explain}
اسی جوڑی فنکشناں نوں، مثلاً، قدرتی عدداں دیاں جوڑیاں کوڈ کرن
(encode) لئی ورت سکدے آں؛ یا دوجے لفظاں وچ، اسی عنصراں (elements) دی
ہر \emph{جوڑی} (pair) نوں اک \emph{اکلے} (single) عدد نال پیش کر سکدے
آں۔ جوڑی فنکشن دا اُلٹ ورت کے اسی عدد نوں \emph{ڈی کوڈ} (decode) کر
سکدے آں، یعنی پتا لا سکدے آں کہ اوہ کیہڑی جوڑی نوں پیش کردا اے۔
\end{explain}

\begin{prob}
سارے غیر منفی ناطق عدداں \pnbLatin{(non-negative rational numbers)} دے سیٹ دی
اک گنتی وار فہرست (enumeration) دیو۔
\end{prob}

\begin{prob}
وکھاؤ کہ $\Rat$ دی گنتی وار فہرست بن سکدی اے، یعنی ایہہ
!!{enumerable} اے۔ یاد رکھو کہ ہر ناطق عدد نوں
$z/m$ کسر دے طور تے لکھیا جا سکدا اے، جتھے $z \in \Int$ تے
$m \in \Nat^+$۔
\end{prob}

\begin{prob}
$\Bin^*$ دی اک گنتی وار فہرست (enumeration) تعریف کرو۔
\end{prob}

\begin{prob}
اپنے ابتدائی منطق کورس توں یاد کرو کہ ہر ممکن صداقتی جدول \pnbLatin{(truth table)} اک صداقتی فنکشن \pnbLatin{(truth function)} نوں ظاہر کردا اے۔ دوجے لفظاں
وچ، صداقتی فنکشن
$\Bin^k \to \Bin$ والے سارے فنکشن نیں، کسے~$k$ لئی۔ ثابت کرو کہ سارے صداقتی فنکشناں
دا سیٹ گنتی وار فہرست بنن جوگ (enumerable) اے۔
\end{prob}

\begin{prob}
وکھاؤ کہ کسے وی لامتناہی، گنتی وار فہرست بنن جوگ
(!!{enumerable}) سیٹ دے سارے متناہی ذیلی سیٹاں دا سیٹ وی گنتی وار
فہرست بنن جوگ (!!{enumerable}) اے۔
\end{prob}

\begin{prob}
$\Nat$ دے کسے ذیلی سیٹ نوں \emph{ہم متناہی} (cofinite) عین اوس ویلے
آکھیا جاندا اے جدوں اوہ قدرتی عدداں دے کسے متناہی ذیلی سیٹ دا $\Nat$
دے اندر متمم ہووے؛ یعنی $A \subseteq \Nat$ عین اوس ویلے ہم متناہی اے جدوں
$\Nat\setminus A$ متناہی ہووے۔
% OLSIZ-002 ماخذ دی درستی: متمم قدرتی عدداں دے سیٹ دے اندر کسے متناہی ذیلی سیٹ دا اے؛ انگریزی عبارت وچ ایہ نسبت رہ گئی سی۔
$I$ اوہ سیٹ ہووے جیہدے !!{element}s، یعنی عناصر، ٹھیک $\Nat$ دے
متناہی تے ہم متناہی ذیلی سیٹ نیں۔ وکھاؤ کہ $I$ گنتی وار فہرست بنن
جوگ، یعنی !!{enumerable}، اے۔
\end{prob}

\begin{prob}
وکھاؤ کہ گنتی وار فہرست بنن جوگ (!!{enumerable}) سیٹاں دا
!!{enumerable} اتحاد \pnbLatin{(enumerable union)} وی گنتی وار فہرست بنن جوگ
(!!{enumerable}) اے۔ یعنی جدوں وی $A_1$، $A_2$، \dots{} سیٹ ہون، تے
ہر $A_i$، !!{enumerable} ہووے، تاں اوہناں ساریاں دا اتحاد
$\bigcup_{i=1}^\infty A_i$ وی !!{enumerable} اے۔ [نوٹ: ایہ اوکھا اے!]
\end{prob}

\begin{prob}
$f \colon A \times B \to \Nat$ کوئی وی جوڑی فنکشن \pnbLatin{(pairing function)}
ہووے۔ وکھاؤ کہ $f$ دا اُلٹ (inverse)، $A \times B$ دی اک گنتی وار
فہرست (enumeration) اے۔
\end{prob}

\begin{prob}
ایہو جہا فنکشن صاف صاف دسو جیہڑا $\Nat^3$ نوں کوڈ (encodes) کردا ہووے۔
\end{prob}
\endgroup
% /OLP-0031

% OLP-0032 content/sets-functions-relations/size-of-sets/pairing-alt.tex
\begingroup
\def\olfilename{pairing-alt}
\expandafter\def\csname ol@sectioncs\endcsname{section}

\olfileid{sfr}{siz}{pai-alt}

\olsection{اک متبادل جوڑی فنکشن \pnbLatin{(An Alternative Pairing Function)}}

\begin{explain}
$\Nat^2$ دیاں ہور وی گنتی وار فہرستاں (enumerations) نیں جیہناں دے
اُلٹ فنکشن \pnbLatin{(inverse functions)} لبھنے ودھ سوکھے نیں۔ ایہناں وچوں اک
ایہ اے۔ گنتی وار فہرست نوں جدول (array) وچ ویکھن دی بجائے، مثبت صحیح
عدداں \pnbLatin{(positive integers)} دی اوس فہرست توں شروع کرو جیہدے نال مُڈھ وچ
خالی تھاواں جوڑیاں ہوئیاں نیں۔ تصور کرو کہ ایہ تھاواں ترتیب وار
جوڑیاں \pnbLatin{(ordered pairs)} $\tuple{n,m}$ نال اگلے ڈھنگ نال بھریاں جا رہیاں نیں۔
پہلاں اوہ جوڑیاں لو جیہناں دے پہلے مقام \pnbLatin{(first coordinate)} اُتے~$0$ اے (یعنی جوڑیاں
$\tuple{0,m}$)؛ پہلی جوڑی (یعنی $\tuple{0,0}$) پہلی خالی تھاں اُتے
رکھو، فیر اک خالی تھاں چھڈ دیو، دوجی جوڑی (یعنی $\tuple{0,1}$) اگلی
خالی تھاں اُتے رکھو، فیر اک ہور خالی تھاں چھڈ دیو، تے ایویں اگے۔
% PNBSIZ-002 ماخذ دی درستی: ایتھے دوجی جوڑی (0,1) اے؛ انگریزی عبارت وچ غلطی نال (0,2) لکھیا سی، جدوں کہ نال دی پہلی قطار وی (0,1) ای وکھاؤندی اے۔
ساڈی گنتی وار فہرست دا (ادھورا) مُڈھ ہُن ایویں دِسدا اے:
\[\small
\begin{array}{@{}c c c c c c c c c c c@{}}
\mathbf 1 & \mathbf 2 & \mathbf 3 & \mathbf 4 & \mathbf 5 & \mathbf 6 & \mathbf 7 & \mathbf 8 & \mathbf 9 & \mathbf{10} & \dots \\ \\
\tuple{0,0} &  & \tuple{0,1} &  & \tuple{0,2} &  & \tuple{0,3} & & \tuple{0,4} &  & \dots \\
\end{array}
\]
ہُن جیہڑیاں تھاواں ہلے وی خالی نیں اوہناں لئی جوڑیاں $\tuple{1,m}$
نال ایہ کم دہراؤ، تے فیر ہر دوجی خالی تھاں چھڈدے جاؤ:
\[\small
\begin{array}{@{}c c c c c c c c c c c@{}}
\mathbf 1 & \mathbf 2 & \mathbf 3 & \mathbf 4 & \mathbf 5 & \mathbf 6 & \mathbf 7 & \mathbf 8 & \mathbf 9 & \mathbf{10} & \dots \\ \\
\tuple{0,0} & \tuple{1,0} & \tuple{0,1} &  & \tuple{0,2} & \tuple{1,1} &
\tuple{0,3} & & \tuple{0,4} &  \tuple{1,2} & \dots \\
\end{array}
\]
اوسے ڈھنگ نال جوڑیاں $\tuple{2,m}$، $\tuple{3,m}$ وغیرہ درج کرو۔
% OLSIZ-003 ماخذ دی درستی: لگاتار قطاراں (2,m)، (3,m) وغیرہ نیں؛ انگریزی عبارت وچ (2,m) دو واری آ گیا سی۔
ساڈی مکمل گنتی وار فہرست ایس طرح شروع ہوندی اے:
\[\small
\begin{array}{@{}cc c c c c c c c c c@{}}
\mathbf 1 & \mathbf 2 & \mathbf 3 & \mathbf 4 & \mathbf 5 & \mathbf 6 & \mathbf 7 & \mathbf 8 & \mathbf 9 & \mathbf{10} & \dots \\ \\
\tuple{0,0} & \tuple{1,0} & \tuple{0,1} & \tuple{2,0}  & \tuple{0,2} &
\tuple{1,1} & \tuple{0,3} & \tuple{3,0}  & \tuple{0,4} &  \tuple{1,2} & \dots \\
\end{array}
\]
جے اسی اُتلے جدول دے خانیاں (cells) نوں ایس گنتی وار فہرست (enumeration)
مطابق نمبر دیئیے، تاں سانوں صاف زگ زیگ لکیر \pnbLatin{(zig-zag line)} نہیں لبھے
گی، بلکہ ایہ ترتیب لبھے گی:
\[
\begin{array}{ c | c | c | c | c | c | c | c }
& \mathbf 0 & \mathbf 1 & \mathbf 2 & \mathbf 3 & \mathbf 4 & \mathbf 5 & \dots \\
\hline
\mathbf 0 & 1 & 3 & 5 & 7 & 9 & 11 & \dots \\
\hline
\mathbf 1 & 2 & 6 & 10 & 14 & 18 & \dots & \dots \\
\hline
\mathbf 2 & 4 & 12 & 20 & 28 & \dots & \dots & \dots \\
\hline
\mathbf 3 & 8 & 24 & 40 & \dots & \dots & \dots & \dots \\
\hline
\mathbf 4 & 16 & 48 & \dots & \dots & \dots & \dots & \dots \\
\hline
\mathbf 5 & 32 & \dots & \dots & \dots & \dots & \dots & \dots \\
\hline
\vdots & \vdots & \vdots & \vdots & \vdots & \vdots & \vdots & \ddots\\
\end{array}
\]
اسی ویکھ سکدے آں کہ قطار (row)~$0$ دیاں جوڑیاں ساڈی گنتی وار فہرست دے
طاق نمبر والے مقاماں \pnbLatin{(odd-numbered positions)} اُتے نیں، یعنی جوڑی $\tuple{0,m}$، مقام $2m+1$
اُتے اے؛ دوجی قطار دیاں جوڑیاں $\tuple{1,m}$ اوہناں مقاماں اُتے نیں
جیہناں دا نمبر کسے طاق عدد دا دُگنا اے، خاص طور تے $2 \cdot (2m+1)$؛
تیجی قطار دیاں جوڑیاں $\tuple{2,m}$ اوہناں مقاماں اُتے نیں جیہناں دا
نمبر کسے طاق عدد دا چوگنا اے، $4 \cdot (2m+1)$؛ تے ایویں اگے۔ ہر قطار
لئی $(2m+1)$ دے عامل (factors) $1$، $2$، $4$، $8$، \dots، ٹھیک~$2$ دیاں
قوتاں (powers) نیں: $1= 2^0$، $2 = 2^1$، $4 = 2^2$، $8 = 2^3$، \dots\@
اصل وچ متعلقہ قوت نما (exponent) ہمیشہ زیرِ غور جوڑی دا پہلا رکن اے۔ ایس لئی جوڑی
$\tuple{n,m}$ لئی عامل $2^n$ اے۔ ایس توں سانوں عام فارمولا ملدا اے:
$2^n \cdot (2m+1)$۔ پر ایہ جوڑیاں توں \emph{مثبت} صحیح عدداں ول نقشہ
بندی (mapping) اے، یعنی $\tuple{0,0}$ دا مقام~$1$ اے۔ جے اسی مقام~$0$ توں شروع
کرنا چاہندے آں تاں نتیجے وچوں~$1$ گھٹانا پئے گا۔ ایس توں سانوں ایہ
ملدا اے:
\end{explain}

\begin{ex}
اگلے فارمولے نال دتا گیا فنکشن (function) $h\colon \Nat^2 \to \Nat$
\[
h(n,m) = 2^n (2m+1) - 1
\]
قدرتی عدداں دیاں جوڑیاں دے سیٹ~$\Nat^2$ لئی اک جوڑی فنکشن \pnbLatin{(pairing function)} اے۔
\end{ex}

\begin{explain}
ایس مطابق، $\Nat^2$ دی ساڈی دوجی گنتی وار فہرست (enumeration) وچ جوڑی
$\tuple{0,0}$ دا کوڈ (code) $h(0,0) = 2^0(2\cdot 0+1) - 1 = 0$ اے؛
$\tuple{1,2}$ دا کوڈ $2^{1} \cdot (2 \cdot 2 + 1) - 1 = 2
\cdot 5 - 1 = 9$ اے؛ تے $\tuple{2,6}$ دا کوڈ $2^{2} \cdot (2
\cdot 6 + 1) - 1 = 51$ اے۔
\end{explain}

کدی کدی قدرتی عدداں دیاں جوڑیاں $\Nat^2$ نوں کوڈ کرنا (encode) ای کافی
ہوندا اے، تے ایہ شرط ضروری نہیں ہوندی کہ کوڈ بندی ہر ہدف قدر لین والی
(surjective) ہووے۔ ایہو جہیاں کوڈ بندیاں دے اُلٹ فنکشن \pnbLatin{(inverse functions)}
صرف جزوی فنکشن \pnbLatin{(partial functions)} ہوندے نیں۔

\begin{ex}
اگلے فارمولے نال دتا گیا فنکشن (function) $j\colon \Nat^2 \to \Nat^+$
\[
j(n,m) = 2^n3^m
\]
اک بہ اک؛ !!a{injective} فنکشن $\Nat^2 \to \Nat$ اے۔
\end{ex}
\endgroup
% /OLP-0032

% OLP-0033 content/sets-functions-relations/size-of-sets/non-enumerability.tex
\begingroup
\def\olfilename{non-enumerability}
\expandafter\def\csname ol@sectioncs\endcsname{section}

\olfileid{sfr}{siz}{nen}

\olsection{\printtoken{S}{nonenumerable} سیٹ \pnbLatin{(Non-enumerable Sets)}}

\begin{editorial}
  ایہ حصہ \olref[enm]{sec} وچ دتی تعریف نوں ورت کے $\Bin^\omega$ تے
  $\Pow{\PosInt}$ دے گنتی وار فہرست نہ بنن جوگ ہون (non-enumerability)
  نوں ثابت کردا اے۔
  اینوں \olref[enm-alt]{sec} والے روپ نالوں تھوڑا ودھ ابتدائی تے تھوڑا
  ودھ تفصیلی بناؤن دا مقصد اے۔
\end{editorial}

کجھ سیٹ، جیویں مثبت صحیح عدداں دا سیٹ $\PosInt$، لامتناہی (infinite)
نیں۔ ہُن تک اسی لامتناہی سیٹاں دیاں جیہڑیاں مثالاں ویکھیاں، اوہ سارے
گنتی وار فہرست بنن جوگ (!!{enumerable}) سن۔ پر کجھ لامتناہی سیٹ ایہ
خاصیت نہیں رکھدے۔ ایہو جہے سیٹاں نوں \emph{گنتی وار فہرست نہ بنن جوگ
(!!{nonenumerable})} آکھیا جاندا اے۔

سب توں پہلاں، شاید ایہ گل پہلے ای حیران کرن والی اے کہ
گنتی وار فہرست نہ بنن جوگ (!!{nonenumerable}) سیٹ موجود نیں۔ ہر گنتی
وار فہرست بنن جوگ (!!{enumerable}) سیٹ~$A$ لئی کوئی ہر ہدف قدر لین
والا (!!a{surjective}) فنکشن $f \colon \PosInt \to A$ ہوندا اے۔ جے
کوئی سیٹ !!{nonenumerable} اے، تاں ایہو جہا کوئی فنکشن نہیں ہوندا۔ یعنی
$\PosInt$ دے لامتناہی عنصراں (!!{element}s) نوں~$A$ ول نقشہ بند کرن
والا کوئی وی فنکشن سارے~$A$ نوں پورا نہیں کر سکدا۔ سو~$A$ دے
!!{element}s، لامتناہی مثبت صحیح عدداں \pnbLatin{(positive integers)} نالوں وی
”ودھ“ نیں۔

اسی ایہ کِویں ثابت کراں گے کہ کوئی سیٹ !!{nonenumerable} اے؟ تہانوں
وکھاؤنا پئے گا کہ ایہو جہا کوئی ہر ہدف قدر لین والا فنکشن موجود نہیں
ہو سکدا۔ ایس دے برابر، تہانوں وکھاؤنا پئے گا کہ~$A$ دے عنصراں نوں اک
پاسے ول لامتناہی فہرست \pnbLatin{(one-way infinite list)} وچ گنتی وار نہیں رکھیا
جا سکدا۔ ایہ کرن دا سب
توں چنگا طریقہ ایہ وکھاؤنا اے کہ~$A$ دے !!{element}s دی ہر فہرست
گھٹ توں گھٹ اک عنصر ضرور چھڈ دیندی اے؛ یا کوئی فنکشن
$f\colon \PosInt \to A$، !!{surjective} نہیں ہو سکدا۔ اسی ایہ کم
کانتور دے \emph{قطری طریقے} \pnbLatin{(Cantor's diagonal method)} نال کر سکدے
آں۔~$A$ دے !!{element}s دی
کوئی فہرست، مثلاً $x_1$، $x_2$، \dots، دتی ہووے، تاں اسی~$A$ دا اک
ہور عنصر بناؤندے آں جیہڑا اپنی بناوٹ کر کے اوس فہرست وچ ہو ای نہیں سکدا۔

ساڈی پہلی مثال سیٹ~$\Bin^\omega$ اے، جیہڑا $0$ تے $1$ دیاں ساریاں
لامتناہی ثنائی لڑیاں \pnbLatin{(infinite binary sequences)}، یعنی بنا کھلوت دیاں
لڑیاں، دا سیٹ اے۔

\begin{thm}
\ollabel{thm:nonenum-bin-omega}
$\Bin^\omega$، گنتی وار فہرست نہ بنن جوگ (!!{nonenumerable}) اے۔
\end{thm}

\begin{proof}
تضاد لئی فرض کرو کہ $\Bin^\omega$، گنتی وار فہرست بنن جوگ
(!!{enumerable}) اے، یعنی فرض کرو کہ اک فہرست (list) $s_{1}$،
$s_{2}$، $s_{3}$، $s_{4}$، \dots{}، ~$\Bin^\omega$ دے سارے
عنصراں (!!{element}s) دی اے۔ ایہناں وچوں ہر $s_i$ آپ $0$ تے~$1$
دی اک لامتناہی لڑی \pnbLatin{(infinite sequence)} اے۔ آؤ $j$ویں عنصر نوں، جیہڑا
ایس فہرست دی
$i$ویں لڑی وچ اے، $s_i(j)$ آکھیے۔ تاں $i$ویں لڑی~$s_i$ ایہ اے:
\[
s_i(1), s_i(2), s_i(3), \dots
\]

اسی ایس فہرست نوں، تے ایس وچ ہر لڑی $s_i$ دے عنصراں نوں، اک جدول
(array) وچ سجا سکدے آں:
\[
\begin{array}{c|c|c|c|c|c}
& 1 & 2 & 3 & 4 & \dots \\\hline
1 & \mathbf{s_{1}(1)} & s_{1}(2) & s_{1}(3) & s_1(4) & \dots \\\hline
2 & s_{2}(1)& \mathbf{s_{2}(2)} & s_2(3) & s_2(4) & \dots \\\hline
3 & s_{3}(1)& s_{3}(2) & \mathbf{s_3(3)} & s_3(4) & \dots \\\hline
4 & s_{4}(1)& s_{4}(2) & s_4(3) & \mathbf{s_4(4)} & \dots \\\hline
\vdots & \vdots & \vdots & \vdots & \vdots & \mathbf{\ddots}
\end{array}
\]
کھبے پاسے تھلے ول لیبل فہرست دیاں لڑیاں $s_1$، $s_2$، \dots دے
نمبر دسدے نیں؛ اُتے آر پار دتے نمبر ہر وکھری لڑی دے عنصراں
(!!{element}s) نوں
لیبل کردے نیں۔ مثال لئی، $s_{1}(1)$ اوس عدد دا ناں اے جیہڑا $0$ یا
~$1$ اے تے لڑی $s_{1}$ دا پہلا !!{element} اے؛ تے ایویں اگے۔

ہُن اسی اک اُلٹی قطری لڑی \pnbLatin{(anti-diagonal sequence)} $\overline{s}$،
جیہڑی $0$ تے $1$ دی لامتناہی لڑی اے، بناؤندے آں، جیہڑی ایس فہرست وچ
ہو ای نہیں سکدی۔ $\overline{s}$ دی تعریف فہرست
$s_1$، $s_2$، \dots اُتے منحصر ہووے گی۔ $0$ تے $1$ دیاں لامتناہی
لڑیاں دی ہر لامتناہی فہرست اک ایہو جہی لامتناہی لڑی~$\overline{s}$
پیدا کردی اے جیہڑی پکی گل اے کہ فہرست وچ نہیں آوے گی۔

$\overline{s}$ نوں تعریف کرن لئی اسی اوہدے سارے !!{element}s دسدے آں،
یعنی اسی $\overline{s}(n)$ دسدے آں، ہر $n \in \PosInt$ لئی۔ اسی اُتلے
جدول دے قطر (diagonal) نوں اُتوں تھلے پڑھ کے (ایسے لئی ناں ”قطری
طریقہ“ اے)، تے فیر ہر رکن نوں اُلٹ کے \pnbLatin{(flipping each bit)}، یعنی ہر
$1$ نوں $0$ تے ہر $0$ نوں~$1$ وچ بدل کے، ایہ کم کردے آں۔ ہور
مجرد طور تے، اسی $\overline{s}(n)$ نوں $0$ یا $1$ ایس مطابق تعریف
کردے آں کہ قطر دا $n$واں !!{element}، $s_n(n)$، $1$ اے یا $0$:
\[
\overline{s}(n) =
\begin{cases}
1 & \text{جے $s_{n}(n) = 0$ اے}\\
0 & \text{جے $s_{n}(n) = 1$ اے}.
\end{cases}
\]
جے تہانوں حالت وار تعریف \pnbLatin{(definition by cases)} نالوں فارمولے چنگے
لگدے نیں، تاں تسیں
$\overline{s}(n) = 1 - s_n(n)$ وی تعریف کر سکدے او۔

صاف اے کہ $\overline{s}$، $0$ تے $1$ دی اک لامتناہی لڑی اے، کیوں جے
ایہ ساڈے جدول دے قطر اُتے آون والے $0$ تے $1$ دی لڑی دا بٹ-اُلٹ عکس
\pnbLatin{(bit-flipped diagonal sequence)}
اے۔ سو $\overline{s}$،~$\Bin^\omega$ دا !!a{element} اے۔ پر ایہ
فہرست $s_1$، $s_2$، \dots{} وچ نہیں ہو سکدی۔ کیوں؟

ایہ فہرست دی پہلی لڑی $s_1$ نہیں ہو سکدی، کیوں جے ایہ پہلے مقام
(coordinate) دے
!!{element} وچ $s_1$ نالوں وکھ اے۔ $s_1(1)$ جو وی ہووے، اسی
$\overline{s}(1)$ نوں اوہدے اُلٹ تعریف کیتا اے۔ ایہ فہرست دی دوجی
لڑی نہیں ہو سکدی، کیوں جے $\overline{s}$ دوجے عنصر وچ $s_2$ نالوں
وکھ اے: جے $s_2(2)$، $0$ اے، تاں $\overline{s}(2)$، $1$ اے، تے
اُلٹ صورت وچ وی ایویں۔ تے ایس طرح اگے۔

ہور ٹھیک طور تے: جے $\overline{s}$ فہرست وچ ہوندی، تاں کوئی $k$ ایہو
جہا ہوندا کہ $\overline{s} = s_{k}$۔ دو لڑیاں عین اکّو عین
(identical) اوس ویلے
نیں جدوں اوہ ہر مقام اُتے اکّو قدر رکھن، یعنی ہر~$n$ لئی
$\overline{s}(n) = s_{k}(n)$۔ سو خاص طور تے $n = k$ والی حالت لیئیے،
تاں $\overline{s}(k) = s_{k}(k)$ ہونا لازمی سی۔ $s_k(k)$ یا $0$ اے
یا~$1$۔ جے ایہ $0$ اے، تاں $\overline{s}(k)$ لازماً~$1$ اے---اسی
$\overline{s}$ نوں ایویں ای تعریف کیتا اے۔ پر جے $s_k(k) = 1$ اے،
تاں فیر، $\overline{s}$ دی ساڈی تعریف کر کے، $\overline{s}(k) = 0$
اے۔ دوہاں صورتاں وچ $\overline{s}(k) \neq s_{k}(k)$۔

اسی ایہ فرض کر کے شروع کیتا سی کہ~$\Bin^\omega$ دے !!{element}s دی
اک فہرست $s_1$، $s_2$، \dots{} اے۔ ایس فہرست توں اسی اک لڑی
~$\overline{s}$ بنائی جیہڑی، جیویں اسی ثابت کیتا، فہرست وچ نہیں ہو
سکدی۔ پر ایہ وی پکی گل اے کہ اوہ $0$ تے $1$ دی اک لڑی اے جے سارے
$s_i$، $0$ تے $1$ دیاں لڑیاں نیں، یعنی
$\overline{s} \in \Bin^\omega$۔ ایہ خاص طور تے وکھاؤندا اے کہ
~$\Bin^\omega$ دے \emph{سارے} !!{element}s دی کوئی فہرست نہیں ہو
سکدی، کیوں جے ہر ایہو جہی فہرست توں اسی اک لڑی~$\overline{s}$ وی بنا
سکدے آں جیہڑی پکی گل اے کہ اوس فہرست وچ نہیں؛ سو $\Bin^\omega$ دیاں
ساریاں لڑیاں دی فہرست ہون دا مفروضہ تضاد (contradiction) ول لے جاندا
اے۔
\end{proof}

\begin{explain}
ایس ثبوت دے طریقے نوں ”قطری بناوٹ“ (diagonalization) آکھیا جاندا اے، کیوں جے ایہ
~$\overline{s}$ نوں تعریف کرن لئی جدول دا قطر ورتدا اے۔ قطری بناوٹ
وچ جدول دا موجود ہونا لازمی نہیں: اسی اوسے ورگا خیال ورت کے سیٹاں دا
!!{enumerable} نہ ہونا وی وکھا سکدے آں، بھاویں کوئی جدول تے حقیقی قطر
شامل نہ ہووے۔
\end{explain}

\begin{thm}
\ollabel{thm:nonenum-pownat}
$\Pow{\PosInt}$، گنتی وار فہرست بنن جوگ (!!{enumerable}) نہیں اے۔
\end{thm}

\begin{proof}
اسی اوسے ڈھنگ نال کم کردے آں:~$\PosInt$ دے ذیلی سیٹاں دے سیٹ
\pnbLatin{(power set)} وچوں ذیلی سیٹاں دی ہر
فہرست لئی وکھاؤندے آں کہ~$\PosInt$ دا کوئی ایہو جہا ذیلی سیٹ موجود اے
جیہڑا فہرست وچ نہیں ہو سکدا۔ فرض کرو کہ~$\PosInt$ دے ذیلی سیٹاں دی
اگلی فہرست دتی ہوئی اے:
\[
Z_{1}, Z_{2}, Z_{3}, \dots
\]
ہُن اسی قطری ذیلی سیٹ \pnbLatin{(diagonal subset)} $\overline{Z}$ نوں رکنیت اُلٹ
کے \pnbLatin{(membership flip)} ایویں تعریف کردے آں کہ ہر
$n \in \PosInt$ لئی، $n \in \overline{Z}$ عین اوس ویلے جدوں
$n \notin Z_{n}$:
\[
\overline{Z} = \Setabs{n \in \PosInt}{n \notin Z_n}
\]
صاف اے کہ $\overline{Z}$ مثبت صحیح عدداں دا سیٹ اے، کیوں جے مفروضے
مطابق ہر~$Z_n$ وی ایہو جہا سیٹ اے، تے ایس لئی
$\overline{Z} \in \Pow{\PosInt}$۔ پر $\overline{Z}$ فہرست وچ نہیں
ہو سکدا۔ ایہ وکھاؤن لئی اسی ثابت کراں گے کہ ہر
$k \in \PosInt$ لئی، $\overline{Z} \neq Z_k$۔

سو $k \in \PosInt$ کوئی وی ہووے۔ اسی $\overline{Z}$ نوں رکنیت دا
فیصلہ اُلٹ کے \pnbLatin{(membership flip)} ایویں تعریف
کیتا اے کہ ہر $n \in \PosInt$ لئی، $n \in \overline{Z}$ عین اوس
ویلے جدوں $n \notin Z_n$۔ خاص طور تے $n=k$ رکھ کے،
$k \in \overline{Z}$ عین اوس ویلے جدوں $k \notin Z_k$۔ پر ایہ
وکھاؤندا اے کہ $\overline{Z} \neq Z_k$، کیوں جے $k$ اک دا
!!a{element} اے پر دوجے دا نہیں، تے ایس لئی $\overline{Z}$ تے $Z_k$
دے !!{element}s وکھ نیں۔ کیوں جے $k$ کوئی وی سی، ایس لئی
$\overline{Z}$ فہرست $Z_1$، $Z_2$، \dots وچ نہیں۔
\end{proof}

\begin{explain}
پچھلے ثبوت وچ قطر دا ذکر نہیں آیا، پر تسیں اینوں قطری ثبوت سمجھ سکدے
او جے ایس دی تصویر ایویں بناؤ: تصور کرو کہ سیٹ $Z_1$، $Z_2$، \dots
اک جدول (array) وچ لکھے نیں، جتھے ہر !!{element}~$j \in Z_i$ نوں
~$j$ویں کالم (column)
وچ لکھیا گیا اے۔ من لو کہ اوس فہرست دے پہلے چار سیٹ
$\{1,2,3,\dots\}$، $\{2, 4, 6, \dots\}$، $\{1,2,5\}$، تے
$\{3,4,5,\dots\}$ نیں۔ تاں جدول ایویں شروع ہووے گا:
\[
\begin{array}{r@{}rrrrrrr}
  Z_1 = \{ & \mathbf{1}, & 2, & 3, & 4, & 5, & 6, & \dots\}\\
  Z_2 = \{ &  & \mathbf{2}, &  & 4, &  & 6, & \dots\}\\
  Z_3 = \{ & 1, & 2, &  &  & 5\phantom{,} &  & \}\\
  Z_4 = \{ &  &  & 3, & \mathbf{4}, & 5, & 6, & \dots\}\\
  \vdots & & & & & \ddots
\end{array}
\]
فیر $\overline{Z}$ اوہ سیٹ اے جیہڑا قطر (diagonal) دے نال تھلے جا کے، قطر اُتے
آون والے عدد چھڈ کے، تے اوہ $j$ شامل کر کے ملدا اے جتھے جدول دی
$j$ویں قطار/کالم (row/column) وچ کھلوت اے۔ اُتلی صورت وچ اسی $1$ تے $2$ چھڈ
دیاں گے،~$3$ شامل کراں گے،~$4$ چھڈ دیاں گے، تے ایویں اگے۔
\end{explain}

\begin{prob}
قطری دلیل \pnbLatin{(diagonal argument)} نال وکھاؤ کہ $\Pow{\Nat}$،
!!{nonenumerable} اے۔
\end{prob}

\begin{prob}\label{sfr:siz:nen:prob:f-posint}
کھلی قطری دلیل \pnbLatin{(explicit diagonal argument)} نال وکھاؤ کہ فنکشناں
$f \colon \PosInt \to \PosInt$
دا سیٹ !!{nonenumerable} اے۔ یعنی وکھاؤ کہ جے $f_1$، $f_2$، \dots
فنکشناں دی فہرست اے تے ہر $f_i\colon \PosInt \to \PosInt$، تاں کوئی
$\overline{f}\colon \PosInt \to \PosInt$ ایہو جہا اے جیہڑا ایس فہرست
وچ نہیں۔
\end{prob}
\endgroup
% /OLP-0033

% OLP-0034 content/sets-functions-relations/size-of-sets/reduction.tex
\begingroup
\def\olfilename{reduction}
\expandafter\def\csname ol@sectioncs\endcsname{section}

\olfileid{sfr}{siz}{red}

\olsection{تقلیل \pnbLatin{(Reduction)}}

\begin{editorial}
  ایہ حصہ تقلیل (reduction) نال گنتی وار فہرست نہ بنن جوگ ہون
  (non-enumerability) نوں ثابت کردا اے،
  تے \olref[nen]{sec} دے نتیجیاں نال میل کھاندا اے۔ اک متبادل، ذرا
  مختصر روپ جیہڑا \olref[nen-alt]{sec} دے نتیجیاں نال میل کھاندا اے،
  \olref[red-alt]{sec} وچ دتا گیا اے۔
\end{editorial}

اسی قطری دلیل \pnbLatin{(diagonal argument)} نال وکھایا سی کہ $\Pow{\PosInt}$،
یعنی مثبت صحیح عدداں دے ذیلی سیٹاں دا سیٹ \pnbLatin{(power set)}، گنتی وار فہرست
نہ بنن جوگ (!!{nonenumerable}) اے۔ ساڈے کول پہلے ای اک ثبوت سی کہ
$\Bin^\omega$، یعنی $0$ تے $1$ دیاں ساریاں لامتناہی لڑیاں
\pnbLatin{(infinite sequences)} دا سیٹ، !!{nonenumerable} اے۔ $\Pow{\PosInt}$
دے !!{nonenumerable} ہون نوں ثابت کرن دا اک ہور طریقہ ایہ اے: وکھاؤ
کہ \emph{جے $\Pow{\PosInt}$، گنتی وار فہرست بنن جوگ
(!!{enumerable}) اے تاں $\Bin^\omega$ وی !!{enumerable} اے}۔ کیوں جے
اسی جاندے آں کہ $\Bin^\omega$، !!{enumerable} نہیں، ایس لئی
$\Pow{\PosInt}$ وی نہیں ہو سکدا۔ اینوں اک مسئلے نوں دوجے مسئلے ول
\emph{تقلیل دینا} \pnbLatin{(reducing one problem to another)} آکھیا
جاندا اے---ایتھے اسی $\Bin^\omega$ دی گنتی وار فہرست بناؤن دے مسئلے
نوں $\Pow{\PosInt}$ دی گنتی وار فہرست بناؤن دے مسئلے ول تقلیل دیندے
آں۔ دوجے مسئلے دا حل---$\Pow{\PosInt}$ دی اک گنتی وار فہرست---پہلے
مسئلے دا حل---$\Bin^\omega$ دی اک گنتی وار فہرست---دے دیوے گا۔

اسی سیٹ~$B$ دی گنتی وار فہرست بناؤن دے مسئلے نوں سیٹ~$A$ دی گنتی
وار فہرست بناؤن دے مسئلے ول کِویں گھٹائیے؟ اسی~$A$ دی اک گنتی وار
فہرست نوں~$B$ دی گنتی وار فہرست وچ بدلّن دا طریقہ دیندے آں۔ ایہ کرن
دا سب توں سوکھا طریقہ اک ہر ہدف قدر لین والا (!!a{surjective}) فنکشن
(function)
$f\colon A \to B$ تعریف کرنا اے۔ جے $x_1$، $x_2$، \dots{}،
~$A$ دی گنتی وار فہرست اے، تاں $f(x_1)$، $f(x_2)$، \dots{}،
~$B$ دی گنتی وار فہرست (enumeration) ہووے گی۔ ساڈی صورت وچ اسی ہر
ہدف قدر لین والا فنکشن $f\colon \Pow{\PosInt} \to \Bin^\omega$ لبھ
رہے آں؛ ایہ ڈومین توں ہدف ول تقلیل دی سمت \pnbLatin{(direction of reduction)}
نوں پکا کردا اے۔

\begin{prob}
وکھاؤ کہ جے کوئی اک بہ اک (!!{injective}) فنکشن $g\colon B \to A$
موجود اے، تے~$B$، گنتی وار فہرست نہ بنن جوگ (!!{nonenumerable}) اے،
تاں~$A$ وی ایویں ای اے۔ ایہ وکھا کے
کرو کہ تسیں~$g$ نوں~$A$ دی اک گنتی وار فہرست توں~$B$ دی گنتی وار
فہرست بناؤن لئی کِویں ورت سکدے او۔
\end{prob}

\begin{proof}[{\olref[nen]{thm:nonenum-pownat}} دا تقلیلی ثبوت]
ایہ تقلیلی ثبوت \pnbLatin{(proof by reduction)} اے۔ فرض کرو کہ
$\Pow{\PosInt}$، !!{enumerable} ہوندا، تے ایس لئی اوہدی اک
گنتی وار فہرست $Z_{1}$، $Z_{2}$، $Z_{3}$، \dots ہوندی۔

فنکشن $f \colon \Pow{\PosInt} \to \Bin^\omega$ ایویں تعریف کرو کہ
$f(Z)$ اوہ خاصیتی ثنائی لڑی \pnbLatin{(characteristic binary sequence)} $s$ ہووے
جیہدے لئی $s(n) = 1$ عین اوس ویلے جدوں
$n \in Z$، تے ہور صورت وچ $s(n) = 0$ ہووے۔
% OLSIZ-004 ماخذ دی درستی: نتیجے والی لڑی دا ناں s اے؛ انگریزی عبارت وچ غیر متعین k والے s_k تے s_k(n) ورتے گئے سن۔
ایہ صاف طور تے اک فنکشن تعریف کردا اے، کیوں جے جدوں وی
$Z \subseteq \PosInt$ ہووے، ہر $n \in \PosInt$ یا تاں $Z$ دا
عنصر (!!a{element}) اے یا نہیں۔ مثال لئی، مثبت جفت عدداں \pnbLatin{(positive even numbers)} دے سیٹ
$2\PosInt = \{2, 4, 6, \dots\}$ نوں لڑی $010101\dots$ ول، خالی سیٹ
\pnbLatin{(empty set)}
نوں $0000\dots$ ول، تے خود سیٹ $\PosInt$ نوں $1111\dots$ ول بھیجیا
جاندا اے۔

ایہ ہر ہدف قدر لین والا (!!{surjective}) وی اے: $0$ تے $1$ دی ہر لڑی
مثبت صحیح عدداں \pnbLatin{(positive integers)} دے کسے سیٹ نال میل کھاندی اے،
یعنی اوس سیٹ نال جیہدے رکن (members) اوہ صحیح عدد نیں جیہڑے لڑی
وچ~$1$ والے مقاماں (positions) دے مطابق نیں۔ ہور ٹھیک طور تے،
فرض کرو $s \in \Bin^\omega$۔ $Z \subseteq \PosInt$ نوں ایویں تعریف کرو:
\[
Z = \Setabs{n \in \PosInt}{s(n) = 1}
\]
فیر $f(Z) = s$، جیویں~$f$ دی تعریف نوں ویکھ کے جانچیا جا سکدا اے۔

ہُن فہرست اُتے غور کرو:
\[
f(Z_1), f(Z_2), f(Z_3), \dots
\]
کیوں جے $f$، ہر ہدف قدر لین والا (!!{surjective}) اے، ایس لئی
$\Bin^\omega$ دا ہر رکن ~$f$ دی کسے دلیل (argument) لئی قدر لازماً اے،
تے سو فہرست وچ لازماً آوے گا۔
ایس لئی ایہ فہرست~$\Bin^\omega$ دے ہر رکن دی گنتی وار
فہرست اے۔

سو جے $\Pow{\PosInt}$، !!{enumerable} ہوندا، تاں $\Bin^\omega$ وی
!!{enumerable} ہوندا۔ پر $\Bin^\omega$، گنتی وار فہرست نہ بنن جوگ
(!!{nonenumerable}) اے
(\olref[nen]{thm:nonenum-bin-omega})۔ ایس لئی $\Pow{\PosInt}$،
!!{nonenumerable} اے۔
\end{proof}

\begin{explain}
تقلیل دی سمت \pnbLatin{(direction of reduction)}، یعنی تقلیل کس پاسے جاندی اے،
ایس بارے بھلیکھا پینا سوکھا اے۔ مثال لئی، کوئی ہر ہدف قدر لین والا
(!!a{surjective}) فنکشن $g \colon \Bin^\omega \to B$ ایہ
\emph{نہیں} ثابت کردا کہ $B$، !!{nonenumerable} اے۔ (فنکشن
$g \colon \Bin^\omega \to \Bin$ اُتے غور کرو جیہڑا
$g(s) = s(1)$ نال تعریف اے، یعنی اوہ فنکشن جیہڑا $0$ تے $1$ دی
اک لڑی نوں اوہدے پہلے عنصر (!!{element}) ول بھیجدا اے۔ ایہ
!!{surjective}
اے، کیوں جے کجھ لڑیاں $0$ توں تے کجھ $1$ توں شروع ہوندیاں نیں۔ پر
$\Bin$ متناہی (finite) اے۔) ایہ وی یاد رکھو کہ فنکشن~$f$ دا
!!{surjective}
ہونا لازمی اے، نہیں تاں دلیل نہیں چلدی: $f(x_1)$، $f(x_2)$،
\dots{} وچ~$B$ دے سارے !!{element}s دے شامل ہون دی ضمانت نہیں
ہو سکدی۔ مثال لئی،
\[
h(n) = \underbrace{000\dots0}_{\text{$n$ $0$'s}}111\dots
\]
% OLSIZ-005 ماخذ دی درستی: ہر نتیجہ n صفراں توں بعد اکاں دی لامتناہی پونچھ رکھدا اے؛ انگریزی نمائش صرف متناہی صفر لکھ کے اوہنوں لامتناہی لڑی والے ہدف وچ رکھدی سی۔
ایتھے اکاں دی لامتناہی ثنائی پونچھ \pnbLatin{(infinite binary tail)} ہر نتیجے نوں
ہدف سیٹ دی لامتناہی لڑی بناؤندی اے۔ ایہ اک فنکشن $h\colon \PosInt \to
\Bin^\omega$ تعریف کردا اے، پر $\PosInt$، !!{enumerable} اے۔
\end{explain}

\begin{prob}
تقلیلی دلیل \pnbLatin{(reduction argument)} نال وکھاؤ کہ مثبت صحیح عدداں دیاں
جوڑیاں \pnbLatin{(pairs of positive integers)} دے سارے \emph{سیٹاں دے} سیٹ،
!!{nonenumerable} اے۔
\end{prob}

\begin{prob}\label{sfr:siz:red:prob:nat-nat}
  تقلیلی دلیل \pnbLatin{(reduction argument)} نال وکھاؤ کہ سیٹ~$X$، جیہڑا سارے
  فنکشناں (functions)
  $f\colon \Nat \to \Nat$ دا سیٹ اے، !!{nonenumerable} اے (اشارا:
  $X$ توں~$\Bin^\omega$ ول
  اک ہر ہدف قدر لین والا فنکشن \pnbLatin{(surjective function)} دیو۔)
\end{prob}

\begin{prob}
وکھاؤ کہ $\Nat^\omega$، یعنی قدرتی عدداں \pnbLatin{(natural numbers)} دیاں
لامتناہی لڑیاں \pnbLatin{(infinite sequences)} دا سیٹ، تقلیلی دلیل \pnbLatin{(reduction argument)} نال !!{nonenumerable} اے۔
\end{prob}

\begin{prob}
$P$ مثبت صحیح عدداں دے سیٹ توں سیٹ $\{0\}$ ول فنکشناں (functions) دا سیٹ ہووے،
تے $Q$ مثبت صحیح عدداں دے سیٹ توں سیٹ $\{0\}$ ول
\emph{جزوی} فنکشناں \pnbLatin{(partial functions)} دا سیٹ ہووے۔ وکھاؤ کہ~$P$،
!!{enumerable} اے
تے~$Q$ نہیں۔ (اشارا: $\Bin^\omega$ دی گنتی وار فہرست بناؤن دے مسئلے
نوں~$Q$ دی گنتی وار فہرست بناؤن ول گھٹاؤ۔)
\end{prob}

\begin{prob}
$S$ مثبت صحیح عدداں \pnbLatin{(positive integers)} دے سیٹ توں سیٹ \{0,1\} ول
سارے ہر ہدف قدر لین والے (!!{surjective}) فنکشناں دا سیٹ ہووے،
یعنی~$S$ سارے !!{surjective}
~$f\colon \PosInt \to \Bin$ اُتے مشتمل اے۔ وکھاؤ کہ~$S$،
!!{nonenumerable} اے۔
\end{prob}

\begin{prob}
وکھاؤ کہ سارے حقیقی عدداں \pnbLatin{(real numbers)} دا سیٹ~$\Real$،
گنتی وار فہرست نہ بنن جوگ (!!{nonenumerable}) اے۔
\end{prob}
\endgroup
% /OLP-0034

% OLP-0035 content/sets-functions-relations/size-of-sets/equinumerous-sets.tex
\begingroup
\def\olfilename{equinumerous-sets}
\expandafter\def\csname ol@sectioncs\endcsname{section}

\olfileid{sfr}{siz}{equ}

\olsection{ہم تعدادیت \pnbLatin{(Equinumerosity)}}

ساڈے کول سیٹاں دے سائز (size) دا اک بدیہی تصور \pnbLatin{(intuitive notion)} اے،
جیہڑا متناہی (finite) سیٹاں لئی چنگی طرح کم کردا اے۔ پر لامتناہی
(infinite) سیٹاں دا کی؟ جے اسی \emph{کسے وی} سائز دے دو سیٹاں دے
سائز دا رسمی طور تے موازنہ کرن دا طریقہ بناؤنا چاہنے آں، تاں ایہ
تعریف کرن توں شروع کرنا چنگا خیال اے کہ سیٹ کدوں اکّو سائز دے ہوندے
نیں۔ فریگے دی مثال ویکھو:
\begin{quote}
  جے اک بیرا ایہ پکا کرنا چاہوے کہ اوہنے میز اُتے عین اوہنیاں چھریاں
  رکھیاں نیں جِنیاں پلیٹاں، تاں اوہنوں دوہاں وچوں کسے ٹولی نوں وی گنّن
  دی لوڑ نہیں۔ اوہ بس ہر پلیٹ دے سجے پاسے اک چھری رکھ دیوے، ایس طرح
  کہ میز اُتے ہر چھری کسے نہ کسے پلیٹ دے سجے پاسے پئی ہووے۔ ایس طرح
  پلیٹاں تے چھریاں اکّو اک نال میل کھاندیاں \pnbLatin{(uniquely correlated)}
  نیں، تے اوہ وی اوسے مکانی تعلق راہیں۔ \citep[\S70]{Frege1884}
\end{quote}
ایس عبارت دی سوجھ نوں اک رسمی تعریف \pnbLatin{(formal definition)} راہیں صاف
کیتا جا سکدا اے:

\begin{defn}\ollabel{comparisondef}
  $A$،~$B$ دے نال \emph{ہم تعداد (equinumerous)} اے، جیہنوں
  $\cardeq{A}{B}$ لکھدے آں، عین اوس ویلے جدوں کوئی اک بہ اک تے ہر
  ہدف قدر لین والا نقشہ (!!a{bijection}) $f \colon A \to B$ موجود
  ہووے۔
\end{defn}

\begin{prop}\ollabel{equinumerosityisequi}
ہم تعدادیت (equinumerosity) اک ہم ارزیت دا تعلق \pnbLatin{(equivalence relation)}
اے۔
\end{prop}

\begin{proof}
سانوں وکھاؤنا اے کہ ہم تعدادیت انعکاسی (reflexive)، متناظر
(symmetric)، تے متعدی (transitive) اے۔ من لو $A, B$ تے~$C$ سیٹ نیں۔

\emph{انعکاسیت (Reflexivity)۔} شناختی نقشہ \pnbLatin{(identity map)}
$\Id{A} \colon A \to A$، جتھے $\Id{A} (x) = x$ ہر $x \in A$ لئی،
اک بہ اک تے ہر ہدف قدر لین والا نقشہ
(!!a{bijection}) اے۔ ایس لئی $\cardeq{A}{A}$۔

\emph{تناظر (Symmetry)۔} فرض کرو $\cardeq{A}{B}$، یعنی کوئی اک بہ
اک تے ہر ہدف قدر لین والا نقشہ (!!a{bijection})
$f\colon A \to B$ موجود اے۔ کیوں جے $f$، !!{bijective} اے، ایس لئی
اوہدا اُلٹا فنکشن \pnbLatin{(inverse function)} $f^{-1}$ موجود اے تے اوہ وی
!!{bijective} اے۔ ایس لئی $f^{-1}\colon B \to A$ وی اک بہ اک تے ہر
ہدف قدر لین والا نقشہ (!!a{bijection}) اے، سو $\cardeq{B}{A}$۔

\emph{تعدیت (Transitivity)۔} فرض کرو کہ $\cardeq{A}{B}$ تے
$\cardeq{B}{C}$، یعنی اک بہ اک تے ہر ہدف قدر لین والے نقشے
(!!{bijection}s) $f\colon A \to B$ تے $g\colon B \to C$ موجود
نیں۔ فیر فنکشناں دی ترکیب (composition)
$\comp{f}{g}\colon A \to C$، !!{bijective} اے، سو
$\cardeq{A}{C}$۔
\end{proof}

\begin{prop}
جے $\cardeq{A}{B}$، تاں $A$ گنتی وار فہرست بنن جوگ
(!!{enumerable}) عین اوس ویلے اے جدوں~$B$ وی ایویں اے۔
\end{prop}

\begin{editorial}
اگلا ثبوت \olref[enm]{defn:enumerable} والی تعریف ورتدا اے، جے
\olref[enm]{sec} شامل کیتا گیا ہووے؛ نہیں تاں اوہ
\olref[enm-alt]{defn:enumerable} والی متبادل تعریف ورتدا اے۔
\end{editorial}

\begin{proof}
فرض کرو $\cardeq{A}{B}$، سو کوئی اک بہ اک تے ہر ہدف قدر لین والا
نقشہ (!!{bijection}) $f \colon A \to B$ موجود اے، تے فرض کرو کہ
$A$ گنتی وار فہرست بنن جوگ (!!{enumerable}) اے۔
\oliflabeldef{sfr:siz:enm:defn:enumerable}{
  فیر یا تاں خالی سیٹ \pnbLatin{(empty set)} والی حالت $A = \emptyset$ اے، یا
  کوئی ہر ہدف قدر لین والا (!!a{surjective}) فنکشن (function)
  $g\colon \PosInt \to A$ موجود اے۔ جے $A = \emptyset$، تاں
  $B = \emptyset$ وی اے (نہیں تاں کوئی عنصر (!!a{element})
  ~$y \in B$ ہوندا، پر کوئی $x \in A$ نہ ہوندا جیہدے لئی
  $f(x) = y$)۔ دوجے پاسے جے $g\colon \PosInt \to A$،
  !!{surjective} اے، تاں فنکشناں دی ترکیب (composition)
  $\comp{g}{f} \colon \PosInt \to B$ وی !!{surjective} اے۔ ایہ
  ویکھن لئی $y \in B$ لو۔ کیوں جے $f$، !!{surjective} اے، ایس لئی
  کوئی $x \in A$ اے جیہدے لئی $f(x) = y$۔ کیوں جے $g$،
  !!{surjective} اے، ایس لئی کوئی مثبت صحیح عدد \pnbLatin{(positive integer)}
  $n \in \PosInt$ اے جیہدے لئی $g(n) = x$۔ ایس لئی،
\[
(\comp{g}{f})(n) = f(g(n)) = f(x) = y
\]
تے سو $\comp{g}{f}$، !!{surjective} اے۔ ایس طرح $\comp{g}{f}$،
~$B$ دی اک گنتی وار فہرست (enumeration) اے، تے ایس لئی $B$ وی
گنتی وار فہرست بنن جوگ (!!{enumerable}) اے۔}
{
  فیر یا تاں $A = \emptyset$ اے یا کوئی اک بہ اک تے ہر ہدف قدر لین
  والا نقشہ (!!a{bijection})~$g$ موجود اے، جیہدی حاصل قدراں دا سیٹ
  (range)~$A$ اے تے جیہدا دائرۂ تعریف (domain) یا $\Nat$ اے یا قدرتی
  عدداں دا اک مُڈھلا سلسلہ \pnbLatin{(initial sequence)} اے۔ جے
  $A = \emptyset$، تاں $B = \emptyset$ وی اے (نہیں تاں کوئی
  ~$y \in B$ ہوندا جیہدے لئی کوئی $x \in A$ ای نہ ہوندا کہ
  $f(x) = y$)۔ سو فرض کرو ساڈے کول ایہ !!{bijection}~$g$ اے۔ فیر
  $\comp{g}{f}$ اک !!a{bijection} اے، جیہدی حاصل قدراں دا سیٹ~$B$
  اے تے جیہدا دائرۂ تعریف~$g$ والے دے برابر اے (یعنی یا $\Nat$ یا
  اوہدا اک مُڈھلا ٹکڑا \pnbLatin{(initial segment)})، سو $B$، !!{enumerable}
  اے۔}
% OLSIZ-006 ماخذ دی درستی: خالی سیٹ والیاں دوہاں مشروط شاخاں وچ f ہی زیرِ بحث بائجیکشن اے؛ انگریزی ماخذ اوتھے g(x)=y لکھدا اے، حالاں کہ g ہن تک متعارف نہیں ہویا یا B ول نہیں جاندا۔

جے $B$ گنتی وار فہرست بنن جوگ (!!{enumerable}) اے، تاں اسی حاصل
کردے آں کہ $A$ وی !!{enumerable} اے، دلیل نوں اُلٹی بائجیکشن
\pnbLatin{(inverse bijection)} !!{bijection} $f^{-1}\colon B \to A$ نال دُہرا
کے،~$f$ دی بجائے۔
\end{proof}

\begin{prob}
وکھاؤ کہ جے $\cardeq{A}{C}$ تے $\cardeq{B}{D}$ دتے ہوون، تے
$A \cap B = C \cap D = \emptyset$، یعنی دوہاں پاسیاں دے سیٹ باہم
الگ (disjoint) ہون، تاں $\cardeq{A \cup B}{C \cup D}$۔
\end{prob}

\begin{prob}
وکھاؤ کہ جے $A$ لامتناہی (infinite) تے گنتی وار فہرست بنن جوگ
(!!{enumerable}) اے، تاں اوہ قدرتی عدداں دے سیٹ نال ہم تعداد
(equinumerous) اے، یعنی $\cardeq{A}{\Nat}$۔
\end{prob}
\endgroup
% /OLP-0035

% OLP-0036 content/sets-functions-relations/size-of-sets/comparing-size.tex
\begingroup
\def\olfilename{comparing-size}
\expandafter\def\csname ol@sectioncs\endcsname{section}

\olfileid{sfr}{siz}{car}

\olsection{وکھ وکھ سائزاں دے سیٹ، تے کانتور دا قضیہ \pnbLatin{(Cantor's Theorem)}}

\begin{explain}
اسی پہلے ایہ خیال ٹھیک ٹھیک بیان کیتا اے کہ دو سیٹاں دا سائز اکّو
اے۔ اسی ایہ خیال وی ٹھیک ٹھیک بیان کر سکدے آں کہ اک سیٹ دوجے نالوں
چھوٹا اے۔ ”سائز وچ چھوٹا یا ہم تعداد \pnbLatin{(smaller than or equinumerous)}“
ہون دی ساڈی تعریف وچ، سیٹاں دے وچکار !!a{bijection} دی تھاں پہلے
سیٹ توں دوجے سیٹ ول !!a{injection} چاہیدی ہووے گی۔ جے ایہو جہا
فنکشن موجود ہووے، تاں پہلے سیٹ دا سائز دوجے سیٹ دے سائز توں گھٹ یا
اوہدے برابر اے۔ بدیہی طور اُتے، اک سیٹ توں دوجے سیٹ ول
!!a{injection} ایہ ضمانت دیندا اے کہ فنکشن دے حاصل قدراں دے سیٹ
(range) وچ گھٹ توں گھٹ اوہنے !!{element}s ہوندے نیں جِنّے دائرۂ
تعریف (domain) وچ، کیوں جے دائرۂ تعریف دے کوئی دو !!{element}s
حاصل قدراں دے اکو !!{element} ول نہیں بھیجے جاندے۔
\end{explain}

\begin{defn}
$A$،~$B$ نالوں \emph{سائز وچ وڈا نہیں \pnbLatin{(no larger than)}} اے، جیہنوں
$\cardle{A}{B}$ لکھدے آں، عین اوس ویلے جدوں کوئی !!a{injection}
$f \colon A \to B$ موجود ہووے۔
\end{defn}

ایہ تعلق صاف طور اُتے انعکاسی (reflexive) تے متعدی (transitive) اے،
پر متناظر (symmetric) نہیں؛ ایہ گل مشق لئی چھڈی گئی اے۔ اسی اک ہور
تصور وی متعارف کرا سکدے آں، جیہڑا دسدا اے کہ اک سیٹ دوجے نالوں سخت
معنی وچ چھوٹا اے۔

\begin{defn}
$A$،~$B$ نالوں \emph{سائز وچ چھوٹا \pnbLatin{(smaller than)}} اے، جیہنوں
$\cardless{A}{B}$ لکھدے آں، عین اوس ویلے جدوں کوئی
!!a{injection}~$f\colon A \to B$ موجود ہووے، پر کوئی
!!{bijection}~$g\colon A \to B$ موجود نہ ہووے؛ یعنی
$\cardle{A}{B}$ تے $\cardneq{A}{B}$۔
\end{defn}

ایہ تعلق صاف طور اُتے غیر انعکاسی (irreflexive) تے متعدی
(transitive) اے؛ ایہ گل مشق لئی چھڈی گئی اے۔ ایس علامت نال اسی آکھ
سکدے آں کہ سیٹ~$A$، !!{enumerable} (گنتی وار فہرست بنن جوگ) عین اوس
ویلے اے جدوں $\cardle{A}{\Nat}$، تے~$A$، !!{nonenumerable} (گنتی
وار فہرست نہ بنن جوگ) عین اوس ویلے اے جدوں
$\cardless{\Nat}{A}$۔ ایس نال اسی
\oliflabeldef{sfr:siz:nen-alt:thm:nonenum-pownat}{%
\olref[sfr][siz][nen-alt]{thm:nonenum-pownat}
نوں ایہ مشاہدہ آکھ کے فیر بیان کر سکدے آں کہ
$\cardless{\Nat}{\Pow{\Nat}}$}{%
\olref[sfr][siz][nen]{thm:nonenum-pownat}
نوں ایہ مشاہدہ آکھ کے فیر بیان کر سکدے آں کہ
$\cardless{\PosInt}{\Pow{\PosInt}}$}۔ اصل وچ، \citet{Cantor1892} نے
ثابت کیتا سی کہ ایہ آخری گل \emph{پوری طرح عام} اے:

\begin{thm}[کانتور (Cantor)]\ollabel{thm:cantor}
$\cardless{A}{\Pow{A}}$، کسے وی سیٹ~$A$ لئی۔
\end{thm}

\begin{proof}
نقشہ $f(x) = \{x\}$ اک !!a{injection} $f \colon A \to \Pow{A}$
اے، کیوں جے جے $x \neq y$، تاں عنصراں نال تعیّن دے اصول
(extensionality) موجب $\{x\} \neq \{y\}$ وی اے، تے ایس لئی
$f(x) \neq f(y)$۔ سو $\cardle{A}{\Pow{A}}$۔

\begin{editorial}
جے \olref[nen]{sec} موجود ہووے تاں اسی تفصیل نال چلن والا ثبوت پیش
کردے آں؛ نہیں تاں \olref[nen-alt]{sec} نال میل کھان والا ودھ تیز،
مختصر ثبوت۔
\end{editorial}

\oliflabeldef{sfr:siz:nen:sec}{%
  ہُن اسی ثابت کراں گے کہ کوئی !!a{surjective} فنکشن، یعنی ہر ہدف قدر
  لین والا فنکشن،~$g\colon A \to \Pow{A}$ موجود نہیں ہو سکدا؛
  !!a{bijective} ہونا تاں دور دی گل اے۔ ایس لئی
  $\cardneq{A}{\Pow{A}}$۔ فرض کرو $g\colon A \to \Pow{A}$ اک فنکشن
  اے۔ کیوں جے $g$ کُل (total) اے، ہر $x \in A$ نوں اک ذیلی سیٹ
  $g(x) \subseteq A$ ول بھیجیا جاندا اے۔ اسی ثابت کر سکدے آں کہ $g$
  ہر ہدف قدر لین والا نہیں ہو سکدا۔ ایس لئی اسی اک ذیلی سیٹ
  ~$\overline{A} \subseteq A$ تعریف کردے آں، جیہڑا اپنی تعریف موجب
  ~$g$ دے حاصل قدراں دے سیٹ (range) وچ نہیں ہو سکدا۔ رکھو
  \[
  \overline{A} = \Setabs{x \in A}{x \notin g(x)}.
  \]
  کیوں جے $g(x)$ ہر $x \in A$ لئی تعریف شدہ اے، $\overline{A}$ صاف
  طور اُتے~$A$ دا اک درست تعریف شدہ ذیلی سیٹ اے۔ پر ایہ~$g$ دے حاصل
  قدراں وچ نہیں ہو سکدا۔ کوئی من مانا $x \in A$ لو؛ اسی ثابت کراں
  گے کہ $\overline{A} \neq g(x)$۔ جے $x \in g(x)$، تاں اوہ
  $x \notin g(x)$ شرط پوری نہیں کردا، تے~$\overline{A}$ دی تعریف
  موجب $x \notin \overline{A}$۔ جے $x \in \overline{A}$، تاں اوہ
  ~$\overline{A}$ دی تعریفی خاصیت، یعنی $x \in A$ تے
  $x \notin g(x)$، پوری کردا اے۔ کیوں جے $x$ من مانا سی، ایہ ثابت
  کردا اے کہ ہر $x \in A$ لئی $x \in g(x)$ عین اوس ویلے اے جدوں
  $x \notin \overline{A}$، تے ایس لئی $g(x) \neq \overline{A}$۔
  ہور لفظاں وچ، $\overline{A}$،~$g$ دے حاصل قدراں وچ نہیں ہو سکدا؛
  ایہ~$g$ دے ہر ہدف قدر لین والے ہون دے مفروضے نال تضاد اے۔
  % OLSIZ-007 ماخذ دی درستی: من مانا x، سیٹ A وچوں چُنیا گیا اے؛ نتیجہ A دے ہر x لئی اے، نہ کہ صرف خط والے A دے رکناں لئی۔
}{ایہ ثابت کرنا باقی اے کہ $\cardneq{A}{\Pow{A}}$۔ برہانِ خلف
  (reductio) لئی فرض کرو $\cardeq{A}{\Pow{A}}$، یعنی کوئی
  !!{bijection} $g \colon A \to \Pow{A}$ موجود اے۔ ہُن غور کرو:
  \[
    D = \Setabs{x \in A}{x \notin g(x)}
  \]
  دھیان دیو کہ $D \subseteq A$، ایس لئی $D \in \Pow{A}$۔ کیوں جے
  $g$ اک !!a{bijection} اے، کوئی $y \in A$ موجود اے جیہدے لئی
  $g(y) = D$۔ پر فیر:
  \[
    y \in g(y) \text{ عین اوس ویلے جدوں } y \in D
    \text{ عین اوس ویلے جدوں } y \notin g(y).
  \]
  ایہ تضاد اے؛ ایس لئی $\cardneq{A}{\Pow{A}}$۔}{}
\end{proof}

\begin{explain}
\oliflabeldef{sfr:siz:nen:thm:nonenum-pownat}{\olref{thm:cantor} دے
  ثبوت نوں \olref[nen]{thm:nonenum-pownat} دے ثبوت نال ملا کے ویکھنا
  سبق آموز اے۔ اوتھے اسی ثابت کیتا سی کہ ہر فہرست $Z_1$، $Z_2$،
  \dots{}، جیہڑی~$\PosInt$ دے ذیلی سیٹاں دی اے، اوہدے توں عدداں دا
  اک سیٹ~$\overline{Z}$ بنایا جا سکدا اے جیہڑا لازماً اوس فہرست وچ
  نہیں۔
  ایس دی ضمانت ایہ گل دیندی سی کہ ہر $n \in \PosInt$ لئی،
  $n \in Z_n$ عین اوس ویلے اے جدوں $n \notin \overline{Z}$۔ ایس
  طرح ہمیشہ کوئی عدد ہوندا اے جیہڑا $Z_n$ تے~$\overline{Z}$ وچوں
  اک دا !!a{element} اے، پر دوجے دا نہیں۔ ایتھے اسی اوسے خیال اُتے
  چلدے آں، پر ہُن اشاریے (indices)~$n$،~$A$ دے !!{element}s نیں،
  ~$\PosInt$ دی تھاں۔ سیٹ~$\overline{A}$ نوں ایویں تعریف کیتا گیا
  اے کہ اوہ~$g(x)$ توں وکھرا اے، ہر $x \in A$ لئی، کیوں جے
  $x \in g(x)$ عین اوس ویلے اے جدوں $x \notin \overline{A}$۔ فیر
  ہمیشہ~$A$ دا کوئی !!a{element} ہوندا اے جیہڑا $g(x)$ تے
  ~$\overline{A}$ وچوں اک دا !!a{element} اے، پر دوجے دا نہیں۔
  جیویں~$\overline{Z}$ فہرست $Z_1$، $Z_2$، \dots{} وچ نہیں ہو سکدا،
  اویویں~$\overline{A}$،~$g$ دے حاصل قدراں وچ نہیں ہو سکدا۔}{}

\oliflabeldef{sfr:siz:nen-alt:thm:nonenum-pownat}{\olref{thm:cantor}
  دے ثبوت نوں \olref[nen-alt]{thm:nonenum-pownat} دے ثبوت نال ملا کے
  ویکھنا سبق آموز اے۔ اوتھے اسی ثابت کیتا سی کہ ہر فہرست $N_0$،
  $N_1$، $N_2$، \dots{}، جیہڑی~$\Nat$ دے ذیلی سیٹاں دی اے، اوہدے
  توں عدداں دا اک سیٹ~$D$ بنایا جا سکدا اے جیہڑا لازماً اوس فہرست وچ
  نہیں۔ ایس دی ضمانت ایہ گل دیندی سی کہ $n \in N_n$ عین اوس ویلے اے
  جدوں $n \notin D$، ہر $n \in \Nat$ لئی۔ ایتھے اسی اوسے خیال اُتے
  چلدے آں، پر ہُن اشاریے (indices)~$n$،~$A$ دے !!{element}s نیں،
  ~$\Nat$ دی تھاں۔ سیٹ~$D$ نوں ایویں تعریف کیتا گیا اے کہ اوہ
  ~$g(x)$ توں وکھرا اے، ہر $x \in A$ لئی، کیوں جے $x \in g(x)$ عین
  اوس ویلے اے جدوں
  $x \notin D$۔}{}

ایس ثبوت نوں رسل دے پیراڈوکس \pnbLatin{(Russell's Paradox)} دے ثبوت،
\olref[sfr][set][rus]{thm:russells-paradox}، نال ملا کے ویکھنا وی فائدہ
مند اے۔ اصل وچ کانتور دے قضیے نے رسل نوں اپنا پیراڈوکس پیش کرن دی
سُوجھ دتی سی۔
\end{explain}

\begin{prob}
  ثابت کرو کہ کوئی !!a{injection} $g\colon \Pow{A} \to A$ موجود نہیں
  ہو سکدا، کسے وی سیٹ~$A$ لئی۔ اشارہ: فرض کرو کہ
  $g\colon \Pow{A} \to A$، !!{injective}، یعنی اک بہ اک، اے۔
  $D = \Setabs{g(B)}{B \subseteq A \text{ تے } g(B) \notin B}$
  اُتے غور کرو تے $x = g(D)$ رکھو۔ ایس حقیقت نوں ورتو کہ~$g$،
  !!{injective} اے، تے تضاد حاصل کرو۔
\end{prob}
\endgroup
% /OLP-0036

% OLP-0037 content/sets-functions-relations/size-of-sets/schroder-bernstein.tex
\begingroup
\def\olfilename{schroder-bernstein}
\expandafter\def\csname ol@sectioncs\endcsname{section}

\olfileid{sfr}{siz}{sb}

\olsection{سائز \pnbLatin{(size)} دا تصور تے شروڈر-برنسٹین \pnbLatin{(Schröder-Bernstein)}}

\begin{explain}
ایہ اک بدیہی خیال اے: جے $A$،~$B$ نالوں سائز وچ وڈا نہیں
\pnbLatin{(no larger than)} اے، تے $B$،~$A$ نالوں سائز وچ وڈا نہیں اے، تاں
$A$ تے~$B$ ہم تعداد (equinumerous) نیں۔ صاف گل ایہ اے کہ جے
ایہ خیال \emph{غلط} ہوندا، تاں اسی ایہ سوچ مشکل نال ای جائز ٹھہرا
سکدے کہ ہم تعدادیت (equinumerosity) بارے ساڈا تعریف کیتا تصور سیٹاں
دے ”سائز“ دے موازنے نال کوئی واسطہ رکھدا اے!{} پر خوش قسمتی نال ایہ
بدیہی خیال درست اے۔ ایس گل دا جواز شروڈر-برنسٹین دے قضیے نال ملدا اے۔
\end{explain}

\begin{thm}[شروڈر-برنسٹین (Schröder-Bernstein)]
	\ollabel{thm:schroder-bernstein}
	جے $\cardle{A}{B}$ تے $\cardle{B}{A}$،
	تاں $\cardeq{A}{B}$۔
\end{thm}

\begin{explain}
ہور لفظاں وچ، جے $A$ توں~$B$ ول اک بہ اک نقشہ
(!!a{injection}) موجود اے، تے $B$ توں~$A$ ول اک بہ اک نقشہ
(!!a{injection}) موجود اے، تاں $A$ توں~$B$ ول اک بہ اک تے ہر
ہدف قدر لین والا نقشہ (!!a{bijection}) موجود اے۔

پر ایس نتیجے نوں ثابت کرنا واقعی خاصا \emph{اوکھا} اے۔ واقعی، بھاویں
کانتور نے ایہ نتیجہ بیان کیتا سی، پر اوہنوں ثابت ہورناں نے کیتا۔
\footnote{تاریخ بارے ہور جانن لئی، مثلاً،
\citet[pp.~165--6]{Potter2004} ویکھو۔}
\oliflabeldef{sfr:cardinals:card-sb:sec}{اسی شروڈر-برنسٹین دا قضیہ
\olref[sfr][cardinals][card-sb]{sec} وچ جا کے ای \emph{ثابت} کر سکاں
گے۔}{}%
فی الحال تسی ایس نتیجے نوں بھروسے اُتے منّ سکدے او (تے منّنا وی پئے
گا)۔

خوش قسمتی نال شروڈر-برنسٹین دا قضیہ \emph{درست} اے، تے ایہ ساڈی ایس
سوچ نوں جائز ٹھہراؤندا اے کہ ساڈے تعریف کیتے تعلق، یعنی
$\cardeq{A}{B}$ تے $\cardle{A}{B}$، سیٹاں دے ”سائز“ نال کوئی واسطہ
رکھدے نیں۔ نال ای، شروڈر-برنسٹین دا قضیہ بہت \emph{کارآمد} اے۔ دو
ہم تعداد (equinumerous) سیٹاں دے وچکار اک بہ اک تے ہر ہدف قدر لین
والا نقشہ (!!a{bijection}) سوچنا اوکھا ہو سکدا اے۔ شروڈر-برنسٹین دا
قضیہ سانوں موازنے نوں دو رخاں وچ ونڈن دیندا اے، ایس لئی سانوں صرف
پہلے سیٹ توں دوجے ول اک بہ اک نقشہ (!!a{injection}) تے، اُلٹ رخ وچ،
دوجے توں پہلے ول اک بہ اک نقشہ سوچنا پیندا اے۔
\end{explain}
\endgroup
% /OLP-0037

\begin{editorial}
اگلے \olref[sfr][siz][enm-alt]{sec}،
\olref[sfr][siz][nen-alt]{sec} تے \olref[sfr][siz][red-alt]{sec}،
\olref[sfr][siz][enm]{sec}، \olref[sfr][siz][nen]{sec} تے
\olref[sfr][siz][red]{sec} دے متبادل روپ \pnbLatin{(alternative versions)} نیں۔
ایہ ٹم بٹن نے اپنی
اوپن سیٹ تھیوری \pnbLatin{(Open Set Theory)} دی کتاب وچ ورتن لئی لکھے سن۔ ایہ
تھوڑے ہور اگے دے درجے دے \pnbLatin{(slightly more advanced)} نیں تے گنتی وار
فہرست بنن جوگ ہون (enumerability) دی اک وکھری تعریف ورتدے نیں، جیہڑی
سیٹ تھیوری دے سیاق \pnbLatin{(set-theory context)} وچ ہور مناسب اے: یعنی $\Nat$
نال یا اوہدے کسے مُڈھلے ٹکڑے \pnbLatin{(initial segment)} نال اک بہ اک مطابقت
(bijection)، نہ کہ فہرست بن سکنا \pnbLatin{(being listable)} یا $\PosInt$ توں ہر
ہدف قدر لین والے فنکشن \pnbLatin{(surjective function)} دا حاصل قدراں دا سیٹ
(range) ہونا۔
% PNBSIZ-001 مقامی تدوینی درستی: ماخذ دے "a difference definition" نوں سیاق مطابق "a different definition" (اک وکھری تعریف) کیتا؛ منجمد ماخذ نہیں بدلیا۔
\end{editorial}

% OLP-0038 content/sets-functions-relations/size-of-sets/enumerability-alt.tex
\begingroup
\def\olfilename{enumerability-alt}
\expandafter\def\csname ol@sectioncs\endcsname{section}

\olfileid{sfr}{siz}{enm-alt}

\olsection{گنتی وار فہرستاں \pnbLatin{(enumerations)} تے گنتی وار فہرست بنن جوگ
(\usetoken{S}{enumerable}) سیٹ}

\begin{editorial}
  ایہ حصہ گنتی وار فہرستاں (enumerations) نوں~$\Nat$ دے مُڈھلے
  ٹکڑیاں \pnbLatin{(initial segments)} نال اک بہ اک تے ہر ہدف قدر لین والیاں
  مطابقتاں (bijections) دے طور اُتے تعریف کردا اے، جیویں سیٹ تھیوری
  وچ کیتا جاندا اے۔ ایس لئی ایہ \olref[enm]{sec} دیاں تعریفاں نال ذرا
  ٹکراؤندا اے، تے اوتھے دیاں ساریاں مثالاں دُہراؤندا اے۔ ایہ اوس
  حصے نالوں کجھ مختصر وی اے۔
\end{editorial}

اسی کسے متناہی (finite) سیٹ دے عنصراں (!!{element}s) نوں بس ترتیب
وار لکھ کے اوہ سیٹ دس سکدے آں۔ اسی ایویں کردے آں جدوں کوئی سیٹ ایس
طرح تعریف کریے:
\[
  A = \{a_1, a_2, \ldots, a_n\}.
\]
% PNBSIZ-003 ماخذ دی درستی: انگریزی ماخذ دا پہلا جملہ نحوی طور تے ٹُٹیا ہویا اے؛ اگلی نمائش تے وضاحت موجب مراد متناہی سیٹ دے رکناں نوں ترتیب وار لکھ کے اوہنوں دسنا اے۔
جے فرض کریے کہ عنصر (!!{element}s) $a_1$، \dots، $a_n$ سارے وکھرے
نیں، تاں ایہ~$A$ تے پہلے $n$ قدرتی عدداں \pnbLatin{(natural numbers)} $0$،
\dots، $n-1$ دے وچکار اک بہ اک تے ہر ہدف قدر لین والی مطابقت
(!!a{bijection}) بناؤندا اے۔ الٹ پاسے، کیوں جے ہر متناہی سیٹ وچ
صرف متناہی گنتی دے عنصر (!!{element}s) ہوندے نیں، ہر متناہی سیٹ نوں
ایہو جہی مطابقت وچ رکھیا جا سکدا اے۔ ہور لفظاں وچ، جے~$A$ متناہی
اے، تاں~$A$ تے $\{0, \dots, n-1\}$ دے وچکار اک بہ اک تے ہر ہدف
قدر لین والی مطابقت (!!a{bijection}) موجود اے، جتھے~$n$،~$A$ دے
عنصراں (!!{element}s) دی گنتی اے۔

جے اسی کجھ قسماں دے لامتناہی (infinite) سیٹاں نوں وی آون دیواں، تاں
اسی کجھ لامتناہی سیٹاں دی وی گنتی وار فہرست بنن دی اجازت دیواں گے۔
اسی ایہ گل ایویں ٹھیک طور اُتے کہہ سکدے آں کہ اک لامتناہی سیٹ دی
گنتی وار فہرست (enumeration) اوس سیٹ تے سارے~$\Nat$ دے وچکار اک بہ
اک تے ہر ہدف قدر لین والی مطابقت (!!a{bijection}) بناؤندی اے۔

\begin{defn}[گنتی وار فہرست (Enumeration)، سیٹ تھیوری دے لحاظ نال]
سیٹ~$A$ دی اک \emph{گنتی وار فہرست (enumeration)} اوہ اک بہ اک تے ہر
ہدف قدر لین والی مطابقت (!!a{bijection}) اے جیہدی حاصل قدراں دا سیٹ
(range)~$A$ ہووے تے جیہدا دائرۂ تعریف (domain) یا قدرتی عدداں دا اک
مُڈھلا سیٹ \pnbLatin{(initial set)} $\{0, 1, \ldots, n\}$ {یا} قدرتی عدداں دا
پورا سیٹ~$\Nat$ ہووے۔
\end{defn}

\begin{explain}
\emph{گنتی وار فہرست (enumeration)} لفظ دے ایس ورتاوے تھلے اک بدیہی خیال اے۔ ایہ
آکھنا کہ اسی سیٹ~$A$ دی گنتی وار فہرست بنا لئی اے، ایہ آکھنا اے کہ
کوئی اک بہ اک تے ہر ہدف قدر لین والی مطابقت (!!a{bijection})~$f$
موجود اے جیہڑی سانوں سیٹ~$A$ دے عنصراں نوں گن کے کڈھن دیندی اے۔
$0$واں عنصر $f(0)$ اے، پہلا $f(1)$ اے،
\ldots{}، $n$واں $f(n)$ اے، \ldots۔\footnote{ہاں، اسی گنتی
$0$ توں شروع کردے آں۔ بے شک اسی~$1$ توں وی شروع کر سکدے ساں۔ ایس
نال کوئی وڈا فرق نہ پیندا۔ سانوں بس~$\Nat$ دی تھاں~$\PosInt$ رکھنا
پیندا۔} ایس دی وجہ اگلی گل نوں شامل کر کے ہور وی صاف کیتی جا سکدی اے:
\end{explain}

\begin{defn}
  \ollabel{defn:enumerable}
  سیٹ~$A$، گنتی وار فہرست بنن جوگ (!!{enumerable}) عین اوس ویلے اے جدوں یا تاں
  $A = \emptyset$ ہووے یا~$A$ دی کوئی گنتی وار فہرست موجود ہووے۔
  اسی~$A$ نوں گنتی وار فہرست نہ بنن جوگ (!!{nonenumerable}) عین اوس
  ویلے آکھدے آں جدوں~$A$، گنتی وار فہرست بنن جوگ (!!{enumerable}) نہ
  ہووے۔
\end{defn}

\begin{explain}
سو اک سیٹ گنتی وار فہرست بنن جوگ (!!{enumerable}) عین اوس ویلے اے
جدوں اوہ خالی ہووے یا تسی اک گنتی وار فہرست نال اوہدے عنصر
(!!{element}s) گن کے کڈھ سکو۔
\end{explain}

\begin{ex}
قدرتی عدداں دی گنتی وار فہرست بناؤن والا فنکشن بس شناختی فنکشن
\pnbLatin{(identity function)}
$\Id{\Nat} \colon \Nat \to \Nat$ اے، جیہڑا $\Id{\Nat}(n) = n$ نال
دتا گیا اے۔ \emph{مثبت} قدرتی عدداں، $\Nat^+ =
\Nat \setminus \{0\}$، دی گنتی وار فہرست بناؤن والا فنکشن
$g(n) = n + 1$ اے، یعنی جانشین فنکشن \pnbLatin{(successor function)}۔
\end{ex}

\begin{prob}
وکھاؤ کہ سیٹ~$A$، گنتی وار فہرست بنن جوگ (!!{enumerable}) عین اوس
ویلے اے جدوں یا تاں $A = \emptyset$ ہووے یا کوئی ہر ہدف قدر لین والا
نقشہ (!!a{surjection}) $f\colon \Nat \to A$ موجود ہووے۔ وکھاؤ کہ
~$A$، گنتی وار فہرست بنن جوگ (!!{enumerable}) عین اوس ویلے اے جدوں
کوئی اک بہ اک نقشہ (!!a{injection}) $g\colon A \to \Nat$ موجود ہووے۔
\end{prob}

\begin{ex}
فنکشن $f\colon \Nat \to \Nat$ تے $g \colon \Nat \to \Nat$، جیہڑے
ایویں دتے گئے نیں:
\begin{align*}
  f(n) & = 2n \text{ تے}\\
  g(n) & = 2n+1
\end{align*}
ترتیب وار جفت (even) قدرتی عدداں تے طاق (odd) قدرتی عدداں دیاں گنتی
وار فہرستاں بناؤندے نیں۔ پر دوہاں وچوں کوئی وی ہر ہدف قدر لین والا
(!!{surjective}) نہیں، سو کوئی وی~$\Nat$ دی گنتی وار فہرست نہیں۔
\end{ex}

\begin{prob}
مربع عدداں \pnbLatin{(square numbers)} $1$، $4$، $9$، $16$، \dots دی گنتی وار
فہرست تعریف کرو۔
\end{prob}

\begin{ex}
$\lceil x \rceil$ نوں \emph{چھت فنکشن \pnbLatin{(ceiling function)}} منّو، جیہڑا
$x$ نوں اوس قدر توں گھٹ نہ ہون والے سب توں چھوٹے صحیح عدد تک اُتے پورا
کردا اے۔ فیر فنکشن
$f \colon \Nat \to \Int$، جیہڑا ایویں دتا گیا اے:
\[
  f(n) = (-1)^{n} \left\lceil\tfrac{n}{2}\right\rceil
\]
صحیح عدداں (integers) دے سیٹ~$\Int$ دی گنتی وار فہرست اگلے ڈھنگ نال
بناؤندا اے:
\[
\begin{array}{c c c c c c c c}
f(0) & f(1) & f(2) & f(3) & f(4) & f(5) & f(6) & \dots \\ \\
\big\lceil \tfrac{0}{2} \big\rceil & -\big\lceil \tfrac{1}{2}\big\rceil &  \big\lceil \tfrac{2}{2} \big\rceil & -\big\lceil \tfrac{3}{2} \big\rceil & \big\lceil \tfrac{4}{2} \big\rceil  & -\big\lceil \tfrac{5}{2}\big\rceil & \big\lceil \tfrac{6}{2} \big\rceil & \dots \\ \\
0 & -1 & 1 & -2 & 2 & -3 & 3& \dots
\end{array}
\]
دھیان دیو کہ $f$، مثبت تے منفی صحیح عدداں دے وچکار ”ٹپدے“ ہوئے
$\Int$ دیاں قدراں پیدا کردا اے۔ تسیں~$f$ نوں اگلی صورتاں نال تعریف
کیتا ہویا وی سمجھ سکدے او:
\[
f(n) = \begin{cases}
  \frac{n}{2} & \text{جے $n$ جفت (even) اے}\\
  -\frac{n+1}{2} & \text{جے $n$ طاق (odd) اے}
  \end{cases}
\]
\end{ex}

\begin{prob}
وکھاؤ کہ جے $A$ تے~$B$ گنتی وار فہرست بنن جوگ (!!{enumerable}) نیں،
تاں $A \cup B$ وی اے۔
\end{prob}

\begin{prob}
$n$ اُتے ریاضیاتی استقرا \pnbLatin{(mathematical induction)} نال وکھاؤ کہ جے
$A_1$، $A_2$، \dots، $A_n$ سارے گنتی وار فہرست بنن جوگ
(!!{enumerable}) نیں، تاں $A_1 \cup \dots \cup A_n$ وی اے۔
\end{prob}
\endgroup
% /OLP-0038
% OLP-0039 content/sets-functions-relations/size-of-sets/non-enumerability-alt.tex
\begingroup
\def\olfilename{non-enumerability-alt}
\expandafter\def\csname ol@sectioncs\endcsname{section}

\olfileid{sfr}{siz}{nen-alt}

\olsection{\printtoken{S}{nonenumerable} سیٹ \pnbLatin{(Non-enumerable Sets)}}

\begin{editorial}
  ایہ حصہ \olref[enm-alt]{sec} وچ دتیاں تعریفاں ورت کے $\Bin^\omega$ تے
  $\Pow{\Nat}$ دے گنتی وار فہرست نہ بنن جوگ ہون (non-enumerability)
  نوں ثابت کردا اے؛ یعنی~$\Nat$ نال اک بہ اک مطابقت (bijection) دی
  شرط رکھدا اے،~$\PosInt$ توں ہر ہدف قدر لین والے نقشے (surjection)
  دی تھاں۔
\end{editorial}

\begin{explain}
قدرتی عدداں دا سیٹ~$\Nat$ لامتناہی (infinite) اے۔ ایہ صاف طور تے
گنتی وار فہرست بنن جوگ (!!{enumerable}) وی اے۔ پر حیران کرن والی
حقیقت ایہ اے کہ \emph{گنتی وار فہرست نہ بنن جوگ
(!!{nonenumerable})} سیٹ وی نیں، یعنی ایہو جہے سیٹ جیہڑے گنتی وار
فہرست بنن جوگ (!!{enumerable}) نہیں
(ویکھو \olref[sfr][siz][enm-alt]{defn:enumerable})۔

ایہ گل حیران کر سکدی اے۔ آخرکار، لامتناہی سیٹ~$A$ لئی ایہ آکھنا کہ
اوہ گنتی وار فہرست نہ بنن جوگ (!!{nonenumerable}) اے، ایہ آکھنا اے
کہ کوئی اک بہ اک مطابقت
(!!{bijection}) $f \colon \Nat \to A$ \emph{موجود نہیں}؛ یعنی
~$\Nat$ دے لامتناہی عنصراں (!!{element}s) نوں~$A$ ول نقشہ بند کرن
والا کوئی فنکشن سارے~$A$ نوں پورا نہیں کردا۔ سو جے~$A$ گنتی وار
فہرست نہ بنن جوگ (!!{nonenumerable}) اے، تاں~$A$ دے
!!{element}s قدرتی عدداں نالوں ”ودھ“ نیں۔
% PNBSIZ-004 ماخذ دی درستی: قدرتی عدداں نال اک بہ اک مطابقت نہ ہون والی مساوات ایتھے صرف لامتناہی سیٹ A لئی اے؛ متناہی سیٹ وی گنتی وار فہرست بنن جوگ نیں۔

کسی سیٹ دے گنتی وار فہرست نہ بنن جوگ (!!{nonenumerable}) ہون نوں
ثابت کرن لئی تہانوں وکھاؤنا پئے گا کہ کوئی مناسب اک بہ اک مطابقت
(!!{bijection}) موجود نہیں ہو سکدی۔ ایہ کرن دا سب توں چنگا طریقہ ایہ
وکھاؤنا اے کہ~$A$ دے عنصراں (!!{element}s) نوں گنتی وار ترتیب دین دی
ہر کوشش گھٹ توں گھٹ اک !!{element} ضرور چھڈ دیندی اے؛ ایہ وکھاؤندا اے
کہ کوئی فنکشن $f\colon \Nat \to A$ ہر ہدف قدر لین والا
(!!{surjective}) نہیں۔ اینوں قائم کرن دی اک عام حکمتِ عملی کانتور دا
\emph{قطری طریقہ} \pnbLatin{(Cantor's diagonal method)} ورتنا اے۔~$A$ دے
!!{element}s دی کوئی فہرست (list)، مثلاً $x_1$، $x_2$، \dots، دتی
ہووے، تاں اسی~$A$ دا اک ہور !!{element} بناؤندے آں جیہڑا اپنی بناوٹ
کر کے اوس فہرست وچ ہو ای نہیں سکدا۔

پر ایہ ساری گل مثال نال سب توں چنگی طرح سمجھ آندی اے۔ سو ساڈی پہلی
مثال سیٹ~$\Bin^\omega$ اے، جیہڑا $0$ تے $1$ دیاں ساریاں لامتناہی
ثنائی لڑیاں \pnbLatin{(infinite binary strings)} دا سیٹ اے۔ ($\Bin$ توں مراد
ثنائی (binary) اے، تے اسی اینوں بس دو عنصراں والے سیٹ
$\{0,1\}$ دے طور تے سمجھ سکدے آں۔)\oliflabeldef{sfr:card-arithmetic:card-opps:sec}{\footnote{ہور
ٹھیک طور تے، سانوں ایہ شرط لانی چاہیدی اے کہ~$\Bin^\omega$،
$\omega$-لڑیاں (omega-sequences) دا اوہ سیٹ اے جیہڑیاں $0$ تے $1$
دیاں نیں، یعنی سیٹ $\funfromto{\omega}{\{0,1\}}$۔ پر ایس دا مطلب
\olref[sfr][card-arithmetic][card-opps]{sec} وچ جا کے صاف ہووے گا۔}}{}
پر فی الحال ایہ ذرا ڈھیلی عبارت کوئی بھلیکھا نہیں پاؤنی چاہیدی۔
\end{explain}

\begin{thm}
\ollabel{thm:nonenum-bin-omega}
$\Bin^\omega$، گنتی وار فہرست نہ بنن جوگ (!!{nonenumerable}) اے۔
\end{thm}

\begin{proof}
$\Bin^\omega$ دے کسے ذیلی سیٹ دی کوئی وی گنتی وار فہرست
(enumeration) لو۔ سو ساڈے کول کوئی فہرست (list) $s_{0}$، $s_{1}$،
$s_{2}$، \dots{} اے، جتھے ہر $s_n$، $0$ تے~$1$ دی اک لامتناہی
ثنائی لڑی \pnbLatin{(infinite binary string)} اے۔ آؤ $s_n(m)$ نوں ایس فہرست دی
$n$ویں لڑی دا $m$واں ہندسہ (digit) آکھیے۔ ہُن اسی اپنی فہرست نوں اک
جدول (array) سمجھ سکدے آں، جتھے $s_n(m)$ نوں $n$ویں قطار (row) تے
$m$ویں کالم (column) وچ رکھیا گیا اے:
% OLSIZ-008 ماخذ دی درستی: انگریزی نثر اشارییاں دے کردار اُلٹ دَسدی اے؛ اگلا جملہ تے جدول s_n(m) نوں nویں لڑی دا mواں ہندسہ مقرر کردے نیں۔
\[
\begin{array}{c|c|c|c|c|c}
& 0 & 1 & 2 & 3 & \dots \\\hline
0 & \mathbf{s_{0}(0)} & s_{0}(1) & s_{0}(2) & s_0(3) & \dots \\\hline
1 & s_{1}(0)& \mathbf{s_{1}(1)} & s_1(2) & s_1(3) & \dots \\\hline
2 & s_{2}(0)& s_{2}(1) & \mathbf{s_2(2)} & s_2(3) & \dots \\\hline
3 & s_{3}(0)& s_{3}(1) & s_3(2) & \mathbf{s_3(3)} & \dots \\\hline
\vdots & \vdots & \vdots & \vdots & \vdots & \mathbf{\ddots}
\end{array}
\]
ہُن اسی اک لامتناہی لڑی~$d$، جیہڑی $0$ تے $1$ دی اے، بناؤندے آں
جیہڑی ایس فہرست وچ نہیں۔ اسی اوہدی ہر قدر (entry) دَس کے ایہ کم کردے
آں، یعنی اسی $d(n)$ مقرر کردے آں، ہر $n \in \Nat$ لئی۔ سمجھ دے
پکھوں، اسی اُتلے جدول
دا قطر (diagonal) تھلے ول پڑھدے آں---ایسے لئی ناں ”قطری طریقہ“
\pnbLatin{(diagonal method)} اے---تے فیر ہر $1$ نوں $0$ تے ہر $0$ نوں~$1$ وچ
بدل دیندے آں۔
% OLSIZ-009 ماخذ دی درستی: انگریزی نثر 1 توں 0 والی تبدیلی دو واری آکھدی اے؛ اگلی حالت وار تعریف موجب دوجی تبدیلی 0 توں 1 اے۔
ہور مجرد طور تے، اسی $d(n)$ نوں $0$ یا $1$ ایس مطابق تعریف کردے آں
کہ قطر دا $n$واں عنصر (!!{element})، $s_n(n)$، $1$ اے یا $0$، یعنی:
\[
d(n) =
\begin{cases}
1 & \text{جے $s_{n}(n) = 0$ اے}\\
0 & \text{جے $s_{n}(n) = 1$ اے}
\end{cases}
\]
صاف اے کہ $d \in \Bin^\omega$، کیوں جے ایہ $0$ تے $1$ دی اک
لامتناہی لڑی اے۔ پر اسی~$d$ نوں ایویں بنایا اے کہ
$d(n) \neq s_n(n)$، ہر $n \in \Nat$ لئی۔ یعنی~$d$، $s_n$ نال
اوہدی $n$ویں قدر وچ وکھ اے۔ سو $d \neq s_n$، ہر $n\in \Nat$ لئی۔
ایس لئی $d$ فہرست $s_0$، $s_1$، $s_2$، \dots وچ نہیں ہو سکدا۔

اسی وکھا دتا اے کہ~$\Bin^\omega$ دے کسے ذیلی سیٹ دی من مانی گنتی
وار فہرست (enumeration) دتی ہووے، تاں اوہ~$\Bin^\omega$ دا کوئی عنصر
(!!{element}) چھڈ دیوے گی۔ سو سیٹ~$\Bin^\omega$ دی کوئی گنتی وار
فہرست نہیں، یعنی $\Bin^\omega$ گنتی وار فہرست نہ بنن جوگ
(!!{nonenumerable}) اے۔
\end{proof}

\begin{explain}
ایس ثبوت دے طریقے نوں ”قطری بناوٹ“ (diagonalization) آکھیا جاندا اے،
کیوں جے ایہ~$d$ نوں تعریف کرن لئی جدول (array) دا قطر (diagonal)
ورتدا اے۔ پر قطری بناوٹ وچ جدول دا موجود ہونا لازمی نہیں۔ اصل وچ،
اسی اوسے ورگا خیال ورت کے کسی سیٹ دا گنتی وار فہرست نہ بنن جوگ
(!!{nonenumerable}) ہونا وی وکھا سکدے آں، بھاویں کوئی جدول تے حقیقی
قطر شامل نہ ہووے۔ اگلا نتیجہ دَسدا اے کہ ایہ کِویں ہوندا اے۔
\end{explain}

\begin{thm}
\ollabel{thm:nonenum-pownat}
$\Pow{\Nat}$، گنتی وار فہرست بنن جوگ (!!{enumerable}) نہیں اے۔
\end{thm}

\begin{proof}
اسی اوسے ڈھنگ نال کم کردے آں: ایہ وکھاؤندے آں کہ~$\Nat$ دے ذیلی
سیٹاں (subsets) دی ہر فہرست (list)~$\Nat$ دا کوئی ذیلی سیٹ چھڈ دیندی
اے۔ سو فرض کرو ساڈے کول کوئی فہرست $N_0, N_1, N_2, \ldots$ اے،
جیہڑی~$\Nat$ دے ذیلی سیٹاں دی اے۔ اسی سیٹ~$D$ ایویں تعریف کردے آں
کہ $n \in D$ عین اوس ویلے جدوں $n \notin N_{n}$:
\[
D = \Setabs{n \in \Nat}{n \notin N_n}
\]
صاف اے کہ $D\subseteq \Nat$۔ پر~$D$ فہرست وچ نہیں ہو سکدا۔ بناوٹ
موجب $n \in D$ عین اوس ویلے جدوں $n\notin N_n$، ایس لئی
$D \neq N_n$، ہر $n \in \Nat$ واسطے۔
\end{proof}

\begin{explain}
پچھلے ثبوت وچ قطر دا ذکر نہیں آیا۔ فیر وی تسیں اینوں قطری دلیل
\pnbLatin{(diagonal argument)} سمجھ سکدے او جے ایس دی تصویر ایویں بناؤ: تصور
کرو کہ سیٹ $N_0$، $N_1$، \dots، اک جدول (array) وچ لکھے نیں، جتھے
~$N_n$ نوں $n$ویں قطار (row) وچ لکھن لئی $m$ نوں $m$ویں کالم
(column) وچ عین اوس ویلے لکھدے آں جدوں $m \in N_n$۔ مثال لئی، من لو
کہ ایس فہرست دے پہلے چار سیٹ
$\{0,1,2,\dots\}$، $\{1, 3, 5, \dots\}$، $\{0,1,4\}$، تے
$\{2,3,4,\dots\}$ نیں؛ تاں ساڈا جدول ایویں شروع ہووے گا:
\[
\begin{array}{r@{}rrrrrrr}
  N_0 = \{ & \mathbf{0}, & 1, & 2, & & & & \dots\}\\
  N_1 = \{ &  & \mathbf{1}, &  & 3, &  & 5, & \dots\}\\
  N_2 = \{ & 0, & 1, &  &  & 4\phantom{,} &  & \}\\
  N_3 = \{ &  &  & 2, & \mathbf{3}, & 4, & & \dots\}\\
  &\vdots & & & & & & \ddots\phantom{\}}
  \end{array}
\]
فیر~$D$ اوہ سیٹ اے جیہڑا قطر (diagonal) دے نال تھلے جا کے تے
$n \in D$ عین اوس ویلے رکھ کے ملدا اے جدوں~$n$ قطر اُتے
\emph{نہیں}۔ سو اُتلی صورت وچ اسی $0$ تے $1$ چھڈ دیاں گے،~$2$ شامل
کراں گے،~$3$ چھڈ دیاں گے، تے ایویں اگے۔
\end{explain}

\begin{prob}
کھلی قطری دلیل \pnbLatin{(explicit diagonal argument)} نال وکھاؤ کہ سارے فنکشناں
$f \colon \Nat \to \Nat$ دا سیٹ گنتی وار فہرست نہ بنن جوگ
(!!{nonenumerable}) اے۔ یعنی وکھاؤ کہ جے $f_1$، $f_2$، \dots،
فنکشناں دی فہرست اے تے ہر $f_i\colon \Nat \to \Nat$، تاں کوئی
$g \colon \Nat \to \Nat$ ایہو جہا اے جیہڑا ایس فہرست وچ نہیں۔
\end{prob}
\endgroup
% /OLP-0039
% OLP-0040 content/sets-functions-relations/size-of-sets/reduction-alt.tex
\begingroup
\def\olfilename{reduction-alt}
\expandafter\def\csname ol@sectioncs\endcsname{section}

\olfileid{sfr}{siz}{red-alt}

\olsection{تقلیل \pnbLatin{(Reduction)}}

\begin{editorial}
  ایہ حصہ تقلیل (reduction) نال گنتی وار فہرست نہ بنن جوگ ہون
  (non-enumerability) نوں ثابت کردا اے، تے \olref[nen-alt]{sec} دے
  نتیجیاں نال میل کھاندا اے۔ اک متبادل، ذرا ودھ تفصیلی روپ جیہڑا
  \olref[nen]{sec} دے نتیجیاں نال میل کھاندا اے، \olref[red]{sec}
  وچ دتا گیا اے۔
\end{editorial}

اسی قطری دلیل \pnbLatin{(diagonal argument)} نال وکھایا سی کہ $\Bin^\omega$،
گنتی وار فہرست نہ بنن جوگ (!!{nonenumerable}) اے۔ اسی اوسے ورگی قطری
دلیل نال وکھایا سی کہ $\Pow{\Nat}$، یعنی قدرتی عدداں دے ذیلی سیٹاں
دا سیٹ \pnbLatin{(power set)}، !!{nonenumerable} اے۔ پر $\Pow{\Nat}$ دے
!!{nonenumerable} ہون نوں ثابت کرن دا اک ہور طریقہ ایہ اے: وکھاؤ کہ
\emph{جے $\Pow{\Nat}$، گنتی وار فہرست بنن جوگ (!!{enumerable}) اے
تاں $\Bin^\omega$ وی !!{enumerable} اے}۔ کیوں جے اسی جاندے آں کہ
$\Bin^\omega$، !!{nonenumerable} اے، ایس لئی $\Pow{\Nat}$ وی ایویں
ای اے۔

اینوں اک مسئلے نوں دوجے مسئلے ول \emph{تقلیل دینا} \pnbLatin{(reducing one problem to another)} آکھیا جاندا اے۔ ایس صورت وچ اسی $\Bin^\omega$ دی
گنتی وار فہرست بناؤن دے مسئلے نوں $\Pow{\Nat}$ دی گنتی وار فہرست
بناؤن دے مسئلے ول تقلیل دیندے آں۔ دوجے مسئلے دا حل---$\Pow{\Nat}$ دی
اک گنتی وار فہرست---پہلے مسئلے دا حل---$\Bin^\omega$ دی اک گنتی وار
فہرست---دے دیوے گا۔

سیٹ~$B$ دی گنتی وار فہرست بناؤن دے مسئلے نوں سیٹ~$A$ دی گنتی وار
فہرست بناؤن دے مسئلے ول تقلیل دین لئی، اسی~$A$ دی اک گنتی وار فہرست
نوں~$B$ دی گنتی وار فہرست وچ بدلّن دا طریقہ دیندے آں۔ ایہ کرن دا سب
توں سوکھا طریقہ اک ہر ہدف قدر لین والا (!!a{surjection}) فنکشن
(function) $f\colon A \to B$ تعریف کرنا اے۔ جے $x_1$، $x_2$،
\dots{}،~$A$ دی گنتی وار فہرست اے، تاں $f(x_1)$، $f(x_2)$،
\dots{}،~$B$ دی گنتی وار فہرست (enumeration) ہووے گی۔ ساڈی صورت وچ
اسی اک ہر ہدف قدر لین والا (!!a{surjection}) فنکشن
$f\colon \Pow{\Nat} \to \Bin^\omega$ لبھ رہے آں؛ ایہ ڈومین توں ہدف
ول تقلیل دی سمت \pnbLatin{(direction of reduction)} نوں پکا کردا اے۔

\begin{prob}
وکھاؤ کہ جے کوئی اک بہ اک (!!{injective}) فنکشن $g\colon B \to A$
موجود اے، تے~$B$ گنتی وار فہرست نہ بنن جوگ (!!{nonenumerable}) اے،
تاں~$A$ وی ایویں ای اے۔ ایہ وکھا کے کرو کہ تسیں~$g$ نوں~$A$ دی اک
گنتی وار فہرست توں~$B$ دی گنتی وار فہرست بناؤن لئی کِویں ورت سکدے او۔
\end{prob}

\begin{proof}[{\olref[nen-alt]{thm:nonenum-pownat}} دا تقلیلی ثبوت]
ایہ تقلیلی ثبوت \pnbLatin{(proof by reduction)} اے۔ فرض کرو کہ $\Pow{\Nat}$،
!!{enumerable} ہوندا، تے ایس لئی اوہدی اک گنتی وار فہرست $N_{1}$،
$N_{2}$، $N_{3}$، \dots ہوندی۔

فنکشن $f \colon \Pow{\Nat} \to \Bin^\omega$ ایویں تعریف کرو کہ
$f(N)$ اوہ خاصیتی ثنائی لڑی \pnbLatin{(characteristic binary sequence)} $s$
ہووے جیہدے لئی $s(n) = 1$ عین اوس ویلے جدوں $n \in N$، تے ہور
صورت وچ $s(n) = 0$ ہووے۔
% OLSIZ-010 ماخذ دی درستی: نتیجے والی لڑی دا ناں s اے؛ انگریزی عبارت وچ غیر متعین k والے s_k تے s_k(n) ورتے گئے سن۔

ایہ صاف طور تے اک فنکشن تعریف کردا اے، کیوں جے جدوں وی
$N \subseteq \Nat$ ہووے، ہر $n \in \Nat$ یا تاں~$N$ دا عنصر
(!!a{element}) اے یا نہیں۔ مثال لئی، جفت قدرتی عدداں \pnbLatin{(even natural numbers)} دا سیٹ
$2\Nat = \Setabs{2n}{n \in \Nat} = \{0,2, 4, 6,
\dots\}$ لڑی $1010101\dots$ ول، خالی سیٹ \pnbLatin{(empty set)} $\emptyset$
لڑی $0000\dots$ ول، تے خود سیٹ~$\Nat$ لڑی $1111\dots$ ول بھیجیا
جاندا اے۔

ایہ ہر ہدف قدر لین والا (!!{surjective}) وی اے: $0$ تے $1$ دی ہر
لڑی قدرتی عدداں \pnbLatin{(natural numbers)} دے کسے سیٹ نال میل کھاندی اے، یعنی
اوس سیٹ نال جیہدے رکن (members) اوہ قدرتی عدد نیں جیہڑے لڑی وچ~$1$
والے مقاماں (positions) دے مطابق نیں۔ ہور ٹھیک طور تے، جے
$s \in \Bin^\omega$، تاں $N \subseteq \Nat$ نوں ایویں تعریف کرو:
\[
N = \Setabs{n \in \Nat}{s(n) = 1}
\]
فیر $f(N) = s$، جیویں~$f$ دی تعریف نوں ویکھ کے جانچیا جا سکدا اے۔

ہُن فہرست اُتے غور کرو:
\[
f(N_1), f(N_2), f(N_3), \dots
\]
کیوں جے~$f$ ہر ہدف قدر لین والا (!!{surjective}) اے، ایس لئی
$\Bin^\omega$ دا ہر رکن~$f$ دی کسے دلیل (argument) لئی قدر دے طور
تے لازماً آوے گا، تے سو فہرست وچ لازماً آوے گا۔ ایس لئی ایہ فہرست
$\Bin^\omega$ دے ہر رکن دی گنتی وار فہرست اے۔

سو جے $\Pow{\Nat}$، !!{enumerable} ہوندا، تاں $\Bin^\omega$ وی
!!{enumerable} ہوندا۔ پر $\Bin^\omega$ گنتی وار فہرست نہ بنن جوگ
(!!{nonenumerable}) اے (\olref[nen-alt]{thm:nonenum-bin-omega})۔ ایس
لئی $\Pow{\Nat}$، !!{nonenumerable} اے۔
\end{proof}

%\begin{explain}
%تقلیل دی سمت (direction of reduction)، یعنی تقلیل کس پاسے جاندی اے،
%ایس بارے بھلیکھا پینا سوکھا اے۔ مثال لئی، کوئی ہر ہدف قدر لین والا
%(!!a{surjective}) فنکشن $g \colon \Bin^\omega \to X$ ایہ
%\emph{نہیں} ثابت کردا کہ~$X$، !!{nonenumerable} اے۔ (فنکشن $g
%\colon \Bin^\omega \to \Bin$ اُتے غور کرو جیہڑا $g(s) = s(1)$
%نال تعریف اے، یعنی اوہ فنکشن جیہڑا $0$ تے $1$ دی اک لڑی نوں اوہدے
%پہلے عنصر (!!{element}) ول بھیجدا اے۔ ایہ ہر ہدف قدر لیندا اے، کیوں
%جے کجھ لڑیاں $0$ توں تے کجھ $1$ توں شروع ہوندیاں نیں۔ پر~$\Bin$
%متناہی (finite) اے۔) ایہ وی یاد رکھو کہ فنکشن~$f$ دا ہر ہدف قدر
%لینا لازمی اے، نہیں تاں دلیل نہیں چلدی: $f(x_1)$، $f(x_2)$،
%\dots{} وچ~$Y$ دے سارے !!{element}s دے شامل ہون دی ضمانت نہیں ہو
%سکدی۔ مثال لئی،
%$h\colon \Nat \to
%\Bin^\omega$ جیہڑا ایویں تعریف اے:
%\[
%h(n) = \underbrace{000\dots0}_{\text{$n$ $0$'s}}111\dots
%\]
%ایتھے اکاں دی لامتناہی ثنائی پونچھ (infinite binary tail) ہر نتیجے
%نوں ہدف سیٹ دی لامتناہی لڑی بناؤندی اے۔ ایہ اک فنکشن اے، پر~$\Nat$،
%!!{enumerable} اے۔
% PNBSIZ-005 ماخذ دی درستی: ہر نتیجہ n صفراں توں بعد اکاں دی لامتناہی پونچھ رکھدا اے؛ غیر فعال انگریزی وضاحت صرف متناہی صفر لکھ کے اوہنوں لامتناہی لڑی والے ہدف وچ رکھدی سی۔
%\end{explain}

\begin{prob}\label{sfr:siz:red-alt:prob:nat-nat}
  تقلیلی دلیل \pnbLatin{(reduction argument)} نال وکھاؤ کہ سیٹ~$X$، جیہڑا سارے
  فنکشناں (functions) $f\colon \Nat \to \Nat$ دا سیٹ اے،
  !!{nonenumerable} اے (اشارا: $X$ توں~$\Bin^\omega$ ول اک ہر ہدف
  قدر لین والا فنکشن \pnbLatin{(surjective function)} دیو۔)
\end{prob}

\begin{prob}
تقلیلی دلیل \pnbLatin{(reduction argument)} نال وکھاؤ کہ قدرتی عدداں دیاں
جوڑیاں \pnbLatin{(pairs of natural numbers)} دے سارے \emph{سیٹاں دے} سیٹ، یعنی
$\Pow{\Nat \times \Nat}$، !!{nonenumerable} اے۔
\end{prob}

\begin{prob}
وکھاؤ کہ~$\Nat^\omega$، یعنی قدرتی عدداں \pnbLatin{(natural numbers)} دیاں
لامتناہی لڑیاں \pnbLatin{(infinite sequences)} دا سیٹ، تقلیلی دلیل \pnbLatin{(reduction argument)} نال !!{nonenumerable} اے۔
\end{prob}

%\begin{prob}
%$P$، $\Nat$ توں سیٹ~$\{0\}$ ول فنکشناں (functions) دا سیٹ ہووے،
%تے~$Q$ مثبت صحیح عدداں (positive integers) دے سیٹ توں سیٹ~$\{0\}$
%ول \emph{جزوی} فنکشناں (partial functions) دا سیٹ ہووے۔ وکھاؤ کہ
%~$P$، !!{enumerable} اے تے~$Q$ نہیں۔ (اشارا: $\Bin^\omega$ دی گنتی
%وار فہرست بناؤن دے مسئلے نوں~$Q$ دی گنتی وار فہرست بناؤن ول گھٹاؤ۔)
%\end{prob}

\begin{prob}
$S$، $\Nat$ توں سیٹ~$\{0,1\}$ ول سارے ہر ہدف قدر لین والے
(!!{surjection}s) فنکشناں دا سیٹ ہووے، یعنی~$S$ سارے
!!{surjection}s~$f \colon \Nat \to \Bin$ اُتے مشتمل اے۔ وکھاؤ کہ
~$S$، !!{nonenumerable} اے۔
\end{prob}

\begin{prob}
وکھاؤ کہ سارے حقیقی عدداں \pnbLatin{(real numbers)} دا سیٹ~$\Real$، گنتی وار
فہرست نہ بنن جوگ (!!{nonenumerable}) اے۔
\end{prob}
\endgroup
% /OLP-0040

\OLEndChapterHook



\par\noindent\emph{متبادل ترتیب۔} سیٹاں دے سائز دا مُڈھلا باب (\texttt{size-of-sets.tex})۔ ایہ ترتیب تدوین جوگ ماخذ فائل وچ موجود اے۔
\endgroup
% /OLP-0027

% OLP-0041 content/sets-functions-relations/arithmetization/arithmetization.tex
\begingroup
\def\olfilename{arithmetization}
\expandafter\def\csname ol@sectioncs\endcsname{section}

\olchapter{sfr}{arith}{عددی نظاماں دی تشکیل \pnbLatin{(Arithmetization)}}

\begin{editorial}
ایس باب وچ سادہ سیٹ تھیوری \pnbLatin{(naive set theory)} دے اندر عددی نظاماں دی
تشکیل \pnbLatin{(construction of the number systems)} پیش کیتی گئی اے۔
ایہ مواد ٹم بٹن دی لکھت Open Set Theory توں لیا گیا اے۔
\end{editorial}

% OLP-0042 content/sets-functions-relations/arithmetization/integers.tex
\begingroup
\def\olfilename{integers}
\expandafter\def\csname ol@sectioncs\endcsname{section}

\olfileid{sfr}{arith}{int}
\olsection{\texorpdfstring{$\Nat$ توں $\Int$ ول: قدرتی عدداں توں صحیح اعداد}{ℕ توں ℤ ول: قدرتی عدداں توں صحیح اعداد}}

ایتھے دو بنیادی مشاہدے نیں:
	\begin{enumerate}
		\item ہر صحیح عدد (integer) نوں $n - m$ دی صورت وچ لکھیا جا سکدا اے، جتھے $n, m \in\Nat$۔
		\item اظہار $n - m$ وچ کوڈ کیتی معلومات نوں ترتیب وار جوڑے \pnbLatin{(ordered pair)} $\tuple{n, m}$ نال وی اوہنی ای طرح کوڈ کیتا جا سکدا اے۔
	\end{enumerate}
اسی پہلے ای جان دے آں کہ قدرتی عدداں دیاں ترتیب وار جوڑیاں $\Nat^2$
دے !!{element}s (elements) نیں، تے اسی فرض کردے آں کہ $\Nat$ نوں
سمجھدے آں۔ ایس لئی ایہناں دوہاں مشاہدیاں توں اک ابتدائی، سادی تجویز
ملدی اے: \emph{صحیح عدداں نوں قدرتی عدداں دیاں ترتیب وار جوڑیاں سمجھ لئیے}۔

پر ایہ تجویز حقیقت وچ حدوں ودھ سادی اے۔ صحیح عدداں دے حساب وچ اسی
$0- 2 = 4 - 6$ چاہنے آں، حالانکہ ترتیب وار جوڑیاں دے طور اُتے
$\tuple{0, 2 }\neq \tuple{4, 6}$ اے۔ ایس لئی $\Nat^2$ نوں سیدھا
صحیح عدداں دا سیٹ نہیں آکھیا جا سکدا۔

پچھلی مشکل نوں عام صورت وچ لکھئیے تاں ساڈی لوڑ ایہ اے:
	$$a - b = c - d \text{ اُدوں تے صرف اُدوں جدوں }a + d = c + b$$
(صحیح عدداں نے \emph{ایویں ای ورتاؤ کرنا} اے: مساوات دے دونیں پاسیاں
وچ $b$ تے $d$ جمع کر دیو۔) ایس ورتاؤ نوں یقینی بناؤن لئی ترتیب وار
جوڑیاں اُتے ہم ارزیت دا تعلق \pnbLatin{(equivalence relation)} $\Intequiv$ ایویں
تعریف کردے آں:
$$\tuple{ a, b } \Intequiv \tuple{c, d}\text{ اُدوں تے صرف اُدوں جدوں }a + d = c + b$$
ہُن ثابت کرنا اے کہ ایہ واقعی ہم ارزیت دا تعلق اے۔
\begin{prop} $\Intequiv$ ہم ارزیت دا تعلق \pnbLatin{(equivalence relation)} اے۔
\end{prop}
	\begin{proof}
	ساڈے نوں وکھاؤنا اے کہ $\Intequiv$ انعکاسی (reflexive)، متناظر
	(symmetric) تے متعدی (transitive) اے۔

	\emph{انعکاسیت (Reflexivity):} $\tuple{a, b} \Intequiv \tuple{a, b}$،
کیوں جے $a + b = b + a$۔

	\emph{تناظر (Symmetry):} فرض کرو $\tuple{a, b} \Intequiv \tuple{c, d}$،
یعنی $a + d = c + b$۔ فیر $c + b = a + d$، ایس لئی
$\tuple{c, d} \Intequiv \tuple{a, b}$۔

	\emph{تعدیت (Transitivity):} فرض کرو $\tuple{a, b} \Intequiv \tuple{c, d}\Intequiv
\tuple{m, n}$۔ فیر $a + d = c + b$ تے $c + n = m + d$۔ لہٰذا
$a + d + c + n = c + b + m + d$، جِتھوں $a + n = m + b$۔
سو $\tuple{a, b } \Intequiv \tuple{m, n}$۔
\end{proof}

ہُن ایس ہم ارزیت دے تعلق نال ہم ارزیت دے طبقے \pnbLatin{(equivalence classes)}
بنا سکدے آں:

\begin{defn}
		صحیح اعداد (integers)، قدرتی عدداں دیاں ترتیب وار جوڑیاں دے
		$\Intequiv$ تحت ہم ارزیت دے طبقے نیں؛ یعنی
		$\Int = \equivclass{\Nat^2}{\Intequiv}$۔
\end{defn}

ایس مقرر کیتی تعریف \pnbLatin{(stipulative definition)} اُتے کئی طرح دے
\emph{فلسفیانہ} ردعمل ہو سکدے نیں۔ اوہناں ول مڑن توں پہلاں کجھ فنی
باریکیاں مکمل کر لینا فائدے مند اے۔

صحیح اعداد دی تعریف توں بعد اوہناں اُتے بنیادی فنکشن (functions) تے
تعلق (relations) وی تعریف کرنے نیں۔ $\equivrep{m, n}{\Intequiv}$ نال
اوہ $\Intequiv$-ہم ارزیت طبقہ مراد لاواں گے جیہدے وچ $\tuple{m, n}$
!!a{element} (element) اے۔\footnote{نوٹ: \olref[sfr][rel][eqv]{def:equivalenceclass}
وچ متعارف علامت مطابق ایسے طبقے نوں
$\equivrep{\tuple{m,n}}{\Intequiv}$ لکھیا جانا سی، پر اوہ ذرا اوکھا
پڑھدا اے۔} یعنی:
	$$\equivrep{m, n}{\Intequiv} = \Setabs{\tuple{a, b} \in \Nat^2}{\tuple{a, b}\Intequiv \tuple{m, n}}$$
ہُن طبقیاں اُتے جمع، ضرب تے ترتیب ایویں تعریف کردے آں:
	\begin{align*}
		\equivrep{a, b}{\Intequiv} + \equivrep{c, d}{\Intequiv} &= \equivrep{a + c, b + d}{\Intequiv}\\
		\equivrep{a, b}{\Intequiv} \times \equivrep{c, d}{\Intequiv} &= \equivrep{a c + b  d, a  d + b c}{\Intequiv}\\
		\equivrep{a, b}{\Intequiv} \leq \equivrep{c, d}{\Intequiv} &\text{ اُدوں تے صرف اُدوں جدوں }a + d \leq b + c
	\end{align*}
(عام رواج مطابق، ”$ab$“ نوں ”$(a \times b)$“ دا مختصر روپ ورتیا اے،
تاں جو مسلّمے سوکھے پڑھے جان۔)
% PNBARI-001 ماخذ دی نحوی درستی: انگریزی عبارت وچ “using ab to stand for” والا to چھڈیا ہویا اے؛ پنجابی صرف فوراً واضح کیتی مراد بیان کردی اے۔
ایہ وی جانچنا اے کہ نمائندہ جوڑا بدلنے نال جواب نہ بدلے تے ایہ تعریفاں
اوہو جیہا ورتاؤ کرن جیہو جیہا اوہناں نوں \emph{کرنا چاہیدا اے}۔ پوری
جانچ لمّی اے، ایس لئی تفصیل \olref[check]{sec} وچ اے۔ مختصر نتیجہ:
ساریاں تعریفاں درست کم کردیاں نیں۔

اک آخری مسئلہ رہندا اے۔ اسی قدرتی عدداں نال صحیح اعداد تعمیر کیتے نیں،
ایس لئی تعمیر کیتے طبقیاں دے اندر قدرتی عدد \emph{خود صحیح عدد نہیں}۔
ایس دی فلسفیانہ اہمیت ول \olref[ref]{sec} وچ مڑاں گے۔ فنی طور اُتے
قدرتی عدداں نوں صحیح اعداد \emph{دے روپ وچ} ورتن دا طریقہ چاہیدا اے
\pnbLatin{(treat natural numbers as integers)}۔ ہر $n \in \Nat$ لئی مقرر کرو:
$n_\Int = \equivrep{n, 0}{\Intequiv}$۔ ہُن تصدیق کرنی اے کہ ایہ
نمائندگی جمع، ضرب تے ترتیب محفوظ رکھدی اے؛ یعنی ہر $m, n \in \Nat$ لئی
\begin{align*}
	(m + n)_\Int &= m_\Int + n_\Int\\
	(m \times n)_\Int &= m_\Int \times n_\Int\\
	m \leq n &\liff m_\Int \leq n_\Int
\end{align*}
ایہ تصدیق سیدھی اے۔ مثال لئی دوجی شرط، یعنی ضرب محفوظ رہن دی گل،
قدرتی عدداں دے ورتاؤ توں ایویں نکلدی اے:
%\begin{proof}
%	قدرتی عدداں دے ورتاؤ توں صاف اے:
	\begin{align*}
%		(m + n)_\Int  = \equivrep{m + n, 0}{\Intequiv} = \equivrep{m + n, 0 + 0}{\Intequiv} = \equivrep{m, 0}{\Intequiv} + \equivrep{n, 0}{\Intequiv} = m_\Int + n_\Int\\
		(m \times n)_\Int  &= \equivrep{m \times n, 0}{\Intequiv} \\
		&= \equivrep{m \times n + 0 \times 0, m \times 0 + 0 \times n}{\Intequiv} \\
		&= \equivrep{m, 0}{\Intequiv} \times \equivrep{n, 0}{\Intequiv} \\
		&= m_\Int \times n_\Int
\end{align*}
%\end{proof}
باقی دو شرطاں دی تصدیق مشق اے۔
\begin{prob}
	وکھاؤ کہ $(m + n)_\Int = m_\Int + n_\Int$ تے
	$m \leq n \liff m_\Int \leq n_\Int$، ہر $m, n \in \Nat$ لئی۔
\end{prob}
\endgroup
% /OLP-0042
% OLP-0043 content/sets-functions-relations/arithmetization/rationals.tex
\begingroup
\def\olfilename{rationals}
\expandafter\def\csname ol@sectioncs\endcsname{section}

\olfileid{sfr}{arith}{rat}
\olsection{\texorpdfstring{$\Int$ توں $\Rat$ تک: ناطق اعداد \pnbLatin{(Rational Numbers)}}{ℤ توں ℚ تک: ناطق اعداد (Rational Numbers)}}

اسی ہُنّے ویکھیا اے کہ سادی سیٹ تھیوری \pnbLatin{(naive set theory)} ورت کے قدرتی
عدداں توں صحیح اعداد کِویں بنائے جا سکدے نیں۔ ہُن اسی بالکل رلدے ملدے
طریقے نال صحیح عدداں توں ناطق اعداد \pnbLatin{(rational numbers)} بناؤنا ویکھاں گے۔
ساڈے مُڈھلے مشاہدے ایہ نیں:
\begin{enumerate}
	\item ہر ناطق عدد نوں کسر (fraction) $\nicefrac{i}{j}$ دی صورت وچ
	لکھیا جا سکدا اے، جتھے $i$ تے $j$ دونیں صحیح عدد نیں، پر مخرج
	(denominator) $j$ صفر نہیں۔
	\item اظہار $\nicefrac{i}{j}$ وچ کوڈ کیتی معلومات نوں ترتیب وار جوڑے
	\pnbLatin{(ordered pair)} $\tuple{ i, j}$ وچ وی اوہو جیہے طور اُتے کوڈ کیتا جا
	سکدا اے۔
\end{enumerate}
ظاہری طریقہ ایہ ہووے گا کہ ناطق عدداں نوں $\Int \times (\Int \setminus
\{0_\Int\})$ توں چُنے ترتیب وار جوڑے \emph{ای سمجھیا} جاوے۔ پر پہلے
وانگ ایہ گل حدوں ودھ سادی اے، کیوں جو اسی چاہنے آں کہ
$\nicefrac{3}{2} = \nicefrac{6}{4}$، پر $\tuple{ 3, 2}\neq \tuple{ 6,
4}$۔ ہور عام طور تے ساڈی لوڑ ایہ اے:
\[
	\nicefrac{a}{b} = \nicefrac{c}{d} \text{ اُدوں تے صرف اُدوں جدوں } a \times d = b \times c
\]
ایہ مطلوبہ برابری آڑی ضرب دے معیار نال ملدی اے۔ ایس لئی اسی $\Int
\times (\Int \setminus \{0_\Int\})$ اُتے اک {ہم ارزیت دا تعلق
\pnbLatin{(equivalence relation)}} ایویں تعریف کردے آں:
\[
	\tuple{ a, b }\Ratequiv \tuple{ c, d} \text{ اُدوں تے صرف اُدوں جدوں }a \times d = b \times c
\]
ساڈے نوں جانچنا پینا اے کہ ایہ واقعی ہم ارزیت دا تعلق اے۔ ایہ جانچ
$\Intequiv$ والے معاملے نال بہت رلدی اے، ایس لئی اسی ایہنوں مشق لئی
چھڈدے آں۔
\begin{prob}
وکھاؤ کہ $\Ratequiv$ ہم ارزیت دا تعلق اے۔
\end{prob}
ایہ تعلق ساڈے نوں برابر کسراں دے سارے نمائندیاں نوں اک طبقے وچ رکھن
دی اجازت دیندا اے:
\begin{defn}
ناطق اعداد، صحیح عدداں دے اوہناں جوڑیاں دے $\Ratequiv$ تحت ہم ارزیت
دے طبقے \pnbLatin{(equivalence classes)} نیں جیہناں دا دوجا رکن صفر نہیں۔ یعنی $\Rat =
\equivclass{(\Int \times (\Int\setminus \{0_\Int\}))}{\Ratequiv}$۔
\end{defn}

صحیح عدداں وانگ، اسی ناطق عدداں اُتے کجھ بنیادی عمل وی تعریف کرنا
چاہنے آں۔ جتھے $\equivrep{i,j}{\Ratequiv}$، $\Ratequiv$ تحت اوہ ہم
ارزیت دا طبقہ اے جیہدے وچ $\tuple{i, j}$ اک !!a{element} اے، اسی
جمع تے ضرب ایویں مقرر کردے آں:
\begin{align*}
	\equivrep{a, b}{\Ratequiv} + \equivrep{c, d}{\Ratequiv} &= \equivrep{ad + bc,  bd}{\Ratequiv}\\
	\equivrep{a, b}{\Ratequiv} \times \equivrep{c, d}{\Ratequiv} &= \equivrep{a  c, b d}{\Ratequiv}.
\intertext{ایہناں ناطق عدداں اُتے $r \leq s$ ترتیب تعریف کرن لئی اسی
ایہ گل ورتدے آں کہ $r \le s$ اُدوں تے صرف اُدوں جدوں $s - r$ منفی
نہیں؛ یعنی $s - r$ نوں $\nicefrac{i}{j}$ دی صورت وچ لکھیا جا سکدا
اے، جتھے $i$ غیر منفی (non-negative) تے $j$ مثبت اے:}
	\equivrep{a, b}{\Ratequiv} \leq \equivrep{c, d}{\Ratequiv} &\text{ اُدوں تے صرف اُدوں جدوں }
	\equivrep{c, d}{\Ratequiv} - \equivrep{a, b}{\Ratequiv} =
	\equivrep{i_\Int, j_\Int}{\Ratequiv}
\end{align*}
% PNBARI-002 ماخذ دی درستی: انگریزی وضاحت اک واری فرق نوں r-s آکھدی اے؛ اگلی تعریف تے پچھلا جملہ دونیں s-r مقرر کردے نیں۔
کِسے $i \in \Nat$ تے $0 \neq j \in \Nat$ لئی۔

فیر ساڈے نوں جانچنا پینا اے کہ ایہ تعریفاں اوہو جیہا ورتاؤ کردیاں نیں
جیہو جیہا اوہناں نوں \emph{کرنا چاہیدا اے}؛ نمائندہ بدلّن نال جواب
نہیں بدلدا تے عمل تے ترتیب درست رہندے نیں۔ ایہ جانچ اسی
\olref[check]{sec} ول بھیج دے آں، تے اوتھے ایہ واقعی درست نکلدی اے!{}
آخر وچ اسی صحیح عدداں نوں ناطق اعداد دے روپ وچ \emph{ورتن} دا طریقہ
چاہنے آں۔ ایس لئی ہر $i \in \Int$ واسطے اسی مقرر کردے آں کہ $i_\Rat =
\equivrep{i, 1_\Int}{\Ratequiv}$۔ فیر اسی \olref[check]{sec} وچ جانچدے
آں کہ ایہ نمائندگی جمع، ضرب تے ترتیب نوں قائم رکھدی اے۔

\begin{prob}
ایہ ثابت کرن لئی کہ صحیح عدداں دی ایہ نمائندگی جمع، ضرب تے ترتیب نوں
قائم رکھدی اے، وکھاؤ کہ $(i + j)_\Rat = i_\Rat+ j_\Rat$ تے $(i \times
j)_\Rat = i_\Rat \times j_\Rat$ تے $i \leq j \liff i_\Rat \leq
j_\Rat$، ہر $i, j \in \Int$ لئی۔
\end{prob}
\endgroup
% /OLP-0043
% OLP-0044 content/sets-functions-relations/arithmetization/reals.tex
\begingroup
\def\olfilename{reals}
\expandafter\def\csname ol@sectioncs\endcsname{section}

\olfileid{sfr}{arith}{real}
\olsection{حقیقی عددی خط \pnbLatin{(Real Line)}}

ہُن اگلا قدم ایہ وکھاؤنا اے کہ ناطق عدداں توں حقیقی عدد کِویں بنائے
جاندے نیں۔ اُس توں پہلاں سانوں ایہ سمجھنا پینا اے کہ حقیقی عدداں نوں
\emph{وکھرا} کیہڑی گل بناؤندی اے۔

حقیقی عدد بہت سارے پاسیاں توں ناطق عدداں وانگ ورتاؤ کردے نیں۔ (فنی
لفظاں وچ، دونیں \emph{ترتیب دار میدان \pnbLatin{(ordered fields)}} نیں؛ ایہدی
تعریف لئی \olref[check]{orderedfield} ویکھو۔) ہُن جے تسیں
\olref[sfr][siz][]{chap} دیاں مشقاں کیتیاں نیں، تاں تہانوں پتا ہووے گا
کہ حقیقی عدداں دی تعداد ناطق عدداں نالوں سختی نال ودھ اے، یعنی
$\cardless{\Rat}{\Real}$۔ ایہ نتیجہ سب توں پہلاں کانتور نے ثابت کیتا
سی۔ پر لگ بھگ ڈھائی ہزار ورھیاں توں ایہ معلوم اے کہ غیر ناطق عدد
\pnbLatin{(irrational numbers)} موجود نیں، یعنی اوہ حقیقی عدد جیہڑے ناطق نہیں۔
بلکہ:
%\footnote{ہُن: ”$\sqrt{2}$“ صرف غیر منفی اصل جذر (nonnegative principal square root) نوں دَسدا اے؛ دوجا جذر اوہدا منفی اے۔ اگلے کجھ ورقیاں وچ اسی صرف ایس اصل جذر اُتے غور کراں گے۔}
% PNBARI-004 ماخذ دی درستی، غیر فعال نوٹ: ریڈیکل دی علامت صرف غیر منفی اصل جذر دَسدی اے، مثبت تے منفی دونیں جذر نہیں۔
\begin{thm}\ollabel{root2irrational}
$\sqrt{2}$ ناطق عدد نہیں، یعنی $\sqrt{2} \notin \Rat$
\end{thm}

\begin{proof}
تضاد دے طریقے نال فرض کرو کہ $\sqrt{2}$ ناطق اے۔ سو
$\sqrt{2} = \nicefrac{m}{n}$، کجھ قدرتی عدد $m$ تے $n$ لئی۔ بلکہ اسی
$m$ تے $n$ ایویں چُن سکدے آں کہ ایہ کسر ہور گھٹ نہ ہو سکے۔ ترتیب بدلیے تاں
$m^{2} = 2n^{2}$ ملدا اے۔ ایتھوں ثبوت نوں دو طریقیاں نال پورا کیتا جا
سکدا اے:

\emph{پہلا، ہندسی طور تے} (ٹیننبام دے طریقے اُتے چل کے)۔\footnote{ایہ
ثبوت \cite{Conway2006} وچ نقل ہویا اے۔} ایہناں مربعواں اُتے غور کرو:
\begin{center}
	\begin{tikzpicture}
		\draw[thick] (0,0) rectangle (3,3);
		\draw[thick, fill=red!50] (0,0) rectangle (2.3,2.3);
		\draw[thick, fill=yellow!50] (0.7,0.7) rectangle (3, 3);
		\draw[thick, fill=orange!50] (0.7,0.7) rectangle (2.3, 2.3);
		\draw[<->] (4, 0.7)--(4, 3);
		\node at (4.25, 1.85) (n) {$n$};
		\draw[<->] (5, 0)--(5, 3);
		\node at (5.25, 1.5) (m) {$m$};
	\end{tikzpicture}
\end{center}
چونکہ $m^2 = 2n^2$ اے، ایس لئی ضلع $n$ والے دونیں مربع جتھے اک دوجے
اُتے چڑھدے نیں، اوس حصے دا رقبہ اوہناں حصیاں دے کُل رقبے دے برابر اے
جیہڑے دوناں مربعواں وچوں کوئی وی نہیں ڈھکدا؛ یعنی نارنجی مربع دا رقبہ
دونیں بے رنگ مربعواں دے رقبیاں دے جوڑ دے برابر اے۔ ہُن جے نارنجی مربع
دا ضلع $p$ تے ہر بے رنگ مربع دا ضلع $q$ ہووے، تاں $p^2 = 2q^2$۔ پر
ہُن $\sqrt{2} = \nicefrac{p}{q}$، جتھے $p < m$ تے $q < n$ تے $p, q
\in \Nat$۔ ایہ اُس گل دے خلاف اے کہ $m$ تے $n$ نوں ممکن حد تک چھوٹا
چُنیا گیا سی۔

\emph{دوجا، رسمی طور تے۔} چونکہ $m^{2} = 2n^{2}$ اے، ایس توں نکلدا اے
کہ $m$ جفت (even) اے۔ (ایہ وکھاؤنا سوکھا اے کہ جے $x$ طاق (odd)
ہووے، تاں $x^2$ وی طاق ہوندا اے۔) سو $m = 2r$، کسے $r \in \Nat$
لئی۔ ترتیب بدلیے تاں $2r^2 = n^2$،
		%				\begin{align*}
		%					(2r)^{2} &= 2n^{2}\\
		%					4r^{2} &= 2n^{2}\\
		%					2r^{2} &= n^{2}
		%				\end{align*}
		% PNBARI-005 ماخذ دی درستی، غیر فعال اخذ: m=2r توں بعد انگریزی تبصرے وچ غیر متعین q آ گیا سی؛ تیناں تھاواں اُتے پابند r ورتیا اے۔
سو $n$ وی جفت اے۔ ایس طرح $m$ تے $n$ دونیں جفت نیں، تے ایس کرکے
کسر $\nicefrac{m}{n}$ نوں \emph{ہور} گھٹ کیتا جا سکدا اے۔ ایہ تضاد اے!
\end{proof}

گل چل پئی اے تاں ایہ تصویری ثبوت سانوں
\olref[his][set][mythology]{sec} دے مواد ول مڑ کے ویکھن دا موقع دیندا
اے۔ ٹیننبام (1927--2006) سراسر جدید ریاضی دان سی؛ پر ایہ ثبوت بے شک
سوہنا، پوری طرح کڑا، تے ہندسی وجدان اُتے ٹیک رکھن والا اے!

بہرحال، حقیقی عدد ناطق عدداں نالوں ”ہور پھیلے ہوئے“ نیں۔ اک معنٰی وچ
ناطق عدداں دے اندر ”خلا“ نیں، تے حقیقی عدد اوہ خلا بھر دے نیں۔
وائرسٹراس نے سمجھیا کہ ایہ گل حقیقی عدداں دی اکّو خاصیت بیان کردی اے
جیہڑی اوہناں نوں ناطق عدداں توں وکھرا کردی اے، یعنی تکمیلیت دی خاصیت
\pnbLatin{(Completeness Property)}: \emph{حقیقی عدداں دے ہر غیر خالی سیٹ، جیہدی
کوئی بالائی حد \pnbLatin{(upper bound)} ہووے، دی اک سب توں چھوٹی بالائی حد \pnbLatin{(least upper bound)} ہوندی اے۔}

ایہ ویکھنا سوکھا اے کہ ناطق عدداں کول تکمیلیت دی خاصیت نہیں اے۔ مثال
لئی $\sqrt{2}$ توں چھوٹے ناطق عدداں دے ایس سیٹ اُتے غور کرو:
\[
	\Setabs{p \in \Rat}{p^2 < 2 \text{ یا }p < 0}
\]
ایس سیٹ دی ناطق عدداں وچ اک بالائی حد اے؛ مثال لئی ایہدے !!{element}s
سارے $3$ نالوں چھوٹے نیں۔
% PNBARI-003 ماخذ دی درستی: element ٹوکن دے فوراً بعد اک بے محل < سی؛ کوئی کھبا operand نہیں، تے جملہ آپ دَسدا اے کہ سارے رکن 3 توں چھوٹے نیں۔
پر ایہدی سب توں چھوٹی بالائی حد کی اے؟ اسی ”$\sqrt{2}$“ آکھنا چاہندے
آں؛ پر اسی ہُنّے ویکھیا اے کہ $\sqrt{2}$ \emph{ناطق نہیں}۔ نالے
$\sqrt{2}$ توں وڈا کوئی \emph{سب توں چھوٹا} ناطق عدد نہیں اے۔ سو ایس
سیٹ دی بالائی حد تاں اے، پر کوئی سب توں چھوٹی بالائی حد نہیں۔ ایس لئی
ناطق عدداں کول تکمیلیت دی خاصیت نہیں اے۔

ایہدے اُلٹ، استمراری خط (continuum) کول ”اصولی طور تے“ تکمیلیت دی
خاصیت ہونی چاہیدی اے۔ اسی صرف ایہ نہیں چاہندے کہ $\sqrt{2}$ اک حقیقی
عدد ہووے؛ اسی ناطق عددی خط دے سارے ”خلا“ بھرنے چاہندے آں۔ بلکہ اسی ایہ
وی چاہندے آں کہ استمراری خط دے اپنے اندر کوئی ”خلا“ نہ ہووے۔ تکمیلیت
دے ذریعے سانوں عین ایہو ملے گا۔
\endgroup
% /OLP-0044
% OLP-0045 content/sets-functions-relations/arithmetization/cuts.tex
\begingroup
\def\olfilename{cuts}
\expandafter\def\csname ol@sectioncs\endcsname{section}

\olfileid{sfr}{arith}{cuts}
\olsection{\texorpdfstring{$\Rat$ توں $\Real$ ول \pnbLatin{(From Rationals to Reals)}}{ℚ توں ℝ ول (From Rationals to Reals)}}

اصل وچ تکمیلیت دی خاصیت \pnbLatin{(Completeness Property)} ایہ وکھاؤندی اے کہ حقیقی
خط \pnbLatin{(real line)} دا ہر نقطہ
$\alpha$ اوس خط نوں بالکل دو حصیاں وچ ونڈدا اے: اک اوہ جیہدی سب توں
چھوٹی بالائی حد \pnbLatin{(least upper bound)} $\alpha$ اے، تے دوجا اوہ جیہدی سب توں
وڈی زیریں حد \pnbLatin{(greatest lower bound)}
$\alpha$ اے۔ ناطق عدداں توں حقیقی اعداد \emph{بنان} لئی، ڈیڈیکائنڈ نے
سجھایا کہ حقیقی عدداں نوں صرف اوہ \emph{کٹاں (cuts)} سمجھ لیا جاوے جیہڑے ناطق
عدداں نوں ونڈدے نیں۔ یعنی اسی $\sqrt{2}$ نوں اوس \emph{کٹ} نال
شناخت کردے آں جیہڑا $< \sqrt{2}$ ناطق عدداں نوں $> \sqrt{2}$ ناطق
عدداں توں وکھ کردا اے۔

آؤ ایس گل نوں صاف کر لئیے۔ جے اسی ناطق عدداں نوں دو حصیاں وچ ونڈئیے،
تاں صرف اوہدے \emph{نچلے حصے \pnbLatin{(lower half)}} اُتے غور کر کے اسی اپنی کیتی ونڈ دی
اکّی شناخت کر سکدے آں۔ سو پوری درستی نال اسی ایہ تعریف دیندے آں:

\begin{defn}[کٹ (Cut)]
اک \emph{کٹ (cut)} $\alpha$ ناطق عدداں دا کوئی وی غیر خالی (non-empty)،
حقیقی یعنی پورے سیٹ توں چھوٹا (proper) مُڈھلا ٹکڑا \pnbLatin{(initial segment)}
اے جیہدا کوئی سب توں وڈا رکن \pnbLatin{(greatest element)} نہیں۔ یعنی $\alpha$ اک کٹ اُدوں تے
صرف اُدوں اے جدوں:
\begin{enumerate}
	\item \emph{غیر خالی، حقیقی \pnbLatin{(non-empty, proper)}}: $\emptyset \neq \alpha \subsetneq \Rat$
	\item \emph{مُڈھلا (initial)}: ہر $p,q \in \Rat$ لئی: جے $p < q \in \alpha$ تاں $p \in \alpha$
	\item \emph{کوئی سب توں وڈا رکن نہیں \pnbLatin{(no maximum)}}: ہر $p \in \alpha$ لئی کوئی $q \in \alpha$ اے جیہدے لئی $p < q$
\end{enumerate}
فیر $\Real$ کٹاں دا سیٹ اے۔
\end{defn}

سو ہُن اسی آکھ سکدے آں کہ $\sqrt{2} = \Setabs{p \in \Rat}{p^2 < 2\text{
یا }p < 0}$۔ بے شک ساڈے نوں جانچنا پینا اے کہ ایہ واقعی اک کٹ
\emph{(cut) اے}، پر اسی ایہ کم \olref[check]{sec} ول بھیج دے آں۔

پہلے وانگ، کجھ ہستیاں نوں تعریف کرن توں بعد اگلا کم اوہناں اُتے بنیادی
فنکشن تے تعلق تعریف کرنا اے۔ اسی اک سوکھے تعلق توں شروع کردے آں:
\begin{align*}
	\alpha \leq \beta \text{ اُدوں تے صرف اُدوں جدوں }\alpha \subseteq \beta
\end{align*}
ترتیب دی ایہ تعریف سانوں مرکزی نتیجہ \emph{بیان} کرن دی سہولت دیندی اے:
کٹاں دے سیٹ کول تکمیلیت دی خاصیت \pnbLatin{(Completeness Property)} اے۔ پوری طرح
کھول کے بیان ایہ شکل رکھدا اے۔ جے $S$ کٹاں دا اک غیر خالی سیٹ اے جیہدی
کوئی بالائی حد \pnbLatin{(upper bound)} اے، تاں $S$ دی اک سب توں چھوٹی بالائی حد
\pnbLatin{(least upper bound)} اے۔ ہور تفصیل نال: اک کٹ
$\lambda$ اے جیہڑا $S$ دی بالائی حد \pnbLatin{(upper bound)} اے، یعنی $(\forall \alpha \in S)\alpha \subseteq
\lambda$، تے $\lambda$ ایسا سب توں چھوٹا کٹ اے، یعنی $(\forall \beta \in \Real)((\forall \alpha \in S)\alpha \subseteq \beta \lif \lambda \subseteq \beta)$۔ ہُن
نتیجے دا ثبوت ایہ اے:

\begin{thm}\ollabel{realcompleteness}
کٹاں دے سیٹ کول تکمیلیت دی خاصیت \pnbLatin{(Completeness Property)} اے۔
\end{thm}

\begin{proof}
$S$ کٹاں دا کوئی وی غیر خالی سیٹ ہووے جیہدی بالائی حد اے۔ مقرر کرو کہ
$\lambda = \bigcup S$؛ ایہ ایس خاندان دا اتحاد (union) اے۔
%\Setabs{p \in \Rat}{(\exists \alpha \in S)p \in \alpha}$$
پہلاں اسی دعویٰ کردے آں کہ $\lambda$ اک کٹ اے:
\begin{enumerate}
\item کیوں جو $S$ غیر خالی اے، ایس لئی گھٹ توں گھٹ اک کٹ $S$ وچ اے، سو
$\emptyset \neq \lambda$۔
% PNBARI-006 ماخذ دی درستی: ایتھے کٹ دی موجودگی بالائی حد توں نہیں بلکہ S دے غیر خالی ہون توں نکلدی اے۔
کیوں جو $S$ کٹاں دا سیٹ اے، $\lambda \subseteq \Rat$۔ کیوں جو $S$ دی
بالائی حد اے، کوئی $p \in \Rat$ ہر کٹ $\alpha \in S$ وچوں غائب اے۔
سو $p\notin \lambda$، تے ایس لئی $\lambda \subsetneq \Rat$۔
\item فرض کرو $p < q \in \lambda$۔ سو کوئی $\alpha \in S$ اے جیہدے
لئی $q \in \alpha$۔ کیوں جو $\alpha$ اک کٹ اے، $p \in \alpha$۔ سو
$p \in \lambda$۔
\item فرض کرو $p \in \lambda$۔ سو کوئی $\alpha \in S$ اے جیہدے لئی
$p \in \alpha$۔ کیوں جو $\alpha$ اک کٹ اے، کوئی $q \in \alpha$ اے
جیہدے لئی $p < q$۔ سو $q \in \lambda$۔
\end{enumerate}
ایہ دعویٰ ثابت کردا اے۔ نال ای صاف اے کہ $(\forall \alpha \in S)\alpha
\subseteq \bigcup S = \lambda$، یعنی $\lambda$، $S$ دی بالائی حد
\pnbLatin{(upper bound)} اے۔ سو ہُن فرض کرو کہ $\beta \in \mathbb{R}$ وی اک بالائی حد اے، یعنی
$(\forall \alpha \in S)\alpha \subseteq \beta$۔ ہر $p \in \Rat$ لئی،
جے $p \in \lambda$، تاں کوئی $\alpha \in S$ اے جیہدے لئی $p \in
\alpha$، سو $p \in \beta$۔ عام کر کے $\lambda \subseteq \beta$۔ سو
$\lambda$، $S$ دی \emph{سب توں چھوٹی} بالائی حد \pnbLatin{(least upper bound)} اے۔
\end{proof}

سو ساڈے کول ہُن ہستیاں دا اک جتھا اے جیہڑا تکمیلیت دی خاصیت
\pnbLatin{(Completeness Property)} پوری کردا اے۔ ایس گل نوں بیان کرن دا اک طریقہ
ایہ اے: ساڈے کٹاں وچ کوئی ”خلا“ (gaps)
نہیں۔ (سو ناطق عدداں دی بجائے حقیقی عدداں دے ہور ”کٹ“ لین نال کوئی
دلچسپ نویں شے نہیں ملے گی۔)

اگے ساڈے نوں حقیقی عدداں اُتے کجھ عمل (operations) تعریف کرنے پینے نیں۔
اسی ایہ مقرر کر کے ناطق عدداں نوں حقیقی عدداں وچ سماؤندے (embed) آں کہ $p_\Real =
\Setabs{q \in \Rat}{q < p}$، ہر $p \in \Rat$ لئی۔ فیر اسی تعریف کردے آں:
\begin{align*}
	\alpha + \beta &= \Setabs{p + q}{p \in \alpha \land q \in \beta}\\
%	\alpha - \beta &\defis \Setabs{p - q}{p \in \alpha \land q \in \Rat \setminus \beta}\\
	\alpha \times \beta &=
	\Setabs{p \times q}{0 \leq p \in \alpha \land 0 \leq q \in \beta} \cup 0_\Real & \text{جے }\alpha, \beta \geq 0_\Real
%	\alpha \div \beta &\defis
%	\Setabs{\nicefrac{p}{q}}{p \in \alpha \land q \in \Rat \setminus \beta}  & \text{جے }\alpha \geq 0_\Real, \beta > 0_\Real
\end{align*}
% PNBARI-007 ماخذ دی درستی: ناطق p دی جاگذاری p_Real اے؛ ایس لئی ضرب وچ حقیقی صفر وی 0_Real اے، اک واری آیا 0^R نہیں۔
ضرب دے دوجے معاملے سنبھالن لئی، پہلاں مقرر کرو: %کہ $0_\Real \times \alpha = 0_\Real = \alpha \times 0_\Real$، تے شامل کرو:
\begin{align*}
	-\alpha &= \Setabs{p - q}{p < 0 \land q \notin \alpha}
\end{align*}
تے فیر ایہ مقرر کرو:
\begin{align*}
	\alpha \times \beta &\defis
	\begin{cases}
		\mathord{-}\alpha \times \mathord{-}\beta &\text{جے }\alpha < 0_\Real\text{ تے }\beta < 0_\Real\\
		\mathord{-}(\mathord{-}\alpha \times \beta) &\text{جے }\alpha < 0_\Real \text{ تے }\beta > 0_\Real\\
		\mathord{-}(\alpha \times \mathord{-}\beta) &\text{جے }\alpha > 0_\Real \text{ تے }\beta < 0_\Real
	\end{cases}
\end{align*}
فیر ساڈے نوں جانچنا پینا اے کہ ایہناں وچوں ہر تعریف ہمیشہ اک کٹ دیندی
اے۔ آخر وچ ساڈے نوں اک سوکھا (پر لمبا چوڑا) مظاہرہ کرنا پینا اے کہ
ایویں تعریف کیتے کٹ عین اوہو جیہا ورتاؤ کردے نیں جیہو جیہا اوہناں نوں
کرنا چاہیدا اے۔ ایتھے کٹ توں مراد cut اے۔ پر اسی ایہ سارا کم
\olref[check]{sec} ول بھیج دے آں۔
\endgroup
% /OLP-0045
% OLP-0046 content/sets-functions-relations/arithmetization/reflections.tex
\begingroup
\def\olfilename{reflections}
\expandafter\def\csname ol@sectioncs\endcsname{section}

\olfileid{sfr}{arith}{ref}
\olsection{کجھ فلسفیانہ غور فکر \pnbLatin{(Some Philosophical Reflections)}}

فنی تفصیلاں (technicalities) لئی اینا کافی اے۔ پر اوہناں نے حاصل کی کیتا؟

بھلا، لگ بھگ بے اختلاف طور تے اوہناں نے سانوں کجھ {سوہنی} خالص ریاضی
\pnbLatin{(pure mathematics)} دِتی۔ نال ای کجھ ڈونگھیاں تصوری کامیابیاں
\pnbLatin{(conceptual achievements)} وی سن۔ ایہ ویکھنا اک گہری بصیرت سی کہ تکمیلیت
دی خاصیت \pnbLatin{(Completeness Property)} حقیقی تے ناطق عدداں وچکار فیصلہ کن فرق
بیان کردی اے۔ نال ای ڈیڈیکائنڈ کٹاں دے طور تے حقیقی عدداں دی صاف
بناوٹ، تجزیے دے موضوع نوں پکی بنیاد دیندی اے۔ اسی جان دے آں کہ
\emph{مکمل ترتیب دار میدان \pnbLatin{(complete ordered field)}} دا تصور مربوط اے، کیوں جو کٹ عین ایسا
میدان بناؤندے نیں۔

ایس سب دے باوجود، سانوں ایہناں کامیابیاں بارے کجھ تحفظات (reservations)
سامنے لیانے چاہیدے نیں۔

پہلاں، ایہ صاف نہیں کہ حقیقی عدداں نوں کٹاں دے حوالے نال سوچنا، اوہناں
دیاں واقف (ممکنہ طور تے لامتناہی) اعشاری توسیعاں دے حوالے نال سوچن توں
کوئی \emph{ودھ} کڑائی والا \pnbLatin{(more rigorous)} اے۔ حقیقی عدداں دی ایہ پچھلی
”بناوٹ“، کوشی لڑی \pnbLatin{(Cauchy sequence)}
راہی حقیقی عدداں دی بناوٹ نال کجھ رل دی اے؛ پر اصل وچ ایہ سترھویں صدی
دے مُڈھ توں اگے ریاضی داناں نوں بنیادی طور تے معلوم سی (ویکھو
\olref[cauchy]{sec})۔ کڑائی (rigour) وچ اصل اضافہ ایس ادراک توں آیا کہ
حقیقی عدداں کول تکمیلیت دی خاصیت اے؛ حقیقی عدداں نوں خاص سیٹاں دے طور تے
بنا سکنا شاید آپنے آپ وچ اینا دلچسپ نہیں۔

ایہ گل ہور وی گھٹ صاف اے کہ صحیح عدداں یا ناطق عدداں دی (بہت سوکھی)
اعدادی تشکیل (arithmetization)، اوہناں شعبیاں وچ کڑائی ودھاؤندی اے۔ ایتھے اک سادہ گل
کرنی فائدہ مند اے۔ صحیح عدداں نوں قدرتی عدداں دیاں ترتیب وار جوڑیاں دے
ہم ارزیت طبقیاں دے طور تے \emph{بنا} کے، فیر ناطق عدداں نوں صحیح عدداں
دیاں ترتیب وار جوڑیاں دے ہم ارزیت طبقیاں دے طور تے بنا کے، تے فیر
حقیقی عدداں نوں ناطق عدداں دے سیٹاں دے طور تے بنا کے، اسی فوراً ایہناں
بناوٹاں نوں \emph{بھل جاندے} آں۔ خاص طور تے: کوئی وی ریاضیاتی ثبوت
وچ ایہناں بناوٹاں نوں \emph{ورتنا (invoke)} نہیں چاہے گا (بے شک اوس ثبوت توں
سوا جیہڑا وکھاوے کہ بناوٹاں نے اوہو جیہا ورتاؤ کیتا جیہو جیہا اوہناں
نوں کرنا چاہیدا سی)۔ سیٹاں دے سیٹاں دے سیٹاں دے سیٹاں دے سیٹاں دے
سیٹاں دے سیٹاں دے قدرتی عدداں بارے گل کرن نالوں، سیدھا حقیقی عدد بارے
گل کرنا بہت سوکھا اے۔
%(کیوں جو حقیقی عدد ناطق عدداں دے سیٹاں دے طور تے بنائے گئے سن؛ تے ایہ صحیح عدداں دیاں ترتیب وار جوڑیاں دے ہم ارزیت طبقیاں، یعنی صحیح عدداں دے سیٹاں دے سیٹاں دے سیٹاں، دے طور تے بنائے گئے سن؛ وغیرہ۔)
%2/3 = [2,3]
%	یعنی {<2,3>, ... }
%	{ {{2},{2,3}}, ... }
%
%حقیقی عدد
%	ناطق عدداں دے سیٹ نیں
%
%ناطق عدد
%	صحیح عدداں دے سیٹاں دے سیٹاں دے سیٹ نیں
%
%صحیح عدد
%	قدرتی عدداں دے سیٹاں دے سیٹاں دے سیٹ نیں

سب توں ودھ شک ایس گل اُتے اے کہ ایہ تعریفاں سانوں دسدیاں نیں کہ صحیح
عدد، ناطق عدد یا حقیقی عدد \emph{مابعدالطبیعیاتی طور تے \pnbLatin{(metaphysically speaking)}} \emph{کی نیں}۔
یعنی ایس گل اُتے شک اے کہ حقیقی عدد (مثال لئی) کجھ خاص سیٹ (سیٹاں دے
سیٹاں دے سیٹ\ldots) \emph{ای نیں}۔ ایسے نظریے دے راہ وچ وڈی رکاوٹ
ایہ اے کہ بناوٹ کئی وکھرے طریقیاں نال کیتی جا سکدی سی۔ حقیقی عدداں دے
معاملے وچ کجھ واقعی دلچسپ طور تے وکھریاں بناوٹاں نیں (ویکھو
\olref[cauchy]{sec})۔ پر کجھ وکھریاں بناوٹاں حاصل کرن دا بالکل معمولی
طریقہ ایہ اے: \olref[sfr][rel][ref]{sec} وانگ، اسی ترتیب وار جوڑیاں نوں
ذرا وکھرے ڈھنگ نال تعریف کر سکدے ساں؛ جے اسی ترتیب وار جوڑے دا ایہ
متبادل تصور ورتدے، تاں ساڈیاں بناوٹاں عین اوہنی کامیابی نال کم کردیاں،
پر آخر وچ سانوں وکھریاں شےواں ملدیاں۔ ایس لئی صحیح، ناطق تے حقیقی
عدداں دیاں کئی مقابل سیٹ-تھیوری بناوٹاں \pnbLatin{(rival set-theoretic constructions)} نیں۔ تے ہُن ایہ دعویٰ کرنا محض من مانا (arbitrary) تے
شرمندگی والا ہووے گا کہ صحیح عدد (مثال لئی) \emph{ایہ}
سیٹ نیں، \emph{اوہ} نہیں۔ (جیویں \olref[sfr][rel][ref]{sec} وچ، ایہ
\citealt{Benacerraf1965} نال مشہور ہوئی دلیل دی اک مثال اے۔)

اک ہور نکتہ وی چکنا فائدہ مند اے: ساڈیاں بناوٹاں بارے کجھ بڑا
\emph{انجوکھا} اے۔ اسی قدرتی عدداں توں شروع کیتا۔ فیر اسی صحیح عدد
بناؤندے آں، تے ”صحیح عدداں دا $0$“ بناؤندے آں، یعنی
$ \equivrep{0,0}{\Intequiv}$۔ پر $0 \neq
\equivrep{0,0}{\Intequiv}$۔ بلکہ ساڈیاں بناوٹاں دے حساب نال کوئی وی
قدرتی عدد صحیح عدد \emph{نہیں}۔ پر ایہ بہت خلافِ وجدان (counter-intuitive)
لگدا اے۔ بلکہ
\olref[sfr][set][imp]{sec} وچ اسی بہت دلیل دِتے بغیر دعویٰ کیتا سی کہ
$\Nat \subseteq \Rat$۔ جے بناوٹاں سانوں عین \emph{دسدیاں} نیں کہ عدد
کی نیں، تاں ایہ دعویٰ صاف طور تے جھوٹا سی۔

سو کجھ پچھے ہٹ کے ویکھئیے، تاں اسی کتھے پہنچدے آں؟ سادی سیٹ تھیوری وچ
کم کردیاں تے قدرتی عدداں نوں پہلے توں ملیا من کے، اسی صحیح، ناطق تے
حقیقی عدداں نوں کجھ خاص سیٹاں دے طور تے \emph{ورت} سکدے آں۔ ایس معنٰی
وچ اسی ایہناں ہستیاں دے نظرییاں نوں سیٹ تھیوری دے اندر \emph{سما
(embed)} سکدے آں۔ پر ایس سماوٹ دی فلسفیانہ اہمیت \pnbLatin{(philosophical import)} بالکل سیدھی سادی نہیں۔

بے شک ایہناں وچوں کوئی گل آخری فیصلہ نہیں!{} نکتہ صرف ایہ اے۔ ایہ
وکھاؤن لئی کہ حقیقی عدداں دی اعدادی تشکیل دی \emph{واقعی} گہری فلسفیانہ
اہمیت \pnbLatin{(philosophical significance)} اے، کوئی اضافی \emph{فلسفیانہ}
دلیل \pnbLatin{(philosophical argument)} دینا ضروری ہووے گا۔

%ہور ایہ کہ: ایس پورے باب دوران اسی قدرتی عدداں نوں سادہ طور تے پہلے توں
%\emph{ملیا (given)} منیا اے۔ پر اسی اوہناں نال کی کر سکدے آں؟ ایہ اگلے باب دا
%موضوع اے۔
\endgroup
% /OLP-0046
% OLP-0047 content/sets-functions-relations/arithmetization/checking-details.tex
\begingroup
\def\olfilename{checking-details}
\expandafter\def\csname ol@sectioncs\endcsname{section}

\olfileid{sfr}{arith}{check}
\olsection{ترتیب دار حلقے تے میدان \pnbLatin{(Ordered Rings and Fields)}}

ایس پورے باب وچ اسی ایہ دعویٰ کردے آئے آں کہ کجھ تعریفاں ”جیویں اوہناں نوں کرنا
چاہیدا اے“ اوویں ورتاؤ کردیاں نیں۔ ایس فنی ضمیمے وچ اسی کھول کے دساں
گے کہ ساڈا مطلب کی اے، تے (ایہدا خاکہ دیاں گے کہ) تعریفاں ”درست“
ورتاؤ کردیاں نیں۔

\olref[int]{sec} وچ اسی $\Int$ اُتے جمع تے ضرب تعریف کیتے سن۔ اسی
وکھاؤنا چاہنے آں کہ جیویں اوہ تعریف کیتے گئے نیں، اوہ $\Int$ نوں اوہ
ساخت دیندے نیں جیہڑی اسی ایہدے کول ”چاہواں گے“۔ خاص طور تے متعلقہ ساخت
تبادلی حلقے \pnbLatin{(commutative ring)} دی اے۔

\begin{defn}
	اک \emph{تبادلی حلقہ \pnbLatin{(commutative ring)}} اک سیٹ $S$ اے، جیہدے نال خاص رکن $0$ تے $1$
	تے عمل $+$ تے $\times$ دِتے گئے نیں، تے جیہڑا ایہ اَٹھ فارمولے پورے کردا اے:
	\begin{align*}
		\text{\textbf{شرکتیت (Associativity)}}&&a + (b+ c) & = (a + b) + c \\
		&& (a \times b) \times c & = a \times (b\times c)\\
		\text{\textbf{تبادلیت (Commutativity)}}&&a + b &= b+ a  \\
		&&  a \times b&= b\times a\\
		\text{\textbf{شناختی رکن (Identities)}}&&a + 0 &= a \\
		&& a \times 1 &= a\\
		\text{\textbf{جمعی اُلٹ (Additive Inverse)}}&&(\exists b\in S)0&=a + b\\
		\text{\textbf{تقسیمیت (Distributivity)}}&&a \times (b+ c ) &= (a \times b) + (a \times c)
	\end{align*}
	لُکّویں طور تے ایہ سارے $S$ تک محدود کلی مقداریوں \pnbLatin{(universal quantifiers)} نال بندھے ہوئے نیں۔
	تے یاد رکھو کہ ایتھے $0$ تے~$1$ ناں والے رکن لازمی نہیں کہ ایسے ناں
	والے قدرتی عدد ہون۔
\end{defn}

سو ایہ جانچن لئی کہ صحیح عدد اک تبادلی حلقہ \pnbLatin{(commutative ring)} بناؤندے نیں، سانوں صرف ایہ
جانچنا اے کہ اسی ایہ اَٹھ شرطاں پوری کردے آں۔ کوئی شرط ثابت کرنا
{اوکھا} نہیں، پر کم ذرا لمبا اے۔ مثال لئی جمع دے معاملے وچ
\emph{شرکتیت (Associativity)} ایویں ثابت ہوندی اے:

\begin{proof}
کوئی وی $i, j, k \in \Int$ مقرر کرو۔ سو کجھ $a_1, b_1, a_2, b_2, a_3, b_3 \in
\Nat$ نیں جیہناں لئی $i = \equivrep{a_1, b_1}{}$ تے $j =
\equivrep{a_2,b_2}{}$ تے $k = \equivrep{a_3, b_3}{}$۔ (پڑھن دی
سوکھ لئی اسی ”$\equivrep{x, y}{}$“ لکھدے آں، نہ کہ
”$\equivrep{x, y}{\Intequiv}$“؛ ایس پورے حصے وچ ایویں کراں گے۔) ہُن:
\begin{align*}
	i + (j + k) &= \equivrep{a_1, b_1}{}+(\equivrep{a_2, b_2}{} + \equivrep{a_3, b_3}{}) \\
	&= \equivrep{a_1,  b_1}{} + \equivrep{a_2+a_3, b_2+b_3}{}\\
	&= \equivrep{a_1 + (a_2 + a_3), b_1 + (b_2 + b_3)}{}\\
	&= \equivrep{(a_1 + a_2) + a_3, (b_1 + b_2) + b_3}{}\\
	&= \equivrep{a_1 + a_2, b_1 + b_2}{} + \equivrep{a_3, b_3}{}\\
	&= (\equivrep{a_1, b_1}{} + \equivrep{a_2, b_2}{}) + \equivrep{a_3, b_3}{}\\
	&= (i+j) + k
\end{align*}
ایس حساب وچ اسی $\Nat$ اُتے جمع دے معلوم ورتاؤ نوں بے تکلف ورتیا اے۔
\end{proof}

ایسے طرح \emph{جمعی اُلٹ \pnbLatin{(Additive Inverse)}} دا ثبوت ایہ اے:

\begin{proof}
$i \in \Int$ مقرر کرو، سو $i = \equivrep{a,b}{}$، کجھ $a,b \in
\Nat$ لئی۔ $j = \equivrep{b,a}{} \in \Int$ مقرر کرو۔
قدرتی عدداں دے ورتاؤ نوں ورت کے $(a+b) + 0 = 0 + (a+b)$، ایس لئی
تعریف دی روئیں $\tuple{a+b, b+a} \sim_\Int \tuple{0,0}$، تے لہٰذا
$\equivrep{a+b, b+a}{} = \equivrep{0, 0}{} = 0_\Int$۔ سو ہُن $i + j =
\equivrep{a,b}{}+\equivrep{b,a}{}=\equivrep{a+b, b+a}{}= \equivrep{0,
0}{} = 0_\Int$۔
\end{proof}

تے \emph{تقسیمیت (Distributivity)} دا ثبوت ایہ اے:

\begin{proof}
پہلاں وانگ $i = \equivrep{a_1, b_1}{}$ تے $j =
\equivrep{a_2,b_2}{}$ تے $k = \equivrep{a_3, b_3}{}$ مقرر کرو۔ ہُن:
\begin{align*}
	i \times (j + k)
	&= \equivrep{a_1, b_1}{} \times (\equivrep{a_2,b_2}{} + \equivrep{a_3, b_3}{})\\
	&= \equivrep{a_1, b_1}{} \times \equivrep{a_2 + a_3,b_2+b_3}{}\\
	&= \equivrep{a_1  (a_2 + a_3) + b_1  (b_2+b_3), a_1  (b_2 + b_3) + b_1 (a_2 + a_3)}{}\\
	&= \equivrep{a_1 a_2 + a_1a_3 + b_1 b_2+b_1b_3, a_1 b_2 + a_1b_3 + a_2b_1 + a_3b_1}{}\\
%		&= \equivrep{a_1 a_2 + b_1b_2 + a_1a_3 + b_1 b_2+b_1b_3, a_1 b_2 + a_1b_3 + a_2b_1 + a_3b_1}{}\\
	&= \equivrep{a_1a_2 + b_1b_2, a_1b_2 + a_2b_1}{} + \equivrep{a_1a_3 + b_1b_3, a_1b_3 + a_3b_1}{}\\
	&= (\equivrep{a_1, b_1}{} \times \equivrep{a_2,b_2}{}) + (\equivrep{a_1, b_1}{} \times  \equivrep{a_3, b_3}{})\\
	&= (i \times j) + (i \times k)
\end{align*}
\end{proof}

باقی پنج شرطاں دا ثبوت اسی مشق لئی چھڈدے آں۔ اوہ کر لین توں بعد اسی
ایہو وکھا دِتا اے کہ $\Int$ اک تبادلی حلقہ \pnbLatin{(commutative ring)} بناؤندا اے، یعنی جمع تے ضرب
(جیویں اوہ تعریف کیتے گئے نیں) اوویں ورتاؤ کردے نیں جیویں اوہناں نوں
کرنا چاہیدا اے۔

\begin{prob}
ثابت کرو کہ $\Int$ اک تبادلی حلقہ \pnbLatin{(commutative ring)} اے۔
\end{prob}

پر ساڈا کم ہُنّے پورا نہیں ہویا۔ $\Int$ اُتے جمع تے ضرب دے نال نال
اسی اک ترتیبی تعلق $\leq$ وی تعریف کیتا سی، تے سانوں جانچنا اے کہ ایہ
وی اوویں ورتاؤ کردا اے جیویں ایہنوں کرنا چاہیدا اے۔ ہور تفصیل نال سانوں
وکھاؤنا اے کہ $\Int$ اک \emph{ترتیب دار (ordered)} حلقہ بناؤندا اے۔\footnote{
	\olref[sfr][rel][ord]{def:linearorder} توں یاد کرو کہ کُل ترتیب \pnbLatin{(total order)} اوہ
	تعلق اے جیہڑا انعکاسی (reflexive)، متعدی (transitive)، ضد متقارن
	(anti-symmetric) تے متصل (connected) ہووے۔ ترتیبی تعلقاں دے سیاق وچ
	ایس اتصال نوں کدی کدی \emph{سہ شقیت (trichotomy)} آکھیا جاندا اے، کیوں جو ہر
	$a$ تے $b$ لئی ساڈے کول $a \leq b \lor a
	= b \lor a \geq b$ اے۔}

\begin{defn}
اک \emph{ترتیب دار حلقہ \pnbLatin{(ordered ring)}} اوہ تبادلی حلقہ \pnbLatin{(commutative ring)} اے جیہدے نال اک کُل ترتیبی
تعلق \pnbLatin{(total order relation)} $\leq$ وی دِتا گیا اے، تے جیہڑا ایہ شرطاں پوری کردا اے:
\begin{align*}
	a \leq b &\lif a + c \leq b + c\\
	(a \leq b \land 0 \leq c) &\lif a \times c \leq b \times c
\end{align*}
\end{defn}

\begin{prob}
ثابت کرو کہ $\Int$ اک ترتیب دار حلقہ \pnbLatin{(ordered ring)} اے۔
\end{prob}

پہلے وانگ، ایہ وکھاؤنا لمبا پر معمول دا کم اے کہ ساڈی بناوٹ مطابق $\Int$ اک
ترتیب دار حلقہ \pnbLatin{(ordered ring)} اے۔ اسی ایہ کم تہاڈے لئی چھڈدے آں۔

ایس نال صحیح عدداں دا کم پورا ہو گیا۔ پر ہُن سانوں ناطق عدداں بارے
ایہو جیہیاں گلاں وکھاؤنیاں نیں۔ خاص طور تے $+$، $\times$ تے $\leq$ دیاں
ساڈیاں دِتیاں تعریفاں تحت ہُن سانوں وکھاؤنا اے کہ ناطق عدد اک ترتیب دار
\emph{میدان (field)} بناؤندے نیں:
\begin{defn}\ollabel{orderedfield}
اک \emph{ترتیب دار میدان \pnbLatin{(ordered field)}} اوہ ترتیب دار حلقہ \pnbLatin{(ordered ring)} اے جیہڑا ایہ شرط وی پوری کردا اے:
\begin{align*}
	\text{\textbf{ضربی اُلٹ (Multiplicative Inverse)}}& & (\forall a \in S \setminus \{0\})(\exists b \in S) a\times b& = 1
\end{align*}
\end{defn}

جدوں تسیں وکھا دِتا کہ $\Int$ اک ترتیب دار حلقہ \pnbLatin{(ordered ring)} بناؤندا اے، تاں ایہ
وکھاؤنا سوکھا پر لمبا اے کہ $\Rat$ اک ترتیب دار میدان \pnbLatin{(ordered field)} بناؤندا اے۔

\begin{prob}
ثابت کرو کہ $\Rat$ اک ترتیب دار میدان \pnbLatin{(ordered field)} اے۔
\end{prob}

صحیح تے ناطق عدداں نال نمٹ کے صرف حقیقی عدد رہ جاندے نیں۔ خاص طور تے
سانوں وکھاؤنا اے کہ $\Real$ اک \emph{مکمل (complete)} ترتیب دار میدان بناؤندا اے،
یعنی تکمیلیت دی خاصیت \pnbLatin{(Completeness Property)} والا ترتیب دار میدان
\pnbLatin{(ordered field)}۔ ہُن \olref[cuts]{realcompleteness} نے ثابت کیتا کہ
$\Real$ کول تکمیلیت دی خاصیت \pnbLatin{(Completeness Property)} اے۔ پر ہُنّے ایہ
جانچن دا (اُکتاؤنا) کم کرنا باقی اے کہ $\Real$ اک ترتیب دار میدان
\pnbLatin{(ordered field)} اے۔
% PNBARI-008 ماخذ دی درستی: انگریزی جملے وچ ”the (tedious) of checking“ دے وچکار کم والا ناں رہ گیا سی؛ ایتھے صرف جانچن دا کم آکھیا اے۔

اوس لمبی مشق ول \emph{دوڑن} توں پہلاں سانوں کجھ ہور ”فوری“ (immediate) گلاں
جانچنیاں پینیاں نیں۔ مثال لئی سانوں ضمانت چاہیدی اے کہ $\alpha +
\beta$، جیویں ایہ تعریف کیتا گیا اے، ہر کٹ $\alpha$ تے $\beta$ لئی
واقعی اک \emph{کٹ (cut)} اے۔ ایس حقیقت دا ثبوت ایہ اے:

\begin{proof}
کیوں جو $\alpha$ تے $\beta$ دونیں کٹ نیں، $\alpha + \beta = \Setabs{p
+ q}{p \in \alpha \land q \in \beta}$، $\Rat$ دا اک غیر خالی (non-empty)،
حقیقی یعنی پورے سیٹ توں چھوٹا (proper) ذیلی سیٹ اے۔ ہُن فرض کرو کہ $x < p + q$، کجھ $p \in \alpha$ تے
$q \in \beta$ لئی۔ فیر $x - p < q$، سو $x - p \in \beta$، تے
$x = p + (x - p) \in \alpha + \beta$۔ سو $\alpha + \beta$، $\Rat$
دا اک مُڈھلا ٹکڑا \pnbLatin{(initial segment)} اے۔ آخر وچ ہر $p + q \in \alpha + \beta$ لئی، کیوں
جو $\alpha$ تے $\beta$ دونیں کٹ نیں، کجھ $p_1 \in \alpha$ تے
$q_1 \in \beta$ نیں جیہناں لئی $p < p_1$ تے $q < q_1$؛ سو
$p + q < p_1 + q_1 \in \alpha + \beta$؛ سو $\alpha + \beta$ دا
کوئی سب توں وڈا رکن، یعنی maximum، نہیں۔
\end{proof}

ایہو جیہی محنت نال تسیں جانچ سکدے او کہ $\alpha - \beta$ تے
$\alpha \times \beta$ تے $\alpha \div \beta$ کٹ (cuts) نیں (آخری معاملے وچ
$\beta$ دے صفر-کٹ (zero-cut) ہون والے معاملے نوں چھڈ کے)۔ پر اک واری فیر
اسی ایہ کم تہاڈے لئی چھڈدے آں۔
% PNBARI-012 ماخذ دی وضاحت: cuts.tex وچ تفریق تے تقسیم دیاں دکھائیاں تعریفاں غیر فعال تبصری سطراں نیں؛ ایتھے جانچ دی درخواست محفوظ اے، پر اوہ تعریفاں فعال نہیں کیتیاں گئیاں۔

\begin{prob}
ثابت کرو کہ $\Real$ اک ترتیب دار میدان \pnbLatin{(ordered field)} اے۔
\end{prob}

پر ایتھے رہیا ہویا اک چھوٹا سرا سمیٹنا اے۔ \olref[cuts]{sec} وچ اسی
سجھایا سی کہ $\sqrt{2} =
\Setabs{p \in \Rat}{p < 0 \text{ یا }p^2 < 2}$ لیا جا سکدا اے۔ پر سانوں
ایہ وکھاؤنا پینا اے کہ ایہ سیٹ اک \emph{کٹ (cut)} اے۔ ایس حقیقت دا ثبوت ایہ اے:

\begin{proof}
صاف اے کہ ایہ ناطق عدداں دا اک غیر خالی (nonempty)، حقیقی یعنی پورے سیٹ
توں چھوٹا مُڈھلا ٹکڑا \pnbLatin{(proper initial segment)} اے؛ سو اینا وکھاؤنا کافی اے
کہ ایہدا کوئی سب توں وڈا رکن، یعنی maximum، نہیں۔ خاص طور تے اینا
وکھاؤنا کافی اے کہ جتھے $p$ اک مثبت ناطق عدد اے جیہدے لئی $p^2 < 2$
تے $q = \frac{2p+2}{p+2}$، اوتھے $p < q$ تے $q^2 < 2$ دونیں۔
$p < q$ ویکھن لئی صرف ایہ نوٹ کرو:
\begin{align*}
	p^2 &< 2\\
	p^2 + 2p &< 2 + 2p\\
	p(p + 2) &< 2 + 2p\\
	p &< \tfrac{2+2p}{p+2} = q
\end{align*}
$q^2 < 2$ ویکھن لئی صرف ایہ نوٹ کرو:
\begin{align*}
	p^2 &< 2\\
	2p^2 + 4p + 2 &< p^2 + 4p+ 4\\
	4p^2 + 8p + 4 &< 2(p^2 + 4p + 4)\\
	(2p+2)^2 & <2(p+2)^2\\
	\tfrac{(2p+2)^2}{(p+2)^2} &< 2\\
	q^2 &< 2
\end{align*}
\end{proof}
\endgroup
% /OLP-0047
% OLP-0048 content/sets-functions-relations/arithmetization/cauchy.tex
\begingroup
\def\olfilename{cauchy}
\expandafter\def\csname ol@sectioncs\endcsname{section}

\olfileid{sfr}{arith}{cauchy}

\olsection{ضمیمہ: حقیقی عدد بطور کوشی لڑیاں \pnbLatin{(Appendix: the Reals as Cauchy Sequences)}}

\olref[cuts]{sec} وچ اسی حقیقی عدداں نوں ڈیڈیکائنڈ کٹاں دے طور تے
بنایا سی۔ ایس حصے وچ اسی اک متبادل بناوٹ سمجھاؤندے آں۔ ایہ کوشی دی
(اوس شے دی جیہنوں ہُن اسی آکھدے آں) کوشی لڑی \pnbLatin{(Cauchy sequence)} والی تعریف اُتے ٹکی اے؛
پر ایس تعریف نال حقیقی عدد \emph{بنان} دا کم انیسویں صدی دے ہور مصنّفاں،
خاص طور تے وائرشٹراس، ہائن، مے رے تے کانتور، دا اے۔ (چنگی تاریخی روداد
لئی \citeauthor{OConnorRobertson:RN} \citeyear{OConnorRobertson:RN} ویکھو۔)

انیسویں صدی ول جان توں پہلاں سائمن سٹیون (1548--1620) اُتے غور کرنا
فائدے مند اے۔ مختصر ایہ کہ سٹیون نے سمجھیا کہ اسی ہر حقیقی عدد نوں
اوہدی اعشاری پھیلاؤ \pnbLatin{(decimal expansion)} راہیں سوچ سکدے آں۔ ایس لئی حتیٰ کہ $\sqrt{2}$
ورگے غیر ناطق عدد دا وی اک سوہنا اعشاری پھیلاؤ اے، جیہڑا ایویں شروع
ہوندا اے:
\[
	1.41421356237\ldots
\]
سیٹ تھیوری وچ اعشاری پھیلاؤ دا نمونہ بنانا بہت سوکھا اے: اوہناں نوں
بس فنکشن (function) $d \colon \Nat \to \Nat$ سمجھو، جتھے $d(n)$ اوہ $n$واں
اعشاری مقام اے جیہدے وچ ساڈی دلچسپی اے۔ فیر اعشاری نقطے توں پہلاں
آؤندے حقیقی عدد دے حصے (ایتھے صرف $1$) نوں سنبھالن لئی اک چھوٹی تبدیلی
لوڑیں دی اے۔ مثال لئی $0.999\ldots = 1$ دی ضمانت لئی اک ہور تبدیلی
(اک ہم ارزیتی تعلق، equivalence relation) وی لوڑیں دی اے۔ پر سٹیون دے
ڈھنگ نال سیٹ تھیوری اندر حقیقی عدداں دی پوری طرح کڑی (rigorous) بناوٹ دینا اوکھا نہیں۔

سٹیون ساڈا مرکزی موضوع نہیں۔ (سٹیون بارے ہور لئی
\citealt{KatzKatz2012} ویکھو۔) پر ایتھے اک بہت نیڑا خیال اے۔
$\sqrt{2}$ دے اعشاری پھیلاؤ نوں سدھا لین دی بجائے اسی ودھدی درستی
والے پھیلاؤ لے کے، $\sqrt{2}$ دیاں ہور تے ہور درست ناطق تقریبیاں
\pnbLatin{(rational approximations)} دی اک \emph{لڑی (sequence)} اُتے غور کر سکدے آں:
\[
1, 1.4, 1.414, 1.4142, 1.41421,\ldots
\]
حقیقی عدداں نوں ”ہور تے ہور چنگیاں تقریبیاں“ راہیں سمجھ سکّن دا خیال
سانوں بصیرتاں دی اک ہور لڑی دا مُڈھ دیندا اے (اوہناں بصیرتاں ورگی
جیہڑیاں اسی $\Rat$ دی $\Int$ توں، یا $\Int$ دی $\Nat$ توں بناوٹ
ویلے ورتیاں سن)۔ بنیادی بصیرتاں ایہ نیں:
\begin{enumerate}
	\item ہر حقیقی عدد نوں اک (شاید لامتناہی) اعشاری پھیلاؤ دے طور تے لکھیا جا سکدا اے۔
	\item اک (شاید لامتناہی) اعشاری پھیلاؤ وچ بند معلومات نوں ناطق عدداں دی لڑی وی اوہناں ہی پورا بند کر سکدی اے۔
	\item ناطق عدداں دی لڑی نوں $\Nat$ توں $\Rat$ ول فنکشن سمجھیا جا سکدا اے؛ بس $f(n)$ نوں لڑی دا $n$واں ناطق عدد لو۔
\end{enumerate}
بے شک $\Nat$ توں $\Rat$ ول ہر \emph{کوئی وی} فنکشن سانوں حقیقی
عدد نہیں دیوے گا۔ مثال لئی ایہ فنکشن ویکھو:
\[
	f(n) =\begin{cases}
			1 & \text{جے }n\text{ طاق اے}\\
			0 &\text{جے }n\text{ جفت اے}
		\end{cases}
\]
اصل فکر ایہ اے کہ لڑی $0,1,0,1,0,1,0,\ldots$ کسے حقیقی عدد ول
”سمٹدی“ نہیں لگدی۔ سو اوہ لڑیواں لین دی ضمانت لئی جیہڑیاں کسے حقیقی
عدد ول سمٹدیاں نیں، سانوں اپنی توجہ اوہناں لڑیواں تک محدود کرنی پینی
اے جیہناں دی کوئی \emph{حد (limit)} ہووے۔

اسی \olref[his][set][limits]{sec} وچ حد دا خیال پہلے ای ویکھ چکے آں۔
پر اسی عین اوہی تعریف نہیں ورت سکدے جیہڑی اوتھے ورتی سی۔ اوتھے اظہار
”$(\forall \epsilon>0)$“ وچ لُکّویں طور تے حقیقی عدداں اُتے مقداری
بندھن سی؛ نالے اسی حقیقی عدداں اُتے فنکشناں دیاں حداں اُتے غور کر رہے
ساں۔ سو اوس تعریف نوں ورتنا حقیقی عدداں نوں پہلے توں ہی من لینا ہوندا؛
تے اسی عین اوہناں نوں \emph{بنان} لگے آں۔ خوش قسمتی نال اسی حد دا اک
بہت نیڑا خیال ورت سکدے آں۔
\begin{defn}\ollabel{def:CauchySequence}
  فنکشن $f: \Nat \to \Rat$ اک \emph{کوشی لڑی \pnbLatin{(Cauchy sequence)}} اُدوں تے صرف اُدوں
  اے جدوں ہر مثبت $\epsilon \in \Rat$ (ایپسلون) لئی
  $(\exists \ell \in \Nat)(\forall m, n > \ell)|f(m) - f(n)| < \epsilon$ ہووے۔
\end{defn}

حد دا عام خیال پہلے والا ای اے: جے تسیں درستی (precision) دی اک خاص سطح
(جیہڑی~$\epsilon$ نال نپی اے) چاہندے او، تاں ویکھن لئی اک ”علاقہ“
(~$\ell$ توں وڈا ہر انپٹ) اے۔ نالے ایہ ویکھنا سوکھا اے کہ ساڈی لڑی
$1$، $1.4$، $1.414$، $1.4142$، $1.41421$\ldots دی اک حد اے: جے تسیں
$\sqrt{2}$ دی تقریب $\nicefrac{1}{10^{n}}$ توں گھٹ غلطی (error) نال چاہندے او،
تاں بس لڑی وچ $n$ویں رکن توں بعد دا کوئی وی رکن ویکھ لو۔

فیر پہلا صاف خیال ایہ ہووے گا کہ حقیقی عدد بس کوئی وی کوشی لڑی \pnbLatin{(Cauchy sequence)}
\emph{اے}۔ پر $\Int$ تے $\Rat$ دیاں بناوٹاں وانگ ایہ گل بہت سادی
ہووے گی: ہر مقرر حقیقی عدد نوں کئی وکھ وکھ کوشی لڑیاں دَسدیاں نیں۔
ایہ ویکھن دا اک سادہ طریقہ ایہ اے۔ کوئی کوشی لڑی~$f$ دِتی ہووے تاں
$g$ نوں عین $f$ والا فنکشن تعریف کرو، بس $g(0)\neq f(0)$ ہووے۔ کیوں
جو دونیں لڑیاں پہلے عدد توں بعد ہر تھاں اکّو قدر رکھدیاں نیں، اسی آخرکار ایہ
آکھنا چاہواں گے کہ \olref{def:CauchySequence} وچ ورتے معنی مطابق
اوہناں دی حد اکّو اے، تے ایس لئی اوہناں نوں اکّو حقیقی عدد ”تعریف
کردیاں“ سمجھنا چاہیدا اے۔ سو سانوں واقعی ایہناں کوشی لڑیواں نوں اکّو
حقیقی عدد سمجھنا چاہیدا اے۔

ایس لئی سانوں فیر کوشی لڑیواں اُتے اک ہم ارزیتی تعلق \pnbLatin{(equivalence relation)} تعریف کرنا، تے
حقیقی عدداں دی ہم ارزیت \emph{طبقیاں \pnbLatin{(equivalence classes)}} نال شناخت کرنی پینی اے۔ پہلے
سانوں اوس فنکشن دا خیال چاہیدا اے جیہڑا حد وچ $0$ ول جاندا اے۔ ہر
فنکشن $h : \Nat \to \Rat$ لئی آکھو کہ \emph{$h$ حد وچ $0$ ول
جاندا اے \pnbLatin{(tends to zero)}} اُدوں تے صرف اُدوں جدوں ہر مثبت $\epsilon \in \Rat$ لئی
$(\exists \ell \in \Nat)(\forall n > \ell)|h(n)| <
\epsilon$ ہووے۔\footnote{ایہدی \olref[his][set][limits]{sec} وچ
$\lim_{x \mathord{\rightarrow}\infty}f(x) = 0$ دی تعریف نال
تُلا کرو۔} ہور، جتھے $f$ تے $g$، $\Nat \to \Rat$ فنکشن نیں،
مقرر کرو کہ $(f-g)(n) = f(n) - g(n)$۔ ہُن تعریف کرو:
\[
	f \Realequiv g \text{ اُدوں تے صرف اُدوں جدوں $(f-g)$ حد وچ $0$ ول جاندا اے}
.
\]
سانوں جانچنا پینا اے کہ $\Realequiv$ اک ہم ارزیتی تعلق \pnbLatin{(equivalence relation)} اے؛ تے ایہ
اے۔ فیر جے اسی چاہواں، تاں حقیقی عدداں نوں $\Realequiv$ تحت، $\Nat
\to \Rat$ دیاں ساریاں کوشی لڑیواں دے ہم ارزیت طبقے \pnbLatin{(equivalence classes)} تعریف کر سکدے آں۔
% PNBARI-009 ماخذ دی درستی: حقیقی عدداں دی شناخت ہم ارزیتی تعلق نال نہیں، اوس تعلق تحت کوشی لڑیواں دے طبقیاں نال ہوندی اے۔

\begin{prob}
$f(n) = 0$، ہر $n$ لئی، لو۔ $g(n) = \frac{1}{(n+1)^2}$ لو۔ وکھاؤ
کہ دونیں کوشی لڑیاں \pnbLatin{(Cauchy sequences)} نیں، تے واقعی دونیں فنکشناں دی حد $0$ اے؛ ایس لئی
$f \sim_\Real g$ وی اے۔
\end{prob}

ایہ کر لین توں بعد اسی عام وانگ اوس ہم ارزیت طبقے لئی
$\equivrep{f}{\Realequiv}$ لکھاں گے جیہدے وچ $f$، !!a{element} یعنی رکن اے۔
پر لکھت نوں پڑھّن جوگ رکھن لئی اگے اسی زیریں اشاریہ چھڈ کے صرف
$\equivrep{f}{}$ لکھاں گے۔ نالے اسی مقرر کردے آں کہ ہر $q \in \Rat$
لئی $q_{\Real} = \equivrep{c_{q}}{}$، جتھے $c_{q}$ مستقل فنکشن
(constant function، یعنی مستقل لڑی) اے،
یعنی $c_q(n) = q$، ہر $n \in \Nat$ لئی۔ فیر اسی حقیقی عدداں اُتے
بنیادی تعلق تے عمل تعریف کردے آں، مثال لئی:
\begin{align*}
	\equivrep{f}{} + \equivrep{g}{} &= 	\equivrep{(f + g)}{} \\
	\equivrep{f}{} \times 	\equivrep{g}{} &= \equivrep{(f \times g)}{}
\end{align*}
جتھے $(f + g)(n) = f(n) + g(n)$ تے $(f \times g)(n) = f(n) \times
g(n)$۔ بے شک سانوں ایہ وی جانچنا پینا اے کہ $(f + g)$، $(f-g)$ تے
$(f\times g)$ وچوں ہر اک وی کوشی لڑی اے، جدوں $f$ تے $g$ کوشی لڑیاں
ہون؛ واقعی اوہ نیں، تے اسی ایہ جانچ تہاڈے لئی چھڈدے آں۔

آخر وچ اسی ترتیب (order) دا اک تصور تعریف کردے آں۔ آکھو کہ
$\equivrep{f}{}$ \emph{مثبت (positive)} اے اُدوں تے صرف اُدوں جدوں
$\equivrep{f}{}\neq 0_\Real$ تے
$(\exists \ell \in \Nat)(\forall n > \ell)0 < f(n)$ دونیں ہون۔
فیر آکھو کہ $\equivrep{f}{} < \equivrep{g}{}$ اُدوں تے صرف اُدوں
جدوں $\equivrep{(g - f)}{}$ مثبت اے۔ سانوں جانچنا پینا اے کہ ایہ
صحیح طور تے تعریف شدہ (well-defined) اے (یعنی ایہ ہم ارزیت طبقے وچوں ”نمائندہ“ (representative) فنکشن
دی چُون اُتے منحصر نہیں)۔
% PNBARI-010 ماخذ دی درستی: حقیقی ہم ارزیت طبقے دی غیر صفری شرط اندر صحیح ٹائپ والا صفر 0_Real اے، 0_Rat نہیں۔
%\begin{align*}
%	\equivrep{f}{} \leq \equivrep{g}{} \text{ اُدوں تے صرف اُدوں جدوں }&\equivrep{f}{} = \equivrep{g}{} \text{ یا }\\
%	&(\exists \ell \in \Nat)(\forall n \in \Nat)(\ell < n \lif f(n) \leq g(n))
%\end{align*}
پر ایہ کر لین توں بعد ایہ وکھاؤنا کافی سوکھا اے کہ ایہ تعریفاں درست
الجبری خاصیتاں دیندیاں نیں؛ یعنی:
\begin{thm}\ollabel{thm:cauchyorderedfield}
کوشی لڑیواں دے ہم ارزیت طبقے \pnbLatin{(equivalence classes)} اک ترتیب دار میدان
\pnbLatin{(ordered field)} بناؤندے نیں۔
% PNBARI-011 ماخذ دی درستی: ایتھے، اگلی مشق تے تکمیلیت دے قضیے وچ میدان دیاں ہستیاں خام کوشی لڑیاں نہیں بلکہ اوہناں دے ہم ارزیت طبقے نیں؛ اگلے ثبوت وچ نمائندہ لڑیاں والا نشان محفوظ اے۔
\end{thm}

\begin{proof}
ایہ مشق اے۔
\end{proof}

\begin{prob}
ثابت کرو کہ کوشی لڑیواں \pnbLatin{(Cauchy sequences)} دے ہم ارزیت طبقے
\pnbLatin{(equivalence classes)} اک ترتیب دار میدان \pnbLatin{(ordered field)} بناؤندے نیں۔
\end{prob}

ایویں بنائے حقیقی عدداں کول تکمیلیت دی خاصیت \pnbLatin{(Completeness Property)} اے،
ایہ ثابت کرنا ہور اوکھا اے، سو اسی ایہدا ثبوت دیندے آں۔

\begin{thm}
کوشی لڑیواں دے ہم ارزیت طبقیاں \pnbLatin{(equivalence classes)} دا ہر غیر خالی
(non-empty) سیٹ، جیہدی بالائی حد \pnbLatin{(upper bound)} ہووے، اک سب توں چھوٹی
بالائی حد \pnbLatin{(least upper bound)} رکھدا اے۔
\end{thm}

\begin{proof}[ثبوت دا خاکہ \pnbLatin{(Proof sketch)}] $S$ کوشی لڑیواں دے ہم ارزیت
طبقیاں دا کوئی وی غیر خالی سیٹ ہووے جیہدی اک بالائی حد اے؛ اگلے نشان
وچ ہر طبقے لئی اوہدی نمائندہ لڑی لکھی جا رہی اے۔ سو کوئی $p \in \Rat$
اے جیہدے لئی $p_{\Real}$، $S$ دی بالائی حد اے۔ $r \in S$ لو؛ فیر کوئی $q \in \Rat$ اے جیہدے لئی
$q_{\Real} < r$۔ سو جے $S$ دی سب توں چھوٹی بالائی حد موجود اے، تاں
اوہ $q_\Real$ تے $p_\Real$ دے وچکار، دونیں سرے شامل کر کے (inclusive)، اے۔

اسی بازگشتی نصف بندی \pnbLatin{(recursive bisection)} نال سب توں چھوٹی بالائی حد
\pnbLatin{(least upper bound)} ول اکّو ویلے ہِٹھاں تے اُتّوں نزدیک آ کے سمٹاں گے۔
خاص طور تے اسی دو فنکشن $f, g \colon \Nat \to \Rat$
تعریف کردے آں؛ مقصد ایہ اے کہ $f$ اُتّوں اوس حد ول سمٹے، تے $g$ ہِٹھاں
توں اوس ول سمٹے۔ اسی ایہ تعریفاں توں شروع کردے آں:
\begin{align*}
	f(0) &= p \\
	g(0) &= q
\end{align*}
فیر جتھے $a_n = \frac{f(n) + g(n)}{2}$، مقرر کرو:\footnote{ایہ اک
بازگشتی تعریف اے۔ پر اسی ہُن \emph{تک} ایہ سوچن دی کوئی وجہ نہیں
دِتی کہ بازگشتی تعریفاں درست نیں۔}
\begin{align*}
	f(n+1) &=
	\begin{cases}
		a_n &\text{جے }(\forall h \in S)\equivrep{h}{} \leq (a_n)_\Real\\
		f(n)&\text{نہیں تاں}
	\end{cases}\\
	g(n+1) &=
	\begin{cases}
		a_n &\text{جے }(\exists h \in S)\equivrep{h}{} \geq (a_n)_\Real\\
	 	g(n) &\text{نہیں تاں}
	\end{cases}
\end{align*}
$f$ تے $g$ دونیں کوشی لڑیاں \pnbLatin{(Cauchy sequences)} نیں۔ (ایہ کافی سوکھ نال جانچیا جا سکدا
اے، پر اسی ایہنوں مشق لئی چھڈدے آں۔) نوٹ کرو کہ فنکشن $(f-g)$ حد وچ
$0$ ول جاندا اے \pnbLatin{(tends to zero)}، کیوں جو ہر قدم اُتے $f$ تے $g$ وچلا فرق آدھا ہو جاندا
اے۔ ایس لئی $\equivrep{f}{} = \equivrep{g}{}$۔

پہلاں اسی وکھاؤندے آں کہ $\equivrep{f}{}$، $S$ دی بالائی حد \pnbLatin{(upper bound)} اے، یعنی
$(\forall h \in S)\equivrep{h}{} \leq \equivrep{f}{}$۔
(اگے ودھدے اسی \olref{thm:cauchyorderedfield} نوں ورتاں گے۔) $h \in S$
لو تے تضاد لئی فرض کرو کہ $\equivrep{f}{} < \equivrep{h}{}$، سو
$0_\Real < \equivrep{(h-f)}{}$۔ کیوں جو $f$ اک اک رُخ نال گھٹدی
\pnbLatin{(monotonically decreasing)} کوشی لڑی اے، کوئی $n \in \Nat$ اے جیہدے لئی
$\equivrep{(c_{f(n)} - f)}{} < \equivrep{(h-f)}{}$۔ سو:
\[
	(f(n))_\Real = \equivrep{c_{f(n)}}{} < \equivrep{f}{} + \equivrep{(h-f)}{} = \equivrep{h}{},
\]
پر ایہ اوس حقیقت دے خلاف اے کہ بناوٹ مطابق
$\equivrep{h}{} \leq (f(n))_\Real$۔

اگے اسی وکھاؤندے آں کہ $\equivrep{f}{} = \equivrep{g}{}$، $S$ دی
\emph{سب توں چھوٹی (least)} بالائی حد اے۔ سو $j$ کوئی وی کوشی لڑی
\pnbLatin{(Cauchy sequence)} ہووے تے
فرض کرو کہ $\equivrep{j}{} < \equivrep{g}{}$۔ اوپر وانگ دلیل نال
(ایہ حقیقت ورت کے کہ $g$ \emph{ودھدی (increasing)} اے)، کوئی $n \in \Nat$ اے
جیہدے لئی $\equivrep{j}{} < (g(n))_\Real$۔ پر بناوٹ مطابق کوئی
$h \in S$ اے جیہدے لئی $(g(n))_\Real \leq \equivrep{h}{}$، سو
$\equivrep{j}{} < \equivrep{h}{}$، تے ایس لئی $\equivrep{j}{}$،
$S$ دی بالائی حد نہیں۔
\end{proof}
\endgroup
% /OLP-0048

\OLEndChapterHook
\endgroup
% /OLP-0041

% OLP-0049 content/sets-functions-relations/infinite/infinite.tex
\begingroup
\def\olfilename{infinite}
\expandafter\def\csname ol@sectioncs\endcsname{section}

\olchapter{sfr}{infinite}{لامتناہی سیٹ \pnbLatin{(Infinite Sets)}}

\begin{editorial}
لامتناہی سیٹاں بارے ایہ باب ٹِم بٹن دی \emph{اوپن سیٹ تھیوری
\pnbLatin{(Open Set Theory)}} توں لیا گیا اے۔
\end{editorial}

% OLP-0050 content/sets-functions-relations/infinite/hilberts-hotel.tex
\begingroup
\def\olfilename{hilberts-hotel}
\expandafter\def\csname ol@sectioncs\endcsname{section}

\olfileid{sfr}{infinite}{hilbert}
\olsection{ہلبرٹ دا ہوٹل \pnbLatin{(Hilbert's Hotel)}}

قدرتی عدداں دا سیٹ بدیہی طور تے لامتناہی (infinite) اے۔ ایس لئی، جے
اسی قدرتی عدداں نوں \emph{پہلاں ای ورت} نہیں لینا چاہندے، تاں ساڈا
پہلا قدم ایہ ہونا چاہیدا اے کہ لامتناہی سیٹ دی پہچان ایسے لفظاں وچ
بیان کریے جیہڑے خود قدرتی عدداں دا ذکر نہ منگن۔ 1924 دے اک لیکچر وچ
ہلبرٹ نے ایس دا اک سوہنا طریقہ پیش کیتا۔ اوہ سانوں ایہ تصور کرن لئی
آکھدا اے:
% PNBINF-001 ماخذ دی درستی: ”our first step must be characterize“ وچ مصدر دا نشان رہ گیا سی؛ ایتھے پہلا قدم خاصیت دَسنا آکھیا اے۔
\begin{quote}
[\ldots] اک ہوٹل جیہدے کمراں دی گنتی متناہی (finite) اے۔ سارے کمراں
وچ عین اک اک مہمان ٹھہریا ہووے۔ جے مہمان ہُن کسے وی ڈھنگ نال اپنے
کمرے وٹا لَین، [پر] وٹاندرا ایس طرح ہووے کہ ہر کمرے وچ ہجے وی اک توں
ودھ بندہ نہ ہووے، تاں کوئی وی کمرہ خالی نہیں ہووے گا، تے ہوٹل دا مالک
ایس طریقے نال ہُن آؤن والے مہمان لئی نویں تھاں نہیں بنا سکدا [\ldots
\textparagraph \ldots]

ہُن اسی فرض کردے آں کہ ہوٹل وچ لامتناہی گنتی دے \pnbLatin{(infinitely many)}
نمبر لائے کمرے $1$، $2$، $3$، $4$، $5$، \dots، ہون، تے ہر کمرے وچ
عین اک مہمان ٹھہریا ہووے۔ جیویں ای کوئی نواں مہمان آوے، مالک نوں بس
ہر پرانے مہمان نوں اوس کمرے وچ بھیجنا اے جیہدا نمبر اک ودھ اے
\pnbLatin{(successor-numbered room)}، تے کمرہ~$1$ نویں آؤندے مہمان لئی خالی ہو
جاوے گا۔
	\begin{center}
		\begin{tikzpicture}[scale = .75]
		\foreach \x in {1, 2, 3, 4, 5, 6, 7, 8, 9}
		{
			\node (\x a) at (\x, 1) {\small{\x}};
			\node (\x b) at (\x, 2) {\small{\x}};
		}
		\node (dotsa) at (10, 1) {\small{\ldots}};
		\node (dotsb) at (10, 2) {\small{\ldots}};
		\draw[->] (1b)--(2a);
		\draw[->] (2b)--(3a);
		\draw[->] (3b)--(4a);
		\draw[->] (4b)--(5a);
		\draw[->] (5b)--(6a);
		\draw[->] (6b)--(7a);
		\draw[->] (7b)--(8a);
		\draw[->] (8b)--(9a);
		\draw[->] (9b)--(dotsa);
		\draw (1,1) circle (.4);
		\end{tikzpicture}
	\end{center}
	(\citealt[730]{EwaldSieg2013} وچ چھپیا؛ ماخذ مصنّفاں دا ترجمہ
	\pnbLatin{(source authors' translation)})
\end{quote}
فیصلہ کن گل ایہ اے کہ ہلبرٹ دے ہوٹل وچ کمرے لامتناہی نیں؛ اوہدی
سمجھاوت نوں اسی ایس دعوے دے معنی دی تعریف بنا سکدے آں۔ اصل وچ ایہ
ڈیڈیکائنڈ دا طریقہ سی (بے شک ایتھے ایہنوں بہت وڈی زمانوی بے ترتیبی
(anachronism) نال پیش کیتا گیا اے؛ ڈیڈیکائنڈ دی تعریف
\citeyear{Dedekind1888} دی اے):

\begin{defn}\ollabel{defn:DedekindInfinite}
سیٹ $A$ نوں \emph{ڈیڈیکائنڈ لامتناہی \pnbLatin{(Dedekind infinite)}} اُدوں تے
صرف اُدوں آکھدے آں جدوں $A$ توں $A$ دے اک حقیقی، یعنی پورے سیٹ توں
چھوٹے، ذیلی سیٹ \pnbLatin{(proper subset)} ول کوئی !!a{injection} \pnbLatin{(injective map)}
موجود ہووے۔ برابری نال، کوئی $o \in A$ تے !!a{injection} \pnbLatin{(injective map)} $f \colon A \to A$ ایہو جیہے ہون کہ $o \notin \ran{f}$؛ یعنی اوہ
رکن نقشے دی رینج (range) وچ نہ آوے۔
\end{defn}
\endgroup
% /OLP-0050
% OLP-0051 content/sets-functions-relations/infinite/dedekind-algebra.tex
\begingroup
\def\olfilename{dedekind-algebra}
\expandafter\def\csname ol@sectioncs\endcsname{section}

\olfileid{sfr}{infinite}{dedekind}
\olsection{ڈیڈیکائنڈ الجبرے \pnbLatin{(Dedekind Algebras)}}

ساڈا مقصد صرف ایہ نہیں کہ قدرتی عدد لامتناہی ہون؛ اوہناں وچ کجھ
(الجبری) خاصیتاں وی ہونیاں چاہیدیاں نیں، تاں جو جمع، ضرب تے ہور
عمل اوہناں اُتے ٹھیک ڈھنگ نال ورتن۔

ڈیڈیکائنڈ دا خیال سی کہ \emph{جانشین فنکشن \pnbLatin{(successor function)}} نوں
بنیادی من کے عدداں دی پہچان ہِیٹھ دِتیاں شرطاں نال کیتی جاوے:
\begin{enumerate}
	\item اک عدد $0$ اے جیہڑا کسے وی عدد دا جانشین نہیں
	\\i.e., $0 \notin \ran{s}$
	\\i.e., $\forall x\ s(x) \neq 0$
	\item وکھ وکھ عدداں دے جانشین وی وکھ وکھ ہوندے نیں
	\\i.e., $s$ اک !!a{injection} اے
	\\i.e., $\forall x \forall y (s(x) = s(y) \lif x = y)$
	\item\ollabel{repeatedapplication} ہر عدد، جانشین فنکشن دے وار وار
	اطلاق \pnbLatin{(repeated application)} نال $0$ توں حاصل ہوندا اے۔
\end{enumerate}
پہلیاں دو شرطاں نوں پہلے درجے دی منطق \pnbLatin{(first-order logic)} وچ بیان
کرنا سوکھا اے (اُتّے ویکھو)، پر صرف ایس منطق نال
\olref{repeatedapplication} نوں نہیں پھڑیا جا سکدا۔ ڈیڈیکائنڈ دی وڈی
پیش رفت ایہ سی کہ شرط \olref{repeatedapplication} نوں سیٹ تھیوری دے
لفظاں وچ ایس طرح نواں بیان کیتا جاوے:
\begin{enumerate}
	\item[3$'$.] قدرتی عدد اوہ سب توں چھوٹا سیٹ نیں جیہڑا
	جانشین فنکشن \emph{تحت بند \pnbLatin{(closed under)}} اے: یعنی سیٹ دے کسے
	وی !!{element} اُتے $s$ لاگو کریے، تاں نتیجے وچ ایس سیٹ دا اک
	ہور !!{element} ملدا اے۔
\end{enumerate}
ہُن ایس گل نوں قدم بہ قدم کھولنا ضروری اے۔

\begin{defn}\ollabel{Closure}
	ہر ایسے فنکشن $f$ لئی جیہڑا اپنے حامل سیٹ \pnbLatin{(carrier set)} دا خود نقشہ
	(self-map/endofunction) اے، حامل دا ذیلی سیٹ $X$، $f$-\emph{بند}
	{اُدوں تے صرف اُدوں} اے جدوں $(\forall x \in X)f(x) \in X$۔ ہُن
	اوسے حامل دے ہر $o$ لئی تعریف کرو:
	$$\closureofunder{f}{o} = \bigcap\Setabs{X}{o \in X\text{ تے }X
\text{، $f$-بند اے}}$$
\end{defn}
% PNBINF-002 ماخذ دی درستی: تعریف وچ f نوں اپنے حامل سیٹ دا خود نقشہ، X نوں حامل دا ذیلی سیٹ تے o نوں اوس حامل دا رکن مقرر کیتا اے؛ ایس دائرے بناں اگلی سطر دا A بے تعارف رہندا تے حاصل قدراں دا گواہ لازماً بند نہیں ہوندا۔

ایس لئی $\closureofunder{f}{o}$ اوہناں سارے $f$-بند سیٹاں دا اشتراک
(intersection) اے جیہناں وچ $o$، !!a{element} اے۔ بدیہی طور تے
$\closureofunder{f}{o}$، \emph{سب توں چھوٹا} $f$-بند سیٹ، یعنی سب توں
چھوٹی پیدا کیتی بندش \pnbLatin{(least generated closure)} اے جیہدے وچ $o$،
!!a{element} اے۔ اگلا نتیجہ ایس بدیہی خیال نوں باقاعدہ درست ثابت کردا اے:
\begin{lem}\ollabel{closureproperties}
	ہر ایسے فنکشن $f$ تے ہر $o \in A$ لئی، جتھے A اوس خود نقشے دا حامل سیٹ اے:
	\begin{enumerate}
		\item\ollabel{closurehaselem} $o \in \closureofunder{f}{o}$؛ تے
		\item\ollabel{closureclosed} $\closureofunder{f}{o}$، $f$-بند اے؛ تے
		\item\ollabel{closuresmallest} جے $X$، $f$-بند اے تے $o \in
		X$، تاں $\closureofunder{f}{o} \subseteq X$
	\end{enumerate}
\end{lem}

\begin{proof}
پہلاں ویکھو کہ گھٹ توں گھٹ اک $f$-بند سیٹ ایسا اے جیہدے وچ $o$،
!!a{element} اے: ایہ رینج (range) نال خاص رکن ملا کے بنیا سیٹ
$\ran{f}\cup \{o\}$ اے۔ ایس لئی ایسے \emph{سارے} سیٹاں دا اشتراک
$\closureofunder{f}{o}$ موجود اے۔ ہُن \olref{closurehaselem}--\olref{closuresmallest}
دیاں تین خاصیتاں جانچدے آں۔

\olref{closurehaselem} لئی، $o \in \closureofunder{f}{o}$؛ کیوں جو
اشتراک وچ شامل ہر سیٹ اندر $o$، !!a{element} اے۔

\olref{closureclosed} لئی فرض کرو $x \in \closureofunder{f}{o}$۔
جدوں وی $o \in X$ تے $X$، $f$-بند اے، اوہدوں $x \in X$؛ تے
$f(x) \in X$، کیوں جو $X$، $f$-بند اے۔ ایس لئی
$f(x) \in \closureofunder{f}{o}$۔

\olref{closuresmallest} لئی اشتراک دی عام خاصیت ورتو: جے $X
\in C$، تاں لازماً $\bigcap C \subseteq X$۔
\end{proof}

ایس نتیجے دی مدد نال ہُن ایہ تعریف دِتی جا سکدی اے:

\begin{defn}
اک \emph{ڈیڈیکائنڈ الجبرا \pnbLatin{(Dedekind algebra)}} اک سیٹ $A$، اوس اُتے
خود نقشہ (self-map/endofunction) $f \colon A \to A$، تے اک خاص رکن
$o \in A$ توں بندا اے، جیہڑے ایہ شرطاں پوری کردے نیں:
	\begin{enumerate}
		\item \ollabel{ded:proper} $o \notin \ran{f}$
		\item \ollabel{ded:injection} $f$ اک !!a{injection} \pnbLatin{(injective map)} اے
		\item \ollabel{ded:closure} $A = \closureofunder{f}{o}$
	\end{enumerate}
\end{defn}

چونکہ $A = \closureofunder{f}{o}$ اے، پچھلا نتیجہ دَسدا اے کہ $A$
اوہ سب توں چھوٹا $f$-بند سیٹ اے جیہدے وچ $o$، !!a{element} اے۔
تعریف دیاں شرطاں \olref{ded:proper} تے \olref{ded:injection} توں فوراً
صاف اے کہ ہر ڈیڈیکائنڈ الجبرا، ڈیڈیکائنڈ لامتناہی \pnbLatin{(Dedekind infinite)} اے۔ ہور دلچسپ گل ایہ اے کہ الٹ پاسے ہر ڈیڈیکائنڈ
لامتناہی سیٹ توں اک ڈیڈیکائنڈ الجبرا بنایا جا سکدا اے۔

\begin{thm}\ollabel{thm:DedekindInfiniteAlgebra}
جے ڈیڈیکائنڈ لامتناہی \pnbLatin{(Dedekind infinite)} سیٹ موجود اے، تاں اک
ڈیڈیکائنڈ الجبرا \pnbLatin{(Dedekind algebra)} وی موجود اے۔
\end{thm}

\begin{proof}
$D$ نوں ڈیڈیکائنڈ لامتناہی منو۔ فیر اک انجیکٹو نقشہ \pnbLatin{(injective map)}
$g \colon D \to D$ تے اک ایسا رکن $o  \in D \setminus \ran{g}$ اے
جیہڑا رینج وچ نہیں۔ ہُن $A = \closureofunder{g}{o}$ لو؛
\olref{closureproperties} مطابق $A$ موجود اے تے $o \in A$۔ اگے فنکشن
دی تحدید \pnbLatin{(function restriction)} $f = \funrestrictionto{g}{A}$ لو۔
اسی ثابت کراں گے کہ $A, f, o$ رل کے اک ڈیڈیکائنڈ الجبرا بناؤندے نیں۔

\olref{ded:proper} لئی: $o \notin \ran{g}$ تے $\ran{f}
\subseteq \ran{g}$؛ ایس رینج شمولیت توں $o\notin \ran{f}$۔

\olref{ded:injection} لئی: $g$، $D$ اُتے انجیکٹو نقشہ اے؛ ایس لئی
$f \subseteq g$ وی لازماً اک انجیکشن (injection) اے۔

\olref{ded:closure} لئی: \olref{closureproperties} مطابق $A$،
$g$-بند اے، تے ایس کرکے $A$، $f$-بند وی اے۔ لہٰذا
$\closureofunder{f}{o} \subseteq A$، \olref{closureproperties} مطابق۔
دوجے پاسے $\closureofunder{f}{o}$ خود $f$-بند اے، تے چونکہ
$f = \funrestrictionto{g}{A}$ اے، ایس لئی $\closureofunder{f}{o}$،
$g$-بند وی اے۔
پیدا کیتی بندش دی سب توں چھوٹی ہون والی خاصیت نال
$A \subseteq \closureofunder{f}{o}$، \olref{closureproperties} مطابق۔
\end{proof}
\endgroup
% /OLP-0051
% OLP-0052 content/sets-functions-relations/infinite/dedekind-induction.tex
\begingroup
\def\olfilename{dedekind-induction}
\expandafter\def\csname ol@sectioncs\endcsname{section}

\olfileid{sfr}{infinite}{induction}
\olsection[حسابی استقرا \pnbLatin{(Arithmetical Induction)}]{ڈیڈیکائنڈ الجبرے تے حسابی استقرا \pnbLatin{(Dedekind Algebras and Arithmetical Induction)}}

ہُن اہم گل ایہ اے کہ اک ڈیڈیکائنڈ الجبرا---بلکہ
\emph{ہر} ڈیڈیکائنڈ الجبرا---قدرتی عدداں دے قائم مقام (surrogate)
وجوں کم کر سکدا اے۔ ایہ اگلے سیدھے نتیجے توں ملدا اے:

\begin{thm}[حسابی استقرا \pnbLatin{(Arithmetical induction)}]\ollabel{thm:dedinfiniteinduction}
$N, s, o$ رل کے اک ڈیڈیکائنڈ الجبرا بناؤن۔ تاں ہر سیٹ $X$ لئی:
	\begin{center}
		جے $o \in X$ تے $(\forall {n} \in N \cap X){s}({n}) \in X$، {تاں} نتیجتاً $N \subseteq X$۔
	\end{center}
\end{thm}

\begin{proof}
ڈیڈیکائنڈ الجبرے دی تعریف نال $N = \closureofunder{s}{o}$۔ جے
${o} \in X$ تے $(\forall {n} \in N)(n \in X \lif
{s}({n}) \in X)$ دونیں شرطاں پوری ہون، تاں حامل سیٹ تے دِتے سیٹ دا
اشتراک، خاص رکن نوں رکھن والا جانشین-بند حامل ذیلی سیٹ اے۔ سب توں
چھوٹی پیدا کیتی بندش \pnbLatin{(least generated closure)} دی خاصیت نال
$N = \closureofunder{s}{o} \subseteq X$۔
% PNBINF-003 ماخذ دی درستی: قضیہ X نوں ہر سیٹ مندا اے، ایس لئی بندش دا چھوٹاپن N نال X دے اشتراک اُتے لایا اے؛ ماخذ دا ثبوت X نوں سدھا s-بند آکھ کے حامل توں باہر دے ممکنہ رکناں نوں نظر انداز کردا اے۔
\end{proof}

استقرا قدرتی عدداں دی خاص پہچان اے، سو نکتہ ایہ بن دا اے: کوئی وی
ڈیڈیکائنڈ لامتناہی سیٹ دِتا ہووے، اوس توں اک ڈیڈیکائنڈ الجبرا بنا کے
اس الجبرے نوں قدرتی عدداں دے قائم مقام (surrogate) وجوں ورتیا جا
سکدا اے۔

ایہ درست اے کہ \olref{thm:dedinfiniteinduction} استقرا نوں
\emph{سیٹ-تھیوری والے (set-theoretic)} لفظاں وچ بیان کردا اے۔ پر
ایسے اصول نوں ہور جانے پچھانے روپ وچ وی سوکھ نال لکھیا جا سکدا اے:

\begin{cor}\ollabel{natinductionschema}
$N, s, o$ رل کے اک ڈیڈیکائنڈ الجبرا بناؤن۔ فیر استقرائی خاکے
\pnbLatin{(induction schema)} مطابق ہر فارمولے $\phi(x)$ لئی، جیہدے وچ پیرامیٹر
(parameters) ہو سکدے نیں:
\begin{center}
	جے $\phi(o)$ تے $(\forall {n} \in N)(\phi(n)\lif
	\phi({s}({n})))$، {تاں} نتیجتاً $(\forall n \in N)\phi(n)$
\end{center}
\end{cor}

\begin{proof}
$X = \Setabs{n \in N}{\phi(n)}$ مقرر کرو، تے فیر
\olref{thm:dedinfiniteinduction} لاگو کرو۔
\end{proof}

ایس نتیجے وچ فارمولے دے ”پیرامیٹر (parameters) رکھن“ دی گل ہوئی اے۔
موٹے طور تے مطلب ایہ اے کہ کوئی وی شےواں $c_1, \ldots, c_k$ لئی
اسی $\phi(x, c_1, \ldots, c_k)$ نال کم کر سکدے آں۔ ہور ٹھیک بیان
وچ ”پیرامیٹر“ دا ناں لین دی لوڑ نہیں: ہر فارمولے
$\phi(x, v_1, \ldots, v_k)$ لئی، جیہدے منطق والے سارے آزاد متغیر
\pnbLatin{(free variables)} سامنے لکھے گئے نیں---ایتھے آزاد متغیر دا مطلب
differential equations والا independent variable نہیں---ساڈے کول ایہ اے:
	\begin{align*}
			\forall v_1 \ldots \forall v_k((&\phi(o, v_1,\ldots, v_k) \land {}\\
			&	(\forall x \in N)(\phi(x,v_1, \ldots, v_k) \lif \phi(s(x), v_1,\ldots, v_k))) \lif {}\\
			&\hspace{3em} (\forall x \in N)\phi(x, v_1,\ldots, v_k))
	\end{align*}
صاف طور تے ”پیرامیٹر رکھن“ والی مختصر گل لکھت نوں بہت سوکھا بنا
دیندی اے۔ (\olref[sth][][]{part} وچ اسی ایہ طریقہ کئی وار ورتاں گے۔)

ہُن ڈیڈیکائنڈ الجبریاں ول واپس آؤ۔ کوئی وی ڈیڈیکائنڈ الجبرا دِتا
ہووے، اسی جمع، ضرب تے قوت چڑھاؤ (exponentiation) دے عام حسابی فنکشن
وی تعریف کر سکدے آں۔ ایہ کم معمولی نہیں: ایس وچ
\emph{بازگشتی تعریف \pnbLatin{(recursive definition)}} دی ترکیب ورتی جاندی اے۔
اسی ایس ترکیب نوں بہت بعد وچ، اک ہور عام سیاق اندر، متعارف کرا کے
جائز ثابت کراں گے۔ (دلچسپی رکھن والے
\olref[sth][ord-arithmetic][]{chap} توں بعد ایتھے مڑ سکدے نیں، یا
\citealt[pp.~95--8]{Potter2004} ورگا کوئی متبادل بیان پڑھ سکدے نیں۔)
پر جدوں $N, s, o$ رل کے اک ڈیڈیکائنڈ الجبرا بناؤن، آخرکار اسی ہِیٹھ
دِتیاں چھے بازگشتی مساوات مقرر کر سکاں گے:
\begin{align*}
	{a} + {o} &= {a} & & & {a} \times {o} &= {o} & & & {a}^{o} &= s(o)\\
{a} + {s}({b}) &= {s}({a}+{b}) &&& {a} \times {s}({b}) &= ({a}\times {b}) + {a}  & & & {a}^{{s}({b})} &= {a}^{b} \times {a}
\end{align*}
تے فیر ثابت کر سکاں گے کہ ایہ اوویں ای ورتاؤ کردیاں نیں جیویں امید اے۔
\endgroup
% /OLP-0052
% OLP-0053 content/sets-functions-relations/infinite/dedekinds-proof.tex
\begingroup
\def\olfilename{dedekinds-proof}
\expandafter\def\csname ol@sectioncs\endcsname{section}

\olfileid{sfr}{infinite}{dedekindsproof}

\olsection[ڈیڈیکائنڈ دا ”ثبوت“ \pnbLatin{(Dedekind's ``Proof'')}]{لامتناہی سیٹ دے وجود دا ڈیڈیکائنڈ دا
”ثبوت“ (Dedekind's ``Proof'' of the Existence of an Infinite Set)}

ایس باب وچ اسی ڈیڈیکائنڈ الجبریاں دے ذریعے قدرتی عدداں دا سیٹ-تھیوری
والا بیان \pnbLatin{(set-theoretic treatment)} دِتا اے۔ \olref[arith][ref]{sec}
وچ اسی، ہور گلاں دے نال، تجزیے دی اعدادی تشکیل \pnbLatin{(arithmetization of analysis)} دی فلسفیانہ اہمیت اُتے غور کیتا سی۔ ہُن ایتھے حاصل کیتے
نتیجے دی اہمیت جانچنی چاہیدی اے۔

پورے \olref[sfr][arith][]{chap} وچ قدرتی عدد پہلے توں دِتے ہوئے من کے
اوہناں نال صحیح، ناطق تے حقیقی عدد صاف طور تے بنائے گئے سن۔ ایس باب
وچ قدرتی عدداں دی کوئی صریح بناوٹ \pnbLatin{(explicit construction)} نہیں دِتی؛
صرف ایہ ثابت کیتا اے کہ \emph{کوئی وی ڈیڈیکائنڈ لامتناہی سیٹ دِتا
ہووے}، تاں اک ایسا سیٹ تعریف ہو سکدا اے جیہڑا اوویں ورتاؤ کرے جیویں
اسی چاندے آں کہ~$\Nat$ ورتاؤ کرے۔

ایس کرکے اسی کسے مابعدالطبیعیاتی (metaphysical) سوال، مثلاً
\emph{قدرتی عدد اصل وچ کیہڑیاں شےواں نیں}، دا جواب دین دا دعویٰ نہیں
کر سکدے---تے ایہ چنگی گل اے۔ \olref[sfr][arith][ref]{sec} وچ زور دِتا
گیا سی کہ $\Real$ نوں ڈیڈیکائنڈ کٹاں دا سیٹ تعریف کرنا سہولت والی
مقرری اے، \emph{دریافت} نہیں۔ اصل نکتہ ایہ اے کہ ڈیڈیکائنڈ کٹ حقیقی
عدداں دیاں اہم ریاضیاتی خاصیتاں دی مثال پیش کردے نیں۔ ایسے طرح ایتھے
\emph{ہر} ڈیڈیکائنڈ الجبرا قدرتی عدداں دیاں اہم ریاضیاتی خاصیتاں دی
مثال پیش کردا اے۔ (ڈیڈیکائنڈ نے ایہ ثابت کرکے گل ہور پکی کیتی کہ سارے
ڈیڈیکائنڈ الجبرے \emph{ہم شکل (isomorphic)} نیں؛ متشاکل اک ہور
دستاویزی متبادل اے (\citeyear[Theorems
132--3]{Dedekind1888})۔ ایس لئی کئی اجوکے ”ساختیت پسند
(structuralists)“ ڈیڈیکائنڈ نوں اپنا پیش رو آکھدے نیں۔)

نال ای، اسی وکھایا اے کہ قدرتی عدداں دے نظریے نوں اک سادی، ناپختی
سیٹ تھیوری \pnbLatin{(naive simple set theory)} وچ کیویں سمایا جا سکدا اے۔ ایہ
سیٹ تھیوری آپ ہجے خاصی غیر رسمی اے، پر بظاہر قدرتی عدداں نوں پہلے توں
دِتا ہویا نہیں مندی۔ ایس حد تک شاید اسی ڈیڈیکائنڈ دے ولولے بھرے
منطقیت والے منصوبے \pnbLatin{(logicist project)} نوں پورا کرن دے راہ اُتے آں،
جیہنوں اوہنے ایویں بیان کیتا:
\begin{quote}
	سائنس وچ جیہڑی گل ثبوت دے قابل اے، اوہنوں ثبوت توں بناں منّنا نہیں
	چاہیدا۔ بھاویں ایہ تقاضا معقول اے، میں ایہ نہیں من سکدا کہ سادہ
	ترین سائنس دیاں بنیاداں رکھن دے سب توں نویں طریقیاں وچ وی ایہ پورا
	ہویا اے؛ یعنی منطق دے اوس حصے وچ جیہڑا عدداں دے نظریے نال واسطہ
	رکھدا اے۔ حساب (الجبرا، تجزیہ) نوں محض منطق دا اک حصہ آکھ کے میرا
	مطلب ایہ اے کہ عدد دا تصور تھاں تے سمے دے تصورات یا وجدان توں پوری
	طرح آزاد اے---بلکہ اوہ فکر دے خالص قانوناں \pnbLatin{(pure laws of thought)}
	دی فوری پیداوار اے۔
	\citep[preface]{Dedekind1888}
\end{quote}
ڈیڈیکائنڈ دا دلیر خیال ایہ اے: ہُن اسی صرف (سادی) سیٹ تھیوری نال
قدرتی عدد بناؤن دا طریقہ وکھایا اے، تے \olref[sfr][arith][]{chap} وچ
ویکھیا سی کہ قدرتی عدد تے کجھ سیٹ تھیوری دِتے ہوون، تاں حقیقی عدد
کیویں بنائے جا سکدے نیں۔ سو شاید ”حساب (الجبرا، تجزیہ)“، ڈیڈیکائنڈ
دے وسیع معنٰی وچ، ”محض منطق دا اک حصہ“ نکلے۔

خیال ایہو اے، پر اک پل رکو۔ ڈیڈیکائنڈ الجبرے، یعنی قدرتی عدداں دے
ساڈے قائم مقام (surrogate)، دی بناوٹ ڈیڈیکائنڈ لامتناہی سیٹ دے وجود
اُتے مشروط (conditional) اے۔ (پچھے
\olref[sfr][infinite][dedekind]{thm:DedekindInfiniteAlgebra} ویکھو۔)
جے ایسے سیٹ دا وجود ”منطق“ یا ”فکر دے خالص قانوناں“ راہیں قائم نہ ہو
سکے، تاں منصوبہ ایتھے ای رک جاندا اے۔

سو سوال ایہ اے: کی ڈیڈیکائنڈ لامتناہی سیٹ دا وجود ”فکر دے خالص
قانوناں \pnbLatin{(pure laws of thought)}“ نال \emph{قائم} کیتا جا سکدا اے؟
ڈیڈیکائنڈ دی کوشش ایہ سی:
\begin{quote}
  میرے خیالاں دی اپنی قلمرو، یعنی اوہناں ساریاں شےواں دی کُلّیت $S$
  جیہڑیاں میرے خیال دا موضوع بن سکدیاں نیں، لامتناہی اے۔ کیوں جو جے
  $s$،~$S$ دے اک رکن نوں ظاہر کردا اے، تاں ایہ خیال $s'$ کہ~$s$ میرے
  خیال دا موضوع بن سکدا اے، آپ~$S$ دا اک رکن اے۔ جے ایہنوں عکس
  $\phi(s)$، یعنی رکن~$s$ دا عکس، سمجھیا جاوے، تاں \dots~$S$ [ڈیڈیکائنڈ]
  لامتناہی اے، تے ایہی ثابت کرنا سی۔
	\citep[\S66]{Dedekind1888}
\end{quote}
اک ایسی کتاب دے عین وچکار، جیہڑی زیادہ تر نہایت کڑے ریاضیاتی ثبوتاں
اُتے مشتمل اے، ایہ دلیل ملنی خاصی حیران کن اے۔ ایتھے دو گلّاں قابلِ
ذکر نیں۔

پہلی گل: ایس ”ثبوت“ وچ اوہ کردار شاید ہی اے جیہنوں اسی اج
”ریاضیاتی“ منّاں گے۔ ایہ نفسیاتی شےواں \pnbLatin{(psychological objects)}، یعنی
خیالاں، دی گل کردا اے، تے اوہ وی محض \emph{ممکن} خیالاں دی۔

دوجی گل: ڈیڈیکائنڈ الجبرے دی ساڈی پیشکش مطابق ایس ”ثبوت“ وچ اک صاف
فنی کمی اے۔ جے ڈیڈیکائنڈ دی دلیل کامیاب وی ہووے، تاں اوہ صرف اینا
قائم کردی اے کہ شےواں لامتناہی گنتی وچ نیں، خاص طور تے خیال۔ پر اوہنوں
ایہ وجہ وی دینی پینی اے کہ~$S$ نوں اک اکیلا \emph{سیٹ} سمجھیا جاوے
جیہدے لامتناہی !!{element}s نیں، نہ کہ~$S$ نوں بس \emph{کجھ شےواں}
سمجھیا جاوے۔ ایہو بہتیاں شےواں توں اک سیٹ تک دا خلا
\pnbLatin{(plurality-to-set gap)} اے۔

ڈیڈیکائنڈ نے ایتھے خلا نہ ویکھیا؛ ایس توں شاید پتا لگدا اے کہ اوہدا
لفظ ”کُلّیت (totality)“ عین \emph{ساڈے} لفظ ”سیٹ“ دے نال نہیں رلدا۔
\footnote{اصل وچ ساڈے کول ایہ سوچن دیاں ہور وجوہ وی نیں کہ ایہ نہیں
رلدا سی؛ \citet[p.~23]{Potter2004} ویکھو۔} ایہ کوئی وڈی حیرانی نہیں۔
پچھلے دو باباں دا منصوبہ---قدرتی عدداں دی اک ”بناوٹ“، تے اوہناں توں
صحیح، حقیقی تے ناطق عدداں دی اک ”بناوٹ“---سارا سادے (naive) ڈھنگ نال
چلایا گیا اے۔ جدوں لوڑ پئی، اسی ایہ سیٹ یا اوہ سیٹ اپنے کم لئی لے
لیا، پر عام اصول بہت گھٹ دِتے کہ عین کیہڑے سیٹ کدوں موجود ہوندے نیں۔
سانوں \emph{کجھ} عام اصول ضرور چاہیدے نیں، نہیں تاں اسی رَسل دے تضاد
\pnbLatin{(Russell's Paradox)} وچ جا ڈِگاں گے۔

ہُن ایس سادہ لوح سیٹ تھیوری توں اگے ودھن دا ویلا آ گیا اے۔
\endgroup
% /OLP-0053
% OLP-0054 content/sets-functions-relations/infinite/card-sb.tex
\begingroup
\def\olfilename{card-sb}
\expandafter\def\csname ol@sectioncs\endcsname{section}

\olfileid{sfr}{infinite}{card-sb}

\olsection{ضمیمہ: شروڈر-برنسٹین \pnbLatin{(Schröder--Bernstein)} دا ثبوت}

سادی (naive) سیٹ تھیوری توں وداع ہون توں پہلاں اک آخری سادہ، پر خاصا
نفیس، ثبوت ویکھنا باقی اے۔ ایہ شروڈر-برنسٹین
(\olref[sfr][siz][sb]{thm:schroder-bernstein}) دا ثبوت اے: جے
$\cardle{A}{B}$ تے $\cardle{B}{A}$، تاں $\cardeq{A}{B}$۔ یعنی
!!{injection}s \pnbLatin{(injective maps)} $f \colon A \to B$ تے
$g \colon B \to A$ دِتے ہوون، تاں !!a{bijection} \pnbLatin{(bijective map)}
$h \colon A \to B$ موجود اے۔

ایس باب وچ اسی ڈیڈیکائنڈ دیاں \emph{بندشاں (closures)} دا تصور ورتیا
اے۔ ڈیڈیکائنڈ نے ایسے تصور نال شروڈر-برنسٹین دا اک سوہنا ثبوت دِتا،
تے اسی اوہ ایتھے پیش کراں گے۔ ذرا وکھرے پر بنیادی طور تے رلدے ملدے
بیان لئی ایہ ثبوت \citet[pp.~157--8]{Potter2004} دے بہت نیڑے چلدا اے۔
تھوڑی آن لائن تلاش وی وکھا دیندی اے کہ $\sqrt{2}$ دی غیر ناطقیت وانگ
ایس قضیے دے \emph{کئی} دلچسپ تے وکھرے ثبوت موجود نیں۔

\olref[sfr][infinite][dedekind]{Closure} ورگی علامت بندی ورتو۔ حامل
سیٹ \pnbLatin{(carrier set)} C اُتے خود نقشے (self-map/endofunction) f تے اوس
حامل دے ذیلی سیٹ B لئی رکھو
\[
\Closureofunder{f}{B} = \bigcap \Setabs{X}{B \subseteq X
\text{ تے $X$، $f$-بند اے}}
\]
ہر ایسے سیٹ $B$ تے فنکشن~$f$ لئی۔ ایس طرح
$\Closureofunder{f}{B}$، B نال پیدا کیتی سب توں چھوٹی بندش \pnbLatin{(least generated closure)}، یعنی سب توں چھوٹا $f$-بند سیٹ اے جیہدے وچ~$B$
شامل اے، ایس معنٰی وچ کہ:
% PNBINF-004 ماخذ دی درستی: f نوں حامل C دا خود نقشہ، B نوں C دا ذیلی سیٹ تے X نوں حامل دے ذیلی سیٹاں وچ چلدا مقرر کیتا اے؛ ماخذ دی بے قید ”ہر فنکشن“ والی عبارت بندش دے وجود تے اگلے اطلاق لئی لازم حامل دائرہ نہیں دیندی۔

\begin{lem}\ollabel{Closureprops}
	ہر ایسے حامل خود نقشے $f$ تے ہر حامل ذیلی سیٹ $B$ لئی:
	\begin{enumerate}
		\item\ollabel{Closurehaselem} $B \subseteq \Closureofunder{f}{B}$، یعنی بندش وچ اوہ ذیلی سیٹ شامل اے؛ تے
		\item\ollabel{Closureclosed} $\Closureofunder{f}{B}$، $f$-بند اے؛ تے
		\item\ollabel{Closuresmallest} جے $X$، $f$-بند اے تے $B
		\subseteq X$، تاں لازماً $\Closureofunder{f}{B} \subseteq X$۔
	\end{enumerate}
\end{lem}

\begin{proof}
ثبوت عین \olref[sfr][infinite][dedekind]{closureproperties} والے ثبوت
ورگا اے۔
\end{proof}

شروڈر-برنسٹین ثابت کرن لئی ہُن صرف اک آخری معاون حقیقت چاہیدی اے:

\begin{prop}\ollabel{sbhelper}
درمیانے سیٹ والے معاون نتیجے \pnbLatin{(sandwich helper)} مطابق، جے
$A \subseteq B \subseteq C$ تے $A \approx C$، تاں $\cardeq{\cardeq{A}{B}}{C}$۔
\end{prop}

\begin{proof}
!!a{bijection} \pnbLatin{(bijective map)} $f \colon C \to A$ دِتا ہووے، تاں
$F = \Closureofunder{f}{C \setminus B}$ لو۔ ہُن فنکشن $g$، جیہدا
دائرۂ تعریف $C$ اے، دو صورتاں (piecewise) وچ ایویں تعریف کرو:
\[
	g(x) =
	\begin{cases}
			f(x) &\text{جے $x \in F$}\\
			x & \text{نہیں تاں}
	\end{cases}
\]
اسی ثابت کراں گے کہ $g$، نقشہ بندی $C \to B$ دے طور تے
!!a{bijection} اے۔ ایس توں $\comp{f^{-1}}{g} \colon A \to B$ دا
!!a{bijection} ہونا نکلے گا تے ثبوت مکمل ہو جاوے گا۔

پہلا دعویٰ ایہ اے کہ جے $x \in F$ پر $y\notin F$، تاں $g(x) \neq
g(y)$۔ برہانِ خلف (reductio) لئی الٹ منو، سو $y = g(y) = g(x) =
f(x)$۔ چونکہ $x \in F$ تے \olref{Closureprops} مطابق $F$، $f$-بند
اے، ایس لئی $y = f(x) \in F$؛ ایہ تضاد اے۔

ہُن $g(x) = g(y)$ فرض کرو۔ اُتلے دعوے نال $x \in F$ اُدوں تے صرف
اُدوں اے جدوں $y \in F$۔ جے $x, y \in F$، تاں $f(x) = g (x) =
g(y) = f(y)$، سو $x = y$، کیوں جو $f$ اک !!a{bijection} اے۔ جے
$x, y \notin F$، تاں $x = g(x) = g(y) = y$۔ لہٰذا $g$ اک
!!a{injection} \pnbLatin{(injective map)} اے۔

ہُن رینج (range) دی برابری $\ran{g} = B$ ثابت کرنی باقی اے۔ کوئی
$x \in B \subseteq C$ پکا کرو۔ جے $x \notin F$، تاں $g(x) = x$۔
جے $x \in F$، تاں $x = f(y)$، کسے $y \in F$ لئی؛ نہیں تاں $F
\setminus \{x\}$، $f$-بند رہندا تے $C\setminus B$ نوں شامل رکھدا،
جیہڑا \olref{Closureprops} مطابق ناممکن اے۔ ہُن $g(y) = f(y) = x$۔
\end{proof}

آخر وچ مرکزی نتیجہ ثابت کردے آں۔ فنکشن دے عکس سیٹ \pnbLatin{(image set)} دی
تعریف یاد کرو: فنکشن $h$ تے سیٹ $D$ دِتے ہوون، تاں اسی
$\funimage{h}{D} = \Setabs{h(x)}{x \in D}$ تعریف کردے آں۔

\begin{proof}[شروڈر-برنسٹین دا ثبوت] $f \colon A \to B$ تے
$g \colon B \to A$ نوں !!{injection}s منو۔ چونکہ $\funimage{f}{A}
\subseteq B$، ایس لئی عکس سیٹاں دی زنجیر
$\funimage{g}{\funimage{f}{A}} \subseteq g[B] \subseteq A$ اے۔ نال
ای $\comp{f}{g} \colon A \to \funimage{g}{\funimage{f}{A}}$ اک
!!{injection} اے، کیوں جو $g$ تے $f$ دونیں انجیکشن نیں؛ بلکہ ہدف سیٹ
دی تعریف دے باعث $\comp{f}{g}$ اک !!a{bijection} اے۔ سو
$\cardeq{\funimage{g}{\funimage{f}{A}}}{A}$، تے \olref{sbhelper}
مطابق !!a{bijection} $h \colon A \to \funimage{g}{B}$ موجود اے۔ ہور
$g^{-1}$ اک !!a{bijection} $\funimage{g}{B} \to B$ اے۔ آخر
$\comp{h}{g^{-1}} \colon A \to B$ اک !!a{bijection} اے۔
\end{proof}
\endgroup
% /OLP-0054
\endgroup
% /OLP-0049

\OLEndPartHook


\par\noindent\emph{متبادل ترتیب۔} سیٹ، تعلق تے فنکشن دا مُڈھلا حصہ (\texttt{sets-functions-relations.tex})۔ ایہ ترتیب تدوین جوگ ماخذ فائل وچ موجود اے۔
\endgroup
% /OLP-0003

% OLP-0055 content/propositional-logic/propositional-logic.tex
\begingroup
\def\olfilename{propositional-logic}
\expandafter\def\csname ol@sectioncs\endcsname{section}

\olpart{pl}{قضیاتی منطق \pnbLatin{(Propositional Logic)}}

\begin{editorial}
  ایس حصے وچ کلاسیکی قضیاتی منطق \pnbLatin{(classical propositional logic)} دا
  مواد شامل اے۔ پہلا باب نسبتاً
  ابتدائی اے تے صرف تعریفاں تے نتیجے درج کردا اے؛ کئی ثبوت پورے نہیں
  دِتے گئے بلکہ مشقاں لئی چھڈّے گئے نیں۔ ثبوتی نظاماں \pnbLatin{(proof systems)}
  تے تمامیت دے قضیے \pnbLatin{(completeness theorem)} دا مواد پہلے درجے دی منطق
  \pnbLatin{(first-order logic)} والے حصے توں شامل کیتا گیا اے، تے
  ``FOL'' ٹیگ نوں false مقرر کیتا گیا اے۔ ایس طرح محمولاں، اصطلاحاں
  تے مسوّراں، یعنی predicates، terms تے quantifiers، نال متعلق ہر گل
  شامل نہیں رہندی، تے !!{structure}s~$\Struct{M}$ یعنی ساختاں دی گل
  نوں !!{valuation}s~$\pAssign{v}$ یعنی قدر بندیاں (valuations) دی گل
  نال بدل دِتا جاندا اے۔

  منصوبہ اے کہ ایس حصے نوں ہور تفصیل نال ودھایا جاوے، تے ہور موضوع
  تے نتیجے، مثلاً صداقتی-فنکشنی تمامیت \pnbLatin{(truth-functional completeness)}، شامل کیتے جان۔
\end{editorial}

% OLP-0056 content/propositional-logic/syntax-and-semantics/syntax-and-semantics.tex
\begingroup
\def\olfilename{syntax-and-semantics}
\expandafter\def\csname ol@sectioncs\endcsname{section}

\olchapter{pl}{syn}{نحو تے معنیات \pnbLatin{(Syntax and Semantics)}}

\begin{editorial}
  ایہ صرف تعریفاں دا نہایت مختصر خلاصہ \pnbLatin{(very quick summary)} اے۔ ایہنوں
  ودھایا جانا چاہیدا اے تاں جو فارمولیاں اُتے استقرا \pnbLatin{(induction on formulas)} نال ثبوتاں دا آسان تعارف \pnbLatin{(gentle introduction)}، تے بہت
  ساریاں ہور مثالاں، دِتیاں جا سکن۔
\end{editorial}

% OLP-0057 content/propositional-logic/syntax-and-semantics/introduction.tex
\begingroup
\def\olfilename{introduction}
\expandafter\def\csname ol@sectioncs\endcsname{section}

\olfileid{pl}{syn}{int}

\olsection{تعارف \pnbLatin{(Introduction)}}

قضیاتی منطق \pnbLatin{(propositional logic)} اوہناں !!{formula}s نال واسطہ رکھدی
اے جیہڑے !!{propositional variable}s، یعنی قضیاتی متغیراں
\pnbLatin{(propositional variables)}، توں قضیاتی رابطیاں \pnbLatin{(propositional connectives)}
$\lnot$، $\land$، $\lor$، $\lif$ تے $\liff$ نال بنائے جاندے نیں۔ سہج
فہم طور تے، !!a{propositional variable}~$p$ اک جملے یا قضیے دی
نمائندگی کردا اے جیہڑا سچ یا جھوٹ ہوندا اے۔ جدوں وی !!a{formula} وچ
!!{propositional variable} دی ``صداقتی قدر \pnbLatin{(truth value)}'' متعین ہو
جاوے، تاں
قضیاتی رابطیاں نال اوہناں توں بنے !!{formula}s وچوں ہر اک دی صداقتی قدر وی متعین
ہو جاندی اے۔ اسیں آکھدے آں کہ قضیاتی منطق
\emph{صداقتی-فنکشنی \pnbLatin{(truth functional)}} اے، کیوں جے ایہدی معنیات
صداقتی قدراں دے فنکشناں
راہیں دِتی جاندی اے۔ خاص طور تے، قضیاتی منطق وچ اسیں صدق تے کذب دی
ہور کوئی تعیین زیرِ غور نہیں لیاؤندے، مثلاً ایہ کہ کوئی گل محض اتفاقی
طور تے سچ ہون دی بجائے لازماً سچ اے، یا ایہ کہ کوئی گل سچ معلوم اے، یا
ایہ کہ کوئی گل ہُن سچ اے نہ کہ پہلے سچ سی یا اگوں سچ ہووے گی۔ اسیں صرف
دو صداقتی قدراں، سچ ($\True$) تے جھوٹ ($\False$)، اُتے غور کردے آں، تے
ایس لئی ایہ امکان بحث توں باہر رکھدے آں کہ کوئی بیان نہ سچ ہووے نہ
جھوٹ، یا صرف ادھا سچ ہووے۔ اسیں صرف اوہناں رابطیاں اُتے وی دھیان
دیندے آں جنہاں توں بنے !!a{formula} دی صداقتی قدر اوہدے حصیاں دیاں
صداقتی قدراں نال پوری طرح متعین ہوندی اے (نہ کہ، مثلاً، اوہدے معنی
نال)۔ خاص طور تے، انگریزی شرطیہ جملیاں دی صداقتی قدر ایس معنی وچ
صداقتی-فنکشنی اے یا نہیں، ایہ اختلافی گل اے۔ مادی شرطیہ~$\lif$، یعنی
material conditional، صداقتی-فنکشنی اے؛ ہور منطقاں اوہناں شرطیاں نال
واسطہ رکھدیاں نیں
جیہڑیاں صداقتی-فنکشنی نہیں ہُندیاں۔

صداقتی-فنکشنی قضیاتی منطق دا نظریہ (theory) تے مابعد نظریہ
(metatheory) تیار کرن لئی، سانوں پہلاں ایہدیاں عبارتاں دی نحو (syntax)
تے معنیات (semantics) دی تعریف کرنی پئے گی۔ اسیں
رابطیاں نال !!{propositional variable}s توں !!{formula}s بناؤن دا اک
طریقہ بیان کراں گے۔ ہور تعریفاں وی ممکن نیں۔ ہور نظام وکھریاں علامات
چُنن گے، رابطیاں دے وکھرے مجموعیاں نوں ابتدائی منّن گے، تے قوسین نوں
وکھرے ڈھنگ نال ورتن گے (یا نام نہاد پولش علامت نگاری وانگ بالکل نہیں
ورتن گے)۔ پر سارے طریقیاں وچ سانجھی گل ایہ اے کہ تشکیل دے قاعدے
\pnbLatin{(formation rules)}
!!{formula}s دے سیٹ دی
\emph{استقرائی طور تے (inductively)} تعریف کردے نیں۔ جے ایہ کم ٹھیک
طرح کیتا جاوے،
تاں ہر عبارت بنیادی طور تے تشکیل دے قاعدیاں مطابق صرف اکّو طریقے نال
بن سکدی اے۔ فارمولیاں نوں
\emph{اکّو طرح نال پڑھن جوگ \pnbLatin{(uniquely readable)}} بناؤن والی استقرائی
تعریف دا مطلب اے کہ
اسیں ایسّے طریقے، یعنی استقرائی تعریف، نال ایہناں عبارتاں دے معنی وی
دے سکدے آں۔

عبارتاں نوں معنی دینا معنیات (semantics) دا دائرہ اے۔ قضیاتی منطق دی
معنیات وچ مرکزی تصور !!a{valuation}، یعنی قدر بندی (valuation)، وچ
پورا اترن (satisfaction) دا اے۔
!!^a{valuation}~$\pAssign{v}$، !!{propositional variable}s نوں صداقتی
قدراں $\True$، $\False$ مقرر کردی اے۔ ہر !!{valuation} اک صداقتی قدر
$\pValue{v}(!A)$، کسی وی !!{formula}~$!A$ لئی، متعین کردی اے۔
!!^a{formula}، !!a{valuation}~$\pAssign{v}$ وچ اُدوں تے صرف اُدوں پورا
اتردا اے جدوں $\pValue{v}(!A) = \True$ ہووے—اسیں ایہنوں
$\pSat{v}{!A}$ لکھدے آں۔ ایس تعلق دی تعریف $!A$ دی ساخت اُتے استقرا
نال وی کیتی جا سکدی اے؛ منطقی رابطیاں دے صداقتی فنکشن ورت کے، مثلاً،
$!A \land !B$ دے پورا اترن نوں $!A$ تے~$!B$ دے پورا اترن (یا نہ
اترن) دے لحاظ نال تعریف کیتا جا سکدا اے۔

جملیاں لئی پورا اترن دے تعلق $\pSat{v}{!A}$ دی بنیاد اُتے اسیں فیر
ہمیشہ صادق ہون (tautology)، منطقی نتیجے (entailment) تے پورا ہون جوگتا
(satisfiability) دے بنیادی معنوی تصوراں
دی تعریف کر سکدے آں۔ !!^a{formula} اک ہمیشہ صادق فارمولا،
$\Entails !A$، اے جے ہر !!{valuation} اوہنوں پورا کردی اے، یعنی
$\pValue{v}(!A) = \True$ ہر~$\pAssign{v}$ لئی۔ اوہی فارمولا
!!{formula}s دے سیٹ دا منطقی نتیجہ ہوندا اے، $\Gamma \Entails !A$، جے ہر اوہ
!!{valuation} جیہڑی~$\Gamma$ دے سارے !!{formula}s نوں پورا کردی اے،
اوہ~$!A$ نوں وی پورا کردی اے۔ تے !!{formula}s دا سیٹ پورا ہون جوگ اے
جے کوئی !!{valuation} اوس وچلے سارے !!{formula}s نوں اکّو ویلے پورا
کردی اے۔ کیوں جے !!{formula}s دی تعریف استقرائی طور تے کیتی گئی اے،
تے پورا اترن دی تعریف اگوں !!{formula}s دی ساخت اُتے استقرا نال کیتی
گئی اے، ایس لئی اسیں اپنی معنیات دیاں خاصیتاں ثابت کرن تے تعریف کیتے
معنوی تصوراں نوں آپس وچ جوڑن لئی استقرا ورت سکدے آں۔
\endgroup
% /OLP-0057

% OLP-0058 content/propositional-logic/syntax-and-semantics/formulas.tex
\begingroup
\def\olfilename{formulas}
\expandafter\def\csname ol@sectioncs\endcsname{section}

\olfileid{pl}{syn}{fml}

\olsection{قضیاتی \usetoken{P}{formula} \pnbLatin{(Propositional Formula)}}

قضیاتی منطق \pnbLatin{(propositional logic)} دے !!^{formula}s
\emph{!!{propositional variable}s}، یعنی قضیاتی متغیراں \pnbLatin{(propositional variables)}، توں
\iftag{prvFalse}{\iftag{prvTrue}{،}{ تے} قضیاتی مستقل~$\lfalse$}{}
\iftag{prvTrue}{\iftag{prvFalse}{ تے}{} قضیاتی مستقل~$\ltrue$}{}
\emph{منطقی رابطیاں \pnbLatin{(logical connectives)}} نال بنائے جاندے نیں۔

\begin{enumerate}
\item !!^a{denumerable} سیٹ \pnbLatin{(denumerable set)}~$\PVar$، جیہڑا
  !!{propositional variable}s
  $\Obj p_0$، $\Obj p_1$، \dots دا اے۔
\tagitem{prvFalse}{!!{falsity} دا قضیاتی مستقل \pnbLatin{(propositional constant)}~$\lfalse$۔}{}
\tagitem{prvTrue}{!!{truth} دا قضیاتی مستقل \pnbLatin{(propositional constant)}~$\ltrue$۔}{}
\item منطقی رابطے \pnbLatin{(logical connectives)}:
  \startycommalist
  \iftag{prvNot}{\ycomma $\lnot$ (نفی؛ negation)}{}%
  \iftag{prvAnd}{\ycomma $\land$ (عطف؛ conjunction)}{}%
  \iftag{prvOr}{\ycomma $\lor$ (فصل؛ disjunction)}{}%
  \iftag{prvIf}{\ycomma $\lif$ (شرطیہ؛ conditional، !!{conditional})}{}%
  \iftag{prvIff}{\ycomma $\liff$ (دو شرطیہ؛ biconditional،
    !!{biconditional})}{}%
\item رموزِ اوقاف \pnbLatin{(punctuation marks)}: (، )، تے کوما۔
\end{enumerate}

اسیں قضیاتی منطق دی ایس زبان (language) نوں $\Lang L_0$ نال ظاہر کردے
آں۔

\iftag{defNot,defOr,defAnd,defIf,defIff,defTrue,defFalse,defEx,defAll}{%
اُتے متعارف کرائے ابتدائی رابطیاں \pnbLatin{(primitive connectives)} دے نال نال،
اسیں ایہ \emph{تعریف شدہ (defined)} علامات وی ورتدے آں:
\startycommalist
  \iftag{defNot}{\ycomma $\lnot$ (نفی)}{}%
  \iftag{defAnd}{\ycomma $\land$ (عطف)}{}%
  \iftag{defOr}{\ycomma $\lor$ (فصل)}{}%
  \iftag{defIf}{\ycomma $\lif$ (شرطیہ، !!{conditional})}{}%
  \iftag{defIff}{\ycomma $\liff$ (دو شرطیہ، !!{biconditional})}{}%
  \iftag{defFalse}{\ycomma $\lfalse$ (کذب، !!{falsity})}{}%
% PNBPL-001 ماخذ دی درستی: اندرلے defTrue شرطی حکم دی خالی تیسری شاخ ماخذ وچ رہ گئی اے؛ ایتھے اوہ عین خالی شاخ شامل اے۔
  \iftag{defTrue}{\ycomma $\ltrue$ (صدق، !!{truth})}{}}{}۔

\begin{tagblock}{defNot,defOr,defAnd,defIf,defIff,defTrue,defFalse,defEx,defAll}
\begin{explain}
تعریف شدہ علامت \pnbLatin{(defined symbol)} باقاعدہ طور تے زبان دا حصہ \pnbLatin{(official part of the language)} نہیں ہُندی، بلکہ اک غیر رسمی مخفف \pnbLatin{(informal abbreviation)} دے طور تے متعارف کرائی جاندی اے: ایہ سانوں اوہناں فارمولیاں نوں
مختصر لکھن دی سہولت دیندی اے جیہڑے صرف ابتدائی علامات نال لکھے جان تاں
خاصے لمّے ہو جان۔ ایہ صاف طور تے اک فائدہ اے۔ پر ایہ توں وڈا فائدہ ایہ
اے کہ ثبوت چھوٹے ہو جاندے نیں۔ جے کوئی علامت ابتدائی (primitive) ہووے،
تاں ثبوتاں
وچ اوہنوں وکھرے طور تے ورتنا پیندا اے۔ ایس لئی جِنّیاں زیادہ ابتدائی
علامتاں ہون، ساڈے ثبوت اونّے ای لمّے ہون گے۔
\end{explain}
\end{tagblock}

% Alternate symbols

\begin{intro}
ہو سکدا اے تُسی اُتّے ساڈی ورتی اصطلاح (terminology) تے علامات
(symbols) توں وکھری اصطلاح تے علامتاں نال واقف ہووو۔ منطق دیاں کتاباں
(تے استاد) ``نفی (negation)'' لئی عام طور تے $\sim$، $\neg$ تے~!\
ورتدے نیں، تے ``عطف (conjunction)'' لئی $\wedge$، $\cdot$ تے $\&$۔
``شرطیہ (conditional)'' یا ``لزوم (implication)'' لئی عام ورتیاں جان والیاں
علامتاں $\rightarrow$، $\Rightarrow$ تے $\supset$ نیں۔
\iftag{prvIff,defIff}{``دو شرطیہ''، ``دو طرفہ لزوم'' یا ``(مادی) ہم
  ارزی'' دیاں علامات $\leftrightarrow$، $\Leftrightarrow$ تے
  $\equiv$ نیں۔}{}
\iftag{prvFalse,defFalse}{$\lfalse$ دی علامت نوں وکھ وکھ تھاواں اُتے
  ``کذب''، ``کاذب''، ``محال'' یا ``زیریں'' آکھیا جاندا اے۔}{}
\iftag{prvTrue,defTrue}{$\ltrue$ دی علامت نوں وکھ وکھ تھاواں اُتے
  ``صدق''، ``صادق'' یا ``بالائی'' آکھیا جاندا اے۔}{}
\end{intro}

\begin{defn}[فارمولا (Formula)]
\ollabel{defn:formulas}
قضیاتی منطق دے \emph{!!{formula}s} دا سیٹ~$\Frm[L_0]$ استقرائی طور تے
(inductively) اینج تعریف کیتا جاندا اے:
\begin{enumerate}
\tagitem{prvFalse}{$\lfalse$ اک ایٹمی (atomic) !!{formula} اے۔}{}

\tagitem{prvTrue}{$\ltrue$ اک ایٹمی (atomic) !!{formula} اے۔}{}

\item ہر !!{propositional variable}~$\Obj p_i$، یعنی ہر قضیاتی متغیر
  \pnbLatin{(propositional variable)}، اک ایٹمی !!{formula} اے۔

\tagitem{prvNot}{جے $!A$، !!a{formula} اے، تاں نفی والی تشکیل
  \pnbLatin{(negation formation)} $\lnot !A$ وی
  !!a{formula} اے۔}{}

\tagitem{prvAnd}{جے $!A$ تے $!B$، !!{formula}s نیں، تاں عطف والی
  تشکیل \pnbLatin{(conjunction formation)} $(!A \land
  !B)$ وی !!a{formula} اے۔}{}

\tagitem{prvOr}{جے $!A$ تے $!B$، !!{formula}s نیں، تاں فصل والی
  تشکیل \pnbLatin{(disjunction formation)} $(!A \lor !B)$
  وی !!a{formula} اے۔}{}

\tagitem{prvIf}{جے $!A$ تے $!B$، !!{formula}s نیں، تاں شرطیہ والی
  تشکیل \pnbLatin{(conditional formation)} $(!A \lif !B)$
  وی !!a{formula} اے۔}{}

\tagitem{prvIff}{جے $!A$ تے $!B$، !!{formula}s نیں، تاں دو شرطیہ والی
  تشکیل \pnbLatin{(biconditional formation)} $(!A \liff !B)$
  وی !!a{formula} اے۔}{}

\tagitem{limitClause}{حد والی شق \pnbLatin{(limiting clause)}: ہور کوئی شے
  !!a{formula} نہیں اے۔}{}
\end{enumerate}
\end{defn}

\begin{explain}
!!{formula}s دی تعریف اک \emph{استقرائی تعریف \pnbLatin{(inductive definition)}}
اے۔ بنیادی طور تے اسیں !!{formula}s دا سیٹ لامتناہی مرحلیاں \pnbLatin{(infinitely many stages)} وچ بناؤندے آں۔ پہلے مرحلے \pnbLatin{(initial stage)} وچ اسیں سارے
ایٹمی
فارمولیاں نوں فارمولے قرار دیندے آں؛ ایہ تعریف دیاں پہلیاں کجھ صورتاں،
یعنی
\iftag{prvTrue}{$\ltrue$، }{}%
\iftag{prvFalse}{$\lfalse$، }{}%
$\Obj p_i$ والیاں صورتاں دے مطابق اے۔ ایس طرح ``ایٹمی (atomic)
!!{formula}''
دا مطلب ایس روپ دا کوئی وی !!{formula} اے۔

تعریف دیاں ہور صورتاں پہلے بنائے !!{formula}s توں نویں !!{formula}s
بناؤن دے تشکیل قاعدے \pnbLatin{(formation rules)} دیندیاں نیں۔ دوجے مرحلے \pnbLatin{(second stage)} وچ اسیں ایہناں قاعدیاں نال ایٹمی
!!{formula}s توں !!{formula}s بنا سکدے آں۔ تیجے مرحلے \pnbLatin{(third stage)}
وچ اسیں ایٹمی
فارمولیاں تے دوجے مرحلے وچ ملے فارمولیاں توں نویں فارمولے بناؤندے آں،
تے ایویں ای اگے۔ !!{formula} اوہو شے اے جیہڑی آخرکار کسے ایہو جہے
مرحلے اُتے بنائی جاوے، تے ہور کوئی شے نہیں۔
\end{explain}

فارمولا $(!B \ast !C)$، جیہڑا $!B$، $!C$ توں دو-تھاں رابطے
\pnbLatin{(two-place connective)}~$\ast$ نال بنیا اے، لکھدے ہویاں، اسیں اکثر سب
توں باہری قوسین دی جوڑی \pnbLatin{(outermost pair of parentheses)} چھڈ کے صرف
$!B \ast !C$ لکھاں گے۔

\begin{tagblock}{defNot,defOr,defAnd,defIf,defIff,defTrue,defFalse,defEx,defAll}
\begin{defn}
تعریف شدہ عامل \pnbLatin{(defined operators)} ورت کے بنائے فارمولیاں نوں اینج
سمجھیا جاوے:

\begin{tagenumerate}{defTrue,defFalse,defNot,defOr,defAnd,defIf,defIff,defEx,defAll}

\tagitem{defTrue}{$\ltrue$،
  \iftag{prvFalse}{$\lnot\lfalse$}{$(!A \lor \lnot !A)$، کسے مقرر
    ایٹمی !!{formula}~$!A$ لئی} دا مخفف (abbreviation) اے۔}{}

\tagitem{defFalse}{$\lfalse$،
  \iftag{prvTrue}{$\lnot\ltrue$}{$(!A \land \lnot !A)$، کسے مقرر
    ایٹمی !!{formula}~$!A$ لئی} دا مخفف (abbreviation) اے۔}{}

\tagitem{defNot}{$\lnot !A$، $!A \lif \lfalse$ دا مخفف
  (abbreviation) اے۔}{}

\tagitem{defOr}{$!A \lor !B$،
  \iftag{prvAnd}{$\lnot(\lnot !A \land \lnot !B)$}{$\lnot !A \lif
    !B$} دا مخفف (abbreviation) اے۔}{}

\tagitem{defAnd}{$!A \land !B$،
  \iftag{prvOr}{$\lnot(\lnot !A \lor \lnot !B)$}{$\lnot (!A \lif
    \lnot !B)$} دا مخفف (abbreviation) اے۔}{}

% PNBPL-002 ماخذ دی درستی: پہلی متبادل عبارت دے آخری قوس دی شروع والی قوس ماخذ وچ رہ گئی اے؛ ایتھے پوری جوڑی شامل اے۔
\tagitem{defIf}{$!A \lif !B$،
  \iftag{prvOr}{$(\lnot !A \lor !B)$}{$\lnot (!A \land \lnot !B)$}
  دا مخفف (abbreviation) اے۔}{}

\tagitem{defIff}{$!A \liff !B$، $(!A \lif !B) \land (!B
  \lif !A)$ دا مخفف (abbreviation) اے۔}{}
\end{tagenumerate}
\end{defn}
\end{tagblock}

\begin{defn}[نحوی یکسانیت \pnbLatin{(Syntactic identity)}]
علامت $\ident$، علامتاں دیاں لڑیاں \pnbLatin{(strings of symbols)} وچ نحوی
یکسانیت ظاہر کردی اے، یعنی
$!A \ident !B$ اُدوں تے صرف اُدوں جدوں $!A$ تے $!B$ علامتاں دیاں
اکّو لمبائی دیاں لڑیاں ہون تے ہر مقام اُتے اکّو علامت رکھدیاں ہون۔
\end{defn}

علامت $\ident$ دے دونیں پاسے جوڑ (concatenation) نال ملیاں لڑیاں آ
سکدیاں نیں؛ مثلاً
$!A \ident (!B \lor !C)$ دا مطلب اے: علامتاں دی لڑی~$!A$ اوہی لڑی
اے جیہڑی، ایس ترتیب نال، اک کھلدی قوس، لڑی $!B$، علامت $\lor$،
لڑی~$!C$ تے اک بند ہُندی قوس نوں ایس ترتیب نال \pnbLatin{(in this order)} جوڑ کے
ملدی اے۔ جے ایہ گل اے، تاں
سانوں پتا اے کہ $!A$ دی پہلی علامت اک کھلدی قوس اے، $!A$ وچ $!B$ اک
ذیلی لڑی (substring) دے طور تے شامل اے (دوجی علامت توں شروع ہوندی اے)،
اوس ذیلی
لڑی دے بعد $\lor$ آؤندا اے، وغیرہ۔
\endgroup
% /OLP-0058

% OLP-0059 content/propositional-logic/syntax-and-semantics/preliminaries.tex
\begingroup
\def\olfilename{preliminaries}
\expandafter\def\csname ol@sectioncs\endcsname{section}

\olfileid{pl}{syn}{pre}

\olsection{مُڈھلیاں گلاں \pnbLatin{(Preliminaries)}}

\begin{thm}[\emph{!!{formula}s اُتے استقرا دا اصول \pnbLatin{(Principle of induction on formulas)}}]
  \ollabel{thm:induction}
  جے کوئی خاصیت~$P$ سارے ایٹمی (atomic) !!{formula}s لئی پوری ہُندی اے تے
  اینج اے کہ
  \begin{enumerate}
    \tagitem{prvNot}{ایہ $\lnot !A$ لئی پوری ہُندی اے جدوں وی ایہ
      ~$!A$ لئی پوری ہُندی اے؛}{}
    \tagitem{prvAnd}{ایہ $(!A \land !B)$ لئی پوری ہُندی اے
      جدوں وی ایہ $!A$ تے~$!B$ لئی پوری ہُندی اے؛}{}
    \tagitem{prvOr}{ایہ $(!A \lor !B)$ لئی پوری ہُندی اے
      جدوں وی ایہ $!A$ تے~$!B$ لئی پوری ہُندی اے؛}{}
    \tagitem{prvIf}{ایہ $(!A \lif !B)$ لئی پوری ہُندی اے
      جدوں وی ایہ $!A$ تے~$!B$ لئی پوری ہُندی اے؛}{}
    \tagitem{prvIff}{ایہ $(!A \liff !B)$ لئی پوری ہُندی اے
      جدوں وی ایہ $!A$ تے~$!B$ لئی پوری ہُندی اے؛}{}
  \end{enumerate}
  تاں $P$ سارے !!{formula}s لئی پوری ہُندی اے۔
\end{thm}

\begin{proof}
  $S$ نوں اوہناں سارے !!{formula}s دی جماعت (class) منّو جنہاں وچ
  خاصیت~$P$ اے۔ صاف اے کہ $S \subseteq \Frm[L_0]$۔ $S$،
  \olref[fml]{defn:formulas} دیاں ساریاں شرطاں پوری کردا اے: ایہدے وچ
  سارے ایٹمی !!{formula}s نیں تے ایہ !!{operator}s تحت بند \pnbLatin{(closed under the operators)} اے۔ $\Frm[L_0]$ ایہو جہی سب توں چھوٹی
  (smallest) جماعت اے، سو
  $\Frm[L_0] \subseteq S$۔ ایس لئی $\Frm[L_0] = S$، تے ہر فارمولے
  وچ خاصیت~$P$ اے۔
\end{proof}

\begin{prop}\ollabel{prop:balanced}
   $\Frm[L_0]$ وچ ہر !!{formula} \emph{متوازن (balanced)} اے، یعنی اوہدے وچ
   کھبیاں قوساں \pnbLatin{(left parentheses)} اونّیاں ای نیں جِنّیاں سجیاں
   قوساں \pnbLatin{(right parentheses)}۔
\end{prop}

\begin{prob}
\olref[pl][syn][pre]{prop:balanced} ثابت (Prove) کرو۔
\end{prob}

\begin{prop} \ollabel{prop:noinit}
کسی !!a{formula} دا پورے توں چھوٹا مُڈھلا ٹکڑا \pnbLatin{(proper initial segment)} !!a{formula} نہیں ہُندا۔
\end{prop}

\begin{prob}
  \olref[pl][syn][pre]{prop:noinit} ثابت (Prove) کرو۔
\end{prob}

\begin{prop}[اکّو طرح نال پڑھت \pnbLatin{(Unique Readability)}]
ہر !!{formula}~$!A$، جیہڑا $ \Frm[L_0]$ وچ اے، اوہدا اگلیاں صورتاں
وچوں کسے اک دے طور تے تجزیہ (parsing) ٹھیک اکّو طریقے نال ہُندا اے:
\begin{enumerate}
\tagitem{prvFalse}{$\lfalse$۔}{}

\tagitem{prvTrue}{$\ltrue$ (truth)۔}{}

\item $\Obj p_n$، کسے قضیاتی متغیر \pnbLatin{(propositional variable)}
$\Obj p_n \in  \PVar$ لئی۔

\tagitem{prvNot}{$\lnot !B$ (negation)، کسے !!{formula}~$!B$ لئی۔}{}

\tagitem{prvAnd}{$(!B \land !C)$ (conjunction)، کسے !!{formula}s $!B$ تے~$!C$ لئی۔}{}

\tagitem{prvOr}{$(!B \lor !C)$ (disjunction)، کسے !!{formula}s $!B$ تے~$!C$ لئی۔}{}

\tagitem{prvIf}{$(!B \lif !C)$ (conditional)، کسے !!{formula}s $!B$ تے~$!C$ لئی۔}{}

\tagitem{prvIff}{$(!B \liff !C)$ (biconditional)، کسے !!{formula}s $!B$ تے~$!C$ لئی۔}{}
\end{enumerate}
ہور ایہ تجزیاتی پڑھت (parsing) \emph{اکّو ای (unique)} ہُندی اے۔
\end{prop}

\begin{proof}
$!A$ اُتے استقرا (induction) نال۔ مثال لئی، منّو کہ $!A$ دیاں دو وکھریاں پڑھتاں
$(!B \lif !C)$ تے $(!B' \lif !C')$ نیں۔ تاں $!B$ تے $!B'$ لازماً
اکّو ہون گے (نہیں تاں اک دوجے دا proper initial segment، یعنی پورے توں
چھوٹا مُڈھلا ٹکڑا، ہُندا)؛ سو
جے $!A$ دیاں دونیں پڑھتاں وکھریاں نیں تاں وجہ ایہو ہو سکدی اے کہ
$!C$ تے $!C'$ علامتاں دی اکّو لڑی \pnbLatin{(same sequence of symbols)} دیاں
وکھریاں پڑھتاں نیں، جیہڑی گل استقرائی مفروضے \pnbLatin{(inductive hypothesis)}
موجب ناممکن اے۔
\end{proof}

\begin{defn}[یکساں بدل \pnbLatin{(Uniform Substitution)}]
جے $!A$ تے $!B$، !!{formula}s نیں، تے $\Obj p_i$ اک !!{propositional variable} اے، تاں $\Subst{!A}{!B}{\Obj p_i}$ اوہ نتیجہ ظاہر کردا اے
جیہڑا $\Obj p_i$ دی ہر واری آمد \pnbLatin{(each occurrence)} نوں $!B$ دی اک آمد
نال $!A$ وچ بدل کے ملدا اے؛ ایسّے طرح $\Obj p_1$، \dots،~$\Obj p_n$
نوں اکّو ویلے \pnbLatin{(simultaneous substitution)}
!!{formula}s $!B_1$، \dots،~$!B_n$ نال بدلنا
$\SSubst{!A}{\subst{!B_1}{\Obj p_1},\dots,\subst{!B_n}{\Obj p_n}}$
نال ظاہر کیتا جاندا اے۔
\end{defn}

\begin{prob} تھلے دتے پنجاں !!{formula}s وچوں ہر اک لئی طے کرو کہ کی
 اوہ !!{formula} بدل (substitution) \( \Subst{!A}{!B}{\Obj p_i} \) دے روپ وچ لکھیا
 جا سکدا اے، جتھے \( !A \) ایہ اے: (i) \( \Obj p_0 \)؛ (ii)
 \( ( \lnot \Obj p_0 \land \Obj p_1) \)؛ تے (iii)
 \( ( ( \lnot \Obj p_0 \lif \Obj p_1 ) \land \Obj p_2 ) \)۔ ہر
 صورت وچ متعلق بدل صاف صاف دسو۔
  \begin{enumerate}
    \item \( \Obj p_1 \)
    \item \( ( \lnot \Obj p_0 \land \Obj p_0 ) \)
    \item \( ( ( \Obj p_0 \lor \Obj p_1 ) \land \Obj p_2 )  \)
    \item \( \lnot ( ( \Obj p_0 \lif \Obj p_1 ) \land \Obj p_2 ) \)
    \item \( (( \lnot ( \Obj p_0 \lif \Obj p_1 ) \lif ( \Obj p_0 \lor \Obj p_1 )) \land \lnot ( \Obj p_0 \land \Obj p_1 )) \)
  \end{enumerate}
\end{prob}

\begin{prob}
استقرا (induction) نال $\Subst{!A}{!B}{p}$ دی ریاضیاتی طور تے ٹھوس
\pnbLatin{(mathematically rigorous)} تعریف دیو۔
\end{prob}
\endgroup
% /OLP-0059

% OLP-0060 content/propositional-logic/syntax-and-semantics/formation-sequences.tex
\begingroup
\def\olfilename{formation-sequences}
\expandafter\def\csname ol@sectioncs\endcsname{section}

\olfileid{pl}{syn}{fseq}

\olsection{تشکیل دے سلسلے \pnbLatin{(Formation Sequences)}}

!!{formula}s نوں استقرائی تعریف \pnbLatin{(inductive definition)} راہیں تعریف
کرنا، تے نال دی نال !!{formula}s دیاں خاصیتاں نوں استقرا (induction)
راہیں ثابت کرن دی تکنیک، اک سوہنا
تے کارگر طریقہ اے۔ پر !!{formula}s دی بناوٹ نوں ہِٹھوں اُتّے ول، قدم
بہ قدم طریقے \pnbLatin{(bottom-up, step-by-step approach)} نال ویکھنا وی فائدہ
مند ہو سکدا اے؛ ایتھے اسیں ایہ کم \emph{تشکیل دے سلسلے \pnbLatin{(formation sequence)}} دے تصور نال کردے آں۔

\begin{defn}[فارمولیاں لئی تشکیل دے سلسلے \pnbLatin{(Formation sequences for formulas)}]
\ollabel{defn:fseq-frm}
علامتاں دیاں لڑیاں \pnbLatin{(strings of symbols)} دا متناہی متوالیہ \pnbLatin{(finite sequence)} $\tuple{!A_0,\dotsc,!A_n}$، جیہڑا زبان~$\Lang L_0$ توں اے،
$!A$ لئی \emph{تشکیل دا سلسلہ \pnbLatin{(formation sequence)}} اے جے
$!A \ident !A_n$، تے ہر $i \leq n$ لئی یا تاں $!A_i$ ایٹمی فارمولا
 \pnbLatin{(atomic formula)} اے، یا ایہو جہے $j,k < i$ موجود نیں کہ اگلیاں وچوں کوئی اک گل پوری
ہُندی اے:
\begin{enumerate}
    \tagitem{prvNot}{$!A_i \ident \lnot !A_j$۔}{}%
    \tagitem{prvAnd}{$!A_i \ident (!A_j \land !A_k)$۔}{}%
    \tagitem{prvOr}{$!A_i \ident (!A_j \lor !A_k)$۔}{}%
    \tagitem{prvIf}{$!A_i \ident (!A_j \lif !A_k)$۔}{}%
    \tagitem{prvIff}{$!A_i \ident (!A_j \liff !A_k)$۔}{}%
\end{enumerate}
\end{defn}

\begin{ex}
\[
\tuple{
    \Obj p_0,
    \Obj p_1,
    (\Obj p_1 \land \Obj p_0),
    \lnot (\Obj p_1 \land \Obj p_0)
}
\]
$\lnot (\Obj p_1 \land \Obj p_0)$ لئی تشکیل دا اک سلسلہ \pnbLatin{(formation sequence)} اے، تے ایہ وی اے:
\[
\tuple{
    \Obj p_0,
    \Obj p_1,
    \Obj p_0,
    (\Obj p_1 \land \Obj p_0),
    (\Obj p_0 \lif \Obj p_1),
    \lnot (\Obj p_1 \land \Obj p_0)
}.
\]
%
دوجی مثال توں ویکھیا جا سکدا اے کہ تشکیل دے سلسلیاں وچ `فالتو مواد
(junk)' وی آ سکدا اے: اوہ فارمولے جیہڑے وادھو (redundant) ہون یا
بناوٹ وچ کوئی حصہ نہ پاؤن \pnbLatin{(do not contribute to the construction)}۔
\end{ex}

\begin{prop}
\ollabel{prop:formed}
ہر !!{formula}~$!A$، جیہڑا~$\Frm[L_0]$ وچ اے، اوہدے لئی تشکیل دا کوئی
سلسلہ \pnbLatin{(formation sequence)} موجود اے۔
\end{prop}

\begin{proof}
منّو $!A$ ایٹمی (atomic) اے۔ تاں سلسلہ~$\tuple{!A}$، $!A$ لئی تشکیل
دا سلسلہ \pnbLatin{(formation sequence)} اے۔
%
ہُن منّو کہ $!B$ تے~$!C$ دے تشکیل دے سلسلے ترتیب وار
$\tuple{!B_0,\dotsc,!B_n}$ تے $\tuple{!C_0,\dotsc,!C_m}$ نیں۔
%
\begin{enumerate}
  \tagitem{prvNot}{جے $!A \ident \lnot !B$،
      تاں $\tuple{!B_0,\dotsc,!B_n,\lnot !B_n}$
      $!A$ لئی تشکیل دا سلسلہ اے۔}{}
  \tagitem{prvAnd}{جے $!A \ident (!B \land !C)$،
      تاں $\tuple{!B_0,\dotsc,!B_n,!C_0,\dotsc,!C_m,(!B_n \land !C_m)}$
      $!A$ لئی تشکیل دا سلسلہ اے۔}{}
  \tagitem{prvOr}{جے $!A \ident (!B \lor !C)$،
      تاں $\tuple{!B_0,\dotsc,!B_n,!C_0,\dotsc,!C_m,(!B_n \lor !C_m)}$
      $!A$ لئی تشکیل دا سلسلہ اے۔}{}
  \tagitem{prvIf}{جے $!A \ident (!B \lif !C)$،
      تاں $\tuple{!B_0,\dotsc,!B_n,!C_0,\dotsc,!C_m,(!B_n \lif !C_m)}$
      $!A$ لئی تشکیل دا سلسلہ اے۔}{}
  \tagitem{prvIff}{جے $!A \ident (!B \liff !C)$،
      تاں $\tuple{!B_0,\dotsc,!B_n,!C_0,\dotsc,!C_m,(!B_n \liff !C_m)}$
      $!A$ لئی تشکیل دا سلسلہ اے۔}{}
\end{enumerate}
!!{formula}s اُتے استقرا دے اصول \pnbLatin{(principle of induction on formulas)}
نال، ہر !!{formula} دا کوئی تشکیل دا سلسلہ اے۔
\end{proof}

اسیں اُلٹا رخ (converse) وی ثابت کر سکدے آں۔ ایہ گل اہم اے، کیوں جے
ایہ وکھاؤندی اے کہ فارمولیاں نوں تعریف کرن دے ساڈے دونیں طریقے ہم ارز
(equivalent) نیں: دونیں اکّو نتیجے دیندے نیں۔ ایہدا مطلب ایہ وی اے کہ
اسیں تشکیل دے سلسلیاں دی لمبائی اُتے عام استقرا \pnbLatin{(ordinary induction)}
ورت کے فارمولیاں بارے قضیے ثابت کر سکدے آں۔

\begin{lem}
\ollabel{lem:fseq-init}
منّو کہ $\tuple{!A_0,\dotsc,!A_n}$،~$!A_n$ لئی تشکیل دا سلسلہ اے، تے
$k \leq n$۔ تاں $\tuple{!A_0,\dotsc,!A_k}$،~$!A_k$ لئی مُڈھلا ذیلی
سلسلہ \pnbLatin{(initial subsequence)} تے تشکیل دا سلسلہ اے۔
\end{lem}

\begin{proof}
مشق (Exercise)۔
\end{proof}

\begin{thm}
\ollabel{thm:fseq-frm-equiv}
$\Frm[L_0]$، زبان~$\Lang L_0$ وچ علامتاں دیاں اوہناں ساریاں لڑیواں
دا سیٹ اے جنہاں دا کوئی تشکیل دا سلسلہ \pnbLatin{(formation sequence)} اے۔
\end{thm}

\begin{proof}
$F$ نوں زبان~$\Lang L_0$ وچ علامتاں دیاں اوہناں ساریاں لڑیواں \pnbLatin{(strings of symbols)} دا سیٹ منّو جنہاں دا کوئی تشکیل دا سلسلہ اے۔ اسیں
\olref[pl][syn][fseq]{prop:formed} وچ ویکھ چکے آں کہ
$\Frm[L_0] \subseteq F$، سو ہُن اسیں اُلٹا رخ ثابت کردے آں۔

منّو $!A$ دا تشکیل دا سلسلہ $\tuple{!A_0,\dotsc,!A_n}$ اے۔ اسیں ثابت
کردے آں کہ $!A \in \Frm[L_0]$، تے ایہ کم~$n$ اُتے مضبوط استقرا
\pnbLatin{(strong induction)} نال کردے آں۔ ساڈا استقرائی مفروضہ \pnbLatin{(induction hypothesis)} اے کہ علامتاں دی ہر اوہ لڑی جیہدے تشکیل دے سلسلے دی
لمبائی $m < n$ اے،~$\Frm[L_0]$ وچ اے۔ تشکیل دے سلسلے دی تعریف موجب
یا تاں $!A_n$ ایٹمی اے، یا ایہو جہے $j,k < n$ لازماً موجود نیں کہ
اگلیاں وچوں کوئی اک صورت اے:
\begin{enumerate}
    \tagitem{prvNot}{$!A_n \ident \lnot !A_j$۔}{}%
    \tagitem{prvAnd}{$!A_n \ident (!A_j \land !A_k)$۔}{}%
    \tagitem{prvOr}{$!A_n \ident (!A_j \lor !A_k)$۔}{}%
    \tagitem{prvIf}{$!A_n \ident (!A_j \lif !A_k)$۔}{}%
    \tagitem{prvIff}{$!A_n \ident (!A_j \liff !A_k)$۔}{}%
\end{enumerate}
ہُن اسیں صورتاں نال دلیل دیندے آں۔ جے $!A_n$ ایٹمی اے تاں
$!A_n \in \Frm[L_0]$۔ ایس دی بجائے منّو کہ $!A \equiv
(!A_j \land !A_k)$۔ \olref[pl][syn][fseq]{lem:fseq-init} موجب،
$\tuple{!A_0,\dotsc,!A_j}$ تے $\tuple{!A_0,\dotsc,!A_k}$ ترتیب وار
$!A_j$ تے~$!A_k$ لئی تشکیل دے سلسلے نیں۔ کیوں جے ایہ~$!A$ دے تشکیل دے
سلسلے دے پورے توں چھوٹے مُڈھلے ذیلی سلسلے \pnbLatin{(proper initial subsequences)} نیں، ایس لئی دونیں دی
لمبائی~$n$ توں گھٹ اے۔ سو استقرائی مفروضے موجب $!A_j$ تے~$!A_k$،
$\Frm[L_0]$ وچ نیں، تے ایس لئی !!a{formula} دی تعریف موجب
$(!A_j \land !A_k)$ وی اوہدے وچ اے۔ ہور صورتاں ایسّے برابر دلیل
\pnbLatin{(parallel reasoning)} نال ملدیاں نیں۔
\end{proof}
\endgroup
% /OLP-0060

% OLP-0061 content/propositional-logic/syntax-and-semantics/valuations-sat.tex
\begingroup
\def\olfilename{valuations-sat}
\expandafter\def\csname ol@sectioncs\endcsname{section}

\olfileid{pl}{syn}{val}

\olsection{\usetoken{P}{valuation} \pnbLatin{(Valuation)} تے پورا اترنا \pnbLatin{(Satisfaction)}}

\begin{defn}[!!^{valuation}s]
$\{\True, \False\}$ نوں دو صداقتی قدراں \pnbLatin{(truth values)}، ``سچ''
(true) تے ``جھوٹ'' (false)، دا سیٹ منّو۔ $\Lang{L_0}$ لئی
\emph{!!{valuation} (valuation)} اک فنکشن (function)~$\pAssign{v}$ اے
جیہڑا زبان دے !!{propositional variable}s نوں $\True$ یا $\False$
وچوں کوئی اک مقرر کردا اے، یعنی $\pAssign{v} \colon
\PVar \to \{\True, \False \}$۔
\end{defn}

\begin{defn}
  دتی ہوئی !!a{valuation}~$\pAssign{v}$ لئی، قدر معلوم کرن والے فنکشن
  \pnbLatin{(evaluation function)} $\pValue{v} \colon \Frm[L_0] \to
  \{\True, \False \}$ دی استقرائی تعریف \pnbLatin{(inductive definition)} ایویں کرو:
  \begin{align*}
    \iftag{prvFalse}{\pValue{v}(\lfalse) & = \False; \\ }{}
    \iftag{prvTrue}{\pValue{v}(\ltrue) & = \True; \\ }{}
    \pValue{v}(\Obj p_n) & = \pAssign{v}(\Obj p_n); \\
    \iftag{prvNot}{\pValue{v}(\lnot !A) & = \begin{cases}
        \True & \text{جے } \pValue{v}(!A) = \False;\\
        \False & \text{ہور صورت وچ۔}
      \end{cases} \\ }{}
    \iftag{prvAnd}{\pValue{v}(!A \land !B) & = \begin{cases}
        \True &
        \text{جے $\pValue{v}(!A) = \True$ تے $\pValue{v}(!B) = \True$؛}\\
        \False &
        \text{جے $\pValue{v}(!A) = \False$ یا $\pValue{v}(!B) = \False$}.
      \end{cases}\\ }{}
    \iftag{prvOr}{\pValue{v}(!A \lor !B) & = \begin{cases}
        \True &
        \text{جے $\pValue{v}(!A) = \True$ یا $\pValue{v}(!B) = \True$؛}\\
        \False &
        \text{جے $\pValue{v}(!A) = \False$ تے $\pValue{v}(!B) = \False$}.
      \end{cases}\\ }{}
    \iftag{prvIf}{\pValue{v}(!A \lif !B) & = \begin{cases}
        \True &
        \text{جے $\pValue{v}(!A) = \False$ یا $\pValue{v}(!B) = \True$؛}\\
        \False &
        \text{جے $\pValue{v}(!A) = \True$ تے $\pValue{v}(!B) = \False$}.
      \end{cases}\\ }{}
    \iftag{prvIff}{\pValue{v}(!A \liff !B) & = \begin{cases}
        \True &
        \text{جے $\pValue{v}(!A) = \pValue{v}(!B)$؛}\\
        \False &
        \text{جے $\pValue{v}(!A) \neq \pValue{v}(!B)$}.
      \end{cases}}{}
\end{align*}
\end{defn}

\begin{explain}
ایہ شرطاں اگلے صداقتی جدولاں \pnbLatin{(truth tables)} دے مطابق نیں:
\begin{center}
\iftag{prvNot}{
  \begin{tabular}{|c||c|} \hline
    $!A$ & $ \lnot !A$ \\
    \hline \hline
    $\True$ & $\False$ \\
    $\False$ & $\True$ \\
    \hline
  \end{tabular}
}{}
\iftag{prvAnd}{
  \begin{tabular}{|cc||c|} \hline
    $!A$ & $!B$ & $!A \land !B$ \\
    \hline \hline
    $\True$ & $\True$ & $\True$ \\
    $\True$ & $\False$ & $\False$ \\
    $\False$ & $\True$ & $\False$ \\
    $\False$ & $\False$ & $\False$ \\
    \hline
  \end{tabular}
}{}
\iftag{prvOr}{
  \begin{tabular}{|cc||c|} \hline
    $!A$ & $!B$ & $!A \lor !B$ \\
    \hline \hline
    $\True$ & $\True$ & $\True$ \\
    $\True$ & $\False$ & $\True$ \\
    $\False$ & $\True$ & $\True$ \\
    $\False$ & $\False$ & $\False$ \\
    \hline
  \end{tabular}
}{}%
\iftag{notprvNot,notprvAnd,notprvOr,notprvIf}{}{\\[1em]}
\iftag{prvIf}{
  \begin{tabular}{|cc||c|} \hline
    $!A$ & $!B$ & $!A \lif !B$ \\
    \hline \hline
    $\True$ & $\True$ & $\True$ \\
    $\True$ & $\False$ & $\False$ \\
    $\False$ & $\True$ & $\True$ \\
    $\False$ & $\False$ & $\True$ \\
    \hline
  \end{tabular}
}{}
\iftag{prvIff}{
  \begin{tabular}{|cc||c|} \hline
    $!A$ & $!B$ & $!A \liff !B$ \\
    \hline \hline
    $\True$ & $\True$ & $\True$ \\
    $\True$ & $\False$ & $\False$ \\
    $\False$ & $\True$ & $\False$ \\
    $\False$ & $\False$ & $\True$ \\
    \hline
  \end{tabular}
}{}
\end{center}
\end{explain}

\begin{prob}
$\Lang{L_0}$ وچ اک سہ رُکنی رابطہ \pnbLatin{(ternary connective)}~$\diamondsuit$
شامل کرن اُتے غور کرو، جیہدی قدر (evaluation) اگلے طریقے نال دتی گئی اے:
\begin{gather*}
  \pValue{v}(\diamondsuit( !A , !B , !C ) ) = \begin{cases}
    \pValue{v}( !B ) &
    \text{جے $\pValue{v}(!A) = \True$؛}\\
    \pValue{v}( !C ) &
    \text{جے $\pValue{v}(!A) = \False $}.
  \end{cases}
\end{gather*}
ایس رابطے دا صداقتی جدول \pnbLatin{(truth table)} لکھو۔
\end{prob}

\begin{thm}[مقامی تعیین \pnbLatin{(Local Determination)}]
\ollabel{thm:LocalDetermination} منّو کہ $\pAssign{v_1}$ تے
$\pAssign{v_2}$ ایہو جہیاں !!{valuation}s نیں جیہڑیاں مقرر فارمولے
~$!A$ وچ آون
والے !!{propositional variable}s اُتے متفق نیں، یعنی
$\pAssign{v_1}(\Obj p_n) = \pAssign{v_2}(\Obj p_n)$ جدوں وی
$\Obj p_n$، !!{formula}~$!A$ وچ آندا اے۔ تاں $\pValue{v_1}$ تے
$\pValue{v_2}$،~$!A$ اُتے وی متفق نیں، یعنی
$\pValue{v_1}(!A) = \pValue{v_2}(!A)$۔
\end{thm}

\begin{proof}
$!A$ اُتے استقرا (induction) نال۔
\end{proof}

\begin{defn}[پورا اترنا (Satisfaction)]
\ollabel{defn:satisfaction} اسیں
  !!a{formula}~$!A$ دے !!a{valuation}~$\pAssign{v}$ نال
  \emph{پورا اترن (satisfaction)} دے تصور، $\pSat{v}{!A}$، دی استقرائی
  تعریف ایویں کر سکدے آں۔ (اسیں $\pSat/{v}{!A}$ نوں منفی پورا اترنا
  \pnbLatin{(negative satisfaction)}، یعنی ``$\pSat{v}{!A}$ نہیں''، دے معنی وچ
  لکھدے آں۔)
\begin{enumerate}
\tagitem{prvFalse}{%
  \indcase{!A}{\lfalse}{$\pSat/{v}{\indfrm}$۔}}{}

\tagitem{prvTrue}{%
  \indcase{!A}{\ltrue}{$\pSat{v}{\indfrm}$؛ ایہ صدق مستقل
  \pnbLatin{(truth constant)} دا معاملہ اے۔}}{}

\item \indcase{!A}{\Obj p_i}{$\pSat{v}{\indfrm}$
  اُدوں تے صرف اُدوں (iff) جدوں $\pAssign{v}(\Obj p_i) = \True$؛
  ایہ قضیاتی متغیر \pnbLatin{(propositional variable)} دا معاملہ اے۔}

\tagitem{prvNot}{%
  \indcase{!A}{\lnot !B}{$\pSat{v}{\indfrm}$ اُدوں تے صرف اُدوں جدوں
    $\pSat/{v}{!B}$؛ ایہ نفی (negation) دا معاملہ اے۔}}{}

\tagitem{prvAnd}{%
  \indcase{!A}{(!B \land !C)}{$\pSat{v}{\indfrm}$ اُدوں تے صرف اُدوں جدوں
    $\pSat{v}{!B}$ تے $\pSat{v}{!C}$؛ ایہ عطف (conjunction) دا معاملہ اے۔}}{}

\tagitem{prvOr}{%
  \indcase{!A}{(!B \lor !C)}{$\pSat{v}{\indfrm}$ اُدوں تے صرف اُدوں جدوں
    $\pSat{v}{!B}$ یا $\pSat{v}{!C}$ (یا دونیں؛ or both)؛ ایہ فصل
    (disjunction) دا معاملہ اے۔}}{}

\tagitem{prvIf}{%
  \indcase{!A}{(!B \lif !C)}{$\pSat{v}{\indfrm}$ اُدوں تے صرف اُدوں جدوں
    $\pSat/{v}{!B}$ یا $\pSat{v}{!C}$ (یا دونیں؛ or both)؛ ایہ شرطیہ
    (conditional) دا معاملہ اے۔}}{}

\tagitem{prvIff}{%
  \indcase{!A}{(!B \liff !C)}{$\pSat{v}{\indfrm}$ اُدوں تے صرف اُدوں جدوں
    یا تاں $\pSat{v}{!B}$ تے $\pSat{v}{!C}$ دونیں (both)، یا نہ
    $\pSat{v}{!B}$ تے نہ $\pSat{v}{!C}$ (neither)؛ ایہ دو شرطیہ
    (biconditional) دا معاملہ اے۔}}{}
\end{enumerate}
جے $\Gamma$، !!{formula}s دا سیٹ \pnbLatin{(set of formulas)} اے، تاں $\pSat{v}{\Gamma}$ اُدوں تے
صرف اُدوں جدوں $\pSat{v}{!A}$ ہر~$!A \in \Gamma$ لئی۔
\end{defn}

\begin{prop}\ollabel{prop:sat-value}
  $\pSat{v}{!A}$ اُدوں تے صرف اُدوں (iff) جدوں
  $\pValue{v}(!A) = \True$ \pnbLatin{(evaluates to True)}۔
\end{prop}

\begin{proof}
  $!A$ اُتے استقرا (induction) نال۔
\end{proof}

\begin{prob}
  \olref[pl][syn][val]{prop:sat-value} ثابت (Prove) کرو۔
\end{prob}
\endgroup
% /OLP-0061

% OLP-0062 content/propositional-logic/syntax-and-semantics/semantic-notions.tex
\begingroup
\def\olfilename{semantic-notions}
\expandafter\def\csname ol@sectioncs\endcsname{section}

\olfileid{pl}{syn}{sem}

\olsection{معنوی تصورات \pnbLatin{(Semantic Notions)}}

اسیں اگلے معنوی تصورات \pnbLatin{(semantic notions)} دی تعریف کردے آں:

\begin{defn}
\begin{enumerate}
\item !!^a{formula}~$!A$
  \emph{پورا ہون جوگ (satisfiable)} اے جے کسے (some)
  ~$\pAssign{v}$ لئی $\pSat{v}{!A}$ ہووے؛ ایہ
  \emph{پورا ہون جوگ نہیں (unsatisfiable)} جے کسے وی (no)
  $\pAssign{v}$ لئی $\pSat{v}{!A}$ نہ ہووے؛
\item !!^a{formula}~$!A$ اک \emph{ہمیشہ صادق فارمولا (tautology)} اے جے
  $\pSat{v}{!A}$ ساریاں (all) !!{valuation}s~$\pAssign{v}$ لئی ہووے؛
\item !!^a{formula}~$!A$ \emph{اتفاقی (contingent)} اے جے اوہ پورا ہون جوگ اے پر
  ہمیشہ صادق فارمولا نہیں؛
\item جے $\Gamma$، !!{formula}s دا سیٹ اے، تاں $\Gamma \Entails !A$
  (``$\Gamma$ توں $!A$ منطقی طور تے نکلدا اے''؛ semantic entailment)
  اُدوں تے صرف اُدوں
  جدوں $\pSat{v}{!A}$ ہر اوہ !!{valuation}~$\pAssign{v}$ لئی ہووے
  جیہدے لئی $\pSat{v}{\Gamma}$۔
\item جے $\Gamma$، !!{formula}s دا سیٹ اے، تاں $\Gamma$
  \emph{پورا ہون جوگ (satisfiable)} اے جے کوئی اک ایہو جہی
  !!a{valuation}~$\pAssign{v}$ اے جیہدے لئی $\pSat{v}{\Gamma}$ ہووے؛
  نہیں تاں $\Gamma$ \emph{پورا ہون جوگ نہیں (unsatisfiable)}۔
\end{enumerate}
\end{defn}

\begin{prob}
 اگلے چار !!{formula}s وچوں ہر اک لئی طے کرو کہ کی اوہ
 (a)~پورا ہون جوگ (satisfiable)، (b)~ہمیشہ صادق فارمولا (tautology)،
 تے (c)~اتفاقی (contingent) اے۔
\begin{enumerate}
  \item \( ( \Obj p_0 \lif ( \lnot \Obj p_1 \lif \lnot \Obj p_0 ) ) \).
  \item \( ( ( \Obj p_0 \land \lnot \Obj p_1 ) \lif ( \lnot \Obj p_0 \land \Obj p_2 )) \liff ( ( \Obj p_2 \lif \Obj p_0 ) \lif ( \Obj p_0 \lif \Obj p_1 )) \).
  \item \( ( \Obj p_0 \liff \Obj p_1 ) \lif ( \Obj p_2 \liff \lnot \Obj p_1 ) \).
  \item \( (( \Obj p_0 \liff ( \lnot \Obj p_1 \land \Obj p_2 )) \lor ( \Obj p_2 \lif ( \Obj p_0 \liff \Obj p_1 ))) \).
\end{enumerate}
\end{prob}

\begin{prop}
\ollabel{prop:semanticalfacts}
\begin{enumerate}
\item $!A$ اُدوں تے صرف اُدوں ہمیشہ صادق فارمولا (tautology) اے جدوں
  $\emptyset \Entails !A$؛
\item جے $\Gamma \Entails !A$ تے $\Gamma \Entails !A \lif !B$ تاں
  $\Gamma \Entails !B$؛ ایہ معنوی modus ponens اے؛
\item جے $\Gamma$ پورا ہون جوگ اے تاں $\Gamma$ دا ہر محدود ذیلی سیٹ
  \pnbLatin{(finite subset)} وی
  پورا ہون جوگ اے؛
\item \ollabel{def:monotonicity} یک رُخیت (Monotonicity): جے $\Gamma \subseteq \Delta$
  تے $\Gamma \Entails !A$ تاں $\Delta \Entails !A$ وی؛
\item \ollabel{def:Cut} تعدیت (Transitivity؛ Cut): جے $\Gamma \Entails !A$ تے
  $\Delta \cup \{ !A\} \Entails !B$ تاں $\Gamma \cup \Delta \Entails
  !B$۔
\end{enumerate}
\end{prop}

\begin{proof}
مشق (Exercise)۔
\end{proof}

\begin{prob}
\olref[pl][syn][sem]{prop:semanticalfacts} ثابت (Prove) کرو۔
\end{prob}

\begin{prop}\ollabel{prop:entails-unsat}
  $\Gamma \Entails !A$ اُدوں تے صرف اُدوں جدوں
  $\Gamma \cup \{\lnot !A\}$ پورا ہون جوگ نہیں (unsatisfiable)۔
\end{prop}

\begin{proof}
مشق (Exercise)۔
\end{proof}

\begin{prob}
\olref[pl][syn][sem]{prop:entails-unsat} ثابت (Prove) کرو۔
\end{prob}

\begin{thm}[معنوی استنباط دا قضیہ \pnbLatin{(Semantic Deduction Theorem)}]
  \ollabel{thm:sem-deduction} $\Gamma \Entails !A \lif !B$ اُدوں تے صرف
  اُدوں جدوں $\Gamma \cup \{!A\} \Entails !B$۔
\end{thm}

\begin{proof}
مشق (Exercise)۔
\end{proof}

\begin{prob}
\olref[pl][syn][sem]{thm:sem-deduction} ثابت (Prove) کرو۔
\end{prob}
\endgroup
% /OLP-0062

% OLP-0714 content/propositional-logic/syntax-and-semantics/completeness.tex
\begingroup
\def\olfilename{completeness}
\expandafter\def\csname ol@sectioncs\endcsname{section}

\olfileid{pl}{prp}{com}

\olsection{قضیاتی منطق دی تمامیت}

\begin{defn}
\ollabel{def:MaxCon}
سیٹ $\Gamma$، جیہڑا !!{formula}s دا اے، \emph{ودھ توں ودھ بے تضاد} اے جے اوہ بے تضاد اے تے جے $\Delta$ ایہو جیہا بے تضاد سیٹ اے کہ $\Gamma
\subseteq \Delta$، تاں $\Gamma = \Delta$۔
\end{defn}

\begin{lem}[سچائی دا لیما]
  \ollabel{lem:truth}
منّ لو $\Gamma$ ودھ توں ودھ بے تضاد اے۔ فیر:
\begin{enumerate}
\item \ollabel{lem:truth-1} $!A \in \Gamma$ اوہدوں تے صرف اوہدوں جدوں $\lnot !A
  \notin \Gamma$۔
\item \ollabel{lem:truth-2} $\Gamma \Proves !A$ اوہدوں تے صرف اوہدوں جدوں $!A \in \Gamma$۔
\item \ollabel{lem:truth-3} $!A \lif !B \in \Gamma$ اوہدوں تے صرف اوہدوں جدوں یا $!A \notin \Gamma$ یا $!B \in \Gamma$۔
\end{enumerate}
\end{lem}

\begin{proof}
\begin{enumerate}
\item منّ لو $!A \in \Gamma$۔ جے $\lnot !A \in \Gamma$ وی اے، تاں $\Gamma$ تضاد والا اے۔ تے جے نہ $!A$ تے نہ $\lnot!A$، $ \Gamma$ وچ اے، تاں \olref[fol][axd][prv]{prop:provability-exhaustive} نال $\Gamma\cup\{!A\}$ تے $ \Gamma \cup \{\lnot !A\}$ وچوں اک بے تضاد اے، یعنی $\Gamma$ ودھ توں ودھ نہیں۔ دوجی سمت دا ثبوت وی ایویں اے۔
\item جے $!A \in \Gamma$، تاں صاف اے کہ $\Gamma \Proves !A$۔ دوجے پاسے، منّ لو $\Gamma \Proves !A$۔ جے $!A \notin \Gamma$، تاں \olref{lem:truth-1} نال $\lnot !A \in \Gamma$۔ سو $\Gamma \Proves
  !A$ تے $\Gamma \Proves \lnot !A$، یعنی تضاد۔
\item مشق۔
\end{enumerate}
\end{proof}

\olref{lem:truth}\olref{lem:truth-2} دی کھبّے توں سجّے ول سمت وکھاؤندی اے کہ $\Gamma$ \emph{اخذ دے لحاظ نال بند} اے، یعنی جے $\Gamma \Proves !A$، تاں $!A
\in \Gamma$ ہر~$!A$ لئی۔

\begin{prob}
  \olref{lem:truth} دا ثبوت پورا کرو۔
\end{prob}

\begin{thm}[تمامیت]
\ollabel{thm:completeness}
جے $\Gamma$ بے تضاد اے، تاں $\Gamma$ پورا ہون جوگا اے۔
\end{thm}

\begin{proof}
منّ لو $!A_0$، $!A_1$،~\dots، زبان دے سارے !!{formula}s دی پوری فہرست اے۔ اسیں !!{formula}s دے سیٹاں دی ودھدی لڑی $\Gamma_0$، $\Gamma_1$،~\dots، بازگشتی طور اُتے ایویں متعین کردے آں:
\begin{align*}
\Gamma_0 & = \Gamma\\
\Gamma_{n+1} &  =
\begin{cases}
  \Gamma_n \cup \{!A_n\} & \text{جے $\Gamma_n \cup \{!A_n\}$ بے تضاد اے،}\\
  \Gamma_n \cup \{\lnot !A_n\} & \text{نئیں تاں}.
\end{cases}
\end{align*}
فیر ایہ متعین کرو:
\[
\Gamma^* = \bigcup_{0\le n}\Gamma_n.
\]
ہُن ثبوت ہِٹھلیاں گلاں واری واری ثابت کر کے اگے ودھدا اے:
\begin{enumerate}
\item ہر $n$ لئی سیٹ $\Gamma_n$ بے تضاد اے ($n$ اُتے استقرا نال، \olref[fol][axd][prv]{prop:provability-exhaustive} ورت کے)۔
\item $\Gamma^*$ بے تضاد اے۔
\item $\Gamma^*$ ودھ توں ودھ اے۔
\end{enumerate}
فیر قدر بندی $v$ ایویں متعین کرو کہ $v(p_i) = \True$ اوہدوں تے صرف اوہدوں جدوں $p_i \in \Gamma^*$۔ فیر $!A$ اُتے استقرا نال وکھایا جاندا اے کہ $\Gamma^*$ وچ رکنیت، $v$ دے حساب نال سچائی دے برابر اے، ہور گنجل دار !!{sentence}s لئی وی۔ $ \overline{v}(!A) = \True$ اوہدوں تے صرف اوہدوں جدوں $!A \in \Gamma^*$۔ خاص کر، $v$، $\Gamma^*$ اُتے پورا اُتردا اے، تے کیونکہ $\Gamma \subseteq \Gamma^*$، $\Gamma$ وی پورا ہون جوگا اے، جیویں سانوں چاہیدا سی۔
\end{proof}

\begin{cor}
جے $\Gamma \Entails !A$، تاں $\Gamma \Proves !A$۔
\end{cor}

\begin{proof}
جے $\Gamma\Proves/ !A$، تاں $\Gamma \cup \{\lnot !A\}$ بے تضاد اے، \olref[fol][axd][prv]{prop:prov-incons} نال۔ سو تمامیت دے قضیے نال $\Gamma \cup
\{\lnot !A\}$ پورا ہون جوگا اے، تے $\Gamma \Entails/ !A$۔
\end{proof}

\begin{prop}[ٹھوس پن دا قضیہ]
\ollabel{prop:compactness}
$\Gamma$ پورا ہون جوگا اے اوہدوں تے صرف اوہدوں جدوں ہر \emph{محدود} ذیلی سیٹ $\Gamma_0$، $ \Gamma$ دا، پورا ہون جوگا اے۔
\end{prop}

\begin{proof}
$\Gamma$ پورا ہون جوگا نہیں اوہدوں تے صرف اوہدوں جدوں اوہ تضاد والا اے، اوہدوں تے صرف اوہدوں جدوں کوئی محدود ذیلی سیٹ $\Gamma_0$، $ \Gamma$ دا، تضاد والا اے، اوہدوں تے صرف اوہدوں جدوں کوئی محدود ذیلی سیٹ $\Gamma_0$، $ \Gamma$ دا، پورا ہون جوگا نہیں۔
\end{proof}

\begin{cor}
$\Gamma \Entails !A$ اوہدوں تے صرف اوہدوں جدوں کسے محدود ذیلی سیٹ $\Gamma_0$، $ \Gamma$ دے لئی، $\Gamma_0 \Entails !A$۔
\end{cor}
\endgroup
% /OLP-0714
% OLP-0715 content/propositional-logic/syntax-and-semantics/soundness.tex
\begingroup
\def\olfilename{soundness}
\expandafter\def\csname ol@sectioncs\endcsname{section}

\olfileid{pl}{prp}{snd}

\olsection{قضیاتی منطق دی صحت مندی}

\begin{thm}[صحت مندی]
\ollabel{thm:soundness}
جے $\Gamma\Proves!A$، تاں $\Gamma \Entails !A$۔
\end{thm}

\begin{proof}
قضیاں اُتے استقرا نال۔ جے $!A$ مسلّمہ اے، تاں $\Entails
!A$ تے ایس لئی بدرجہ اولیٰ $\Gamma \Entails!A$۔ ایویں ای جے $!A \in \Gamma$، تاں $\Gamma \Entails!A$۔ استقرائی قدم لئی منّ لو $!A$، \emph{شرطیہ توں نتیجہ کڈھن دے قاعدے} نال $!C$ تے $!C\lif!A$ توں ملیا اے۔ استقرا دے مفروضے دے طور اُتے ایہ وی منّ لو کہ قضیہ $!C\lif !A$ تے $!C$ لئی ٹھیک اے۔ \olref[pl][syn][sem]{prop:semanticalfacts} دی تِیجی مدّ توں مطلوبہ نتیجہ ملدا اے۔
\end{proof}

\begin{cor}
جے $\Gamma$ پورا ہون جوگا اے، تاں $\Gamma$ بے تضاد اے۔ سو قضیاتی منطق بے تضاد اے۔
\end{cor}
\endgroup
% /OLP-0715
\OLEndChapterHook
\endgroup
% /OLP-0056


\chapter{قضیاتی منطق دے ثبوتی نظام}
قضیاتی ثبوتاں دے قاعدے تے تمامیت دے ثبوت اِنہاں باباں وچ ویکھو۔
\begin{itemize}
\item \hyperref[fol:prf::chap]{اخذ نظام}
\item \hyperref[fol:seq::chap]{سیکوئنٹ حساب}
\item \hyperref[fol:ntd::chap]{فطری استنباط}
\item \hyperref[fol:tab::chap]{تابلو}
\item \hyperref[fol:axd::chap]{مسلّمی اخذ}
\item \hyperref[fol:com::chap]{تمامیت دا قضیہ}
\end{itemize}
ماخذ فائل توں اِنہاں باباں دا صرف قضیاتی روپ وی تیار ہو سکدا اے۔


\OLEndPartHook
\endgroup
% /OLP-0055

% OLP-0138 content/first-order-logic/first-order-logic.tex
\begingroup
\def\olfilename{first-order-logic}
\expandafter\def\csname ol@sectioncs\endcsname{section}

\olpart{fol}{پہلے درجے دی منطق}

\begin{editorial}
  ایہ حصہ پہلے درجے دی منطق دے مابعد نظریے نوں
  تمامیت تک بیان کردا اے۔ فی الحال ایہ قضیاتی
  منطق دی الگ بحث اُتے منحصر نہیں؛ ہر نتیجہ ایتھے
  ثابت کیتا جاندا اے۔ جے ”FOL“ ٹیگ غلط ہووے
  تاں ماخذ فائلیں مسوّراں والا مواد کڈھ دیندیاں نیں
  (تے ”!!{structure}“ دی تھاں
  ”!!{valuation}“، $\Struct{M}$ دی تھاں
  $\pAssign{v}$، وغیرہ رکھ دیندیاں نیں)۔
  اصل وچ قضیاتی منطق والے حصے دا زیادہ تر
  مواد، پہلے درجے والے مواد وچ ”FOL“ ٹیگ
  بند کر کے ای بن جاندا اے۔

  جے قضیاتی منطق والا حصہ وی شامل ہووے
  تاں بہت دہراؤ ہوندا اے۔ البتہ منصوبہ ایہ
  اے کہ اگے چل کے ایہ حصہ قضیاتی منطق
  والے مواد نوں وی حساب وچ رکھ سکے:
  جیہڑا مواد پہلے ای آ چکیا اے اوہنوں
  اتھوں کڈھ دے، تے قضیاتی منطق والے
  حصے دے متعلقہ مقاماں دے حوالیاں نال
  ثبوت مختصر کرے۔
\end{editorial}

% OLP-0139 content/first-order-logic/introduction/introduction.tex
\begingroup
\def\olfilename{introduction}
\expandafter\def\csname ol@sectioncs\endcsname{section}

\olchapter{fol}{int}{پہلے درجے دی منطق دا تعارف}

% OLP-0140 content/first-order-logic/introduction/first-order-logic.tex
\begingroup
\def\olfilename{first-order-logic}
\expandafter\def\csname ol@sectioncs\endcsname{section}

\olfileid{fol}{int}{fol}

\olsection{پہلے درجے دی منطق}

رسمی منطق دے پہلے تعارف توں تسی غالباً پہلے درجے
دی منطق نال واقف او۔\footnote{اصل وچ اسیں کم و بیش
ایہی فرض کردے آں!{} جے نہیں، تاں تسی ہور بنیادی
درجے دی کتاب، مثلاً \emph{forall x}
\citep{Magnus2021}، ویکھ سکدے او۔}
تسی اینوں ”مسوّراں والی منطق“ یا ”محمولی منطق“
وی آکھدے ہووگے۔ پہلے درجے دی منطق سب توں پہلے
اک رسمی زبان اے۔ یعنی ایس دی اک خاص لفظیات اے،
تے ایس دے اظہارات اوس لفظیات دیاں علامتاں دیاں
لڑیاں نیں۔ پر ہر لڑی جائز نہیں۔ جائز اظہارات
دیاں وکھ وکھ قسماں نیں: ٹرم، !!{formula}s تے
!!{sentence}s۔ ساڈی بنیادی دلچسپی پہلے درجے
دی منطق دے !!{sentence}s وچ اے: ایہ انگریزی
دے جملیاں دا رسمی ہم روپ دیندے نیں، تے انہاں
بارے اسیں اوہ سوال پچھ سکدے آں جنہاں وچ منطقی
ماہر عموماً دلچسپی رکھدا اے۔ مثال لئی:
\begin{itemize}
    \item کی $!B$،~$!A$ دا منطقی نتیجہ اے؟
    \item کی $!A$ منطقی طور تے سچا، منطقی طور تے
غلط، یا حالت اُتے منحصر اے؟
    \item کی $!A$ تے $!B$ ہم ارز نیں؟
\end{itemize}

ایہ سوال بنیادی طور تے پہلے درجے دی منطق دے
!!{sentence}s دے ”معنی“ بارے نیں۔ مثال لئی،
فلسفی ایس سوال نوں، کی $!B$ منطقی طور تے~$!A$
توں نکلدا اے، ایویں پرکھے گا: کی کوئی ایسی حالت
اے جتھے $!A$ سچا اے پر~$!B$ غلط
($!B$،~$!A$ توں نہیں نکلدا)، یا ہر اوہ حالت
جیہڑی $!A$ نوں سچا بناندی اے،~$!B$ نوں وی
سچا بناندی اے ($!B$،~$!A$ توں نکلدا اے)؟
پر ہن تک اسیں ایہ نہیں دسیا کہ ”حالت“ کیہ اے؛
ایہ کم \emph{معنویات} دا اے۔ پہلے درجے دی منطق
دیاں معنویات فلسفی دے ”حالت“ والے فطری خیال
دا ریاضیاتی طور تے ٹھیک نمونہ دیندیاں نیں،
تے نال ایہ وی طے کردیاں نیں---تے ایہ اہم اے---
کہ کسے !!a{sentence}~$!A$ دا کسے حالت
\emph{وچ سچا ہونا} کیہ اے۔ اسیں جیہڑا
ریاضیاتی طور تے ٹھیک نمونہ بناؤن لگے آں، اوہنوں
!!a{structure} آکھدے آں۔ جیہڑا تعلق
”وچ سچا“ نوں ٹھیک معنی دیندا اے، اوہنوں
\emph{پورا اترن} دا تعلق آکھدے نیں۔ سو اسیں
ایہ تعریف کراں گے کہ $!A$،
ساخت~$\Struct{M}$ وچ پورا اتردا اے
(علامتاں وچ: $\Sat{M}{!A}$)،
جتھے !!{sentence}s~$!A$ تے
!!{structure}s~$\Struct{M}$ ہون۔
ایہ ہو جان توں بعد اسیں ”توں نتیجہ نکلنا“
یا ”منطقی طور تے سچا ہونا“ ورگیاں ہور
معنوی اصطلاحاں دیاں ٹھیک تعریفاں وی دے سکاں
گے۔ انہاں تعریفاں نال ریاضیاتی طور تے
ایہ طے ہو سکے گا کہ، مثال لئی،
% PNB-INT-001: چیونٹی/کیڑے دی دلیل دے پانچ bracket slips عین نشان زدہ درست کیتے۔
$\lforall[x][(!A(x) \lif !B(x))],
\lexists[x][!A(x)] \Entails \lexists[x][!B(x)]$
درست اے یا نہیں۔ جواب یقیناً ”ہاں“ اے۔
جے تسی انگریزی دے جملیاں نوں پہلے درجے
دی منطق وچ علامتاں نال لکھنا سکھ چکے او،
تاں ایہ، مثال لئی، ”ساریاں چیونٹیاں کیڑے
نیں؛ چیونٹیاں موجود نیں؛ ایس لئی کیڑے موجود
نیں“ دی علامتی صورت اے۔ ایہ صاف طور تے
درست دلیل اے؛ ایس لئی ساڈی رسمی زبان
وچ ”توں نتیجہ نکلنا“ دا ریاضیاتی نمونہ
وی ایہی جواب دینا چاہیدا اے۔

رسمی منطق دے پہلے تعارف توں تہانوں ایہ
وی یاد ہووے گا کہ \emph{!!{derivation}s}
ہوندے نیں۔ جے تسی رسمی منطق دا ابتدائی
کورس پڑھیا اے، تاں استاد نے تہانوں ایہوجے
!!{derivation}s لبھن دی مشق کرائی ہووے گی،
شاید اک ایسا !!a{derivation} وی لبھن
لئی آکھیا ہووے جیہڑا اُتّے والا منطقی
نتیجہ وکھاندا اے۔ !!{derivation}s دین
دے کئی طریقے نیں: ممکن اے تسی ”فطری
استنباط“ یا ”سچائی دے درخت“ پڑھے ہون،
پر ہور وی بہت نیں۔ !!{derivation} نظاماں
دا مقصد اوہ اوزار دینا اے جنہاں نال اُتّے
والے منطقی سوالاں دے جواب لبھن۔ مثال لئی،
فطری استنباط دا اک !!{derivation} جیہدے
% PNB-INT-001: universal premise تے existential conclusion دے source bracket slips درست۔
مقدمے
$\lforall[x][(!A(x) \lif !B(x))]$
تے $\lexists[x][!A(x)]$ ہون تے نتیجہ
(آخری سطر) $\lexists[x][!B(x)]$ ہووے،
ایہ \emph{تصدیق کردا اے} کہ
$\lexists[x][!B(x)]$ منطقی طور تے
$\lforall[x][(!A(x) \lif !B(x))]$
تے $\lexists[x][!A(x)]$ توں نکلدا اے۔

پر ایہ کیوں؟ بظاہر !!{derivation} نظاماں
دا معنویات نال کوئی تعلق نہیں: رسمی
!!{derivation} دینا صرف علامتاں نوں
کجھ قاعدیاں مطابق رکھنا اے؛ اوس وچ
”حالتاں“ یا ”وچ سچا“ دا ذکر وی نہیں۔
!!{derivation} نظاماں تے معنویات دا
تعلق مابعد منطقی تحقیق نال قائم کرنا
پینا اے۔ لوڑ اک ریاضیاتی ثبوت دی اے،
% PNB-INT-001: آخری biconditional مثال دے source bracket slips درست۔
مثلاً ایہ کہ
$\lexists[x][!B(x)]$ دا رسمی
!!{derivation} مقدمیاں
$\lforall[x][(!A(x) \lif !B(x))]$
تے $\lexists[x][!A(x)]$ توں
اُدوں تے صرف اُدوں
ممکن اے جدوں
$\lforall[x][(!A(x) \lif !B(x))]$
تے $\lexists[x][!A(x)]$ مل کے
~$\lexists[x][!B(x)]$ دا منطقی
نتیجہ دیندے ہون۔ ایہ کم کرنے توں
پہلاں، البتہ، نحو تے معنویات دیاں
تعریفاں ٹھیک بٹھان لئی کافی باریک
کم کرنا پینا اے۔
\endgroup
% /OLP-0140

% OLP-0141 content/first-order-logic/introduction/syntax.tex
\begingroup
\def\olfilename{syntax}
\expandafter\def\csname ol@sectioncs\endcsname{section}

\olfileid{fol}{int}{syn}

\olsection{نحو}

سب توں پہلاں سانوں ٹھیک ٹھیک طے کرنا پینا اے کہ علامتاں دیاں
کیہڑیاں لڑیاں پہلے درجے دی منطق دے !!{sentence}s شمار ہوندیاں
نیں۔ ایہ کم اسیں بعد وچ کراں گے؛ فی الحال مثالاں نال اگے
ودھدے آں۔ پہلے درجے دی منطق دیاں بنیادی اینٹاں---یعنی
لفظیات---دے دو حصے نیں۔ پہلا حصہ اوہ علامتاں نیں جنہاں
نال اسیں کسے خاص شے ول اشارہ کردے آں یا اوہدے بارے کوئی
خاص گل کہندے آں۔ کسے شے ول اشارہ کرن لئی اسیں
!!{constant}s ورتدے آں، تے اوہدی بابت گل کرن لئی
!!{predicate}s۔ مثال لئی، کسے اک شے ول اشارہ کرن
لئی اسیں $\Obj a$ نوں !!a{constant} بنا سکدے آں،
تے اوس شے بارے !!{sentence}~$\Atom{\Obj P}{\Obj
a}$ دے ذریعے کجھ کہہ سکدے آں۔ جے تہاڈے ذہن وچ
”$\Obj a$“ تے ”$\Obj P$“ دے کجھ معنی ہون،
تاں $\Atom{\Obj P}{\Obj a}$ نوں انگریزی دے جملے
وانگ پڑھ سکدے او (تے جدوں تسی پہلی واری رسمی منطق
سکھی ہووے گی، غالباً ایویں پڑھیا وی ہووے گا)۔
جدوں پہلے درجے دی منطق دے ایہوجے سادے !!{sentence}s
تہاڈے کول ہون، تاں لفظیات دے دوجے حصے، یعنی منطقی
علامتاں (رابطے تے مسوّر)، نال ہور گنجل دار جملے
بنا سکدے او۔ مثلاً اسیں
$(\Atom{\Obj P}{\Obj a} \land \Atom{\Obj Q}{\Obj b})$
یا~$\lexists[\Obj x][\Atom{\Obj P}{\Obj x}]$
ورگے اظہار بنا سکدے آں۔

ٹھیک معنوی تعریفاں تے اپنے !!{derivation} نظاماں
دے قاعدے دین لئی، جیہڑے سخت مابعد منطقی مطالعے
لئی ضروری نیں، سانوں سب توں پہلاں ایہ ٹھیک
تعریف دینا پینی اے کہ پہلے درجے دی منطق دا
!!a{sentence} کیہنوں آکھدے آں۔ بنیادی خیال سمجھنا
اوکھا نہیں: کجھ سادے !!{sentence}s اسیں صرف
!!{predicate}s تے !!{constant}s نال بنا سکدے آں،
مثلاً~$\Atom{\Obj P}{\Obj a}$۔ فیر انہاں توں
رابطیاں تے مسوّراں نال ہور گنجل دار جملے
بناندے آں۔ پر اوہ عین کیہڑے قاعدے نیں جنہاں
مطابق اسیں ایہ گنجل دار !!{sentence}s بنا سکدے
آں؟ انہاں قاعدیاں نوں دَسنا ضروری اے؛ ورنہ
”پہلے درجے دی منطق دا !!{sentence}“ ساڈی تعریف
وچ کافی ٹھیک نہیں ہووے گا۔ ایتھے کجھ مسئلے
نیں۔ پہلا، عین اوہ لڑیاں چننا جیہڑیاں
!!{sentence}s شمار ہون۔ دوجا، ایہ کم ایس طرح
کرنا کہ اسیں \emph{سارے} !!{sentence}s بارے
ریاضیاتی ثبوت دے سکّیے۔ آخر وچ سانوں
!!{sentence}s اُتے کجھ بنیادی عملاں دی وی ٹھیک
تعریف دینی پینی اے، مثلاً ”ہر $\Obj x$،
جیہڑا~$!A$ وچ آوندا اے، اوہدی تھاں~$\Obj b$
رکھو“۔
اُکّھیڑ ایہ اے کہ پہلے درجے دی منطق دے
مسوّر تے !!{variable}s ایہ کم کیویں کرنا اے،
ایہنوں بالکل صاف نہیں رہن دیندے۔ مثال لئی،
کی $\lexists[\Obj x][\Atom{\Obj P}{\Obj a}]$
نوں !!a{sentence} شمار کرنا چاہیدا اے؟
تے $\lexists[\Obj x][\lexists[\Obj x][\Atom{\Obj P}{\Obj x}]]$
دا کیہ؟ تے ”$\Obj x$ دی تھاں~$\Obj b$ رکھو،
جیہڑا $(\Atom{\Obj P}{\Obj x} \land \lexists[\Obj x][\Atom{\Obj P}{\Obj x}])$
وچ آوندا اے“ دا نتیجہ کیہ ہونا چاہیدا اے؟
\endgroup
% /OLP-0141

% OLP-0142 content/first-order-logic/introduction/formulas.tex
\begingroup
\def\olfilename{formulas}
\expandafter\def\csname ol@sectioncs\endcsname{section}

\olfileid{fol}{int}{fml}

\section{\usetoken{P}{formula}}

پہلے درجے دی منطق دے !!{sentence}s نوں پوری ٹھیکائی نال
مقرر کرن تے !!{variable}s دی ورتوں توں پیدا ہون والے
مسئلیاں نال نمٹن لئی اسیں ایہ طریقہ ورتاں گے۔
سب توں پہلاں اسیں اظہارات دا اک \emph{وکھرا} سیٹ
تعریف کراں گے: !!{formula}s۔ ایس توں بعد
اسیں انہاں وچ !!{variable}s دے کردار نوں ویکھ
سکاں گے---تے کجھ ہور تصورات، یعنی ”آزاد“ تے
”پابند“ !!{variable}s، دے سہارے ایہ تعریف
کر سکاں گے کہ !!a{sentence} کیہ اے (یعنی ایسا
!!a{formula} جیہدے وچ آزاد !!{variable}s نہ
ہون)۔ اسیں ایہ صرف ایس لئی نہیں کردے کہ
”!!{sentence}“ دی تعریف سنبھالن جوگی ہو جائے،
بلکہ ایس لئی وی کہ معنوی تصور، یعنی پورا اترن،
دی تعریف دے طریقے وچ ایہ فرق نہایت اہم ہووے گا۔

آؤ اک سادہ پہلے درجے دی زبان لئی ”!!{formula}“
تعریف کریے، جیہدے وچ صرف اک !!{predicate}~$\Obj P$
تے اک !!{constant}~$\Obj a$ ہون، تے منطقی
علامتاں وچ صرف $\lnot$، $\land$ تے~$\lexists$
ہون۔ ساڈیاں پوری تعریفاں ایہناں توں کہیں عام
ہون گیاں: اسیں لامتناہی تعداد وچ !!{predicate}s
تے !!{constant}s دی اجازت دیاں گے۔ اصل وچ
اسیں !!{function}s نوں وی ویکھاں گے، جیہڑے
!!{constant}s تے !!{variable}s نال مل کے
”ٹرم“ بنا سکدے نیں۔ فی الحال صرف $\Obj a$
تے متغیر ساڈے ٹرم ہون گے۔ البتہ سانوں
لامتناہی !!{variable}s ضرور چاہیدے نیں۔
باقاعدہ طور تے اسیں $\Obj v_0$، $\Obj
v_1$، \dots، نوں متغیر علامتاں ورتاں گے۔

\begin{defn}
\emph{!!{formula}s} دا سیٹ~$\Frm$ ایویں تعریف ہوندا اے:
\begin{enumerate}
\item\ollabel{fmls-atom} $\Atom{\Obj P}{\Obj a}$ تے $\Atom{\Obj
  P}{\Obj v_i}$ !!{formula}s نیں ($i \in \Nat$)۔

\tagitem{prvNot}{\ollabel{fmls-not}جے $!A$ !!a{formula} اے، تاں $\lnot !A$
  وی !!{formula} اے۔}{}

\tagitem{prvAnd}{جے $!A$ تے $!B$ !!{formula}s نیں، تاں $(!A \land
  !B)$ وی !!a{formula} اے۔}{}

\tagitem{prvEx}{\ollabel{fmls-ex}جے $!A$ !!a{formula} اے تے $x$ !!a{variable}
  اے، تاں $\lexists[x][!A]$ وی !!a{formula} اے۔}{}

\tagitem{limitClause}{\ollabel{fmls-limit}ایہناں توں سوا ہور کجھ وی !!a{formula} نہیں۔}{}
\end{enumerate}
\end{defn}

\olref{fmls-atom} سانوں دسدا اے کہ
$\Atom{\Obj P}{\Obj a}$ تے
$\Atom{\Obj P}{\Obj v_i}$، ہر $i \in
\Nat$ لئی، !!{formula}s نیں۔ ایہناں نوں
\emph{ایٹمی} !!{formula}s آکھدے نیں۔
ایہ ابتدا لئی مواد دیندے نیں۔ باقی شقّاں
پہلاں بنے فارمولیاں توں نویں !!{formula}s
بنان دے طریقے دیندیاں نیں۔ مثال لئی،
\olref{fmls-not} توں سانوں
$\lnot \Atom{\Obj P}{\Obj v_2}$ اک
!!a{formula} ملدا اے، کیوں جے
$\Atom{\Obj P}{\Obj v_2}$ پہلے ہی
\olref{fmls-atom} مطابق !!a{formula}
اے۔ فیر \olref{fmls-ex} توں
$\lexists[\Obj v_2][\lnot \Atom{\Obj P}{\Obj v_2}]$
وی اک !!{formula} ملدا اے، تے ایس
طرح اگے۔ \olref{fmls-limit} دسدا
اے کہ \emph{صرف} ایس طریقے نال
بنائیاں لڑیاں ہی !!{formula}s شمار
ہوندیاں نیں۔ خاص طور تے
$\lexists[\Obj v_0][\Atom{\Obj P}{\Obj a}]$
تے $\lexists[\Obj
  v_0][\lexists[\Obj v_0][\Atom{\Obj P}{\Obj a}]]$
!!{formula}s \emph{ہوندے} نیں؛ پر
$(\lnot \Atom{\Obj P}{\Obj a})$
فارمولا نہیں، کیوں جے اوہدے
باہر وادھو قوسین نیں۔

!!{formula}s نوں ایس طریقے نال تعریف کرن
نوں \emph{استقرائی تعریف} آکھدے نیں۔
ایہ سانوں !!{formula}s بارے استقرا
والے ثبوت دا اک روپ، یعنی
\emph{ساختی استقرا}، ورت کے قضیے
ثابت کرن دی اجازت دیندا اے۔
انہاں تصورات اُتے عمومی بحث
\olref[mth][ind][idf]{sec} تے
\olref[mth][ind][sti]{sec} وچ اے؛
اگے دے ثبوتاں وچ جان توں پہلاں
تہانوں اوہ حصہ مُڑ ویکھ لینا چاہیدا اے۔
بنیادی خیال ایہ اے: جے تسی ثابت کرنا
چاہندے او کہ کوئی گل سارے !!{formula}s
لئی سچی اے، تاں سب توں پہلاں اوہ
ایٹمی !!{formula}s لئی ثابت کرو؛
فیر دکھاؤ کہ \emph{جے} اوہ کسے
!!{formula}~$!A$ (تے~$!B$) لئی سچی
اے، تاں $\lnot !A$، $(!A \land !B)$
تے $\lexists[x][!A]$ لئی \emph{وی}
سچی اے۔ مثال لئی، ایس نال ایہ گل
$\lexists[\Obj v_2][\lnot \Atom{\Obj P}{\Obj v_2}]$
لئی ثابت ہوندی اے: پہلے حصے توں
پتا لگدا اے کہ ایٹمی !!{formula}~$\Atom{\Obj
P}{\Obj v_2}$ لئی ایہ سچی اے۔
فیر دوجے حصے توں
$\lnot \Atom{\Obj P}{\Obj v_2}$
لئی سچی ہوندی اے؛ تے اوسے حصے نوں
دوبارہ لا کے
$\lexists[\Obj v_2][\lnot \Atom{\Obj P}{\Obj v_2}]$
لئی وی۔ چونکہ سارے !!{formula}s
ایٹمی !!{formula}s توں استقرائی
طریقے نال بندے نیں، ایس لئی
ایہ دلیل انہاں وچوں ہر اک لئی
چلدی اے۔
\endgroup
% /OLP-0142

% OLP-0143 content/first-order-logic/introduction/satisfaction.tex
\begingroup
\def\olfilename{satisfaction}
\expandafter\def\csname ol@sectioncs\endcsname{section}

\olfileid{fol}{int}{sat}

\olsection{پورا اترنا}

جدوں سانوں پتا لگ جائے کہ !!{formula}s کیہ نیں،
تاں اسیں پہلے درجے دی منطق دیاں معنویات ول ہن
ہی ودھ سکدے آں: ایتھے بنیادی تعریف
!!a{structure} دی اے۔ ساڈی سادہ زبان لئی
!!a{structure}~$\Struct M$ دے صرف تِن حصے
نیں: اک غیر خالی سیٹ $\Domain{M}$ جیہنوں
\emph{!!{domain}} آکھدے نیں، $\Obj a$
جس شے ول~$\Struct M$ وچ اشارہ کردا اے،
تے $\Obj P$ جنہاں شےواں بارے
~$\Struct M$ وچ سچا اے۔~$\Obj a$ دی چُنی
شے نوں~$\Assign{\Obj a}{M}$ لکھدے آں،
تے جنہاں شےواں بارے $\Obj P$ سچا اے
اوہناں دے سیٹ نوں~$\Assign{\Obj P}{M}$۔
!!^a{structure}~$\Struct{M}$ ایہناں
تِن شےواں توں ہی بندی اے:
$\Domain{M}$، $\Assign{\Obj a}{M} \in
\Domain{M}$، تے $\Assign{\Obj P}{M} \subseteq \Domain{M}$۔
عمومی صورت ہور گنجل دار ہووے گی، کیوں جے اوہدے
وچ کئی !!{predicate}s تے !!{constant}s
ہون گے،
% PNB-INT-002: frozen source نے کئی تھاواں دی اہلیت مستقل نوں دتی؛ ایتھے محمول مراد اے۔
!!{predicate}s اک توں ودھ تھاواں والے وی
ہو سکدے نیں، تے !!{function}s وی ہون گے۔

ایہناں گلّاں نال اوہناں !!{formula}s لئی
پورا اترن دی تعریف دے سکدے آں جنہاں وچ
!!{variable}s نہیں۔ خیال ایہ اے کہ
تعریف استقرائی ہووے، ٹھیک اوس طرح جیویں
اسیں !!{formula}s دی تعریف کیتی اے۔
اسیں دساں گے کہ اک ایٹمی فارمولا
~$\Struct{M}$ وچ کدوں پورا اتردا اے،
تے فیر، مثال لئی، $\lnot !A$ دا
~$\Struct{M}$ وچ پورا اترنا ایس اُتے
کیویں منحصر اے کہ $!A$،
~$\Struct{M}$ وچ پورا
اتردا اے یا نہیں۔ مثال لئی، اسیں
ایویں تعریف دے سکدے آں:
\begin{enumerate}
  \item $\Atom{\Obj P}{\Obj a}$،
  ~$\Struct{M}$ وچ اُدوں تے صرف اُدوں
  پورا اتردا اے جدوں
  $\Assign{\Obj a}{M} \in \Assign{\Obj P}{M}$۔
  \item $\lnot !A$،~$\Struct{M}$ وچ
  اُدوں تے صرف اُدوں پورا اتردا اے
  جدوں $!A$،~$\Struct{M}$ وچ پورا
  نہیں اتردا۔
  \item $(!A \land !B)$،
  ~$\Struct{M}$ وچ اُدوں تے صرف
  اُدوں پورا اتردا اے جدوں
  $!A$،~$\Struct{M}$ وچ پورا
  اتردا اے، تے $!B$ وی~$\Struct{M}$
  وچ پورا اتردا اے۔
\end{enumerate}
فرض کرو $\Domain{M} = \{0, 1, 2\}$،
$\Assign{\Obj a}{M} = 1$،
تے $\Assign{\Obj P}{M} = \{1, 2\}$۔
ایس تعریف مطابق
$\Atom{\Obj P}{\Obj a}$،
~$\Struct{M}$ وچ پورا اتردا اے
(کیوں جے $\Assign{\Obj a}{M} =  1 \in \{1,2\} = \Assign{\Obj P}{M}$)۔
ایس توں ایہ وی پتا لگدا اے کہ
$\lnot \Atom{\Obj P}{\Obj a}$،
~$\Struct{M}$ وچ پورا نہیں اتردا؛
ایس دے الٹ
$\lnot\lnot \Atom{\Obj P}{\Obj
a}$ پورا اتردا اے، تے
$(\lnot \Atom{\Obj P}{\Obj a} \land \Atom{\Obj P}{\Obj a})$
پورا نہیں اتردا؛ تے ایس طرح اگے۔

اُکّھیڑ اوہدوں آندی اے جدوں اسیں
مسوّراں لئی تعریف دینا چاہندے آں:
اسیں کجھ ایہ کہنا چاہواں گے کہ
”$\lexists[\Obj
v_0][\Atom{\Obj P}{\Obj v_0}]$
اُدوں تے صرف اُدوں پورا اتردا اے
جدوں $\Atom{\Obj P}{\Obj
v_0}$ پورا اتردا اے“۔
پر !!{structure}~$\Struct{M}$
سانوں نہیں دسدی کہ !!{variable}s
نال کیہ کرنا اے۔ اسیں اصل وچ
ایہ کہنا چاہندے آں کہ
$\Atom{\Obj P}{\Obj v_0}$
\emph{$\Obj v_0$ دی کسے قدر لئی}
پورا اتردا اے۔ ایس نوں ٹھیک طور
تے بیان کرن لئی سانوں اک طریقہ
چاہیدا اے جیہڑا
~$\Domain{M}$ دے !!{element}s
صرف $\Obj a$ نوں ہی نہیں بلکہ
~$\Obj v_0$ نوں وی دے سکے۔
ایس مقصد لئی اسیں !!{variable}
\emph{قدر-بندیاں} متعارف کراندے آں۔
!!^a{variable} قدر-بندی سِرف
اک فنکشن~$s$ اے جیہڑا
!!{variable}s نوں~$\Domain{M}$
دے !!{element}s ول بھیجدا اے
(ساڈی مثال وچ
% PNB-INT-003: frozen source دے 1، 2، 3 نوں مثال دے دائرے 0، 1، 2 مطابق درست کیتا۔
$0$، $1$، یا~$2$ وچوں اک ول)۔
چونکہ سانوں پہلے توں نہیں پتا
کہ کسے !!a{formula} وچ کیہڑے
!!{variable}s آ سکدے نیں،
اسیں $s$ دی قدر-بندی دا
دائرہ انہاں وچوں کجھ !!{variable}s
تک محدود نہیں کر سکدے۔
سادہ حل ایہ اے کہ $s$
\emph{سارے} !!{variable}s
~$\Obj v_0$، $\Obj v_1$،
\dots\@ نوں قدراں دے؛
اسیں انہاں وچوں صرف اوہی
ورتاں گے جنہاں دی لوڑ اے۔

!!{formula}s دے پورا اترن نوں
صرف !!a{structure} دے لحاظ نال
تعریف کرن دی بجائے، اسیں اوہنوں
!!a{structure}~$\Struct{M}$
\emph{تے} !!a{variable}
قدر-بندی~$s$ دونہاں دے
لحاظ نال تعریف کراں گے،
تے مختصراً $\Sat{M}{!A}[s]$
لکھاں گے۔ ساڈی تعریف وچ
ہن !!{variable}s والے
ایٹمی !!{formula}s لئی
اک وادھو شقّ وی آوے گی:
\begin{enumerate}
  \item $\Sat{M}{\Atom{\Obj P}{\Obj a}}[s]$
  اُدوں تے صرف اُدوں جدوں
  $\Assign{\Obj a}{M} \in \Assign{\Obj P}{M}$۔
  \item $\Sat{M}{\Atom{\Obj P}{\Obj v_i}}[s]$
  اُدوں تے صرف اُدوں جدوں
  $s(\Obj v_i) \in \Assign{\Obj P}{M}$۔
  \item $\Sat{M}{\lnot !A}[s]$
  اُدوں تے صرف اُدوں جدوں
  $\Sat{M}{!A}[s]$ نہیں۔
  \item $\Sat{M}{(!A \land !B)}[s]$
  اُدوں تے صرف اُدوں جدوں
  $\Sat{M}{!A}[s]$ تے
  $\Sat{M}{!B}[s]$ دونیں۔
\end{enumerate}
ایس نال اک مسئلہ حل ہو گیا:
ہن اسیں دَس سکدے آں کہ
~$\Struct{M}$ وچ
$\Atom{\Obj P}{\Obj v_0}$
قدر~$s(\Obj v_0)$ لئی
کدوں پورا اتردا اے۔
پر $\lexists[\Obj v_0][\Atom{\Obj P}{\Obj
v_0}]$ دی تعریف ٹھیک بناؤن
لئی اک ہور کم کرنا لازمی اے:
اسیں چاہندے آں کہ
$\Sat{M}{\lexists[\Obj v_0][\Atom{\Obj P}{\Obj v_0}]}[s]$
اُدوں تے صرف اُدوں ہووے
جدوں $\Sat{M}{\Atom{\Obj P}{\Obj v_0}}[s']$
\emph{کسے} قدر-بندی $s'$
ولوں~$\Obj v_0$ نوں قدر دین
نال ہووے۔
پر~$\Obj v_0$ نوں دتی
قدر لازمی نہیں اوہی ہووے
جیہڑی $s(\Obj v_0)$
چُندا اے۔ ایس لئی اک
علامت متعارف کراندے آں:
جے $m \in \Domain{M}$،
تاں $\Subst{s}{m}{\Obj v_0}$
اوہ قدر-بندی اے جیہڑی
~$s$ ورگی ہی اے
(~$\Obj v_0$ توں سوا
سارے !!{variable}s لئی)،
پر $\Obj v_0$ نوں~$m$
دیندی اے۔ ہن تعریف ایویں
پوری ہوندی اے:
\begin{enumerate}\setcounter{enumi}{4}
  \item $\Sat{M}{\lexists[\Obj v_i][!A]}[s]$
  اُدوں تے صرف اُدوں جدوں
  $\Sat{M}{!A}[\Subst{s}{m}{\Obj v_i}]$
  کسے~$m \in \Domain{M}$ لئی ہووے۔
\end{enumerate}
کی ایہ ٹھیک چلدا اے؟
فرض کرو اسیں $s(\Obj v_i) = 0$
سارے~$i \in \Nat$ لئی رکھدے آں۔
$\Sat{M}{\lexists[\Obj v_0][\Atom{\Obj P}{\Obj v_0}]}[s]$
اُدوں تے صرف اُدوں
جدوں کوئی $m \in \Domain{M}$
ہووے جیہدے لئی
$\Sat{M}{\Atom{\Obj P}{\Obj
v_0}}[\Subst{s}{m}{\Obj v_0}]$۔
ایسا اک موجود اے:
اسیں $m = 1$ یا
$m = 2$ چُن سکدے آں۔
غور کرو، ایہ گل اوہدوں
وی سچی اے جدوں~$s(\Obj v_0)$
دی~$\Obj v_0$ نوں خود~$s$
ولوں دتی قدر---ایتھے
$0$---ایہ کم نہیں کردی۔
اسیں
$\Sat{M}{\Atom{\Obj P}{\Obj v_0}}[\Subst{s}{1}{\Obj
v_0}]$ حاصل کردے آں،
پر $\Sat{M}{\Atom{\Obj P}{\Obj v_0}}[s]$
نہیں۔

جے ایہ بیان تہانوں
الجھن بھریا تے بھاری
لگدا اے، تاں واقعی
ایہ ایسا اے۔ پر ایہ
وادھو پیچیدگی سارے
!!{formula}s دے پورا
اترن دی ٹھیک استقرائی
تعریف دین لئی ضروری اے،
تے معنوی تصورات دیاں
ٹھیک تعریفاں لئی وی
سانوں ایہوجی شے
چاہیدی اے۔ ہور طریقے
وی نیں، پر سب ایہنے
ہی (بے) ڈھنگے نیں۔
\endgroup
% /OLP-0143

% OLP-0144 content/first-order-logic/introduction/sentences.tex
\begingroup
\def\olfilename{sentences}
\expandafter\def\csname ol@sectioncs\endcsname{section}

\olfileid{fol}{int}{snt}

\olsection{\usetoken{P}{sentence}}

چنگا، ہن !!{structure}s تے !!{formula}s
لئی پورا اترن (یعنی ”وچ سچا ہون“)
دی تعریف دا اک خاکہ ساڈے کول اے۔
پر اوہنوں اک وادھو شے---!!a{variable}
قدر-بندی---دی لوڑ اے، جدوں کہ اسیں
!!{sentence}s دی تعریف چاہندے ساں۔
قدر-بندیاں توں چھٹکارا کیویں
پائیے، تے !!{sentence}s اصل وچ
کیہ نیں؟

رسمی منطق دے اپنے پہلے تعارف وچ
تہانوں شاید !!{variable}s تے
مسوّراں دے رشتے بارے بحث یاد
ہووے۔ مسوّر دے فوراً بعد
ہمیشہ !!a{variable} آؤندا اے؛
فیر !!{sentence} دے اوس حصے
وچ جتھے مسوّر لاگو ہوندا اے
(اوہدے ”دائرۂ اثر“ وچ)،
اوہ !!{variable} اوس مسوّر
نال ”پابند“ سمجھیا جاندا اے۔
!!{formula}s وچ ایہ ضروری
نہیں سی کہ ہر !!{variable}
لئی کوئی متعلقہ مسوّر وی ہووے؛
جیہڑے !!{variable}s دے
متعلقہ مسوّر نہ ہون، اوہ
”آزاد“ یا ”غیر پابند“ ہوندے
نیں۔ اسیں !!{sentence}s
اوہناں سارے !!{formula}s
نوں آکھاں گے جنہاں وچ
کوئی آزاد !!{variable}s
نہ ہون۔

اک واری فیر، !!a{formula}~$!A$
وچ !!a{variable} دی کسے
آمد بارے ایہ سمجھنا اوکھا
نہیں کہ اوہ کدوں پابند اے،
کیہڑا مسوّر اوہنوں پابند
کردا اے، تے کدوں آزاد اے۔
ممکن اے تسی ایس دی جانچ
دا کوئی طریقہ وی سکھیا
ہووے، شاید قوسین گن کے۔
پر سانوں ٹھیک تعریف
چاہیدی اے---تے چونکہ
اسیں !!{formula}s نوں
استقرائی طور تے تعریف
کیتا اے، ایس لئی
!!a{variable}~$x$
دیاں آزاد تے پابند
آمدواں دی تعریف
!!a{formula}~$!A$ وچ وی
استقرائی طور تے دے
سکدے آں۔ ساڈی سادہ
زبان لئی ایہ کجھ
ایس طرح ہو سکدی اے:
\begin{enumerate}
  \item جے $!A$ ایٹمی اے،
  تاں اوہدے وچ $x$ دیاں
  ساریاں آمدواں آزاد نیں
  (یعنی $x$ دی آمد
  ~$\Atom{\Obj P}{x}$ وچ آزاد اے)۔
  \item جے $!A$ دا روپ
  $\lnot !B$ اے، تاں
  $x$ دی کوئی آمد
  ~$\lnot !B$ وچ اُدوں
  تے صرف اُدوں آزاد
  اے جدوں اوہدے مطابق
  آمد~$x$،
  ~$!B$ وچ آزاد اے
  (یعنی~$
  !B$ وچ متغیراں
  دیاں آزاد آمدواں
  عین~$\lnot !B$
  وچ انہاں دے مطابق
  آمدواں نیں)۔
  \item جے $!A$ دا روپ
  $(!B \land !C)$ اے،
  تاں~$x$ دی کوئی آمد
  $(!B \land !C)$ وچ
  اُدوں تے صرف اُدوں
  آزاد اے جدوں اوہدے
  مطابق آمد~$x$،
  ~$!B$ یا~$!C$
  وچ آزاد اے۔
  \item جے $!A$ دا روپ
  $\lexists[x][!B]$ اے،
  تاں $x$ دی کوئی
  وی آمد~$!A$ وچ
  آزاد نہیں؛ جے روپ
  $\lexists[y][!B]$
  اے، جتھے $y$،
  ~$x$ توں وکھرا
  !!{variable} اے،
  تاں $x$ دی کوئی
  آمد $\lexists[y][!B]$
  وچ اُدوں تے صرف
  اُدوں آزاد اے جدوں
  اوہدے مطابق آمد~$x$،
  ~$!B$ وچ آزاد اے۔
\end{enumerate}

جدوں متغیراں دیاں
آزاد تے پابند آمدواں
دی ٹھیک تعریف ساڈے
کول ہووے، تاں اسیں
سادہ طور تے کہہ سکدے
آں: !!a{sentence} اوہ
ہر !!{formula} اے
جیہدے وچ !!{variable}s
دیاں آزاد آمدواں
نہ ہون۔
\endgroup
% /OLP-0144

% OLP-0145 content/first-order-logic/introduction/semantic-notions.tex
\begingroup
\def\olfilename{semantic-notions}
\expandafter\def\csname ol@sectioncs\endcsname{section}

\olfileid{fol}{int}{sem}

\olsection{معنوی تصورات}

اسیں اُتّے ذکر کیتا سی کہ جدوں اسیں ایہ ویکھدے آں کہ
$\Sat{M}{!A}[s]$ قائم اے یا نہیں، تاں سہولت لئی
$s$ نوں سارے !!{variable}s نوں قدراں دین دیندے آں؛
پر ورتیاں صرف اوہ قدراں جاندیاں نیں جیہڑیاں اوہ
~$!A$ وچ آؤن والے !!{variable}s نوں دیندا اے۔
اصل وچ~$!A$ دے صرف \emph{آزاد} متغیراں دیاں
قدراں اہم نیں۔ اسیں احتیاط نال کم کردے آں،
ایس لئی ایہ گل ثابت وی کراں گے۔ چونکہ
!!{sentence}s وچ کوئی آزاد متغیر نہیں ہوندا،
!!a{structure} وچ انہاں دے پورا اترن یا نہ
اترن لئی $s$ دا کوئی اثر نہیں۔ ایس لئی
جدوں~$!A$ اک !!a{sentence} ہووے، تاں
$\Sat{M}{!A}$ دی تعریف ایہ دے سکدے آں کہ
”$\Sat{M}{!A}[s]$ سارے~$s$ لئی“؛
تے ایہ اُدوں تے صرف اُدوں سچ اے جدوں
$\Sat{M}{!A}[s]$ گھٹ توں گھٹ اک~$s$ لئی ہووے۔
!!{formula}s دے پورا اترن دی قابلِ عمل تعریف
لئی سانوں !!{variable} قدر-بندیاں متعارف
کرانیاں پئیاں، پر !!{sentence}s لئی پورا
اترنا !!{variable} قدر-بندیاں توں آزاد اے۔

جدوں ”$\Sat{M}{!A}$“ دی تعریف ساڈے کول آ
جاندی اے، تاں پہلے درجے دی منطق دے
!!{sentence}s بارے ”حالت“ تے ”وچ سچا ہون“
دا مطلب وی واضح ہو جاندا اے۔ !!{sentence}s
لئی $\Sat{M}{!A}$ دی تعریف دی بنیاد اُتے
فیر اسیں منطقی اعتبار، منطقی نتیجہ تے پورا
ہون جوگتا دے بنیادی معنوی تصورات تعریف
کر سکدے آں۔ اک جملہ معتبر اے،
$\Entails !A$، جے ہر !!{structure}
اوہنوں پورا کردی اے۔ !!{sentence}s دے کسے سیٹ
توں اوہ منطقی طور تے نکلدا اے،
$\Gamma \Entails !A$، جے ہر اوہ
!!{structure} جیہڑی~$\Gamma$ دے سارے
!!{sentence}s نوں پورا کردی اے،
~$!A$ نوں وی پورا کردی ہووے۔ تے
!!{sentence}s دا اک سیٹ پورا ہون جوگا
اے جے کوئی !!{structure} اوہدے سارے
!!{sentence}s نوں اکّو ویلے پورا کردی ہووے۔

چونکہ !!{formula}s دی تعریف استقرائی اے،
تے پورا اترن دی تعریف وی !!{formula}s دی
بناوٹ اُتے استقرا نال دتی جاندی اے،
اسیں استقرا ورت کے اپنی معنویات دیاں
خاصیتاں ثابت کر سکدے آں تے تعریف شدہ
معنوی تصورات دا آپسی تعلق وی واضح
کر سکدے آں۔ اسیں ایہناں وچوں کجھ
خاصیتاں اکٹھیاں کر کے ثابت کراں گے:
کجھ ایس لئی کہ اوہ آپ دلچسپ نیں،
پر زیادہ ایس لئی کہ جدوں اسیں معنویات
تے !!{derivation} نظاماں دے رشتے دی
جانچ کراں گے، تاں ایہ کم آؤن گیاں۔
ایس کم لئی سانوں کجھ اوہ نحوی تصورات
تے عمل وی ٹھیک ٹھیک، یعنی استقرائی
طور تے، تعریف کرنے پین گے جنہاں دا
ہن تک ذکر نہیں کیتا۔
\endgroup
% /OLP-0145

% OLP-0146 content/first-order-logic/introduction/substitution.tex
\begingroup
\def\olfilename{substitution}
\expandafter\def\csname ol@sectioncs\endcsname{section}

\olfileid{fol}{int}{sub}

\olsection{بدل}

اسیں اک مثال نال ویکھاں گے کہ ایہ گلّاں
آپس وچ کیویں جڑدیاں نیں، تے نحو تے معنویات
دی بناوٹ ساڈے بعد دے ودھ پیچیدہ مطالعے
لئی بنیاد کیویں رکھدی اے۔ ساڈے
!!{derivation} نظاماں نوں سانوں
$\Atom{\Obj P}{\Obj a}$،
% PNB-INT-004: frozen source وچ Atom دا دوسرا argument quantifier دے body توں باہر رہ گیا سی۔
$\lforall[\Obj v_0][\Atom{\Obj P}{\Obj v_0}]$
توں !!{derive} کرن دینا چاہیدا اے۔
شاید اسیں ایس نوں استنباط دے اک قاعدے
وجوں وی لکھنا چاہیئے۔ پر ایہ کرن لئی
سانوں اوہ قاعدہ سب توں عمومی روپ وچ
بیان کرنا پئے گا: صرف $\Obj P$،
$\Obj a$ تے $\Obj v_0$ لئی نہیں،
بلکہ کسے وی !!{formula}~$!A$،
ٹرم~$t$ تے !!{variable}~$x$ لئی۔
(یاد رکھو، !!{constant}s ٹرم ہوندے
نیں؛ اسیں !!{constant}s تے
!!{function}s توں بنے ہور گنجل دار
ٹرماں نوں وی ویکھاں گے۔)
ایس لئی اسیں کجھ ایس طرح کہنا
چاہندے آں: ”جدوں وی تسی
$\lforall[x][!A(x)]$ نوں
!!{derive}d کر چکے ہو، تاں
~$!A(t)$ دا نتیجہ کڈھن دا جواز اے“؛
یعنی $\lforall[x]$ ہٹا کے
~$x$ دی تھاں~$t$ رکھن نال
ملن والا نتیجہ۔ پر ”$x$ دی
تھاں~$t$ رکھنا“ ٹھیک ٹھیک کیہ
ہوندا اے؟ $!A(x)$ تے
~$!A(t)$ دا آپس وچ کیہ
تعلق اے؟ کی ایہ ہر ویلے
درست رہندا اے؟

ایہ گل ٹھیک ٹھیک کرن لئی اسیں
\emph{بدل} دے عمل دی تعریف کردے آں۔
بدل اصل وچ ذرا اوکھا اے، کیوں جے
اسیں~$x$ دی ہر آمد
~$!A$ وچ صرف~$t$ نال نہیں بدل سکدے،
تے ہر~$t$ نوں ہر~$x$ دی
تھاں بدلنا جائز نہیں۔ اسیں فیر
استقرائی تعریفاں نال ایہ مسئلہ
حل کراں گے۔ پر ایہ کم ہو جان
توں بعد ”$!A(t)$ دا نتیجہ
$\lforall[x][!A(x)]$ توں کڈھو“
اک ٹھیک ٹھیک استنباطی قاعدہ
ہووے گا۔ ایس توں علاوہ اسیں
دکھا سکاں گے کہ ایہ چنگا
استنباطی قاعدہ اے، ایس معنی
وچ کہ $\lforall[x][!A(x)]$
توں~$!A(t)$ منطقی طور تے
نکلدا اے۔ پر ایہ ثابت کرن
لئی اسیں فیر اک ایہوجی گل
ثابت کرنی پئے گی جیہڑی پہلی
نظر وچ تہانوں پچھن اُتے
مجبور کر سکدی اے: ”اسیں
ایہ کیوں کر رہے آں؟“
$\lforall[x][!A(x)]$ توں
~$!A(t)$ دا منطقی نتیجہ
ہونا ایس حقیقت اُتے منحصر
اے کہ $\Sat{M}{!A(t)}$
% PNB-INT-005: اصل دی implicit context واضح؛ایہ نئے source-error دا دعویٰ نہیں۔
قائم اے یا نہیں، ایہ جائز بدل لئی،ساخت تے اصل فارمولا
اوہی رکھ کے تے باقی آزاد متغیراں دیاں قدراں نہ بدل کے،صرف
ٹرم~$t$ دی قدر اُتے
منحصر اے۔ یعنی جے
$m$ اوہ !!{element} ہووے
~$\Domain{M}$ دا جیہنوں
~$t$ اصل قدر-بندی دے لحاظ نال چُندا اے، تاں
$\Sat{M}{!A(t)}[s]$ اُدوں
تے صرف اُدوں جدوں
$\Sat{M}{!A(x)}[\Subst{s}{m}{x}]$۔
ایہ اوہدوں وی درست اے
جدوں $t$ وچ !!{variable}s
ہون، پر نتیجہ ٹھیک کیویں
بیان کرنا اے، ایس بارے
سانوں احتیاط کرنی پئے گی۔
\endgroup
% /OLP-0146

% OLP-0147 content/first-order-logic/introduction/models-theories.tex
\begingroup
\def\olfilename{models-theories}
\expandafter\def\csname ol@sectioncs\endcsname{section}

\olfileid{fol}{int}{mod}

\olsection{نمونے تے نظریے}

جدوں پہلے درجے دی منطق دی نحو تے معنویات
تعریف ہو جان، اسیں !!{structure}s تے معنوی
تصورات دیاں خاصیتاں دی جانچ شروع کر سکدے
آں۔ اسیں !!{derivation} نظام وی تعریف
کر سکدے آں تے انہاں دی پڑتال وی کر
سکدے آں۔ !!{sentence}s دے کسے سیٹ
لئی اسیں پچھ سکدے آں: کیہڑیاں
!!{structure}s ایس سیٹ دے سارے
!!{sentence}s نوں سچا کردیاں نیں؟
جے !!{sentence}s دا سیٹ~$\Gamma$
دتا ہووے، تاں اوہناں نوں پورا کرن والی
!!a{structure}~$\Struct{M}$ نوں
\emph{$\Gamma$ دا نمونہ} آکھدے آں۔
اسیں~$\Gamma$ توں شروع کر کے
اوہدے نمونے لبھن دی کوشش کر سکدے
آں---اوہ کیہو جیہے ہوندے نیں؟
اوہناں دا سائز کِنّا وڈا یا چھوٹا
ہونا لازمی اے؟ پر اسیں کسے اک
!!{structure} یا !!{structure}s دے
اک مجموعے توں شروع کر کے ایہ وی
پچھ سکدے آں: انہاں وچ کیہڑے
!!{sentence}s سچے نیں؟ کی کوئی
!!{sentence}s ایہناں !!{structure}s
دی \emph{عین شناخت} ایس معنی وچ
کردے نیں کہ اوہ تے صرف اوہی
اوہناں وچ سچے ہوندے نیں؟
ایہ سوال \emph{نمونیاں دے نظریے}
دا موضوع نیں۔ ایہی سوال
\emph{مسلّمی طریقے} دی بنیاد
وی نیں: !!{structure}s دے اک
مجموعے نوں کسے نظریے دے
مسلّماں، یعنی !!{sentence}s
دے سیٹ، نال بیان کرنا۔
ایہ ایس مشاہدے نال ممکن اے
کہ مسلّماں توں پہلے درجے دی
منطق وچ منطقی طور تے نکلن
والے عین اوہی !!{sentence}s
مسلّماں دے سارے نمونیاں وچ
سچے ہوندے نیں۔

اک بہت سادہ مثال لئی پیش ترتیب
ویکھو۔ پیش ترتیب اوہ تعلق~$R$
اے جیہڑا کسے سیٹ~$A$ اُتے
انعکاسی وی اے تے متعدی وی۔
% PNB-INT-006: غیر خالی دائرہ تے معیاری زبان دی ضمنی روایت واضح؛نواں source-error دعویٰ نہیں۔
اک غیر خالی سیٹ~$A$، جیہدے اُتے دو
تھاواں والا تعلق
$R \subseteq A \times A$
ہووے، اوہی کچھ اے جیہڑا
سانوں اک دو تھاواں والی
تعلق-علامت~$\Obj P$ والی
پہلے درجے دی زبان لئی
!!a{structure} دین لئی چاہیدا
اے: اسیں $\Domain{M} = A$
تے $\Assign{\Obj P}{M} = R$
رکھاں گے۔ چونکہ $R$ اک
پیش ترتیب اے، ایہ انعکاسی
تے متعدی اے؛ اسیں پہلے
درجے دی منطق دے
!!{sentence}s دا اک
سیٹ~$\Gamma$ لبھ سکدے
آں جیہڑا ایہ گل بیان کردا اے:
\begin{align*}
  & \lforall[\Obj v_0][\Atom{\Obj P}{\Obj v_0,\Obj v_0}]\\
  & \lforall[\Obj v_0][\lforall[\Obj v_1][\lforall[\Obj v_2][((\Atom{\Obj P}{\Obj v_0,\Obj v_1} \land \Atom{\Obj P}{\Obj v_1,\Obj v_2}) \lif \Atom{\Obj P}{\Obj v_0,\Obj v_2})]]]
\end{align*}
ایہ !!{sentence}s بس ”ہر~$x$ لئی
$Rxx$“ ($R$ انعکاسی اے)
تے ”جدوں $Rxy$ تے $Ryz$
ہون تاں $Rxz$ وی“ ($R$
متعدی اے) دیاں علامتی
لکھتاں نیں۔ اسیں ویکھدے
آں کہ !!a{structure}~$\Struct{M}$
ایہناں دو !!{sentence}s دے
سیٹ~$\Gamma$ دا نمونہ
اُدوں تے صرف اُدوں اے
جدوں $R$ (یعنی
$\Assign{\Obj P}{M}$)،
~$A$ اُتے پیش ترتیب
اے (یعنی $\Domain{M}$)۔
دوجے لفظاں وچ، $\Gamma$
دے نمونے عین پیش ترتیباں
نیں۔ ساریاں پیش ترتیباں
دی کوئی وی خاصیت جیہڑی
صرف~$\Obj P$ نوں
!!{predicate} منن والی
پہلے درجے دی زبان وچ
بیان ہو سکدی ہووے
(جیویں اُتّے انعکاسیت
تے تعدیت)، اوہ
~$\Gamma$ دے دو
!!{sentence}s توں منطقی
طور تے نکلدی اے،
تے اُلٹا وی۔ ایس لئی
اسیں~$\Gamma$ دے
نمونیاں بارے جو کجھ
ثابت کردے آں، اوہ
ساریاں پیش ترتیباں
بارے ثابت کردے آں۔

ہر خاص نظریے تے اوہدے
نمونیاں دے طبقے لئی
(جیویں $\Gamma$ تے
ساریاں پیش ترتیباں)،
ایہ دلچسپ سوال ہون
گے کہ مطابق پہلے درجے
دی زبان وچ کیہ بیان
ہو سکدا اے تے کیہ
نہیں۔ !!{structure}s
دیاں کجھ خاصیتاں
ایہوجیہیاں نیں جیہڑیاں
ہر زبان تے نمونیاں دے
ہر طبقے لئی اہم نیں،
خاص کر !!{domain}
دے سائز نال متعلق
خاصیتاں۔ ایتھے برابری والی،زیادہ توں زیادہ قابلِشمار
علامتی ذخیرے دی معیاری زبان تے اوہدیاں ذیلی زباناں مراد نیں۔ مثال لئی،
!!{domain} وچ عین
~$n$~!!{element}s
ہون دی گل کسے وی
~$n \in \PosInt$
لئی ہمیشہ بیان
کیتی جا سکدی اے۔
لامتناہی !!{sentence}s
دا اک سیٹ ورت کے
ایہ وی بیان کیتا
جا سکدا اے کہ
!!{domain} لامتناہی
اے۔ پر ایہ بیان
نہیں کیتا جا سکدا
کہ دائرہ محدود اے
یا دائرہ
!!{nonenumerable}
اے۔ پہلے درجے دیاں
زباناں دیاں حدّاں
بارے ایہ نتیجے
محدود گواہی تے
لووینہائم--سکولم
دے قضیاں توں نکلدے
نیں۔
\endgroup
% /OLP-0147

% OLP-0148 content/first-order-logic/introduction/soundness-completeness.tex
\begingroup
\def\olfilename{soundness-completeness}
\expandafter\def\csname ol@sectioncs\endcsname{section}

\olfileid{fol}{int}{scp}

\olsection{صحت مندی تے تمامیت}

اسیں پہلے درجے دی منطق لئی
!!{derivation} نظام وی متعارف
کراں گے۔ منطق داناں نے کئی
!!{derivation} نظام بنائے نیں،
پر اوہ سارے !!{sentence}s دے
درمیان اکّو !!{derivability}
تعلق تعریف کردے نیں۔ اسیں کہندے
آں کہ $\Gamma$ توں
~$!A$ \emph{!!{derive}s}
اے، $\Gamma \Proves !A$،
جے اک مخصوص، ٹھیک ٹھیک
تعریف کیتی قسم دا
!!a{derivation} موجود ہووے۔
!!^{derivation}s ہمیشہ
علامتاں دیاں محدود
بناوٹاں ہوندیاں نیں---
ہو سکدا اے !!{sentence}s
دی فہرست، یا کوئی ہور
گنجل دار ساخت۔ !!{derivation}
نظاماں دا مقصد ایہ طے
کرن دا وسیلہ دینا اے کہ
کی کوئی !!a{sentence}
کسے سیٹ~$\Gamma$ توں
منطقی طور تے نکلدا اے۔
ایس مقصد لئی ضروری اے
کہ $\Gamma \Entails !A$
اُدوں تے صرف اُدوں
ہووے جدوں $\Gamma
\Proves !A$ ہووے۔

جے $\Gamma \Proves !A$
ہووے پر $\Gamma \Entails !A$
نہ ہووے، تاں ساڈا
!!{derivation} نظام حد توں
ودھ طاقت ور ہووے گا:
اوہ بہت کجھ ثابت کر دے
گا۔ ایہ خاصیت کہ جے
$\Gamma
\Proves !A$ تاں
$\Gamma \Entails !A$
وی، \emph{صحت مندی}
کہلاندی اے، تے چنگے
!!{derivation} نظام لئی
ایہ گھٹ توں گھٹ شرط
اے۔ دوجے پاسے، جے
$\Gamma \Entails !A$
ہووے پر $\Gamma \Proves !A$
نہ ہووے، تاں ساڈا
!!{derivation} نظام حد توں
گھٹ طاقت ور ہووے گا:
اوہ کافی کچھ ثابت نہیں
کردا۔ ایہ خاصیت کہ جے
$\Gamma \Entails !A$
تاں $\Gamma \Proves !A$
وی، \emph{تمامیت}
کہلاندی اے۔ صحت مندی
دا ثبوت عموماً نسبتاً
آسان اے (!!{derivation}s
دی بناوٹ اُتے استقرا نال،
کیوں جے انہاں دی تعریف
استقرائی اے)۔ تمامیت دا
ثبوت اوکھا اے۔

صحت مندی تے تمامیت دے
کئی اہم نتیجے نیں۔ جے
!!{sentence}s دا کوئی
سیٹ~$\Gamma$ تضاد
!!{derive}s کرے (مثلاً
$!A \land \lnot !A$)،
تاں اوہ \emph{متضاد}
کہلاندا اے۔ متضاد
$\Gamma$s دے کوئی
نمونے نہیں ہو سکدے؛
اوہ پورا ہون جوگ
نہیں۔ تمامیت توں
اُلٹی گل نکلدی اے:
ہر اوہ $\Gamma$
جیہڑا متضاد نہیں---
یا جیویں اسیں آکھاں
گے، \emph{بے تضاد}---
اوہدا اک نمونہ ہوندا
اے۔ اصل وچ ایہ تمامیت
دے ہم ارز اے، تے
اسی روپ وچ اسیں
تمامیت ثابت کراں گے۔
ایہ ڈونگھا، شاید
حیران کن نتیجہ اے:
صرف ایس گل توں کہ
تسی $!A \land \lnot !A$
نوں $\Gamma$ توں
ثابت نہیں کر سکدے،
ایہ یقینی بن جاندا
اے کہ کوئی
!!a{structure} موجود
اے جیہڑی~$\Gamma$
دے بیان مطابق اے۔
ایس لئی تمامیت
ایس سوال دا جواب
دیندی اے: !!{sentence}s
دے کیہڑے سیٹاں دے
نمونے ہوندے نیں؟
جواب: اوہ سارے
تے صرف اوہی سیٹ
جیہڑے بے تضاد نیں۔

صحت مندی تے تمامیت
دے قضیاں دے دو
اہم نتیجے نیں: محدود
گواہی تے لووینہائم--سکولم
دا قضیہ۔ نمونیاں دے
نظریے وچ ایہ اہم
نتیجے نیں تے انہاں
نال کئی دلچسپ نتیجے
قائم کیتے جا سکدے
نیں۔ اسیں دو دا
ذکر پہلے کر چکے آں:
پہلے درجے دی منطق
ایہ بیان نہیں کر
سکدی کہ کسے
!!a{structure} دا
!!{domain} محدود
اے، یا اوہدا دائرہ
!!{nonenumerable}
اے۔

تاریخی طور تے ایہ
سب کجھ---پہلے درجے
دی منطق دی نحو تے
معنویات دی تعریف،
چنگے !!{derivation}
نظاماں دی تعریف،
انہاں دی صحت مندی
تے تمامیت دا ثبوت،
تے پہلے درجے دیاں
زباناں وچ کیہ بیان
ہو سکدا اے تے کیہ
نہیں، ایس دی وضاحت---
سمجھ کے درست کرن
وچ لمّا سماں لگیا۔
ہن سانوں ایہ کرنا
آؤندا اے، پر ساریاں
تفصیلاں وچوں لنگھنا
اج وی الجھن بھریا
تے تھکا دینے والا
ہو سکدا اے۔ پر
ایہ اہم وی اے،
کیوں جے پہلے درجے
دی منطق دی رسمی
زبان لئی ایتھے
بنائے طریقے منطق،
کمپیوٹر سائنس
تے لسانیات وچ
ہر پاسے ورتے
جاندے نیں۔ ایس لئی
تفصیلاں اُتے محنت
دا فائدہ لمبے
سمے وچ ملدا اے۔
\endgroup
% /OLP-0148

\OLEndChapterHook
\endgroup
% /OLP-0139

% OLP-0149 content/first-order-logic/syntax-and-semantics/syntax.tex
\begingroup
\def\olfilename{syntax}
\expandafter\def\csname ol@sectioncs\endcsname{section}

\olchapter{fol}{syn}{پہلے درجے دی منطق دی نحو}

% OLP-0150 content/first-order-logic/syntax-and-semantics/intro-syntax.tex
\begingroup
\def\olfilename{intro-syntax}
\expandafter\def\csname ol@sectioncs\endcsname{section}

\olfileid{fol}{syn}{itx}

\olsection{تعارف}

پہلے درجے دی منطق دا نظریہ تے
ماورائے نظریہ ودھاؤن لئی سانوں
پہلاں اوہدی لکھتاں دی نحو تے
معنویات تعریف کرنیاں پین
گیاں۔ پہلے درجے دی منطق
دیاں لکھتاں ٹرم تے
!!{formula}s ہوندیاں نیں۔
ٹرم !!{variable}s،
!!{constant}s تے
!!{function}s توں
بندے نیں۔ دوجے پاسے،
!!^{formula}s،
!!{predicate}s نوں
ٹرماں نال ملا کے
بندے نیں (ایہ سب
توں چھوٹے، ”ایٹمی“
!!{formula}s ہوندے
نیں)، تے فیر ایٹمی
!!{formula}s توں
منطقی رابطے تے
مسوّر ورت کے ہور
گنجل دار فارمولے
بنا سکدے آں۔
تشکیل دے قاعدے
بیان کرن دے کئی
طریقے نیں؛ اسیں
صرف اک ممکن
طریقہ دیندے آں۔
ہور نظام وکھریاں
علامتاں چُنن گے،
وکھرے منطقی
رابطیاں نوں
ابتدائی منّن
گے، تے قوسین
وکھرے ڈھنگ
نال ورتن گے
(یا بالکل نہیں
ورتن گے، جیویں
پولش علامت
نگاری وچ)۔
پر سارے طریقیاں
وچ سانجھی گل
ایہ اے کہ تشکیل
دے قاعدے ٹرماں
تے !!{formula}s
دے سیٹ دی تعریف
\emph{استقرائی طور تے}
کردے نیں۔ جے
ایہ کم ٹھیک کیتا
جائے، تاں ہر لکھت
تشکیل دے قاعدیاں
مطابق بنیادی طور
تے صرف اک ہی
طریقے نال بن
سکدی اے۔ لکھتاں
نوں \emph{اکّو
طرح نال پڑھن
جوگ} بناؤن والی
استقرائی تعریف
دا مطلب ایہ اے
کہ اسیں اوہناں
لکھتاں نوں
معنی وی اوسے
طریقے نال، یعنی
استقرائی تعریف
نال، دے سکدے آں۔
\endgroup
% /OLP-0150

% OLP-0151 content/first-order-logic/syntax-and-semantics/first-order-languages.tex
\begingroup
\def\olfilename{first-order-languages}
\expandafter\def\csname ol@sectioncs\endcsname{section}

\olfileid{fol}{syn}{fol}

\olsection{پہلے درجے دیاں زباناں}


پہلے درجے دی منطق دیاں لکھتاں اک بنیادی علامتی ذخیرے توں بندیاں نیں،
جیہدے وچ \emph{!!{variable}s}، \emph{!!{constant}s}،
\emph{!!{predicate}s} تے کدی کدی \emph{!!{function}s} ہوندے نیں۔
ایہناں نال منطقی رابطے، مسوّر تے قوساں تے وِرمیاں ورگیاں
رموزی علامتاں ملا کے \emph{ٹرم} تے \emph{!!{formula}s}
بنائے جاندے نیں۔

\begin{explain}
غیر رسمی طور تے !!{predicate}s خاصیتاں تے تعلقاں دے ناں نیں،
!!{constant}s وکھ وکھ شےواں دے ناں نیں، تے !!{function}s
نقشیاں دے ناں نیں۔ !!{identity}~$\eq$ نوں چھڈ کے ایہ
\emph{غیر منطقی علامتاں} نیں تے مل کے اک زبان بناؤندیاں نیں۔
پہلے درجے دی ہر زبان~$\Lang L$ اوہدیاں غیر منطقی
علامتاں توں متعین ہوندی اے۔ سب توں عام صورت وچ
$\Lang L$ وچ ہر قسم دیاں بے انت علامتاں ہوندیاں نیں۔
\end{explain}

عام صورت وچ اسیں پہلے درجے دی منطق وچ ایہ علامتاں ورتدے آں:

\begin{enumerate}
\item منطقی علامتاں
\begin{enumerate}
\item منطقی رابطے:
  \startycommalist
  \iftag{prvNot}{\ycomma $\lnot$ (نفی)}{}%
  \iftag{prvAnd}{\ycomma $\land$ (عطف)}{}%
  \iftag{prvOr}{\ycomma $\lor$ (فصل)}{}%
  \iftag{prvIf}{\ycomma $\lif$ (!!{conditional})}{}%
  \iftag{prvIff}{\ycomma $\liff$ (!!{biconditional})}{}%
  \iftag{prvAll}{\ycomma $\lforall$ (کُلّی مسوّر)}{}%
  \iftag{prvEx}{\ycomma $\lexists$ (وجودی مسوّر)}{}.
\tagitem{prvFalse}{!!{falsity} لئی قضیاتی مستقل~$\lfalse$۔}{}
\tagitem{prvTrue}{!!{truth} لئی قضیاتی مستقل~$\ltrue$۔}{}
\item دو-مقامی !!{identity}~$\eq$۔
\item !!{variable}s دا اک !!{denumerable}s سیٹ: $\Obj v_0$, $\Obj v_1$, $\Obj
  v_2$, \dots
\end{enumerate}
\item غیر منطقی علامتاں، جیہڑیاں پہلے درجے دی منطق دی
\emph{معیاری زبان} بناؤندیاں نیں
\begin{enumerate}
\item $n$-مقامی !!{predicate}s دا ہر $n>0$ لئی اک !!{denumerable}s سیٹ: $\Obj
  A^n_0$, $\Obj A^n_1$, $\Obj A^n_2$, \dots
\item !!{constant}s دا اک !!{denumerable}s سیٹ: $\Obj c_0$, $\Obj c_1$, $\Obj
  c_2$, \dots.
\item $n$-مقامی !!{function}s دا ہر $n>0$ لئی اک !!{denumerable}s سیٹ:
  $\Obj f^n_0$, $\Obj f^n_1$, $\Obj f^n_2$, \dots
\end{enumerate}
\item رموزِ اوقاف: (, )، تے وِرما۔
\end{enumerate}

ساڈیاں زیادہ تر تعریفاں تے نتیجے پہلے درجے دی منطق دی پوری
معیاری زبان لئی بیان کیتے جان گے۔ پر ورتوں دے مطابق اسیں
زبان نوں صرف کجھ !!{predicate}s، !!{constant}s تے
!!{function}s تک محدود وی کر سکدے آں۔

\begin{ex}
حساب دی زبان~$\Lang L_A$ وچ اک دو-مقامی !!{predicate}~$<$،
اک !!{constant}~$\Obj 0$، اک یک-مقامی !!{function}~$\prime$،
تے دو دو-مقامی !!{function}s~$+$ تے~$\times$ ہوندے نیں۔
\end{ex}

\begin{ex}
سیٹ نظریے دی زبان~$\Lang L_Z$ وچ صرف اک دو-مقامی
!!{predicate}~$\in$ ہوندا اے۔
\end{ex}

\begin{ex}
ترتیباں دی زبان~$\Lang L_\le$ وچ صرف اک دو-مقامی
!!{predicate}~$\le$ ہوندا اے۔
\end{ex}

فیر وی ایہ رواجی علامتاں نیں: باقاعدہ طور تے ایہ صرف قائم مقام ناں نیں؛
مثلاً $<$، $\in$ تے $\le$، $\Obj A^2_0$ دے قائم مقام ناں نیں؛
$\Obj 0$، $\Obj c_0$ دا؛ $\prime$، $\Obj f^1_0$ دا؛
$+$، $\Obj f^2_0$ دا؛ تے $\times$، $\Obj f^2_1$ دا۔

\iftag{defNot,defOr,defAnd,defIf,defIff,defTrue,defFalse,defEx,defAll}{%
اُتے متعارف کرائے ابتدائی رابطیاں تے
\iftag{notprvEx,notprvAll}{مسوّر}{مسوّراں} دے نال نال،
اسیں ایہ \emph{تعریف شدہ} علامتاں وی ورتدے آں:
\startycommalist
  \iftag{defNot}{\ycomma $\lnot$ (نفی)}{}%
  \iftag{defAnd}{\ycomma $\land$ (عطف)}{}%
  \iftag{defOr}{\ycomma $\lor$ (فصل)}{}%
  \iftag{defIf}{\ycomma $\lif$ (!!{conditional})}{}%
  \iftag{defIff}{\ycomma $\liff$ (!!{biconditional})}{}%
  \iftag{defAll}{\ycomma $\lforall$ (کُلّی مسوّر)}{}%
  \iftag{defEx}{\ycomma $\lexists$ (وجودی مسوّر)}{}%
  \iftag{defFalse}{\ycomma !!{falsity}~$\lfalse$}{}%
  \iftag{defTrue}{\ycomma !!{truth}~$\ltrue$}}{}.

\begin{tagblock}{defNot,defOr,defAnd,defIf,defIff,defTrue,defFalse,defEx,defAll}
\begin{explain}
تعریف شدہ علامت باقاعدہ طور تے زبان دا حصہ نہیں ہوندی،
بلکہ غیر رسمی مخفف دے طور تے متعارف ہوندی اے: ایہ
اوہناں فارمولیاں نوں مختصر لکھن دی سہولت دیندی اے
جیہڑے صرف ابتدائی علامتاں نال کافی لمّے ہو جاندے۔
ایہ اک صاف فائدہ اے۔ پر اس توں وڈا فائدہ ایہ اے کہ ثبوت
چھوٹے ہو جاندے نیں۔ جے کوئی علامت ابتدائی ہووے، تاں
ثبوتاں وچ اوہدا وکھرا معاملہ کرنا پیندا اے۔ اس لئی
ابتدائی علامتاں جِنیاں ودھ ہون، ساڈے ثبوت اُنّے لمّے ہوندے نیں۔
\end{explain}
\end{tagblock}

% Alternate symbols

\begin{intro}
ممکن اے تہانوں اُتے ورتی اصطلاحات تے علامتاں توں
وکھریاں اصطلاحات تے علامتاں دا پتا ہووے۔ منطق دیاں کتاباں
(تے استاد) ”نفی“ لئی عام طور تے $\sim$، $\neg$ یا~!\
تے ”عطف“ لئی $\wedge$، $\cdot$ یا $\&$ ورتدے نیں۔
”شرطیہ“ یا ”استلزام“ لئی عام علامتاں $\rightarrow$،
$\Rightarrow$ تے $\supset$ نیں۔
\iftag{prvIff,defIff}{”دو طرفہ شرطیہ“، ”دو طرفہ استلزام“
  یا ”(مادی) ہم ارزی“ لئی $\leftrightarrow$،
  $\Leftrightarrow$ تے $\equiv$ علامتاں نیں۔}{}
\iftag{prvFalse,defFalse}{$\lfalse$ دی علامت نوں ”کذب“،
  ”فالسَم“، ”محال“ یا ”تھلّہ“ وی آکھیا جاندا اے۔}{}
\iftag{prvTrue,defTrue}{$\ltrue$ دی علامت نوں ”صدق“،
  ”وِیرَم“ یا ”چوٹی“ وی آکھیا جاندا اے۔}{}

رواج اے کہ !!{constant}s (جیہناں نوں کدی ناں وی آکھدے نیں)
لئی لاطینی حروفِ تہجی دے مُڈھ دے چھوٹے حرف (مثلاً $a$، $b$، $c$)
تے !!{variable}s لئی آخر دے چھوٹے حرف (مثلاً $x$، $y$، $z$)
ورتے جان۔ مسوّر !!{variable}s نال ملدے نیں، مثلاً $x$؛
علامت نگاری دیاں وکھریاں صورتاں وچ کُلّی مسوّر لئی
$\forall x$، $(\forall x)$، $(x)$، $\Pi x$، $\bigwedge_x$
تے وجودی مسوّر لئی $\exists x$، $(\exists
x)$، $(Ex)$، $\Sigma x$، $\bigvee_x$ شامل نیں۔
\end{intro}

\begin{explain}
اسیں سارے قضیاتی عملگر تے دونیں مسوّر زبان دیاں ابتدائی
علامتاں من سکدے آں۔ یا اسیں ابتدائی علامتاں دا چھوٹا
ذخیرہ چُن کے باقی !!{operator}s نوں تعریف شدہ من سکدے آں۔
بولیائی عملگراں دے ”صداقتی-فنکشنی طور تے مکمل“ سیٹاں وچ
$\{ \lnot, \lor \}$، $\{ \lnot, \land \}$ تے $\{ \lnot,
\lif\}$ شامل نیں—ایہناں وچوں ہر سیٹ نوں دوہاں وچوں کسی وی
اک مسوّر نال ملا کے اظہار دے لحاظ نال مکمل پہلے درجے دی
زبان بنائی جا سکدی اے۔

تہانوں دو ہور !!{operator}s دا پتا ہو سکدا اے: شیفر دی لکیر~$|$
(ہنری شیفر دے ناں اُتے)، تے پیئرس دا تیر~$\downarrow$،
جیہنوں کوائن دا خنجر وی آکھدے نیں۔ جدوں انہاں نوں بالترتیب
”نینڈ (nand)“ تے ”نور (nor)“ دے عام معنیاں وچ لیا جائے،
تاں ایہ عملگر اکیلے اکیلے صداقتی-فنکشنی طور تے مکمل نیں۔
\end{explain}
\endgroup
% /OLP-0151

% OLP-0152 content/first-order-logic/syntax-and-semantics/terms-formulas.tex
\begingroup
\def\olfilename{terms-formulas}
\expandafter\def\csname ol@sectioncs\endcsname{section}

\olfileid{fol}{syn}{frm}

\olsection{ٹرم تے \printtoken{P}{formula}}

جدوں پہلے درجے دی اک زبان~$\Lang L$ دتی جاوے، تاں اسیں
زبان~$\Lang L$ دی بنیادی علامتاں توں بندیاں لکھتاں دی تعریف کر سکدے آں۔
ایہناں وچ خاص طور تے \emph{ٹرم} تے \emph{!!{formula}s} شامل نیں۔

\begin{defn}[ٹرم]
\ollabel{defn:terms}
\emph{ٹرماں} دا سیٹ~$\Trm[L]$، $\Lang L$ لئی، اینج
استقرائی طور تے تعریف کیتا جاندا اے:
\begin{enumerate}
\item ہر !!{variable} اک ٹرم اے۔
\item $\Lang L$ دا ہر !!{constant} اک ٹرم اے۔
\item جے $f$ اک $n$-مقامی !!{function} اے تے $t_1$, \dots, $t_n$
  ٹرم نیں، تاں $\Atom{f}{t_1, \ldots, t_n}$ وی اک ٹرم اے۔
\tagitem{limitClause}{
  ہور کوئی شے ٹرم نہیں اے۔}{}
\end{enumerate}
جیہڑے ٹرم وچ کوئی !!{variable}s نہ ہون، اوہ \emph{بند ٹرم} اے۔
\end{defn}

\begin{explain}
ساڈی زبان تے ٹرماں دی وضاحت وچ !!{constant}s علامتاں دی
وکھری قسم دے طور تے آؤندے نیں، پر اسیں انہاں نوں
صفر-مقامی !!{function}s وچ وی شامل کر سکدے ساں۔ فیر
ٹرماں دی تعریف دی دوجی شق دی لوڑ نہ رہندی۔ بس
$\Atom{f}{t_1, \ldots, t_n}$ نوں اکیلا $f$
سمجھنا پیندا اے، جے $n =
0$ ہووے۔
\end{explain}

\begin{defn}[فارمولے]
\ollabel{defn:formulas}
\emph{!!{formula}s} دا سیٹ~$\Frm[L]$، زبان~$\Lang L$ لئی،
اینج استقرائی طور تے تعریف کیتا جاندا اے:
\begin{enumerate}
\tagitem{prvFalse}{$\lfalse$ اک ایٹمی !!{formula} اے۔}{}

\tagitem{prvTrue}{$\ltrue$ اک ایٹمی !!{formula} اے۔}{}

\item جے $R$ اک $n$-مقامی !!{predicate} اے جو~$\Lang L$ دا اے، تے $t_1$, \dots,
  $t_n$، $\Lang L$ دے ٹرم نیں، تاں $\Atom{R}{t_1,\ldots, t_n}$ اک
  ایٹمی !!{formula} اے۔

\item جے $t_1$ تے $t_2$، $\Lang L$ دے ٹرم نیں، تاں $\Atom{\eq}{t_1, t_2}$
  اک ایٹمی !!{formula} اے۔

\tagitem{prvNot}{جے $!A$، !!a{formula} اے، تاں $\lnot !A$ وی
  !!a{formula} اے۔}{}

\tagitem{prvAnd}{جے $!A$ تے $!B$، !!{formula}s نیں، تاں $(!A \land
  !B)$ وی !!a{formula} اے۔}{}

\tagitem{prvOr}{جے $!A$ تے $!B$، !!{formula}s نیں، تاں $(!A \lor !B)$
  وی !!a{formula} اے۔}{}

\tagitem{prvIf}{جے $!A$ تے $!B$، !!{formula}s نیں، تاں $(!A \lif !B)$
  وی !!a{formula} اے۔}{}

\tagitem{prvIff}{جے $!A$ تے $!B$، !!{formula}s نیں، تاں $(!A \liff !B)$
  وی !!a{formula} اے۔}{}

\tagitem{prvAll}{جے $!A$، !!a{formula} اے تے $x$، !!a{variable} اے،
  تاں $\lforall[x][!A]$ وی !!a{formula} اے۔}{}

\tagitem{prvEx}{جے $!A$، !!a{formula} اے تے $x$، !!a{variable} اے،
  تاں $\lexists[x][!A]$ وی !!a{formula} اے۔}{}

\tagitem{limitClause}{ہور کوئی شے !!a{formula} نہیں اے۔}{}
\end{enumerate}
\end{defn}

\begin{explain}
ٹرماں دے سیٹ تے !!{formula}s دے سیٹ دیاں تعریفاں
\emph{استقرائی تعریفاں} نیں۔ بنیادی طور تے اسیں
!!{formula}s دا سیٹ لامتناہی مرحلیاں وچ بناؤندے آں۔
پہلے مرحلے وچ اسیں سارے ایٹمی فارمولیاں نوں فارمولے
قرار دیندے آں؛ ایہ تعریف دیاں پہلیاں کجھ صورتاں، یعنی
\iftag{prvTrue}{$\ltrue$, }{}%
\iftag{prvFalse}{$\lfalse$, }{}%
$\Atom{R}{t_1,\dots,t_n}$ تے $\Atom{\eq}{t_1,t_2}$
والیاں صورتاں دے مطابق اے۔ ایس لئی ”ایٹمی
!!{formula}“ دا مطلب ایس روپ دا کوئی وی !!{formula} اے۔

تعریف دیاں ہور صورتاں پہلے بنائے !!{formula}s توں
نویں !!{formula}s بناؤن دے قاعدے دیندیاں نیں۔
دوجے مرحلے وچ اسیں اوہناں نال ایٹمی !!{formula}s توں
!!{formula}s بنا سکدے آں۔ تیجے مرحلے وچ اسیں ایٹمی
فارمولیاں تے دوجے مرحلے وچ ملے فارمولیاں توں نویں
فارمولے بناؤندے آں، تے ایویں اگے۔ !!{formula} صرف
اوہ شے اے جیہڑی آخرکار کسے ایہو جہے مرحلے اُتے
بنائی جاوے؛ ہور کوئی شے نہیں۔
\end{explain}

رواجاً اسیں $\eq$ نوں اوہدیاں دو دلیلاں دے وچ لکھدے تے
قوساں چھڈ دیندے آں: $\eq[t_1][t_2]$،
$\Atom{\eq}{t_1,t_2}$ دا مخفف اے۔ اس توں علاوہ
$\lnot \Atom{\eq}{t_1,t_2}$ نوں $\eq/[t_1][t_2]$
لکھ کے مختصر کردے آں۔ جدوں فارمولا $(!B \ast !C)$،
جیہڑا $!B$ تے $!C$ توں دو-مقامی رابطے~$\ast$
نال بنیا اے، لکھدے آں تاں اکثر سب توں باہری قوساں دی جوڑی
چھڈ کے صرف~$!B \ast !C$ لکھدے آں۔

\begin{intro}
منطق دیاں کجھ کتاباں وچ !!{variable}~$x$ دا~$!A$
وچ آنا لازمی منیا جاندا اے، تاں جو
\iftag{prvEx}{$\lexists[x][!A]$ }{}%
\iftag{notprvEx,notprvAll}{}{تے }%
\iftag{prvAll}{$\lforall[x][!A]$ }{}%
نوں
\iftag{notprvEx,notprvAll}{!!a{formula}}{!!{formula}s}
منیا جائے۔ ایہ شرط نہ رکھن نال کوئی خرابی نہیں آؤندی،
تے کم آسان ہو جاندا اے۔
\end{intro}

\begin{tagblock}{defNot,defOr,defAnd,defIf,defIff,defTrue,defFalse,defEx,defAll}
\begin{defn}
تعریف شدہ عملگراں نال بنے فارمولیاں نوں اینج سمجھیا جاوے:

\begin{tagenumerate}{defTrue,defFalse,defNot,defOr,defAnd,defIf,defIff,defEx,defAll}

\tagitem{defTrue}{$\ltrue$،
  \iftag{prvFalse}{$\lnot\lfalse$}{$(!A \lor \lnot !A)$، کسے
    مقرر ایٹمی !!{formula}~$!A$ لئی} دا مخفف اے۔}{}

\tagitem{defFalse}{$\lfalse$،
  \iftag{prvTrue}{$\lnot\ltrue$}{$(!A \land \lnot !A)$، کسے
    مقرر ایٹمی !!{formula}~$!A$ لئی} دا مخفف اے۔}{}

\tagitem{defNot}{$\lnot !A$، $!A \lif \lfalse$ دا مخفف اے۔}{}

\tagitem{defOr}{$!A \lor !B$،
  \iftag{prvAnd}{$\lnot(\lnot !A \land \lnot !B)$}{$\lnot !A \lif
    !B$} دا مخفف اے۔}{}

\tagitem{defAnd}{$!A \land !B$،
  \iftag{prvOr}{$\lnot(\lnot !A \lor \lnot !B)$}{$\lnot (!A \lif
    \lnot !B)$} دا مخفف اے۔}{}

% PNB-SYN-001 ماخذ دی درستی: defIf دی پہلی متبادل عبارت وچ بند ہُندی قوس توں پہلاں کھلدی قوس ماخذ وچ رہ گئی؛ ایتھے پچھلے PNBPL-002 وانگوں جوڑی پوری کیتی اے۔
\tagitem{defIf}{$!A \lif !B$،
  \iftag{prvOr}{$(\lnot !A \lor !B)$}{$\lnot (!A \land \lnot !B)$} دا مخفف اے۔}{}

\tagitem{defIff}{$!A \liff !B$، $(!A \lif !B) \land (!B
  \lif !A)$ دا مخفف اے۔}{}

\tagitem{defAll}{$\lforall[x][!A]$، $\lnot\lexists[x][\lnot !A]$ دا مخفف اے۔}{}

\tagitem{defEx}{$\lexists[x][!A]$، $\lnot\lforall[x][\lnot !A]$ دا مخفف اے۔}{}
\end{tagenumerate}
\end{defn}
\end{tagblock}

جے اسیں کسے خاص ورتوں دی زبان وچ کم کریے، تاں دو-مقامی
!!{predicate}s تے !!{function}s نوں اکثر اوہناں دے متعلقہ
ٹرماں دے وچ لکھدے آں، مثلاً حساب دی زبان وچ $t_1 < t_2$
تے $(t_1 + t_2)$، تے سیٹ نظریے دی زبان وچ $t_1 \in t_2$۔
حساب دی زبان دا جانشین فنکشن رواجاً اپنی دلیل دے
\emph{بعد} وی لکھیا جاندا اے:~$t'$۔ پر باقاعدہ طور تے
ایہ صرف $\Atom{\Obj A^2_0}{t_1, t_2}$،
$\Obj f^2_0(t_1, t_2)$، $\Atom{\Obj A^2_0}{t_1, t_2}$
تے $f^1_0(t)$ دے رواجی مخفف نیں، بالترتیب۔

\begin{defn}[نحوی یکسانیت]
علامت $\ident$، علامتاں دیاں لڑیاں وچ نحوی یکسانیت
ظاہر کردی اے، یعنی $!A \ident !B$ اُدوں تے صرف
اُدوں جدوں $!A$ تے $!B$ اکّو لمبائی دیاں علامتی
لڑیاں ہون تے ہر تھاں اُتے اکّو علامت رکھدیاں ہون۔
\end{defn}

علامت $\ident$ دے دونیں پاسے جوڑ نال ملیاں لڑیاں
آ سکدیاں نیں؛ مثلاً $!A \ident (!B \lor !C)$ دا
مطلب اے کہ علامتاں دی لڑی~$!A$ اوہی لڑی اے
جیہڑی اک کھلدی قوس، لڑی $!B$، علامت $\lor$،
لڑی~$!C$ تے اک بند ہُندی قوس نوں ایس ترتیب نال
جوڑ کے ملدی اے۔ جے ایہ گل ہووے، تاں اسیں جاندے
آں کہ~$!A$ دی پہلی علامت کھلدی قوس اے، $!A$
وچ $!B$ اک ذیلی لڑی دے طور تے شامل اے (دوجی
علامت توں شروع ہو کے)، اس ذیلی لڑی دے بعد~$\lor$
آؤندا اے، وغیرہ۔

جیویں ٹرم تے !!{formula}s بنیادی عنصراں توں
استقرائی تعریفاں راہیں بندے نیں، اسیں انہاں بارے
گلاں ثابت کرن لئی ایہ استقرائی اصول ورت سکدے آں۔

\begin{lem}[\emph{ٹرماں اُتے استقرا دا اصول}]
    \ollabel{lem:trmind}
    فرض کرو $\Lang L$ پہلے درجے دی زبان اے۔
    جے کوئی خاصیت~$P$ اینج ہووے کہ
    %
    \begin{enumerate}
        \item اوہ ہر !!{variable}~$v$ لئی درست ہووے،
        %
        \item اوہ ہر !!{constant}~$a$ لئی درست ہووے جو~$\Lang L$ دا اے، تے
        %
        \item اوہ $f(t_1,\dotsc,t_n)$ لئی وی درست ہووے جدوں
            اوہ $t_1$,~\dots, $t_n$ لئی درست ہووے تے $f$
            اک $n$-مقامی !!{function} ہووے جو~$\Lang L$ دا اے
    \end{enumerate}
    (ایہ من کے کہ $t_1$,~\dots, $t_n$، $\Lang{L}$ دے ٹرم نیں)،
    تاں $P$، $\Trm[L]$ دے ہر ٹرم لئی درست اے۔
\end{lem}

\begin{prob}
    \olref[fol][syn][frm]{lem:trmind} ثابت کرو۔
\end{prob}

\begin{lem}[\emph{!!{formula}s اُتے استقرا دا اصول}]
    \ollabel{thm:frmind}
    فرض کرو $\Lang L$ پہلے درجے دی زبان اے۔
    جے کوئی خاصیت~$P$ سارے ایٹمی !!{formula}s لئی
    درست ہووے، تے اینج وی ہووے کہ
    %
    \begin{enumerate}
        \tagitem{prvNot}{اوہ $\lnot !A$ لئی وی درست ہووے
            جدوں اوہ~$!A$ لئی درست ہووے؛}{}
        \tagitem{prvAnd}{اوہ $(!A \land !B)$ لئی وی درست ہووے
            جدوں اوہ $!A$ تے~$!B$ لئی درست ہووے؛}{}
        \tagitem{prvOr}{اوہ $(!A \lor !B)$ لئی وی درست ہووے
            جدوں اوہ $!A$ تے~$!B$ لئی درست ہووے؛}{}
        \tagitem{prvIf}{اوہ $(!A \lif !B)$ لئی وی درست ہووے
            جدوں اوہ $!A$ تے~$!B$ لئی درست ہووے؛}{}
        \tagitem{prvIff}{اوہ $(!A \liff !B)$ لئی وی درست ہووے
            جدوں اوہ $!A$ تے~$!B$ لئی درست ہووے؛}{}
        \tagitem{prvEx}{اوہ $\lexists[x][!A]$ لئی وی درست ہووے
            جدوں اوہ~$!A$ لئی درست ہووے؛}{}
        \tagitem{prvAll}{اوہ $\lforall[x][!A]$ لئی وی درست ہووے
            جدوں اوہ~$!A$ لئی درست ہووے؛}{}
  \end{enumerate}
  (ایہ من کے کہ $!A$ تے $!B$، $\Lang{L}$ دے !!{formula}s نیں)،
  تاں $P$، $\Frm[L]$ دے سارے فارمولیاں لئی درست اے۔
\end{lem}
\endgroup
% /OLP-0152

% OLP-0153 content/first-order-logic/syntax-and-semantics/unique-readability.tex
\begingroup
\def\olfilename{unique-readability}
\expandafter\def\csname ol@sectioncs\endcsname{section}

\olfileid{fol}{syn}{unq}

\olsection{اکّو طرح نال پڑھت}

\begin{explain}
جس طرح اسیں !!{formula}s دی تعریف کیتی اے، اوہ ضمانت
دیندا اے کہ ہر !!{formula} دی \emph{اکّو پڑھت} ہووے؛
یعنی !!{formula}s دے تشکیل قاعدیاں مطابق اوہنوں بناؤن
دا بنیادی طور تے صرف اک طریقہ اے، تے اوہدی نحوی ”تعبیر“
دا وی صرف اک طریقہ اے۔ جے ایہ نہ ہوندا، تاں مبہم
!!{formula}s ہوندے، یعنی اوہ !!{formula}s جیہناں
دیاں اک توں ودھ پڑھتاں یا نحوی تعبیراں ہوندیاں—ایہ صاف
اے کہ اسیں اس توں بچنا چاہندے آں۔ پر اس توں وی
اہم گل ایہ اے کہ ایس خاصیت دے بغیر اگلیاں
زیادہ تر تعریفاں تے ثبوت کم نہیں کرن گے۔

ایہ گل واضح کرن دا شاید سب توں چنگا طریقہ ایہ ویکھنا
اے کہ جے اسیں !!{formula}s دے تشکیل لئی خراب قاعدے
دتے ہوندے، جیہڑے اکّو طرح نال پڑھت دی ضمانت نہ
دیندے، تاں کی ہوندا۔ مثال لئی، اسیں رابطیاں دے
تشکیل قاعدیاں وچ قوساں بھل سکدے ساں؛ مثلاً اسیں
ایہ اجازت دے سکدے ساں:
\begin{quote}
جے $!A$ تے $!B$، !!{formula}s نیں، تاں $!A \lif !B$ وی اے۔
\end{quote}
ایٹمی فارمولا $!D$ توں شروع کر کے ایہ سانوں $!D
\lif !D$ بناؤن دی اجازت دیندا۔ ایس توں، $!D$ دے نال،
سانوں $!D \lif !D
\lif !D$ مل جاندا۔ پر ایہ کرن دے دو طریقے نیں:
\begin{enumerate}
\item اسیں $!D$ نوں $!A$ منّیئے تے $!D \lif !D$ نوں $!B$۔
\item اسیں $!A$ نوں $!D \lif !D$ منّیئے تے $!B$،~$!D$ اے۔
\end{enumerate}
ایس دے مطابق !!{formula}~$!D \lif !D \lif !D$
نوں ”پڑھن“ دے دو طریقے نیں۔ ایہ $!B
\lif !C$ دے روپ وچ اے، جدوں $!B$، $!D$ اے تے $!C$،
$!D \lif !D$ اے؛ پر \emph{ایہ ایس روپ وچ وی اے}:
$!B \lif !C$، جتھے $!B$، $!D \lif !D$ اے تے $!C$
اوہو~$!D$ اے۔

جے ایہ ہووے، تاں ساڈیاں تعریفاں ہر وار کم نہیں کرن
گیاں۔ مثال لئی، جدوں اسیں کسے فارمولے دے
!!{main operator} دی تعریف کردے آں، تاں آکھدے آں:
$!B \lif !C$ روپ دے فارمولے وچ !!{main operator}،
~$\lif$ دی دکھائی گئی آمد اے۔ پر جے اسیں فارمولا $!D
\lif !D \lif !D$ نوں اُتے دسے دو وکھرے طریقیاں نال
$!B \lif !C$ دے روپ وچ رکھ سکدے آں، تاں پہلی
صورت وچ $\lif$ دی پہلی آمد !!{main operator} بنے گی
تے دوجی صورت وچ دوجی آمد۔ پر ساڈا مقصد
!!{main operator} نوں !!{formula} دا اک \emph{فنکشن}
% PNB-SYN-006: اگلی main-operator تعریف مطابق ایٹمی فارمولیاں دا مرکزی عامل نہیں؛ضمنی غیرایٹمی دائرہ صاف۔
منّنا اے؛ یعنی ہر غیر ایٹمی !!{formula} وچ !!{main operator} دی ٹھیک اک آمد ہونی چاہیدی اے۔
\end{explain}

\begin{lem}
!!a{formula}~$!A$ وچ کھبیاں تے سجیاں قوساں دی
گنتی برابر اے۔
\end{lem}

\begin{proof}
اسیں ایہ $!A$ دی تشکیل اُتے استقرا نال ثابت کردے آں۔
ایس لئی دو گلاں چاہیدیاں نیں: (الف) پہلاں ثابت کریے
کہ سارے ایٹمی فارمولیاں وچ متعلقہ خاصیت اے
(استقرا دی بنیاد)۔ (ب) فیر ثابت کریے کہ جدوں دتے
فارمولیاں توں نویں فارمولے بنائے جان، تاں جے پرانیاں
وچ خاصیت اے، نویاں وچ وی ہووے۔

$l(!A)$ نوں کھبیاں قوساں دی گنتی تے $r(!A)$ نوں
~$!A$ وچ سجیاں قوساں دی گنتی منّو؛ ایس طرح
$l(t)$ تے $r(t)$، ٹرم~$t$ وچ کھبیاں تے سجیاں
قوساں دی گنتی نیں۔

\begin{prob}
    ہر ٹرم~$t$ لئی $l(t) = r(t)$ ثابت کرو۔
\end{prob}

\begin{enumerate}
\tagitem{prvFalse}{\indcase{!A}{\lfalse}{$\indfrm$ وچ $0$ کھبیاں تے $0$
    سجیاں قوساں نیں۔}}{}

\tagitem{prvTrue}{\indcase{!A}{\ltrue}{$\indfrm$ وچ $0$ کھبیاں تے $0$
    سجیاں قوساں نیں۔}}{}

\item \indcase{!A}{\Atom{R}{t_1,\dots,t_n}}{$l(\indfrm) = 1 + l(t_1) +
  \dots + l(t_n) = 1 + r(t_1) + \dots + r(t_n) = r(\indfrm)$۔ ایتھے
  اسیں اوہ گل ورتدے آں جیہڑی مشق لئی چھڈی گئی اے:
  $l(t) = r(t)$، ہر ٹرم~$t$ لئی۔}

\item \indcase{!A}{\eq[t_1][t_2]}{$l(\indfrm) = l(t_1) + l(t_2) =
  r(t_1) + r(t_2) = r(\indfrm)$۔}

\tagitem{prvNot}{\indcase{!A}{\lnot !B}{استقرائی مفروضے نال
    $l(!B) = r(!B)$۔ اس لئی $l(\indfrm) = l(!B) = r(!B) =
    r(\indfrm)$۔}}{}

\item \indcase{!A}{(!B \ast !C)}{استقرائی مفروضے نال $l(!B) =
  r(!B)$ تے $l(!C) = r(!C)$۔ اس لئی $l(\indfrm) = 1 + l(!B) + l(!C) =
  1 + r(!B) + r(!C) = r(\indfrm)$۔}

\tagitem{prvAll}{\indcase{!A}{\lforall[x][!B]}{استقرائی
    مفروضے نال $l(!B) = r(!B)$۔ اس لئی $l(\indfrm) = l(!B) = r(!B) =
    r(\indfrm)$۔}}{}

% Print case for \lexists if prvEx and not prvAll
\tagitem{prvAll,notprvEx}{}{\indcase{!A}{\lexists[x][!B]}{استقرائی
    مفروضے نال $l(!B) = r(!B)$۔ اس لئی $l(\indfrm) = l(!B) = r(!B) =
    r(\indfrm)$۔}}

% Just say similarly if prvEx and prvAll
\tagitem{notprvAll,notprvEx}{}{\indcase{!A}{\lexists[x][!B]}{ایویں ای۔}}
\end{enumerate}
\end{proof}

\begin{defn}[پورے توں چھوٹا مُڈھلا ٹکڑا]
علامتاں دی لڑی $!B$، علامتاں دی لڑی~$!A$ دا
\emph{پورے توں چھوٹا مُڈھلا ٹکڑا} اے، جے $!B$
نال علامتاں دی اک غیر خالی لڑی جوڑ کے~$!A$ ملے۔
\end{defn}

\begin{lem}\ollabel{lem:no-prefix}
جے $!A$، !!a{formula} اے تے $!B$، $!A$ دا
پورے توں چھوٹا مُڈھلا ٹکڑا اے، تاں $!B$،
!!a{formula} نہیں اے۔
\end{lem}

\begin{proof}
مشق۔
\end{proof}

\begin{prob}
\olref[fol][syn][unq]{lem:no-prefix} ثابت کرو۔
\end{prob}

\begin{prop}
\ollabel{prop:unique-atomic}
جے $!A$ ایٹمی !!{formula} اے، تاں اوہ ہیٹھ
دیتیاں شرطاں وچوں ٹھیک اک شرط پوری کردا اے۔
\begin{enumerate}
\tagitem{prvFalse}{$!A \ident \lfalse$۔}{}
\tagitem{prvTrue}{$!A \ident \ltrue$۔}{}
\item $!A \ident \Atom{R}{t_1,\dots,t_n}$، جتھے $R$ اک $n$-مقامی
  !!{predicate} اے، $t_1$, \dots, $t_n$ ٹرم نیں، تے $R$،
  $t_1$, \dots, $t_n$ وچوں ہر اک اکّو طرح متعین اے۔
\item $!A \ident \eq[t_1][t_2]$، جتھے $t_1$ تے $t_2$
  اکّو طرح متعین ٹرم نیں۔
\end{enumerate}
\end{prop}

\begin{proof}
مشق۔
\end{proof}

\begin{prob}
\olref[fol][syn][unq]{prop:unique-atomic} ثابت کرو (اشارہ:
ٹرماں لئی \olref[fol][syn][unq]{lem:no-prefix} دا اک روپ
بیان کرو تے ثابت کرو۔)
\end{prob}

\begin{prop}[اکّو طرح نال پڑھت]
ہر !!{formula} ہیٹھ دیتیاں شرطاں وچوں ٹھیک اک
شرط پوری کردا اے۔
\begin{enumerate}
\item $!A$ ایٹمی اے۔

\tagitem{prvNot}{$!A$، $\lnot !B$ دے روپ وچ اے۔}{}

\tagitem{prvAnd}{$!A$، $(!B \land !C)$ دے روپ وچ اے۔}{}

\tagitem{prvOr}{$!A$، $(!B \lor !C)$ دے روپ وچ اے۔}{}

\tagitem{prvIf}{$!A$، $(!B \lif !C)$ دے روپ وچ اے۔}{}

\tagitem{prvIff}{$!A$، $(!B \liff !C)$ دے روپ وچ اے۔}{}

\tagitem{prvAll}{$!A$، $\lforall[x][!B]$ دے روپ وچ اے۔}{}

\tagitem{prvEx}{$!A$، $\lexists[x][!B]$ دے روپ وچ اے۔}{}
\end{enumerate}
ہور ایہ کہ ہر صورت وچ $!B$، یا $!B$ تے $!C$،
اکّو طرح متعین نیں۔ ایس دا مطلب اے کہ، مثال لئی،
وکھرے جوڑے $!B$, $!C$ تے $!B'$, $!C'$ نہیں ہو
سکدے جیہناں لئی $!A$ دونیں روپاں وچ ہووے:
\iftag{prvIf}{$(!B \lif !C)$ تے $(!B' \lif !C')$۔}{
    \iftag{prvOr}{$(!B \lor !C)$ تے $(!B' \lor !C')$۔}{
      \iftag{prvAnd}{$(!B \land !C)$ تے $(!B' \land !C')$۔}{}}}
\end{prop}

\begin{proof}
تشکیل دے قاعدے لازمی کردے نیں کہ جے !!a{formula}
ایٹمی نہیں، تاں اوہ کھلدی قوس~(،
\iftag{prvNot}{$\lnot$،}{} یا مسوّر نال شروع ہووے۔
دوجے پاسے، جیہڑا وی !!{formula} ہیٹھ دیاں
علامتاں وچوں کسے اک نال شروع ہووے، اوہ ایٹمی
ہووے گا: !!a{predicate}، !!a{function}، !!a{constant}\iftag{prvFalse}{، $\lfalse$}{}\iftag{prvTrue}{، $\ltrue$}{}۔

اس لئی سانوں اصل وچ صرف ایہ دکھانا اے کہ جے $!A$
اک پاسے $(!B \ast
!C)$ دے روپ وچ اے تے دوجے پاسے $(!B' \mathbin{\ast'} !C')$،
تاں $!B \ident
!B'$، $!C \ident !C'$، تے $\ast = {\ast'}$۔

فرض کرو $!A \ident (!B \ast !C)$ وی اے تے $!A \ident (!B'
\mathbin{\ast'} !C')$ وی۔ تاں یا $!B \ident !B'$، یا ایہ
نہیں۔ جے ایہ اے، تاں صاف اے کہ $\ast = {\ast'}$ تے
$!C \ident !C'$، کیوں جے ایہ فیر $!A$ دیاں
اوہ ذیلی لڑیاں نیں جیہڑیاں اکّو مقام توں شروع ہوندیاں
تے اکّو لمبائی رکھدیاں نیں۔ دوجی صورت $!B \not\ident !B'$
اے۔ کیوں جے $!B$ تے
$!B'$ دونیں $!A$ دیاں ذیلی لڑیاں نیں تے اکّو مقام
توں شروع ہوندیاں نیں، تاں اک دوجی دا پورے توں
چھوٹا مُڈھلا ٹکڑا لازماً ہووے گی۔ پر ایہ
\olref{lem:no-prefix} دے مطابق ناممکن اے۔
\end{proof}
\endgroup
% /OLP-0153

% OLP-0154 content/first-order-logic/syntax-and-semantics/main-operator.tex
\begingroup
\def\olfilename{main-operator}
\expandafter\def\csname ol@sectioncs\endcsname{section}

\olfileid{fol}{syn}{mai}

\olsection{فارمولے دا \printtoken{S}{main operator}}

\begin{explain}
!!a{formula}~$!A$ بناؤن وچ سب توں آخر ورتے عملگر
دی گل کرنا اکثر فائدہ مند ہوندا اے۔ ایس عملگر نوں
~$!A$ دا \emph{مرکزی عامل} آکھدے آں۔ بدیہی طور
تے ایہ $!A$ دا ”سب توں باہرا“ عملگر اے۔ مثال
لئی، $\lnot !A$ دا مرکزی عامل $\lnot$ اے،
تے $(!A \lor !B)$ دا مرکزی عامل $\lor$ اے، وغیرہ۔
\end{explain}


\begin{defn}[!!^{main operator}]
\ollabel{def:main-op}
!!a{formula}~$!A$ دا \emph{!!{main operator}}
اینج تعریف کیتا جاندا اے:
\begin{enumerate}
\item \indcase*{!A}{!A}{$\indfrm$ دا کوئی !!{main operator} نہیں۔}

\tagitem{prvNot}{\indcase{!A}{\lnot !B}{$\indfrm$ دا !!{main operator}
  ~$\lnot$ اے۔}}{}

\tagitem{prvAnd}{\indcase{!A}{(!B \land !C)}{$\indfrm$ دا !!{main operator}
  ~$\land$ اے۔}}{}

\tagitem{prvOr}{\indcase{!A}{(!B \lor !C)}{$\indfrm$ دا !!{main operator}
  ~$\lor$ اے۔}}{}

\tagitem{prvIf}{\indcase{!A}{(!B \lif !C)}{$\indfrm$ دا !!{main operator}
  ~$\lif$ اے۔}}{}

\tagitem{prvIff}{\indcase{!A}{(!B \liff !C)}{$\indfrm$ دا !!{main operator}
  ~$\liff$ اے۔}}{}

\tagitem{prvAll}{\indcase{!A}{\lforall[x][!B]}{$\indfrm$ دا !!{main operator}
  ~$\lforall$ اے۔}}{}

\tagitem{prvEx}{\indcase{!A}{\lexists[x][!B]}{$\indfrm$ دا !!{main operator}
  ~$\lexists$ اے۔}}{}
\end{enumerate}
\end{defn}

ہر صورت وچ ساڈی مراد فارمولے دے !!{main operator}
دی دکھائی گئی خاص \emph{آمد} اے۔ مثال لئی، کیوں جے
فارمولا $((!D \lif !E) \lif (!E \lif !D))$، $(!B \lif !C)$
دے روپ وچ اے، جتھے $!B$، $(!D \lif !E)$ اے تے
$!C$، $(!E \lif !D)$ اے، ایس لئی $\lif$ دی
دوجی آمد !!{main operator} اے۔

\begin{explain}
% PNB-SYN-007: اگلے حصے دے صریح فرق مطابق ایتھے چھوٹے فارمولیاں اُتے ایس فنکشن دیاں قدراں نہیں ورتیاں؛ ایہ صورتاں دے لحاظ نال سیدھی تعریف اے، بازگشتی تعریف نہیں۔
ایہ اک فنکشن دی \emph{صورتاں دے لحاظ نال سیدھی} تعریف اے، جیہڑا سارے
غیر ایٹمی !!{formula}s نوں اوہناں دے !!{main operator}
دی آمد نال جوڑدا اے۔ جس طرح !!{formula}s
استقرائی طور تے تعریف کیتے گئے نیں، ہر
!!{formula}~$!A$، \olref{def:main-op} دیاں
صورتاں وچوں اک پوری کردا اے۔ ایہ ضمانت دیندا
اے کہ ہر غیر ایٹمی !!{formula}~$!A$ لئی
!!a{main operator} موجود اے۔ کیوں جے ہر !!{formula}
ایہناں شرطاں وچوں صرف اک پوری کردا اے، تے
ہر صورت وچ اوہ چھوٹے !!{formula}s وی اکّو طرح
متعین نیں جیہناں توں $!A$ بنیا اے، اس لئی
~$!A$ دے !!{main operator} دی آمد اکّو اے، تے ایس طرح
اسیں اک فنکشن تعریف کر دتا اے۔
\end{explain}

اسیں !!{formula}s نوں \olref{tab:main-op} وچ
دتے ناں اوہناں دے !!{main operator} دی علامت
مطابق دیندے آں۔\iftag{defNot,defOr,defAnd,defIf,defIff,defTrue,defFalse,defEx,defAll}
{ پر یاد رکھو کہ تعریف شدہ عملگر باقاعدہ طور تے
!!{formula}s وچ نہیں آؤندے۔ اوہ صرف مخفف نیں،
اس لئی باقاعدہ طور تے کسے فارمولے دے مرکزی عامل
نہیں بن سکدے۔ لہٰذا سارے !!{formula}s بارے ثبوتاں
وچ اوہناں دا وکھرا معاملہ کرن دی لوڑ نہیں۔}{}

% PNB-SYN-002 ماخذ دی درستی: جدول دے عطف، فصل تے شرطیہ مثالاں دیاں بند ہُندیاں قوساں ماخذ وچ ریاضی موڈ توں باہر نیں؛ ایتھے تیناں مکمل فارمولے ریاضی موڈ وچ رکھے نیں۔
\begin{table}[!h]
\centering
\begin{tabular}{c | c | c}
!!^{main operator} & !!{formula} دی قسم & مثال\\
\hline
کوئی نہیں & ایٹمی (!!{formula}) &
\iftag{prvFalse}{$\lfalse$,}{}
\iftag{prvTrue}{$\ltrue$,}{}
$\Atom{R}{t_1, \dots, t_n}$,
$\eq[t_1][t_2]$\\
$\lnot$ & نفی & $\lnot !A$ \\
$\land$ & عطف & $(!A \land !B)$ \\
$\lor$ & فصل & $(!A \lor !B)$ \\
$\lif$ & !!{conditional} & $(!A \lif !B)$ \\
$\liff$ & !!{biconditional} & $(!A \liff !B)$ \\
$\lforall[][]$ & کُلّی (!!{formula}) & $\lforall[x][!A]$ \\
$\lexists[][]$ & وجودی (!!{formula}) & $\lexists[x][!A]$
\end{tabular}
\caption{!!{formula}s دے ناں تے مرکزی عامل}
\ollabel{tab:main-op}
\end{table}
\endgroup
% /OLP-0154

% OLP-0155 content/first-order-logic/syntax-and-semantics/subformulas.tex
\begingroup
\def\olfilename{subformulas}
\expandafter\def\csname ol@sectioncs\endcsname{section}

\olfileid{fol}{syn}{sbf}

\olsection{\printtoken{P}{subformula}}

\begin{explain}
کسے دتے !!{formula} نوں ”بناؤن“ والے !!{formula}s
دی گل کرنا اکثر فائدہ مند ہوندا اے۔ اسیں اوہناں نوں
اوہدے \emph{!!{subformula}s} آکھدے آں۔ ہر
!!{formula} اپنے آپ دا !!a{subformula} وی اے؛
$!A$ دا اوہدے اپنے آپ $!A$ توں وکھرا ذیلی
فارمولا \emph{پورے توں چھوٹا !!{subformula}} اے۔
\end{explain}

\begin{defn}[براہِ راست !!^{subformula}]
جے $!A$، !!a{formula} اے، تاں $!A$ دے
\emph{براہِ راست !!{subformula}s} اینج استقرائی
طور تے تعریف کیتے جاندے نیں:
\begin{enumerate}
\item ایٹمی !!{formula}s دے کوئی براہِ راست !!{subformula}s نہیں۔

\tagitem{prvNot}{\indcase{!A}{\lnot !B}{$\indfrm$ دا اکلوتا براہِ راست
    !!{subformula}،~$!B$ اے۔}}{}

\item \indcase{!A}{(!B \ast !C)}{$\indfrm$ دے براہِ راست !!{subformula}s
  $!B$ تے $!C$ نیں ($\ast$ دو-مقامی رابطیاں وچوں کوئی وی اک اے)۔}

\tagitem{prvAll}{\indcase{!A}{\lforall[x][!B]}{$\indfrm$ دا اکلوتا براہِ راست
    !!{subformula}،~$!B$ اے۔}}{}

\tagitem{prvEx}{\indcase{!A}{\lexists[x][!B]}{$\indfrm$ دا اکلوتا براہِ راست
    !!{subformula}،~$!B$ اے۔}}{}
\end{enumerate}
\end{defn}

\begin{defn}[پورے توں چھوٹا !!^{subformula}]
جے $!A$، !!a{formula} اے، تاں $!A$ دے
\emph{پورے توں چھوٹے !!{subformula}s} اینج بازگشتی
طور تے تعریف کیتے جاندے نیں:
\begin{enumerate}
\item ایٹمی !!{formula}s دے کوئی پورے توں چھوٹے !!{subformula}s نہیں۔

\tagitem{prvNot}{\indcase{!A}{\lnot !B}{$\indfrm$ دے پورے توں چھوٹے
    !!{subformula}s،~$!B$ تے~$!B$ دے سارے پورے توں
    چھوٹے !!{subformula}s نیں۔}}{}

\item \indcase{!A}{(!B \ast !C)}{$\indfrm$ دے پورے توں چھوٹے !!{subformula}s
  $!B$، $!C$، نالے $!B$ دے سارے پورے توں چھوٹے !!{subformula}s
  تے~$!C$ دے سارے ایہو جہے ذیلی فارمولے نیں۔}

\tagitem{prvAll}{\indcase{!A}{\lforall[x][!B]}{$\indfrm$ دے پورے توں
    چھوٹے !!{subformula}s،~$!B$ تے~$!B$ دے سارے پورے توں
    چھوٹے !!{subformula}s نیں۔}}{}

\tagitem{prvEx}{\indcase{!A}{\lexists[x][!B]}{$\indfrm$ دے پورے توں
    چھوٹے !!{subformula}s،~$!B$ تے~$!B$ دے سارے پورے توں
    چھوٹے !!{subformula}s نیں۔}}{}
\end{enumerate}
\end{defn}

\begin{defn}[!!^{subformula}]
$!A$ دے !!{subformula}s وچ $!A$ خود تے اوہدے
سارے پورے توں چھوٹے !!{subformula}s شامل نیں۔
\end{defn}

\begin{explain}
دھیان دیو کہ اسیں براہِ راست !!{subformula}s تے پورے
توں چھوٹے !!{subformula}s دی تعریف دو ذرا وکھرے
طریقیاں نال کیتی اے۔ پہلی صورت وچ اسیں فارمولا~$!A$
دے ہر ممکن روپ~$!A$ لئی اوہدے براہِ راست !!{subformula}s
دی سیدھی تعریف دتی اے۔ ایہ صورتاں دے لحاظ نال
صاف تعریف اے، تے صورتاں !!{formula}s دے سیٹ
دی استقرائی تعریف دے مطابق نیں۔ دوجی صورت وچ
اسیں سارے !!{formula}s دے سیٹ دی تعریف والا
ڈھنگ فیر اپنایا اے، پر ہر صورت وچ چھوٹے
!!{formula}s $!B$، $!C$ دے پورے توں چھوٹے
!!{subformula}s وی انہاں !!{formula}s دے نال شامل
کیتے نیں۔ ایس نال تعریف \emph{بازگشتی}
بن جاندی اے۔ عام طور تے استقرائی طور تے تعریف
کیتے سیٹ (ساڈے معاملے وچ، !!{formula}s) اُتے
کسے فنکشن دی تعریف بازگشتی اے جے فنکشن دی
تعریف دیاں صورتاں وچ خود اوہ فنکشن ورتیا جائے۔
پر تعریف ٹھیک ہون لئی سانوں یقین کرنا پیندا اے
کہ فنکشن دیاں قدراں صرف اوہناں دلیلاں لئی ورتیاں
جان جیہڑیاں زیرِ تعریف دلیل توں ”پہلاں“ آؤندیاں نیں:
ساڈے معاملے وچ، جدوں $(!B \ast !C)$ لئی
”پورے توں چھوٹے !!{subformula}“ تعریف کردے آں،
تاں صرف ”پہلاں“ والے !!{formula}s $!B$ تے $!C$
دے پورے توں چھوٹے !!{subformula}s ورتدے آں۔
\end{explain}

\begin{prop}
\ollabel{prop:subfrm-trans}
فرض کرو $!B$، $!A$ دا ذیلی فارمولا اے تے
$!C$، $!B$ دا ذیلی فارمولا اے۔ تاں $!C$، $!A$
دا ذیلی فارمولا اے۔ دوجے لفظاں وچ ذیلی فارمولا
تعلق متعدی اے۔
\end{prop}

\begin{prob}
\olref[fol][syn][sbf]{prop:subfrm-trans} ثابت کرو۔
\end{prob}

\begin{prop}
\ollabel{prop:count-subfrms}
فرض کرو $!A$ اک فارمولا اے جیہدے وچ $n$ رابطے
تے مسوّر نیں۔ تاں $!A$ دے زیادہ توں زیادہ $2n+1$
ذیلی فارمولے نیں۔
\end{prop}

\begin{prob}
\olref[fol][syn][sbf]{prop:count-subfrms} ثابت کرو۔
\end{prob}
\endgroup
% /OLP-0155

% OLP-0156 content/first-order-logic/syntax-and-semantics/formation-sequences.tex
\begingroup
\def\olfilename{formation-sequences}
\expandafter\def\csname ol@sectioncs\endcsname{section}

\olfileid{fol}{syn}{fseq}

\olsection{تشکیل دے سلسلے}

!!{formula}s نوں استقرائی تعریف نال تعریف کرنا، تے !!{formula}s
دیاں خاصیتاں استقرا نال ثابت کرن دا نال دا طریقہ، اک سوہنا تے
کارگر ڈھنگ اے۔ پر !!{formula}s دی بناوٹ نوں ہِٹھوں اُتّے ول،
قدم بہ قدم وی ویکھیا جا سکدا اے۔ ایتھے اسیں ایہ کم
\emph{تشکیل دے سلسلے} دے تصور نال کردے آں۔
%
ایس طریقے نال ٹرم تے !!{formula}s نوں اوہناں دیاں استقرائی
تعریفاں دا حوالہ دتے بغیر پیش کرن لئی، پہلے اسیں کسے
زبان~$\Lang L$ دیاں علامتاں توں بنی من مانی لڑی دا تصور دیندے آں۔

\begin{defn}[لڑیاں]
\ollabel{defn:string}
فرض کرو $\Lang L$ پہلے درجے دی زبان اے۔ \emph{$\Lang
L$-لڑی}، $\Lang L$ دیاں علامتاں دا محدود سلسلہ اے۔ جدوں
سیاق توں زبان~$\Lang L$ صاف ہووے، تاں اسیں
$\Lang L$-لڑی نوں صرف \emph{لڑی} وی آکھاں گے۔
\end{defn}

\begin{ex}
کسے وی پہلے درجے دی زبان $\Lang L$ لئی، سارے
$\Lang L$-!!{formula}s، $\Lang L$-لڑیاں نیں؛ پر
لازم نہیں کہ ہر ایسی لڑی فارمولا ہووے۔ مثال لئی،
\[)(\Obj v_0\lif\lexists\] اک $\Lang L$-لڑی اے،
پر $\Lang L$-!!{formula} نہیں۔
\end{ex}

\begin{defn}[ٹرماں لئی تشکیل دے سلسلے]
\ollabel{defn:fseq-trm}
$\Lang L$-لڑیاں دا محدود سلسلہ $\tuple{t_0,\dotsc,t_n}$،
ٹرم $t$ لئی \emph{تشکیل دا سلسلہ} اے جے $t \ident t_n$،
تے ہر $i \leq n$ لئی یا تاں $t_i$ !!a{variable} یا
!!a{constant} اے، یا $\Lang L$ وچ اک $k$-مقامی
!!{function}~$f$ اے تے ایہو جہے $m_1,\dotsc,m_k < i$
موجود نیں کہ $t_i \ident f(t_{m_1},\dotsc,t_{m_k})$۔
% PNB-SYN-003: منجمد ماخذ وچ k-مقامی فنکشن لئی m_0 توں m_k تک k+1 دلیلاں سن؛ ایتھے k دلیلاں رکھیاں نیں۔
جدوں فرق واضح کرن دی لوڑ ہووے، اسیں ٹرماں دے تشکیل دے سلسلیاں نوں
\emph{ٹرم تشکیل دے سلسلے} آکھاں گے۔
\end{defn}

\begin{ex}
ایہ سلسلہ
\[
    \tuple{\Obj c_0, \Obj v_0, \Atom{\Obj f^2_0}{\Obj c_0, \Obj v_0}, \Atom{\Obj f^1_0}{\Atom{\Obj f^2_0}{\Obj c_0, \Obj v_0}}}
\]
ٹرم $\Atom{\Obj f^1_0}{\Atom{\Obj
f^2_0}{\Obj c_0, \Obj v_0}}$ لئی تشکیل دا سلسلہ اے،
تے ایہ وی:
\[
    \tuple{\Obj v_0, \Obj c_0, \Atom{\Obj f^2_0}{\Obj c_0, \Obj v_0}, \Atom{\Obj f^1_0}{\Atom{\Obj f^2_0}{\Obj c_0, \Obj v_0}}}.
\]
\end{ex}

\begin{defn}[فارمولیاں لئی تشکیل دے سلسلے]
\ollabel{defn:fseq-frm}
$\Lang L$-لڑیاں دا محدود سلسلہ $\tuple{!A_0,\dotsc,!A_n}$
$!A$ لئی \emph{تشکیل دا سلسلہ} اے جے $!A \ident !A_n$،
تے ہر $i \leq n$ لئی یا تاں $!A_i$ ایٹمی !!{formula} اے،
یا ایہو جہے $j,k < i$ تے !!a{variable}~$x$ موجود نیں
کہ اگلیاں وچوں کوئی اک صورت پوری ہُندی اے:
\begin{enumerate}
    \tagitem{prvNot}{$!A_i \ident \lnot !A_j$.}{}%
    \tagitem{prvAnd}{$!A_i \ident (!A_j \land !A_k)$.}{}%
    \tagitem{prvOr}{$!A_i \ident (!A_j \lor !A_k)$.}{}%
    \tagitem{prvIf}{$!A_i \ident (!A_j \lif !A_k)$.}{}%
    \tagitem{prvIff}{$!A_i \ident (!A_j \liff !A_k)$.}{}%
    \tagitem{prvAll}{$!A_i \ident \lforall[x][!A_j]$.}{}%
    \tagitem{prvEx}{$!A_i \ident \lexists[x][!A_j]$.}{}%
\end{enumerate}
جدوں فرق واضح کرن دی لوڑ ہووے، اسیں فارمولیاں دے تشکیل دے سلسلیاں نوں
\emph{فارمولا تشکیل دے سلسلے} آکھاں گے۔
\end{defn}

\begin{ex}
\[
\tuple{
    \Atom{\Obj A^1_0}{\Obj v_0},
    \Atom{\Obj A^1_1}{\Obj c_1},
    (\Atom{\Obj A^1_1}{\Obj c_1} \land \Atom{\Obj A^1_0}{\Obj v_0}),
    \lexists[\Obj v_0][(\Atom{\Obj A^1_1}{\Obj c_1} \land \Atom{\Obj A^1_0}{\Obj v_0})]
}
\]
$\lexists[\Obj v_0][(\Atom{\Obj A^1_1}{\Obj
c_1} \land \Atom{\Obj A^1_0}{\Obj v_0})]$ لئی
تشکیل دا سلسلہ اے؛ ایہ وی اک اے:
\begin{multline*}
\tuple{
    \Atom{\Obj A^1_0}{\Obj v_0},
    \Atom{\Obj A^1_1}{\Obj c_1},
    (\Atom{\Obj A^1_1}{\Obj c_1} \land \Atom{\Obj A^1_0}{\Obj v_0}),
    \Atom{\Obj A^1_1}{\Obj c_1},\\
    \lforall[\Obj v_1][\Atom{\Obj A^1_0}{\Obj v_0}],
    \lexists[\Obj v_0][(\Atom{\Obj A^1_1}{\Obj c_1} \land \Atom{\Obj A^1_0}{\Obj v_0})]
}.
\end{multline*}
%
دوجی مثال وکھاؤندی اے کہ تشکیل دے سلسلیاں وچ ”فالتو مواد“ وی
آ سکدا اے: اوہ !!{formula}s جیہڑے دہرائے گئے ہون یا
بناوٹ وچ کوئی حصہ نہ پاؤندے ہون۔
\end{ex}

\begin{prop}\ollabel{prop:formed}
ہر !!{formula}~$!A$ جیہڑا~$\Frm[L]$ وچ اے،
اوہدے لئی تشکیل دا سلسلہ موجود اے۔
\end{prop}

\begin{proof}
فرض کرو $!A$ ایٹمی اے۔ تاں سلسلہ $\tuple{!A}$،
$!A$ لئی تشکیل دا سلسلہ اے۔
%
ہُن فرض کرو $!B$ تے~$!C$ دے تشکیل دے سلسلے
ترتیب وار $\tuple{!B_0,\dotsc,!B_n}$ تے
$\tuple{!C_0,\dotsc,!C_m}$ نیں۔
%
\begin{enumerate}
  \tagitem{prvNot}{جے $!A \ident \lnot !B$،
      تاں $\tuple{!B_0,\dotsc,!B_n,\lnot !B_n}$
      $!A$ لئی تشکیل دا سلسلہ اے۔}{}
  \tagitem{prvAnd}{جے $!A \ident (!B \land !C)$،
      تاں $\tuple{!B_0,\dotsc,!B_n,!C_0,\dotsc,!C_m,(!B_n \land !C_m)}$
      $!A$ لئی تشکیل دا سلسلہ اے۔}{}
  \tagitem{prvOr}{جے $!A \ident (!B \lor !C)$،
      تاں $\tuple{!B_0,\dotsc,!B_n,!C_0,\dotsc,!C_m,(!B_n \lor !C_m)}$
      $!A$ لئی تشکیل دا سلسلہ اے۔}{}
  \tagitem{prvIf}{جے $!A \ident (!B \lif !C)$،
      تاں $\tuple{!B_0,\dotsc,!B_n,!C_0,\dotsc,!C_m,(!B_n \lif !C_m)}$
      $!A$ لئی تشکیل دا سلسلہ اے۔}{}
  \tagitem{prvIff}{جے $!A \ident (!B \liff !C)$،
      تاں $\tuple{!B_0,\dotsc,!B_n,!C_0,\dotsc,!C_m,(!B_n \liff !C_m)}$
      $!A$ لئی تشکیل دا سلسلہ اے۔}{}
  \tagitem{prvAll}{جے $!A \ident \lforall[x][!B]$،
      تاں $\tuple{!B_0,\dotsc,!B_n,\lforall[x][!B_n]}$
      $!A$ لئی تشکیل دا سلسلہ اے۔}{}
  \tagitem{prvEx}{جے $!A \ident \lexists[x][!B]$،
      تاں $\tuple{!B_0,\dotsc,!B_n,\lexists[x][!B_n]}$
      $!A$ لئی تشکیل دا سلسلہ اے۔}{}
\end{enumerate}
!!{formula}s اُتے استقرا دے اصول موجب ہر !!{formula}
دا تشکیل دا سلسلہ اے۔
\end{proof}

اسیں اُلٹا رخ وی ثابت کر سکدے آں۔ ایہ اہم اے،
کیوں جے ایہ وکھاؤندا اے کہ فارمولیاں نوں تعریف
کرن دے ساڈے دونیں طریقے ہم ارز نیں: اوہ اکّو
نتیجہ دیندے نیں۔ ایس دا مطلب ایہ وی اے کہ
اسیں تشکیل دے سلسلے دی لمبائی اُتے عام استقرا
لا کے فارمولیاں بارے قضیے ثابت کر سکدے آں۔

\begin{lem}
\ollabel{lem:fseq-init}
فرض کرو $\tuple{!A_0,\dotsc,!A_n}$،~$!A_n$ لئی
تشکیل دا سلسلہ اے، تے $k \leq n$۔ تاں
$\tuple{!A_0,\dotsc,!A_k}$،~$!A_k$ لئی
تشکیل دا سلسلہ اے۔
\end{lem}

\begin{proof}
مشق۔
\end{proof}

\begin{prob}
\olref[fol][syn][fseq]{lem:fseq-init} ثابت کرو۔
\end{prob}

\begin{thm}
\ollabel{thm:fseq-frm-equiv}
$\Frm[L]$ اوہناں ساریاں $\Lang L$-لڑیاں~$!A$
دا سیٹ اے جیہناں دی شرط ایہ اے کہ $!A$ لئی فارمولا تشکیل دا سلسلہ موجود اے۔
\end{thm}

\begin{proof}
$F$ نوں زبان~$\Lang L$ دیاں اوہناں ساریاں علامتی
لڑیاں دا سیٹ منّو جنہاں دا تشکیل دا سلسلہ اے۔
اسیں \olref[fol][syn][fseq]{prop:formed} وچ
ویکھ چکے آں کہ $\Frm[L] \subseteq F$؛ سو ہُن
اسیں اُلٹا رخ ثابت کردے آں۔

فرض کرو $!A$ دا تشکیل دا سلسلہ
$\tuple{!A_0,\dotsc,!A_n}$ اے۔ اسیں ثابت
کردے آں کہ $!A \in \Frm[L]$؛ ایہ کم
$n$ یعنی آخری رکن دے اشاریے اُتے مضبوط استقرا نال کردے آں۔
ساڈا استقرائی مفروضہ اے کہ جیہڑی وی علامتی لڑی
کسے ایسے تشکیل دے سلسلے دا آخری رکن اے
جیہدا آخری اشاریہ $m < n$ اے، اوہ $\Frm[L]$ وچ اے۔
% PNB-SYN-004: ماخذ n اُتے استقرا کردا پر n+1 رکناں نوں لمبائی n آکھدا اے؛ ایتھے آخری اشاریہ مستقل رکھیا اے۔
تشکیل دے سلسلے دی تعریف موجب یا تاں
$!A \ident !A_n$ ایٹمی اے، یا ایہو جہے
$j,k < n$ موجود نیں کہ اگلیاں وچوں کوئی اک صورت اے:
\begin{enumerate}
    \tagitem{prvNot}{$!A \ident \lnot !A_j$.}{}%
    \tagitem{prvAnd}{$!A \ident (!A_j \land !A_k)$.}{}%
    \tagitem{prvOr}{$!A \ident (!A_j \lor !A_k)$.}{}%
    \tagitem{prvIf}{$!A \ident (!A_j \lif !A_k)$.}{}%
    \tagitem{prvIff}{$!A \ident (!A_j \liff !A_k)$.}{}%
    \tagitem{prvAll}{$!A \ident \lforall[x][!A_j]$.}{}%
    \tagitem{prvEx}{$!A \ident \lexists[x][!A_j]$.}{}%
\end{enumerate}
ہُن اسیں صورتاں دے لحاظ نال دلیل دیندے آں۔
جے $!A$ ایٹمی اے تاں $!A_n \in \Frm[L]$۔
ایس دی بجائے فرض کرو کہ $!A \ident (!A_j \land !A_k)$۔
% PNB-SYN-005: ماخذ دے دو Frm[L_0] حوالیاں نوں پہلے درجے دے Frm[L] نال تے equiv نوں نحوی ident نال درست کیتا اے۔
\olref[fol][syn][fseq]{lem:fseq-init} موجب
$\tuple{!A_0,\dotsc,!A_j}$ تے
$\tuple{!A_0,\dotsc,!A_k}$ ترتیب وار
$!A_j$ تے~$!A_k$ لئی تشکیل دے سلسلے نیں۔
ایہ $!A$ دے سلسلے دے پورے توں چھوٹے مُڈھلے
ذیلی سلسلے نیں؛ اوہناں دے آخری اشاریے
اصل سلسلے دے آخری اشاریے~$n$ توں گھٹ نیں۔
ایس لئی استقرائی مفروضے
موجب $!A_j$ تے~$!A_k$،~$\Frm[L]$ وچ نیں،
تے !!a{formula} دی تعریف موجب
$(!A_j \land !A_k)$ وی اوہدے وچ اے۔
باقی صورتاں لئی وی ایسّے طرح دی دلیل اے۔
\end{proof}

ٹرماں لئی تشکیل دے سلسلیاں دیاں خاصیتاں
!!{formula}s دے سلسلیاں دیاں خاصیتاں ورگیاں نیں۔

\begin{prop}
\ollabel{prop:fseq-trm-equiv}
$\Trm[L]$ اوہناں ساریاں $\Lang L$-لڑیاں $t$
دا سیٹ اے جیہناں دی شرط ایہ اے کہ $t$ لئی ٹرم تشکیل دا سلسلہ موجود اے۔
\end{prop}

\begin{proof}
مشق۔
\end{proof}

\begin{prob}
\olref[fol][syn][fseq]{prop:fseq-trm-equiv} ثابت کرو۔
اشارہ: \olref[fol][syn][fseq]{thm:fseq-frm-equiv}
دے ثبوت ورگی حکمتِ عملی ورتو۔
\end{prob}

تشکیل دے سلسلیاں وچ دو طرح دا ”فالتو مواد“ آ
سکدا اے: دہرائے گئے رکن، تے اوہ رکن جیہڑے
فارمولے یا ٹرم دی بناوٹ لئی غیر متعلق ہون۔
اسیں گھٹ توں گھٹ تشکیل دے سلسلیاں نوں ویکھ
کے دونیں طرح دے فالتو رکناں نوں ہٹا سکدے آں۔

\begin{defn}[گھٹ توں گھٹ تشکیل دے سلسلے]
\ollabel{defn:minimal-fseq}
تشکیل دا سلسلہ $\tuple{!A_0, \ldots, !A_n}$،
جیہڑا فارمولا~$!A$ لئی اے،~$!A$ لئی
\emph{گھٹ توں گھٹ تشکیل دا سلسلہ} اے جے
ہر ہور تشکیل دے سلسلے~$s$، جیہڑا~$!A$ لئی ہووے،
نوں لے کے ایس~$s$ دی لمبائی~$n+1$ توں وڈّی
یا برابر ہووے۔

ایسّے طرح تشکیل دا سلسلہ $\tuple{t_0, \ldots, t_n}$،
جیہڑا ٹرم~$t$ لئی اے،~$t$ لئی
\emph{گھٹ توں گھٹ تشکیل دا سلسلہ} اے جے
ہر ہور تشکیل دے سلسلے~$s$، جیہڑا~$t$ لئی ہووے،
نوں لے کے ایس~$s$ دی لمبائی~$n+1$ توں وڈّی
یا برابر ہووے۔
\end{defn}

دھیان دیو کہ کسے فارمولے یا ٹرم لئی اک توں
ودھ گھٹ توں گھٹ تشکیل دے سلسلے ہو سکدے نیں،
پر اوہناں وچ عین اوہی لڑیاں شامل ہون گیاں۔

\begin{prop}
\ollabel{prop:subformula-equivs}
اگلیاں گلاں ہم ارز نیں:
\begin{enumerate}
    \item $!B$،~$!A$ دا ذیلی !!{formula} اے۔
    \item $!B$،~$!A$ دے ہر تشکیل دے سلسلے وچ آؤندا اے۔
    \item $!B$،~$!A$ دے کسے گھٹ توں گھٹ تشکیل دے سلسلے وچ آؤندا اے۔
\end{enumerate}
\end{prop}

\begin{proof}
مشق۔
\end{proof}

\begin{prob}
\olref[fol][syn][fseq]{prop:subformula-equivs} ثابت کرو۔
\end{prob}

\begin{history}
تشکیل دے سلسلے Raymond Smullyan نے اپنی
کتاب \emph{First-Order Logic} \citep{Smullyan1968}
وچ پیش کیتے سن۔ تشکیل دے سلسلیاں دیاں ہور خاصیتاں
\citet{Zuckerman1973} نے ثابت کیتیاں۔
\end{history}
\endgroup
% /OLP-0156

% OLP-0157 content/first-order-logic/syntax-and-semantics/free-vars-sentences.tex
\begingroup
\def\olfilename{free-vars-sentences}
\expandafter\def\csname ol@sectioncs\endcsname{section}

\olfileid{fol}{syn}{fvs}

\olsection{آزاد \printtoken{P}{variable} تے \printtoken{P}{sentence}}

\begin{defn}[!!a{variable} دیاں آزاد آمدواں]
\ollabel{defn:free-occ}
!!a{formula} وچ !!a{variable} دیاں \emph{آزاد}
آمدواں اینج استقرائی طور تے تعریف کیتیاں جاندیاں نیں:
\begin{enumerate}
\item \indcase*{!A}{$!A$ ایٹمی اے}{$\indfrm$ وچ !!{variable}
  دیاں ساریاں آمدواں آزاد نیں۔}

\tagitem{prvNot}{\indcase{!A}{\lnot !B}{$\indfrm$ وچ !!{variable}
  دیاں آزاد آمدواں عین اوہو نیں جیہڑیاں $!B$ وچ آزاد نیں۔}}{}

\item \indcase{!A}{(!B \ast !C)}{$\indfrm$ وچ !!{variable}
  دیاں آزاد آمدواں $!B$ دیاں آزاد آمدواں
  تے~$!C$ دیاں آزاد آمدواں نوں ملا کے بندیاں نیں۔}

\tagitem{prvAll}{\indcase{!A}{\lforall[x][!B]}{$\indfrm$ وچ !!{variable}
  دیاں آزاد آمدواں~$!B$ دیاں ساریاں آزاد آمدواں نیں،
  سوائے~$x$ دیاں آمدواں دے۔}}{}

\tagitem{prvEx}{\indcase{!A}{\lexists[x][!B]}{$\indfrm$ وچ !!{variable}
  دیاں آزاد آمدواں~$!B$ دیاں ساریاں آزاد آمدواں نیں،
  سوائے~$x$ دیاں آمدواں دے۔}}{}
\end{enumerate}
\end{defn}

\begin{defn}[پابند متغیر]
فارمولا~$!A$ وچ !!a{variable} دی کوئی آمد
\emph{پابند} اے جے اوہ آزاد نہیں۔
\end{defn}

\begin{prob}
\olref[fol][syn][fvs]{defn:free-occ} دے طریقے تے
پابند متغیراں دیاں آمدواں دی استقرائی تعریف دیو۔
\end{prob}

\begin{defn}[دائرۂ اثر]
\iftag{prvAll}{جے $\lforall[x][!B]$، فارمولا~$!A$ وچ کسے ذیلی
  فارمولے دی اک آمد اے، تاں~$!B$ دی اوہدے مطابق آمد نوں~$!A$ وچ
  $\lforall[x]$ دی اوہدے مطابق آمد دا
  \emph{دائرۂ اثر} آکھدے آں۔ \iftag{prvEx}{$\lexists[x]$ لئی وی
  ایہو گل اے۔}{}}{جے
  $\lexists[x][!B]$، فارمولا~$!A$ وچ کسے ذیلی فارمولے
  دی اک آمد اے، تاں~$!B$ دی اوہدے مطابق آمد نوں~$!A$ وچ
  $\lexists[x]$ دی اوہدے مطابق آمد دا
  \emph{دائرۂ اثر} آکھدے آں۔}

جے $!B$ مسوّر دی کسے آمد
\iftag{prvAll}{$\lforall[x]$\iftag{prvEx}{ یا
    $\lexists[x]$}{}}{$\lexists[x]$} دا دائرۂ اثر~$!A$ وچ اے،
تاں $x$ دیاں~$!B$ وچ آزاد آمدواں
\iftag{prvAll}{$\lforall[x][!B]$\iftag{prvEx}{ تے
$\lexists[x][!B]$}{}}{$\lexists[x][!B]$} وچ پابند نیں۔
اسیں آکھدے آں کہ ایہ آمدواں مذکورہ مسوّر دی آمد
\emph{نال پابند} نیں۔
\end{defn}

\begin{ex}
ایہ فارمولا ویکھو:
\[
\lexists[\Obj v_0][\underbrace{\Atom{\Obj A^2_0}{\Obj v_0,\Obj v_1}}_{!B}] 
\]
$!B$، $\lexists[\Obj v_0]$ دا دائرۂ اثر اے۔
ایہ مسوّر $\Obj v_0$ دی آمد $!B$ وچ
پابند کردا اے، پر $\Obj v_1$ دی آمد نوں نہیں۔
ایس لئی ایس مثال وچ $\Obj v_1$ آزاد متغیر اے۔


ہُن ویکھو کہ کسے ہور پیچیدہ !!{formula}~$!A$
وچ ایہ گل کیویں کم کردی اے:
\[
\lforall[\Obj v_0][\underbrace{(\Atom{\Obj A^1_0}{\Obj v_0} \lif
    \Atom{\Obj A^2_0}{\Obj v_0, \Obj v_1})}_{!B}] \lif \lexists[\Obj
  v_1][\underbrace{(\Atom{\Obj A^2_1}{\Obj v_0, \Obj v_1} \lor \lforall[\Obj v_0][\overbrace{\lnot \Atom{\Obj A^1_1}{\Obj v_0}}^{!D}])}_{!C}]
\]
$!B$ پہلے $\lforall[\Obj v_0]$ دا دائرۂ اثر اے،
$!C$، $\lexists[\Obj v_1]$ دا دائرۂ اثر اے، تے
$!D$ دوجے $\lforall[\Obj v_0]$ دا دائرۂ اثر اے۔
پہلا $\lforall[\Obj v_0]$، $\Obj v_0$
دیاں آمدواں~$!B$ وچ پابند کردا اے؛
$\lexists[\Obj v_1]$، $\Obj v_1$ دی
آمد~$!C$ وچ پابند کردا اے؛ تے دوجا
$\lforall[\Obj v_0]$، $\Obj v_0$
دی آمد~$!D$ وچ پابند کردا اے۔
$\Obj v_1$ دی پہلی آمد تے $\Obj v_0$ دی چوتھی
آمد~$!A$ وچ آزاد نیں۔ $\Obj v_0$ دی آخری آمد
$!D$ وچ آزاد اے، پر $!C$ تے~$!A$ وچ پابند اے۔
\end{ex}

\begin{defn}[جملہ]
!!^a{formula}~$!A$، \article{sentence}
\emph{!!{sentence}} اُدوں تے صرف اُدوں اے جدوں
اوہدے وچ !!{variable}s دیاں کوئی آزاد آمدواں نہ ہون۔
\end{defn}

% منجمد ماخذ دا غیر مطبوعہ اداراتی نوٹ: ایتھے ہور مثالاں شامل کرن دا اشارہ اے، پر اصل وچ اوہ مثالاں موجود نہیں۔
\endgroup
% /OLP-0157

% OLP-0158 content/first-order-logic/syntax-and-semantics/substitution.tex
\begingroup
\def\olfilename{substitution}
\expandafter\def\csname ol@sectioncs\endcsname{section}

\olfileid{fol}{syn}{sub}

\olsection{تھاں بدل}

\begin{defn}[ٹرم وچ تھانواں بدلنا]
اسیں $\Subst{s}{t}{x}$ نوں، یعنی $t$ نوں~$x$ دی
ہر آمد دی تھاں~$s$ وچ \emph{رکھن} دا نتیجہ، بازگشتی
طور تے تعریف کردے آں:
\begin{enumerate}
\item \indcase{s}{c}{$\Subst{\indfrm}{t}{x}$ بس $s$ اے۔}

\item \indcase{s}{y}{$\Subst{\indfrm}{t}{x}$ وی بس~$s$ اے،
  بشرطیکہ $y$ متغیر ہووے تے $y \not\ident x$ ہووے۔}

\item \indcase{s}{x}{$\Subst{\indfrm}{t}{x}$،~$t$ اے۔}

\item\indcase{s}{\Atom{f}{t_1, \dots, t_n}}{$\Subst{\indfrmp}{t}{x}$،
  $\Atom{f}{\Subst{t_1}{t}{x}, \dots, \Subst{t_n}{t}{x}}$ اے۔}
\end{enumerate}
\end{defn}

\begin{defn}
ٹرم~$t$، $x$ لئی~$!A$ وچ \emph{!!{free for}} اے
جے~$x$ دی کوئی آزاد آمد~$!A$ وچ کسے ایسے مسوّر دے
دائرۂ اثر وچ نہ آوے جیہڑا~$t$ دے کسے متغیر نوں پابند کردا ہووے۔
\end{defn}

\begin{ex} ~
\begin{enumerate}
\item $\Obj v_8$، $\Obj v_1$ لئی
$\lexists[\Obj
  v_3]\Atom{\Obj A^2_4}{\Obj v_3,\Obj v_1}$ وچ بدل لئی آزاد اے۔

\item $\Obj f^2_1(\Obj v_1, \Obj v_2)$، $\Obj
  v_0$ لئی $\lforall[\Obj v_2]\Atom{\Obj A^2_4}{\Obj v_0,\Obj v_2}$
  وچ بدل لئی \emph{آزاد نہیں} اے۔
\end{enumerate}
\end{ex}

\begin{defn}[!!a{formula} وچ تھانواں بدلنا]
جے $!A$، !!a{formula} اے، $x$~!!a{variable} اے،
تے $t$~اوہ ٹرم اے جیہڑا~$x$ لئی~$!A$ وچ
!!{free for} اے، تاں $\Subst{!A}{t}{x}$،
$t$ نوں~$x$ دیاں ساریاں آزاد آمدواں دی تھاں
$!A$ وچ رکھن دا نتیجہ اے۔
\begin{enumerate}
\tagitem{prvFalse}{\indcase{!A}{\lfalse}{$\Subst{\indfrm}{t}{x}$،
    $\lfalse$ اے۔}}{}

\tagitem{prvTrue}{\indcase{!A}{\ltrue}{$\Subst{\indfrm}{t}{x}$،
    $\ltrue$ اے۔}}{}

\item \indcase{!A}{\Atom{P}{t_1,\dots,
    t_n}}{$\Subst{\indfrm}{t}{x}$،
    $\Atom{P}{\Subst{t_1}{t}{x},
    \dots, \Subst{t_n}{t}{x}}$ اے۔}

\item \indcase{!A}{\eq[t_1][t_2]}{$\Subst{\indfrmp}{t}{x}$،
  $\Subst{t_1}{t}{x} = \Subst{t_2}{t}{x}$ اے۔}

\tagitem{prvNot}{\indcase{!A}{\lnot !B}{$\Subst{\indfrmp}{t}{x}$،
    $\lnot \Subst{!B}{t}{x}$ اے۔}}{}

\tagitem{prvAnd}{\indcase{!A}{(!B \land
    !C)}{$\Subst{\indfrmp}{t}{x}$،
    $(\Subst{!B}{t}{x} \land
    \Subst{!C}{t}{x})$ اے۔}}{}

\tagitem{prvOr}{\indcase{!A}{(!B \lor
    !C)}{$\Subst{\indfrmp}{t}{x}$،
    $(\Subst{!B}{t}{x} \lor
    \Subst{!C}{t}{x})$ اے۔}}{}

\tagitem{prvIf}{\indcase{!A}{(!B \lif
    !C)}{$\Subst{\indfrmp}{t}{x}$،
    $(\Subst{!B}{t}{x} \lif
    \Subst{!C}{t}{x})$ اے۔}}{}

\tagitem{prvIff}{\indcase{!A}{(!B \liff
    !C)}{$\Subst{\indfrmp}{t}{x}$،
    $(\Subst{!B}{t}{x} \liff
    \Subst{!C}{t}{x})$ اے۔}}{}

\tagitem{prvAll}{
\indcase{!A}{\lforall[y][!B]}{$\Subst{\indfrmp}{t}{x}$
    دا نتیجہ $\lforall[y][\Subst{!B}{t}{x}]$ اے،
    بشرطیکہ $y$، $x$ توں وکھرا متغیر ہووے؛
    نہیں تاں $\Subst{\indfrmp}{t}{x}$
    بس $\indfrm$ اے۔}}{}

\tagitem{prvEx}{
\indcase{!A}{\lexists[y][!B]}{$\Subst{\indfrmp}{t}{x}$
    دا نتیجہ $\lexists[y][\Subst{!B}{t}{x}]$ اے،
    بشرطیکہ $y$، $x$ توں وکھرا متغیر ہووے؛
    نہیں تاں $\Subst{\indfrmp}{t}{x}$
    بس $\indfrm$ اے۔}}{}
\end{enumerate}
\end{defn}

\begin{explain}
دھیان دیو کہ تھانواں بدلنا بے اثر وی ہو سکدا اے:
جے $x$، $!A$ وچ کِتے وی نہ آوے، تاں
$\Subst{!A}{t}{x}$ بس~$!A$ اے۔

$t$ دا~$x$ لئی~$!A$ وچ !!{free for} ہونا
اگلی مثال ورگیاں صورتاں کڈھن لئی ضروری اے۔
جے $!A \ident \lexists[y][x < y]$
تے $t \ident y$، تاں $\Subst{!A}{t}{x}$
دا نتیجہ $\lexists[y][y <
  y]$ بن جاندا۔ ایتھے آزاد متغیر $y$،
تھانواں بدلّن توں بعد مسوّر $\lexists[y]$ دی
پابندی وچ ”پھس“ جاندا اے، جیہڑی گل ساڈے لئی
ناپسندیدہ اے۔ مثال لئی، اسیں چاہندے آں کہ
جدوں وی $\lforall[x][!B]$ پورا اترے، تاں
$\Subst{!B}{t}{x}$ وی پورا اترے۔ پر
$\lforall[x][\lexists[y][x < y]]$ ویکھو
(ایتھے $!B$، $\lexists[y][x <
  y]$ اے)۔ ایہ ایسا جملہ اے جیہڑا، مثال لئی،
قدرتی عدداں بارے سچ اے: ہر عدد~$x$ لئی اوہدے
توں وڈّا کوئی عدد~$y$ موجود اے۔ جے اسیں
$y$ نوں~$x$ دی تھاں رکھن دی اجازت دیندے،
تاں $\Subst{!B}{y}{x} \ident \lexists[y][y < y]$
ملدا، جیہڑا جھوٹ اے۔ اسیں ایہ صورت
ایہ شرط رکھ کے روکدے آں کہ~$t$ وچ موجود
کوئی وی آزاد متغیر~$!A$ دے کسے مسوّر نال
پابند نہ ہو جائے۔
\end{explain}

بوجھل علامت نگاری توں بچن لئی اسیں اکثر ایہ رواج
ورتدے آں: جے $!A$ اک !!a{formula} اے جیہدے وچ
!!{variable}~$x$ آزاد ہو سکدا اے، تاں
اسیں ایہ وکھاؤن لئی~$!A(x)$ وی لکھدے آں۔
جدوں ایہ صاف ہووے کہ کیہڑا $!A$ تے~$x$
مراد نیں، تے $t$ اک ٹرم ہووے (فرض کیتا
جاندا اے کہ اوہ $x$ لئی~$!A(x)$ وچ بدل لئی
آزاد اے)، تاں اسیں $!A(t)$ نوں
$\Subst{!A}{t}{x}$ دا مختصر روپ لکھدے آں۔
مثال لئی، اسیں آکھ سکدے آں، ”$!A(t)$ نوں
$\lforall[x][!A(x)]$ دی اک مخصوص صورت آکھدے آں۔“
ایس دا مطلب اے کہ جے $!A$ کوئی وی
!!{formula} ہووے، $x$~!!a{variable} ہووے،
تے $t$ اک ایسا ٹرم ہووے جیہڑا~$x$ لئی~$!A$ وچ
بدل لئی آزاد ہووے، تاں $\Subst{!A}{t}{x}$
فارمولا~$\lforall[x][!A]$ دی اک مخصوص صورت اے۔
\endgroup
% /OLP-0158

% OLP-0645 content/first-order-logic/syntax-and-semantics/introduction.tex
\begingroup
\def\olfilename{introduction}
\expandafter\def\csname ol@sectioncs\endcsname{section}

\olfileid{fol}{syn}{int}

\olsection{جان پچھان}

پہلے درجے دی منطق دا نظریہ تے اوہدے بارے اُتلا نظریہ بناؤن لئی سانوں پہلاں ایہدیاں عبارتاں دی بناوٹ تے معنی دی تعریف کرنی پئے گی۔ پہلے درجے دی منطق دیاں عبارتاں ٹرماں تے !!{formula}s نیں۔ ٹرم !!{variable}s، !!{constant}s تے !!{function}s توں بندے نیں۔ اگّوں !!^{formula}s، !!{predicate}s نوں ٹرماں نال جوڑ کے بندے نیں (ایہ سبھ توں نِکّے، «جوہری» !!{formula}s نیں)، تے فیر جوہری !!{formula}s توں منطقی جوڑاں تے مسوّراں دی مدد نال ہور گنجل دار فارمولا بنا سکدے آں۔ بناوٹ دے قاعدے دَسن دے کئی وکھ وکھ ڈھنگ نیں۔ اسیں بس اک ممکن ڈھنگ دِندے آں۔ دوجے نظام وکھ وکھ نشان چُنن گے، بنیادی جوڑاں دے وکھ وکھ سیٹ رکھن گے، قوساں وکھرے طریقے نال ورتن گے (یا بالکل نہیں ورتن گے، جیویں پولش لکھت وچ)۔ پر سارے طریقیاں دی سانجھی گل ایہ اے کہ بناوٹ دے قاعدے ٹرماں تے !!{formula}s دے سیٹ دی \emph{استقرائی} تعریف کردے نیں۔ جے ایہ ٹھیک طرح ہووے، تاں بناوٹ دے قاعدیاں مطابق ہر عبارت بنیادی طور اُتے بس اک ای ڈھنگ نال بن سکدی اے۔ عبارت دی ایہ استقرائی تعریف، جیہدے نال اوہ \emph{اکّو ڈھنگ نال پڑھی جا سکدی} اے، سانوں ایہناں عبارتاں دے معنی وی اوسے ڈھنگ نال، یعنی استقرائی تعریف نال، دین دی کھُل دیندی اے۔

عبارتاں نوں معنی دینا علمِ معنی دا کم اے۔ علمِ معنی دا مرکزی خیال !!a{structure} وچ پورا اترنا اے۔ !!^a{structure} زبان دیاں بنیادی اکائیاں نوں معنی دیندی اے۔ !!a{domain} شےواں دا غیر خالی سیٹ اے۔ مسوّراں دی تعبیر ایویں ہوندی اے کہ اوہ ایس دائرے دیاں شےواں اُتے پھردے نیں۔ !!{constant}s نوں دائرے دے رکن دِتّے جاندے نیں، !!{function}s نوں !!{domain} توں اوسے دائرے ول فنکشن دِتّے جاندے نیں، تے !!{predicate}s نوں !!{domain} اُتے تعلق دِتّے جاندے نیں۔ !!{domain} تے بنیادی لفظاں لئی دِتّیاں تعبیراں رَل کے !!a{structure} بناؤندیاں نیں۔ !!^{variable}s، !!{formula}s وچ آ سکدے نیں، تے معنی دین لئی سانوں !!{element}s وی !!{domain} توں ایہناں متغیراں نوں دین پَین گے۔ ایہ متغیر قدر بندی اے۔ آخر پورا اترن دا تعلق ایہناں گلّاں نوں جوڑدا اے۔ !!^a{formula}، !!a{structure}~$\Struct{M}$ وچ متغیر قدر بندی~$s$ دے لحاظ نال پورا اُتر سکدا اے، جیہنوں $\Sat{M}{!A}[s]$ لکھدے آں۔ ایس تعلق دی وی~$!A$ دی بناوٹ اُتے استقرا نال تعریف ہوندی اے۔ منطقی جوڑاں دیاں سچائی سارنیاں ورت کے، مثال لئی، $!A \land !B$ دا پورا اترنا، $!A$ تے ~$!B$ دے پورا اترن (یا نہ اترن) نال تعریف ہوندا اے۔ فیر پتا لگدا اے کہ جے !!{formula}~$!A$ اک !!a{sentence} اے، یعنی اوہدے وچ کوئی آزاد متغیر نہیں، تاں متغیر قدر بندی دا کوئی اثر نہیں۔ سو اسیں !!{sentence}s دے !!{structure}s وچ بس پورا اترن (یا نہ اترن) دی گل کر سکدے آں۔

جملیاں لئی پورا اترن دے تعلق $\Sat{M}{!A}$ دی بنیاد اُتے اسیں معنی دیاں بنیادی اصطلاحاں، صحت، لازمی نتیجہ تے پورا اتر سکنا، دی تعریف کر سکدے آں۔ اک جملہ صحیح اے، $\Entails !A$، جے ہر ساخت وچ پورا اُترے۔ اوہ !!{sentence}s دے کسے سیٹ دا لازمی نتیجہ اے، $\Gamma \Entails !A$، جے ہر !!{structure} جیہدے وچ سارے !!{sentence}s، جیہڑے~$\Gamma$ وچ نیں، پورے اُترن، اوہدے وچ~$!A$ وی پورا اُترے۔ جملیاں دا سیٹ پورا اُتر سکدا اے جے کوئی !!{structure} ایہو جیہی ہووے جیہدے وچ اوہدے سارے !!{sentence}s اکٹھے پورے اُترن۔ کیوں جے !!{formula}s دی استقرائی تعریف ہوندی اے تے اگّوں پورا اترن دی تعریف !!{formula}s دی بناوٹ اُتے استقرا نال ہوندی اے، اسیں استقرا ورت کے اپنے علمِ معنی دیاں خاصیتاں ثابت کر سکدے آں تے ایہناں معنی دیاں اصطلاحاں دا آپس دا تعلق وکھا سکدے آں۔
\endgroup
% /OLP-0645
\OLEndChapterHook



\par\noindent\emph{متبادل ترتیب۔} نحو تے معنویات اکّو باب وچ (\texttt{syntax-and-semantics.tex})۔ ایہ ترتیب تدوین جوگ ماخذ فائل وچ موجود اے۔
\endgroup
% /OLP-0149

% OLP-0159 content/first-order-logic/syntax-and-semantics/semantics.tex
\begingroup
\def\olfilename{semantics}
\expandafter\def\csname ol@sectioncs\endcsname{section}

\olchapter{fol}{sem}{پہلے درجے دی منطق دیاں معنویات}

% OLP-0160 content/first-order-logic/syntax-and-semantics/intro-semantics.tex
\begingroup
\def\olfilename{intro-semantics}
\expandafter\def\csname ol@sectioncs\endcsname{section}

\olfileid{fol}{syn}{its}

\olsection{تعارف}

اظہاراں دے معنی متعین کرنا معنویات دا کم اے۔ معنویات دا مرکزی
تصور !!a{structure} وچ پورا اترن دا اے۔ !!^a{structure} زبان دے
بنیادی حصیاں نوں معنی دیندی اے: !!a{domain} اشیاء دا اک غیر خالی
سیٹ اے۔ مسوّراں دی تعبیر ایس دائرے اُتے چلن دے طور تے ہوندی اے،
!!{constant}s نوں دائرے دے عنصر دتے جاندے نیں، !!{function}s نوں
اوہناں دی مقامیت مطابق !!{domain} دے عناصر دے ترتیب وار ٹپل توں
ایس دائرے دے عناصر ول جان والے فنکشن دتے جاندے نیں، تے
!!{predicate}s نوں !!{domain} اُتے تعلق دتے جاندے نیں۔
% PNB-SEM-001: ماخذ دا دائرے توں اوہدے اپنے آپ ول والا بیان صرف اک-مقامی فنکشن لئی اے؛ ایتھے مقامیت مطابق ٹپل توں دائرے ول فنکشن دتے نیں۔
!!{domain} تے بنیادی لفظیات نوں دتیاں قدراں مل کے !!a{structure}
بناندیاں نیں۔ !!^{variable}s، !!{formula}s وچ آ سکدے نیں؛ معنوی
تعبیر لئی سانوں اوہناں نوں وی !!{domain} دے !!{element}s دینے پیندے نیں---
ایہ متغیر قدر بندی اے۔ پورا اترن دا تعلق آخرکار ایہ ساریاں گلاں
جوڑدا اے۔ !!^a{formula}، !!a{structure}~$\Struct{M}$ وچ
!!a{variable} قدر بندی~$s$ دے لحاظ نال پورا اتر سکدا اے؛
اس نوں $\Sat{M}{!A}[s]$ لکھدے آں۔ ایہ تعلق وی~$!A$ دی ساخت
اُتے استقرا نال تعریف ہوندا اے: منطقی رابطیاں دے صدقی جدول
ورت کے، مثال لئی، $(!A \land !B)$ دا پورا اترنا $!A$ تے~$!B$
دے پورا اترن (یا نہ اترن) دے لحاظ نال تعریف ہوندا اے۔ فیر
ایہ نکلدا اے کہ جے !!{formula}~$!A$، !!a{sentence} ہووے،
یعنی اوہدے وچ کوئی آزاد متغیر نہ ہووے، تاں !!{variable} قدر بندی
دا فرق نہیں پیندا؛ ایس لئی اسیں !!{sentence}s دے !!{structure}s
وچ بس پورا اترن (یا نہ اترن) دی گل کر سکدے آں۔

!!{sentence}s لئی پورا اترن دے تعلق $\Sat{M}{!A}$ دی بنیاد اُتے
اسیں منطقی اعتبار، منطقی نتیجے تے پورا ہون جوگتا دے بنیادی معنوی
تصورات تعریف کر سکدے آں۔ !!^a{sentence} معتبر اے، $\Entails !A$،
جے ہر !!{structure} اوہنوں پورا کرے۔ !!{sentence}s دے کسے سیٹ توں
اوہ منطقی نتیجے دے طور تے نکلدا اے، $\Gamma \Entails !A$، جے ہر
اوہ !!{structure} جیہڑی~$\Gamma$ دے سارے !!{sentence}s پورے کردی
اے،~$!A$ نوں وی پورا کرے۔ تے !!{sentence}s دا کوئی سیٹ پورا ہون
جوگ اے جے کوئی اک !!{structure} اوہدے سارے !!{sentence}s نوں
اکو ویلے پورا کرے۔ کیوں جے !!{formula}s استقرائی طور تے تعریف
ہوندے نیں، تے پورا اترنا وی !!{formula}s دی ساخت اُتے استقرا
نال تعریف ہوندا اے، اسیں استقرا ورت کے اپنی معنویات دیاں خاصیتاں
ثابت کر سکدے آں تے متعین کردہ معنوی تصورات دے باہمی تعلق وکھا
سکدے آں۔
\endgroup
% /OLP-0160

% OLP-0161 content/first-order-logic/syntax-and-semantics/structures.tex
\begingroup
\def\olfilename{structures}
\expandafter\def\csname ol@sectioncs\endcsname{section}

\olfileid{fol}{syn}{str}

\olsection{پہلے درجے دیاں زباناں لئی \printtoken{P}{structure}}

\begin{explain}
پہلے درجے دیاں زباناں اپنے آپ وچ \emph{بنا تعبیر دے} ہوندیاں نیں:
!!{constant}s، !!{function}s تے !!{predicate}s نال کوئی خاص
معنی پہلے توں جڑے نہیں ہوندے۔ معنی \article{structure}
\emph{!!{structure}} متعین کرکے دتے جاندے نیں۔ اوہ
\emph{دائرہ} متعین کردی اے، یعنی اوہ اشیاء جنہاں ول
!!{constant}s اشارہ کردے نیں، جنہاں اُتے !!{function}s کم کردے
نیں، تے جنہاں اُتے مسوّرے چلدے نیں۔ نال ایہ وی متعین ہوندا اے
کہ کیہڑے !!{constant}s کیہڑیاں اشیاء ول اشارہ کردے نیں،
!!a{function} اشیاء نوں دوجیاں اشیاء ول کیویں بھیجدا اے، تے
!!{predicate}s کیہڑیاں اشیاء اُتے لاگو ہوندے نیں۔
!!^{structure}s منطق وچ \emph{معنوی} تصورات، مثلاً منطقی نتیجے،
اعتبار تے پورا ہون جوگتا دی بنیاد نیں۔ علمی عبارت وچ انہاں نوں
کدے ``ساختاں،'' کدے ``تعبیراں،'' تے کدے ``نمونے'' آکھیا جاندا اے۔
\end{explain}

\begin{defn}[!!^{structure}s]
\Article{structure} \emph{!!{structure}}~$\Struct M$، پہلے درجے دی
زبان $\Lang{L}$ لئی، ہیٹھ لکھیاں چیزاں توں بندی اے:
\begin{enumerate}
\item \emph{دائرہ:} اک غیر خالی سیٹ، $\Domain M$
\item \emph{!!{constant}s دی تعبیر:} ہر !!{constant}~$c$ لئی
  جیہڑا $\Lang{L}$ وچ اے، !!a{element} $\Assign{c}{M} \in \Domain M$
\item \emph{!!{predicate}s دی تعبیر:} ہر $n$-مقامی
  !!{predicate}~$R$ لئی جیہڑا $\Lang{L}$ وچ اے ($\eq$ توں علاوہ)، اک $n$-مقامی
  تعلق $\Assign{R}{M} \subseteq \Domain{M}^n$
\item \emph{!!{function}s دی تعبیر:} ہر $n$-مقامی
  !!{function}~$f$ لئی جیہڑا $\Lang{L}$ وچ اے، اک $n$-مقامی فنکشن $\Assign{f}{M}
  \colon \Domain{M}^n \to \Domain{M}$
\end{enumerate}
\end{defn}

\begin{ex}
حساب دی زبان لئی !!^a{structure}~$\Struct M$ وچ اک سیٹ،
$\Domain M$ دا اک عنصر $\Assign{\Obj 0}{M}$ جیہڑا
!!{constant}~$\Obj 0$ دی تعبیر اے، اک یک-مقامی فنکشن
$\Assign{\Obj \prime}{M} \colon \Domain{M} \to \Domain M$،
دو دو-مقامی فنکشن $\Assign{\Obj +}{M}$ تے $\Assign{\Obj
    \times}{M}$ (دونیں $\Domain M^2 \to \Domain M$)، تے اک
دو-مقامی تعلق $\Assign{\Obj <}{M} \subseteq \Domain{M}^2$ ہوندے نیں۔

ایہو جہی ساخت دی اک صاف مثال ایہ اے:
\begin{enumerate}
\item $\Domain N = \Nat$
\item $\Assign{\Obj 0}{N} = 0$
\item $\Assign{\Obj \prime}{N}(n) = n + 1$ ہر $n \in \Nat$ لئی
\item $\Assign{\Obj +}{N}(n, m) = n + m$ ہر $n, m \in \Nat$ لئی
\item $\Assign{\Obj \times}{N}(n, m) = n\cdot m$ ہر $n, m \in \Nat$ لئی
\item $\Assign{\Obj <}{N} = \Setabs{\tuple{n, m}}{n \in \Nat, m \in
  \Nat, n < m}$
\end{enumerate}
ایس طرح تعریف کیتی ساخت~$\Struct N$ نوں، $\Lang L_A$ لئی،
\emph{حساب دا معیاری نمونہ} آکھدے آں، کیوں جے ایہ~$\Lang L_A$
دیاں غیر منطقی ثابتاں دی تعبیر ٹھیک اوہو جیہی کردی اے جیہڑی
تُسی توقع کردے او۔

پر~$\Lang L_A$ لئی ہور بہت ساریاں !!{structure}s وی ممکن نیں۔
مثلاً اسیں دائرے لئی صحیح عدداں دا سیٹ~$\Int$ لے سکدے آں،
قدرتی اعداد~$\Nat$ دی تھاں، تے $\Obj
0$، $\Obj \prime$، $\Obj +$، $\Obj \times$ تے $\Obj <$
دیاں تعبیراں اوہدے مطابق متعین کر سکدے آں۔ بلکہ~$\Lang L_A$
لئی ایسی ساختاں وی بنا سکدے آں جنہاں دا عدداں نال دور دا وی
واسطہ نہ ہووے۔
\end{ex}

\begin{ex}
ساخت~$\Struct M$ نوں، سیٹ نظریے دی زبان~$\Lang L_Z$ لئی،
بس اک سیٹ تے اک دو-مقامی تعلق چاہیدا اے۔ ایس لئی، خالص فنی
لحاظ نال، مثلاً لوکاں دا سیٹ تے تعلق ``$x$، $y$ توں وڈی عمر
دا اے''، $\Lang L_Z$ لئی !!a{structure} بن سکدا اے؛ ایویں
$\Nat$ تے $n \ge m$ دا تعلق وی، جدوں $n, m \in \Nat$ ہون۔

$\Lang L_Z$ لئی اک خاص دلچسپ !!{structure} اوہ اے جیہدے دائرے دے
!!{element}s خود سیٹ نیں تے $\Obj \in$ دی تعبیر واقعی تعلق
``$x$،~$y$ دا !!a{element} اے'' اے۔ ایہ \emph{وراثتاً محدود سیٹاں}
دی !!{structure}~$\Struct{{HF}}$ اے:
\begin{enumerate}
\item $\Domain{{{HF}}} = \emptyset \cup \Pow{\emptyset} \cup
  \Pow{\Pow{\emptyset}} \cup \Pow{\Pow{\Pow{\emptyset}}} \cup \dots$;
\item $\Assign{\Obj \in}{{{HF}}} = \Setabs{\tuple{x, y}}{x, y \in
  \Domain{{{HF}}}, x \in y}$.
\end{enumerate}
\end{ex}

\begin{digress}
کیہڑی شے نوں !!a{structure} منن دے ضابطے ساڈی منطق اُتے اثر
پاندے نیں۔ مثال لئی، خالی دائریاں دی ممانعت، مسوّر والے جملیاں
دے پورا اترن (یا سچ ہون) دی عام تعبیر ہِٹھ، ایہ یقینی بناندی اے
کہ $\lexists[x][(!A(x) \lor \lnot !A(x))]$ معتبر اے---یعنی
منطقی سچائی اے۔ تے ایہ شرط کہ سارے !!{constant}s دائرے دی
کسے شے ول اشارہ کرن، وجودی تعمیم نوں استدلال دا درست نمونہ
بناندی اے: $!A(a)$، ایس لئی $\lexists[x][!A(x)]$۔ جے اسیں
ناواں نوں دائرے توں باہر دیاں اشیاء ول اشارہ کرن دی اجازت
دیندے، یا اشارہ نہ کرن دی، تاں اسیں \emph{آزاد منطق} ول جا
رہے ہوندے؛ اوتھے وجودی تعمیم لئی اک وادھو مقدمہ چاہیدا اے:
$!A(a)$ تے $\lexists[x][\eq[x][a]]$، ایس لئی
$\lexists[x][!A(x)]$۔
\end{digress}
\endgroup
% /OLP-0161

% OLP-0162 content/first-order-logic/syntax-and-semantics/covered-structures.tex
\begingroup
\def\olfilename{covered-structures}
\expandafter\def\csname ol@sectioncs\endcsname{section}

\olfileid{fol}{syn}{cov}

\olsection{پہلے درجے دیاں زباناں لئی ڈھکیاں ہوئیاں \printtoken{P}{structure}}

\begin{explain}
یاد رکھو، اک ٹرم \emph{بند} اے جے اوہدے وچ کوئی !!{variable}s نہیں۔
\end{explain}

\begin{defn}[بند ٹرماں دی !!^{value}]
جے $t$، زبان~$\Lang L$ دا بند ٹرم اے تے $\Struct M$،
~$\Lang L$ لئی !!{structure} اے، تاں اوہدی
\emph{!!{value}}~$\Value{t}{M}$ ایویں تعریف ہوندی اے:
\begin{enumerate}
\item جے $t$ بس !!{constant}~$c$ اے، تاں $\Value{c}{M} = \Assign{c}{M}$۔
\item جے $t$ دی صورت $\Atom{f}{t_1, \ldots, t_n}$ اے، تاں
  \[
  \Value{t}{M} = \Assign{f}{M}(\Value{t_1}{M}, \ldots,
  \Value{t_n}{M}).
  \]
\end{enumerate}
\end{defn}

\begin{defn}[ڈھکی ہوئی !!{structure}]
!!{structure} \emph{ڈھکی ہوئی} اے جے دائرے دا ہر عنصر
کسے بند ٹرم دی !!{value} ہووے۔
\end{defn}

\begin{ex}
فرض کرو~$\Lang L$ ایسی زبان اے جیہدے وچ !!{constant}s
$\Obj{zero}$، $\Obj{one}$، $\Obj{two}$، \dots، دو-مقامی
!!{predicate}~$<$، تے دو-مقامی !!{function}s $+$ تے $\times$
نیں۔ فیر !!a{structure}~$\Struct
M$،~$\Lang L$ لئی، اوہ اے جیہدا دائرہ $\Domain M = \{0, 1, 2, \ldots
\}$ اے، تے تعبیراں $\Assign{\Obj{zero}}{M} = 0$،
$\Assign{\Obj{one}}{M} = 1$، $\Assign{\Obj{two}}{M} = 2$،
تے ایس طرح اگے نیں۔ دو-مقامی تعلق دی علامت $<$ لئی سیٹ
$\Assign{<}{M}$ اوہناں ساریاں جوڑیاں دا سیٹ اے
$\tuple{c_1, c_2} \in \Domain{M}^2$ جنہاں وچ
$c_1$،~$c_2$ توں گھٹ اے: مثال لئی، $\tuple{1, 3} \in
\Assign{<}{M}$ پر $\tuple{2, 2} \notin \Assign{<}{M}$۔
دو-مقامی !!{function} $+$ لئی، $\Assign{+}{M}$ نوں معمول
مطابق تعریف کرو---مثلاً، $\Assign{+}{M}(2,3)$ دا حاصل~$5$
اے؛ دو-مقامی !!{function}~$\times$ لئی وی ایویں۔ ایس لئی
$\Obj{four}$ دی !!{value} بس~$4$ اے، تے
$\times(\Obj{two},
+(\Obj{three},\Obj{zero}))$ دی !!{value} (یا درمیان-علامتی
لکھت وچ، $\Obj{two} \times
(\Obj{three} + \Obj{zero})$) ایہ اے:
% PNB-SEM-002: منجمد ماخذ وچ پہلی مساوات دی علامت دو واری اے؛ ایتھے صرف اگلی سطر دی اک علامت رکھی اے۔
\begin{multline*}
\Value{\times(\Obj{two}, +(\Obj{three},\Obj{zero}))}{M} \\
\begin{aligned}
& =\Assign{\times}{M}(\Value{\Obj{two}}{M}, \Value{+(\Obj{three}, \Obj{zero})}{M})\\
& = \Assign{\times}{M}(\Value{\Obj{two}}{M}, \Assign{+}{M}(\Value{\Obj{three}}{M},
\Value{\Obj{zero}}{M})) \\
& = \Assign{\times}{M}(\Assign{\Obj{two}}{M}, \Assign{+}{M}(\Assign{\Obj{three}}{M},
\Assign{\Obj{zero}}{M})) \\
& = \Assign{\times}{M}(2, \Assign{+}{M}(3, 0)) \\
& = \Assign{\times}{M}(2, 3) \\
& = 6
\end{aligned}
\end{multline*}
\end{ex}

\begin{prob}
کی حساب دا معیاری نمونہ $\Struct N$ ڈھکیا ہویا اے؟
اپنا جواب سمجھاؤ۔
\end{prob}
\endgroup
% /OLP-0162

% OLP-0163 content/first-order-logic/syntax-and-semantics/satisfaction.tex
\begingroup
\def\olfilename{satisfaction}
\expandafter\def\csname ol@sectioncs\endcsname{section}

\olfileid{fol}{syn}{sat}

\olsection{\article{formula} \printtoken{S}{formula} دا
  \article{structure} \printtoken{S}{structure} وچ پورا اترنا}

\begin{explain}
ٹرم تے !!{formula}s اک پاسے، تے !!{structure}s دوجے پاسے
ہون، تاں انہاں نوں جوڑن والے بنیادی تصور ٹرم دی \emph{!!{value}}
تے !!a{formula} دا \emph{پورا اترنا} نیں۔ عام لفظاں وچ ٹرم دی
!!{value} کسے !!a{structure} دا اک !!{element} اے---جے ٹرم صرف
ثابتہ ہووے تاں اوہدی !!{value} اوہ شے اے جیہڑی !!{structure}
اُس ثابتہ نوں دیندی اے؛ جے ٹرم !!{function}s نال بنیا ہووے تاں
اوہدی !!{value} ثابتیاں دیاں !!{value}s، متغیراں نوں قدر بندی
وچ دتیاں قدراں، تے ٹرم وچلے فنکشناں نوں دتے فنکشناں توں
حساب کیتی جاندی اے۔
% PNB-SEM-017: کھلے فنکشن-ٹرماں لئی متغیراں دیاں قدراں وی چاہیدیاں نیں؛ اگلی مکمل تعریف مطابق ایہ شرط صریح کیتی اے۔
!!^a{formula} کسے
!!a{structure} وچ \emph{پورا اتردا} اے جے محمولاں نوں دتی
تعبیر !!{formula} نوں اُس !!{structure} دے دائرے وچ سچ بناوے۔
پورا اترن دا ایہ تصور استقرائی طور تے متعین ہوندا اے: !!{structure} دی
تعیین سیدھا دسدی اے کہ ایٹمی !!{formula}s کدوں پورے اتردے نیں؛
فیر پیچیدہ !!{formula} دا پورا اترنا اوہدے مرکزی رابطے یا مسوّر
تے فوری !!{subformula}s دے پورا اترن یا نہ اترن مطابق تعریف کردے آں۔

% PNB-SEM-016: مسوّر والے فارمولے دا اندرلا فارمولا آزاد متغیر رکھ سکدا اے، پر ضروری نہیں؛ بے اثر مسوّر وی جائز نیں۔
مسوّراں دی صورت ذرا نازک اے: مسوّر والے !!{formula} دے فوری
!!{subformula} وچ !!{variable} آزاد ہو سکدا اے، پر !!{structure}s
!!{variable}s دیاں !!{value}s متعین نہیں کردیاں۔ ایس مشکل لئی
اسیں \emph{متغیر قدر بندیاں} وی لیاندے آں تے پورا اترنا اکیلی
!!a{structure} دے لحاظ نال نہیں، بلکہ !!a{structure} تے
!!a{variable} قدر بندی دوناں دے لحاظ نال تعریف کردے آں۔
\end{explain}

\begin{defn}[متغیر قدر بندی]
\emph{متغیر قدر بندی}~$s$، !!a{structure}~$\Struct{M}$ لئی،
اک فنکشن اے جیہڑا ہر !!{variable} نوں~$\Domain M$ دے
!!a{element} ول بھیجدا اے؛ یعنی $s\colon \Var \to \Domain M$۔
\end{defn}

\begin{explain}
!!^a{structure} ہر !!{constant} نوں !!a{value} دیندی اے،
تے متغیر قدر بندی ہر متغیر نوں قدر دیندی اے۔ پر اسیں انہاں توں
بنے ٹرماں نال وی !!{domain} دے !!{element}s دا ناں لینا چاہندے
آں۔ ایس لئی اسیں ٹرماں دی !!{value} استقرائی طور تے تعریف
کردے آں۔ !!{constant}s تے متغیراں لئی قدر بالکل اوہ اے جیہڑی
!!{structure} یا متغیر قدر بندی متعین کردی اے؛ ودھ پیچیدہ
ٹرماں لئی اوہ قدر !!{structure} ولوں !!{function}s نوں دتے
فنکشناں نال بازگشتی طور تے حساب ہوندی اے۔
\end{explain}

\begin{defn}[ٹرماں دی !!^{value}]
جے $t$ زبان~$\Lang L$ دا ٹرم اے، $\Struct M$ اُس~$\Lang L$
لئی !!{structure} اے، تے $s$،~$\Struct M$ لئی !!a{variable}
قدر بندی اے، تاں \emph{!!{value}}~$\Value{t}{M}[s]$ ایویں
تعریف ہوندی اے:
\begin{enumerate}
\item \indcase{t}{c}{$\Value{\indfrm}{M}[s] = \Assign{\indcomplex}{M}$۔}
\item \indcase{t}{x}{$\Value{\indfrm}{M}[s] = s(\indcomplex)$۔}
\item \indcase{t}{\Atom{f}{t_1, \ldots, t_n}}{
\[
\Value{\indfrm}{M}[s] = \Assign{f}{M}(\Value{t_1}{M}[s], \ldots,
\Value{t_n}{M}[s]).
\]}
\end{enumerate}
\end{defn}

\begin{defn}[$x$-بدل قدر بندی]
جے $s$، !!a{structure}~$\Struct M$ لئی !!a{variable} قدر بندی
اے، تاں ہر ایہو جیہی !!{variable} قدر بندی~$s'$، جیہڑی~$\Struct M$
لئی ہووے تے~$s$ توں ودھ توں ودھ~$x$ نوں دتی قدر وچ وکھری
ہووے، اوہنوں~$s$ دی \emph{$x$-بدل قدر بندی} آکھدے آں۔
جے $s'$، $x$ دے لحاظ نال~$s$ دی بدل قدر بندی اے، تاں اسیں
$\varAssign{s'}{s}{x}$ لکھدے آں۔
\end{defn}

\begin{explain}
دھیان دیو: $x$-بدل قدر بندی لئی اصل قدر بندی~$s$ دے مقابل
~$x$ نوں وکھری قدر دینا \emph{ضروری نہیں}۔ اصل وچ ہر قدر بندی
اپنے آپ دی وی اک $x$-بدل قدر بندی اے۔
\end{explain}

\begin{defn}
  جے $s$، !!a{structure}~$\Struct M$ لئی !!a{variable} قدر بندی
  اے تے $m \in \Domain{M}$، تاں~$\Subst{s}{m}{x}$ اوہ متغیر
  قدر بندی اے جیہڑی ایویں تعریف ہوندی اے:
  \[\Subst{s}{m}{x}(y) = \begin{cases}
    m & \text{جے } y \ident x\\
    s(y) & \text{ہور صورت وچ}.
  \end{cases}\]
\end{defn}

دوجے لفظاں وچ، $\Subst{s}{m}{x}$ اوہ خاص $x$-بدل
قدر بندی اے جیہڑی~$s$ توں بندی اے: ایہ دائرے دا !!{element}~$m$،~$x$ نوں دیندی اے؛
تے $x$ توں ہور !!{variable}s نوں اوہی قدراں دیندی اے جیہڑیاں
$s$ دیندی اے۔

\begin{defn}[پورا اترنا]
\ollabel{defn:satisfaction}
!!a{formula}~$!A$ دا !!a{structure}~$\Struct M$ وچ
!!a{variable} قدر بندی~$s$ دے لحاظ نال پورا اترنا، علامتاں وچ
$\Sat{M}{!A}[s]$، ہیٹھاں بازگشتی طور تے تعریف ہوندا اے۔
(اسیں $\Sat/{M}{!A}[s]$ نوں ``$\Sat{M}{!A}[s]$ نہیں'' دے
معنی وچ لکھدے آں۔)
\begin{enumerate}
\tagitem{prvFalse}{%
  \indcase{!A}{\lfalse}{$\Sat/{M}{\indfrm}[s]$۔}}{}

\tagitem{prvTrue}{%
  \indcase{!A}{\ltrue}{$\Sat{M}{\indfrm}[s]$۔}}{}

\item \indcase{!A}{\Atom{R}{t_1, \dots, t_n}}{$\Sat{M}{\indfrm}[s]$
  اُدوں تے صرف اُدوں، جدوں $\langle \Value{t_1}{M}[s], \dots, \Value{t_n}{M}[s] \rangle \in
  \Assign{R}{M}$۔}

\item \indcase{!A}{\eq[t_1][t_2]}{$\Sat{M}{\indfrm}[s]$ اُدوں تے صرف اُدوں، جدوں
  $\Value{t_1}{M}[s] = \Value{t_2}{M}[s]$۔}

\tagitem{prvNot}{%
  \indcase{!A}{\lnot !B}{$\Sat{M}{\indfrm}[s]$ اُدوں تے صرف اُدوں، جدوں
    $\Sat/{M}{!B}[s]$۔}}{}

\tagitem{prvAnd}{%
  \indcase{!A}{(!B \land !C)}{$\Sat{M}{\indfrm}[s]$ اُدوں تے صرف اُدوں، جدوں $\Sat{M}{!B}[s]$
    تے $\Sat{M}{!C}[s]$ دونیں۔}}{}

\tagitem{prvOr}{%
  \indcase{!A}{(!B \lor !C)}{$\Sat{M}{\indfrm}[s]$ اُدوں تے صرف اُدوں، جدوں
    $\Sat{M}{!B}[s]$ یا $\Sat{M}{!C}[s]$، یا دونیں۔}}{}

\tagitem{prvIf}{%
  \indcase{!A}{(!B \lif !C)}{$\Sat{M}{\indfrm}[s]$ اُدوں تے صرف اُدوں، جدوں $\Sat/{M}{!B}[s]$
    یا $\Sat{M}{!C}[s]$، یا دونیں۔}}{}

\tagitem{prvIff}{%
  \indcase{!A}{(!B \liff !C)}{$\Sat{M}{\indfrm}[s]$ اُدوں تے صرف اُدوں، جدوں یا تاں
    $\Sat{M}{!B}[s]$ تے $\Sat{M}{!C}[s]$ دونیں ہون، یا نہ $\Sat{M}{!B}[s]$
    ہووے نہ $\Sat{M}{!C}[s]$۔}}{}

\tagitem{prvAll}{%
  \indcase{!A}{\lforall[x][!B]}{$\Sat{M}{\indfrm}[s]$ اُدوں تے صرف اُدوں، جدوں ہر
    !!{element}~$m \in \Domain M$ لئی $\Sat{M}{!B}[\Subst{s}{m}{x}]$۔}}{}

\tagitem{prvEx}{%
  \indcase{!A}{\lexists[x][!B]}{$\Sat{M}{\indfrm}[s]$ اُدوں تے صرف اُدوں، جدوں گھٹ توں گھٹ
  اک !!{element}~$m \in \Domain M$ لئی $\Sat{M}{!B}[\Subst{s}{m}{x}]$۔}}{}
\end{enumerate}
\end{defn}

\begin{explain}
متغیر قدر بندیاں آخری
\iftag{notprvEx,notprvAll}{شق}{دو شقّاں} وچ اہم نیں۔\iftag{prvAll}{ اسیں
$\lforall[x][!B(x)]$ دا پورا اترنا ایہ آکھ کے تعریف نہیں کر سکدے:
``ہر $m \in \Domain{M}$ لئی $\Sat{M}{!B(m)}$۔''}{}\iftag{prvEx}{ اسیں
$\lexists[x][!B(x)]$ دا پورا اترنا ایہ آکھ کے تعریف نہیں کر سکدے:
``گھٹ توں گھٹ اک $m \in \Domain{M}$ لئی $\Sat{M}{!B(m)}$۔''}{} وجہ
ایہ اے کہ جے $m \in \Domain M$، تاں اوہ زبان دی علامت نہیں؛
ایس لئی $!B(m)$~کوئی !!a{formula} نہیں (یعنی
$\Subst{!B}{m}{x}$ غیر تعریف شدہ اے)۔ اسیں ایہ وی فرض نہیں
کر سکدے کہ ساڈے کول ہر !!{element}~$\Struct{M}$ دا ناں رکھن
لئی !!{constant}s یا ٹرم موجود نیں، کیوں جے !!{structure}s دی
تعریف وچ ایسی کوئی شرط نہیں۔ معیاری زبان وچ !!{constant}s دا
سیٹ !!{denumerable} اے؛ ایس لئی جے $\Domain{M}$
!!{enumerable} نہ ہووے تاں ہر شے دا ناں رکھن لئی کافی
!!{constant}s ہو ہی نہیں سکدے۔

ایس مسئلے نوں اسیں !!{variable} قدر بندیاں لیان نال حل کردے
آں؛ ایہ متغیراں نوں دائرے دے !!{element}s نال سیدھا جوڑدیاں
نیں۔ فیر مثال لئی ایہ آکھن دی تھاں کہ
\iftag{prvEx}{$\lexists[x][!B(x)]$}{$\lforall[x][!B(x)]$}،
~$\Struct M$ وچ اُدوں تے صرف اُدوں پورا اتردا اے جدوں
% PNB-SEM-009: ماخذ دی prvEx شاخ وچ گھٹ توں گھٹ اک عنصر آکھن توں بعد اندرلے بیان دی شرط رہ گئی اے۔
\iftag{prvEx}{گھٹ توں گھٹ اک $m \in
\Domain{M}$ لئی $\Sat{M}{!B(m)}$}{ہر $m \in \Domain{M}$ لئی $\Sat{M}{!B(m)}$}،
اسیں ایہ آکھدے آں کہ اوہ~$\Struct M$ وچ~$s$ دے
\emph{لحاظ نال} اُدوں تے صرف اُدوں پورا اتردا اے جدوں
$!B(x)$،~$\Subst{s}{m}{x}$ دے لحاظ نال
\iftag{prvEx}{گھٹ توں گھٹ اک}{ہر} $m \in \Domain M$ لئی پورا اتردا اے۔
\end{explain}

\begin{ex}
فرض کرو $\Lang{L} = \{a, b, f, R\}$، جتھے $a$ تے $b$
!!{constant}s نیں، $f$ اک دو-مقامی !!{function} اے، تے $R$
اک دو-مقامی !!{predicate} اے۔ ہیٹھ دتی !!{structure}~$\Struct{M}$ ویکھو:
\begin{enumerate}
\item $\Domain M = \{1, 2, 3, 4\}$
\item $\Assign{a}{M} = 1$
\item $\Assign{b}{M} = 2$
\item $\Assign{f}{M}(x, y) = x+y$ جے $x+y \le 3$ ہووے، تے
  ہور صورت وچ $= 3$۔
\item $\Assign{R}{M} = \{\tuple{1, 1}, \tuple{1, 2}, \tuple{2, 3}, \tuple{2, 4}\}$
\end{enumerate}
فنکشن $s(x) = 1$ ہر !!{variable} نوں $1 \in \Domain{M}$
دیندا اے، ایس لئی ایہ~$\Struct{M}$ لئی متغیر قدر بندی اے۔

فیر
\begin{align*}
\Value{f(a,b)}{M}[s] & = \Assign{f}{M}(\Value{a}{M}[s], \Value{b}{M}[s]).
\intertext{کیوں جے $a$ تے $b$ !!{constant}s نیں، $\Value{a}{M}[s]
  = \Assign{a}{M} = 1$ تے $\Value{b}{M}[s] = \Assign{b}{M} = 2$۔ ایس لئی}
\Value{f(a,b)}{M}[s] & = \Assign{f}{M}(1, 2) = 1+2 = 3.
\intertext{$f(f(a,b),a)$ دی قدر حساب کرن لئی اسیں ایہ ویکھدے آں:}
\Value{f(f(a,b),a)}{M}[s] & = \Assign{f}{M}(\Value{f(a, b)}{M}[s],
\Value{a}{M}[s]) = \Assign{f}{M}(3, 1) = 3,
\intertext{کیوں جے $3+1 > 3$۔ تے کیوں جے $s(x) = 1$ تے $\Value{x}{M}[s] =
  s(x)$، ساڈے کول ایہ وی اے:}
\Value{f(f(a,b),x)}{M}[s] & = \Assign{f}{M}(\Value{f(a, b)}{M}[s],
\Value{x}{M}[s]) = \Assign{f}{M}(3, 1) = 3,
\end{align*}

ایٹمی !!{formula}~$R(t_1, t_2)$ اُدوں پورا اتردا اے جدوں
اوہدیاں دلیلاں دیاں قدراں دا ٹپل، یعنی $\tuple{\Value{t_1}{M}[s],
  \Value{t_2}{M}[s]}$،~$\Assign{R}{M}$ دا !!a{element} ہووے۔
ایس لئی، مثال لئی، $\Sat{M}{R(b,f(a,b))}[s]$، کیوں جے
$\tuple{\Value{b}{M},
  \Value{f(a,b)}{M}} = \tuple{2, 3} \in \Assign{R}{M}$؛ پر
% PNB-SEM-003: محمول دی تعبیر اک تعلق اے؛ ایہ متغیر قدر بندی مطابق بدلدا تعلق نہیں۔
$\Sat/{M}{R(x, f(a,b))}[s]$، کیوں جے $\tuple{1, 3} \notin \Assign{R}{M}$۔

جے اک غیر ایٹمی فارمولا~$!A$ پورا اتردا اے یا نہیں جاننا ہووے،
تاں استقرائی تعریف دی اوہ شق لاگو کرو جیہڑی اوہدے مرکزی رابطے
نال متعلق اے۔ مثال لئی، $R(a, a) \lif (R(b,
x) \lor R(x, b))$ دا مرکزی رابطہ~$\lif$ اے، تے
\begin{align*}
  & \Sat{M}{R(a, a) \lif (R(b, x) \lor R(x, b))}[s] \text{ اُدوں تے صرف اُدوں جدوں }\\
  & \qquad
  \Sat/{M}{R(a,a)}[s] \text{ یا } \Sat{M}{R(b, x) \lor R(x, b)}[s]
  \intertext{کیوں جے $\Sat{M}{R(a,a)}[s]$ (کیوں جے $\tuple{1,1} \in
    \Assign{R}{M}$)، ہُن تک جواب پکا نہیں؛ پہلے ایہ جاننا پئے گا کہ
    $\Sat{M}{R(b, x) \lor R(x, b)}[s]$ ہوندا اے یا نہیں:}
  & \Sat{M}{R(b, x) \lor R(x, b)}[s] \text{ اُدوں تے صرف اُدوں جدوں }\\
  & \qquad \Sat{M}{R(b,x)}[s] \text{ یا } \Sat{M}{R(x, b)}[s]
  \intertext{ایہ صورت واقعی اے، کیوں جے $\Sat{M}{R(x, b)}[s]$
    (کیوں جے $\tuple{1,2} \in \Assign{R}{M}$)۔}
\end{align*}

یاد رکھو کہ $x$-بدل قدر بندی~$s$ دی اوہ متغیر قدر بندی اے
جیہڑی~$s$ توں ودھ توں ودھ~$x$ نوں دتی قدر وچ وکھری اے۔
~$\Domain{M}$ دے ہر !!{element} لئی اک $x$-بدل قدر بندی~$s$ دی اے:
\begin{align*}
  s_1 & = \Subst{s}{1}{x}, &
  s_2 & = \Subst{s}{2}{x},\\
  s_3 & = \Subst{s}{3}{x}, &
  s_4 & = \Subst{s}{4}{x}.
\end{align*}
مثال لئی، $s_2(x) = 2$ تے $s_2(y) = s(y) = 1$ ہر اوہ
متغیر~$y$ لئی جیہڑا~$x$ نہیں۔ ایہ ساریاں $x$-بدل قدر بندیاں~$s$
دیاں نیں، ساخت~$\Struct{M}$ لئی، کیوں جے
$\Domain{M} = \{1, 2, 3, 4\}$۔ خاص طور تے $s_1 = s$
(~$s$ ہمیشہ اپنے آپ دی اک $x$-بدل قدر بندی اے)۔

\iftag{prvEx}{جے وجودی مسوّر والا
  !!{formula}~$\lexists[x][!A(x)]$ پورا اتردا اے یا نہیں
  جاننا ہووے، تاں ایہ ویکھنا پئے گا کہ
  $\Sat{M}{!A(x)}[\Subst{s}{m}{x}]$ گھٹ توں گھٹ اک
  $m \in \Domain M$ لئی ہوندا اے یا نہیں۔ ایس لئی
  \[
  \Sat{M}{\lexists[x][(R(b,x) \lor R(x,b))]}[s],
  \]
  کیوں جے $\Sat{M}{R(b,x) \lor R(x, b)}[\Subst{s}{1}{x}]$
  ($\Subst{s}{3}{x}$ وی گواہ بن سکدا سی)۔ پر
  \[
  \Sat/{M}{\lexists[x][(R(b,x) \land R(x,b))]}[s]
  \]
  کیوں جے اسیں جیہڑا وی $m \in \Domain{M}$ چنیے،
  $\Sat/{M}{R(b,x) \land R(x,b)}[\Subst{s}{m}{x}]$۔}{}

\iftag{prvAll}{جے کلی مسوّر والا
  !!{formula}~$\lforall[x][!A(x)]$ پورا اتردا اے یا نہیں
  جاننا ہووے، تاں ایہ ویکھنا پئے گا کہ
  $\Sat{M}{!A(x)}[\Subst{s}{m}{x}]$ ہر $m \in \Domain M$
  لئی ہوندا اے یا نہیں۔ ایس لئی
  \[
  \Sat{M}{\lforall[x][(R(x,a) \lif R(a,x))]}[s],
  \]
  کیوں جے $\Sat{M}{R(x,a) \lif R(a,x)}[\Subst{s}{m}{x}]$ ہر
  $m \in \Domain M$ لئی۔
  $m = 1$ لئی $\Sat{M}{R(a,x)}[\Subst{s}{1}{x}]$ اے،
  ایس لئی نتیجہ سچ اے؛ $m = 2$، $3$ تے~$4$ لئی
  $\Sat/{M}{R(x,a)}[\Subst{s}{m}{x}]$ اے، ایس لئی مقدم جھوٹا اے۔
  پر
  \[
  \Sat/{M}{\lforall[x][(R(a,x) \lif R(x,a))]}[s]
  \]
  کیوں جے $\Sat/{M}{R(a,x) \lif R(x,a)}[\Subst{s}{2}{x}]$
  (کیوں جے $\Sat{M}{R(a, x)}[\Subst{s}{2}{x}]$ تے
  $\Sat/{M}{R(x,
  a)}[\Subst{s}{2}{x}]$)۔}{}

\iftag{defEx}{جے وجودی مسوّر والا
  !!{formula}~$\lexists[x][!A(x)]$ پورا اتردا اے یا نہیں
  جاننا ہووے، تاں ایہ ویکھنا پئے گا کہ
  $\Sat{M}{\lnot \lforall[x][\lnot !A(x)]}[s]$ ہوندا اے یا نہیں۔
  مثال لئی
  \[
  \Sat{M}{\lexists[x][(R(b,x) \lor R(x,b))]}[s].
  \]
  پہلے $\Sat{M}{R(b,x) \lor R(x, b)}[\Subst{s}{1}{x}]$
  ($\Subst{s}{3}{x}$ وی گواہ بن سکدا سی)۔ ایس لئی
  $\Sat/{M}{\lnot(R(b,x) \lor R(x,b))}[\Subst{s}{1}{x}]$؛ لہٰذا
  % PNB-SEM-010: دو اخذاں وچ نفی توں بعد اک وادھو کھلدی قوس اے؛ ایتھے قوساں برابر نیں۔
  $\Sat/{M}{\lforall[x][\lnot(R(b,x) \lor R(x,b))]}[s]$، تے
  نتیجتاً $\Sat{M}{\lnot\lforall[x][\lnot(R(b,x) \lor
  R(x,b))]}[s]$۔ دوجے پاسے
  % PNB-SEM-004: منجمد ماخذ وچ منفی وجودی فارمولے اندر اک وادھو کاما اے؛ ایتھے ہٹایا اے۔
  \[
  \Sat/{M}{\lexists[x][(R(b,x) \land R(x,b))]}[s].
  \]
  ایہ اس لئی کہ $\Sat{M}{\lforall[x][\lnot(R(b,x) \land
  R(x,b))]}[s]$؛ کیوں جے کسے وی $m \in \Domain M$ لئی
  $\Sat{M}{R(b,x) \land R(x,b)}[\Subst{s}{m}{x}]$ نہیں۔ انہاں
  مثالاں توں واضح اے کہ $\Sat{M}{\lexists[x][!A(x)]}[s]$ اُدوں تے
  صرف اُدوں جدوں $\Sat{M}{!A(x)}[\Subst{s}{m}{x}]$ گھٹ توں گھٹ
  اک $m \in \Domain M$ لئی۔}{}

\iftag{defAll}{جے کلی مسوّر والا
  !!{formula}~$\lforall[x][!A(x)]$ پورا اتردا اے یا نہیں
  جاننا ہووے، تاں ایہ ویکھنا پئے گا کہ
  $\Sat{M}{\lnot\lexists[x][\lnot !A(x)]}[s]$ ہوندا اے یا نہیں۔
  مثال لئی
  \[
  \Sat{M}{\lforall[x][(R(x,a) \lif R(a,x))]}[s],
  \]
  پہلے $\Sat{M}{R(x,a) \lif R(a,x)}[\Subst{s}{m}{x}]$ ہر
  $m \in \Domain M$ لئی
  ($\Sat{M}{R(a,x)}[\Subst{s}{1}{x}]$ تے
  % PNB-SEM-005: دوجے عنصر لئی جھوٹا مقدم R(x,a) اے، R(a,x) نہیں۔
  $\Sat/{M}{R(x,a)}[\Subst{s}{m}{x}]$ جدوں $m = 2$، $3$ یا~$4$)۔
  ایس لئی کسے $m \in \Domain M$ لئی
  $\Sat{M}{\lnot(R(x,a) \lif R(a,x))}[\Subst{s}{m}{x}]$ نہیں، لہٰذا
  $\Sat/{M}{\lexists[x][\lnot(R(x,a) \lif R(a,x))]}[s]$۔ نتیجتاً
  $\Sat{M}{\lnot\lexists[x][\lnot(R(x,a) \lif R(a,x))]}[s]$۔
  دوجے پاسے
  \[
  \Sat/{M}{\lforall[x][(R(a,x) \lif R(x,a))]}[s],
  \]
  کیوں جے $\Sat/{M}{R(a,x) \lif R(x,a)}[\Subst{s}{2}{x}]$،
  تے ایس لئی $\Sat{M}{\lexists[x][\lnot(R(a,x) \lif R(x,a))]}[s]$۔
  انہاں مثالاں توں واضح اے کہ $\Sat{M}{\lforall[x][!A(x)]}[s]$
  اُدوں تے صرف اُدوں جدوں $\Sat{M}{!A(x)}[\Subst{s}{m}{x}]$
  ہر $m \in \Domain M$ لئی۔}{}

ذرا ودھ پیچیدہ صورت ویکھو:
\[
\lforall[x][(R(a,x) \lif \lexists[y][R(x,y)])].
\]
کیوں جے $\Sat/{M}{R(a,x)}[\Subst{s}{3}{x}]$ تے
$\Sat/{M}{R(a,x)}[\Subst{s}{4}{x}]$، اس لئی شرطیہ دے نتیجے
بارے سوچن والیاں صورتاں صرف $m = 1$ تے
% PNB-SEM-006: ماخذ وچ دوجی قدر توں پہلاں m رہ گیا اے؛ ایتھے اشاریہ صریح اے۔
$m = 2$ نیں۔ کی $\Sat{M}{\lexists[y][R(x,y)]}[\Subst{s}{1}{x}]$
درست اے؟ ایہ اُدوں درست اے جے گھٹ توں گھٹ اک $n \in \Domain M$
ہووے جیہدے لئی
$\Sat{M}{R(x,y)}[\Subst{\Subst{s}{1}{x}}{n}{y}]$۔ واقعی جے
$n = 1$ چنیے تاں $\Subst{\Subst{s}{1}{x}}{n}{y} = \Subst{s}{1}{y} =
s$۔ کیوں جے $s(x) = 1$، $s(y) = 1$ تے
$\tuple{1,1} \in \Assign{R}{M}$، جواب ہاں اے۔

ایہ جانن لئی کہ $\Sat{M}{\lexists[y][R(x,y)]}[\Subst{s}{2}{x}]$
درست اے یا نہیں، اسیں !!{variable} قدر بندیاں
$\Subst{\Subst{s}{2}{x}}{n}{y}$ ویکھدے آں۔ ایتھے $n = 1$ لئی
ایہ قدر بندی~$s_2 = \Subst{s}{2}{x}$ اے، جیہڑی $R(x,y)$ نوں
پورا نہیں کردی ($s_2(x) =
2$، $s_2(y) = 1$ تے $\tuple{2,1}\notin \Assign{R}{M}$)۔ پر
$\Subst{\Subst{s}{2}{x}}{3}{y} = \Subst{s_2}{3}{y}$ ویکھو۔
$\Sat{M}{R(x,y)}[\Subst{s_2}{3}{y}]$، کیوں جے
$\tuple{2,3} \in
\Assign{R}{M}$؛ ایس لئی $\Sat{M}{\lexists[y][R(x,y)]}[s_2]$۔

% PNB-SEM-007: ایتھے باہرلے کلی بیان دا متغیر m اے، n نہیں۔
سو ہر $m \in \Domain M$ لئی یا تاں
$\Sat/{M}{R(a,x)}[\Subst{s}{m}{x}]$ (جے $m = 3$، $4$) یا
$\Sat{M}{\lexists[y][R(x,y)]}[\Subst{s}{m}{x}]$ (جے $m = 1$، $2$)؛
ایس لئی
\[
\Sat{M}{\lforall[x][(R(a,x) \lif \lexists[y][R(x,y)])]}[s].
\]
دوجے پاسے
\[
\Sat/{M}{\lexists[x][(R(a,x) \land \lforall[y][R(x,y)])]}[s].
\]
$\Sat{M}{R(a,x)}[\Subst{s}{m}{x}]$ صرف $m = 1$ تے $m =
2$ لئی اے۔ پر~$m$ دیاں دوناں ایہناں قدراں لئی اک
$n \in
\Domain M$ ایسا اے جیہدے نال کلی شق ناکام ہوندی اے:
% PNB-SEM-008: پہلی قدر لئی چوتھا عنصر ناکامی دا گواہ اے؛ دوجی قدر لئی دائرے دا پہلا عنصر ورتیا اے۔
پہلی قدر لئی $n = 4$ چنو، تے دوجی لئی دائرے دا پہلا عنصر، یعنی ۱۔
ایس طرح $\Sat/{M}{R(x,y)}[\Subst{\Subst{s}{m}{x}}{n}{y}]$
تے نتیجتاً $\Sat/{M}{\lforall[y][R(x,y)]}[\Subst{s}{m}{x}]$
دونیں $m = 1$ تے $m =
2$ لئی حاصل ہوندے نیں۔ کُل ملا کے کوئی
$m \in \Domain M$ ایسا نہیں جیہدے لئی $\Sat{M}{R(a,x)
\land \lforall[y][R(x,y)]}[\Subst{s}{m}{x}]$۔
\end{ex}

\iftag{defEx}{%
\begin{prop}\ollabel{prop:sat-ex}
$\Sat{M}{\lexists[x][!B(x)]}[s]$ اُدوں تے صرف اُدوں جدوں
کوئی $x$-بدل قدر بندی $s'$ موجود ہووے جیہڑی~$s$ دی ہووے تے
    $\Sat{M}{!B(x)}[s']$۔
\end{prop}

\begin{proof}
  مشق۔
\end{proof}
}{}

\tagprob[defEx]{notdefEx}
\begin{prob}
  \olref[fol][syn][sat]{prop:sat-ex} ثابت کرو۔
\end{prob}
\tagendprob

\iftag{defAll}{%
\begin{prop}\ollabel{prop:sat-all}
$\Sat{M}{\lforall[x][!B(x)]}[s]$ اُدوں تے صرف اُدوں جدوں
ہر $x$-بدل قدر بندی~$s'$ لئی، جیہڑی~$s$ دی ہووے،
  $\Sat{M}{!B(x)}[s']$۔
\end{prop}

\begin{proof}
  مشق۔
\end{proof}
}{}

\tagprob[defAll]{notdefAll}
\begin{prob}
  \olref[fol][syn][sat]{prop:sat-all} ثابت کرو۔
\end{prob}
\tagendprob

\begin{prob}
فرض کرو $\Lang L = \{c, f, A\}$، جیہدے وچ اک !!{constant}،
اک یک-مقامی !!{function} تے اک دو-مقامی !!{predicate} اے، تے
!!{structure}~$\Struct{M}$ ایویں دتی اے:
\begin{enumerate}
\item $\Domain M = \{1, 2, 3\}$
\item $\Assign{c}{M} = 3$
\item $\Assign{f}{M}(1) = 2, \Assign{f}{M}(2) = 3, \Assign{f}{M}(3) = 2$
\item $\Assign{A}{M} = \{\tuple{1, 2}, \tuple{2, 3}, \tuple{3, 3}\}$
\end{enumerate}
(الف) $s(v) = 1$ رکھو سارے !!{variable}s~$v$ لئی۔ معلوم کرو کہ
\[
\Sat{M}{\lexists[x][(A(f(z), c) \lif \lforall[y][(A(y, x) \lor A(f(y),
      x))])]}[s]
\]
درست اے یا نہیں۔ وجہ سمجھاؤ۔

(ب) اک وکھری ساخت تے !!{variable} قدر بندی دیو جیہدے وچ
ایہ !!{formula} پورا نہ اترے۔
\end{prob}
\endgroup
% /OLP-0163

% OLP-0164 content/first-order-logic/syntax-and-semantics/assignments.tex
\begingroup
\def\olfilename{assignments}
\expandafter\def\csname ol@sectioncs\endcsname{section}

\olfileid{fol}{syn}{ass}

\olsection{متغیر قدر بندیاں}

\begin{explain}
!!{variable} قدر بندی~$s$ \emph{ہر} متغیر نوں قدر دیندی اے---تے
متغیر لامتناہی نیں۔ ظاہر اے کہ ایہ لازم نہیں۔ اسیں !!{variable}
قدر بندیاں توں سارے !!{variable}s نوں قدراں دیون دا تقاضا صرف
ایس لئی کردے آں کہ گل بہت سوکھی ہو جاندی اے۔ ٹرم~$t$ دی قدر
تے !!a{structure} وچ !!a{formula}~$!A$ دا~$s$ دے لحاظ نال
پورا اترنا یا نہ اترنا صرف انہاں قدراں اُتے منحصر اے جیہڑیاں
قدر بندی~$s$،~$t$ دے !!{variable}s تے~$!A$ دے آزاد
!!{variable}s نوں دیندی اے۔ اگلے دو قضیے ایہی گل دسدے نیں۔
``منحصر ہون'' نوں پکا کرن لئی اسیں ثابت کردے آں کہ دو متغیر
قدر بندیاں جے~$t$ دے سارے متغیراں اُتے متفق ہون تاں اک ہی
قدر دیندیاں نیں؛ تے جے اوہ~$!A$ دے سارے آزاد متغیراں اُتے
متفق ہون تاں $!A$ پہلی دے لحاظ نال اُدوں تے صرف اُدوں
پورا اتردا اے جدوں دوجی دے لحاظ نال پورا اتردا اے۔
\end{explain}

\begin{prop}\ollabel{prop:valindep}
جے ٹرم~$t$ دے !!{variable}s، $x_1$، \dots،~$x_n$ وچوں ہون، تے
$s_1(x_i) = s_2(x_i)$، $i = 1$، \dots،~$n$ لئی، تاں
$\Value{t}{M}[s_1] = \Value{t}{M}[s_2]$۔
\end{prop}

\begin{proof}
~$t$ دی پیچیدگی اُتے استقرا کردے آں۔ بنیادی صورت وچ $t$ یا
!!a{constant} ہو سکدا اے یا~$x_1$، \dots،~$x_n$ وچوں کوئی
متغیر۔ جے $t = c$، تاں $\Value{t}{M}[s_1] = \Assign{c}{M} =
\Value{t}{M}[s_2]$۔ جے $t = x_i$، تاں قضیے دی مفروضہ شرط
مطابق $s_1(x_i) = s_2(x_i)$، ایس لئی
$\Value{t}{M}[s_1] = s_1(x_i) = s_2(x_i) = \Value{t}{M}[s_2]$۔

استقرائی قدم لئی فرض کرو $t = \Atom{f}{t_1, \dots, t_k}$، تے
دعویٰ $t_1$، \dots، $t_k$ لئی درست اے۔ فیر
\begin{align*}
  \Value{t}{M}[s_1] & = \Value{\Atom{f}{t_1, \dots, t_k}}{M}[s_1] \\
  & = \Assign{f}{M}(\Value{t_1}{M}[s_1], \dots, \Value{t_k}{M}[s_1]).
\intertext{$j = 1$، \dots،~$k$ لئی~$t_j$ دے !!{variable}s،
    $x_1$، \dots،~$x_n$ وچوں نیں۔ استقرائی مفروضے مطابق
    $\Value{t_j}{M}[s_1] = \Value{t_j}{M}[s_2]$۔ ایس لئی}
\Value{t}{M}[s_1] & = \Value{\Atom{f}{t_1, \dots, t_k}}{M}[s_1] \\
  & = \Assign{f}{M}(\Value{t_1}{M}[s_1], \dots, \Value{t_k}{M}[s_1]) \\
  & = \Assign{f}{M}(\Value{t_1}{M}[s_2], \dots, \Value{t_k}{M}[s_2]) \\
  & = \Value{\Atom{f}{t_1, \dots, t_k}}{M}[s_2] = \Value{t}{M}[s_2].
\end{align*}
\end{proof}

\begin{prop}\ollabel{prop:satindep}
جے $!A$ دے آزاد !!{variable}s، $x_1$، \dots،~$x_n$ وچوں ہون، تے
$s_1(x_i) = s_2(x_i)$، $i = 1$، \dots،~$n$ لئی، تاں
$\Sat{M}{!A}[s_1]$ اُدوں تے صرف اُدوں جدوں $\Sat{M}{!A}[s_2]$۔
\end{prop}

\begin{proof}
اسیں $!A$ دی پیچیدگی اُتے استقرا کردے آں۔ بنیادی صورت وچ، جدوں
$!A$ ایٹمی ہووے، $!A$ ایہناں وچوں ہو سکدا اے:
\iftag{prvTrue}{$\ltrue$،}{}
\iftag{prvFalse}{$\lfalse$،}{}
$\Atom{R}{t_1, \dots, t_k}$، جیہڑا $k$-مقامی محمول $R$ تے
ٹرماں $t_1$، \dots،~$t_k$ نال بندا اے؛ یا $\eq[t_1][t_2]$،
جیہڑا ٹرماں $t_1$ تے~$t_2$ نال بندا اے۔ پچھلی دو صورتاں وچ اسیں !!{biconditional}
دی صرف اگلی سمت دکھاندے آں، کیوں جے اُلٹی سمت دا ثبوت
متناظر اے۔

\begin{enumerate}
\tagitem{prvTrue}{%
  \indcase{!A}{\ltrue}{$\Sat{M}{\indfrm}[s_1]$ تے
    $\Sat{M}{\indfrm}[s_2]$ دونیں۔}}{}

\tagitem{prvFalse}{%
  \indcase{!A}{\lfalse}{$\Sat/{M}{!A}[s_1]$ تے
    $\Sat/{M}{!A}[s_2]$ دونیں۔}}{}

\item
  \indcase{!A}{\Atom{R}{t_1, \ldots, t_k}}{فرض کرو
    $\Sat{M}{\indfrm}[s_1]$۔ تاں
    \[
    \langle \Value{t_1}{M}[s_1], \ldots, \Value{t_k}{M}[s_1] \rangle
    \in \Assign{R}{M}.
    \]
    $i = 1$، \dots،~$k$ لئی \olref{prop:valindep} مطابق
    $\Value{t_i}{M}[s_1] = \Value{t_i}{M}[s_2]$۔ ایس لئی
    % PNB-SEM-011: ٹپل دی پہلی جگہ پہلا ٹرم ای رہندا اے، چلدا اشاریہ نہیں۔
    $\langle \Value{t_1}{M}[s_2], \ldots, \Value{t_k}{M}[s_2] \rangle
    \in \Assign{R}{M}$، تے نتیجتاً $\Sat{M}{\indfrm}[s_2]$۔}
\item
  \indcase{!A}{\eq[t_1][t_2]}{فرض کرو $\Sat{M}{\indfrm}[s_1]$۔
    تاں $\Value{t_1}{M}[s_1] = \Value{t_2}{M}[s_1]$۔ ایس لئی
    \begin{align*}
      \Value{t_1}{M}[s_2] & = \Value{t_1}{M}[s_1]
      & \text{(\olref{prop:valindep} مطابق)} \\
      & = \Value{t_2}{M}[s_1]
      & \text{(کیوں جے $\Sat{M}{\eq[t_1][t_2]}[s_1]$)}\\
      &= \Value{t_2}{M}[s_2]
      & \text{(\olref{prop:valindep} مطابق)،}
    \end{align*}
    ایس لئی $\Sat{M}{\eq[t_1][t_2]}[s_2]$۔}
\end{enumerate}

ہُن فرض کرو کہ $\Sat{M}{!B}[s_1]$ اُدوں تے صرف اُدوں جدوں
$\Sat{M}{!B}[s_2]$، سارے اوہناں !!{formula}s~$!B$ لئی
جیہڑے~$!A$ نالوں گھٹ پیچیدہ نیں۔ استقرائی قدم~$!A$ دے مرکزی عامل مطابق
صورتاں وچ ونڈیا جاندا اے۔ ہر صورت وچ اسیں !!{biconditional}
دی صرف اگلی سمت دکھاندے آں؛ اُلٹی سمت دا ثبوت متناظر اے۔
مسوّراں والیاں صورتاں توں چھٹ اسیں استقرائی مفروضہ ذیلی-!!{formula}s~$!B$
اُتے لاگو کردے آں، جیہڑے~$!A$ دے نیں۔~$!B$ دے آزاد متغیر،
~$!A$ دے آزاد متغیراں وچوں نیں۔ ایس لئی جے $s_1$ تے $s_2$،~$!A$
دے آزاد متغیراں اُتے متفق ہون، تاں اوہ~$!B$ دے آزاد متغیراں
اُتے وی متفق ہون گے، تے استقرائی مفروضہ~$!B$ اُتے لاگو ہووے گا۔

\begin{enumerate}
\tagitem{defNot}{}{%
  \iftag{probNot}{%
    \indcase!{!A}{\lnot !B}{}}{%
    \indcase{!A}{\lnot !B}{جے $\Sat{M}{\indfrm}[s_1]$، تاں
      $\Sat/{M}{!B}[s_1]$؛ استقرائی مفروضے مطابق
      $\Sat/{M}{!B}[s_2]$، لہٰذا $\Sat{M}{\indfrm}[s_2]$۔}}}

\tagitem{defAnd}{}{%
  \iftag{probAnd}{%
    \indcase!{!A}{!B \land !C}{}}{%
    \indcase{!A}{!B \land !C}{جے $\Sat{M}{\indfrm}[s_1]$، تاں
      $\Sat{M}{!B}[s_1]$ تے $\Sat{M}{!C}[s_1]$؛ استقرائی
      مفروضے مطابق $\Sat{M}{!B}[s_2]$ تے $\Sat{M}{!C}[s_2]$۔ ایس لئی
      $\Sat{M}{\indfrm}[s_2]$۔}}}

\tagitem{defOr}{}{%
  \iftag{probOr}{%
    \indcase!{!A}{!B \lor !C}{}}{%
    \indcase{!A}{!B \lor !C}{جے $\Sat{M}{\indfrm}[s_1]$، تاں
      $\Sat{M}{!B}[s_1]$ یا $\Sat{M}{!C}[s_1]$۔ استقرائی مفروضے
      مطابق $\Sat{M}{!B}[s_2]$ یا $\Sat{M}{!C}[s_2]$؛ لہٰذا
      $\Sat{M}{\indfrm}[s_2]$۔}}}

\tagitem{defIf}{}{%
  \iftag{probIf}{%
    \indcase!{!A}{!B \lif !C}{}}{%
    \indcase{!A}{!B \lif !C}{جے $\Sat{M}{\indfrm}[s_1]$، تاں
    $\Sat/{M}{!B}[s_1]$ یا $\Sat{M}{!C}[s_1]$۔ استقرائی مفروضے
    مطابق $\Sat/{M}{!B}[s_2]$ یا $\Sat{M}{!C}[s_2]$؛ ایس لئی
    $\Sat{M}{\indfrm}[s_2]$۔}}}

\tagitem{defIff}{}{%
  \iftag{probIff}{%
    \indcase!{!A}{!B \liff !C}{}}{%
    \indcase{!A}{!B \liff !C}{جے $\Sat{M}{\indfrm}[s_1]$، تاں یا تاں
      $\Sat{M}{!B}[s_1]$ تے $\Sat{M}{!C}[s_1]$ دونیں، یا $\Sat/{M}{!B}[s_1]$ تے
      $\Sat/{M}{!C}[s_1]$ دونیں۔ استقرائی مفروضے مطابق یا تاں
      $\Sat{M}{!B}[s_2]$ تے $\Sat{M}{!C}[s_2]$ دونیں، یا $\Sat/{M}{!B}[s_2]$
      تے $\Sat/{M}{!C}[s_2]$ دونیں۔ دونوں صورتاں وچ $\Sat{M}{\indfrm}[s_2]$۔}}}

\tagitem{defEx}{}{%
  \iftag{probEx}{%
    \indcase!{!A}{\lexists[x][!B]}{}}{%
    \indcase{!A}{\lexists[x][!B]}{جے $\Sat{M}{\indfrm}[s_1]$، تاں
      کوئی $m \in \Domain{M}$ ایسا اے جیہدے لئی
      $\Sat{M}{!B}[\Subst{s_1}{m}{x}]$۔ رکھو
      $s_1' = \Subst{s_1}{m}{x}$ تے $s_2' =
      \Subst{s_2}{m}{x}$۔~$!B$ دے آزاد متغیر $x_1$،
      \dots، $x_n$ تے $x$ وچوں نیں۔ $s_1'(x_i) = s_2'(x_i)$، کیوں جے
      $s_1'$ تے $s_2'$، $x$-بدل قدر بندیاں نیں، بالترتیب $s_1$ تے~$s_2$ دیاں،
      تے مفروضے مطابق $s_1(x_i) = s_2(x_i)$۔
      % PNB-SEM-018: درج متغیر مسوّر دے متغیر دے برابر وی ہو سکدا اے؛ دونوں صورتاں دی برابری صریح اے۔
      ایہ سبب مسوّر دے متغیر توں ہور درج متغیراں لئی اے؛ جے کوئی
      درج متغیر مسوّر دے متغیر دے برابر ہووے، تاں دوہاں وچ رکھی
      اکّو قدر نال برابری ملدی اے۔ انہاں دی تعریف توں
      $s_1'(x) = s_2'(x) = m$؛ $s_1'$ تے~$s_2'$ دی تعریف وچ ایہی
      قدر رکھی گئی سی۔ فیر استقرائی
      مفروضہ $!B$ تے $s_1'$، $s_2'$ اُتے لاگو ہوندا اے، ایس لئی
      $\Sat{M}{!B}[s_2']$۔ ہُن چونکہ $s_2' = \Subst{s_2}{m}{x}$،
      کوئی $m \in \Domain{M}$ ایسا اے جیہدے لئی
      $\Sat{M}{!B}[\Subst{s_2}{m}{x}]$، لہٰذا
      $\Sat{M}{\indfrm}[s_2]$۔}}}

\tagitem{defAll}{}{%
  \iftag{probAll}{%
    \indcase!{!A}{\lforall[x][!B]}{}}{%
    \indcase{!A}{\lforall[x][!B]}{جے $\Sat{M}{\indfrm}[s_1]$، تاں ہر
      $m \in \Domain{M}$ لئی $\Sat{M}{!B}[\Subst{s_1}{m}{x}]$۔ اسیں
      دکھانا چاہندے آں کہ ہر $m \in \Domain{M}$ لئی
      $\Sat{M}{!B}[\Subst{s_2}{m}{x}]$ وی اے۔ سو کوئی وی
      $m \in \Domain{M}$ چنو، تے
      % PNB-SEM-012: ماخذ وچ اک غیر متعین قدر بندی دو وار بدل دی اے؛ دونوں بدل صورتاں اپنی اپنی اصل قدر بندی توں بندیاں نیں۔
      $s_1' = \Subst{s_1}{m}{x}$ تے $s_2' =
      \Subst{s_2}{m}{x}$ ویکھو۔ ساڈے کول $\Sat{M}{!B}[s_1']$ اے۔
      ~ $!B$ دے آزاد متغیر $x_1$، \dots، $x_n$ تے $x$ وچوں نیں۔
      $s_1'(x_i) = s_2'(x_i)$، کیوں جے $s_1'$ تے $s_2'$
      $x$-بدل قدر بندیاں نیں، بالترتیب $s_1$ تے~$s_2$ دیاں، تے مفروضے مطابق
      $s_1(x_i) = s_2(x_i)$۔ ایہ سبب مسوّر دے متغیر توں ہور
      درج متغیراں لئی اے؛ جے کوئی درج متغیر مسوّر دے متغیر دے
      برابر ہووے، تاں دوہاں وچ رکھی اکّو قدر نال برابری ملدی اے۔
      تعریف توں $s_1'(x) = s_2'(x) = m$؛
      ایہ~$s_1'$ تے~$s_2'$ نوں دتی گئی قدر اے۔ فیر استقرائی مفروضہ~$!B$
      تے~$s_1'$، $s_2'$ اُتے لاگو ہوندا اے، تے
      $\Sat{M}{!B}[s_2']$ ملدا اے۔ ایہ ہر $m \in \Domain{M}$ لئی اے؛
      یعنی $\Sat{M}{!B}[\Subst{s_2}{m}{x}]$ سارے~$m \in
      \Domain{M}$ لئی اے۔ ایس لئی $\Sat{M}{\indfrm}[s_2]$۔}}}
\end{enumerate}
استقرا نال نتیجہ ملدا اے کہ $\Sat{M}{!A}[s_1]$ اُدوں تے صرف اُدوں
جدوں $\Sat{M}{!A}[s_2]$، جے $!A$ دے آزاد !!{variable}s، $x_1$،
\dots، $x_n$ وچوں ہون تے $s_1(x_i)=s_2(x_i)$،
$i = 1$، \dots،~$n$ لئی۔
\end{proof}

\begin{probtag}{probNot,probOr,probAnd,probIf,probIff,probEx,probAll}
\olref[fol][syn][ass]{prop:satindep} دا ثبوت مکمل کرو۔
\end{probtag}

\begin{explain}
!!^{sentence}s وچ آزاد متغیر نہیں ہوندے؛ ایس لئی کوئی وی دو متغیر
قدر بندیاں کسے وی جملے دے سارے آزاد متغیراں (جنہاں دی گنتی صفر
اے) نوں اوہی قدراں دیندیاں نیں۔ پچھلا ثابت کیتا قضیہ ایہ دسدا اے
کہ !!a{sentence} دا کسے ساخت وچ متغیر قدر بندی دے لحاظ نال
پورا اترنا یا نہ اترنا اُس قدر بندی توں بالکل آزاد اے۔ اسیں ایہ
حقیقت درج کردے آں۔ ایہی اگلی تعریف نوں جائز بناندی اے، جیہڑی
!!a{structure} وچ !!a{sentence} دے پورا اترن لئی متغیر قدر
بندی دا ذکر نہیں کردی۔
\end{explain}

\begin{cor}
\ollabel{cor:sat-sentence}
جے $!A$ اک !!a{sentence} اے تے $s$ متغیر قدر بندی اے، تاں
$\Sat{M}{!A}[s]$ اُدوں تے صرف اُدوں جدوں
$\Sat{M}{!A}[s']$، ہر متغیر قدر بندی~$s'$ لئی۔
\end{cor}

\begin{proof}
  $s'$ کوئی وی متغیر قدر بندی لو۔ چونکہ $!A$ اک !!a{sentence} اے،
  اوہدے کوئی آزاد متغیر نہیں؛ ایس لئی ہر متغیر قدر بندی~$s'$
  بدیہی طور تے~$!A$ دے سارے آزاد متغیراں نوں اوہی قدراں دیندی
  اے جیہڑیاں~$s$ دیندی اے۔ سو \olref{prop:satindep} دی شرط پوری
  اے، تے $\Sat{M}{!A}[s]$ اُدوں تے صرف اُدوں جدوں
  $\Sat{M}{!A}[s']$۔
\end{proof}

\begin{defn}
\ollabel{defn:satisfaction}
جے $!A$ اک !!a{sentence} اے، تاں اسیں آکھدے آں کہ
!!a{structure}~$\Struct M$،~$!A$ نوں \emph{پورا کردی اے}،
$\Sat{M}{!A}$، اُدوں تے صرف اُدوں جدوں
$\Sat{M}{!A}[s]$ ہر متغیر قدر بندی~$s$ لئی۔
\end{defn}

جے $\Sat{M}{!A}$، تاں اسیں سادہ لفظاں وچ ایہ وی آکھدے آں کہ
\emph{$!A$،~$\Struct{M}$ وچ سچ اے۔} پورا اترن دا تصور
انفرادی !!{sentence}s توں !!{sentence}s دے سیٹاں تک فطری طور تے
پھیل جاندا اے۔

\begin{defn}
\ollabel{defn:sat}
جے $\Gamma$، !!{sentence}s دا سیٹ~$\Gamma$ اے، تاں اسیں
آکھدے آں کہ !!a{structure}~$\Struct M$،~$\Gamma$ نوں
\emph{پورا کردی اے}، $\Sat{M}{\Gamma}$، اُدوں تے صرف اُدوں
جدوں $\Sat{M}{!A}$ ہر $!A \in \Gamma$ لئی۔
\end{defn}

\begin{prop}\ollabel{prop:sentence-sat-true}
  فرض کرو $\Struct{M}$ اک !!a{structure}، $!A$ اک
  !!a{sentence} تے $s$ متغیر قدر بندی اے۔ $\Sat{M}{!A}$
  اُدوں تے صرف اُدوں جدوں $\Sat{M}{!A}[s]$۔
\end{prop}

\begin{proof}
مشق۔
\end{proof}

\begin{prob}
\olref[fol][syn][ass]{prop:sentence-sat-true} ثابت کرو۔
\end{prob}

\begin{prop}\ollabel{prop:sat-quant}
  فرض کرو $!A(x)$ وچ صرف $x$ آزاد اے، تے $\Struct M$ اک
  !!a{structure} اے۔ فیر:
  \begin{tagenumerate}{prvEx,prvAll}
    \tagitem{prvEx}{$\Sat{M}{\lexists[x][!A(x)]}$ اُدوں تے صرف اُدوں جدوں
      $\Sat{M}{!A(x)}[s]$ گھٹ توں گھٹ اک متغیر قدر بندی~$s$ لئی۔}{}
    \tagitem{prvAll}{$\Sat{M}{\lforall[x][!A(x)]}$ اُدوں تے صرف اُدوں جدوں
      $\Sat{M}{!A(x)}[s]$ ہر متغیر قدر بندی~$s$ لئی۔}{}
\end{tagenumerate}
\end{prop}

\begin{proof}
مشق۔
\end{proof}

\begin{prob}
\olref[fol][syn][ass]{prop:sat-quant} ثابت کرو۔
\end{prob}

\begin{prob}
\DeclareRobustCommand{\VDash}{\mathrel{||}\joinrel\Relbar}
فرض کرو $\Lang L$ ایسی زبان اے جیہدے وچ !!{function}s نہیں۔
% PNB-SEM-021: تازہ ثابتہ نشان چنن لئی معیاری زبان دا لامتناہی ذخیرہ ایتھے صریح مفروضہ اے۔
ایہ مشق ثابتہ نشاناں دا لامتناہی ذخیرہ من کے کیتی جاندی اے،
تاں جو ہر قدم اُتے تازہ ناں چنیا جا سکے۔
!!{structure}~$\Struct M$، !!a{constant}~$c$ تے
$a \in \Domain M$ دتے ہون، تاں $\Struct M[a/c]$ اوہ
!!{structure} تعریف کرو جیہڑی~$\Struct M$ ورگی اے، سوائے ایس دے
کہ $\Assign{c}{M[a/c]} = a$۔ $\Struct M \VDash !A$،
!!{sentence}s~$!A$ لئی ایویں تعریف کرو:
\begin{enumerate}
\tagitem{prvFalse}{%
  \indcase{!A}{\lfalse}{$\Struct{M} \VDash \indfrm$ نہیں۔}}{}

\tagitem{prvTrue}{%
  \indcase{!A}{\ltrue}{$\Struct{M} \VDash \indfrm$۔}}{}

\item \indcase{!A}{\Atom{R}{d_1, \dots, d_n}}{$\Struct M \VDash \indfrm$
  اُدوں تے صرف اُدوں جدوں $\langle \Assign{d_1}{M}, \dots, \Assign{d_n}{M} \rangle \in
  \Assign{R}{M}$۔}

\item \indcase{!A}{\eq[d_1][d_2]}{$\Struct M \VDash \indfrm$ اُدوں تے صرف اُدوں جدوں
  $\Assign{d_1}{M} = \Assign{d_2}{M}$۔}

\tagitem{prvNot}{%
  \indcase{!A}{\lnot !B}{$\Struct{M} \VDash \indfrm$ اُدوں تے صرف اُدوں جدوں
    $\Struct{M} \VDash {!B}$ نہ ہووے۔}}{}

\tagitem{prvAnd}{%
  \indcase{!A}{(!B \land !C)}{$\Struct{M} \VDash
    \indfrm$ اُدوں تے صرف اُدوں جدوں $\Struct{M} \VDash!B$ تے
    $\Struct{M} \VDash !C$ دونیں۔}}{}

\tagitem{prvOr}{%
  \indcase{!A}{(!B \lor !C)}{$\Struct{M} \VDash \indfrm$ اُدوں تے صرف اُدوں جدوں
    $\Struct{M} \VDash !B$ یا $\Struct{M} \VDash !C$، یا دونیں۔}}{}

\tagitem{prvIf}{%
  \indcase{!A}{(!B \lif !C)}{$\Struct{M} \VDash \indfrm$ اُدوں تے صرف اُدوں جدوں
    $\Struct {M} \VDash !B$ نہ ہووے یا $\Struct M \VDash !C$ ہووے، یا دونیں۔}}{}

\tagitem{prvIff}{%
  \indcase{!A}{(!B \liff !C)}{$\Struct{M} \VDash {\indfrm}$ اُدوں تے صرف اُدوں جدوں یا تاں
    $\Struct{M} \VDash {!B}$ تے $\Struct{M} \VDash {!C}$ دونیں، یا
    نہ $\Struct{M} \VDash {!B}$ ہووے نہ $\Struct{M} \VDash {!C}$۔}}{}

\tagitem{prvAll}{%
  \indcase{!A}{\lforall[x][!B]}{$\Struct{M} \VDash {\indfrm}$ اُدوں تے صرف اُدوں جدوں
    ہر $a \in \Domain{M}$ لئی $\Struct{M[a/c]} \VDash \Subst{!B}{c}{x}$،
    جے $c$،~$!B$ وچ نہ آوے۔}}{}

\tagitem{prvEx}{%
  \indcase{!A}{\lexists[x][!B]}{$\Struct{M} \VDash
  {\indfrm}$ اُدوں تے صرف اُدوں جدوں کوئی $a \in \Domain M$ ایسا ہووے
  جیہدے لئی $\Struct{M[a/c]} \VDash \Subst{!B}{c}{x}$، جے
  $c$،~$!B$ وچ نہ آوے۔}}{}
\end{enumerate}
$x_1$، \dots، $x_n$،~$!A$ دے سارے آزاد !!{variable}s ہون؛
$c_1$، \dots، $c_n$،~$!A$ وچ نہ آون والے ثابتہ نشان ہون؛
انہاں تازہ ثابتہ نشاناں نوں آپس وچ وکھرا چنو؛
$a_1$، \dots، $a_n \in \Domain M$ ہون؛ تے $s(x_i) = a_i$۔

ثابت کرو کہ $\Sat{M}{!A}[s]$ اُدوں تے صرف اُدوں جدوں
$\Struct M[a_1/c_1,\dots,a_n/c_n]
\VDash \Subst{\Subst{!A}{c_1}{x_1}\dots}{c_n}{x_n}$۔

(ایہ مسئلہ دکھاندا اے کہ پہلے درجے دی منطق لئی متغیر قدر بندیاں
توں بغیر وی معنویات دتی جا سکدی اے۔)
\end{prob}

\begin{prob}
فرض کرو $f$ اک یک-مقامی فنکشن نشان اے جیہڑا~$!A(x,y)$ وچ نہیں آندا۔
% PNB-SEM-022: قدر بندی توں بغیر صدق تے سکولم بدل لئی جملہ ہون تے جائز بدل دیاں پچھلی شرطاں صریح نیں۔
ایتھے فارمولے دے آزاد متغیر صرف دتے دوناں متغیراں وچوں نیں،
تے دتے فنکشن ٹرم دا بدل پچھلی تعریف مطابق بدل لئی آزاد اے۔
ثابت کرو کہ کوئی !!a{structure}~$\Struct{M}$ ایسی موجود اے جیہدے لئی
$\Sat{M}{\lforall[x][\lexists[y][!A(x,y)]]}$، اُدوں تے صرف اُدوں
جدوں کوئی~$\Struct M'$ ایسی موجود ہووے جیہدے لئی
$\Sat{M'}{\lforall[x][!A(x,f(x))]}$۔

(ایہ مسئلہ سکولم دے قضیے دی اک خاص صورت اے؛
$\lforall[x][!A(x,f(x))]$ نوں
$\lforall[x][\lexists[y][!A(x,y)]]$ دی \emph{سکولم معیاری صورت}
آکھدے آں۔)
\end{prob}
\endgroup
% /OLP-0164

% OLP-0165 content/first-order-logic/syntax-and-semantics/extensionality.tex
\begingroup
\def\olfilename{extensionality}
\expandafter\def\csname ol@sectioncs\endcsname{section}

\olfileid{fol}{syn}{ext}

\olsection{غیر متعلقہ علامتاں توں آزادی}

\begin{explain}
غیر متعلقہ علامتاں توں آزادی، جیہنوں کدے متعلقہ ہون دا اصول وی
آکھدے نیں، عام لفظاں وچ ایویں دسی جا سکدی اے: !!{formula}~$!A$
دا !!a{structure}~$\Struct M$ وچ !!a{variable} قدر بندی~$s$
% PNB-SEM-020: قضیے وچ عین اکّو دائرہ فرض اے؛ اکیلے حجم دی برابری اوہ شرط نہیں۔
دے لحاظ نال پورا اترنا صرف خود !!{domain} تے اُن
تعبیراں اُتے منحصر اے جیہڑیاں~$\Struct M$ تے~$s$ زبان دے
اوہناں اجزا نوں دیندے نیں جیہڑے واقعی~$!A$ وچ آندے نیں۔

ایس اصول دا اک سیدھا نتیجہ ایہ اے: جے دو !!{structure}s~$\Struct M$
تے~$\Struct M'$ کسے جملے~$!A$ وچ آون والے زبان دے سارے اجزا
دی تعبیر اُتے متفق ہون تے اوہناں دا دائرہ اک ہووے، تاں
~$\Struct M$ تے~$\Struct M'$ نوں ایہ گل وی اکّو جیہی دسنی
پئے گی کہ $!A$ سچ اے یا نہیں۔
\end{explain}

\begin{prop}[غیر متعلقہ علامتاں توں آزادی]
  \ollabel{prop:extensionality}
  فرض کرو $!A$ اک !!a{formula} اے، $\Struct M_1$ تے $\Struct M_2$
  دو !!{structure}s نیں جیہناں لئی $\Domain{M_1} = \Domain{M_2}$،
  تے $s$ ایس مشترک $\Domain{M_1} = \Domain{M_2}$ اُتے متغیر
  قدر بندی اے۔ جے $\Assign{c}{M_1} = \Assign{c}{M_2}$،
  $\Assign{R}{M_1}=\Assign{R}{M_2}$ تے
  $\Assign{f}{M_1} = \Assign{f}{M_2}$ ہر اُس !!{constant}~$c$،
  تعلقی نشان~$R$، تے !!{function}~$f$ لئی جیہڑا~$!A$ وچ آندا اے،
  تاں $\Sat{M_1}{!A}[s]$ اُدوں تے صرف اُدوں جدوں
  $\Sat{M_2}{!A}[s]$۔
\end{prop}

\begin{proof}
  پہلاں~$t$ اُتے استقرا نال ثابت کرو کہ فارمولے وچ آون والے ہر ٹرم لئی
  % PNB-SEM-013: ماخذ ہر ٹرم آکھدا اے، پر تعبیراں دی برابری صرف فارمولے وچ آون والیاں علامتاں لئی فرض اے۔
  $\Value{t}{M_1}[s] = \Value{t}{M_2}[s]$۔ فیر~$!A$ اُتے
  استقرا کرکے قضیہ ثابت کرو؛ بنیادی استقرائی صورت وچ، جدوں $!A$
  ایٹمی ہووے، ہُن ثابت کیتا دعویٰ ورتو۔
\end{proof}

\begin{prob}
\olref[fol][syn][ext]{prop:extensionality} دا ثبوت پوری تفصیل نال
دیو۔
\end{prob}

\begin{cor}[!!^{sentence}s لئی غیر متعلقہ علامتاں توں آزادی]
  \ollabel{cor:extensionality-sent}
  فرض کرو $!A$ اک !!a{sentence} اے تے $\Struct{M_1}$، $\Struct{M_2}$
  اوہ نیں جیہڑے \olref{prop:extensionality} وچ دتے نیں۔ تاں
  $\Sat{M_1}{!A}$ اُدوں تے صرف اُدوں جدوں $\Sat{M_2}{!A}$۔
\end{cor}

\begin{proof}
\olref[ass]{cor:sat-sentence} دے ذریعے ایہ \olref{prop:extensionality}
توں نکلدا اے۔
\end{proof}

ایس توں ودھ، ٹرم دی قدر تے !!a{structure} دا !!a{formula}
نوں پورا کرنا یا نہ کرنا، اوہدے ذیلی ٹرماں دیاں قدراں اُتے
ہی منحصر ہوندا اے۔

\begin{prop}\ollabel{prop:ext-terms}
فرض کرو $\Struct M$ اک !!a{structure} اے، $t$ تے $t'$ ٹرم نیں،
تے $s$ متغیر قدر بندی اے۔ تاں
$\Value{\Subst{t}{t'}{x}}{M}[s] =
\Value{t}{M}[\Subst{s}{\Value{t'}{M}[s]}{x}]$۔
\end{prop}

\begin{proof}
~$t$ اُتے استقرا نال۔
\begin{enumerate}
\item جے $t$ ثابتہ اے، مثلاً $t\ident c$، تاں $\Subst{t}{t'}{x} =
  c$، تے $\Value{c}{M}[s] = \Assign{c}{M} =
  \Value{c}{M}[\Subst{s}{\Value{t'}{M}[s]}{x}]$۔

\item جے $t$،~$x$ توں وکھرا متغیر اے، مثلاً $t \ident y$، تاں
  $\Subst{t}{t'}{x} = y$، تے $\Value{y}{M}[s] =
  \Value{y}{M}[\Subst{s}{\Value{t'}{M}[s]}{x}]$، کیوں جے
  $\varAssign{s}{\Subst{s}{\Value{t'}{M}[s]}{x}}{x}$۔

\item جے $t \ident x$، تاں $\Subst{t}{t'}{x} = t'$۔ پر
  $\Value{x}{M}[\Subst{s}{\Value{t'}{M}[s]}{x}] = \Value{t'}{M}[s]$،
  ~ $\Subst{s}{\Value{t'}{M}[s]}{x}$ دی تعریف مطابق۔

\item جے $t \ident \Atom{f}{t_1,\dots,t_n}$، تاں ساڈے کول ایہ اے:
% PNB-SEM-014: شروع دا دہرایا برابر نشان ہٹایا اے؛ اگلی سطر دا برابر نشان برقرار اے۔
\begin{multline*}
  \Value{\Subst{t}{t'}{x}}{M}[s]  \\
  \begin{aligned}[b]
& = \Value{\Atom{f}{\Subst{t_1}{t'}{x}, \dots, \Subst{t_n}{t'}{x}}}{M}[s]\\
& \qquad    \text{ $\Subst{t}{t'}{x}$ دی تعریف مطابق}\\
& = \Assign{f}{M}(\Value{\Subst{t_1}{t'}{x}}{M}[s], \dots,
    \Value{\Subst{t_n}{t'}{x}}{M}[s])\\
    & \qquad  \text{ $\Value{\Atom{f}{\dots}}{M}[s]$ دی تعریف مطابق}\\
& = \Assign{f}{M}(\Value{t_1}{M}[\Subst{s}{\Value{t'}{M}[s]}{x}], \dots,
   \Value{t_n}{M}[\Subst{s}{\Value{t'}{M}[s]}{x}])\\
& \qquad    \text{ استقرائی مفروضے مطابق}\\
& = \Value{t}{M}[\Subst{s}{\Value{t'}{M}[s]}{x}]
    \text{ $\Value{\Atom{f}{\dots}}{M}[\Subst{s}{\Value{t'}{M}[s]}{x}]$ دی تعریف مطابق}
  \end{aligned}
\end{multline*}
\end{enumerate}
\end{proof}

\begin{prop}\ollabel{prop:ext-formulas} فرض کرو $\Struct M$ اک
!!a{structure} اے، $!A$ اک !!a{formula}، $t'$ اک ٹرم تے $s$
متغیر قدر بندی اے۔
% PNB-SEM-019: فارمولے وچ بدل پچھلی تعریف مطابق اوہدوں ای جائز اے جدوں دتا ٹرم بدل لئی آزاد ہووے۔
ایہ شرط وی اے کہ دتے فارمولے وچ منتخب متغیر لئی دتا ٹرم بدل لئی
آزاد ہووے۔ تاں $\Sat{M}{\Subst{!A}{t'}{x}}[s]$ اُدوں تے
صرف اُدوں جدوں $\Sat{M}{!A}[\Subst{s}{\Value{t'}{M}[s]}{x}]$۔
\end{prop}

\begin{proof}
مشق۔
\end{proof}

\begin{prob}
\olref[fol][syn][ext]{prop:ext-formulas} ثابت کرو۔
\end{prob}

\begin{explain}
  \cref{fol:syn:ext:prop:ext-terms,fol:syn:ext:prop:ext-formulas} دی
  گل ایہ اے: فرض کرو ساڈے کول ٹرم $t$ یا !!a{formula}~$!A$ اے
  تے کوئی ہور ٹرم~$t'$ اے، تے اسیں $\Subst{t}{t'}{x}$ دی قدر
  یا $\Subst{!A}{t'}{x}$ دا !!a{structure}~$\Struct M$
  وچ !!a{variable} قدر بندی~$s$ دے لحاظ نال پورا اترنا جاننا
  چاہندے آں۔ فارمولے والے بدل لئی اوپر دتی جائز ہون دی شرط ضروری
  اے۔ تاں اسیں پہلے بدل کرکے فیر~$\Struct{M}$ تے~$s$
  دے لحاظ نال قدر یا پورا اترنا ویکھ سکدے آں؛ یا پہلے قدر~$m = \Value{t'}{M}[s]$
  معلوم کر سکدے آں، جیہڑی~$t'$ دی~$\Struct{M}$ وچ~$s$ دے لحاظ
  نال قدر اے؛ ایہ حساب کرکے !!{variable} قدر بندی نوں~$\Subst{s}{m}{x}$
  بنا سکدے آں، تے فیر~$t$ دی قدر~$\Struct{M}$ تے
  ~$\Subst{s}{m}{x}$ دے لحاظ نال، یا
  $\Sat{M}{!A}[\Subst{s}{m}{x}]$ ویکھ سکدے آں۔
  \Cref{fol:syn:ext:prop:ext-terms,fol:syn:ext:prop:ext-formulas}
  ضمانت دیندے نیں کہ دوناں راہواں نال اک ہی جواب آوے گا۔
\end{explain}
\endgroup
% /OLP-0165

% OLP-0166 content/first-order-logic/syntax-and-semantics/semantic-notions.tex
\begingroup
\def\olfilename{semantic-notions}
\expandafter\def\csname ol@sectioncs\endcsname{section}

\olfileid{fol}{syn}{sem}
\olsection{معنوی تصورات}

\begin{explain}
پہلے درجے دیاں زباناں لئی !!{structure}s دی تعریف مل جان توں بعد
اسیں جملیاں دیاں کجھ بنیادی معنوی خاصیتاں تے انہاں دے آپسی
تعلقاں دی تعریف کر سکدے آں۔ انہاں وچ سبھ توں سادہ تصور
جملے دا \emph{منطقی اعتبار} اے۔ جملہ معتبر اے جے اوہ ہر
!!{structure} وچ پورا اترے۔ معتبر جملے اوہ نیں جیہڑے اپنے
غیر منطقی نشاناں دی تعبیر چاہے جیویں وی کیتی جائے، پورے اترن۔
ایس لئی معتبر جملیاں نوں \emph{منطقی سچائیاں} وی آکھدے نیں:
اوہ ہر !!{structure} وچ سچے، یعنی پورے اترن والے، ہوندے نیں؛
لہٰذا انہاں دا سچ صرف انہاں وچ آون والے منطقی نشاناں تے انہاں
دی نحوی !!{structure} اُتے منحصر اے، غیر منطقی نشاناں یا انہاں دی تعبیر
اُتے نہیں۔
\end{explain}

\begin{defn}[منطقی اعتبار]
اک جملہ $!A$ \emph{معتبر} اے، $\Entails !A$، اُدوں تے صرف اُدوں
جدوں $\Sat{M}{!A}$ ہر !!{structure}~$\Struct M$ لئی ہووے۔
\end{defn}

\begin{defn}[منطقی نتیجہ]
جملیاں دے سیٹ~$\Gamma$ توں جملہ~$!A$ \emph{منطقی طور تے نکلدا اے}،
$\Gamma \Entails !A$، اُدوں تے صرف اُدوں جدوں ہر اوہ !!{structure}~$\Struct M$
جیہڑے لئی $\Sat{M}{\Gamma}$ ہووے، اوہدے لئی $\Sat{M}{!A}$ وی ہووے۔
\end{defn}


\begin{defn}[پورا ہون جوگتا]
جملیاں دا سیٹ~$\Gamma$ \emph{پورا ہون جوگا} اے جے
$\Sat{M}{\Gamma}$ کسے !!{structure}~$\Struct M$ لئی ہووے۔ جے
$\Gamma$ پورا ہون جوگا نہ ہووے، تاں اوہنوں \emph{پورا ہون جوگا نہیں} آکھدے نیں۔
\end{defn}

\begin{prop}
جملہ $!A$ معتبر اے اُدوں تے صرف اُدوں جدوں $\Gamma \Entails !A$
جملیاں دے ہر سیٹ~$\Gamma$ لئی ہووے۔
\end{prop}

\begin{proof}
اگلی سمت لئی فرض کرو $!A$ معتبر اے، تے $\Gamma$ جملیاں دا
اک سیٹ اے۔ کوئی !!a{structure}~$\Struct M$ لو جیہڑے لئی
$\Sat{M}{\Gamma}$ ہووے۔ چونکہ $!A$ معتبر اے، $\Sat{M}{!A}$؛
لہٰذا $\Gamma \Entails !A$۔

اُلٹی سمت دا متقابلِ عکس ثابت کرن لئی فرض کرو $!A$ معتبر نہیں،
یعنی کوئی !!a{structure}~$\Struct M$ اے جیہڑے لئی $\Sat/{M}{!A}$۔
جے $\Gamma = \{ \ltrue \}$ رکھئے، تاں $\ltrue$ دے معتبر ہون کرکے
$\Sat{M}{\Gamma}$۔ ایس طرح اک !!a{structure}~$\Struct M$
ایہو جیہی اے کہ $\Sat{M}{\Gamma}$ پر $\Sat/{M}{!A}$؛ لہٰذا
$\Gamma$ توں $!A$ منطقی طور تے نہیں نکلدا۔
\end{proof}

\begin{prop}
\ollabel{prop:entails-unsat}
$\Gamma \Entails !A$ اُدوں تے صرف اُدوں جدوں $\Gamma \cup \{\lnot !A\}$
پورا ہون جوگا نہ ہووے۔
\end{prop}

\begin{proof}
اگلی سمت لئی فرض کرو $\Gamma \Entails !A$، تے خلافِ فرض کوئی
!!a{structure}~$\Struct M$ ایہو جیہی اے کہ $\Sat{M}{\Gamma
  \cup \{ \lnot !A \}}$۔ چونکہ $\Sat{M}{\Gamma}$ تے $\Gamma \Entails
!A$، اسیں $\Sat{M}{!A}$ لیندے آں۔ نال ہی، چونکہ
$\Sat{M}{\Gamma\cup \{ \lnot !A \}}$، اسیں $\Sat{M}{\lnot !A}$
وی لیندے آں۔ ایس طرح $\Sat{M}{!A}$ تے $\Sat/{M}{!A}$ دونے
ملدے نیں، جو تضاد اے۔ لہٰذا ایہو جیہی !!{structure}~$\Struct M$
ہو ہی نہیں سکدی، تے $\Gamma \cup \{ \lnot !A \}$ پورا ہون جوگا نہیں۔

الٹی سمت لئی فرض کرو $\Gamma \cup \{ \lnot !A \}$ پورا ہون جوگا نہیں۔
ایس لئی ہر !!{structure}~$\Struct M$ لئی یا تاں
$\Sat/{M}{\Gamma}$ ہووے گا یا $\Sat{M}{!A}$۔ لہٰذا ہر اوہ
!!{structure}~$\Struct M$ جیہڑے لئی $\Sat{M}{\Gamma}$ ہووے،
اوہدے لئی $\Sat{M}{!A}$ وی ہووے گا؛ سو $\Gamma \Entails !A$۔
\end{proof}

\begin{prob}
% PNB-SEM-023: تیجے حصے دا مسوّری نتیجہ تے ثابتہ-بدل جملے ہون لئی صرف چنیا متغیر آزاد ہو سکدا اے۔
تیجے حصے وچ دتے فارمولے دے آزاد متغیر صرف چنے مسوّری متغیر
وچوں نیں، تاں جو معنوی نتیجے وچ آون والیاں عبارتاں جملے ہون۔
\begin{enumerate}
\item دکھاؤ کہ $\Gamma \Entails \bot$ اُدوں تے صرف اُدوں جدوں $\Gamma$ پورا ہون جوگا نہیں۔
\item دکھاؤ کہ $\Gamma \cup \{!A\} \Entails \bot$ اُدوں تے صرف اُدوں جدوں $\Gamma \Entails \lnot !A$۔
\item فرض کرو $c$، $!A$ یا $\Gamma$ وچ نہیں آندا۔ دکھاؤ کہ
  $\Gamma \Entails \lforall[x][!A]$ اُدوں تے صرف اُدوں جدوں $\Gamma \Entails
  \Subst{!A}{c}{x}$۔
\end{enumerate}
\end{prob}

\begin{prop}
جے $\Gamma \subseteq \Gamma'$ تے $\Gamma \Entails !A$، تاں $\Gamma'
\Entails !A$۔
\end{prop}

\begin{proof}
فرض کرو $\Gamma \subseteq \Gamma'$ تے $\Gamma \Entails !A$۔
$\Struct M$ کوئی ایہو جیہی ساخت لو جیہڑے لئی $\Sat{M}{\Gamma'}$؛
تاں $\Sat{M}{\Gamma}$، تے چونکہ $\Gamma \Entails !A$،
سانوں $\Sat{M}{!A}$ ملدا اے۔ لہٰذا جدوں وی $\Sat{M}{\Gamma'}$
ہووے، $\Sat{M}{!A}$ وی ہووے گا؛ ایس لئی $\Gamma' \Entails !A$۔
\end{proof}


\begin{thm}[معنوی استنباط دا قضیہ]
\ollabel{thm:sem-deduction}
$\Gamma \cup \{!A\} \Entails !B$ اُدوں تے صرف اُدوں جدوں
$\Gamma \Entails !A \lif !B$۔
\end{thm}

\begin{proof}
اگلی سمت لئی فرض کرو $\Gamma \cup \{ !A \} \Entails !B$، تے
$\Struct M$ کوئی !!a{structure} اے جیہڑے لئی $\Sat{M}{\Gamma}$۔
جے $\Sat{M}{!A}$، تاں $\Sat{M}{\Gamma \cup \{ !A \} }$؛ تے چونکہ
$\Gamma \cup \{ !A \}$ توں $!B$ منطقی طور تے نکلدا اے،
سانوں $\Sat{M}{!B}$ ملدا اے۔ ایس لئی $\Sat{M}{!A \lif !B}$،
سو $\Gamma \Entails !A \lif !B$۔

الٹی سمت لئی فرض کرو $\Gamma \Entails !A \lif !B$، تے
$\Struct M$ کوئی !!a{structure} اے جیہڑے لئی $\Sat{M}{\Gamma \cup \{ !A
  \}}$۔ فیر $\Sat{M}{\Gamma}$، لہٰذا $\Sat{M}{!A \lif !B}$؛
تے چونکہ $\Sat{M}{!A}$، اسیں $\Sat{M}{!B}$ لیندے آں۔ ایس لئی
جدوں وی $\Sat{M}{\Gamma \cup \{ !A \} }$ ہووے، $\Sat{M}{!B}$
وی ہووے گا؛ سو $\Gamma \cup \{ !A \} \Entails !B$۔
\end{proof}

\begin{prop}\ollabel{prop:quant-terms}
  فرض کرو $\Struct{M}$ اک !!a{structure} اے، $!A(x)$ اک ایہو جیہا
  !!a{formula} اے جیہدا اک آزاد متغیر~$x$ اے، تے $t$~اک بند ٹرم اے۔ تاں:
  \begin{tagenumerate}{prvEx,prvAll}
    \tagitem{prvEx}{$!A(t) \Entails \lexists[x][!A(x)]$}{}
    \tagitem{prvAll}{$\lforall[x][!A(x)] \Entails !A(t)$}{}
  \end{tagenumerate}
\end{prop}

\begin{proof}
  \begin{tagenumerate}{prvEx,prvAll}
    \tagitem{prvEx}{%
      \iftag{probEx}{مشق۔}{فرض کرو $\Sat{M}{!A(t)}$۔ $s$ اک
        ایہو جیہی متغیر قدر بندی لو جیہڑے لئی $s(x) = \Value{t}{M}$۔ فیر
        $\Sat{M}{!A(t)}[s]$، کیوں جے $!A(t)$ اک !!a{sentence} اے۔
        \olref[ext]{prop:ext-formulas} دے ذریعے $\Sat{M}{!A(x)}[s]$۔
        \olref[ass]{prop:sat-quant} دے ذریعے
        $\Sat{M}{\lexists[x][!A(x)]}$۔}}{}
    \tagitem{prvAll}{%
      \iftag{probAll}{مشق۔}{فرض کرو
        $\Sat{M}{\lforall[x][!A(x)]}$۔ $s$ اک ایہو جیہی متغیر
        قدر بندی لو جیہڑے لئی $s(x) = \Value{t}{M}$۔
        \olref[ass]{prop:sat-quant} دے ذریعے $\Sat{M}{!A(x)}[s]$۔
        \olref[ext]{prop:ext-formulas} دے ذریعے $\Sat{M}{!A(t)}[s]$۔
        \olref[ass]{prop:sentence-sat-true} دے ذریعے
        $\Sat{M}{!A(t)}$، کیوں جے $!A(t)$ اک !!a{sentence} اے۔}}{}
  \end{tagenumerate}
\end{proof}

\begin{probtag}{probEx,probAll}
  \olref[fol][syn][sem]{prop:quant-terms} دا ثبوت پورا کرو۔
\end{probtag}
\endgroup
% /OLP-0166

\OLEndChapterHook



\par\noindent\emph{متبادل ترتیب۔} نحو تے معنویات اکّو باب وچ (\texttt{syntax-and-semantics.tex})۔ ایہ ترتیب تدوین جوگ ماخذ فائل وچ موجود اے۔
\endgroup
% /OLP-0159

% OLP-0167 content/first-order-logic/models-theories/models-theories.tex
\begingroup
\def\olfilename{models-theories}
\expandafter\def\csname ol@sectioncs\endcsname{section}

\olchapter{fol}{mat}{نظریے تے انہاں دے نمونے}

% OLP-0168 content/first-order-logic/models-theories/introduction.tex
\begingroup
\def\olfilename{introduction}
\expandafter\def\csname ol@sectioncs\endcsname{section}

\olfileid{fol}{mat}{int}
\olsection{تعارف}

\begin{explain}
مسلّمی طریقے دی ترقی سائنس دی تاریخ دی اک وڈی کامیابی اے، تے
ریاضی دی تاریخ وچ خاص اہمیت رکھدی اے۔ کسے شعبے نوں مسلّماں
دی بنیاد اُتے قائم کردیاں بہت سارے سوال واضح کرن پیندے نیں:
ایہ شعبہ کس بارے اے؟ ایس دے سبھ توں بنیادی تصور کیہڑے نیں؟
اوہ اک دوجے نال کس طرح جڑے نیں؟ کی شعبے دے سارے تصورات
ایہناں بنیادی تصورات دے ذریعے تعریف کیتے جا سکدے نیں؟ ایہ
تصورات کیہڑے قاعدیاں اُتے عمل کردے نیں، تے کیہڑے قاعدیاں اُتے
لازماً عمل کرنا چاہیدا اے؟

مسلّمی طریقہ تے منطق اک دوجے لئی بہت مناسب نیں۔ رسمی منطق
مسلّمی نظریے بنانے، نظریے دے مسلّماں توں ٹھیک ٹھیک متعین
طریقے نال قضیے ثابت کرن، تے مسلّماں نوں پورا کرن والے سارے
نظاماں دیاں خاصیتاں دی باقاعدہ جانچ لئی وسیلے دیندی اے۔
\end{explain}

\begin{defn}
!!{sentence}s دا سیٹ~$\Gamma$ \emph{بند} اے اُدوں تے صرف اُدوں
جدوں $\Gamma \Entails !A$ ہون اُتے $!A \in \Gamma$ وی ہووے۔
!!{sentence}s دے سیٹ~$\Gamma$ دی \emph{بندش}
$\Setabs{!A}{\Gamma \Entails !A}$ اے۔

اسیں آکھدے آں کہ~$\Gamma$ نوں جملیاں دے سیٹ~$\Delta$ نے
\emph{مسلّماں نال متعین کیتا اے} جے $\Gamma$،~$\Delta$ دی بندش اے۔
\end{defn}

\begin{explain}
اسیں مسلّمی نظریے نوں اوہناں جملیاں دے سیٹ دے طور تے سمجھ سکدے آں
جیہڑا اوہدے مسلّماں دے سیٹ~$\Delta$ توں متعین ہوندا اے۔ دوجے
لفظاں وچ، فرض کرو اسیں اک ایہو جیہی پہلے درجے دی زبان ورتدے
آں جیہدے غیر منطقی نشان اُس سائنس دے بنیادی اجزا لئی نیں
جیہنوں اسیں مسلّمی طریقے نال پڑھنا چاہندے آں، تے نال اوہناں
بنیادی قاعدیاں نوں بیان کرن والے !!{sentence}s دا سیٹ اے۔ تاں
اسیں نظریے نوں ایس زبان دے اوہناں سارے !!{sentence}s نال ظاہر
سمجھ سکدے آں جیہڑے مسلّماں توں منطقی طور تے نکلدے نیں۔
ایہ دائرہ اک ہی بنیادی جزو تے سادہ مسلّماں والے نمونیاں، جیویں
جزوی ترتیباں دے نظریے، توں لے کے نیوٹنی میکانیات ورگے
پیچیدہ نظریاں تک پھیلیا ہویا اے۔

مسلّمی طریقے ول ایس رسمی پہنچ نوں اہم بنان والے منطقی حقائق
ہیٹھاں دتے نیں۔ فرض کرو $\Gamma$ کسے نظریے دا مسلّمی نظام،
یعنی جملیاں دا اک سیٹ، اے۔
\begin{enumerate}
\item اسیں ٹھیک ٹھیک دس سکدے آں کہ مسلّمی نظام کدوں مقصودہ
  !!{structure}s دی جماعت نوں پکڑدا اے۔ یعنی جے ساڈی دلچسپی
  !!{structure}s دی کسے خاص جماعت وچ اے، تاں مسلّمی نظام~$\Gamma$
  اوہ جماعت اُدوں تے صرف اُدوں پکڑے گا جدوں !!{structure}s
  ٹھیک اوہی~$\Struct M$ ہون جیہناں لئی $\Sat{M}{\Gamma}$۔
\item اس کم وچ ناکامی ایس کرکے ہو سکدی اے کہ کجھ $\Struct M$
  ایہو جیہے ہون جیہناں لئی $\Sat{M}{\Gamma}$، پر $\Struct M$
  ساڈے مقصودہ !!{structure}s وچ شامل نہ ہووے۔ ایس توں اسیں
  اوہ مسلّمے شامل کرن ول جا سکدے آں جیہڑے~$\Struct M$ وچ سچے نہ ہون۔
\item جے گھٹ توں گھٹ ایس حد تک کامیابی ہووے کہ $\Gamma$ ساریاں
  مقصودہ !!{structure}s وچ سچا ہووے، تاں جملہ~$!A$ وی ساریاں
  مقصودہ !!{structure}s وچ سچا ہووے گا جدوں وی $\Gamma \Entails !A$۔
  ایس لئی اسیں منطقی وسیلے، جیویں !!{derivation} دے طریقے، ورت کے
  جملیاں دا ساریاں مقصودہ !!{structure}s وچ سچا ہونا صرف ایہ
  دکھا کے ثابت کر سکدے آں کہ اوہ مسلّماں توں منطقی طور تے نکلدے نیں۔
\item کدے ساڈے ذہن وچ مقصودہ !!{structure}s نہیں ہوندیاں، بلکہ
  اسیں خود مسلّماں توں شروع کردے آں: کجھ بنیادی اجزا لیندے
  آں جیہناں لئی اسیں چاہندے آں کہ اوہ کجھ قاعدیاں اُتے عمل کرن، تے انہاں
  قاعدیاں نوں مسلّمی نظام وچ درج کردے آں۔ سبھ توں پہلاں اسیں
  ایہ جانچنا چاہندے آں کہ مسلّمے آپس وچ متضاد نہ ہون؛ جے ہون،
  تاں کوئی تصور انہاں قاعدیاں اُتے عمل نہیں کر سکدا تے ساڈا
  بنایا نظریہ غیر موافق ہووے گا۔~$\Gamma$ دا کوئی نمونہ لبھ کے
  اسیں تصدیق کر سکدے آں کہ ایہ مسئلہ نہیں۔ تے جے ساڈے نظریے
  دے نمونے ہون، تاں اسیں منطقی طریقیاں نال انہاں دی جانچ وی
  کر سکدے آں تے نمونے بنا وی سکدے آں۔
\item مسلّماں دا اک دوجے توں آزاد ہونا وی اہم سوال اے۔ ہو سکدا
  اے کہ کوئی مسلّمہ اصل وچ باقیاں دا منطقی نتیجہ ہووے، تے ایس
  لئی فالتو ہووے۔ اسیں ثابت کر سکدے آں کہ مسلّمہ $!A$،
  جیہڑا $\Gamma$ وچ اے، فالتو اے جے ثابت کر دیئے کہ $\Gamma \setminus
  \{!A\} \Entails !A$۔ ایس دے الٹ، ایہ دکھا کے کہ
  $(\Gamma \setminus \{!A\}) \cup \{\lnot
  !A\}$ پورا ہون جوگا اے، اسیں ثابت کر سکدے آں کہ ایہ مسلّمہ
  فالتو نہیں۔ مثال لئی، ایس طریقے نال دکھایا گیا سی کہ متوازی
  خطاں دا مسلّمہ جیومیٹری دے ہور مسلّماں توں آزاد اے۔
\item اک ہور اہم سوال نظریے وچ تصورات دے تعریف دے قابل ہون
  دا اے: زبان دی چُون ایہ طے کردی اے کہ نظریے دے نمونے کیہناں
  چیزاں توں بندے نیں۔ پر نظریے دے ہر پہلو نوں اوہدے نمونیاں
  وچ الگ توں ظاہر کرنا ضروری نہیں۔ مثال لئی، ہر ترتیب
  $\le$ اک متعلقہ سخت ترتیب~$<$ طے کردی اے: اک مل جائے، تاں
  دوجی تعریف کیتی جا سکدی اے۔ ایس لئی ایہ ضروری نہیں کہ
  ایہو جیہی ترتیب والے نظریے دے نمونے وچ متعلقہ سخت ترتیب
  \emph{وی} الگ شامل ہووے۔ عمومی طور تے اک تعلق دوجیاں دے
  ذریعے کدوں تعریف کیتا جا سکدا اے؟ تے کدوں کسے تعلق نوں
  دوجیاں دے ذریعے تعریف کرنا ناممکن اے، جیہدے نتیجے وچ اوہنوں
  زبان دے بنیادی اجزا وچ شامل کرنا پیندا اے؟
\end{enumerate}
\end{explain}
\endgroup
% /OLP-0168

% OLP-0169 content/first-order-logic/models-theories/expressing-props-of-structures.tex
\begingroup
\def\olfilename{expressing-props-of-structures}
\expandafter\def\csname ol@sectioncs\endcsname{section}

\olfileid{fol}{mat}{exs}

\olsection{\printtoken{P}{structure} دیاں خاصیتاں دا اظہار}

\begin{explain}
اکثر فنکشناں تے تعلقاں اُتے شرطاں بیان کرنا مفید تے ضروری ہوندا
اے؛ یا زیادہ عمومی طور تے ایہ بیان کرنا کہ کسے ساخت دے فنکشن
تے تعلق اوہ شرطاں پوریاں کردے نیں۔ مثال لئی اسیں اوہناں
!!{structure}s نوں، جیہڑیاں کسے زبان دے !!{predicate}s دا
ساڈا مطلوبہ ”مطلب“ ٹھیک ”پکڑدیاں“ نیں، اوہناں توں وکھ کرنا
چاہندے آں جیہڑیاں ایہ کم نہیں کردیاں۔ بے شک اسیں پوری آزادی
نال طے کر سکدے آں کہ ساڈیاں ”مقصودہ“ !!{structure}s کیہڑیاں
نیں: مثال لئی، !!{predicate}~$\le$ دی تعبیر اک ترتیب ہونی
چاہیدی اے؛ یا اسیں صرف اوہناں تعبیراں وچ دلچسپی
رکھدے آں جیہناں وچ~$\Lang L$ دی تعبیر لئی دائرہ سیٹاں توں
بنے تے $\Obj \in$ دی تعبیر ”!!a{element} ہون“ دے تعلق نال
کیتی جائے۔ پر کی اسیں ایہ کم خود زبان دے !!{sentence}s نال
کر سکدے آں؟ دوجے لفظاں وچ، !!a{structure}~$\Struct M$ اُتے
کیہڑیاں شرطاں اسیں !!a{sentence} (یا شاید
!!{sentence}s دے سیٹ) نال،~$\Struct M$ دی زبان وچ، بیان کر سکدے آں؟ کجھ شرطاں اسیں
بیان نہیں کر سکاں گے۔ مثال لئی،~$\Lang L_A$ وچ ایہو جیہا
% PNB-MAT-001: ناممکن شرط والا جملہ کوئی نمونہ نہیں رکھدا؛ دائرے نوں عین قدرتی اعداد تک پابند کرن دا دعویٰ پورا ہون جوگے جملے بارے اے۔
کوئی پورا ہون جوگا جملہ نہیں جیہڑا !!{structure}~$\Struct M$ وچ صرف ایس
صورت سچا ہووے جے $\Domain M = \Nat$۔ اسیں ایہ بیان نہیں کر
سکدے کہ ”دائرے وچ صرف قدرتی عدد نیں“۔ ایتھے گل دائرے دیاں
اصل چیزاں نوں عین ایہناں عدداں تک پابند کرن دی اے، اوہدے
حسابی ڈھانچے بارے شرطاں بیان کرن دی نہیں۔ پر !!{structure}s دیاں
کجھ ”ساختی خاصیتاں“ شاید اسیں بیان کر سکئیے۔ !!{structure}s
دیاں کیہڑیاں خاصیتاں اسیں !!{sentence}s راہیں بیان کر سکدے آں؟
یا دوجے لفظاں وچ، !!{structure}s دیاں کیہڑیاں جماعتاں نوں
اس طرح بیان کر سکدے آں کہ اوہ !!a{sentence} (یا
!!{sentence}s دے سیٹ) نوں سچا بناندیاں نیں؟
\end{explain}

\begin{defn}[سیٹ دا نمونہ]
فرض کرو $\Gamma$، زبان~$\Lang L$ وچ !!{sentence}s دا سیٹ
اے۔ اسیں آکھدے آں کہ !!a{structure}~$\Struct M$،~$\Gamma$
دا \emph{نمونہ اے} جے $\Sat{M}{!A}$ ہر $!A \in \Gamma$ لئی ہووے۔
\end{defn}

\begin{ex}
جملہ $\lforall[x][x \le x]$،~$\Struct M$ وچ اُدوں تے صرف
اُدوں سچا اے جدوں $\Assign{\le}{M}$ انعکاسی تعلق ہووے۔ جملہ
$\lforall[x][\lforall[y][((x \le y \land y \le x) \lif x = y)]]$
،~$\Struct M$ وچ اُدوں تے صرف اُدوں سچا اے جدوں
$\Assign{\le}{M}$ ضد تناظری ہووے۔ جملہ
$\lforall[x][\lforall[y][\lforall[z][((x \le y \land y \le z)
      \lif x \le z)]]]$،~$\Struct M$ وچ اُدوں تے صرف اُدوں
سچا اے جدوں $\Assign{\le}{M}$ متعدی ہووے۔ ایس لئی
\begin{align*}
\{\quad &\lforall[x][x \le x], \\
   & \lforall[x][\lforall[y][((x \le y \land y \le
    x) \lif x = y)]], \\
   &\lforall[x][\lforall[y][\lforall[z][((x \le y
      \land y \le z) \lif x \le z)]]] \quad \}
\end{align*}
دے نمونے ٹھیک اوہی ساختاں نیں جیہناں وچ~$\Assign{\le}{M}$
انعکاسی، ضد تناظری تے متعدی اے؛ یعنی جزوی ترتیب اے۔ ایس لئی
اسیں انہاں نوں \emph{جزوی ترتیباں دے پہلے درجے دے نظریے} دے
مسلّمے بنا سکدے آں۔
\end{ex}
\endgroup
% /OLP-0169

% OLP-0170 content/first-order-logic/models-theories/theories.tex
\begingroup
\def\olfilename{theories}
\expandafter\def\csname ol@sectioncs\endcsname{section}

\olfileid{fol}{mat}{the}

\olsection{پہلے درجے دے نظریاں دیاں مثالاں}

\begin{ex}
زبان~$\Lang L_<$ وچ سخت خطی ترتیباں دا نظریہ ہیٹھلے سیٹ
دے مسلّماں نال متعین کیتا جاندا اے:
\begin{align*}
\{\quad & \lforall[x][\lnot x < x], \\
& \lforall[x][\lforall[y][((x < y \lor y <
    x) \lor x = y)]], \\
& \lforall[x][\lforall[y][\lforall[z][((x < y
      \land y < z) \lif x < z)]]] \quad \}
\end{align*}
ایہ ساڈیاں مقصودہ !!{structure}s نوں پوری طرح پکڑدا اے:
ہر سخت خطی ترتیب ایس مسلّمی نظام دا نمونہ اے؛ تے الٹ وی، جے $R$
% PNB-SEM-015: الٹی سمت وچ سخت خطی ترتیب چاہیدی اے، صرف انعکاسی خطی ترتیب نہیں۔
سیٹ $X$ اُتے سخت خطی ترتیب اے، تاں اوہ ساخت $\Struct M$ جیہدے لئی
$\Domain M = X$ تے $\Assign{<}{M} = R$، ایس نظریے دا نمونہ اے۔
\end{ex}

\begin{ex}
گروپاں دا نظریہ، جیہدی زبان وچ $\Obj 1$ (!!{constant}) تے $\cdot$
(دو تھانواں والا !!{function}) نیں، ہیٹھلے مسلّماں نال متعین اے:
\begin{align*}
& \lforall[x][\eq[(x \cdot \Obj 1)][x]]\\
& \lforall[x][\lforall[y][\lforall[z][\eq[(x \cdot (y \cdot z))][((x
          \cdot y) \cdot z)]]]]\\
& \lforall[x][\lexists[y][\eq[(x \cdot y)][\Obj 1]]]
\end{align*}
\end{ex}

\begin{ex}
پیانو حساب دا نظریہ، حساب دی زبان~$\Lang L_A$ وچ ہیٹھلے
جملیاں نال مسلّمی بنیاد پاؤندا اے۔
\begin{align*}
& \lforall[x][\lforall[y][(\eq[x'][y'] \lif \eq[x][y])]]\\
& \lforall[x][\eq/[\Obj 0][x']]\\
& \lforall[x][\eq[(x + \Obj 0)][x]]\\
& \lforall[x][\lforall[y][\eq[(x + y')][(x + y)']]]\\
& \lforall[x][\eq[(x \times \Obj 0)][\Obj 0]]\\
& \lforall[x][\lforall[y][\eq[(x \times y')][((x \times y) + x)]]]\\
& \lforall[x][\lforall[y][(x < y \liff \lexists[z][\eq[(z' + x)][y])]]]\\
\intertext{نال ہی ایس روپ دے سارے جملے وی}
& (!A(\Obj 0) \land \lforall[x][(!A(x) \lif !A(x'))]) \lif \lforall[x][!A(x)]
\end{align*}
چونکہ پچھلے روپ دے بے انتہا جملے نیں، ایس لئی ایہ مسلّمی
نظام لامحدود اے۔ پچھلے روپ نوں \emph{استقرائی خاکہ} آکھدے نیں۔
(اصل وچ، استقرائی خاکہ ایتھے دسے گئے توں تھوڑا ہور پیچیدہ اے۔)

% PNB-SEM-024: ایہ اشارہ استقرائی خاکے توں پہلاں والی آخری حسابی شق ول اے، خود خاکے ول نہیں۔
استقرائی خاکے توں پہلاں والا آخری مسلّمہ~$<$ دی اک \emph{واضح تعریف} اے۔
\end{ex}

\begin{ex}
خالص سیٹاں دا نظریہ ریاضی دیاں بنیاداں وچ، تے ریاضی دے فلسفے
وچ وی، اہم کردار رکھدا اے۔ سیٹ خالص اے جے اوہدے سارے
!!{element}s وی خالص سیٹ ہون۔ ایس لئی خالی سیٹ وی خالص اے،
پر جیہڑے سیٹ دا کوئی !!a{element} خود سیٹ نہ ہووے، اوہ خالص
نہیں۔ سو خالص سیٹ اوہ نیں جیہڑے صرف خالی سیٹ توں بنے ہون
تے انہاں وچ کوئی ”غیر سیٹ بنیادی شے“، یعنی اوہ چیز جیہڑی
آپ سیٹ نہ ہووے، شامل نہ ہووے۔

ہیٹھلے نوں خالص سیٹاں دے نظریے دا مسلّمی نظام سمجھیا جا سکدا اے:
\begin{align*}
& \lexists[x][\lnot \lexists[y][y \in x]]\\
& \lforall[x][\lforall[y][(\lforall[z](z \in x \liff z \in y) \lif 
\eq[x][y])]]\\
& \lforall[x][\lforall[y][\lexists[z][\lforall[u][(u \in z \liff
(\eq[u][x] \lor \eq[u][y]))]]]]\\
& \lforall[x][\lexists[y][\lforall[z][(z \in y \liff \lexists[u][(z \in
        u \land u \in x)])]]]\\
\intertext{نال ہی ایس روپ دے سارے جملے وی} &
\lexists[x][\lforall[y][(y \in x \liff !A(y))]]
\end{align*}
پہلا مسلّمہ آکھدا اے کہ اک ایہو جیہا سیٹ اے جیہدے کوئی
!!{element}s نہیں (یعنی $\emptyset$ موجود اے)؛ دوجا آکھدا اے
کہ سیٹاں دی برابری انہاں دے عنصراں نال طے ہوندی اے؛ تریجا
آکھدا اے کہ ہر دو سیٹاں $X$ تے $Y$ لئی سیٹ $\{X, Y\}$ موجود
اے؛ تے چوتھا کہ ہر سیٹ $X$ لئی سیٹ $\cup X$ موجود اے، جتھے
$\cup X$،~$X$ دے سارے عنصراں دا اتحاد اے۔

آخر وچ دتے گئے !!{sentence}s نوں اکٹھے
\emph{سادہ خاصیت توں سیٹ بنان دا خاکہ} آکھدے نیں۔ ایہ بنیادی
طور تے آکھدا اے کہ ہر $!A(x)$ لئی سیٹ $\Setabs{x}{!A(x)}$
موجود اے؛ پہلی نظر وچ ایہ سچا، مفید، تے شاید ضروری وی لگدا
اے۔ ایہنوں ”سادہ“ اس لئی آکھدے نیں کہ حقیقت وچ ایہ نظریے
نوں پورا ہون جوگا نہیں رہن دیندا: جے $!A(y)$ نوں $\lnot y \in y$
لو، تاں ایہ !!{sentence} ملدا اے:
\[
\lexists[x][\lforall[y][(y \in x \liff \lnot y \in y)]]
\]
تے ایہ !!{sentence} کسے وی !!{structure} وچ پورا نہیں اتردا۔
\end{ex}

\begin{ex}
\emph{حصہ-نظریے} (mereology) وچ \emph{حصہ ہونا} (parthood)
اک بنیادی تعلق اے۔ سیٹاں دے نظریاں وانگ حصے ہون دے وی
نظریے نیں جیہڑے ایس تعلق دیاں وکھ وکھ، تے کدے آپس وچ
ٹکراندیاں، سمجھاں نوں مسلّماں راہیں متعین کردے نیں۔

حصہ-نظریے دی زبان وچ اک ہی دو تھانواں والا محمولی نشان
~$\Obj P$ اے، تے $\Atom{\Obj P}{x, y}$ دا ”مطلب“ ایہ اے
کہ $x$،~$y$ دا حصہ اے۔ ایس تعبیر نوں ذہن وچ رکھ کے، ایس
زبان دی !!a{structure} نوں \emph{حصیاں دی ساخت} آکھدے نیں۔
بے شک اک دو تھانواں والے محمول دی ہر ساخت ایس ناں دی حق دار
نہیں۔ ”حصہ ہون“ نوں پکڑن لئی $\Assign{\Obj P}{M}$ نوں کجھ
شرطاں پوریاں کرنیاں چاہیدیاں نیں، جیہڑیاں اسیں حصہ-نظریے دے
مسلّمے بنا سکدے آں۔ مثال لئی چیزاں وچ حصہ ہون دا تعلق جزوی
ترتیب اے: ہر چیز اپنے آپ دا حصہ اے، بھاویں \emph{غیر سخت}
حصہ؛ دو وکھریاں چیزاں اک دوجے دے حصے نہیں ہو سکدیاں؛
تے کسے چیز دے حصے دا حصہ خود وی اُس چیز دا حصہ اے۔ غور کرو
کہ ایس معنی وچ ”حصہ ہون“ ”ذیلی سیٹ ہون“ نال ملدا اے، پر
”رکن ہون“ نال نہیں، کیوں جے رکن ہون نہ انعکاسی اے نہ متعدی۔
\begin{align*}
& \lforall[x][\Atom{\Obj P}{x,x}] \\
& \lforall[x][\lforall[y][((\Part{x}{y} \land \Part{y}{x})
      \lif \eq[x][y])]] \\
& \lforall[x][\lforall[y][\lforall[z][((\Part{x}{y} \land
        \Part{y}{z}) \lif \Part{x}{z})]]]\\
\intertext{ایس توں ودھ، ہر دو چیزاں لئی حصیاں دا اک مجموعہ ہوندا اے (اک ایہو جیہی چیز جیہدے
  حصہ اوہ دونے ہون، تے ایس لحاظ نال اوہ سبھ توں چھوٹی ہووے)۔}  &
\lforall[x][\lforall[y][\lexists[z][\lforall[u][(\Part{z}{u} \liff
        (\Part{x}{u} \land \Part{y}{u}))]]]]
\end{align*}
ایہ حصہ ہون دے اوہناں بنیادی اصولاں وچوں صرف کجھ نیں جنہاں
اُتے مابعدالطبیعیات دے ماہر غور کردے نیں۔ پر ہور اصول، کجھ
تعریف شدہ تعلق پہلے متعارف کرائے بغیر، جلد ہی بیان یا لکھن
وچ اوکھے ہو جاندے نیں۔ مثال لئی حصہ-نظریے وچ دلچسپی رکھن
والے زیادہ تر ماہر ایہنوں وی معتبر اصول منندے نیں: جدوں وی
کسے چیز~$x$ دا اک سخت حصہ~$y$ ہووے، تاں اوہدا اک ہور
حصہ~$z$ وی ہووے جیہدا~$y$ نال کوئی مشترک حصہ نہ ہووے،
تے $y$ تے $z$ دے مجموعے نال~$x$ بنے۔
\end{ex}
\endgroup
% /OLP-0170

% OLP-0171 content/first-order-logic/models-theories/expressing-relations.tex
\begingroup
\def\olfilename{expressing-relations}
\expandafter\def\csname ol@sectioncs\endcsname{section}

\olfileid{fol}{mat}{exr}

\olsection{\article{structure} \printtoken{S}{structure} وچ تعلقاں دا اظہار}

\begin{explain}
!!{formula}s دا اک وڈا ورتارا ایہ اے کہ !!a{structure}~$\Struct M$ دیاں خاصیتاں
تے تعلقاں نوں زبان~$\Lang L$ دیاں بنیادی علامتاں راہیں ظاہر
کیتا جا سکدا اے، جیہڑی~$\Struct M$ دی زبان اے۔ ایہدے نال ساڈی مراد ایہ اے:
$\Struct M$ دا !!{domain} اشیا دا سیٹ اے۔ !!{constant}s،
!!{function}s تے !!{predicate}s دی~$\Struct M$ وچ تشریح بالترتیب
$\Domain M$ دے کجھ عناصر، $\Domain M$ اُتے فنکشناں تے
$\Domain M$ اُتے تعلقاں راہیں ہوندی اے۔ مثال لئی، جے
$\Obj A^2_0$، $\Lang L$ وچ ہووے، تاں $\Struct M$ ایہدے نال
تعلق~$R = \Assign{{\Obj A^2_0}}{M}$ جوڑدی اے۔ فیر فارمولا
$\Atom{\Obj A^2_0}{\Obj v_1, \Obj v_2}$ ہی اس تعلق دا
\emph{اظہار کردا اے}، ایس معنی وچ: جے متغیراں دی قدر بندی~$s$،
$\Obj v_1$ نوں $a \in \Domain{M}$ تے
$\Obj v_2$ نوں $b \in \Domain M$ ول بھیجے، تاں
\[
Rab \text{\quad جے تے صرف جے\quad} \Sat{M}{\Atom{\Obj A^2_0}{\Obj v_1, \Obj v_2}}[s].
\]
ایتھے متغیراں دی قدر بندی لانی ضروری اے: اسیں ایہ نہیں کہہ
سکدے کہ ”$Rab$ جے تے صرف جے
$\Sat{M}{\Atom{\Obj A^2_0}{a, b}}$“، کیونکہ $a$ تے
  $b$ ساڈی زبان دیاں علامتاں نہیں، بلکہ
  $\Domain{M}$ دے !!{element}s نیں۔

ساڈے کول صرف ایٹمی !!{formula}s ای نہیں؛ اسیں انہاں نوں منطقی
رابطیاں تے مسوّراں نال جوڑ وی سکدے آں۔ ایس لئی ودھ پیچیدہ
!!{formula}s اوہ ہور تعلق تعریف کر سکدے نیں جیہڑے
$\Struct M$ وچ سدھے بنیادی تعلقاں دے طور اُتے شامل نہیں۔
ساڈی دلچسپی ایس وچ اے کہ ایہ کم کیویں ہوندا اے، تے خاص طور اُتے
ایہ کہ کسے !!a{structure} وچ کیہڑے تعلق تعریف کیتے جا سکدے نیں۔
\end{explain}

\begin{defn}
فرض کرو $!A(\Obj v_1,\dots, \Obj v_n)$، $\Lang L$ دا
!!a{formula} ہووے جس وچ صرف $\Obj v_1$،\dots،
$\Obj v_n$ آزاد آون، تے $\Struct M$، $\Lang L$ لئی
!!a{structure} ہووے۔ $!A(\Obj v_1,\dots, \Obj v_n)$
تعلق~$R \subseteq \Domain M^n$ دا \emph{اظہار کردا اے} جے تے
صرف جے
\[
Ra_1\dots a_n \text{\quad جے تے صرف جے\quad} \Sat{M}{\Atom{!A}{\Obj
    v_1,\dots, \Obj v_n}}[s]
\]
ہر اوہ متغیر قدر بندی~$s$ لئی، جیہدے وچ
$s(\Obj v_i) = a_i$ ($i = 1, \dots, n$)، ایہ برابری قائم ہووے۔
\end{defn}

\begin{ex}
حساب دے معیاری نمونے~$\Struct N$ وچ !!{formula}
$\Obj v_1 < \Obj v_2 \lor \eq[\Obj v_1][\Obj v_2]$،
تعلق $\le$ دا اظہار کردا اے، جیہڑا $\Nat$ اُتے اے۔ !!{formula}
$\eq[\Obj v_2][\Obj v_1']$ جانشین تعلق دا اظہار کردا اے؛
یعنی اوہ تعلق $R \subseteq \Nat^2$ جتھے $Rnm$ ایس صورت وچ
قائم اے جدوں $m$، $n$ دا جانشین ہووے۔ فارمولا
$\eq[\Obj v_1][\Obj v_2']$ پیشرو تعلق دا اظہار کردا اے۔
!!{formula}s
$\lexists[\Obj v_3][(\eq/[\Obj v_3][\Obj 0] \land
  \eq[\Obj v_2][(\Obj v_1 + \Obj v_3)])]$ تے $\lexists[\Obj
  v_3][\eq[(\Obj v_1 + {\Obj v_3}')][v_2]]$ دونے
$\Obj <$ تعلق دا اظہار کردے نیں۔ ایس دا مطلب اے کہ حساب دی زبان
وچ محمول علامت~$<$ اصل وچ لازمی نہیں؛ ایہنوں تعریف کیتا جا سکدا اے۔
\end{ex}

\begin{explain}
ایہ خیال صرف خاص !!{structure}s وچ ای دلچسپ نہیں، بلکہ عام طور اُتے
وی اوہدوں اہم اے جدوں اسیں کسے زبان نال اک مطلوب نمونے یا نمونیاں
دا بیان کردے آں، یعنی جدوں اسیں نظرییاں اُتے غور کردے آں۔
ایہ نظریے اکثر صرف تھوڑے جیہے !!{predicate}s نوں بنیادی
علامتاں دے طور اُتے رکھدے نیں، پر جیہڑے دائرے دا اوہ بیان کردے
نیں، اوہدے وچ ہور بہت سارے تعلق وی اہم ہوندے نیں۔ جے ایہ ہور
تعلق، زبان دے بنیادی !!{predicate}s دی تشریح کرن والے تعلقاں
راہیں باقاعدہ طور اُتے ظاہر کیتے جا سکن، تاں اسیں کہندے آں کہ
انہاں نوں اس زبان وچ \emph{تعریف} کیتا جا سکدا اے۔
\end{explain}

\begin{prob}
$\Lang L_A$ وچ اوہ !!{formula}s لبھو جیہڑے ہیٹھلے تعلقاں
دی تعریف کردے نیں:
\begin{enumerate}
\item $n$، $i$ تے $j$ دے درمیان اے؛
\item $n$، $m$ نوں پورا ونڈدا اے (یعنی $m$، $n$ دا مضرب اے)؛
% PNB-MAT-002: اصل دی وضاحت اک نوں وی شامل کردی سی؛ اولی عدد اک توں وڈا ہونا ضروری اے۔
\item $n$ اولی عدد اے (یعنی اوہ اک توں وڈا اے تے $1$ تے $n$ توں علاوہ ہور کوئی عدد
  $n$ نوں پورا نہیں ونڈدا)۔
\end{enumerate}
\end{prob}

\begin{prob}
فرض کرو فارمولا $!A(\Obj v_1, \Obj v_2)$، تعلق
$R \subseteq \Domain M^2$ دا اظہار !!a{structure}~$\Struct M$
وچ کردا اے۔ اوہ
فارمولے لبھو جیہڑے ہیٹھلے تعلقاں دا اظہار کردے نیں:
\begin{enumerate}
\item اُلٹا تعلق $R^{-1}$، جیہڑا $R$ دا اُلٹا اے؛
\item نسبتی حاصل ضرب $R \mid R$؛
\end{enumerate}
کی تسی $R^+$، یعنی $R$ دی تعدی تکمیل، دا اظہار کرن دا کوئی
طریقہ لبھ سکدے او؟
\end{prob}

\begin{prob}
فرض کرو $\Lang{L}$ اوہ زبان اے جیہدے وچ صرف دو تھانواں
والی محمول علامت $<$ ہووے (ہور کوئی !!{constant}s،
!!{function}s یا !!{predicate}s نہ ہون---البتہ
$\eq$ ضرور ہووے)۔ $\Struct{N}$ اوہ ساخت ہووے جیہدے لئی
$\Domain{N} = \Nat$ تے
$\Assign{<}{N} = \Setabs{\tuple{n,m}}{n <
  m}$۔ ہیٹھلے دعویاں دا ثبوت دیو:
\begin{enumerate}
\item $\{ 0 \}$، $\Struct{N}$ وچ تعریف دے قابل اے؛
\item $\{ 1 \}$، $\Struct{N}$ وچ تعریف دے قابل اے؛
\item $\{ 2 \}$، $\Struct{N}$ وچ تعریف دے قابل اے؛
\item ہر $n \in \Nat$ لئی سیٹ $\{ n \}$،
  $\Struct{N}$ وچ تعریف دے قابل اے؛
\item $\Domain{N}$ دا ہر محدود ذیلی سیٹ
  $\Struct{N}$ وچ تعریف دے قابل اے؛
\item $\Domain{N}$ دا ہر ہم متناہی ذیلی سیٹ
  $\Struct{N}$ وچ تعریف دے قابل اے (جتھے
  $X \subseteq \Nat$ ہم متناہی اے جے تے صرف جے
  $\Nat \setminus X$ محدود ہووے)۔
\end{enumerate}
\end{prob}
\endgroup
% /OLP-0171

% OLP-0172 content/first-order-logic/models-theories/set-theory.tex
\begingroup
\def\olfilename{set-theory}
\expandafter\def\csname ol@sectioncs\endcsname{section}

\olfileid{fol}{mat}{set}

\olsection{سیٹاں دا نظریہ}

تقریباً ساری ریاضی سیٹاں دے نظریے وچ تشکیل دتی جا سکدی اے۔
ایس نظریے وچ ریاضی دی تشکیل لئی کئی کم کرن پینڈے نیں۔ سب توں
پہلاں تعلق~$\in$ لئی مسلّماں دا اک سیٹ چاہیدا اے۔ کجھ وکھرے
مسلّمی نظام بنائے گئے نیں، جنہاں وچ کدی کدی~$\in$ دیاں
خاصیتاں وی اک دوجے نال ٹکرا جاندیاں نیں۔ $\Log{ZFC}$ ناں دا
مسلّمی نظام، یعنی چناؤ دے مسلّمے نال زرمیلو--فرینکل سیٹ نظریہ،
انہاں وچ نمایاں اے: ایہ سب توں ودھ ورتیا تے پڑھیا گیا نظام اے،
کیونکہ ایہدے مسلّمے تقریباً اوہ ساریاں گلاں ثابت کرن لئی کافی
نکلدے نیں جیہڑیاں ریاضی دان ثابت کر سکن دی امید رکھدے نیں۔
پر ایہ قائم کرن توں پہلاں ایہ واضح کرن ضروری اے کہ ریاضی دان
جیہڑیاں گلاں کہنا چاہندے نیں، انہاں دا \emph{اظہار} کیویں
کیتا جا سکدا اے۔ شروع وچ ای زبان وچ کوئی !!{constant}s یا
!!{function}s نہیں، ایس لئی پہلی نظر وچ ایہ صاف نہیں کہ اسیں
کسے خاص سیٹ (مثلاً $\emptyset$ یا $\Nat$) بارے گل کیویں کریے،
سیٹاں اُتے عملاں (مثلاً $X \cup Y$ تے $\Pow{X}$) بارے کیویں
گل کریے، یا سیٹاں توں علاوہ ہور چیزاں، جیویں تعلقاں تے
فنکشناں، نوں شامل کرن والیاں بناوٹاں بارے گل کیویں کریے۔

شروع وچ ”دا رکن ہونا“ واحد تعلق نہیں جیہدے نال ساڈی دلچسپی اے:
”دا ذیلی سیٹ ہونا“ وی تقریباً اونا ای اہم لگدا اے۔ پر اسیں
”دا ذیلی سیٹ ہونا“ نوں ”دا رکن ہونا“ دے ذریعے \emph{تعریف}
کر سکدے آں۔ ایہدے لئی سیٹ نظریے دی زبان وچ اک
!!a{formula}~$!A(x, y)$ لبھنا اے جیہڑا سیٹاں دی جوڑی
$\tuple{X, Y}$ لئی جے تے صرف جے پورا اترے جدوں
$X \subseteq Y$ ہووے۔ پر $X$، $Y$ دا ذیلی سیٹ اوہدوں تے
صرف اوہدوں اے جدوں~$X$ دے سارے !!{element}s،
~$Y$ دے وی !!{element}s ہون۔ ایس لئی اسیں
$\subseteq$ نوں ہیٹھلے فارمولے نال تعریف کر سکدے آں:
\[
\lforall[z][(z \in x \lif z \in y)]
\]
ہن جتھے وی کسے فارمولے وچ تعلق~$\subseteq$ ورتنا ہووے،
اسیں اوہدی تھاں ایہ فارمولا ورت سکدے آں ($x$ تے $y$ نوں
مناسب طرح بدل کے، تے لوڑ پئے تاں پابند متغیر~$z$ دا ناں
بدل کے)۔ مثال لئی، سیٹاں دے عنصراں نال تعیّن دا مطلب ایہ
اے کہ جے کوئی دو سیٹ~$x$ تے $y$ اک دوجے وچ شامل ہون،
تاں $x$ تے $y$ اکّو سیٹ ہون گے۔ ایہ گل
$\lforall[x][\lforall[y][((x \subseteq y \land y
    \subseteq x) \lif x = y)]]$ نال ظاہر کیتی جا سکدی اے،
یا جے $\subseteq$ دی تھاں اُپرلی تعریف رکھیے، تاں
\[
\lforall[x][\lforall[y][((\lforall[z][(z \in x \lif z \in y)] \land
    \lforall[z][(z \in y \lif z \in x)]) \lif x = y)]].
\]
ایہ حقیقت وچ $\Log{ZFC}$ دے مسلّماں وچوں اک اے، یعنی
”عنصراں نال تعیّن دا مسلّمہ“۔

$\emptyset$ لئی کوئی !!{constant} نہیں، پر ”$x$ خالی اے“
نوں $\lnot \lexists[y][y \in x]$ نال ظاہر کیتا جا سکدا اے۔
فیر ”$\emptyset$ موجود اے“ !!{sentence}~$\lexists[x][\lnot \lexists[y][y \in
    x]]$ بن جاندا اے۔ ایہ~$\Log{ZFC}$ دا اک ہور مسلّمہ اے۔
(دھیان دیو کہ عنصراں نال تعیّن دا مسلّمہ ایہ لازم کردا اے
کہ خالی سیٹ صرف اک ای اے۔) سیٹ نظریے دی زبان وچ جدوں وی
$\emptyset$ بارے گل کرنی ہووے، اسیں ایویں لکھاں گے:
”اک سیٹ اے جیہڑا خالی اے تے \dots“۔ مثال دے طور اُتے،
ایہ ظاہر کرن لئی کہ $\emptyset$ ہر سیٹ دا ذیلی سیٹ اے،
اسیں لکھ سکدے آں
\[
\lexists[x][(\lnot\lexists[y][y \in x] \land \lforall[z][x \subseteq
    z])]
\]
جتھے بے شک $x \subseteq z$ نوں وی اوہدی تعریف نال
بدلنا پئے گا۔

سیٹاں اُتے عملاں، جیویں $X \cup Y$ تے $\Pow{X}$، بارے گل
کرن لئی وی اسیں ایہو جیہا طریقہ ورتنا اے۔ سیٹ نظریے دی
زبان وچ فنکشن دیاں علامتاں نہیں، پر فنکشنی تعلق
$X \cup Y = Z$ تے $\Pow{X} = Y$ نوں اسیں ایویں ظاہر
کر سکدے آں:
\begin{align*}
& \lforall[u][((u \in x \lor u \in y) \liff u \in z)]\\
& \lforall[u][(u \subseteq x \liff u \in y )]
\end{align*}
کیونکہ $X \cup Y$ دے !!{element}s عین اوہ سیٹ نیں جیہڑے
یا~$X$ دے !!{element}s ہون یا~$Y$ دے !!{element}s،
تے $\Pow{X}$ دے !!{element}s عین~$X$ دے ذیلی سیٹ نیں۔
پر ایہدے نال اسیں $x \cup y$ یا $\Pow{x}$ نوں ٹرماں وانگ
نہیں ورت سکدے؛ اسیں صرف اوہ پورے !!{formula}s ورت سکدے
آں جیہڑے $X \cup Y = Z$ تے $\Pow{X} = Y$ تعلقاں دی
تعریف کردے نیں۔ حقیقت وچ اسیں اجے ایہ وی نہیں جاندے کہ
ایہ تعلق کدے پورے وی اتردے نیں یا نہیں؛ یعنی ایہ نہیں
جاندے کہ اتحاد سیٹ تے قوت سیٹ ہمیشہ موجود ہوندے نیں۔
مثال لئی !!{sentence}
$\lforall[x][\lexists[y][\Pow{x} = y]]$،
~$\Log{ZFC}$ دا اک ہور مسلّمہ اے (قوت سیٹ دا مسلّمہ)۔

ہن ترتیب وار جوڑیاں یا فنکشناں بارے کیویں گل کریے؟ ایتھے
سانوں دسنا پئے گا کہ ترتیب وار جوڑیاں تے فنکشناں نوں
خاص قسم دے سیٹ کیویں سمجھیا جا سکدا اے۔ ترتیب وار جوڑی
$\tuple{x, y}$ دی اک تعریف سیٹ $\{\{x\}, \{x, y\}\}$
اے۔ پر پہلاں وانگ اسیں ایس سیٹ دا ناں رکھن والا
!!a{function} شامل نہیں کر سکدے؛ اسیں صرف تعلق
$\tuple{x, y} = z$، یعنی $\{\{x\}, \{x, y\}\} = z$،
دی تعریف کر سکدے آں:
\[
\lforall[u][(u \in z \liff (\lforall[v][(v \in u \liff v = x)] \lor
  \lforall[v][(v \in u \liff (v = x \lor v = y))]))]
\]
ایہ کہندا اے کہ !!{element}s~$u$، جیہڑے~$z$ وچ نیں، عین اوہ سیٹ
نیں جنہاں دا اکلوتا !!{element} $x$ اے، یا جنہاں دے
صرف !!{element}s $x$ تے~$y$ نیں؛ دوجے لفظاں وچ،
اوہ سیٹ جیہڑے یا $\{x\}$ دے برابر نیں یا
$\{x, y\}$ دے برابر نیں۔ ایہ بناوٹ ملن توں بعد اسیں
ہور گل وی کہہ سکدے آں، مثلاً $X \times Y = Z$:
\[
\lforall[z][(z \in Z \liff \lexists[x][\lexists[y][(x \in
        X \land y \in Y \land \tuple{x, y} = z)]])]
\]

فنکشن $f \colon X \to Y$ نوں تعلق $f(x)
= y$ دے روپ وچ، یعنی جوڑیاں دے سیٹ
$\Setabs{\tuple{x,y}}{f(x) = y}$ دے روپ وچ سمجھیا جا سکدا اے۔
فیر اسیں کہہ سکدے آں کہ سیٹ~$f$، $X$ توں $Y$ ول فنکشن
اے جے (الف) ایہ $\subseteq X \times Y$ تعلق اے، (ب) ایہ
کُل اے، یعنی ہر $x \in X$ لئی کوئی $y \in Y$ اے
جیہدے لئی $\tuple{x, y} \in f$ ہووے، تے (ج) ایہ
فنکشنی اے، یعنی جدوں وی $\tuple{x, y}, \tuple{x, y'} \in f$
ہوون، $y = y'$ ہووے (کیونکہ اکّو انپٹ لئی فنکشن دی
قدر اکّو ہونی چاہیدی اے)۔ ایس لئی ”$f$، $X$ توں $Y$ ول
فنکشن اے“ ایویں لکھیا جا سکدا اے:
% PNB-SYN-008: فنکشن دے بیان وچ باہرلے مشروط بیان دی رہ گئی قوس بند کیتی اے۔
\begin{align*}
 \lforall[u][(u \in f \lif {}] & \lexists[x][\lexists[y][(x \in X \land y \in
      Y \land \tuple{x, y} = u)]]) \land {}\\
 \lforall[x][(x \in X \lif {}] &
      (\lexists[y][(y \in Y \land \mathrm{maps}(f, x, y))] \land {}\\
& (\lforall[y][\lforall[y'][((\mathrm{maps}(f, x, y) \land
    \mathrm{maps}(f, x, y')) \lif y = y')]])))
\end{align*}
جتھے $\mathrm{maps}(f, x, y)$، $\lexists[v][(v \in f \land
  \tuple{x, y} = v)]$ دا مخفف اے (ایہ !!{formula}،
”$f(x) = y$“ دا اظہار کردا اے)۔

ہن ایہ ظاہر کرن وی اوکھا نہیں کہ $f\colon X \to Y$
!!{injective} اے، مثلاً:
\begin{multline*}
 f \colon X \to Y \land \lforall[x][\lforall[x'][((x \in X \land x' \in
     X \land {}]] \\
 \lexists[y][(\mathrm{maps}(f, x, y) \land \mathrm{maps}(f,
        x', y))]) \lif x = x')
\end{multline*}
فنکشن~$f\colon X \to Y$ !!{injective} اے جے تے صرف
جے، جدوں وی $f$، $x, x' \in X$ نوں اکّو~$y$
ول بھیجے، $x = x'$ ہووے۔ جے اسیں ایس فارمولے
دا مخفف $\mathrm{inj}(f, X, Y)$ رکھ لئیے، تاں ہن
سیٹ نظریے دی زبان وچ کانتور دے قضیے ورگی اہم گل
وی کہہ سکدے آں: $\Pow{X}$ توں $X$ ول کوئی
!!{injective} فنکشن نہیں:
\[
\lforall[X][\lforall[Y][(\Pow{X} = Y \lif
    \lnot\lexists[f][\mathrm{inj}(f, Y, X)])]]
\]

کسے نوں لگ سکدا اے کہ سیٹ نظریے نوں اک ہور مسلّمہ چاہیدا
اے جیہڑا ہر وصف لئی اک سیٹ دی موجودگی یقینی بناوے۔
% PNB-SEM-025: ایتھے دکھائے گئے جملے لئی ہور آزاد متغیر نہیں؛ پیرامیٹر والے عمومی خاکے دی بندش ہیٹھاں واضح اے۔
جے $!A(x)$ سیٹ نظریے دا اوہ فارمولا اے جیہدے وچ
صرف متغیر~$x$ آزاد ہو سکدا اے، تاں اسیں !!{sentence}
\[
\lexists[y][\lforall[x][(x \in y \liff !A(x))]].
\]
اُتے غور کر سکدے آں۔ ایہ !!{sentence} کہندا اے کہ
سیٹ~$y$ موجود اے جیہدے !!{element}s اوہ سارے تے
صرف اوہ $x$ نیں جیہڑے~$!A(x)$ نوں پورا کردے نیں۔
ایس خاکے نوں ”خاصیت توں سیٹ بنان دا اصول“ کہندے نیں۔
ایہ بڑا کارآمد لگدا اے؛ پر بدقسمتی نال ایہ غیر موافق اے۔
$!A(x) \ident \lnot x \in x$ رکھو؛ فیر خاصیت توں
سیٹ بنان دا اصول کہندا اے:
\[
\lexists[y][\lforall[x][(x \in y \liff x \notin x)]],
\]
یعنی اوہ سارے سیٹاں دا اک سیٹ موجود اے جیہڑے
اپنے آپ دے !!{element}s نہیں۔ ایسا کوئی سیٹ موجود
نہیں ہو سکدا---ایہ رسل دا پیراڈوکس اے۔ حقیقت وچ
$\Log{ZFC}$ ایس اصول دا اک محدود---تے موافق---روپ
رکھدا اے، جیہنوں وکھ کرن دا اصول کہندے نیں:
% PNB-SYN-008: وکھ کرن دے خاکے وچ باہرلی دو طرفہ شرط دی رہ گئی قوس بند کیتی اے۔
\[
\lforall[z][\lexists[y][\lforall[x][(x \in y \liff (x \in z \land
      !A(x)))]]].
\]
عمومی خاکے وچ وادھو آزاد متغیراں نوں پیرامیٹر من کے پورے
بیان نوں انہاں اُتے ہمہ گیر مسوّراں نال بند کیتا جاندا اے؛
باہرلے مسوّراں دے متغیر نام تازہ رکھے جاندے نیں۔ ”موافق“
ایتھے وکھ کرن دے اصول بارے اے، پورے سیٹ نظریے دی بےتضادی
دا ثبوت نہیں۔

\begin{prob}
اک ایداں دا !!a{derivation} دے کے دکھاؤ جیہدے نال ثابت ہووے
کہ خاصیت توں سیٹ بنان دا اصول غیر موافق اے، یعنی
\[
\lexists[y][\lforall[x][(x \in y \liff x \notin x)]] \Proves \lfalse.
\]
پہلاں ایہ ثابت کرنا مددگار ہو سکدا اے کہ
$(A \lif \lnot A) \land (\lnot A \lif A)
\Proves \lfalse$۔
\end{prob}
\endgroup
% /OLP-0172

% OLP-0173 content/first-order-logic/models-theories/size-of-structures.tex
\begingroup
\def\olfilename{size-of-structures}
\expandafter\def\csname ol@sectioncs\endcsname{section}

\olfileid{fol}{mat}{siz}

\olsection{\printtoken{P}{structure} دا سائز ظاہر کرنا}

\begin{explain}
ساختاں دیاں کجھ خاصیتاں اسیں زبان دیاں غیر منطقی علامتاں
ورتے بغیر وی ظاہر کر سکدے آں۔ مثال لئی کجھ
!!{sentence}s ایداں دے نیں جیہڑے کسے !!a{structure}
وچ جے تے صرف جے سچے ہوندے نیں جدوں اس !!{structure}
دے !!{domain} وچ !!{element}s دی گنتی کسے خاص
مثبت عدد~$n$ توں گھٹ نہ ہووے، ودھ نہ ہووے، یا عین اوہنی ہووے۔
% PNB-SEM-026: نقطیاں والے خاکے مثبت گنتی مطابق محدود جوڑاں دے مخفف نیں؛ اک لئی خالی جوڑ سچ اے۔
ہیٹھاں نقطیاں والیاں عبارتیں مطلوب گنتی مطابق محدود جوڑاں
تے متبادلاں دے مخفف نیں۔ اک عنصر لئی جوڑی وار وکھرے ہون
دیاں شرطاں دی خالی جوڑ-فہرست نوں سچ سمجھیا جاندا اے۔
\end{explain}

\begin{prop}
!!{sentence}
\begin{multline*}
  \formula A_{\ge n} \ident \lexists[x_1][\lexists[x_2][\dots\lexists[x_n][{}]]]\\
  \begin{aligned}
  (\eq/[x_1][x_2] \land {}
  \eq/[x_1][x_3] \land \eq/[x_1][x_4] \land \dots \land \eq/[x_1][x_n] \land {}\\
\eq/[x_2][x_3] \land \eq/[x_2][x_4] \land \dots \land {} \eq/[x_2][x_n] \land {} \\
\vdots\\
\eq/[x_{n-1}][x_n])
\end{aligned}
\end{multline*}
!!a{structure}~$\Struct M$ وچ جے تے صرف جے سچا اے
جدوں $\Domain M$ وچ گھٹ توں گھٹ $n$ !!{element}s
ہون۔ ایس لئی $\Sat{M}{\lnot !A_{\ge n+1}}$ جے تے
صرف جے جدوں $\Domain M$ وچ ودھ توں ودھ~$n$
!!{element}s ہون۔

\end{prop}

\begin{prop}
!!{sentence}
\begin{multline*}
\formula A_{= n} \ident \lexists[x_1][\lexists[x_2][\dots\lexists[x_n][{}]]] \\
\begin{aligned}
  (\eq/[x_1][x_2] \land {}
  \eq/[x_1][x_3] \land \eq/[x_1][x_4] \land \dots \land \eq/[x_1][x_n] \land {}\\
\eq/[x_2][x_3] \land \eq/[x_2][x_4] \land \dots \land {} \eq/[x_2][x_n] \land {} \\
\vdots\\
\eq/[x_{n-1}][x_n] \land {} \\
\lforall[y][(\eq[y][x_1] \lor \dots \lor \eq[y][x_n]]))
\end{aligned}
\end{multline*}
!!a{structure}~$\Struct M$ وچ جے تے صرف جے سچا اے
جدوں $\Domain M$ وچ عین $n$ !!{element}s ہون۔
\end{prop}

\begin{prop}
اک !!{structure} لامتناہی اے جے تے صرف جے اوہ
\[
\{!A_{\ge 1}, !A_{\ge 2}, !A_{\ge 3}, \dots \}.
\]
دا نمونہ اے۔
\end{prop}

ایداں دا کوئی اک خالص منطقی جملہ نہیں جیہڑا~$\Struct
M$ وچ جے تے صرف جے سچا ہووے جدوں $\Domain M$
لامتناہی ہووے۔ البتہ غیر منطقی !!{predicate}s ورتن
والے ایداں دے !!{sentence}s دتے جا سکدے نیں جنہاں
دے صرف لامتناہی نمونے ہون (بھاویں ہر لامتناہی
!!{structure} انہاں دا نمونہ نہ ہووے)۔ محدود ساخت
ہون دی خاصیت !!{sentence}s دے لامتناہی سیٹ نال وی ظاہر
نہیں کیتی جا سکدی۔
% PNB-SEM-027: ناقابلِ شمار ہون بارے ایہ دعویٰ ایتھے ورتی گننے جوگی زبان لئی اے۔
ایتھے ورتی گننے جوگی زبان وچ !!{nonenumerable} ساخت ہون دی
خاصیت وی جملیاں دے لامتناہی سیٹ نال ظاہر
نہیں کیتی جا سکدی۔ ایہ حقیقتاں محدود گواہی دے
قضیے تے لووینہائم--سکولم (L\"owenheim--Skolem)
دے قضیاں توں نکلدیاں نیں۔
\endgroup
% /OLP-0173

\OLEndChapterHook
\endgroup
% /OLP-0167

% OLP-0063 content/first-order-logic/proof-systems/proof-systems.tex
\begingroup
\def\olfilename{proof-systems}
\expandafter\def\csname ol@sectioncs\endcsname{section}

\iftag{FOL}
      {\olchapter{fol}{prf}{\usetoken{S}{derivation} نظام \pnbLatin{(Derivation Systems)}}}
      {\olchapter{pl}{prf}{\usetoken{S}{derivation} نظام \pnbLatin{(Derivation Systems)}}}

\begin{editorial}
ایہ باب !!{derivation} نظاماں \pnbLatin{(derivation systems)} بارے عام مواد اکٹھا کردا اے۔ کوئی درسی
کتاب جیہڑی کسے خاص نظام نوں ورتدی ہووے، اوس نظام دا تعارف کراؤن والے
باب دے شروع وچ تعارفی حصہ \pnbLatin{(introduction section)} تے نال متعلقہ جائزے
والا حصہ \pnbLatin{(relevant survey section)} شامل کر سکدی اے۔
\end{editorial}

% OLP-0064 content/first-order-logic/proof-systems/introduction.tex
\begingroup
\def\olfilename{introduction}
\expandafter\def\csname ol@sectioncs\endcsname{section}

\iftag{FOL}
      {\olfileid{fol}{prf}{int}}
      {\olfileid{pl}{prf}{int}}

\olsection{تعارف \pnbLatin{(Introduction)}}

منطقاں وچ عام طور تے معنیات (semantics) وی ہُندی اے تے !!a{derivation}
نظام \pnbLatin{(derivation system)} وی۔ معنیات دا واسطہ صدق (truth)، پورا ہون
جوگتا (satisfiability)، منطقی اعتبار (validity) تے منطقی نتیجے
(entailment) ورگے تصورات نال اے۔ !!{derivation} نظاماں دا مقصد منطقی
نتیجہ تے اعتبار قائم کرن دا سراسر نحوی طریقہ \pnbLatin{(purely syntactic method)}
دینا اے۔ ایہ ایس معنی وچ سراسر نحوی نیں کہ ایہو جہے نظام وچ کوئی
!!{derivation} اک محدود نحوی شے \pnbLatin{(finite syntactic object)}
ہُندی اے، عام طور تے !!{sentence}s یا !!{formula}s دا کوئی سلسلہ (یا
ہور محدود بندوبست)۔ چنگے !!{derivation} نظاماں وچ ایہ خاصیت ہُندی اے
کہ !!{sentence}s یا !!{formula}s دے دِتے ہوئے ہر سلسلے یا بندوبست نوں
مشینی طریقے نال جانچیا جا سکدا اے کہ اوہ ``درست'' اے (mechanically
verified as ``correct'')۔

پہلے درجے دی منطق لئی سب توں سادے (تے تاریخی طور تے پہلے)
!!{derivation} نظام \emph{مسلّمی (axiomatic)} سن۔ ایہو جہے نظام وچ !!{formula}s
دا کوئی سلسلہ !!a{derivation} اُدوں گِنیا جاندا اے جے اوہدے وچ ہر
وکھرا !!{formula} یا تاں ``مسلّماں'' (axioms) دے کسے مقرر سیٹ وچ شامل ہووے،
یا مقرر گنتی دے ``استنباطی قاعدیاں'' \pnbLatin{(inference rules)} وچوں کسے اک راہیں سلسلے وچ اپنے
توں پہلاں آون والے !!{formula}s توں نکلدا ہووے---تے ایہ مشینی طریقے
نال جانچیا جا سکدا اے کہ کی !!a{formula} مسلّمہ اے تے کی اوہ استنباطی
قاعدیاں وچوں کسے اک راہیں ہور !!{formula}s توں ٹھیک طرح نکلدا اے۔
مسلّمی !!{derivation} نظاماں نوں بیان کرنا---تے ماورائے نظریاتی طور
تے (metatheoretically) سنبھالنا وی---سَوکھا اے، پر اوہناں وچ !!{derivation}s نوں پڑھنا تے
سمجھنا اوکھا اے، تے اوہناں نوں بنانا وی اوکھا اے۔

ہور !!{derivation} نظام ایس مقصد نال بنائے گئے نیں کہ !!{derivation}s
نوں بنانا سَوکھا ہووے یا مکمل ہو جان توں بعد !!{derivation}s نوں
سمجھنا سَوکھا ہووے۔ مثالاں فطری استنباط \pnbLatin{(natural deduction)}، صداقتی
درخت \pnbLatin{(truth trees)}، جنہاں نوں تابلو ثبوت \pnbLatin{(tableaux proofs)} وی آکھیا
جاندا اے، تے سیکوئنٹ حساب \pnbLatin{(sequent calculus)} نیں۔ کجھ !!{derivation}
نظام خاص طور تے مشینی عمل (mechanization) نوں سامنے رکھ کے تیار کیتے
گئے نیں، مثلاً ریزولیوشن (resolution) دا طریقہ سافٹ ویئر وچ نافذ کرنا سَوکھا اے (پر اوہدیاں
!!{derivation}s نوں سمجھنا لگ بھگ ناممکن اے)۔ ایہناں ہور !!{derivation}
نظاماں وچوں زیادہ تر، !!{derivation}s نوں سلسلیاں دی بجائے !!{formula}s
دے درختاں \pnbLatin{(trees of formulas)} دے طور تے پیش کردے نیں۔ ایس نال حصیاں
دا انحصار (dependency) ویکھنا سَوکھا ہو جاندا اے: !!a{derivation} دے
کیہڑے حصے ہور کیہڑے حصیاں اُتے نربھر نیں۔

سو کسی دِتّی منطق، جیویں پہلے درجے دی منطق، لئی وکھ وکھ
!!{derivation} نظام ایس گل دیاں وکھ وکھ واضح تعریفاں دیندے نیں کہ
!!a{sentence} دا \emph{قضیہ (theorem)} ہون دا کی مطلب اے تے !!a{sentence} دا
کجھ ہور جملیاں توں !!{derivable} (derivable) ہون دا کی مطلب اے۔ ایہ کم جیویں وی
کیتا جاوے (مسلّمی !!{derivation}s، فطری استنباطاں، سیکوئنٹ
!!{derivation}s، صداقتی درختاں یا ریزولیوشن نال ردّاں راہیں)، اسیں
چاہندے آں کہ ایہ تعلق منطقی اعتبار (validity) تے منطقی نتیجے
(entailment) دے معنوی تصورات نال
میل کھان۔ آؤ $\Proves !A$ نوں ``$!A$ اک قضیہ اے'' لئی تے
``$\Gamma \Proves !A$'' نوں ``$!A$،~$\Gamma$ توں !!{derivable} اے''
لئی لکھئیے۔ $\Proves$ دی تعریف جیویں وی کیتی جاوے، اسیں چاہندے آں کہ
اوہ $\Entails$ نال میل کھاوے، یعنی:
\begin{enumerate}
\item $\Proves !A$ اُدوں تے صرف اُدوں جدوں $\Entails !A$
\item $\Gamma \Proves !A$ اُدوں تے صرف اُدوں جدوں $\Gamma \Entails !A$
\end{enumerate}
اُپر دتے ``صرف اُدوں'' والے رخ نوں \emph{صحت مندی (soundness)} آکھیا جاندا اے۔
!!^a{derivation} نظام صحت مند اے جے !!{derivability} (derivability)، منطقی نتیجے (یا
اعتبار) دی ضمانت دیندی اے۔ ہر ڈھنگ دے !!{derivation} نظام دا صحت مند
ہونا لازم اے؛ غیر صحت مند !!{derivation} نظام بالکل فائدہ مند نہیں۔
آخرکار، !!a{derivation} دا پورا مقصد اعتبار یا منطقی نتیجے دی نحوی
ضمانت دینا اے۔ اسیں اپنے پیش کیتے !!{derivation} نظاماں دی صحت مندی
ثابت کراں گے۔

اُلٹا ``اُدوں'' والا رخ وی اہم اے: ایہنوں \emph{تمامیت (completeness)} آکھیا جاندا
اے۔ تمامیت والا !!{derivation} نظام اینا طاقتور ہُندا اے کہ جدوں وی $!A$
منطقی طور تے معتبر ہووے، اوہ $!A$ نوں قضیہ وکھا سکے، تے
$\Gamma \Proves !A$ وی وکھا سکے جدوں وی $\Gamma \Entails !A$۔ تمامیت
قائم کرنا ودھ اوکھا اے، تے کجھ منطقاں دا کوئی تمامیت والا !!{derivation} نظام
نہیں۔ پہلے درجے دی منطق دا اے۔ کرٹ گوڈل (Kurt Gödel) اوہ پہلا شخص سی جہنے 1929 دے
اپنے ڈاکٹری مقالے وچ پہلے درجے دی منطق دے !!a{derivation} نظام لئی
تمامیت ثابت کیتی۔

!!{derivation} نظاماں نال جڑیا اک ہور تصور \emph{بے تضادی (consistency)} دا اے۔
!!{sentence}s دے سیٹ نوں متضاد (inconsistent) آکھیا جاندا اے جے اوس توں
جو کجھ وی !!{derive} کیتا جا سکدا ہووے، تے نہیں تاں بے تضاد (consistent)۔
متضاد ہونا، پورا ہون جوگ نہ ہون (unsatisfiability) دا نحوی ہم منصب
\pnbLatin{(syntactic counterpart)} اے: پورا ہون جوگ نہ ہون والے سیٹاں
وانگ، !!{sentence}s دے متضاد سیٹ چنگے نظریے نہیں بناؤندے؛ اوہ بنیادی
طور تے خراب ہُندے نیں۔ !!{sentence}s دے بے تضاد سیٹ ضروری نہیں کہ سچے
یا فائدہ مند ہون، پر گھٹّو گھٹ منطقی فائدے دی ایہ سب توں مُڈھلی حد
پار کردے نیں۔ وکھ وکھ !!{derivation} نظاماں لئی !!{sentence}s دے سیٹ
دی بے تضادی دی خاص تعریف وکھری ہو سکدی اے، پر~$\Proves$ وانگ اسیں
چاہندے آں کہ بے تضادی اپنے معنوی ہم منصب، پورا ہون جوگتا
(satisfiability)، نال اکّو
ہووے۔ اسیں چاہندے آں کہ ہمیشہ ایہ گل ہووے کہ $\Gamma$ اُدوں تے صرف
اُدوں بے تضاد اے جدوں اوہ پورا ہون جوگ اے۔ ایتھے ``صرف اُدوں'' والا
رخ تمامیت دے برابر اے (بے تضادی، پورا ہون جوگتا دی ضمانت دیندی اے)،
تے ``اُدوں'' والا رخ صحت مندی دے برابر اے (پورا ہون جوگتا، بے تضادی
دی ضمانت دیندی اے)۔ اصل وچ، کلاسیکی پہلے درجے دی منطق لئی صحت مندی
تے تمامیت دے دونیں روپ ہم ارز نیں۔
\endgroup
% /OLP-0064

% OLP-0065 content/first-order-logic/proof-systems/sequent-calculus.tex
\begingroup
\def\olfilename{sequent-calculus}
\expandafter\def\csname ol@sectioncs\endcsname{section}

\iftag{FOL}
      {\olfileid{fol}{prf}{seq}}
      {\olfileid{pl}{prf}{seq}}

\olsection{سیکوئنٹ حساب \pnbLatin{(Sequent Calculus)}}

بھاویں بہت سارے !!{derivation} نظام !!{sentence}s دے بندوبستاں نال کم
کردے نیں، سیکوئنٹ حساب \pnbLatin{(sequent calculus)} \emph{سیکوئنٹاں (sequents)}
نال کم کردا اے۔ سیکوئنٹ ایس
روپ دی اک عبارت اے
\[
!A_1, \dots, !A_m \Sequent !B_1, \dots, !B_m,
\]
یعنی !!{sentence}s دے سلسلیاں دی اک جوڑی، جہنوں سیکوئنٹ دی علامت
\pnbLatin{(sequent symbol)}~$\Sequent$ وکھ کردی اے۔ دونیں سلسلیاں وچوں کوئی وی
خالی ہو سکدا اے۔ سیکوئنٹ حساب وچ !!^a{derivation} (derivation)
سیکوئنٹاں دا اک درخت ہُندا اے،
جتھے سب توں اُتّے والے سیکوئنٹ خاص روپ دے ہُندے نیں (اوہناں نوں
``ابتدائی سیکوئنٹ \pnbLatin{(initial sequents)}'' یا ``مسلّماں (axioms)'' آکھیا
جاندا اے) تے ہر ہور سیکوئنٹ استنباط دے قاعدیاں \pnbLatin{(rules of inference)}
وچوں کسے اک راہیں اپنے عین اُتّے والے سیکوئنٹاں توں
نکلدا اے۔ استنباط دے قاعدے یا تاں سیکوئنٹاں اندر !!{sentence}s نوں
بدلدے نیں (کھبے یا سجے پاسے اوہناں نوں شامل، خارج یا اگے پچھے کر کے)،
یا قاعدے دے نتیجے وچ اک مرکب !!{formula} متعارف کراؤندے نیں۔ مثال لئی،
$\LeftR{\land}$ دا قاعدہ $!A,
\Gamma \Sequent \Delta$ توں $!A \land !B, \Gamma \Sequent \Delta$ دا
استنباط کرن دی اجازت دیندا اے، تے $\RightR{\lif}$ دا قاعدہ $!A, \Gamma \Sequent
\Delta, !B$ توں $\Gamma \Sequent \Delta, !A \lif !B$ دا، ہر $\Gamma$،
$\Delta$، $!A$ تے~$!B$ لئی۔ (خاص طور تے، $\Gamma$ تے~$\Delta$ خالی
وی ہو سکدے نیں۔)

سیکوئنٹ حساب اُتے مبنی $\Proves$ تعلق \pnbLatin{(derivability relation)} دی تعریف
ایویں کیتی جاندی اے:
$\Gamma \Proves !A$ اُدوں تے صرف اُدوں جدوں کوئی سلسلہ $\Gamma_0$ ایہو
جہا ہووے کہ ہر $!A$ جیہڑا $\Gamma_0$ وچ اے،~$\Gamma$ وچ ہووے، تے اک
!!{derivation} ہووے جہدی جڑ (root) اُتے سیکوئنٹ~$\Gamma_0 \Sequent !A$
ہووے۔ $!A$ سیکوئنٹ حساب وچ اک قضیہ (theorem) اے جے
سیکوئنٹ~$\Sequent !A$ دا
!!a{derivation} ہووے۔ مثال لئی، ایہ !!a{derivation} وکھاؤندا اے کہ
$\Proves (!A \land !B) \lif !A$:
\begin{prooftree}
  \Axiom$!A \fCenter !A$
  \RightLabel{\LeftR{\land}}
  \UnaryInf$!A \land !B \fCenter !A$
  \RightLabel{\RightR{\lif}}
  \UnaryInf$\fCenter (!A \land !B) \lif !A$
\end{prooftree}

سیکوئنٹ حساب وچ سیٹ $\Gamma$ متضاد (inconsistent) اے جے
$\Gamma_0 \Sequent$ دا
!!a{derivation} ہووے (جتھے ہر $!A \in \Gamma_0$،~$\Gamma$
وچ ہووے تے سیکوئنٹ دا سجا پاسا خالی \pnbLatin{(empty right side)} ہووے)۔
قاعدہ \RightR{\Weakening} \pnbLatin{(right weakening)} ورت کے، کسے متضاد سیٹ توں ہر !!{sentence} نوں
!!{derive} کیتا جا سکدا اے۔

سیکوئنٹ حساب نوں 1930 دی دہائی وچ Gerhard Gentzen نے ایجاد کیتا۔
اپنی منظم تے متناظر بناوٹ \pnbLatin{(systematic and symmetric design)} پاروں ایہ
!!{derivation}s دا نظریہ تیار کرن لئی بہت فائدہ مند رسمی ڈھانچا
(formalism) اے۔ سیکوئنٹ حساب وچ !!{derivation}s لبھنے
نسبتاً سَوکھے نیں، پر ایہ !!{derivation}s اکثر پڑھن وچ اوکھے ہُندے نیں
تے ثبوتاں نال اوہناں دا تعلق ویکھنا کدی کدی سَوکھا نہیں ہُندا۔ ایس دے باوجود
ایہ !!{derivation} نظاماں دا بہت سلیقے والا طریقہ ثابت ہویا اے، تے بہت
ساریاں منطقاں دے سیکوئنٹ حسابی نظام نیں۔
\endgroup
% /OLP-0065

% OLP-0066 content/first-order-logic/proof-systems/natural-deduction.tex
\begingroup
\def\olfilename{natural-deduction}
\expandafter\def\csname ol@sectioncs\endcsname{section}

\iftag{FOL}
      {\olfileid{fol}{prf}{ntd}}
      {\olfileid{pl}{prf}{ntd}}

\olsection{فطری استنباط \pnbLatin{(Natural Deduction)}}

فطری استنباط \pnbLatin{(natural deduction)} اک !!a{derivation} نظام اے جہدا مقصد
اصل استدلال (actual reasoning، خاص طور تے ریاضی داناں دے ورتے ہوئے
باقاعدہ استدلال) دی صورت نوں پیش کرنا اے۔
اصل استدلال کئی ``فطری'' نمونیاں راہیں اگے ودھدا اے۔ مثال لئی، صورتاں
راہیں ثبوت \pnbLatin{(proof by cases)} سانوں اک فصلی مقدمے \pnbLatin{(disjunctive premise)}
دی بنیاد اُتے کوئی نتیجہ قائم کرن دی
اجازت دیندا اے، ایس گل نوں قائم کر کے کہ اوہ نتیجہ ہر اک فصلی جز توں
نکلدا اے۔ بالواسطہ ثبوت \pnbLatin{(indirect proof)} سانوں کوئی نتیجہ ایس گل نال
قائم کرن دی اجازت
دیندا اے کہ اوہدی نفی اک تضاد تک لے جاندی اے۔ شرطیہ ثبوت ``جے \dots
تاں \dots'' شرطیہ دعوا \pnbLatin{(conditional claim)} ایس گل نال قائم کردا اے کہ
تالی (consequent)، مقدم (antecedent) توں نکلدا اے۔ فطری استنباط ایہناں
فطری استنباطاں وچوں کجھ دی رسمی تشکیل اے۔ ہر منطقی رابطے تے مسوّر دے
نال دو قاعدے، اک تعارفی قاعدہ \pnbLatin{(introduction rule)} تے اک اخراجی قاعدہ
\pnbLatin{(elimination rule)}، آؤندے
نیں تے ہر قاعدہ ایہو جہے کسے فطری استنباطی نمونے دے مطابق ہُندا اے۔
مثال لئی، $\Intro{\lif}$ شرطیہ ثبوت دے مطابق اے تے $\Elim{\lor}$
صورتاں راہیں ثبوت دے۔ اک خاص طور تے سادہ قاعدہ $\Elim{\land}$ اے،
جیہڑا $!A \land !B$ توں~$!A$ (یا $!B$) دا استنباط کرن دی اجازت دیندا
اے۔

فطری استنباط نوں ہور !!{derivation} نظاماں توں وکھ کرن والی اک خاصیت
اوس ولوں مفروضیاں (assumptions) دی ورتوں اے۔ فطری استنباط وچ
!!^a{derivation}،
!!{formula}s دا اک درخت ہُندا اے۔ اک !!{formula}، !!{formula}s دے درخت
دی جڑ (root) اُتے کھلوتا ہُندا اے، تے درخت دے ``پتے'' (leaves) اوہ
!!{formula}s ہُندے
نیں جنہاں توں نتیجہ اخذ کیتا جاندا اے۔ فطری استنباط وچ کجھ پتے والے
!!{formula}s، !!{derivation} دے اندر کردار نبھاؤندے نیں پر جدوں
!!{derivation} نتیجے تک پہنچدا اے تاں اوہ ``ورت کے مکائے'' جا چکے ہُندے
نیں۔ ایہ اصل استدلال وچ ایہو جہے مفروضے متعارف کراؤن دے عمل دے مطابق
اے جیہڑے تھوڑھے چِر لئی ای لاگو رہندے نیں۔ مثال لئی، صورتاں راہیں ثبوت
وچ اسیں ہر فصلی جز نوں سچ منّ لیندے آں؛ شرطیہ ثبوت وچ اسیں مقدم نوں
سچ منّ لیندے آں؛ بالواسطہ ثبوت وچ اسیں نتیجے دی نفی نوں سچ منّ لیندے
آں۔ وقتی مفروضے متعارف کراؤن تے فیر درمیانی قدم قائم کرن لئی اوہناں
نوں ختم کر دین دا ایہ طریقہ فطری استنباط دی نمایاں خاصیت اے۔ فطری
استنباطی !!{derivation} دے پتیاں اُتے فارمولیاں نوں مفروضے آکھیا جاندا
اے، تے استنباط دے کجھ قاعدے اوہناں نوں ``!!{discharge}'' (discharge)
کر سکدے نیں۔
مثال لئی، جے ساڈے کول کجھ مفروضیاں توں~$!B$ دا !!a{derivation} اے تے
اوہناں مفروضیاں وچ~$!A$ وی شامل اے، تاں $\Intro{\lif}$ دا قاعدہ سانوں
~$!A \lif !B$ دا استنباط کرن تے~$!A$ دے روپ والے ہر مفروضے نوں
خارج کرن دی اجازت دیندا اے۔ (ایہ حساب رکھن لئی کہ کیہڑے
استنباط اُتے کیہڑے مفروضے خارج کیتے جاندے نیں، اسیں استنباط
تے اوس ولوں خارج کیتے مفروضیاں اُتے اک نمبر دا نشان لاؤندے
آں۔) جیہڑے مفروضے !!{derivation} دے آخر اُتے !!{undischarged}
(undischarged) رہندے
نیں، اوہ رل کے نتیجے دے صدق لئی کافی ہُندے نیں، ایس لئی
!!a{derivation} قائم کردا اے کہ اوہدے !!{undischarged} مفروضیاں توں
اوہدا نتیجہ منطقی طور تے نکلدا اے۔

فطری استنباط اُتے مبنی اخذ ہون جوگتا دا تعلق \pnbLatin{(derivability relation)}
$\Gamma \Proves !A$ اُدوں تے صرف اُدوں
قائم ہُندا اے جدوں کوئی !!a{derivation} ایہو جہا ہووے کہ $!A$ درخت دا
آخری !!{sentence} ہووے تے ہر !!{undischarged} پتا~$\Gamma$ وچ ہووے۔
$!A$ فطری استنباط وچ اُدوں تے صرف اُدوں اک قضیہ (theorem) اے جدوں کوئی
!!a{derivation} ایہو جہا ہووے کہ $!A$ آخری !!{sentence} ہووے تے سارے
مفروضے !!{discharged} ہون۔ مثال لئی، ایہ !!a{derivation} وکھاؤندا اے
کہ $\Proves (!A
\land !B) \lif !A$:
\begin{prooftree}
  \AxiomC{$\Discharge{!A \land !B}{1}$}
  \RightLabel{\Elim{\land}}
  \UnaryInfC{$!A$}
  \DischargeRule{\Intro{\lif}}{1}
  \UnaryInfC{$(!A \land !B) \lif !A$}
\end{prooftree}
نشان~$1$ دسدا اے کہ مفروضہ $!A \land !B$ نوں \Intro{\lif}
استنباط اُتے !!{discharged} (discharged) کیتا جاندا اے۔

فطری استنباط وچ سیٹ~$\Gamma$ اُدوں تے صرف اُدوں متضاد (inconsistent)
اے جدوں $\Gamma \Proves \lfalse$۔ کذب دا قاعدہ \pnbLatin{(falsity rule)}
\FalseInt{} ایس گل نوں یقینی بناؤندا
اے کہ کسے متضاد سیٹ توں ہر !!{sentence} نوں !!{derive} کیتا جا سکدا اے۔

فطری استنباطی نظام Gerhard Gentzen تے Stanis\l{}aw Ja\'skowski ولوں
1930 دی دہائی وچ بنائے گئے، تے بعد وچ Dag Prawitz تے Frederic Fitch
ولوں ہور اگے ودھائے گئے۔ چونکہ ایہدے استنباط ثبوت دے فطری طریقیاں دی
صورت نوں پیش کردے نیں، ایس لئی فلسفی ایہنوں پسند کردے نیں۔ Fitch ولوں
تیار کیتے روپ اکثر منطق دیاں تعارفی درسی کتاباں وچ ورتے جاندے نیں۔
منطق دے فلسفے وچ، فطری استنباط دے قاعدیاں نوں کدی کدی منطقی رابطیاں دے
معنی دین والا منّیا گیا اے (``ثبوتی-نظریاتی معنیات
\pnbLatin{(proof-theoretic semantics)}'')۔
\endgroup
% /OLP-0066

% OLP-0067 content/first-order-logic/proof-systems/tableaux.tex
\begingroup
\def\olfilename{tableaux}
\expandafter\def\csname ol@sectioncs\endcsname{section}

\iftag{FOL}
      {\olfileid{fol}{prf}{tab}}
      {\olfileid{pl}{prf}{tab}}

\olsection{\usetoken{P}{tableau}}

بھاویں بہت سارے !!{derivation} نظام !!{sentence}s دے بندوبستاں نال کم
کردے نیں، !!{tableau}s (tableaux) !!{signed formula}s \pnbLatin{(signed formulas)}
نال کم کردے نیں۔ !!^a{signed formula} اک جوڑی اے جیہڑی صداقتی قدر دی
اک علامت \pnbLatin{(truth-value sign)}
($\True$ یا $\False$) تے !!a{sentence} اُتے مشتمل ہُندی اے
\[
\sFmla{\True}{!A} \text{ یا } \sFmla{\False}{!A}.
\]
!!^a{tableau} وچ !!{signed formula}s نوں ہِٹھاں ول شاخاں کڈھن والے
درخت \pnbLatin{(downward-branching tree)} دے روپ وچ ترتیب دتا جاندا اے۔ ایہ کجھ
\emph{مفروضیاں (assumptions)} نال شروع
ہُندا اے تے اوہناں !!{signed formula}s نال اگے ودھدا اے جیہڑے استنباط
دے قاعدیاں \pnbLatin{(rules of inference)} وچوں کوئی اک لاگو کر کے اپنے توں اُتّے والے
!!{signed formula}s وچوں کسے اک توں حاصل ہُندے نیں۔ ہر قاعدہ سانوں کسے شاخ دے سرے
اُتے اک یا ودھ !!{signed formula}s، یا نال نال دو !!{signed formula}s
شامل کرن دی اجازت دیندا اے---ایس دوجی صورت وچ شاخ دو حصیاں وچ ونڈی
\pnbLatin{(splits into two)}
جاندی اے تے شامل کیتے دو !!{signed formula}s اوہناں دو شاخاں دے سرے بن جاندے نیں۔

کسی مرکب !!{signed formula} اُتے لاگو کیتا قاعدہ اوہ !!{signed formula}s
شامل کردا اے جیہڑے عین ذیلی !!{formula}s \pnbLatin{(immediate subformulas)} ہُندے
نیں۔ ایہ
قاعدے جوڑیاں وچ آؤندے نیں، دوہاں علامتاں وچوں ہر اک لئی اک قاعدہ۔
مثال لئی، $\TRule{\True}{\land}$ دا قاعدہ
$\sFmla{\True}{!A \land !B}$ اُتے لاگو ہُندا اے، تے ایہ
دونیں !!{signed formula}s $\sFmla{\True}{!A}$ تے~$\sFmla{\True}{!B}$
نوں $\sFmla{\True}{!A \land !B}$ والی کسے وی شاخ دے سرے اُتے شامل
کرن دی اجازت دیندا اے؛ تے $\TRule{\False}{!A \land !B}$ دا قاعدہ
$\sFmla{\False}{!A}$ تے $\sFmla{\False}{!B}$ نوں نال نال شامل کر کے
شاخ نوں دو حصیاں وچ ونڈن دی اجازت دیندا اے۔ !!^a{tableau} بند (closed)
اے جے اوہدی ہر اک شاخ وچ !!{signed formula}s دی اک میل کھاندی جوڑی
\pnbLatin{(matching pair)}
$\sFmla{\True}{!A}$ تے $\sFmla{\False}{!A}$ شامل ہووے۔

!!{tableau}s اُتے مبنی $\Proves$ تعلق \pnbLatin{(tableau derivability)} دی تعریف
ایویں کیتی جاندی اے: $\Gamma \Proves !A$ اُدوں تے صرف اُدوں جدوں کوئی
محدود سیٹ \pnbLatin{(finite set)}~$\Gamma_0 = \{!B_1,
\dots, !B_n\} \subseteq \Gamma$ ایہو جہا ہووے کہ مفروضیاں
\[
\{\sFmla{\False}{!A}, \sFmla{\True}{!B_1}, \dots, \sFmla{\True}{!B_n}\}
\]
لئی کوئی بند !!{tableau} \pnbLatin{(closed tableau)} ہووے۔ مثال لئی، ایہ بند
!!{tableau} وکھاؤندا
اے کہ $\Proves (!A \land !B) \lif !A$:
\begin{oltableau}
  [\sFmla{\False}{(\formula{A} \land \formula{B}) \lif \formula{A}}, just = \TAss
    [\sFmla{\True}{\formula{A} \land \formula{B}}, just = {\TRule{\False}{\lif}[1]}
      [\sFmla{\False}{\formula{A}}, just = {\TRule{\False}{\lif}[1]}
        [\sFmla{\True}{\formula{A}}, just = {\TRule{\True}{\lif}[2]}
          [\sFmla{\True}{\formula{B}}, just = {\TRule{\True}{\lif}[2]}, close]
        ]
      ]
    ]
  ]
\end{oltableau}

!!{tableau} حساب \pnbLatin{(tableau calculus)} وچ سیٹ $\Gamma$ متضاد
(inconsistent) اے جے مفروضیاں
\[
\{\sFmla{\True}{!B_1}, \dots, \sFmla{\True}{!B_n}\}
\]
لئی، کسے $!B_i \in \Gamma$ واسطے، کوئی بند !!{tableau} ہووے۔

!!^{tableau}s نوں 1950 دی دہائی وچ Evert Beth تے Jaakko Hintikka
نے اک دوجے توں آزاد ہو کے ایجاد کیتا، تے Raymond Smullyan نے اوہناں
نوں سادہ بنا کے عام کیتا۔ ایہ ورتن وچ بہت سَوکھے نیں، کیوں جو
!!a{tableau} بنانا بہت منظم کارروائی \pnbLatin{(systematic procedure)} اے۔
!!{tableau}s دی منظم نوعیت پاروں ایہناں نوں کمپیوٹر راہیں لاگو
\pnbLatin{(implementation by computer)} کرنا وی سَوکھا اے۔ پر
!!a{tableau} اکثر پڑھن وچ اوکھا ہُندا اے تے اوہناں دا ثبوتاں نال تعلق
ویکھنا کدی کدی سَوکھا نہیں ہُندا۔ ایہ طریقہ کافی عام وی اے تے بہت
ساریاں منطقاں دے !!{tableau} نظام نیں۔ !!^{tableau}s سانوں اوہ
!!{structure}s لبھن وچ وی مدد دیندے نیں جیہڑے دتے ہوئے !!{sentence}s
(یا اوہناں دے سیٹاں) نوں پورا کردے نیں: جے سیٹ پورا ہون جوگ اے تاں
اوہدے لئی کوئی بند !!{tableau} نہیں ہووے گا، یعنی ہر !!{tableau} وچ
اک کھلی شاخ \pnbLatin{(open branch)} ہووے گی۔ پورا کرن والی !!{structure} نوں کسے
کھلی شاخ توں ``پڑھ کے کڈھیا'' \pnbLatin{(read off)} جا سکدا اے، بشرطیکہ اوس شاخ اُتے ہر اوہ قاعدہ
لاگو کیتا جا چکیا ہووے جیہڑا لاگو ہو سکدا اے۔ سیکوئنٹ حساب نال وی
ایس دا بہت نیڑا تعلق اے: بنیادی طور تے، بند !!{tableau} سیکوئنٹ حساب
\pnbLatin{(sequent calculus)} دا اک مختصر کیتا !!{derivation} اے جیہڑا اُلٹا لکھیا
\pnbLatin{(written upside-down)} ہُندا اے۔
\endgroup
% /OLP-0067

% OLP-0068 content/first-order-logic/proof-systems/axiomatic-deduction.tex
\begingroup
\def\olfilename{axiomatic-deduction}
\expandafter\def\csname ol@sectioncs\endcsname{section}

\iftag{FOL}
      {\olfileid{fol}{prf}{axd}}
      {\olfileid{pl}{prf}{axd}}

\olsection{مسلّمی \usetoken{P}{derivation} \pnbLatin{(Axiomatic Derivation)}}

مسلّمی !!{derivation}s \pnbLatin{(axiomatic derivations)} سب توں پرانے تے سب توں سادہ منطقی
!!{derivation} نظام نیں۔ ایہناں دے !!{derivation}s بس !!{sentence}s دے
سلسلے ہُندے نیں۔ !!{sentence}s دا کوئی سلسلہ اک درست !!{derivation}
اُدوں منّیا جاندا اے جدوں اوہدے وچ ہر !!{sentence}~$!A$ ہِٹھ لکھیاں
شرطاں وچوں کوئی اک پوری کرے:
\begin{enumerate}
\item $!A$ اک مسلّمہ (axiom) اے، یا
\item $!A$، !!{sentence}s دے دتے ہوئے سیٹ~$\Gamma$ دا !!a{element} اے، یا
\item $!A$ نوں استنباط دے کسے قاعدے \pnbLatin{(rule of inference)} توں جواز ملدا اے۔
\end{enumerate}
مسلّمہ ہون لئی $!A$ نوں کجھ مقررہ !!{sentence} خاکیاں \pnbLatin{(sentence schemas)} وچوں کسے اک دا روپ رکھنا پیندا اے۔ مسلّمہ خاکیاں \pnbLatin{(axiom schemas)} دے بہت سارے ایہو جہے سیٹ نیں جیہڑے پہلے درجے دی منطق لئی اک
تسلی بخش (صحت مند تے مکمل؛ sound and complete) !!{derivation} نظام
دیندے نیں۔ کجھ نوں اوہناں رابطیاں دے لحاظ نال ترتیب دتا جاندا اے
جنہاں اُتے اوہ لاگو ہُندے نیں؛ مثال لئی، خاکے
\[
!A \lif (!B \lif !A) \qquad !B \lif (!B \lor !C) \qquad (!B \land !C) \lif !B
\]
عام مسلّمے نیں جیہڑے $\lif$، $\lor$ تے~$\land$ اُتے لاگو ہُندے نیں۔
کجھ مسلّمی نظاماں دا مقصد مسلّماں دی گنتی گھٹ توں گھٹ رکھنا ہُندا
اے۔ ایہ گل اوس اُتے منحصر اے کہ کیہڑے رابطیاں نوں ابتدائی (primitive)
منّیا جاندا اے؛ ایسے مسلّمی نظام لبھنے وی ممکن نیں جیہڑے صرف اک
مسلّمے \pnbLatin{(single axiom)} اُتے مشتمل ہون۔

استنباط دا قاعدہ \pnbLatin{(rule of inference)} اک شرطیہ عبارت اے جیہڑی
!!a{derivation} وچ !!a{sentence} نوں جواز ملن لئی کافی شرط \pnbLatin{(sufficient condition)} دیندی اے۔ قاعدۂ اثباتِ مقدم \pnbLatin{(modus ponens)}
ایہو جہا اک بہت عام قاعدہ اے: ایہ آکھدا اے کہ جے $!A$ تے $!A \lif
!B$ نوں پہلاں ای جواز مل چکیا اے، تاں $!B$ نوں وی جواز ملدا اے۔ ایس
دا مطلب اے کہ !!a{derivation} وچ !!{sentence}~$!B$ والی سطر نوں جواز
ملدا اے، بشرطیکہ $!A$ تے $!A \lif !B$ دونیں (کسے
!!{sentence}~$!A$ لئی) اوس !!{derivation} وچ~$!B$ توں پہلاں (earlier)
آ چکے ہون۔

مسلّمی !!{derivation}s اُتے مبنی $\Proves$ تعلق \pnbLatin{(axiomatic derivability)} دی تعریف ایویں کیتی جاندی اے: $\Gamma \Proves !A$ اُدوں
تے صرف اُدوں جدوں کوئی
!!a{derivation} ایہو جہا ہووے کہ !!{sentence}~$!A$ اوہدا آخری فارمولا
\pnbLatin{(last formula)}
ہووے (تے $\Gamma$ نوں اوس !!{derivation} دے اوہناں !!{sentence}s دا
سیٹ منّیا جاوے جنہاں نوں اُتّے دتی شرط~(2) توں جواز ملدا اے)۔ $!A$
اک قضیہ (theorem) اے جے $!A$ دا کوئی ایہو جہا !!a{derivation} ہووے جتھے
$\Gamma$ خالی ہووے، یعنی !!{derivation} وچ ہر !!{sentence} نوں یا
(1) توں یا~(3) توں جواز ملدا ہووے۔ مثال لئی، ایہ !!a{derivation}
وکھاؤندا اے کہ $\Proves !A  \lif (!B \lif (!B \lor !A))$:
\begin{derivation}
  1. & $!B \lif (!B \lor !A)$ \\
  2. & $(!B \lif (!B \lor !A)) \lif (!A  \lif (!B \lif (!B \lor !A)))$\\
  3. & $!A  \lif (!B \lif (!B \lor !A))$
\end{derivation}
سطر~1 دا !!{sentence} مسلّمے $!A \lif (!A \lor !B)$ دے روپ دا اے
($!A$ تے $!B$ دے کردار الٹے کر کے)۔ سطر~2 دا جملہ مسلّمے $!A \lif
(!B \lif !A)$ دے روپ دا اے۔ ایس لئی دوناں سطراں نوں جواز ملدا اے۔
سطر~3 نوں قاعدۂ اثباتِ مقدم توں جواز ملدا اے: جے اسیں ایہنوں مختصر
کر کے $!D$ لکھیے، تاں سطر~2 دا روپ $!C \lif !D$ اے، جتھے $!C$،
$!B \lif (!B \lor !A)$ یعنی سطر~1 اے۔

سیٹ $\Gamma$ متضاد (inconsistent) اے جے $\Gamma \Proves \lfalse$۔ اک
مکمل مسلّمی نظام \pnbLatin{(complete axiom system)} $\lfalse \lif !A$ نوں وی،
ہر~$!A$ لئی، ثابت کرے گا، ایس لئی جے
$\Gamma$ متضاد اے تاں $\Gamma \Proves !A$، ہر~$!A$ لئی۔

منطق لئی مسلّمی !!{derivation}s نظام سب توں پہلاں Gottlob Frege نے
اپنی 1879 دی \emph{Begriffsschrift} وچ دتے سن؛ ایس کارن اوہ کم اکثر
جدید منطق دا پہلا کم منّیا جاندا اے۔ Alfred North Whitehead تے
Bertrand Russell دی \emph{Principia Mathematica} وچ، تے 1920 دی دہائی
وچ David Hilbert تے اوہدے شاگرداں ولوں، ایہناں نظاماں نوں کمال تک
پہنچایا گیا۔ ایس لئی ایہناں نوں اکثر ``فریگ نظام \pnbLatin{(Frege systems)}'' یا
``ہلبرٹ نظام \pnbLatin{(Hilbert systems)}'' آکھیا جاندا اے۔ ایہ بہت لچکدار
(versatile) نیں، ایس معنی وچ کہ کسے منطق لئی
مسلّمی نظام لبھنا اکثر سَوکھا ہُندا اے۔ چونکہ !!{derivation}s دی بناوٹ
بہت سادہ اے تے اوہناں وچ استنباط دے صرف اک یا دو قاعدے ہُندے نیں، ایس
لئی اوہناں دے \emph{بارے} گلاں ثابت کرنا \pnbLatin{(metatheoretical reasoning)}
وی نسبتاً سَوکھا اے۔ پر عملی ورتوں \pnbLatin{(in practice)} وچ ایہ بہت اوکھے نیں،
یعنی ثبوت لبھنے تے لکھنے اوکھے نیں۔
\endgroup
% /OLP-0068

\OLEndChapterHook
\endgroup
% /OLP-0063

% OLP-0069 content/first-order-logic/sequent-calculus/sequent-calculus.tex
\begingroup
\def\olfilename{sequent-calculus}
\expandafter\def\csname ol@sectioncs\endcsname{section}

\iftag{FOL}
      {\olchapter{fol}{seq}{سیکوئنٹ حساب \pnbLatin{(The Sequent Calculus)}}}
      {\olchapter{pl}{seq}{سیکوئنٹ حساب \pnbLatin{(The Sequent Calculus)}}}

\begin{editorial}
  ایہ باب کلاسیکی پہلے درجے دی منطق \pnbLatin{(classical first-order logic)} لئی
  Gentzen دا معیاری سیکوئنٹ حساب \pnbLatin{(standard sequent calculus)} LK پیش
  کردا اے۔ ایہنوں ہور مثالاں تے مشقاں دی لوڑ اے۔ سیکوئنٹ حساب نوں اک
  ثبوتی نظام \pnbLatin{(proof system)} وجوں متعلقہ مواد شامل یا خارج کرن لئی
  ``prfLK'' ٹیگ ورتو۔
\end{editorial}

% OLP-0070 content/first-order-logic/sequent-calculus/rules-and-proofs.tex
\begingroup
\def\olfilename{rules-and-proofs}
\expandafter\def\csname ol@sectioncs\endcsname{section}

\iftag{FOL}
      {\olfileid{fol}{seq}{rul}}
      {\olfileid{pl}{seq}{rul}}

\olsection{قاعدے تے \usetoken{P}{derivation} \pnbLatin{(Rules and Derivations)}}

اگلی بحث لئی، منّ لو کہ $\Gamma, \Delta, \Pi, \Lambda$، !!{sentence}s
دے محدود سلسلے \pnbLatin{(finite sequences)} ظاہر کردے نیں۔

\begin{defn}[سیکوئنٹ (Sequent)]
\emph{سیکوئنٹ (sequent)} ایس روپ دی اک عبارت اے
\[
\Gamma \Sequent \Delta
\]
جتھے $\Gamma$ تے $\Delta$، زبان $\Lang L$ دے !!{sentence}s دے محدود
(ممکنہ طور تے خالی؛ possibly empty) سلسلے نیں۔ $\Gamma$ نوں
\emph{مقدم (antecedent)} آکھیا جاندا اے، جد کہ $\Delta$ نوں
\emph{تالی (succedent)} آکھیا جاندا اے۔
\end{defn}

\begin{explain}
سیکوئنٹ دے پچھے بدیہی خیال \pnbLatin{(intuitive idea)} ایہ اے: جے مقدم وچ سارے
!!{sentence}s قائم ہون تاں تالی دے !!{sentence}s وچوں گھٹ توں گھٹ اک
\pnbLatin{(at least one)} قائم ہووے۔ یعنی، جے
$\Gamma = \tuple{!A_1, \dots, !A_m}$ تے
$\Delta = \tuple{!B_1, \dots, !B_n}$، تاں $\Gamma \Sequent \Delta$
اُدوں تے صرف اُدوں قائم ہُندا اے جدوں
\[
(!A_1 \land \cdots \land !A_m) \lif (!B_1 \lor \cdots \lor
!B_n)
\]
قائم ہُندا اے۔ دو خاص صورتاں \pnbLatin{(special cases)} نیں: اک جتھے $\Gamma$ خالی اے تے دوجی
جدوں $\Delta$~خالی اے۔ جدوں $\Gamma$ خالی اے، یعنی $m = 0$، $\quad
\Sequent \Delta$ اُدوں تے صرف اُدوں قائم ہُندا اے جدوں $!B_1 \lor
\dots \lor !B_n$ قائم ہُندا اے۔ جدوں $\Delta$ خالی اے، یعنی $n = 0$،
$\Gamma \Sequent \quad$ اُدوں تے صرف اُدوں قائم ہُندا اے جدوں
$\lnot(!A_1 \land \dots \land !A_m)$ قائم ہُندا اے۔ اسیں آکھدے آں
کہ کوئی سیکوئنٹ اُدوں تے صرف اُدوں معتبر (valid) اے جدوں اوہدے مطابق
!!{sentence} \pnbLatin{(corresponding sentence)} معتبر اے۔
\end{explain}

جے $\Gamma$، !!{sentence}s دا اک سلسلہ اے، تاں اسیں $\Gamma, !A$ اوس
نتیجے نوں لکھدے آں جیہڑا $!A$ نوں~$\Gamma$ دے سجے سرے اُتے شامل
\pnbLatin{(append to the right end)} کرن
نال ملدا اے (تے $!A, \Gamma$ اوس نتیجے نوں جیہڑا $!A$ نوں~$\Gamma$
دے کھبے سرے اُتے شامل \pnbLatin{(append to the left end)} کرن نال ملدا اے)۔ جے $\Delta$ وی
!!{sentence}s دا اک سلسلہ اے، تاں $\Gamma, \Delta$ ایس لکھی ترتیب نال
دونیں سلسلیاں نوں جوڑ (concatenation) کے بنیا سلسلہ اے۔

\begin{defn}[ابتدائی سیکوئنٹ \pnbLatin{(Initial Sequent)}]
\emph{ابتدائی سیکوئنٹ \pnbLatin{(initial sequent)}} اک ایہو جہا سیکوئنٹ اے
\iftag{prvFalse,prvTrue}{جیہڑا اگلے روپاں وچوں کسے اک دا ہووے:
  \begin{enumerate}
    \item $!A \Sequent !A$
    \tagitem{prvTrue}{$\quad \Sequent \ltrue$}{}
    \tagitem{prvFalse}{$\lfalse \Sequent \quad$}{}
  \end{enumerate}}
{جیہڑا $!A \Sequent !A$ دے روپ دا ہووے}، زبان دے کسے وی
!!{sentence} $!A$ لئی۔
\end{defn}

سیکوئنٹ حساب وچ !!^{derivation}s سیکوئنٹاں دے کجھ خاص درخت \pnbLatin{(trees of sequents)} نیں، جتھے سب توں اُتّے والے سیکوئنٹ \pnbLatin{(topmost sequents)}
ابتدائی سیکوئنٹ ہُندے نیں، تے جے کوئی
سیکوئنٹ اک یا دو ہور سیکوئنٹاں دے تھلّے آوے تاں اوہ استنباط دے کسے
قاعدے \pnbLatin{(rule of inference)} راہیں درست طور تے اوہناں توں نکلنا چاہیدا اے۔ $\Log{LK}$ دے
قاعدیاں دیاں دو وڈیاں قسماں نیں:
\emph{منطقی (logical)} قاعدے تے \emph{ساختی (structural)} قاعدے۔ منطقی قاعدیاں دے ناں ہِٹھلے
سیکوئنٹ وچ $!A$ تے/یا $!B$ والے !!{sentence} دے !!{main operator} \pnbLatin{(main operator)} دے ناں اُتے رکھے جاندے نیں۔ ہر اک دے دو روپ ہُندے نیں:
اک اوس سیکوئنٹ دا استنباط کرن لئی جیہدے کھبے پاسے !!{operator} والا
!!{sentence} ہووے، تے اک اوس لئی جیہدے سجے پاسے ایہ !!{sentence} ہووے۔
\endgroup
% /OLP-0070

% OLP-0071 content/first-order-logic/sequent-calculus/propositional-rules.tex
\begingroup
\def\olfilename{propositional-rules}
\expandafter\def\csname ol@sectioncs\endcsname{section}

\iftag{FOL}
      {\olfileid{fol}{seq}{prl}}
      {\olfileid{pl}{seq}{prl}}

\olsection{قضیاتی قاعدے \pnbLatin{(Propositional Rules)}}

\subsection{$\lnot$ دے قاعدے \pnbLatin{(negation rules)}}

\begin{defish}
\Axiom$ \Gamma \fCenter \Delta, !A $
\RightLabel{\LeftR{\lnot}}
\UnaryInf$ \lnot !A, \Gamma \fCenter \Delta$
\DisplayProof
\hfill
\Axiom$!A, \Gamma \fCenter \Delta$
\RightLabel{\RightR{\lnot}}
\UnaryInf$ \Gamma \fCenter \Delta, \lnot !A $
\DisplayProof
\end{defish}

\subsection{$\land$ دے قاعدے \pnbLatin{(conjunction rules)}}

\begin{defish}\noindent
\begin{tabular}{l}
\Axiom$ !A, \Gamma \fCenter \Delta$
\RightLabel{\LeftR{\land}}
\UnaryInf$ !A \land !B, \Gamma \fCenter \Delta$
\DisplayProof
\\[3ex]
\Axiom$!B, \Gamma \fCenter \Delta$
\RightLabel{\LeftR{\land}}
\UnaryInf$!A \land !B, \Gamma \fCenter \Delta$
\DisplayProof
\end{tabular}
\hfill
\Axiom$\Gamma \fCenter \Delta, !A$
\Axiom$ \Gamma \fCenter \Delta, !B$
\RightLabel{\RightR{\land}}
\BinaryInf$ \Gamma \fCenter \Delta, !A \land !B $
\DisplayProof
\end{defish}

\subsection{$\lor$ دے قاعدے \pnbLatin{(disjunction rules)}}

\begin{defish}
\Axiom$!A, \Gamma \fCenter \Delta$
\Axiom$!B, \Gamma \fCenter \Delta$
\RightLabel{\LeftR{\lor}}
\BinaryInf$!A \lor !B, \Gamma \fCenter \Delta$
\DisplayProof
\hfill
\begin{tabular}{r}
\Axiom$\Gamma \fCenter \Delta, !A$
\RightLabel{\RightR{\lor}}
\UnaryInf$ \Gamma \fCenter \Delta, !A \lor !B$
\DisplayProof
\\[3ex]
\Axiom$ \Gamma \fCenter \Delta, !B$
\RightLabel{\RightR{\lor}}
\UnaryInf$ \Gamma \fCenter \Delta, !A \lor !B$
\DisplayProof
\end{tabular}
\end{defish}

\subsection{$\lif$ دے قاعدے \pnbLatin{(conditional rules)}}

\begin{defish}
\Axiom$ \Gamma \fCenter \Delta, !A$
\Axiom$ !B, \Pi \fCenter \Lambda$
\RightLabel{\LeftR{\lif}}
\BinaryInf$ !A \lif !B, \Gamma, \Pi \fCenter \Delta, \Lambda$
\DisplayProof
\hfill
\Axiom$ !A, \Gamma \fCenter \Delta, !B$
\RightLabel{\RightR{\lif}}
\UnaryInf$ \Gamma \fCenter \Delta, !A \lif !B $
\DisplayProof
\end{defish}
\endgroup
% /OLP-0071

% OLP-0072 content/first-order-logic/sequent-calculus/quantifier-rules.tex
\begingroup
\def\olfilename{quantifier-rules}
\expandafter\def\csname ol@sectioncs\endcsname{section}

\olfileid{fol}{seq}{qrl}

\olsection{مسوّراں دے قاعدے \pnbLatin{(Quantifier Rules)}}

\subsection{$\lforall$ دے قاعدے \pnbLatin{(universal-quantifier rules)}}

\begin{defish}
\Axiom$ !A(t), \Gamma \fCenter \Delta$
\RightLabel{\LeftR{\lforall}}
\UnaryInf$ \lforall[x][!A(x)],\Gamma \fCenter \Delta$
\DisplayProof
\hfill
\Axiom$ \Gamma \fCenter \Delta, !A(a) $
\RightLabel{\RightR{\lforall}}
\UnaryInf$ \Gamma \fCenter \Delta, \lforall[x][!A(x)]$
\DisplayProof
\end{defish}

\LeftR{\lforall} وچ، $t$ اک بند ٹرم \pnbLatin{(closed term)} اے (یعنی، جیہدے وچ متغیر نہ
ہون)۔ \RightR{\lforall} وچ، $a$ اک !!a{constant} اے جیہڑا
\RightR{\lforall} قاعدے دے ہِٹھلے سیکوئنٹ \pnbLatin{(lower sequent)} وچ کتے وی نہیں آنا چاہیدا۔
اسیں $a$ نوں \RightR{\forall} استنباط دا \emph{ایگن متغیر (eigenvariable)} آکھدے
آں۔\footnote{اسیں ``ایگن متغیر (eigenvariable)'' اصطلاح ورتدے آں، بھاویں اُپرلے قاعدے
وچ $a$ اک !!a{constant} (constant) اے۔ ایس دے تاریخی سبب \pnbLatin{(historical reasons)} نیں۔}

\subsection{$\lexists$ دے قاعدے \pnbLatin{(existential-quantifier rules)}}

\begin{defish}
\Axiom$ !A(a), \Gamma \fCenter \Delta $
\RightLabel{\LeftR{\lexists}}
\UnaryInf$ \lexists[x][!A(x)], \Gamma \fCenter \Delta$
\DisplayProof
\hfill
\Axiom$ \Gamma \fCenter \Delta, !A(t) $
\RightLabel{\RightR{\lexists}}
\UnaryInf$ \Gamma \fCenter \Delta, \lexists[x][!A(x)]$
\DisplayProof
\end{defish}

ایتھے فیر، $t$ اک بند ٹرم \pnbLatin{(closed term)} اے، تے $a$ اک !!a{constant} (constant) اے جیہڑا
\LeftR{\lexists} قاعدے دے ہِٹھلے سیکوئنٹ \pnbLatin{(lower sequent)} وچ نہیں آؤندا۔ اسیں $a$ نوں
\LeftR{\lexists} استنباط دا \emph{ایگن متغیر (eigenvariable)} آکھدے آں۔

ایہ شرط کہ کوئی ایگن متغیر \RightR{\lforall} یا
\LeftR{\lexists} استنباط دے ہِٹھلے سیکوئنٹ وچ نہ آوے، \emph{ایگن
متغیر دی شرط \pnbLatin{(eigenvariable condition)}} آکھوَندی اے۔

\begin{explain}
یاد کرو اوہ رواج کہ جدوں $!A$ اک ایہو جہا !!a{formula} اے جیہدے وچ
!!{variable}~$x$ آزاد \pnbLatin{(free variable)} اے، تاں اسیں ایہ گل~$!A(x)$ لکھ کے ظاہر کردے
آں۔ ایسے سیاق وچ، $!A(t)$ فیر~$\Subst{!A}{t}{x}$ دا مختصر روپ (shorthand) اے۔ سو
اسیں \RightR{\lexists} دا قاعدہ ایویں وی لکھ سکدے آں:
\begin{prooftree}
  \Axiom$\Gamma \fCenter \Delta, \Subst{!A}{t}{x}$
  \RightLabel{\RightR{\lexists}}
  \UnaryInf$\Gamma \fCenter \Delta, \lexists[x][!A]$
\end{prooftree}
غور کرو کہ $t$ پہلے ای~$!A$ وچ آ سکدا اے؛ مثال لئی، $!A$ دا روپ
~$\Atom{\Obj P}{t,x}$ ہو سکدا اے۔ ایس لئی،~$\Gamma \Sequent \Delta,
\lexists[x][\Atom{\Obj P}{t,x}]$ دا استنباط~$\Gamma \Sequent \Delta,
\Atom{\Obj P}{t,t}$ توں کرنا \RightR{\lexists} دی درست ورتوں \pnbLatin{(correct application)} اے---تسیں
مقدمے وچ~$t$ دیاں اک یا ودھ، تے لازمی نہیں کہ ساریاں \pnbLatin{(one or more, and not necessarily all)}، آؤن دیاں تھاواں
نوں پابند !!{variable}~$x$ \pnbLatin{(bound variable)} نال ``بدل'' سکدے او۔ پر
\RightR{\lforall} تے~\LeftR{\lexists} دیاں ایگن متغیر شرطاں \pnbLatin{(eigenvariable conditions)} دا تقاضا
اے کہ !!{constant}~$a$،~$!A$ وچ نہ آوے۔ ایس لئی، تسیں
$\Gamma \Sequent \Delta, \lforall[x][\Atom{\Obj P}{a,x}]$ دا
استنباط $\Gamma \Sequent \Delta, \Atom{\Obj P}{a,a}$ توں
$\RightR{\lforall}$ ورت کے درست طور تے نہیں کر سکدے۔
\end{explain}

\begin{explain}
\RightR{\lexists} تے \LeftR{\lforall} وچ ٹرم~$t$ اُتے کوئی ہور پابندی \pnbLatin{(no further restrictions)}
نہیں۔ دوجے پاسے، \LeftR{\lexists} تے \RightR{\lforall} دے قاعدیاں
وچ، ایگن متغیر دی شرط دا تقاضا اے کہ !!{constant}~$a$ اُتّے والے
سیکوئنٹ \pnbLatin{(upper sequent)} وچ~$!A(a)$ توں باہر کتے وی نہ آوے \pnbLatin{(nowhere outside)}۔ ایہ نظام دی صحت مندی (soundness) یقینی
بناؤن لئی ضروری اے، یعنی ایہ صرف معتبر سیکوئنٹ \pnbLatin{(only valid sequents)} !!{derive}s کرے۔ ایس
شرط توں بغیر، ہِٹھ لکھے استنباط جائز ہُندے:
\begin{prooftree}
  \Axiom$!A(a) \fCenter !A(a)$
  \RightLabel{*\LeftR{\lexists}}
  \UnaryInf$\lexists[x][!A(x)] \fCenter !A(a)$
  \RightLabel{\RightR{\lforall}}
  \UnaryInf$\lexists[x][!A(x)] \fCenter \lforall[x][!A(x)]$
  \DisplayProof\bottomAlignProof
  \qquad
  \Axiom$!A(a) \fCenter !A(a)$
  \RightLabel{*\RightR{\lforall}}
  \UnaryInf$!A(a) \fCenter \lforall[x][!A(x)]$
  \RightLabel{\LeftR{\lexists}}
  \UnaryInf$\lexists[x][!A(x)] \fCenter \lforall[x][!A(x)]$
\end{prooftree}
پر $\lexists[x][!A(x)] \Sequent \lforall[x][!A(x)]$ معتبر نہیں \pnbLatin{(not valid)}۔
\end{explain}
\endgroup
% /OLP-0072

% OLP-0073 content/first-order-logic/sequent-calculus/structural-rules.tex
\begingroup
\def\olfilename{structural-rules}
\expandafter\def\csname ol@sectioncs\endcsname{section}

\iftag{FOL}
      {\olfileid{fol}{seq}{srl}}
      {\olfileid{pl}{seq}{srl}}

\olsection{ساختی قاعدے \pnbLatin{(Structural Rules)}}

سانوں کجھ ایہو جیہے ساختی قاعدیاں \pnbLatin{(structural rules)} دی وی لوڑ اے جیہڑے سیکوئنٹ دے کھبے تے
سجے پاسے !!{sentence}s نوں مُڑ ترتیب دین (rearrange) دی اجازت دیندے نیں۔ چونکہ
منطقی قاعدیاں \pnbLatin{(logical rules)} دا تقاضا اے کہ مقدمے (premise) وچ اوہ !!{sentence}s، جِنہاں اُتے
قاعدہ لاگو ہُندا اے، یا تاں بالکل کھبے سرے یا بالکل سجے سرے \pnbLatin{(far left or far right)} اُتے ہون،
ایس لئی سانوں اک ``تبادلے'' دے قاعدے \pnbLatin{(Exchange rule)} دی لوڑ اے جیہڑا !!{sentence}s
نوں ٹھیک تھاں \pnbLatin{(required position)} اُتے لے جان دی اجازت دیندا اے۔ کدی کدی دو اکّو جیہے
!!{sentence}s \pnbLatin{(two identical sentences)} نوں اک وچ ملا سکنا، تے کسے وی پاسے اک !!a{sentence}
شامل کر سکنا \pnbLatin{(add on either side)} وی ضروری اے۔

\subsection{کمزور کرنا \pnbLatin{(Weakening)}}

\begin{defish}
\Axiom$ \Gamma \fCenter \Delta $
\RightLabel{\LeftR{\Weakening}}
\UnaryInf$ !A, \Gamma \fCenter \Delta$
\DisplayProof
\hfill
\Axiom$ \Gamma \fCenter \Delta$
\RightLabel{\RightR{\Weakening}}
\UnaryInf$ \Gamma \fCenter \Delta, !A$
\DisplayProof
\end{defish}

\subsection{سکیڑنا \pnbLatin{(Contraction)}}

\begin{defish}
\Axiom$ !A, !A, \Gamma \fCenter \Delta $
\RightLabel{\LeftR{\Contraction}}
\UnaryInf$ !A, \Gamma \fCenter \Delta$
\DisplayProof
\hfill
\Axiom$ \Gamma \fCenter \Delta, !A, !A$
\RightLabel{\RightR{\Contraction}}
\UnaryInf$ \Gamma \fCenter \Delta, !A$
\DisplayProof
\end{defish}

\subsection{تبادلہ \pnbLatin{(Exchange)}}

\begin{defish}
\Axiom$ \Gamma, !A, !B, \Pi \fCenter \Delta $
\RightLabel{\LeftR{\Exchange}}
\UnaryInf$ \Gamma, !B, !A, \Pi \fCenter \Delta$
\DisplayProof
\hfill
\Axiom$ \Gamma \fCenter \Delta, !A, !B, \Lambda$
\RightLabel{\RightR{\Exchange}}
\UnaryInf$ \Gamma \fCenter \Delta, !B, !A, \Lambda$
\DisplayProof
\end{defish}

کمزور کرن، سکیڑن تے تبادلے دے کئی لگاتار استنباط \pnbLatin{(series of inferences)} اکثر دُہریاں
استنباطی لیکاں \pnbLatin{(double inference lines)} نال وکھائے جان گے۔

ہِٹھلا قاعدہ، جہنوں ``کٹ'' (Cut) آکھیا جاندا اے، سختی نال ویکھیا جاوے تاں
ضروری نہیں \pnbLatin{(not strictly necessary)}، پر !!{derivation}s نوں مُڑ ورتن تے اکٹھیاں کرن \pnbLatin{(reuse and combine)} نوں بہت
سَوکھا بنا دیندا اے۔
\begin{defish}
\[
\Axiom$ \Gamma \fCenter \Delta, !A$
\Axiom$ !A, \Pi \fCenter \Lambda $
\RightLabel{\Cut}
\BinaryInf$ \Gamma, \Pi \fCenter \Delta, \Lambda$
\DisplayProof
\]
\end{defish}
\endgroup
% /OLP-0073

% OLP-0074 content/first-order-logic/sequent-calculus/derivations.tex
\begingroup
\def\olfilename{derivations}
\expandafter\def\csname ol@sectioncs\endcsname{section}

\iftag{FOL}
      {\olfileid{fol}{seq}{der}}
      {\olfileid{pl}{seq}{der}}

\olsection{\usetoken{P}{derivation}}

\begin{explain}
اسیں دَس چکے آں کہ ابتدائی سیکوئنٹ \pnbLatin{(initial sequent)} کیہو جہا ہُندا اے، تے استنباط دے
قاعدے \pnbLatin{(rules of inference)} وی دے چکے آں۔ سیکوئنٹ حساب \pnbLatin{(sequent calculus)} وچ !!^{derivation}s ایہناں توں
استقرائی طور تے بنائے جاندے نیں \pnbLatin{(inductively generated)}: ہر !!{derivation} یا تاں اپنے آپ وچ
اک ابتدائی سیکوئنٹ ہُندا اے، یا اک یا دو !!{derivation}s تے اوہناں دے
بعد اک استنباط \pnbLatin{(one or two derivations followed by an inference)} اُتے مشتمل ہُندا اے۔
\end{explain}

\begin{defn}[$\Log{LK}$ !!{derivation}]
اک \emph{$\Log{LK}$-!!{derivation}} \pnbLatin{(LK derivation)} جیہڑی سیکوئنٹ~$S$ دی اے،
سیکوئنٹاں دا اک محدود درخت \pnbLatin{(finite tree of sequents)} اے جیہڑا ہِٹھ دتیاں شرطاں پوری کردا اے:
\begin{enumerate}
\item درخت دے سب توں اُتلے سیکوئنٹ \pnbLatin{(topmost sequents)} ابتدائی سیکوئنٹ ہُندے نیں۔
\item درخت دا سب توں ہِٹھلا سیکوئنٹ \pnbLatin{(bottommost sequent)}~$S$ اے۔
\item درخت وچ $S$ توں سوا ہر سیکوئنٹ، استنباط دے قاعدے دی کسے درست
  ورتوں دا مقدمہ \pnbLatin{(premise of a correct application)} اے، جیہدا نتیجہ درخت وچ اوس سیکوئنٹ دے عین ہِٹھاں
  کھلوتا اے \pnbLatin{(conclusion directly below)}۔
\end{enumerate}
فیر اسیں $S$ نوں !!{derivation} دا \emph{آخرلا سیکوئنٹ (end-sequent)} آکھدے آں تے
ایہ وی آکھدے آں کہ $S$، \emph{$\Log{LK}$ وچ !!{derivable} اے} (derivable in LK؛ یا
$\Log{LK}$-!!{derivable} اے)۔
\end{defn}

\begin{ex}
ہر ابتدائی سیکوئنٹ \pnbLatin{(initial sequent)}، مثلاً $!C \Sequent !C$، اک !!a{derivation} اے۔
اس اُتے، من لو، $\LeftR{\Weakening}$ قاعدہ \pnbLatin{(Weakening rule)} لا کے اسیں اک نویں
!!{derivation} حاصل کر سکدے آں،
\begin{prooftree}
\Axiom$ \Gamma \fCenter \Delta $
\RightLabel{\LeftR{\Weakening}}
\UnaryInf$ !A, \Gamma \fCenter \Delta$
\end{prooftree}
پر ایہ قاعدہ عام روپ \pnbLatin{(general rule schema)} وچ ورتن لئی اے: قاعدے وچلے $!A$ دی تھاں اسیں
کوئی وی !!{sentence} رکھ سکدے آں \pnbLatin{(replace with any sentence)}، مثلاً~$!D$ وی۔ جے مقدمہ ساڈے
ابتدائی سیکوئنٹ $!C \Sequent !C$ نال میل کھاندا اے \pnbLatin{(premise matches)}، تاں ایس دا مطلب
اے کہ $\Gamma$ تے $\Delta$ دونیں صرف~$!C$ نیں، تے نتیجہ فیر
$!D, !C \Sequent !C$ ہووے گا (conclusion)۔ سو ہِٹھ لکھی شے اک درست !!a{derivation} \pnbLatin{(correct derivation)} اے:
\begin{prooftree}
\Axiom$ !C \fCenter !C $
\RightLabel{\LeftR{\Weakening}}
\UnaryInf$ !D, !C \fCenter !C$
\end{prooftree}
ہُن اسیں کوئی ہور قاعدہ، من لو $\LeftR{\Exchange}$ (Exchange)، لا سکدے آں،
جیہڑا سانوں کھبے پاسے دے دو !!{sentence}s نوں آپو وچ بدل دین \pnbLatin{(switch two sentences on the left)} دی
اجازت دیندا اے۔ سو ہِٹھ لکھی شے وی اک درست !!{derivation} اے:
\begin{prooftree}
\Axiom$ !C \fCenter !C $
\RightLabel{\LeftR{\Weakening}}
\UnaryInf$ !D, !C \fCenter !C$
\RightLabel{\LeftR{\Exchange}}
\UnaryInf$ !C, !D \fCenter !C$
\end{prooftree}
ایس قاعدے دی ورتوں \pnbLatin{(application of the rule)} وچ—قاعدہ ایویں دِتا گیا سی—
\begin{prooftree}
\Axiom$ \Gamma, !A, !B, \Pi \fCenter \Delta $
\RightLabel{\LeftR{\Exchange}}
\UnaryInf$ \Gamma, !B, !A, \Pi \fCenter \Delta,$
\end{prooftree}
$\Gamma$ تے $\Pi$ دونیں خالی (empty) سن، $\Delta$، $!C$ اے، تے $!A$ تے
$!B$ دے کردار (roles) ترتیب وار (respectively) $!D$ تے~$!C$ نے نبھائے نیں۔ لگ بھگ
ایسے طریقے نال اسیں ایہ وی ویکھدے آں کہ
\begin{prooftree}
\Axiom$ !D \fCenter !D $
\RightLabel{\LeftR{\Weakening}}
\UnaryInf$ !C, !D \fCenter !D$
\end{prooftree}
اک !!a{derivation} اے۔ ہُن اسیں ایہناں دوناں !!{derivation}s نوں لے
کے $\RightR{\land}$ نال اکٹھیاں کر سکدے آں \pnbLatin{(combine using right conjunction)}۔ اوہ قاعدہ ایہ سی
\begin{prooftree}
\Axiom$\Gamma \fCenter \Delta, !A$
\Axiom$ \Gamma \fCenter \Delta, !B$
\RightLabel{\RightR{\land}}
\BinaryInf$ \Gamma \fCenter \Delta, !A \land !B $
\end{prooftree}
ساڈے معاملے وچ، مقدمیاں نوں اوہناں !!{derivation}s دے آخرلے سیکوئنٹاں
نال میل کھانا چاہیدا اے \pnbLatin{(premises match the last sequents)} جیہڑے ایہناں مقدمیاں اُتے ختم ہُندے نیں۔ ایس دا مطلب
اے کہ $\Gamma$، $!C, !D$ اے، $\Delta$ خالی اے، $!A$، $!C$ اے تے
$!B$، $!D$ اے۔ سو جے استنباط درست ہونا اے، تاں نتیجہ $!C, !D \Sequent !C
\land !D$ اے (conclusion)۔
\begin{prooftree}
\Axiom$ !C \fCenter !C $
\RightLabel{\LeftR{\Weakening}}
\UnaryInf$ !D, !C \fCenter !C$
\RightLabel{\LeftR{\Exchange}}
\UnaryInf$ !C, !D \fCenter !C$
\Axiom$ !D \fCenter !D $
\RightLabel{\LeftR{\Weakening}}
\UnaryInf$ !C, !D \fCenter !D$
\RightLabel{\RightR{\land}}
\BinaryInf$ !C, !D \fCenter !C \land !D $
\end{prooftree}
یقیناً، اسیں مقدمیاں دی ترتیب اُلٹ وی سکدے آں \pnbLatin{(reverse the premises)}؛ فیر $!A$، $!D$
ہووے گا تے $!B$،~$!C$ ہووے گا۔
\begin{prooftree}
\Axiom$ !D \fCenter !D $
\RightLabel{\LeftR{\Weakening}}
\UnaryInf$ !C, !D \fCenter !D$
\Axiom$ !C \fCenter !C $
\RightLabel{\LeftR{\Weakening}}
\UnaryInf$ !D, !C \fCenter !C$
\RightLabel{\LeftR{\Exchange}}
\UnaryInf$ !C, !D \fCenter !C$
\RightLabel{\RightR{\land}}
\BinaryInf$ !C, !D \fCenter !D \land !C $
\end{prooftree}
\end{ex}
\endgroup
% /OLP-0074

% OLP-0075 content/first-order-logic/sequent-calculus/proving-things.tex
\begingroup
\def\olfilename{proving-things}
\expandafter\def\csname ol@sectioncs\endcsname{section}

\iftag{FOL}
      {\olfileid{fol}{seq}{pro}}
      {\olfileid{pl}{seq}{pro}}

\olsection{\usetoken{P}{derivation} دیاں مثالاں \pnbLatin{(Examples of derivation)}}

\begin{ex}
اک $\Log{LK}$-!!{derivation} سیکوئنٹ $!A \land !B \Sequent !A$ لئی دیو \pnbLatin{(Give an LK derivation)}۔

اسیں مطلوب آخرلا سیکوئنٹ \pnbLatin{(desired end-sequent)} !!{derivation} دے سب توں ہِٹھاں لکھ کے شروع کردے آں۔
\begin{prooftree}
\AxiomC{}
\UnaryInf$!A\land !B \fCenter !A$
\end{prooftree}
اگلا کم ایہ لبھنا اے کہ ایس روپ دے ہِٹھلے سیکوئنٹ \pnbLatin{(lower sequent)} والی استنباط
(inference) کیہڑی ہو سکدی اے۔ ایہ کوئی ساختی قاعدہ \pnbLatin{(structural rule)} ہو سکدا اے،
پر شروع وچ منطقی قاعدہ \pnbLatin{(logical rule)} لبھنا چنگا رہندا اے۔ ہِٹھلے سیکوئنٹ وچ
اکّو منطقی رابط \pnbLatin{(logical connective)}، یعنی $\land$، آندا اے، سو اسیں $\land$ دا
قاعدہ لبھ رہے آں۔ نالے $\land$ دی علامت مقدم (antecedent) وچ اے، ایس لئی اسیں
\LeftR{\land} قاعدہ ویکھدے آں۔
\begin{prooftree}
\AxiomC{}
\RightLabel{\LeftR{\land}}
\UnaryInf$!A\land !B \fCenter !A$
\end{prooftree}
\LeftR{\land} استنباط دے اوپری سیکوئنٹ \pnbLatin{(upper sequent)} لئی دو چوناں نیں: اوپری
سیکوئنٹ $!A \Sequent !A$ ہو سکدا اے، یا $!B \Sequent !A$۔ صاف اے کہ $!A
\Sequent !A$ اک ابتدائی سیکوئنٹ \pnbLatin{(initial sequent)} اے—جیہڑا ساڈے لئی چنگا اے—جد کہ
$!B \Sequent !A$ عام طور تے اخذ ہون جوگ نہیں \pnbLatin{(not derivable in general)}۔ سو اسیں اوپری سیکوئنٹ بھر دیندے آں:
\begin{prooftree}
\Axiom$!A \fCenter !A$
\RightLabel{\LeftR{\land}}
\UnaryInf$!A\land !B \fCenter !A$
\end{prooftree}
ہُن ساڈے کول اک درست $\Log{LK}$-!!{derivation} اے جیہڑی سیکوئنٹ $!A
\land !B \Sequent !A$ دی اے۔
\end{ex}

\begin{ex}
اک $\Log{LK}$-!!{derivation} سیکوئنٹ $\lnot !A \lor !B
\Sequent !A \lif !B$ لئی دیو \pnbLatin{(Give an LK derivation)}۔

مطلوب آخرلا سیکوئنٹ \pnbLatin{(desired end-sequent)} !!{derivation} دے سب توں ہِٹھاں لکھ کے شروع کرو۔
\begin{prooftree}
\AxiomC{}
\UnaryInf$\lnot !A \lor !B \fCenter !A \lif !B$
\end{prooftree}
اوہ منطقی قاعدہ لبھن لئی جیہڑا سانوں ایہ آخرلا سیکوئنٹ دے سکدا اے، اسیں آخرلے
سیکوئنٹ دے منطقی رابطیاں \pnbLatin{(logical connectives)}، یعنی $\lnot$، $\lor$ تے $\lif$،
نوں ویکھدے آں۔ ایس ویلے ساڈا سروکار صرف $\lor$ تے $\lif$ نال اے، کیوں جے
اوہ آخرلے سیکوئنٹ دیاں !!{sentence}s دے !!{main operator}s نیں؛ $\lnot$ کسے ہور
رابط دے دائرے (scope) اندر اے، سو اوہدے نال اسیں بعد وچ نمٹاں گے۔ ایس لئی آخری
استنباط دے منطقی قاعدے لئی ساڈے کول \LeftR{\lor} تے \RightR{\lif} دیاں چوناں
نیں۔ اصل وچ کوئی وی چُن سکدے آں، پر آؤ \RightR{\lif} قاعدہ چُنیے—جے ہور کوئی
سبب نہ وی ہووے، تاں ایہ سانوں ہُنّے دو شاخاں وچ ونڈن \pnbLatin{(splitting into two branches)}
نوں ملتوی کرن دیندا اے۔ جیہڑی \RightR{\lif} استنباط ایہ ہِٹھلا سیکوئنٹ دے سکدی
اے، اوہدے روپ مطابق ایہ اینج ہونی چاہیدی اے:
\begin{prooftree}
\AxiomC{}
\UnaryInf$ !A, \lnot !A \lor !B \fCenter !B $
\RightLabel{\RightR{\lif}} \UnaryInf$ \lnot !A \lor !B \fCenter !A \lif !B $
\end{prooftree}

جے اسیں $\lnot !A \lor !B$ نوں مقدم دے باہرلے سرے ول لے جائیے، تاں
\LeftR{\lor} قاعدہ لا سکدے آں۔ خاکے (schema) مطابق ایہ دو اوپری سیکوئنٹاں وچ
اینج ونڈیا جانا چاہیدا اے:
% PNBSEQ-001 ماخذ دی درستی: ہِٹھ دِتیاں چار تھاواں اُتے فارمولے صرف مقدم دے اندر بدلدے نیں، ایس لئی Right Exchange دی تھاں Left Exchange اے؛ نتیجے والے پاسے دیاں دو Right Exchange ورتوں جوں دیاں توں نیں۔
\begin{prooftree}
\AxiomC{}
\UnaryInf$\lnot !A, !A \fCenter !B$
\AxiomC{}
\UnaryInf$!B, !A \fCenter !B$
\RightLabel{\LeftR{\lor}}
\BinaryInf$ \lnot !A \lor !B, !A \fCenter !B $
\RightLabel{\LeftR{\Exchange}}
\UnaryInf$ !A, \lnot !A \lor !B \fCenter !B $
\RightLabel{\RightR{\lif}}
\UnaryInf$ \lnot !A \lor !B \fCenter !A \lif !B $
\end{prooftree}
چیتا رکھو کہ اسیں اُتّے ول اپنا رستہ بناندے ہوے \pnbLatin{(wind our way up)} ابتدائی
سیکوئنٹاں تک پہنچن دی کوشش کر رہے آں؛ لگدا اے اسیں کافی نیڑے آں!{} سجی شاخ
\pnbLatin{(right branch)} ابتدائی سیکوئنٹ توں صرف اک Weakening تے اک Exchange دور اے، تے
فیر اوہ پوری ہو جاندی اے:
\begin{prooftree}
\AxiomC{}
\UnaryInf$\lnot !A, !A \fCenter !B$
\Axiom$!B \fCenter !B$
\RightLabel{\LeftR{\Weakening}}
\UnaryInf$!A, !B \fCenter !B$
\RightLabel{\LeftR{\Exchange}}
\UnaryInf$!B, !A \fCenter !B$
\RightLabel{\LeftR{\lor}}
\BinaryInf$\lnot !A \lor !B, !A \fCenter !B $
\RightLabel{\LeftR{\Exchange}}
\UnaryInf$ !A, \lnot !A \lor !B \fCenter !B $
\RightLabel{\RightR{\lif}}
\UnaryInf$ \lnot !A \lor !B \fCenter !A \lif !B $
\end{prooftree}

ہُن کھبی شاخ \pnbLatin{(left branch)} ول ویکھو۔ کسے وی !!{sentence} وچ اکّو منطقی رابط،
یعنی مقدم والیاں !!{sentence}s وچ $\lnot$ دی علامت، اے؛ سو اسیں \LeftR{\lnot}
قاعدے دی اک ورتوں (instance) لبھ رہے آں۔
\begin{prooftree}
\AxiomC{}
\UnaryInf$ !A \fCenter !B, !A$
\RightLabel{\LeftR{\lnot}}
\UnaryInf$\lnot !A, !A \fCenter !B$
\Axiom$!B \fCenter !B$
\RightLabel{\LeftR{\Weakening}}
\UnaryInf$!A, !B \fCenter !B$
\RightLabel{\LeftR{\Exchange}}
\UnaryInf$!B, !A \fCenter !B$
\RightLabel{\LeftR{\lor}}
\BinaryInf$\lnot !A \lor !B, !A \fCenter !B $
\RightLabel{\LeftR{\Exchange}}
\UnaryInf$ !A, \lnot !A \lor !B \fCenter !B $
\RightLabel{\RightR{\lif}}
\UnaryInf$ \lnot !A \lor !B \fCenter !A \lif !B $
\end{prooftree}
جیویں اسیں سجی شاخ پوری کیتی سی، اویویں ایس کھبی شاخ نوں وی پورا کرن لئی صرف اک
Weakening تے اک Exchange دی لوڑ اے۔
\begin{prooftree}
\Axiom$!A \fCenter !A$
\RightLabel{\RightR{\Weakening}}
\UnaryInf$ !A \fCenter !A, !B$
\RightLabel{\RightR{\Exchange}}
\UnaryInf$ !A \fCenter !B, !A$
\RightLabel{\LeftR{\lnot}}
\UnaryInf$\lnot !A, !A \fCenter !B$
\Axiom$!B \fCenter !B$
\RightLabel{\LeftR{\Weakening}}
\UnaryInf$!A, !B \fCenter !B$
\RightLabel{\LeftR{\Exchange}}
\UnaryInf$!B, !A \fCenter !B$
\RightLabel{\LeftR{\lor}}
\BinaryInf$\lnot !A \lor !B, !A \fCenter !B $
\RightLabel{\LeftR{\Exchange}}
\UnaryInf$ !A, \lnot !A \lor !B \fCenter !B $
\RightLabel{\RightR{\lif}}
\UnaryInf$ \lnot !A \lor !B \fCenter !A \lif !B $
\end{prooftree}
\end{ex}

\begin{ex}
اک $\Log{LK}$-!!{derivation} سیکوئنٹ $\lnot !A \lor \lnot !B
\Sequent \lnot (!A \land !B)$ لئی دیو \pnbLatin{(Give an LK derivation)}۔

اُتّے والے طریقے ورت کے اسیں مطلوب آخرلا سیکوئنٹ \pnbLatin{(desired end-sequent)} سب توں
ہِٹھاں لکھ کے شروع کردے آں۔
\begin{prooftree}
\AxiomC{}
\UnaryInf$ \lnot !A \lor \lnot !B \fCenter \lnot (!A \land !B) $
\end{prooftree}
آخرلے سیکوئنٹ دیاں !!{sentence}s دے دستیاب مرکزی رابط \pnbLatin{(available main connectives)} $\lor$ تے $\lnot$ نیں۔ ایتھے \LeftR{\lor} یا \RightR{\lnot}،
دونواں وچوں کوئی وی قاعدہ لگ سکدا اے، پر اسیں \RightR{\lnot} نال شروع کردے آں،
کیوں جے ایہ اک پل لئی دو شاخاں وچ ونڈن توں بچاندا اے:
\begin{prooftree}
\AxiomC{}
\UnaryInf$!A \land !B, \lnot !A \lor \lnot !B \fCenter $
\RightLabel{\RightR{\lnot}}
\UnaryInf$\lnot !A \lor \lnot !B \fCenter \lnot (!A \land !B)$
\end{prooftree}
% PNBSEQ-002 ماخذ دی درستی: ہِٹھلے دونیں متبادل سیکوئنٹاں وچ B دی نفی ماخذ توں رہ گئی سی؛ مثال دے آخرلے سیکوئنٹ تے اگلے سارے ثبوتی درخت not-B ورتدے نیں۔
ہُن ساڈے کول \LeftR{\land} یا \LeftR{\lor} قاعدہ ویکھن دی چون اے۔ آؤ ویکھئیے
کہ \LeftR{\land} لاؤن نال کی ہُندا اے: اسیں سیکوئنٹ $!A,
\lnot !A \lor \lnot !B \Sequent \quad$ یا سیکوئنٹ $!B, \lnot !A
\lor \lnot !B \Sequent \quad$، دونواں وچوں کسے نال شروع کر سکدے آں۔ کیوں جے
!!{derivation}، $!A$ تے $!B$ دے لحاظ نال متناظر (symmetric) اے، سو پہلا چُن لئیے:
\begin{prooftree}
\AxiomC{}
\UnaryInf$!A, \lnot !A \lor \lnot !B \fCenter $
\RightLabel{\LeftR{\land}}
\UnaryInf$!A \land !B, \lnot !A \lor \lnot !B \fCenter $
\RightLabel{\RightR{\lnot}}
\UnaryInf$\lnot !A \lor \lnot !B \fCenter \lnot (!A \land !B)$
\end{prooftree}
!!{derivation} نوں اگے بھردے ہوے سانوں پتا لگدا اے کہ اک مسئلہ آ پیندا اے:
\begin{prooftree}
\Axiom$!A \fCenter !A$
\RightLabel{\LeftR{\lnot}}
\UnaryInf$ \lnot !A, !A \fCenter$
\AxiomC{}
\RightLabel{?}
\UnaryInf$!A \fCenter !B$
\RightLabel{\LeftR{\lnot}}
\UnaryInf$ \lnot !B, !A \fCenter$
\RightLabel{\LeftR{\lor}}
\BinaryInf$\lnot !A \lor \lnot !B, !A \fCenter $
\RightLabel{\LeftR{\Exchange}}
\UnaryInf$!A, \lnot !A \lor \lnot !B \fCenter $
\RightLabel{\LeftR{\land}}
\UnaryInf$!A \land !B, \lnot !A \lor \lnot !B \fCenter $
\RightLabel{\RightR{\lnot}}
\UnaryInf$\lnot !A \lor \lnot !B \fCenter \lnot (!A \land !B)$
\end{prooftree}
سجی شاخ دے اُتلے سیکوئنٹ نوں ہور گھٹایا (reduced) نہیں جا سکدا، تے ساختی
استنباطاں راہیں وی اوہنوں ابتدائی سیکوئنٹ نہیں بنایا جا سکدا۔ سو ایہ چُنن والا
ثبوتی رستہ \pnbLatin{(proof path)} درست نہیں۔ صاف اے کہ اُتّے \LeftR{\land} قاعدہ لانا
غلتی سی۔ پچھلی حالت ول مُڑ کے (backtracking) اوہدی تھاں \LeftR{\lor} قاعدہ
لاواں، تاں سانوں ایہ ملدا اے:
\begin{prooftree}
\AxiomC{}
\UnaryInf$\lnot !A, !A \land !B \fCenter $

\AxiomC{}
\UnaryInf$\lnot !B, !A \land !B \fCenter $

\RightLabel{\LeftR{\lor}}
\BinaryInf$\lnot !A \lor \lnot !B, !A \land !B \fCenter $
\RightLabel{\LeftR{\Exchange}}
\UnaryInf$!A \land !B, \lnot !A \lor \lnot !B \fCenter $
\RightLabel{\RightR{\lnot}}
\UnaryInf$\lnot !A \lor \lnot !B \fCenter \lnot (!A \land !B)$
\end{prooftree}
ہر شاخ نوں پہلاں وانگ پورا کر کے سانوں ایہ ملدا اے:
\begin{prooftree}
\Axiom$ !A \fCenter!A$
\RightLabel{\LeftR{\land}}
\UnaryInf$!A \land !B \fCenter !A$
\RightLabel{\LeftR{\lnot}}
\UnaryInf$\lnot !A, !A \land !B \fCenter $

\Axiom$ !B \fCenter !B$
\RightLabel{\LeftR{\land}}
\UnaryInf$!A \land !B \fCenter !B$
\RightLabel{\LeftR{\lnot}}
\UnaryInf$\lnot !B, !A \land !B \fCenter $

\RightLabel{\LeftR{\lor}}
\BinaryInf$\lnot !A \lor \lnot !B, !A \land !B \fCenter $
\RightLabel{\LeftR{\Exchange}}
\UnaryInf$!A \land !B, \lnot !A \lor \lnot !B \fCenter $
\RightLabel{\RightR{\lnot}}
\UnaryInf$\lnot !A \lor \lnot !B \fCenter \lnot (!A \land !B)$
\end{prooftree}
(ایہناں قدماں وچ اسیں $\land$ دے قاعدے $\lnot$ دے قاعدیاں توں ہِٹھاں وی لا
سکدے ساں، تے فیر وی اک درست !!{derivation} ملدی۔)
\end{ex}

\begin{ex}
ہُن تک اسیں Contraction دا قاعدہ نہیں ورتیا، پر کدے کدے اوہدی لوڑ پیندی اے۔
ایہ اوہو جیہی اک مثال اے۔ من لو اسیں $\quad \Sequent !A \lor \lnot !A$ ثابت
کرنا چاہنے آں۔ $\RightR{\lor}$ نوں پچھلے پاسے ول (backwards) لاؤن نال سانوں
ہِٹھ دِتیاں دوناں !!{derivation}s وچوں اک ملے گی:
\begin{prooftree}
\AxiomC{}
\UnaryInf$ \fCenter !A$
\RightLabel{\RightR{\lor}}
\UnaryInf$ \fCenter !A \lor \lnot !A$
\DisplayProof\qquad\bottomAlignProof
\AxiomC{}
\UnaryInf$!A \fCenter $
\RightLabel{\RightR{\lnot}}
\UnaryInf$ \fCenter \lnot !A$
\RightLabel{\RightR{\lor}}
\UnaryInf$ \fCenter !A \lor \lnot !A$
\end{prooftree}
ظاہر اے کہ ایہناں وچوں کوئی وی ابتدائی سیکوئنٹ اُتے ختم نہیں ہُندی۔ ترکیب ایہ
سمجھنا اے کہ Contraction دا قاعدہ سانوں !!a{sentence} دیاں دو نقلاں نوں اک وچ
سکیڑن \pnbLatin{(combine two copies into one)} دیندا اے—تے جدوں اسیں ثبوت لبھ رہے آں، یعنی
ہِٹھوں اُتّے ول \pnbLatin{(bottom to top)} جا رہے آں، تاں مقدمے وچ $!A \lor \lnot !A$ دی
اک نقل سنبھال کے رکھ سکدے آں، مثلاً:
\begin{prooftree}
\AxiomC{}
\UnaryInf$ \fCenter !A \lor \lnot !A, !A$
\RightLabel{\RightR{\lor}}
\UnaryInf$ \fCenter !A \lor \lnot !A, !A \lor \lnot !A$
\RightLabel{\RightR{\Contraction}}
\UnaryInf$ \fCenter !A \lor \lnot !A$
\end{prooftree}
ہُن اسیں $\RightR{\lor}$ دوسری واری لا سکدے آں تے~$\lnot !A$ وی حاصل کر
لیندے آں، جیہڑا اک مکمل !!{derivation} تک لے جاندا اے۔
\begin{prooftree}
\Axiom$!A \fCenter !A$
\RightLabel{\RightR{\lnot}}
\UnaryInf$\fCenter !A, \lnot !A$
\RightLabel{\RightR{\lor}}
\UnaryInf$\fCenter !A, !A \lor \lnot !A$
\RightLabel{\RightR{\Exchange}}
\UnaryInf$ \fCenter !A \lor \lnot !A, !A$
\RightLabel{\RightR{\lor}}
\UnaryInf$ \fCenter !A \lor \lnot !A, !A \lor \lnot !A$
\RightLabel{\RightR{\Contraction}}
\UnaryInf$ \fCenter !A \lor \lnot !A$
\end{prooftree}
\end{ex}

\begin{prob}
ہِٹھلے سیکوئنٹاں دیاں !!{derivation}s دیو \pnbLatin{(Give derivations)}:
\begin{enumerate}
\item $!A \land (!B \land !C) \Sequent (!A \land !B) \land !C$.
\item $!A \lor (!B \lor !C) \Sequent (!A \lor !B) \lor !C$.
\item $!A \lif (!B \lif !C) \Sequent !B \lif (!A \lif !C)$.
\item $!A \Sequent \lnot\lnot !A$.
\end{enumerate}
\end{prob}

\begin{prob}
ہِٹھلے سیکوئنٹاں دیاں !!{derivation}s دیو \pnbLatin{(Give derivations)}:
\begin{enumerate}
\item $(!A \lor !B) \lif !C \Sequent !A \lif !C$.
\item $(!A \lif !C) \land (!B \lif !C) \Sequent (!A \lor !B) \lif !C$.
\item $\Sequent \lnot(!A \land \lnot !A)$.
\item $!B \lif !A \Sequent \lnot !A \lif \lnot !B$.
\item $\Sequent (!A \lif \lnot !A) \lif \lnot !A$.
\item $\Sequent \lnot(!A \lif !B) \lif \lnot !B$.
\item $!A \lif !C \Sequent \lnot (!A \land \lnot !C)$.
\item $!A \land \lnot !C \Sequent \lnot (!A \lif !C)$.
\item $!A \lor !B, \lnot !B \Sequent !A$.
\item $\lnot !A \lor \lnot !B \Sequent \lnot(!A \land !B)$.
\item $\Sequent (\lnot !A \land \lnot !B) \lif\lnot(!A \lor !B)$.
\item $\Sequent \lnot(!A \lor !B) \lif (\lnot !A \land \lnot !B)$.
\end{enumerate}
\end{prob}

\begin{prob}
ہِٹھلے سیکوئنٹاں دیاں !!{derivation}s دیو \pnbLatin{(Give derivations)}:
\begin{enumerate}
\item $\lnot(!A \lif !B) \Sequent !A$.
\item $\lnot(!A \land !B) \Sequent \lnot !A \lor \lnot !B$.
\item $!A \lif !B \Sequent \lnot !A \lor !B$.
\item $\Sequent \lnot \lnot !A \lif !A$.
\item $!A \lif !B, \lnot !A \lif !B \Sequent !B$.
\item $(!A \land !B) \lif !C \Sequent (!A \lif !C) \lor (!B \lif !C)$.
\item $(!A \lif !B) \lif !A \Sequent !A$.
\item $\Sequent (!A \lif !B) \lor (!B \lif !C)$.
\end{enumerate}
(ایہناں ساریاں لئی $\RightR{\Contraction}$ قاعدے دی لوڑ اے \pnbLatin{(These all require the rule)}۔)
\end{prob}
\endgroup
% /OLP-0075

% OLP-0076 content/first-order-logic/sequent-calculus/proving-things-quant.tex
\begingroup
\def\olfilename{proving-things-quant}
\expandafter\def\csname ol@sectioncs\endcsname{section}

\olfileid{fol}{seq}{prq}

\olsection{مسوّراں نال \usetoken{P}{derivation} \pnbLatin{(Derivation with Quantifiers)}}

\begin{ex}
اک $\Log{LK}$-!!{derivation} سیکوئنٹ $\lexists[x][\lnot !A(x)]
\Sequent \lnot \lforall[x][!A(x)]$ لئی دیو \pnbLatin{(Give an LK derivation)}۔

مسوّراں نال کم کردیاں سانوں ایہ یقینی بنانا پیندا اے کہ ایگن متغیر دی شرط
\pnbLatin{(eigenvariable condition)} دی خلاف ورزی نہ ہووے، تے کدی کدی ایس لئی کجھ استنباط
کرن دی ترتیب نال ردّوبدل کرنا پیندا اے۔ عام طور تے، ایگن متغیر دی شرط والے
قاعدیاں نوں پہلاں نمٹاؤن دی کوشش فائدہ دیندی اے (مکمل ثبوت وچ اوہ ہور ہِٹھاں
ہون گے)۔ نالے، اگوں ویکھ کے ایہ اندازہ لاؤن دی کوشش کرنا چنگا اے کہ ابتدائی
سیکوئنٹ \pnbLatin{(initial sequent)} کیہو جہا ہو سکدا اے۔ ساڈی صورت وچ اوہنوں $!A(a)
\Sequent !A(a)$ ورگا ہونا پئے گا۔ ایس دا مطلب اے کہ جدوں اسیں مسوّراں دے
قاعدیاں نوں ``اُلٹے پاسے'' لاندے آں، تاں اکّو ٹرم \pnbLatin{(same term)}—جیہنوں اسیں
$a$ آکھاں گے—$\lforall$ تے $\lexists$ دے دونیں قاعدیاں لئی چُننا پئے گا۔ جے
اسیں ہر قاعدے لئی وکھرا ٹرم چُنیے، تاں آخر نوں $!A(a) \Sequent !A(b)$ ورگا
کجھ ملے گا، جیہڑا ظاہر اے اخذ ہون جوگ نہیں۔

معمول وانگ شروع کردیاں اسیں لکھدے آں
\begin{prooftree}
\AxiomC{}
\UnaryInf$\lexists[x][\lnot !A(x)] \fCenter \lnot \lforall[x][!A(x)]$
\end{prooftree}
اسیں یا \LeftR{\exists} قاعدہ لا سکدے آں یا \RightR{\lnot} قاعدہ۔ کیوں جے
\LeftR{\exists} قاعدہ ایگن متغیر دی شرط دے تابع اے، ایس نوں بعد دی تھاں
ہُنّے نمٹا لینا چنگا اے؛ سو اسیں پہلاں ایہو قاعدہ لاواں گے۔
\begin{prooftree}
\AxiomC{}
\UnaryInf$ \lnot !A(a) \fCenter \lnot \lforall[x][!A(x)]$
\RightLabel{\LeftR{\lexists}}
\UnaryInf$ \lexists[x][\lnot !A(x)] \fCenter \lnot \lforall[x][!A(x)]$
\end{prooftree}
\LeftR{\lnot} تے \RightR{\lnot} قاعدے پچھلے پاسے ول (backwards) لاؤن نال سانوں
ایہ ملدا اے:
\begin{prooftree}
\AxiomC{}
\UnaryInf$\lforall[x][!A(x)] \fCenter !A(a)$
\RightLabel{\LeftR{\lnot}}
\UnaryInf$\lnot !A(a), \lforall[x][!A(x)] \fCenter $
\RightLabel{\LeftR{\Exchange}}
\UnaryInf$\lforall[x][!A(x)], \lnot !A(a) \fCenter $
\RightLabel{\RightR{\lnot}}
\UnaryInf$ \lnot !A(a) \fCenter \lnot \lforall[x] !A(x)$
\RightLabel{\LeftR{\lexists}}
\UnaryInf$ \lexists[x] \lnot !A(x) \fCenter \lnot \lforall[x] !A(x)$
\end{prooftree}
ایس مقام اُتے ساڈی اکّو چون \LeftR{\forall} قاعدہ لانا اے۔ کیوں جے ایہ قاعدہ
ایگن متغیر والی پابندی دے تابع نہیں، سو ایتھے کوئی رکاوٹ نہیں۔ یاد رکھو، اسیں
$!A(a) \Sequent !A(a)$ دے روپ دا ابتدائی سیکوئنٹ حاصل کرنا چاہنے آں؛ ایس لئی
قاعدہ لاندیاں سانوں $a$ نوں $!A$ وچ رکھن والا ٹرم چُننا چاہیدا اے۔
\begin{prooftree}
\Axiom$!A(a) \fCenter !A(a)$
\RightLabel{\LeftR{\lforall}}
\UnaryInf$\lforall[x][!A(x)] \fCenter !A(a)$
\RightLabel{\LeftR{\lnot}}
\UnaryInf$\lnot !A(a), \lforall[x][!A(x)] \fCenter $
\RightLabel{\LeftR{\Exchange}}
\UnaryInf$\lforall[x][!A(x)], \lnot !A(a) \fCenter $
\RightLabel{\RightR{\lnot}}
\UnaryInf$ \lnot !A(a) \fCenter \lnot \lforall[x][!A(x)]$
\RightLabel{\LeftR{\lexists}}
\UnaryInf$ \lexists[x][ \lnot !A(x)] \fCenter \lnot \lforall[x][!A(x)]$
\end{prooftree}
خاص کر مسوّراں نال کم کردیاں، ایس مقام اُتے دوبارہ جانچنا ضروری اے کہ ایگن
متغیر دی شرط دی خلاف ورزی تے نہیں ہوئی۔ اسیں جیہڑا اکّو قاعدہ لایا اے جیہڑا
ایگن متغیر دی شرط دے تابع اے، اوہ \LeftR{\exists} سی؛ تے ایگن متغیر~$a$
اوہدے ہِٹھلے سیکوئنٹ \pnbLatin{(lower sequent)}، یعنی آخرلے سیکوئنٹ (end-sequent)، وچ
نہیں آندا۔ ایس لئی ایہ اک درست !!{derivation} اے۔
\end{ex}

\begin{prob}
ہِٹھلے سیکوئنٹاں دیاں !!{derivation}s دیو \pnbLatin{(Give derivations)}:
\begin{enumerate}
\item $\Sequent (\lforall[x][!A(x)] \land \lforall[y][!B(y)]) \lif
\lforall[z][(!A(z) \land !B(z))]$.
\item $\Sequent (\lexists[x][!A(x)] \lor \lexists[y][!B(y)]) \lif
\lexists[z][(!A(z) \lor !B(z))]$.
\item $\lforall[x][(!A(x) \lif !B)] \Sequent \lexists[y][!A(y)] \lif !B$.
\item $\lforall[x][\lnot !A(x)] \Sequent \lnot\lexists[x][!A(x)]$.
\item $\Sequent \lnot\lexists[x][!A(x)] \lif \lforall[x][\lnot !A(x)]$.
\item $\Sequent \lnot\lexists[x][\lforall[y][((!A(x,y) \lif \lnot
!A(y,y)) \land (\lnot !A(y,y) \lif !A(x,y)))]]$.
\end{enumerate}
\end{prob}

\begin{prob}
ہِٹھلے سیکوئنٹاں دیاں !!{derivation}s دیو \pnbLatin{(Give derivations)}:
\begin{enumerate}
\item $\Sequent \lnot\lforall[x][!A(x)] \lif \lexists[x][\lnot!A(x)]$.
\item $(\lforall[x][!A(x)] \lif !B) \Sequent \lexists[y][(!A(y) \lif !B)]$.
\item $\Sequent \lexists[x][(!A(x) \lif \lforall[y][!A(y)])]$.
\end{enumerate}
(ایہناں ساریاں لئی $\RightR{\Contraction}$ قاعدے دی لوڑ اے \pnbLatin{(These all require the rule)}۔)
\end{prob}
\endgroup
% /OLP-0076

% OLP-0077 content/first-order-logic/sequent-calculus/proof-theoretic-notions.tex
\begingroup
\def\olfilename{proof-theoretic-notions}
\expandafter\def\csname ol@sectioncs\endcsname{section}

\begin{editorial}
% PNBSEQ-003: frozen source says natural deduction; the complete LK context requires sequent calculus.
  ایہ حصہ سیکوئنٹ حساب \pnbLatin{(sequent calculus)} لئی ثابت پذیری دے تعلق
  \pnbLatin{(provability relation)} تے بے تضادی دیاں تعریفاں اکٹھیاں کردا اے۔
\end{editorial}

\iftag{FOL}
      {\olfileid{fol}{seq}{ptn}}
      {\olfileid{pl}{seq}{ptn}}

\olsection{ثبوتی-نظریاتی تصورات \pnbLatin{(Proof-Theoretic Notions)}}

\begin{explain}
جیویں اسیں کئی اہم معنوی تصورات (منطقی اعتبار، منطقی نتیجہ تے پورا ہون
جوگتا) دی تعریف دِتّی اے، ہُن اسیں اوہناں دے مطابق
\emph{ثبوتی-نظریاتی تصورات \pnbLatin{(proof-theoretic notions)}} دی تعریف دیندے آں۔
ایہناں دی تعریف !!{sentence}s دے !!{structure}s وچ پورا اترن ول رجوع کر کے
نہیں، بلکہ کجھ سیکوئنٹاں دے !!{derivability} یا !!{nonderivability} ول رجوع
کر کے دِتّی جاندی اے۔ ایہ اک اہم دریافت سی کہ ایہ تصور آپو وچ میل کھاندے
نیں۔ اوہناں دے میل کھان دی گل \emph{صحت مندی (soundness)} تے
\emph{تمامیت دے قضیے \pnbLatin{(completeness theorem)}} دا مضمون اے۔
\end{explain}

\begin{defn}[قضیے (Theorems)]
!!^a{sentence}~$!A$ اک \emph{قضیہ (theorem)} اے جے~$\Log{LK}$ وچ
سیکوئنٹ $\quad \Sequent !A$ دی !!a{derivation} ہووے۔ اسیں $\Proves !A$
لکھدے آں جے $!A$ قضیہ اے، تے جے ایہ قضیہ نہیں تاں
$\Proves/ !A$ لکھدے آں۔
\end{defn}

\begin{defn}[!!^{derivability}]
!!^a{sentence}~$!A$، !!{sentence}s دے سیٹ~$\Gamma$ توں
\emph{!!{derivable}} اے، $\Gamma \Proves !A$، اُدوں تے صرف اُدوں جدوں
کوئی محدود ذیلی سیٹ~$\Gamma_0 \subseteq \Gamma$ تے کوئی سلسلہ
$\Gamma_0'$ ایہو جہا ہووے جیہڑا~$\Gamma_0$ وچلے !!{sentence}s دا ہووے،
تے $\Log{LK}$، $\Gamma_0' \Sequent !A$ نوں !!{derive}s کرے۔ جے $!A$، $\Gamma$ توں
!!{derivable} نہیں تاں اسیں $\Gamma \Proves/ !A$ لکھدے آں۔
\end{defn}

سکیڑن، کمزور کرن تے تبادلے دے قاعدیاں کرکے~$\Gamma_0'$ وچ
!!{sentence}s دی ترتیب تے گنتی نال فرق نہیں پیندا: جے سیکوئنٹ
$\Gamma_0' \Sequent !A$ !!{derivable} اے، تاں $\Gamma_0'' \Sequent !A$
وی اے، ہر اوہ~$\Gamma_0''$ لئی جیہدے وچ اوہی !!{sentence}s نیں جیہڑے
~$\Gamma_0'$ وچ نیں۔ مثال لئی، جے
$\Gamma_0 = \{!B, !C\}$ اے، تاں $\Gamma_0' = \tuple{!B, !B, !C}$ تے
$\Gamma_0'' = \tuple{!C, !C, !B}$ دونیں ایہو جہے سلسلے نیں جنہاں وچ
صرف~$\Gamma_0$ دے !!{sentence}s نیں۔ جے ایہناں وچوں اک والا سیکوئنٹ
!!{derivable} اے، تاں دوجا وی اے، مثال دے طور تے:
\begin{prooftree}
  \AxiomC{}
  \Deduce$!B, !B, !C \fCenter !A$
  \RightLabel{\LeftR{\Contraction}}
  \UnaryInf$!B, !C \fCenter !A$
  \RightLabel{\LeftR{\Exchange}}
  \UnaryInf$!C, !B \fCenter !A$
  \RightLabel{\LeftR{\Weakening}}
  \UnaryInf$!C, !C, !B \fCenter !A$
\end{prooftree}
ہُن توں اسیں آکھاں گے کہ جے~$\Gamma_0$، !!{sentence}s دا محدود سیٹ اے،
تاں $\Gamma_0 \Sequent !A$ ہر اوہ سیکوئنٹ اے جیہدے مقدم وچ~$\Gamma_0$
دے !!{sentence}s دا سلسلہ ہووے؛ تے جے لوڑ پئے تاں سکیڑن، تبادلے تے کمزور
کرن دے قدم بنا دَسے شامل سمجھاں گے۔

\begin{defn}[بے تضادی (Consistency)]
جملیاں دا سیٹ~$\Gamma$ \emph{متضاد (inconsistent)} اے اُدوں تے صرف
اُدوں جدوں کوئی محدود ذیلی سیٹ~$\Gamma_0 \subseteq \Gamma$ ایہو جہا ہووے
کہ $\Log{LK}$، $\Gamma_0 \Sequent \quad$ نوں !!{derive}s کرے۔ جے~$\Gamma$
متضاد نہیں، یعنی جے ہر محدود~$\Gamma_0 \subseteq \Gamma$ لئی $\Log{LK}$،
$\Gamma_0 \Sequent \quad$ نوں !!{derive} نہیں کردا، تاں اسیں اوہنوں
\emph{بے تضاد (consistent)} آکھدے آں۔
\end{defn}

\begin{prop}[انعکاسیت (Reflexivity)]
\ollabel{prop:reflexivity}
جے $!A \in \Gamma$ اے، تاں $\Gamma \Proves !A$۔
\end{prop}

\begin{proof}
ابتدائی سیکوئنٹ $!A \Sequent !A$ !!{derivable} اے، تے $\{!A\}
\subseteq \Gamma$۔
\end{proof}
  
\begin{prop}[یک رُخیت (Monotonicity)]
\ollabel{prop:monotonicity}
جے $\Gamma \subseteq \Delta$ تے $\Gamma \Proves !A$، تاں $\Delta
\Proves !A$۔
\end{prop}

\begin{proof}
منّ لو $\Gamma \Proves !A$، یعنی کوئی محدود~$\Gamma_0
\subseteq \Gamma$ ایہو جہا اے کہ $\Gamma_0 \Sequent !A$
!!{derivable} اے۔ چونکہ $\Gamma \subseteq \Delta$ اے، ایس لئی~$\Gamma_0$
،~$\Delta$ دا وی محدود ذیلی سیٹ اے۔ سو $\Gamma_0 \Sequent !A$ دی
!!{derivation} ایہ وی وکھاؤندی اے کہ $\Delta \Proves !A$۔
\end{proof}

\begin{prop}[تعدیت (Transitivity)]
\ollabel{prop:transitivity}
جے $\Gamma \Proves !A$ تے $\{!A\} \cup \Delta \Proves
!B$، تاں $\Gamma \cup \Delta \Proves !B$۔
\end{prop}

\begin{proof}
جے $\Gamma \Proves !A$ اے، تاں کوئی محدود~$\Gamma_0 \subseteq \Gamma$
تے کوئی !!a{derivation}~$\pi_0$، $\Gamma_0 \Sequent !A$ دی اے۔ جے
$\{!A\} \cup \Delta \Proves !B$ اے، تاں کسے محدود ذیلی سیٹ~$\Delta_0
\subseteq \Delta$ لئی کوئی !!a{derivation}~$\pi_1$، $!A, \Delta_0
\Sequent !B$ دی اے۔ ہِٹھلی !!{derivation} ویکھو:
\begin{prooftree}
  \AxiomC{}
  \RightLabel{$\pi_0$}
  \Deduce$\Gamma_0 \fCenter !A$
  \AxiomC{}
  \RightLabel{$\pi_1$}
  \Deduce$!A, \Delta_0 \fCenter !B$
  \RightLabel{\Cut}
  \BinaryInf$\Gamma_0, \Delta_0 \fCenter !B$
\end{prooftree}
چونکہ $\Gamma_0 \cup \Delta_0 \subseteq \Gamma \cup \Delta$ اے، ایس لئی
ایہ وکھاؤندا اے کہ $\Gamma \cup \Delta \Proves !B$۔
\end{proof}

دھیان دیو کہ ایس دا خاص مطلب ایہ اے کہ جے $\Gamma \Proves !A$ تے $!A
\Proves !B$، تاں $\Gamma \Proves !B$۔ ایہ وی نکلدا اے کہ جے $!A_1,
\dots, !A_n \Proves !B$ تے $\Gamma \Proves !A_i$ ہر~$i$ لئی، تاں
$\Gamma \Proves !B$۔

\begin{prop}
\ollabel{prop:incons}
$\Gamma$ اُدوں تے صرف اُدوں متضاد اے جدوں $\Gamma \Proves {!A}$، ہر
جملے~$!A$ لئی۔
\end{prop}

\begin{proof}
مشق۔
\end{proof}

\begin{prob}
\olref[fol][seq][ptn]{prop:incons} ثابت کرو۔
\end{prob}

\begin{prop}[محدود گواہی (Compactness)]
  \ollabel{prop:proves-compact}
  \begin{enumerate}
  \item جے $\Gamma \Proves !A$ اے، تاں کوئی محدود ذیلی سیٹ $\Gamma_0
    \subseteq \Gamma$ ایہو جہا اے کہ $\Gamma_0 \Proves !A$۔
  \item جے~$\Gamma$ دا ہر محدود ذیلی سیٹ بے تضاد اے، تاں~$\Gamma$ بے تضاد اے۔
  \end{enumerate}
\end{prop}

\begin{proof}
  \begin{enumerate}
    \item جے $\Gamma \Proves !A$ اے، تاں کوئی محدود ذیلی سیٹ
      $\Gamma_0 \subseteq \Gamma$ ایہو جہا اے کہ سیکوئنٹ $\Gamma_0
      \Sequent !A$ دی کوئی !!a{derivation} اے۔ ایس لئی $\Gamma_0
      \Proves !A$۔
    \item جے~$\Gamma$ متضاد اے، تاں کوئی محدود ذیلی سیٹ~$\Gamma_0
      \subseteq \Gamma$ ایہو جہا اے کہ $\Log{LK}$، $\Gamma_0 \Sequent
      \quad$ نوں !!{derive}s کردا اے۔ پر فیر~$\Gamma_0$،~$\Gamma$ دا اک
      متضاد محدود ذیلی سیٹ اے۔
  \end{enumerate}
\end{proof}
\endgroup
% /OLP-0077

% OLP-0078 content/first-order-logic/sequent-calculus/provability-consistency.tex
\begingroup
\def\olfilename{provability-consistency}
\expandafter\def\csname ol@sectioncs\endcsname{section}

\iftag{FOL}
      {\olfileid{fol}{seq}{prv}}
      {\olfileid{pl}{seq}{prv}}

\olsection{\usetoken{S}{derivability} تے بے تضادی \pnbLatin{(Derivability and Consistency)}}

ہُن اسیں !!{derivability} دے تعلق دیاں کئی خاصیتاں قائم کراں گے۔ ایہ
اپنے آپ وچ وی دلچسپ نیں، پر ہر خاصیت تمامیت دے قضیے دے ثبوت وچ اک کردار
نبھاوے گی۔

\begin{prop}\ollabel{prop:provability-contr}
  جے $\Gamma \Proves !A$ تے $\Gamma \cup \{!A\}$ متضاد اے، تاں
  $\Gamma$ متضاد اے۔
\end{prop}

\begin{proof}
ایہو جہے محدود~$\Gamma_0$ تے~$\Gamma_1 \subseteq \Gamma$ نیں کہ
$\Log{LK}$، $\Gamma_0 \Sequent !A$ تے $!A, \Gamma_1
\Sequent \quad$ نوں !!{derive}s کردا اے۔ پہلی~$\Log{LK}$-!!{derivation}،
یعنی $\Gamma_0 \Sequent !A$ والی، نوں~$\pi_0$ آکھو، تے دوجی
~$\Log{LK}$-!!{derivation}، یعنی $\Gamma_1, !A
\Sequent \quad$ والی، نوں~$\pi_1$ آکھو۔ فیر اسیں ایہ !!{derive} کر سکدے آں:
\begin{prooftree}
\AxiomC{}
\RightLabel{$\pi_0$}
\Deduce$ \Gamma_0 \fCenter !A $
\AxiomC{}
\RightLabel{$\pi_1$}
\Deduce$!A, \Gamma_1 \fCenter $
\RightLabel{\Cut}
\BinaryInf$ \Gamma_0,\Gamma_1 \fCenter $
\end{prooftree}

چونکہ $\Gamma_0 \subseteq \Gamma$ تے $\Gamma_1 \subseteq \Gamma$، سو
$\Gamma_0 \cup \Gamma_1 \subseteq \Gamma$؛ ایس لئی~$\Gamma$ متضاد اے۔
\end{proof}

\begin{prop}
\ollabel{prop:prov-incons}
$\Gamma \Proves !A$ اُدوں تے صرف اُدوں جدوں $\Gamma \cup \{\lnot !A\}$
متضاد اے۔
\end{prop}

\begin{proof}
% PNBSEQ-004: finite witnesses replace the source's unrestricted set-valued sequent antecedents.
پہلاں منّ لو $\Gamma \Proves !A$، یعنی کوئی محدود~$\Gamma_0
\subseteq \Gamma$ تے کوئی !!a{derivation}~$\pi_0$ ایہو جہے نیں کہ
$\Gamma_0 \Sequent !A$ اوہدا آخرلا سیکوئنٹ اے۔ اک
$\LeftR{\lnot}$ قاعدہ شامل کر کے، سانوں~$\lnot !A, \Gamma_0
\Sequent \quad$ دی !!a{derivation} مل جاندی اے؛ یعنی
$\Gamma \cup \{\lnot !A\}$ متضاد اے۔

جے $\Gamma \cup \{\lnot !A\}$ متضاد اے، تاں کوئی محدود~$\Gamma_1
\subseteq \Gamma$ تے کوئی !!a{derivation}~$\pi_1$ ایہو جہے نیں کہ
$\lnot !A, \Gamma_1 \Sequent \quad$ اوہدا آخرلا سیکوئنٹ اے۔ ہِٹھلی
!!a{derivation}، $\Gamma_1 \Sequent !A$ دی اے:
\begin{prooftree}
  \Axiom$!A \fCenter !A$
  \RightLabel{\RightR{\lnot}}
  \UnaryInf$\fCenter !A, \lnot !A$
  \AxiomC{}
  \RightLabel{$\pi_1$}
  \Deduce$\lnot !A, \Gamma_1 \fCenter$
  \RightLabel{\Cut}
  \BinaryInf$\Gamma_1 \fCenter !A$
\end{prooftree}
ایہ محدود گواہی مطلوبہ اخذ نوں قائم کردی اے۔
\end{proof}

\begin{prob}
ثابت کرو کہ $\Gamma \Proves \lnot !A$ اُدوں تے صرف اُدوں جدوں
$\Gamma \cup \{!A\}$ متضاد اے۔
\end{prob}

\begin{prop}\ollabel{prop:explicit-inc}
  جے $\Gamma \Proves !A$ تے $\lnot !A \in \Gamma$، تاں $\Gamma$ متضاد اے۔
\end{prop}

\begin{proof}
  منّ لو $\Gamma \Proves !A$ تے $\lnot !A \in \Gamma$۔ تاں کوئی
  !!a{derivation}~$\pi$ ایہو جہی اے جہدا سیکوئنٹ $\Gamma_0 \Sequent !A$ اے۔
  سیکوئنٹ $\lnot !A, \Gamma_0 \Sequent \quad$ وی !!{derivable} اے:
  \begin{prooftree}
    \AxiomC{}
    \RightLabel{$\pi$}
    \Deduce$\Gamma_0 \fCenter !A$
    \Axiom$!A \fCenter !A$
    \RightLabel{\LeftR{\lnot}}
    \UnaryInf$\lnot !A, !A \fCenter$
    \RightLabel{\LeftR{\Exchange}}
    \UnaryInf$!A, \lnot !A \fCenter$
    \RightLabel{\Cut}
    \BinaryInf$\Gamma_0, \lnot !A \fCenter$
  \end{prooftree}
  چونکہ $\lnot !A \in \Gamma$ تے $\Gamma_0 \subseteq \Gamma$، ایس توں
  وکھائی دیندا اے کہ $\Gamma$ متضاد اے۔
\end{proof}

\begin{prop}\ollabel{prop:provability-exhaustive}
  جے $\Gamma \cup \{!A\}$ تے $\Gamma \cup \{\lnot !A\}$ دونیں متضاد
  نیں، تاں $\Gamma$ متضاد اے۔
\end{prop}

\begin{proof}
ایہو جہے محدود سیٹ $\Gamma_0 \subseteq \Gamma$ تے $\Gamma_1
\subseteq \Gamma$، تے $\Log{LK}$-!!{derivation}s $\pi_0$ تے $\pi_1$
نیں، جیہڑے ترتیب وار $!A, \Gamma_0 \Sequent \quad$ تے $\lnot !A,
\Gamma_1 \Sequent \quad$ دے نیں۔ فیر اسیں ایہ !!{derive} کر سکدے آں:
\begin{prooftree}
\AxiomC{}
\RightLabel{$\pi_0$}
\Deduce$ !A, \Gamma_0 \fCenter $
\RightLabel{\RightR{\lnot}}
\UnaryInf$ \Gamma_0 \fCenter \lnot !A$
\AxiomC{}
\RightLabel{$\pi_1$}
\Deduce$\lnot !A, \Gamma_1 \fCenter  $
\RightLabel{\Cut}
\BinaryInf$ \Gamma_0, \Gamma_1 \fCenter $
\end{prooftree}
چونکہ $\Gamma_0 \subseteq \Gamma$ تے $\Gamma_1 \subseteq \Gamma$، سو
$\Gamma_0 \cup \Gamma_1 \subseteq \Gamma$۔ ایس لئی $\Gamma$ متضاد اے۔
\end{proof}
\endgroup
% /OLP-0078

% OLP-0079 content/first-order-logic/sequent-calculus/provability-propositional.tex
\begingroup
\def\olfilename{provability-propositional}
\expandafter\def\csname ol@sectioncs\endcsname{section}

\iftag{FOL}
      {\olfileid{fol}{seq}{ppr}}
      {\olfileid{pl}{seq}{ppr}}

\olsection{\usetoken{S}{derivability} تے قضیاتی رابطے \pnbLatin{(Derivability and the Propositional Connectives)}}

\begin{explain}
  اسیں قائم کردے آں کہ سیکوئنٹ حساب دا !!{derivability} تعلق~$\Proves$
  قضیاتی رابطیاں نال جڑے کجھ بنیادی حقائق قائم کرن لئی اینا طاقتور اے،
  جیویں کہ $!A \land !B \Proves !A$ تے $!A, !A \lif !B \Proves !B$
  (قاعدۂ اثباتِ مقدم؛ modus ponens)۔ تمامیت دے قضیے دے ثبوت لئی ایہناں
  حقائق دی لوڑ اے۔
\end{explain}

\begin{prop}\ollabel{prop:provability-land}
  \begin{enumerate}
  \item \ollabel{prop:provability-land-left} دونیں $!A \land !B \Proves
    !A$ تے $!A \land !B \Proves !B$۔
  \item \ollabel{prop:provability-land-right} $!A, !B \Proves !A \land
    !B$۔
  \end{enumerate}
\end{prop}

\begin{proof}
  \begin{enumerate}
  \item دونیں سیکوئنٹ $!A \land !B \Sequent !A$ تے $!A \land !B \Sequent
    !B$ !!{derivable} نیں:
    \begin{prooftree}
      \Axiom$!A \fCenter !A$
      \RightLabel{\LeftR{\land}}
      \UnaryInf$!A \land !B \fCenter !A$
      \DisplayProof\qquad\bottomAlignProof
      \Axiom$!B \fCenter !B$
      \RightLabel{\LeftR{\land}}
      \UnaryInf$!A \land !B \fCenter !B$
    \end{prooftree}
    \item سیکوئنٹ $!A, !B \Sequent !A \land !B$ دی !!a{derivation} ایہ اے:
    \begin{prooftree}
      \Axiom$!A \fCenter !A$
      \Axiom$!B \fCenter !B$
      \RightLabel{\RightR{\land}}
      \BinaryInf$!A, !B \fCenter !A \land !B$
    \end{prooftree}
  \end{enumerate}
\end{proof}

\begin{prop}\ollabel{prop:provability-lor}
  \begin{enumerate}
  \item $!A \lor !B, \lnot !A, \lnot !B$ متضاد اے۔
  \item دونیں $!A \Proves !A \lor !B$ تے $!B \Proves !A \lor !B$۔
  \end{enumerate}
\end{prop}

\begin{proof}
  \begin{enumerate}
  \item اسیں سیکوئنٹ $!A \lor !B, \lnot !A,
    \lnot !B \Sequent$ دی !!a{derivation} دیندے آں:
    \begin{prooftree}
      \Axiom$!A \fCenter !A$
      \RightLabel{\LeftR{\lnot}}
      \UnaryInf$\lnot !A, !A \fCenter$
      \doubleLine
      \UnaryInf$!A, \lnot !A, \lnot !B \fCenter$
      \Axiom$!B \fCenter !B$
      \RightLabel{\LeftR{\lnot}}
      \UnaryInf$\lnot !B, !B \fCenter$
      \doubleLine
      \UnaryInf$!B, \lnot !A, \lnot !B \fCenter$
      \RightLabel{\LeftR{\lor}}
      \BinaryInf$ !A \lor !B, \lnot !A, \lnot !B \fCenter $
    \end{prooftree}
    (چیتا رکھو کہ دُہریاں استنباطی لیکاں کمزور کرن، سکیڑن تے تبادلے دے
    کئی استنباطاں نوں دسدیاں نیں۔)
  \item دونیں سیکوئنٹ $!A \Sequent !A \lor !B$ تے $!B \Sequent !A
    \lor !B$ دیاں !!{derivation}s نیں:
    \begin{prooftree}
      \Axiom$!A \fCenter !A$
      \RightLabel{\RightR{\lor}}
      \UnaryInf$!A \fCenter !A \lor !B$
      \DisplayProof\qquad\bottomAlignProof
      \Axiom$!B \fCenter !B$
      \RightLabel{\RightR{\lor}}
      \UnaryInf$!B \fCenter !A \lor !B$
    \end{prooftree}
  \end{enumerate}
\end{proof}

\begin{prop}\ollabel{prop:provability-lif}
  \begin{enumerate}
  \item \ollabel{prop:provability-lif-left}  $!A, !A \lif !B \Proves !B$۔
  \item \ollabel{prop:provability-lif-right}
    دونیں $\lnot !A \Proves !A \lif !B$ تے $!B \Proves !A \lif !B$۔
  \end{enumerate}
\end{prop}

\begin{proof}
  \begin{enumerate}
    \item سیکوئنٹ $!A \lif !B, !A \Sequent !B$ !!{derivable} اے:
      \begin{prooftree}
        \Axiom$!A \fCenter !A$
        \Axiom$!B \fCenter !B$
        \RightLabel{\LeftR{\lif}}
        \BinaryInf$!A \lif !B, !A  \fCenter !B$
      \end{prooftree}
    \item دونیں سیکوئنٹ $\lnot !A \Sequent !A \lif !B$ تے $!B
      \Sequent !A \lif !B$ !!{derivable} نیں:
      \begin{prooftree}
        \Axiom$!A \fCenter !A$
        \RightLabel{\LeftR{\lnot}}
        \UnaryInf$\lnot !A, !A \fCenter$
        \RightLabel{\LeftR{\Exchange}}
        \UnaryInf$!A, \lnot !A \fCenter$
        \RightLabel{\RightR{\Weakening}}
        \UnaryInf$!A, \lnot !A \fCenter !B$
        \RightLabel{\RightR{\lif}}
        \UnaryInf$\lnot !A \fCenter !A \lif !B$
        \DisplayProof\qquad\bottomAlignProof
        \Axiom$!B \fCenter !B$
        \RightLabel{\LeftR{\Weakening}}
        \UnaryInf$!A, !B \fCenter !B$
        \RightLabel{\RightR{\lif}}
        \UnaryInf$!B \fCenter !A \lif !B$
      \end{prooftree}
  \end{enumerate}
\end{proof}
\endgroup
% /OLP-0079

% OLP-0080 content/first-order-logic/sequent-calculus/provability-quantifiers.tex
\begingroup
\def\olfilename{provability-quantifiers}
\expandafter\def\csname ol@sectioncs\endcsname{section}

\olfileid{fol}{seq}{qpr}
\olsection{\usetoken{S}{derivability} تے مسوّر \pnbLatin{(Derivability and the Quantifiers)}}

\begin{explain}
  تمامیت دے قضیے دا تقاضا ایہ وی اے کہ سیکوئنٹ حساب دے قاعدے~$\Proves$
  بارے اوہ حقائق دیندے ہون جیہڑے ایس حصے وچ قائم کیتے گئے نیں۔
\end{explain}

\begin{thm}
\ollabel{thm:strong-generalization} جے $c$ اک ایہو جہا مستقل (constant) اے جیہڑا
$\Gamma$ یا $!A(x)$ وچ نہیں آؤندا تے $\Gamma \Proves !A(c)$، تاں $\Gamma
\Proves \lforall[x][!A(x)]$۔
\end{thm}

\begin{proof}
$\pi_0$، $\Log{LK}$ وچ $\Gamma_0 \Sequent !A(c)$ دی اک !!{derivation} ہووے، جتھے
$\Gamma_0 \subseteq \Gamma$ اک محدود ذیلی سیٹ اے۔ اک $\RightR{\lforall}$ استنباط جوڑ کے سانوں
$\Gamma_0 \Sequent \lforall[x][!A(x)]$ دی !!a{derivation} ملدی اے، کیوں جے $c$،
$\Gamma$ یا $!A(x)$ وچ نہیں آؤندا تے ایس لئی ایگن متغیر دی شرط
\pnbLatin{(eigenvariable condition)} پوری ہُندی اے۔
\end{proof}

\begin{prop}
\ollabel{prop:provability-quantifiers}
\begin{tagenumerate}{prvEx,prvAll}
\tagitem{prvEx}{$!A(t) \Proves \lexists[x][!A(x)]$.}{}
\tagitem{prvAll}{$\lforall[x][!A(x)] \Proves !A(t)$.}{}
\end{tagenumerate}
\end{prop}

\begin{proof}
\begin{tagenumerate}{prvEx,prvAll}
\tagitem{prvEx}{سیکوئنٹ $!A(t) \Sequent \lexists[x][!A(x)]$
  !!{derivable} اے:
  \begin{prooftree}
    \Axiom$!A(t) \fCenter !A(t)$
    \RightLabel{\RightR{\lexists}}
    \UnaryInf$!A(t) \fCenter \lexists[x][!A(x)]$
    \end{prooftree}}{}
\tagitem{prvAll}{سیکوئنٹ $\lforall[x][!A(x)] \Sequent !A(t)$
  !!{derivable} اے:
  \begin{prooftree}
    \Axiom$!A(t) \fCenter !A(t)$
    \RightLabel{\LeftR{\lforall}}
    \UnaryInf$\lforall[x][!A(x)] \fCenter !A(t)$
    \end{prooftree}}{}
\end{tagenumerate}
\end{proof}
\endgroup
% /OLP-0080

% OLP-0081 content/first-order-logic/sequent-calculus/soundness.tex
\begingroup
\def\olfilename{soundness}
\expandafter\def\csname ol@sectioncs\endcsname{section}

\iftag{FOL}
      {\olfileid{fol}{seq}{sou}}
      {\olfileid{pl}{seq}{sou}}

\olsection{صحت مندی \pnbLatin{(Soundness)}}

\begin{explain}
!!^a{derivation} نظام، جیویں سیکوئنٹ حساب \pnbLatin{(sequent calculus)}، 
\emph{صحت مند (sound)} اے جے اوہ اوہ گلاں !!{derive} نہیں کر سکدا جیہڑیاں
اصل وچ قائم نہیں۔ ایس لئی صحت مندی، !!{derivation} نظاماں دی اک طرح دی
ضمانت شدہ حفاظتی خاصیت \pnbLatin{(guaranteed safety property)} اے۔ جیہڑی
ثبوتی-نظریاتی خاصیت دی گل ہو رہی ہووے، اوہدے مطابق اسیں مثال لئی ایہ جاننا
چاہندے آں کہ
\begin{enumerate}
\item ہر !!{derivable}~$!A$، \iftag{FOL}{معتبر}{ہمیشہ صادق فارمولا} اے؛
\item جے !!a{sentence} کجھ ہور جملیاں توں !!{derivable} اے، تاں اوہ اوہناں دا
  منطقی نتیجہ وی اے؛
\item جے !!{sentence}s دا کوئی سیٹ متضاد اے، تاں اوہ پورا ہون جوگ نہیں۔
\end{enumerate}
ایہ !!a{derivation} نظام دیاں اہم خاصیتاں نیں۔ جے ایہناں وچوں کوئی وی
قائم نہ ہووے، تاں !!{derivation} نظام ناقص اے---اوہ لوڑ توں ودھ
!!{derive} کرے گا۔ ایس کرکے !!{derivation} نظام دی صحت مندی قائم کرنا
انتہائی اہم اے۔

چونکہ ایہ ساریاں ثبوتی-نظریاتی خاصیتاں سیکوئنٹ حساب وچ کجھ سیکوئنٹاں دی
!!{derivability} راہیں تعریف کیتیاں گئیاں نیں، ایس لئی اُپرلے (1)--(3)
ثابت کرن لئی !!{derivable} سیکوئنٹاں دیاں معنوی خاصیتاں بارے کجھ ثابت
کرنا پیندا اے۔ اسیں پہلاں ایہ تعریف کراں گے کہ کسے سیکوئنٹ دے
\emph{معتبر (valid)} ہون دا کی مطلب اے، تے فیر وکھاواں گے کہ ہر
!!{derivable} سیکوئنٹ معتبر اے۔ فیر (1)--(3) ایس نتیجے توں ذیلی قضیاں
دے طور تے نکل آون گے۔
\end{explain}

\begin{defn}
\iftag{FOL}{!!^a{structure}~$\Struct
  M$}{!!^a{valuation}~$\pAssign{v}$} سیکوئنٹ
$\Gamma \Sequent \Delta$ نوں \emph{پورا کردا اے (satisfies)} اُدوں تے صرف اُدوں جدوں یا تاں
$\iftag{FOL}{\Sat/{M}}{\pSat/{v}}{!A}$ کسے $!A \in \Gamma$ لئی ہووے یا
$\iftag{FOL}{\Sat{M}}{\pSat{v}}{!A}$ کسے $!A \in \Delta$ لئی ہووے۔

کوئی سیکوئنٹ \emph{معتبر (valid)} اے اُدوں تے صرف اُدوں جدوں ہر
\iftag{FOL}{!!{structure}~$\Struct M$}{!!{valuation}~$\pAssign{v}$}
اوہنوں پورا کردا ہووے۔
\end{defn}

\begin{thm}[صحت مندی (Soundness)]
\ollabel{thm:sequent-soundness} جے $\Log{LK}$، $\Theta
\Sequent \Xi$ نوں !!{derive}s کردا اے، تاں $\Theta \Sequent \Xi$ معتبر اے۔
\end{thm}

\begin{proof}
$\pi$، $\Theta \Sequent \Xi$ دی !!a{derivation} ہووے۔ اسیں استنباطاں دی
گنتی~$n$ اُتے استقرا کردے آں، جیہڑی~$\pi$ وچ اے۔

جے استنباطاں دی گنتی~$0$ اے، تاں~$\pi$ صرف اک ابتدائی سیکوئنٹ اُتے مشتمل
اے۔ ہر ابتدائی سیکوئنٹ $!A \Sequent !A$ ظاہر اے معتبر اے، کیوں جے ہر
\iftag{FOL}{$\Struct M$}{$\pAssign{v}$} لئی یا تاں
$\iftag{FOL}{\Sat/{M}}{\pSat/{v}}{!A}$ یا
$\iftag{FOL}{\Sat{M}}{\pSat{v}}{!A}$۔

جے استنباطاں دی گنتی~0 توں ودھ اے، تاں اسیں سب توں ہِٹھلے استنباط دی
قسم مطابق صورتاں وکھ کردے آں۔ استقرائی مفروضے کرکے اسیں من سکدے آں کہ
اوس استنباط دے مقدمے معتبر نیں، کیوں جے کسے وی مقدمے دی !!{derivation}
وچ استنباطاں دی گنتی~$n$ توں گھٹ اے۔

پہلاں اسیں صرف اک مقدمے والے ممکن استنباطاں اُتے غور کردے آں۔
\begin{enumerate}
\item آخرلا استنباط کمزور کرنا (weakening) اے۔ فیر $\Theta \Sequent \Xi$
  یا تاں $!A, \Gamma \Sequent \Delta$ اے (جے آخرلا استنباط
  \LeftR{\Weakening} اے) یا $\Gamma \Sequent \Delta, !A$ اے (جے اوہ
  \RightR{\Weakening} اے)، تے !!{derivation} ایہناں وچوں کسے اک اُتے
  ختم ہُندی اے:
  \begin{prooftree}
    \AxiomC{}
    \Deduce$\Gamma \fCenter \Delta$
    \RightLabel{\LeftR{\Weakening}}
    \UnaryInf$!A, \Gamma \fCenter \Delta$
    \DisplayProof\qquad\bottomAlignProof
    \AxiomC{}
    \Deduce$\Gamma \fCenter \Delta$
    \RightLabel{\RightR{\Weakening}}
    \UnaryInf$\Gamma \fCenter \Delta, !A$
  \end{prooftree}
  استقرائی مفروضے مطابق $\Gamma \Sequent \Delta$ معتبر اے، یعنی ہر
  \iftag{FOL}{!!{structure}~$\Struct{M}$}{!!{valuation}~$\pAssign{v}$}
  لئی یا تاں کوئی $!C \in \Gamma$ اے جیہدے لئی
  $\iftag{FOL}{\Sat/{M}}{\pSat/{v}}{!C}$، یا کوئی $!C \in
  \Delta$ اے جیہدے لئی $\iftag{FOL}{\Sat{M}}{\pSat{v}}{!C}$۔
  
  جے $\iftag{FOL}{\Sat/{M}}{\pSat/{v}}{!C}$ کسے $!C \in \Gamma$ لئی
  ہووے، تاں $!C \in \Theta$ وی اے: کھبے کمزور کرنے دے قاعدے وچ $\Theta = !A, \Gamma$،
  تے سجے کمزور کرنے دے قاعدے وچ نتیجے دا کھبا پاسا مقدمے دے کھبے پاسے
  ورگا ای رہندا اے۔ دوناں صورتاں وچ مقدمے دے کھبے پاسے دا ہر فارمولا
  نتیجے دے کھبے پاسے وچ رہندا اے۔ سو
  $\iftag{FOL}{\Sat/{M}}{\pSat/{v}}{!C}$ کسے $!C \in \Theta$ لئی اے۔
  ایسے طرح، جے $\iftag{FOL}{\Sat{M}}{\pSat{v}}{!C}$ کسے $!C \in
  \Delta$ لئی اے، تاں $!C \in \Xi$ ہون کرکے
  $\iftag{FOL}{\Sat{M}}{\pSat{v}}{!C}$ کسے $!C \in \Xi$ لئی اے۔ سو
  $\Theta \Sequent \Xi$ معتبر اے۔
\item آخرلا استنباط \LeftR{\lnot} اے۔ فیر آخرلے استنباط دا مقدمہ
  $\Gamma \Sequent \Delta, !A$ تے نتیجہ $\lnot !A, \Gamma \Sequent
  \Delta$ اے، یعنی !!{derivation} دا آخر ایہ اے:
  \begin{prooftree}
    \AxiomC{}
    \Deduce$\Gamma \fCenter \Delta, !A$
    \RightLabel{\LeftR{\lnot}}
    \UnaryInf$\lnot !A, \Gamma \fCenter \Delta$
  \end{prooftree}
  تے $\Theta = \lnot !A, \Gamma$، جدکہ $\Xi = \Delta$۔

  استقرائی مفروضہ سانوں دَسدا اے کہ $\Gamma \Sequent \Delta, !A$
  معتبر اے، یعنی ہر \iftag{FOL}{$\Struct{M}$}{$\pAssign{v}$} لئی یا
  (a) کسے $!C \in \Gamma$ لئی
  $\iftag{FOL}{\Sat/{M}}{\pSat/{v}}{!C}$، یا (b) کسے $!C \in
  \Delta$ لئی $\iftag{FOL}{\Sat{M}}{\pSat{v}}{!C}$، یا (c)
  $\iftag{FOL}{\Sat{M}}{\pSat{v}}{!A}$۔ اسیں وکھانا چاہندے آں کہ
  $\Theta \Sequent \Xi$ وی معتبر اے۔ من لو
  \iftag{FOL}{$\Struct{M}$ اک !!a{structure} اے}{$\pAssign{v}$ اک
    !!a{valuation} اے}۔ جے (a) قائم اے، تاں کوئی $!C \in \Gamma$ اے
  جیہدے لئی $\iftag{FOL}{\Sat/{M}}{\pSat/{v}}{!C}$، پر $!C \in
  \Theta$ وی اے۔ جے (b) قائم اے، تاں کوئی $!C \in \Delta$ اے جیہدے
  لئی $\iftag{FOL}{\Sat{M}}{\pSat{v}}{!C}$، پر $!C \in \Xi$ وی اے۔
  آخر وچ، جے $\iftag{FOL}{\Sat{M}}{\pSat{v}}{!A}$، تاں
  $\iftag{FOL}{\Sat/{M}}{\pSat/{v}}{\lnot !A}$۔ چونکہ $\lnot !A \in
  \Theta$، ایس لئی کوئی $!C \in \Theta$ اے جیہدے لئی
  $\iftag{FOL}{\Sat/{M}}{\pSat/{v}}{!C}$۔ سو $\Theta
  \Sequent \Xi$ معتبر اے۔
\item آخرلا استنباط \RightR{\lnot} اے: مشق۔
\item آخرلا استنباط \LeftR{\land} اے۔ ایس دے دو روپ نیں: $!A \land !B$
  مقدمے دے کھبے پاسے موجود $!A$ یا $!B$ توں کڈھیا جا سکدا اے۔
  پہلی صورت وچ~$\pi$ دا آخر ایہ اے:
  \begin{prooftree}
    \AxiomC{}
    \Deduce$!A, \Gamma \fCenter \Delta$
    \RightLabel{\LeftR{\land}}
    \UnaryInf$!A \land !B, \Gamma \fCenter \Delta$
  \end{prooftree}
  تے $\Theta = !A \land !B, \Gamma$، جدکہ $\Xi = \Delta$۔ من لو
  \iftag{FOL}{!!a{structure}~$\Struct
    M$}{!!a{valuation}~$\pAssign{v}$}۔ استقرائی مفروضے مطابق
  $!A, \Gamma \Sequent \Delta$ معتبر اے، سو یا (a)
  $\iftag{FOL}{\Sat/{M}}{\pSat/{v}}{!A}$، (b)
  $\iftag{FOL}{\Sat/{M}}{\pSat/{v}}{!C}$ کسے $!C \in \Gamma$ لئی، یا
  (c)~$\iftag{FOL}{\Sat{M}}{\pSat{v}}{!C}$ کسے $!C \in \Delta$ لئی۔ صورت (a) وچ
  $\iftag{FOL}{\Sat/{M}}{\pSat/{v}}{!A \land !B}$، سو کوئی $!C \in
  \Theta$ (یعنی $!A \land !B$) اے جیہدے لئی
  $\iftag{FOL}{\Sat/{M}}{\pSat/{v}}{!C}$۔ صورت (b) وچ کوئی $!C
  \in \Gamma$ اے جیہدے لئی $\iftag{FOL}{\Sat/{M}}{\pSat/{v}}{!C}$، تے
  $!C \in \Theta$ وی اے۔ صورت (c) وچ کوئی $!C \in \Delta$ اے جیہدے
  لئی $\iftag{FOL}{\Sat{M}}{\pSat{v}}{!C}$، تے $!C \in \Xi$ وی اے،
  کیوں جے $\Xi = \Delta$۔ سو ہر صورت وچ
  \iftag{FOL}{$\Struct M$}{$\pAssign{v}$}، $!A \land !B, \Gamma
  \Sequent \Delta$ نوں پورا کردا اے۔
  % PNBSEQ-005: frozen source concludes Gamma -> Delta, but this rule case proves the displayed conclusion A and B, Gamma -> Delta.
  چونکہ \iftag{FOL}{$\Struct{M}$}{$\pAssign{v}$} من مانا سی،
  $!A \land !B, \Gamma \Sequent \Delta$ معتبر اے۔ $!A \land !B$ نوں
  $!B$ توں کڈھن والی صورت وی ایسے طرح، $!A$ نوں $!B$ نال بدل کے،
  سنبھالی جاندی اے۔
\item آخرلا استنباط \RightR{\lor} اے۔ ایس دے دو روپ نیں: $!A \lor !B$
  مقدمے دے سجے پاسے موجود $!A$ یا $!B$ توں کڈھیا جا سکدا اے۔
  پہلی صورت وچ~$\pi$ دا آخر ایہ اے:
  \begin{prooftree}
    \AxiomC{}
    \Deduce$\Gamma \fCenter \Delta, !A$
    \RightLabel{\RightR{\lor}}
    \UnaryInf$\Gamma \fCenter \Delta, !A \lor !B$
  \end{prooftree}
  ہُن $\Theta = \Gamma$ تے $\Xi = \Delta, !A \lor !B$۔ من لو
  \iftag{FOL}{!!a{structure}~$\Struct{M}$}{!!a{valuation}~$\pAssign{v}$}۔
  چونکہ $\Gamma \Sequent \Delta, !A$ معتبر اے، سو یا (a)
  $\iftag{FOL}{\Sat{M}}{\pSat{v}}{!A}$، (b)
  $\iftag{FOL}{\Sat/{M}}{\pSat/{v}}{!C}$ کسے $!C \in \Gamma$ لئی، یا
  (c)~$\iftag{FOL}{\Sat{M}}{\pSat{v}}{!C}$ کسے $!C \in \Delta$ لئی۔ صورت (a) وچ
  $\iftag{FOL}{\Sat{M}}{\pSat{v}}{!A \lor !B}$۔ صورت (b) وچ کوئی
  $!C \in \Gamma$ اے جیہدے لئی
  $\iftag{FOL}{\Sat/{M}}{\pSat/{v}}{!C}$۔ صورت (c) وچ کوئی $!C
  \in \Delta$ اے جیہدے لئی $\iftag{FOL}{\Sat{M}}{\pSat{v}}{!C}$۔ سو
  ہر صورت وچ \iftag{FOL}{$\Struct{M}$}{$\pAssign{v}$}،
  $\Gamma \Sequent \Delta, !A \lor !B$، یعنی $\Theta \Sequent \Xi$،
  نوں پورا کردا اے۔ چونکہ \iftag{FOL}{$\Struct{M}$}{$\pAssign{v}$}
  من مانا سی، $\Theta \Sequent \Xi$ معتبر اے۔ $!A \lor !B$ نوں $!B$
  توں کڈھن والی صورت وی ایسے طرح، $!A$ نوں $!B$ نال بدل کے، سنبھالی جاندی اے۔
\item آخرلا استنباط \RightR{\lif} اے۔ فیر~$\pi$ دا آخر ایہ اے:
  \begin{prooftree}
    \AxiomC{}
    \Deduce$!A, \Gamma \fCenter \Delta, !B$
    \RightLabel{\RightR{\lif}}
    \UnaryInf$\Gamma \fCenter \Delta, !A \lif !B$
  \end{prooftree}
  فیر، استقرائی مفروضہ دَسدا اے کہ مقدمہ معتبر اے؛ اسیں وکھانا چاہندے آں
  کہ نتیجہ وی معتبر اے۔ من لو
  \iftag{FOL}{$\Struct{M}$}{$\pAssign{v}$} من مانا اے۔ چونکہ $!A,
  \Gamma \Sequent \Delta, !B$ معتبر اے، ہِٹھلیاں صورتاں وچوں گھٹّو گھٹ
  اک قائم اے: (a) $\iftag{FOL}{\Sat/{M}}{\pSat/{v}}{!A}$، (b)
  $\iftag{FOL}{\Sat{M}}{\pSat{v}}{!B}$، (c)
  $\iftag{FOL}{\Sat/{M}}{\pSat/{v}}{!C}$ کسے~$~!C \in \Gamma$ لئی، یا
  (d) $\iftag{FOL}{\Sat{M}}{\pSat{v}}{!C}$ کسے $!C \in \Delta$ لئی۔ صورتاں (a) تے (b) وچ
  $\iftag{FOL}{\Sat{M}}{\pSat{v}}{!A \lif !B}$، سو کوئی $!C \in
  \Delta, !A \lif !B$ اے جیہدے لئی
  $\iftag{FOL}{\Sat{M}}{\pSat{v}}{!C}$۔ صورت (c) وچ کسے $!C \in
  \Gamma$ لئی $\iftag{FOL}{\Sat/{M}}{\pSat/{v}}{!C}$۔ صورت (d) وچ
  کسے $!C \in \Delta$ لئی $\iftag{FOL}{\Sat{M}}{\pSat{v}}{!C}$۔ ہر
  صورت وچ \iftag{FOL}{$\Struct{M}$}{$\pAssign{v}$}، $\Gamma
  \Sequent \Delta, !A \lif !B$ نوں پورا کردا اے۔ چونکہ
  \iftag{FOL}{$\Struct{M}$}{$\pAssign{v}$} من مانا سی، $\Gamma
  \Sequent \Delta, !A \lif !B$ معتبر اے۔ \iftag{FOL}{%
\item آخرلا استنباط \LeftR{\lforall} اے۔ فیر کوئی !!a{formula}~$!A(x)$
  تے کوئی بند ٹرم~$t$ ایہو جہا اے کہ~$\pi$ دا آخر ایہ اے:
  \begin{prooftree}
    \AxiomC{}
    \Deduce$!A(t), \Gamma \fCenter \Delta$
    \RightLabel{\LeftR{\lforall}}
    \UnaryInf$\lforall[x][!A(x)], \Gamma \fCenter \Delta$
  \end{prooftree}
  اسیں وکھانا چاہندے آں کہ نتیجہ $\lforall[x][!A(x)], \Gamma
  \Sequent \Delta$ معتبر اے۔ من لو !!a{structure}~$\Struct M$۔ چونکہ
  مقدمہ $!A(t), \Gamma \Sequent \Delta$ معتبر اے، سو (a)
  $\Sat/{M}{!A(t)}$، (b) $\Sat/{M}{!C}$ کسے $!C \in \Gamma$ لئی، یا
  (c)~$\Sat{M}{!C}$ کسے $!C \in \Delta$ لئی۔ صورت (a) وچ
  \olref[syn][sem]{prop:quant-terms} مطابق، جے
  $\Sat{M}{\lforall[x][!A(x)]}$ تاں $\Sat{M}{!A(t)}$۔ چونکہ
  $\Sat/{M}{!A(t)}$، ایس لئی $\Sat/{M}{\lforall[x][!A(x)]}$۔ صورتاں (b)
  تے (c) وچ $\Struct{M}$، $\lforall[x][!A(x)], \Gamma
  \Sequent \Delta$ نوں وی پورا کردا اے۔ چونکہ $\Struct M$ من مانا سی،
  $\lforall[x][!A(x)], \Gamma \Sequent \Delta$ معتبر اے۔
\item آخرلا استنباط \RightR{\lexists} اے: مشق۔
\item آخرلا استنباط \RightR{\lforall} اے۔ فیر کوئی !!a{formula}~$!A(x)$
  تے کوئی !!a{constant}~$a$ ایہو جہا اے کہ~$\pi$ دا آخر ایہ اے:
  \begin{prooftree}
    \AxiomC{}
    \Deduce$\Gamma \fCenter \Delta, !A(a)$
    \RightLabel{\RightR{\lforall}}
    \UnaryInf$\Gamma \fCenter \Delta, \lforall[x][!A(x)]$
  \end{prooftree}
  جتھے ایگن متغیر دی شرط \pnbLatin{(eigenvariable condition)} پوری ہُندی اے، یعنی
  $a$، $!A(x)$، $\Gamma$ یا $\Delta$ وچ نہیں آؤندا۔ استقرائی مفروضے
  مطابق آخرلے استنباط دا مقدمہ معتبر اے۔ سانوں وکھانا اے کہ نتیجہ وی معتبر
  اے، یعنی ہر !!{structure}~$\Struct M$ لئی یا (a)
  $\Sat{M}{\lforall[x][!A(x)]}$، (b) $\Sat/{M}{!C}$ کسے $!C \in \Gamma$
  لئی، یا (c)~$\Sat{M}{!C}$ کسے $!C \in \Delta$ لئی۔

  من لو $\Struct{M}$ اک من مانا !!{structure} اے۔ جے (b) یا (c) قائم اے
  تاں کم ہو گیا، سو من لو دوناں وچوں کوئی وی قائم نہیں: ہر $!C \in
  \Gamma$ لئی $\Sat{M}{!C}$، تے ہر $!C \in \Delta$ لئی
  $\Sat/{M}{!C}$۔ سانوں وکھانا اے کہ (a) قائم اے، یعنی
  $\Sat{M}{\lforall[x][!A(x)]}$۔
  \olref[syn][ass]{prop:sat-quant} مطابق ایہ وکھانا کافی اے کہ
  $\Sat{M}{!A(x)}[s]$ ہر متغیر تعیّن~$s$ لئی۔ سو $s$ نوں کوئی من مانا متغیر تعیّن
  من لو۔ اوس ساخت~$\Struct{M'}$ اُتے غور کرو جیہڑی~$\Struct{M}$ ورگی اے،
  پر $\Assign{a}{M'} = s(x)$۔
  \olref[syn][ext]{cor:extensionality-sent} مطابق ہر $!C \in \Gamma$ لئی
  $\Sat{M'}{!C}$، کیوں جے $a$،~$\Gamma$ وچ نہیں آؤندا، تے ہر $!C \in
  \Delta$ لئی $\Sat/{M'}{!C}$۔ پر مقدمہ معتبر اے، سو
  $\Sat{M'}{!A(a)}$۔ \olref[syn][ass]{prop:sentence-sat-true} مطابق
  $\Sat{M'}{!A(a)}[s]$، کیوں جے $!A(a)$ اک جملہ اے۔ ہُن
  $\varAssign{s}{s}{x}$ تے $s(x) = \Value{a}{M'}[s]$، کیوں جے اسیں
  $\Struct{M'}$ دی تعریف ایویں ای کیتی اے۔ سو
  \olref[syn][ext]{prop:ext-formulas} لاگو ہُندا اے تے سانوں
  $\Sat{M'}{!A(x)}[s]$ ملدا اے۔ چونکہ $a$،~$!A(x)$ وچ نہیں آؤندا،
  \olref[syn][ext]{prop:extensionality} مطابق $\Sat{M}{!A(x)}[s]$۔ چونکہ
  $s$ من مانا سی، اسیں ہر متغیر تعیّن لئی $\Sat{M}{!A(x)}[s]$ ثابت کر
  دِتّا اے۔
\item آخرلا استنباط \LeftR{\lexists} اے: مشق۔
}{}
\end{enumerate}
ہُن دو مقدمیاں والے ممکن استنباطاں اُتے غور کریے۔
\begin{enumerate}
\item آخرلا استنباط کٹ (cut) اے؛ فیر~$\pi$ دا آخر ایہ اے:
  \begin{prooftree}
    \AxiomC{}
    \Deduce$\Gamma \fCenter \Delta, !A$
    \AxiomC{}
    \Deduce$!A, \Pi \fCenter \Lambda$
    \RightLabel{\Cut}
    \BinaryInf$\Gamma, \Pi \fCenter \Delta, \Lambda$
  \end{prooftree}
  من لو \iftag{FOL}{$\Struct{M}$ اک !!a{structure} اے}{$\pAssign{v}$ اک
    !!a{valuation} اے}۔ استقرائی مفروضے مطابق مقدمے معتبر نیں، سو
  \iftag{FOL}{$\Struct{M}$}{$\pAssign{v}$} دوناں مقدمیاں نوں پورا کردا اے۔
  اسیں دو صورتاں وکھ کردے آں: (a)
  $\iftag{FOL}{\Sat/{M}}{\pSat/{v}}{!A}$ تے (b)
  $\iftag{FOL}{\Sat{M}}{\pSat{v}}{!A}$۔ صورت (a) وچ، کھبے مقدمے نوں پورا
  کرن لئی \iftag{FOL}{$\Struct{M}$}{$\pAssign{v}$} نوں $\Gamma
  \Sequent \Delta$ پورا کرنا پئے گا۔ پر فیر اوہ نتیجے نوں وی پورا کردا اے۔
  صورت (b) وچ، سجے مقدمے نوں پورا کرن لئی
  % PNBSEQ-006: frozen source prints Pi setminus Lambda; the right Cut premise requires the sequent Pi -> Lambda.
  \iftag{FOL}{$\Struct{M}$}{$\pAssign{v}$} نوں $\Pi \Sequent \Lambda$
  پورا کرنا پئے گا۔ فیر وی \iftag{FOL}{$\Struct{M}$}{$\pAssign{v}$}
  نتیجے نوں پورا کردا اے۔
\item آخرلا استنباط \RightR{\land} اے۔ فیر~$\pi$ دا آخر ایہ اے:
  \begin{prooftree}
    \AxiomC{}
    \Deduce$\Gamma \fCenter \Delta, !A$
    \AxiomC{}
    \Deduce$\Gamma \fCenter \Delta, !B$
    \RightLabel{\RightR{\land}}
    \BinaryInf$\Gamma \fCenter \Delta, !A \land !B$
  \end{prooftree}
  من لو \iftag{FOL}{!!a{structure}~$\Struct
    M$}{!!a{valuation}~$\pAssign{v}$}۔ جے
  \iftag{FOL}{$\Struct{M}$}{$\pAssign{v}$}، $\Gamma \Sequent
  \Delta$ نوں پورا کردا اے، تاں کم ہو گیا۔ سو من لو اوہ نہیں کردا۔
  % PNBSEQ-007: frozen prose uses proof-tree-only fCenter here; the valid premise is the sequent Gamma -> Delta, A.
  چونکہ $\Gamma \Sequent \Delta, !A$ استقرائی مفروضے مطابق معتبر اے،
  ایس لئی $\iftag{FOL}{\Sat{M}}{\pSat{v}}{!A}$۔ ایسے طرح، چونکہ $\Gamma
  \Sequent \Delta, !B$ معتبر اے، ایس لئی
  $\iftag{FOL}{\Sat{M}}{\pSat{v}}{!B}$۔ پر فیر
  $\iftag{FOL}{\Sat{M}}{\pSat{v}}{!A \land !B}$۔
\item آخرلا استنباط \LeftR{\lor} اے: مشق۔
\item آخرلا استنباط \LeftR{\lif} اے۔ فیر~$\pi$ دا آخر ایہ اے:
  \begin{prooftree}
    \AxiomC{}
    \Deduce$\Gamma \fCenter \Delta, !A$
    \AxiomC{}
    \Deduce$!B, \Pi \fCenter \Lambda$
    \RightLabel{\LeftR{\lif}}
    \BinaryInf$!A \lif !B, \Gamma, \Pi \fCenter \Delta, \Lambda$
  \end{prooftree}
  فیر، من لو \iftag{FOL}{!!a{structure}~$\Struct{M}$}{!!a{valuation}~$\pAssign{v}$}
  تے فرض کرو کہ \iftag{FOL}{$\Struct{M}$}{$\pAssign{v}$}،
  $\Gamma, \Pi \Sequent \Delta, \Lambda$ نوں پورا نہیں کردا۔ سانوں
  وکھانا اے کہ $\iftag{FOL}{\Sat/{M}}{\pSat/{v}}{!A \lif !B}$۔ جے
  \iftag{FOL}{$\Struct{M}$}{$\pAssign{v}$}، $\Gamma, \Pi \Sequent
  \Delta, \Lambda$ نوں پورا نہیں کردا، تاں اوہ نہ $\Gamma \Sequent
  \Delta$ نوں تے نہ $\Pi \Sequent \Lambda$ نوں پورا کردا اے۔ چونکہ
  $\Gamma \Sequent \Delta, !A$ معتبر اے، ایس لئی
  $\iftag{FOL}{\Sat{M}}{\pSat{v}}{!A}$۔ چونکہ $!B, \Pi \Sequent
  \Lambda$ معتبر اے، ایس لئی $\iftag{FOL}{\Sat/{M}}{\pSat/{v}}{!B}$۔
  پر فیر $\iftag{FOL}{\Sat/{M}}{\pSat/{v}}{!A \lif !B}$، تے ایہی اسیں
  وکھانا چاہندے ساں۔
\end{enumerate}
\end{proof}

\tagprob{FOL}
\begin{prob}
\olref[fol][seq][sou]{thm:sequent-soundness} دا ثبوت پورا کرو۔
\end{prob}
\tagendprob

\tagprob{notFOL}
\begin{prob}
\olref[pl][seq][sou]{thm:sequent-soundness} دا ثبوت پورا کرو۔
\end{prob}
\tagendprob

\begin{cor}
\ollabel{cor:weak-soundness}
جے $\Proves !A$، تاں $!A$، \iftag{FOL}{معتبر}{ہمیشہ صادق فارمولا} اے۔
\end{cor}

\begin{cor}
\ollabel{cor:entailment-soundness}
جے $\Gamma \Proves !A$، تاں $\Gamma \Entails !A$۔
\end{cor}

\begin{proof}
جے $\Gamma \Proves !A$، تاں کسے محدود ذیلی سیٹ $\Gamma_0 \subseteq
\Gamma$ لئی $\Gamma_0 \Sequent !A$ دی !!a{derivation} اے۔
\olref{thm:sequent-soundness} مطابق ہر
\iftag{FOL}{!!{structure}~$\Struct{M}$}{!!{valuation}~$\pAssign{v}$}
یا تاں کسے $!B \in \Gamma_0$ نوں کاذب بناؤندا اے یا $!A$ نوں صادق۔ سو
جے $\iftag{FOL}{\Sat{M}}{\pSat{v}}{\Gamma}$، تاں
$\iftag{FOL}{\Sat{M}}{\pSat{v}}{!A}$ وی۔
\end{proof}

\begin{cor}
\ollabel{cor:consistency-soundness}
جے $\Gamma$ پورا ہون جوگ اے، تاں اوہ بے تضاد اے۔
\end{cor}

\begin{proof}
اسیں مقابل نقیض ثابت کردے آں۔ فرض کرو کہ $\Gamma$ بے تضاد نہیں۔ فیر کوئی
محدود $\Gamma_0 \subseteq \Gamma$ تے $\Gamma_0 \Sequent \quad$ دی
!!a{derivation} اے۔ \olref{thm:sequent-soundness} مطابق
$\Gamma_0 \Sequent \quad$ معتبر اے۔ دوجے لفظاں وچ، ہر
\iftag{FOL}{!!{structure}~$\Struct{M}$}{!!{valuation}~$\pAssign{v}$}
لئی کوئی $!C \in \Gamma_0$ اے جیہدے لئی
$\iftag{FOL}{\Sat/{M}}{\pSat/{v}}{!C}$، تے چونکہ $\Gamma_0
\subseteq \Gamma$، اوہ $!C$،~$\Gamma$ وچ وی اے۔ ایس لئی کوئی وی
\iftag{FOL}{$\Struct{M}$}{$\pAssign{v}$}،~$\Gamma$ نوں پورا نہیں کردا،
تے $\Gamma$ پورا ہون جوگ نہیں۔
\end{proof}
\endgroup
% /OLP-0081

% OLP-0082 content/first-order-logic/sequent-calculus/identity.tex
\begingroup
\def\olfilename{identity}
\expandafter\def\csname ol@sectioncs\endcsname{section}

\olfileid{fol}{seq}{ide}

\olsection{\usetoken{P}{derivation} وچ \usetoken{S}{identity} دی شمولیت}

!!^{derivation}s وچ !!{identity} شامل کرن لئی وادھو ابتدائی سیکوئنٹ
تے استنباطی قاعدے چاہیدے نیں۔

\begin{defn}[$\eq$ لئی ابتدائی سیکوئنٹ]
جے $t$ بند ٹرم اے، تاں ${} \Sequent \eq[t][t]$ اک ابتدائی سیکوئنٹ اے۔
\end{defn}

$\eq$ دے قاعدے ایہ نیں ($t_1$ تے $t_2$ بند ٹرم نیں):

\begin{defish}
\Axiom$ \eq[t_1][t_2], \Gamma \fCenter \Delta, !A(t_1) $
\RightLabel{$\eq$}
\UnaryInf$\eq[t_1][t_2], \Gamma \fCenter \Delta, !A(t_2)$
\DisplayProof
\hfill
\Axiom$\eq[t_1][t_2], \Gamma \fCenter \Delta, !A(t_2) $
\RightLabel{$\eq$}
\UnaryInf$\eq[t_1][t_2], \Gamma  \fCenter \Delta, !A(t_1)$
\DisplayProof
\end{defish}

\begin{ex}
جے $s$ تے $t$ بند ٹرم نیں، تاں $\eq[s][t], !A(s)
\Proves !A(t)$:
\begin{prooftree}
\Axiom$ !A(s) \fCenter !A(s)$
\RightLabel{\LeftR{\Weakening}}
\UnaryInf$\eq[s][t], !A(s)  \fCenter !A(s)$
\RightLabel{$\eq$}
\UnaryInf$\eq[s][t], !A(s)  \fCenter !A(t)$
\end{prooftree}
تساں ایس نوں اکو ہون والیاں چیزاں دی اک دوجے دی تھاں ورتوں دا اصول
\pnbLatin{(substitutability of identicals)}، یا لائبنز دا قانون \pnbLatin{(Leibniz' Law)}،
دے ناں نال جان دے ہوؤ گے۔

$\Log{LK}$ ثابت کردا اے کہ $\eq$ متناظر (symmetric) تے متعدی (transitive) اے:
\begin{prooftree}
\Axiom$ \fCenter \eq[t_1][t_1] $
\RightLabel{\LeftR{\Weakening}}
\UnaryInf$ \eq[t_1][t_2] \fCenter \eq[t_1][t_1] $
\RightLabel{$\eq$}
\UnaryInf$ \eq[t_1][t_2] \fCenter \eq[t_2][t_1]$
\DisplayProof\qquad\bottomAlignProof
\Axiom$ \eq[t_1][t_2] \fCenter \eq[t_1][t_2] $
\RightLabel{\LeftR{\Weakening}}
\UnaryInf$\eq[t_2][t_3], \eq[t_1][t_2]  \fCenter \eq[t_1][t_2] $
\RightLabel{$\eq$}
\UnaryInf$\eq[t_2][t_3], \eq[t_1][t_2]  \fCenter \eq[t_1][t_3]$
\RightLabel{\LeftR{\Exchange}}
\UnaryInf$\eq[t_1][t_2], \eq[t_2][t_3]  \fCenter \eq[t_1][t_3]$
\end{prooftree}
کھبے پاسے والے !!{derivation} وچ !!{formula}~$\eq[x][t_1]$ ساڈا
$!A(x)$ اے۔ سجے پاسے اسیں $!A(x)$ نوں~$\eq[t_1][x]$ لیندے آں۔
\end{ex}

\begin{prob}
ہیٹھ دتے سیکوئنٹاں دے !!{derivation}s بناؤ:
\begin{enumerate}
\item $\Sequent \lforall[x][\lforall[y][((x = y \land !A(x)) \lif !A(y))]]$
\item $\lexists[x][!A(x)] \land \lforall[y][\lforall[z][((!A(y) \land
    !A(z)) \lif y = z)]] \Sequent
\lexists[x][(!A(x) \land \lforall[y][(!A(y) \lif y = x)])]$
\end{enumerate}
\end{prob}
\endgroup
% /OLP-0082
% OLP-0083 content/first-order-logic/sequent-calculus/soundness-identity.tex
\begingroup
\def\olfilename{soundness-identity}
\expandafter\def\csname ol@sectioncs\endcsname{section}

\olfileid{fol}{seq}{sid}

\olsection{\usetoken{S}{identity} نال صحت مندی}

\begin{prop}
$\Log{LK}$، اکو ہون لئی ابتدائی سیکوئنٹاں تے قاعدیاں سمیت، صحت مند اے۔
\end{prop}

\begin{proof}
${} \Sequent \eq[t][t]$ روپ دے ابتدائی سیکوئنٹ معتبر نیں، کیوں جے
ہر !!{structure}~$\Struct M$ لئی $\Sat{M}{\eq[t][t]}$۔ (چیتے رہے کہ
اسیں ٹرم $t$ نوں بند منّیا اے، یعنی ایس وچ کوئی متغیر نہیں، ایس لئی
متغیراں دیاں قدر بندیاں غیر متعلق نیں۔)

فرض کرو کہ کسے !!a{derivation} دا آخری استنباط $=$ اے۔ تاں مقدمہ
$\eq[t_1][t_2], \Gamma \Sequent \Delta, !A(t_1)$ تے نتیجہ
$\eq[t_1][t_2], \Gamma \Sequent \Delta, !A(t_2)$ اے۔
!!a{structure}~$\Struct M$ لو۔ اسیں وکھانا اے کہ نتیجہ معتبر اے، یعنی جے
$\Sat{M}{\eq[t_1][t_2]}$ تے $\Sat{M}{\Gamma}$، تاں یا
$\Sat{M}{!C}$ کسے $!C \in \Delta$ لئی، یا $\Sat{M}{!A(t_2)}$۔

استقرا دے مفروضے مطابق مقدمہ معتبر اے۔ ایس دا مطلب اے کہ جے
$\Sat{M}{\eq[t_1][t_2]}$ تے $\Sat{M}{\Gamma}$، تاں یا (a) کسے $!C
\in \Delta$ لئی $\Sat{M}{!C}$، یا (b) $\Sat{M}{!A(t_1)}$۔ (a) وچ
گل ثابت ہو گئی۔ (b) نوں ویکھو۔ $s$ اک متغیر قدر بندی ہووے جیہدے لئی
$s(x) = \Value{t_1}{M}$۔ \olref[syn][ass]{prop:sentence-sat-true} دے مطابق
$\Sat{M}{!A(t_1)}[s]$۔ چونکہ $\varAssign{s}{s}{x}$، ایس لئی
\olref[syn][ext]{prop:ext-formulas} دے مطابق $\Sat{M}{!A(x)}[s]$۔ چونکہ
$\Sat{M}{\eq[t_1][t_2]}$، ایس لئی $\Value{t_1}{M} = \Value{t_2}{M}$،
تے ایس کرکے $s(x) = \Value{t_2}{M}$۔
\olref[syn][ext]{prop:ext-formulas} فیر لاون نال
$\Sat{M}{!A(t_2)}[s]$ وی ملدا اے۔
\olref[syn][ass]{prop:sentence-sat-true} دے مطابق
$\Sat{M}{!A(t_2)}$۔
\end{proof}
\endgroup
% /OLP-0083

\OLEndChapterHook
\endgroup
% /OLP-0069

% OLP-0084 content/first-order-logic/natural-deduction/natural-deduction.tex
\begingroup
\def\olfilename{natural-deduction}
\expandafter\def\csname ol@sectioncs\endcsname{section}

\iftag{FOL}
      {\olchapter{fol}{ntd}{فطری استنباط \pnbLatin{(Natural Deduction)}}}
      {\olchapter{pl}{ntd}{فطری استنباط \pnbLatin{(Natural Deduction)}}}

\begin{editorial}
ایہ باب Gentzen/Prawitz دے انداز وچ فطری استنباط دا اک نظام پیش کردا اے۔

ثبوتی نظام وجوں فطری استنباط نال متعلقہ مواد شامل یا خارج کرن لئی
``prfND'' ٹیگ ورتو۔
\end{editorial}

% OLP-0085 content/first-order-logic/natural-deduction/rules-and-proofs.tex
\begingroup
\def\olfilename{rules-and-proofs}
\expandafter\def\csname ol@sectioncs\endcsname{section}

\iftag{FOL}
      {\olfileid{fol}{ntd}{rul}}
      {\olfileid{pl}{ntd}{rul}}

\olsection{قاعدے تے \usetoken{P}{derivation}}

\begin{explain}
فطری استنباط دے نظام ریاضیاتی ثبوت وچ ورتے جاندے غیر رسمی استدلال دے
بہت نیڑے ہون لئی بنائے جاندے نیں (ایس کرکے ایہ کجھ حد تک
``فطری'' نیں)۔ فطری استنباط دے ثبوت مفروضیاں توں شروع ہوندے نیں۔
فیر استنباطی قاعدے لاگو کیتے جاندے نیں۔ \Intro{\lnot}, \Intro{\lif},
\iftag{FOL}{\Elim{\lor} تے
  \Elim{\lexists}}{ تے \Elim{\lor}} استنباطی قاعدیاں راہیں مفروضے
``!!{discharged}'' کیتے جاندے نیں، تے وضاحت لئی !!{discharged}
مفروضے دا لیبل اوہ استنباط دے نال لکھیا جاندا اے۔
\end{explain}

\begin{defn}[مفروضہ]
\emph{مفروضہ} اوہ !!{sentence} اے جیہڑا کسے وی شاخ دے سب توں اُپرلے
مقام اُتے ہووے۔
\end{defn}

فطری استنباط وچ !!^{derivation}s، !!{sentence}s دے کجھ خاص درخت نیں،
جنہاں وچ سب توں اُپرلے !!{sentence}s مفروضے ہوندے نیں، تے جے
% PNBNTD-001: source calls these upper sentence nodes "sequents"; the target uses sentences.
!!a{sentence} اک، دو یا تِن ہور جملیاں دے ہِٹھاں آوے، تاں اوہ لازمی طور
تے استنباطی قاعدے مطابق ٹھیک نکلے۔ استنباط دے اُتّے والے
!!{sentence}s نوں \emph{مقدمے} تے ہِٹھلے !!{sentence} نوں
\emph{نتیجہ} آکھدے نیں۔ قاعدے جوڑیاں وچ آؤندے نیں: ہر !!{operator}
لئی اک تعارفی تے اک اخراجی قاعدہ۔ ایہ نتیجے وچ !!a{operator} شامل
کردے یا قاعدے دے مقدمے وچوں !!a{operator} ہٹا دیندے نیں۔ کجھ قاعدے
خاص قسم دے مفروضے نوں \emph{!!{discharged}} کرن دی اجازت دیندے نیں۔
کیہڑا مفروضہ کیہڑے استنباط نال !!{discharged} ہوندا اے، ایہ دسن لئی
اسیں مفروضے تے استنباط دوناں نوں لیبل وی دیندے آں۔ ایہ مفروضہ
``$\Discharge{!A}{n}$.'' لکھ کے وکھایا جاندا اے۔

% Only include this sentence is on of \land, lor, lif or lnot is defined.

\iftag{notprvNot,notprvAnd,notprvOr,notprvIf}{}%
	{رواج اے کہ سارے !!{operator}s $\land$, $\lor$, $\lif$, $\lnot$، تے $\lfalse$ لئی قاعدیاں نوں ویکھیا جائے، بھاویں ایہناں وچوں کجھ تعریف شدہ وی ہون۔}
\endgroup
% /OLP-0085

% OLP-0086 content/first-order-logic/natural-deduction/propositional-rules.tex
\begingroup
\def\olfilename{propositional-rules}
\expandafter\def\csname ol@sectioncs\endcsname{section}

\iftag{FOL}
      {\olfileid{fol}{ntd}{prl}}
      {\olfileid{pl}{ntd}{prl}}

\olsection{قضیاتی قاعدے}


\subsection{$\land$ لئی قاعدے}

\begin{defish}
\AxiomC{$!A$}
\AxiomC{$!B$}
\RightLabel{\Intro{\land}}
\BinaryInfC{$!A \land !B$}
\DisplayProof
\hfill
\begin{tabular}{r}
\AxiomC{$!A \land !B$}
\RightLabel{\Elim{\land}}
\UnaryInfC{$!A$}
\DisplayProof
\\[3ex]
\AxiomC{$!A \land !B$}
\RightLabel{\Elim{\land}}
\UnaryInfC{$!B$}
\DisplayProof
\end{tabular}
\end{defish}

\subsection{$\lor$ لئی قاعدے}

\begin{defish}
\begin{tabular}{r}
\AxiomC{$!A$}
\RightLabel{\Intro{\lor}}
\UnaryInfC{$!A \lor !B$}
\DisplayProof
\\[3ex]
\AxiomC{$!B$}
\RightLabel{\Intro{\lor}}
\UnaryInfC{$!A \lor !B$}
\DisplayProof
\end{tabular}
\hfill
\AxiomC{$!A \lor !B$}
\AxiomC{$\Discharge{!A}{n}$}
\DeduceC{$!C$}
\AxiomC{$\Discharge{!B}{n}$}
\DeduceC{$!C$}
\DischargeRule{\Elim{\lor}}{n}
\TrinaryInfC{$!C$}
\DisplayProof
\end{defish}

\subsection{$\lif$ لئی قاعدے}

\begin{defish}
\AxiomC{$\Discharge{!A}{n}$}
\DeduceC{$!B$}
\DischargeRule{\Intro{\lif}}{n}
\UnaryInfC{$!A \lif !B$}
\DisplayProof
\hfill
\AxiomC{$!A \lif !B$}
\AxiomC{$!A$}
\RightLabel{\Elim{\lif}}
\BinaryInfC{$!B$}
\DisplayProof
\end{defish}

\subsection{$\lnot$ لئی قاعدے}

\begin{defish}
\AxiomC{$\Discharge{!A}{n}$}
\noLine
\DeduceC{$\lfalse$}
\DischargeRule{\Intro{\lnot}}{n}
\UnaryInfC{$\lnot !A$}
\DisplayProof
\hfill
\AxiomC{$\lnot !A$}
\AxiomC{$!A$}
\RightLabel{\Elim{\lnot}}
\BinaryInfC{$\lfalse$}
\DisplayProof
\end{defish}

\subsection{$\lfalse$ لئی قاعدے}

\begin{defish}
\AxiomC{$\lfalse$}
\RightLabel{\FalseInt}
\UnaryInfC{$!A$}
\DisplayProof
\hfill
\AxiomC{$\Discharge{\lnot !A}{n}$}
\DeduceC{$\lfalse$}
\DischargeRule{\FalseCl}{n}
\UnaryInfC{$!A$}
\DisplayProof
\end{defish}

ایہ گل نوٹ کرو کہ $\Intro{\lnot}$ تے $\FalseCl$ بہت رلدے ملدے
نیں: فرق ایہ اے کہ $\Intro{\lnot}$ نفی والا
!!{sentence}~$\lnot !A$ اخذ کردا اے، پر $\FalseCl$ اثباتی
!!{sentence}~$!A$ اخذ کردا اے۔

جدوں وی کوئی قاعدہ ایہ دسے کہ کسے مفروضے نوں خارج کیتا جا سکدا اے،
اسیں ایہنوں اجازت سمجھدے آں، لازمی شرط نہیں۔ مثلاً
$\Intro{\lif}$ قاعدے وچ اسیں شکل~$!A$ دے جتنے مرضی مفروضے
مقدمہ~$!B$ دے !!{derivation} وچ خارج کر سکدے آں، صفر وی۔
\endgroup
% /OLP-0086

% OLP-0087 content/first-order-logic/natural-deduction/quantifier-rules.tex
\begingroup
\def\olfilename{quantifier-rules}
\expandafter\def\csname ol@sectioncs\endcsname{section}

\olfileid{fol}{ntd}{qrl}

\olsection{مسوّراں دے قاعدے}

\subsection{$\lforall$ لئی قاعدے}

\begin{defish}
\AxiomC{$!A(a)$}
\RightLabel{\Intro{\lforall}}
\UnaryInfC{$\lforall[x][\Atom{!A}{x}]$}
\DisplayProof
\hfill
\AxiomC{$\lforall[x][\Atom{!A}{x}]$}
\RightLabel{\Elim{\lforall}}
\UnaryInfC{$!A(t)$}
\DisplayProof
\end{defish}

$\lforall$ دے قاعدیاں وچ $t$ اک بند ٹرم اے (ایہوجیہا ٹرم جیہدے
وچ کوئی متغیر نہ ہووے)، تے $a$~اک !!a{constant} اے جیہڑا
نتیجے~$\lforall[x][!A(x)]$ وچ یا کسے اوہ مفروضے وچ نہیں آؤندا
جیہڑا مقدمے~$!A(a)$ اُتے ختم ہون والے !!{derivation} وچ
!!{undischarged} ہووے۔ اسیں $a$ نوں \Intro{\lforall} استنباط دا
\emph{ایگن متغیر} آکھدے آں۔\footnote{اسیں ``ایگن متغیر'' دا ناں
ورتدے آں، بھاویں اُتّے والے قاعدے وچ $a$ اک مستقل اے۔ ایہدا
تاریخی سبب اے۔}

\subsection{$\lexists$ لئی قاعدے}

\begin{defish}
\AxiomC{$\Atom{!A}{t}$}
\RightLabel{\Intro{\lexists}}
\UnaryInfC{$\lexists[x][\Atom{!A}{x}]$}
\DisplayProof
\hfill
\AxiomC{$\lexists[x][\Atom{!A}{x}]$}
\AxiomC{[$\Atom{!A}{a}$]$^n$}
\DeduceC{$!C$}
\DischargeRule{\Elim{\lexists}}{n}
\BinaryInfC{$!C$}
\DisplayProof
\end{defish}

ایتھے وی $t$ بند ٹرم اے، تے $a$ اک !!a{constant} اے جیہڑا
مقدمے $\lexists[x][!A(x)]$ وچ، نتیجے~$!C$ وچ، یا دوہاں
مقدمیاں اُتے ختم ہون والیاں !!{derivation}s دے کسے وی
!!{undischarged} مفروضے وچ نہیں آؤندا ($!A(a)$ والے مفروضیاں
توں سوا)۔ اسیں $a$ نوں \Elim{\lexists} استنباط دا
\emph{ایگن متغیر} آکھدے آں۔

% PNBNTD-002: source says absent from premises, but universal introduction has premise A(a); use conclusion and undischarged assumptions.
ایگن متغیر دا متعلقہ نتیجے وچ، وجودی اخراج دے وجودی مقدمے وچ،
تے مقدمیاں تک جان والیاں
!!{derivation}s دے اوہناں مفروضیاں وچ نہ آؤنا جیہڑے
!!{undischarged} ہون، \Intro{\lforall} یا \Elim{\lexists}
استنباط لئی \emph{ایگن متغیر دی شرط} آکھیا جاندا اے۔

\begin{explain}
ایہ رواج یاد کرو: جے $!A$ اک !!a{formula} اے جیہدے وچ
!!{variable}~$x$ آزاد ہووے، اسیں ایہنوں~$!A(x)$ لکھ کے ظاہر کردے
آں۔ اوسے سیاق وچ $!A(t)$،~$\Subst{!A}{t}{x}$ دا مختصر روپ اے۔
اس لئی اسیں $\Intro\lexists$ دا قاعدہ ایویں وی لکھ سکدے آں:
\begin{prooftree}
\AxiomC{$\Subst{!A}{t}{x}$}
\RightLabel{\Intro{\lexists}}
\UnaryInfC{$\lexists[x][!A]$}
\end{prooftree}
یاد رکھو کہ $t$ پہلے ای~$!A$ وچ آ سکدا اے؛ مثلاً $!A$،
~$\Atom{\Obj P}{t,x}$ ہو سکدا اے۔ ایس لئی
$\lexists[x][\Atom{\Obj P}{t,x}]$ دا
$\Atom{\Obj P}{t,t}$ توں استنباط~$\Intro\lexists$ دی درست ورتوں اے: مقدمے
وچ~$t$ دیاں اک یا ودھ آمدواں، پر لازمی نہیں کہ ساریاں،
پابند !!{variable}~$x$ نال ``بدل'' سکدے او۔ پر
$\Intro\lforall$ تے~$\Elim\lexists$ دیاں ایگن متغیر شرطاں منگدیاں
نیں کہ !!{constant}~$a$،~$!A$ وچ نہ آوے۔ ایس لئی
$\lforall[x][\Atom{\Obj P}{a,x}]$ نوں
$\Atom{\Obj P}{a,a}$ توں~$\Intro\lforall$ نال درست طور تے اخذ نہیں کر سکدے۔
\end{explain}

\begin{explain}
% PNBNTD-003: source's "anything" is constrained by its prior closed-term requirement; no further restriction on closed t.
\Intro{\lexists} تے \Elim{\lforall} وچ بند ٹرم~$t$ اُتے ہور
کوئی پابندی نہیں، ایس لئی کسے اضافی شرط دی فکر نہیں کرنی پیندی۔
دوجے پاسے، \Elim{\lexists} تے \Intro{\lforall} دے قاعدیاں وچ
ایگن متغیر دی شرط منگدی اے کہ !!{constant}~$a$ نتیجے یا کسے
!!{undischarged} مفروضے وچ کتے وی نہ آوے۔ ایہ شرط نظام دی
صحت مندی یقینی بناؤن لئی ضروری اے؛ یعنی نظام صرف اوہ
!!{sentence}s، !!{undischarged} مفروضیاں توں !!{derive}s کرے
جیہڑے اوہناں توں منطقی طور تے نکلدے ہون۔ ایس شرط توں بغیر
ہِٹھلا استنباط جائز ہو جاندا:
\begin{prooftree}
  \AxiomC{$\lexists[x][!A(x)]$}
  \AxiomC{$\Discharge{!A(a)}{1}$}
  \RightLabel{*\Intro{\lforall}}
  \UnaryInfC{$\lforall[x][!A(x)]$}
  \RightLabel{\Elim{\lexists}}
  \BinaryInfC{$\lforall[x][!A(x)]$}
\end{prooftree}
پر $\lexists[x][!A(x)] \Entails/ \lforall[x][!A(x)]$۔

چونکہ مسوّراں دے اخراجی قاعدے !!{variable}s دی تھان بند ٹرم
رکھن دی ای اجازت دیندے نیں، ایس توں نکلدا اے کہ !!{sentence}s دے
کسے سیٹ توں اخذ ہو سکدا ہر !!{formula} خود وی اک
!!a{sentence} اے۔
\end{explain}
\endgroup
% /OLP-0087

% OLP-0088 content/first-order-logic/natural-deduction/derivations.tex
\begingroup
\def\olfilename{derivations}
\expandafter\def\csname ol@sectioncs\endcsname{section}

\iftag{FOL}
      {\olfileid{fol}{ntd}{der}}
      {\olfileid{pl}{ntd}{der}}

\olsection{اخذ (\usetoken{P}{derivation})}

\begin{explain}
اسیں دس چکے آں کہ مفروضہ کیہ ہوندا اے، تے استنباط دے قاعدے وی
دے چکے آں۔ فطری استنباط وچ !!^{derivation}s ایہناں توں
استقرائی طور تے بنائے جاندے نیں: ہر !!{derivation} یا تاں اکیلا
مفروضہ ہوندا اے، یا اک، دو، یا تِن !!{derivation}s دے بعد
درست استنباط لا کے بندا اے۔
\end{explain}

\begin{defn}[!!^{derivation}]
\Article{derivation} \emph{!!{derivation}}، !!a{sentence}~$!A$ دا
مفروضیاں~$\Gamma$ توں اک محدود درخت اے جیہڑا !!{sentence}s دا
بنیا ہووے تے ہِٹھ لکھیاں شرطاں پوری کردا ہووے:
\begin{enumerate}
\item درخت دے سب توں اُپرلے !!{sentence}s یا تاں $\Gamma$ وچ ہوندے
  نیں یا درخت دے کسے استنباط نال !!{discharged} ہوندے نیں۔
\item درخت دا سب توں ہِٹھلا !!{sentence}~$!A$ اے۔
\item درخت دا ہر !!{sentence}، ہِٹھلے آخرلے جملے~$!A$ توں سوا،
  استنباطی قاعدے دی درست ورتوں دا مقدمہ اے، تے اوس قاعدے دا نتیجہ
  درخت وچ اوس !!{sentence} دے ٹھیک ہِٹھاں آؤندا اے۔
\end{enumerate}
فیر اسیں $!A$ نوں !!{derivation} دا \emph{نتیجہ} تے $\Gamma$
نوں ایہدے !!{undischarged} مفروضے آکھدے آں۔

جے $!A$ دا~$\Gamma$ توں !!a{derivation} موجود ہووے، اسیں آکھدے
آں کہ $!A$،~$\Gamma$ توں \emph{!!{derivable}} اے، یا علامتاں وچ:
$\Gamma \Proves
!A$۔ جے~$!A$ دا !!a{derivation} ہووے جیہدے ہر مفروضے نوں
!!{discharged} کیتا گیا ہووے، اسیں~$\Proves !A$ لکھدے آں۔
\end{defn}

\begin{ex}
ہر مفروضہ اپنے آپ وچ !!a{derivation} اے۔ مثلاً $!A$ اپنے آپ وچ
!!a{derivation} اے، تے $!B$ وی۔ اسیں ایہناں توں، مثال لئی
$\Intro{\land}$ دا قاعدہ لا کے، نواں !!{derivation} بنا سکدے آں:
\begin{prooftree}
\AxiomC{$!A$}
\AxiomC{$!B$}
\RightLabel{\Intro{\land}}
\BinaryInfC{$!A \land !B$}
\end{prooftree}
ایہ قاعدے عام ورتوں لئی نیں: اسیں ایہدے وچ $!A$ تے~$!B$ دی
تھاں کوئی وی !!{sentence}s رکھ سکدے آں، مثلاً $!C$ تے~$!D$۔
فیر نتیجہ $!C \land !D$ ہووے گا، اس لئی
\begin{prooftree}
\AxiomC{$!C$}
\AxiomC{$!D$}
\RightLabel{\Intro{\land}}
\BinaryInfC{$!C \land !D$}
\end{prooftree}
اک درست !!{derivation} اے۔ یقیناً اسیں مفروضیاں دی ترتیب وی
بدل سکدے آں: $!D$،~$!A$ والا کردار تے $!C$،~$!B$ والا
کردار ادا کرے۔ ایس لئی
\begin{prooftree}
\AxiomC{$!D$}
\AxiomC{$!C$}
\RightLabel{\Intro{\land}}
\BinaryInfC{$!D \land !C$}
\end{prooftree}
وی اک درست !!{derivation} اے۔

ہن اسیں ہور قاعدہ، مثلاً $\Intro{\lif}$، لا سکدے آں۔ ایہ
شرطیہ نتیجہ کڈھن دی اجازت دیندا اے، تے اوہناں سارے مفروضیاں نوں
!!{discharge} کرن دی وی جیہڑے اوس شرطیہ دے مقدم دے عین برابر ہون۔
ایس لئی ہِٹھلے دونیں !!{derivation}s درست ہون گے:
\begin{prooftree}
\AxiomC{$\Discharge{!C}{1}$}
\AxiomC{$!D$}
\RightLabel{\Intro{\land}}
\BinaryInfC{$!C \land !D$}
\DischargeRule{\Intro{\lif}}{1}
\UnaryInfC{$!C \lif (!C \land !D)$}
\DisplayProof\bottomAlignProof
\AxiomC{$!C$}
\AxiomC{$\Discharge{!D}{1}$}
\RightLabel{\Intro{\land}}
\BinaryInfC{$!C \land !D$}
\DischargeRule{\Intro{\lif}}{1}
\UnaryInfC{$!D \lif (!C \land !D)$}
\end{prooftree}
ایہ بالترتیب وکھاندے نیں کہ $!D \Proves !C \lif (!C \land !D)$
تے $!C \Proves !D \lif (!C \land !D)$۔

یاد رکھو، مفروضے خارج کرنا اجازت اے، لازمی شرط نہیں: ساڈے لئی
مفروضے خارج کرنا ضروری نہیں۔ خاص طور تے، جے اوہ مفروضے
!!{derivation} وچ موجود وی نہ ہون، فیر وی اسیں قاعدہ لا سکدے آں۔
مثلاً ہِٹھلا استنباط جائز اے، بھاویں کسے مفروضے~$!A$ نوں
!!{discharged} کرن لئی موجود نہیں:
\begin{prooftree}
  \AxiomC{$!B$}
  \DischargeRule{\Intro{\lif}}{1}
  \UnaryInfC{$!A \lif !B$}
\end{prooftree}
\end{ex}
\endgroup
% /OLP-0088

% OLP-0089 content/first-order-logic/natural-deduction/proving-things.tex
\begingroup
\def\olfilename{proving-things}
\expandafter\def\csname ol@sectioncs\endcsname{section}

\iftag{FOL}
      {\olfileid{fol}{ntd}{pro}}
      {\olfileid{pl}{ntd}{pro}}

\olsection{اخذ (\usetoken{P}{derivation}) دیاں مثالاں}

\begin{ex}
آؤ !!{sentence} $(!A \land !B) \lif !A$ دا !!a{derivation} دیئے۔

پہلاں اسیں مطلوبہ نتیجہ !!{derivation} دے سب توں ہِٹھ لکھدے آں۔
\begin{prooftree}
\AxiomC{}
\UnaryInfC{$(!A\land !B) \lif !A$}
\end{prooftree}

اگلا سوال ایہ اے کہ ایس روپ دا !!a{sentence} کس قسم دے استنباط
توں نکل سکدا اے۔ نتیجے دا !!{main operator}~$\lif$ اے، ایس لئی
اسیں \Intro{\lif} دا قاعدہ لا کے نتیجے تک پہنچن دی کوشش کراں گے۔
نال نال ورتے ہوئے مفروضے لکھدے جاؤ تے استنباطی قاعدیاں دے ناں
لاؤ؛ ایس طرح ثبوت دے آخر وچ ویکھنا سوکھا رہندا اے کہ سارے مفروضے
!!{discharged} ہو چکے نیں یا نہیں۔
\begin{prooftree}
\AxiomC{$\Discharge{!A \land !B}{1}$}
\DeduceC{$!A$}
\DischargeRule{\Intro{\lif}}{1} 
\UnaryInfC{$(!A\land !B) \lif !A$}
\end{prooftree}

ہن مفروضے $!A \land !B$ توں $!A$ تک دے قدم بھرنے نیں۔ اتھے ساڈے
کول صرف اک منطقی رابط، $\land$، اے؛ ایس لئی $\land$ دا اخراجی
قاعدہ ورتنا پئے گا۔ ایہدے نال ہِٹھلا ثبوت ملدا اے:
\begin{prooftree}
\AxiomC{$\Discharge{!A \land !B}{1}$}
\RightLabel{\Elim{\land}}
\UnaryInfC{$!A$}
\DischargeRule{\Intro{\lif}}{1} 
\UnaryInfC{$(!A\land !B) \lif !A$}
\end{prooftree}
ہن ساڈے کول $(!A \land
!B) \lif !A$ دا درست !!{derivation} اے۔
\end{ex}

\begin{ex}
ہن آؤ $(\lnot !A \lor !B) \lif (!A \lif !B)$ دا !!a{derivation}
دیئے۔

پہلاں اسیں مطلوبہ نتیجہ !!{derivation} دے سب توں ہِٹھ لکھدے آں۔
\begin{prooftree}
\AxiomC{}
\UnaryInfC{$(\lnot !A \lor !B) \lif (!A \lif !B)$}
\end{prooftree}
ایہ نتیجہ دین والا منطقی قاعدہ لبھن لئی اسیں نتیجے دے منطقی
رابطے ویکھدے آں: $\lnot$، $\lor$ تے $\lif$۔ فی الحال ساڈا دھیان
$\lif$ دی پہلی ورتوں اُتے اے، کیوں جے اخذ دے آخری !!{sentence}
% PNBNTD-011: the ND tree ends in one sentence, not an end-sequent; retain the main-operator scope.
دا !!{main operator} ایہی اے۔ $\lnot$، $\lor$ تے $\lif$ دی دوجی
ورتوں کسے ہور رابطے دے دائرے وچ نیں، ایس لئی اوہناں نوں بعد وچ
سنبھالاں گے۔ سو اسیں \Intro{\lif} دے قاعدے توں شروع کردے آں۔
ایس دی درست ورتوں دا روپ ایہ ہووے گا:
\begin{prooftree}
\AxiomC{$\Discharge{\lnot !A \lor !B}{1}$}
\DeduceC{$!A \lif !B$}
\DischargeRule{\Intro{\lif}}{1}
\UnaryInfC{$(\lnot !A \lor !B) \lif (!A \lif !B)$}
\end{prooftree}
ہن اگے ودھن دے دو رستے نیں۔ یا اسیں ہِٹھوں اُپر ول کم جاری رکھ کے
\Intro{\lif} دی اک ہور ورتوں لبھیے، یا اُپروں ہِٹھاں ول کم کر کے
\Elim{\lor} دا قاعدہ لائیے۔ آؤ دوجا رستہ چنیے۔ مفروضے
$\lnot !A \lor !B$ نوں \Elim{\lor} دے سب توں کھبے مقدمے وجوں
ورتاں گے۔ \Elim{\lor} دی درست ورتوں وچ باقی دو مقدمے نتیجے
$!A \lif !B$ دے عین برابر ہونے چاہیدے نیں؛ البتہ ہر اک نوں
$\lnot !A \lor !B$ دے دو فصلی جزاں وچوں متعلقہ جز نوں مفروضہ
من کے وکھرا اخذ کیتا جا سکدا اے۔ ایس لئی ساڈا !!{derivation}
ایویں لگے گا:
\begin{prooftree}
\AxiomC{$\Discharge{\lnot !A \lor !B}{1}$}
\AxiomC{$\Discharge{\lnot !A}{2}$}
\DeduceC{$!A \lif !B$}
\AxiomC{$\Discharge{!B}{2}$}
\DeduceC{$!A \lif !B$}
\DischargeRule{\Elim{\lor}}{2}
\TrinaryInfC{$!A \lif !B$}
\DischargeRule{\Intro{\lif}}{1} 
\UnaryInfC{$(\lnot !A \lor !B) \lif (!A \lif !B)$}
\end{prooftree}

سجے پاسے دیاں دو ذیلی شاخاں وچ اسیں $!A \lif !B$ نوں !!{derive}
کرنا چاہندے آں؛ ایس لئی \Intro{\lif} ورتنا سب توں چنگا اے۔
\begin{prooftree}
\AxiomC{$\Discharge{\lnot !A \lor !B}{1}$}
\AxiomC{$\Discharge{\lnot !A}{2}, \Discharge{!A}{3}$}
\DeduceC{$!B$}
\DischargeRule{\Intro{\lif}}{3}
\UnaryInfC{$!A \lif !B$}
\AxiomC{$\Discharge{!B}{2}, \Discharge{!A}{4}$}
\DeduceC{$!B$}
\DischargeRule{\Intro{\lif}}{4}
\UnaryInfC{$!A \lif !B$}
\DischargeRule{\Elim{\lor}}{2}
\TrinaryInfC{$!A \lif !B$}
\DischargeRule{\Intro{\lif}}{1} 
\UnaryInfC{$(\lnot !A \lor !B) \lif (!A \lif !B)$}
\end{prooftree}

!!{derivation} دے باقی دو خالی حصیاں لئی سانوں $!B$ دے
!!{derivation}s چاہیدے نیں: وچکارلی شاخ وچ $\lnot !A$ تے $!A$
توں، تے سب توں سجی شاخ وچ $!A$ تے $!B$ توں۔
% PNBNTD-004: source calls the second subsidiary branch "on the left"; the tree makes it rightmost.
پہلاں وچکارلا حصہ لیندے آں۔ $\lnot
!A$ تے $!A$، \Elim{\lnot} دے دو مقدمے نیں:
\begin{prooftree}
\AxiomC{$\Discharge{\lnot !A}{2}$}
\AxiomC{$\Discharge{!A}{3}$}
\RightLabel{\Elim{\lnot}}
\BinaryInfC{$\lfalse$}
\DeduceC{$!B$}
\end{prooftree}
\FalseInt لا کے اسیں $!B$ نوں نتیجے وجوں حاصل کر سکدے آں
تے ایہ شاخ مکمل ہو جاندی اے۔
\begin{prooftree}
\AxiomC{$\Discharge{\lnot !A \lor !B}{1}$}
\AxiomC{$\Discharge{\lnot !A}{2}$}
\AxiomC{$\Discharge{!A}{3}$}
% PNBNTD-005: source intermediate label Intro(false) contradicts its Elim(not) rule and finished tree.
\RightLabel{\Elim{\lnot}}
\BinaryInfC{$\lfalse$}
\RightLabel{\FalseInt}
\UnaryInfC{$!B$}
\DischargeRule{\Intro{\lif}}{3}
\UnaryInfC{$!A \lif !B$}
\AxiomC{$\Discharge{!B}{2}, \Discharge{!A}{4}$}
\DeduceC{$!B$}
\DischargeRule{\Intro{\lif}}{4}
\UnaryInfC{$!A \lif !B$}
\DischargeRule{\Elim{\lor}}{2}
\TrinaryInfC{$!A \lif !B$}
\DischargeRule{\Intro{\lif}}{1} 
\UnaryInfC{$(\lnot !A \lor !B) \lif (!A \lif !B)$}
\end{prooftree}

ہن سب توں سجی شاخ ویکھو۔ اتھے ایہ سمجھنا ضروری اے کہ
!!{derivation} دی تعریف مفروضے \emph{خارج کرن دی اجازت دیندی اے}،
پر اوہناں نوں خارج کرنا \emph{لازمی نہیں بناؤندی}۔ دوجے لفظاں وچ،
جے اسیں $!B$ نوں مفروضیاں $!A$ تے $!B$ وچوں اک توں
اخذ کر لئیے، دوجے نوں ورتے بنا، تاں ایہ ٹھیک اے۔ $!B$ نوں خود $!B$ توں
!!{derive} کرنا بالکل سادہ اے: اکیلا $!B$ ایہو جیہا
!!a{derivation} اے تے کوئی استنباط نہیں چاہیدا۔ سو اسیں
مفروضہ~$!A$ ہٹا سکدے آں۔
\begin{prooftree}
\AxiomC{$\Discharge{\lnot !A \lor !B}{1}$}
\AxiomC{$\Discharge{\lnot !A}{2}$}
\AxiomC{$\Discharge{!A}{3}$}
\RightLabel{\Elim{\lnot}}
\BinaryInfC{$\lfalse$}
\RightLabel{\FalseInt}
\UnaryInfC{$!B$}
\DischargeRule{\Intro{\lif}}{3}
\UnaryInfC{$!A \lif !B$}
\AxiomC{$\Discharge{!B}{2}$}
\RightLabel{\Intro{\lif}}
\UnaryInfC{$!A \lif !B$}
\DischargeRule{\Elim{\lor}}{2}
\TrinaryInfC{$!A \lif !B$}
\DischargeRule{\Intro{\lif}}{1} 
\UnaryInfC{$(\lnot !A \lor !B) \lif (!A \lif !B)$}
\end{prooftree}
دھیان دیو کہ مکمل !!{derivation} وچ سب توں سجی \Intro{\lif}
والی استنباطی ورتوں حقیقت وچ کوئی مفروضہ خارج نہیں کردی۔
\end{ex}

\begin{ex}
ہن تک سانوں \FalseCl{} دے قاعدے دی لوڑ نہیں پئی۔ ایہ ایس گل
وچ خاص اے کہ ایہ اوہ مفروضہ وی خارج کرن دی اجازت دیندا اے جیہڑا
قاعدے دے نتیجے دا ذیلی-!!{formula} نہیں ہوندا۔ ایہ \FalseInt{}
دے قاعدے نال نیڑے جڑیا ہویا اے۔ اصل وچ \FalseInt{} دا قاعدہ،
\FalseCl{} دے قاعدے دی اک خاص صورت اے---اک منطق، ”وجدانیت پسند
منطق“ \pnbLatin{(intuitionistic logic)}، وی اے جیہدے وچ صرف \FalseInt{}
دی اجازت اے۔ جدوں ہور کوئی طریقہ کم نہ کرے، \FalseCl{} آخری
چارہ اے۔ مثلاً فرض کرو اسیں $!A \lor \lnot !A$ نوں !!{derive}
کرنا چاہندے آں۔ ساڈی عام تدبیر $!A \lor \lnot !A$ نوں
$\Intro{\lor}$ لا کے !!{derive} کرن دی ہووے گی۔ پر ایس لئی
بنا کسے مفروضے دے یا $!A$ یا $\lnot !A$ نوں !!{derive} کرنا
پئے گا، تے ایہ نہیں ہو سکدا۔ ہن \FalseCl{} کم آؤندا اے:
\begin{prooftree}
  \AxiomC{$\Discharge{\lnot(!A \lor \lnot !A)}{1}$}
  \DeduceC{$\lfalse$}
  \DischargeRule{\FalseCl}{1}
  \UnaryInfC{$!A \lor \lnot !A$}
\end{prooftree}
ہن سانوں $\lfalse$ دا !!a{derivation}،
$\lnot(!A \lor \lnot !A)$ توں چاہیدا اے۔ کیوں جے $\lfalse$،
$\Elim{\lnot}$ دا نتیجہ اے، اسیں اوہ قاعدہ آزما سکدے آں:
\begin{prooftree}
  \AxiomC{$\Discharge{\lnot(!A \lor \lnot !A)}{1}$}
  \DeduceC{$\lnot !A$}
  \AxiomC{$\Discharge{\lnot(!A \lor \lnot !A)}{1}$}
  \DeduceC{$!A$}
  \RightLabel{\Elim{\lnot}}
  \BinaryInfC{$\lfalse$}
  \DischargeRule{\FalseCl}{1}
  \UnaryInfC{$!A \lor \lnot !A$}
\end{prooftree}
$\lnot !A$ دا !!a{derivation} لبھن دی ساڈی تدبیر
$\Intro{\lnot}$ لاون نوں آکھدی اے:
\begin{prooftree}
  \AxiomC{$\Discharge{\lnot(!A \lor \lnot !A)}{1}, \Discharge{!A}{2}$}
  \DeduceC{$\lfalse$}
  \DischargeRule{\Intro{\lnot}}{2}
  \UnaryInfC{$\lnot !A$}
  \AxiomC{$\Discharge{\lnot(!A \lor \lnot !A)}{1}$}
  \DeduceC{$!A$}
  \RightLabel{\Elim{\lnot}}
  \BinaryInfC{$\lfalse$}
  \DischargeRule{\FalseCl}{1}
  \UnaryInfC{$!A \lor \lnot !A$}
\end{prooftree}
اتھے اسیں $\lfalse$ سوکھ نال حاصل کر سکدے آں: $\Elim{\lnot}$
نوں مفروضے $\lnot(!A \lor \lnot !A)$ تے
$!A \lor \lnot !A$ اُتے لا کے۔ دوجا مقدمہ ساڈے نویں مفروضے $!A$
توں $\Intro{\lor}$ راہیں نکلدا اے:
\begin{prooftree}
  \AxiomC{$\Discharge{\lnot(!A \lor \lnot !A)}{1}$}
  \AxiomC{$\Discharge{!A}{2}$}
  \RightLabel{\Intro{\lor}}
  \UnaryInfC{$!A \lor \lnot !A$}
  \RightLabel{\Elim{\lnot}}
  \BinaryInfC{$\lfalse$}
  \DischargeRule{\Intro{\lnot}}{2}
  \UnaryInfC{$\lnot !A$}
  \AxiomC{$\Discharge{\lnot(!A \lor \lnot !A)}{1}$}
  \DeduceC{$!A$}
  \RightLabel{\Elim{\lnot}}
  \BinaryInfC{$\lfalse$}
  \DischargeRule{\FalseCl}{1}
  \UnaryInfC{$!A \lor \lnot !A$}
\end{prooftree}
سجے پاسے اسیں اوہی تدبیر ورتدے آں، پر اتھے $!A$ نوں
\FalseCl راہیں حاصل کردے آں:
\begin{prooftree}
  \AxiomC{$\Discharge{\lnot(!A \lor \lnot !A)}{1}$}
  \AxiomC{$\Discharge{!A}{2}$}
  \RightLabel{\Intro{\lor}}
  \UnaryInfC{$!A \lor \lnot !A$}
  \RightLabel{\Elim{\lnot}}
  \BinaryInfC{$\lfalse$}
  \DischargeRule{\Intro{\lnot}}{2}
  \UnaryInfC{$\lnot !A$}
  \AxiomC{$\Discharge{\lnot(!A \lor \lnot !A)}{1}$}
  \AxiomC{$\Discharge{\lnot !A}{3}$}
  \RightLabel{\Intro{\lor}}
  \UnaryInfC{$!A \lor \lnot !A$}
  \RightLabel{\Elim{\lnot}}
  \BinaryInfC{$\lfalse$}
  \DischargeRule{\FalseCl}{3}
  \UnaryInfC{$!A$}
  \RightLabel{\Elim{\lnot}}
  \BinaryInfC{$\lfalse$}
  \DischargeRule{\FalseCl}{1}
  \UnaryInfC{$!A \lor \lnot !A$}
\end{prooftree}
\end{ex}

\begin{prob}
ہِٹھ لکھیاں گلّاں وکھاؤن والے !!{derivation}s دیو:
\begin{enumerate}
\item $!A \land (!B \land !C) \Proves (!A \land !B) \land !C$.
\item $!A \lor (!B \lor !C) \Proves (!A \lor !B) \lor !C$.
\item $!A \lif (!B \lif !C) \Proves !B \lif (!A \lif !C)$.
\item $!A \Proves \lnot\lnot !A$.
\end{enumerate}
\end{prob}

\begin{prob}
ہِٹھ لکھیاں گلّاں وکھاؤن والے !!{derivation}s دیو:
\begin{enumerate}
\item $(!A \lor !B) \lif !C \Proves !A \lif !C$.
\item $(!A \lif !C) \land (!B \lif !C) \Proves (!A \lor !B) \lif !C$.
\item $\Proves \lnot(!A \land \lnot !A)$.
\item $!B \lif !A \Proves \lnot !A \lif \lnot !B$.
\item $\Proves (!A \lif \lnot !A) \lif \lnot !A$.
\item $\Proves \lnot(!A \lif !B) \lif \lnot !B$.
\item $!A \lif !C \Proves \lnot (!A \land \lnot !C)$.
\item $!A \land \lnot !C \Proves \lnot (!A \lif !C)$.
\item $!A \lor !B, \lnot !B \Proves !A$.
\item $\lnot !A \lor \lnot !B \Proves \lnot(!A \land !B)$.
\item $\Proves (\lnot !A \land \lnot !B) \lif\lnot(!A \lor !B)$.
\item $\Proves \lnot(!A \lor !B) \lif (\lnot !A \land \lnot !B)$.
\end{enumerate}
\end{prob}

\begin{prob}
ہِٹھ لکھیاں گلّاں وکھاؤن والے !!{derivation}s دیو:
\begin{enumerate}
\item $\lnot(!A \lif !B) \Proves !A$.
\item $\lnot(!A \land !B) \Proves \lnot !A \lor \lnot !B$.
\item $!A \lif !B \Proves \lnot !A \lor !B$.
\item $\Proves \lnot \lnot !A \lif !A$.
\item $!A \lif !B, \lnot !A \lif !B \Proves !B$.
\item $(!A \land !B) \lif !C \Proves (!A \lif !C) \lor (!B \lif !C)$.
\item $(!A \lif !B) \lif !A \Proves !A$.
\item $\Proves (!A \lif !B) \lor (!B \lif !C)$.
\end{enumerate}
(ایہناں ساریاں لئی $\FalseCl$~دا قاعدہ ضروری اے۔)
\end{prob}
\endgroup
% /OLP-0089

% OLP-0090 content/first-order-logic/natural-deduction/proving-things-quant.tex
\begingroup
\def\olfilename{proving-things-quant}
\expandafter\def\csname ol@sectioncs\endcsname{section}

\olfileid{fol}{ntd}{prq}

\olsection{مسوّراں نال \usetoken{P}{derivation}}

\begin{ex}
مسوّراں نال کم کردے ہوئیاں اسیں دھیان رکھنا اے کہ ایگن متغیر
دی شرط نہ ٹُٹے۔ کدی کدی ایس لئی کجھ استنباط لاون دی ترتیب نوں
بدل کے ویکھنا پیندا اے۔ عام طور اُتے پہلے اوہ قاعدے سنبھالن نال
مدد ملدی اے جیہڑے ایگن متغیر دی شرط دے تابع نیں؛ مکمل ثبوت وچ
اوہ ہِٹھلے پاسے ہون گے۔

آؤ ویکھئیے کہ !!{formula}
$\lexists[x][\lnot !A(x)] \lif \lnot \lforall[x][!A(x)]$
دا !!a{derivation} کیویں دیاں گے۔ معمول وانگ، شروع وچ اسیں ایہ
لکھدے آں:
\begin{prooftree}
\AxiomC{}
\UnaryInfC{$\lexists[x][\lnot !A(x)]\lif \lnot \lforall[x][!A(x)]$}
\end{prooftree}
فیر اسیں لکھدے آں کہ آخری قدم نوں \Intro{\lif} دے قاعدے نال
جائز ٹھہراون لئی کیہ چاہیدا اے:
\begin{prooftree}
\AxiomC{$\Discharge{\lexists[x][\lnot !A(x)]}{1}$}
\DeduceC{$\lnot \lforall[x][!A(x)]$}
\DischargeRule{\Intro{\lif}}{1}
\UnaryInfC{$\lexists[x][\lnot !A(x)]\lif \lnot \lforall[x][!A(x)]$}
\end{prooftree}
کیوں جے $\lnot \lforall[x][!A(x)]$ اُتے لاون لئی کوئی صاف
ظاہر قاعدہ نہیں لبھدا، اسیں !!{derivation} نوں ایس طرح ترتیب
دیواں گے کہ \Elim{\lexists} دا قاعدہ ورت سکئیے۔ اتھے ایگن
متغیر دی شرط دھیان وچ رکھنی اے، تے ایہو جیہا مستقل چننا اے جیہڑا
% PNBNTD-006: the source cites existential A, but its major premise is existential not-A.
$\lexists[x][\lnot !A(x)]$ وچ یا اوہناں مفروضیاں وچ نہ آوے
جنہاں اُتے ایہ اخذ نربھر اے۔ (پر اتھے کوئی !!{constant}s آؤندے
ہی نہیں، ایس لئی کوئی وی انتخاب ٹھیک اے۔)
\begin{prooftree}
\AxiomC{$\Discharge{\lexists[x][\lnot !A(x)]}{1}$}
\AxiomC{$\Discharge{\lnot !A(a)}{2}$}
\DeduceC{$\lnot \lforall[x][!A(x)]$}
\DischargeRule{\Elim{\lexists}}{2}
\BinaryInfC{$\lnot \lforall[x][!A(x)]$}
\DischargeRule{\Intro{\lif}}{1}
\UnaryInfC{$\lexists[x][\lnot !A(x)] \lif \lnot \lforall[x][!A(x)]$}
\end{prooftree}
$\lnot \lforall[x][!A(x)]$ اخذ کرن لئی اسیں \Intro{\lnot}
دا قاعدہ آزماواں گے: ایس لئی اک تضاد اخذ کرنا ہووے گا، ممکن اے
کہ $\lforall[x][!A(x)]$ نوں وادھو مفروضہ بنا کے۔ ظاہر اے، ایس
تضاد وچ مفروضہ $\lnot !A(a)$ وی کم آ سکدا اے، جیہڑا
\Elim{\lexists} دے استنباط نال خارج ہووے گا۔ اسیں ایہدا روپ
ایویں بنا سکدے آں:
\begin{prooftree}
\AxiomC{$\Discharge{\lexists[x][\lnot !A(x)]}{1}$}
\AxiomC{$\Discharge{\lnot !A(a)}{2}, \Discharge{\lforall[x][!A(x)]}{3}$}
\DeduceC{$\lfalse$}
\DischargeRule{\Intro{\lnot}}{3}
\UnaryInfC{$\lnot \lforall[x][!A(x)]$}
\DischargeRule{\Elim{\lexists}}{2}
\BinaryInfC{$\lnot \lforall[x][!A(x)]$}
\DischargeRule{\Intro{\lif}}{1}
\UnaryInfC{$\lexists[x][\lnot !A(x)]\lif \lnot \lforall[x][!A(x)]$}
\end{prooftree}
لگدا اے کہ ہن اسیں تضاد دے نیڑے آ گئے آں۔ لاون لئی سب توں
سوکھا قاعدہ \Elim{\lforall} اے، جیہدے نال ایگن متغیر دی
کوئی شرط نہیں۔ کلی مسوّر والے~$x$ دی تھاں اسیں کوئی وی بند ٹرم
رکھ سکدے آں؛ ایس لئی تضاد تک پہنچن واسطے اوہی مستقل~$a$
ورتنا سب توں مناسب اے۔
% PNBNTD-007: source says "any term"; the governing universal-elimination rule requires a closed term.
\begin{prooftree}
\AxiomC{$\Discharge{\lexists[x][\lnot !A(x)]}{1}$}
\AxiomC{$\Discharge{\lnot !A(a)}{2}$}
\AxiomC{$\Discharge{\lforall[x][!A(x)]}{3}$}
\RightLabel{\Elim{\lforall}}
\UnaryInfC{$!A(a)$}
\RightLabel{\Elim{\lnot}}
\BinaryInfC{$\lfalse$}
\DischargeRule{\Intro{\lnot}}{3}
\UnaryInfC{$\lnot \lforall[x][!A(x)]$}
\DischargeRule{\Elim{\lexists}}{2}
\BinaryInfC{$\lnot \lforall[x][!A(x)]$}
\DischargeRule{\Intro{\lif}}{1}
\UnaryInfC{$\lexists[x][\lnot !A(x)]\lif \lnot \lforall[x][!A(x)]$}
\end{prooftree}

خاص کر مسوّراں نال کم کردے ہوئیاں، ایس مرحلے اُتے دُوبارہ جانچنا
ضروری اے کہ ایگن متغیر دی شرط نہیں ٹُٹی۔ اسیں جیہڑا واحد قاعدہ
ایگن متغیر دی شرط دے تابع لایا، اوہ \Elim{\exists} سی؛ تے
ایگن متغیر~$a$ کسے وی اوہ مفروضے وچ نہیں آؤندا جیہدے
اُتے ایہ اخذ نربھر اے۔ ایس لئی ایہ درست !!{derivation} اے۔
\end{ex}

\begin{ex}
کدی اسیں ہور !!{formula}s توں !!a{formula} اخذ کردے آں۔
ایہناں صورتاں وچ کجھ مفروضے خارج نہ ہوئے وی رہ سکدے نیں۔
ساڈے لئی مفروضیاں تے آخری مقصد، دونیں دا حساب رکھنا ضروری اے۔

آؤ ویکھئیے کہ !!{formula} $\lexists[x][!C(x,b)]$ دا
!!a{derivation} مفروضیاں $\lexists[x][(!A(x)
\land !B(x))]$ تے $\lforall[x][(!B(x) \lif !C(x,b))]$
توں کیویں دیاں گے۔ معمول وانگ شروع کر کے نتیجہ سب توں ہِٹھ
لکھدے آں۔
\begin{prooftree}
\AxiomC{}
\UnaryInfC{$\lexists[x][!C(x,b)]$}
\end{prooftree}

ساڈے کول ورتن لئی دو مقدمے نیں۔ پہلے مقدمے نوں ورتن لئی، یعنی
$\lexists[x][!C(x, b)]$ دا !!a{derivation}
$\lexists[x][(!A(x) \land !B(x))]$ توں لبھن دی کوشش
وچ، اسیں \Elim{\lexists} دا قاعدہ ورتاں گے۔ ایس قاعدے اُتے
ایگن متغیر دی شرط لگدی اے، ایس لئی اسیں ایہ قاعدہ پہلے
لاواں گے۔ نتیجہ ایہ اے:
\begin{prooftree}
\AxiomC{$\lexists[x][(!A(x) \land !B(x))]$}
\AxiomC{$\Discharge{!A(a) \land !B(a)}{1}$}
\DeduceC{$\lexists[x][!C(x,b)]$}
\DischargeRule{\Elim{\lexists}}{1}
\BinaryInfC{$\lexists[x][!C(x,b)]$}
\end{prooftree}
اسیں جنہاں دو مفروضیاں نال کم کر رہے آں، اوہناں وچ~$!B$
سانجھا اے۔ ایس ویلے $!B(a)$ نوں وکھ کرن لئی \Elim{\land}
لاونا فائدہ مند ہو سکدا اے۔
\begin{prooftree}
% PNBNTD-008: source intermediate tree reverses the final conjunction parenthesis and argument bracket.
\AxiomC{$\lexists[x][(!A(x) \land !B(x))]$}
\AxiomC{$\Discharge{!A(a) 
\land !B(a)}{1}$}
\RightLabel{\Elim{\land}}
\UnaryInfC{$!B(a)$}
\DeduceC{$\lexists[x][!C(x,b)]$}
\DischargeRule{\Elim{\lexists}}{1}
\BinaryInfC{$\lexists[x][!C(x,b)]$}
\end{prooftree}

ساڈا دوجا مفروضہ~$\lforall[x][(!B(x) \lif
  !C(x,b))]$ اے۔ اتھے ایگن متغیر دی شرط نہیں، سو اسیں
\Elim{\lforall} لا کے $x$ دی تھاں !!{constant}~$a$ رکھ سکدے
آں؛ ایس نال $!B(a) \lif !C(a, b)$ ملدا اے۔ ہن ساڈے کول
$!B(a) \lif !C(a,b)$ تے $!B(a)$، دونیں نیں۔ اگلا قدم
\Elim{\lif} دے قاعدے دی سیدھی ورتوں ہونا چاہیدا اے۔
\begin{prooftree}
\AxiomC{$\lexists[x][(!A(x) \land !B(x))]$}
\AxiomC{$\lforall[x][(!B(x) \lif !C(x,b))]$}
\RightLabel{\Elim{\lforall}}
\UnaryInfC{$!B(a) \lif !C(a,b)$}
\AxiomC{$\Discharge{!A(a) 
\land !B(a)}{1}$}
\RightLabel{\Elim{\land}}
\UnaryInfC{$!B(a)$}
\RightLabel{\Elim{\lif}}
\BinaryInfC{$!C(a,b)$}
\DeduceC{$\lexists[x][!C(x,b)]$}
\DischargeRule{\Elim{\lexists}}{1}
\insertBetweenHyps{\hspace{-5em}}
\BinaryInfC{$\lexists[x][!C(x,b)]$}
\end{prooftree}
ہن منزل بہت نیڑے اے!{} \Intro{\lexists} دی اک ورتوں
نال مقصد حاصل ہو جاندا اے۔
\begin{prooftree}
\AxiomC{$\lexists[x][(!A(x) \land !B(x))]$}
\AxiomC{$\lforall[x][(!B(x) \lif !C(x,b))]$}
\RightLabel{\Elim{\lforall}}
\UnaryInfC{$!B(a) \lif !C(a,b)$}
\AxiomC{$\Discharge{!A(a) 
\land !B(a)}{1}$}
\RightLabel{\Elim{\land}}
\UnaryInfC{$!B(a)$}
\RightLabel{\Elim{\lif}}
\BinaryInfC{$!C(a,b)$}
\RightLabel{\Intro{\lexists}}
\UnaryInfC{$\lexists[x][!C(x,b)]$}
\DischargeRule{\Elim{\lexists}}{1}
\insertBetweenHyps{\hspace{-5em}}
\BinaryInfC{$\lexists[x][!C(x,b)]$}
\end{prooftree}
ہر قدم اُتے اسیں ایگن متغیر دی شرط نہ ٹُٹن دی جانچ کیتی،
ایس لئی اسیں بھروسے نال آکھ سکدے آں کہ ایہ درست
!!{derivation} اے۔
\end{ex}

\begin{ex}
!!{formula} $\lnot\lforall[x][!A(x)]$ دا !!a{derivation}
مفروضیاں $\lforall[x][!A(x)] 
\lif \lexists[y][!B(y)]$ تے $\lnot\lexists[y][!B(y)]$
توں دیو۔ معمول وانگ اسیں مطلوبہ !!{formula}
سب توں ہِٹھ لکھدے آں۔
\begin{prooftree}
\AxiomC{}
\UnaryInfC{$\lnot\lforall[x][!A(x)]$}
\end{prooftree}
!!{derivation} دی آخری سطر اک نفی اے، سو آؤ
\Intro{\lnot} آزما کے ویکھئیے۔ ایس لئی سانوں لبھنا ہووے گا
کہ تضاد کیویں !!{derive} کیتا جاوے۔
\begin{prooftree}
\AxiomC{$\Discharge{\lforall[x][!A(x)]}{1}$}
\DeduceC{$\lfalse$}
\DischargeRule{\Intro{\lnot}}{1}
\UnaryInfC{$\lnot\lforall[x][!A(x)]$}
\end{prooftree}
ہن تک ٹھیک اے۔ اسیں \Elim{\lforall} لا سکدے آں، پر ایہ
صاف نہیں کہ ایس نال ساڈے مقصد تک پہنچن وچ مدد ملے گی۔ ایس دی
تھاں اپنا اک مفروضہ ورتدے آں۔ $\lforall[x][!A(x)] \lif
\lexists[y][!B(y)]$ نوں $\lforall[x][!A(x)]$ نال ملا کے
اسیں \Elim{\lif} دا قاعدہ لا سکدے آں۔
\begin{prooftree}
\AxiomC{$\lforall[x][!A(x)] \lif \lexists[y][!B(y)]$}
\AxiomC{$\Discharge{\lforall[x][!A(x)]}{1}$}
\RightLabel{\Elim{\lif}}
\BinaryInfC{$\lexists[y][!B(y)]$}
\DeduceC{$\lfalse$}
\DischargeRule{\Intro{\lnot}}{1}
\UnaryInfC{$\lnot\lforall[x][!A(x)]$}
\end{prooftree}
ہن اک آخری مفروضہ ورتنا رہ گیا اے، تے لگدا اے کہ
\Elim{\lnot} لا کے ایہ سانوں تضاد تک پہنچا دے گا۔
\begin{prooftree}
\AxiomC{$\lnot\lexists[y][!B(y)]$}
\AxiomC{$\lforall[x][!A(x)] \lif \lexists[y][!B(y)]$}
\AxiomC{$\Discharge{\lforall[x][!A(x)]}{1}$}
\RightLabel{\Elim{\lif}}
\BinaryInfC{$\lexists[y][!B(y)]$}
\RightLabel{\Elim{\lnot}}
\BinaryInfC{$\lfalse$}
\DischargeRule{\Intro{\lnot}}{1}
\UnaryInfC{$\lnot\lforall[x][!A(x)]$}
\end{prooftree}
\end{ex}

\begin{prob}
ہِٹھ لکھیاں گلّاں وکھاؤن والے !!{derivation}s دیو:
\begin{enumerate}
\item $\Proves (\lforall[x][!A(x)] \land \lforall[y][!B(y)]) \lif
\lforall[z][(!A(z) \land !B(z))]$.
\item $\Proves (\lexists[x][!A(x)] \lor \lexists[y][!B(y)]) \lif
\lexists[z][(!A(z) \lor !B(z))]$.
\item $\lforall[x][(!A(x) \lif !B)] \Proves \lexists[y][!A(y)] \lif !B$.
\item $\lforall[x][\lnot !A(x)] \Proves \lnot\lexists[x][!A(x)]$.
\item $\Proves \lnot\lexists[x][!A(x)] \lif \lforall[x][\lnot !A(x)]$.
\item $\Proves \lnot\lexists[x][\lforall[y][((!A(x,y) \lif \lnot
!A(y,y)) \land (\lnot !A(y,y) \lif !A(x,y)))]]$.
\end{enumerate}
\end{prob}

\begin{prob}
ہِٹھ لکھیاں گلّاں وکھاؤن والے !!{derivation}s دیو:
\begin{enumerate}
\item $\Proves \lnot\lforall[x][!A(x)] \lif \lexists[x][\lnot!A(x)]$.
\item $(\lforall[x][!A(x)] \lif !B) \Proves \lexists[y][(!A(y) \lif !B)]$.
\item $\Proves \lexists[x][(!A(x) \lif \lforall[y][!A(y)])]$.
\end{enumerate}
(ایہناں ساریاں لئی $\FalseCl$~دا قاعدہ ضروری اے۔)
\end{prob}
\endgroup
% /OLP-0090

% OLP-0091 content/first-order-logic/natural-deduction/proof-theoretic-notions.tex
\begingroup
\def\olfilename{proof-theoretic-notions}
\expandafter\def\csname ol@sectioncs\endcsname{section}

\iftag{FOL}
      {\olfileid{fol}{ntd}{ptn}}
      {\olfileid{pl}{ntd}{ptn}}

\olsection{ثبوتی-نظریاتی تصورات \pnbLatin{(Proof-Theoretic Notions)}}

\begin{editorial}
ایہ حصہ فطری استنباط لئی ثابت پذیری دے تعلق تے بے تضادی دیاں
تعریفاں اکٹھیاں کردا اے۔
\end{editorial}

\begin{explain}
جیویں اسیں کئی اہم معنوی تصورات (منطقی اعتبار، منطقی نتیجہ تے پورا ہون
جوگتا) دی تعریف دِتّی اے، ہُن اسیں اوہناں دے مطابق
\emph{ثبوتی-نظریاتی تصورات \pnbLatin{(proof-theoretic notions)}} دی تعریف دیندے آں۔
ایہناں دی تعریف !!{sentence}s دے !!{structure}s وچ پورا اترن ول رجوع کر کے
نہیں، بلکہ کجھ !!{sentence}s دا ہور جملیاں توں !!{derivability}
یا !!{nonderivability} ول رجوع کر کے دِتّی جاندی اے۔ ایہ اک اہم دریافت سی
کہ ایہ تصور آپو وچ میل کھاندے نیں۔ اوہناں دے میل کھان دی گل
\emph{صحت مندی (soundness)} تے \emph{تمامیت دے قضیاں \pnbLatin{(completeness theorems)}}
دا مضمون اے۔
\end{explain}


\begin{defn}[قضیے (Theorems)]
!!{sentence}~$!A$ اک \emph{قضیہ (theorem)} اے جے فطری استنباط وچ
~$!A$ دی کوئی !!a{derivation} ہووے جیہدے سارے مفروضے !!{discharged} ہون۔
اسیں $\Proves !A$ لکھدے آں جے $!A$ قضیہ اے، تے جے ایہ قضیہ نہیں تاں
$\Proves/ !A$ لکھدے آں۔
\end{defn}

\begin{defn}[!!^{derivability}]
!!^a{sentence} $!A$، !!{sentence}s دے سیٹ~$\Gamma$ توں
\emph{!!{derivable}} اے، $\Gamma \Proves !A$، جے کوئی
!!{derivation} ایہو جہی ہووے جیہدا نتیجہ~$!A$ ہووے تے جیہدا ہر مفروضہ
یا تاں !!{discharged} ہووے یا~$\Gamma$ وچ ہووے۔ جے $!A$، $\Gamma$ توں
!!{derivable} نہیں تاں اسیں $\Gamma \Proves/ !A$ لکھدے آں۔
\end{defn}

\begin{defn}[بے تضادی (Consistency)]
!!{sentence}s دا سیٹ~$\Gamma$ \emph{متضاد (inconsistent)} اے اُدوں تے صرف
اُدوں جدوں $\Gamma
\Proves \lfalse$۔ جے $\Gamma$ متضاد نہیں، یعنی جے
$\Gamma \Proves/ \lfalse$، تاں اسیں اوہنوں \emph{بے تضاد (consistent)}
آکھدے آں۔
\end{defn}

\begin{prop}[انعکاسیت (Reflexivity)]
\ollabel{prop:reflexivity}
جے $!A \in \Gamma$ اے، تاں $\Gamma \Proves !A$۔
\end{prop}

\begin{proof}
مفروضہ $!A$ اپنے آپ وچ~$!A$ دی اک !!a{derivation} اے جیہدے سارے
!!{undischarged} مفروضے (یعنی $!A$)~$\Gamma$ وچ نیں۔
\end{proof}
  

\begin{prop}[یک رُخیت (Monotonicity)]
\ollabel{prop:monotonicity}
جے $\Gamma \subseteq \Delta$ تے $\Gamma \Proves !A$، تاں $\Delta
\Proves !A$۔
\end{prop}

\begin{proof}
~$!A$ دی~$\Gamma$ توں ہر !!{derivation}،~$!A$ دی~$\Delta$ توں وی
!!a{derivation} اے۔
\end{proof}

\begin{prop}[تعدیت (Transitivity)]
\ollabel{prop:transitivity}
جے $\Gamma \Proves !A$ تے $\{!A\} \cup \Delta \Proves
!B$، تاں $\Gamma \cup \Delta \Proves !B$۔
\end{prop}

\begin{proof}
جے $\Gamma \Proves !A$ اے، تاں اک !!a{derivation}~$\delta_0$،~$!A$
دی اے، جیہدے سارے !!{undischarged} مفروضے~$\Gamma$ وچ نیں۔ جے $\{!A\} \cup
\Delta \Proves !B$ اے، تاں اک !!a{derivation}~$\delta_1$،~$!B$
دی اے، جیہدے سارے !!{undischarged} مفروضے~$\{!A\} \cup \Delta$ وچ نیں۔
ہُن ایہ ویکھو:
\begin{prooftree}
  \AxiomC{$\Delta, \Discharge{!A}{1}$}
  \RightLabel{$\delta_1$}
  \DeduceC{$!B$}
  \DischargeRule{\Intro{\lif}}{1}
  \UnaryInfC{$!A \lif !B$}
  \AxiomC{$\Gamma$}
  \RightLabel{$\delta_0$}
  \DeduceC{$!A$}
  \RightLabel{\Elim{\lif}}
  \BinaryInfC{$!B$}
\end{prooftree}
ہُن سارے !!{undischarged} مفروضے $\Gamma \cup
\Delta$ وچ نیں؛ ایس لئی ایہ وکھاؤندا اے کہ $\Gamma \cup \Delta \Proves !B$۔
\end{proof}

جدوں $\Gamma = \{!A_1, !A_2, \ldots, !A_k\}$ کوئی محدود سیٹ ہووے، تاں اسیں مختصر علامت $!A_1,!A_2,\ldots,!A_k \Proves !B$، $\Gamma \Proves !B$ دی تھاں ورت سکدے آں؛ خاص طور تے $!A \Proves !B$ دا مطلب اے کہ $\{!A\} \Proves !B$۔

دھیان دیو کہ جے $\Gamma \Proves !A$ تے $!A
\Proves !B$، تاں $\Gamma \Proves !B$۔ ایہ وی نکلدا اے کہ جے $!A_1,
\dots, !A_n \Proves !B$ تے $\Gamma \Proves !A_i$ ہر~$i$ لئی، تاں
$\Gamma \Proves !B$۔

\begin{prop}
\ollabel{prop:incons}
ہِٹھلیاں گلاں اک دوجے دے برابر نیں:
    \begin{enumerate}
      \item \( \Gamma \) متضاد اے۔
      \item \( \Gamma \Proves {!A} \)، ہر !!{sentence}~\( {!A} \) لئی۔
      \item \( \Gamma \Proves {!A} \) تے \( \Gamma \Proves \lnot {!A} \)، کسے !!{sentence}~\( {!A} \) لئی۔
    \end{enumerate}
\end{prop}

\begin{proof}
مشق۔
\end{proof}

\tagprob{FOL}
\begin{prob}
\olref[fol][ntd][ptn]{prop:incons} ثابت کرو۔
\end{prob}
\tagendprob

\tagprob{notFOL}
\begin{prob}
\olref[pl][ntd][ptn]{prop:incons} ثابت کرو۔
\end{prob}
\tagendprob

\begin{prop}[محدود گواہی (Compactness)]
  \ollabel{prop:proves-compact}
  \begin{enumerate}
  \item جے $\Gamma \Proves !A$، تاں کوئی محدود ذیلی سیٹ $\Gamma_0
    \subseteq \Gamma$ ایہو جہا اے کہ $\Gamma_0 \Proves !A$۔
  \item جے~$\Gamma$ دا ہر محدود ذیلی سیٹ
    بے تضاد اے، تاں $\Gamma$ بے تضاد اے۔
  \end{enumerate}
\end{prop}

\begin{proof}
  \begin{enumerate}
    \item جے $\Gamma \Proves !A$، تاں
      اک !!a{derivation}~$\delta$،~$!A$ دی~$\Gamma$ توں اے۔ منّ لو $\Gamma_0$،~$\delta$ دے
      !!{undischarged} مفروضیاں دا سیٹ اے۔ چونکہ ہر !!{derivation} محدود
      ہوندی اے، ایس لئی~$\Gamma_0$ وچ صرف محدود گنتی دے
      !!{sentence}s ہو سکدے نیں۔ سو~$\delta$،~$!A$ دی اک !!a{derivation} محدود~$\Gamma_0 \subseteq \Gamma$
      توں اے۔
    \item ایہ (1) دا متقابلِ عکس اے، خاص حالت $!A
      \ident \lfalse$ لئی۔
  \end{enumerate}
\end{proof}
\endgroup
% /OLP-0091

% OLP-0092 content/first-order-logic/natural-deduction/provability-consistency.tex
\begingroup
\def\olfilename{provability-consistency}
\expandafter\def\csname ol@sectioncs\endcsname{section}

\iftag{FOL}
      {\olfileid{fol}{ntd}{prv}}
      {\olfileid{pl}{ntd}{prv}}
      
\olsection{\usetoken{S}{derivability} تے بے تضادی \pnbLatin{(Consistency)}}

ہُن اسیں !!{derivability} دے تعلق دیاں کئی خاصیتاں ثابت کراں گے۔ ایہ
اپنے آپ وچ وی دلچسپ نیں، تے ہر اک تمامیت دے قضیے دے ثبوت وچ اپنا کردار
ادا کرے گی۔

\begin{prop}\ollabel{prop:provability-contr}
  جے $\Gamma \Proves !A$ تے $\Gamma \cup \{!A\}$ متضاد اے،
  تاں $\Gamma$ متضاد اے۔
\end{prop}

\begin{proof}
منّ لو~$!A$ دی~$\Gamma$ توں !!{derivation}~$\delta_1$ اے، تے
~$\lfalse$ دی $\Gamma \cup \{!A\}$ توں
!!{derivation}~$\delta_2$ اے۔ فیر اسیں ایویں !!{derive} کر سکدے آں:
\begin{prooftree}
\AxiomC{$\Gamma, \Discharge{!A}{1}$}
\RightLabel{$\delta_2$}
\DeduceC{$\lfalse$}
\DischargeRule{\Intro{\lnot}}{1}
\UnaryInfC{$\lnot !A$}
\AxiomC{$\Gamma$}
\RightLabel{$\delta_1$}
\DeduceC{$!A$}
\RightLabel{\Elim{\lnot}}
\BinaryInfC{$\lfalse$}
\end{prooftree}
نویں !!{derivation} وچ مفروضہ~$!A$ !!{discharged} اے، ایس لئی ایہ
~$\Gamma$ توں !!a{derivation} اے۔
\end{proof}

\begin{prop}
\ollabel{prop:prov-incons}
$\Gamma \Proves !A$ اُدوں تے صرف اُدوں جدوں $\Gamma \cup \{\lnot !A\}$ متضاد اے۔
\end{prop}

\begin{proof}
پہلاں منّ لو $\Gamma \Proves !A$، یعنی اک
!!a{derivation}~$\delta_0$،~$!A$ دی اے جیہدے !!{undischarged}
مفروضے~$\Gamma$ وچ نیں۔ اس طرح اسیں $\lfalse$ دی
$\Gamma \cup \{\lnot !A\}$ توں !!a{derivation} حاصل کردے آں:
\begin{prooftree}
  \AxiomC{$\lnot !A$}
  \AxiomC{$\Gamma$}
  \RightLabel{$\delta_0$}
  \DeduceC{$!A$}
  \RightLabel{\Elim{\lnot}}
  \BinaryInfC{$\lfalse$}
\end{prooftree}

ہُن منّ لو $\Gamma \cup \{\lnot !A\}$ متضاد اے، تے
$\delta_1$ اوہدی مطابق~$\lfalse$ دی !!{derivation} اے جیہدے
!!{undischarged} مفروضے~$\Gamma \cup \{\lnot !A\}$ وچ نیں۔
اسیں~$!A$ دی صرف~$\Gamma$ توں !!a{derivation}،
~$\FalseCl$ ورت کے حاصل کردے آں:
\begin{prooftree}
  \AxiomC{$\Gamma, \Discharge{\lnot !A}{1}$}
  \RightLabel{$\delta_1$}
  \DeduceC{$\lfalse$}
  \RightLabel{\FalseCl}
  \DischargeRule{\FalseCl}{1}
  \UnaryInfC{$!A$}
\end{prooftree}
\end{proof}

\begin{prob}
ثابت کرو کہ $\Gamma \Proves \lnot !A$ اُدوں تے صرف اُدوں جدوں $\Gamma \cup \{!A\}$
متضاد اے۔
\end{prob}

\begin{prop}\ollabel{prop:explicit-inc}
  جے $\Gamma \Proves !A$ تے $\lnot !A \in \Gamma$، تاں $\Gamma$
  متضاد اے۔
\end{prop}

\begin{proof}
  منّ لو $\Gamma \Proves !A$ تے $\lnot !A \in \Gamma$۔ فیر
  اک !!a{derivation}~$\delta$،~$!A$ دی~$\Gamma$ توں اے۔
  $\Elim{\lnot}$ قاعدے دی ایہ سادی ورتوں ویکھو:
  \begin{prooftree}
    \AxiomC{$\lnot !A$}
    \AxiomC{$\Gamma$}
    \RightLabel{$\delta$}
    \DeduceC{$!A$}
    \RightLabel{\Elim{\lnot}}
    \BinaryInfC{$\lfalse$}
  \end{prooftree}
  چونکہ $\lnot !A \in \Gamma$، سارے !!{undischarged} مفروضے
  ~$\Gamma$ وچ نیں؛ ایس لئی ایہ وکھاؤندا اے کہ $\Gamma \Proves \lfalse$۔
\end{proof}

\begin{prop}\ollabel{prop:provability-exhaustive}
  جے $\Gamma \cup \{!A\}$ تے $\Gamma \cup \{\lnot !A\}$ دونیں
  متضاد نیں، تاں $\Gamma$ متضاد اے۔
\end{prop}

\begin{proof}
!!{derivation}s~$\delta_1$ تے~$\delta_2$، بالترتیب~$\lfalse$ دیاں
  $\Gamma \cup \{ !A \}$ توں تے $\lfalse$ دیاں $\Gamma \cup \{ \lnot !A
  \}$ توں نیں۔ فیر اسیں ایویں !!{derive} کر سکدے آں:
\begin{prooftree}
\AxiomC{$\Gamma, \Discharge{\lnot !A}{2}$}
\RightLabel{$\delta_2$}
\DeduceC{$\lfalse$}
\DischargeRule{\Intro{\lnot}}{2}
\UnaryInfC{$\lnot \lnot !A$}
\AxiomC{$\Gamma, \Discharge{!A}{1}$}
\RightLabel{$\delta_1$}
\DeduceC{$\lfalse$}
\DischargeRule{\Intro{\lnot}}{1}
\UnaryInfC{$\lnot !A$}
\RightLabel{\Elim{\lnot}}
\BinaryInfC{$\lfalse$}
\end{prooftree}
چونکہ مفروضے $!A$ تے $\lnot !A$ دونیں !!{discharged} نیں، ایہ
~$\lfalse$ دی صرف~$\Gamma$ توں !!a{derivation} اے۔ ایس لئی $\Gamma$
متضاد اے۔
\end{proof}
\endgroup
% /OLP-0092

% OLP-0093 content/first-order-logic/natural-deduction/provability-propositional.tex
\begingroup
\def\olfilename{provability-propositional}
\expandafter\def\csname ol@sectioncs\endcsname{section}

\iftag{FOL}
      {\olfileid{fol}{ntd}{ppr}}
      {\olfileid{pl}{ntd}{ppr}}

\olsection{\usetoken{S}{derivability} تے قضیاتی رابطے \pnbLatin{(Propositional Connectives)}}

\begin{explain}
  اسیں ثابت کردے آں کہ فطری استنباط دا !!{derivability} تعلق~$\Proves$
  قضیاتی رابطیاں نال جڑے کجھ بنیادی حقائق قائم کرن لئی اینا طاقتور اے،
  جیویں کہ $!A \land !B
  \Proves !A$ تے $!A, !A \lif !B \Proves !B$ (قاعدۂ اثباتِ مقدم؛ modus ponens)۔
  تمامیت دے قضیے دے ثبوت لئی ایہناں حقائق دی لوڑ اے۔
\end{explain}

\begin{prop}\ollabel{prop:provability-land}
  \begin{enumerate}
  \item \ollabel{prop:provability-land-left} دونیں $!A \land !B \Proves
    !A$ تے $!A \land !B \Proves !B$۔
  \item \ollabel{prop:provability-land-right} $!A, !B \Proves !A \land !B$۔
  \end{enumerate}
\end{prop}

\begin{proof}
  \begin{enumerate}
  \item اسیں دونیں !!{derive} کر سکدے آں:
    \begin{prooftree}
      \AxiomC{$!A \land !B$}
      \RightLabel{\Elim{\land}}
      \UnaryInfC{$!A$}
      \DisplayProof\qquad\bottomAlignProof
      \AxiomC{$!A \land !B$}
      \RightLabel{\Elim{\land}}
      \UnaryInfC{$!B$}
    \end{prooftree}
  \item اسیں ایہ !!{derive} کر سکدے آں:
    \begin{prooftree}
      \AxiomC{$!A$}
      \AxiomC{$!B$}
      \RightLabel{\Intro{\land}}
      \BinaryInfC{$!A \land !B$}
    \end{prooftree}
  \end{enumerate}
\end{proof}


\begin{prop}\ollabel{prop:provability-lor}
  \begin{enumerate}
  \item $!A \lor !B, \lnot !A, \lnot !B$ متضاد اے۔
  \item دونیں $!A \Proves !A \lor !B$ تے $!B \Proves !A \lor !B$۔
  \end{enumerate}
\end{prop}

\begin{proof}
  \begin{enumerate}
  \item ایہ !!{derivation} ویکھو:
    \begin{prooftree}
      \AxiomC{$!A \lor !B$}
      \AxiomC{$\lnot !A$}
      \AxiomC{$\Discharge{!A}{1}$}
      \RightLabel{\Elim{\lnot}}
      \BinaryInfC{$\lfalse$}
      \AxiomC{$\lnot !B$}
      \AxiomC{$\Discharge{!B}{1}$}
      \RightLabel{\Elim{\lnot}}
      \BinaryInfC{$\lfalse$}
      \DischargeRule{\Elim{\lor}}{1}
      \TrinaryInfC{$\lfalse$}
    \end{prooftree}
    ایہ~$\lfalse$ دی !!a{derivation} اے جیہدے !!{undischarged}
    مفروضے $!A \lor !B$، $\lnot !A$ تے $\lnot !B$ نیں۔
  \item اسیں دونیں !!{derive} کر سکدے آں:
    \begin{prooftree}
      \AxiomC{$!A$}
      \RightLabel{\Intro{\lor}}
      \UnaryInfC{$!A \lor !B$}
      \DisplayProof\qquad\bottomAlignProof
      \AxiomC{$!B$}
      \RightLabel{\Intro{\lor}}
      \UnaryInfC{$!A \lor !B$}
    \end{prooftree}
  \end{enumerate}
\end{proof}

\begin{prop}\ollabel{prop:provability-lif}
  \begin{enumerate}
  \item \ollabel{prop:provability-lif-left}  $!A, !A \lif !B \Proves !B$۔
  \item \ollabel{prop:provability-lif-right}
    دونیں $\lnot !A \Proves !A \lif !B$ تے $!B \Proves !A \lif !B$۔
  \end{enumerate}
\end{prop}

\begin{proof}
  \begin{enumerate}
  \item اسیں ایہ !!{derive} کر سکدے آں:
    \begin{prooftree}
      \AxiomC{$!A \lif !B$}
      \AxiomC{$!A$}
      \RightLabel{\Elim{\lif}}
      \BinaryInfC{$!B$}
    \end{prooftree}
    
  \item ہِٹھلیاں دو !!{derivation}s ایہ وکھاؤندیاں نیں:
    \begin{prooftree}
      \AxiomC{$\lnot !A$}
      \AxiomC{$\Discharge{!A}{1}$}
      \RightLabel{\Elim{\lnot}}
      \BinaryInfC{$\lfalse$}
      \RightLabel{\FalseInt}
      \UnaryInfC{$!B$}
      \DischargeRule{\Intro{\lif}}{1}
      \UnaryInfC{$!A \lif !B$}
      \DisplayProof\qquad\bottomAlignProof
      \AxiomC{$!B$}
      \RightLabel{\Intro{\lif}}
      \UnaryInfC{$!A \lif !B$}
    \end{prooftree}
    دھیان دیو کہ $\Intro{\lif}$ مفروضہ~$!A$ نوں !!{discharge} کر سکدا
    اے، پر ایہ کرنا لازمی نہیں۔
  \end{enumerate}
\end{proof}
\endgroup
% /OLP-0093

% OLP-0094 content/first-order-logic/natural-deduction/provability-quantifiers.tex
\begingroup
\def\olfilename{provability-quantifiers}
\expandafter\def\csname ol@sectioncs\endcsname{section}

\olfileid{fol}{ntd}{qpr}

\olsection{\usetoken{S}{derivability} تے مسوّر \pnbLatin{(Quantifiers)}}

\begin{explain}
  تمامیت دے قضیے لئی ایہ وی لوڑیںدا اے کہ فطری استنباط دے قاعدے
  ایس حصے وچ قائم کیتے~$\Proves$ بارے حقائق وی دین۔
\end{explain}

\begin{thm}
\ollabel{thm:strong-generalization} جے $c$ اک ایہو جہا مستقل اے جیہڑا
$\Gamma$ یا $!A(x)$ وچ نہیں آؤندا، تے $\Gamma \Proves !A(c)$، تاں $\Gamma
\Proves \lforall[x][!A(x)]$۔
\end{thm}

\begin{proof}
منّ لو $\delta$، $!A(c)$ دی $\Gamma$ توں !!a{derivation} اے۔
اک \Intro{\lforall} استنباط شامل کر کے اسیں
$\lforall[x][!A(x)]$ دی !!a{derivation} حاصل کردے آں۔ چونکہ $c$،
$\Gamma$ یا $!A(x)$ وچ نہیں آؤندا، ایگن متغیر والی شرط پوری ہوندی اے۔
\end{proof}

\begin{prop}
\ollabel{prop:provability-quantifiers}
% PNBNTD-009: OLP-0087's displayed quantifier rules require t to be closed in both branches.
ایہ دونیں بیان ہر بند ٹرم لئی نیں:
\begin{tagenumerate}{prvEx,prvAll}
\tagitem{prvEx}{$!A(t) \Proves \lexists[x][!A(x)]$.}{}

\tagitem{prvAll}{$\lforall[x][!A(x)] \Proves !A(t)$.}{}
\end{tagenumerate}
\end{prop}

\begin{proof}
\begin{tagenumerate}{prvEx,prvAll}
\tagitem{prvEx}{ہِٹھلی~$\lexists[x][!A(x)]$ دی~$!A(t)$ توں !!a{derivation} اے:
\begin{prooftree}
\AxiomC{$!A(t)$}
\RightLabel{\Intro{\lexists}}
\UnaryInfC{$\lexists[x][!A(x)]$}
\end{prooftree}}{}

\tagitem{prvAll}{ہِٹھلی~$!A(t)$ دی
  ~$\lforall[x][!A(x)]$ توں !!a{derivation} اے:
\begin{prooftree}
\AxiomC{$\lforall[x][!A(x)]$}
\RightLabel{\Elim{\lforall}}
\UnaryInfC{$!A(t)$}
\end{prooftree}}{}
\end{tagenumerate}
\end{proof}
\endgroup
% /OLP-0094

% OLP-0095 content/first-order-logic/natural-deduction/soundness.tex
\begingroup
\def\olfilename{soundness}
\expandafter\def\csname ol@sectioncs\endcsname{section}

\iftag{FOL}
      {\olfileid{fol}{ntd}{sou}}
      {\olfileid{pl}{ntd}{sou}}
      
\olsection{صحت مندی \pnbLatin{(Soundness)}}

\begin{explain}
!!^a{derivation} نظام، جیویں فطری استنباط، \emph{صحت مند (sound)}
اُدوں اے جے اوہ اوہ گلاں !!{derive} نہیں کر سکدا جیہڑیاں اصل وچ
نتیجے دے طور تے نہیں نکلدیاں۔ ایس لئی صحت مندی، !!{derivation} نظاماں
لئی ضمانت والی حفاظتی خاصیت اے۔ جیہڑی ثبوتی-نظریاتی خاصیت دا سوال ہووے،
اوہدے مطابق اسیں مثال دے طور تے جاننا چاہواں گے کہ
\begin{enumerate}
\item ہر !!{derivable} !!{sentence} \iftag{FOL}{معتبر}{ہمیشہ صادق فارمولا} اے؛
\item جے کوئی !!a{sentence} کجھ ہور جملیاں توں !!{derivable} اے، تاں اوہ
  اوہناں دا منطقی نتیجہ وی اے؛
\item جے !!{sentence}s دا کوئی سیٹ متضاد اے، تاں اوہ پورا ہون جوگ نہیں۔
\end{enumerate}
ایہ !!a{derivation} نظام دیاں اہم خاصیتاں نیں۔ جے ایہناں وچوں کوئی وی
قائم نہ ہووے، تاں !!{derivation} نظام ناقص ہووے گا---اوہ لوڑ توں ودھ
!!{derive} کرے گا۔ ایس لئی !!{derivation} نظام دی صحت مندی ثابت کرنا
نہایت اہم اے۔
\end{explain}

\begin{thm}[صحت مندی (Soundness)]
\ollabel{thm:soundness}
جے $!A$، !!{undischarged} مفروضیاں
$\Gamma$ توں !!{derivable} اے، تاں $\Gamma \Entails !A$۔
\end{thm}

\begin{proof}
$\delta$، $!A$ دی !!a{derivation} ہووے۔ اسیں~$\delta$ وچ
استنباطاں دی گنتی اُتے استقرا کراں گے۔

استقرا دی بنیاد لئی اسیں اوہ صورت ثابت کردے آں جدوں استنباطاں دی گنتی
~$0$ اے۔ ایس صورت وچ $\delta$ صرف اک
!!{sentence}~$!A$، یعنی اک مفروضے، اُتے مشتمل اے۔ ایہ مفروضہ
!!{undischarged} اے، کیوں جے مفروضے صرف استنباط دے ذریعے
!!{discharged} ہوندے نیں تے ایتھے کوئی استنباط نہیں۔ ایس لئی ہر
\iftag{FOL}{!!{structure}~$\Struct{M}$}{!!{valuation}~$\pAssign{v}$}
جیہڑا ثبوت دے سارے !!{undischarged} مفروضیاں نوں پورا کردا اے،
~$!A$ نوں وی پورا کردا اے۔

ہُن استقرائی قدم ویکھو۔ منّ لو $\delta$ وچ~$n$ استنباط نیں۔ سب توں
ہِٹھلے استنباط دے اک یا ودھ مقدمے ذیلی !!{derivation}s نال
!!{derive} ہوندے نیں، تے ہر اک وچ~$n$ توں گھٹ
استنباط نیں۔ اسیں استقرائی مفروضہ مندے آں: سب توں ہِٹھلے استنباط دے
مقدمے اوہناں ذیلی !!{derivation}s دے !!{undischarged} مفروضیاں
توں منطقی طور تے نکلدے نیں جیہڑیاں انہاں مقدمیاں اُتے ختم ہوندیاں نیں۔
اسیں وکھاؤنا اے کہ نتیجہ~$!A$ سارے ثبوت دے !!{undischarged} مفروضیاں
توں منطقی طور تے نکلدا اے۔

اسیں سب توں ہِٹھلے استنباط دی قسم مطابق صورتاں وکھ کردے آں۔ پہلاں
اک مقدمے والے ممکنہ استنباط ویکھدے آں۔

\begin{enumerate}
\item منّ لو آخری استنباط \Intro{\lnot} اے: !!{derivation} دی صورت ایہ اے
  \begin{prooftree}
    \AxiomC{$\Gamma, \Discharge{!A}{n}$}
    \RightLabel{$\delta_1$}
    \DeduceC{$\lfalse$}
    \DischargeRule{\Intro{\lnot}}{n}
    \UnaryInfC{$\lnot !A$}
  \end{prooftree}
  استقرائی مفروضے مطابق $\lfalse$، !!{undischarged}
  مفروضیاں $\Gamma \cup \{!A\}$ توں نکلدا اے جیہڑے~$\delta_1$
  دے نیں۔ کوئی
  \iftag{FOL}{!!a{structure}~$\Struct{M}$}{!!a{valuation}~$\pAssign{v}$}
  ویکھو۔ اسیں ثابت کرنا اے کہ جے
  $\iftag{FOL}{\Sat{M}}{\pSat{v}}{\Gamma}$، تاں
  $\iftag{FOL}{\Sat{M}}{\pSat{v}}{\lnot !A}$۔ تضاد لئی منّ لو کہ
  $\iftag{FOL}{\Sat{M}}{\pSat{v}}{\Gamma}$، پر
  $\iftag{FOL}{\Sat/{M}}{\pSat/{v}}{\lnot !A}$، یعنی
  $\iftag{FOL}{\Sat{M}}{\pSat{v}}{!A}$۔ ایس دا مطلب ہووے گا کہ
  $\iftag{FOL}{\Sat{M}}{\pSat{v}}{\Gamma \cup \{!A\}}$۔ ایہ
  ساڈے استقرائی مفروضے دے خلاف اے۔ ایس لئی
  $\iftag{FOL}{\Sat{M}}{\pSat{v}}{\lnot !A}$۔

\item آخری استنباط \Elim{\land} اے: دو صورتاں نیں: $!A$ یا $!B$
  نوں مقدمے $!A \land !B$ توں اخذ کیتا جا سکدا اے۔ پہلی صورت ویکھو۔
  !!{derivation}~$\delta$ ایویں اے:
  \begin{prooftree}
    \AxiomC{$\Gamma$}
    \RightLabel{$\delta_1$}
    \DeduceC{$!A \land !B$}
    \RightLabel{\Elim{\land}}
    \UnaryInfC{$!A$}
  \end{prooftree}
  استقرائی مفروضے مطابق $!A \land !B$، !!{undischarged}
  مفروضیاں~$\Gamma$ توں نکلدا اے جیہڑے~$\delta_1$ دے نیں۔ اک
  % PNBNTD-010: PL branch needs a valuation, not a first-order structure.
  \iftag{FOL}{!!a{structure}}{!!a{valuation}}~\iftag{FOL}{$\Struct{M}$}{$\pAssign{v}$} ویکھو۔
  اسیں ثابت کرنا اے کہ جے
  $\iftag{FOL}{\Sat{M}}{\pSat{v}}{\Gamma}$، تاں
  $\iftag{FOL}{\Sat{M}}{\pSat{v}}{!A}$۔ منّ لو
  $\iftag{FOL}{\Sat{M}}{\pSat{v}}{\Gamma}$۔ استقرائی مفروضے
  ($\Gamma \Entails !A \land !B$) توں اسیں جان دے آں کہ
  $\iftag{FOL}{\Sat{M}}{\pSat{v}}{!A \land !B}$۔ تعریف مطابق،
  $\iftag{FOL}{\Sat{M}}{\pSat{v}}{!A \land !B}$ اُدوں تے صرف اُدوں
  جدوں $\iftag{FOL}{\Sat{M}}{\pSat{v}}{!A}$ تے
  $\iftag{FOL}{\Sat{M}}{\pSat{v}}{!B}$۔ (جتھے $!B$،
  $!A \land !B$ توں اخذ ہوندا اے، اوہ صورت وی ایویں حل ہوندی اے۔)
  
\item آخری استنباط \Intro{\lor} اے: دو صورتاں نیں: $!A
  \lor !B$، مقدمے~$!A$ توں یا مقدمے~$!B$ توں اخذ ہو سکدا اے۔
  پہلی صورت ویکھو۔ !!{derivation} دی صورت ایہ اے
  \begin{prooftree}
    \AxiomC{$\Gamma$}
    \RightLabel{$\delta_1$}
    \DeduceC{$!A$}
    \RightLabel{\Intro{\lor}}
    \UnaryInfC{$!A \lor !B$}
  \end{prooftree}
  استقرائی مفروضے مطابق $!A$، !!{undischarged}
  مفروضیاں~$\Gamma$ توں نکلدا اے جیہڑے~$\delta_1$ دے نیں۔ کوئی
  \iftag{FOL}{!!a{structure}~$\Struct{M}$}{!!a{valuation}~$\pAssign{v}$}
  ویکھو۔ اسیں ثابت کرنا اے کہ جے
  $\iftag{FOL}{\Sat{M}}{\pSat{v}}{\Gamma}$، تاں
  $\iftag{FOL}{\Sat{M}}{\pSat{v}}{!A \lor !B}$۔ منّ لو
  $\iftag{FOL}{\Sat{M}}{\pSat{v}}{\Gamma}$؛ فیر
  $\iftag{FOL}{\Sat{M}}{\pSat{v}}{!A}$، کیوں جے
  $\Gamma \Entails !A$ (استقرائی مفروضہ)۔ ایس لئی
  $\iftag{FOL}{\Sat{M}}{\pSat{v}}{!A \lor !B}$ وی ہووے گا۔
  (جتھے $!A \lor !B$،~$!B$ توں اخذ ہوندا اے، اوہ صورت وی
  ایویں حل ہوندی اے۔)
  
\item آخری استنباط \Intro{\lif} اے: $!A \lif !B$ اک ایس ذیلی ثبوت
  توں اخذ ہوندا اے جیہدے وچ مفروضہ~$!A$ تے نتیجہ~$!B$ اے، یعنی
  \begin{prooftree}
    \AxiomC{$\Gamma, \Discharge{!A}{n}$}
    \RightLabel{$\delta_1$}
    \DeduceC{$!B$}
    \DischargeRule{\Intro{\lif}}{n}
    \UnaryInfC{$!A \lif !B$}
  \end{prooftree}
  استقرائی مفروضے مطابق $!B$،~$\delta_1$ دے !!{undischarged}
  مفروضیاں توں نکلدا اے، یعنی $\Gamma \cup \{!A\} \Entails
  !B$۔ کوئی
  \iftag{FOL}{!!a{structure}~$\Struct{M}$}{!!a{valuation}~$\pAssign{v}$}
  ویکھو۔~$\delta$ دے !!{undischarged} مفروضے صرف $\Gamma$ نیں،
  کیوں جے آخری استنباط وچ $!A$ خارج ہو جاندا اے۔ ایس لئی اسیں ثابت
  کرنا اے کہ $\Gamma \Entails !A \lif !B$۔ تضاد لئی منّ لو کہ کسے
  \iftag{FOL}{!!{structure}~$\Struct{M}$}{!!{valuation}~$\pAssign{v}$}
  لئی $\iftag{FOL}{\Sat{M}}{\pSat{v}}{\Gamma}$ پر
  $\iftag{FOL}{\Sat/{M}}{\pSat/{v}}{!A \lif !B}$۔ فیر
  $\iftag{FOL}{\Sat{M}}{\pSat{v}}{!A}$ تے
  $\iftag{FOL}{\Sat/{M}}{\pSat/{v}}{!B}$۔ پر مفروضے مطابق $!B$،
  $\Gamma \cup \{!A\}$ دا منطقی نتیجہ اے، یعنی
  $\iftag{FOL}{\Sat{M}}{\pSat{v}}{!B}$، جو تضاد اے۔ ایس لئی
  $\Gamma \Entails !A \lif !B$۔

\item آخری استنباط \FalseInt اے: ایتھے $\delta$ ایویں ختم ہوندی اے
  \begin{prooftree}
    \AxiomC{$\Gamma$}
    \RightLabel{$\delta_1$}
    \DeduceC{$\lfalse$}
    \RightLabel{\FalseInt}
    \UnaryInfC{$!A$}
  \end{prooftree}
  استقرائی مفروضے مطابق $\Gamma \Entails \lfalse$۔ اسیں ثابت کرنا
  اے کہ $\Gamma \Entails !A$۔ جے ایہ نہ ہووے، تاں کسے
  ~$\iftag{FOL}{\Struct{M}}{\pAssign{v}}$ لئی
    $\iftag{FOL}{\Sat{M}}{\pSat{v}}{\Gamma}$ تے
    $\iftag{FOL}{\Sat/{M}}{\pSat/{v}}{!A}$ ہووے گا۔ پر ہر حال وچ
    $\iftag{FOL}{\Sat/{M}}{\pSat/{v}}{\lfalse}$، ایس لئی ایہ
    $\Gamma \Entails/ \lfalse$ دا مطلب دیندا اے، جیہڑا استقرائی
    مفروضے دے خلاف اے۔

\item آخری استنباط \FalseCl اے: مشق۔

\iftag{FOL}{
\item آخری استنباط \Intro{\lforall} اے: فیر $\delta$ دی صورت ایہ اے
  \begin{prooftree}
    \AxiomC{$\Gamma$}
    \RightLabel{$\delta_1$}
    \DeduceC{$!A(a)$}
    \RightLabel{\Intro{\lforall}}
    \UnaryInfC{$\lforall[x][!A(x)]$}
  \end{prooftree}
  استقرائی مفروضے مطابق مقدمہ $!A(a)$، !!{undischarged}
  مفروضیاں $\Gamma$ دا منطقی نتیجہ اے۔ کوئی ساخت
  $\Struct{M}$ ویکھو جیہدے لئی $\Sat{M}{\Gamma}$۔ اسیں ثابت
  کرنا اے کہ $\Sat{M}{\lforall[x][!A(x)]}$۔ چونکہ
  $\lforall[x][!A(x)]$ !!a{sentence} اے، ایس دا مطلب اے کہ
  ہر متغیر-قدر بندی~$s$ لئی $\Sat{M}{!A(x)}[s]$ ثابت کیتا جاوے
  (\olref[syn][ass]{prop:sat-quant})۔ چونکہ $\Gamma$ سارے دا سارا
  جملیاں اُتے مشتمل اے، $\Sat{M}{!B}[s]$ ہر
  $!B \in \Gamma$ لئی قائم اے، بحوالہ
  \olref[syn][sat]{defn:satisfaction}۔ $\Struct{M'}$ اوہو جہی ساخت
  ہووے جیہڑی $\Struct{M}$ توں صرف ایس گل وچ وکھری اے کہ
  $\Assign{a}{M'} = s(x)$۔ چونکہ $a$،~$\Gamma$ وچ نہیں آؤندا،
  \olref[syn][ext]{cor:extensionality-sent} مطابق
  $\Sat{M'}{\Gamma}$۔ چونکہ $\Gamma \Entails
  !A(a)$، ایس لئی $\Sat{M'}{!A(a)}$۔ چونکہ $!A(a)$ !!a{sentence}
  اے، \olref[syn][ass]{prop:sentence-sat-true} مطابق
  $\Sat{M'}{!A(a)}[s]$۔ \olref[syn][ext]{prop:ext-formulas} مطابق
  $\Sat{M'}{!A(x)}[s]$ اُدوں تے صرف اُدوں جدوں
  $\Sat{M'}{!A(a)}$ (چیتے رکھو کہ $!A(a)$، صرف
  $\Subst{!A(x)}{a}{x}$ اے)۔ ایس لئی
  $\Sat{M'}{!A(x)}[s]$۔ چونکہ $a$،~$!A(x)$ وچ نہیں آؤندا،
  \olref[syn][ext]{prop:extensionality} مطابق
  $\Sat{M}{!A(x)}[s]$۔ پر $s$ کوئی وی متغیر-قدر بندی سی،
  ایس لئی $\Sat{M}{\lforall[x][!A(x)]}$۔
  
\item آخری استنباط \Intro{\lexists} اے: مشق۔

\item آخری استنباط \Elim{\forall} اے: مشق۔
}{}
\end{enumerate}

ہُن کئی مقدمیاں والے ممکنہ استنباط ویکھدے آں:
\Elim{\lor}، \Intro{\land}، \iftag{FOL}{\Elim{\lif}، تے
  \Elim{\lexists}}{تے \Elim{\lif}}۔
\begin{enumerate}

\item آخری استنباط \Intro{\land} اے۔ $!A \land !B$ مقدمیاں $!A$
  تے $!B$ توں اخذ ہوندا اے، تے $\delta$ دی صورت ایہ اے
  \begin{prooftree}
    \AxiomC{$\Gamma_1$}
    \RightLabel{$\delta_1$}
    \DeduceC{$!A$}
    \AxiomC{$\Gamma_2$}
    \RightLabel{$\delta_2$}
    \DeduceC{$!B$}
    \RightLabel{\Intro{\land}}
    \BinaryInfC{$!A \land !B$}
  \end{prooftree}
  استقرائی مفروضے مطابق $!A$، !!{undischarged}
  مفروضیاں~$\Gamma_1$ توں نکلدا اے جیہڑے~$\delta_1$ دے نیں،
  تے $!B$، !!{undischarged} مفروضیاں~$\Gamma_2$ توں نکلدا اے
  جیہڑے~$\delta_2$ دے نیں۔
  ~$\delta$ دے !!{undischarged} مفروضے $\Gamma_1 \cup
  \Gamma_2$ نیں، ایس لئی اسیں ثابت کرنا اے کہ
  $\Gamma_1 \cup \Gamma_2 \Entails
  !A \land !B$۔ کوئی
  \iftag{FOL}{!!a{structure}~$\Struct{M}$}{!!a{valuation}~$\pAssign{v}$}
  ویکھو جیہدے لئی
  $\iftag{FOL}{\Sat{M}}{\pSat{v}}{\Gamma_1 \cup \Gamma_2}$۔
  چونکہ $\iftag{FOL}{\Sat{M}}{\pSat{v}}{\Gamma_1}$،
  $\iftag{FOL}{\Sat{M}}{\pSat{v}}{!A}$ ہووے گا، کیوں جے
  $\Gamma_1 \Entails !A$؛ تے چونکہ
  $\iftag{FOL}{\Sat{M}}{\pSat{v}}{\Gamma_2}$،
  $\iftag{FOL}{\Sat{M}}{\pSat{v}}{!B}$ ہووے گا، کیوں جے
  $\Gamma_2 \Entails
  !B$۔
  دونیں ملا کے
  $\iftag{FOL}{\Sat{M}}{\pSat{v}}{!A \land !B}$۔
  
\item آخری استنباط \Elim{\lor} اے: مشق۔

\item آخری استنباط \Elim{\lif} اے۔ $!B$، مقدمیاں
  $!A \lif !B$ تے~$!A$ توں اخذ ہوندا اے۔
  !!{derivation}~$\delta$ ایویں اے:
  \begin{prooftree}
    \AxiomC{$\Gamma_1$}
    \RightLabel{$\delta_1$}
    \DeduceC{$!A \lif !B$}
    \AxiomC{$\Gamma_2$}
    \RightLabel{$\delta_2$}
    \DeduceC{$!A$}
    \RightLabel{\Elim{\lif}}
    \BinaryInfC{$!B$}
  \end{prooftree}
  استقرائی مفروضے مطابق $!A \lif !B$، !!{undischarged}
  مفروضیاں~$\Gamma_1$ توں نکلدا اے جیہڑے~$\delta_1$ دے نیں،
  تے $!A$، !!{undischarged} مفروضیاں~$\Gamma_2$ توں نکلدا اے
  جیہڑے~$\delta_2$ دے نیں۔ کوئی
  \iftag{FOL}{!!a{structure}~$\Struct{M}$}{!!a{valuation}~$\pAssign{v}$}
  ویکھو۔ اسیں ثابت کرنا اے کہ جے
  $\iftag{FOL}{\Sat{M}}{\pSat{v}}{\Gamma_1 \cup
    \Gamma_2}$، تاں $\iftag{FOL}{\Sat{M}}{\pSat{v}}{!B}$۔
  منّ لو $\iftag{FOL}{\Sat{M}}{\pSat{v}}{\Gamma_1 \cup \Gamma_2}$۔
  چونکہ $\Gamma_1 \Entails !A \lif !B$،
  $\iftag{FOL}{\Sat{M}}{\pSat{v}}{!A
    \lif !B}$۔ چونکہ $\Gamma_2 \Entails !A$، اسیں
  $\iftag{FOL}{\Sat{M}}{\pSat{v}}{!A}$ حاصل کردے آں۔ ایس دا
  مطلب اے کہ $\iftag{FOL}{\Sat{M}}{\pSat{v}}{!B}$ (کیوں جے جے
  $\iftag{FOL}{\Sat/{M}}{\pSat/{v}}{!B}$ ہوندا، تاں
  $\iftag{FOL}{\Sat{M}}{\pSat{v}}{!A}$ دے نال
  $\iftag{FOL}{\Sat/{M}}{\pSat/{v}}{!A \lif !B}$ ہوندا،
  جیہڑا~$\iftag{FOL}{\Sat{M}}{\pSat{v}}{!A \lif !B}$
  دے خلاف اے)۔

\item آخری استنباط \Elim{\lnot} اے: مشق۔
  
\tagitem{FOL}{آخری استنباط \Elim{\lexists} اے: مشق۔}{}
\end{enumerate}
\end{proof}

\tagprob{FOL}
\begin{prob}
\olref[fol][ntd][sou]{thm:soundness} دا ثبوت مکمل کرو۔
\end{prob}
\tagendprob

\tagprob{notFOL}
\begin{prob}
\olref[pl][ntd][sou]{thm:soundness} دا ثبوت مکمل کرو۔
\end{prob}
\tagendprob

\begin{cor}
\ollabel{cor:weak-soundness}
جے $\Proves !A$، تاں $!A$ \iftag{FOL}{معتبر}{ہمیشہ صادق فارمولا} اے۔
\end{cor}

\begin{cor}
\ollabel{cor:consistency-soundness}
جے $\Gamma$ پورا ہون جوگ اے، تاں ایہ بے تضاد اے۔
\end{cor}

\begin{proof}
اسیں متقابلِ عکس ثابت کردے آں۔ منّ لو $\Gamma$ بے تضاد نہیں۔
فیر $\Gamma \Proves \lfalse$، یعنی~$\lfalse$ دی !!a{derivation}
موجود اے جیہدے !!{undischarged} مفروضے~$\Gamma$ وچ نیں۔
\olref{thm:soundness} مطابق ہر
\iftag{FOL}{!!{structure}~$\Struct{M}$}{!!{valuation}~$\pAssign{v}$}
جیہڑا $\Gamma$ نوں پورا کردا اے، اوہنوں $\lfalse$ وی پورا کرنا
پئے گا۔ چونکہ
$\iftag{FOL}{\Sat/{M}}{\pSat/{v}}{\lfalse}$ ہر
\iftag{FOL}{!!{structure}~$\Struct{M}$}{!!{valuation}~$\pAssign{v}$}
لئی قائم اے،
کوئی وی \iftag{FOL}{$\Struct{M}$}{$\pAssign{v}$}،
$\Gamma$ نوں پورا نہیں کر سکدا؛ یعنی $\Gamma$ پورا ہون جوگ نہیں۔
\end{proof}
\endgroup
% /OLP-0095

% OLP-0096 content/first-order-logic/natural-deduction/identity.tex
\begingroup
\def\olfilename{identity}
\expandafter\def\csname ol@sectioncs\endcsname{section}

\olfileid{fol}{ntd}{ide}

\olsection{\usetoken{P}{derivation} وچ \usetoken{S}{identity}}

!!^{derivation}s وچ !!{identity} لئی وادھو استنباطی قاعدے چاہیدے نیں۔

\begin{defish}
\AxiomC{}
\RightLabel{\Intro{\eq}}
\UnaryInfC{$\eq[t][t]$}
\DisplayProof
\hfill
\begin{tabular}{r}
\AxiomC{$\eq[t_1][t_2]$}
\AxiomC{$!A(t_1)$}
\RightLabel{\Elim{\eq}}
\BinaryInfC{$!A(t_2)$}
\DisplayProof
\\[3ex]
\AxiomC{$\eq[t_1][t_2]$}
\AxiomC{$!A(t_2)$}
\RightLabel{\Elim{\eq}}
\BinaryInfC{$!A(t_1)$}
\DisplayProof
\end{tabular}
\end{defish}

اُپر دتے قاعدیاں وچ $t$، $t_1$ تے $t_2$ بند ٹرم نیں۔
\Intro{\eq} قاعدہ بنا کسے مفروضے دے سِدھا $\eq[t][t]$ دے روپ دا
کوئی وی برابری دا بیان !!{derive} کرن دی اجازت دیندا اے۔

\begin{ex}
جے $s$ تے $t$ بند ٹرم نیں، تاں $!A(s), \eq[s][t] \Proves !A(t)$:
\begin{prooftree}
\AxiomC{$\eq[s][t]$}
\AxiomC{$!A(s)$}
\RightLabel{$\Elim{\eq}$}
\BinaryInfC{$!A(t)$}
\end{prooftree}
ایہ تہاڈے لئی ”اکو ہون والیاں چیزاں دی اک دوجے دی تھاں ورتوں دا اصول
\pnbLatin{(substitutability of identicals)}“، یا لائبنز دا قانون \pnbLatin{(Leibniz' Law)}،
دے ناں نال جانیا پہچانیا ہو سکدا اے۔
\end{ex}

\begin{prob}
ثابت کرو کہ $=$ متناظر (symmetric) تے متعدی (transitive) دونیں اے،
یعنی $\lforall[x][\lforall[y][(\eq[x][y] \lif
    \eq[y][x])]]$ تے $\lforall[x][\lforall[y][\lforall[z]((\eq[x][y]
    \land \eq[y][z]) \lif \eq[x][z])]]$ دیاں
!!{derivation}s بناؤ۔
\end{prob}

\begin{ex}
اسیں ایہ !!{sentence} !!{derive} کردے آں
\begin{align*}
& \lforall[x][\lforall[y][((!A(x) \land !A(y)) \lif \eq[x][y])]]
\intertext{!!{sentence} توں}
& \lexists[x][\lforall[y][(!A(y) \lif \eq[y][x])]]
\end{align*}
اسیں !!{derivation} نوں اُلٹے پاسے توں بناندے آں:
\begin{prooftree}
\AxiomC{$\lexists[x][\lforall[y][(!A(y) \lif \eq[y][x])]]
\quad \Discharge{!A(a) \land !A(b)}{1}$}
\DeduceC{$\eq[a][b]$}
\DischargeRule{\Intro{\lif}}{1}
\UnaryInfC{$((!A(a) \land !A(b)) \lif \eq[a][b])$}
\RightLabel{\Intro{\lforall}}
\UnaryInfC{$\lforall[y][((!A(a) \land !A(y)) \lif \eq[a][y])]$}
\RightLabel{\Intro{\lforall}}
\UnaryInfC{$\lforall[x][\lforall[y][((!A(x) \land !A(y)) \lif \eq[x][y])]]$}
\end{prooftree}
ہُن سانوں مرکزی مفروضہ ورتنا پئے گا: چونکہ ایہ اک وجودی
!!{formula} اے، اسیں \Elim{\lexists} ورت کے وچلا نتیجہ
$\eq[a][b]$ !!{derive} کردے آں۔
\begin{prooftree}
\AxiomC{$\lexists[x][\lforall[y][(!A(y) \lif \eq[y][x])]]$}
\AxiomC{$\Discharge{\lforall[y][(!A(y) \lif \eq[y][c]])}{2}$}
\noLine
\UnaryInfC{$\Discharge{!A(a) \land !A(b)}{1}$}
\DeduceC{$\eq[a][b]$}
\DischargeRule{\Elim{\lexists}}{2}
\BinaryInfC{$\eq[a][b]$}
\DischargeRule{\Intro{\lif}}{1}
\UnaryInfC{$((!A(a) \land !A(b)) \lif \eq[a][b])$}
\RightLabel{\Intro{\lforall}}
\UnaryInfC{$\lforall[y][((!A(a) \land !A(y)) \lif \eq[a][y])]$}
\RightLabel{\Intro{\lforall}}
\UnaryInfC{$\lforall[x][\lforall[y][((!A(x) \land !A(y)) \lif \eq[x][y])]]$}
\end{prooftree}
اُتّے سجے پاسے والی ذیلی !!{derivation} اپنے مفروضے ورت کے
$\eq[a][c]$ تے $\eq[b][c]$ وکھا کے مکمل ہوندی اے۔ ایس لئی دو
وکھریاں !!{derivation}s چاہیدیاں نیں۔ $\eq[a][c]$ دی
!!{derivation} ایویں اے:
\begin{prooftree}
\AxiomC{$\Discharge{\lforall[y][(!A(y) \lif \eq[y][c]])}{2}$}
\RightLabel{\Elim{\lforall}}
\UnaryInfC{$!A(a) \lif \eq[a][c]$}
\AxiomC{$\Discharge{!A(a) \land !A(b)}{1}$}
\RightLabel{\Elim{\land}}
\UnaryInfC{$!A(a)$}
\RightLabel{\Elim{\lif}}
\BinaryInfC{$\eq[a][c]$}
\end{prooftree}
$\eq[a][c]$ تے $\eq[b][c]$ توں اسیں $\eq[a][b]$
\Elim{\eq} نال !!{derive} کردے آں۔
\end{ex}

\begin{prob}
ہیٹھ دتے !!{formula}s دیاں !!{derivation}s بناؤ:
\begin{enumerate}
\item $\lforall[x][\lforall[y][((\eq[x][y] \land !A(x)) \lif !A(y))]]$
\item $\lexists[x][!A(x)] \land \lforall[y][\lforall[z][((!A(y) \land
    !A(z)) \lif \eq[y][z])]] \lif \lexists[x][(!A(x) \land
  \lforall[y][(!A(y) \lif \eq[y][x])])]$
\end{enumerate}
\end{prob}
\endgroup
% /OLP-0096
% OLP-0097 content/first-order-logic/natural-deduction/soundness-identity.tex
\begingroup
\def\olfilename{soundness-identity}
\expandafter\def\csname ol@sectioncs\endcsname{section}

\olfileid{fol}{ntd}{sid}

\olsection{\usetoken{S}{identity} نال صحت مندی}

\begin{prop}
$\eq$ دے قاعدیاں سمیت فطری استنباط صحت مند اے۔
\end{prop}

\begin{proof}
$\eq[t][t]$ دے روپ دا ہر !!{formula} منطقی طور تے معتبر اے، کیوں جے
ہر !!{structure}~$\Struct M$ لئی $\Sat{M}{\eq[t][t]}$۔ (چیتے رہے کہ
اسیں ٹرم $t$ نوں بند منّیا اے، یعنی ایس وچ کوئی متغیر نہیں، ایس لئی
متغیراں دیاں قدر بندیاں غیر متعلق نیں۔)

فرض کرو کہ کسے !!a{derivation} دا آخری استنباط \Elim{\eq} اے، یعنی
!!{derivation} دی ہیٹھ دتی صورت اے:
\begin{prooftree}
  \AxiomC{$\Gamma_1$}
  \RightLabel{$\delta_1$}
  \DeduceC{$\eq[t_1][t_2]$}
  \AxiomC{$\Gamma_2$}
  \RightLabel{$\delta_2$}
  \DeduceC{$!A(t_1)$}
  \RightLabel{\Elim{\eq}}
  \BinaryInfC{$!A(t_2)$}
\end{prooftree}
$\eq[t_1][t_2]$ تے $!A(t_1)$ دے مقدمے ترتیب وار !!{undischarged}
مفروضیاں~$\Gamma_1$ تے $\Gamma_2$ توں !!{derive} کیتے گئے نیں۔
اسیں وکھانا چاہندے آں کہ $!A(t_2)$، $\Gamma_1 \cup \Gamma_2$ توں
نکلدا اے۔ اوہ !!a{structure}~$\Struct{M}$ لو جیہدے لئی
$\Sat{M}{\Gamma_1 \cup
  \Gamma_2}$۔ استقرا دے مفروضے مطابق، $\Sat{M}{!A(t_1)}$ تے
$\Sat{M}{\eq[t_1][t_2]}$۔ ایس لئی
$\Value{t_1}{M} = \Value{t_2}{M}$۔ $s$ کوئی وی متغیر قدر بندی ہووے،
تے $m = \Value{t_1}{M} = \Value{t_2}{M}$۔
\olref[fol][syn][ext]{prop:ext-formulas} دے مطابق
$\Sat{M}{!A(t_1)}[s]$ تب تے صرف تب، جدوں
$\Sat{M}{!A(x)}[\Subst{s}{m}{x}]$، تے ایہ تب تے صرف تب، جدوں
$\Sat{M}{!A(t_2)}[s]$۔ چونکہ $\Sat{M}{!A(t_1)}$، ایس لئی
$\Sat{M}{!A(t_2)}$۔
\end{proof}
\endgroup
% /OLP-0097

\OLEndChapterHook
\endgroup
% /OLP-0084

% OLP-0098 content/first-order-logic/tableaux/tableaux.tex
\begingroup
\def\olfilename{tableaux}
\expandafter\def\csname ol@sectioncs\endcsname{section}

\iftag{FOL}
      {\olchapter{fol}{tab}{تابلو (Tableaux)}}
      {\olchapter{pl}{tab}{تابلو (Tableaux)}}

\begin{editorial}
ایہ باب صداقتی قدراں دیاں علامتاں والے تجزیاتی تابلو دا اک نظام پیش کردا اے۔
  
% PNBTAB-001: اصل عبارت وچ natural deduction اے؛ prfTab ایتھے تابلو لئی اے۔
ثبوتی نظام وجوں تابلو نال متعلقہ مواد شامل یا خارج کرن لئی
``prfTab'' ٹیگ ورتو۔
\end{editorial}

% OLP-0099 content/first-order-logic/tableaux/rules-and-proofs.tex
\begingroup
\def\olfilename{rules-and-proofs}
\expandafter\def\csname ol@sectioncs\endcsname{section}

\iftag{FOL}
      {\olfileid{fol}{tab}{rul}}
      {\olfileid{pl}{tab}{rul}}

\olsection{قاعدے تے \usetoken{P}{tableau}}

!!^a{tableau}، !!a{sentence} دے کسے !!a{structure} وچ سچ یا جھوٹ
ہون دے ممکنہ طریقیاں دا منظم جائزہ اے۔ تابلو دیاں بنیادی اکائیاں
!!{signed formula}s نیں: !!{sentence}s دے نال صداقتی قدر دی اک
``علامت''، یعنی $\True$ یا~$\False$۔ ایہ علامت لائے ہوئے
!!{formula}s اک (ہیٹھاں ول ودھدے) درخت وچ ترتیب دتے جاندے نیں۔

\begin{defn}
  اک \emph{!!{signed formula}} صداقتی قدر تے !!a{sentence} دی جوڑی اے، یعنی:
  \[
  \sFmla{\True}{!A} \text{ یا } \sFmla{\False}{!A}.
  \]
\end{defn}

معنی دے لحاظ نال اسیں $\sFmla{\True}{!A}$ نوں ``$!A$ سچ ہو
سکدا اے'' تے $\sFmla{\False}{!A}$ نوں ``$!A$ جھوٹ ہو سکدا اے''
پڑھ سکدے آں (کسے !!{structure} وچ)۔

درخت وچ ہر !!{signed formula} یا تاں \emph{مفروضہ} اے (جیہڑے درخت
دے سب توں اُتّے درج ہوندے نیں)، یا اوہدے اُتّے موجود کسے
!!a{signed formula} توں استنباط دے کئی قاعدیاں وچوں اک ورت کے حاصل
کیتا جاندا اے۔ پچھلے !!{formula} دے ہر ممکن !!{main operator} لئی دو
قاعدے نیں: اک جدوں علامت~$\True$ ہووے، تے اک جدوں علامت~$\False$
ہووے۔ کجھ قاعدے درخت نوں شاخاں وچ ونڈدے نیں، تے کجھ شاخ وچ صرف
!!{signed formula}s شامل کردے نیں۔ قاعدہ فوراً پچھلے
!!{signed formula} اُتے ہی لاگو کرنا ضروری نہیں؛ اوہ جڑ توں اوس جگہ تک والی
شاخ دے کسے وی !!{signed formula} اُتے لاگو ہو سکدا اے (تے اکثر
کرنا پیندا اے)۔

شاخ \emph{بند} اُدوں ہوندی اے جدوں اوہدے وچ $\sFmla{\True}{!A}$
تے $\sFmla{\False}{!A}$ دونیں ہون۔ بند !!{tableau} اوہ اے جیہدی ہر
شاخ بند ہووے۔ بدیہی تشریح مطابق، ہر شاخ اک اکٹھے ممکن ہون والی صورت
دسدی اے، پر $\sFmla{\True}{!A}$ تے $\sFmla{\False}{!A}$ اکٹھے
ممکن نہیں۔ ہور لفظاں وچ، جے شاخ بند اے، تاں اوہ امکان جیہڑا اوہ بیان کردی اے،
ردّ ہو چکیا اے۔ خاص طور تے ایس دا مطلب اے کہ بند !!{tableau} اوہ
سارے امکان ردّ کردا اے جنہاں وچ $\sFmla{\True}{!A}$ دے روپ دے
ہر مفروضے نوں سچ تے $\sFmla{\False}{!A}$ دے روپ دے ہر مفروضے نوں
جھوٹ، اکّو ویلے، بنایا جا سکدا ہووے۔

\emph{$!A$ لئی} بند !!{tableau} اوہ بند !!{tableau} اے جیہدی
جڑ~$\sFmla{\False}{!A}$ ہووے۔ جے ایہو جہا بند !!{tableau} موجود
ہووے، تاں $!A$ دے جھوٹ ہون دے سارے امکان ردّ ہو جاندے نیں؛ یعنی
$!A$ ہر !!{structure} وچ سچ ہونا لازمی اے۔
\endgroup
% /OLP-0099

% OLP-0100 content/first-order-logic/tableaux/propositional-rules.tex
\begingroup
\def\olfilename{propositional-rules}
\expandafter\def\csname ol@sectioncs\endcsname{section}

\iftag{FOL}
      {\olfileid{fol}{tab}{prl}}
      {\olfileid{pl}{tab}{prl}}

\olsection{قضیاتی قاعدے}

\subsection{$\lnot$ دے قاعدے}

\begin{defish}
\AxiomC{\sFmla{\True}{\lnot !A}}
\RightLabel{\TRule{\True}{\lnot}}
\UnaryInfC{\sFmla{\False}{!A}}
\DisplayProof
\hfill
\AxiomC{\sFmla{\False}{\lnot !A}}
\RightLabel{\TRule{\False}{\lnot}}
\UnaryInfC{\sFmla{\True}{!A}}
\DisplayProof
\end{defish}

\subsection{$\land$ دے قاعدے}

\begin{defish}\noindent
\AxiomC{\sFmla{\True}{!A \land !B}}
\RightLabel{\TRule{\True}{\land}}
\UnaryInfC{\sFmla{\True}{!A}}
\noLine
\UnaryInfC{\sFmla{\True}{!B}}
\DisplayProof
\hfill
\AxiomC{\sFmla{\False}{!A \land !B}}
\RightLabel{\TRule{\False}{\land}}
\UnaryInfC{$\sFmla{\False}{!A} \quad \mid \quad \sFmla{\False}{!B}$}
\DisplayProof
\end{defish}

\subsection{$\lor$ دے قاعدے}

\begin{defish}
\AxiomC{\sFmla{\True}{!A \lor !B}}
\RightLabel{\TRule{\True}{\lor}}
\UnaryInfC{$\sFmla{\True}{!A} \quad \mid \quad \sFmla{\True}{!B}$}
\DisplayProof
\hfill
\AxiomC{\sFmla{\False}{!A \lor !B}}
\RightLabel{\TRule{\False}{\lor}}
\UnaryInfC{\sFmla{\False}{!A}}
\noLine
\UnaryInfC{\sFmla{\False}{!B}}
\DisplayProof
\end{defish}

\subsection{$\lif$ دے قاعدے}

\begin{defish}
\AxiomC{\sFmla{\True}{!A \lif !B}}
\RightLabel{\TRule{\True}{\lif}}
\UnaryInfC{$\sFmla{\False}{!A} \quad \mid \quad \sFmla{\True}{!B}$}
\DisplayProof
\hfill
\AxiomC{\sFmla{\False}{!A \lif !B}}
\RightLabel{\TRule{\False}{\lif}}
\UnaryInfC{\sFmla{\True}{!A}}
\noLine
\UnaryInfC{\sFmla{\False}{!B}}
\DisplayProof
\end{defish}

\subsection{کٹ دا قاعدہ}

\begin{defish}
\AxiomC{}
\RightLabel{\Cut}
\UnaryInfC{$\sFmla{\True}{!A} \quad \mid \quad \sFmla{\False}{!A}$}
\DisplayProof
\end{defish}

\Cut{} دا قاعدہ پچھلے کسے !!{signed formula} ``اُتے'' لاگو نہیں
کیتا جاندا؛ بلکہ ایہ !!a{tableau} دی ہر شاخ نوں دو حصیاں وچ ونڈن دی
اجازت دیندا اے: اک شاخ وچ $\sFmla{\True}{!A}$ ہُندا اے، دوجی
وچ~$\sFmla{\False}{!A}$۔ ایہ ضروری نہیں---!!{signed formula}s
دے جس وی سیٹ لئی بند !!{tableau} ہووے، اوہدے لئی \Cut{} ورتے
بغیر وی اک بند تابلو موجود ہُندا اے---پر ایہ سانوں !!{tableau}s
نوں سَوکھے طریقے نال اکٹھے کرن دی اجازت دیندا اے۔
\endgroup
% /OLP-0100

% OLP-0101 content/first-order-logic/tableaux/quantifier-rules.tex
\begingroup
\def\olfilename{quantifier-rules}
\expandafter\def\csname ol@sectioncs\endcsname{section}

\olfileid{fol}{tab}{qrl}

\olsection{مسوّراں دے قاعدے}

\subsection{$\lforall$ لئی قاعدے}

\begin{defish}
\AxiomC{\sFmla{\True}{\lforall[x][!A(x)]}}
\RightLabel{\TRule{\True}{\forall}}
\UnaryInfC{\sFmla{\True}{!A(t)}}
\DisplayProof
\hfill
\AxiomC{\sFmla{\False}{\lforall[x][!A(x)]}}
\RightLabel{\TRule{\False}{\lforall}}
\UnaryInfC{\sFmla{\False}{!A(a)}}
\DisplayProof
\end{defish}

\TRule{\True}{\lforall} وچ $t$ اک بند ٹرم اے (یعنی جیہدے وچ
کوئی متغیر نہ ہووے)۔ \TRule{\False}{\lforall} وچ $a$~اک
!!a{constant} اے جیہڑا \TRule{\False}{\lforall} قاعدے توں اُتّے
والی شاخ وچ کتے وی نہیں آؤنا چاہیدا۔ اسیں $a$ نوں
\TRule{\False}{\forall} استنباط دا \emph{ایگن متغیر} آکھدے
آں۔\footnote{اسیں ``ایگن متغیر'' اصطلاح ورتدے آں، بھاویں اُتّے
والے قاعدے وچ $a$ اک !!a{constant} اے۔ ایس دے تاریخی سبب نیں۔}

\subsection{$\lexists$ لئی قاعدے}

\begin{defish}
\AxiomC{\sFmla{\True}{\lexists[x][!A(x)]}}
\RightLabel{\TRule{\True}{\lexists}}
\UnaryInfC{\sFmla{\True}{!A(a)}}
\DisplayProof
\hfill
\AxiomC{\sFmla{\False}{\lexists[x][!A(x)]}}
\RightLabel{\TRule{\False}{\lexists}}
\UnaryInfC{\sFmla{\False}{!A(t)}}
\DisplayProof
\end{defish}

ایتھے فیر، $t$ اک بند ٹرم اے، تے $a$ اک !!a{constant} اے جیہڑا
\TRule{\True}{\lexists} قاعدے توں اُتّے والی شاخ وچ نہیں آؤندا۔
اسیں $a$ نوں \TRule{\True}{\lexists} استنباط دا
\emph{ایگن متغیر} آکھدے آں۔

ایہ شرط کہ ایگن متغیر \TRule{\False}{\lforall} یا
\TRule{\True}{\lexists} استنباط توں اُتّے والی شاخ وچ نہ آوے،
\emph{ایگن متغیر دی شرط} آکھوَندی اے۔

\begin{explain}
ایہ رواج یاد کرو: جدوں $!A$ اک !!a{formula} اے جیہدے وچ
!!{variable}~$x$ آزاد ہووے، اسیں ایہ گل~$!A(x)$ لکھ کے ظاہر کردے
آں۔ اوسے سیاق وچ $!A(t)$،~$\Subst{!A}{t}{x}$ دا مختصر روپ اے۔
سو اسیں \TRule{\False}{\lexists} دا قاعدہ ایویں وی لکھ سکدے آں:
\begin{prooftree}
  \AxiomC{\sFmla{\False}{\lexists[x][!A]}}
  \RightLabel{\TRule{\False}{\lexists}}
  \UnaryInfC{\sFmla{\False}{\Subst{!A}{t}{x}}}
\end{prooftree}
یاد رکھو کہ $t$ پہلے ای~$!A$ وچ آ سکدا اے؛ مثلاً $!A$ دا روپ
~$\Atom{\Obj P}{t,x}$ ہو سکدا اے۔ ایس لئی
$\sFmla{\False}{\Atom{\Obj P}{t,t}}$ دا
~$ \sFmla{\False}{\lexists[x][\Atom{\Obj P}{t,x}]}$ توں
استنباط~\TRule{\False}{\lexists} دی درست ورتوں اے۔ پر
\TRule{\False}{\lforall} تے~\TRule{\True}{\lexists} دیاں
ایگن متغیر شرطاں منگدیاں نیں کہ !!{constant}~$a$،~$!A$ وچ نہ
آوے۔ ایس لئی
$\sFmla{\False}{\Atom{\Obj P}{a,a}}$ دا
$\sFmla{\False}{\lforall[x][\Atom{\Obj P}{a,x}]}$
توں~$\TRule{\False}{\lforall}$ ورت کے درست استنباط نہیں کر سکدے۔
\end{explain}

\begin{explain}
% PNBTAB-002: source's "no restrictions" is constrained by the earlier closed-term requirement; no further restriction on closed t.
\TRule{\True}{\lforall} تے \TRule{\False}{\lexists} وچ بند
ٹرم~$t$ اُتے ہور کوئی پابندی نہیں۔ دوجے پاسے،
\TRule{\True}{\lexists} تے \TRule{\False}{\lforall} دے قاعدیاں
وچ ایگن متغیر دی شرط منگدی اے کہ !!{constant}~$a$ متعلقہ استنباط
توں اُتّے والی شاخ وچ کتے وی نہ آوے۔ ایہ نظام دی صحت مندی یقینی
بناؤن لئی ضروری اے۔ ایس شرط توں بغیر ہِٹھلا
$\lexists[x][\formula{A}(x)] \lif \lforall[x][\formula{A}(x)]$
لئی بند !!{tableau} ہو جاندا:
\begin{center}
\begin{tableau}{}
  [\sFmla{\False}{\lexists[x][\formula{A}(x)] \lif \lforall[x][\formula{A}(x)]}, just=\TAss
    [\sFmla{\True}{\lexists[x][\formula{A}(x)]},
      just={\TRule{\False}{\lif}[1]}
      [\sFmla{\False}{\lforall[x][\formula{A}(x)]},
        just={\TRule{\False}{\lif}[1]}
        [\sFmla{\True}{\formula{A}(a)},
          just={\TRule{\True}{\lexists}[2]}
          [\sFmla{\False}{\formula{A}(a)},
            just={\TRule{\False}{\lforall}[3]}, close]
        ]
      ]
    ]
  ]
\end{tableau}
\end{center}
پر $\lexists[x][\formula{A}(x)] \lif
\lforall[x][\formula{A}(x)]$ معتبر نہیں۔
\end{explain}
\endgroup
% /OLP-0101

% OLP-0102 content/first-order-logic/tableaux/derivations.tex
\begingroup
\def\olfilename{derivations}
\expandafter\def\csname ol@sectioncs\endcsname{section}

\iftag{FOL}
      {\olfileid{fol}{tab}{der}}
      {\olfileid{pl}{tab}{der}}

\olsection{\usetoken{P}{tableau}}

\begin{explain}
اسیں دس چکے آں کہ مفروضہ کیہ اے، تے استنباط دے قاعدے وی دے چکے
آں۔ !!^{tableau}s ایہناں توں استقرائی طور تے بندے نیں: ہر
!!{tableau} یا تاں اک ایہو جیہی شاخ اے جیہڑی اک یا ودھ مفروضیاں
اُتے مشتمل ہووے، یا پھر شاخ اُتے استنباطی قاعدہ لا کے پہلے
!!a{tableau} توں حاصل ہوندا اے۔
\end{explain}

\begin{defn}[!!^{tableau}]
مفروضیاں $\sFmla{S_1}{!A_1}$، \dots،
$\sFmla{S_n}{!A_n}$ لئی !!^a{tableau} (جتھے ہر $S_i$ یا تاں
$\True$ اے یا~$\False$) !!{signed formula}s دا اک محدود درخت
اے جیہڑا ہِٹھ لکھیاں شرطاں پوری کردا اے:
\begin{enumerate}
\item درخت دے سب توں اُتّے والے $n$ عدد !!{signed formula}s،
  $\sFmla{S_i}{!A_i}$ نیں، اک دوجے دے ہِٹھاں۔
\item درخت وچ ہر اوہ !!{signed formula} جیہڑا مفروضیاں وچوں نہیں،
  اوہدے اُتّے والی شاخ دے کسے !!a{signed formula} اُتے استنباطی
  قاعدے دی درست ورتوں توں حاصل ہوندا اے۔
\end{enumerate}
!!a{tableau} دی شاخ اُدوں تے صرف اُدوں \emph{بند} اے جدوں اوہدے
وچ $\sFmla{\True}{!A}$ تے~$\sFmla{\False}{!A}$ دونیں
ہون، تے نہیں تاں \emph{کھلی} اے۔ جیہڑے !!^a{tableau} دی ہر شاخ
بند ہووے، اوہ (اپنے مفروضیاں دے سیٹ لئی) \emph{بند !!{tableau}}
اے۔ جے !!a{tableau} بند نہیں، یعنی اوہدے وچ گھٹ توں گھٹ اک کھلی
شاخ اے، تاں اوہ \emph{کھلا} اے۔
\end{defn}

\begin{ex}
  مفروضیاں دا ہر سیٹ اپنے آپ وچ !!a{tableau} اے، پر عام طور تے
  بند نہیں ہووے گا۔ (ظاہر اے کہ ایہ صرف اُدوں بند اے جدوں مفروضیاں
  وچ پہلے ای !!{signed formula}s دی جوڑی
  $\sFmla{\True}{!A}$ تے~$\sFmla{\False}{!A}$ موجود ہووے۔)

  % PNBTAB-003: bare A is the formula component of the signed formula; the sign remains present in the tableau.
  !!a{tableau} (کھلا یا بند) وچ موجود !!a{signed formula}، جیہدا
  فارمولا~$!A$ اے، اُتے استنباط دے قاعدیاں وچوں اک لا کے اس توں
  نواں، وڈا تابلو بنا سکدے آں۔ ایہ قاعدہ اوہدے لاگو ہون والی
  $!A$ دی آمد رکھن والی ہر شاخ دے سرے اُتے اک یا ودھ
  !!{signed formula}s شامل کرے گا۔

  مثال لئی، مفروضہ $\sFmla{\True}{!A \land \lnot
    !A}$ لو۔ صرف اوس مفروضے اُتے مشتمل (کھلا) !!{tableau} ایہ اے:
  \begin{center}
    \begin{tableau}{}
      [\sFmla{\True}{\formula{A} \land \lnot \formula{A}}, just=\TAss]
    \end{tableau}{}
  \end{center}
  ایس مفروضے اُتے $\TRule{\True}{\land}$ دا قاعدہ لا کے نواں
  !!{tableau} ملدا اے۔ ایہ قاعدہ سانوں !!{tableau} وچ دو نویاں
  لیکاں شامل کرن دی اجازت دیندا اے، $\sFmla{\True}{!A}$ تے
  $\sFmla{\True}{\lnot !A}$:
  \begin{center}
    \begin{tableau}{}
      [\sFmla{\True}{\formula{A} \land \lnot \formula{A}}, just=\TAss
        [\sFmla{\True}{\formula{A}}, just={\TRule{\True}{\land}[1]},
          [\sFmla{\True}{\lnot \formula{A}}, just={\TRule{\True}{\land}[1]}
          ]
        ]
      ]
    \end{tableau}{}
  \end{center}
  جدوں اسیں !!{tableau}s لکھدے آں تاں لاگو کیتے قاعدے سجے پاسے
  درج کردے آں (مثلاً $\TRule{\True}{\land} 1$ دا مطلب اے کہ
  اوس لکیر دا !!{signed formula}، $\TRule{\True}{\land}$ قاعدہ
  لکیر~$1$ دے !!{signed formula} اُتے لا کے
  حاصل ہویا اے)۔ ایس نویں !!{tableau} وچ ہور
  !!{signed formula}s نیں، پر صرف اک
  ($\sFmla{\True}{\lnot !A}$) اُتے اسیں قاعدہ لا سکدے آں
  (ایس صورت وچ $\TRule{\True}{\lnot}$ دا قاعدہ)۔ ایس توں بند
  !!{tableau} ملدا اے:
  \begin{center}
    \begin{tableau}{}
      [\sFmla{\True}{\formula{A} \land \lnot \formula{A}}, just=\TAss
        [\sFmla{\True}{\formula{A}}, just={\TRule{\True}{\land}[1]}
          [\sFmla{\True}{\lnot \formula{A}}, just={\TRule{\True}{\land}[1]}
            [\sFmla{\False}{\formula{A}}, just={\TRule{\True}{\lnot}[3]}, close]
          ]
        ]
      ]
    \end{tableau}{}
  \end{center}
\end{ex}
\endgroup
% /OLP-0102

% OLP-0103 content/first-order-logic/tableaux/proving-things.tex
\begingroup
\def\olfilename{proving-things}
\expandafter\def\csname ol@sectioncs\endcsname{section}

\iftag{FOL}
      {\olfileid{fol}{tab}{pro}}
      {\olfileid{pl}{tab}{pro}}

\olsection{\usetoken{P}{tableau} دیاں مثالاں}

\begin{ex}
آؤ !!{sentence} $(!A \land !B) \lif !A$ لئی اک بند !!{tableau} لبھیے۔

شروع وچ اس نال متعلق مفروضہ !!{tableau} دے سب توں اُتّے لکھدے آں۔
\begin{oltableau}
  [\sFmla{\False}{(\formula{A} \land \formula{B}) \lif \formula{A}},
    just = \TAss]
\end{oltableau}

مفروضہ صرف اک اے، ایس لئی قاعدہ لاگو کرن لئی وی صرف اک
!!{signed formula} اے۔ (ہر !!{signed formula} اُتے زیادہ توں
زیادہ اک ای قاعدہ لاگو ہو سکدا اے: اوہ !!{sentence} دی متعلقہ
علامت تے !!{main operator} دا قاعدہ اے۔) ایس صورت وچ سانوں
$\TRule{\False}{\lif}$ لاگو کرنا پینا اے۔
\begin{oltableau}
  [\sFmla{\False}{(\formula{A} \land \formula{B}) \lif \formula{A}},
    checked, just = \TAss
    [\sFmla{\True}{\formula{A} \land \formula{B}},
      just={\TRule{\False}{\lif}[1]}
      [\sFmla{\False}{\formula{A}}, just={\TRule{\False}{\lif}[1]}]
    ]
  ]
\end{oltableau}
ایہ یاد رکھن لئی کہ کِہڑے !!{signed formula}s اُتے اوہناں دے
متعلقہ قاعدے لا چکے آں، اسیں جملے دے کول اک صحیح دا نشان لکھدے
آں۔ پر صحیح دا نشان \emph{صرف} اُدوں لاؤ جدوں قاعدہ ساریاں
کھلی شاخاں اُتے لاگو ہو چکیا ہووے۔ جدوں !!a{signed formula}
دا متعلقہ قاعدہ ہر کھلی شاخ وچ لا دتا جاوے، سانوں مُڑ اوہدے
اُتے اوہ قاعدہ لاون دی لوڑ نہیں رہندی۔ ایتھے شاخ صرف اک اے،
ایس لئی قاعدہ وی اک واری لاونا پینا اے۔ (دھیان رہے کہ صحیح دے
نشان تابلو بنان وچ سہولت لئی نیں؛ اوہ تابلو دی باقاعدہ نحوی
ساخت دا حصہ نہیں۔)

ہن اک نواں !!{signed formula} اے جِہدے اُتے قاعدہ لا سکدے آں:
$\sFmla{\True}{!A \land !B}$، جیہڑا لکیر~$2$ اُتے اے۔
$\TRule{\True}{\land}$ دا قاعدہ لاون نال ایہ ملدا اے:
\begin{oltableau}
  [\sFmla{\False}{(\formula{A} \land \formula{B}) \lif \formula{A}},
    checked, just = \TAss
    [\sFmla{\True}{\formula{A} \land \formula{B}},
      just={\TRule{\False}{\lif}[1]}, checked
      [\sFmla{\False}{\formula{A}}, just={\TRule{\False}{\lif}[1]}
        [\sFmla{\True}{\formula{A}}, just={\TRule{\True}{\land}[2]}
          [\sFmla{\True}{\formula{B}}, just={\TRule{\True}{\land}[2]}, close
          ]
        ]
      ]
    ]
  ]
\end{oltableau}
کیوں جے ہن شاخ وچ $\sFmla{\True}{!A}$ (لکیر~$4$ اُتے)
تے $\sFmla{\False}{!A}$ (لکیر~$3$ اُتے) دونیں نیں، ایس لئی
شاخ بند اے۔ ایہو اک شاخ اے، ایس لئی !!{tableau} وی بند اے۔
اسیں $(!A \land !B) \lif !A$ لئی اک بند !!{tableau} لبھ لیا اے۔
\end{ex}

\begin{ex}
ہن $(\lnot !A \lor !B) \lif (!A
\lif !B)$ لئی اک بند !!{tableau} لبھیے۔

اس نال متعلق مفروضے توں شروع کردے آں:
\begin{oltableau}
  [\sFmla{\False}{(\lnot \formula{A} \lor \formula{B}) \lif
      (\formula{A} \lif \formula{B})}, just=\TAss]
\end{oltableau}
ایس !!{tableau} دے اِکو !!{signed formula} دا !!{main operator}~$\lif$
تے علامت~$\False$ اے، ایس لئی اوہدے اُتے $\TRule{\False}{\lif}$
دا قاعدہ لا کے ایہ حاصل کردے آں:
\begin{oltableau}
  [\sFmla{\False}{(\lnot \formula{A} \lor \formula{B})
      \lif (\formula{A} \lif \formula{B})}, just=\TAss, checked
    [\sFmla{\True}{\lnot \formula{A} \lor \formula{B}},
      just={\TRule{\False}{\lif}[1]}
      [\sFmla{\False}{(\formula{A} \lif \formula{B})},
        just={\TRule{\False}{\lif}[1]}
      ]
    ]
  ]
\end{oltableau}
ہن سانوں چُننا اے کہ~$\TRule{\True}{\lor}$ لکیر~$2$ اُتے
لائیے یا $\TRule{\False}{\lif}$ لکیر~$3$ اُتے۔ حقیقت وچ
ترتیب نال فرق نہیں پیندا، بشرطیکہ ہر !!{signed formula} دا
متعلقہ قاعدہ ہر شاخ وچ لاگو ہووے۔ سو پہلا قاعدہ چُندے آں۔
$\TRule{\True}{\lor}$ دا قاعدہ !!{tableau} نوں دو شاخاں وچ
ونڈن دی اجازت دیندا اے؛ قاعدے دے دو نتیجے ایہناں دو نویاں
شاخاں وچ شامل ہون والے نویں !!{signed formula}s ہون گے۔
ایس توں ایہ ملدا اے:
\begin{oltableau}
  [\sFmla{\False}{(\lnot \formula{A} \lor \formula{B}) \lif
      (\formula{A} \lif \formula{B})}, just=\TAss, checked
    [\sFmla{\True}{\lnot \formula{A} \lor \formula{B}},
      just={\TRule{\False}{\lif}[1]}, checked
      [\sFmla{\False}{(\formula{A} \lif \formula{B})},
        just={\TRule{\False}{\lif}[1]}
        [\sFmla{\True}{\lnot \formula{A}}, just={\TRule{\True}{\lor}[2]}]
        [\sFmla{\True}{\formula{B}}, just={\TRule{\True}{\lor}[2]}]
      ]
    ]
  ]
\end{oltableau}
اسیں $\TRule{\False}{\lif}$ دا قاعدہ لکیر~$3$ اُتے ہجے
نہیں لایا؛ ہن ایہ کردے آں۔ وقت بچان لئی دونوں شاخاں اُتے
اکو واری لا دندے آں۔ یاد رکھو کہ !!a{signed formula} دے کول
صحیح دا نشان صرف اُدوں لکھدے آں جدوں اوہدا متعلقہ قاعدہ ہر
کھلی شاخ وچ لا چکے ہوئیے۔ ایس لئی جِس جِس شاخ وچ اوہ
!!{signed formula} موجود ہووے، اوہدی انتہا اُتے قاعدہ لاونا
چنگا اے۔ ایس طرح بعد وچ وکھ وکھ شاخاں دے ہِٹھاں جا کے اوس
!!{signed formula} ول مُڑن دی لوڑ نہیں پَندی۔
\begin{oltableau}
  [\sFmla{\False}{(\lnot \formula{A} \lor \formula{B}) \lif
      (\formula{A} \lif \formula{B})}, just=\TAss, checked
    [\sFmla{\True}{\lnot \formula{A} \lor \formula{B}},
      just={\TRule{\False}{\lif}[1]}, checked
      [\sFmla{\False}{(\formula{A} \lif \formula{B})},
        just={\TRule{\False}{\lif}[1]}, checked
        [\sFmla{\True}{\lnot \formula{A}}, just={\TRule{\True}{\lor}[2]}
          [\sFmla{\True}{\formula{A}}, just={\TRule{\False}{\lif}[3]}
            [\sFmla{\False}{\formula{B}}, just={\TRule{\False}{\lif}[3]}]
          ]
        ]
        [\sFmla{\True}{\formula{B}}, just={\TRule{\True}{\lor}[2]}
          [\sFmla{\True}{\formula{A}}, just={\TRule{\False}{\lif}[3]}
            [\sFmla{\False}{\formula{B}}, just={\TRule{\False}{\lif}[3]}, close]
          ]
        ]
      ]
    ]
  ]
\end{oltableau}
سجی شاخ ہن بند اے۔ کھبی شاخ اُتے اسیں ہجے
$\TRule{\True}{\lnot}$ دا قاعدہ لکیر~$4$ نوں لا سکدے آں۔ ایس نال
$\sFmla{\False}{!A}$ ملدا اے تے کھبی شاخ وی بند ہو جاندی اے:
\begin{oltableau}
  [\sFmla{\False}{(\lnot \formula{A} \lor \formula{B}) \lif
      (\formula{A} \lif \formula{B})}, just=\TAss, checked
    [\sFmla{\True}{\lnot \formula{A} \lor \formula{B}},
      just={\TRule{\False}{\lif}[1]}, checked
      [\sFmla{\False}{(\formula{A} \lif \formula{B})},
        just={\TRule{\False}{\lif}[1]}, checked
        [\sFmla{\True}{\lnot \formula{A}}, just={\TRule{\True}{\lor}[2]}
          [\sFmla{\True}{\formula{A}}, just={\TRule{\False}{\lif}[3]}
            [\sFmla{\False}{\formula{B}}, just={\TRule{\False}{\lif}[3]}
              [\sFmla{\False}{\formula{A}},
                just={\TRule{\True}{\lnot}[4]}, close]
            ]
          ]
        ]
        [\sFmla{\True}{\formula{B}}, just={\TRule{\True}{\lor}[2]}
          [\sFmla{\True}{\formula{A}}, just={\TRule{\False}{\lif}[3]}
            [\sFmla{\False}{\formula{B}},
              just={\TRule{\False}{\lif}[3]}, close]
          ]
        ]
      ]
    ]
  ]
\end{oltableau}
\end{ex}

\begin{ex}
اسیں مفروضیاں دے طور اُتے کِنّے وی !!{signed formula}s لئی
!!{tableau}s بنا سکدے آں۔ اکثر اک توں ودھ شاخاں بنان والے
قاعدے لاون دی وی لوڑ ہوندی اے؛ تے عام صورت وچ !!a{tableau}
دیاں کِنّیاں وی شاخاں ہو سکدیاں نیں۔ مثال لئی،
$\{\sFmla{\True}{!A \lor (!B \land !C)}, \sFmla{\False}{(!A \lor !B) \land (!A
  \lor !C)}\}$ لئی !!a{tableau} ویکھو۔ پہلے مفروضے اُتے
$\TRule{\True}{\lor}$ لا کے شروع کردے آں:
\begin{oltableau}
  [\sFmla{\True}{\formula{A} \lor (\formula{B} \land \formula{C})},
    just=\TAss, checked
    [\sFmla{\False}{(\formula{A} \lor \formula{B}) \land
        (\formula{A} \lor \formula{C})}, just=\TAss
      [\sFmla{\True}{\formula{A}}, just={\TRule{\True}{\lor}[1]}]
      [\sFmla{\True}{\formula{B} \land \formula{C}},
        just={\TRule{\True}{\lor}[1]}]
    ]
  ]
\end{oltableau}
ہن $\TRule{\False}{\land}$ دا قاعدہ لکیر~$2$ اُتے لا
سکدے آں۔ ایہ دونوں شاخاں اُتے اکو واری لا کے لکیر~$2$ دے
کول صحیح دا نشان لا سکدے آں:
\begin{oltableau}
  [\sFmla{\True}{\formula{A} \lor (\formula{B} \land \formula{C})},
    just=\TAss, checked
    [\sFmla{\False}{(\formula{A} \lor \formula{B}) \land
        (\formula{A} \lor \formula{C})}, just=\TAss, checked
      [\sFmla{\True}{\formula{A}}, just={\TRule{\True}{\lor}[1]}
        [\sFmla{\False}{\formula{A} \lor \formula{B}},
          just={\TRule{\False}{\land}[2]}]
        [\sFmla{\False}{\formula{A} \lor \formula{C}},
          just={\TRule{\False}{\land}[2]}]
      ]
      [\sFmla{\True}{\formula{B} \land \formula{C}},
        just={\TRule{\True}{\lor}[1]}
        [\sFmla{\False}{\formula{A} \lor \formula{B}},
          just={\TRule{\False}{\land}[2]}]
        [\sFmla{\False}{\formula{A} \lor \formula{C}},
          just={\TRule{\False}{\land}[2]}]
      ]
    ]
  ]
\end{oltableau}
ہن $\TRule{\False}{\lor}$ اوہناں ساریاں شاخاں اُتے لا سکدے
آں جِہڑیاں وچ $!A \lor !B$ اے:
\begin{oltableau}
  [\sFmla{\True}{\formula{A} \lor (\formula{B} \land \formula{C})},
    just=\TAss, checked
    [\sFmla{\False}{(\formula{A} \lor \formula{B}) \land
        (\formula{A} \lor \formula{C})}, just=\TAss, checked
      [\sFmla{\True}{\formula{A}}, just={\TRule{\True}{\lor}[1]}
        [\sFmla{\False}{\formula{A} \lor \formula{B}},
          just={\TRule{\False}{\land}[2]}, checked
          [\sFmla{\False}{\formula{A}}, just={\TRule{\False}{\lor}[4]}
            [\sFmla{\False}{\formula{B}},
              just={\TRule{\False}{\lor}[4]}, close]
          ]
        ]
        [\sFmla{\False}{\formula{A} \lor \formula{C}},
          just={\TRule{\False}{\land}[2]}]
      ]
      [\sFmla{\True}{\formula{B} \land \formula{C}},
        just={\TRule{\True}{\lor}[1]}
        [\sFmla{\False}{\formula{A} \lor \formula{B}},
          just={\TRule{\False}{\land}[2]}, checked
          [\sFmla{\False}{\formula{A}}, just={\TRule{\False}{\lor}[4]}
            [\sFmla{\False}{\formula{B}},
              just={\TRule{\False}{\lor}[4]}]
          ]
        ]
        [\sFmla{\False}{\formula{A} \lor \formula{C}},
          just={\TRule{\False}{\land}[2]}]
      ]
    ]
  ]
\end{oltableau}
سب توں کھبّی شاخ ہن بند اے۔ ہن $\TRule{\False}{\lor}$
$!A \lor !C$ اُتے لائیے:
\begin{oltableau}
  [\sFmla{\True}{\formula{A} \lor (\formula{B} \land \formula{C})},
    just=\TAss, checked
    [\sFmla{\False}{(\formula{A} \lor \formula{B}) \land
        (\formula{A} \lor \formula{C})}, just=\TAss, checked
      [\sFmla{\True}{\formula{A}}, just={\TRule{\True}{\lor}[1]}
        [\sFmla{\False}{\formula{A} \lor \formula{B}},
          just={\TRule{\False}{\land}[2]}, checked
          [\sFmla{\False}{\formula{A}},
            just={\TRule{\False}{\lor}[4]}
            [\sFmla{\False}{\formula{B}},
              just={\TRule{\False}{\lor}[4]}, close]
          ]
        ]
        [\sFmla{\False}{\formula{A} \lor \formula{C}},
          just={\TRule{\False}{\land}[2]}, checked
          [\sFmla{\False}{\formula{A}},
            just={\TRule{\False}{\lor}[4]}, move by=2
            [\sFmla{\False}{\formula{C}},
              just={\TRule{\False}{\lor}[4]}, close]
          ]
        ]
      ]
      [\sFmla{\True}{\formula{B} \land \formula{C}},
        just={\TRule{\True}{\lor}[1]}
        [\sFmla{\False}{\formula{A} \lor \formula{B}},
          just={\TRule{\False}{\land}[2]}, checked
          [\sFmla{\False}{\formula{A}},
            just={\TRule{\False}{\lor}[4]}
            [\sFmla{\False}{\formula{B}},
              just={\TRule{\False}{\lor}[4]}
            ]
          ]
        ]
        [\sFmla{\False}{\formula{A} \lor \formula{C}},
          just={\TRule{\False}{\land}[2]}, checked
          [\sFmla{\False}{\formula{A}},
            just={\TRule{\False}{\lor}[4]}, move by=2
            [\sFmla{\False}{\formula{C}},
              just={\TRule{\False}{\lor}[4]}]
          ]
        ]
      ]
    ]
  ]
\end{oltableau}
دھیان دو کہ وضاحت لئی اسیں $\TRule{\False}{\lor}$ نوں دوجی
واری لاون دا نتیجہ ہِٹھاں منتقل کیتا اے۔ ایتھے ایہ ضروری نہیں
سی، کیوں جے قاعدے دے حوالے اوہی رہنَے سن۔

دو شاخاں ہجے کھلیاں نیں، تے
$\sFmla{\True}{!B \land !C}$، جیہڑا لکیر~$3$ اُتے اے،
اوہدے کول ہجے صحیح دا نشان نہیں۔
اوس اُتے $\TRule{\True}{\land}$ لا کے اک بند !!{tableau}
حاصل کردے آں:
\begin{oltableau}
  [\sFmla{\True}{\formula{A} \lor (\formula{B} \land \formula{C})},
    just=\TAss, checked
    [\sFmla{\False}{(\formula{A} \lor \formula{B}) \land
        (\formula{A} \lor \formula{C})}, just=\TAss, checked
      [\sFmla{\True}{\formula{A}}, just={\TRule{\True}{\lor}[1]}
        [\sFmla{\False}{\formula{A} \lor \formula{B}},
          just={\TRule{\False}{\land}[2]}, checked
          [\sFmla{\False}{\formula{A}},
            just={\TRule{\False}{\lor}[4]}
            [\sFmla{\False}{\formula{B}},
              just={\TRule{\False}{\lor}[4]}, close]
          ]
        ]
        [\sFmla{\False}{\formula{A} \lor \formula{C}},
          just={\TRule{\False}{\land}[2]}, checked
          [\sFmla{\False}{\formula{A}},
            just={\TRule{\False}{\lor}[4]}
            [\sFmla{\False}{\formula{C}},
              just={\TRule{\False}{\lor}[4]}, close]
          ]
        ]
      ]
      [\sFmla{\True}{\formula{B} \land \formula{C}},
        just={\TRule{\True}{\lor}[1]}, checked
        [\sFmla{\False}{\formula{A} \lor \formula{B}},
          just={\TRule{\False}{\land}[2]}, checked
          [\sFmla{\False}{\formula{A}},
            just={\TRule{\False}{\lor}[4]}
            [\sFmla{\False}{\formula{B}},
              just={\TRule{\False}{\lor}[4]}
              [\sFmla{\True}{\formula{B}},
                just={\TRule{\True}{\land}[3]}
                [\sFmla{\True}{\formula{C}},
                  just={\TRule{\True}{\land}[3]},close
                ]
              ]
            ]
          ]
        ]
        [\sFmla{\False}{\formula{A} \lor \formula{C}},
          just={\TRule{\False}{\land}[2]}, checked
          [\sFmla{\False}{\formula{A}},
            just={\TRule{\False}{\lor}[4]}
            [\sFmla{\False}{\formula{C}},
              just={\TRule{\False}{\lor}[4]}
              [\sFmla{\True}{\formula{B}},
                just={\TRule{\True}{\land}[3]}
                [\sFmla{\True}{\formula{C}},
                  just={\TRule{\True}{\land}[3]},close
                ]
              ]
            ]
          ]
        ]
      ]
    ]
  ]
\end{oltableau}

مقابلے لئی، مفروضیاں دے اوسے سیٹ دا ایہ بند !!{tableau}
ویکھو، جِہدے وچ قاعدے وکھری ترتیب نال لائے گئے نیں:
\begin{oltableau}
  [\sFmla{\True}{\formula{A} \lor (\formula{B} \land \formula{C})},
    just=\TAss, checked
    [\sFmla{\False}{(\formula{A} \lor \formula{B}) \land
        (\formula{A} \lor \formula{C})}, just=\TAss, checked
      [\sFmla{\False}{\formula{A} \lor \formula{B}},
        just={\TRule{\False}{\land}[2]}, checked
        [\sFmla{\False}{\formula{A}},
          just={\TRule{\False}{\lor}[3]}
          [\sFmla{\False}{\formula{B}},
            just={\TRule{\False}{\lor}[3]}
            [\sFmla{\True}{\formula{A}},
              just={\TRule{\True}{\lor}[1]},close
            ]
            [\sFmla{\True}{\formula{B} \land \formula{C}},
              just={\TRule{\True}{\lor}[1]}, checked
              [\sFmla{\True}{\formula{B}},
                just={\TRule{\True}{\land}[6]}
                [\sFmla{\True}{\formula{C}},
                  just={\TRule{\True}{\land}[6]},close
                ]
              ]
            ]
          ]
        ]
      ]
      [\sFmla{\False}{\formula{A} \lor \formula{C}},
        just={\TRule{\False}{\land}[2]}, checked
        [\sFmla{\False}{\formula{A}},
          just={\TRule{\False}{\lor}[3]}
          [\sFmla{\False}{\formula{C}},
            just={\TRule{\False}{\lor}[3]}
            [\sFmla{\True}{\formula{A}},
              just={\TRule{\True}{\lor}[1]},close
            ]
            [\sFmla{\True}{\formula{B} \land \formula{C}},
              just={\TRule{\True}{\lor}[1]}, checked
              [\sFmla{\True}{\formula{B}},
                just={\TRule{\True}{\land}[6]}
                [\sFmla{\True}{\formula{C}},
                  just={\TRule{\True}{\land}[6]},close
                ]
              ]
            ]
          ]
        ]
      ]
    ]
  ]
\end{oltableau}
\end{ex}

\begin{prob}
ہِٹھ دتیاں صورتاں لئی بند !!{tableau}s بناؤ:
\begin{enumerate}
\item $\sFmla{\True}{!A \land (!B \land !C)}, \sFmla{\False}{(!A \land !B) \land !C}$.
\item $\sFmla{\True}{!A \lor (!B \lor !C)}, \sFmla{\False}{(!A \lor !B) \lor !C}$.
\item $\sFmla{\True}{!A \lif (!B \lif !C)}, \sFmla{\False}{!B \lif (!A \lif !C)}$.
\item $\sFmla{\True}{!A}, \sFmla{\False}{\lnot\lnot !A}$.
\end{enumerate}
\end{prob}

\begin{prob}
ہِٹھ دتیاں صورتاں لئی بند !!{tableau}s بناؤ:
\begin{enumerate}
\item $\sFmla{\True}{(!A \lor !B) \lif !C}, \sFmla{\False}{!A \lif !C}$.
\item $\sFmla{\True}{(!A \lif !C) \land (!B \lif !C)}, \sFmla{\False}{(!A \lor !B) \lif !C}$.
\item $\sFmla{\False}{\lnot(!A \land \lnot !A)}$.
\item $\sFmla{\True}{!B \lif !A}, \sFmla{\False}{\lnot !A \lif \lnot !B}$.
\item $\sFmla{\False}{(!A \lif \lnot !A) \lif \lnot !A}$.
\item $\sFmla{\False}{\lnot(!A \lif !B) \lif \lnot !B}$.
\item $\sFmla{\True}{!A \lif !C}, \sFmla{\False}{\lnot (!A \land \lnot !C)}$.
\item $\sFmla{\True}{!A \land \lnot !C}, \sFmla{\False}{\lnot (!A \lif !C)}$.
% PNBTAB-004: the frozen source puts two premises inside one signed-formula argument; each is separately signed True here.
\item $\sFmla{\True}{!A \lor !B}, \sFmla{\True}{\lnot !B}, \sFmla{\False}{!A}$.
\item $\sFmla{\True}{\lnot !A \lor \lnot !B}, \sFmla{\False}{\lnot(!A \land !B)}$.
\item $\sFmla{\False}{(\lnot !A \land \lnot !B) \lif\lnot(!A \lor !B)}$.
\item $\sFmla{\False}{\lnot(!A \lor !B) \lif (\lnot !A \land \lnot !B)}$.
\end{enumerate}
\end{prob}

\begin{prob}
ہِٹھ دتیاں صورتاں لئی بند !!{tableau}s بناؤ:
\begin{enumerate}
\item $\sFmla{\True}{\lnot(!A \lif !B)}, \sFmla{\False}{!A}$.
\item $\sFmla{\True}{\lnot(!A \land !B)}, \sFmla{\False}{\lnot !A \lor \lnot !B}$.
\item $\sFmla{\True}{!A \lif !B}, \sFmla{\False}{\lnot !A \lor !B}$.
\item $\sFmla{\False}{\lnot \lnot !A \lif !A}$.
\item $\sFmla{\True}{!A \lif !B}, \sFmla{\True}{\lnot !A \lif !B}, \sFmla{\False}{!B}$.
\item $\sFmla{\True}{(!A \land !B) \lif !C}, \sFmla{\False}{(!A \lif !C) \lor (!B \lif !C)}$.
\item $\sFmla{\True}{(!A \lif !B) \lif !A}, \sFmla{\False}{!A}$.
\item $\sFmla{\False}{(!A \lif !B) \lor (!B \lif !C)}$.
\end{enumerate}
\end{prob}
\endgroup
% /OLP-0103

% OLP-0104 content/first-order-logic/tableaux/proving-things-quant.tex
\begingroup
\def\olfilename{proving-things-quant}
\expandafter\def\csname ol@sectioncs\endcsname{section}

\olfileid{fol}{tab}{prq}

\olsection{مسوّراں والے \usetoken{P}{tableau}}

\begin{ex}
مسوّراں نال کم کردیاں سانوں ایگن متغیر دی شرط نہ ٹُٹن دا خیال
رکھنا پیندا اے؛ کدی کدی ایس لئی کجھ استنباطاں دی ترتیب
بدلنی پیندی اے۔ عام طور تے ایگن متغیر دی شرط والے قاعدے
پہلاں لاون دی کوشش فائدہ دیندی اے (مکمل !!{tableau} وچ
اوہ اُتّے ول ہون گے)۔

آؤ ویکھیے کہ !!{sentence}
$\lexists[x][\lnot !A(x)] \lif \lnot \lforall[x][!A(x)]$
لئی !!a{tableau} کیویں دیندے آں۔ ہمیشہ وانگ پہلے مفروضہ لکھدے آں،
\begin{oltableau}
[\sFmla{\False}{\lexists[x][\lnot \formula{A}(x)]\lif \lnot
    \lforall[x][\formula{A}(x)]}, just=\TAss]
\end{oltableau}
کیوں جے !!{main operator} $\lif$ اے، اسیں
$\TRule{\False}{\lif}$ لاوندے آں:
\begin{oltableau}
[\sFmla{\False}{\lexists[x][\lnot \formula{A}(x)]\lif \lnot
    \lforall[x][\formula{A}(x)]}, just=\TAss, checked
  [\sFmla{\True}{\lexists[x][\lnot \formula{A}(x)]},
    just={\TRule{\False}{\lif}[1]}
    [\sFmla{\False}{\lnot\lforall[x][\formula{A}(x)]},
      just={\TRule{\False}{\lif}[1]}]
  ]
]
\end{oltableau}
اگلی لکیر جِہدے نال نبڑنا اے، اوہ~$2$ اے۔ اسیں
$\TRule{\True}{\lexists}$ ورتدے آں۔ ایس لئی اک نواں
!!{constant} چاہیدا اے؛ کیونکہ ہجے کوئی !!{constant}s موجود
نہیں، اسیں کوئی وی چُن سکدے آں، مثلاً~$a$۔
\begin{oltableau}
[\sFmla{\False}{\lexists[x][\lnot \formula{A}(x)]\lif \lnot
    \lforall[x][\formula{A}(x)]}, just=\TAss, checked
  [\sFmla{\True}{\lexists[x][\lnot \formula{A}(x)]},
    just={\TRule{\False}{\lif}[1]}, checked
    [\sFmla{\False}{\lnot\lforall[x][\formula{A}(x)]},
      just={\TRule{\False}{\lif}[1]}
      [\sFmla{\True}{\lnot\formula{A}(a)}, just={\TRule{\True}{\lexists}[2]}]
    ]
  ]
]
\end{oltableau}
ہن $\TRule{\False}{\lnot}$ لکیر~$3$ اُتے لاوندے آں:
\begin{oltableau}
[\sFmla{\False}{\lexists[x][\lnot \formula{A}(x)]\lif \lnot
    \lforall[x][\formula{A}(x)]}, just=\TAss, checked
  [\sFmla{\True}{\lexists[x][\lnot \formula{A}(x)]},
    just={\TRule{\False}{\lif}[1]}, checked
    [\sFmla{\False}{\lnot\lforall[x][\formula{A}(x)]},
      just={\TRule{\False}{\lif}[1]}, checked
      [\sFmla{\True}{\lnot\formula{A}(a)}, just={\TRule{\True}{\lexists}[2]}
        [\sFmla{\True}{\lforall[x][\formula{A}(x)]},
          just = {\TRule{\False}{\lnot}[3]}
        ]
      ]
    ]
  ]
]
\end{oltableau}
پہلاں $\TRule{\True}{\lnot}$ لکیر~$4$ اُتے، تے ایس توں بعد
$\TRule{\True}{\lforall}$ لکیر~$5$ اُتے لا کے اک بند
!!{tableau} حاصل کردے آں۔
\begin{oltableau}
[\sFmla{\False}{\lexists[x][\lnot \formula{A}(x)]\lif \lnot
    \lforall[x][\formula{A}(x)]}, just=\TAss, checked
  [\sFmla{\True}{\lexists[x][\lnot \formula{A}(x)]},
    just={\TRule{\False}{\lif}[1]}, checked
    [\sFmla{\False}{\lnot\lforall[x][\formula{A}(x)]},
      just={\TRule{\False}{\lif}[1]}, checked
      [\sFmla{\True}{\lnot\formula{A}(a)}, just={\TRule{\True}{\lexists}[2]}
        [\sFmla{\True}{\lforall[x][\formula{A}(x)]},
          just = {\TRule{\False}{\lnot}[3]}
          [\sFmla{\False}{\formula{A}(a)}, just={\TRule{\True}{\lnot}[4]}
            [\sFmla{\True}{\formula{A}(a)},
              just={\TRule{\True}{\lforall}[5]}, close
            ]
          ]
        ]
      ]
    ]
  ]
]
\end{oltableau}
\end{ex}

\begin{ex}
آؤ ویکھیے کہ ہِٹھلے سیٹ لئی !!a{tableau} کیویں دیندے آں:
\[
\sFmla{\False}{\lexists[x][!C(x,b)]},
\sFmla{\True}{\lexists[x][(!A(x) \land !B(x))]},
\sFmla{\True}{\lforall[x][(!B(x) \lif !C(x,b))]}.
\]
ہمیشہ وانگ مفروضیاں توں شروع کردے آں:
\begin{oltableau}
  [\sFmla{\False}{\lexists[x][\formula{C}(x,b)]}, just=\TAss
    [\sFmla{\True}{\lexists[x][(\formula{A}(x) \land \formula{B}(x))]},
      just=\TAss
      [\sFmla{\True}{\lforall[x][(\formula{B}(x) \lif \formula{C}(x,b))]},
        just=\TAss]
    ]
  ]
\end{oltableau}
ایگن متغیر دی شرط والا قاعدہ ہمیشہ پہلاں لاونا چاہیدا اے؛
ایتھے اوہ $\TRule{\True}{\lexists}$ اے جیہڑا
لکیر~$2$ اُتے لاگو ہووے گا۔ مفروضیاں وچ !!{constant}~$b$
موجود اے، ایس لئی وکھرا مستقل ورتنا پینا اے؛ مُڑ~$a$ چُن لیندے آں۔
\begin{oltableau}
  [\sFmla{\False}{\lexists[x][\formula{C}(x,b)]}, just=\TAss
    [\sFmla{\True}{\lexists[x][(\formula{A}(x) \land \formula{B}(x))]},
      just=\TAss, checked
      [\sFmla{\True}{\lforall[x][(\formula{B}(x) \lif \formula{C}(x,b))]},
        just=\TAss
        [\sFmla{\True}{\formula{A}(a) \land \formula{B}(a)},
          just={\TRule{\True}{\lexists}[2]}]
      ]
    ]
  ]
\end{oltableau}
% PNBTAB-006: rules without the eigenvariable condition still require the substitution term t to be closed.
جے ہن $\TRule{\False}{\lexists}$ لکیر~$1$ اُتے یا
$\TRule{\True}{\lforall}$ لکیر~$3$ اُتے لاوندے آں، تاں
فیصلہ کرنا پیندا اے کہ کِہڑا ٹرم~$t$، $x$ دی تھاں رکھیے۔
ایہناں قاعدیاں لئی ایگن متغیر دی شرط نہیں، ایس لئی کوئی وی
بند ٹرم چُن سکدے آں۔ کجھ صورتاں وچ قاعدہ وکھ وکھ~$t$s نال
کئی وار لاونا وی پی سکدا اے۔ پر عام طور تے !!{tableau} وچ
پہلے توں موجود ٹرماں وچوں اک چُننا فائدہ دیندا اے---ایتھے
$a$ تے~$b$---تے ایتھے اسیں اندازہ لا سکدے آں کہ $a$
بند شاخ تک ودھ آسانی نال لے جاوے گا۔
\begin{oltableau}
  [\sFmla{\False}{\lexists[x][\formula{C}(x,b)]}, just=\TAss
    [\sFmla{\True}{\lexists[x][(\formula{A}(x) \land \formula{B}(x))]},
      just=\TAss, checked
      [\sFmla{\True}{\lforall[x][(\formula{B}(x) \lif \formula{C}(x,b))]}, just=\TAss
        [\sFmla{\True}{\formula{A}(a) \land \formula{B}(a)}, just={\TRule{\True}{\lexists}[2]}
          [\sFmla{\False}{\formula{C}(a,b)}, just={\TRule{\False}{\lexists}[1]}
            [\sFmla{\True}{\formula{B}(a) \lif \formula{C}(a,b)},
              just={\TRule{\True}{\lforall}[3]}
            ]
          ]
        ]
      ]
    ]
  ]
\end{oltableau}
اسیں لکیر $1$ تے $3$ دے !!{signed formula}s کول صحیح دا نشان
نہیں لاوندے، کیوں جے اوہناں نوں مُڑ ورتنا پی سکدا اے۔ ہن
$\TRule{\True}{\land}$ لکیر~$4$ اُتے لاؤ:
\begin{oltableau}
  [\sFmla{\False}{\lexists[x][\formula{C}(x,b)]}, just=\TAss
    [\sFmla{\True}{\lexists[x][(\formula{A}(x) \land \formula{B}(x))]},
      just=\TAss, checked
      [\sFmla{\True}{\lforall[x][(\formula{B}(x) \lif \formula{C}(x,b))]}, just=\TAss
        [\sFmla{\True}{\formula{A}(a) \land \formula{B}(a)},
          just={\TRule{\True}{\lexists}[2]}, checked
          [\sFmla{\False}{\formula{C}(a,b)}, just={\TRule{\False}{\lexists}[1]}
            [\sFmla{\True}{\formula{B}(a) \lif \formula{C}(a,b)},
              just={\TRule{\True}{\lforall}[3]}
              [\sFmla{\True}{\formula{A}(a)}, just={\TRule{\True}{\land}[4]}
                [\sFmla{\True}{\formula{B}(a)}, just={\TRule{\True}{\land}[4]}
                ]
              ]
            ]
          ]
        ]
      ]
    ]
  ]
\end{oltableau}
جے ہن $\TRule{\True}{\lif}$ لکیر~$6$ اُتے لائیے،
تاں !!{tableau} بند ہو جاندا اے:
\begin{oltableau}
  [\sFmla{\False}{\lexists[x][\formula{C}(x,b)]}, just=\TAss
    [\sFmla{\True}{\lexists[x][(\formula{A}(x) \land \formula{B}(x))]},
      just=\TAss, checked
      [\sFmla{\True}{\lforall[x][(\formula{B}(x) \lif \formula{C}(x,b))]},
        just=\TAss
        [\sFmla{\True}{\formula{A}(a) \land \formula{B}(a)},
          just={\TRule{\True}{\lexists}[2]}, checked
          [\sFmla{\False}{\formula{C}(a,b)},
            just={\TRule{\False}{\lexists}[1]}
            [\sFmla{\True}{\formula{B}(a) \lif \formula{C}(a,b)},
              just={\TRule{\True}{\lforall}[3]}, checked
              [\sFmla{\True}{\formula{A}(a)},
                just={\TRule{\True}{\land}[4]}
                [\sFmla{\True}{\formula{B}(a)},
                  just={\TRule{\True}{\land}[4]}
                  [\sFmla{\False}{\formula{B}(a)},
                    just={\TRule{\True}{\lif}[6]},close
                  ]
                  [\sFmla{\True}{\formula{C}(a,b)},
                    just={\TRule{\True}{\lif}[6]},close
                  ]
                ]
              ]
            ]
          ]
        ]
      ]
    ]
  ]
\end{oltableau}


\end{ex}

\begin{ex}
ہِٹھلے سیٹ لئی !!a{tableau} بناوندے آں:
\[
\sFmla{\True}{\lforall[x][!A(x)]}, \sFmla{\True}{\lforall[x][!A(x)] 
  \lif \lexists[y][!B(y)]}, \sFmla{\True}{\lnot\lexists[y][!B(y)]}.
\]
ہمیشہ وانگ مفروضے لکھدے آں:
\begin{oltableau}
  [\sFmla{\True}{\lforall[x][\formula{A}(x)]}, just=\TAss
    [\sFmla{\True}{\lforall[x][\formula{A}(x)] \lif
        \lexists[y][\formula{B}(y)]}, just=\TAss
      [\sFmla{\True}{\lnot\lexists[y][\formula{B}(y)]}, just=\TAss
      ]
    ]
  ]
\end{oltableau}
شروع وچ $\TRule{\True}{\lnot}$ دا قاعدہ لکیر~$3$ اُتے
لاوندے آں۔ ``ایگن متغیر دی شرط والے قاعدے ہمیشہ پہلاں لاؤ''
دی اک نتیجہ وار ہدایت ایہ اے کہ ``ایگن متغیر دی شرط توں
بغیر مسوّراں دے قاعدے لوڑ پین تک ملتوی رکھو''۔ شاخاں وچ
ونڈن والے قاعدے وی ملتوی رکھو۔
\begin{oltableau}
  [\sFmla{\True}{\lforall[x][\formula{A}(x)]}, just=\TAss
    [\sFmla{\True}{\lforall[x][\formula{A}(x)] \lif
        \lexists[y][\formula{B}(y)]}, just=\TAss
      [\sFmla{\True}{\lnot\lexists[y][\formula{B}(y)]}, just=\TAss, checked
        [\sFmla{\False}{\lexists[y][\formula{B}(y)]}, just={\TRule{\True}{\lnot}[3]}]
      ]
    ]
  ]
\end{oltableau}
نویں لکیر~$4$ اُتے $\TRule{\False}{\lexists}$ چاہیدا اے؛
ایہ ایگن متغیر دی شرط توں بغیر مسوّر دا قاعدہ اے۔ سو اسیں
ایہنوں ملتوی کر کے $\TRule{\True}{\lif}$ لکیر~$2$ اُتے
لاوندے آں۔
\begin{oltableau}
  [\sFmla{\True}{\lforall[x][\formula{A}(x)]}, just=\TAss
    [\sFmla{\True}{\lforall[x][\formula{A}(x)] \lif
        \lexists[y][\formula{B}(y)]}, just=\TAss, checked
      [\sFmla{\True}{\lnot\lexists[y][\formula{B}(y)]}, just=\TAss, checked
        [\sFmla{\False}{\lexists[y][\formula{B}(y)]},
          just={\TRule{\True}{\lnot}[3]},
          [\sFmla{\False}{\lforall[x][\formula{A}(x)]}, just={\TRule{\True}{\lif}[2]}]
          [\sFmla{\True}{\lexists[y][\formula{B}(y)]}, just={\TRule{\True}{\lif}[2]}]
        ]
      ]
    ]
  ]
\end{oltableau}
دونیں نویں !!{signed formula}s اُتے ایگن متغیر دی شرط والے
قاعدے لاگو ہوندے نیں، ایس لئی اگلا کم ایہو اے:
\begin{oltableau}
  [\sFmla{\True}{\lforall[x][\formula{A}(x)]}, just=\TAss
    [\sFmla{\True}{\lforall[x][\formula{A}(x)] \lif
        \lexists[y][\formula{B}(y)]}, just=\TAss, checked
      [\sFmla{\True}{\lnot\lexists[y][\formula{B}(y)]}, just=\TAss, checked
        [\sFmla{\False}{\lexists[y][\formula{B}(y)]},
          just={\TRule{\True}{\lnot}[3]}
          [\sFmla{\False}{\lforall[x][\formula{A}(x)]}, just={\TRule{\True}{\lif}[2]},checked
            [\sFmla{\False}{\formula{A}(b)}, just={\TRule{\False}{\lforall}[5]}]
          ]
          [\sFmla{\True}{\lexists[y][\formula{B}(y)]}, just={\TRule{\True}{\lif}[2]},checked
            [\sFmla{\True}{\formula{B}(c)}, just={\TRule{\True}{\lexists}[5]}]
          ]
        ]
      ]
    ]
  ]
\end{oltableau}
% PNBTAB-005: the second closing rule is applied to line 4, not line 3 as the frozen source prose says.
شاخاں بند کرن لئی سانوں لکیر $1$ تے~$4$ دے !!{signed formula}s
ورتنے پین گے۔ اوہناں دے متعلقہ قاعدیاں (\TRule{\True}{\lforall}
تے \TRule{\False}{\lexists}) اُتے ایگن متغیر دی شرط نہیں،
ایس لئی مناسب بند ٹرم چُن سکدے آں۔ ایتھے اوہ ترتیب وار
$b$ تے~$c$ نیں۔
\begin{oltableau}
  [\sFmla{\True}{\lforall[x][\formula{A}(x)]}, just=\TAss
    [\sFmla{\True}{\lforall[x][\formula{A}(x)] \lif
        \lexists[y][\formula{B}(y)]}, just=\TAss, checked
      [\sFmla{\True}{\lnot\lexists[y][\formula{B}(y)]}, just=\TAss, checked
        [\sFmla{\False}{\lexists[y][\formula{B}(y)]},
          just={\TRule{\True}{\lnot}[3]}
          [\sFmla{\False}{\lforall[x][\formula{A}(x)]},
            just={\TRule{\True}{\lif}[2]},checked
            [\sFmla{\False}{\formula{A}(b)},
              just={\TRule{\False}{\lforall}[5]}
              [\sFmla{\True}{\formula{A}(b)},
                just={\TRule{\True}{\lforall}[1]},close
              ]
            ]
          ]
          [\sFmla{\True}{\lexists[y][\formula{B}(y)]},
            just={\TRule{\True}{\lif}[2]},checked
            [\sFmla{\True}{\formula{B}(c)},
              just={\TRule{\True}{\lexists}[5]}
              [\sFmla{\False}{\formula{B}(c)},
                just={\TRule{\False}{\lexists}[4]},close
              ]
            ]
          ]
        ]
      ]
    ]
  ]
\end{oltableau}
\end{ex}

\begin{prob}
ہِٹھ دتیاں صورتاں لئی بند !!{tableau}s بناؤ:
\begin{enumerate}
\item $\sFmla{\False}{(\lforall[x][!A(x)] \land \lforall[y][!B(y)])
\lif \lforall[z][(!A(z) \land !B(z))]}$.
\item $\sFmla{\False}{(\lexists[x][!A(x)] \lor \lexists[y][!B(y)])
\lif \lexists[z][(!A(z) \lor !B(z))]}$.
\item $\sFmla{\True}{\lforall[x][(!A(x) \lif !B)]},
\sFmla{\False}{\lexists[y][!A(y)] \lif !B}$.
\item $\sFmla{\True}{\lforall[x][\lnot !A(x)]},
\sFmla{\False}{\lnot\lexists[x][!A(x)]}$.
\item $\sFmla{\False}{\lnot\lexists[x][!A(x)] \lif \lforall[x][\lnot
!A(x)]}$.
\item $\sFmla{\False}{\lnot\lexists[x][\lforall[y][((!A(x,y) \lif
\lnot !A(y,y)) \land (\lnot !A(y,y) \lif !A(x,y)))]]}$.
\end{enumerate}
\end{prob}

\begin{prob}
ہِٹھ دتیاں صورتاں لئی بند !!{tableau}s بناؤ:
\begin{enumerate}
\item $\sFmla{\False}{\lnot\lforall[x][!A(x)] \lif \lexists[x][\lnot!A(x)]}$.
\item $\sFmla{\True}{(\lforall[x][!A(x)] \lif !B)}, \sFmla{\False}{\lexists[y][(!A(y) \lif !B)]}$.
\item $\sFmla{\False}{\lexists[x][(!A(x) \lif \lforall[y][!A(y)])]}$.
\end{enumerate}
\end{prob}
\endgroup
% /OLP-0104

% OLP-0105 content/first-order-logic/tableaux/proof-theoretic-notions.tex
\begingroup
\def\olfilename{proof-theoretic-notions}
\expandafter\def\csname ol@sectioncs\endcsname{section}

\iftag{FOL}
      {\olfileid{fol}{tab}{ptn}}
      {\olfileid{pl}{tab}{ptn}}
      
\olsection{ثبوتی-نظریاتی تصورات \pnbLatin{(Proof-Theoretic Notions)}}

\begin{editorial}
ایہ حصہ تابلواں لئی ثابت پذیری دے تعلق تے بے تضادی دیاں
تعریفاں اکٹھیاں کردا اے۔
\end{editorial}

\begin{explain}
جیویں اسیں کئی اہم معنوی تصورات (منطقی اعتبار، منطقی نتیجہ تے پورا ہون
جوگتا) دی تعریف دِتّی اے، ہُن اسیں اوہناں دے مطابق
\emph{ثبوتی-نظریاتی تصورات \pnbLatin{(proof-theoretic notions)}} دی تعریف دیندے آں۔
ایہناں دی تعریف !!{sentence}s دے !!{structure}s وچ پورا اترن ول رجوع
کر کے نہیں، بلکہ مخصوص بند !!{tableau}x دے موجود ہون ول رجوع کر کے
دِتّی جاندی اے۔ ایہ اک اہم دریافت سی کہ ایہ تصورات آپو وچ میل کھاندے
نیں۔ اوہناں دے میل کھان دی گل \emph{صحت مندی (soundness)} تے
\emph{تمامیت دے قضیاں \pnbLatin{(completeness theorems)}} دا مضمون اے۔
\end{explain}


\begin{defn}[قضیے (Theorems)]
!!{sentence}~$!A$ اک \emph{قضیہ (theorem)} اے جے
$\sFmla{\False}{!A}$ لئی کوئی بند !!{tableau} موجود اے۔ اسیں
$\Proves !A$ لکھدے آں جے $!A$ قضیہ اے، تے جے نہیں تاں
$\Proves/ !A$ لکھدے آں۔
\end{defn}

\begin{defn}[!!^{derivability}]
!!^a{sentence} $!A$، !!{sentence}s دے سیٹ~$\Gamma$ توں
\emph{!!{derivable}} اے، یعنی $\Gamma \Proves !A$، اُدوں تے صرف
اُدوں جدوں کوئی محدود سیٹ $\{!B_1, \dots, !B_n\} \subseteq \Gamma$
موجود ہووے تے ہیٹھ دِتّے سیٹ لئی بند !!{tableau} موجود ہووے:
\[
\{
  \sFmla{\False}{!A},
  \sFmla{\True}{!B_1}, \dots,
  \sFmla{\True}{!B_n}
  \}.
\]
جے $!A$، $\Gamma$ توں !!{derivable} نہیں اے، تاں اسیں
$\Gamma \Proves/ !A$ لکھدے آں۔
\end{defn}

\begin{defn}[بے تضادی (Consistency)]
!!{sentence}s دا سیٹ~$\Gamma$ \emph{متضاد (inconsistent)} اے اُدوں تے
صرف اُدوں جدوں کوئی محدود سیٹ $\{!B_1, \dots, !B_n\} \subseteq
\Gamma$ موجود ہووے تے ہیٹھ دِتّے سیٹ لئی بند !!{tableau} موجود ہووے:
\[
\{
  \sFmla{\True}{!B_1}, \dots,
  \sFmla{\True}{!B_n}  
  \}.
\]
جے $\Gamma$ متضاد نہیں اے، تاں اسیں اوہنوں \emph{بے تضاد
(consistent)} آکھدے آں۔
\end{defn}

\begin{prop}[انعکاسیت (Reflexivity)]
\ollabel{prop:reflexivity}
جے $!A \in \Gamma$ اے، تاں $\Gamma \Proves !A$۔
\end{prop}

\begin{proof}
جے $!A \in \Gamma$ اے، تاں $\{!A\}$، $\Gamma$ دا محدود ذیلی سیٹ
اے تے ایہ !!{tableau}
\begin{oltableau}
  [\sFmla{\False}{\formula{A}}, just = \TAss
    [\sFmla{\True}{\formula{A}}, just = \TAss,close]
  ]
\end{oltableau}
بند اے۔
\end{proof}
  

\begin{prop}[یک رُخیت (Monotonicity)]
\ollabel{prop:monotonicity}
جے $\Gamma \subseteq \Delta$ تے $\Gamma \Proves !A$، تاں $\Delta
\Proves !A$۔
\end{prop}

\begin{proof}
$\Gamma$ دا ہر محدود ذیلی سیٹ، $\Delta$ دا وی محدود ذیلی سیٹ اے۔
\end{proof}

\begin{prop}[تعدیت (Transitivity)]
\ollabel{prop:transitivity}
جے $\Gamma \Proves !A$ تے $\{!A\} \cup
\Delta \Proves !B$، تاں $\Gamma \cup \Delta \Proves !B$۔
\end{prop}

\begin{proof}
جے $\{!A\} \cup \Delta \Proves !B$، تاں کوئی محدود ذیلی سیٹ $\Delta_0 =
\{!C_1, \dots, !C_n\} \subseteq \Delta$ موجود اے جیہدے لئی
% PNBTAB-007: each D_i is a sentence in Gamma; the frozen source incorrectly uses D_m subseteq Gamma.
\begin{align*}
\{\sFmla{\False}{!B}, & \sFmla{\True}{!A}, \sFmla{\True}{!C_1},
\dots, \sFmla{\True}{!C_n}\}
\intertext{اک بند !!{tableau} موجود اے۔ جے $\Gamma \Proves !A$، تاں
  ایہو جیہے $!D_1$، \dots، $!D_m \in \Gamma$ موجود نیں کہ}
\{\sFmla{\False}{!A}, & \sFmla{\True}{!D_1},
\dots, \sFmla{\True}{!D_m}\}
\end{align*}
لئی بند !!{tableau} موجود اے۔

ہُن اوہ !!{tableau} ویکھو جیہدے مفروضے ایہ نیں:
\[
\sFmla{\False}{!B},
\sFmla{\True}{!C_1}, \dots, \sFmla{\True}{!C_n},
\sFmla{\True}{!D_1}, \dots, \sFmla{\True}{!D_m}.
\]
$!A$ اُتے \Cut{} قاعدہ لاؤ۔ ایس نال دو شاخاں بندیاں نیں: اک وچ
$\sFmla{\True}{!A}$ تے دوجی وچ $\sFmla{\False}{!A}$ اے۔ ایس طرح،
پہلی شاخ اُتے
\[
\{\sFmla{\False}{!B}, \sFmla{\True}{!A},
\sFmla{\True}{!C_1}, \dots, \sFmla{\True}{!C_n}\}
\]
دے سارے علامت لائے فارمولے موجود نیں۔ ایہناں مفروضیاں لئی بند !!{tableau}
موجود ہون کرکے اسیں اوہنوں ایس شاخ نال جوڑ سکدے آں؛
$\sFmla{\True}{!A}$ وچوں لنگھن والی ہر شاخ بند ہو جاندی اے۔ دوجی شاخ
اُتے
\[
\{\sFmla{\False}{!A}, \sFmla{\True}{!D_1}, \dots,
\sFmla{\True}{!D_m}\}
\]
دے سارے علامت لائے فارمولے موجود نیں، ایس لئی اوہدے نال وی متعلقہ بند
!!{tableau} جوڑ کے پوری دوسری شاخ بند کیتی جا سکدی اے۔ ایس توں
$\Gamma \cup \Delta \Proves !B$ ثابت ہوندا اے۔
\end{proof}

خیال کرو کہ ایس دا خاص مطلب ایہ وی اے: جے $\Gamma \Proves !A$ تے $!A
\Proves !B$، تاں $\Gamma \Proves !B$۔ ایہ وی نتیجہ نکلدا اے کہ جے $!A_1,
\dots, !A_n \Proves !B$ تے $\Gamma \Proves !A_i$ ہر~$i$ لئی، تاں
$\Gamma \Proves !B$۔

\begin{prop}
\ollabel{prop:incons}
$\Gamma$ متضاد اے اُدوں تے صرف اُدوں جدوں $\Gamma \Proves !A$
  ہر !!{sentence}~$!A$ لئی ہووے۔
\end{prop}

\begin{proof}
مشق۔
\end{proof}

\tagprob{FOL}
\begin{prob}
\olref[fol][tab][ptn]{prop:incons} ثابت کرو۔
\end{prob}
\tagendprob

\tagprob{notFOL}
\begin{prob}
\olref[pl][tab][ptn]{prop:incons} ثابت کرو۔
\end{prob}
\tagendprob

\begin{prop}[محدود گواہی (Compactness)]
  \ollabel{prop:proves-compact}
  \begin{enumerate}
  \item جے $\Gamma \Proves !A$، تاں کوئی محدود ذیلی سیٹ $\Gamma_0
    \subseteq \Gamma$ موجود اے جیہدے لئی $\Gamma_0 \Proves !A$۔
  \item جے $\Gamma$ دا ہر محدود ذیلی سیٹ بے تضاد اے، تاں $\Gamma$
    وی بے تضاد اے۔
  \end{enumerate}
\end{prop}

\begin{proof}
  \begin{enumerate}
    \item جے $\Gamma \Proves !A$، تاں کوئی محدود ذیلی سیٹ
      $\Gamma_0 = \{!B_1, \dots, !B_n\}$ موجود اے تے ایہناں لئی بند !!{tableau}
      موجود اے:
      \[
      \{\sFmla{\False}{!A}, \sFmla{\True}{!B_1}, \dots, \sFmla{\True}{!B_n}\}
      \]
      ایہ !!{tableau} $\Gamma_0 \Proves !A$ وی دکھاندا اے۔
    \item جے $\Gamma$ متضاد اے، تاں کسے محدود ذیلی سیٹ
      $\Gamma_0 = \{!B_1, \dots, !B_n\}$ لئی ہیٹھ دِتّے سیٹ دا بند
      !!{tableau} موجود اے:
      \[
      \{\sFmla{\True}{!B_1}, \dots, \sFmla{\True}{!B_n}\}
      \]
      ایہ بند !!{tableau} دکھاندا اے کہ $\Gamma_0$ متضاد اے۔
  \end{enumerate}
\end{proof}
\endgroup
% /OLP-0105

% OLP-0106 content/first-order-logic/tableaux/provability-consistency.tex
\begingroup
\def\olfilename{provability-consistency}
\expandafter\def\csname ol@sectioncs\endcsname{section}

\iftag{FOL}
      {\olfileid{fol}{tab}{prv}}
      {\olfileid{pl}{tab}{prv}}

\olsection{\usetoken{S}{derivability} تے بے تضادی \pnbLatin{(Derivability and Consistency)}}

ہُن اسیں !!{derivability} دے تعلق دیاں کئی خاصیتاں قائم کراں گے۔ ایہ
اپنے آپ وچ وی دلچسپ نیں، پر ہر اک تمامیت دے قضیے دے ثبوت وچ کردار
نبھاوے گی۔

\begin{prop}\ollabel{prop:provability-contr}
  جے $\Gamma \Proves !A$ تے $\Gamma \cup \{!A\}$ متضاد اے، تاں
  $\Gamma$ متضاد اے۔
\end{prop}

\begin{proof}
% PNBTAB-008: Gamma_1 is indexed by C_m in the displayed tableau, not C_n as the source definition says.
ایہو جیہے دو محدود ذیلی سیٹ $\Gamma_0 = \{!B_1, \dots, !B_n\}$ تے $\Gamma_1
  =\{!C_1, \dots, !C_m\} \subseteq \Gamma$ موجود نیں کہ
  \begin{align*}
    \{\sFmla{\False}{!A}, &
    \sFmla{\True}{!B_1}, \dots, \sFmla{\True}{!B_n}\} \\
    \{\sFmla{\True}{!A}, &
    \sFmla{\True}{!C_1}, \dots, \sFmla{\True}{!C_m}\}
  \end{align*}
  دوناں لئی بند !!{tableau}s موجود نیں۔ $!A$ اُتے \Cut{} قاعدہ لا کے
  اسیں ایہناں نوں اک بند !!{tableau} وچ جوڑ سکدے آں، جیہڑا دکھاندا اے کہ $\Gamma_0
  \cup \Gamma_1$ متضاد اے۔ چونکہ $\Gamma_0
  \subseteq \Gamma$ تے $\Gamma_1 \subseteq \Gamma$، سو $\Gamma_0 \cup
  \Gamma_1 \subseteq \Gamma$؛ ایس لئی $\Gamma$~متضاد اے۔
\end{proof}

\begin{prop}
\ollabel{prop:prov-incons}
$\Gamma \Proves !A$ اُدوں تے صرف اُدوں جدوں $\Gamma \cup \{\lnot !A\}$ متضاد اے۔
\end{prop}

\begin{proof}
پہلاں منّ لو $\Gamma \Proves !A$، یعنی ہیٹھ دِتّے سیٹ لئی بند
!!{tableau} موجود اے:
\[
\{\sFmla{\False}{!A},
\sFmla{\True}{!B_1}, \dots, \sFmla{\True}{!B_n}\}
\]
$\TRule{\True}{\lnot}$ قاعدہ ورت کے ایس نوں ہیٹھ دِتّے سیٹ دے
بند !!{tableau} وچ بدلیا جا سکدا اے:
\[
\{\sFmla{\True}{\lnot !A},
\sFmla{\True}{!B_1}, \dots, \sFmla{\True}{!B_n}\}.
\]

دوجے پاسے، جے پچھلے سیٹ لئی بند !!{tableau} موجود اے، تاں اوہنوں
پہلے سیٹ دے بند !!{tableau} وچ ایویں بدل سکدے آں: پہلے
مفروضے~$\sFmla{\True}{\lnot !A}$ اُتے \TRule{\True}{\lnot}
لا کے نکلے سارے فارمولے تے اوہ مفروضہ ہٹا دِیو، تے
$\sFmla{\False}{!A}$ مفروضے دے طور تے شامل کر دِیو۔ جے پہلے
تابلو دی کوئی شاخ $\sFmla{\True}{\lnot !A}$ اُتے
\TRule{\True}{\lnot} لاون دے نتیجے، یعنی $\sFmla{\False}{!A}$،
کرکے بند سی، تاں نویں !!{tableau} دی اوہدے مطابق شاخ وی بند اے۔
جے پراݨے تابلو دی کوئی شاخ مفروضہ $\sFmla{\True}{\lnot !A}$ تے
$\sFmla{\False}{\lnot !A}$ دوناں ہون کرکے بند سی، تاں
$\TRule{\False}{\lnot}$ نوں $\sFmla{\False}{\lnot !A}$ اُتے لا کے
$\sFmla{\True}{!A}$ حاصل کرو۔ شامل کیتے مفروضے
$\sFmla{\False}{!A}$ نال ایہ شاخ بند ہو جاندی اے۔
\end{proof}

\begin{prob}
ثابت کرو کہ $\Gamma \Proves \lnot !A$ اُدوں تے صرف اُدوں جدوں
$\Gamma \cup \{!A\}$ متضاد اے۔
\end{prob}

\begin{prop}\ollabel{prop:explicit-inc}
  جے $\Gamma \Proves !A$ تے $\lnot !A \in \Gamma$، تاں $\Gamma$
  متضاد اے۔
\end{prop}

\begin{proof}
  منّ لو $\Gamma \Proves !A$ تے $\lnot !A \in \Gamma$۔ تاں ایہو جیہے
  $!B_1$، \dots، $!B_n \in \Gamma$ موجود نیں کہ
  \[
  \{\sFmla{\False}{!A}, \sFmla{\True}{!B_1}, \dots, \sFmla{\True}{!B_n}\}
  \]
  دا بند تابلو موجود اے۔ مفروضہ \sFmla{\False}{!A} دی تھاں
  \sFmla{\True}{\lnot !A} رکھو، تے مفروضیاں توں بعد
% PNBTAB-009: True-negation applies to True not-A, not to False A as the frozen prose says.
  \sFmla{\True}{\lnot !A} اُتے \TRule{\True}{\lnot} لا کے اوہدا نتیجہ
  درج کرو۔ پراݨے !!{tableau} وچ جیہڑا وی !!{sentence} لکیر~$1$ دے
  حوالے نال جائز سی، اوہ ہُن نویں !!{tableau} وچ
% PNBTAB-010: n+1 initial assumptions followed by one derived line place False A at n+2.
  لکیر~$n+2$ دے حوالے نال جائز اے۔ ایس لئی جے پراݨا !!{tableau}
  بند سی تاں نواں وی بند اے۔ ایہ دکھاندا اے کہ $\Gamma$ متضاد اے،
  کیوں جے سارے مفروضے~$\Gamma$ وچ نیں۔
\end{proof}

\begin{prop}\ollabel{prop:provability-exhaustive}
  جے $\Gamma \cup \{!A\}$ تے $\Gamma \cup \{\lnot !A\}$ دونیں
  متضاد نیں، تاں $\Gamma$ متضاد اے۔
\end{prop}

\begin{proof}
  جے ایہو جیہے $!B_1$، \dots، $!B_n \in \Gamma$ تے $!C_1$، \dots،
  $!C_m \in \Gamma$ موجود نیں کہ
  \begin{align*}
    \{\sFmla{\True}{!A}, &
    \sFmla{\True}{!B_1}, \dots, \sFmla{\True}{!B_n}\} \text{ تے }\\
    \{\sFmla{\True}{\lnot !A}, &
    \sFmla{\True}{!C_1}, \dots, \sFmla{\True}{!C_m}\}
  \end{align*}
  دوناں لئی بند !!{tableau}s موجود نیں، تاں اسیں اک جوڑیا ہویا
  !!{tableau} بنا سکدے آں جیہڑا $\Gamma$ دا تضاد دکھاندا اے۔ اوہدے
  مفروضے $\sFmla{\True}{!B_1}$، \dots، $\sFmla{\True}{!B_n}$
  دے نال $\sFmla{\True}{!C_1}$، \dots، $\sFmla{\True}{!C_m}$
  رکھو تے فیر \Cut{} قاعدہ لاؤ۔ ایس نال دو شاخاں بندیاں نیں: اک
  $\sFmla{\True}{!A}$ نال شروع ہوندی اے، دوجی
  $\sFmla{\False}{!A}$ نال۔
  
  کھبّے پاسے پہلے !!{tableau} دا اوہ حصہ جوڑو جیہڑا اوہدے مفروضیاں توں
  ہِٹھاں اے۔ ایتھے ہر قاعدے دا اطلاق ہن وی ٹھیک اے، کیوں جے پہلے
  !!{tableau} دے سارے مفروضے، $\sFmla{\True}{!A}$ سمیت، موجود نیں۔
  ایس لئی $\sFmla{\True}{!A}$ دے ہِٹھاں ہر شاخ بند ہو جاندی اے۔
  
  سَجّے پاسے دوجے !!{tableau} دا مفروضیاں توں ہِٹھلا حصہ جوڑو، پر
  $\TRule{\True}{\lnot}$ نوں $\sFmla{\True}{\lnot !A}$ اُتے
  لاون دے سارے نتیجے ہٹا دِیو۔ $\TRule{\True}{\lnot}$ نوں
  $\sFmla{\True}{\lnot !A}$ اُتے لاون دا نتیجہ
  $\sFmla{\False}{!A}$ اے؛ ایہ وی موجود اے، کیوں جے جوڑے ہوئے
  !!{tableau} دے سَجّے پاسے \Cut{} قاعدے دا ایہی نتیجہ اے۔

  جے دوجے تابلو دی کوئی شاخ ایس کرکے بند سی کہ اوہدے وچ مفروضہ
  $\sFmla{\True}{\lnot !A}$ (جیہڑا جوڑے ہوئے !!{tableau} وچ ہن
  مفروضہ نہیں) تے $\sFmla{\False}{\lnot !A}$ دونیں سن، تاں
  $\TRule{\False}{\lnot}$ نوں $\sFmla{\False}{\lnot !A}$ اُتے لا کے
  $\sFmla{\True}{!A}$ حاصل کرو۔ ہُن جوڑے ہوئے !!{tableau} دی
  متعلقہ شاخ وی بند اے، کیوں جے اوہدے وچ \Cut{} قاعدے دا سَجّا نتیجہ
  $\sFmla{\False}{!A}$ وی موجود اے۔ جے دوجے !!{tableau} دی شاخ
  کسے ہور سبب توں بند سی، تاں جوڑے ہوئے !!{tableau} دی متعلقہ شاخ
  وی بند اے، کیوں جے پراݨے دوجے !!{tableau} دی شاخ اُتے
  $\sFmla{\True}{\lnot !A}$ توں وکھرے !!{signed formula}s وچوں ہر اک جوڑے
  ہوئے !!{tableau} دی متعلقہ شاخ اُتے وی موجود اے۔
  \end{proof}
\endgroup
% /OLP-0106

% OLP-0107 content/first-order-logic/tableaux/provability-propositional.tex
\begingroup
\def\olfilename{provability-propositional}
\expandafter\def\csname ol@sectioncs\endcsname{section}

\iftag{FOL}
      {\olfileid{fol}{tab}{ppr}}
      {\olfileid{pl}{tab}{ppr}}

\olsection{\usetoken{S}{derivability} تے قضیاتی رابطے \pnbLatin{(Derivability and the Propositional Connectives)}}

\begin{explain}
  اسیں قائم کردے آں کہ تابلو حساب دا !!{derivability} تعلق~$\Proves$
  قضیاتی رابطیاں نال جڑے کجھ بنیادی حقائق قائم کرن لئی اینا طاقتور اے،
  جیویں کہ $!A \land !B \Proves !A$ تے $!A, !A \lif !B \Proves !B$
  (قاعدۂ اثباتِ مقدم؛ modus ponens)۔ تمامیت دے قضیے دے ثبوت لئی ایہناں
  حقائق دی لوڑ اے۔
\end{explain}

\begin{prop}\ollabel{prop:provability-land}
  \begin{enumerate}
  \item \ollabel{prop:provability-land-left} دونیں $!A \land !B \Proves
    !A$ تے $!A \land !B \Proves !B$۔
  \item \ollabel{prop:provability-land-right} $!A, !B \Proves !A \land
    !B$۔
  \end{enumerate}
\end{prop}

\begin{proof}
  \begin{enumerate}
  \item دوناں $\{\sFmla{\False}{!A}, \sFmla{\True}{!A \land !B}\}$ تے
    $\{\sFmla{\False}{!B}, \sFmla{\True}{!A \land !B}\}$ لئی بند
    !!{tableau}s موجود نیں:
% PNBTAB-011: normalize eight malformed sFmla sign/formula argument pairs in the four conjunction/disjunction trees.
    \begin{oltableau}{}
      [\sFmla{\False}{\formula{A}}, just=\TAss
        [\sFmla{\True}{\formula{A} \land \formula{B}}, just=\TAss
          [\sFmla{\True}{\formula{A}},just={\TRule{\True}{\land}[2]}
            [\sFmla{\True}{\formula{B}},just={\TRule{\True}{\land}[2]}, close
            ]
          ]
        ]
      ]
    \end{oltableau}
    \begin{oltableau}{}
      [\sFmla{\False}{\formula{B}}, just=\TAss
        [\sFmla{\True}{\formula{A} \land \formula{B}}, just=\TAss
          [\sFmla{\True}{\formula{A}},just={\TRule{\True}{\land}[2]}
            [\sFmla{\True}{\formula{B}},just={\TRule{\True}{\land}[2]}, close
            ]
          ]
        ]
      ]
    \end{oltableau}
    \item $\{\sFmla{\True}{!A},
      \sFmla{\True}{!B}, \sFmla{\False}{!A \land !B}\}$ لئی بند
      !!{tableau} ایہ اے:
      \begin{oltableau}
        [\sFmla{\False}{\formula{A} \land \formula{B}}, just = \TAss
          [\sFmla{\True}{\formula{A}}, just = \TAss
            [\sFmla{\True}{\formula{B}}, just=\TAss
              [\sFmla{\False}{\formula{A}}, just = {\TRule{\False}{\land}[1]}, close]
              [\sFmla{\False}{\formula{B}}, just = {\TRule{\False}{\land}[1]}, close]
            ]
          ]
        ]
      \end{oltableau}
  \end{enumerate}
\end{proof}

\begin{prop}\ollabel{prop:provability-lor}
  \begin{enumerate}
  \item $\{!A \lor !B, \lnot !A, \lnot !B\}$ متضاد اے۔
  \item دونیں $!A \Proves !A \lor !B$ تے $!B \Proves !A \lor !B$۔
  \end{enumerate}
\end{prop}

\begin{proof}
  \begin{enumerate}
  \item اسیں $\{\sFmla{\True}{!A \lor !B},
    \sFmla{\True}{\lnot !A}, \sFmla{\True}{\lnot !B}\}$ لئی اک بند
    !!{tableau} دیندے آں:
    \begin{oltableau}
      [\sFmla{\True}{\formula{A} \lor \formula{B}}, just = \TAss
        [\sFmla{\True}{\lnot \formula{A}}, just = \TAss
          [\sFmla{\True}{\lnot \formula{B}}, just = \TAss
            [\sFmla{\False}{\formula{A}}, just = {\TRule{\True}{\lnot}[2]}
              [\sFmla{\False}{\formula{B}}, just = {\TRule{\True}{\lnot}[3]}
                [\sFmla{\True}{\formula{A}}, just = {\TRule{\True}{\lor}[1]}, close]
                [\sFmla{\True}{\formula{B}}, just = {\TRule{\True}{\lor}[1]}, close]
              ]
            ]
          ]
        ]
      ]
    \end{oltableau}
  \item دوناں $\{\sFmla{\False}{!A \lor !B}, \sFmla{\True}{!A}\}$ تے
    $\{\sFmla{\False}{!A \lor !B}, \sFmla{\True}{!B}\}$ لئی بند
    !!{tableau}s موجود نیں:
    \begin{oltableau}{}
      [\sFmla{\False}{\formula{A} \lor \formula{B}}, just=\TAss
        [\sFmla{\True}{\formula{A}}, just=\TAss
          [\sFmla{\False}{\formula{A}},just={\TRule{\False}{\lor}[1]}
            [\sFmla{\False}{\formula{B}},just={\TRule{\False}{\lor}[1]}, close
            ]
          ]
        ]
      ]
    \end{oltableau}
    \begin{oltableau}{}
      [\sFmla{\False}{\formula{A} \lor \formula{B}}, just=\TAss
        [\sFmla{\True}{\formula{B}}, just=\TAss
          [\sFmla{\False}{\formula{A}},just={\TRule{\False}{\lor}[1]}
            [\sFmla{\False}{\formula{B}},just={\TRule{\False}{\lor}[1]}, close
            ]
          ]
        ]
      ]
    \end{oltableau}
  \end{enumerate}
\end{proof}

\begin{prop}\ollabel{prop:provability-lif}
  \begin{enumerate}
  \item \ollabel{prop:provability-lif-left} $!A, !A \lif !B \Proves !B$۔
  \item \ollabel{prop:provability-lif-right}
    دونیں $\lnot !A \Proves !A \lif !B$ تے $!B \Proves !A \lif !B$۔
  \end{enumerate}
\end{prop}

\begin{proof}
  \begin{enumerate}
    \item $\{\sFmla{\False}{!B}, \sFmla{\True}{!A \lif !B},
      \sFmla{\True}{!A}\}$ لئی بند !!{tableau} ایہ اے:
      \begin{oltableau}
        [\sFmla{\False}{\formula{B}}, just=\TAss
          [\sFmla{\True}{\formula{A} \lif \formula{B}}, just=\TAss
            [\sFmla{\True}{\formula{A}}, just=\TAss
              [\sFmla{\False}{\formula{A}}, just = {\TRule{\True}{\lif}[2]}, close]
              [\sFmla{\True}{\formula{B}}, just = {\TRule{\True}{\lif}[2]}, close]
            ]
          ]
        ]
      \end{oltableau}
    \item دوناں $\{\sFmla{\False}{!A \lif !B},
      \sFmla{\True}{\lnot !A}\}$ تے $\{\sFmla{\False}{!A \lif !B},
      \sFmla{\True}{!B}\}$ لئی بند !!{tableau}s موجود نیں:
      \begin{oltableau}
        [\sFmla{\False}{\formula{A} \lif \formula{B}}, just = \TAss
          [\sFmla{\True}{\lnot \formula{A}}, just = \TAss
            [\sFmla{\True}{\formula{A}}, just = {\TRule{\False}{\lif}[1]}
              [\sFmla{\False}{\formula{B}}, just = {\TRule{\False}{\lif}[1]}
                [\sFmla{\False}{\formula{A}}, just = {\TRule{\True}{\lnot}[2]}, close]
              ]
            ]
          ]
        ]
      \end{oltableau}
      \begin{oltableau}
        [\sFmla{\False}{\formula{A} \lif \formula{B}}, just = \TAss
          [\sFmla{\True}{\formula{B}}, just = \TAss
            [\sFmla{\True}{\formula{A}}, just = {\TRule{\False}{\lif}[1]}
              [\sFmla{\False}{\formula{B}}, just = {\TRule{\False}{\lif}[1]},close
              ]
            ]
          ]
        ]
      \end{oltableau}
  \end{enumerate}
\end{proof}
\endgroup
% /OLP-0107

% OLP-0108 content/first-order-logic/tableaux/provability-quantifiers.tex
\begingroup
\def\olfilename{provability-quantifiers}
\expandafter\def\csname ol@sectioncs\endcsname{section}

\olfileid{fol}{tab}{qpr}
\olsection{\usetoken{S}{derivability} تے مسوّر \pnbLatin{(Derivability and the Quantifiers)}}

\begin{explain}
  تمامیت دے قضیے دا تقاضا ایہ وی اے کہ تابلو دے قاعدے~$\Proves$
  بارے اوہ حقائق دیندے ہون جیہڑے ایس حصے وچ قائم کیتے گئے نیں۔
\end{explain}

\begin{thm}
\ollabel{thm:strong-generalization} جے $c$ اک ایہو جہا مستقل (constant) اے جیہڑا
$\Gamma$ یا~$!A(x)$ وچ نہیں آؤندا تے $\Gamma \Proves !A(c)$، تاں $\Gamma
\Proves \lforall[x][!A(x)]$۔
\end{thm}

\begin{proof}
فرض کرو $\Gamma \Proves !A(c)$، یعنی $!B_1$، \dots، $!B_n
\in \Gamma$ موجود نیں تے ہِٹھلے سیٹ لئی اک بند !!{tableau} اے:
\begin{align*}
\{ \sFmla{\False}{!A(c)}, &
\sFmla{\True}{!B_1}, \dots, \sFmla{\True}{!B_n} \}.
\intertext{سانوں وکھانا اے کہ ہِٹھلے سیٹ لئی وی اک بند !!{tableau} موجود اے:}
\{ \sFmla{\False}{\lforall[x][!A(x)]}, &
\sFmla{\True}{!B_1}, \dots, \sFmla{\True}{!B_n} \}.
\end{align*}
ایہ بند !!{tableau} لو تے پہلا مفروضہ
$\sFmla{\False}{\lforall[x][!A(x)]}$ نال بدلو؛ فیر سارے
مفروضیاں توں بعد $\sFmla{\False}{!A(c)}$ شامل کرو۔
\begin{center}
\begin{tableau}{not line numbering}
  [\sFmla{\False}{\formula{A}(c)}
    [\sFmla{\True}{\formula{B}_1}
      [\vdots
      [\sFmla{\True}{\formula{B}_n}, 
        [][][]]]]]
\end{tableau}\qquad
\begin{tableau}{not line numbering}
  [\sFmla{\False}{\lforall[x][\formula{A}(x)]}
    [\sFmla{\True}{\formula{B}_1}
      [\vdots
        [\sFmla{\True}{\formula{B}_n},
          [\sFmla{\False}{\formula{A}(c)}
        [][][]]]]]]
\end{tableau}
\end{center}
ایہ تابلو ہن وی بند اے، کیوں جے جیہڑے سارے
!!{sentence}s پہلے مفروضیاں دے طور تے دستیاب سن، اوہ ہن وی
!!{tableau} دے شروع وچ موجود نیں۔ شامل کیتی لکیر
$\TRule{\False}{\lforall}$ دی درست ورتوں نال ملی اے، کیوں جے
!!{constant}~$c$، $!B_1$، \dots، $!B_n$ یا
$\lforall[x][!A(x)]$ وچ نہیں آؤندا؛ یعنی نویں !!{tableau}
وچ شامل کیتی لکیر توں اُتّے کہیں وی نہیں آؤندا۔
\end{proof}

\begin{prop}
\ollabel{prop:provability-quantifiers}
\begin{tagenumerate}{prvEx,prvAll}
\tagitem{prvEx}{$!A(t) \Proves \lexists[x][!A(x)]$.}{}
\tagitem{prvAll}{$\lforall[x][!A(x)] \Proves !A(t)$.}{}
\end{tagenumerate}
\end{prop}

\begin{proof}
\begin{tagenumerate}{prvEx,prvAll}
\tagitem{prvEx}{$\sFmla{\False}{\lexists[x][!A(x)]}, \sFmla{\True}{!A(t)}$
  لئی اک بند !!{tableau} ایہ اے:
  \begin{oltableau}
    [\sFmla{\False}{\lexists[x][\formula{A}(x)]}, just = \TAss
      [\sFmla{\True}{\formula{A}(t)}, just = \TAss
        [\sFmla{\False}{\formula{A}(t)}, just = {\TRule{\False}{\lexists}[1]}, close]]]
    \end{oltableau}}{}
% PNBTAB-012: source line 79's trailing comma is punctuation, not part of the signed premise.
\tagitem{prvAll}{$\sFmla{\False}{\formula{A}(t)}, \sFmla{\True}{\lforall[x][!A(x)]}$
  لئی اک بند !!{tableau} ایہ اے:
  \begin{oltableau}
    [\sFmla{\False}{\formula{A}(t)}, just = \TAss
      [\sFmla{\True}{\lforall[x][\formula{A}(x)]}, just = \TAss
        [\sFmla{\True}{\formula{A}(t)}, just = {\TRule{\True}{\lforall}[2]}, close]]]
    \end{oltableau}}{}
\end{tagenumerate}
\end{proof}
\endgroup
% /OLP-0108

% OLP-0109 content/first-order-logic/tableaux/soundness.tex
\begingroup
\def\olfilename{soundness}
\expandafter\def\csname ol@sectioncs\endcsname{section}

\iftag{FOL}
      {\olfileid{fol}{tab}{sou}}
      {\olfileid{pl}{tab}{sou}}
      
\olsection{صحت مندی \pnbLatin{(Soundness)}}

\begin{explain}
!!^a{derivation} نظام، جیویں تابلو، \emph{صحت مند (sound)} اے جے اوہ
اوہ گلاں !!{derive} نہیں کر سکدا جیہڑیاں اصل وچ قائم نہیں۔ ایس لئی
صحت مندی !!{derivation} نظاماں دی ضمانت شدہ حفاظتی خاصیت دی اک قسم اے۔
جیہڑی ثبوتی-نظریاتی خاصیت دی گل ہو رہی ہووے، اوہدے مطابق اسیں، مثال
لئی، ایہ جاننا چاہندے آں کہ
\begin{enumerate}
\item ہر !!{derivable}~$!A$، \iftag{FOL}{معتبر}{ہمیشہ صادق فارمولا} اے؛
\item جے کوئی !!a{sentence} کجھ ہور جملیاں توں !!{derivable} اے، تاں اوہ
  اوہناں دا منطقی نتیجہ وی اے؛
\item جے !!{sentence}s دا اک سیٹ متضاد اے، تاں اوہ پورا ہون جوگ نہیں۔
\end{enumerate}
ایہ !!a{derivation} نظام دیاں اہم خاصیتاں نیں۔ جے ایہناں وچوں کوئی
وی خاصیت قائم نہ ہووے، تاں !!{derivation} نظام ناقص اے---اوہ لوڑ
توں ودھ !!{derive} کرے گا۔ ایس کرکے !!a{derivation} نظام دی صحت
مندی قائم کرنا نہایت ضروری اے۔

چونکہ ایہ ساریاں ثبوتی-نظریاتی خاصیتاں کسے نہ کسے قسم دے بند
!!{tableau}s راہیں تعریف کیتیاں گئیاں نیں، ایس لئی اُپرلے (1)--(3)
ثابت کرن لئی بند !!{tableau}s دیاں معنوی خاصیتاں بارے کجھ ثابت کرنا
پیندا اے۔ اسیں پہلے تعریف کراں گے کہ کسے ساخت وچ !!a{signed formula} دے پورا ہون دا کی مطلب اے، تے فیر وکھاواں گے کہ جے کوئی
!!{tableau} بند ہووے تاں کوئی وی ساخت اوہدے سارے مفروضیاں نوں پورا
نہیں کردی۔ (1)--(3) فیر ایس نتیجے دے ذیلی قضیاں دے طور تے نکلن گے۔
\end{explain}

\begin{defn}
\iftag{FOL}{!!^a{structure}}{!!^a{valuation}}~$\iftag{FOL}{\Struct{M}}{\pAssign{v}}$
!!a{signed formula} $\sFmla{\True}{!A}$ نوں \emph{پورا کردی اے (satisfies)}
اُدوں تے صرف اُدوں جدوں
$\iftag{FOL}{\Sat{M}}{\pSat{v}}{!A}$؛ تے اوہ
$\sFmla{\False}{!A}$ نوں اُدوں تے صرف اُدوں پورا کردی اے جدوں
$\iftag{FOL}{\Sat/{M}}{\pSat/{v}}{!A}$۔ $\iftag{FOL}{\Struct{M}}{\pAssign{v}}$
!!{signed formula}s دے سیٹ~$\Gamma$ نوں اُدوں تے صرف اُدوں پورا
کردی اے جدوں ہر $\sFmla{S}{!A} \in \Gamma$ نوں پورا کردی اے۔
$\Gamma$، \emph{پورا ہون جوگ (satisfiable)} اے جے کوئی
\iftag{FOL}{!!a{structure}}{!!a{valuation}} اوہنوں پورا کردی ہووے،
ورنہ اوہ \emph{پورا ہون جوگ نہیں (unsatisfiable)}۔
\end{defn}

\begin{thm}[صحت مندی (Soundness)]
  \ollabel{thm:tableau-soundness}
  جے $\Gamma$ لئی اک بند !!{tableau} اے، تاں $\Gamma$ پورا ہون جوگ نہیں۔
\end{thm}

\begin{proof}
اک !!a{tableau} دی شاخ نوں پورا ہون جوگ آکھاں گے اُدوں تے صرف اُدوں
جدوں اوس اُتے موجود !!{signed formula}s دا سیٹ پورا ہون جوگ ہووے؛
تے !!a{tableau} نوں پورا ہون جوگ آکھاں گے جے اوہدے وچ گھٹّو گھٹ
اک پورا ہون جوگ شاخ ہووے۔

اسیں ہِٹھلی گل وکھاندے آں: پورا ہون جوگ !!{tableau} نوں استنباط دے
کسے وی اک قاعدے نال ودھاؤن دا نتیجہ ہمیشہ اک پورا ہون جوگ
!!{tableau} ہوندا اے۔ ایس نال قضیہ ثابت ہو جاوے گا: ہر بند
!!{tableau}، صرف~$\Gamma$ دے مفروضیاں اُتے مشتمل ابتدائی
!!{tableau} اُتے استنباط دے قاعدے لا کے بندا اے۔ سو جے $\Gamma$
پورا ہون جوگ ہوندا، تاں اوہدے لئی ہر !!{tableau} وی پورا ہون جوگ
ہوندا۔ پر بند !!{tableau} صاف طور تے پورا ہون جوگ نہیں: ہر شاخ
وچ $\sFmla{\True}{!A}$ تے $\sFmla{\False}{!A}$ دونیں ہوندے نیں،
تے کسے ساخت لئی ایہ ممکن نہیں کہ اوہ اکّو ویلے~$!A$ نوں پورا
وی کرے تے نہ وی کرے۔

فرض کرو ساڈے کول اک پورا ہون جوگ !!{tableau} اے، یعنی ایہو جہا
!!a{tableau} جیہدے وچ گھٹّو گھٹ اک پورا ہون جوگ شاخ اے۔ استنباط دا
قاعدہ لاون نال یا تاں کسے شاخ وچ !!{signed formula}s شامل ہوندے
نیں، یا اوہ شاخ دو حصیاں وچ ونڈدی اے۔ جے !!{tableau} دی کوئی
پورا ہون جوگ شاخ ایس قاعدے دی ورتوں نال ودھائی نہیں گئی، تاں اوہ
ودھے ہوئے !!{tableau} وچ وی پورا ہون جوگ رہندی اے؛ ایس لئی ودھیا
ہویا تابلو پورا ہون جوگ اے۔ سو سانوں صرف اوہ صورت ویکھنی اے جتھے
قاعدہ اک پورا ہون جوگ شاخ اُتے لاگو ہوندا اے۔

منّ لو $\Gamma$ اوس شاخ دے !!{signed formula}s دا سیٹ اے، تے
$\sFmla{S}{!A} \in \Gamma$ اوہ !!{signed formula} اے جیہدے اُتے
قاعدہ لاگو ہوندا اے۔ جے قاعدہ شاخ نوں نہیں ونڈدا، تاں سانوں
وکھانا اے کہ ودھی ہوئی شاخ، یعنی $\Gamma$ نال قاعدے دے نتیجے
جوڑ کے، ہن وی پورا ہون جوگ اے۔ جے قاعدہ شاخ نوں ونڈدا اے، تاں
سانوں وکھانا اے کہ بنیاں دو شاخاں وچوں گھٹّو گھٹ اک پورا ہون جوگ اے۔

پہلے اسیں اوہ ممکن استنباط ویکھدے آں جیہڑے شاخ نوں نہیں ونڈدے۔
\begin{enumerate}
\item $\TRule{\True}{\lnot}$ نوں
  $\sFmla{\True}{\lnot !B} \in \Gamma$ اُتے لا کے شاخ ودھائی جاندی اے۔ فیر ودھی ہوئی
  شاخ وچ $\Gamma \cup \{\sFmla{\False}{!B}\}$ دے !!{signed formula}s
  ہوندے نیں۔ فرض کرو
  $\iftag{FOL}{\Sat{M}}{\pSat{v}}{\Gamma}$۔ خاص طور تے
  $\iftag{FOL}{\Sat{M}}{\pSat{v}}{\lnot !B}$۔ ایس لئی
  $\iftag{FOL}{\Sat/{M}}{\pSat/{v}}{!B}$، یعنی
  $\iftag{FOL}{\Struct{M}}{\pAssign{v}}$،
  $\sFmla{\False}{!B}$ نوں پورا کردی اے۔
\item $\TRule{\False}{\lnot}$ نوں
  $\sFmla{\False}{\lnot !B} \in \Gamma$ اُتے لا کے شاخ ودھائی جاندی اے: مشق۔
\item $\TRule{\True}{\land}$ نوں
  $\sFmla{\True}{!B \land !C} \in \Gamma$ اُتے لا کے شاخ ودھائی جاندی اے، جیہدے نال
  شاخ اُتے دو نویں !!{signed formula}s آ جاندے نیں:
  $\sFmla{\True}{!B}$ تے $\sFmla{\True}{!C}$۔ فرض کرو
  $\iftag{FOL}{\Sat{M}}{\pSat{v}}{\Gamma}$، خاص طور تے
  $\iftag{FOL}{\Sat{M}}{\pSat{v}}{!B \land !C}$۔ فیر
  $\iftag{FOL}{\Sat{M}}{\pSat{v}}{!B}$ تے
  $\iftag{FOL}{\Sat{M}}{\pSat{v}}{!C}$۔ ایس دا مطلب اے کہ
  $\iftag{FOL}{\Struct{M}}{\pAssign{v}}$ دوناں
  $\sFmla{\True}{!B}$ تے $\sFmla{\True}{!C}$ نوں پورا کردی اے۔
\item $\TRule{\False}{\lor}$ نوں
  $\sFmla{\False}{!B \lor !C} \in \Gamma$ اُتے لا کے شاخ ودھائی جاندی اے: مشق۔
\item $\TRule{\False}{\lif}$ نوں
  $\sFmla{\False}{!B \lif !C} \in \Gamma$ اُتے لا کے شاخ ودھائی جاندی اے: ایس نال شاخ
  اُتے دو نویں !!{signed formula}s، $\sFmla{\True}{!B}$ تے
  $\sFmla{\False}{!C}$، آ جاندے نیں۔ فرض کرو
  $\iftag{FOL}{\Sat{M}}{\pSat{v}}{\Gamma}$، خاص طور تے
  $\iftag{FOL}{\Sat/{M}}{\pSat/{v}}{!B \lif !C}$۔ فیر
  $\iftag{FOL}{\Sat{M}}{\pSat{v}}{!B}$ تے
  $\iftag{FOL}{\Sat/{M}}{\pSat/{v}}{!C}$۔ ایس دا مطلب اے کہ
  $\iftag{FOL}{\Struct{M}}{\pAssign{v}}$ دوناں
  $\sFmla{\True}{!B}$ تے $\sFmla{\False}{!C}$ نوں پورا کردی اے۔

% PNBTAB-013: five B(x) occurrences in the FOL universal cases are corrected to A(x); see the source-correction analysis.
\iftag{FOL}{%
\item $\TRule{\True}{\lforall}$ نوں
  $\sFmla{\True}{\lforall[x][!A(x)]} \in \Gamma$ اُتے لا کے شاخ ودھائی جاندی اے: ایس نال
  شاخ اُتے نواں !!{signed formula}~$\sFmla{\True}{!A(t)}$ آ جاندا
  اے۔ فرض کرو $\Sat{M}{\Gamma}$، خاص طور تے
  $\Sat{M}{\lforall[x][!A(x)]}$۔
  \olref[syn][sem]{prop:quant-terms} مطابق $\Sat{M}{!A(t)}$۔
  ایس لئی $\Struct{M}$، $\sFmla{\True}{!A(t)}$ نوں پورا کردی اے۔

\item $\TRule{\False}{\lforall}$ نوں
  $\sFmla{\False}{\lforall[x][!A(x)]} \in \Gamma$ اُتے لا کے شاخ ودھائی جاندی اے: ایس نال
  نواں !!{signed formula}~$\sFmla{\False}{!A(a)}$ آ جاندا اے،
  جتھے $a$ اک !!a{constant} اے جیہڑا~$\Gamma$ وچ نہیں آؤندا۔
  چونکہ $\Gamma$ پورا ہون جوگ اے، کوئی ساخت~$\Struct{M}$ موجود اے
  جیہدے لئی $\Sat{M}{\Gamma}$، خاص طور تے
  $\Sat/{M}{\lforall[x][!A(x)]}$۔ سانوں وکھانا اے کہ
  $\Gamma \cup \{\sFmla{\False}{!A(a)}\}$ پورا ہون جوگ اے۔
  ایس لئی اسیں مناسب ساخت~$\Struct{M'}$ ایویں تعریف کردے آں۔

  \olref[syn][ass]{prop:sat-quant} مطابق
  $\Sat/{M}{\lforall[x][!A(x)]}$ اُدوں تے صرف اُدوں جدوں کسے
  متغیر تعیّن~$s$ لئی $\Sat/{M}{!A(x)}[s]$۔ ہن $\Struct{M'}$
  نوں $\Struct{M}$ ورگی ساخت لو، سوائے ایس دے کہ
  $\Assign{a}{M'} = s(x)$۔
  \olref[syn][ext]{cor:extensionality-sent} مطابق ہر
  $\sFmla{\True}{!C} \in \Gamma$ لئی $\Sat{M'}{!C}$، تے ہر
  $\sFmla{\False}{!C} \in \Gamma$ لئی $\Sat/{M'}{!C}$، کیوں جے
  $a$،~$\Gamma$ وچ نہیں آؤندا۔

  \olref[syn][ext]{prop:extensionality} مطابق
  $\Sat/{M'}{!A(x)}[s]$۔
  \olref[syn][ext]{prop:ext-formulas} مطابق
  $\Sat/{M'}{!A(a)}[s]$۔ چونکہ $!A(a)$ اک !!a{sentence} اے،
  \olref[syn][ass]{prop:sentence-sat-true} مطابق
  $\Sat/{M'}{!A(a)}$؛ یعنی $\Struct{M'}$،
  $\sFmla{\False}{!A(a)}$ نوں پورا کردی اے۔
\item $\TRule{\True}{\lexists}$ نوں
  $\sFmla{\True}{\lexists[x][!B(x)]} \in \Gamma$ اُتے لا کے شاخ ودھائی جاندی اے: مشق۔
\item $\TRule{\False}{\lexists}$ نوں
  $\sFmla{\False}{\lexists[x][!B(x)]} \in \Gamma$ اُتے لا کے شاخ ودھائی جاندی اے: مشق۔}{}
\end{enumerate}
ہن شاخ نوں ونڈن والے ممکن استنباط ویکھدے آں۔
\begin{enumerate}
\item $\TRule{\False}{\land}$ نوں
  $\sFmla{\False}{!B \land !C} \in \Gamma$ اُتے لا کے شاخ ودھائی جاندی اے۔ ایہدے نال
  دو شاخاں بندیاں نیں: کھبّی شاخ $\sFmla{\False}{!B}$ راہیں اگے
  ودھدی اے، تے سجّی $\sFmla{\False}{!C}$ راہیں۔ فرض کرو
  $\iftag{FOL}{\Sat{M}}{\pSat{v}}{\Gamma}$، خاص طور تے
  $\iftag{FOL}{\Sat/{M}}{\pSat/{v}}{!B \land !C}$۔ فیر
  $\iftag{FOL}{\Sat/{M}}{\pSat/{v}}{!B}$ یا
  $\iftag{FOL}{\Sat/{M}}{\pSat/{v}}{!C}$۔ پہلی صورت وچ
  $\iftag{FOL}{\Struct{M}}{\pAssign{v}}$،
  $\sFmla{\False}{!B}$ نوں پورا کردی اے، یعنی
  $\iftag{FOL}{\Struct{M}}{\pAssign{v}}$ کھبّی شاخ دے
  سارے فارمولیاں نوں پورا کردی اے۔ دوجی صورت وچ
  $\iftag{FOL}{\Struct{M}}{\pAssign{v}}$،
  $\sFmla{\False}{!C}$ نوں پورا کردی اے، یعنی
  $\iftag{FOL}{\Struct{M}}{\pAssign{v}}$ سجّی شاخ دے
  سارے فارمولیاں نوں پورا کردی اے۔
\item $\TRule{\True}{\lor}$ نوں
  $\sFmla{\True}{!B \lor !C} \in \Gamma$ اُتے لا کے شاخ ودھائی جاندی اے: مشق۔
\item $\TRule{\True}{\lif}$ نوں
  $\sFmla{\True}{!B \lif !C} \in \Gamma$ اُتے لا کے شاخ ودھائی جاندی اے: مشق۔
\item شاخ \Cut{} لا کے ودھائی جاندی اے: ایس نال دو شاخاں بندیاں
  نیں، اک وچ $\sFmla{\True}{!B}$ تے دوجی وچ
  $\sFmla{\False}{!B}$۔ چونکہ
  $\iftag{FOL}{\Sat{M}}{\pSat{v}}{\Gamma}$ تے یا تاں
  $\iftag{FOL}{\Sat{M}}{\pSat{v}}{!B}$ یا
  $\iftag{FOL}{\Sat/{M}}{\pSat/{v}}{!B}$، ایس لئی
  $\iftag{FOL}{\Struct{M}}{\pAssign{v}}$ کھبّی یا سجّی شاخ
  وچوں اک نوں پورا کردی اے۔
\end{enumerate}
\end{proof}

\tagprob{FOL}
\begin{prob}
\olref[fol][tab][sou]{thm:tableau-soundness} دا ثبوت مکمل کرو۔
\end{prob}
\tagendprob

\tagprob{notFOL}
\begin{prob}
\olref[pl][tab][sou]{thm:tableau-soundness} دا ثبوت مکمل کرو۔
\end{prob}
\tagendprob

\begin{cor}
\ollabel{cor:weak-soundness}
جے $\Proves !A$، تاں $!A$ \iftag{FOL}{معتبر}{ہمیشہ صادق فارمولا} اے۔
\end{cor}

\begin{cor}
\ollabel{cor:entailment-soundness}
جے $\Gamma \Proves !A$، تاں $\Gamma \Entails !A$۔
\end{cor}

\begin{proof}
  جے $\Gamma \Proves !A$، تاں کجھ $!B_1$، \dots، $!B_n \in
  \Gamma$ ایہو جہے نیں کہ
  $\{\sFmla{\False}{!A}, \sFmla{\True}{!B_1}, \dots,
  \sFmla{\True}{!B_n}\}$ لئی اک بند !!{tableau} اے۔
  \olref{thm:tableau-soundness} مطابق ہر
  \iftag{FOL}{!!{structure}}{!!{valuation}}~$\iftag{FOL}{\Struct{M}}{\pAssign{v}}$
  یا تاں کسے $!B_i$ نوں کاذب بناندی اے یا $!A$ نوں صادق۔
  ایس لئی جے $\iftag{FOL}{\Sat{M}}{\pSat{v}}{\Gamma}$، تاں
  $\iftag{FOL}{\Sat{M}}{\pSat{v}}{!A}$ وی۔
\end{proof}

\begin{cor}
\ollabel{cor:consistency-soundness}
جے $\Gamma$ پورا ہون جوگ اے، تاں ایہ بے تضاد اے۔
\end{cor}

\begin{proof}
اسیں متقابلِ عکس ثابت کردے آں۔ فرض کرو $\Gamma$ بے تضاد نہیں۔
فیر کجھ $!B_1$، \dots، $!B_n \in \Gamma$ نیں تے
$\{\sFmla{\True}{!B_1}, \dots, \sFmla{\True}{!B_n}\}$ لئی
اک بند !!{tableau} اے۔ \olref{thm:tableau-soundness} مطابق کوئی
$\iftag{FOL}{\Struct{M}}{\pAssign{v}}$ ایہو جہی نہیں کہ
$\iftag{FOL}{\Sat{M}}{\pSat{v}}{!B_i}$ ہر
$i=1$، \dots،~$n$ لئی قائم ہووے۔ پر ایس دا
مطلب اے کہ $\Gamma$ پورا ہون جوگ نہیں۔
\end{proof}
\endgroup
% /OLP-0109

% OLP-0110 content/first-order-logic/tableaux/identity.tex
\begingroup
\def\olfilename{identity}
\expandafter\def\csname ol@sectioncs\endcsname{section}

\olfileid{fol}{tab}{ide}

\olsection{\usetoken{P}{tableau} وچ \usetoken{S}{identity} دی شمولیت}

!!^{tableau}s وچ !!{identity} لئی وادھو استنباطی قاعدے چاہیدے نیں۔
$\eq$ دے قاعدے ایہ نیں ($t$، $t_1$ تے $t_2$ بند ٹرم نیں):

\begin{defish}
\AxiomC{}
\RightLabel{$\eq$}
\UnaryInfC{\sFmla{\True}{\eq[t][t]}}
\DisplayProof
\hfill
\AxiomC{\sFmla{\True}{\eq[t_1][t_2]}}
\noLine
\UnaryInfC{\sFmla{\True}{!A(t_1)}}
\RightLabel{$\TRule{\True}{\eq}$}
\UnaryInfC{\sFmla{\True}{!A(t_2)}}
\DisplayProof
\hfill
\AxiomC{\sFmla{\True}{\eq[t_1][t_2]}}
\noLine
\UnaryInfC{\sFmla{\False}{!A(t_1)}}
\RightLabel{$\TRule{\False}{\eq}$}
\UnaryInfC{\sFmla{\False}{!A(t_2)}}
\DisplayProof
\end{defish}
دھیان دیو کہ باقی سارے قاعدیاں دے برعکس، $\TRule{\True}{\eq}$ تے
$\TRule{\False}{\eq}$ لئی لازم اے کہ شاخ اُتے پہلے ای \emph{دو}
علامت لائے !!{formula}s موجود ہون، یعنی
$\sFmla{\True}{\eq[t_1][t_2]}$ تے $\sFmla{S}{!A(t_1)}$ دونیں۔

\begin{ex}
جے $s$ تے $t$ بند ٹرم نیں، تاں $\eq[s][t], !A(s)
\Proves !A(t)$:
\begin{oltableau}
  [\sFmla{\False}{\formula{A}(t)}, just = \TAss
    [\sFmla{\True}{\eq[s][t]}, just = \TAss
      [\sFmla{\True}{\formula{A}(s)}, just = \TAss
        [\sFmla{\True}{\formula{A}(t)}, just={\TRule{\True}{\eq}[2, 3]}, close]
      ]
    ]
  ]
\end{oltableau}
تساں ایس نوں اکو ہون والیاں چیزاں دی اک دوجے دی تھاں ورتوں دا اصول
\pnbLatin{(substitutability of identicals)}، یا لائبنز دا قانون \pnbLatin{(Leibniz' Law)}،
دے ناں نال جان دے ہوؤ گے۔

% PNBTAB-014: two explanatory identity-rule premise formulas are normalized; the three proof trees are unchanged.
!!^{tableau}s ایہ ثابت کردے نیں کہ $\eq$ متناظر اے، یعنی
$\eq[s_1][s_2]
\Proves \eq[s_2][s_1]$:
\begin{oltableau}
  [\sFmla{\False}{\eq[s_2][s_1]}, just = \TAss
    [\sFmla{\True}{\eq[s_1][s_2]}, just = \TAss
      [\sFmla{\True}{\eq[s_1][s_1]}, just = {$\eq$}
        [\sFmla{\True}{\eq[s_2][s_1]}, just = {\TRule{\True}{\eq}[2, 3]}, close]
      ]
    ]
  ]
\end{oltableau}
ایتھے سطر~$2$ پہلا لازمی !!{formula}~$\sFmla{\True}{\eq[s_1][s_2]}$
اے جیہڑا $\TRule{\True}{\eq}$ لئی چاہیدا اے۔ سطر~$3$ دوجا
لازمی فارمولا اے، جیہدا روپ $\sFmla{\True}{!A(s_1)}$ اے۔
$!A(x)$ نوں $\eq[x][s_1]$ سمجھو؛ فیر $!A(s_1)$،
$\eq[s_1][s_1]$ اے تے $!A(s_2)$، $\eq[s_2][s_1]$ اے۔

ایہ وی ثابت کردے نیں کہ $\eq$ متعدی اے، یعنی
$\eq[s_1][s_2],
\eq[s_2][s_3] \Proves \eq[s_1][s_3]$:
\begin{oltableau}
  [\sFmla{\False}{\eq[s_1][s_3]}, just = \TAss
    [\sFmla{\True}{\eq[s_1][s_2]}, just = \TAss
      [\sFmla{\True}{\eq[s_2][s_3]}, just = \TAss
        [\sFmla{\True}{\eq[s_1][s_3]}, just = {\TRule{\True}{\eq}[3, 2]}, close]
      ]
    ]
  ]
\end{oltableau}
ایس !!{tableau} وچ $\TRule{\True}{\eq}$ دا پہلا لازمی
!!{formula} سطر~$3$ دا $\sFmla{\True}{\eq[s_2][s_3]}$ اے
($s_2$،~$t_1$ دا کردار ادا کردا اے، تے $s_3$،~$t_2$ دا)۔
دوجا لازمی فارمولا، جیہدا روپ~$\sFmla{\True}{!A(s_2)}$ اے،
سطر~$2$ وچ اے۔ ایتھے $!A(x)$ نوں $\eq[s_1][x]$ سمجھو؛
ایس طرح $!A(s_2)$، $\eq[s_1][s_2]$ بندا اے (یعنی سطر~$2$)
تے $!A(s_3)$ نتیجے دا !!{formula}~$\eq[s_1][s_3]$ بندا اے۔
\end{ex}

\begin{prob}
ہِٹھ دتیاں لئی بند !!{tableau}s بناؤ:
\begin{enumerate}
\item $\sFmla{\False}{\lforall[x][\lforall[y][((x = y \land !A(x))
      \lif !A(y))]]}$
\item $\sFmla{\False}{\lexists[x][(!A(x) \land
    \lforall[y][(!A(y) \lif y = x)])]}$,\\
  $\sFmla{\True}{\lexists[x][!A(x)] \land
    \lforall[y][\lforall[z][((!A(y) \land !A(z)) \lif y = z)]]}$
\end{enumerate}
\end{prob}
\endgroup
% /OLP-0110
% OLP-0111 content/first-order-logic/tableaux/soundness-identity.tex
\begingroup
\def\olfilename{soundness-identity}
\expandafter\def\csname ol@sectioncs\endcsname{section}

\olfileid{fol}{tab}{sid}

\olsection{\usetoken{S}{identity} نال صحت مندی}

\begin{prop}
!!^{tableau}s جیہناں وچ اکو ہون دے قاعدے وی شامل نیں، صحت مند نیں:
کوئی بند !!{tableau} پورا ہون جوگ نہیں۔
\end{prop}

\begin{proof}
ساڈے لئی، پہلے وانگ، صرف ایہ وکھانا کافی اے کہ جے !!a{tableau} دی
کوئی شاخ پورا ہون جوگ اے، تاں اوس شاخ اُتے~$\eq$ دے قاعدیاں وچوں
کوئی اک لا کے بنن والی شاخ وی پورا ہون جوگ اے۔ $\Gamma$ نوں شاخ اُتے
موجود علامت لائے !!{formula}s دا سیٹ منّ لو، تے $\Struct{M}$
!!a{structure} ہووے جیہڑی~$\Gamma$ نوں پورا کردی اے۔

فرض کرو شاخ نوں~$\eq$ نال ودھایا جاندا اے، یعنی علامت لایا
!!{formula}~$\sFmla{\True}{\eq[t][t]}$ شامل کیتا جاندا اے۔
ظاہر اے کہ $\Sat{M}{\eq[t][t]}$، ایس لئی $\Struct{M}$،
$\Gamma \cup \{\sFmla{\True}{\eq[t][t]}\}$ نوں وی پورا کردی اے۔

% PNBTAB-015: the explicitly True equality-rule conclusion has its True sign; all other formulas are unchanged.
جے شاخ نوں $\TRule{\True}{\eq}$ نال ودھایا جائے، تاں اسیں علامت لایا
!!{formula}~$\sFmla{\True}{!A(t_2)}$ شامل کردے آں، پر $\Gamma$ وچ
دونیں~$\sFmla{\True}{\eq[t_1][t_2]}$ تے $\sFmla{\True}{!A(t_1)}$
موجود نیں۔ ایس لئی $\Sat{M}{\eq[t_1][t_2]}$ تے
$\Sat{M}{!A(t_1)}$۔ اک متغیراں دی قدر بندی $s$ لو جیہدے لئی
$s(x) = \Value{t_1}{M}$۔ \olref[syn][ass]{prop:sentence-sat-true}
مطابق $\Sat{M}{!A(t_1)}[s]$۔ چونکہ $\varAssign{s}{s}{x}$،
\olref[syn][ext]{prop:ext-formulas} مطابق $\Sat{M}{!A(x)}[s]$۔
چونکہ $\Sat{M}{\eq[t_1][t_2]}$، اسیں جان دے آں کہ
$\Value{t_1}{M} = \Value{t_2}{M}$، ایس لئی
$s(x) = \Value{t_2}{M}$۔ \olref[syn][ext]{prop:ext-formulas} نوں
دوبارہ لا کے $\Sat{M}{!A(t_2)}[s]$ وی ملدا اے۔
\olref[syn][ass]{prop:sentence-sat-true} مطابق
$\Sat{M}{!A(t_2)}$۔ $\TRule{\False}{\eq}$ دی صورت وی ایس طرح
حل ہوندی اے۔
\end{proof}
\endgroup
% /OLP-0111

\OLEndChapterHook
\endgroup
% /OLP-0098

% OLP-0112 content/first-order-logic/axiomatic-deduction/axiomatic-deduction.tex
\begingroup
\def\olfilename{axiomatic-deduction}
\expandafter\def\csname ol@sectioncs\endcsname{section}

\iftag{FOL}
      {\olchapter{fol}{axd}{مسلّمی \usetoken{P}{derivation}}}
      {\olchapter{pl}{axd}{مسلّمی \usetoken{P}{derivation}}}

\begin{editorial}
  ہن تک ایہ پکا کرن دی کوشش نہیں کیتی گئی کہ ایس باب دا مواد اوہناں
  وکھ وکھ ٹیگاں دے مطابق اے جیہڑے دسدے نیں کہ کیہڑے منطقی رابطے تے
  مسوّر ابتدائی نیں تے کیہڑے تعریف شدہ: سارے رابطے تے مسوّر ابتدائی
  منّے گئے نیں، صرف $\liff$ نوں تعریف شدہ منّیا گیا اے۔ جے FOL
  ٹیگ سچا ہووے، تاں اسیں مسوّراں والا روپ تیار کردے آں، ورنہ
  مسوّراں توں بنا روپ۔
\end{editorial}

% OLP-0113 content/first-order-logic/axiomatic-deduction/rules-and-proofs.tex
\begingroup
\def\olfilename{rules-and-proofs}
\expandafter\def\csname ol@sectioncs\endcsname{section}

\iftag{FOL}
      {\olfileid{fol}{axd}{rul}}
      {\olfileid{pl}{axd}{rul}}

\olsection{قاعدے تے \usetoken{P}{derivation}}

\begin{explain}
  مسلّمی !!{derivation}s شاید منطق لئی سب توں سادہ !!{derivation} نظام
  نیں۔ !!^a{derivation} بس !!{formula}s دا اک سلسلہ اے۔ کسے سلسلے نوں
  !!a{derivation} منّن لئی اوہدے ہر !!{formula} دا یا تاں کسے مسلّمے
  دی اک مثال ہونا لازم اے، یا اوہدے توں پہلاں سلسلے وچ آ چکے اک یا ودھ
  !!{formula}s توں کسے استنباطی قاعدے نال نکلنا لازم اے۔
  !!^a{derivation} اپنے آخری !!{formula} نوں !!{derive}s کردا اے۔
\end{explain}

\begin{defn}[!!^{derivability}]
جے $\Gamma$، $\Lang L$ دے !!{formula}s دا سیٹ اے، تاں~$\Gamma$ توں
اک \emph{!!{derivation}}، !!{formula}s دا اک محدود سلسلہ $!A_1$،
\dots،~$!A_n$ اے، جیہدے ہر $i \le n$ لئی ہِٹھ دتیاں وچوں کوئی اک
گل درست ہووے:
\begin{enumerate}
\item $!A_i \in \Gamma$؛ یا
\item $!A_i$ اک مسلّمہ اے؛ یا
\item $!A_i$، کسے استنباطی قاعدے نال، کسے $!A_j$ (تے $!A_k$) توں
  نکلدا اے جتھے $j < i$ (تے $k < i$)۔
\end{enumerate}
\end{defn}

درست !!{derivation} کِنوں منّنا اے، ایہہ اس گل اُتے منحصر اے کہ اسیں
کیہڑے استنباطی قاعدے منّ دے آں (تے ظاہر اے کیہڑیاں گلاں نوں مسلّمے
منّ دے آں)۔ استنباطی قاعدہ اک جے-تاں بیان اے جیہڑا دسدَا اے کہ کِنہاں
شرطاں ہِٹھ !!a{derivation} دا قدم~$A_i$ اک درست استنباطی قدم اے۔

\begin{defn}[استنباط دا قاعدہ]
\emph{استنباط دا قاعدہ}،~$\Gamma$ توں !!a{derivation} وچ درست
استنباطی قدم ہون لئی اک کافی شرط دیندا اے۔
\end{defn}

مثال لئی، صرف اک عنصر دا سلسلہ $!A$، جتھے $!A \in \Gamma$، ظاہر اے
!!a{derivation} منّیا جاندا اے، ایس لئی ہِٹھ دتا بیان اک بہت سادہ
استنباطی قاعدہ ہو سکدا اے:
\begin{quote}
  جے $!A \in \Gamma$، تاں $!A$،~$\Gamma$ توں کسے وی
  !!{derivation} وچ ہمیشہ اک درست استنباطی قدم اے۔
\end{quote}
ایسے طرح جے $!A$ مسلّماں وچوں اک اے، تاں اکیلا $!A$ وی
!!a{derivation} اے؛ ایس لئی ایہ وی استنباطی قاعدہ اے:
\begin{quote}
  جے $!A$ اک مسلّمہ اے، تاں $!A$ اک درست استنباطی قدم اے۔
\end{quote}
قاعدہ ہور دلچسپ اُدوں ہوندا اے جدوں اوہ زیرِ غور قدم توں پہلاں آ چکے
!!{formula}s دا حوالہ دیندا اے۔ ہِٹھ دتے قاعدے نوں \emph{قاعدۂ اثباتِ
مقدم \pnbLatin{(modus ponens)}} آکھدے نیں:
\begin{quote}
  جے $!B \lif !A$ تے $!B$، !!{derivation} وچ اِس توں اوپر آ چکے ہون،
  تاں~$!A$ اک درست استنباطی قدم اے۔
\end{quote}
جے ایہو اکّو استنباطی قاعدہ ہووے، تاں اُتّے دتی !!{derivation} دی
تعریف دا مطلب ایہ اے: $!A_1$، \dots،~$!A_n$ اُدوں تے صرف اُدوں
!!a{derivation} اے جدوں ہر $i \le n$ لئی ہِٹھ دتیاں وچوں کوئی اک
گل درست ہووے:
\begin{enumerate}
\item $!A_i \in \Gamma$؛ یا
\item $!A_i$ اک مسلّمہ اے؛ یا
\item کسے $j < i$ لئی $!A_j$، $!B \lif !A_i$ ہووے، تے کسے $k < i$
  لئی $!A_k$،~$!B$ ہووے۔
\end{enumerate}
آخری شق آکھدی اے کہ $!A_i$، قاعدۂ اثباتِ مقدم نال~$!A_j$
($!B \lif !A_i$) تے $!A_k$ ($!B$) توں نکلدا اے۔ جے اسیں $1$
توں~$n$ تک جا سکیے تے ہر واری ساڈے کول !!a{formula}~$!A_i$
ایہو جہا ہووے جیہڑا یا تاں~$\Gamma$ وچ ہووے، یا مسلّمہ ہووے، یا
جیہنوں کوئی استنباطی قاعدہ درست قدم دسے، تاں پورا سلسلہ اک درست
!!{derivation} منّیا جاندا اے۔

\begin{defn}[!!^{derivability}]
!!{formula}~$!A$، $\Gamma$ توں \emph{!!{derivable}} اے، جیہنوں
$\Gamma \Proves !A$ لکھدے آں، جے $\Gamma$ توں !!a{derivation}
موجود ہووے جیہڑا $!A$ اُتے ختم ہووے۔
\end{defn}

\begin{defn}[قضیے]
!!^a{formula}~$!A$ اک \emph{قضیہ} اے جے خالی سیٹ توں $!A$ دا
!!a{derivation} موجود ہووے۔ اسیں $\Proves !A$ لکھدے آں جے $!A$
قضیہ اے، تے جے نہیں تاں $\Proves/ !A$ لکھدے آں۔
\end{defn}
\endgroup
% /OLP-0113

% OLP-0114 content/first-order-logic/axiomatic-deduction/axioms-rules-propositional.tex
\begingroup
\def\olfilename{axioms-rules-propositional}
\expandafter\def\csname ol@sectioncs\endcsname{section}

\iftag{FOL}
      {\olfileid{fol}{axd}{prp}}
      {\olfileid{pl}{axd}{prp}}

\olsection{قضیاتی رابطیاں لئی مسلّمے تے قاعدے}

\begin{defn}[مسلّمے]
قضیاتی رابطیاں لئی \emph{مسلّماں} دا سیٹ~$\PAx$ ہِٹھ دتے روپاں والے
سارے !!{formula}s اُتے مشتمل اے:
\begin{align}
  & (!A \land !B) \lif !A \ollabel{ax:land1}\\
  & (!A \land !B) \lif !B \ollabel{ax:land2}\\
  & \formula A \lif (!B \lif (!A \land !B)) \ollabel{ax:land3}\\
  & \formula A \lif (!A \lor !B) \ollabel{ax:lor1}\\
  & \formula A \lif (!B \lor !A) \ollabel{ax:lor2}\\
  & (!A \lif !C) \lif ((!B \lif !C) \lif ((!A \lor !B) \lif !C)) \ollabel{ax:lor3}\\
  & \formula A \lif (!B \lif !A) \ollabel{ax:lif1}\\
  & (!A \lif (!B \lif !C)) \lif ((!A \lif !B) \lif (!A \lif !C)) \ollabel{ax:lif2}\\
  & (!A \lif !B) \lif ((!A \lif \lnot !B) \lif \lnot !A) \ollabel{ax:lnot1}\\
  & \lnot !A \lif (!A \lif !B) \ollabel{ax:lnot2}\\
  & \ltrue \ollabel{ax:ltrue}\\
  & \lfalse \lif !A \ollabel{ax:lfalse1}\\
  & (!A \lif \lfalse) \lif \lnot !A \ollabel{ax:lfalse2}\\
  & \lnot\lnot !A \lif !A \ollabel{ax:dne}
\end{align}
\end{defn}

\begin{defn}[قاعدۂ اثباتِ مقدم]
  جے $!B$ تے $!B \lif !A$ پہلے ای !!a{derivation} وچ آ چکے ہون، تاں $!A$
  اک درست استنباطی قدم اے۔
\end{defn}

اسیں قاعدۂ اثباتِ مقدم نوں مختصر کر کے ``\MP'' لکھاں گے۔
\endgroup
% /OLP-0114

% OLP-0115 content/first-order-logic/axiomatic-deduction/axioms-rules-quantifiers.tex
\begingroup
\def\olfilename{axioms-rules-quantifiers}
\expandafter\def\csname ol@sectioncs\endcsname{section}

\olfileid{fol}{axd}{qua}

\olsection{مسوّراں لئی مسلّمے تے قاعدے}

\begin{defn}[مسوّراں لئی مسلّمے]
مسوّراں اُتے لاگو \emph{مسلّمے} ہِٹھ دتے روپاں دیاں ساریاں مثالاں نیں:
\begin{align}
\ollabel{ax:q1} & \lforall[x][!B] \lif !B(t), \\
\ollabel{ax:q2} & \formula B(t) \lif \lexists[x][!B].
\end{align}
کسے وی بند ٹرم~$t$ لئی۔
\end{defn}

\begin{defn}[مسوّراں لئی قاعدے]
% PNBAXD-013: خاص مستقل تیار ہون والے مسوّری فارمولے وچ وی غیر حاضر ہووے،جیویں صحت مندی دے ثبوت وچ لازم اے۔
\item جے $!B \lif !A(a)$ پہلے ای !!{derivation} وچ آ چکیا ہووے تے
  $a$ نہ~$\Gamma$ وچ آؤندا ہووے، نہ~$!B$ وچ، تے نہ تیار ہون والے
  مسوّری فارمولے وچ، تاں $!B \lif
  \lforall[x][!A(x)]$ اک درست استنباطی قدم اے۔
\item جے $!A(a) \lif !B$ پہلے ای !!{derivation} وچ آ چکیا ہووے تے
  $a$ نہ~$\Gamma$ وچ آؤندا ہووے، نہ~$!B$ وچ، تے نہ تیار ہون والے
  مسوّری فارمولے وچ، تاں $\lexists[x][!A(x)] \lif
  !B$ اک درست استنباطی قدم اے۔
\end{defn}

اسیں ایہناں وچوں ہر اک قاعدے نوں مختصر کر کے ``\QR'' لکھاں گے۔
\endgroup
% /OLP-0115

% OLP-0116 content/first-order-logic/axiomatic-deduction/proving-things.tex
\begingroup
\def\olfilename{proving-things}
\expandafter\def\csname ol@sectioncs\endcsname{section}

\iftag{FOL}
      {\olfileid{fol}{axd}{pro}}
      {\olfileid{pl}{axd}{pro}}

\olsection{\usetoken{P}{derivation} دیاں مثالاں}


\begin{ex}
فرض کرو اسیں $(\lnot !D \lor !E) \lif (!D \lif
!E)$ ثابت کرنا چاہندے آں۔ ظاہر اے ایہ ساڈے کسے وی مسلّمے دی مثال نہیں؛
ایس لئی ایہنوں !!{derive} کرن لئی \MP{} دا قاعدہ ورتنا پئے گا۔ ساڈا
اکّو قاعدہ~MP اے، جیہڑا $!A$ تے $!A \lif !B$ دتے ہون اُتے~$!B$
نوں درست قدم ٹھہران دی اجازت دیندا اے۔ اک طریقہ ایہ ہو سکدا اے کہ
\olref[prp]{ax:lor3} وچ $!A$ دی تھاں $\lnot
!D$، $!B$ دی تھاں $!E$ تے $!C$ دی تھاں $!D \lif !E$ رکھیے،
یعنی ایہ مثال:
\[
(\lnot !D \lif (!D \lif !E)) \lif ((!E \lif (!D \lif !E)) \lif ((\lnot
!D \lor !E) \lif (!D \lif !E))).
\]
کیوں؟ MP نوں دو واری ورتن نال آخری حصہ ملدا اے، تے اوہی اسیں
چاہندے آں۔ اسیں سوکھ نال ویکھ سکدے آں کہ $\lnot !D \lif (!D \lif !E)$،
\olref[prp]{ax:lnot2} دی مثال اے، تے $!E \lif (!D \lif !E)$،
\olref[prp]{ax:lif1} دی مثال اے۔ ایس لئی ساڈا !!{derivation} ایہ اے:
\begin{derivation}
  1. &  $\lnot !D \lif (!D \lif !E)$ & \olref[prp]{ax:lnot2} \\
  2. & $(\lnot !D \lif (!D \lif !E)) \lif {}$\\
  &\qquad $((!E \lif (!D \lif !E)) \lif ((\lnot !D \lor !E) \lif (!D \lif !E)))$ & \olref[prp]{ax:lor3}\\
  3. & $(!E \lif (!D \lif !E)) \lif ((\lnot !D \lor !E) \lif (!D \lif !E))$ & 1, 2, \MP\\
  4. & $!E \lif (!D \lif !E)$ & \olref[prp]{ax:lif1}\\
  5. & $(\lnot !D \lor !E) \lif (!D \lif !E)$ & 3, 4, \MP
\end{derivation}
\end{ex}

\begin{ex}\ollabel{ex:identity}
آؤ $!D \lif !D$ دا !!a{derivation} لبھن دی کوشش کریے۔ ایہ کسے
مسلّمے دی مثال نہیں، ایس لئی ایہنوں !!{derive} کرن لئی \MP{} ورتنا
پئے گا۔ \olref[prp]{ax:lif1}،~$!A \lif !B$ دے روپ دا اک مسلّمہ
اے جیہدے اُتے اسیں~\MP لا سکدے آں۔ پر ایس حالت وچ \MP{} جیہڑے
$!B$ نوں درست قدم ٹھہرائے، اوہ لازماً~$!D \lif !D$ ہونا چاہیدا
اے، کیوں جے اسیں ایہو !!{derive} کرنا چاہندے آں۔ ایس دا مطلب اے
کہ $!A$ وی $!D$ ہونا چاہیدا اے؛ یعنی اسیں
\olref[prp]{ax:lif1} دی ایہ مثال ویکھ سکدے آں:
\[
!D \lif (!D \lif !D)
\]
\MP ورتن لئی ساڈے کول اوہدے مطابق دوجے مقدمے، یعنی~$!A$، دا
جواز وی ہونا چاہیدا اے۔ پر ایتھے اوہ~$!D$ ہووے گا، تے اسیں
اکیلا~$!D$ !!{derive} نہیں کر سکاں گے۔ ایس لئی سانوں ہور
طریقہ چاہیدا اے۔

صرف~$\lif$ نال متعلق دوجا مسلّمہ \olref[prp]{ax:lif2} اے، یعنی
\[
(!A \lif (!B \lif !C)) \lif ((!A \lif !B) \lif (!A \lif !C))
\]
\MP{} نوں دو واری ورت کے اسیں آخری اندرلا شرطیہ حاصل کر سکدے آں۔
فیر ایہدا مطلب اے کہ سانوں \olref[prp]{ax:lif2} دی ایسی مثال
چاہیدی اے جیہدے وچ $!A \lif !C$، اوہ !!{formula}~$!D \lif !D$
ہووے جیہنوں اسیں ثابت کرنا چاہندے آں۔ ایس لئی $!A$ تے $!C$
دوویں~$!D$ نیں۔ ہُن $!B$ کیہڑا چُنیے تاں جو $!A \lif (!B \lif !C)$
تے $!A \lif !B$، یعنی ایتھے $!D \lif (!B \lif !D)$ تے $!D \lif !B$،
دوویں !!{derivable} ہون؟ پہلا تاں $!B$ دی ہر چُون لئی
\olref[prp]{ax:lif1} دی مثال اے۔ تے $!D \lif !B$ وی
\olref[prp]{ax:lif1} دی مثال ہووے گا جے $!B$، $(!D \lif !D)$ ہووے۔
ایس لئی ساڈا !!{derivation} ایہ اے:
\begin{derivation}
  1. & $!D \lif ((!D \lif !D) \lif !D)$ & \olref[prp]{ax:lif1}\\
  2. & $(!D \lif ((!D \lif !D) \lif !D)) \lif {}$\\
  & \qquad $((!D \lif (!D \lif !D)) \lif (!D \lif !D))$ & \olref[prp]{ax:lif2}\\
  3. & $(!D \lif (!D \lif !D)) \lif (!D \lif !D)$ & 1, 2, \MP\\
  4. & $!D \lif (!D \lif !D)$ & \olref[prp]{ax:lif1}\\
  5. & $!D \lif !D$ & 3, 4, \MP
\end{derivation}
\end{ex}

\begin{ex}\ollabel{ex:chain}
کدے اسیں دکھانا چاہندے آں کہ کجھ ہور !!{formula}s~$\Gamma$ توں
کسے !!{formula} دا !!a{derivation} موجود اے۔ مثال لئی، آؤ
دکھائیے کہ اسیں $!A \lif !C$، $\Gamma = \{!A \lif !B,
!B \lif !C\}$ توں !!{derive} کر سکدے آں۔
\begin{derivation}
  1. & $!A \lif !B$ & \Hyp\\
  2. & $!B \lif !C$ & \Hyp\\
  3. & $(!B \lif !C) \lif (!A \lif (!B \lif !C))$ & \olref[prp]{ax:lif1} \\
  4. & $!A \lif (!B \lif !C)$ & 2, 3, \MP\\
  5. & $(!A \lif (!B \lif !C)) \lif {}$\\
  & \qquad $((!A \lif !B) \lif (!A \lif !C))$ & \olref[prp]{ax:lif2}\\
  6. & $((!A \lif !B) \lif (!A \lif !C))$ & 4, 5, \MP\\
  7. & $!A \lif !C$ & 1, 6, \MP
\end{derivation}
``\Hyp'' نال نشان لگیاں سطراں (``مفروضہ'' لئی) دسدیاں
نیں کہ اس سطر دا !!{formula}،~$\Gamma$ دا !!a{element} اے۔
\end{ex}

\begin{prop}
  \ollabel{prop:chain} جے $\Gamma \Proves !A \lif !B$ تے $\Gamma
  \Proves !B \lif !C$، تاں $\Gamma \Proves !A \lif !C$
\end{prop}

\begin{proof}
  فرض کرو $\Gamma \Proves !A \lif !B$ تے $\Gamma \Proves !B \lif
  !C$۔ تاں $!A \lif !B$ دا !!a{derivation}~$\Gamma$ توں موجود
  اے، تے $!B \lif !C$ دا !!a{derivation}~$\Gamma$ توں وی۔ اوہناں
  نوں اک دوجے دے پچھے جوڑ کے اکّو !!{derivation} بنا لو۔ ہُن پچھلی
  مثال دے !!{derivation} دیاں سطراں 3--7 شامل کرو۔ ایس نال
  $!A \lif !C$ دا !!a{derivation}، جیہدے وچ ایہ فارمولا نویں
  !!{derivation} دی آخری سطر اے،~$\Gamma$ توں ملدا اے۔ دھیان دیو
  کہ سطراں 4 تے~7 دے جواز درست رہندے نیں جے سطر نمبر~1 دے
  حوالے نوں $!A \lif !B$ دے !!{derivation} دی آخری سطر دے
  حوالے نال، تے سطر نمبر~2 دے حوالے نوں $!B \lif !C$ دے
  !!{derivation} دی آخری سطر دے حوالے نال بدل دتا جائے۔
\end{proof}

\begin{prob}
مسلّماں توں !!{derivation}s پیش کر کے دکھاؤ کہ ہِٹھ دتیاں گلاں
درست نیں:
\begin{enumerate}
\item $(!A \land !B) \lif (!B \land !A)$
\item $((!A \land !B) \lif !C) \lif (!A \lif (!B \lif !C))$
\item $\lnot(!A \lor !B) \lif \lnot !A$
\end{enumerate}
\end{prob}
\endgroup
% /OLP-0116

% OLP-0117 content/first-order-logic/axiomatic-deduction/proving-things-quant.tex
\begingroup
\def\olfilename{proving-things-quant}
\expandafter\def\csname ol@sectioncs\endcsname{section}

\olfileid{fol}{axd}{prq}

\olsection{مسوّراں نال \usetoken{P}{derivation}}

\begin{ex}
آؤ $(\lforall[x][!A(x)] \land
\lforall[y][!B(y)]) \lif \lforall[x][(!A(x) \land !B(x))]$ دا
!!a{derivation} پیش کریے۔

پہلاں، غور کرو کہ
% PNBAXD-011: make the omitted composition and second conditional axiom explicit; all displayed formulas and citations remain unchanged.
\begin{align*}
  (\lforall[x][!A(x)] \land \lforall[y][!B(y)]) & \lif \lforall[x][!A(x)]\\
  \intertext{ایہ \olref[prp]{ax:land1} دی مثال اے، تے}
  \lforall[x][!A(x)] & \lif !A(a) \\
  \intertext{\olref[qua]{ax:q1} دی مثال اے۔ ایس لئی \olref[pro]{prop:chain} توں سانوں پتا لگدا اے کہ}
  (\lforall[x][!A(x)] \land \lforall[y][!B(y)]) & \lif !A(a) \\
  \intertext{!!{derivable} اے۔ ایسے طرح، چونکہ}
  (\lforall[x][!A(x)] \land \lforall[y][!B(y)]) & \lif \lforall[y][!B(y)] \qquad\text{تے}\\
    \lforall[y][!B(y)] & \lif !B(a)\\
    \intertext{ترتیب نال \olref[prp]{ax:land2} تے \olref[qua]{ax:q1} دیاں مثالاں نیں،}
    (\lforall[x][!A(x)] \land \lforall[y][!B(y)]) & \lif !B(a)\\
    \intertext{\olref[pro]{prop:chain} نال اخذ ہو سکدا اے۔ پہلے جزو لئی حاصل شرطیہ نتیجے نوں \olref[prp]{ax:land3} دی مناسب مثال نال ایس زنجیری نتیجے راہیں ملاؤ۔ فیر شرطیہ رابطے دے دوجے مسلّمے تے دوجے جزو لئی حاصل شرطیہ نتیجے نال~\MP دی دو وار ورتوں توں اسیں ویکھدے آں کہ}
    (\lforall[x][!A(x)] \land \lforall[y][!B(y)]) & \lif (!A(a) \land !B(a))\\
    \intertext{اخذ ہو سکدا اے۔ ہُن \QR{} ورت کے اسیں ایہ حاصل کر سکدے آں:}
(\lforall[x][!A(x)] \land \lforall[y][!B(y)]) & \lif \lforall[x][(!A(x) \land !B(x))].
\end{align*}
\end{ex}
\endgroup
% /OLP-0117

% OLP-0118 content/first-order-logic/axiomatic-deduction/proof-theoretic-notions.tex
\begingroup
\def\olfilename{proof-theoretic-notions}
\expandafter\def\csname ol@sectioncs\endcsname{section}

\iftag{FOL}
      {\olfileid{fol}{axd}{ptn}}
      {\olfileid{pl}{axd}{ptn}}

\olsection{ثبوتی-نظریاتی تصورات \pnbLatin{(Proof-Theoretic Notions)}}

\begin{explain}
جیویں اسیں کئی اہم معنوی تصورات
(\iftag{FOL}{منطقی اعتبار}{ہمیشہ صادق ہونا}، منطقی نتیجہ، پورا ہون
جوگتا) دی تعریف دتی اے، ہُن اسیں اوہناں دے مطابق
\emph{ثبوتی-نظریاتی تصورات} دی تعریف دیندے آں۔ ایہناں دی تعریف
!!{sentence}s دے !!{structure}s وچ پورا اترن ول رجوع کر کے نہیں،
بلکہ کجھ فارمولیاں دے !!{derivability} یا !!{nonderivability} ول
رجوع کر کے دتی جاندی اے۔ ایہناں دوہاں طرح دے تصورات دا اک دوجے
نال میل کھانا اک اہم دریافت سی۔ ایس میل دا بیان \emph{صحت مندی}
تے \emph{تمامیت دے قضیاں} وچ اے۔
\end{explain}

\begin{defn}[!!^{derivability}]
!!^a{formula}~$!A$، $\Gamma$ توں \emph{!!{derivable}} اے، جیہنوں
$\Gamma \Proves !A$ لکھدے آں، جے~$\Gamma$ توں !!a{derivation}
موجود ہووے جیہڑا $!A$ اُتے ختم ہووے۔
\end{defn}

\begin{defn}[قضیے]
!!^a{formula}~$!A$ اک \emph{قضیہ} اے جے خالی سیٹ توں $!A$ دا
!!a{derivation} موجود ہووے۔ اسیں $\Proves !A$ لکھدے آں جے $!A$
قضیہ اے، تے $\Proves/ !A$ لکھدے آں جے نہیں۔
\end{defn}

\begin{defn}[بے تضادی]
!!{formula}s دا سیٹ~$\Gamma$ \emph{بے تضاد} اے اُدوں تے صرف اُدوں
جدوں $\Gamma\Proves/ \lfalse$؛ نہیں تاں ایہ \emph{متضاد} اے۔
\end{defn}

\begin{prop}[انعکاسیت]
\ollabel{prop:reflexivity}
جے $!A \in \Gamma$، تاں $\Gamma \Proves !A$۔
\end{prop}

\begin{proof}
  !!{formula}~$!A$ اپنے آپ وچ~$!A$ دا~$\Gamma$ توں
  !!a{derivation} اے۔
\end{proof}

\begin{prop}[یک رُخیت]
\ollabel{prop:monotonicity}
جے $\Gamma \subseteq \Delta$ تے $\Gamma \Proves !A$، تاں $\Delta
\Proves !A$۔
\end{prop}

\begin{proof}
$!A$ دا~$\Gamma$ توں کوئی وی !!{derivation}، $!A$ دا~$\Delta$
توں وی !!a{derivation} اے۔
\end{proof}

\begin{prop}[تعدیت]
\ollabel{prop:transitivity}
جے $\Gamma \Proves !A$ تے $\{!A\} \cup \Delta \Proves
!B$، تاں $\Gamma \cup \Delta \Proves !B$۔
\end{prop}

\begin{proof}
  فرض کرو $\{!A\} \cup \Delta \Proves !B$۔ تاں اک
  !!a{derivation} $!B_1$، \dots، $!B_l = !B$،~$\{!A\} \cup
  \Delta$ توں موجود اے۔ اوس !!{derivation} دے کجھ قدم ایسے
  قاعدے دی وجہ توں درست ہون گے جیہڑا پہلے آ چکی سطر~$!B_i = !A$
  دا حوالہ دیندا اے۔ مفروضے مطابق $!A$ دا~$\Gamma$ توں
  !!a{derivation} موجود اے، یعنی !!a{derivation}~$!A_1$، \dots،
  $!A_k = !A$، جیہدے ہر $!A_i$ یا تاں مسلّمہ اے، یا~$\Gamma$
  دا !!a{element} اے، یا استنباط دے قاعدے نال درست اے۔ ہُن
  ایہ سلسلہ ویکھو:
  \[
  !A_1, \dots, !A_k = !A, !B_1, \dots, !B_l = !B.
  \]
% PNBAXD-001: آخری وضاحت وچ B_i دے چھڈے ہوے فارمولا نشان نوں بحال کیتا۔
% PNBAXD-002: تین ممکن جوازاں لئی ماخذ دے صرف قاعدہ آکھن نوں جواز کیتا۔
  ایہ~$!B$ دا $\Gamma \cup \Delta$ توں درست !!{derivation} اے،
  کیوں جے ہر $!B_i = !A$ ہُن اوہو جواز رکھدا اے جیہڑا~$!A_k = !A$
  نوں درست ٹھہراندا اے۔
\end{proof}

دھیان دیو کہ ایس دا خاص مطلب ایہ اے: جے $\Gamma \Proves !A$ تے $!A
\Proves !B$، تاں $\Gamma \Proves !B$۔ ایہ وی نکلدا اے کہ جے $!A_1,
\dots, !A_n \Proves !B$ تے $\Gamma \Proves !A_i$ ہر~$i$ لئی، تاں
$\Gamma \Proves !B$۔

\begin{prop}
\ollabel{prop:incons}
$\Gamma$ اُدوں تے صرف اُدوں متضاد اے جدوں $\Gamma \Proves !A$ ہر~$!A$ لئی۔
\end{prop}

\begin{proof}
مشق۔
\end{proof}

\tagprob{FOL}
\begin{prob}
\olref[fol][axd][ptn]{prop:incons} ثابت کرو۔
\end{prob}
\tagendprob

\tagprob{notFOL}
\begin{prob}
\olref[pl][axd][ptn]{prop:incons} ثابت کرو۔
\end{prob}
\tagendprob

\begin{prop}[محدود گواہی (Compactness)]
\ollabel{prop:proves-compact}
  \begin{enumerate}
  \item جے $\Gamma \Proves !A$، تاں اک محدود ذیلی سیٹ $\Gamma_0
    \subseteq \Gamma$ ایہو جہا اے کہ $\Gamma_0 \Proves !A$۔
  \item جے~$\Gamma$ دا ہر محدود ذیلی سیٹ بے تضاد اے، تاں
    $\Gamma$ بے تضاد اے۔
  \end{enumerate}
\end{prop}

\begin{proof}
  \begin{enumerate}
% PNBAXD-012: پہلی درجے دے اخذ وچ مسوّری قاعدیاں تے اوہناں دی غیر حاضری شرط نوں وی شامل کیتا۔
    \item جے $\Gamma \Proves !A$، تاں !!{formula}s دا اک محدود
      سلسلہ $!A_1$، \dots،~$!A_n$ موجود اے جیہدے لئی
      $!A \ident !A_n$ اے؛ ہر $!A_i$ یا تاں منطقی مسلّمہ اے، یا
      ~$\Gamma$ دا !!a{element} اے، یا پچھلے !!{formula}s توں
      قاعدۂ اثباتِ مقدم نال، یا پہلی درجے دی منطق وچ مسوّراں دے
      مقررہ استنباطی قاعدے نال، نکلدا اے۔ $\Gamma_0$ اوہناں $!A_i$
      دا سیٹ لو جیہڑے~$\Gamma$ وچ نیں۔ تاں اوہ !!{derivation}
      ایسے طرح~$\Gamma_0$ توں وی !!a{derivation} اے: مسوّری قاعدے
      دا خاص مستقل اصل مفروضیاں وچ غیر حاضر سی، ایس لئی چُنے ہوے
      ذیلی سیٹ وچ وی غیر حاضر رہندا اے؛قاعدے دیاں باقی غیر حاضری
      شرطاں وی نہیں بدلّدیاں۔ تے ایس لئی
      $\Gamma_0 \Proves !A$۔
    \item ایہہ (1) دا برعکس بالمخالف اے، اوس خاص حالت وچ جدوں $!A
      \ident \lfalse$۔
\end{enumerate}
\end{proof}
\endgroup
% /OLP-0118

% OLP-0119 content/first-order-logic/axiomatic-deduction/deduction-theorem.tex
\begingroup
\def\olfilename{deduction-theorem}
\expandafter\def\csname ol@sectioncs\endcsname{section}

\iftag{FOL}
      {\olfileid{fol}{axd}{ded}}
      {\olfileid{pl}{axd}{ded}}

\olsection{استنباط دا قضیہ \pnbLatin{(Deduction Theorem)}}

جیویں اسیں ویکھیا اے، مسلّمی نظام وچ !!{derivation}s لکھنا اوکھا اے
تے !!{derivation}s لبھنا وی مشکل ہو سکدا اے۔ !!{formula}s دیاں
لمیاں فہرستاں عملاً لکھن دی تھاں عموماً کجھ سادے نتیجے ورت کے ایہہ
دسنا سوکھا ہوندا اے کہ ایسے !!{derivation}s موجود نیں۔ اسیں پہلے
ایہو جیہے تین نتیجے ثابت کر چکے آں:
\olref[ptn]{prop:reflexivity} دسدا اے کہ اسیں ہمیشہ $\Gamma
\Proves !A$ آکھ سکدے آں، جدوں ایہہ معلوم ہووے کہ $!A \in
\Gamma$۔ \olref[ptn]{prop:monotonicity} دسدا اے کہ جے $\Gamma \Proves !A$
تاں $\Gamma \cup \{!B\} \Proves !A$ وی اے۔ تے
\olref[ptn]{prop:transitivity} توں نکلدا اے کہ جے $\Gamma \Proves !A$
تے $!A \Proves !B$، تاں $\Gamma \Proves !B$۔ ہُن اک ہور سادا
نتیجہ ویکھو: قاعدۂ اثباتِ مقدم دا اک ``میٹا'' روپ:

\begin{prop}
  \ollabel{prop:mp} جے $\Gamma \Proves !A$ تے $\Gamma \Proves !A \lif
  !B$، تاں $\Gamma \Proves !B$۔
\end{prop}

\begin{proof}
  ساڈے کول $\{!A, !A \lif !B\} \Proves !B$ اے:
  \begin{derivation}
    1. & $!A$ & مفروضہ\\
    2. & $!A \lif !B$ & مفروضہ\\
    3. & $!B$ & 1, 2, MP
  \end{derivation}
  \olref[ptn]{prop:transitivity} نال $\Gamma \Proves !B$ ملدا اے۔
\end{proof}

ایس سیاق وچ ساڈا سب توں اہم نتیجہ استنباط دا قضیہ اے:

\begin{thm}[استنباط دا قضیہ]
\ollabel{thm:deduction-thm} $\Gamma \cup \{!A\} \Proves !B$ اُدوں تے
صرف اُدوں جدوں $\Gamma \Proves !A \lif !B$۔
\end{thm}

\begin{proof}
``جے'' والا رخ فوراً مل جاندا اے۔ جے $\Gamma \Proves !A \lif !B$
تاں \olref[ptn]{prop:monotonicity} دے مطابق $\Gamma \cup \{!A\} \Proves !A \lif !B$
وی اے۔ \olref[ptn]{prop:reflexivity} توں $\Gamma \cup \{!A\} \Proves !A$
وی اے۔ ایس لئی \olref{prop:mp} نال $\Gamma \cup
\{!A\} \Proves !B$ ملدا اے۔

``صرف اُدوں'' والے رخ لئی اسیں $!B$ دے $\Gamma \cup \{!A\}$
توں !!{derivation} دی لمبائی اُتے استقرا کردے آں۔

استقرا دی بنیاد وچ لمبائی~$1$ دے ہر !!{derivation} لئی دعویٰ
ثابت کردے آں۔ $!B$ دا $\Gamma \cup \{!A\}$ توں لمبائی~$1$
والا !!^a{derivation} اکّلے $!B$ اُتے مشتمل اے؛ جے ایہہ درست اے
تاں $!B$ یا $\in \Gamma \cup \{!A\}$ اے یا مسلّمہ اے۔ جے
$!B \in \Gamma$ اے یا مسلّمہ اے، تاں $\Gamma \Proves !B$۔
نالے \olref[prp]{ax:lif1} توں $\Gamma
\Proves !B \lif (!A \lif !B)$ وی اے؛ فیر \olref{prop:mp} توں
$\Gamma \Proves !A \lif !B$ ملدا اے۔ جے $!B \in \{ !A\}$
اے، تاں $\Gamma \Proves !A \lif !B$، کیوں جے ایس حالت وچ آخری
!!{sentence}~$!A \lif !B$ اوہی $!A \lif !A$ اے، جیہڑا اسیں
\olref[pro]{ex:identity} وچ !!{derive} کیتا اے۔

استقرائی قدم وچ فرض کرو $!B$ دا $\Gamma \cup \{!A\}$ توں
!!a{derivation} ایسے قدم~$!B$ اُتے ختم ہوندا اے جیہڑا قاعدۂ
اثباتِ مقدم نال درست ٹھہردا اے۔ (جے آخری قدم ایس قاعدے نال درست
نہیں ٹھہردا، تاں $!B \in \Gamma$ اے، $!B
\ident !A$ اے، یا $!B$ مسلّمہ اے؛ تے استقرا دی بنیاد والی
اوہی دلیل لاگو ہوندی اے۔) تاں !!{derivation} دے کجھ پچھلے قدم
$!C \lif !B$ تے $!C$ نیں، کسے !!{formula}~$!C$ لئی؛ یعنی
$\Gamma \cup \{!A\} \Proves !C \lif !B$ تے $\Gamma \cup \{!A\}
\Proves !C$۔ ایہناں دے متعلقہ !!{derivation}s چھوٹے نیں، ایس لئی
استقرائی مفروضہ اوہناں اُتے لاگو ہوندا اے۔ سو ساڈے کول دوویں نیں:
\begin{align*}
  & \Gamma \Proves !A \lif (!C \lif !B); \\
  & \Gamma \Proves !A \lif !C.
\end{align*}
نالے \olref[prp]{ax:lif2} توں
\[
\Gamma \Proves (!A \lif (!C \lif !B)) \lif
((!A\lif !C)  \lif (!A \lif !B)),
\]
تے \olref{prop:mp} دو واری لا کے مطلوبہ $\Gamma \Proves !A \lif !B$
ملدا اے۔
\end{proof}

دھیان دیو کہ \olref[prp]{ax:lif1} تے \olref[prp]{ax:lif2}
خاص ایس طرح چُنے گئے سن کہ استنباط دا قضیہ درست رہے۔

!!{derivability} بارے ہٹھ لکھے کجھ کارآمد حقائق ثبوت لئی مشق
وجوں چھڈدے آں۔

\begin{prop}
  \ollabel{prop:derivfacts}
\begin{enumerate}
% PNBAXD-003: پہلے نتیجے دے دو کھلے رہ گئے بریکٹ بند کیتے۔
\item $\Proves (!A \lif !B) \lif ((!B \lif !C)
  \lif (!A \lif !C))$؛ \ollabel{derivfacts:a}
\item جے $\Gamma \cup \{ \lnot !A\}
  \Proves \lnot !B$، تاں $\Gamma \cup \{ !B\} \Proves
  !A$ (برعکس بالمخالف؛ Contraposition)؛ \ollabel{derivfacts:b}
\item $\{ !A, \lnot!A\} \Proves
    !B$ (تضاد توں ہر نتیجہ؛ Ex Falso Quodlibet, Explosion)؛ \ollabel{derivfacts:c}
\item $\{ \lnot\lnot!A\} \Proves
  !A$ (دوہری نفی دا خاتمہ؛ Double Negation Elimination)؛\ollabel{derivfacts:d}
\item جے $\Gamma \Proves \lnot\lnot!A$، تاں $\Gamma \Proves
  !A$؛\ollabel{derivfacts:e}
\end{enumerate}
\end{prop}

\tagprob{FOL}
\begin{prob}
  \olref[fol][axd][ded]{prop:derivfacts} ثابت کرو۔
\end{prob}
\tagendprob

\tagprob{notFOL}
\begin{prob}
  \olref[pl][axd][ded]{prop:derivfacts} ثابت کرو۔
\end{prob}
\tagendprob
\endgroup
% /OLP-0119

% OLP-0120 content/first-order-logic/axiomatic-deduction/deduction-theorem-quantifiers.tex
\begingroup
\def\olfilename{deduction-theorem-quantifiers}
\expandafter\def\csname ol@sectioncs\endcsname{section}

\olfileid{fol}{axd}{ddq}

\olsection{مسوّراں سمیت استنباط دا قضیہ}

\begin{thm}[استنباط دا قضیہ]
\ollabel{thm:deduction-thm-q} جے $\Gamma \cup \{!A\} \Proves !B$،
تاں $\Gamma \Proves !A \lif !B$۔
\end{thm}

\begin{proof}
اسیں فیر $!B$ دے $\Gamma \cup \{!A\}$ توں !!{derivation} دی
لمبائی اُتے استقرا کردے آں۔

استقرا دی بنیاد دا ثبوت بالکل اوہی اے جیہڑا
\olref[ded]{thm:deduction-thm} دے ثبوت وچ اے۔

استقرائی قدم لئی فیر فرض کرو کہ $!B$ دا !!{derivation}،
$\Gamma \cup \{!A\}$ توں، ایسے قدم~$!B$ اُتے ختم ہوندا اے جیہڑا
استنباط دے قاعدے نال درست ٹھہردا اے۔ جے ایہ قاعدہ اثباتِ مقدم اے،
تاں اسیں \olref[ded]{thm:deduction-thm} دے ثبوت وانگ اگے ودھدے آں۔
جے قاعدہ \QR اے، تاں اسیں جاندے آں کہ $!B \ident !C \lif \lforall[x][!D(x)]$
تے $!C \lif !D(a)$ دے روپ دا !!a{formula} پہلے !!{derivation}
وچ آ چکیا اے؛ ایتھے $a$ نہ~$!C$ وچ، نہ~$!A$ وچ، تے نہ~$\Gamma$
وچ آؤندا اے۔ ایس لئی ساڈے کول اے:
% PNBAXD-005: دوجی میٹا شرطیہ وچ اک کھلا بریکٹ بند کیتا۔
\begin{align*}
  \Gamma \cup \{!A\} & \Proves !C \lif !D(a),\\
  \intertext{تے استقرائی مفروضہ لاگو ہوندا اے، یعنی ساڈے کول اے}
    \Gamma & \Proves !A \lif (!C \lif !D(a)).\\
  \intertext{ایس توں}
  & \Proves (!A \lif (!C \lif !D(a))) \lif ((!A \land !C) \lif !D(a))\\
  \intertext{تے قاعدۂ اثباتِ مقدم نال ملدا اے}
  \Gamma & \Proves (!A \land !C) \lif !D(a).\\
  \intertext{کیوں جے خاص مستقل دی غیر حاضری والی شرط ہُن وی پوری اے، اسیں ایس !!{derivation} وچ \QR نال درست ہون والا اک قدم ودھا سکدے آں، تے حاصل کردے آں}
    \Gamma & \Proves (!A \land !C) \lif \lforall[x][!D(x)].\\
    \intertext{ساڈے کول ایہ وی اے}
    & \Proves ((!A \land !C) \lif \lforall[x][!D(x)]) \lif (!A \lif (!C \lif \lforall[x][!D(x)])),\\
    \intertext{لہٰذا قاعدۂ اثباتِ مقدم نال}
    \Gamma & \Proves !A \lif (!C \lif \lforall[x][!D(x)]),
\end{align*}
% PNBAXD-004: آخری نتیجے وچ ماخذ توں رہ گیا A والا شرطیہ بحال کیتا۔
یعنی $\Gamma \Proves !A \lif !B$۔

اوہ حالت جتھے $!B$ قاعدے \QR نال درست ٹھہردا اے، پر اوہدا روپ
$\lexists[x][!D(x)] \lif !C$ اے، مشق لئی چھڈدے آں۔
\end{proof}

\begin{prob}
  \olref[fol][axd][ddq]{thm:deduction-thm-q} دا ثبوت پورا کرو۔
\end{prob}
\endgroup
% /OLP-0120

% OLP-0121 content/first-order-logic/axiomatic-deduction/provability-consistency.tex
\begingroup
\def\olfilename{provability-consistency}
\expandafter\def\csname ol@sectioncs\endcsname{section}

\iftag{FOL}
      {\olfileid{fol}{axd}{prv}}
      {\olfileid{pl}{axd}{prv}}

\olsection{\usetoken{S}{derivability} تے بے تضادی}

ہُن اسیں !!{derivability} دے تعلق دیاں کئی خاصیتاں ثابت کراں گے۔
ایہ اپنے آپ وچ وی دلچسپ نیں، پر تمامیت دے قضیے دے ثبوت وچ
ہر اک دا اپنا کردار ہووے گا۔

\begin{prop}\ollabel{prop:provability-contr}
  جے $\Gamma \Proves !A$ تے $\Gamma \cup \{!A\}$ متضاد اے،
  تاں $\Gamma$ متضاد اے۔
\end{prop}

\begin{proof}
  جے $\Gamma \cup \{!A\}$ متضاد اے، تاں $\Gamma \cup \{!A\}
  \Proves \lfalse$۔ \olref[ptn]{prop:reflexivity} توں
  $\Gamma \Proves !B$ ہر اس $!B \in \Gamma$ لئی۔ مفروضے مطابق
  $\Gamma \Proves !A$ وی اے، ایس لئی $\Gamma \Proves !B$
  ہر اس $!B \in \Gamma \cup
  \{!A\}$ لئی۔ \olref[ptn]{prop:transitivity}
  توں $\Gamma \Proves
  \lfalse$، یعنی $\Gamma$ متضاد اے۔
\end{proof}

\begin{prop}
\ollabel{prop:prov-incons}
$\Gamma \Proves !A$ اُدوں تے صرف اُدوں جدوں $\Gamma \cup \{\lnot !A\}$ متضاد اے۔
\end{prop}

\begin{proof}
پہلاں فرض کرو $\Gamma \Proves !A$۔ تاں \olref[ptn]{prop:monotonicity}
توں $\Gamma \cup \{\lnot !A\}
\Proves !A$۔ \olref[ptn]{prop:reflexivity} توں $\Gamma \cup \{\lnot
!A\} \Proves \lnot !A$۔ ساڈے کول \olref[prp]{ax:lnot2} توں
$\Proves \lnot !A \lif (!A \lif \lfalse)$ وی اے۔ ایس لئی
\olref[ded]{prop:mp} دو واری لا کے $\Gamma \cup \{\lnot
!A\} \Proves \lfalse$ ملدا اے۔

ہُن فرض کرو $\Gamma \cup \{\lnot !A\}$ متضاد اے، یعنی $\Gamma
\cup \{\lnot !A\} \Proves \lfalse$۔ استنباط دے قضیے نال $\Gamma
\Proves \lnot !A \lif \lfalse$۔ \olref[prp]{ax:lfalse2} توں $\Gamma
\Proves (\lnot !A \lif
\lfalse) \lif \lnot\lnot !A$، ایس لئی \olref[ded]{prop:mp} نال $\Gamma
\Proves \lnot\lnot !A$۔ چونکہ $\Gamma \Proves
\lnot\lnot !A \lif !A$ (\olref[prp]{ax:dne})، ایس لئی فیر
\olref[ded]{prop:mp} نال $\Gamma
\Proves !A$۔
\end{proof}

\begin{prob}
ثابت کرو کہ $\Gamma \Proves \lnot !A$ اُدوں تے صرف اُدوں جدوں
$\Gamma \cup \{!A\}$ متضاد اے۔
\end{prob}

\begin{prop}\ollabel{prop:explicit-inc}
  جے $\Gamma \Proves !A$ تے $\lnot !A \in \Gamma$، تاں $\Gamma$
  متضاد اے۔
\end{prop}

\begin{proof}
  \olref[prp]{ax:lnot2} توں $\Gamma \Proves \lnot !A \lif (!A \lif \lfalse)$۔
  \olref[ded]{prop:mp} دو واری لا کے $\Gamma \Proves \lfalse$۔
\end{proof}


\begin{prop}\ollabel{prop:provability-exhaustive}
  جے $\Gamma \cup \{!A\}$ تے $\Gamma \cup \{\lnot !A\}$
  دوویں متضاد نیں، تاں $\Gamma$ متضاد اے۔
\end{prop}

\begin{proof}
مشق۔
\end{proof}

\tagprob{FOL}
\begin{prob}
  \olref[fol][axd][prv]{prop:provability-exhaustive} ثابت کرو۔
\end{prob}
\tagendprob

\tagprob{notFOL}
\begin{prob}
  \olref[pl][axd][prv]{prop:provability-exhaustive} ثابت کرو۔
\end{prob}
\tagendprob
\endgroup
% /OLP-0121

% OLP-0122 content/first-order-logic/axiomatic-deduction/provability-propositional.tex
\begingroup
\def\olfilename{provability-propositional}
\expandafter\def\csname ol@sectioncs\endcsname{section}

\iftag{FOL}
      {\olfileid{fol}{axd}{ppr}}
      {\olfileid{pl}{axd}{ppr}}

\olsection{\usetoken{S}{derivability} تے قضیاتی رابطے}

\begin{explain}
  اسیں ثابت کراں گے کہ مسلّمی اخذ دا !!{derivability} تعلق~$\Proves$
  قضیاتی رابطیاں نال متعلق کجھ بنیادی حقیقتاں ثابت کرن لئی کافی مضبوط اے؛
  مثلاً $!A \land !B
  \Proves !A$ تے $!A, !A \lif !B \Proves !B$ (قاعدۂ اثباتِ مقدم)۔
  ایہ حقیقتاں تمامیت دے قضیے دے ثبوت لئی ضروری نیں۔
\end{explain}

\begin{prop}\ollabel{prop:provability-land}
  \begin{enumerate}
  \item \ollabel{prop:provability-land-left} $!A \land !B \Proves
    !A$ تے $!A \land !B \Proves !B$، دوویں۔
  \item \ollabel{prop:provability-land-right} $!A, !B \Proves !A \land !B$۔
  \end{enumerate}
\end{prop}

\begin{proof}
  \begin{enumerate}
    % PNBAXD-006: دوجے اخذ دا حوالہ ax:land2 اے؛ منجمد ماخذ وچ ax:land1 دُہرایا گیا اے۔
    \item \olref[prp]{ax:land1} تے \olref[prp]{ax:land2} توں
      قاعدۂ اثباتِ مقدم نال۔
  \item \olref[prp]{ax:land3} توں قاعدۂ اثباتِ مقدم دو واری لا کے۔
  \end{enumerate}
\end{proof}


\begin{prop}\ollabel{prop:provability-lor}
  \begin{enumerate}
  \item $!A \lor !B, \lnot !A, \lnot !B$ متضاد اے۔
  \item $!A \Proves !A \lor !B$ تے $!B \Proves !A \lor !B$، دوویں۔
  \end{enumerate}
\end{prop}

\begin{proof}
  \begin{enumerate}
  % PNBAXD-007: ایتھے درکار مسلّمہ ax:lnot2 اے؛ منجمد ماخذ وچ ax:lnot1 دا حوالہ اے۔
  \item \olref[prp]{ax:lnot2} توں $\Proves \lnot !A \lif (!A
    \lif \lfalse)$ تے $\Proves \lnot !B \lif (!B \lif \lfalse)$
    ملدے نیں۔ ایس لئی استنباط دے قضیے نال ساڈے کول
    $\{\lnot !A\} \Proves !A \lif
    \lfalse$ تے $\{\lnot !B\} \Proves !B \lif \lfalse$ نیں۔
    \olref[prp]{ax:lor3} توں $\{\lnot !A, \lnot !B\} \Proves (!A
    \lor !B) \lif \lfalse$ ملدا اے۔ استنباط دے قضیے نال
    $\{!A \lor !B, \lnot !A, \lnot !B\} \Proves \lfalse$۔
  \item \olref[prp]{ax:lor1} تے \olref[prp]{ax:lor2} توں
    قاعدۂ اثباتِ مقدم نال۔
  \end{enumerate}
\end{proof}

\begin{prop}\ollabel{prop:provability-lif}
  \begin{enumerate}
  \item \ollabel{prop:provability-lif-left}  $!A, !A \lif !B \Proves !B$۔
  \item \ollabel{prop:provability-lif-right}
    $\lnot !A \Proves !A \lif !B$ تے $!B \Proves !A \lif !B$، دوویں۔
  \end{enumerate}
\end{prop}

\begin{proof}
  \begin{enumerate}
  \item اسیں !!{derive} کر سکدے آں:
    \begin{derivation}
      1. & $!A$ & \Hyp\\
      2. & $!A \lif !B$ & \Hyp\\
      3. & $!B$ & 1, 2, \MP
    \end{derivation}
    
  \item پہلا نتیجہ \olref[prp]{ax:lnot2} توں، تے دوجا
    \olref[prp]{ax:lif1} تے استنباط دے قضیے توں، بالترتیب۔
  \end{enumerate}
\end{proof}
\endgroup
% /OLP-0122

% OLP-0123 content/first-order-logic/axiomatic-deduction/provability-quantifiers.tex
\begingroup
\def\olfilename{provability-quantifiers}
\expandafter\def\csname ol@sectioncs\endcsname{section}

\olfileid{fol}{axd}{qpr}

\olsection{\usetoken{S}{derivability} تے مسوّرے}

\begin{explain}
  تمامیت دے قضیے لئی ایہ وی ضروری اے کہ مسلّمی اخذاں توں
  $\Proves$ بارے اوہ حقیقتاں حاصل ہون جیہڑیاں ایس حصے وچ
  ثابت کیتیاں جاندیاں نیں۔
\end{explain}

\begin{thm}
\ollabel{thm:strong-generalization} جے $c$ اک !!a{constant} اے جیہڑا
نہ $\Gamma$ وچ آؤندا اے تے نہ $!A(x)$ وچ، تے $\Gamma \Proves !A(c)$،
تاں $\Gamma \Proves \lforall[x][!A(x)]$۔
\end{thm}

\begin{proof}
% PNBAXD-015: ہمیشہ صادق والی مسلّمی سطر تے اثباتِ مقدم دا آخری جواز واضح کیتا۔
استنباط دے قضیے نال $\Gamma \Proves \ltrue \lif !A(c)$۔ چونکہ
$c$ نہ $\Gamma$ وچ آؤندا اے تے نہ~$\top$ وچ، ساڈے کول $\Gamma \Proves
\ltrue \lif \lforall[x][!A(x)]$ ملدا اے۔ استنباط دا قضیہ فیر لا کے،
تے ہمیشہ صادق والی مسلّمی سطر نوں قاعدۂ اثباتِ مقدم نال ورت کے،
$\Gamma \Proves \lforall[x][!A(x)]$۔
\end{proof}

\begin{prop}
\ollabel{prop:provability-quantifiers}
\begin{tagenumerate}{prvEx,prvAll}
\tagitem{prvEx}{$!A(t) \Proves \lexists[x][!A(x)]$.}{}

\tagitem{prvAll}{$\lforall[x][!A(x)] \Proves !A(t)$.}{}
\end{tagenumerate}
\end{prop}

\begin{proof}
% PNBAXD-014: q1/q2 دی اصل بند ٹرم والی حد صاف کیتی۔
دونوں حالتاں وچ مثال لئی ورتیا ٹرم بند اے۔
\begin{tagenumerate}{prvEx,prvAll}
\tagitem{prvEx}{\olref[qua]{ax:q2} تے استنباط دے قضیے توں۔}{}
\tagitem{prvAll}{\olref[qua]{ax:q1} تے استنباط دے قضیے توں۔}{}
\end{tagenumerate}
\end{proof}
\endgroup
% /OLP-0123

% OLP-0124 content/first-order-logic/axiomatic-deduction/soundness.tex
\begingroup
\def\olfilename{soundness}
\expandafter\def\csname ol@sectioncs\endcsname{section}

\iftag{FOL}
      {\olfileid{fol}{axd}{sou}}
      {\olfileid{pl}{axd}{sou}}
      
\olsection{صحت مندی}

\begin{explain}
!!^a{derivation} نظام، جیویں مسلّمی اخذ، \emph{صحت مند} اے جے ایہ
اوہ دعوے !!{derive} نہیں کر سکدا جیہڑے حقیقت وچ درست نہیں۔ ایس لئی
صحت مندی !!{derivation} نظاماں لئی اک قسم دی یقینی تحفظ والی خاصیت اے۔
جیہڑی ثبوتی-نظریاتی خاصیت دا سوال ہووے، اوس مطابق اسیں مثال دے طور تے
ایہ جاننا چاہواں گے کہ
\begin{enumerate}
\item ہر !!{derivable}~$!A$ منطقی طور تے معتبر ہووے؛
\item جے $!A$ کجھ ہور مفروضیاں~$\Gamma$ توں !!{derivable} ہووے،
  تاں اوہ اوہناں دا منطقی نتیجہ وی ہووے؛
\item جے !!{formula}s دا سیٹ~$\Gamma$ متضاد ہووے، تاں اوہ پورا ہون جوگ نہ ہووے۔
\end{enumerate}
ایہ !!a{derivation} نظام دیاں اہم خاصیتاں نیں۔ جے انہاں وچوں کوئی
وی خاصیت نہ ہووے، تاں !!{derivation} نظام وچ کمی اے---ایہ حد توں ودھ
!!{derive} کرے گا۔ ایس لئی کسے !!a{derivation} نظام دی صحت مندی
ثابت کرنا نہایت اہم اے۔
\end{explain}

\begin{prop}
  جے $!A$ اک مسلّمہ اے، تاں
  \iftag{FOL}
    {$\Sat{M}{!A}[s]$ ہر !!{structure}~$\Struct{M}$ تے ہر تعیین~$s$ لئی۔}
    {$\pSat{v}{!A}$ ہر !!{valuation}~$\pAssign{v}$ لئی۔}
\end{prop}

\begin{proof}
\iftag{FOL}{اسیں ایہ جانچنا اے کہ سارے مسلّمے منطقی طور تے معتبر نیں۔
  مثال لئی \olref[qua]{ax:q1} دی حالت ویکھو: فرض کرو $t$، $x$ دی تھاں
  $!A$ وچ !!{free for} اے، تے فرض کرو
  $\Sat{M}{\lforall[x][!A]}[s]$۔ تاں پورا اترن دی تعریف مطابق
  ہر $\varAssign{s'}{s}{x}$ لئی $\Sat{M}{!A}[s']$ وی اے؛ خاص طور تے
  جے $s'(x) = \Value{t}{M}[s]$ ہووے۔
  \olref[syn][ext]{prop:ext-formulas} توں
  $\Sat{M}{\Subst{!A}{t}{x}}[s]$۔ ایس توں
  $\Sat{M}{(\lforall[x][!A] \lif \Subst{!A}{t}{x})}[s]$
  ثابت ہوندا اے۔}{ہر مسلّمے لئی صداقتی جدول بنا کے جانچو کہ اوہ
  ہمیشہ صادق اے۔}
\end{proof}

\begin{thm}[صحت مندی]
\ollabel{thm:soundness}
جے $\Gamma \Proves !A$ تاں $\Gamma \Entails !A$۔
\end{thm}

\begin{proof}
$!A$ دے~$\Gamma$ توں !!{derivation} دی لمبائی اُتے استقرا کردے آں۔
جے استنباطی قاعدیاں نال درست ٹھہرے کوئی قدم نہ ہون، تاں !!{derivation}
دے سارے !!{formula}s یا مسلّماں دیاں مثالاں نیں یا~$\Gamma$ وچ نیں۔
پچھلے قضیے توں سارے مسلّمے
\iftag{FOL}{منطقی طور تے معتبر}{ہمیشہ صادق} نیں، ایس لئی جے $!A$
مسلّمہ اے، تاں $\Gamma \Entails !A$۔ جے $!A \in \Gamma$، تاں
صاف ظاہر اے کہ $\Gamma \Entails !A$۔

جے~$!A$ دے !!{derivation} دا آخری قدم قاعدۂ اثباتِ مقدم نال درست
ٹھہردا اے، تاں !!{derivation} وچ !!{formula}s $!B$ تے $!B \lif !A$
موجود نیں، تے استقرائی مفروضہ انہاں !!{formula}s اُتے ختم ہون والے
!!{derivation} دے حصے اُتے لاگو ہوندا اے (کیوں جے اوہناں وچ
استنباطی قاعدے نال درست ٹھہرے گھٹ توں گھٹ اک قدم کم اے)۔ ایس لئی
استقرائی مفروضے توں $\Gamma \Entails !B$ تے $\Gamma \Entails !B \lif
!A$۔ فیر
\iftag{FOL}
      {\olref[syn][sem]{thm:sem-deduction}}
      {\olref[pl][syn][sem]{thm:sem-deduction}}
توں $\Gamma \Entails !A$۔

\iftag{FOL}{ہُن فرض کرو آخری قدم \QR نال درست ٹھہردا اے۔
% PNBAXD-008: دو مسوّری صورتاں وچ ماخذ توں رہ گیا ! فارمولا ٹوکن بحال کیتا۔
تاں ایس قدم دا روپ $!C \lif \lforall[x][!B(x)]$ اے تے ایس توں
پہلاں $!C \lif !B(c)$ دا قدم موجود اے؛ ایتھے $c$ نہ $\Gamma$ وچ،
نہ~$!C$ وچ، تے نہ $\lforall[x][!B(x)]$ وچ آؤندا اے۔ استقرائی مفروضے
توں $\Gamma \Entails !C \lif !B(c)$۔
\olref[syn][sem]{thm:sem-deduction} توں $\Gamma \cup \{!C\} \Entails
!B(c)$۔

ایسی ساخت~$\Struct{M}$ لو جیہڑی $\Sat{M}{\Gamma \cup
  \{!C\}}$ پوری کردی ہووے۔ اسیں دکھانا اے کہ
$\Sat{M}{\lforall[x][!B(x)]}$۔ چونکہ
$\lforall[x][!B(x)]$ اک !!a{sentence} اے، ایس دا مطلب ایہ دکھانا اے
کہ ہر متغیر تعیین~$s$ لئی $\Sat{M}{!B(x)}[s]$
(\olref[syn][ass]{prop:sat-quant})۔ چونکہ $\Gamma \cup \{!C\}$
سارے جملیاں اُتے مشتمل اے، $\Sat{M}{!D}[s]$ ہر $!D \in
\Gamma$ لئی اے،
\olref[syn][sat]{defn:satisfaction} مطابق۔ ہُن $\Struct{M'}$ نوں
$\Struct{M}$ ورگی ساخت لو، صرف ایہ فرق ہووے کہ
$\Assign{c}{M'} = s(x)$۔ چونکہ $c$ نہ~$\Gamma$ وچ آؤندا اے تے
نہ~$!C$ وچ، \olref[syn][ext]{cor:extensionality-sent} توں
$\Sat{M'}{\Gamma \cup \{!C\}}$۔ چونکہ $\Gamma \cup \{!C\}
\Entails !B(c)$،
% PNBAXD-009: ماخذ دے اک satisfaction argument وچ رہ گیا ! ٹوکن بحال کیتا۔
$\Sat{M'}{!B(c)}$۔ چونکہ $!B(c)$ اک !!a{sentence} اے،
\olref[syn][ass]{prop:sentence-sat-true} توں
$\Sat{M'}{!B(c)}[s]$۔ \olref[syn][ext]{prop:ext-formulas} توں
$\Sat{M'}{!B(x)}[s]$ اُدوں تے صرف اُدوں جدوں
$\Sat{M'}{!B(c)}[s]$ (یاد رکھو کہ
$!B(c)$ دراصل $\Subst{!B(x)}{c}{x}$ اے)۔ ایس لئی
$\Sat{M'}{!B(x)}[s]$۔ چونکہ $c$،~$!B(x)$ وچ نہیں آؤندا،
\olref[syn][ext]{prop:extensionality} توں
$\Sat{M}{!B(x)}[s]$۔ پر $s$ کوئی وی متغیر تعیین سی، ایس لئی
$\Sat{M}{\lforall[x][!B(x)]}$۔ لہٰذا
$\Gamma \cup \{!C\} \Entails \lforall[x][!B(x)]$۔
\olref[syn][sem]{thm:sem-deduction} توں $\Gamma
\Entails !C \lif \lforall[x][!B(x)]$۔

اوہ حالت جتھے $!A$، \QR{} نال درست ٹھہردا اے پر اوہدا روپ
$\lexists[x][!B(x)] \lif !C$ اے، مشق لئی چھڈدے آں۔}{}
\end{proof}

\tagprob{FOL}
\begin{prob}
  \olref[fol][axd][sou]{thm:soundness} دا ثبوت پورا کرو۔
\end{prob}
\tagendprob

\begin{cor}
\ollabel{cor:weak-soundness}
جے $\Proves !A$، تاں $!A$ \iftag{FOL}{منطقی طور تے معتبر}{ہمیشہ صادق} اے۔
\end{cor}

\begin{cor}
\ollabel{cor:consistency-soundness}
جے $\Gamma$ پورا ہون جوگ اے، تاں اوہ بے تضاد اے۔
\end{cor}

\begin{proof}
اسیں متقابلِ عکس ثابت کردے آں۔ فرض کرو $\Gamma$ بے تضاد نہیں۔
تاں $\Gamma \Proves \lfalse$، یعنی~$\lfalse$ دا
~$\Gamma$ توں !!a{derivation} موجود اے۔ \olref{thm:soundness} توں
ہر \iftag{FOL}{!!{structure}~$\Struct{M}$}{!!{valuation}~$\pAssign{v}$}
جیہڑی $\Gamma$ نوں پورا کردی اے، اوہنوں~$\lfalse$ وی پورا کرنا پوے گا۔
چونکہ
\iftag{FOL}{$\Sat/{M}{\lfalse}$ ہر
  !!{structure}~$\Struct{M}$ لئی}{$\pSat/{v}{\lfalse}$ ہر
  !!{valuation}~$\pAssign{v}$ لئی}، کوئی وی
\iftag{FOL}{$\Struct{M}$}{$\pAssign{v}$}، $\Gamma$ نوں پورا
نہیں کر سکدا؛ یعنی $\Gamma$ پورا ہون جوگ نہیں۔
\end{proof}
\endgroup
% /OLP-0124

% OLP-0125 content/first-order-logic/axiomatic-deduction/identity.tex
\begingroup
\def\olfilename{identity}
\expandafter\def\csname ol@sectioncs\endcsname{section}

\olfileid{fol}{axd}{ide}

\olsection{\usetoken{P}{derivation} نال \usetoken{S}{identity}}

!!{derivation}s وچ $\eq$ نوں شامل کرن لئی اسیں بس نویں مسلّمے دے
خاکے شامل کردے آں۔ !!{derivation} تے $\Proves$ دیاں تعریفاں
اوہو رہندیاں نیں؛ ہُن نویں مسلّماں نوں وی اجازت اے۔

\begin{defn}[!!{identity} دے مسلّمے]
\begin{align}
\ollabel{ax:id1} & \eq[t][t], \\
\ollabel{ax:id2} & \eq[t_1][t_2] \lif (!B(t_1) \lif
  \formula B(t_2)),
\end{align}
% PNBAXD-010: ماخذ دے ’’بند ٹرم‘‘ نوں اگلیاں بے قید قضیاں نال ملا کے ایتھے عام ٹرم کیتا۔
کسے وی ٹرماں $t$, $t_1$, $t_2$ لئی۔
\end{defn}

\begin{prop}\ollabel{prop:sound}
  مسلّمے \olref{ax:id1} تے \olref{ax:id2} منطقی طور تے معتبر نیں۔
\end{prop}

\begin{proof}
  مشق۔
\end{proof}

\begin{prob}
\olref[fol][axd][ide]{prop:sound} ثابت کرو۔
\end{prob}

\begin{prop}\ollabel{prop:iden1}
 $\Gamma \Proves \eq[t][t]$، کسے وی ٹرم $t$ تے سیٹ~$\Gamma$ لئی۔
\end{prop}

\begin{prop}\ollabel{prop:iden2}
  جے $\Gamma \Proves !A(t_1)$ تے $\Gamma \Proves
  \eq[t_1][t_2]$، تاں $\Gamma \Proves !A(t_2)$۔
\end{prop}

\begin{proof}
!!{formula}
\[
(\eq[t_1][t_2] \lif (!A(t_1) \lif !A(t_2)))
\]
~\olref{ax:id2} دی اک مثال اے۔ نتیجہ \MP نوں دو واری
لاون نال نکلدا اے۔
\end{proof}
\endgroup
% /OLP-0125

% OLP-0643 content/first-order-logic/axiomatic-deduction/provability.tex
\begingroup
\def\olfilename{provability}
\expandafter\def\csname ol@sectioncs\endcsname{section}

\iftag{FOL}
      {\olfileid{fol}{axd}{prv-alt}}
      {\olfileid{pl}{axd}{prv-alt}}

\olsection{\usetoken{S}{derivability} دیاں خاصیتاں}

\begin{prop}[وڈے سیٹ ول برقرار رہنا] جے $\Gamma \subseteq \Delta$ تے $\Gamma \Proves !A$، تاں $\Delta \Proves !A$ وی اے۔ \end{prop}

\begin{proof} ہر محدود $\Gamma_0 \subseteq \Gamma$،~$\Delta$ دا وی محدود ذیلی سیٹ اے۔ ایس لئی $\Gamma_0 \Proves !A$ دا !!a{derivation} ایہ وی وکھاؤندا اے کہ
$\Delta \Proves !A$۔ \end{proof}

\begin{prop}\ollabel{prop:provability} $\Gamma \Proves !A$ اُدوں تے صرف اُدوں جدوں $\Gamma\cup \{\lnot!A \}$ متضاد اے۔ \end{prop}

\begin{prop} \begin{enumerate}
\item \ollabel{prop:provability-contr} جے $\Gamma \Proves !A$ تے $\Gamma
\cup \{ !A\} \Proves \lfalse$، تاں $\Gamma$ متضاد اے۔

\item \ollabel{prop:provability-lnot} جے $\Gamma \cup \{!A\}
\Proves \lfalse$، تاں $\Gamma \Proves \lnot !A$۔

\item \ollabel{prop:provability-exhaustive} جے $\Gamma \cup \{!A\}
\Proves \lfalse$ تے $\Gamma \cup \{\lnot !A\} \Proves
\lfalse$، تاں $\Gamma \Proves \lfalse$۔

\tagitem{prvOr}{\ollabel{prop:provability-lor-left} جے $\Gamma \cup \{!A\}
\Proves \lfalse$ تے $\Gamma \cup \{!B\} \Proves
\lfalse$، تاں $\Gamma \cup \{!A \lor !B\} \Proves \lfalse$۔}{}

\tagitem{prvOr}{\ollabel{prop:provability-lor-right} جے $\Gamma
\Proves !A$ یا $\Gamma \Proves !B$، تاں $\Gamma
\Proves !A \lor !B$۔}{}

\tagitem{prvAnd}{\ollabel{prop:provability-land-left} جے $\Gamma
\Proves !A \land !B$، تاں $\Gamma \Proves !A$ تے
$\Gamma \Proves !B$۔}{}

\tagitem{prvAnd}{\ollabel{prop:provability-land-right} جے $\Gamma
\Proves !A$ تے $\Gamma \Proves !B$، تاں $\Gamma
\Proves !A \land !B$۔}{}

\tagitem{prvIf}{\ollabel{prop:provability-mp} جے $\Gamma \Proves
!A$ تے $\Gamma \Proves !A \lif !B$، تاں $\Gamma
\Proves !B$۔}{}

\tagitem{prvIf}{\ollabel{prop:provability-lif} جے $\Gamma \Proves
\lnot !A$ یا $\Gamma \Proves !B$، تاں $\Gamma \Proves
!A \lif !B$۔}{} \end{enumerate} \end{prop}

\begin{proof}
\begin{enumerate}
\item منّ لو $\Gamma_0 \cup \{!A\} \Proves \lfalse$۔

\begin{derivation} 
1. & $\Gamma_0 \cup \{!A\} \Proves \lfalse$ & فرض \\
2. & $\Gamma_1 \Proves !A$ & فرض\\ 
3. & $\Gamma_0 \cup \Gamma_1 \Proves \lfalse$ & قطع \\ 
\end{derivation}

چونکہ $\Gamma_0 \subseteq \Gamma$ تے $\Gamma_1 \subseteq \Gamma$،
ایس لئی $\Gamma_0 \cup \Gamma_1 \subseteq \Gamma$، سو $\Gamma \Proves \lfalse$۔

\item اسیں منّنے آں کہ $\Gamma_0 \cup\{\lnot !A\} \Proves \lfalse$۔

\begin{derivation} 
1. & $\Gamma_0 \cup\{\lnot !A\} \Proves \lfalse$ & فرض \\ 
2. & $\Gamma_0 \Proves \lnot !A \lif \lfalse$ & اخذ دا قضیہ\\
3. & $\Gamma_0 \Proves (\lnot !A \lif \lfalse) \lif !A$ & Ax0\\ 
4. & $\Gamma_0 \Proves !A$ & MP 2, 3 \\ 
\end{derivation}

\item اسیں منّنے آں کہ $\Gamma \cup \{ !A \} \Proves \lfalse$ تے
$\Gamma_1 \cup \{\lnot !A \} \Proves \lfalse$۔

\begin{derivation} 
1. & $\Gamma_0 \cup\{!A\} \Proves \lfalse$ & فرض \\
2. & $\Gamma_0 \Proves !A \lif \lfalse$ & اخذ دا قضیہ \\ 
3. & $\Gamma_1 \cup \{\lnot !A \} \Proves \lfalse$ & فرض\\ 
4. & $\Gamma_0 \Proves (!A \lif \lfalse) \lif \lnot !A$ & Ax0\\ 
5. & $\Gamma_0 \Proves \lnot !A$ & MP 2, 4\\ 
6. & $\Gamma_0 \cup \Gamma_1 \Proves \lfalse$ & قطع 3, 5 \\
\end{derivation}

چونکہ $\Gamma_0 \subseteq \Gamma$ تے $\Gamma_1 \subseteq \Gamma$،
ایس لئی $\Gamma_0 \cup \Gamma_1 \subseteq \Gamma$۔ سو $\Gamma \Proves \lfalse$۔

% prop:provability-lor-left 
\tagitem{defOr}{}{
\iftag{probOr}{مشق۔}{مشق۔}}

% prop:provability-lor-right 
\tagitem{defOr}{}{ \iftag{probOr}{مشق۔}{
اسیں منّنے آں کہ $\Gamma_0 \Proves !A$۔

\begin{derivation} 
1. & $\Gamma_0 \Proves !A$ & فرض \\ 
2. & $\Gamma_0 \Proves !A \lif (!A \lor !B)$ & Ax0\\ 
3. & $\Gamma_0 \Proves !A \lor !B$ & MP 1, 2 \\
\end{derivation}

ایس لئی $\Gamma \Proves !A \lor !B$۔ جدوں $\Gamma \Proves
!B$ ہووے، ثبوت ایہو جیہا اے۔}}

% prop:provability-land-left 
\tagitem{defAnd}{}{ \iftag{probAnd}{مشق}{
منّ لو $\Gamma_0 \Proves !A \land !B$

\begin{derivation} 
1. & $\Gamma_0 \Proves !A \land !B$ & فرض \\ 
2. & $\Gamma_0 \Proves (!A\land !B) \lif !A $ & Ax0\\ 
3. & $\Gamma_0 \Proves !A \lor !B$ & MP 1, 2 \\
\end{derivation}

سو $\Gamma \Proves !A$۔ ایہو جیہا !!{derivation} شروع کر کے وکھا سکدے آں کہ
$\Gamma \Proves !B$۔}}

% prop:provability-land-right
\tagitem{defAnd}{}{\iftag{probAnd}{مشق۔}{مشق۔}}

% prop:provability-mp 
\tagitem{defIf}{}{ \iftag{probIf}{مشق۔}{
اسیں دوویں منّنے آں، $\Gamma_0 \Proves !A$ تے $\Gamma_0 \Proves !A \lif !B $۔

\begin{derivation} 
1. & $\Gamma_0 \Proves !A$ & فرض \\ 
2. & $\Gamma_0 \Proves !A \lif !B $ & فرض\\ 
3. & $\Gamma_0 \Proves !B$ & MP 1, 2\\
\end{derivation}

چونکہ $\Gamma_0 \cup \Gamma_1 \subseteq \Gamma$، ایس دا مطلب اے کہ $\Gamma
\Proves !B$۔}}

% prop:provability-lif 
\tagitem{defIf}{}{\iftag{probIf}{مشق۔}{
پہلاں منّ لو $\Gamma \Proves \lnot !A$۔

\begin{derivation} 
1. & $\Gamma_0 \Proves \lnot !A$ & فرض \\ 
2. & $\Gamma_0 \Proves \lnot !A \lif (!A \lif !B) $ & Ax0\\ 
3. & $\Gamma_0 \Proves !A \lif !B$ & MP 1,
2\\ \end{derivation}

$\Gamma \Proves !B$ والا ثبوت وی ایہو جیہا اے:

\begin{derivation} 
1. & $\Gamma_0 \Proves !B$ & فرض \\ 
2. & $\Gamma_0 \Proves !B \lif (!A \lif !B) $ & Ax0\\ 
3. & $\Gamma_0 \Proves !A \lif !B$ & MP 1, 2\\
\end{derivation}}}

\end{enumerate} 
\end{proof}

\begin{probtag}{probOr,probAnd,probIf} \olref{prop:provability} دا ثبوت پورا کرو۔
\end{probtag}

\begin{thm}[عام کرن دا قضیہ] \ollabel{thm:Generalization} جے $\Gamma \Proves
!A$ تے $x$، $\Gamma$ دے کسے وی !!{formula} وچ آزاد نہیں، تاں
$\Gamma \Proves \lforall[x][!A]$۔
\end{thm}

\begin{proof} $\mathsf{Thm}_\Gamma$ اُتے استقرا نال۔ جے $!A$ مسلّمہ اے، تاں $\lforall[x][!A]$ وی مسلّمہ اے۔ جے $!A\in \Gamma$، تاں $x$، $!A$ وچ آزاد نہیں۔ سو $!A \lif \lforall[x][!A]$ مسلّمہ اے (\textbf{Ax3}) تے MP نال $\Gamma
\Proves \lforall[x][!A]$۔ ہُن منّ لو $!A$، MP نال نکلدا اے کیوں جے
$\Gamma \Proves !B$ تے $\Gamma \Proves !B \lif !A$۔ استقرا دی فرضیت نال $\Gamma \Proves \lforall[x][!B]$ تے $\Gamma \Proves
\lforall[x][( !B \lif !A)]$۔ پر \textbf{Ax2} نال ایہ وی اے کہ \[ \Gamma \Proves
\lforall[x][( !B \lif !A)] \lif (\lforall[x][ !B] \lif \lforall[x][!A]), \]
تے MP دو واری لا کے لوڑ مطابق $\Gamma \Proves \lforall[x][!A]$ ملدا اے۔ \end{proof}

\begin{thm}\ollabel{thm:WeakGen} (\emph{ثابتاں اُتے کمزور عام کرن دا قضیہ})
جے $\Gamma \Proves !A$ تے $c$ اک !!a{constant} اے جیہڑا $\Gamma$ وچ نہیں آؤندا، تاں اک متغیر $x$ اے جیہڑا $!A$ وچ نہیں آؤندا، ایویں کہ $\Gamma
\Proves \lforall[x][\Subst{!A}{x}{c}]$، تے ثبوت وچ $c$ نہیں آؤندا۔
\end{thm}

\begin{proof} منّ لو $!A_1,\ldots,!A_n$، $!A$ دا $\Gamma$ توں ثبوت اے، سو $!A=!A_n$۔ کوئی متغیر $x$ چُنو جیہڑا $!A_1,\ldots,!A_n$ وچ نہیں آؤندا، تے نویں لڑی ویکھو: $\Subst{!A_1}{x}{c},\ldots,\Subst{!A_n}{x}{c}.$ ایہو جیہی لڑی $\Subst{!A}{x}{c}$ دا $\Gamma$ توں ثبوت اے۔ اصل وچ، ہر $i=1,\ldots,n$ لئی:
\begin{itemize} \item جے $!A$ مسلّمہ اے، تاں $\Subst{!A_i}{x}{c}$ وی مسلّمہ اے۔
\item جے $!A_i \in \Gamma$، تاں چونکہ $c$، $\Gamma$ وچ نہیں، ایس لئی $\Subst{!A_i}{x}{c} = !A_i \in \Gamma$۔
\item جے $!A_i$، $!A_j$ تے $!A_j \lif !A_i$ توں ملدا اے، تاں $\Subst{!A_i}{x}{c}$، $\Subst{!A_j}{x}{c}$ تے $\Subst{(!A_j \lif !A_i)}{x}{c}
= \Subst{!A_j}{x}{c} \lif \Subst{!A_i}{x}{c}$ توں MP نال نکلدا اے۔
\end{itemize} صاف اے کہ !!{constant} $c$ ہُن نویں لڑی وچ نہیں آؤندا۔ ہُن $\Gamma'$ وچ $\Gamma$ دے اوہ !!{formula}s رکھو جیہڑے $\Subst{!A}{x}{c}$ دے !!{derivation} وچ آؤندے نیں۔ فیر $x$، $\Gamma'$ دے کسے !!{formula} وچ آزاد نہیں تے $\Gamma' \Proves \Subst{!A}{x}{c}$۔
عام کرن دے قضیے (\olref{thm:Generalization}) نال $\Gamma'
\Proves \lforall[x][\Subst{!A}{x}{c}]$۔ وڈے سیٹ ول برقرار رہن نال لوڑ مطابق $\Gamma \Proves
\lforall[x][\Subst{!A}{x}{c}]$۔
\end{proof}

\begin{explain}
ساڈا مقصد \olref{thm:WeakGen} دی شرط، کہ $x$، $!A$ وچ (بالکل) نہیں آؤندا، نوں ایہناں گھٹ سخت شرطاں نال بدلنا اے کہ $x$، $!A$ وچ آزاد نہیں تے $c$ دی تھاں پاؤن لئی $!A$ وچ !!{free for} اے۔ صاف اے کہ بندھے !!{variable} نوں بدل کے ایہ ہو سکدا اے۔ سو پہلاں اسیں ایہو ثابت کرنا اے۔
\end{explain}

\begin{lem}
\ollabel{lem:free-for}
جے $x$، $c$ دی تھاں پاؤن لئی $!A$ وچ آزاد اے تے $y$، $!A$ وچ آزاد نہیں، تاں $x$، $y$ دی تھاں پاؤن لئی $\Subst{!A}{y}{c}$ وچ !!{free for} اے۔
\end{lem}

\begin{proof} جے $x$، $y$ دی تھاں پاؤن لئی $\Subst{!A}{y}{x}$ وچ !!{free for} نہیں، تاں $y$ دی کوئی آزاد موجودگی $\Subst{!A}{y}{c}$ وچ مسوّر $\lforall[x]$ دی حد وچ آؤندی اے۔ پر فرض مطابق $y$، $!A$ وچ آزاد نہیں، سو ایہ ساری موجودگیاں بدل $\Subst{}{y}{c}$ توں آؤندیاں نیں۔ ایس لئی $c$ دی کوئی موجودگی $\lforall[x]$ دی حد وچ آؤندی اے تے $x$، $c$ دی تھاں پاؤن لئی $!A$ وچ !!{free for} نہیں۔
\end{proof}

\begin{lem}[بندھا متغیر بدلنا]
\ollabel{lem:ChangeBdVar}
جے $x$ تے $y$، $!A$ وچ آزاد نہیں تے دوویں $c$ دی تھاں پاؤن لئی $!A$ وچ !!{free for} نیں، تاں $\Proves \lforall[x][\Subst{!A}{x}{c}] \equiv
\lforall[y][\Subst{!A}{y}{c}]$ (تے ثبوت وچ $c$ نہیں آؤندا)۔
\end{lem}

\begin{proof}
شرطاں توں $y$، $c$ دی تھاں پاؤن لئی $!A$ وچ !!{free for} اے تے $x$، $!A$ وچ آزاد نہیں۔ سو پچھلے \olref{lem:free-for} نال $y$، $x$ دی تھاں پاؤن لئی $\Subst{!A}{x}{c}$ وچ !!{free for} اے۔ ایس توں نکلدا اے کہ
$\lforall[x][\Subst{!A}{x}{c} \lif \Subst{!A}{x}{c} \Subst{}{y}{x}]$
اک مسلّمہ اے (\textbf{Ax1})۔ چونکہ $x$، $!A$ وچ آزاد نہیں، سانوں
$\Subst{!A}{x}{c} \Subst{}{y}{x} = \Subst{!A}{y}{c}$ ملدا اے، ایس لئی
\[
\Proves \lforall[x][\Subst{!A}{x}{c}] \lif \Subst{!A}{y}{c},
\]
تے اخذ دے قضیے نال $\lforall[x][\Subst{!A}{x}{c}] \Proves
\Subst{!A}{y}{c}$۔ چونکہ $y$، $!A$ وچ آزاد نہیں، اوہ
$\lforall[x][\Subst{!A}{x}{c}]$ وچ وی آزاد نہیں۔ سو عام کرن دے قضیے نال
$\lforall[x][\Subst{!A}{x}{c}] \Proves \lforall[y][\Subst{!A}{y}{c}]$،
تے اخذ دا قضیہ فیر لا کے $\Proves
\lforall[x][\Subst{!A}{x}{c}] \lif \lforall[y][\Subst{!A}{y}{c}]$ ملدا اے۔
$\Proves \lforall[y][\Subst{!A}{y}{c}] \lif \lforall[x]
[\Subst{!A}{x}{c}]$ دا ثبوت عین اُلٹے پاسے ایویں ای اے، سو نتیجہ Taut نال نکلدا اے۔
\end{proof}

\begin{thm}[ثابتاں اُتے تکڑا عام کرن دا قضیہ]
\ollabel{thm:StrongGen}
جے $\Gamma \Proves !A$، !!{constant} $c$، $\Gamma$ وچ نہیں آؤندا، تے $x$، $!A$ وچ آزاد نہیں پر $c$ دی تھاں پاؤن لئی $!A$ وچ !!{free for} اے، تاں $\Gamma \Proves \lforall[x][\Subst{!A}{x}{c}]$۔
\end{thm}

\begin{proof}
چونکہ $\Gamma \Proves !A$ تے $c$، $\Gamma$ وچ نہیں، کمزور عام کرن دے قضیے نال اک !!a{variable} $y$ اے جیہڑا $!A$ وچ نہیں آؤندا تے
$\Gamma \Proves \lforall[y][\Subst{!A}{y}{c}]$۔ چونکہ $y$، $!A$ وچ (بالکل) نہیں، ایہ $!A$ وچ آزاد نہیں تے $c$ دی تھاں پاؤن لئی $!A$ وچ !!{free for} وی اے۔ جے نالے (فرض مطابق) $x$، $!A$ وچ آزاد نہیں تے $c$ دی تھاں پاؤن لئی $!A$ وچ !!{free for} اے، تاں بندھا !!{variable} بدلّن دیاں شرطاں پوریاں نیں۔ سو
$\Proves \lforall[x][\Subst{!A}{x}{c}] \equiv
\lforall[y][\Subst{!A}{y}{c}]$، ایس توں $\Gamma \Proves
\lforall[x][\Subst{!A}{x}{c}]$۔
\end{proof}

\begin{thm}
\ollabel{thm:provability-quantifiers}
\begin{tagenumerate}{prvEx,prvAll}
\tagitem{prvEx}{جے $\Gamma \Proves !A(t)$، تاں $\Gamma \Proves
  \lexists[x][!A(x)]$۔}{}
\tagitem{prvAll}{جے $\Gamma \Proves \lforall[x][!A(x)]$، تاں $\Gamma
  \Proves !A(t)$۔}{}
\end{tagenumerate}
\end{thm}

\begin{proof}
\begin{tagenumerate}{prvEx,prvAll}
%%Need axioms for exists
\tagitem{prvEx}{مشق۔}{}
\tagitem{prvAll} {منّ لو $\Gamma_0 \Proves \lforall[x][!A(x)]$۔

\begin{derivation} 
1. & $\Gamma_0 \Proves \lforall[x][!A(x)]$ & فرض \\
2. & $\Gamma_0 \Proves \lforall[x][!A(x)] \lif \Subst{!A}{t}{x}$ & Ax1 \\ 
3. & $\Gamma_0 \Proves \Subst{!A}{t}{x}$ & MP 1, 2 \\ 
\end{derivation}}{}

\end{tagenumerate} 
\end{proof}
\endgroup
% /OLP-0643
\OLEndChapterHook
\endgroup
% /OLP-0112

% OLP-0126 content/first-order-logic/completeness/completeness.tex
\begingroup
\def\olfilename{completeness}
\expandafter\def\csname ol@sectioncs\endcsname{section}

\iftag{FOL}
      {\olchapter{fol}{com}{تمامیت دا قضیہ}}
      {\olchapter{pl}{com}{تمامیت دا قضیہ}}

% OLP-0127 content/first-order-logic/completeness/introduction.tex
\begingroup
\def\olfilename{introduction}
\expandafter\def\csname ol@sectioncs\endcsname{section}

\iftag{FOL}
      {\olfileid{fol}{com}{int}}
      {\olfileid{pl}{com}{int}}

\olsection{تعارف}

تمامیت دا قضیہ منطق دے سب توں بنیادی نتائج وچوں اک اے۔ ایس دے
دو بیان نیں، جنہاں دی ہم ارزیت اسیں ثابت کراں گے۔ پہلے بیان
مطابق ایہ معنوی نتیجے تے ساڈے !!{derivation} نظام دے تعلق بارے
بنیادی گل کہندا اے: جے کوئی !!a{sentence}~$!A$ کجھ
!!{sentence}s~$\Gamma$ دا نتیجہ ہووے، تاں اک
!!a{derivation} وی موجود اے جیہڑا $\Gamma
\Proves !A$ قائم کردا اے۔ ایس طرح !!{derivation} نظام
اُنا طاقت ور اے جِنا اوہ انہاں گلاں نوں ثابت کیتے بنا ہو سکدا اے
جیہڑیاں حقیقت وچ نتیجہ نہیں بندیاں۔

دوجے بیان مطابق ایہ نمونے دے وجود دا نتیجہ اے: !!{sentence}s
دا ہر بے تضاد سیٹ پورا ہون جوگ اے۔ بے تضادی اک ثبوتی-نظریاتی
تصور اے: ایہ کہندا اے کہ ساڈا !!{derivation} نظام کسے خاص
طرح دیاں !!{derivation}s نہیں بنا سکدا۔ پر صرف ایس گل توں کہ
~$\Gamma$ توں ایس طرح دیاں !!{derivation}s نہیں نکلدیاں، ایہ
کیویں پکا ہویا کہ
% PNB-COM-012: FOL شاخ وچ وی دتے مفروضیاں نوں پورا کرن والی ساخت دی شرط لازمی اے۔
\iftag{FOL}{!!a{structure}~$\Struct{M}$ موجود ہووے تے دتے سیٹ نوں پورا کرے}{!!{valuation}{~$\pAssign{v}$}
موجود ہووے جیہدی لئی $\iftag{FOL}{\Sat{M}}{\pSat{v}}{\Gamma}$}؟
تمامیت دا قضیہ پہلی واری ثابت ہون توں پہلاں---بلکہ ساڈے ہُن
والے !!{derivation} نظاماں توں وی پہلاں---جرمنی دے عظیم ریاضی
دان ڈیوڈ ہلبرٹ دی رائے سی کہ ریاضیاتی نظرییاں دی بے تضادی
انہاں دے موضوع بنن والیاں چیزاں دے وجود دی ضمانت دیندی اے۔
انہاں نے گوٹلوب فریگے نوں اک خط وچ ایویں لکھیا:
\begin{quote}
  جے من مانی طرح دتے مسلّمے، اپنے سارے نتیجیاں سمیت، اک دوجے
  نال تضاد نہ کردے ہون، تاں اوہ سچے نیں تے مسلّماں نال تعریف
  کیتیاں چیزاں موجود نیں۔ میرے لئی سچائی تے وجود دا معیار ایہی اے۔
\end{quote}
فریگے نے ایس نال سخت اختلاف کیتا۔ تمامیت دے قضیے دا دوجا
بیان دکھاندا اے کہ ہلبرٹ گھٹ توں گھٹ ایس معنی وچ درست سن کہ
جے مسلّمے بے تضاد ہون، تاں \emph{کوئی نہ کوئی}
\iftag{FOL}{!!{structure}}{!!{valuation}} موجود اے جیہڑی
انہاں ساریاں نوں سچا بناندی اے۔

ایہو وجوہاں تمامیت دے قضیے---یا خاص طور تے اوس دے ثبوت---دی
اہمیت دیاں اکلوتیاں وجوہاں نہیں۔ ایس دے کئی اہم نتائج نیں،
جنہاں وچوں کجھ اُتے اسیں الگ گل کراں گے۔ مثال لئی، کوئی وی
!!{derivation} جیہڑا $\Gamma \Proves !A$ وکھاندا اے محدود
ہوندا اے، ایس لئی~$\Gamma$ دے صرف محدود گنتی دے
!!{sentence}s ورت سکدا اے۔ اس توں تمامیت دے قضیے نال ایہ
نکلدا اے کہ جے $!A$،~$\Gamma$ دا نتیجہ اے، تاں اوہ پہلے ای
~$\Gamma$ دے کسے محدود ذیلی سیٹ دا نتیجہ اے۔ ایس نوں
\emph{محدود گواہی (compactness)} آکھدے نیں۔ برابر دا بیان
ایہ اے کہ جے $\Gamma$ دا ہر محدود ذیلی سیٹ بے تضاد اے، تاں
$\Gamma$ وی بے تضاد اے۔

بھاویں محدود گواہی دا قضیہ، تمامیت دے قضیے توں
!!{derivation}s راہیں گھوم کے نکلدا اے، پر تمامیت دے قضیے دے
\emph{ثبوت} نوں ورت کے ایس نوں سِدھا ثابت کرنا وی ممکن اے۔
ثبوت اک ایسی خاصیت---بے تضادی---والا !!{sentence}s دا سیٹ
لیندا اے تے ایس سیٹ توں !!a{structure} بناندا اے جیہدیاں
کجھ خاصیتاں ہوندیاں نیں (ایتھے ایہ کہ اوہ سیٹ نوں پورا کردی اے)۔
تقریباً ایہو تعمیر، بے تضاد سیٹاں دی تھاں ”محدود طور تے پورے
ہون جوگ“ !!{sentence}s دے سیٹاں توں شروع کر کے، محدود گواہی
سِدھی قائم کر سکدی اے۔
\iftag{FOL}{ایس تعمیر دے ہور نتیجے وی نیں؛ مثال لئی،
  !!{sentence}s دے ہر پورے ہون جوگ سیٹ دا اک محدود یا
  !!{denumerable}s نمونہ موجود اے۔ (ایس نتیجے نوں
  L\"owenheim--Skolem دا قضیہ آکھدے نیں۔) عام طور تے
  !!{sentence}s دے سیٹاں توں !!{structure}s بنان دی
  تعمیر منطق وچ اکثر ورتی جاندی اے، تے کدی فلسفے وچ وی۔}{}
\endgroup
% /OLP-0127

% OLP-0128 content/first-order-logic/completeness/outline.tex
\begingroup
\def\olfilename{outline}
\expandafter\def\csname ol@sectioncs\endcsname{section}

\iftag{FOL}
      {\olfileid{fol}{com}{out}}
      {\olfileid{pl}{com}{out}}

\olsection{ثبوت دا خاکہ}

تمامیت دے قضیے دا ثبوت ذرا گنجل دار اے، تے پہلی واری پڑھدیاں
راہ گواچ جانا سوکھا اے۔ ایس لئی آؤ ثبوت دا خاکہ ویکھئیے۔ پہلا
قدم نظر دا رخ بدلنا اے، جیہڑا ثبوت ول جان دا رستہ دکھاندا اے۔
جے تمامیت نوں ”جدوں وی $\Gamma
\Entails !A$ ہووے، اوہدوں $\Gamma \Proves !A$ وی ہووے“ دے روپ
وچ سوچئیے، تاں کوئی خیال آنا وی اوکھا لگ سکدا اے: $\Gamma
\Proves !A$ وکھاون لئی سانوں !!a{derivation} لبھنا پئے گا،
تے مفروضہ $\Gamma \Entails !A$ ایس کم وچ ظاہری طور تے
کوئی مدد نہیں کردا۔ کجھ ثبوتی نظاماں لئی !!a{derivation} سِدھا
بنایا جا سکدا اے، پر اسیں ذرا وکھرا رستہ لواں گے۔ نظر دا
لازمی بدلاؤ ایہ اے کہ تمامیت ایویں وی بیان ہو سکدی اے: ”جے
$\Gamma$ بے تضاد اے، تاں اوہ پورا ہون جوگ اے۔“ شاید اسیں
~$\Gamma$ دی معلومات تے اوہدے بے تضاد ہون دے مفروضے نوں
ورت کے \iftag{FOL}{!!a{structure}}{!!a{valuation}} بنا
سکئیے جیہڑی~$\Gamma$ دے ہر
\iftag{FOL}{!!{sentence}}{!!{formula}} نوں پورا کرے۔ آخر
سانوں پتہ اے کہ سانوں کیہو جہی
\iftag{FOL}{!!{structure}}{!!{valuation}} چاہیدی اے:
جیہڑی~$\Gamma$ دے بیان دے مطابق ہووے!{}

جے $\Gamma$ وچ صرف \iftag{FOL}{جوہری
  !!{sentence}s}{!!{propositional variable}s} ہون، تاں
اوس لئی نمونہ بنانا سوکھا اے۔\iftag{FOL}{ فرض کرو سارے جوہری
  !!{sentence}s دا روپ $\Atom{P}{a_1,\dots,a_n}$ اے، جتھے
  $a_i$ !!{constant}s نیں۔}{} سانوں بس
\iftag{FOL}{%
  !!a{domain}~$\Domain{M}$ تے~$P$ لئی ایسی تعبیر چننی اے کہ
  $\Sat{M}{\Atom{P}{a_1,\ldots,a_n}}$ ہووے۔ ایہ اوکھا نہیں:
  $\Domain{M} = \Nat$ رکھو، $\Assign{\Obj c_i}{M} = i$ رکھو،
  تے ہر $\Atom{P}{a_1,\ldots,a_n} \in \Gamma$ لئی
  ٹپل $\tuple{k_1,
    \dots, k_n}$ نوں $\Assign{P}{M}$ وچ پاؤ، جتھے $k_i$
  مستقل علامت~$a_i$ دا اشاریہ اے (یعنی $a_i \ident \Obj c_{k_i}$)۔}{%
  !!a{valuation}~$\pAssign{v}$ چننی اے جیہدی لئی $\pSat{v}{p}$
  ہر $p \in \Gamma$ لئی ہووے۔ ایس لئی
  $\pAssign{v}(p) = \True$ اُدوں تے صرف اُدوں رکھو جدوں $p \in \Gamma$۔}

ہُن فرض کرو $\Gamma$ وچ کوئی !!{formula}~$\lnot !B$ اے، جتھے
$!B$ جوہری اے۔ سانوں فکر ہو سکدی اے کہ
$\iftag{FOL}{\Struct{M}}{\pAssign{v}}$ دی تعمیر، $\lnot !B$
نوں سچا بنان دے امکان نال ٹکراوے گی۔ پر ایتھے~$\Gamma$ دی
بے تضادی کم آؤندی اے: جے $\lnot !B \in \Gamma$، تاں
$!B \notin \Gamma$؛ نہیں تاں $\Gamma$ متضاد ہووے گا۔
تے جے $!B \notin \Gamma$، تاں ساڈی تعمیر مطابق
~$\iftag{FOL}{\Struct{M}}{\pAssign{v}}$ وچ
$\iftag{FOL}{\Sat/{M}}{\pSat/{v}}{!B}$، ایس لئی
$\iftag{FOL}{\Sat{M}}{\pSat{v}}{\lnot !B}$۔
ایتھوں تک گل ٹھیک اے۔

پر جے $\Gamma$ وچ پیچیدہ، غیر جوہری فارمولے ہون تاں؟
مثال لئی جے $!A \land !B$ اوس وچ ہووے، تاں اوہنوں سچا
بنان لئی اسیں ایویں چلنا اے جیویں $!A$ تے $!B$ دونیں
~$\Gamma$ وچ ہون۔ تے جے $!A \lor !B \in \Gamma$، تاں
اوہناں وچوں گھٹ توں گھٹ اک نوں سچا بنانا پئے گا، یعنی
ایویں چلنا پئے گا جیویں اوہناں وچوں اک~$\Gamma$ وچ ہووے۔

ایس توں ایہ خیال نکلدا اے: اسیں~$\Gamma$ وچ ہور
!!{formula}s شامل کریے تاکہ (a)~بنیا سیٹ بے تضاد رہے تے
% PNB-COM-001: مکمل سیٹ دی اگلی عین تعریف مطابق صرف جوہری نہیں، ہر جملہ طے ہونا چاہیدا اے۔
(b)~ہر ممکن !!{sentence}~$!A$ لئی بنے سیٹ وچ یا تاں
$!A$ ہووے یا $\lnot !A$، تے (c)~جدوں وی
$!A \land !B$ سیٹ وچ ہووے، تاں $!A$ تے $!B$ دونیں
وی ہون؛ جے $!A \lor !B$ سیٹ وچ ہووے، تاں $!A$ یا $!B$
وچوں گھٹ توں گھٹ اک وی ہووے، وغیرہ۔ اسیں ایہ کم جاری رکھدے
آں (ممکن اے ہمیشہ لئی)۔ ایس طرح شامل کیتے سارے
!!{formula}s دے سیٹ نوں~$\Gamma^*$ آکھو۔ فیر اُتّے دتی
تعمیر سانوں \iftag{FOL}{!!a{structure}~$\Struct{M}$}{!!a{valuation}~$\pAssign{v}$}
دے گی، جیہدے بارے اسیں استقرا نال ثابت کر سکدے آں کہ اوہ
~$\Gamma^*$ دے سارے جملے پورے کردی اے، تے ایس لئی~$\Gamma$
دے سارے جملے وی، کیوں جے $\Gamma \subseteq \Gamma^*$۔
معلوم ہوندا اے کہ (a) تے~(b) یقینی بنانا کافی اے۔ جملیاں دا
اوہ سیٹ جیہدے لئی (b) پوری ہووے \emph{مکمل} کہلاندا اے۔
ایس لئی ساڈا کم بے تضاد سیٹ~$\Gamma$ نوں بے تضاد تے مکمل
سیٹ~$\Gamma^*$ تک ودھانا ہووے گا۔

\iftag{FOL}{%
ایس منصوبے وچ اک رکاوٹ اے: جے $\lexists[x][!A(x)] \in \Gamma$
ہووے، تاں اسیں چاہواں گے کہ کوئی !!{constant}~$c$ چن کے
ایس عمل وچ $!A(c)$ شامل کریے۔ پر ایہ ہمیشہ کیویں ممکن اے؟
ہو سکدا اے ساڈی زبان وچ صرف تھوڑے !!{constant}s ہون، تے
اوہناں وچوں ہر اک لئی $\lnot !A(c) \in \Gamma$ ہووے۔
اسیں نال $!A(c)$ وی شامل نہیں کر سکدے، کیوں جے فیر
سیٹ متضاد ہو جائے گا تے سانوں پتہ نہیں لگے گا کہ
$\Struct{M}$ نے $!A(c)$ نوں سچا بنانا اے یا
$\lnot !A(c)$ نوں۔ ایہ وی ہو سکدا اے کہ $\Gamma$ وچ
ایسی زبان دے جملے ہون جیہدے وچ کوئی مستقل علامت سرے توں
ہے ای نہیں (مثلاً سیٹ تھیوری دی زبان)۔

ایس مسئلے دا حل ایہ اے کہ شروع وچ لامتناہی گنتی دے
مستقل ودھا لئیے، نال ایسے جملے وی جیہڑے اوہناں نوں
مسوّراں نال درست طرح جوڑن۔ (ظاہر اے اسیں ایہ وی
جانچنا اے کہ ایس نال کوئی تضاد پیدا نہیں ہوندا۔)

ساڈی پہلی تعمیر اوہدوں چنگی چل دی اے جدوں جوہری جملیاں
وچ صرف !!{constant}s ہون۔ پر زبان وچ !!{function}s وی
ہو سکدے نیں۔ فیر~$\Nat$ اُتے ایہناں !!{function}s لئی
مناسب فنکشن لبھنا اوکھا ہو سکدا اے تاکہ ساری گل ٹھیک بنے۔ ایس لئی
اک ہور ترکیب اے: $i$ نوں $\Obj c_i$ دی تعبیر دے طور تے ورتن دی
تھاں !!{constant}s دا سیٹ آپ میدان بنا لو۔ فیر
$\Struct M$ ہر !!{constant} نوں اوہ آپ دے سکدا اے:
$\Assign{\Obj c_i}{M} = \Obj c_i$۔ پر ہور اگے
کیوں نہ جائیے؟
% PNB-COM-002: اگلی رسمی term-model تعریف مطابق میدان تے فنکشن دے اندراج بند ٹرم نیں۔
$\Domain{M}$ نوں زبان دے سارے \emph{بند ٹرماں} دا سیٹ
بنا لو!{} ایس صورت وچ !!{function}s لئی فنکشن دین دا
واضح طریقہ اے (جیہڑے بند ٹرم اندراج لے کے بند ٹرم قدر
دیندے نیں): !!{function}~$\Obj f^n_i$ نوں اوہ فنکشن
دیو جیہڑا $n$ بند ٹرم $t_1$،
\dots،~$t_n$ اندراج وچ لے کے $\Obj f^n_i(t_1, \dots, t_n)$
ٹرم قدر وچ کڈھے۔

باقی رہ گئی~$\eq$ دی گل۔ !!{predicate}~$\eq$ دی
تعبیر مقرر اے: $\Sat{M}{\eq[t][t']}$ اُدوں تے صرف
اُدوں جدوں $\Value{t}{M} = \Value{t'}{M}$۔ ہُن
جے اسیں ایہ انتظام کریے کہ ٹرم~$t$ دی !!{value}
$t$ آپ ہووے، تاں ایہ !!{structure} $\eq[t][t']$
دے روپ دے \emph{کسے وی} جملے نوں سچا نہیں بنائے گی
جد تک $t$ تے $t'$ عین اکّو ٹرم نہ ہون۔ ایہ یقیناً
مسئلہ اے، کیوں جے !!{function}s والی زبان وچ تقریباً
ہر دلچسپ نظریے دے قضیاں وچ $\eq[t][t']$ دے روپ
دے جملے ہون گے جتھے $t$ تے $t'$ وکھ ٹرم نیں
(مثلاً حساب دے نظرییاں وچ: $\eq[(\Obj 0+ \Obj 0)][\Obj 0]$)۔
ایس مسئلے نوں حل کرن لئی اسیں~$\Struct M$ دا میدان
بدل دیندے آں:~$\Domain{M}$ وچ اشیا دے طور تے ٹرماں دی
تھاں ٹرماں دے سیٹ ورتدے آں، تے ہر سیٹ وچ اوہ سارے
ٹرم ہوندے نیں جنہاں نوں~$\Gamma$ دے جملے برابر ہون
دا تقاضا کردے نیں۔ مثال لئی جے $\Gamma$ حساب دا
نظریہ اے، تاں انہاں وچوں اک سیٹ وچ $\Obj 0$،
$(\Obj 0 + \Obj 0)$، $(\Obj
0 \times \Obj 0)$ وغیرہ شامل ہون گے۔ ایہو سیٹ
اسیں $\Obj 0$ نوں دیواں گے؛ ثابت ہووے گا کہ ایہ
سیٹ اپنے اندر ہر ٹرم دی قدر وی اے، مثال لئی
$(\Obj 0 + \Obj 0)$ دی قدر وی۔ ایس لئی
$\eq[(\Obj 0+ \Obj 0)][\Obj 0]$ دا جملہ ایس
بدلی ہوئی !!{structure} وچ سچا ہووے گا۔}{}

اسیں ایویں کراں گے۔ پہلے !!{complete} بے تضاد سیٹاں
دیاں خاصیتاں جانچاں گے؛ خاص طور تے ثابت کراں گے کہ
!!a{complete} بے تضاد سیٹ وچ $!A \land !B$ اُدوں
تے صرف اُدوں ہوندا اے جدوں $!A$ تے~$!B$ دونیں
اوس وچ ہون؛ تے $!A \lor !B$ اُدوں تے صرف اُدوں
ہوندا اے جدوں اوہناں وچوں گھٹ توں گھٹ اک ہووے،
وغیرہ (\olref[ccs]{prop:ccs})۔\iftag{FOL}{ فیر اسیں
  جملیاں دے ”مسوّری گواہاں والے“ سیر شدہ سیٹاں
  \pnbLatin{(saturated sets)} دی تعریف تے جانچ کراں گے۔ ایسے سیٹ
  وچ شرطیہ جملے ہوندے نیں جیہڑے ہر مسوّر والے
  !!{sentence} نوں اوہدی مناسب مثالاں نال جوڑدے نیں
  (\olref[hen]{defn:henkin-exp})۔ اسیں وکھاواں گے کہ
  کوئی وی بے تضاد سیٹ~$\Gamma$ اک سیر شدہ
  سیٹ~$\Gamma'$ تک ودھایا جا سکدا اے
  (\olref[hen]{lem:henkin})۔ جے سیٹ بے تضاد، سیر شدہ
  تے !!{complete} ہووے، تاں اوہدی ایہ خاصیت وی اے:
  اوہدے وچ \iftag{prvEx}{$\lexists[x][!A(x)]$ اُدوں تے
    صرف اُدوں اے جدوں $!A(t)$ کسے بند ٹرم~$t$ لئی
    اوہدے وچ اے}{}\iftag{defEx,defAll}{}{ تے
  }\iftag{prvAll}{$\lforall[x][!A(x)]$ اُدوں تے صرف
    اُدوں اے جدوں $!A(t)$ ہر بند ٹرم~$t$ لئی
    اوہدے وچ اے}{} (\olref[hen]{prop:saturated-instances})۔}{}
فیر اسیں\iftag{FOL}{ سیر شدہ}{} بے تضاد
سیٹ~$\iftag{FOL}{\Gamma'}{\Gamma}$ لواں گے تے
وکھاواں گے کہ اوہنوں اک \iftag{FOL}{سیر شدہ،
  بے تضاد تے}{بے تضاد تے}
!!{complete} سیٹ~$\Gamma^*$ تک ودھایا جا سکدا اے
(\olref[lin]{lem:lindenbaum})۔ ایہ سیٹ $\Gamma^*$
ساڈے \iftag{FOL}{ٹرماں توں بنے نمونے
  \pnbLatin{(term model)}~$\Struct M(\Gamma^*)$}{!!{valuation}~$\pAssign v(\Gamma^*)$}
دی تعریف لئی ورتیا جائے گا۔
\iftag{FOL}{ٹرماں توں بنے نمونے دا میدان بند ٹرماں
  دا سیٹ اے، تے اوس دے !!{predicate}s دی تعبیر
  جوہری !!{sentence}s توں طے ہوندی اے}{قدر بندی
  !!{propositional variable}s توں طے ہوندی اے} جیہڑے
~$\Gamma^*$ وچ نیں (\olref[mod]{defn:termmodel})۔
اسیں\iftag{FOL}{ سیر شدہ،}{} مکمل بے تضاد
سیٹاں دیاں خاصیتاں ورت کے دکھاواں گے کہ واقعی
$\iftag{FOL}{\Sat{M(\Gamma^*)}}{\pSat{v(\Gamma^*)}}{!A}$
اُدوں تے صرف اُدوں جدوں $!A \in
\Gamma^*$ (\olref[mod]{lem:truth})؛ ایس لئی خاص طور تے
$\iftag{FOL}{\Sat{M(\Gamma^*)}}{\pSat{v(\Gamma^*)}}{\Gamma}$۔
\iftag{FOL}{آخر وچ اسیں ویکھاں گے کہ جے
  $\Gamma$ وچ~$\eq$ وی ہووے تاں ٹرماں توں بنے
  نمونے دی تعریف کیویں کرنی اے
  (\olref[ide]{defn:term-model-factor})، تے ثابت
  کراں گے کہ اوہ~$\Gamma^*$ نوں پورا کردا اے
  (\olref[ide]{lem:truth})۔}{}
\endgroup
% /OLP-0128

% OLP-0129 content/first-order-logic/completeness/complete-consistent-sets.tex
\begingroup
\def\olfilename{complete-consistent-sets}
\expandafter\def\csname ol@sectioncs\endcsname{section}

\iftag{FOL}
      {\olfileid{fol}{com}{ccs}}
      {\olfileid{pl}{com}{ccs}}
      
\olsection{\usetoken{P}{sentence} دے مکمل بے تضاد سیٹ}

\begin{defn}[مکمل سیٹ]
\ollabel{def:complete-set} !!{sentence}s دا سیٹ~$\Gamma$
\emph{!!{complete}} اُدوں تے صرف اُدوں اے جدوں ہر
!!{sentence}~$!A$ لئی یا تاں $!A \in
\Gamma$ ہووے یا $\lnot !A \in \Gamma$ ہووے۔
\end{defn}

\begin{explain}
!!^{complete} جملیاں دے سیٹ کوئی سوال بے جواب نہیں چھڈدے۔
ہر !!{sentence}~$!A$ لئی $\Gamma$ ایہ ”دسدا“ اے کہ $!A$
سچا اے یا جھوٹھا۔ !!{complete} سیٹاں دی اہمیت تمامیت دے
قضیے دے ثبوت توں اگے وی اے۔ مثال لئی، جیہڑا نظریہ
!!{complete} ہووے تے مسلّماں راہیں بیان کرن جوگ ہووے،
اوہدے ہر جملے بارے مؤثر طریقے نال فیصلہ کیتا جا سکدا اے۔
\end{explain}

\begin{explain}
!!^{complete} بے تضاد سیٹ تمامیت دے ثبوت وچ اہم نیں،
کیوں جے اسیں یقین نال کہہ سکدے آں کہ !!{sentence}s دا
ہر بے تضاد سیٹ~$\Gamma$، کسے !!a{complete} بے تضاد
سیٹ~$\Gamma^*$ وچ شامل اے۔ !!^a{complete} بے تضاد
سیٹ وچ ہر !!{sentence}~$!A$ لئی یا $!A$ ہوندا اے یا
اوہدی نفی $\lnot !A$، پر دونیں نہیں۔ ایہ خاص طور تے
\iftag{FOL}{جوہری !!{sentence}s}{قضیاتی متغیراں} لئی وی
سچ اے۔ ایس لئی !!a{complete} بے تضاد سیٹ توں
\iftag{FOL}{، ایسی زبان وچ جیہڑی مناسب !!{constant}s نال
ودھائی گئی ہووے،}{} اسیں
\iftag{FOL}{!!a{structure}}{!!a{valuation}} بنا سکدے
آں جیہدے وچ \iftag{FOL}{!!{predicate}s دی تعبیر}{قضیاتی
متغیراں نوں دتی صداقتی قدر} دا تعین ایس گل مطابق
ہووے گا کہ کِہڑے \iftag{FOL}{جوہری
!!{sentence}s}{!!{propositional variable}s}
~$\Gamma^*$ وچ نیں۔ فیر ایہ ثابت کیتا جا سکدا اے کہ ایہ
\iftag{FOL}{!!{structure}}{!!{valuation}}،~$\Gamma^*$
دے سارے !!{sentence}s نوں سچا بناندی اے (تے ایس لئی
~$\Gamma$ دے سارے جملیاں نوں وی)۔ ایس آخری گل دے ثبوت
لئی لوڑ اے کہ $\lnot !A \in \Gamma^*$ اُدوں تے صرف اُدوں
جدوں $!A \notin \Gamma^*$، تے $(!A \lor !B) \in
\Gamma^*$ اُدوں تے صرف اُدوں جدوں $!A \in \Gamma^*$
یا $!B \in \Gamma^*$، وغیرہ۔
\end{explain}

اگے اسیں $\Proves$ دی انعکاسیت، یک رُخیت تے تعدیت دیاں
خاصیتاں اکثر بنا وکھرا ذکر کیتیاں ورتاں گے (ویکھو
\iftag{FOL}{%
  \tagrefs{prfSC/{fol:seq:ptn:sec},prfND/{fol:ntd:ptn:sec},prfAX/{fol:axd:ptn:sec},prfTab/{fol:tab:ptn:sec}}}{%
  \tagrefs{prfSC/{pl:seq:ptn:sec},prfND/{pl:ntd:ptn:sec},prfAX/{pl:axd:ptn:sec},prfTab/{pl:tab:ptn:sec}}})۔

\begin{prop}
\ollabel{prop:ccs}
فرض کرو $\Gamma$ !!{complete} تے بے تضاد اے۔ تاں:
\begin{enumerate}
\item \ollabel{prop:ccs-prov-in} جے $\Gamma \Proves !A$،
  تاں $!A \in \Gamma$۔

\tagitem{prvAnd}{\ollabel{prop:ccs-and} $!A \land !B \in \Gamma$
  اُدوں تے صرف اُدوں جدوں $!A \in \Gamma$ تے $!B \in \Gamma$
  دونیں ہون۔}{}

\tagitem{prvOr}{\ollabel{prop:ccs-or} $!A \lor !B \in \Gamma$ اُدوں
  تے صرف اُدوں جدوں $!A \in \Gamma$ یا $!B \in \Gamma$
  وچوں گھٹ توں گھٹ اک ہووے۔}{}

\tagitem{prvIf}{\ollabel{prop:ccs-if} $!A \lif !B \in \Gamma$ اُدوں
  تے صرف اُدوں جدوں $!A \notin \Gamma$ یا $!B \in \Gamma$
  وچوں اک شرط پوری ہووے۔}{}
\end{enumerate}
\end{prop}

\begin{proof}
اگلیاں ساریاں گلاں لئی فرض کرو کہ $\Gamma$ !!{complete} تے
بے تضاد اے۔
\begin{enumerate}
\item جے $\Gamma \Proves !A$، تاں $!A \in \Gamma$۔

فرض کرو $\Gamma \Proves !A$۔ برہانِ خلف لئی فرض کرو کہ $!A
\notin \Gamma$۔ چونکہ $\Gamma$ !!{complete} اے، ایس لئی
$\lnot !A \in \Gamma$۔ ہُن
\iftag{FOL}{%
  \tagrefs{prfSC/{fol:seq:prv:prop:explicit-inc},
    prfND/{fol:ntd:prv:prop:explicit-inc},
    prfAX/{fol:axd:prv:prop:explicit-inc},
    prfTab/{fol:tab:prv:prop:explicit-inc}}}{%
  \tagrefs{prfSC/{pl:seq:prv:prop:explicit-inc},
    prfND/{pl:ntd:prv:prop:explicit-inc},
    prfAX/{pl:axd:prv:prop:explicit-inc},
    prfTab/{pl:tab:prv:prop:explicit-inc}}}
مطابق $\Gamma$ متضاد اے۔ ایہ $\Gamma$ دے بے تضاد ہون
دے مفروضے دے خلاف اے۔ ایس لئی $!A \notin
\Gamma$ نہیں ہو سکدا، سو $!A \in \Gamma$۔

\tagitem{defAnd}{}{%
\iftag{probAnd}{مشق۔}{%
$!A \land !B \in \Gamma$ اُدوں تے صرف اُدوں جدوں
$!A \in \Gamma$ تے $!B \in \Gamma$ دونیں ہون:

اگلے رخ لئی فرض کرو $!A \land !B \in \Gamma$۔ فیر
\iftag{FOL}{%
  \tagrefs{prfSC/{fol:seq:ppr:prop:provability-land},
    prfND/{fol:ntd:ppr:prop:provability-land},
    prfAX/{fol:axd:ppr:prop:provability-land},
    prfTab/{fol:tab:ppr:prop:provability-land}}}{%
  \tagrefs{prfSC/{pl:seq:ppr:prop:provability-land},
    prfND/{pl:ntd:ppr:prop:provability-land},
    prfAX/{pl:axd:ppr:prop:provability-land},
    prfTab/{pl:tab:ppr:prop:provability-land}}}، شق~(1)
توں $\Gamma \Proves !A$ تے $\Gamma \Proves !B$۔
\olref{prop:ccs-prov-in} مطابق $!A \in \Gamma$ تے
$!B \in \Gamma$، جیویں چاہیدا سی۔

اُلٹے رخ لئی فرض کرو $!A \in \Gamma$ تے $!B \in
\Gamma$۔ ہُن
\iftag{FOL}{%
  \tagrefs{prfSC/{fol:seq:ppr:prop:provability-land},
    prfND/{fol:ntd:ppr:prop:provability-land},
    prfAX/{fol:axd:ppr:prop:provability-land},
    prfTab/{fol:tab:ppr:prop:provability-land}}}{%
  \tagrefs{prfSC/{pl:seq:ppr:prop:provability-land},
    prfND/{pl:ntd:ppr:prop:provability-land},
    prfAX/{pl:axd:ppr:prop:provability-land},
    prfTab/{pl:tab:ppr:prop:provability-land}}}، شق~(2)
توں $\Gamma \Proves !A \land !B$۔ فیر
\olref{prop:ccs-prov-in} مطابق $!A \land
!B \in \Gamma$، جیویں چاہیدا سی۔}}

\tagitem{defOr}{}{%
\iftag{probOr}{مشق۔}{%
پہلے اسیں وکھاواں گے کہ جے $!A \lor !B \in \Gamma$،
تاں یا $!A \in \Gamma$ یا $!B \in \Gamma$۔ فرض کرو
$!A \lor !B \in \Gamma$، پر $!A
\notin \Gamma$ تے $!B \notin \Gamma$۔ چونکہ $\Gamma$
!!{complete} اے، ایس لئی $\lnot !A \in \Gamma$ تے
$\lnot !B \in \Gamma$۔ ہُن
\iftag{FOL}{%
  \tagrefs{prfSC/{fol:seq:ppr:prop:provability-lor},
    prfND/{fol:ntd:ppr:prop:provability-lor},
    prfAX/{fol:axd:ppr:prop:provability-lor},
    prfTab/{fol:tab:ppr:prop:provability-lor}}}{%
  \tagrefs{prfSC/{pl:seq:ppr:prop:provability-lor},
    prfND/{pl:ntd:ppr:prop:provability-lor},
    prfAX/{pl:axd:ppr:prop:provability-lor},
    prfTab/{pl:tab:ppr:prop:provability-lor}}}، شق (1)
مطابق $\Gamma$ متضاد اے؛ ایہ تضاد اے۔ سو یا $!A
\in \Gamma$ یا $!B \in \Gamma$۔

اُلٹے رخ لئی فرض کرو کہ $!A \in \Gamma$ یا $!B \in
\Gamma$۔ ہُن
\iftag{FOL}{%
  \tagrefs{prfSC/{fol:seq:ppr:prop:provability-lor},
    prfND/{fol:ntd:ppr:prop:provability-lor},
    prfAX/{fol:axd:ppr:prop:provability-lor},
    prfTab/{fol:tab:ppr:prop:provability-lor}}}{%
  \tagrefs{prfSC/{pl:seq:ppr:prop:provability-lor},
    prfND/{pl:ntd:ppr:prop:provability-lor},
    prfAX/{pl:axd:ppr:prop:provability-lor},
    prfTab/{pl:tab:ppr:prop:provability-lor}}}، شق (2)
توں $\Gamma \Proves !A \lor !B$۔ فیر
\olref{prop:ccs-prov-in} مطابق $!A \lor
!B \in \Gamma$، جیویں چاہیدا سی۔}}

\tagitem{defIf}{}{%
\iftag{probIf}{مشق۔}{%
اگلے رخ لئی فرض کرو $!A \lif !B \in \Gamma$، تے
برہانِ خلف لئی فرض کرو کہ $!A \in \Gamma$ تے
$!B \notin \Gamma$۔ انہاں مفروضیاں اُتے
$!A \lif !B \in \Gamma$ تے $!A \in \Gamma$۔
\iftag{FOL}{%
  \tagrefs{prfSC/{fol:seq:ppr:prop:provability-lif},
    prfND/{fol:ntd:ppr:prop:provability-lif},
    prfAX/{fol:axd:ppr:prop:provability-lif},
    prfTab/{fol:tab:ppr:prop:provability-lif}}}{%
  \tagrefs{prfSC/{pl:seq:ppr:prop:provability-lif},
    prfND/{pl:ntd:ppr:prop:provability-lif},
    prfAX/{pl:axd:ppr:prop:provability-lif},
    prfTab/{pl:tab:ppr:prop:provability-lif}}}، شق (1)
توں $\Gamma \Proves !B$۔ پر فیر
\olref{prop:ccs-prov-in} مطابق $!B \in
\Gamma$، جیہڑا $!B \notin \Gamma$ دے مفروضے دے
خلاف اے۔

اُلٹے رخ لئی پہلے اوہ صورت ویکھو جتھے $!A \notin
\Gamma$۔ چونکہ $\Gamma$ !!{complete} اے، ایس لئی
$\lnot !A \in \Gamma$۔ ہُن
\iftag{FOL}{%
  \tagrefs{prfSC/{fol:seq:ppr:prop:provability-lif},
    prfND/{fol:ntd:ppr:prop:provability-lif},
    prfAX/{fol:axd:ppr:prop:provability-lif},
    prfTab/{fol:tab:ppr:prop:provability-lif}}}{%
  \tagrefs{prfSC/{pl:seq:ppr:prop:provability-lif},
    prfND/{pl:ntd:ppr:prop:provability-lif},
    prfAX/{pl:axd:ppr:prop:provability-lif},
    prfTab/{pl:tab:ppr:prop:provability-lif}}}، شق (2)
توں $\Gamma \Proves !A \lif !B$۔ فیر
\olref{prop:ccs-prov-in} مطابق $!A \lif !B \in
\Gamma$، جیویں چاہیدا سی۔

ہُن اوہ صورت ویکھو جتھے $!B \in \Gamma$۔ ہُن
\iftag{FOL}{%
  \tagrefs{prfSC/{fol:seq:ppr:prop:provability-lif},
    prfND/{fol:ntd:ppr:prop:provability-lif},
    prfTab/{fol:tab:ppr:prop:provability-lif},
    prfAX/{fol:axd:ppr:prop:provability-lif}}}{%
  \tagrefs{prfSC/{pl:seq:ppr:prop:provability-lif},
    prfND/{pl:ntd:ppr:prop:provability-lif},
    prfTab/{pl:tab:ppr:prop:provability-lif},
    prfAX/{pl:axd:ppr:prop:provability-lif}}}
دی شق (2) توں فیر $\Gamma \Proves !A \lif !B$۔
\olref{prop:ccs-prov-in} مطابق $!A \lif !B \in
\Gamma$۔}}
\end{enumerate}
\end{proof}

\tagprob{FOL}
\begin{prob}
\olref[fol][com][ccs]{prop:ccs} دا ثبوت مکمل کرو۔
\end{prob}
\tagendprob

\tagprob{notFOL}
\begin{prob}
\olref[pl][com][ccs]{prop:ccs} دا ثبوت مکمل کرو۔
\end{prob}
\tagendprob
\endgroup
% /OLP-0129

% OLP-0130 content/first-order-logic/completeness/henkin-expansions.tex
\begingroup
\def\olfilename{henkin-expansions}
\expandafter\def\csname ol@sectioncs\endcsname{section}

\olfileid{fol}{com}{hen}

\olsection{ہینکن دی توسیع}

\begin{explain}
تمامیت دے قضیے نوں ثابت کرن وچ اک مشکل ایہ اے کہ مکمل بے تضاد
سیٹ~$\Gamma$ توں جیہڑا ماڈل اسیں بناندے آں، اوہنوں
~$\Gamma$ وچ موجود سارے مسوّر !!{formula}s سچے بنانے پینے نیں۔
ایہ یقینی بناؤن لئی اسیں لیون ہینکن دی دتی اک ترکیب ورتدے آں۔
بنیادی طور اُتے، اس ترکیب وچ زبان نوں لامحدود تعداد وچ نویں
!!{constant}s نال ودھایا جاندا اے، تے اک آزاد !!{variable}
والے ہر !!{formula}~$!A(x)$ لئی ایس صورت دا اک فارمولا
شامل کیتا جاندا اے:
\iftag{prvEx}
      {$\lexists[x][!A(x)] \lif !A(c)$}
      {$\lnot\lforall[x][!A(x)] \lif \lnot !A(c)$},
جتھے $c$ نویں !!{constant}s وچوں اک اے۔ جدوں اسیں
~$\Gamma$ نوں پورا کرن والی !!{structure} بناواں گے،
ایہ ایس گل دی ضمانت دیوے گا کہ
\iftag{prvEx}
{ہر سچے وجودی جملے دا اک گواہ}
{ہر جھوٹھے کلی جملے دی اک ردّ مثال}
انہاں نویں مستقل علامتاں وچوں موجود ہووے۔
\end{explain}

\begin{prop}
\ollabel{prop:lang-exp}
جے $\Gamma$، $\Lang L$ وچ بے تضاد اے تے $\Lang L'$
نوں $\Lang L$ وچ نویں !!{constant}s دا
!!a{denumerable} سیٹ $\Obj d_0$،
$\Obj d_1$، \dots شامل کر کے بنایا گیا اے، تاں
$\Gamma$،~$\Lang L'$ وچ وی بے تضاد اے۔
\end{prop}

\begin{defn}[سیر شدہ سیٹ]
!!{formula}s دا سیٹ $\Gamma$، زبان $\Lang {L}$ وچ
\emph{سیر شدہ} اُدوں تے صرف اُدوں اے جدوں ہر
!!{formula}~$!A(x) \in \Frm[L]$ لئی، جیہدے وچ اک
آزاد !!{variable}~$x$ اے، کوئی
!!a{constant}~$c \in \Lang{L}$ ایسا ہووے کہ
\iftag{prvEx}
      {$\lexists[x][!A(x)] \lif !A(c) \in \Gamma$}
      {$\lnot\lforall[x][!A(x)] \lif \lnot !A(c) \in \Gamma$}۔
\end{defn}

اگلی تعریف اگلے قضیے دے ثبوت وچ ورتی جائے گی۔

\begin{defn}
\ollabel{defn:henkin-exp}
فرض کرو $\Lang L'$ اوہو جیہی اے جیویں
\olref{prop:lang-exp} وچ اے۔ اک ترتیب
$!A_0(x_0)$، $!A_1(x_1)$، \dots مقرر کرو۔
ایہ انہاں سارے !!{formula}s~$!A_i(x_i)$ دی ترتیب
اے جیہڑے~$\Lang L'$ دے نیں تے جِنہاں وچ صرف اک
متغیر ($x_i$) آزاد آوندا اے۔ اسیں
!!{sentence}s~$!D_n$ نوں $n$ اُتے استقرا نال متعین کردے آں۔

$c_0$، انہاں $\Obj d_i$ وچوں پہلی ایسی
!!{constant} لو جیہڑی اسیں~$\Lang{L}$ وچ شامل کیتی
اے تے جیہڑی~$!A_0(x_0)$ وچ نہیں آوندی۔ فرض کرو
$!D_0$، \dots،~$!D_{n-1}$ پہلے ہی متعین ہو چکے نیں؛
ہُن $c_n$ نویاں !!{constant}s~$\Obj d_i$ وچوں پہلی
ایسی لو جیہڑی نہ $!D_0$، \dots،~$!D_{n-1}$ وچ
آوے تے نہ~$!A_n(x_n)$ وچ۔

ہُن $!D_{n}$ نوں ایہ !!{formula} مقرر کرو:
\iftag{prvEx}
{$\lexists[x_{n}][!A_{n}(x_{n})] \lif
  !A_{n}(c_{n})$}{$\lnot\lforall[x_{n}][!A_{n}(x_{n})] \lif \lnot
  !A_{n}(c_{n})$}۔
\end{defn}

\begin{lem}
\ollabel{lem:henkin}
ہر بے تضاد سیٹ~$\Gamma$ نوں اک سیر شدہ بے تضاد
سیٹ~$\Gamma'$ تک ودھایا جا سکدا اے۔
\end{lem}

\begin{proof}
جملیاں دا اک بے تضاد سیٹ~$\Gamma$ زبان~$\Lang{L}$
وچ دتا ہووے، تاں نویں !!{constant}s دا اک
!!a{denumerable} سیٹ شامل کر کے زبان نوں
~$\Lang{L'}$ تک ودھاؤ۔ \olref{prop:lang-exp} مطابق
$\Gamma$ ایس ودھی ہوئی زبان وچ وی بے تضاد اے۔ اگے
$!D_i$ اوہو جیہڑے \olref{defn:henkin-exp} وچ نیں، لو۔
ایہ مقرر کرو:
\begin{align*}
\Gamma_0 & = \Gamma \\
\Gamma_{n+1} & = \Gamma_n \cup \{!D_n \}
\end{align*}
یعنی $\Gamma_{n+1} = \Gamma \cup \{ !D_0, \dots, !D_n \}$،
تے $\Gamma' = \bigcup_{n} \Gamma_n$ مقرر کرو۔
$\Gamma'$ دا سیر شدہ ہونا صاف اے۔

جے $\Gamma'$ متضاد ہوندا، تاں کسے نہ کسے $n$ لئی
$\Gamma_n$ متضاد ہوندا (مشق: دسو کیوں)۔ ایس لئی
$\Gamma'$ دے بے تضاد ہون نوں وکھاؤن لئی $n$ اُتے
استقرا نال ایہ وکھانا کافی اے کہ ہر سیٹ~$\Gamma_n$
بے تضاد اے۔

استقرا دی بنیاد بس ایہ دعویٰ اے کہ $\Gamma_0 = \Gamma$
بے تضاد اے، جیہڑا قضیے دا مفروضہ اے۔ استقرا دے
اگلے قدم لئی فرض کرو $\Gamma_{n}$ بے تضاد اے پر
$\Gamma_{n+1} =
\Gamma_n \cup \{!D_n\}$ متضاد اے۔ یاد کرو
$!D_n$ ایہ اے:
\iftag{prvEx}
{$\lexists[x_{n}][!A_{n}(x_n)] \lif !A_{n}(c_{n})$}
{$\lnot\lforall[x_{n}][!A_{n}(x_n)] \lif \lnot !A_{n}(c_{n})$}،
جتھے $!A_n(x_n)$، $\Lang{L'}$ دا !!a{formula} اے
جیہدے وچ صرف متغیر~$x_n$ آزاد اے۔ $c_n$ چُنن دے
طریقے مطابق (ویکھو \olref{defn:henkin-exp})، $c_n$
نہ~$!A_n(x_n)$ وچ آوندا اے تے نہ~$\Gamma_n$ وچ۔

جے $\Gamma_n \cup \{!D_n\}$ متضاد اے، تاں
$\Gamma_n \Proves \lnot !D_n$، تے ایس لئی
اگلیاں دونیں گلاں درست نیں:
\iftag{prvEx}{
\[
\Gamma_n \Proves \lexists[x_n][!A_n(x_n)]
\qquad
\Gamma_n \Proves \lnot !A_n(c_n)
\]}{
\[
\Gamma_n \Proves \lnot\lforall[x_n][!A_n(x_n)]
\qquad
\Gamma_n \Proves !A_n(c_n)
\]}
چونکہ $c_n$ نہ
$\Gamma_n$ وچ آوندا اے نہ~$!A_n(x_n)$ وچ،
\tagrefs{prfAX/{fol:axd:qpr:thm:strong-generalization},prfSC/{fol:seq:qpr:thm:strong-generalization},prfND/{fol:ntd:qpr:thm:strong-generalization},prfTab/{fol:tab:qpr:thm:strong-generalization}}
لاگو ہوندا اے۔
\iftag{prvEx}
{$\Gamma_n \Proves \lnot !A_n(c_n)$}
{$\Gamma_n \Proves !A_n(c_n)$}
توں سانوں ایہ ملدا اے:
\iftag{prvEx}
{$\Gamma_n \Proves \lforall[x_n][\lnot !A_n(x_n)]$}
{$\Gamma_n \Proves \lforall[x_n][!A_n(x_n)]$}۔
ایس طرح ساڈے کول اک پاسے
\iftag{prvEx}
{$\Gamma_n \Proves \lexists[x_n][!A_n(x_n)]$ تے
$\Gamma_n \Proves \lforall[x_n][\lnot !A_n(x_n)]$}
{$\Gamma_n \Proves \lnot\lforall[x_n][!A_n(x_n)]$ تے
$\Gamma_n \Proves \lforall[x_n][!A_n(x_n)]$}
دونیں نیں؛ سو $\Gamma_n$ آپ متضاد اے۔
\iftag{prvEx}
{(خیال کرو کہ
\iftag{prvAll}
{$\lforall[x_n][\lnot !A_n(x_n)] \Proves
  \lnot\lexists[x_n][!A_n(x_n)]$}
{$\lforall[x_n][\lnot !A_n]$ نوں
  $\lnot\lexists[x_n][\lnot\lnot !A_n(x_n)]$ دے طور اُتے
  متعین کیتا گیا اے، تے
  $\lnot\lexists[x_n][\lnot\lnot !A_n(x_n)] \Proves
  \lnot\lexists[x_n][!A_n(x_n)]$}۔)}{}
ایہ تضاد اے: $\Gamma_n$ نوں بے تضاد فرض کیتا سی۔
لہٰذا $\Gamma_n \cup \{ !D_n\}$ بے تضاد اے۔
\end{proof}

\begin{explain}
% PNB-COM-013: وجودی خاصیت دا tag prvEx اے،prvAll نہیں؛اگلی رسمی شق ایہو فرق رکھدی اے۔
ہُن اسیں وکھاواں گے کہ جیہڑے سیٹ
\emph{مکمل}، بے تضاد تے سیر شدہ نیں، اوہناں دی
خاصیت ایہ اے کہ \iftag{prvAll}{اوہناں وچ اک
  کلی طور اُتے مسوّر !!{sentence} اُدوں تے صرف
  اُدوں ہوندا اے جدوں اوہدیاں ساریاں مثالاں
  اوہناں وچ ہون}{}\iftag{defAll,defEx}{}{ تے
  }\iftag{prvEx}{اوہناں وچ اک وجودی طور اُتے
  مسوّر !!{sentence} اُدوں تے صرف اُدوں ہوندا
  اے جدوں اوہدی گھٹ توں گھٹ اک مثال اوہناں
  وچ ہووے}{}۔ اسیں ایہ گل وکھاؤن لئی ورتاں گے
کہ مکمل، بے تضاد، سیر شدہ سیٹ توں بنائی
!!{structure} اوہدے سارے مسوّر جملیاں نوں سچا
بناندی اے۔
\end{explain}

\begin{prop}\ollabel{prop:saturated-instances}
فرض کرو $\Gamma$ مکمل، بے تضاد تے سیر شدہ اے۔
\begin{tagenumerate}{prvEx,prvAll}
\tagitem{prvEx}{$\lexists[x][!A(x)] \in \Gamma$
  اُدوں تے صرف اُدوں جدوں $!A(t) \in \Gamma$
  ہووے، گھٹ توں گھٹ اک بند ٹرم~$t$ لئی۔}{}
\tagitem{prvAll}{$\lforall[x][!A(x)] \in \Gamma$
  اُدوں تے صرف اُدوں جدوں $!A(t) \in \Gamma$
  ہووے، ہر بند ٹرم~$t$ لئی۔}{}
\end{tagenumerate}
\end{prop}

\begin{proof}
  \begin{tagenumerate}{prvEx,prvAll}
  \tagitem{prvEx}{%
    \iftag{probEx}{مشق۔}{پہلے فرض کرو کہ $\lexists[x][!A(x)]
      \in \Gamma$۔ چونکہ $\Gamma$ سیر شدہ اے،
      $(\lexists[x][!A(x)] \lif !A(c)) \in \Gamma$
      کسے !!{constant}~$c$ لئی۔
      \tagrefs{prfAX/{fol:axd:ppr:prop:provability-lif},prfSC/{fol:seq:ppr:prop:provability-lif},prfND/{fol:ntd:ppr:prop:provability-lif},prfTab/{fol:tab:ppr:prop:provability-lif}}
      دی شق~(1) تے
      \olref[ccs]{prop:ccs}\olref[ccs]{prop:ccs-prov-in}
      مطابق $!A(c) \in \Gamma$۔

    اُلٹے رخ لئی سیر شدگی دی لوڑ نہیں: فرض کرو
    $!A(t) \in \Gamma$۔ تاں
    \tagrefs{prfAX/{fol:axd:qpr:prop:provability-quantifiers},prfSC/{fol:seq:qpr:prop:provability-quantifiers},prfND/{fol:ntd:qpr:prop:provability-quantifiers},prfTab/{fol:tab:qpr:prop:provability-quantifiers}}
    دی شق~(1) مطابق $\Gamma \Proves \lexists[x][!A(x)]$۔
    \olref[ccs]{prop:ccs}\olref[ccs]{prop:ccs-prov-in}
    مطابق $\lexists[x][!A(x)] \in \Gamma$۔}}{}
  \tagitem{prvAll}{%
    \iftag{probAll}{مشق۔}{فرض کرو کہ $!A(t) \in \Gamma$
      ہر بند ٹرم~$t$ لئی۔ برہانِ خلف لئی فرض کرو
      $\lforall[x][!A(x)] \notin \Gamma$۔ چونکہ
      $\Gamma$ مکمل اے، ایس لئی
      $\lnot\lforall[x][!A(x)] \in \Gamma$۔
      سیر شدگی مطابق
      \iftag{prvEx}{$(\lexists[x][\lnot !A(x)] \lif \lnot !A(c)) \in
        \Gamma$}{$(\lnot\lforall[x][!A(x)] \lif \lnot !A(c)) \in
        \Gamma$} کسے !!{constant}~$c$ لئی۔ مفروضے
      مطابق، چونکہ $c$ اک بند
      ٹرم اے، $!A(c) \in \Gamma$۔ پر ایہ
      $\Gamma$ نوں متضاد بنا دیندا اے۔ (مشق:
      اوہ !!{derivation} دیو جیہڑی وکھاوے کہ
      \[
      \lnot \lforall[x][!A(x)],
      \iftag{prvEx}{\lexists[x][\lnot !A(x)] \lif \lnot
        !A(c)}{\lnot\lforall[x][!A(x)] \lif \lnot
        !A(c)}, !A(c)
      \]
      متضاد نیں۔)

      اُلٹے رخ لئی وی سیر شدگی دی لوڑ نہیں:
      فرض کرو $\lforall[x][!A(x)] \in \Gamma$۔
      تاں
      \tagrefs{prfAX/{fol:axd:qpr:prop:provability-quantifiers},prfSC/{fol:seq:qpr:prop:provability-quantifiers},prfND/{fol:ntd:qpr:prop:provability-quantifiers},prfTab/{fol:tab:qpr:prop:provability-quantifiers}}
      دی شق~(2) مطابق $\Gamma \Proves !A(t)$۔
      \olref[ccs]{prop:ccs} مطابق سانوں
      $!A(t) \in \Gamma$ ملدا اے۔}}{}
  \end{tagenumerate}
\end{proof}
\endgroup
% /OLP-0130

% OLP-0131 content/first-order-logic/completeness/lindenbaums-lemma.tex
\begingroup
\def\olfilename{lindenbaums-lemma}
\expandafter\def\csname ol@sectioncs\endcsname{section}

\iftag{FOL}
      {\olfileid{fol}{com}{lin}}
      {\olfileid{pl}{com}{lin}}

\olsection{لنڈنبام دا معاون قضیہ \pnbLatin{(Lindenbaum's lemma)}}

\begin{explain}
ہُن اسیں اک معاون قضیہ ثابت کراں گے جیہڑا وکھاندا اے کہ
!!{sentence}s دا ہر بے تضاد سیٹ، جملیاں دے اک ایسے
سیٹ وچ شامل ہوندا اے جیہڑا صرف بے تضاد ہی نہیں بلکہ
!!{complete} وی اے۔ ثبوت وچ اسیں ہر واری اک
!!{sentence} شامل کردے آں تے ہر قدم اُتے سیٹ دی
بے تضادی برقرار رکھدے آں۔ ایہ اس طرح کردے آں کہ
ہر $!A$ لئی کسے نہ کسے مرحلے اُتے یا $!A$ یا
$\lnot !A$ شامل ہو جائے۔ ایس تعمیر دے سارے
مرحلیاں دا اتحاد فیر یا $!A$ یا اوہدی نفی~$\lnot !A$
رکھدا اے، ایس لئی مکمل اے۔ اوہ بے تضاد وی اے،
کیوں جے ہر مرحلے اُتے اسیں تضاد پیدا نہ ہون دی
ضمانت کردے آں۔
\end{explain}

\begin{lem}[لنڈنبام دا معاون قضیہ \pnbLatin{(Lindenbaum's lemma)}]
\ollabel{lem:lindenbaum} ہر بے تضاد سیٹ~$\Gamma$
نوں زبان~$\Lang{L}$ وچ اک !!a{complete} تے بے تضاد
سیٹ~$\Gamma^*$ تک ودھایا جا سکدا اے۔
\end{lem}

\begin{proof}
فرض کرو $\Gamma$ بے تضاد اے۔ $!A_0$، $!A_1$،
\dots{} نوں~$\Lang L$ دے سارے !!{sentence}s دی اک
ترتیب منو۔ $\Gamma_0 = \Gamma$ مقرر کرو، تے
\[
\Gamma_{n+1} =
\begin{cases}
\Gamma_n \cup \{ !A_n \} & \textrm{جے $\Gamma_n \cup \{!A_n\}$ بے
  تضاد ہووے؛} \\
\Gamma_n \cup \{ \lnot !A_n \} & \textrm{ورنہ۔}
\end{cases}
\]
$\Gamma^* = \bigcup_{n \geq 0} \Gamma_n$ مقرر کرو۔

ہر $\Gamma_n$ بے تضاد اے: $\Gamma_0$ تعریف مطابق
بے تضاد اے۔ جے $\Gamma_{n+1} = \Gamma_n \cup \{!A_n\}$،
تاں ایس دی وجہ ایہ اے کہ ایہ ودھایا سیٹ بے تضاد اے۔
جے اوہ بے تضاد نہیں، تاں
$\Gamma_{n+1} = \Gamma_n \cup \{\lnot !A_n\}$۔
سانوں جانچنا اے کہ $\Gamma_n \cup \{\lnot !A_n\}$
بے تضاد اے۔ فرض کرو ایہ بے تضاد نہیں۔ تاں
$\Gamma_n \cup \{!A_n\}$ تے
$\Gamma_n \cup \{\lnot !A_n\}$ \emph{دونیں}
متضاد نیں۔ ایس دا مطلب اے کہ
\iftag{FOL}{%
  \tagrefs{prfAX/{fol:axd:prv:prop:provability-exhaustive},
    prfSC/{fol:seq:prv:prop:provability-exhaustive},
    prfND/{fol:ntd:prv:prop:provability-exhaustive},
    prfTab/{fol:tab:prv:prop:provability-exhaustive}}}{%
  \tagrefs{prfAX/{pl:axd:prv:prop:provability-exhaustive},
    prfSC/{pl:seq:prv:prop:provability-exhaustive},
    prfND/{pl:ntd:prv:prop:provability-exhaustive},
    prfTab/{pl:tab:prv:prop:provability-exhaustive}}}
مطابق $\Gamma_n$ متضاد ہووے گا، جیہڑا استقرا دے
مفروضے دے خلاف اے۔

ہر~$n$ تے ہر $i < n$ لئی $\Gamma_i \subseteq \Gamma_n$۔
ایہ $n$ اُتے سادہ استقرا توں ملدا اے۔ $n=0$ لئی
کوئی $i < 0$ نہیں، ایس لئی دعویٰ آپے درست اے۔
استقرائی قدم لئی فرض کرو ایہ~$n$ لئی درست اے۔
اسیں وکھاواں گے کہ جے $i < n+1$، تاں
$\Gamma_i \subseteq \Gamma_{n+1}$۔
تعمیر مطابق $\Gamma_{n+1} = \Gamma_n \cup \{!A_n\}$
یا $=
\Gamma_n \cup \{\lnot !A_n\}$۔
سو $\Gamma_n \subseteq \Gamma_{n+1}$۔
جے $i < n+1$، تاں استقرا دے مفروضے مطابق
$\Gamma_i \subseteq \Gamma_n$ (جے $i < n$)،
یا صاف گل $\Gamma_n
\subseteq \Gamma_n$ (جے $i = n$) مطابق۔
$\Gamma_i \subseteq
\Gamma_{n+1}$ ایس توں~$\subseteq$ دی تعدیت نال ملدا اے۔

ایس توں $\Gamma^*$ دی بے تضادی وی نکلدی اے۔
وجہ ایہ اے: فرض کرو $\Gamma' \subseteq \Gamma^*$
اک محدود سیٹ اے۔
% PNB-COM-017: خالی محدود ذیلی سیٹ لئی largest stage موجود نہیں؛بے تضاد اصل مرحلے دا ذیلی سیٹ ہون کافی اے۔
جے ایہ سیٹ خالی اے، تاں اوہ اصل بے تضاد مرحلے دا ذیلی سیٹ اے، سو بے تضاد اے۔ ورنہ ہر $!B \in \Gamma'$ کسے
مرحلے~$\Gamma_i$ وچ وی اے، جیہدا اشاریہ~$i$ اے۔
انہاں اشارییاں وچوں
$n$ سب توں وڈا لو۔ چونکہ
$\Gamma_i \subseteq \Gamma_n$ جے $i \le n$ ہووے،
ایس لئی ہر
$!B \in \Gamma'$ وی $\in \Gamma_n$ اے،
یعنی $\Gamma' \subseteq \Gamma_n$، تے
$\Gamma_n$~بے تضاد اے۔ ایس طرح ہر محدود ذیلی
سیٹ $\Gamma' \subseteq
\Gamma^*$ بے تضاد اے۔
\iftag{FOL}{%
  \tagrefs{prfAX/{fol:axd:ptn:prop:proves-compact},
    prfSC/{fol:seq:ptn:prop:proves-compact},
    prfND/{fol:ntd:ptn:prop:proves-compact},
    prfTab/{fol:tab:ptn:prop:proves-compact}}}{%
  \tagrefs{prfAX/{pl:axd:ptn:prop:proves-compact},
    prfSC/{pl:seq:ptn:prop:proves-compact},
    prfND/{pl:ntd:ptn:prop:proves-compact},
    prfTab/{pl:tab:ptn:prop:proves-compact}}}
مطابق $\Gamma^*$ بے تضاد اے۔

$\Frm[L]$ دا ہر !!{sentence} اوہدی فہرست وچ
آوندا اے جیہڑی~$\Gamma^*$ نوں متعین کرن لئی
ورتی گئی۔ جے $!A_n \notin \Gamma^*$، تاں
ایس دی وجہ ایہ اے کہ $\Gamma_n \cup \{!A_n\}$
متضاد سی۔ پر فیر $\lnot !A_n
\in \Gamma^*$، ایس لئی $\Gamma^*$
!!{complete} اے۔
\end{proof}
\endgroup
% /OLP-0131

% OLP-0132 content/first-order-logic/completeness/construction-of-model.tex
\begingroup
\def\olfilename{construction-of-model}
\expandafter\def\csname ol@sectioncs\endcsname{section}

\iftag{FOL}
      {\olfileid{fol}{com}{mod}}
      {\olfileid{pl}{com}{mod}}

\olsection{اک نمونے دی تعمیر}

\begin{explain}
\iftag{FOL}{فی الحال ساڈی توجہ $\eq$ اُتے نہیں، یعنی اسیں
  صرف ایہ وکھانا چاہندے آں کہ !!{sentence}s دا اک بے
  تضاد سیٹ~$\Gamma$ جیہدے وچ~$\eq$ نہیں، پورا ہون
  جوگ اے۔ پہلے اسیں~$\Gamma$ نوں اک بے تضاد،
  !!{complete} تے سیر شدہ سیٹ~$\Gamma^*$ تک ودھاندے
  آں۔ ایس صورت وچ نمونے~$\Struct{M(\Gamma^*)}$ دی
  تعریف سادہ اے: اسیں~$\Lang{L'}$ دے بند ٹرماں دا
  سیٹ میدان لَیندے آں۔ ہر !!{constant} دی تعبیر
  اوہ آپ اے، تے عام طور اُتے ہر بند ٹرم~$t$ لئی
  $\Value{t}{M(\Gamma^*)} = t$ یقینی بناندے آں۔
  !!{predicate}s دی تعبیری توسیع ایس طرح مقرر
  کردے آں کہ جوہری !!{sentence}،
  $\Struct{M(\Gamma^*)}$ وچ اُدوں تے صرف اُدوں
  سچا ہووے جدوں اوہ~$\Gamma^*$ وچ ہووے۔ ظاہر اے،
  ایس نال~$\Gamma^*$ دے سارے جوہری
  !!{sentence}s، $\Struct{M(\Gamma^*)}$ وچ سچے
  ہون گے۔ باقی جملے وی سچے ہون گے جے شروع دا
  $\Gamma^*$ بے تضاد، مکمل تے سیر شدہ اے۔}{ہُن
  اسیں ایسی !!a{valuation} متعین کرن لئی تیار آں
  جیہڑی سارے $!A \in \Gamma$ نوں سچا بناندی اے۔
  ایس لئی پہلے لنڈنبام دا معاون قضیہ ورت کے
  سانوں مکمل بے تضاد $\Gamma^* \supseteq
  \Gamma$ ملدا اے۔ فیر~$\Gamma^*$ دے
  !!{propositional variable}s توں
  ~$\pAssign v(\Gamma^*)$ متعین کردے آں۔}
\end{explain}

\iftag{FOL}{%
\begin{defn}[ٹرماں توں بنا نمونہ \pnbLatin{(term model)}]
\ollabel{defn:termmodel}
فرض کرو $\Gamma^*$، زبان~$\Lang L$ دے
!!{sentence}s دا !!a{complete}، بے تضاد تے سیر
شدہ سیٹ اے۔ \emph{ٹرماں توں بنا
  نمونہ}~$\Struct M(\Gamma^*)$،~$\Gamma^*$ لئی
اگلی !!{structure} اے:
\begin{enumerate}
\item !!{domain}~$\Domain{M(\Gamma^*)}$،~$\Lang L$ دے
  سارے بند ٹرماں دا سیٹ اے۔
\item !!a{constant} $c$ دی تعبیر $c$ آپ اے:
  $\Assign{c}{M(\Gamma^*)} = c$۔
\item !!{function}~$f$ دی تعبیر اوہ فنکشن اے
  جیہڑا بند ٹرماں $t_1$، \dots، $t_n$ نوں
  اندراجاں دے طور اُتے لے کے بند ٹرم
  $f(t_1, \dots, t_n)$ قدر دیندا اے:
\[
\Assign{f}{M(\Gamma^*)}(t_1, \dots, t_n) = f(t_1,\dots, t_n)
\]
\item جے $R$ اک $n$-جگہ !!{predicate} اے، تاں
  \[
  \tuple{t_1, \dots,
  t_n} \in \Assign{R}{M(\Gamma^*)} \text{ اُدوں تے صرف اُدوں } \Atom{R}{t_1, \dots,
    t_n} \in \Gamma^*.
  \]
\end{enumerate}
\end{defn}
}{%
\begin{defn}
\ollabel{defn:termmodel}
فرض کرو $\Gamma^*$، !!{formula}s دا اک مکمل
بے تضاد سیٹ اے۔ فیر اسیں ایہ مقرر کردے آں:
\[
\pAssign v(\Gamma^*)(p) = \begin{cases}
  \True & \text{جے $p \in \Gamma^*$}\\
  \False & \text{جے $p \notin \Gamma^*$}
  \end{cases}
\]
\end{defn}
}

\iftag{FOL}{%
ہُن اسیں بند ٹرماں لئی جانچاں گے کہ واقعی $\Value{t}{M(\Gamma^*)} = t$۔

\begin{lem}
\ollabel{lem:val-in-termmodel} فرض کرو $\Struct M(\Gamma^*)$،
\olref{defn:termmodel} دا ٹرماں توں بنا نمونہ اے؛
% PNB-COM-014: بغیر variable assignment دے value تے ایس induction دا دائرہ بند ٹرم اے۔
تاں ہر بند ٹرم لئی $\Value{t}{M(\Gamma^*)} = t$۔
\end{lem}
\begin{proof}
ثبوت $t$ اُتے استقرا نال اے۔ بنیاد وچ، جدوں $t$
!!a{constant} اے، دعویٰ ٹرماں توں بنے نمونے
دی تعریف توں سیدھا ملدا اے۔ استقرائی قدم لئی
فرض کرو $t_1, \ldots, t_n$ بند ٹرم نیں
جِنہاں لئی $\Value{t_i}{M(\Gamma^*)} = t_i$،
تے $f$ اک $n$-مقامی !!{function} اے۔ تاں
\begin{align*}
\Value{f(t_1,\ldots,t_n)}{M(\Gamma^*)} &= \Assign{f}{M(\Gamma^*)}(\Value{t_1}
 {M(\Gamma^*)},
\ldots, \Value{t_n}{M(\Gamma^*)}) \\
&= \Assign{f}{M(\Gamma^*)}(t_1, \dots, t_n) \\
&= f(t_1,\dots, t_n),
\end{align*}
تے ایس لئی استقرا نال ایہ ہر بند ٹرم~$t$
لئی درست اے۔
\end{proof}
}{}

\iftag{FOL}{%
\begin{explain}
  \iftag{prvEx}{کوئی !!^a{structure}~$\Struct{M}$ اک وجودی
    طور اُتے مسوّر !!{sentence}~$\lexists[x][!A(x)]$
    سچا بنا سکدی اے، بھاویں اوہدی کوئی ایسی
    مثال~$!A(t)$ نہ ہووے جیہنوں اوہ سچا بناندی اے۔}{}
  \iftag{prvAll}{کوئی !!^a{structure}~$\Struct{M}$ ساریاں
    مثالاں~$!A(t)$ سچیاں بنا سکدی اے اک کلی طور
    اُتے مسوّر !!{sentence}~$\lforall[x][!A(x)]$ دیاں،
    پر~$\lforall[x][!A(x)]$ نوں سچا نہ بنائے۔}{}
  وجہ ایہ اے کہ عام طور اُتے~$\Domain{M}$ دا ہر
  !!{element} کسے بند ٹرم دی قدر نہیں ہوندا
  ($\Struct{M}$ پوری طرح بند ٹرماں نال ڈھکی نہ وی
  ہو سکدی اے)۔ ایس لئی پورا ہون دا تعلق
  متغیراں دی قدر بندیاں راہیں متعین کیتا جاندا اے۔
  پر ساڈے ٹرماں توں بنے نمونے~$\Struct{M(\Gamma^*)}$
  لئی ایہ لوڑ نہیں، کیونکہ اوہ ڈھکیا ہویا اے۔
  اگلا نتیجہ ایہی وکھاندا اے۔
\end{explain}

\begin{prop}
\ollabel{prop:quant-termmodel}
فرض کرو $\Struct M(\Gamma^*)$،
\olref{defn:termmodel} دا ٹرماں توں بنا نمونہ اے۔
\begin{tagenumerate}{prvEx,prvAll}
\tagitem{prvEx}{$\Sat{M(\Gamma^*)}{\lexists[x][!A(x)]}$
  اُدوں تے صرف اُدوں جدوں
  $\Sat{M(\Gamma^*)}{!A(t)}$ ہووے گھٹ توں گھٹ اک
  بند ٹرم~$t$ لئی۔}{}
\tagitem{prvAll}{$\Sat{M(\Gamma^*)}{\lforall[x][!A(x)]}$
  اُدوں تے صرف اُدوں جدوں
  $\Sat{M(\Gamma^*)}{!A(t)}$ ہووے ہر بند
  ٹرم~$t$ لئی۔}{}
\end{tagenumerate}
\end{prop}

\begin{proof}
\begin{tagenumerate}{prvEx,prvAll}
\tagitem{prvEx}{%
  \iftag{probEx}{مشق۔}{\olref[syn][ass]{prop:sat-quant}
    مطابق $\Sat{M(\Gamma^*)}{\lexists[x][!A(x)]}$
    اُدوں تے صرف اُدوں جدوں گھٹ توں گھٹ اک
    متغیر-قدر بندی~$s$ لئی
    $\Sat{M(\Gamma^*)}{!A(x)}[s]$۔
    چونکہ $\Domain{M(\Gamma^*)}$،~$\Lang{L}$ دے
    بند ٹرماں توں بنیا اے، ایس لئی ایہ اُدوں
    تے صرف اُدوں ہووے گا جدوں کوئی بند ٹرم~$t$
    ایسا ہووے کہ $s(x) = t$ تے
    $\Sat{M(\Gamma^*)}{!A(x)}[s]$۔
    \olref[fol][syn][ext]{prop:ext-formulas} مطابق
    $\Sat{M(\Gamma^*)}{!A(x)}[s]$ اُدوں تے صرف
    اُدوں جدوں $\Sat{M(\Gamma^*)}{!A(t)}[s]$،
    جتھے $s(x) = t$۔ ہُن
    \olref[fol][syn][ass]{prop:sentence-sat-true} مطابق
    $\Sat{M(\Gamma^*)}{!A(t)}[s]$ اُدوں تے صرف
    اُدوں جدوں $\Sat{M(\Gamma^*)}{!A(t)}$،
    کیوں جے $!A(t)$ اک جملہ اے۔}}{}
\tagitem{prvAll}{%
  \iftag{probAll}{مشق۔}{\olref[syn][ass]{prop:sat-quant}
    مطابق $\Sat{M(\Gamma^*)}{\lforall[x][!A(x)]}$
    اُدوں تے صرف اُدوں جدوں ہر متغیر-قدر بندی
    $s$ لئی $\Sat{M(\Gamma^*)}{!A(x)}[s]$۔
    یاد کرو $\Domain{M(\Gamma^*)}$،
    ~$\Lang{L}$ دے بند ٹرماں توں بنیا اے؛
    ایس لئی ہر بند ٹرم~$t$ لئی $s(x) = t$
    والی متغیر-قدر بندی ملدی اے، تے ہر متغیر-قدر
    بندی لئی $s(x)$ کوئی بند ٹرم~$t$ اے۔
    \olref[fol][syn][ext]{prop:ext-formulas} مطابق
    $\Sat{M(\Gamma^*)}{!A(x)}[s]$ اُدوں تے
    صرف اُدوں جدوں $\Sat{M(\Gamma^*)}{!A(t)}[s]$،
    جتھے $s(x) = t$۔
    \olref[fol][syn][ass]{prop:sentence-sat-true}
    مطابق $\Sat{M(\Gamma^*)}{!A(t)}[s]$
    اُدوں تے صرف اُدوں جدوں
    $\Sat{M(\Gamma^*)}{!A(t)}$، کیوں جے
    $!A(t)$ اک جملہ اے۔}}{}
\end{tagenumerate}
\end{proof}
}{}

\tagprob[FOL]{probEx,probAll}
\begin{prob}
\olref[fol][com][mod]{prop:quant-termmodel} دا ثبوت مکمل کرو۔
\end{prob}
\tagendprob

\begin{lem}[صدق دا معاون قضیہ]
\ollabel{lem:truth} \iftag{FOL}{فرض کرو $!A$ وچ~$\eq$ نہیں آوندا۔ تاں}{}
$\iftag{FOL}{\Sat{M(\Gamma^*)}}{\pSat{v(\Gamma^*)}}{!A}$
اُدوں تے صرف اُدوں جدوں $!A \in \Gamma^*$۔
\end{lem}

\begin{proof}
اسیں دونیں رخ اکٹھے، $!A$ اُتے استقرا نال ثابت کردے آں۔
\begin{enumerate}
\tagitem{prvFalse}{\indcase{!A}{\lfalse}
  {$\iftag{FOL}{\Sat/{M(\Gamma^*)}{\lfalse}}{\pSat/{v(\Gamma^*)}{\lfalse}}$
    پورا ہون دی تعریف مطابق۔ دوجے پاسے
    $\lfalse \notin
    \Gamma^*$، کیونکہ $\Gamma^*$ بے تضاد اے۔}}{}

\tagitem{prvTrue}{\indcase{!A}{\ltrue}
  {$\iftag{FOL}{\Sat{M(\Gamma^*)}}{\pSat{v(\Gamma^*)}}{\ltrue}$
    پورا ہون دی تعریف مطابق۔ دوجے پاسے
    $\ltrue \in
    \Gamma^*$، کیونکہ $\Gamma^*$ بے تضاد تے
    !!{complete} اے، تے
    $\Gamma^* \Proves \ltrue$۔}}{}

\tagitem{FOL}{\indcase{!A}{R(t_1, \dots, t_n)}
  {$\Sat{M(\Gamma^*)}{\Atom{R}{t_1, \dots, t_n}}$
    اُدوں تے صرف اُدوں جدوں $\tuple{t_1,
      \dots, t_n} \in \Assign{R}{M(\Gamma^*)}$
    (پورا ہون دی تعریف مطابق)، تے ایہ
    اُدوں تے صرف اُدوں جدوں
    $R(t_1, \dots, t_n) \in \Gamma^*$
    ($\Struct
    M(\Gamma^*)$ دی تعمیر مطابق)۔}}{\indcase{!A}{p}
  {$\pSat{v(\Gamma^*)}{p}$ اُدوں تے صرف
    اُدوں جدوں $\pAssign v(\Gamma^*)(p) = \True$
    (پورا ہون دی تعریف مطابق)، تے ایہ
    اُدوں تے صرف اُدوں جدوں
    $p \in \Gamma^*$ ($\pAssign v(\Gamma^*)$
    دی تعمیر مطابق)۔}}

\tagitem{prvNot}{%
  \indcase{!A}{\lnot !B}
    {$\iftag{FOL}{\Sat{M(\Gamma^*)}}{\pSat{v(\Gamma^*)}}{\indfrm}$
    اُدوں تے صرف اُدوں جدوں
    $\iftag{FOL}{\Sat/{M(\Gamma^*)}{!B}}{\pSat/{v(\Gamma^*)}{!B}}$
    (پورا ہون دی تعریف مطابق)۔ استقرا دے مفروضے
    توں $\iftag{FOL}{\Sat/{M(\Gamma^*)}{!B}}{\pSat/{v(\Gamma^*)}{!B}}$
    اُدوں تے صرف اُدوں جدوں
    $!B \notin \Gamma^*$۔
    چونکہ $\Gamma^*$ بے تضاد تے !!{complete}
    اے، $!B
    \notin \Gamma^*$ اُدوں تے صرف اُدوں جدوں
    $\lnot !B \in \Gamma^*$۔}}{}

\tagitem{prvAnd}{%
\iftag{probAnd}{%
  \indcase!{!A}{!B \land !C}{}}{%
  \indcase{!A}{!B \land
    !C}{$\iftag{FOL}{\Sat{M(\Gamma^*)}}{\pSat{v(\Gamma^*)}}{\indfrm}$
    اُدوں تے صرف اُدوں جدوں
    $\iftag{FOL}{\Sat{M(\Gamma^*)}}{\pSat{v(\Gamma^*)}}{!B}$
    تے $\iftag{FOL}{\Sat{M(\Gamma^*)}}{\pSat{v(\Gamma^*)}}{!C}$
    دونیں ہون (پورا ہون دی تعریف مطابق)؛
    استقرا دے مفروضے مطابق ایہ اُدوں تے صرف
    اُدوں جدوں $!B \in \Gamma^*$ تے
    $!C \in
    \Gamma^*$ دونیں ہون۔
    \olref[ccs]{prop:ccs}\olref[ccs]{prop:ccs-and}
    مطابق ایہ اُدوں تے صرف اُدوں جدوں
    $(!B \land !C) \in \Gamma^*$۔}}}{}

\tagitem{prvOr}{%
\iftag{probOr}{%
  \indcase!{!A}{!B \lor !C}{}}{%
  \indcase{!A}{!B \lor !C}
    {$\iftag{FOL}{\Sat{M(\Gamma^*)}}{\pSat{v(\Gamma^*)}}{\indfrm}$
    اُدوں تے صرف اُدوں جدوں
    $\iftag{FOL}{\Sat{M(\Gamma^*)}}{\pSat{v(\Gamma^*)}}{!B}$
    یا $\iftag{FOL}{\Sat{M(\Gamma^*)}}{\pSat{v(\Gamma^*)}}{!C}$
    ہووے (پورا ہون دی تعریف مطابق)؛
    استقرا دے مفروضے مطابق ایہ اُدوں تے صرف
    اُدوں جدوں $!B \in \Gamma^*$ یا
    $!C \in
    \Gamma^*$۔ ایہ اُدوں تے صرف اُدوں جدوں
    $(!B
    \lor !C) \in \Gamma^*$
    (\olref[ccs]{prop:ccs}\olref[ccs]{prop:ccs-or}
    مطابق)۔}}}{}

\tagitem{prvIf}{%
\iftag{probIf}{%
  \indcase!{!A}{!B \lif !C}{}}{%
  \indcase{!A}{!B \lif !C}
    {$\iftag{FOL}{\Sat{M(\Gamma^*)}}{\pSat{v(\Gamma^*)}}{\indfrm}$
    اُدوں تے صرف اُدوں جدوں
    $\iftag{FOL}{\Sat/{M(\Gamma^*)}}{\pSat/{v(\Gamma^*)}}{!B}$
    یا $\iftag{FOL}{\Sat{M (\Gamma^*)}}{\pSat{v(\Gamma^*)}}{!C}$
    ہووے (پورا ہون دی تعریف مطابق)؛
    استقرا دے مفروضے مطابق ایہ اُدوں تے صرف
    اُدوں جدوں $!B \notin \Gamma^*$ یا
    $!C \in
    \Gamma^*$۔ ایہ اُدوں تے صرف اُدوں جدوں
    $(!B
    \lif !C) \in \Gamma^*$
    (\olref[ccs]{prop:ccs}\olref[ccs]{prop:ccs-if}
    مطابق)۔}}}{}

\iftag{FOL}{%
\tagitem{prvAll}{%
  \iftag{probAll}{%
    \indcase!{!A}{\lforall[x][!B(x)]}{}}{%
    % PNB-COM-003: بند ٹرم دی قید ماڈل دی تعریف تے prop:quant-termmodel مطابق۔
    \indcase{!A}{\lforall[x][!B(x)]}{$\Sat{M(\Gamma^*)}{\indfrm}$
      اُدوں تے صرف اُدوں جدوں
      $\Sat{M(\Gamma^*)}{!B(t)}$ ہر بند ٹرم~$t$ لئی
      (\olref{prop:quant-termmodel})۔ استقرا دے
      مفروضے مطابق ایہ اُدوں تے صرف اُدوں جدوں
      $!B(t) \in \Gamma^*$ ہر بند ٹرم~$t$ لئی؛
      \olref[hen]{prop:saturated-instances} مطابق
      ایہ اگوں اُدوں تے صرف اُدوں جدوں
      % PNB-COM-004: موجودہ صورت B(x) اے؛ منجمد A(x) نال تصادم۔
      $\lforall[x][!B(x)] \in \Gamma^*$۔}}}{}

\tagitem{prvEx}{%
  \iftag{probEx}{%
    \indcase!{!A}{\lexists[x][!B(x)]}{}}{%
    % PNB-COM-003: وجودی صورت وچ وی بند ٹرم دی قید۔
    \indcase{!A}{\lexists[x][!B(x)]}{$\Sat{M(\Gamma^*)}{\indfrm}$
      اُدوں تے صرف اُدوں جدوں
      $\Sat{M(\Gamma^*)}{!B(t)}$ گھٹ توں
      گھٹ اک بند ٹرم~$t$ لئی
      (\olref{prop:quant-termmodel})۔ استقرا دے
      مفروضے مطابق ایہ اُدوں تے صرف اُدوں جدوں
      $!B(t) \in \Gamma^*$ گھٹ توں گھٹ اک
      بند ٹرم~$t$ لئی۔ ہُن
      \olref[hen]{prop:saturated-instances} مطابق
      ایہ اگوں اُدوں تے صرف اُدوں جدوں
      $\lexists[x][!B(x)] \in \Gamma^*$۔}}}{}
}{}
\end{enumerate}
\end{proof}


\tagprob[FOL]{probOr,probAnd,probIf,probEx,probAll}

\begin{prob}
\olref[fol][com][mod]{lem:truth} دا ثبوت مکمل کرو۔
\end{prob}

\tagendprob

\tagprob[notFOL]{probOr,probAnd,probIf}

\begin{prob}
\olref[pl][com][mod]{lem:truth} دا ثبوت مکمل کرو۔
\end{prob}

\tagendprob
\endgroup
% /OLP-0132

% OLP-0133 content/first-order-logic/completeness/identity.tex
\begingroup
\def\olfilename{identity}
\expandafter\def\csname ol@sectioncs\endcsname{section}

\olfileid{fol}{com}{ide}
\olsection{برابری}

\begin{explain}
پچھلے حصے وچ ٹرماں توں بنے نمونے دی تعمیر،
اوہناں سیٹاں~$\Gamma$ لئی پہلے درجے دی منطق
دی تمامیت قائم کرن واسطے کافی اے جِنہاں وچ~$\eq$
نہیں آندا۔ ٹرماں توں بنا نمونہ
$!A \in \Gamma^*$ دے ہر اوہ جملے نوں پورا کردا
اے جیہدے وچ~$\eq$ نہیں؛ ایس لئی
$!A \in \Gamma$ دے سارے جملے وی پورے ہوندے
نیں۔ پر جے~$\eq$ موجود ہووے تاں ایہ نمونہ کم
نہیں کردا۔ کیوں جے فیر~$\Gamma^*$ وچ
!!a{sentence}~$\eq[t][t']$ ہو سکدا اے، جد کہ
ٹرماں توں بنے نمونے وچ ہر ٹرم دی قدر اوہ ٹرم
آپ اے۔ ایس لئی جے $t$ تے $t'$ وکھرے ٹرم
ہون، تاں نمونے وچ اوہناں دیاں قدراں، یعنی
بالترتیب $t$ تے $t'$، وی وکھریاں ہون گیاں؛
چنانچہ $\eq[t][t']$ جھوٹا ہووے گا۔ ایس
مشکل نوں اسیں ہم ارز طبقیاں راہیں ونڈ
(factoring) دی تعمیر نال دور کر سکدے آں۔
\end{explain}

\begin{defn}
  فرض کرو $\Gamma^*$، زبان~$\Lang L$ دے جملیاں
  دا بے تضاد تے !!{complete} سیٹ اے۔
  تعلق~$\approx$، زبان~$\Lang L$ دے بند
  ٹرماں دے سیٹ اُتے ایویں متعین کردے آں:
  \[
  t \approx t' \text{\quad اُدوں تے صرف اُدوں \quad} \eq[t][t'] \in \Gamma^*
  \]
\end{defn}

\begin{prop}
\ollabel{prop:approx-equiv}
تعلق~$\approx$ دیاں اگلیاں خاصیتاں نیں:
\begin{enumerate}
\item $\approx$ انعکاسی اے۔
\item $\approx$ متناظر اے۔
\item $\approx$ متعدی اے۔
\item جے $t \approx t'$، $f$ اک
  !!a{function} اے، تے $t_1$، \dots،
  $t_{i-1}$، $t_{i+1}$، \dots، $t_n$
  بند ٹرم نیں، تاں
\[
\Atom{f}{t_1,\dots, t_{i-1}, t, t_{i+1}, \dots, t_n} \approx
\Atom{f}{t_1,\dots, t_{i-1}, t', t_{i+1}, \dots, t_n}.
\]
\item جے $t \approx t'$، $R$ اک
  !!a{predicate} اے، تے $t_1$، \dots،
  $t_{i-1}$، $t_{i+1}$، \dots، $t_n$
  بند ٹرم نیں، تاں
\begin{multline*}
\Atom{R}{t_1,\dots, t_{i-1}, t, t_{i+1}, \dots, t_n} \in \Gamma^* \text{ اُدوں تے صرف اُدوں } \\
\Atom{R}{t_1,\dots, t_{i-1}, t', t_{i+1}, \dots, t_n} \in \Gamma^*.
\end{multline*}
\end{enumerate}
\end{prop}

\begin{proof}
چونکہ~$\Gamma^*$ بے تضاد تے
!!{complete} اے، ایس لئی
$\eq[t][t'] \in \Gamma^*$ اُدوں تے صرف اُدوں
اے جدوں $\Gamma^* \Proves \eq[t][t']$۔
لہٰذا اگلیاں گلاں وکھانا کافی اے:
\begin{enumerate}
\item $\Gamma^* \Proves \eq[t][t]$
  ہر بند ٹرم~$t$ لئی۔
\item جے $\Gamma^* \Proves \eq[t][t']$، تاں
  $\Gamma^* \Proves \eq[t'][t]$۔
\item جے $\Gamma^* \Proves \eq[t][t']$ تے
  $\Gamma^* \Proves \eq[t'][t'']$، تاں
  $\Gamma^* \Proves \eq[t][t'']$۔
\item جے $\Gamma^* \Proves \eq[t][t']$، تاں
% PNB-COM-005: frozen source has an extra comma after t_{i+1} in the first function term; only that comma is removed here.
\[
\Gamma^* \Proves
\eq[\Atom{f}{t_1,\dots,t_{i-1},t,t_{i+1},\dots,t_n}][\Atom{f}{t_1,\dots,t_{i-1},t',t_{i+1},\dots,t_n}]
\]
ہر $n$-مقامی !!{function}~$f$ تے بند
ٹرماں $t_1$، \dots، $t_{i-1}$،
$t_{i+1}$، \dots،~$t_n$ لئی۔
\item جے $\Gamma^* \Proves \eq[t][t']$
تے
$\Gamma^* \Proves
\Atom{R}{t_1,\dots,t_{i-1},t,t_{i+1},\dots,t_n}$،
تاں
$\Gamma^* \Proves \Atom{R}{t_1,\dots,t_{i-1},t',t_{i+1},\dots,t_n}$
ہر $n$-مقامی !!{predicate}~$R$ تے بند
ٹرماں $t_1$، \dots، $t_{i-1}$،
$t_{i+1}$، \dots،~$t_n$ لئی۔
\end{enumerate}
\end{proof}

\begin{prob}
\olref[fol][com][ide]{prop:approx-equiv} دا ثبوت مکمل کرو۔
\end{prob}

\begin{defn}
فرض کرو $\Gamma^*$، زبان~$\Lang L$ وچ اک
بے تضاد تے !!{complete} سیٹ اے، $t$ اک
بند ٹرم اے، تے $\approx$ پچھلی تعریف والا
تعلق اے۔
% PNB-COM-015: اصل Trm[L] وچ کھلے ٹرم وی نیں؛quotient-domain صرف بند ٹرماں دا اے۔
بند ٹرماں دا سیٹ ایویں وکھ لکھو:
$\mathcal{T} = \Setabs{t \in \Trm[L]}{t\text{ بند اے}}$۔ فیر:
\[
\equivrep{t}{\approx} = \Setabs{t'}{t'\in \mathcal{T}, t \approx t'}
\]
تے $\equivclass{\mathcal{T}}{\approx} = \Setabs{\equivrep{t}{\approx}}{t \in \mathcal{T}}$۔
\end{defn}

\begin{defn}
\ollabel{defn:term-model-factor}
فرض کرو $\Struct M = \Struct M(\Gamma^*)$،
\olref[mod]{defn:termmodel} والا
$\Gamma^*$ لئی ٹرماں توں بنا نمونہ اے۔
فیر $\Struct{\equivclass{M}{\approx}}$
اگلی !!{structure} اے:
\begin{enumerate}
\item $\Domain{\equivclass{M}{\approx}} = \equivclass{\mathcal{T}}{\approx}$۔
\item $\Assign{c}{\equivclass{M}{\approx}} = \equivrep{c}{\approx}$
\item $\Assign{f}{\equivclass{M}{\approx}}(\equivrep{t_1}{\approx}, \dots,
  \equivrep{t_n}{\approx}) = \equivrep{\Atom{f}{t_1,\dots, t_n}}{\approx}$
\item $\tuple{\equivrep{t_1}{\approx}, \dots, \equivrep{t_n}{\approx}} \in
  \Assign{R}{\equivclass{M}{\approx}}$ اُدوں تے
  صرف اُدوں جدوں
  $\Sat{M}{\Atom{R}{t_1,\dots, t_n}}$؛ یعنی
  اُدوں تے صرف اُدوں جدوں
  $\Atom{R}{t_1,\dots, t_n} \in \Gamma^*$۔
\end{enumerate}
\end{defn}

\begin{explain}
غور کرو کہ
$\Assign{f}{\equivclass{M}{\approx}}$
تے $\Assign{R}{\equivclass{M}{\approx}}$
دی تعریف، $\equivclass{\mathcal{T}}{\approx}$
دے عناصر نوں $\equivrep{t}{\approx}$ دی
صورت وچ لکھ کے کیتی اے؛ یعنی ہر طبقے دے
\emph{نمائندے}~$t \in \equivrep{t}{\approx}$
راہیں۔ سانوں یقینی بنانا پئے گا کہ ایہ
تعریفاں نمائندیاں دے انتخاب اُتے منحصر نہ
ہون۔ مطلب، جے کوئی ہور انتخاب~$t'$ اوہی
ہم ارزیت طبقہ دیندا اے
($\equivrep{t}{\approx} = \equivrep{t'}{\approx}$)،
تاں تعریفاں دا نتیجہ وی اوہی رہے۔ مثال
لئی جے $R$ اک جگہ والا !!{predicate}
اے، تاں تعریف دی آخری شق آکھدی اے کہ
$\equivrep{t}{\approx} \in
\Assign{R}{\equivclass{M}{\approx}}$ اُدوں
تے صرف اُدوں جدوں
% PNB-COM-006: the frozen source negates R(t) here; representative independence requires the other representative t' in this hypothetical counterexample.
$\Sat{M}{\Atom{R}{t}}$۔ پر جے کسے ہور
ٹرم~$t'$ لئی $t \approx t'$ ہووے تے
$\Sat/{M}{\Atom{R}{t'}}$، تاں تعریف
$\equivrep{t'}{\approx} \notin
\Assign{R}{\equivclass{M}{\approx}}$
منگے گی۔ جدوں $t \approx t'$ اے، تاں
$\equivrep{t}{\approx} =
\equivrep{t'}{\approx}$؛ پر
$\equivrep{t}{\approx}
\in \Assign{R}{\equivclass{M}{\approx}}$
تے $\equivrep{t}{\approx}
\notin \Assign{R}{\equivclass{M}{\approx}}$
اکّو ویلے نہیں ہو سکدے۔
\olref{prop:approx-equiv} یقین دیندا اے کہ
ایہ تضاد نہیں آ سکدا۔
\end{explain}

\begin{prop}
$\Struct{\equivclass{M}{\approx}}$
صحیح طرح متعین اے؛ یعنی جے $t_1$،
  \dots، $t_n$، $t_1'$، \dots، $t_n'$
  بند ٹرم نیں، تے $t_i \approx t_i'$
  ہووے، تاں
\begin{enumerate}
\item $\equivrep{\Atom{f}{t_1,\dots, t_n}}{\approx} =
    \equivrep{\Atom{f}{t_1',\dots, t_n'}}{\approx}$؛ یعنی
  \[
  \Atom{f}{t_1,\dots, t_n} \approx \Atom{f}{t_1',\dots, t_n'}
  \]
  تے
\item $\Sat{M}{\Atom{R}{t_1,\dots, t_n}}$
  اُدوں تے صرف اُدوں جدوں
  $\Sat{M}{\Atom{R}{t_1',\dots, t_n'}}$؛ یعنی
  \[
    \Atom{R}{t_1,\dots, t_n} \in \Gamma^* \text{ اُدوں تے صرف اُدوں }
    \Atom{R}{t_1',\dots, t_n'} \in \Gamma^*.
  \]
\end{enumerate}
\end{prop}

\begin{proof}
$n$ اُتے استقرا نال ایہ
\olref{prop:approx-equiv} توں نکلدا اے۔
\end{proof}

ٹرماں توں بنے نمونے وانگ، صدق والا
معاون قضیہ ثابت کرن توں پہلاں سانوں
اگلا معاون قضیہ چاہیدا اے۔

\begin{lem}
\ollabel{lem:val-in-termmodel-factored}
فرض کرو $\Struct M = \Struct M
(\Gamma^*)$؛ فیر ہر بند ٹرم لئی
% PNB-COM-014: factored-value statement وی بند ٹرم بارے اے۔
$\Value{t}{\equivclass{M}{\approx}} = \equivrep{t}
{\approx}$۔
\end{lem}
\begin{proof}
ایہدا ثبوت \olref[mod]{lem:val-in-termmodel}
دے ثبوت ورگا اے۔
\end{proof}

\begin{prob}
\olref[fol][com][ide]{lem:val-in-termmodel-factored}
دا ثبوت مکمل کرو۔
\end{prob}

\begin{lem}
\ollabel{lem:truth}
$\Sat{\equivclass{M}{\approx}}{!A}$
اُدوں تے صرف اُدوں جدوں $!A \in
\Gamma^*$، ہر جملے~$!A$ لئی۔
\end{lem}

\begin{proof}
$!A$ اُتے استقرا نال، ٹھیک
\olref[mod]{lem:truth} دے ثبوت وانگ۔
اکّو صورت نوں وادھو دھیان دی لوڑ اے:
جدوں $!A \ident \eq[t][t']$ ہووے۔
\begin{align*}
\Sat{\equivclass{M}{\approx}}{\eq[t][t']} & \text{ اُدوں تے صرف اُدوں } \equivrep{t}{\approx} = \equivrep{t'}{\approx}
\text{ ($\Struct{\equivclass{M}{\approx}}$ دی تعریف نال)}\\
& \text{ اُدوں تے صرف اُدوں } t \approx t' \text{ ($\equivrep{t}{\approx}$ دی تعریف نال)}\\
& \text{ اُدوں تے صرف اُدوں } \eq[t][t'] \in \Gamma^* \text{ ($\approx$ دی تعریف نال)۔}
\end{align*}
\end{proof}

\begin{digress}
غور کرو، جتھے $\Struct{M(\Gamma^*)}$
ہمیشہ !!{enumerable} تے لامحدود اے،
اُتھے $\Struct{\equivclass{M}{\approx}}$
محدود وی ہو سکدا اے: ہو سکدا اے کہ
ہم ارزیت دے طبقے
$\equivrep{t}{\approx}$ گنتی وچ
صرف محدود ہون۔ ایہ امید دے مطابق اے،
کیوں جے~$\Gamma$ وچ ایسے
!!{sentence}s ہو سکدے نیں جیہڑے ہر
اوہ !!{structure} محدود ہون دی شرط
رکھدے نیں جیہدے وچ ایہ سچے ہون۔
مثال لئی
$\lforall[x][\lforall[y][x =
y]]$ اک بے تضاد !!{sentence} اے، پر
ایہ صرف اوہناں !!{structure}s وچ
سچا اے جِنہاں دے !!a{domain} وچ
بالکل اک !!{element} ہووے۔
\end{digress}
\endgroup
% /OLP-0133

% OLP-0134 content/first-order-logic/completeness/completeness-thm.tex
\begingroup
\def\olfilename{completeness-thm}
\expandafter\def\csname ol@sectioncs\endcsname{section}

\iftag{FOL}
      {\olfileid{fol}{com}{cth}}
      {\olfileid{pl}{com}{cth}}

\olsection{تمامیت دا قضیہ}

\begin{explain}
آؤ ہُن اپنے نتائج اکٹھے کریے: ایس توں تمامیت دا قضیہ ملدا اے۔
\end{explain}

\begin{thm}[تمامیت دا قضیہ]
\ollabel{thm:completeness}
فرض کرو $\Gamma$، !!{sentence}s دا اک سیٹ اے۔
جے $\Gamma$ بے تضاد اے، تاں اوہ پورا ہون جوگا اے۔
\end{thm}

\begin{proof}
فرض کرو $\Gamma$ بے تضاد اے۔\iftag{FOL}{
\olref[hen]{lem:henkin} مطابق اک سیر شدہ،
بے تضاد سیٹ $\Gamma' \supseteq \Gamma$
موجود اے۔}{} \olref[lin]{lem:lindenbaum}
مطابق اک $\Gamma^* \supseteq
\iftag{FOL}{\Gamma'}{\Gamma}$ موجود اے
جیہڑا بے تضاد تے !!{complete} اے۔
\iftag{FOL}{
  چونکہ $\Gamma' \subseteq \Gamma^*$،
  ہر !!{formula}~$!A(x)$ لئی
  $\Gamma^*$ وچ ایس صورت دا !!a{sentence}
  موجود اے:
  \iftag{prvEx}
        {$\lexists[x][!A(x)] \lif !A(c)$}
        {$\lnot\lforall[x][!A(x)] \lif \lnot !A(c)$}
  ایس لئی $\Gamma^*$ سیر شدہ اے۔ جے
  $\Gamma$ وچ~$\eq$ نہیں، تاں}{ایس صورت وچ}
\olref[mod]{lem:truth} مطابق،
$\iftag{FOL}{\Sat{M(\Gamma^*)}}{\pSat{v(\Gamma^*)}}{!A}$
اُدوں تے صرف اُدوں جدوں $!A \in
\Gamma^*$۔ ایس توں خاص طور اُتے ہر
$!A \in \Gamma$ لئی
$\iftag{FOL}{\Sat{M(\Gamma^*)}}{\pSat{v(\Gamma^*)}}{!A}$
ملدا اے؛ لہٰذا $\Gamma$ پورا ہون جوگا اے۔
\iftag{FOL}{
  جے $\Gamma$ وچ~$\eq$ موجود اے،
  تاں \olref[ide]{lem:truth} مطابق
  سارے !!{sentence}s~$!A$ لئی
  $\Sat{\equivclass{M}{\approx}}{!A}$
  اُدوں تے صرف اُدوں جدوں $!A \in \Gamma^*$۔
  خاص طور اُتے
  $\Sat{\equivclass{M}{\approx}}{!A}$
  ہر $!A \in \Gamma$ لئی
  اے، ایس لئی $\Gamma$ پورا ہون جوگا اے۔}{}
\end{proof}

\begin{cor}[تمامیت دا قضیہ، دوجا بیان]
\ollabel{cor:completeness}
ہر $\Gamma$ تے !!{sentence}s~$!A$ وچوں ہر اک لئی:
جے $\Gamma \Entails !A$، تاں
$\Gamma \Proves !A$۔
\end{cor}

\begin{proof}
غور کرو کہ \olref{cor:completeness} تے
\olref{thm:completeness} وچ $\Gamma$
دونیں وار کلی طور اُتے مسوّر اے۔ الجھن
توں بچن لئی \olref{thm:completeness}
نوں ہور متغیر نال فیر بیان کریے:
!!{sentence}s دے ہر سیٹ~$\Delta$
لئی، جے $\Delta$ بے تضاد اے، تاں
اوہ پورا ہون جوگا اے۔ اِس دا
عکسِ نقیض ایہ اے: جے $\Delta$
پورا ہون جوگا نہیں، تاں
$\Delta$ متضاد اے۔ ایس نوں اسیں
نتیجے دا ثبوت دین لئی ورتاں گے۔

فرض کرو $\Gamma \Entails !A$۔ فیر
$\Gamma \cup \{\lnot !A\}$
\olref[syn][sem]{prop:entails-unsat}
مطابق پورا ہون جوگا نہیں۔
$\Gamma \cup \{\lnot !A\}$
نوں اپنا $\Delta$ لے کے،
\olref{thm:completeness} دے پچھلے
بیان توں ملدا اے کہ
$\Gamma \cup \{\lnot !A\}$
متضاد اے۔ ہُن
\iftag{FOL}{%
  \tagrefs{prfAX/{fol:axd:prv:prop:prov-incons},
    prfSC/{fol:seq:prv:prop:prov-incons},
    prfND/{fol:ntd:prv:prop:prov-incons},
    prfTab/{fol:tab:prv:prop:prov-incons}}}{%
  \tagrefs{prfAX/{pl:axd:prv:prop:prov-incons},
    prfSC/{pl:seq:prv:prop:prov-incons},
    prfND/{pl:ntd:prv:prop:prov-incons},
    prfTab/{pl:tab:prv:prop:prov-incons}}}
مطابق $\Gamma \Proves !A$۔
\end{proof}

\tagprob{FOL}
\begin{prob}
\olref[fol][com][cth]{cor:completeness}
ورت کے \olref[fol][com][cth]{thm:completeness}
ثابت کرو، تے ایس طرح وکھاؤ کہ تمامیت دے
قضیے دے دونیں بیان ہم ارز نیں۔
\end{prob}
\tagendprob

\tagprob{notFOL}
\begin{prob}
\olref[pl][com][cth]{cor:completeness}
ورت کے \olref[pl][com][cth]{thm:completeness}
ثابت کرو، تے ایس طرح وکھاؤ کہ تمامیت دے
قضیے دے دونیں بیان ہم ارز نیں۔
\end{prob}
\tagendprob

\tagprob{FOL}
\begin{prob}
!!a{derivation} نظام دے مکمل ہون لئی
اوہدے قاعدے اینے زورآور ہونے چاہیدے نیں
کہ ہر پورا ہون جوگ نہ ہون والا سیٹ
متضاد ثابت ہو سکے۔ تمامیت ثابت کرن
لئی !!{derivation} دے کہڑے قاعدے
لازمی سن؟ کی انہاں وچوں کوئی قاعدہ
ثبوت وچ کتھے وی نہیں ورتیا گیا؟
ایہناں سوالاں دا جواب دین لئی اک فہرست
یا خاکہ بناؤ جیہڑا وکھاوے کہ
!!{derivation} دے کہڑے قاعدے
\olref[fol][com][cth]{thm:completeness}
دے ثبوت تک لے جان والے کہڑے
نتیجیاں وچ ورتے گئے۔ اِنہاں ثبوتاں
وچ قاعدیاں دی غیر صریح ورتوں وی
ضرور نشان زد کرو۔
\end{prob}
\tagendprob

\tagprob{notFOL}
\begin{prob}
!!a{derivation} نظام دے مکمل ہون لئی
اوہدے قاعدے اینے زورآور ہونے چاہیدے نیں
کہ ہر پورا ہون جوگ نہ ہون والا سیٹ
متضاد ثابت ہو سکے۔ تمامیت ثابت کرن
لئی !!{derivation} دے کہڑے قاعدے
لازمی سن؟ کی انہاں وچوں کوئی قاعدہ
ثبوت وچ کتھے وی نہیں ورتیا گیا؟
ایہناں سوالاں دا جواب دین لئی اک فہرست
یا خاکہ بناؤ جیہڑا وکھاوے کہ
!!{derivation} دے کہڑے قاعدے
\olref[pl][com][cth]{thm:completeness}
دے ثبوت تک لے جان والے کہڑے
نتیجیاں وچ ورتے گئے۔ اِنہاں ثبوتاں
وچ قاعدیاں دی غیر صریح ورتوں وی
ضرور نشان زد کرو۔
\end{prob}
\tagendprob
\endgroup
% /OLP-0134

% OLP-0135 content/first-order-logic/completeness/compactness.tex
\begingroup
\def\olfilename{compactness}
\expandafter\def\csname ol@sectioncs\endcsname{section}

\iftag{FOL}
      {\olfileid{fol}{com}{com}}
      {\olfileid{pl}{com}{com}}

\olsection{محدود گواہی دا قضیہ}

تمامیت دے قضیے دا اک اہم نتیجہ محدود گواہی دا قضیہ اے۔
ایہ قضیہ کہندا اے کہ جے !!{sentence}s دے کسے سیٹ دا ہر
\emph{محدود} ذیلی سیٹ پورا ہون جوگا ہووے، تاں سارا سیٹ
پورا ہون جوگا اے---بھاویں اوہ سیٹ آپ لامتناہی ہووے۔
ایہ گل پہلی نظر وچ بالکل واضح نہیں۔ پہلے پہل ایہ امکان
رد ہوندا نظر نہیں آؤندا کہ !!{sentence}s دا کوئی لامتناہی
سیٹ متضاد ہووے، پر تضاد، ایویں سمجھو، صرف لامتناہی تعداد
وچوں پیدا ہوندا ہووے۔ محدود گواہی دا قضیہ ایہ امکان رد
کردا اے: !!{sentence}s دا کوئی ایسا پورا ہون جوگا نہ ہون والا
لامتناہی سیٹ نہیں جیہدا ہر محدود ذیلی سیٹ پورا ہون جوگا
ہووے۔ تمامیت دے قضیے وانگ، ایس دا منطقی نتیجے نال
متعلق وی اک بیان اے: جے !!{sentence}s دا لامتناہی سیٹ
کجھ نتیجے وچ لیاندا اے، تاں اوہ نتیجہ اوہدے کسے محدود
ذیلی سیٹ توں ہی نکلدا اے۔

\begin{defn}
  !!{formula}s دا سیٹ~$\Gamma$ \emph{محدود طور تے پورا ہون جوگا}
  اُدوں تے صرف اُدوں اے جدوں ہر محدود $\Gamma_0 \subseteq \Gamma$
  پورا ہون جوگا ہووے۔
\end{defn}

\begin{thm}[محدود گواہی دا قضیہ]
\ollabel{thm:compactness}
ہر جملیاں دے سیٹ~$\Gamma$ تے ہر جملے~$!A$ لئی
ہٹھ لکھیاں گلاں درست نیں:
\begin{enumerate}
  \item $\Gamma \Entails !A$ اُدوں تے صرف اُدوں جدوں کوئی محدود
    $\Gamma_0 \subseteq \Gamma$ موجود ہووے جیہدے لئی
    $\Gamma_0 \Entails !A$۔
  \item $\Gamma$ پورا ہون جوگا اُدوں تے صرف اُدوں اے جدوں
    اوہ محدود طور تے پورا ہون جوگا اے۔
\end{enumerate}
\end{thm}

\begin{proof}
اسیں (2) ثابت کردے آں۔ جے $\Gamma$ پورا ہون جوگا اے، تاں
\iftag{FOL}{!!a{structure}~$\Struct{M}$}{!!a{valuation}~$\pAssign{v}$}
موجود اے جیہدی لئی
$\iftag{FOL}{\Sat{M}}{\pSat{v}}{!A}$
ہر $!A \in \Gamma$ اُتے درست اے۔
ظاہر اے کہ ایہی $\iftag{FOL}{\Struct{M}}{\pAssign{v}}$
،~$\Gamma$ دے ہر محدود ذیلی سیٹ نوں وی پورا کردی اے؛
ایس لئی $\Gamma$ محدود طور تے پورا ہون جوگا اے۔

ہُن فرض کرو $\Gamma$ محدود طور تے پورا ہون جوگا اے۔
فیر ہر محدود $\Gamma_0 \subseteq \Gamma$ پورا ہون جوگا
اے۔ صحت مطابق
(\iftag{FOL}{%
  \tagrefs{prfAX/{fol:axd:sou:cor:consistency-soundness},
    prfSC/{fol:seq:sou:cor:consistency-soundness},
    prfND/{fol:ntd:sou:cor:consistency-soundness},
    prfTab/{fol:tab:sou:cor:consistency-soundness}}}{%
  \tagrefs{prfAX/{pl:axd:sou:cor:consistency-soundness},
    prfSC/{pl:seq:sou:cor:consistency-soundness},
    prfND/{pl:ntd:sou:cor:consistency-soundness},
    prfTab/{pl:tab:sou:cor:consistency-soundness}}}),
ہر محدود ذیلی سیٹ بے تضاد اے۔ فیر
\iftag{FOL}{%
  \tagrefs{prfAX/{fol:axd:ptn:prop:proves-compact},
    prfSC/{fol:seq:ptn:prop:proves-compact},
    prfND/{fol:ntd:ptn:prop:proves-compact},
    prfTab/{fol:tab:ptn:prop:proves-compact}}}{%
  \tagrefs{prfAX/{pl:axd:ptn:prop:proves-compact},
    prfSC/{pl:seq:ptn:prop:proves-compact},
    prfND/{pl:ntd:ptn:prop:proves-compact},
    prfTab/{pl:tab:ptn:prop:proves-compact}}}
مطابق $\Gamma$ آپ وی بے تضاد ہووے گا۔
تمامیت مطابق (\olref[cth]{thm:completeness})،
چونکہ $\Gamma$ بے تضاد اے، ایس لئی اوہ پورا ہون جوگا اے۔
\end{proof}

\tagprob{FOL}
\begin{prob}
\olref[fol][com][com]{thm:compactness}
دا حصہ (1) ثابت کرو۔
\end{prob}
\tagendprob

\tagprob{notFOL}
\begin{prob}
\olref[pl][com][com]{thm:compactness}
دا حصہ (1) ثابت کرو۔
\end{prob}
\tagendprob

\iftag{FOL}{%
\begin{ex}
% PNB-COM-008: The source calls equality instances sentences but Trm[L] includes variables; restrict t to closed old-language terms, as required by its own term semantics.
ہر نمونے~$\Struct{M}$ وچ، جیہڑا نظریے~$\Gamma$ دا نمونہ اے،
ہر بند ٹرم~$t$،~$\Domain{M}$ دے کسے نہ کسے
!!a{element} نوں ظاہر کردا اے۔ پر کی اسیں ایہ ضمانت
دے سکدے آں کہ $\Domain{M}$ دا ہر !!{element} وی کسے
بند ٹرم توں ظاہر ہووے؟ ہور لفظاں وچ، کی ایسے نظریے~$\Gamma$
نیں جنہاں دے سارے نمونے بند ٹرماں نال پوری طرح ڈھکے ہوئے
ہون؟ محدود گواہی دا قضیہ وکھاندا اے کہ جے $\Gamma$
دے لامتناہی نمونے ہون، تاں ایہ گل درست نہیں۔ وجہ ایہ اے:
فرض کرو $\Struct{M}$،~$\Gamma$ دا اک لامتناہی نمونہ اے،
تے $c$ اک !!a{constant} اے جیہڑا $\Gamma$ دی زبان وچ نہیں۔
% PNB-COM-016: پچھلا PNB-COM-007 غلط سی؛eq/ اصل وچ inequality اے،اضافی lnot ہٹا کے اصل polarity بحال۔
$\Delta$ نوں $\eq/[c][t]$ جملیاں دا سیٹ لو،
جتھے~$t$،~$\Lang{L}$، یعنی~$\Gamma$ دی زبان
دا بند ٹرم اے؛ یعنی
  \[
  \Delta = \Setabs{\eq/[c][t]}{t \in \Trm[L],\ t\text{ بند اے}}.
  \]
$\Gamma \cup \Delta$ دا کوئی وی محدود ذیلی سیٹ
$\Gamma' \cup \Delta'$ دے روپ وچ لکھیا جا سکدا اے،
جتھے $\Gamma' \subseteq \Gamma$ تے
$\Delta' \subseteq \Delta$۔ چونکہ $\Delta'$ محدود اے،
ایس وچ صرف محدود گنتی دے ٹرم آ سکدے نیں۔ کوئی
$a \in \Domain{M}$ لو،~$\Domain{M}$ دا ایسا
!!a{element} جیہڑا اوہناں وچوں کسے ٹرم توں ظاہر نہ
ہووے، تے $\Struct{M'}$ نوں اوہ
!!{structure} لو جیہڑی $\Struct{M}$ ورگی ہی اے،
پر جیہدے وچ $\Assign{c}{M'} = a$ وی اے۔
چونکہ $a \neq \Value{t}{M}$ ہر~$t$ لئی
جیہڑا~$\Delta'$ وچ آؤندا اے، ایس لئی
$\Sat{M'}{\Delta'}$۔ چونکہ
$\Sat{M}{\Gamma}$، $\Gamma' \subseteq
\Gamma$، تے $c$،~$\Gamma$ وچ نہیں آؤندا،
ایس لئی $\Sat{M'}{\Gamma'}$ وی۔
دونیں ملا کے،
$\Sat{M'}{\Gamma' \cup \Delta'}$
ہر محدود ذیلی سیٹ $\Gamma' \cup \Delta'$
لئی اے۔
لہٰذا $\Gamma \cup \Delta$ دا ہر محدود ذیلی
سیٹ پورا ہون جوگا اے۔ محدود گواہی مطابق،
$\Gamma \cup \Delta$ آپ وی پورا ہون جوگا
اے۔ ایس لئی $\Gamma \cup \Delta$ دے ایسے نمونے
موجود نیں جنہاں لئی
$\Sat{M}{\Gamma \cup \Delta}$۔
ایسا ہر $\Struct{M}$،~$\Gamma$ دا نمونہ
اے۔
% PNB-COM-018: نہ ڈھکے ہون دا دعویٰ اصل زبان دے reduct بارے اے؛نواں c خود نویں عنصر دا ناں اے۔
اصل زبان لئی، یعنی نئے مستقل دی تعبیر ہٹا کے، اوہ بند ٹرماں نال پوری طرح ڈھکیا ہویا
نہیں، کیونکہ $\Value{c}{M} \neq
\Value{t}{M}$ ہر بند ٹرم~$t$
لئی اے جیہڑا~$\Lang{L}$ وچ اے۔
\end{ex}

\begin{ex}
ایسی !!a{language}~$\Lang{L}$ اُتے غور کرو جیہدے
وچ !!{predicate}~$<$، !!{constant}s~$\Obj{0}$ تے
$\Obj{1}$، تے !!{function}s~$+$، $\times$ تے
$-$ ہون۔ $\Gamma$ نوں ایس !!{language} دے اوہناں
سارے !!{sentence}s دا سیٹ لو جیہڑے
!!{structure}~$\Struct{Q}$ وچ سچے نیں،
جیہدا دائرہ~$\Rat$ اے تے جیہدے نشاناں دیاں
تعبیراں معمول والیاں نیں۔ ایس طرح
$\Gamma$،~$\Lang{L}$ دے ناطق اعداد بارے
سارے سچے !!{sentence}s دا سیٹ اے۔
یقیناً~$\Rat$ وچ (بلکہ~$\Real$ وچ وی)
کوئی ایسا عدد~$r$ نہیں جیہڑا~$0$ توں وڈا
پر $1/k$ توں، ہر $k \in \PosInt$ لئی،
چھوٹا ہووے۔ ایسا عدد، جے موجود ہووے،
\emph{لامتناہی طور تے چھوٹا عدد (infinitesimal)}
ہووے گا: صفر نہیں، پر بے حد چھوٹا۔
محدود گواہی دا قضیہ ورت کے ایہ وکھایا
جا سکدا اے کہ $\Gamma$ دے ایسے نمونے
موجود نیں جنہاں وچ لامتناہی طور تے
چھوٹے عدد موجود ہون۔ ساڈی زبان وچ
تقسیم دا کوئی !!a{function} نہیں (صفر
نال تقسیم غیر معرّف اے، تے !!{function}s
دی تعبیر ہمہ جا فنکشناں نال ہونی چاہیدی
اے)۔ پر اسیں پھر وی ایہ بیان کر سکدے
آں کہ $r < 1/k$، کیونکہ ایہ اُدوں تے
صرف اُدوں اے جدوں $r \cdot k < 1$۔
ہُن $c$ اک نواں !!{constant} لو تے
$\Delta$ نوں ایہ سیٹ لو:
\[
\{\Obj{0} < c\}
\cup \Setabs{ c \times \num k < \Obj{1} }{k \in \PosInt}
\]
(جتھے $\num{k} = (\Obj{1} + (\Obj{1} + \dots + (\Obj{1} +
\Obj{1})\dots))$ وچ $k$ عدد~$\Obj{1}$ نیں)۔
$\Delta_0$ نوں~$\Delta$ دا کوئی وی محدود ذیلی سیٹ
لو۔ کوئی~$K$ موجود اے ایسا کہ سارے
!!{sentence}s~$c \times \num{k} < \Obj{1}$
جیہڑے~$\Delta_0$ وچ آؤندے نیں،
اوہناں دا $k < K$ ہووے۔ جے اسیں $\Struct{Q}$ نوں
$\Struct{Q'}$ تک ودھائیے تے
$\Assign{c}{Q'} = 1/K$ رکھیے،
تاں $\Sat{Q'}{\Gamma_0 \cup \Delta_0}$
ہر محدود $\Gamma_0 \subseteq \Gamma$ لئی؛
ایس لئی $\Gamma \cup \Delta$ محدود
طور تے پورا ہون جوگا اے (مشق:
ایہ تفصیل نال ثابت کرو)۔ محدود گواہی
مطابق $\Gamma \cup \Delta$ پورا ہون
جوگا اے۔ ہر نمونے~$\Struct{S}$ وچ،
جیہڑا~$\Gamma \cup \Delta$ دا نمونہ اے، اک
لامتناہی طور تے چھوٹا عدد، یعنی
$\Assign{c}{S}$، موجود اے۔
\end{ex}
}{}

\tagprob{FOL}
\begin{prob}
حساب دے معیاری نمونے~$\Struct{N}$
وچ کوئی !!{element}~$k \in \Domain{N}$
ایسا نہیں جیہڑا ہر فارمولے~$\num{n}
< x$ نوں پورا کرے (جتھے $\num{n}$،
$\Obj{0}^{\prime\dots\prime}$ اے،
جیہدے وچ $n$ عدد $\prime$ نیں)۔
محدود گواہی دا قضیہ ورت کے وکھاؤ کہ
حساب دی زبان دے اوہ !!{sentence}s
جیہڑے حساب دے معیاری نمونے
$\Struct{N}$ وچ سچے نیں،
!!a{structure}~$\Struct{N'}$ وچ
وی سچے نیں جیہدے وچ ایسا
!!a{element} موجود اے جیہڑا
ہر فارمولے~$\num{n} < x$
نوں \emph{پورا کردا} اے۔
\end{prob}
\tagendprob

\iftag{FOL}{
\begin{ex}
اسیں جان دے آں کہ !!{identity} والی
پہلے درجے دی منطق ایہ بیان کر سکدی
اے کہ !!{domain} دا حجم گھٹ توں گھٹ
کنا ہووے: جملہ~$!A_{\ge n}$
(جیہڑا کہندا اے ”گھٹ توں گھٹ $n$
وکھریاں چیزاں نیں“) صرف اوہناں
ساختاں وچ سچا اے جنہاں دے
$\Domain{M}$ وچ گھٹ توں گھٹ~$n$
چیزاں ہون۔ سو جے اسیں
  \[
  \Delta = \Setabs{!A_{\ge n}}{n \ge 1}
  \]
لئیے، تاں $\Delta$ دا ہر نمونہ
لامتناہی ہونا لازم اے۔ ایس طرح
کسی نظریے وچ~$\Delta$ شامل کر کے
اسیں ضمانت دے سکدے آں کہ اوہدے
صرف لامتناہی نمونے ہون: $\Gamma \cup
\Delta$ دے نمونے ٹھیک اوہی
$\Gamma$ دے لامتناہی نمونے نیں۔

ایس لئی پہلے درجے دی منطق
لامتناہیت بیان کر سکدی اے۔ پر
محدود گواہی دا قضیہ وکھاندا اے
کہ ایہ محدودیت بیان نہیں کر
سکدی۔ فرض کرو کہ جملیاں
دا کوئی سیٹ~$\Lambda$ ٹھیک
اوہناں ساریاں تے صرف اوہناں
!!{structure}s وچ پورا ہووے
جیہڑیاں محدود نیں۔ فیر
$\Delta \cup \Lambda$
محدود طور تے پورا ہون جوگا
اے۔ کیوں؟ فرض کرو
$\Delta' \cup \Lambda' \subseteq \Delta \cup
\Lambda$ محدود اے، جتھے
$\Delta' \subseteq \Delta$ تے
$\Lambda' \subseteq \Lambda$۔
% PNB-COM-009: The source's maximum does not exist when Delta-prime is empty; in that case take the bound to be one.
جے ایہ پہلا ذیلی سیٹ خالی نہ ہووے،
تاں $n$ اوہ سب توں وڈا عدد
لو جیہدے لئی $!A_{\ge n}
\in \Delta'$؛ جے اوہ خالی
ہووے تاں ایہ حد اک رکھ لو۔
چونکہ $\Lambda$ ہر محدود
ساخت وچ پورا ہوندا اے،
ایس دا اک نمونہ~$\Struct{M}$
اے جیہدے !!{element}s
محدود تعداد وچ، پر $\ge n$
نیں۔ فیر
$\Sat{M}{\Delta' \cup \Lambda'}$۔
محدود گواہی مطابق
$\Delta \cup \Lambda$
دا اک لامتناہی نمونہ اے،
جیہڑا ایس مفروضے نال
ٹکراؤندا اے کہ $\Lambda$
صرف محدود !!{structure}s
وچ پورا ہوندا اے۔
\end{ex}
}{}
\endgroup
% /OLP-0135

% OLP-0136 content/first-order-logic/completeness/compactness-direct.tex
\begingroup
\def\olfilename{compactness-direct}
\expandafter\def\csname ol@sectioncs\endcsname{section}

\iftag{FOL}
      {\olfileid{fol}{com}{cpd}}
      {\olfileid{pl}{com}{cpd}}

\olsection{محدود گواہی دے قضیے دا سِدھا ثبوت}

تمامیت دے قضیے نوں ورتے بغیر، اوہدے ثبوت دے ای
خیالاں نال محدود گواہی دا قضیہ سِدھا ثابت
کیتا جا سکدا اے۔ تمامیت دے قضیے دے ثبوت وچ
اسیں !!{sentence}s دے بے تضاد سیٹ~$\Gamma$
توں شروع کیتا، اوہنوں اک بے تضاد\iftag{FOL}{،
  سیر شدہ،}{} تے !!{complete}
!!{sentence}s دے سیٹ~$\Gamma^*$ تک ودھایا،
تے فیر وکھایا کہ
% PNB-COM-019: برابری والی صورت وچ نحوی ٹرماں دی بجائے ہم ارز طبقیاں دا نمونہ۔
\iftag{FOL}{ٹرم
  نمونے~$\Struct{M(\Gamma^*)}$، یا برابری ہون دی صورت وچ اوہدے
  ہم ارز طبقیاں دے نمونے}{!!{valuation}~$\pAssign{v(\Gamma^*)}$}
وچ، جیہنوں~$\Gamma^*$ توں بنایا گیا سی،
~$\Gamma$ دے سارے !!{sentence}s
سچے نیں؛ ایس لئی~$\Gamma$ پورا ہون جوگا اے۔

ایہی طریقہ ورت کے اسیں وکھا سکدے آں کہ جملیاں
دا کوئی محدود طور تے پورا ہون جوگا سیٹ، پورا
ہون جوگا اے۔ بس صدق دے معاون قضیے تک لے جان والے نتائج
دے اوہ بیان ثابت کرنے پین گے جنہاں وچ
”بے تضاد“ دی تھاں ”محدود طور تے پورا
ہون جوگا“ آ جاندا اے۔

\begin{prop}
\ollabel{prop:fsat-ccs}
فرض کرو $\Gamma$،~!!{complete} تے محدود طور تے
پورا ہون جوگا اے۔ فیر:
\begin{enumerate}
\tagitem{prvAnd}{$(!A \land !B) \in \Gamma$
  اُدوں تے صرف اُدوں جدوں $!A \in \Gamma$
  تے $!B \in \Gamma$ دونیں ہون۔}{}

\tagitem{prvOr}{$(!A \lor !B) \in \Gamma$
  اُدوں تے صرف اُدوں جدوں یا $!A \in \Gamma$
  ہووے یا $!B \in \Gamma$ ہووے۔}{}

\tagitem{prvIf}{$(!A \lif !B) \in \Gamma$
  اُدوں تے صرف اُدوں جدوں یا $!A \notin \Gamma$
  ہووے یا $!B \in \Gamma$ ہووے۔}{}
\end{enumerate}
\end{prop}

\tagprob{FOL}
\begin{prob}
\olref[fol][com][cpd]{prop:fsat-ccs}
ثابت کرو۔ $\Proves$ نوں نہ ورتو۔
\end{prob}
\tagendprob

\tagprob{notFOL}
\begin{prob}
\olref[pl][com][cpd]{prop:fsat-ccs}
ثابت کرو۔ $\Proves$ نوں نہ ورتو۔
\end{prob}
\tagendprob

\iftag{FOL}{%
\begin{lem}
\ollabel{lem:fsat-henkin}
ہر محدود طور تے پورا ہون جوگے
سیٹ~$\Gamma$ نوں اک سیر شدہ،
محدود طور تے پورا ہون جوگے
سیٹ~$\Gamma'$ تک ودھایا جا سکدا اے۔
\end{lem}
}{}

\tagprob{FOL}
\begin{prob}
\olref[fol][com][cpd]{lem:fsat-henkin}
ثابت کرو۔ (اشارہ: اہم قدم ایہ وکھانا اے کہ جے
$\Gamma_n$ محدود طور تے پورا ہون جوگا اے،
تاں $\Gamma_n \cup \{!D_n\}$ وی
محدود طور تے پورا ہون جوگا اے؛
!!{derivation}s یا بے تضادی دا سہارا نہ لو۔)
\end{prob}
\tagendprob

\iftag{FOL}{
\begin{prop}\ollabel{prop:fsat-instances}
فرض کرو $\Gamma$ مکمل، محدود طور تے
پورا ہون جوگا تے سیر شدہ اے۔
\begin{tagenumerate}{prvEx,prvAll}
\tagitem{prvEx}{$\lexists[x][!A(x)] \in \Gamma$
  اُدوں تے صرف اُدوں جدوں $!A(t) \in \Gamma$
  ہووے کسے اک بند ٹرم~$t$ لئی۔}{}
\tagitem{prvAll}{$\lforall[x][!A(x)] \in \Gamma$
  اُدوں تے صرف اُدوں جدوں $!A(t) \in \Gamma$
  ہووے ہر بند ٹرم~$t$ لئی۔}{}
\end{tagenumerate}
\end{prop}
}{}

\tagprob{FOL}
\begin{prob}
\olref[fol][com][cpd]{prop:fsat-instances}
ثابت کرو۔
\end{prob}
\tagendprob

\begin{lem}
\ollabel{lem:fsat-lindenbaum}
ہر محدود طور تے پورا ہون جوگے
سیٹ~$\Gamma$ نوں !!a{complete} تے
محدود طور تے پورا ہون جوگے
سیٹ~$\Gamma^*$ تک ودھایا جا سکدا اے۔
\end{lem}

\tagprob{FOL}
\begin{prob}
\olref[fol][com][cpd]{lem:fsat-lindenbaum}
ثابت کرو۔ (اشارہ: اہم قدم ایہ وکھانا اے
کہ جے $\Gamma_n$ محدود طور تے پورا
ہون جوگا اے، تاں یا
$\Gamma_n \cup \{!A_n\}$ یا
$\Gamma_n \cup \{\lnot !A_n\}$
محدود طور تے پورا ہون جوگا اے۔)
\end{prob}
\tagendprob

\tagprob{notFOL}
\begin{prob}
\olref[pl][com][cpd]{lem:fsat-lindenbaum}
ثابت کرو۔ (اشارہ: اہم قدم ایہ وکھانا اے
کہ جے $\Gamma_n$ محدود طور تے پورا
ہون جوگا اے، تاں یا
$\Gamma_n \cup \{!A_n\}$ یا
$\Gamma_n \cup \{\lnot !A_n\}$
محدود طور تے پورا ہون جوگا اے۔)
\end{prob}
\tagendprob

\begin{thm}[محدود گواہی]
\ollabel{thm:compactness-direct}
$\Gamma$ پورا ہون جوگا اُدوں تے
صرف اُدوں اے جدوں اوہ محدود
طور تے پورا ہون جوگا ہووے۔
\end{thm}

\begin{proof}
جے $\Gamma$ پورا ہون جوگا اے، تاں
\iftag{FOL}{!!a{structure}~$\Struct{M}$}{!!a{valuation}~$\pAssign{v}$}
موجود اے جیہدی لئی
$\iftag{FOL}{\Sat{M}}{\pSat{v}}{!A}$
ہووے، ہر $!A \in \Gamma$ لئی۔
ظاہر اے کہ ایہی
\iftag{FOL}{$\Struct{M}$}{$\pAssign{v}$}
،~$\Gamma$ دے ہر محدود ذیلی سیٹ
نوں وی پورا کردی اے؛ ایس لئی
$\Gamma$ محدود طور تے پورا ہون جوگا اے۔

ہُن فرض کرو $\Gamma$ محدود طور تے پورا
ہون جوگا اے۔\iftag{FOL}{ \olref{lem:fsat-henkin}
  مطابق اک محدود طور تے پورا ہون جوگا، سیر شدہ
  سیٹ $\Gamma' \supseteq \Gamma$ موجود اے۔}{}
\olref{lem:fsat-lindenbaum} مطابق،
$\iftag{FOL}{\Gamma'}{\Gamma}$
نوں اک !!a{complete} تے محدود طور تے
پورا ہون جوگے سیٹ~$\Gamma^*$ تک ودھایا
جا سکدا اے\iftag{FOL}{، تے $\Gamma^*$ ہن وی
  سیر شدہ اے}{}۔
\olref[mod]{defn:termmodel} وانگ
\iftag{FOL}{ٹرم
  نمونہ~$\Struct{M(\Gamma^*)}$}{!!{valuation}~$\pAssign{v(\Gamma^*)}$}
بناؤ۔\iftag{FOL}{ غور کرو کہ
\olref[mod]{prop:quant-termmodel} اس گل اُتے
منحصر نہیں سی کہ $\Gamma^*$ بے تضاد اے
(یا، ایس حوالے نال، !!{complete} یا سیر شدہ
اے)، بلکہ صرف ایس گل اُتے کہ
$\Struct{M(\Gamma^*)}$ پوری طرح بند
ٹرماں نال ڈھکیا ہویا اے۔}{}
% PNB-COM-010: دونوں شاخاں وچ finite-satisfiability proposition دا حوالہ صاف کیتا۔
صدق دے معاون قضیے دا ثبوت (\olref[mod]{lem:truth})
اوہدوں وی چلدا اے جے اسیں
\olref[ccs]{prop:ccs} دے حوالے دی تھاں
\olref{prop:fsat-ccs} دا حوالہ رکھ دیئے
\iftag{FOL}{، تے
\olref[hen]{prop:saturated-instances}
دے حوالے دی تھاں \olref{prop:fsat-instances}
دا حوالہ رکھ دیئے۔
% PNB-COM-019: برابری موجود ہووے تاں بند ٹرماں دا quotient ضروری؛finite satisfaction congruence دی بنیاد اے۔
برابری موجود ہووے تاں \olref[ide]{defn:term-model-factor}
والا طبقیاں دا نمونہ ورتو،نہ الگ الگ نحوی ٹرماں والا نمونہ۔
برابری دی انعکاسیت،تناظر،تعدیت تے فنکشن/محمول دی بدل پذیری
دے ہر قدم لئی صرف محدود برابریاں درکار نیں؛اوہناں دے معنوی
نتیجے مکمل سیٹ وچ شامل ہوندے نیں۔ ایس لئی تعبیر نمائندے دے
انتخاب توں آزاد اے تے \olref[ide]{lem:truth} دی برابری والی
صورت وی صدق دے استقرائی ثبوت وچ چل دی اے}{}۔
% PNB-COM-019: صدق دے معاون قضیے دے اخذی جواز دی تھاں محدود معنوی بندش دا مکمل سبب۔
ایتھے کسی اخذی قاعدے دا سہارا نہیں: مکمل،محدود طور تے پورا ہون
جوگے سیٹ دے کسی محدود ذیلی سیٹ دا ہر معنوی نتیجہ اُس وڈے سیٹ
وچ شامل اے۔ ورنہ مکملیت نال اوہدی نفی شامل ہووے گی،تے اوہ
محدود ذیلی سیٹ نال نفی ملا کے اک محدود،پورا نہ ہون جوگا ذیلی
سیٹ بن جائے گا،جو مفروضے دے خلاف اے۔ خاص طور اُتے کذب دا
مستقل شامل نہیں ہو سکدا،صدق دا مستقل شامل ہوندا اے،تے ہر
جملے دی نفی اُدوں تے صرف اُدوں شامل اے جدوں اوہ جملہ شامل
نہ ہووے۔ ایہ محدود معنوی جواز اصل صدق دے معاون قضیے وچ اخذ تے بے تضادی
والیاں وجوہاں دی تھاں ورتو۔
\end{proof}

\tagprob{FOL}
\begin{prob}
صدق دے معاون قضیے دے اوہ روپ دا پورا ثبوت
لکھو (\olref[fol][com][mod]{lem:truth})
جیہڑا \olref[fol][com][cpd]{thm:compactness-direct}
دے ثبوت لئی درکار اے۔
\end{prob}
\tagendprob

\tagprob{notFOL}
\begin{prob}
صدق دے معاون قضیے دے اوہ روپ دا پورا ثبوت
لکھو (\olref[pl][com][mod]{lem:truth})
جیہڑا \olref[pl][com][cpd]{thm:compactness-direct}
دے ثبوت لئی درکار اے۔
\end{prob}
\tagendprob
\endgroup
% /OLP-0136

% OLP-0137 content/first-order-logic/completeness/downward-ls.tex
\begingroup
\def\olfilename{downward-ls}
\expandafter\def\csname ol@sectioncs\endcsname{section}

\olfileid{fol}{com}{dls}
\olsection{لووینہائم--سکولم دا قضیہ}

لووینہائم--سکولم دا قضیہ کہندا اے کہ جے کسے نظریے
دا لامتناہی نمونہ اے، تاں اوہدا اک ایہو جیہا نمونہ
وی اے جیہڑا زیادہ توں زیادہ !!{denumerable} اے۔
ایس گل دا فوری نتیجہ ایہ اے کہ پہلے درجے دی منطق
ایہ بیان نہیں کر سکدی کہ کسے !!a{structure}
دا سائز !!{nonenumerable} اے: کوئی وی
!!{sentence} یا !!{sentence}s دا سیٹ جیہڑا
ساریاں !!{nonenumerable} !!{structure}s وچ پورا
ہوندا اے، کسے !!{enumerable} ساخت وچ وی پورا
ہوندا اے۔

\begin{thm}
\ollabel{thm:downward-ls}
جے $\Gamma$ بے تضاد اے، تاں اوہدا
!!a{enumerable} نمونہ اے؛ یعنی اوہ اک ایسی
!!a{structure} وچ پورا ہون جوگا اے جیہدا
دائرہ یا محدود اے یا !!{denumerable}۔
\end{thm}

\begin{proof}
جے $\Gamma$ بے تضاد اے، تاں تمامیت دے قضیے
دے ثبوت توں ملی !!{structure}~$\Struct M$
% PNB-COM-020: تعمیر Henkin دی ودھی زبان دے بند ٹرماں یا اوہناں دے طبقیاں اُتے اے۔
دا دائرہ $\Domain{M}$ ودھی زبان~$\Lang L'$
دے بند ٹرماں دے سیٹ توں وڈا نہیں؛طبقیاں دی گنتی وی نمائندیاں
دی گنتی توں وڈّی نہیں۔ اصل شمار جوگ علامات وچ شمار جوگ نویاں
مستقل علامتاں شامل کرن نال ایہ بند ٹرماں دا سیٹ شمار جوگ ہی
رہندا اے۔ ایس لئی
$\Struct M$ زیادہ توں زیادہ !!{denumerable}
اے۔
\end{proof}

\begin{thm}
\ollabel{noidentity-ls}
جے $\Gamma$ برابری دے بغیر پہلے درجے دی
منطق دی زبان وچ !!{sentence}s دا اک بے تضاد
سیٹ اے، تاں اوہدا !!a{denumerable} نمونہ اے؛
یعنی اوہ اک ایسی !!a{structure} وچ پورا ہون
جوگا اے جیہدا دائرہ لامتناہی تے
!!{enumerable} اے۔
\end{thm}

\begin{proof}
جے $\Gamma$ بے تضاد اے تے اوہدے وچ کوئی
ایہو جیہا جملہ نہیں جیہدے وچ برابری آوے،
تاں تمامیت دے قضیے دے ثبوت توں ملی
!!{structure}~$\Struct M$ دا دائرہ
% PNB-COM-011: ٹرم نمونے دا اصل دائرہ کھلے ٹرماں دی بجائے بند ٹرم نیں۔
$\Domain{M}$ زبان~$\Lang L'$ دے بند
ٹرماں دے سیٹ دے عین برابر اے۔ سو
$\Struct{M}$ !!{denumerable} اے، کیونکہ
$\Trm[L']$ گنتی وار فہرست بنن جوگ اے،
تے اوس پھیلی زبان وچ شامل لامتناہی نویاں مستقل
علامتاں بند ٹرماں دے سیٹ نوں لامتناہی
بناندیاں نیں۔
\end{proof}

\begin{ex}[سکولم دا پیراڈوکس]
زرمیلو--فرینکل سیٹ تھیوری~$\Log{ZFC}$
اک بہت طاقت ور ڈھانچہ اے جیہدے وچ
تقریباً سارے ریاضیاتی بیان کیتے جا سکدے
نیں، سیٹاں دے سائز بارے حقیقتاں وی۔
مثلاً $\Log{ZFC}$ ایہ ثابت کر سکدا اے
کہ حقیقی اعداد دا سیٹ~$\Real$
!!{nonenumerable} اے؛ ایہ کانتور دا
قضیہ وی ثابت کر سکدا اے کہ کسے وی
سیٹ دے ذیلی سیٹاں دا سیٹ اوس سیٹ
توں وڈا ہوندا اے، وغیرہ۔ جے
$\Log{ZFC}$ بے تضاد اے، تاں اوہدے
سارے نمونے لامتناہی نیں۔ ہور ایہ کہ
اوہناں ساریاں وچ ایہوجے !!{element}s
وی نیں جنہاں بارے نظریہ کہندا اے کہ
اوہ !!{nonenumerable} نیں؛ مثال لئی
اوہ عنصر جیہڑا قدرتی اعداد دے ذیلی
سیٹاں دے سیٹ دے وجود بارے
~$\Log{ZFC}$ دے قضیے نوں سچا کردا
اے۔ لووینہائم--سکولم دے قضیے مطابق
$\Log{ZFC}$ دے !!{enumerable}
نمونے وی نیں---ایہوجے نمونے جنہاں
وچ ”!!{nonenumerable}“ سیٹ ہوندے
نیں، حالانکہ نمونے آپ !!{enumerable}
ہوندے نیں۔
\end{ex}
\endgroup
% /OLP-0137

% OLP-0644 content/first-order-logic/completeness/maximally-consistent-sets.tex
\begingroup
\def\olfilename{maximally-consistent-sets}
\expandafter\def\csname ol@sectioncs\endcsname{section}

\olfileid{fol}{com}{mcs}
\olsection{\usetoken{P}{sentence} دے زیادہ توں زیادہ بے تضاد سیٹ}

\begin{defn}[زیادہ توں زیادہ بے تضاد سیٹ]
\ollabel{mcs}
!!{sentence}s دا سیٹ~$\Gamma$ اُدوں تے صرف اُدوں \emph{زیادہ توں زیادہ بے تضاد} اے جدوں
\begin{enumerate}
\item $\Gamma$ بے تضاد اے، تے
\item جے $\Gamma \subsetneq \Gamma'$ ہووے، تاں $\Gamma'$ متضاد اے۔
\end{enumerate}
\end{defn}

\begin{explain}
اُتلی تعریف دے برابر اک ہور تعریف ایہ اے۔ جملیاں دا سیٹ $\Gamma$ اُدوں تے صرف اُدوں \emph{زیادہ توں زیادہ بے تضاد} اے جدوں
\begin{enumerate}
\item $\Gamma$ بے تضاد اے، تے
\item جے $\Gamma \cup \{ !A \}$ بے تضاد اے، تاں $!A \in \Gamma$۔
\end{enumerate}
دوجے لفظاں وچ، جیہڑے !!{sentence}s پہلاں ای~$\Gamma$ وچ نہیں، اوہناں نوں زیادہ توں زیادہ بے تضاد سیٹ~$\Gamma$ وچ شامل کرن نال بنیا وڈا سیٹ متضاد ہو جاندا اے۔
\end{explain}

\begin{explain}
زیادہ توں زیادہ بے تضاد سیٹ تکمیل دے ثبوت وچ اہم نیں، کیوں جے اسیں ضمانت دے سکدے آں کہ !!{sentence}s دا ہر بے تضاد سیٹ~$\Gamma$ کسے زیادہ توں زیادہ بے تضاد سیٹ~$\Gamma^*$ وچ شامل اے، تے ہر !!{sentence}~$!A$ لئی زیادہ توں زیادہ بے تضاد سیٹ وچ یا $!A$ اے یا اوہدا انکار $\lnot !A$ اے۔ ایہ خاص کر جوہری !!{sentence}s لئی سچ اے، سو !!{constant}s نال ڈھکدے طریقے نال ودھائی زبان دے زیادہ توں زیادہ بے تضاد سیٹ توں اسیں اک !!a{structure} بنا سکدے آں جیہدے وچ !!{predicate}s دی تعبیر ایس حساب نال دِتّی جاندی اے کہ کیہڑے جوہری !!{sentence}s~$\Gamma^*$ وچ نیں۔ فیر وکھایا جا سکدا اے کہ ایہ !!{structure}وچلے سارے !!{sentence}s~$\Gamma^*$ وچ (تے ایس لئی $\Gamma$ وچ وی) سچ کردی اے۔ ایس آخری گل دے ثبوت لئی چاہیدا اے کہ $\lnot !A \in
\Gamma^*$ اُدوں تے صرف اُدوں جدوں $!A \notin \Gamma^*$، $(!A \lor !B) \in \Gamma^*$ اُدوں تے صرف اُدوں جدوں $!A
\in \Gamma^*$ یا $!B \in \Gamma^*$، تے ایس طرح اگے۔
\end{explain}

\begin{prop}
\ollabel{prop:mcs}
منّ لو $\Gamma$ زیادہ توں زیادہ بے تضاد اے۔ فیر:
\begin{enumerate}
\item \ollabel{prop:mcs-prov-in} جے $\Gamma \Proves !A$، تاں $!A \in
  \Gamma$۔

\item \ollabel{prop:mcs-either-or} ہر $!A$ لئی، یا $!A \in
  \Gamma$ یا $\lnot !A \in \Gamma$۔

\tagitem{prvAnd}{\ollabel{prop:mcs-and} $(!A \land !B) \in \Gamma$
  اُدوں تے صرف اُدوں جدوں دوویں $!A \in \Gamma$ تے $!B \in \Gamma$۔}{}

\tagitem{prvOr}{\ollabel{prop:mcs-or} $(!A \lor !B) \in \Gamma$ اُدوں تے صرف اُدوں جدوں
  یا $!A \in \Gamma$ یا $!B \in \Gamma$۔}{}

\tagitem{prvIf}{\ollabel{prop:mcs-if} $(!A \lif !B) \in \Gamma$ اُدوں تے صرف اُدوں جدوں
  یا $!A \notin \Gamma$ یا $!B \in \Gamma$۔}{}
\end{enumerate}
\end{prop}

\begin{proof}
ایہناں ساریاں گلاں لئی اسیں منّ لینے آں کہ $\Gamma$ زیادہ توں زیادہ بے تضاد اے۔
\begin{enumerate}
\item جے $\Gamma \Proves !A$، تاں $!A \in \Gamma$۔

منّ لو $\Gamma \Proves !A$۔ اُلٹا فرض کرو کہ $!A
\notin \Gamma$۔ فیر چونکہ $\Gamma$ زیادہ توں زیادہ بے تضاد اے، $\Gamma
\cup \{!A\}$ متضاد اے، سو $\Gamma \cup \{!A\} \Proves
\lfalse$۔ \tagrefs{prfSC/{fol:seq:prv:prop:provability-contr},prfND/{fol:ntd:prv:prop:provability-contr}} مطابق
$\Gamma$ متضاد اے۔ ایہ اوس فرض دے اُلٹ اے کہ $\Gamma$ بے تضاد اے۔ سو ایہ نہیں ہو سکدا کہ $!A \notin
\Gamma$، ایس لئی $!A \in \Gamma$۔

\item ہر $!A$ لئی، یا $!A \in \Gamma$ یا $\lnot !A \in \Gamma$۔

اُلٹا فرض کرو کہ کسے~$!A$ لئی دوویں $!A \notin \Gamma$ تے $\lnot !A \notin \Gamma$۔ چونکہ $\Gamma$ زیادہ توں زیادہ بے تضاد اے، $\Gamma \cup \{!A\}$ تے $\Gamma \cup \{\lnot !A\}$ دوویں متضاد نیں، سو $\Gamma \cup \{!A\} \Proves \lfalse$ تے $\Gamma \cup
\{\lnot !A\} \Proves \lfalse$۔ \tagrefs{prfSC/{fol:seq:prv:prop:provability-exhaustive},prfND/{fol:ntd:prv:prop:provability-exhaustive}} مطابق
$\Gamma$ متضاد اے، سو تضاد نکلدا اے۔ ایس لئی کوئی ایہو جیہا !!a{sentence}~$!A$ نہیں ہو سکدا تے ہر $!A$ لئی، یا $!A \in \Gamma$ یا $\lnot !A
\in \Gamma$۔

\tagitem{defAnd}{}{%
\iftag{probAnd}{مشق۔}{%
$(!A \land !B) \in \Gamma$ اُدوں تے صرف اُدوں جدوں دوویں $!A \in \Gamma$ تے $!B \in \Gamma$۔

سدھے پاسے لئی منّ لو $(!A \land !B) \in \Gamma$۔ فیر
$\Gamma \Proves !A \land !B$۔ \tagrefs{prfSC/{fol:seq:ppr:prop:provability-land-left},prfND/{fol:ntd:ppr:prop:provability-land-left}} مطابق
$\Gamma \Proves !A$ تے $\Gamma \Proves !B$۔ \olref{prop:mcs-prov-in} مطابق $!A \in \Gamma$ تے $!B \in \Gamma$، جیویں چاہیدا سی۔

اُلٹے پاسے لئی منّ لو $!A \in \Gamma$ تے $!B \in
\Gamma$۔ فیر $\Gamma \Proves !A$ تے $\Gamma \Proves !B$۔ \tagrefs{prfSC/{fol:seq:ppr:prop:provability-land-right},prfND/{fol:ntd:ppr:prop:provability-land-right}} مطابق
$\Gamma \Proves !A \land !B$۔ \olref{prop:mcs-prov-in} مطابق $(!A \land
!B) \in \Gamma$۔}}

\tagitem{defOr}{}{%
\iftag{probOr}{مشق۔}{%
$(!A \lor !B) \in \Gamma$ اُدوں تے صرف اُدوں جدوں یا $!A \in \Gamma$ یا $!B \in \Gamma$۔

سدھے پاسے دی اُلٹی نفی لئی منّ لو $!A
\notin \Gamma$ تے $!B \notin \Gamma$۔ اسیں وکھاؤنا اے کہ $(!A \lor
!B) \notin \Gamma$۔ چونکہ $\Gamma$ زیادہ توں زیادہ بے تضاد اے، $\Gamma
\cup \{!A\} \Proves \lfalse$ تے $\Gamma \cup \{!B\} \Proves \lfalse$۔
\tagrefs{prfSC/{fol:seq:ppr:prop:provability-lor},prfND/{fol:ntd:ppr:prop:provability-lor}} مطابق
$\Gamma \cup \{(!A \lor !B)\}$ متضاد اے۔ سو $(!A \lor !B)
\notin \Gamma$، جیویں چاہیدا سی۔

اُلٹے پاسے لئی منّ لو $!A \in \Gamma$ یا $!B \in
\Gamma$۔ فیر $\Gamma \Proves !A$ یا $\Gamma \Proves !B$۔ \tagrefs{prfSC/{fol:seq:ppr:prop:provability-lor},prfND/{fol:ntd:ppr:prop:provability-lor}} مطابق
$\Gamma \Proves !A \lor !B$۔ \olref{prop:mcs-prov-in} مطابق $(!A \lor
!B) \in \Gamma$، جیویں چاہیدا سی۔}}

\tagitem{defIf}{}{%
\iftag{probIf}{مشق۔}{%
$(!A \lif !B) \in \Gamma$ اُدوں تے صرف اُدوں جدوں یا $!A \notin \Gamma$ یا $!B \in \Gamma$۔

سدھے پاسے لئی منّ لو $(!A \lif !B) \in \Gamma$، تے اُلٹا فرض کرو کہ $!A \in \Gamma$ تے $!B \notin \Gamma$۔ ایہناں فرضاں اُتے $\Gamma \Proves !A \lif !B$ تے $\Gamma \Proves !A$۔ \tagrefs{prfSC/{fol:seq:ppr:prop:provability-lif-left},prfND/{fol:ntd:ppr:prop:provability-lif-left}} مطابق
$\Gamma \Proves !B$۔ پر فیر \olref{prop:mcs-prov-in} مطابق $!B \in
\Gamma$، جیہڑا اوس فرض دے اُلٹ اے کہ $!B \notin \Gamma$۔

اُلٹے پاسے لئی پہلاں اوہ صورت ویکھو جتھے $!A \notin
\Gamma$۔ \olref{prop:mcs-either-or} مطابق $\lnot !A \in \Gamma$ تے سو $\Gamma \Proves \lnot !A$۔ \tagrefs{prfSC/{fol:seq:ppr:prop:provability-lif-right},prfND/{fol:ntd:ppr:prop:provability-lif-right}} مطابق
$\Gamma \Proves !A \lif !B$۔ \olref{prop:mcs-prov-in} فیر لا کے سانوں $(!A \lif !B) \in \Gamma$ ملدا اے، جیویں چاہیدا سی۔

ہُن اوہ صورت ویکھو جتھے $!B \in \Gamma$۔ فیر $\Gamma \Proves !B$ تے \tagrefs{prfSC/{fol:seq:ppr:prop:provability-lif-right},prfND/{fol:ntd:ppr:prop:provability-lif-right}} مطابق
$\Gamma \Proves !A \lif !B$۔ \olref{prop:mcs-prov-in} مطابق $(!A \lif
!B) \in \Gamma$۔}}
\end{enumerate}
\end{proof}

\begin{prob}
\olref[fol][com][mcs]{prop:mcs} دا ثبوت پورا کرو۔
\end{prob}
\endgroup
% /OLP-0644
\OLEndChapterHook
\endgroup
% /OLP-0126

% OLP-0174 content/first-order-logic/beyond/beyond.tex
\begingroup
\def\olfilename{beyond}
\expandafter\def\csname ol@sectioncs\endcsname{section}

\olchapter{fol}{byd}{پہلے درجے دی منطق توں اگے}

\begin{editorial}
  ایہ باب، Jeremy Avigad دے منطق والے نوٹاں توں ڈھالیا گیا،
  ہور منطقی نظاماں دی نہایت مختصر جھلک پیش کردا اے۔ ایہ
  ایداں دے کورس لئی اگے مطالعے دے موضوع تجویز کرن والا باب
  اے جیہڑا انہاں موضوعاں نوں خود نہیں پڑھاندا۔ ایتھے ذکر
  کیتے ہر موضوع نوں---امید اے---آخرکار اوپن لاجک پراجیکٹ
  وچ اپنے وکھرے حصے دی سطح اُتے تفصیلی جگہ ملے گی۔
\end{editorial}

% OLP-0175 content/first-order-logic/beyond/introduction.tex
\begingroup
\def\olfilename{introduction}
\expandafter\def\csname ol@sectioncs\endcsname{section}

\olfileid{fol}{byd}{int}

\olsection{جائزہ}

پہلے درجے دی منطق دلچسپی دا واحد منطقی نظام نہیں:
ایہدیاں کئی توسیعاں تے وکھرے روپ نیں۔ کسے منطق وچ
عام طور اُتے زبان دی رسمی وضاحت ہوندی اے؛ عموماً،
پر ہمیشہ نہیں، اک استنباطی نظام ہوندا اے؛ تے عموماً،
پر ہمیشہ نہیں، اک مقصود معنویات ہوندی اے۔ پر اصطلاح
”منطق“ دی ایس فنی ورتوں نال اک واضح سوال اٹھدا اے:
جیہڑیاں منطقاں پہلے درجے دی منطق نہیں، انہاں دا
”منطق“ لفظ نال کیہڑا رشتہ اے جدوں ایہ لفظ عام فہم
یا فلسفیانہ معنی وچ ورتیا جاوے؟ ہیٹھاں بیان کیتے سارے
نظام کسے نہ کسے طرح دی استدلالی سوچ نوں نمونہ بنان
لئی تیار کیتے گئے نیں؛ کی اسیں دس سکدے آں کہ انہاں
نوں منطقی کیہڑی گل بناؤندی اے؟

ایہدے آسان جواب سامنے نہیں آندے۔ ”منطق“ لفظ وکھ وکھ
طریقیاں تے سیاقاں وچ ورتیا جاندا اے، تے ”سچائی“ وانگ
ایس تصور دا وی کئی فلسفیانہ نقطۂ نظراں توں تجزیہ
ہویا اے۔ مثال لئی کوئی منطقی استدلال دا مقصد ایہ طے
کرنا سمجھ سکدا اے کہ کیہڑے بیان لازماً سچے نیں، تجربے
توں آزاد سچے نیں، غیر منطقی اصطلاحاں دی تشریح توں
آزاد سچے نیں، اپنی صورت دی وجہ توں سچے نیں، یا
لسانی رواج دی وجہ توں سچے نیں؛ تے انہاں ہر اک
تصور نوں ہور خاصی وضاحت دی لوڑ اے۔ جے توجہ صرف
ریاضی وچ ورتی جاندی منطق تک وی محدود رکھیے، تاں
وی ایہدی حد بارے بہت گھٹ اتفاق اے۔ مثال لئی
\textit{Principia Mathematica} وچ رسل تے وائٹ ہیڈ
نے فریگے توں شروع ہوئی {\em منطقیت} دی روایت وچ
ریاضی نوں منطق دی بنیاد اُتے تشکیل دین دی کوشش
کیتی۔ انہاں دا منطقی نظام اعلیٰ قسم دی منطق دا
اک روپ سی، جیہڑا ہیٹھاں بیان کیتے نظام نال ملدا
جلدا اے۔ آخرکار انہاں نوں ایداں دے مسلّمے شامل
کرن پئے جیہڑے اکثر معیاراں مطابق خالص منطقی نہیں
لگدے (خاص طور اُتے لامتناہیت دا مسلّمہ تے
تقلیل پذیری دا مسلّمہ)؛ پر فیر وی کوئی ایہ موقف
رکھ سکدا اے کہ اعلیٰ درجے دے استدلال دیاں کجھ
صورتاں منطقی منّیاں جانیاں چاہیدیاں نیں۔ ایس دے
برعکس، کوائن، جیہدی وجودیات ”قضیہاں“
(propositions) نوں گفتگو دیاں جائز شےواں نہیں
منّدی، دلیل دیندا اے کہ دوجے درجے دی تے
اعلیٰ درجے دی منطق اصل وچ سیٹ نظریے دیاں
ای صورتاں نیں، بس منطق دا بھیس پایا ہویا اے؛
دوجے لفظاں وچ، محمولاں اُتے مسوّریت والے نظام
خالص منطقی نہیں۔

فی الحال ایداں دے فلسفیانہ سوال کسے ہور ویلے
لئی چھڈ دینا بہتر اے، تے ہیٹھلے نظاماں نوں بس
استدلال دیاں وکھ وکھ قسماں دی رسمی مثالی
صورتاں سمجھنا اے، چاہے اوہ منطقی ہون یا نہ ہون۔
\endgroup
% /OLP-0175

% OLP-0176 content/first-order-logic/beyond/many-sorted-logic.tex
\begingroup
\def\olfilename{many-sorted-logic}
\expandafter\def\csname ol@sectioncs\endcsname{section}

\olfileid{fol}{byd}{msl}

\olsection{کئی قسماں والی منطق}

پہلے درجے دی منطق وچ متغیر تے مسوّر اک ہی
!!{domain} دے عناصر اُتے چلدے نیں۔ پر اکثر کئی
وکھ وکھ، اک دوجے توں جدا !!{domain}s رکھنا فائدہ مند
ہوندا اے: مثال لئی اعداد دا !!a{domain}،
ہندسی شےواں دا !!a{domain}، اعداد توں اعداد ول جان
والے فنکشناں دا !!a{domain}، ابیلی گروپاں دا
!!a{domain}، تے ایس طرح ہور۔

کئی قسماں والی منطق ایس کم لئی اک ڈھانچا دیندی اے۔
شروع وچ ”قسماں“ (sorts) دی فہرست لئی جاندی اے: کسے شے
دی ”قسم“ (sort) ایہ دسدی اے کہ اوہ کیہڑے
”!!{domain}“ وچ آندی اے۔ فیر ہر قسم لئی وکھرے
!!{variable}s تے مسوّر ہوندے نیں، تے (عموماً) ہر
قسم لئی اپنی !!{identity} وی ہوندی اے۔ فنکشناں تے
تعلقاں دے دلیل والے مقام وی قسماں دے حساب نال
مقرر ہوندے نیں، یعنی ہر مقام اُتے کیہڑی قسم دی شے
آ سکدی اے۔ ایس توں علاوہ پہلے درجے دی منطق دے
معمول دے قاعدے برقرار رہندے نیں؛ مسوّراں دے
قاعدیاں دی ہر قسم لئی اپنی صورت دہرائی جاندی اے۔

مثال لئی بین الاقوامی تعلقاں دے مطالعے واسطے اسیں
ایسی زبان چن سکدے آں جیہدیاں شےواں دیاں دو قسماں
ہون: فرانسیسی شہری تے جرمن شہری۔ پہلی قسم دیاں
شےواں لئی اک اک-جگہی تعلق ”انگوری شراب پیندا اے“ رکھ
سکدے آں؛ دوجی قسم دیاں شےواں لئی دوجا اک-جگہی
تعلق ”ورسٹ کھاندا اے“؛ تے دو-جگہی تعلق ”دو
ملکاں توں بنیا ویاہتا جوڑا بناندے نیں“، جیہدے
دو دلیل نیں: پہلی دلیل پہلی قسم دی، تے دوجی
دوجی قسم دی۔ جے متغیر $a$, $b$, $c$ فرانسیسی
شہریاں اُتے تے $x$, $y$, $z$ جرمن شہریاں اُتے
چلدے نیں، تاں
% PNB-SYN-009: اندرلے مسوّر دے متغیر تے ماتحت فارمولے نوں اوہناں دے اپنے اختیاری آرگومنٹاں وچ رکھیا۔
\[
\lforall[a][\lforall[x][(\Atom{\Obj{MarriedTo}}{a,x} \lif
(\Atom{\Obj{DrinksWine}}{a} \lor \lnot \Atom{\Obj{EatsWurst}}{x})]]
\]
ایہ بیان کردا اے کہ جے کوئی فرانسیسی شخص کسے
جرمن نال ویاہیا اے، تاں یا اوہ فرانسیسی شخص
انگوری شراب پیندا اے یا اوہ جرمن شخص ورسٹ نہیں کھاندا۔

کئی قسماں والی منطق نوں قدرتی طریقے نال پہلے درجے
دی منطق وچ سمویا جا سکدا اے: کئی قسماں دے
!!{domain}s دیاں ساریاں شےواں نوں اک پہلے درجے دے
!!{domain} وچ اکٹھا کر دیو، قسماں دی پہچان رکھن
لئی اک-جگہی !!{predicate}s ورتو، تے مسوّراں دا
دائرہ متعلقہ قسم تک محدود کرو۔ مثال لئی اوپرلی
مثال نال ملدی پہلے درجے دی زبان وچ پہلے بیان
کیتے تعلقاں دے نال اک-جگہی !!{predicate}s
”$\Obj{German}$“ تے ”$\Obj{French}$“ وی ہون گے،
پر تعلقاں دی قسم-بندی دیاں شرطاں علامتاں توں
ہٹا دتیاں جان گیاں۔ اک قسم-بند مسوّر
$\lforall[x][!A]$ دا ترجمہ، جتھے $x$ جرمن
قسم دا !!a{variable} اے، ایہ بنے گا:
\[
\lforall[x][(\Atom{\Obj{German}}{x} \lif !A)].
\]
اسیں ایداں دے مسلّمے شامل کرن ہون گے جیہڑے
قسماں دے اک دوجے توں جدا ہون دی ضمانت دین،
مثلاً $\lforall[x][\lnot (\Atom{\Obj{German}}{x} \land
  \Atom{\Obj{French}}{x})]$؛ نال ای ایداں دے
مسلّمے وی جیہڑے ضمانت دین کہ ”انگوری شراب پیندا اے“
صرف انہاں شےواں لئی درست ہووے جیہڑیاں
$\Atom{\Obj{French}}{x}$ محمول پورا کردیاں نیں،
وغیرہ۔
% PNB-SEM-028: معمول دے غیر خالی قسمی دائرے، ایس دو-قسمی مثال دی مکمل ڈھکائی، تے مستقل/فنکشن دی قسم-بندی واضح کیتی۔
ایتھے معمول دے قاعدیاں لئی ہر قسم دا دائرہ غیر خالی
منیا گیا اے۔ ایس دو-قسمی مثال وچ اکٹھے دائرے دا ہر
عنصر فرانسیسی یا جرمن قسم وچوں کسے اک دا ہونا چاہیدا
اے۔ جے زبان وچ مستقل یا فنکشن ہون، تاں مستقل اپنی
مقرر قسم وچ ہوون، تے فنکشن مقرر انپٹ قسماں دے عناصر
نوں اپنی مقرر آؤٹ پٹ قسم دا عنصر دین۔ فنکشناں دیاں
مقرر انپٹ قسماں توں باہر دیاں قدراں ایس ترجمے وچ
بےاثر نیں۔ انہاں رواجاں تے مسلّمیاں نال ایہ وکھانا
اوکھا نہیں کہ کئی قسماں والے !!{sentence}s دا
ترجمہ پہلے درجے دے !!{sentence}s وچ ہوندا اے،
تے کئی قسماں والے !!{derivation}s دا ترجمہ
پہلے درجے دے !!{derivation}s وچ۔ ایس طرح کئی
قسماں والیاں !!{structure}s دا وی متعلقہ پہلے
درجے دیاں !!{structure}s وچ ”ترجمہ“ ہوندا اے،
تے الٹ پاسے وی؛ ایس لئی کئی قسماں والی منطق
لئی تمامیت دا قضیہ وی مل جاندا اے۔
\endgroup
% /OLP-0176

% OLP-0177 content/first-order-logic/beyond/second-order-logic.tex
\begingroup
\def\olfilename{second-order-logic}
\expandafter\def\csname ol@sectioncs\endcsname{section}

\olfileid{fol}{byd}{sol}

\olsection{دوجے درجے دی منطق}

دوجے درجے دی منطق دی زبان وچ صرف افراد دے
!!a{domain} اُتے ہی نہیں، ایس !!{domain} اُتے
بنے تعلقاں اُتے وی مسوّریت ہو سکدی اے۔ پہلے درجے
دی دتی ہوئی زبان $\Lang{L}$ لئی، ہر $k$ واسطے
ایدے وچ !!{variable}s $R$ شامل کیتے جاندے نیں
جیہڑے $k$-جگہی تعلقاں اُتے چل دے نیں، تے انہاں
متغیراں اُتے مسوّریت دی اجازت ہوندی اے۔ جے $R$
اک $k$-جگہی تعلق دا !!a{variable} ہوئے، تے
$t_1$, \dots,~$t_k$ عام (پہلے درجے دیاں)
اصطلاحاں ہون، تاں
$\Atom{R}{t_1,\dots,t_k}$ اک جوہری !!{formula}
اے۔ ہور !!{formula}s دا مجموعہ پہلے درجے دی
منطق وانگ ہی متعین ہوندا اے، بس دوجے درجے دی
مسوّریت لئی وادھو شقّاں ہوندیاں نیں۔ چیتا رہے
کہ !!{identity} صرف پہلے درجے دیاں اصطلاحاں
لئی اے: جے $R$ تے $S$ اکّو مقامیت~$k$ دے
تعلق-!!{variable}s ہون، تاں
$\eq[R][S]$ نوں ایس گل دا مختصر نشان بنا
سکدے آں:
\[
\lforall[x_1 \dots][\lforall[x_k][(\Atom{R}{x_1, \dots, x_k} \liff
  \Atom{S}{x_1, \dots, x_k})]].
\]

دوجے درجے دی منطق دے قاعدے صرف مسوّراں دے
قاعدیاں نوں نویں دوجے درجے دے متغیراں تک
ودھاندے نیں۔ پر ایتھے ذرا دھیان نال سمجھانا
پیندا اے کہ ایہ متغیر~$\Lang{L}$ دے
!!{predicate}s نال، تے عمومی طور اُتے
~$\Lang{L}$ دے !!{formula}s نال، کیویں
مل کے کم کردے نیں۔ گھٹ توں گھٹ تعلقی متغیر
اصطلاحاں ونجوں گنے جاندے نیں، ایس لئی
ایس صورت دے استنباط ہوندے نیں:
\[
!A(R) \Proves \lexists[R][!A(R)]
\]
پر جے $\Lang{L}$ حساب دی زبان ہوئے جیہدے وچ
اک مستقل تعلقی علامت~$<$ ہوئے، تاں ایہ
استنباط وی درست ہون دی توقع ہوئے گی:
\[
x < y \Proves \lexists[R][\Atom{R}{x,y}]
\]
یا دتے ہوئے !!{formula}~$!A$ لئی:
\[
!A(x_1, \dots, x_k) \Proves \lexists[R][\Atom{R}{x_1,\dots,x_k}]
\]
ہور عام طور اُتے اسیں ایس صورت دے استنباط
دی اجازت چاہ سکدے آں:
\[
\Subst{!A}{\lambd[\vec x][!B(\vec x)]}{R} \Proves
\lexists[R][!A]
\]
جتھے $\Subst{!A}{\lambd[\vec x][!B(\vec x)]}{R}$
دا مطلب اے کہ
$\Obj{R}{t_1,\dots,t_k}$ صورت دے ہر جوہری
!!{formula} نوں،~$!A$ وچ جتھے ایس تعلقی متغیر دی آزاد آمد ہووے،
$!B(t_1, \dots, t_k)$ نال بدل
دتا جاوے۔
% PNB-SEM-029: بدل صرف آزاد تعلقی آمد لئی اے؛ پابند ناں تازہ رکھ کے پیرامیٹر گرفت توں بچایا، مشق وچ آزاد آمد واضح کیتی۔
ایہ بدل پابند متغیراں دے نام ضرورت مطابق تازہ رکھ کے
کیتی جاندی اے، تاں جو فارمولے دے آزاد پیرامیٹر غلطی
نال پابند نہ ہو جان؛ عام اصطلاحاں دی بدل وی ایس
گرفت توں بچن والی شرط دے تابع اے۔ ایہ پچھلا قاعدہ {\em خاصیت توں
تعلق بنان دے خاکے} دے برابر اے، یعنی ایس صورت
دا اک مسلّمہ:
\[
\lexists[R][\lforall[x_1, \dots, x_k][(!A(x_1, \dots, x_k) \liff
\Atom{R}{x_1, \dots, x_k})]],
\]
دوجے درجے دی زبان دے ہر !!{formula}~$!A$
لئی اک، جیہدے وچ $R$ آزاد متغیر نہ ہوئے۔
(مشق: وکھاؤ کہ جے $R$ نوں~$!A$ وچ آزاد آون دی
اجازت دتی جاوے، تاں ایہ خاکہ غیر موافق
ہو جاندا اے!)

منطق دان جدوں ”دوجے درجے دی منطق دے
مسلّمے“ آکھدے نیں، تاں عموماً انہاں دی مراد
پہلے درجے دی منطق دے مسوّر-قاعدیاں نوں
دوجے درجے تک گھٹ توں گھٹ ودھانا، تے نال
خاصیت توں تعلق بنان دا خاکہ شامل کرنا ہوندا
اے۔ پر ایدے مسلّمیاں تے قاعدیاں دے ہور
کمزور ذیلی نظاماں دا مطالعہ وی دلچسپ
ہوندا اے۔ مثال لئی چیتا رہے کہ اپنی پوری
عمومیت وچ ایہ خاکہ \emph{غیر اسنادی}
(impredicative) اے: ایہ
$\Atom{R}{x_1, \dots, x_k}$ تعلق دی موجودگی
دا دعویٰ کرن دی اجازت دیندا اے جیہڑا دوجے
درجے دے مسوّراں والے !!a{formula} نال
”متعین“ ہوندا اے؛ تے ایہ مسوّر ایداں دے
سارے تعلقاں دے سیٹ اُتے چل دے نیں---جیہدے
وچ خود $R$ وی شامل اے!{} ویہویں صدی دے
موڑ دے نیڑے رسل دے پیراڈوکس اُتے اک عام
ردّعمل ایہ سی کہ قصور ایداں دیاں تعریفاں
دا اے، تے ریاضی دیاں بنیاداں بناندے ہوئیاں
انہاں توں بچیا جاوے۔ جے
!!{formula}~$!A$ وچ دوجے درجے دے مسوّراں
دی ورتوں منع کر دتی جاوے، تاں خاصیت
توں تعلق بنان دا اک {\em اسنادی}
(predicative) روپ ملدا اے، جیہڑا کجھ
کمزور اے۔

معنیات دے لحاظ نال دوجے درجے دی
!!{structure} نوں ایس طرح سمجھیا جا سکدا
اے: زبان لئی اک پہلے درجے دی
!!{structure}، تے نال ایہو جہے تعلقاں دا
اک سیٹ جیہڑے ایس دے !!{domain} اُتے بنے
ہون تے جیہناں اُتے دوجے درجے دے مسوّر
چل دے ہون (زیادہ عین طور اُتے، ہر $k$
لئی $k$-جگہی تعلقاں دا اک سیٹ)۔
ظاہر اے، جے خاصیت توں تعلق بنان دا خاکہ
!!{derivation} نظام وچ شامل ہوئے، تاں وادھو
شرط ایہ اے کہ ”دوجے درجے دے حصے“ وچ
اینے تعلق ہون کہ اوہ مسلّمے پورے ہون؛
ورنہ !!{derivation} نظام درست نہیں رہندا!{}
کافی تعلق مہیا کرن دا اک آسان طریقہ ایہ
اے کہ دوجے درجے دے حصے وچ پہلے درجے دے
حصے اُتے بن سکدے \emph{سارے} تعلق رکھے
جان۔ ایسی !!a{structure} نوں
\emph{مکمل} (full) آکھدے نیں، تے اک معنی
وچ ایہی زبان دی ”مقصود !!{structure}“
اے۔ جے توجہ صرف مکمل !!{structure}s تک
محدود رکھیے، تاں ایہنوں دوجے درجے دی
\emph{مکمل} معنیات آکھیا جاندا اے۔
ایس صورت وچ !!{structure} متعین کرنا
اصل وچ پہلے درجے دا حصہ متعین کرنا ہی
اے، کیوں جے دوجے درجے دے حصے دے عناصر
اوتھوں آپے متعین ہو جاندے نیں۔

مختصر ایہ کہ دوجے درجے دی منطق دی گل
کردیاں کجھ ابہام رہندا اے۔ !!{derivation}
نظام دے لحاظ نال دو مراداں ہو سکدیاں نیں:
\begin{enumerate}
\item اک ”گھٹ توں گھٹ“ دوجے درجے دا
  !!{derivation} نظام، جیہدے نال خاصیت
  توں تعلق بنان دے کجھ مسلّمے ہون۔
\item ”معیاری“ دوجے درجے دا
  !!{derivation} نظام، جیہدے وچ
  خاصیت توں تعلق بنان دا پورا خاکہ ہوئے۔
\end{enumerate}
معنیات دے لحاظ نال وی دو مراداں نیں:
\begin{enumerate}
\item ”کمزور“ معنیات، جتھے !!a{structure}
  وچ پہلے درجے دا حصہ تے نال دوجے درجے
  دا اینا وڈا حصہ ہوئے جو خاصیت توں
  تعلق بنان دے مسلّمے پورے کرے۔
\item ”معیاری“ دوجے درجے دی معنیات،
  جتھے صرف مکمل !!{structure}s وچاریاں
  جاندیاں نیں۔
\end{enumerate}
جے منطق دان اپنی مراد والا
!!{derivation} نظام یا معنیات واضح نہ کرن،
تاں عموماً ہر فہرست دی دوجی شے مراد
ہوندی اے۔ ایس معنیات دا فائدہ ایہ اے
کہ، جداں اسیں ویکھاں گے، ایہ بہت ساریاں
عام ریاضیاتی ساختاں دا
تماثل تک یکتا بیان دیندی اے؛
نال ای !!{derivation} نظام کافی طاقتور
تے ایس معنیات لئی درست اے۔ نقصان ایہ اے
کہ !!{derivation} نظام ایس معنیات لئی
\emph{تمام} نہیں؛ اصل وچ مکمل دوجے درجے
دی معنیات لئی \emph{کوئی وی} مؤثر طریقے
نال دتا !!{derivation} نظام تمام نہیں۔
دوجے پاسے، اسیں ویکھاں گے کہ
!!{derivation} نظام کمزور معنیات لئی
\emph{تمام} اے؛ یعنی جے کوئی جملہ
ثابت نہیں ہو سکدا، تاں \emph{کوئی نہ
کوئی} !!{structure}، ضروری نہیں کہ
مکمل ہوئے، ایسی اے جیہدے وچ اوہ
جھوٹا اے۔

دوجے درجے دی منطق دی زبان کافی
اظہار-طاقت رکھدی اے۔ اک-جگہی تعلقاں
نوں !!{domain} دے ذیلی سیٹاں نال
پہچانیا جا سکدا اے؛ ایس لئی انہاں
سیٹاں اُتے وی مسوّریت ہو سکدی اے۔
مثال لئی قدرتی اعداد دا استقرا اک
ہی مسلّمے نال بیان کیتا جا سکدا اے:
\[
\lforall[R][((\Atom{R}{\Obj{0}} \land \lforall[x][(\Atom{R}{x} \lif
    \Atom{R}{x'})]) \lif \lforall[x][\Atom{R}{x}])].
\]
جے حساب دی زبان دیاں علامتاں
$\Obj 0, \Obj \prime, +,
\times$ تے $<$ ہون، تاں انہاں دے
رویے لئی ہیٹھلے مسلّمے شامل
کیتے جا سکدے نیں:
\begin{enumerate}
\item $\lforall[x][\lnot x' = \Obj 0]$
% PNB-SYN-010: اصل دا غیر مقرر s اوپر دتی جانشین علامت دی بجائے آ گیا سی؛ اوہی پرائم علامت ورتی۔
\item $\lforall[x][\lforall[y][(x' = y' \lif x = y)]]$
\item $\lforall[x][(x + \Obj 0) = x]$
\item $\lforall[x][\lforall[y][(x + y') = (x + y)']]$
\item $\lforall[x][(x \times \Obj 0) = \Obj 0]$
\item $\lforall[x][\lforall[y][(x \times y') = ((x \times y) + x)]]$
\item $\lforall[x][\lforall[y][(x < y \liff \lexists[z][y = (x + z')])]]$
\end{enumerate}
ایہ وکھانا اوکھا نہیں کہ ایہ مسلّمے،
اوپرلے استقرا دے مسلّمے نال مل کے،
حساب دے معیاری نمونے
!!{structure}~$\Struct{N}$ دا
تماثل تک یکتا بیان دیندے نیں،
بشرطیکہ مکمل دوجے درجے دی معنیات
ورتی جا رہی ہوئے۔ ایداں دی کوئی وی
!!{structure}~$\Struct{M}$ لو جیہدے
وچ ایہ مسلّمے سچے نیں۔ اک
فنکشن~$f$، $\Nat$ توں
$\Struct{M}$ دے !!{domain} تک،
$\Nat$ اُتے عام بازگشت نال
ایس طرح متعین کرو کہ
$f(0) = \Assign{\Obj 0}{M}$ تے
$f(x+1) = \Assign{\prime}{M}(f(x))$
ہوئے۔ $\Nat$ اُتے عام استقرا نال،
تے ایس گل نال کہ مسلّمے (1)
تے~(2)~$\Struct M$ وچ پورے ہوندے
نیں، ویکھدے آں کہ $f$ !!{injective}
اے۔ $f$ دے !!{surjective} ہون
لئی~$P$ نوں~$\Domain{M}$ دے انہاں
عناصر دا سیٹ لو جیہڑے~$f$ دی
حاصل قدراں وچ نیں۔ چونکہ
$\Struct M$ مکمل اے، $P$ دوجے درجے
دے !!{domain} وچ شامل اے۔
~$f$ دی بناوٹ توں پتا اے کہ
$\Assign{\Obj
  0}{M}$ سیٹ~$P$ وچ اے، تے
$P$،~$\Assign{\prime}{M}$ دے
عمل ہیٹھ بند اے۔ استقرا دا مسلّمہ
$\Struct M$ وچ (خاص طور اُتے~$P$
لئی) پورا ہون دی وجہ توں
$P$~پہلے درجے دے پورے !!{domain}
دے برابر اے، جیہڑا~$\Struct M$ دا
دائرہ اے۔ ایس نال $f$ دا
!!a{bijection} ہون ثابت ہوندا اے۔
$f$ دا ہومومارفزم ہون وکھانا وی
اوداں ای آسان اے: $\Nat$ اُتے
عام استقرا وار وار ورتو۔

سیٹ نظریے دی زبان وچ فنکشن تعلق
دی اک خاص قسم ہوندا اے؛ مثال لئی
اک-جگہی فنکشن~$f$ نوں ایداں دے
دو-جگہی تعلق~$R$ نال پہچانیا
جا سکدا اے جیہڑا
$\lforall[x][\lexists![y][R(x,y)]]$
پورا کرے۔ ایس لئی فنکشناں اُتے
وی مسوّریت ہو سکدی اے۔ مکمل
معنیات ورت کے لامتناہی
!!{structure}s دی جماعت نوں
انہاں !!{structure}s $\Struct M$
دی جماعت ونجوں متعین کیتا جا سکدا
اے جیہناں لئی $\Struct M$ دے
!!{domain} توں ایدے اپنے کسے
حقیقی ذیلی سیٹ ول اک
انجیکٹو فنکشن ہوئے:
\[
\lexists[f][(\lforall[x][\lforall[y][(\eq[f(x)][f(y)] \lif \eq[x][y])]]
  \land \lexists[y][\lforall[x][\eq/[f(x)][y]]])].
\]
ایس جملے دی نفی فیر متناہی
!!{structure}s دی جماعت متعین
کردی اے۔
% PNB-SEM-030: اپنے حقیقی ذیلی سیٹ نال برابر سائز دا بیان ہر لامتناہی دائرے لئی آکھن وچ معمول دا انتخابی مفروضہ ورتیا گیا اے۔
ایتھے ہر لامتناہی سیٹ وچ اک گننے جوگا لامتناہی
ذیلی سیٹ ہون دا معمول دا مفروضہ وی ورتیا گیا اے،
مثال لئی انتخاب دے مسلّمے والے سیٹ نظریاتی پس منظر
وچ۔ ایس مفروضے توں بغیر اوپرلا فارمولا بالخصوص
اپنے حقیقی ذیلی سیٹ نال برابر سائز رکھن والے دائرے
بیان کردا اے؛ ایہدی نفی نوں صرف متناہی دائرے ہون
دا عین بیان نہیں آکھ سکدے۔

ایس توں علاوہ، خطی ترتیب دی
تعریف وچ ہیٹھلا اضافہ کرکے
حُسنِ ترتیب والیاں ساختاں دی
جماعت متعین کیتی جا سکدی اے:
\[
\lforall[P][(\lexists[x][\Atom{P}{x}] \lif \lexists[x][(\Atom{P}{x}
    \land \lforall[y][(y < x \lif \lnot \Atom{P}{y})])])].
\]
ایہ آکھدا اے کہ ہر غیر خالی
سیٹ دا سب توں چھوٹا عنصر ہوندا
اے، جے ”سیٹ“ نوں ”اک-جگہی
تعلق“ نال اک من لیا جاوے۔ اک
ہور مثال وچ گرافاں دی جڑت دا
تصور ایس طرح بیان ہو سکدا اے
کہ راساں نوں دو غیر خالی
حصیاں وچ ایداں وکھ نہیں
کیتا جا سکدا کہ انہاں
وچکار کوئی کنارہ نہ ہوئے۔ ایتھے گراف دے کنارے
بغیر سمت دے منے گئے نیں، یعنی دو راساں دا تعلق
دوواں پاسیاں توں اکّو اے:
% PNB-SEM-031: گراف دی عام جڑت لئی بغیر سمت والا تعلق واضح کیتا؛ سمت والے گراف دی کمزور جڑت دا دعویٰ نہیں۔
\[
\lnot \lexists[A][(\lexists[x][A(x)] \land \lexists[y][\lnot A(y)]
  \land \lforall[w][\lforall[z][((\Atom{A}{w} \land \lnot \Atom{A}{z})
      \lif \lnot \Atom{R}{w,z})]])].
\]
ہور مشق ونجوں تسی انہاں متناہی
!!{structure}s دی جماعت متعین کرن
دی کوشش کر سکدے او جیہناں دے
!!{domain} دا سائز جفت اے۔
ایدے توں وی ودھ نمایاں گل ایہ
اے کہ ناطق اعداد نوں شامل کرن
والے مکمل ترتیب دار میدان ونجوں
حقیقی اعداد دا
تماثل تک یکتا بیان دتا
جا سکدا اے۔

مختصر ایہ کہ دوجے درجے دی
منطق پہلے درجے دی منطق نالوں
کیتے ودھ اظہار-طاقت رکھدی
اے۔ ایہ چنگی خبر اے؛ ہن بری
خبر: اسیں دس چکے آں کہ
مکمل دوجے درجے دی معنیات
لئی کوئی مؤثر !!{derivation}
نظام تمام نہیں۔ چاہے ایہ
چنگا لگے یا برا، پہلے درجے
دی منطق دیاں کئی خاصیتاں
ایتھے نہیں ملتیاں، جنہاں وچ
محدود گواہی (compactness) تے
L\"owenheim--Skolem دے
قضیے وی شامل نیں۔

دوجے پاسے، جے کوئی مکمل
دوجے درجے دی معنیات چھڈ
کے کمزور معنیات منّن نوں
تیار ہوئے، تاں گھٹ توں
گھٹ دوجے درجے دا
!!{derivation} نظام ایس
معنیات لئی تمام اے۔
دوجے لفظاں وچ جے
$\Proves$ نوں ”گھٹ توں گھٹ
نظام وچ ثابت کردا اے“
تے $\Entails$ نوں
”کمزور معنیات وچ منطقی
نتیجہ اے“ پڑھیا جاوے،
تاں جدوں وی
$\Gamma \Entails !A$،
اوہدوں
$\Gamma \Proves !A$
وکھایا جا سکدا اے۔ جے
!!{derivation} نظام وچ
خاص خاص خاصیت توں تعلق
بنان دے مسلّمے شامل
کرنے ہون، تاں معنیات
نوں انہاں مسلّمیاں پورے
کرن والیاں دوجے درجے
دیاں ساختاں تک محدود
کرنا پیندا اے: مثال
لئی جے $\Delta$ ایداں
دے مسلّمیاں دا سیٹ
ہوئے (ممکن اے سارے
ای ہون)، تاں
$\Gamma
\cup \Delta \Entails !A$
توں
$\Gamma \cup \Delta \Proves !A$
ملدا اے۔ خاص طور
اُتے، جے $!A$ انہاں
خاصیت توں تعلق بنان
دے مسلّمیاں نال
ثابت نہ ہو سکدا
ہوئے، تاں
$\lnot !A$ دا اک
نمونہ موجود اے
جیہدے وچ ایہ
مسلّمے فیر وی
پورے ہوندے نیں۔

کمزور معنیات لئی
تمامیت دا قضیہ
کیوں اے، ایہ ویکھن
دا سب توں آسان
طریقہ دوجے درجے
دی منطق نوں کئی
قسماں والی منطق
سمجھنا اے۔ اک قسم
عام ”پہلے درجے“
دے !!{domain}
ونجوں سمجھی جاندی
اے، تے فیر ہر $k$
لئی ”$k$-جگہی
تعلقاں“ دا
!!a{domain}
ہوندا اے۔ زبان
وچ پہلے توں
شامل تعلقی
علامتاں
”$\Atom{\Obj{true}_k}{R,x_1,\dots,x_k}$“
رکھدے آں؛ ایس
دا مطلب ایہ
اے کہ $R$،
$x_1$, \dots,~$x_k$
لئی پورا
اتردا اے،
جتھے $R$
”$k$-جگہی تعلق“
قسم دا متغیر
اے تے
$x_1$, \dots,~$x_k$
پہلے درجے
دی قسم دیاں
شےواں نیں۔

ایس پہچان
نال کمزور
دوجے درجے
دی معنیات
اصل وچ کئی
قسماں والی
منطق دی عام
معنیات ہی
اے؛ تے اسیں
پہلاں ویکھ
چکے آں کہ
کئی قسماں
والی منطق
پہلے درجے
دی منطق وچ
سموئی جا
سکدی اے۔
دوواں پاسیاں
دے ترجمیاں
نوں پیش نظر
رکھ کے، دوجے
درجے دی منطق
دا کمزور تصور
اصل وچ پہلے
درجے دی منطق
دا ہی اک
بھیس بدلیا
روپ اے،
جتھے
!!{domain}
وچ ”شےواں“
تے ”تعلق“
دوویں ہوندے
نیں تے مناسب
مسلّمے انہاں
اُتے لاگو
ہوندے نیں۔
\endgroup
% /OLP-0177

% OLP-0178 content/first-order-logic/beyond/higher-order-logic.tex
\begingroup
\def\olfilename{higher-order-logic}
\expandafter\def\csname ol@sectioncs\endcsname{section}

\olfileid{fol}{byd}{hol}

\olsection{اعلیٰ درجے دی منطق}

پہلے درجے دی منطق توں دوجے درجے دی منطق
ول آ کے رسمی زبان دے اندر پہلے درجے دے
!!{domain} دیاں شےواں دے سیٹاں بارے
گل کرنا ممکن ہویا۔ فیر ایتھے ہی کیوں
رکیے؟ مثال لئی تیجے درجے دی منطق نال
شےواں دے سیٹاں دے سیٹ وچارے جا سکدے
نیں، یا شاید ایسے سیٹ وی جیہناں وچ
شےواں تے شےواں دے سیٹ دوویں شامل ہون۔
چوتھے درجے دی منطق فیر ایس قسم دیاں
شےواں دے سیٹاں بارے گل کرن دی
اجازت دیوے گی۔ جداں تسی شاید
اندازہ لا لیا اے، ایہ خیال
من چاہے درجے تک دہرایا جا سکدا اے۔

عمل وچ اعلیٰ درجے دی منطق دی
!!{formula} بندی اکثر تعلقاں دی بجائے فنکشناں
دے ذریعے کیتی جاندی اے۔
(قدرتی مطابقتاں نوں ملحوظ رکھ کے
ایہ فرق بنیادی نہیں۔) کجھ بنیادی
”قسماں“ $A$, $B$, $C$,~\dots
دتی ہوون (جیہناں نوں ہن ”انواع“
یا types آکھاں گے)، تاں اسیں
ہیٹھلی شرط نال نویاں انواع بنا
سکدے آں:
\begin{quote}
جے $\sigma$ تے $\tau$ محدود
تشکیل والیاں انواع ہون،
تاں $\sigma \to \tau$ وی ایسی
نوع اے۔
\end{quote}
انواع نوں نحوی ”نشان“ سمجھو
جیہڑے ساڈے !!{domain} دیاں
شےواں دی قسم بندی کردے نیں؛
$\sigma \to \tau$ انہاں شےواں
نوں بیان کردا اے جیہڑیاں نوع~$\sigma$
دیاں شےواں توں نوع~$\tau$ دیاں
شےواں ول فنکشن نیں۔ مثال لئی
سچائی دیاں قدراں ”سچ“ تے ”جھوٹ“
دی اک نوع $\Omega$، تے قدرتی
اعداد دی اک نوع $\Nat$ رکھ
سکدے آں۔ فیر نوع $\Nat \to
\Omega$ دیاں شےواں نوں اک-جگہی
تعلق یا $\Nat$ دے ذیلی سیٹ
سمجھیا جا سکدا اے؛ نوع
$\Nat \to \Nat$ دیاں شےواں
قدرتی اعداد توں قدرتی اعداد
ول فنکشن نیں؛ تے نوع
$(\Nat \to \Nat) \to \Nat$
دیاں شےواں ”فنکشنل“
نیں، یعنی اعلیٰ نوع دے
اوہ فنکشن جیہڑے فنکشناں نوں
اعداد ول بھیج دے نیں۔

دوجے درجے دی منطق وانگ،
اعلیٰ درجے دی منطق نوں
کئی قسماں والی منطق دا
اک روپ سمجھیا جا سکدا اے،
جتھے ہر مطلوبہ نوع دیاں
شےواں لئی اک قسم ہوئے۔
پر عموماً اعلیٰ نوع دی
منطق دی نحو نوں شروع
توں سیدھا متعین کرنا
زیادہ صاف ہوندا اے۔
مثال لئی، محدود تشکیل
والیاں انواع دا مجموعہ
استقرائی طور تے ایس
طرح متعین کر سکدے آں:
\begin{enumerate}
\item $\Nat$ اک محدود تشکیل والی نوع اے۔
\item جے $\sigma$ تے $\tau$ محدود تشکیل
  والیاں انواع ہون، تاں $\sigma
  \to \tau$ وی محدود تشکیل والی نوع اے۔
\item جے $\sigma$ تے $\tau$ محدود تشکیل
  والیاں انواع ہون، تاں $\sigma \times
  \tau$ وی محدود تشکیل والی نوع اے۔
\end{enumerate}
ایتھے محدود ہون دا مطلب نوع دی تشکیل دے محدود
قدم نیں؛ نوع دیاں شےواں دی گنتی متناہی ہون دا
دعویٰ نہیں۔
مطلب ایہ اے کہ $\Nat$
قدرتی اعداد دی نوع،
$\sigma \to \tau$ نوع
$\sigma$ توں نوع $\tau$
ول فنکشناں دی نوع،
تے $\sigma \times \tau$
جوڑیاں دی نوع اے
جیہناں دا اک رکن
$\sigma$ توں تے دوجا
$\tau$ توں ہوئے۔
فیر اصطلاحاں دا
مجموعہ استقرائی طور
تے ایس طرح متعین
کر سکدے آں:
\begin{enumerate}
\item ہر نوع $\sigma$ لئی
  متغیراں $x$, $y$,
  $z$, \dots دا
  ذخیرہ موجود اے؛
  ایہ نوع $\sigma$
  دے متغیر نیں۔
\item $\Obj 0$ نوع $\Nat$
  دی اصطلاح اے۔
\item $\Obj S$ (جانشین)
  نوع $\Nat \to \Nat$
  دی اصطلاح اے۔
\item جے $s$ نوع $\sigma$
  دی اصطلاح، تے $t$
  نوع $\Nat \to (\sigma \to \sigma)$
  دی اصطلاح ہوئے، تاں
  $\Obj{R}_{st}$ نوع
  $\Nat \to \sigma$
  دی اصطلاح اے۔
\item جے $s$ نوع
  $\tau \to \sigma$ دی
  اصطلاح تے $t$
  نوع~$\tau$ دی اصطلاح
  ہوئے، تاں $s(t)$
  نوع $\sigma$ دی
  اصطلاح اے۔
\item جے $s$ نوع~$\sigma$
  دی اصطلاح تے $x$
  نوع~$\tau$ دا متغیر
  ہوئے، تاں $\lambd[x][s]$
  نوع $\tau \to \sigma$
  دی اصطلاح اے۔
\item جے $s$ نوع~$\sigma$
  دی اصطلاح تے $t$
  نوع~$\tau$ دی اصطلاح
  ہوئے، تاں $\tuple{s, t}$
  نوع $\sigma \times
  \tau$ دی اصطلاح اے۔
\item جے $s$ نوع~$\sigma \times \tau$
  دی اصطلاح ہوئے، تاں
  $p_1(s)$ نوع~$\sigma$
  دی تے $p_2(s)$
  نوع~$\tau$ دی اصطلاح
  اے۔
\end{enumerate}
مطلب ایہ اے کہ
$\Obj{R}_{st}$ ایس
بازگشتی تعریف والے
فنکشن نوں بیان کردا اے:
\begin{align*}
\Obj{R}_{st}(0) & = s \\
\Obj{R}_{st}(x+1) & = t(x, R_{st}(x)),
\end{align*}
% PNB-SEM-032: arrow-type دے مطابق recursor دی دو دلیل والی لکھت دو لگاتار اطلاقاں دا مختصر نشان اے۔
بازگشت دی دوجی مساوات وچ دو دلیل والی لکھت
دو لگاتار اطلاقاں دا مختصر نشان اے: پہلے قدرتی عدد
اُتے اطلاق نال اک فنکشن ملدا اے، فیر اوہ پچھلی
حاصل قدر اُتے لاگو ہوندا اے۔
$\tuple{s, t}$ ایس جوڑے
نوں بیان کردا اے جیہدا
پہلا رکن~$s$ تے دوجا
رکن~$t$ اے؛ تے $p_1(s)$
تے~$p_2(s)$،~$s$ دے
پہلے تے دوجے رکن
(”پروجیکشن“) نوں بیان
کردے نیں۔ آخرکار
$\lambd[x][s]$ اوہ
فنکشن~$f$ بیان کردا
اے جیہڑا ایویں
متعین اے:
\[
f(x) = s
\]
ہر~$x$ لئی، جیہڑا
% PNB-SYN-011: شق6 وچ lambda دا input tau تے body/result sigma اے؛ اصل دی وضاحت وچ غلط sigma نوں tau کیتا۔
نوع~$\tau$ دا ہوئے؛
ایس طرح شق (6)
خاصیت توں فنکشن بنان
دا اک روپ دیندی اے،
جیہدے نال اصطلاحاں
راہی فنکشن متعین
ہو سکدے نیں۔
!!^{formula}s اکّو
نوع دیاں اصطلاحاں
وچکار !!{identity}
دے بیان $\eq[s][t]$،
عام قضیاتی رابطیاں
تے اعلیٰ نوع دے
مسوّراں توں بنائے
جاندے نیں۔ فیر
نظام دے مسلّمے
اوپر متعین اصطلاحاں
دیاں بنیادی مساوات،
تے نال مسوّراں
تے !!{identity}
دے عام منطقی
قاعدے رکھ سکدے
آں۔

جے محدود تشکیل والیاں انواع دے
نظام وچ سچائی دیاں
قدراں دی نوع
$\Omega$ شامل کر
دتی جاوے، تاں
ایہدی ورتوں نوں
چلاؤن والے مسلّمے
وی شامل کرنے
پین گے۔ اصل وچ
ذرا ہنر نال
پیچیدہ !!{formula}s
نوں بالکل ہٹا
کے انہاں دی
تھاں نوع $\Omega$
دیاں اصطلاحاں
رکھیاں جا
سکدیاں نیں!{}
ثبوتی نظام
فیر اوس حساب
نال بدل
سکدا اے۔
نتیجہ بنیادی
طور تے
\emph{انواع دا
سادہ نظریہ}
اے جیہڑا
الونزو چرچ
نے 1930 دی
دہائی وچ
پیش کیتا۔

دوجے درجے
دی منطق
وانگ، اعلیٰ
نوع دی معنیات
دے وی وکھ
وکھ روپ
ورتے جا
سکدے نیں۔
مکمل روپ
وچ نوع
$\sigma \to \tau$
دے متغیر
نوع~$\sigma$
دیاں شےواں
توں نوع~$\tau$
دیاں شےواں
تک دے
\emph{سارے}
فنکشناں
اُتے چل
دے نیں۔
جداں امید
ہو سکدی
اے، ایہ
معنیات
اینی طاقتور
اے کہ
ایدے لئی
کوئی تمام،
مؤثر
!!{derivation}
نظام نہیں۔
پر اک
کمزور معنیات
وی وچار
سکدے آں،
جیہدے وچ
!!a{structure}
عناصر دے
سیٹ
$T_\tau$،
ہر نوع
$\tau$ لئی
اک، تے نال
اطلاق،
پروجیکشن
وغیرہ دے
مناسب عمل
رکھدی اے۔
جے تفصیل
ٹھیک طرح
پوری کیتی
جاوے، تاں
اوپر بیان
کیتے
!!{derivation}
نظاماں
واسطے
تمامیت دے
قضیے مل
سکدے نیں۔

اعلیٰ نوع
دی منطق
پرکشش اے
کیوں جے
ایہ ایسا
ڈھانچا دیندی
اے جیہدے
وچ ریاضی
دا وڈا حصہ
قدرتی طریقے
نال سمویا
جا سکدا اے:
$\Nat$
توں شروع
کرکے حقیقی
اعداد،
متصل فنکشن
وغیرہ متعین
ہو سکدے
نیں۔ ایہ
وجدانیت
پسند منطق
دے سیاق
وچ وی خاص
طور تے
پرکشش اے،
کیوں جے
انواع دیاں
”تعمیری“
تشریحاں
صاف نیں۔
بلکہ اعلیٰ
نوع دی
معنیات دے
تعمیری روپ
وی تیار
کیتے جا
سکدے نیں
(جیہڑے
کلاسیکی
منطق دی
بجائے
وجدانیت
پسند منطق
اُتے مبنی
ہون)؛
ایہ روپ
انہاں
تعمیری
تشریحاں
نوں چنگی
طرح واضح
کردے نیں
تے کئی
لحاظاں
نال کلاسیکی
ہم منصباں
نسبت ودھ
دلچسپ نیں۔
\endgroup
% /OLP-0178

% OLP-0179 content/first-order-logic/beyond/intuitionistic-logic.tex
\begingroup
\def\olfilename{intuitionistic-logic}
\expandafter\def\csname ol@sectioncs\endcsname{section}

\olfileid{fol}{byd}{il}

\olsection{وجدانیت پسند منطق}

دوجے درجے تے اعلیٰ درجے دی منطق دے
برعکس، وجدانیت پسند پہلے درجے دی منطق
کلاسیکی روپ نوں محدود کرکے بندی اے؛
ایہ ہور ”تعمیری“ قسم دے استدلال دا
نمونہ بنان لئی اے۔ ہیٹھلی مثالاں
ایدے پیچھے دیاں کجھ وجوہاں واضح
کر سکدیاں نیں۔

فرض کرو اک دن کوئی بندہ تہاڈے کول
آ کے اعلان کرے کہ اوہنے اک قدرتی
عدد~$x$ معلوم کر لیا اے جیہدی
خاصیت ایہ اے: جے $x$ اولی عدد
ہوئے تاں ریمان دا مفروضہ سچ اے،
تے جے $x$ مرکب عدد ہوئے تاں
ریمان دا مفروضہ جھوٹا اے۔
کمال دی خبر!{} ریمان دا مفروضہ
سچ اے یا نہیں، ایہ ریاضی دے
وڈے کھلے سوالاں وچوں اک اے؛
ایتھے لگدا اے کہ سوال نوں
حسابی جانچ تک گھٹا دتا گیا
اے، یعنی ایہ معلوم کرنا رہ
گیا اے کہ اک خاص عدد اولی
اے یا نہیں۔

$x$ دی جادوئی قدر کیہڑی اے؟
اوہ ایویں بیان کردا اے:
$x$ اوہ قدرتی عدد اے
جیہڑا، جے ریمان دا
مفروضہ سچ ہوئے، $7$
دے برابر اے، تے نہیں
تاں $9$ دے برابر۔

تسی غصے وچ اپنے پیسے
واپس منگدے او۔ کلاسیکی
نقطۂ نظر توں اوپرلی
وضاحت واقعی $x$ دی
اک منفرد قدر متعین
کردی اے؛ پر تہاڈے
لئی مطلوب $x$ دی
اوہ قدر اے جیہڑی
\emph{صاف صاف} دتی
گئی ہوئے۔

ہور، شاید گھٹ بناوٹی،
مثال ویکھو۔ اسیں
جاندے آں کہ غیر
ناطق عدد نوں ناطق
قوت تک چڑھا کے
ناطق نتیجہ مل
سکدا اے؛ مثال
لئی
$\sqrt{2}^2 = 2$۔
پر ایہ گل اینی
صاف نہیں کہ غیر
ناطق عدد نوں
\emph{غیر ناطق}
قوت تک چڑھا کے
ناطق نتیجہ مل
سکدا اے یا نہیں۔
ہیٹھلا قضیہ ہاں
وچ جواب دیندا اے:

\begin{thm}
ایسے غیر ناطق
اعداد $a$ تے
$b$ موجود نیں
جیہناں لئی
$a^b$ ناطق
اے۔
\end{thm}

\begin{proof}
$\sqrt{2}^{\sqrt{2}}$
لو۔ جے ایہ
ناطق ہوئے، تاں
کم ہو گیا:
$a = b = \sqrt{2}$
رکھ لو۔ ورنہ
ایہ غیر ناطق
اے۔ ایس صورت
وچ
\[
(\sqrt{2}^{\sqrt{2}})^{\sqrt{2}} = \sqrt{2}^{\sqrt{2} \cdot
  \sqrt{2}} = \sqrt{2}^2 = 2,
\]
جیہڑا یقیناً
ناطق اے۔
اس لئی ایس
صورت وچ
$a$ نوں
$\sqrt{2}^{\sqrt{2}}$
تے $b$ نوں~$\sqrt 2$
رکھ لو۔
\end{proof}

کی ایہ درست
ثبوت اے؟
زیادہ تر ریاضی
دان سمجھدے نیں
کہ اے۔ پر
فیر وی ایتھے
اک گل کجھ
تسلی نہیں دیندی
اے: اسیں
اک خاصیت
رکھن والے
حقیقی اعداد
دے جوڑے
دی موجودگی
ثابت کر
دتی، پر
ایہ نہیں
دس سکدے
کہ اوہ
\emph{کیہڑا}
جوڑا اے۔
اوہی نتیجہ
ایس طرح وی
ثابت ہو سکدا
اے کہ ثبوت
وچ جوڑا
$a$, $b$
\emph{صاف
دتا} جاوے:
$a = \sqrt{3}$
تے
$b = \log_3 4$
لو۔
% PNB-SEM-033: explicit pair لئی دوواں گواہاں دی irrationality دی source وچ مختصر چھڈی جانچ واضح کیتی۔
ایہ دوویں عدد غیر ناطق نیں۔ تین دے مربع جذر نوں
گھٹ توں گھٹ صورت دی صحیح عددی نسبت منّن نال
نسبت دے دوناں عدد تین نال ونڈن جوگ نکلدے نیں،
جیہڑا ایس گھٹ توں گھٹ صورت دے خلاف اے۔ دوجے
عدد نوں مثبت صحیح اعداد دی نسبت منّن نال تین دی
اک مثبت صحیح عددی گھات، چار دی اک مثبت صحیح عددی
گھات دے برابر ہونی پئے گی؛ پہلی تین نال ونڈن جوگ
اے تے دوجی نہیں، ایس لئی ایہ وی ممکن نہیں۔ فیر
\[
a^b = \sqrt{3}^{\log_3 4} = 3^{1/2 \cdot \log_3 4} = (3^{\log_3
  4})^{1/2} = 4^{1/2}= 2,
\]
کیوں جے
$3^{\log_3 x} = x$۔

وجدانیت پسند
منطق ایہو
جہے استدلال
دا نمونہ
بنان لئی
اے جیہدے
وچ پہلے
ثبوت والی
چال دی
اجازت نہیں۔
ایہ ثابت
کرنا کہ
کوئی $x$
موجود اے
جو~$!A(x)$
پورا کردا
اے، ایس
دا مطلب
اے کہ
تسی اک
خاص~$x$
تے ثبوت
دیو کہ
اوہ $!A$
پورا کردا
اے، جداں
دوجے ثبوت
وچ۔ ایہ
ثابت کرن
لئی کہ
$!A$ یا
$!B$ درست
اے، انہاں
وچوں اک
دا ثبوت
دینا پیندا
اے۔

رسمی طور
اُتے وجدانیت
پسند پہلے
درجے دی
منطق، پہلے
درجے دے
!!a{derivation}
نظام نوں
خاص طریقے
نال محدود
کرن نال
ملدی اے۔
ایس طرح
دوجے تے
اعلیٰ درجے
دی منطق
دے وجدانیت
پسند روپ
وی نیں۔
ریاضی دے
حساب نال
ایہ بس
رسمی استنباطی
نظام نیں؛
پر جداں
پہلاں دسیا،
انہاں دا
مقصد ریاضیاتی
استدلال دی
اک قسم
دا نمونہ
بنانا اے۔
کوئی ایہ
سمجھ سکدا
اے کہ ایہ
اوہ استدلال
اے جیہڑا
ریاضی بارے
کسی خاص
فلسفیانہ
موقف (مثلاً
برووَر دی
وجدانیت)
نال جائز
ٹھہردا اے؛
یا ایہ
اک زیادہ
”ٹھوس“
تے اطمینان
بخش ریاضیاتی
استدلال اے
(بشپ دے
تعمیریت
والے انداز
وچ)؛
تے ایس
اُتے بحث
وی ہو
سکدی اے
کہ رسمی
بیان غیر
رسمی محرک
نوں صحیح
پکڑدا اے
یا نہیں۔
پر ساڈا
فلسفیانہ
موقف کوئی
وی ہوئے،
اسیں وجدانیت
پسند منطق
دا رسمی
طور اُتے
پیش کیتی
منطق ونجوں
مطالعہ کر
سکدے آں؛
تے کئی
ریاضیاتی
منطق دان،
اپنیاں
وجوہاں نال،
ایہ مطالعہ
دلچسپ سمجھدے
نیں۔

وجدانیت پسند
رابطیاں دی
اک غیر رسمی
تعمیری تشریح
اے، جیہنوں
عموماً BHK
تشریح آکھدے
نیں؛ ایہ
نام برووَر،
ہیٹنگ تے
کولموگوروف
دے ناماں
توں بنیا
اے۔ تشریح
ایویں اے:
$!A
\land !B$
دا ثبوت،
$!A$ دے
ثبوت تے
$!B$ دے
ثبوت دی
جوڑی اے؛
$!A \lor !B$
دا ثبوت،
$!A$ یا
$!B$ وچوں
اک دا ثبوت
اے، تے
صاف خبر
ہوئے کہ
کیہڑا؛
$!A \lif !B$
دا ثبوت
اک طریقہ
اے جو
$!A$ دے
ثبوت نوں
~$!B$ دے
ثبوت وچ
بدلدا اے؛
$\lforall[x][!A(x)]$
دا ثبوت
اک طریقہ
اے جو
$!A(x)$ دا
ثبوت دیندا
اے،~$x$
دی ہر قدر
لئی؛
تے
$\lexists[x][!A(x)]$
دا ثبوت
~$x$ دی
اک قدر
تے ایس
قدر دے
~$!A$ پورا
کرن دا
ثبوت اے۔
ایہ تشریح
حسابی اصطلاحاں
وچ ”کری--ہاورڈ
تماثل“ یا
”!!{formula}s
بطور انواع
دا اصول“
وی کہلا
سکدی اے:
!!a{formula}
نوں اک
خاص قسم
دے ڈیٹا
ٹائپ دی
تعیین
سمجھو، تے
ثبوتاں نوں
انہاں ڈیٹا
ٹائپاں دیاں
حسابی شےواں،
جیہڑیاں
متعلقہ
!!{formula}
دی سچائی
سامنے لے
آندیاں
نیں۔

وجدانیت پسند
منطق نوں
اکثر کلاسیکی
منطق ”وچوں“
تیسرے امکان
دے اخراج دا
قانون کڈ
کے باقی
رہن والی
منطق سمجھیا
جاندا اے۔
ہیٹھلا قضیہ
ایہ گل
ہور عین
کردا اے۔
\begin{thm}
وجدانیت پسند
منطق وچ
ہیٹھلے مسلّمی
خاکے اک
دوجے دے
برابر نیں:
\begin{enumerate}
\item $(\lnot !A \lif \lfalse) \lif !A$.
\item $!A \lor \lnot !A$
\item $\lnot \lnot !A \lif !A$
\end{enumerate}
\end{thm}
ایہ برابری سارے خاکیاں نوں اضافی مسلّمے منّن بارے
اے؛ ایہ نہیں کہ دکھائے ہوئے فارمولے وجدانیت پسند
منطق وچ پہلے توں ثابت نیں۔ سارے خاکے دا فرض
کرنا تے اک دتے ہوئے فارمولے لئی اوہدی صرف اک
مخصوص صورت دا فرض کرنا وکھ وکھ گلّاں نیں۔
اک خاکے دی
مخصوص صورت
باقی دوواں
وچوں کسے
اک توں
کڈنا وجدانیت
پسند منطق
دی چنگی
مشق اے۔

وجدانیت پسند
قضیاتی منطق
دے پہلے
استنباطی نظام،
برووَر دی
وجدانیت
دی رسمی
صورتاں ونجوں،
کولموگوروف،
گلیوینکو
تے ہیٹنگ
نے اک دوجے
توں آزاد
پیش کیتے۔
وجدانیت پسند
پہلے درجے
دی منطق
(تے وجدانیت
پسند ریاضی
دے کجھ حصیاں)
دی پہلی
رسمی صورت
ہیٹنگ نے
دتی۔
اگرچہ کئی
کلاسیکی
طور اُتے
درست خاکے
وجدانیت پسند
طور اُتے
درست نہیں،
پر کئی
ہوندے نیں۔

\emph{دوہری نفی دا ترجمہ}
کلاسیکی تے
وجدانیت پسند
منطق دے
اک اہم رشتے
نوں بیان
کردا اے۔
ایہ ہیٹھاں
استقرائی
طور اُتے
متعین اے
(~$!A^N$
نوں کلاسیکی
!!{formula}~$!A$
دا ”وجدانیت
پسند“ ترجمہ
سمجھو):
\begin{align*}
\formula A^N & \ident \lnot\lnot !A \quad \text{جوہری !!{formula}s $!A$ لئی} \\
(!A \land !B)^N & \ident (!A^N \land !B^N) \\
(!A \lor !B)^N & \ident  \lnot\lnot (!A^N \lor !B^N) \\
(!A \lif !B)^N & \ident (!A^N \lif !B^N) \\
(\lforall[x][!A])^N & \ident \lforall[x][!A^N] \\
(\lexists[x][!A])^N & \ident \lnot\lnot\lexists[x][!A^N]
\end{align*}
کولموگوروف تے
گلیوینکو نے
قضیاتی منطق
لئی ایس ترجمے
دے روپ دتے
سن؛ محمولی
منطق لئی
ایہ گوڈل
تے گینٹزن
نے اک دوجے
توں آزاد
دتا۔
فیر ساڈے
کول ایہ
اے:

\begin{thm}
\begin{enumerate}
\item $!A \liff !A^N$ کلاسیکی طور اُتے ثابت ہوندا اے۔
\item جے $!A$ کلاسیکی طور اُتے ثابت ہوئے، تاں
  $!A^N$ وجدانیت پسند طور اُتے ثابت ہوندا اے۔
\end{enumerate}
\end{thm}

ہن ایہ
گفتگو تصور
کرو۔ کلاسیکی
ریاضی دان:
”میں $!A$
ثابت کر
دتا!“
وجدانیت پسند
ریاضی دان:
”تہاڈا ثبوت
درست نہیں۔
اصل وچ تسی
$!A^N$ ثابت
کیتا اے۔“
کلاسیکی
ریاضی دان:
”مینوں ایہ
وی منظور!“
کلاسیکی
ریاضی دان
لئی وجدانیت
پسند بندہ
بس باریک
نکتے اُتے
اڑیا ہویا
اے، کیوں
جے اوہدے
حساب نال
دوویں ہم
ارز نیں۔
پر وجدانیت
پسند اصرار
کردا اے
کہ فرق
اے۔

چیتا رہے
کہ اوپرلا
ترجمہ صرف
خالص منطق
بارے اے؛
ایہ اس
سوال دا
جواب نہیں
دیندا کہ
کلاسیکی
تے وجدانیت
پسند ریاضی
لئی مناسب
\emph{غیر منطقی}
مسلّمے کیہڑے
نیں، یا
انہاں دا
باہمی رشتہ
کی اے۔
پر اوپرلے
قضیے دی
ہیٹھلی
چھوٹی
توسیع کجھ
فائدہ مند
خبر دیندی
اے:

\begin{thm}
جے $\Gamma$ کلاسیکی طور اُتے $!A$ ثابت کرے،
تاں $\Gamma^N$ وجدانیت پسند طور اُتے
$!A^N$ ثابت کردا اے۔
\end{thm}

دوجے لفظاں
وچ جے
$!A$ کجھ
مفروضیاں
توں کلاسیکی
طور اُتے
ثابت ہوئے،
تاں $!A^N$
انہاں مفروضیاں
دے دوہری
نفی والے
ترجمیاں
توں ثابت
ہوندا اے۔

اک جملے
یا قضیاتی
!!{formula}
دا وجدانیت
پسند اعتبار
وکھان لئی
ثبوت دینا
کافی اے۔
پر ایہ
کیویں وکھاؤ
گے کہ
اوہ معتبر
نہیں؟
ایہدے لئی
ایسی معنیات
چاہیدی اے
جیہڑی درست
ہوئے،
تے بہتر
اے تمام
وی ہوئے۔
کرپکے دی
معنیات ایس
لوڑ لئی
بہت ڈھکدی
اے۔

اسیں اوہی
کم کر
سکدے آں
جیہڑا کلاسیکی
منطق لئی
کیتا سی:
معنیات متعین
کرو، تے
درستگی تے
تمامیت ثابت
کرو۔
پر اک
فرق چیتے
رکھن جوگا
اے۔ کلاسیکی
منطق وچ
معنیات اک
لحاظ نال
”فطری“
سی، جیویں
منطقی رابطیاں
دے مطلب
وچ پہلے
توں موجود
ہوئے۔
کرپکے معنیات
لئی کجھ
بدیہی محرک
دتے جا
سکدے نیں،
پر اوہ
پہلی معنیات
وانگ ناگزیر
محسوس نہیں
ہوندی۔
ایس توں
علاوہ کلاسیکی
!!{structure}
اک فطری
ریاضیاتی تصور
اے؛ اسیں
یا تاں
!!a{structure}
نوں کلاسیکی
پہلے درجے
دی منطق
دے مطالعے
دا آلہ
سمجھ سکدے
آں، یا
کلاسیکی پہلے
درجے دی
منطق نوں
ریاضیاتی
!!{structure}s
دے مطالعے
دا آلہ۔
ایس دے
مقابلے وچ
کرپکے دیاں
!!{structure}s
صرف منطقی
بناوٹ ونجوں
دیکھیاں جا
سکدیاں
نیں؛ انہاں
دی کوئی
آزاد ریاضیاتی
دلچسپی نظر
نہیں آندی۔

قضیاتی زبان
لئی کرپکے
دی !!{structure}~
$\mModel M = \tuple{W, R, V}$
وچ اک
سیٹ~$W$،
جزوی ترتیب~$R$
جیہڑی~$W$
اُتے ہوئے
تے جیہدا
سب توں
چھوٹا !!{element}
ہوئے،
تے قضیاتی
متغیراں نوں
~$W$ دے
!!{element}s
نال ”سچ نوں برقرار رکھن والی“
(monotone) طور
تے جوڑن
والی قدر
بندی ہوندی
اے۔ خیال
ایہ اے
کہ $W$
دے !!{element}s
”جہاناں“
یا ”علمی
حالَتاں“ دی
نمائندگی کردے
نیں؛ اک
عنصر $v \geq
u$،~$u$
دی ”ممکنہ
اگلی حالت“
دسدا اے؛
تے~$u$
نال جوڑے
قضیاتی متغیر
اوہ قضیے
نیں جیہڑے
~$u$ دی
حالت وچ
سچے معلوم
نیں۔ ایس شرط دا مطلب اے کہ جیہڑا قضیاتی متغیر
اک حالت وچ سچا اے، اوہ اوس دی ہر ممکنہ اگلی
حالت وچ وی سچا رہندا اے؛ صرف متغیراں نوں
حالَتاں نال من چاہے طریقے نال جوڑنا کافی نہیں۔ فیر
تحمیلی صدق
(forcing)
دا تعلق
$\mSat{M}{!A}[w]$
ایس رشتے
نوں زبان
دے سارے
!!{formula}s
تک ودھاؤندا
اے؛
$\mSat{M}{!A}[w]$
نوں ”$!A$
حالت~$w$
وچ سچ
اے“ پڑھو۔
ایہ رشتہ
استقرائی
طور اُتے
ایویں متعین
اے:
\begin{enumerate}
\item $\mSat{M}{\Obj p_i}[w]$ جے تے صرف جے
  $\Obj p_i$ اوہناں قضیاتی متغیراں وچوں
  ہوئے جیہڑے~$w$ نال جوڑے گئے نیں۔
\item $\mSat/{M}{\lfalse}[w]$۔
\item $\mSat{M}{(!A \land !B)}[w]$ جے تے صرف جے
  $\mSat{M}{!A}[w]$ تے $\mSat{M}{!B}[w]$۔
\item $\mSat{M}{(!A \lor !B)}[w]$ جے تے صرف جے
  $\mSat{M}{!A}[w]$ یا $\mSat{M}{!B}[w]$۔
\item $\mSat{M}{(!A \lif !B)}[w]$ جے تے صرف جے
  جدوں وی $w' \geq w$ تے
  $\mSat{M}{!A}[w']$ ہون، تاں
  $\mSat{M}{!B}[w']$ وی ہوئے۔
\end{enumerate}
چنگی مشق
ایہ اے
کہ ایداں
دی کرپکے
!!{structure}
بنا کے،
جیہڑی مخالف
مثال دیوے،
وکھاؤ کہ
$\lnot (p \land q) \lif
(\lnot p \lor \lnot q)$
وجدانیت پسند
طور اُتے
معتبر نہیں۔
\endgroup
% /OLP-0179

% OLP-0180 content/first-order-logic/beyond/modal-logics.tex
\begingroup
\def\olfilename{modal-logics}
\expandafter\def\csname ol@sectioncs\endcsname{section}

\olfileid{fol}{byd}{mod}

\olsection{امکانی منطق دیاں قسماں}

ہیٹھلے شرطیہ جملے دی مثال ویکھو:
\begin{quote}
  جے جیریمی اوس کمرے وچ اکیلا اے،
  تاں اوہ نشے وچ، ننگا، تے کرسیاں
  اُتے نچدا پیا اے۔
\end{quote}
ایہ ایداں دا شرطیہ بیان اے جیہڑا
مادی معنی وچ سچ ہو سکدا اے،
پر فیر وی غلط تاثر دیندا اے:
الفاظ ایس گل دا اشارہ کردے نیں
کہ مقدم تے نتیجے وچ اوس توں
مضبوط رشتہ اے جو محض ایہ ہووے
کہ یا مقدم جھوٹا اے یا تالی
سچی اے۔ یعنی جملہ ایہ ظاہر
کردا اے کہ دعویٰ صرف ایس
خاص جہان وچ سچ نہیں (جتھے
ایہ معمولی طور تے سچ ہو
سکدا اے، کیوں جے جیریمی
کمرے وچ اکیلا نہیں)، بلکہ
جے مقدم سچ \emph{ہوندا}
تاں تالی وی \emph{سچ ہوندی}۔
دوجے لفظاں وچ بیان نوں
ایہ سمجھیا جا سکدا اے
کہ دعویٰ صرف ایس جہان وچ
نہیں بلکہ ہر ”ممکن“
جہان وچ سچ اے؛ یا صرف
ایس خاص جہان وچ سچ
ہون دی بجائے
\emph{لازماً} سچ اے۔

امکانی منطق ایس قسم
دی لازمی سچائی دا
مطلب صاف کرن لئی
بنائی گئی اے۔ عام
قضیاتی منطق وچ اک
چوکور عامل شامل
کرکے امکانی قضیاتی
منطق ملدی اے؛ یعنی
جے $!A$ اک
!!a{formula} اے، تاں
$\Box !A$ وی
فارمولا اے۔
بدیہی طور اُتے
$\Box !A$ ایہ
آکھدا اے کہ
$!A$ \emph{لازماً}
سچ اے، یا ہر
ممکن جہان
وچ سچ اے۔
$\Diamond !A$ نوں
عموماً
$\lnot \Box \lnot !A$
دا مختصر نشان
منیا جاندا اے؛ ایس
دا مطلب ایہ
پڑھیا جا سکدا
اے کہ $!A$
\emph{ممکنہ طور
اُتے} سچ اے۔
ظاہر اے
کہ محمولی
منطق وچ
وی امکانی
عامل شامل
کیتے جا
سکدے نیں۔

کرپکے دیاں
!!{structure}s
امکانی منطق
دی معنیات
دین لئی
ورتی جا
سکدیاں نیں؛
اصل وچ
کرپکے نے
ایہ معنیات
پہلاں امکانی
منطق لئی
سوچی سی۔
ایتھے جزوی
ترتیباں تک
محدود رہن
دی بجائے،
عمومی طور
اُتے ”ممکن
جہاناں“ دا
سیٹ $P$
تے جہاناں
وچ دو-جگہی
”رسائی“
دا تعلق
$\Atom{R}{x,y}$
ہوندا اے۔
بدیہی طور
اُتے
$\Atom{R}{p,q}$
آکھدا اے
کہ جہان~$q$
جہان~$p$
نال موافق
اے؛ یعنی
جے اسیں
جہان~$p$
”وچ“ آں،
تاں ایہ
امکان وی
سوچنا پیندا
اے کہ
جہان~$q$
ورگا ہو
سکدا سی۔
کسی جہان وچ چوکور عامل والا بیان اوہدوں سچ ہوندا اے جدوں
اندرلا بیان اوس جہان توں رسائی والے ہر جہان وچ سچ ہووے؛
امکان والے بیان لئی کم از کم اک ایسا جہان کافی اے۔
سو عمومی صورت وچ سارے ممکن جہان نہیں، بلکہ متعلقہ جہان توں
رسائی والے جہان ویکھنے نیں۔ پچھلی وجدانیت پسند معنیات دے برعکس،
ایتھے ایٹمی قضیے دا سچ رسائی دے تعلق نال آگے لازماً برقرار
رہن دی شرط نہیں۔

امکانی منطق
نوں کدی
”مفہومی“
(intensional)
منطق وی
آکھدے نیں،
”امتدادی“
(extensional)
منطق دے
مقابلے وچ۔
کلاسیکی
منطق ورگی
امتدادی
منطق دی
مقصود معنیات
صرف اک
جہان، یعنی
”واقعی“
جہان، دا
حوالہ دیندی
اے؛ جدکہ
مفہومی منطق
دی معنیات
وڈھیری پیچیدہ
وجودیات اُتے
ٹکدی اے۔
لازمیت نوں
!!{structure} دے
روپ وچ پیش
کرن توں علاوہ
امکانی عامل
ہور لسانی
بناوٹاں نوں
!!{structure} دے
روپ وچ پیش
کرن لئی وی
ورتے جا
سکدے نیں:
ورتوں دے
مطابق
$\Box$ تے
$\Diamond$ دا
مطلب نواں
لے لو۔
مثال لئی:
\begin{enumerate}
\item ثابت پذیری
  دی منطق
  وچ $\Box !A$
  نوں ”$!A$
  ثابت ہو
  سکدا اے“
  تے
  $\Diamond !A$
  نوں ”$!A$
  موافق اے“
  پڑھدے نیں۔
  ایہ دوویں پڑھتاں اکّو مقررہ پس منظر دے نظریے دے لحاظ نال نیں:
  ثابت کرن دے قاعدے بدل بدل کے معنی نہیں لینے؛ موافقت دا مطلب
  اوہ نظریہ ایس بیان نال تضاد پیدا نہ کرے۔
\item علم دی
  منطق وچ
  $\Box !A$
  نوں ”میں
  $!A$ جاندا
  آں“ یا
  ”میں $!A$
  منّدا آں“
  پڑھ سکدے
  آں۔
\item زمانی
  منطق وچ
  $\Box !A$
  نوں ”$!A$
  ہمیشہ سچ
  اے“ تے
  $\Diamond !A$
  نوں ”$!A$
  کدی کدی
  سچ اے“
  پڑھیا جا
  سکدا اے۔
\end{enumerate}

اسیں چاہندے
آں کہ منطق
وچ امکانی
عامل بارے
قاعدے تے
مسلّمے شامل
ہون۔ مثال
لئی نظام
$\Log{S4}$
وچ عام
قضیاتی منطق
دے مسلّمے
تے قاعدے،
نال ہیٹھلے
مسلّمے ہوندے
نیں:
% PNB-SEM-034: لازمیت دا قاعدہ بے کھلے مقامی مفروضیاں ثابت نتیجے لئی اے۔
\begin{align*}
& \Box (!A \lif !B) \lif (\Box !A \lif \Box !B)\\
& \Box !A \lif !A \\
& \Box !A \lif \Box \Box !A
\intertext{نال اک قاعدہ وی: جے $!A$ کوئی کھلا مقامی مفروضہ ورتے بغیر
  ثابت شدہ قضیہ اے، تاں $\Box !A$ کڈو۔ ایہ قاعدہ محض فرض کیتے
  ہوئے بیان اُتے نہیں لا سکدے۔
  $\Log{S5}$ ہیٹھلا مسلّمہ شامل کردا اے:}
& \Diamond !A \lif \Box \Diamond !A
\end{align*}
ایہناں مسلّمیاں
دیاں وکھ
صورتاں وکھ
ورتوں لئی
مناسب ہو
سکدیاں نیں؛
مثال لئی
S5 نوں
عموماً منطقی
لازمیت دا
تصور بیان
کرن والا
منیا جاندا
اے۔ چنگی
گل ایہ
اے کہ
کرپکے دیاں
!!{structure}s
وچ رسائی
دے تعلق
نوں فطری
طریقے نال
محدود کرکے
عموماً
ایسی معنیات
مل جاندی
اے جیہدے
لئی !!{derivation}
نظام درست
تے تمام
ہوئے۔
مثال لئی
$\Log{S4}$
انہاں کرپکے
!!{structure}s
دی جماعت
نال ملدا
اے جیہناں
وچ رسائی
دا تعلق
انعکاسی
تے متعدی
اے۔
$\Log{S5}$
انہاں کرپکے
!!{structure}s
دی جماعت
نال ملدا
اے جیہناں
وچ رسائی
دا تعلق
{\em آفاقی}
اے، یعنی
ہر جہان
توں ہر
جہان
تک رسائی
ہوئے؛ ایس
لئی
$\Box !A$
پورا اتردا اے
جے تے
صرف جے
$!A$ ہر
جہان وچ
پورا اترے۔
\endgroup
% /OLP-0180

% OLP-0181 content/first-order-logic/beyond/other-logics.tex
\begingroup
\def\olfilename{other-logics}
\expandafter\def\csname ol@sectioncs\endcsname{section}

\olfileid{fol}{byd}{oth}

\olsection{ہور منطق دیاں قسماں}

ہُن تک شاید تہاڈے ذہن وچ ایہ گل آ گئی ہووے کہ نویں منطق بنانی اوکھی نہیں۔
تُسی وی اپنی نحوی ساخت بنا سکدے او، اک استدلالی نظام گھڑ سکدے او، تے اوہدے
نال ملدی معنیات ترتیب دے سکدے او۔ جے تسی چاہندے او کہ !!{derivation}
دا نظام انہاں معنیات دے لحاظ نال تمام ہوئے، تاں کجھ سوجھ بوجھ درکار ہو
سکدی اے؛ تے عام لوکاں نوں ایہ منوان لئی وی محنت لگ سکدی اے کہ تہاڈی
منطق واقعی دلچسپ اے۔ پر ایس دے بدلے تُسی اپنی منطق دیاں ریاضیاتی تے
کمپیوٹیشنل خاصیتاں دی کھوج کردیاں کئی گھنٹے چنگی، بے ضرر تفریح وچ
گزار سکدے او۔

پچھلیاں کئی دہاکیاں وچ رسمی منطق دیاں قسماں حیران کن حد تک ودھیاں نیں۔
دھندلی منطق \pnbLatin{(fuzzy logic)} مبہم خاصیتاں بارے استدلال دا نمونہ بناؤن لئی
بنائی جاندی اے۔ احتمالاتی منطق \pnbLatin{(probabilistic logic)} غیر یقینی صورتاں
بارے استدلال دا نمونہ بناندی اے۔ ڈیفالٹ منطق دے نظام \pnbLatin{(default logics)}
تے نویں معلومات نال نتیجے پلٹ سکدے ہون والی منطق دے نظام
\pnbLatin{(nonmonotonic logics)} ایسے استدلال دا نمونہ بناندے نیں: یعنی اوہ ”معقول“ نتیجے جو نویں
معلومات آون اُتے رد ہو سکدے نیں۔ علم بارے استدلال لئی علمی منطق دیاں
قسماں \pnbLatin{(epistemic logics)}، سبب تے نتیجے دے رشتیاں بارے استدلال لئی
سببیتی منطق دیاں قسماں \pnbLatin{(causal logics)}، تے اخلاقی فرضاں تے ذمہ واریاں
بارے استدلال لئی ”اخلاقی فرض دی منطق“ \pnbLatin{(deontic logics)} وی نیں۔ ایہناں
نظاماں نوں متعارف کراؤن دی اصل ترغیب فلسفیانہ، ریاضیاتی یا
کمپیوٹیشنل ہووے، ایس دے مطابق تہاڈے سامنے انہاں دا مطالعہ ریاضیاتی
منطق، فلسفیانہ منطق، مصنوعی ذہانت، ادراکی سائنس یا کسے ہور شعبے دے
تحت آ سکدا اے۔

ایہ فہرست لمی ہوندی جاندی اے، تے امکان بے انت لگدے نیں۔ سارے انسانی
استدلال نوں حسابی عمل وچ ڈھال دیون دا لائبنز دا خواب شاید اسیں کدے
پورا نہ کر سکیے---پر ایہ گل سانوں کوشش کرن توں نہیں روک سکدی۔
\endgroup
% /OLP-0181

\OLEndChapterHook
\endgroup
% /OLP-0174

\OLEndPartHook
\endgroup
% /OLP-0138

% OLP-0182 content/model-theory/model-theory.tex
\begingroup
\def\olfilename{model-theory}
\expandafter\def\csname ol@sectioncs\endcsname{section}

\olpart{mod}{نمونیاں دا نظریہ}

\begin{editorial}
  نمونیاں دے نظریے دا ایہ مواد نامکمل تے تجرباتی اے۔ فی الحال ایہ
  آلدو انتونیلی دے نمونیاں دے نظریے بارے نوٹس نوں ورت کے ڈھالیا گیا مواد اے؛
  پہلے درجے دی منطق والے حصے وچ شامل موضوع (نظریے، تمامیت، محدود
  گواہی دا قضیہ) ایس وچوں کڈ دتے گئے نیں۔ نصابی کتاب لئی کارآمد
  ہون واسطے ایس نوں ہور بہت زیادہ تعارف، ترغیب تے وضاحت، نال مشقاں،
  درکار نیں۔ اینڈی ارانا خاص طور تے ایس حصے اُتے کم کرن دی
  منصوبہ بندی کر رہیا اے (\gitissue{65})۔
\end{editorial}

% OLP-0183 content/model-theory/basics/basics.tex
\begingroup
\def\olfilename{basics}
\expandafter\def\csname ol@sectioncs\endcsname{section}

\olchapter{mod}{bas}{نمونیاں دے نظریے دیاں بنیادی گلّاں}

% OLP-0184 content/model-theory/basics/reducts-and-expansions.tex
\begingroup
\def\olfilename{reducts-and-expansions}
\expandafter\def\csname ol@sectioncs\endcsname{section}

\olfileid{mod}{bas}{red}
\section{ساختاں دے گھٹائے تے ودھائے روپ}

اکثر ایہناں زباناں دا مقابلہ کرنا فائدہ مند یا ضروری ہوندا اے جنہاں
دیاں کجھ علامتاں سانجھیاں ہوون، تے انہاں زباناں دیاں !!{structure}s
دا وی۔ سبھ توں عام صورت ایہ اے کہ !!a{language}~$\Lang{L}$ دیاں
ساریاں علامتاں !!a{language}~$\Lang{L'}$ وچ وی شامل ہون، یعنی
$\Lang{L} \subseteq
\Lang{L'}$۔ فیر $\Lang{L}$-!!{structure}~$\Struct{M}$ نوں
وادھو علامتاں دی تشریح شامل کرکے ہمیشہ اک $\Lang{L'}$-!!{structure}
بنایا جا سکدا اے، جدکہ سانجھیاں علامتاں دی تشریح اوہی رکھی جاندی اے۔
دوجے پاسے، $\Lang{L'}$-ساخت~$\Struct{M'}$ توں اوہناں علامتاں دیاں
تشریحاں ”بھلا“ کے، جیہڑیاں~$\Lang{L}$ وچ نہیں آندیاں، اک
$\Lang{L}$-ساخت حاصل کیتی جا سکدی اے۔
ایتھے زبان دیاں علامتاں تے اوہناں دیاں تشریحاں گھٹدیاں یا ودھدیاں نیں؛
دائرہ اوہی رہندا اے، دائرے دے عنصر کڈن یا شامل کرن دی گل نہیں۔

\begin{defn}
\ollabel{defn:reduct}
فرض کرو $\Lang L \subseteq \Lang L'$، $\Struct M$ اک
$\Lang L$-!!{structure} اے تے $\Struct M'$ اک $\Lang L'$-!!{structure}
اے۔ $\Struct M$، $\Struct M'$ دا $\Lang L$ تک \emph{گھٹایا روپ}
(reduct) اے، تے $\Struct M'$، $\Struct M$ دا $\Lang L'$ تک
\emph{ودھایا روپ} (expansion) اے، اُدوں تے صرف اُدوں جدوں
\begin{enumerate}
\item $\Domain{M} = \Domain{M'}$
\item ہر !!{constant}~$c \in \Lang L$ لئی، $\Assign{c}{M} =
  \Assign{c}{M'}$۔
\item ہر !!{function}~$f \in \Lang L$ لئی، $\Assign{f}{M} =
  \Assign{f}{M'}$۔
\item ہر !!{predicate}~$P \in \Lang L$ لئی، $\Assign{P}{M} =
  \Assign{P}{M'}$۔
\end{enumerate}
\end{defn}

\begin{prop}
\ollabel{prop:reduct}
جے اک $\Lang{L}$-!!{structure}~$\Struct{M}$، اک
$\Lang{L'}$-!!{structure}~$\Struct{M'}$ دا گھٹایا روپ اے، تاں
ہر $\Lang{L}$-!!{sentence}s~$!A$ لئی
\[
\Sat{M}{!A} \text{ اُدوں تے صرف اُدوں } \Sat{M'}{!A}.
\]
\end{prop}

\begin{proof}
  مشق۔
\end{proof}

\begin{prob}
\olref[mod][bas][red]{prop:reduct} ثابت کرو۔
\end{prob}

\begin{defn}
جدوں ساڈے کول اک $\Lang{L}$-ساخت $\Struct{M}$ ہووے، تے
$\Lang{L'} =
\Lang{L} \cup \{P\}$، $\Lang{L}$ وچ صرف اک $n$-تھانواں والا
!!{predicate}~$P$ شامل کرکے ملی ودھی زبان ہووے، تے
$R \subseteq \Domain{M}^n$ اک $n$-تھانواں والا تعلق ہووے، تاں
$\Expan{M}{R}$ اوس ودھائے روپ~$\Struct{M'}$ لئی لکھدے آں
جو~$\Struct{M}$ دا اے تے جس وچ $\Assign{P}{M'} = R$۔
\end{defn}
\endgroup
% /OLP-0184

% OLP-0185 content/model-theory/basics/substructures.tex
\begingroup
\def\olfilename{substructures}
\expandafter\def\csname ol@sectioncs\endcsname{section}

\olfileid{mod}{bas}{sub}
\olsection{ذیلی \printtoken{p}{structure}}

اک !!a{structure}~$\Struct{M}$ دا !!{domain} دوجی ساخت~$\Struct{M'}$
دے دائرے دا ذیلی سیٹ ہو سکدا اے۔ پر $\Struct{M}$ نوں
$\Struct{M'}$ دا ”حصہ“ اُدوں ای سمجھنا چاہیدا اے جدوں نہ صرف
$\Domain{M}
\subseteq \Domain{M'}$ ہووے، بلکہ $\Struct{M}$ تے $\Struct{M'}$
زبان دیاں علامتاں دی تشریح وچ گھٹ توں گھٹ سانجھے حصے~$\Domain{M}$
اُتے اک دوجے نال ”متفق“ وی ہون۔
ایتھے زبان اوہی رہندی اے، پر دائرہ چھوٹا ہو سکدا اے؛
پچھلے گھٹائے روپ وچ زبان گھٹدی سی، دائرہ نہیں۔

\begin{defn}
\ollabel{defn:substructure}
!!{structure}s $\Struct M$ تے $\Struct M'$ اکّو زبان~$\Lang L$
لئی دتیاں ہون۔ اسیں $\Struct M$ نوں $\Struct M'$ دی
\emph{ذیلی !!{structure}} (substructure)، تے $\Struct M'$ نوں
$\Struct M$ دی \emph{توسیع} (extension) آکھدے آں، جس نوں
$\Struct M \substruct \Struct M'$ لکھدے آں، اُدوں تے صرف اُدوں جدوں
\begin{enumerate}
\item $\Domain{M} \subseteq \Domain{M'}$،
\item ہر مستقل $c \in \Lang L$ لئی، $\Assign{c}{M} =
    \Assign{c}{M'}$؛
\item ہر $n$-تھانواں والے !!{function} $f \in \Lang L$ لئی
  $\Assign{f}{M}(a_1, \dots, a_n) = \Assign{f}{M'}(a_1, \dots, a_n)$
  سارے $a_1$، \dots، $a_n \in \Domain{M}$ لئی۔
\item ہر $n$-تھانواں والے !!{predicate} $R \in \Lang L$ لئی، $\langle
  a_1, \dots, a_n\rangle \in \Assign{R}{M}$ اُدوں تے صرف اُدوں جدوں
  $\langle a_1, \dots,
  a_n\rangle \in \Assign{R}{M'}$ ہووے، سارے $a_1$، \dots، $a_n \in
  \Domain{M}$ لئی۔
\end{enumerate}
دائرہ غیر خالی اے، ہر مستقل دی قدر ایس وچ اے تے فنکشن دی چھوٹے
دائرے دیاں قدراں اُتے ورتوں دا نتیجہ وی ایتھے ای رہندا اے۔
ایہ گلاں دتی ہوئی ساخت دا دائرہ تے کل فنکشن ہون نال لازم نیں۔
\end{defn}

% PNB-SEM-035: پہلے درجے دی ساخت دا دائرہ غیر خالی اے؛خالی ذیلی سیٹ لئی ایہ دعویٰ نہیں۔
\begin{rem}
\ollabel{rem:substructure}
جے زبان وچ کوئی مستقل یا !!{function}s نہ ہون، تاں ہر غیر خالی $N
\subseteq \Domain{M}$ اک ذیلی !!{structure}~$\Struct{N}$
متعین کردا اے جو $\Struct M$ دی اے تے جس دا !!{domain}~$\Domain{N} = N$
اے؛ ایس لئی $\Assign{R}{N} =
\Assign{R}{M} \cap N^n$ رکھدے آں۔
\end{rem}

% ایس بارے کجھ ثابت کریے؟ مثالاں؟
\endgroup
% /OLP-0185

% OLP-0186 content/model-theory/basics/overspill.tex
\begingroup
\def\olfilename{overspill}
\expandafter\def\csname ol@sectioncs\endcsname{section}

\olfileid{mod}{bas}{ove}
\olsection{حد توں اگے پھیلاؤ \pnbLatin{(Overspill)}}

\begin{thm}
\ollabel{overspill} جے جملیاں دے سیٹ $\Gamma$ دے من مانے وڈے
محدود نمونے ہون، تاں اوہدا اک لامتناہی نمونہ وی ہوندا اے۔
\end{thm}

\begin{proof}
$\Gamma$ دی زبان وچ قابلِ شمار لامتناہی تعداد دے نویں مستقل
$c_0$، $c_1$، \dots شامل کرو تے سیٹ $\Gamma \cup \{c_i \neq c_j :
i \neq j\}$ ویکھو۔ ایہ کہنا کہ $\Gamma$ دے من مانے وڈے محدود نمونے
نیں، ایس دا مطلب اے کہ ہر $m >0$ لئی کوئی $n\ge m$ اے جیہدے لئی
$\Gamma$ دا ایداں دا نمونہ اے جس دے دائرے وچ $n$ عنصر نیں۔ ودھے سیٹ دا
کوئی وی محدود حصہ چنو۔ اوہدے وچ صرف محدود تعداد دے نویں مستقل آندے
نیں؛ اصل سیٹ دا اِنّے عنصر رکھن والا یا اوہدوں وڈا محدود نمونہ چن کے
اوہناں مستقل نشانیاں نوں وکھ وکھ عنصر دیو۔ باقی نویں مستقل اک من مرضی
دے عنصر نوں نام دے سکدے نیں؛ اصل زبان دی تشریح نہیں بدلنی۔ ایس لئی
چنے محدود حصے دے اصل جملے تے عدمِ برابریاں اکٹھیاں سچ ہون گیاں۔ ایس توں
ایہ نکلدا اے کہ $\Gamma \cup \{c_i \neq
c_j : i \neq j\}$ محدود طور تے پورا ہون جوگا اے۔ محدود گواہی دے
قضیے مطابق $\Gamma
\cup \{c_i \neq c_j : i \neq j\}$ دا اک نمونہ $\Struct M$
اے جس دا دائرہ لازماً لامتناہی اے، کیوں جے ایہ ساریاں عدمِ برابریاں
$c_i \neq c_j$ پوری کردا اے۔
آخر وچ نویں مستقل نشانیاں دی تشریح بھلا کے اصل زبان والی ساخت لے لو؛
دائرہ اوہی لامتناہی رہندا اے تے اصل سیٹ دے سارے جملے وی سچے رہندے نیں۔
\end{proof}

\begin{prop}
\ollabel{inf-not-fo}
کسی وی پہلے درجے دی زبان دا کوئی جملہ $!A$ ایسا نہیں کہ اوہ
!!a{structure}~$\Struct M$ وچ اُدوں تے صرف اُدوں سچ ہووے جدوں
اوس !!{structure} دا دائرہ $\Domain{M}$ لامتناہی ہووے۔
ایتھے اکّو جملے نال ساریاں تے صرف لامتناہی ساختاں دی پہچان دی گل اے؛
جملیاں دے پورے لامتناہی سیٹ نال ایہ شرط بیان کر سکّن دی نفی نہیں۔
\end{prop}

\begin{proof}
جے ایداں دا کوئی $!A$ ہوندا، تاں اوہدی نفی $\lnot !A$ ہر تے
صرف محدود !!{structure}s وچ سچی ہوندی۔ ایس لئی اوہدے من مانے وڈے
محدود نمونے ہوندے، پر کوئی لامتناہی نمونہ نہ ہوندا؛ ایہ
\olref{overspill} دے خلاف اے۔
\end{proof}
\endgroup
% /OLP-0186

% OLP-0187 content/model-theory/basics/isomorphism.tex
\begingroup
\def\olfilename{isomorphism}
\expandafter\def\csname ol@sectioncs\endcsname{section}

\olfileid{mod}{bas}{iso}
\olsection{ہم شکل ساختاں}

پہلے درجے دیاں !!{structure}s دو طرحاں نال اک دوجے ورگیاں ہو
سکدیاں نیں۔ اک طریقہ ایہ اے کہ اوہ اکّو !!{sentence}s نوں سچ
بناندیاں ہون۔ ایہو جیہیاں !!{structure}s نوں اسیں \emph{جملہ وار
ہم ارز} \pnbLatin{(elementarily equivalent)} آکھدے آں۔ پر دو ساختاں بہت وکھ
ہو کے وی اکّو !!{sentence}s سچ بنا سکدیاں نیں---مثال لئی اک
!!{enumerable} ہو سکدی اے تے دوجی نہیں۔ ایس دی وجہ ایہ اے کہ
!!a{structure} دیاں بہت ساریاں خاصیتاں پہلے درجے دیاں زباناں وچ
بیان نہیں ہو سکدیاں: یا زبان کافی وسیع نہیں ہوندی، یا پہلے درجے
دی منطق دیاں بنیادی حدّاں، جیویں L\"owenheim--Skolem قضیہ، راہ وچ
ہوندیاں نیں۔ ایس لئی !!{structure}s دے اک دوجے ورگیاں ہون دا اک
ہور، زیادہ سخت مفہوم ایہ اے کہ اوہ بنیادی طور اُتے اکّو جیہیاں ہون:
انہاں نوں بناؤن والیاں چیزاں وکھ ہوون، پر ساختی خاصیتاں وکھ نہ
ہون۔ ایس مفہوم نوں ٹھیک ٹھیک بیان کرن دا طریقہ \emph{ہم شکلی}
(isomorphism) دا تصور اے۔

\begin{defn}
\ollabel{defn:elem-equiv}
دو !!{structure}s $\Struct{M}$ تے $\Struct M'$، جیہڑیاں اکو
!!{language}~$\Lang{L}$ دیاں ہون، دتیاں ہون۔ اسیں $\Struct{M}$ نوں
$\Struct M'$ دے نال \emph{جملہ وار ہم ارز} آکھدے آں، جس نوں
$\Struct{M} \equiv \Struct M'$ لکھدے آں، اُدوں تے صرف اُدوں جدوں
ہر !!{sentence}~$!A$، جیہڑا~$\Lang{L}$ دا ہووے، اوہدے لئی $\Sat{M}{!A}$ اُدوں تے
صرف اُدوں جدوں $\Sat{M'}{!A}$۔
\end{defn}

\begin{defn}
\ollabel{defn:isomorphism} دو !!{structure}s $\Struct{M}$ تے
$\Struct M'$، جیہڑیاں اکو !!{language}~$\Lang L$ دیاں ہون، دتیاں ہون۔ اسیں
$\Struct{M}$ نوں~$\Struct M'$ دے نال \emph{ہم شکل} آکھدے آں، جس
نوں $\Struct{M}
\simeq \Struct M'$ لکھدے آں، اُدوں تے صرف اُدوں جدوں اک فنکشن
$h \colon
\Domain{M} \to \Domain{M'}$ ہووے جیہڑا ایہ شرطاں پوری کرے:
\begin{enumerate}
\item $h$ !!{injective} اے: جے $h(x) =
  h(y)$ تاں $x = y$؛
\item $h$ !!{surjective} اے: ہر $y \in \Domain{M'}$ لئی کوئی
  $x \in \Domain{M}$ اے جیہدے لئی $h(x) = y$؛
\item \ollabel{defn:iso-const}ہر !!{constant} $c$ لئی:
  $h(\Assign{c}{M}) = \Assign{c}{M'}$؛
\item \ollabel{defn:iso-pred}ہر $n$-تھانواں والے !!{predicate}~$P$ لئی:
  \[
  \tuple{a_1, \dots, a_n}\in \Assign{P}{M} \quad\text{اُدوں تے صرف اُدوں}\quad
  \tuple{h(a_1), \dots, h(a_n)} \in \Assign{P}{M'};
  \]
\item \ollabel{defn:iso-func}ہر $n$-تھانواں والے !!{function} $f$ لئی:
  \[
  h(\Assign{f}{M}(a_1, \dots, a_n)) =
  \Assign{f}{M'}(h(a_1), \dots, h(a_n)).
  \]
\end{enumerate}
\end{defn}

\begin{thm}
\ollabel{thm:isom}
جے $\Struct{M} \iso \Struct M'$ تاں $\Struct{M} \elemequiv
\Struct{M'}$۔
\end{thm}

\begin{proof}
فرض کرو $h$، $\Struct{M}$ توں $\Struct M'$ اُتے ہم شکلی اے۔ ہر
متغیر قدر بندی~$s$ لئی، $h \circ s$، $h$ تے $s$ دا مرکب اے؛ یعنی
$\Struct{M'}$ وچ اوہ متغیر قدر بندی جیہدے لئی $(h \circ s)(x) = h(s(x))$۔
$t$ تے $!A$ اُتے استقرا نال ایہ زیادہ مضبوط دعوے ثابت ہو سکدے نیں:
\begin{enumerate}
  \item[a.] $h(\Value{t}{M}[s]) = \Value{t}{M'}[h\circ s]$۔
  \item[b.] $\Sat{M}{!A}[s]$ اُدوں تے صرف اُدوں جدوں $\Sat{M'}{!A}[h \circ s]$۔
\end{enumerate}
پہلا دعویٰ~$t$ دی پیچیدگی اُتے استقرا نال ثابت ہوندا اے۔
\begin{enumerate}
\item جے $t \ident c$، تاں $\Value{c}{M}[s] = \Assign{c}{M}$ تے
  $\Value{c}{M'}[h \circ s] = \Assign{c}{M'}$۔ ایس لئی
  $h(\Value{t}{M}[s]) = h(\Assign{c}{M}) = \Assign{c}{M'}$ (
  \olref{defn:isomorphism} دی \olref{defn:iso-const} مطابق) $=
  \Value{t}{M'}[h \circ s]$۔
\item جے $t \ident x$، تاں $\Value{x}{M}[s] = s(x)$ تے
  $\Value{x}{M'}[h \circ s] = h(s(x))$۔ ایس لئی $h(\Value{x}{M}[s]) =
  h(s(x)) = \Value{x}{M'}[h \circ s]$۔
\item جے $t \ident f(t_1, \dots, t_n)$، تاں
  \begin{align*}
    \Value{t}{M}[s] & = \Assign{f}{M}(\Value{t_1}{M}[s], \dots, \Value{t_n}{M}[s]) \quad\text{تے}\\
  \Value{t}{M'}[h \circ s] & = \Assign{f}{M'}(\Value{t_1}{M'}[h \circ
    s], \dots, \Value{t_n}{M'}[h \circ s]).
  \end{align*}
  استقرائی مفروضہ ایہ اے کہ ہر $i$ لئی، $h(\Value{t_i}{M}[s])
  = \Value{t_i}{M'}[h\circ s]$۔ ایس لئی
  \begin{align}
    h(\Value{t}{M}[s])
    & = h(\Assign{f}{M}(\Value{t_1}{M}[s], \dots, \Value{t_n}{M}[s]) \notag\\
    & = \Assign{f}{M'}(h(\Value{t_1}{M}[s]), \dots,
    h(\Value{t_n}{M}[s])) \ollabel{iso-1}\\
    & = \Assign{f}{M'}(\Value{t_1}{M'}[h \circ s], \dots,
    \Value{t_n}{M'}[h \circ s]) \ollabel{iso-2}\\
    & = \Value{t}{M'}[h\circ s] \notag
  \end{align}
  ایتھے \olref{iso-1}، \olref{defn:isomorphism} دی
  \olref{defn:iso-func} توں تے \olref{iso-2} استقرائی مفروضے توں نکلدا اے۔
\end{enumerate}
حصہ (b) مشق دے طور اُتے چھڈیا گیا اے۔

جے $!A$ اک جملہ اے، تاں متغیر قدر بندیاں~$s$ تے $h \circ s$
غیر متعلقہ نیں، تے سانوں $\Sat{M}{!A}$ اُدوں تے صرف اُدوں جدوں
$\Sat{M'}{!A}$ ملدا اے۔
\end{proof}

\begin{prob}
\olref[mod][bas][iso]{thm:isom} دے حصہ (b) دا ثبوت پوری تفصیل نال
کرو۔ خاص طور اُتے دسو کہ \olref[mod][bas][iso]{defn:isomorphism}
وچ ہم شکلی دیاں پنج خاصیتاں وچوں ہر اک کتھے ورتی جاندی اے۔
\end{prob}

\begin{defn}
ساخت $\Struct{M}$ دی \emph{خود ہم شکلی} (automorphism)،
$\Struct{M}$ توں اوہدے اپنے اُتے اک ہم شکلی اے۔
\end{defn}

\begin{prob}
دکھاؤ کہ ہر ساخت $\Struct{M}$ لئی، جے $X$، $\Struct{M}$ دا
قابلِ تعریف ذیلی سیٹ ہووے، جیہدے تعریفی فارمولے وچ زبان توں باہر دے
اضافی پیرامیٹر نہ ہون، تے $h$، $\Struct{M}$ دی خود ہم شکلی ہووے،
تاں $X
= \Setabs{h(x)}{x \in X}$ (یعنی $X$، $h$ دے تحت قائم رہندا اے)۔
\end{prob}
\endgroup
% /OLP-0187

% OLP-0188 content/model-theory/basics/theory-of-m.tex
\begingroup
\def\olfilename{theory-of-m}
\expandafter\def\csname ol@sectioncs\endcsname{section}

\olfileid{mod}{bas}{thm}

\section{اک \printtoken{S}{structure} دا نظریہ}

ہر !!{structure}~$\Struct{M}$ کجھ !!{sentence}s نوں سچ تے کجھ نوں
جھوٹا بناؤندی اے۔ اوہدے سچ بنائے سارے !!{sentence}s دے سیٹ نوں
اوہدا \emph{نظریہ} آکھیا جاندا اے۔ ایہ سیٹ واقعی اک نظریہ اے، کیوں جے
اوہدے توں جیہڑی وی گل منطقی طور اُتے نکلے، اوہ اوہدے سارے نمونیاں
وچ سچ ہونی چاہیدی اے، جنہاں وچ~$\Struct{M}$ وی شامل اے۔ ایس لئی
اوہ نتیجہ وی ایس سارے سچ جملیاں دے سیٹ وچ آ جاندا اے۔

\begin{defn}
  !!a{structure}~$\Struct M$ دتی ہووے۔ $\Struct{M}$ دا
  \emph{نظریہ}، اوہناں !!{sentence}s دا سیٹ $\Theory{M}$ اے
  جیہڑے $\Struct{M}$ وچ سچ نیں؛ یعنی $\Theory{M} =
  \Setabs{!A}{\Sat{M}{!A}}$۔
\end{defn}

اسیں ``نظریہ'' دا لفظ غیر رسمی طور اُتے اوہناں !!{sentence}s دے
سیٹاں لئی وی ورتدے آں جنہاں دی کوئی مِتھی ہوئی تعبیر ہووے، چاہے اوہ
استنتاج دے لحاظ نال بند ہوون یا نہ ہوون۔

\begin{prop}
ہر $\Struct{M}$ لئی، $\Theory{M}$ مکمل اے۔
\end{prop}

\begin{proof}
ہر !!{sentence}~$!A$ لئی یا $\Sat{M}{!A}$ یا $\Sat{M}{\lnot !A}$،
ایس لئی یا $!A \in \Theory{M}$ یا $\lnot !A \in \Theory{M}$۔
\end{proof}

\begin{prop}\ollabel{prop:equiv}
  جے $\Struct{N} \models !A$ ہر $!A \in \Theory{M}$ لئی، تاں
  $\Struct{M} \elemequiv \Struct{N}$۔ ایتھے دوہاں ساختاں دی زبان اکّو اے۔
\end{prop}

\begin{proof}
کیوں جے $\Sat{N}{!A}$ ہر $!A \in \Theory{M}$ لئی، ایس لئی
$\Theory{M} \subseteq \Theory{N}$۔ جے $\Sat{N}{!A}$، تاں
$\Sat/{N}{\lnot !A}$؛ ایس لئی $\lnot !A
\notin \Theory{M}$۔ جے نفی پہلے نظریے وچ ہوندی، تاں دوسری ساخت وچ وی
سچ ہونی سی، پر ایتھے اوہ سچ نہیں۔ کیوں جے $\Theory{M}$ مکمل اے، ایس لئی $!A \in
\Theory{M}$۔ سو $\Theory{N} \subseteq \Theory{M}$، تے سانوں
$\Struct{M} \elemequiv \Struct{N}$ ملدا اے۔
\end{proof}

\begin{rem}\ollabel{remark:R}
  $\Struct{R} = \langle\Real, <\rangle$ نوں ویکھو: ایہ اوہ
  !!{structure} اے جیہدا دائرہ حقیقی اعداد دا سیٹ $\Real$ اے،
  تے اوہ !!{language} وچ اے جیہدے وچ صرف اک دو-تھانواں والا
  !!{predicate} اے، جس دی تعبیر حقیقی اعداد اُتے $<$ تعلق اے۔
  ظاہر اے کہ $\Struct{R}$ !!{nonenumerable} اے؛ پر چونکہ
  $\Theory{R}$ واضح طور اُتے موافق اے، L\"owenheim--Skolem قضیے
  مطابق ایہدا اک !!a{enumerable} نمونہ اے، مثلاً $\Struct{S}$؛
  تے \olref{prop:equiv} توں $\Struct{R}
  \equiv \Struct{S}$۔ اک دا دائرہ قابلِ شمار اے تے دوجے دا نہیں،
  ایس لئی اوہناں وچ ہم شکلی لئی درکار اک بہ اک مطابقت نہیں ہو سکدی۔
  ہور ایہ کہ $\Struct{R}$ تے $\Struct{S}$
  ہم شکل نہیں، ایس لئی ایہ مثال دکھاؤندی اے کہ
  \olref[iso]{thm:isom} دا الٹا دعویٰ عام طور اُتے درست نہیں۔
\end{rem}
\endgroup
% /OLP-0188

% OLP-0189 content/model-theory/basics/partial-iso.tex
\begingroup
\def\olfilename{partial-iso}
\expandafter\def\csname ol@sectioncs\endcsname{section}

\olfileid{mod}{bas}{pis}
\section{جزوی ہم شکلیاں}

\begin{defn}
  دو !!{structure}s $\Struct{M}$ تے $\Struct{N}$ دتیاں ہون۔
  $\Struct{M}$ توں $\Struct{N}$ ول اک \emph{جزوی ہم شکلی}
  \pnbLatin{(partial isomorphism)} اک محدود جزوی فنکشن $p$ اے، جیہڑا
  $\Domain M$ دے عناصر نوں دلیل بنا کے $\Domain N$ وچ قدراں دیندا
  اے، تے اپنے دائرے اُتے \olref[iso]{defn:isomorphism} دیاں
  ہم شکلی والیاں شرطاں پوری کردا اے:
  \begin{enumerate}
  \item $p$ !!{injective} اے؛
  \item ہر !!{constant}~$c$ لئی: جے $p(\Assign{c}{M})$ دی قدر
    متعین ہووے، تاں $p(\Assign{c}{M}) = \Assign{c}{N}$؛
  \item ہر $n$-تھانواں والے !!{predicate} $P$ لئی: جے $a_1$، \dots، $a_n$
    سارے $p$ دے دائرے وچ ہوون، تاں $\langle a_1, \dots, a_n\rangle \in
    \Assign P M$ اُدوں تے صرف اُدوں جدوں $\langle p(a_1), \dots, p(a_n) \rangle
    \in \Assign P N$؛
  \item ہر $n$-تھانواں والے !!{function} $f$ لئی: جے $a_1$، \dots، $a_n$
    سارے $p$ دے دائرے وچ ہوون، تاں $p(\Assign f M (a_1, \dots,a_n))
    = \Assign f N (p(a_1), \dots, p(a_n))$۔
  \end{enumerate}
  $p$ دے محدود ہون دا مطلب ایہ اے کہ $\dom{p}$ محدود اے۔
  دوہاں ساختاں دی زبان اکّو اے۔ فنکشن والی شرط وچ، دائرے دے عناصر
  اُتے فنکشن دی قدر وی جزوی نقشے دے دائرے وچ ہونی چاہیدی اے، نہیں تاں
  کھبے پاسے دی قدر متعین نہیں ہوئے گی۔ پورے ہدف دائرے اُتے ہون دی شرط نہیں۔
\end{defn}

دھیان دیو کہ خالی فنکشن~$\emptyset$ ہر دو !!{structure}s وچکار
ہمیشہ اک جزوی ہم شکلی اے۔
ایتھے کتاب دی زبان مطابق ہر محمول تے فنکشن دے تھانواں دی تعداد مثبت اے؛
مستقل دی شرط وی صرف اوہدوں لگدی اے جدوں اوہدی قدر نقشے دے دائرے وچ ہووے۔

\begin{defn}\ollabel{defn:partialisom}
  دو !!{structure}s $\Struct{M}$ تے $\Struct{N}$
  \emph{جزوی طور اُتے ہم شکل} نیں، جس نوں $\Struct{M} \iso[p]
  \Struct{N}$ لکھدے آں، اُدوں تے صرف اُدوں جدوں جزوی ہم شکلیاں
  دا اک غیر خالی سیٹ $I$ ہووے، جیہڑیاں $\Struct{M}$ تے $\Struct{N}$
  وچکار ہون، تے ایہ سیٹ \emph{اگے-پیچھے} (back-and-forth) خاصیت پوری کرے:
  \begin{enumerate}
  \item (\emph{اگے}) ہر $p \in I$ تے $a \in \Domain M$ لئی
    کوئی $q \in I$ اے جیہدے لئی $p \subseteq q$، تے $a$،
    $q$ دے دائرے وچ اے؛
  \item (\emph{پیچھے}) ہر $p \in I$ تے $b \in \Domain N$ لئی
    کوئی $q \in I$ اے جیہدے لئی $p \subseteq q$، تے $b$،
    $q$ دیاں حاصل قدراں وچ اے۔
  \end{enumerate}
\end{defn}

\begin{thm}\ollabel{thm:p-isom1}
  جے $\Struct{M} \iso[p] \Struct{N}$ تے $\Struct{M}$ تے
  $\Struct{N}$ !!{enumerable} ہون، تاں $\Struct{M} \iso
  \Struct{N}$۔
\end{thm}

\begin{proof}
  کیوں جے $\Struct{M}$ تے $\Struct{N}$ !!{enumerable} نیں، فرض کرو
  $\Domain{M} =
  \{a_0, a_1, \ldots \}$ تے $\Domain{N} = \{b_0, b_1, \ldots \}$۔
  جے کوئی دائرہ محدود ہووے، تاں سارے عنصر لسٹ کرکے آخری عنصر نوں دہراؤ؛
  کتاب مطابق دائرہ غیر خالی اے۔ لسٹ دے ہر اصل عنصر نوں ہن وی اک مرحلہ ملدا اے۔
  کسے وی $p_0 \in I$ توں شروع کرکے اسیں جزوی ہم شکلیاں دا
  اک ودھدا سلسلہ $p_0 \subseteq p_1 \subseteq p_2
  \subseteq \cdots$ ایویں مقرر کردے آں:
  \begin{enumerate}
  \item جے $n+1$ طاق ہووے، یعنی $n = 2r$، تاں اگے والی خاصیت
    نال کوئی $p_{n+1} \in I$ لبھو جیہدے لئی $p_n \subseteq p_{n+1}$
    تے $a_r$، $p_{n+1}$ دے دائرے وچ ہووے؛
  \item جے $n+1$ جفت ہووے، یعنی $n =2r+1$، تاں پیچھے والی خاصیت
    نال کوئی $p_{n+1} \in I$ لبھو جیہدے لئی $p_n \subseteq p_{n+1}$
    تے $b_r$، $p_{n+1}$ دیاں حاصل قدراں وچ ہووے۔
  \end{enumerate}
ایس طرح پہلی اگلی چال پہلی لسٹ دے شروع والے عنصر نوں، تے پہلی پچھلی چال
دوجی لسٹ دے شروع والے عنصر نوں شامل کردی اے؛ کوئی شروع والا عنصر چھٹدا نہیں۔
ہن جے اسیں ایہ رکھئیے:
\[
p = \bigcup_{n\ge 0} p_n,
\]
تاں $p$، $\Struct{M}$ تے $\Struct{N}$ وچکار اک ہم شکلی اے۔
\end{proof}

\begin{prob}
  پوری تفصیل نال دکھاؤ کہ
  \olref[mod][bas][pis]{thm:p-isom1} وچ مقرر کیتا $p$ واقعی اک
  ہم شکلی اے۔
\end{prob}

\begin{thm}\ollabel{thm:p-isom2}
  فرض کرو $\Struct{M}$ تے $\Struct{N}$ اک خالص تعلقی
  !!{language} دیاں !!{structure}s نیں (یعنی !!a{language} جیہدے
  وچ صرف !!{predicate}s ہون، کوئی !!{function}s یا مستقل نہ ہون)۔
  تاں جے $\Struct{M} \iso[p] \Struct{N}$، ایہ وی ہووے گا کہ
  $\Struct{M} \elemequiv
  \Struct{N}$۔
\end{thm}

\begin{proof}
  ایس دعوے وچ جزوی ہم شکلی اوپر دتے اگے-پیچھے والے غیر خالی سیٹ وچوں
  ہونی چاہیدی اے۔ فارمولے دے سارے آزاد متغیر تھلے دتی متغیر لسٹ وچ ہون۔
  !!{formula}s اُتے استقرا کرکے دکھایا جاندا اے کہ جے $a_1$، \dots، $a_n$ تے
  $b_1$، \dots، $b_n$ ایسے ہون کہ اگے-پیچھے والے سیٹ دی اک جزوی ہم شکلی $p$ ہر
  $a_i$ نوں $b_i$ ول بھیجے تے $s_1(x_i) =a_i$ تے $s_2(x_i) =b_i$
  ($i =1$، \dots،~$n$ لئی)، تاں $\Sat{M}{!A}[s_1]$ اُدوں تے
  صرف اُدوں جدوں $\Sat{N}{!A}[s_2]$۔ $n=0$ دا حال
  $\Struct{M} \elemequiv \Struct{N}$ دیندا اے۔
  اکیلی جزوی ہم شکلی، جیہڑی ایس توسیع والے سیٹ دا حصہ نہ ہووے، کافی نہیں۔
\end{proof}

\begin{rem}
جے !!{function}s وی ہون، تاں پچھلا نتیجہ فیر وی درست اے، پر
$p$ توں پیدا ہون والی اوہ ہم شکلی ویکھنی پیندی اے جیہڑی
$\Struct{M}$ دی، $a_1$، \dots، $a_n$ توں پیدا ہوئی
ذیلی!!{structure} تے $\Struct{N}$ دی، $b_1$، \dots، $b_n$
توں پیدا ہوئی ذیلی!!{structure} وچکار ہووے۔
ایتھے وی اگے-پیچھے والا پورا مفروضہ قائم اے، صرف اکیلا جزوی نقشہ کافی نہیں۔
جملیاں دے خالی متغیر لسٹ والے حال وچ غیر خالی نقشے تک توسیع لے کے کم کیتا
جا سکدا اے؛ فارمولے وچ آؤن والیاں مستقل قدراں وی متعلقہ توسیع وچ شامل کرو۔
\end{rem}

پچھلے نتیجے نوں مرحلیاں وچ وی وکھایا جا سکدا اے: !!a{formula}
وچ اک دوجے دے اندر آئے مقدار نماواں دی تعداد تے متعلقہ جزوی
ہم شکلیاں نوں کنی واری ودھایا جا سکدا اے، انہاں وچ ربط قائم کرکے۔

\begin{defn}
  ہر !!{formula}~$!A$ لئی $!A$ دا \emph{مقدار نمائی درجہ}
  \pnbLatin{(quantifier rank)}، جیہنوں $\QuantRank{!A} \in \Nat$ لکھدے آں،
  استقرائی طور اُتے $!A$ وچ اک دوجے دے اندر آئے مقدار نماواں دی
  سبھ توں وڈی تعداد دے طور اُتے مقرر ہوندا اے۔ دو
  !!{structure}s $\Struct{M}$ تے $\Struct{N}$ \emph{$n$-ہم ارز}
  نیں، جیہنوں $\Struct{M} \elemequiv[n] \Struct{N}$ لکھدے آں،
  جے اوہ مقدار نمائی درجہ~$n$ توں گھٹ یا اوہدے برابر رکھن والے
  سارے جملیاں بارے متفق ہون۔
\end{defn}

\begin{prop}\ollabel{prop:qr-finite}
  $\Lang{L}$ اک محدود خالص تعلقی !!{language} ہووے، یعنی ایسی
  !!{language} جیہدے وچ محدود گنتی دے !!{predicate}s تے
  !!{constant}s ہون، تے کوئی !!{function}s نہ ہوون۔ تاں ہر
  $n \in \Nat$ لئی !!{language} $\Lang{L}$ وچ پہلے درجے دے
  !!{sentence}s، جنہاں دا مقدار نمائی درجہ $n$ توں ودھ نہ ہووے،
  منطقی ہم ارزی تک صرف محدود گنتی وچ نیں۔
  ایس قضیے وچ مستقل نشانیاں دی محدود تعداد دی اجازت اے؛ پچھلے قضیے والی
  بنا مستقل والی قید ایتھے نہیں لائی گئی۔ جملیاں دے نال ایہ نتیجہ اوہناں
  فارمولیاں لئی وی درست اے جیہڑے اک مقرر محدود لسٹ دے متغیراں توں باہر
  کوئی آزاد متغیر نہیں رکھدے۔ محدود گنتی ہم ارزی جماعتاں دی اے، لفظی فارمولیاں دی نہیں۔
\end{prop}

\begin{proof}
  $n$ اُتے استقرا نال؛ فارمولیاں والے روپ وچ آزاد متغیراں دی محدود لسٹ مقرر رکھو۔
\end{proof}

\begin{defn}
  !!a{structure}~$\Struct{M}$ دتی ہووے؛ $\Domain M^{<\omega}$،
  $\Domain{M}$ دے عناصر اُتے بنے سارے محدود سلسلیاں دا سیٹ ہووے۔
  اسیں $\mathbf{a},
  \mathbf{b}, \mathbf{c}, \ldots$ نوں عناصر دے محدود سلسلیاں
  لئی ورتدے آں۔ جے $\mathbf{a} \in \Domain{M}^{<\omega}$ تے
  $a \in \Domain{M}$، تاں $\mathbf{a}a$، $\mathbf{a}$ نال
  $a$ دا \emph{جوڑ} (concatenation) ظاہر کردا اے۔
\end{defn}

\begin{defn}
  !!{structure}s $\Struct{M}$ تے $\Struct{N}$ دتیاں ہون۔
  اسیں برابر لمبائی والے سلسلیاں وچکار تعلقات
  $I_n \subseteq \Domain M^{<\omega} \times \Domain N^{<\omega}$
  نوں $n$ اُتے استقرا نال ایویں مقرر کردے آں:
   \begin{enumerate}
   \item $I_0(\mathbf{a},\mathbf{b})$ اُدوں تے صرف اُدوں جدوں
     $\mathbf{a}$ تے $\mathbf{b}$، $\Struct{M}$ تے
     $\Struct{N}$ وچ اکّو جوہری !!{formula}s پورے کردے ہون؛ یعنی
     جے $s_1(x_i) = a_i$ تے $s_2(x_i) =
     b_i$، تے $!A$ جوہری ہووے تے اوہدے سارے !!{variable}s،
     $x_1$، \dots،~$x_n$ وچوں ہون، تاں $\Sat{M}{!A}[s_1]$ اُدوں تے
     صرف اُدوں جدوں~$\Sat{N}{!A}[s_2]$۔
   \item $I_{n+1} (\mathbf{a},\mathbf{b})$ اُدوں تے صرف اُدوں جدوں ہر
     $a\in \Domain M$ لئی کوئی $b\in \Domain N$ ہووے جیہدے لئی $I_n
     (\mathbf{a}a,\mathbf{b}b)$، تے الٹی سمت وچ وی ایویں ہووے۔
   \end{enumerate}
  جوہری شرط وچ متغیر لسٹ دی آخری گنتی سلسلے دی لمبائی اے؛ اوہنوں
  تعلق دے استقرائی درجے دی گنتی نال نہ ملاؤ۔ دوہاں سلسلیاں دی لمبائی برابر اے۔
\end{defn}


\begin{defn}
  $\Struct{M} \approx_n \Struct{N}$ لکھو جے
  $I_n(\emptyseq,\emptyseq)$، $\Struct{M}$ تے
  $\Struct{N}$ لئی درست ہووے (جتھے $\emptyseq$ خالی سلسلہ اے)۔
\end{defn}

\begin{thm}\ollabel{thm:b-n-f}
  $\Lang{L}$ اک خالص تعلقی !!{language} ہووے۔ تاں $I_n
  (\mathbf{a},\mathbf{b})$ توں نکلدا اے کہ ہر $!A$، جیہدے لئی
  $\QuantRank{!A} \le n$، بارے $\Sat{M}{!A}[\mathbf{a}]$ اُدوں تے
  صرف اُدوں جدوں $\Sat{N}{!A}[\mathbf{b}]$ (ایتھے وی $\mathbf{a}$،
  $!A$ نوں پورا کردا اے جے کوئی وی $s$، جیہدے لئی
  $s(x_i) = a_i$، $!A$ نوں پورا کرے)۔ ہور ایہ کہ جے $\Lang{L}$
  محدود ہووے، تاں الٹا دعویٰ وی درست اے۔
  ایتھے ہر فارمولے دے سارے آزاد متغیر دتی متغیر لسٹ وچ ہون؛ ایس لئی
  لسٹ نال ملدیاں قدر بندیاں باقی متغیراں اُتے کیہ قدراں دیندیاں نیں، فرق نہیں پیندا۔
\end{thm}

\begin{proof}
  $I_n(\mathbf{a},\mathbf{b})$ توں ایہ نکلن دا ثبوت کہ
  $\mathbf{a}$ تے $\mathbf{b}$، مقدار نمائی درجہ $n$ توں
  ودھ نہ رکھن والے اکّو !!{formula}s پورے کردے نیں،
  $!A$ اُتے سادے استقرا نال ملدا اے۔ الٹے دعوے لئی اسیں
  $n$ اُتے استقرا کردے آں، تے \olref{prop:qr-finite} ورتدے آں؛
  ایہ یقینی بناؤندا اے کہ ہر $n$ لئی ایس مقدار نمائی درجے دے
  غیر ہم ارز !!{formula}s دی زیادہ توں زیادہ محدود گنتی اے۔
  ایتھے اوپر ثابت کیتا مقرر آزاد متغیر لسٹ والا روپ ورتیا گیا اے۔

  $n=0$ لئی ایہ مفروضہ کہ $\mathbf{a}$ تے $\mathbf{b}$ اکّو
  بے مقدار نما !!{formula}s پورے کردے نیں، ایہ وی دیندا اے کہ اوہ
  اکّو جوہری فارمولے پورے کردے نیں؛ ایس لئی
  $I_0(\mathbf{a},\mathbf{b})$۔

  $n+1$ دے حال وچ فرض کرو $\mathbf{a}$ تے $\mathbf{b}$،
  مقدار نمائی درجہ $n+1$ توں ودھ نہ رکھن والے اکّو
  !!{formula}s پورے کردے نیں۔ $I_{n+1}(\mathbf{a},\mathbf{b})$
  دکھاؤن لئی دونوں سمتاں دکھاؤنیاں نیں۔ اگلی سمت لئی ہر $a \in \Domain M$ لئی
  کوئی $b \in
  \Domain N$ ہووے جیہدے لئی $I_n(\mathbf{a}a,\mathbf{b}b)$؛
  تے استقرائی مفروضے نال ایہ دکھانا کافی اے کہ ہر $a \in
  \Domain M$ لئی کوئی $b \in \Domain N$ ہووے جیہدے لئی
  $\mathbf{a}a$ تے $\mathbf{b}b$ مقدار نمائی درجہ
  $n$ توں ودھ نہ رکھن والے اکّو !!{formula}s پورے کردے ہون۔

  کوئی $a \in \Domain M$ لو؛ $!T^a_n$ لئی اوہناں !!{formula}s
  $!B(x,\mathbf{y})$ دیاں منطقی ہم ارزی جماعتاں وچوں اک اک نمائندہ چنو
  جنہاں دا درجہ $n$ توں ودھ نہیں تے جنہاں نوں $\mathbf{a}a$، $\Struct{M}$ وچ پورا کردا اے۔
  $!T^a_n$ ہن چنے ہوئے سچ نمائندیاں دا محدود سیٹ اے، سارے لفظی فارمولیاں دا نہیں۔
  ایس لئی اوہناں ساریاں دی عطف نال پہلے درجے دا اک
  !!{formula} بنا لو؛ عطف مقدار نمائی درجہ نہیں ودھاؤندی۔ ایس توں نکلدا اے کہ
  $\mathbf{a}$، $\lexists[x][!T^a_n(x,\mathbf{y})]$ نوں پورا
  کردا اے، جیہدا مقدار نمائی درجہ $n+1$ توں ودھ نہیں۔
  مفروضے مطابق $\mathbf{b}$، $\Struct{N}$ وچ ایہو
  !!{formula} پورا کردا اے، ایس لئی کوئی $b \in \Domain N$ اے
  جیہدے نال $\mathbf{b}b$، $!T^a_n$ نوں پورا کردا اے؛ خاص طور
  اُتے $\mathbf{b}b$، مقدار نمائی درجہ $n$ توں ودھ نہ رکھن
  والے اوہو !!{formula}s پورے کردا اے جیہڑے $\mathbf{a}a$ پورے
  کردا اے۔ دونوں رخ ایس لئی ملدے نیں کہ اصل لسٹ لئی جھوٹے فارمولے
  دی نفی وی ایس درجے وچ سچی اے؛ اوہدی ہم ارزی جماعت دا نمائندہ وی
  چنے سیٹ وچ اے، اس لئی نئی لسٹ اوہ جھوٹا فارمولا سچ نہیں کر سکدی۔
  ایس طرح ہر $b \in \Domain N$ لئی کوئی
  $a\in \Domain M$ لبھدا اے جیہدے نال $\mathbf{a}a$ تے
  $\mathbf{b}b$ مقدار نمائی درجہ $n$ توں ودھ نہ رکھن والے
  اکّو !!{formula}s پورے کردے نیں؛ ایس نال ثبوت پورا ہوندا اے۔
\end{proof}

\begin{cor}\ollabel{cor:b-n-f}
  جے $\Struct{M}$ تے $\Struct{N}$ اک محدود !!{language} دیاں
  خالص تعلقی !!{structure}s نیں، تاں $\Struct{M} \approx_n\Struct{N}$
  اُدوں تے صرف اُدوں جدوں $\Struct{M} \elemequiv[n] \Struct{N}$۔
  خاص طور اُتے $\Struct{M} \elemequiv \Struct{N}$ اُدوں تے
  صرف اُدوں جدوں ہر $n$ لئی $\Struct{M} \approx_n \Struct{N}$۔
\end{cor}
\endgroup
% /OLP-0189

% OLP-0190 content/model-theory/basics/dlo.tex
\begingroup
\def\olfilename{dlo}
\expandafter\def\csname ol@sectioncs\endcsname{section}

\olfileid{mod}{bas}{dlo}
\section{گھنیاں خطی ترتیباں}

\begin{defn}
  \emph{سِراں توں خالی گھنی خطی ترتیب}
  \pnbLatin{(dense linear ordering without endpoints)} اوہ !!a{structure}
  $\Struct{M}$ اے جیہڑی اکّو دو-تھانواں والا
  !!{predicate}~$<$ رکھن والی !!{language} دی اے تے تھلے دتے
  جملے پورے کردی اے:
  \begin{enumerate}
  \item $\lforall[x][\lnot x < x]$؛
  \item $\lforall[x][\lforall[y][\lforall[z][(x < y \lif (y < z \lif x
    <z ))]]]$؛
  \item $\lforall[x][\lforall[y][(x< y \lor \eq[x][y] \lor y < x)]]$؛
  \item $\lforall[x][\lexists[y][x < y]]$؛
  \item $\lforall[x][\lexists[y][y < x]]$؛
  \item $\lforall[x][\lforall[y][(x < y \lif \lexists[z][(x < z \land
        z < y))]]]$۔
 \end{enumerate}
\end{defn}

\begin{thm}\ollabel{thm:cantorQ}
  سِراں توں خالی ہر دو !!{enumerable} گھنیاں خطی ترتیباں
  ہم شکل ہوندیاں نیں۔
\end{thm}

\begin{proof}
  $\Struct{M_1}$ تے $\Struct{M_2}$ سِراں توں خالی
  !!{enumerable} گھنیاں خطی ترتیباں ہون، جتھے
  ${<_1} = \Assign{<}{M_1}$ تے ${<_2} =
  \Assign{<}{M_2}$؛ انہاں دے وچکار ساریاں جزوی ہم شکلیاں دا
  سیٹ $\PIso{I}$ ہووے۔ $\PIso{I}$ غیر خالی اے کیوں جے گھٹ توں گھٹ
  $\emptyset \in \PIso{I}$۔ اسیں دکھاندے آں کہ $\PIso{I}$،
  اگے-پیچھے خاصیت پوری کردا اے۔ تاں
  $\Struct{M_1} \iso[p] \Struct{M_2}$، تے
  \olref[pis]{thm:p-isom1} توں قضیہ نکلدا اے۔

  $\PIso{I}$ دی اگے والی خاصیت دکھاؤن لئی $p \in
  \PIso{I}$ لو۔ جے نقشہ خالی ہووے، تاں دوجی ساخت دے غیر خالی دائرے
  وچوں کوئی وی عنصر چن کے دتے عنصر نال اوہدا جوڑا شامل کر لو۔ جے دتا
  عنصر پہلے ای نقشے دے دائرے وچ اے، تاں نقشہ آپ ہی مطلوبہ توسیع اے۔
  اگے غیر خالی نقشے تے دائرے توں باہر دتے عنصر دا حال اے؛ لسٹ وچ
  نقشے دے دائرے دے سارے عناصر لو۔ $p(a_i) = b_i$، $i = 1$، \dots،~$n$ لئی
  ہووے؛ دائرے دی لسٹ نوں ترتیب دے کے فرض کرو $a_1 <_1 a_2 <_1 \cdots <_1
  a_n$۔ کوئی $a \in \Domain{M_1}$ دتا ہووے؛ $b \in \Domain{M_2}$
  ایویں لبھو:
  \begin{enumerate}
  \item جے $a <_1 a_1$، تاں $b \in \Domain{M_2}$ ایسا لو
    کہ $b <_2
    b_1$؛
  \item جے $a_n <_1 a$، تاں $b \in \Domain{M_2}$ ایسا لو
    کہ $b_n <_2 b$؛
 \item جے $a_i <_1 a <_1 a_{i+1}$ کسے $i$ لئی، تاں
   $b \in
   \Domain{M_2}$ ایسا لو کہ $b_i <_2 b <_2 b_{i+1}$۔
  \end{enumerate}
  مطلوبہ خاصیت والا $b$ ہمیشہ لبھ جاندا اے، کیوں جے
  $\Struct{M_2}$ سِراں توں خالی گھنی خطی ترتیب اے۔
  $q = p \cup \{ \langle a, b \rangle \}$ مقرر کرو؛ ایویں
  $q \in \PIso{I}$، $p$ دا مطلوبہ ودھاؤ اے۔ ایس نال اگے
  والی خاصیت ثابت ہوندی اے۔ پیچھے والی خاصیت دا ثبوت ایہو
  جیہا اے۔ ایس لئی $\Struct{M_1} \iso[p]
  \Struct{M_2}$؛ \olref[pis]{thm:p-isom1} نال
  $\Struct{M_1} \iso
  \Struct{M_2}$۔
\end{proof}

\begin{prob}
  \olref[mod][bas][dlo]{thm:cantorQ} دا ثبوت ایہ جانچ کے پورا
  کرو کہ $\PIso{I}$ پیچھے والی خاصیت پوری کردا اے۔
\end{prob}

\begin{rem}
  $\Struct{S}$ کوئی وی سِراں توں خالی !!{enumerable} گھنی
  خطی ترتیب ہووے۔ تاں (\olref{thm:cantorQ} مطابق)
  $\Struct{S} \iso
  \Struct{Q}$، جتھے $\Struct{Q} = (\Rat, <)$ اوہ
  !!{enumerable} گھنی خطی ترتیب اے جیہدا دائرہ ناطق اعداد دا
  سیٹ $\Rat$ اے۔ ہن \olref[thm]{remark:R} وچ آئی
  !!{structure}~$\Struct{R} =
  (\Real, <)$ نوں فیر ویکھو۔ اسیں ویکھیا سی کہ اک
  !!a{enumerable} !!{structure}~$\Struct{S}$ اے جیہدے لئی
  $\Struct{R}
  \elemequiv \Struct{S}$۔ پر $\Struct{S}$ سِراں توں خالی
  !!a{enumerable} گھنی خطی ترتیب اے، ایس لئی اوہ
  !!{structure}~$\Struct{Q}$ دے نال ہم شکل (تے ایس سبب
  جملہ وار ہم ارز) اے۔ جملہ وار ہم ارزی دی تعدی خاصیت توں
  $\Struct{R} \elemequiv
  \Struct{Q}$۔ (ایہ نتیجہ اسیں اوہو اگے-پیچھے استدلال نال
  $\Struct{R} \iso[p] \Struct{Q}$ قائم کرکے براہِ راست وی
  دکھا سکدے ساں۔)
\end{rem}
\endgroup
% /OLP-0190

%\olimport{nonstandard-arithmetic}

% OLP-0653 content/model-theory/basics/nonstandard-arithmetic.tex
\begingroup
\def\olfilename{nonstandard-arithmetic}
\expandafter\def\csname ol@sectioncs\endcsname{section}

\olfileid{mod}{bas}{nsa}
\section{حساب دے غیر معیاری نمونے}

\begin{defn}
منّ لو $\Lang{L_N}$ حساب دی !!{language} اے، جیہدے وچ اک !!{constant} $\Obj{0}$، اک دو تھاواں والا !!{predicate} $<$، اک اک تھاں والا !!{function} $\prime$، تے دو تھاواں والے !!{function}s $+$ تے $\times$ نیں۔
\begin{enumerate}
\item حساب دا \emph{معیاری نمونہ}، !!{structure} $\Struct{N}$ اے، $\Lang{L_N}$ لئی، جیہدے وچ $\Nat = \{ 0, 1, 2, \dots\}$ اے تے $\Obj{0}$ دی تعبیر $0$، $<$ دی تعبیر~$\Nat$ اُتے چھوٹا ہون دا تعلق، تے $\prime$، $+$ تے $\times$ دیاں تعبیراں بالترتیب $\Nat$ اُتے جانشین، جمع تے ضرب نیں۔
\item \emph{سچّا حساب} نظریہ $\Theory{N}$ اے۔
\end{enumerate}
\end{defn}

$\Lang{L_N}$ وچ کم کردیاں، $\Obj{0}^{\prime\dots\prime}$ دی شکل والا ہر ٹرم، جیہدے وچ جانشین فنکشن $n$ واری $\Obj{0}$ اُتے لگدا اے، مختصر~$\num{n}$ لکھدے آں۔

\begin{defn}
!!{structure}~$\Struct{M}$، $\Lang{L_N}$ لئی، اُدوں تے صرف اُدوں \emph{معیاری} اے جدوں $\Struct{N} \iso \Struct{M}$۔
\end{defn}

\begin{thm}\ollabel{thm:non-std}
سچّے حساب دے غیر معیاری !!{enumerable} نمونے موجود نیں۔
\end{thm}

\begin{proof}
$\Lang{L_N}$ نوں اک نواں !!{constant}~$c$ شامل کر کے ودھاؤ، تے ایہ نظریہ ویکھو:
\[
\Theory{N} \cup \Setabs{\num{n} < c}{n \in \Nat}.
\]
نظریہ محدود طور اُتے پورا اتر سکدا اے۔ ایس لئی محدود گواہی دے قضیے نال ایہدا اک نمونہ $\Struct{M}$ اے، جیہنوں تھلّے ول لووَن ہائیم--سکولم دے قضیے نال !!{enumerable} منّ سکدے آں۔ جتھے $\Domain{M}$، $\Struct{M}$ دا دائرہ اے، منّ لو $\Struct{M}$، $\Lang{L}$ دے غیر منطقی ثابت نشاناں دی تعبیر ایویں کردی اے:
$\mathbf{z} = \Assign{\Obj{0}}{M} \in \Domain{M}$، ${\prec} =
\Assign{<}{M} \subseteq M^2$، $* = \Assign{\prime}{M} \colon
\Domain{M} \to \Domain{M}$، تے $\oplus = \Assign{+}{M}, \otimes =
\Assign{\times}{M} \colon \Domain{M}^2 \to \Domain{M}$۔ ہر $x
\in \Domain{M}$ لئی اسیں $x^*$ اوہ رکن لکھدے آں جو $\Domain{M}$ وچ $x$ اُتے~$*$ لا کے ملدا اے۔

ہُن جے $h$، $\Struct{N}$ تے $\Struct{M}$ دی ہم ساختی ہووے، تاں کوئی $n \in \Nat$ ہووے گا جیہدے لئی $h(n) = \Assign{c}{M}$۔ سو منّ لو $s$، $\Struct{N}$ وچ کوئی وی قدر بندی اے جیہدے لئی $s(x) = n$۔ تاں $\Sat{N}{\eq[\num{n}][x]}[s]$۔ \olref[iso]{thm:isom} دے ثبوت نال $\Sat{M}{\eq[\num{n}][x]}[h \circ s]$ وی، ایس لئی $\Assign{c}{M} =
\mathbf{z}^{*\cdots *}$ ($*$ نوں $n$ واری لا کے)۔ پر ایہ ناممکن اے، کیوں جے مفروضے نال $\Sat{M}{\num{n} < c}$ اے تے $\prec$ اپنے آپ نال لاگو نہیں ہوندا۔ سو $\Struct{M}$ غیر معیاری اے۔
\end{proof}

\begin{prob}
تعلق $R$، سیٹ $X$ اُتے، اُدوں تے صرف اُدوں \emph{حُسنِ بنیاد والا} اے جدوں~$R$ وچ تھلّے ول جاندی کوئی لامتناہی لڑی نہیں، یعنی جے $x_0$، $x_1$، $x_2$،~\dots سیٹ $X$ وچ ایہو جیہے نہیں کہ $\dots x_2Rx_1Rx_0$۔ فرض کرو زرمَیلو--فرینکل سیٹ نظریہ~$ZF$ بے تضاد اے۔ وکھاؤ کہ $ZF$ دے حُسنِ بنیاد توں سکھنے نمونے نیں، یعنی ایہو جیہے نمونے $\mathfrak{M}$ جیہناں لئی $\dots x_2 \in x_1 \in x_0$۔
\end{prob}

چونکہ غیر معیاری نمونہ $\Struct{M}$، جیہڑا \olref{thm:non-std} وچ اے، معیاری نمونے دے نال ابتدائی طور اُتے ہم ارز اے، $\Struct{M}$ دیاں کئی خاصیتاں اخذ ہو سکدیاں نیں۔ ایس حصے دا باقی بیان ایس کم لئی اے۔ ایس نال سانوں !!{enumerable} غیر معیاری نمونیاں، جیہڑے $\Theory{N}$ دے نیں، دی دقیق پہچان ملے گی۔

\begin{enumerate}
\item $\Domain{M}$ دا کوئی رکن $\prec$ دے لحاظ نال اپنے آپ توں چھوٹا نہیں۔ جملہ $\lforall[x][\lnot x < x]$، $\Struct{N}$ وچ سچ اے تے ایس لئی $\Struct{M}$ وچ وی۔
\item ایہو جیہی دلیل نال ملدا اے کہ $\prec$، $\Domain{M}$ دی اک \emph{خطی ترتیب} اے، یعنی $\Domain{M}$ اُتے ساریاں وکھریاں جوڑیاں نوں ترتیب دین والا، اپنے آپ نال لاگو نہ ہون والا تے تعدی والا تعلق اے۔
\item رکن $\mathbf{z}$، $\prec$ دے لحاظ نال $\Domain{M}$ دا سبھ توں چھوٹا رکن اے۔
\item $\Domain{M}$ دا ہر رکن $\prec$ دے لحاظ نال اپنے $*$-جانشین توں چھوٹا اے، تے $x^*$، $\prec$ دے لحاظ نال $\Domain{M}$ دا سبھ توں چھوٹا رکن اے جیہڑا $x$ توں وڈا اے۔
\item $\Struct{M}$ وچ اک مُڈھلا حصہ ($\prec$ دا) اے جیہڑا $\Nat$ نال ہم ساخت اے:
$\mathbf{z}, \mathbf{z}^*, \mathbf{z}^{**}, \dots$۔ اسیں ایہنوں $\Domain{M}$ دا \emph{معیاری حصہ} آکھدے آں۔ $\Domain{M}$ دا باقی ہر رکن \emph{غیر معیاری} اے۔ $\Domain{M}$ وچ غیر معیاری رکن ضرور ہون گے، نہیں تاں \olref{thm:non-std} دے ثبوت دا فنکشن $h$ ہم ساختی ہووے گا۔ اسیں $n, m, \dots$ نوں !!{variable}s وجوں ورتدے آں جیہڑے $\Struct{M}$ دے ایس معیاری حصے اُتے پھردے نیں۔
\item ہر غیر معیاری رکن ہر معیاری رکن توں وڈا اے۔ ایہ ایس لئی اے کہ ہر $n \in \Nat$ لئی،
  \[
  \Sat{N}{\lforall[z][(\lnot(\eq[z][\Obj{0}] \lor
  \dots \lor \eq[z][\num{n}]) \lif \num{n} < z)]},
  \]
سو جے $z \in \Domain{M}$ سارے معیاری رکناں توں وکھرا اے، تاں اوہ اوہناں ساریاں توں \emph{وڈا} اے۔
\item $\Domain{M}$ دا ہر رکن، $\mathbf{z}$ توں چھُٹ، $*$-جانشین اے، $\Domain{M}$ دے کسے اکّو رکن دا، جیہنوں $^*x$ لکھدے آں۔ جے $x = y^*$ تاں $x$ تے $y$ دوویں معیاری نیں جے اوہناں وچوں اک معیاری اے (تے جے اک غیر معیاری اے تاں دوویں غیر معیاری نیں)۔
\item ہم ارزی تعلق $\approx$، $\Domain{M}$ اُتے، ایویں تعریف کرو کہ $x \approx y$ اُدوں تے صرف اُدوں جدوں کسے \emph{معیاری} $n$ لئی یا $x \oplus n = y$ یا $y \oplus n =x$۔ دوجے لفظاں وچ $x
  \approx y$ اُدوں تے صرف اُدوں جدوں $x$ تے $y$ دی دُوری محدود اے۔ جے $n$ تے $m$ معیاری نیں تاں $n \approx m$۔ $x$ دا \emph{بلاک} ہم ارزی جماعت $[x] = \Setabs{y}{x
  \approx y}$ تعریف کرو۔
\item منّ لو $x \prec y$ جتھے $x \not\approx y$۔ چونکہ
  $\Sat{N}{\lforall[x][\lforall[y][(x < y \lif (x' < y \lor x' =
      y))]]}$، یا $x^* \prec y$ یا $x^* = y$۔ پچھلی صورت ناممکن اے کیوں جے اوہدے توں $x \approx y$ نکلدا اے، سو $x \prec
  y$۔ ایویں ای جے $x \prec y$ تے $x \not\approx y$، تاں $x \prec
  {^*y}$۔ ایس لئی جے $x \prec y$ تے $x \not\approx y$، تاں ہر $w \approx x$، $\prec$ دے لحاظ نال ہر $v \approx y$ توں چھوٹا اے۔ ایس حساب نال ہر بلاک $[x]$ دوہاں پاسیاں ول لامتناہی لڑی بناؤندا اے:
  \[
  \cdots \prec  {^{**}x} \prec {^*}x \prec x \prec x^* \prec x^{**}
  \prec \cdots
  \]
جیہنوں $Z$-لڑی آکھدے نیں کیوں جے ایہدی ترتیبی قسم صحیح عدداں والی اے۔
\item $\prec$ دی ترتیب بلاکاں تک چُک سکدے آں۔ جے $x \prec y$، تاں $x$ دا بلاک $y$ دے بلاک توں چھوٹا اے۔ جے بلاک وچ کوئی غیر معیاری رکن ہووے تاں اوہ \emph{غیر معیاری} اے۔ معیاری بلاک سبھ توں چھوٹا بلاک اے۔
\item کوئی سبھ توں چھوٹا غیر معیاری بلاک نہیں۔ جے $y$ غیر معیاری اے تاں اک $x \prec y$ اے جتھے $x$ وی غیر معیاری اے تے $x
  \not\approx y$۔ ثبوت۔ معیاری نمونے $\Struct{N}$ وچ ہر عدد نوں دو نال تقسیم کیتا جا سکدا اے، ممکن اے باقی اک ہووے، یعنی
  $\Sat{N}{\lforall[y][\lforall[x][(y = x + x \lor y = x + x +
        \Obj{0}')]]}$۔ ابتدائی ہم ارزی نال ہر $y \in
  \Domain{M}$ لئی کوئی $x \in \Domain{M}$ اے ایویں کہ یا $x \oplus x
  = y$ یا $x \oplus x \oplus \mathbf{z}^*= y$۔ جے $x$ معیاری ہووے، تاں $y$ وی، ایس لئی $x$ غیر معیاری اے۔ ہور، $x$ تے $y$ وکھرے بلاکاں وچ نیں، یعنی $x \not\approx y$۔ ایہ ویکھن لئی منّ لو اوہ اِکّو بلاک وچ نیں، یعنی $x \oplus n =
  y$ کسے معیاری $n$ لئی۔ جے $y = x \oplus x$، تاں $x \oplus n = x
  \oplus x$، جیہدے توں $x = n$، جمع دے ردّ کرن والے قاعدے نال (جیہڑا $\Struct{N}$ وچ لگدا اے تے ایس لئی $\Struct{M}$ وچ وی)، تے فیر $x$ معیاری ہووے گا۔ جے $y = x \oplus x
  \oplus \mathbf{z}^*$ تاں وی ایہو جیہی دلیل اے۔
\item ایہو جیہی دلیل نال، کوئی سبھ توں وڈا بلاک نہیں۔
\item بلاکاں دی ترتیب گنجان اے۔ جے $[x]$، $[y]$ توں چھوٹا اے (جتھے $x \not\approx y$)، تاں اک بلاک $[z]$ اے جیہڑا دوہاں توں وکھرا تے اوہناں دے وچکار اے۔ منّ لو $x \prec y$۔ پہلے وانگ $x \oplus
  y$ نوں دو نال تقسیم کیتا جا سکدا اے (ممکن اے کجھ باقی ہووے)، سو کوئی $u
  \in \Domain{M}$ اے ایویں کہ یا $x \oplus y = u \oplus u$ یا $x
  \oplus y = u \oplus u \oplus \mathbf{z}^*$۔ رکن $u$، $x$ تے $y$ دی اوسط اے، سو دوہاں دے وچکار اے۔ منّ لو $x \oplus y =
  u \oplus u$ (دوجی صورت ایہو جیہی اے)۔ جے $u \approx x$، تاں کسے معیاری $n$ لئی:
  \[
  x \oplus y = x \oplus n \oplus x \oplus n,
  \]
ایس لئی $y = x \oplus n \oplus n$ تے سانوں $x \approx y$ ملدا، جیہڑا فرض دے اُلٹ اے۔ سو $u \not\approx x$۔ ایہو جیہی دلیل $u \not\approx y$ دیندی اے۔
\end{enumerate}
ایس لئی غیر معیاری بلاک ناطق عدداں وانگ ترتیب وچ نیں۔ اوہ دوہاں سِریاں توں بناں !!a{enumerable} خطی ترتیب بناؤندے نیں۔ ایس توں نکلدا اے کہ سچّے حساب دے کسے وی دو !!{enumerable} غیر معیاری نمونے $\Struct{M}_1$ تے $\Struct{M_2}$ جے بس $<$ تے $=$ والی زبان تک گھٹا دِتّے جان، تاں اوہ ہم ساخت نیں۔ ہم ساختی $h$ دی تعریف ایویں کیتی جا سکدی اے۔ $\Struct{M_1}$ تے $\Struct{M_2}$ دے معیاری حصے معیاری نمونے $\Struct{N}$ نال ہم ساخت نیں، سو اک دوجے نال وی۔ غیر معیاری حصے والے بلاک آپ ناطق عدداں وانگ ترتیب وچ نیں، سو \olref[dlo]{thm:cantorQ} نال ہم ساخت نیں۔ بلاکاں دی ہم ساختی نوں بلاکاں دے \emph{اندر} تک ودھا سکدے آں۔ ہر بلاک وچوں کوئی وی رکن چُن کے جوڑی بناؤ، فیر $x$ دے جانشین دا عکس، $\Struct{M_1}$ توں، ایویں لو کہ اوہ $x$ دے عکس دا جانشین ہووے، $\Struct{M_2}$ وچ۔ دھیان دیو، ایس توں ایہ \emph{نہیں} نکلدا کہ $\mathfrak{M}_1$ تے $\mathfrak{M}_2$ حساب دی پوری زبان وچ ہم ساخت نیں (ہم ساختی ہمیشہ دِتّے نشاناں دے مجموعے دے لحاظ نال ہوندی اے)، کیوں جے $\Domain{M_1}$ تے $\Domain{M_2}$ اُتے جمع تے ضرب تعریف کرن دے ایہو جیہے ڈھنگ نیں جو ہم ساخت نہیں۔ (ایہ واٹ دے مشہور قضیے توں وی نکلدا اے کہ کسے پورے نظریے دے قابلِ شمار نمونیاں دی گنتی~2 نہیں ہو سکدی۔)

\begin{prob}
وکھاؤ کہ حساب دے غیر معیاری نمونے وچ کوئی سبھ توں وڈا بلاک نہیں ہو سکدا۔
\end{prob}

\begin{prob}
منّ لو $\Lang{L}$ پہلے درجے دی !!{language} اے جیہدے وچ $<$ اکّو اک !!{predicate} اے ($\eq$ توں چھُٹ)، تے منّ لو $\Struct{N} = (\Nat,
<)$۔ $\Struct{N}$ دے سارے محدود یا محدود متمم والے ذیلی سیٹ تعریف کیتے جا سکدے نیں۔ وکھاؤ کہ $\Struct{N}$ دے \emph{بس ایہو} ذیلی سیٹ تعریف کیتے جا سکدے نیں۔

(اشارہ۔ پہلاں $\Obj{prc}(x,y)$، $\Lang{L}$ دا اوہ فارمولا ہووے جیہڑا «$x$، $y$ دا نال دا پیشرو اے» دا مختصر روپ اے:
\[
x<y \land \lnot \lexists[z][(x<z \land z < y)].
\]
ہُن $\Struct{N}$ دے ہر تعریف جوگ ذیلی سیٹ نال فارمولا $!A(x)$، $\Lang{L}$ وچ، جُڑدا اے۔ کسے وی ایہو جیہے $!A$ لئی جملہ $!D$ ویکھو:
\[
\lexists[x][\lforall[y][\lforall[z][((x<y \land x<z \land \Obj{prc}(y,z)
  \land !A(y)) \lif !A(z))]]].
\]
وکھاؤ کہ $\Sat{N}{!D}$ اُدوں تے صرف اُدوں جدوں $\Struct{N}$ دا $!A$ نال تعریف کیتا ذیلی سیٹ یا محدود اے یا اوہدا متمم محدود اے۔

ہُن منّ لو $\Struct{M}$ اک غیر معیاری نمونہ اے جیہڑا $\Struct{N}$ نال ابتدائی طور اُتے ہم ارز اے۔ جے $a \in \Domain{M}$ غیر معیاری اے، تاں منّ لو $b, c \in \Domain{M}$، $a$ توں وڈے نیں تے $b$،~$c$ دا نال دا پیشرو اے۔ تاں اک خود ہم ساختی $h$، $\Domain{M}$ دی، اے جیہدے لئی $h(b)=c$ (کیوں؟)۔ ایس لئی جے $b$، $!A$ پورا کردا اے، تاں $c$ وی (کیوں؟)۔ ایس توں نکلدا اے کہ $!D$، $\Struct{M}$ وچ سچ اے تے ایس لئی $\Struct{N}$ وچ وی۔ پر ایس توں نکلدا اے کہ $\Struct{N}$ دا $!A$ نال تعریف کیتا ذیلی سیٹ یا محدود اے یا اوہدا متمم محدود اے۔
\end{prob}
\endgroup
% /OLP-0653
\OLEndChapterHook
\endgroup
% /OLP-0183

% OLP-0191 content/model-theory/models-of-arithmetic/models-of-arithmetic.tex
\begingroup
\def\olfilename{models-of-arithmetic}
\expandafter\def\csname ol@sectioncs\endcsname{section}

\olchapter{mod}{mar}{حساب دے نمونے}

% OLP-0192 content/model-theory/models-of-arithmetic/introduction.tex
\begingroup
\def\olfilename{introduction}
\expandafter\def\csname ol@sectioncs\endcsname{section}

\olfileid{mod}{mar}{int}
\section{تعارف}

حساب دا \emph{معیاری نمونہ} اوہ
!!{structure}~$\Struct{N}$ اے جیہدے وچ $\Domain{N} = \Nat$
تے $\Obj{0}$، $\prime$، $+$، $\times$ تے $<$ دیاں تعبیراں
اُویں نیں جیویں توقع ہوندی اے۔ یعنی $\Obj{0}$ دی تعبیر $0$ اے،
$\prime$ جانشین فنکشن اے، $+$ قدرتی اعداد دی جمع تے $\times$
اوہناں دی ضرب اے، جتھے اعداد~$\Nat$ وچوں نیں۔ خاص طور اُتے،
\begin{align*}
  \Assign{\Obj{0}}{N} & = 0\\
  \Assign{\prime}{N}(n) & = n + 1\\
  \Assign{+}{N}(n, m) & = n + m\\
  \Assign{\times}{N}(n, m) & = nm
\end{align*}
بے شک $\Lang{L_A}$ دیاں ایسی ساختاں وی نیں جنہاں دے
دائرے~$\Nat$ توں وکھرے نیں۔ مثال لئی، اسیں
$\Struct{M}$ لے سکدے آں جیہدا دائرہ
$\Domain{M} = \{a\}^*$ اے۔
ایہ صرف اک علامت~$a$ توں بنیاں محدود ترتیب وار لڑیواں
دا سیٹ اے، یعنی $\emptyset$، $a$، $aa$، $aaa$، \dots؛
تعبیراں ایہ نیں:
\begin{align*}
  \Assign{\Obj{0}}{M} & = \emptyset\\
  \Assign{\prime}{M}(s) & = s \concat a\\
  \Assign{+}{M}(a^n, a^m) & = a^{n + m}\\
  \Assign{\times}{M}(a^n, a^m) & = a^{nm}
\end{align*}
چھوٹے ہون دا تعلق لڑیواں دیاں لمبائیاں دی قدرتی ترتیب دے مطابق اے۔
ایہ دو ساختاں اک خاص معنی وچ ”بنیادی طور اُتے اکّو جیہیاں“
نیں: فرق صرف اوہناں دے !!{domain}s دے !!{element}s وچ اے،
ایس گل وچ نہیں کہ تعبیر کرن والے فنکشن اوہناں !!{domain}s دے
!!{element}s نوں آپس وچ کِویں جوڑدے نیں۔ اسیں آکھدے آں کہ
ایہ دو !!{structure}s \emph{ہم شکل} نیں۔

محدود گواہی دے قضیے توں آسانی نال نکلدا اے کہ
$\Struct{N}$ وچ سچ کوئی وی نظریہ ایسے نمونے وی رکھدا اے
جو~$\Struct{N}$ نال ہم شکل نہیں۔ ایسی ساختاں نوں
\emph{غیر معیاری} آکھدے نیں۔ اوہناں دی دلچسپ گل ایہ اے کہ
معیاری نمونے دے !!{element}s، یعنی $\Struct{N}$ دے
یا اوہدے نال ہم شکل ساریاں !!{structure}s دے عنصر،
معیاری عددی عبارتاں~$\num{n}$ دیاں قدراں نال پورے ہو جاندے نیں؛
یعنی
\[
\Domain{N} = \Setabs{\Value{\num{n}}{N}}{n \in \Nat}
\]
پر رابنسن حساب (Q) دے غیر معیاری نمونے وچ ایہ گل درست نہیں:
جے $\Struct{M}$ غیر معیاری اے تے رابنسن حساب دے مسلّمے پورے کردا اے،
تاں گھٹ توں گھٹ اک
$x \in \Domain{M}$ ایسا اے کہ
$x \neq \Value{\num{n}}{M}$ ہر~$n$ لئی۔

ایہ غیر معیاری عنصر بڑے دلچسپ نیں: ایہ ”لامتناہی
قدرتی اعداد“ نیں۔ پر اک معنی وچ اوہناں دا وجود
نامکملیت دے مظاہر دی وجہ وی سمجھاؤندا اے۔ مثال لئی
پیانو حساب دی بے تضادی دا بیان~$\OCon[\Th{PA}]$
ویکھو، یعنی $\lnot
\lexists[x][\OPrf[\Th{PA}](x, \gn{\lfalse})]$۔
کیوں جے $\Th{PA}$ نہ $\OCon[\Th{PA}]$ نوں ثابت کردا اے
نہ $\lnot \OCon[\Th{PA}]$ نوں، ایس لئی دوواں وچوں
جیہڑا مرضی بیان $\Th{PA}$ وچ بے تضادی برقرار رکھ کے
شامل کیتا جا سکدا اے۔ کیوں جے $\Th{PA}$ بے تضاد اے،
$\Sat{N}{\OCon[\Th{PA}]}$، تے نتیجتاً
$\Sat/{N}{\lnot
  \OCon[\Th{PA}]}$۔ ایس لئی $\Struct{N}$،
$\Th{PA} \cup \{\lnot \OCon[\Th{PA}]\}$ دا
\emph{نمونہ نہیں}؛ ایس نظریے دے سارے نمونے لازماً
غیر معیاری نیں۔ $\Th{PA} \cup \{\lnot \OCon[\Th{PA}]\}$
دے نمونیاں وچ کوئی ایسا !!{element} ہونا چاہیدا اے
جو $\lexists[x][\OPrf[\Th{PA}](x, \gn{\lfalse})]$
نوں سچ بناؤن والا گواہ ہووے، یعنی نمونے دے اندر~$\Th{PA}$ توں
تضاد دی !!a{derivation} دا گوڈل عدد۔
ایہ !!a{element} معیاری نہیں ہو سکدا، کیوں جے
$\Th{PA} \Proves \lnot
\OPrf[\Th{PA}](\num{n}, \gn{\lfalse})$
ہر~$n$ لئی ہوندا اے۔
\endgroup
% /OLP-0192

% OLP-0193 content/model-theory/models-of-arithmetic/standard-models.tex
\begingroup
\def\olfilename{standard-models}
\expandafter\def\csname ol@sectioncs\endcsname{section}

\olfileid{mod}{mar}{stm}
\olsection{حساب دے معیاری نمونے}

حساب دی زبان~$\Lang{L_A}$ دا مقصد ظاہر اے کہ اعداد،
خاص طور اُتے قدرتی اعداد، بارے گل کرنا اے۔ ایس لئی ”اوہ“
معیاری نمونہ~$\Struct{N}$ خاص اے: ایہ اوہ نمونہ اے
جیہدے بارے اسیں گل کرنا چاہندے آں۔ پر منطق وچ ساڈی
دلچسپی اکثر ساختی خاصیتاں وچ ہوندی اے، تے ہر دو
ہم شکل !!{structure}s ایہ خاصیتاں سانجھیاں رکھدیاں نیں۔
ایس لئی اسیں تعریف ذرا کھلی رکھ سکدے آں تے ہر اوہ
!!{structure} ”معیاری“ من سکدے آں جیہڑی~$\Struct{N}$
نال ہم شکل ہووے۔

\begin{defn}
$\Lang{L_A}$ لئی !!{structure} \emph{معیاری} اے
جے اوہ~$\Struct{N}$ نال ہم شکل ہووے۔
\end{defn}

\begin{prop}
\ollabel{prop:standard-domain} جے !!a{structure}~$\Struct{M}$
معیاری اے، تاں اوہدا دائرہ معیاری عددی عبارتاں دیاں
قدراں دا سیٹ اے، یعنی
\[
\Domain{M} = \Setabs{\Value{\num{n}}{M}}{n \in \Nat}
\]
\end{prop}

\begin{proof}
واضح اے کہ ہر $\Value{\num{n}}{M} \in \Domain{M}$۔
بس ایہ وکھانا اے کہ ہر $x \in \Domain{M}$
$\Value{\num{n}}{M}$ دے برابر اے، کسے~$n$ لئی۔
کیوں جے $\Struct{M}$ معیاری اے، ایہ~$\Struct{N}$
نال ہم شکل اے۔ فرض کرو $g\colon \Nat \to \Domain{M}$
اک ہم شکلی اے۔ فیر $g(n) = g(\Value{\num{n}}{N}) =
\Value{\num{n}}{M}$۔ پر ہر $x \in \Domain{M}$
لئی کوئی~$n \in \Nat$ اے جیہدے لئی $g(n) = x$،
کیوں جے $g$ ہر ہدف قدر لیندا اے (!!{surjective})۔
\end{proof}

\begin{explain}
جے ساخت~$\Struct{M}$، $\Lang{L_A}$ لئی، معیاری اے،
تاں اوہدے !!{domain} دے سارے عنصراں نوں
معیاری عددی عبارتاں $\num{0}$، $\num{1}$،
$\num{2}$، \dots نال ناں دتا جا سکدا اے؛
یعنی ٹرماں $\Obj{0}$، $\Obj{0}'$،
$\Obj{0}''$، وغیرہ نال۔ بے شک ایس دا ایہ مطلب نہیں کہ
$\Domain{M}$ دے !!{element}s خود اعداد \emph{نیں}؛
صرف ایہ کہ اسیں اوہناں نوں اوہو طرح چن سکدے آں
جیویں~$\Domain{N}$ وچ اعداد نوں چن سکدے آں۔
\end{explain}

\begin{prob}
وکھاؤ کہ \olref[mod][mar][stm]{prop:standard-domain}
دا الٹا دعویٰ غلط اے؛ یعنی ایسی !!a{structure}~$\Struct{M}$
دی مثال دیو جیہدے لئی
$\Domain{M} = \Setabs{\Value{\num{n}}{M}}{n \in \Nat}$
ہووے، پر اوہ~$\Struct{N}$ نال ہم شکل نہ ہووے۔
\end{prob}

\begin{prop}
\ollabel{prop:thq-standard}
جے $\Sat{M}{\Th{Q}}$ ہووے تے
$\Domain{M} = \Setabs{\Value{\num{n}}{M}}{n
  \in \Nat}$، تاں $\Struct{M}$ معیاری اے۔
\end{prop}

\begin{proof}
اسیں وکھانا اے کہ $\Struct{M}$،
~$\Struct{N}$ نال ہم شکل اے۔ فنکشن
$g\colon \Nat \to \Domain{M}$ ویکھو جیہدی تعریف
$g(n) = \Value{\num{n}}{M}$ اے۔ مفروضے مطابق
$g$ ہر ہدف قدر لیندا اے (!!{surjective})۔ ایہ
اک بہ اک (!!{injective}) وی اے:
$\Th{Q} \Proves
\eq/[\num{n}][\num{m}]$ جدوں $n \neq m$
ہووے۔ ایس لئی، کیوں جے
$\Sat{M}{\Th{Q}}$، تاں
$\Sat{M}{\eq/[\num{n}][\num{m}]}$،
جدوں $n \neq m$۔ ایس لئی جے $n \neq m$، تاں
$\Value{\num{n}}{M} \neq
\Value{\num{m}}{M}$، یعنی $g(n) \neq g(m)$۔

ایہ وی ثابت کرنا اے کہ $g$ ہم شکلی اے۔
\begin{enumerate}
\item ساڈے کول $g(\Assign{\Obj{0}}{N}) = g(0)$ اے،
  کیوں جے $\Assign{\Obj{0}}{N} = 0$۔ $g$ دی تعریف توں
  $g(0) = \Value{\num{0}}{M}$۔ پر $\num{0}$
  بس $\Obj{0}$ اے، تے جیہڑا ٹرم !!a{constant} ہووے
  اوہدی قدر اوہ اے جیہڑی !!{structure} اوس
  !!{constant} نوں دندی اے، یعنی
  $\Value{\Obj{0}}{M} = \Assign{\Obj{0}}{M}$۔
  ایس لئی مطلوبہ برابری
  $g(\Assign{\Obj{0}}{N}) = \Assign{\Obj{0}}{M}$ مل گئی۔
\item $g(\Assign{\prime}{N}(n)) = g(n+1)$،
  کیوں جے $\prime$، $\Struct{N}$ وچ
  ~$\Nat$ اُتے جانشین فنکشن اے۔ فیر
  $g(n+1) = \Value{\num{n+1}}{M}$،
  $g$ دی تعریف توں۔
  پر $\num{n+1}$ اوہو ٹرم اے جو $\num{n}'$،
  ایس لئی $\Value{\num{n+1}}{M} =
  \Value{\num{n}'}{M}$۔ قدر والے فنکشن دی تعریف توں
  ایہ $= \Assign{\prime}{M}(\Value{\num{n}}{M})$ اے۔
  کیوں جے $\Value{\num{n}}{M} = g(n)$،
  اسیں $g(\Assign{\prime}{N}(n)) =
  \Assign{\prime}{M}(g(n))$ حاصل کردے آں۔
\item $g(\Assign{+}{N}(n,m)) = g(n+m)$،
  کیوں جے $+$، $\Struct{N}$ وچ
  ~$\Nat$ اُتے جمع دا فنکشن اے۔ فیر
  $g(n+m) = \Value{\num{n+m}}{M}$،
  $g$ دی تعریف توں۔
  پر $\Th{Q} \Proves
  \num{n+m} = (\num{n} + \num{m})$،
  ایس لئی $\Value{\num{n+m}}{M} =
  \Value{\num{n}+\num{m}}{M}$۔ قدر والے فنکشن دی
  تعریف توں ایہ $=
  \Assign{+}{M}(\Value{\num{n}}{M},\Value{\num{m}}{M})$
  اے۔ کیوں جے $\Value{\num{n}}{M} = g(n)$ تے
  $\Value{\num{m}}{M} = g(m)$، اسیں
  $g(\Assign{+}{N}(n, m)) = \Assign{+}{M}(g(n), g(m))$
  حاصل کردے آں۔
\item $g(\Assign{\times}{N}(n, m)) =
  \Assign{\times}{M}(g(n), g(m))$: مشق۔
\item $\tuple{n,m} \in \Assign{<}{N}$ اُدوں تے صرف اُدوں
  جدوں $n < m$۔ جے $n < m$، تاں
  $\Th{Q} \Proves \num{n} < \num{m}$،
  تے $\Sat{M}{\num{n} <
  \num{m}}$ وی۔ ایس لئی
  $\tuple{\Value{\num{n}}{M}, \Value{\num{m}}{M}} \in
  \Assign{<}{M}$، یعنی
  $\tuple{g(n), g(m)} \in \Assign{<}{M}$۔
  جے $n \not< m$، تاں
  $\Th{Q} \Proves \lnot \num{n} < \num{m}$،
  تے نتیجتاً $\Sat/{M}{\num{n} < \num{m}}$۔
  ایس لئی، پہلے وانگ،
  $\tuple{g(n), g(m)} \notin \Assign{<}{M}$۔
  دوہاں حالاں نوں ملا کے
  $\tuple{n,m} \in \Assign{<}{N}$
  اُدوں تے صرف اُدوں جدوں
  $\tuple{g(n), g(m)} \in
  \Assign{<}{M}$۔
\end{enumerate}
\end{proof}

\begin{explain}
فنکشن~$g$، $\Nat$ توں کسے وی ہور
!!{structure}~$\Struct{M}$ دے دائرے ول، $\Lang{L_A}$
لئی، نقشہ متعین کرن
دا سب توں سیدھا طریقہ اے، کیوں جے ہر ایسی
$\Struct{M}$ وچ $\num{0}$، $\num{1}$،
$\num{2}$ وغیرہ نال ناں ملن والے !!{element}s
ہوندے نیں۔ ایس لئی حیرانی نہیں کہ جے $\Struct{M}$،
$\num{n}$ دے بارے گھٹ توں گھٹ کجھ بنیادی بیان
~$\Struct{N}$ وانگ سچے بناؤندی اے، تے $g$ اک بہ اک
مطابقت وی اے، تاں $g$ ہم شکلی بن جاندا اے۔
ایہناں مفروضیاں تحت، جے $\Domain{M}$ وچ $\num{n}$
دے ناں والے عنصراں توں سوا کوئی !!{element}s نہیں،
تاں ایہی واحد ہم شکلی اے۔
\end{explain}

\begin{prop}
\ollabel{prop:thq-unique-iso} جے $\Struct{M}$ معیاری اے،
تاں \olref{prop:thq-standard} دے ثبوت والا $g$،
$\Struct{N}$ توں~$\Struct{M}$ ول واحد ہم شکلی اے۔
\end{prop}

\begin{proof}
فرض کرو $h\colon \Nat \to \Domain{M}$،
$\Struct{N}$ تے~$\Struct{M}$ وچ ہم شکلی اے۔
اسیں $g = h$ نوں~$n$ اُتے استقرا نال وکھاؤندے آں۔
جے $n = 0$، تاں
$g(0) = \Assign{\Obj{0}}{M}$، $g$ دی تعریف توں۔
پر کیوں جے $h$
ہم شکلی اے، $h(0) =
h(\Assign{\Obj{0}}{N}) =\Assign{\Obj{0}}{M}$،
ایس لئی $g(0) = h(0)$۔

ہُن $n+1$ دی حالت ویکھو۔ ساڈے کول
\begin{align*}
  g(n+1) & = \Value{\num{n+1}}{M} \text{ $g$ دی تعریف توں}\\
  & = \Value{\num{n}'}{M} \text{ کیوں جے $\num{n+1}\ident \num{n}'$}\\
  & = \Assign{\prime}{M}(\Value{\num{n}}{M})
    \text{ $\Value{t'}{M}$ دی تعریف توں}\\
  & = \Assign{\prime}{M}(g(n))  \text{ $g$ دی تعریف توں}\\
  & = \Assign{\prime}{M}(h(n)) \text{ استقرائی مفروضے توں}\\
  & = h(\Assign{\prime}{N}(n)) \text{ کیوں جے $h$ ہم شکلی اے}\\
  & = h(n+1)
\end{align*}
\end{proof}

\begin{explain}
ہر !!{denumerable} سیٹ~$M$ لئی $\Nat$ تے $M$
وچکار !!a{bijection} ہوندی اے، ایس لئی ہر ایسا سیٹ~$M$
کسی معیاری نمونے~$\Struct{M}$ دا !!{domain}
بن سکدا اے۔ بلکہ جدوں اک شے $z \in
M$ تے اک مناسب فنکشن $s$، بالترتیب
$\Assign{\Obj{0}}{M}$ تے
$\Assign{\prime}{M}$ دے طور اُتے چن لو،
تاں $+$، $\times$ تے $<$ دیاں تعبیراں
پہلے ای مقرر ہو جاندیاں نیں۔ صرف اوہ فنکشن
$s\colon M \to M \setminus \{z\}$،
معیاری نمونے وچ~$\Assign{\prime}{M}$ لئی مناسب نیں
جو اک بہ اک (!!{injective}) تے ہر ہدف قدر لین والے
(!!{surjective}) دوویں ہون۔ $s$ دے حاصل قدراں دے سیٹ وچ~$z$
نہیں آ سکدا؛ نہیں تاں
$\lforall[x][\eq/[\Obj 0][x']]$ غلط ہوندا۔
ایہ !!{sentence}، $\Struct{N}$ وچ سچ اے،
تے ایس لئی $\Struct{M}$ وچ وی سچ ہونا چاہیدا اے۔
فنکشن~$s$ دا اک بہ اک (!!{injective}) ہونا لازم اے،
کیوں جے جانشین فنکشن~$\Assign{\prime}{N}$،
$\Struct{N}$ وچ ایداں دا اے، تے
~$\Assign{\prime}{N}$ دا اک بہ اک (!!{injective}) ہونا
!!a{sentence} نال بیان کیتا جاندا اے جیہڑا
~$\Struct{N}$ وچ سچ اے۔ اوہدا ہر ہدف قدر لین والا
(!!{surjective}) ہونا وی لازم اے، نہیں تاں کوئی
$x \in M \setminus \{z\}$،
~$s$ دے حاصل قدراں دے سیٹ وچ نہ ہوندا؛ یعنی
!!{sentence} $\lforall[x][(\eq[x][\Obj 0] \lor
\lexists[y][\eq[y'][x]])]$،
~$\Struct{M}$ وچ غلط ہوندا، حالانکہ
~$\Struct{N}$ وچ سچ اے۔
\end{explain}
\endgroup
% /OLP-0193

% OLP-0194 content/model-theory/models-of-arithmetic/non-standard-models.tex
\begingroup
\def\olfilename{non-standard-models}
\expandafter\def\csname ol@sectioncs\endcsname{section}

\olfileid{mod}{mar}{nst}
\section{غیر معیاری نمونے}

\begin{explain}
اسیں $\Lang{L_A}$ لئی !!a{structure} نوں معیاری
آکھدے آں جے اوہ~$\Struct{N}$ نال ہم شکل ہووے۔
جے کوئی !!a{structure}، $\Struct{N}$ نال ہم شکل
نہ ہووے، تاں اوہنوں غیر معیاری آکھیا جاندا اے۔
\end{explain}

\begin{defn}
!!{structure}~$\Struct{M}$، $\Lang{L_A}$ لئی،
\emph{غیر معیاری} اے جے اوہ~$\Struct{N}$ نال ہم
شکل نہیں۔ اوہ !!{element}s $x \in
\Domain{M}$ جیہڑے $\Value{\num{n}}{M}$
دے برابر ہون کسے $n \in \Nat$ لئی،
~$\Struct{M}$ دے \emph{معیاری اعداد} کہلاندے نیں،
تے باقی \emph{غیر معیاری اعداد}۔
\end{defn}

\begin{explain}
\olref[stm]{prop:standard-domain} مطابق، ہر معیاری
!!{structure}، $\Lang{L_A}$ لئی، صرف معیاری
!!{element}s رکھدی اے۔ ایس لئی غیر معیاری
!!{structure}، جیہڑی رابنسن حساب (Q) دے مسلّمے پورے کردی ہووے،
وچ گھٹ توں گھٹ اک غیر معیاری
عنصر ہونا لازم اے۔ بلکہ غیر معیاری !!{element}
دا وجود خود ضمانت دیندا اے کہ !!{structure}
غیر معیاری اے۔
\end{explain}

\begin{prop}
جے !!a{structure}~$\Struct{M}$، $\Lang{L_A}$ لئی،
وچ غیر معیاری عدد موجود اے، تاں $\Struct{M}$
غیر معیاری اے۔
\end{prop}

\begin{proof}
الٹ فرض کرو: $\Struct{M}$ معیاری اے، پر
اوہدے وچ غیر معیاری عدد~$x$ اے۔
$g\colon \Nat \to \Domain{M}$ اک ہم شکلی ہووے۔
~$n$ اُتے استقرا نال آسانی نال ویکھیا جا سکدا اے کہ
$g(\Value{\num{n}}{N}) = \Value{\num{n}}{M}$۔
دوجے لفظاں وچ، $g$، $\Struct{N}$ دے معیاری
اعداد نوں~$\Struct{M}$ دے معیاری اعداد ول بھیجدا اے۔
جے $\Struct{M}$ وچ غیر معیاری عدد ہووے، تاں $g$
ہر ہدف قدر نہیں لے سکدا (!!{surjective})؛ ایہ مفروضے
دے خلاف اے۔
\end{proof}

\begin{prob}
یاد کرو کہ $\Th{Q}$ وچ ایہ مسلمات نیں:
\begin{align*}
& \lforall[x][\lforall[y][(\eq[x'][y'] \lif \eq[x][y])]] \tag{$!Q_1$}\\
& \lforall[x][\eq/[\Obj 0][x']] \tag{$!Q_2$}\\
& \lforall[x][(\eq[x][\Obj 0] \lor \lexists[y][\eq[x][y']])] \tag{$!Q_3$}
\end{align*}
ایس طرح دیاں !!{structure}s~$\Struct{M_1}$،
$\Struct{M_2}$، $\Struct{M_3}$ دیو کہ:
\begin{enumerate}
\item $\Sat{M_1}{!Q_1}$، $\Sat{M_1}{!Q_2}$، $\Sat/{M_1}{!Q_3}$؛
\item $\Sat{M_2}{!Q_1}$، $\Sat/{M_2}{!Q_2}$، $\Sat{M_2}{!Q_3}$؛ تے
\item $\Sat/{M_3}{!Q_1}$، $\Sat{M_3}{!Q_2}$، $\Sat{M_3}{!Q_3}$؛
\end{enumerate}
صاف اے کہ ہر اک لئی بس~$\Assign{\Obj{0}}{M_i}$
تے $\Assign{\prime}{M_i}$ متعین کرنے نیں۔
\end{prob}

\begin{explain}
$\Lang{L_A}$ لئی غیر معیاری !!{structure}s
متعین کرنا کافی آسان اے۔ مثال لئی
!!{domain}~$\Int$ والی ساخت لے لو تے ساریاں
غیر منطقی علامتاں دی معمول مطابق تعبیر کرو۔
کیوں جے منفی اعداد $\num{n}$ دیاں قدراں
نہیں، کسے وی~$n$ لئی، ایہ ساخت غیر معیاری اے۔
بے شک ایہ حساب دا \emph{نمونہ} اوس معنی وچ نہیں
ہووے گی جیہدے وچ اوہ اوہو جملے سچے کرے
جو~$\Struct{N}$ سچے کردا اے۔ مثال لئی،
$\lforall[x][\eq/[x'][\Obj{0}]]$ غلط اے۔
پر محدود گواہی دا قضیہ ورت کے اسیں آسانی نال
ثابت کر سکدے آں کہ حساب دے غیر معیاری نمونے موجود نیں۔
\end{explain}

\begin{prop}
فرض کرو $\Th{TA} = \Setabs{!A}{\Sat{N}{!A}}$،
~$\Struct{N}$ دا نظریہ اے۔ $\Th{TA}$ دا اک
!!a{enumerable} غیر معیاری نمونہ موجود اے۔
\end{prop}

\begin{proof}
$\Lang{L_A}$ وچ نواں !!{constant}~$c$ شامل کرو
تے !!{sentence}s دے ایس سیٹ نوں ویکھو:
\[
\Gamma = \Th{TA} \cup \{\eq/[c][\num{0}], \eq/[c][\num{1}],
\eq/[c][\num{2}], \dots\}
\]
ہر نمونہ~$\Struct{M^c}$، $\Gamma$ دا، کوئی
!!a{element}~$x =
\Assign{c}{M^c}$ رکھے گا جو غیر معیاری اے،
کیوں جے $x \neq \Value{\num{n}}{M}$
ہر $n \in \Nat$ لئی۔
ایہ وی صاف اے کہ $\Sat{M^c}{\Th{TA}}$،
کیوں جے $\Th{TA} \subseteq \Gamma$۔
جے اسیں $\Struct{M^c}$ توں
!!a{structure}~$\Struct{M}$، $\Lang{L_A}$ لئی،
صرف~$c$ نوں بھلا کے بنا لئیے،
تاں اوہدے دائرے وچ غیر معیاری~$x$
ہن وی موجود اے، تے~$\Sat{M}{\Th{TA}}$ وی۔
پچھلی گل ایس لئی یقینی اے کہ $c$،
$\Th{TA}$ وچ نہیں آؤندا۔ ودھائی ہوئی زبان گنتی وار اے؛
لووینہائم--سکولم دے نزولی قضیے مطابق، جے $\Gamma$ دا نمونہ ہووے،
تاں اوہدا قابلِ شمار نمونہ وی ہووے گا۔ ایس لئی صرف ایس سیٹ دا
نمونہ ہون دا ثبوت کافی اے۔

اسیں محدود گواہی دے قضیے نال وکھاؤندے آں
کہ~$\Gamma$ دا نمونہ اے۔ جے~$\Gamma$ دا
ہر محدود ذیلی سیٹ پورا ہون جوگا اے، تاں
~$\Gamma$ وی پورا ہون جوگا اے۔
من مانا محدود ذیلی سیٹ $\Gamma_0 \subseteq
\Gamma$ ویکھو۔ $\Gamma_0$ وچ~$\Th{TA}$ دے کجھ
!!{sentence}s تے~$\eq/[c][\num{n}]$
دی صورت دے کجھ جملے نیں، پر صرف محدود تعداد وچ۔
فرض کرو $k$ صفر اے جے ایہ ذیلی سیٹ عددی عبارتاں نال برابر نہ
ہون والی کوئی شرط شامل نہ کرے؛ ہور حال وچ اوہ سب توں وڈا عدد اے جیہدے لئی
$\eq/[c][\num{k}] \in \Gamma_0$۔
$\Struct{N_k}$ نوں~$\Struct{N}$ وچ
تعبیر~$\Assign{c}{N_k} = k+1$
شامل کرکے متعین کرو۔
$\Sat{N_k}{\Gamma_0}$: جے $!A
\in \Th{TA}$، تاں $\Sat{N_k}{!A}$،
کیوں جے $\Struct{N_k}$، $\Struct{N}$ ورگی اے
ہر لحاظ نال سوائے~$c$ دے،
تے $c$، $!A$ وچ نہیں آؤندا۔
تے $\Sat{N_k}{\eq/[c][\num{n}]}$،
کیوں جے $n \le k$، تے
$\Value{c}{N_k} = k+1$۔
ایس لئی~$\Gamma$ دا ہر محدود ذیلی سیٹ
پورا ہون جوگا اے۔
\end{proof}
\endgroup
% /OLP-0194

% OLP-0195 content/model-theory/models-of-arithmetic/models-of-q.tex
\begingroup
\def\olfilename{models-of-q}
\expandafter\def\csname ol@sectioncs\endcsname{section}

\olfileid{mod}{mar}{mdq}
\section{\texorpdfstring{$\Th{Q}$ دے نمونے}{Q دے نمونے}}

\begin{explain}
اسیں جاندے آں کہ غیر معیاری !!{structure}s وی نیں جنہاں وچ اوہی
!!{sentence}s سچ نیں جو~$\Struct{N}$ وچ سچ نیں؛ یعنی اوہ
~$\Th{TA}$ دے نمونے نیں۔ کیوں جے $\Sat{N}{\Th{Q}}$، ایس لئی
~$\Th{TA}$ دا ہر نمونہ~$\Th{Q}$ دا وی نمونہ اے۔ $\Th{Q}$،
~$\Th{TA}$ توں بہت کمزور اے؛ مثال دے طور اُتے
$\Th{Q} \Proves/ \lforall[x][\lforall[y][\eq[(x +
      y)][(y+x)]]]$۔ کمزور نظرییاں نوں پورا کرنا سوکھا ہوندا اے:
اوہناں دے نمونے ودھیرے ہوندے نیں۔ مثال دے طور اُتے $\Th{Q}$ دے
ایہو جئے نمونے نیں جنہاں وچ
$\lforall[x][\lforall[y][\eq[(x + y)][(y+x)]]]$ غلط اے، پر
اوہ~$\Th{TA}$ دے نمونے نہیں ہو سکدے؛ ایس معاملے وچ
~$\Th{PA}$ دے وی نہیں۔ $\Th{Q}$ دے نمونے نسبتاً سادے وی نیں:
اسیں اوہناں نوں صاف صاف متعین کر سکدے آں۔
\end{explain}

\begin{ex}
\ollabel{ex:model-K-of-Q}
اوہ !!{structure}~$\Struct{K}$ لو جہدا دائرہ $\Domain{K} = \Nat
\cup \{a\}$ اے تے علامتاں دیاں تعبیراں ایہ نیں:
\begin{align*}
  \Assign{\Obj{0}}{K} & = 0\\
  \Assign{\prime}{K}(x) & =
  \begin{cases}
    x+1 & \text{جے $x\in \Nat$}\\
    a & \text{جے $x = a$}
  \end{cases}\\
  \Assign{+}{K}(x, y) & =
  \begin{cases}
    x+y & \text{جے $x$, $y \in\Nat$}\\
    a & \text{باقی حالتاں وچ}
  \end{cases}\\
  \Assign{\times}{K}(x, y) & =
  \begin{cases}
    xy & \text{جے $x$, $y \in\Nat$}\\
    0 & \text{جے $x = 0$ یا $y = 0$}\\
    a & \text{باقی حالتاں وچ}\\
  \end{cases}\\
  \Assign{<}{K} & =
  \Setabs{\tuple{x,y}}{x, y \in \Nat \text{ تے } x<y} \cup
  \Setabs{\tuple{x,a}}{x \in \Domain{K}}
\end{align*}
ایہ دکھاؤن لئی کہ $\Sat{K}{\Th{Q}}$، سانوں تصدیق کرنی پئے گی کہ
~$\Th{Q}$ دے سارے مسلمات~$\Struct{K}$ وچ سچ نیں۔ سہولت لئی
$x^\nssucc$ نوں $\Assign{\prime}{K}(x)$ لئی لکھیے
(ایہ $x$ دا~$\Struct{K}$ وچ جانشین اے)،
$x \nsplus y$ نوں $\Assign{+}{K}(x, y)$ لئی
(ایہ $x$ تے $y$ دا~$\Struct{K}$ وچ جوڑ اے)،
$x \nstimes y$ نوں $\Assign{\times}{K}(x, y)$ لئی
(ایہ $x$ تے~$y$ دا~$\Struct{K}$ وچ ضرب دا حاصل اے)، تے
$x \nsless y$ نوں $\tuple{x,y} \in \Assign{<}{K}$ لئی لکھیے۔
ایہناں مختصر علامتاں نال~$\Struct{K}$ دیاں تعبیراں نوں ہور
صاف جدولاں وچ ایویں لکھ سکدے آں:
\[
\begin{array}{c|c}
  x & x^\nssucc \\
  \hline
  n & n+1 \\
  a & a
\end{array}
\qquad
\begin{array}{c|ccc}
  x \nsplus y & 0 & m & a \\
  \hline
  0 & 0 & m & a \\
  n & n & n+m & a \\
  a & a & a & a \\
\end{array}
\qquad
\begin{array}{c|ccc}
  x \nstimes y & 0 & m>0 & a \\
  \hline
  0 & 0 & 0 & 0 \\
  n>0 & 0 & nm & a \\
  a & 0 & a & a \\
\end{array}
\]
$n \nsless m$ اوہدوں تے صرف اوہدوں جدوں $n<m$، جدوں
$n$، $m \in \Nat$ ہون؛ تے $x \nsless a$ ہر
~$x \in \Domain{K}$ لئی اے۔

$\Sat{K}{\lforall[x][\lforall[y][(\eq[x'][y'] \lif \eq[x][y])]]}$
کیوں جے $\nssucc$ اک بہ اک (!!{injective}) اے۔
$\Sat{K}{\lforall[x][\eq/[\Obj 0][x']]}$ کیوں جے
$0$ کسے دا $\nssucc$-جانشین~$\Struct{K}$ وچ نہیں۔
$\Sat{K}{\lforall[x][(\eq[x][\Obj 0] \lor \lexists[y][\eq[x][y']])]}$
کیوں جے ہر $n>0$ لئی $n = (n-1)^\nssucc$، تے
$a = a^\nssucc$۔

$\Sat{K}{\lforall[x][\eq[(x + \Obj 0)][x]]}$ کیوں جے
$n \nsplus 0 = n+0 = n$، تے $a\nsplus 0 = a$،
~$\nsplus$ دی تعریف مطابق۔ ایہ ثابت کرنا کہ
$\Sat{K}{\lforall[x][\lforall[y][\eq[(x + y')][(x + y)']]]}$
ذرا اوکھا اے۔ جے $n$ تے $m$ دونے معیاری ہون تاں:
\begin{align*}
(n \nsplus m^\nssucc) & = (n+(m+1)) = (n+m)+1 = (n \nsplus m)^\nssucc
\intertext{کیوں جے معیاری اعداد اُتے $\nsplus$ تے $^\nssucc$،
  $+$ تے $\prime$ دے مطابق نیں۔ ہن من لو $x \in \Domain{K}$۔ تاں}
(x \nsplus a^\nssucc) & = (x \nsplus a) = a = a^\nssucc = (x \nsplus a)^\nssucc
\intertext{باقی حالت ایہ اے کہ $y \in \Domain{K}$ ہووے پر $x =
  a$۔ اتھے ایہناں حالتاں نوں وی وکھ کرنا پئے گا کہ $y =
  n$ معیاری اے یا $y = a$:}
(a \nsplus n^\nssucc) & = (a \nsplus (n+1)) = a = a^\nssucc = (a \nsplus n)^\nssucc\\
(a \nsplus a^\nssucc) & = (a \nsplus a) = a = a^\nssucc = (a \nsplus a)^\nssucc
\end{align*}
یقیناً ایہ ضرورت توں کجھ ودھ تفصیل اے۔ مثال دے طور اُتے،
جدوں $a \nsplus z = a$، چاہے $z$ جو وی ہووے، تاں اسیں فوراً
نتیجہ کڈھ سکدے آں کہ $a \nsplus a^\nssucc = a$۔
باقی مسلمات دی تصدیق وی ایس طرح کیتی جا سکدی اے۔

ایس طرح~$\Struct{K}$،~$\Th{Q}$ دا نمونہ اے۔ ایتھے
جمع~$\nsplus$ بدل ترتیب نال وی اوہی قدر دیندی اے۔
پر کجھ ہور جملے~$\Struct{N}$ وچ سچ تے~$\Struct{K}$ وچ
غلط نیں، تے کجھ ایس دے اُلٹ۔ مثال دے طور اُتے
$a \nsless a$، ایس لئی $\Sat{K}{\lexists[x][x < x]}$ تے
$\Sat/{K}{\lforall[x][\lnot x<x]}$۔ ایس توں ظاہر ہوندا اے
کہ $\Th{Q} \Proves/ \lforall[x][\lnot x < x]$۔
\end{ex}

\begin{prob}
ثابت کرو کہ \olref[mod][mar][mdq]{ex:model-K-of-Q} والی
~$\Struct{K}$،~$\Th{Q}$ دے باقی مسلمات وی پورے کردی اے:
\begin{align*}
  & \lforall[x][\eq[(x \times \Obj 0)][\Obj 0]] \tag{$!Q_6$}\\
  & \lforall[x][\lforall[y][\eq[(x \times y')][((x \times y) + x)]]] \tag{$!Q_7$}\\
  & \lforall[x][\lforall[y][(x < y \liff \lexists[z][\eq[(z' + x)][y]])]] \tag{$!Q_8$}
\end{align*}
صرف~$\prime$ والی کوئی !!a{sentence} لبھو جو~$\Struct{N}$ وچ
سچ پر~$\Struct{K}$ وچ غلط ہووے۔
\end{prob}

\begin{ex}
\ollabel{ex:model-L-of-Q} اوہ !!{structure}~$\Struct{L}$ لو جہدا
دائرہ $\Domain{L} = \Nat \cup \{a, b\}$ اے تے
$\Assign{\prime}{L} = \nssucc$، $\Assign{+}{L} = \nsplus$
دیاں تعبیراں ایہناں جدولاں نال دتیاں گئیاں نیں:
\[
\begin{array}{c|c}
  x & x^\nssucc \\
  \hline
  n & n+1 \\
  a & a\\
  b & b
\end{array}
\qquad
\begin{array}{c|ccc}
  x \nsplus y & m & a & b\\
  \hline
  n & n+m & b & a\\
  a & a & b & a\\
  b & b & b & a
\end{array}
\]
کیوں جے $\nssucc$ اک بہ اک (!!{injective}) اے، $0$ اوہدے
حاصل قدراں دے سیٹ وچ نہیں، تے ہر $x \in \Domain{L}$ جو
~$0$ نہیں، اوہدے حاصل قدراں دے سیٹ وچ اے، ایس لئی مسلمات
$!Q_1$--$!Q_3$،~$\Struct{L}$ وچ سچ نیں۔ ہر $x$ لئی
$x \nsplus 0 = x$، سو $!Q_4$ وی سچ اے۔ $!Q_5$ لئی
$x \nsplus y^\nssucc$ تے $(x \nsplus y)^\nssucc$ نوں
ویکھو۔ جے $x$ تے $y$ دونے معیاری ہون، ایہ برابر نیں؛ کیوں جے
ایس صورت وچ $\nssucc$ تے $\nsplus$، $\prime$ تے $+$
دے مطابق نیں۔ جے $x$ غیر معیاری تے $y$ معیاری ہووے،
تاں $x \nsplus y^\nssucc = x =
x^\nssucc = (x \nsplus y)^\nssucc$۔
جے $x$ تے $y$ دونے غیر معیاری ہون تاں چار حالتاں نیں:
\begin{align*}
&  a \nsplus a^\nssucc  = b = b^\nssucc = (a \nsplus a)^\nssucc\\
&  b \nsplus b^\nssucc  = a = a^\nssucc = (b \nsplus b)^\nssucc\\
&  b \nsplus a^\nssucc  = b = b^\nssucc = (b \nsplus a)^\nssucc\\
&  a \nsplus b^\nssucc  = a = a^\nssucc = (a \nsplus b)^\nssucc\\
\intertext{جے $x$ معیاری ہووے پر $y$ غیر معیاری ہووے، تاں}
&  n \nsplus a^\nssucc  = n \nsplus a = b = b^\nssucc = (n \nsplus a)^\nssucc\\
&  n \nsplus b^\nssucc  = n \nsplus b = a = a^\nssucc = (n \nsplus b)^\nssucc
\end{align*}
ایس لئی $\Sat{L}{!Q_5}$۔ پر $a \nsplus 0 \neq 0 \nsplus a$،
سو $\Sat/{L}{\lforall[x][\lforall[y][\eq[(x+y)][(y+x)]]]}$۔
\end{ex}

\begin{prob}
\olref[mod][mar][mdq]{ex:model-L-of-Q} والی~$\Struct{L}$ وچ
~$\nstimes$ تے $\nsless$ وی شامل کرو جو~$\times$ تے $<$
دیاں تعبیراں بنن۔ دکھاؤ کہ تہاڈی ساخت~$\Th{Q}$ دے باقی
مسلمات پورے کردی اے:
\begin{align*}
& \lforall[x][\eq[(x \times \Obj 0)][\Obj 0]] \tag{$!Q_6$}\\ &
  \lforall[x][\lforall[y][\eq[(x \times y')][((x \times y) + x)]]]
  \tag{$!Q_7$}\\ & \lforall[x][\lforall[y][(x < y \liff \lexists[z][\eq[(z'+x)][y]])]] \tag{$!Q_8$}
\end{align*}
\end{prob}

\begin{prob}
\olref[mod][mar][mdq]{ex:model-L-of-Q} والی~$\Struct{L}$ وچ
$a^\nssucc = a$ تے $b^\nssucc = b$۔ کی~$\Th{Q}$ دا کوئی
ایہو جیا نمونہ اے جس وچ $a^\nssucc = b$ تے
$b^\nssucc = a$ ہووے؟
\end{prob}

\begin{explain}
اسیں~$\Th{Q}$ دے ایہو جئے نمونے صاف صاف بنائے نیں جنہاں وچ
غیر معیاری !!{element}s، معیاری عنصراں توں پرلے پاسے
نیں۔ اصل وچ مسلمات ایہ گل لازم کردے نیں۔ غیر معیاری
!!{element}~$x$، ${} \nsless 0$ نہیں ہو سکدا، کیوں جے
$\Th{Q} \Proves \lforall[x][\lnot x<0]$
(\olref[inc][req][min]{lem:less-zero} ویکھو)۔ نالے ہر $n$
لئی $\Th{Q} \Proves \lforall[x][(x < \num{n}' \lif
(\eq[x][\num{0}] \lor \eq[x][\num{1}] \lor \dots \lor
\eq[x][\num{n}]))]$
(\olref[inc][req][min]{lem:less-nsucc})، سو
$a \nsless n$ کسے وی~$n>0$ لئی نہیں ہو سکدا۔
\end{explain}
\endgroup
% /OLP-0195

% OLP-0196 content/model-theory/models-of-arithmetic/models-of-pa.tex
\begingroup
\def\olfilename{models-of-pa}
\expandafter\def\csname ol@sectioncs\endcsname{section}

\olfileid{mod}{mar}{mpa}
\section{\texorpdfstring{$\Th{PA}$ دے نمونے}{PA دے نمونے}}

\begin{explain}
~$\Th{TA}$ دا ہر غیر معیاری نمونہ~$\Th{PA}$ دا وی نمونہ اے۔
اسیں جاندے آں کہ~$\Th{TA}$ دے، تے ایس لئی~$\Th{PA}$ دے،
غیر معیاری نمونے موجود نیں۔ ایہ وی جان دے آں کہ ایہناں
غیر معیاری نمونیاں وچ غیر معیاری عدد، یعنی دائرے دے
اوہ !!{element}s ہوندے نیں جو سارے معیاری اعداد توں
پرلے پاسے نیں۔ پر اوہ کس ترتیب وچ نیں؟ اوہ کنے نیں؟
اسیں ویکھ چکے آں کہ کمزور نظریے~$\Th{Q}$ دے نمونے
صرف اک غیر معیاری عدد وی رکھ سکدے نیں۔ پر ایہ سادیاں
!!{structure}s، $\Th{PA}$ یا $\Th{TA}$ دے نمونے نہیں۔

~$\Th{PA}$ یا $\Th{TA}$ دے نمونیاں دی ساخت سمجھن دی
کُنجی ایہ ویکھنا اے کہ ایہناں نظرییاں وچ کیہڑے حقائق
!!{derivable} نیں۔ مثال دے طور اُتے خود $\Th{PA}$،
$\lforall[x][\eq/[x][x']]$ تے
$\lforall[x][\lforall[y][\eq[(x+y)][(y+x)]]]$
ثابت کردا اے؛ ایس لئی جن سادیاں ساختاں وچ ایہ
!!{sentence}s غلط نیں، اوہ~$\Th{PA}$ دے نمونے نہیں ہو سکدیاں۔

من لو~$\Struct{M}$،~$\Th{PA}$ دا نمونہ اے۔ تاں جے
$\Th{PA} \Proves !A$، $\Sat{M}{!A}$۔ فیر
$\nszero$ نوں $\Assign{\Obj{0}}{M}$ لئی،
$\nssucc$ نوں $\Assign{\prime}{M}$ لئی،
$\nsplus$ نوں $\Assign{+}{M}$ لئی،
$\nstimes$ نوں $\Assign{\times}{M}$ لئی، تے
$\nsless$ نوں $\Assign{<}{M}$ لئی ورتیے۔
ہن ہر !!{sentence}~$!A$، $\nszero$، $\nssucc$،
$\nsplus$، $\nstimes$ تے $\nsless$ بارے کوئی شرط بیان
کردا اے، تے جے $\Sat{M}{!A}$ ہووے تاں اوہ شرط پوری
ہونی چاہیدی اے۔ مثال دے طور اُتے، جے $\Sat{M}{!Q_1}$،
یعنی
$\Sat{M}{\lforall[x][\lforall[y][(\eq[x'][y'] \lif \eq[x][y])]]}$،
تاں $\nssucc$ اک بہ اک (!!{injective}) ہونا چاہیدا اے۔
\end{explain}

\begin{prop}
~$\Struct{M}$ وچ $\nsless$ اک سخت خطی ترتیب اے، یعنی ایہناں
شرطاں نوں پورا کردی اے:
\begin{enumerate}
\item $x \nsless x$ کسے وی~$x \in \Domain{M}$ لئی نہیں۔
\item جے $x \nsless y$ تے $y \nsless z$، تاں $x \nsless z$۔
\item ہر $x \neq y$ لئی $x \nsless y$ یا $y \nsless x$۔
\end{enumerate}
\end{prop}

\begin{proof}
$\Th{PA}$ ایہ ثابت کردا اے:
\begin{enumerate}
\item $\lforall[x][\lnot x < x]$
\item $\lforall[x][\lforall[y][\lforall[z][((x < y \land y < z) \lif x < z)]]]$
\item $\lforall[x][\lforall[y][((x < y \lor y < x) \lor \eq[x][y]))]]$
\end{enumerate}
\end{proof}

\begin{prop}
\ollabel{prop:M-discrete} $\nszero$،~$\Domain{M}$ دا
$\nsless$-ترتیب وچ سب توں چھوٹا !!{element} اے۔
ہر $x$ لئی $x \nsless x^\nssucc$، تے
$x^\nssucc$ ایہ خاصیت رکھن والا $\nsless$ دے لحاظ نال
سب توں چھوٹا !!{element} اے۔ ہر غیر صفر $x$ لئی اک منفرد
$y$ اے جس لئی $y^\nssucc = x$۔
(اسیں $y$ نوں $x$ دا~$\Struct{M}$ وچ پیشرو آکھدے آں
تے اوہنوں~$^\nssucc x$ نال ظاہر کردے آں۔)
\end{prop}

\begin{proof}
مشق۔
\end{proof}

\begin{prob}
~$\Lang{L_A}$ وچ ایہو جئے !!{sentence}s لبھو جو
~$\Th{PA}$ وچ !!{derivable} ہون (تے ایس لئی
~$\Struct{N}$ وچ سچ ہون) تے
\olref[mod][mar][mpa]{prop:M-discrete} وچ دتیاں
$\nszero$، $\nssucc$ تے $\nsless$ دیاں خاصیتاں
لازم کرن۔
\end{prob}

\begin{prop}
~$\Struct{M}$ دے سارے معیاری !!{element}s،
~$\nsless$ دے لحاظ نال، سارے غیر معیاری !!{element}s
توں چھوٹے نیں۔
\end{prop}

\begin{proof}
اسیں $n$ نوں $\Value{\num{n}}{M}$، یعنی~$\Struct{M}$
دے معیاری !!{element}، لئی مختصر لکھاں گے۔
خود $\Th{Q}$ ہر~$n \in \Nat$ لئی
$\lforall[x][(x < \num{n}' \lif (\eq[x][\num{0}] \lor
  \eq[x][\num{1}] \lor \dots \lor \eq[x][\num{n}]))]$
ثابت کردا اے۔ کوئی !!{element}s، $\nsless \nszero$
نہیں۔ سو جے $n$ معیاری تے $x$ غیر معیاری ہووے،
تاں $x \nsless n$ نہیں ہو سکدا۔ تعریف مطابق غیر
معیاری عنصر اوہ اے جو $\Value{\num{n}}{M}$ نہیں،
کسے وی~$n \in \Nat$ لئی؛ ایس لئی $x \neq n$ وی۔
کیوں جے $\nsless$ خطی ترتیب اے، تاں لازماً
$n \nsless x$۔
\end{proof}

\begin{prop}
ہر غیر معیاری !!{element}~$x$،~$\Domain{M}$ وچ، ایس ذیلی سیٹ
دا عنصر اے:
\[
\dots ^{\nssucc\nssucc\nssucc}x \nsless ^{\nssucc\nssucc}x \nsless
^{\nssucc}x \nsless x \nsless x^{\nssucc} \nsless x^{\nssucc\nssucc}
\nsless x^{\nssucc\nssucc\nssucc} \nsless \dots
\]
ایس ذیلی سیٹ نوں اسیں $x$ دا \emph{بلاک} آکھدے آں تے
$[x]$ لکھدے آں۔ ایس دا نہ کوئی سب توں چھوٹا تے نہ کوئی
سب توں وڈا !!{element} اے۔ ایہ اوہناں
$y \in \Domain{M}$ دا سیٹ اے جنہاں لئی کوئی معیاری~$n$
ایہو جیا ہووے کہ $x \nsplus n = y$ یا
$y \nsplus n = x$۔
\end{prop}

\begin{proof}
ایہو جیا سیٹ~$[x]$ ہمیشہ موجود اے، کیوں جے
ہر غیر معیاری !!{element}~$y$ دا،~$\Domain{M}$ وچ، اک منفرد
جانشین~$y^\nssucc$ تے اک منفرد پیشرو~$^\nssucc y$
ہوندا اے۔ اک دوجے دے لگاتار !!{element}s $y$،
$y^\nssucc$ لئی $y \nsless y^\nssucc$، تے
$y^\nssucc$، $\nsless$ دے لحاظ نال
~$\Domain{M}$ دا اوہ سب توں چھوٹا !!{element} اے
جس توں $y$، $\nsless$ دے لحاظ نال چھوٹا اے۔
کیوں جے ہمیشہ $^\nssucc y \nsless y$ تے
$y \nsless y^\nssucc$، ایس لئی $[x]$ دا نہ کوئی
سب توں چھوٹا تے نہ کوئی سب توں وڈا !!{element} اے۔
جے $y \in [x]$ تاں $x \in [y]$؛ کیوں جے فیر
یا $y^{\nssucc\dots\nssucc} = x$ ہووے گا یا
$x^{\nssucc\dots\nssucc} = y$۔ جے
$y^{\nssucc\dots\nssucc} = x$ ہووے (جتھے
$n$ واری $\nssucc$ لگیا اے)، تاں
$y \nsplus n = x$، تے الٹ رخ وی؛ کیوں جے
$\Th{PA} \Proves \lforall[x][\eq[x^{\prime\dots\prime}][(x +
    \num{n})]]$ (جے $n$، $\prime$ دیاں وارییاں دی تعداد اے)۔
\end{proof}

\begin{prop}
جے $[x] \neq [y]$ تے $x \nsless y$، تاں ہر
$u \in [x]$ تے ہر $v \in [y]$ لئی $u \nsless v$۔
\end{prop}

\begin{proof}
ایہ گل نوٹ کرو کہ
$\Th{PA} \Proves \lforall[x][\lforall[y][(x < y \lif (x' < y
    \lor x' = y))]]$۔ ایس لئی جے $u \nsless v$
ہووے، تاں $u \nsplus n^\nssucc \nsless v$
وی اے، ہر~$n$ لئی، جدوں $[u] \neq [v]$ ہووے۔

ہر $u \in [x]$، $\nsless y$ اے: مفروضے مطابق
$x \nsless y$۔ جے $u \nsless x$، تاں تعدیت نال
$u \nsless y$۔ تے جے $x \nsless u$ پر
$u \in [x]$، تاں
$u = x\nsplus n^\nssucc$ کسے~$n$ لئی؛ ایس لئی ہُنھی ثابت کیتی
گل نال $u \nsless y$۔

ہن من لو کہ $v \in [y]$، $\nsless y$ اے؛ یعنی
$v \nsplus m^\nssucc = y$ کسے معیاری~$m$ لئی۔
ایس توں $v \nsless x$ ممکن نہیں، ورنہ
$y = v \nsplus m^\nssucc \nsless x$ ہوندا۔
صاف اے کہ $x \neq v$ وی، ورنہ
$x \nsplus m^\nssucc = v \nsplus m^\nssucc = y$
ہوندا تے $[x] = [y]$ ملدا۔ سو $x \nsless v$۔
فیر $x \nsplus n^\nssucc \nsless v$ ہر~$n$ لئی
وی اے۔ ایس لئی جے $x \nsless u$ تے
$u \in [x]$، تاں $u \nsless v$۔ جے
$u \nsless x$، تاں تعدیت نال $u \nsless v$۔

آخر وچ، جے $y \nsless v$، تاں $u \nsless v$،
کیوں جے جیویں اسیں دکھایا اے $u \nsless y$ تے
$y \nsless v$۔
\end{proof}

\begin{cor}
جے $[x] \neq [y]$، تاں $[x] \cap [y] = \emptyset$۔
\end{cor}

\begin{proof}
من لو $z \in [x]$ تے $x \nsless y$۔ تاں
$z \nsless u$ ہر $u \in [y]$ لئی۔ جے
$z \in [y]$ وی ہووے، تاں $z \nsless z$ ملدا۔
جے $y \nsless x$ ہووے تاں وی ایہو دلیل اے۔
\end{proof}

\begin{explain}
ایس دا مطلب اے کہ بلاکاں نوں خود ایس طرح ترتیب دتی جا سکدی
اے جو $\nsless$ دے مطابق ہووے: دو وکھرے بلاکاں لئی $[x] \nsless [y]$
اوہدوں تے صرف اوہدوں جدوں $x \nsless y$؛ یا ایس دے
برابر، جدوں $u \nsless v$ ہووے، کسے وی
$u \in [x]$ تے $v \in [y]$ لئی۔
صاف اے کہ معیاری بلاک $[0]$
سب توں چھوٹا بلاک اے۔ اوہ کسے غیر معیاری بلاک نال
نہیں ملدا، تے کوئی دو غیر معیاری بلاک وی آپس وچ
نہیں ملدے۔ خاص طور اُتے جانشین یا پیشرو وار وار
لا کے دوجے بلاک تک نہیں پہنچ سکدے۔
\end{explain}

\begin{prop}
جے $x$ تے $y$ غیر معیاری ہون، تاں
$x \nsless x \nsplus y$ تے $x
\nsplus y \notin [x]$۔
\end{prop}

\begin{proof}
جے $y$ غیر معیاری اے تاں $y \neq \nszero$۔
$\Th{PA} \Proves
\lforall[x][(y \neq \Obj{0} \lif x < (x+y))]$۔
ہن من لو $x \nsplus y \in [x]$۔ کیوں جے
$x \nsless x \nsplus y$، ایس لئی
$x \nsplus n^\nssucc = x \nsplus y$ ہوندا۔
پر $\Th{PA} \Proves
\lforall[x][\lforall[y][\lforall[z][(\eq[(x+y)][(x+z)] \lif y = z)]]]$
(جمع دی سانجھی دلیل کڈھن دا قانون)۔ ایس دا مطلب اے کہ
$y = n^\nssucc$ کسے معیاری~$n$ لئی؛ پر $y$ نوں
غیر معیاری منیا گیا سی۔
\end{proof}

\begin{prop}
کوئی سب توں چھوٹا غیر معیاری بلاک نہیں۔
\end{prop}

\begin{proof}
$\Th{PA} \Proves \lforall[x][\lexists[y][(\eq[(y+y)][x] \lor
      \eq[(y+y)'][x])]]$؛ یعنی ہر $x$ نوں~$2$ نال
ونڈیا جا سکدا اے (ممکن اے باقی~$1$ بچ جائے)۔
جے $x$ غیر معیاری ہووے تاں $y$ وی غیر معیاری اے۔
پچھلے ثابت کیتے بیان مطابق $y \nsless y \nsplus y$ تے
$y \nsplus y \notin [y]$۔ تاں
$y \nsless (y \nsplus y)^\nssucc$ تے
$(y \nsplus y)^\nssucc \notin [y]$ وی۔
پر $x = y \nsplus y$ یا
$x = (y \nsplus y)^\nssucc$، سو
$y \nsless x$ تے $y \notin [x]$۔
\end{proof}

\begin{prop}
غیر معیاری نمونے وچ کوئی سب توں وڈا بلاک نہیں۔
\end{prop}

\begin{proof}
مشق۔
\end{proof}

\begin{prob}
دکھاؤ کہ~$\Th{PA}$ دے غیر معیاری نمونے وچ
کوئی سب توں وڈا بلاک نہیں۔
\end{prob}

\begin{prop}
\ollabel{prop:blocks-dense}
بلاکاں دی ترتیب گھنی اے۔ یعنی جے $x \nsless y$
تے $[x] \neq [y]$، تاں اوہناں دوناں توں وکھرا
کوئی بلاک $[z]$ اوہناں دے وچکار اے۔
\end{prop}

\begin{proof}
من لو $x \nsless y$۔ پہلاں وانگ
$x \nsplus y$ نوں دو نال ونڈیا جا سکدا اے
(ممکن اے کجھ باقی بچ جائے): کوئی
$z \in \Domain{M}$ ایہو جیا اے کہ یا
$x \oplus y = z \oplus z$، یا
$x \oplus y = (z \oplus
z)^\nssucc$۔ عنصر $z$، $x$ تے~$y$ دا
“درمیانی قدر” اے، تے $x \nsless z$ تے
$z \nsless y$۔
\end{proof}

\begin{prob}
\olref[mod][mar][mpa]{prop:blocks-dense} دا تفصیلی
ثبوت لکھو۔ کیہڑا !!{sentence} $\Th{PA}$ نوں
!!{derive} کرنا پئے گا تاں جو~$z$ دی موجودگی
یقینی ہووے؟ $x \nsless z$ تے $z \nsless y$ کیوں،
تے $[x] \neq [z]$ تے $[z] \neq [y]$ کیوں؟
\end{prob}

\begin{explain}
قابلِ شمار غیر معیاری نمونے وچ، غیر معیاری بلاک ناطق اعداد وانگ ترتیب وچ
نیں: اوہ !!a{denumerable}، سِراں توں خالی،
گھنی خطی ترتیب بناؤندے نیں۔ ایہ دکھایا جا سکدا اے کہ
ایہو جیہیاں کوئی دو !!{denumerable} ترتیباں
ہم شکل ہوندیاں نیں۔ ایس توں نتیجہ نکلدا اے کہ
سچے حساب دے کسے وی دو !!{enumerable} غیر معیاری
نمونیاں $\Struct{M}_1$ تے $\Struct{M_2}$ دے،
صرف $<$ تے $=$ والی زبان تک گھٹائے روپ ہم شکل
ہوندے نیں۔ واقعی ہم شکلی $h$ نوں ایویں متعین
کیتا جا سکدا اے: $\Struct{M_1}$ تے
$\Struct{M_2}$ دے معیاری حصے، معیاری نمونے
$\Struct{N}$ نال ہم شکل نیں، ایس لئی آپس وچ
وی ہم شکل نیں۔ غیر معیاری حصے بناؤن والے بلاک
خود ناطق اعداد وانگ مرتب نیں، تے ایس لئی
ہم شکل نیں؛ بلاکاں دی ہم شکلی نوں ہر بلاک
\emph{دے اندر} وی ہم شکلی تک ایویں ودھایا جا
سکدا اے کہ ہر بلاک وچ من چاہے عنصراں نوں
ملایا جائے، تے فیر $x$ دے~$\Struct{M_1}$ وچ
جانشین دی تصویر، $x$ دی تصویر
دا~$\Struct{M_2}$ وچ جانشین رکھیا جائے۔
دھیان دیو کہ ایس توں ایہ \emph{نہیں} نکلدا کہ
$\mathfrak{M}_1$ تے $\mathfrak{M}_2$ حساب
دی پوری زبان وچ ہم شکل نیں (اصل وچ ہم شکلی
ہمیشہ کسے !!a{language} دے لحاظ نال ہوندی اے)؛
کیوں جے $\Domain{M_1}$ تے
$\Domain{M_2}$ اُتے جمع تے ضرب متعین کرن
دے ایسے طریقے وی نیں جو ہم شکل نہیں۔
(ایہ واٹ دے اک مشہور قضیے توں وی نکلدا اے
کہ مکمل نظریے دے قابل شمار نمونیاں دی تعداد~2
نہیں ہو سکدی۔)
\end{explain}
\endgroup
% /OLP-0196

% OLP-0197 content/model-theory/models-of-arithmetic/computable-models.tex
\begingroup
\def\olfilename{computable-models}
\expandafter\def\csname ol@sectioncs\endcsname{section}

\olfileid{mod}{mar}{cmp}
\section{حساب دے قابلِ حساب نمونے}

\begin{explain}
معیاری نمونے~$\Struct{N}$ دیاں دو چنگیاں خاصیتاں نیں۔
اوہدا دائرہ قدرتی اعداد~$\Nat$ نیں، یعنی اوہدے عنصر
اوہی چیزاں نیں جنہاں بارے اسیں حساب دی زبان وچ گل
کرنا چاہندے آں؛ تے معیاری عددی عبارت~$\num{n}$،
اصل وچ~$n$ نوں ہی ظاہر کردی اے۔ دوجی چنگی خاصیت
ایہ اے کہ~$\Lang{L_A}$ دیاں غیر منطقی علامتاں
دیاں ساریاں تعبیراں \emph{قابلِ حساب} نیں۔
جانشین، جمع تے ضرب دے فنکشن، جو
$\Assign{\prime}{N}$، $\Assign{+}{N}$ تے
$\Assign{\times}{N}$ دا کردار ادا کردے نیں،
اعداد دے قابلِ حساب فنکشن نیں۔ (مثال دے طور اُتے
ٹیورنگ مشیناں نال حساب کیتے جا سکدے نیں، یا
ابتدائی بازگشت نال متعین ہو سکدے نیں۔)
تے~$\Struct{N}$ اُتے چھوٹا ہون دا تعلق، یعنی
$\Assign{<}{N}$، فیصلہ کرن جوگ اے۔

$\Th{Q}$ تے $\Th{PA}$ ورگے حسابی نظرییاں دے غیر
معیاری نمونیاں وچ لازماً غیر معیاری عنصر ہوندے نیں۔
ایس لئی اوہناں دے دائریاں وچ عموماً~$\Nat$ دے
نال ہور !!{element}s وی ہوندے نیں۔ پر کوئی وی
گنتی وار لامتناہی !!{structure} کسے وی !!{denumerable}
سیٹ اُتے بنایا جا سکدا اے، جیویں~$\Nat$ اُتے۔
سو دائرہ~$\Nat$ رکھن والے غیر معیاری نمونے وی
موجود نیں۔ ایسے نمونے~$\Struct{M}$ وچ یقیناً
کجھ اعداد نوں اوہ کردار نہیں مل سکدا جو عام طور
اُتے اوہناں دا اے؛ کیوں جے کوئی~$k$ ایسا ہونا
لازمی اے جو $\Value{\num{n}}{M}$ توں
وکھرا ہووے، ہر~$n \in \Nat$ لئی۔
\end{explain}

\begin{defn}
!!^a{structure}~$\Struct{M}$،~$\Lang{L_A}$ لئی،
\emph{قابلِ حساب (computable)} اوہدوں تے صرف
اوہدوں اے جدوں $\Domain{M} = \Nat$ ہووے، تے
$\Assign{\prime}{M}$، $\Assign{+}{M}$،
$\Assign{\times}{M}$ قابلِ حساب فنکشن ہون،
تے $\Assign{<}{M}$ فیصلہ کرن جوگ تعلق ہووے۔
\end{defn}

\begin{ex}\ollabel{ex:comp-model-q}
\olref[mdq]{ex:model-K-of-Q} والی ساخت
$\Struct{K}$ یاد کرو۔ اوہدا دائرہ
$\Domain{K} = \Nat
\cup \{a\}$ سی تے تعبیراں ایہ سن:
\begin{align*}
  \Assign{\Obj{0}}{K} & = 0\\
  \Assign{\prime}{K}(x) & =
  \begin{cases}
    x+1 & \text{جے $x\in \Nat$}\\
    a & \text{جے $x = a$}
  \end{cases}\\
  \Assign{+}{K}(x, y) & =
  \begin{cases}
    x+y & \text{جے $x$, $y \in\Nat$}\\
    a & \text{باقی حالتاں وچ}
  \end{cases}\\
  \Assign{\times}{K}(x, y) & =
  \begin{cases}
    xy & \text{جے $x$, $y \in\Nat$}\\
    0 & \text{جے $x=0$ یا $y=0$}\\
    a & \text{باقی حالتاں وچ}\\
  \end{cases}\\
  \Assign{<}{K} & =
  \Setabs{\tuple{x,y}}{x, y \in \Nat \text{ تے } x<y} \cup
  \Setabs{\tuple{x,a}}{x \in \Domain{K}}
\end{align*}
پر $\Domain{K}$ !!{denumerable} اے تے ایس لئی
~$\Nat$ دے ہم تعداد اے۔ مثال دے طور اُتے
$g\colon \Nat \to \Domain{K}$ جس لئی $g(0) =
a$ تے $g(n) = n-1$، $n>0$ لئی، !!a{bijection}
اے۔ ایس نوں نویں نمونے~$\Struct{K'}$، جو
~$\Th{Q}$ دا نمونہ اے، تے $\Struct{K}$ وچکار
ہم شکلی وچ بدل سکدے آں۔ $\Struct{K'}$ وچ
~$\Lang{L_A}$ دیاں علامتاں لئی مختلف فنکشن
تے تعلق مقرر کرنے پین گے، کیوں جے
~$\Nat$ دے وکھرے !!{element}s معیاری تے
غیر معیاری اعداد دے کردار ادا کردے نیں۔

خاص طور اُتے $0$ ہن~$a$ دا کردار ادا کردا اے،
سب توں چھوٹے معیاری عدد دا نہیں۔ سب توں چھوٹا
معیاری عدد ہن~$1$ اے۔ ایس لئی اسیں
$\Assign{\Obj{0}}{K'} = 1$ مقرر کردے آں۔
ہن جانشین فنکشن وی وکھرا اے: جے معیاری عدد،
یعنی $n > 0$، دتا جاوے، تاں ایہ فیر وی
$n+1$ دیندا اے۔ پر $0$ ہن~$a$ دا کردار
ادا کردا اے، جو آپ ہی اپنا جانشین اے۔ سو
$\Assign{\prime}{K'}(0) = 0$۔
جمع تے ضرب لئی وی ایویں اے:
\begin{align*}
\Assign{+}{K'}(x, y) & =
  \begin{cases}
    x+y-1 & \text{جے $x$, $y >0$}\\
    0 & \text{باقی حالتاں وچ}
  \end{cases}\\
  \Assign{\times}{K'}(x, y) & =
  \begin{cases}
    1 & \text{جے $x = 1$ یا $y = 1$}\\
    xy - x - y + 2 & \text{جے $x$, $y > 1$}\\
    0 & \text{باقی حالتاں وچ}\\
  \end{cases}
\end{align*}
تے $\tuple{x, y} \in \Assign{<}{K'}$ اوہدوں
تے صرف اوہدوں جدوں $x < y$ تے $x > 0$ تے
$y > 0$، یا جے $y = 0$ ہووے۔

ایہ سارے فنکشن قدرتی اعداد دے قابلِ حساب
فنکشن نیں تے $\Assign{<}{K'}$،~$\Nat$
اُتے فیصلہ کرن جوگ تعلق اے؛ پر ایہ
~$\Nat$ اُتے جانشین، جمع تے ضرب والے اوہی
فنکشن نہیں، تے $\Assign{<}{K'}$ وی
$<$ والا~$\Nat$ اُتے اوہی تعلق نہیں۔
\end{ex}

\begin{prob}
ایہو جی !!a{structure}~$\Struct{L'}$ دیو جس دا
$\Domain{L'} = \Nat$ ہووے تے جو
\olref[mod][mar][mdq]{ex:model-L-of-Q} والی
~$\Struct{L}$ نال ہم شکل ہووے۔
\end{prob}

\begin{explain}
\Olref{ex:comp-model-q} توں پتا چلدا اے کہ
~$\Th{Q}$ دے دائرہ~$\Nat$ والے قابلِ حساب
غیر معیاری نمونے نیں۔ پر اگلا نتیجہ دکھاؤندا
اے کہ~$\Th{PA}$ دے نمونیاں لئی ایہ گل سچ
نہیں (تے ایس لئی~$\Th{TA}$ دے نمونیاں لئی وی نہیں)۔
\end{explain}

\begin{thm}[ٹیننبام دا قضیہ]
ہم شکلی دے لحاظ نال~$\Struct{N}$،~$\Th{PA}$ دا اکلوٹا قابلِ حساب
نمونہ اے۔
\end{thm}
\endgroup
% /OLP-0197

\OLEndChapterHook
\endgroup
% /OLP-0191

% OLP-0198 content/model-theory/interpolation/interpolation.tex
\begingroup
\def\olfilename{interpolation}
\expandafter\def\csname ol@sectioncs\endcsname{section}

\olchapter{mod}{int}{درمیانی جملے دا قضیہ (انٹرپولیشن)}

% OLP-0199 content/model-theory/interpolation/introduction.tex
\begingroup
\def\olfilename{introduction}
\expandafter\def\csname ol@sectioncs\endcsname{section}

\olfileid{mod}{int}{int}

\olsection{تعارف}

درمیانی جملے دے قضیے دا بیان ایہ اے: فرض کرو دتے دو
فارمولے جملے نیں تے
$\Entails !A \lif !B$۔ تاں اک !!a{sentence} $!C$
موجود اے جس لئی $\Entails !A \lif !C$ تے
$\Entails !C \lif !B$۔ ایس توں ودھ، ہر
!!{constant}، !!{function} تے !!{predicate}
($\eq$ توں علاوہ) جو $!C$ وچ اے، $!A$
تے~$!B$ دوہاں وچ وی آؤندا اے۔
!!{sentence} $!C$ نوں
$!A$ تے~$!B$ دا \emph{درمیانی جملہ
(interpolant)} آکھیا جاندا اے۔

درمیانی جملے دا قضیہ اپنی تھاں دلچسپ اے، پر
اوہدی بنیادی اہمیت ایہ اے کہ ایس نال کسے نظریے
وچ قابلِ تعریف ہون بارے نتائج ثابت کیتے جا سکدے
نیں، تے ایہ وی پتہ لگایا جا سکدا اے کہ کنہاں
شرطاں تحت دو بے تضاد نظرییاں نوں ملاؤن توں
بے تضاد نظریہ بندا اے۔ پہلے نتیجے نوں بیتھ دا
قابلِ تعریف ہون دا قضیہ آکھدے نیں؛ دوجے نوں
رابنسن دا مشترکہ بے تضادی دا قضیہ۔
\endgroup
% /OLP-0199

% OLP-0200 content/model-theory/interpolation/separation.tex
\begingroup
\def\olfilename{separation}
\expandafter\def\csname ol@sectioncs\endcsname{section}

\olfileid{mod}{int}{sep}

\olsection{\printtoken{P}{sentence} نوں وکھ کرنا}

درمیانی جملے دے قضیے دے ثبوت توں پہلے کجھ
بنیاد رکھنی پئے گی۔ $!A$ تے $!B$ لئی
درمیانی جملہ اوہ !!a{sentence}~$!C$ اے
جس لئی $!A \Entails !C$ تے
$!C \Entails !B$۔ تقابلِ نقیض راہیں
پچھلی گل اوہدوں تے صرف اوہدوں سچ اے جدوں
$\lnot !B \Entails \lnot
!C$۔ ایس خاصیت والے !!^a{sentence}~$!C$
بارے آکھدے آں کہ اوہ \emph{وکھ کردا} اے
$!A$ تے $\lnot !B$ نوں۔ سو $!A$ تے
$!B$ لئی درمیانی جملہ لبھنا، ایسا
!!a{sentence} لبھن دے برابر اے جو $!A$
تے $\lnot !B$ نوں وکھ کرے۔ عام طور وانگ
اک وسیع صورت فائدے والی ہووے گی: ایسا
جملہ جو !!{sentence}s دے دو \emph{سیٹاں}
نوں وکھ کرے۔

\begin{defn}
اک جملہ $!C$، جملیاں دے سیٹاں $\Gamma$
تے $\Delta$ نوں \emph{وکھ کردا} اے اوہدوں
تے صرف اوہدوں جدوں $\Gamma \Entails !C$
تے $\Delta \Entails
\lnot !C$۔ جے ایسا کوئی !!{sentence} موجود
نہ ہووے، تاں $\Gamma$ تے $\Delta$
\emph{وکھ نہ کیتے جا سکّن والے} نیں۔
\end{defn}

$\Gamma$، $\Delta$ تے $!C$ دے نمونیاں
دیاں جماعتاں وچ شمولیت دے تعلق ہیٹھاں
وکھائے گئے نیں:
\begin{figure}[h]
  \centering
  \begin{tikzpicture}[node distance=2cm, auto, thick]
    \draw [rounded corners] (0,0) -- (6,0) -- (6,3) -- (0,3) --  cycle;
    \draw (1.5,1.5) circle (0.9cm);
    \draw (4.5,1.5) circle (0.9cm);
    \path node at (1.5,1.5) {\Large $\Gamma$}; 
    \path node at (4.5,1.5) {\Large $\Delta$}; 
    \path node at (0.4,2.6) {\Large $\formula{C}$};
    \path node at (3.4,0.4) {\Large $\lnot \formula{C}$};
    \draw (2.5,0) .. controls (2.5,1.5) and (3.5,1.5) .. (3.5,3);
  \end{tikzpicture}
  \caption{$!C$، $\Gamma$ تے $\Delta$ نوں وکھ کردا اے}
  \ollabel{fig:sep}
\end{figure}


\begin{lem}\ollabel{lem:sep1}
فرض کرو $\Lang{L}_0$ اوہ زبان اے جس وچ
ہر اوہ !!{constant}، !!{function} تے
!!{predicate} ($\doteq$ توں علاوہ) اے جو
$\Gamma$ تے $\Delta$ \emph{دوہاں}
وچ آؤندا اے؛ تے $\Lang{L}'_0$،
لامتناہی نویں !!{constant}s $\Obj c_n$،
$n \ge 0$ لئی، شامل کرکے ملن والی
زبان ہووے۔ تاں جے $\Gamma$ تے $\Delta$،
$\Lang{L}_0$ وچ وکھ نہ کیتے جا سکّن،
اوہ $\Lang{L}'_0$ وچ وی وکھ نہ کیتے
جا سکّن گے۔
\end{lem}

\begin{proof}
اسیں بالواسطہ ثبوت دیندے آں: تضاد لئی فرض کرو
کہ $\Gamma$ تے $\Delta$،
$\Lang{L}'_0$ وچ وکھ کیتے جا سکدے نیں۔
تاں $\Gamma \Entails
\Subst{!C}{c}{x}$ تے
$\Delta \Entails \lnot \Subst{!C}{c}{x}$
کسے $!C \in \Lang{L}_0$ لئی ہون گے
(ایتھے $c$ اک نواں !!{constant}
اے؛ جے $!C$ وچ ایسے اک توں ودھ نویں
!!{constant} ہون، اوہ صورت وی ایس ورگی اے)۔
محدود گواہی دے قضیے نال \emph{محدود}
ذیلی سیٹ $\Gamma_0$، $\Gamma$ وچوں،
تے $\Delta_0$، $\Delta$ وچوں ملدے نیں،
جنہاں لئی $\Gamma_0 \Entails \Subst{!C}{c}{x}$
تے $\Delta_0 \Entails \lnot \Subst{!C}{c}{x}$۔
$!G$ نوں $\Gamma_0$ دے سارے !!{formula}s
دا منطقی جوڑ، تے $!H$ نوں $\Delta_0$ دے
سارے !!{formula}s دا منطقی جوڑ من لو۔ تاں
\begin{align*}
  \formula G & \Entails \Subst{!C}{c}{x}, & \formula H  \Entails \lnot \Subst{!C}{c}{x}.
\end{align*}
پہلے نتیجے توں عمومیت دے قاعدے راہیں
$!G \Entails
\lforall[x][!C]$ ملدا اے؛ دوجے
توں تقابلِ نقیض راہیں
$\Subst{!C}{c}{x} \Entails \lnot !H$
ملدا اے، جس توں
$\lforall[x][!C]
\Entails \lnot !H$ وی۔ فیر
تقابلِ نقیض نال $!H \Entails
\lnot \lforall[x][!C]$ ملدا اے۔
یک رُخیت نال
\begin{align*}
  \Gamma &\Entails \lforall[x][!C], & 
  \Delta & \Entails \lnot \lforall[x][!C],
\end{align*}
ایس لئی $\lforall[x][!C]$،
$\Gamma$ تے $\Delta$ نوں
$\Lang{L}_0$ وچ وکھ کردا اے۔
\end{proof}

\begin{lem}\ollabel{lem:sep2}
فرض کرو $\Gamma \cup \{ \lexists[x][!S] \}$
تے $\Delta$ وکھ نہ کیتے جا سکّن، تے
$c$ اک نواں !!{constant} ہووے جو
$\Gamma$، $\Delta$ یا $!S$ وچ نہ ہووے۔
تاں $\Gamma \cup \{ \lexists[x][!S], \Subst{!S}{c}{x} \}$
تے $\Delta$ وی وکھ نہ کیتے جا سکّن گے۔
\end{lem}

\begin{proof}
تضاد لئی فرض کرو کہ $!C$،
$\Gamma \cup \{
\lexists[x][!S], \Subst{!S}{c}{x}\}$
تے $\Delta$ نوں وکھ کردا اے، جدوں کہ
نال ہی $\Gamma \cup \{\lexists[x]{!S} \}$
تے $\Delta$ وکھ نہ کیتے جا سکدے ہون۔
دو صورتاں وکھ کردے آں:
\begin{enumerate}
\item جے $c$، $!C$ وچ نہ آوے، تاں
  $\Gamma \cup
  \{\lexists[x][!S], \lnot!C \}$
  پورا ہون جوگا اے (ورنہ $!C$،
  $\Gamma \cup \{\lexists[x][!S] \}$
  تے $\Delta$ نوں وکھ کردا)۔
  $\Subst{!S}{c}{x}$ شامل کرن توں بعد وی
  ایہ پورا ہون جوگا رہندا اے؛ ایس لئی
  $!C$، $\Gamma \cup \{ \lexists[x][!S], \Subst{!S}{c}{x} \}$
  تے $\Delta$ نوں آخرکار وکھ نہیں کردا۔
\item جے $c$، $!C$ وچ آوے، تاں
  $!C$ دی صورت
  $\Subst{!C}{c}{x}$ اے۔ ایس توں
  \[
  \Gamma \cup \{ \lexists[x][!S], \Subst{!S}{c}{x}\} \Entails \Subst{!C}{c}{x},
  \]
  ملدا اے؛ فیر استنباط دے قضیے
  تے عمومیت نال
  $\Gamma, \lexists[x][!S] \Entails \lforall[x][(!S \lif !C)]$
  تے آخر $\Gamma
  \cup \{ \lexists[x][!S] \} \Entails \lexists[x][!C]$۔
  دوجے پاسے $\Delta \Entails \lnot \Subst{!C}{c}{x}$
  اے؛ سو عمومیت نال
  $\Delta \Entails \lnot \lexists[x][!C]$۔
  ایس لئی $\Gamma
  \cup \{\lexists[x][!S] \}$
  تے $\Delta$ وکھ کیتے جا سکدے نیں،
  جو تضاد اے۔
\end{enumerate}
\end{proof}
\endgroup
% /OLP-0200

% OLP-0201 content/model-theory/interpolation/interpolation-proof.tex
\begingroup
\def\olfilename{interpolation-proof}
\expandafter\def\csname ol@sectioncs\endcsname{section}

\olfileid{mod}{int}{prf}

\olsection{کریگ دا درمیانی جملے دا قضیہ}

\begin{thm}[کریگ دا درمیانی جملے دا قضیہ]
\ollabel{thm:interpol} 
فرض کرو دتے دو فارمولے جملے نیں تے
$\Entails !A \lif !B$، تاں اک !!a{sentence}
$!C$ موجود اے جس لئی
$\Entails !A \lif !C$ تے
$\Entails !C \lif !B$، تے ہر
!!{constant}، !!{function} تے
!!{predicate} ($\eq$ توں علاوہ)
جو $!C$ وچ اے، $!A$
تے~$!B$ دوہاں وچ وی آؤندا اے۔
!!{sentence} $!C$ نوں $!A$ تے~$!B$
دا \emph{درمیانی جملہ (interpolant)}
آکھیا جاندا اے۔
\end{thm}

\begin{proof}
فرض کرو $\Lang{L}_1$، $!A$ دی زبان،
تے $\Lang{L}_2$، $!B$ دی زبان اے۔
$\Lang{L}_0 = \Lang{L}_1 \cap \Lang{L}_2$
رکھو۔ ہر $i \in \{0, 1, 2 \}$ لئی
$\Lang{L}'_i$، $\Lang{L}_i$ وچ
لامتناہی نویں !!{constant}s $\Obj c_0,
\Obj c_1, \Obj c_2, \dots$ شامل کرکے بناؤ۔

جے $!A$ پورا ہون جوگا نہ ہووے،
$\lexists[x][\eq/[x][x]]$ اک درمیانی
جملہ اے۔ جے $\lnot !B$ پورا ہون جوگا
نہ ہووے (تے ایس لئی $!B$ منطقی طور
اُتے معتبر ہووے)، تاں
$\lexists[x][\eq[x][x]]$ اک درمیانی
جملہ اے۔ سو اسیں ایہ وی فرض کر سکدے آں
کہ $!A$ تے $\lnot !B$ دوہاں
پورے ہون جوگے نیں۔

درمیانی جملے دے قضیے دا تقابلِ نقیض
ثابت کرن لئی فرض کرو کہ $!A$ تے
$!B$ لئی کوئی درمیانی جملہ نہیں۔
دوجے لفظاں وچ فرض کرو کہ $\{!A\}$
تے $\{\lnot !B\}$،
$\Lang{L}_0$ وچ وکھ نہ کیتے
جا سکدے نیں۔

ساڈا مقصد جوڑی $(\{ !A \}, \{\lnot!B\})$
نوں اک ایسی غیر قابلِ تفریق جوڑی
$(\Gamma^*, \Delta^*)$ تک ودھانا اے
جس نوں ہور نہ ودھایا جا سکے۔
$!A_0$، $!A_1$، $!A_2$، \dots،
$\Lang{L}'_1$ دے !!{sentence}s دی
گنتی وار فہرست ہون؛ تے $!B_0$،
$!B_1$، $!B_2$، \dots،
~$\Lang{L}'_2$ دے !!{sentence}s دی۔
اسیں !!{sentence}s دے سیٹاں دیاں دو
ودھدیاں سلسلیاں $(\Gamma_n, \Delta_n)$،
$n \ge 0$ لئی، ایویں متعین کردے آں۔
$\Gamma_0 = \{ !A\}$ تے
$\Delta_0 = \{\lnot !B \}$ رکھو۔
فرض کرو $(\Gamma_n, \Delta_n)$ پہلے
توں متعین نیں؛ $\Gamma_{n+1}$ تے
$\Delta_{n+1}$ ایویں متعین کرو:
ہر اگلا سیٹ پچھلا سیٹ رکھدا اے؛ ہیٹھاں دتی شرط
پوری ہووے تاں صرف دتے اضافے کرو، نہیں تاں اوہ سیٹ
جیویں اے رہن دیو۔ گواہ لئی شامل کیتے مشترک نویں
مستقلاں وچوں اک چنو:
\begin{enumerate}
\item جے $\Gamma_n \cup \{!A_n \}$ تے
  $\Delta_n$، $\Lang{L}'_0$ وچ
  وکھ نہ کیتے جا سکّن، تاں $!A_n$ نوں
  $\Gamma_{n+1}$ وچ رکھو۔ ہور ایہ کہ جے
  $!A_n$ وجودی !!{formula}
  $\lexists[x][!S]$ ہووے، تاں ایسا نواں
  !!{constant} $c$ چنو جو
  $\Gamma_n$، $\Delta_n$، $!A_n$ یا
  $!B_n$ وچ نہ آؤندا ہووے، تے
  $\Subst{!S}{c}{x}$ نوں
  $\Gamma_{n+1}$ وچ رکھو۔
\item جے $\Gamma_{n+1}$ تے
  $\Delta_n \cup \{!B_n \}$،
  $\Lang{L}'_0$ وچ وکھ نہ کیتے جا سکّن،
  تاں $!B_n$ نوں $\Delta_{n+1}$ وچ رکھو۔
  ہور ایہ کہ جے $!B_n$ وجودی
  !!{formula} $\lexists[x][!S]$ ہووے،
  تاں ایسا نواں !!{constant} $c$
  چنو جو $\Gamma_{n+1}$، $\Delta_n$،
  $!A_n$ یا $!B_n$ وچ نہ آؤندا ہووے،
  تے $\Subst{!S}{c}{x}$ نوں
  $\Delta_{n+1}$ وچ رکھو۔
\end{enumerate}
آخر وچ ایہ متعین کرو:
\begin{align*}
  \Gamma^* & = \bigcup_{n\ge 0} \Gamma_n, & 
  \Delta^* & = \bigcup_{n\ge 0} \Delta_n.
\end{align*}
ہن $n$ اُتے دوہری استقرا نال
اسیں ثابت کر سکدے آں:
\begin{enumerate}
\item\ollabel{part-a} $\Gamma_n$ تے
  $\Delta_n$، $\Lang{L}'_0$
  وچ وکھ نہ کیتے جا سکّن؛
\item\ollabel{part-b} $\Gamma_{n+1}$
  تے $\Delta_n$، $\Lang{L}'_0$
  وچ وکھ نہ کیتے جا سکّن۔
\end{enumerate}
\olref{part-a} دی ابتدا
\olref[sep]{lem:sep1} توں ملدی اے۔
\olref{part-b} لئی تین صورتاں
وکھ کرنیاں پین گیاں:
\begin{enumerate}
\item جے $\Gamma_0 \cup \{!A_0 \}$
  تے $\Delta_0$ وکھ کیتے جا سکّن،
  تاں $\Gamma_1 = \Gamma_0$ تے
  \olref{part-b} صرف \olref{part-a}
  اے؛
\item جے $\Gamma_1 = \Gamma_0 \cup\{ !A_0\}$،
  تاں بناوٹ دے مطابق $\Gamma_1$ تے
  $\Delta_0$ وکھ نہیں کیتے جا سکدے۔
\item ہن اوہ صورت رہندی اے جدوں $!A_0$
  وجودی ہووے، سو
  $\Gamma_1 = \Gamma_0 \cup \{ \lexists[x][!S], \Subst{!S}{c}{x}
  \}$۔ بناوٹ دے مطابق
  $\Gamma_0 \cup \{ \lexists[x][!S]\}$
  تے $\Delta_0$ وکھ نہیں کیتے جا سکدے؛
  سو \olref[sep]{lem:sep2} نال
  $\Gamma_0 \cup \{ \lexists[x][!S], \Subst{!S}{c}{x} \}$
  تے $\Delta_0$ وی وکھ نہیں کیتے جا سکدے۔
\end{enumerate}
ایس نال اُتے \olref{part-a} تے
\olref{part-b} لئی استقرا دی ابتدا
مکمل ہوندی اے۔ ہن استقرائی قدم۔
\olref{part-a} لئی، جے
$\Delta_{n+1} \supseteq \Delta_n \cup \{ !B_n \}$
تاں $\Gamma_{n+1}$ تے
$\Delta_{n+1}$ بناوٹ دے مطابق وکھ
نہیں کیتے جا سکدے (جے $!B_n$
وجودی وی ہووے تاں
\olref[sep]{lem:sep2} نال)؛ جے
$\Delta_{n+1} = \Delta_n$ ہووے
(کیوں جے $\Gamma_{n+1}$ تے
$\Delta_n \cup \{!B_n\}$ وکھ کیتے
جا سکدے نیں)، تاں
\olref{part-b} دا استقرائی مفروضہ ورتدے
آں۔ \olref{part-b} دے استقرائی قدم لئی،
جے $\Gamma_{n+2} \supseteq \Gamma_{n+1} \cup
\{!A_{n+1} \}$ تاں $\Gamma_{n+2}$
تے $\Delta_{n+1}$ بناوٹ دے مطابق
وکھ نہیں کیتے جا سکدے (جے $!A_{n+1}$
وجودی وی ہووے تاں
\olref[sep]{lem:sep2} نال)؛ تے جے
$\Gamma_{n+2} = \Gamma_{n+1}$ ہووے،
تاں \olref{part-a} دا ہن ثابت ہویا
استقرائی قدم ورتدے آں۔ ایس نال
\olref{part-a} تے \olref{part-b}
اُتے استقرا مکمل ہوندی اے۔

ایس توں $\Gamma^*$ تے $\Delta^*$
وکھ نہیں کیتے جا سکدے؛ جے ایسا نہ
ہووے، تاں محدود گواہی دے قضیے نال
کوئی $n \ge 0$ ایسا ہووے گا کہ
$\Gamma_n$ تے $\Delta_n$ نوں
وکھ کیتا جا سکے، \olref{part-a} دے خلاف۔
خاص طور اُتے $\Gamma^*$ تے
$\Delta^*$ بے تضاد نیں: کیوں جے جے
پہلا یا دوجا سیٹ متضاد ہووے، تاں
اوہناں نوں بالترتیب
$\lexists[x][\eq/[x][x]]$ یا
  $\lforall[x][\eq[x][x]]$
نال وکھ کیتا جا سکدا اے۔

ہن اسیں دکھاؤندے آں کہ $\Gamma^*$،
$\Lang{L}'_1$ وچ مکمل بے تضاد اے
تے ایویں $\Delta^*$،
$\Lang{L}'_2$ وچ۔ پہلے لئی فرض کرو
کہ $!A_n \notin \Gamma^*$ تے
$\lnot !A_n
\notin \Gamma^*$، کسے
$n \ge 0$ لئی۔ جے
$!A_n \notin \Gamma^*$،
تاں $\Gamma_n \cup \{!A_n \}$
نوں $\Delta_n$ توں وکھ کیتا
جا سکدا اے، تے ایس لئی کوئی
$!C \in \Lang{L}'_0$ اے جس لئی
دوہاں گلاں ایہ نیں:
\begin{align*}
  \Gamma^* & \Entails !A_n \lif !C, & 
  \Delta^* & \Entails \lnot !C.
\end{align*}
ایویں جے $\lnot !A_n \notin \Gamma^*$،
تاں کوئی $!C' \in
\Lang{L}'_0$ اے جس لئی دوہاں گلاں:
\begin{align*}
  \Gamma^* & \Entails \lnot !A_n \lif !C', & 
  \Delta^* & \Entails \lnot !C'.
\end{align*}
قضیاتی منطق نال
$\Gamma^* \Entails !C \lor !C'$ تے
$\Delta^* \Entails \lnot (!C \lor !C')$؛
سو $!C \lor
!C'$، $\Gamma^*$ تے $\Delta^*$
نوں وکھ کردا اے۔ ایس ورگی دلیل
دکھاؤندی اے کہ $\Delta^*$ وی مکمل اے۔

آخر وچ اسیں دکھاؤندے آں کہ
$\Gamma^* \cap \Delta^*$،
$\Lang{L}'_0$ وچ مکمل بے تضاد اے۔
ایہ ظاہر طور اُتے بے تضاد اے، کیوں
جے ایہ بے تضاد سیٹاں دا مشترک حصہ اے۔
مکمل پن دکھاؤن لئی $!S \in
\Lang{L}'_0$ نوں اک جملہ من لو۔ ہن $\Gamma^*$،
$\Lang{L'_1}
\supseteq \Lang{L'_0}$ وچ مکمل اے،
تے ایویں $\Delta^*$،
$\Lang{L'_2} \supseteq \Lang{L'_0}$
وچ۔ ایس لئی یا $!S \in \Gamma^*$
یا $\lnot !S \in \Gamma^*$، تے
یا $!S \in \Delta^*$ یا
$\lnot !S \in \Delta^*$۔ جے
$!S \in
\Gamma^*$ تے $\lnot !S \in \Delta^*$
تاں $!S$، $\Gamma^*$ تے
$\Delta^*$ نوں وکھ کردا؛ تے جے
$\lnot !S \in
\Gamma^*$ تے $!S \in \Delta^*$
تاں $\Gamma^*$ تے $\Delta^*$
نوں $\lnot !S$ وکھ کردا۔ ایس لئی
یا $!S \in
\Gamma^* \cap \Delta^*$ یا
$\lnot !S \in \Gamma^* \cap \Delta^*$،
تے $\Gamma^* \cap \Delta^*$
مکمل اے۔

بناوٹ وچ ہر قبول کیتے وجودی جملے نال گواہ شامل کیتا گیا اے۔
ہر بند ٹرم دے برابر ہون والا عنصر موجود ہون دی شرط ہمیشہ
صادق اے؛ ودھائی زبان دی فہرست وچ ایہ شرط وی آؤندی اے
تے اوہدے لئی شامل کیتا مشترک نواں مستقل گواہ ملدا اے۔
سو بند ٹرماں دے برابری والے ہر طبقے دا نمائندہ ایسا مستقل اے۔
کیوں جے $\Gamma^*$ مکمل بے تضاد تے گواہاں والا اے،
ایس دا اک نمونہ $\Struct{M}'_1$ اے
جس دا !!{domain} $\Domain{M'_1}$
عین اوہناں عنصراں اُتے مشتمل اے
جو شامل کیتے مشترک نویں !!{constant}s دی تعبیر
$\Assign{c}{M'_1}$ نال ملدے نیں؛
بالکل تمامیت دے قضیے دے ثبوت وانگ
(\olref[fol][com][cth]{thm:completeness})۔
ایویں $\Delta^*$ دا اک نمونہ
$\Struct{M}'_2$ اے جس دا !!{domain}
$\Domain{M'_2}$،
شامل کیتے مشترک نویں !!{constant}s دیاں تعبیراں
$\Assign{c}{M'_2}$ نال ملدا اے۔

$\Struct{M_1}$، $\Struct{M'_1}$
توں $\Lang{L'_1} \setminus \Lang{L'_0}$
وچلی !!{constant}s، !!{function}s تے
!!{predicate}s دیاں تعبیراں ہٹا کے
بناؤ؛ تے $\Struct{M_2}$ وی ایسے طرح۔
ایتھے نقشے دی تعریف وچ مستقل شامل کیتے مشترک نویں
مستقلاں وچوں ورتیا جاندا اے؛ پچھلے قدم نال ایہ دوہاں
دائریاں دے ہر عنصر نوں ناں دیندے نیں۔
تاں نقشہ $h \colon M_1 \to M_2$
جو $h(\Assign{c}{M'_1}) = \Assign{c}{M'_2}$
نال متعین اے، $\Lang{L}'_0$
وچ ہم شکلی اے، کیوں جے
$\Gamma^* \cap \Delta^*$،
$\Lang{L}'_0$ وچ مکمل بے تضاد اے،
جیویں دکھایا۔ ایہ ایس لئی نکلدا
اے کہ کوئی وی $\Lang{L}'_0$-!!{sentence}
یا $\Gamma^*$ تے $\Delta^*$ دوہاں
وچ اے یا کسے وچ وی نہیں: سو
$\Assign{c}{M'_1} \in
\Assign{P}{M'_1}$ اوہدوں تے صرف
اوہدوں جدوں $\Atom{P}{c} \in \Gamma^*$،
اوہدوں تے صرف اوہدوں جدوں
$\Atom{P}{c} \in \Delta^*$، اوہدوں
تے صرف اوہدوں جدوں
$\Assign{c}{M'_2} \in
\Assign{P}{M'_2}$۔ ہم شکلی دیاں
باقی شرطاں وی ایویں ثابت کیتیاں
جا سکدیاں نیں۔

ہن اک نمونہ $\Struct{M}$،
!!{language}
$\Lang{L'_1} \cup \Lang{L'_2}$ لئی
ایویں متعین کریے:
\begin{enumerate}
\item !!{domain} $\Domain{M}$،
  عین $\Domain{M_2}$ اے، یعنی سارے
  عنصر $\Assign{c}{M'_2}$ دا سیٹ؛
\item جے !!a{predicate}~$P$،
  $\Lang{L_2} \setminus
  \Lang{L_1}$ وچ اے، تاں
  $\Assign{P}{M} = \Assign{P}{M'_2}$؛
\item جے محمول $P$،
  $\Lang{L}_1\setminus \Lang{L}_2$
  وچ اے، تاں
  $\Assign{P}{M} = h(\Assign{P}{M'_1})$،
  یعنی $\tuple{\Assign{c_1}{M'_2}, \dots, \Assign{c_n}{M'_2}} \in
  \Assign{P}{M}$ اوہدوں تے صرف اوہدوں
  جدوں $\tuple{\Assign{c_1}{M'_1}, \dots,
  \Assign{c_n}{M'_1}} \in \Assign{P}{M'_1}$۔
\item جے !!a{predicate} $P$،
  $\Lang{L}_0$ وچ اے، تاں
  $\Assign{P}{M} =
  \Assign{P}{M'_2} = h(\Assign{P}{M'_1})$۔
\item $\Lang{L}'_1 \cup \Lang{L}'_2$
  دے !!^{function}s، جنہاں وچ
  !!{constant}s وی نیں، ایویں ہی
  سمجھے جاندے نیں۔
\end{enumerate}

آخر وچ !!{formula}s اُتے استقرا
نال دکھایا جاندا اے کہ
$\Struct{M}$، $\Struct{M'_1}$
نال $\Lang{L'_1}$ دے سارے
!!{formula}s بارے، تے
$\Struct{M'_2}$ نال
$\Lang{L'_2}$ دے سارے
!!{formula}s بارے متفق اے؛ پہلے مقابلے وچ ہم شکلی
راہیں ملدیاں متغیر قدر-بندیاں، تے دوجے وچ اکو قدر-بندی
ورتی جاندی اے۔ جملیاں لئی قدر-بندی دا فرق نہیں پیندا۔
خاص طور اُتے
$\Struct{M} \Entails \Gamma^* \cup \Delta^*$،
جس توں $\Struct{M}
\Entails !A$ تے
$\Struct{M} \Entails \lnot!B$،
تے $\not\Entails !A \lif !B$۔
ایس نال کریگ دے درمیانی جملے دے
قضیے دا ثبوت مکمل ہوندا اے۔
\end{proof}
\endgroup
% /OLP-0201

% OLP-0202 content/model-theory/interpolation/definability.tex
\begingroup
\def\olfilename{definability}
\expandafter\def\csname ol@sectioncs\endcsname{section}

\olfileid{mod}{int}{def}

\olsection{قابلِ تعریف ہون دا قضیہ}

درمیانی جملے دے قضیے دا اک اہم اطلاق
بیتھ دا قابلِ تعریف ہون دا قضیہ اے۔
اک $n$-جگہی تعلق~$R$ متعین کرن لئی
اسیں ایسا !!a{formula}~$!C$ دے سکدے
آں جس وچ $n$ آزاد !!{variable}s ہون
تے جس وچ~$R$ نہ آوے۔ ایہ
~$R$ دی \emph{صریح} تعریف
~$!C$ راہیں ہووے گی۔ فیر اسیں ایہ وی آکھ
سکدے آں کہ اک نظریہ~$\Sigma(P)$،
ایسی !!{language} وچ جس وچ
$n$-جگہی !!{predicate}~$P$ ہووے،
~$P$ نوں صریح طور اُتے متعین کردا
اے جے ایس وچ اک باقاعدہ صریح تعریف
موجود ہووے (یا گھٹ توں گھٹ اوہ
ایس توں منطقی طور اُتے نکلدی ہووے)،
یعنی
\[
\Sigma(P) \Entails \lforall[x_1][\dots
  \lforall[x_n][(\Atom{P}{x_1,\dots, x_n} \liff !C(x_1, \dots,
    x_n))]].
\]
پر صریح تعریف، تعلق نوں پوری طرح
متعین کرن دے مفہوم وچ، تعریف دا صرف
اک طریقہ اے۔ کوئی نظریہ ایسا وی ہو
سکدا اے کہ ہر نمونے وچ باقی
!!{language} دی تعبیر،~$P$ دی تعبیر
مقرر کر دیوے۔ قابلِ تعریف ہون دا
قضیہ آکھدا اے کہ جدوں کوئی نظریہ
~$P$ دی تعبیر ایس طرح مقرر کرے،
یعنی~$P$ نوں \emph{ضمنی طور اُتے
متعین کرے}، تاں اوہ اوہنوں صریح طور
اُتے وی متعین کردا اے۔

\begin{defn}
فرض کرو $\Lang{L}$ ایسی !!a{language}
اے جس وچ !!{predicate}~$P$ نہیں۔
$\Sigma(P)$، $\Lang{L}
\cup \{P\}$ دے !!{sentence}s دا سیٹ،
~$P$ نوں
\emph{صریح طور اُتے متعین کردا} اے
اوہدوں تے صرف اوہدوں جدوں
کوئی !!a{formula}~$!C(x_1, \dots, x_n)$،
$\Lang{L}$ دا، ایسا ہووے
کہ
\[
\Sigma(P) \Entails \lforall[x_1][\dots
  \lforall[x_n][(\Atom{P}{x_1,\dots, x_n} \liff !C(x_1, \dots,
    x_n))]].
\]
\end{defn}

\begin{defn}
فرض کرو $\Lang{L}$ ایسی !!a{language}
اے جس وچ !!{predicate}s~$P$ تے~$P'$
نہیں۔ $\Sigma(P)$،
$\Lang{L} \cup \{P\}$ دے
!!{sentence}s دا سیٹ،
$P$ نوں \emph{ضمنی طور اُتے
متعین کردا} اے اوہدوں تے صرف
اوہدوں جدوں
\[
\Sigma(P) \cup \Sigma(P') \Entails \lforall[x_1][\dots
  \lforall[x_n][(\Atom{P}{x_1,\dots, x_n} \liff \Atom{P'}{x_1,\dots,
      x_n})]],
\]
جتھے $\Sigma(P')$، $P$ دی جگہ
$P'$ ہر تھاں اکّو طرح رکھ کے
$\Sigma(P)$ توں ملدا اے۔
\end{defn}

دوجے لفظاں وچ، کسے وی نمونے
$\Struct{M}$ تے $R, R' \subseteq
\Domain{M}^n$ لئی، جے
$\Expan{M}{R} \Entails \Sigma(P)$
تے $\Expan{M}{R'} \Entails \Sigma(P')$
دوہاں ہون، تاں $R=R'$؛ جتھے
$\Expan{M}{R}$ اوہ !!{structure}
~$\Struct{M'}$ اے جو
$\Lang{L}$ نوں
$\Lang{L} \cup \{P\}$ تک ودھا کے
ملدی اے تے جس وچ
$\Assign{P}{M'} = R$؛ تے
$\Expan{M}{R'}$ لئی وی ایویں؛ دوجی توسیع وچ
نویں ناں والا محمول دوجے تعلق نال تعبیر ہوندا اے۔

\begin{thm}[بیتھ دا قابلِ تعریف ہون دا قضیہ]
$\Sigma(P)$، $\Lang{L}
  \cup\{P\}$-!!{sentence}s دا سیٹ،
$P$ نوں ضمنی طور
اُتے متعین کردا اے اوہدوں تے صرف
اوہدوں جدوں $\Sigma(P)$،
$P$ نوں صریح طور اُتے متعین کرے۔
\end{thm}

\begin{proof}
جے $\Sigma(P)$، $P$ نوں صریح
طور اُتے متعین کرے، تاں دوہاں
\begin{align*}
  \Sigma(P) & \Entails & \lforall[x_1][\dots \lforall[x_n]
    [(\Atom{P}{x_1,\dots, x_n} \liff !C(x_1,\dots,x_n))]]\\
  \Sigma(P') & \Entails & \lforall[x_1][\dots \lforall[x_n]
    [(\Atom{P'}{x_1,\dots, x_n} \liff !C(x_1,\dots,x_n))]]
\end{align*}
ہون گے، تے نتیجہ نکل آؤندا اے۔
الٹے رخ لئی فرض کرو کہ $\Sigma(P)$،
$P$ نوں ضمنی طور اُتے متعین کردا
اے۔ پہلے اسیں نویں !!{constant}s
$c_1$، \dots،~$c_n$،
$\Lang{L}$ وچ شامل کردے آں۔ تاں
\[
\Sigma(P) \cup \Sigma(P') \Entails
\Atom{P}{c_1, \dots, c_n} \to  \Atom{P'}{c_1, \dots, c_n}.
\]
محدود گواہی دے قضیے نال محدود سیٹ
$\Delta_0 \subseteq \Sigma(P)$
تے $\Delta_1 \subseteq \Sigma(P')$
ایسے ملدے نیں کہ
\[
\Delta_0 \cup \Delta_1 \Entails
\Atom{P}{c_1, \dots, c_n} \to \Atom{P'}{c_1, \dots, c_n}.
\]
$!D(P)$ نوں اوہناں سارے
!!{sentence}s $!A(P)$ دا منطقی جوڑ من لو
جنہاں لئی یا $!A(P) \in \Delta_0$
ہووے یا $!A(P') \in \Delta_1$؛
تے $!D(P')$ نوں اوہناں سارے
!!{sentence}s $!A(P')$ دا منطقی جوڑ من لو
جنہاں لئی یا $!A(P) \in \Delta_0$
ہووے یا $!A(P') \in \Delta_1$۔
تاں $!D(P) \land
!D(P') \Entails \Atom{P}{c_1, \dots, c_n} \to P'c_1\dots c_n$۔
ایس نوں ایویں مُڑ ترتیب دے
سکدے آں کہ دتے خاص !!{predicate} دے اصل تے نویں ناں،
$\Entails$ دے آپو اپنے پاسے آ جان؛ باقی زبان مشترک رہ سکدی اے:
\[
!D(P) \land \Atom{P}{c_1, \dots, c_n} \Entails
!D(P') \to \Atom{P'}{c_1, \dots, c_n}.
\]
کریگ دے درمیانی جملے دے قضیے نال
اک !!a{sentence} $!C(c_1,\dots, c_n)$
اے جس وچ $P$ یا $P'$ نہیں، تے
جس لئی:
\begin{align*}
  \formula D(P) \land \Atom{P}{c_1, \dots, c_n} & \Entails !C(c_1,\dots, c_n); \\
  \formula C(c_1,\dots, c_n) & \Entails !D(P') \to \Atom{P'}{c_1, \dots, c_n}.
\end{align*}
ایہناں وچوں پہلے منطقی نتیجے
توں ساڈے کول
$!D(P) \Entails
\Atom{P}{c_1,\dots, c_n} \lif !C(c_1,\dots, c_n)$
آؤندا اے۔ تے دوجے توں، کیوں جے
اک $\Lang{L} \cup \{P\}$-نمونہ
$\Expan{M}{R}
\Entails !A(P)$ اوہدوں تے صرف
اوہدوں جدوں متعلقہ
$\Lang{L} \cup
\{P'\}$-نمونہ
$\Expan{M}{R} \models !A(P')$
ہووے (باقی زبان، دائرہ تے تعلق دی تعبیر اکو اے، صرف
محمول دا ناں بدلدا اے)، اسیں
$!C(c_1,\dots, c_n) \Entails !D(P) \lif \Atom{P}{c_1,\dots, c_n}$
لیندے آں، جس توں:
\[
!D(P) \Entails !C(c_1,\dots,c_n) \to \Atom{P}{c_1,\dots, c_n}.
\]
دوہاں نوں ملا کے
$!D(P) \Entails \Atom{P}{c_1,\dots, c_n}
\liff !C(c_1, \dots, c_n)$،
تے یک رُخیت تے عمومیت نال ایہ وی۔ ایتھے نویں مستقلاں
دی ہر آمد نوں متعلقہ آزاد متغیر نال بدل کے اوہناں دے
ناواں توں پاک فارمولا بناؤ؛ پہلاں پابند متغیراں دے ناں
بدلو جے لوڑ ہووے، تاں جو نواں آزاد متغیر کسے مسوّر دے
دائرۂ اثر وچ نہ پھس جائے:
\[
\Sigma(P) \Entails
\lforall[x_1][\dots\lforall[x_n][(\Atom{P}{x_1,\dots, x_n} \liff
    !C(x_1,\dots, x_n))]].
\]
\end{proof}
\endgroup
% /OLP-0202

\OLEndChapterHook
\endgroup
% /OLP-0198

% OLP-0203 content/model-theory/lindstrom/lindstrom.tex
\begingroup
\def\olfilename{lindstrom}
\expandafter\def\csname ol@sectioncs\endcsname{section}

\olchapter{mod}{lin}{لنڈسٹروم (Lindstr\"om) دا قضیہ}

% OLP-0204 content/model-theory/lindstrom/introduction.tex
\begingroup
\def\olfilename{introduction}
\expandafter\def\csname ol@sectioncs\endcsname{section}

\olfileid{mod}{lin}{int}
\section{تعارف}

ایس باب وچ ساڈا مقصد Lindstr\"om دی اوہ خصوصیت بندی ثابت کرنا اے جس
دے مطابق پہلے درجے دی منطق زیادہ توں زیادہ وسیع منطق اے جس لئی (کجھ ہور
شرطاں دے تحت) محدود گواہی دا قضیہ تے نزولی L\"owenheim--Skolem دا
قضیہ قائم رہندے نیں
(\olref[fol][com][com]{thm:compactness} تے
\olref[fol][com][dls]{thm:downward-ls})۔ پہلے سانوں منطقاں دے
اُس عمومی طبقے دی ہور وسیع خصوصیت بندی چاہیدی اے جس اُتے ایہ
قضیہ لاگو ہوندا اے۔ اسیں اپنے آپ نوں \emph{تعلقی} زباناں تک
محدود رکھاں گے، یعنی ایہو جیہیاں زباناں جنہاں وچ صرف
!!{predicate}s تے انفرادی مستقلاں ہوندے نیں، پر !!{function}s
نہیں ہوندے۔
\endgroup
% /OLP-0204

% OLP-0205 content/model-theory/lindstrom/abstract-logics.tex
\begingroup
\def\olfilename{abstract-logics}
\expandafter\def\csname ol@sectioncs\endcsname{section}

\olfileid{mod}{lin}{alg}
\olsection{تجریدی منطقاں}

\begin{defn}
اک \emph{تجریدی منطق} اک جوڑا $\tuple{L, \models_L}$ اے؛ اِتھے $L$
اوہ فنکشن اے جو ہر !!{language}~$\Lang{L}$ نوں
!!{sentence}s دا سیٹ $L(\Lang{L})$ دیندا اے، تے $\models_L$
اُس !!{language}~$\Lang{L}$ دیاں !!{structure}s تے
$L(\Lang{L})$ دے !!{element}s وچکار اک رشتہ اے۔ خاص طور اُتے
$\tuple{F, \models}$ عام پہلے درجے دی منطق اے: $F$ اوہ فنکشن اے
جو !!{language}~$\Lang{L}$ نوں $\Lang{L}$ دیاں غیر منطقی علامتاں توں بنے
پہلے درجے دے !!{sentence}s دا سیٹ دیندا اے، تے $\models$
پہلے درجے دی منطق دا صدقی رشتہ اے۔
\end{defn}

دھیان دیو کہ کسے مقرر !!{language} لئی اسیں ہن وی
!!{structure} دا اوہو مفہوم ورتدے آں جو پہلے درجے دی منطق وچ
سی۔ البتہ اسیں ایہ پہلے توں نہیں مندے کہ !!{sentence}s عام
طریقے نال $\Lang{L}$ دیاں بنیادی علامتاں توں بنے ہون، نہ ایہ
کہ رشتہ $\models_L$ پہلے درجے دی منطق وانگ تکراری تعریف نال
مقرر ہوندا ہووے۔ مثلاً ایہ تعریف پوری طرح عمومی اے؛ ایس دا
مقصد $\tuple{L,\models_L}$ وچ ایسے !!{sentence}s نوں وی شامل
کرنا اے جنہاں وچ لامتناہی لمبے
عطف یا فصل، یا $\lexists$ تے $\lforall$ توں ہٹ کے مسوّر
(مثلاً «لامتناہی طور تے بہت سارے~$x$ ایس طرح دے نیں کہ \dots»)،
یا شاید مسوّراں دے لامتناہی لمبے ابتدائی سلسلے ہون۔ ایہ گل
صاف رکھن لئی کہ $L(\Lang{L})$ دے «!!{sentence}s» لازماً عام
پہلے درجے دے !!{sentence}s نہیں، ایس باب وچ اسیں اوہناں لئی
!!{variable}s $!E$, $!F$,~\dots ورتاں گے، تے $!A$, $!B$,~\dots
صرف عام پہلے درجے دے !!{formula}s لئی رکھاں گے۔

\begin{defn}
$\Mod(L){!E}$ توں مراد طبقہ $\Setabs{\Struct{M}}{\Struct{M}
  \models_L !E}$ لوو۔ جے !!{language} صراحت نال دسنی ہووے،
تاں اسیں $\Mod[L](L){!E}$ لکھدے آں۔ دو
!!{structure}s $\Struct{M}$ تے $\Struct{N}$،
$\Lang{L}$ لئی،
$\tuple{L, \models_L}$ وچ \emph{جملہ وار ہم ارز} ہوندیاں نیں،
جِنہاں نوں $\Struct{M} \elemequiv[L] \Struct{N}$ لکھدے آں،
جے $L(\Lang{L})$ دے اوہو !!{sentence}s دوہاں وچ سچ ہون۔
\end{defn}

\begin{defn}
اک تجریدی منطق $\tuple{L,\models_L}$،
!!{language} $\Lang{L}$ لئی \emph{معمول مطابق (normal)} اے جے ایہ
ہیٹھاں لکھیاں خاصیتاں پوری کرے:
\begin{enumerate}
\item (\emph{$L$-یک رُخیت}) !!{language}s $\Lang{L}$ تے
  $\Lang{L'}$ لئی، جے $\Lang{L} \subseteq \Lang{L'}$، تاں
  $L(\Lang{L}) \subseteq L(\Lang{L'})$۔
\item (\emph{توسیع دی خاصیت}) ہر $!E \in L(\Lang{L})$ لئی
  اک \emph{محدود} ذیلی سیٹ $\Lang{L'}$ موجود اے، جو
  $\Lang{L}$ دا اے،
  ایس طرح کہ رشتہ $\Struct{M} \models_L !E$ صرف
  $\Struct{M}$ دے $\Lang{L'}$ تک گھٹائے روپ اُتے منحصر ہووے؛
  یعنی جے $\Struct{M}$ تے $\Struct{N}$ دا $\Lang{L'}$ تک
  گھٹایا روپ اکّو ہووے، تاں $\Struct{M}
  \models_L !E$ جے تے صرف جے $\Struct{N} \models_L !E$۔
\item (\emph{ہم شکلی دی خاصیت}) جے $\Struct{M} \models_L !E$
  تے $\Struct{M} \simeq \Struct{N}$، تاں
  $\Struct{N} \models_L !E$ وی۔
\item (\emph{نواں نام رکھن دی خاصیت}) رشتہ $\models_L$ نام
  بدلنے نال محفوظ رہندا اے: جے !!{language} $\Lang{L}'$
  $\Lang{L}$ دے ہر نشان $P$ دی تھاں اوہو تعدادِ دلائل
  والا نشان $P'$ تے ہر مستقل $c$ دی تھاں وکھرا مستقل $c'$
  رکھ کے بنے، تاں ہر !!{structure}~$\Struct{M}$ تے
  !!{sentence}~$!E$ لئی $\Struct{M}
  \models_L !E$ جے تے صرف جے $\Struct{M}' \models_L !E'$،
  جتھے $\Struct{M}'$ اک $\Lang{L}'$-!!{structure} اے جو
  $\Lang{L}$ دی اصل ساخت دے اکو دائرے تے تعبیراں نال بن دی اے؛
  صرف علامتاں دے ناں بدل دے نیں، تے سارے نویں ناں آپس وچ وکھرے نیں۔
  ایتھے $!E' \in L(\Lang{L}')$۔
\item (\emph{بولین خاصیت}) تجریدی منطق $\tuple{L,
  \models_L}$ بولین رابطیاں دے تحت ایس معنی وچ بند اے
  کہ ہر $!E \in L(\Lang{L})$ لئی کوئی~$!F \in
  L(\Lang{L})$ موجود اے جس لئی $\Struct{M} \models_L !F$
  جے تے صرف جے $\Struct{M} \not\models_L !E$، تے ہر
  $!E$ تے $!F$ لئی کوئی $!G$ موجود اے جس لئی
  $\Mod(L){!G} = \Mod(L){!E} \cap
  \Mod(L){!F}$۔ ایس طرح بند جوہری !!{formula}s تے ہور
  رابطیاں لئی وی۔
\item (\emph{مسوّر دی خاصیت}) ہر مستقل $c$ جو $\Lang{L}$ وچ اے
  تے $!E \in L(\Lang{L})$ لئی کوئی $!F \in L(\Lang{L'})$
  موجود اے ایس طرح کہ
  \[
  \Mod[L'](L){!F} = \Setabs{\Struct{M}}{\Expan{M}{a} \in
  \Mod[L](L){!E} \text{ کسے } a \in \Domain{M}\text{ لئی}},
  \]
  جتھے $\Lang{L'} = \Lang{L} \setminus \{c\}$، تے
  $\Expan{M}{a}$، $\Struct{M}$ دی $\Lang{L}$ تک اوہ توسیع
  اے جس وچ $a$، $c$ نوں دتی قدر اے۔
\item (\emph{نسبتی بنانے دی خاصیت}) جے !!a{sentence} $!E \in
  L(\Lang{L})$ تے نشان $R$, $c_1$, \dots, $c_n$،
  $\Lang{L}$ وچ نہ ہون، تاں کوئی !!a{sentence} $!F \in
  L(\Lang{L} \cup \{R,c_1,\ldots,c_n\})$ موجود اے، جس نوں
  $!E$ دا $\Atom{R}{x, c_1, \dots c_n}$ تک
  \emph{نسبتی روپ} آکھدے نیں، ایس طرح کہ ہر
  !!{structure}~$\Struct{M}$ لئی، جدوں منتخب دائرہ غیر خالی ہووے
  تے اصل زبان دے ہر مستقل دی قدر ایس وچ ہووے۔ ایس باب دیاں
  تعلقی زباناں لئی ایہ شرطاں اوہ دائرہ ذیلی ساخت بناون لئی لازم نیں:
  \[
  \Expan{M}{X, b_1, \ldots, b_n} \models_L !F \text{ اوہدوں تے
    صرف اوہدوں جدوں } \Struct{N} \models_L !E,
  \]
  جتھے $\Struct{N}$، $\Struct{M}$ دی اوہ ذیلی ساخت اے
  جس دا !!{domain}
  $\Domain{N} = \Setabs{a\in \Domain{M}}{\tuple{a, b_1, \dots, b_n} \in X}$
  اے (ویکھو \olref[bas][sub]{rem:substructure})، تے
  $\Expan{M}{X, b_1, \ldots,
    b_n}$، $\Struct{M}$ دی اوہ توسیع اے جو بالترتیب
  $R$, $c_1$, \dots, $c_n$ دی تعبیر
  $X$, $b_1,$ \dots, $b_n$ نال کردی اے (جتھے
  $X \subseteq M^{n+1}$)۔
\end{enumerate}
\end{defn}
 
\begin{defn}
دو تجریدی منطقاں $\tuple{L_1, \models_{L_1}}$ تے
$\tuple{L_2, \models_{L_2}}$ دتیاں ہون۔ اسیں آکھدے آں کہ
دوجی منطق پہلی دے \emph{گھٹ توں گھٹ برابر اظہار جوگی} اے،
جس نوں $\tuple{L_1, \models_{L_1}}
\leq \tuple{L_2, \models_{L_2}}$ لکھدے آں، جے ہر
!!{language} $\Lang{L}$ تے ہر !!{sentence}
$!E \in L_1(\Lang{L})$ لئی کوئی !!a{sentence} $!F
\in L_2(\Lang{L})$ موجود ہووے جس لئی
$\Mod[L](L_1){!E} =
\Mod[L](L_2){!F}$۔ منطقاں
$\tuple{L_1, \models_{L_1}}$ تے
$\tuple{L_2, \models_{L_2}}$ \emph{ہم ارز} نیں جے
$\tuple{L_1,
  \models_{L_1}} \leq \tuple{L_2, \models_{L_2}}$ تے
$\tuple{L_2,
  \models_{L_2}} \leq \tuple{L_1, \models_{L_1}}$۔
\end{defn}

\begin{rem}
  پہلے درجے دی منطق، یعنی تجریدی منطق $\tuple{F, \models}$،
  معمول مطابق اے۔ دراصل اُتّے دتیاں خاصیتاں پہلے درجے دی
  منطق لئی زیادہ تر سیدھیاں نیں۔ صرف ایہ یاد رکھو کہ توسیع
  دی خاصیت غیر متعلقہ علامتاں توں آزادی
  (Extensionality/Relevance) اُتے آ ٹکدی اے، تے !!a{sentence} $!E$
  دا $\Atom{R}{x, c_1, \dots, c_n}$ تک نسبتی روپ لبھن لئی
  پہلاں وجودی مسوّراں نوں نفی تے عمومی مسوّر دے مساوی روپ وچ لکھو؛ فیر
  ہر !!{subformula} $\lforall[x][!F]$ دی تھاں
  $\lforall[x][(\Atom{R}{x, c_1, \dots, c_n} \lif \ !F)]$
  رکھ کے لبھدا اے۔ ہور ایہ کہ جے
  $\tuple{L, \models_L}$ معمول مطابق ہووے، تاں
  $\tuple{F, \models} \leq \tuple{L, \models_{L}}$،
  جیویں پہلے درجے دے !!{formula}s اُتے استقرا نال
  دکھایا جا سکدا اے۔ ایس لئی عمومیت گنوائے بغیر اسیں
  فرض کر سکدے آں کہ ہر پہلے درجے دا !!{sentence} ہر
  معمول مطابق منطق وچ شامل اے۔
\end{rem}
\endgroup
% /OLP-0205

% OLP-0206 content/model-theory/lindstrom/ls-property.tex
\begingroup
\def\olfilename{ls-property}
\expandafter\def\csname ol@sectioncs\endcsname{section}

\olfileid{mod}{lin}{lsp}

\olsection{محدود گواہی تے L\"owenheim--Skolem دیاں خاصیتاں}

ہن اسیں تجریدی منطقاں لئی محدود گواہی تے
L\"owenheim--Skolem دیاں سیدھیاں توسیعاں دیندے آں۔

\begin{defn}
اک تجریدی منطق $\tuple{L, \models_L}$ وچ \emph{محدود گواہی
  دی خاصیت} اے جے ہر سیٹ $\Gamma$، جو
$L(\Lang{L})$ دے !!{sentence}s دا ہووے، اوہدوں قابلِ اطمینان ہووے جدوں اوہدا
ہر محدود $\Gamma_0 \subseteq \Gamma$ قابلِ اطمینان ہووے۔
\end{defn}

\begin{defn}
$\tuple{L, \models_L}$ وچ \emph{نزولی L\"owenheim--Skolem
  دی خاصیت} اے جے ہر قابلِ اطمینان $\Gamma$ دا
!!a{enumerable} نمونہ موجود ہووے، یعنی محدود یا گنتی وار
لامتناہی نمونہ۔
\end{defn}


\olref[bas][pis]{defn:partialisom} وچ جزوی ہم شکلی دا تصور خالص
«الجبری» اے، یعنی اوہ زبان دے !!{sentence}s نوں
نہیں، بلکہ صرف !!{structure}s دی !!{language}~$\Lang{L}$
وچ دتیاں غیر منطقی علامتاں نوں ویکھ کے بیان ہوندا اے؛ ایس لئی ایہ
تجریدی منطقاں اُتے وی لاگو ہوندا اے۔ پہلے درجے دی منطق لئی
\olref[bas][pis]{thm:p-isom2} توں سانوں پتہ اے کہ جے دو
!!{structure}s جزوی طور اُتے ہم شکل ہون تاں اوہ جملہ وار
ہم ارز نیں۔ اوہ ثبوت تجریدی منطقاں تک نہیں پہنچدا، کیوں جے
من مانے $!E \in L(\Lang{L})$ لئی !!{formula}s اُتے
استقرا دستیاب ہونا ضروری نہیں؛ پر جے L\"owenheim--Skolem دی
خاصیت ہووے، تاں قضیہ پھر وی درست اے۔

\begin{thm}
\ollabel{thm:abstract-p-isom}
فرض کرو $\tuple{L, \models_L}$ اک معمول مطابق منطق اے جس وچ
L\"owenheim--Skolem دی خاصیت اے۔ تاں ہر دو !!{structure}s
جو جزوی طور اُتے ہم شکل ہون،
$\tuple{L,
  \models_L}$ وچ جملہ وار ہم ارز نیں۔
\end{thm}

\begin{proof}
فرض کرو $\Struct{M} \simeq_p \Struct{N}$، پر کسے $!E$ لئی
$\Struct{M} \models_L !E$ جدوں کہ
$\Struct{N} \not\models_L !E$۔ ہم شکلی دی خاصیت توں اسیں
ہم شکل نقلاں ورت سکدے آں۔ پر مستقل نشانیاں ہون تاں
$\Domain{M}$ تے $\Domain{N}$ نوں پوری طرح وکھرا نہیں رکھنا:
پرانیاں مستقل نشانیاں دیاں متناظر قدراں نوں اکّو عنصر بنا کے پہچانو،
تے باقی عناصر نوں وکھرے نشان دیو۔ اوہ سانجھا سیٹ $K$ آکھو۔
اگے-پیچھے خاصیت نال ہر محدود لسٹ دیاں نامی قدراں کسے جزوی نقشے
وچ آ سکدیاں نیں؛ سو دوہاں ساختاں دے مستقل برابر ہون دے نمونے
اکّو نیں، تے سانجھیاں نامی قدراں اُتے ہر محمول دی تعبیر وی اکّو اے۔
ایس لئی ایہ پہچان اک درست سانجھا حصہ بناندی اے۔
اسیں
$!E \in
L(\Lang{L})$ لئی اصل !!{language}~$\Lang{L}$ برقرار رکھاں گے؛
توسیع دی خاصیت صرف محدود علامتاں اُتے صدق دی وابستگی دیندی اے،
ایس توں جملے دی محدود زبان وچ آمد اپنے آپ نہیں نکلدی۔ ہِٹھاں
پہلے درجے دے جملیاں دا سیٹ ورتاں گے، جس لئی نزولی خاصیت کافی اے۔
$\mathcal{I}$، $\Struct{M}$ تے $\Struct{N}$ وچکار جزوی
ہم شکلیاں دا اک غیر خالی، اگے-پیچھے خاصیت والا سیٹ ہووے۔
عمومیت گنوائے بغیر ایہ وی فرض
کرو کہ جے $p
\in \mathcal{I}$ تے $q \subseteq p$، تاں
$q \in \mathcal{I}$ وی۔ ہر ایسے نقشے نوں تمام سانجھیاں نامی
قدراں اُتے شناختی نقشے نال ملا کے $j_p=p\cup\operatorname{id}_K$
لکھو۔ کسے نامی قدر نوں اصل نقشے وچ شامل کرن دی توسیع اوہنوں
اوہدی متناظر قدر اُتے ای بھیجدی اے؛ یک بہ یک ہون کارن کوئی ہور
عنصر اوس قدر اُتے نہیں جا سکدا۔ سو ایہ ملاپ وی جزوی یک بہ یک
نقشہ اے تے سارے محمول محفوظ رکھدا اے:
کسے محمول نوں پرکھن لئی صرف محدود ٹپل دی لوڑ اے، تے اوس ٹپل
دیاں نامی قدراں نوں اصل اگے-پیچھے سیٹ دے نقشے وچ شامل کیتا جا سکدا اے۔
نویں نقشے محدود ہون ضروری نہیں؛ پر اوہناں دا غیر خالی خاندان ہن وی
دونوں سمتاں وچ ہر نواں عنصر شامل کر سکدا اے۔

$\Domain{M}^{<\omega}$، $\Domain{M}$ دے !!{element}s دیاں
محدود ترتیب وار لڑیّاں دا سیٹ اے؛ ساخت بناندیاں لڑیّاں دیاں
وکھرے نشان لائیاں نقلاں ورتو، تاکہ عنصر تے لڑی دے کوڈ آپس وچ
نہ رلن، تے دوہاں پاسیاں دیاں خالی لڑیّاں دے کوڈ وی وکھرے ہون۔ $S$،
$\Domain{M}^{<\omega}$ اُتے تین مقامی رشتہ ہووے جو لڑیّاں
جوڑن نوں ظاہر کردا اے؛ یعنی جے $\mathbf{a}, \mathbf{b},
\mathbf{c} \in \Domain{M}^{<\omega}$، تاں
$S(\mathbf{a}, \mathbf{b},
\mathbf{c})$ اوہدوں تے صرف اوہدوں درست اے جدوں
$\mathbf{c}$، $\mathbf{a}$ تے $\mathbf{b}$ نوں جوڑ
کے بنی ہووے۔ تے $T$ اوہ تین مقامی رشتہ ہووے جس لئی
$T(\mathbf{a}, b, \mathbf{c})$، $b \in M$ تے
$\mathbf{a}, \mathbf{c} \in \Domain{M}^{<\omega}$ لئی
اوہدوں تے صرف اوہدوں درست اے جدوں $\mathbf{a} =
a_1, \dots a_n$ تے $\mathbf{c} = a_1, \dots a_n, b$۔
نویں تین مقامی !!{predicate}s $P$ تے $Q$ چنو، تے
!!{structure} $\Struct{M}^*$ بناؤ جس دا دائرہ
$\Domain{M} \cup \Domain{M}^{<\omega}$ اے،
جس وچ $\Struct{M}$ ذیلی ساخت اے، تے جس وچ $P$ تے $Q$
دی تعبیر بالترتیب جوڑن دے رشتے $S$ تے $T$ نیں
(اس لئی $\Struct{M}^*$،
!!{language} $\Lang{L} \cup \{ P, Q\}$ وچ اے)۔

$\Domain{N}^{<\omega}$، $S'$، $T'$، $P'$، $Q'$ تے
$\Struct{N}^*$ اوہو جیہے طریقے نال مقرر کرو۔ چونکہ مفروضے
مطابق $\Struct{M} \simeq_p
\Struct{N}$، ایس لئی اک رشتہ $I$،
$\Domain{M}^{<\omega}$ تے
$\Domain{N}^{<\omega}$ وچکار اے جس لئی
$I(\mathbf{a}, \mathbf{b})$ اوہدوں تے صرف اوہدوں درست اے
جدوں برابر لمبائی دیاں لڑیّاں $\mathbf{a}$ تے $\mathbf{b}$
دے متناظر عناصر نوں کسے نقشے $p\in\mathcal{I}$ نال ملایا جا سکدا ہووے؛
\olref[bas][pis]{defn:partialisom} دی اگے-پیچھے شرط
پورے خاندان دی شرط اے، اکیلے ٹپل دی نہیں۔ ہن $\Struct{C}$
نوں اوہ !!{structure} آکھو جس دا
!!{domain}، $\Struct{M}^*$ تے $\Struct{N}^*$ دے
!!{domain}s دا اتحاد، نال ہر نقشے $j_p$ دا اک وکھرا کوڈ، اے۔
اصل زبان دے محمول دوہاں پاسیاں اُتے اصل تعبیر رکھن، تے باقی
مخلوط ٹپلاں اُتے جھوٹے رکھے جا سکدے نیں؛ سانجھا حصہ $K$ اُتے
تعبیراں پہلے ای رلدیاں نیں۔ اپنے اپنے متعلقہ زبان والے گھٹائے روپ
وچ $\Struct{M}^*$ تے $\Struct{N}^*$ ذیلی !!{structure}s نیں۔
ایس ساخت دی
!!{language} وچ اک وادھو دو مقامی !!{predicate}~$R$ اے
جس دی تعبیر رشتہ~$I$ اے، نالے اوہ !!{predicate}s
جو !!{domain}s $\Domain{M}^*$ تے~$\Domain{N}^*$
نوں ظاہر کردے نیں۔ ہور نویں یک مقامی محمول $U$ تے $V$
صرف اصل دو دائرے چنّن، یک مقامی محمول $H$ نقشیاں دے کوڈ چنّے،
تے تین مقامی محمول $G$ لئی $G(h,a,b)$ دا مطلب ہووے کہ
کوڈ $h$ والا نقشہ عنصر $a$ نوں $b$ نال ملاندا اے۔

\begin{figure}[h]
  \centering
  \begin{tikzpicture}[node distance=2cm, auto, thick, >=stealth']
    \draw [rounded corners] (0,0) -- (8,0) -- (8,4) -- (0,4) --  cycle;
    \draw (2,2) circle (0.5cm);
    \draw (2,2) circle (1.25cm);
    \draw (6,2) circle (0.5cm);
    \draw (6,2) circle (1.25cm);
    \path node at (0.75,3.5) {\large $\Struct{C}$};
    \path node at (2,2) {\large $\Struct{M}$};
    \path node at (6,2) {\large $\Struct{N}$};
    \path node at (3.5,1) {\large $\Struct{M}^*$};
    \path node at (7.5,1) {\large $\Struct{N}^*$};
    \node (Idom) at (2.8,2) {};
    \node (Irng) at (5.2,2) {};
    \draw[<->, bend left] (Idom) to node {\large $I$} (Irng) ;
  \end{tikzpicture} 
  \caption{!!{structure}~$\Struct{C}$ تے اوہدے اندرلی
    جزوی ہم شکلی دا خاکہ۔ وکھرے دائرے مستقل نشانیاں توں خالی حال
    دکھاندے نیں؛ مستقل ہون تاں دوہاں پاسیاں دیاں سانجھیاں نامی قدراں
    نوں اکّو عنصر سمجھنا اے۔ نقشیاں دے وادھو کوڈ بیرلے دائرے وچ نیں،
    پر ایس خاکے وچ الگ نہیں دکھائے گئے۔}
\end{figure}

بنیادی مشاہدہ ایہ اے کہ !!{structure}~$\Struct{C}$ دی
!!{language} وچ اک \emph{تجریدی} !!{sentence} $!D_1$
موجود اے، جو $\Struct{C}$ وچ سچ اے تے آکھدا اے کہ
$\Struct{M} \models_L !E$ تے
$\Struct{N} \not\models_L !E$ (ایس لئی نسبتی بنانے
دی خاصیت تے بولین بندش درکار نیں)۔ ایہ جملہ پہلے نسبتی روپ
توں بنے گا: پہلا روپ $U$ تک، تے نفی والا دوجا روپ $V$ تک۔
دونوں چنے دائرے غیر خالی نیں تے ہر پرانے مستقل دی قدر دوہاں وچ اے۔
توسیع دی خاصیت وچ صدق صرف اوہ محدود گھٹایا روپ ویکھدا اے؛
لہٰذا جملہ جس زبان وچ دستیاب ہووے، اوہدے ہور نویں علامتاں نال
ودھائے روپ اُتے وی اوہو صدق رہندا اے۔ ایہ صدقی بات اے، جملے نوں
محدود زبان دا قرار دین دی بات نہیں۔ من مانے تجریدی اصل جملے لئی
ایس بنے جملے نوں پہلے درجے دا نہیں کہہ سکدے۔
اصل دے واحد \emph{پہلے درجے دے} !!{sentence} $!D_2$ دی تھاں
پہلے درجے دے جملیاں دا سیٹ $\Delta_2$ ورتو۔ ایس وچ چنے دائرے
تے کوڈ غیر خالی ہون، نقشے جزوی یک بہ یک ہون، تے اگے-پیچھے
توسیع لئی ایہ جملے شامل نیں:
\[
\begin{gathered}
\exists a\,U(a),\qquad \exists b\,V(b),\qquad \exists h\,H(h),\\
\forall h,a,b\,(G(h,a,b)\to(H(h)\land U(a)\land V(b))),\\
\forall h,a,b,b'\,((G(h,a,b)\land G(h,a,b'))\to b=b'),\\
\forall h,a,a',b\,((G(h,a,b)\land G(h,a',b))\to a=a'),\\
\forall h,a\,((H(h)\land U(a))\to
\exists k,b\,(H(k)\land V(b)\land G(k,a,b)\land
\forall u,v\,(G(h,u,v)\to G(k,u,v)))),\\
\forall h,b\,((H(h)\land V(b))\to
\exists k,a\,(H(k)\land U(a)\land G(k,a,b)\land
\forall u,v\,(G(h,u,v)\to G(k,u,v)))).
\end{gathered}
\]
ہر پرانے مستقل $c$ لئی جملے $U(c)\land V(c)$ تے
$\forall h\,(H(h)\to G(h,c,c))$ شامل کرو۔ ہر پرانے
$n$-مقامی محمول $A$ لئی ایہ جملہ شامل کرو:
\[
\forall h,a_1,\ldots,a_n,b_1,\ldots,b_n\,
\bigl((H(h)\land\bigwedge_{i=1}^{n}G(h,a_i,b_i))\to
(A(a_1,\ldots,a_n)\leftrightarrow A(b_1,\ldots,b_n))\bigr).
\]
ایہ صرف پہلے درجے دے جملیاں دا سیٹ اے، زبان محدود ہون دی شرط
نہیں؛ شناختی رشتے دی حفاظت نقشے دی یک بہ یک ہون دی شرط توں آندی اے۔
عام پہلے درجے دے ایہ جملے معمول مطابق تجریدی منطق وچ دستیاب نیں۔
ایہ سارا سیٹ $\Struct{C}$ وچ سچ اے: ہر کوڈ اصل $j_p$ دا اے،
توسیع اصل اگے-پیچھے خاندان دی توسیع توں ملدی اے۔
L\"owenheim--Skolem دی خاصیت جملیاں دے پورے قابلِ اطمینان سیٹ
$\{!D_1\}\cup\Delta_2$ نوں !!a{enumerable}
نمونہ $\Struct{C}_0$ دیندی اے۔ اوہدے $U$ تے $V$ والے دائرے
چن کے اصل زبان دیاں ذیلی ساختاں $\Struct{M}_0$ تے
$\Struct{N}_0$ لوو؛ اوہ غیر خالی تے قابلِ شمار نیں، تے
$\Struct{M}_0 \models_L !E$ تے
$\Struct{N}_0
\not\models_L !E$۔ کوڈاں دے نقشیاں دیاں ساریاں محدود پابندیاں
اک غیر خالی جزوی ہم شکلی خاندان بناندیاں نیں: کسے محدود نقشے
دا اصل کوڈ چن کے اُتّے لکھی توسیع نال نواں عنصر شامل کرو، تے
فیر صرف لوڑ جوگی محدود پابندی لوو۔ محمولاں تے مستقلاں دی حفاظت
درج جملیاں توں آندی اے۔ ایس طرح $\Struct{M}_0 \simeq_p \Struct{N}_0$
دا اصل دلیل والا اگے-پیچھے طریقہ ہن نئی دو ساختاں لئی وی حاصل اے؛
رشتہ~$I$ دی اصل کوڈ بندی نوں نقشیاں دے گراف نال صریح کیتا گیا اے۔
پر !!{enumerable}، جزوی طور اُتے
ہم شکل !!{structure}s دراصل
\olref[bas][pis]{thm:p-isom1} توں ہم شکل ہوندیاں نیں؛
ایہ معمول مطابق تجریدی منطقاں دی ہم شکلی دی خاصیت دے خلاف اے۔
\end{proof}
\endgroup
% /OLP-0206

% OLP-0207 content/model-theory/lindstrom/lindstrom-proof.tex
\begingroup
\def\olfilename{lindstrom-proof}
\expandafter\def\csname ol@sectioncs\endcsname{section}

\olfileid{mod}{lin}{prf}

\olsection{لنڈسٹروم (Lindstr\"om) دا قضیہ}

\begin{lem}
\ollabel{lem:lindstrom}
فرض کرو تجریدی منطق معمول مطابق اے، تے $!E \in L(\Lang{L})$، جتھے $\Lang{L}$ محدود
اے، تے ایہ وی فرض کرو کہ کوئی $n \in \Nat$ موجود اے
ایس طرح کہ ہر دو !!{structure}s $\Struct{M}$ تے~$\Struct{N}$
لئی جے $\Struct{M} \equiv_n
\Struct{N}$ تے $\Struct{M} \models_L !E$ ہووے، تاں
$\Struct{N}
\models_L !E$ وی ہووے۔ تاں $!E$ پہلے درجے دے اک
!!{sentence} دے ہم ارز اے؛ یعنی پہلے درجے دا کوئی $!D$
موجود اے جس لئی $\Mod[L](L){!E} = \Mod[L](L){!D}$؛
ایہ دونوں نمونہ طبقے دتی زبان دیاں ساختاں لئی نیں۔
\end{lem}

\begin{proof} 
$n$ اوہ ہووے جس لئی ہر دو $n$-ہم ارز !!{structure}s
$\Struct{M}$ تے $\Struct{N}$، $!E$ نوں دتی قدر اُتے
متفق ہون۔ \olref[bas][pis]{prop:qr-finite} یاد کرو:
محدود !!{language} وچ مقدار نمائی درجہ~$n$ توں
ودھ نہ ہون والے پہلے درجے دے !!{sentence}s، منطقی
ہم ارزی تک، صرف محدود تعداد وچ نیں۔ ایہ لفظی جملیاں دی محدود
تعداد نہیں: ہر منطقی ہم ارزی جماعت وچوں اک نمائندہ چنو۔ معمول
مطابق منطق وچ، پچھلی وضاحت مطابق، عام پہلے درجے دے جملے اوہناں
دے معمول دے صدق نال شامل سمجھے جا سکدے نیں۔ ہن ہر مقرر
!!{structure}~$\Struct{M}$ لئی $!D_{\Struct{M}}$،
چنے پہلے درجے دے !!{sentence}s~$!A$ دا اوہ عطف
ہووے جو $\Struct{M}$ وچ سچ ہون تے جنہاں لئی
$\QuantRank{!A} \le n$؛ ایہ عطف ہن واقعی محدود اے۔ ایس طرح
$\Struct{N} \models !D_{\Struct{M}}$ جے تے صرف جے
$\Struct{N}
\equiv_n \Struct{M}$۔ فیر $!D = \textstyle\bigvee
\Setabs{!D_{\Struct{M}}}{\Struct{M} \models_L !E}$
رکھو؛ ایس فصل وچ وی ہر حاصل درجہ وار قسم لئی صرف اک
نمائندہ رکھو، سو ایہ واقعی محدود پہلے درجے دا جملہ اے۔ خالی
عطف نوں ہمیشہ سچا جملہ تے خالی فصل نوں ہمیشہ جھوٹا جملہ سمجھو۔

نتیجہ $\Mod[L](L){!E} = \Mod[L](L){!D}$ نکلدا اے۔
واقعی، جے $\Struct{N} \models_L !D$، تاں کسے
$\Struct{M} \models_L
!E$ لئی $\Struct{N} \models !D_{\Struct{M}}$ ہووے گا،
جس توں $\Struct{N} \models_L !E$ وی نکلدا اے
(لمے دے مفروضے نال)۔ الٹے پاسے، جے
$\Struct{N} \models_L !E$، تاں $!D_\Struct{N}$
$!D$ دا اک جزوِ انفصال اے، تے چونکہ
$\Struct{N}
\models !D_\Struct{N}$، ایس لئی
$\Struct{N} \models_L !D$ وی۔
\end{proof}

\begin{thm}[Lindstr\"om دا قضیہ]
  \ollabel{thm:lindstrom} فرض کرو $\tuple{L, \models_L}$ وچ
  معمول مطابق منطق دیاں شرطاں، محدود گواہی تے نزولی
  L\"owenheim--Skolem دیاں خاصیتاں نیں۔ تاں
  $\tuple{L, \models_L} \le \tuple{F, \models}$ (اس لئی
  $\tuple{L, \models_L}$ پہلے درجے دی منطق دے ہم ارز اے)۔
\end{thm}

\begin{proof}
ہر $!E \in L(\Lang{L})$ لوو۔ توسیع دی خاصیت توں اک
محدود !!{language} $\Lang{L_0}\subseteq\Lang{L}$ ملدی اے
جس دے گھٹائے روپ اُتے $!E$ دا صدق منحصر اے۔ ایہ صدقی وابستگی
اے؛ اسیں $!E$ نوں چھوٹی زبان دا جملہ نہیں منّاں گے۔ ہر چھوٹی
زبان دی ساخت نوں اصل زبان تک ودھایا جا سکدا اے، تے اکّو گھٹایا
روپ رکھن والیاں ایس طرح دیاں ساختاں وچ $!E$ دا صدق اکّو اے۔
سو \olref{lem:lindstrom} دے محدود نمائندیاں والے استدلال نوں
ایس صدقی خاصیت اُتے ورتنا کافی اے: کوئی $n\in\Nat$ ایسا ہووے
کہ دو اصل !!{structure}s $\Struct{M}$ تے~$\Struct{N}$ دے
چھوٹے زبان والے روپ $\Struct{M}\equiv_n\Struct{N}$ ہون
تاں $!E$ بارے متفق ہون۔ ایتھے ایہ ہم ارزی صرف انہاں چھوٹے
روپاں لئی لکھی گئی اے۔ محدود درجہ وار قسم دے اوہو عطف تے فصل
اک پہلے درجے دا !!{sentence} بنان گے جو چھوٹی زبان وچ اے
تے اصل زبان وچ وی $!E$ دے ہم ارز اے۔ ایس توں
$\tuple{L, \models_L} \le \tuple{F, \models}$ نکلے گا؛
معمول مطابق منطق وچ پہلے درجے دی شمولیت الٹی سمت تے قوساں وچ
لکھی ہم ارزی وی دیندی اے۔ ہِٹھاں محدود خالص تعلقی زبان لئی
\olref[bas][pis]{thm:b-n-f} مطابق مقدار نمائی درجے دی ہم ارزی نوں
اگے-پیچھے دی محدود الجبری شرط $I_n(\emptyseq,\emptyseq)$
نال ورتاں گے۔

تضاد لئی فرض کرو کہ ہر $n$ لئی اصل !!{structure}s
$\Struct{M}_n$ تے $\Struct{N}_n$ موجود نیں جنہاں دے
چھوٹے روپ مقدار نمائی درجہ $n+1$ تک ہم ارز نیں، پر
$\Struct{M}_n\models_L !E$ تے
$\Struct{N}_n\not\models_L !E$۔ اک نواں مستقل $c_\circ$
چنو۔ پہلی ساخت وچ اک عنصر چن کے، محدود درجے دی ہم ارزی دی
اگے والی شرط نال دوسری وچ ایسا متناظر عنصر چن سکدے آں کہ
انہاں نوں نواں مستقل نام دین پچھوں محدود زبان
$\Lang{L_0}\cup\{c_\circ\}$ دے روپ مقدار نمائی درجہ $n$ تک ہم ارز
ہون۔ اصل زبان دے باقی مستقلاں نوں ایس نامی عنصر دی قدر دیو،
تے باقی محمول خالی رکھو۔ چھوٹا روپ نہیں بدلیا، سو $!E$ دا
مختلف صدق وی نہیں بدلیا۔ ہن اسی محدود ودھی زبان لئی
$I_n(\emptyseq,\emptyseq)$ حاصل اے۔ اوہدیاں بند جوہری
!!{sentence}s صرف محدود تعداد وچ نیں؛ اس لئی کسے بے حد
ذیلی سلسلے وچ سارے جوڑیاں دے اوہو جوہری !!{sentence}s سچ ہون گے۔ جوڑی
دے دونوں پاسے درجہ صفر اُتے وی متفق نیں۔ ایہ ذیلی سلسلہ
ایویں دوبارہ گنو کہ نئی جوڑی نمبر $n$ دا پرانا درجہ کم از کم
$n$ ہووے؛ اصل $I_n$ شرط قائم رہندی اے۔ ہم شکل نقلاں وچ
مستقلاں دیاں متناظر قدراں نوں مشترک نامی حصہ بنا کے پہچانو،
تے ہر جوڑی دے باقی عناصر نوں الگ نشان دیو۔ بند جوہری صدق دا
اتفاق مشترک نامی حصے اُتے برابری تے محمول تعبیر دا اتفاق یقینی
کردا اے؛ سو اتحاد وچ اصل محمول دوہاں پاسیاں نال رل سکدے نیں۔
$\Struct{M}$ اوہ !!{structure} ہووے جو ساریاں $\Struct{M}_n$ ساختاں دا اتحاد اے،
جس وچ انہاں دیاں اصل محمول تعبیراں دا اتحاد اے تے ہور مخلوط
ٹپلاں اُتے محمول جھوٹے نیں؛ ہر $\Struct{M}_n$ اصل زبان
دے روپ وچ ذیلی ساخت اے۔ \olref[lsp]{thm:abstract-p-isom}
دے درست کیتے ثبوت وانگ، $\Struct{M}^*$ دی
!!{domain} $\Domain{M}\cup\Domain{M}^{<\omega}$ وچ
لڑیّاں دے الگ نشان والے کوڈ ورتو۔ ودھی !!{language} وچ
جوڑن دے محمول $P$ تے~$Q$ اوہو نیں۔

اسی طرح $\Struct{N}_n$، $\Struct{N}$ تے
$\Struct{N}^*$ مقرر کرو۔ بڑی !!{structure} نوں
$\Struct{C}$ آکھو، تاکہ اصل پاسے والی ساخت دا ناں نہ رلے۔
اوہدی !!{domain} وچ $\Struct{M}^*$ تے $\Struct{N}^*$
دے !!{domain}s، فطری عدد~$\Nat$ دے وکھرے کوڈ تے اوہناں
دی فطری ترتیب~$\le$، نال جزوی نقشیاں دے الگ کوڈ شامل ہون۔
یک مقامی نواں محمول $Z$ عددی کوڈ چنّے۔ اصل نامی حصہ مشترک
اے؛ لڑیّاں، عددی کوڈ تے نقشہ کوڈ اصل عناصر توں وکھرے نیں۔
اِس دی !!{language} وچ وادھو محمول !!{domain}s
$\Domain{M}$، $\Domain{N}$، $\Domain{M}^{<\omega}$ تے
$\Domain{N}^{<\omega}$ نوں چنّن، تے $R$ تے $S$،
$\Struct{M}_n$ تے $\Struct{N}_n$ دے دائرے تے متعلقہ لڑی
کوڈ ایس معنی وچ چنّن:
\begin{align*}
  \Domain{M_n} & = \Setabs{a \in \Domain{M}}{R(a, n)}; & 
  \Domain{N_n} & = \Setabs{a \in \Domain{N}}{S(a,n)}; \\
  \Domain{M}^{<\omega}_n & = \Setabs{a \in \Domain{M}^{<\omega}}{R(a,n)}; &
  \Domain{N}^{<\omega}_n & = \Setabs{a \in \Domain{N}^{<\omega}}{S(a,n)}. 
\end{align*}
ایس !!{structure} وچ ہن اگے-پیچھے دی باقی گہرائی تے اصل جوڑی دا اشاریہ وکھ رکھنا
لازم اے۔ اصل عنصر دائرے چنّن والے محمول $U_0$ تے $V_0$
آکھو، تے دو مقامی محمول ایویں مقرر کرو:
$U(a,m)\leftrightarrow(U_0(a)\land R(a,m))$ تے
$V(b,m)\leftrightarrow(V_0(b)\land S(b,m))$۔ ہن لڑی
کوڈ انہاں چنے حصیاں وچ نہیں آندے۔ نواں تین مقامی رشتہ $J(d,m,h)$ دا مطلب ہووے:
کوڈ $h$ والا نقشہ جوڑی نمبر $m$ وچ گھٹ توں گھٹ $d$
ہور قدم اگے-پیچھے توسیع برداشت کردا اے۔ تین مقامی محمول
$G(h,a,b)$ اوس نقشے دا گراف اے۔ اصل $I_d$ وچ آنے والے
محدود tuple-نقشیاں نال مشترک نامی حصے دا شناختی نقشہ ملا کے
ایسے گراف بناؤ، جیویں پچھلے ثبوت وچ کیتا۔ محدود درجے دی شرط توں
ہر $m$ لئی $J(m,m,h)$ والا کوئی کوڈ ملدا اے؛ گہرائی صفر
تک ہر توسیع اصل جوڑی $m$ نوں قائم رکھدی اے۔ ایہ کوڈ بندی
اصل لڑیّاں دی تعمیر نوں نہیں بدل دی، پر غلط اکّو اشاریے والے
رشتے دی تھاں اصل جزوی نقشیاں نوں صریح رکھدی اے۔

ہن !!{language}~$\Lang{L}$ نوں $R$، $S$، $J$، $G$،
$Z$ وغیرہ نال ودھاؤ؛ تمام کوڈ بندی دیاں علامتاں اصل زبان
توں نئی چنّیاں نیں۔ نسبتی روپ، بولین بندش تے مسوّر دی
خاصیت نال اک تجریدی !!a{sentence}~$!D$ حاصل اے جو ہر
عدد کوڈ $m$ لئی $\Struct{M}_m\models_L !E$ تے
$\Struct{N}_m\not\models_L !E$ آکھدا اے۔ ایس نوں
بنان لئی عارضی نواں مستقل جوڑی دا اشاریہ بنے؛ پہلے $U$
تے $V$ دے اوس پیرامیٹر والے حصیاں تک $!E$ تے اوہدی نفی
دا نسبتی روپ لوو، فیر «جے اشاریہ عددی اے» دی شرط نال ساریاں
قدراں اُتے عمومی مسوّر لوو۔ ایہ پیرامیٹر ذیلی ساخت اندر ہون
ضروری نہیں؛ ہر پرانے مستقل دی قدر اندر ہونا ضروری اے۔ نال
پہلے درجے دے جملیاں دا سیٹ $\Delta$ ورتو۔ ایہ $Z$ اُتے
$\le$ نوں خطی ترتیب بناندا اے، اک پہلا عددی رکن چنّدا اے،
ہر عدد نوں فوری جانشین تے ہر غیر پہلے عدد نوں فوری پیش رو
دیندا اے۔ $B(e,d)$ فوری پیش رو دا رشتہ ہووے، جس دی تعریف
سیٹ وچ ایہ جملہ کرے:
\[
\forall e,d\,(B(e,d)\leftrightarrow
(Z(e)\land Z(d)\land e<d\land
\neg\exists t\,(Z(t)\land e<t\land t<d))).
\]
خطی ترتیب اُتے فوری پیش رو تے جانشین ایس تعریف نال اکیلے
مقرر ہوندے نیں۔ ہر عددی جوڑی
لئی $U$ تے $V$ دے اصل-عنصر حصے غیر خالی ہون، تے ہر
پرانے مستقل تے $c_\circ$ دی قدر دوہاں حصیاں وچ ہووے۔ ہر
$J(d,m,h)$ والا گراف انہاں حصیاں وچ جزوی یک بہ یک نقشہ
ہووے، ساریاں نامی قدراں نوں اپنی قدر اُتے ای ملائے، تے ہر
اصل محمول دے ٹپلاں دا صدق دونوں رخ محفوظ رکھے۔ ایہ گراف
شرطاں پچھلے ثبوت والی پہلے درجے دی شرطاں دا جوڑی-اشاریے
والا روپ نیں۔ صریح طور اُتے ایہ جملے شامل کرو:
\[
\begin{gathered}
\forall m\,(Z(m)\to(\exists a\,U(a,m)\land\exists b\,V(b,m))),\\
\forall d,m,h\,(J(d,m,h)\to(Z(d)\land Z(m))),\\
\forall d,m,h,a,b\,((J(d,m,h)\land G(h,a,b))\to(U(a,m)\land V(b,m))),\\
\forall h,a,b,b'\,((G(h,a,b)\land G(h,a,b'))\to b=b'),\\
\forall h,a,a',b\,((G(h,a,b)\land G(h,a',b))\to a=a').
\end{gathered}
\]
ہر پرانے مستقل $c$، تے $c_\circ$، لئی ایہ جملے شامل کرو:
\[
\forall m\,(Z(m)\to(U(c,m)\land V(c,m))),\qquad
\forall d,m,h\,(J(d,m,h)\to G(h,c,c)).
\]
ہر پرانے $r$-مقامی محمول $A$ لئی ایہ جملہ شامل کرو:
\[
\forall d,m,h,a_1,\ldots,a_r,b_1,\ldots,b_r\,
\bigl((J(d,m,h)\land\bigwedge_{i=1}^{r}G(h,a_i,b_i))\to
(A(a_1,\ldots,a_r)\leftrightarrow A(b_1,\ldots,b_r))\bigr).
\]
لڑی کوڈ حصیاں نوں اصل-عنصر حصیاں نال نہ ملاؤ۔ سیٹ وچ ایہ
ابتدائی کوڈ تے دونوں توسیع شرطاں وی شامل ہون:
\[
\begin{aligned}
&\forall m\,(Z(m)\to\exists h\,J(m,m,h)),\\
&\forall d,m,h,e,a\,((J(d,m,h)\land B(e,d)\land U(a,m))\to \\
&\quad\exists k,b\,(J(e,m,k)\land V(b,m)\land G(k,a,b)\land \\
&\quad\forall u,v\,(G(h,u,v)\to G(k,u,v)))),\\
&\forall d,m,h,e,b\,((J(d,m,h)\land B(e,d)\land V(b,m))\to \\
&\quad\exists k,a\,(J(e,m,k)\land U(a,m)\land G(k,a,b)\land \\
&\quad\forall u,v\,(G(h,u,v)\to G(k,u,v)))).
\end{aligned}
\]
ایتھے $U$ تے $V$ صرف اصل عناصر چنّن، تے $J$ دی شرط وچ
پہلے دو دلائل عددی ہون۔ $\Delta$ وچ ہر غیر متعلقہ پرانا محمول خالی تے ہر
غیر متعلقہ پرانا مستقل $c_\circ$ دے برابر ہون دے جملے وی
شامل کرو۔ ایہ سارے جملے تے $!D$، $\Struct{C}$ وچ سچ
نیں۔ نال عددی مستقل $\nu_k$، معیاری عدد $k\in\Nat$ دے
کوڈ نوں نام دیندے ہون؛ $\Delta$ وچ پہلے رکن توں انہاں دی
فوری جانشینی لڑی دے جملے رکھو۔ اسیں صرف پہلے درجے دا سیٹ
آکھ رہے آں، اک واحد محدود جملہ یا حقیقی معیاری عددی دائرہ
پہلے درجے نال مجبور کرن دا دعویٰ نہیں۔

اک نواں مستقل $c_*$ چن کے جملہ-سیٹ
$\{!D\}\cup\Delta\cup\{Z(c_*),\ \nu_k<c_*:k\in\Nat\}$
لوو، جتھے سخت ترتیب $<$، $\le$ تے نا برابری نال تعریف
ہوندی اے۔ ہر محدود ذیلی سیٹ $\Struct{C}$ وچ پورا ہوندا
اے: $c_*$ نوں اوہدے اندر آئے سارے محدود معیاری عددی کوڈاں
توں وڈا معیاری عدد دیو؛ باقی سیٹ پہلے ای سچ اے۔ زبان ودھاون
نال $!D$ دا صدق وی قائم اے۔ محدود گواہی دی خاصیت توں
ایس پورے سیٹ دا نمونہ $\Struct{C}^{\dagger}$ ملدا اے۔
اوہدے اندر $c_*$ دی قدر $n^*$ سارے معیاری عددی کوڈاں توں
وڈی اے۔ اصل زبان والے، غیر خالی ذیلی !!{structure}s
$\Struct{M_{n^*}}$ تے $\Struct{N_{n^*}}$، متعلقہ
اصل-عنصر حصیاں توں لوو؛ $!D$ توں
$\Struct{M_{n^*}}\models_L !E$ تے
$\Struct{N_{n^*}}\not\models_L !E$۔ ہر معیاری محدود
$k$ لئی $n^*-k$ دا مطلب فوری پیش رو دے $k$ قدم اے؛
ایہ موجود اے تے $n^*-k>0$، کیوں جے پہلی قدر تک محدود
قدم وچ پہنچنا $n^*$ نوں کسے معیاری کوڈ دے برابر بنا دیندا۔
ہن $\mathcal{I}$ اوہناں محدود نقشیاں $p$ دا سیٹ ہووے جنہاں
دا دائرہ $k$ عناصر دا اے تے جو کسے گراف $G(h,-,-)$ دی
پابندی نیں، جتھے $J(n^*-k,n^*,h)$ اے۔ اصل جوڑی دا
اشاریہ ہن ہمیشہ $n^*$ اے؛ صرف باقی گہرائی گھٹدی اے۔
ابتدائی کوڈ والی شرط خالی نقشہ دیندی اے، سو خاندان غیر خالی
اے۔ جے مطلوب عنصر پہلے ای دائرے یا حاصل قدراں وچ اے، اوہو
$p$ رکھو۔ نہیں تاں لکھی توسیع شرط نال پورا پرانا گراف قائم
رکھ کے مطلوب نواں جوڑا لوو، تے صرف $p$ تے اوہ جوڑا رکھو؛
ایہ $k+1$ عناصر والا محدود نقشہ اے تے اگلی پیش رو گہرائی
دا کوڈ اوہنوں رکھدا اے۔ گراف دیاں شرطاں یک بہ یک پن، محمول تے
مستقل حفاظت دیندیاں نیں۔ ایس طرح دونوں اگے-پیچھے رخ حاصل
نیں تے $\Struct{M_{n^*}}\simeq_p\Struct{N_{n^*}}$۔
پر \olref[lsp]{thm:abstract-p-isom}، جس دے معمول مطابق ہون تے
نزولی مفروضے ایتھے دونوں موجود نیں، انہاں نوں $L$ دے لحاظ
نال جملہ وار ہم ارز کردا اے؛ ایہ $!E$ دے مختلف صدق دے
خلاف اے۔ سو مقدار نمائی درجے دی محدود حد موجود اے، تے اوپر دا محدود
نمائندیاں والا استدلال ہر اصل زبان دے ہر جملے لئی مطلوب
پہلے درجے دا مساوی جملہ دیندا اے۔
\end{proof}
\endgroup
% /OLP-0207

\OLEndChapterHook
\endgroup
% /OLP-0203

\OLEndPartHook
\endgroup
% /OLP-0182

% OLP-0208 content/computability/computability.tex
\begingroup
\def\olfilename{computability}
\expandafter\def\csname ol@sectioncs\endcsname{section}

\olpart{cmp}{قابلِ حسابیت}

\begin{editorial}
  ایہ حصہ قابلِ حسابیت دے نظریے بارے Jeremy Avigad دے نوٹاں
  اُتے مبنی اے۔ بازگشتی فنکشناں والے باب تے قابلِ حسابیت
  دے نظریے والے باب، دوہاں وچ مشقاں نیں؛ تے مجموعی طور اُتے ساری پیشکش نوں مقصد دی
  وضاحت، مثالاں، تفصیلاں تے مشقاں ہور شامل کرکے ودھایا جا سکدا اے۔
\end{editorial}

% OLP-0209 content/computability/recursive-functions/recursive-functions.tex
\begingroup
\def\olfilename{recursive-functions}
\expandafter\def\csname ol@sectioncs\endcsname{section}

\olchapter{cmp}{rec}{بازگشتی فنکشناں}

\begin{editorial}
  ایہ بازگشتی فنکشناں بارے Jeremy Avigad دے نوٹ نیں۔
  Richard Zach نے انہاں دی نظرثانی کیتی تے انہاں نوں ودھایا اے۔
  ایس باب وچ کجھ مشقاں وی موجود نیں، تے نحو نوں عددی روپ دینے
  \pnbLatin{(arithmetization of syntax)} بارے بحث دی بنیاد فراہم کرن لئی
  ایہ باب وکھرا وی شامل کیتا جا سکدا اے۔
\end{editorial}



% OLP-0210 content/computability/recursive-functions/introduction.tex
\begingroup
\def\olfilename{introduction}
\expandafter\def\csname ol@sectioncs\endcsname{section}

\olfileid{cmp}{rec}{int}
\olsection{تعارف}

قابلِ حسابیت دا ریاضیاتی نظریہ بنانے لئی سب توں پہلاں
قابلِ حسابیت دا اک \emph{نمونہ} بنانا پیندا اے۔ ہن اسی
قابلِ حسابیت نوں اوہ کم سمجھدے آں جو کمپیوٹر کردے نیں،
تے کمپیوٹر علامتاں نال کم کردے نیں۔ پر جدوں قابلِ حسابیت
دے نظریے بننے شروع ہوئے، حساب دی نمائندہ مثال \emph{عددی}
حساب سی۔ ریاضی دان ہمیشہ اعداد دے نظریے والے فنکشناں،
یعنی ایہو جیہے قابلِ حساب فنکشناں $f\colon \Nat^n \to
\Nat$، وچ دلچسپی رکھدے آئے نیں۔ ایس لئی حیرانی نہیں کہ
قابلِ حسابیت دے نظریے دے آغاز وچ انہاں ہی فنکشناں دا
مطالعہ ہویا۔ قابلِ حساب عددی فنکشناں دیاں سب توں جانے
پچھانے مثالاں، جیویں قدرتی عدداں اُتے جمع، ضرب تے قوت
چڑھاؤ، اک دلچسپ خاصیت سانجھی رکھدیاں نیں: انہاں نوں
\emph{بازگشتی طور اُتے} تعریف کیتا جا سکدا اے۔ ایس لئی
بازگشتی تعریفاں دے آسرے \emph{قابلِ حساب فنکشن} دی عام
تعریف دین دی کوشش فطری اے۔ اعداد دے نظریے والے فنکشناں
نوں بازگشتی طور اُتے تعریف کرن دے بہت سارے ممکنہ طریقیاں
وچوں اک خاصا سادہ تعریف دا نمونہ اتھے مرکزی بن جاندا اے:
جس نوں \emph{بنیادی بازگشت \pnbLatin{(primitive recursion)}} آکھدے نیں۔

قابلِ حساب فنکشناں دے نال نال، اسی قابلِ حساب سیٹاں تے
تعلقاں وچ وی دلچسپی رکھ سکدے آں۔ کوئی سیٹ قابلِ حساب اے
جے اسی ایہ جواب حساب کر سکّیے کہ دتا ہویا عدد سیٹ دا
!!a{element} اے یا نہیں؛ تے کوئی تعلق اوہدوں تے صرف اوہدوں
قابلِ حساب اے جدوں اسی ایہ حساب کر سکّیے کہ ٹپل
$\tuple{n_1, \dots, n_k}$ تعلق دا !!a{element} اے یا نہیں۔
کسے سیٹ یا تعلق دے \emph{خصوصیتی فنکشن \pnbLatin{(characteristic function)}} نوں ویکھ کے قابلِ حساب سیٹاں تے تعلقاں دی
بحث نوں قابلِ حساب فنکشناں دی بحث وچ سمویا جا سکدا اے۔
ایس طرح اسی بنیادی بازگشتی تعلق وی تعریف کر سکدے آں؛
مثال لئی تعلق ``$n$، $m$ نوں بغیر باقی دے تقسیم کردا اے''
اک بنیادی بازگشتی تعلق اے۔

بنیادی بازگشتی فنکشن---یعنی اوہ جنہاں نوں صرف بنیادی
بازگشت نال تعریف کیتا جا سکدا اے---پر اعداد دے نظریے
دے سارے قابلِ حساب فنکشن نہیں نیں۔ بنیادی بازگشت دیاں
کئی توسیعاں اُتے غور ہویا اے، پر فنکشناں نوں حساب کرن
دا سب توں طاقت ور تے بڑے پیمانے اُتے تسلیم شدہ اضافی
طریقہ بغیر مقررہ حد دی تلاش \pnbLatin{(unbounded search)} اے۔ اِس
نال \emph{جزوی بازگشتی فنکشناں \pnbLatin{(partial recursive functions)}} دی تعریف ملدی اے، تے اِس نال جڑی اک
تعریف \emph{عمومی بازگشتی فنکشناں \pnbLatin{(general recursive functions)}} دی وی ملدی اے۔ عمومی بازگشتی فنکشن قابلِ
حساب تے کُل (total)، یعنی ہر ان پٹ اُتے تعریف شدہ، ہوندے
نیں؛ تے ایہ تعریف ٹھیک انہاں
جزوی بازگشتی فنکشناں دی نشان دہی کردی اے جو کُل ہون۔
بازگشتی فنکشن حساب دے ہر ہور نمونے (Turing machines،
lambda calculus وغیرہ) دی کارگزاری دی نقل کر سکدے نیں؛
ایس لئی اوہ حساب دے کئی تسلیم شدہ نمونیاں وچوں اک نیں۔
\endgroup
% /OLP-0210

% OLP-0211 content/computability/recursive-functions/primitive-recursion.tex
\begingroup
\def\olfilename{primitive-recursion}
\expandafter\def\csname ol@sectioncs\endcsname{section}

\olfileid{cmp}{rec}{pre}
\olsection{بنیادی بازگشت}

قدرتی عدداں دی اک خاصیت ایہ اے کہ ہر قدرتی عدد تک~$0$
توں جانشین بنانے دا عمل~$+1$ محدود واری لا کے پہنچیا جا
سکدا اے---ہر قدرتی عدد یا تاں~$0$ اے یا جانشین اُتے
جانشین \dots{} لا کے~$0$ توں بنیا اے۔ ایس خاصیت
نوں ورت کے فنکشن
~$h\colon\Nat \to \Nat$ متعین کرن دا اک طریقہ ایہ اے:
(a)~دسّو کہ $h$ دی دلیل~$0$ لئی قیمت کیہ اے، تے
(b)~ایہ وی دسّو کہ $h(x)$ دی قیمت معلوم ہووے تاں
$h(x+1)$ دی قیمت کِویں حساب کرنی اے۔ (a) ساڈے نوں
سِدھا $h(0)$ دَسدا اے، ایس لئی $h$~،~$0$ اُتے
متعین اے۔ ہن (b) دی ہدایت $x=0$ لئی لا کے اسی
$h(1) = h(0+1)$ نوں~$h(0)$ توں حساب کر سکدے آں۔
اوہی ہدایت $x=1$ لئی ورت کے اسی $h(2) = h(1+1)$
نوں~$h(1)$ توں حساب کردے آں، تے ایویں اگے ودھدے آں۔
ہر قدرتی عدد~$x$ لئی آخر اوہ قدم آ جاندا اے جتھے
$h(x)$ توں $h(x+1)$ دی قیمت متعین کردے آں؛ ایس
لئی $h(x)$ ہر $x \in \Nat$ لئی متعین اے۔

مثال لئی، فرض کرو اسی $h\colon \Nat \to \Nat$ نوں
ہیٹھ لکھیاں دو مساواتاں نال متعین کردے آں:
\begin{align*}
h(0) & =  1\\
h(x+1) & =  2 \cdot h(x)
\end{align*}
جے ضرب کرنی پہلے توں آؤندی ہووے، تاں ایہ مساواتاں
اوپرلے (a) تے~(b) لئی لوڑ دی معلومات دیندیاں نیں۔
دوجی مساوات نوں وار وار لا کے سانوں ملدا اے:
\begin{align*}
  h(1) & = 2\cdot h(0) = 2,\\
  h(2) & = 2\cdot h(1) = 2\cdot 2,\\
  h(3) & = 2 \cdot h(2) = 2\cdot 2 \cdot 2,\\
  & \vdots
\end{align*}
اسی ویکھدے آں کہ ایہ متعین کیتا فنکشن~$h$،
$h(x) = 2^x$ اے۔

قدرتی عدداں دی ایس خاصیت نال یقین ہو جاندا اے کہ
ایہ دو شرطاں پوری کرن والا صرف اک ہی فنکشن~$h$ اے۔
ایہو جیہیاں دو مساواتاں نوں فنکشن~$h$ دی
\emph{بنیادی بازگشت دے ذریعے تعریف \pnbLatin{(definition by primitive recursion)}} آکھدے نیں۔ اِس نوں «بازگشتی» ایس لئی
آکھدے نیں کہ اسی~$h$ نوں ایس طرح تعریف کردے آں
کہ تعریف، خاص طور اُتے دوجی مساوات، دے سجے پاسے
$h$ آپ وی آؤندا اے۔ اِس نوں «بنیادی» ایس لئی
آکھدے نیں کہ $h(x+1)$ نوں متعین کرن لئی اسی
صرف~$h(x)$، یعنی اوہدے عین پہلے والی قیمت، ورتدے
آں۔~$\Nat$ اُتے فنکشن نوں بازگشتی طریقے نال تعریف
کرن دا ایہ سب توں سادہ طریقہ اے۔

اسی جمع تے ضرب ورگے ہور وی بنیادی فنکشن بنیادی
بازگشت نال تعریف کر سکدے آں۔ پر انہاں صورتاں وچ
فنکشن $2$-جگہ والے نیں۔ اسی دلیل دی اک جگہ مقرر
کر لَیندے آں تے دوجی اُتے بازگشت چلاوندے آں۔ مثال
لئی $\Add(x, y)$ دی تعریف لئی اسی~$x$ مقرر کرکے
پہلاں $y=0$ لئی قیمت، تے فیر $y+1$ لئی قیمت
~$y$ دے حوالے نال متعین کردے آں۔ چونکہ $x$ مقرر
اے، ایس لئی اوہ تعریف کرن والیاں مساواتاں دے کھبے
تے سجے دوہاں پاسیاں آوے گا:
\begin{align*}
\Add(x,0) & =  x\\
\Add(x,y+1) & =  \Add(x,y)+1
\end{align*}
ایہ مساواتاں $\Add$ دی قیمت سارے $x$
\emph{تے}~$y$ لئی متعین کردیاں نیں۔ مثال لئی
$\Add(2,3)$ کڈھن واسطے اسی $x = 2$ رکھ کے
تعریفی مساواتاں ورتدے آں: پہلی توں
$\Add(2,0) = 2$ لبھدا اے، فیر دوجی وار وار
لا کے $\Add(2,1) = 2 + 1 = 3$،
$\Add(2, 2) = 3 + 1 = 4$،
$\Add(2, 3) = 4 + 1 = 5$ لبھدا اے۔

$\Add$ دی تعریف وچ اسی دوجی مساوات دے سجے
پاسے $+$ ورتیا، پر صرف~$1$ جوڑن لئی۔ دوجے لفظاں
وچ اسی جانشین فنکشن $\Succ(z) = z+1$ نوں پچھلی
قیمت $\Add(x,y)$ اُتے لا کے $\Add(x,y+1)$
متعین کیتا۔ ایس لئی بازگشتی تعریف نوں اکّو فنکشن
دے ذریعے سمجھ سکدے آں، جیہڑا اسیں پچھلی قیمت اُتے
لاؤندے آں۔ پر فنکشن نوں صرف پچھلی قیمت اُتے ہی
نہیں، بلکہ $x$ تے~$y$ اُتے وی منحصر ہون دینا
نقصان دہ نہیں---تے کدے ضروری ہوندا اے۔ ویکھو:
\begin{align*}
  \Mult(x,0) & =  0 \\
  \Mult(x,y+1) & =  \Add(\Mult(x,y),x)
\end{align*}
ایہ فنکشن $\Mult$ دی بنیادی بازگشتی تعریف اے:
فنکشن $\Add$ نوں پچھلی قیمت $\Mult(x,y)$
تے پہلی دلیل~$x$ دوہاں اُتے لایا جاندا اے۔
ایہ فنکشن~$\Mult(x,y)$ نوں سارے $x$ تے~$y$
لئی وی متعین کردا اے۔ مثال لئی $\Mult(2,3)$،
$\Mult(2,0)$، $\Mult(2,1)$،
$\Mult(2,2)$ تے~$\Mult(2,3)$ وار وار حساب
کرن نال متعین ہوندا اے:
\begin{align*}
  \Mult(2,0) & = 0\\
  \Mult(2,1) & = \Mult(2,0+1) =
  \Add(\Mult(2,0), 2) = \Add(0, 2) = 2\\
  \Mult(2,2) & = \Mult(2,1+1) =
  \Add(\Mult(2,1), 2) = \Add(2, 2) = 4\\
  \Mult(2,3) & = \Mult(2,2+1) =
  \Add(\Mult(2,2), 2) = \Add(4, 2) = 6
\end{align*}

تاں عام نمونہ ایہ اے: فنکشن~$h(x_0, \dots, x_{k-1}, y)$
دی بنیادی بازگشتی تعریف لئی اسی دو مساواتاں دیندے
آں۔ پہلی $h(x_0, \dots, x_{k-1}, 0)$ دی قیمت
~$h$ دا حوالہ دتے بغیر متعین کردی اے۔ دوجی،
$h(x_0, \dots, x_{k-1}, y+1)$
دی قیمت $h(x_0, \dots, x_{k-1}, y)$،
باقی دلیلاں $x_0$، \dots،~$x_{k-1}$ تے~$y$
دے حوالے نال متعین کردی اے۔ ایس دوجی مساوات وچ
~$h$ دی صرف عین پہلے والی قیمت ورتی جا سکدی اے۔
جے اسی انہاں دو مساواتاں دے سجے پاسے دتے عملاں
نوں خود فنکشن $f$ تے~$g$ منّیے، تاں بنیادی بازگشت
نال نواں فنکشن~$h$ تعریف کرن دا عام نمونہ ایہ اے:
\begin{align*}
  h(x_0, \dots, x_{k-1}, 0) & = f(x_0, \dots, x_{k-1})\\
  h(x_0, \dots, x_{k-1}, y+1) & =
  g(x_0, \dots, x_{k-1}, y, h(x_0, \dots, x_{k-1}, y))
\end{align*}
$\Add$ دی صورت وچ $k=1$ تے $f(x_0) = x_0$
(شناختی فنکشن)، تے $g(x_0, y, z) = z + 1$
(اوہ $3$-جگہ والا فنکشن جیہڑا اپنی تیسری دلیل
دا جانشین واپس کردا اے):
\begin{align*}
  \Add(x_0, 0) & = f(x_0) = x_0\\
  \Add(x_0, y+1) & = g(x_0, y, \Add(x_0, y)) =
  \Succ(\Add(x_0, y))
\end{align*}
$\Mult$ دی صورت وچ $f(x_0) = 0$ (اوہ مستقل
فنکشن جیہڑا ہمیشہ~$0$ واپس کردا اے) تے
$g(x_0, y, z) = \Add(z,x_0)$ (اوہ
$3$-جگہ والا فنکشن جیہڑا اپنی آخری تے پہلی
دلیل دا جوڑ واپس کردا اے):
\begin{align*}
  \Mult(x_0,0) & =  f(x_0) = 0 \\
  \Mult(x_0,y+1) & =  g(x_0, y, \Mult(x_0,y)) =
  \Add(\Mult(x_0,y), x_0)
\end{align*}
\endgroup
% /OLP-0211

% OLP-0212 content/computability/recursive-functions/composition.tex
\begingroup
\def\olfilename{composition}
\expandafter\def\csname ol@sectioncs\endcsname{section}

\olfileid{cmp}{rec}{com}
\olsection{ترکیب}

جے $f$ تے $g$ قدرتی عدداں دے دو اک-جگہ فنکشن
ہون، تاں اسی انہاں دی ترکیب بنا سکدے آں:
$h(x) = f(g(x))$۔ نواں فنکشن~$h(x)$، فنکشناں
$f$ تے~$g$ توں \emph{ترکیب (composition)}
راہیں متعین ہوندا اے۔ اسی ایس نوں اک توں ودھ
دلیلاں والے فنکشناں تک عام کرنا چاہندے آں۔

ایہ کرن دا اک طریقہ ایہ اے: فرض کرو $f$ اک
$k$-جگہ فنکشن اے، تے $g_0$، \dots،
$g_{k-1}$، کل $k$ فنکشن نیں، جیہڑے سارے
$n$-جگہ نیں۔ فیر اسی نواں $n$-جگہ فنکشن
$h$ ایویں متعین کر سکدے آں:
\[
h(x_0, \dots, x_{n-1}) =
f(g_0(x_0, \dots, x_{n-1}), \dots, g_{k-1}(x_0, \dots, x_{n-1}))
\]
جے $f$ تے سارے $g_i$ قابلِ حساب ہون، تاں
$h$ وی قابلِ حساب اے: $h(x_0,
\dots, x_{n-1})$ حساب کرن لئی پہلاں
$y_i = g_i(x_0, \dots, x_{n-1})$ دیاں
قیمتاں ہر $i = 0$، \dots،~$k-1$
واسطے حساب کرو۔ فیر ایہ قیمتاں $f$ وچ
پا کے $h(x_0, \dots, x_{n-1}) =
f(y_0, \dots, y_{k-1})$ حساب کرو۔

موجودہ فنکشناں توں نواں فنکشن حساب کرن دی ایہ
وضاحت شاید حد توں ودھ پابند لگے۔ اک گل ایہ
کہ کدے اسی فنکشن دیاں ساریاں دلیلاں نہیں
ورتدے، جیویں اسی $g(x, y, z) = \Succ(z)$
نوں~$\Add$ دی بنیادی بازگشتی تعریف وچ
ورتن لئی متعین کیتا سی۔ فرض کرو ساڈے کول ہیٹھ لکھے
فنکشن ورتن دی اجازت اے:
\[
\Proj{n}{i}(x_0, \dots, x_{n-1}) = x_i
\]
فنکشناں~$\Proj{k}{i}$ نوں
\emph{پروجیکشن (projection)} فنکشن آکھدے نیں:
$\Proj{n}{i}$ اک $n$-جگہ فنکشن اے۔ فیر
$g$ نوں ایویں متعین کیتا جا سکدا اے:
\[
g(x, y, z) = \Succ(\Proj{3}{2}(x, y, z)).
\]
اتھے $f$ دا کردار $1$-جگہ فنکشن $\Succ$
نبھاؤندا اے، ایس لئی $k=1$۔ تے ساڈے کول
اک $3$-جگہ فنکشن $\Proj{3}{2}$ اے،
جیہڑا $g_0$ دا کردار نبھاؤندا اے۔ نتیجہ
اک $3$-جگہ فنکشن اے، جیہڑا تیسری دلیل
دا جانشین واپس کردا اے۔

پروجیکشن فنکشن ساڈے نوں دلیلاں دی ترتیب
بدل کے یا اکّو دلیل نوں دوبارہ ورت کے وی
نویں فنکشن تعریف کرن دی سہولت دیندے نیں۔
مثال لئی فنکشن $h(x) = \Add(x, x)$ نوں
ایویں متعین کیتا جا سکدا اے:
\[
h(x_0) = \Add(\Proj{1}{0}(x_0),\Proj{1}{0}(x_0)).
\]
اتھے $k=2$، $n=1$، $f(y_0,y_1)$ دا
کردار $\Add$ نبھاؤندا اے، تے
$g_0(x_0)$ تے $g_1(x_0)$ دوہاں دا کردار
~$\Proj{1}{0}(x_0)$، یعنی اک-جگہ
پروجیکشن فنکشن (شناختی فنکشن وی آکھدے نیں)،
نبھاؤندا اے۔

جے $f(y_0, y_1)$ پہلے توں ساڈے کول فنکشن
ہووے، تاں فنکشن $h(x_0, x_1) = f(x_1, x_0)$
نوں ایویں متعین کر سکدے آں:
\[
h(x_0, x_1) = f(\Proj{2}{1}(x_0, x_1),\Proj{2}{0}(x_0, x_1)).
\]
اتھے $k=2$، $n = 2$، تے $g_0$ تے
$g_1$ دے کردار، بالترتیب، $\Proj{2}{1}$
تے~$\Proj{2}{0}$ نبھاؤندے نیں۔

تہانوں ایہ فکر وی ہو سکدی اے کہ $g_0$،
\dots،~$g_{k-1}$ ساریاں دی دلیل-گنتی
اکّو~$n$ ہونی ضروری اے۔ (یاد رکھو کہ
فنکشن دی \emph{دلیل-گنتی (arity)} اوہدیاں
دلیلاں دی تعداد اے؛ $n$-جگہ فنکشن دی
دلیل-گنتی~$n$ اے۔) پر پروجیکشن فنکشن
شامل کرن نال مطلوبہ لچک مل جاندی اے۔
مثال لئی فرض کرو $f$ تے~$g$، $3$-جگہ
فنکشن نیں، تے $h$~اوہ $2$-جگہ فنکشن اے
جیہڑا ایویں متعین اے:
\[
h(x,y) = f(x,g(x,x,y),y).
\]
فنکشن~$h$ دی تعریف پروجیکشن فنکشناں نال
ایویں دوبارہ لکھ سکدے آں:
\[
h(x,y) = f(\Proj{2}{0}(x,y), g(\Proj{2}{0}(x,y), \Proj{2}{0}(x,y),
\Proj{2}{1}(x,y)), \Proj{2}{1}(x,y)).
\]
فیر $h$، $f$ دی $\Proj{2}{0}$،
$l$ تے $\Proj{2}{1}$ نال ترکیب اے، جتھے
\[
l(x,y) = g(\Proj{2}{0}(x,y),\Proj{2}{0}(x,y),\Proj{2}{1}(x,y)),
\]
یعنی $l$، $g$ دی $\Proj{2}{0}$،
$\Proj{2}{0}$ تے~$\Proj{2}{1}$
نال ترکیب اے۔
\endgroup
% /OLP-0212

% OLP-0213 content/computability/recursive-functions/pr-functions.tex
\begingroup
\def\olfilename{pr-functions}
\expandafter\def\csname ol@sectioncs\endcsname{section}

\olfileid{cmp}{rec}{prf}
\olsection{بنیادی بازگشتی فنکشن}

آؤ فیر درج کر لَیئے کہ موجودہ فنکشناں توں
بنیادی بازگشت تے ترکیب ورت کے نویں فنکشن
کِویں تعریف کیتے جا سکدے نیں۔

\begin{defn}
  \ollabel{defn:primitive-recursion} فرض کرو $f$ اک $k$-جگہ
  فنکشن اے ($k\ge 1$) تے $g$ اک $(k+2)$-جگہ فنکشن اے۔
  $f$ تے $g$ توں \emph{بنیادی بازگشت راہیں متعین کیتا فنکشن}
  اوہ $(k+1)$-جگہ فنکشن~$h$ اے جیہڑا ایہناں مساواتاں نال
  متعین ہوندا اے:
  \begin{align*}
    h(x_0,\dots,x_{k-1},0) & =  f(x_0,\dots,x_{k-1}) \\
    h(x_0,\dots,x_{k-1},y+1) & = g(x_0,\dots,x_{k-1}, y, h(x_0,\dots,x_{k-1}, y))
  \end{align*}
\end{defn}

\begin{defn}
  \ollabel{defn:composition} فرض کرو $f$ اک $k$-جگہ فنکشن اے،
  تے $g_0$، \dots، $g_{k-1}$ کل $k$ فنکشن نیں، جیہڑے
  سارے $n$-جگہ نیں۔ $f$ تے
  \emph{$g_0$، \dots،~$g_{k-1}$ توں ترکیب راہیں متعین کیتا فنکشن}
  اوہ $n$-جگہ فنکشن~$h$ اے جیہڑا ایویں متعین ہوندا اے:
  \[
  h(x_0, \dots, x_{n-1}) =
  f(g_0(x_0, \dots, x_{n-1}), \dots, g_{k-1}(x_0, \dots, x_{n-1})).
  \]
\end{defn}

$\Succ$ تے پروجیکشن فنکشناں
\[
\Proj{n}{i}(x_0,\dots,x_{n-1}) = x_i,
\]
دے نال، ہر قدرتی عدد $n$ تے $i < n$ لئی،
اسی فنکشن~$\Zero(x) = 0$ نوں وی بنیادی
بازگشتی فنکشناں وچ شامل کراں گے۔

\begin{defn}
  بنیادی بازگشتی فنکشناں دا سیٹ، $\Nat^n$ توں
  $\Nat$ ول فنکشناں دا اوہ سیٹ اے جیہڑا
  ہیٹھ لکھیاں شرطاں نال استقرائی طور اُتے
  متعین ہوندا اے:
  \begin{enumerate}
  \item $\Zero$ بنیادی بازگشتی اے۔
  \item $\Succ$ بنیادی بازگشتی اے۔
  \item ہر پروجیکشن فنکشن $\Proj{n}{i}$ بنیادی بازگشتی اے۔
  \item جے $f$ اک $k$-جگہ بنیادی بازگشتی فنکشن اے
    تے $g_0$، \dots،~$g_{k-1}$، $n$-جگہ
    بنیادی بازگشتی فنکشن نیں، تاں $f$ دی
    $g_0$، \dots،~$g_{k-1}$ نال ترکیب وی
    بنیادی بازگشتی اے۔
  \item جے $f$ اک $k$-جگہ بنیادی بازگشتی فنکشن اے
    تے $g$ اک $k+2$-جگہ بنیادی بازگشتی فنکشن اے،
    تاں $f$ تے $g$ توں بنیادی بازگشت راہیں
    متعین کیتا فنکشن وی بنیادی بازگشتی اے۔
\end{enumerate}
\end{defn}

\begin{explain}
مختصر لفظاں وچ، بنیادی بازگشتی فنکشناں دا سیٹ
اوہ سب توں چھوٹا سیٹ اے جیہڑا $\Zero$،
$\Succ$ تے پروجیکشن فنکشناں~$\Proj{n}{j}$
نوں رکھدا اے، تے ترکیب تے بنیادی بازگشت
تحت بند اے۔

بنیادی بازگشتی فنکشناں دے سیٹ نوں بیان کرن
دا اک ہور طریقہ اوہنوں ``مرحلیاں'' دے
حوالے نال متعین کرنا اے۔ $S_0$ شروعاتی
فنکشناں دا سیٹ ہووے: $\Zero$، $\Succ$
تے پروجیکشن فنکشن۔ ایہ مرحلے~$0$ دے
بنیادی بازگشتی فنکشن نیں۔ جے مرحلہ $S_i$
متعین ہو چکیا ہووے، تاں $S_{i+1}$ اوہناں
سارے فنکشناں دا سیٹ ہووے جیہڑے پہلے توں~$S_i$
وچ موجود فنکشناں اُتے ترکیب یا بنیادی بازگشت
دی اکّو مثال لا کے ملدے نیں۔ فیر
\[
S = \bigcup_{i \in \Nat} S_i
\]
سارے بنیادی بازگشتی فنکشناں دا سیٹ اے۔
\end{explain}

آؤ تصدیق کریے کہ $\Add$ بنیادی بازگشتی
فنکشن اے۔

\begin{prop}
  جمع دا فنکشن $\Add(x,y) = x+y$ بنیادی بازگشتی اے۔
\end{prop}

\begin{proof}
ساڈے کول $\Add$ دی بنیادی بازگشتی تعریف
پہلے ہی دو فنکشناں $f$ تے~$g$ دے
ذریعے اے، جیہڑی \olref{defn:primitive-recursion}
دے نمونے نال ملدی اے:
\begin{align*}
  \Add(x_0, 0) & = f(x_0) = x_0\\
  \Add(x_0, y+1) & = g(x_0, y, \Add(x_0, y)) =
  \Succ(\Add(x_0, y))
\end{align*}
ایس لئی $\Add$ بنیادی بازگشتی اے،
بشرطیکہ $f$ تے $g$ وی بنیادی بازگشتی ہون۔
$f(x_0) = x_0 = \Proj{1}{0}(x_0)$،
تے پروجیکشن فنکشن بنیادی بازگشتی گنے جاندے
نیں، ایس لئی $f$ بنیادی بازگشتی اے۔
فنکشن $g$ اوہ تِن-جگہ فنکشن
$g(x_0, y, z)$ اے جیہڑا ایویں متعین اے:
\[
g(x_0, y, z) = \Succ(z).
\]
ایس توں ہن تک ایہ نتیجہ نہیں نکلدا کہ
$g$ بنیادی بازگشتی اے، کیوں جے $g$
تے $\Succ$ بالکل اکّو فنکشن نہیں:
$\Succ$ اک-جگہ اے تے $g$ نوں تِن-جگہ
ہونا پیندا اے۔ پر اسی $g$ نوں ترکیب
راہیں باقاعدہ ایویں تعریف کر سکدے آں:
\[
g(x_0, y, z) = \Succ(\Proj{3}{2}(x_0, y, z))
\]
چونکہ $\Succ$ تے $\Proj{3}{2}$
بنیادی بازگشتی فنکشن گنے جاندے نیں،
ایس لئی $g$ وی بنیادی بازگشتی اے،
کیونکہ اوہ بنیادی بازگشتی فنکشناں توں
ترکیب راہیں متعین ہو سکدا اے۔
\end{proof}

\begin{prop}
  \ollabel{prop:mult-pr}
  ضرب دا فنکشن $\Mult(x,y) = x \cdot y$ بنیادی بازگشتی اے۔
\end{prop}

\begin{proof}
  مشق۔
\end{proof}

\begin{prob}
  \olref[cmp][rec][prf]{prop:mult-pr} ثابت کرو:
  دکھاؤ کہ $\Mult$ دی بنیادی بازگشتی تعریف
  نوں \olref[cmp][rec][prf]{defn:primitive-recursion}
  دے مطلوبہ روپ وچ لکھیا جا سکدا اے، تے
  متعلقہ فنکشن $f$ تے~$g$ بنیادی بازگشتی نیں۔
\end{prob}

\begin{ex}
ایہ ساڈی بنیادی بازگشتی تعریف دی سب توں پہلی
مثال اے:
\begin{align*}
h(0) & =  1 \\
h(y+1) & =  2 \cdot h(y).
\end{align*}
ایہ فنکشن \olref{defn:primitive-recursion} دے
مطلوبہ روپ وچ نہیں آ سکدا، کیونکہ $k=0$۔
تعریف وچ مستقل عدد $1$ تے $2$ وی آؤندے
نیں۔ پہلے مسئلے نوں حل کرن لئی اک اضافی
غیر مستعمل دلیل شامل کریے تے فنکشن~$h'$
ایویں تعریف کریے:
\begin{align*}
h'(x_0, 0) & =  f(x_0) = 1 \\
h'(x_0, y+1) & =  g(x_0, y, h'(x_0, y)) = 2 \cdot h'(x_0, y).
\end{align*}
فنکشن $f(x_0) = 1$ نوں $\Succ$
تے $\Zero$ دی ترکیب نال تعریف کر
سکدے آں: $f(x_0) = \Succ(\Zero(x_0))$۔
فنکشن $g$ نوں $g'(z) = 2 \cdot z$
تے پروجیکشن فنکشناں توں ترکیب راہیں
تعریف کر سکدے آں:
\begin{align*}
g(x_0, y, z) & = g'(\Proj{3}{2}(x_0, y, z))
\intertext{تے $g'$ نوں اگے ترکیب راہیں ایویں تعریف کر سکدے آں}
g'(z) & = \Mult(g''(z), \Proj{1}{0}(z))
\intertext{تے}
g''(z) & = \Succ(f(z)),
\end{align*}
جتھے $f$ اوہی اُپر والا اے:
$f(z) = \Succ(\Zero(z))$۔ ہن جدوں
ساڈے کول~$h'$ اے، تاں اسی فیر ترکیب
ورّت کے $h(y) =
h'(\Proj{1}{0}(y),\Proj{1}{0}(y))$
رکھ سکدے آں۔ ایہ دکھاؤندا اے کہ $h$
نوں بنیادی فنکشناں توں ترکیباں تے
بنیادی بازگشتاں دی اک ترتیب نال تعریف
کیتا جا سکدا اے، ایس لئی $h$~بنیادی
بازگشتی اے۔
\end{ex}
\endgroup
% /OLP-0213

% OLP-0214 content/computability/recursive-functions/notation-pr-functions.tex
\begingroup
\def\olfilename{notation-pr-functions}
\expandafter\def\csname ol@sectioncs\endcsname{section}

\olfileid{cmp}{rec}{not}
\olsection{بنیادی بازگشت دیاں علامتی لکھتاں}

بنیادی بازگشتی فنکشناں دی ٹھیک استقرائی تعریف
دا اک فائدہ ایہ اے کہ اسی اوہناں نوں
منظم طریقے نال بیان کر سکدے آں۔ مثال لئی،
ہر ایہو جیہے فنکشن نوں اک «علامتی لکھت»
دے سکدے آں، اِس طرح: صفر، جانشین تے
پروجیکشن فنکشناں لئی بالترتیب $\Zero$،
$\Succ$ تے $\Proj{n}{i}$ دیاں علامتاں
ورتو۔ ہن فرض کرو $h$ نوں اک $k$-جگہ
فنکشن~$f$ تے $n$-جگہ فنکشناں $g_0$،
\dots،~$g_{k-1}$ توں ترکیب راہیں
متعین کیتا گیا اے، تے ایہناں فنکشناں نوں
بالترتیب $F$، $G_0$، \dots،~$G_{k-1}$
دیاں علامتی لکھتاں دتیاں گئیاں نیں۔
تاں اک نویں علامت $\fn{Comp}_{k,n}$
ورت کے اسی فنکشن $h$ نوں
$\fn{Comp}_{k,n}[F,G_0,\dots,G_{k-1}]$
نال ظاہر کر سکدے آں۔

بنیادی بازگشت راہیں متعین فنکشناں لئی
وی اِس ورگیاں علامتی لکھتاں ورت سکدے
آں۔ فرض کرو $(k+1)$-جگہ فنکشن~$h$
نوں $k$-جگہ فنکشن~$f$ تے
$(k+2)$-جگہ فنکشن~$g$ توں بنیادی
بازگشت راہیں متعین کیتا گیا اے،
تے $f$ تے~$g$ نوں بالترتیب $F$
تے~$G$ دیاں علامتی لکھتاں دتیاں
گئیاں نیں۔ تاں~$h$ نوں دتی
علامتی لکھت $\fn{Rec}_k[F,G]$ اے۔

یاد کرو کہ جمع دا فنکشن بنیادی بازگشت
راہیں ایویں متعین ہوندا اے:
\begin{align*}
  \Add(x_0, 0) & = \Proj{1}{0}(x_0) = x_0\\
  \Add(x_0, y+1) & = \Succ(\Proj{3}{2}(x_0, y, \Add(x_0, y))) = \Add(x_0, y) +1
\end{align*}
ایتھے~$f$ دا کردار $\Proj{1}{0}$
ادا کردا اے، تے~$g$ دا کردار
$\Succ(\Proj{3}{2}(x_0, y, z))$
ادا کردا اے۔ اِس نوں
$\fn{Comp}_{1,3}[\Succ,\Proj{3}{2}]$
دی علامتی لکھت دتی جاندی اے، کیونکہ
ایہ $1$-جگہ فنکشن~$\Succ$ تے
$3$-جگہ فنکشن~$\Proj{3}{2}$ توں
ترکیب راہیں فنکشن متعین کرن دا نتیجہ
اے۔ اِس ترتیب نال اسی جمع دے فنکشن
نوں ایویں ظاہر کر سکدے آں:
\[
\fn{Rec}_1[\Proj{1}{0},\fn{Comp}_{1,3}[\Succ,\Proj{3}{2}]].
\]
ایہ علامتی لکھتاں بعض ویلے کم آؤندیاں
نیں، مثلاً بنیادی بازگشتی فنکشناں نوں
گِن کے ترتیب وار درج کرن ویلے۔

\begin{prob}
  $\Mult$ لئی پوری بنیادی بازگشتی علامتی لکھت دیو۔
\end{prob}
\endgroup
% /OLP-0214

% OLP-0215 content/computability/recursive-functions/pr-functions-computable.tex
\begingroup
\def\olfilename{pr-functions-computable}
\expandafter\def\csname ol@sectioncs\endcsname{section}

\olfileid{cmp}{rec}{cmp}
\olsection{بنیادی بازگشتی فنکشن قابلِ حساب نیں}

فرض کرو فنکشن $h$ نوں بنیادی بازگشت راہیں
متعین کیتا گیا اے:
\begin{eqnarray*}
h(\vec x, 0)     & = & f(\vec x) \\
h(\vec x, y + 1) & = & g(\vec x, y, h(\vec x, y))
\end{eqnarray*}
تے فرض کرو فنکشن $f$ تے $g$ قابلِ حساب
نیں۔ (ایتھے اسی $\vec x$ نوں $x_0$،
\dots، $x_{k-1}$ دا مختصر روپ ورتدے
آں۔) تاں $h(\vec x, 0)$ صاف طور اُتے
حساب کیتا جا سکدا اے، کیونکہ ایہ بس
$f(\vec x)$ اے جیہڑا فرض مطابق قابلِ
حساب اے۔ فیر $h(\vec x, 1)$ وی حساب
کیتا جا سکدا اے، کیونکہ $1 = 0 + 1$
تے ایس لئی $h(\vec x, 1)$ بس ایہ اے:
\begin{align*}
 h(\vec x, 1) & = g(\vec x, 0, h(\vec x, 0)) =  g(\vec x, 0, f(\vec x)).
\intertext{اسی ایس طرح اگے ودھ کے حساب کر سکدے آں}
h(\vec x, 2) & = g(\vec x, 1, h(\vec x, 1)) = g(\vec x, 1, g(\vec x, 0, f(\vec x)))\\
h(\vec x, 3) & = g(\vec x, 2, h(\vec x, 2)) = g(\vec x, 2, g(\vec x, 1, g(\vec x, 0, f(\vec x))))\\
h(\vec x, 4) & = g(\vec x, 3, h(\vec x, 3)) = g(\vec x, 3, g(\vec x, 2, g(\vec x, 1, g(\vec x, 0, f(\vec x)))))\\
& \vdots
\end{align*}
لہٰذا عام طور اُتے $h(\vec x, y)$
حساب کرن لئی پہلے $h(\vec x, 0)$،
فیر $h(\vec x, 1)$، \dots، باری باری
حساب کردے جاواں، جد تک $h(\vec x, y)$
تک نہ پہنچ جاواں۔

ایس لئی جے فنکشن $f$ تے $g$ قابلِ
حساب ہون، تاں بنیادی بازگشتی تعریف
اک نواں قابلِ حساب فنکشن دیندی اے۔
اِس طرح جے فنکشن $f$ تے $g_i$ قابلِ
حساب ہون، تاں اوہناں دی ترکیب نال
ملن والا فنکشن وی قابلِ حساب ہوندا
اے۔

چونکہ بنیادی فنکشن $\Zero$، $\Succ$
تے $\Proj{n}{i}$ قابلِ حساب نیں،
تے ترکیب تے بنیادی بازگشت قابلِ حساب
فنکشناں توں قابلِ حساب فنکشن دیندیاں
نیں، اِس لئی ہر بنیادی بازگشتی فنکشن
قابلِ حساب اے۔
\endgroup
% /OLP-0215

% OLP-0216 content/computability/recursive-functions/examples.tex
\begingroup
\def\olfilename{examples}
\expandafter\def\csname ol@sectioncs\endcsname{section}

\olfileid{cmp}{rec}{exa}
\olsection{بنیادی بازگشتی فنکشناں دیاں مثالاں}


ساڈے کول بنیادی بازگشتی فنکشناں دیاں
کجھ مثالاں پہلے توں نیں: جمع تے ضرب
دے فنکشن~$\Add$ تے $\Mult$۔
شناختی فنکشن $\fn{id}(x) = x$ بنیادی
بازگشتی اے، کیونکہ ایہ بس~$\Proj{1}{0}$
اے۔ مستقل فنکشن $\fn{const}_n(x) = n$
وی بنیادی بازگشتی نیں، کیونکہ اوہناں نوں
$\Zero$ تے $\Succ$ توں لگاتار ترکیب
راہیں متعین کیتا جا سکدا اے۔ ایہ اوہدوں
کم آؤندا اے جدوں بنیادی بازگشتی تعریفاں
وچ مستقل عدد ورتنے ہوون۔ مثال لئی، جے
اسی فنکشن $f(x) = 2 \cdot x$
متعین کرنا چاہیئے، تاں اوہنوں
$\fn{const}_n(x)$ تے ضرب دی ترکیب توں
$f(x) =
\Mult(\fn{const}_2(x), \Proj{1}{0}(x))$
دے روپ وچ حاصل کر سکدے آں۔
اگوں اسی ایہ طریقہ ورتاں گے۔

\begin{prop}
  قوت چڑھاؤ دا فنکشن $\fn{exp}(x, y) = x^y$ بنیادی بازگشتی اے۔
\end{prop}

\begin{proof}
  اسی $\fn{exp}$ نوں بنیادی بازگشت
  راہیں ایویں متعین کر سکدے آں:
  \begin{align*}
    \fn{exp}(x, 0) & = 1\\
    \fn{exp}(x, y+1) & = \Mult(x, \fn{exp}(x,y)).
    \intertext{سختی نال آکھیے تاں ایہ بنیادی بازگشتی فنکشناں توں بازگشتی تعریف نہیں۔ باقاعدہ روپ ایہ اے:}
    \fn{exp}(x, 0) & = f(x)\\
    \fn{exp}(x, y+1) & = g(x, y, \fn{exp}(x,y)).
    \intertext{جتھے}
    f(x) & = \Succ(\Zero(x)) = 1\\
    g(x, y, z) & = \Mult(\Proj{3}{0}(x, y, z), \Proj{3}{2}(x, y, z)) = x \cdot z
  \end{align*}
  تے ایس لئی $f$ تے $g$ بنیادی
  بازگشتی فنکشناں توں ترکیب راہیں
  متعین ہوندے نیں۔
\end{proof}

\begin{prop}
  پیشرو فنکشن $\fn{pred}(y)$، جیہڑا
  ایویں متعین اے،
  \[
  \fn{pred}(y) = \begin{cases}
    0 & \text{جے $y=0$}\\
    y-1 & \text{ورنہ}
  \end{cases}
  \]
  بنیادی بازگشتی اے۔
\end{prop}

\begin{proof}
  غور کرو:
 \begin{align*}
   \fn{pred}(0) & = 0 \text{ تے}\\
   \fn{pred}(y+1) & = y.
 \end{align*}
 ایہ بنیادی بازگشتی تعریف دے لگ بھگ
 اے۔ سختی نال آکھیے تاں ایہ بنیادی
 بازگشت راہیں تعریف دے نمونے وچ پورا
 نہیں آؤندا، کیونکہ اوہ نمونہ گھٹّو گھٹ
 اک ہور دلیل~$x$ منگدا اے۔ ایہ گل وی
 وکھری اے کہ $\fn{pred}(y)$ نوں
 $\fn{pred}(y+1)$ دی تعریف وچ اصل وچ
 ورتیا ہی نہیں جاندا۔ پر اسی پہلے
 $\fn{pred}'(x, y)$ ایویں متعین کر
 سکدے آں:
 \begin{align*}
   \fn{pred}'(x, 0) & = \Zero(x) = 0,\\
   \fn{pred}'(x, y+1) & = \Proj{3}{1}(x, y, \fn{pred'}(x, y)) = y.
 \end{align*}
تے فیر اوہدے توں ترکیب راہیں
$\fn{pred}$ متعین کر سکدے آں،
مثلاً $\fn{pred}(x) = \fn{pred}'(\Zero(x), \Proj{1}{0}(x))$۔
\end{proof}

\begin{prop}
  عاملی فنکشن $\fn{fac}(x) = \fact{x} = 1 \cdot 2 \cdot 3
  \cdot \dots \cdot x$ بنیادی بازگشتی اے۔
\end{prop}

\begin{proof}
  بنیادی بازگشتی تعریف دا سیدھا
  روپ ایہ اے:
  \begin{align*}
    \fn{fac}(0) &= 1\\
    \fn{fac}(y+1) & = \fn{fac}(y) \cdot (y+1).
    \intertext{باقاعدہ طور اُتے سانوں پہلے اک دو-جگہ فنکشن $h$ متعین کرنا پیندا اے}
    h(x, 0) & = \fn{const}_1(x)\\
    h(x, y+1) & = g(x, y, h(x, y))
    \intertext{جتھے $g(x, y, z) = \Mult(\Proj{3}{2}(x, y, z),
      \Succ(\Proj{3}{1}(x, y, z)))$ تے فیر رکھو}
    \fn{fac}(y) & = h(\Proj{1}{0}(y), \Proj{1}{0}(y)) = h(y,y).
  \end{align*}
  ہن توں اگے اسی تھوڑی آسانی
  رکھاں گے تے ترکیب تے بنیادی
  بازگشت دیاں باقاعدہ تعریفاں
  ہر وار نہیں دیواں گے۔
\end{proof}

\begin{prop}
  صفر اُتے رک جان والی تفریق، $x \tsub y$،
  جیہڑی ایویں متعین اے،
  \[
  x \tsub y = \begin{cases}
    0 & \text{جے $x < y$}\\
    x-y & \text{ورنہ}
  \end{cases}
  \]
  بنیادی بازگشتی اے۔
\end{prop}

\begin{proof}
  ساڈے کول ایہ اے:
  \begin{align*}
    x \tsub 0 & = x\\
    x \tsub (y+1) & = \fn{pred}(x \tsub y) 
  \end{align*}
\end{proof}

\begin{prop}
 $x$ تے $y$ وچلا فاصلہ، $\left|x-y\right|$، بنیادی بازگشتی اے۔
\end{prop}

\begin{proof}
  ساڈے کول $\left| x-y \right| = (x \tsub y) + (y \tsub x)$
  اے؛ ایس لئی فاصلے نوں $+$ تے $\tsub$
  توں ترکیب راہیں متعین کیتا جا سکدا
  اے، تے ایہ دوہاں بنیادی بازگشتی نیں۔
\end{proof}

\begin{prop}
  $x$ تے $y$ وچوں وڈا عدد، $\fn{max}(x,y)$، بنیادی بازگشتی اے۔
\end{prop}

\begin{proof}
  اسی $\fn{max}(x,y)$ نوں $+$ تے
  $\tsub$ توں ترکیب راہیں ایویں
  متعین کر سکدے آں:
  \[
  \fn{max}(x,y) \defis x + (y \tsub x).
  \]
  جے $x$ وڈا اے، یعنی $x \ge y$،
  تاں $y \tsub x = 0$، ایس لئی $x
  + (y \tsub x) = x + 0 = x$۔
  جے $y$ وڈا اے، تاں $y \tsub x =
  y - x$، تے ایس لئی $x + (y \tsub x) = x + (y - x) = y$۔
\end{proof}

\begin{prop}
  \ollabel{prop:min-pr}
  $x$ تے $y$ وچوں چھوٹا عدد، $\fn{min}(x,y)$، بنیادی بازگشتی اے۔
\end{prop}

\begin{proof}
  مشق۔
\end{proof}

\begin{prob}
  \olref[cmp][rec][exa]{prop:min-pr} ثابت کرو۔
\end{prob}

\begin{prob}
دکھاؤ کہ \[
f(x, y) =
2^{(2^{\iddots^{2^{x}}})}\raisebox{1ex}{\bigg\rbrace}
\raisebox{1ex}{\text {$y$ وار $2$}}\] بنیادی بازگشتی اے۔
\end{prob}

\begin{prob}
دکھاؤ کہ صحیح عددی تقسیم $d(x, y) = \lfloor x/y \rfloor$
(یعنی تقسیم جتھے اعشاری نقطے توں اگے
دا سارا حصہ نظرانداز کیتا جائے) بنیادی
بازگشتی اے۔ جدوں $y = 0$ ہووے، اسی
$d(x, y) = 0$ مقرر کردے آں۔
بنیادی بازگشت تے ترکیب ورت کے~$d$
دی اک صاف تعریف دیو۔
\end{prob}

\begin{prop}
بنیادی بازگشتی فنکشناں دا سیٹ
ہیٹھلے دو عملاں تحت بند اے:
\begin{enumerate}
\item محدود جمع: جے $f(\vec x, z)$ بنیادی بازگشتی اے،
تاں ایہ فنکشن وی بنیادی بازگشتی اے:
\[
g(\vec x, y) \defis \sum_{z = 0}^y f(\vec x, z).
\]
\item محدود حاصلِ ضرب: جے $f(\vec x, z)$ بنیادی بازگشتی اے،
تاں ایہ فنکشن وی بنیادی بازگشتی اے:
\[
h(\vec x, y) \defis \prod_{z = 0}^y f(\vec x, z).
\]
\end{enumerate}
\end{prop}

\begin{proof}
مثال لئی، محدود جمع نوں ایہناں
مساواتاں نال بازگشتی طور اُتے
متعین کیتا جاندا اے:
\begin{align*}
  g(\vec x, 0) & = f(\vec x, 0)\\
  g(\vec x, y+1) & = g(\vec x, y) + f(\vec x, y+1).
\end{align*}
\end{proof}
\endgroup
% /OLP-0216

% OLP-0217 content/computability/recursive-functions/pr-relations.tex
\begingroup
\def\olfilename{pr-relations}
\expandafter\def\csname ol@sectioncs\endcsname{section}

\olfileid{cmp}{rec}{prr}
\olsection{بنیادی بازگشتی تعلق}


\begin{defn}
تعلق $R(\vec x)$ نوں بنیادی بازگشتی
اوہدوں کہندے نیں جدوں اوہدا خصوصیتی
فنکشن
\[
\Char{R}(\vec x) = \left\{
  \begin{array}{ll}
  1 & \mbox{جے $R(\vec x)$ درست ہووے} \\
  0 & \mbox{ورنہ}
  \end{array}
\right.
\]
بنیادی بازگشتی ہووے۔
\end{defn}

دوجے لفظاں وچ، جدوں اسی بنیادی
بازگشتی تعلق $R(\vec x)$ دی گل
کردے آں، تاں مراد اوہ تعلق اے
جیہڑا $\Char{R}(\vec x) = 1$
دے روپ وچ لکھیا جا سکدا اے،
جتھے $\Char{R}$ اک بنیادی
بازگشتی فنکشن اے جیہڑا ہر انپٹ
اُتے یا 1 یا 0 واپس دیندا اے۔
مثال لئی، تعلق $\fn{IsZero}(x)$
اوہدوں تے صرف اوہدوں درست اے
جدوں $x = 0$ ہووے؛ ایہ فنکشن
$\Char{\fn{IsZero}}$ دے مطابق اے،
جیہڑا بنیادی بازگشت راہیں
ایویں متعین اے:
\begin{align*}
\Char{\fn{IsZero}}(0) & = 1,\\
\Char{\fn{IsZero}}(x+1) & = 0.
\end{align*}

ایہ صاف ہونا چاہیدا اے کہ تعلقاں
نوں ہور بنیادی بازگشتی فنکشناں
نال ترکیب دے روپ وچ ورتیا جا
سکدا اے۔ ایس لئی ہیٹھلے تعلق
وی بنیادی بازگشتی نیں:
\begin{enumerate}
\item برابری دا تعلق، $x = y$،
  جیہڑا $\fn{IsZero}(\left|x -
  y\right|)$ راہیں متعین اے۔
\item چھوٹا یا برابر ہوون دا تعلق،
  $x \leq y$، جیہڑا $\fn{IsZero}(x
  \tsub y)$ راہیں متعین اے۔
\end{enumerate}


\begin{prop}
  بنیادی بازگشتی تعلقاں دا سیٹ
  بولیائی عملاں تحت بند اے؛ یعنی
  جے $P(\vec x)$ تے $Q(\vec x)$
  بنیادی بازگشتی نیں، تاں
  ایہ وی نیں:
  \begin{enumerate}
  \item $\lnot P(\vec x)$
  \item $P(\vec x) \land Q(\vec x)$
  \item $P(\vec x) \lor Q(\vec x)$
  \item $P(\vec x) \lif Q(\vec x)$
  \end{enumerate}
\end{prop}

\begin{proof}
  فرض کرو $P(\vec x)$ تے $Q(\vec x)$
  بنیادی بازگشتی نیں، یعنی اوہناں دے
  خصوصیتی فنکشن $\Char{P}$ تے
  $\Char{Q}$ ایہو جیہے نیں۔ سانوں
  دکھانا اے کہ $\lnot P(\vec x)$
  وغیرہ دے خصوصیتی فنکشن وی
  بنیادی بازگشتی نیں:
  \[
  \Char{\lnot P}(\vec x) = \begin{cases}
    0 & \text{جے $\Char{P}(\vec x) = 1$ ہووے}\\
    1 & \text{ورنہ}
  \end{cases}
  \]
  اسی $\Char{\lnot P}(\vec x)$
  نوں $1 \tsub \Char{P}(\vec x)$
  دے روپ وچ متعین کر سکدے آں:
  \[
  \Char{P \land Q}(\vec x) = \begin{cases}
    1 & \text{جے $\Char{P}(\vec x) = \Char{Q}(\vec x) = 1$ ہووے}\\
    0 & \text{ورنہ}
  \end{cases}
  \]
  اسی $\Char{P \land Q}(\vec x)$
  نوں $\Char{P}(\vec x) \cdot
  \Char{Q}(\vec x)$ یا
  $\fn{min}(\Char{P}(\vec x), \Char{Q}(\vec
  x))$ دے روپ وچ متعین کر سکدے آں۔
  اِس طرح،
  \begin{align*}
    \Char{P \lor Q}(\vec x) & = \fn{max}(\Char{P}(\vec x), \Char{Q}(\vec x)) \text{ تے}\\
    \Char{P \lif Q}(\vec x) & = \fn{max}(1 \tsub
  \Char{P}(\vec x), \Char{Q}(\vec x)).
  \end{align*}
\end{proof}

\begin{prop}
  بنیادی بازگشتی تعلقاں دا سیٹ
  محدود مسوّریت تحت بند اے؛ یعنی
  جے $R(\vec x, z)$ بنیادی
  بازگشتی تعلق اے، تاں ایہ تعلق
  وی بنیادی بازگشتی نیں:
  \begin{align*}
    & \bforall{z < y}{R(\vec x, z)} \text{ تے}\\
    & \bexists{z < y}{R(\vec x, z)}.
  \end{align*}
  $\bforall{z < y}{R(\vec x, z)}$
  ، $\vec x$ تے $y$ لئی اوہدوں
  تے صرف اوہدوں درست اے جدوں
  $R(\vec x, z)$ ہر اوہ~$z$
  لئی درست ہووے جیہڑا~$y$
  توں چھوٹا اے؛ اِسے طرح
  $\bexists{z < y}{R(\vec x, z)}$
  لئی وی متعلقہ وجودی شرط اے۔
\end{prop}

\begin{proof}
  رواجاً اسی $\bforall{z < 0}{R(\vec x, z)}$
  نوں درست منّدے آں (سادہ وجہ
  ایہ اے کہ کوئی $z$،~$0$
  توں چھوٹا نہیں) تے
  $\bexists{z < 0}{R(\vec x, z)}$
  نوں غلط منّدے آں۔ محدود کلی
  مسوّر محدود حاصلِ ضرب یا
  وار وار گھٹ قدر وانگ کم کردا
  اے؛ یعنی جے $P(\vec x, y) \defiff \bforall{z <
  y}{R(\vec x, z)}$، تاں
  $\Char{P}(\vec x, y)$ ایویں
  متعین کیتا جا سکدا اے:
  \begin{align*}
    \Char{P}(\vec x, 0) & = 1\\
    \Char{P}(\vec x, y+1) & =
    \fn{min}(\Char{P}(\vec x, y), \Char{R}(\vec x, y)).
  \end{align*}
  محدود وجودی مسوّریت نوں
  اِسے طرح $\fn{max}$ ورت کے
  متعین کیتا جا سکدا اے۔
  یا فیر محدود کلی مسوّریت
  توں ایہ ہم ارزی ورت کے:
  $\bexists{z < y}{R(\vec x, z)}
  \liff \lnot \bforall{z < y}{\lnot R(\vec x, z)}$۔
  غور کرو کہ، مثال لئی،
  $\bexists{x \leq y}{\dots
  x\dots}$ دے روپ دا محدود
  مسوّر $\bexists{x < y+1}{\dots x \dots}$
  دے برابر اے۔
\end{proof}

\begin{prob}
  دکھاؤ کہ تِن-جگہ تعلق $x \equiv y \mod n$
  (یعنی modulo~$n$ دے حساب نال مطابقت)
  بنیادی بازگشتی اے۔
\end{prob}

اک ہور کم آؤن والا بنیادی بازگشتی
فنکشن شرطی فنکشن $\fn{cond}(x,y,z)$
اے، جیہڑا ایویں متعین اے:
\begin{align*}
  \fn{cond}(x,y,z) & = \begin{cases}
  y & \text{جے $x = 0$ ہووے} \\
  z & \text{ورنہ}.
\end{cases}
\intertext{ایہ بازگشتی طور اُتے ایویں متعین اے}
\fn{cond}(0,y,z) & = y,\\
\fn{cond}(x+1,y,z) & = z.
\end{align*}
ایس نوں ورت کے بنیادی بازگشتی
تعلقاں راہیں صورتاں دے حساب
نال بنیادی بازگشتی فنکشناں
دیاں تعریفاں دا جواز دتا جا
سکدا اے:

\begin{prop}
جے $g_0(\vec x)$، \dots،~$g_m(\vec x)$
بنیادی بازگشتی فنکشن نیں، تے
$R_0(\vec x)$، \dots، $R_{m-1}(\vec x)$
بنیادی بازگشتی تعلق نیں، تاں
ہیٹھ تعریف کیتا فنکشن $f$
\[
f(\vec x) = \begin{cases}
    g_0(\vec x) & \text{جے $R_0(\vec{x})$ ہووے} \\
    g_1(\vec x) & \text{جے $R_1(\vec{x})$ ہووے تے $R_0(\vec{x})$ نہ ہووے} \\
    \vdots & \\
    g_{m-1}(\vec x) & \text{جے $R_{m-1}(\vec{x})$ ہووے تے پچھلیاں وچوں کوئی شرط پوری نہ ہووے}
    \\
    g_m(\vec x) & \mbox{ورنہ}
\end{cases}
\]
وی بنیادی بازگشتی اے۔
\end{prop}

\begin{proof}
  جدوں $m = 1$ ہووے، ایہ بس
  ایس طرح متعین فنکشن اے:
  \[
  f(\vec x) = \fn{cond}(\Char{\lnot R_0}(\vec x),g_0(\vec x),g_1(\vec
  x)).
  \]
  جے $m$، $1$ توں وڈا ہووے،
  تاں اِس روپ دیاں تعریفاں
  دی ترکیب کر سکدے آں۔
\end{proof}
\endgroup
% /OLP-0217

% OLP-0218 content/computability/recursive-functions/bounded-minimization.tex
\begingroup
\def\olfilename{bounded-minimization}
\expandafter\def\csname ol@sectioncs\endcsname{section}

\olfileid{cmp}{rec}{bmi}
\olsection{محدود کمترین قدر دی تلاش}

\begin{explain}
کئی وار کسے خاصیت یا تعلق~$P$
نوں پورا کرن والا سب توں چھوٹا
عدد لے کے فنکشن متعین کرنا کم
آؤندا اے۔ جے $P$ دا فیصلہ
کیتا جا سکدا ہووے، تاں اسی
ممکنہ اعداد $0$، $1$،
$2$، \dots، اک اک کرکے
آزمائیے، جد تک~$P$ نوں
پورا کرن والا سب توں چھوٹا
عدد نہ لبھ جائے۔ ایس قسم دی
بغیر حد تلاش سانوں بنیادی
بازگشتی فنکشناں دے دائرے توں
باہر لے جاندی اے۔ پر جے ساڈی
دلچسپی صرف اوہ سب توں چھوٹا
عدد لبھن وچ ہووے جیہڑا
\emph{الگ توں دتی گئی کسے حد توں چھوٹا}
ہووے، تاں اسی بنیادی بازگشتی
دائرے وچ رہندے آں۔ دوجے لفظاں
وچ، تھوڑا ہور عمومی طور اُتے،
فرض کرو ساڈے کول بنیادی
بازگشتی تعلق~$R(x,z)$ اے۔
اوہ فنکشن ویکھو جیہڑا
$x$ تے~$y$ نوں سب توں
چھوٹا $z < y$ دیندا اے جیہدے
لئی $R(x, z)$ درست ہووے۔
ایہ وی $R(x, 0)$،
$R(x, 1)$، \dots،
$R(x, y-1)$ نوں جانچ کے
حساب کیتا جا سکدا اے۔
پر ایہ بنیادی بازگشتی کیوں اے؟
\end{explain}

\begin{prop}
جے $R(\vec x, z)$ بنیادی
بازگشتی اے، تاں فنکشن
$m_R(\vec{x}, y)$ وی بنیادی
بازگشتی اے۔ ایہ اوہ سب
توں چھوٹا~$z$ دیندا اے
جیہڑا~$y$ توں چھوٹا ہووے
تے جیہدے لئی $R(\vec x, z)$
درست ہووے، جے ایہو جیہا
کوئی ہووے؛ ورنہ $y$ دیندا
اے۔ اسی فنکشن~$m_R$ نوں
ایویں لکھاں گے:
\[
\bmin{z < y}{R(\vec{x}, z)},
\]
\end{prop}

\begin{proof}
غور کرو کہ کوئی~$z < 0$
ایہو جیہا نہیں ہو سکدا
جیہدے لئی $R(\vec{x}, z)$
درست ہووے، کیونکہ کوئی
$z < 0$ ہے ہی نہیں۔
ایس لئی $m_R(\vec x, 0) = 0$۔

جے حد $y + 1$ دے روپ
وچ ہووے، تاں تِن صورتاں نیں:
\begin{enumerate}
\item کوئی $z < y$ اے جیہدے
لئی $R(\vec{x}, z)$ درست
اے؛ اِس صورت وچ
$m_R(\vec{x}, y+1) = m_R(\vec{x}, y)$۔
\item ایہو جیہا کوئی~$z<y$
نہیں، پر $R(\vec{x}, y)$
درست اے؛ فیر
$m_R(\vec{x}, y+1) = y$۔
\item کوئی $z < y+1$ نہیں
جیہدے لئی $R(\vec{x}, z)$
درست ہووے؛ فیر
$m_R(\vec{x}, y+1) = y+1$۔
\end{enumerate}
ایس لئی اسی $m_R(\vec x, 0)$
نوں بنیادی بازگشت راہیں
ایویں متعین کر سکدے آں:
\begin{align*}
m_R(\vec{x}, 0) & = 0\\
m_R(\vec{x}, y+1) & =
\begin{cases}
m_R(\vec{x}, y) & \text{جے $m_R(\vec x, y) \neq y$ ہووے}\\
y & \text{جے $m_R(\vec x, y) = y$ تے $R(\vec{x}, y)$ درست ہووے}\\
y+1 & \text{ورنہ۔}
\end{cases}
\end{align*}
غور کرو کہ کوئی $z<y$
جیہدے لئی $R(\vec x, z)$
درست ہووے، اوہدوں تے
صرف اوہدوں اے جدوں
$m_R(\vec x,
y) \neq y$ ہووے۔
\end{proof}

\begin{prob}
فرض کرو $R(\vec x, z)$
بنیادی بازگشتی اے۔ فنکشن
$m'_R(\vec{x}, y)$ نوں
ایویں متعین کرو: ایہ سب
توں چھوٹا~$z$ واپس دیوے
جیہڑا~$y$ توں چھوٹا ہووے
تے جیہدے لئی $R(\vec{x}, z)$
درست ہووے، جے کوئی ہووے،
تے ورنہ $0$ واپس دیوے۔
ایہ تعریف بنیادی بازگشت
راہیں~$\Char{R}$ توں کرو۔
\end{prob}
\endgroup
% /OLP-0218

% OLP-0219 content/computability/recursive-functions/primes.tex
\begingroup
\def\olfilename{primes}
\expandafter\def\csname ol@sectioncs\endcsname{section}

\olfileid{cmp}{rec}{pri}
\olsection{اولی اعداد}

محدود مسوّریت تے محدود کمترین قدر دی تلاش سانوں ایہہ وکھاون
لئی کافی اوزار دیندیاں نیں کہ قدرتی طور سامنے آون والے
فنکشن تے تعلق بنیادی بازگشتی نیں۔ مثال لئی ``$x$
، $y$ نوں پورا ونڈدا اے'' والا تعلق ویکھو، جیہنوں
$x \mid y$ لکھدے آں۔ $x \mid y$ اوہدوں درست اے
جدوں $y$،~$x$ دا صحیح عددی مضرب ہووے؛ غیر صفر
تقسیم کرن والے دی صورت وچ، ایہہ $y$ نوں~$x$ نال
بغیر باقی تقسیم کرن دے برابر اے۔ (جے ایہہ تعلق درست
نہ ہووے تے تقسیم کرن والا غیر صفر ہووے، یعنی $x$
نال $y$ دی تقسیم وچ باقی $> 0$ ہووے، تاں اسی
$x \nmid y$ لکھدے آں۔) دوجے لفظاں وچ، $x \mid y$
اوہدوں تے صرف اوہدوں اے جدوں کسے~$z$ لئی
$x \cdot z = y$ ہووے۔ جے ایہو جیہا شاہد موجود ہووے،
تاں اک شاہد~$z$ ایسا چنیا جا سکدا اے جیہڑا
$\leq y$ ہووے؛ صفر نال صفر دی صورت وچ صفر ای
چن لو۔ ایس لئی $x \mid y$ اوہدوں تے صرف اوہدوں
اے جدوں کسے $z \le y$ لئی $x \cdot z = y$
ہووے۔ اسی تعلق $x \mid y$ نوں $=$ تے ضرب توں
محدود وجودی مسوّریت راہیں ایویں متعین کر سکدے آں:
\[
x \mid y \defiff \bexists{z \leq y}{(x \cdot z) = y}.
\]
ایس طرح اسی وکھا دتا اے کہ $x \mid y$ بنیادی
بازگشتی اے۔

قدرتی عدد~$x$ نوں \emph{اولی} کہندے آں جے اوہ
نہ $0$ ہووے نہ $1$، تے اوہنوں صرف $1$ تے اوہ
آپ پورا ونڈن۔ دوجے لفظاں وچ، اولی عدد ایہو جیہے
نیں کہ جدوں وی $y \mid x$ ہووے، تاں یا $y = 1$
ہووے یا~$y=x$۔ ایہہ جانچن لئی کہ $x$~اولی اے
یا نہیں، سانوں صرف ایہہ جانچنا پیندا اے کہ
$y \mid x$ اے یا نہیں، اوہناں سارے $y \le x$
لئی؛ کیونکہ ایتھے حد
دے تحت جے $y > x$ ہووے، تاں آپے~$y \nmid x$
ہووے گا۔ ایس لئی تعلق $\fn{Prime}(x)$، جیہڑا
اوہدوں تے صرف اوہدوں درست اے جدوں $x$ اولی ہووے،
ایویں متعین ہو سکدا اے:
\[
\fn{Prime}(x) \defiff x \geq 2 \land \bforall{y \leq x}{(y \mid x \lif y
  = 1 \lor y = x)}
\]
تے ایس لئی ایہہ بنیادی بازگشتی اے۔

اولی اعداد $2$، $3$، $5$، $7$، $11$ وغیرہ نیں۔
فنکشن $p(x)$ ویکھو، جیہڑا ایس سلسلے دا $x$واں
اولی عدد واپس دیندا اے؛ یعنی $p(0) =
2$، $p(1) = 3$، $p(2) = 5$ وغیرہ۔ (سہولت
لئی اسی اکثر $p(x)$ نوں $p_x$ لکھاں گے
($p_0=2$، $p_1=3$ وغیرہ)۔)

جے ساڈے کول فنکشن
$\fn{nextPrime(x)}$ ہووے، جیہڑا~$x$ توں وڈا
پہلا اولی عدد واپس دیندا اے، تاں $p$ نوں بنیادی
بازگشت راہیں سوکھے نال متعین کیتا جا سکدا اے:
\begin{align*}
  p(0) & = 2\\
  p(x+1) & = \fn{nextPrime}(p(x))
\end{align*}
چونکہ $\fn{nextPrime}(x)$ اوہ سب توں چھوٹا
$y$ اے جیہڑا $y > x$ ہووے تے $y$~اولی ہووے،
ایہہ بغیر حد تلاش نال سوکھے حساب کیتا جا سکدا اے۔
پر اقلیدس دے اک نتیجے کر کے ایہہ محدود کمترین
قدر دی تلاش نال وی متعین ہو سکدا اے: $x$ توں
وڈا تے $\fact{x}+1$ توں ودھ نہ ہوون والا کوئی
اولی عدد ضرور ہوندا اے۔
\[
  \fn{nextPrime(x)} =
  \bmin{y \leq \fact{x}+1}{(y > x \land \fn{Prime}(y))}.
\]
ایہہ وکھاؤندا اے کہ $\fn{nextPrime}(x)$ تے ایس
لئی $p(x)$ صرف قابلِ حساب ای نہیں، بلکہ بنیادی
بازگشتی نیں۔

(جے دلچسپی ہووے، تاں اقلیدس دے نتیجے دا چھوٹا
ثبوت ایہہ اے۔ دو توں گھٹ حد دے دو کیس سیدھے
جانچے جا سکدے نیں؛ ہُن حد گھٹ توں گھٹ دو ہووے۔
فرض کرو $p_n$ سب توں وڈا اولی عدد اے جیہڑا
$\le x$ ہووے، تے ضرب حاصل $p =
p_0\cdot p_1 \cdot \dots \cdot p_n$ ویکھو؛ ایہہ
اوہناں سارے اولی اعداد~$\le x$ دا ضرب حاصل اے۔ یا تاں
$p+1$ خود اولی اے یا $x$ تے~$p+1$ دے وچ
کوئی اولی عدد اے۔ کیوں؟ فرض کرو $p+1$ اولی
نہیں۔ تاں کوئی اولی عدد $q \mid p+1$ اے جیہڑا
$q < p+1$ ہووے۔ $\le x$ والے اولی اعداد وچوں
کوئی وی $p+1$ نوں پورا نہیں ونڈدا۔ ($p$ دی تعریف
مطابق، ہر اولی عدد $p_i \le x$،~$p$ نوں پورا
ونڈدا اے، یعنی باقی~$0$ نال۔ ایس لئی ہر
$p_i \le x$، $p+1$ نوں تقسیم کرن اُتے باقی~$1$
دیندا اے، تے ایس کرکے $p_i \nmid p+1$۔) لہٰذا
$q$ اولی اے جیہڑا $> x$ تے $< p+1$ اے۔
نالے $p \le \fact{x}$، ایس لئی کوئی اولی عدد
$> x$ تے $\le
\fact{x}+1$ ضرور اے۔)

\begin{prob}
محدود کمترین قدر دی تلاش ورت کے صحیح عددی تقسیم
$d(x, y)$ متعین کرو۔
\end{prob}
\endgroup
% /OLP-0219

% OLP-0220 content/computability/recursive-functions/sequences.tex
\begingroup
\def\olfilename{sequences}
\expandafter\def\csname ol@sectioncs\endcsname{section}

\olfileid{cmp}{rec}{seq}
\olsection{سلسلے}

بنیادی بازگشتی فنکشناں دا سیٹ خاصا مضبوط اے۔ پر جدوں ساڈے کول
\emph{سلسلیاں} نال کم کرن دا مناسب طریقہ ہووے گا، تاں اسی ہور
وی زیادہ کم کر سکاں گے۔ اسی قدرتی عدداں دے محدود سلسلیاں
نوں قدرتی عدداں نال ایویں جوڑاں گے: سلسلہ
$\langle a_0, a_1, a_2, \dots, a_k \rangle$ ایس عدد دے
مطابق اے:
\[
p_0^{a_0+1} \cdot p_1^{a_1+1} \cdot p_2^{a_2+1} \cdot \dots \cdot
p_k^{a_k+1}.
\]
اسی ہر اسّ وچ اک شامل کردے آں تاکہ، مثال لئی، سلسلے
$\langle 2, 7, 3\rangle$ تے $\langle 2, 7, 3, 0, 0 \rangle$
دے عددی کوڈ وکھرے ہون۔ $0$ تے~$1$ دونوں نوں خالی
سلسلے دا کوڈ منیا جا سکدا اے؛ پکی چُون لئی
$\emptyseq$ توں مراد~$0$ لاں گے۔

سلسلیاں دی ایہہ کوڈکاری اس کرکے کم کردی اے کہ حساب دا
بنیادی قضیہ کہندا اے: ہر قدرتی عدد $n \ge 2$ نوں اکّو
اک طریقے نال ایس روپ وچ لکھیا جا سکدا اے:
\[
n = p_0^{a_0} \cdot p_1^{a_1} \cdot \dots \cdot p_k^{a_k}
\]
جتھے $a_k \ge 1$ ہووے۔ ایس توں پکا ہوندا اے کہ نقشہ
$\tuple{}(a_0,
\dots, a_k) = \tuple{a_0, \dots, a_k}$ اک بہ اک اے:
وکھرے سلسلے وکھرے عدداں نال جُڑدے نیں؛ ہر عدد نال
زیادہ توں زیادہ اک سلسلہ جُڑدا اے۔

ہُن اسی وکھاواں گے کہ سلسلے دی لمبائی لبھنا، اوہدا
$i$واں رکن لبھنا، اک رکن سلسلے دے آخر شامل کرنا، تے
دو سلسلے جوڑنا، سب بنیادی بازگشتی عمل نیں۔

\begin{prop}
فنکشن $\len{s}$، جیہڑا سلسلے $s$ دی لمبائی واپس
دیندا اے، بنیادی بازگشتی اے۔
\end{prop}

\begin{proof}
فرض کرو $R(i, s)$ ایہہ تعلق اے:
\[
R(i, s) \text{ اوہدوں تے صرف اوہدوں جدوں }
p_i \mid s \land p_{i+1} \nmid s.
\]
$R$ واضح طور اُتے بنیادی بازگشتی اے۔ جدوں وی $s$
غیر خالی سلسلے دا کوڈ ہووے، یعنی
\[
s = p_0^{a_0+1} \cdot \dots \cdot p_{k}^{a_{k}+1},
\]
تاں $R(i,s)$ اوہدوں درست اے جدوں $p_i$ سب توں
وڈا اولی عدد ہووے جیہدے لئی $p_i \mid s$؛
یعنی $i = k$۔ ایس لئی
$s$ دی لمبائی $i+1$ اوہدوں تے صرف اوہدوں اے
جدوں $p_i$ سب توں وڈا اولی عدد ہووے جیہڑا~$s$
نوں پورا ونڈدا اے۔ ایس کرکے اسی ایویں رکھ سکدے آں:
\[
\len{s} =
\begin{cases}
  0 & \text{جے $s = 0$ یا $s = 1$ ہووے} \\
  1 + \bmin{i < s}{R(i, s)} & \text{ورنہ}
\end{cases}
\]
ایتھے محدود کمترین قدر دی تلاش ورت سکدے آں؛ کیونکہ
صرف اک $i$ ای $R(i,s)$ نوں پورا کردا اے، جے
$s$~سلسلے دا کوڈ ہووے؛ تے جے $i$ موجود
ہووے تاں اوہ خود~$s$ توں چھوٹا ہوندا اے۔
\end{proof}

\begin{prop}
فنکشن $\fn{append}(s,a)$، جیہڑا $a$ نوں سلسلے
$s$ دے آخر شامل کرن دا نتیجہ واپس دیندا اے، بنیادی بازگشتی اے۔
\end{prop}

\begin{proof}
$\fn{append}$ نوں ایویں متعین کیتا جا سکدا اے:
\[
\fn{append}(s,a) =
\begin{cases}
  2^{a+1} & \text{جے $s = 0$ یا $s = 1$ ہووے} \\
  s \cdot p_{\len{s}}^{a+1} & \text{ورنہ}
\end{cases}
\]
\end{proof}

\begin{prop}
فنکشن $\fn{element}(s,i)$، جیہڑا $i$واں رکن
سلسلے $s$ وچوں
واپس دیندا اے (شروع دے رکن نوں $0$واں آکھدے آں)،
یا $0$ واپس دیندا اے جے $i$، $s$ دی لمبائی دے
برابر یا اوہدے توں وڈا ہووے، بنیادی بازگشتی اے۔
\end{prop}

\begin{proof}
غور کرو کہ $a$، $i$واں رکن~$s$ دا اوہدوں تے
صرف اوہدوں اے جدوں $p_i^{a+1}$،~$p_i$ دی اوہ
سب توں وڈی قوت ہووے جیہڑی~$s$ نوں پورا ونڈدی
اے؛ یعنی $p_i^{a+1} \mid s$ پر
$p_i^{a+2} \nmid s$۔ ایس لئی:
\[
\fn{element}(s,i) =
\begin{cases}
  0 & \mbox{جے $i \geq \len{s}$ ہووے} \\
  \bmin{a < s}{(p_i^{a+2} \nmid s)} & \text{ورنہ}
\end{cases}
\]
\end{proof}

اوپر متعین فنکشناں دے باقاعدہ ناں ہر وار ورتن دی بجائے
اسی مختصر علامتاں لیانیاں آں۔ $(s)_i$ نوں
$\fn{element}(s,i)$ دی جگہ لکھاں گے، تے $\tuple{s_0, \dots, s_k}$
نال ایہہ مختصر کراں گے:
\[
\fn{append}(\fn{append}(\dots \fn{append}(\emptyseq,s_0)
\dots),s_k).
\]
غور کرو کہ جے $s$ دی لمبائی~$k$ اے، تاں
$s$ دے رکن $(s)_0$، \dots،~$(s)_{k-1}$ نیں۔

\begin{prop}
فنکشن $\fn{concat}(s,t)$، جیہڑا دو سلسلے
جوڑدا اے، بنیادی بازگشتی اے۔
\end{prop}

\begin{proof}
سانوں ایہو جیہا فنکشن $\fn{concat}$ چاہیدا اے جیہدے لئی
\[
  \fn{concat}(\tuple{a_0, \dots, a_k}, \tuple{b_0, \dots, b_l}) =
  \tuple{a_0, \dots, a_k, b_0, \dots, b_l}.
\]
اسی اک ``مددگار'' فنکشن
$\fn{hconcat}(s,t,n)$ ورتاں گے جیہڑا پہلے $n$ رکن
$t$ توں لے کے~$s$ دے آخر جوڑدا اے۔ ایہہ فنکشن بنیادی
بازگشت راہیں ایویں متعین ہو سکدا اے:
\begin{align*}
  \fn{hconcat}(s,t,0) & = s\\
  \fn{hconcat}(s,t,n+1) & = \fn{append}(\fn{hconcat}(s,t,n),(t)_n)
  \intertext{فیر اسی $\fn{concat}$ نوں ایویں متعین کر سکدے آں:}
  \fn{concat}(s,t) & = \fn{hconcat}(s,t,\len{t}).
\end{align*}
\end{proof}

اسی $s \concat t$ نوں $\fn{concat}(s,t)$ دی بجائے لکھاں گے۔

ساڈے لئی سلسلے دے عددی کوڈ دی اک حد، اوہدی لمبائی
تے سب توں وڈے رکن دے حساب نال لبھ سکنا مفید ہووے گا۔
فرض کرو $s$ اک سلسلہ اے جیہدی لمبائی~$k$ اے، تے
ہر رکن کسے عدد~$x$ توں چھوٹا یا اوہدے برابر اے۔
تاں $s$ دے زیادہ توں زیادہ $k$ اولی عوامل نیں،
ہر اک زیادہ توں زیادہ~$p_{k-1}$ اے، تے ہر اک دا
اسّ زیادہ توں زیادہ $x+1$ ہوندا اے، جدوں~$s$ دی اولی عامل
بندی لکھی جاندی اے۔ دوجے لفظاں وچ، مثبت لمبائی لئی اسی
ایویں متعین کردے آں؛ خالی سلسلے لئی حد اک مقرر اے:
\[
\fn{sequenceBound}(x,k) = p_{k-1}^{k \cdot (x+1)},
\]
تاں اوپر دسے سلسلے~$s$ دا عددی کوڈ زیادہ توں
زیادہ~$\fn{sequenceBound}(x,k)$ اے۔

سلسلیاں دی ایہہ حد سانوں محدود تلاش نال نویں فنکشن
متعین کرن دا طریقہ دیندی اے۔ مثال لئی، اسی
$\fn{concat}$ نوں محدود تلاش نال متعین کر سکدے آں۔
سانوں صرف لبھّن والی شے (جوڑے گئے سلسلے دے عدد)
دی اک بنیادی بازگشتی \emph{تعیین} تے ایہہ حد دینی
اے کہ کتھے تک لبھنا اے۔ ایہہ تعریف کم کردی اے:
\begin{align*}
  \fn{concat}(s,t) = {} & \bmin{v < \fn{sequenceBound}(s+t,\len{s} +
    \len{t})}{} \\
  & \quad(\len{v} = \len{s} + \len{t} \land {}\\
    & \qquad \bforall{i < \len{s}}{((v)_i = (s)_i)} \land {} \\
      & \qquad \bforall{j < \len{t}}{((v)_{\len{s}+j} = (t)_j)})
\end{align*}

\begin{prob}
وکھاؤ کہ اک بنیادی بازگشتی فنکشن~$\fn{sconcat}(s)$
موجود اے جیہدے لئی
\[
\fn{sconcat}(\tuple{s_0, \dots, s_k}) = s_0 \concat \dots \concat s_k.
\]
\end{prob}

\begin{prob}
وکھاؤ کہ اک بنیادی بازگشتی فنکشن~$\fn{tail}(s)$
موجود اے جیہدے لئی
\begin{align*}
  \fn{tail}(\emptyseq) & = 0 \text{ تے}\\
  \fn{tail}(\tuple{s_0, \dots, s_{k}}) & = \tuple{s_1, \dots, s_{k}}.
\end{align*}
\end{prob}

\begin{prop}
  \ollabel{prop:subseq}
فنکشن $\fn{subseq}(s, i, n)$، جیہڑا $s$ دا اوہ
ذیلی سلسلہ واپس دیندا اے جیہدی لمبائی~$n$ اے تے
جیہڑا $i$ویں رکن توں شروع ہوندا اے، بنیادی بازگشتی اے۔
\end{prop}

\begin{proof}
مشق۔
\end{proof}

\begin{prob}
  \olref[cmp][rec][seq]{prop:subseq} ثابت کرو۔
\end{prob}
  
\endgroup
% /OLP-0220

% OLP-0221 content/computability/recursive-functions/trees.tex
\begingroup
\def\olfilename{trees}
\expandafter\def\csname ol@sectioncs\endcsname{section}

\olfileid{cmp}{rec}{tre}
\olsection{درخت}

کدی کدی درختاں نوں قدرتی عدداں دے روپ وچ ظاہر کرنا
مفید ہوندا اے؛ جداں اسی سلسلیاں نوں عدداں دے روپ وچ،
تے اوہناں دیاں خاصیتاں تے عملاں نوں اوہناں دے کوڈاں
اُتے بنیادی بازگشتی تعلقاں تے فنکشناں دے روپ وچ ظاہر
کر سکدے آں۔ ایتھے اسی سلسلے تے اوہناں دے کوڈ ورتاں گے۔
درخت یا تاں اکّلی گرہ ہو سکدا اے (ممکن اے اوہدے
نال کوئی لیبل ہووے)، یا اک گرہ جیہدے نال ذیلی درخت
جُڑے ہون (اوہدے نال وی لیبل ہو سکدا اے)۔
ایس گرہ نوں درخت دی \emph{جڑ} کہندے آں، تے اوہدے نال
سِدھے جُڑے ذیلی درختاں نوں \emph{براہِ راست ذیلی درخت}۔

اسی درختاں نوں بازگشتی طور اُتے سلسلہ
$\tuple{k, d_1, \dots, d_k}$ نال کوڈ کردے آں،
جتھے $k$ براہِ راست ذیلی درختاں دی تعداد اے،
تے $d_1$، \dots،~$d_k$ اوہناں دے کوڈ نیں۔
جے گرہواں دے لیبل وی ہون، تاں اوہناں نوں براہِ راست
ذیلی درختاں توں بعد شامل کیتا جا سکدا اے۔ ایس لئی
صرف اک گرہ تے لیبل~$l$ والا درخت
$\tuple{0,l}$ نال کوڈ ہووے گا؛ تے اک جڑ (لیبل~$l_1$)
نال جُڑیاں دو اکّلی گرہواں (لیبل $l_2$، $l_3$)
والا درخت $\tuple{2,
  \tuple{0, l_2}, \tuple{0, l_3}, l_1}$ نال کوڈ ہووے گا۔

\begin{prop}
  \ollabel{prop:subtreeseq}
فنکشن $\fn{SubtreeSeq}(t)$، جیہڑا اک سلسلے دا کوڈ
واپس دیندا اے جیہدے رکن کوڈ~$t$ والے درخت دے سارے
ذیلی درختاں دے کوڈ ہون، بنیادی بازگشتی اے۔
\end{prop}

\begin{proof}
پہلاں غور کرو کہ $\fn{ISubtrees}(t) = \fn{subseq}(t, 1, (t)_0)$
بنیادی بازگشتی اے تے درخت~$t$ دے براہِ راست
ذیلی درختاں دے کوڈ واپس دیندا اے۔ ہُن اسی اک مددگار
فنکشن $\fn{hSubtreeSeq}(t,n)$ متعین کر سکدے آں
جیہڑا جڑ توں $n$ قدم تک دے سارے ذیلی درختاں دا
سلسلہ حساب کردا اے۔ درخت~$t$ دی جڑ توں $0$
قدم تک دے ذیلی درختاں دا سلسلہ—یعنی اوہ جیہڑا
درخت~$t$ دی جڑ توں شروع ہوندا اے—صرف~$t$
اُتے مشتمل اے۔ سطح~$n+1$ تک درخت $t$ دے سارے
ذیلی درختاں دا سلسلہ لین لئی، سطح $n$ تک جمع
ذیلی درختاں نال سطح $n$ تک دے انہاں ساریاں دے
براہِ راست ذیلی درختاں دا سلسلہ جوڑ دیندے آں۔
ایہہ سارے اکٹھے کرن لئی غور کرو کہ جے $f(x)$
بنیادی بازگشتی فنکشن اے جیہڑا سلسلیاں دے کوڈ واپس
دیندا اے، تاں $g_f(s, k) = f((s)_0) \concat
\dots \concat f((s)_k)$ وی بنیادی بازگشتی اے:
    \begin{align*}
      g(s, 0) & = f((s)_0)\\
      g(s, k+1) & = g(s, k) \concat f((s)_{k+1})
    \end{align*}
    مثال لئی، جے $s$ درختاں دا سلسلہ اے، تاں
    $h(s) = g_{\fn{ISubtrees}}(s, \len{s})$ سلسلے~$s$
    دے رکناں دے براہِ راست ذیلی درختاں دا سلسلہ دیندا
    اے۔ اسی ایہہ ورت کے~$\fn{hSubtreeSeq}$ نوں
    ایویں متعین کر سکدے آں:
    \begin{align*}
      \fn{hSubtreeSeq}(t, 0) & = \tuple{t} \\
      \fn{hSubtreeSeq}(t, n+1) & = \fn{hSubtreeSeq}(t, n) \concat
      h(\fn{hSubtreeSeq}(t, n)).
    \end{align*}
    کوڈ~$t$ والے درخت وچ ذیلی درختاں دی زیادہ توں
    زیادہ سطح، یعنی جڑ تے کسے پتے والی گرہ دے
    وچکار زیادہ توں زیادہ فاصلہ، کوڈ~$t$ توں ودھ نہیں۔
    ایس لئی کوڈ~$t$ والے درخت دے سارے ذیلی درختاں
    دے کوڈاں دا اک سلسلہ $\fn{hSubtreeSeq}(t, t)$
    نال مل جاندا اے۔
\end{proof}

\begin{prob}
  ثبوت وچ $\fn{hSubtreeSeq}$ دی تعریف،
  \olref[cmp][rec][tre]{prop:subtreeseq} وچ، عام طور
  اُتے دہرائیاں شامل کردی اے۔ اک وکھری تعریف دیو
  جیہڑی یقینی بنائے کہ نتیجے والی فہرست وچ کسے
  ذیلی درخت دا کوڈ صرف اک وار آوے۔
\end{prob}
\endgroup
% /OLP-0221

% OLP-0222 content/computability/recursive-functions/other-recursions.tex
\begingroup
\def\olfilename{other-recursions}
\expandafter\def\csname ol@sectioncs\endcsname{section}

\olfileid{cmp}{rec}{ore}
\olsection{بازگشت دے ہور روپ}

جوڑیاں تے سلسلے بناون دے طریقے ورت کے، بنیادی
بازگشت دے ہور وکھرے (تے مفید) روپاں دا جواز دتا جا
سکدا اے۔ مثال لئی، اکثر دو فنکشناں نوں اکّو ویلے
متعین کرنا مفید ہوندا اے، جداں ہیٹھلی تعریف وچ:
\begin{align*}
h_0(\vec x, 0) & = f_0(\vec x) \\
h_1(\vec x, 0) & = f_1(\vec x) \\
h_0(\vec x, y+1) & = g_0(\vec x, y, h_0(\vec x, y), h_1(\vec x, y)) \\
h_1(\vec x, y+1) & = g_1(\vec x, y, h_0(\vec x, y), h_1(\vec x, y))
\end{align*}
ایہہ \emph{اکٹھّی بازگشت \pnbLatin{(simultaneous recursion)}}
دی اک مثال اے۔ فنکشن متعین کرن دا اک ہور مفید طریقہ
ایہہ اے کہ $h(\vec x, y+1)$ دی قیمت نوں
$h(\vec x, 0)$، \dots،~$h(\vec x, y)$ دیاں
\emph{ساریاں} قیمتاں دے حوالے نال دتا جاوے، جداں
ہیٹھلی تعریف وچ:
\begin{align*}
h(\vec x, 0) & = f(\vec x) \\
h(\vec x, y+1) & = g(\vec x, y, \tuple{h(\vec x, 0), \dots, h(\vec x, y)}).
\end{align*}
ہیٹھلا خاکہ ایس خیال نوں ہور مختصر طور اُتے ظاہر
کردا اے:
\[
h(\vec x, y) = g(\vec x, y, \tuple{h(\vec x, 0), \dots, h(\vec x, y-1)})
\]
ایتھے $g$ دی آخری دلیل بارے ایہہ سمجھنا اے کہ
جدوں $y$، $0$ ہووے، تاں اوہ صرف خالی سلسلہ اے۔ دوہاں
روپاں وچ خیال ایہہ اے کہ «جانشین والا قدم» حساب
کردے ویلے فنکشن $h$ ہن تک حساب کیتیاں قیمتاں
دا پورا سلسلہ ورت سکدا اے۔ ایس نوں
\emph{پچھلی ساریاں قیمتاں والی} بازگشت
\pnbLatin{(course-of-values recursion)} کہندے نیں۔ اک خاص
مثال لئی، ایس نال ہیٹھلی قسم دی تعریف دا جواز دتا
جا سکدا اے:
\begin{align*}
h(\vec x, y) & = \begin{cases}
  g(\vec x, y, h(\vec x, k(\vec x, y))) & \text{جے $k(\vec x, y) < y$} \\
  f(\vec x) & \text{نہیں تاں}
\end{cases}
\end{align*}
دوجے لفظاں وچ، $h$ دی قیمت $y$ اُتے
\emph{کسے وی} پچھلی دلیل اُتے $h$ دی قیمت ورت کے
لبھی جا سکدی اے؛ اوہ پچھلی دلیل
~$k$ دیندا اے۔

\begin{prob}
  باقی دا فنکشن $r(x,y)$ پچھلی ساریاں قیمتاں والی
  بازگشت نال متعین کرو۔ (جے $x$، $y$ قدرتی عدد
  ہون تے $y > 0$، تاں $r(x,y)$ اوہ عدد اے
  جیہڑا~$y$ توں چھوٹا اے تے
  $x = z\times y + r(x,y)$ کسے~$z$ لئی ہووے۔ گل نوں مقرر
  رکھن لئی، جے $y=0$، تاں $r(x,0) = 0$
  رکھ لَیندے آں۔)
\end{prob}

ایہہ سوچو کہ انہاں فنکشناں نوں عام بنیادی بازگشت
نال کِویں حاصل کیتا جا سکدا اے۔ بنیادی بازگشت دا
اک آخری روپ ہور لچکدار اے، کیوں جے ایس وچ
\emph{پیرامیٹر} (نال چلن والیاں قیمتاں) نوں عمل
دے دوران بدلّن دی اجازت اے:
\begin{align*}
h(\vec x, 0) & = f(\vec x) \\
h(\vec x, y+1) & = g(\vec x, y, h(k(\vec x), y))
\end{align*}
ایس دا عمل وی عام بنیادی بازگشت راہیں سرانجام دتا
جا سکدا اے۔ (ایہہ کرنا سوکھا نہیں۔ اشارے لئی، حساب
دے قدم ہتھ نال کھول کے ویکھو۔)
\endgroup
% /OLP-0222

% OLP-0223 content/computability/recursive-functions/non-pr-functions.tex
\begingroup
\def\olfilename{non-pr-functions}
\expandafter\def\csname ol@sectioncs\endcsname{section}

\olfileid{cmp}{rec}{npr}
\olsection{بنیادی بازگشتی نہ ہون والے فنکشن}

بنیادی بازگشتی فنکشن، بدیہی سمجھ مطابق قابلِ حساب
سارے فنکشناں نوں نہیں سمیٹدے۔ بدیہی طور اُتے ایہہ
صاف ہونا چاہیدا اے کہ اسی سارے اک-جگہ والے بنیادی
بازگشتی فنکشناں دی فہرست
$f_0$، $f_1$، $f_2$،~\dots ایویں بنا سکدے آں
کہ $f_x$ دی قیمت انپٹ $y$ اُتے موثر طریقے نال
حساب کر سکّیے؛ دوجے لفظاں وچ، فنکشن $g(x,y)$،
جیہدی تعریف ایہہ اے:
\[
g(x,y) = f_x(y)
\]
قابلِ حساب اے۔ پر فیر ایہہ فنکشن وی قابلِ حساب اے:
\begin{eqnarray*}
h(x) & = & g(x,x) + 1 \\
& = & f_x(x) +1.
\end{eqnarray*}
ہر بنیادی بازگشتی فنکشن $f_i$ لئی، $h$ تے
$f_i$ دی قیمت $i$ اُتے وکھری اے۔ ایس لئی
$h$ قابلِ حساب اے، پر بنیادی بازگشتی نہیں؛ تے
$g$ بارے وی ایہو گل کہی جا سکدی اے۔ ایہہ کانتور
دی قطری دلیل دا اک «موثر» روپ اے۔

قابلِ حساب پر بنیادی بازگشتی نہ ہون والے فنکشناں
دیاں ہور صاف مثالاں وی دتیاں جا سکدیاں نیں۔ مثال
لئی، علامت $g^n(x)$ نوں $g(g(\dots g(x)))$
لئی ورتو، جیہدے وچ کل $n$ وار $g$ آؤندا اے؛
تے فنکشناں دا سلسلہ $g_0,g_1,\dots$ ایویں
متعین کرو:
\begin{eqnarray*}
g_0(x) & = & x+1 \\
g_{n + 1}(x) & = & g_n^x(x)
\end{eqnarray*}
تسی جانچ سکدے او کہ ہر فنکشن $g_n$ بنیادی
بازگشتی اے۔ ہر اگلا فنکشن پچھلے توں بہت تیز
ودھدا اے؛ $g_1(x)$، $2x$ دے برابر اے،
$g_2(x)$، $2^x \cdot x$ دے برابر اے، تے
$g_3(x)$ لگ بھگ $x$ واری $2$ نوں اک دوجے
دی قوت بنا کے اُتے اُتے رکھن والے مینار وانگ
ودھدا اے۔ Ackermann--P\'eter فنکشن بنیادی طور
اُتے $G(x) = g_x(x)$ اے، تے دکھایا جا سکدا اے
کہ ایہہ ہر بنیادی بازگشتی فنکشن توں تیز ودھدا اے۔

ہُن بنیادی بازگشتی فنکشناں دی گنتی وار فہرست
بناون دے مسئلے ول واپس آؤ۔ یاد رکھو کہ ہر
بنیادی بازگشتی فنکشن نوں اسی اک علامتی لکھت
دِتّی اے؛ ایس لئی علامتی لکھتاں دی گنتی وار فہرست
بناونا کافی اے۔ اسی قدرتی عدد $\#(F)$ نوں
ہر علامتی لکھت $F$ نال بازگشتی طور اُتے ایویں
جوڑ سکدے آں:
\begin{eqnarray*}
\#(0) & = & \langle 0 \rangle \\
\#(S) & = & \langle 1 \rangle \\
\#(\Proj{n}{i}) & = & \langle 2, n, i \rangle \\
\#(\fn{Comp}_{k,l}[H,G_0,\dots,G_{k-1}]) & = & \langle
3,k,l,\#(H),\#(G_0),\dots,\#(G_{k-1}) \rangle \\
\#(\fn{Rec}_l[G,H]) & = & \langle 4, l, \#(G), \#(H) \rangle
\end{eqnarray*}
ایتھے اسی ایہہ حقیقت ورت رہے آں کہ پچھلے حصے دے
کوڈاں نال عدداں دے ہر سلسلے نوں قدرتی عدد دے
روپ وچ ویکھیا جا سکدا اے۔ نتیجہ ایہہ اے کہ ہر
کوڈ نوں اک قدرتی عدد ملدا اے۔ بے شک، کجھ سلسلے
(تے ایس لئی کجھ عدد) کسے علامتی لکھت دے مطابق
نہیں ہندے؛ پر اسی $f_i$ نوں اوہ اک-جگہ والا
بنیادی بازگشتی فنکشن رکھ سکدے آں جیہدی علامتی
لکھت دا کوڈ $i$ اے، جے $i$ واقعی ایہو جیہی
لکھت دا کوڈ ہووے؛ نہیں تاں مستقل $0$ فنکشن
رکھ لَیندے آں۔ ایس طرح اک-جگہ والے بنیادی
بازگشتی فنکشناں دی گنتی وار فہرست بناون دا اک
واضح طریقہ مل گیا اے۔

(اصل وچ، کجھ فنکشن، جداں مستقل صفر فنکشن، فہرست
وچ اک توں ودھ وار آون گے۔ ایہہ صرف ساڈی کوڈکاری
دی وجہ توں نہیں، بلکہ ایس وجہ توں وی اے کہ
مستقل صفر فنکشن دی اک توں ودھ علامتی لکھت اے۔
اگے اسی ویکھاں گے کہ ایہناں دہرائیاں توں قابلِ
حساب طریقے نال بچیا نہیں جا سکدا؛ مثال لئی، کوئی
قابلِ حساب فنکشن نہیں جیہڑا فیصلہ کرے کہ دِتّی
علامتی لکھت مستقل صفر فنکشن نوں ظاہر کردی اے
یا نہیں۔)

ہُن اسی فنکشن $g(x,y)$ نوں $f_x(y)$ نال
متعین کر سکدے آں، جتھے $f_x$ اوہی گنتی وار
فہرست دا حوالہ اے جیہڑی ہُن دسی اے۔ سانوں کِویں
پتا اے کہ $g(x,y)$ قابلِ حساب اے؟ بدیہی طور اُتے
ایہہ صاف اے: $g(x,y)$ حساب کرن لئی، پہلاں
$x$ دا کوڈ «کھولو» تے ویکھو کہ ایہہ اک-جگہ
والے فنکشن دی علامتی لکھت اے یا نہیں۔ جے اے،
تاں اوہ فنکشن دی قیمت انپٹ~$y$ اُتے حساب کرو۔

\begin{tagblock}{TMs}
\begin{digress}
ممکن اے تسی پہلاں ہی قائل ہوو کہ (کجھ محنت نال!)
ایہہ کم کرن والا پروگرام (مثال لئی، Java یا C++
وچ) لکھیا جا سکدا اے؛ تے ہُن اسی چرچ--ٹیورنگ
دعویٰ \pnbLatin{(Church--Turing thesis)} دا سہارا لے سکدے
آں، جیہڑا کہندا اے کہ بدیہی سمجھ مطابق جیہڑی وی
شے قابلِ حساب اے، اوہ ٹیورنگ مشین نال حساب کیتی
جا سکدی اے۔

بے شک، $g(x,y)$ دے قابلِ حساب ہون نوں دکھاون
دا ہور سِدھا طریقہ ایہہ اے کہ اوہنوں حساب کرن
والی ٹیورنگ مشین صاف صاف بیان کیتی جاوے۔ خاص
طور اُتے ایس نال چرچ--ٹیورنگ دعوے تے بدیہی سمجھ
دا سہارا نہیں لینا پئے گا۔ چھیتی ہی ساڈے کول
اینیاں بناوتاں ہو جان گیاں کہ حساب دے اک نمونے
دا سہارا لے کے $g(x,y)$ دے قابلِ حساب ہون نوں
دکھا سکّیے؛ اوہ نمونہ ایسا اے جیہدے حسابی عمل دی ٹیورنگ مشین اُتے
\emph{نقل کیتی} جا سکدی اے: یعنی بازگشتی فنکشن۔
\end{digress}
\end{tagblock}
\endgroup
% /OLP-0223

% OLP-0224 content/computability/recursive-functions/partial-functions.tex
\begingroup
\def\olfilename{partial-functions}
\expandafter\def\csname ol@sectioncs\endcsname{section}

\olfileid{cmp}{rec}{par}
\olsection{جزوی بازگشتی فنکشن}

بازگشتی فنکشناں دی تعریف دا جواز سمجھّن لئی، غور
کرو کہ قابلِ حساب پر بنیادی بازگشتی نہ ہون والے
فنکشناں دا ساڈا ثبوت اصل وچ ہور وی بہت کجھ ثابت
کردا اے۔ دلیل سادہ سی: اسی صرف ایہہ حقیقت ورتی
کہ فنکشناں $f_0,f_1,\dots$ دی گنتی وار فہرست
ایویں بن سکدی اے کہ $x$ تے $y$ دے فنکشن دے
روپ وچ $f_x(y)$ قابلِ حساب ہووے۔ ایس لئی دلیل
\emph{فنکشناں دی ہر اوہ جماعت اُتے لاگو ہوندی اے
جیہدی فہرست ایس طرح بنائی جا سکدی اے}۔ ایہہ
سانوں مشکل وچ پا دیندا اے: اسی قابلِ حساب فنکشناں
نوں صریح طور اُتے بیان کرنا چاہندے آں؛ پر قابلِ
حساب فنکشناں دے کسے مجموعے دا ایہو جیہا صریح
بیان ساریاں نوں نہیں سمیٹ سکدا!

راہ ایہہ اے کہ \emph{جزوی} فنکشناں نوں وی شامل
کیتا جاوے۔ اسی ویکھاں گے کہ جزوی قابلِ حساب
فنکشناں دی گنتی وار فہرست بناونا \emph{ممکن اے}۔
\iftag{TMs}{ اصل وچ، سانوں کافی حد تک پہلاں ہی
ایہہ پتا اے، کیوں جے ٹیورنگ مشیناں دی فہرست اک
باقاعدہ طریقے نال بنائی جا سکدی اے۔}{} اسی اگے
اپنی قطری دلیل ول واپس آواں گے تے ویکھاں گے کہ
جزوی فنکشن شامل کرن نال اوہ دلیل کیوں نہیں چلدی۔

ہُن سوال ایہہ اے: سارے جزوی بازگشتی فنکشن حاصل
کرن لئی بنیادی بازگشتی فنکشناں وچ کیہ شامل کرنا
پئے گا؟ دو کم کرنے نیں:
\begin{enumerate}
\item بنیادی بازگشتی فنکشناں دی تعریف نوں بدلنا،
تاں کہ جزوی فنکشناں دی وی گنجائش ہووے۔
\item تعریف وچ کجھ \emph{شامل} کرنا، تاں کہ
کجھ نویں جزوی فنکشن وی آ جان۔
\end{enumerate}

پہلا کم سوکھا اے۔ پہلے وانگ، صفر، جانشین تے
پروجیکشن فنکشناں توں شروع کراں گے، تے ترکیب تے
بنیادی بازگشت تحت بندش لواں گے۔ فرق صرف ایہہ
اے کہ ترکیب تے بنیادی بازگشت دی تعریف بدلنی
پئے گی، تاں کہ تعریف وچ آون والے کجھ ٹرم غیر
تعریف شدہ وی ہو سکّن۔ جے $f$ تے $g$ جزوی
فنکشن ہون، تاں $f(x) \fdefined$ دا مطلب ایہہ
لواں گے کہ $f$، $x$ اُتے تعریف شدہ اے، یعنی
$x$، $f$ دے دائرۂ تعریف وچ اے؛ تے $f(x)
\fundefined$ دا اُلٹ مطلب اے، یعنی $f$،~$x$
اُتے تعریف شدہ نہیں۔ $f(x) \simeq g(x)$ دا
مطلب ایہہ لواں گے کہ یا تاں $f(x)$ تے $g(x)$
دونیں غیر تعریف شدہ نیں، یا دونیں تعریف شدہ تے
برابر نیں۔ ہور پیچیدہ ٹرماں لئی وی ایہو علامتاں
ورتاں گے۔ ایہہ طے کر لَیندے آں کہ جے
$h$ تے $g_0$، \dots،~$g_k$ سارے جزوی
فنکشن ہون، تاں
\[
h(g_0(\vec x),\dots,g_k(\vec x))
\]
اُدوں تے صرف اُدوں تعریف شدہ اے جدوں ہر $g_i$،
$\vec x$ اُتے تعریف شدہ ہووے، تے $h$،
$g_0(\vec x)$، \dots،~$g_k(\vec x)$ دلیلاں
دے اِس ٹپل اُتے تعریف شدہ ہووے۔ ایس سمجھ نال جزوی فنکشناں لئی
ترکیب تے بنیادی بازگشت دی تعریف اوپر وانگ ہی اے،
بس «$=$» دی جگہ «$\simeq$» رکھنا پئے گا۔

جزوی فنکشن حاصل کرن لئی بنیادی بازگشتی فنکشناں
دی تعریف وچ جیہڑی شے شامل کراں گے، اوہ
\emph{غیر محدود تلاش دا عامل \pnbLatin{(unbounded search operator)}} اے۔ جے $f(x,\vec z)$ قدرتی عدداں
اُتے کوئی وی جزوی فنکشن ہووے، تاں
$\mu x \; f(x,\vec z)$ دی تعریف ایہہ اے:
\begin{quote}
  اوہ کمترین $x$ جیہدے لئی $f(0,\vec z), f(1,\vec z), \dots, f(x,\vec
  z)$ سارے تعریف شدہ ہون، تے $f(x,\vec z) = 0$
  ہووے، جے ایہو جیہا $x$ موجود ہووے
\end{quote}
تے ایہہ سمجھنا اے کہ ہور ہر صورت وچ
$\mu x \; f(x,\vec z)$ غیر تعریف شدہ اے۔ ایس
نال $\mu x \; f(x,\vec z)$ اکّو طرح متعین
ہوندا اے۔

\begin{explain}
غور کرو کہ ساڈی تعریف \iftag{TMs}{ٹیورنگ
مشیناں، یا }{} الگورتھماں یا حساب دے کسے خاص
نمونے دا حوالہ نہیں دیندی۔ پر ترکیب تے بنیادی
بازگشت وانگ، غیر محدود تلاش دے پچھے وی عمل تے
حساب بارے اک بدیہی سمجھ اے۔ جزوی فنکشن دی
حساب پذیری وچ، جیہڑیاں دلیلاں اُتے فنکشن غیر
تعریف شدہ اے، اوہ انہاں انپٹاں دے مطابق نیں
جیہناں لئی حساب نہیں رکدا۔ $\mu x \;
f(x,\vec z)$ حساب کرن دا طریقہ ایہہ ہووے گا:
$f(0,\vec z), f(1,\vec z),
f(2,\vec z)$ حساب کردے جاؤ، جدوں تک صفر
واپس نہ آوے۔ پر جے وچکار دا کوئی حساب نہ رکے،
تاں $\mu x \; f(x,\vec z)$ دا حساب وی نہیں
رکدا۔
\end{explain}

جے $R(x,\vec z)$ کوئی وی تعلق ہووے، تاں
$\mu x \; R(x,\vec z)$ دی تعریف
$\mu x \; (1 \tsub \Char{R}(x,\vec z))$
رکھدے آں۔ دوجے لفظاں وچ، $\mu x \;
R(x,\vec z)$، $x$ دی اوہ کمترین قیمت واپس
دیندا اے جیہدے لئی $R(x,\vec z)$ پورا اترے۔
ایس لئی جے $f(x,\vec z)$ کُل فنکشن ہووے، تاں
$\mu x \; f(x,\vec
z)$، $\mu x \; (f(x,\vec z) = 0)$ دے
برابر اے۔ پر یاد رکھو کہ ساڈی پہلی تعریف ہور
عام اے، کیوں جے اوہ ایہہ گنجائش وی رکھدی اے
کہ $f(x,\vec z)$ ہر تھاں تعریف شدہ نہ ہووے
(ایہدے اُلٹ، تعلق دا خصوصیتی فنکشن ہمیشہ کُل
ہوندا اے)۔

\begin{defn}
\emph{جزوی بازگشتی فنکشناں} دا سیٹ قدرتی عدداں
توں قدرتی عدداں ول مختلف دلیل-گنتیاں والے جزوی
فنکشناں دا سب توں چھوٹا سیٹ اے، جیہدے وچ صفر،
جانشین تے پروجیکشن فنکشن شامل نیں، تے جیہڑا
ترکیب، بنیادی بازگشت تے غیر محدود تلاش تحت بند اے۔
\end{defn}

بے شک، کجھ جزوی بازگشتی فنکشن کُل وی ہون گے،
یعنی ہر دلیل اُتے تعریف شدہ ہون گے۔

\begin{defn}
\ollabel{defn:recursive-fn}
\emph{بازگشتی فنکشناں} دا سیٹ اوہناں جزوی
بازگشتی فنکشناں دا سیٹ اے جیہڑے کُل نیں۔
\end{defn}

بازگشتی فنکشن نوں کدی کدی «کُل بازگشتی»
\pnbLatin{(total recursive)} وی کہندے نیں، تاں کہ اوہدے
ہر تھاں تعریف شدہ ہون اُتے زور دتا جاوے۔
\endgroup
% /OLP-0224

% OLP-0225 content/computability/recursive-functions/normal-form.tex
\begingroup
\def\olfilename{normal-form}
\expandafter\def\csname ol@sectioncs\endcsname{section}

\olfileid{cmp}{rec}{nft}
\olsection{نارمل روپ دا قضیہ}

\begin{thm}[کلینی دا نارمل روپ دا قضیہ \pnbLatin{(Kleene's Normal Form Theorem)}]
\ollabel{thm:kleene-nf}
اک بنیادی بازگشتی تعلق $T(e, x, s)$ تے اک بنیادی
بازگشتی فنکشن $U(s)$ موجود نیں، جیہناں دی ایہہ
خاصیت اے: جے $f$ کوئی وی جزوی بازگشتی فنکشن
ہووے، تاں کسے~$e$ لئی،
\[
f(x) \simeq U(\umin{s}{T(e, x, s)})
\]
ہر $x$ لئی اے۔
\end{thm}

\begin{explain}
نارمل روپ دے قضیے دا ثبوت پیچیدہ اے، پر بنیادی
خیال سادہ اے۔ ہر جزوی بازگشتی فنکشن دا اک
\emph{اشاریہ (index)}~$e$ ہوندا اے؛ بدیہی طور
اُتے ایہہ اک عدد اے جیہڑا اوہدے پروگرام یا تعریف
دا کوڈ اے۔ جے $f(x)
\fdefined$، تاں حساب دا ریکارڈ باقاعدہ طریقے نال
رکھیا تے کسے عدد~$s$ وچ کوڈ کیتا جا سکدا اے؛
تے ایہہ گل کہ $s$،~$f$ دے اُس حساب دا کوڈ
اے جیہدا انپٹ~$x$ اے، صرف $x$ تے تعریف~$e$ ورت
کے بنیادی بازگشتی طریقے نال چیک کیتی جا سکدی
اے۔ ایس لئی تعلق~$T$، یعنی «اشاریہ~$e$ والے
فنکشن دا انپٹ~$x$ لئی اک حساب اے، تے $s$
اُس حساب دا کوڈ اے»، بنیادی بازگشتی اے۔ حساب
دا پورا ریکارڈ~$s$ دتا ہووے، تاں~$s$ دا
«نتیجہ»~$f(x)$ دی قیمت اے، تے اوہ وی~$s$
توں بنیادی بازگشتی طریقے نال حاصل کیتی جا سکدی اے۔

نارمل روپ دا قضیہ دسدا اے کہ کسے وی جزوی
بازگشتی فنکشن دی تعریف لئی صرف اک غیر محدود
تلاش دی لوڑ اے۔ بنیادی طور اُتے اسی سارے عدداں
وچ تلاش کر سکدے آں، جدوں تک اک ایہو جیہا عدد
نہ لبھے جیہڑا اشاریہ~$e$ والے فنکشن دے
انپٹ~$x$ لئی حساب دا کوڈ ہووے۔ اسی عدداں~$e$
نوں جزوی بازگشتی فنکشناں دے «ناں» دے طور اُتے
ورت سکدے آں، تے $\cfind{e}$ اوس فنکشن~$f$
لئی لکھ سکدے آں جیہڑا قضیے دی مساوات نال
تعریف ہوندا اے۔ یاد رکھو کہ کسے وی جزوی بازگشتی
فنکشن دے اک توں ودھ اشاریے ہو سکدے نیں—اصل
وچ، ہر جزوی بازگشتی فنکشن دے لامتناہی اشاریے نیں۔
\end{explain}
\endgroup
% /OLP-0225

% OLP-0226 content/computability/recursive-functions/halting-problem.tex
\begingroup
\def\olfilename{halting-problem}
\expandafter\def\csname ol@sectioncs\endcsname{section}

\olfileid{cmp}{rec}{hlt}
\olsection{حساب رُکّن دا مسئلہ}

عام طور اُتے \emph{حساب رُکّن دا مسئلہ \pnbLatin{(halting problem)}} ایہہ فیصلہ کرنے دا مسئلہ اے کہ قابلِ
حساب فنکشن دا بیان~$e$ (مثلاً پروگرام) تے اک
عدد~$n$ دتے ہون، تاں کی انپٹ~$n$ اُتے اوس
فنکشن دا حساب رُکدا اے، یعنی کوئی نتیجہ دیندا
اے۔ ایلن ٹیورنگ دا مشہور ثبوت اے کہ ایہہ مسئلہ
خود کسے قابلِ حساب فنکشن نال حل نہیں ہو سکدا؛
یعنی فنکشن
\[
h(e, n) =
\begin{cases}
1 & \text{جے حساب $e$ انپٹ $n$ اُتے رُکے}\\
0 & \text{ہور صورت وچ،}
\end{cases}
\]
قابلِ حساب نہیں۔

جزوی بازگشتی فنکشناں دے سیاق وچ، پروگرام دے
بیان دی جگہ کلینی دے نارمل روپ دے قضیے وچ
دتا اشاریہ~$e$ آ سکدا اے۔ جے $f$ اک جزوی
بازگشتی فنکشن ہووے، تاں ہر اوہ $e$~جیہدے لئی
نارمل روپ دے قضیے دی مساوات پوری اُترے،~$f$
دا اشاریہ اے۔ اک عدد~$e$ دتا ہووے، تاں نارمل
روپ دا قضیہ دسدا اے کہ
\[
\cfind{e}(x) \simeq U(\mu s \; T(e, x, s))
\]
جزوی بازگشتی اے، تے ہر جزوی بازگشتی
$f\colon \Nat
\to \Nat$ لئی اک $e \in \Nat$ موجود اے
جیہدے لئی $\cfind{e}(x) \simeq
f(x)$ ہر~$x \in \Nat$ اُتے اے۔ اصل وچ،
ہر ایہو جیہے $f$ لئی اک نہیں، بلکہ لامتناہی
ایہو جیہے~$e$ نیں۔ \emph{حساب رُکّن دی خبر
دین والے فنکشن \pnbLatin{(halting function)}}~$h$ دی
تعریف ایہہ اے:
\[
h(e, x) =
\begin{cases}
1 & \text{جے $\cfind{e}(x) \fdefined$}\\
0 & \text{ہور صورت وچ۔}
\end{cases}
\]
یاد رکھو کہ $h(e, x) = 0$ جے
$\cfind{e}(x) \fundefined$، پر اُدوں وی
جدوں~$e$ کسے جزوی بازگشتی فنکشن دا اشاریہ
ہی نہ ہووے۔

\begin{thm}
\ollabel{thm:halting-problem}
حساب رُکّن دی خبر دین والا فنکشن~$h$ جزوی
بازگشتی نہیں۔
\end{thm}

\begin{proof}
جے $h$
جزوی بازگشتی ہوندا، تاں اسی ایہہ تعریف کر سکدے:
\[
d(y) =
\begin{cases}
1 & \text{جے $h(y, y) = 0$}\\
\umin{x}{x \neq x} & \text{ہور صورت وچ۔}
\end{cases}
\]
کیوں جے کوئی وی عدد $x$، $x \neq x$ پورا
نہیں کردا، ایس لئی کوئی $\umin{x}{x \neq
x}$ نہیں، تے $d(y) \fundefined$ اُدوں
تے صرف اُدوں جدوں $h(y,y) \neq 0$۔ ایس
تعریف توں ایہہ نتیجے نکلدے نیں:
\begin{enumerate}
\item $d(y) \fdefined$ اُدوں تے صرف اُدوں
جدوں $\cfind{y}(y) \fundefined$ یا $y$
کسے جزوی بازگشتی فنکشن دا اشاریہ نہ ہووے۔
\item $d(y) \fundefined$ اُدوں تے صرف اُدوں
جدوں $\cfind{y}(y) \fdefined$۔
\end{enumerate}
جے $h$ جزوی بازگشتی ہوندا، تاں $d$ وی جزوی
بازگشتی ہوندا۔ ایس لئی کلینی دے نارمل روپ دے
قضیے مطابق، اوہدا اک اشاریہ~$e_d$ ہوندا۔
$h(e_d, e_d)$ دی قیمت ویکھو۔ دو ممکن
صورتاں نیں: $0$ تے~$1$۔
\begin{enumerate}
\item جے $h(e_d, e_d) = 1$، تاں
$\cfind{e_d}(e_d) \fdefined$۔ پر
$\cfind{e_d} \simeq d$، تے $d(e_d)$
اُدوں تے صرف اُدوں تعریف شدہ اے جدوں
$h(e_d, e_d) = 0$۔ ایس لئی
$h(e_d, e_d) \neq 1$۔
\item جے $h(e_d, e_d) = 0$، تاں یا
$e_d$ کسے جزوی بازگشتی فنکشن دا اشاریہ
نہیں، یا اوہ اشاریہ اے تے $\cfind{e_d}(e_d)
\fundefined$۔ پر فیر وی $\cfind{e_d} \simeq d$،
تے $d(e_d)$ اُدوں تے صرف اُدوں غیر تعریف
شدہ اے جدوں $\cfind{e_d}(e_d) \fdefined$۔
\end{enumerate}
نتیجہ ایہہ اے کہ آخرکار $e_d$ کسے جزوی
بازگشتی فنکشن دا اشاریہ نہیں ہو سکدا۔ پر جے
$h$ جزوی بازگشتی ہوندا، تاں $d$ وی ہوندا؛
ایس لئی اوہدے اشاریے دے روپ وچ $e_d$ دی
ساڈی تعریف جائز ہوندی۔ لازماً ایہہ نتیجہ کڈھنا
پَینا اے کہ $h$ جزوی بازگشتی نہیں ہو سکدا۔
\end{proof}
\endgroup
% /OLP-0226

% OLP-0227 content/computability/recursive-functions/general-recursive-functions.tex
\begingroup
\def\olfilename{general-recursive-functions}
\expandafter\def\csname ol@sectioncs\endcsname{section}

\olfileid{cmp}{rec}{gen}
\olsection{عمومی بازگشتی فنکشن}

کُل فنکشناں دا اک سیٹ حاصل کرن دا ہور وی
طریقہ اے۔ اک کُل فنکشن $f(x,\vec z)$ نوں
\emph{باقاعدہ (regular)} کہیئے جے قدرتی عدداں
دے ہر سلسلے $\vec z$ لئی اک $x$ موجود
ہووے جیہدے لئی $f(x,\vec z) = 0$۔ دوجے
لفظاں وچ، باقاعدہ فنکشن عین اوہ فنکشن نیں
جیہناں اُتے غیر محدود تلاش لا کے آخرکار کُل
فنکشن حاصل ہووے۔ احتیاط نال غیر محدود تلاش
نوں صرف باقاعدہ فنکشناں تک محدود رکھ سکدے آں:

\begin{defn}
\ollabel{defn:general-recursive}
\emph{عمومی بازگشتی فنکشناں \pnbLatin{(general recursive functions)}} دا سیٹ قدرتی عدداں توں قدرتی
عدداں ول مختلف دلیل-گنتیاں والے فنکشناں دا
سب توں چھوٹا سیٹ اے، جیہدے وچ صفر، جانشین
تے پروجیکشن فنکشن شامل نیں، تے جیہڑا ترکیب،
بنیادی بازگشت تے \emph{باقاعدہ} فنکشناں اُتے
غیر محدود تلاش لاون تحت بند اے۔
\end{defn}

صاف اے کہ ہر عمومی بازگشتی فنکشن کُل اے۔
\olref{defn:general-recursive} تے
\olref[par]{defn:recursive-fn} دا فرق ایہہ اے
کہ دوجی تعریف وچ راہ وچ جزوی بازگشتی فنکشن
ورتن دی اجازت اے؛ شرط صرف ایہہ اے کہ آخر
وچ حاصل ہون والا فنکشن کُل ہووے۔ ایس لئی لفظ
«عمومی»، جیہڑا تاریخ توں بچیا ہویا ناں اے،
معنی نال نہیں ملدا: ظاہری طور اُتے
\olref{defn:general-recursive}،
\olref[par]{defn:recursive-fn} نالوں \emph{گھٹ}
عمومی اے۔ پر خوش قسمتی نال ایہہ فرق صرف
ظاہری اے؛ تعریفاں وکھ نیں، پر عمومی بازگشتی
فنکشناں دا سیٹ تے بازگشتی فنکشناں دا سیٹ
اکّو نیں۔
\endgroup
% /OLP-0227

\OLEndChapterHook
\endgroup
% /OLP-0209

% OLP-0228 content/computability/computability-theory/computability-theory.tex
\begingroup
\def\olfilename{computability-theory}
\expandafter\def\csname ol@sectioncs\endcsname{section}

\olchapter{cmp}{thy}{قابلِ حسابیت دا نظریہ}

\begin{editorial}
  ایس باب دے مواد اُتے نظرثانی ہونی چاہیدی اے
  تے ایہنوں ہور ودھایا جانا چاہیدا اے۔ ایس باب وچ
  مشقاں وی موجود نیں۔
\end{editorial}

% OLP-0229 content/computability/computability-theory/introduction.tex
\begingroup
\def\olfilename{introduction}
\expandafter\def\csname ol@sectioncs\endcsname{section}

\olfileid{cmp}{thy}{int}
\olsection{تعارف}

منطق دی جیہڑی شاخ \emph{قابلِ حسابیت دا نظریہ} کہلاندی اے،
اوہ فنکشناں تے سیٹاں دی قابلِ حسابیت یا نسبتی قابلِ حسابیت
\pnbLatin{(relative computability)} نال متعلق سوالاں نوں گھوکھدی اے۔
کلینی دے اثر دی اک نشانی ایہ اے کہ پہلے ایس موضوع دا ناں
\emph{بازگشت دا نظریہ} سی، تے اج وی دونوں ناں عام ورتے جاندے نیں۔

آؤ فنکشن~$f\colon \Nat \pto \Nat$ نوں \emph{جزوی قابلِ حساب
  \pnbLatin{(partial computable)}} آکھئیے جے اوہنوں حساب دے کسے نمونے وچ
حسابیا جا سکدا ہووے۔ جے $f$ کُل اے، تاں اسی صرف ایہ آکھاں گے
کہ $f$~\emph{قابلِ حساب} اے۔ اک تعلق~$R$ وی قابلِ حساب کہلاندا اے
جے اوہدا خصوصیتی فنکشن~$\Char{R}$ قابلِ حساب ہووے۔ جے $f$
تے~$g$ جزوی فنکشن نیں، تاں اسی $f(x) \fdefined$ ایہ دسّن لئی
لکھاں گے کہ $f$،~$x$ اُتے تعریف شدہ اے، یعنی $x$،~$f$ دے
دائرۂ تعریف وچ اے؛ تے $f(x) \fundefined$ ایس دے اُلٹ معنی لئی،
یعنی ایہ دسّن لئی کہ $f$،~$x$ اُتے تعریف شدہ نہیں۔ اسی
$f(x) \simeq g(x)$ دا مطلب ایہ لاں گے کہ یا تاں $f(x)$ تے
$g(x)$ دونوں غیر تعریف شدہ نیں، یا دونوں تعریف شدہ تے برابر نیں۔

حساب دے کسے خاص نمونے دا حوالہ دِتے بغیر وی قابلِ حسابیت دے
نظریے نوں گھوکھیا جا سکدا اے۔ ایس لئی ایہ وکھایا جاندا اے کہ
اک ہمہ گیر جزوی قابلِ حساب فنکشن~$\fn{Un}(k, x)$ موجود اے۔
ایہ ساڈے لئی جزوی قابلِ حساب فنکشناں دی فہرست بنانا ممکن کردا
اے۔ اسی علامت~$\cfind{k}$ نال $k$-ویں اک آرگومنٹ والے جزوی
قابلِ حساب فنکشن نوں ظاہر کراں گے، جیہدی تعریف
$\cfind{k}(x) \simeq \fn{Un}(k, x)$ اے۔ (کلینی ایس مقصد لئی
$\{ k \}$ ورتدا سی، پر حالیہ وقتاں وچ ایہ علامت اینی عام نہیں
ورتی گئی۔) تھوڑا ہور عام طور اُتے، اسی کسے وی تعداد دے
آرگومنٹاں والے جزوی قابلِ حساب فنکشناں دی فہرست اکّو یکساں
طریقے نال بنا سکدے آں، تے $\cfind{k}[n]$ نال $k$-ویں،
$n$ آرگومنٹاں والے جزوی بازگشتی فنکشن نوں ظاہر کراں گے۔

جے $f(\vec x, y)$ اک کُل یا جزوی فنکشن اے، تاں $\umin{y}{f
(\vec x, y)}$،~$\vec x$ دا اوہ فنکشن اے جیہڑا سب توں
چھوٹا~$y$ واپس دیندا اے جیہدے لئی $f(\vec x, y) = 0$
ہووے، بشرطیکہ $f(\vec x, 0)$، \dots، $f(\vec x, y-1)$
سارے تعریف شدہ ہون؛ جے ایہوجیہا کوئی $y$ نہیں، تاں
$\umin{y}{f (\vec x, y)}$ غیر تعریف شدہ اے۔ جے $R(\vec x, y)$
اک تعلق اے، تاں $\umin{y}{R(\vec x, y)}$ دی تعریف سب توں
چھوٹے~$y$ دے طور اُتے کیتی جاندی اے جیہدے لئی $R(\vec x, y)$
سچ ہووے؛ ہور لفظاں وچ، سب توں چھوٹا~$y$ جیہدے لئی
$1 \tsub \Char{R}(\vec x, y) = 0$ ہووے۔

ایہ وکھاؤن لئی کہ اک فنکشن قابلِ حساب اے، دو طریقیاں نال
اگے ودھیا جا سکدا اے:
\begin{enumerate}
\item سخت ریاضیاتی طریقے نال: اک ٹیورنگ مشین یا جزوی بازگشتی
  فنکشن نوں صاف صاف بیان کرو، تے وکھاؤ کہ اوہ تہاڈے مطلوبہ
  فنکشن دا حساب کردا اے؛
\item غیر رسمی طریقے نال: ایہدا حساب کرن والا اک الگورتھم بیان
  کرو، تے چرچ دے دعوے \pnbLatin{(Church's thesis)} دا سہارا لو۔
\end{enumerate}
دوہاں طریقیاں وچ کوئی تِکھّی حد نہیں؛ الگورتھم دی تفصیلی وضاحت
وچ اینی معلومات ہونی چاہیدی اے کہ ایہ نسبتاً صاف ہووے کہ اصولی
طور اُتے مطلوبہ ٹیورنگ مشین یا جزوی بازگشتی تعریفاں دا سلسلہ
کیویں بنایا جا سکدا اے۔ پوری ریاضیاتی سختی والیاں تعریفاں توں
وڈی وضاحت ملّن دی امید گھٹ اے، تے اسی ایہناں دو طریقیاں وچ
اک مناسب درمیانی راہ لبھن دی کوشش کراں گے؛ مُکدی گل، اسی
ایہوجیہے ثبوت لبھن دی کوشش کراں گے جیہڑے سوجھ نال سمجھ آون
تے ریاضیاتی طور اُتے سخت وی ہون، تے وکھاؤن کہ ٹھیک
ٹھیک تعریفاں حاصل کیتیاں جا سکدیاں نیں۔
\endgroup
% /OLP-0229

% OLP-0230 content/computability/computability-theory/coding-computations.tex
\begingroup
\def\olfilename{coding-computations}
\expandafter\def\csname ol@sectioncs\endcsname{section}

\olfileid{cmp}{thy}{cod}
\olsection{حساباں نوں کوڈ دینا}

حساب دے ہر نمونے وچ ایہ کم کیتے جا سکدے نیں:
\begin{enumerate}
\item قابلِ حساب فنکشناں دیاں \emph{تعریفاں} نوں اک منظم طریقے
  نال بیان کرنا۔ مثلاً ٹیورنگ مشین دیاں تعییناں، بازگشتی تعریفاں
  یا کسے پروگرامنگ زبان دے پروگرام ایہ تعریفاں دے سکدے نیں۔
\item کسے تعریف نال دِتے فنکشن دے کسے دِتے انپٹ اُتے حساب دا
  پورا ریکارڈ بیان کرنا۔ مثلاً ٹیورنگ مشین دا حساب اوہناں
  ترتیب وار حالتاں دے سلسلے نال بیان ہو سکدا اے جیہڑیاں حساب
  دے ہر قدم اُتے ہوندیاں نیں: مشین دی حالت تے ٹیپ دا مواد۔
\item ایہ پرکھنا کہ حساب دا دِتا ہویا ممکنہ ریکارڈ واقعی اوہ
  ریکارڈ اے جیہدے مطابق دِتی تعریف والا قابلِ حساب فنکشن دِتے
  انپٹ اُتے حسابیا جاندا اے؛ ایس انپٹ اُتے فنکشن تعریف شدہ
  ہوندا اے، یعنی حساب رُک جاندا اے۔
\item حساب دے پورے ریکارڈ دے ایہوجیہے بیان توں دِتے انپٹ لئی
  فنکشن دی قدر کڈھنا۔ مثلاً ٹیورنگ مشین دے رُکّن والے حساب
  دے بالکل آخری قدم اُتے ٹیپ دا مواد ایہ قدر ہوندا اے۔
\end{enumerate}

کوڈنگ راہیں قابلِ حساب فنکشن دے ہر بیان نوں اک عددی
\emph{اشاریہ (index)} دِتا جا سکدا اے، تے ایہ ایداں کیتا
جا سکدا اے کہ اشاریے توں ہدایتاں اک قابلِ حساب طریقے نال
واپس کڈھیاں جا سکّن۔ ایس طرح حساب دے پورے ریکارڈ نوں وی
اکّو عدد دے کوڈ وچ رکھیا جا سکدا اے۔ ایس توں بنن والا حسابی
تعلق ``$s$، اشاریہ~$e$ والے فنکشن دے انپٹ~$x$ اُتے حساب
دے ریکارڈ دا کوڈ اے'' تے فنکشن ``کوڈ~$s$ والے حسابی سلسلے
دا نتیجہ'' دونوں قابلِ حساب ہوندے نیں؛ بلکہ اوہ بنیادی
بازگشتی نیں۔

ایہ بنیادی حقیقت بڑی طاقتور اے، تے ایس دی مدد نال اسی
قابلِ حسابیت بارے کئی حیران کن تے اہم نتیجے ثابت کر سکدے
آں، بھاویں حساب دا جیہڑا وی نمونہ چنیا گیا ہووے۔
\endgroup
% /OLP-0230

% OLP-0231 content/computability/computability-theory/normal-form.tex
\begingroup
\def\olfilename{normal-form}
\expandafter\def\csname ol@sectioncs\endcsname{section}

\olfileid{cmp}{thy}{nfm}
\olsection{نارمل روپ دا قضیہ}

فرض کرو اسی قابلِ حساب فنکشناں دیاں تعریفاں بیان کر سکدے
آں، تے پرکھ سکدے آں کہ کسے انپٹ اُتے اوس فنکشن دا حساب کرن دے
پورے ریکارڈ دا دِتا ہویا ممکنہ بیان درست اے یا نہیں۔ فیر
ایہ گل معقول لگدی اے کہ حساب دے نمونے توں آزاد رہ کے، اسی
کسے وی قابلِ حساب فنکشن~$f$ دی کسے وی انپٹ~$x$ اُتے قدر
ایویں لبھ سکدے آں:
\begin{enumerate}
  \item حسابی ریکارڈاں دے سارے ممکنہ بیان تلاش کرو۔
  \item پرکھو کہ دِتا ریکارڈ $f(x)$ دے حساب دا ریکارڈ اے یا نہیں۔
  \item جے درست ریکارڈ لبھ جائے، تاں اوہدے وچوں $f(x)$ دی قدر کڈھو۔
\end{enumerate}
ایہ گل واقعی درست اے: کلینی دے نارمل روپ دے قضیے دا مضمون
ایہو اے۔

\begin{thm}[کلینی دے نارمل روپ دا قضیہ]
\ollabel{thm:normal-form}
اک بنیادی بازگشتی تعلق~$T(e, x, s)$ تے اک بنیادی بازگشتی
فنکشن~$U(s)$ موجود نیں، جیہناں دی ایہ خاصیت اے: جے $f$
کوئی وی جزوی قابلِ حساب فنکشن اے، تاں کسے~$e$ لئی
\[
f(x) \simeq U(\umin{s}{T(e, x, s)})
\]
ہر~$x$ اُتے۔
\end{thm}

\begin{proof}[ثبوت دا خاکہ]
حساب دے کسے وی نمونے وچ قابلِ حساب فنکشن~$f$ دا بیان پوری
ریاضیاتی سختی نال تعریف کیتا جا سکدا اے، تے ایس بیان نوں
اک قدرتی عدد~$e$ دا کوڈ دِتا جا سکدا اے۔ ``حسابی سلسلہ''
دا تصور وی پوری ریاضیاتی سختی نال تعریف کیتا جا سکدا اے،
جیہڑا اشاریہ~$e$ والے فنکشن دے انپٹ~$x$ اُتے حساب دا
عمل درج کردا اے۔ ایہوجیہے حسابی سلسلے نوں وی اک عدد~$s$
دا کوڈ دِتا جا سکدا اے۔ ایہ کم ایداں ہو سکدا اے کہ
\begin{enumerate}
  \item تعلق $T(e, x, s)$، جیہڑا صرف تے صرف اُدوں قائم ہوندا اے جدوں
  عدد~$s$، اشاریہ~$e$ والے فنکشن دے انپٹ~$x$ اُتے حسابی
  سلسلے دا کوڈ ہووے، تے
  \item فنکشن $U(s)$، جیہڑا کوڈ~$s$ والے حسابی سلسلے نوں
  اوس حساب دے آخری نتیجے ول نقشہ بند کردا اے،
\end{enumerate}
دونوں قابلِ حساب ہون۔ بلکہ تعلق~$T$ تے فنکشن~$U$ بنیادی
بازگشتی نیں۔
\end{proof}

\begin{explain}
نارمل روپ دے قضیے دا پوری ریاضیاتی سختی نال ثبوت دین لئی
سانوں حساب دا اک نمونہ مقرر کرنا، قابلِ حساب فنکشناں دے
بیاناں تے حسابی سلسلیاں دی کوڈنگ پوری تفصیل نال کرنا، تے
ایہ تصدیق کرنا پوے گا کہ تعلق~$T$ تے فنکشن~$U$ بنیادی
بازگشتی نیں۔ زیادہ تر اطلاقاں لئی ایہ کافی اے کہ $T$
تے~$U$ قابلِ حساب ہون تے $U$~کُل ہووے۔

نارمل روپ دے قضیے دا ثبوت شاید اک مختصر یاد رکھن والی گل
دے روپ وچ چیتے رکھنا سب توں چنگا اے: $\umin{s}{T(e, x, s)}$
اشاریہ~$e$ والے فنکشن دے انپٹ~$x$ اُتے حسابی سلسلے نوں
لبھدا اے، تے جے کوئی ایہوجیہا سلسلہ لبھ جائے تاں $U$
اوہدا نتیجہ واپس دیندا اے۔
\end{explain}

جے حساب دا نمونہ جزوی بازگشتی فنکشناں والا ہووے، جیہڑا
کلینی نے شروع وچ ورتیا سی، تاں ایہ نتیجہ دسّدا اے کہ
کسے وی فنکشن دی تعریف لئی غیر محدود تلاش صرف اک واری
ورتی جان دی لوڑ اے، یعنی اکّو $\umin{y}{f(\vec x, y)}$
آپریٹر کافی اے۔ ایس معنی وچ ایہ وکھاندا اے کہ ہر جزوی
بازگشتی فنکشن دا اک نارمل روپ اے۔

$T$ تے $U$ نال فہرست $\cfind{0}$، $\cfind{1}$، $\cfind{2}$،
\dots دی تعریف کیتی جا سکدی اے۔ ہن توں اسی فرض کراں گے
کہ $T$ تے~$U$ دی اک مناسب چُون مقرر کر لئی اے، تے مساوات
\[
\cfind{e}(x) \simeq U(\umin{s}{T(e,x,s)})
\]
نوں $\cfind{e}$ دی \emph{تعریف} منّاں گے۔

ایہ اک ہور کم دی حقیقت اے:

\begin{thm}
ہر جزوی قابلِ حساب فنکشن دے لامتناہی اشاریے نیں۔
\end{thm}

ایہ گل وی سوجھ نال صاف سمجھ آؤندی اے۔ کسے قابلِ حساب
فنکشن دا بیان ملے، تاں اک ہور بیان بنایا جا سکدا اے
جیہڑا اوہو فنکشن، یعنی اوہو انپٹ--آؤٹ پٹ جوڑا، حسابے،
پر تھوڑے وکھرے طریقے نال۔ مثلاً اوہ پہلے کوئی ایہوجیہا
کم کرے جیہدا حساب اُتے کوئی اثر نہ ہووے: پرکھے کہ
$0 = 0$ اے، یا $5$ تک گنے، وغیرہ۔ بدلے ہوئے بیان دا
اشاریہ اصل اشاریے توں ہمیشہ وکھرا ہووے گا۔ دونوں اوہو
فنکشن دے اشاریے نیں؛ بس اوہدا حساب تھوڑے وکھرے طریقے
نال کیتا گیا اے۔
\endgroup
% /OLP-0231

% OLP-0232 content/computability/computability-theory/s-m-n.tex
\begingroup
\def\olfilename{s-m-n}
\expandafter\def\csname ol@sectioncs\endcsname{section}

\olfileid{cmp}{thy}{smn}
\olsection{\texorpdfstring{$s$-$m$-$n$ دا قضیہ}{s-m-n دا قضیہ}}

\begin{explain}
اگلا قضیہ ``$s$-$m$-$n$ دا قضیہ'' کہلاندا اے؛ ایس ناں دی
وجہ ہُنّے صاف ہو جائے گی۔ اوکھا کم ایہ سمجھنا اے کہ
قضیہ ٹھیک ٹھیک کیہ کہندا اے؛ جے بیان سمجھ آ جائے، تاں
نتیجہ خاصا ظاہر لگے گا۔
\end{explain}

\begin{thm}
\ollabel{thm:s-m-n}
  قدرتی عدداں دے ہر جوڑے $n$ تے~$m$ لئی اک بنیادی
  بازگشتی فنکشن~$s^m_n$ موجود اے، ایداں کہ ہر سلسلے
  $e$، $a_0$، \dots، $a_{m-1}$، $y_0$، \dots، $y_{n-1}$ لئی
  \[
  \cfind{s^m_n(e, a_0, \dots, a_{m-1})}[n](y_0, \dots, y_{n-1}) \simeq
  \cfind{e}[m+n](a_0, \dots, a_{m-1}, y_0, \dots, y_{n-1}).
\]
\end{thm}

\begin{explain}
% PNB-CMP-011: source B7 دے دو explanatory x پروگرام دے اشاریہ e نال بدلے؛ theorem unchanged۔
$s^m_n$ نوں \emph{پروگراماں} اُتے کم کرن والا سمجھنا مددگار
اے۔ یعنی $s^m_n$ اک پروگرام~$e$ لیندا اے، جیہڑا $(m+n)$
آرگومنٹاں والے فنکشن لئی اے، تے نال مقرر انپٹ $a_0$،
\dots، $a_{m-1}$ وی؛ تے اوہ پروگرام
$s^m_n(e, a_0, \dots, a_{m-1})$ واپس دیندا اے، جیہڑا باقی
آرگومنٹاں دے $n$-آرگومنٹ والے فنکشن لئی اے۔ \iftag{TMs}{جے
تسی $e$ نوں اک ٹیورنگ مشین دا بیان سمجھو، تاں
$s^m_n(e, a_0, \dots, a_{m-1})$ اوہ ٹیورنگ مشین اے جیہڑی
انپٹ $y_0$، \dots،~$y_{n-1}$ ملّن اُتے، $a_0$، \dots،~$a_{m-1}$
نوں انپٹ لڑی دے مُڈھ وچ لا دیندی اے، تے~$e$ چلا دیندی
اے۔ فیر ہر $s^m_n$ صرف اک بنیادی بازگشتی فنکشن اے جیہڑا
مناسب ٹیورنگ مشین دا کوڈ لبھدا اے۔}{}
\end{explain}
\endgroup
% /OLP-0232

% OLP-0233 content/computability/computability-theory/universal-part-function.tex
\begingroup
\def\olfilename{universal-part-function}
\expandafter\def\csname ol@sectioncs\endcsname{section}

\olfileid{cmp}{thy}{uni}
\olsection{ہمہ گیر جزوی قابلِ حساب فنکشن}

\begin{thm}
\ollabel{thm:univ-comp}
اک ہمہ گیر جزوی قابلِ حساب فنکشن~$\fn{Un}(e,x)$ موجود اے۔
ہور لفظاں وچ، اک فنکشن $\fn{Un}(e,x)$ موجود اے جیہدے لئی:
\begin{enumerate}
\item $\fn{Un}(e,x)$ جزوی قابلِ حساب اے۔
\item جے $f(x)$ کوئی وی جزوی قابلِ حساب فنکشن اے، تاں
اک قدرتی عدد $e$ موجود اے ایداں کہ $f(x) \simeq \fn{Un}(e,x)$
ہر~$x$ لئی۔
\end{enumerate}
\end{thm}

\begin{proof}
$\fn{Un}(e,x) \simeq U(\umin{s}{T(e,x,s)})$ رکھو، جتھے $U$ تے
$T$ کلینی دے نارمل روپ دے قضیے (\olref[nfm]{thm:normal-form})
والے نیں۔
\end{proof}

\begin{explain}
ایہ صرف ایس گل دا ٹھیک ٹھیک بیان اے کہ ساڈے کول جزوی
قابلِ حساب فنکشناں دی اک موثر گنتی وار فہرست اے۔ خیال
ایہ اے کہ جے اسی $f_e$ اوس فنکشن لئی لکھئیے جیہدی تعریف
$f_e(x) = \fn{Un}(e,x)$ اے، تاں سلسلے $f_0$، $f_1$، $f_2$،
\dots وچ سارے جزوی قابلِ حساب فنکشن شامل نیں، تے خاصیت
ایہ اے کہ $f_e(x)$ نوں $e$ تے~$x$ دے لحاظ نال ``یکساں''
طریقے اُتے حسابیا جا سکدا اے۔ سادگی لئی اسی دو آرگومنٹاں
والا فنکشن ورت رہے آں جیہڑا اک آرگومنٹ والے فنکشناں لئی
ہمہ گیر اے، پر عدداں دے سلسلیاں دی کوڈنگ نال ایہنوں
ہور ودھ آرگومنٹاں تک آسانی نال عام کیتا جا سکدا اے۔
% PNB-CMP-012: اصل final g نوں recursive کہندا؛ partial f توں صرف partial g لازم،total نہیں۔
مثلاً، جے $f(x,y,z)$ اک $3$-آرگومنٹ والا جزوی بازگشتی
فنکشن اے، تاں فنکشن $g(x) \simeq f((x)_0, (x)_1, (x)_2)$
اک آرگومنٹ والا جزوی بازگشتی فنکشن اے۔
\end{explain}
\endgroup
% /OLP-0233

% OLP-0234 content/computability/computability-theory/no-universal-function.tex
\begingroup
\def\olfilename{no-universal-function}
\expandafter\def\csname ol@sectioncs\endcsname{section}

\olfileid{cmp}{thy}{nou}
\olsection{کوئی ہمہ گیر کُل قابلِ حساب فنکشن نہیں}

% PNB-CMP-013: اصل 'total for the partial computable functions'،پچھلے قضیے تے اگلے تقابل مطابق 'universal for' اے،کُل فنکشن نہیں۔
اگرچہ اک جزوی قابلِ حساب فنکشن موجود اے جیہڑا سارے جزوی
قابلِ حساب فنکشناں لئی ہمہ گیر اے، پر کوئی کُل قابلِ حساب
فنکشن نہیں جیہڑا سارے کُل قابلِ حساب فنکشناں لئی ہمہ گیر ہووے۔

\begin{thm}
\ollabel{thm:no-univ}
کوئی ہمہ گیر قابلِ حساب فنکشن موجود نہیں۔ ہور لفظاں وچ،
کوئی وی فنکشن $\fn{Un}'(k, x)$ جیہدے بارے ایہ ہووے کہ جے
$f(x)$ اک کُل قابلِ حساب فنکشن اے، تاں اک قدرتی عدد~$k$
موجود ہووے ایداں کہ $f(x) = \fn{Un}'(k,x)$ ہر~$x$ لئی،
اوہ قابلِ حساب نہیں ہو سکدا۔
\end{thm}

\begin{proof}
ثبوت اک سادہ قطری دلیل اے: جے $\fn{Un}'(k,x)$ کُل تے
قابلِ حساب ہووے، تاں
\[
d(x) = \fn{Un}'(x, x) + 1
\]
وی کُل تے قابلِ حساب ہووے گا۔ پر تعریف مطابق، $d(k)$،
$\fn{Un}'(k,k)$ دے برابر نہیں۔ ایس لئی ہر $k$ لئی،
$d(x)$ تے~$\fn{Un}'(k, x)$ دیاں قدراں گھٹ توں گھٹ
اک~$x$ اُتے وکھریاں نیں، یعنی $x = k$ اُتے۔
\end{proof}

\begin{explain}
اُتّے \olref[uni]{thm:univ-comp} وکھاندا اے کہ ایس قطری
دلیل توں بچیا جا سکدا اے، پر صرف ایس قیمت اُتے کہ
ہمہ گیر فنکشن نوں جزوی ہون دِتا جائے۔ یعنی $\fn{Un}$
کُل قابلِ حساب فنکشناں لئی ہمہ گیر اے، بس اوہ خود کُل
نہیں۔ جزوی صورت وچ قطری دلیل کم نہیں کردی۔
\end{explain}

\begin{prob}
  ایہ سمجھّن لئی کہ \olref{thm:no-univ} دے ثبوت والی
  قطری دلیل جزوی صورت وچ کیوں کم نہیں کردی، فنکشن
  $f(x) \simeq \fn{Un}(x,x)+1$ اُتے غور کرو۔ کیہ ایہ
  جزوی قابلِ حساب اے؟ جے اے، تاں ایس دا اک اشاریہ~$e$
  اے، یعنی $f(x) \simeq \fn{Un}(e,x)$۔ فیر~$f(e)$
  بارے تسی کیہ کہہ سکدے او؟
\end{prob}
\endgroup
% /OLP-0234

% OLP-0235 content/computability/computability-theory/halting-problem.tex
\begingroup
\def\olfilename{halting-problem}
\expandafter\def\csname ol@sectioncs\endcsname{section}

\olfileid{cmp}{thy}{hlt}
\olsection{حساب رُکّن دا مسئلہ}

تعمیر مطابق، ہمہ گیر جزوی قابلِ حساب فنکشن~$\fn{Un}(e,x)$
اُدوں تے صرف اُدوں تعریف شدہ اے جدوں~$e$ دے کوڈ والا فنکشن
انپٹ~$x$ اُتے حساب کرکے کوئی قدر دیوے۔ فطری سوال ایہ اے کہ
کیہ اسی فیصلہ کر سکدے آں کہ ایہ حالت اے یا نہیں۔ اصل وچ
ایہ فیصلہ نہیں کیتا جا سکدا۔ حساب دے ٹیورنگ مشین نمونے وچ
ایس دا مطلب اے کہ کوئی دِتّی ٹیورنگ مشین کسے دِتّے انپٹ
اُتے رُکے گی یا نہیں، ایہ حسابی طور اُتے ناقابلِ فیصلہ اے۔
ایس لئی اگلے قضیے نوں ``حساب رُکّن دے مسئلے دی ناقابلِ
فیصلہ حیثیت'' آکھیا جاندا اے۔ تھلّے اسی دو ثبوت دیاں گے۔
پہلا پچھلی بحث دی کڑی اگّے ودھاندا اے، دوجا ودھ سِدھا اے۔

\begin{thm}
\ollabel{thm:halting-problem}
من لو
\[
h(e, x)  =
\begin{cases}
1 & \text{جے\/ $\fn{Un}(e, x)$ تعریف شدہ ہووے} \\
0 & \text{ہور صورت وچ۔}
\end{cases}
\]
تاں $h$ قابلِ حساب نہیں۔
\end{thm}

\begin{proof}
فرض کرو $h$ قابلِ حساب اے۔ اسی وکھانا اے کہ ایس نال
اک ہمہ گیر قابلِ حساب فنکشن دی تعریف ہو سکدی اے۔ من لو
\[
\fn{Un'}(e,x) =
\begin{cases}
\fn{Un}(e,x) & \text{جے $h(e,x) = 1$} \\
0 & \text{ہور صورت وچ۔}
\end{cases}
\]
ہن $\fn{Un'}(e, x)$ اک کُل فنکشن اے، تے جے~$h$
قابلِ حساب اے تاں ایہ وی قابلِ حساب اے۔ مثال دے طور اُتے،
اسی بنیادی بازگشت نال $g$ دی تعریف کر سکدے آں:
\begin{align*}
g(0, e, x) & \simeq 0 \\
g(y+1, e, x) & \simeq \fn{Un}(e,x);
\end{align*}
فیر
\[
\fn{Un'}(e,x) \simeq g(h(e,x),e,x).
\]
$\fn{Un'}(e,x)$، $\fn{Un}(e,x)$ دے برابر اے
جتھے وی دوجا تعریف شدہ اے؛ ایس لئی $\fn{Un'}$
اوہناں جزوی قابلِ حساب فنکشناں لئی ہمہ گیر اے جیہڑے
اتفاق نال کُل نیں۔ پر ایہ \olref[nou]{thm:no-univ}
نال ٹکراندا اے۔
\end{proof}

\begin{proof}
فرض کرو $h(e,x)$ قابلِ حساب ہووے۔ $g$ دی تعریف ایویں کرو:
\[
g(x) =
\begin{cases}
  0                & \text{جے $h(x,x) = 0$} \\
  \fundefined & \text{ہور صورت وچ۔}
\end{cases}
\]
فنکشن $g$ جزوی قابلِ حساب اے۔ مثال دے طور اُتے، ایس نوں
$\umin{y}{h(x,x) = 0}$ دے طور اُتے تعریف کیتا جا سکدا اے۔
ایس لئی کسے~$e$ لئی $g(x) \simeq \fn{Un}(e,
x)$ ہر~$x$ اُتے۔ کیہ $g$، $e$ اُتے تعریف شدہ اے؟
جے اے، تاں $g$ دی تعریف مطابق $h(e,e) = 0$
(تعریف شدہ ہوئے بغیر~$h$ دی قدر~$0$ نہیں ہو سکدی)۔
$h$ دی تعریف مطابق ایس دا مطلب اے کہ $\fn{Un}(e, e)$
غیر تعریف شدہ اے۔ پر ساڈی فرض کیتی برابری
$g(x) \simeq \fn{Un}(e, x)$ ہر~$x$ لئی اے،
ایس لئی $g(e)$ وی غیر تعریف شدہ اے: تضاد۔
دوجے پاسے جے $g(e)$ غیر تعریف شدہ اے، تاں
$h(e,e) \neq 0$، تے ایس لئی $h(e,e) = 1$۔
ایس توں $\fn{Un}(e, e)$ تعریف شدہ ہوندا اے۔
پر $g(x) \simeq \fn{Un}(e, x)$ ہون کرکے $g(e)$
وی تعریف شدہ ہوندا۔ فیر تضاد۔
\end{proof}

\begin{tagblock}{TMs}
\begin{explain}
ایس دلیل نوں ٹیورنگ مشیناں دے لفظاں وچ وی بیان کر سکدے آں۔
فرض کرو کوئی ٹیورنگ مشین~$H$ ہووے جیہڑی ٹیورنگ مشین~$E$
دا بیان تے اک انپٹ~$x$ لے کے فیصلہ کرے کہ $E$ انپٹ~$x$
اُتے رُکدی اے یا نہیں۔ فیر اسی اک ہور ٹیورنگ مشین~$G$
بنا سکدے آں جیہڑی اکّو انپٹ~$x$ لے، $H$ نوں چلا کے
فیصلہ کرے کہ مشین $M_x$، جیہدا اشاریہ~$x$ اے، انپٹ~$x$
اُتے رُکدی اے یا نہیں، تے اوہدے اُلٹ کرے۔ ہور لفظاں وچ،
جے $H$ دسے کہ $M_x$ انپٹ~$x$ اُتے رُکدی اے، تاں $G$
لامتناہی چکر وچ پے جاندی اے؛ تے جے $H$ دسے کہ $M_x$
انپٹ~$x$ اُتے نہیں رُکدی، تاں $G$ بس رُک جاندی اے۔
کیہ $G$ اپنا اشاریہ انپٹ بنا کے رُکدی اے؟ اُتّے والی
دلیل دسدی اے کہ اوہ اُدوں تے صرف اُدوں رُکدی اے جدوں
اوہ نہیں رُکدی؛ ایہ تضاد اے۔ ایس لئی ساڈا ایہ فرض غلط
ہونا لازمی اے کہ ایداں دی کوئی ٹیورنگ مشین~$H$ موجود اے۔
\end{explain}
\end{tagblock}
\endgroup
% /OLP-0235

% OLP-0236 content/computability/computability-theory/russells-paradox.tex
\begingroup
\def\olfilename{russells-paradox}
\expandafter\def\csname ol@sectioncs\endcsname{section}

\olfileid{cmp}{thy}{rus}
\olsection{رسل دے پیراڈوکس نال تقابل}

ایس حصے دیاں دلیلاں دا رسل دے پیراڈوکس نال موازنہ کرنا،
تے اوہناں دیاں سانجھاں تے فرق ویکھنا، فائدہ مند اے:
\begin{enumerate}
% PNB-CMP-014: اصل پہلی شق وچ X غیر متعین اے؛سیٹ S دی تعریف تے رسل دے ثبوت مطابق S ہی اپنے آپ نال تقابل وچ آؤندا اے۔
\item رسل دا پیراڈوکس: من لو $S = \Setabs{x}{x \notin x}$۔
فیر $S
  \in S$ اُدوں تے صرف اُدوں جدوں $S \notin S$،
  تے ایہ تضاد اے۔

  \emph{نتیجہ:} ایہو جہا کوئی سیٹ~$S$ نہیں۔ ``سارے سیٹاں
  دا سیٹ'' موجود منّنا سیٹ نظریے دے ہور مسلّماں نال
  غیر موافق اے۔

\item رسل دے پیراڈوکس دا اک بدلیا روپ: من لو $F$ اوہ ``فنکشن''
  اے جیہڑا سارے فنکشناں دے سیٹ توں $\{ 0, 1 \}$ تک اے،
  تے ایس دی تعریف ایویں اے:
  \[
  F(f) =
  \begin{cases}
    1 & \text{جے $f$، $f$ دے دائرۂ تعریف وچ ہووے، تے $f(f) = 0$} \\
    0 & \text{ہور صورت وچ}
  \end{cases}
  \]
  ایداں دی دلیل وکھاندی اے کہ $F(F) = 0$ اُدوں تے صرف
  اُدوں جدوں $F(F) = 1$، تے ایہ تضاد اے۔

  \emph{نتیجہ:} $F$ فنکشن نہیں۔ ``سارے فنکشناں دا سیٹ''
  فنکشن دا دائرۂ تعریف بنن لئی بہت وڈا اے۔

\item قطری دلیل: من لو $f_0$، $f_1$، \dots جزوی قابلِ حساب
  فنکشناں دی فہرست اے، تے $G \colon \Nat \to
  \{ 0, 1 \}$ دی تعریف ایویں اے:
  \[
  G(x) =
  \begin{cases}
    1 & \text{جے $f_x(x)\downarrow = 0$} \\
    0 & \text{ہور صورت وچ}
  \end{cases}
  \]
  جے $G$ قابلِ حساب اے، تاں ایہ فنکشن $f_k$ اے،
  کسے $k$ لئی۔ پر فیر $G(k) = 1$ اُدوں تے صرف اُدوں
  جدوں $G(k) = 0$، تے ایہ تضاد اے۔

  \emph{نتیجہ:} $G$ قابلِ حساب نہیں۔ یاد رکھو کہ سیٹ نظریے
  دے مسلّماں مطابق $G$ فیر وی اک فنکشن اے؛ایتھے کوئی
  پیراڈوکس نہیں، صرف فرق دی وضاحت اے۔
\end{enumerate}

جزوی فنکشناں، قابلِ حساب فنکشناں، جزوی قابلِ حساب فنکشناں
وغیرہ دی ایہ گل الجھا سکدی اے۔ $\Nat$ توں $\Nat$ تک
سارے جزوی فنکشناں دا سیٹ چیزاں دا اک وڈا اکٹھ اے۔
اوہناں وچوں کجھ کُل نیں، کجھ قابلِ حساب نیں، کجھ کُل
تے قابلِ حساب دونوں نیں، تے کجھ نہ کُل نیں نہ قابلِ حساب۔
چیتے رکھو کہ جدوں اسی صرف ``فنکشن'' آکھدے آں، تاں
عام طور اُتے اسی کُل فنکشن دی گل کر رہے ہوندے آں۔ ایس طرح
چار قسمّاں نیں:
\begin{enumerate}
\item قابلِ حساب فنکشن
\item جزوی قابلِ حساب فنکشن جیہڑے کُل نہیں
\item فنکشن جیہڑے قابلِ حساب نہیں
\item جزوی فنکشن جیہڑے نہ کُل نیں نہ قابلِ حساب
\end{enumerate}
ایہ فرق صاف کرن لئی $\Nat$ توں $\Nat$ تک سارے جزوی
فنکشناں دا اک وڈا چوکور بنا کے ویکھو، تے اوہدے اندر
دو اک دوجے نال سانجھ رکھدے حصے نشان لاؤ: اک کُل
فنکشناں لئی، دوجا جزوی قابلِ حساب فنکشناں لئی۔
ایس خاکے توں بنن والے ہر حصے وچوں اک اک شے دی مثال
دے سکنا چنگی مشق اے۔
\endgroup
% /OLP-0236

% OLP-0237 content/computability/computability-theory/computable-sets.tex
\begingroup
\def\olfilename{computable-sets}
\expandafter\def\csname ol@sectioncs\endcsname{section}

\olfileid{cmp}{thy}{cps}
\olsection{قابلِ حساب سیٹ}

قابلِ حساب فنکشناں توں حساب پذیری دا خیال قابلِ حساب
سیٹاں تک ودھایا جا سکدا اے:

\begin{defn}
  من لو $S$ قدرتی عدداں دا اک سیٹ اے۔ تاں $S$ اُدوں تے
  صرف اُدوں \emph{قابلِ حساب} اے جدوں اوہدا خصوصیتی
  فنکشن~$\Char{S}$ قابلِ حساب ہووے۔ ہور لفظاں وچ،
  $S$~اُدوں تے صرف اُدوں قابلِ حساب اے جدوں فنکشن
\[
\Char{S}(x) =
\begin{cases}
1 & \text{جے $x \in S$} \\
0 & \text{ہور صورت وچ}
\end{cases}
\]
قابلِ حساب ہووے۔ ایداں ای تعلق $R(x_0, \dots, x_{k-1})$
اُدوں تے صرف اُدوں قابلِ حساب اے جدوں اوہدا خصوصیتی
فنکشن قابلِ حساب ہووے۔

قابلِ حساب سیٹاں تے تعلقاں نوں \emph{فیصلہ کرن جوگ
(decidable)} وی آکھیا جاندا اے۔
\end{defn}

\begin{explain}
دھیان دیو کہ ہن ساڈے کول حساب پذیری دے کئی خیال نیں:
جزوی فنکشناں لئی، فنکشناں لئی، تے سیٹاں لئی۔ ایہناں
نوں آپس وچ نہ رلاؤ!{} \iftag{TMs}{جزوی فنکشن دا
حساب کرن والی ٹیورنگ مشین اوہناں انپٹ قدراں لئی
فنکشن دا آؤٹ پٹ دیندی اے جنہاں اُتے فنکشن تعریف شدہ
اے؛ سیٹ دا حساب کرن والی ٹیورنگ مشین ایہ فیصلہ کرن
مگروں کہ دِتّی انپٹ قدر سیٹ وچ اے یا نہیں، یا $1$
یا~$0$ واپس دیندی اے۔}{}
\end{explain}
\endgroup
% /OLP-0237

% OLP-0238 content/computability/computability-theory/ce-sets.tex
\begingroup
\def\olfilename{ce-sets}
\expandafter\def\csname ol@sectioncs\endcsname{section}

\olfileid{cmp}{thy}{ces}
\olsection{قابلِ حساب گنتی وار فہرست بنن جوگ سیٹ}

\begin{defn}
کوئی سیٹ \emph{قابلِ حساب گنتی وار فہرست بنن جوگ
\pnbLatin{(computably enumerable)}} اُدوں اے جدوں اوہ خالی
ہووے یا کسے قابلِ حساب فنکشن دیاں حاصل قدراں دا
سیٹ ہووے۔
\end{defn}

\begin{history}
قابلِ حساب گنتی وار فہرست بنن جوگ سیٹاں نوں
\emph{بازگشتی طور اُتے گنتی وار فہرست بنن جوگ
\pnbLatin{(recursively enumerable)}} سیٹ وی آکھیا جاندا اے۔
ایہ اصلی اصطلاح اے؛ اج کل دوہاں اصطلاحاں دے نال
نال اوہناں دے مخفف ``c.e.'' تے ``r.e.'' وی عام
ورتے جاندے نیں۔
\end{history}

\begin{explain}
تعریف دا مطلب کی اے تے ایہ اصطلاح کیوں مناسب اے،
ایس اُتے غور کرو۔ خیال ایہ اے کہ جے $S$ قابلِ حساب
فنکشن~$f$ دیاں حاصل قدراں دا سیٹ ہووے، تاں
\[
S = \{ f(0), f(1), f(2), \dots \},
\]
ایس لئی $f$ نوں $S$ دے رکناں دی گنتی وار فہرست
بناؤن والا فنکشن سمجھیا جا سکدا اے۔ دھیان دیو،
تعریف مطابق $f$ دا ودھدا جان والا فنکشن ہونا ضروری
نہیں؛ یعنی فہرست ودھدی ترتیب وچ ہونی لازمی نہیں۔
بلکہ $f$ دا اک بہ اک (injective) ہونا وی ضروری نہیں؛
فہرست $f(0)$، $f(1)$، $f(2)$، \dots{} وچ $S$ دے
رکناں دی دہرائی ہو سکدی اے۔ مثال لئی مستقل فنکشن
$f(x) = 0$ سیٹ $\{ 0
\}$ دی گنتی وار فہرست بناندا اے۔
\end{explain}

ہر قابلِ حساب سیٹ قابلِ حساب گنتی وار فہرست بنن جوگ
ہوندا اے۔ ایہ ویکھن لئی من لو $S$ قابلِ حساب اے۔
جے $S$ خالی اے، تاں تعریف مطابق اوہ قابلِ حساب
گنتی وار فہرست بنن جوگ اے۔ نہیں تاں $a$ نوں $S$
دا کوئی وی رکن چن لو۔ $f$ دی تعریف ایویں کرو:
\[
f(x) =
\begin{cases}
x & \text{جے $\Char{S}(x) = 1$} \\
a & \text{ہور صورت وچ۔}
\end{cases}
\]
تاں $f$ قابلِ حساب فنکشن اے، تے $S$ فنکشن~$f$
دیاں حاصل قدراں دا سیٹ اے۔
\endgroup
% /OLP-0238

% OLP-0239 content/computability/computability-theory/equiv-ce-defs.tex
\begingroup
\def\olfilename{equiv-ce-defs}
\expandafter\def\csname ol@sectioncs\endcsname{section}

\olfileid{cmp}{thy}{eqc}

\olsection[قابلِ حساب فہرست بنن جوگ سیٹاں دیاں تعریفاں]{قابلِ حساب
  گنتی وار فہرست بنن جوگ سیٹاں دیاں ہم ارز تعریفاں}


ہِٹھاں دتے بیان ایہ سمجھاؤن دے کئی اہم تے ہم ارز
طریقے نیں کہ سیٹ دے قابلِ حساب گنتی وار فہرست
بنن جوگ ہون دا مطلب کی اے۔

\begin{thm}
\ollabel{thm:ce-equiv}
من لو $S$ قدرتی عدداں دا سیٹ اے۔ تاں ہِٹھلیاں
شرطاں آپس وچ ہم ارز نیں:
\begin{enumerate}
\item\ollabel{case:ce} $S$ قابلِ حساب گنتی وار فہرست بنن جوگ اے۔
\item\ollabel{case:ran-pc} $S$ کسے \emph{جزوی} قابلِ حساب فنکشن دیاں حاصل قدراں دا سیٹ اے۔
\item\ollabel{case:ran-prim} $S$ خالی اے یا کسے بنیادی بازگشتی فنکشن دیاں حاصل قدراں دا سیٹ اے۔
\item\ollabel{case:ce-domain} $S$ کسے جزوی قابلِ حساب فنکشن دا \emph{دائرۂ تعریف} اے۔
\end{enumerate}
\end{thm}

\begin{explain}
پہلیاں تِن شرطاں دا مطلب اے کہ کسے وی غیر خالی
قابلِ حساب گنتی وار فہرست بنن جوگ سیٹ دی گنتی وار
فہرست قابلِ حساب فنکشن، جزوی قابلِ حساب فنکشن یا
بنیادی بازگشتی فنکشن نال بنائی جا سکدی اے۔ چوتھی
شرط دسدی اے کہ جے $S$ قابلِ حساب گنتی وار فہرست
بنن جوگ اے، تاں کسے اشاریے~$e$ لئی
\[
S = \Setabs{x}{\cfind{e}(x) \fdefined}.
\]
ہور لفظاں وچ، $S$ اوہناں انپٹاں دا سیٹ اے جنہاں
لئی $\cfind{e}$ دی حسابی کارروائی رُک جاندی اے۔
ایس واسطے قابلِ حساب گنتی وار فہرست بنن جوگ سیٹاں
نوں کدی کدی \emph{نیم فیصلہ کرن جوگ
(semi-decidable)} وی آکھیا جاندا اے: جے کوئی عدد
سیٹ وچ ہووے تاں آخرکار ``ہاں'' مل جاندی اے، پر
جے اوہ سیٹ وچ نہ ہووے تاں ``نہیں'' کدی نہیں ملدی!{}
\end{explain}

\begin{proof}
ہر بنیادی بازگشتی فنکشن قابلِ حساب ہوندا اے تے ہر
قابلِ حساب فنکشن جزوی قابلِ حساب وی ہوندا اے، ایس
لئی \olref{case:ran-prim} توں \olref{case:ce}
نکلدا اے، تے \olref{case:ce} توں
\olref{case:ran-pc} نکلدا اے۔ (دھیان دیو: جے $S$
خالی ہووے، تاں $S$ اوس جزوی قابلِ حساب فنکشن
دیاں حاصل قدراں دا سیٹ اے جو کسے وی انپٹ اُتے
تعریف شدہ نہیں۔) جے اسی ایہ وکھا دئیے کہ
\olref{case:ran-pc} توں \olref{case:ran-prim}
نکلدا اے، تاں پہلیاں تِن شرطاں دی ہم ارزیت ثابت
ہو جائے گی۔

ہن من لو $S$ جزوی قابلِ حساب فنکشن
$\cfind{e}$ دیاں حاصل قدراں دا سیٹ اے۔ جے $S$
خالی اے تاں گل ثابت ہو گئی۔ نہیں تاں $a$ نوں
$S$ دا کوئی وی رکن چن لو۔ کلینی دے نارمل روپ
والے قضیے مطابق اسی لکھ سکدے آں:
\[
\cfind{e}(x) = U(\umin{s}{T(e, x, s)}).
\]
خاص طور اُتے، $\cfind{e}(x) \fdefined$ تے $= y$
اُدوں تے صرف اُدوں ہوندا اے جدوں کوئی $s$ ہووے
جس لئی $T(e, x, s)$ درست ہووے تے $U(s) = y$
ہووے۔ $f(z)$ دی تعریف ایویں کرو:
\[
f(z) = \begin{cases}
  U((z)_1) & \text{جے $T(e, (z)_0, (z)_1)$} \\
  a        & \text{ہور صورت وچ۔}
\end{cases}
\]
تاں $f$ بنیادی بازگشتی اے، کیوں جے $T$ تے $U$
بنیادی بازگشتی نیں۔ \iftag{TMs}{ٹیورنگ مشیناں
دے لحاظ نال، جے $z$ جوڑی
$\tuple{(z)_0, (z)_1}$ دا کوڈ ہووے تے $(z)_1$،
مشین~$M_e$ دی انپٹ $(z)_0$ اُتے رُک جان والی
حسابی کارروائی ہووے، تاں $f$ اوس کارروائی دا
آؤٹ پٹ دیندا اے؛ نہیں تاں اوہ~$a$ واپس دیندا اے۔}{}

ساڈے لئی ایہ ثابت کرنا باقی اے کہ $S$ عین~$f$
دیاں حاصل قدراں دا سیٹ اے؛ یعنی ہر قدرتی
عدد~$y$ لئی، $y \in S$ اُدوں تے صرف اُدوں
جدوں اوہ~$f$ دیاں حاصل قدراں وچ ہووے۔ پہلے
رخ وچ من لو $y \in S$۔ تاں $y$،
$\cfind{e}$ دیاں حاصل قدراں وچ اے؛ ایس لئی
کسے $x$ تے~$s$ لئی $T(e,x,s)$ درست اے
تے $U(s) = y$۔ پر فیر $y = f(\tuple{x,s})$۔
اُلٹے رخ وچ من لو $y$، فنکشن~$f$ دیاں
حاصل قدراں وچ اے۔ تاں یا $y = a$ اے،
یا کسے~$z$ لئی $T(e,(z)_0,(z)_1)$ تے
$U((z)_1) = y$۔ % PNB-CMP-015: latter case uses input (z)_0, not unbound x.
دوجی صورت وچ $\cfind{e}((z)_0) \fdefined =
y$ ہون کرکے، دوناں صورتاں وچ $y$،
$S$ وچ اے۔

(علامت $\cfind{e}(x) \fdefined = y$ دا مطلب
اے کہ ``$\cfind{e}(x)$ تعریف شدہ اے تے
$y$ دے برابر اے۔'' اسی $\cfind{e}(x) =
y$ وی لکھ سکدے ساں، پر وادھو تیر کدی کدی
ایہ یاد کران لئی فائدہ مند اے کہ ایتھے فنکشن
جزوی اے۔)

\olref{thm:ce-equiv} دا ثبوت پورا کرن لئی
ایہ وکھانا کافی اے کہ \olref{case:ce} تے~
\olref{case:ce-domain} ہم ارز نیں۔ پہلاں
وکھائیے کہ \olref{case:ce} توں~
\olref{case:ce-domain} نکلدا اے۔ % PNB-CMP-016: source omits empty CE case in this direction.
جے زیرِ غور سیٹ خالی اے تاں اوہ کسے وی انپٹ
اُتے تعریف نہ ہون والے جزوی قابلِ حساب فنکشن
دا دائرۂ تعریف اے۔ ہن من لو $S$
قابلِ حساب فنکشن~$f$ دیاں حاصل قدراں دا سیٹ
اے، یعنی
\[
S = \Setabs{y}{\text{کسے $x$ لئی، } f(x) = y}.
\]
ایہ تعریف کرو:
\[
g(y) = \umin{x}{(f(x) = y)}.
\]
تاں $g$ جزوی قابلِ حساب فنکشن اے، تے $g(y)$
اُدوں تے صرف اُدوں تعریف شدہ اے جدوں
کسے~$x$ لئی $f(x) = y$ ہووے۔ ہور
لفظاں وچ، $g$ دا دائرۂ تعریف~$f$
دیاں حاصل قدراں دا سیٹ اے۔
\iftag{TMs}{ٹیورنگ مشیناں دے لحاظ نال: جے
ٹیورنگ مشین~$F$، $S$ دے رکناں دی
گنتی وار فہرست بناندی اے، تاں مشین $G$
نوں ایداں بناؤ کہ اوہ $S$ بارے نیم فیصلہ کرے۔
اوہ~$F$ دے آؤٹ پٹاں وچ دِتّا رکن لبھدی
رہندی اے؛ رکن لبھ جائے تاں رُک جاندی اے،
نہیں تاں ہمیشہ تلاش جاری رکھدی اے۔}{}

آخر وچ \olref{case:ce-domain} توں~
\olref{case:ce} نکلنا وکھاؤن لئی من لو
$S$ جزوی قابلِ حساب فنکشن~$\cfind{e}$
دا دائرۂ تعریف اے، یعنی
\[
S = \Setabs{x}{\cfind{e}(x) \fdefined}.
\]
جے $S$ خالی اے تاں گل ثابت اے؛ نہیں تاں
$a$ نوں~$S$ دا کوئی وی رکن چن لو۔
$f$ دی تعریف ایویں کرو:
\[
f(z) = \begin{cases}
(z)_0 & \text{جے $T(e,(z)_0,(z)_1)$} \\
a & \text{ہور صورت وچ۔}
\end{cases}
\]
تاں اُتّے وانگ کوئی عدد $x$، فنکشن~$f$
دیاں حاصل قدراں وچ اُدوں تے صرف اُدوں
اے جدوں $\cfind{e}(x) \fdefined$، یعنی
اُدوں تے صرف اُدوں جدوں $x \in S$۔
\iftag{TMs}{ٹیورنگ مشیناں دے لحاظ نال:
جے مشین $M_e$، $S$ بارے نیم فیصلہ
کردی اے، تاں ٹیورنگ مشیناں دیاں ساریاں ممکنہ
حسابی کارروائیاں دی فہرست نوں اک اک کرکے ویکھو تے
اوہناں رُک جان والیاں کارروائیاں دے
متعلقہ انپٹ واپس دیو؛ ایس طرح $S$
دے رکناں دی فہرست بن جاندی اے۔}{}
\end{proof}

\olref{thm:ce-equiv} دی شرط~
\olref{case:ce-domain}، قابلِ حساب گنتی وار
فہرست بنن جوگ سیٹاں دی فہرست بناؤن دا اک
سَوکھا طریقہ دیندی اے: ہر~$e$ لئی $W_e$
نوں $\cfind{e}$ دا دائرۂ تعریف من لو،
یعنی
\[
W_e = \Setabs{x}{\cfind{e}(x) \fdefined}.
\]
تاں جے $A$ کوئی وی قابلِ حساب گنتی وار
فہرست بنن جوگ سیٹ اے، تاں $A = W_e$
کسے~$e$ لئی ہوندا اے۔

اگلا قضیہ قابلِ حساب گنتی وار فہرست بنن
جوگ سیٹاں دی اک ہور تعریف دیندا اے۔

\begin{thm}
\ollabel{thm:exists-char}
سیٹ $S$ اُدوں تے صرف اُدوں قابلِ حساب
گنتی وار فہرست بنن جوگ اے جدوں کوئی
قابلِ حساب تعلق $R(x,y)$ ہووے جس لئی
\[
S = \Setabs{ x }{ \lexists[y][R(x,y)] }.
\]
\end{thm}

\begin{proof}
پہلے رخ وچ من لو $S$ قابلِ حساب گنتی وار
فہرست بنن جوگ اے۔ تاں کسے $e$ لئی،
$S = W_e$۔ ایس~$e$ لئی اسی $S$
نوں ایویں لکھ سکدے آں:
\[
S = \Setabs{ x }{ \lexists[y][T(e, x, y)] }.
\]
اُلٹے رخ وچ من لو $S = \Setabs{ x }{
  \lexists[y][R(x, y)] }$۔ $f$~دی
تعریف ایویں کرو:
\[
f(x) \simeq \umin{y}{R(x, y)}.
\]
تاں $f$ جزوی قابلِ حساب اے، تے $S$
فنکشن~$f$ دا دائرۂ تعریف اے۔
\end{proof}
\endgroup
% /OLP-0239

% OLP-0240 content/computability/computability-theory/non-comp-set.tex
\begingroup
\def\olfilename{non-comp-set}
\expandafter\def\csname ol@sectioncs\endcsname{section}

\olfileid{cmp}{thy}{ncp}
\olsection{ناقابلِ حساب سیٹ موجود نیں}


اُتّے اسی ویکھیا سی کہ ہر قابلِ حساب سیٹ، قابلِ
حساب گنتی وار فہرست بنن جوگ ہوندا اے۔ کی اِس
دا اُلٹ وی درست اے؟ اگلا نتیجہ وکھاندا اے کہ
عام طور اُتے ایہ درست نہیں۔

\begin{thm}\ollabel{thm:K-0}
من لو $K_0$ سیٹ
$\Setabs{\tuple{e, x}}{\cfind{e}(x) \fdefined}$
اے۔ تاں $K_0$ قابلِ حساب گنتی وار فہرست بنن
جوگ اے، پر قابلِ حساب نہیں۔
\end{thm}

\begin{proof}
ایہ ویکھن لئی کہ $K_0$ قابلِ حساب گنتی وار
فہرست بنن جوگ اے، دھیان دیو کہ اوہ ہِٹھاں تعریف کیتے
فنکشن~$f$ دا دائرۂ تعریف اے:
\[
f(z) = \umin{y}{(\len{z} = 2 \land T((z)_0, (z)_1, y))}.
\]
کیوں جے جے $\cfind{e}(x)$ تعریف شدہ اے،
تاں $f(\tuple{e, x})$ رُک جان والی حسابی
کارروائی دا سلسلہ لبھ لیندا اے؛ جے
$\cfind{e}(x)$ تعریف شدہ نہیں، تاں
$f(\tuple{e, x})$ وی تعریف شدہ نہیں؛
تے جے $z$ جوڑی دا کوڈ ای نہیں،
تاں $f(z)$ وی تعریف شدہ نہیں۔

$K_0$ دے قابلِ حساب نہ ہون دی وجہ بس
حساب رُکّن دے مسئلے دا ناقابلِ فیصلہ ہونا
اے، \olref[hlt]{thm:halting-problem}۔
\end{proof}

سیٹ $K_0$ اوہناں جوڑیاں $\tuple{e,x}$
دا سیٹ اے جنہاں لئی $\cfind{e}(x) \fdefined$؛
یعنی $\tuple{e,x} \in K_0$ اُدوں تے
صرف اُدوں جدوں $\cfind{e}$ انپٹ~$x$
اُتے تعریف شدہ ہووے، یعنی رُک جائے۔
ایس لئی ایس نوں ``رُکّن والا سیٹ''
وی آکھیا جاندا اے۔ سیٹ
$K = \Setabs{e}{\cfind{e}(e) \fdefined}$
نوں ``اپنے آپ اُتے رُکّن والا سیٹ''
آکھیا جاندا اے۔ ایہ اکثر اک معیاری
ناقابلِ فیصلہ سیٹ دے طور اُتے ورتیا جاندا اے۔

\begin{thm}\ollabel{thm:K}
اپنے آپ اُتے رُکّن والا سیٹ
$K = \Setabs{e}{\cfind{e}(e) \fdefined}$
، !!{c.e.} اے پر فیصلہ کرن جوگ نہیں۔
\end{thm}

\begin{proof}
من لو $K$ فیصلہ کرن جوگ اے، یعنی اوہدا
خصوصیتی فنکشن $\Char{K}$ قابلِ حساب اے۔
ایہ تعریف کرو:
\[d(e) = \begin{cases}
1 & \text{جے\/ $\Char{K}(e) = 0$}\\
\fundefined & \text{ہور صورت وچ۔}
\end{cases}
\]
من لو $k$، $d$ دا اشاریہ اے، یعنی
$d \simeq \cfind{k}$۔ تاں $d(k)
\simeq \cfind{k}(k)$۔ ایہ اوس حقیقت
دے خلاف اے کہ $d(k) \fdefined$
اُدوں تے صرف اُدوں جدوں
$\cfind{k}(k) \fundefined$، جیہڑی
$d$ دی تعریف توں نکلدی اے۔

$K$، $f(x) = \umin{y}{T(x,x,y)}$ دا
دائرۂ تعریف اے، ایس لئی !!{c.e.} اے۔
\end{proof}
\endgroup
% /OLP-0240

% OLP-0241 content/computability/computability-theory/ce-closed-cup-cap.tex
\begingroup
\def\olfilename{ce-closed-cup-cap}
\expandafter\def\csname ol@sectioncs\endcsname{section}

\olfileid{cmp}{thy}{clo}

\olsection[CE سیٹاں دا اتحاد تے اشتراک]{قابلِ حساب گنتی وار
  فہرست بنن جوگ سیٹ اتحاد تے اشتراک تحت بند نیں}

اگلا قضیہ قابلِ حساب گنتی وار فہرست بنن جوگ
سیٹاں دے مجموعے دیاں کجھ بندش دیاں خاصیتاں
دسدا اے۔

\begin{thm}
من لو $A$ تے $B$ قابلِ حساب گنتی وار
فہرست بنن جوگ نیں۔ تاں $A \cap B$
تے $A \cup B$ وی ایداں دے نیں۔
\end{thm}

\begin{proof}
\olref[eqc]{thm:ce-equiv} سانوں قابلِ حساب گنتی وار
فہرست بنن جوگ سیٹاں دیاں وکھ وکھ تعریفاں
ورتن دی اجازت دیندا اے۔ مثال دے طور اُتے
اسی کئی وکھرے ثبوت دیاں گے۔
% PNB-CMP-017: the first two source constructions assume nonempty input sets.
جے دوہاں وچوں کوئی اک سیٹ خالی ہووے تاں
اوہناں دا اتحاد دوجا سیٹ اے تے اشتراک خالی؛
ایس صورت وچ دعویٰ فوراً نکلدا اے۔ پہلے
دو ثبوتاں دیاں تعمیرات وچ اس لئی غیر خالی
صورت نوں ویکھنا کافی اے۔

پہلے ثبوت لئی من لو $A$ دی گنتی وار فہرست
قابلِ حساب فنکشن~$f$ بناندا اے، تے $B$
دی گنتی وار فہرست قابلِ حساب فنکشن~$g$
بناندا اے۔ ایہ رکھو:
\begin{align*}
h(x) & = \umin{y}{(f(y) = x \lor g(y) = x)} \text{ تے}\\
j(x) & = \umin{y}{(f((y)_0) = x \land g((y)_1) = x)}.
\end{align*}
تاں $A \cup B$، $h$ دا دائرۂ تعریف
اے، تے $A \cap B$، فنکشن~$j$
دا دائرۂ تعریف اے۔

\begin{explain}
حسابی لحاظ نال ایہ ہو رہیا اے: جے ساڈے کول
$A$ تے $B$ دیاں گنتی وار فہرستاں بناؤن والے
طریقے ہون، تاں کسے رکن $x$ دے
$A \cup B$ وچ ہون بارے نیم فیصلہ
کرن لئی دوناں وچوں کسے وی فہرست وچ $x$
لبھ سکدے آں؛ تے $x$ دے $A \cap B$
وچ ہون بارے نیم فیصلہ کرن لئی دوناں
فہرستاں وچ اکّو ویلے $x$ لبھ سکدے آں۔
\end{explain}

دوجے ثبوت لئی فیر من لو $A$ دی
فہرست~$f$ بناندا اے تے $B$ دی
فہرست~$g$۔ ایہ رکھو:
\[
k(x) = \begin{cases}
f(x/2) & \text{جے $x$ جفت ہووے} \\
g((x-1)/2) & \text{جے $x$ طاق ہووے۔}
\end{cases}
\]
تاں $k$، $A \cup B$ دی فہرست
بناندا اے؛ خیال ایہ اے کہ $k$،
$f$ تے~$g$ دیاں دتیاں فہرستاں
وارو واری ورتدا اے۔ $A \cap
B$ دی فہرست بنانا ذرا اوکھا اے۔
جے $A \cap B$ خالی اے، تاں ایہ
تعریف مطابق ای گنتی وار فہرست بنن جوگ
اے۔ نہیں تاں $c$ نوں $A \cap B$
دا کوئی وی رکن چن لو، تے $l$ دی
تعریف ایویں کرو:
\[
l(x) = \begin{cases}
f((x)_0) & \text{جے $f((x)_0) = g((x)_1)$} \\
c & \text{ہور صورت وچ۔}
\end{cases}
\]
حسابی لحاظ نال $l$، $f$ تے $g$
دیاں فہرستاں دے رکناں دیاں جوڑیاں
ویکھدا اے تے جتھے دوویں ملن، اوہ
قدر واپس کردا اے؛ نہیں تاں صرف $c$
واپس کردا اے، یعنی اوہی مقرر قدر دُہراندا اے۔

آخری ثبوت لئی من لو $A$، جزوی
فنکشن $m(x)$ دا \emph{دائرۂ تعریف}
اے تے $B$، جزوی فنکشن $n(x)$
دا دائرۂ تعریف اے۔ تاں $A \cap B$،
جزوی فنکشن $m(x) +
n(x)$ دا دائرۂ تعریف اے۔

\begin{explain}
حسابی لحاظ نال جے $A$ اوہ قدراں
دا سیٹ اے جنہاں اُتے $m$ رُکدا اے،
تے $B$ اوہ قدراں دا سیٹ اے جنہاں
اُتے $n$ رُکدا اے، تاں $A \cap B$
اوہ قدراں دا سیٹ اے جنہاں اُتے
دونوں طریقے رُکدے نیں۔
\end{explain}

$A \cup B$ نوں رُکّن والیاں قدراں
دے سیٹ دے روپ وچ لکھنا ودھ اوکھا اے،
کیوں جے $m$ تے $n$ دے حسابی عملاں دی
متوازی طور اُتے مشینی نقل کرنی پیندی اے۔ من لو
$d$، $m$ دا اشاریہ اے تے $e$،
$n$ دا اشاریہ اے؛ ہور لفظاں وچ،
$m = \cfind{d}$ تے $n = \cfind{e}$۔
تاں $A \cup B$ اِس فنکشن دا دائرۂ
تعریف اے:
\[
p(x) = \umin{y}{(T(d,x,y) \lor T(e,x,y))}.
\]
\begin{explain}
حسابی لحاظ نال، انپٹ $x$ اُتے $p$،
یا $m$ لئی یا $n$ لئی رُک جان والی
حسابی کارروائی لبھدا اے، تے دوہاں
وچوں کوئی وی لبھ جائے تاں رُک جاندا اے۔
\end{explain}
\end{proof}
\endgroup
% /OLP-0241

% OLP-0242 content/computability/computability-theory/complement-ce.tex
\begingroup
\def\olfilename{complement-ce}
\expandafter\def\csname ol@sectioncs\endcsname{section}

\olfileid{cmp}{thy}{cmp}
\olsection[سیٹاں دا متمم]{قابلِ حساب گنتی وار فہرست بنن جوگ
  سیٹ متمم تحت بند نہیں}

من لو $A$ قابلِ حساب گنتی وار فہرست بنن جوگ اے۔
کی~$A$ دا متمم،
$\Complement{A} = \Nat \setminus A$،
ہمیشہ ایداں دا ای گنتی وار فہرست بنن جوگ
ہوندا اے؟ اگلا قضیہ تے نتیجہ دسّن گے
کہ جواب نہیں اے۔

\begin{thm}
\ollabel{thm:ce-comp}
من لو $A$ قدرتی اعداد دا کوئی وی سیٹ اے۔
تاں $A$ قابلِ حساب اے جے تے صرف جے
$A$ تے $\Complement{A}$ دوویں قابلِ حساب
گنتی وار فہرست بنن جوگ ہون۔
\end{thm}

\begin{proof}
اک پاسے دا ثبوت سوکھا اے: جے $A$ قابلِ حساب
اے تاں $\Complement{A}$ وی قابلِ حساب اے
($\Char{A} = 1 \tsub
\Char{\Complement{A}}$)؛ ایس لئی دوویں
قابلِ حساب گنتی وار فہرست بنن جوگ نیں۔

دوجے پاسے، من لو $A$ تے~$\Complement{A}$
دوویں قابلِ حساب گنتی وار فہرست بنن جوگ نیں۔
من لو $A$، $\cfind{d}$ دا دائرۂ تعریف اے،
تے $\Complement{A}$، $\cfind{e}$ دا دائرۂ
تعریف اے۔ $h$ دی تعریف ایویں کرو:
\[
h(x) = \umin{s}{(T(d,x,s) \lor T(e,x,s))}.
\]
ہور لفظاں وچ، انپٹ~$x$ اُتے $h$،
$\cfind{d}$ یا $\cfind{e}$ دی رُک جان والی
حسابی کارروائی لبھدا اے۔ ہن جے $x \in A$
اے تاں پہلی صورت وچ، تے جے $x \in
\Complement{A}$ اے تاں دوجی صورت وچ،
ایہ تلاش ضرور کامیاب ہووے گی۔ ایس لئی
$h$ اک کُل قابلِ حساب فنکشن اے۔
% PNB-CMP-018: A is the d-domain, not the e-domain in the frozen source.
ہن ہر~$x$ لئی $x \in A$ جے تے صرف
جے $T(d, x, h(x))$، یعنی جے
$\cfind{d}$ ایتھے تعریف شدہ اے۔
چونکہ $T(d, x, h(x))$ اک قابلِ حساب
تعلق اے، ایس لئی $A$ قابلِ حساب اے۔
\end{proof}

\begin{explain}
ایس گل نوں غیر رسمی حسابی لفظاں وچ سمجھنا
سوکھا اے: $A$ بارے فیصلہ کرن لئی،
انپٹ $x$ اُتے $\cfind{d}$ تے $\cfind{e}$
دیاں رُک جان والیاں حسابی کارروائیاں
لبھو۔ دوہاں وچوں اک ضرور رُکے گا؛
% PNB-CMP-019: search the d/e domains; d, not e, witnesses membership in A.
جے $\cfind{d}$ رُکے تاں $x$،~$A$
وچ اے، تے نہیں تاں $x$،~$\Complement{A}$
وچ اے۔
\end{explain}

\begin{cor}
\ollabel{cor:comp-k}
$\Complement{K_0}$ قابلِ حساب گنتی وار
فہرست بنن جوگ نہیں اے۔
\end{cor}

\begin{proof}
سانوں پتا اے کہ $K_0$ قابلِ حساب گنتی وار
فہرست بنن جوگ اے، پر قابلِ حساب نہیں۔
جے $\Complement{K_0}$ وی قابلِ حساب
گنتی وار فہرست بنن جوگ ہوندا تاں
\olref{thm:ce-comp} مطابق $K_0$ قابلِ
حساب ہوندا۔ ایہ \olref[ncp]{thm:K-0}
نال متصادم اے۔
\end{proof}
\endgroup
% /OLP-0242

% OLP-0243 content/computability/computability-theory/reducibility.tex
\begingroup
\def\olfilename{reducibility}
\expandafter\def\csname ol@sectioncs\endcsname{section}

\olfileid{cmp}{thy}{red}
\olsection{تقلیل}

\begin{explain}
ہن سانوں پتا اے کہ گھٹّو گھٹ اک سیٹ،
$K_0$، قابلِ حساب گنتی وار فہرست بنن جوگ
اے پر قابلِ حساب نہیں۔ ایہ گل صاف ہونی
چاہیدی اے کہ ایہو جیہے ہور سیٹ وی نیں۔ تقلیل دا طریقہ
سانوں ہر وار بنیادی اصولاں ول واپس گئے
بغیر دوجے سیٹاں دیاں ایہ خاصیتاں ثابت
کرن دا اک طاقتور طریقہ دیندا اے۔

عام طور اُتے سیٹ $A$ دی سیٹ~$B$ ول
تقلیل اک ایداں دا طریقہ اے جیڑا
!!{element}s دے~$B$ وچ ہون یا نہ ہون
دے جواباں نوں !!{element}s دے~$A$ وچ ہون
یا نہ ہون دے جواباں وچ بدل دیندا اے۔
اسی کئی توں اک تقلیل اُتے توجہ
دیاں گے، پر تقلیل دیاں ہور وی کئی قسماں
نیں، جنہاں دیاں خاصیتاں وکھ وکھ نیں۔
تقلیل دے تصور حسابی پیچیدگی دے مطالعے
وچ وی مرکزی نیں، جتھے کارکردگی دا سوال
وی ویکھنا پیندا اے۔ مثال دے طور اُتے
کسے سیٹ نوں NP-مکمل اوہدوں آکھیا
جاندا اے جے اوہ NP وچ ہووے تے ہر NP
مسئلے نوں اوہدے ول تقلیل دتی جا سکدی
ہووے؛ اِس لئی تھلّے دسی تقلیل ورگا
تصور ورتیا جاندا اے، پر اضافی شرط ایہ
اے کہ تقلیل دا حساب کثیر حدّی وقت
\pnbLatin{(polynomial time)} وچ کیتا جا سکّے۔

اسی تقلیل دا تصور پہلے وی بالواسطہ
ورت چکے آں۔ سیٹ~$K$ دی تعریف کرو:
\[
K = \Setabs{x}{\cfind{x}(x) \fdefined},
\]
یعنی $K = \Setabs{x}{x \in W_x}$۔
حساب رُکّن دا مسئلہ حل نہ ہو سکّن
والا ساڈا ثبوت، \olref[hlt]{thm:halting-problem}،
سبھ توں سدھا ایہ دسدا اے کہ $K$ قابلِ
حساب نہیں۔ یاد کرو کہ $K_0$ ایہ سیٹ اے:
\[
K_0 = \Setabs{\tuple{e, x}}{\cfind{e}(x) \fdefined },
\]
% PNB-CMP-020: W_e membership requires the same (e,x) pair order as K_0 above.
یعنی $K_0 = \Setabs{\tuple{e,x}}{x \in W_e}$۔
$K$ دے ناقابلِ حساب ہون دا کوئی وی
ثبوت $K_0$ دے ناقابلِ حساب ہون دے
ثبوت تک سوکھے نال ودھایا جا سکدا اے:
جے $K_0$ قابلِ حساب ہوندا تاں
!!a{element}~$x$ دے $K$ وچ ہون
یا نہ ہون دا فیصلہ صرف ایہ پچھ کے
ہو جاندا کہ جوڑا $\tuple{x, x}$،
$K_0$ وچ اے یا نہیں۔ فنکشن~$f$
جیڑا $x$ نوں $\tuple{x, x}$ ول
لے جاندا اے، $K$ دی $K_0$ ول
\emph{تقلیل} دی اک مثال اے۔
\end{explain}

\begin{defn}
من لو $A$ تے $B$ قدرتی اعداد دے
سیٹ نیں۔ قابلِ حساب فنکشن~$f\colon
\Nat \to \Nat$، $A$ دی~$B$ ول
\emph{کئی توں اک تقلیل} اے جے تے
صرف جے ہر قدرتی عدد~$x$ لئی:
\[
x \in A \quad \text{جے تے صرف جے} \quad f(x) \in B.
\]
جے ایہو جیہا $f$ موجود ہووے تاں
اسی آکھدے آں کہ $A$ دی~$B$ ول
\emph{کئی توں اک تقلیل ہو سکدی اے}؛
ایہنوں $A \leq_m B$ لکھدے آں۔
جے $A$ دی $B$ ول کئی توں اک تقلیل
ہو سکدی اے تے الٹی سمت
وی ایداں ای ہووے، تاں $A$ تے $B$ نوں
\emph{کئی توں اک ہم ارز} آکھدے آں؛
ایہنوں $A \equiv_m B$ لکھدے آں۔
\end{defn}

جے اُتّے دتی تعریف والا فنکشن $f$
اک بہ اک ہووے، تاں اسی آکھدے آں کہ $A$ دی~$B$ ول
\emph{اک توں اک تقلیل ہو سکدی اے}۔ ہیٹھاں دتیاں زیادہ تر
تقلیلاں ایہ ودھ سخت شرط پوری کردیاں
نیں، پر اسی ایس حقیقت نوں نہیں ورتاں گے۔

\begin{digress}
ایہ گل سچ اے، پر بالکل ظاہر نہیں،
کہ اک توں اک تقلیل واقعی کئی توں اک
تقلیل نالوں ودھ سخت شرط اے۔ ہور
لفظاں وچ، ایداں دے لامتناہی سیٹ
$A$ تے~$B$ موجود نیں کہ $A$ دی
~$B$ ول کئی توں اک تقلیل ہو
سکدی اے پر~$B$ ول اک توں اک
تقلیل نہیں ہو سکدی۔
\end{digress}
\endgroup
% /OLP-0243

% OLP-0244 content/computability/computability-theory/prop-reduce.tex
\begingroup
\def\olfilename{prop-reduce}
\expandafter\def\csname ol@sectioncs\endcsname{section}

\olfileid{cmp}{thy}{ppr}
\olsection{تقلیل دیاں خاصیتاں}

اسی $A \leq_m B$ لکھدے آں جے $A$
دی~$B$ ول تقلیل ہووے۔ ایہ علامت
ایہ خیال دیندی اے کہ جے $A \leq_m B$
تاں $A$،~$B$ نالوں ودھ اوکھا نہیں،
تے $B$،~$A$ جِنا یا اوہدے نالوں
ودھ اوکھا اے؛ گویا $\le_m$،
اعداد اُتے عام $\le$ وانگوں،
سیٹاں یا فیصلے دے مسئلیاں نوں اوہناں
دی پیچیدگی دے حساب نال ترتیب دیندا اے۔
اگلے دو قضیے ایس خیال دی تائید کردے
نیں۔ پہلا دسدا اے کہ $\le_m$ متعدی اے۔

\begin{prop}
\ollabel{prop:trans-red}
جے $A \leq_m B$ تے $B \leq_m C$
تاں $A \leq_m C$۔
\end{prop}

\begin{proof}
$A$ دی~$B$ ول تقلیل تے $B$ دی
$C$ ول تقلیل دی ترکیب نال
$A$ دی~$C$ ول تقلیل ملدی اے۔
\end{proof}

\begin{prob}
\olref{prop:trans-red} ایہ وکھا کے ثابت کرو:
جے $f$ تے~$g$ بالترتیب $A$ دی~$B$
ول تے $B$ دی~$C$ ول کئی توں اک
تقلیلاں ہون، تاں $\comp{f}{g}$،
$A$ دی~$C$ ول کئی توں اک
تقلیل اے۔
\end{prob}

\begin{prop}
\ollabel{prop:reduce}
من لو $A$ تے $B$ کوئی وی سیٹ نیں
تے $A \leq_m B$۔
\begin{enumerate}
\item جے $B$ قابلِ حساب گنتی وار فہرست
بنن جوگ اے تاں $A$ وی ایداں دا اے۔
\item جے $B$ قابلِ حساب اے تاں
$A$ وی قابلِ حساب اے۔
\end{enumerate}
\end{prop}

\begin{proof}
من لو $f$، $A$ دی~$B$ ول کئی
توں اک تقلیل اے۔ پہلے دعوے لئی
ایہ ویکھنا کافی اے: جے $B$ کسے
جزوی قابلِ حساب فنکشن~$g$ دا
% PNB-CMP-022: CE-domain characterization requires g partial computable.
دائرۂ تعریف اے تاں $A$،
$\comp{f}{g}$ دا دائرۂ تعریف اے:
\begin{align*}
x \in A & \text{ جے تے صرف جے } f(x) \in B \\
& \text{ جے تے صرف جے }  g(f(x)) \fdefined.
\end{align*}

دوجے دعوے لئی یاد کرو کہ جے $B$
قابلِ حساب اے تاں $B$
تے~$\Complement{B}$ دوویں قابلِ حساب
گنتی وار فہرست بنن جوگ نیں
(\olref[cmp]{thm:ce-comp})۔ ایہ وی
سوکھے نال ویکھیا جا سکدا اے کہ $f$،
$\Complement{A}$ دی
$\Complement{B}$ ول کئی توں اک
تقلیل وی اے؛ ایس لئی ثبوت دے پہلے
حصے مطابق $A$ تے~$\Complement{A}$
دوویں !!{computably enumerable} نیں۔ ہن \olref[cmp]{thm:ce-comp}
مطابق $A$ وی قابلِ حساب اے۔
(دوجا رستہ ایہ ویکھنا اے کہ
$\Char{A} = \comp{f}{\Char{B}}$؛
ایس لئی جے $\Char{B}$ قابلِ حساب
اے تاں $\Char{A}$ وی قابلِ حساب اے۔)
\end{proof}

\begin{prob}
من لو $f$، $A$ دی~$B$ ول کئی
توں اک تقلیل اے۔ وکھاؤ کہ $f$،
$\Complement{A}$ دی
$\Complement{B}$ ول وی کئی توں
اک تقلیل اے۔
\end{prob}

\begin{prob}
% PNB-CMP-021: a many-one reduction is total on Nat, not only defined on A.
وکھاؤ کہ جے $f\colon \Nat \to \Nat$
پہلے زیرِ غور سیٹ دی دوجے ول
کئی توں اک تقلیل اے، تاں $\Char{A}
= \comp{f}{\Char{B}}$۔
\end{prob}

\begin{digress}
تقلیل دا اک ہور عمومی تصور،
\emph{ٹیورنگ تقلیل}، ہور حالتاں وچ
خاص کر ناقابلِ فیصلہ ہون دے نتیجے
ثابت کرن لئی مفید اے۔ دھیان دیو
کہ \olref[cmp]{cor:comp-k} مطابق
$K_0$ دے متمم دی~$K_0$ ول کئی توں اک
تقلیل نہیں ہو سکدی، کیوں جے
اوہ قابلِ حساب گنتی وار فہرست
بنن جوگ نہیں۔ پر وجدانی طور اُتے
جے تہانوں $K_0$ بارے سوالاں دے
جواب پتا ہون تاں اوہدے متمم بارے
سوالاں دے جواب وی پتا لگ جان گے۔
اسی آکھدے آں کہ سیٹ $A$ دی~$B$ ول
ٹیورنگ تقلیل ہو سکدی اے جدوں اک قابلِ حساب
طریقہ $A$ بارے سوالاں دے جواب
طے کرن لئی~$B$ بارے سوال
پچھ سکّے۔ ایہ کئی توں
اک تقلیل نالوں کھلّا تصور اے:
اوتھے (1)~$B$ بارے صرف اک سوال
پچھیا جا سکدا اے؛ تے (2)~اوہدے
ہاں دا جواب $A$ بارے ہاں دے
جواب وچ، تے ناں دا جواب ناں دے
جواب وچ بدلنا لازمی اے۔ فیر وی
جے $A$ دی~$B$ ول ٹیورنگ تقلیل
ہو سکدی اے تے $B$ قابلِ
حساب اے تاں $A$ وی قابلِ حساب
اے۔ پر، جداں اسی ویکھ چکے آں،
قابلِ حساب گنتی وار فہرست بنن
جوگ ہون لئی ایداں دا مشابہ دعویٰ
درست نہیں۔

تقلیل دیاں ایس باب وچ وچاریاں
وکھ وکھ قسماں بارے سوچو تے
اوہناں دا فرق سمجھو۔ البتہ اسی
ایس باب وچ صرف کئی توں اک
تقلیل نال ای کم کراں گے۔ ضمنی
طور اُتے، پچھلے پیرا وچ دتیاں
دونوں قسمیں حسابی پیچیدگی وچ
وی ملدیاں نیں، اِس اضافی شرط
نال کہ ٹیورنگ مشیناں کثیر حدّی
وقت \pnbLatin{(polynomial time)} وچ چلّن:
کئی توں اک تقلیل دا پیچیدگی
والا روپ \emph{کارپ تقلیل}
تے ٹیورنگ تقلیل دا پیچیدگی
والا روپ \emph{کُک تقلیل}
کہلاندا اے۔
\end{digress}
\endgroup
% /OLP-0244

% OLP-0245 content/computability/computability-theory/complete-ce-sets.tex
\begingroup
\def\olfilename{complete-ce-sets}
\expandafter\def\csname ol@sectioncs\endcsname{section}

\olfileid{cmp}{thy}{cce}
\olsection{مکمل قابلِ حساب گنتی وار فہرست بنن جوگ سیٹ}

\begin{defn}
سیٹ $A$، کئی توں اک تقلیل دے لحاظ
نال \emph{مکمل !!{computably enumerable} سیٹ}
اے جے:
\begin{enumerate}
\item $A$ قابلِ حساب گنتی وار فہرست
بنن جوگ ہووے، تے
\item ہر ہور قابلِ حساب گنتی وار
فہرست بنن جوگ سیٹ $B$ لئی
$B \leq_m A$ ہووے۔
\end{enumerate}
\end{defn}

ہور لفظاں وچ، مکمل قابلِ حساب گنتی وار
فہرست بنن جوگ سیٹ ایس جماعت دے
سبھ توں اوکھے سیٹ نیں۔ انہاں دی مدد نال
کسے \emph{وی} قابلِ حساب گنتی وار فہرست
بنن جوگ سیٹ دی رکنیت بارے سوالاں دے
جواب دِتے جا سکدے نیں۔

\begin{thm}
$K$، $K_0$ تے $K_1$ تینوں مکمل
قابلِ حساب گنتی وار فہرست بنن جوگ
سیٹ نیں۔
\end{thm}

\begin{proof}
ایہ ویکھن لئی کہ $K_0$ مکمل اے، من لو $B$
کوئی وی قابلِ حساب گنتی وار فہرست
بنن جوگ سیٹ اے۔ تاں کسے اشاریے
$e$ لئی:
\[
B = W_e = \Setabs{x}{\cfind{e}(x) \fdefined}.
\]
من لو $f$ اوہ فنکشن اے جیڑا
$f(x) = \tuple{e, x}$ نال دتا اے۔
تاں ہر قدرتی عدد $x$ لئی
$x \in B$ جے تے صرف جے
$f(x) \in K_0$۔ یعنی $f$،
$B$ دی~$K_0$ ول تقلیل اے۔

ایہ ویکھن لئی کہ $K_1$ مکمل اے، دھیان دیو
کہ \olref[k1]{prop:k1} دے ثبوت وچ
اسی $K_0$ دی اوس ول تقلیل بنائی
سی۔ ایس لئی \olref[ppr]{prop:trans-red}
مطابق ہر قابلِ حساب گنتی وار
فہرست بنن جوگ سیٹ دی~$K_1$
ول وی تقلیل ہو سکدی اے۔

% PNB-CMP-023: completeness of K requires K_0 to reduce to K, not conversely.
$K_0$ دی $K$ ول تقلیل وی ایس
طرح بن سکدی اے۔ جوڑے دے درست
کوڈ توں اک پروگرام دا اشاریہ
قابلِ حساب طریقے نال بناؤ جیڑا
اپنا انپٹ نظر انداز کر کے اوس
جوڑے وچ دتے پروگرام نوں اوہدے
دتے انپٹ اُتے چلاوے؛ s-m-n
قضیے نال ایہ اشاریہ ہر جوڑے توں
اکّو حسابی طریقے نال بن سکدا اے۔
جے کوڈ جوڑے دا نہ ہووے تاں
کدی نہ رُکّن والے پروگرام دا
مقرر اشاریہ لو۔ ایس طرح بنیا
پروگرام اپنے ہی اشاریے اُتے
اُدوں تے صرف اُدوں رُکے گا
جدوں اصل کوڈ رُکّن والیاں جوڑیاں
دے سیٹ وچ ہووے۔
\end{proof}

\begin{prob}
$K$ دی $K_0$ ول تقلیل دیو۔
\end{prob}

\begin{digress}
ایویں پتا لگدا اے کہ ہن تک
ویکھے گئے قابلِ حساب گنتی وار
فہرست بنن جوگ سیٹاں دیاں
ساریاں مثالاں یا تاں قابلِ
حساب نیں یا مکمل۔ ایہ حیرانی
والی گل لگنی چاہیدی اے! کی
ایدان دے قابلِ حساب گنتی وار
فہرست بنن جوگ سیٹ وی نیں
جیڑے نہ قابلِ حساب نیں نہ
مکمل؟ جواب ہاں اے؛ پر ایہ
گل 1950ء دے دہاکے دے
وچکار جا کے فریڈبرگ تے
مُچنک نے وکھ وکھ ثابت کیتی۔
\end{digress}
\endgroup
% /OLP-0245

% OLP-0246 content/computability/computability-theory/k-1.tex
\begingroup
\def\olfilename{k-1}
\expandafter\def\csname ol@sectioncs\endcsname{section}

\olfileid{cmp}{thy}{k1}
\olsection{تقلیل دی اک مثال}

آؤ \olref[ppr]{prop:reduce} دی اک
ورتوں ویکھئیے۔

\begin{prop}
\ollabel{prop:k1}
ایہ سیٹ من لو:
\[
K_1 = \Setabs{e}{\cfind{e}(0) \fdefined}.
\]
تاں $K_1$ قابلِ حساب گنتی وار
فہرست بنن جوگ اے پر قابلِ حساب
نہیں۔
\end{prop}

\begin{proof}
چونکہ $K_1 = \Setabs{e}{\lexists[s][T(e,0,s)]}$،
ایس لئی \olref[eqc]{thm:exists-char}
مطابق $K_1$ قابلِ حساب گنتی وار
فہرست بنن جوگ اے۔

ایہ ثابت کرن لئی کہ $K_1$ قابلِ حساب نہیں،
اسی وکھاواں گے کہ $K_0$ دی اوس ول تقلیل ہو سکدی اے۔

\begin{explain}
ایہ تھوڑا اوکھا اے، کیوں جے $K_1$
ورت کے اسی صرف ایداں دیاں
حسابی کارروائیاں بارے پچھ سکدے
آں جو اک مقرر انپٹ، $0$،
نال شروع ہوندیاں نیں۔ من لو
تہاڈا اک ذہین دوست ایہناں سوالاں
دے جواب دے سکدا اے؛ ایداں
دے دوستاں نوں اوریکل
آکھدے نیں۔ ہن کوئی تہاڈے کول
آ کے پچھدا اے کہ $\tuple{e,
  x}$، $K_0$ وچ اے یا نہیں،
یعنی مشین $e$ انپٹ $x$ اُتے
رُکدی اے یا نہیں۔ تسی اک ہور
مشین $e_x$ بنا سکدے او جیڑی
\emph{ہر} انپٹ نوں نظر انداز
کر کے مشین~$e$ نوں انپٹ
$x$ اُتے چلاوے۔ ہن مشین $e$
دے انپٹ $x$ اُتے رُکّن دا
سوال عین مشین $e_x$ دے
انپٹ $0$ یا کسے وی ہور انپٹ
اُتے رُکّن دے سوال دے برابر
اے۔ تاں تسی دوست نوں پچھو
کہ نویں مشین $e_x$ انپٹ
$0$ اُتے رُکدی اے یا نہیں؛
بدلے ہوئے سوال دا جواب اصل
سوال دا جواب دے دیندا اے۔
ایہ $K_0$ دی~$K_1$ ول
مطلوبہ تقلیل اے۔
\end{explain}

ہمہ گیر جزوی قابلِ حساب فنکشن
ورت کے $f$ نوں ہِٹھاں دتا
تِن آرگومنٹاں والا فنکشن
من لو:
\[
f(x,y,z) \simeq \cfind{x}(y).
\]
دھیان دیو کہ $f$ اپنا تِیجا
انپٹ بالکل نظر انداز کردا اے۔
ایہو جیہا اشاریہ $e$ چنو کہ
$f = \cfind{e}[3]$؛ تاں:
\[
\cfind{e}[3](x,y,z) \simeq \cfind{x}(y).
\]
$s$-$m$-$n$ قضیے مطابق اک
فنکشن $s(e,x,y)$ موجود اے
ایداں کہ ہر $z$ لئی:
\begin{align*}
\cfind{s(e,x,y)}(z) & \simeq \cfind{e}[3](x,y,z) \\
& \simeq \cfind{x}(y).
\end{align*}

\begin{explain}
اُتّے دتی غیر رسمی دلیل
دے لفظاں وچ، $s(e,x,y)$
اوس مشین دا اشاریہ اے جیڑی
ہر انپٹ $z$ نوں نظر انداز
کر کے $\cfind{x}(y)$
دی حسابی کارروائی کردی اے۔
\end{explain}

خاص طور اُتے:
\[
\cfind{s(e,x,y)}(0) \fdefined \quad \text{جے تے صرف جے} \quad
\cfind{x}(y) \fdefined.
\]
ہور لفظاں وچ، $\tuple{x, y} \in K_0$
جے تے صرف جے $s(e,x,y) \in
K_1$۔ ہن $d$ نوں ایداں دے
پروگرام دا مقرر اشاریہ چن لو
جیڑا کدی نہیں رُکدا۔ تاں
ہِٹھاں تعریف کیتا فنکشن
$g$:
% PNB-CMP-024: total reduction must send non-pair codes outside K_1.
\[
g(w) = \begin{cases}
s(e,(w)_0,(w)_1) & \text{جے } \len{w}=2,\\
d & \text{ہور صورت وچ۔}
\end{cases}
\]
$K_0$ دی~$K_1$ ول
تقلیل اے۔ غیر جوڑا کوڈ
پہلے سیٹ وچ نہیں، تے اوہدے
لئی چنیا مستقل اشاریہ وی
دوجے سیٹ وچ نہیں۔
\end{proof}
\endgroup
% /OLP-0246

% OLP-0247 content/computability/computability-theory/total.tex
\begingroup
\def\olfilename{total}
\expandafter\def\csname ol@sectioncs\endcsname{section}

\olfileid{cmp}{thy}{tot}
\olsection{کُل ہون دا فیصلہ ممکن نہیں}

آؤ $s$-$m$-$n$ قضیے دی اک ہور
ورتوں ویکھئیے، جیہدے نال اسی
وکھاواں گے کہ اک سیٹ قابلِ
حساب نہیں۔ $\fn{Tot}$ کُل
قابلِ حساب فنکشناں دے
اشارییاں دا سیٹ اے، یعنی:
\[
\fn{Tot} = \Setabs{x}{\text{ہر $y$ لئی، $\cfind{x}(y)\fdefined$}}.
\]

\begin{prop}
\ollabel{prop:total}
$\fn{Tot}$ قابلِ حساب نہیں۔
\end{prop}

\begin{proof}
ایہ ویکھن لئی کہ $\fn{Tot}$ قابلِ حساب نہیں،
ایہ وکھانا کافی اے کہ $K$ دی ایس
ول تقلیل ہو سکدی اے۔ $h(x,y)$
نوں ایداں تعریف کرو:
\[
h(x,y) \simeq
\begin{cases}
0 & \text{جے $x \in K$} \\
\fundefined & \text{ہور صورت وچ}
\end{cases}
\]
دھیان دیو کہ $h(x,y)$
دا $y$ نال بالکل کوئی
واسطہ نہیں۔ ایہ ویکھنا اوکھا نہیں کہ $h$ جزوی
قابلِ حساب اے: انپٹ
$x, y$ اُتے $h$ دی قدر
کڈھن لئی پہلے $\cfind{x}$ دی انپٹ~$x$ اُتے
حسابی کارروائی دی مشینی نقل کرو؛ جے
ایہ حسابی کارروائی
رُک جائے، تاں $h(x,y)$
دا آؤٹ پٹ $0$ آؤندا اے
تے اوہ رُک جاندا اے۔
سو $h(x,y)$ عین
$\Zero(\umin{s}{T(x,x,s)})$
اے، جتھے $\Zero$
مستقل صفر فنکشن اے۔

$s$-$m$-$n$ قضیے
مطابق اک بنیادی بازگشتی
فنکشن~$k(x)$ موجود اے
ایداں کہ ہر $x$ تے~$y$
لئی:
\[
\cfind{k(x)}(y) =
\begin{cases}
0 & \text{جے $x \in K$} \\
\fundefined & \text{ہور صورت وچ}
\end{cases}
\]
ایس لئی $\cfind{k(x)}$
کُل اے جے $x \in K$
ہووے، تے نہیں تاں
ہر انپٹ اُتے غیر تعریف
شدہ اے۔ لہٰذا $k$،
$K$ دی~$\fn{Tot}$ ول
تقلیل اے۔
\end{proof}

\begin{digress}
اصل وچ $\fn{Tot}$
قابلِ حساب گنتی وار
فہرست بنن جوگ وی
نہیں؛ ایہدی پیچیدگی
«حسابی درجہ بندی»
وچ ہور اُتّے اے۔
پر ایتھے اسی ایس
مزید مضبوط نتیجے
دی بحث نہیں کراں گے۔
\end{digress}
\endgroup
% /OLP-0247

% OLP-0248 content/computability/computability-theory/rice-theorem.tex
\begingroup
\def\olfilename{rice-theorem}
\expandafter\def\csname ol@sectioncs\endcsname{section}

\olfileid{cmp}{thy}{rce}
\olsection{رائس دا قضیہ}

سوچو تاں پتا لگدا اے کہ
\olref[tot]{prop:total} دے
ثبوت وچ $\fn{Tot}$ دیاں
خاص تفصیلاں دا کوئی کردار
نہیں۔ اسی $h(x,y)$ نوں
ایداں بنایا سی کہ ایہ
مستقل فنکشن
$j(y) = 0$ وانگ کم کرے
عین اُدوں جدوں
$x$، $K$ وچ ہووے؛
پر اوس حال وچ اسی ایہنوں
کسے وی ہور جزوی قابلِ حساب
فنکشن وانگ کم کروا سکدے
ساں۔ ایس گل توں اک ہور
عام قضیہ ملدا اے۔ موٹے طور اُتے
ایہ آکھدا اے کہ قابلِ حساب فنکشناں
دی جیہڑی خاصیت کجھ فنکشناں وچ ہووے
تے کجھاں وچ نہ ہووے، اوہدا فیصلہ
قابلِ حساب طریقے نال نہیں ہو سکدا۔

یاد رکھو کہ $\cfind{0}$،
$\cfind{1}$، $\cfind{2}$،
~\dots جزوی قابلِ حساب
فنکشناں دی ساڈی معیاری
گنتی وار فہرست اے۔

\begin{thm}[رائس دا قضیہ]
  $C$ جزوی قابلِ حساب فنکشناں
  دا کوئی وی سیٹ ہووے،
  تے $A =
  \Setabs{n}{\cfind{n} \in C}$
  من لو۔ جے $A$ قابلِ حساب
  اے، تاں یا $C$ خالی اے
  یا $C$ سارے جزوی قابلِ حساب
  فنکشناں دا سیٹ اے۔
\end{thm}

{\em اشارییاں دا سیٹ} اوہ
سیٹ $A$ اے جیہدی خاصیت
ایہ اے: جے $n$ تے $m$
ایکو فنکشن «حساب» کرن
والے اشاریے ہون، تاں
یا دونوں $n$ تے $m$،
$A$ وچ ہوندے نیں، یا
کوئی وی نہیں ہوندا۔
ایہ ویکھنا اوکھا نہیں کہ قضیے وچ
دتے سیٹ~$A$ وچ ایہ خاصیت اے۔
الٹ پاسے، جے $A$~
اشارییاں دا سیٹ اے تے
$C$~اوہناں اشارییاں
نال حساب کیتے فنکشناں
دا سیٹ اے، تاں
$A = \Setabs{n}{\cfind{n} \in C}$۔

\begin{explain}
ایس اصطلاح وچ رائس دا
قضیہ ایہو آکھدا اے کہ
اشارییاں دے ایداں دے
سیٹ دا فیصلہ ممکن نہیں
جیڑا نہ خالی ہووے، نہ
سارے اشارییاں دا سیٹ۔
قضیہ سمجھݨ لئی
\emph{پروگراماں}، مثلاً
تہاڈی پسند دی پروگرامنگ
زبان وچ لکھے پروگراماں،
تے اوہناں ولوں حساب کیتے
فنکشناں وچ فرق یاد رکھنا
مدد کردا اے۔ پروگرام
(اشاریے) نحوی چیزاں نیں؛
اوہناں بارے کئی سوالاں
دا جواب قابلِ حساب اے:
کی ایس پروگرام وچ 150
توں ودھ نشان نیں؟ کی
22 توں ودھ سطراں نیں؟
کی ایس وچ «while»
بیان اے؟ کی «hello world»
لڑی کدے «print» بیان
دا آرگومنٹ بن کے آؤندی اے؟
رائس دا قضیہ آکھدا اے
کہ پروگرام دے
\emph{روَیّے} بارے
ایہو جیہے کسے وی سوال دا جواب
قابلِ حساب نہیں جیہدا جواب کجھ
پروگراماں لئی ہاں تے کجھ لئی ناں ہووے۔
مثلاً: کی پروگرام انپٹ
$0$ اُتے رُکدا اے؟
کی ایہ کدے رُکدا اے؟
کی ایہ کدے کوئی جفت
عدد آؤٹ پٹ کردا اے؟
\end{explain}

\begin{proof}[رائس دے قضیے دا ثبوت]
من لو $C$ نہ خالی اے
تے نہ سارے جزوی قابلِ
حساب فنکشناں دا سیٹ؛
$A$ اوہناں فنکشناں دے
اشارییاں دا سیٹ اے
جو~$C$ وچ نیں۔ اسی
وکھاواں گے کہ جے $A$
قابلِ حساب ہووے تاں
رُکّن دا مسئلہ حل ہو
جاوے؛ ایس لئی $A$~
قابلِ حساب نہیں۔

عامیت برقرار رکھدیاں اسی
من سکدے آں کہ اوہ فنکشن
$f$، جیڑا کسے وی انپٹ اُتے
تعریف شدہ نہیں،
$C$ وچ نہیں۔ جے
ایہ اوس سیٹ وچ ہووے تاں
ہِٹھلی دلیل وچ $C$ تے
اوس دا متمم آپس وچ
بدل دیو۔ $g$، $C$~
وچلا کوئی وی فنکشن
ہووے۔ خیال ایہ اے کہ
جے اسی~$A$ دا فیصلہ
کر سکدے، تاں $f$ نوں
حساب کرن والے اشاریے
تے $g$ نوں حساب کرن
والے اشاریے وکھ کر
سکدے؛ فیر ایس نال
رُکّن دا مسئلہ حل
کر سکدے۔

ایہو طریقہ اے۔ ہمہ گیر جزوی
قابلِ حساب فنکشناں دی مدد نال
اسی اک فنکشن ایداں تعریف کر سکدے آں:
\[
h(x,y) \simeq
\begin{cases}
\text{غیر تعریف شدہ} & \text{جے $\cfind{x}(x) \fundefined$} \\
g(y) & \text{ہور صورت وچ۔}
\end{cases}
\]
$h$ دی قدر کڈھن لئی
پہلے $\cfind{x}(x)$
حساب کرن دی کوشش
کرو؛ جے ایہ حسابی
کارروائی رُک جائے
تاں $g(y)$ حساب
کرو؛ تے جے
{\em اوہ} کارروائی
وی رُک جائے تاں
اوہدا آؤٹ پٹ دے
دیو۔ ہور رسمی طور
اُتے ایہ لکھ سکدے
آں:
\[
h(x,y) \simeq \Proj{2}{0}(g(y),\fn{Un}(x,x)).
\]
جتھے $\Proj{2}{0}(z_0, z_1) = z_0$
$2$ آرگومنٹاں والا
پروجیکشن فنکشن اے؛
ایہ $0$-واں آرگومنٹ
واپس دیندا اے تے
قابلِ حساب اے۔

تاں $h$ جزوی قابلِ
حساب فنکشناں دی
ترکیب اے؛ سجّا پاسا
اُدوں تے صرف اُدوں
تعریف شدہ اے تے~$g(y)$
دے برابر اے جدوں
$\fn{Un}(x,x)$ تے
$g(y)$ دونوں تعریف
شدہ ہون۔

دھیان دیو، مقرر~$x$
لئی جے $\cfind{x}(x)$
غیر تعریف شدہ اے
تاں $h(x,y)$ ہر~$y$
لئی غیر تعریف شدہ
اے؛ تے جے
$\cfind{x}(x)$ تعریف
شدہ اے تاں
$h(x,y) \simeq g(y)$۔
سو $x$ دی ہر مقرر
قدر لئی یا $h(x,y)$
عین $f$ وانگ کم
کردا اے یا عین $g$
وانگ۔ $\cfind{x}(x)$
تعریف شدہ اے یا نہیں،
ایہ فیصلہ کرنا ایہناں
دو حالتاں وچوں اک
دا فیصلہ کرنا اے۔
پر ایہ اوہو سوال
اے کہ $h_x(y) \simeq h(x,y)$
نال دتا فنکشن~$C$
وچ اے یا نہیں؛ جے
$A$ قابلِ حساب ہووے
تاں اسی ایہ فیصلہ
کر سکدے۔

ہور رسمی طور اُتے،
چونکہ $h$~دو آرگومنٹاں
والا جزوی قابلِ حساب
فنکشن اے، ایس لئی
ایہ $\cfind{e}$ دے
برابر اے، جتھے $e$
اوہدا اشاریہ اے۔
$s$-$m$-$n$
قضیے مطابق اک بنیادی
بازگشتی فنکشن~$s$
اے ایداں کہ ہر~$x$
لئی $\cfind{s(e,x)}(y) = h_x(y)$۔
ہن ہر $x$ لئی،
جے $\cfind{x}(x) \fdefined$
تاں $\cfind{s(e,x)}$
اوہو فنکشن اے جو~$g$،
اس لئی $s(e,x)$،
$A$ وچ اے۔ دوجے
پاسے، جے
$\cfind{x}(x)
\uparrow$ تاں
$\cfind{s(e,x)}$
اوہو فنکشن اے جو~$f$،
اس لئی $s(e,x)$،
~$A$ وچ نہیں۔
ہور لفظاں وچ،
ہر~$x$ لئی،
$x \in K$ جے تے
صرف جے $s(e,x) \in A$۔
جے $A$ قابلِ حساب
ہووے تاں $K$~وی
قابلِ حساب ہووے گا،
جو تضاد اے۔ ایس
لئی $A$ قابلِ حساب
نہیں۔
\end{proof}

رائس دا قضیہ بہت
طاقتور اے۔ ہِٹھلا
فوری نتیجہ ایس دیاں
کجھ مثالاں دیندا اے۔
\begin{cor}
ہِٹھاں دتے سیٹ
ناقابلِ فیصلہ نیں۔
\begin{enumerate}
\item $\Setabs{x}{\text{$17$، $\cfind{x}$ دے حاصل قدراں وچ اے}}$
\item $\Setabs{x}{\text{$\cfind{x}$ مستقل اے}}$
\item $\Setabs{x}{\text{$\cfind{x}$ کُل اے}}$
\item $\Setabs{x}{\text{جدوں وی $y < y'$، $\cfind{x}(y) \fdefined$ ہووے،
    تے جے $\cfind{x}(y') \fdefined$ ہووے تاں $\cfind{x}(y) < \cfind{x}(y')$}}$
\end{enumerate}
\end{cor}

\begin{proof}
  ایہ سارے اشارییاں
  دے اوہ سیٹ نیں جیڑے
  نہ خالی نیں، نہ سارے
  اشارییاں اُتے مشتمل۔
\end{proof}
\endgroup
% /OLP-0248

% OLP-0249 content/computability/computability-theory/fixed-point-thm.tex
\begingroup
\def\olfilename{fixed-point-thm}
\expandafter\def\csname ol@sectioncs\endcsname{section}

\olfileid{cmp}{thy}{fix}
\olsection{مقرر نقطے دا قضیہ}

آؤ رُکّن دے مسئلے
نوں فیر ویکھئیے۔
عارضی علامت وجوں
$\gn{\cfind{x}(y)}$
نوں $\tuple{x, y}$
لئی لکھو؛ ایہنوں
$\cfind{x}(y)$ دی
قدر دے اک «ناں»
وجوں سمجھو۔ ایس
علامت نال اسی اپنے اوہناں ثبوتاں وچوں اک
ہور لفظاں وچ بیان کر سکدے آں جیہڑے
وکھاؤندے نیں کہ رُکّن دا مسئلہ ناقابلِ فیصلہ اے۔

سوال: کی اک قابلِ حساب
فنکشن $h$ موجود اے
جیہدی ہِٹھلی خاصیت
ہووے؟ ہر $x$ تے
$y$ لئی:
\[
h(\gn{\cfind{x}(y)}) =
\begin{cases}
1 & \text{جے $\cfind{x}(y) \fdefined$} \\
0 & \text{ہور صورت وچ۔}
\end{cases}
\]

جواب: نہیں؛ ورنہ
ہِٹھاں دتا جزوی فنکشن
\[
g(x) \simeq
\begin{cases}
0 & \text{جے $h(\gn{\cfind{x}(x)}) = 0$} \\
\text{غیر تعریف شدہ} & \text{ہور صورت وچ}
\end{cases}
\]
قابلِ حساب ہووے گا
تے اوس دا کوئی
اشاریہ~$e$ ہووے گا۔
پر فیر:
\[
\cfind{e}(e) \simeq
\begin{cases}
0 & \text{جے $h(\gn{\cfind{e}(e)}) = 0$} \\
\text{غیر تعریف شدہ} & \text{ہور صورت وچ،}
\end{cases}
\]
ایس صورت وچ
$\cfind{e}(e)$
اُدوں تے صرف
اُدوں تعریف شدہ
اے جدوں ایہ
تعریف شدہ نہیں؛
ایہ تضاد اے۔

ہن $\cfind{e}$ والی
مساوات ویکھو۔
اک لحاظ نال ایتھے
اپنے آپ دا حوالہ
اے: اسی $\cfind{e}(e)$
دی قدر نوں اک خاص
طریقے نال
$\gn{\cfind{e}(e)}$
اُتے منحصر بنایا
اے۔ مقرر نقطے دا
قضیہ آکھدا اے کہ
عام طور اُتے اسی
ایہ {\em کر سکدے}
آں؛ صرف تضاد ثابت
کرن لئی نہیں۔

\olref{lem:fixed-equiv}
مقرر نقطے دا قضیہ
بیان کرن دے دو
ہم ارز طریقے دیندا
اے۔ منطقی لحاظ
نال ایہ بیان دونوں
سچ ہون کارن
ہم ارز ہوندے نیں؛
پر ساڈی مراد ایہ
اے کہ ہر اک دوجے
توں سدھا نکلدا اے،
ایس لئی دونوں اکّو قضیے نوں
بیان کرن دے وکھو وکھرے طریقے نیں۔

\begin{lem}
\ollabel{lem:fixed-equiv}
ہِٹھلے بیان ہم ارز نیں:
\begin{enumerate}
\item ہر جزوی قابلِ
  حساب فنکشن $g(x,y)$
  لئی اک اشاریہ~$e$
  اے ایداں کہ ہر~$y$
  لئی:
\[
\cfind{e}(y) \simeq g(e,y).
\]
\item ہر قابلِ حساب
  فنکشن~$f(x)$ لئی
  اک اشاریہ~$e$ اے
  ایداں کہ ہر~$y$
  لئی:
\[
\cfind{e}(y) \simeq \cfind{f(e)}(y).
\]
\end{enumerate}
\end{lem}

\begin{proof}
$(1) \Rightarrow (2)$:
$f$ دتا ہووے تاں
$g$ نوں $g(x,y) \simeq
\fn{Un}(f(x),y)$
دی مدد نال تعریف کرو۔
(1) نال ایداں دا
اشاریہ~$e$ لو کہ
ہر~$y$ لئی:
\begin{align*}
\cfind{e}(y) & = \fn{Un}(f(e),y) \\
& = \cfind{f(e)}(y).
\end{align*}

$(2) \Rightarrow (1)$:
$g$ دتا ہووے تاں
$s$-$m$-$n$ قضیہ
ورت کے $f$ لو
ایداں کہ ہر $x$
تے~$y$ لئی
$\cfind{f(x)}(y) \simeq g(x,y)$۔
(2) نال ایداں دا
اشاریہ~$e$ لو کہ:
\begin{align*}
\cfind{e}(y) & = \cfind{f(e)}(y) \\
& = g(e,y).
\end{align*}
ایس نال ثبوت
مکدا اے۔
\end{proof}

\begin{explain}
بیان (1) ثابت کرن توں پہلاں، جیہدے توں
(2) وی نکلدا اے، سوچو ایہ کِنّا حیران کن اے۔
$e$ نوں کمپیوٹر
پروگرام سمجھو۔
بیان (1) آکھدا
اے کہ کوئی وی
جزوی قابلِ حساب
$g(x,y)$ دتا
ہووے تاں تسی
اک ایداں دا
کمپیوٹر پروگرام
$e$ لبھ سکدے
او جیڑا
$g_e(y) \simeq g(e,y)$
حساب کرے۔
یعنی تسی ایداں
دا پروگرام لبھ
سکدے او جیڑا اوس فنکشن دی قدر
کڈھدا اے جیہڑا ایس پروگرام دا ای حوالہ دیندا اے۔
\end{explain}

\begin{thm}
\olref{lem:fixed-equiv}
دے دونوں بیان
سچ نیں۔ خاص طور
اُتے، ہر جزوی
قابلِ حساب فنکشن
$g(x,y)$ لئی اک
اشاریہ~$e$ اے
ایداں کہ ہر~$y$
لئی:
\[
\cfind{e}(y) \simeq g(e,y).
\]
\end{thm}

\begin{proof}
درکار اجزا اُتّے
رُکّن دے مسئلے
دی بحث وچ پہلے
ای موجود نیں۔
$\fn{diag}(x)$
اک قابلِ حساب
فنکشن ہووے جیڑا
ہر $x$ لئی
$f_x(y) \simeq \cfind{x}(x,y)$
دا اشاریہ دیندا
اے، یعنی:
\[
\cfind{\fn{diag}(x)}(y) \simeq \cfind{x}(x,y).
\]
$\fn{diag}$ نوں
ایداں دا فنکشن
سمجھو جیڑا دو
آرگومنٹاں والے
فنکشن دے پروگرام
نوں اک آرگومنٹ
والے فنکشن دے
پروگرام وچ بدلدا
اے۔ ایس لئی اصل فنکشن دے پہلے آرگومنٹ
وجوں اوہدا اپنا پروگرام، یعنی اوہدا اشاریہ،
مقرر کردے آں۔ $\fn{diag}$
دی رسمی تعریف
ایداں بن سکدی
اے: پہلے $s$
نوں تعریف کرو:
\[
s(x,y) \simeq \fn{Un}^2(x,x,y),
\]
جتھے $\fn{Un}^2$
تِن آرگومنٹاں
والا فنکشن اے
جو دو آرگومنٹاں
والے جزوی قابلِ
حساب فنکشناں
لئی ہمہ گیر اے۔
فیر $s$-$m$-$n$
قضیے نال بنیادی
بازگشتی فنکشن
$\fn{diag}$ لبھدا
اے جیہڑا ایہ
شرط پوری کردا
اے:
\[
\cfind{\fn{diag}(x)}(y) \simeq s(x,y).
\]

ہن فنکشن $l$
نوں ایداں
تعریف کرو:
\[
l(x,y) \simeq g(\fn{diag}(x),y).
\]
تے $\gn{l}$،
$l$ دا اک
اشاریہ ہووے۔
آخر $e = \fn{diag}(\gn{l})$
من لو۔ فیر
ہر $y$ لئی:
\begin{align*}
\cfind{e}(y) & \simeq \cfind{\fn{diag}(\gn{l})}(y) \\
& \simeq \cfind{\gn{l}}(\gn{l}, y) \\
& \simeq l(\gn{l}, y) \\
& \simeq g(\fn{diag}(\gn{l}),y) \\
& \simeq g(e, y),
\end{align*}
جیویں چاہیدا سی۔
\end{proof}

\begin{explain}
ایتھے ہو کی رہیا
اے؟ من لو تہاڈا
کم اک ایداں دا
کمپیوٹر پروگرام
لکھنا اے جیڑا
اپنا متن آپ چھاپے۔
ہور من لو تسی
ایہو جیہی پروگرامنگ
زبان ورتدے او جیہدی
لڑیاں والے فنکشناں
دی لائبریری بھرپور
تے انوکھی اے۔
خاص کر، فرض کرو
ایس زبان وچ
اک فنکشن $\fn{diag}$
اے جیڑا ایداں کم
کردا اے: انپٹ
لڑی~$s$ دتی ہووے
تاں $\fn{diag}$ اوس~$s$ وچ «x» نشان
دیاں ساریاں تھاواں لبھدا اے تے ہر وار
اوس دی تھاں اصل لڑی دا اقتباسی
نشانیاں وچ بند روپ پا دیندا اے۔ مثلاً،
ہِٹھلی لڑی انپٹ
ہووے:
\begin{quote}
\begin{verbatim}
hello x world
\end{verbatim}
\end{quote}
تاں فنکشن ہِٹھلا
نتیجہ دیندا اے:
\begin{quote}
\begin{verbatim}
hello 'hello x world' world
\end{verbatim}
\end{quote}
ایس صورت وچ
مطلوبہ پروگرام
لکھنا سوکھا اے؛
تسی جانچ سکدے
او کہ:
\begin{quote}
\begin{verbatim}
print(diag('print(diag(x))'))
\end{verbatim}
\end{quote}
ایہ کم کردا اے۔
C++ تے Java
ورگیاں عام
پروگرامنگ زباناں
وچ وی اوہی خیال
کم کردا اے،
بھاویں اوہدی
بناوٹ ہور اوکھی
ہووے۔

مقرر نقطے دے
قضیے دے ثبوت
توں اسی ہن
صرف دو کدم
دور آں۔
من لو پرنٹ
فنکشن دی اک
صورت $\fn{print}(x,y)$
لڑی $x$ تے
اک ہور عددی
آرگومنٹ $y$
لیندی اے، تے
لڑی $x$ نوں
$y$ واری
چھاپدی اے۔
تاں ہِٹھلا
«پروگرام»:
\begin{quote}
\begin{verbatim}
getinput(y);
print(diag('getinput(y); print(diag(x), y)'), y)
\end{verbatim}
\end{quote}
انپٹ $y$ اُتے
اپنا متن $y$
واری چھاپدا
اے۔ $\fn{getinput}$---
$\fn{print}$---
$\fn{diag}$ دے
ڈھانچے دی تھاں
کوئی وی فنکشن
$g(x,y)$ رکھو
تاں ملدا اے:
\begin{quote}
\begin{verbatim}
g(diag('g(diag(x), y)'), y)
\end{verbatim}
\end{quote}
ایہ پروگرام
انپٹ $y$
اُتے $g$
نوں ایس پروگرام دے اپنے متن
تے $y$ اُتے
چلاوندا اے۔
«اقتباسی نشانیاں
وچ بند کرن»
نوں «ایس دا
اشاریہ ورتن»
وجوں سمجھو
تاں اُتّے
والا ثبوت
ملدا اے۔

فی الحال جے
تسی ایس ثبوت
نوں رسمی چال
یا کالا جادو
سمجھنا چاہو
تاں کوئی گل
نہیں۔ پر تہانوں اُتّے دتی دلیل دیاں
تفصیلاں آپ دوبارہ بنا سکنا چاہیدا اے۔
جدوں اسی
نامکملیت دے
قضیے، تے
اوہناں نال
متعلق «مقرر
نقطے دا قضیہ»،
ثابت کراں گے
تاں اسی ایس گل نوں سمجھن دے ہور
طریقیاں اُتے گل کراں گے کہ ایہ کم کیوں کردا اے۔
\end{explain}

\begin{tagblock}{lambda}
\begin{digress}
ایہو خیال
«مقرر نقطہ»
کمبینیٹر وی
دیندا اے۔
من لو تہاڈے
کول لیمبڈا اصطلاح
$g$ اے، تے
تسی ہور اصطلاح
$k$ چاہندے
او جیہدی
خاصیت ایہ
ہووے کہ $k$، $\beta$ دے لحاظ نال
$gk$ دے ہم ارز ہووے۔
ساڈے علامتی
قاعدیاں مطابق
ایہ اصطلاحاں
تعریف کرو:
\[
\fn{diag}(x) = xx
\]
تے
\[
l(x) = g(\fn{diag}(x))
\]
ہور لفظاں وچ،
$l$ اصطلاح
$\lambd[x][g(xx)]$
اے۔ $k$ نوں
$ll$ اصطلاح
من لو۔ فیر:
\begin{align*}
k & = (\lambd[x][g(xx)])(\lambd[x][g(xx)]) \\
& \red  g((\lambd[x][g(xx)])(\lambd[x][g(xx)])) \\
& = gk.
\end{align*}
جے ایہ لو:
\[
Y = \lambd[g][((\lambd[x][g(xx)])(\lambd[x][g(xx)]))]
\]
تاں $Yg$ تے
$g(Yg)$ اک
سانجھی اصطلاح
تک گھٹدے نیں؛
ایس لئی $Yg \equiv_\beta
g(Yg)$۔ ایہنوں
«کری دا کمبینیٹر»
آکھدے نیں۔
جے ایس دی
تھاں ایہ لو:
\[
Y = (\lambd[xg][g(xxg)])(\lambd[xg][g(xxg)])
\]
تاں واقعی $Yg$،
$g(Yg)$ ول
گھٹدا اے، جیڑی
ہور مضبوط گل
اے۔ $Y$ دی
ایس دوجی
صورت نوں
«ٹیورنگ دا
کمبینیٹر»
آکھدے نیں۔
\end{digress}
\end{tagblock}
\endgroup
% /OLP-0249

% OLP-0250 content/computability/computability-theory/application-fixed-point.tex
\begingroup
\def\olfilename{application-fixed-point}
\expandafter\def\csname ol@sectioncs\endcsname{section}

\olfileid{cmp}{thy}{apf}
\olsection{مقرر نقطے دے قضیے دی ورتوں}

مقرر نقطے دے قضیے نال اسی
جزوی قابلِ حساب فنکشناں نوں اوہناں
دے اپنے اشارییاں دے حوالے نال
تعریف کر سکدے آں۔ مثلاً،
اک اشاریہ $e$ لبھ
سکدا اے ایداں کہ
ہر $y$ لئی:
\[
\cfind{e}(y) = e + y.
\]
ہور مثال وچ،
مقرر نقطے دے قضیے
دا ثبوت ورت کے Java
یا C++ وچ ایداں
دا پروگرام بنایا
جا سکدا اے جیڑا
اپنا متن آپ چھاپے۔

یاد رکھو: جے ہر
$e$ لئی $W_e$
نوں $\cfind{e}$
دا دائرۂ تعریف
منئیے، تاں سلسلہ
$W_0$، $W_1$،
$W_2$،~\dots
قابلِ حساب گنتی
وار فہرست بنن جوگ
سیٹاں دی فہرست
دیندا اے۔ ایہناں
وچوں کجھ سیٹ
قابلِ حساب نیں۔
سوال ایہ اے:
کی کوئی الگورتھم
اک قدر $x$
انپٹ لے کے،
جے $W_x$ قابلِ
حساب ہووے،
اوس دے خصوصیتی
فنکشن دا اشاریہ
واپس دے سکدا اے؟
جواب «نہیں» اے؛
ایداں دا کوئی
الگورتھم نہیں:

\begin{thm}
کوئی جزوی قابلِ حساب
فنکشن $f$ نہیں جیہدی ایہ خاصیت ہووے:
جدوں وی $W_e$
قابلِ حساب ہووے،
تاں $f(e)$
تعریف شدہ ہووے
تے $\cfind{f(e)}$
اوس دا خصوصیتی
فنکشن ہووے۔
\end{thm}

\begin{proof}
% PNB-CMP-025: proof must cover partial computable f and undefined-index cases.
من لو $f$ کوئی وی
جزوی قابلِ حساب
فنکشن اے؛ اسی اک
$e$ بناواں گے
ایداں کہ $W_e$
قابلِ حساب ہووے،
پر $\cfind{f(e)}$
اوس دا خصوصیتی
فنکشن نہ ہووے،
یا مطلوبہ اشاریہ
ہی نہ ملے۔
مقرر نقطے دا
قضیہ ورت کے
اک اشاریہ $e$
لبھ سکدا اے
ایداں کہ:
\[
\cfind{e}(y) \simeq
\begin{cases}
0 & \text{جے $y=0$ تے $\cfind{f(e)}(0) \fdefined = 0$} \\
\text{غیر تعریف شدہ} & \text{ہور صورت وچ۔}
\end{cases}
\]
یعنی $e$، مقرر
نقطے دا قضیہ
ہِٹھاں دتے
فنکشن اُتے
لا کے ملدا
اے:
\[
g(x,y) \simeq
\begin{cases}
0 & \text{جے $y=0$ تے $\cfind{f(x)}(0) \fdefined = 0$} \\
\text{غیر تعریف شدہ} & \text{ہور صورت وچ۔}
\end{cases}
\]
غیر رسمی طور اُتے ایویں ویکھ سکدے آں
کہ $g$ جزوی قابلِ حساب اے:
انپٹ $x$ تے
$y$ اُتے الگورتھم
پہلے جانچدا اے
کہ $y$،~$0$
دے برابر اے
یا نہیں۔ جے
اے، تاں اوہ
$f(x)$ حساب
کردا اے،
فیر ہمہ گیر مشین
نال $\cfind{f(x)}(0)$
حساب کردا اے۔
جے ایہ آخری
حساب رُک کے
~$0$ واپس
دیوے، تاں
الگورتھم~$0$
دے کے رُکدا
اے؛ ہور صورت
وچ نہیں رُکدا۔
جے اشاریہ کڈھن والا پہلا حساب
ہی نہ رُکے،
تاں وی پورا
الگورتھم نہیں
رُکدا۔

ہن دھیان دیو:
جے $\cfind{f(e)}(0)$
تعریف شدہ اے
تے $0$ دے
برابر اے، تاں
$\cfind{e}(y)$
عین اُدوں
تعریف شدہ
اے جدوں $y$،
$0$ دے
برابر ہووے؛
اس لئی
$W_e = \{ 0 \}$۔
جے $\cfind{f(e)}(0)$
غیر تعریف شدہ
اے، یا تعریف
شدہ اے پر
$0$ دے برابر
نہیں، تاں
$W_e = \emptyset$۔
جے ایہ اشاریہ موجود ہووے، دونوں
صورتاں وچ $\cfind{f(e)}$،
$W_e$ دا
خصوصیتی فنکشن
نہیں: انپٹ $0$
اُتے صفر آوے
تاں جواب اک
ہونا چاہیدا؛
غیر صفر آوے
تاں جواب صفر
ہونا چاہیدا۔
جے اشاریہ
دین والا حساب
ہی تعریف شدہ
نہ ہووے، تاں
مطلوبہ اشاریہ
نہیں ملدا؛
تے جے اوہ
اشاریہ ملے
پر اوس دا
فنکشن صفر
اُتے غیر تعریف
شدہ ہووے،
تاں اوہ کُل
خصوصیتی فنکشن
نہیں۔
\end{proof}
\endgroup
% /OLP-0250

% OLP-0251 content/computability/computability-theory/def-functions-self-reference.tex
\begingroup
\def\olfilename{def-functions-self-reference}
\expandafter\def\csname ol@sectioncs\endcsname{section}

\olfileid{cmp}{thy}{slf}
\olsection{اپنے آپ دے حوالے نال فنکشناں دی تعریف}

عام طور اُتے فنکشناں نوں
اپنے آپ دے حوالے نال
تعریف کر سکنا فائدہ مند
ہوندا اے۔ مثال دے طور
اُتے، جے قابلِ حساب
فنکشن $k$، $l$ تے~$m$
دتے ہون، تاں مقرر نقطے
دا معاون قضیہ دسدا اے
کہ اک جزوی قابلِ حساب
فنکشن~$f$ موجود اے
جیڑا ہر~$y$ لئی ہِٹھلی
مساوات پوری کردا اے:
\[
f(y) \simeq
\begin{cases}
k(y) & \text{جے $l(y) = 0$} \\
f(m(y)) & \text{ہور صورت وچ۔}
\end{cases}
\]
ہور خاص طور اُتے، $f$ لبھن لئی
پہلاں ایہ تعریف کرو:
\[
g(x,y) \simeq
\begin{cases}
k(y) & \text{جے $l(y) = 0$} \\
\cfind{x}(m(y)) & \text{ہور صورت وچ}
\end{cases}
\]
تے فیر مقرر نقطے دا
معاون قضیہ ورت کے اک
اشاریہ~$e$ لبھیا جاندا
اے ایداں کہ
$\cfind{e}(y) = g(e,y)$۔

اک ٹھوس مثال ویکھو۔
«سب توں وڈا مشترک پورا
ونڈن والا عدد» \pnbLatin{(greatest common divisor)} دا فنکشن
$\fn{gcd}(u,v)$ ایداں
تعریف کیتا جا سکدا اے:
\[
\fn{gcd}(u,v) \simeq
\begin{cases}
v & \text{جے $u = 0$} \\
\fn{gcd}(\fn{mod}(v, u), u) & \text{ہور صورت وچ}
\end{cases}
\]
جتھے $\fn{mod}(v, u)$
توں مراد $v$ نوں~$u$
نال ونڈن مگروں بچیا
باقی اے۔ مقرر نقطے
دے معاون قضیے نوں
لا کے پتا لگدا اے کہ
$\fn{gcd}$ جزوی قابلِ
حساب اے۔ (حقیقت وچ،
$y$ نوں جوڑے
$\tuple{u, v}$ دا کوڈ
من کے ایہنوں اُتلی
صورت وچ لکھ سکدے آں۔)
ایس توں بعد~$u$ اُتے
استقرا نال ثابت ہوندا اے کہ
$\fn{gcd}$ دراصل کُل اے۔

بے شک، اپنے آپ دے
حوالے والیاں تعریفاں
پچھلیاں مثالاں توں
کِتے ودھ پیچیدہ وی
بنا سکدے آں۔ بہتیاں
پروگرامنگ زباناں اک
نہ اک طریقے نال فنکشناں
نوں اپنے آپ دے حوالے
نال تعریف کرن دی اجازت
دیندیاں نیں۔ پر دھیان
دیو: ایہ گل، اوہدے
مقابلے وچ قدرے گھٹ
حیران کن اے کہ اسی
اک فنکشن نوں اوس فنکشن
دا حساب کرن والے الگورتھم
دے \emph{اشاریے} دے
حوالے نال تعریف کر
سکدے آں؛ مقرر نقطے
دا قضیہ اپنے پورے عموم
وچ ایہ دوجی سہولت
دیندا اے۔
\endgroup
% /OLP-0251

\OLEndChapterHook
\endgroup
% /OLP-0228

\OLEndPartHook
\endgroup
% /OLP-0208

% OLP-0252 content/turing-machines/turing-machines.tex
\begingroup
\def\olfilename{turing-machines}
\expandafter\def\csname ol@sectioncs\endcsname{section}

\olpart{tur}{ٹیورنگ مشیناں}

% OLP-0253 content/turing-machines/machines-computations/machines-computations.tex
\begingroup
\def\olfilename{machines-computations}
\expandafter\def\csname ol@sectioncs\endcsname{section}

\olchapter{tur}{mac}{ٹیورنگ مشیناں دیاں حسابی کارروائیاں}

% OLP-0254 content/turing-machines/machines-computations/introduction.tex
\begingroup
\def\olfilename{introduction}
\expandafter\def\csname ol@sectioncs\endcsname{section}

\olfileid{tur}{mac}{int}
\olsection{تعارف}

اک فنکشن، مثلاً $\Nat$ توں $\Nat$ تک،
\emph{قابلِ حساب} ہون دا کی مطلب اے؟
ابتدائی جواباں وچوں اک، تے سب توں معروف
جواب، ایہ اے کہ جے اک ٹیورنگ مشین
اوس فنکشن دا حساب کر سکدی اے تاں
اوہ قابلِ حساب اے۔ ایہ تصور ایلن ٹیورنگ
\pnbLatin{(Alan Turing)} نے 1936 وچ پیش کیتا۔
ٹیورنگ مشیناں \emph{حساب دے اک نمونے}
دی مثال نیں: ایہ «حساب کرن دے طریقے»
دے خیال نوں ریاضیاتی طور اُتے ٹھیک ٹھیک
تعریف کرن دا اک طریقہ نیں۔ اوس خیال
دا بالکل کی مطلب اے، ایس اُتے بحث
ہوندی اے؛ پر عام اتفاق اے کہ ٹیورنگ
مشیناں حساب کرن دے طریقے مقرر کرن
دا اک طریقہ نیں۔ بھاویں «ٹیورنگ مشین»
سن کے حرکت کرن والے پرزیاں دی مادی
مشین دا خیال آندا اے، پر ٹھیک ٹھیک آکھیے
تاں ٹیورنگ مشین خالص ریاضیاتی بناوٹ اے۔
ایہ حساب کرن دے طریقے دا اوہ ریاضیاتی روپ اے
جیہدے اُتے اصلی مشین دیاں عملی بندشاں نہیں لگدیاں۔ مثال لئی،
اسی ٹیورنگ مشین دے وقت یا یادداشت
دی ضرورت اُتے کوئی حد نہیں رکھدے:
اوہ ایداں دا حساب وی کر سکدی اے
جیدے لئی کائنات دے ایٹماں دی گنتی
توں ودھ ذخیرہ کرن دی تھاں یا قدم
درکار ہون۔

\begin{explain}
شاید ٹیورنگ مشین نوں اک خاص قسم
دے فرضی آلے دا پروگرام سمجھنا
سب توں چنگا اے۔ ایس آلے وچ اک
\emph{پٹی} (tape) تے اک
\emph{پڑھن-لکھن والا سر} (head)
ہوندا اے۔ ساڈی ٹیورنگ مشین دی
صورت وچ پٹی اک پاسے، یعنی سجّے
پاسے، ول لامتناہی اے۔ ایہ
\emph{خانیاں} وچ ونڈی ہوئی اے؛
ہر خانے وچ کوئی نشان ہو سکدا اے
جیڑا اک محدود \emph{حرفاں دے مجموعے}
(alphabet) وچوں لیا جاندا اے۔
ایس مجموعے وچ وکھو وکھرے
نشان کوئی وی محدود گنتی وچ
رکھ سکدے آں،
پر زیادہ تر ساڈے لئی تِن ہی کافی
ہون گے: $\TMendtape$،
$\TMblank$ تے $\TMstroke$۔
جدوں آلہ شروع ہوندا اے،
پٹی خالی ہوندی اے، یعنی
ہر خانے وچ $\TMblank$ ہوندا اے؛
بس سب توں کھبے خانے وچ
$\TMendtape$ ہوندا اے تے
محدود گنتی دے کجھ خانے
\emph{انپٹ} رکھدے نیں۔
ہر ویلے آلہ محدود گنتی
دیاں \emph{حالتاں} وچوں
کسے اک حالت وچ ہوندا اے۔
شروع وچ سر سب توں کھبے
خانے نوں پڑھدا اے تے آلہ
اک مقرر کیتی \emph{ابتدائی حالت}
وچ ہوندا اے۔ آلے دے چلن
دے ہر قدم اُتے، سر ہِٹھ
والے خانے دا نشان، آلے
دی موجودہ حالت تے ٹیورنگ
مشین دا پروگرام رل کے
اگلا عمل طے کردے نیں۔
ٹیورنگ مشین دا پروگرام
اک جزوی فنکشن نال دتا جاندا
اے: ایہ اک حالت~$q$
تے اک نشان~$\sigma$
انپٹ لے کے اک تِکا
$\tuple{q', \sigma',
  D}$ واپس دیندا اے۔
جدوں آلہ حالت $q$
وچ ہووے تے نشان $\sigma$
پڑھے، تاں اوہ موجودہ
خانے دا نشان $\sigma'$
نال بدلدا اے؛ $D$
دی قدر $\TMleft$،
$\TMright$ یا $\TMstay$
ہون مطابق سر کھبے ول سرکدا اے،
سجّے ول سرکدا اے یا اوتھے ای رہندا اے؛
تے آلہ حالت~$q'$
وچ چلا جاندا اے۔

مثال لئی، \olref{fig:tm}
وچ دکھائی حالت ویکھو۔
\begin{figure}
  \olasset{\olpath/assets/diagrams/turing-machine.tikz}
  \caption{اک ٹیورنگ مشین اپنا پروگرام چلا رہی اے۔}
  \ollabel{fig:tm}
\end{figure}
ٹیورنگ مشین دی پٹی دے
دکھائی دیندے حصے وچ
سب توں کھبے خانے اُتے
پٹی دے اختتام دا نشان
$\TMendtape$ اے؛ اوس دے
بعد تِن $1$، اک $0$
تے ہور چار $1$ نیں۔
سر کھبے پاسے توں تیجے
خانے نوں پڑھ ریا اے؛
اوتھے اک~$1$ اے تے
آلہ حالت~$q_1$ وچ اے۔
اسی آکھدے آں، «مشین
اک $1$ پڑھ رہی اے
جدوں اوہ حالت~$q_1$
وچ اے۔» جے ٹیورنگ
مشین دا پروگرام انپٹ
$\tuple{q_1, 1}$ لئی
تِکا $\tuple{q_2,
0, \TMstay}$ واپس دیوے،
تاں مشین تیجے خانے
والے~$1$ نوں اک~$0$
نال بدل دیوے گی،
پڑھن-لکھن والے سر نوں
اوتھے ہی رکھے گی تے
حالت~$q_2$ وچ چلی
جاوے گی۔ فیر جے پروگرام
تِکا $\tuple{q_3, 0, \TMright}$،
انپٹ $\tuple{q_2, 0}$
لئی واپس دیوے، تاں مشین
موجودہ~$0$ اُتے فیر
اک~$0$ لکھے گی
(یعنی سر ہِٹھ پٹی دا نشان
اوہو رہے گا)، سر نوں اک خانہ
سجّے ول سرکا دے گی تے حالت~$q_3$
وچ آ جاوے گی۔ ایس
طرح اگے چلدا اے۔

اسی آکھدے آں کہ مشین
\emph{رُک جاندی اے} جدوں
اوہ کسے حالت $q_n$ وچ
ایہو جیہا نشان $\sigma$ پڑھے جیہدے لئی
$\tuple{q_n, \sigma}$
دا کوئی حکم موجود نہ ہووے؛
یعنی انپٹ
$\tuple{q_n,\sigma}$
لئی حالتی تبدیلی
دا فنکشن غیر تعریف شدہ
ہووے۔ دوجے لفظاں وچ،
مشین کول کرن لئی
کوئی حکم نہیں رہندا،
تے ایس ویلے اوہ
چلنا بند کر دیندی اے۔
کدی کدی رُکّن نوں
اک خاص رُکّن والی
حالت~$h$ نال وی
دکھایا جاندا اے۔
ایہ اگے ودھ تفصیل
نال وکھایا جاوے گا۔
\end{explain}

\begin{digress}
ٹیورنگ دے مقالے
«قابلِ حساب اعداد بارے» \pnbLatin{(On computable numbers)}
دی خوبی ایہ اے کہ
اوہ صرف رسمی تعریف
نہیں دیندا، بلکہ ایہ
دلیل وی دیندا اے
کہ تعریف قابلِ حساب ہون دے
بدیہی تصور نوں بیان کردی اے۔ تعریف توں
صاف ہونا چاہیدا اے
کہ ٹیورنگ مشین نال
قابلِ حساب ہر فنکشن
بدیہی معنی وچ قابلِ
حساب اے۔ ٹیورنگ
ایہ اُلٹی گل لئی،
یعنی جیہڑا وی فنکشن
اسی فطری طور اُتے
قابلِ حساب سمجھدے
آں اوہ وی ایداں
دی مشین نال قابلِ
حساب اے، تِن طرح
دیاں دلیلاں دیندا
اے۔ اوس دے اپنے
لفظاں وچ اوہ ایہ نیں:
\begin{enumerate}
\item بدیہی سمجھ نوں سدھا بنیاد بنانا۔
\item دو تعریفاں دی ہم ارزیت
  دا ثبوت (جے نویں تعریف
  بدیہی سمجھ نال زیادہ میل کھاندی ہووے)۔
\item قابلِ حساب اعداد دیاں وڈیاں جماعتاں دیاں مثالاں دینا۔
\end{enumerate}
ساڈا مقصد قابلِ حساب ہون
دے تصور نوں «اصولی
طور اُتے» تعریف کرن
دی کوشش اے، یعنی
وقت تے تھاں دیاں
عملی حداں نوں لحاظ
وچ لیائے بغیر۔ بے شک،
قابلِ حساب ہون دی سب توں وسیع
تعریف مقرر ہو جائے تاں محدود
وسیلیاں نال حساب نوں
وی ویکھ سکدے آں؛
ایہ «حسابی پیچیدگی»
\pnbLatin{(computational complexity)}
نوں کہے جاندے
مضمون دا مرکز اے۔
\end{digress}

\begin{history}
ایلن ٹیورنگ \pnbLatin{(Alan Turing)}
نے 1936 وچ ٹیورنگ
مشیناں ایجاد کیتیاں۔
اوس ویلے اوس دی
دلچسپی پہلے درجے دی
منطق دے قابلِ فیصلہ
ہون وچ سی، پر ایس
مقالے نوں کمپیوٹر
ڈیزائن دیاں بنیاداں
بابت فیصلہ کن مقالہ
وی آکھیا گیا اے۔
مقالے وچ ٹیورنگ دی
توجہ قابلِ حساب
حقیقی اعداد اُتے
اے، یعنی اوہ حقیقی
اعداد جِنہاں دے
اعشاری پھیلاؤ قابلِ
حساب نیں؛ پر اوہ
ایہ وی دسدا اے
کہ اپنے تصورات نوں
قدرتی اعداد اُتے
قابلِ حساب فنکشناں
وغیرہ لئی ڈھالنا
اوکھا نہیں۔ یاد رکھو
کہ ایہ پہلا کم کرن
والا عام مقصد دا
کمپیوٹر 1941 وچ
بنن توں پورے پنج
سال پہلے دی گل اے؛
اوہ کمپیوٹر جرمن
کونراڈ زُوزے \pnbLatin{(Konrad Zuse)}
نے اپنے ماں پیو دے گھر دی
بیٹھک وچ بنایا سی۔ ایہ
ٹیورنگ تے اوس دے
بلیچلے پارک \pnbLatin{(Bletchley Park)}
دے ساتھیواں ولوں
کوڈ توڑن والا کولوسس
(Colossus) بناؤن
(1943) توں ست سال،
امریکی ENIAC (1945)
توں نو سال، مانچسٹر
وچ پہلا برطانوی
عام مقصد دا کمپیوٹر
— Manchester Small-Scale
Experimental Machine —
بنن (1948) توں باراں
سال، تے امریکیواں
ولوں BINAC دی پہلی
آزمائش (1949) توں
تیرہ سال پہلے سی۔
مانچسٹر SSEM نوں
پہلا ایسا کمپیوٹر
ہون دا امتیاز حاصل
اے جیڑا پروگرام
اپنی یادداشت وچ
رکھدا سی؛ ایس توں
پہلیاں مشیناں نوں
ہر نویں کم لئی
ہتھ نال تارا دوبارہ
جوڑݨے پیندے سن۔
\end{history}
\endgroup
% /OLP-0254

% OLP-0255 content/turing-machines/machines-computations/representing-tms.tex
\begingroup
\def\olfilename{representing-tms}
\expandafter\def\csname ol@sectioncs\endcsname{section}

\olfileid{tur}{mac}{rep}
\olsection{ٹیورنگ مشیناں دی نمائندگی}

\begin{explain}
ٹیورنگ مشیناں نوں \emph{حالتی خاکیاں}
\pnbLatin{(state diagrams)} نال تصویری شکل وچ
دکھایا جا سکدا اے۔ ایہناں خاکیاں
وچ حالتاں دے خانے تیراں نال جڑے
ہوندے نیں۔ جیویں توقع اے، ہر حالتی خانہ مشین دی
اک حالت دکھاندا اے۔ ہر تیر اوہ
حکم دکھاندا اے جیڑا اوس حالت توں
چلایا جا سکدا اے؛ حکم دی تفصیل
متعلقہ تیر دے اُتے یا تھلّے
لکھی ہوندی اے۔ ہِٹھلی مشین ویکھو:
ایس دیاں صرف دو اندرونیاں
حالتاں $q_0$ تے $q_1$ نیں،
تے اک حکم اے:
\[
\begin{tikzpicture}[->,>=stealth',shorten >=1pt,auto,node distance=2.8cm,
                    semithick]
  \tikzstyle{every state}=[fill=none,draw=black,text=black]

  \node[initial,state]   (A)              {$q_0$};
  \node[state]   (B) [right of=A] {$q_1$};

  \path (A) edge  node {\TMtrans{\TMblank}{\TMstroke}{\TMright}} (B);
\end{tikzpicture}
\]
یاد کرو کہ ٹیورنگ مشین کول
پڑھن-لکھن والا سر تے اک پٹی ہوندی اے
جیہدے اُتے انپٹ لکھیا ہوندا اے۔ ایس
حکم نوں ایداں پڑھو:
\emph{جے اک~$\TMblank$
حالت $q_0$ وچ پڑھیا جاوے،
تاں اک~$\TMstroke$ لکھو،
سر نوں سجّے ول سرکاؤ، تے حالت
$q_1$ وچ چلے جاؤ}۔
ایہ اوہو گل اے کہ
حالتی تبدیلی دا فنکشن
$\tuple{q_0, \TMblank}$
نوں $\tuple{q_1, \TMstroke,
\TMright}$ ول بھیجدا اے۔
\end{explain}

\begin{ex}
\emph{جفت مشین}: ہِٹھلی ٹیورنگ
مشین اُدوں تے صرف اُدوں
رُکدی اے جدوں پٹی اُتے
$\TMstroke$ دی گنتی جفت
ہووے (ایس مفروضے ہِٹھ کہ
سارے $\TMstroke$ پہلے
$\TMblank$ توں پہلاں آندے
نیں)۔
\[
\begin{tikzpicture}[->,>=stealth',shorten >=1pt,auto,node distance=2.8cm,
                    semithick]
  \tikzstyle{every state}=[fill=none,draw=black,text=black]

  \node[initial,state]         (A)              {$q_0$};
  \node[state]         (B) [right of=A] {$q_1$};

  \path (A) edge [bend left] node {\TMtrans{\TMstroke}{\TMstroke}{\TMright}} (B)
        (B) edge [loop above] node {\TMtrans{\TMblank}{\TMblank}{\TMright}} (B)
            edge [bend left] node {\TMtrans{\TMstroke}{\TMstroke}{\TMright}} (A);
\end{tikzpicture}
\]

ایہ حالتی خاکہ ہِٹھلے
حالتی تبدیلی دے فنکشن
دے مطابق اے:
\begin{align*}
\delta(q_0, \TMstroke) & = \tuple{q_1, \TMstroke, \TMright},\\
\delta(q_1, \TMstroke) & = \tuple{q_0, \TMstroke, \TMright},\\
\delta(q_1, \TMblank)  & = \tuple{q_1, \TMblank, \TMright}
\end{align*}
\end{ex}

\begin{explain}
اُتلی مشین صرف اُدوں
رُکدی اے جدوں انپٹ وچ
لکیراں دی گنتی جفت ہووے۔
ہور صورت وچ مشین نظری طور اُتے
ہمیشہ چلدی رہندی اے۔ کسے وی
مشین تے انپٹ لئی مشین
دیاں \emph{مکمل حالتاں}
(configurations) دا سلسلہ
ویکھ کے آؤٹ پٹ دا پتا
لایا جا سکدا اے۔ مکمل
حالت دی رسمی تعریف
اسی اگے دیواں گے۔
فی الحال ایہناں نوں چلن دے
وکھو وکھرے ویلیاں اُتے مشین دی حالت
دکھان والے خاکیاں دا
سلسلہ سمجھ لو۔ ہر مکمل
حالت پٹی دا متن، مشین
دی حالت تے پڑھن-لکھن
والے سر دی جگہ دکھاندی
اے۔

آؤ جفت مشین نوں چار
$\TMstroke$ دا انپٹ دے
کے اوس دیاں مکمل حالتاں
دا سلسلہ ویکھئیے۔ ایس
صورت وچ مشین دے رُکّن
دی توقع اے۔ فیر اسی
مشین نوں تِن $\TMstroke$
دے انپٹ اُتے چلاواں گے۔
اوس صورت وچ اوہ ہمیشہ چلدی
رہے گی۔

% PNB-CMP-026: source says “state one” although the initial index is zero.
مشین حالت~$q_0$ وچ
شروع ہوندی اے تے
سب توں کھبے پاسے والا~$\TMstroke$
پڑھدی اے۔ مشین دی
ابتدائی مکمل حالت
ایداں دکھا سکدے آں:
\[
\TMendtape \TMstroke_0 \TMstroke \TMstroke \TMstroke \TMblank \ldots
\]
اُتلی مکمل حالت
سمجھنی سوکھی اے۔
جیویں ویکھ سکدے آں،
مشین صفر نمبر والی
ابتدائی حالت وچ
سب توں کھبے پاسے والا~$\TMstroke$
پڑھدی اے۔ پہلی~$\TMstroke$
دے تھلّے حالتی ناں
دا ذیلی ہندسہ ایس
گل نوں دکھاندا اے۔
ایتھے لاگو ہون
والا حکم
$\delta(q_0, \TMstroke) = \tuple{q_1, \TMstroke, \TMright}$
اے؛ ایس لئی مشین دا سر
پٹی اُتے سجّے ول سرکدا اے تے مشین
حالت~$q_1$ وچ
چلی جاندی اے۔
\[
\TMendtape \TMstroke \TMstroke_1 \TMstroke \TMstroke \TMblank \ldots
\]
ہن مشین حالت~$q_1$
وچ اک~$\TMstroke$
پڑھ رہی اے؛ ایس
لئی اسی حکم
$\delta(q_1, \TMstroke) = \tuple{q_0,
  \TMstroke, \TMright}$
دی پیروی کرنی اے۔
ایس توں ہِٹھلی
مکمل حالت ملدی اے:
\[
\TMendtape \TMstroke \TMstroke \TMstroke_0 \TMstroke \TMblank \ldots
\]
مشین اگے چلدی اے
تاں اوہو قاعدے
اوہو ترتیب وچ
فیر لگدے نیں؛
اگلیاں دو مکمل
حالتاں ایہ نیں:
\[
\TMendtape \TMstroke \TMstroke \TMstroke \TMstroke_1 \TMblank \ldots
\]
\[
\TMendtape \TMstroke \TMstroke \TMstroke \TMstroke \TMblank_0 \ldots
\]
ہن مشین حالت~$q_0$
وچ اک~$\TMblank$
پڑھ رہی اے۔
حالتی تبدیلی دے
خاکے توں صاف اے
کہ ایتھے کوئی
حکم موجود نہیں؛
ایس لئی مشین
رُک گئی اے۔
ایہدا مطلب اے
کہ انپٹ قبول
ہو گیا اے۔

ہن فرض کرو اسی
مشین نوں تِن
$\TMstroke$ دا انپٹ
دے کے شروع کریے۔
پہلیاں کجھ مکمل
حالتاں پچھلیاں
ورگیاں نیں: اوہو
حکم چلدے نیں،
صرف پٹی دا انپٹ
تھوڑا وکھرا اے:
\[
\TMendtape \TMstroke_0 \TMstroke \TMstroke \TMblank \ldots
\]
\[
\TMendtape \TMstroke \TMstroke_1 \TMstroke \TMblank \ldots
\]
\[
\TMendtape \TMstroke \TMstroke \TMstroke_0 \TMblank \ldots
\]
\[
\TMendtape \TMstroke \TMstroke \TMstroke \TMblank_1 \ldots
\]
ہن مشین سارے
$\TMstroke$ لنگھ
چکی اے تے
اک~$\TMblank$
حالت~$q_1$ وچ
پڑھ رہی اے۔
خاکے مطابق
ایس لئی حکم
$\delta(q_1, \TMblank) =\tuple{q_1,
\TMblank, \TMright}$
موجود اے۔ چونکہ
پٹی سجّے پاسے
ول بے انتہا
$\TMblank$ نال
بھری اے، مشین
ایہو حکم
\emph{ہمیشہ}
چلاندی رہے گی،
حالت~$q_1$ وچ رہندی اے،
تے اوہدا سر ہور ہور سجّے ول سرکدا رہندا اے۔
اوہ کدی
نہیں رُکے گی
تے انپٹ قبول
نہیں کرے گی۔
\end{explain}

\begin{explain}
ایہ یاد رکھنا ضروری
اے کہ کجھ مشیناں رُکدیاں ای نہیں۔ جے
رُکّن دا مطلب ایہ اے
کہ مشین کول چلان لئی
حکم مک جان، تاں ہر
حالت وچ ہر نشان لئی
باہر جان والا تیر
رکھ کے اسی ایداں
دی مشین بنا سکدے آں
جیڑی کدی نہ رُکے۔
جفت مشین نوں وی
$\TMblank$ حالت~$q_0$
وچ پڑھݨ
لئی اک حکم پا کے
ہمیشہ چلن والی
بنا سکدے آں۔
\end{explain}

\begin{ex}
\[
\begin{tikzpicture}[->,>=stealth',shorten >=1pt,auto,node distance=2.8cm,
                    semithick]
  \tikzstyle{every state}=[fill=none,draw=black,text=black]

  \node[initial,state] (A)              {$q_0$};
  \node[state]         (B) [right of=A] {$q_1$};

  \path
  (A) edge [bend left] node {\TMtrans{\TMstroke}{\TMstroke}{\TMright}} (B)
      edge [loop above] node {\TMtrans{\TMblank}{\TMblank}{\TMright}} (A)
  (B) edge [loop above] node {\TMtrans{\TMblank}{\TMblank}{\TMright}} (B)
      edge [bend left] node {\TMtrans{\TMstroke}{\TMstroke}{\TMright}} (A);
\end{tikzpicture}
\]
\end{ex}

\begin{explain}
ٹیورنگ مشیناں دی نمائندگی
مشینی جدولاں \pnbLatin{(machine tables)} نال وی ہو سکدی اے۔
ایہناں وچ پٹی دے
حرفاں دا مجموعہ
$x$-محور اُتے
تے مشین دیاں
حالتاں دا سیٹ
$y$-محور اُتے
دکھایا جاندا اے۔
جدول وچ ہر حالت دی قطار تے ہر نشان دے کالم
دے ملاپ والے خانے وچ حکم دا باقی حصہ،
یعنی نویں حالت،
نواں نشان تے
حرکت دی سمت،
لکھی ہوندی اے۔
مشینی جدول توں
ایہ لبھنا سوکھا
اے کہ مشین کس
حالت تے کس نشان
اُتے رُکدی اے۔
جدول وچ جتھے
خانہ خالی ہووے،
اوتھے مشین رُک
سکدی اے۔ حالتی
خاکیاں تے حکم دی
فہرست وچ رُکّن
والیاں جگہاں فوراً
ظاہر نہیں ہوندیاں؛
جدول دے خالی
خانے ویکھ کے
اوہ جلدی لبھ
جاندیاں نیں۔
\end{explain}

\begin{ex}
جفت مشین دا جدول ایہ اے:
\[
\centering
\begin{tabular}{lllll}
\cline{2-4}
\multicolumn{1}{l|}{}      & \multicolumn{1}{c|}{$\TMblank$}                
& \multicolumn{1}{c|}{$\TMstroke$}     & \multicolumn{1}{c|}{$\TMendtape$}   &  \\ \cline{1-4}
\multicolumn{1}{|l|}{$q_0$} & \multicolumn{1}{l|}{}                          
& \multicolumn{1}{l|}{\TMtrans{\TMstroke}{q_1}{\TMright}} & 
\multicolumn{1}{l|}{\phantom{\TMtrans{\TMstroke}{q_1}{\TMright}}} &  \\ \cline{1-4}
\multicolumn{1}{|l|}{$q_1$} & \multicolumn{1}{l|}{\TMtrans{\TMblank}{q_1}{\TMright}} 
& \multicolumn{1}{l|}{\TMtrans{\TMstroke}{q_0}{\TMright}} & 
\multicolumn{1}{l|}{\phantom{\TMtrans{\TMstroke}{q_1}{\TMright}}} &  \\ \cline{1-4}
\end{tabular}
\]
جیویں ویکھ سکدے آں،
مشین اک~$\TMblank$
حالت~$q_0$ وچ پڑھدے
ویلے رُکدی اے۔
\end{ex}

\begin{explain}
ہُن تک اسی صرف
اوہ مشیناں ویکھیاں
نیں جیہڑیاں انپٹ
پڑھدیاں تے قبول
کردیاں نیں۔ پر
ٹیورنگ مشیناں
پڑھ وی سکدیاں
تے لکھ وی سکدیاں
نیں۔ ایہدی اک مثال
(ہور بہت مثالاں وی
نیں) \emph{دوگنا کرن
والی مشین} (doubler)
اے۔ جے پٹی اُتے~$n$
$\TMstroke$ دا
اک ٹکڑا دے کے
ایہ مشین شروع
کریے، تاں ایہ
$2n$ $\TMstroke$
دا ٹکڑا آؤٹ پٹ
دیندی اے۔
\end{explain}

\begin{ex}
  \ollabel{ex:doubler} دوگنا
  کرن والی مشین بناؤن
  توں پہلاں مسئلہ
  حل کرن دی اک
  \emph{تدبیر} بنانی
  ضروری اے۔ جیویں
  اسی مشین نوں
  ہُن تک بیان کیتا
  اے، ایہ یاد نہیں
  رکھ سکدی کہ
  اوس نے کنے
  $\TMstroke$ پڑھے
  نیں؛ ایس لئی
  پٹی اُتے موجود
  سارے $\TMstroke$
  دا حساب رکھن
  دا کوئی طریقہ
  لبھنا پئے گا۔
  اک طریقہ ایہ
  اے کہ آؤٹ پٹ
  نوں انپٹ توں
  اک~$\TMblank$ نال
  وکھ کر دئیے۔
  مشین فیر انپٹ
  وچوں پہلا
  $\TMstroke$ مٹا
  دیوے، باقی انپٹ
  لنگھے، اک
  $\TMblank$ چھڈے
  تے دو نویں
  $\TMstroke$ لکھے۔
  فیر مشین واپس
  جا کے انپٹ وچ
  دوجا $\TMstroke$
  لبھے تے اوس
  نوں وی دوگنا
  کرے۔ انپٹ دے
  ہر اک $\TMstroke$
  لئی آؤٹ پٹ دے
  دو $\TMstroke$
  لکھے گی۔ چلن
  نال نال انپٹ
  مٹاندی جاوے تاں
  ایہ یقین رہندا
  اے کہ کوئی
  $\TMstroke$ نہ
  رہ جاوے تے نہ اوہنوں دو واری دوگنا
  کیتا جاوے۔ جدوں
  سارا انپٹ مٹ
  جاوے گا، پٹی
  اُتے $2n$
  $\TMstroke$ رہ
  جان گے۔ نتیجے
  والی ٹیورنگ
  مشین دا حالتی
  خاکہ \olref{fig:doubler}
  وچ دکھایا گیا اے۔
  \begin{figure}
\[
\begin{tikzpicture}[->,>=stealth',shorten >=1pt,auto,node distance=2.8cm,
                    semithick]
  \tikzstyle{every state}=[fill=none,draw=black,text=black]

  \node[initial,state] (1)              {$q_0$};
  \node[state]         (2) [right of=1] {$q_1$};
  \node[state]         (3) [right of=2] {$q_2$};
  \node[state]         (4) [below of=3] {$q_3$};
  \node[state]         (5) [left of=4]  {$q_4$};
  \node[state]         (6) [left of=5]  {$q_5$};

  \path (1) edge node {\TMtrans{\TMstroke}{\TMblank}{\TMright}} (2)
    (2) edge [loop above] node {\TMtrans{\TMstroke}{\TMstroke}{\TMright}} (2)
      edge node {\TMtrans{\TMblank}{\TMblank}{\TMright}} (3)
    (3) edge [loop above] node {\TMtrans{\TMstroke}{\TMstroke}{\TMright}} (3)
        edge node {\TMtrans{\TMblank}{\TMstroke}{\TMright}} (4)
    (4) edge [loop below] node {\TMtrans{\TMblank}{\TMstroke}{\TMleft}} (4)
        edge node {\TMtrans{\TMstroke}{\TMstroke}{\TMleft}} (5)
    (5) edge [loop below]  node {\TMtrans{\TMstroke}{\TMstroke}{\TMleft}} (5)
        edge              node {\TMtrans{\TMblank}{\TMblank}{\TMleft}} (6)
    (6) edge [loop below] node {\TMtrans{\TMstroke}{\TMstroke}{\TMleft}} (6)
        edge              node {\TMtrans{\TMblank}{\TMblank}{\TMright}} (1);
\end{tikzpicture}
\]
\caption{اک دوگنا کرن والی مشین}
\ollabel{fig:doubler}
\end{figure}
\end{ex}

\begin{prob}
کوئی وی انپٹ
چنو تے
\olref[tur][mac][rep]{ex:doubler}
وچ دتی دوگنا
کرن والی مشین
دیاں مکمل حالتاں
دا سلسلہ قدم بہ
قدم ویکھو۔
\end{prob}

\begin{prob}
حرفاں دے مجموعے
$\{\TMendtape,\TMblank, A, B\}$
والی ٹیورنگ مشین
بناؤ جیڑی $A$
تے $B$ دیاں اوہ
لڑیاں قبول کرے،
یعنی اوہناں اُتے
رُک جائے، جِنہاں
وچ $A$ دی گنتی
$B$ دی گنتی دے
برابر ہووے
\emph{تے} سارے
$A$، سارے $B$
توں پہلاں آندے
ہون۔ جِنہاں لڑیاں
وچ $A$ دی گنتی
$B$ دی گنتی
دے برابر نہ ہووے
یا سارے $A$،
سارے~$B$ توں
پہلاں نہ آندے
ہون، اوہناں نوں
رد کرے، یعنی
اوہناں اُتے نہ
رُکے۔ مثال لئی،
مشین $AABB$
تے $AAABBB$
قبول کرے، پر $AAB$
تے $AABBAABB$
دونوں نوں رد کرے۔
\end{prob}

\begin{prob}
حرفاں دے مجموعے
$\{\TMendtape,\TMblank, A, B\}$
والی ٹیورنگ مشین
بناؤ جیڑی کوئی
وی لڑی $\alpha$
انپٹ لے؛ ایہ
لڑی $A$ تے $B$
دی بنی ہووے۔
مشین اوس دی اک ہور نقل
اوہدے نال جوڑ کے
$\alpha\alpha$
دی صورت وچ
آؤٹ پٹ بناوے۔
مثال لئی، انپٹ
$ABBA$ توں
آؤٹ پٹ
$ABBAABBA$
بننا چاہیدا اے۔
\end{prob}

\begin{prob}
\emph{حرف وار ترتیب؟:}
حرفاں دے مجموعے
$\{\TMendtape,\TMblank, A, B\}$
والی ٹیورنگ مشین
بناؤ جیڑی
$A$ تے $B$
دا اک محدود
سلسلہ انپٹ
لے کے جانچے
کہ سارے $A$
سارے $B$
دے کھبے پاسے
نیں یا نہیں۔
مشین انپٹ
دی لڑی پٹی
اُتے رہن دیوے؛
جے لڑی «حرف
وار ترتیب وچ»
ہووے تاں رُکے،
نہ ہووے تاں
ہمیشہ چلدی
رہے۔
\end{prob}

\begin{prob}
\emph{حرف وار ترتیب
دین والی مشین:}
حرفاں دے مجموعے
$\{\TMendtape,\TMblank, A, B\}$
والی ٹیورنگ مشین
بناؤ جیڑی
$A$ تے $B$
دا اک محدود
سلسلہ انپٹ
لے کے اوس
دے حرف ایداں
فیر ترتیب دے
کہ سارے $A$
سارے~$B$ دے
کھبے پاسے
ہون۔ مثال لئی،
لڑی $BABAA$
لڑی $AAABB$
بن جائے، تے
لڑی $ABBABB$
لڑی $AABBBB$
بن جائے۔
\end{prob}
\endgroup
% /OLP-0255

% OLP-0256 content/turing-machines/machines-computations/turing-machines.tex
\begingroup
\def\olfilename{turing-machines}
\expandafter\def\csname ol@sectioncs\endcsname{section}

\olfileid{tur}{mac}{tur}
\olsection{ٹیورنگ مشیناں}

\begin{explain}
ٹیورنگ مشین دی
رسمی تعریف ویکھن
نوں تجریدی لگدی اے،
پر اصل وچ سادی اے:
ایہ ٹیورنگ مشین
دا عمل مقرر کرن
لئی درکار ساری
معلومات اک ریاضیاتی
بناوٹ وچ اکٹھیاں
کردی اے۔ ایہناں
وچ (1) مشین دیاں
ممکنہ حالتاں،
(2) پٹی اُتے
آ سکّن والے
نشان، (3) مشین
دی ابتدائی حالت
تے (4) مشین دے
حکماں دا مجموعہ
شامل نیں۔
\end{explain}

\begin{defn}[ٹیورنگ مشین]
اک \emph{ٹیورنگ مشین}
$M$ چار رکناں
دا اک ترتیب وار ٹپل
$\langle Q, \Sigma, q_0,
\delta\rangle$ اے،
جیہدے وچ:
\begin{enumerate}
\item \emph{حالتاں} دا اک
  محدود سیٹ~$Q$،
\item اک محدود
  \emph{حرفاں دا مجموعہ}
  $\Sigma$، جیدے وچ
  $\TMendtape$ تے
  $\TMblank$ شامل نیں،
\item اک \emph{ابتدائی
  حالت}~$q_0 \in Q$،
\item اک محدود
  \emph{حکماں دا مجموعہ}~$\delta\colon Q \times \Sigma
  \pto Q \times \Sigma \times \{\TMleft, \TMright, \TMstay\}$۔
\end{enumerate}
جزوی فنکشن~$\delta$
نوں $M$ دا
\emph{حالتی تبدیلی
دا فنکشن} وی
آکھدے نیں۔
\end{defn}

\begin{explain}
اسی مندے آں کہ
پٹی صرف اک پاسے
ول لامتناہی اے۔
ایس لئی اک خاص
نشان~$\TMendtape$
نوں پٹی دے کھبے
سِرے دی نشانی
بناؤنا فائدہ مند اے۔
ایس نال ٹیورنگ
مشین دے پروگرام
نوں پتا لاؤن وچ
سہولت رہندی اے
کہ اوہ کدوں پٹی
توں باہر نکلن
دے خطرے وچ اے۔
اسی ایہ شرط لا
سکدے ساں کہ ایس
نشان اُتے کدی
نواں نشان نہ لکھیا
جاوے، یعنی
$\delta(q,\TMendtape) = \tuple{q', \TMendtape, x}$
جے $\delta(q,\TMendtape)$
تعریف شدہ ہووے۔
کجھ درسی کتاباں
ایہ شرط رکھدیاں
نیں؛ اسی نہیں
رکھدے۔ ٹیورنگ
مشین بناؤندے
ویلے تسی صرف
ایہ خیال رکھ
سکدے او کہ
اوہ~$\TMendtape$
نوں کدی نہ بدلے۔
ہور ایہ وی
اے کہ کجھ
صورتاں وچ ایس
نشان نوں بدلݨ
دی اجازت دینا
مفید لچک دیندا
اے۔
\end{explain}

\begin{ex}
\emph{جفت مشین:}
جفت مشین رسمی
طور اُتے چار رکناں
دا ترتیب وار ٹپل
$\tuple{Q, \Sigma, q_0, \delta}$
اے، جتھے:
\begin{align*}
Q & = \{ q_0, q_1 \} \\
\Sigma & = \{ \TMendtape, \TMblank, \TMstroke \}, \\
\delta(q_0, \TMstroke) & = \tuple{q_1, \TMstroke, \TMright},\\
\delta(q_1, \TMstroke) & = \tuple{q_0, \TMstroke, \TMright},\\
\delta(q_1, \TMblank)  & = \tuple{q_1, \TMblank, \TMright}.
\end{align*}
\end{ex}
\endgroup
% /OLP-0256

% OLP-0257 content/turing-machines/machines-computations/configuration.tex
\begingroup
\def\olfilename{configuration}
\expandafter\def\csname ol@sectioncs\endcsname{section}

\olfileid{cmp}{tur}{con}
\olsection{مکمل حالتاں تے حسابی کارروائیاں}

\begin{explain}
پہلاں جفت مشین دیاں مکمل
حالتاں وچوں قدم بہ قدم
لنگھنا یاد کرو۔ پٹی،
پڑھن-لکھن والے سر تے
ٹیورنگ مشین دے پروگرام
توں بنیا تصوراتی نظام
دراصل ٹیورنگ مشین دی
حسابی کارروائی نوں سمجھݨ
دا اک بدیہی طریقہ اے۔
رسمی طور اُتے، دتے ہوئے
انپٹ لئی مشین دی حسابی
کارروائی نوں
\emph{مکمل حالتاں} دے سلسلے
وجوں تعریف کر سکدے آں۔
ہر مکمل حالت نشاناں
دے اک سلسلے (جیڑا
حسابی کارروائی دے اوس
ویلے پٹی دا متن دکھاندا
اے)، پڑھن-لکھن والے
سر دی جگہ دسن والا
اک عدد، تے مشین دی
اک حالت توں بندی اے۔ ایہناں دی
مدد نال اسی تعریف کر
سکدے آں کہ ٹیورنگ
مشین~$M$ دتے ہوئے
انپٹ اُتے کیہ حساب
کردی اے۔
\end{explain}

\begin{defn}[مکمل حالت]
ٹیورنگ مشین $M = \tuple{Q, \Sigma, q_0,
\delta}$ دی اک \emph{مکمل حالت}
تِن رکناں دا ترتیب وار
ٹپل $\tuple{C, m, q}$ اے،
جتھے:
\begin{enumerate}
\item $C \in \Sigma^*$، $\Sigma$
  دے نشاناں دا اک محدود
  سلسلہ اے،
\item $m \in \Nat$ اک عدد اے
  تے $< \len{C}$ اے، تے
\item $q \in Q$۔
\end{enumerate}
بدیہی طور اُتے سلسلہ~$C$
پٹی دا متن اے۔ ایہ سب توں کھبے خانے توں
آخری متعلقہ خانے تک دے سارے نشان نیں۔ آخری
متعلقہ خانہ اوہناں خانیاں وچوں سب توں سجّے پاسے اے
جیہڑے خالی نہیں یا جیہناں تک سر پہلے پہنچ چکیا اے۔ $m$~اوس
خانے دا نمبر اے جیہنوں
پڑھن-لکھن والا سر
ایس ویلے پڑھ رہیا اے؛
سب توں کھبے خانے دا
نمبر $0$ اے۔ $q$ مشین دی
موجودہ حالت اے۔
\end{defn}

\begin{explain}
ٹیورنگ مشین دا ممکنہ
انپٹ نشاناں دا اک سلسلہ
اے، عموماً اوہ سلسلہ
جیڑا کسے نہ کسے روپ
وچ اک عدد دا کوڈ اے۔ مشین دی
ابتدائی مکمل حالت
اوہ حالت اے جتھوں اسی
اُس انپٹ اُتے مشین
دا عمل شروع کردے آں:
پٹی اُتے پہلے پٹی دے
سِرے دا نشان تے اوہدے
فوراً مگروں سجے پاسے
دے خانیاں وچ انپٹ
لکھیا ہوندا اے؛
پڑھن-لکھن والا سر
انپٹ دے سب توں کھبے
خانے نوں پڑھ رہیا
ہوندا اے (یعنی کھبے
سِرے دے نشان توں
اگلے سجے خانے نوں)،
تے مشین مقرر کیتی
ابتدائی حالت~$q_0$
وچ ہوندی اے۔
\end{explain}

\begin{defn}[ابتدائی مکمل حالت]
$M$ دی انپٹ $I \in \Sigma^*$
لئی \emph{ابتدائی مکمل
حالت} ایہ اے:
\[
\tuple{\TMendtape \frown I, 1, q_0}.
\]
\end{defn}

\begin{explain}
نشان~$\frown$ دا مطلب
\emph{جوڑنا} (concatenation)
اے۔
% PNB-CMP-027: source B9 says left; B7 and the initial tuple put input right of the endmarker.
انپٹ دی لڑی پٹی دے
کھبے سِرے دے نشان
توں فوراً سجے پاسے
شروع ہوندی اے۔
\end{explain}

\begin{defn}
اسی آکھدے آں کہ مکمل
حالت $\tuple{C, m, q}$
\emph{اک قدم وچ مکمل حالت
$\tuple{C', m', q'}$ پیدا
کردی اے} (مشین~$M$ دے
مطابق)، اُدوں تے صرف
اُدوں جدوں:
\begin{enumerate}
\item $m$-ویں تھاں اُتے $C$ دا نشان
  $\sigma$ ہووے،
\item $M$ دے حکماں دا
  مجموعہ $\delta(q, \sigma) =
  \tuple{q', \sigma', D}$
  مقرر کردا ہووے،
\item $m$-ویں تھاں اُتے $C'$ دا
  نشان $\sigma'$ ہووے، تے
\item
\begin{enumerate}
\item جے $D = L$ ہووے تاں
  $m' = m - 1$، بشرطیکہ
  $m>0$ ہووے؛ نہیں
  تاں $m'=0$، یا
\item $D = R$ تے $m' = m + 1$، یا
\item $D = N$ تے $m' = m$،
\end{enumerate}
\item جے $m' = \len{C}$
  ہووے تاں $\len{C'} = \len{C} + 1$
  تے $m'$-ویں تھاں اُتے $C'$ دا
  نشان~$\TMblank$ ہووے۔
  نہیں تاں $\len{C'}=\len{C}$۔
\item ہر اوس $i$ لئی
  جدے لئی $i < \len{C}$
  تے $i \neq m$ ہووے،
  $C'(i) = C(i)$ ہووے،
\end{enumerate}
\end{defn}

\begin{defn}\ollabel{defn:run-output}
$M$ دا \emph{انپٹ~$I$
اُتے چلن} مکمل حالتاں
دا سلسلہ $C_i$ اے،
تے ایہ ساریاں مکمل حالتاں $M$ دیاں نیں۔
ایتھے $C_0$، $M$ دی
انپٹ~$I$ لئی ابتدائی مکمل
حالت اے، تے ہر $C_i$
اک قدم وچ $C_{i+1}$
پیدا کردی اے۔

اسی آکھدے آں کہ $M$
انپٹ $I$ اُتے $k$ قدماں
مگروں \emph{رُکدی اے}
جے $C_k =
\tuple{C, m, q}$ ہووے،
~$m$-ویں تھاں اُتے $C$ دا
نشان~$\sigma$ ہووے، تے
$\delta(q,
\sigma)$ غیر تعریف شدہ
ہووے۔ ایس صورت وچ
مشین~$M$ دے انپٹ~$I$ دا
\emph{آؤٹ پٹ}~$O$ اے؛
$O$ نشاناں دی ایداں دی
لڑی اے جیڑی~$\TMblank$
اُتے ختم نہیں ہوندی،
تے $C = \TMendtape \concat O \concat \TMblank^j$
کسے~$j \in \Nat$
لئی ہووے۔ ($\TMblank^j$،
$j$ عدد~$\TMblank$ نشاناں
دا سلسلہ اے۔)
\end{defn}

\begin{explain}
ایس تعریف دے مطابق
آؤٹ پٹ~$O$، جیہڑا مشین~$M$ دا اے،
ہمیشہ~$\TMblank$ توں
وکھرے نشان اُتے ختم
ہوندا اے؛ یا جے ویلے~$k$
پوری پٹی~$\TMblank$
نال بھری ہووے (سب توں
کھبے~$\TMendtape$
توں سوا)، تاں $O$
خالی لڑی اے۔
\end{explain}
\endgroup
% /OLP-0257

% OLP-0258 content/turing-machines/machines-computations/unary-numbers.tex
\begingroup
\def\olfilename{unary-numbers}
\expandafter\def\csname ol@sectioncs\endcsname{section}

\olfileid{tur}{mac}{una}
\olsection{اعداد دی اک نشانی نمائندگی}

\begin{explain}
ٹیورنگ مشیناں اپنی پٹی
اُتے لکھے نشاناں دے
سلسلّیاں اُتے کم کردیاں
نیں۔ مشین دے ورتے
حرفاں دے مجموعے دے
مطابق ایہ سلسلے کئی
طرح دے انپٹ تے آؤٹ پٹ
دکھا سکدے نیں۔ خاص طور
اُتے، \emph{حسابی}
فنکشن، یعنی قدرتی
اعداد دے فنکشن، دا
حساب کرن والیاں ٹیورنگ
مشیناں وچ ساڈی دلچسپی
اے۔ مثبت صحیح اعداد
نوں دکھان دا اک سادہ
طریقہ ایہ اے کہ ہر عدد
نوں اکّو نشان~$\TMstroke$
دی لڑی نال کوڈ کریے۔
جے $n \in \Nat$ ہووے
تاں $\TMstroke^n$، $n = 0$
ہون اُتے خالی سلسلہ
اے، تے باقی صورت وچ
ٹھیک $n$ عدد~$\TMstroke$
نشاناں دا سلسلہ اے۔
\end{explain}

\begin{defn}[حسابی کارروائی]
ٹیورنگ مشین~$M$ فنکشن
$f\colon \Nat^k \to \Nat$
دا \emph{حساب کردی اے}
اُدوں تے صرف اُدوں جدوں
$M$ ہِٹھلے انپٹ اُتے
رُکے:
\[
\TMstroke^{n_1} \TMblank \TMstroke^{n_2} \TMblank \dots \TMblank \TMstroke^{n_k}
\]
تے آؤٹ پٹ
$\TMstroke^{f(n_1, \dots, n_k)}$
ہووے۔
\end{defn}

\begin{prob}
ایہدی تعریف دیو کہ
ٹیورنگ مشین~$M$ فنکشن
$f\colon \Nat^k \to \Nat^m$
دا حساب کدوں کردی اے۔
\end{prob}

\begin{ex}\ollabel{ex:adder}
\emph{جمع:}
آؤ فنکشن $f(n,m) = n + m$
دا حساب کرن والی مشین
بنائیے۔ مشین دی پٹی
اُتے شروع وچ $\TMstroke$
نشاناں دے دو ٹکڑے
ہون، جنہاں دیاں لمبائیاں
$n$ تے~$m$ ہون، تے
رُکّن ویلے $n+m$~$\TMstroke$
نشاناں دا اک ٹکڑا رہ
جاوے۔ انپٹ دے ایہہ دو
$\TMstroke$ ٹکڑے اک~$\TMblank$
نال وکھ کیتے ہوندے
نیں۔ ایس لئی اک طریقہ
ایہ اے کہ $\TMblank$
والے خانے اُتے اک لکیر
لکھی جاوے تے آخری~$\TMstroke$
مٹا دتا جاوے۔
\begin{figure}\[
\begin{tikzpicture}[->,>=stealth',shorten >=1pt,auto,node distance=2.8cm,
                    semithick]
  \tikzstyle{every state}=[fill=none,draw=black,text=black]

  \node[initial,state]         (A)              {$q_0$};
  \node[state]         (B) [right of=A] {$q_1$};
  \node[state]         (C) [right of=B] {$q_2$};

  \path (A) edge node {\TMtrans{\TMblank}{\TMstroke}{\TMstay}} (B)
           edge [loop above] node {\TMtrans{\TMstroke}{\TMstroke}{\TMright}} (B)
        (B) edge [loop above] node {\TMtrans{\TMstroke}{\TMstroke}{\TMright}} (B)
            edge node {\TMtrans{\TMblank}{\TMblank}{\TMleft}} (C)
        (C) edge [loop above] node {\TMtrans{\TMstroke}{\TMblank}{\TMstay}} (C);
\end{tikzpicture}
\]\caption{$f(x,y) = x+y$ دا حساب کرن والی مشین}
\ollabel{fig:adder}
\end{figure}
\end{ex}

\begin{prob}
\olref[tur][mac][una]{ex:adder}
وچلی مشین دیاں مکمل
حالتاں دا سلسلہ
انپٹ~$\tuple{3,2}$ لئی
قدم بہ قدم لکھو۔ جے
مشین $0+0$ دا حساب
کرے تاں کیہ ہوندا اے؟
\end{prob}

\begin{explain}
\olref[rep]{ex:doubler}
وچ اسی اک ٹیورنگ مشین
دی مثال دتی سی جیڑی
انپٹ وجوں~$\TMstroke$
نشاناں دی لڑی لیندی
اے تے پٹی اُتے دوگنے
~$\TMstroke$ نشاناں دی
لڑی نال رُکدی اے:
دوگنا کرن والی مشین۔
پر آؤٹ پٹ وچ
$\TMblank$ نشان دوگنے
~$\TMstroke$ ٹکڑے دے
کھبے پاسے رہ جاندے
نیں۔ ایس لئی ایہ مشین
فنکشن $f(x) = 2x$ دا حساب
نہیں کردی، بھاویں تسی شاید
ایہ سمجھیا ہووے کہ کردی اے۔ اسی ایس
مسئلے نوں ٹھیک کرن دے
دو طریقے بیان کراں گے۔
\end{explain}

\begin{ex}
\olref{fig:doubler-disc}
وچلی مشین فنکشن $f(x) =
2x$ دا حساب کردی اے۔
\olref[rep]{ex:doubler}
وچلی مشین ہر انپٹ
$\TMstroke$ مٹا کے
دور سجے پاسے دو
$\TMstroke$ لکھدی اے؛
ایس دے برعکس ایہ
مشین انپٹ دی ہر
~$\TMstroke$ لئی
سجّے پاسے صرف اک
~$\TMstroke$ ودھاؤندی
اے۔ ایہنوں انپٹ
دے ختم ہون دی جگہ
یاد رکھنی پیندی اے،
ایس لئی انپٹ تے
نواں لکھیاں لکیراں
دے وچ اک~$\TMblank$
چھڈدی اے، جیہنوں
آخر وچ~$\TMstroke$
نال بھر دیندی اے۔
نالے انپٹ وچ اپنی
موجودہ جگہ یاد
رکھن لئی اسی انپٹ
دے ٹکڑے دا اک
~$\TMstroke$ عارضی
طور اُتے~$\TMblank$
نال بدل دیندے آں۔
  \begin{figure}
\[
\begin{tikzpicture}[->,>=stealth',shorten >=1pt,auto,node distance=2.8cm,
                    semithick]
  \tikzstyle{every state}=[fill=none,draw=black,text=black]

  \node[initial,state] (0)              {$q_0$};
  \node[state]         (1) [above of=0] {$q_1$};
  \node[state]         (2) [above of=1] {$q_2$};
  \node[state]         (3) [right of=2] {$q_3$};
  \node[state]         (4) [below of=3] {$q_4$};
  \node[state]         (5) [below of=4] {$q_5$};
  \node[state]         (6) [above right of=4] {$q_6$};
  \node[state]         (7) [below of=6] {$q_7$};
  \node[state]         (8) [below of=7] {$q_8$};

  \path (0) edge node {\TMtrans{\TMstroke}{\TMblank}{\TMright}} (1)
%    (0) edge node {\TMtrans{\TMblank}{\TMblank}{\TMstay}} (h)
    (1) edge [loop left] node {\TMtrans{\TMstroke}{\TMstroke}{\TMright}} (1)
      edge node {\TMtrans{\TMblank}{\TMblank}{\TMright}} (2)
    (2) edge [loop above] node {\TMtrans{\TMstroke}{\TMstroke}{\TMright}} (2)
        edge node {\TMtrans{\TMblank}{\TMstroke}{\TMleft}} (3)
    (3) edge [loop above] node {\TMtrans{\TMstroke}{\TMstroke}{\TMleft}} (3)
        edge node[left] {\TMtrans{\TMblank}{\TMblank}{\TMleft}} (4)
    (4) edge        node[left] {\TMtrans{\TMstroke}{\TMstroke}{\TMleft}} (5)
    (5) edge [loop below] node {\TMtrans{\TMstroke}{\TMstroke}{\TMleft}} (5)
        edge              node {\TMtrans{\TMblank}{\TMstroke}{\TMright}} (0)
    (4) edge      node[sloped] {\TMtrans{\TMblank}{\TMstroke}{\TMright}} (6)
    (6) edge              node {\TMtrans{\TMblank}{\TMstroke}{\TMright}} (7)
    (7) edge [loop right] node {\TMtrans{\TMstroke}{\TMstroke}{\TMright}} (7)
        edge              node {\TMtrans{\TMblank}{\TMblank}{\TMleft}} (8)
    (8) edge [loop right] node {\TMtrans{\TMstroke}{\TMblank}{\TMstay}} (8);
    \end{tikzpicture}
\]
\caption{$f(x) = 2x$ دا حساب کرن والی مشین}
\ollabel{fig:doubler-disc}
\end{figure}
\end{ex}

\begin{ex}\ollabel{ex:mover}
$f(x) = 2x$ دا حساب
کرن دا دوجا طریقہ
ایہ اے کہ اصل
دوگنا کرن والی مشین
نوں برقرار رکھیا جاوے،
پر آخر وچ ہور حالتاں
تے حکم شامل کیتے
جان جیڑے دوگنی لکیراں
دا ٹکڑا پٹی دے بالکل
کھبے پاسے لے جان۔
\olref{fig:mover} وچلی
مشین صرف ایہ آخری
کم کردی اے۔ جے شروع
وچ پٹی اُتے پہلے
$\TMblank$ نشاناں دا
ٹکڑا، فیر~$\TMstroke$
نشاناں دا ٹکڑا ہووے
(تے سر~$\TMblank$
دے ٹکڑے وچ کسے وی
جگہ ہووے)، تاں ایہ
مشین~$\TMstroke$
نشاناں نوں اک اک کر
کے مٹا کے پٹی دے
شروع وچ لکھدی اے۔
کم پورا ہون دا پتا
لاؤن لئی ایہ پہلے
$\TMstroke$ ٹکڑے دے
آخر اُتے اک~$\TMendtape$
نشان لاؤندی اے،
جیہڑا آخر وچ مٹا
دتا جاندا اے۔ اسی
حالتاں نوں نمبر دینا~$q_6$
توں شروع کیتا اے
تاں جو ایہناں نوں
دوگنا کرن والی مشین
نال جوڑ سکئیے۔ صرف
اک ہور حکم
$\delta(q_0,
\TMblank) = \tuple{q_6,\TMblank,\TMstay}$
دی لوڑ اے؛ یعنی
~$q_0$ توں~$q_6$
تک اک تیر جس اُتے
~$\TMtrans{\TMblank}{\TMblank}{\TMstay}$
لکھیا ہووے۔ (اک
باریک مسئلہ اے:
نتیجے والی مشین
انپٹ~$x=0$ لئی
کم نہیں کردی۔ ایہنوں
اسی مشق لئی چھڈدے
آں۔)
\begin{figure}
  \[
  \begin{tikzpicture}[->,>=stealth',shorten >=1pt,auto,node distance=2.8cm,
                      semithick]
    \tikzstyle{every state}=[fill=none,draw=black,text=black]
    \node[initial,state] (6)              {$q_6$};
    \node[state]         (7) [right of=6] {$q_7$};
    \node[state]         (8) [right of=7] {$q_8$};
    \node[state]         (9) [below of=8] {$q_9$};
    \node[state]         (10) [left of=9] {$q_{10}$};
    \node[state]         (11) [left of=10]  {$q_{11}$};
    \node[state]         (12) [below of=11]  {$q_{12}$};
    \node[state]         (13) [right of=12] {$q_{13}$};
    \node[state]         (14) [right of=13] {$q_{14}$};
    \path
    (6)  edge node {\TMtrans{\TMblank}{\TMblank}{\TMright}} (7)
    (7)  edge [loop above] node {\TMtrans{\TMblank}{\TMblank}{\TMright}} (7)
         edge node {\TMtrans{\TMstroke}{\TMstroke}{\TMright}} (8)
    (8)  edge [loop above] node {\TMtrans{\TMstroke}{\TMstroke}{\TMright}} (8)
         edge node {\TMtrans{\TMblank}{\TMendtape}{\TMleft}} (9)
    (9)  edge [loop right] node {\TMtrans{\TMstroke}{\TMstroke}{\TMleft}} (9)
         edge node {\TMtrans{\TMblank}{\TMblank}{\TMright}} (10)
    (10) edge node[above] {\TMtrans{\TMstroke}{\TMblank}{\TMleft}} (11)
    (11) edge [loop above] node {\TMtrans{\TMblank}{\TMblank}{\TMleft}} (11)
         edge node[left] {\begin{tabular}{@{}l@{}}
          \TMtrans{\TMendtape}{\TMendtape}{\TMright}\\
          \TMtrans{\TMstroke}{\TMstroke}{\TMright}
         \end{tabular}} (12)
    (12) edge node[] {\TMtrans{\TMblank}{\TMstroke}{\TMright}} (13)
    (13) edge [loop below] node {\TMtrans{\TMblank}{\TMblank}{\TMright}} (13)
         edge node[sloped] {\TMtrans{\TMstroke}{\TMblank}{\TMleft}} (11)
    (13) edge node {\TMtrans{\TMendtape}{\TMblank}{\TMstay}} (14);
  \end{tikzpicture}
  \]
  \caption{$\TMstroke$ نشاناں دا ٹکڑا کھبے پاسے لے جانا}
  \ollabel{fig:mover}
\end{figure}
\end{ex}

\begin{prob}
\olref[tur][mac][una]{ex:mover}
وچ اسی اوہ مشین بیان کیتی اے
جیڑی
\olref[tur][mac][una]{fig:doubler-disc}
وچلی دوگنا کرن والی
مشین تے
\olref[tur][mac][una]{fig:mover}
وچلی ٹکڑا سرکاؤن
والی مشین نوں
ملا کے بندی اے۔
جے ایہ جوڑی ہوئی
مشین انپٹ~$x=0$
اُتے، یعنی خالی پٹی
اُتے، شروع کیتی
جاوے تاں کیہ
ہوندا اے؟ مشین
نوں کِویں بدلو گے
تاں جو ایس صورت
وچ اوہ آؤٹ پٹ~$2x=0$
نال رُکے؟ (ایہ
کم اک حالت تے
اک حالتی تبدیلی
ودھا کے ہو جانا
چاہیدا اے۔)
\end{prob}

\begin{prob}
\emph{تفریق:} اک
ٹیورنگ مشین بناؤ
جیڑی لکیراں دیاں دو
غیر خالی لڑیاں
انپٹ وجوں لے،
جنہاں دیاں لمبائیاں
$n$ تے~$m$ ہون،
جتھے $n >
m$، تے فنکشن
$f(n,m) = n - m$
دا حساب کرے۔
\end{prob}

\begin{prob}
\emph{برابری:} اک
ٹیورنگ مشین بناؤ
جیڑی ہِٹھلے فنکشن
دا حساب کرے:
\[
\fn{equality}(n,m) = 
\begin{cases}
  \text{1} & \text{جے~$n = m$} \\
  \text{0} & \text{جے~$n \neq m$}
\end{cases}
\]
جتھے~$n$ تے~$m \in \PosInt$
نیں۔
\end{prob}

\begin{prob}
فنکشن $\min(x,y)$
دا حساب کرن لئی
ٹیورنگ مشین بناؤ،
جتھے $x$ تے $y$
مثبت صحیح اعداد
نیں؛ پٹی اُتے اوہ
$\TMstroke$ نشاناں
دیاں لڑیاں نال
دکھائے جاندے نیں
جنہاں دے وچ
اک $\TMblank$ ہوندا
اے۔ مشین دے
حرفاں دے مجموعے
وچ ہور نشان وی
ورت سکدے او۔

فنکشن $\min$
اپنے آرگومنٹاں وچوں
سب توں چھوٹی
قدر چُنندا اے؛
ایس لئی
$\min(3,5)=3$،
$\min(20,16)=16$،
$\min(4,4)=4$
تے ایس طرح اگے۔
\end{prob}

\begin{defn}
ٹیورنگ مشین~$M$
جزوی فنکشن
$f\colon \Nat^k
  \pto \Nat$ دا
حساب کردی اے
اُدوں تے صرف
اُدوں جدوں:
  \begin{enumerate}
    \item $M$ انپٹ
    $\TMstroke^{n_1}\concat\TMblank\concat
    \dots \concat\TMblank\concat\TMstroke^{n_k}$
    اُتے آؤٹ پٹ
    $\TMstroke^{m}$
    نال رُکدی اے، جے
    $f(n_1, \dots, n_k) = m$
    ہووے۔
    \item $M$ یا تاں بالکل
    نہیں رُکدی،
    یا ایہو جیہے آؤٹ پٹ
    نال رُکدی اے
    جیڑا صرف
    ~$\TMstroke$
    نشاناں دا اک
    ٹکڑا نہیں، جے
    $f(n_1, \dots, n_k)$
    غیر تعریف شدہ
    ہووے۔
  \end{enumerate}
\end{defn}
\endgroup
% /OLP-0258

% OLP-0259 content/turing-machines/machines-computations/halting-states.tex
\begingroup
\def\olfilename{halting-states}
\expandafter\def\csname ol@sectioncs\endcsname{section}

\olfileid{tur}{mac}{hal}
\olsection{رُکّن والیاں حالتاں}

\begin{explain}
بھاویں اسی اپنی مشیناں
لئی ایہ تعریف رکھی اے
کہ اوہ صرف اُدوں رُکدیاں
نیں جدوں چلاؤن لئی
کوئی حکم نہ ہووے، ٹیورنگ
مشیناں دی عام نمائندگی
وچ اک خاص \emph{رُکّن
والی حالت}~$h$ ہوندی
اے، تے $h \in Q$ اے۔

رُکّن والی حالت دا
خیال سادا اے: جدوں
مشین اپنا کم مکمل
کر لیندی اے (انپٹ
قبول کرن لئی تیار
ہو جاوے، یا آؤٹ پٹ
لکھ چکی ہووے)، اوہ
حالت~$h$ وچ چلی جاندی
اے تے اوس حالت وچ رُک جاندی اے۔
کجھ مشیناں دیاں
رُکّن والیاں دو حالتاں
ہوندیاں نیں: اک انپٹ
قبول کرن لئی تے
دوجی انپٹ رد کرن لئی۔
\end{explain}

\begin{ex}\emph{رُکّن والیاں حالتاں}۔
ایس خیال نوں کھولن لئی
اسی جفت مشین نوں تھوڑھا
بدل کے شروع کردے آں۔
جے انپٹ دیاں لکیراں دی گنتی جفت ہووے تاں
مشین نوں حالت~$q_0$
وچ رُکّان دی بجائے
اسی اک حکم ودھا سکدے
آں جیڑا مشین نوں رُکّن
والی حالت وچ لے جاوے۔
\[
\begin{tikzpicture}[->,>=stealth',shorten >=1pt,auto,node distance=2.8cm,
                    semithick]
  \tikzstyle{every state}=[fill=none,draw=black,text=black]

  \node[initial, state]         (A)                     {$q_0$};
  \node[state]         (B) [right of=A] {$q_1$};
  \node[state]         (C) [below of=A] {$h$};

  \path (A) edge [bend left] node {\TMtrans{\TMstroke}{\TMstroke}{R}} (B)
  	    edge node {\TMtrans{\TMblank}{\TMblank}{N}} (C)
        (B) edge [loop above] node {\TMtrans{\TMblank}{\TMblank}{R}} (B)
            edge [bend left] node {\TMtrans{\TMstroke}{\TMstroke}{R}} (A);
\end{tikzpicture}
\]

ہن مثال نوں ہور ودھائیے۔
جدوں مشین نوں پتا لگدا اے
کہ انپٹ دیاں لکیراں دی گنتی طاق اے،
اوہ کدی نہیں رُکدی۔ لوپ والے حکم نوں
رد کرن والی حالت~$r$
تک جان والے حکم نال
بدل کے اسی مشین وچ
اک \emph{رد} والی حالت
شامل کر سکدے آں۔
\[
\begin{tikzpicture}[->,>=stealth',shorten >=1pt,auto,node distance=2.8cm,
                    semithick]
  \tikzstyle{every state}=[fill=none,draw=black,text=black]

  \node[initial,state]         (A)                     {$q_0$};
  \node[state]         (B) [right of=A] {$q_1$};
  \node[state]         (C) [below of=A] {$h$};
  \node[state]         (D) [below of=B] {$r$};

  \path (A) edge [bend left] node {\TMtrans{\TMstroke}{\TMstroke}{R}} (B)
  	    edge node {\TMtrans{\TMblank}{\TMblank}{N}} (C)
        (B) edge node {\TMtrans{\TMblank}{\TMblank}{N}} (D)
            edge [bend left] node {\TMtrans{\TMstroke}{\TMstroke}{R}} (A);
\end{tikzpicture}
\]
\end{ex}

\begin{explain}
ایہو جہیاں صورتاں وچ
خاص رُکّن والی حالت
شامل کرنی فائدہ مند
ہو سکدی اے، کیوں جے
ایس نال صاف پتا لگدا
اے کہ مشین کدوں کسے
انپٹ نوں قبول یا رد
کردی اے۔ پر ایہ گل
یاد رکھنی ضروری اے
کہ خاص رُکّن والی
حالت ودھاؤن نال مشین
دی حساب کرن دی طاقت
نہیں ودھدی۔ ایس طرح
رُکّن دے گھٹ رسمی
خیال دے وی اپنے فائدے
نیں۔ ایس باب وچ ہن
تک ورتی گئی رُکّن
دی تعریف نال
\emph{رُکّن دے مسئلے}
بارے ثبوت سمجھنا تے
وکھانا سوکھا ہو جاندا اے۔ ایس
لئی اسی اپنی اصل
تعریف ہی ورتدے رہندے
آں۔
\end{explain}
\endgroup
% /OLP-0259

% OLP-0260 content/turing-machines/machines-computations/disciplined-machines.tex
\begingroup
\def\olfilename{disciplined-machines}
\expandafter\def\csname ol@sectioncs\endcsname{section}

\olfileid{tur}{mac}{dis}
\olsection{پابند مشیناں}

\begin{explain}
\olref[hal]{sec} وچ
اسی اوہ ٹیورنگ مشیناں
ویکھیاں سن جنہاں دی
اکّو مقرر رُکّن والی
حالت~$h$ اے۔ جے ایہ
مشیناں کدی رُکّن،
تاں لازماً حالت~$h$
وچ ہی رُکّن گیاں۔
ایس طرح اکّو
رُکّن والی حالت
رکھن والیاں مشیناں
عام ٹیورنگ مشیناں
دے مقابلے وچ زیادہ
پابند نیں۔ مشیناں
دے عمل اُتے ہور
پابندیاں وی لا سکدے
آں۔ مثال دے طور
اُتے اسی مشین نوں
خانہ~$0$ اُتے پٹی
دے سِرے دا نشان
مٹاؤن توں، یا
خانہ~$0$ توں
کھبے پاسے جان دی
کوشش توں، منع
نہیں کیتا۔ (ساڈی
تعریف مطابق ایس
صورت وچ سر بس
خانہ~$0$ اُتے
ہی رہندا اے۔ کجھ ہور تعریفاں مطابق
ایس صورت وچ مشین رُک جاندی اے۔) ایہ وی کدے
کدے مفید ہوندا اے
کہ جے مشین رُکے
تاں خانہ~$1$
اُتے ہی رُکے۔
\end{explain}

\begin{defn}\ollabel{defn:disciplined}
ٹیورنگ مشین~$M$
\emph{پابند} اے اُدوں
تے صرف اُدوں جدوں:
\begin{enumerate}
    \item اوس دی اکّو
    مقرر رُکّن والی
    حالت~$h$ ہووے،
    \item جے اوہ رُکے
    تاں اوس دا سر خانہ~$1$
    پڑھ رہیا ہووے،
    \item اوہ $\TMendtape$
    نشان خانہ~$0$
    اُتے کدی نہ مٹاوے،
    تے
    \item اوہ خانہ~$0$
    توں کھبے پاسے
    جان دی کوشش کدی
    نہ کرے۔
\end{enumerate}
\end{defn}

\begin{explain}
اسی پہلے بیان کر
چکے آں کہ کسے وی
ٹیورنگ مشین نوں
اوہو عمل کرن والی اک
مشین وچ بدلیا
جا سکدا اے جیہدی
مقرر رُکّن والی حالت
ہووے۔ ایہ کم بس
نئی حالت~$h$
شامل کر کے، تے
ہر اوہو جیہے جوڑے لئی
$\delta(q, \sigma) = \tuple{h, \sigma, N}$
حکم شامل کر کے
ہوندا اے جتھے
$\tuple{q,\sigma}$ اُتے
اصل $\delta$~غیر
تعریف شدہ ہووے۔
ایہ سچ اے، اگرچہ
ثابت کرنا لمبا کم
اے، کہ ہر ٹیورنگ
مشین~$M$ نوں اک
پابند ٹیورنگ
مشین~$M'$ وچ
بدلیا جا سکدا اے
جیڑی اوہناں ہی
انپٹاں اُتے رُکے
تے اوہی آؤٹ پٹ
دیوے۔ مثال دے
طور اُتے، جے مشین
رُکّن لگے تے
خانہ~$1$ اُتے
نہ ہووے، تاں اسی
ہور حکم شامل کر
سکدے آں تاں جو
سر کھبے پاسے
چلدا رہے، جد تک
پٹی دے سِرے دا
نشان نہ لبھ لوے؛
فیر سر اک خانہ سجّے ول سرکے
تے مشین رُک جاوے۔ باقی شرطاں
کِویں پوری کیتیاں
جان، ایہ سوچݨ
لئی اسی تہانوں
چھڈدے آں۔
\end{explain}

\begin{ex}
    \olref{fig:adder-disc}
    وچ اسی
    \olref[una]{ex:adder}
    والی جمع دی مشین
    نوں پابند مشین
    وچ بدلدے آں۔
    \begin{figure}\[
\begin{tikzpicture}[->,>=stealth',shorten >=1pt,auto,node distance=2.8cm,
                    semithick]
  \tikzstyle{every state}=[fill=none,draw=black,text=black]

  \node[initial,state]         (A)              {$q_0$};
  \node[state]         (B) [right of=A] {$q_1$};
  \node[state]         (C) [below right of=B] {$q_2$};
  \node[state]         (D) [below left of=C] {$q_3$};
  \node[state]         (H) [left of=D] {$h$};

  \path (A) edge node {\TMtrans{\TMblank}{\TMstroke}{\TMstay}} (B)
           edge [loop below] node {\TMtrans{\TMstroke}{\TMstroke}{\TMright}} (B)
        (B) edge [loop below] node {\TMtrans{\TMstroke}{\TMstroke}{\TMright}} (B)
            edge node {\TMtrans{\TMblank}{\TMblank}{\TMleft}} (C)
        (C) edge node {\TMtrans{\TMstroke}{\TMblank}{\TMleft}} (D)
        (D) edge [loop above] node {\TMtrans{\TMstroke}{\TMstroke}{\TMleft}} (D)
        edge node {\TMtrans{\TMendtape}{\TMendtape}{\TMright}} (H);
\end{tikzpicture}
\]\caption{پابند جمع والی مشین}
\ollabel{fig:adder-disc}
\end{figure}
\end{ex}

\begin{prop}\ollabel{prop:disciplined}
ہر ٹیورنگ مشین~$M$
لئی اک پابند ٹیورنگ
مشین~$M'$ موجود اے
جیڑی آؤٹ پٹ~$O$
نال رُکدی اے جے
$M$~آؤٹ پٹ~$O$
نال رُکدی ہووے،
تے نہیں رُکدی
جے $M$~نہ رُکے۔
خاص طور اُتے،
جے فنکشن
$f\colon\Nat^n \to \Nat$
دا حساب کوئی
ٹیورنگ مشین
کردی اے، تاں
اوہدا حساب اک
پابند ٹیورنگ
مشین وی کردی
اے۔
\end{prop}

\begin{prob}\label{tur:mac:dis:prob:disc-succ}
فنکشن $f(x) = x+1$
دا حساب کرن والی
پابند مشین دیو۔
\end{prob}


\begin{prob}\label{tur:mac:dis:prob:copier}
ایسی پابند مشین
لبھو جیڑی
انپٹ $\TMstroke^n$
نال شروع ہو کے
آؤٹ پٹ
$\TMstroke^n \concat \TMblank \concat \TMstroke^n$
دیوے۔
\end{prob}
\endgroup
% /OLP-0260

% OLP-0261 content/turing-machines/machines-computations/combining-machines.tex
\begingroup
\def\olfilename{combining-machines}
\expandafter\def\csname ol@sectioncs\endcsname{section}

\olfileid{tur}{mac}{cmb}
\olsection{ٹیورنگ مشیناں نوں جوڑنا}

\begin{explain}
ٹیورنگ مشیناں دیاں جیہڑیاں مثالاں اسی ہن تک ویکھیاں نیں،
اوہ خاصیاں سادیاں سن۔ پر حقیقت وچ جیہڑا مسئلہ وی
کسے جدید پروگرامنگ زبان نال حل کیتا جا سکدا اے،
اوہ ٹیورنگ مشیناں نال وی حل کیتا جا سکدا اے۔
ہور پیچیدہ ٹیورنگ مشیناں بنانے لئی ایہ سمجھنا
اہم اے کہ اسی مشیناں نوں آپس وچ جوڑ سکدے آں:
عمل نوں سادے حصیاں وچ ونڈ کے وڈے مسئلے
حل کرن والی مشین بنا سکدے آں۔
جے پیچیدہ مسئلے نوں فطری انداز وچ حصیاں وچ
ونڈیا جا سکدا ہووے، تاں اسی کئی مرحلیاں وچ
ہر حصے لئی سادی مشین بنا کے اوہناں نوں اک
مشین وچ جوڑ سکدے آں جیڑی پورا مسئلہ حل کرے۔
اگلے حصے وچ رُکّن دے مسئلے نال نبڑدیاں ایہ گل
خاص طور اُتے اہم ہووے گا۔

ٹیورنگ مشیناں $M = \tuple{Q, \Sigma, q_0, \delta}$
تے~$M' = \tuple{Q', \Sigma', q_0', \delta'}$
نوں اسی کِویں جوڑئیے؟ ہن $M$~دے رُکّن مگروں
پٹی دی مکمل حالت نوں مشین~$M'$ دے چلن دی
شروع والی مکمل حالت بنا لیندے آں۔
ایہ کم کرن والی اک ٹیورنگ مشین $M \frown M'$
بنان لئی ایہ قدم چکو:
\begin{enumerate}
    \item حالتاں~$Q'$ نوں، جیڑیاں $M'$ دیاں نیں،
    نویں نمبر یا ناں دیو، تاں جو $M$
    تے~$M'$ دی کوئی حالت سانجھی نہ ہووے
    ($Q \cap Q' = \emptyset$)۔
    \item $M \frown M'$ دیاں حالتاں $Q \cup Q'$ نیں۔
    \item پٹی دے نشاناں دا مجموعہ $\Sigma \cup \Sigma'$ اے۔
    \item ابتدائی حالت~$q_0$ اے۔
    \item حالتی تبدیلی والا فنکشن $\delta''$
    ایہ اے:
    \[\delta''(q,\sigma) =
    \begin{cases}
      \delta(q,\sigma) & \text{جے $q \in Q$ ہووے تے $\delta(q,\sigma)$ تعریف شدہ ہووے}\\
      \delta'(q,\sigma) & \text{جے $q \in Q'$ ہووے}\\
      \tuple{q_0', \sigma, \TMstay} & \text{جے $q \in Q$ ہووے تے
      $\delta(q,\sigma)$ غیر تعریف شدہ ہووے}
    \end{cases}\]
\end{enumerate}
جوڑی ہوئی مشین وچ $M$ دا حکم ورتیا جاندا اے
جدوں حالت $q \in Q$ ہووے تے اوس ویلے کوئی حکم تعریف شدہ ہووے؛
$M'$ دے حکم حالت~$q \in Q'$ وچ ورتے جاندے نیں۔
جے حالت $q \in Q$ ہووے تے سر اوہ
نشان~$\sigma$ پڑھدا ہووے جیہدے لئی $M$
دی حالتی تبدیلی تعریف شدہ نہیں
(یعنی $M$ ایتھے رُک جاندی)،
تاں مشین~$M'$ دی ابتدائی حالت وچ
چلی جاندی اے۔ پٹی دے نشان
تے سر دی جگہ اوہو رہندے نیں۔

یاد رکھو: جے مشین~$M$ پابند نہ ہووے،
تاں $M$~دے رُکّن ویلے سر دی جگہ دا
سانوں پتا نہیں ہوندا۔ ایس لئی
$M$~دی آخری مکمل حالت وچ سر دا
خانہ~$1$ پڑھنا ضروری نہیں۔
مشیناں جوڑدیاں ایہ گل ذہن وچ رکھنی چاہیدی اے۔
\end{explain}

\begin{ex}
\emph{مشیناں نوں جوڑنا:} اسی ایہو جیہی اک مشین بناواں گے
جیڑی $\TMstroke$ دیاں دو لڑیاں،
اک دی لمبائی $n$ تے دوجی دی $m$،
انپٹ وچ لے کے رُکے تاں پٹی اُتے
$2(m+n)$ $\TMstroke$ دی اکّو لڑی ہووے۔
ایہ بناؤن لئی اسی دو پہلاں توں
جانیاں مشیناں جوڑ سکدے آں:
جمع والی مشین تے دوگنا کرن والی مشین۔
پہلاں جمع والی مشین دا حالتی خاکہ بناندے آں۔
\[
\begin{tikzpicture}[->,>=stealth',shorten >=1pt,auto,node distance=2.8cm,
                    semithick]
  \tikzstyle{every state}=[fill=none,draw=black,text=black]

  \node[initial,state] (A)              {$q_0$};
  \node[state]         (B) [right of=A] {$q_1$};
  \node[state]         (C) [right of=B] {$q_2$};

  \path (A) edge node {\TMtrans{\TMblank}{\TMstroke}{\TMstay}} (B)
            edge [loop above] node {\TMtrans{\TMstroke}{\TMstroke}{\TMright}} (B)
        (B) edge [loop above] node {\TMtrans{\TMstroke}{\TMstroke}{\TMright}} (B)
            edge node {\TMtrans{\TMblank}{\TMblank}{\TMleft}} (C)
        (C) edge [loop above] node {\TMtrans{\TMstroke}{\TMblank}{\TMstay}} (C);
\end{tikzpicture}
\]
حالت~$q_2$ وچ رُکّن دی بجائے
اسی عمل جاری رکھنا چاہندے آں
تاں جو حاصل نوں دوگنا کیتا جا سکے۔
یاد کرو کہ دوگنا کرن والی مشین
انپٹ دی پہلی لکیر مٹا کے وکھرے
آؤٹ پٹ وچ دو لکیراں لکھدی اے۔
اک حکم ہور شامل کریے تاں جو
سر جمع والی مشین دے آؤٹ پٹ
دی پہلی لکیر پڑھ رہیا ہووے۔
\[
\begin{tikzpicture}[->,>=stealth',shorten >=1pt,auto,node distance=2.8cm,
                    semithick]
  \tikzstyle{every state}=[fill=none,draw=black,text=black]

  \node[initial,state] (A)              {$q_0$};
  \node[state]         (B) [right of=A] {$q_1$};
  \node[state]         (C) [right of=B] {$q_2$};
  \node[state]         (D) [below left of=C] {$q_3$};
  \node[state]         (E) [below left of=D] {$q_4$};

  \path (A) edge node {\TMtrans{\TMblank}{\TMstroke}{\TMstay}} (B)
            edge [loop above] node {\TMtrans{\TMstroke}{\TMstroke}{\TMright}} (B)
        (B) edge [loop above] node {\TMtrans{\TMstroke}{\TMstroke}{\TMright}} (B)
            edge node {\TMtrans{\TMblank}{\TMblank}{\TMleft}} (C)
        (C) edge node {\TMtrans{\TMstroke}{\TMblank}{\TMleft}} (D)
        (D) edge [loop left] node {\TMtrans{\TMstroke}{\TMstroke}{\TMleft}} (D)
            edge node[left, xshift=-2mm] {\TMtrans{\TMendtape}{\TMendtape}{\TMright}} (E);
\end{tikzpicture}
\]
ہن انپٹ نوں دوگنا کرنا سوکھا اے:
دوگنا کرن والی مشین نوں حالت~$q_4$
نال جوڑ دیو۔ ایس لئی اوس مشین دیاں
حالتیں نویں ناں نال لکھنیاں پین گیاں
تاں جو اوہ~$q_4$ توں شروع ہوون،
نہ کہ~$q_0$ توں؛ ایس طرح دو
ابتدائی حالتیں نہیں بنن گیاں۔
آخری خاکہ \olref{fig:combined} ورگا ہووے گا۔
\begin{figure}
\[
\begin{tikzpicture}[->,>=stealth',shorten >=1pt,auto,node distance=2.8cm,
                    semithick]
  \tikzstyle{every state}=[fill=none,draw=black,text=black]
  \node[initial,state] (A)              {$q_0$};
  \node[state]         (B) [right of=A] {$q_1$};
  \node[state]         (C) [right of=B] {$q_2$};
  \node[state]         (D) [below left of=C] {$q_3$};
  \node[state]         (E) [below left of=D] {$q_4$};
  \node[state]         (2) [right of=E] {$q_5$};
  \node[state]         (3) [right of=2] {$q_6$};
  \node[state]         (4) [below of=3] {$q_7$};
  \node[state]         (5) [left of=4]  {$q_8$};
  \node[state]         (6) [left of=5]  {$q_9$};

  \path (A) edge node {\TMtrans{\TMblank}{\TMstroke}{\TMstay}} (B)
            edge [loop above] node {\TMtrans{\TMstroke}{\TMstroke}{\TMright}} (B)
        (B) edge [loop above] node {\TMtrans{\TMstroke}{\TMstroke}{\TMright}} (B)
            edge node {\TMtrans{\TMblank}{\TMblank}{\TMleft}} (C)
        (C) edge node {\TMtrans{\TMstroke}{\TMblank}{\TMleft}} (D)
        (D) edge [loop left] node {\TMtrans{\TMstroke}{\TMstroke}{\TMleft}} (D)
            edge node[left, xshift=-2mm] {\TMtrans{\TMendtape}{\TMendtape}{\TMright}} (E)
    (E) edge node {\TMtrans{\TMstroke}{\TMblank}{\TMright}} (2)
    (2) edge [loop above] node {\TMtrans{\TMstroke}{\TMstroke}{\TMright}} (2)
      edge node {\TMtrans{\TMblank}{\TMblank}{\TMright}} (3)
    (3) edge [loop above] node {\TMtrans{\TMstroke}{\TMstroke}{\TMright}} (3)
        edge node {\TMtrans{\TMblank}{\TMstroke}{\TMright}} (4)
    (4) edge [loop below] node {\TMtrans{\TMblank}{\TMstroke}{\TMleft}} (4)
        edge node {\TMtrans{\TMstroke}{\TMstroke}{\TMleft}} (5)
    (5) edge [loop below]  node {\TMtrans{\TMstroke}{\TMstroke}{\TMleft}} (5)
        edge              node {\TMtrans{\TMblank}{\TMblank}{\TMleft}} (6)
    (6) edge [loop below] node {\TMtrans{\TMstroke}{\TMstroke}{\TMleft}} (6)
        edge              node {\TMtrans{\TMblank}{\TMblank}{\TMright}} (E);
\end{tikzpicture}
\]
\caption{جمع تے دوگنا کرن والیاں مشیناں نوں جوڑنا}
\ollabel{fig:combined}
\end{figure}
\end{ex}

\begin{prop}
    جے $M$ تے $M'$ پابند مشیناں
    بالترتیب فنکشن $f\colon
    \Nat^k \to \Nat$ تے $f'\colon \Nat \to \Nat$
    دا حساب کردیاں نیں، تاں
    $M \frown M'$ وی پابند اے
    تے~$\comp{f}{f'}$ دا حساب کردی اے۔
\end{prop}

\begin{proof}
    چونکہ $M$ پابند اے، اس لئی
    جدوں اوہ آؤٹ پٹ~$f(n_1,\dots,n_k) = m$
    دے کے رُکدی اے تاں سر خانہ~$1$
    پڑھ رہیا ہوندا اے۔ ہن جے اسی
    مشین~$M'$ دی ابتدائی حالت وچ
    داخل ہوئیے، تاں ایہ مشین آؤٹ پٹ
    $f'(m)$ دے کے رُکے گی تے سر فیر
    خانہ~$1$ پڑھدا ہووے گا۔
    \olref[dis]{defn:disciplined} دیاں
    باقی شرطاں وی پوری ہوندیاں نیں۔
\end{proof}

\begin{prob}
    فنکشن $f(x) = x+2$ دا حساب کرن
    والی پابند ٹیورنگ مشین دیو:
    \cref{tur:mac:dis:prob:disc-succ} والی
    مشین~$M$ لے کے $M \frown M$ بناؤ۔
\end{prob}
\endgroup
% /OLP-0261

% OLP-0262 content/turing-machines/machines-computations/variants.tex
\begingroup
\def\olfilename{variants}
\expandafter\def\csname ol@sectioncs\endcsname{section}

\olfileid{tur}{mac}{var}
\olsection{ٹیورنگ مشیناں دے وکھرے روپ}

ٹیورنگ مشیناں دی تعریف کرن دے کئی طریقے
ہو سکدے نیں؛ ساڈی تعریف اوہناں وچوں
صرف اک اے۔ کجھ گلاں وچ ساڈی تعریف
ہور تعریفاں توں زیادہ کھلی اے۔ اسی
نشاناں دا کوئی وی محدود مجموعہ پٹی
لئی منّ لیندے آں؛ اک زیادہ پابند تعریف
صرف دو نشان، $\TMstroke$
تے~$\TMblank$، منّ سکدی اے۔ اسی
مشین نوں پٹی اُتے نشان لکھن تے
نال ہی سر ہلاؤن دی اجازت دیندے آں؛
ہور تعریفاں وچ اک قدم اُتے
یا نشان لکھیا جا سکدا اے یا سر ہلایا جا سکدا اے۔
ساڈی تعریف وچ سر نوں ہلائے بنا
نشان لکھنا وی ممکن اے؛ کجھ ہور
تعریفاں $\TMstay$ والے “حکم” نوں
شامل نہیں کردیاں۔ دوجے پاسے کجھ
گلاں وچ ساڈی تعریف زیادہ پابند اے۔
اسی پٹی نوں صرف اک پاسے ول
لامتناہی منّیا اے؛ ہور تعریفاں وچ
پٹی کھبے تے سجے، دوہاں پاسیاں
ول لامتناہی ہو سکدی اے۔ بلکہ
وکھ وکھ پٹیاں دی کوئی وی تعداد،
یا خانیاں دی لامتناہی جالی وی
منّی جا سکدی اے۔ اسی ٹیورنگ
مشین دے حکماں نوں حالتی تبدیلی
والے فنکشن نال ظاہر کردے آں؛
ہور تعریفاں حالتی تبدیلی والا
تعلق ورتدیاں نیں، جیہدے وچ اکّو
صورتِ حال وچ مشین لئی اک توں
ودھ ممکن حکم ہو سکدے نیں۔

تعریف دی ایہ آخری ڈھیل خاص دلچسپ اے۔
ساڈی تعریف وچ جدوں مشین حالت~$q$
وچ نشان~$\sigma$ پڑھدی اے،
تاں $\delta(q, \sigma)$ نواں نشان،
نئی حالت تے سر دی اگلی جگہ
طے کردا اے۔ پر جے حکماں
دا مجموعہ موجودہ حالت-نشان
جوڑیاں $\tuple{q,
  \sigma}$ تے نویں حالت-نشان-سمت
تِکاں $\tuple{q', \sigma',
  D}$ دے وچکار تعلق ہووے،
تاں ٹیورنگ مشین دا عمل
اکّو طرح طے ہونا ضروری
نہیں۔ حکماں دے تعلق وچ
$\tuple{q,
  \sigma, q', \sigma', D}$ تے
$\tuple{q, \sigma, q'', \sigma'', D'}$
دوویں شامل ہو سکدے نیں۔
ایس صورت وچ مشین
\emph{غیر قطعی} (non-deterministic)
ٹیورنگ مشین ہوندی اے۔
حسابی پیچیدگی دے نظریے وچ
ایہناں دا اہم کردار اے۔

ٹیورنگ مشین کدوں رُکدی اے،
ایس بارے وی وکھرے رواج نیں:
ساڈے ہاں جدوں حالتی تبدیلی
والا فنکشن غیر تعریف شدہ ہووے
تاں مشین رُکدی اے؛ ہور تعریفاں
خاص مقرر رُکّن والی حالت
لازم کردیاں نیں۔ \olref[hal]{sec}
وچ اسی سمجھایا اے کہ مقرر
رُکّن والی حالت دی شرط
ٹیورنگ مشیناں دے حساب کرن
دے دائرے نوں نہیں گھٹاندی۔
ساڈیاں مشیناں دی پٹی صرف
اک پاسے ول لامتناہی اے،
ایس لئی کدی مشین کوئی
حکم ٹھیک طرح پورا نہیں
کر سکدی: جے سر سب توں
کھبے خانے نوں پڑھ رہیا
ہووے تے کھبے ول جانا
ہووے۔ ساڈی تعریف وچ سر
پٹی توں “ڈگّن” دی بجائے
اوہتھے ہی رہندا اے؛ پر
اسی ایہ وی تعریف کر سکدے
سن کہ ایس صورت وچ مشین
رُک جاوے۔ ایہ تعریف وی
ہم ارز اے: خانہ~$0$ توں
کھبے جان دی کوشش اُتے
رُکّن والی مشین دے عمل
دی نقل کرن لئی ہر اوہ
حالتی تبدیلی
$\delta(q,\TMendtape) =
\tuple{q',\sigma,\TMleft}$
حذف کر دیو۔ فیر
$\TMendtape$ اُتے کھبے
جان دی کوشش کرن دی
بجائے مشین رُک جاوے
گی۔\footnote{ایہ گل
\emph{پوری طرح} درست نہیں،
کیوں جے ساڈے قاعدیاں
وچ $\TMendtape$ نشان
خانہ~$0$ توں
علاوہ ہور خانیاں اُتے
وی لکھن
تے پڑھݨ توں کوئی
نہیں روکدا
(ویکھو \olref[una]{ex:mover})۔
ایس مشکل توں بچن لئی
ایس کم واسطے اک دوجا
نشان~$\TMendtape'$
شامل کیتا جا سکدا اے۔}

اعداد نوں ظاہر کرن دے
طریقے وی وکھرے ہو سکدے
نیں، ایس لئی ٹیورنگ مشین
دا حساب کیتا ہویا انپٹ-آؤٹ پٹ
فنکشن وی وکھرے روپ وچ
پیش ہو سکدا اے۔ اسی
اعداد دی اک نشانی نمائندگی
ورتدے آں؛ ثنائی نمائندگی
وی ورتی جا سکدی اے۔
ایس لئی $\TMblank$
تے~$\TMendtape$ توں
علاوہ دو نشان لوڑیں دے نیں۔

ہن اک دلچسپ حقیقت ویکھو:
ایہناں وکھرے روپاں وچوں
کسی روپ نال وی ایہ گل نہیں
بدلدی کہ کیہڑے فنکشن
ٹیورنگ مشین نال حساب ہو
سکدے نیں۔ \emph{جے کوئی
فنکشن اک تعریف مطابق
ٹیورنگ-حسابی اے، تاں اوہ
ایہناں ساریاں تعریفاں
مطابق ٹیورنگ-حسابی اے۔}

اسی ایس گل دی پوری
تصدیق دیاں تفصیلاں وچ
نہیں جاواں گے۔ بس اک
مثال ویکھو: دوہاں پاسیاں
ول لامتناہی پٹی یا کئی
پٹیاں دی اجازت نال
حسابی طاقت نہیں ودھدی۔
وجہ، موٹے طور اُتے،
ایہ اے کہ اکّو، اک
پاسے ول لامتناہی پٹی
والی ٹیورنگ مشین کئی
پٹیاں یا دوہاں پاسیاں
ول لامتناہی پٹی دی
نقل کر سکدی اے۔ مثال
دے طور اُتے، ودھ
حالتیں تے حکم ورت کے
کئی پٹیاں والی یا
دوہاں پاسیاں ول
لامتناہی پٹی والی
مشین دے پروگرام نوں
اک، اک پاسے ول
لامتناہی پٹی والی
مشین دے پروگرام وچ
“ترجمہ” کیتا جا
سکدا اے۔ نئی مشین
پٹی~$1$ دے خانیاں
(یا دوہاں پاسیاں
والی پٹی دے “مثبت”
خانیاں) لئی جُفت
خانے، تے پٹی~$2$
دے خانیاں (یا
“منفی” خانیاں)
لئی طاق خانے
ورت سکدی اے۔
\endgroup
% /OLP-0262

% OLP-0263 content/turing-machines/machines-computations/church-turing-thesis.tex
\begingroup
\def\olfilename{church-turing-thesis}
\expandafter\def\csname ol@sectioncs\endcsname{section}

\olfileid{tur}{mac}{ctt}
\olsection{چرچ--ٹیورنگ دا دعویٰ}

ٹیورنگ مشیناں دا مقصد
مؤثر طریقۂ کار دے تصور
دی جگہ اک درست
ریاضیاتی تعریف دینا اے۔
ٹیورنگ دا خیال سی کہ
جیڑھا بندہ مؤثر طریقۂ کار
تے ٹیورنگ مشین، دوہاں
دے تصور سمجھدا ہووے،
اوہنوں ایہ گل فطری
لگے گی کہ مؤثر
طریقۂ کار نال کیتا
جا سکّن والا ہر
کم ٹیورنگ مشین وی
کر سکدی اے۔
ایس دعوے نوں ایہ گل
وی سہارا دیندی اے
کہ مؤثر طریقۂ کار
دی جیہڑیاں ہور
درست ریاضیاتی تعریفاں
پیش کیتیاں گئیاں،
اوہ ٹیورنگ مشین
دے تصور نال
حساب ہو سکّن والے فنکشناں دے لحاظ
نال ہم ارز نکلدیاں نیں۔ یعنی اوہ بالکل
اوہناں ہی فنکشناں
دا حساب کر سکدیاں
نیں۔ ایس دعوے
نوں \emph{چرچ--ٹیورنگ
دا دعویٰ} کہندے نیں۔

\begin{defn}[چرچ--ٹیورنگ دا دعویٰ]
\emph{چرچ--ٹیورنگ دا دعویٰ}
ایہ اے کہ مؤثر طریقۂ کار
نال جیہڑا وی حساب ہو سکدا اے،
اوہ ٹیورنگ مشین وی کر سکدی اے۔
\end{defn}

چرچ--ٹیورنگ دے دعوے
نوں دو طریقیاں نال
ورتیا جاندا اے۔
پہلا طریقہ سستی
دا بہانہ بن سکدا اے۔
فرض کرو ساڈے کول
کسے شے دا حساب کرن
لئی مؤثر طریقۂ کار
دا بیان ہووے، مثلاً
“فرضی کوڈ” وچ۔
فیر اسی چرچ--ٹیورنگ
دے دعوے دا حوالہ
دے کے ایہ کہہ
سکدے آں کہ کوئی
ٹیورنگ مشین وی
اوہی فنکشن دا حساب
کردی اے، بھاویں
اسی حقیقت وچ اوہ
مشین بنائی نہ ہووے۔

چرچ--ٹیورنگ دے دعوے
دی دوجی ورتوں فلسفیانہ
طور اُتے زیادہ دلچسپ
اے۔ ایہ ثابت کیتا
جا سکدا اے کہ
کجھ فنکشن ٹیورنگ
مشیناں نال حساب
نہیں ہو سکدے۔
ایس نتیجے تے
چرچ--ٹیورنگ دے
دعوے نوں ملا کے
ایہ نتیجہ نکلدا اے
کہ اوہناں دا حساب
کسے وی مؤثر طریقے نال نہیں
ہو سکدا۔ جے کوئی ایہو جیہا مؤثر
طریقۂ کار ہوندا،
تاں چرچ--ٹیورنگ
دے دعوے مطابق
اوہدے لئی ٹیورنگ
مشین وی ہوندی۔
ایس لئی جے اسی
ثابت کر دیئے کہ
کسے فنکشن دا حساب
کرن والی کوئی
ٹیورنگ مشین نہیں،
تاں اوہدے لئی
مؤثر طریقۂ کار
وی نہیں ہو سکدا۔
خاص طور اُتے،
چرچ--ٹیورنگ دے
دعوے نوں ایہ کہن
لئی ورتیا جاندا اے
کہ رُکّن دے مسئلے دا حل
کسے ٹیورنگ مشین کول نہیں،
تے ایہ مسئلہ کسے وی مؤثر
طریقے نال حل نہیں ہو سکدا۔
\endgroup
% /OLP-0263

\OLEndChapterHook
\endgroup
% /OLP-0253

% OLP-0264 content/turing-machines/undecidability/undecidability.tex
\begingroup
\def\olfilename{undecidability}
\expandafter\def\csname ol@sectioncs\endcsname{section}

\olchapter{tur}{und}{ناقابلِ فیصلہ ہونا}

% OLP-0265 content/turing-machines/undecidability/introduction.tex
\begingroup
\def\olfilename{introduction}
\expandafter\def\csname ol@sectioncs\endcsname{section}

\olfileid{tur}{und}{int} 
\olsection{تعارف}

ایہ گل شاید صاف لگے
کہ کجھ فنکشن، ایتھوں تک کہ
کجھ حسابی فنکشن وی،
حساب کرن جوگ نہیں۔
فنکشن بہت
ودھ نیں تے اوہناں
دا چلن بہت پیچیدہ
ہو سکدا اے۔ مثال
دے طور اُتے
تابکار ذریاں دے
ٹُٹّن، یا ہور
بے قاعدہ یا اتفاقی
چلن نال تعریف
ہون والے فنکشن
ویکھو۔ فرض کرو
اسی کسے مقرر ویلے
توں اک اک سیکنڈ
دے وقفے گننا
شروع کریے،
تے فنکشن $f(n)$
نوں کائنات دے
اوہناں ذریاں دی
گنتی رکھئیے جیڑے
اُس ابتدائی ویلے
توں بعد $n$-ویں
اک سیکنڈ دے
وقفے وچ ٹُٹدے
نیں۔ ایہ ایسا
فنکشن لگدا اے
جس دا حساب
کرن دی امید
اسی کدی نہیں
رکھ سکدے۔

پر ایہ تصور نہ کر
سکنا کہ ایہو جیہے فنکشن
دا حساب کِویں
کیتا جاوے، اوہنوں
واقعی ناقابلِ حساب
ثابت کرن توں
بالکل وکھ گل اے۔
حقیقت وچ بظاہر
انتہائی پیچیدہ
فنکشن وی تجریدی
معنیٰ وچ حساب
ہو سکدا اے۔
مثلاً فرض کرو
کائنات دا وقت
محدود اے:
کجھ کائناتی
نظریاں مطابق
بہت دور مستقبل
وچ کائنات اک
نکتے وچ سکڑ
جاوے گی۔ فیر
ابتدائی ویلے توں
بعد صرف محدود،
بھاویں ناقابلِ
تصور حد تک
وڈی، تعداد وچ
سیکنڈ نیں جنہاں
لئی $f(n)$
تعریف شدہ اے۔
تے جیہڑا فنکشن
صرف محدود
انپٹاں اُتے
تعریف شدہ ہووے،
اوہ حساب کرن
جوگ اے: اسی
اوہدے آؤٹ پٹ
اک وڈے جدول
وچ لکھ سکدے
آں، یا اک بہت
وڈے ٹیورنگ
مشین دے حالتی
تبدیلی والے
خاکے وچ کوڈ
کر سکدے آں۔

اکثر سانوں اوہناں
فنکشناں دیاں
خاص صورتاں وچ
دلچسپی ہوندی اے
جنہاں دیاں قدراں
ہاں یا ناں والے
سوالاں دے جواب
دیندیاں نیں۔
مثلاً “کی $n$
اولی عدد اے؟”
سوال نال ایہ
فنکشن جڑیا اے:
\[
\fn{isprime}(n) = \begin{cases}
  1 & \text{جے $n$ اولی عدد ہووے}\\
  0 & \text{ہور صورت وچ۔}
  \end{cases}
\]
اسی کہندے آں کہ
ہاں/ناں والے
سوال دا \emph{مؤثر
طریقے نال فیصلہ}
ہو سکدا اے جے
اوہدے نال جڑیا
$1/0$-قدری
فنکشن مؤثر
طریقے نال حساب
ہو سکدا ہووے۔

ریاضیاتی طور اُتے
ایہ ثابت کرن لئی
کہ کجھ فنکشن
مؤثر طریقے نال
حساب نہیں ہو
سکدے، یا کجھ
مسئلیاں دا مؤثر
طریقے نال فیصلہ
نہیں ہو سکدا،
حساب دا اک خاص
نمونہ طے کرنا
ضروری اے۔ فیر
ایہ دکھانا پیندا
اے کہ کجھ فنکشن
اوہ نمونہ حساب
نہیں کر سکدا،
یا کجھ مسئلیاں
دا فیصلہ نہیں
کر سکدا۔ مثلاً
اسی دکھا سکدے
آں کہ کجھ فنکشن
ٹیورنگ مشین نال حساب
نہیں ہو سکدے تے کجھ
مسئلیاں دا فیصلہ
ٹیورنگ مشین نہیں کر سکدی۔
فیر چرچ--ٹیورنگ
دے دعوے دا حوالہ
دے کے ایہ نتیجہ
کڈھ سکدے آں کہ
صرف ٹیورنگ
مشیناں ہی ہر
فنکشن دا حساب
کرن توں قاصر
نہیں، بلکہ کوئی
وی مؤثر طریقۂ
کار ایہ کم نہیں
کر سکدا۔

ایہناں منفی نتیجیاں
نوں ثابت کرن دی
اہم گل ایہ اے
کہ اسی خود ٹیورنگ
مشیناں نوں عدد
دے سکدے آں۔
ایس دا سوکھا
طریقہ اوہناں دی
گنتی وار فہرست
بناؤنا اے: مثلاً
ٹیورنگ مشیناں
تے اوہناں دے
پروگرام لکھن دا
اک مقرر طریقہ
چن کے فیر
اوہناں نوں
باقاعدہ ترتیب
نال فہرست وچ
لکھو۔ جدوں
ایہ ممکن ہونا
صاف ہو جاوے،
تاں سیٹاں دے سائز
دی سادی دلیل نال
ٹیورنگ مشین
توں ناقابلِ
حساب فنکشناں
دا وجود نکلدا
اے: $\Nat$ توں~$\Nat$
فنکشناں دا سیٹ
(بلکہ صرف
$\Nat$ توں
$\{0, 1\}$
فنکشناں دا سیٹ
وی) !!{nonenumerable}
اے؛ پر ساریاں
ٹیورنگ مشیناں
دی فہرست بن
سکدی اے، ایس
لئی ٹیورنگ
مشین نال حساب
ہو سکّن والے
فنکشناں دا
سیٹ صرف
!!{denumerable}
اے۔

اسی \emph{خاص}
فنکشن تے مسئلے
وی تعریف کر
سکدے آں جنہاں
نوں بالترتیب
ناقابلِ حساب
تے ناقابلِ فیصلہ
ثابت کیتا جا
سکدا اے۔ اک
مسئلہ اوہ اے جیہنوں
\emph{رُکّن دا مسئلہ}
کہندے نیں۔ ٹیورنگ
مشین دے حکماں
دی فہرست نال
اوس دا محدود
بیان دتا جا
سکدا اے۔ ٹیورنگ
مشین دا ایسا
بیان، یعنی
مشین دا پروگرام،
یقیناً دوجی
ٹیورنگ مشین
لئی انپٹ
بن سکدا اے۔
ایس لئی اسی
ایہناں مشیناں
بارے سوالاں
دا فیصلہ کرن
والیاں ٹیورنگ
مشیناں اُتے
غور کر سکدے
آں۔ اک خاص
دلچسپ سوال
ایہ اے:
“کی دتی گئی
ٹیورنگ مشین
انپٹ~$n$
نال شروع ہو
کے آخرکار
رُکے گی؟”
چنگا ہوندا جے
ایس سوال دا
فیصلہ کرن
والی ٹیورنگ
مشین ہوندی:
ایہنوں معیار
جانچ والی
مشین سمجھو
جیڑی دیکھدی
کہ ہور مشیناں
لامتناہی چکّر
وچ نہ پھسن۔
دلچسپ حقیقت،
جیڑی ٹیورنگ
نے ثابت کیتی،
ایہ اے کہ
ایسی مشین
موجود نہیں۔
کوئی اک ٹیورنگ
مشین ایسی
نہیں ہو سکدی
جیڑی ٹیورنگ
مشین~$M$
دے بیان تے
کسے عدد~$n$
نوں انپٹ
وچ لے کے
ہمیشہ رُکے
تے آؤٹ پٹ
$1$ یا $0$
ایس حساب نال
دیوے کہ
$M$ مشین
انپٹ $n$
نال شروع ہو
کے رُکدی
اے یا نہیں۔

جدوں ساڈے کول
خاص ناقابلِ
فیصلہ مسئلیاں
دیاں مثالاں
ہون، تاں اسی
اوہناں نال ہور
مسئلیاں دا وی
ناقابلِ فیصلہ
ہونا دکھا سکدے
آں۔ مثال دے
طور اُتے اک
مشہور ناقابلِ
فیصلہ سوال
ایہ اے: “کی
پہلے درجے دا
!!{formula}~$!A$
درست اے؟”
کوئی ٹیورنگ مشین ایہو جیہی نہیں
جیہنوں انپٹ وجوں پہلے درجے دا
!!{formula}~$!A$ دتا جاوے
تے اوہ ہر صورت رُکے تے
آؤٹ پٹ $1$
یا $0$ ایس
حساب نال
دیوے کہ $!A$
درست اے
یا نہیں۔
تاریخی طور
اُتے ایس
مسئلے نوں
مؤثر طریقے
نال حل کرن
والا طریقۂ
کار لبھݨ
دے سوال
نوں بس
“فیصلے دا
مسئلہ”
کہندے سن؛
ایس لئی
اسی کہندے
آں کہ
فیصلے دا
مسئلہ حل
نہیں ہو
سکدا۔
ٹیورنگ
تے چرچ
نے ایہ
نتیجہ لگ
بھگ اکّو
ویلے وکھ
وکھ ثابت
کیتا سی،
ایس لئی
ایہنوں
چرچ--ٹیورنگ
دا قضیہ
وی کہندے
نیں۔
\endgroup
% /OLP-0265

% OLP-0266 content/turing-machines/undecidability/enumerating-tms.tex
\begingroup
\def\olfilename{enumerating-tms}
\expandafter\def\csname ol@sectioncs\endcsname{section}

\olfileid{tur}{und}{enu}
\olsection{ٹیورنگ مشیناں دی گنتی وار فہرست}

\begin{explain}
اسی دکھا سکدے آں کہ ساریاں ٹیورنگ مشیناں دا مجموعہ !!{enumerable}
اے۔ ایہ ایس گل توں نکلدا اے کہ ہر ٹیورنگ مشین دا محدود بیان
دتا جا سکدا اے۔ حالتاں دا مجموعہ تے پٹی دے نشاناں دا مجموعہ
دونوں محدود نیں۔ حالتی تبدیلی دا فنکشن $Q \times \Sigma$
توں $Q \times \Sigma \times \{\TMleft, \TMright, \TMstay\}$
تک اک جزوی فنکشن اے؛ ایس لئی جن محدود دلیلی جوڑیاں اُتے
ایہ تعریف شدہ اے، اوہناں لئی ایس دیاں قدراں فہرست وچ لکھ
کے ایہنوں وی متعین کیتا جا سکدا اے۔

ایہ گل اپنی حد تک ٹھیک اے، پر اتھے اک باریک فرق اے۔ ٹیورنگ
مشین دی تعریف نے ایس گل اُتے کوئی پابندی نہیں لائی کہ حالتاں
دے مجموعے تے پٹی دے حروف دے !!{element}s کیہڑے ہو سکدے نیں۔
مثلاً ہر حقیقی عدد لئی، تکنیکی طور اُتے، اک ایسی ٹیورنگ مشین
اے جیڑی اُس عدد نوں اپنی حالت بنا لَیندی اے۔ پر مشین دا
\emph{عمل} ایس گل توں آزاد اے کہ حالت تے حرف دے طور اُتے کیہڑیاں
اشیا ورتیاں گئیاں نیں۔ \olref{fig:variants} وچ دتیاں دو ٹیورنگ
مشیناں ویکھو۔
\begin{figure}
\begin{center}
\begin{tikzpicture}[->,>=stealth',shorten >=1pt,auto,node distance=2.8cm,
                    semithick]
  \tikzstyle{every state}=[fill=none,draw=black,text=black]

  \node[initial,state]         (A)              {$q_0$};
  \node[state]         (B) [right of=A] {$q_1$};

  \path (A) edge [bend left] node {\TMtrans{\TMstroke}{\TMstroke}{\TMright}} (B)
        (B) edge [loop above] node {\TMtrans{\TMblank}{\TMblank}{\TMright}} (B)
            edge [bend left] node {\TMtrans{\TMstroke}{\TMstroke}{\TMright}} (A);
\end{tikzpicture}\\
\begin{tikzpicture}[->,>=stealth',shorten >=1pt,auto,node distance=2.8cm,
  semithick]
\tikzstyle{every state}=[fill=none,draw=black,text=black]

\node[initial,state]         (A)              {$s$};
\node[state]         (B) [right of=A] {$h$};

\path (A) edge [bend left] node {\TMtrans{A}{A}{\TMright}} (B)
(B) edge [loop above] node {\TMtrans{\TMblank}{\TMblank}{\TMright}} (B)
edge [bend left] node {\TMtrans{A}{A}{\TMright}} (A);
\end{tikzpicture}
\end{center}
\caption{\emph{جُفت} مشین دے وکھرے روپ}
\ollabel{fig:variants}
\end{figure}
ایہ دو خاکے دو مشیناں نوں ظاہر کردے نیں: $M$ دی پٹی دا
حرفی مجموعہ $\Sigma = \{\TMendtape,\TMblank,\TMstroke\}$
تے حالتاں دا مجموعہ $\{q_0,q_1\}$ اے؛ $M'$ دا حرفی
مجموعہ $\Sigma' =
\{\TMendtape,\TMblank,A\}$ تے حالات $\{s,h\}$ نیں۔ پر
باقی ہدایات اکّو جیہیاں نیں: $M$، $n$ عدد $\TMstroke$
نشاناں دی لڑی اُتے اُدوں تے صرف اُدوں رُکدی اے جدوں $n$ جُفت
ہووے؛ تے $M'$، $n$ عدد $A$ نشاناں دی لڑی اُتے، اُدوں تے
صرف اُدوں رُکدی اے جدوں $n$ جُفت ہووے۔ اسی صرف $\TMstroke$
دا ناں بدل کے~$A$، $q_0$ دا بدل کے~$s$، تے $q_1$ دا
بدل کے~$h$ رکھیا اے۔ ایہ مثال عام اصول دِندی اے: جے
اک مشین دوجی توں حالتاں تے نشاناں دے ناں ایس طرح بدلݨ نال
ملدی اے، تاں اسی اوہناں نوں اکّو جیہی سمجھ سکدے آں۔
بلکہ اسی مشین دے نشاناں تے حالتاں نوں بس مثبت صحیح اعداد
ہی سمجھ سکدے آں: $\sigma_0$ دی تھاں~$1$، $\sigma_1$
دی تھاں~$2$، وغیرہ؛ $\TMendtape$ نوں~$1$، $\TMblank$
نوں~$2$، وغیرہ۔ ایس طرح \emph{جُفت} مشین
\olref{fig:standard-even} وچ وکھائی مشین بن جاندی اے۔
\begin{figure}
\[\begin{tikzpicture}[->,>=stealth',shorten >=1pt,auto,node distance=2.8cm,
  semithick]
\tikzstyle{every state}=[fill=none,draw=black,text=black]

\node[initial,state]         (A)              {$1$};
\node[state]         (B) [right of=A] {$2$};

\path (A) edge [bend left] node {\TMtrans{3}{3}{\TMright}} (B)
(B) edge [loop above] node {\TMtrans{2}{2}{\TMright}} (B)
edge [bend left] node {\TMtrans{3}{3}{\TMright}} (A);
\end{tikzpicture}
\]
\caption{اک \emph{معیاری} جُفت مشین}
\ollabel{fig:standard-even}
\end{figure}
جس ٹیورنگ مشین دیاں حالتاں تے نشان مثبت صحیح اعداد ہون،
اسی اوہنوں \emph{معیاری} مشین کہہ سکدے آں تے ہن توں صرف
ایہناں مشیناں اُتے غور کر سکدے آں۔\footnote{“معیاری
مشین” اک رائج معیاری اصطلاح نہیں۔}

ساڈا مقصد ایہ دکھانا سی کہ ٹیورنگ مشیناں دا مجموعہ
!!{enumerable} اے۔ اُپرلی گل دے پیشِ نظر ایس لئی ایہ
دکھانا کافی اے کہ معیاری ٹیورنگ مشیناں دا مجموعہ
!!{enumerable} اے۔ فرض کرو ساڈے کول اک معیاری مشین
$M = \tuple{Q, \Sigma, q_0, \delta}$ اے۔ اسی مثبت صحیح
اعداد دی اک محدود لڑی نال ایہدا بیان کِویں دیواں؟ پہلے
حالتاں دی تعداد، فیر خود حالات، نشاناں دی تعداد، خود نشان
تے ابتدائی حالت فہرست وچ لکھانگے۔ (چیتے رکھو کہ ایہ
سارے مثبت صحیح اعداد نیں، کیوں جے $M$ معیاری مشین اے۔)
ہن~$\delta$ دا کیہ کریے؟ ممکنہ آرگومنٹاں، یعنی
$\tuple{q,\sigma}$ جوڑیاں، دا مجموعہ محدود اے، کیوں جے
$Q$ تے~$\Sigma$ محدود نیں۔ ایس لئی~$\delta$ دی ساری
معلومات بس اوہناں $5$-رُکنی سلسلیاں
$\tuple{q, \sigma, q', \sigma', d}$ دی محدود فہرست اے
جنہاں لئی $\delta(q,\sigma) = \tuple{q', \sigma', D}$
ہووے۔ اتھے $d$ اک عدد اے جیڑا سمت~$D$ نوں کوڈ کردا
اے (مثلاً $1$ نال~$\TMleft$، $2$ نال~$\TMright$، تے
$3$ نال~$\TMstay$)۔

ایس طرح ہر معیاری ٹیورنگ مشین دا بیان مثبت صحیح اعداد دی
اک محدود فہرست، یعنی اک سلسلے $s_M \in
(\PosInt)^*$، نال دتا جا سکدا اے۔ مثلاً معیاری
\emph{جُفت} مشین نوں ایہ سلسلہ کوڈ کردا اے:
\[
2, \underbrace{1, 2}_Q, 3, \overbrace{1, 2, 3}^\Sigma, 1, \underbrace{1, 3, 2, 3, 2}_{\delta(1,3) = \tuple{2,3,R}}, 
\overbrace{2, 2, 2, 2, 2}^{\delta(2,2) = \tuple{2,2,R}},
\underbrace{2, 3, 1, 3, 2}_{\delta(2,3) = \tuple{1,3,R}}.
\]
\end{explain}

\begin{thm}
$\Nat$ توں~$\Nat$ تک ایسے فنکشن موجود نیں جیڑے ٹیورنگ
مشین نال حساب نہیں ہو سکدے۔
\end{thm}

\begin{proof}
سانوں پتا اے کہ مثبت صحیح اعداد دے محدود سلسلیاں
~$(\PosInt)^*$ دا مجموعہ !!{enumerable} اے
(\cref{sfr:siz:zigzag:prob:posint-star})۔ ایس توں نکلدا
اے کہ معیاری ٹیورنگ مشیناں دے بیان، جو~$(\PosInt)^*$
دا ذیلی مجموعہ نیں، خود وی گنتی وار فہرست بنن جوگ نیں۔
ہر ٹیورنگ-حسابی فنکشن $\Nat$ توں~$\Nat$ تک کسے نہ
کسے (اصل وچ کئی) ٹیورنگ مشین نال حساب ہوندا اے۔ اوس
مشین دیاں حالتاں تے نشانیاں دے ناں مثبت صحیح اعداد
رکھ کے، خاص طور اُتے $\TMendtape$ نوں~$1$، $\TMblank$
نوں~$2$، تے $\TMstroke$ نوں~$3$ رکھ کے، اسی ویکھدے آں
کہ ہر ٹیورنگ-حسابی فنکشن دا حساب اک معیاری ٹیورنگ مشین
وی کردِی اے۔ ایس دا مطلب اے کہ $\Nat$ توں~$\Nat$ تک
سارے ٹیورنگ-حسابی فنکشناں دا مجموعہ وی گنتی وار فہرست
بنن جوگ اے۔

دوجے پاسے، $\Nat$ توں~$\Nat$ تک سارے فنکشناں دا
مجموعہ !!{enumerable} نہیں اے
(\cref{sfr:siz:red:prob:nat-nat})۔ جے سارے فنکشن
کسی نہ کسی ٹیورنگ مشین نال حساب ہو سکدے ہوندے، تاں
اوہناں نوں حساب کرن والیاں مشیناں دے بیان فہرست وچ
لکھ کے اسی سارے فنکشناں دی فہرست بنا سکدے ہوندے۔
ایس لئی کجھ فنکشن ٹیورنگ مشین نال حساب نہیں ہو سکدے۔
\end{proof}

\begin{prob}
  کی تسی ٹیورنگ مشیناں دا ایسا بیان سوچ سکدے او جس وچ
  حالات تے حرفی نشاناں نوں صاف صاف فہرست وچ لکھن
  دی لوڑ نہ ہووے؟ تسی “معیاری” مشین دا اپنا تصور
  تعریف کر سکدے او، پر ایہ وی دسو کہ تہاڈے نویں معنیٰ
  وچ ہر ٹیورنگ مشین دی حسابی کارروائی اک “معیاری”
  مشین کیوں کر سکدی اے۔
\end{prob}
\endgroup
% /OLP-0266

% OLP-0267 content/turing-machines/undecidability/universal-tm.tex
\begingroup
\def\olfilename{universal-tm}
\expandafter\def\csname ol@sectioncs\endcsname{section}

\olfileid{tur}{und}{uni}
\olsection{ہمہ گیر ٹیورنگ مشیناں}

\olref[enu]{sec} وچ اسی گل کیتی سی کہ ہر ٹیورنگ مشین دا
بیان صحیح اعداد دے اک محدود سلسلے نال دتا جا سکدا اے۔
ایہ سلسلہ مشین دیاں حالتاں، حرفی نشاناں، ابتدائی حالت
تے ہدایات نوں کوڈ کردا اے۔ اسی ایہ وی دَسیا سی کہ
ایہناں سارے بیانواں دا مجموعہ !!{enumerable} اے۔ کیوں
جے ایہناں بیانواں دا مجموعہ !!{denumerable} اے، ایس دا
مطلب اے کہ~$\Nat$ توں ایہناں بیانواں تک اک
!!a{surjective} فنکشن موجود اے۔ ایہو جیہا !!a{surjective}
فنکشن، مثلاً، کینٹر دے ٹیڑھے-میڑھے طریقے نال لبھ
سکدا اے۔ ایہ سانوں ساریاں ٹیورنگ مشیناں (دے بیانواں)
دی گنتی وار فہرست بناؤن دا طریقہ دیندا اے۔ ایہو جیہی اک
فہرست مقرر کر لئیے، تاں $1$ نمبر، $2$ نمبر، \dots، $e$ نمبر
والی ٹیورنگ مشین دی گل کر سکدے آں۔ ایہناں
عدداں نوں \emph{اشاریے} کہندے نیں۔

\begin{defn}
جے $M$~ساڈی مقررہ فہرست وچ $e$ نمبر والی ٹیورنگ مشین اے،
تاں اسی $e$~نوں~$M$ دا اک \emph{اشاریہ} کہندے آں۔
اسی $M_e$ لکھ کے $e$ نمبر والی ٹیورنگ مشین مراد لیندے آں۔
\end{defn}

اک مشین دے کئی اشاریے ہو سکدے نیں۔ مثلاً~$M$ دے
دو بیانواں وچ اوہدیاں ہدایات فہرست وچ لکھن دی ترتیب
وکھ ہو سکدی اے، تے ایہناں وکھرے بیانواں دے اشاریے
وی وکھ ہون گے۔

اہم گل ایہ اے کہ ٹیورنگ مشیناں دے بیانواں دی فہرست
ایس طرح بنائی جا سکدی اے کہ اشاریے توں~$M$ دا بیان
مؤثر طریقے نال حساب ہو جاوے، تے مشین~$M$ دے بیان توں
اوس دا اک اشاریہ وی مؤثر طریقے نال حساب ہو جاوے۔
فیر چرچ--ٹیورنگ دے دعوے مطابق ایسی ٹیورنگ مشین
لبھ سکدی اے جیڑی اشاریے~$e$ والی مشین دا بیان
حاصل کرے تے اوہ بیان اپنی پٹی اُتے آؤٹ پٹ دے طور
اُتے لکھ دے۔ ایہ بیان~$\TMstroke$ دیاں لڑیوں
دے سلسلے دی صورت وچ ہووے گا (ہر لڑی~$M_e$
دے بیان والے سلسلے دے اک مثبت صحیح عدد دی نمائندگی کرے گی)۔

ایس توں بعد ایہ سوال فطری اے: ٹیورنگ مشین دے اشاریے
اُتے تعریف ہون والے کیہڑے فنکشن خود ٹیورنگ مشین نال
حساب ہو سکدے نیں؟ ٹیورنگ مشین دے اشارییاں دیاں
کیہڑیاں خاصیتاں دا فیصلہ ٹیورنگ مشین کر سکدی اے؟
اک مثال: جیہڑا فنکشن اشاریے~$e$ نوں اوس مشین دیاں
حالتاں دی تعداد ول بھیجدا اے جیہدا اشاریہ~$e$ اے،
اوہ ٹیورنگ مشین نال حساب ہو سکدا اے۔ ایسی مشین ایہ کم
کرے گی: جے اوہنوں پٹی اُتے $e$~$\TMstroke$ نشاناں دی
اک لڑی دتی جاوے، تاں پہلے اوہ $e$ دا کوڈ کھول کے
مشین دا بیان حاصل کرے گی۔ ہن ایہ بیان پٹی اُتے
~$\TMstroke$ دیاں لڑیوں دے سلسلے نال ظاہر اے۔ ایس
سلسلے دا پہلا !!{element} حالتاں دی تعداد اے۔ ایس لئی
ہن مشین نوں پہلی $\TMstroke$ لڑی توں علاوہ باقی سب
مٹانا تے فیر رُکنا اے۔

اک قابلِ توجہ نتیجہ ایہ اے:

\begin{thm}\ollabel{thm:universal-tm} اک \emph{ہمہ گیر
  ٹیورنگ مشین}~$U$ موجود اے جیڑی انپٹ $\tuple{e,n}$
  نال شروع ہووے تاں
  \begin{enumerate}
    \item اُدوں تے صرف اُدوں رُکدی اے جدوں $M_e$ انپٹ~$n$
    اُتے رُکدی اے، تے
    \item جے $M_e$ آؤٹ پٹ $m$ نال رُکے، تاں~$U$ وی
    اوہو آؤٹ پٹ لے کے رُکدی اے۔
  \end{enumerate}
  ایس طرح $U$ اوہ فنکشن $f\colon \Nat \times \Nat \pto \Nat$
  دا حساب کردی اے جیہدے لئی $f(e,n) = m$ ہووے جے $M_e$
  انپٹ~$n$ نال شروع ہو کے آؤٹ پٹ~$m$ دے کے رُکے،
  تے دوجی صورت وچ ایہ تعریف شدہ نہ ہووے۔
\end{thm}

\begin{proof}
  $U$ نوں حقیقت وچ مکمل بنا کے پیش کرنا عملاً ناممکن
  جیہا اے، کیوں جے ایہ بے حد پیچیدہ مشین اے۔ پر اسی
  اوہدے عمل دا خاکہ بیان کر سکدے آں تے فیر چرچ--ٹیورنگ
  دے دعوے نوں لاگو کر سکدے آں۔ شروع وچ $U$ دی پٹی
  اُتے $e$ عدد وچ $\TMstroke$ نشاناں دی اک لڑی اے،
  تے اوس دے بعد $n$~$\TMstroke$ نشاناں دی لڑی اے۔
  مشین پہلے اشاریے~$e$ دا کوڈ انپٹ~$n$ دے سجّے
  پاسے کھولدی اے۔ ایس نال اعداد دی اک فہرست ملدی
  اے، یعنی~$\TMstroke$ دیاں لڑیوں جیڑیاں~$\TMblank$
  نال وکھ نیں، تے ایہ مشین~$M_e$ دیاں ہدایات بیان
  کردیاں نیں۔ فیر $U$، مشین~$M_e$ دی ابتدائی حالت
  دا عدد تے عدد~$1$، مشین~$M_e$ دے بیان دے سجّے
  پاسے لکھدی اے۔ (ایہ وی اک نشان والی نمائندگی وچ
  $\TMstroke$ دیاں لڑیوں نیں۔) اگّے اوہ انپٹ،
  یعنی $n$~$\TMstroke$ نشاناں والی لڑی، سجّے
  پاسے نقل کردی اے؛ پر ہر $\TMstroke$ دی تھاں
  تین $\TMstroke$ نشاناں دی لڑی رکھدی اے
  (چیتے رکھو کہ $\TMstroke$ نشان دا عدد~$3$ اے؛
  $1$،~$\TMendtape$ دا عدد، تے $2$،~$\TMblank$
  دا عدد اے)۔ ایہناں لڑیوں دے سلسلے دے کھبے
  سِرے اُتے، جیڑیاں~$\TMblank$ نال $U$ دی پٹی
  اُتے وکھ نیں، اوہ اکّو~$\TMstroke$ لکھدی اے؛
  ایہ~$\TMendtape$ دا کوڈ اے۔

  ہن $U$ دی پٹی اُتے ایہ سب کجھ اے: اشاریہ~$e$،
  عدد~$n$، ابتدائی حالت دا کوڈ عدد (یعنی “موجودہ
  حالت”)، سر دی ابتدائی جگہ~$1$ دا عدد (یعنی
  “سر دی موجودہ جگہ”)، تے مصنوعی “پٹی” دا ابتدائی
  مواد۔ آخری چیز~$\TMstroke$ دیاں لڑیوں دا اک سلسلہ
  اے جو~$M_e$ دے نشاناں دے کوڈ عدداں، یعنی
  مصنوعی “نشاناں”، دی نمائندگی کردا اے؛ لڑیوں نوں
  ~$\TMblank$ وکھ کردا اے۔

  ہن ایہ مشین اوہ عمل نقل کردی اے جیڑا $M_e$
  انپٹ~$n$ اُتے شروع ہو کے کردِی؛ اوہ ایہ قدم چکدی اے:
  \begin{enumerate}
    \item “سر دی موجودہ جگہ” دا عدد~$k$ لبھو
    (شروع وچ ایہ~$1$ اے)؛
    \item “پٹی” دے $k$-ویں ٹکڑے ول جا کے اُتھوں دا
    “نشان” ویکھو؛
    \item\ollabel{find-inst}%
    موجودہ “حالت” تے “نشان” نال میل کھاندی ہدایت لبھو؛
    \item “پٹی” دے $k$-ویں ٹکڑے ول واپس جا کے اوتھے
    دے “نشان” دی تھاں اُس نشان دا کوڈ عدد لکھو
    جیڑا $M_e$ لکھدی؛
    \item سر نوں اُس تھاں لے جاؤ جتھے موجودہ
    “حالت” درج اے، تے اوتھے دا عدد نویں حالت
    دے عدد نال بدل دو؛
    \item جتھے “پٹی دی جگہ” درج اے اُتھے جا کے
    اک~$\TMstroke$ مٹا دو یا اک~$\TMstroke$ ودھا دو
    (ترتیب نال، جے ہدایت کھبّے یا سجّے ہلّن لئی
    آکھے)؛
    \item ایہ کم دوبارہ کرو۔\footnote{ایتھے اسی کجھ باریک
    اوکھائیاں دا پورا حساب نہیں دِتا۔ مثلاً “سر دی
    موجودہ جگہ” دا حساب رکھن والی گنتی ودھے تاں
    $U$~نوں ہور جگہ چاہیدی ہو سکدی اے؛ اوس صورت
    وچ اوہنوں ساری مصنوعی “پٹی” سجّے پاسے
    سرکاؤنی پَوے گی۔}
  \end{enumerate}
  جے $M_e$ انپٹ~$n$ اُتے شروع ہو کے کدی نہیں
  رُکدی، تاں $U$ وی نہیں رُکدی؛ ایس لئی اوہدا
  آؤٹ پٹ تعریف شدہ نہیں۔

  جے قدم~\olref{find-inst} وچ پتا لگے کہ~$M_e$ دے
  بیان وچ موجودہ “حالت”/“نشان” جوڑی لئی کوئی ہدایت
  نہیں، تاں $M_e$ رُک جاندی۔ ایسا ہووے تاں $U$
  اپنی پٹی دا مصنوعی “پٹی” دے کھبے پاسے والا حصہ
  مٹا دیندی اے۔ تین~$\TMstroke$ نشاناں دی ہر لڑی
  لئی، جیڑی اک~$\TMstroke$ دی نمائندگی کردی اے
  (ایہ نشان $M_e$ دی پٹی اُتے اے)، اوہ اپنی پٹی دے
  کھبے سِرے اُتے اک $\TMstroke$ لکھدی تے نال نال مصنوعی
  “پٹی” نوں مٹاندی جاندی اے۔ ایہ کم مکّے تاں
  $U$ دی پٹی اُتے~$\TMstroke$ نشاناں دی
  لمبائی~$m$ والی اکّو لڑی رہ جاندی اے۔
  
  جے $U$ نوں مصنوعی “پٹی” اُتے تین~$\TMstroke$
  نشاناں دی لڑی توں وکھ کجھ ملے، تاں اوہ فوراً
  رُک جاندی اے۔ ایس صورت وچ $U$ دی پٹی اُتے
  ~$\TMstroke$ نشاناں دی اکّو لڑی نہیں ہوندی؛
  ایس لئی اوہدا آؤٹ پٹ قدرتی عدد نہیں، یعنی
  ایس صورت وچ $f(e,n)$ تعریف شدہ نہیں۔
\end{proof}
\endgroup
% /OLP-0267

% OLP-0268 content/turing-machines/undecidability/halting-problem.tex
\begingroup
\def\olfilename{halting-problem}
\expandafter\def\csname ol@sectioncs\endcsname{section}

\olfileid{tur}{und}{hal}
\olsection{رُکّن دا مسئلہ}

\begin{explain}
فرض کرو اسی ٹیورنگ مشیناں دے بیانواں دی اک گنتی وار
فہرست مقرر کر لئی اے۔ ایس طرح ہر ٹیورنگ مشین نوں اک
\emph{اشاریہ} ملدا اے: مشین دے بیانواں دی فہرست
$M_1$، $M_2$، $M_3$، \dots{} وچ اوہدے بیان دی جگہ۔

سانوں پتا اے کہ کجھ فنکشناں دا حساب ٹیورنگ مشین نال
نہیں ہو سکدا: مشیناں دے بیانواں دا مجموعہ، تے ایس لئی
ٹیورنگ مشیناں دا مجموعہ، !!{enumerable} اے، پر
$\Nat$ توں $\Nat$ تک سارے فنکشناں دا مجموعہ
ایہو جیہا نہیں۔ پر اسی ناقابلِ حساب فنکشناں دیاں خاص
مثالاں وی لبھ سکدے آں۔ اک مثال رُکّن دا فنکشن اے۔
\end{explain}

\begin{defn}[رُکّن دا فنکشن]
 \emph{رُکّن دا فنکشن}~$h$ ایویں تعریف ہوندا اے:
\[
h(e,n) =
\begin{cases}
  \text{0} & \text{جے مشین~$M_e$ انپٹ $n$ اُتے نہ رُکے} \\
  \text{1} & \text{جے مشین~$M_e$ انپٹ $n$ اُتے رُکے}
\end{cases}
\]
\end{defn}

\begin{defn}[رُکّن دا مسئلہ]
\emph{رُکّن دا مسئلہ} ایہ طے کرن دا مسئلہ اے کہ کسے وی
$e$، $n$ لئی ٹیورنگ مشین~$M_e$ انپٹ وچ~$n$ لکیراں
لے کے رُکدی اے یا نہیں۔
\end{defn}

\begin{explain}
اسی~$h$ دا ٹیورنگ-ناقابلِ حساب ہونا اک متعلقہ فنکشن
~$s$ دا ٹیورنگ-ناقابلِ حساب ہونا دکھا کے ثابت کراں گے۔
ایہ ثبوت ایس گل اُتے ٹکیا اے کہ جیہڑا وی کم ٹیورنگ
مشین حساب کر سکدی اے، اوہ پابند ٹیورنگ مشین وی کر
سکدی اے (\olref[mac][dis]{sec})؛ تے ایس گل اُتے
کہ دو ٹیورنگ مشیناں جوڑ کے اک مشین بنائی جا سکدی اے
(\olref[mac][cmb]{sec})۔
\end{explain}

\begin{defn} فنکشن~$s$ ایویں تعریف ہوندا اے:
\[
s(e) =
\begin{cases}
  \text{0} & \text{جے مشین~$M_e$ انپٹ $e$ اُتے نہ رُکے} \\
  \text{1} & \text{جے مشین~$M_e$ انپٹ $e$ اُتے رُکے}
\end{cases}
\]
\end{defn}

\begin{lem}
فنکشن~$s$ ٹیورنگ مشین نال حساب نہیں ہو سکدا۔
\end{lem}

\begin{proof}
تضاد لئی فرض کرو کہ فنکشن~$s$ ٹیورنگ مشین نال حساب
ہو سکدا اے۔ فیر اک ٹیورنگ مشین~$S$ ہووے گی جیڑی
~$s$ دا حساب کرے۔ عمومیّت گھٹائے بغیر اسی فرض کر
سکدے آں کہ جدوں $S$ رُکدی اے، اوہدے سر تھلّے
پٹی دا پہلا خانہ ہوندا اے (یعنی اسی اوہنوں پابند روپ وچ
لے سکدے آں)۔ ایہ مشین اک دوجی مشین~$J$ نال
“جوڑی” جا سکدی اے۔ جے اوہ انپٹ~$0$ نال شروع
ہووے، یعنی پٹی دے سِرے دے نشان دے سجّے پاسے
پہلے خانے اُتے ابتدائی حالت وچ~$\TMblank$ پڑھے،
تاں اوہ رُک جاندی اے؛ دوجی صورت وچ سجّے پاسے
چلدی رہندی اے تے کدی نہیں رُکدی۔ $S \concat J$،
یعنی $S$ نوں~$J$ نال جوڑ کے بنی مشین، خود ٹیورنگ
مشین اے؛ ایس لئی ایہ $M_e$ اے، کسے نہ کسے~$e$ لئی
(یعنی فہرست وچ کسے نہ کسے تھاں آندی اے)۔ ہن $M_e$
نوں~$e$ عدد $\TMstroke$ لکیراں دے انپٹ اُتے
شروع کرو۔ دو صورتاں نیں: یا تاں $M_e$ رُکدی اے
یا نہیں رُکدی۔
\begin{enumerate}
\item فرض کرو $M_e$ انپٹ وچ $e$ عدد $\TMstroke$
  لکیراں لے کے رُکدی اے۔ فیر $s(e)
  = 1$ اے۔ ایس لئی $S$، انپٹ~$e$ اُتے شروع
  ہو کے، پٹی اُتے اکّو $\TMstroke$ آؤٹ پٹ دے
  کے رُکدی اے۔ فیر $J$ پٹی اُتے اک $\TMstroke$
  دے نال شروع ہوندی اے۔ ایس صورت وچ $J$ نہیں
  رُکدی۔ پر $M_e$، $S
  \concat J$ مشین اے؛ ایس لئی اوہنوں عین اوہو
  کرنا چاہیدا اے جیڑا $S$ توں بعد $J$ کردیاں نیں
  (یعنی ایتھے سجّے پاسے چلدے رہنا تے کدی نہ رُکنا)۔
  ایس لئی $M_e$ انپٹ وچ $e$ عدد $\TMstroke$
  لکیراں لے کے رُک نہیں سکدی۔

\item ہن فرض کرو $M_e$ انپٹ وچ $e$ عدد
  $\TMstroke$ لکیراں لے کے نہیں رُکدی۔ فیر
  $s(e) = 0$ اے، تے $S$، انپٹ~$e$ اُتے
  شروع ہو کے، خالی پٹی دے نال رُکدی اے۔ $J$
  خالی پٹی اُتے شروع ہو کے فوراً رُکدی اے۔ فیر
  $M_e$ وی اوہو کردی اے جیڑا $S$ توں بعد~$J$
  کردیاں نیں؛ ایس لئی $M_e$ انپٹ وچ $e$ عدد
  $\TMstroke$ لکیراں لے کے ضرور رُکدی اے۔
\end{enumerate}
دونوں صورتاں وچ ساڈے فرض دے خلاف نتیجہ نکلیا۔
ایس توں پتا لگدا اے کہ ایہو جیہی ٹیورنگ مشین~$S$
نہیں ہو سکدی: $s$~ٹیورنگ مشین نال حساب نہیں ہو سکدا۔
\end{proof}

\begin{thm}[رُکّن دے مسئلے دا ناقابلِ حل ہونا]
\ollabel{thm:halting-problem} رُکّن دا مسئلہ حل نہیں ہو
سکدا؛ یعنی فنکشن~$h$ ٹیورنگ مشین نال حساب نہیں
ہو سکدا۔
\end{thm}

\begin{proof}
فرض کرو $h$ ٹیورنگ مشین نال حساب ہو سکدا ہووے،
مثلاً اک مشین~$H$ نال۔ اسی $H$ ورت کے~$s$ دا
حساب کرن والی مشین بنا سکدے آں: پہلے انپٹ دی
اک نقل بنا لو (تے دوہاں نوں اک~$\TMblank$ نشان
نال وکھ رکھو)۔ فیر شروع ول واپس آ کے~$H$ چلاؤ۔
پہلا کم کرن والی مشین صاف طور اُتے بن سکدی اے
(\cref{tur:mac:dis:prob:copier})؛ جے $H$ موجود
ہوندی، تاں اسی اوہنوں ایہو جیہی نقل کرن والی مشین نال
“جوڑ” کے نویں مشین بنا لیندے جیڑی طے کردی کہ
$M_e$ انپٹ~$e$ اُتے رُکدی اے یا نہیں، یعنی~$s$
دا حساب کردی۔ پر اسی پہلے ہی دکھا چکّے آں کہ
ایہو جیہی کوئی مشین نہیں ہو سکدی۔ ایس لئی $h$~وی
ٹیورنگ مشین نال حساب نہیں ہو سکدا۔
\end{proof}

\begin{prob}
تِن-رُکّن (3-Halt) دا مسئلہ ایہ اے کہ کسے وی چنی
ہوئی ٹیورنگ مشین لئی ایہ طے کرن دا فیصلہ جاتی
طریقہ دِتا جاوے کہ اوہ باقی خالی پٹی اُتے تِن
$\TMstroke$ لکیراں دے انپٹ نال رُکدی اے یا نہیں۔
ثابت کرو کہ تِن-رُکّن دا مسئلہ حل نہیں ہو سکدا۔
\end{prob}

\begin{prob}
دکھاؤ کہ جے مشین/انپٹ جوڑیاں $M_e$ تے~$n$
لئی، جتھے $e \neq n$، رُکّن دا مسئلہ حل ہو سکدا
ہووے، تاں $e = n$ والی صورتاں لئی وی حل ہو
سکدا اے۔
\end{prob}

\begin{prob}
اسی ثابت کیتا کہ رُکّن دا مسئلہ حل نہیں ہو سکدا جدوں
انپٹ اک عدد~$e$ ہووے، جیڑا ساریاں ٹیورنگ
مشیناں دی فہرست راہیں اک مشین~$M_e$ دی نشاندہی
کردا اے۔ پر جے اسی \olref[tur][und][enu]{sec} وچ دتے
مشیناں دے بیان ای سدھے
انپٹ بنن دی اجازت دیئیے تاں کیہ ہووے گا؟ کی کوئی
ٹیورنگ مشین مشین دے اشاریے دی تھاں اوہدا بیان
انپٹ وچ لے کے رُکّن دا مسئلہ حل کر سکدی اے؟
دسو کیوں یا کیوں نہیں۔
\end{prob}

\begin{prob} دکھاؤ کہ ایویں تعریف شدہ \emph{جزوی}
فنکشن~$s'$ ٹیورنگ مشین نال حساب \emph{ہو سکدا} اے:
  \[
  s'(e) =
  \begin{cases}
    \text{1} & \text{جے مشین~$M_e$ انپٹ $e$ اُتے رُکے}\\
    \text{غیر تعریف شدہ} & \text{جے مشین~$M_e$ انپٹ $e$ اُتے نہ رُکے}
  \end{cases}
  \]
\end{prob}
\endgroup
% /OLP-0268

% OLP-0269 content/turing-machines/undecidability/decision-problem.tex
\begingroup
\def\olfilename{decision-problem}
\expandafter\def\csname ol@sectioncs\endcsname{section}

\olfileid{tur}{und}{dec}
\olsection{فیصلے دا مسئلہ}

اسی پہلے درجے دی منطق نوں \emph{قابلِ فیصلہ} اُدوں تے
صرف اُدوں کہندے آں جدوں کوئی مؤثر طریقہ ایہ طے کردا
ہووے کہ دتا ہویا !!{sentence} درست اے یا نہیں۔ پر ایہو
جیہا طریقہ نہیں: پہلے درجے دے جملیاں دی درستی دا فیصلہ کرن
والا مسئلہ حل نہیں ہو سکدا۔

ایہ اہم منفی نتیجہ قائم کرن لئی اسی ثابت کراں گے کہ
ٹیورنگ مشین فیصلے دا مسئلہ حل نہیں کر سکدی۔ یعنی
ایہو جیہی کوئی ٹیورنگ مشین نہیں جیڑی، جدوں وی اوہ
ایس پٹی اُتے شروع ہووے جیہدے اُتے پہلے درجے دا !!{sentence}
لکھیا ہووے، آخرکار رُکے تے ایہ ویکھ کے کہ اوہ !!{sentence}
درست اے یا نہیں، $1$ یا~$0$ آؤٹ پٹ دے۔ چرچ--ٹیورنگ دے
دعوے مطابق ہر حساب کرن جوگ فنکشن ٹیورنگ مشین نال
وی حساب ہو سکدا اے۔ ایس لئی جے ایہ “درستی دا
فنکشن” کسے مؤثر طریقے نال حساب ہو سکدا ہووے،
تاں ٹیورنگ مشین نال وی حساب ہو سکدا اے۔ جے ایہ
ٹیورنگ مشین نال حساب نہیں ہو سکدا، تاں کوئی مؤثر
طریقہ وی ایہنوں حساب نہیں کر سکدا۔

فیصلے دے مسئلے دا ناقابلِ حل ہونا ثابت کرن لئی
اسی رُکّن دے مسئلے نوں فیصلے دے مسئلے تک گھٹاواں گے۔
ایس دا مطلب ایہ اے: اسی ثابت کر چکّے آں کہ فنکشن~$h(e,w)$،
جیہدی قدر $1$ اے جے~$e$ دے بیان والی ٹیورنگ مشین
انپٹ~$w$ اُتے رُکے تے دوجی صورت وچ~$0$ اے،
ٹیورنگ مشین نال حساب نہیں ہو سکدا۔ اسی دکھاواں گے کہ
جے پہلے درجے دے جملیاں دی درستی دا فیصلہ کرن والی ٹیورنگ مشین ہوندی،
تاں~$h$ دا حساب کرن والی ٹیورنگ مشین وی ہوندی۔
چونکہ $h$ ٹیورنگ مشین نال حساب نہیں ہو سکدا،
ایس لئی درستی دا فیصلہ کرن والی ٹیورنگ مشین
وی نہیں ہو سکدی۔

ایس حکمتِ عملی دا پہلا قدم ایہ دکھانا اے کہ ہر
انپٹ~$w$ تے ٹیورنگ مشین~$M$ لئی اسی مؤثر
طریقے نال اک !!a{sentence} $!T(M, w)$ بیان
کر سکدے آں، جیڑا~$M$ دیاں ہدایات تے
انپٹ~$w$ نوں ظاہر کرے، تے اک !!a{sentence}~$!E(M, w)$
وی بیان کر سکدے آں، جیڑا ایہ آکھے کہ “$M$
آخرکار رُکدی اے”، ایس طرح کہ:
\begin{quote}
  $\Entails !T(M, w) \lif !E(M,w)$ اُدوں تے صرف اُدوں جدوں $M$
  انپٹ~$w$ اُتے رُکدی اے۔
\end{quote}
ساڈے ثبوت دا وڈا حصہ ایہ جملے $!T(M, w)$
تے~$!E(M, w)$ بیان کرن تے ایہ جانچن وچ
لگے گا کہ $!T(M, w) \lif !E(M, w)$ اُدوں تے صرف
اُدوں درست اے جدوں $M$~انپٹ~$w$ اُتے رُکدی اے۔
\endgroup
% /OLP-0269

% OLP-0270 content/turing-machines/undecidability/representing-tms.tex
\begingroup
\def\olfilename{representing-tms}
\expandafter\def\csname ol@sectioncs\endcsname{section}

\olfileid{tur}{und}{rep}
\olsection{ٹیورنگ مشیناں دی منطقی نمائندگی}

\begin{explain}
ٹیورنگ مشیناں تے اوہناں دے چلن نوں پہلے درجے دی منطق دے
اک !!a{sentence} نال ظاہر کرن لئی سانوں اک موزوں زبان دی
تعریف کرنی پئے گی۔ ایس زبان دے دو حصے نیں: مشین دیاں مکمل
حالتاں بیان کرن لئی !!{predicate}s، تے چال دے قدماں (“ویلیاں”)
تے پٹی اُتے تھاماں نوں نمبر دین لئی عبارتاں۔

اسی دو قسماں دے !!{predicate}s ورتاں گے؛ دونیں دو-جگہی
نیں: ہر حالت~$q$ لئی اک !!a{predicate}~$\Obj Q_q$، تے پٹی دے
ہر نشان~$\sigma$ لئی اک !!a{predicate}~$\Obj S_\sigma$۔ پہلے
والے سانوں~$M$ دی حالت تے اوہدے پڑھن لکھن والے سر دی تھاں
دسن دیندے نیں؛ دوجے والے دسدے نیں کہ پٹی اُتے کیہڑے نشان نیں۔

پٹی اُتے سر دیاں تھاماں تے پورے ہو چکّے قدماں دی گنتی ظاہر
کرن لئی سانوں عدد ظاہر کرن دا طریقہ چاہیدا اے۔ ایس لئی
اسی اک !!a{constant}~$\Obj 0$ تے اک $1$-جگہی فنکشن~$\prime$،
یعنی جانشین فنکشن، ورتدے آں۔ رواج مطابق ایہ اپنی دلیل
توں \emph{بعد} لکھیا جاندا اے (تے اسی بریکٹ چھڈ دیندے آں)۔

ہر عدد~$n$ لئی اک مقررہ ٹرم~$\num{n}$ اے، یعنی~$n$ دی
\emph{عددی عبارت}، جیڑی~$\Lang L_M$ وچ اوہنوں ظاہر کردی اے۔
$\num{0}$، $\Obj 0$ اے؛ $\num{1}$، $\Obj 0'$ اے؛
$\num{2}$، $\Obj 0''$ اے؛ تے ایداں اگے۔ زیادہ رسمی طور اُتے:
\begin{align*}
\num{0} & = \Obj 0 \\
\num{n+1} &= \num{n}'
\end{align*}
ٹرم~$\num{0}$، یعنی~$\Obj 0$، پٹی دے سب توں کھبے خانے
دا ناں وی اے تے پہلا قدم چُکن توں پہلاں دے ویلے
(ابتدائی مکمل حالت) دا ناں وی۔ ٹرم~$\num{1}$، یعنی~$\Obj 0'$،
سب توں کھبے خانے دے سجّے پاسے والے خانے دا ناں اے، تے
پہلا قدم چُکن توں بعد دے ویلے دا وی؛ ایداں ہی اگے۔

اسی پٹی اُتے تھاماں دی ترتیب (جدوں ایس دا مطلب “دے کھبے
پاسے” ہووے) تے چلن دے قدماں دی ترتیب (جدوں مطلب “توں پہلاں”
ہووے) دونیں ظاہر کرن لئی اک !!a{predicate}~$<$ وی ورتاں گے۔

زبان مقرر کرن توں بعد اسی~$!T(M,
w)$ دے “مسلّمے” لکھاں گے، یعنی اوہ !!{sentence}s جیہڑے
اکٹھے ہو کے مشین~$M$ دا انپٹ~$w$ اُتے چلن بیان
کردے نیں۔ کجھ !!{sentence}s، $\Obj 0$، $\prime$ تے~$<$
دیاں شرطاں مقرر کرن گے؛ کجھ !!{sentence}s ابتدائی مکمل
حالت بیان کرن گے؛ تے کجھ !!{sentence}s بیان کرن گے کہ
مشین~$M$ کوئی خاص حکم چلاؤن توں بعد کس مکمل حالت وچ اے۔
\end{explain}

\begin{defn}
  \ollabel{defn:tm-descr}
جے ٹیورنگ مشین $M = \tuple{Q, \Sigma, q_0, \delta}$ دتی
ہووے، تاں زبان~$\Lang L_M$ وچ ایہ چیزاں ہوندیاں نیں:
\begin{enumerate}
\item اک دو-جگہی !!{predicate} $\Obj Q_q(x, y)$ ہر حالت~$q \in
  Q$ لئی۔ مطلب ایہ
  اے کہ $\Obj Q_q(\num{m}, \num{n})$ آکھدا اے: “$n$
  قدم چُکن توں بعد~$M$ حالت~$q$ وچ اے تے~$m$ویں خانہ
  پڑھ رہی اے۔”
\item اک دو-جگہی !!{predicate} $\Obj S_\sigma(x, y)$ ہر
  نشان~$\sigma\in \Sigma$ لئی۔ مطلب ایہ اے کہ $\Obj S_\sigma(\num{m},
  \num{n})$ آکھدا اے: “$n$ قدم چُکن توں بعد~$m$ویں خانے
  وچ نشان~$\sigma$ اے۔”
\item اک !!{constant} $\Obj 0$
\item اک یک-مقامی !!{function} $\prime$
\item اک دو-جگہی !!{predicate} $<$
\end{enumerate}
\end{defn}

ٹیورنگ مشین~$M$ دا انپٹ $w = \sigma_{i_1}\dots\sigma_{i_k}$
اُتے چلن بیان کرن والے !!{sentence}s ایہ نیں:
\begin{enumerate}
\item عدداں تے~$<$ بارے مسلّمے:
\begin{enumerate}
\item اک !!^a{sentence} جیڑا آکھدا اے کہ ہر عدد اپنے جانشین توں
  گھٹ اے:
\[
\lforall[x][x < x']
\]
\item اک !!^a{sentence} جیڑا یقینی بناندا اے کہ~$<$ متعدی اے:
\[
\lforall[x][\lforall[y][\lforall[z][
      ((x < y \land y < z) \lif x < z)]]]
\]
\end{enumerate}
\item انپٹ دی ابتدائی مکمل حالت بیان کرن والے مسلّمے:
\begin{enumerate}
\item $0$ قدم چُکن توں بعد، یعنی مشین چلن توں پہلاں،
  $M$ ابتدائی حالت~$q_0$ وچ اے تے خانہ~$1$ پڑھ رہی اے:
\[
\Obj Q_{q_0}(\num{1}, \num{0})
\]
\item پہلے $k+1$ خانیاں وچ ترتیب وار نشان $\TMendtape$،
  $\sigma_{i_1}$، \dots، $\sigma_{i_k}$ نیں:
\[
\Obj S_\TMendtape(\num{0}, \num{0}) \land
\Obj S_{\sigma_{i_1}}(\num{1}, \num{0}) \land
\dots \land
\Obj S_{\sigma_{i_k}}(\num{k}, \num{0})
\]
\item باقی پٹی خالی اے:
\[
\lforall[x][(\num{k} < x \lif \Obj S_\TMblank(x, \num{0}))]
\]
\end{enumerate}
\item اک مکمل حالت توں اگلی مکمل حالت وچ تبدیلی بیان کرن
  والے مسلّمے:

اگلے حصے لئی~$!A(x, y)$ نوں اوہ عطف منّو جیہڑا ایس روپ
والے سارے !!{sentence}s جوڑ کے بندا اے:
\[
\lforall[z][
  (((z < x \lor x < z) \land \Obj S_\sigma(z, y))
  \lif \Obj S_\sigma(z, y'))]
\]
جتھے~$\sigma \in \Sigma$ اے۔ اسی~$!A(\num{m},\num{n})$
ایہ آکھن لئی ورتدے آں: “خانہ~$m$ توں علاوہ پٹی
$n+1$ قدم چُکن توں بعد اوہو جیہی اے جیہی~$n$ قدم
چُکن توں بعد سی۔”
\begin{enumerate}
\item \ollabel{rep-right} ہر حکم $\delta(q_i, \sigma) =
  \tuple{q_j, \sigma', \TMright}$ لئی ایہ !!{sentence}:
\begin{align*}
& \lforall[x][\lforall[y][(
   (\Obj Q_{q_i}(x, y) \land \Obj S_{\sigma}(x, y)) \lif {}]] \\
&\qquad   (\Obj Q_{q_j}(x', y') \land \Obj S_{\sigma'}(x, y') \land
\formula A(x, y)))
\end{align*}
ایہ آکھدا اے: جے~$y$ قدم چُکن توں بعد مشین حالت~$q_i$
وچ خانے~$x$ نوں پڑھ رہی اے، جتھے نشان~$\sigma$ اے، تاں
$y+1$ قدم چُکن توں بعد اوہ خانہ~$x+1$ پڑھ رہی اے،
حالت~$q_j$ وچ اے، خانے~$x$ وچ ہن~$\sigma'$ اے، تے
خانہ~$x$ توں علاوہ ہر خانے وچ اوہو نشان اے جیڑا
$y$ قدم چُکن توں بعد سی۔

\item \ollabel{rep-left} ہر حکم $\delta(q_i, \sigma) =
  \tuple{q_j, \sigma', \TMleft}$ لئی ایہ !!{sentence}:
% PNB-CMP-030: the old scanned/written square is x', so the frame condition exempts x', not x.
\begin{align*}
& \lforall[x][\lforall[y][
    ((\Obj Q_{q_i}(x', y) \land \Obj S_{\sigma}(x', y)) \lif {}]]\\
& \qquad   (\Obj Q_{q_j}(x, y') \land \Obj S_{\sigma'}(x', y') \land
\formula A(x', y))) \land {}\\
& \lforall[y][((\Obj Q_{q_i}(\num{0}, y) \land \Obj S_{\sigma}(\num{0},
    y)) \lif {}]\\
& \qquad (\Obj Q_{q_j}(\num{0}, y') \land \Obj S_{\sigma'}(\num{0},
  y') \land !A(\num{0}, y)))
\end{align*}
ذرا سوچو ایہ کم کیویں کردا اے: ہن اسی “جے مشین خانہ~$x$
پڑھ رہی ہووے \dots” توں نہیں، بلکہ “جے مشین خانہ~$x+1$
پڑھ رہی ہووے \dots” توں شروع کردے آں۔ کھبے ول جان دا
مطلب اے کہ اگلے قدم وچ مشین خانہ~$x$ پڑھدی اے؛ پر
نشان خانہ~$x+1$ وچ لکھدی اے۔ اسی ایویں ایس لئی کردے
آں کہ ساڈے کول تفریق یا پیشرو فنکشن دا نشان نہیں۔

دھیان دیو کہ $x+1$ روپ دے عدد $1$، $2$، \dots{} نیں؛
ایس وچ اوہ صورت شامل نہیں جدوں مشین خانہ~$0$ پڑھ رہی
ہووے تے اوہنوں کھبے ول جانا ہووے (ایہ کر نہیں سکدی،
ایس لئی اپنی تھاں رہندی اے)۔ اوہ خاص صورت عطف دے دوجے
جزو نال بیان ہوئی اے: جے~$y$ قدم چُکن توں بعد مشین
خانہ~$0$ حالت~$q_i$ وچ پڑھ رہی اے تے خانہ~$0$
وچ نشان~$\sigma$ اے، تاں~$y+1$ قدم چُکن توں بعد
اوہ ہن وی خانہ~$0$ پڑھ رہی اے، ہن حالت~$q_j$
وچ اے، خانے~$0$ دا نشان~$\sigma'$ اے، تے
خانہ~$0$ توں علاوہ خانے اوہو نشان رکھدے نیں
جیہڑے~$y$ قدم چُکن توں بعد اوہناں وچ سن۔
\item \ollabel{rep-stay} ہر حکم $\delta(q_i, \sigma) =
  \tuple{q_j, \sigma', \TMstay}$ لئی ایہ !!{sentence}:
\begin{align*}
& \lforall[x][\lforall[y][(
   (\Obj Q_{q_i}(x, y) \land \Obj S_{\sigma}(x, y)) \lif {}]] \\
&\qquad   (\Obj Q_{q_j}(x, y') \land \Obj S_{\sigma'}(x, y') \land
\formula A(x, y)))
\end{align*}
\end{enumerate}
\end{enumerate}
اُتلے سارے !!{sentence}s جوڑ کے بنے عطف نوں~$!T(M, w)$
آکھو، جیہڑا ٹیورنگ مشین~$M$ تے انپٹ~$w$ لئی اے۔

ایہ ظاہر کرن لئی کہ~$M$ آخرکار رُکدی اے، سانوں
اک ایہو جیہا !!{sentence} لبھنا پئے گا جیڑا آکھے: “کجھ
قدم چُکن توں بعد حالتی تبدیلی دا فنکشن غیر تعریف
شدہ ہووے گا۔” $X$~اوہناں ساریاں جوڑیاں
$\tuple{q, \sigma}$ دا سیٹ ہووے جِنہاں لئی
~$\delta(q, \sigma)$ غیر تعریف شدہ اے۔ فیر
$!E(M, w)$ ایہ !!{sentence} ہووے:
\[
\lexists[x][\lexists[y][(\bigvee_{\tuple{q, \sigma} \in
      X}(\Obj Q_q(x, y) \land \Obj S_\sigma(x, y)))]]
\]

جے اسی مقررہ رُکّن والی حالت~$h$ رکھنے والی ٹیورنگ
مشین ورتیے، تاں ایہ ہور وی سوکھا اے: فیر
!!{sentence}~$!E(M, w)$
\[
\lexists[x][\lexists[y][\Obj Q_h(x, y)]]
\]
ایہ ظاہر کردا اے کہ مشین آخرکار رُکدی اے۔

\begin{prop}
\ollabel{prop:mlessk}
جے $m < k$ ہووے، تاں $!T(M, w) \Entails \num{m} < \num{k}$۔
\end{prop}

\begin{proof}
مشق۔
\end{proof}

\begin{prob}
\olref[tur][und][rep]{prop:mlessk} ثابت کرو۔
(اشارہ: $k-m$ اُتے استقرا ورتو۔)
\end{prob}
\endgroup
% /OLP-0270

% OLP-0271 content/turing-machines/undecidability/verification.tex
\begingroup
\def\olfilename{verification}
\expandafter\def\csname ol@sectioncs\endcsname{section}

\olfileid{tur}{und}{ver} 
\olsection{نمائندگی دی تصدیق}

\begin{explain}
ساڈی نمائندگی دے درست کم کرن دی تصدیق لئی سانوں دو
گلاں ثابت کرنیاں پین گیاں۔ پہلے، ایہ دکھانا پئے گا
کہ جے~$M$ انپٹ~$w$ اُتے رُکدی اے، تاں
$!T(M, w) \lif !E(M, w)$ درست اے۔ فیر ایہدا اُلٹ
دکھانا پئے گا: جے~$!T(M, w) \lif !E(M, w)$ درست
اے، تاں~$M$ انپٹ~$w$ اُتے واقعی آخرکار
رُکدی اے۔

ایہ دو نتیجے ثابت کرن دیاں تدبیراں بہت وکھریاں نیں۔
پہلے نتیجے لئی سانوں دکھانا اے کہ پہلے درجے دی منطق
دا اک !!a{sentence}، یعنی~$!T(M, w) \lif !E(M, w)$،
درست اے۔ ایہ دکھان دا سب توں سوکھا طریقہ اک
!!a{derivation} دینا اے۔ پر ساڈا ثبوت ہر~$M$ تے~$w$
لئی ٹھیک ہووے؛ ایس لئی ایتھے لامتناہی جملے نیں۔ جیہڑا وی
!!{sentence} ہووے، اوہدا اک !!a{derivation} دینا پئے گا۔
ایس لئی اسی استقرا نال ایہ
ثابت کر سکدے آں کہ~$M$ تے~$w$ جیہو جیہے وی ہون،
تے~$M$ نوں انپٹ~$w$ اُتے رُکّن لئی جنے وی قدم
لگّن، $!T(M, w) \lif !E(M, w)$ دا اک
!!a{derivation} موجود ہووے گا۔

فطری طور اُتے ساڈا استقرا اُنہاں قدماں دی گنتی اُتے
ہووے گا جیہڑے~$M$ رُکّن والی مکمل حالت تک پہنچن
توں پہلاں چُکدی اے۔ استقرائی ثبوت وچ اسی قائم
کراں گے کہ ہر قدم~$n$ لئی، جیہڑا~$M$ انپٹ~$w$
اُتے چلن وچ چُکدی اے، $!T(M, w) \Entails !C(M, w, n)$،
جتھے~$!C(M, w, n)$، $M$ دے انپٹ~$w$ اُتے
چلدیاں~$n$ قدم چُکن توں بعد دی مکمل حالت ٹھیک
بیان کردا اے۔ ہن جے~$M$ انپٹ~$w$ اُتے، مثلاً،
$n$ قدم چُکن توں بعد رُکدی اے، تاں~$!C(M, w, n)$
رُکّن والی مکمل حالت بیان کرے گا۔ اسی ایہ وی
دکھاواں گے کہ~$!C(M, w, n)
\Entails !E(M, w)$ جدوں وی~$!C(M, w, n)$ ایسی رُکّن
والی مکمل حالت بیان کردا اے۔ ایس لئی جے~$M$ انپٹ~$w$
اُتے رُکدی اے، تاں کسے~$n$ لئی~$M$، $n$ قدم
چُکن توں بعد رُکّن والی مکمل حالت وچ ہووے گی۔
فیر~$!T(M, w) \Entails !C(M, w, n)$، جتھے
$!C(M, w, n)$ رُکّن والی مکمل حالت بیان کردا اے؛
تے ایس صورت وچ~$!C(M, w, n) \Entails !E(M, w)$
ہون کرکے سانوں~$T(M, w) \Entails !E(M, w)$ ملدا
اے، یعنی~$\Entails !T(M, w) \lif !E(M, w)$۔

اُلٹے رخ دی تدبیر بالکل وکھری اے۔ ایتھے اسی
$\Entails !T(M, w) \lif !E(M, w)$ فرض کرکے
ثابت کردے آں کہ~$M$ انپٹ~$w$ اُتے رُکدی اے۔
مفروضے توں~$!T(M, w) \Entails !E(M,
w)$ ملدا اے، یعنی ہر اُس !!{structure} وچ~$!E(M,
w)$ سچ اے جتھے~$!T(M, w)$ سچ اے۔ ایس
لئی اسی اک !!a{structure}~$\Struct{M}$ بیان کراں
گے جتھے~$!T(M, w)$ سچ اے: اوہدا دائرہ~$\Nat$
ہووے گا، تے سارے~$\Obj Q_q$ تے~$\Obj S_\sigma$
دیاں تعبیراں~$M$ دے انپٹ~$w$ اُتے چلن دیاں
مکمل حالتاں توں ملن گیاں۔ مثال لئی،
$\Sat{M}{\Obj Q_q(\num{m}, \num{n})}$ اُدوں تے صرف اُدوں
% PNB-CMP-031: source names T here; the running machine is M.
جدوں~$M$ انپٹ~$w$ اُتے~$n$ قدم چلن توں بعد
حالت~$q$ وچ اے تے خانہ~$m$ پڑھ رہی اے۔
ہن چونکہ مفروضے مطابق~$!T(M, w) \Entails !E(M, w)$،
تے تعمیر مطابق~$\Sat{M}{!T(M,
  w)}$، ایس لئی~$\Sat{M}{!E(M, w)}$۔ پر
$\Sat{M}{!E(M, w)}$ اُدوں تے صرف اُدوں جدوں کوئی
$n \in \Domain{M} = \Nat$ ایہو جیہا ہووے کہ~$M$
انپٹ~$w$ اُتے~$n$ قدم چُکن توں بعد رُکّن
والی مکمل حالت وچ ہووے۔
\end{explain}

\begin{defn} 
$!C(M, w, n)$ ایہ !!{sentence} ہووے:
\[ 
\Obj Q_q(\num{m}, \num{n}) \land \Obj S_{\sigma_0}(\num{0}, \num{n})
\land \dots \land \Obj S_{\sigma_k}(\num{k}, \num{n}) \land
\lforall[x][(\num{k} < x \lif \Obj S_\TMblank(x, \num{n}))]
\] 
جتھے~$q$، $M$ دی ویلے~$n$ والی حالت اے،
$M$ خانہ~$m$ ویلے~$n$ اُتے پڑھ رہی اے، خانے~$i$
وچ نشان~$\sigma_i$ ویلے~$n$ اُتے اے، $0
\le i \le k$ لئی؛ تے~$k$ یا تاں ویلے~$0$
اُتے پٹی دے سب توں سجّے غیر خالی خانے دا نمبر
اے یا~$n$ قدماں تک سر جیہڑے خانیاں تک پہنچ چکیا
اے، اوہناں وچ سب توں سجّے خانے دا نمبر، جیہڑا وی وڈا ہووے۔
\end{defn}

\begin{lem}\ollabel{lem:halt-config-implies-halt}
جے~$M$ انپٹ~$w$ اُتے~$n$ قدم چُکن
توں بعد رُکّن والی مکمل حالت وچ اے، تاں
$!C(M, w, n) \Entails !E(M, w)$۔
\end{lem}

\begin{proof}
فرض کرو~$M$ انپٹ~$w$ اُتے~$n$ قدم چُکن توں
بعد رُکدی اے۔ کوئی حالت~$q$، خانہ~$m$ تے نشان
~$\sigma$ ایداں دے نیں کہ:
\begin{enumerate} 
\item $n$ قدم چُکن توں بعد~$M$ حالت~$q$ وچ
  خانہ~$m$ پڑھ رہی اے، جتھے~$\sigma$ لکھیا اے۔
\item حالتی تبدیلی دا فنکشن~$\delta(q, \sigma)$
  غیر تعریف شدہ اے۔
\end{enumerate}
$!C(M, w, n)$ ایس مکمل حالت دا بیان اے تے ایس
وچ اجزا~$\Obj Q_{q}(\num{m}, \num{n})$ تے~$\Obj
S_{\sigma}(\num{m}, \num{n})$ شامل ہون گے۔ ایہ
اجزا مل کے~$!E(M,
w)$ لازم کردے نیں:
\[
\lexists[x][\lexists[y][(\bigvee_{\tuple{q, \sigma} \in
      X}(\Obj Q_q(x, y) \land \Obj S_\sigma(x, y)))]]
\]
% PNB-CMP-032: use the stated halting pair and the S predicate symbol.
کیوں جے~$\Obj Q_{q}(\num{m}, \num{n}) \land \Obj S_{\sigma}(\num{m},
\num{n}) \Entails \bigvee_{\tuple{q, \sigma} \in X} (\Obj Q_q(\num{m},
\num{n}) \land \Obj S_{\sigma}(\num{m}, \num{n}))$،
ایس لئی کہ~$\tuple{q,\sigma} \in X$ اے۔
\end{proof}

\begin{explain}
ایس لئی جے~$M$ انپٹ $w$ اُتے رُکدی اے، تاں
کوئی~$n$ ایہو جیہا اے کہ~$!C(M,
w, n) \Entails !E(M,w)$۔ ہن اسی دکھاواں گے
کہ ہر اُس ویلے~$n$ لئی، جتھے مشین پہلاں نہ رُکی
ہووے، $!T(M, w) \Entails !C(M, w, n)$۔
\end{explain}

\begin{lem}
\ollabel{lem:config}
% PNB-CMP-033: include the configuration at the first halting step.
ہر~$n$ لئی، جے~$M$، $n$ قدم چُکن توں پہلاں
نہ رُکی ہووے، تاں~$!T(M, w)
\Entails !C(M, w, n)$۔
\end{lem}

\begin{proof}
استقرا دی بنیاد: جے~$n = 0$، تاں~$!C(M, w, 0)$
دے عطف دے سارے جزو~$!T(M, w)$ دے عطف وچ وی نیں۔ ایس
لئی اوہ ایس توں منطقی طور اُتے نکلدے نیں۔

استقرائی مفروضہ: جے~$M$، $n$ویں قدم توں پہلاں
نہ رُکی ہووے، تاں~$!T(M,w) \Entails !C(M, w, n)$۔
سانوں دکھانا اے کہ (جے~$!C(M, w, n)$ رُکّن والی
مکمل حالت بیان نہ کردا ہووے) تاں~$!T(M, w) \Entails
!C(M, w, n+1)$۔

% PNB-CMP-034: induction step also needs n=0.
فرض کرو~$n \ge 0$، تے~$n$ قدم چُکن توں بعد
$M$، انپٹ~$w$ اُتے شروع ہو کے، حالت~$q$ وچ
خانہ~$m$ پڑھ رہی اے۔ چونکہ~$M$، $n$ قدم چُکن
توں بعد نہیں رُکدی، ہِٹھلیاں تِن صورتاں وچوں
کسے اک صورت دا حکم~$M$ دے
پروگرام وچ موجود ہووے گا:

\begin{enumerate} 
\item \ollabel{right} $\delta(q, \sigma) = \tuple{q', \sigma', \TMright}$

\item \ollabel{left} $\delta(q, \sigma) = \tuple{q', \sigma', \TMleft}$

\item \ollabel{stay} $\delta(q, \sigma) = \tuple{q', \sigma', \TMstay}$
\end{enumerate}

اسی ایہناں تِن صورتاں نوں واری واری ویکھاں گے۔

\begin{enumerate} 
\item فرض کرو~\olref{right} والے روپ دا حکم موجود اے۔
  \olref[rep]{defn:tm-descr}\olref[rep]{rep-right} مطابق ایہدے
  معنی ایہ نیں کہ
% PNB-CMP-035: omit the old scanned square m from retained-symbol conjuncts.
\begin{align*} 
& \lforall[x][\lforall[y][((\Obj Q_{q}(x,
  y) \land \Obj S_\sigma(x, y)) \lif {}]]\\
& \qquad (\Obj Q_{q'}(x',
  y') \land \Obj S_{\sigma'}(x, y') \land !A(x, y)))  
\intertext{، $!T(M,w)$ دا اک جزو اے۔ ایس توں ہِٹھلا
  !!{sentence} نکلدا اے (کُلی تخصیص نال~$\num{m}$ نوں
  $x$ دی تھاں تے~$\num{n}$ نوں~$y$ دی تھاں رکھ کے):}
& (\Obj Q_{q}(\num{m}, \num{n}) \land \Obj S_{\sigma}(\num{m},
\num{n})) \lif {}\\
& \qquad (\Obj Q_{q'}(\num{m}', \num{n}') \land
\Obj S_{\sigma'}(\num{m}, \num{n}') \land !A(\num{m}, \num{n})).
\intertext{استقرائی مفروضے مطابق~$!T(M, w) \Entails !C(M, w, n)$،
  یعنی:}
& \Obj Q_q(\num{m}, \num{n}) \land \Obj S_{\sigma_0}(\num{0}, \num{n})
\land \dots \land \Obj S_{\sigma_k}(\num{k}, \num{n}) \land \\
& \qquad\lforall[x][(\num{k} < x \lif \Obj S_\TMblank(x, \num{n}))]\\
\intertext{چونکہ~$n$ قدم چُکن توں بعد خانہ~$m$ وچ~$\sigma$
  اے، ایس لئی متعلقہ جزو~$\Obj S_\sigma(\num{m}, \num{n})$
  اے؛ لہٰذا ایہ لازم آندا اے:}
& \Obj Q_{q}(\num{m}, \num{n}) \land \Obj S_{\sigma}(\num{m},
   \num{n})
\intertext{ہن سانوں ایہ ملدا اے:}
& \Obj Q_{q'}(\num{m}', \num{n}') \land \Obj S_{\sigma'}(\num{m},
  \num{n}') \land {}\\
& \qquad\bigwedge_{\substack{0\le i\le k\\ i\ne m}}
  \Obj S_{\sigma_i}(\num{i}, \num{n}') \land {}\\
& \qquad \lforall[x][(\num{k} < x \lif \Obj S_\TMblank(x, \num{n}'))]
\end{align*}
ایہ نتیجہ ایویں ملدا اے: پہلی لکیر پچھلے شرطیہ دے
تالی توں اثباتِ مقدم نال سدھی نکلدی اے۔ وچلی
لکیر وچ~$S_{\sigma_m}(\num{m},\num{n}')$ شامل نہیں۔
ایس لکیر دا ہر جزو $!C(M, w, n)$ دے متعلقہ جزو تے
$!A(\num{m}, \num{n})$ توں نکلدا اے۔

جے~$m < k$، تاں~$!T(M,w) \Proves \num{m} < \num {k}$
(\olref[rep]{prop:mlessk})، تے~$<$ دی تعدیت توں
سانوں~$\lforall[x][(\num{k} < x \lif \num{m} < x)]$
ملدا اے۔ جے~$m = k$، تاں ایہو
$\lforall[x][(\num{k} < x \lif \num{m} < x)]$
صرف منطق نال نکلدا اے۔ آخری لکیر فیر
$!C(M, w,
n)$ دے متعلقہ جزو، $\lforall[x][(\num{k} < x \lif \num{m} < x)]$،
تے~$!A(\num{m},
\num{n})$ توں نکلدی اے۔ جے~$m<k$، تاں ایہ
پہلاں ہی~$!C(M, w, n+1)$ اے۔

ہن فرض کرو~$m=k$۔ ایس صورت وچ~$n+1$ قدم
چُکن توں بعد پٹی دا سر خانہ~$k+1$ تک وی پہنچ چکیا
اے۔ ہن سر جیہڑے خانیاں تک پہنچیا اے، اوہناں وچوں ایہ سب توں سجّا اے۔
ایس لئی~$!C(M, w, n+1)$ وچ اک نواں جزو
$\Obj S_\TMblank(\num{k}',\num{n}')$ اے، تے آخری جزو
$\lforall[x][(\num{k}' < x \lif \Obj S_\TMblank(x, \num{n}'))]$
اے۔ سانوں جانچنا اے کہ ایہ دونیں !!{sentence}s
وی منطقی طور اُتے نکلدے نیں۔

ساڈے کول پہلے ہی
$\lforall[x][(\num{k} < x \lif \Obj S_\TMblank(x,
  \num{n}'))]$ اے۔ خاص طور اُتے ایہدے توں
$\num{k} < \num{k}' \lif
\Obj S_\TMblank(\num{k}', \num{n}')$ ملدا اے۔ مسلّمے
$\lforall[x][x <
  x']$ توں~$\num{k} < \num{k}'$ ملدا اے۔ اثباتِ
مقدم نال~$\Obj
S_\TMblank(\num{k}',\num{n}')$ نکلدا اے۔

نالے، چونکہ~$!T(M,w) \Proves \num{k} < \num{k}'$،
تاں~$<$ دی تعدیت دا مسلّمہ سانوں
$\lforall[x][(\num{k}' < x \lif \Obj
  S_\TMblank(x, \num{n}'))]$ دیندا اے۔
(ایس دی جانچ اسی مشق لئی چھڈدے آں۔)

\item فرض کرو~\olref{left} والے روپ دا حکم موجود اے۔
  فیر~\olref[rep]{defn:tm-descr}\olref[rep]{rep-left} مطابق:
% PNB-CMP-035: omit the written square from the left-case retained-symbol conjunction.
% PNB-CMP-036: left-frame exemption is the written x-prime; edge states match q and q-prime.
\begin{align*} 
& \lforall[x][\lforall[y][((\Obj Q_{q}(x', y) \land \Obj
    S_{\sigma}(x', y)) \lif {}]]\\
& \qquad (\Obj Q_{q'}(x, y') \land \Obj
  S_{\sigma'}(x', y') \land !A(x', y))) \land {}\\
& \lforall[y][((\Obj Q_{q}(\num{0}, y) \land \Obj S_{\sigma}(\num{0},
    y)) \lif {}]\\
& \qquad (\Obj Q_{q'}(\num{0}, y') \land \Obj S_{\sigma'}(\num{0}, y')
  \land !A(\num{0}, y)))
\intertext{، $!T(M,w)$ دا اک جزو اے۔ جے~$m>0$،
  تاں~$l = m - 1$ رکھو (یعنی~$m = l+1$)۔ اُتلے
  !!{sentence} دے پہلے جزو توں ایہ نتیجہ نکلدا اے:}
& (\Obj Q_{q}(\num{l}', \num{n}) \land \Obj S_{\sigma}(\num{l}', \num{n}))
\lif {} \\
& \qquad (\Obj Q_{q'}(\num{l}, \num{n}') \land \Obj S_{\sigma'}(\num{l}',
\num{n}') \land !A(\num{l}', \num{n}))
\intertext{دوجی صورت وچ~$l = m = 0$ رکھو تے دوجے
  جزو توں نکلن والا ہِٹھلا !!{sentence} ویکھو:}
& ((\Obj Q_{q}(\num{0}, \num{n}) \land \Obj S_{\sigma}(\num{0}, \num{n})) \lif {}\\
& \qquad   (\Obj Q_{q'}(\num{0}, \num{n}') \land \Obj S_{\sigma'}(\num{0}, \num{n}') \land
\formula A(\num{0}, \num{n}))) 
\intertext{دوہاں وچوں ہر اک جملے نال ایہ نتیجہ نکلدا اے:}
&  \Obj Q_{q'}(\num{l}, \num{n}') \land \Obj S_{\sigma'}(\num{m},
  \num{n}') \land {}\\
&\qquad  \bigwedge_{\substack{0\le i\le k\\ i\ne m}}
  \Obj S_{\sigma_i}(\num{i}, \num{n}') \land {} \\
& \qquad  \lforall[x][(\num{k} < x
  \lif \Obj S_\TMblank(x, \num{n}'))]
\end{align*}
پہلاں وانگ۔ (دھیان دیو کہ پہلی صورت وچ
$\num{l}' \ident \num{l+1}
\ident \num{m}$ تے دوجی صورت وچ
$\num{l} \ident \num{0}$۔) پر ایہ عین
$!C(M, w, n+1)$ اے۔

\item صورت~\olref{stay} مشق لئی چھڈدے آں۔
\end{enumerate}
اسی دکھا دتا کہ ہر اُس~$n$ لئی، جیہدے توں پہلاں
مشین نہ رُکی ہووے، $!T(M, w) \Entails !C(M, w, n)$۔
\end{proof}

\begin{prob}
\olref[tur][und][ver]{lem:config} دے ثبوت وچ
صورت~\olref[tur][und][ver]{stay} مکمل کرو۔
\end{prob}

\begin{prob}
$\Obj S_{\sigma_i}(\num{i}, \num{n}')$ دا اک
!!a{derivation}، $\Obj S_{\sigma_i}(\num{i}, \num{n})$
تے~$!A(m, n)$ توں دیو (ایس مفروضے ہِٹھ کہ~$i \neq
m$، یعنی یا~$i < m$ یا~$m < i$)۔
\end{prob}

\begin{prob}
$\lforall[x][(\num{k}' < x \lif \Obj
  S_\TMblank(x, \num{n}'))]$ دا اک !!a{derivation}
$\lforall[x][(\num{k} < x \lif \Obj
  S_\TMblank(x, \num{n}'))]$، $\lforall[x][x < x']$ تے
$\lforall[x][\lforall[y][\lforall[z][ ((x < y \land y < z) \lif x <
      z)]]]$ توں دیو۔)
\end{prob}


\begin{lem}
\ollabel{lem:valid-if-halt}
جے~$M$ انپٹ~$w$ اُتے رُکدی اے، تاں
$!T(M, w) \lif
!E(M, w)$ درست اے۔
\end{lem}

\begin{proof}
\olref{lem:config} توں سانوں پتا اے کہ رُکّن توں
پہلاں یا رُکّن دے ویلے~$n$
اُتے مکمل حالت دا بیان~$!C(M, w, n)$، جیہڑا~$M$
دی ویلے~$n$ اُتے مکمل حالت دسدے، $!T(M, w)$
توں نکلدا اے۔ فرض کرو~$M$، $k$ قدم چُکن توں
بعد رُکدی اے۔ اوس ویلے اوہ خانہ~$m$ پڑھ رہی
ہووے گی، کسے~$m \in \Nat$ لئی۔ فیر
$!C(M, w, k)$، $M$ دی رُکّن والی مکمل حالت
بیان کردا اے؛ یعنی ایس وچ~$\Obj Q_q(\num{m}, \num{k})$
تے~$\Obj
S_\sigma(\num{m}, \num{k})$ دونیں جزو شامل نیں،
جدوکہ~$\delta(q,\sigma)$ غیر تعریف شدہ اے۔
ایس لئی~\olref{lem:halt-config-implies-halt} مطابق
$!C(M, w, k) \Entails !E(M,
w)$۔ پر چونکہ~$!T(M, w) \Entails !C(M, w, k)$،
ایس لئی~$!T(M, w)
\Entails !E(M, w)$، تے نتیجتاً
$!T(M, w) \lif !E(M, w)$ درست اے۔
\end{proof}


\begin{explain} 
اپنے دعوے دی تصدیق مکمل کرن لئی سانوں اُلٹا
رخ وی قائم کرنا پئے گا: جے~$!T(M, w) \lif !E(M, w)$
درست اے، تاں~$M$ انپٹ~$w$ اُتے واقعی
رُکدی اے۔
\end{explain}

\begin{lem}
\ollabel{lem:halt-if-valid}
جے~$\Entails !T(M, w) \lif !E(M, w)$، تاں~$M$
انپٹ~$w$ اُتے رُکدی اے۔
\end{lem}

\begin{proof}
$\Lang L_M$ دی !!{structure}~$\Struct M$ ویکھو
جیہدا دائرہ~$\Nat$ اے۔ ایہ~$\Obj 0$ دی تعبیر~$0$،
$\prime$ دی تعبیر جانشین فنکشن تے~$<$ دی تعبیر “توں گھٹ”
والے تعلق دے طور اُتے کردی اے؛ تے محمولاں~$\Obj Q_q$
تے~$\Obj S_\sigma$ دی ایویں:
\begin{align*}
  \Assign{\Obj Q_q}{M} & =
\Setabs{\tuple{m, n}}{\begin{array}{ll}\text{انپٹ~$w$ اُتے شروع ہو کے، $n$ قدم بعد،}\\ \text{$M$ حالت $q$
  وچ خانہ~$m$ پڑھ رہی اے}\end{array}} \\
\Assign{\Obj S_\sigma}{M} & = \Setabs{\tuple{m, n}}{\begin{array}{ll}
\text{انپٹ~$w$ اُتے شروع ہو کے، $n$ قدم بعد،}\\ \text{خانہ~$m$، $M$ دا، وچ
  نشان~$\sigma$ اے}\end{array}}
\end{align*}
ہور لفظاں وچ، اسی !!{structure}~$\Struct{M}$
ایویں بناندے آں کہ اوہ~$M$ دا انپٹ~$w$
اُتے حقیقی چلن قدم بہ قدم بیان کرے۔
% PNB-CMP-037: extend a finite run by retaining its final halting configuration.
جے مشین رُک جائے تاں بعدلے ویلّیاں لئی آخری
رُکّن والی مکمل حالت اوہو جیہی رکھو؛ ایویں
تعبیراں سارے قدرتی عددی ویلّیاں اُتے مقرر ہو جاندیاں نیں۔
صاف اے کہ~$\Sat{M}{!T(M, w)}$۔ جے
$\Entails !T(M, w) \lif !E(M, w)$،
تاں~$\Sat{M}{!E(M, w)}$ وی، یعنی:
\[
\Sat{M}{\lexists[x][\lexists[y][(\bigvee_{\tuple{q, \sigma} \in
      X}(\Obj Q_q(x, y) \land \Obj S_\sigma(x, y)))]]}.
\]
چونکہ~$\Domain{M} = \Nat$، ایس لئی کوئی
$m$، $n \in \Nat$ ضرور ہون گے جنہاں لئی
$\Sat{M}{\Obj Q_q(\num{m}, \num{n}) \land \Obj S_\sigma(\num{m},
\num{n})}$، کسے~$q$ تے~$\sigma$ لئی، جتھے~$\delta(q, \sigma)$ غیر تعریف
شدہ اے۔ $\Struct M$ دی تعریف مطابق ایہدے معنی
ایہ نیں کہ~$M$، انپٹ~$w$ اُتے شروع ہو کے، $n$
قدم چُکن توں بعد حالت~$q$ وچ اے تے نشان~$\sigma$
پڑھ رہی اے، تے حالتی تبدیلی دا فنکشن غیر تعریف
شدہ اے؛ یعنی~$M$ رُک چکی اے۔
\end{proof}
\endgroup
% /OLP-0271

% OLP-0272 content/turing-machines/undecidability/unsolvability-decision-problem.tex
\begingroup
\def\olfilename{unsolvability-decision-problem}
\expandafter\def\csname ol@sectioncs\endcsname{section}

\olfileid{tur}{und}{uns} 
\olsection{فیصلے دا مسئلہ حل نہیں ہو سکدا}

\begin{thm}
\ollabel{thm:decision-prob}
فیصلے دا مسئلہ حل نہیں ہو سکدا۔ کوئی ایہو جیہی ٹیورنگ
مشین~$D$ نہیں کہ جدوں پٹی اُتے پہلے درجے دی منطق
دا اک !!a{sentence}~$!B$ انپٹ وجوں رکھ کے اوہنوں
چلایا جاوے، تاں~$D$ آخرکار رُکے تے~$1$ آؤٹ پٹ اُدوں تے صرف
اُدوں دے جدوں~$!B$ درست ہووے، تے دوجی صورت وچ~$0$ دے۔
\end{thm}

\begin{proof}
فرض کرو فیصلے دا مسئلہ حل ہو سکدا ہووے، یعنی
ایہو جیہی ٹیورنگ مشین~$D$ موجود ہووے۔ فیر اسی
رُکّن دا مسئلہ ایویں حل کر سکدے آں۔
اسی اک ٹیورنگ مشین~$E$ بناندے آں جیہڑی
ٹیورنگ مشین دا اشاریہ~$e$، یعنی~$M_e$ دا نمبر، تے انپٹ~$w$
ملن اُتے متعلقہ !!{sentence}~
$!T(M_e, w) \lif !E(M_e, w)$ حساب کردی
اے تے پٹی دا سب توں کھبّا خانہ پڑھدی ہوئی
رُک جاندی اے۔ فیر مشین~$E \concat D$،
انپٹ~$e$ تے~$w$ ملن اُتے، پہلے~
$!T(M_e, w) \lif
!E(M_e, w)$ حساب کرے گی
تے فیر اوس انپٹ اُتے فیصلے دے مسئلے والی
مشین~$D$ چلائے گی۔ $D$ آؤٹ پٹ~$1$ نال
اُدوں تے صرف اُدوں رُکدی اے جدوں
$!T(M_e, w) \lif !E(M_e, w)$ درست ہووے،
ورنہ~$0$ دیندی اے۔
\olref[ver]{lem:halt-if-valid} تے
\olref[ver]{lem:valid-if-halt} مطابق
$!T(M_e, w) \lif !E(M_e, w)$ اُدوں تے صرف اُدوں
درست اے جدوں~$M_e$ انپٹ~$w$ اُتے رُکدی اے۔
ایس لئی~$E\concat D$، انپٹ~$e$ تے~$w$
ملن اُتے، آؤٹ پٹ~$1$ نال اُدوں تے صرف اُدوں رُکدی
اے جدوں~$M_e$ انپٹ~$w$ اُتے رُکدی اے،
تے دوجی صورت وچ آؤٹ پٹ~$0$ نال رُکدی اے۔ دوجے لفظاں
وچ~$E \concat D$ رُکّن دا مسئلہ حل
کر دیوے گی۔ پر~\olref[hal]{thm:halting-problem}
توں سانوں پتا اے کہ ایہو جیہی کوئی ٹیورنگ مشین
موجود نہیں ہو سکدی۔
\end{proof}

\begin{cor}\ollabel{cor:undecidable-sat}%
پہلے درجے دی منطق دا دتا ہویا !!{sentence}
پورا ہو سکدا اے یا نہیں، ایہدا فیصلہ نہیں کیتا جا سکدا۔
\end{cor}

\begin{proof}
  فرض کرو اک ٹیورنگ مشین~$S$ ایہ فیصلہ کر
  سکدی ہووے کہ جملہ پورا ہو سکدا اے یا نہیں۔
  فیر اسی فیصلے دا مسئلہ ایویں حل کر سکدے
  آں۔ اک !!a{sentence}~$B$ انپٹ وجوں ملے تاں
  $!B$ نوں اک خانہ سجّے ول سرکا دیو۔
  خانہ~$1$ اُتے واپس آ کے نشان~$\lnot$ لکھ دیو۔

  ہن ٹیورنگ مشین~$S$ چلاؤ۔ ایہ آخرکار پٹی
  اُتے یا آؤٹ پٹ~$1$ نال رُکے گی (جے
  $\lnot !B$ پورا ہو سکدا ہووے) یا~$0$
  نال (جے~$\lnot !B$ پورا نہ ہو سکدا ہووے)۔
  جے~$\TMstroke$ خانہ~$1$ اُتے ہووے تاں
  اوہ مٹا دیو؛ جے خانہ~$1$ خالی ہووے تاں
  اوتھے~$\TMstroke$ لکھ دیو، فیر رُک جاؤ۔

  ایہ ٹیورنگ مشین ہمیشہ رُکدی اے، تے
  آؤٹ پٹ~$1$ اُدوں تے صرف اُدوں دیندی اے جدوں
  $\lnot
  !B$ پورا نہ ہو سکدا ہووے، ورنہ~$0$۔
  کیوں جے~$!B$ اُدوں تے صرف اُدوں درست اے جدوں
  $\lnot !B$~پورا نہ ہو سکدا ہووے،
  ایس لئی مشین~$1$ اُدوں تے صرف اُدوں دیندی اے
  جدوں~$!B$ درست ہووے، ورنہ~$0$؛ یعنی
  ایہ فیصلے دا مسئلہ حل کر دیوے گی۔
\end{proof}

\begin{explain}
ایس لئی کوئی ایہو جیہی ٹیورنگ مشین نہیں جیہڑی
سوال «کی~$!B$ پہلے درجے دی منطق دا
درست !!{sentence} اے؟» دا ہمیشہ ٹھیک
«ہاں» یا «نہیں» جواب دیوے۔ پر اک ایہو جیہی
ٹیورنگ مشین \emph{موجود اے} جیہڑی جدوں صحیح
جواب «ہاں» ہووے تاں «ہاں» آکھدی اے، تے جدوں
صحیح جواب «نہیں» ہووے تاں رُکدی نہیں۔
ایہ پہلے درجے دی منطق دے صحت مندی تے
تمامیت دے قضیے، تے ایس حقیقت توں نکلدا اے
کہ سارے !!{derivation}s مؤثر طریقے نال
اک اک کر کے فہرست وچ لکھے جا سکدے نیں۔
\end{explain}

\begin{thm}
  \ollabel{thm:valid-ce}%
  پہلے درجے دے !!{sentence}s درست نیں یا نہیں، ایہ
  طے کرن دا مسئلہ نیم فیصلہ کرن جوگ اے۔ اک ٹیورنگ مشین~$E$
  موجود اے جیہڑی پٹی اُتے پہلے درجے دی
  منطق دا اک !!a{sentence}~$!B$ انپٹ وجوں
  رکھ کے شروع کیتی جائے تاں~$E$ آخرکار رُک
  کے آؤٹ پٹ~$1$ اُدوں تے صرف اُدوں دیندی اے
  جدوں~$!B$ درست ہووے۔ دوجی صورت وچ رُکدی نہیں۔
\end{thm}

\begin{proof}
  پہلے درجے دی منطق دے سارے ممکن
  !!{derivation}s اک مؤثر طریقے نال اک
  توں بعد اک بنائے جا سکدے نیں۔ مشین~$E$
  ایہی کردی اے، تے جدوں ایہ اک ایہو جیہا
  !!a{derivation} لبھدی اے جیہڑا
  $\Proves
  !B$ دکھاندا اے، تاں آؤٹ پٹ~$1$ نال
  رُک جاندی اے۔ صحت مندی دے قضیے مطابق
  جے~$E$ آؤٹ پٹ~$1$ نال رُکدی اے تاں
  ایس دی وجہ~$\Entails !B$ اے۔
  تمامیت دے قضیے مطابق جے~$\Entails !B$
  ہووے تاں اک !!a{derivation} موجود اے
  جیہڑا~$\Proves !B$ دکھاندا اے۔ چونکہ
  $E$ سارے ممکن !!{derivation}s ترتیب
  نال بناندی اے، ایس لئی اوہ آخرکار
  $\Proves !B$ دکھان والا اخذ لبھ لوے گی،
  تے آؤٹ پٹ~$1$ نال رُک جائے گی۔
\end{proof}
\endgroup
% /OLP-0272

% OLP-0273 content/turing-machines/undecidability/trakhtenbrot.tex
\begingroup
\def\olfilename{trakhtenbrot}
\expandafter\def\csname ol@sectioncs\endcsname{section}

\olfileid{tur}{und}{tra}
\olsection{تراختنبرُوت دا قضیہ}

\begin{explain}
  \olref[rep]{sec} وچ اسی !!{sentence}s~$!T(M,w)$
  تے~$!E(M,w)$ دی تعریف کیتی سی،
  ٹیورنگ مشین~$M$ تے انپٹ لڑی~$w$ لئی۔ فیر
  \olref[ver]{lem:valid-if-halt} تے
  \olref[ver]{lem:halt-if-valid} وچ دکھایا
  کہ~$!T(M,w) \lif !E(M,w)$ اُدوں تے صرف اُدوں
  درست اے جدوں~$M$ انپٹ~$w$ اُتے شروع
  ہو کے آخرکار رُکدی اے۔ رُکّن دا مسئلہ
  ناقابلِ فیصلہ اے۔ ایس توں نکلدا اے کہ
  پہلے درجے دی منطق وچ !!{sentence}s درست نیں یا
  نہیں، ایہ فیصلہ نہیں ہو سکدا۔ ایہ وی فیصلہ نہیں
  ہو سکدا کہ اوہ پورے ہو سکدے نیں یا نہیں
  (\cref{tur:und:uns:thm:decision-prob,tur:und:uns:cor:undecidable-sat})۔

  پر جملیاں دی درستی تے پورا ہو سکّن دیاں تعریفاں
  ہر قسم دیاں !!{structure}s لئی دتیاں
  جاندیاں نیں، چاہے اوہ محدود ہون یا لامتناہی۔
  تُہانوں شاید لگے کہ ایہ فیصلہ کرنا سوکھا
  ہووے گا کہ اک دتا ہویا !!a{sentence} کسے محدود
  !!{structure} وچ پورا ہو سکدا اے یا نہیں
  (یا ساریاں محدود !!{structure}s وچ سچ
  اے یا نہیں)۔ اسی فیصلے دے مسئلے دی
  حل نہ ہو سکّن والا ثبوت ایویں ڈھال سکدے
  آں کہ اوہ ایہناں صورتاں وچ وی فیصلے دا ناممکن ہونا دکھاوے۔

  پہلے~\olref[ver]{lem:halt-if-valid} دے
  ثبوت ول ویکھو۔ اوتھے اسی اک
  نمونہ~$\Struct M$ بنایا سی جو~$!T(M,w)$
  نوں پورا کردا اے تے جیہڑا
  عین اوہ کم بیان کردا اے جیہڑا مشین~$M$
  انپٹ~$w$ اُتے شروع ہو کے کردی اے۔
  اوس نمونے دا دائرہ~$\Nat$ سی، یعنی لامتناہی۔
  پر جے~$M$ انپٹ~$w$ اُتے واقعی رُکدی
  ہووے، تاں اسی ایس طرح اک محدود
  نمونہ~$\Struct M'$ بنا سکدے آں۔ فرض کرو
  $M$ انپٹ~$w$ اُتے شروع ہو کے~$k$ قدم
  بعد رُکدی اے۔ دائرہ~$\Domain{M'}$ لئی
  سیٹ~$\{0, \dots, n\}$ لو، جتھے
% PNB-CMP-038: leave a point strictly beyond the last executed time and needed tape position.
  $n$، $k$ تے~$w$ دی لمبائی وچوں وڈے
  عدد توں اک ودھ اے؛ تے ایہ تعبیر رکھو:
  \[
    \Assign{\prime}{M'}(x) = 
    \begin{cases}
      x + 1 &\text{جے $x < n$ ہووے}\\
      n &\text{ورنہ،}
    \end{cases}
  \]
  تے~$\tuple{x,y} \in \Assign{<}{M'}$
  اُدوں تے صرف اُدوں جدوں~$x < y$ یا~$x = y = n$۔
  باقی~$\Struct{M'}$ نوں~$\Struct{M}$ وانگ
  تعریف کرو۔ $\Struct{M'}$ دی تعریف
  مطابق، عین~\olref[ver]{lem:halt-if-valid}
  دے ثبوت وانگ، $\Sat{M'}{!T(M,w)}$۔
  تے چونکہ اسی فرض کیتا اے کہ~$M$
  انپٹ~$w$ اُتے رُکدی اے، ایس لئی
  $\Sat{M'}{!E(M,w)}$۔ لہٰذا
  $\Struct{M'}$، $!T(M,w) \land !E(M,w)$
  دا محدود نمونہ اے (دھیان دیو کہ
  اسی~$\lif$ دی تھاں~$\land$ رکھیا اے)۔
  
  ثبوت دا ادھا کم ہو گیا: اسی دکھا دتا
  کہ جے~$M$ انپٹ~$w$ اُتے رُکے، تاں
  $!T(M,w) \land !E(M,w)$ دا محدود نمونہ
  اے۔ پر ایہدی اُلٹ گل سچ نہیں۔ کجھ
  ٹیورنگ مشیناں کسے انپٹ~$w$
  اُتے نہیں رُکدیاں، پر~$!T(M,w)
  \land !E(M,w)$ دا فیر وی محدود نمونہ
  ہو سکدا اے۔ مثال لئی مشین~$M$ لو
  جیہدی اکّو حالت~$q_0$ تے اکّو
  حکم~$\delta(q_0,\TMblank) =
  \tuple{q_0,\TMblank,\TMstay}$ اے۔
  خالی انپٹ~$w = \emptyseq$
  اُتے شروع ہو کے ایہ مشین کدی نہیں
  رُکدی: ایہ لامتناہی چکر وچ اے، پر نہ
  پٹی بدلدی اے نہ سر ہلاندی اے۔
  ساریاں مکمل حالتاں اکّو جہیاں نیں
  (اوہی حالت، سر دی اوہی جگہ، پٹی دے
  اوہی نشان)۔ اسی اک محدود
  !!{structure}~$\Struct{M''}$ تعریف کر
  سکدے آں جیہڑی~$!T(M,\emptyseq) \land
  !E(M,\emptyseq)$ پورا کردی اے (مشق)۔
  البتہ اسی~$!T(M,w)$ نوں مناسب طریقے
  نال بدل سکدے آں تاکہ ایداں دیاں
  !!{structure}s خارج ہو جان۔

  \begin{prob}
    ٹیورنگ مشین~$M$ لو جیہدی
    اکّو حالت~$q_0$ تے اکّو
% PNB-CMP-039: the sole state and both successor interpretations are M-double-prime.
    حکم~$\delta(q_0,\TMblank) =
    \tuple{q_0,\TMblank,\TMstay}$ اے۔
    $\Domain{M''} = \{0, 1, 2\}$،
    $\Assign{\prime}{M''}(0) =
    \Assign{\prime}{M''}(1) = 1$
    تے~$\Assign{\prime}{M''}(2) = 2$ رکھو؛
    تے~$\Assign{<}{M''} = \{\tuple{0,1},
    \tuple{1,1}, \tuple{2,2}\}$۔
    $\Assign{\Obj Q_{q_0}}{M''}$،
    $\Assign{\Obj
    S_{\TMblank}}{M''}$ تے
    $\Assign{\Obj S_{\TMendtape}}{M''}$
    ایویں تعریف کرو کہ~$!T(M,\emptyseq)$
    تے~$!E(M,\emptyseq)$ دونیں سچ
    ہو جان، تے سمجھاؤ کیوں۔ اشارہ:
    $\delta(q_0, \TMendtape)$
    غیر تعریف شدہ اے۔ یقین کرو کہ
    \begin{align*}
    & \Obj Q_{q_0}(\num{1}, \num{n}) \land \Obj S_{\TMendtape}(\num{0}, \num{n})
      \land 
      \lforall[x][(\num{0} < x \lif \Obj S_\TMblank(x, \num{n}))] \text{\quad ہر $n \in \Nat$ لئی}\\
    & \lexists[y][(\Obj Q_{q_0}(\num{0}, y) \land \Obj S_{\TMendtape}(\num{0}, y))]
    \end{align*}
    دونیں~$\Struct{M''}$ وچ سچ ہون۔
  \end{prob}
\end{explain}

ٹیورنگ مشین~$M$ دے انپٹ
$w = \sigma_{i_1}\dots\sigma_{i_k}$ اُتے
چلن نوں بیان کرن والے !!{sentence}s ویکھو:
\begin{enumerate}
\item عدداں تے~$<$ نوں بیان کرن والے
مسلّمے (بالکل~\olref[rep]{sec} وچ
$!T(M,w)$ دی تعریف وانگ)۔
\item ابتدائی مکمل حالت نوں بیان کرن والے
مسلّمے: $!T(M,w)$ دی تعریف وانگ۔
\item اک مکمل حالت توں اگلی مکمل حالت
ول تبدیلی نوں بیان کرن والے مسلّمے:

اگے~$!A(x, y)$ پہلے وانگ ہووے، تے
\[
  !B(y) \ident \lforall[x][(x < y \lif \eq/[x][y])].
\]
\begin{enumerate}
\item \ollabel{rep-right} ہر حکم
$\delta(q_i, \sigma) =
  \tuple{q_j, \sigma', \TMright}$
لئی ایہ !!{sentence}:
\begin{align*}
& \lforall[x][\lforall[y][(
   (\Obj Q_{q_i}(x, y) \land \Obj S_{\sigma}(x, y)) \lif {}]] \\
&\qquad   (\Obj Q_{q_j}(x', y') \land \Obj S_{\sigma'}(x, y') \land
\formula A(x, y) \land !B(y')))
\end{align*}
\item \ollabel{rep-left} ہر حکم
$\delta(q_i, \sigma) =
  \tuple{q_j, \sigma', \TMleft}$
لئی ایہ !!{sentence}:
% PNB-CMP-040: non-edge left writes x-prime and must also force B at the next time.
\begin{align*}
& \lforall[x][\lforall[y][
    ((\Obj Q_{q_i}(x', y) \land \Obj S_{\sigma}(x', y)) \lif {}]]\\
& \qquad   (\Obj Q_{q_j}(x, y') \land \Obj S_{\sigma'}(x', y') \land
\formula A(x', y) \land !B(y'))) \land {}\\
& \lforall[y][((\Obj Q_{q_i}(\num{0}, y) \land \Obj S_{\sigma}(\num{0},
    y)) \lif {}]\\
& \qquad (\Obj Q_{q_j}(\num{0}, y') \land \Obj S_{\sigma'}(\num{0},
  y') \land !A(\num{0}, y) \land !B(y')))
\end{align*}
\item \ollabel{rep-stay} ہر حکم
$\delta(q_i, \sigma) =
  \tuple{q_j, \sigma', \TMstay}$
لئی ایہ !!{sentence}:
\begin{align*}
& \lforall[x][\lforall[y][(
   (\Obj Q_{q_i}(x, y) \land \Obj S_{\sigma}(x, y)) \lif {}]] \\
&\qquad (\Obj Q_{q_j}(x, y') \land \Obj S_{\sigma'}(x, y') \land
\formula A(x, y) \land !B(y')))
\end{align*}
\end{enumerate}
ویکھو کہ~$M$ دیاں حالتی تبدیلیاں بیان
کرن والے !!{sentence}s وچوں ہر اک
$!T(M,w)$ دے متعلقہ !!{sentence}
ورگا اے۔ ہر اک دے آخر وچ
$!B(y')$ جوڑیا اے۔ $!B(y')$ ایہ
یقینی بناندا اے کہ «اگلی» مکمل حالت
دا نمبر~$y'$، پچھلے سارے نمبراں
$\Obj 0$، $\Obj 0'$، \dots{} توں وکھرا اے۔
% \item Sentences that express consistency conditions:
% \begin{enumerate}
%   \item\ollabel{rep-con-symb}%
%   !!^a{sentence} that says that no square can have more than one
%   symbol on it at any given time:
%   \[\lforall[x][\lforall[y][(\Obj S_{\sigma}(x, y) \lif \lnot \Obj
%   S_{\sigma'}(x, y))]],\]
%   for every pair $\sigma \neq \sigma'$ of tape symbols of~$T$.
%   \item\ollabel{rep-con-state}%
%   !!^a{sentence} that says that $M$ cannot be in more than one state
%   at any given time:
%   \[\lforall[x][\lforall[y][(\Obj Q_{q}(x, y) \lif
%   \lforall[z][\lnot \Obj
%   Q_{q'}(z, y))]],\]
%   for every pair $q \neq q'$ of states of~$T$.
%   \item\ollabel{rep-con-tape}%
%   !!^a{sentence} that says that $M$ cannot be on more than one tape
%   square at any given time:
%   \[\lforall[x][\lforall[y][(\Obj Q_{q}(x, y) \lif
%   \lforall[z][\eq/[z][y]\lif \lnot \Obj
%   Q_{q'}(z, y))]],\]
%   for every pair $q$, $q'$ of states of~$T$.
% \end{enumerate}
\end{enumerate}
اُتلے سارے !!{sentence}s جوڑ کے بنے عطف نوں
$!T'(M, w)$ آکھو، ٹیورنگ
مشین~$M$ تے انپٹ~$w$ لئی۔

\begin{lem}\ollabel{lem:halts-sat}
  جے~$M$ انپٹ~$w$ اُتے شروع ہو کے
  رُکدی اے، تاں~$!T'(M,w) \land !E(M,w)$
  دا اک محدود نمونہ اے۔
\end{lem}

\begin{proof}
  $\Struct{M'}$ نوں
  \olref[ver]{lem:halt-if-valid} دے ثبوت
  وانگ لو، پر ایہ رکھو:
  \begin{align*}
    \Domain{M'} & = \{0, \dots, n\},\\
    \Assign{\prime}{M'}(x) & = 
    \begin{cases}
      x + 1 &\text{جے $x < n$ ہووے}\\
      n &\text{ورنہ،}
    \end{cases} \\
    \tuple{x,y} \in \Assign{<}{M'} &\text{عین اوہدوں جدوں $x < y$ یا $x = y = n$،}
  \end{align*}
% PNB-CMP-038: cap strictly beyond the last transition and input extent.
  جتھے~$n = \max(k,\len{w})+1$،
  تے~$k$ اوہ سب توں گھٹ عدد اے جیہدے لئی
  مشین~$M$ انپٹ~$w$ اُتے شروع ہو
  کے~$k$ قدم چُکّن توں بعد رُک
  چکی ہووے۔
  اسی ایہ جانچ مشق لئی چھڈدے آں کہ
% PNB-CMP-041: E is a named formula, and the modified theory is T-prime.
  $\Sat{M'}{!T'(M,w) \land !E(M,w)}$۔
\end{proof}

\begin{prob}
  \olref[tur][und][tra]{lem:halts-sat}
  دا ثبوت مکمل کرو: ایہ ثابت کرو کہ
% PNB-CMP-041: exercise must establish the new T-prime conjunction with !E.
  $\Sat{M'}{!T'(M,w) \land !E(M,w)}$۔
\end{prob}

\begin{lem}\ollabel{lem:sat-halts}
  جے~$!T'(M,w) \land !E(M,w)$ دا
  اک محدود نمونہ ہووے، تاں~$M$
  انپٹ~$w$ اُتے شروع ہو کے رُکدی اے۔
\end{lem}

\begin{proof}
  اسی متضادِ نقیض ثابت کراں گے۔
  فرض کرو~$M$ انپٹ~$w$ اُتے شروع ہو
  کے نہیں رُکدی۔ جے~$!T'(M,w) \land
  !E(M,w)$ دا سرے توں کوئی نمونہ
  نہ ہووے تاں کم پورا اے۔ ایس لئی
% PNB-CMP-042: the assumed model belongs to the lemma's T-prime theory.
  فرض کرو~$\Struct{M}$،
  $!T'(M,w) \land !E(M, w)$ دا نمونہ
  اے۔ سانوں دکھانا اے کہ ایہ محدود
  نہیں ہو سکدا۔
  
  اسی~\olref[ver]{lem:config} وانگ
  ثابت کر سکدے آں کہ جے~$M$،
  انپٹ~$w$ اُتے شروع ہو کے
  $n$ قدم بعد ہن تک نہ رُکی ہووے
  تے ایہ قدم-گنتی مثبت ہووے،
% PNB-CMP-043: B is forced at positive successor times, not the zero base.
  تاں~$!T'(M,w)
  \Entails !C(M, w, n) \land
  !B(\num{n})$؛ صفر اُتے ابتدائی
  مکمل حالت سِدھی ابتدائی مسلّماں
  توں ملدی اے۔ چونکہ~$M$ انپٹ~$w$
  اُتے شروع ہو کے کدی نہیں رُکدی،
  ایس لئی~$!T'(M,w) \Entails !C(M, w, n) \land
  !B(\num{n})$ ہر مثبت~$n \in \Nat$
  لئی۔ دھیان دیو کہ
  \olref[rep]{prop:mlessk} مطابق
  $!T'(M,w) \Entails \num{k} < \num{n}$
  ہر~$k < n$ لئی۔ نالے
  $!B(\num{n}) \Entails \num{k} < \num{n} \lif
  \eq/[\num{k}][\num{n}]$۔
  ایس لئی~$\Sat{M}{\eq/[\num{k}][\num{n}]}$
  ہر~$k < n$ لئی،
  یعنی ٹرماں~$\num{k}$ دی لامتناہی
  لڑی دیاں قدراں~$\Struct{M}$ وچ
  آپس وچ وکھ وکھ ہون گیاں۔ ایس لئی
  $\Domain{M}$ لامتناہی ہونا
  ضروری اے؛ ایس لئی~$\Struct{M}$،
  $!T'(M,w) \land !E(M, w)$
  دا محدود نمونہ نہیں ہو سکدا۔
\end{proof}

\begin{prob}
  \olref[tur][und][tra]{lem:sat-halts}
  دا ثبوت مکمل کرو: ایہ ثابت کرو کہ
  جے~$M$ انپٹ~$w$ اُتے شروع ہو کے
  $n$ قدم بعد نہ رُکی ہووے، تاں
% PNB-CMP-043: the exercise needs the positive-time qualification.
  $!T'(M,w) \Entails !B(\num{n})$
  ہر مثبت عدد لئی۔ صفر دی
  مکمل حالت ابتدائی مسلّماں توں لو۔
\end{prob}

\begin{thm}[تراختنبرُوت دا قضیہ]
  \ollabel{thm:trakhtenbrodt}
  ایہ فیصلہ نہیں کیتا جا سکدا کہ
  پہلے درجے دی منطق دا دتا ہویا
  !!{sentence} کوئی محدود نمونہ رکھدا اے
  یا نہیں (یعنی اوہ کسے محدود
  ساخت وچ پورا ہو سکدا اے یا نہیں)۔
\end{thm}

\begin{proof}
  فرض کرو کوئی ٹیورنگ مشین~$F$
  محدود نمونہ ہون دا فیصلہ
  کردی ہووے۔ فیر کسے وی
  ٹیورنگ مشین~$M$ تے انپٹ~$w$
  توں اسی~$!T'(M,w) \land !E(M,w)$
  جملہ حساب کرکے~$F$ نال فیصلہ
  کر سکدے آں کہ ایہدا محدود
  نمونہ اے یا نہیں۔
  \cref{tur:und:tra:lem:halts-sat,tur:und:tra:lem:sat-halts}
  مطابق ایہدا محدود نمونہ اُدوں تے
  صرف اُدوں اے جدوں~$M$ انپٹ~$w$
  اُتے شروع ہو کے رُکدی اے۔
  ایس لئی~$F$ نال رُکّن دا
  مسئلہ حل ہو جائے گا۔ سانوں
  پتا اے کہ ایہ حل نہیں ہو سکدا۔
\end{proof}

\begin{cor}\ollabel{cor:fproof-incomp}%
% PNB-CMP-044: the contradiction uses effective enumeration of derivations.
  ہر \emph{محدود} ساخت وچ درستی لئی
  کوئی صحت مند تے کامل !!{derivation} دا نظام،
  جیہدے اخذ مؤثر طور اُتے گنتی وار
  فہرست وچ لکھے جا سکن، نہیں ہو سکدا۔
  ایہ ایہو جیہا !!a{derivation} دا نظام ہووے گا
  جیہدے وچ~$\Proves !B$
  اُدوں تے صرف اُدوں ہووے جدوں
  $\Sat{M}{!B}$ ہر محدود
  !!{structure}~$\Struct{M}$ وچ ہووے۔
\end{cor}

\begin{proof}
  مشق۔
\end{proof}

\begin{prob}
  \olref[tur][und][tra]{cor:fproof-incomp}
  ثابت کرو۔ دھیان دیو کہ
  $!B$ ہر محدود !!{structure}
  وچ اُدوں تے صرف اُدوں پورا اتردا
  اے جدوں~$\lnot !B$ کسے وی محدود
  ساخت وچ پورا نہیں ہو سکدا۔
  سمجھاؤ کہ محدود نمونہ ہون دا مسئلہ،
  \olref[tur][und][uns]{thm:valid-ce}
  والے معنی وچ نیم فیصلہ
  کرن جوگ کیوں اے۔ ایس
  توں دلیل دیو کہ جے محدود
% PNB-CMP-044: effective enumeration is needed for the corollary's decision procedure.
  درستی لئی کوئی مؤثر
  !!a{derivation} دا نظام
  ہووے، تاں محدود نمونہ ہون دا
  فیصلہ ہو سکدا ہووے گا۔
\end{prob}
\endgroup
% /OLP-0273

\OLEndChapterHook
\endgroup
% /OLP-0264

\OLEndPartHook
\endgroup
% /OLP-0252

% OLP-0274 content/incompleteness/incompleteness.tex
\begingroup
\def\olfilename{incompleteness}
\expandafter\def\csname ol@sectioncs\endcsname{section}

\olpart{inc}{نامکملیت}

\begin{editorial}
 ایہ حصّہ نامکملیت دے قضیاں بارے اے۔ ایہدے لئی پہلے درجے دی
 منطق والے حصیاں دا مواد (خاص کر ثبوتی نظام) تے قابلِ حسابیت والے
 حصے دے بازگشتی فنکشناں دا مواد چاہیدا اے۔ ایہ Jeremy Avigad
 دے نوٹس اُتے مبنی اے، جنہاں وچ Richard Zach نے ترمیماں کیتیاں نیں۔
\end{editorial}

% OLP-0275 content/incompleteness/introduction/introduction.tex
\begingroup
\def\olfilename{introduction}
\expandafter\def\csname ol@sectioncs\endcsname{section}

\olchapter{inc}{int}{نامکملیت دا تعارف}

% OLP-0276 content/incompleteness/introduction/historical-background.tex
\begingroup
\def\olfilename{historical-background}
\expandafter\def\csname ol@sectioncs\endcsname{section}

\olfileid{inc}{int}{bgr}

\olsection{تاریخی پس منظر}

ایس حصے وچ اسی مختصر طور اُتے اوہناں تاریخی پیش رفتاں دی گل
کراں گے جیہڑیاں نامکملیت دے قضیاں دا پس منظر سمجھّن وچ مدد
دیندیاں نیں۔ خاص طور اُتے، پہلے ریاضیاتی منطق دی تاریخ دا
سرسری خاکہ دیواں گے۔ فیر ریاضی دیاں بنیاداں دی تاریخ
بارے کجھ گلّاں کراں گے۔

\begin{digress}
  «ریاضیاتی منطق» دی ترکیب دے دو مطلب ہو سکدے نیں۔ «ریاضیاتی»
  نوں موضوع دی صفت سمجھو تاں مراد «ریاضی دی منطق» اے، یعنی
  ریاضی وچ استدلال دے اصول۔ جے ایہنوں طریقے دی صفت سمجھو
  تاں مراد «منطق دی ریاضی» اے، یعنی استدلال دے اصولاں دا
  ریاضیاتی مطالعہ۔ اگلے بیان وچ ریاضیاتی منطق دے ایہ دونیں
  مطلب آؤندے نیں، کئی وار اکّو ویلے۔
\end{digress}

منطق دا باقاعدہ مطالعہ ارسطو توں شروع ہویا منیا جاندا اے۔
اوہ لگ بھگ 384--322 \textsc{قبلِ مسیح} وچ جیویا۔ اوہ دیاں
\emph{زمرے (Categories)}، \emph{پہلے تجزیے \pnbLatin{(Prior analytics)}} تے \emph{بعد دے تجزیے \pnbLatin{(Posterior analytics)}}
علمی استدلال دے اصولاں دا باقاعدہ مطالعہ کردیاں نیں۔
ایہناں وچ قیاس، یعنی مقرر مقدمیاں توں نتیجہ کڈھن والی
صورت، دا وی پورا تے منظم مطالعہ اے۔

ارسطو دی منطق قرونِ وسطیٰ وچ مدرسی فلسفے اُتے حاوی رہی۔
ایتھوں تک کہ اٹھارہویں صدی وچ وی کانٹ دا خیال سی کہ
ارسطو دی منطق مکمل اے تے ایہنوں بدلّن دی لوڑ نہیں۔
پر قیاس دا نظریہ ریاضیاتی استدلال دے سطحی ترین پہلو توں
اگے کوئی ٹھیک نمونہ نہیں بنا سکدا۔ اوہدے توں اک صدی
پہلاں، نیوٹن دے ہم عصر لائبنز نے منطقی استدلال لئی
اک مکمل «حساب» دا تصور کیتا سی۔ اوہنے ایداں دا حساب
تیار کرن ول کجھ ابتدائی قدم وی چُکے، تے بنیادی طور اُتے
قضیاتی منطق دی اک صورت بیان کیتی۔

انیویں صدی منطق لئی فیصلہ کن موڑ سی۔ 1854 وچ
جارج بُول نے \emph{سوچ دے قانوناں \pnbLatin{(The Laws of Thought)}}
لکھی، جیہدے وچ قضیاتی منطق دا مفصل الجبری مطالعہ اے۔
ایہ مطالعہ اج دی پیش کش توں بہتا وکھرا نہیں۔ 1879 وچ
گوٹلوب فریگے نے اپنی \emph{Begriffsschrift}
(«تصوری تحریر») شائع کیتی۔ اوہنے قضیاتی منطق وچ
مسوّراں تے تعلق جوڑے۔ ایس طرح ایہدے وچ
پہلے درجے دی منطق وی شامل سی۔ حقیقت وچ فریگے دے
منطقی نظاماں وچ اعلیٰ درجے دی منطق تے ہور وی کجھ
شامل سی۔ اپنی \emph{حساب دے بنیادی قانوناں
  \pnbLatin{(Basic Laws of Arithmetic)}} وچ فریگے نے ایہ
دکھاون دی کوشش کیتی کہ سارا حساب اوہدی
Begriffsschrift وچ صرف منطقی مفروضیاں توں اخذ
ہو سکدا اے۔ بدقسمتی نال اوہ مفروضے متضاد نکلے،
جیویں رسل نے 1902 وچ دکھایا۔ پر متضاد مسلمے
نوں پاسے رکھ دیئیے تاں فریگے نے جدید منطق لگ بھگ
اکیلے ہی ایجاد کر دتی؛ ایہ حیران کن کارنامہ سی۔
بُول توں بعد الجبری سوچ رکھنے والے ہور مفکراں،
جنہاں وچ پئرس تے شروڈر شامل سن، نے وی
مسوّری منطق آزاد طور اُتے تیار کیتی۔

ہن ریاضی دیاں بنیاداں وچ پیش رفت ول آئیے۔
ظاہر اے کہ منطق ریاضی وچ اہم کم کردی اے، ایس لئی
ایہ پیش رفت منطق دیاں اوہناں تبدیلیاں نال
جیہڑیاں اسیں ہُن تک بیان کیتیاں نیں،
کافی جڑی ہوئی اے۔ مثال لئی فریگے نے اپنی
منطق ایس واضح مقصد نال بنائی سی کہ ساری
ریاضی صرف اوس دے منطقی ڈھانچے اُتے کھڑی
ہو سکدی اے۔ خاص کر اوہ دکھانا چاہندا سی
کہ ریاضی تجربے توں آزاد \emph{تجزیاتی}
سچائیاں اُتے مشتمل اے، نہ کہ تجربے توں
آزاد \emph{ترکیبی} سچائیاں اُتے، جیویں
کانٹ سمجھدا سی۔

بہت سارے لوک ریاضی دے باقاعدہ جنم نوں یونانیاں نال
جوڑدے نیں۔ لگ بھگ 300 قبلِ مسیح لکھی اقلیدس دی
\emph{عناصر (Elements)} پہلے ہی یونانی ریاضی دی
اک پختہ مثال اے، جیہدے وچ دلیل دی پختگی تے ٹھیک پن اُتے
زور اے۔ اقلیدس دی \emph{عناصر} وچ تعریفاں تے
ثبوت اج وی ہائی سکولاں لئی ہندسے دیاں درسی کتاباں وچ لگ بھگ
اوہو جہے موجود نیں (جتھے ہائی سکولاں وچ ہندسہ ہن وی
پڑھایا جاندا اے)۔ ریاضیاتی استدلال دے ایس نمونے
نوں صرف ریاضی وچ نہیں بلکہ فلسفے دیاں شاخاں وچ وی
سخت استدلال دی مثال سمجھیا گیا اے۔ (اسپینوزا نے
تاں اخلاقی تے مذہبی دلیلیں وی اقلیدسی طرز وچ پیش
کیتیاں؛ اوہ ویکھ کے حیرانی ہوندی اے!)

سترہویں صدی وچ نیوٹن تے لائبنز نے حسابان ایجاد کیتا۔
(ایہ جھگڑا صدیاں تک چلدا رہیا کہ پہلاں ایہ کم کس نے کیتا۔ اج
اکثر محقق سمجھدے نیں کہ دوناں دی پیش رفت وڈی حد
تک اک دوجے توں آزاد سی۔) حسابان وچ، مثال لئی،
لامتناہی حد تک چھوٹیاں مقداراں دے لامتناہی جوڑاں بارے
استدلال ہوندا اے۔ ایہناں خاصیتاں دی وجہ توں
بشپ برکلے نے تنقید کیتی: اوہدے مطابق رب اُتے
یقین اوہدے ویلے دی ریاضی توں گھٹ معقول نہیں سی۔
اٹھارہویں صدی وچ حسابان دے طریقے وڈے پیمانے
اُتے ورتے گئے؛ مثال لئی لیونارڈ اویلر نے
لامتناہی جوڑاں والے حساب نال حیران کن
نتیجے کڈھے۔

انیویں صدی وچ ریاضی داناں نے حسابان نوں ہور مضبوط
بنیاد اُتے رکھ کے برکلے دی تنقید دا جواب دین دی
کوشش کیتی۔ کوشی، وائرشٹراس، بولزانو تے ہوراں دی
کوششاں توں حد، استمرار، مشتق گیری تے تکامل دیاں
اج والیاں تعریفاں «ایپسلون تے ڈیلٹا» دی مدد نال
بنیاں؛ یعنی اوہناں وچ لامتناہی حد تک چھوٹی
مقداراں دا کوئی حوالہ نہ رہا۔ ایس صدی دے پچھلے
حصے وچ ریاضی داناں نے ہور اگے ودھ کے حسابان
دے سارے پہلو، حقیقی اعداد سمیت، قدرتی عدداں
دی مدد نال سمجھاون دی کوشش کیتی۔ (کرونیکر:
«رب نے پورے عدد بنائے؛ باقی سب انسان دا کم اے۔»)
1872 وچ ڈیڈیکنڈ نے «استمرار تے غیر ناطق اعداد»
\pnbLatin{(Continuity and the irrational numbers)} لکھی،
جیہدے وچ اوہنے ناطق عدداں دے سیٹاں راہیں
حقیقی اعداد دی «تعمیر» دکھائی (تے تُسی جان دے
او کہ ناطق عدداں نوں قدرتی عدداں دے جوڑیاں
وجوں سمجھیا جا سکدا اے)۔ 1888 وچ اوہنے
«Was sind und was sollen die Zahlen» (موٹے
طور اُتے، «قدرتی عدد کیہ نیں تے اوہناں نوں
کیہ ہونا چاہیدا اے؟») لکھی، جیہدا مقصد
قدرتی عدداں نوں خالص «منطقی» لفظاں وچ
سمجھانا سی۔ 1887 وچ کرونیکر نے
«\"Uber den Zahlbegriff» («عدد دے تصور بارے»)
لکھی، جیہدے وچ اوہنے ساریاں ریاضیاتی
اشیاں نوں صحیح عدداں دے ذریعے بیان کرن
دی گل کیتی۔ 1889 وچ جوزیپے پیانو نے
قدرتی عدداں لئی رسمی، علامتی مسلّمے دتے۔

انیویں صدی دے آخر وچ لامتناہی نال برتاؤ
وچ وی اک نویں دلیری آئی۔ اوہدے توں پہلاں
لامتناہی اشیاں تے ساختاں، مثلاً قدرتی عدداں
دا سیٹ، احتیاط نال ورتیاں جاندیاں سن۔
«لامتناہی تعداد وچ» دا مطلب «جنے چاہو اونے»
سمجھیا جاندا سی، تے «حد ول جانا» دا مطلب
«جنا چاہو اونا نیڑے آنا»۔ پر جارج کانٹور
نے دکھایا کہ لامتناہی نوں اپنے پورے معنی
وچ وی لیا جا سکدا اے۔ کانٹور، ڈیڈیکنڈ
تے ہوراں دے کم نے ریاضی نوں سیٹاں دے
عام نظریے راہیں سمجھّن دا اوہ طریقہ متعارف
کراؤن وچ مدد دتی جیہڑا اج وڈے پیمانے
اُتے قبول اے۔

ایہ گلاں سانوں ویہویں صدی دی منطق تے
ریاضی دیاں بنیاداں دی پیش رفت تک لے
آندیاں نیں۔
% PNB-CMP-045: the dated 1902 event is Russell's letter to Frege, not the date of discovery.
1902 وچ رسل نے فریگے نوں اوس دے منطقی
نظام وچ ملے تضاد بارے خط لکھ کے آگاہ
کیتا۔ 1904 وچ زرمیلو نے «انتخاب دے
مسلّمے» دی مدد نال کانٹور دا حُسنِ ترتیب
دا اصول ثابت کیتا۔ ایس مسلّمے دے جائز
ہون بارے کافی بحث ہوئی۔ 1910 تے 1913
دے وچکار رسل تے وائٹ ہیڈ دی
\emph{Principia Mathematica} دیاں تین جلدّاں
شائع ہوئیاں، جنہاں نے ریاضی نوں منطقی
بنیاداں اُتے کھڑا کرن دا فریگے والا
منصوبہ اگے ودھایا۔ پر رسل تے وائٹ ہیڈ
نوں دو ایہو جہے اصول مننّے پئے جنہاں نوں
خالص منطقی کہہ کے جائز ٹھہرانا اوکھا سی:
لامتناہیت دا اک مسلمہ تے «قابلِ تقلیل
ہون» دا اک مسلمہ۔ 1900 دی پہلی دہائی
وچ پوانکارے نے ریاضی وچ «پورے مجموعے
اُتے انحصار کرن والیاں تعریفاں»
\pnbLatin{(impredicative definitions)}
دی ورتوں اُتے تنقید کیتی۔ 1910 دی
دہائی وچ براوَر نے ساری ریاضی نوں
«وجدانیت پسند» بنیاد اُتے دوبارہ
کھڑا کرن دی تجویز پیش کرنی شروع
کیتی، جیہڑی تیسرے امکان دے اخراج
دا قاعدہ~($!A \lor \lnot !A$)
نہیں ورتدی۔

سچ مچ عجیب زمانہ سی!{} ساری ریاضی
نوں منطق وچ گھٹا دین والے منصوبے
نوں ہن «منطقیت» آکھیا جاندا اے؛
اُپر بیان کیتیاں اوکھاں کارن ایہنوں
عام طور اُتے ناکام سمجھیا جاندا اے۔
ریاضی نوں وجدانیت پسند ذہنی تعمیراں
راہیں بنانے والے منصوبے نوں «وجدانیت»
آکھدے نیں، تے عام ریاضی اُتے ایہدیاں
پابندیاں حد توں ودھ سخت سمجھیاں
جاندیاں نیں۔ صدی دے موڑ اُتے ڈیوڈ
ہلبرٹ، ہر ویلے دے سب توں اثر
رکھن والے ریاضی داناں وچوں اک،
کانٹور تے ڈیڈیکنڈ دے نویں تجریدی
طریقیاں دا پکا حامی سی: «کانٹور نے
ساڈے لئی جیہڑی جنت بنائی اے،
کوئی سانوں اوتھوں نہیں کڈھ سکے گا۔»
نال ہی اوہ ایہناں نویں طریقیاں
(جنہاں نوں ہن عجیب طور اُتے
«کلاسیکی» آکھیا جاندا اے) اُتے
بنیادی تنقید بارے حساس سی۔ اوہنے
دونیں فائدے اکٹھے رکھن دا اک راہ
سجھایا:
\begin{enumerate}
\item کلاسیکی طریقیاں نوں رسمی
  مسلّماں تے قاعدیاں نال بیان کرو؛
  ریاضیاتی سوالاں نوں مسلّمی
  نظام وچ !!{formula}s بنا کے پیش کرو۔
\item بھروسے جوگ «متناہی» طریقیاں نال
  ثابت کرو کہ ایہ رسمی اخذی نظام
  بے تضاد نیں۔
\end{enumerate}

پہلا مقصد حاصل کرن ول ہلبرٹ دا
کم بہت اگے ودھیا۔ 1899 وچ اوہنے
اپنی مشہور کتاب
\emph{ہندسہ دیاں بنیاداں \pnbLatin{(Foundations of geometry)}}
وچ ہندسہ لئی ایہ کم کیتا سی۔ بعد
دے سالاں وچ اوہ تے اوس دے کئی
شاگرداں تے ساتھیّاں نے ریاضی
دے ہور میداناں وچ وی ہندسہ
والا کم کرن دی کوشش کیتی۔
ہلبرٹ نے خود حساب تے تجزیہ
لئی مسلّماتی نظام دتے۔ زرمیلو
نے سیٹ نظریے دے مسلّمے بنائے،
جنہاں نوں بعد وچ فرینکل،
سکولم، فان نوئمان تے ہوراں
نے اگے ودھایا۔ 1920 دی
دہائی دے وچکار «ساری» ریاضی
نوں مسلّماتی شکل دین دا دعویٰ
کرن والے دو طریقے موجود سن:
رسل تے وائٹ ہیڈ دی
\emph{Principia mathematica}،
تے اوہ سیٹ نظریہ جیہڑا
زرمیلو--فرینکل سیٹ نظریے
دے ناں نال مشہور ہویا۔

1921 وچ ہلبرٹ نے ایہناں
نظاماں دی بے تضادی ثابت
کرن دا ہدف رکھ کے تحقیقی
منصوبہ شروع کیتا۔ ایس
کم وچ اوس دے کئی شاگرداں،
خاص کر برنائز، ایکرمان
تے بعد وچ گینٹزن، نے
اوہدی مدد کیتی۔ بنیادی
خیال ایہ سی کہ ریاضی وچ
تضاد دا !!{derivation} ہو سکّن
دا سوال علامتاں دیاں ممکن
لڑیاں بارے اک امتزاجی مسئلے
دی صورت وچ رکھیا جاوے۔ یعنی
جملیاں دیاں اوہ لڑیاں ویکھیاں
جان جیہڑیاں حساب،
تجزیے یا سیٹ نظریے دے
مسلّماتی نظام توں،
مثلاً، $!A \land
\lnot !A$ دے صحیح
!!{derivation} ہون دیاں
شرطاں پوریاں کردیاں
ہون۔ ایسی علامتی لڑی
ناممکن ہون دا ثبوت
خود اک ریاضیاتی ثبوت
ہووے گا، ایس لئی
ایہناں مسلّماتی نظاماں
وچ رسمی طور اُتے
بیان کیتا جا سکدا
ہووے گا۔ دوجے لفظاں
وچ، کوئی جملہ~$\OCon$
ہووے گا جیہڑا کہے
گا کہ، مثال لئی،
حساب بے تضاد اے۔
ہور ایہ کہ متعلقہ
نظاماں وچ ایہ جملہ
ثابت وی ہو سکنا
چاہیدا اے، خاص کر
جے اوس دا ثبوت
بہت محدود «متناہی»
وسیلے ہی منگدا ہووے۔

دوجا ہدف، یعنی بنائے
ہوئے مسلّماتی نظاماں
نال ہر ریاضیاتی سوال
دا فیصلہ ہونا، دو طرحاں
نال ٹھیک ٹھیک بیان
ہو سکدا اے۔ پہلا بیان
ایہ اے: ریاضی دے مسلّماتی
نظام دی زبان دے ہر
جملے~$!A$ لئی یا
تاں~$!A$ یا~$\lnot !A$
اوس دے مسلّماں توں
ثابت ہو سکدا ہووے۔
جے ایہ درست ہووے
تاں کوئی جملہ ایہو
جیہا نہ ہووے گا
جیہنوں انہاں مسلّماں
توں نہ ثابت کیتا
جا سکے نہ رد؛
ایہ مسلّمے ہر سوال
دا فیصلہ کردے
ہون گے۔ ایس خاصیت
والے مسلّماتی نظام
نوں \emph{کامل}
آکھدے نیں۔ بے شک
کسے دتے جملے لئی
ایہ طے کرنا ہن وی
اوکھا ہو سکدا اے
کہ دوناں وچوں
کیہڑا متبادل ثابت ہو سکدا اے۔
پر اصولی طور اُتے
ایہ معلوم کرن دا
کوئی طریقہ ہونا
چاہیدا اے۔ ہلبرٹ
دے زیرِ غور مسلّماتی
تے !!{derivation}
نظاماں لئی حقیقت
وچ کاملیت توں
ایہو جیہا طریقہ
موجود ہونا نکلدا
سی، بھاویں ہلبرٹ
نوں ایس دا احساس
نہیں سی۔ سوال
سمجھّن دا دوجا
طریقہ ایس توں
ودھ زورآور شرط
اے: کوئی مشینی،
حسابی طریقہ ہووے
جیہڑا دتے جملے~$!A$
لئی طے کرے کہ
اوہ مسلّماں توں
اخذ ہون جوگ اے
یا نہیں۔

% PNB-CMP-046: restrict Gödel's claims to effective arithmetic-rich consistent theories; own consistency is unprovable, not automatically independent.
1931 وچ گوڈل نے
«نامکملیت دے دو قضیے»
ثابت کیتے، جنہاں
نے دکھایا کہ ایہ
منصوبہ اپنی مطلوبہ
صورت وچ کامیاب
نہیں ہو سکدا۔
ریاضی دا ہر مسلّماتی
نظام نہیں بلکہ
ہر اوہ بے تضاد،
مؤثر طور اُتے
مسلّماں نال پیش
کیتا جا سکّن
والا تے حساب
لئی کافی طاقتور
نظام نامکمل ہوندا
اے۔ مناسب اخذی
شرطاں ہیٹھ ایسا
نظام اپنی بے
تضادی بیان کرن
والا جملہ آپ
ثابت نہیں کر
سکدا۔ اوس جملے
دی نفی وی ثابت
نہ ہو سکّن دا
دعویٰ وادھو
مفروضیاں اُتے
منحصر اے۔

ایہ نتیجہ ہلبرٹ دے
اصل منصوبے لئی
سخت ضرب سی۔ پر
جیویں ریاضی وچ
اکثر ہوندا اے،
ایہدے نال دلچسپ
نویں تحقیق دے
راہ وی کھُل گئے۔
جے ساری ریاضی
نوں سمیٹن والا
اکّو رسمی نظام
نہیں، تاں محدود
دائرے والے نظام
بنانے تے اوہناں
وچ کیہ ثابت ہو
سکدا اے، ایہ
پرکھّن دا کم
معنی رکھدا اے۔
ایہ وی فائدے دی گل اے
کہ انہاں نظاماں دی بے
تضادی ثابت کرن لئی
گھٹ پابندی والے طریقے
بنائیے تے ناپیئے کہ
اوہناں دی بے تضادی
ثابت کرنی کِنّی اوکھی اے۔
گوڈل نے دکھایا
کہ (تقریباً)
ہر رسمی نظام
کجھ سوالاں دا
فیصلہ نہیں
کردا؛ ایس
لئی کسے خاص
نظام دے
ایسے «دلچسپ»
غیر فیصلہ شدہ
سوال لبھنا،
تے ایہ جاننا
کہ اوہناں
دا فیصلہ کرن
لئی نظام
کِنّا طاقتور
ہونا لوڑیں دا
اے، اہم
اے۔ اج تک
منطق دان
ایہ سوال
ریاضی دی
اک نویں شاخ،
ثبوتاں دے
نظریے، وچ
کھوج رہے نیں۔
\endgroup
% /OLP-0276

% OLP-0277 content/incompleteness/introduction/definitions.tex
\begingroup
\def\olfilename{definitions}
\expandafter\def\csname ol@sectioncs\endcsname{section}

\olfileid{inc}{int}{def}

\olsection{تعریفاں}

ہلبرٹ دے منصوبے وچ ریاضی نوں رسمی شکل دین،
تے فیر ایس رسمی نظام دی بے تضادی تے کاملیت
دکھاون لئی، پہلا کم اک زبان، منطقی ڈھانچہ تے
مسلّماں دا نظام چُننا ہووے گا۔ ایتھے اسی فرض
کر لئیے کہ ریاضی نوں پہلے درجے دی زبان وچ
رسمی شکل دتی جا سکدی اے۔ مطلب ایہ کہ
کوئی ایہو جیہا سیٹ اے جیہدے وچ !!{constant}s،
!!{function}s تے !!{predicate}s ہون۔ ایہ نشان
پہلے درجے دی منطق دے رابطیاں تے مسوّراں نال
مل کے ریاضی دے دعوے بیان کر سکدے نیں۔ اکثر لوک
مندے نیں کہ ایڈی زبان موجود اے: سیٹ نظریے
دی زبان، جیہدے وچ $\in$ واحد غیر منطقی نشان
اے۔ پہلی نظر وچ ایڈی سادی زبان دے اِنّی
وڈی اظہار دی طاقت رکھن دا دعویٰ ناقابلِ
یقین لگدا اے؛ ایہ ثابت کرن لئی بہت کم
کرنا پیا کہ عملاً ساری ریاضی ایس بہت
محدود لفظیات وچ بیان ہو سکدی اے۔ گل
سادی رکھن لئی فی الحال بحث نوں حساب
تک محدود کریے، یعنی ریاضی دا اوہ
حصہ جیہڑا صرف قدرتی عدداں~$\Nat$ نال
سروکار رکھدا اے۔ حساب دے حقائق بیان
کرن دی فطری زبان~$\Lang L_A$ اے۔
$\Lang L_A$ وچ اک دو تھانواں
والا !!{predicate}~$<$، اک
!!{constant}~$\Obj 0$،
اک یک-مقامی !!{function}~$\prime$،
تے دو دو تھانواں والے
!!{function}s~$+$ تے~$\times$ نیں۔

\begin{defn}
!!{sentence}s نال بنیا سیٹ~$\Gamma$ اک
\emph{نظریہ} اے جے اوہ منطقی نتیجیاں
دے لحاظ نال بند ہووے، یعنی جے
$\Gamma = \Setabs{!A}{\Gamma \Entails
!A}$۔
\end{defn}

نظریہ متعین کرن دے دو سوکھے طریقے
نیں۔ اک ایہ کہ جیہڑے !!{sentence}s کسے
!!{structure} وچ سچ ہون، اوہناں دا سیٹ لے لئیے۔
مثال لئی $\Lang L_A$ دی
اوہ !!{structure} ویکھو جیہدے
!!{domain}~$\Nat$ اے تے
ہر غیر منطقی نشان دی تعبیر
اوہو اے جیہڑی توقع ہوندی اے۔

\begin{defn}\ollabel{def:standard-model}
\emph{حساب دا معیاری نمونہ}
اوہ !!{structure}~$\Struct
N$ اے جیہدی تعریف ایویں اے:
\begin{enumerate}
\item $\Domain N = \Nat$
\item $\Assign{\Obj 0}{N} = 0$
\item $\Assign{\Obj \prime}{N}(n) = n + 1$ ہر $n \in \Nat$ لئی
\item $\Assign{\Obj +}{N}(n, m) = n + m$ ہر $n, m \in \Nat$ لئی
\item $\Assign{\Obj \times}{N}(n, m) = n\cdot m$ ہر $n, m \in \Nat$ لئی
\item $\Assign{\Obj <}{N} = \Setabs{\tuple{n, m}}{n \in \Nat, m \in
  \Nat, n < m}$
\end{enumerate}
\end{defn}

'$\times$' تے '$\cdot$' دا فرق دھیان
نال ویکھو: $\times$ حساب دی زبان
دا اک نشان اے۔ بے شک اسی ایہ نشان
ضرب یاد کراؤن لئی چُنیا اے، پر
$\times$ خود ضرب دا عمل نہیں؛ ایہ
دو تھانواں والا فنکشن نشان اے
(باقاعدہ ناں~$\Obj f^2_1$)۔
ایہدے برعکس $\cdot$ عام ضرب
دا \emph{فنکشن اے}۔ جدوں
$n \cdot m$ ویکھو تاں مراد
عدداں~$n$ تے~$m$ دا حاصلِ
ضرب اے؛ جدوں $x \times y$
ویکھو تاں گل حساب دی زبان
دے اک ٹرم بارے اے۔ معیاری
نمونے وچ فنکشن نشان~$\times$
دی تعبیر قدرتی عدداں اُتے
فنکشن~$\cdot$ نال کیتی
جاندی اے۔ جمع لئی اسی $+$
نوں حساب دی زبان دے فنکشن
نشان تے قدرتی عدداں اُتے
جمع دے فنکشن، دوناں لئی
ورتدے آں۔ ایتھے سیاق
توں طے کرنا پیندا اے کہ
مراد کیہ اے۔

\begin{defn}
\emph{سچّے حساب} دا نظریہ
اوہ سیٹ اے جیہدے وچ اوہی !!{sentence}s
نیں جیہڑے حساب دے معیاری
نمونے وچ پورے اتردے نیں،
یعنی
\[
\Th{TA} = \Setabs{!A}{\Sat{N}{!A}}.
\]
\end{defn}

$\Th{TA}$ اک نظریہ اے:
جے $\Th{TA} \Entails !A$
ہووے تاں $!A$ ہر اُس
!!{structure} وچ پورا
اتردا اے جیہڑی~$\Th{TA}$
نوں پورا کردی اے۔
$\Sat{N}{\Th{TA}}$ ہون
کارن~$\Sat{N}{!A}$
وی اے، تے ایس لئی
$!A \in \Th{TA}$۔

نظریہ~$\Gamma$ متعین
کرن دا دوجا طریقہ ایہ
اے کہ !!{sentence}s نال
بنے کسے سیٹ~$\Gamma_0$
دے سارے منطقی نتیجیاں
دا سیٹ لے لئیے۔ ایس
صورت وچ $\Gamma$،
$\Gamma_0$ دی منطقی
نتیجیاں دے لحاظ نال
«بندش» اے۔ ایہ طریقہ
تاں ہی کارآمد اے جے
$\Gamma_0$ صاف طور اُتے
متعین ہووے، مثلاً
اوہ~$\Gamma_0$ دے !!{element}s دی
فہرست دتی گئی ہووے۔
گھٹ توں گھٹ~$\Gamma_0$
دے بارے فیصلہ ہو سکنا
چاہیدا اے: کسے
!!a{sentence} دے
$\Gamma_0$ دا رکن
ہون یا نہ ہون دی
قابلِ حساب جانچ ہووے۔
$\Gamma_0$ دے جیہڑے
!!{sentence}s نیں، اوہناں نوں
$\Gamma$ دے \emph{مسلّمے}
آکھدے آں، تے آکھدے
آں کہ $\Gamma$،
$\Gamma_0$ نال
\emph{مسلّماتی طور متعین}
اے۔

\begin{defn}
نظریہ~$\Gamma$،
$\Gamma_0$ نال
\emph{مسلّماتی طور متعین}
اے جے تے صرف جے
\[
\Gamma = \Setabs{!A}{\Gamma_0 \Entails !A}
\]
\end{defn}

\begin{defn}
اگلے جملیاں نال مسلّماتی
طور متعین نظریے~$\Th{Q}$
نوں «رابنسن دا~$\Th{Q}$»
آکھدے نیں؛ ایہ حساب
دا بہت سادہ نظریہ اے۔
\begin{align*}
& \lforall[x][\lforall[y][(\eq[x'][y'] \lif \eq[x][y])]] \tag{$!Q_1$}\\
& \lforall[x][\eq/[\Obj 0][x']] \tag{$!Q_2$}\\
& \lforall[x][(\eq[x][\Obj 0] \lor \lexists[y][\eq[x][y']])] \tag{$!Q_3$}\\
& \lforall[x][\eq[(x + \Obj 0)][x]] \tag{$!Q_4$}\\
& \lforall[x][\lforall[y][\eq[(x + y')][(x + y)']]] \tag{$!Q_5$}\\
& \lforall[x][\eq[(x \times \Obj 0)][\Obj 0]] \tag{$!Q_6$}\\
& \lforall[x][\lforall[y][\eq[(x \times y')][((x \times y) + x)]]] \tag{$!Q_7$}\\
& \lforall[x][\lforall[y][(x < y \liff \lexists[z][\eq[(z' + x)][y]])]] \tag{$!Q_8$}
\end{align*}
!!{sentence}s نال بنیا سیٹ
$\{!Q_1, \dots, !Q_8\}$
ہی $\Th{Q}$ دے مسلّمے
نیں؛ ایس لئی $\Th{Q}$
اوہناں توں نکلن والے
سارے !!{sentence}s
دا سیٹ اے:
\[
\Th{Q} = \Setabs{!A}{\{!Q_1, \dots, !Q_8\} \Entails !A}.
\]
\end{defn}

\begin{defn}
فرض کرو $!A(x)$،
$\Lang L_A$ دا
!!a{formula} اے جیہدے
آزاد متغیر~$x$ تے
$y_1$, \dots, $y_n$
نیں۔ فیر ایس شکل دا
ہر !!{sentence}
\[
\lforall[y_1][\dots\lforall[y_n][((!A(\Obj 0) \land \lforall[x][(!A(x)
\lif !A(x'))]) \lif \lforall[x][!A(x)])]]
\]
\emph{استقرائی خاکے}
دی اک مثال اے۔

\emph{پیانو حساب}~$\Th{PA}$
اوہ نظریہ اے جیہڑا
$\Th{Q}$ دے مسلّماں
تے استقرائی خاکے
دیاں ساریاں مثالاں
نال مسلّماتی طور
متعین ہوندا اے۔
\end{defn}

\begin{explain}
استقرائی خاکے دی ہر
مثال~$\Struct{N}$
وچ سچ اے۔ ایہ گل
سب توں سوکھی اوہدوں
ویکھی جا سکدی اے
جدوں !!{formula}~$!A$
دا صرف اک آزاد
!!{variable}~$x$
ہووے۔ فیر $!A(x)$
اک ذیلی سیٹ~$X_{!A}$
متعین کردا اے؛
ایہ $\Nat$ دا
ذیلی سیٹ اے،
تے تعریف
$\Struct{N}$ وچ اے۔
$X_{!A}$ وچ
اوہ سارے~$n \in \Nat$
نیں جنہاں لئی
$\Sat{N}{!A(x)}[s]$
ہووے، جدوں
$s(x) = n$۔
استقرائی
خاکے دی متعلقہ
مثال ایہ اے:
\[
((!A(\Obj 0) \land \lforall[x][(!A(x) \lif !A(x'))]) \lif 
  \lforall[x][!A(x)]).
\]
جے ایس دا مقدمہ
$\Struct{N}$ وچ
سچ ہووے تاں
$0 \in X_{!A}$
تے جدوں وی
$n \in X_{!A}$
ہووے، $n+1$
وی اوہدے وچ
ہووے گا۔
$0 \in X_{!A}$
توں~$1 \in X_{!A}$
ملدا اے؛
$1 \in X_{!A}$
توں~$2 \in X_{!A}$
ملدا اے؛ تے
ایویں اگے۔
ایس لئی ہر
$n \in \Nat$
لئی~$n \in X_{!A}$
اے۔ ایہو
گل اے کہ
$\lforall[x][!A(x)]$
، $\Struct{N}$
وچ پورا اتردا اے۔
\end{explain}

$\Th{Q}$ تے $\Th{PA}$
دونیں مسلّماتی طور
متعین نظریے نیں۔
وڈا سوال ایہ اے
کہ اوہ کِنّے طاقتور
نیں؟ مثال لئی،
کیہ $\Th{PA}$،
$\Nat$ بارے
$\Lang L_A$ وچ
بیان ہو سکّن
والیاں ساریاں
سچائیاں ثابت کر
سکدا اے؟ ٹھیک
ٹھیک پچھیے تاں
کیہ $\Th{PA}$
دے مسلّمے
$\Lang L_A$ وچ
پوچھے جا سکّن
والے ہر سوال
دا فیصلہ کردے
نیں؟

ایہی سوال ہور
طرح ایویں پچھ
سکدے آں: کیہ
$\Th{PA} = \Th{TA}$
اے؟ ظاہر اے کہ
$\Th{TA}$،
$\Nat$ بارے
ساریاں سچائیاں
ثابت کردا اے
(یعنی اوہناں
نوں شامل کردا
اے) تے
$\Lang L_A$
وچ بیان ہو
سکّن والے
سارے سوالاں
دا فیصلہ کردا
اے۔ وجہ ایہ
اے کہ جے
$!A$،
$\Lang L_A$
دا !!a{sentence}
ہووے، تاں یا
$\Sat{N}{!A}$
یا
$\Sat{N}{\lnot !A}$
؛ ایس لئی
یا
$\Th{TA} \Entails !A$
یا
$\Th{TA} \Entails \lnot !A$۔
ایہو جیہے
نظریے نوں
\emph{!!{complete}}
آکھدے نیں۔

\begin{defn}
نظریہ~$\Gamma$
\emph{!!{complete}}
اے جے تے صرف
جے اوہدی زبان
دے ہر
!!{sentence}~$!A$
لئی یا
$\Gamma \Entails !A$
یا
$\Gamma \Entails \lnot !A$
ہووے۔
\end{defn}

\begin{explain}
تمامیت دے قضیے
موجب،
$\Gamma \Entails !A$
جے تے صرف جے
$\Gamma \Proves
!A$؛ ایس لئی
$\Gamma$ کامل
اے جے تے
صرف جے اوہدی
زبان دے ہر
!!{sentence}~$!A$
لئی یا
$\Gamma \Proves !A$
یا
$\Gamma \Proves \lnot !A$
ہووے۔
\end{explain}

اک ہور سوال وی
اگے آندا اے:
کیہ کوئی حسابی
طریقہ اے جیہدے
نال اسی پرکھ
سکیے کہ کوئی
!!a{sentence}،
$\Th{TA}$ وچ
اے، $\Th{PA}$
وچ اے، یا
صرف~$\Th{Q}$
وچ اے؟
کسے سیٹ
(مثلاً
!!{sentence}s نال بنے سیٹ)
دے !!{decidable}
ہون دی تعریف
دے کے ایس
سوال نوں ہور
ٹھیک بنا
سکدے آں۔

\begin{defn}
سیٹ~$X$
\emph{!!{decidable}}
اے جے تے
صرف جے اک
حسابی طریقہ
ہووے جیہڑا
انپٹ~$x$
اُتے~$1$
واپس دیوے
جے $x \in X$
ہووے، تے نہیں
تاں~$0$ دیوے۔
\end{defn}

ہن سوال ایہ
بندا اے:
کیہ $\Th{TA}$
($\Th{PA}$،
$\Th{Q}$)
!!{decidable}
اے؟

ایہناں سارے
سوالاں دا
جواب «نہیں»
ہووے گا۔
ایہناں وچوں
کوئی وی نظریہ
فیصلہ کرن جوگ
نہیں۔ پر ایہ
واقعہ صرف ایہناں
خاص نظریاں
تک محدود نہیں۔
حقیقت وچ،
کجھ شرطاں
پوری کرن والا
\emph{ہر}
نظریہ ایہناں
نتیجیاں ہیٹھ
آندا اے۔ اک
شرط، جیہڑی
$\Th{Q}$ تے
$\Th{PA}$ پوری
کردے نیں،
ایہ اے کہ
اوہ مسلّماں
دے اک سیٹ نال
متعین نیں،
جیہڑا !!a{decidable}
اے۔

\begin{defn}
نظریہ
\emph{!!{axiomatizable}} اے جے اوہ
مسلّماں دے اک سیٹ نال
مسلّماتی طور متعین ہووے،
تے اوہ سیٹ !!a{decidable} ہووے۔
\end{defn}

\begin{ex}
جیہڑا نظریہ
!!{sentence}s نال بنے
محدود سیٹ
نال مسلّماتی
طور متعین اے،
اوہ
!!{axiomatizable}
اے؛ وجہ ایہ
کہ ہر محدود
سیٹ
!!{decidable}
ہوندا اے۔
ایس لئی، مثال
دے طور اُتے،
$\Th{Q}$
!!{axiomatizable}
اے۔

$\Th{PA}$
ورگے خاکیاں
راہیں متعین
نظریے وی
!!{axiomatizable}
نیں۔ ایہ
پرکھّن لئی
کہ $!B$،
$\Th{PA}$
دے مسلّماں
وچ شامل اے
یا نہیں، یعنی
اوہ فنکشن
$\Char{X}$
حساب کرن لئی
جیہدے وچ
$\Char{X}(!B) = 1$
ہووے جے
$!B$،
$\Th{PA}$
دا مسلمہ
اے؛ نہیں
تاں~$= 0$
ہووے،
ایہ طریقہ
ورت سکدے آں۔
پہلاں پرکھو
کہ $!B$،
$\Th{Q}$
دے مسلّماں
وچوں اک اے
یا نہیں۔
جے اے،
تاں جواب
«ہاں» تے
$\Char{X}(!B) = 1$
اے۔ جے
نہیں، تاں
پرکھو کہ
ایہ استقرائی
خاکے دی
اک مثال
اے یا نہیں۔
ایہ کم
باقاعدہ طریقے
نال ہو
سکدا اے؛
ایس معاملے
وچ شاید
ایویں ویکھنا
سب توں سوکھا
اے: استقرائی
خاکے دی
ہر مثال
چند کلی
مسوّراں توں
شروع ہوندی
اے، جنہاں
مگروں اک
شرطیے والا
ذیلی
!!{formula}
آندا اے۔
اُس شرطیے
دا تالی
$\lforall[x][!A(x, y_1, \dots, y_n)]$
اے۔ ایتھے
$x$ تے
$y_1$, \dots, $y_n$
سارے~$!A$
دے آزاد
متغیر نیں،
تے $!B$
دے شروع والے
مسوّرے
$y_1$, \dots,~$y_n$
نوں باندھدے
نیں۔ ایہ~$!A$
کڈھن تے ایہدے
آزاد متغیراں
نوں شروع دے کلی
مسوّراں تے
$\lforall[x]$
نال بندھے
متغیراں
نال ملاؤن
مگروں، اسی
پرکھدے آں
کہ شرطیے
دا مقدمہ
ایس شکل
نال ملدا
اے یا نہیں:
\[
!A(\Obj 0, y_1, \dots, y_n) \land \lforall[x][(!A(x, y_1, \dots, y_n)
\lif !A(x', y_1, \dots, y_n))]
\]
جے ملدا اے
تاں $!B$
استقرائی
خاکے دی
مثال اے؛
جے نہیں ملدا
تاں $!B$
ایہو جیہی مثال
نہیں۔
\end{ex}

ایس سوال،
تے زیادہ
عمومی سوال
کہ کیہڑے
نظریے
کامل یا
فیصلہ کرن
جوگ نیں،
دے جواب
لئی اگلی
تعریف وی
فائدہ مند
ہووے گی۔
یاد کرو
کہ سیٹ~$X$
!!{enumerable}
اے جے تے
صرف جے اوہ
خالی اے
یا کوئی
!!a{surjective}
فنکشن~$f \colon \Nat
\to X$ موجود
اے۔ ایسے
فنکشن نوں
$X$ دی
گنتی وار
فہرست آکھدے
نیں۔

\begin{defn}
سیٹ~$X$
\emph{!!{computably enumerable}}
(مختصراً
!!{c.e.}) اے
جے تے
صرف جے اوہ
خالی ہووے
یا اوہدی
قابلِ حساب
گنتی وار
فہرست
موجود ہووے۔
\end{defn}

!!{axiomatizability}
دے نال
اک ہور
شرط جیہڑی
نامکملیت
دے قضیاں
لئی متعلقہ
نظریاں اُتے
لاگو ہووے
گی، ایہ
اے کہ اوہ
قابلِ حساب
فنکشناں
تے
!!{decidable}
تعلقّاں بارے
بنیادی حقائق
ثابت کرن
جوگے ہون۔
«بنیادی
حقائق»
توں مراد
اوہ
!!{sentence}s
نیں جیہڑے
قابلِ حساب
فنکشناں
دے
مدخلاں
لئی قدراں
بیان کردے
نیں۔
«کافی طاقتور»
توں مراد
ایہ اے
کہ متعلقہ
نظریے
ایہ جملے
اپنے قضیاں
وچ شامل
کردے ہون۔
مثال لئی
اک بنیادی
قابلِ حساب
فنکشن، جمع،
نوں ویکھو۔
$+$ دے
مدخلاں
$2$ تے
$3$ لئی
قدر~$5$
اے، یعنی
$2+3 = 5$۔
حساب دی
زبان وچ
$+$ دے
مدخلاں
$2$ تے
$3$ لئی
قدر~$5$
ہون دا
بیان ایہ
جملہ کردا اے:
$(\num{2} + \num{3}) = \num{5}$۔
مثلاً
$\Th{Q}$
ایہ جملہ
ثابت کردا
اے۔ عام
طور اُتے
اسی چاہندے
آں کہ
ہر قابلِ
حساب فنکشن
$f(x_1, x_2)$
لئی کوئی
!!a{formula}~$!A_f(x_1, x_2, y)$
زبان~$\Lang{L_A}$
وچ ایویں ہووے کہ
$\Th{Q}
\Proves !A_f(\num{n_1}, \num{n_2}, \num{m})$
جدوں وی
$f(n_1, n_2) =
m$ ہووے۔
ایس طرح
$\Th{Q}$
ثابت کردا
اے کہ
$f$ دی
قدر مدخلاں
$n_1$،
$n_2$
لئی~$m$
اے۔ حقیقت
وچ اسی
ایہدے توں
کجھ ودھ
وی منگدے
آں: نظریہ
ایہ وی
ثابت کرے
کہ~$f$ دی
قدر مدخلاں
$n_1$،
$n_2$ لئی
کوئی
ہور عدد
نہیں ہو
سکدی۔
!!{decidable}
تعلقّاں
لئی وی
ایویں
اے۔ اگلی
دو تعریفاں
ایہ گل
ٹھیک
بناندیاں
نیں۔

\begin{defn}
!!{formula}~$!A(x_1, \dots, x_k, y)$
فنکشن~$f\colon \Nat^k \to \Nat$
دی نظریے~$\Gamma$ وچ
\emph{!!{represents}} کردا اے
جے تے
صرف جے،
جدوں وی
$f(n_1,
\dots, n_k) = m$
ہووے، تاں
\begin{enumerate}
\item $\Gamma \Proves !A(\num{n_1}, \dots, \num{n_k}, \num{m})$، تے
\item $\Gamma \Proves \lforall[y](!A(\num{n_1}, \dots, \num{n_k},
y) \lif y = \num{m})$۔
\end{enumerate}
\end{defn}

\begin{defn}
% PNB-CMP-047: source relation definition omits the theory Gamma used in its two conditions.
!!{formula}~$!A(x_1, \dots, x_k)$
نظریے~$\Gamma$ وچ
تعلق~$R \subseteq \Nat^k$
دی \emph{!!{represents}} کردا اے
جے تے
صرف جے،
\begin{enumerate}
\item جدوں وی $R(n_1, \dots, n_k)$ ہووے،
  $\Gamma \Proves !A(\num{n_1},
\dots, \num{n_k})$، تے
\item جدوں وی $R(n_1, \dots, n_k)$ نہ ہووے،
  $\Gamma \Proves \lnot
!A(\num{n_1}, \dots, \num{n_k})$۔
\end{enumerate}
\end{defn}

نامکملیت دے
قضیے لاگو
ہون لئی
نظریہ «کافی
طاقتور» اے
جے اوہ
سارے قابلِ
حساب فنکشناں
تے سارے
!!{decidable}
تعلقّاں دی
!!{represents}
کردا اے۔
$\Th{Q}$~
تے اوہدے
توسیعی نظریے
ایہ شرط پوری
کردے نیں،
پر ایس
گل دا ثبوت
دین لئی
سانوں
ہن کجھ ویلا
لگے گا۔
ایہ
$\Th{Q}$
دے ثبوتی
اختیار بارے
غیر معمولی
حقیقت اے؛
ایہ ثابت
کرنی اوکھی
اے کیوں
جے
$\Th{Q}$
دے مسلّمے
گھٹ نیں
تے سارے
حقائق اوہناں
توں اخذ
کرنے پینے
نیں۔ پر
$\Th{Q}$
خود بہت
کمزور نظریہ
اے۔ ایس لئی
بھاویں ایہ
ثابت کرنا
اوکھا اے
کہ
$\Th{Q}$
سارے قابلِ
حساب فنکشناں
دی نمائندگی
کردا اے،
اکثر دلچسپ
نظریے
$\Th{Q}$
توں طاقتور
نیں، یعنی
اوہ
$\Th{Q}$
توں ودھ
کجھ ثابت
کردے نیں۔
جے
$\Th{Q}$
کجھ ثابت
کردا اے
تاں ہر
زیادہ طاقتور
نظریہ وی
اوہ ثابت
کردا اے۔
ایس لئی
$\Th{Q}$
سارے قابلِ
حساب فنکشناں
دی نمائندگی
کردا ہووے
تاں ہر
زیادہ طاقتور
نظریہ وی
ایہ کردا
اے۔ ایہ
دکھاون نال
نامکملیت
دے قضیے
بہت سارے
دلچسپ نظریاں
اُتے لاگو
ہوندے نیں۔
پوری ریاضی
نوں رسمی
شکل دین
دا امیدوار
نظریہ یقیناً
$\Th{Q}$
دا ثابت
کیتا سب کجھ
ثابت کرے؛
گھٹ توں
گھٹ اوہنوں
ابتدائی
حساباں دے
نتیجے سنبھال
سکنے چاہیدے
نیں۔ سو
«ساری ریاضی»
دے نظریے
دا کوئی وی
امیدوار اوہناں
وچوں ہووے
گا جنہاں
اُتے نامکملیت
دے قضیے
لاگو ہوندے
نیں۔
\endgroup
% /OLP-0277

% OLP-0278 content/incompleteness/introduction/overview.tex
\begingroup
\def\olfilename{overview}
\expandafter\def\csname ol@sectioncs\endcsname{section}

\olfileid{inc}{int}{ovr}

\olsection{نامکملیت دے نتیجیاں دا اجمالی جائزہ}

ہلبرٹ نوں امید سی کہ ریاضی نوں اک ایہو جیہے !!{axiomatizable}
نظریے وچ رسمی شکل دتی جا سکدی اے جیہدا
!!{complete} تے !!{decidable} ہونا ثابت
کیتا جا سکے۔ اوہ ایہ وی چاہندا سی
کہ بہت کمزور، «متناہی» طریقیاں نال ایس نظریے
دی بے تضادی ثابت ہووے؛ ایس نال وجدانی مکتبِ
فکر دے اعتراضاں دے مقابلے وچ کلاسیکی ریاضی
دا دفاع ہووے گا۔ گوڈل دے نامکملیت دے قضیاں
نے دکھایا کہ ایہ سارے مقصد اکٹھے پورے نہیں
ہو سکدے۔

گوڈل دے پہلے نامکملیت قضیے نے دکھایا کہ
رسل تے وائٹ ہیڈ دی \emph{Principia Mathematica}
دا اک روپ !!{complete} نہیں۔ پر ثبوت خاص
کتاب تک محدود نہیں سی۔ اوہ بہت طرح دے
نظریاں اُتے لاگو ہون والی عام دلیل سی۔
ایس لئی مسئلہ صرف ایہ نہیں کہ
\emph{Principia Mathematica} ریاضی نوں پوری
طرح نہیں سمیٹ سکی، بلکہ ہلبرٹ دے مطلوبہ
خواص رکھن والا \emph{کوئی وی} نظریہ ایہ
کم نہیں کر سکدا۔ ایہ طے کرن وچ کجھ ویلا
لگیا کہ نظریے دیاں کیہڑیاں خاصیتاں نامکملیت
دے قضیاں لئی کافی نیں، تے گوڈل دی دلیل
نوں اینج عام بنایا جائے کہ اوہ صرف ایہناں
خاصیتاں اُتے منحصر ہووے۔ ہن اسی حساب دی
زبان~$\Lang L_A$ وچ نظریاں لئی پہلے
نامکملیت قضیے دا اک بہت عام بیان دے
سکدے آں۔

\begin{thm}
جے $\Gamma$ زبان~$\Lang L_A$ دا اک بے تضاد
تے !!{axiomatizable} نظریہ اے جیہڑا سارے
قابلِ حساب فنکشناں تے سارے !!{decidable}
تعلقّاں دی !!{represents} کردا اے، تاں
$\Gamma$ !!{complete} نہیں۔
\end{thm}

% PNB-CMP-048: move the closing parenthesis after Gamma outside inline math.
% PNB-CMP-049: distinguish ordinary Goedel self-unprovability from Rosser's consistency-only independent witness.
$\Gamma$ دے !!{complete} نہ ہون دا مطلب
ایہ اے کہ گھٹ توں گھٹ اک !!{sentence}~$!A$
ایہو جیہا اے جیہدے لئی $\Gamma \Proves/ !A$ تے
$\Gamma \Proves/ \lnot !A$، دونیں گلاں سچ نیں۔
ایہو جیہے !!a{sentence} نوں نظریے
(یعنی~$\Gamma$) دے لحاظ نال \emph{آزاد}
کہندے نیں۔ آزاد جملیاں دا ہونا اسی نسبتاً
چھیتی ثابت کر سکدے آں۔ گوڈل دی دلیل دی
خاص طاقت ایہ اے کہ اوہ اک \emph{متعین مثال}
بنا کے دیندی اے؛ پر اتھے دو دعویاں نوں
وکھ رکھنا ضروری اے۔ اک !!a{sentence}~$!G_\Gamma$
بنایا جاندا اے، جیہنوں~$\Gamma$ لئی
\emph{گوڈل جملہ} کہندے نیں۔ ایہ
ثابت نہیں ہو سکدا، کیوں جے
$\Gamma$ وچ~$!G_\Gamma$ ایس دعوے دے
برابر اے کہ~$!G_\Gamma$، نظریے
$\Gamma$ وچ ثابت نہیں ہو سکدا۔
صرف بے تضادی توں ایہ نکلدا اے کہ
ایہ جملہ ثابت نہیں ہو سکدا۔
ایس عام گوڈل جملے دی نفی وی ثابت نہ ہون
لئی ہور مضبوط شرط، مثلاً اومیگا-بے تضادی،
چاہیدی اے۔ صرف بے تضادی دی شرط ہیٹھ
اک متعین آزاد !!{sentence} روسر دی بدلی ہوئی
بناوٹ توں ملدا اے۔ گوڈل جملے دی بناوٹ
\emph{تعمیری} اے۔ جے~$\Gamma$ دے مسلّماں
دا بیان تے !!{derivation} دے نظام دی تفصیل
دتی ہووے، ثبوت~$!G_\Gamma$ نوں حقیقت وچ
لکھن دا طریقہ دیندا اے۔

گوڈل دی بناوٹ لئی ضروری اے کہ
$\Lang L_A$ دے اپنے ٹرم تے
!!{formula}s جیہڑیاں خاصیتاں رکھدے نیں، تے
اوہناں اُتے جیہڑے عمل ہوندے نیں، اوہ~$\Lang L_A$ وچ بیان ہو سکن۔
مثلاً «$!A$ اک !!a{sentence} اے»،
«$\delta$،~$!A$ دا اک !!a{derivation} اے»،
تے $\Subst{!A}{t}{x}$ ورگا عمل۔
ایہ طریقہ، اول، نشاناں تے اوہناں دیاں
لڑیواں نوں قدرتی عدداں دے کوڈاں وچ
بدل کے ایہ خاصیتاں تے تعلق بیان کرے۔
ٹرم، !!{formula}s تے !!{derivation}s
ایہو جیہیاں نشان-لڑیاں نیں، جد کہ
$\Lang L_A$ قدرتی عدداں بارے گل
کر سکدی اے۔ دوجا، کوڈ ایویں بنے کہ
$\Gamma$ متعلقہ حقائق ثابت کر سکے۔
ایس لئی دکھانا پینا اے کہ کوڈ شدہ
خاصیتاں قدرتی عدداں دیاں
!!{decidable} خاصیتاں نیں تے
کوڈاں اُتے عمل، قدرتی عدداں اُتے
قابلِ حساب فنکشناں نال ملدے نیں۔
ایہنوں «نحو نوں عددی روپ دینا»
\pnbLatin{(arithmetization of syntax)} کہندے نیں۔

پر نحو نوں عددی روپ دین دی پڑتال
توں پہلے اسی ایہ شرط ویکھاں گے کہ
نظریہ~$\Gamma$ «کافی طاقتور» ہووے،
یعنی سارے قابلِ حساب فنکشناں تے
!!{decidable} تعلقّاں دی
!!{represents} کردا ہووے۔ ایس لئی
«قابلِ حساب» دی ٹھیک تعریف دینی
پینی اے۔ ایہدے کئی طریقے نیں، مثلاً
ٹیورنگ مشین دے نمونے نال، یا کسے
عام مقصد والی پروگرامنگ زبان دے
پروگراماں راہیں حساب کیتے جان والے
فنکشناں نال۔ پر ساڈا مقصد ایہناں
فنکشناں تے تعلقّاں دی زبان~$\Lang L_A$
دے نظریے وچ !!{represents} کرنا اے،
ایس لئی صرف عددی فنکشناں دی
قابلِ حسابیت دی سادی تعریف چُننا
چنگا اے۔ ایہ \emph{بازگشتی فنکشن}
دا تصور اے۔ سو پہلے اسی بازگشتی
فنکشناں نوں ویکھاں گے۔ فیر دکھاواں
گے کہ $\Th{Q}$ پہلے ہی سارے
بازگشتی فنکشناں تے تعلقّاں دی
!!{represents} کردا اے۔ ایس توں
نامکملیت دا قضیہ $\Th{Q}$ تے
$\Th{PA}$ ورگے مخصوص نظریاں
اُتے لاگو کر سکاں گے، کیوں جے
ایہ «کافی طاقتور» نظریاں دیاں
مثالاں ثابت ہو چکیاں ہون گیاں۔
% PNB-CMP-050: remove the frozen source's malformed glossary plural suffix.

نحو نوں عددی روپ دین دا آخری
نتیجہ اک !!a{formula}~$\OProv[\Gamma](x)$
اے۔ ایہ !!{formula}s نوں عدداں
نال کوڈ کر کے~$\Gamma$ دے مسلّماں
توں ثبوت ہون دی گل بیان کردا اے۔
خاص طور تے، جے~$!A$ دا کوڈ
عدد~$n$ اے تے $\Gamma \Proves
!A$، تاں $\Gamma \Proves
\Prov[\Gamma](\num{n})$۔
$\Gamma$ لئی ایہ «ثبوت پذیری
دا محمول» اک خاص مفہوم وچ
$\Gamma$ دی بے تضادی نوں وی
$\Lang L_A$ دے اک !!a{sentence}
وچ بیان کرن دی سہولت دیندا اے:
$\Gamma$ دا «بے تضادی دا جملہ»
!!{sentence}~$\lnot \Prov[\Gamma](\num{n})$
ہووے، جتھے~$n$ کسے تضاد،
مثلاً~$\lfalse$، دا کوڈ اے۔
% PNB-CMP-051: the source's unconditional second-theorem sentence is qualified by its following sentence.
دوجا نامکملیت قضیہ کہندا اے
کہ بے تضاد !!{axiomatizable}
نظریے، جدوں ایس قضیے دیاں
ہور مطلوبہ شرطاں وی پوری کرن،
اپنی بے تضادی دے جملے دا ثبوت
نہیں دے سکدے۔ ایہ شرطاں
صرف سارے قابلِ حساب فنکشناں
تے !!{decidable} تعلقّاں دی
نمائندگی توں کجھ ودھ سخت
نیں؛ اسی دکھاواں گے کہ
$\Th{PA}$ اوہناں نوں پورا
کردا اے۔
\endgroup
% /OLP-0278

% OLP-0279 content/incompleteness/introduction/undecidability.tex
\begingroup
\def\olfilename{undecidability}
\expandafter\def\csname ol@sectioncs\endcsname{section}

\olfileid{inc}{int}{dec}

\olsection{فیصلہ نہ ہو سکنا تے نامکملیت}

گوڈل دے نامکملیت قضیاں دے ثبوت لئی
نحو نوں عددی روپ دینا ضروری اے۔
پر ایہدے بغیر وی، صرف ایہ فرض کر
کے کہ اک نظریہ سارے !!{decidable}
تعلقّاں دی !!{represents} کردا اے،
کجھ اہم نتیجے حاصل ہو سکدے نیں۔
ایہ قطری دلیل رُکّن دے مسئلے دا
فیصلہ ناممکن ہون والی دلیل ورگی اے۔

\begin{thm}
جے $\Gamma$ اک بے تضاد نظریہ اے
جیہڑا ہر !!{decidable} تعلق دی
!!{represents} کردا اے، تاں
$\Gamma$ !!{decidable} نہیں۔
\end{thm}

\begin{proof}
فرض کرو $\Gamma$ !!{decidable}
ہووے۔ اسی دکھاواں گے کہ جے
$\Gamma$ ہر !!{decidable}
تعلق دی !!{represents} کردا اے،
تاں اوہ متضاد ہووے گا۔

% PNB-CMP-052: restore the indexed formula in the two diagonal membership-test fragments.
!!^{decidable} خاصیتاں، یعنی
یک-مقامی تعلق، اک آزاد
متغیر والے !!{formula}s نال
ظاہر کیتیاں جاندیاں نیں۔ فرض
کرو $!A_0(x)$، $!A_1(x)$، \dots،
ایہو جیہے سارے !!{formula}s ایس ترتیب نال
اک قابلِ حساب گنتی وار فہرست وچ آندے نیں۔
ہن ایہ سیٹ $D \subseteq \Nat$
ویکھو:
\[
D = \Setabs{n}{\Gamma \Proves \lnot !A_n(\num{n})}
\]
سیٹ~$D$ !!{decidable} اے، کیوں جے
اسی $n \in D$ جانچن لئی پہلے
$!A_n(x)$ نکال سکدے آں تے
اوہدے توں $\lnot !A_n(\num{n})$
بنا سکدے آں۔ ٹرم~$\num{n}$
نال $x$ دی ہر آزاد آمد
بدلنی، جیہڑی $!A_n(x)$ وچ
اے، تے
$!A_n(\num{n})$ دے اگے
$\lnot$ لانا صاف مشینی کم اے۔
فرض مطابق $\Gamma$ !!{decidable}
اے؛ سو اسی جانچ سکدے آں کہ
$\lnot !A_n(\num{n}) \in \Gamma$
اے یا نہیں۔ جے اے تاں
$n \in D$، تے جے نہیں تاں
$n \notin D$۔ ایس طرح
$D$ وی !!{decidable} اے۔

% PNB-CMP-053: normalize four malformed verbal glossary suffixes in this unit.
کیوں جے $\Gamma$ ساریاں
!!{decidable} خاصیتاں دی
!!{represents} کردا اے،
اوہ~$D$ دی وی
!!{represents} کردا اے۔
$D$ دی !!{represents} نظریے~$\Gamma$ وچ کرن والے
سارے !!{formula}s
فہرست $!A_0(x)$، $!A_1(x)$،
\dots، وچ شامل نیں۔ سو ایہو جیہا
عدد~$d$ چنو کہ $!A_d(x)$،
$D$ دی، نظریے~$\Gamma$ وچ،
!!{represents}
کردا ہووے۔ جے
$d \notin D$، تاں
$!A_d(x)$ ولوں~$D$
دی !!{represents} توں
$\Gamma \Proves \lnot !A_d(\num{d})$
ملدا اے۔ پر ایس دا مطلب
ایہ اے کہ $d$،
سیٹ~$D$ دی تعریف والی شرط
پوری کردا اے؛ یعنی
$d \in D$۔ ایہ
$d \notin D$ دے خلاف اے۔
ایس لئی برہانِ خلف نال
$d \in D$۔

جدوں $d \in D$ اے، تاں
سیٹ~$D$ دی تعریف توں
$\Gamma \Proves \lnot
!A_d(\num{d})$۔ دوجی
طرف، کیوں جے $!A_d(x)$،
$D$ دی، نظریے~$\Gamma$ وچ،
!!{represents} کردا اے،
$\Gamma \Proves !A_d(\num{d})$
وی اے۔ ایس لئی
$\Gamma$ متضاد اے۔
\end{proof}

\begin{explain}
پچھلا قضیہ دکھاندا اے کہ کوئی
بے تضاد نظریہ جیہڑا سارے
!!{decidable} تعلقّاں دی
!!{represents} کردا اے،
!!{decidable} نہیں ہو سکدا۔
اسی دکھاواں گے کہ
$\Th{Q}$ واقعی سارے
!!{decidable} تعلقّاں دی
!!{represents} کردا اے۔
ایس لئی $\Th{Q}$
نوں شامل کرن والے سارے
نظریے، مثلاً $\Th{PA}$
تے $\Th{TA}$، وی
ایہناں تعلقّاں دی نمائندگی
کردے نیں، تے ایس واسطے
اوہ وی !!{decidable}
نہیں۔ (ایہ سارے نظریے
معیاری نمونے وچ سچ نیں،
ایس لئی بے تضاد نیں۔)

ایس نتیجے توں پہلے
نامکملیت قضیے دا اک کمزور
روپ وی ملدا اے۔ ہر نظریہ
جیہڑا !!{axiomatizable}
تے !!{complete} اے،
!!{decidable} وی اے۔
ایس لئی بے تضاد نظریے
جیہڑے !!{axiomatizable}
نیں تے ساریاں
!!{decidable} خاصیتاں
دی !!{represents} کردے
نیں، !!{complete} نہیں
ہو سکدے۔
\end{explain}

\begin{thm}
جے $\Gamma$ !!{axiomatizable}
تے !!{complete} اے، تاں
اوہ !!{decidable} اے۔
\end{thm}

\begin{proof}
ہر متضاد نظریہ
!!{decidable} اے، کیوں جے
ایہو جیہے نظریے وچ سارے
!!{sentence}s ہوندے نیں۔
ایس لئی سوال «کیہ
$!A \in \Gamma$ اے؟»
دا جواب ہمیشہ «ہاں»
ہوندا اے، یعنی فیصلہ
کیتا جا سکدا اے۔

ہن فرض کرو $\Gamma$
بے تضاد اے، تے
!!{axiomatizable} تے
!!{complete} وی اے۔
!!{axiomatizable} ہون
دی وجہ توں $\Gamma$
!!{computably enumerable}
اے۔ نظریے~$\Gamma$ دے مسلّماں توں
جیہڑے درست !!{derivation}s نکلدے نیں،
اوہناں دی فہرست اک قابلِ حساب فنکشن
دے سکدا اے۔
ہر درست
!!{derivation} توں اوہ
!!{sentence} نکال سکدے
آں جیہڑا اوہ
!!{derive}s کردا اے؛
سو اک قابلِ حساب فنکشن
$\Gamma$ دے سارے
قضیاں دی فہرست بنا
دیندا اے۔ اک
!!{sentence}،
نظریے~$\Gamma$
دا قضیہ اُدوں تے صرف
اُدوں اے جدوں $\lnot !A$
قضیہ نہیں؛ ایہ
$\Gamma$ دی بے
تضادی تے
!!{complete} ہون
توں نکلدا اے۔
ہن $!A \in \Gamma$
دا فیصلہ ایویں کرو:
$\Gamma$ دے سارے
قضیے فہرست وچ آندے
جاؤن۔ جدوں $!A$
فہرست وچ آوے،
پتہ لگ جاندا اے
کہ $\Gamma \Proves !A$۔
جدوں $\lnot
!A$ آوے،
پتہ لگ جاندا اے
کہ $\Gamma \Proves/ !A$۔
$\Gamma$ دے
!!{complete} ہون
دی وجہ توں ایہناں
دوناں وچوں اک صورت
ضرور آوے گی؛
سو طریقہ آخرکار
جواب دے دیندا اے۔
\end{proof}

\begin{cor}
\ollabel{cor:incompleteness}
جے $\Gamma$ بے تضاد
تے !!{axiomatizable}
اے، تے ہر
!!{decidable}
خاصیت دی
!!{represents}
کردا اے، تاں
اوہ !!{complete}
نہیں۔
\end{cor}

\begin{proof}
جے $\Gamma$
!!{complete} ہوندا،
پچھلے قضیے توں
!!{decidable} ہوندا،
کیوں جے اوہ
!!{axiomatizable}
تے بے تضاد اے۔
پر $\Gamma$
ہر !!{decidable}
خاصیت دی
!!{represents}
کردا اے؛ پہلے
قضیے توں اوہ
!!{decidable}
نہیں۔ ایہ تضاد اے۔
\end{proof}

\begin{prob}
دکھاؤ کہ
$\Th{TA} = \Setabs{!A}{\Sat{N}{!A}}$
!!{axiomatizable}
نہیں۔ تُسی
ایہ فرض کر
سکدے او کہ
$\Th{TA}$
ساریاں فیصلہ
کرن جوگ
خاصیتاں دی
نمائندگی کردا اے۔
\end{prob}

جدوں اسی ثابت
کر لواں گے
کہ مثلاً
$\Th{Q}$
ساریاں
!!{decidable}
خاصیتاں دی
!!{represents}
کردا اے،
تاں نتیجہ
دسّے گا کہ
$\Th{Q}$
لازماً نامکمل
اے۔ پر ایس
دلیل وچ کسے
آزاد !!{sentence}
دی مثال نہیں
ملدی؛ صرف
ایہ ثابت
ہوندا اے کہ
ایہو جیہا !!a{sentence}
موجود اے۔
ایہو جیہی مثال
بناؤن لئی
نحو نوں عددی
روپ دینا تے
گوڈل دے
اصل ثبوت
دا خیال
استعمال کرنا
پینا اے۔
تے ظاہر اے
کہ سانوں
پہلا دعویٰ وی
ثابت کرنا
اے: $\Th{Q}$
واقعی ساریاں
!!{decidable}
خاصیتاں دی
!!{represents}
کردا اے۔

ایہ وی یاد
رکھنا چاہیدا
اے کہ ہر
\emph{دلچسپ}
نظریہ نامکمل
یا ناقابلِ
فیصلہ نہیں۔
بہت نظریے
دلچسپ ریاضیاتی
حقائق بیان
کرن جوگے
نیں، پر گوڈل
دے نتیجے
دیاں مطلوبہ
شرطاں پوری
نہیں کردے۔
مثلاً
$\Th{Pres} = \Setabs{!A \in \Lang{L_{A^+}}}{\Sat{N}{!A}}$
اوہ سیٹ اے جیہدے وچ حساب دی
اوہ زبان جیہدے وچ~$\times$ نہیں،
دے اوہی !!{sentence}s نیں
جیہڑے معیاری نمونے وچ سچ نیں۔
ایہ نظریہ
کامل وی اے
تے فیصلہ کرن
جوگ وی۔
ایہنوں
پریسبورگر
حساب کہندے
نیں؛ ایہ
قدرتی عدداں
بارے اوہ
سارے سچ
ثابت کردا اے
جیہڑے صرف
$\Obj{0}$،
$\prime$ تے~$+$
نال بیان
ہو سکدے نیں۔
\endgroup
% /OLP-0279

\OLEndChapterHook
\endgroup
% /OLP-0275

% OLP-0280 content/incompleteness/arithmetization-syntax/arithmetization-syntax.tex
\begingroup
\def\olfilename{arithmetization-syntax}
\expandafter\def\csname ol@sectioncs\endcsname{section}

\olchapter{inc}{art}{نحو نوں عددی روپ دینا}

\begin{editorial}
  دھیان رہے کہ صداقتی قدراں دیاں علامتاں والے تابلواں لئی
  نحو نوں عددی روپ دین دا طریقہ ہن تک دستیاب نہیں۔
\end{editorial}

% OLP-0281 content/incompleteness/arithmetization-syntax/introduction.tex
\begingroup
\def\olfilename{introduction}
\expandafter\def\csname ol@sectioncs\endcsname{section}

\olfileid{inc}{art}{int}
\olsection{تعارف}

قابلِ حسابیت تے منطق نوں جوڑن لئی سانوں ایسا طریقہ چاہیدا اے
جیہدے نال منطق دیاں چیزاں (علامتاں، ٹرم، !!{formula}s تے
!!{derivation}s)، اوہناں اُتے ہوندے عمل، تے اوہناں دیاں خاصیتاں
تے تعلقاں نوں حسابی طریقیاں نال ورت سکّیے۔ ایہ کم سیدھا وی ہو سکدا
اے، جے اسیں علامتاں، علامتاں دے سلسلیاں تے اوہناں توں بنیاں ہور
چیزاں اُتے قابلِ حساب فنکشن تے تعلق ویکھیے۔ منطقی نحو دی ہر شے
متناہی اے تے علامتاں دے !!a{enumerable} مجموعیاں توں بندی اے،
ایس لئی حساب دے کجھ نمونیاں وچ ایہ طریقہ چل سکدا اے۔ پر حساب دے
ہور نمونے، جیویں بازگشتی فنکشن، صرف عدداں، عدداں وچلے تعلقاں
تے عدداں اُتے فنکشناں تک محدود نیں۔ آخرکار اسیں کجھ نظرییاں دے
اندر وی نحو نال کم کرنا چاہندے آں، خاص کر اوہناں نظرییاں وچ
جیہڑے حساب دی زبان وچ لکھے گئے نیں۔ اوتھے نحو نوں \emph{عددی
روپ دینا} ضروری اے۔ یعنی نحوی چیزاں نوں عدداں دے روپ وچ، اوہناں
اُتے ہوندے عملاں نوں حسابی فنکشناں دے روپ وچ تے اوہناں دے تعلقاں
نوں حسابی تعلقاں دے روپ وچ پیش کرنا پیندا اے۔ نحوی چیزاں نوں
عدد دین دا ایہ خیال لائبنز دے ویلے توں چلیا آؤندا اے۔

علامتاں نوں «کوڈ» وجوں عدد دینا نسبتاً سوکھا اے۔ کجھ علامتاں
لئی ذرا اوکھ اے، کیوں جے، مثلاً، !!{variable}s لامتناہی نیں تے
دلیلاں دی ہر مقررہ تعداد~$n$ والے !!{function}s وی لامتناہی نیں۔
فیر وی علامتاں نوں ترتیب نال عدد دِتے جا سکدے نیں، جیویں
$\Obj v_2$ تے $\Obj v_3$ نوں وکھو وکھ کوڈ۔ علامتاں دے سلسلیاں
(جیویں ٹرماں تے !!{formula}s) نوں کوڈ دینا ہور اوکھا اے۔ پر جے
عدداں دے سلسلیاں نال خالص حسابی طریقے نال کم ہو سکدا اے
(مثلاً اولی عدداں دیاں قوتاں نال سلسلیاں نوں کوڈ دے کے)، تاں
اسیں اک اک علامت دی کوڈکاری توں اگے ودھ کے علامتاں دے سلسلیاں
نوں وی کوڈ دے سکدے آں۔ ایس توں اگے !!{formula}s دے سلسلیاں یا
ہور طرح دیاں بناوٹاں، جیویں !!{derivation}s، نوں وی کوڈ دے سکدے
آں۔ ایس وادھو کوڈکاری نوں «گوڈل نمبر دینا» (Gödel numbering)
آکھدے آں۔ ہر ٹرم، !!{formula} تے !!{derivation} نوں اک گوڈل
نمبر دِتا جاندا اے۔

جے اسیں علامتاں دے سلسلیاں دی تھاں اوہناں دیاں علامتاں دے
کوڈاں دے سلسلے لئیے، تے سلسلیاں نوں کوڈ دین دا ایسا طریقہ چنیے
جیہڑے نال قابلِ حساب فنکشناں راہیں کم ہو سکے، تاں گوڈل نمبراں
نال وی قابلِ حساب فنکشناں راہیں کم ہو سکدا اے۔ عمل وچ سارے
متعلقہ فنکشن بنیادی بازگشتی ہون گے۔ مثلاً سلسلے دے کوڈ توں
اوہدی لمبائی لبھن والا فنکشن تے اوہدا $i$-واں رکن کڈھن والا
فنکشن، دونوں بنیادی بازگشتی نیں۔ جے سلسلے دا کوڈ، مثلاً،
!!a{formula}~$!A$ دا گوڈل نمبر ہووے، تاں ایس توں فوراً پتا لگدا
اے کہ !!a{formula} دی لمبائی تے !!a{formula} دے $i$-ویں نشان
دا کوڈ وی $!A$ دے گوڈل نمبر توں کڈھیا جا سکدا اے۔ ایہ ثابت
کرنا ذرا اوکھا اے کہ، مثلاً، صحیح طرح بنے ٹرم یا درست
!!{derivation} دا گوڈل نمبر ہون دی خاصیت بنیادی بازگشتی اے۔
فیر وی ایہ ثابت ہو سکدا اے، کیوں جے متعلقہ سلسلیاں (ٹرم،
!!{formula}s تے !!{derivation}s) دیاں تعریفاں استقرائی نیں۔

مثال لئی بدل دا عمل ویکھو۔ جے $!A$ اک فارمولا، $x$ اک متغیر
تے $t$ اک ٹرم ہووے، تاں $\Subst{!A}{t}{x}$ اوہ نتیجہ اے جیہڑا
$x$ دیاں~$!A$ وچ ساریاں آزاد آمدواں دی تھاں $t$ رکھ کے ملدا اے۔
ہن فرض کرو کہ $!A$، $x$ تے $t$ نوں ایس ترتیب نال گوڈل نمبر
$k$، $l$ تے $m$ دِتے گئے نیں۔ ایہو کوڈکاری $\Subst{!A}{t}{x}$
نوں وی اک گوڈل نمبر دیندی اے، من لو~$n$۔ $k$، $l$ تے $m$ نوں
$n$ ول لیجان والا ایہ فنکشن بدل دے عمل دا حسابی روپ اے۔ جدوں
بدل دا عمل $!A$، $x$ تے $t$ توں $\Subst{!A}{t}{x}$ دیندا اے،
تاں بدل دا عددی روپ گوڈل نمبراں $k$، $l$ تے $m$ توں گوڈل
نمبر~$n$ دیندا اے۔ اگے اسیں ویکھاں گے کہ ایہ فنکشن بنیادی
بازگشتی اے۔

نحو نوں عددی روپ دینا صرف نظری دلچسپی دی گل نہیں۔ البتہ شروع
وچ ایہ سمجھنا سوکھا نہیں سی کہ حساب دی زبان ورگیاں زباناں وی
عبارتاں دیاں پیچیدہ خاصیتاں نوں رسمی روپ دے سکدیاں نیں، حالانکہ
اوہناں وچ زباناں بارے «گل کرن» دا کوئی اپنا وسیلہ نہیں۔ ایس توں
اگلا سوال ایہ اے کہ ایس زبان وچ لکھیا کوئی نظریہ، جیویں پیانو
حساب، اپنی زبان بارے کیہ \emph{ثابت} کر سکدا اے (مثلاً، کیہ
!!{sentence}s ثابت کیتے جا سکدے نیں، یا کیہ اوہ سچ نیں)۔ ایس
سوال توں گوڈل تے تارسکی دے مشہور قضیاں تک پہنچدے آں، جیہڑے
ثبوت تے تعریف دیاں حداں دسدے نیں۔ گوڈل دے قضیے ثابت نہ ہو
سکّن بارے نیں تے تارسکی دا قضیہ ایس گل بارے اے کہ سچ دی تعریف
نہیں دِتّی جا سکدی۔ نحو نوں عددی روپ دین دا طریقہ قابلِ حسابیت
دے نظریے دے کجھ اہم نتیجے ثابت کرن لئی وی کم آؤندا اے، جیویں
نظرییاں دی حسابی طاقت یا حساب دے وکھو وکھ نمونیاں دے آپسی تعلق
بارے نتیجے۔ ایہ طریقہ ہور چیزاں تے خاصیتاں نوں عددی روپ دین دا
نمونہ وی اے۔ مثلاً ٹیورنگ مشیناں دیاں مکمل حالتاں تے اوہناں
دیاں حسابی کارروائیاں نوں وی ایویں عددی روپ دِتا جا سکدا اے۔
ایس نال اک نمونے (مثلاً ٹیورنگ مشیناں) دیاں حسابی کارروائیاں
دوجے نمونے (مثلاً بازگشتی فنکشناں) وچ نقل کر کے چلائیاں جا
سکدیاں نیں۔
\endgroup
% /OLP-0281

% OLP-0282 content/incompleteness/arithmetization-syntax/coding-symbols.tex
\begingroup
\def\olfilename{coding-symbols}
\expandafter\def\csname ol@sectioncs\endcsname{section}

\olfileid{inc}{art}{cod}
\olsection{علامتاں نوں کوڈ دینا}

پہلے درجے دی منطق دی بنیادی زبان~$\Lang L$ ایہ علامتاں ورتدی اے:
\[
\lfalse \quad \lnot \quad \lor \quad \land \quad \lif \quad \lforall
\quad \lexists \quad \eq \quad ( \quad ) \quad ,
\]
نال ایہ متغیراں تے !!{constant}s دے !!{enumerable} مجموعے،
تے دلیلاں دی ہر ممکن تعداد لئی !!{function}s تے !!{predicate}s
دے !!{enumerable} مجموعے وی ورتدی اے۔ اسیں ہر علامت نوں
\emph{کوڈ} وجوں اک عدد ایس طرح دے سکدے آں کہ ہر علامت دا
اپنا وکھرا کوڈ ہووے تے دو وکھریاں علامتاں دا کوڈ اکّو نہ ہووے۔
ایہ ممکن اے، کیوں جے ساریاں علامتاں دا مجموعہ !!{enumerable}
اے، تے ایس مجموعے تے قدرتی عدداں دے مجموعے وچ !!a{bijection}
اے۔ پر اسیں ایہ وی یقینی بنانا چاہندے آں کہ کوڈ توں علامت تے
اوہدے بارے کجھ معلومات (مثلاً !!a{function} دیاں دلیلاں دی
تعداد) قابلِ حساب طریقے نال واپس لبھ سکّیے۔ ایہ کرن دے کئی
طریقے نیں۔ ایتھے اسیں اک طریقہ دیندے آں جیہڑا بنیادی بازگشتی
فنکشن ورتدا اے۔ (یاد رکھو کہ $\tuple{n_0, \dots, n_k}$ اوہ عدد
اے جیہڑا $n_0$, \dots, $n_k$ دے سلسلے دا کوڈ اے۔)

\begin{defn}
جے $s$ زبان~$\Lang L$ دی اک علامت ہووے، تاں
اوہدا \emph{علامتی کوڈ}~$\scode s$ اینج مقرر کردے آں:
\begin{enumerate}
\item جے $s$ منطقی علامتاں وچوں ہووے، تاں
  $\scode s$ دی قدر ایہ جدول دندا اے:
\[
\begin{array}{cccccccccc}
  \lfalse & \lnot & \lor & \land & \lif & \lforall \\
\tuple{0, 0} & \tuple{0, 1} & \tuple{0, 2} & \tuple{0, 3} &
\tuple{0, 4} & \tuple{0, 5} \\
\lexists & \eq & ( & ) & ,\\
\tuple{0, 6} & \tuple{0, 7} &
\tuple{0, 8} & \tuple{0, 9} & \tuple{0, 10}
\end{array}
\]
\item جے $s$، $i$-واں متغیر $\Obj v_i$ ہووے،
  تاں $\scode s = \tuple{1, i}$۔
\item جے $s$، $i$-واں !!{constant}~$\Obj c_i$
  ہووے، تاں $\scode s = \tuple{2, i}$۔
\item جے $s$، $i$-واں $n$-دلیلی
  !!{function}~$\Obj f_i^n$ ہووے، تاں
  $\scode s = \tuple{3, n, i}$۔
\item جے $s$، $i$-واں $n$-دلیلی
  !!{predicate}~$\Obj P_i^n$ ہووے، تاں
  $\scode s = \tuple{4, n, i}$۔
\end{enumerate}
\end{defn}

\begin{prop}
ایہ تعلق بنیادی بازگشتی نیں:
\begin{enumerate}
\item $\fn{Fn}(x, n)$ اُدوں تے صرف اُدوں جدوں $x$،
  $\Obj f^n_i$ دا کوڈ ہووے، کسے~$i$ لئی۔ سو $x$ اک $n$-دلیلی !!{function} دا کوڈ اے۔
\item $\fn{Pred}(x, n)$ اُدوں تے صرف اُدوں جدوں $x$،
  $\Obj P^n_i$ دا کوڈ ہووے، کسے~$i$ لئی، یا $x$، $\eq$ دا کوڈ ہووے تے $n = 2$۔ سو $x$ اک
  $n$-دلیلی !!{predicate} دا کوڈ اے۔
\end{enumerate}
\end{prop}

\begin{defn}
جے $s_0, \dots, s_{n-1}$ علامتاں دا سلسلہ
ہووے، تاں اوہدا \emph{گوڈل نمبر}
$\tuple{\scode{s_0}, \dots, \scode{s_{n-1}}}$ اے۔
\end{defn}

\begin{explain}
دھیان رکھو کہ \emph{کوڈ} تے \emph{گوڈل نمبر} وکھ وکھ
چیزاں نیں۔ مثلاً متغیر~$\Obj v_5$ دا کوڈ~$\scode{\Obj
  v_5} = \tuple{1, 5} = 2^2\cdot 3^6$ اے۔ پر متغیر~$\Obj v_5$
نوں ٹرم وجوں ویکھیے تاں اوہ علامتاں دا اک سلسلہ وی اے، جیہدی
لمبائی~$1$ اے۔ ایس \emph{گوڈل نمبر}~$\Gn{\Obj v_5}$ دی قدر،
جدوں \emph{ٹرم}~$\Obj v_5$ دی گل ہووے، ایہ اے:
$\tuple{\scode{\Obj v_5}} = 2^{\scode{\Obj
    v_5} + 1} = 2^{2^2\cdot 3^6 + 1}$۔
\end{explain}

\begin{ex}
یاد کرو کہ جے $k_0$, \dots, $k_{n-1}$
عدداں دا سلسلہ ہووے، تاں اولی عدداں دیاں قوتاں
والی کوڈکاری وچ سلسلے دا کوڈ
$\tuple{k_0, \dots, k_{n-1}}$ ایہ اے:
\[
2^{k_0+1}\cdot3^{k_1+1}\cdot \dots \cdot p_{n-1}^{k_{n-1}+1},
\]
ایتھے $p_i$، $i$-واں اولی عدد اے
(گنتی $p_0 = 2$ توں شروع ہوندی اے)۔
سو مثلاً فارمولا $\eq[\Obj v_0][\Obj 0]$،
یا ہور کھول کے ${\eq}(\Obj v_0,\Obj c_0)$،
دا گوڈل نمبر ایہ اے:
\[
\tuple{\scode{\eq},\scode{(},\scode{\Obj v_0},\scode{,}, \scode{\Obj
    c_0},\scode{)}}.
\]
ایتھے $\scode{\eq}$، $\tuple{0,7} = 2^{0+1}\cdot
3^{7+1}$ اے، $\scode{\Obj v_0}$،
$\tuple{1,0} = 2^{1+1}\cdot3^{0+1}$ اے،
تے ایداں ہی اگے۔ ایس لئی $\Gn{=(\Obj v_0,\Obj c_0)}$
ایہ بندا اے:
\begin{multline*}
2^{\scode{=} + 1}\cdot 3^{\scode{(}+1}\cdot 5^{\scode{\Obj v_0}+1}
\cdot 7^{\scode{,} + 1} \cdot 11^{\scode{\Obj c_0}+1} \cdot
13^{\scode{)}+1} = \\
2^{2^1\cdot 3^8 + 1}\cdot 3^{2^1\cdot 3^9+1}\cdot 5^{2^2\cdot 3^1+1}
\cdot 7^{2^1\cdot 3^{11} + 1} \cdot 11^{2^3\cdot3^1+1} \cdot
13^{2^1\cdot3^{10}+1} = \\
2^{13\,123}\cdot 3^{39\,367}\cdot 5^{13}\cdot 7^{354\,295}\cdot11^{25}\cdot13^{118\,099}.
\end{multline*}
\end{ex}
\endgroup
% /OLP-0282

% OLP-0283 content/incompleteness/arithmetization-syntax/coding-terms.tex
\begingroup
\def\olfilename{coding-terms}
\expandafter\def\csname ol@sectioncs\endcsname{section}

\olfileid{inc}{art}{trm}
\olsection{ٹرماں نوں کوڈ دینا}

\begin{explain}
ٹرم علامتاں دا اک خاص قسم دا سلسلہ اے۔ ایہ ٹرماں دی تشکیل
دے قاعدیاں مطابق ثابتہ نشاناں تے متغیراں توں استقرائی ڈھنگ
نال بندا اے۔ علامتاں دے سلسلیاں نوں عدداں دے روپ وچ کوڈ کیتا
جا سکدا اے۔ ایس لئی علامتاں نوں کوڈ دین دا طریقہ تے عدداں دے
سلسلیاں نوں کوڈ دین دا طریقہ اکٹھے ورتدے آں۔ ایس کرکے ٹرماں
نوں گوڈل نمبر دینا اوکھا نہیں۔ اصل اوکھ ایہ وکھاؤن وچ اے
کہ کسے عدد دا صحیح طرح بنے ٹرم دا گوڈل نمبر ہون والی خاصیت
قابلِ حساب اے، بلکہ بنیادی بازگشتی اے۔
\end{explain}

!!^{variable}s تے !!{constant}s سبھ توں سادے ٹرم نیں۔ ایہ پرکھنا
سوکھا اے کہ $x$ ایسے کسے ٹرم دا گوڈل نمبر اے یا نہیں۔
$\fn{Var}(x)$ اُدوں قائم اے جدوں $x$، $\Gn{\Obj v_i}$ ہووے،
کسے~$i$ لئی۔ دوجے لفظاں وچ، $x$ لمبائی~$1$ والے سلسلے دا کوڈ
اے تے اوس سلسلے دا اکّو رکن $(x)_0$ کسے !!{variable}~$\Obj v_i$
دا کوڈ اے۔ یعنی $x$، $\tuple{\tuple{1, i}}$ ہووے، کسے~$i$ لئی۔
ایویں $\fn{Const}(x)$ اُدوں قائم اے جدوں $x$، $\Gn{\Obj c_i}$
ہووے، کسے~$i$ لئی۔ ایہ دونوں تعلق بنیادی بازگشتی نیں، کیوں جے
ایسا کوئی $i$ موجود ہووے تاں اوہ لازماً $< x$ اے:
\begin{align*}
  \fn{Var}(x) & \defiff \bexists{i<x}{x = \tuple{\tuple{1, i}}}\\
  \fn{Const}(x) & \defiff \bexists{i<x}{x = \tuple{\tuple{2, i}}}
\end{align*}

\begin{prop}
\ollabel{prop:term-primrec}
تعلق $\fn{Term}(x)$ تے $\fn{ClTerm}(x)$ بنیادی بازگشتی نیں۔
ایہ بالترتیب اُدوں تے صرف اُدوں قائم ہوندے نیں جدوں $x$
کسے ٹرم یا بند ٹرم دا گوڈل نمبر ہووے۔
\end{prop}

\begin{proof}
علامتاں دا سلسلہ~$s$ اُدوں تے صرف اُدوں ٹرم اے جدوں ٹرماں دا
کوئی سلسلہ~$s_0$، \dots، $s_{k-1} = s$ موجود ہووے جیہڑا وکھاوے
کہ ٹرماں دی تشکیل دے قاعدیاں مطابق !!{constant}s تے !!{variable}s
توں ٹرم~$s$ کِویں بنیا۔ تشکیل دا کوئی دِتّا سلسلہ قاعدیاں مطابق ہون لئی ہر $i < k$ لئی
ایہناں وچوں کوئی اک شرط پوری ہونی چاہیدی اے:
\begin{enumerate}
\item $s_i$ کوئی !!a{variable}~$\Obj v_j$ ہووے، یا
\item $s_i$ کوئی !!a{constant}~$\Obj c_j$ ہووے، یا
\item $s_i$، $n$ ٹرماں $t_1$، \dots، $t_n$ توں بنیا ہووے۔
  ایہ ٹرم تھاں~$i$ توں پہلاں آئے ہون تے انہاں اُتے اک
  $n$-دلیلی !!{function}~$\Obj f^n_j$ لایا گیا ہووے۔
\end{enumerate}
ایہ وکھاؤن لئی کہ گوڈل نمبراں اُتے متعلقہ تعلق بنیادی بازگشتی
اے، سانوں ایس شرط نوں بنیادی بازگشتی فنکشناں، تعلقاں تے
محدود مسوّریت نال بیان کرنا پینا اے۔

فرض کرو $y$ اوہ عدد اے جیہڑا سلسلے $s_0$، \dots، $s_{k-1}$
دا کوڈ اے، یعنی $y = \tuple{\Gn{s_0}, \dots, \Gn{s_{k-1}}}$۔
ایہ گوڈل نمبر~$x$ والے ٹرم دی تشکیل دے سلسلے دا کوڈ اُدوں
تے صرف اُدوں اے جدوں ہر $i < k$ لئی:
\begin{enumerate}
\item $\fn{Var}((y)_i)$ قائم ہووے، یا
\item $\fn{Const}((y)_i)$ قائم ہووے، یا
\item کوئی $n$ تے کوئی عدد~$z = \tuple{z_1, \dots, z_n}$ ایسے
  ہون کہ ہر $z_l$ کسے $(y)_{i'}$ دے برابر ہووے، جتھے
  $i' < i$، تے
\[
(y)_i = \Gn{\Obj f^n_j(} \concat \fn{flatten}(z) \concat \Gn{)},
\]
\end{enumerate}
نال ایہ وی ضروری اے کہ $(y)_{k-1} = x$ ہووے۔
(فنکشن $\fn{flatten}(z)$ سلسلے $\tuple{\Gn{t_1}, \dots, \Gn{t_n}}$
نوں $\Gn{t_1, \dots, t_n}$ وچ بدلدا اے تے بنیادی بازگشتی اے۔)

شق (3) وچ اشاریے $j$، $n$، گوڈل نمبر $z_l$ جیہڑے ٹرماں $t_l$
دے نیں، تے کوڈ~$z$ جیہڑا سلسلے~$\tuple{z_1, \dots, z_n}$ دا اے،
ایہ سارے~$y$ توں چھوٹے نیں۔ اُتلے $k$ دی تھاں
$\len{y}$ رکھ سکدے آں۔ ایس لئی «$y$، گوڈل نمبر~$x$ والے ٹرم دی
تشکیل دے سلسلے دا کوڈ اے» والی شرط ایس طرح بیان کر سکدے آں کہ
تعلق دا بنیادی بازگشتی ہونا ظاہر ہو جاوے۔

ہن ایہ وکھانا رہ گیا اے کہ $y$ لئی بنیادی بازگشتی حد موجود اے۔
جے $x$ کسے ٹرم دا گوڈل نمبر اے، تاں اوہدی تشکیل دا ایسا
سلسلہ ضرور ہووے گا جیہدے وچ ودھ توں ودھ $\len{x}$ ٹرم ہون گے۔
ایس لئی کہ~$s$ دی تشکیل دے سلسلے دا ہر ٹرم~$s$ وچ کسے تھاں
توں شروع ہوندا اے تے دو ذیلی ٹرم اکّو تھاں توں شروع نہیں ہو سکدے۔
ظاہر اے کہ $s$ دے ہر ذیلی ٹرم دا گوڈل نمبر $\le x$ اے۔ ایس لئی ہمیشہ
ایسا تشکیل دا سلسلہ موجود اے جیہدا کوڈ
$\le p_{k-1}^{k(x+1)}$ ہووے، جتھے $k=\len{x}$ اے۔

$\fn{ClTerm}$ لئی !!{variable}s والی شق چھڈ دیو۔
\end{proof}

\begin{prob}
وکھاؤ کہ فنکشن $\fn{flatten}(z)$، جیہڑا سلسلے
$\tuple{\Gn{t_1}, \dots, \Gn{t_n}}$ نوں $\Gn{t_1, \dots, t_n}$
وچ بدلدا اے، بنیادی بازگشتی اے۔
\end{prob}

\begin{prop}\ollabel{prop:num-primrec}
  فنکشن $\fn{num}(n) = \Gn{\num{n}}$
  بنیادی بازگشتی اے۔
\end{prop}

\begin{proof}
  اسیں $\fn{num}(n)$ دی تعریف بنیادی بازگشت نال کردے آں:
  \begin{align*}
    \fn{num}(0) & = \Gn{\Obj 0}\\
    \fn{num}(n+1) & = \Gn{\prime(} \concat \fn{num}(n) \concat \Gn{)}.
  \end{align*}
\end{proof}
\endgroup
% /OLP-0283

% OLP-0284 content/incompleteness/arithmetization-syntax/coding-formulas.tex
\begingroup
\def\olfilename{coding-formulas}
\expandafter\def\csname ol@sectioncs\endcsname{section}

\olfileid{inc}{art}{frm}
\olsection{\printtoken{P}{formula} نوں کوڈ دینا}

تعلق $\fn{Term}(x)$ دی تعریف بنیادی بازگشتی ڈھنگ نال ہو جان
توں بعد اسیں اوہنوں ورت کے !!{formula}s لئی متعلقہ تعلق
$\fn{Frm}(x)$ دی تعریف وی بنیادی بازگشتی ڈھنگ نال کر سکدے آں۔

\begin{prop}
تعلق $\fn{Atom}(x)$ اُدوں تے صرف اُدوں قائم اے جدوں $x$ کسے
ایٹمی !!{formula} دا گوڈل نمبر ہووے۔ ایہ تعلق بنیادی بازگشتی اے۔
\end{prop}

\begin{proof}
عدد $x$ کسے ایٹمی !!{formula} دا گوڈل نمبر اُدوں تے صرف اُدوں
اے جدوں ایہناں وچوں کوئی اک شرط پوری ہووے:
\begin{enumerate}
\item ایسے $n$، $j < x$ تے $z < x$ موجود ہون کہ ہر $i < n$
  لئی $\fn{Term}((z)_i)$ قائم ہووے تے $x = $
\[
\Gn{\Obj P^n_j(} \concat \fn{flatten}(z) \concat \Gn{)}.
\]
\item ایسے $z_1, z_2 < x$ موجود ہون کہ $\fn{Term}(z_1)$
  تے $\fn{Term}(z_2)$ دونوں قائم ہون، تے $x = $
\[
\Gn{{\eq}(} \concat z_1 \concat \Gn{,} \concat z_2 \concat{\Gn{)}}.
\]
  \tagitem{prvFalse}{$x = \Gn{\lfalse}$.}{}
  \tagitem{prvTrue}{$x = \Gn{\ltrue}$.}{}
\end{enumerate}
\end{proof}

\begin{prop}
\ollabel{prop:frm-primrec}
تعلق $\fn{Frm}(x)$ اُدوں تے صرف اُدوں قائم اے جدوں $x$
!!a{formula} دا گوڈل نمبر ہووے۔ ایہ تعلق بنیادی بازگشتی اے۔
\end{prop}

\begin{proof}
علامتاں دا سلسلہ~$s$ اُدوں تے صرف اُدوں !!a{formula} اے جدوں
!!{formula} دی تشکیل دا کوئی سلسلہ~$s_0$، \dots، $s_{k-1} = s$
موجود ہووے جیہڑا وکھاوے کہ تشکیل دے قاعدیاں مطابق ایٹمی
!!{formula}s توں~$s$ کِویں بنیا۔ ہر $s_i$ دا کوڈ (تے پورے
سلسلے $\tuple{s_0, \dots, s_{k-1}}$ دا کوڈ وی) کوڈ~$x$ توں
چھوٹا اے، جیہڑا~$s$ دا کوڈ اے۔
\end{proof}

\begin{prob}
\olref[inc][art][trm]{prop:term-primrec} دے پہلے ثبوت دی طرز اُتے
\olref[inc][art][frm]{prop:frm-primrec} دا تفصیلی ثبوت دیو۔
\end{prob}

\begin{prop}
\ollabel{prop:freeocc-primrec}
تعلق $\fn{FreeOcc}(x, z, i)$ اُدوں تے صرف اُدوں قائم اے جدوں
فارمولے دی $i$-ویں علامت، جتھے اوس فارمولے دا گوڈل نمبر~$x$
اے، گوڈل نمبر~$z$ والے متغیر دی آزاد آمد ہووے۔
ایہ تعلق بنیادی بازگشتی اے۔
\end{prop}

\begin{proof}
مشق۔
\end{proof}

\begin{prob}
\olref[inc][art][frm]{prop:freeocc-primrec} ثابت کرو۔ تسیں ایہ
گل ورت سکدے او کہ !!a{formula} دی ہر اوہ لگاتار ذیلی لڑی
جیہڑی آپ !!a{formula} ہووے، اوس دی ذیلی !!{formula} اے۔
\end{prob}

\begin{prop}
خاصیت $\fn{Sent}(x)$ اُدوں تے صرف اُدوں قائم اے جدوں $x$
!!a{sentence} دا گوڈل نمبر ہووے۔ ایہ خاصیت بنیادی بازگشتی اے۔
\end{prop}

\begin{proof}
!!{sentence} اوہ !!a{formula} اے جیہدے وچ !!{variable}s دی
کوئی آزاد آمد نہیں ہوندی۔ ایس لئی $\fn{Sent}(x)$ اُدوں تے
صرف اُدوں قائم اے جدوں
\begin{multline*}
\bforall{i<\len{x}}{\bforall{z<x}{}}\\
(\bexists{j<z}{z=\Gn{\Obj v_j}} \lif \lnot\fn{FreeOcc}(x,z,i)).
\end{multline*}
\end{proof}
\endgroup
% /OLP-0284

% OLP-0285 content/incompleteness/arithmetization-syntax/substitution.tex
\begingroup
\def\olfilename{substitution}
\expandafter\def\csname ol@sectioncs\endcsname{section}

\olfileid{inc}{art}{sub}
\olsection{تھاں بدل}

یاد رکھو کہ تھانواں بدلنا اوہ عمل اے جیہدے وچ
!!a{variable}~$u$ دیاں !!a{formula}~$!A$ وچ ساریاں آزاد
آمدواں دی تھاں ٹرم~$t$ رکھدے آں۔ ایہنوں $\Subst{!A}{t}{u}$
لکھدے آں۔ ایہ عمل !!{variable}s، !!{formula}s تے ٹرماں
دے گوڈل نمبراں اُتے کیتا جاوے، تاں بنیادی بازگشتی اے۔

\begin{prop}\ollabel{prop:subst-primrec}
اک بنیادی بازگشتی فنکشن $\fn{Subst}(x, y, z)$ موجود اے
جیہدے لئی ایہ مساوات قائم اے:
\[
\fn{Subst}(\Gn{!A}, \Gn{t}, \Gn{u}) = \Gn{\Subst{!A}{t}{u}}.
\]
\end{prop}

\begin{proof}
فیر فنکشن $\fn{hSubst}$ دی تعریف بنیادی بازگشت نال
ایویں کر سکدے آں:
\begin{multline*}
\begin{aligned}
\fn{hSubst}(x, y, z, 0) & = \emptyseq \\
\fn{hSubst}(x, y, z, i+1) & =
\end{aligned}\\
\begin{cases}
\fn{hSubst}(x, y, z, i) \concat y & \text{جے $\fn{FreeOcc}(x, z, i)$ ہووے} \\
\fn{append}(\fn{hSubst}(x, y, z, i), (x)_{i}) & \text{نہیں تاں۔}
\end{cases}
\end{multline*}
ہن $\fn{Subst}(x, y, z)$ دی تعریف
$\fn{hSubst}(x, y, z, \len{x})$ دے برابر رکھ کے کر سکدے آں۔
\end{proof}

\begin{prop}
\ollabel{prop:free-for}
تعلق $\fn{FreeFor}(x, y, z)$ اُدوں تے صرف اُدوں قائم
اے جدوں گوڈل نمبر~$y$ والا ٹرم، گوڈل نمبر~$z$ والے متغیر
لئی، گوڈل نمبر~$x$ والے فارمولے وچ !!{free for} ہووے۔
ایہ تعلق بنیادی بازگشتی اے۔
\end{prop}

\begin{proof} مشق۔ \end{proof}

\begin{prob}
\olref[inc][art][sub]{prop:free-for}
ثابت کرو۔
\end{prob}
\endgroup
% /OLP-0285

% OLP-0286 content/incompleteness/arithmetization-syntax/proofs-in-lk.tex
\begingroup
\def\olfilename{proofs-in-lk}
\expandafter\def\csname ol@sectioncs\endcsname{section}

\olfileid{inc}{art}{plk}
\olsection{\texorpdfstring{$\Log{LK}$ وچ \usetoken{P}{derivation}}{LK وچ اخذ}}

\begin{explain}
!!{derivation}s نوں عددی روپ دین لئی سانوں !!{derivation}s نوں
عدداں نال ظاہر کرنا پینا اے۔ !!{derivation}s سیکوئنٹاں دے درخت نیں
تے ہر استنباط دا اک لیبل وی اے۔ ایس لئی بازگشتی طریقہ سِدھا اے۔
اسیں ہر !!{derivation} نوں اک ٹپل نال ظاہر کردے آں، جیہدے جز
آخرلا سیکوئنٹ، استنباط دا لیبل تے آخری استنباط دے مقدمیاں تک
پہنچن والے ذیلی !!{derivation}s دے عددی روپ نیں۔
\end{explain}

\begin{defn}
جے $\Gamma$، !!{sentence}s دا متناہی سلسلہ اے، $\Gamma =
\tuple{!A_1, \dots, !A_n}$، تاں
$\Gn{\Gamma} = \tuple{\Gn{!A_1},
  \dots, \Gn{!A_n}}$ اے۔

جے $\Gamma \Sequent \Delta$
اک سیکوئنٹ اے، تاں
$\Gamma \Sequent \Delta$ دا
گوڈل نمبر ایہ اے:
\[
\Gn{\Gamma \Sequent \Delta} = \tuple{\Gn{\Gamma}, \Gn{\Delta}}
\]

جے $\pi$، $\Log{LK}$ وچ
!!a{derivation} اے، تاں
$\Gn{\pi}$ دی تعریف ایویں کردے آں:
\begin{enumerate}
\item جے $\pi$ صرف ابتدائی
  سیکوئنٹ $\Gamma \Sequent \Delta$
  اُتے مشتمل ہووے، تاں
  $\Gn{\pi}$ ایہ اے:
  \[
    \tuple{0, \Gn{\Gamma \Sequent \Delta}}.
  \]
\item جے $\pi$ دا آخرلا استنباط
  اک یا دو مقدمیاں والا ہووے،
  اوہدا نتیجہ
  $\Gamma \Sequent \Delta$
  ہووے، تے $\pi_1$ تے
  $\pi_2$ اوہ براہِ راست ذیلی ثبوت ہون جیہڑے آخری استنباط دے
  مقدمیاں اُتے ختم ہوندے نیں، تاں $\Gn{\pi}$ دی قدر بالترتیب ایہ اے:
  \begin{align*}
    & \tuple{1, \Gn{\pi_1}, \Gn{\Gamma \Sequent \Delta}, k} \text{ یا}\\
    & \tuple{2, \Gn{\pi_1}, \Gn{\pi_2}, \Gn{\Gamma \Sequent \Delta}, k},
  \end{align*}
  ایتھے $k$ دی قدر، آخری استنباط وچ ورتے قاعدے مطابق،
  ہِٹھلے جدول توں ملدی اے:

  \begin{tabular}{lccccccc}
    \text{قاعدہ:} & \LeftR{\Weakening} & \RightR{\Weakening} &
    \LeftR{\Contraction} & \RightR{\Contraction} &
    \LeftR{\Exchange} & \RightR{\Exchange} \\
    $k$: & 1 & 2 & 3 & 4 & 5 & 6 \\[2ex]
    \text{قاعدہ:} & \LeftR{\lnot} & \RightR{\lnot} &
    \LeftR{\land} & \RightR{\land} &
    \LeftR{\lor} & \RightR{\lor} \\
    $k$: & 7 & 8 & 9 & 10 & 11 & 12 \\[2ex]
    \text{قاعدہ:} & \LeftR{\lif} & \RightR{\lif} &
    \LeftR{\lforall} & \RightR{\lforall} &
    \LeftR{\lexists} & \RightR{\lexists} \\
    $k$: & 13 & 14 & 15 & 16 & 17 & 18 \\[2ex]
    \text{قاعدہ:} & \Cut & = \\
    $k$: & 19 & 20
  \end{tabular}
\end{enumerate}
\end{defn}

\begin{ex}
ایہ سوکھا جیہا !!{derivation} ویکھو:
\begin{prooftree}
  \Axiom$!A \fCenter !A$
  \RightLabel{\LeftR{\land}}
  \UnaryInf$!A \land !B \fCenter !A$
  \RightLabel{\RightR{\lif}}
  \UnaryInf$\fCenter (!A \land !B) \lif !A$
\end{prooftree}
ابتدائی سیکوئنٹ دا گوڈل نمبر
$p_0 = \tuple{0,
  \Gn{!A \Sequent !A}}$
ہوئے گا۔ $\LeftR{\land}$ دے
نتیجے اُتے ختم ہون والے
!!{derivation} دا گوڈل نمبر
$p_1 =
    \tuple{1, p_0, \Gn{!A \land !B \Sequent !A}, 9}$
ہووے گا ($1$ ایس لئی کہ
$\LeftR{\land}$ دا اک مقدمہ اے؛
نتیجے~$!A \land !B \Sequent !A$
دا گوڈل نمبر؛ تے $9$،
$\LeftR{\land}$ نوں کوڈ کرن
والا نمبر اے)۔ فیر پورے
!!{derivation} دا گوڈل نمبر
$\tuple{1, p_1, \Gn{\Sequent (!A \land !B) \lif !A)}, 14}$
اے؛ یعنی
\[
\tuple{1, \tuple{1, \tuple{0, \Gn{!A \Sequent !A)}}, \Gn{!A \land !B \Sequent !A}, 9},
  \Gn{\Sequent (!A \land !B) \lif !A}, 14}.
\]
\end{ex}

\begin{explain}
!!{derivation}s دی عددی روپ طے کر کے سانوں ایہ وی دکھانا
پینا اے کہ اسیں ایہناں
!!{derivation}s اُتے بنیادی بازگشتی
ڈھنگ نال عمل کر سکدے آں تے انہاں
دیاں ضروری خاصیتاں تے تعلقاں
نوں وی ایس طرح بیان کر سکدے آں۔
کجھ عمل سادے نیں: مثال لئی
!!a{derivation} دے گوڈل
نمبر~$p$ توں
$\fn{EndSeq}(p) = (p)_{(p)_0+1}$
اوس دے آخرلے سیکوئنٹ دا
گوڈل نمبر دیندا اے تے
$\fn{LastRule}(p) = (p)_{(p)_0+2}$
اوس دے آخری قاعدے دا
کوڈ دیندا اے۔ خاصیت
$\fn{Sequent}(s)$، جیہدی
تعریف ایہ اے،
\[
  \len{s} = 2 \land \bforall{i<\len{(s)_0} + \len{(s)_1}}{\fn{Sent}(((s)_0 \concat (s)_1)_i)}
\]
$s$ لئی عین اوس ویلے
سچ اے جدوں $s$،
!!{sentence}s توں بنے
سیکوئنٹ دا گوڈل نمبر
ہووے۔ کجھ ہور عمل کِتے ودھ اوکھے نیں۔ ایتھے اسیں گھٹ توں گھٹ اوہناں دا خاکہ دساں گے۔
مقصد ایہ دکھانا اے کہ
«$\pi$، $!A$ دا
!!a{derivation} $\Gamma$
توں اے» والا تعلق، $\pi$ تے
$!A$ دے گوڈل نمبراں
دا بنیادی بازگشتی تعلق اے۔
\end{explain}

\begin{prop}\ollabel{prop:followsby}
خاصیت $\fn{Correct}(p)$ اُدوں تے صرف اُدوں قائم اے جدوں
!!{derivation}~$\pi$ دا آخری استنباط درست ہووے، جتھے اوس اخذ دا
گوڈل نمبر~$p$ اے۔ ایہ خاصیت بنیادی بازگشتی اے۔
\end{prop}

\begin{proof}
$\Gamma \Sequent \Delta$
ابتدائی سیکوئنٹ اے جے یا تاں
کوئی !!a{sentence}~$!A$
ایہوجیہا ہووے کہ
$\Gamma \Sequent \Delta$،
$!A \Sequent !A$ ہووے؛
یا کوئی ٹرم~$t$ ایہوجیہا
ہووے کہ
$\Gamma \Sequent \Delta$،
$\emptyset \Sequent \eq[t][t]$
ہووے۔ گوڈل نمبراں دے لحاظ نال
$\fn{InitSeq}(s)$ عین
اوس ویلے سچ اے جدوں
\begin{align*}
\bexists{x < s}{} (\fn{Sent}(x) & \land
s = \tuple{\tuple{x},\tuple{x}})
\lor {}\\
\bexists{t<s}{} (\fn{Term}(t) & \land
s = \tuple{0, \tuple{\Gn{{\eq}(} \concat t \concat \Gn{,} \concat t \concat \Gn{)}}}).
\end{align*}

سانوں ایہ وی دکھانا
پینا اے کہ استنباط دے
ہر قاعدے~$R$ لئی
تعلق $\fn{FollowsBy}_R(p)$
بنیادی بازگشتی اے۔
ایتھے
$\fn{FollowsBy}_R(p)$
عین اوس ویلے سچ اے
جدوں $p$، کسے
!!{derivation}~$\pi$
دا گوڈل نمبر ہووے،
تے~$\pi$ دا آخرلا
سیکوئنٹ، قاعدہ~$R$ ٹھیک طرح لاؤن نال~$\pi$ دے براہِ راست ذیلی
!!{derivation}s دے نتیجیاں توں نکلدا ہووے۔

اک سوکھی مثال \RightR{\land} دا قاعدہ اے۔ جے $\pi$ اک درست
$\RightR{\land}$ استنباط اُتے ختم ہووے، تاں اوہ ایویں دِسے گا:
\begin{prooftree}
  \AxiomC{}
  \RightLabel{$\pi_1$}
  \Deduce$\Gamma \fCenter \Delta, !A$

    \AxiomC{}
    \RightLabel{$\pi_2$}
    \Deduce$\Gamma \fCenter \Delta, !B$

  \RightLabel{\RightR\land}
  \BinaryInf$\Gamma \fCenter \Delta, !A \land !B$
\end{prooftree}
ایس لئی
!!{derivation} $\pi$
وچ آخری استنباط
$\RightR{\land}$ دا
درست اطلاق اُدوں تے صرف اُدوں اے جدوں
!!{sentence}s دے
سلسلے $\Gamma$
تے $\Delta$، نالے
دو !!{sentence}s
$!A$ تے~$!B$
ایہوجیہے ہون کہ
$\pi_1$ دا آخرلا
سیکوئنٹ
$\Gamma \Sequent \Delta, !A$
ہووے، $\pi_2$ دا
آخرلا سیکوئنٹ
$\Gamma \Sequent \Delta, !B$
ہووے، تے $\pi$
دا آخرلا سیکوئنٹ
$\Gamma \Sequent \Delta, !A \land !B$
ہووے۔ ہن بس ایس نوں
گوڈل نمبراں وچ
بیان کرنا اے۔ جے
$s = \Gn{\Gamma \Sequent \Delta}$
ہووے، تاں
$(s)_0 = \Gn{\Gamma}$
تے $(s)_1 = \Gn{\Delta}$
ہوندے نیں۔ ایس لئی
$\fn{FollowsBy}_{\RightR{\land}}(p)$
عین اوس ویلے سچ اے
جدوں
\begin{align*}
& \bexists{g < p}{\bexists{d < p}{\bexists{a < p}{\bexists{b < p}{\quad}}}} \\
& \qquad \fn{EndSequent}(p) =
  \tuple{g, d \concat
    \tuple{\Gn{(} \concat a \concat \Gn{\land} \concat b \concat \Gn{)}}} \land {} \\
& \qquad \fn{EndSequent}((p)_1) = \tuple{g, d \concat \tuple{a}} \land {}\\
& \qquad \fn{EndSequent}((p)_2) = \tuple{g, d \concat \tuple{b}} \land {}\\
& \qquad (p)_0 = 2 \land \fn{LastRule}(p) = 10.
\end{align*}
ہر سطر دی گل بالترتیب ایہ اے کہ
«سلسلہ~($\Gamma$)،
جیہدا گوڈل نمبر~$g$ اے،
سلسلہ~($\Delta$)،
جیہدا گوڈل نمبر~$d$ اے،
!!a{formula}~($!A$)،
جیہدا گوڈل نمبر~$a$ اے،
تے !!a{formula}~($!B$)،
جیہدا گوڈل نمبر~$b$ اے،
موجود نیں»، تے
نال ایہ کہ
«$\pi$ دا آخرلا
سیکوئنٹ
$\Gamma \Sequent \Delta, !A
\land !B$ اے»،
«$\pi_1$ دا آخرلا
سیکوئنٹ
$\Gamma \Sequent \Delta,
!A$ اے»،
«$\pi_2$ دا آخرلا
سیکوئنٹ
$\Gamma \Sequent \Delta, !B$
اے»، تے
«$\pi$ دیاں دو
براہِ راست ذیلی
اخذاں نیں تے آخری
استنباطی قاعدہ
$\RightR\land$
اے (جیہدا نمبر~$10$
اے)»۔

$\pi$ وچ آخری استنباط
$\RightR\lexists$
دا درست اطلاق اُدوں تے صرف اُدوں اے جدوں
سلسلے $\Gamma$
تے $\Delta$، اک
!!a{formula}~$!A$،
اک متغیر~$x$،
تے اک ٹرم~$t$
ایہوجیہے ہون کہ
$\pi$ دا آخرلا
سیکوئنٹ
$\Gamma \Sequent \Delta, \lexists[x][!A]$
تے $\pi_1$ دا
آخرلا سیکوئنٹ
$\Gamma \Sequent \Delta,
\Subst{!A}{t}{x}$
ہووے۔ ایس لئی گوڈل
نمبراں دے لحاظ نال
$\fn{FollowsBy}_{\RightR\lexists}(p)$
عین اوس ویلے سچ اے
جدوں
\begin{align*}
  & \bexists{g<p}{\bexists{d<p}{\bexists{a<p}{\bexists{x<p}{\bexists{t<p}{\quad}}}}}\\
  & \qquad \fn{EndSequent}(p) =
  \tuple{
    g,
    d \concat \tuple{\Gn{\lexists} \concat x \concat a}
  } \land {}\\
  & \qquad \fn{EndSequent}((p)_1) =
  \tuple{
    g,
    d \concat \tuple{\fn{Subst}(a, t, x)}
  } \land {}\\
  & \qquad (p)_0 = 1 \land \fn{LastRule}(p) = 18.
\end{align*}
فیر اسیں $\fn{Correct}(p)$ دی تعریف ایویں کردے آں:
\begin{multline*}
  \fn{Sequent}(\fn{EndSequent}(p)) \land {}\\
  [(\fn{LastRule}(p) = 1 \land
  \fn{FollowsBy}_{\LeftR\Weakening}(p)) \lor \dots \lor {}\\
  (\fn{LastRule}(p) = 20 \land \fn{FollowsBy}_{\eq}(p)) \lor {}\\
  (p)_0 = 0 \land \fn{InitialSeq}(\fn{EndSequent}(p))]
\end{multline*}
پہلی سطر یقینی بناندی
اے کہ~$d$ دا آخرلا
سیکوئنٹ واقعی
!!{sentence}s دا
سیکوئنٹ اے۔ آخری
سطر اوس صورت لئی اے
جدوں $p$ بس اک
ابتدائی سیکوئنٹ ہووے۔
\end{proof}

\begin{prob}
ہِٹھلیاں خاصیتاں
\olref[inc][art][plk]{prop:followsby}
دی طرز اُتے تعریف کرو:
\begin{enumerate}
\item $\fn{FollowsBy}_{\Cut}(p)$،
\item $\fn{FollowsBy}_{\LeftR\lif}(p)$،
\item $\fn{FollowsBy}_{\eq}(p)$،
\item $\fn{FollowsBy}_{\RightR{\lforall}}(p)$۔
\end{enumerate}
آخری صورت لئی ایہ وی وکھاؤ کہ
گوڈل نمبر~$p$
والے !!{derivation}
دا آخری استنباط
ایگن متغیر والی
شرط پوری کردا
اے یا نہیں، ایہ
بنیادی بازگشتی
ڈھنگ نال پرکھیا
جا سکدا اے۔ یعنی
ایگن متغیر~$a$،
جیہڑا $\RightR{\lforall}$
استنباط دا اے،
آخرلے سیکوئنٹ
وچ نہ آوے۔
\end{prob}

\begin{prop}
\ollabel{prop:deriv}
تعلق $\fn{Deriv}(p)$ بنیادی بازگشتی اے۔ ایہ قائم اے جے $p$
کسے درست !!{derivation}~$\pi$ دا گوڈل نمبر ہووے۔
\end{prop}

\begin{proof}
!!^a{derivation}~$\pi$
اوس ویلے درست اے
جدوں اوہدا ہر استنباط کسے قاعدے
دا درست اطلاق
ہووے؛ یعنی اوہدے ذیلی !!{derivation}s
وچوں ہر اک
درست استنباط اُتے
ختم ہووے۔ ایس لئی
$\fn{Deriv}(d)$
عین اوس ویلے سچ
اے جدوں
\[
\bforall{i<\len{\fn{SubtreeSeq}(p)}}{\fn{Correct}((\fn{SubtreeSeq}(p))_i}.
\]
\end{proof}

\begin{prop}
فرض کرو $\Gamma$،
!!{sentence}s دا
بنیادی بازگشتی
سیٹ اے۔ تاں تعلق
$\Prf[\Gamma](x, y)$،
جیہڑا ایہ بیان
کردا اے کہ
«$x$،
!!a{derivation}~$\pi$
دا کوڈ اے جیہدا
آخرلا سیکوئنٹ
$\Gamma_0 \Sequent !A$
اے، جتھے متناہی $\Gamma_0 \subseteq \Gamma$
ہووے، تے
$y$،~$!A$ دا
گوڈل نمبر اے»،
بنیادی بازگشتی
اے۔
\end{prop}

\begin{proof}
فرض کرو
«$y \in \Gamma$»
بنیادی بازگشتی
محمول~$R_\Gamma(y)$
نال دتا گیا اے۔
سانوں ایہ دکھانا
اے کہ
$\Prf[\Gamma](x, y)$،
جیہڑا عین اوس
ویلے سچ اے جدوں
$y$ کسے
جملے~$!A$ دا
گوڈل نمبر ہووے
تے $x$، $\Log{LK}$ وچلے کسے !!{derivation} دا کوڈ ہووے جیہدا
آخرلا سیکوئنٹ
$\Gamma_0 \Sequent !A$
ہووے،
بنیادی بازگشتی
اے۔

پچھلے قضیے مطابق خاصیت
$\fn{Deriv}(x)$،
جیہڑی عین اوس
ویلے سچ اے جدوں
$x$ کسے درست
!!{derivation}~$\pi$
دا کوڈ ہووے
جیہڑا $\Log{LK}$
وچ اے،
بنیادی بازگشتی
اے۔ جے $x$
ایہوجیہا کوڈ
اے، تاں
$\fn{EndSequent}(x)$،
~$\pi$ دے
آخرلے سیکوئنٹ
دا کوڈ اے، تے
ایس لئی
$(\fn{EndSequent}(x))_0$
آخرلے سیکوئنٹ
دے کھبے پاسے
دا کوڈ، تے
$(\fn{EndSequent}(x))_1$
سجے پاسے دا
کوڈ اے۔ ایس
لئی اسیں
«$\pi$ دے
آخرلے سیکوئنٹ
دا سجا پاسا
صرف~$!A$ اے»
نوں
$\len{(\fn{EndSequent}(x))_1} = 1 \land ((\fn{EndSequent}(x))_1)_0 =
x$
نال بیان کر
سکدے آں۔
$\pi$ دے
آخرلے سیکوئنٹ
دا کھبا پاسا آپے متناہی اے۔ سانوں صرف ایہ بیان کرنا اے کہ
اوہدا ہر جملہ~$\Gamma$ وچ اے۔ ایس لئی
$\Prf[\Gamma](x, y)$
نوں ایس طرح
تعریف کر سکدے
آں:
\begin{align*}
\Prf[\Gamma](x, y) \defiff {}&
\fn{Deriv}(x) \land {} \\
& \bforall{i <
  \len{(\fn{EndSequent}(x))_0}}{R_\Gamma(((\fn{EndSequent}(x))_0)_i)} \land {}\\
& \len{(\fn{EndSequent}(x))_1} = 1 \land ((\fn{EndSequent}(x))_1)_0 = y.
\end{align*}
\end{proof}
\endgroup
% /OLP-0286

% OLP-0287 content/incompleteness/arithmetization-syntax/proofs-in-nd.tex
\begingroup
\def\olfilename{proofs-in-nd}
\expandafter\def\csname ol@sectioncs\endcsname{section}

\olfileid{inc}{art}{pnd}
\olsection{فطری استنباط وچ \usetoken{P}{derivation}}

\begin{explain}
!!{derivation}s نوں عددی روپ دین لئی سانوں !!{derivation}s نوں
عدداں نال ظاہر کرنا پینا اے۔ کیوں جے !!{derivation}s،
!!{formula}s دے درخت نیں تے ہر استنباط نال اک یا دو لیبل وی
ہوندے نیں، ایس لئی بازگشتی طریقہ سِدھا اے۔
اسیں !!{derivation} نوں اک ٹپل نال ظاہر کردے آں۔
اوس ٹپل دے جز آخری استنباط دے مقدمیاں تک پہنچن والیاں
براہِ راست ذیلی !!{derivation}s دی گنتی، اوہناں ذیلی !!{derivation}s دے عددی روپ، تے آخرلا !!{formula}،
آخری استنباط دا مفروضہ خارج کرن والا لیبل، تے آخری
استنباط دی قسم دسن والا اک عدد نیں۔
\end{explain}

\begin{defn}
  جے $\delta$ فطری استنباط وچ !!a{derivation} اے،
  تاں $\Gn{\delta}$ دی استقرائی تعریف ایویں کردے آں:
  \begin{enumerate}
  \item
    جے $\delta$ صرف مفروضے~$!A$ اُتے مشتمل اے،
    تاں $\Gn{\delta}$، $\tuple{0, \Gn{!A}, n}$ اے۔
    جے ایہ !!{undischarged} مفروضہ اے تاں عدد~$n$،
    $0$ اے؛ ورنہ اوس دا عددی لیبل اے۔
  \item
    جے $\delta$ دا آخرلا استنباط صفر، اک، دو یا
    تن مقدمیاں والا اے، تاں $\Gn{\delta}$ بالترتیب
    ایہ اے:
    \begin{align*}
      & \tuple{0, \Gn{!A}, n, k}, \\
      & \tuple{1, \Gn{\delta_1}, \Gn{!A}, n, k}, \\
      & \tuple{2, \Gn{\delta_1}, \Gn{\delta_2}, \Gn{!A}, n, k}, \text{ یا}\\
      & \tuple{3, \Gn{\delta_1}, \Gn{\delta_2}, \Gn{\delta_3}, \Gn{!A}, n,
        k},
    \end{align*}
    ایتھے $\delta_1$، $\delta_2$ تے $\delta_3$،
    $\delta$ دے آخری استنباط دے مقدمے یا مقدمیاں
    اُتے ختم ہون والے ذیلی !!{derivation}s نیں۔
    $!A$، $\delta$ وچ آخری استنباط دا نتیجہ اے؛
    $n$ اوس استنباط دا مفروضہ خارج کرن والا لیبل
    اے (جے استنباط کوئی مفروضہ خارج نہیں کردا
    تاں $0$)؛ تے آخری استنباط وچ ورتے قاعدے مطابق $k$ دی قدر ہِٹھلے جدول توں ملدی اے۔

    \begin{tabular}{lccccccc}
      \text{قاعدہ:} & \Intro{\land} & \Elim{\land} & \Intro{\lor} & \Elim{\lor} \\
      $k$: & 1 & 2 & 3 & 4  \\[1.5ex]
      \text{قاعدہ:} &  \Intro{\lif} & \Elim{\lif} & \Intro{\lnot} & \Elim{\lnot} \\
      $k$: & 5 & 6 & 7 & 8 \\[1.5ex]
      \text{قاعدہ:}  &  \FalseInt & \FalseCl & \Intro{\lforall} & \Elim{\lforall} \\
      $k$: & 9 & 10 & 11 & 12 \\[1.5ex]
      \text{قاعدہ:} & \Intro{\lexists} & \Elim{\lexists} & \Intro{\eq} & \Elim{\eq} \\
      $k$: & 13 & 14 & 15 & 16
    \end{tabular}
  \end{enumerate}
\end{defn}

\begin{ex}
  ایہ سوکھا جیہا !!{derivation} ویکھو:
  \begin{prooftree}
    \AxiomC{$\Discharge{!A \land !B}{1}$}
    \RightLabel{\Elim{\land}}
    \UnaryInfC{$!A$}
    \DischargeRule{\Intro{\lif}}{1}
    \UnaryInfC{$(!A \land !B) \lif !A$}
  \end{prooftree}
  مفروضے دا گوڈل نمبر
  $d_0 = \tuple{0, \Gn{!A \land !B}, 1}$ ہوئے گا۔
  $\Elim{\land}$ دے نتیجے اُتے ختم ہون والے
  !!{derivation} دا گوڈل نمبر
  $d_1 = \tuple{1, d_0, \Gn{!A}, 0, 2}$
  ہوئے گا ($1$ ایس لئی کہ $\Elim{\land}$ دا اک مقدمہ
  اے؛ نتیجے~$!A$ دا گوڈل نمبر؛ $0$ ایس لئی کہ
  کوئی مفروضہ خارج نہیں ہوندا؛ تے $2$،
  $\Elim{\land}$ نوں کوڈ کرن والا عدد اے)۔
  فیر پورے !!{derivation} دا گوڈل نمبر
  $\tuple{1, d_1, \Gn{((!A \land !B) \lif
      !A)}, 1, 5}$ اے؛ یعنی
  \[
  \tuple{1, \tuple{1, \tuple{0, \Gn{(!A \land !B)}, 1}, \Gn{!A}, 0, 2},
    \Gn{((!A \land !B) \lif !A)}, 1, 5}.
  \]
\end{ex}

\begin{explain}
!!{derivation}s دی عددی روپ طے کر کے
سانوں ایہ وی دکھانا پینا اے کہ اسیں انہاں
!!{derivation}s دے گوڈل نمبراں اُتے بنیادی
بازگشتی ڈھنگ نال عمل کر سکدے آں، تے اوہناں دیاں ضروری خاصیتاں تے تعلقاں نوں وی ایس طریقے نال بیان کر سکدے آں۔ کجھ عمل سادے نیں: مثال لئی،
!!a{derivation} دے گوڈل نمبر~$d$ توں
$\fn{EndFmla}(d) = (d)_{(d)_0+1}$
اوس دے آخرلے !!{formula} دا گوڈل نمبر
دیندا اے؛
$\fn{DischargeLabel}(d) = (d)_{(d)_0 + 2}$
مفروضہ خارج کرن والا لیبل دیندا اے؛ تے
$\fn{LastRule}(d) = (d)_{(d)_0 + 3}$
آخری استنباط دی قسم دسن والا عدد دیندا اے۔
کجھ ہور عمل کافی اوکھے نیں۔ اسیں گھٹ توں
گھٹ انہاں دا طریقہ خاکہ وار دساں گے۔
مقصد ایہ دکھانا اے کہ
«$\delta$،~$!A$ دا
!!a{derivation} $\Gamma$ توں اے» والا تعلق،
$\delta$ تے~$!A$ دے گوڈل نمبراں
دا بنیادی بازگشتی تعلق اے۔
\end{explain}

\begin{prop}
  ہِٹھلے تعلق بنیادی بازگشتی نیں:
  \begin{enumerate}
  \item $!A$، $\delta$ وچ لیبل~$n$ والا
    مفروضے وجوں آؤندا اے۔
  \item $\delta$ وچ لیبل~$n$ والے سارے
    مفروضے $!A$ دے روپ وچ نیں
    (یعنی اسیں مفروضہ~$!A$ نوں
    لیبل~$n$ ورت کے $\delta$ وچ
    !!{discharge} کر سکدے آں)۔
  \end{enumerate}
\end{prop}

\begin{proof}
  سانوں دکھانا اے کہ !!{formula}s دے گوڈل
  نمبراں تے !!{derivation}s دے گوڈل نمبراں
  دے وچکار متعلقہ تعلق بنیادی بازگشتی نیں۔
  \begin{enumerate}
  \item اسیں دکھانا چاہندے آں کہ
    $\fn{Assum}(x, d, n)$ بنیادی بازگشتی اے۔
    ایہ اوس ویلے سچ ہوندا اے جدوں $x$،
    گوڈل نمبر~$d$ والے !!{derivation}
    دے لیبل~$n$ والے کسے مفروضے دا گوڈل
    نمبر ہووے۔ ایہ اوس ویلے اے جدوں
    گوڈل نمبر~$\tuple{0, x, n}$ والا
    !!{derivation}، $d$ دا ذیلی
    !!{derivation} ہووے۔ دھیان دیو کہ
    !!{derivation}s دی ساڈی کوڈکاری،
    \olref[cmp][rec][tre]{sec} وچ دتی
    درختاں دی کوڈکاری دی اک خاص صورت اے۔
    ایس لئی بنیادی بازگشتی فنکشن
    $\fn{SubtreeSeq}(d)$، $d$ دے سارے
    ذیلی !!{derivation}s دے گوڈل نمبراں
    دا سلسلہ دیندا اے (جیہدی لمبائی ودھ توں ودھ $d$ اے)۔
    سو اسیں ایہ تعریف کر سکدے آں:
    \[
    \fn{Assum}(x, d, n) \defiff \bexists{i<d}{(\fn{SubtreeSeq}(d))_i =
      \tuple{0, x, n}}.
    \]
  \item اسیں دکھانا چاہندے آں کہ
    $\fn{Discharge}(x, d, n)$ بنیادی
    بازگشتی اے۔ ایہ اوس ویلے سچ ہوندا
    اے جدوں لیبل~$n$ والے،
    گوڈل نمبر~$d$ دے
    !!{derivation} وچلے سارے
    مفروضے، گوڈل نمبر~$x$
    والے !!{formula} ہی ہون۔
    ایہ تعلق عین اوس ویلے سچ اے جدوں
    $\bforall{y<d}{(\fn{Assum}(y, d, n) \lif y
      = x)}$ ہووے۔
  \end{enumerate}
\end{proof}

\begin{prop}
  \ollabel{prop:followsby}
  خاصیت $\fn{Correct}(d)$ اُدوں تے صرف اُدوں قائم اے جدوں
!!{derivation}~$\delta$ دا آخری استنباط درست ہووے، جتھے اوس اخذ
دا گوڈل نمبر~$d$ اے۔ ایہ خاصیت بنیادی بازگشتی اے۔
\end{prop}

\begin{proof}
  ایتھے سانوں استنباط دے ہر قاعدے~$R$
  لئی ایہ دکھانا پینا اے کہ تعلق
  $\fn{FollowsBy}_R(d)$ بنیادی بازگشتی
  اے۔ $\fn{FollowsBy}_R(d)$ عین اوس
  ویلے سچ ہوندا اے جدوں $d$ کسے
  !!{derivation}~$\delta$ دا گوڈل نمبر
  ہووے، تے $\delta$ دا آخرلا
  !!{formula}، قاعدہ~$R$ ٹھیک طرح لاؤن نال $\delta$ دے براہِ راست ذیلی
!!{derivation}s دے نتیجیاں توں نکلدا ہووے۔

  اک سوکھی مثال \Intro{\land} دا قاعدہ اے۔ جے $\delta$ اک درست
$\Intro{\land}$ استنباط اُتے ختم ہووے، تاں اوہ ایویں دِسے گا:
  \begin{prooftree}
    \AxiomC{}
    \RightLabel{$\delta_1$}
    \DeduceC{$!A$}

    \AxiomC{}
    \RightLabel{$\delta_2$}
    \DeduceC{$!B$}

    \RightLabel{\Intro\land}
    \BinaryInfC{$!A \land !B$}
  \end{prooftree}
  فیر گوڈل نمبر~$d$، جیہڑا~$\delta$ نوں دِتا گیا اے،
  $\tuple{2, d_1, d_2,
    \Gn{(!A \land !B)}, 0, k}$ اے؛ جتھے
  $\fn{EndFmla}(d_1) = \Gn{!A}$،
  $\fn{EndFmla}(d_2) = \Gn{!B}$،
  $n=0$ تے $k=1$ نیں۔ سو اسیں
  $\fn{FollowsBy}_{\Intro\land}(d)$
  نوں ایس طرح تعریف کر سکدے آں:
  \begin{multline*}
    (d)_0 = 2 \land \fn{DischargeLabel}(d) = 0 \land \fn{LastRule}(d) = 1 \land {}\\
    \fn{EndFmla}(d) = {}\\
    \Gn{(} \concat \fn{EndFmla}((d)_1) \concat \Gn{\land}
    \concat \fn{EndFmla}((d)_2) \concat \Gn{)}.
  \end{multline*}

  اک ہور سادہ مثال $\Intro\eq$ دا قاعدہ
  اے۔ ایس دے کوئی مقدمے نہیں، ایس لئی
  $(d)_0 = 0$ اے، مفروضیاں وانگ۔
  ایس دا مفروضہ خارج کرن والا لیبل وی
  نہیں؛ یعنی $n=0$۔ پر $!A$ ضرور
  $\eq[t][t]$ دے روپ وچ ہووے،
  جتھے~$t$ بند ٹرم اے۔
  ایتھے اک بنیادی بازگشتی تعریف ایہ اے:
  \begin{multline*}
    (d)_0 = 0 \land \fn{DischargeLabel}(d) = 0 \land {}\\
    \bexists{t<d}{(\fn{ClTerm}(t) \land
      \fn{EndFmla}(d) = {}\\
      \Gn{{\eq}(} \concat t \concat \Gn{,} \concat t \concat \Gn{)})}.
  \end{multline*}

  ذرا اوکھی مثال وچ
  $\fn{FollowsBy}_{\Intro{\lif}}(d)$
  عین اوس ویلے سچ ہوندا اے جدوں
  $\delta$ دا آخرلا !!{formula}،
  $(!A \lif !B)$ دے روپ وچ ہووے،
  $\delta_1$ دا آخرلا !!{formula}
  $!B$ ہووے، تے $\delta$ وچ
  لیبل~$n$ والا ہر مفروضہ
  $!A$ دے روپ وچ ہووے۔
  اسیں ایس نوں بنیادی بازگشتی
  طریقے نال ایویں بیان کر
  سکدے آں:
  \begin{multline*}
    (d)_0 = 1 \land {}\\
    \bexists{a<d}{(\fn{Discharge}(a, (d)_1, \fn{DischargeLabel}(d)) \land {}}\\
      \fn{EndFmla}(d) = (\Gn{(} \concat a \concat \Gn{\lif}
      \concat \fn{EndFmla}((d)_1) \concat \Gn{)}))
  \end{multline*}
  ($a$ نوں~$!A$ دا گوڈل نمبر سمجھو)۔

  اک ہور مثال لئی \Intro{\lexists}
  نوں ویکھو۔ ایتھے $\delta$
  دا آخری استنباط عین اوس
  ویلے درست اے جدوں کوئی
  !!a{formula}~$!A$، بند ٹرم~$t$
  تے !!a{variable}~$x$ ایہوجیہے
  ہون کہ $\Subst{!A}{t}{x}$،
  !!{derivation}~$\delta_1$ دا
  آخرلا !!{formula} ہووے تے
  $\lexists[x][!A]$ آخری
  استنباط دا نتیجہ ہووے۔
  سو $\fn{FollowsBy}_{\Intro{\lexists}}(d)$
  عین اوس ویلے سچ ہوندا اے جدوں
  \begin{multline*}
    (d)_0 = 1 \land \fn{DischargeLabel}(d) = 0 \land {} \\
    \bexists{a < d}{\bexists{x<d}{\bexists{t<d}{
          (\fn{ClTerm}(t) \land \fn{Var}(x)  \land {}}}}\\
    \fn{Subst}(a,t,x) = \fn{EndFmla}((d)_1) \land
    \fn{EndFmla}(d) = (\Gn{\lexists} \concat x \concat a)).
  \end{multline*}

  فیر اسیں $\fn{Correct}(d)$ دی تعریف ایویں کردے آں:
  \begin{multline*}
    \fn{Sent}(\fn{EndFmla}(d)) \land {}\\
    (\fn{LastRule}(d) = 1 \land
    \fn{FollowsBy}_{\Intro\land}(d)) \lor \dots \lor {}\\
    (\fn{LastRule}(d) = 16 \land \fn{FollowsBy}_{\Elim\eq}(d)) \lor {}\\
    \bexists{n<d}{\bexists{x<d}{(d = \tuple{0, x, n})}}.
  \end{multline*}
  پہلی سطر یقینی بناندی اے کہ~$d$
  دا آخرلا !!{formula} واقعی
  جملہ اے۔ آخری سطر اوس صورت
  لئی اے جدوں $d$ بس اک
  مفروضہ ہووے۔
\end{proof}

\begin{prob}
  ہِٹھلیاں خاصیتاں
  \olref[inc][art][pnd]{prop:followsby}
  دی طرز اُتے تعریف کرو:
  \begin{enumerate}
  \item $\fn{FollowsBy}_{\Elim{\lif}}(d)$،
  \item $\fn{FollowsBy}_{\Elim{\eq}}(d)$،
  \item $\fn{FollowsBy}_{\Elim{\lor}}(d)$،
  \item $\fn{FollowsBy}_{\Intro{\lforall}}(d)$۔
  \end{enumerate}
  آخری صورت لئی ایہ وی وکھاؤ کہ گوڈل
  نمبر~$d$ والے !!{derivation}
  دا آخری استنباط ایگن متغیر
  والی شرط پوری کردا اے یا
  نہیں، ایہ بنیادی بازگشتی
  ڈھنگ نال پرکھیا جا سکدا اے۔
  یعنی ایگن متغیر~$a$،
  جیہڑا $\Intro{\lforall}$
  استنباط دا اے،
  نہ~$d$ دے آخرلے
  !!{formula} وچ آوے،
  نہ~$d$ دے کسے
  خارج نہ کیتے مفروضے
  وچ۔ ایس لئی تسیں
  \olref[inc][art][pnd]{prop:openassum}
  والا بنیادی بازگشتی
  محمول $\fn{OpenAssum}$
  ورت سکدے او۔
\end{prob}

\begin{prop}
  \ollabel{prop:deriv}
  تعلق $\fn{Deriv}(d)$،
  جیہڑا اوس ویلے سچ
  ہوندا اے جدوں $d$
  کسے درست
  !!{derivation}~$\delta$
  دا گوڈل نمبر ہووے،
  بنیادی بازگشتی اے۔
\end{prop}

\begin{proof}
  !!^a{derivation}~$\delta$
  اوس ویلے درست اے جدوں
  اوہدا ہر استنباط کسے
  قاعدے دا درست اطلاق
  ہووے؛ یعنی اوہدے ذیلی
  !!{derivation}s وچوں
  ہر اک درست استنباط
  اُتے ختم ہووے۔ ایس
  لئی $\fn{Deriv}(d)$
  عین اوس ویلے سچ
  اے جدوں
  \[
  \bforall{i<\len{\fn{SubtreeSeq}(d)}}{\fn{Correct}((\fn{SubtreeSeq}(d))_i)}
  \]
\end{proof}

\begin{prop}
  \ollabel{prop:openassum}
  تعلق $\fn{OpenAssum}(z, d)$ بنیادی بازگشتی اے۔ ایہ قائم اے جے
$z$ کسے !!a{undischarged} مفروضے~$!A$ دا گوڈل نمبر ہووے جیہڑا
!!{derivation}~$\delta$ وچ اے، تے اوس اخذ دا گوڈل نمبر~$d$ ہووے۔
\end{prop}

\begin{proof}
  کسے مفروضے دی اک
  آمد اوس ویلے
  !!{discharged} ہوندی
  اے جدوں اوہ
  لیبل~$n$ نال
  $\delta$ دے کسے
  ذیلی !!{derivation}
  وچ آوے تے اوہ
  ذیلی اخذ اوسے لیبل~$n$ نال
  مفروضہ خارج کرن
  والے قاعدے اُتے
  ختم ہووے۔ سو
  $!A$، $\delta$
  دا !!a{undischarged}
  مفروضہ اوس ویلے
  اے جدوں اوہدیاں
  آمدواں وچوں گھٹ
  توں گھٹ اک
  $\delta$ وچ
  !!{discharged} نہ ہوئی
  ہووے۔ ایتھے ایہ گل یاد رکھو کہ
  $\delta$ وچ
  $!A$ دیاں
  !!{discharged} تے
  !!{undischarged}
  دونوں قسم دیاں
  آمدواں ہو سکدیاں
  نیں۔

  اک سلسلہ
  $\delta_0$، \dots،
  $\delta_k$ لو،
  جتھے $\delta_0 =
  \delta$، $\delta_k$
  مفروضہ
  $\Discharge{!A}{n}$
  ہووے (کسے~$n$ لئی)، تے
  $\delta_{i+1}$،
  $\delta_i$ دا
  براہِ راست ذیلی
  !!{derivation}
  ہووے۔ جے ایسا سلسلہ ملے جیہدے وچ کوئی
  وی~$\delta_i$
  لیبل~$n$ نال
  مفروضہ خارج
  کرن والے
  استنباط اُتے
  ختم نہ ہووے،
  تاں $!A$،
  $\delta$ دا
  !!a{undischarged}
  مفروضہ اے۔

  بنیادی بازگشتی فنکشن
  $\fn{SubtreeSeq}(d)$
  سانوں $\delta$ دے
  سارے ذیلی
  !!{derivation}s دے
  گوڈل نمبراں دا
  سلسلہ دیندا اے۔
  $\delta$ دے ذیلی
  !!{derivation}s دے
  گوڈل نمبراں دا
  ہر سلسلہ اوس دا
  ذیلی سلسلہ اے۔
  ذیلی سلسلہ ہون
  دا تعلق بنیادی
  بازگشتی اے:
  $\fn{Subseq}(s, s')$
  عین اوس ویلے
  سچ اے جدوں
  $\bforall{i<\len{s}}{\lexists[j<\len{s'}][(s)_i = (s')_j]}$۔
  براہِ راست ذیلی
  !!{derivation}
  ہون دا تعلق وی
  ایہوجیہا اے:
  $\fn{Subderiv}(d, d')$
  عین اوس ویلے
  سچ اے جدوں
  $\bexists{j<(d')_0}{d = (d')_j}$
  ہووے۔ سو $\fn{OpenAssum}(z, d)$ دی تعریف ایویں کر سکدے آں:
  \begin{multline*}
    \bexists{s<\fn{SubtreeSeq}(d)}{(\fn{Subseq}(s, \fn{SubtreeSeq}(d))
      \land (s)_0 = d \land {}} \\
    \bexists{n<d}{((s)_{\len{s} \tsub 1} = \tuple{0, z, n} \land {}}\\
      \bforall{i<(\len{s} \tsub 1)}{(\fn{Subderiv}((s)_{i+1}, (s)_i) \land {}}\\
      \fn{DischargeLabel}((s)_i) \neq n))).
  \end{multline*}
\end{proof}

\begin{prop}
  \ollabel{prop:prf-prim-rec}
  فرض کرو $\Gamma$،
  !!{sentence}s دا
  بنیادی بازگشتی
  سیٹ اے۔ تاں
  تعلق
  $\Prf[\Gamma](x, y)$،
  جیہڑا ایہ بیان
  کردا اے کہ
  «$x$، کسے
  !!a{derivation}~$\delta$
  دا کوڈ اے،
  جیہڑا $!A$ نوں $\Gamma$ وچلے !!{undischarged} مفروضیاں توں کڈھدا اے، تے
  $y$،~$!A$
  دا گوڈل نمبر
  اے»، بنیادی
  بازگشتی اے۔
\end{prop}

\begin{proof}
  فرض کرو
  «$y \in \Gamma$»
  بنیادی بازگشتی محمول~$R_\Gamma(y)$ نال دِتّا گیا اے۔
  سانوں دکھانا اے
  کہ تعلق
  $\Prf[\Gamma](x, y)$
  بنیادی بازگشتی
  اے۔ ایہ عین
  اوس ویلے سچ
  ہوندا اے جدوں
  $y$ کسے
  جملے~$!A$ دا
  گوڈل نمبر ہووے،
  تے $x$ کسے
  فطری استنباطی
  !!{derivation}
  دا کوڈ ہووے
  جیہدا آخرلا
  !!{formula}~$!A$
  اے، تے جیہدے
  سارے !!{undischarged}
  مفروضے
  $\Gamma$ وچ
  ہون۔

  \olref{prop:deriv}
  دے مطابق
  خاصیت
  $\fn{Deriv}(x)$،
  جیہڑی عین اوس
  ویلے سچ اے
  جدوں $x$
  فطری استنباط
  وچ کسے درست
  !!{derivation}~$\delta$
  دا گوڈل نمبر
  ہووے، بنیادی
  بازگشتی اے۔
  ایس لئی اسیں
  $\Prf[\Gamma](x, y)$ دی تعریف ایویں کر سکدے آں:
  \begin{align*}
    \Prf[\Gamma](x, y) \defiff {}
    & \fn{Deriv}(x) \land \fn{EndFmla}(x) = y \land {} \\
    & \bforall{z < x}{(\fn{OpenAssum}(z, x) \lif R_\Gamma(z))}.
  \end{align*}
\end{proof}
\endgroup
% /OLP-0287

% OLP-0288 content/incompleteness/arithmetization-syntax/proofs-in-ax.tex
\begingroup
\def\olfilename{proofs-in-ax}
\expandafter\def\csname ol@sectioncs\endcsname{section}

\olfileid{inc}{art}{pax}
\olsection{مسلّمی \usetoken{P}{derivation}}

\begin{explain}
مسلّمی !!{derivation}s نوں عددی روپ دین لئی
سانوں !!{derivation}s نوں عدداں دے روپ وچ
ظاہر کرنا پینا اے۔ کیوں جے !!{derivation}s
بس !!{formula}s دے سلسلے نیں، ایس لئی
واضح طریقہ ایہ اے کہ ہر !!{derivation}
نوں اوس وچلے !!{formula}s دے کوڈاں دے
سلسلے دا کوڈ دیئیے۔
\end{explain}

\begin{defn}
جے $\delta$ اک مسلّمی !!{derivation} اے،
جیہڑا !!{formula}s $!A_1$، \dots،
~$!A_n$ توں بنیا اے، تاں
$\Gn{\delta}$ ایہ اے:
\[
\tuple{\Gn{!A_1}, \dots, \Gn{!A_n}}.
\]
\end{defn}

\begin{ex}
  ایہ سوکھا جیہا !!{derivation} ویکھو:
  \begin{derivation}
  1. & $!B \lif (!B \lor !A)$ \\
  2. & $(!B \lif (!B \lor !A)) \lif (!A  \lif (!B \lif (!B \lor !A)))$\\
  3. & $!A  \lif (!B \lif (!B \lor !A))$
  \end{derivation}
  ایس !!{derivation} دا گوڈل نمبر
  ایہ ہووے گا:
  \begin{align*}
  \openTuple\,
    & \Gn{!B \lif (!B \lor !A)}, \\
    &\Gn{(!B \lif (!B \lor !A)) \lif (!A  \lif (!B \lif (!B \lor !A)))},\\
    & \Gn{!A  \lif (!B \lif (!B \lor !A))} \,\closeTuple.
  \end{align*}
\end{ex}

\begin{explain}
!!{derivation}s دی عددی روپ طے کر کے سانوں ایہ وی دکھانا پینا اے
کہ اسیں انہاں !!{derivation}s اُتے
بنیادی بازگشتی ڈھنگ نال عمل کر
سکدے آں تے انہاں دیاں ضروری
خاصیتاں تے تعلقاں نوں وی ایس
طرح بیان کر سکدے آں۔ کجھ عمل
سوکھے نیں۔ مثال لئی، کسے
!!a{derivation} دے گوڈل نمبر~$d$
توں $(d)_{\len{d}-1}$ اوس دے
آخرلے !!{formula} دا گوڈل نمبر
دیندا اے۔ کجھ ہور عمل کافی
اوکھے نیں۔ اسیں گھٹ توں گھٹ
انہاں دا طریقہ خاکہ وار دساں
گے۔ مقصد ایہ دکھانا اے کہ
«$\delta$،~$!A$ دا
!!a{derivation} $\Gamma$ توں
اے» والا تعلق، $\delta$
تے~$!A$ دے گوڈل نمبراں دا بنیادی بازگشتی تعلق اے۔
\end{explain}

\begin{prop}
\ollabel{prop:followsby}
ایہ تعلق بنیادی بازگشتی نیں:
\begin{enumerate}
\item $!A$ اک مسلّمہ اے۔
\item $i$ویں سطر، جیہڑی $\delta$ وچ اے،
  دا جواز قاعدۂ اثباتِ مقدم نال اے۔
\item $i$ویں سطر، جیہڑی $\delta$ وچ اے،
  دا جواز \QR نال اے۔
\item $\delta$ اک درست
  !!{derivation} اے۔
\end{enumerate}
\end{prop}

\begin{proof}
سانوں دکھانا اے کہ !!{formula}s دے
گوڈل نمبراں تے !!{derivation}s دے
گوڈل نمبراں دے وچکار متعلقہ
تعلق بنیادی بازگشتی نیں۔
\begin{enumerate}
\item ساڈے کول مسلّمہ خاکیاں دی
  اک مقرر فہرست اے؛ تے $!A$
  مسلّمہ اوس ویلے اے جدوں اوہ
  کسے اک خاکے دے دتے روپ وچ
  ہووے۔ کیوں جے خاکیاں دی فہرست متناہی اے، ایس لئی
  ہر مسلّمہ خاکے لئی ایہ
  دکھانا کافی اے کہ اسیں
  بنیادی بازگشتی ڈھنگ نال
  پرکھ سکدے آں کہ $!A$
  اوس روپ وچ اے یا نہیں۔
  مثال لئی ایہ مسلّمہ
  خاکہ ویکھو:
  \[
  !B \lif (!C \lif !B).
  \]
  $!A$ ایس مسلّمہ خاکے دی
  مثال اوس ویلے اے جدوں
  !!{formula}s $!B$ تے $!C$
  ایہوجیہے ہون کہ $!A$
  ایویں ملے: '$($' نال
  $!B$، فیر '$\lif$'،
  فیر '$($'، فیر $!C$،
  فیر '$\lif$'، فیر $!B$
  تے آخر~'$))$' جوڑیا
  جائے۔ اسیں گوڈل نمبر $n$ دی،
  جیہڑا~$!A$ دا اے،
  متعلقہ خاصیت پرکھ سکدے
  آں، کیوں جے سلسلیاں
  نوں جوڑنا بنیادی
  بازگشتی اے، تے $!B$
  تے $!C$ دے گوڈل نمبر
  ضرور~$!A$ دے گوڈل
  نمبر توں چھوٹے نیں۔
  ایس دی وجہ ایہ اے کہ
  جدوں ایہ تعلق سچ
  ہووے، تاں $!B$
  تے $!C$ دونوں
  ~$!A$ دے ذیلی
  !!{formula}s نیں۔
  لہٰذا اسیں ایہ
  تعریف کر سکدے آں:
  \begin{multline*}
  \fn{IsAx}_{!B \lif (!C \lif !B)}(n) \defiff \bexists{b< n}{
    \bexists{c < n}{(\fn{Sent}(b) \land \fn{Sent}(c) \land {}}}\\ n =
  \Gn{(} \concat b \concat \Gn{\lif} \concat \Gn{(} \concat c \concat
  \Gn{\lif} \concat b \concat \Gn{))}).
  \end{multline*}
  جے ہر مسلّمہ خاکے
  لئی ایہوجیہی تعریف
  ہووے، تاں اوہناں دا منطقی «یا»
  خاصیت $\fn{IsAx}(n)$
  دی تعریف کردا اے:
  «$n$ کسے مسلّمے
  دا گوڈل نمبر اے»۔
\item $i$ویں سطر، جیہڑی $\delta$ وچ اے،
  دا جواز قاعدۂ
  اثباتِ مقدم نال
  عین اوس ویلے اے
  جدوں سطراں $j$
  تے $k < i$
  موجود ہون؛ سطر~$j$
  دا !!{sentence}
  کوئی فارمولا $!A$
  ہووے، سطر~$k$
  دا جملہ
  $!A \lif !B$
  ہووے، تے
  سطر~$i$ دا
  جملہ~$!B$
  ہووے۔
  \begin{multline*}
    \fn{MP}(d, i) \defiff \bexists{j < i}{\bexists{k < i}{}}\\
    (d)_k = \Gn{(} \concat (d)_j
      \concat \Gn{\lif} \concat (d)_i \concat \Gn{)}
  \end{multline*}
  کیوں جے محدود
  مسوّریت، سلسلیاں
  دا جوڑ تے $=$
  بنیادی بازگشتی
  نیں، ایہ اک
  بنیادی بازگشتی
  تعلق دی تعریف اے۔
\item $\delta$ دی کسے
  سطر دا جواز \QR{}
  نال اوس ویلے اے
  جدوں اوہ
  $!B \lif \lforall[x][!A(x)]$
  دے روپ وچ ہووے،
  پہلاں آ چکی کوئی سطر $!B \lif !A(c)$ ہووے، جتھے کوئی
  !!{constant}~$c$
  ہووے، تے $c$
  ~$!B$ وچ نہ
  آوندا ہووے۔
  ایہ عین اوس
  ویلے اے جدوں
  \begin{enumerate}
  \item کوئی !!a{sentence}~$!B$
    ہووے، تے
  \item !!a{formula}~$!A(x)$
    ہووے جیہدے وچ
    صرف متغیر~$x$
    آزاد ہووے؛ نال
    ایہ کہ
  \item سطر $i$ وچ
    $!B \lif \lforall[x][!A(x)]$
    ہووے،
  \item کوئی سطر
    $j < i$،
    $!B \lif \Subst{!A}{c}{x}$
    رکھدی ہووے،
    جتھے $c$ مستقل
    ہووے،
  \item تے اوہ
    ~$!B$ وچ نہ
    آوندا ہووے۔
  \end{enumerate}
  ایہ ساریاں گلّاں
  بنیادی بازگشتی
  ڈھنگ نال پرکھیاں
  جا سکدیاں نیں،
  کیوں جے $!B$،
  $!A(x)$ تے $x$
  دے گوڈل نمبر،
  سطر~$i$ دے
  فارمولے دے
  گوڈل نمبر توں
  چھوٹے نیں؛ تے
  $a$ دا گوڈل
  نمبر، سطر~$j$
  دے فارمولے دے
  گوڈل نمبر توں
  چھوٹا اے:
  \begin{multline*}
    \fn{QR}_1(d, i) \defiff \bexists{b < (d)_i}{\bexists{x <
        (d)_i}{\bexists{a < (d)_i}{\bexists{c < (d)_j}{(}}}}
    \\ \fn{Var}(x) \land \fn{Const}(c) \land {} \\
    (d)_i = \Gn{(} \concat
    b \concat \Gn{\lif} \concat \Gn{\lforall} \concat x \concat a
    \concat \Gn{)} \land {}\\
    (d)_j = \Gn{(} \concat b \concat
    \Gn{\lif} \concat \fn{Subst}(a,c,x) \concat \Gn{)} \land {}\\ \fn{Sent}(b)
    \land \fn{Sent}(\fn{Subst}(a,c,x)) \land {} \bforall{k <
      \len{b}}{(b)_k \neq (c)_0})
  \end{multline*}
  ایتھے اسیں $c$ نوں ثابتہ نشان دا تے $x$ نوں متغیر دا
  گوڈل نمبر سمجھدے آں، جدوں دوناں نوں ٹرم وجوں ویکھیے
  (یعنی انہاں دیاں علامتاں دے کوڈ نہیں)۔
  اسیں ایہ
  پرکھدے آں
  کہ $x$ ہی
  ~$!A(x)$ دا
  اکیلا آزاد
  متغیر اے؛
  ایس لئی
  $\Subst{!A(x)}{c}{x}$
  دا !!a{sentence}
  ہونا پرکھدے
  آں۔ تے ایہ
  یقینی بناندے
  آں کہ $c$،
  $!B$ وچ نہ
  آوے؛ ایس لئی
  منگدے آں
  کہ~$!B$ دی
  ہر علامت
  ~$c$ توں
  وکھری ہووے۔

  \QR{} دا
  دوجا روپ
  مشق وجوں چھڈدے آں۔
\item $d$ کسے
  درست !!{derivation}
  دا گوڈل نمبر
  عین اوس ویلے
  اے جدوں اوہدی
  ہر سطر یا
  تاں مسلّمہ
  ہووے، یا
  قاعدۂ اثباتِ
  مقدم نال، یا
  ~\QR نال
  جواز رکھدی
  ہووے۔ ایس لئی:
  \[
  \fn{Deriv}(d) \defiff \bforall{i < \len{d}}{(\fn{IsAx}((d)_i) \lor
    \fn{MP}(d,i) \lor \fn{QR}(d, i))}
  \]
\end{enumerate}
\end{proof}

\begin{prob}
ایہ تعلق
\olref[inc][art][pax]{prop:followsby}
دی طرز اُتے
تعریف کرو:
\begin{enumerate}
\item $\fn{IsAx}_{!A \lif (!B \lif (!A \land !B))}(n)$،
\item $\fn{IsAx}_{\lforall[x][!A(x)] \lif !A(t)}(n)$،
\item $\fn{QR}_{2}(d, i)$
  (\QR دے دوجے روپ لئی)۔
\end{enumerate}
\end{prob}

\begin{prop}
فرض کرو $\Gamma$،
!!{sentence}s دا
بنیادی بازگشتی
سیٹ اے۔ تاں
تعلق
$\Prf[\Gamma](x, y)$،
جیہڑا ایہ بیان
کردا اے کہ
«$x$ کسے
!!a{derivation}~$\delta$
دا کوڈ اے
جیہڑا $!A$ نوں~$\Gamma$ توں کڈھدا اے، تے
$y$،~$!A$ دا
گوڈل نمبر اے»،
بنیادی بازگشتی
اے۔
\end{prop}

\begin{proof}
فرض کرو
«$y \in \Gamma$»
بنیادی بازگشتی
محمول~$R_\Gamma(y)$ نال دِتّا گیا اے۔
سانوں دکھانا اے
کہ تعلق
$\Prf[\Gamma](x, y)$
بنیادی بازگشتی
اے؛ ایتھے
$\Prf[\Gamma](x,
y)$ عین اوس
ویلے سچ اے
جدوں $y$
کسے
!!a{sentence}~$!A$
دا گوڈل نمبر
ہووے تے $x$،
~$!A$ دے
!!a{derivation}
دا کوڈ ہووے
جیہڑا $\Gamma$
توں اے۔

پچھلے قضیے مطابق خاصیت
$\fn{Deriv}(x)$،
جیہڑی عین اوس
ویلے سچ اے
جدوں $x$
کسے درست
!!{derivation}~$\delta$
دا کوڈ ہووے،
بنیادی بازگشتی
اے۔ پر اوس
تعریف وچ سیٹ
$\Gamma$ نوں
!!{derivation}
دیاں سطراں
دا جواز دین
دا اک ہور
طریقہ نہیں
گِنّیا گیا۔
ساڈے بنیادی بازگشتی پرکھ وچ، جیہڑی ویکھدی اے
کہ سطر دا
جواز \QR{}
نال اے یا
نہیں، ایہ
شرط وی
نظرانداز سی
کہ مستقل~$c$
نوں~$\Gamma$
وچ آون دی
اجازت نہیں۔
ساڈی تعریف
نوں بدل کے
براہِ راست
$\Gamma$ دا
لحاظ رکھیا
جا سکدا اے،
پر $\fn{Deriv}$
تے استنباط
دا قضیہ
ورتنا سوکھا
اے۔
$\Gamma \Proves !A$
عین اوس
ویلے اے
جدوں
!!{sentence}s
دی کوئی متناہی فہرست $!B_1$،
\dots،
$!B_n \in \Gamma$
ایہوجیہی ہووے
کہ
$\{!B_1, \dots, !B_n\} \Proves !A$۔
تے استنباط
دے قضیے
مطابق ایہ
اوس ویلے اے
جدوں
$\Proves (!B_1 \lif (!B_2 \lif
\cdots (!B_n \lif !A)\cdots))$
ہووے۔
گوڈل نمبر~$z$
والا
!!a{sentence}
ایس روپ وچ
اے یا نہیں،
ایہ بنیادی
بازگشتی ڈھنگ
نال پرکھیا
جا سکدا اے۔
سو $x$ نوں
گوڈل نمبر~$y$
والے
!!{sentence}
دے
!!a{derivation}
دا گوڈل نمبر
\emph{$\Gamma$ توں}
سمجھّن دی
تھاں، اسیں
$x$ نوں
اُتّے دتے
روپ دے تہہ
در تہہ
شرطیہ دا
!!a{derivation}
دا گوڈل نمبر
سمجھدے آں،
جیہڑا~$\emptyset$
توں اے۔

پہلاں، جے
ساڈے کول
!!{sentence}s
دا سلسلہ
ہووے، تاں
اسیں بنیادی
بازگشتی ڈھنگ
نال ایہ
شرطیہ بنا
سکدے آں
جیہدے سارے
مقدم ایہ
جملے ہون
تے تالی
اک دتا
!!{sentence}
ہووے:
\begin{align*}
  \fn{hCond}(s, y, 0) & = y \\
  \fn{hCond}(s, y, n+1) & = \Gn{(} \concat (s)_{n} \concat \Gn{\lif}
  \concat \fn{Cond}(s, y, n) \concat \Gn{)}\\
  \fn{Cond}(s, y) & = \fn{hCond}(s, y, \len{s})\\
  \intertext{سو اسیں $\Prf[\Gamma](x, y)$ دی تعریف ایویں کر سکدے آں}
  \Prf[\Gamma](x, y) & \defiff \bexists{s < \fn{sequenceBound}(x,x)}{(} \\
  &\qquad (x)_{\len{x}-1} = \fn{Cond}(s,y) \land {} \\
  &\qquad\bforall{i<\len{s}}{(s)_i \in \Gamma} \land {} \\
  &\qquad\fn{Deriv}(x)).
\end{align*}
~$s$ دی حد
ایس گل توں ملدی اے کہ ہر
$(s)_i$،
!!{derivation}
دی آخری
سطر دے کسے
ذیلی
!!{formula}
دا گوڈل نمبر
اے؛ یعنی
اوہ
$(x)_{\len{x}-1}$
توں چھوٹا اے۔
مقدمیاں
$!B \in \Gamma$
دی گنتی،
یعنی~$s$
دی لمبائی،
~$x$ دی
آخری سطر
دی لمبائی
توں گھٹ اے۔
\end{proof}
\endgroup
% /OLP-0288

\OLEndChapterHook
\endgroup
% /OLP-0280

% OLP-0289 content/incompleteness/representability-in-q/representability-in-q.tex
\begingroup
\def\olfilename{representability-in-q}
\expandafter\def\csname ol@sectioncs\endcsname{section}

\olchapter{inc}{req}{\texorpdfstring{$\Th{Q}$ وچ قابلِ نمائندگی ہونا}{Q وچ قابلِ نمائندگی ہونا}}

% OLP-0290 content/incompleteness/representability-in-q/introduction.tex
\begingroup
\def\olfilename{introduction}
\expandafter\def\csname ol@sectioncs\endcsname{section}

\olfileid{inc}{req}{int}
\olsection{تعارف}

نامکملیت دے قضیے اوہناں نظرییاں اُتے لاگو ہوندے نیں جیہناں وچ
قابلِ حساب فنکشناں بارے بنیادی گلاں بیان تے ثابت کیتیاں جا سکدیاں
نیں۔ اسیں ایہو جیہا اک نہایت کمزور نظریہ دساں گے جیہنوں «$\Th{Q}$»
آکھدے نیں (یا کدی کدی رافائل رابنسن دے ناں اُتے «رابنسن دا $Q$»)۔
اسیں دساں گے کہ کسے فنکشن دا $\Th{Q}$ وچ \emph{قابلِ نمائندگی}
ہون دا کیہ مطلب اے، تے فیر ایہ گل ثابت کراں گے:
\begin{quote}
  کسے فنکشن دی $\Th{Q}$ وچ نمائندگی اُدوں تے صرف اُدوں ہو سکدی اے
جدوں اوہ قابلِ حساب ہووے۔
\end{quote}
ایس نال سانوں قابلِ حسابیت دا اک ہور نمونہ ملدا اے۔ ایہ وی
وکھاواں گے کہ سیٹ $\Setabs{!A}{\Th{Q} \Proves !A}$ فیصلہ پذیر نہیں۔
ایس لئی رُکّن دے مسئلے نوں ایس سیٹ دی رکنیت دے فیصلے وچ
گھٹاؤاں گے۔ آخر تک اسیں ایس توں کِتے ودھ طاقتور نتیجے ثابت
کر چُکے ہواں گے۔

$\Th{Q}$ دی زبان حساب دی زبان اے۔ $\Th{Q}$ دے مسلّمے ہِٹھ دِتّے نیں۔
اوہناں نوں !!{identity} والی پہلے درجے دی منطق دے ہور مسلّمیاں
تے قاعدیاں نال رلا کے ورتنا اے۔
\begin{align*}
& \lforall[x][\lforall[y][(\eq[x'][y'] \lif \eq[x][y])]] \tag{$!Q_1$}\\
& \lforall[x][\eq/[\Obj 0][x']] \tag{$!Q_2$}\\
& \lforall[x][(\eq[x][\Obj 0] \lor \lexists[y][\eq[x][y']])] \tag{$!Q_3$}\\
& \lforall[x][\eq[(x + \Obj 0)][x]] \tag{$!Q_4$}\\
& \lforall[x][\lforall[y][\eq[(x + y')][(x + y)']]] \tag{$!Q_5$}\\
& \lforall[x][\eq[(x \times \Obj 0)][\Obj 0]] \tag{$!Q_6$}\\
& \lforall[x][\lforall[y][\eq[(x \times y')][((x \times y) + x)]]] \tag{$!Q_7$}\\
& \lforall[x][\lforall[y][(x < y \liff \lexists[z][\eq[(z' + x)][y]])]] \tag{$!Q_8$}
\end{align*}
ہر قدرتی عدد $n$ لئی عددی عبارت $\num{n}$ اوہ ٹرم
$\Obj{0}^{\prime\prime\ldots\prime}$ اے جیہدے وچ کل $n$ پرائم نشان
ہون۔ سو $\num{0}$ صرف
!!{constant}~$\Obj{0}$ اے، $\num{1}$،
$\Obj{0}'$ اے، $\num{2}$، $\Obj{0}''$ اے، تے ایویں اگّے۔

حساب دے نظریے وجوں $\Th{Q}$ \emph{بہت ای} کمزور اے۔ مثال لئی،
ایس وچ $\lforall[x][\eq/[x][x']]$ یا
$\lforall[x][\lforall[y][(x + y) = (y + x)]]$ ورگیاں سادیاں
گلاں وی ثابت نہیں ہو سکدیاں۔ اسیں ویکھاں گے کہ $\Th{Q}$ دے
دلچسپ ہون دی وڈی وجہ \emph{ایہو} کمزوری اے۔ اوہ بس اینا کُ
طاقتور اے کہ نامکملیت دا قضیہ لاگو ہو سکے۔ $\Th{Q}$ دے
دلچسپ ہون دی اک ہور وجہ ایہ اے کہ اوہدے مسلّمیاں دا سیٹ
\emph{متناہی} اے۔

$\Th{Q}$ توں طاقتور نظریہ، جیہنوں \emph{پیانو حساب} $\Th{PA}$
کہندے نیں،~$\Th{Q}$ وچ استقرا دا اک خاکہ شامل کرن نال ملدا اے۔
\[
(!A(\Obj 0) \land \lforall[x][(!A(x) \lif !A(x'))]) \lif \lforall[x][!A(x)]
\]
جتھے $!A(x)$ کوئی وی فارمولا اے۔ جے $!A(x)$ وچ $x$ توں علاوہ ہور !!{variable}s آزاد ہون،
تاں اسیں شروع وچ کلی مسوّر لا کے اوہناں ساریاں نوں پابند کر
دیندے آں تاں جے استقرا دے خاکے دی ایہ صورت !!a{sentence} ہووے۔
مثال لئی، جے $!A(x, y)$ وچ !!{variable}~$y$ وی آزاد
ہووے، تاں متعلقہ مثال ایہ اے:
\[
\lforall[y][((!A(\Obj 0) \land \lforall[x][(!A(x) \lif !A(x'))]) \lif
  \lforall[x][!A(x)])]
\]
استقرا دے خاکے دیاں ایہو جہیاں صورتاں ورت کے~$\Th{PA}$ دے
مسلّمیاں توں، $\Th{Q}$ دے مسلّمیاں نالوں، کِتے ودھ کجھ ثابت
کیتا جا سکدا اے۔ حقیقت وچ، قدرتی عدداں بارے
ایہوجیہے «قدرتی» بیان لبھّن لئی کافی محنت پینی اے
جیہڑے پیانو حساب وچ ثابت نہ ہو سکّن!{}

\begin{defn}
\ollabel{defn:representable-fn}
  قدرتی عدداں توں قدرتی عدداں ول فنکشن $f(x_0,\ldots,x_k)$ نوں
{\em $\Th{Q}$ وچ قابلِ نمائندگی} آکھدے آں جے کوئی فارمولا
$!A_f(x_0,\dots,x_k,y)$ ایہو جیہا ہووے کہ جدوں وی
$f(n_0,\dots,n_k) = m$ ہووے، تاں $\Th{Q}$ ایہ دونوں گلاں ثابت کرے:
\begin{enumerate}
\item\ollabel{defn:rep:a} $!A_f(\num{n_0}, \dots, \num{n_k}, \num{m})$
\item\ollabel{defn:rep:b} $\lforall[y][(!A_f(\num{n_0}, \dots,
\num{n_k}, y) \lif \num{m} = y)]$۔
\end{enumerate}
\end{defn}

ایس تعریف نوں ہور طرح وی بیان کر سکدے آں۔ مثال لئی ایس
تعریف دے برابر ایہ شرط وی رکھی جا سکدی اے کہ $\Th{Q}$ ایہ ثابت کرے: $\lforall[y][(!A_f(\num{n_0}, \dots,
\num{n_k}, y) \liff \eq[y][\num{m}])]$۔

\begin{thm}
\ollabel{thm:representable-iff-comp}
کسے فنکشن دی $\Th{Q}$ وچ نمائندگی اُدوں تے صرف اُدوں ہو سکدی اے
جدوں اوہ قابلِ حساب ہووے۔
\end{thm}

ایس قضیے دیاں دوہاں سمتاں ثابت کرنیاں نیں۔ نحو نوں عددی روپ
دے چکے ہوئیے تاں قابلِ نمائندگی توں قابلِ حساب ہون والی سمت
کافی سِدھی اے۔ دوجی سمت لئی ودھ کم کرنا پینا اے۔ خیال ایہ اے کہ اسیں «قابلِ حساب» نوں
ٹھیک معنی دین لئی «عمومی بازگشتی» نوں چُناں گے تے وکھاواں گے کہ ہر عمومی بازگشتی فنکشن دی~$\Th{Q}$ وچ
نمائندگی ہو سکدی اے۔ یاد کرو، فنکشن عمومی بازگشتی
ہوندا اے جے اوہ $\Zero$، جانشین فنکشن~$\Succ$ تے
پروجیکشن فنکشناں~$\Proj{n}{i}$ توں، ترکیب، بنیادی
بازگشت تے باقاعدہ کمترین قدر دی تلاش ورت کے، تعریف کیتا جا سکے۔
سو ایہ دکھاون دا اک طریقہ کہ ہر عمومی بازگشتی
فنکشن دی~$\Th{Q}$ وچ نمائندگی ہو سکدی اے، ایہ اے
کہ بنیادی فنکشناں دی نمائندگی ہو سکدی اے تے جدوں
کجھ فنکشناں دی نمائندگی ہو سکدی ہووے، تاں انہاں
توں ترکیب، بنیادی بازگشت تے باقاعدہ کمترین قدر دی تلاش نال
تعریف کیتے فنکشناں دی وی ہو سکے۔ دوجے لفظاں وچ، اسیں ایہ وکھا سکدے آں کہ بنیادی فنکشن قابلِ
نمائندگی نیں تے ترکیب، بنیادی بازگشت تے باقاعدہ کمترین قدر
دی تلاش نال بنے فنکشن وی قابلِ نمائندگی فنکشناں دے طبقے
وچ رہندے نیں۔ ایہ یقینی بناندا اے کہ ہر عمومی بازگشتی
فنکشن دی نمائندگی ہو سکدی اے۔

ایہ وکھانا اوکھا اے کہ قابلِ نمائندگی فنکشناں توں بنیادی
بازگشت نال بنے فنکشن وی قابلِ نمائندگی نیں۔ ایس قدم توں
بچن لئی پہلاں وکھاواں گے کہ بنیادی بازگشت توں بغیر وی
کم چل سکدا اے۔ یعنی ہر عمومی بازگشتی فنکشن بنیادی فنکشناں
توں صرف ترکیب تے باقاعدہ کمترین قدر دی تلاش نال تعریف
کیتا جا سکدا اے۔ ایس لئی وکھاواں گے کہ اک خاص باقاعدہ کمترین
قدر دی تلاش بنیادی بازگشت دا کم کر سکدی اے۔ ایہ طریقہ
ورتن لئی کجھ ہور بنیادی فنکشن شامل کرن پین گے۔ اوہ جمع، ضرب
تے برابری دے تعلق دا امتیازی فنکشن~$\Char{=}$ نیں۔ فیر قضیہ
ایہ وکھا کے ثابت کر سکدے آں کہ \emph{ایہناں} سارے بنیادی
فنکشناں دی~$\Th{Q}$ وچ نمائندگی ہو سکدی اے تے ترکیب تے
باقاعدہ کمترین قدر دی تلاش نال قابلِ نمائندگی فنکشناں توں
بنے فنکشن وی قابلِ نمائندگی نیں۔
\endgroup
% /OLP-0290

% OLP-0291 content/incompleteness/representability-in-q/representable-comp.tex
\begingroup
\def\olfilename{representable-comp}
\expandafter\def\csname ol@sectioncs\endcsname{section}

\olfileid{inc}{req}{rpc}
\olsection{\texorpdfstring{$\Th{Q}$ وچ قابلِ نمائندگی فنکشن قابلِ حساب نیں}{Q وچ قابلِ نمائندگی فنکشن قابلِ حساب نیں}}

اسیں ثابت کراں گے کہ جیہڑے فنکشن دی~$\Th{Q}$ وچ نمائندگی
ہو سکدی اے، اوہ قابلِ حساب اے۔ پہلاں سانوں~$\Th{Q}$ وچ
قابلِ نمائندگی فنکشناں بارے اک معاون قضیہ ثابت کرنا پینا اے۔

\begin{lem}\ollabel{lem:rep-q}
  جے $f(x_0,
  \dots, x_k)$ دی~$\Th{Q}$ وچ نمائندگی ہو سکدی اے،
  تاں کوئی !!a{formula}~$!A(x_0, \dots, x_k, y)$ ایہو جیہا
موجود اے کہ
  \[
  \Th{Q} \Proves !A_f(\num{n_0}, \dots, \num{n_k}, \num{m})
  \quad\text{جے تے صرف جے}\quad m = f(n_0, \dots, n_k).
  \]
\end{lem}

\begin{proof}
  «جے» والی سمت
\olref[int]{defn:representable-fn}\olref[int]{defn:rep:a}
اے۔ «صرف جے» والی سمت لئی فرض
کرو $\Th{Q} \Proves !A_f(\num{n_0},
\dots, \num{n_k}, \num{m})$، پر
$m \neq f(n_0, \dots, n_k)$۔ من لو $l =
f(n_0, \dots, n_k)$۔ شرط
\olref[int]{defn:representable-fn}\olref[int]{defn:rep:a}
مطابق، $\Th{Q}
\Proves !A_f(\num{n_0}, \dots, \num{n_k}, \num{l})$۔
شرط
\olref[int]{defn:representable-fn}\olref[int]{defn:rep:b}
مطابق، ایہ وی ثابت ہوندا اے:
$\lforall[y][(!A_f(\num{n_0}, \dots, \num{n_k}, y) \lif \num{l} =
y)]$۔ منطق تے ایس مفروضے نوں ورت کے کہ
$\Th{Q} \Proves
!A_f(\num{n_0}, \dots, \num{n_k}, \num{m})$، سانوں
ایہ ملدا اے کہ $\Th{Q}
\Proves \eq[\num{l}][\num{m}]$۔ دوجے پاسے،
\olref[bre]{lem:q-proves-neq} مطابق، $\Th{Q} \Proves
\eq/[\num{l}][\num{m}]$۔ ایس توں نکلدا اے کہ $\Th{Q}$ بے تضاد نہیں۔ پر ایہ ہو نہیں سکدا،
کیوں جے معیاری نمونہ $\Th{Q}$ اُتے پورا اُتردا اے (ویکھو
\olref[int][def]{def:standard-model})،
$\Sat{N}{\Th{Q}}$؛ تے صحت مندی دے قضیے مطابق جیہڑے نظریے کسے نمونے وچ پورے اتردے نیں، اوہ ہمیشہ بے تضاد ہوندے نیں
(\tagrefs{prfAX/{fol:axd:sou:cor:consistency-soundness},
prfSC/{fol:seq:sou:cor:consistency-soundness},
prfND/{fol:ntd:sou:cor:consistency-soundness},
prfTab/{fol:tab:sou:cor:consistency-soundness}})۔
\end{proof}

\begin{lem}
جیہڑے فنکشن دی $\Th{Q}$ وچ نمائندگی ہو سکدی اے، اوہ قابلِ حساب اے۔
\end{lem}

\begin{proof}
پہلاں ایس گل دی سِدھی وجہ ویکھیے۔~$f$ دی قدر کڈھن لئی
اسیں حساب دی زبان وچ سارے ممکن !!{derivation}s~$\delta$ دی
فہرست بناواں گے۔ ایہ مشینی طریقے نال ہو سکدا اے۔ ہر اک نوں
پرکھاں گے کہ اوہ ایس نظریے دے مسلّمیاں توں !!a{derivation} اے
جیہدا آخرلا !!a{formula} ایس روپ وچ ہووے یا نہیں:
~$!A_f(\num{n_0}, \dots,
\num{n_k}, \num{m})$ (یعنی
\olref{lem:rep-q} والا !!{formula}، جیہڑا $f$
دی~$\Th{Q}$ وچ نمائندگی کردا اے)۔ جے ہووے،
تاں \olref{lem:rep-q} مطابق
$m = f(n_0, \dots, n_k)$، تے سانوں~$f$
دی قدر مل گئی۔ تلاش ضرور رُکے گی، کیوں جے
$\Th{Q} \Proves !A_f(\num{n_0}, \dots, \num{n_k},
\num{f(n_0, \dots, n_k)})$؛ سو آخرکار سانوں
ٹھیک قسم دا اک $\delta$ مل جائے گا۔

ایہ طریقہ ہلے پوری طرح متعین نہیں، کیوں جے ایہ نرے عدداں دی تھاں !!{derivation}s
تے !!{formula}s اُتے عمل کردا اے، تے اسیں
ایہ وی ٹھیک ٹھیک نہیں دسّیا کہ «سارے ممکن
!!{derivation}s دی فہرست بنانا» مشینی طریقے
نال کیوں ممکن اے۔ پر، جیویں اسیں ویکھ چکے آں، ٹرماں، !!{formula}s تے !!{derivation}s
نوں گوڈل نمبراں نال کوڈ کیتا جا سکدا اے۔
اسیں حساب دا اک ٹھیک طرح متعین نمونہ، یعنی
عمومی بازگشتی فنکشن، وی متعارف کرا چُکے
آں۔ تے اسیں ویکھیا اے کہ تعلق
$\Prf[\Th{Q}](d,y)$ (جیہڑا عین اوس ویلے
سچ اے جدوں $d$ اوس !!a{derivation} دا گوڈل نمبر ہووے جیہڑا گوڈل
نمبر~$y$ والے !!{formula} نوں~$\Th{Q}$ دے مسلّمیاں توں کڈھدا اے)
(بنیادی) بازگشتی اے۔ ہور بنیادی
بازگشتی فنکشن جیہڑے سانوں چاہیدے
ہون گے، اوہ $\fn{num}$
(\olref[art][trm]{prop:num-primrec}) تے
$\fn{Subst}$
(\olref[art][sub]{prop:subst-primrec})
نیں۔ انہاں توں~$f$ نوں کمترین قدر
دی تلاش راہیں تعریف کیتا جا سکدا
اے؛ سو $f$ بازگشتی اے۔

پہلاں ایہ تعریف کرو:
\begin{multline*}
  A(n_0, \dots, n_k, m) = \\
  \fn{Subst}(\fn{Subst}(\dots\fn{Subst}(\Gn{!A_f}, \fn{num}(n_0), \Gn{x_0}),\\ \dots),
  \fn{num}(n_k),  \Gn{x_k}), \fn{num}(m), \Gn{y})
\end{multline*}
ایہ لکھت اوکھی لگدی اے۔ ایہ فنکشن
$A(n_0, \dots, n_k,
m) = \Gn{!A_f(\num{n_0}, \dots, \num{n_k}, \num{m})}$
اے۔

ہن تعلق~$R(n_0, \dots, n_k, s)$ ویکھو،
جیہڑا اُدوں قائم اے جدوں $(s)_0$ اوس !!a{derivation} دا گوڈل نمبر
ہووے جیہڑا~$\Th{Q}$ توں
$!A_f(\num{n_0}, \dots, \num{n_k}, \num{(s)_1})$
دا اخذ اے:
\[
R(n_0, \dots, n_k, s) \quad\text{جے تے صرف جے}\quad \Prf[\Th{Q}]((s)_0, A(n_0,
\dots, n_k, (s)_1))
\]
جے اسیں اک~$s$ لبھ سکّیے جیہدے لئی
$R(n_0, \dots, n_k, s)$ سچ ہووے،
تاں سانوں عدداں دا اک جوڑا مل گیا---$(s)_0$ تے~$(s)_1$---ایہو جیہا کہ $(s)_0$،~$A_f(\num{n_0}, \dots,
\num{n_k}, (s)_1)$ دے
!!a{derivation} دا گوڈل نمبر اے۔
سو~$s$ دی تلاش غیر رسمی ثبوت
وچ $d$ تے $m$ دی جوڑی لبھّن
وانگ اے۔ تے ایہوجیہا~$s$
«لبھّن» والا قابلِ حساب فنکشن
باقاعدہ کمترین قدر دی تلاش
راہیں تعریف کیتا جا سکدا اے۔
دھیان دیو کہ $R$ باقاعدہ اے۔ ہر $n_0$، \dots، $n_k$ لئی
!!a{derivation}~$\delta$ موجود
اے جیہڑا ایہ دکھاندا اے کہ
$\Th{Q} \Proves !A_f(\num{n_0}, \dots,
\num{n_k}, \num{f(n_0, \dots, n_k)})$؛
سو $R(n_0, \dots, n_k, s)$،
$s = \tuple{\Gn{\delta}, f(n_0, \dots, n_k)}$
لئی سچ اے۔ ایس لئی اسیں $f$
نوں ایویں لکھ سکدے آں:
\[
f(n_0,\dots,n_{k}) = (\umin{s}{R(n_0, \dots, n_k, s)})_1.
\]
\end{proof}
\endgroup
% /OLP-0291

% OLP-0292 content/incompleteness/representability-in-q/beta-function.tex
\begingroup
\def\olfilename{beta-function}
\expandafter\def\csname ol@sectioncs\endcsname{section}

\olfileid{inc}{req}{bet}
\olsection{بیٹا فنکشن دا معاون قضیہ}

سانوں وکھانا اے کہ جمع، ضرب تے $\Char{=}$ موجود ہون تاں
اسیں بنیادی بازگشت دا کم کر سکدے آں۔ ایس لئی سانوں سلسلیاں
اُتے کم کرن والے فنکشن بنانے پین گے۔ (جے قوت چڑھاؤن دا عمل وی
ہووے، تاں کم سوکھا ہووے۔) جدوں اسیں بنیادی بازگشت ورت سکدے ساں،
تاں «$n$-واں اولی عدد» ورگیاں شےواں دی تعریف کر سکدے ساں تے
کوڈ دین دا سِدھا طریقہ چُن سکدے ساں۔ ایتھے بنیادی بازگشت
پہلے توں ساڈے کول نہیں۔ اسیں وکھانا چاہندے آں کہ کمترین قدر
دی تلاش نال بنیادی بازگشت دا کم ہو سکدا اے۔
ایس لئی سانوں سلسلیاں دے کوڈاں لئی ہور طریقہ لبھنا پینا اے۔

\begin{lem}
\ollabel{lem:beta}
اک فنکشن $\beta(d,i)$ موجود اے جیہدے لئی ایہ شرط پوری ہوندی
اے۔ ہر سلسلے $a_0$، \dots،~$a_n$ لئی کوئی عدد~$d$ ایسا ملدا اے
کہ ہر $i \le n$ لئی $\beta(d,i) = a_i$ ہووے۔ نالے $\beta$ دی
تعریف بنیادی فنکشناں توں صرف ترکیب تے باقاعدہ کمترین قدر دی
تلاش نال کیتی جا سکدی اے۔
\end{lem}

$d$ نوں سلسلے $\tuple{a_0, \dots, a_n}$ دا کوڈ سمجھو۔ فنکشن
$\beta(d,i)$ ایس سلسلے دا $i$-واں رکن دیندا اے۔ (ایہ اوہ کوڈکاری
\emph{نہیں} جیہڑی اسیں اولی عدداں دیاں قوتاں نال پہلے کر چکے آں۔)
معاون قضیہ صرف تھوڑی جیہی گل ثابت کردا اے۔ ایہ سلسلے جوڑن یا
آخر وچ رکن پاؤن دا طریقہ نہیں دیندا۔ ایہ وی نہیں آکھدا کہ ترکیب
تے باقاعدہ کمترین قدر دی تلاش نال بنے فنکشناں نوں ورت کے~$d$
دا $a_0$، \dots،~$a_n$ توں \emph{حساب} لایا جا سکدا اے۔ ایہ صرف
ایہ آکھدا اے کہ اک «کوڈ کھولّن» والا فنکشن موجود اے جیہدے لئی
ہر سلسلے دا کوئی نہ کوئی «کوڈ» موجود اے۔

$\beta$ دی لکھت گوڈل توں آئی اے۔ فیر دھیان دیو کہ معاون قضیہ
ثابت کرن دا اوکھا کم بظاہر گھٹ وسائلاں نال مناسب~$\beta$
دی تعریف کرنا اے۔ ساڈے کول
صرف ترکیب تے کمترین قدر دی تلاش دے طریقے نیں، تے بنیادی فنکشناں
وچ جمع، ضرب تے~$\Char{=}$ وی ورت سکدے آں۔ ایس معاون قضیے دے کئی
ثبوت نیں۔ اج وی سبھ توں صاف طریقیاں وچوں اک گوڈل دا اپنا طریقہ
اے۔ اوس نے نظریۂ اعداد دا اک نتیجہ ورتیا جیہنوں سن زی دا قضیہ
آکھدے نیں (روایتی ناں «چینی باقیات دا قضیہ» اے)۔

\begin{defn}
دو قدرتی عدد $a$ تے $b$ اُدوں تے صرف اُدوں \emph{آپس وچ اوّل}
ہوندے نیں جدوں اوہناں دونوں نوں پورا ونڈن والا سبھ توں وڈا
عدد~$1$ ہووے۔ یعنی اوہناں دونوں نوں پورا ونڈن والا ہور کوئی
عدد نہیں۔
\end{defn}

\begin{defn}
قدرتی عدد $a$ تے $b$ \emph{مودولو~$c$ ہم باقی} ہوندے نیں،
$a \equiv b \mod c$، اُدوں تے صرف اُدوں جدوں $c \mid (a-b)$۔
یعنی $a$ تے $b$ نوں~$c$ نال ونڈیے تاں دونوں وار اکّو باقی
بچدی اے۔
\end{defn}

سن زی دا قضیہ ایہ اے:
\begin{thm}
فرض کرو $x_0$، \dots،~$x_n$ وچ ہر جوڑے دے عدد آپس وچ اوّل
نیں، تے $y_0$، \dots،~$y_n$ کوئی وی عدد نیں۔ تاں کوئی
عدد $z$ ایسا موجود اے کہ
\begin{align*}
z & \equiv y_0 \mod x_0 \\
z & \equiv y_1 \mod x_1 \\
& \vdots  \\
z & \equiv y_n \mod x_n.
\end{align*}
\end{thm}

سن زی دا قضیہ اسیں ایویں ورتاں گے۔ جے $x_0$، \dots،~$x_n$
بالترتیب $y_0$، \dots،~$y_n$ توں وڈے ہون، تاں $z$ نوں
سلسلے $\tuple{y_0, \dots,y_n}$ دا کوڈ بنا سکدے آں۔~$y_i$
کڈھن لئی $z$ نوں~$x_i$ نال ونڈ کے باقی لینی اے۔ ایس طریقے
لئی سانوں $x_0$، \dots، ~$x_n$ دیاں مناسب قدراں لبھنیاں پین گیاں۔

ایس کم وچ دو گلاں ساڈی مدد کرن گیاں۔
$y_0$، \dots،~$y_n$ دتے ہون،
تاں ایہ من لو:
\begin{align*}
j &= \max(n, y_0 + 1, \dots, y_n + 1), \\
m &= \lcm(1,\dots,j),
\end{align*}
تے ایہ من لو:
\begin{align*}
x_0 & = 1 + m \\
x_1 & = 1 + 2 \cdot m \\
x_2 & = 1 + 3 \cdot m \\
& \vdots  \\
x_n & = 1 + (n+1) \cdot m
\end{align*}
تاں دو گلّاں سچ نیں:
\begin{enumerate}
\item\ollabel{rel-prime} $x_0,\dots,x_n$ وچ ہر جوڑے دے عدد آپس وچ اوّل نیں۔
\item\ollabel{less} ہر $i$ لئی $y_i < x_i$۔
\end{enumerate}
\olref{rel-prime} لئی دو وکھ اشاریے لو۔ جے $p$ اولی عدد اے
تے $p \mid x_i$ تے $p \mid x_k$، تاں
$p \mid 1 + (i+1) m$ تے
$p \mid 1 + (k+1) m$۔
پر فیر $p$ اوہناں دے فرق
نوں پورا ونڈدا اے:
\[
(1 + (i+1)m) - (1+ (k+1)m) = (i-k) m.
\]
کیوں جے $p$، $1 + (i+1)m$ نوں پورا ونڈدا اے، ایس لئی
$m$ نوں پورا نہیں ونڈ سکدا۔ جے ونڈدا ہووے تاں پہلی ونڈ وچ
باقی~$1$ بچدی۔ ایس لئی $p$، $i-k$ نوں پورا ونڈدا اے، کیوں
جے $p$، $(i-k)m$ نوں پورا ونڈدا اے۔ پر $\left|i-k\right|$
ودھ توں ودھ~$n$ اے، تے اسیں $j \geq n$ چُنیا اے۔ ایس توں
$p \mid m$ نکلدا اے، جیہڑا تضاد اے۔ سو کوئی اولی عدد $x_i$
تے $x_k$ دونوں نوں پورا نہیں ونڈدا۔ شرط~\olref{less} سِدھی اے:
$y_i < j \leq m < x_i$۔

ہن $\beta$ فنکشن دا معاون قضیہ ثابت کریے۔ یاد رکھو کہ اسیں $0$،
جانشین، جمع، ضرب، $\Char{=}$ تے پروجیکشن ورت سکدے آں۔ اوہناں
توں ترکیب تے باقاعدہ فنکشناں اُتے کمترین قدر دی تلاش نال بنیا
کوئی وی فنکشن ورت سکدے آں۔ جے کسے تعلق دے امتیازی فنکشن دی
تعریف ایویں ہو سکے تاں اوہ تعلق وی ورت سکدے آں۔ پہلے وانگ
وکھا سکدے آں کہ بولیائی عملاں تے محدود مسوّریت نال بنے تعلق
وی ایس جماعت وچ رہندے نیں۔ مثال لئی:
\begin{align*}
\fn{not}(x) & \defis \Char{=}(x,0)\\
\bmin{x \leq z}{R(x,y)} & \defis \umin{x}{(R(x,y) \lor x = z)}\\
\bexists{x \leq z}{R(x,y)} & \defiff R(\bmin{x \leq z}{R(x,y)}, y)
\end{align*}
فیر اسیں دکھا سکدے
آں کہ ہِٹھلیاں
ساریاں شےواں وی
بنیادی بازگشت توں
بغیر تعریف ہو
سکدیاں نیں:
\begin{enumerate}
\item جوڑی فنکشن،
$J(x,y) = \frac{1}{2}[(x+y)(x+y+1)] + x$؛
% شاید ایتھے ہور دَسّنا چاہیدا اے کہ کیہ ہو رہیا اے؛گل کجھ الجھی اے۔
\item پروجیکشن فنکشن
\begin{align*}
K(z) & = \bmin{x \leq z}{\bexists{y \leq z}{z = J(x,y)}},\\
L(z) & = \bmin{y \leq z}{\bexists{x \leq z}{z = J(x,y)}};
\end{align*}
\item چھوٹا ہون دا تعلق $x < y$؛
\item پورا ونڈن دا تعلق $x \mid y$؛
% \item $x \tsub y$
% \item $\fn{Prime}(x)$
% \item فرض کرو $p$ اولی اے؛تعلق «$x$،$p$ دی اک قوت اے»:
% \[
% \bforall{y \leq x}{(y \mid x \lif y = 1 \lor y = x)}.
% \]
\item فنکشن $\fn{rem}(x,y)$، جیہڑا $y$ نوں~$x$ نال ونڈن توں بچی
  باقی دیندا اے۔
\end{enumerate}
ہن ایہ تعریف کرو:
\begin{align*}
\beta^*(d_0,d_1,i) & = \fn{rem}(1+(i+1) d_1,d_0) \text{ تے}\\
\beta(d,i) & = \beta^*(K(d),L(d),i).
\end{align*}
ایہو فنکشن سانوں چاہیدا اے۔
$a_0,\dots,a_n$ اُتّے وانگ
دتے ہون،تاں ایہ من لو:
\[
j = \max(n,a_0+1,\dots,a_n+1),
\]
تے $d_1 = \lcm(1,\dots,j)$
من لو۔ اُتّے دی
\olref{rel-prime}
مطابق سانوں پتا
اے کہ $1+d_1$،
$1+2 d_1$، \dots،
$1+(n+1) d_1$ وچ ہر جوڑے دے عدد آپس وچ اوّل نیں۔ نالے
~\olref{less} مطابق اوہ سارے $a_0,\dots,a_n$ توں وڈے نیں۔
سن زی دے قضیے
مطابق کوئی قدر~$d_0$ موجود اے جیہڑی ہر~$i$ لئی ایہ شرط پوری کردی اے۔
\[
d_0 \equiv a_i \mod (1+(i+1)d_1)
\]
تے سو (کیوں جے
$d_1$،~$a_i$
توں وڈّا اے):
\[
a_i = \fn{rem}(1+(i+1)d_1,d_0).
\]
$d = J(d_0,d_1)$
من لو۔ تاں ہر
$i \le n$
لئی ساڈے کول
ایہ اے:
\begin{align*}
\beta(d,i) & =  \beta^*(d_0,d_1,i) \\
& =  \fn{rem}(1+(i+1) d_1,d_0) \\
& =  a_i
\end{align*}
ایہو سانوں چاہیدا
سی۔ ایس نال
$\beta$-فنکشن دے
معاون قضیے دا
ثبوت مکمل اے۔

\begin{prob}
  وکھاؤ کہ تعلق $x < y$، $x \mid y$ تے فنکشن~$\fn{rem}(x,y)$
دیاں تعریفاں بنیادی بازگشت توں بغیر دِتیاں جا سکدیاں نیں۔ تسیں
$0$، جانشین، جمع، ضرب، $\Char{=}$، پروجیکشن، محدود کمترین قدر
دی تلاش تے محدود مسوّریت ورت سکدے او۔
\end{prob}
\endgroup
% /OLP-0292

% OLP-0293 content/incompleteness/representability-in-q/prim-rec.tex
\begingroup
\def\olfilename{prim-rec}
\expandafter\def\csname ol@sectioncs\endcsname{section}

\olfileid{inc}{req}{pri}
\olsection{بنیادی بازگشت دی نقل}

ہن اسیں وکھا سکدے آں کہ بیٹا فنکشن ورت کے بنیادی بازگشتی
تعریف دا کم باقاعدہ کمترین قدر دی تلاش نال ہو سکدا اے۔ فرض کرو ساڈے کول
$f(\vec x)$ تے $g(\vec x, y, z)$
نیں۔ تاں فنکشن~$h(x,\vec z)$ دی تعریف $f$ تے~$g$ توں بنیادی بازگشت
نال ایویں ہوندی اے۔
\begin{align*}
h(\vec x, 0) & =  f(\vec x) \\
h(\vec x, y+1) & =  g(\vec x, y, h(\vec x, y)).
\end{align*}
سانوں وکھانا اے کہ $h$ دی تعریف $f$ تے~$g$ توں صرف ترکیب
تے باقاعدہ کمترین قدر دی تلاش نال ہو سکدی اے۔ ایس کم وچ
بنیادی فنکشن تے اوہناں توں ایس طرح بنے فنکشن، جیویں~$\beta$،
ورت سکدے آں۔

\begin{lem}
\ollabel{lem:prim-rec}
جے $h$ دی تعریف $f$ تے $g$ توں بنیادی بازگشت نال ہو سکدی اے،
تاں $f$، $g$ تے فنکشناں $\Zero$، $\Succ$، $\Proj{n}{i}$،
$\Add$، $\Mult$، $\Char{=}$ توں صرف ترکیب تے باقاعدہ کمترین قدر
دی تلاش نال وی اوہدی تعریف ہو سکدی اے۔
\end{lem}

\begin{proof}
پہلاں مددگار فنکشن $\hat h(\vec x, y)$ دی تعریف کرو۔ ایہ
سبھ توں چھوٹا عدد~$d$ دیندا اے جیہدے لئی $d$ ایسے سلسلے دا
کوڈ اے جیہڑا ایہ شرطاں پوریاں کردا اے۔
\begin{enumerate}
\item $(d)_0 = f(\vec x)$،تے
\item ہر $i < y$ لئی
  $(d)_{i+1} = g(\vec x, i, (d)_i)$،
\end{enumerate}
ایتھے $(d)_i$، $\beta(d,i)$ دی مختصر لکھت اے۔ یعنی $\hat h$
سلسلے $\tuple{h(\vec x, 0), h(\vec x, 1), \dots, h(\vec x, y)}$
دا کوڈ دیندا اے۔
اسیں $\hat h$
نوں ایویں لکھ
سکدے آں:
\[
\hat h(\vec x, y) = \umin{d}{(\beta(d,0) = f(\vec x) \land \bforall{i <
  y}{\beta(d,i+1) = g(\vec x, i,\beta(d,i)})}.
\]
ایتھے بنیادی بازگشت دی لوڑ نہیں۔ صرف کمترین قدر دی تلاش ورتی
اے۔ جیہڑے فنکشن اُتے ایہ
تلاش لاندے آں، اوہ باقاعدہ اے، کیوں جے بیٹا فنکشن دے معاون
قضیے \olref[bet]{lem:beta} نال ہر دِتّی دلیل لئی ایسا کوڈ ملدا اے۔

ہن ساڈے کول ایہ مساوات اے۔
\[
h(\vec x, y) = \beta(\hat h(\vec x, y), y),
\]
ایس لئی $h$ دی تعریف دِتّے فنکشناں تے بنیادی فنکشناں توں صرف
ترکیب تے باقاعدہ کمترین قدر دی تلاش نال ہو سکدی اے۔
\end{proof}
\endgroup
% /OLP-0293

% OLP-0294 content/incompleteness/representability-in-q/basic-representable.tex
\begingroup
\def\olfilename{basic-representable}
\expandafter\def\csname ol@sectioncs\endcsname{section}

\olfileid{inc}{req}{bre}
\olsection{\texorpdfstring{بنیادی فنکشن~$\Th{Q}$ وچ قابلِ نمائندگی نیں}{بنیادی فنکشن Q وچ قابلِ نمائندگی نیں}}

پہلاں سانوں وکھانا اے کہ سارے بنیادی فنکشن~$\Th{Q}$ وچ
قابلِ نمائندگی نیں۔ ایس لئی سانوں وکھانا اے کہ ہر $k$-دلیلی
بنیادی فنکشن $f(x_0,\dots,x_{k-1})$ نال !!a{formula}
$!A_f(x_0,\dots,x_{k-1},y)$ کِویں جوڑیے جیہڑا اوہدی نمائندگی کرے۔

اسیں صفر، جانشین، جمع، ضرب، برابری دے امتیازی فنکشن تے
پروجیکشن فنکشناں دی نمائندگی کر سکاں گے۔ ہر صورت وچ نمائندگی
کرن والا فارمولا سِدھا اے۔ مثلاً صفر دی نمائندگی فارمولا
$y = \Obj 0$ کردا اے، جانشین دی نمائندگی !!{formula} $x_0' = y$
کردا اے، تے جمع دی نمائندگی !!{formula} $(x_0 + x_1) = y$
کردا اے۔ اصل کم ایہ وکھانا اے کہ $\Th{Q}$ متعلقہ !!{sentence}s
ثابت کر سکدا اے۔ مثلاً اُتلا !!{formula} جمع دی نمائندگی
کردا اے جے ایہ وکھائیے کہ قدرتی عدداں $m$ تے $n$ دے ہر جوڑے
لئی $\Th{Q}$ ایہ دونوں گلاں ثابت کردا اے۔
\begin{align*}
& \eq[\num n + \num m][\num {n+m}] \text{ تے}\\
& \lforall[y][(\eq[(\num n + \num m)][y] \lif \eq[y][\num{n+m}])].
\end{align*}

\begin{prop}
\ollabel{prop:rep-zero}
صفر فنکشن $\Zero(x) = 0$ دی~$\Th{Q}$ وچ نمائندگی فارمولا
$!A_{\Zero}(x,y) \ident \eq[y][\Obj 0]$ کردا اے۔
\end{prop}

\begin{prop}
\ollabel{prop:rep-succ}
جانشین فنکشن $\Succ(x) = x+1$ دی~$\Th{Q}$ وچ نمائندگی فارمولا
$!A_{\Succ}(x,y) \ident \eq[y][x']$ کردا اے۔
\end{prop}

\begin{prop}
\ollabel{prop:rep-proj}
پروجیکشن فنکشن $\Proj{n}{i}(x_0, \dots, x_{n-1}) = x_i$
دی~$\Th{Q}$ وچ نمائندگی ہِٹھلا فارمولا کردا اے۔
\[!A_{\Proj{n}{i}}(x_0, \dots, x_{n-1}, y) \ident \eq[y][x_i].\]
\end{prop}

\begin{prob}
ثابت کرو کہ $\eq[y][\Obj 0]$، $\eq[y][x']$ تے $\eq[y][x_i]$ بالترتیب
$\Zero$، $\Succ$ تے $\Proj{n}{i}$ دی نمائندگی کردے نیں۔
\end{prob}

\begin{prop}
\ollabel{prop:rep-id}
% PNB-CMP-054: case دی عام عبارت لئی اک text wrapper؛فارمولے قائم۔
برابری~$=$ دا امتیازی فنکشن،
\[
\Char{=}(x_0, x_1) =
\begin{cases}
  1 & \text{جے } x_0 =x_1\\
  0 & \text{نہیں تاں}
\end{cases}
\]
دی~$\Th{Q}$ وچ نمائندگی ہِٹھلا فارمولا کردا اے۔
\[
  !A_{\Char{=}}(x_0, x_1, y) \ident (\eq[x_0][x_1] \land \eq[y][\num{1}]) \lor (\eq/[x_0][x_1] \land
\eq[y][\num{0}]).
\]
\end{prop}

ثبوت لئی ہِٹھلا معاون قضیہ چاہیدا اے۔

\begin{lem}
\ollabel{lem:q-proves-neq} قدرتی عدداں $n$ تے $m$ لئی جے $n \neq m$ ہووے، تاں $\Th{Q} \Proves \eq/[\num n][\num m]$۔
\end{lem}

\begin{proof}
$n$ اُتے استقرا نال وکھاؤ کہ ہر $m$ لئی، جے $n \neq m$، تاں $Q \Proves \eq/[\num n][\num m]$۔

بنیادی صورت وچ $n = 0$ اے۔ جے $m$، $0$ دے برابر نہیں، تاں
$m = k + 1$، جتھے $k$ کوئی قدرتی عدد اے۔ ساڈے کول اک مسلّمہ اے جیہڑا کہندا اے
$\lforall[x][\eq/[0][x']]$۔ مسوّر دے اک مسلّمے نال، $x$ دی تھاں $\num k$ رکھ کے، اسیں $\eq/[0][\num k']$ کڈھ سکدے آں۔
پر $\num k'$ بس $\num m$ ای اے۔

استقرائی قدم وچ اسیں $n$ لئی دعویٰ درست فرض کر سکدے آں،
تے $n+1$ نوں ویکھدے آں۔ $m$ کوئی وی قدرتی عدد ہووے۔ دو
صورتاں نیں:یا $m = 0$،یا کسے $k$ لئی $m = k+1$۔ پہلی صورت دا ثبوت اُتلے طریقے نال ہو جاندا اے۔ دوجی صورت وچ
فرض کرو $n+1 \neq
k+1$۔ تاں $n \neq k$۔ $n$ لئی استقرائی مفروضے نال ساڈے
کول $\Th{Q} \Proves \eq/[\num n][\num k]$ اے۔ ساڈے کول
اک مسلّمہ اے جیہڑا کہندا اے
$\lforall[x][\lforall[y][\eq[x'][y'] \lif \eq[x][y]]]$۔ مسوّر دا اک مسلّمہ ورت کے ساڈے کول
$\eq[\num n'][\num k'] \lif \eq[\num n][\num
  k]$ اے۔ قضیاتی منطق نال اسیں $\Th{Q}$ وچ ایہ کڈھ سکدے آں:
$\eq/[\num n][\num k] \lif \eq/[\num n'][\num k']$۔
اثباتِ مقدم دے قاعدے نال اسیں $\eq/[\num n'][\num k']$ کڈھ سکدے آں؛
ایہو سانوں چاہیدا اے، کیوں جے $\num k'$،$\num m$ اے۔
\end{proof}

\begin{explain}
دھیان دیو کہ معاون قضیہ بنیادی گل ایہ آکھدا اے کہ $\Th{Q}$ ثابت کر سکدا اے کہ وکھریاں
عددی عبارتاں وکھریاں شےواں دسدیاں نیں۔ مثال لئی،$\Th{Q}$
ایہ ثابت کردا اے: $0'' \neq 0'''$۔ پر ایہ دکھاؤن لئی کہ
گل ہر صورت وچ درست اے،ذرا احتیاط چاہیدی اے۔ ایہ وی دھیان رکھو کہ اسیں جیہڑا استقرا ورت رہے آں، اوہ
$\Th{Q}$ دے \emph{باہر} اے۔
\end{explain}

\begin{proof}[\olref{prop:rep-id} دا ثبوت]
جے $n = m$ ہووے، تاں $\num{n}$ تے $\num{m}$ اکّو ٹرم نیں،
تے $\Char{=}(n, m) = 1$۔ پر
$\Th{Q} \Proves (\eq[\num{n}][\num{m}] \land
\eq[\num{1}][\num{1}])$،سو اوہ
$!A_=(\num{n}, \num{m},
\num{1})$ ثابت کردا اے۔ جے $n \neq m$ ہووے، تاں
$\Char=(n, m) = 0$۔ \olref{lem:q-proves-neq} نال،
$\Th{Q} \Proves \eq/[\num{n}][\num{m}]$،تے ایس لئی
$(\eq/[\num{n}][\num{m}] \land \Obj 0 = \Obj 0)$ وی۔
سو $\Th{Q}
\Proves !A_=(\num{n}, \num{m}, \num{0})$۔

دوجے حصے لئی وی دو صورتاں نیں۔ جے $n = m$،سانوں
دکھانا اے کہ $\Th{Q} \Proves \lforall[y][(!A_=(\num{n}, \num{m}, y)
  \lif \eq[y][\num{1}])]$۔ غیر رسمی دلیل وچ فرض کرو
$!A_=(\num{n},
\num{m}, y)$،یعنی:
\[
(\eq[\num{n}][\num{n}] \land \eq[y][\num{1}]) \lor
(\eq/[\num{n}][\num{n}] \land \eq[y][\num{0}])
\]
فصل دا کھبا جز منطق نال $\eq[y][\num{1}]$ دا نتیجہ دیندا اے؛
اوہدا سجا جز $\eq[\num{n}][\num{n}]$ نال ٹکراندا اے،جیہڑا
منطق نال ثابت ہو سکدا اے۔

دوجی صورت وچ فرض کرو $n \neq m$۔ تاں
$!A_=(\num{n},
\num{m}, y)$ ایہ اے:
\[
(\eq[\num{n}][\num{m}] \land \eq[y][\num{1}]) \lor
(\eq/[\num{n}][\num{m}] \land \eq[y][\num{0}])
\]
ایتھے فصل دا کھبا جز $\eq/[\num{n}][\num{m}]$ نال ٹکراندا اے،
جیہڑا \olref{lem:q-proves-neq} نال $\Th{Q}$ وچ ثابت ہو سکدا اے؛
فصل دا سجا جز $\eq[y][\num{0}]$ دا نتیجہ دیندا اے۔
\end{proof}

\begin{prop}
\ollabel{prop:rep-add}
جمع فنکشن $\Add(x_0, x_1) = x_0+x_1$ دی~$\Th{Q}$ وچ نمائندگی ہِٹھلا فارمولا کردا اے۔
\[
  !A_{\Add}(x_0, x_1, y) \ident \eq[y][(x_0 + x_1)].
\]
\end{prop}

\begin{lem}
\ollabel{lem:q-proves-add}
$\Th{Q} \Proves \eq[(\num{n} + \num{m})][\num{n+m}]$
\end{lem}

\begin{proof}
اسیں ایہ $m$ اُتے استقرا نال ثابت کردے آں۔ جے $m = 0$،
دعویٰ ایہ اے کہ
$\Th{Q} \Proves \eq[(\num{n} + \Obj 0)][\num{n}]$۔ ایہ
مسلّمے~$!Q_4$ توں نکلدا اے۔ ہن $m$ لئی دعویٰ فرض کرو؛
ہن $m+1$ لئی دعویٰ ثابت کریے۔ یعنی وکھائیے کہ
$\Th{Q} \Proves \eq[(\num{n} +
  \num{m+1})][\num{n+m+1}]$۔ دھیان دیو کہ $\num{m+1}$
بس $\num{m}'$ اے، تے $\num{n+m+1}$ بس $\num{n+m}'$ اے۔
مسلّمے $!Q_5$ نال،$\Th{Q}
\Proves \eq[(\num{n} + \num{m}')][(\num{n}+\num{m})']$۔
استقرائی مفروضے نال،
$\Th{Q} \Proves \eq[(\num{n} + \num{m})][\num{n+m}]$۔ سو
$\Th{Q} \Proves \eq[(\num{n} + \num{m}')][\num{n+m}']$۔
\end{proof}

\begin{proof}[\olref{prop:rep-add} دا ثبوت]
!!{formula}~$!A_\Add(x_0, x_1, y)$ جیہڑا $\Add$ دی نمائندگی کردا اے، ایہ اے
$\eq[y][(x_0 + x_1)]$۔ پہلاں اسیں دکھاندے آں کہ جے
$\Add(n, m) = k$،تاں
$\Th{Q} \Proves !A_\Add(\num{n}, \num{m}, \num{k})$،
یعنی $\Th{Q}
\Proves \eq[\num{k}][(\num{n} + \num{m})]$۔ پر کیوں جے
$k = n + m$، ایس لئی $\num{k}$ عین $\num{n+m}$ اے،تے اسیں
\olref{lem:q-proves-add} وچ ایہ دکھا چُکے آں:
$\Th{Q} \Proves \eq[(\num{n} +
  \num{m})][\num{n+m}]$۔

سانوں ایہ وی دکھانا اے کہ جے $\Add(n, m) = k$،تاں
\[
\Th{Q} \Proves \lforall[y][(!A_\Add(\num{n}, \num{m}, y) \lif
  \eq[y][\num{k}])].
\]
فرض کرو ساڈے کول $\eq[(\num{n} + \num{m})][y]$ اے۔
کیوں جے
\[
\Th{Q} \Proves \eq[(\num{n}+\num{m})][\num{n+m}],
\]
اسیں کھبے پاسے دی تھاں $\num{n+m}$ رکھ کے
$\eq[\num{n+m}][y]$ حاصل کر سکدے آں،ہر~$y$ لئی۔
\end{proof}

\begin{prop}
\ollabel{prop:rep-mult}
ضرب فنکشن $\Mult(x_0, x_1) = x_0 \cdot x_1$ دی~$\Th{Q}$
وچ نمائندگی ہِٹھلا فارمولا کردا اے۔
\[
  !A_{\Mult}(x_0, x_1, y) \ident y = (x_0 \times x_1).
\]
\end{prop}

\begin{proof} مشق۔ \end{proof}

\begin{lem}
\ollabel{lem:q-proves-mult}
$\Th{Q} \Proves \eq[(\num{n} \times \num{m})][\num{n \cdot m}]$
\end{lem}

\begin{proof} مشق۔ \end{proof}

\begin{prob}
\olref[inc][req][bre]{lem:q-proves-mult} ثابت کرو۔
\end{prob}

\begin{prob}
\olref[inc][req][bre]{lem:q-proves-mult} ورت کے
\olref[inc][req][bre]{prop:rep-mult} ثابت کرو۔
\end{prob}

\begin{explain}
  یاد رکھو کہ اسیں حساب دی زبان دے فنکشن نشان لئی $\times$
ورتدے آں تے عدداں اُتے عام ضرب دے عمل لئی $\cdot$۔ ایس لئی
$\cdot$ عدد دسن والیاں عبارتاں وچکار آ سکدا اے، جیویں
$m \cdot n$ وچ۔ $\times$ صرف حساب دی زبان دے ٹرماں وچکار
آندا اے، جیویں $(\num{m} \times \num{n})$ وچ۔ ہور وی اُلجھن
والی گل ایہ اے کہ $+$، !!{function} تے جمع دے عمل، دونوں لئی
ورتیا جاندا اے۔ جدوں ایہ ٹرماں وچکار آوے، جیویں
$(\num{n} + \num{m})$ وچ، تاں ایہ حساب دی زبان دا $2$-دلیلی
!!{function} اے۔ جدوں ایہ عدداں وچکار آوے، جیویں $n+m$ وچ،
تاں ایہ جمع دا عمل اے۔ $\num{n+m}$ وچ وی ایہو جمع اے۔
ایہ عدد~$n+m$ دی معیاری عددی عبارت اے۔
\end{explain}
\endgroup
% /OLP-0294

% OLP-0295 content/incompleteness/representability-in-q/composition-representable.tex
\begingroup
\def\olfilename{composition-representable}
\expandafter\def\csname ol@sectioncs\endcsname{section}

\olfileid{inc}{req}{cmp}
\olsection{\texorpdfstring{ترکیب $\Th{Q}$ وچ قابلِ نمائندگی اے}{ترکیب Q وچ قابلِ نمائندگی اے}}

فرض کرو $h$ دی تعریف ایہ اے:
\[
h(x_0,\dots,x_{l-1}) = f(g_0(x_0,\dots,x_{l-1}), \dots,
g_{k-1}(x_0,\dots,x_{l-1})).
\]
اسیں پہلے ای !!{formula}s $!A_f, !A_{g_0}, \dots, !A_{g_{k-1}}$
لبھ چکے آں جیہڑے بالترتیب فنکشناں $f$ تے $g_0$، \dots،~$g_{k-1}$
دی نمائندگی کردے نیں۔ سانوں !!a{formula}~$!A_h$ لبھنا اے جیہڑا
$h$ دی نمائندگی کرے۔

پہلاں اوہ سوکھی صورت ویکھیے جتھے سارے فنکشن $1$-دلیلی نیں،
یعنی $h(x) = f(g(x))$۔ جے $!A_f(y, z)$،$f$
دی نمائندگی کردا اے،تے $!A_g(x, y)$،$g$ دی نمائندگی کردا اے،
تاں سانوں !!a{formula}~$!A_h(x, z)$ چاہیدا اے جیہڑا $h$ دی
نمائندگی کرے۔ دھیان دیو کہ $h(x) = z$ اُدوں تے صرف اُدوں اے جدوں
کوئی~$y$ ایسا ہووے جیہدے لئی $z = f(y)$ تے $y = g(x)$
دونوں درست ہون۔ جے $h(x) = z$ ہووے، تاں $g(x)$ ایسا~$y$ اے۔
جے کوئی ایسا $y$ موجود ہووے، تاں $y = g(x)$ تے $z = f(y)$
توں $z = f(g(x))$ نکلدا اے۔ ایس توں لگدا اے کہ
$\lexists[y][(!A_g(x, y) \land !A_f(y,
  z))]$،$!A_h(x, z)$ لئی اک مناسب فارمولا لگدا اے۔ سانوں بس ایہ پرکھنا اے کہ $\Th{Q}$ متعلقہ !!{formula}s ثابت کردا اے۔

\begin{prop}
\ollabel{prop:rep1}
جے $h(n) = m$ ہووے، تاں $\Th{Q} \Proves !A_h(\num{n}, \num{m})$۔
\end{prop}

\begin{proof}
فرض کرو $h(n) = m$، یعنی $f(g(n)) = m$۔ $k = g(n)$ من لو۔ ایس توں ایہ ملدا اے۔
\begin{align*}
  \Th{Q} & \Proves !A_g(\num{n}, \num{k})
  \intertext{کیوں جے $!A_g$، $g$ دی نمائندگی کردا اے، تے}
  \Th{Q} & \Proves !A_f(\num{k}, \num{m})
  \intertext{کیوں جے $!A_f$، $f$ دی نمائندگی کردا اے۔ ایس لئی}
  \Th{Q} & \Proves !A_g(\num{n}, \num{k}) \land !A_f(\num{k}, \num{m})
  \intertext{تے ایس دے نتیجے وچ ایہ وی:}
  \Th{Q} & \Proves \lexists[y][(!A_g(\num{n}, y) \land !A_f(y, \num{m}))],
\end{align*}
یعنی $\Th{Q} \Proves !A_h(\num{n}, \num{m})$۔
\end{proof}

\begin{prop}
\ollabel{prop:rep2}
جے $h(n) = m$ ہووے، تاں $\Th{Q} \Proves \lforall[z][(!A_h(\num{n}, z) \lif
  z = \num{m})]$۔
\end{prop}

\begin{proof}
فرض کرو $h(n) = m$، یعنی $f(g(n)) = m$۔ $k = g(n)$ من لو۔ ایس توں ایہ ملدا اے۔
\begin{align*}
  \Th{Q} & \Proves \lforall[y][(!A_g(\num{n}, y) \lif \eq[y][\num{k}])]
  \intertext{کیوں جے $!A_g$، $g$ دی نمائندگی کردا اے، تے}
  \Th{Q} & \Proves \lforall[z][(!A_f(\num{k}, z) \lif \eq[z][\num{m}])]
  \intertext{کیوں جے $!A_f$،$f$ دی نمائندگی کردا اے۔ تھوڑی جیہی منطق نال ایہ وی وکھا سکدے آں۔}
  \Th{Q} & \Proves \lforall[z][(\lexists[y][(!A_g(\num{n}, y) \land
      \formula A_f(y, z))] \lif \eq[z][\num{m}])].
\end{align*}
یعنی $\Th{Q} \Proves \lforall[y][(!A_h(\num n, y) \lif \eq[y][\num m])]$۔
\end{proof}

ایہو خیال اوکھی صورت وچ وی چلدا اے جتھے $f$ تے~$g_i$ دیاں
دلیلاں دی تعداد $1$ توں ودھ اے۔

\begin{prop}
\ollabel{prop:rep-composition}
جے $!A_f(y_0, \dots, y_{k-1}, z)$،$f(y_0, \dots, y_{k-1})$
دی~$\Th{Q}$ وچ نمائندگی کردا اے،تے
$!A_{g_i}(x_0, \dots, x_{l-1}, y)$،
$g_i(x_0, \dots, x_{l-1})$ دی~$\Th{Q}$ وچ نمائندگی کردا اے،تاں
% PNB-CMP-055: دونیں quantifier scopes مکمل conjunction دے آخر تک۔
\begin{multline*}
  \lexists[y_0\dots][\lexists[y_{k-1}][(!A_{g_0}(x_0,\dots,x_{l-1},y_0) \land
      \dots \land {}\\
  \formula A_{g_{k-1}}(x_0,\dots,x_{l-1},y_{k-1}) \land !A_f(y_0,\dots,y_{k-1},z))]]
\end{multline*}
ہِٹھ لکھے فنکشن دی نمائندگی کردا اے۔
\[
h(x_0, \dots, x_{l-1}) = f(g_0(x_0, \dots, x_{l-1}), \dots, g_{k-1}(x_0,
\dots, x_{l-1})).
\]
\end{prop}

\begin{proof}
مشق۔
\end{proof}

\begin{prob}
% PNB-CMP-056: قدر تے یکتائی دے دونوں ثبوت guide نیں،اک دہرایا حوالہ نہیں۔
\olref[inc][req][cmp]{prop:rep1} تے
\olref[inc][req][cmp]{prop:rep2} دے ثبوت سامنے رکھ کے
\olref[inc][req][cmp]{prop:rep-composition} دا پورا ثبوت دیو۔
\end{prob}
\endgroup
% /OLP-0295

% OLP-0296 content/incompleteness/representability-in-q/minimization-representable.tex
\begingroup
\def\olfilename{minimization-representable}
\expandafter\def\csname ol@sectioncs\endcsname{section}

\olfileid{inc}{req}{min}
% PNB-CMP-057: step6 وچ درمیانی برابری؛آخری سطر وچ مطلوبہ نتیجہ۔
% PNB-CMP-058: وجودی تعمیم دا گواہ c اے،c' نہیں۔
% PNB-CMP-059: آخری بیان وچ شروع دا باقاعدگی دا فرض صاف کیتا۔
\olsection{\texorpdfstring{باقاعدہ کمترین قدر دی تلاش $\Th{Q}$ وچ قابلِ نمائندگی اے}{باقاعدہ کمترین قدر دی تلاش Q وچ قابلِ نمائندگی اے}}

بغیر مقررہ حد دی تلاش اُتے غور کریے۔ فرض کرو $g(x, z)$ باقاعدہ
اے تے $\Th{Q}$ وچ قابلِ نمائندگی اے، مثلاً !!{formula}~$!A_g(x, z, y)$
نال۔ $f$ دی تعریف $f(z) = \umin{x}{[g(x, z) = 0]}$ نال کریے۔
اسیں کوئی !!a{formula}~$!A_f(z, y)$ لبھنا چاہندے آں جیہڑا~$f$
دی نمائندگی کرے۔~$f(z)$ دی قدر اوہ عدد~$x$ اے جیہڑا (a) شرط
$g(x, z) = 0$ پوری کردا اے تے (b) ایہ شرط پوری کرن والے عدداں
وچوں سبھ توں چھوٹا اے۔ یعنی ہر $w < x$ لئی $g(w, z) \neq 0$۔
ایس لئی ہِٹھلا فارمولا چُننا سِدھی گل اے۔
\[
!A_f(z,y) \ident !A_g(y, z, \Obj 0) \land \lforall[w][(w < y \lif \lnot
  !A_g(w, z, \Obj 0))].
\]
عام صورت وچ سانوں $z$ دی تھاں
$z_0$، \dots، $z_k$ رکھنے پین گے۔

ثبوت لئی فیر کجھ ایہو جیہے معاون قضیے چاہیدے ہون گے جیہڑے
$\Th{Q}$ ثابت کر سکدا اے۔

\begin{lem}
\ollabel{lem:succ} ہر !!{constant}~$a$ تے ہر قدرتی
عدد~$n$ لئی،
\[
\Th{Q} \Proves \eq[(a' + \num n)][(a + \num n)'].
\]
\end{lem}

\begin{proof}
پہلاں وانگ ثبوت $n$ اُتے استقرا نال اے۔ بنیادی صورت $n =
0$ وچ سانوں وکھانا اے کہ $\Th{Q}$، $\eq[(a' + \Obj 0)][(a + \Obj
0)']$ ثابت کردا اے۔پر ساڈے کول:
\begin{align}
  \Th{Q} & \Proves \eq[(a' + \Obj 0)][a'] \quad \text{مسلّمے $Q_4$ نال}
  \ollabel{step1}\\
  \Th{Q} & \Proves  \eq[(a + \Obj 0)][a] \quad \text{مسلّمے $Q_4$ نال}
  \ollabel{step2} \\
  \Th{Q} & \Proves \eq[(a + \Obj 0)'][a'] \quad \text{\olref{step2} نال}
  \ollabel{step3} \\
  \Th{Q} & \Proves \eq[(a' + \Obj 0)][(a + \Obj 0)'] \quad
  \text{\olref{step1} تے \olref{step3} نال}\notag
\end{align}
استقرائی قدم وچ اسیں فرض کر سکدے آں کہ $\Th{Q}
\Proves \eq[(a' + \num n)][(a + \num n)']$ پہلاں ثابت ہو چکیا اے۔ کیوں جے
$\num{n+1}$، $\num{n}'$ اے،سانوں وکھانا اے کہ $\Th{Q}$،
$\eq[(a' + \num{n}')][(a + \num{n}')']$ ثابت کردا اے۔ساڈے کول:
\begin{align}
  \Th{Q} & \Proves \eq[(a' + \num n')][(a' + \num n)'] \quad
  \text{مسلّمے $!Q_5$ نال} \ollabel{step5}\\
  \Th{Q} & \Proves \eq[(a' + \num n)'][(a + \num n')'] \quad
  \text{استقرائی مفروضے،برابری دی ہم آہنگی تے جمع دے جانشینی مسلّمے نال} \ollabel{step6}\\
  \Th{Q} & \Proves \eq[(a' + \num n')][(a + \num n')'] \quad
  \text{\olref{step5} تے \olref{step6} نال۔} \notag
\end{align}
\end{proof}

یاد رکھو کہ ایہ نتیجہ ایس گل نالوں کمزور اے کہ
$\Th{Q}$، $\lforall[x][\lforall[y][(x' + y) = (x + y)']]$ ثابت کردا اے۔
ایہ !!{sentence}~$\Struct{N}$ وچ سچ اے، پر $\Th{Q}$ اینوں
ثابت نہیں کردا۔

\begin{lem}
\ollabel{lem:less-zero}
$\Th{Q} \Proves \lforall[x][\lnot x < \Obj 0]$.
\end{lem}

\begin{proof}
اسیں ثبوت غیر رسمی ڈھنگ نال دیندے آں، یعنی رسمی !!{derivation} بنان دے صرف اشارے دیندے آں۔

سانوں $\lnot a < \Obj 0$ کسے وی~$a$ لئی ثابت کرنا اے۔ $<$
دی تعریف مطابق سانوں $\lnot
\lexists[y][\eq[(y'+a)][\Obj 0]]$ نوں $\Th{Q}$ وچ ثابت کرنا اے۔اسیں $\lexists[y][\eq[(y'+a)][\Obj 0]]$ فرض کر کے تضاد کڈھاں گے۔فرض کرو
$\eq[(b'+a)][\Obj 0]$۔$!Q_3$ ورت کے سانوں
$\eq[a][\Obj 0] \lor \lexists[y][\eq[a][y']]$ ملدا اے۔ہن دوہاں صورتاں نوں وکھ وکھ ویکھدے آں۔

پہلی صورت وچ $\eq[a][\Obj 0]$ قائم اے۔$\eq[(b'+a)][\Obj 0]$ توں سانوں
$\eq[(b' + \Obj 0)][\Obj 0]$ ملدا اے۔مسلّمے~$!Q_4$ نال،جیہڑا $\Th{Q}$ دا اے،ساڈے کول
$\eq[(b' + \Obj 0)][b']$ اے،تے ایس لئی $\eq[b'][\Obj 0]$۔پر
مسلّمے~$!Q_2$ نال ساڈے کول $\eq/[b'][\Obj 0]$ وی اے۔ ایہ تضاد اے۔

دوجی صورت وچ کسے $c$ لئی $\eq[a][c']$۔پر فیر ساڈے کول $\eq[(b' +
c')][\Obj 0]$ اے۔مسلّمے~$!Q_5$ نال،سانوں $\eq[(b' + c)'][\Obj 0]$ ملدا اے۔ ایہ فیر مسلّمے~$Q_2$ دے خلاف اے۔
\end{proof}

\begin{lem}
\ollabel{lem:less-nsucc}
ہر قدرتی عدد~$n$ لئی،
  \[
  \Th{Q} \Proves
  \lforall[x][(x < \num {n+1} \lif (\eq[x][\Obj 0] \lor \dots \lor
    \eq[x][\num n]))].
  \]
\end{lem}

\begin{proof}
اسیں~$n$ اُتے استقرا ورتدے آں۔ پہلاں $n = 0$ والی بنیادی صورت ویکھیے۔
ایس صورت وچ،سانوں $a < \num 1 \lif \eq[a][\Obj 0]$ کسے وی~$a$ لئی
وکھانا اے۔فرض کرو $a < \num 1$۔فیر $<$ دے تعریفی مسلّمے نال
ساڈے کول $\lexists[y][\eq[(y'+a)][\Obj 0']]$ اے،کیوں جے $\num 1 \ident
\Obj 0'$۔

فرض کرو $b$ ایہ شرط پوری کردا اے، یعنی ساڈے کول $\eq[(b'+a)][\Obj 0']$ اے۔
سانوں $\eq[a][\Obj 0]$ وکھانا اے۔مسلّمے~$!Q_3$ نال،یا تاں
$\eq[a][\Obj 0]$ اے،یا کسے $c$ لئی $\eq[a][c']$۔پہلی صورت وچ نتیجہ مل گیا۔ ہن فرض کرو $\eq[a][c']$۔فیر
ساڈے کول $\eq[(b' + c')][\Obj 0']$ اے۔مسلّمے~$!Q_5$ نال،جیہڑا $\Th{Q}$ دا اے،سانوں
$\eq[(b'+c)'][\Obj 0']$ ملدا اے۔مسلّمے~$!Q_1$ نال،سانوں $\eq[(b' + c)][\Obj
0]$ ملدا اے۔پر مسلّمے~$!Q_8$ نال،ایس دا مطلب $c < \Obj 0$ اے،جیہڑا
\olref{lem:less-zero} دے خلاف اے۔

ہن استقرائی قدم لئی $n+1$ دی صورت ثابت کریے،~$n$ دی صورت نوں من کے۔ فرض کرو $a < \num {n+2}$۔فیر $!Q_3$ نال دو صورتاں ملدیاں نیں۔ $\eq[a][\Obj 0]$ تے کسے $b$ لئی
$\eq[a][b']$۔پہلی صورت وچ $\eq[a][\Obj 0] \lor \dots \lor
\eq[a][\num{n+1}]$ فوراً نکلدا اے۔دوجی صورت وچ ساڈے کول $b'
< \num {n+2}$،یعنی $b' < \num{n+1}'$ اے۔مسلّمے~$!Q_8$ نال،کسے $c$ لئی
$\eq[(c'+b')][\num{n+1}']$۔مسلّمے $!Q_5$ نال
$\eq[(c'+b)'][\num{n+1}']$۔مسلّمے~$!Q_1$ نال $\eq[(c'+b)][\num{n+1}]$،
تے ایس لئی $b < \num{n+1}$،مسلّمے~$!Q_8$ نال۔استقرائی مفروضے نال
$\eq[b][\Obj 0] \lor \dots \lor \eq[b][\num{n}]$۔منطق نال ایس توں
$\eq[b'][\Obj 0'] \lor \dots \lor \eq[b'][\num{n}']$ ملدا اے،تے ایس لئی
$\eq[a][\num{1}] \lor \dots \lor \eq[a][\num{n+1}]$،کیوں جے
$\eq[a][b']$۔
\end{proof}

\begin{lem}
  \ollabel{lem:trichotomy} ہر قدرتی عدد~$m$ لئی،
  \[
  \Th{Q} \Proves
  \lforall[y][((y < \num{m} \lor \num{m} < y) \lor \eq[y][\num{m}])].
  \]
\end{lem}

\begin{proof}
$m$ اُتے استقرا ورتدے آں۔ پہلاں $m=0$ دی صورت اُتے غور کریے۔$\Th{Q} \Proves
\lforall[y][(\eq[y][\Obj 0] \lor \lexists[z][\eq[y][z']])]$،~$!Q_3$ نال۔
کوئی وی $a$ لو۔ فیر یا تاں $\eq[a][\Obj 0]$ اے،یا کسے~$b$ لئی
$\eq[a][b']$۔پہلی صورت وچ،ساڈے کول $(a < \Obj 0 \lor \Obj
0 < a) \lor \eq[a][\Obj 0]$ وی اے۔پر جے $\eq[a][b']$،تاں
$\eq[(b' + \Obj 0)][(a + \Obj 0)]$،~$\eq$ دی منطق نال۔$!Q_4$ نال $\eq[(a +
\Obj 0)][a]$،سو ساڈے کول $\eq[(b' + \Obj 0)][a]$ اے،تے ایس لئی
$\lexists[z][\eq[(z' + \Obj 0)][a]]$۔$<$ دی
تعریف،جیہڑی $!Q_8$ وچ اے،نال $\Obj 0 < a$۔جے $\Obj 0 < a$،تاں $(\Obj 0 < a \lor a
< \Obj 0) \lor \eq[a][\Obj 0]$ وی۔

ہن استقرائی مفروضہ من لو۔ ساڈے کول
\begin{align*}
  \Th{Q} & \Proves \lforall[y][((y < \num{m} \lor \num{m} < y) \lor
    \eq[y][\num{m}])]
  \intertext{اے تے اسیں وکھانا چاہندے آں کہ}
  \Th{Q} & \Proves \lforall[y][((y < \num{m+1} \lor \num{m+1} < y) \lor
    \eq[y][\num{m+1}])]
\end{align*}
کوئی وی $a$ لو۔ $!Q_3$ نالیا تاں $\eq[a][\Obj 0]$ اے،یا
کسے~$b$ لئی $\eq[a][b']$۔پہلی صورت وچ،ساڈے کول $\eq[\num{m}' +
a][\num{m+1}]$ اے،~$!Q_4$ نال،تے ایس لئی $a < \num{m+1}$،~$!Q_8$ نال۔

ہن دوجی صورت $\eq[a][b']$ ویکھیے۔ استقرائی مفروضے نال $(b < \num{m} \lor \num{m} < b) \lor
    \eq[b][\num{m}]$۔

فصل دا پہلا جز $b < \num{m}$،$!Q_8$ نال،
$\lexists[z][\eq[(z' + b)][\num{m}]]$ دے برابر اے۔ فرض کرو $c$ وچ ایہ خاصیت اے۔
جے $\eq[(c' + b)][\num{m}]$،تاں $\eq[(c' + b)'][\num{m}']$ وی۔$!Q_5$
نال $\eq[(c' + b)'][(c' + b')]$۔ایس لئی $\eq[(c' +
b')][\num{m}']$۔اسیں $\lexists[u][\eq[(u' + b')][\num{m+1}]]$،~$c$ نوں وجودی تعمیم دا گواہ بنا کے حاصل کردے آں،تے ایہ یاد رکھدے آں کہ $\num{m}'
\ident \num{m+1}$۔ایس لئی جے $b < \num{m}$،تاں $b' < \num{m+1}$،
تے ایس لئی $a < \num{m+1}$۔

دوجے جز لئی فرض کرو $\num{m} < b$،یعنی $\lexists[z][\eq[(z' +
\num{m})][b]]$۔فرض کرو $c$ ایسا~$z$ اے، یعنی $\eq[(c' +
\num{m})][b]$۔منطق نال $\eq[(c' + \num{m})'][b']$۔$!Q_5$ نال
$\eq[(c' + \num{m}')][b']$۔کیوں جے $\eq[a][b']$ تے $\num{m}' \ident
\num{m+1}$،ایس لئی $\eq[(c' + \num{m+1})][a]$۔$!Q_8$ نال $\num{m+1} < a$۔

آخری جز لئی فرض کرو $\eq[b][\num{m}]$۔فیر منطق نال
$\eq[b'][\num{m}']$،تے ایس لئی $\eq[a][\num{m+1}]$۔

ایس لئی~$m$ تے~$b$ دی صورت وچ فصل دے ہر جز توں اسیں
~$m+1$ تے~$a$ لئی اوسے نال ملدا جز کڈھ سکدے آں۔
\end{proof}
  
\begin{prop}
\ollabel{prop:rep-minimization}
جے $!A_g(x, z, y)$ باقاعدہ فنکشن $g(x, z)$ دی~$\Th{Q}$ وچ نمائندگی کردا اے، تاں
\[
!A_f(z,y) \ident !A_g(y, z, \Obj 0) \land \lforall[w][(w < y \lif \lnot
  !A_g(w, z, \Obj 0))]
\]
$f(z) = \umin{x}{[g(x, z) = 0]}$ دی نمائندگی کردا اے۔
\end{prop}

\begin{proof}
پہلاں وکھائیے کہ جے $f(n) = m$ ہووے، تاں $\Th{Q} \Proves !A_f(\num n, \num m)$،
یعنی
\begin{align}
  \Th{Q} & \Proves !A_g(\num{m}, \num{n}, \Obj 0) \land \lforall[w][(w <
    \num{m} \lif \lnot !A_g(w, \num{n}, \Obj 0))]. \notag
  \intertext{کیوں جے $!A_g(x, z, y)$، $g(x, z)$ دی نمائندگی کردا اے تے $g(m, n) =
    0$ جدوں $f(n) = m$،ایس لئی ساڈے کول}
\Th{Q} & \Proves !A_g(\num{m}, \num{n}, \Obj 0). \notag
\intertext{جے $f(n) = m$،تاں ہر $k < m$ لئی $g(k, n) \neq 0$۔سو}
\Th{Q} & \Proves \lnot !A_g(\num{k}, \num{n}, \Obj 0). \notag
\intertext{سانوں}
\Th{Q} & \Proves \lforall[w][(w < \num{m} \lif \lnot
  \formula A_g(w, \num{n}, \Obj 0))]. \ollabel{rep-less}
\end{align}
ملدا اے۔ $m = 0$ دی صورت وچ \olref{lem:less-zero} نال تے دوجی صورت وچ \olref{lem:less-nsucc} نال۔

ہن وکھائیے کہ جے $f(n) = m$ ہووے، تاں $\Th{Q} \Proves
\lforall[y][(!A_f(\num{n}, y) \lif \eq[y][\num{m}])]$۔اسیں فیر دلیل دے اشارے دیندے آں۔ رسمی ثبوت بناؤن دا کم پڑھن والے لئی چھڈدے آں۔

فرض کرو $!A_f(\num{n}, b)$۔ایس توں سانوں (a) $!A_g(b, \num{n},
\Obj 0)$ تے (b) $\lforall[w][(w < b \lif \lnot !A_g(w, \num{n}, \Obj
  0))]$ ملدے نیں۔\olref{lem:trichotomy} نال $(b < \num{m} \lor \num{m} < b)
\lor \eq[b][\num{m}]$۔اسیں وکھاواں گے کہ $b < \num{m}$ تے $\num{m}
< b$وچوں ہر اک صورت وچ تضاد نکلدا اے۔

جے $\num{m} < b$، تاں~(b) توں $\lnot !A_g(\num{m}, \num{n}, \Obj 0)$
ملدا اے۔پر $m = f(n)$،سو $g(m, n) = 0$،تے ایس لئی $\Th{Q} \Proves
!A_g(\num{m}, \num{n}, \Obj 0)$،کیوں جے $!A_g$،~$g$ دی نمائندگی کردا اے۔ایس توں تضاد نکلدا اے۔

ہن فرض کرو $b < \num{m}$۔فیر،کیوں جے \olref{rep-less} نال $\Th{Q} \Proves \lforall[w][(w <
  \num{m} \lif \lnot !A_g(w, \num{n}, \Obj 0))]$،سانوں
$\lnot !A_g(b, \num{n}, \Obj 0)$ ملدا اے۔ایہ وی~(a) نال ٹکراندا اے۔
\end{proof}
\endgroup
% /OLP-0296

% OLP-0297 content/incompleteness/representability-in-q/comp-representable.tex
\begingroup
\def\olfilename{comp-representable}
\expandafter\def\csname ol@sectioncs\endcsname{section}

\olfileid{inc}{req}{crq}
% PNB-CMP-060: converse تے stronger-theory مثالاں لئی ضروری موافقت دا فرض صاف۔
\olsection{\texorpdfstring{قابلِ حساب فنکشن $\Th{Q}$ وچ قابلِ نمائندگی نیں}{قابلِ حساب فنکشن Q وچ قابلِ نمائندگی نیں}}

\begin{thm}
ہر قابلِ حساب فنکشن~$\Th{Q}$ وچ قابلِ نمائندگی اے۔
\end{thm}

\begin{proof}
گل نوں ٹھیک طرح متعین کرن لئی چرچ--ٹیورنگ دا دعویٰ ورتیے۔
اسیں کسے فنکشن نوں اُدوں تے صرف اُدوں قابلِ حساب آکھاں گے
جدوں اوہ عمومی بازگشتی ہووے۔ عمومی بازگشتی فنکشناں دی تعریف صفر
فنکشن~$\Zero$،جانشین فنکشن~$\Succ$ تے پروجیکشن
فنکشن~$\Proj{n}{i}$ توں ترکیب، بنیادی بازگشت تے باقاعدہ کمترین قدر دی تلاش نال
کیتی جا سکدی اے۔\olref[pri]{lem:prim-rec} نال،ہر فنکشن~$h$
جیہڑا $f$ تے~$g$ توں بنیادی بازگشت نال تعریف کیتا جا سکدا اے،اوہ
ترکیب تے باقاعدہ کمترین قدر دی تلاش نال $f$، $g$ تے $\Zero$،
$\Succ$، $\Proj{n}{i}$، $\Add$، $\Mult$، $\Char{=}$ توں وی تعریف کیتا جا سکدا اے۔ایس لئی کوئی فنکشن اُدوں تے صرف اُدوں عمومی بازگشتی اے جدوں اوہ $\Zero$،
$\Succ$، $\Proj{n}{i}$، $\Add$، $\Mult$، $\Char{=}$ توں ترکیب
تے باقاعدہ کمترین قدر دی تلاش نال تعریف کیتا جا سکدا اے۔

اسیں ایہ وی وکھا چکے آں کہ متعلقہ بنیادی فنکشن
~$\Th{Q}$ وچ قابلِ نمائندگی نیں
(\cref{inc:req:bre:prop:rep-zero,inc:req:bre:prop:rep-succ,inc:req:bre:prop:rep-proj,inc:req:bre:prop:rep-id,inc:req:bre:prop:rep-add,inc:req:bre:prop:rep-mult})،
تے قابلِ نمائندگی فنکشناں توں
ترکیب یا باقاعدہ کمترین قدر دی تلاش نال تعریف کیتا گیا ہر فنکشن
(\olref[cmp]{prop:rep-composition}،
\olref[min]{prop:rep-minimization}) وی قابلِ نمائندگی اے۔ ایس لئی ہر
عمومی بازگشتی فنکشن~$\Th{Q}$ وچ قابلِ نمائندگی اے۔
\end{proof}

\begin{explain}
اسیں وکھایا اے کہ قابلِ حساب فنکشناں دا سیٹ اوہناں سارے
فنکشناں دا سیٹ اے جیہڑے $\Th{Q}$ وچ قابلِ نمائندگی نیں۔ ایہ ثبوت ہور وی عام
اے۔ نمائندگی دی تعریف توں سوکھا ای دِسدا اے کہ ہر اوہ نظریہ
جیہڑا $\Th{Q}$ نوں ودھاؤندا اے یا جیہدے وچ $\Th{Q}$ دی تعبیر
ہو سکدی اے، قابلِ حساب فنکشناں دی نمائندگی کر سکدا اے۔ اُلٹی
سمت وچ، ایہو جیہے کسے بے تضاد نظریے دے !!{derivation} نظام وچ
جے !!{derivation} ہون دی پرکھ قابلِ حساب اے، تاں ہر قابلِ
نمائندگی فنکشن قابلِ حساب اے۔ بے تضادی دا ایہ فرض رکھ کے،
قابلِ حساب فنکشناں دا سیٹ پیانو حساب یا زرمیلو--فرینکل سیٹ
نظریے وچ قابلِ نمائندگی فنکشناں دا سیٹ وی اے۔ جیویں گوڈل نے
دَسیا، ایہ گل ذرا حیران کرن والی اے۔ اسیں ویکھاں گے کہ کیہ
ثابت ہو سکدا اے، ایہ چُنے نظریے اُتے بہت انحصار کردا اے۔
عام طور اُتے مسلّمے جِنّے طاقتور ہون، اوِنّا ودھ ثابت ہو سکدا
اے۔ پر مسلّماتی نظرییاں دے اک وڈے دائرے وچ قابلِ نمائندگی
فنکشن عین قابلِ حساب فنکشن نیں۔ جے اوہ نظریے بے تضاد تے مسلّماتی
طور متعین ہون، تاں ہور طاقتور نظریے ہور فنکشناں دی نمائندگی
نہیں کردے۔
\end{explain}
\endgroup
% /OLP-0297

% OLP-0298 content/incompleteness/representability-in-q/representing-relations.tex
\begingroup
\def\olfilename{representing-relations}
\expandafter\def\csname ol@sectioncs\endcsname{section}

\olfileid{inc}{req}{rel}

\olsection{تعلقاں دی نمائندگی}

دسیے کہ کسے \emph{تعلق} دا قابلِ نمائندگی ہون دا کیہ مطلب اے۔

\begin{defn}
\ollabel{defn:representing-relations} قدرتی عدداں اُتے تعلق $R(x_0,\dots,x_k)$ نوں
{\em $\Th{Q}$ وچ قابلِ نمائندگی} آکھدے آں جے کوئی فارمولا
$!A_R(x_0,\dots,x_k)$ ایہو جیہا ہووے۔ جدوں وی $R(n_0,\dots,n_k)$
سچ ہووے، $\Th{Q}$، $!A_R(\num{n_0},\dots,\num{n_k})$ ثابت
کرے، تے جدوں وی $R(n_0,\dots,n_k)$ جھوٹ ہووے، $\Th{Q}$،
$\lnot !A_R(\num{n_0}, \dots, \num{n_k})$ ثابت کرے۔
\end{defn}

\begin{thm}
\ollabel{thm:representing-rels} کوئی تعلق $\Th{Q}$ وچ اُدوں تے صرف اُدوں قابلِ نمائندگی اے
جدوں اوہ قابلِ حساب اے۔
\end{thm}

\begin{proof}
پہلی سمت لئی فرض کرو $R(x_0,\dots,x_k)$ دی
نمائندگی فارمولا $!A_R(x_0,\dots,x_k)$ کردا اے۔$R$ دا حساب کرن لئی دلیلاں $n_0$، \dots،~$n_k$ اُتے ایویں کرو۔
$!A_R(\num{n_0}, \dots, \num{n_k})$ دے ثبوت تے
$\lnot !A_R(\num{n_0}, \dots, \num{n_k})$ دے ثبوت لئی دو تلاشاں اکٹھیاں چلا کے باری باری اگے ودھاؤ۔
ساڈے مفروضے نال اوہناں وچوں اک ثبوت ضرور لبھے گا۔ جے پہلا
لبھے تاں «ہاں» دسو، نہیں تاں «نہیں» دسو۔

دوجی سمت لئی فرض کرو $R(x_0, \dots, x_k)$ قابلِ حساب اے۔تعریف نال،
ایس دا مطلب اے کہ امتیازی فنکشن $\Char{R}(x_0, \dots, x_k)$
قابلِ حساب اے۔\olref[int]{thm:representable-iff-comp} نال،$\Char{R}$
دی نمائندگی کوئی فارمولا کردا اے،مثلاً $!A_{\Char{R}}(x_0, \dots, x_k,
y)$۔مان لو $!A_R(x_0, \dots, x_k)$ اوہ فارمولا اے جیہڑا $!A_{\Char{R}}(x_0,
\dots, x_k, \num{1})$ اے۔ ہن کسے وی $n_0$، \dots،~$n_k$ لئی،جے $R(n_0,
\dots, n_k)$ سچ اے،تاں $\Char{R}(n_0, \dots, n_k) = 1$،تے ایس
صورت وچ $\Th{Q}$، $!A_{\Char{R}}(\num{n_0}, \dots, \num{n_k},
\num{1})$ ثابت کردا اے،تے ایس لئی $\Th{Q}$، $!A_R(\num{n_0}, \dots,
\num{n_k})$ ثابت کردا اے۔ جے $R(n_0, \dots, n_k)$ جھوٹ اے،تاں
$\Char{R}(n_0, \dots, n_k) = 0$۔ایس دا مطلب اے کہ $\Th{Q}$،
\[
\lforall[y][(!A_{\Char{R}}(\num{n_0}, \dots, \num{n_k}, y) \lif y =
  \num{0})].
\]
ثابت کردا اے۔کیوں جے $\Th{Q}$، $\eq/[\num{0}][\num{1}]$ ثابت کردا اے،ایس لئی $\Th{Q}$،
$\lnot !A_{\Char{R}}(\num{n_0}, \dots, \num{n_k}, \num{1})$ ثابت کردا اے،تے ایس لئی
اوہ $\lnot !A_R(\num{n_0}, \dots, \num{n_k})$ ثابت کردا اے۔
\end{proof}

\begin{prob}
وکھاؤ کہ جے $R$،~$\Th{Q}$ وچ قابلِ نمائندگی اے، تاں~$\Char{R}$ وی قابلِ نمائندگی اے۔
\end{prob}
\endgroup
% /OLP-0298

% OLP-0648 content/incompleteness/representability-in-q/c.tex
\begingroup
\def\olfilename{c}
\expandafter\def\csname ol@sectioncs\endcsname{section}

\olfileid{inc}{req}{cee}
\olsection{\texorpdfstring{فنکشاں دا سیٹ $C$}{فنکشاں دا سیٹ C}}

منّ لو $C$ فنکشاں دا سبھ توں نِکّا سیٹ اے جیہدے وچ ایہ شامل نیں:
\begin{enumerate}
\item $0$،
\item جانشین فنکشن،
\item جمع،
\item ضرب،
\item پروجیکشن فنکشن، تے
\item برابری دا امتیازی فنکشن، $\Char{=}$۔
\end{enumerate}
نالے ایہ سیٹ ایہناں عملّاں ہیٹھ بند اے:
\begin{enumerate}
\item فنکشاں دی ترکیب، تے
\item باقاعدہ فنکشاں اُتے غیر محدود تلاش۔
\end{enumerate}
چیتے رکھو، آخری پابندی دا مطلب بس ایہ اے کہ تسی $\mu$ عمل اوہدوں ای ورت سکدے او جدوں نتیجہ کُل فنکشن ہووے۔ ایہدی \emph{عام بازگشتی} فنکشاں دی تعریف نال ٹاکرا کرو۔ ایتھے جمع، ضرب تے $\Char{=}$ شامل کیتے نیں، تے بنیادی بازگشت چھڈ دِتّی اے۔

صاف اے کہ $C$ وچلا ہر فنکشن بازگشتی اے، کیوں جے جمع، ضرب تے $\Char{=}$ بازگشتی نیں۔ اسیں وکھاواں گے کہ اُلٹی گل وی سچ اے۔ ایہدا مطلب اے کہ $C$ وچلیاں باقی شےواں نال اسیں بنیادی بازگشت کر سکدے آں۔
\endgroup
% /OLP-0648
% OLP-0647 content/incompleteness/representability-in-q/c-representable.tex
\begingroup
\def\olfilename{c-representable}
\expandafter\def\csname ol@sectioncs\endcsname{section}

\olfileid{inc}{req}{cre}
\olsection{\texorpdfstring{$C$ وچلے فنکشن $\Th{Q}$ وچ قابلِ نمائندگی نیں}{C وچلے فنکشن Q وچ قابلِ نمائندگی نیں}}

سانوں وکھاؤنا اے کہ $C$ وچلا ہر فنکشن $\Th{Q}$ وچ قابلِ نمائندگی اے۔ آخر سانوں وکھاؤنا پئے گا کہ ہر $k$-دلیلی فنکشن $f(x_0,\dots,x_{k-1})$، جیہڑا $C$ وچ اے، نال فارمولا $!A_f(x_0,\dots,x_{k-1},y)$ کِویں جوڑیے جیہڑا اوہدی نمائندگی کرے۔

\begin{lem}
قدرتی عدد $n$ تے $m$ دِتّے ہون، تاں جے $n \neq m$، فیر $Q \Proves \num
n \neq \num m$۔
\end{lem}

\begin{proof}
$n$ اُتے استقرا نال وکھاؤ کہ ہر $m$ لئی، جے $n \neq m$، تاں
$Q \Proves \eq/[\num n][\num m]$۔

مُڈھلی صورت وچ $n = 0$۔ جے $m$، $0$ دے برابر نہیں، تاں $m = k +
1$ اے، کسے قدرتی عدد $k$ لئی۔ ساڈے کول مسلّمہ $\lforall[x][\eq/[0][x']]$ اے۔ مسوّر دے اک مسلّمے نال، $x$ دی تھاں $\num k$ رکھ کے، اسیں $\eq/[0][\num k']$ کڈھ سکدے آں۔ پر $\num k'$ تاں $\num m$ ای اے۔

استقرا دے قدم وچ اسیں منّ سکدے آں کہ دعویٰ $n$ لئی سچ اے، تے $n+1$ نوں ویکھدے آں۔ کوئی وی قدرتی عدد $m$ لو۔ دو صورتاں نیں۔ یا $m = 0$، یا کسے $k$ لئی $m = k+1$۔ پہلی صورت اُتلے طریقے نال ہو جاندی اے۔ دوجی صورت وچ منّ لو $n+1 \neq
k+1$۔ تاں $n \neq k$۔ $n$ لئی استقرا دے مفروضے نال
$\Th{Q} \Proves \eq/[\num n][\num k]$۔ ساڈے کول مسلّمہ اے
$\lforall[x][\lforall[y][\eq[x'][y'] \lif \eq[x][y]]]$۔ مسوّر دے اک مسلّمے نال $\eq[\num n'][\num k'] \lif \eq[\num n][\num
  k]$ ملدا اے۔ قضیائی منطق نال $\Th{Q}$ وچ
$\eq/[\num n][\num k] \lif \eq/[\num n'][\num k']$ کڈھ سکدے آں۔ اثباتِ مقدم نال $\eq/[\num n'][\num k']$ کڈھ سکدے آں۔ سانوں ایہو چاہیدا سی، کیوں جے $\num k'$ دراصل $\num m$ اے۔
\end{proof}

دھیان دیو، معاون قضیہ کوئی وڈی گل نہیں کردا۔ ایہ بنیادی طور اُتے کہندا اے کہ $\Th{Q}$ ثابت کر سکدا اے کہ وکھ وکھ عددی عبارتاں وکھ وکھ شےواں نوں دسدیاں نیں۔ مثلاً $\Th{Q}$، $0'' \neq 0'''$ ثابت کردا اے۔ پر ایہ عام طور اُتے وکھاؤن لئی کجھ دھیان چاہیدا اے۔ ایہ وی دھیان دیو کہ بھاویں اسیں استقرا ورتدے آں، ایہ استقرا \emph{باہر} اے، $\Th{Q}$ توں۔

اسیں صفر، جانشین، جمع، ضرب، برابری دے امتیازی فنکشن تے پروجیکشن فنکشاں دی نمائندگی کر سکاں گے۔ ہر صورت وچ نمائندگی کرن والا ڈھکدا فنکشن بالکل سدھا اے۔ مثلاً صفر دی نمائندگی ایہ فارمولا کردا اے:
\[
y = 0,
\]
جانشین دی نمائندگی ایہ فارمولا کردا اے:
\[
x_0' = y,
\]
تے جمع دی نمائندگی ایہ فارمولا کردا اے:
\[
x_0 + x_1 = y.
\]
اصل کم ایہ وکھاؤنا اے کہ $\Th{Q}$ متعلقہ بیان ثابت کر سکدا اے۔ مثلاً ایہ آکھن لئی کہ اُتلا فارمولا جمع دی نمائندگی کردا اے، سانوں وکھاؤنا پئے گا کہ قدرتی عدداں $m$ تے $n$ دے ہر جوڑے لئی $\Th{Q}$ ایہ ثابت کردا اے:
\[
\eq[\num n + \num m][\num {n+m}]
\]
تے
\[
\lforall[y][(\eq[(\num n + \num m)][y] \lif \eq[y][\num {n+m}])].
\]

ترکیب بارے کیہ؟ منّ لو $h$ دی تعریف ایویں اے:
\[
h(x_0,\dots,x_{l-1}) = f(g_0(x_0,\dots,x_{l-1}), \dots,
g_{k-1}(x_0,\dots,x_{l-1})).
\]
ایتھے اسیں پہلاں ای فارمولے $!A_f, !A_{g_0}, \dots,
!A_{g_{k-1}}$ لبھ چکے آں جیہڑے بالترتیب فنکشاں $f,g_0,\dots,g_{k-1}$ دی نمائندگی کردے نیں۔ فیر اسیں فارمولا $!A_h$، جیہڑا $h$ دی نمائندگی کردا اے، ایویں تعریف کر سکدے آں کہ $!A_h(x_0,\dots,x_{l-1},y)$ ایہ ہووے:
\begin{multline*}
  \lexists[z_0\dots][\lexists[z_{k-1}][(!A_{g_0}(x_0,\dots,x_{l-1},z_0) \land
  \dots \land !A_{g_{k-1}}(x_0,\dots,x_{l-1},z_{k-1}) \land]]\\
  \formula A_f(z_0,\dots,z_{k-1},y)).
\end{multline*}

آخر غیر محدود تلاش ویکھو۔ منّ لو $g(x,\vec z)$ باقاعدہ اے تے $\Th{Q}$ وچ قابلِ نمائندگی اے، مثلاً فارمولا $!A_g(x,\vec
z,y)$ اوہدی نمائندگی کردا اے۔ $f$ دی تعریف $f(\vec z) = \mu x \; g(x,\vec z)$ نال کرو۔ اسیں فارمولا $!A_f(\vec z,y)$ لبھنا چاہندے آں جیہڑا $f$ دی نمائندگی کرے۔ اک قدرتی چُون ایہ اے:
\[
!A_f(\vec z,y) \ident !A_g(y,\vec z,0) \land \lforall[w][(w < y
\lif \lnot !A_g(w,\vec z,0))].
\]

کجھ ہور معاون قضیاں نال وکھایا جا سکدا اے کہ ایہ کم کردا اے، مثلاً ایہناں نال۔

\begin{lem}
  ہر متغیر $x$ تے ہر قدرتی عدد $n$ لئی، $\Th{Q}$،
  $\eq[(x'
  + \num n)][(x + \num n)']$ ثابت کردا اے۔
\end{lem}

فیر دَس دیئے کہ ایہ ایس گل توں کمزور اے کہ $\Th{Q}$،
$\lforall[x][\lforall[y][(x' + y) = (x +y)']]$ ثابت کردا اے (جیہڑی غلط اے)۔

\begin{proof}
ثبوت، جیویں عام ورتدے آں، $n$ اُتے استقرا نال اے۔ مُڈھلی صورت وچ $n =
0$، سانوں وکھاؤنا اے کہ $\Th{Q}$، $\eq[(x'+0)][(x+0)']$ ثابت کردا اے۔ پر ساڈے کول ایہ اے:
\begin{eqnarray*}
& & \eq[(x' + 0)][x'] \quad \text{مسلّمہ 4 توں} \\
& & \eq[(x + 0)][x] \quad \text{مسلّمہ 4 توں} \\
& & \eq[(x+0)'][x'] \quad \text{سطر 2 نال} \\
& & \eq[(x' + 0)][(x + 0)'] \quad \text{سطراں 1 تے 3}
\end{eqnarray*}
استقرا دے قدم وچ منّ سکدے آں کہ اسیں $\eq[(x' +
  \num n)][(x + \num n)']$، $\Th{Q}$ وچ اخذ کر چکے آں۔ چونکہ $\num{n+1}$،
$\num
n'$ اے، سانوں وکھاؤنا اے کہ $\Th{Q}$، $\eq[(x' + \num n')][(x + \num
n')']$ ثابت کردا اے۔ ایہ برابریاں دی لڑی $\Th{Q}$ وچ اخذ ہو سکدی اے:
\begin{eqnarray*}
(x' + \num n') & \eq & (x' + \num n)' \quad \text{مسلّمہ 5} \\
& \eq & (x + \num n')' \quad \text{استقرا دے مفروضے توں}
\end{eqnarray*}
\end{proof}

\begin{lem}
\begin{enumerate}
\item $\Th{Q}$، $\lnot (x < \num 0)$ ثابت کردا اے۔
\item ہر قدرتی عدد $n$ لئی، $\Th{Q}$ ایہ ثابت کردا اے:
\[
x < \num {n+1} \lif (\eq[x][\num 0] \lor \dots \lor \eq[x][\num n]).
\]
\end{enumerate}
\end{lem}

\begin{proof}
آؤ 1 تے 2 دا کجھ حصہ غیر رسمی طریقے نال کریئے (یعنی بس ایہ اشارے دیئے کہ رسمی !!{derivation} کِویں بنایا جا سکدا اے)۔

پہلے حصے لئی، $<$ دی تعریف مطابق، سانوں $\lnot
\lexists[y][\eq[(y'+x)][0]]$، $\Th{Q}$ وچ ثابت کرنا اے۔ پہلے درجے دی منطق دے مسلّمے تے قاعدے ورتیاں ایہ $\lforall[y][\eq/[(y' +
    x)][0]]$ دے برابر اے۔ خیال ایہ اے۔ منّ لو $\eq[(y'+x)][0]$۔ جے $x$ صفر اے، تاں $\eq[(y' + 0)][0]$۔ پر $\Th{Q}$ دے مسلّمہ 4 نال
$\eq[(y' + 0)][y']$ تے مسلّمہ 2 نال $\eq/[y'][0]$، سو تضاد اے۔ ایس لئی
$\lforall[y][\eq/[(y' +x)][0]]$۔ جے $x$، $0$ نہیں، تاں مسلّمہ 3 نال کوئی $z$ اے جیہدے لئی $\eq[x][z']$۔ پر فیر
$\eq[(y' + z')][0]$۔ مسلّمہ~5 نال $\eq[(y' + z)'][0]$، سو فیر مسلّمہ~2 نال تضاد اے۔

دوجے حصے لئی $n$ اُتے استقرا ورتو۔ مُڈھلی صورت ویکھو، جدوں $n =0$۔ اوس صورت وچ سانوں $x < \num 1 \lif \eq[x][\num
  0]$ وکھاؤنا اے۔ منّ لو $x < \num 1$۔ تاں $<$ دی تعریف والے مسلّمے نال $\lexists[y][\eq[(y'+x)][0']]$۔ منّ لو $y$ وچ ایہ خاصیت اے، سو $\eq[y'+x][0']$۔

سانوں $\eq[x][0]$ وکھاؤنا اے۔ مسلّمہ 3 نال جے $x$ صفر نہیں تاں $z'$ دے برابر اے، کسے $z$ لئی۔ فیر $\eq[(y' + z')][0']$۔ مسلّمہ~5، جیہڑا
$\Th{Q}$ دا اے، نال $\eq[(y'+z)'][0']$۔ مسلّمہ 1 نال $\eq[(y' +
    z)][0]$۔ پر تعریف مطابق ایہدا مطلب $z < 0$ اے، جیہڑا حصّہ~1 نال ٹکراؤندا اے۔
\end{proof}

\begin{explain}
اسیں وکھایا اے کہ قابلِ حساب فنکشاں دا سیٹ $\Th{Q}$ وچ قابلِ نمائندگی فنکشاں دا سیٹ اے۔ اصل وچ ثبوت ہور عام اے۔ نمائندگی دی تعریف توں سوکھا دِسدا اے کہ ہر نظریہ جیہڑا $\Th{Q}$ دا وادھا اے (یا جیہدے وچ $\Th{Q}$ دی تعبیر کیتی جا سکدی اے)، قابلِ حساب فنکشاں دی نمائندگی کر سکدا اے۔ اُلٹے پاسے، ہر !!{derivation} نظام وچ جیہدے وچ ثبوت دا تصور قابلِ حساب اے، ہر قابلِ نمائندگی فنکشن قابلِ حساب اے۔ سو، مثلاً، قابلِ حساب فنکشاں دا سیٹ پیانو حساب، یا زرمَیلو–فرینکل سیٹ نظریے، وچ قابلِ نمائندگی فنکشاں دا سیٹ وی اے۔ جیویں گوڈل نے دَسیا، ایہ کجھ حیرانی والی گل اے۔ اسیں ویکھاں گے کہ ثابت پذیری دے سوال ایس گل اُتے بہت نربھر کردے نیں کہ تسی کیہڑا نظریہ ویکھ رہے او۔ موٹے طور اُتے، مسلّمے جِنّے تکڑے ہون، اُنا ودھ ثابت کر سکدے او۔ پر مسلّماتی نظرییاں دی اک وڈی حد وچ قابلِ نمائندگی فنکشن عین قابلِ حساب فنکشن ای نیں۔
\end{explain}
\endgroup
% /OLP-0647
% OLP-0299 content/incompleteness/representability-in-q/undecidability.tex
\begingroup
\def\olfilename{undecidability}
\expandafter\def\csname ol@sectioncs\endcsname{section}

\olfileid{inc}{req}{und}
\olsection{ناقابلِ فیصلہ ہونا}

اسیں نظریے $\Th{T}$ نوں \emph{ناقابلِ فیصلہ} آکھدے آں جے کوئی
ایہو جیہا حسابی طریقۂ کار موجود نہ ہووے جیہڑا ہر واری محدود گنتی
دے قدماں اندر ایس سوال دا ٹھیک جواب دے سکے: “کی $\Th{T}$،
$!A$ ثابت کردا اے؟” ایہ شرط ہر اُس جملے~$!A$ لئی اے جیہڑا
نظریے~$\Th{T}$ دی زبان وچ ہووے۔ ایس لئی $\Th{Q}$ اُدوں تے صرف
اُدوں قابلِ فیصلہ ہووے گا جدوں کوئی حسابی طریقۂ کار، حساب دی
زبان دا جملہ~$!A$ دتے جان اُتے، فیصلہ کر سکے کہ
$\Th{Q} \Proves !A$ اے یا نہیں۔ ایہ گل ہور ٹھیک طرح ایویں پچھ
سکدے آں: کی تعلق~$\Prov[\Th{Q}](y)$ بازگشتی اے؟ ایہ تعلق~$y$
لئی اُدوں تے صرف اُدوں قائم اے جدوں $y$ اُس جملے دا گوڈل عدد
ہووے جیہڑا~$\Th{Q}$ وچ ثابت ہو سکدا اے۔ جواب اے، نہیں۔

\begin{thm}
$\Th{Q}$ ناقابلِ فیصلہ اے۔ یعنی ایہ تعلق
\[
\Prov[\Th{Q}](y) \defiff \fn{Sent}(y) \land
\lexists[x][\Prf[\Th{Q}](x, y)]
\]
بازگشتی نہیں۔
\end{thm}

\begin{proof}
فرض کرو ایہ تعلق بازگشتی اے۔ فیر اسیں حساب رُکّن دا مسئلہ
ایویں حل کر سکدے آں۔ $e$ تے $n$ دتے جان اُتے، سانوں پتا اے کہ
$\cfind{e}(n) \fdefined$ اُدوں تے صرف اُدوں جدوں کوئی~$s$ ہووے
جیہدے لئی $T(e, n, s)$ قائم ہووے۔ ایتھے $T$، کلینی دا اوہ
محمول اے جیہڑا \olref[cmp][rec][nft]{thm:kleene-nf} وچ اے۔
کیوں جے $T$ بنیادی بازگشتی اے، ایس لئی~$\Th{Q}$ وچ اک
فارمولے $!B_T$ نال ایہدی نمائندگی ہو سکدی اے۔ یعنی
$\Th{Q} \Proves !B_T(\num{e}, \num{n}, \num{s})$ اُدوں تے
صرف اُدوں جدوں $T(e, n, s)$ قائم ہووے۔ جے
$\Th{Q} \Proves !B_T(\num{e}, \num{n}, \num{s})$، تاں
$ \Th{Q} \Proves \lexists[y][!B_T(\num{e}, \num{n}, y)]$ وی۔
جے کوئی ایہو جیہا $s$ موجود نہ ہووے، تاں
$\Th{Q} \Proves \lnot !B_T(\num{e}, \num{n}, \num{s})$
ہر~$s$ لئی۔ پر $\Th{Q}$، $\omega$-بے تضاد اے، یعنی جے
$\Th{Q} \Proves \lnot !A(\num{n})$ ہر~$n \in \Nat$ لئی،
تاں $\Th{Q} \Proves/ \lexists[y][!A(y)]$۔ سانوں ایہ پتا اے
کیوں جے $\Th{Q}$ دے مسلّمے معیاری نمونے~$\Struct{N}$ وچ
سچے نیں۔ ایس لئی $\Th{Q}
\Proves/ \lexists[y][!B_T(\num{e}, \num{n}, y)]$۔
دوجے لفظاں وچ،
$\Th{Q} \Proves \lexists[y][!B_T(\num{e}, \num{n}, y)]$
اُدوں تے صرف اُدوں جدوں کوئی $s$ ہووے جیہدے لئی $T(e, n, s)$
قائم ہووے، یعنی اُدوں تے صرف اُدوں جدوں $\cfind{e}(n) \fdefined$۔
$e$ تے~$n$ توں اسیں $\Gn{\lexists[y][!B_T(\num{e},
    \num{n}, y)]}$ کڈھ سکدے آں۔ ایہ عدد کڈھن والے بنیادی
بازگشتی فنکشن نوں $g(e, n)$ آکھ لئیے۔ فیر
\[
h(e, n) =
\begin{cases}
1 & \text{جے $\Prov[\Th{Q}](g(e, n))$}\\
0 & \text{نہیں تاں}.
\end{cases}
\]
ایس طرح $h$ وی بازگشتی ہوندا جے $\Prov[\Th{Q}]$ بازگشتی
ہوندا۔ پر~$h$ بازگشتی نہیں، جیویں
\olref[cmp][rec][hlt]{thm:halting-problem} توں پتا لگدا اے۔
ایس لئی $\Prov[\Th{Q}]$ وی بازگشتی نہیں ہو سکدا۔
\end{proof}

\begin{cor}
پہلے درجے دی منطق ناقابلِ فیصلہ اے۔
\end{cor}

\begin{proof}
جے پہلے درجے دی منطق قابلِ فیصلہ ہوندی، تاں~$\Th{Q}$ وچ
ثبوت پذیری دا فیصلہ وی کیتا جا سکدا سی۔ ایہ ایس لئی کہ
$\Th{Q} \Proves !A$ اُدوں تے صرف اُدوں جدوں
$\Proves !T \lif !A$، جتھے $!T$، $\Th{Q}$ دے مسلّمیاں دا
اقتران اے۔
\end{proof}
\endgroup
% /OLP-0299

% OLP-0300 content/incompleteness/representability-in-q/sigma1-completeness.tex
\begingroup
\def\olfilename{sigma1-completeness}
\expandafter\def\csname ol@sectioncs\endcsname{section}

\olfileid{inc}{inp}{s1c}
\olsection{\texorpdfstring{$\Sigma_1$}{Sigma-1} دی تمامیت}

بھاویں $\Th{Q}$ تے ایس دیاں بے تضاد تے مسلّمہ بند ہو سکّن والیاں
توسیعاں نامکمل نیں، اسیں ویکھ چکے آں کہ $\Th{Q}$ عددی ٹرماں بارے
کئی بنیادی گلاں ثابت کردا اے۔ ایہ نتیجہ ہور وی اگے ودھایا جا
سکدا اے۔ ایہ سمجھّن لئی کہ~$\Th{Q}$ وچ کجھ ثابت ہو سکّن دی حد
کِتھے تک اے، اسیں $\Delta_0$، $\Sigma_1$ تے $\Pi_1$ !!{formula}s
دے تصور متعارف کراندے آں۔ موٹے طور اُتے، $\Sigma_1$ !!{formula}
اوہ اے جیہدی صورت $\lexists[x][!B(x)]$ ہووے، جتھے $!B$ صرف
قضیاتی رابطیاں تے محدود مسوّراں نال بنیا ہووے۔ اسیں وکھاواں گے
کہ جے $!A$ اک $\Sigma_1$ !!{sentence} اے جیہڑا $\Struct{N}$
وچ سچ اے، تاں $\Th{Q} \Proves !A$ (\olref{thm:sigma1-completeness})۔

\begin{defn}
\ollabel{defn:bd-quant}
% PNB-CMP-061: حد والے ٹرم وچ ایس مسوّر دا متغیر نہیں؛ہور آزاد متغیر آ سکدے نیں۔
\emph{محدود وجودی !!{formula}} دی صورت
$\lexists[x][(x < t \land !A(x))]$ ہوندی اے، جتھے $t$ کوئی وی
ایہو جیہا ٹرم اے جیہدے وچ ایس مسوّر دا متغیر نہ آوے۔ رواج
مطابق اسیں ایہنوں $\bexists{x < t}{!A(x)}$ لکھدے آں۔
%
\emph{محدود کلی !!{formula}} دی صورت
$\lforall[x][(x < t \lif !A(x))]$ ہوندی اے، جتھے $t$ کوئی وی
ایہو جیہا ٹرم اے جیہدے وچ ایس مسوّر دا متغیر نہ آوے۔ رواج
مطابق اسیں ایہنوں $\bforall{x < t}{!A(x)}$ لکھدے آں۔
\end{defn}

\begin{defn}
\ollabel{defn:delta0-sigma1-pi1-frm}
!!{formula} $!B$، $\Delta_0$ اے جے اوہ جوہری !!{formula}s توں
صرف قضیاتی رابطے تے محدود مسوّر ورت کے بنیا ہووے۔
%
!!{formula} $!A$، $\Sigma_1$ اے جے $!A \ident \lexists[x][!B(x)]$،
جتھے $!B$، $\Delta_0$ اے۔
%
!!{formula} $!A$، $\Pi_1$ اے جے $!A \ident \lforall[x][!B(x)]$،
جتھے $!B$، $\Delta_0$ اے۔
\end{defn}

\begin{lem}
\ollabel{lem:q-proves-clterm-id} فرض کرو $t$ اک بند ٹرم اے
جیہدے لئی $\Value{t}{N} = n$۔ تاں $\Th{Q} \Proves \eq[t][\num n]$۔
\end{lem}

\begin{proof}
اسیں ایہ گل~$t$ دی پیچیدگی اُتے استقرا نال ثابت کردے آں۔
بنیادی صورت وچ $\Value{\Obj 0}{N} = 0$ تے
$\Th{Q} \Proves \eq[\Obj 0][\num 0]$، کیوں جے $\num 0 \ident \Obj 0$۔
%
استقرائی قدم لئی، $t_1$ تے $t_2$ ایہو جیہے ٹرم منّو جنہاں
لئی $\Value{t_1}{N} = n_1$، $\Value{t_2}{N} = n_2$،
$\Th{Q} \Proves \eq[t_1][\num n_1]$ تے
$\Th{Q} \Proves \eq[t_2][\num n_2]$۔

فیر $\Value{(t_1')}{N} = n_1 + 1$ اے۔ سانوں $\Th{Q} \Proves
\eq[t_1'][{\num n_1}']$ ملدا اے جدوں اسیں پہلے درجے دی منطق
وچ برابری دے قاعدے استقرائی مفروضے تے !!{formula}
$\eq[\num{n_1}'][\num{n_1}']$ اُتے لاندے آں۔ ایس لئی عددی
ٹرماں دی تعریف نال $\Th{Q} \Proves \eq[t_1'][\num{n_1 + 1}]$ ملدا اے۔

جمع دی صورت وچ
\[
      \Value{(t_1 + t_2)}{N}
    = \Value{t_1}{N} + \Value{t_2}{N}
    = n_1 + n_2.
\]
استقرائی مفروضے تے برابری دے قاعدیاں نال
$\Th{Q} \Proves \eq[t_1 + t_2][\num{n_1} + t_2]$ ملدا اے۔
برابری دے قاعدے دوجی واری لائیے تاں
$\Th{Q} \Proves \eq[t_1 + t_2][\num{n_1} + \num{n_2}]$ ملدا اے۔
\olref[inc][req][bre]{lem:q-proves-add} نال
$\Th{Q} \Proves \eq[\num{n_1} + \num{n_2}][\num{n_1 + n_2}]$
ملدا اے، ایس لئی $\Th{Q} \Proves \eq[t_1 + t_2][\num{n_1 + n_2}]$۔

~$\times$ دی صورت وچ وی ایہو جیہی دلیل ورتی جا سکدی اے، جتھے
اسیں \olref[inc][req][bre]{lem:q-proves-mult} ورتدے آں۔
%
ایس طرح حساب دے سارے بند ٹرم شامل ہو گئے۔ ایس لئی
$\Th{Q} \Proves \eq[t][\num n]$ ہر اُس بند ٹرم~$t$ لئی
جیہدے واسطے $\Value{t}{N} = n$۔
\end{proof}

\begin{prob}
\olref{lem:q-proves-clterm-id} دے جس حصے وچ~$\times$ آؤندا اے،
اوہدا تفصیلی ثبوت دیو۔
\end{prob}

\begin{lem}
\ollabel{lem:atomic-completeness}
فرض کرو $t_1$ تے $t_2$ بند ٹرم نیں۔تاں
\begin{enumerate}
\item جے $\Value{t_1}{N} = \Value{t_2}{N}$،
    تاں $\Th{Q} \Proves \eq[t_1][t_2]$۔
\item جے $\Value{t_1}{N} \neq \Value{t_2}{N}$،
    تاں $\Th{Q} \Proves \eq/[t_1][t_2]$۔
\item جے $\Value{t_1}{N} < \Value{t_2}{N}$،
    تاں $\Th{Q} \Proves t_1 < t_2$۔
\item جے $\Value{t_2}{N} \leq \Value{t_1}{N}$،
    تاں $\Th{Q} \Proves \lnot(t_1 < t_2)$۔
\end{enumerate}
\end{lem}

\begin{proof}
دتے ہوئے ٹرماں $t_1$ تے $t_2$ لئی اسیں $n = \Value{t_1}{N}$
تے $m = \Value{t_2}{N}$ مقرر کردے آں۔

% PNB-CMP-062: دوجے ٹرم دی معیاری قدر m اے؛بلاشرط n نہیں۔
فرض کرو $!A \ident t_1 = t_2$۔ \olref{lem:q-proves-clterm-id}
نال $\Th{Q} \Proves \eq[t_1][\num n]$ تے $\Th{Q} \Proves \eq[t_2][\num m]$۔
جے $n = m$، تاں $\Th{Q} \Proves \eq[\num n][\num m]$، تے برابری
دی تعدّیت نال $\Th{Q} \Proves \eq[t_1][t_2]$ ملدا اے۔
جے $n \neq m$، تاں $\Th{Q} \Proves \eq/[\num n][\num m]$، تے
فیر برابری دی تعدّیت ورت کے $\Th{Q} \Proves \eq/[t_1][t_2]$ ملدا اے۔

ہن $!A \ident t_1 < t_2$ منّو۔ دونیں صورتاں وچ اسیں
مسلّمے~$!Q_8$ نوں ورتاں گے۔ ایہ مسلّمہ آکھدا اے کہ
$x < y \liff \lexists[z][\eq[z' + x][y]]$ سارے $x,y$ لئی۔

% PNB-CMP-063: بند عددی جمع Q8 دے گواہ دے ترتیب وار روپ وچ؛Q دی عام تبادلی خاصیت نہیں۔
فرض کرو $\Sat{N}{t_1 < t_2}$۔ تاں کوئی $k \in \Nat$ اے
جیہدے لئی $n + k + 1 = m$۔ \olref{lem:q-proves-clterm-id} نال
$\Th{Q} \Proves \eq[t_1][\num n]$ تے $\Th{Q} \Proves \eq[t_2][\num m]$۔
فیر ایس لیما دا پہلا حصہ تے بند عددی جمع دی معیاری قدر ورت کے
$\Th{Q} \Proves \eq[{\num k}' + \num n][\num m]$ ملدا اے۔
برابری دی تعدّیت تے برابر ٹرماں دی بدل نال
$\Th{Q} \Proves \eq[{\num k}' + t_1][t_2]$ ملدا اے۔
ایس لئی $\Th{Q} \Proves \lexists[z][\eq[z' + t_1][t_2]]$۔
~$!Q_8$ نوں سجے پاسے توں کھبے پاسے ول ورت کے
$\Th{Q} \Proves t_1 < t_2$ ملدا اے۔

% PNB-CMP-064: صفر نال مثبت برابری Q2 نال تضاد؛Q3 سابق دا مسلّمہ اے۔
ہن فرض کرو $\Sat/{N}{t_1 < t_2}$، یعنی $m \leq n$۔
%
اسیں~$\Th{Q}$ وچ کم کردے آں تے $t_1 < t_2$ فرض کردے آں۔
~$!Q_8$ نوں کھبے پاسے توں سجے پاسے ول ورتن نال کوئی~$z$
ملدا اے جیہدے لئی $\eq[z' + t_1][t_2]$۔
کیوں جے $\Th{Q} \Proves \eq[t_1][\num n]$ تے
$\Th{Q} \Proves \eq[t_2][\num m]$، ایس لئی $\eq[z' + \num n][\num m]$۔
%
~$m$ اُتے باہرونی استقرا نال، ~$!Q_5$ تے جانشیناں دی برابری
توں اصل ٹرماں دی برابری والا مسلّمہ ورت کے،
$\eq[z' + \num{n - m}][\Obj 0]$ ملدا اے۔
جے $m = n$، تاں جمع وچ صفر دے مسلّمے نال $\eq[z'][\Obj 0]$
ملدا اے، جیہڑا~$!Q_2$ نال تضاد دیندا اے۔
جے $m < n$، تاں $\eq[(z' + \num{n - m - 1})'][\Obj 0]$
فیر~$!Q_5$ نال ملدا اے۔ ایہ وی~$!Q_2$ نال تضاد دیندا اے۔
ایس لئی $\Th{Q} \Proves \lnot(t_1 < t_2)$۔
\end{proof}

\begin{lem}
\ollabel{lem:bounded-quant-equiv}
فرض کرو $!A$ اک !!a{formula} اے، $t$ اک بند ٹرم اے تے
$k=\Value{t}{N}$۔ تاں
\begin{enumerate}
\item $\Th{Q} \Proves \bforall{x<t}{!A(x)}$ اُدوں تے صرف اُدوں جدوں $\Th{Q} \Proves
    !A(\num 0) \land \dots \land !A(\num{k-1})$۔
\item $\Th{Q} \Proves \bexists{x<t}{!A(x)}$ اُدوں تے صرف اُدوں جدوں $\Th{Q} \Proves
    !A(\num 0) \lor \dots \lor !A(\num{k-1})$۔
\end{enumerate}
\end{lem}

\begin{proof}
    اسیں محدود کلی مسوّر دی صورت ثابت کردے آں۔
    جے $\Value{t}{N} = 0$، تاں ہم ارزی دا کھبا پاسا~$\Th{Q}$
    وچ ثابت ہو سکدا اے، کیوں جے \olref[inc][req][min]{lem:less-zero}
    مطابق کوئی $x<\num 0$ نہیں۔
    % PNB-CMP-065: کلی مسوّر دی صفر حد والی صورت وچ سچا خالی اقتران ورتیا اے۔
    ایس طرح اسیں خالی اقتران نوں $\ltrue$ منّ سکدے آں، جیہڑا
    ~$\Th{Q}$ وچ وی ثابت ہو سکدا اے۔
    %
    ہن فرض کرو $\Value{t}{N} = k+1$، جتھے~$k$ کوئی قدرتی عدد اے۔
    \olref{lem:q-proves-clterm-id} نال اسیں منّ سکدے آں کہ
    ساڈے کول $\bforall{x<\num{k+1}}{!A(x)}$ دی صورت دا !!a{formula} اے۔
    
    فرض کرو $\Th{Q} \Proves \bforall{x<\num{k+1}}{!A(x)}$،
    تے $n \leq k$ منّو۔ \olref{lem:atomic-completeness} نال
    $\Th{Q} \Proves \num n < \num{k+1}$، ایس لئی منطق نال
    $\Th{Q} \Proves !A(\num n)$ ملدا اے۔ ایس دلیل نوں ہر عدد
    لئی اک واری، کل $k+1$ وار ورتیئے، یعنی ہر $n \leq k$ لئی،
    تاں مطلوبہ $\Th{Q} \Proves !A(\num 0)
    \land \dots \land !A(\num k)$ ملدا اے۔
    
    دوجے رخ لئی فرض کرو $\Th{Q} \Proves
    !A(\num 0) \land \dots \land !A(\num k)$۔ اسیں
    $\Th{Q}$ وچ کم کردیاں $x < \num{k+1}$ فرض کردے آں۔
    \olref[inc][req][min]{lem:less-nsucc} نال
    $x = \num 0 \lor \dots \lor x = \num k$ ملدا اے۔ ایس توں
    منطق نال~$!A(x)$ تے فیر کلی دعویٰ
    $\bforall{x<\num{k+1}}{!A(x)}$ ملدا اے۔
    
    محدود وجودی مسوّر والے !!{formula}s لئی ہم ارزی دا ثبوت
    وی ایس طرح دا اے۔
\end{proof}

\begin{prob}
\olref{lem:bounded-quant-equiv} دی وجودی صورت دا تفصیلی ثبوت دیو۔
\end{prob}

\begin{lem}
\ollabel{lem:delta0-completeness}
جے $!A$ اک $\Delta_0$ !!{sentence} اے جیہڑا
$\Struct{N}$ وچ سچ اے،تاں $\Th{Q} \Proves !A$۔
\end{lem}

\begin{proof}
اسیں ایہ گل !!{formula} دی پیچیدگی اُتے استقرا نال ثابت کردے آں۔
%
بنیادی صورت \olref{lem:atomic-completeness} توں ملدی اے۔
ہن اسیں استقرائی قدم ول آندے آں۔ سوکھ لئی نفی دی صورت نوں
اوس !!{formula} دی بناوٹ مطابق اگے ونڈدے آں جیہدے اُتے
نفی لائی گئی اے۔

\begin{enumerate}
\item فرض کرو $(!A \land !B)$، $\Struct{N}$ وچ سچ اے۔
ایس لئی $!A$ تے $!B$،~$\Struct{N}$ وچ سچے نیں۔
استقرائی مفروضے نال $\Th{Q} \Proves !A$ تے
$\Th{Q} \Proves !B$، ایس لئی منطق نال $\Th{Q} \Proves (!A \land !B)$۔
%
\item فرض کرو $\lnot (!A \land !B)$، $\Struct{N}$ وچ سچ اے۔
تاں $\lnot !A$ یا $\lnot !B$، $\Struct{N}$ وچ سچ اے۔
عمومیت قائم رکھدیاں پہلی صورت منّ لئیے۔ استقرائی مفروضے نال
$\Th{Q} \Proves \lnot !A$، تے منطق نال
$\Th{Q} \Proves \lnot (!A \land !B)$۔
%
\item فرض کرو $(!A \lor !B)$، $\Struct{N}$ وچ سچ اے۔
تاں یا $!A$، $\Struct{N}$ وچ سچ اے، یا $!B$، $\Struct{N}$
وچ سچ اے۔ عمومیت قائم رکھدیاں پہلی صورت منّ لئیے۔ استقرائی
مفروضے نال $\Th{Q} \Proves !A$، تے منطق نال
$\Th{Q} \Proves (!A \lor !B)$۔
%
\item فرض کرو $\lnot(!A \lor !B)$، $\Struct{N}$ وچ سچ اے۔
ایس لئی $\lnot !A$ تے $\lnot !B$، $\Struct{N}$ وچ سچے نیں۔
استقرائی مفروضے نال $\Th{Q} \Proves \lnot !A$ تے
$\Th{Q} \Proves \lnot !B$، تے منطق نال
$\Th{Q} \Proves \lnot(!A \lor !B)$۔
%
\item فرض کرو $\bforall{x<t}{!A(x)}$،~$\Struct{N}$ وچ سچ اے،
جتھے $t$ اک بند ٹرم اے تے $k=\Value{t}{N}$۔
استقرائی مفروضے تے منطق نال، جے $!A(\num n)$،~$\Struct{N}$ وچ
ہر $n < \Value{t}{N}$ لئی سچ اے، تاں
$\Th{Q} \Proves !A(\num 0) \land \dots \land !A(\num{k-1})$۔
\olref{lem:bounded-quant-equiv} نال $\Th{Q} \Proves \bforall{x<t}{!A(x)}$ ملدا اے۔
%
\item محدود وجودی مسوّر دی صورت وچ ساڈے کول
$\bexists{x < t}{!A(x)}$ دی صورت دا !!a{sentence} اے۔
ایہ صورت محدود کلی مسوّر والی صورت ورگی اے۔
%
\item فرض کرو $\lnot \bforall{x<t}{!A(x)}$،~$\Struct{N}$ وچ
سچ اے، جتھے $t$ اک بند ٹرم اے۔ ایہ !!{sentence} اوس
!!{sentence} دے ہم ارز اے جیہدی صورت $\bexists{x<t}{\lnot !A(x)}$
اے۔ ایہ ہم ارزی~$\Th{Q}$ وچ اخذ کیتی جا سکدی اے۔ ایس لئی
اسیں محدود وجودی مسوّراں والی دلیل ورت سکدے آں۔
%
\item ایس طرح فرض کرو $\lnot \bexists{x<t}!A(x)$، $\Struct{N}$
وچ سچ اے، جتھے $t$ اک بند ٹرم اے۔ ایہ !!{sentence}،
$\Th{Q}$ وچ $\bforall{x<t}{\lnot!A(x)}$ دے ہم ارز اے۔
ایس لئی اسیں محدود کلی مسوّراں والی دلیل ورت سکدے آں۔
%
\item آخر اُتے فرض کرو $\lnot !A$، $\Struct{N}$ وچ سچ اے۔
ہن دو صورتاں رہ جاندیاں نیں۔ پہلی وچ $!A$ جوہری اے تے
دوجی وچ $\lnot !A \ident \lnot\lnot !B$ اے، کسے
$\Delta_0$ !!{sentence} $!B$ لئی۔
جے $!A$ جوہری اے، تاں \olref{lem:atomic-completeness} نال
$\Th{Q} \Proves \lnot !A$۔
جے $\lnot !A \ident \lnot\lnot !B$، تاں منطق نال~$\Th{Q}$
وچ ایہدی~$!B$ نال ہم ارزی ثابت ہو سکدی اے۔ ہم ارز جملہ~$\Struct{N}$
وچ سچ اے، کیوں جے $\lnot !A$،~$\Struct{N}$ وچ سچ اے۔
ایس لئی استقرائی مفروضے نال $\Th{Q} \Proves \lnot !A$۔
\end{enumerate}
\end{proof}

\begin{prob}
\olref{lem:delta0-completeness} دی وجودی صورت دا تفصیلی ثبوت دیو۔
\end{prob}

\begin{thm}
\ollabel{thm:sigma1-completeness}
جے $!A$ اک $\Sigma_1$ !!{sentence} اے جیہڑا
~$\Struct{N}$ وچ سچ اے،تاں $\Th{Q} \Proves !A$۔
\end{thm}

\begin{proof}
جے $\lexists{x}!A(x)$ اک $\Sigma_1$ !!{sentence} اے جیہڑا
~$\Struct{N}$ وچ سچ اے، تاں کوئی قدرتی عدد~$n$ تے متغیر
قدر بندی~$s$ موجود نیں جنہاں لئی $s(x) = n$ تے
$\Sat{N}{!A(x)}[s]$۔ پورا اترن دے تعلق دیاں معیاری خاصیتاں
توں $\Sat{N}{!A(\num n)}$ ملدا اے۔ ہن $!A(\num n)$ اک
$\Delta_0$ !!{formula} اے، ایس لئی \olref{lem:delta0-completeness}
نال $\Th{Q} \Proves !A(\num n)$، تے فیر منطق نال
$\Th{Q} \Proves \lexists[x][!A(x)]$۔
\end{proof}
\endgroup
% /OLP-0300

\OLEndChapterHook
\endgroup
% /OLP-0289

% OLP-0301 content/incompleteness/theories-computability/theories-computability.tex
\begingroup
\def\olfilename{theories-computability}
\expandafter\def\csname ol@sectioncs\endcsname{section}

\begin{editorial}
  ایس باب دا دارومدار قابلِ حسابیت دے نظریے والے باب دے مواد اُتے اے۔
  جے اوہ باب شامل نہ ہووے، تاں ایہ باب چھڈیا جا سکدا اے۔ ہن ایہ
  Jeremy Avigad دے نوٹس نوں ڈھال کے تیار کیتا ابتدائی روپ اے۔
  ایہدی اجے نظرثانی نہیں ہوئی تے ایس وچ مشقاں دی گھاٹ اے۔
\end{editorial}

\olchapter{inc}{tcp}{نظریے تے قابلِ حسابیت}

% OLP-0302 content/incompleteness/theories-computability/introduction.tex
\begingroup
\def\olfilename{introduction}
\expandafter\def\csname ol@sectioncs\endcsname{section}

\olfileid{inc}{tcp}{int}

\olsection{تعارف}

\begin{editorial}
ایس حصے نوں فیر توں لکھن دی لوڑ اے۔
\end{editorial}

ساڈے کول ایہ تعریفاں تے قضیے نیں۔
\begin{enumerate}
\item ایس گل دی تعریف کہ کسے فنکشن دا $\Th{Q}$ وچ قابلِ نمائندگی
  ہونا کیہہ اے (\olref[req][int]{defn:representable-fn})
\item ایس گل دی تعریف کہ کسے تعلق دا $\Th{Q}$ وچ قابلِ نمائندگی
  ہونا کیہہ اے (\olref[req][rel]{defn:representing-relations})
\item اک قضیہ جیہڑا آکھدا اے کہ $\Th{Q}$ وچ قابلِ نمائندگی فنکشن
  عین اوہی نیں جیہڑے قابلِ حساب نیں (\olref[req][int]{thm:representable-iff-comp})
\item اک قضیہ جیہڑا آکھدا اے کہ $\Th{Q}$ وچ قابلِ نمائندگی تعلق
  عین اوہی نیں جیہڑے قابلِ حساب نیں
  % PNB-CMP-066: اصل چوتھے حوالے دی غائب کھلّن والی قوس درست۔
  (\olref[req][rel]{thm:representing-rels})
\end{enumerate}

{\em نظریہ} جملیاں دا اوہ سیٹ اے جیہڑا اخذ دے لحاظ نال بند ہووے۔
یعنی جدوں وی $T$، $!A$ نوں ثابت کردا اے، تاں $!A$، $T$ وچ ہوندا اے۔
نظریے نوں سمجھّن دا سوکھا طریقہ ایہ اے کہ جملیاں دا اک مجموعہ
تے اوہ ساریاں گلاں وی منّ لئیے جیہڑیاں انہاں جملیاں توں لازم آندیاں نیں۔
اگے اسیں $\Th{Q}$ نال اوہ {\em نظریہ} مراد لواں گے جیہڑا
\olref[req][int]{sec} دے اٹھ مسلّمیاں توں اخذ ہو سکّن والے سارے
جملیاں دا سیٹ اے۔ یاد رکھو کہ اسیں $\Th{Q}$ دے فارمولیاں نوں عدداں
نال کوڈ کر سکدے آں۔ جے $!A$ ایہو جیہا فارمولا اے، تاں $\Gn{!A}$
نال اوہ عدد مراد لئیے جیہڑا $!A$ نوں کوڈ کردا اے۔ ایس کوڈ بندی
دے حساب نال ہن اسیں پچھ سکدے آں کہ فارمولیاں دے وکھ وکھ سیٹ
قابلِ حساب نیں یا نہیں۔
\endgroup
% /OLP-0302

% OLP-0303 content/incompleteness/theories-computability/q-is-ce.tex
\begingroup
\def\olfilename{q-is-ce}
\expandafter\def\csname ol@sectioncs\endcsname{section}

\olfileid{inc}{tcp}{qce}

\olsection{\texorpdfstring{$\Th{Q}$ \printtoken{S}{c.e.}-مکمل اے}{Q قابلِ حساب گنتی وار فہرست بنن جوگ-مکمل اے}}

\begin{thm}
$\Th{Q}$ !!{c.e.} اے، پر قابلِ فیصلہ نہیں۔ اصل وچ ایہ اک مکمل
!!{c.e.} سیٹ اے۔
\end{thm}

\begin{proof}
ایہ ویکھنا سوکھا اے کہ $\Th{Q}$ !!{c.e.} اے۔ ایہ اوہناں جملیاں
دے کوڈاں $y$ دا سیٹ اے جنہاں لئی کوئی ثبوت $x$ موجود اے جیہڑا
$y$ نوں $\Th{Q}$ وچ ثابت کردا اے۔
\[
Q = \Setabs{y}{\lexists[x][\Prf[\Th{Q}](x,y)]}.
\]
سانوں پتا اے کہ $\Prf[\Th{Q}](x,y)$ قابلِ حساب اے، بلکہ ایہ
بنیادی بازگشتی اے۔ ہر سیٹ جیہڑا اُپرلی وجودی صورت وچ لکھیا
جا سکے، قابلِ حساب گنتی وار فہرست بنن جوگ اے۔

ایہ آکھنا کہ ایہ اک مکمل قابلِ حساب گنتی وار فہرست بنن جوگ سیٹ
اے، ایہ آکھن دے برابر اے کہ $K
\leq_m Q$، جتھے $K = \Setabs{x}{!A_x(x) \downarrow}$۔
سو اسیں وکھاواں گے کہ $K$ دی $\Th{Q}$ ول کئی توں اک تقلیل
اے۔ کلینی دا محمول $T(e,x,s)$ بنیادی بازگشتی اے، ایس لئی
$\Th{Q}$ وچ اوہدی نمائندگی ہو سکدی اے۔ منّ لئیے کہ اوہدی
نمائندگی $!A_T$ کردا اے۔ فیر ہر $x$ لئی
\begin{align*}
x \in K & \lif \lexists[s][T(x,x,s)] \\
& \lif \lexists[s][(\Th{Q} \Proves !A_T(\num x,\num x, \num s))]
  \\
& \lif \Th{Q} \Proves \lexists[s][!A_T(\num x, \num x, s)].
\end{align*}
دوجے پاسے، جے $\Th{Q} \Proves \lexists[s][!A_T(\num x, \num x, s)]$،
تاں معیاری نمونے وچ کوئی قدرتی عدد $n$ اے جیہدے لئی فارمولا
$!A_T(\num x,
\num x, \num n)$ سچ اے۔ ایہ ایس لئی کہ Q دے اٹھ مسلّمے
معیاری نمونے وچ سچے نیں تے منطقی اخذ سچائی محفوظ رکھدا اے۔
ہن جے $T(x,x,n)$ جھوٹا ہوندا، تاں $\Th{Q}$،
$\lnot !A_T(\num x, \num x, \num n)$ ثابت کردا، کیوں جے
$!A_T$، تعلق $T$ دی نمائندگی کردا اے۔ پر فیر $\Th{Q}$
اک جھوٹا فارمولا ثابت کردا، جیہڑا تضاد اے۔ ایس لئی
$T(x,x,n)$ سچ اے، جیہدے توں $!A_x(x) \downarrow$ نکلدا اے۔

% PNB-CMP-067: آخری دو جملہ/کوڈ خلاصیاں وچ عددی تعلق T دی جگہ اوہدا نمائندہ فارمولا !A_T؛پہلے T استعمال عین۔
ایس طرح ہر $x$ لئی، $x$ سیٹ $K$ وچ اُدوں تے صرف اُدوں اے
جدوں $\Th{Q}$، جملہ $\lexists[s][!A_T(\num x,\num x,s)]$ ثابت
کردا اے۔ سو فنکشن $f$، جیہڑا $x$ نوں جملے
$\lexists[s][!A_T(\num x,
  \num x, s)]$ دے کوڈ اُتے بھیجدا اے، $K$ دی $\Th{Q}$ ول
کئی توں اک تقلیل اے۔
\end{proof}
\endgroup
% /OLP-0303

% OLP-0304 content/incompleteness/theories-computability/oconsis-ext-of-q-undec.tex
\begingroup
\def\olfilename{oconsis-ext-of-q-undec}
\expandafter\def\csname ol@sectioncs\endcsname{section}

\olfileid{inc}{tcp}{oqn}

\olsection{\texorpdfstring{$\omega$-بے تضاد توسیعاں جیہڑیاں $\Th{Q}$ نوں شامل کردیاں نیں، ناقابلِ فیصلہ نیں}{ω-بے تضاد توسیعاں جیہڑیاں Q نوں شامل کردیاں نیں، ناقابلِ فیصلہ نیں}}

\begin{explain}
$\Th{Q}$ دی قابلِ حساب گنتی وار فہرست بنن جوگ سیٹاں وچ مکمل ہون
دا ثبوت ایس گل اُتے ٹکیا سی کہ $\Th{Q}$ وچ ہر ثابت ہو سکّن والا
جملہ قدرتی عدداں بارے “سچ” اے۔ اگلی تعریف تے قضیہ سچائی دے اوہناں خاص
پہلوآں دی نشان دہی کردے نیں جیہڑے پچھلے ثبوت لئی لوڑیںدے سن۔
ایس طرح نتیجہ ہور مضبوط ہو جاندا اے۔ ایس قضیے اُتے بہتی دیر
نہ لاؤ، کیوں جے اسیں چھیتی ایہنوں ہور مضبوط کراں گے۔ ایہنوں
شامل کرن دا خاص سبب تاریخی اے۔ گوڈل نے اپنے اصل مقالے وچ
$\omega$-بے تضادی دا تصور ورتیا سی۔ تھوڑے چِر پچھوں
$\omega$-بے تضادی دی جگہ عام بے تضادی رکھ کے اوہدا نتیجہ
ہور مضبوط کر دتا گیا۔
\end{explain}

\begin{defn}
\ollabel{thm:oconsis-q}
نظریہ $\Th{T}$، $\omega$-بے تضاد اے جے ہِٹھلی شرط پوری ہووے۔
جے $\lexists[x][!A(x)]$ کوئی وی جملہ اے تے $\Th{T}$،
$\lnot
!A(\num 0)$، $\lnot !A(\num 1)$، $\lnot !A(\num 2)$، \dots
سبھ ثابت کردا اے، تاں $\Th{T}$، جملہ
$\lexists[x][!A(x)]$ ثابت نہیں کردا۔
\end{defn}

\begin{thm}
فرض کرو $\Th{T}$ اک $\omega$-بے تضاد نظریہ اے جیہڑا $\Th{Q}$
نوں شامل کردا اے۔ تاں $\Th{T}$ قابلِ فیصلہ نہیں۔
\end{thm}

\begin{proof}
جے $\Th{T}$ وچ $\Th{Q}$ شامل اے، تاں $\Th{T}$ قابلِ حساب
فنکشناں تے تعلقاں دی نمائندگی کردا اے۔ سانوں پچھلے ثبوت وچ
بس تھوڑی جیہی تبدیلی کرنی پئے گی۔ پہلاں وانگ جے $x \in K$،
تاں $\Th{T}$، جملہ $\lexists[s][!A_T(\num
  x,\num x, s)]$ ثابت کردا اے۔ دوجے پاسے فرض کرو $\Th{T}$
جملہ $\lexists[s][!A_T(\num x, \num x, s)]$ ثابت کردا اے۔
فیر $x$ سیٹ $K$ وچ ضرور اے۔ جے ایہ اوہدے وچ نہ ہووے، تاں
مشین $x$ دی انپٹ $x$ اُتے کوئی رُکّن والی حسابی کارروائی نہیں۔
کیوں جے $!A_T$، کلینی دے تعلق $T$ دی نمائندگی کردا اے،
ایس لئی $\Th{T}$، فارمولے $\lnot !A_T(\num x, \num x, \num 0)$،
$\lnot !A_T(\num x, \num
x, \num 1)$، \dots سبھ ثابت کردا اے۔ ایس طرح $\Th{T}$،
$\omega$-بے تضاد نہیں رہندا۔
\end{proof}
\endgroup
% /OLP-0304

% OLP-0305 content/incompleteness/theories-computability/extensions-of-q-not-decidable.tex
\begingroup
\def\olfilename{extensions-of-q-not-decidable}
\expandafter\def\csname ol@sectioncs\endcsname{section}

\olfileid{inc}{tcp}{cqn}

\olsection{\texorpdfstring{$\Th{Q}$ دیاں بے تضاد توسیعاں ناقابلِ فیصلہ نیں}{Q دیاں بے تضاد توسیعاں ناقابلِ فیصلہ نیں}}

\begin{explain}
یاد رکھو کہ اک نظریہ \emph{بے تضاد} اے جے اوہ
$!A$ تے $\lnot !A$ اکٹھے ثابت نہ کرے، کسے وی فارمولے~$!A$ لئی۔
تضاد توں ہر شے اخذ ہو سکدی اے۔ ایس لئی تضاد والے نظریے وچ ہر
جملہ ثابت ہو سکدا اے۔ ظاہر اے کہ جے نظریہ $\omega$-بے تضاد
ہووے، تاں اوہ بے تضاد وی اے۔ عام بے تضادی دی شرط کمزور اے۔
یعنی کجھ نظریے بے تضاد ہوندے نیں پر $\omega$-بے تضاد نہیں
ہوندے۔ اسیں \olref[oqn]{thm:oconsis-q} وچ مفروضے دی جگہ
عام بے تضادی رکھ کے ہور مضبوط قضیہ حاصل کر سکدے آں۔
\end{explain}

\begin{lem}
کوئی “ہمہ گیر قابلِ حساب تعلق” موجود نہیں۔ یعنی کوئی دو رکنی
قابلِ حساب تعلق $R(x,y)$ موجود نہیں جیہدے وچ ایہ خاصیت ہووے۔
جدوں وی $S(y)$ اک رکنی قابلِ حساب تعلق ہووے، کوئی ایہو جیہا $k$
ہووے جیہدے لئی ہر $y$ اُتے، $S(y)$ اُدوں تے صرف اُدوں سچ
ہووے جدوں $R(k,y)$ سچ ہووے۔
\end{lem}

\begin{proof}
فرض کرو $R(x,y)$ ہمہ گیر قابلِ حساب تعلق اے۔ $S(y)$ نوں
تعلق $\lnot R(y,y)$ منّ لئیے۔ کیوں جے $S(y)$ قابلِ حساب اے،
ایس لئی کسے $k$ لئی $S(y)$ تے $R(k,y)$ ہم ارز نیں۔
فیر $S(k)$، $R(k,k)$ تے $\lnot R(k,k)$ دونہاں دے ہم ارز
اے، جیہڑا تضاد اے۔
\end{proof}

\begin{thm}
فرض کرو $\Th{T}$ کوئی وی بے تضاد نظریہ اے جیہڑا $\Th{Q}$
نوں شامل کردا اے۔ تاں $\Th{T}$ قابلِ فیصلہ نہیں۔
\end{thm}

\begin{proof}
فرض کرو $\Th{T}$، $\Th{Q}$ دی اک بے تضاد تے قابلِ فیصلہ
توسیع اے۔ اسیں $\Th{T}$ دی مدد نال اک ہمہ گیر قابلِ حساب
تعلق متعین کراں گے۔ ایس توں تضاد نکلے گا۔

تعلق $R(x,y)$ اُدوں تے صرف اُدوں قائم اے جدوں
\begin{quote}
$x$، فارمولے $!D(u)$ دا کوڈ اے تے $\Th{T}$،
فارمولا $!D(\num y)$ ثابت کردا اے۔
\end{quote}
فرض مطابق $\Th{T}$ قابلِ فیصلہ اے، ایس لئی $R$ قابلِ حساب اے۔
ہن وکھائیے کہ $R$ ہمہ گیر اے۔ جے $S(y)$ کوئی وی قابلِ حساب
تعلق اے، تاں $\Th{Q}$ وچ، تے ایس لئی $\Th{T}$ وچ وی، اوہدی
نمائندگی فارمولا $!D_S(u)$ کر سکدا اے۔ فیر ہر $n$ لئی
\begin{eqnarray*}
S(\num n) & \lif & T \vdash !D_S(\num n) \\
& \lif & R(\Gn{!D_S(u)}, n)
\end{eqnarray*}
تے
\begin{eqnarray*}
\lnot S(\num n) & \lif & T \vdash \lnot !D_S(\num n) \\
& \lif & T \not\vdash !D_S(\num n) \quad \text{(کیوں جے $\Th{T}$
  بے تضاد اے)} \\
& \lif & \lnot R(\Gn{!D_S(u)}, n).
\end{eqnarray*}
یعنی ہر $y$ لئی $S(y)$ اُدوں تے صرف اُدوں سچ اے جدوں
$R(\Gn{!D_S(u)}, y)$ سچ اے۔ ایس لئی $R$ ہمہ گیر اے،
تے ایہی اوہ تضاد اے جیہڑا اسیں لبھ رہے ساں۔
\end{proof}

“سچا حساب” اوہ نظریہ اے جیہڑا $\Setabs{!A}{ \Sat{\Nat}{!A}}$
نال دتا جاندا اے، یعنی حساب دی زبان دے اوہناں جملیاں دا سیٹ
جیہڑے معیاری تعبیر وچ سچے نیں۔

\begin{cor}
سچا حساب قابلِ فیصلہ نہیں۔
\end{cor}

%حصے~\olref{undefinability:truth:section} وچ اسیں تارسکی دا
%ہور مضبوط نتیجہ بیان کراں گے۔
\endgroup
% /OLP-0305

% OLP-0306 content/incompleteness/theories-computability/computably-axiomatizable.tex
\begingroup
\def\olfilename{computably-axiomatizable}
\expandafter\def\csname ol@sectioncs\endcsname{section}

\olfileid{inc}{tcp}{cax}

\olsection{\printtoken{S}{axiomatizable} نظریے}

نظریہ $\Th{T}$، \emph{!!{axiomatizable}} اے جے اوہدے
مسلّمیاں دا کوئی قابلِ حساب سیٹ~$A$ ہووے۔ ($A$ نوں $\Th{T}$ لئی
مسلّمیاں دا سیٹ آکھن دا مطلب اے کہ $T = \Setabs{!A}{A \Proves !A}$۔)
قدرتی عدداں دی کوئی وی “معقول” مسلّمہ بندی ایہ شرط پوری کرے گی۔
خاص طور اُتے ہر اوہ نظریہ !!{axiomatizable} اے جیہدے مسلّمیاں
دا سیٹ محدود اے۔

\begin{lem}
فرض کرو $\Th{T}$ !!{axiomatizable} اے۔ تاں $\Th{T}$
!!{computably enumerable} اے۔
\end{lem}

\begin{proof}
فرض کرو $A$، $\Th{T}$ لئی مسلّمیاں دا قابلِ حساب سیٹ اے۔
$!A \in T$ دی تصدیق لئی مسلّمیاں توں $!A$ دا !!a{derivation}
لبھو۔ ثبوت مل جائے تاں رکنیت ثابت ہو جاندی اے۔ جے ثبوت موجود
نہ ہووے، تاں ایہ تلاش ضروری نہیں کہ مُک جائے۔

دوجے لفظاں وچ $!A$، $\Th{T}$ وچ اُدوں تے صرف اُدوں اے
جدوں مسلّمیاں دی اک محدود فہرست $!B_1$، \dots، $!B_k$
موجود ہووے، جیہدے سارے مسلّمے سیٹ $A$ وچ ہون، تے پہلے درجے
دی منطق وچ $(!B_1 \land \dots \land !B_k) \lif !A$ دا
!!a{derivation} ہووے۔ اسیں پہلاں جان دے آں کہ “کوئی \dots
موجود اے جیہدے لئی \dots” دے روپ وچ تعریف ہو سکّن والا ہر
سیٹ !!{c.e.} اے، جے دوجے “\dots” والی شرط قابلِ حساب ہووے۔
\end{proof}
\endgroup
% /OLP-0306

% OLP-0307 content/incompleteness/theories-computability/complete-decidable.tex
\begingroup
\def\olfilename{complete-decidable}
\expandafter\def\csname ol@sectioncs\endcsname{section}

\olfileid{inc}{tcp}{cdc}

\olsection{\printtoken{S}{axiomatizable} تے مکمل نظریے قابلِ فیصلہ نیں}

اک نظریے نوں {\em مکمل} آکھدے آں جے ہر جملے $!A$ لئی،
$!A$ تے $\lnot !A$ وچوں گھٹ توں گھٹ اک ثابت ہو سکدا ہووے۔

\begin{lem}
فرض کرو نظریہ $\Th{T}$ مکمل تے !!{axiomatizable} اے۔
تاں $\Th{T}$ قابلِ فیصلہ اے۔
\end{lem}

\begin{proof}
فرض کرو $\Th{T}$ مکمل اے تے $A$ مسلّمیاں دا قابلِ حساب سیٹ اے۔
جے $\Th{T}$ بے تضاد نہیں، تاں اوہدا فیصلہ کرنا سوکھا اے۔
(الگورتھم ہر درست جملے لئی ایہو جواب دیندا اے، “ہاں۔”)
ہن اسیں اوہ صورت لیندے آں جیہدے وچ $\Th{T}$ بے تضاد اے۔
اسیں ایتھے نظریے دیاں دو ممکن صورتاں دا ثبوت دے رہے آں۔
مسلّمیاں دی فہرست توں بے تضادی دا فیصلہ کرن دا کوئی سانجھا
طریقہ فرض نہیں کیتا۔ عددی کوڈاں اُتے کم کردیاں پہلاں اوہناں
کوڈاں نوں رد کرو جیہڑے کسے جملے دے کوڈ نہیں۔

ایہ فیصلہ کرن لئی کہ جملہ $!A$، $\Th{T}$ وچ اے یا نہیں،
اوہ !!a{derivation} لبھو جیہڑا~$!A$ نوں~$\Th{T}$ توں اخذ
کرے، تے نال ای اوہ !!a{derivation} لبھو جیہڑا~$\lnot !A$ نوں
اخذ کرے۔ مسلّمیاں دے قابلِ حساب سیٹ نال ثبوت جانچدیاں دونوں
تلاشاں نوں واری واری اگے ودھاؤ۔ کیوں جے $\Th{T}$ مکمل اے،
دونہاں وچوں گھٹ توں گھٹ اک ثبوت ضرور لبھے گا۔ تے کیوں جے
$\Th{T}$ بے تضاد اے، جے $\lnot !A$ دا !!a{derivation} لبھ
جائے، تاں $!A$ دا کوئی !!{derivation} نہیں ہو سکدا۔

دوجے لفظاں وچ اسیں پہلاں جان دے آں کہ $\Th{T}$ !!{c.e.} اے۔
پہلاں ثابت کیتے قضیے مطابق، ایہ وکھانا کافی اے کہ $\Th{T}$ دا
متمم وی !!{c.e.} اے۔ پر اُتلے جملے دے سیاق وچ، !!a{formula}~$!A$،
~$\bar{\Th{T}}$ وچ اُدوں تے صرف اُدوں اے جدوں $\lnot !A$،
~$\Th{T}$ وچ اے۔ ایس لئی $\bar{\Th{T}} \leq_m \Th{T}$۔
درست جملیاں دے کوڈاں اُتے ایہ تبدیلی نفی دا کوڈ بناؤندی اے۔
جے متمم سارے قدرتی عدداں دے اندر لیندے آں، تاں نحوی جانچ نال
غلط جملہ کوڈ فوراً پچھان کے متمم دی گنتی وچ شامل کر لیندے آں۔
تقلیل اوہناں کوڈاں نوں کسے اک مقررہ ثابت جملے دے کوڈ اُتے
بھیج دیندی اے۔
\end{proof}
\endgroup
% /OLP-0307

% OLP-0308 content/incompleteness/theories-computability/first-incompleteness.tex
\begingroup
\def\olfilename{first-incompleteness}
\expandafter\def\csname ol@sectioncs\endcsname{section}

\olfileid{inc}{tcp}{inc}

\olsection{\texorpdfstring{$\Th{Q}$ دی کوئی مکمل، بے تضاد، !!^{axiomatizable}
  توسیع نہیں}{Q دی کوئی مکمل، بے تضاد، مسلّماتی طور متعین ہون جوگ توسیع نہیں}}

\begin{thm}
\ollabel{thm:first-incompleteness}
$\Th{Q}$ دی کوئی مکمل، بے تضاد، !!{axiomatizable} توسیع نہیں۔
\end{thm}

\begin{proof}
اسیں پہلاں ایہ ثابت کر چکے آں کہ $\Th{Q}$ دی کوئی بے تضاد
توسیع قابلِ فیصلہ نہیں۔ پر جے $\Th{T}$ مکمل تے !!{axiomatized}
اے، تاں اوہ قابلِ فیصلہ اے۔
\end{proof}

\begin{explain}
ایہ قضیہ گوڈل دے 1931 والے پہلے نامکملیت قضیے دی اصل صورت
دے نیڑے اے۔ نویں اصطلاحاں توں چھٹ، دو وڈے فرق نیں۔ گوڈل دی
شرط عام بے تضادی دی تھاں $\omega$-بے تضادی سی۔ اوہ “!!{axiomatizable}”
نوں پورے عموم نال وی نہیں ورت سکدا سی، کیوں جے قابلِ حسابیت دا
رسمی تصور اجے نہیں سی بنیا۔ (قابلِ حسابیت دے رسمی نمونے اگلے
دس ورھیاں وچ بنے۔ گوڈل نے وی ایس کم وچ حصہ پایا۔ ایہ نمونے
بنان دا اک وڈا مقصد ایہ پہچاننا سی کہ کہڑے نظرییاں اُتے گوڈل
دا نتیجہ لاگو ہو سکدا اے۔)

قضیہ آکھدا اے کہ توسیع وچ مکمل ہون، بے تضادی تے !!{axiomatizability}
تینوں کٹھیاں نہیں ہو سکدیاں۔ پر جے کوئی اک شرط چھڈ دیو، تاں
باقی دو پوریاں ہو سکدیاں نیں۔ $\Th{Q}$ بے تضاد اے تے اوہدی
مسلّمہ بندی قابلِ حساب اے، پر اوہ مکمل نہیں۔ تضاد والا نظریہ
مکمل اے تے اوہدی مسلّمہ بندی قابلِ حساب اے (مثال لئی
$\{ \eq/[0][0] \}$ نال)، پر اوہ بے تضاد نہیں۔ حساب دے سچے
جملیاں دا سیٹ مکمل تے بے تضاد اے، پر اوہدی مسلّمہ بندی قابلِ
حساب نہیں۔
\end{explain}
\endgroup
% /OLP-0308

% OLP-0309 content/incompleteness/theories-computability/inseparability.tex
\begingroup
\def\olfilename{inseparability}
\expandafter\def\csname ol@sectioncs\endcsname{section}

\olfileid{inc}{tcp}{ins}

\olsection{\texorpdfstring{$\Th{Q}$ وچ ثابت ہو سکّن والے تے قابلِ رد جملیاں دے
  سیٹ قابلِ حساب طریقے نال وکھ نہیں کیتے جا سکدے}{Q وچ ثابت ہو سکّن والے تے قابلِ رد جملیاں دے سیٹ قابلِ حساب طریقے نال وکھ نہیں کیتے جا سکدے}}

منّو $\bar{\Th{Q}}$ اوہناں جملیاں دا سیٹ اے جنہاں دیاں
{\em نفیاں} $\Th{Q}$ وچ ثابت ہو سکدیاں نیں، یعنی
$\bar{\Th{Q}} = \Setabs{!A}{\Th{Q} \Proves
  \lnot !A}$۔ ایہ قابلِ رد جملیاں دا سیٹ اے، نظریے دا پورا
متمم نہیں۔ یاد رکھو کہ بے سانجھے سیٹ $A$ تے $B$ قابلِ حساب
طریقے نال وکھ نہیں کیتے جا سکدے جے کوئی ایہو جیہا قابلِ حساب
سیٹ $C$ نہ ہووے جیہدے لئی $A
\subseteq C$ تے $B \subseteq \Complement{C}$۔

\begin{lem}
$\Th{Q}$ تے $\bar{\Th{Q}}$ قابلِ حساب طریقے نال وکھ نہیں کیتے
جا سکدے۔
\end{lem}

\begin{proof}
فرض کرو $C$ اک قابلِ حساب سیٹ اے جیہدے لئی $\Th{Q} \subseteq C$
تے $\bar{\Th{Q}} \subseteq \Complement{C}$۔ $R(x,y)$ نوں ایہ
تعلق منّ لئیے۔
\begin{quote}
$x$، فارمولے $!D(u)$ دا کوڈ اے تے $!D(\num y)$، $C$ وچ اے۔
\end{quote}
اسیں وکھاواں گے کہ $R(x,y)$ ہمہ گیر قابلِ حساب تعلق اے،
تے ایس توں تضاد نکلے گا۔ فارمولے دا کوڈ جانچنا، اوہدے وچ عددی
نام رکھنا تے دتے ہوئے قابلِ حساب سیٹ دی رکنیت جانچنا، تینوں کم
مؤثر طریقے نال ہو سکدے نیں۔ ایس لئی پورا تعلق قابلِ حساب اے۔
جے پہلا انپٹ فارمولے دا صحیح کوڈ نہ ہووے، تاں تعلق قائم نہیں ہوندا۔

فرض کرو $S(y)$ قابلِ حساب اے تے $!D_S(u)$، $\Th{Q}$ وچ
اوہدی نمائندگی کردا اے۔ فیر
\begin{eqnarray*}
S(\num n) & \lif & \Th{Q} \Proves !D_S(\num n) \\
& \lif & \formula D_S(\num n) \in C
\end{eqnarray*}
تے
\begin{eqnarray*}
\lnot S(\num n) & \lif & \Th{Q} \Proves \lnot !D_S(\num n) \\
& \lif & \formula D_S(\num n) \in \bar{\Th{Q}} \\
& \lif & \formula D_S(\num n) \not\in C
\end{eqnarray*}
ایس لئی $S(y)$ تے $R(\#(!D_S(u)),y)$ ہم ارز نیں۔ کوڈ آزاد
متغیر والے اصل فارمولے دا اے۔ ہر انپٹ لئی اوہدی عددی مثال
سیٹ دی رکنیت جانچن لئی ورتی جاندی اے۔ ایس طرح ہر اک رکنی
قابلِ حساب تعلق لئی اک مقرر قطار ملدی اے۔ فیر ہمہ گیر تعلق
فرض کرن نال پچھلا قطری تضاد نکلدا اے۔
% PNB-CMP-068: اصل آخری کوڈ وچ عددی مثال دی بجائے آزاد متغیر والا فارمولا؛واحد محدود ماخذی تصحیح۔
\end{proof}
\endgroup
% /OLP-0309

% OLP-0310 content/incompleteness/theories-computability/consis-with-q.tex
\begingroup
\def\olfilename{consis-with-q}
\expandafter\def\csname ol@sectioncs\endcsname{section}

\olfileid{inc}{tcp}{con}

\olsection{\texorpdfstring{$\Th{Q}$ نال کٹھے بے تضاد نظریے ناقابلِ فیصلہ نیں}{Q نال کٹھے بے تضاد نظریے ناقابلِ فیصلہ نیں}}

اگلے قضیے مطابق $\Th{Q}$ ناقابلِ فیصلہ اے، تے ہر اوہ نظریہ
وی ناقابلِ فیصلہ اے جیہڑا $\Th{Q}$ نال کٹھا بے تضاد ہووے۔

\begin{thm}
فرض کرو $\Th{T}$ حساب دی زبان دا کوئی نظریہ اے جیہڑا
$\Th{Q}$ نال کٹھا بے تضاد اے (یعنی $\Th{T} \cup \Th{Q}$
بے تضاد اے)۔ تاں $\Th{T}$ ناقابلِ فیصلہ اے۔
\end{thm}

\begin{proof}
یاد رکھو کہ $\Th{Q}$ دے مسلّمیاں دا سیٹ محدود اے، یعنی
$!Q_1$، \dots، $!Q_8$۔ انہاں دی جگہ اسیں اکّو مسلّمہ وی
رکھ سکدے آں، $!E = !Q_1 \land
\dots \land !Q_8$۔

فرض کرو $\Th{T}$ اک قابلِ فیصلہ نظریہ اے جیہڑا~$\Th{Q}$ نال
کٹھا بے تضاد اے۔ منّ لئیے کہ
\[
C = \Setabs{!A}{\Th{T} \Proves !E \lif !A}.
\]
اسیں وکھاواں گے کہ $C$، $\Th{Q}$ تے $\bar{\Th{Q}}$ نوں
وکھ کرن والا قابلِ حساب سیٹ اے۔ ایس توں تضاد نکلے گا۔
دتے ہوئے جملے توں مقرر شرطیہ بنانا مؤثر کم اے۔ فرض کیتے
فیصلہ کرن والے طریقے نال پتا لگ سکدا اے کہ ایہ شرطیہ ثابت
ہو سکدا اے یا نہیں۔ ایس لئی سیٹ قابلِ حساب اے۔ پہلاں جے
$!A$، $\Th{Q}$ وچ اے، تاں $!A$، $\Th{Q}$ دے مسلّمیاں توں
ثابت ہو سکدا اے۔ اخذ دے قضیے مطابق، پہلے درجے دی منطق وچ
$!E \lif !A$ دا !!a{derivation} اے۔ ایس لئی $!A$، $C$ وچ اے۔

دوجے پاسے، جے $!A$، $\bar{\Th{Q}}$ وچ اے، تاں پہلے درجے دی
منطق وچ $!E \lif \lnot !A$ دا ثبوت اے۔ جے $\Th{T}$،
$!E \lif !A$ وی ثابت کرے، تاں $\Th{T}$، $\lnot !E$ ثابت
کردا اے۔ اوس صورت وچ $\Th{T}
\cup \Th{Q}$ متضاد ہووے گا۔ پر فرض مطابق
$\Th{T} \cup \Th{Q}$ بے تضاد اے۔ ایس لئی $\Th{T}$،
$!E \lif !A$ ثابت نہیں کردا تے $!A$، $C$ وچ نہیں۔

اسیں وکھایا اے کہ جے $!A$، $\Th{Q}$ وچ اے، تاں اوہ $C$
وچ اے۔ جے $!A$، $\bar{\Th{Q}}$ وچ اے، تاں اوہ
$\Complement{C}$ وچ اے۔ ایس طرح $C$ دونہاں سیٹاں نوں
قابلِ حساب طریقے نال وکھ کردا اے۔ ایہی اوہ تضاد اے جیہڑا
اسیں لبھ رہے ساں۔
\end{proof}

ایہ قضیہ بڑا مضبوط اے۔ مثال لئی، ایس توں ایہ نتیجہ نکلدا اے۔

\begin{cor}
حساب دی زبان دی پہلے درجے دی منطق، یعنی سیٹ
\[
\Setabs{!A}{\text{$!A$ پہلے درجے دی منطق وچ ثابت ہو سکدا اے}}
\]
ناقابلِ فیصلہ اے۔
\end{cor}

\begin{proof}
پہلے درجے دی منطق $\emptyset$ دے منطقی نتیجیاں دا سیٹ اے۔
ایہ سیٹ~$\Th{Q}$ نال کٹھا بے تضاد اے۔
\end{proof}
\endgroup
% /OLP-0310

% OLP-0311 content/incompleteness/theories-computability/interpretability.tex
\begingroup
\def\olfilename{interpretability}
\expandafter\def\csname ol@sectioncs\endcsname{section}

\olfileid{inc}{tcp}{itp}

\olsection{\texorpdfstring{$\Th{Q}$ دی تعبیر جنہاں نظرییاں وچ ہو سکدی اے، اوہ ناقابلِ فیصلہ نیں}{Q دی تعبیر جنہاں نظرییاں وچ ہو سکدی اے، اوہ ناقابلِ فیصلہ نیں}}

اسیں ایہناں نتیجیاں نوں ہور مضبوط بنا سکدے آں۔ غیر رسمی
طور اُتے اک زبان $\Lang{L_1}$ دی دوجی زبان $\Lang{L_2}$ وچ
تعبیر وچ $\Lang{L_1}$ دے دائرے، تعلق دیاں علامتاں تے فنکشن
دیاں علامتاں دی تعریف $\Lang{L_2}$ دے !!{formula}s نال کرنی
پیندی اے۔ اسیں ایتھے اوہ تفصیل نہیں دیواں گے، پر ایس تعریف
نوں پوری طرح ٹھیک ٹھیک بیان کیتا جا سکدا اے۔

\begin{thm}
فرض کرو $\Th{T}$ ایہو جیہی زبان دا نظریہ اے جیہدے وچ
حساب دی زبان دی تعبیر ہو سکدی اے، تے تعبیر ایہو جیہی اے کہ
$\Th{T}$، $\Th{Q}$ دی تعبیر نال کٹھا بے تضاد اے۔ تاں
$\Th{T}$ ناقابلِ فیصلہ اے۔ جے $\Th{T}$، $\Th{Q}$ دے
مسلّمیاں دیاں تعبیراں ثابت کردا اے، تاں $\Th{T}$ دی کوئی
بے تضاد توسیع قابلِ فیصلہ نہیں۔
\end{thm}

پچھلے قضیے دے ثبوت وچ تھوڑی جیہی تبدیلی نال ایہ ثبوت مل
جاندا اے۔ کسے مخالف مثال نال $\Th{Q}$ تے $\bar{\Th{Q}}$ نوں
قابلِ حساب طریقے نال وکھ کرن والا سیٹ مل جائے گا۔ $\Th{ZFC}$،
یعنی انتخاب دے مسلّمے والا زرمیلو--فرینکل سیٹ نظریہ، ایسی
مسلّمی بنیاد اے جیہڑی عام ریاضی دا خاصا حصہ کرن لئی کافی اے۔
$\Th{ZFC}$ وچ قدرتی اعداد دی تعریف ہو سکدی اے، تے ایس تعبیر
دے حساب نال $\Th{Q}$ دے مسلّمے سچے نیں۔ ایس توں ایہ نتیجہ
نکلدا اے۔

\begin{cor}
% PNB-CMP-069: اصل نتیجے وچ بے تضادی دی غائب شرط بحال؛اُتلے قضیے دی شرط مطابق۔
$\Th{ZFC}$ دی کوئی بے تضاد تے قابلِ فیصلہ توسیع نہیں۔
\end{cor}

\begin{cor}
$\Th{ZFC}$ دی کوئی توسیع ایہو جیہی نہیں جیہڑی مکمل، بے تضاد تے
قابلِ حساب طریقے نال !!{axiomatizable} ہووے۔
\end{cor}

$\Th{ZFC}$ دی زبان وچ اکّو دو رکنی تعلق دی علامت، $\in$، اے۔
(حقیقت وچ برابری دی وی لوڑ نہیں۔) ایس توں ایہ نتیجہ نکلدا اے۔

\begin{cor}
دو رکنی تعلق دی علامت والی کسے وی زبان دی پہلے درجے دی
منطق ناقابلِ فیصلہ اے۔
\end{cor}

ایہ نتیجہ دو اک رکنی فنکشناں دیاں علامتاں والی کسے وی زبان
لئی وی سچ اے، کیوں جے ایہناں دوہاں نال دو رکنی تعلق دی علامت
دی نقل کیتی جا سکدی اے۔ انہاں نتیجیاں دیاں حدّاں ٹھیک ایتھے
ای نیں۔ جے زبان وچ ساریاں تعلق دیاں علامتاں \emph{اک رکنی} ہون،
تے ودھ توں ودھ اک \emph{اک رکنی} فنکشن دی علامت ہووے، تاں
اوہدی پہلے درجے دی منطق قابلِ فیصلہ اے۔

اک ہور دلچسپ گل۔ اسیں جان دے آں کہ $\Obj 0$، $'$، $+$،
$\times$، $<$ والیاں علامتاں دی زبان دے اوہناں جملیاں دا سیٹ
ناقابلِ فیصلہ اے جیہڑے معیاری نمونے وچ سچے نیں۔ اصل وچ $<$
نوں باقی علامتاں نال تعریف کیتا جا سکدا اے، تے فیر $+$ نوں
$\times$ تے $'$ نال۔ ایس لئی $\Obj 0$، $'$، $\times$ والی
زبان دے سچے جملیاں دا سیٹ ناقابلِ فیصلہ اے۔ دوجے پاسے،
پریسبورگر نے وکھایا کہ $\Obj 0$، $'$، $+$ والی حسابی زبان
دے اوہناں جملیاں دا سیٹ قابلِ فیصلہ اے جیہڑے معیاری نمونے وچ
سچے نیں۔ پر ایہ فیصلہ کرن لئی اینا حساب کرنا پیندا اے کہ
طریقہ عملی طور اُتے ورتن جوگ نہیں۔
\endgroup
% /OLP-0311

\OLEndChapterHook
\endgroup
% /OLP-0301

% OLP-0312 content/incompleteness/incompleteness-provability/incompleteness-provability.tex
\begingroup
\def\olfilename{incompleteness-provability}
\expandafter\def\csname ol@sectioncs\endcsname{section}

\olchapter{inc}{inp}{نامکملیت تے ثبوت پذیری}

% OLP-0313 content/incompleteness/incompleteness-provability/introduction.tex
\begingroup
\def\olfilename{introduction}
\expandafter\def\csname ol@sectioncs\endcsname{section}

\olfileid{inc}{inp}{int}

\olsection{تعارف}

ہلبرٹ دا خیال سی کہ قدرتی اعداد ورگی کسے ریاضیاتی ساخت دے
مسلّمیاں دا نظام اوہدوں ای کافی اے جدوں اوہدے توں اوس ساخت
بارے ہر سچا بیان اخذ ہو سکے۔ پچھوں رسمی نظاماں دے ذریعے اخذ وچ
اوہدی دلچسپی توں لگدا اے کہ اوہ ایہ ضمانت چاہندا سی کہ مثلاً
قدرتی اعداد بارے دلیل دین لئی ورتے رسمی نظام بے تضاد وی ہون
تے \emph{مکمل} وی۔ یعنی اوہناں دی زبان دا ہر بیان یا اوہدی نفی
!!{derivable} ہووے۔ گوڈل دا پہلا نامکملیت قضیہ وکھاؤندا اے کہ
مسلّمیاں دا ایہو جیہا کوئی نظام موجود نہیں۔ حساب دا کوئی رسمی
نظام مکمل، بے تضاد تے !!{axiomatizable} تینوں اکّو ویلے نہیں ہو سکدا۔
بلکہ کوئی وی “کافی طاقتور” بے تضاد تے !!{axiomatizable}
ریاضیاتی نظریہ مکمل نہیں۔

ہلبرٹ دا اک ہور، ودھ اہم مقصد، جدید (“کلاسیکی”) ریاضی دا
جواز قائم کرن دے پروگرام دا مرکز سی۔ اوہ کلاسیکی استدلال دی
نمائندگی کرن والے رسمی نظاماں دی بے تضادی دے ثبوت محدود
طریقیاں نال لبھنا چاہندا سی۔ ایس لئی گوڈل دا دوجا نامکملیت
قضیہ ہلبرٹ دے پروگرام لئی کہیں وڈا دھچکا سی۔ پہلے قضیے وانگ
ایہنوں وی موٹے طور اُتے بیان کیتا جا سکدا اے۔ حساب دا کوئی
کافی طاقتور نظریہ اپنی بے تضادی ثابت نہیں کر سکدا۔
“کافی طاقتور” آکھدیاں سانوں~$\Th{Q}$ توں کجھ ودھ طاقت
مراد لینی پئے گی۔

گوڈل دے نامکملیت قضیے دے اصل ثبوت دا خیال ایپیمنیدس دے
پیراڈوکس وچ لبھدا اے۔ کریٹ دا ایپیمنیدس آکھدا سی کہ کریٹ دے
سارے لوک جھوٹ بولدے نیں۔ ایس پیراڈوکس دی سدھی صورت ایہ
اے، “ایہ جملہ جھوٹا اے۔” سچائی دی تھاں !!{derivability} رکھ
کے گوڈل نے !!a{sentence} نوں رسمی روپ دتا جیہڑا بالواسطہ
اپنے بارے آکھدا اے کہ اوہ !!{derivable} نہیں۔ جے اوہ
!!{sentence} !!{derivable} ہووے، تاں نظریہ متضاد ہو جائے۔
گوڈل نے ایہ وی وکھایا کہ اوس !!{sentence} دی نفی زیرِ غور
مسلّمیاں دے نظام توں !!{derivable} نہیں۔ (ایس دوجے حصے لئی
گوڈل نوں فرض کرنا پیا کہ نظریہ~$\Th{T}$ “$\omega$-بے تضاد”
اے۔ $\omega$-بے تضادی، عام بے تضادی نال جڑی ہوئی پر ودھ
مضبوط خاصیت اے۔\footnote{یعنی ہر $\omega$-بے تضاد نظریہ
بے تضاد ہوندا اے، پر الٹی شرط لازم نہیں۔} گوڈل توں کجھ سال
پچھوں روسر نے وکھایا کہ نظریے~$\Th{T}$ دی عام بے تضادی
فرض کرنی کافی اے۔)

پہلا سوال ایہ اے کہ اپنے آپ ول حوالہ کرن والا !!a{sentence}
کِویں بنایا جائے۔ ہر !!{formula}~$!A$ لئی جیہڑا~$\Th{Q}$
دی زبان دا اے، منّ لئیے کہ~$\gn{!A}$ اوہ عددی نام اے جیہڑا
~$\Gn{!A}$ دے مطابق اے۔ ذرا غور کرو۔ $!A$،~$\Th{Q}$
دی زبان دا !!a{formula} اے۔ $\Gn{!A}$ اک قدرتی عدد اے،
تے $\gn{!A}$،~$\Th{Q}$ دی زبان دا اک \emph{ٹرم} اے۔
ایس لئی ہر !!{formula}~$!A$، جیہڑا~$\Th{Q}$ دی زبان دا اے،
دا اک \emph{ناں}،~$\gn{!A}$، اے۔ ایہ ناں~$\Th{Q}$ دی زبان
دا ٹرم اے۔ ایس طرح سانوں اوہ ڈھانچا ملدا اے جیہدے وچ
~$\Th{Q}$ دی زبان دے !!{formula}s دوجے !!{formula}s بارے
گل “کر” سکدے نیں۔ اگلے لمے نوں مقرر نقطے دا لمہ آکھدے آں۔

\begin{lem}
فرض کرو $\Th{T}$،~$\Th{Q}$ دی کوئی وی توسیع اے، تے $!B(x)$
کوئی ایہو جیہا !!{formula} اے جیہدے وچ صرف متغیر~$x$ آزاد اے۔
تاں اک !!a{sentence}~$!A$ موجود اے جیہدے لئی
$\Th{T} \Proves!A \liff !B(\gn{!A})$۔
\end{lem}

لمہ آکھدا اے کہ کوئی وی خاصیت $!B(x)$ دتی ہووے تاں اک
!!a{sentence}~$!A$ اے جیہڑا آکھدا اے، “$!B(x)$ میرے بارے
سچ اے۔” $\Th{T}$ ایہ گل “جاندا” اے۔

ایہو جیہا !!a{sentence} اسیں کِویں بناواں؟ ایپیمنیدس دے
پیراڈوکس دی ایہ صورت ویکھو جیہڑی کوائن نے دتی۔
\begin{quote}
“اپنا اقتباس اگے رکھن نال جھوٹ نکلدا اے” اپنا اقتباس اگے
رکھن نال جھوٹ نکلدا اے۔
\end{quote}
ایہ جملہ سِدھا اپنے آپ ول حوالہ نہیں کردا۔ ایہ بس اقتباس
دیاں نشانیاں وچ دتّی عبارت بارے اک نحوی دعویٰ کردا اے۔
ایس لحاظ نال ایہ ہِٹھلے جملیاں ورگا اے۔
\begin{enumerate}
\item “رابرٹ” اک سوہنا ناں اے۔
\item “میں دوڑیا۔” اک چھوٹا جملہ اے۔
\item “تِن لفظ نیں” وچ تِن لفظ نیں۔
\end{enumerate}
پر جے “اپنا اقتباس اگے رکھن نال جھوٹ نکلدا اے” والی عبارت
لئیے تے اوہدے اگے اوسے عبارت دا اقتباس رکھ دیئے، تاں کیہ
ہوندا اے؟ اوہی اصل جملہ بن جاندا اے!{} ایس طرح جملہ اپنے
بارے آکھدا اے کہ ایہ جھوٹا اے۔
\endgroup
% /OLP-0313

% OLP-0314 content/incompleteness/incompleteness-provability/fixed-point-lemma.tex
\begingroup
\def\olfilename{fixed-point-lemma}
\expandafter\def\csname ol@sectioncs\endcsname{section}

\olfileid{inc}{inp}{fix}

\olsection{مقرر نقطے دا لمہ}

\begin{explain}
مقرر نقطے دا لمہ آکھدا اے کہ کسے وی !!{formula}~$!B(x)$
لئی اک !!a{sentence}~$!A$ اے جیہدے لئی
$\Th{T} \Proves !A \liff !B(\gn{!A})$، بشرطیکہ $\Th{T}$،
~$\Th{Q}$ دی توسیع ہووے۔ جھوٹے جملے لئی اسیں چاہواں گے
کہ $!A$ دی،~$\Th{T}$ وچ، ایس بیان نال ہم ارزی ثابت ہووے،
“$\gn{!A}$ جھوٹا اے۔” یعنی $\Gn{!A}$ اک جھوٹے
!!{sentence} دا گوڈل نمبر اے۔ ثبوت دا خیال سمجھّن لئی $!A$
دا کوائن والی غیر رسمی صورت نال ملاپ ویکھو،
“اپنا اقتباس اگے رکھن نال جھوٹ نکلدا اے” اپنا اقتباس اگے
رکھن نال جھوٹ نکلدا اے۔ کسے عبارت دے اگے اوسے دا اقتباس
رکھ کے جملہ بناؤن نوں اوس عبارت دی \emph{قطری بناوٹ}
آکھ سکدے آں۔ جیہڑا جملہ بنے اوہ اوس عبارت دی قطری صورت
اے۔ ایس طرح “اپنا اقتباس اگے رکھن نال جھوٹ نکلدا اے” دی
قطری صورت اے، “اپنا اقتباس اگے رکھن نال جھوٹ نکلدا اے”
اپنا اقتباس اگے رکھن نال جھوٹ نکلدا اے۔
دھیان دیو کہ کوائن دا جھوٹا جملہ “جھوٹ نکلدا اے” دی قطری
صورت نہیں۔ ایہ “اپنا اقتباس اگے رکھن نال جھوٹ نکلدا اے”
دی قطری صورت اے۔ سو جھوٹا جملہ بناؤن والی خاصیت وچ خود
قطری بناوٹ شامل اے!{}

حساب دی زبان وچ اک آزاد متغیر والے !!a{formula} دا اقتباس
ایویں بناندے آں۔ اوہدا گوڈل نمبر کڈھدے آں تے آزاد متغیر دی
تھاں اوس نمبر دا معیاری عددی نام رکھ دیندے آں۔
~$!E(x)$ دی قطری صورت $!E(\num{n})$ اے، جتھے
$n = \Gn{!E(x)}$۔ (اگے توں $\num{\Gn{!E(x)}}$ نوں
مختصر کر کے $\gn{!E(x)}$ لکھاں گے۔) سو جے $!B(x)$
“جھوٹا اے” والی خاصیت ہووے، تاں “اپنا اقتباس اگے رکھن نال
جھوٹ نکلدا اے” دا مطلب ہووے گا، “اپنی قطری صورت دے گوڈل
نمبر اُتے لاؤن نال جھوٹ نکلدا اے۔” فرض کرو ساڈے کول
~$\Obj{diag}$ علامت اے، فنکشن $\fn{diag}(n)$ لئی۔
ایہ فنکشن گوڈل نمبر~$n$ والے !!{formula} دی قطری صورت
دا گوڈل نمبر کڈھدا اے۔ تاں اسیں $!E(x)$ نوں
$!B(\Obj{diag}(x))$ لکھ سکدے آں۔ فیر کوائن دے جھوٹے
جملے دا روپ اوہدی قطری صورت ہووے گا، یعنی $!E(\gn{!E(x)})$
یا $!B(\Obj{diag}(\gn{!B(\Obj{diag}(x))}))$۔
~$!B(x)$ کوئی ہور خاصیت وی ہو سکدی اے تے ایہو بناوٹ
اوہدے لئی وی چلے گی۔ نامکملیت دے قضیے لئی اسیں $!B(x)$
ایہ لواں گے، “$x$،~$\Th{T}$ وچ !!{derivable} نہیں۔”
فیر $!E(x)$ دا مطلب ہووے گا، “اپنی قطری صورت دے گوڈل
نمبر اُتے لاؤن نال ایہو جیہا !!a{sentence} ملدا اے جیہڑا
~$\Th{T}$ وچ !!{derivable} نہیں۔”

ایہ گل~$\Th{T}$ وچ رسمی بناؤن لئی سانوں $\fn{diag}$ نوں
رسمی روپ دین دا طریقہ لبھنا پئے گا۔ فنکشن $\fn{diag}(n)$
قابلِ حساب اے، بلکہ بنیادی بازگشتی اے۔ جے $n$ فارمولا
~$!E(x)$ دا گوڈل نمبر اے، تاں $\fn{diag}(n)$،
~$!E(\gn{!E(x)})$ دا گوڈل نمبر دیندا اے۔ (چیتے رکھو کہ
$\gn{!E(x)}$،~$!E(x)$ دے گوڈل نمبر دا معیاری عددی نام
اے، یعنی $\num{\Gn{!E(x)}}$۔) جے $\Obj{diag}$،
$\Th{T}$ وچ فنکشن $\fn{diag}$ دی نمائندگی کرن والی علامت
ہووے، تاں اسیں $!A$ نوں فارمولا
$!B(\Obj{diag}(\gn{!B(\Obj{diag}(x))}))$ منّ سکدے آں۔
دھیان دیو کہ
\begin{align*}
\fn{diag}(\Gn{!B(\Obj{diag}(x))}) & =
\Gn{!B(\Obj{diag}(\gn{!B(\Obj{diag}(x))}))} \\
& = \Gn{!A}.
\end{align*}
فرض کرو $\Th{T}$ ایہ !!{derive} کر سکدا اے،
\[
\Obj{diag}(\gn{!B(\Obj{diag}(x))}) = \gn{!A},
\]
تاں اوہ ایہ وی !!{derive} کر سکدا اے،
$!B(\Obj{diag}(\gn{!B(\Obj{diag}(x))}))
\liff !B(\gn{!A})$۔ تعریف مطابق ایس ہم ارزی دا کھبا
پاسا~$!A$ اے۔

عام طور اُتے $\Obj{diag}$، $\Th{T}$ دی زبان دی کوئی فنکشن
علامت نہیں ہووے گی، تے~$\Th{Q}$ دی تاں یقیناً نہیں۔ پر
$\fn{diag}$ قابلِ حساب اے، ایس لئی~$\Th{Q}$ وچ اک فارمولے
$!D_{\fn{diag}}(x,y)$ نال ایہ \emph{قابلِ نمائندگی} اے۔
ایس لئی $!B(\Obj{diag}(x))$ دی تھاں
$\lexists[y][(!D_{\fn{diag}}(x,y) \land !B(y))]$ لکھ
سکدے آں۔ ایہ بدل کرن نال اُپرلا ثبوت چل پیندا اے، تے ایہ
دلیل~$\Th{Q}$ وچ ای چل پیندی اے۔
\end{explain}

\begin{lem}
\ollabel{lem:fixed-point} فرض کرو $!B(x)$ کوئی فارمولا اے
جیہدے وچ صرف متغیر~$x$ آزاد اے۔ تاں اک جملہ~$!A$ اے
جیہدے لئی $\Th{Q} \Proves
!A \liff !B(\gn{!A})$۔
\end{lem}


\begin{proof}
دتے ہوئے $!B(x)$ لئی، $!E(x)$ نوں ایہ فارمولا منّ لئیے،
$\lexists[y][(!D_{\fn{diag}}(x,y) \land !B(y))]$۔
تے $!A$~اوہدی قطری صورت اے، یعنی فارمولا $!E(\gn{!E(x)})$۔

کیوں جے $!D_{\fn{diag}}$، $\fn{diag}$ دی نمائندگی کردا اے
تے $\fn{diag}(\Gn{!E(x)}) = \Gn{!A}$، ایس لئی $\Th{Q}$
ایہ !!{derive} کر سکدا اے،
\begin{align}
  & \formula D_{\fn{diag}}(\gn{!E(x)}, \gn{!A}) \ollabel{repdiag1} \\
  & \lforall[y][(!D_{\fn{diag}}(\gn{!E(x)},y) \lif
  \eq[y][\gn{!A}])]. \ollabel{repdiag2}
\end{align}
ہن اسیں وکھاواں گے کہ $\Th{Q} \Proves !A \liff !B(\gn{!A})$۔
ایہ دلیل غیر رسمی طور اُتے صرف منطق تے~$\Th{Q}$ وچ
!!{derivable} حقیقتاں نال دیواں گے۔

پہلاں فرض کرو~$!A$، یعنی $!E(\gn{!E(x)})$۔ $!E(x)$ دی
تعریف مطابق $!E(\gn{!E(x)})$ ایہو فارمولا اے،
\[
\lexists[y][(!D_{\fn{diag}}(\gn{!E(x)},y) \land !B(y))].
\]
کوئی ایہو جیہا~$y$ لئیے۔ کیوں جے
$!D_{\fn{diag}}(\gn{!E(x)},y)$، ایس لئی \olref{repdiag2}
مطابق $y = \gn{!A}$۔ سو $!B(y)$ توں~$!B(\gn{!A})$ ملدا اے۔

ہن فرض کرو $!B(\gn{!A})$۔ \olref{repdiag1} نال سانوں
\begin{align*}
& \formula D_{\fn{diag}}(\gn{!E(x)}, \gn{!A}) \land !B(\gn{!A}).
\intertext{ایس توں ایہ نکلدا اے کہ}
& \lexists[y][(!D_{\fn{diag}}(\gn{!E(x)},y) \land !B(y))].
\end{align*}
ایہ $!E(\gn{!E(x)})$ اے، یعنی~$!A$۔
\end{proof}

\begin{digress}
ایس ثبوت دا ملاپ قابلِ حسابیت دے نظریے وچ مقرر نقطے دے لمے
دے ثبوت نال کرو۔ ایتھے اسیں اک \emph{بیان} نوں اوہدے اپنے
حوالے نال متعین کرنا چاہندے آں۔ اُتھے اسیں اک \emph{فنکشن}
نوں اوہدے اپنے حوالے نال متعین کرنا چاہندے ساں۔ ایس فرق توں
بغیر خیال اکّو اے۔
\end{digress}

\begin{prob}
  !!^a{formula}~$!A(x)$ اک \emph{سچائی دی تعریف} اے جے
  $\Th{Q} \Proves !B \liff !A(\gn{!B})$ ہر !!{sentence}s~$!B$
  لئی ہووے۔ مقرر نقطے دا لمہ ورت کے وکھاؤ کہ کوئی !!{formula}
  سچائی دی تعریف نہیں اے۔
\end{prob}
\endgroup
% /OLP-0314

% OLP-0315 content/incompleteness/incompleteness-provability/first-incompleteness-thm.tex
\begingroup
\def\olfilename{first-incompleteness-thm}
\expandafter\def\csname ol@sectioncs\endcsname{section}

\olfileid{inc}{inp}{1in}

\olsection{پہلا نامکملیت قضیہ}

ہن اسیں پہلے نامکملیت قضیے دا گوڈل والا اصل ثبوت بیان
کر سکدے آں۔ فرض کرو $\Th{T}$ ایہو جیہا نظریہ اے جیہدی
مسلّمہ بندی قابلِ حساب اے، تے جیہدی زبان حساب دی زبان دی
توسیع اے۔ ایہ وی فرض کرو کہ $\Th{T}$ وچ $\Th{Q}$ دے
مسلّمے شامل نیں۔ خاص طور اُتے ایس توں نکلدا اے کہ
$\Th{T}$ قابلِ حساب فنکشناں تے تعلقاں دی نمائندگی کردا اے۔

اسیں ویکھیا اے کہ جے فارمولیاں تے ثبوتاں نوں عدداں نال
کوڈ کرن دا معقول طریقہ ہووے، تاں تعلق $\Prf[\Th{T}](x,y)$
قابلِ حساب اے۔ $\Prf[\Th{T}](x,y)$ اُدوں تے صرف اُدوں
قائم اے جدوں $x$ اوس !!a{derivation} دا گوڈل نمبر ہووے
جیہڑا گوڈل نمبر~$y$ والے !!{formula} دا~$\Th{T}$ وچ
اخذ اے۔ گوڈل نے اپنے خاص نظریے لئی اپنے مقالے وچ دتی 45
فنکشناں تے تعلقاں دی فہرست ورت کے وکھایا کہ ایہ تعلق بنیادی
بازگشتی اے۔ اوس فہرست دا 45واں تعلق، $x B y$، اوہی
$\Prf[\Th{T}](x,y)$ اے جیہڑا گوڈل دے خاص نظریے
~$\Th{T}$ لئی ورتیا جاندا اے۔ چیتے رکھو،
گوڈل اپنے مقالے وچ جیہنوں “بازگشتی” آکھدا اے، اسیں ہن
اوہنوں “بنیادی بازگشتی” آکھدے آں۔

کیوں جے $\Prf[\Th{T}](x,y)$ قابلِ حساب اے، ایس لئی
$\Th{T}$ وچ اوہدی نمائندگی ہو سکدی اے۔ اسیں اوہدی
نمائندگی کرن والے فارمولے نوں $\OPrf[\Th{T}](x,y)$ آکھاں
گے۔ $\OProv[\Th{T}](y)$ نوں فارمولا
$\lexists[x][\OPrf[\Th{T}](x,y)]$ منّ لئیے۔ ایہ گوڈل دی
فہرست دے 46ویں تعلق، $\fn{Bew}(y)$، نوں بیان کردا اے۔
گوڈل آکھدا اے کہ ایہ اکّو تعلق اے جیہدے بارے پکے طور اُتے
نہیں آکھیا جا سکدا کہ ایہ بازگشتی اے۔ اوہدا مطلب غالباً ایہ سی
کہ تعریف توں ایہ صاف نہیں کہ تعلق قابلِ حساب اے یا نہیں۔
پچھوں ہون والی پیش رفت توں پتا لگدا اے کہ ایہ قابلِ حساب نہیں۔

فرض کرو $\Th{T}$ اک !!{axiomatizable} نظریہ اے جیہدے وچ
~$\Th{Q}$ شامل اے۔ تاں $\Prf[\Th{T}](x, y)$ قابلِ فیصلہ
اے، ایس لئی~$\Th{Q}$ وچ !!a{formula}~$\OPrf[\Th{T}](x, y)$
نال اوہدی نمائندگی ہو سکدی اے۔ $\OProv[\Th{T}](y)$ اوہی
فارمولا اے جیہڑا اسیں اُتے بیان کیتا۔ مقرر نقطے دے لمے نال
اک فارمولا $!G_\Th{T}$ ملدا اے جیہدے لئی $\Th{Q}$،
تے ایس لئی $\Th{T}$ وی، ایہ !!{derive}s کردا اے،
\begin{equation}
\ollabel{eqn:qpf}
!G_\Th{T} \liff \lnot \OProv[\Th{T}](\gn{!G_\Th{T}}).
\end{equation}
دھیان دیو کہ $!G_\Th{T}$ دا مطلب ایہ اے، “$!G_\Th{T}$،
~$\Th{T}$ وچ !!{derivable} نہیں۔”

\begin{lem}\ollabel{lem:cons-G-unprov}
جے $\Th{T}$،~$\Th{Q}$ دی توسیع والا اک بے تضاد
!!{axiomatizable} نظریہ اے، تاں $\Th{T} \Proves/ !G_\Th{T}$۔
\end{lem}

\begin{proof}
فرض کرو $\Th{T}$، $!G_\Th{T}$ نوں !!{derive}s کردا اے۔
تاں اک !!a{derivation} \emph{ضرور موجود اے}۔ ایس لئی کسے عدد
$m$ لئی تعلق $\Prf[\Th{T}](m,
\Gn{!G_\Th{T}})$ قائم اے۔ فیر $\Th{Q}$، جملہ
$\OPrf[\Th{T}](\num m, \gn{!G_\Th{T}})$ !!{derive}s کردا اے۔
سو $\Th{Q}$ ایہ وی !!{derive}s کردا اے،
$\lexists[x][\OPrf[\Th{T}](x,\gn{!G_\Th{T}})]$، جیہڑا تعریف
مطابق $\OProv[\Th{T}](\gn{!G_\Th{T}})$ اے۔ \olref{eqn:qpf}
نال $\Th{Q}$، $\lnot !G_\Th{T}$ !!{derive}s کردا اے۔
کیوں جے $\Th{T}$ وچ $\Th{Q}$ شامل اے، ایس لئی
~$\Th{T}$ وی ایہ اخذ کردا اے۔ اسیں وکھایا اے کہ جے
$\Th{T}$، $!G_\Th{T}$ !!{derive}s کرے، تاں اوہ
$\lnot !G_\Th{T}$ وی !!{derive}s کرے گا۔ ایس طرح اوہ متضاد
ہو جائے گا۔
\end{proof}

\begin{defn}
\ollabel{thm:oconsis-q}
نظریہ $\Th{T}$، \emph{$\omega$-بے تضاد} اے جے ایہ شرط
پوری ہووے۔ جے $\lexists[x][!A(x)]$ کوئی وی جملہ اے تے
$\Th{T}$، $\lnot
!A(\num 0)$، $\lnot !A(\num 1)$، $\lnot !A(\num 2)$،
\dots سبھ !!{derive}s کردا اے، تاں $\Th{T}$،
$\lexists[x][!A(x)]$ ثابت نہیں کردا۔
\end{defn}

دھیان دیو کہ ہر $\omega$-بے تضاد نظریہ عام بے تضاد وی اے۔
وجہ ایہ اے کہ جے $\Th{T}$ متضاد اے، تاں $\Th{T} \Proves !A$
ہر~$!A$ لئی۔ خاص طور اُتے، جے $\Th{T}$ متضاد اے، تاں
اوہ $\lnot !A(\num n)$ ہر~$n$ لئی !!{derive}s کردا اے
تے~$\lexists[x][!A(x)]$ وی !!{derive}s کردا اے۔ سو جے
$\Th{T}$ متضاد اے، تاں اوہ $\omega$-بے تضاد نہیں۔ ایس
شرط دے عکسِ نقیض توں نکلدا اے کہ جے $\Th{T}$،
$\omega$-بے تضاد اے، تاں اوہ لازماً عام بے تضاد اے۔

\begin{lem}\ollabel{lem:omega-cons-G-unref}
جے $\Th{T}$ اک $\omega$-بے تضاد !!{axiomatizable} نظریہ اے
جیہڑا~$\Th{Q}$ دی توسیع اے، تاں
$\Th{T} \Proves/ \lnot !G_\Th{T}$۔
\end{lem}

\begin{proof}
اسیں وکھاواں گے کہ جے $\Th{T}$، $\lnot !G_\Th{T}$
!!{derive}s کرے، تاں اوہ $\omega$-بے تضاد نہیں۔ فرض کرو
$\Th{T}$، $\lnot !G_\Th{T}$ !!{derive}s کردا اے۔ جے
$\Th{T}$ متضاد اے، تاں اوہ $\omega$-بے تضاد نہیں تے
دلیل پوری اے۔ دوجی صورت وچ $\Th{T}$ بے تضاد اے، ایس لئی
اوہ $!G_\Th{T}$ !!{derive} نہیں کردا، \olref{lem:cons-G-unprov}
مطابق۔ کیوں جے $!G_\Th{T}$ دا کوئی !!{derivation}
$\Th{T}$ وچ موجود نہیں، ایس لئی $\Th{Q}$ ایہ سبھ
!!{derive}s کردا اے،
\[
\lnot \OPrf[\Th{T}](\num 0, \gn{!G_\Th{T}}), \lnot \OPrf[\Th{T}](\num 1,
\gn{!G_\Th{T}}), \lnot \OPrf[\Th{T}](\num 2, \gn{!G_\Th{T}}), \dots
\]
تے~$\Th{T}$ وی ایہ سبھ اخذ کردا اے۔ دوجے پاسے
\olref{eqn:qpf} مطابق $\lnot
!G_\Th{T}$، $\lexists[x][\OPrf[\Th{T}](x,\gn{!G_\Th{T}})]$
دے ہم ارز اے۔ ایس لئی $\Th{T}$، $\omega$-بے تضاد نہیں۔
\end{proof}

\begin{prob}
ہر $\omega$-بے تضاد نظریہ عام بے تضاد اے۔ وکھاؤ کہ الٹی
شرط لازم نہیں، یعنی کجھ نظریے بے تضاد نیں پر
$\omega$-بے تضاد نہیں۔ ایس لئی وکھاؤ کہ $\Th{Q} \cup
\{\lnot !G_\Th{Q}\}$ عام بے تضاد اے پر $\omega$-بے تضاد
نہیں۔
\end{prob}

\begin{thm}
\ollabel{thm:first-incompleteness}
فرض کرو $\Th{T}$ کوئی وی $\omega$-بے تضاد
!!{axiomatizable} نظریہ اے جیہڑا~$\Th{Q}$ دی توسیع اے۔
تاں $\Th{T}$ مکمل نہیں۔
\end{thm}

\begin{proof}
جے $\Th{T}$، $\omega$-بے تضاد اے، تاں اوہ عام بے تضاد
وی اے۔ ایس لئی $\Th{T}
\Proves/ !G_\Th{T}$، \olref{lem:cons-G-unprov} مطابق۔
\olref{lem:omega-cons-G-unref} مطابق
$\Th{T} \Proves/ \lnot !G_\Th{T}$۔ ایس توں نکلدا اے کہ
$\Th{T}$ نامکمل اے، کیوں جے اوہ نہ $!G_\Th{T}$ !!{derive}s
کردا اے تے نہ $\lnot !G_\Th{T}$۔
\end{proof}
\endgroup
% /OLP-0315

% OLP-0316 content/incompleteness/incompleteness-provability/rosser-thm.tex
\begingroup
\def\olfilename{rosser-thm}
\expandafter\def\csname ol@sectioncs\endcsname{section}

\olfileid{inc}{inp}{ros}

\olsection{روسر دا قضیہ}

کی اسیں گوڈل دا ثبوت بدل کے ہور مضبوط نتیجہ لے سکدے آں،
جیہدے وچ “$\omega$-بے تضاد” دی تھاں صرف “بے تضاد” کافی
ہووے؟ روسر دے لبھے طریقے نال جواب “ہاں” اے۔ روسر دی ترکیب
ایہ اے کہ اک بدلیا ہویا !!{derivability} محمول
$\ORProv_T(y)$، $\OProv[\Th{T}](y)$ دی تھاں ورتیا جائے۔

\begin{thm}
\ollabel{thm:rosser}
فرض کرو $\Th{T}$ کوئی وی بے تضاد !!{axiomatizable} نظریہ
اے جیہڑا $\Th{Q}$ دی توسیع اے۔ تاں $\Th{T}$ مکمل نہیں۔
\end{thm}

\begin{proof}
چیتے رکھو کہ $\OProv[\Th{T}](y)$ دی تعریف
$\lexists[x][\OPrf[\Th{T}](x,
  y)]$ اے۔ ایتھے $\OPrf[\Th{T}](x, y)$ اوس قابلِ فیصلہ
تعلق دی نمائندگی کردا اے جیہڑا اُدوں تے صرف اُدوں قائم اے
جدوں $x$، اوس !!a{derivation} دا گوڈل نمبر ہووے جیہڑا گوڈل
نمبر~$y$ والے !!{sentence} دا اخذ اے۔ ایس طرح $x$ تے~$y$
وچلا اوہ تعلق وی قابلِ فیصلہ اے جیہڑا اُدوں قائم اے جدوں $x$،
گوڈل نمبر~$y$ والے جملے دی \emph{تردید} دا گوڈل نمبر ہووے۔
$\fn{not}(x)$ نوں ایہو جیہا بنیادی بازگشتی فنکشن منّ لئیے کہ
جے $x$، فارمولا $!A$ دا کوڈ ہووے، تاں $\fn{not}(x)$،
$\lnot
!A$ دا کوڈ اے۔ تاں $\Refut[\Th{T}](x, y)$ اُدوں تے صرف
اُدوں قائم اے جدوں $\Prf[\Th{T}](x, \fn{not}(y))$ قائم اے۔
منّ لئیے کہ $\ORefut[\Th{T}](x, y)$ اوہدی نمائندگی کردا اے۔
تاں جے $\Th{T} \Proves \lnot !A$ تے $\delta$ متعلقہ
!!{derivation} اے، تاں $\Th{Q} \Proves
\ORefut[\Th{T}](\gn{\delta}, \gn{!A})$۔ اسیں
$\ORProv[\Th{T}](y)$ دی تعریف ایویں کردے آں،
\[
\lexists[x][(\OPrf[\Th{T}](x,y) \land \lforall[z][(z < x \lif \lnot
  \ORefut[\Th{T}](z,y))])].
\]
موٹے طور اُتے $\ORProv[\Th{T}](y)$ آکھدا اے کہ $y$ دا
$\Th{T}$ وچ اک ثبوت اے، تے $y$ دی کوئی تردید ایس ثبوت
توں چھوٹے کوڈ والی نہیں۔ جے $\Th{T}$ بے تضاد اے، تاں سچائی
دے لحاظ نال $\ORProv[\Th{T}](y)$ اوہناں ای عدداں لئی قائم اے
جنہاں لئی $\OProv[\Th{T}](y)$ قائم اے۔ پر $\Th{T}$ وچ
\emph{ثبوت پذیری} دے لحاظ نال دونہاں دیاں خاصیتاں وکھریاں نیں۔
ہن اسیں سچائی تے ثبوت پذیری دا فرق جان دے آں۔ جے $\Th{T}$
\emph{مت}ضاد اے، تاں ایہ دونیں اوہناں ای عدداں لئی
قائم \emph{نہیں} ہوندے!{} (اکثر $\ORProv[\Th{T}](y)$
نوں “$y$ روسر دے لحاظ نال ثابت ہو سکدا اے” پڑھیا جاندا اے۔
پر ایس نال بھلیکھا پے سکدا اے۔ روسر والی خاصیت ثبوت پذیری
دی کوئی خاص قسم نہیں۔ متضاد نظرییاں وچ کجھ !!{sentence}s
ثابت ہو سکدے نیں پر روسر والی شرط پوری نہیں کردے۔ بھلیکھے
توں بچن لئی اسیں بس “$y$ شمووَبل (shmovable) اے” وی پڑھ
سکدے آں۔)

مقرر نقطے دے لمے نال اک فارمولا $!R_\Th{T}$ ملدا اے جیہدے لئی
\begin{equation}
  \Th{Q} \Proves !R_\Th{T} \liff \lnot \ORProv[\Th{T}](\gn{!R_\Th{T}}).
  \ollabel{RT}
\end{equation}
\olref[1in]{thm:first-incompleteness} دے ثبوت توں وکھ،
ایتھے اسیں ایہ وکھاواں گے کہ جے $\Th{T}$ بے تضاد اے، تاں
$\Th{T}$، $!R_\Th{T}$ نوں !!{derive} نہیں کردا تے
$\Th{T}$، $\lnot !R_\Th{T}$ نوں وی !!{derive} نہیں کردا۔
(یعنی ایتھے $\omega$-بے تضادی دا مفروضہ ضروری نہیں۔)

پہلاں وکھائیے کہ $\Th{T} \Proves/ !R_{\Th{T}}$۔ فرض کرو
ایہ ثابت ہوندا اے۔ تاں $!R_{\Th{T}}$ دا $T$ توں
!!a{derivation} اے۔ اوہدا گوڈل نمبر $n$ منّ لئیے۔
تاں $\Th{Q} \Proves \OPrf[\Th{T}](\num{n}, \gn{!R_{\Th{T}}})$،
کیوں جے $\OPrf[\Th{T}]$، $\Prf[\Th{T}]$ دی~$\Th{Q}$
وچ نمائندگی کردا اے۔ نالے ہر $k < n$ لئی، $k$،
$\lnot !R_{\Th{T}}$ دے کسے !!a{derivation} دا گوڈل
نمبر نہیں، کیوں جے $\Th{T}$ بے تضاد اے۔ سو ہر $k < n$ لئی
$\Th{Q} \Proves \lnot
\ORefut[\Th{T}](\num{k}, \gn{!R_{\Th{T}}})$۔
\olref[req][min]{lem:less-nsucc} مطابق
$\Th{Q} \Proves \lforall[z][(z < \num{n} \lif \lnot \ORefut[\Th{T}](z,
  \gn{!R_{\Th{T}}}))]$۔ ایس لئی
\[
\Th{Q} \Proves \lexists[x][(\OPrf[\Th{T}](x,\gn{!R_{\Th{T}}}) \land \lforall[z][(z
    < x \lif \lnot \ORefut[\Th{T}](z,\gn{!R_{\Th{T}}}))])],
\]
تے ایہ اوہی $\ORProv[\Th{T}](\gn{!R_{\Th{T}}})$ اے۔
\olref{RT} نال $\Th{Q} \Proves \lnot !R_{\Th{T}}$۔
کیوں جے $\Th{T}$، $\Th{Q}$ دی توسیع اے، ایس لئی
$\Th{T} \Proves \lnot !R_{\Th{T}}$ وی۔ پر اسیں فرض کیتا اے کہ
$\Th{T} \Proves !R_{\Th{T}}$۔ تاں $\Th{T}$ متضاد ہو جائے گا،
جیہڑا قضیے دے مفروضے دے خلاف اے۔

ہن وکھائیے کہ $\Th{T} \Proves/ \lnot !R_{\Th{T}}$۔ فرض کرو
ایہ ثابت ہوندا اے، تے $n$، $\lnot !R_{\Th{T}}$ دے کسے
!!a{derivation} دا گوڈل نمبر اے۔ تاں
$\Refut[\Th{T}](n, \Gn{!R_{\Th{T}}})$ قائم اے۔ کیوں جے
$\ORefut[\Th{T}]$، $\Refut[\Th{T}]$ دی $\Th{Q}$ وچ
نمائندگی کردا اے، ایس لئی $\Th{Q} \Proves
\ORefut[\Th{T}](\num{n}, \gn{!R_{\Th{T}}})$۔ اسیں فیر
وکھاواں گے کہ $\Th{T}$ متضاد ہو جائے گا، کیوں جے اوہ
$!R_{\Th{T}}$ وی !!{derive} کرے گا۔ ساڈے کول اے کہ
\begin{align*}
\Th{Q} & \Proves !R_{\Th{T}} \liff \lnot \ORProv[\Th{T}](\gn{!R_{\Th{T}}}),
\intertext{تے کیوں جے $\Th{T}$،~$\Th{Q}$ دی توسیع اے، ایہ وکھانا کافی اے کہ}
\Th{Q} & \Proves \lnot \ORProv[\Th{T}](\gn{!R_{\Th{T}}}).
\end{align*}
!!{sentence} $\lnot
\ORProv[\Th{T}](\gn{!R_{\Th{T}}})$، یعنی
\begin{align*}
  \lnot & \lexists[x][(\OPrf[\Th{T}](x,\gn{!R_{\Th{T}}}) \land
    \lforall[z][(z < x \lif \lnot \ORefut[\Th{T}](z,\gn{!R_{\Th{T}}}))])],
  \intertext{منطقی طور اُتے ایس دے ہم ارز اے}
  & \lforall[x][(\OPrf[\Th{T}](x,\gn{!R_{\Th{T}}}) \lif
    \lexists[z][(z < x \land \ORefut[\Th{T}](z,\gn{!R_{\Th{T}}}))])].
\end{align*}
اسیں منطق تے $\Th{Q}$ دے !!{derive}s کیتے حقائق نال
غیر رسمی دلیل دیواں گے۔ کوئی وی $x$ لئیے تے فرض کریے کہ
$\OPrf[\Th{T}](x,
\gn{!R_{\Th{T}}})$۔ اسیں پہلاں جان دے آں کہ
$\Th{T} \Proves/ !R_{\Th{T}}$، ایس لئی ہر مقرر $k$ لئی
$\Th{Q} \Proves \lnot \OPrf[\Th{T}](\num{k}, \gn{!R_{\Th{T}}})$۔
سو ہر $k$ لئی $\eq/[x][\num{k}]$۔ خاص طور اُتے (a)
$\eq/[x][\num{n}]$۔ نالے
$\lnot(\eq[x][\num{0}]
\lor \eq[x][\num{1}] \lor \dots \lor \eq[x][\num{n-1}])$۔
ایس توں \olref[req][min]{lem:less-nsucc} نال (b)
$\lnot(x < \num{n})$ ملدا اے۔
\olref[req][min]{lem:trichotomy} نال $\num{n} < x$ ملدا اے۔
کیوں جے $\Th{Q} \Proves
\ORefut[\Th{T}](\num{n}, \gn{!R_{\Th{T}}})$، ایس لئی
$\num{n} < x \land
\ORefut[\Th{T}](\num{n}, \gn{!R_{\Th{T}}})$، تے ایس توں
$\lexists[z][(z < x
\land \ORefut[\Th{T}](z,\gn{!R_{\Th{T}}}))]$ ملدا اے۔
کیوں جے $x$ کوئی وی ہو سکدا سی، ایس لئی مطلوب نتیجہ ملدا اے،
\[
\lforall[x][(\OPrf[\Th{T}](x,\gn{!R_{\Th{T}}}) \lif
  \lexists[z][(z < x \land \ORefut[\Th{T}](z,\gn{!R_{\Th{T}}}))])].
\]
\end{proof}

\begin{prob}
قدرتی عدداں دے دو سیٹ $A$ تے $B$ نوں
\emph{قابلِ حساب طریقے نال وکھ نہ ہو سکّن والے} آکھدے آں
جے کوئی ایہو جیہا !!{decidable} سیٹ~$X$ نہ ہووے جیہدے لئی
$A \subseteq X$ تے $B \subseteq \Complement{X}$۔
(ایتھے $\Complement{X}$، $\Nat \setminus X$ اے، یعنی
$X$ دا متمم۔) فرض کرو $\Th{T}$، $\Th{Q}$ دی بے تضاد
!!{axiomatizable} توسیع اے۔ فرض کرو $A$ اوہناں !!{sentence}s
دے گوڈل نمبراں دا سیٹ اے جیہڑے~$\Th{T}$ وچ ثابت ہو سکدے
نیں، تے $B$ اوہناں جملیاں دے گوڈل نمبراں دا سیٹ اے جنہاں دی
تردید~$\Th{T}$ وچ ہو سکدی اے۔ ثابت کرو کہ $A$ تے~$B$
قابلِ حساب طریقے نال وکھ نہیں کیتے جا سکدے۔
\end{prob}
\endgroup
% /OLP-0316

% OLP-0317 content/incompleteness/incompleteness-provability/godels-paper.tex
\begingroup
\def\olfilename{godels-paper}
\expandafter\def\csname ol@sectioncs\endcsname{section}

\olfileid{inc}{inp}{gop}

\olsection{گوڈل دے اصل مقالے نال تقابل}

گوڈل دا 1931 والا مقالہ پڑھنا فائدے دا اے۔ اوہدی تمہید وچ
انہاں خیالات دا خاکہ اے جنہاں اُتے اسیں ہن گل کیتی اے۔
مقالے اُتے سرسری نظر نال وی ہر قدم دی گل سمجھ آ جاندی اے۔
پہلاں گوڈل رسمی نظام $P$ دی نحو، مسلّمے تے ثبوت دے قاعدے
دسدا اے۔ فیر بنیادی بازگشتی فنکشناں تے تعلقاں دی تعریف
کردا اے۔ فیر وکھاؤندا اے کہ $x B y$ بنیادی بازگشتی اے،
تے دلیل دیندا اے کہ بنیادی بازگشتی فنکشناں تے تعلقاں دی
$\Th{P}$ وچ نمائندگی ہو سکدی اے۔ ایس توں بعد اوہ اُتے
دسے طریقے وانگ نامکملیت دا قضیہ ثابت کردا اے۔ تیجے حصے وچ
اوہ وکھاؤندا اے کہ ثابت نہ ہو سکّن والے دعوے نوں حساب دی
زبان دا جملہ بنایا جا سکدا اے۔ ایتھوں $\beta$-لمہ نکلیا۔
اسیں وی ایہ لمہ لڑیآں سنبھالن لئی ورتیا سی جدوں اسیں وکھایا
کہ بازگشتی فنکشناں دی $\Th{Q}$ وچ نمائندگی ہو سکدی اے۔
گوڈل کافی ہون والے مسلّمیاں دا سبھ توں چھوٹا سیٹ وکھ نہیں
کردا، پر ہن سانوں پتا اے کہ $\Th{Q}$ کافی اے۔ آخر وچ
چوتھے حصے اندر اوہ دوجے نامکملیت قضیے دے ثبوت دا خاکہ دیندا اے۔
\endgroup
% /OLP-0317

% OLP-0318 content/incompleteness/incompleteness-provability/provability-conditions.tex
\begingroup
\def\olfilename{provability-conditions}
\expandafter\def\csname ol@sectioncs\endcsname{section}

\olfileid{inc}{inp}{prc}

\olsection{\texorpdfstring{پیانو حساب $\Th{PA}$ لئی \usetoken{S}{derivability} دیاں شرطاں}{پیانو حساب PA لئی اخذ ہون دی اہلیت دیاں شرطاں}}

پیانو حساب، یعنی $\Th{PA}$، اوہ نظریہ اے جیہڑا
$\Th{Q}$ وچ سارے !!{formula}s لئی استقرا دے مسلّمے
شامل کردا اے۔ دوجے لفظاں وچ، $\Th{Q}$ وچ ایہو جیہے
مسلّمے شامل کیتے جاندے نیں،
\[
(!A(0) \land \lforall[x][(!A(x) \lif !A(x'))]) \lif \lforall[x][!A(x)]
\]
ہر !!{formula}~$!A$ لئی۔ دھیان دیو کہ ایہ دراصل اک
\emph{خاکہ} اے، یعنی لامتناہی مسلّمے۔ پتا لگدا اے کہ
$\Th{PA}$ نوں محدود مسلّمیاں نال متعین {\em نہیں} کیتا
جا سکدا۔ پر مؤثر طریقے نال ایہ فیصلہ ہو سکدا اے کہ علامتاں
دی دتی ہوئی لڑی استقرا دے مسلّمے دی مثال اے یا نہیں۔
ایس لئی $\Th{PA}$ دے مسلّمیاں دا سیٹ قابلِ حساب اے۔
$\Th{PA}$،~$\Th{Q}$ نالوں کہیں مضبوط نظریہ اے۔ مثال لئی،
استقرا دے عام طریقے نال جمع تے ضرب دی تبادلی خاصیت
سوکھیاں ثابت ہو سکدی اے۔ اصل وچ عدد نظریے تے امتزاجیات
دیاں محدود وسیلیاں والیاں زیادہ تر دلیلیں~$\Th{PA}$ وچ
دِتّیاں جا سکدیاں نیں۔

کیوں جے $\Th{PA}$ دی مسلّمہ بندی قابلِ حساب اے، ایس لئی
!!{derivability} محمول $\Prf[\Th{PA}](x,y)$ قابلِ حساب اے
تے اوہدی~$\Th{Q}$ وچ نمائندگی ہو سکدی اے (تے ایس لئی
~$\Th{PA}$ وچ وی)۔ پہلے وانگ $\OPrf[\Th{PA}](x,y)$ نوں
ایس تعلق دی نمائندگی کرن والا فارمولا منّ لئیے۔
$\OProv[\Th{PA}](y)$ نوں ایہ فارمولا منّ لئیے،
% PNB-CMP-070: عددی Prf دی تھاں رسمی نمائندہ OPrf؛frozen source B7 دا واحد macro correction۔
$\lexists[x][\OPrf[\Th{PA}](x,y)]$۔ سوکھے لفظاں وچ ایہ
آکھدا اے، “$y$، $\Th{PA}$ دے مسلّمیاں توں !!{derivable}
اے۔” سانوں $\Th{Q}$ دے مسلّمیاں توں کجھ ودھ ایس لئی
چاہیدا اے کہ نظریہ اینا مضبوط ہووے جو ایس !!{derivability}
محمول بارے کجھ بنیادی حقائق !!{derive} کر سکے۔ سانوں
ایہ حقائق چاہیدے نیں۔
\begin{enumerate}
\item[P1.] جے $\Th{PA} \Proves !A$، تاں $\Th{PA} \Proves
  \OProv[\Th{PA}](\gn{!A})$۔
\item[P2.] سارے !!{formula}s $!A$ تے $!B$ لئی،
  \[
  \Th{PA} \Proves \OProv[\Th{PA}](\gn{!A \lif !B}) \lif
  (\OProv[\Th{PA}](\gn{!A}) \lif \OProv[\Th{PA}](\gn{!B})).
  \]
\item[P3.] ہر !!{formula}~$!A$ لئی،
  \[
  \Th{PA} \Proves \OProv[\Th{PA}](\gn{!A})
  \lif \OProv[\Th{PA}](\gn{\OProv[\Th{PA}](\gn{!A})}).
  \]
\end{enumerate}
انہاں تنّاں خاصیتاں دی تصدیق لئی !!{formula}
$\OProv[\Th{PA}](y)$ نوں ٹھیک ٹھیک بیان کرنا پئے گا تے
$\Th{PA}$ دے مسلّمیاں نال متعلقہ رسمی !!{derivation}s
دینے پین گے۔ شرطاں (1) تے~(2) سوکھیاں نیں۔ اصلی محنت
شرط~(3) لئی کرنی پیندی اے۔ (ذرا سوچو ایس وچ کیہو جیہا
کم کرنا پئے گا \dots) تفصیل وچ جانا لمبا تے ایتھے
غیر دلچسپ ہووے گا، ایس لئی ہن منّ لئیے کہ
$\Th{PA}$ لئی اُتے دتیاں تنّے خاصیتاں قائم نیں۔
$\OProv[\Th{PA}](y)$ دی کوئی معقول چون ہِٹھلی شرط وی
پوری کرے گی۔
\begin{enumerate}
\item[P4.] جے $\Th{PA} \Proves \OProv[\Th{PA}](\gn{!A})$،
  تاں $\Th{PA} \Proves !A$۔
\end{enumerate}
پر سانوں ایس حقیقت دی لوڑ نہیں پئے گی۔

\begin{digress}
ویسے گوڈل نے وی ساڈے وانگ تفصیل چھڈ دتی سی۔ 1931 دے
مقالے دے آخر وچ اوہ دوجے نامکملیت قضیے دے ثبوت دا خاکہ
دیندا اے تے آکھدا اے کہ پوری تفصیل اگلے مقالے وچ دیوے گا۔
اوہ ایہ کم کرن تک کدے پہنچیا ای نہیں۔ دلیل سمجھّن والے
سارے لوک منّ دے سن کہ ثبوت پورا کیتا جا سکدا اے، ایس لئی
گوڈل نوں تفصیل لکھن دی لوڑ نہ پئی۔
\end{digress}
\endgroup
% /OLP-0318

% OLP-0319 content/incompleteness/incompleteness-provability/second-incompleteness-thm.tex
\begingroup
\def\olfilename{second-incompleteness-thm}
\expandafter\def\csname ol@sectioncs\endcsname{section}

\olfileid{inc}{inp}{2in}

\olsection{نامکملیت دا دوجا قضیہ}

اسیں ایہہ گل کیویں بیان کریے کہ $\Th{PA}$ اپنی بے تضادی نوں
ثابت نہیں کردا؟ $\Th{PA}$ نوں متضاد کہنا ایہہ کہنا اے کہ
$\Th{PA} \Proves \eq[0][1]$۔ ایس لئی بے تضادی دے جملے
$\OCon[\Th{PA}]$ نوں !!{sentence} $\lnot
\OProv[\Th{PA}](\gn{\eq[0][1]})$ دے روپ وچ لے سکدے آں؛فیر
ہِٹھلا قضیہ مطلوبہ نتیجہ دیندا اے:

\begin{thm}
\ollabel{thm:second-incompleteness} 
جے $\Th{PA}$ بے تضاد اے تاں $\Th{PA}$، $\OCon[\Th{PA}]$
نوں !!{derive} نہیں کردا۔
\end{thm}

ایہہ دھیان وچ رکھنا اہم اے کہ قضیہ $\OCon[\Th{PA}]$ دی خاص
نمائندگی اُتے منحصر اے،یعنی $\OProv[\Th{PA}](y)$ دی خاص
نمائندگی اُتے۔ سانوں صرف ایہہ چاہیدا اے کہ
$\OProv[\Th{PA}](y)$ دی نمائندگی تین !!{derivability} شرطاں
پوری کرے؛ایس لئی قضیہ ہر اُس نظریے تک وی پھیلدا اے جس دا
!!a{derivability} محمول ایہناں خاصیتاں نوں پورا کردا ہووے۔

گوڈل دی دلیل دا خاکہ پڑھنا فائدہ مند اے،کیونکہ قضیہ چنگے
اختتامی نکتے وانگ سامنے آ جاندا اے۔ دلیل ایویں اے۔
$!G_\Th{PA}$ اوہ گوڈل جملہ ہووے جیہڑا اسیں
\olref[1in]{thm:first-incompleteness} دے ثبوت وچ بنایا سی۔ اسیں
دکھا چکے آں: «جے $\Th{PA}$ بے تضاد اے تاں $\Th{PA}$،
$!G_\Th{PA}$ نوں !!{derive} نہیں کردا»۔ جے اس گل نوں
$\Th{PA}$ دے \emph{اندر} رسمی بنائیے،تاں ہِٹھلا نتیجہ ثابت
ہوندا اے:
\[
\OCon[\Th{PA}] \lif \lnot \OProv[\Th{PA}](\gn{!G_\Th{PA}}).
\]
% PNB-CMP-071: B9 bare Prov دی تھاں formal OProv؛اصل اگلی تمام صورتاں نال مطابقت۔
ہن منّو $\Th{PA}$، $\OCon[\Th{PA}]$ نوں !!{derive}s کردا
اے۔ فیر اوہ $\lnot
\OProv[\Th{PA}](\gn{!G_\Th{PA}})$ نوں وی !!{derive}s کردا اے۔
پر کیوں جے $!G_\Th{PA}$ گوڈل جملہ اے،ایہہ
$!G_\Th{PA}$ نال ہم ارز اے۔ ایس لئی
$\Th{PA}$، $!G_\Th{PA}$ نوں !!{derive}s کردا اے۔

پر اسیں جاندے آں کہ جے $\Th{PA}$ بے تضاد ہووے تاں اوہ
$!G_\Th{PA}$ نوں !!{derive} نہیں کردا!{} سو جے $\Th{PA}$
بے تضاد اے تاں اوہ $\OCon[\Th{PA}]$ نوں !!{derive} نہیں
کر سکدا۔

دلیل نوں ہور دقیق بناؤن لئی $!G_\Th{PA}$ نوں
$\Th{PA}$ لئی گوڈل جملہ منّاں گے،تے !!{derivability} شرطاں
(P1)--(P3) نال دکھاواں گے کہ $\Th{PA}$،
$\OCon[\Th{PA}] \lif
!G_\Th{PA}$ نوں !!{derive}s کردا اے۔ ایس توں پتا لگے گا کہ
$\Th{PA}$، $\OCon[\Th{PA}]$ نوں !!{derive} نہیں کردا۔
ہِٹھاں $\Th{PA}$ دے اندر ایس ثبوت دا خاکہ اے۔ (سہولت لئی
اسیں $\Th{PA}$ دے ذیلی اشاریے چھڈ دیندے آں۔)
\begin{align}
& \formula G \liff \lnot \OProv(\gn{!G}) \ollabel{G2-1}\\
& \qquad\text{$!G$، گوڈل دا جملہ اے}\notag \\
& \formula G \lif \lnot \OProv(\gn{!G}) \ollabel{G2-2}\\
  & \qquad\text{\olref{G2-1} توں} \notag\\
& \formula G \lif
  (\OProv(\gn{!G}) \lif \lfalse) \ollabel{G2-3}\\
  & \qquad\text{\olref{G2-2} توں منطق نال}\notag\\
& \OProv(\gn{
    \formula G \lif
    (\OProv(\gn{!G}) \lif \lfalse)
  }) \ollabel{G2-4}\\
  & \qquad\text{\olref{G2-3} توں شرط P1 نال} \notag\\
& \OProv(\gn{!G}) \lif
  \OProv(\gn{
    (\OProv(\gn{!G}) \lif \lfalse)
    }) \ollabel{G2-5}\\
  & \qquad\text{\olref{G2-4} توں شرط P2 نال} \notag\\
& \OProv(\gn{!G}) \lif (\OProv(\gn{\OProv(\gn{!G})}) \lif \OProv(\gn{\lfalse})) \ollabel{G2-6}\\
  & \qquad\text{\olref{G2-5} توں شرط P2 تے منطق نال} \notag\\
& \OProv(\gn{!G}) \lif 
  \OProv(\gn{\OProv(\gn{!G})}) \ollabel{G2-7}\\
   & \qquad\text{شرط P3 نال} \notag\\
& \OProv(\gn{!G}) \lif \OProv(\gn{\lfalse}) \ollabel{G2-8}\\
  & \qquad \text{\olref{G2-6} تے \olref{G2-7} توں منطق نال}\notag\\
& \OCon \lif \lnot \OProv(\gn{!G}) \ollabel{G2-9}\\
  & \qquad\text{\olref{G2-8} دے عکسِ نقیض تے $\OCon \ident \lnot \OProv(\gn{\lfalse})$ نال}\notag \\
& \OCon \lif !G \notag\\
  & \qquad\text{\olref{G2-1} تے \olref{G2-9} توں منطق نال}\notag
\end{align}
اُتلی دلیل وچ منطق دی ورتوں صرف قضیاتی منطق دے بنیادی حقائق
نیں۔ مثال لئی \olref{G2-3}،
$\Proves \lnot!A \liff (!A\lif
\lfalse)$ ورتدا اے،تے \olref{G2-8}،
$!A \lif (!B \lif !C), !A \lif !B
\Proves !A \lif !C$ ورتدا اے۔ \olref{G2-5} تے
\olref{G2-6} وچ شرط~P2 دی ورتوں P2 دیاں انہاں مثالاں اُتے
منحصر اے:
$\OProv(\gn{!A \lif !B}) \lif
(\OProv(\gn{!A}) \lif \OProv(\gn{!B}))$۔ پہلی مثال وچ
$!A \ident
!G$ تے $!B \ident \OProv(\gn{!G}) \lif \lfalse$؛
% PNB-CMP-071: B11 آخری P2 instance وچ G دے formula marker ! بحال۔
دوجی وچ $!A
\ident \OProv(\gn{!G})$ تے $!B \ident \lfalse$۔

نامکملیت دے دوجے قضیے دا ہور عمومی روپ ایہ اے:

\begin{thm}
\ollabel{thm:second-incompleteness-gen} منّو $\Th{T}$ کوئی وی
بے تضاد، !!{axiomatized} نظریہ اے جیہڑا $\Th{Q}$ دی توسیع
کردا اے؛تے $\OProv[\Th{T}](y)$ کوئی وی ایسا فارمولا ہووے
جیہڑا $\Th{T}$ لئی !!{derivability} شرطاں P1--P3 پوری کردا
اے۔ فیر $\Th{T}$، $\OCon[T]$ نوں !!{derive} نہیں کردا۔
\end{thm}

\begin{prob}
دکھاؤ کہ $\Th{PA}$، $!G_{\Th{PA}} \lif \OCon[\Th{PA}]$
نوں !!{derive}s کردا اے۔
\end{prob}

\begin{digress}
ایس گل دا حاصل ایہہ اے کہ ریاضی دا کوئی وی «معقول» بے تضاد
نظریہ اپنی بے تضادی دا جملہ !!{derive} نہیں کر سکدا۔ منّو
$\Th{T}$ ریاضی دا نظریہ اے جیہڑے وچ $\Th{Q}$ تے ہلبرٹ
دی «محدود وسیلی» دلیلاں (اوہناں توں جو وی مراد ہووے) شامل نیں۔
تاں پورا $\Th{T}$، $\Th{T}$ دی بے تضادی دا جملہ
!!{derive} نہیں کر سکدا؛سو لازماً اوہدا محدود
وسیلی حصہ وی $\Th{T}$ دی بے تضادی دا جملہ !!{derive} نہیں
کر سکدا۔ ایس معنی وچ «پوری ریاضی» لئی محدود وسیلی بے تضادی
دا ثبوت نہیں ہو سکدا۔

«محدود وسیلی» لفظ دی تعبیر وچ کجھ گنجائش اے،تے گوڈل
1931 دے مقالے وچ ایہہ امکان منّدا اے کہ جیہڑی گل نوں اسیں
«محدود وسیلی» سمجھدے آں اوہ ریاضی دیاں انہاں قسماں توں باہر
ہووے جنہاں نوں ہلبرٹ رسمی بنانا چاہندا سی۔ پر ایتھے گوڈل
نے کھُل دلی دکھائی؛اج ایہہ سمجھنا اوکھا اے کہ کوئی شے معقول
طور تے «محدود وسیلی» وی کہلائے تے،مثلاً $\Th{ZFC}$،یعنی
انتخاب دے مسلّمے سمیت زرمیلو--فرینکل سیٹ نظریے،وچ رسمی طور
بیان وی نہ کیتی جا سکے۔
\end{digress}
\endgroup
% /OLP-0319

% OLP-0320 content/incompleteness/incompleteness-provability/lob-thm.tex
\begingroup
\def\olfilename{lob-thm}
\expandafter\def\csname ol@sectioncs\endcsname{section}

\olfileid{inc}{inp}{lob}

\olsection{لوب دا قضیہ}

نظریے~$\Th{T}$ دا گوڈل جملہ $\lnot
\OProv[\Th{T}](y)$ دا مقرر نقطہ اے،یعنی ایسا !!a{sentence}~$!G$
جیہڑا
\[
\Th{T} \Proves \lnot \OProv[\Th{T}](\gn{!G}) \liff !G.
\]
پورا کردا اے۔ جے نظریہ بے تضاد ہووے تاں ایہہ !!{derivable}
نہیں: جے $\Th{T} \Proves !G$ ہووے،تاں (a) !!{derivability}
شرط~(1) توں $\Th{T} \Proves \OProv[\Th{T}](\gn{!G})$،
تے (b) $\Th{T}
\Proves !G$ نال $\Th{T} \Proves \lnot \OProv[\Th{T}](\gn{!G})
\liff !G$ مل کے $\Th{T} \Proves \lnot
\OProv[\Th{T}](\gn{!G})$ دیندے نیں؛فیر $\Th{T}$
متضاد ہووے گا۔ ہن قدرتی سوال ایہہ اے کہ
$\OProv[\Th{T}](y)$ دے مقرر نقطے دی کی حیثیت اے،یعنی
ایسا !!a{sentence}~$!H$ جیہڑا
\[
\Th{T} \Proves \OProv[\Th{T}](\gn{!H}) \liff !H.
\]
پورا کردا اے۔ جے ایہہ !!{derivable} ہووے،تاں شرط~(1)
توں $\Th{T} \Proves \OProv[\Th{T}](\gn{!H})$؛پر اُتلے
ہم ارزی والے بیان اُتے قیاسِ استثنائی لائیے تاں وی ایہو نتیجہ
ملدا اے۔ ایس لئی گوڈل جملے والی دلیل وانگ اتھوں
$\Th{T}$ دا تضاد نہیں نکلدا،کم از کم اوسے دلیل نال تاں نہیں۔
البتہ ایس توں ایہہ وی ثابت نہیں ہوندا کہ $\Th{T}$
\emph{واقعی} $!H$ نوں !!{derive} کردا اے۔

اس سوال وچ پیش رفت لئی اس نوں کجھ عمومی روپ وچ ویکھدے آں۔
مقرر نقطے دی ہم ارزی دا کھبے توں سجے ول رخ،
$\OProv[\Th{T}](\gn{!H}) \lif !H$،اک عمومی خاکے
$\OProv[\Th{T}](\gn{!A}) \lif !A$ دی مثال اے جیہنوں
\emph{انعکاس دا اصول} کہندے نیں۔ ایہہ ناں ایس لئی اے کہ،
اک معنی وچ،ایہہ $\Th{T}$ دے اپنے !!{derive} کیتے ہوئیاں
جملیاں بارے غور نوں ظاہر کردا اے؛بنیادی گل ایہہ اے:
«جے $\Th{T}$، $!A$ نوں !!{derive} کر سکدا اے،تاں
$!A$ سچ اے»،ہر~$!A$ لئی۔ ایہ گل صرف اوہناں نظریاں لئی سچ اے جنہاں دے
ثابت کیتے سارے جملے سچے ہون۔ ایس توں اندازہ ہوندا اے کہ نظریے عام طور اُتے اس
اصول دی ہر مثال نوں !!{derive} نہیں کرن گے۔ تاں کِہڑیاں
مثالاں کوئی کافی قوی نظریہ،جیہڑا !!{derivability} شرطاں وی
پوری کردا ہووے،!!{derive} کر سکدا اے؟ یقیناً اوہ ساریاں
جنہاں وچ $!A$ خود !!{derivable} اے۔ اگلا نتیجہ دسدَا
اے کہ صرف ایہو مثالاں نیں۔

\begin{thm}\ollabel{thm:lob}
منّو $\Th{T}$، $\Th{Q}$ دی توسیع کرن والا کوئی
!!a{axiomatizable} نظریہ اے،تے
$\OProv[\Th{T}](y)$ اک ایسا فارمولا اے جیہڑا
\olref[2in]{sec} دیاں P1--P3 شرطاں پوری کردا اے۔ جے
$\Th{T}$، $\OProv[\Th{T}](\gn{!A}) \lif !A$ نوں
!!{derive}s کردا اے،تاں واقعی $\Th{T}$، $!A$ نوں
!!{derive}s کردا اے۔
\end{thm}

دوجے لفظاں وچ،جے $\Th{T} \Proves/ !A$،تاں
$\Th{T} \Proves/
\OProv[\Th{T}](\gn{!A}) \lif !A$۔ ایس نتیجے نوں
لوب دا قضیہ کہندے نیں۔

\begin{explain}
لوب دے قضیے دے ثبوت دا خیال سانتا کلاز دے موجود ہون دی
اک چترائی بھری دلیل توں ملدا اے۔ (جے ایہہ نتیجہ پسند نہ
آوے،تاں اپنی پسند دا کوئی ہور نتیجہ رکھ لو۔) دلیل ایہہ اے:
\begin{enumerate}
\item $X$ اوہ جملہ ہووے: «جے $X$ سچ اے،تاں سانتا کلاز
  موجود اے»۔
\item منّو $X$ سچ اے۔
\item فیر اس جملے دی کہی گل درست اے؛یعنی جے $X$ سچ اے
  تاں سانتا کلاز موجود اے۔
\item چونکہ اسیں $X$ نوں سچ منّیا اے،(2) تے~(3) توں
  قیاسِ استثنائی نال ایہہ نتیجہ کڈھ سکدے آں کہ سانتا کلاز
  موجود اے۔
\item اسیں فرض~(2)،«$X$ سچ اے»،توں (4)،«سانتا کلاز
  موجود اے»،اخذ کر لیا۔ شرطیہ ثبوت نال دکھایا: «جے $X$
  سچ اے،تاں سانتا کلاز موجود اے»۔
\item پر ایہہ تاں عین جملہ~$X$ اے۔ ایس لئی اسیں دکھا
  دتا کہ $X$ سچ اے۔
\item فیر اُتلی (2)--(4) والی دلیل نال سانتا کلاز موجود اے۔
\end{enumerate}
ایس خیال نوں رسمی بنا کے «سچ اے» دی تھاں «!!{derivable}
اے» تے «سانتا کلاز موجود اے» دی تھاں~$!A$ رکھیے،تاں
لوب دے قضیے دا ثبوت ملدا اے۔ چال ایہہ اے کہ مقرر نقطے دا
لمہ !!{formula}~$\OProv[\Th{T}](y) \lif !A$ اُتے
لائیے۔ ایس دا مقرر نقطہ پچھلے خاکے دے جملے~$X$ دے
مطابق اے۔
\end{explain}

\begin{proof}[\olref{thm:lob} دا ثبوت]
منّو $!A$ ایسا !!a{sentence} اے جیہڑے لئی $\Th{T}$،
$\OProv[\Th{T}](\gn{!A}) \lif !A$ نوں !!{derive}s
کردا اے۔ $!B(y)$ !!{formula}~$\OProv[\Th{T}](y)
\lif !A$ ہووے؛مقرر نقطے دے لمے نال ایسا !!a{sentence}~$!D$
لبھو جیہڑے لئی $\Th{T}$، $!D \liff !B(\gn{!D})$ نوں
!!{derive}s کردا اے۔ فیر ہِٹھلا ہر بیان
$\Th{T}$ وچ !!{derivable} اے:
\begin{align}
  & \formula D \liff (\OProv[\Th{T}](\gn{!D}) \lif !A) \ollabel{L-1}\\
  & \qquad \text{$!D$، $!B(y)$ دا مقرر نقطہ اے}\notag \\
  & \formula D \lif (\OProv[\Th{T}](\gn{!D}) \lif !A) \ollabel{L-2}\\
  & \qquad\text{\olref{L-1} توں}\notag\\
  & \OProv[\Th{T}](\gn{!D \lif (\OProv[\Th{T}](\gn{!D}) \lif !A)}) \ollabel{L-3}\\
  & \qquad \text{\olref{L-2} توں شرط P1 نال}\notag \\
  & \OProv[\Th{T}](\gn{!D}) \lif \OProv[\Th{T}](\gn{\OProv[\Th{T}](\gn{!D}) \lif !A})
  \ollabel{L-4}\\
  &\qquad \text{\olref{L-3} توں شرط P2 نال}\notag \\
  & \OProv[\Th{T}](\gn{!D}) \lif (\OProv[\Th{T}](\gn{\OProv[\Th{T}](\gn{!D})}) \lif \OProv[\Th{T}](\gn{!A})) \ollabel{L-5}\\
  &\qquad \text{\olref{L-4} توں دوبارہ P2 نال} \notag\\
& \OProv[\Th{T}](\gn{!D}) \lif \OProv[\Th{T}](\gn{\OProv[\Th{T}](\gn{!D})}) \ollabel{L-6}\\
  & \qquad\text{!!{derivability} شرط P3 نال} \notag\\
  & \OProv[\Th{T}](\gn{!D}) \lif \OProv[\Th{T}](\gn{!A}) \ollabel{L-7} \\
  &\qquad\text{\olref{L-5} تے \olref{L-6} توں}\notag\\
  & \OProv[\Th{T}](\gn{!A}) \lif !A \ollabel{L-8}\\
  &\qquad\text{قضیے دے مفروضے توں} \notag\\
  & \OProv[\Th{T}](\gn{!D}) \lif !A \ollabel{L-9}\\
  &\qquad\text{\olref{L-7} تے \olref{L-8} توں}\notag\\
  & (\OProv[\Th{T}](\gn{!D}) \lif !A) \lif !D \ollabel{L-10}\\
  & \qquad \text{\olref{L-1} توں}\notag \\
  & \formula D \ollabel{L-11}\\
  & \qquad\text{\olref{L-9} تے \olref{L-10} توں}\notag \\
  & \OProv[\Th{T}](\gn{!D}) \ollabel{L-12}\\
  & \qquad\text{\olref{L-11} توں شرط~P1 نال}\notag \\
  & \formula A \qquad\qquad\text{\olref{L-8} تے \olref{L-12} توں}\notag
\end{align}
\end{proof}

لوب دے قضیے توں نامکملیت دے دوجے قضیے دا اک مختصر ثبوت
وی ملدا اے،اوہناں نظریاں لئی جنہاں دا !!a{derivability}
محمول P1--P3 شرطاں پوری کردا ہووے: جے $\Th{T} \Proves
\OProv[\Th{T}](\gn{\lfalse}) \lif \lfalse$،تاں
$\Th{T} \Proves \lfalse$۔ جے $\Th{T}$ بے تضاد اے،
تاں $\Th{T} \Proves/ \lfalse$۔ ایس لئی
$\Th{T}
\Proves/ \OProv[\Th{T}](\gn{\lfalse}) \lif \lfalse$،
یعنی $\Th{T} \Proves/
\OCon[\Th{T}]$۔ اسی قضیے نال ایہہ وی دکھا سکدے آں کہ
$!H$،جیہڑا $\OProv[\Th{T}](x)$ دا مقرر نقطہ اے،
!!{derivable} اے۔ کیونکہ
\begin{align*}
  \Th{T} & \Proves \OProv[\Th{T}](\gn{!H}) \liff !H\\
  \intertext{خاص طور اُتے}
    \Th{T} & \Proves \OProv[\Th{T}](\gn{!H}) \lif !H
\end{align*}
تے ایس لئی لوب دے قضیے نال $\Th{T} \Proves !H$۔

% Going in the other direction, for homework I
% may ask you to work through a short proof of L\"ob's theorem, using
% the second incompleteness theorem instead of the fixed-point lemma.

\begin{prob}
منّو $\Th{T}$ اک قابلِ حساب طریقے نال مسلّمہ بند نظریہ
اے،تے $\OProv[\Th{T}]$، $\Th{T}$ دا
!!a{derivability} محمول اے۔ ہِٹھلے چار بیان ویکھو:
\begin{enumerate}
\item جے $T \Proves !A$،تاں $T \Proves \OProv[\Th{T}](\gn{!A})$۔
\item $T \Proves !A \lif \OProv[\Th{T}](\gn{!A})$۔
\item جے $T \Proves \OProv[\Th{T}](\gn{!A})$،تاں $T \Proves !A$۔
\item $T \Proves \OProv[\Th{T}](\gn{!A}) \lif !A$
\end{enumerate}
ہر اک بیان کنہاں شرطاں وچ درست اے؟
\end{prob}
\endgroup
% /OLP-0320

% OLP-0321 content/incompleteness/incompleteness-provability/tarski-thm.tex
\begingroup
\def\olfilename{tarski-thm}
\expandafter\def\csname ol@sectioncs\endcsname{section}

\olfileid{inc}{inp}{tar}

\olsection{سچائی دی ناقابلِ تعریفیت}

\emph{قابلِ تعریف ہون} دا تصور حساب دی زبان لئی رسمی معنیات
ہون اُتے منحصر اے۔ اسیں حساب دی زبان دے فارمولیاں تے جملیاں
دا سیٹ بیان کیتا اے۔ «مرادی تعبیر» وچ ایہہ جملے قدرتی
اعداد بارے دعوے کردے نیں،تے ایسا دعویٰ سچ یا جھوٹ ہو سکدا
اے۔ $\Struct{N}$ اوہ !!{structure} ہووے جیہدا دائرہ
$\Nat$ اے تے حساب دی زبان دیاں علامتاں دی معیاری تعبیر
ورتدا اے۔ فیر $\Sat{N}{!A}$ دا مطلب اے: «$!A$ معیاری
تعبیر وچ سچ اے»۔

\begin{defn}
قدرتی اعداد دا کوئی تعلق $R(x_1,\dots,x_k)$، $\Struct{N}$
وچ اُدوں تے صرف اُدوں \emph{قابلِ تعریف} اے جدوں حساب دی زبان وچ
ایسا فارمولا $!A(x_1,\dots,x_k)$ ہووے کہ ہر
$n_1,\dots,n_k$ لئی $R(n_1,\dots,n_k)$ اُس ویلے تے
صرف اوس ویلے درست ہووے جدوں $\Sat{N}{!A(\num n_1,\dots,\num
  n_k)}$۔
\end{defn}

دوجے لفظاں وچ، کوئی تعلق $\Struct{N}$ وچ اُدوں تے صرف اُدوں
قابلِ تعریف اے جدوں اوہ نظریے $\Th{TA}$ وچ قابلِ نمائندگی
ہووے؛ایتھے $\Th{TA} =
\Setabs{!A}{\Sat{N}{!A}}$ حساب دے سچ جملیاں دا سیٹ
اے۔ (جے ایہہ گل فوراً واضح نہ ہووے،تاں تعریفاں ول
مڑ کے پڑتال کرو تے اپنے آپ نوں اس دی وجہ سمجھاؤ۔)

\begin{lem}
ہر قابلِ حساب تعلق~$\Struct{N}$ وچ قابلِ تعریف اے۔
\end{lem}

\begin{proof}
آسانی نال پڑتال ہو جاندی اے کہ $\Th{Q}$ وچ کسے تعلق
دی نمائندگی کرن والا فارمولا، $\Struct{N}$ وچ اوہی تعلق
متعین کردا اے۔
\end{proof}

ہن پُچھ سکدے آں کہ کی الٹ وی درست اے؟ یعنی
$\Struct{N}$ وچ قابلِ تعریف ہر تعلق قابلِ حساب اے؟
جواب نہیں۔ مثال لئی:

\begin{lem}
حساب رُکّن دا تعلق $\Struct{N}$ وچ قابلِ تعریف اے۔
\end{lem}

\begin{proof}
یاد کرو کلینی دے نارمل روپ دا قضیہ کہندا اے کہ ہر جزوی
قابلِ حساب فنکشن~$f$ دا کوئی اشاریہ~$e$ اے جیہڑے نال
$f(x) =
U(\umin{s}{T(e,x,s)})$ ہر $x \in \Nat$ لئی؛ایتھے
$U$ تے $T$ بنیادی بازگشتی
تے ایس لئی کُل نیں۔ سو $f(x)$ اُدوں تے صرف اُدوں متعین اے، یعنی حساب اُدوں ہی رُکدا
اے، جدوں کوئی ایہو جیہا~$s$ ہووے جیہڑے لئی
$T(e,x,s)$ درست ہووے۔

ہن $H$ حساب رُکّن دا تعلق ہووے،یعنی
\[
H = \Setabs{\tuple{e,x}}{\lexists[s][T(e, x, s)]}.
\]
$!D_T$، $T$ دی $\Struct{N}$ وچ تعریف کردا ہووے۔ فیر
\[
H = \Setabs{\tuple{e,x}}{\Sat{N}{\lexists[s][!D_T(\num e, \num x, s)]}},
\]
ایس لئی $\lexists[s][!D_T(z, x, s)]$،
~$H$ نوں $\Struct{N}$ وچ تعریف کردا اے۔
\end{proof}

\begin{prob}
دکھاؤ کہ $Q(n) \defiff n \in \Setabs{\Gn{!A}}{\Th{Q} \Proves !A}$
حساب وچ قابلِ تعریف اے۔
\end{prob}

خود $\Th{TA}$ دا کیہہ؟ کی اوہ حساب وچ قابلِ تعریف اے؟
یعنی کی سیٹ $\Setabs{\Gn{!A}}{\Sat{N}{!A}}$ حساب
وچ قابلِ تعریف اے؟ تارسکی دے قضیے دا جواب ناں وچ اے۔

\begin{thm}
\ollabel{thm:tarski}
حساب دے سچ !!{sentence}s دا سیٹ حساب وچ قابلِ تعریف
نہیں۔
\end{thm}

\begin{proof}
منّو $!D(x)$ اُس دی تعریف کردا اے؛یعنی
$\Sat{N}{!A}$ اُدوں تے صرف اُدوں جدوں
$\Sat{N}{!D(\gn{!A})}$۔ مقرر نقطے دے لمے نال کوئی
فارمولا $!A$ ملدا اے جیہڑے لئی
$\Th{Q} \Proves !A \liff \lnot !D(\gn{!A})$،
تے ایس توں $\Sat{N}{!A \liff \lnot !D(\gn{!A})}$۔
پر فیر $\Sat{N}{!A}$ اُدوں تے صرف اُدوں جدوں
$\Sat{N}{\lnot !D(\gn{!A})}$،جو ایس مفروضے دے
خلاف اے کہ $!D(y)$ حساب دے سچ بیانواں دا سیٹ تعریف
کردا اے۔
\end{proof}

تارسکی نے ایہہ تجزیہ سچائی دے ہور عمومی فلسفیانہ
تصور اُتے وی لاگو کیتا۔ کوئی وی زبان $L$ منّو؛تارسکی
دے مطابق $L$ لئی سچائی دا مناسب تصور ہر جملے $X$ واسطے
ایہہ شرط پوری کرے گا:
\begin{quote}
«$X$» اُدوں تے صرف اُدوں سچ اے جدوں $X$۔
\end{quote}
انگریزی وچ تارسکی دی بار بار نقل ہون والی مثال
“Snow is white” اے؛پنجابی وچ ایہہ شرط ایویں پڑھیے:
\begin{quote}
«برف چٹی اے» اُدوں تے صرف اُدوں سچ اے جدوں برف چٹی اے۔
\end{quote}
پر ہر ایسی زبان وچ جیہڑی قطری فنکشن دی نمائندگی لئی کافی
قوی ہووے،تے ہر لسانی محمول $T(x)$ لئی،اسیں ایسا جملہ
$X$ بنا سکدے آں جیہڑا ایہہ آکھے: «$X$ اُدوں تے صرف اُدوں جدوں
$T(\text{«$X$»})$ نہیں»۔ چونکہ اسیں نہیں چاہندے کہ
سچائی دا محمول کسے جملے نوں سچ تے جھوٹ دونوں قرار دے،
تارسکی نے نتیجہ کڈھیا کہ ایسی زبان دے سارے جملیاں لئی
سچائی دا محمول متعین کرن لئی کسے نہ کسے طرح اوس زبان
دیاں حداں توں باہر جانا پئے گا۔ دوجے لفظاں وچ،ایسی
زبان دا سچائی محمول خود اوسے زبان وچ قابلِ تعریف نہیں۔
\endgroup
% /OLP-0321

\OLEndChapterHook
\endgroup
% /OLP-0312

\OLEndPartHook
\endgroup
% /OLP-0274

% OLP-0322 content/second-order-logic/second-order-logic.tex
\begingroup
\def\olfilename{second-order-logic}
\expandafter\def\csname ol@sectioncs\endcsname{section}

\olpart{sol}{دوجے درجے دی منطق}

\begin{editorial}
دوجے درجے دی منطق والا حصہ ہجے مُڈھ وچ اے۔
\end{editorial}

% OLP-0323 content/second-order-logic/syntax-and-semantics/syntax-and-semantics.tex
\begingroup
\def\olfilename{syntax-and-semantics}
\expandafter\def\csname ol@sectioncs\endcsname{section}

\olchapter{sol}{syn}{نحو تے معنیات}

\begin{editorial}
ہُن تک دوجے درجے دی منطق دی بنیادی نحو تے معنیات بیان
کیتیاں نیں۔ باب ہون دے لحاظ نال ایہ ہجے بہت چھوٹا اے۔
دوجے درجے دی منطق دے !!{derivation} نظاماں بارے گل کرن لئی
پہلاں ایہ دسنا پئے گا کہ دوجے درجے دے متغیراں دی تھاں ہور
عبارت کیویں رکھی جاندی اے۔ ایس وچ کجھ باریک مسئلے وی نیں۔
\end{editorial}

% OLP-0324 content/second-order-logic/syntax-and-semantics/introduction.tex
\begingroup
\def\olfilename{introduction}
\expandafter\def\csname ol@sectioncs\endcsname{section}

\olfileid{sol}{syn}{int}

\olsection{تعارف}

پہلے درجے دی منطق وچ اسیں کسے زبان دیاں غیر منطقی
علامتاں، یعنی اوہدے !!{constant}s، !!{function}s تے
!!{predicate}s، نوں منطقی علامتاں نال ملا کے پہلے درجے
دیاں !!{structure}s بارے گل ظاہر کردے آں۔ ایس کم لئی پورا اترن دا تصور ورتیدا اے۔ ایہ !!a{structure}~$\Struct{M}$
تے متغیراں دی قدر بندی~$s$ نوں !!a{formula}~$!A$ نال جوڑدا اے۔
$\Sat{M}{!A}[s]$ اُدوں تے صرف اُدوں لاگو اے جدوں $!A$ دی کہی گل
سچ ہووے، جدوں اوہدے
!!{constant}s، !!{function}s تے !!{predicate}s دی
تعبیر $\Struct{M}$ دے مطابق تے آزاد متغیراں دی
تعبیر~$s$ دے مطابق ہووے۔ !!{identity}~$\eq$ دی
تعبیر $\Sat{M}{!A}[s]$ دی تعریف وچ پہلے توں شامل
اے،جیویں~$\lforall$ تے $\lexists$ دی تعبیر وی۔
پہلی علامت نوں ہمیشہ ساخت دے !!{domain}~$\Domain{M}$
اُتے برابری دے تعلق دے طور تے سمجھیا جاندا اے، تے مسوّراں دا دائرہ
پورا !!{domain} ہوندا اے۔ پر خاص گل ایہ اے کہ مسوّریت صرف
!!{domain} دے عناصر اُتے ہو سکدی اے؛ایس لئی مسوّر
دے بعد صرف شے دے !!{variable}s آ سکدے نیں۔

دوجے درجے دی منطق وچ زبان تے پورا اترن دی تعریف
دونوں نوں وسیع کیتا جاندا اے،تاکہ آزاد تے پابند
فنکشنی تے محمولی متغیر،تے اوہناں اُتے مسوّریت،شامل
ہوون۔ ایہہ متغیر !!{function}s تے !!{predicate}s نال اوہو تعلق رکھدے نیں
جیہڑا شے دے متغیر !!{constant}s نال رکھدے نیں۔ دوجے درجے دی
منطق دے ٹرم تے
!!{formula}s بناؤن وچ ایہہ وی اوہی کردار نبھاؤندے
نیں،تے اوہناں اُتے مسوّریت وی ملدی جلُدی طرح
نبھائی جاندی اے۔ \emph{معیاری} معنیات وچ دوجے درجے
دے مسوّر درست نوع دیاں ساریاں ممکنہ چیزاں اُتے
چلدے نیں۔ فنکشنی متغیراں لئی ایہ سارے $n$-رکنی فنکشن نیں،
جیہڑے $\Domain{M}$ دے ترتیب وار ٹپل نوں~$\Domain{M}$ دا اک
عنصر دیندے نیں۔ محمولی متغیراں لئی ایہ سارے $n$-رکنی تعلق نیں۔ مثال لئی
$\lforall[\Obj{v_0}][(\Obj{P^1_0}(\Obj{v_0}) \lor \lnot
  \Obj{P^1_0}(\Obj{v_0}))]$ پہلے تے دوجے دونوں درجیاں
دی منطق دا فارمولا اے؛پر دوجے درجے وچ اسیں
$\lforall[\Obj{V^1_0}][\lforall[\Obj{v_0}][(\Obj{V^1_0}(\Obj{v_0}) \lor
    \lnot \Obj{V^1_0}(\Obj{v_0}))]]$ تے
$\lexists[\Obj{V^1_0}][\lforall[\Obj{v_0}][(\Obj{V^1_0}(\Obj{v_0}) \lor
    \lnot \Obj{V^1_0}(\Obj{v_0}))]]$ وی ویکھ سکدے آں۔
کیوں جے انہاں وچ کوئی آزاد متغیر نہیں، ایس لئی ایہ دوجے درجے
دی منطق دے !!{sentence}s نیں۔ ایتھے $\Obj{V^1_0}$
دوجے درجے دا $1$-رکنی محمولی متغیر اے۔ $\Obj{V^1_0}$
دیاں جائز تعبیراں اوہو نیں جیہڑیاں اسیں کسے
$1$-رکنی !!{predicate}،مثلاً $\Obj{P^1_0}$،نوں دے
سکدے آں،یعنی
~$\Domain{M}$ دے ذیلی سیٹ۔ ایہناں ذیلی سیٹاں اُتے مسوّریت دا مطلب ایہ اے کہ $\lforall[\Obj{v_0}][(\Obj{V^1_0}(\Obj v_0) \lor \lnot
  \Obj{V^1_0}(v_0))]$، $\Domain{M}$ دا ذیلی سیٹ
$\Obj{V^1_0}$ دی قدر دے طور تے دین دے
ہر طریقے لئی، یا گھٹ توں گھٹ اک طریقے لئی، درست ہووے۔
ہر دِتّی شے کسے سیٹ وچ یا شامل ہوندی اے یا شامل نہیں ہوندی۔
ایس لئی دوویں جملے ہر
!!{structure} وچ سچ نیں۔
\endgroup
% /OLP-0324

% OLP-0325 content/second-order-logic/syntax-and-semantics/terms-formulas.tex
\begingroup
\def\olfilename{terms-formulas}
\expandafter\def\csname ol@sectioncs\endcsname{section}

\olfileid{sol}{syn}{frm}

\olsection{ٹرم تے \printtoken{P}{formula}}

پہلے درجے دی منطق وانگ، دوجے درجے دی منطق دیاں لکھتاں
بنیادی علامتاں دی اک فہرست توں بنائیاں جاندیاں نیں۔ ایس وچ
\emph{!!{variable}s}، \emph{!!{constant}s}،
\emph{!!{predicate}s} تے کدے کدے
\emph{!!{function}s} شامل نیں۔ انہاں نوں منطقی رابطیاں،
مسوّراں، تے قوساں تے کومیاں ورگیاں لکھت دیاں نشانیاں
نال ملا کے \emph{ٹرم} تے \emph{!!{formula}s}
بنائے جاندے نیں۔ فرق ایہہ اے کہ دوجے درجے دی منطق وچ
شے دے متغیراں توں علاوہ تعلقاں تے فنکشوناں دے
متغیر وی ہوندے نیں، تے انہاں اُتے مسوّریت دی
اجازت وی ہوندی اے۔ سو دوجے درجے دی منطق دیاں منطقی
علامتاں پہلے درجے دی منطق والیاں نیں، نال ایہہ وی:

\begin{enumerate}
\item ہر رُکن-گنتی~$n$ لئی دوجے درجے دے تعلقی !!{variable}s دا
  اک !!{denumerable}s سیٹ:
  $\Obj V_0^n$, $\Obj V_1^n$, $\Obj V_2^n$, \dots
\item دوجے درجے دے فنکشنی !!{variable}s دا اک
  !!{denumerable}s سیٹ:
  $\Obj u_0^n$, $\Obj u_1^n$, $\Obj u_2^n$, \dots
\end{enumerate}

جیویں اسیں $x$, $y$, $z$ نوں پہلے درجے دے
متغیراں $\Obj v_i$ لئی ماورا متغیراں دے طور تے ورتدے آں،
اُویں $X$, $Y$, $Z$ وغیرہ نوں $\Obj V_i^n$ لئی
تے $u$, $v$ وغیرہ نوں~$\Obj u_i^n$ لئی ماورا متغیراں
دے طور تے ورتاں گے۔

\begin{explain}
دوجے درجے دی زبان دیاں غیر منطقی علامتاں اوہو
طریقے نال متعین کیتیاں جاندیاں نیں جیویں پہلے درجے
دی زبان دیاں: اوہدے !!{constant}s، !!{function}s
تے !!{predicate}s دی فہرست دے کے۔

پہلے درجے دی منطق وچ !!{identity}~$\eq$ نوں عموماً
شامل کیتا جاندا اے۔ پہلے درجے دی منطق وچ
زبان~$\Lang{L}$ دیاں غیر منطقی علامتاں کسے وی
دلچسپ گل نوں بیان کرن لئی بہت ضروری نیں۔ بے شک
کجھ !!{sentence}s غیر منطقی علامتاں توں بناں وی
ہوندے نیں، پر صرف~$\eq$ نال کوئی دلچسپ گل کہنا
اوکھا اے۔ دوجے درجے دی منطق وچ، تعلقی تے فنکشنی متغیر بے حد نیں۔ ایس لئی اسیں
وکھریاں غیر منطقی علامتاں توں بناں وی اوہ کجھ آکھ سکدے آں
جیہڑا پہلے درجے دی کسے زبان وچ آکھ سکدے آں۔
\end{explain}

\begin{defn}[دوجے درجے دے ٹرم]
زبان~$\Lang L$ دے \emph{دوجے درجے دے ٹرماں} دا سیٹ،
$\TrmSOL[L]$، \olref[fol][syn][frm]{defn:terms}
دی تعریف وچ ایہہ شق شامل کر کے متعین ہوندا اے:
\begin{enumerate}
\item جے $u$ اک $n$-رکنی فنکشنی متغیر اے تے $t_1$, \dots, $t_n$
  ٹرم نیں، تاں $\Atom{u}{t_1, \ldots, t_n}$ اک ٹرم اے۔
\end{enumerate}
\end{defn}

\begin{explain}
سو دوجے درجے دا ٹرم پہلے درجے دے ٹرم ورگا اے۔
فرق ایہ اے کہ جتھے پہلے درجے دے ٹرم
وچ !!a{function}~$\Obj{f^n_i}$ آؤندا اے، اُتھے
دوجے درجے دے ٹرم وچ اُس دی تھاں فنکشنی
متغیر~$\Obj{u^n_i}$ آ سکدا اے۔
\end{explain}

\begin{defn}[دوجے درجے دے \usetoken{s}{formula}]
\emph{دوجے درجے دے !!{formula}s} دا سیٹ~$\FrmSOL[L]$،
زبان~$\Lang L$ لئی،
\olref[fol][syn][frm]{defn:formulas} دی تعریف وچ
ایہہ شقّاں شامل کر کے متعین ہوندا اے:
\begin{enumerate}
\item جے $X$ اک $n$-رکنی محمولی متغیر اے تے $t_1$, \dots,
  $t_n$، زبان~$\Lang L$ دے دوجے درجے دے ٹرم نیں، تاں
  $\Atom{X}{t_1,\ldots, t_n}$ اک ایٹمی !!{formula} اے۔

\tagitem{prvAll}{جے $!A$، !!a{formula} اے تے $u$ اک فنکشنی متغیر اے،
  تاں $\lforall[u][!A]$ وی !!a{formula} اے۔}{}

\tagitem{prvAll}{جے $!A$، !!a{formula} اے تے $X$ اک محمولی متغیر اے،
  تاں $\lforall[X][!A]$ وی !!a{formula} اے۔}{}

\tagitem{prvEx}{جے $!A$، !!a{formula} اے تے $u$ اک فنکشنی متغیر اے،
  تاں $\lexists[u][!A]$ وی !!a{formula} اے۔}{}

\tagitem{prvEx}{جے $!A$، !!a{formula} اے تے $X$ اک محمولی متغیر اے،
  تاں $\lexists[X][!A]$ وی !!a{formula} اے۔}{}
\end{enumerate}
\end{defn}
\endgroup
% /OLP-0325

% OLP-0326 content/second-order-logic/syntax-and-semantics/satisfaction.tex
\begingroup
\def\olfilename{satisfaction}
\expandafter\def\csname ol@sectioncs\endcsname{section}

\olfileid{sol}{syn}{sat}

\olsection{پورا اترنا}

\begin{explain}
دوجے درجے دے !!{formula}s لئی پورا اترن دے
تعلق~$\Sat{M}{!A}[s]$ دی تعریف کرن لئی اسیں
دوجے درجے دے !!{variable}s نوں وی تعریف وچ
شامل کرنا اے۔ دوجے درجے دی منطق وچ !!a{structure}
دا تصور اوہی اے جو پہلے درجے دی منطق وچ اے۔
فرق صرف متغیر قدر بندی~$s$ وچ اے۔ ہُن اوہ پہلے درجے دے
!!{variable}s دے نال نال دوجے درجے دے !!{variable}s نوں وی
قدراں دیندی اے۔
\end{explain}

\begin{defn}[متغیر قدر بندی]
\emph{متغیر قدر بندی}~$s$، !!a{structure}~$\Struct{M}$
لئی، اک فنکشن اے جیہڑا ہر
\begin{enumerate}
\item شے دے !!{variable}~$\Obj{v_i}$ نوں~$\Domain M$
  دے اک عنصر ول بھیجدا اے؛ یعنی
  $s(\Obj{v_i}) \in \Domain{M}$؛
\item $n$-رکنی تعلقی متغیر~$\Obj{V_i^n}$ نوں
  اک $n$-رکنی تعلق ول بھیجدا اے، جیہڑا
  ~$\Domain{M}$ اُتے اے؛ یعنی
  $s(\Obj{V_i^n}) \subseteq \Domain{M}^n$؛
\item $n$-رکنی فنکشنی متغیر~$\Obj{u_i^n}$ نوں
  اک $n$-رکنی فنکشن دیندا اے جیہڑا
  $\Domain{M}$ دے ترتیب وار ٹپل نوں $\Domain{M}$ دا اک عنصر
  دیندا اے۔ یعنی
  $s(\Obj{u_i^n})\colon \Domain{M}^n \to \Domain{M}$؛
\end{enumerate}
\end{defn}

\begin{explain}
!!^a{structure} ہر !!{constant} تے !!{function}
نوں !!a{value} دیندی اے، تے دوجے درجے دی
متغیر قدر بندی ہر شے دے متغیر تے فنکشنی متغیر
نوں اوسے ترتیب نال اک شے تے اک فنکشن دیندی اے۔ دوویں رل کے
ہر ٹرم دی قدر متعین کرن دی اجازت دیندے نیں۔
\end{explain}

\begin{defn}[ٹرم دی \usetoken{S}{value}]
جے $t$ زبان~$\Lang L$ دا ٹرم اے، $\Struct M$
اُس~$\Lang L$ لئی !!a{structure} اے، تے $s$
اُس~$\Struct M$ لئی !!a{variable} قدر بندی اے،
تاں \emph{!!{value}}~$\Value{t}{M}[s]$ پہلے
درجے دے ٹرماں وانگ تعریف ہوندی اے، نال ایہہ شق:
\begin{quote}
\indcase{t}{\Atom{u}{t_1, \ldots, t_n}}{
\[
\Value{\indfrm}{M}[s] = s(u)(\Value{t_1}{M}[s], \ldots,
\Value{t_n}{M}[s]).
\]}
\end{quote}
\end{defn}

\begin{defn}[$x$-بدل قدر بندی]
جے $s$، !!a{structure}~$\Struct M$ لئی
!!a{variable} قدر بندی اے، تاں ہر اوہ
!!{variable} قدر بندی~$s'$، جیہڑی~$\Struct M$
لئی ہووے تے~$s$ توں ودھ توں ودھ~$x$ نوں دتی
قدر وچ وکھری ہووے، اوہنوں~$s$ دی \emph{$x$-بدل قدر بندی}
آکھدے آں۔
جے $s'$، $x$ دے لحاظ نال~$s$ دی بدل قدر بندی اے، تاں
$\varAssign{s'}{s}{x}$ لکھدے آں۔
(دوجے درجے دے متغیراں $X$ یا $u$ لئی وی
ایویں ہی۔)
\end{defn}

\begin{defn}
  جے $s$، !!a{structure}~$\Struct M$ لئی
  !!a{variable} قدر بندی اے تے $m \in \Domain{M}$،
  تاں~$\Subst{s}{m}{x}$ اوہ !!{variable}
  قدر بندی اے جیہڑی ایویں تعریف
  ہوندی اے:
  \[\Subst{s}{m}{y} = \begin{cases}
    m & \text{جے } y \ident x\\
    s(y) & \text{ہور صورت وچ},
  \end{cases}\]
  جے $X$ اک $n$-رکنی تعلقی !!{variable} اے تے $M \subseteq
  \Domain{M}^n$، تاں $\Subst{s}{M}{X}$ !!{variable} قدر بندی
  ایویں تعریف ہوندی اے:
  \[\Subst{s}{M}{y} = \begin{cases}
    M & \text{جے } y \ident X\\
    s(y) & \text{ہور صورت وچ}.
  \end{cases}\]
  جے $u$ اک $n$-رکنی فنکشنی !!{variable} اے تے
  $f\colon \Domain{M}^n
  \to \Domain{M}$، تاں $\Subst{s}{f}{u}$ !!{variable} قدر بندی
  ایویں تعریف ہوندی اے:
  \[\Subst{s}{f}{y} = \begin{cases}
    f & \text{جے } y \ident u\\
    s(y) & \text{ہور صورت وچ}.
  \end{cases}\]
  ہر صورت وچ $y$ پہلے یا دوجے درجے دا کوئی وی
  !!{variable} ہو سکدا اے۔
\end{defn}

\begin{defn}[پورا اترنا]
دوجے درجے دے !!{formula}s~$!A$ لئی پورا اترن
دی تعریف \olref[fol][syn][sat]{defn:satisfaction}
ورگی اے، نال ایہہ شقّاں شامل نیں:
\begin{enumerate}
\item \indcase{!A}{\Atom{X^n}{t_1, \dots, t_n}}{$\Sat{M}{\indfrm}[s]$
  اُدوں تے صرف اُدوں جدوں $\langle \Value{t_1}{M}[s], \dots, \Value{t_n}{M}[s] \rangle \in
  s(X^n)$۔}
\tagitem{prvAll}{%
  \indcase{!A}{\lforall[X][!B]}{$\Sat{M}{\indfrm}[s]$ اُدوں تے صرف اُدوں جدوں ہر $M
  \subseteq \Domain{M}^n$ لئی $\Sat{M}{!B}[\Subst{s}{M}{X}]$۔}}{}

\tagitem{prvEx}{%
  \indcase{!A}{\lexists[X][!B]}{$\Sat{M}{\indfrm}[s]$ اُدوں تے صرف اُدوں جدوں
  گھٹ توں گھٹ اک $M \subseteq \Domain{M}^n$ ایسا ہووے جیہدے لئی
  $\Sat{M}{!B}[\Subst{s}{M}{X}]$۔}}{}

\tagitem{prvAll}{%
  \indcase{!A}{\lforall[u][!B]}{$\Sat{M}{\indfrm}[s]$ اُدوں تے صرف اُدوں جدوں ہر
  $f\colon \Domain{M}^n \to \Domain{M}$ لئی
  $\Sat{M}{!B}[\Subst{s}{f}{u}]$۔}}{}

\tagitem{prvEx}{%
  \indcase{!A}{\lexists[u][!B]}{$\Sat{M}{\indfrm}[s]$ اُدوں تے صرف اُدوں جدوں
  گھٹ توں گھٹ اک $f\colon \Domain{M}^n \to \Domain{M}$ ایسا ہووے
  جیہدے لئی $\Sat{M}{!B}[\Subst{s}{f}{u}]$۔}}{}
\end{enumerate}
\end{defn}

\begin{ex}
  !!{formula} $\lforall[z][(\Atom{X}{z} \liff \lnot
    \Atom{Y}{z})]$ ویکھو۔ ایس وچ دوجے درجے دے
  مسوّر کوئی نہیں، پر دوجے درجے دے !!{variable}s
  $X$ تے~$Y$ موجود نیں (ایتھے دوناں نوں
  اک-رکنی منّیا اے)۔ ایس دے مطابق پہلے درجے
  دا !!{sentence}
  $\lforall[z][(\Atom{P}{z} \liff \lnot \Atom{R}{z})]$
  آکھدا اے کہ جو شے~$P$ دی تعبیر وچ آؤندی اے،
  اوہ~$R$ دی تعبیر وچ نہیں آؤندی، تے اُلٹا وی۔
  کسے !!a{structure} وچ !!a{predicate}~$P$
  دی تعبیر~$\Assign{P}{M}$ نال دتی جاندی اے۔
  پر $X$ تے~$Y$ ورگے دوجے درجے دے !!{variable}s
  دی تعبیر !!{structure} خود نہیں، بلکہ
  !!a{variable} قدر بندی دیندی اے۔ چونکہ دوجے
  درجے دا !!{formula} کوئی !!a{sentence} نہیں
  (ایس وچ آزاد !!{variable}s $X$ تے~$Y$ نیں)،
  اس لئی ایہہ صرف !!a{structure}~$\Struct{M}$
  تے نال !!a{variable} قدر بندی~$s$ دے لحاظ نال
  پورا اتردا اے۔

  $\Sat{M}{\lforall[z][(\Atom{X}{z} \liff \lnot \Atom{Y}{z})]}[s]$
  اُدوں تے صرف اُدوں لاگو اے جدوں~$s(X)$ دے !!{element}s اوہو ہون
  جیہڑے دائرے وچ~$s(Y)$ دے !!{element}s نہیں۔
  یعنی اُدوں تے صرف اُدوں جدوں
  $s(Y) = \Domain{M} \setminus s(X)$۔
  مثال لئی $\Domain{M} = \{1, 2, 3\}$ لو۔
  کیوں جے کوئی !!{predicate}s، !!{function}s
  یا !!{constant}s شامل نہیں، اس لئی
  ~$\Struct{M}$ دا دائرہ ہی متعلقہ اے۔
  ہُن $s_1(X) = \{1, 2\}$ تے $s_1(Y) = \{3\}$
  لئی $\Sat{M}{\lforall[z][(\Atom{X}{z}
      \liff \lnot \Atom{Y}{z})]}[s_1]$۔

  ایس دے اُلٹ، جے $s_2(X) = \{1, 2\}$ تے
  $s_2(Y) = \{2, 3\}$ ہون، تاں
  $\Sat/{M}{\lforall[z][(\Atom{X}{z} \liff \lnot
      \Atom{Y}{z})]}[s_2]$۔ وجہ ایہہ اے کہ
  $\Sat{M}{\Atom{X}{z}}[\Subst{s_2}{2}{z}]$
  (کیوں جے $2 \in \Subst{s_2}{2}{z}(X)$)، پر
  $\Sat/{M}{\lnot \Atom{Y}{z}}[\Subst{s_2}{2}{z}]$
  (کیوں جے $2 \in
  \Subst{s_2}{2}{z}(Y)$ وی)۔
\end{ex}

\begin{ex}
  $\Sat{M}{\lexists[Y][(\lexists[y][\Atom{Y}{y}] \land
  \lforall[z][(\Atom{X}{z} \liff \lnot \Atom{Y}{z})])]}[s]$
  اُدوں ہوندا اے جدوں کوئی $N \subseteq \Domain{M}$
  ایسا ہووے کہ
  $\Sat{M}{(\lexists[y][\Atom{Y}{y}] \land \lforall[z][(\Atom{X}{z}
  \liff \lnot \Atom{Y}{z})])}[\Subst{s}{N}{Y}]$۔
  ایہہ اُنہاں صورتاں وچ ہوندا اے جدوں
  $N \neq \emptyset$ ہووے (تاکہ
  $\Sat{M}{\lexists[y][\Atom{Y}{y}]}[\Subst{s}{N}{Y}]$)
  تے پچھلی مثال وانگ
  $N = \Domain{M} \setminus s(X)$ وی ہووے۔
  دوجے لفظاں وچ،
  $\Sat{M}{\lexists[Y][(\lexists[y][\Atom{Y}{y}] \land
  \lforall[z][(\Atom{X}{z} \liff \lnot \Atom{Y}{z})])]}[s]$
  اُدوں تے صرف اُدوں جدوں
  $\Domain{M} \setminus s(X)$ خالی نہ ہووے؛ یعنی
  $s(X) \neq \Domain{M}$۔ سو مثال لئی !!{formula}
  اُس ویلے پورا اتردا اے جدوں $\Domain{M}
  = \{1, 2, 3\}$ تے $s(X) = \{1, 2\}$ ہون؛
  پر اُس ویلے نہیں جدوں $s(X) = \{1, 2, 3\}
  = \Domain{M}$ ہووے۔

  چونکہ جے $s(X) = \Domain{M}$ ہووے تاں
  !!{formula} پورا نہیں اتر دا، اس لئی
  !!{sentence}
  \[
  \lforall[X][\lexists[Y][(\lexists[y][\Atom{Y}{y}] \land
      \lforall[z][(\Atom{X}{z} \liff \lnot \Atom{Y}{z})])]]
  \]
  کدے پورا نہیں اتر دا: ہر !!{structure}~$\Struct{M}$
  لئی اوہ قدر بندی جیہڑی
  $s(X) = \Domain{M}$ رکھے، ایس !!{sentence}
  نوں جھوٹا بنا دیندی اے۔ دوجے پاسے جملہ
  \[
  \lexists[X][\lexists[Y][(\lexists[y][\Atom{Y}{y}] \land
      \lforall[z][(\Atom{X}{z} \liff \lnot \Atom{Y}{z})])]]
  \]
  ہر قدر بندی~$s$ دے لحاظ نال پورا اتردا اے،
  کیوں جے اسیں ہمیشہ کوئی $M \subseteq \Domain{M}$
  لبھ سکدے آں جیہڑا $M \neq \Domain{M}$ ہووے
  (مثلاً $M = \emptyset$)۔
\end{ex}

\begin{ex}
  دوجے درجے دا !!{sentence}~$\lforall[X][\lforall[y][X(y)]]$
  آکھدا اے کہ ہر $1$-رکنی تعلق، یعنی ہر خاصیت،
  ہر شے اُتے صادق اے۔ ایہہ ظاہر اے کہ ایہہ کدے
  سچ نہیں: ہر~$\Struct{M}$ وچ اک متغیر قدر بندی~$s$
  لئی، جیہڑی $s(X) = \emptyset$ تے
  $s(y) = a \in \Domain{M}$ رکھے، ساڈے کول
  $\Sat/{M}{X(y)}[s]$ اے۔ ایس دا مطلب اے کہ
  $!A \lif \lforall[X][\lforall[y][X(y)]]$
  دوجے درجے دی منطق وچ $\lnot !A$ دے ہم ارز اے۔
  یعنی $\Sat{M}{!A \lif
  \lforall[X][\lforall[y][X(y)]]}$ اُدوں تے صرف
  اُدوں جدوں $\Sat{M}{\lnot !A}$۔
  دوجے لفظاں وچ، دوجے درجے دی منطق وچ اسیں
  $\lnot$ نوں $\lforall$ تے~$\lif$ نال
  تعریف کر سکدے آں۔
\end{ex}

\begin{prob}
  دِکھاؤ کہ دوجے درجے دی منطق وچ $\lforall$
  تے~$\lif$ نال اسیں ہور رابطیاں دی تعریف کر سکدے آں:
  \begin{enumerate}
    \item ثابت کرو کہ دوجے درجے دی منطق وچ
    $!A \land !B$،
    $\lforall[X][(!A \lif (!B \lif \lforall[x][X(x)])\lif
    \lforall[x][X(x)])]$ دے ہم ارز اے۔
    \item صرف $\lforall$ تے $\lif$ ورت کے
    $!A \lor !B$ دے ہم ارز دوجے درجے دا
    اک فارمولا لبھو۔
  \end{enumerate}
\end{prob}
\endgroup
% /OLP-0326

% OLP-0327 content/second-order-logic/syntax-and-semantics/semantic-notions.tex
\begingroup
\def\olfilename{semantic-notions}
\expandafter\def\csname ol@sectioncs\endcsname{section}

\olfileid{sol}{syn}{sem}
\olsection{معنوی تصورات}

\begin{explain}
\emph{منطقی اعتبار}، \emph{منطقی نتیجہ} تے
\emph{پورا ہون جوگتا} دیاں تعریفاں دوجے درجے دی منطق وچ
اوہو جیہیاں نیں جیہڑیاں پہلے درجے دی منطق وچ نیں۔ فرق ایہ
اے کہ ہُن پورا اترن دا تعلق دوجے درجے دے !!{formula}s لئی
ورتیدا اے۔ دوجے درجے دا
!!{sentence} اوہ !!a{formula} اے جیہدے سارے
متغیر پابند ہون، محمولی تے فنکشنی متغیر وی۔
\end{explain}

\begin{defn}[منطقی اعتبار]
جملہ $!A$ \emph{معتبر} اے، $\Entails !A$،
اُدوں تے صرف اُدوں جدوں $\Sat{M}{!A}$ ہر
!!{structure}~$\Struct M$ لئی ہووے۔
\end{defn}

\begin{defn}[منطقی نتیجہ]
جملیاں دے سیٹ~$\Gamma$ توں جملہ~$!A$
\emph{منطقی طور تے نکلدا اے}، $\Gamma
\Entails !A$، اُدوں تے صرف اُدوں جدوں ہر اوہ !!{structure}~$\Struct M$ جیہدے وچ
$\Sat{M}{\Gamma}$ ہووے، اوہدے وچ $\Sat{M}{!A}$ وی ہووے۔
\end{defn}

\begin{defn}[پورا ہون جوگتا]
جملیاں دا سیٹ~$\Gamma$ \emph{پورا ہون جوگا} اے
جے $\Sat{M}{\Gamma}$ کسے
!!{structure}~$\Struct M$ لئی ہووے۔ جے
$\Gamma$ پورا ہون جوگا نہ ہووے، تاں اوہنوں
\emph{پورا ہون جوگا نہیں} آکھدے نیں۔
\end{defn}
\endgroup
% /OLP-0327

% OLP-0328 content/second-order-logic/syntax-and-semantics/expressive-power.tex
\begingroup
\def\olfilename{expressive-power}
\expandafter\def\csname ol@sectioncs\endcsname{section}

\olfileid{sol}{syn}{exp}
\olsection{اظہار دی طاقت}

\begin{explain}
دوجے درجے دے متغیراں اُتے مسوّریت زبان دی
اظہار دی طاقت نوں پہلے درجے دی منطق دے مقابل
بہت ودھا دیندی اے۔ دوجے درجے دی وجودی مسوّریت
نال اسیں آکھ سکدے آں کہ کجھ خاصیتاں والے فنکشن یا تعلق موجود نیں۔ پہلے درجے دی منطق
وچ ایہہ کہن دا اکلوتا طریقہ ایہہ اے کہ اوس مقصد
لئی اک غیر منطقی علامت، یعنی !!a{function}
یا !!{predicate}، مقرر کیتی جائے۔ دوجے درجے
دی عمومی مسوّریت نال اسیں آکھ سکدے آں کہ
دائرے دے سارے ذیلی سیٹ، اُس اُتے سارے
تعلق، یا !!{domain} توں !!{domain} ول سارے
فنکشوناں وچ کوئی خاصیت اے۔ پہلے درجے دی
منطق وچ اسیں صرف اوہناں ذیلی سیٹاں، تعلقاں
یا فنکشوناں بارے ایہہ آکھ سکدے آں جیہڑے زبان
دی کسے غیر منطقی علامت نوں دتے گئے ہون۔
تے جدوں اسیں کسے خاصیت والے ذیلی سیٹاں،
تعلقاں یا فنکشوناں دے وجود، یا انہاں سبھناں
دے اوہ خاصیت رکھن، بارے آکھدے آں، تاں اُس
خاصیت نوں بیان کردیاں وی دوجے درجے دی مسوّریت
ورت سکدے آں۔ ایس نال اسیں اوہناں تعلقاں دی تعریف کر سکدے آں
جیہڑے پہلے درجے دی منطق وچ قابلِ تعریف نہیں، تے دائرے دیاں اوہ خاصیتاں
بیان کر سکدے آں جو پہلے درجے دی منطق وچ
بیان نہیں ہو سکدیاں۔
\end{explain}

\begin{defn}
جے $\Struct{M}$ زبان~$\Lang{L}$ لئی
!!a{structure} اے، تاں تعلق~$R \subseteq
\Domain{M}^2$،~$\Lang{L}$ وچ
\emph{قابلِ تعریف} اے جے کوئی !!{formula}~$!A_R(x, y)$
ایہو جیہا ہووے جیہدے وچ صرف $x$
تے~$y$ آزاد ہون، تے $R(a, b)$
(یعنی $\tuple{a,b} \in R$) اُدوں تے صرف
اُدوں لاگو اے جدوں $\Sat{M}{!A_R(x, y)}[s]$
اوہ قدر بندی لئی جیہڑی $s(x) = a$
تے $s(y) = b$ رکھے۔
\end{defn}

\begin{ex}
پہلے درجے دی منطق وچ اسیں برابری دا
تعلق~$\Id{\Domain{M}}$ (یعنی
$\Setabs{\tuple{a,a}}{a \in
  \Domain{M}}$) فارمولے~$\eq[x][y]$
نال متعین کر سکدے آں۔ دوجے درجے دی منطق وچ
ایہہ تعلق \emph{$\eq$ توں بناں} وی متعین ہو سکدا اے۔ جے $a$ تے $b$،~$\Domain{M}$
دے اکّو !!{element} ہون، تاں اوہ
~$\Domain{M}$ دے اوہناں ہی ذیلی سیٹاں دے
!!{element}s ہون گے (کیوں جے سیٹ اپنے
!!{element}s نال متعین ہوندے نیں)۔
اُلٹے پاسے، جے $a$ تے $b$ وکھرے ہون،
تاں کوئی ذیلی سیٹ ایسا اے جیہدے !!{element}s وچ اک اے
تے دوجا نہیں۔ مثال لئی $a \in \{a\}$،
پر $b
\notin \{a\}$ جے $a \neq b$۔
سو «~$\Domain{M}$ دے اکّو ذیلی سیٹاں دے
!!{element}s ہونا» اوہ تعلق اے جیہڑا
$a$ تے $b$ دے لئی اُدوں تے صرف اُدوں
درست اے جدوں $a = b$۔ ایہہ تعلق
دوجے درجے دی منطق وچ بیان ہو سکدا اے،
کیوں جے اسیں~$\Domain{M}$ دے سارے
ذیلی سیٹاں اُتے مسوّریت کر سکدے آں۔
سو ہِٹھلا !!{formula}،
$\Id{\Domain{M}}$ نوں متعین کردا اے:
\[
\lforall[X][(X(x) \liff X(y))]
\]
\end{ex}

\begin{prob}
دکھاؤ کہ $\lforall[X][(X(x) \lif X(y))]$
(دھیان: $\lif$ اے،~$\liff$ نہیں!)
$\Id{\Domain{M}}$ نوں متعین کردا اے۔
\end{prob}

\begin{ex}
جے $R$ اک دو رُکنی !!{predicate} اے، تاں
$\Assign{R}{M}$،~$\Domain{M}$ اُتے
اک دو رُکنی تعلق اے۔ شاید کجھ الجھن
پیدا ہووے، پر اسیں $R$ نوں !!{predicate}~$R$
تے خود تعلق~$\Assign{R}{M}$ دوناں لئی
ورتاں گے۔ \emph{تعدی تکمیل}~$R^*$، جیہڑی~$R$
دی اے، اوہ تعلق اے جو $a$ تے $b$ دے درمیان
اُدوں تے صرف اُدوں ہوندا اے جدوں کجھ
$c_1$, \dots, $c_k$ لئی
$R(a,c_1)$, $R(c_1, c_2)$, \dots,
$R(c_k,b)$ درست ہون۔ ایس وچ $k = 0$
والی صورت وی شامل اے؛ یعنی جے
$R(a,b)$ ہوندا اے، تاں $R^*(a,b)$
وی ہوندا اے۔ ایس دا مطلب
$R \subseteq R^*$ اے۔ اصل وچ
$R^*$ اوہ سب توں چھوٹا متعدی تعلق اے
جیہڑا~$R$ نوں شامل کردا اے۔
دوجے درجے دی منطق وچ
اسیں ایہہ آکھ سکدے آں کہ $X$ اک
متعدی تعلق اے جیہڑا~$R$ نوں شامل
کردا اے:
\begin{multline*}
  \formula B_R(X) \ident \lforall[x][\lforall[y][(R(x,y) \lif X(x, y))]] \land {}\\
\lforall[x][\lforall[y][\lforall[z][((X(x,y) \land X(y,z)) \lif X(x,
      z))]]].
\end{multline*}
پہلا عطف آکھدا اے کہ $R \subseteq X$،
تے دوجا آکھدا اے کہ $X$ متعدی اے۔

ایہہ آکھنا کہ $X$ ایہو جیہا سب توں
چھوٹا تعلق اے، اِس گل دے برابر اے کہ
اوہ خود ہر اوس تعلق وچ شامل اے جیہڑا
$R$ نوں شامل کردا تے متعدی اے۔
سو اسیں~$R$ دی تعدی تکمیل ہیٹھلے
!!{formula} نال متعین کر سکدے آں:
\[
R^*(X) \ident !B_R(X) \land \lforall[Y][(!B_R(Y) \lif
  \lforall[x][\lforall[y][(X(x, y) \lif Y(x,y))]])].
\]
$\Sat{M}{R^*(X)}[s]$ اُدوں تے صرف اُدوں لاگو اے
جدوں $s(X) = R^*$۔
~$R$ دی تعدی تکمیل پہلے درجے دی
منطق وچ بیان نہیں ہو سکدی۔
\end{ex}
\endgroup
% /OLP-0328

% OLP-0329 content/second-order-logic/syntax-and-semantics/inf-count.tex
\begingroup
\def\olfilename{inf-count}
\expandafter\def\csname ol@sectioncs\endcsname{section}

\olfileid{sol}{syn}{siz}

\olsection{لامتناہی تے \usetoken{S}{enumerable}
  \usetoken{P}{domain} دا بیان}

سیٹ~$M$ (ڈیڈیکائنڈ معنیاں وچ) لامتناہی اُدوں تے صرف اُدوں اے جدوں
اک !!a{injective} (اک بہ اک) فنکشن~$f\colon M
\to M$ موجود ہووے جیہڑا !!{surjective} (ہر ہدف قدر لین والا) نہ ہووے؛ یعنی $\ran{f} \neq M$
ہووے۔ پہلے درجے دی منطق وچ اسیں اک اک-رکنی !!{function}~$f$ نوں
ویکھ سکدے آں تے آکھ سکدے آں کہ فنکشن~$\Assign{f}{M}$، جیہڑا
!!a{structure}~$\Struct{M}$ وچ اوس نوں دتا گیا اے، !!{injective}
اے تے $\ran{f} \neq
\Domain{M}$ اے:
\[
\lforall[x][\lforall[y][(\eq[f(x)][f(y)] \to \eq[x][y])]] \land
\lexists[y][\lforall[x][\eq/[y][f(x)]]].
\]
جے $\Struct{M}$ وچ ایہ !!{sentence} پورا اتردا اے، تاں
$\Assign{f}{M}:
\Domain{M} \to \Domain{M}$، !!{injective} اے؛ ایس لئی
$\Domain{M}$ لازماً لامتناہی اے۔ اُلٹے پاسے، جے $\Domain{M}$
لامتناہی اے تے ایس طرح دا فنکشن موجود اے، تاں اسیں
$\Assign{f}{M}$ نوں اوہ فنکشن دے سکدے آں۔ فیر $\Struct{M}$
وچ ایہ !!{sentence} پورا اترے گا۔ پر ایس طریقے لئی ضروری اے کہ ساڈی
زبان وچ اوس مقصد لئی ورتی غیر منطقی علامت~$f$ موجود ہووے۔
دوجے درجے دی منطق وچ اسیں سِدھا آکھ سکدے آں کہ ایہو جیہا فنکشن
\emph{موجود اے}۔ ہُن~$f$ علامت دی لوڑ نہیں رہندی، تے سانوں خالص
دوجے درجے دی منطق دا !!{sentence} ملدا اے:
\[
\fn{Inf} \ident \lexists[u][(\lforall[x][\lforall[y][(\eq[u(x)][u(y)]
      \lif \eq[x][y])]] \land \lexists[y][\lforall[x][\eq/[y][u(x)]]])].
\]
$\Sat{M}{\fn{Inf}}$ اُدوں تے صرف اُدوں جدوں $\Domain{M}$
لامتناہی اے۔ ہُن اسیں $\fn{Fin} \ident \lnot \fn{Inf}$ تعریف
کر سکدے آں؛ $\Sat{M}{\fn{Fin}}$ اُدوں تے صرف اُدوں جدوں
$\Domain{M}$ متناہی اے۔ خالص پہلے درجے دی منطق دا کوئی اک
!!{sentence}، !!{domain} دے لامتناہی ہون نوں بیان نہیں کر سکدا،
بھاویں ایہو جیہے جملیاں دا اک لامتناہی سیٹ ایہہ گل بیان کر سکدا اے۔
خالص پہلے درجے دی منطق دے !!{sentence}s دا کوئی وی سیٹ ایہو
جیہا نہیں جیہڑا کسے !!a{structure} وچ اُدوں تے صرف اُدوں پورا اترے
جدوں اوہدا دائرہ متناہی ہووے۔

\begin{prop}
$\Sat{M}{\fn{Inf}}$ اُدوں تے صرف اُدوں جدوں $\Domain{M}$ لامتناہی اے۔
\end{prop}

\begin{proof}
$\Sat{M}{\fn{Inf}}$ اُدوں تے صرف اُدوں جدوں
  $\Sat{M}{\lforall[x][\lforall[y][(\eq[u(x)][u(y)] \lif \eq[x][y])]]
    \land \lexists[y][\lforall[x][\eq/[y][u(x)]]]}[s]$ کسے~$s$
  لئی درست ہووے۔ جے ایہہ درست اے، تاں $s(u)$ اک !!a{injective}
  فنکشن اے، تے کوئی $y \in \Domain{M}$،~$s(u)$ دی حاصل قدراں
  وچ نہیں آؤندا۔ اُلٹ پاسے، جے اک !!a{injective}
  $f\colon \Domain{M} \to \Domain{M}$ ایہو جیہا اے جیہدے لئی
  $\ran{f}
  \neq \Domain{M}$، تاں $s(u) = f$ رکھ کے ایہو جہی متغیر
  قدر بندی مل جاندی اے۔
\end{proof}

سیٹ $M$، !!{enumerable}
اُدوں اے جدوں اوہدے !!{element}s دی اک گنتی وار فہرست
\[
m_0, m_1, m_2, \dots
\]
موجود ہووے (دہراؤ توں بغیر، پر ممکن اے متناہی ہووے)۔ ایہو جہی
فہرست اُدوں تے صرف اُدوں موجود اے جدوں کوئی !!a{element}
$z \in M$ تے اک فنکشن $f\colon M \to M$ ایہو جیہے ہون کہ
$z$، $f(z)$، $f(f(z))$، \dots،~$M$ دے سارے !!{element}s
ہون۔ جے فہرست موجود اے تاں $z = m_0$ تے
$f(m_k) = m_{k+1}$ (یا $f(m_k) = m_k$ جے $m_k$ فہرست دا آخری
!!{element} اے) مطلوبہ !!{element} تے فنکشن دندے نیں۔
اُلٹ پاسے، جے ایہو جیہے $z$ تے $f$ موجود نیں، تاں
$z$، $f(z)$، $f(f(z))$، \dots،~$M$ دی گنتی وار فہرست اے؛
ایس لئی $M$، !!{enumerable} اے۔ دوجے درجے دی منطق وچ اسیں
$z$ تے~$f$ دے وجود نوں بیان کر سکدے آں، تے ایہو جیہا
!!{sentence} بنا سکدے آں جیہڑا کسے !!a{structure} وچ اُدوں
تے صرف اُدوں درست ہووے جدوں اوہ !!{structure}، !!{enumerable} ہووے:
\[
\fn{Count} \ident
\lexists[z][\lexists[u][\lforall[X][((X(z) \land
      \lforall[x][(X(x) \lif X(u(x)))]) \lif \lforall[x][X(x)])]]]
\]

\begin{prop}
$\Sat{M}{\fn{Count}}$ اُدوں تے صرف اُدوں جدوں $\Domain{M}$،
!!{enumerable} اے۔
\end{prop}

\begin{proof}
فرض کرو $\Domain{M}$، !!{enumerable} اے، تے $m_0$، $m_1$،
\dots، اوہدی گنتی وار فہرست اے۔ دہراؤ ہٹا کے اسی یقین
کر سکدے آں کہ کوئی $m_k$ دو واری
نہ آوے۔ $f(m_k) = m_{k+1}$ رکھو۔ ہُن
$s(z) = m_0$ تے $s(u) = f$ رکھو۔ اسیں وکھاؤندے آں کہ
\[
\Sat{M}{\lforall[X][((X(z) \land \lforall[x][(X(x) \lif X(u(x)))])
    \lif \lforall[x][X(x)])]}[s]
\]
فرض کرو $M \subseteq \Domain{M}$ کوئی وی ذیلی سیٹ اے۔ ایہ وی فرض کرو کہ
$\Sat{M}{(X(z) \land \lforall[x][(X(x) \lif
X(u(x)))])}[\Subst{s}{M}{X}]$۔ تاں $\Subst{s}{M}{X}(z) \in M$
تے جدوں وی $x \in M$ ہووے، $(\Subst{s}{M}{X}(u))(x) \in M$
وی ہووے گا۔ دوجے لفظاں وچ، کیوں جے
$\varAssign{\Subst{s}{M}{X}}{s}{X}$ اے، $m_0 \in M$ اے
تے جے $x \in M$ تاں $f(x) \in M$؛ ایس لئی
$m_0 \in M$، $m_1 = f(m_0) \in M$،
$m_2 = f(f(m_0)) \in M$، وغیرہ۔ لہٰذا $M = \Domain{M}$،
تے ایس لئی
$\Sat{M}{\lforall[x][X(x)]}[\Subst{s}{M}{X}]$۔ کیوں جے
$M \subseteq
\Domain{M}$ کوئی وی ذیلی سیٹ سی، سو:
$\Sat{M}{\fn{Count}}$۔

ہُن فرض کرو کہ $\Sat{M}{\fn{Count}}$، یعنی
\[
\Sat{M}{\lforall[X][((X(z) \land \lforall[x][(X(x) \lif X(u(x)))])
    \lif \lforall[x][X(x)])]}[s]
\]
کسے~$s$ لئی۔ $m = s(z)$ تے $f = s(u)$ رکھو، تے
$M = \{m,
f(m), f(f(m)), \dots\}$ نوں ویکھو۔ ایس طرح تعریف
کیتا~$M$ صاف طور تے !!{enumerable} اے۔ فیر مفروضے دے مطابق
\[
\Sat{M}{(X(z) \land \lforall[x][(X(x) \lif X(u(x)))])
    \lif \lforall[x][X(x)]}[\Subst{s}{M}{X}]
\]
پورا اتردا اے۔ نال ای $\Sat{M}{X(z)}[\Subst{s}{M}{X}]$ وی
پورا اتردا اے، کیوں جے $M \ni m =
\Subst{s}{M}{X}(z)$ اے؛ تے
$\Sat{M}{\lforall[x][(X(x) \lif
X(u(x)))]}[\Subst{s}{M}{X}]$ وی پورا اتردا اے، کیوں جے
جدوں وی $x \in M$ ہووے، $f(x) \in
M$ وی ہوندا اے۔ سو کیوں جے شرط تے پورا شرطیہ دوویں پورے اتردے نیں، اوس دا
نتیجہ وی پورا اترنا لازمی اے:
$\Sat{M}{\lforall[x][X(x)]}[\Subst{s}{M}{X}]$۔ پر ایس
دا مطلب $M = \Domain{M}$ اے، تے ایس لئی $\Domain{M}$،
!!{enumerable} اے، کیوں جے تعریف مطابق $M$ اے۔
\end{proof}

\begin{prob}
!!{sentence} $\fn{Inf} \land \fn{Count}$ عین اوہناں دائریاں وچ
درست اے جیہڑے !!{denumerable} نیں۔ $\fn{Count}$ دی تعریف
ایویں بدلو کہ اک وکھرا !!{sentence} بنے جیہڑا دائرے دے
!!{denumerable} ہون نوں سِدھا بیان کرے، تے ثابت کرو کہ اوہ
ایویں ہی کردا اے۔
\end{prob}
\endgroup
% /OLP-0329

% OLP-0716 content/second-order-logic/syntax-and-semantics/language-of-sol.tex
\begingroup
\def\olfilename{language-of-sol}
\expandafter\def\csname ol@sectioncs\endcsname{section}

\olfileid{sol}{syn}{lan}

\olsection{دوجے درجے دی منطق دی زبان}

پہلے درجے دی منطق وانگ دوجے درجے دی منطق دیاں عبارتاں اک بنیادی لفظ نامے توں بندیاں نیں، جیہدے وچ \emph{!!{variable}s}، \emph{!!{constant}s}، \emph{!!{predicate}s} تے کدے \emph{!!{function}s} ہوندے نیں۔ ایہناں توں، منطقی جوڑاں، مسوّراں تے قوساں تے کامیاں ورگیاں وقفے دیاں نشانیاں نال رلا کے، \emph{ٹرم} تے \emph{!!{formula}s} بندے نیں۔ فرق ایہ اے کہ شیواں دے متغیراں توں علاوہ دوجے درجے دی منطق وچ تعلقاں تے فنکشناں دے متغیر وی نیں، تے ایہ اوہناں اُتے مسوّر لاؤن دی کھُل دیندی اے۔ سو دوجے درجے دی منطق دیاں منطقی نشانیاں پہلے درجے دی منطق والیاں نیں، تے نال ایہ:

\begin{enumerate}
\item اک !!{denumerable}s سیٹ دوجے درجے دے تعلقاں دے !!{variable}s دا، ہر تھاں-گنتی~$n$ لئی: $\Obj V_0^n$، $\Obj V_1^n$، $\Obj V_2^n$، \dots
\item اک !!{denumerable}s سیٹ دوجے درجے دے فنکشناں دے !!{variable}s دا:
  $\Obj u_0^n$، $\Obj u_1^n$، $\Obj u_2^n$، \dots
\end{enumerate}

پہلے درجے دی منطق وچ !!{identity}~$\eq$ عام طور اُتے شامل ہوندی اے۔ پہلے درجے دی منطق وچ زبان~$\Lang{L}$ دیاں غیر منطقی نشانیاں کوئی دلچسپ گل ظاہر کرن لئی ضروری نیں۔ بے شک !!{sentences} وی نیں جیہڑے کوئی غیر منطقی نشانی نہیں ورتدے، پر صرف~$\eq$ نال کوئی دلچسپ گل کہنا اوکھا اے۔ دوجے درجے دی منطق وچ تعلقاں تے فنکشناں دے متغیر بے حد نیں، ایس لئی جیہڑی وی گل پہلے درجے دی زبان وچ کہہ سکدے آں، اوہ ایتھے غیر منطقی نشانیاں دی خاص رسد توں بناں وی کہہ سکدے آں۔
\endgroup
% /OLP-0716
\OLEndChapterHook
\endgroup
% /OLP-0323

% OLP-0330 content/second-order-logic/metatheory/metatheory.tex
\begingroup
\def\olfilename{metatheory}
\expandafter\def\csname ol@sectioncs\endcsname{section}

\olchapter{sol}{met}{دوجے درجے دی منطق دا ماورائے نظریہ}

% OLP-0331 content/second-order-logic/metatheory/introduction.tex
\begingroup
\def\olfilename{introduction}
\expandafter\def\csname ol@sectioncs\endcsname{section}

\olfileid{sol}{met}{int}

\olsection{تعارف}

پہلے درجے دی منطق دیاں کئی چنگیاں خاصیتاں نیں۔ اسیں جاندے آں کہ ایہ
فیصلہ کرن جوگ نہیں، پر گھٹ توں گھٹ ایہ مسلّماتی طور تے متعین کیتی جا سکدی
اے۔ یعنی پہلے درجے دی منطق لئی صحت مند تے تمام ثبوتی نظام موجود نیں۔ اوہناں توں اک !!a{derivability} تعلق~$\Proves$ ملدا اے۔ ایس دی خاصیت ایہ اے کہ !!{sentence}s دے کسے وی سیٹ~$\Gamma$ تے کسے وی
!!{sentence}~$!A$ لئی، $\Gamma \Entails !A$ اُدوں تے صرف اُدوں جدوں
$\Gamma \Proves !A$ ہووے۔ ایس توں خاص طور تے ایہ نکلدا اے کہ پہلے
درجے دی منطق دے منطقی طور تے معتبر جملے !!{computably enumerable} نیں۔ اک قابلِ حساب فنکشن~$f\colon \Nat \to
\Sent[L]$ موجود اے۔ ایس~$f$ دیاں قدراں عین~$\Lang{L}$ دے سارے تے صرف اوہی
!!{sentence}s نیں جیہڑے منطقی طور تے معتبر نیں۔ ایہ ایس لئی اے کہ
!!{derivation}s دی مؤثر گنتی وار فہرست بن سکدی اے۔ ایس فہرست وچ
جیہڑا اخذ اک~!!{sentence} !!{derive} کردا اے، اوہنوں اوسے~!!{sentence}
نال جوڑ دیندے آں۔ دوجے درجے دی منطق دی اظہاری طاقت پہلے درجے نالوں ودھ
اے، ایس لئی عام طور تے اوہدے سارے معتبر جملے لبھنے ودھ اوکھے نیں۔
اصل وچ اسیں دکھاواں گے کہ دوجے درجے دی منطق نہ صرف ناقابلِ فیصلہ اے
بلکہ ایس دے معتبر جملے !!{computably enumerable} وی نہیں۔ ایس دا مطلب اے کہ دوجے درجے دی منطق لئی کوئی
ایہو جیہا صحت مند تے تمام ثبوتی نظام نہیں ہو سکدا جیہدے اخذ مؤثر
طور تے فہرست بن سکدے ہون (بھاویں صحت مند پر نامکمل ثبوتی نظام موجود
نیں تے حقیقت وچ تحقیق دے اہم موضوع نیں)۔

پہلے درجے دی منطق دیاں ہور دو خاصیتاں وی نیں: ایہ محدود گواہی دی
خاصیت رکھدی اے (جے !!{sentence}s دے کسے سیٹ~$\Gamma$ دا ہر محدود
ذیلی سیٹ پورا ہون جوگا اے، تاں خود $\Gamma$ وی پورا ہون جوگا اے)،
تے ایس اُتے لووینہائم--سکولم دا قضیہ لاگو ہوندا اے (جے $\Gamma$
دا اک لامتناہی نمونہ ہووے، تاں اوہدا اک !!a{denumerable} نمونہ وی ہوندا اے)۔ ایہ دونیں نتیجے دوجے درجے دی
منطق لئی ناکام ہوندے نیں۔ وجہ فیر ایہی اے کہ دوجے درجے دی منطق
!!{domain}s دے آکار بارے اوہ حقائق بیان کر سکدی اے جیہڑے پہلے درجے
دی منطق نہیں کر سکدی۔
\endgroup
% /OLP-0331

% OLP-0332 content/second-order-logic/metatheory/second-order-arithmetic.tex
\begingroup
\def\olfilename{second-order-arithmetic}
\expandafter\def\csname ol@sectioncs\endcsname{section}

\olfileid{sol}{met}{spa}

\olsection{دوجے درجے دا حساب}

یاد کرو کہ پیانو حساب دے نظریے~$\Th{PA}$ وچ~$\Th{Q}$ دے آٹھ
مسلّمے شامل نیں:
\begin{align*}
& \lforall[x][\eq/[x'][\Obj 0]]\\
& \lforall[x][\lforall[y][(\eq[x'][y'] \lif \eq[x][y])]]\\
& \lforall[x][(\eq[x][\Obj 0] \lor \lexists[y][\eq[x][y']])]\\ 
& \lforall[x][\eq[(x + \Obj 0)][x]]\\
& \lforall[x][\lforall[y][\eq[(x + y')][(x + y)']]]\\
& \lforall[x][\eq[(x \times \Obj 0)][\Obj 0]]\\
& \lforall[x][\lforall[y][\eq[(x \times y')][((x \times y) + x)]]]\\
& \lforall[x][\lforall[y][(x < y \liff \lexists[z][\eq[(z' + x)][y]])]]\\
\intertext{تے نال ایس روپ دے سارے جملے وی}
& (!A(\Obj 0) \land \lforall[x][(!A(x) \lif !A(x'))]) \lif \lforall[x][!A(x)].
\end{align*}
ایہ آخری بیان اک «خاکہ» اے، یعنی اک روپ جیہدے توں حساب دی زبان دے
لامتناہی !!{sentence}s بندے نیں، ہر !!{formula}~$!A(x)$ لئی اک۔
اسیں ایس نوں (پہلے درجے دا) \emph{استقرا دا مسلّماتی خاکہ} آکھدے آں۔
\emph{دوجے درجے دے} پیانو حساب~$\Th{PA^2}$ وچ استقرا اکّو جملے
نال بیان ہو سکدا اے۔ $\Th{PA^2}$ وچ اُپرلے پہلے آٹھ مسلّماں دے نال
(دوجے درجے دا) \emph{استقرا دا مسلّمہ} شامل اے:
\[
\lforall[X][((X(\Obj 0) \land \lforall[x][(X(x) \lif X(x'))]) \lif \lforall[x][X(x)])].
\]
ایہ کہندا اے کہ جے !!{domain} دا کوئی ذیلی سیٹ~$X$،
~$\Assign{\Obj{0}}{M}$ نوں شامل کرے تے دائرے دے کسے وی~$x \in \Domain{M}$
نوں شامل کرے تاں~$\Assign{\prime}{M}(x)$ نوں وی شامل کرے
(یعنی اوہ «جانشین لین اُتے بند» ہووے)، تاں اوہ پورے !!{domain} نوں
شامل کردا اے (یعنی $X =
\Domain{M})$۔

استقرا دا مسلّمہ ایہ یقینی بناؤندا اے کہ ہر اوہ !!{structure} جیہدے وچ
ایہ مسلّمہ پورا اتردا اے، اوہدے~$\Domain{M}$ وچ صرف اوہی !!{element}s نیں
جنہاں دا موجود ہونا مسلّمے لازم کردے نیں، یعنی~$\num{n}$ دیاں
قدراں، جدوں $n \in \Nat$ ہووے۔ $\Th{PA^2}$ دے کسے نمونے وچ
غیر معیاری عدد نہیں ہوندا۔

\begin{thm}
\ollabel{thm:sol-pa-standard}
جے $\Sat{M}{\Th{PA^2}}$ ہووے، تاں $\Domain{M} =
\Setabs{\Value{\num{n}}{M}}{n \in \Nat}$۔
\end{thm}

\begin{proof}
من لو $N = \Setabs{\Value{\num{n}}{M}}{n \in \Nat}$، تے فرض کرو
$\Sat{M}{\Th{PA^2}}$۔ بے شک کسے وی $n \in \Nat$ لئی
$\Value{\num{n}}{M} \in \Domain{M}$، ایس لئی $N \subseteq \Domain{M}$۔

ہُن دوجے رخ دی شمولیت ویکھدے آں۔ متغیراں دی ایہو جیہی قدر بندی~$s$ لو
جیہدے لئی $s(X) = N$ ہووے۔ مفروضے مطابق:
\begin{align*}
\Sat{M}{\lforall[X][((X(\Obj 0) \land \lforall[x][(X(x)
      \lif X(x'))]) \lif \lforall[x][X(x)])]}, & \text{ ایس لئی}\\
\Sat{M}{(X(\Obj 0) \land \lforall[x][(X(x) \lif X(x'))]) \lif
  \lforall[x][X(x)]}[s]. & 
\end{align*}
ایس شرطیہ دا اگلا حصہ ویکھو۔ $\Value{\Obj{0}}{M} \in
N$، ایس لئی $\Sat{M}{X(\Obj{0})}[s]$۔ عطف دا دوجا جزو،
$\lforall[x][(X(x) \lif X(x'))]$ وی پورا اتردا اے۔ فرض
کرو $x \in N$۔ $N$ دی تعریف مطابق،
$x = \Value{\num{n}}{M}$ ہووے گا، کسے~$n$ لئی۔ فیر
$\Assign{\prime}{M}(x) = \Value{\num{n+1}}{M} \in N$۔ ایس لئی
$\Assign{\prime}{M}(x) \in N$۔

اسیں ثابت کیتا اے کہ
$\Sat{M}{X(\Obj 0) \land \lforall[x][(X(x) \lif X(x'))]}[s]$۔
نتیجتاً $\Sat{M}{\lforall[x][X(x)]}[s]$۔ ایس دا مطلب اے کہ
ہر $x \in \Domain{M}$ لئی $x \in s(X) = N$۔ سو
$\Domain{M} \subseteq N$۔
\end{proof}

\begin{cor}
\ollabel{cor:sol-pa-categorical}
~$\Th{PA^2}$ دے کوئی وی دو نمونے ہم شکل نیں۔
\end{cor}

\begin{proof}
\olref{thm:sol-pa-standard} مطابق $\Th{PA^2}$ دے ہر نمونے دا دائرہ
$\Value{\num{n}}{M}$ دیاں قدراں نال ای پورا ہو جاندا اے۔ ایہو جیہا ہر
نمونہ~$\Th{Q}$ دا وی نمونہ اے۔ \olref[mod][mar][stm]{prop:thq-standard}
مطابق ایسا ہر نمونہ معیاری اے، یعنی~$\Struct{N}$ نال ہم شکل اے۔
\end{proof}

اُپر اسیں $\Th{PA^2}$ نوں پہلے آٹھ حسابی مسلّماں تے دوجے درجے
دے استقرا دے مسلّمے والا نظریہ منّیا سی۔ دوجے درجے دی منطق دی
اظہاری طاقت کر کے دوجے درجے دے پیانو حساب لئی
حساب دے مسلّماں وچوں \emph{صرف پہلے دو} تے استقرا کافی نیں۔

\begin{prop}\ollabel{prop:sol-pa-definable}
من لو $\Th{PA^{2\dagger}}$ اوہ دوجے درجے دا نظریہ اے جیہدے وچ
حساب دے پہلے دو مسلّمے (جانشین دے مسلّمے) تے دوجے درجے دا
استقرا دا مسلّمہ شامل نیں۔ فیر $\le$، $+$ تے $\times$،
$\Th{PA^{2\dagger}}$ وچ قابلِ تعریف نیں۔
\end{prop}

\begin{proof}
$\le$ دے قابلِ تعریف ہون لئی سانوں اک فارمولا~$!A_\le(x, y)$
لبھنا پئے گا جیہدے لئی $\Sat{N}{!A_\le(\num n, \num m)}$ اُدوں
تے صرف اُدوں ہووے جدوں $n \le m$۔ ایہ فارمولا ویکھو:
\[
!B(x, Y) \ident Y(x) \land \lforall[y][(Y(y) \lif Y(y'))]
\]
صاف اے کہ $!B(\num n, Y)$ کسے سیٹ $Y \subseteq \Nat$ لئی اُدوں
تے صرف اُدوں پورا اتردا اے جدوں
$\Setabs{m}{n \le m} \subseteq Y$۔ ایس لئی اسیں
$!A_\le(x, y) \ident
\lforall[Y][(!B(x, Y) \lif Y(y))]$ لے سکدے آں۔

جمع دے قابلِ تعریف ہون لئی ویکھو کہ $k+l=m$ اُدوں تے صرف اُدوں
جدوں کوئی فنکشن $u$ موجود ہووے جیہدے لئی $u(0)=k$،
$u(n')=u(n)'$ ہر $n$ لئی، تے $m = u(l)$ ہووے۔ ایس ہم ارزی
نال $\Th{PA^{2\dagger}}$ وچ جمع دی تعریف ایس فارمولے نال کر سکدے آں:
% PNB-CMP-072: source's quantified w must occur in the recurrence body.
\[
!A_+(x,y,z) \ident \lexists[u][(u(\Obj 0)=x \land \lforall[w][u(w')=u
 (w)'] \land u(y)=z)]
\] 
ایتھوں صاف اے کہ $\Sat{N}{!A_+(\num k, \num l, \num m)}$
اُدوں تے صرف اُدوں ہووے جدوں $k+l=m$۔
\end{proof}

\begin{prob}
\olref[sol][met][spa]{prop:sol-pa-definable} دا ثبوت پورا کرو۔
\end{prob}
\endgroup
% /OLP-0332

% OLP-0333 content/second-order-logic/metatheory/undecidability-and-axiomatizability.tex
\begingroup
\def\olfilename{undecidability-and-axiomatizability}
\expandafter\def\csname ol@sectioncs\endcsname{section}

\olfileid{sol}{met}{nax}

\olsection{دوجے درجے دی منطق مسلّماں نال متعین نہیں ہو سکدی}

\begin{thm}
\ollabel{thm:sol-undecidable}
دوجے درجے دی منطق قابلِ فیصلہ نہیں۔
\end{thm}

\begin{proof}
پہلے درجے دا کوئی !!{sentence} پہلے درجے دی منطق وچ اُدوں تے صرف
اُدوں معتبر اے جدوں اوہ دوجے درجے دی منطق وچ معتبر ہووے۔
پہلے درجے دی منطق قابلِ فیصلہ نہیں۔
\end{proof}

\begin{thm}
\ollabel{cor:sol-not-axiomatizable}
دوجے درجے دی منطق لئی کوئی !!{derivation} نظام ایسا نہیں
جیہڑا صحت مند تے تمام دوویں ہووے۔
\end{thm}

\begin{proof}
حساب دی زبان دا کوئی !!a{sentence}~$!A$ من لو۔
$\Sat{N}{!A}$ اُدوں تے صرف اُدوں جدوں $\Th{PA^2} \Entails !A$۔
من لو $!P$ نظریے~$\Th{PA^2}$ دے نو مسلّماں نوں عطف نال جوڑ کے بنیا جملہ اے۔
$\Th{PA^2} \Entails !A$ اُدوں تے صرف اُدوں جدوں
$\Entails !P \lif !A$، یعنی ہر ساخت وچ
% PNB-CMP-073: source places Sat's closing brace outside inline math.
$\Sat{M}{!P \lif !A}$۔ ہن ایہہ !!{sentence} ویکھو:
$\lforall[z][\lforall[u][\lforall[u'][\lforall[u''][\lforall[L][(!P' \lif !A')]]]]]$۔
ایہ اصل جملے وچ $\Obj{0}$ دی تھاں $z$، $\prime$ دی تھاں
اک-رُکنی فنکشنی متغیر~$u$، $+$ تے $\times$ دی تھاں بالترتیب
دو رُکنی فنکشنی متغیر~$u'$ تے~$u''$، تے $<$ دی تھاں
دو رُکنی تعلقی متغیر~$L$ رکھ کے، فیر انہاں سب اُتے عمومی
مسوّریت لا کے بنیا اے۔ ایہ خالص دوجے درجے دی منطق دا
معتبر جملہ اُدوں تے صرف اُدوں اے جدوں پہلا شرطیہ جملہ معتبر اے؛
ایہ اُدوں تے صرف اُدوں اے جدوں $\Th{PA^2} \Entails !A$،
تے ایہ اُدوں تے صرف اُدوں اے جدوں $\Sat{N}{!A}$۔
سو جے دوجے درجے دی منطق لئی کوئی صحت مند تے تمام
ثبوتی نظام ہووے، تاں اسیں اوس راہیں اک قابلِ حساب گنتی وار
فہرست~$f\colon \Nat \to \Sent[L_A]$ بنا سکدے آں جیہدیاں
قدراں عین اوہی !!{sentence}s نیں جیہڑے $\Struct{N}$ وچ سچے نیں۔ ایہ فنکشن
$\Th{Q}$ وچ کسے پہلے درجے دے فارمولے~$!B_f(x, y)$ نال
قابلِ نمائندگی ہووے گا۔ فیر !!{formula}~$\lexists[x][!B_f(x, y)]$
پہلے درجے دے سچ !!{sentence}s دا سیٹ~$\Struct{N}$ وچ متعین
کرے گا۔ ایہ تارسکی دے قضیے دے خلاف اے۔
\end{proof}
\endgroup
% /OLP-0333

% OLP-0334 content/second-order-logic/metatheory/compactness.tex
\begingroup
\def\olfilename{compactness}
\expandafter\def\csname ol@sectioncs\endcsname{section}

\olfileid{sol}{met}{com}

\olsection{دوجے درجے دی منطق محدود گواہی دی خاصیت نہیں رکھدی}

\begin{explain}
جملیاں دے سیٹ~$\Gamma$ نوں \emph{محدود طور تے پورا ہون جوگا}
آکھدے آں جے اوہدا ہر محدود ذیلی سیٹ پورا ہون جوگا ہووے۔ پہلے درجے دی منطق
وچ جے !!{sentence}s دا سیٹ~$\Gamma$ محدود طور تے پورا ہون جوگا ہووے،
تاں اوہ پورا ہون جوگا وی اے۔ ایس خاصیت نوں \emph{محدود گواہی}
کہندے نیں۔ منطقی نتیجے دی بولی وچ ایہدا اک ہم ارز روپ وی اے: جے
$\Gamma \Entails !A$، تاں $\Gamma_0 \Entails !A$ کسے
محدود ذیلی سیٹ~$\Gamma_0 \subseteq \Gamma$ لئی وی۔ ایس روپ وچ
ایہ تمامیت دے قضیے دا سیدھا نتیجہ اے: جے $\Gamma \Entails !A$،
تاں تمامیت دے مطابق $\Gamma \Proves !A$۔ پر کوئی !!{derivation}،
$\Gamma$ دے صرف محدود گنتی دے !!{sentence}s ای ورت سکدا اے۔

دوجے درجے دی منطق وچ محدود گواہی دی خاصیت نہیں۔ دوجے درجے دے
!!{sentence}s دے ایسے سیٹ نیں جیہڑے محدود طور تے پورے ہون جوگے
نیں پر آپ پورے ہون جوگے نہیں۔ اوہناں توں کوئی~$!A$ منطقی طور
اُتے نکلدا اے، پر اوہناں دے کسے محدود ذیلی سیٹ توں~$!A$
منطقی طور اُتے نہیں نکلدا۔
\end{explain}

% PNB-CMP-074: frozen source repeats the prior undecidability theorem label.
\begin{thm}
\ollabel{thm:sol-not-compact}
دوجے درجے دی منطق محدود گواہی دی خاصیت نہیں رکھدی۔
\end{thm}

\begin{proof}
یاد کرو کہ
\[
\fn{Inf} \ident \lexists[u][(\lforall[x][\lforall[y][(\eq[u(x)][u(y)]
      \lif \eq[x][y])]] \land \lexists[y][\lforall[x][\eq/[y][u(x)]]])]
\]
!!a{structure} وچ اُدوں تے صرف اُدوں پورا اتردا اے جدوں اوہدا
دائرہ لامتناہی ہووے۔ من لو $!A^{\ge n}$ ایسا جملہ اے جیہڑا
آکھدا اے کہ دائرے وچ گھٹ توں گھٹ $n$ !!{element}s نیں، مثلاً
\[
!A^{\ge n} \ident \lexists[x_1][\dots\lexists[x_n][(\eq/[x_1][x_2]
    \land \eq/[x_1][x_3] \land \dots \land \eq/[x_{n-1}][x_n])]].
\]
!!{sentence}s دا ایہ سیٹ ویکھو:
\[
\Gamma = \{\lnot \fn{Inf}, !A^{\ge 1}, !A^{\ge 2}, !A^{\ge 3}, \dots\}.
\]
ایہ محدود طور تے پورا ہون جوگا اے۔ کسے وی محدود ذیلی سیٹ~$\Gamma_0
\subseteq \Gamma$ لئی کوئی مثبت $k$ چن لو جیہڑا اوس ذیلی سیٹ وچ
آئے ہر کم توں کم آکار دے اشاریے توں گھٹ نہ ہووے۔ جے ذیلی سیٹ
وچ ایہو جیہی کوئی شرط نہ ہووے تاں اک چن لو۔ فیر $!A^{\ge k} \in \Gamma$،
پر کوئی
% PNB-CMP-075: the finite subset, not Gamma, has no larger size requirement.
$!A^{\ge n} \in \Gamma_0$ نہیں جیہدے لئی $n>k$ ہووے۔
جے $\Domain{M}$ وچ $k$ !!{element}s ہون، تاں $\Sat{M}{\Gamma_0}$۔
پر $\Gamma$ پورا ہون جوگا نہیں: جے $\Sat{M}{\lnot \fn{Inf}}$،
تاں $\Domain{M}$ متناہی اے، منّ لو آکار~$k$۔ فیر
$\Sat/{M}{!A^{\ge k+1}}$۔
\end{proof}

\begin{prob}
ایسا سیٹ~$\Gamma$ تے !!a{sentence}~$!A$ دی مثال دو جیہدے لئی
$\Gamma \Entails !A$، پر ہر محدود ذیلی سیٹ~$\Gamma_0 \subseteq
\Gamma$ لئی $\Gamma_0 \Entails/ !A$ ہووے۔
\end{prob}
\endgroup
% /OLP-0334

% OLP-0335 content/second-order-logic/metatheory/loewenheim-skolem.tex
\begingroup
\def\olfilename{loewenheim-skolem}
\expandafter\def\csname ol@sectioncs\endcsname{section}

\olfileid{sol}{met}{lst}

\olsection{لووینہائم--سکولم دا قضیہ دوجے درجے دی منطق وچ لاگو نہیں ہوندا}

\begin{explain}
لووینہائم--سکولم (L\"owenheim--Skolem) دا نزولی قضیہ آکھدا اے
کہ !!{sentence}s دے ہر اوس سیٹ دا !!a{enumerable} نمونہ وی ہوندا اے
جیہدا کوئی لامتناہی نمونہ ہووے۔ ایہ وی تمامیت دے قضیے توں نکلدا
اے: تمامیت دا ثبوت !!{sentence}s دے ہر بے تضاد سیٹ لئی اک نمونہ
بناؤندا اے، تے اوہ نمونہ !!{enumerable} ہوندا اے۔ اک صعودی
لووینہائم--سکولم قضیہ وی اے: جے !!{sentence}s دے کسے سیٹ دا
!!a{denumerable} نمونہ ہووے، تاں اوہدا !!a{nonenumerable}
نمونہ وی ہوندا اے۔ ایہ دوویں قضیے دوجے درجے دی منطق وچ لاگو نہیں ہوندے۔
\end{explain}


\begin{thm}
\ollabel{thm:sol-no-ls} لووینہائم--سکولم دا قضیہ دوجے درجے دی منطق وچ
لاگو نہیں ہوندا۔ ایہو جیہے !!{sentence}s
نیں جنہاں دے لامتناہی نمونے نیں، پر کوئی !!{enumerable} نمونہ نہیں۔
\end{thm}

\begin{proof}
یاد کرو کہ
\[
\fn{Count} \ident \lexists[z][\lexists[u][\lforall[X][((X(z) \land
      \lforall[x][(X(x) \lif X(u(x)))]) \lif \lforall[x][X(x)])]]]
\]
!!a{structure}~$\Struct{M}$ وچ اُدوں تے صرف اُدوں درست اے جدوں
$\Domain{M}$، !!{enumerable} ہووے۔ ایس لئی $\lnot\fn{Count}$
ساخت~$\Struct{M}$ وچ اُدوں تے صرف اُدوں درست اے جدوں
$\Domain{M}$، !!{nonenumerable} ہووے۔ ایہو جہیاں !!{structure}s
موجود نیں: کسے وی !!{nonenumerable} سیٹ نوں !!{domain} بنا لو،
مثلاً $\Pow{\Nat}$ یا~$\Real$۔ سو $\lnot \fn{Count}$ دے
لامتناہی نمونے نیں، پر کوئی !!{enumerable} نمونہ نہیں۔
\end{proof}

\begin{thm}
ایہو جیہے !!{sentence}s نیں جنہاں دے !!{denumerable} نمونے نیں،
پر کوئی !!{nonenumerable} نمونہ نہیں۔
\end{thm}

\begin{proof}
$\fn{Count} \land \fn{Inf}$، $\Nat$ وچ درست اے، پر کسے وی
!!{structure}~$\Struct{M}$ وچ درست نہیں جدوں $\Domain{M}$،
!!{nonenumerable} ہووے۔
\end{proof}
\endgroup
% /OLP-0335

\OLEndChapterHook
\endgroup
% /OLP-0330

% OLP-0336 content/second-order-logic/sol-and-set-theory/sol-and-set-theory.tex
\begingroup
\def\olfilename{sol-and-set-theory}
\expandafter\def\csname ol@sectioncs\endcsname{section}

\olchapter{sol}{set}{دوجے درجے دی منطق تے سیٹ نظریہ}

\begin{editorial}
ایہ حصہ دوجے درجے دی منطق وچ قوتی سیٹاں تے استمرار (continuum) دی کوڈ بندی بارے اے۔ نتیجے بیان کیتے گئے نیں، پر
انہاں دے ثبوت ہجے لکھنے باقی نیں۔ ہجے کوئی مشق نہیں، تے خود
تعریفاں تے نتیجیاں وچ وی مسئلے ہو سکدے نیں۔
ایہنوں احتیاط نال ورتو، تے جے کوئی گل غلط یا غیر واضح ہووے تاں دسّو۔
\end{editorial}

% OLP-0337 content/second-order-logic/sol-and-set-theory/introduction.tex
\begingroup
\def\olfilename{introduction}
\expandafter\def\csname ol@sectioncs\endcsname{section}

\olfileid{sol}{set}{int}

\olsection{تعارف}

دوجے درجے دی منطق دائرے دے ذیلی سیٹاں دے نال نال فنکشوناں اُتے وی
مسوّریت کر سکدی اے؛ ایس لئی توقع اے کہ گھٹ توں گھٹ سیٹ نظریہ دا
کجھ حصہ ایس منطق وچ کیتا جا سکدا اے۔ «کیتا جا سکدا اے» توں مراد
ایہ اے کہ سیٹ نظریہ دیاں خاصیتاں تے بیان دوجے درجے دی منطق
وچ ظاہر کیتے جا سکدے نیں، تے ایس لئی سیٹاں واسطے کسے خاص غیر منطقی
علامتاں دے ذخیرے دی لوڑ نہیں، مثلاً سیٹ نظریہ وچ رکنیت دا
!!{predicate}۔ مثال لئی اسیں سیٹاں دے اتحاد تے اشتراک تے ذیلی سیٹ
ہون دا تعلق متعین کر سکدے آں۔ سیٹاں دے سائز دا موازنہ وی کر سکدے آں تے کانتور دے قضیے ورگے نتیجے بیان کر سکدے آں۔
\endgroup
% /OLP-0337

% OLP-0338 content/second-order-logic/sol-and-set-theory/comparing-sets.tex
\begingroup
\def\olfilename{comparing-sets}
\expandafter\def\csname ol@sectioncs\endcsname{section}

\olfileid{sol}{set}{cmp}

\olsection{سیٹاں دا موازنہ}

\begin{prop}
ایہ !!{formula} $\lforall[x][(X(x) \lif Y(x))]$ ذیلی سیٹ ہون دا
تعلق متعین کردا اے، یعنی $\Sat{M}{\lforall[x][(X(x) \lif Y(x))]}[s]$
اُدوں تے صرف اُدوں لاگو اے جدوں $s(X) \subseteq s(Y)$۔
\end{prop}

\begin{prop}
ایہ !!{formula} $\lforall[x][(X(x) \liff Y(x))]$ سیٹاں اُتے
برابری دا تعلق متعین کردا اے، یعنی
$\Sat{M}{\lforall[x][(X(x) \liff Y(x))]}[s]$
اُدوں تے صرف اُدوں لاگو اے جدوں $s(X) = s(Y)$۔
\end{prop}

\begin{prop}
ایہ !!{formula} $\lexists[x][X(x)]$ خالی نہ ہون دی خاصیت متعین
کردا اے، یعنی $\Sat{M}{\lexists[x][X(x)]}[s]$ اُدوں تے صرف اُدوں لاگو اے جدوں
$s(X) \neq \emptyset$۔
\end{prop}

اک سیٹ~$X$، سیٹ~$Y$ نالوں سائز وچ وڈا نہیں، یعنی
$\cardle{X}{Y}$، اُدوں تے صرف اُدوں جدوں اک !!a{injective} فنکشن
$f\colon X \to Y$ موجود ہووے۔ چونکہ اسیں ایہ ظاہر کر سکدے آں کہ
کوئی فنکشن اک بہ اک اے، تے ایہ وی کہ~$X$ دے رکناں اُتے اوہدیاں
قدراں~$Y$ وچ نیں، ایس لئی اسیں دائرے دے ذیلی سیٹاں اُتے سائز وچ وڈا
نہ ہون دا تعلق وی متعین کر سکدے آں۔

\begin{prop}
ایہ فارمولا
\[
\lexists[u][(\lforall[x][(X(x) \lif Y(u(x)))] \land \lforall[x][\lforall[y][((X(x) \land X(y) \land \eq[u(x)][u(y)]) \lif
      \eq[x][y])]])]
\]
سائز وچ وڈا نہ ہون دا تعلق متعین کردا اے۔
% PNB-CMP-076: صرف X دے دو رکناں لئی اک بہ اک شرط؛ منجمد اصل دا پوری دائرے اُتے injectivity والا غیر ضروری تقاضا درست کیتا۔
\end{prop}

دو سیٹاں دا سائز اکّو اے، یا اوہ ”ہم تعداد“ (equinumerous) نیں،
یعنی $\cardeq{X}{Y}$، اُدوں تے صرف اُدوں جدوں اک !!a{bijective} فنکشن
~$f\colon X \to Y$ موجود ہووے۔

\begin{prop}
ایہ فارمولا
\begin{multline*}
  \lexists[u][(\lforall[x][(X(x) \lif Y(u(x)))] \land {}\\
    \lforall[x][\lforall[y][((X(x) \land X(y) \land \eq[u(x)][u(y)]) \lif \eq[x][y])]]
  \land {} \\\lforall[y][(Y(y) \lif \lexists[x][(X(x)
      \land \eq[y][u(x)])])])]
\end{multline*}
ہم تعداد ہون دا تعلق متعین کردا اے۔
% PNB-CMP-077: X اُتے مقامی injectivity؛ منجمد اصل دی global شرط برابر سائز والے بعض لامتناہی سیٹ رد کردندی سی۔
\end{prop}

اسیں ایہناں !!{formula}s نوں ترتیب وار $X \subseteq Y$،
$X = Y$، $X \neq \emptyset$، $\cardle{X}{Y}$، تے
$\cardeq{X}{Y}$ نال مختصر لکھاں گے۔ (ایہ گل تھوڑی الجھن پیدا کر
سکدی اے، کیوں جے سیٹاں~$X$ تے~$Y$ بارے غیر رسمی گل کردیاں اسیں
ایہی علامتاں ورتدے آں؛ پر ایتھے ایہ علامتاں دوجے درجے دی منطق دے
اوہناں !!{formula}s دے مخفف نیں جنہاں وچ اک-رُکنی تعلقی متغیر
$X$ تے~$Y$ شامل نیں۔)

\begin{prop}
ایہ !!{sentence} $\lforall[X][\lforall[Y][((\cardle{X}{Y} \land
    \cardle{Y}{X}) \lif \cardeq{X}{Y})]]$ معتبر اے۔
\end{prop}

\begin{proof}
ایہ !!{sentence} کسے !!a{structure}~$\Struct{M}$ وچ پورا اتردا اے جے
ہر ذیلی سیٹ $X \subseteq \Domain{M}$ تے $Y \subseteq \Domain{M}$ لئی،
$\cardle{X}{Y}$ تے $\cardle{Y}{X}$ ہون توں $\cardeq{X}{Y}$ نکلے۔
پر ایہ گل \emph{ہر} سیٹ~$X$ تے~$Y$ لئی درست اے---ایہی
شروڈر-برنسٹین دا قضیہ اے۔
\end{proof}
\endgroup
% /OLP-0338

% OLP-0339 content/second-order-logic/sol-and-set-theory/cardinalities.tex
\begingroup
\def\olfilename{cardinalities}
\expandafter\def\csname ol@sectioncs\endcsname{section}

\olfileid{sol}{set}{crd}

\olsection{سیٹاں دی کارڈینیلٹی}

\begin{explain}
جیویں اسیں ایہ ظاہر کر سکدے آں کہ دائرہ متناہی یا لامتناہی اے،
!!{enumerable} یا !!{nonenumerable} اے، اویویں اسیں~$\Domain{M}$ دے کسے
ذیلی سیٹ دے متناہی یا لامتناہی، !!{enumerable} یا !!{nonenumerable} ہون دی خاصیت
دی تعریف وی کر سکدے آں۔
\end{explain}

\begin{prop}
ایہ فارمولا $\fn{Inf}(X) \ident$
\begin{multline*}
\lexists[u][(\lforall[x][\lforall[y][((X(x) \land X(y) \land \eq[u(x)][u(y)]) \lif
      \eq[x][y])]] \land \lforall[x][(X(x) \lif X(u(x)))] \land {} \\\lexists[y][(X(y) \land \lforall[x][(X(x)
      \lif \eq/[y][u(x)])])])]
\end{multline*}
متغیر قدر بندی~$s$ دے لحاظ نال اُدوں تے صرف اُدوں پورا اتردا اے جدوں
$s(X)$ لامتناہی ہووے۔
% PNB-CMP-078: X اُتے اک بہ اک خود نقشہ تے X وچ بندش شامل؛اصل global u نال متناہی X نوں وی لامتناہی آکھ سکدا سی۔
\end{prop}

\begin{prop}
ایہ فارمولا $\fn{Count}(X) \ident $
\begin{multline*}
(\lnot\lexists[x][X(x)] \lor {} \\
\lexists[z][\lexists[u][(X(z) \land
    \lforall[x][(X(x) \lif X(u(x)))] \land {} \\ \lforall[Y][((Y(z) \land
      \lforall[x][(Y(x) \lif Y(u(x)))]) \lif X = Y])]])
)
\end{multline*}
متغیر قدر بندی~$s$ دے لحاظ نال اُدوں تے صرف اُدوں پورا اتردا اے جدوں
$s(X)$، !!{enumerable} ہووے۔
% PNB-CMP-079: خالی سیٹ لئی الگ disjunct؛اصل z∈X غیر خالی enumerable سیٹ ہی گنّدا سی۔
\end{prop}

کانتور دے قضیے توں سانوں پتہ اے کہ !!{nonenumerable} سیٹ موجود نیں؛
اصل وچ لامتناہی سائز دے وکھ وکھ درجے نیں، تے اوہناں دی گنتی وی لامتناہی اے۔
سیٹ نظریہ سیٹاں دے سائز دا پورا حساب بناؤندا اے تے سیٹاں نوں
لامتناہی کارڈینیل اعداد دیندا اے۔ قدرتی اعداد متناہی سیٹاں دے سائز ماپن والے
کارڈینیل اعداد دا کم کردے نیں۔ !!{denumerable}
سیٹاں دی کارڈینیلٹی پہلی لامتناہی کارڈینیلٹی اے، جیہنوں
~$\aleph_0$ (”الیف-ناٹ“ یا ”الیف-زیرو“) آکھدے نیں۔ اگلا
لامتناہی سائز~$\aleph_1$ اے۔ ایہ اوہ سب توں چھوٹا سائز اے جیہڑا
کسے غیر قابلِ شمار سیٹ دا ہو سکدا اے؛ایہو جیہا سیٹ نہ متناہی ہوندا
اے نہ سائز~$\aleph_0$ رکھدا اے۔ اسیں ”$X$ دا سائز $\aleph_0$
اے“ ایویں تعریف کر سکدے آں:
$\fn{Aleph}_0(X) \liff \fn{Inf}(X) \land \fn{Count}(X)$۔
سیٹ~$X$ دا سائز~$\aleph_1$ اُدوں تے صرف اُدوں اے جدوں
اوہدا ہر اوہ ذیلی سیٹ جیہڑا سائز وچ اوہدے توں گھٹ اے، یا تاں
متناہی ہووے یا اوہدا سائز~$\aleph_0$ ہووے۔ اصل سیٹ آپ نہ
متناہی ہووے تے نہ اوہدا سائز~$\aleph_0$ ہووے۔
ایس لئی اسیں ایہ فارمولا لکھ سکدے آں:
$\fn{Aleph_1}(X) \ident \lforall[Y][((Y \subseteq X \land
  \lnot\cardeq{Y}{X}) \lif
  (\lnot\fn{Inf}(Y) \lor \fn{Aleph}_0(Y)))] \land \lnot
\fn{Count}(X)$۔ سائز~$\aleph_2$ وی ایس طرح تعریف ہوندا اے، وغیرہ۔
% PNB-CMP-080: Y=X نوں smaller-subset antecedent توں کڈھ کے cardinal strictness،تے X نوں nonenumerable منگیا؛اصل finite X قبول تے aleph_1 X رد کردا سی۔

اک سائز خاص دلچسپی رکھدا اے: استمرار (continuum) دی کارڈینیلٹی۔
ایہ~$\Pow{\Nat}$ دا سائز اے، یعنی اوہی سائز جیہڑا~$\Real$ دا اے۔
ایہ گل کہ کوئی سیٹ استمرار دے سائز دا اے، دوجے درجے دی منطق وچ
وی ظاہر کیتی جا سکدی اے، پر ایس لئی کجھ ہور کم کرنا پیندا اے۔
\endgroup
% /OLP-0339

% OLP-0340 content/second-order-logic/sol-and-set-theory/power-of-continuum.tex
\begingroup
\def\olfilename{power-of-continuum}
\expandafter\def\csname ol@sectioncs\endcsname{section}

\olfileid{sol}{set}{pow}

\olsection{استمرار دا کارڈینیل سائز}

\begin{explain}
دوجے درجے دی منطق وچ اسیں !!{domain} دے ذیلی سیٹاں اُتے مسوّریت کر
سکدے آں، پر !!{domain} دے ذیلی سیٹاں دے سیٹاں اُتے نہیں۔ایہ کم
براہ راست کرن لئی سانوں \emph{تیجے درجے} دی منطق چاہیدی ہووے گی۔
مثال لئی، جے اسیں کانتور دا ایہ قضیہ لکھنا چاہیئے کہ کسے سیٹ دے
قوتی سیٹ توں خود اوس سیٹ ول کوئی !!{injective} (اک بہ اک) فنکشن
نہیں ہوندا،تاں اسیں شاید ایویں کوشش کریے: ”ہر سیٹ~$X$ تے ہر
سیٹ~$P$ لئی،جے $P$،~$X$ دا قوتی سیٹ اے،تاں
$\cardle{P}{X}$ نہیں“۔ایہ آکھن لئی کہ $P$،~$X$ دا قوتی سیٹ اے،
سانوں ایہ رسمی بنانا پئے گا کہ $P$ دے !!{element}s عین~$X$ دے
سارے ذیلی سیٹ نیں؛یعنی کجھ ایس طرح:
$\lforall[Y][(P(Y) \liff Y \subseteq X)]$۔دقت $P(Y)$ وچ اے:
ایہ دوجے درجے دی منطق دا !!a{formula} نہیں،کیونکہ $P$ ورگے
اک-رُکنی تعلقی متغیر دی دلیل دی تھاں صرف ٹرم آ سکدے نیں۔

پر جے دائرہ کافی وڈا ہووے،اسیں سیٹاں دے سیٹاں اُتے مسوّریت دی
\emph{نقل} بنا سکدے آں۔خیال ایہ اے کہ دو رُکنی تعلق~$R$،
!!{domain} دے !!{element}s نوں !!{domain} دے !!{element}s نال
جوڑدا اے۔ایہو جیہا~$R$ ملے تاں کسے~$x$ نال $R$ دے ذریعے
جڑے سارے !!{element}s اکٹھے کیتے جا سکدے نیں:
$\Setabs{y \in \Domain{M}}{R(x,y)}$ اوہ سیٹ اے جیہنوں~$x$
”کوڈ“ کردا اے۔الٹ پاسے،جے
$Z \subseteq \Pow{\Domain{M}}$،~$\Domain{M}$ دے ذیلی سیٹاں دا
کوئی مجموعہ ہووے،تے~$\Domain{M}$ وچ گھٹ توں گھٹ~$Z$ دے سیٹاں
جِنّے !!{element}s ہون،تاں اک تعلق~$R \subseteq \Domain{M}^2$
وی ملدا اے جس نال ہر $Y \in Z$ نوں کوئی~$x$ کوڈ کردا اے۔
ایتھے کوڈ بناؤن والا تعلق~$R$ اے۔
\end{explain}

\begin{defn}
جے $R \subseteq \Domain{M}^2$،تاں
\emph{$x$،$R$ دے لحاظ نال~$\Setabs{y
  \in \Domain{M}}{R(x, y)}$ نوں کوڈ کردا اے}۔
\end{defn}

جے کوئی !!a{element}~$x \in \Domain{M}$،
$R$ دے لحاظ نال سیٹ $Z \subseteq \Domain{M}$ نوں کوڈ کردا اے،
تاں سیٹ $Y \subseteq \Domain{M}$ سیٹاں دے اک سیٹ نوں کوڈ کردا اے:
یعنی اوہ سارے سیٹ جیہڑے~$Y$ دے !!{element}s کوڈ کردے نیں۔ ایس طرح کوئی
سیٹ~$Y$،$R$ دے لحاظ نال~$\Pow{X}$ نوں کوڈ کر سکدا اے۔ایہ
عین اوس ویلے ہوندا اے جدوں ہر $Z \subseteq X$ نوں کوئی $x
\in Y$،$R$ دے لحاظ نال~$Z$ نوں کوڈ کرے،تے ہر $x \in Y$،
$R$ دے لحاظ نال کسے $Z \subseteq X$ نوں کوڈ کرے۔

\begin{prop}
ایہ !!{formula}
  \[
  \fn{Codes}(x, R, Z) \ident \lforall[y][(Z(y) \liff
    R(x, y))]
  \]
ظاہر کردا اے کہ $s(x)$،$s(R)$ دے لحاظ نال~$s(Z)$ نوں کوڈ
کردا اے۔ایہ !!{formula}
\begin{multline*}
  \fn{Pow}(Y, R, X) \ident \\
  \lforall[Z][(Z \subseteq X \lif \lexists[x][(Y(x) \land
    \fn{Codes}(x, R, Z))])] \land {} \\ \lforall[x][(Y(x) \lif
  \lforall[Z][(\fn{Codes}(x, R, Z) \lif Z \subseteq X)]]
\end{multline*}
ظاہر کردا اے کہ $s(Y)$،$s(R)$ دے لحاظ نال~$s(X)$ دا قوتی
سیٹ کوڈ کردا اے۔ یعنی $s(Y)$ دے رکن $s(R)$ دے لحاظ نال عین
$s(X)$ دے ذیلی سیٹاں نوں کوڈ کردے نیں۔
\end{prop}

\begin{explain}
ایس ترکیب نال اسیں ذیلی سیٹاں دی تھاں اوہناں دے کوڈاں اُتے
مسوّر لا کے قوتی سیٹ بارے بیان لکھ سکدے آں۔ مثال لئی، کانتور دا
قضیہ ہُن ایویں لکھیا جا سکدا اے کہ~$X$
دا قوتی سیٹ کوڈ کرن والے کسے وی تعلق دے کوڈاں دے سیٹ توں
خود~$X$ ول کوئی !!{injective} (اک بہ اک) فنکشن نہیں۔
\end{explain}

\begin{prop}
ایہ جملہ
\begin{multline*}
  \lforall[X][\lforall[Y][\lforall[R][(\fn{Pow}(Y, R, X) \lif \\
      \lnot \lexists[u][(\lforall[x][\lforall[y][(\eq[u(x)][u(y)] \lif
            \eq[x][y])]] \land {}\\
        \lforall[x][(Y(x) \lif X(u(x)))])])]]]
\end{multline*}
منطقی طور اُتے معتبر اے۔
\end{prop}

\begin{explain}
!!a{denumerable} سیٹ دا قوتی سیٹ !!{nonenumerable} اے۔ ایس لئی اوہدی
کارڈینیلٹی ہر !!{denumerable} سیٹ دے سائز،یعنی~$\aleph_0$،
توں وڈی اے۔$\Pow{\Nat}$ دے سائز نوں ایتھے ”استمرار دا سائز“
آکھیا گیا اے،کیونکہ ایہ حقیقی عددی خط،~$\Real$،دے نقطیاں
جِنّا اے۔جے !!{domain}،کسی !!a{denumerable} سیٹ دا قوتی سیٹ
کوڈ کرن لئی کافی وڈا ہووے،تاں اسیں ایہ آکھ کے ظاہر کر سکدے آں
کہ اک سیٹ استمرار دے سائز دا اے: اوہ سیٹ~$Y$ دے ہم تعداد اے
جیہڑا سیٹ~$X$،جیہدا سائز~$\aleph_0$ اے،دے قوتی سیٹ نوں
کوڈ کردا اے تے ہر ذیلی سیٹ لئی عین اک وکھرے کوڈ دی شرط پوری کردا اے۔ (جے دائرہ اِنّا وڈا نہ
ہووے،یعنی اس وچ~$\Real$
نال ہم تعداد کوئی ذیلی سیٹ نہ ہووے،تاں~$\Pow{X}$ نوں کوڈ کرن
والا تعلق وی نہیں ہو سکدا۔)
% PNB-CMP-083: اصل any coding Y آکھدا؛duplicate codes نال |Y| وڈا ہو سکدا،B14 دی uniqueness شرط لازم۔
\end{explain}

\begin{prop}
جے $\cardle{\Real}{\Domain{M}}$،تاں ایہ !!{formula}
\begin{multline*}
\fn{Cont}(Y) \ident 
\lexists[X][\lexists[R][((\fn{Aleph}_0(X) \land
      \fn{Pow}(Y, R, X)) \land \\ 
      \lforall[x][\lforall[y][((Y(x) \land {} 
      Y(y) \land \lforall[z][R(x,z) \liff R(y,z)]) \lif \eq[x][y])]])]]
\end{multline*}
ظاہر کردا اے کہ $\cardeq{s(Y)}{\Real}$۔
\end{prop}

\begin{proof}
  $\fn{Pow}(Y, R, X)$ ظاہر کردا اے کہ~$s(Y)$،$s(R)$ دے
  لحاظ نال~$s(X)$ دا قوتی سیٹ کوڈ کردا اے؛$\fn{Aleph}_0(X)$
  آکھدا اے کہ ایہ اصل سیٹ گنتی وار لامتناہی اے۔ ایس لئی~$s(Y)$ دا
  سائز کم از کم استمرار دے سائز جِنّا اے،اگرچہ وڈا وی ہو سکدا اے
  (جے~$s(Y)$ دے کئی رکن~$X$ دا اکّو ذیلی سیٹ کوڈ کرن)۔ آخری
  عطفی جزو ایہ گنجائش ختم کردا اے۔ اوہ منگدا اے کہ
  $s(Y)$ دے رکناں تے~$s(X)$ دے ذیلی سیٹاں وچ نسبت~$s(R)$
  دے ذریعے !!{injective} (اک بہ اک) ہووے۔
% PNB-CMP-081: اصل proof وچ s(Z) بےتعریف؛ایتھے مقرر base s(X) اے۔
\end{proof}

\begin{prop}
$\cardeq{\Domain{M}}{\Real}$ اُدوں تے صرف اُدوں جدوں
\begin{multline*}
  \Sat{M}{\lexists[Y][(\fn{Cont}(Y) \land {} \\
      \lexists[u][(\lforall[x][\lforall[y][(\eq[u(x)][u(y)] \lif
            \eq[x][y])]] \land {} \\
        \lforall[y][(Y(y) \liff \lexists[x][\eq[y][u(x)]])])])]}
        \end{multline*}
  ۔
% PNB-CMP-082: Cont(Y) exact continuum codes،تے u:D→Y bijection؛اصل Y⊆range(u) زائد بڑے D نوں وی قبول کردا سی۔
  \end{prop}

\begin{explain}
استمرار دا مفروضہ ایہ بیان اے کہ استمرار دا سائز پہلی
!!{nonenumerable} کارڈینیلٹی اے؛یعنی~$\Pow{\Nat}$ دا
سائز~$\aleph_1$ اے۔
\end{explain}

\begin{prop}
استمرار دا مفروضہ سچ اے اُدوں تے صرف اُدوں جدوں
\[\fn{CH} \ident
\lforall[X][(\fn{Aleph}_1(X) \liff \fn{Cont}(X))]\]
ہر ساخت وچ درست ہووے۔
\end{prop}

ایہ ہم ارزی درست نہیں کہ استمرار دا مفروضہ جھوٹا ہووے اُدوں تے
صرف اُدوں جدوں $\lnot \fn{CH}$ معتبر ہووے۔ !!a{enumerable} دائرے
وچ نہ~$\aleph_1$ سائز دا کوئی ذیلی سیٹ ہوندا اے تے نہ استمرار
دے سائز دا؛اس لئی ہر !!a{enumerable} دائرے وچ~$\fn{CH}$
ہمیشہ سچ اے۔ پر اسیں اک ہور جملہ دے سکدے آں جیہڑا ہر ساخت وچ
عین اوس ویلے درست ہووے جدوں استمرار دا مفروضہ جھوٹا ہووے:

\begin{prop}
استمرار دا مفروضہ جھوٹا اے اُدوں تے صرف اُدوں جدوں
\[\fn{NCH} \ident
\lforall[X][(\fn{Cont}(X) \lif \lexists[Y][(Y \subseteq X \land \lnot
    \fn{Count}(Y) \land \lnot \cardeq{X}{Y})])]\]
ہر ساخت وچ درست ہووے۔
\end{prop}
\endgroup
% /OLP-0340

\OLEndChapterHook
\endgroup
% /OLP-0336

\OLEndPartHook
\endgroup
% /OLP-0322

% OLP-0341 content/lambda-calculus/lambda-calculus.tex
\begingroup
\def\olfilename{lambda-calculus}
\expandafter\def\csname ol@sectioncs\endcsname{section}

\olpart{lam}{لامبڈا کیلکولس}

\begin{editorial}
ایہ حصہ لامبڈا کیلکولس بارے اے۔ تعارفی باب Jeremy Avigad دے
نوٹاں اُتے مبنی اے۔ اوہدا کجھ مواد ہُن اگلے باباں وچ وی آ گیا اے۔
نحو، چرچ--روسر خاصیت تے لامبڈا دے ذریعے قابلِ تعریف ہون والے باب
Zesen Qian نے Mitacs دی گرمی دی انٹرن شپ دوران تیار کیتے سن۔
اوہناں دی پڑتال تے نظرثانی ہجے باقی اے۔
\end{editorial}

% OLP-0342 content/lambda-calculus/introduction/introduction.tex
\begingroup
\def\olfilename{introduction}
\expandafter\def\csname ol@sectioncs\endcsname{section}

\olchapter{lam}{int}{تعارف}

\begin{editorial}
ایس باب وچ Jeremy دے لامبڈا کیلکولس بارے اصل مختصر نوٹ نیں۔
ایہدے حصے اکٹھے کرن دی لوڑ اے۔ لامبڈا دے ذریعے قابلِ تعریف ہون بارے
ایتھوں دا مواد اوسے موضوع دے وکھرے، ہور تفصیلی باب نال رلانا باقی اے۔
\end{editorial}

% OLP-0343 content/lambda-calculus/introduction/overview.tex
\begingroup
\def\olfilename{overview}
\expandafter\def\csname ol@sectioncs\endcsname{section}

\olfileid{lam}{int}{ovr}
\olsection{جائزہ}

لامبڈا کیلکولس نوں Alonzo Church نے 1930 دی دہائی دے شروع وچ
تعمیری منطق دی بنیاد لئی تیار کیتا سی،قابلِ حساب فنکشوناں دے نمونے
وَجوں \emph{نہیں}۔ پر چھیتی ایہ وکھایا گیا کہ ایہ قابلِ حسابیت دیاں
ہور تعریفاں،مثلاً ٹیورنگ دے ذریعے قابلِ حساب فنکشوناں تے جزوی بازگشتی
فنکشوناں،دے ہم ارز اے۔ایہ گل کہ شروع وچ ایس نتیجے اُتے تھوڑی حیرانی
ہوئی سی، ایس ہم ارزی نوں ہور دلچسپ بنا دیندی اے۔

لامبڈا لکھت وچ فنکشن دا حوالہ سیدھا اوہنوں متعین کرن والی
علامتی عبارت نال دے سکدے آں، ایس لئی اوہدا وکھرا ناں رکھن دی لوڑ نہیں۔ایہ کہن دی بجائے کہ ``$f$ اوہ فنکشن ہووے جو
$f(x) = x + 3$ نال تعریف ہوندا اے،'' اسی کہہ سکدے آں:
``$f$ فنکشن $\lambd[x][(x + 3)]$ ہووے۔'' ہور لفظاں وچ
$\lambd[x][(x+3)]$ محض اوس فنکشن دا \emph{ناں} اے جیڑھا اپنی
دلیل وچ تین جوڑدا اے۔ایس عبارت وچ $x$ صرف تھاں رکھن والا متغیر
\pnbLatin{(dummy variable)} اے:اوہی فنکشن $\lambd[y][(y + 3)]$ نال وی
ظاہر کیتا جا سکدا اے۔ایہ لکھت دوجے پیرامیٹراں دے ہوندیاں وی کم
کردی اے۔مثلاً فرض کرو $g(x, y)$ دو متغیراں دا فنکشن اے،تے $k$
اک قدرتی عدد اے۔فیر $\lambd[x][g(x,k)]$ اوہ فنکشن اے جیڑھا
ہر~$x$ نوں~$g(x, k)$ ول بھیجدا اے۔

علامتی عبارت توں ایس طرح فنکشن متعین کرن نوں \emph{لامبڈا تجرید
\pnbLatin{(lambda abstraction)}} کہندے نیں۔لامبڈا تجرید دے دوجے پاسے
\emph{اطلاق (application)} اے۔ جے ساڈے کول فنکشن $f$ ہووے
(مثلاً قدرتی عدداں اُتے تعریف شدہ)،تاں اسی اوس نوں 2 ورگی کسے
قدر اُتے \emph{لاگو} کر سکدے آں۔عام لکھت وچ نتیجے لئی~$f(2)$
لکھدے آں۔

لامبڈا تجرید تے اطلاق نوں کٹھے کیتا جاوے تاں کی ہوندا اے؟ جیہڑی دلیل اُتے فنکشن لاگو ہویا اے، اوہنوں تجرید
کیتے متغیر دی تھاں رکھ کے نتیجے والی عبارت سادہ بنائی
جا سکدا اے۔مثلاً
\[
(\lambd[x][(x + 3)])(2)
\]
نوں~$2 + 3$ تک سادہ کیتا جا سکدا اے۔

ایتھوں تک اسی روایتی تصورات لئی صرف نویاں لکھتاں متعارف کروائیاں
نیں۔پر لامبڈا کیلکولس سیٹ نظریاتی زاویے توں ہور بنیادی وکھرا رُخ
لیندا اے۔ایس ڈھانچے وچ:
\begin{enumerate}
\item ہر شے کسی فنکشن نوں ظاہر کردی اے۔
\item فنکشن لامبڈا تجرید نال تعریف کیتے جا سکدے نیں۔
\item ہر شے دا اطلاق ہر ہور شے اُتے ہو سکدا اے۔
\end{enumerate}
مثلاً جے $F$ لامبڈا کیلکولس دا اک ٹرم اے، تاں $F(F)$ نوں
ہمیشہ معنی خیز منیا جاندا اے۔ایس کھلی چھوٹ والے ڈھانچے نوں
\emph{بلا نوع بندی} لامبڈا کیلکولس کہندے نیں؛ایتھے ``بلا نوع بندی''
دا مطلب اے کہ کِنہوں کِس اُتے لاگو کیتا جا سکدا اے،ایس اُتے کوئی
پابندی نہیں۔

\begin{digress}
اک \emph{نوع بندی والا} لامبڈا کیلکولس وی اے، جیڑا بلا نوع بندی
والے روپ دی اہم قسم اے۔بھاویں کئی پاسیاں توں دونے ملدے نیں،پر نوع
بندی والے نوں کلاسیکی سیٹ نظریاتی ڈھانچے نال ملانا بہوں سوکھا اے،
تے اوس دیاں کجھ خاصیتاں کافی وکھریاں نیں۔

لامبڈا کیلکولس اُتے تحقیق نظریاتی کمپیوٹر سائنس تے پروگرامنگ
زباناں دی بناوٹ وچ مرکزی ثابت ہوئی اے۔John McCarthy دی 1950 دی
دہائی وچ بنائی LISP اوہناں مُڈھلیاں زباناں وچوں اک اے جنہاں اُتے ایہناں خیالاں دا اثر پیا۔
\end{digress}
\endgroup
% /OLP-0343

% OLP-0344 content/lambda-calculus/introduction/syntax.tex
\begingroup
\def\olfilename{syntax}
\expandafter\def\csname ol@sectioncs\endcsname{section}

\olfileid{lam}{int}{syn}
\olsection{لامبڈا کیلکولس دی نحو}

شروع وچ متغیراں $x$،$y$،$z$،~\dots دی اک لڑی تے کجھ مستقل
نشاناں $a$،$b$،$c$،~\dots لیندے آں۔ ٹرماں دا سیٹ استقرا نال ایویں متعین کردے آں:
\begin{enumerate}
\item ہر متغیر اک ٹرم اے۔
\item ہر مستقل اک ٹرم اے۔
\item جے $M$ تے $N$ ٹرم نیں،تاں $(MN)$ وی ٹرم اے۔
\item جے $M$ ٹرم اے تے $x$ متغیر اے،تاں $(\lambd[x][M])$ وی اک
  ٹرم اے۔
\end{enumerate}
$(MN)$ شکل والے ٹرماں نوں \emph{اطلاق}،تے
$(\lambd[x][M])$ شکل والے ٹرماں نوں \emph{تجرید} کہندے نیں۔

جیہڑے نظام وچ کوئی مستقل نشان نہ ہووے، اوہنوں \emph{خالص}
لامبڈا کیلکولس کہندے نیں۔اسی زیادہ تر خالص $\lambd$-کیلکولس
وچ کم کراں گے،اس لئی سارے چھوٹے لاطینی حروف متغیراں لئی ہون گے۔
وڈے لاطینی حروف ($M$،$N$،وغیرہ) $\lambd$-کیلکولس دے ٹرماں
لئی ورتاں گے۔

اسیں لکھت دے کجھ رواج ورتاں گے:

\begin{conv}
\begin{enumerate}
\item جدوں قوساں چھڈ دئیے، اطلاق کھبے توں سجے ول جُڑدا اے۔
  مثلاً جے $M$،$N$،$P$،تے $Q$ ٹرم نیں،تاں $MNPQ$،
  $(((MN)P)Q)$ دا مختصر روپ اے۔
\item قوساں چھڈن اُتے لامبڈا تجرید نوں ممکنہ سب توں وڈا دائرۂ
  اثر دتا جاندا اے۔مثلاً $\lambd[x][MNP]$ نوں
  $(\lambd[x][((MN)P)])$ پڑھیا جاندا اے۔
\item اک لامبڈا نال کئی متغیراں دی تجرید لکھی جا سکدی اے۔
  مثلاً $\lambd[xyz][M]$،
  $\lambd[x][\lambd[y][\lambd[z][M]]]$ دا مختصر روپ اے۔
\end{enumerate}
\end{conv}

مثلاً
\[
\lambd[xy][xxyx \lambd[z][xz]]
\]
ایس نال مراد اے:
\[
\lambd[x][\lambd[y][((((xx)y)x)(\lambd[z][(xz)]))]].
\]
ایہ رواج زبانی یاد کر لو۔ شروع وچ ایہ تُہانوں چکرا دین گے۔ عادت
پے جائے تاں ایہ قوساں دے ڈھیر نالوں گھٹ چکران گے۔

دو ٹرم جیہڑے صرف پابند متغیراں دے ناماں وچ فرق رکھدے نیں،اوہناں
نوں $\alpha$-ہم ارز کہندے نیں؛مثلاً $\lambd[x][x]$ تے
$\lambd[y][y]$۔ایہناں نوں ``اکّو'' ٹرم سمجھنا سہولت دی گل اے؛
دوجے لفظاں وچ،جدوں اسی کہندے آں کہ $M$ تے $N$ اکّو نیں،تاں
مراد اے کہ پابند متغیراں دے ناں بدل کے اوہ اکّو نیں۔ متغیر دی
جیہڑی آمد اپنے ناں والے $\lambd$ دے دائرۂ اثر وچ ہووے،اوہ
``پابند'' اے؛باقی ``آزاد'' اے۔پچھلی مثال وچ کوئی آزاد متغیر نہیں؛
پر
\[
(\lambd[z][yz])x
\]
وچ $y$ تے $x$ آزاد نیں،تے $z$ پابند اے۔
\endgroup
% /OLP-0344

% OLP-0345 content/lambda-calculus/introduction/reduction.tex
\begingroup
\def\olfilename{reduction}
\expandafter\def\csname ol@sectioncs\endcsname{section}

\olfileid{lam}{int}{red}
\olsection{لامبڈا ٹرماں دی تقلیل}

لامبڈا ٹرماں نال کی کر سکدے آں؟ اوہناں نوں سادہ بنا سکدے آں۔ جے $M$ تے $N$ کوئی وی لامبڈا ٹرم ہون تے $x$ کوئی وی
متغیر ہووے،تاں $\Subst{M}{N}{x}$ نال اوس نتیجے نوں ظاہر کر
سکدے آں جو $N$ نوں~$x$ دی تھاں~$M$ وچ رکھن توں ملدا اے۔ایہ
کرن توں پہلاں $M$ دے اوہناں پابند متغیراں دے ناں بدل لیندے آں
جنہاں کر کے بدل توں بعد~$N$ دے آزاد متغیر پابند ہو سکدے نیں۔ مثلاً
\[
\Subst{(\lambd[w][xxw])}{yyz}{x} = \lambd[w][(yyz)(yyz)w].
\]

\begin{digress}
اک ٹرم دوجے دی تھاں رکھن لئی $[N/x]M$،$[x/N]M$،تے $M[x/N]$ وی
وکھ وکھ لکھتاں نیں۔خبردار رہو!
\end{digress}

بدیہی طور اُتے $(\lambd[x][M])N$ تے $\Subst{M}{N}{x}$ دا
اکّو مطلب اے؛پہلے ٹرم دی تھاں دوجا رکھن دے عمل نوں
\emph{$\beta$-سکڑاؤ} کہندے نیں۔$(\lambd[x][M])N$ نوں
\emph{ریڈیکس} تے $\Subst{M}{N}{x}$ نوں اوس دا
\emph{کانٹریکٹم} کہندے نیں۔عام طور اُتے جے کسے ٹرم $P$ نوں
~$P'$ وچ اوس دے کسی ذیلی ٹرم دا $\beta$-سکڑاؤ کر کے بدلیا جا
سکدا ہووے،تاں آکھدے آں کہ $P$،\emph{اک قدم وچ $\beta$-تقلیل
نال $P'$ بندا اے}،تے $P \redone P'$ لکھدے آں۔جے $P$ توں
اک قدم دیاں کجھ تقلیلاں نال، جِنہاں دی گنتی صفر وی ہو سکدی اے،
~$P'$ حاصل ہووے،تاں $P$،\emph{$\beta$-تقلیل نال $P'$ بندا
اے}؛نشاناں وچ $P \red P'$۔جیہڑا ٹرم ہور
$\beta$-تقلیل نہ ہو سکے،اوس نوں \emph{$\beta$-ناقابلِ تقلیل}
یا \emph{$\beta$-نارمل} کہندے نیں۔جدوں سیاق صاف ہووے،تاں
``$\beta$-تقلیل'' دی تھاں صرف ``تقلیل'' آکھاں گے،تے اینج
ہی ہور تھاواں اُتے۔

آؤ کجھ مثالاں ویکھئیے۔
\begin{enumerate}
\item ساڈے کول
\begin{align*}
(\lambd[x][xxy]) \lambd[z][z] & \redone (\lambd[z][z])(\lambd[z][z]) y \\
& \redone (\lambd[z][z]) y \\
& \redone y.
\end{align*}
\item کسے ٹرم نوں ``سادہ'' کرنا اوس نوں ہور پیچیدہ وی بنا سکدا اے:
\begin{align*}
(\lambd[x][xxy])(\lambd[x][xxy]) & \redone (\lambd[x][xxy])(\lambd[x][xxy])y \\
& \redone (\lambd[x][xxy])(\lambd[x][xxy])yy \\
& \redone \dots
\end{align*}
\item تقلیل ٹرم نوں جیویں دا تیویں وی چھڈ سکدی اے:
\[
(\lambd[x][xx])(\lambd[x][xx]) \redone (\lambd[x][xx])(\lambd[x][xx]).
\]
\item کجھ ٹرماں دی تقلیل اک توں ودھ طریقیاں نال وی ہو سکدی اے؛
  مثلاً
\[
(\lambd[x][(\lambd[y][yx]) z]) v \redone (\lambd[y][yv]) z
\]
  بیرونی ترین اطلاق نوں سکیڑ کے؛تے
\[
(\lambd[x][(\lambd[y][yx]) z]) v \redone (\lambd[x][zx]) v
\]
  اندرونی ترین اطلاق نوں سکیڑ کے۔پر خیال رکھو کہ ایتھے
  دونے حاصل ٹرم اگے تقلیل ہو کے اکّو ٹرم $zv$ بن دے نیں۔
\end{enumerate}

پچھلی مثال دا آخری نتیجہ اتفاق نہیں؛ایہ لامبڈا کیلکولس دی
اک ڈونگی تے اہم خاصیت نوں دکھاندا اے،جس نوں
«چرچ--روسر خاصیت» کہندے نیں۔
\endgroup
% /OLP-0345

% OLP-0346 content/lambda-calculus/introduction/church-rosser.tex
\begingroup
\def\olfilename{church-rosser}
\expandafter\def\csname ol@sectioncs\endcsname{section}

\olfileid{lam}{int}{cr}
\olsection{چرچ--روسر خاصیت}

\begin{thm}
\ollabel{thm:church-rosser}
فرض کرو $M$، $N_1$ تے $N_2$ ٹرم نیں جِنہاں لئی $M \red N_1$
تے $M \red N_2$۔فیر کوئی ٹرم $P$ موجود اے جیہدے لئی
$N_1 \red P$ تے $N_2 \red P$۔
\end{thm}

\begin{cor}
فرض کرو $M$ دی تقلیل کر کے نارمل روپ حاصل کیتا جا سکدا اے۔
فیر ایہ نارمل روپ اکّو اے۔
\end{cor}

\begin{proof}
جے $M \red N_1$ تے $M \red N_2$،تاں پچھلے قضیے دے مطابق
کوئی ٹرم~$P$ موجود اے جیہدے لئی $N_1$ تے $N_2$ دوویں~$P$ ول تقلیل ہو
سکدے نیں۔جے $N_1$ تے~$N_2$ دوویں نارمل روپ وچ ہون،تاں ایہ
صرف اوس ویلے ہو سکدا اے جدوں $N_1 \ident P \ident N_2$۔
\end{proof}

آخر وچ،دو ٹرماں $M$ تے~$N$ نوں
\emph{$\beta$-ہم ارز}،یا صرف \emph{ہم ارز}،اوس ویلے آکھاں گے
جدوں اوہ تقلیل ہو کے اک سانجھے ٹرم تک پُہنچن۔ دوجے لفظاں وچ، جے کوئی
$P$ ہووے جس لئی $M \red P$ تے $N \red P$۔ایہنوں
$M \equal[\beta] N$ لکھدے نیں۔\olref{thm:church-rosser} دی
مدد نال تُسی پرکھ سکدے او کہ $\equal[\beta]$ ہم ارزی دا اک
تعلق اے،تے اوس وچ ایہ ودھ خاصیت اے کہ ہر $M$ تے~$N$ لئی،
جے $M \red N$ یا $N \red M$،تاں $M \equal[\beta] N$۔
(اصل وچ ایہ وی دکھایا جا سکدا اے کہ $\equal[\beta]$ ایس خاصیت
والے ہم ارزیت دے تعلقاں وچ \emph{سب توں چھوٹا} اے۔)
\endgroup
% /OLP-0346

% OLP-0347 content/lambda-calculus/introduction/currying.tex
\begingroup
\def\olfilename{currying}
\expandafter\def\csname ol@sectioncs\endcsname{section}

\olfileid{lam}{rep}{cur}
\olsection{کریئنگ}

اک $\lambd$-تجرید $\lambd[x][M]$ اک دلیل والا فنکشن ظاہر کردی اے۔ جدوں
اسیں کئی دلیلاں والا فنکشن متعین کرنا چاہیئے، ایہ حد بن جاندی اے۔اک طریقہ ایہ ہو سکدا سی کہ $\lambd$-کیلکولس
نوں جوڑے،تِن-ٹپل وغیرہ بناؤن دی اجازت دے کے ودھایا جائے؛فیر،
مثلاً،تِن رُکنی فنکشن $\lambd[x][M]$ اپنی دلیل وچ اک
تِن-ٹپل دی توقع رکھے گا۔پر ایہ کم \emph{کریئنگ} نال کرنا
زیادہ سہل اے۔

اک مثال ویکھئیے۔اک پل لئی فرض کر لیندے آں کہ $+$ دا عمل
$\lambd$-کیلکولس وچ ساڈے کول اے۔جمع دا فنکشن $2$-رُکنی اے،
یعنی دو دلیلاں لیندا اے۔پر $\lambd$-تجرید سانوں صرف اک
دلیل والے فنکشن دیندی اے:ایس دی نحو
$\lambd[(x,y)][(x+y)]$ ورگے اظہاراں دی اجازت نہیں دیندی۔اس
دے باوجود اسی اک دلیل والا فنکشن~$f_x(y)$ ویکھ سکدے آں جو
$\lambd[y][(x+y)]$ نال دتا گیا اے؛ایہ $x$ نوں اپنی اکلوتی
دلیل~$y$ وچ جوڑدا اے۔اصل وچ ایہ اک فنکشن نہیں،سگوں ``$x$ جوڑو''
والے وکھ وکھ فنکشناں دا خاندان اے،ہر عدد~$x$ لئی اک۔ہن اسی
اک ہور اک دلیل والا فنکشن~$g$ نوں $\lambd[x][f_x]$ دے طور اُتے
متعین کر سکدے آں۔دلیل $x$ اُتے لگن نال $g(x)$ فنکشن~$f_x$
واپس کردا اے—ایس لئی اوس دیاں قدراں آپ فنکشن نیں۔ہن جے
$g$ نوں $x$ اُتے لاگو کر کے حاصل نوں فیر~$y$ اُتے لاگو کریے،
تاں $(g(x))y = f_x(y) = x+y$ ملدا اے۔اینج اک دلیل والا
فنکشن~$g$ دو دلیل والے جمع فنکشن والا کم کر سکدا اے۔
``کریئنگ'' بس دو رُکنی فنکشن نوں اک دلیل والے فنکشوناں وچ بدلّن دی ترکیب دا ناں اے،
جنہاں دیاں قدراں آپ اک دلیل والے فنکشن ہوندے نیں۔

ہُن $\lambd$-کیلکولس دی اصل نحو وچ اک مثال ویکھئیے۔فنکشن
$f(x,y) = x$ نوں کیویں ظاہر کریے؟جے دو دلیلاں قبول کر کے پہلی
واپس کرن والا فنکشن تعریف کرنا ہووے،تاں
$\lambd[x][\lambd[y][x]]$ لکھ سکدے آں۔لفظی طور اُتے ایہ
فنکشن اک دلیل~$x$ قبول کر کے فنکشن~$\lambd[y][x]$ واپس کردا
اے۔فنکشن $\lambd[y][x]$ ہور دلیل~$y$ قبول کردا اے،پر اوس نوں
چھڈ دیندا اے تے ہمیشہ~$x$ واپس کردا اے۔ویکھئیے جدوں اسی
$\lambd[x][\lambd[y][x]]$ نوں دو دلیلاں اُتے لاگو کردے آں
تاں کی ہوندا اے:
\begin{align*}
  (\lambd[x][\lambd[y][x]])MN 
  \bredone & (\lambd[y][M])N \\
  \bredone & M
\end{align*}

عام طور اُتے،$x_1$، \dots،~$x_n$ پیرامیٹراں والے،کسی
ٹرم~$N$ نال تعریف کیتے فنکشن نوں اسی
$\lambd[x_1][\lambd[x_2][\ldots\lambd[x_n][N]]]$ لکھ سکدے
آں۔جے اوس نوں $n$ دلیلاں اُتے لاگو کریے تاں ایہ ملدا اے:
\begin{multline*}
  (\lambd[x_1][\lambd[x_2][\ldots\lambd[x_n][N]]]) M_1 \dots M_n \bredone\\
  \begin{aligned}
  \bredone {} & (\Subst{(\lambd[x_2][\ldots\lambd[x_n][N]])}{M_1}{x_1}) M_2
  \dots M_n\\
   \eqs {} & (\lambd[x_2][\ldots\lambd[x_n][\Subst{N}{M_1}{x_1}]]) M_2
            \dots M_n \\
  \vdots & \\
  \bredone {} &\Subst{\Subst{N}{M_1}{x_1}\ldots}{M_n}{x_n}
  \end{aligned}
\end{multline*}
آخری سطر دا لفظی مطلب ایہ اے کہ فنکشن دی تعریف دے اندرلے ٹرم وچ
$M_i$ نوں $x_i$ دی تھاں رکھیا جائے؛کئی دلیلاں والا فنکشن
لاگو کرن ویلے اسی عین ایہو چاہندے آں۔
\endgroup
% /OLP-0347

% OLP-0348 content/lambda-calculus/introduction/lambda-definability.tex
\begingroup
\def\olfilename{lambda-definability}
\expandafter\def\csname ol@sectioncs\endcsname{section}

\olfileid{lam}{int}{rep}
\olsection{\usetoken{S}{lambda definable} حسابی فنکشن}

لامبڈا کیلکولس حساب دا نمونہ کیویں بن سکدا اے؟شروع وچ تاں
ایس گل دا مطلب وی صاف نہیں ہوندا۔قدرتی اعداد اُتے قابلِ حساب ہون
دی گل کرن لئی سانوں ایہناں اعداد دی مناسب نمائندگی لبھنی
پئے گی۔ ہِٹھاں اک ایہو جیہی نمائندگی اے جیہڑی حیرت انگیز
طور اُتے چنگا کم کردی اے۔

\begin{defn}
ہر قدرتی عدد~$n$ لئی \emph{چرچ عدد} $\num{n}$ نوں لامبڈا
ٹرم $\lambd[x][\lambd[y][(x(x(x(\dots
x(y)))))]]$ دے طور اُتے متعین کرو۔ ایتھے کُل $n$ واری $x$ آؤندا اے۔
\end{defn}

ٹرماں $\num{n}$ دا کم عمل دہراؤنا اے۔ انپٹ $f$ اُتے
$\num{n}$ اوہ فنکشن واپس کردا اے جیہڑا $y$ نوں $f^n(y)$
ول بھیجدا اے۔دھیان دیو کہ ہر چرچ عدد نارمل روپ وچ اے۔ہن
اسی دس سکدے آں کہ لامبڈا ٹرم دا قدرتی اعداد اُتے کسے
فنکشن دا «حساب کرنا» کیہہ معنی رکھدا اے۔

\begin{defn}
فرض کرو $f(x_0, \dots, x_{k-1})$ اک $k$-تھاواں والا جزوی
فنکشن اے جیہڑا $\Nat$ توں $\Nat$ ول جاندا اے۔اسی آکھدے آں کہ اک
$\lambd$-ٹرم~$F$ فنکشن~$f$ نوں \emph{!!{lambda define}s}
جے تے صرف جے قدرتی اعداد دے ہر سلسلے $n_0$،
\dots،~$n_{k-1}$ لئی ہِٹھلی تقلیل
\[
F\, \num{n_0}\, \num{n_1} \dots \num{n_{k-1}} \red \num{f(n_0, n_1, \dots,
  n_{k-1})}
\]
اُدوں ہوندی اے جدوں $f(n_0, \dots, n_{k-1})$
متعین ہووے۔ جے اوہ متعین نہ ہووے تاں
$F\, \num{n_0}\, \num{n_1} \dots \num{n_{k-1}}$
دا کوئی نارمل روپ نہ ہووے۔
\end{defn}

\begin{thm}
\ollabel{thm:lambda-def}
اک فنکشن $f$ جزوی قابلِ حساب اے جے تے صرف جے اوہ کسے
لامبڈا ٹرم نال !!{lambda defined} اے۔
\end{thm}

\begin{explain}
ایہ قضیہ کجھ حیران کردا اے۔ حساب دے نمونے دے طور اُتے
لامبڈا کیلکولس کافی سادہ کیلکولس اے؛ایس دے صرف دو عمل
نیں:لامبڈا تجرید تے اطلاق!{} پر ایہناں محدود وسیلیاں
نال وی ہر حسابی طریقۂ کار عمل وچ لیاندا جا سکدا اے۔
\end{explain}
\endgroup
% /OLP-0348

% OLP-0349 content/lambda-calculus/introduction/lambda-computable.tex
\begingroup
\def\olfilename{lambda-computable}
\expandafter\def\csname ol@sectioncs\endcsname{section}

\olfileid{lam}{int}{cmp}
\olsection{\usetoken{S}{lambda definable} فنکشن قابلِ حساب نیں}

\begin{thm}
\ollabel{thm:lambda-computable}
جے اک جزوی فنکشن $f$ کسے لامبڈا ٹرم نال !!{lambda defined}
ہووے، تاں اوہ قابلِ حساب اے۔
\end{thm}

\begin{proof}
فرض کرو فنکشن~$f$ لامبڈا ٹرم~$X$ نال !!{lambda defined}
اے۔ آؤ $f$ دا حساب کرن دا اک غیر رسمی طریقہ بیان کریے۔
انپٹ $m_0$، \dots،~$m_{n-1}$ ملن اُتے ٹرم
$X \num{m_0} \ldots \num{m_{n-1}}$ لکھ لو۔اک درخت بناؤ:
پہلاں اصل ٹرم دیاں ساریاں اک قدم دیاں تقلیلاں لکھو۔ اوہناں دے
ہِٹھاں اوہناں وچوں ہر اک دیاں ساریاں اک قدم دیاں تقلیلاں لکھو
(یعنی اصل ٹرم دیاں دو-قدم تقلیلاں)؛فیر ایہ کم جاری رکھو۔
جے کدے کوئی چرچ عدد لبھ پئے تاں اوہی جواب دے طور اُتے واپس کرو۔
جے کدے نہ لبھے تاں تلاش جاری رہندی اے تے فنکشن اوس انپٹ اُتے متعین نہیں۔

چرچ دے دعوے دی ورتوں نال آکھ سکدے آں کہ ایہ فنکشن قابلِ
حساب اے۔قضیہ ثابت کرن دا بہتر طریقہ ایہ ہووے گا کہ ایس
تلاش دی بازگشتی وضاحت دتی جائے۔مثلاً،بنیادی بازگشتی فنکشناں
تے تعلقاں دا اک سلسلہ تعریف کیتا جا سکدا اے:
``$\fn{IsASubterm}$''، ``$\fn{Substitute}$''،
``$\fn{ReducesToInOneStep}$''، ``$\fn{ReductionSequence}$''،
``$\fn{Numeral}$'' وغیرہ۔$f(m_0, \dots, m_{n-1})$
دا حساب کرن دا جزوی بازگشتی طریقہ فیر
$X \num{m_0} \dots \num{m_{n-1}}$ توں شروع ہو کے
کسے چرچ عدد اُتے ختم ہون والی اک-قدم تقلیلاں دی لڑی
لبھدا اے،تے اوس چرچ عدد دے مطابق عدد واپس کردا اے۔
تفصیل لمبی تے اکتاؤنی اے،پر باقی کم معمول دے مطابق اے۔
\end{proof}
\endgroup
% /OLP-0349

% OLP-0350 content/lambda-calculus/introduction/computable-lambda.tex
\begingroup
\def\olfilename{computable-lambda}
\expandafter\def\csname ol@sectioncs\endcsname{section}

\olfileid{lam}{int}{lrp}
\olsection{قابلِ حساب فنکشن \usetoken{S}{lambda definable} نیں}

\begin{thm}
\ollabel{thm:computable-lambda}
ہر جزوی قابلِ حساب فنکشن !!{lambda definable} اے۔
\end{thm}

\begin{proof}
سانوں دکھانا اے کہ ہر جزوی قابلِ حساب فنکشن~$f$ کسے
لامبڈا ٹرم~$F$ نال !!{lambda defined} اے۔ کلینی دے
نارمل روپ دے قضیے دے مطابق ایہ دکھانا کافی اے کہ ہر
بنیادی بازگشتی فنکشن کسے لامبڈا ٹرم نال !!{lambda defined}
اے۔ فیر ایہ دکھاؤنا اے کہ !!{lambda definable} فنکشن موزوں
ترکیباں تے بغیر مقررہ حد دی تلاش ہِٹھ بند نیں۔ہر بنیادی بازگشتی
فنکشن نوں کسے لامبڈا ٹرم نال !!{lambda defined} دکھاؤن
لئی ایہ دکھانا کافی اے کہ ابتدائی فنکشن !!{lambda definable}
نیں،تے اوہ جزوی فنکشن جیہڑے !!{lambda definable} نیں
ترکیب، بنیادی بازگشت تے بغیر مقررہ حد دی تلاش ہِٹھ بند نیں۔
\end{proof}

باقی ثبوت پڑھنا سہل بناؤن لئی اسیں زیادہ رائج لکھت ورتاں گے۔
مثلاً،$M(x, y, z)$ نوں $Mxyz$ دی تھاں لکھاں گے۔ایہ لکھت
بھاویں اک اشارہ دیندی اے،پر یاد رکھو کہ بلا نوع بندی لامبڈا
کیلکولس دے ٹرم دیاں دلیلاں دی کوئی مقرر گنتی نہیں ہوندی۔ ایس لئی اوسے ٹرم~$M$ لئی $M(x,y)$ تے
$M(x,y,z,w)$ لکھنا وی اکّو طرح جائز اے۔پر ایہ لکھت
دسدی اے کہ اسیں $M$ نوں تِن متغیراں والا فنکشن سمجھ کے ورت رہے آں۔
ایس نال تعریفاں دا مطلب وی صاف رہندا اے۔
ایس طرح اسی ``$M$ نوں $M(x,y,z) = \dots$ نال تعریف کرو''
آکھاں گے،بجائے ایس دے کہ ``$M$ نوں
$M = \lambd[x][\lambd[y][\lambd[z][\dots]]]$ نال تعریف کرو''۔
\endgroup
% /OLP-0350

% OLP-0351 content/lambda-calculus/introduction/basic-pr-lambda.tex
\begingroup
\def\olfilename{basic-pr-lambda}
\expandafter\def\csname ol@sectioncs\endcsname{section}

\olfileid{lam}{int}{bas}
\olsection{مُڈھلے بنیادی بازگشتی فنکشن \usetoken{S}{lambda definable} نیں}

\begin{lem}
فنکشن $\Zero$، $\Succ$ تے $\Proj{n}{i}$، !!{lambda definable} نیں۔
\end{lem}

\begin{proof}
$\Zero$ تاں بس
$\lambd[x][\lambd[y][y]]$ اے۔

جانشین فنکشن~$\Succ$ دی تعریف $\fn{Succ}(u) = \lambd[x][\lambd[y][x(uxy)]]$ اے۔
سوچو ایہ کیوں کم کردا اے۔ ہر چرچ عدد $\num{n}$ نوں وار وار عمل کرن
والا ٹرم سمجھو۔ کسے وی فنکشن~$f$ لئی $\fn{Succ}(\num{n},f)$
اوہ فنکشن اے جیہڑا انپٹ~$y$ لئی $f$ دا $n$ واری اطلاق کردا اے،
شروعات~$y$ توں کر کے۔ فیر اک واری ہور لاگو کردا اے۔

دِتّیاں دلیلاں وچوں اک دلیل چُنن والے فنکشوناں لئی کوئی ہور گل نہیں۔
$\fn{Proj}^n_i(x_0, \dots,
x_{n-1}) = x_i$۔ دوجے لفظاں وچ ساڈے لکھت دے رواج مطابق $\fn{Proj}^n_i$
لامبڈا ٹرم $\lambd[x_0][\dots \lambd[x_{n-1}][x_i]]$ اے۔
\end{proof}
\endgroup
% /OLP-0351

% OLP-0352 content/lambda-calculus/introduction/composition.tex
\begingroup
\def\olfilename{composition}
\expandafter\def\csname ol@sectioncs\endcsname{section}

\olfileid{lam}{int}{com}
\olsection{\usetoken{S}{lambda definable} فنکشن ترکیب ہِٹھ بند نیں}

\begin{lem}
!!{lambda definable} فنکشن ترکیب ہِٹھ بند نیں۔
\end{lem}

\begin{proof}
فرض کرو $f$ دی تعریف $h$، $g_0,$
\dots،~$g_{k-1}$ دی ترکیب نال ہوئی اے۔ منّ لئیے کہ $h$، $g_0$،
\dots،~$g_{k-1}$ بالترتیب $H$، $G_0$، \dots،~$G_{k-1}$ نال
!!{lambda defined} نیں۔ سانوں ٹرم~$F$ لبھنا اے جیہڑا~$f$ نوں
!!{lambda define}s۔ اسیں~$F$ نوں سیدھا ایویں متعین کر سکدے آں:
\[
F(x_0, \dots, x_{l-1}) = H(G_0(x_0, \dots, x_{l-1}),
\dots, G_{k-1}(x_0, \dots, x_{l-1})).
\]
دوجے لفظاں وچ لامبڈا کیلکولس دی زبان فنکشوناں دی ترکیب ظاہر کرن
لئی موزوں اے۔
\end{proof}
\endgroup
% /OLP-0352

% OLP-0353 content/lambda-calculus/introduction/primitive-recursion.tex
\begingroup
\def\olfilename{primitive-recursion}
\expandafter\def\csname ol@sectioncs\endcsname{section}

\olfileid{lam}{int}{pr}
\olsection{\usetoken{S}{lambda definable} فنکشن بنیادی بازگشت ہِٹھ بند نیں}

ہُن بنیادی بازگشت لئی سانوں کجھ کم کرنا پئے گا۔ اسیں ایہ کم
پڑاؤواں وچ کراں گے۔ پہلاں وانگ منّ لئیے کہ ساڈے کول ٹرم $G$ تے $H$
پہلاں ای نیں، جیہڑے فنکشن $g$ تے~$h$ نوں بالترتیب !!{lambda define} کردے نیں۔
سانوں اوہ ٹرم $H$ چاہیدا اے جیہڑا ہِٹھ دِتّے فنکشن~$f$ نوں
!!{lambda define}s:
\begin{align*}
f(0, \vec z) & = g(\vec z) \\
f(x+1, \vec z) & = h(z, f(x,\vec z), \vec z).
\end{align*}
سو عام طور اُتے، جے لامبڈا ٹرم $G'$ تے~$H'$ دِتّے ہون، تاں
ایہو جیہا ٹرم~$F$ لبھنا کافی اے جیہدے لئی
\begin{align*}
F(\num{0}, \vec z) & \equiv G(\vec z) \\
F(\overline{n+1}, \vec z) & \equiv H(\num{n}, F(\num{n}, \vec z), \vec z)
\end{align*}
ہر قدرتی عدد~$n$ لئی ہووے۔ $G'$ تے~$H'$ فنکشن $g$ تے~$h$ نوں
!!{lambda define} کردے نیں، ایس لئی جدوں وی اسیں چرچ عدد $\num{\vec m}$
نوں $\vec z$ دی تھاں رکھاں گے، $F(\num{n+1}, \num{\vec m})$
تقلیل ہو کے صحیح جواب والے نارمل روپ تک پُہنچے گا۔

پر ایس کم لئی ایہو جیہا ٹرم~$F$ لبھنا کافی اے جیہڑا ایہ شرطاں
پوریاں کرے:
\begin{align*}
F(\num 0) & \equiv G \\
F(\overline {n+1}) & \equiv H(\num{n},F(\num{n}))
\intertext{ہر قدرتی عدد $n$ لئی، جتھے}
G & = \lambd[\vec z][G'(\vec z)] \text{ تے}\\
H(u,v) & = \lambd[\vec z][H'(u,v(u,\vec z),\vec z)].
\end{align*}
دوجے لفظاں وچ، لامبڈا دی ایس ترکیب نال وادھو پیرامیٹراں $\vec z$
دی وکھری فکر نہیں کرنی پیندی۔ اوہ لامبڈا لکھت وچ ای سما جاندے نیں۔

ٹرم $F$ دی تعریف توں پہلاں سانوں ترتیب وار جوڑیاں نوں ورتن دا
طریقہ چاہیدا اے۔ اگلا لمہ ایہ طریقہ دیندا اے۔

\begin{lem}
اک لامبڈا ٹرم $D$ موجود اے جیہدے لئی لامبڈا ٹرماں $M$ تے~$N$
دے ہر جوڑے اُتے $D(M,N)(\num{0}) \red M$ تے $D(M,N)(\num{1}) \red N$۔
\end{lem}

\begin{proof}
پہلاں لامبڈا ٹرم $K$ دی تعریف ایویں کرو:
\[
K(y) = \lambd[x][y].
\]
دوجے لفظاں وچ، $K$ اوہ ٹرم $\lambd[y][\lambd[x][y]]$ اے۔ ہور
طرح ویکھیے تاں ہر $M$ لئی، $K(M)$ اک مستقل فنکشن اے جیہڑا
ہر انپٹ اُتے~$M$ واپس کردا اے۔

ہُن $D(x,y,z)$ دی تعریف $D(x,y,z) = z (K(y))x$ نال کرو۔ فیر
ساڈے کول
\begin{align*}
D(M,N,\num 0) & \red \num 0 (K(N)) M \red M \text{ تے}\\
D(M,N,\num 1) & \red \num 1 (K(N)) M \red K(N) M \red N,
\end{align*}
جیویں سانوں چاہیدا سی۔
\end{proof}

خیال ایہ اے کہ $D(M,N)$ جوڑے $\tuple{M, N}$ دی نمائندگی کردا اے۔
جے منّ لئیے کہ $P$ ایہو جیہے جوڑے دی نمائندگی کردا اے، تاں $P(\num 0)$
تے $P(\num 1)$ بالترتیب کھبے تے سجے حصے نوں چُندے نیں، یعنی
$(P)_0$ تے $(P)_1$۔ اسیں ایہ آخری دو لکھتاں ورتاں گے۔

\begin{lem}
!!{lambda definable} فنکشن بنیادی بازگشت ہِٹھ بند نیں۔
\end{lem}

\begin{proof}
سانوں ایہ دکھاؤنا اے کہ کوئی وی ٹرم $G$ تے $H$ دِتّے ہون، اسیں
ایہو جیہا ٹرم $F$ لبھ سکدے آں جیہدے لئی
\begin{align*}
F(\num{0}) & \equiv G \\
F(\num{n+1}) & \equiv H(\num{n}, F(\num{n}))
\end{align*}
ہر قدرتی عدد~$n$ لئی ہووے۔ موٹا خیال ایہ اے کہ چرچ اعداں نوں
وار وار عمل کرن والے ٹرم سمجھ کے \emph{جوڑیاں} دیاں لڑیّاں دا حساب کریے:
\[
\tuple{\num 0, F(\num 0)}, \tuple{\num 1, F(\num 1)}, \dots,
\]
دھیان دیو کہ پہلا جوڑا بس $\tuple{\num 0, G}$ اے۔ جے جوڑا
$\tuple{\num{n}, F(\num{n})}$ دِتّا ہووے، تاں اگلے جوڑے
$\tuple{\num{n+1}, F(\num{n+1})}$ نوں
$\tuple{\num{n+1}, H(\num{n}, F(\num{n}))}$ دے ہم ارز ہونا چاہیدا اے۔
اسیں اک لامبڈا ٹرم~$T$ بناواں گے جیہڑا ایہ اک قدم دا بدل کردا اے۔

تفصیل ایہ اے۔ $T(u)$ دی تعریف ایویں کرو:
\[
T(u) = \tuple{S((u)_0), H((u)_0,(u)_1)}.
\]
ہُن ایہ پرکھنا سوکھا اے کہ کسے وی عدد~$n$ لئی،
\[
T(\tuple{\num{n}, M}) \red \tuple{\num{n+1}, H(\num{n}, M)}.
\]
جیویں اُتے دَسّیا سی، $G$ تے $H$ دِتّے ہون تاں~$F(u)$ دی تعریف
ایویں کرو:
\[
F(u) = (u(T,\tuple{\num 0, G}))_1.
\]
دوجے لفظاں وچ، انپٹ~$\num{n}$ اُتے، $F$، $T$ نوں $n$ واری
$\tuple{\num 0, G}$ توں شروع کر کے لاگو کردا اے، تے فیر دوجا جزو
واپس کردا اے۔ شروع وچ ساڈے کول ایہ اے:
\begin{enumerate}
\item $\num{0} (T, \tuple{\num 0, G}) \equiv \tuple{\num{0}, G}$
\item $F(\num{0}) \equiv G$
\end{enumerate}
~$n$ اُتے استقرا نال اسیں ہر قدرتی عدد لئی ایہ دکھا سکدے آں:
\begin{enumerate}
\item $\num{n+1}(T, \tuple{\num{0}, G}) \equiv \tuple{\num{n+1}, F(\num{n+1})}$
\item $F(\num{n+1}) \equiv H(\num{n}, F(\num{n}))$
\end{enumerate}
دوجی گل لئی ساڈے کول ایہ اے:
\begin{align*}
F(\num{n+1}) & \red (\num{n+1}(T, \tuple{\num 0, G}))_1 \\
& \equiv (T(\num{n} (T, \tuple{\num 0, G})))_1 \\
& \equiv (T(\tuple{\num{n}, F(\num{n})}))_1 \\
& \equiv (\tuple{\num{n+1}, H(\num{n}, F(\num{n}))})_1 \\
& \equiv H(\num{n}, F(\num{n})).
\end{align*}
ایتھے اسیں آخری توں پہلی سطر وچ استقرا دا مفروضہ ورتیا اے۔ پہلی
گل لئی ساڈے کول ایہ اے:
\begin{align*}
\num{n+1} (T, \tuple{\num 0, G}) &
\equiv T(\num{n} (T, \tuple{\num 0, G})) \\
& \equiv T( \tuple{\num{n}, F(\num{n})}) \\
& \equiv \tuple{\num{n+1}, H(\num{n}, F(\num{n}))} \\
& \equiv \tuple{\num{n+1}, F(\num{n+1})}.
\end{align*}
ایتھے اسیں آخری سطر وچ دوجی گل ورتی اے۔ سو اسیں دکھا دِتّا اے کہ
$F(\num 0) \equiv G$ تے ہر $n$ لئی $F(\num {n+1}) \equiv H(\num
n, F(\num{n}))$۔ سانوں عین ایہو چاہیدا سی۔
\end{proof}
\endgroup
% /OLP-0353

% OLP-0354 content/lambda-calculus/introduction/fixed-point-combinator.tex
\begingroup
\def\olfilename{fixed-point-combinator}
\expandafter\def\csname ol@sectioncs\endcsname{section}

\olfileid{lam}{int}{fix}
\olsection{مقرر نقطے دے کمبینیٹر}

فرض کرو تُہاڈے کول لامبڈا ٹرم $g$ اے تے تُہانوں اک ہور ٹرم $k$
چاہیدا اے جیہدے لئی $k$ دی $\beta$-ہم ارزی $gk$ نال ہووے۔
ساڈی لکھت دے رواج ورت کے ایہ ٹرم متعین کرو:
\[
\fn{diag}(x) = xx
\]
تے
\[
l(x) = g(\fn{diag}(x))
\]
دوجے لفظاں وچ، $l$ اوہ ٹرم $\lambd[x][g(xx)]$ اے۔ منّ لو $k$
اوہ ٹرم $ll$ اے۔ فیر ساڈے کول ایہ اے:
\begin{align*}
k & = (\lambd[x][g(xx)])(\lambd[x][g(xx)]) \\
& \red  g((\lambd[x][g(xx)])(\lambd[x][g(xx)])) \\
& = gk.
\end{align*}
جے اسیں
\[
Y = \lambd[g][((\lambd[x][g(xx)])(\lambd[x][g(xx)]))]
\]
لئیے، تاں $Yg$ تے $g(Yg)$ تقلیل ہو کے اکّو ٹرم تک پُہنچدے نیں۔
سو $Yg \equiv_\beta g(Yg)$۔ ایہنوں «کری دا کمبینیٹر» آکھدے نیں۔
جے اس دی تھاں
\[
Y = (\lambd[xg][g(xxg)])(\lambd[xg][g(xxg)])
\]
لئیے، تاں $Yg$ دی تقلیل نال خود $g(Yg)$ ملدا اے۔ ایہ ودھ مضبوط
بیان اے۔ $Y$ دے ایس روپ نوں «ٹیورنگ دا کمبینیٹر» آکھدے نیں۔
\endgroup
% /OLP-0354

% OLP-0355 content/lambda-calculus/introduction/minimization.tex
\begingroup
\def\olfilename{minimization}
\expandafter\def\csname ol@sectioncs\endcsname{section}

\olfileid{lam}{int}{min}
\olsection{\usetoken{S}{lambda definable} فنکشن گھٹ توں گھٹ حل لبھن ہِٹھ بند نیں}

\begin{lem}
منّ لو $f(x,y)$، !!{lambda definable} اے۔ $g$ دی تعریف ایویں ہووے:
\[
g(x) \simeq \umin{y}{f(x,y)}.
\]
فیر $g$، !!{lambda definable} اے۔
\end{lem}

\begin{proof}
موٹا خیال ایہ اے۔ $x$ دِتّا ہووے تاں اسیں مقرر نقطے والا لامبڈا
ٹرم $Y$ ورت کے فنکشن $h_x(n)$ متعین کراں گے، جیہڑا کوئی~$y$
لبھن لئی~$n$ توں تلاش شروع کردا اے۔ فیر $g(x)$ بس $h_x(0)$ اے۔
فنکشن~$h_x$ نوں مقرر نقطے والی مساوات دے حل دے طور اُتے لکھ سکدے آں:
\[
h_x(n) \simeq
\begin{cases}
n & \text{جے $f(x,n) = 0$} \\
h_x(n+1) & \text{ہور صورت وچ۔}
\end{cases}
\]

تفصیل ایہ اے۔ کیوں جے $f$ بنیادی بازگشتی اے، ایہ کسے ٹرم $F$
نال !!{lambda defined} اے۔ یاد کرو ساڈے کول اک لامبڈا ٹرم~$D$
وی اے جیہدے لئی $D(M, N, \bar{0}) \red M$ تے $D(M, N, \bar{1})
\red N$۔ فی الحال $x$ نوں مقرر رکھدے آں۔ $h_x$ نوں !!{lambda define}
کرن لئی سانوں اک ٹرم~$H$ لبھنا اے جیہڑا~$x$ اُتے منحصر ہووے
تے ایہ شرط پوری کرے:
\[
H(\num{n}) \equiv D(\num{n}, H(S(\num{n})), F(x, \num{n})).
\]
ایہ کم اسیں مقرر نقطے والے ٹرم~$Y$ نال کر سکدے آں۔ پہلاں منّ لو
$U$ ایہ ٹرم اے:
\[
\lambd[h][\lambd[z][D(z,(h(Sz)),F(x,z))]],
\]
تے فیر منّ لو $H$ اوہ ٹرم~$YU$ اے۔ دھیان دیو کہ~$H$ دا اکّو
آزاد متغیر~$x$ اے۔ ہُن دکھائیے کہ $H$ اُپرلی مساوات پوری کردا اے۔

$Y$ دی تعریف مطابق ساڈے کول ایہ اے:
\[
H = YU \equiv U(YU) = U(H).
\]
خاص طور اُتے، ہر قدرتی عدد~$n$ لئی،
\begin{align*}
H(\num{n}) & \equiv U(H, \num{n}) \\
& \red D(\num{n}, H(S(\num{n})), F(x, \num{n})),
\end{align*}
جیویں سانوں چاہیدا سی۔ دھیان دیو کہ جے آخری سطر وچ چرچ عدد
$\num{m}$ نوں~$x$ دی تھاں رکھیے، تاں عبارت دی تقلیل نال
$\num{n}$ ملدا اے جے $F(\num{m},
\num{n})$ دی تقلیل نال $\num{0}$ ملے۔ تے جے
$H(S(\num{n}))$ ول تقلیل دی گل کریے، تاں اوہ اُدوں ہوندی اے جدوں
$F(\num{m}, \num{n})$ دی تقلیل نال کوئی ہور چرچ عدد ملے۔

ثبوت پورا کرن لئی منّ لو $G$ اوہ ٹرم $\lambd[x][H(\num 0)]$
اے۔ فیر $G$، $g$ نوں !!{lambda define}s۔ دوجے لفظاں وچ، ہر~$m$
لئی $G(\num m)$ دی تقلیل نال~$\overline {g(m)}$ ملدا اے جے
$g(m)$ متعین ہووے۔ ہور صورت وچ کوئی نارمل روپ نہیں۔
\end{proof}
\endgroup
% /OLP-0355

\OLEndChapterHook
\endgroup
% /OLP-0342

% OLP-0356 content/lambda-calculus/syntax/syntax.tex
\begingroup
\def\olfilename{syntax}
\expandafter\def\csname ol@sectioncs\endcsname{section}

\olchapter{lam}{syn}{نحو}

% OLP-0357 content/lambda-calculus/syntax/terms.tex
\begingroup
\def\olfilename{terms}
\expandafter\def\csname ol@sectioncs\endcsname{section}

\olfileid{lam}{syn}{trm}
\olsection{ٹرم}

لامبڈا کیلکولس دے ٹرم استقرا نال بنائے جاندے نیں۔ ایہناں لئی
لامتناہی متغیر $\Obj{v_0}$، $\Obj{v_1}$، \dots، نشان~«$\lambd$»،
تے قوساں ورتدے آں۔ اسیں $x$، $y$، $z$، \dots{} نوں متغیراں
لئی تے $M$، $N$، $P$، \dots{} نوں ٹرماں لئی ورتاں گے۔

\begin{defn}[ٹرم] \ollabel{defn:term}
لامبڈا کیلکولس دے \emph{ٹرماں} دا سیٹ استقرا نال ایویں متعین کردے آں:
\begin{enumerate}
  \item \ollabel{defn:term-var} جے $x$ متغیر اے، تاں $x$ اک ٹرم اے۔
  \item \ollabel{defn:term-abs} جے $x$ متغیر اے تے $M$ ٹرم اے،
    تاں $(\lambd[x][M])$ اک ٹرم اے۔
  \item \ollabel{defn:term-app} جے $M$ تے $N$ دوویں ٹرم نیں،
    تاں $(MN)$ اک ٹرم اے۔
\end{enumerate}
\end{defn}

جے ٹرم $(\lambd[x][M])$ قاعدے \olref{defn:term-abs} نال
بنیا اے، تاں اسیں آکھدے آں کہ ایہ \emph{تجرید} دا نتیجہ اے۔
$x$ جیہڑا $\lambd[x]$ وچ اے، اوہنوں \emph{!!{parameter}} آکھدے آں۔
ٹرم $(MN)$ جیہڑا قاعدے \olref{defn:term-app} نال بنیا اے،
\emph{اطلاق} دا نتیجہ اے۔

اُپر دِتّی تعریف وچ ٹرماں دیاں ساریاں قوساں لکھی جاندیاں نیں۔
ایہ لکھت بھاری ہو سکدی اے، جیویں ٹرم
$(\lambd[x][((\lambd[x][x])(\lambd[x][(xx)]))])$ توں دِسدا اے۔
اسیں قوساں چھڈن دے رواج اگے دساں گے۔ رسمی تعریف نال پر ایہ
لبھنا سوکھا اے کہ قاعدیاں \olref{defn:term} مطابق ٹرم کیویں بنیا۔
مثال لئی، ٹرم $(\lambd[x][((\lambd[x][x])(\lambd[x][(xx)]))])$
بناؤن دا آخری قدم لازماً تجرید اے، جیہدا !!{parameter}~$x$ اے۔
تجرید ٹرم $((\lambd[x][x])(\lambd[x][(xx)]))$ توں شروع ہوندی اے،
جیہڑا دو ٹرماں دا اطلاق اے۔ اوہناں وچوں ہر ٹرم آپ تجرید دا
نتیجہ اے۔ ایس طرح اسیں اگے ودھ سکدے آں۔

\begin{prob}
دسو کہ $(\lambd[g][(\lambd[x][(g (x x))])
  (\lambd[x][(g (x x))])])$ کیویں بنیا اے۔
\end{prob}
\endgroup
% /OLP-0357

% OLP-0358 content/lambda-calculus/syntax/unique-readability.tex
\begingroup
\def\olfilename{unique-readability}
\expandafter\def\csname ol@sectioncs\endcsname{section}

\olfileid{lam}{syn}{unq}
\olsection{اکّو طرح پڑھیا جانا}

اسیں پُچھ سکدے آں کہ کی ہر ٹرم بناؤن دا اکّو طریقہ اے۔ جواب ہاں
اے۔ ہر لامبڈا ٹرم بناؤن تے سمجھّن دا اکّو طریقہ اے۔ اگلی گل بات وچ
\emph{بناوٹ} توں مراد \olref[trm]{defn:term} دے قاعدے اک یا کئی
وار ورت کے ٹرم بناؤن دی کاروائی اے۔

\begin{lem}\ollabel{lem:term-start}
ٹرم متغیر یا قوس نال شروع ہوندا اے۔
\end{lem}

\begin{proof}
کوئی شے ٹرم اُدوں ای منّی جاندی اے جدوں اوہ \olref[trm]{defn:term}
مطابق بنی ہووے۔ جے ایہ \olref[trm]{defn:term-var} دا نتیجہ اے، تاں
لازماً متغیر اے۔ جے ایہ \olref[trm]{defn:term-abs} یا
\olref[trm]{defn:term-app} دا نتیجہ اے، تاں قوس نال شروع ہوندا اے۔
\end{proof}

\begin{lem}\ollabel{lem:app-start}
اطلاق دا نتیجہ یا تاں دو قوساں نال، یا اک قوس تے اک متغیر نال
شروع ہوندا اے۔
\end{lem}

\begin{proof}
جے $M$ اطلاق دا نتیجہ اے، تاں اوہدی شکل $(PQ)$ اے، سو اوہ قوس نال
شروع ہوندا اے۔ کیوں جے $P$ اک ٹرم اے، ایس لئی \olref{lem:term-start}
مطابق اوہ قوس یا متغیر نال شروع ہوندا اے۔
\end{proof}

\begin{lem}\ollabel{lem:initial}
کسے ٹرم دا کوئی اوہ مُڈھلا حصہ جیہڑا پورے ٹرم توں چھوٹا ہووے،
آپ ٹرم نہیں ہو سکدا۔
\end{lem}

\begin{prob}
ٹرماں دی لمبائی اُتے استقرا نال \olref[lam][syn][unq]{lem:initial}
ثابت کرو۔
\end{prob}

\begin{prop}[اکّو طرح پڑھیا جانا] \ollabel{prop:unq}
ہر ٹرم دی بناوٹ اکّو اے۔ دوجے لفظاں وچ، جے ٹرم $M$ کسے بناوٹ نال
بنیا اے، تاں اوہو اکّو بناوٹ اے جیہڑی ایہ ٹرم بنا سکدی اے۔
\end{prop}

\begin{proof}
اسیں ٹرماں دی بناوٹ اُتے استقرا نال ثابت کردے آں۔
\begin{enumerate}
  \item $M$ دی شکل $x$ اے، جتھے $x$ کوئی متغیر اے۔ کیوں جے
    تجرید تے اطلاق دے نتیجے ہمیشہ قوساں نال شروع ہوندے نیں،
    اوہناں نال $M$ نہیں بن سکدا۔ سو $M$ دی بناوٹ لازماً
    \olref[trm]{defn:term}\olref[trm]{defn:term-var} دا اکّو قدم اے۔
  \item $M$ دی شکل $(\lambd[x][N])$ اے، جتھے $x$ کوئی متغیر
    تے $N$ اک ٹرم اے۔ ایہ \olref[trm]{defn:term}\olref[trm]{defn:term-var}
    نال نہیں بن سکدا، کیوں جے ایہ اکّلا متغیر نہیں۔
    \olref{lem:app-start} مطابق ایہ اطلاق دا نتیجہ وی نہیں۔
    سو $M$ صرف~$N$ دی تجرید نال بنیا ہو سکدا اے۔ استقرا دے مفروضے
    توں پتا اے کہ $N$ دی اپنی بناوٹ وی اکّو اے۔
  \item $M$ دی شکل $(PQ)$ اے، جتھے $P$ تے $Q$ ٹرم نیں۔
    کیوں جے ایہ قوس نال شروع ہوندا اے، ایس لئی
    \olref[trm]{defn:term}\olref[trm]{defn:term-var} نال نہیں
    بن سکدا۔ \olref{lem:term-start} مطابق $P$، $\lambd$ نال
    شروع نہیں ہو سکدا، سو $(PQ)$ تجرید دا نتیجہ نہیں ہو سکدا۔
    ہُن منّ لو $M$ اطلاق نال کسے ہور طریقے وی بن سکدا اے،
    مثلاً ایہ $(P'Q')$ دی شکل دا وی اے۔ فیر $P$، $P'$ دا
    پورے توں چھوٹا مُڈھلا حصہ اے، یا ایہدا اُلٹ اے۔ پر
    \olref{lem:initial} مطابق ایہ ممکن نہیں۔ سو $P$ تے $Q$
    اکّو طرح متعین نیں، تے استقرا دے مفروضے توں پتا اے کہ
    $P$ تے $Q$ دی بناوٹ وی اکّو اے۔
\end{enumerate}
\end{proof}

اُپرلے قضیے نوں سوکھے لفظاں وچ ایویں آکھ سکدے آں:
\begin{prop}
  ٹرم $M$ دی شکل صرف ہِٹھلیاں وچوں اک ہو سکدی اے:
  \begin{enumerate}
    \item $x$، جتھے $x$ اوہ متغیر اے جیہڑا $M$ توں اکّو طرح متعین اے۔
    \item $(\lambd[x][N])$، جتھے $x$ متغیر اے تے $N$ دوجا ٹرم۔
      دوویں $M$ توں اکّو طرح متعین نیں۔
    \item $(PQ)$، جتھے $P$ تے $Q$ دو ٹرم نیں جیہڑے $M$ توں
      اکّو طرح متعین نیں۔
  \end{enumerate}
\end{prop}
\endgroup
% /OLP-0358

% OLP-0359 content/lambda-calculus/syntax/abbreviated-syntax.tex
\begingroup
\def\olfilename{abbreviated-syntax}
\expandafter\def\csname ol@sectioncs\endcsname{section}

\olfileid{lam}{syn}{abb}
\olsection{مختصر نحو}

\olref[trm]{defn:term} وچ تعریف کیتے ٹرم کدے لکھنے اوکھے
ہوندے نیں، ایس لئی ہور مختصر نحو ورتنی سوکھی اے۔ پر ایہ دھیان
رکھنا اے کہ مختصر لکھت وچ وی ہر ٹرم اکّو طرح پڑھیا جا سکے۔
وڈی پدھر اُتے ورتی جاندی اک صورت \emph{مختصر ٹرم} اے، جیہڑی
ایویں لکھی جاندی اے۔

\begin{enumerate}
\item جدوں قوساں چھڈ دئیے، اطلاق کھبے توں سجے ول جُڑدا اے۔
  مثلاً جے $M$، $N$، $P$ تے~$Q$ ٹرم نیں، تاں $MNPQ$،
  $(((MN)P)Q)$ دا مختصر روپ اے۔
\item قوساں چھڈن اُتے لامبڈا تجرید دا دائرۂ اثر جِنّا ہو سکے
  وڈا رکھدے آں۔ مثلاً $\lambd[x][MNP]$ نوں
  $(\lambd[x][MNP])$ پڑھدے آں۔
\item اک لامبڈا نال کئی متغیراں دی تجرید لکھ سکدے آں۔
  مثلاً $\lambd[xyz][M]$،
  $\lambd[x][\lambd[y][\lambd[z][M]]]$ دا مختصر روپ اے۔
\end{enumerate}

مثلاً،
\[
\lambd[xy][xxyx \lambd[z][xz]]
\]
ایہ ہِٹھلی عبارت دا مختصر روپ اے:
\[
(\lambd[x][(\lambd[y][((((xx)y)x)(\lambd[z][(xz)]))])]).
\]

\begin{prob}
مختصر ٹرم $\lambd[g][(\lambd[x][g (x x)])
  \lambd[x][g (x x)]]$ نوں پورا لکھو۔
\end{prob}
\endgroup
% /OLP-0359

% OLP-0360 content/lambda-calculus/syntax/free-variables.tex
\begingroup
\def\olfilename{free-variables}
\expandafter\def\csname ol@sectioncs\endcsname{section}

\olfileid{lam}{syn}{fv}
\olsection{آزاد متغیر}

لامبڈا کیلکولس فنکشوناں بارے اے، تے فنکشن لامبڈا تجرید نال
بندے نیں۔ بدیہی خیال ایہ اے کہ $\lambd[x][M]$ اوہ فنکشن اے
جیہدیاں قدراں~$M$ توں ملدیاں نیں، جدوں فنکشن دی دلیل
نوں~$x$ دی قدر بنایا جاوے۔ پر~$x$ دی~$M$ وچ ہر آمد متعلقہ
نہیں۔ جے $M$ وچ اک ہور تجرید $\lambd[x][N]$ اے، تاں $x$
دیاں~$N$ وچلیاں آمدواں $\lambd[x][N]$ نال متعلق نیں،
$\lambd[x][M]$ نال نہیں۔ سو لامبڈا تجرید $\lambd[x]$،
$\lambd[x][M]$ دے اندر، $x$ دیاں~$M$ وچ اوہناں آمدواں
نوں \emph{پابند کردی اے} جیہڑیاں پہلاں توں کسے ہور
لامبڈا تجرید نال پابند نہیں، یعنی~$x$ دیاں~$M$ وچ
\emph{آزاد} آمدواں۔

\begin{defn}[دائرۂ اثر]
جے $\lambd[x][M]$ ٹرم~$N$ دے اندر آوے، تاں~$N$ دی متعلقہ
آمد~$\lambd[x]$ دا \emph{دائرۂ اثر} اے۔
\end{defn}

\begin{defn}[آزاد تے پابند آمد]
متغیر $x$ دی ٹرم~$M$ وچ آمد \emph{آزاد} اے جے اوہ کسے
$\lambd[x]$ دے دائرۂ اثر وچ نہیں، تے ہور صورت وچ \emph{پابند} اے۔
متغیر~$x$ دی $\lambd[x][M]$ وچ آمد مُڈھلے $\lambd[x]$ نال
اُدوں تے صرف اُدوں پابند اے جدوں $x$ دی اوہ آمد~$M$ وچ آزاد اے۔
\end{defn}

\begin{ex}
$\lambd[x][x y]$ وچ $x$ تے $y$ دوویں $\lambd[x]$ دے دائرۂ
اثر وچ نیں، سو $x$، $\lambd[x]$ نال پابند اے۔ کیوں جے $y$
کسے $\lambd[y]$ دے دائرۂ اثر وچ نہیں، اوہ آزاد اے۔
$\lambd[x][x x]$ وچ $x$ دیاں دوویں آمدواں~$\lambd[x]$ نال پابند
نیں، کیوں جے دوویں $xx$ وچ آزاد نیں۔ $((\lambd[x][xx])x)$ وچ~$x$
دی آخری آمد آزاد اے، کیوں جے اوہ کسے $\lambd[x]$ دے دائرۂ اثر وچ
نہیں۔ $\lambd[x][(\lambd[x][x])x]$ وچ پہلے $\lambd[x]$ دا دائرۂ
اثر $(\lambd[x][x])x$ اے، تے دوجے $\lambd[x]$ دا دائرۂ اثر~$x$
دی آخری توں پہلی آمد اے۔ $(\lambd[x][x])x$ وچ $x$ دی آخری آمد آزاد
اے، تے آخری توں پہلی پابند اے۔ ایس لئی $x$ دی آخری توں پہلی آمد،
$\lambd[x][(\lambd[x][x])x]$ وچ دوجے $\lambd[x]$ نال پابند اے،
تے آخری آمد پہلے $\lambd[x]$ نال پابند اے۔
\end{ex}

ٹرم $P$ وچ متغیراں دیاں ساریاں آمدواں پرکھ کے اسیں آزاد
متغیراں دا سیٹ لبھ سکدے آں۔ ایہ سیٹ $\FV{P}$ نال لکھدے آں، تے
اوہدی قدرتی تعریف ایہ اے:

\begin{defn}[ٹرم دے آزاد متغیر] \ollabel{def:fv}
ٹرم دے \emph{آزاد متغیراں} دا سیٹ استقرا نال ایویں متعین کردے آں:
\begin{enumerate}
  \item $\FV{x} = \{x\}$ \ollabel{def:fv1}
  \item $\FV{\lambd[x][N]} = \FV{N} \setminus \{x\}$ \ollabel{def:fv2}
  \item $\FV{PQ} = \FV{P} \cup \FV{Q}$ \ollabel{def:fv3}
\end{enumerate}
\end{defn}

\begin{prob}
\begin{enumerate}
  \item $\lambd[g]$ تے دوواں $\lambd[x]$ دے دائرۂ اثر ایس ٹرم وچ لبھو:
    $\lambd[g][(\lambd[x][g (x x)]) \lambd[x][g (x x)]]$۔
  \item $\lambd[g][(\lambd[x][g (x x)]) \lambd[x][g (x x)]]$ وچ
    کی متغیراں دیاں ساریاں آمدواں پابند نیں؟ ہر اک نوں کیہڑی تجرید
    پابند کردی اے؟
  \item $\FV{\lambd[x][(\lambd[y][(\lambd[z][x y]) z]) y]}$ دسو۔
\end{enumerate}
\end{prob}

\begin{explain}
آزاد متغیر باہرلی دنیا، یعنی \emph{ماحول}، ول اک حوالے وانگ اے۔
جیہڑے ٹرم وچ آزاد متغیر ہون، اوہنوں ادھورا متعین ٹرم سمجھ سکدے آں،
کیوں جے اوہدا ورتاؤ ماحول نوں کیویں مقرر کردے آں، ایس اُتے منحصر اے۔
مثال لئی، ٹرم $\lambd[x][f x]$ اک دلیل~$x$ لیندا اے تے اوس دلیل
اُتے $f$ دا نتیجہ واپس کردا اے۔ ایتھے متغیر~$f$ آزاد اے۔ ٹرم دی
ایہ قدر اوہدے ماحول اُتے، خاص کر اوس ماحول وچ $f$ دی قدر اُتے،
منحصر اے۔

جے ایس ٹرم اُتے تجرید لائیے، تاں $\lambd[f][\lambd[x][f
    x]]$ ملدا اے۔ ایہ ٹرم ہُن ماحول دے متغیر $f$ اُتے منحصر نہیں۔
ہُن ایہ اوہ فنکشن ظاہر کردا اے جیہڑا دو دلیل لیندا اے تے پہلی
نوں دوجی اُتے لا کے نتیجہ واپس کردا اے۔ ماحول وچ $f$ نوں بدلن
نال ایس ٹرم دے ورتاؤ اُتے کوئی اثر نہیں پئے گا، کیوں جے ٹرم
صرف اوہو قدر ورتے گا جیہڑی دلیل دے طور اُتے دِتّی جاوے گی۔
ماحول وچ $f$ دی قدر نہیں ورتے گا۔
\end{explain}

\begin{defn}[بند ٹرم، کمبینیٹر]
جیہڑے ٹرم وچ کوئی آزاد متغیر نہ ہووے، اوہنوں \emph{بند ٹرم}
یا \emph{کمبینیٹر} آکھدے آں۔
\end{defn}

\begin{lem}\ollabel{lem:fv}
\begin{enumerate}
  \item \ollabel{lem:fv-abs} جے $y \neq x$، تاں $y \in
    \FV{\lambd[x][N]}$ اُدوں تے صرف اُدوں جدوں $y \in \FV{N}$۔
  \item \ollabel{lem:fv-app} $y \in \FV{PQ}$ اُدوں تے صرف اُدوں
    جدوں $y \in \FV{P}$ یا $y \in \FV{Q}$۔
\end{enumerate}
\end{lem}

\begin{proof}
مشق۔
\end{proof}

\begin{prob}
\olref[lam][syn][fv]{lem:fv} ثابت کرو۔
\end{prob}
\endgroup
% /OLP-0360

% OLP-0361 content/lambda-calculus/syntax/substitution.tex
\begingroup
\def\olfilename{substitution}
\expandafter\def\csname ol@sectioncs\endcsname{section}

\olfileid{lam}{syn}{sub}
\olsection{تھانواں بدلنا}

\begin{explain}
آزاد متغیر ماحول دے متغیراں ول حوالے نیں۔ ایس لئی آزاد متغیر
دی تھاں کوئی خاص قدر ورتنی سمجھ آؤن والی گل اے۔ مثلاً اسیں
$f$ نوں $\lambd[x][f x]$ وچ کسے خاص ٹرم نال بدل سکدے آں،
جیویں عین اوہی قدر واپس دین والا فنکشن $\lambd[y][y]$۔
ایس توں $\lambd[x][(\lambd[y][y]) x]$ ملدا اے۔ آزاد
متغیراں دی تھاں لامبڈا ٹرم رکھن نوں تھانواں بدلنا آکھدے آں۔
\end{explain}

\begin{defn}[تھانواں بدلنا] \ollabel{defn:substitution}
  ٹرم $N$ نوں متغیر $x$ دی تھاں ٹرم $M$ وچ رکھن والے \emph{بدل}، یعنی
  $\Subst{M}{N}{x}$، دی استقرائی تعریف ایہ اے۔
  \begin{enumerate}
    \item $\Subst{x}{N}{x} = N$. \ollabel{defn:substitution-1}
    \item $\Subst{y}{N}{x}  = y$ جے $x \neq y$۔ \ollabel{defn:substitution-2}
    \item $\Subst{PQ}{N}{x}  = (\Subst{P}{N}{x}) (\Subst{Q}{N}{x})$.
      \ollabel{def:substitution-3}
    \item $\Subst{(\lambd[y][P])}{N}{x} = \lambd[y][\Subst{P}{N}{x}]$،
      جے $x \neq y$ تے $y \notin \FV{N}$۔ ہور صورت وچ ایہ غیر معین اے۔
      \ollabel{defn:substitution-4}
  \end{enumerate}
\end{defn}

\begin{explain}
\olref{defn:substitution}\olref{defn:substitution-4} وچ اسیں
$x
\neq y$ دی شرط لاندے آں، کیوں جے اسیں متغیر~$x$ دیاں
ٹرم~$M$ وچلیاں \emph{پابند} آمدواں دی تھاں~$N$ نہیں رکھنا چاہندے۔
مثلاً، جے اسیں بدل $\Subst{\lambd[x][x]}{y}{x}$ کڈھیے، تاں
نتیجہ $\lambd[x][y]$ دی تھاں بس $\lambd[x][x]$ ہونا چاہیدا اے۔

جدوں $N$ نوں~$x$ دی تھاں $\lambd[y][P]$ وچ رکھیے، تاں
$y \notin \FV{N}$ دی شرط وی لاندے آں۔ مثلاً، اسیں $y$
نوں $x$ دی تھاں $\lambd[y][x]$ وچ نہیں رکھ سکدے، یعنی
$\Subst{\lambd[y][x]}{y}{x}$ نہیں کر سکدے، کیوں جے ایس
توں $\lambd[y][y]$ ملے گا۔ ایہ ٹرم اوہ فنکشن ظاہر کردا اے
جیہڑا دلیل قبول کردا اے تے اوہو دلیل سیدھی واپس دیندا اے۔
پر ٹرم $\lambd[y][x]$ اوہ فنکشن ظاہر کردا اے جیہڑا ہمیشہ
ٹرم~$x$ واپس دیندا اے، یا اوہ شے جیہدے ول $x$ دا حوالہ اے۔
سو سانوں اصل وچ اوہ فنکشن چاہیدا اے جیہڑا دلیل قبول کر کے
اوہنوں چھڈ دیندا اے تے ماحول دا متغیر~$y$ واپس دیندا اے۔
ایہ صحیح طرح کرن لئی پہلاں پابند متغیر~$y$ دا ناں بدلنا پئے گا۔
\end{explain}

\begin{prob}
  ہِٹھلے بدلاں دا نتیجہ کیہ اے؟
  \begin{enumerate}
  \item $\Subst{\lambd[y][x(\lambd[w][vwx])]}{(uv)}{x}$
  \item $\Subst{\lambd[y][x(\lambd[x][x])]}{(\lambd[y][xy])}{x}$
  \item $\Subst{y(\lambd[v][xv])}{(\lambd[y][vy])}{x}$
  \end{enumerate}
\end{prob}

\begin{thm} \ollabel{thm:notinfv}
  جے $x \notin \FV{M}$، تاں $\FV{\Subst{M}{N}{x}} = \FV{M}$،
  بشرطیکہ مساوات دا کھبا پاسا معین ہووے۔
\end{thm}

\begin{proof}
  $M$ دی بناوٹ اُتے استقرا نال۔
  \begin{enumerate}
  \item $M$ اک متغیر اے۔ مشق۔
  \item $M$ دی شکل $(PQ)$ اے۔ مشق۔
  \item $M$ دی شکل $\lambd[y][P]$ اے۔ کیوں جے
    $\Subst{\lambd[y][P]}{N}{x}$ معین اے، ایس لئی ایہ
    $\lambd[y][\Subst{P}{N}{x}]$ ہونا چاہیدا اے۔ فیر
    $\Subst{P}{N}{x}$ وی معین ہونا چاہیدا اے۔ نالے $x \neq y$
    تے $x \notin \FV{Q}$۔ فیر۔
    \begin{multline*}
      \FV{\Subst{\lambd[y][P]}{N}{x}} = \\
    \begin{aligned}
      & = \FV{\lambd[y][\Subst{P}{N}{x}]} &&
      \text{\olref{defn:substitution-4} مطابق}\\
      & = \FV{\Subst{P}{N}{x}} \setminus \{y\} &&
      \text{\olref[fv]{def:fv}\olref[fv]{def:fv2} مطابق}\\
      & = \FV{P} \setminus \{y\} && \text{استقرائی مفروضے مطابق} \\
      & = \FV{\lambd[y][P]} && \text{\olref[fv]{def:fv}\olref[fv]{def:fv2} مطابق}
    \end{aligned}
    \end{multline*}
  \end{enumerate}
\end{proof}

\begin{prob}
  \olref[lam][syn][sub]{thm:notinfv} دا ثبوت پورا کرو۔
\end{prob}

\begin{thm} \ollabel{thm:infv}
  جے $x \in \FV{M})$، تاں $\FV{\Subst{M}{N}{x}} = (\FV{M} \setminus
  \{x\}) \cup \FV{N}$، بشرطیکہ مساوات دا کھبا پاسا معین ہووے۔
\end{thm}

\begin{proof}
  $M$ دی بناوٹ اُتے استقرا نال۔
  \begin{enumerate}
  \item $M$ اک متغیر اے۔ مشق۔
  \item $M$ دی شکل $PQ$ اے۔ کیوں جے $\Subst{(PQ)}{N}{y}$ معین
    اے، ایس لئی ایہ $(\Subst{P}{N}{x})(\Subst{Q}{N}{x})$
    ہونا چاہیدا اے، تے دوویں بدل معین نیں۔ نالے کیوں جے
    $x \in \FV{PQ}$، ایس لئی $x \in \FV{P}$ یا $x \in \FV{Q}$
    یا دوویں۔ باقی ثبوت مشق لئی چھڈدے آں۔
  \item $M$ دی شکل $\lambd[y][P]$ اے۔ کیوں جے
    $\Subst{\lambd[y][P]}{N}{x}$ معین اے، ایس لئی ایہ
    $\lambd[y][\Subst{P}{N}{x}]$ ہونا چاہیدا اے، جتھے
    $\Subst{P}{N}{x}$ معین اے، $x \neq y$ تے $y \notin \FV{N}$۔
    نالے کیوں جے $y \in
    \FV{\lambd[x][P]}$، ایس لئی $y \in \FV{P}$ وی اے۔ ہُن۔
    \begin{multline*}
      \FV{\Subst{(\lambd[y][P])}{N}{x}} = \\
      \begin{aligned}
      & = \FV{\lambd[y][\Subst{P}{N}{x}]} \\
      & = \FV{\Subst{P}{N}{x}} \setminus \{y\} \\
      & = ((\FV{P} \setminus \{y\}) \cup (\FV{N} \setminus \{x\})
       && \text{استقرائی مفروضے مطابق} \\
      & = (\FV{P} \setminus \{x, y\}) \cup \FV{N}
       && x \notin \FV{N} \\
      & = (\FV{\lambd[y][P]} \setminus \{x\}) \cup \FV{N}
      \end{aligned}
      \end{multline*}
  \end{enumerate}
\end{proof}

\begin{prob}
  \olref[lam][syn][sub]{thm:infv} دا ثبوت پورا کرو۔
\end{prob}

\begin{thm}\ollabel{thm:clr}
  $x \notin \FV{\Subst{M}{N}{x}}$، جے سجا پاسا معین ہووے
  تے $x \notin \FV{N}$۔
\end{thm}

\begin{proof}
  مشق۔
\end{proof}

\begin{prob}
  \olref[lam][syn][sub]{thm:clr} ثابت کرو۔
\end{prob}

\begin{thm}\ollabel{thm:inv}
  جے $\Subst{M}{y}{x}$ معین اے تے $y \notin \FV{M}$،
  تاں $\Subst{\Subst{M}{y}{x}}{x}{y} = M$۔
\end{thm}

\begin{proof}
  $M$ دی بناوٹ اُتے استقرا نال۔
  \begin{enumerate}
  \item $M$ اک متغیر $z$ اے۔ مشق۔
  \item $M$ دی شکل $(PQ)$ اے۔ فیر۔
    \begin{align*}
      \Subst{\Subst{(PQ)}{y}{x}}{x}{y}
      &=\Subst{((\Subst{P}{y}{x})(\Subst{Q}{y}{x}))}{x}{y} \\
      &= (\Subst{\Subst{P}{y}{x}}{x}{y})(\Subst{\Subst{Q}{y}{x}}{x}{y}) \\
      &= (PQ) \text{ استقرائی مفروضے مطابق}
    \end{align*}
  \item $M$ دی شکل $\lambd[z][N]$ اے۔ کیوں جے
    $\Subst{\lambd[z][N]}{y}{x}$ معین اے، ایس لئی سانوں پتا اے
    کہ $z \neq y$۔ سو۔
    \begin{align*}
      \Subst{\Subst{(\lambd[z][N])}{y}{x}}{x}{y}\\
      & = \Subst{(\lambd[z][\Subst{N}{y}{x}])}{x}{y} \\
      & = \lambd[z][\Subst{\Subst{N}{y}{x}}{x}{y}] \\
      &= \lambd[z][N] \text{ استقرائی مفروضے مطابق}
    \end{align*}
  \end{enumerate}
\end{proof}

\begin{prob}
  \olref[lam][syn][sub]{thm:inv} دا ثبوت پورا کرو۔
\end{prob}
\endgroup
% /OLP-0361

% OLP-0362 content/lambda-calculus/syntax/alpha.tex
\begingroup
\def\olfilename{alpha}
\expandafter\def\csname ol@sectioncs\endcsname{section}

\olfileid{lam}{syn}{alp}
\olsection{\texorpdfstring{$\alpha$-روپ بدلنا}{α-روپ بدلنا}}

$\lambd[x][x]$ تے $\lambd[y][y]$ دا آپس وچ کیہ تعلق اے؟
دوویں عین اوہی قدر واپس دین والا فنکشن ظاہر کردے نیں۔ ظاہر اے کہ
نحو دے لحاظ نال ایہ وکھرے ٹرم نیں۔ ایہناں وچ فرق بس پابند
متغیر دے ناں دا اے۔ اک ٹرم دوجے وچ پابند متغیر دا ناں بدل کے
ملدا اے۔ ایس نوں \emph{$\alpha$-روپ بدلنا} آکھدے آں۔

\begin{defn}[پابند متغیر دا ناں بدلنا، $\aconvone$]
  جے ٹرم $M$ وچ $\lambd[x][N]$ دی اک آمد اے، $y \notin
  \FV{N}$ اے، تے $\Subst{N}{y}{x}$ معین اے، تاں ایس آمد دی تھاں
  \begin{equation*}
    \lambd[y][\Subst{N}{y}{x}]
  \end{equation*}
  رکھ کے $M'$ بناؤن نوں \emph{پابند متغیر دا ناں بدلنا} آکھدے
  آں۔ اینوں $M \redone[\alpha] M'$ لکھدے آں۔
\end{defn}

\begin{defn}[تعلق دی موافقت]
  ٹرماں اُتے تعلق $R$ نوں \emph{موافق} آکھدے آں جے ایہ
  ہِٹھلیاں شرطاں پوریاں کردا اے۔
  \begin{enumerate}
  \item جے $R N N'$، تاں $R \lambd[x][N] \lambd[x][N']$
  \item جے $R P P'$، تاں $R (PQ) (P'Q)$
  \item جے $R Q Q'$، تاں $R (PQ) (PQ')$
  \end{enumerate}
\end{defn}

سو تعریف نوں ہور طرح بیان کریے۔
\begin{defn}[پابند متغیر دا ناں بدلنا، $\aconvone$]
  \emph{پابند متغیر دا ناں بدلنا} ($\redone[\alpha]$) ٹرماں
  اُتے سبھ توں چھوٹا موافق تعلق اے جیہڑا ہِٹھلی شرط پوری کردا اے۔
  \begin{align*}
    &\lambd[x][N] \redone[\alpha] \lambd[y][\Subst{N}{y}{x}] && \text{جے
      $x \neq y$، $y \notin \FV{N}$} \\
    & &&\text{تے $\Subst{N}{y}{x}$ معین اے}
  \end{align*}
\end{defn}

ایتھے «سبھ توں چھوٹا» دا مطلب اے کہ تعلق وچ صرف اوہ جوڑے
نیں جیہڑے موافقت تے ودھیک شرط پوری کرن لئی ضروری نیں۔ سو ایس
تعلق دی تعریف ایویں وی ہو سکدی اے:

\begin{defn}[پابند متغیر دا ناں بدلنا، $\aconvone$] \ollabel{defn:aconvone}
  \emph{پابند متغیر دا ناں بدلنا} ($\aconvone$) استقرا نال
  ایویں متعین ہوندا اے۔
  \begin{enumerate}
  \item جے $N \aconvone N'$، تاں $\lambd[x][N] \aconvone
    \lambd[x][N']$۔ \ollabel{defn:aconvone1}
  \item جے $P \aconvone P'$، تاں $(PQ) \aconvone (P'Q)$۔ \ollabel{defn:aconvone2}
  \item جے $Q \aconvone Q'$، تاں $(PQ) \aconvone (PQ')$۔ \ollabel{defn:aconvone3}
  \item جے $x \neq y$، $y \notin \FV{N}$ تے $\Subst{N}{y}{x}$ معین اے، تاں
    $\lambd[x][N] \redone[\alpha] \lambd[y][\Subst{N}{y}{x}]$۔
    \ollabel{defn:aconvone4}
  \end{enumerate}
\end{defn}

دوویں تعریفاں ہم ارز نیں۔ ایہدا ثبوت اسیں مشق لئی چھڈدے آں۔
اگّوں اسیں استقرائی تعریف ورتاں گے۔

\begin{defn}[$\alpha$-روپ بدلنا، $\aconv$]
  \emph{$\alpha$-روپ بدلنا} ($\aconv$) ٹرماں اُتے سبھ توں چھوٹا
  انعکاسی تے تعدّی والا تعلق اے جیہدے اندر $\aconvone$ شامل اے۔
\end{defn}

اُتے وانگ، «سبھ توں چھوٹا» دا مطلب اے کہ تعلق وچ بس اوہی
جوڑے نیں جیہڑے تعدّی تے $\aconvone$ کر کے ضروری نیں۔ ایس توں
ہِٹھلی ہم ارز تعریف ملدی اے:

\begin{defn}[$\alpha$-روپ بدلنا، $\aconv$] \ollabel{defn:aconv}
  \emph{$\alpha$-روپ بدلنا} ($\aconv$) استقرا نال ایویں متعین ہوندا اے۔
  \begin{enumerate}
  \item جے $P \aconv Q$ تے $Q \aconv R$، تاں $P \aconv R$۔
    \ollabel{defn:aconv1}
  \item جے $P \aconvone Q$، تاں $P \aconv Q$۔ \ollabel{defn:aconv2}
  \item $P \aconv P$۔ \ollabel{defn:aconv3}
  \end{enumerate}
\end{defn}

\begin{ex}
  $\lambd[x][f x]$ دا $\alpha$-روپ بدل کے $\lambd[y][f
  y]$ ملدا اے، تے اُلٹا وی۔ سادہ لفظاں وچ، دوویں ایہو جیہے
  فنکشن نیں جیہڑے اک دلیل لیندے نیں تے اوہ دلیل اُتے $f$ لاؤن دا
  نتیجہ دیندے نیں۔ ایتھے حوالہ ماحول دے متغیر~$f$ دا اے۔

  $\lambd[x][f x]$ دا $\alpha$-روپ بدل کے $\lambd[x][g
    x]$ نہیں ملدا۔ سادہ لفظاں وچ، ایہ بالترتیب ماحول دے متغیراں
  $f$ تے~$g$ ول اشارہ کردے نیں، ایس لئی ایہ وکھرے فنکشن نیں۔
  جیہڑے ماحول وچ $f$ تے $g$ وکھرے ہون، اوتھے ایہناں دا ورتاؤ
  وی وکھرا اے۔
\end{ex}

\begin{prob}
  کی ہِٹھلے ٹرماں دے جوڑیاں وچ $\alpha$-روپ بدلیا جا سکدا اے؟
  \begin{enumerate}
  \item $\lambd[x][\lambd[y][x]]$ تے $\lambd[y][\lambd[x][y]]$
  \item $\lambd[x][\lambd[y][x]]$ تے $\lambd[c][\lambd[b][a]]$
  \item $\lambd[x][\lambd[y][x]]$ تے $\lambd[c][\lambd[b][a]]$
  \end{enumerate}
\end{prob}

\begin{lem}\ollabel{lem:fv-one}
  جے $P \aconvone Q$، تاں $\FV{P} = \FV{Q}$۔
\end{lem}

\begin{proof}
  $P \aconvone Q$ دے !!{derivation} اُتے استقرا نال۔
  \begin{enumerate}
  \item جے آخری قاعدہ \olref{defn:aconvone4} اے، تاں $P$ دی
    شکل $\lambd[x][N]$ تے $Q$ دی شکل
    $\lambd[y][\Subst{N}{y}{x}]$ اے، جتھے $x \neq y$، $y \notin \FV{N}$
    تے $\Subst{N}{y}{x}$ معین اے۔ اسیں $x \in \FV{N}$ دے
    ہون یا نہ ہون مطابق صورتاں وکھ کردے آں۔
    \begin{enumerate}
    \item جے $x \in FV(N)$، تاں۔
      \begin{align*}
        \FV{\lambd[y][\Subst{N}{y}{x}]} &=
        \FV{\Subst{N}{y}{x}} \setminus \{y\} \\
        & = ((\FV{N} \setminus \{x\}) \cup \{y\}) \setminus \{y\}
         && \text{ \olref[sub]{thm:infv} مطابق} \\
        & = \FV{N} \setminus \{x\} \\
        & = \FV{\lambd[x][N]}
      \end{align*}
    \item جے $x \notin FV(N)$، تاں۔
      \begin{align*}
        \FV{\lambd[y][\Subst{N}{y}{x}]}
        & = FV{\Subst{N}{y}{x}} \setminus \{y\} \\
        & = \FV{N} \setminus \{x\}
         && \text{\olref[sub]{thm:notinfv} مطابق} \\
        & = \FV{\lambd[x][N]}.
      \end{align*}
    \end{enumerate}
  \item باقی تِن صورتاں مشق لئی چھڈدے آں۔
  \end{enumerate}
\end{proof}

\begin{prob}
  \olref[lam][syn][alp]{lem:fv-one} دا ثبوت پورا کرو۔
\end{prob}

\begin{lem}\ollabel{lem:inv}
  جے $P \aconvone Q$، تاں $Q \aconvone P$۔
\end{lem}

\begin{proof}
  $P \aconvone Q$ دے !!{derivation} اُتے استقرا۔
  \begin{enumerate}
  \item جے آخری قاعدہ \olref{defn:aconvone4} اے، تاں $P$ دی
    شکل $\lambd[x][N]$ تے $Q$ دی شکل
    $\lambd[y][\Subst{N}{y}{x}]$ اے، جتھے $x \neq y$، $y \notin \FV{N}$
    تے $\Subst{N}{y}{x}$ معین اے۔ پہلاں، $y \notin
    \FV{\Subst{N}{y}{x}}$، \olref[sub]{thm:clr} مطابق۔
    \olref[sub]{thm:inv} توں ملدا اے کہ $\Subst{\Subst{N}{y}{x}}{x}{y}$
    معین اے تے~$N$ دے برابر وی اے۔ فیر \olref{defn:aconvone4}
    مطابق $\lambd[y][\Subst{N}{y}{x}]
    \aconvone \lambd[x][\Subst{\Subst{N}{y}{x}}{x}{y}] =
    \lambd[x][N]$۔
  \end{enumerate}
\end{proof}

\begin{prob}
  \olref[lam][syn][alp]{lem:inv} دا ثبوت پورا کرو۔
\end{prob}

\begin{thm}
  $\alpha$-روپ بدلنا ٹرماں اُتے اک ہم ارزیتی تعلق اے، یعنی
  ایہ انعکاسی، متقابل تے تعدّی والا اے۔
\end{thm}

\begin{proof}
  \begin{enumerate}
  \item ہر ٹرم $M$ لئی، $M$ نوں $M$ وچ \emph{صفر} پابند ناں بدل کے
    بدلیا جا سکدا اے۔
  \item جے $P$ دا $\alpha$-روپ پابند ناں بدلن دے اک سلسلے نال
    $Q$ بنے، تاں $Q$ توں اسیں ایہناں بدلاں نوں
    \olref{lem:inv} مطابق اُلٹی ترتیب وچ واپس کر کے~$P$ حاصل
    کر سکدے آں۔
  \item جے $P$ دا $\alpha$-روپ پابند ناں بدلن دے اک سلسلے نال
    $Q$، تے $Q$ دا دوجے سلسلے نال $R$ بنے، تاں پہلا سلسلہ
    تے فیر دوجا سلسلہ ورت کے $P$ نوں $R$ وچ بدل سکدے آں۔
  \end{enumerate}
\end{proof}

اگّوں اسیں آکھاں گے کہ $M$ تے $N$ \emph{$\alpha$-ہم ارز}
نیں، $M
\aeq N$، جے تے صرف جے $M$ دا $\alpha$-روپ بدل کے $N$ ملے۔
جیویں ہُن دکھایا اے، ایہ اُدوں تے صرف اُدوں ہوندا اے جدوں $N$
دا $\alpha$-روپ بدل کے~$M$ ملے۔

\begin{thm}\ollabel{thm:fv}
  جے $M \aeq N$، تاں $\FV{M} = \FV{N}$۔
\end{thm}

\begin{proof}
  \olref{lem:fv-one} توں سیدھا ملدا اے۔
\end{proof}

\begin{lem}\ollabel{lem:sub:R}
  جے $R \aeq R'$ تے $\Subst{M}{R}{y}$ معین اے، تاں
  $\Subst{M}{R'}{y}$ معین اے تے $\alpha$-ہم ارز اے
  $\Subst{M}{R}{y}$ نال۔
\end{lem}

\begin{proof}
  مشق۔
\end{proof}

\begin{prob}
  \olref[lam][syn][alp]{lem:sub:R} ثابت کرو۔
\end{prob}

یاد کرو کہ \olref[sub]{sec} وچ کجھ صورتاں لئی بدل غیر معین اے۔
پر ٹرماں دا $\alpha$-روپ بدل کے اسیں پابند متغیراں دے ناں
ایویں بدل سکدے آں کہ بدل ہمیشہ معین ہووے۔ نتیجہ $\alpha$-ہم ارزی
قائم رکھدا اے، جیویں ایہ قضیہ دسدا اے:

\begin{thm}\ollabel{thm:sub}
  ہر $M$، $R$ تے~$y$ لئی، کوئی $M'$ موجود اے جیہدے لئی $M \aeq M'$
  تے $\Subst{M'}{R}{y}$ معین اے۔ نالے جے کوئی ہور جوڑا
  $M'' \aeq M$ تے $R''$ اے جیہدے لئی $\Subst{M''}{R''}{y}$
  معین اے تے $R'' \aeq R$، تاں $\Subst{M'}{R}{y} \aeq \Subst{M''}{R''}{y}$۔
\end{thm}

\begin{proof}
  $M$ دی بناوٹ اُتے استقرا نال۔
  \begin{enumerate}
  \item $M$ اک متغیر~$z$ اے۔ مشق۔
  \item فرض کرو $M$ دی شکل $\lambd[x][N]$ اے۔ ایہو جیہا متغیر
    $z$ چُنو جیہڑا $x$ تے $y$ توں وکھرا ہووے تے جیہدے لئی
    $z \notin \FV{N}$ تے $z
    \notin \FV{R}$۔ استقرائی مفروضے مطابق کوئی $N'$ اے جیہدے
    لئی $N' \aeq N$ تے $\Subst{N'}{z}{x}$ معین اے۔ فیر
    $\lambd[x][N] \aeq \lambd[x][N']$ وی اے، جیویں
    \olref{defn:aconvone}\olref{defn:aconvone1} توں ملدا اے۔
    ہُن $\lambd[x][N']
    \aeq \lambd[z][\Subst{N'}{z}{x}]$ اے،
    \olref{defn:aconvone}\olref{defn:aconvone4} مطابق۔ ایہ
    کر سکدے آں کیوں جے $z \ne x$، $z \notin FV(N')$ تے
    $\Subst{N'}{z}{x}$ معین اے۔ آخر وچ
    $\Subst{\lambd[z][\Subst{N'}{z}{x}]}{R}{y}$ معین اے،
    کیوں جے $z \neq y$ تے $z \notin FV(R)$۔

    نالے جے کوئی ہور $N''$ تے $R''$ اوہو شرطاں پوریاں کرن، تاں۔
    \begin{multline*}
      \Subst{(\lambd[z][\Subst{N''}{z}{x}])}{R''}{y} =\\
    \begin{aligned}
      &= \lambd[z][\Subst{\Subst{N''}{z}{x}}{R''}{y}] \\
      &= \lambd[z][\Subst{\Subst{N''}{z}{x}}{R}{y}] && \text{
        \olref{lem:sub:R} مطابق}\\
      &=\lambd[z][\Subst{\Subst{N'}{z}{x}}{R}{y}]
      && \text{استقرائی
        مفروضے مطابق}\\
      &=\Subst{(\lambd[z][\Subst{N'}{z}{x}])}{R}{y}
    \end{aligned}
    \end{multline*}
    \item $M$ دی شکل $(PQ)$ اے۔ مشق۔
  \end{enumerate}
\end{proof}

\begin{prob}
  \olref[lam][syn][alp]{thm:sub} دا ثبوت پورا کرو۔
\end{prob}

\begin{cor}\ollabel{cor:sub}
  ہر $M$، $R$ تے~$y$ لئی $M'$ تے $R'$ دا اک جوڑا موجود اے
  جیہدے لئی $M' \aeq M$، $R \aeq R'$ تے $\Subst{M'}{R'}{y}$
  معین اے۔ نالے جے کوئی ہور جوڑا $M'' \aeq M$ تے $R''$ اے
  جیہدے نال $\Subst{M'}{R'}{y}$ معین اے، تاں $\Subst{M'}{R'}{y} \aeq
  \Subst{M''}{R''}{y}$۔
\end{cor}

\begin{proof}
  \olref{thm:sub} توں سیدھا ملدا اے۔
\end{proof}
\endgroup
% /OLP-0362

% OLP-0363 content/lambda-calculus/syntax/de-bruijn.tex
\begingroup
\def\olfilename{de-bruijn}
\expandafter\def\csname ol@sectioncs\endcsname{section}

\olfileid{lam}{syn}{deb}

\olsection{ڈی برائن دا اشاریہ}

$\alpha$-ہم ارزی بڑی سُبھاوک گل اے، کیوں جے
$\alpha$-ہم ارز ٹرماں دا «اکّو مطلب» اے۔ دراصل لامبڈا
ٹرماں دی ایہو جیہی نحو دِتّی جا سکدی اے جیہڑی اوہناں ٹرماں
وچ فرق نہیں کردی جنہاں دا $\alpha$-روپ اک دوجے وچ بدلیا جا
سکدا اے۔ ایس کم دی سبھ توں معروف نحو وچ متغیراں دی تھاں
اوہناں دا \emph{ڈی برائن اشاریہ} ورتیا جاندا اے۔

جدوں اسیں $\lambd[x][M]$ لکھدے آں، تاں صاف دسدے آں کہ
$x$ فنکشن دا پیرامیٹر اے، تاں جو $x$ نوں $M$ وچ اوس
پیرامیٹر دے حوالے لئی ورت سکیے۔ پر ڈی برائن اشاریے وچ
پیرامیٹراں دے ناں نہیں ہوندے۔ فنکشن دے اندر اوہناں دے حوالے
لئی اک عدد لکھیا جاندا اے، جیہڑا حوالے تے پیرامیٹر دے
وچکار تجریداں دیاں تہاں دی گنتی دسدا اے۔ مثال $\lambd[x][\lambd[y][y x]]$
ویکھو۔ باہرلی تجرید متغیر~$x$ نوں تے اندرلی تجرید متغیر~$y$
نوں پابند کردی اے۔ ذیلی ٹرم $y x$ اندرلی تجرید دے دائرۂ اثر
وچ اے۔ $y$ تے اوہدی تجرید~$\lambd[y]$ دے وچکار کوئی تجرید
نہیں، پر $x$ تے اوہدی تجرید~$\lambd[x]$ دے وچکار اک تجرید
اے۔ سو $0\, 1$ نوں $y x$ دی تھاں لکھدے آں، تے سارے ٹرم
لئی $\lambd[][\lambd[][01]]$ لکھدے آں۔

\begin{defn}
  ڈی برائن ٹرم استقرا نال ایویں متعین ہوندے نیں۔
  \begin{enumerate}
  \item $n$، جتھے $n$ کوئی وی قدرتی عدد اے۔
  \item $PQ$، جتھے $P$ تے $Q$ دوویں ڈی برائن ٹرم نیں۔
  \item $\lambd[][N]$، جتھے $N$ اک ڈی برائن ٹرم اے۔
  \end{enumerate}
\end{defn}

عام لامبڈا ٹرماں نوں ڈی برائن اشاریے والے ٹرماں وچ لکھن دا
باقاعدہ طریقہ ایویں اے۔
\begin{defn}
  \begin{align*}
    F_\Gamma(x) &= \Gamma(x) \\
    F_\Gamma(PQ) &= F_\Gamma(P)F_\Gamma(Q) \\
    F_\Gamma(\lambd[x][N]) &= \lambd[][F_{x,\Gamma}(N)]
  \end{align*}
  ایتھے $\Gamma$ متغیراں دی فہرست اے جیہدے اشاریے صفر توں
  شروع ہوندے نیں، تے $\Gamma(x)$ متغیر $x$ دی $\Gamma$ وچلی
  تھاں دسدا اے۔ مثلاً، جے $\Gamma$ دی فہرست $x,y,z$ اے، تاں
  $\Gamma(x)$ دی قدر $0$ تے $\Gamma(z)$ دی قدر $2$ اے۔

  $x,\Gamma$ اوہ فہرست اے جیہڑی $x$ نوں $\Gamma$ دے مُڈھ
  اُتے رکھ کے بندی اے۔ مثلاً، پچھلی مثال دی $\Gamma$ لئی
  $w,\Gamma$ دی فہرست $w,x,y,z$ اے۔
\end{defn}

ڈی برائن ٹرم توں عام لامبڈا ٹرم واپس حاصل کرن دا طریقہ ایویں اے۔

\begin{defn}
  \begin{align*}
    G_\Gamma(n) &= \Gamma[n] \\
    G_\Gamma(PQ) &= G_\Gamma(P) G_\Gamma(Q) \\
    G_\Gamma(\lambd[][N]) &= \lambd[x][G_{x,\Gamma}(N)]
  \end{align*}
  ایتھے $\Gamma$ فیر متغیراں دی فہرست اے جیہدے اشاریے صفر توں
  شروع ہوندے نیں، تے $\Gamma[n]$ تھاں~$n$ اُتے موجود متغیر دسدا اے۔
  مثلاً، جے $\Gamma$ دی فہرست $x,y,z$ اے، تاں $\Gamma[1]$ دا
  متغیر $y$ اے۔

  آخری مساوات وچ متغیر $x$ دی تھاں کوئی وی اوہ متغیر چُندے
  آں جیہڑا $\Gamma$ وچ نہیں۔
\end{defn}

ایتھے اسیں کجھ نتیجے ثبوت توں بغیر دیندے آں۔

\begin{prop}
  جے $M \aconvone M'$ تے $\Gamma$ کوئی وی اوہ فہرست اے جیہدے
  وچ $\FV{M}$ شامل اے، تاں $F_\Gamma(M) \eqs F_\Gamma(M')$۔
\end{prop}
\endgroup
% /OLP-0363

% OLP-0364 content/lambda-calculus/syntax/term-revisited.tex
\begingroup
\def\olfilename{term-revisited}
\expandafter\def\csname ol@sectioncs\endcsname{section}

\olfileid{lam}{syn}{tr}
\olsection{\texorpdfstring{ٹرم بطور $\alpha$-ہم ارزیتی طبقے}{ٹرم بطور α-ہم ارزیتی طبقے}}

اگّوں اسیں ٹرماں نوں $\alpha$-ہم ارزی دے حساب نال ویکھاں گے۔
ایہدا مطلب اے کہ جدوں اسیں کوئی ٹرم لکھدے آں، تاں ساڈی مراد
اوہدا $\alpha$-ہم ارزیتی طبقہ اے۔ مثلاً، $\lambd[a][\lambd[b][a c]]$
اوہناں سارے ٹرماں دا سیٹ دسدا اے جیہڑے اوہدے نال $\alpha$-ہم ارز نیں،
جیویں $\lambd[a][\lambd[b][a
    c]]$، $\lambd[b][\lambd[a][b c]]$ وغیرہ۔

پچھلے حصیاں وچ $N, Q$ ورگے حرف اک ٹرم دسدے سن۔ اگّوں
اسیں اوہی حرف طبقہ دسن لئی ورتاں گے، تے ہِٹھاں ساڈا مطالعہ
ٹرماں دی تھاں ایہناں طبقیاں بارے ہووے گا۔ $x, y$ ورگے حرف
ہلے وی متغیر دسّن گے۔

اسیں $\rep{M}$ دی لکھت طبقہ~$M$ دے کسے وی !!{element}
لئی ورتاں گے، تے جے اک توں ودھ چاہیدے ہون تاں $\rep{M}[0], \rep{M}[1]$ وغیرہ۔

گل نوں سوکھا رکھن لئی ٹرماں والی لکھت فیر ورتدے آں۔
طبقیاں لئی تعریف ہِٹھ دِتّی اے۔
\begin{defn}
  \begin{enumerate}
  \item $\lambd[x][N]$ اوہ طبقہ اے جیہدے وچ
    $\lambd[x][\rep{N}]$ شامل اے۔
  \item $PQ$ اوہ طبقہ اے جیہدے وچ $\rep{P}\rep{Q}$ شامل اے۔
  \end{enumerate}
\end{defn}

ویکھنا اوکھا نہیں کہ ایہ ٹھیک طرح معین نیں، کیوں جے
$\alpha$-روپ بدلن دا تعلق موافق اے۔

\begin{defn} \ollabel{def:fv}
  $\alpha$-ہم ارزیتی طبقہ~$M$ دے \emph{آزاد متغیر}، یعنی
  $FV(M)$، دی تعریف $FV(\rep{M})$ اے۔
\end{defn}

ایہ ٹھیک طرح معین اے، کیوں جے $FV(\rep{M}[0]) =
FV(\rep{M}[1])$، جیویں \olref[alp]{thm:fv} وچ دکھایا اے۔

طبقیاں وچ بدل لئی وی اوہی لکھت ورتدے آں۔
\begin{defn} \ollabel{def:sub}
  $R$ نوں $y$ دی تھاں $M$ وچ رکھن والا \emph{بدل}، یعنی
  $\Subst{M}{R}{y}$، دی تعریف $\Subst{\rep{M}}{\rep{R}}{y}$ اے،
  کسے وی ایہو جیہے $\rep{M}$ تے $\rep{R}$ لئی جیہڑے بدل نوں
  معین کردے ہون۔
\end{defn}

ایہ وی ٹھیک طرح معین اے، جیویں \olref[alp]{cor:sub} وچ دکھایا اے۔

دھیان دیو کہ ایس تعریف نال ساڈا استدلال کِنّا سوکھا ہو جاندا
اے۔ ایہدی اک مثال ہِٹھ اے۔
\begin{align}
  \Subst{\lambd[x][x]}{y}{x} & =\ollabel{eq:1}\\
  &= \Subst{\lambd[z][z]}{y}{x} \ollabel{eq:2}\\
  &= \lambd[z][\Subst{z}{y}{x}] \\
  &= \lambd[z][z]
\end{align}

جے \olref{eq:1} نوں ہلے وی ٹرماں والا بدل منّیئے، تاں ایہ
غیر معین اے۔ پر جیویں پہلاں آکھیا اے، ہُن اسیں اینوں طبقیاں
والا بدل سمجھدے آں۔ ایس لئی \olref{eq:2} ممکن اے۔
$\lambd[x][x]$ دی تھاں $\lambd[z][z]$ رکھ سکدے آں، کیوں جے
دوویں اکّو طبقہ وچ نیں۔

ایسے کارن اگّوں اسیں منّاں گے کہ ساڈے چُنے ہوئے نمائندے بدل
لئی لوڑ دیاں شرطاں ہمیشہ پوریاں کردے نیں۔ مثلاً، جدوں
$\Subst{\lambd[x][N]}{R}{y}$ ویکھیے، تاں منّاں گے کہ نمائندہ
$\lambd[x][N]$ ایویں چُنیا اے کہ $x \neq y$ تے $x \notin FV(R)$۔

کیوں جے $\lambd[x][x]$ نوں «طبقہ» آکھنا ذرا اوپرا لگدا اے،
اگّوں ایہناں نوں $\Lambda$-ٹرم آکھدے آں، یا ایس حصے دے باقی
متن وچ بس «ٹرم»۔ ایس طرح ایہ پہلاں توں جانے ہوئے $\lambda$-ٹرماں
توں وکھرے رہن گے۔

\begin{editorial}
  ہلے اسیں ٹرماں نوں نہیں چھڈ سکدے۔ $\Lambda$-ٹرماں دی ساری
  تعریف $\lambda$-ٹرماں اُتے ٹکی اے، تے اسیں $\Lambda$-ٹرماں
  اُتے فنکشن متعین کرن دا کوئی طریقہ ہلے نہیں دِتّا۔ ایس دا
  مطلب اے کہ ایہو جیہا ہر فنکشن پہلاں $\lambda$-ٹرماں اُتے
  متعین کرنا پینا اے، تے فیر $\Lambda$-ٹرماں اُتے «منتقل» کرنا
  پینا اے، جیویں اسیں بدلاں لئی کیتا سی۔ فیر وی اسیں منّ کے
  چلدے آں کہ پڑھن والا بدیہی طور اُتے سمجھ سکدا اے کہ
  $\Lambda$-ٹرماں اُتے فنکشن کیویں متعین ہو سکدے نیں۔
\end{editorial}
\endgroup
% /OLP-0364

% OLP-0365 content/lambda-calculus/syntax/beta.tex
\begingroup
\def\olfilename{beta}
\expandafter\def\csname ol@sectioncs\endcsname{section}

\olfileid{lam}{int}{bet}
\olsection{\texorpdfstring{$\beta$-تقلیل}{β-تقلیل}}

جدوں $(\lambd[m][(\lambd[y][y]) m])$ ویکھیے، تاں سُبھاوک خیال
آؤندا اے کہ ایہدا $\lambd[m][m]$ نال کوئی جوڑ اے، یعنی پہلا
ٹرم «سوکھا» کرن نال دوجا ٹرم ملنا چاہیدا اے۔ \emph{$\beta$-تقلیل}
ایس خیال نوں باقاعدہ شکل دیندی اے۔

\begin{defn}[$\beta$-سکڑاؤ، $\bredone$] \ollabel{defn:betacontr}
  \emph{$\beta$-سکڑاؤ} ($\bredone$) ٹرماں اُتے سبھ توں چھوٹا
  موافق تعلق اے جیہڑا ہِٹھلی شرط پوری کردا اے۔
  \[
    (\lambd[x][N])Q \bredone \Subst{N}{Q}{x}
  \]
  اسیں آکھدے آں کہ $P$، \emph{$\beta$-سکڑاؤ} نال $Q$ بن جاندا
  اے جے $P \bredone Q$۔ $(\lambd[x][N])Q$ دی شکل والے ٹرم
  نوں \emph{ریڈیکس} آکھدے آں۔
\end{defn}

\begin{prob} \ollabel{prob:def}
  $\beta$-سکڑاؤ دی ہم ارز استقرائی تعریف کھول کے لکھو، جیویں
  اسیں پابند متغیر دا ناں بدلن لئی \olref[lam][syn][alp]{defn:aconvone}
  وچ کیتا سی۔
\end{prob}

\begin{defn}[$\beta$-تقلیل، $\bred$] \ollabel{defn:betared}
  \emph{$\beta$-تقلیل} ($\bred$) ٹرماں اُتے سبھ توں چھوٹا
  انعکاسی تے تعدّی والا تعلق اے جیہدے وچ $\bredone$ شامل اے۔
  اسیں آکھدے آں کہ $P$ دی $\beta$-تقلیل نال $Q$ ملدا اے جے
  $P \bred Q$۔
\end{defn}

جدوں گل دا مطلب صاف ہووے، اسیں $\redone$ لکھاں گے
$\bredone$ دی تھاں، تے $\red$ لکھاں گے $\bred$ دی تھاں۔

سادے لفظاں وچ، $M \bred N$ اُدوں تے صرف اُدوں اے جدوں $M$
نوں $N$ وچ $\beta$-سکڑاؤ دے صفر یا کئی قدماں نال بدلیا جا سکے۔

\begin{defn}[$\beta$-نارمل]
جیہڑا ٹرم $\beta$-سکڑاؤ نال ہور نہیں سکڑ سکدا، اوہنوں
\emph{$\beta$-نارمل} آکھدے آں۔
\end{defn}

جے $M \bred N$ تے $N$، $\beta$-نارمل اے، تاں اسیں
$N$ نوں~$M$ دا \emph{نارمل روپ} آکھدے آں۔ کوئی پُچھ سکدا اے
کہ کی ٹرم دا نارمل روپ اکّو ہوندا اے۔ جواب ہاں اے، جیویں
اگّے ویکھاں گے۔

کجھ مثالاں ویکھیے۔
\begin{enumerate}
\item ساڈے کول ایہ تقلیلاں نیں۔
\begin{align*}
(\lambd[x][xxy]) \lambd[z][z] & \redone (\lambd[z][z])(\lambd[z][z]) y \\
& \redone (\lambd[z][z]) y \\
& \redone y
\end{align*}
\item ٹرم نوں «سوکھا» کرنا اوہنوں ہور گُنجھلدار وی بنا سکدا اے۔
\begin{align*}
(\lambd[x][xxy])(\lambd[x][xxy]) & \redone (\lambd[x][xxy])(\lambd[x][xxy])y \\
& \redone (\lambd[x][xxy])(\lambd[x][xxy])yy \\
& \redone \dots
\end{align*}
\item ایہ ٹرم نوں اویویں وی رہن دے سکدا اے۔
\[
(\lambd[x][xx])(\lambd[x][xx]) \redone (\lambd[x][xx])(\lambd[x][xx])
\]
\item کجھ ٹرماں دی اک توں ودھ طریقیاں نال تقلیل ہو سکدی اے۔
  مثال لئی، سبھ توں باہرلے اطلاق دے سکڑاؤ نال ایہ ملدا اے۔
\[
(\lambd[x][(\lambd[y][yx]) z]) v \redone (\lambd[y][yv]) z
\]
تے سبھ توں اندرلے اطلاق دے سکڑاؤ نال ایہ ملدا اے۔
\[
(\lambd[x][(\lambd[y][yx]) z]) v \redone (\lambd[x][zx]) v
\]
دھیان دیو کہ ایس صورت وچ دوہاں ٹرماں دی اگلی تقلیل نال
اکّو ٹرم $zv$ ملدا اے۔
\end{enumerate}

پچھلی مثال دا آخری نتیجہ اتفاق نہیں۔ ایہ لامبڈا کیلکولس دی
اک ڈونگھی تے اہم خاصیت دی مثال اے جیہنوں چرچ–روسر خاصیت
آکھدے آں۔

\begin{digress}
  عام طور اُتے ٹرم دی $\beta$-تقلیل دا اک توں ودھ طریقہ ہوندا
  اے۔ ایس لئی تقلیل دیاں کئی چالاں بنائیاں گئیاں نیں، جنہاں
  وچوں سبھ توں عام \emph{سُبھاوک چال} اے۔ سُبھاوک چال ہمیشہ
  \emph{سبھ توں کھبے} ریڈیکس نوں سکیڑدی اے۔ ریڈیکس دی تھاں
  ٹرم وچ اوہدی شروعات والی تھاں اے۔ ایس چال دی خاصیت ایہ اے
  کہ کسے چال نال ٹرم دی تقلیل نارمل روپ تک اُدوں تے صرف اُدوں
  ہو سکدی اے جدوں سُبھاوک چال نال اوہدی تقلیل نارمل روپ تک
  ہو سکے۔ اگّے اسیں سُبھاوک چال ورتاں گے، جے کوئی ہور چال نہ
  دسی جاوے۔
\end{digress}

\begin{defn}[$\beta$-ہم ارزی، $\equal$]
  \emph{$\beta$-ہم ارزی} ($\equal$) اوہ تعلق اے جیہڑا استقرا نال
  ایویں متعین ہوندا اے۔
  \begin{enumerate}
  \item $M \equal M$۔
  \item جے $M \equal N$، تاں $N \equal M$۔
  \item جے $M \equal N$ تے $N \equal O$، تاں $M \equal O$۔
  \item جے $M \equal N$، تاں $PM \equal PN$۔
  \item جے $M \equal N$، تاں $MQ \equal NQ$۔
  \item جے $M \equal N$، تاں $\lambd[x][M] \equal \lambd[x][N]$۔
  \item $(\lambd[x][N])Q \equal \Subst{N}{Q}{x}$۔
  \end{enumerate}
\end{defn}

پہلے تِن قاعدے ایس تعلق نوں ہم ارزیتی تعلق بناؤندے نیں۔
اگلے تِن ایہنوں موافق بناؤندے نیں۔ آخری قاعدہ یقین بناؤندا اے
کہ $\beta$-سکڑاؤ ایس وچ شامل اے۔

سادے لفظاں وچ، دو ٹرم $\beta$-ہم ارز اُدوں تے صرف اُدوں
ہوندے نیں جدوں اک نوں دوجے وچ $\beta$-سکڑاؤ یا
$\beta$-سکڑاؤ دے «اُلٹ» دے صفر یا ودھ قدماں نال بدلیا جا
سکے۔ $\beta$-سکڑاؤ دا اُلٹ ایویں متعین اے کہ $M$ دا
اُلٹا-$\beta$-سکڑاؤ $N$ تک اُدوں تے صرف اُدوں اے جدوں $N$
دا $\beta$-سکڑاؤ $M$ تک اے۔

اُپرلے قاعدیاں دے نال اسیں ہور قاعدے پا کے تعلق نوں ودھاواں گے
تے ودھے ہوئے ہم ارزیتی تعلق نوں $\equal[X]$ لکھاں گے، جتھے $X$
وادھا کرن والا قاعدہ اے۔
\endgroup
% /OLP-0365

% OLP-0366 content/lambda-calculus/syntax/eta.tex
\begingroup
\def\olfilename{eta}
\expandafter\def\csname ol@sectioncs\endcsname{section}

\olfileid{lam}{syn}{eta}
\olsection{\texorpdfstring{$\eta$-روپ بدلنا}{η-روپ بدلنا}}

$\lambda$ ٹرماں اُتے اک ہور تعلق اے۔ \olref[fv]{sec} وچ
اسیں $\lambd[x][(fx)]$ دی مثال ورتی سی، جیہڑا اک دلیل قبول
کردا اے تے اوہدے اُتے $f$ لاؤندا اے۔ دوجے لفظاں وچ ایہ اوہی
فنکشن اے جو~$f$ اے۔ $\lambd[x][(fx)]N$ تے $fN$ دوہاں دی
تقلیل نال~$fN$ ملدا اے۔ اسیں ایس خیال نوں $\eta$-تقلیل
تے $\eta$-وادھے نال ظاہر کردے آں۔

\begin{defn}[$\eta$-سکڑاؤ، $\eredone$]
  \ollabel{defn:beredone}
  \emph{$\eta$-سکڑاؤ} ($\eredone$) ٹرماں اُتے سبھ توں چھوٹا
  موافق تعلق اے جیہڑا ہِٹھلی شرط پوری کردا اے۔
  \[
    \lambd[x][M x] \eredone M \text{ بشرطیکہ } x \notin FV(M)
  \]
\end{defn}

\begin{defn}[$\beta\eta$-تقلیل، $\bered$] \ollabel{defn:bered}
  \emph{$\beta\eta$-تقلیل} ($\bered$) ٹرماں اُتے سبھ توں چھوٹا
  انعکاسی تے تعدّی والا تعلق اے جیہدے وچ $\bredone$ تے
  $\eredone$ شامل نیں۔ یعنی ایہدے قاعدے انعکاس تے تعدّی دے
  قاعدے تے ایہ دو ہور قاعدے نیں۔
  \begin{enumerate}
  \item جے $M \bredone N$، تاں $M \bered N$۔ \ollabel{defn:bered3}
  \item جے $M \eredone N$، تاں $M \bered N$۔ \ollabel{defn:bered4}
  \end{enumerate}

\end{defn}

\begin{defn}
  اسیں ہم ارزیتی تعلق $\equal$ نوں $\eta$-روپ بدلن دے قاعدے نال
  ودھاؤندے آں۔
  \[
  \lambd[x][f x] \equal f
  \]
  ودھے ہوئے تعلق نوں $\equal[\eta]$ لکھدے آں۔
\end{defn}

$\eta$-ہم ارزی اہم اے، کیوں جے ایہ لامبڈا ٹرماں دی
امتدادی برابری نال جُڑی اے۔

\begin{defn}[امتدادی برابری]
  اسیں ہم ارزیتی تعلق $\equal$ نوں (\ext) دے قاعدے نال ودھاؤندے آں۔
  \begin{center}
  جے $Mx \equal Nx$، تاں $M \equal N$، بشرطیکہ $x \notin FV(MN)$۔
  \end{center}
  ودھے ہوئے تعلق نوں $\equal[\ext]$ لکھدے آں۔
\end{defn}

موٹے طور اُتے قاعدہ ایہ آکھدا اے کہ جے دو ٹرم فنکشن
وجوں اکّو دلیل اُتے اکّو جیہا ورتاؤ کرن، تاں اوہناں نوں
برابر سمجھنا چاہیدا اے۔

ہُن ثابت کردے آں کہ $\eta$ دا قاعدہ عین امتدادی برابری
دیندا اے تے ایس توں ودھ کجھ نہیں دیندا۔

\begin{thm}
  $M \equal[\ext] N$ اُدوں تے صرف اُدوں اے جدوں $M \equal[\eta] N$۔
\end{thm}

\begin{proof}
  پہلاں اسیں ثابت کردے آں کہ $\equal[\eta]$ امتدادی برابری
  دے قاعدے ہِٹھ بند اے۔ یعنی $ext$ دا قاعدہ $\equal[\eta]$ وچ
  کجھ ہور نہیں پاؤندا۔ ایس لئی $\equal[\eta]$ وچ $\equal[\ext]$
  شامل اے، تے جے $M \equal[\ext] N$، تاں $M \equal[\eta] N$۔

  ایہ ثابت کرن لئی کہ $\equal[\eta]$، \ext ہِٹھ بند اے،
  دھیان دیو کہ جے کوئی $M
  \equal N$، \ext{} دے قاعدے نال اخذ ہویا اے، تاں اوہدا
  مفروضہ $Mx
  \equal[\eta] Nx$ اے۔ فیر $\lambd[x][Mx]
  \equal[\eta] \lambd[x][Nx]$ ملدا اے $\equal$ دے اک قاعدے
  نال۔ دوہاں پاسیاں اُتے $\eta$ لاؤن نال $M \equal[\eta] N$
  ملدا اے۔

  ایس طرح ثابت کردے آں کہ $\eta$ دا قاعدہ $\equal[ext]$ وچ
  شامل اے۔ ہر $\lambd[x][Mx]$ تے $M$ لئی، جتھے $x \notin
  FV(M)$، ساڈے کول $(\lambd[x][Mx])x \equal[\ext] Mx$ اے۔
  \ext{} دے قاعدے نال ایس توں $\lambd[x][Mx] \equal[\ext] M$
  ملدا اے۔
\end{proof}
\endgroup
% /OLP-0366

% OLP-0651 content/lambda-calculus/syntax/conversion.tex
\begingroup
\def\olfilename{conversion}
\expandafter\def\csname ol@sectioncs\endcsname{section}

\olfileid{lam}{int}{con}
\olsection{روپ بدلنا تے تقلیل}

لامبڈا حساب وچ روپ بدلّن یا تقلیل دیاں تِن قسماں نیں۔ روپ بدلّن تے تقلیل دا فرق ایہ اے کہ «تبدیلی» اک پاسے ہوندی اے یا دوہاں پاسیاں ول، جیویں اسیں ویکھاں گے۔

$\alpha$-روپ بدلنا متغیر دا ناں بدلّن دی گل اے۔ ایہ قدرتی اے کہ $\lambda x. x$ تے $\lambda y.y$ نوں اِکّو فنکشن (شناختی فنکشن) دی نمائندگی کردیاں سمجھیا جاوے۔ ایہ خیال اسیں ایویں رسمی بناؤنے آں:
\endgroup
% /OLP-0651
\OLEndChapterHook
\endgroup
% /OLP-0356

% OLP-0367 content/lambda-calculus/church-rosser/church-rosser.tex
\begingroup
\def\olfilename{church-rosser}
\expandafter\def\csname ol@sectioncs\endcsname{section}

\olchapter{lam}{cr}{چرچ–روسر خاصیت}

% OLP-0368 content/lambda-calculus/church-rosser/definitions-and-properties.tex
\begingroup
\def\olfilename{definitions-and-properties}
\expandafter\def\csname ol@sectioncs\endcsname{section}

\olfileid{lam}{cr}{dap}

\olsection{تعریف تے خاصیتاں}

ایس باب وچ اسیں چرچ–روسر خاصیت دا خیال تے ایہدیاں کجھ
عام خاصیتاں پیش کردے آں۔

\begin{defn}[چرچ–روسر خاصیت، CR]
  ٹرماں اُتے تعلق $\xredone$ نوں \emph{چرچ–روسر خاصیت} والا
  اُدوں تے صرف اُدوں آکھدے آں جدوں $M \xredone P$ تے
  $M \xredone Q$ ہون دی ہر صورت وچ کوئی~$N$ موجود ہووے
  جیہدے لئی $P \xredone N$ تے $Q \xredone N$۔
\end{defn}

اسیں لامبڈا کیلکولس نوں حساب دا اک نمونہ سمجھ سکدے آں،
جیہدے وچ نارمل روپ والے ٹرم «قدراں» نیں تے ٹرماں اُتے
تقلیل دا تعلق «حساب دے قاعدے» دیندا اے۔ چرچ–روسر خاصیت
آکھدی اے کہ حساب اگّے ودھاؤن دے کئی طریقے ہون تے وی عبارت
دی اکّو قدر ہوندی اے۔

مُڈھلے الجبرا دی اک مثال ویکھو۔ $4 \times (1+2) + 3$ دا
حساب اک توں ودھ طریقیاں نال ہو سکدا اے۔ ایہدی تقلیل
$4
\times 3+3$ تک ہو سکدی اے، جے پہلاں $1+2$ دی تقلیل
~$3$ تک کریے۔ یا $4 \times 1+4
\times 2+3$ تک ہو سکدی اے، جے پہلاں $4 \times (1+2)$ دی
تقلیل ضرب نوں جمع دے دوہاں حصیاں اُتے ونڈن دے قاعدے نال کریے۔ دوہاں دی اگلی تقلیل
$12+3$ تک ہو سکدی اے۔

جے $\xredone$ دی تھاں $\beta$-تقلیل لئیے، تاں سوکھا ای
دِسدا اے کہ چرچ–روسر خاصیت توں ایہ نتیجہ نکلدا اے کہ جے
کسے ٹرم دا نارمل روپ موجود ہووے، تاں اوہ اکّو اے۔ فرض کرو
$M$ دی تقلیل $P$ تے $Q$ تک ہو سکدی اے، تے دوویں نارمل
روپ نیں۔ چرچ–روسر خاصیت مطابق کوئی $N$ موجود اے جیہدے
تک $P$ تے $Q$ دوہاں دی تقلیل ہوندی اے۔ مفروضے مطابق $P$
تے $Q$ نارمل روپ نیں، ایس لئی $P$ تے $Q$ دی $N$ تک
تقلیل صرف صفر قدماں والی تقلیل ہو سکدی اے۔ یعنی $P$،
$Q$ تے $N$ اکّو ٹرم نیں۔ ایس کرکے ٹرم دے \emph{اکّو}
نارمل روپ دی گل کرنا ٹھیک اے۔

سو جدوں اسیں لامبڈا کیلکولس نوں حساب دا نمونہ سمجھدے آں،
تاں کسے ٹرم دا نارمل روپ اوہ حساب دا «آخری نتیجہ» سمجھیا
جا سکدا اے جیہڑا ایس ٹرم توں شروع ہوندا اے۔ اُپرلا نتیجہ
دسدا اے کہ جے حساب دا کوئی آخری نتیجہ ہووے، تاں اکّو
ہووے گا۔ جیویں $4 \times (1+2)+3$ دا نتیجہ اکّو اے،
یعنی~$15$۔

\begin{thm} \ollabel{thm:str}
  جے تعلق $\xredone$ وچ چرچ–روسر خاصیت اے، تے $\xred$
  سبھ توں چھوٹا تعدّی والا تعلق اے جیہدے وچ $\xredone$
  شامل اے، تاں $\xred$ وچ وی چرچ–روسر خاصیت اے۔
\end{thm}

\begin{proof}
  فرض کرو ہِٹھلیاں زنجیراں نیں۔
  \begin{align*}
    M & \xredone P_1 \xredone \dots \xredone P_m \text{ تے}\\
    M & \xredone Q_1 \xredone \dots \xredone Q_n.
  \end{align*}
  قضیہ ثابت کرن لئی اسیں ٹرماں دا جال $N$ بناواں گے،
  جیہدی اُچائی $m + 1$ تے چوڑائی $n + 1$ اے۔ $N_{i,j}$
  نال $i$ویں قطار تے $j$ویں کالم دا ٹرم لکھدے آں۔

  اسیں $N$ ایویں بناواں گے کہ $N_{i,j} \xredone N_{i+1,j}$ تے
  $N_{i,j} \xredone N_{i,j+1}$۔ تعریف ایویں اے۔
  \begin{align*}
    N_{0,0} &= M \\
    N_{i,0} &= P_i && \text{جے } 1 \le i \le m \\
    N_{0,j} &= Q_j && \text{جے } 1 \le j \le n \\
  \intertext{تے ہور صورت وچ}
    N_{i,j} &= R
  \end{align*}
  ایتھے $R$ اوہ ٹرم اے جیہدے لئی $N_{i-1,j} \xredone R$ تے
  $N_{i,j-1}
  \xredone R$۔ $\xredone$ دی چرچ–روسر خاصیت مطابق ایہو
  جیہا ٹرم ہمیشہ موجود ہوندا اے۔

  ہُن ساڈے کول $N_{m,0} \xredone \dots \xredone N_{m,n}$ تے
  $N_{0,n}
  \xredone \dots \xredone N_{m,n}$ نیں۔ دھیان دیو کہ $N_{m,0}$،
  $P$ اے تے $N_{0,n}$، $Q$ اے۔ $\xred$ دی تعریف توں
  قضیہ نکلدا اے۔
\end{proof}
\endgroup
% /OLP-0368
% OLP-0369 content/lambda-calculus/church-rosser/parallel-beta-reduction.tex
\begingroup
\def\olfilename{parallel-beta-reduction}
\expandafter\def\csname ol@sectioncs\endcsname{section}

\olfileid{lam}{cr}{pb}

\olsection{\texorpdfstring{متوازی $\beta$-تقلیل}{متوازی β-تقلیل}}

اسیں \emph{متوازی $\beta$-تقلیل} دا خیال پیش کردے آں تے
ثابت کردے آں کہ ایہدے وچ چرچ–روسر خاصیت اے۔

\begin{defn}[متوازی $\beta$-تقلیل، $\bredpar$] \ollabel{defn:bredpar}
  ٹرماں دی متوازی تقلیل ($\bredpar$) استقرا نال ایویں متعین ہوندی اے۔
  \begin{enumerate}
    \item \ollabel{defn:bredpar1} $x \bredpar x$۔
    \item \ollabel{defn:bredpar2} جے $N \xrightarrow{\beta} N'$، تاں $ \lambd[x][N] \bredpar
      \lambd[x][N']$۔
    \item \ollabel{defn:bredpar3} جے $P \bredpar P'$ تے $Q \bredpar Q'$، تاں $PQ \bredpar
      P'Q'$۔
    \item \ollabel{defn:bredpar4} جے $N \bredpar N'$ تے $Q \bredpar Q'$، تاں
      $(\lambd[x][N])Q \bredpar \Subst{N'}{Q'}{x}$۔
  \end{enumerate}
\end{defn}

متوازی $\beta$-تقلیل نال ٹرم دے جِنّے مرضی ریڈیکس اکّو قدم
وچ سکیڑ سکدے آں۔ $\beta$-تقلیل نالوں فرق ایہ اے کہ اک قدم
وچ صرف اوہی ریڈیکس سکیڑے جا سکدے نیں جیہڑے اصل ٹرم وچ
موجود نیں۔ جیہڑے متوازی $\beta$-تقلیل دے نتیجے وچ پیدا ہون،
اوہ اوس قدم وچ نہیں سکیڑے جا سکدے۔ مثلاً ٹرم
$(\lambd[f][fx])(\lambd[y][y])$ دی متوازی $\beta$-تقلیل نال
یا اوہو ٹرم ملدا اے یا $(\lambd[y][y])x$۔ ایس قدم وچ
~$x$ تک نہیں پہنچ سکدے، بھاویں $\beta$-تقلیل نال~$x$ مل
سکدا اے۔ وجہ ایہ اے کہ ایہ ریڈیکس متوازی $\beta$-تقلیل
دے اک قدم پچھوں ای پیدا ہوندا اے۔ متوازی $\beta$-تقلیل
دا دوجا قدم لا کے~$x$ مل جاندا اے۔

\begin{thm}\ollabel{thm:refl}
  $M \bredpar M$۔
\end{thm}

\begin{proof}
  مشق۔
\end{proof}

\begin{prob}
  \olref[lam][cr][pb]{thm:refl} ثابت کرو۔
\end{prob}

\begin{defn}[$\beta$-مکمل تقلیل]\ollabel{defn:bcd}
  \emph{$\beta$-مکمل تقلیل}~$\bcd{M}$، جیہڑی $M$ توں ملدی اے،
  استقرا نال ایویں متعین ہوندی اے۔
  \begin{align}
    \bcd{x} &= x \ollabel{defn:bcd1} \\
    \bcd{(\lambd[x][N])} &= \lambd[x][\bcd{N}] \ollabel{defn:bcd2}\\
    \bcd{(PQ)} &= \bcd{P}\bcd{Q} && \text{جے $P$، $\lambd$ تجرید نہیں}
    \ollabel{defn:bcd3} \\
    \bcd{((\lambd[x][N])Q)} &= \Subst{\bcd{N}}{\bcd{Q}}{x} \ollabel{defn:bcd4}
  \end{align}
\end{defn}

جیویں ناں توں پتا لگدا اے، ٹرم دی $\beta$-مکمل تقلیل
«پوری متوازی تقلیل» اے۔ متوازی $\beta$-تقلیل وچ اسیں
کوئی ریڈیکس اویویں چھڈ سکدے آں۔ مکمل تقلیل وچ اوہناں ساریاں
نوں سکیڑنا لازمی اے۔ ایس لئی $(\lambd[f][fx])(\lambd[y][y])$
دی مکمل تقلیل $(\lambd[y][y])x$ اے۔ اصل ٹرم آپ نتیجہ نہیں اے۔

\begin{editorial}
  ایس تعریف وچ دِقّت ایہ اے کہ اسیں ہلے نہیں دسیا کہ
  ($\lambd$-)ٹرماں اُتے فنکشن بازگشت نال کیویں متعین کردے نیں۔
  ایہ اگّے ٹھیک کراں گے۔
\end{editorial}

\begin{lem}\ollabel{lem:comp}
  جے $M \bredpar M'$ تے $R \bredpar R'$، تاں $\Subst{M}{R}{y}
  \bredpar \Subst{M'}{R'}{y}$۔
\end{lem}

\begin{proof}
  $M \bredpar M'$ دے !!{derivation} اُتے استقرا نال۔
  \begin{enumerate}
    \item آخری قدم \olref{defn:bredpar1} اے۔ مشق۔
    \item آخری قدم \olref{defn:bredpar2} اے۔ فیر $M$ دی شکل
      $\lambd[x][N]$ تے $M'$ دی شکل $\lambd[x][N']$ اے،
      جتھے $N \bredpar N'$۔ اسیں ثابت کرنا اے کہ
      $\Subst{(\lambd[x][N])}{R}{y} \bredpar
      \Subst{(\lambd[x][N'])}{R'}{y}$، یعنی
      $\lambd[x][\Subst{N}{R}{y}] \bredpar
      \lambd[x][\Subst{N'}{R}{y}]$۔ ایہ \olref{defn:bredpar2}
      تے استقرائی مفروضے توں سیدھا ملدا اے۔
    \item آخری قدم \olref{defn:bredpar3} اے۔ مشق۔
    \item آخری قدم \olref{defn:bredpar4} اے۔ $M$ دی شکل
      $(\lambd[x][N])Q$ تے $M'$ دی شکل $\Subst{N'}{Q'}{x}$ اے۔
      اسیں ثابت کرنا اے کہ $\Subst{((\lambd[x][N])Q)}{R}{y} \bredpar
      \Subst{\Subst{N'}{Q'}{x}}{R'}{y}$، یعنی
      $(\lambd[x][\Subst{N}{R}{y}])\Subst{Q}{R}{y} \bredpar
      \Subst{\Subst{N'}{R'}{y}}{\Subst{Q'}{R'}{y}}{x}$۔ ایہ
      \olref{defn:bredpar4} تے استقرائی مفروضے توں ملدا اے۔
  \end{enumerate}
\end{proof}

\begin{prob}
  \olref[lam][cr][pb]{lem:comp} دا ثبوت پورا کرو۔
\end{prob}

\begin{lem}\ollabel{lem:cont}
  جے $M \bredpar M'$، تاں $M' \bredpar \bcd{M}$۔
\end{lem}
\begin{proof}
  $M \bredpar M'$ دے !!{derivation} اُتے استقرا نال۔
  \begin{enumerate}
    \item آخری قاعدہ \olref{defn:bredpar1} اے۔ مشق۔
    \item آخری قاعدہ \olref{defn:bredpar2} اے۔ $M$ دی شکل
      $\lambd[x][N]$ تے $M'$ دی شکل $\lambd[x][N']$ اے، جتھے
      $N \bredpar N'$۔ اسیں دکھاؤنا اے کہ $\lambd[x][N'] \bredpar
      \bcd{(\lambd[x][N])}$، یعنی $\lambd[x][N'] \bredpar
      \lambd[x][\bcd{N}]$، جیویں \olref{defn:bcd2} دسدا اے۔
      ایہ \olref{defn:bredpar2} تے استقرائی مفروضے توں ملدا اے۔
    \item آخری قاعدہ \olref{defn:bredpar3} اے۔ $M$ دی شکل $PQ$
      تے $M'$ دی شکل $P'Q'$ اے، کجھ $P$، $Q$، $P'$ تے $Q'$ لئی
      جیہناں وچ $P \bredpar P'$ تے $Q \bredpar Q'$۔
      استقرائی مفروضے مطابق $P'
      \bredpar \bcd{P}$ تے $Q' \bredpar \bcd{Q}$۔
      \begin{enumerate}
        \item جے $P$ دی شکل $\lambd[x][N]$ اے، کجھ $x$ تے $N$ لئی،
          تاں $P'$ دی شکل $\lambd[x][N']$ ہونی چاہیدی اے، کسے
          $N'$ لئی جیہدے نال $N \bredpar N'$۔ استقرائی مفروضے
          مطابق $N' \bredpar \bcd{N}$ تے $Q' \bredpar \bcd{Q}$۔
          فیر $(\lambd[x][N'])Q' \bredpar
          \Subst{\bcd{N}}{\bcd{Q}}{x}$، \olref{defn:bredpar4} مطابق۔
        \item جے $P$ کوئی $\lambd$ تجرید نہیں، تاں $P'Q' \bredpar
          \bcd{P}\bcd{Q}$، \olref{defn:bredpar3} مطابق۔ سجا
          پاسا $ \bcd{PQ}$ اے، \olref{defn:bcd3} مطابق۔
      \end{enumerate}
    \item آخری قاعدہ \olref{defn:bredpar4} اے۔ $M$ دی شکل
      $(\lambd[x][N])Q$ تے $M'$ دی شکل $\Subst{N'}{Q'}{x}$ اے،
      کجھ $x$، $N$، $Q$، $N'$ تے~$Q'$ لئی جیہناں وچ
      $N \bredpar N'$ تے $Q \bredpar
      Q'$۔ استقرائی مفروضے مطابق $N' \bredpar \bcd{N}$ تے
      $Q' \bredpar \bcd{Q}$۔ \olref{lem:comp} مطابق
      $\Subst{N'}{Q'}{x} \bredpar \Subst{\bcd{N}}{\bcd{Q}}{x}$،
      جیہدا سجا پاسا عین $\bcd{((\lambd[x][N])Q)}$ اے۔
  \end{enumerate}
\end{proof}

\begin{prob}
  \olref[lam][cr][pb]{lem:cont} دا ثبوت پورا کرو۔
\end{prob}

\begin{thm}\ollabel{thm:cr}
  $\bredpar$ وچ چرچ–روسر خاصیت اے۔
\end{thm}

\begin{proof}
  \olref{lem:cont} توں سیدھا ملدا اے۔
\end{proof}
\endgroup
% /OLP-0369
% OLP-0370 content/lambda-calculus/church-rosser/beta-reduction.tex
\begingroup
\def\olfilename{beta-reduction}
\expandafter\def\csname ol@sectioncs\endcsname{section}

\olfileid{lam}{cr}{b}

\olsection{\texorpdfstring{$\beta$-تقلیل}{β-تقلیل}}

\begin{lem}\ollabel{lem:one-par}
  جے $M \bredone M'$، تاں $M \bredpar M'$۔
\end{lem}
\begin{proof} جے $M \bredone M'$، تاں $M$ دی شکل
    $(\lambd[x][N])Q$ تے $M'$ دی شکل $\Subst{N}{Q}{x}$ اے،
    کجھ $x$، $N$ تے~$Q$ لئی۔ کیوں جے $N \bredpar N$ تے
    $Q \bredpar Q$، \olref[pb]{thm:refl} مطابق، ایس لئی
    $(\lambd[x][N])Q
    \bredpar \Subst{N}{Q}{x}$ سیدھا
    \olref[pb]{defn:bredpar}\olref[pb]{defn:bredpar4} توں ملدا اے۔
\end{proof}

\begin{lem}\ollabel{lem:par-red}
  جے $M \bredpar M'$، تاں $M \bred M'$۔
\end{lem}

\begin{proof} $M \bredpar M'$ دے !!{derivation} اُتے استقرا نال۔
  \begin{enumerate}
  \item آخری قاعدہ \olref[pb]{defn:bredpar1} اے۔ فیر $M$ تے $M'$
    دوویں $x$ نیں تے $x \bred x$۔
  \item آخری قاعدہ \olref[pb]{defn:bredpar2} اے۔ $M$ دی شکل
    $\lambd[x][N]$ تے $M'$ دی شکل $\lambd[x][N']$ اے، کجھ
    $x$، $N$، $N'$ لئی، جتھے $N \bredpar N'$۔ استقرائی مفروضے
    مطابق $N \bred N'$۔ فیر $\lambd[x][N] \bred \lambd[x][N']$
    اوہناں ای $\bredone$ سکڑاؤواں دے سلسلے نال جیہڑا
    $N \bred N'$ وچ ورتیا اے۔
  \item آخری قاعدہ \olref[pb]{defn:bredpar3} اے۔ $M$ دی شکل
    $PQ$ تے $M'$ دی شکل $P'Q'$ اے، کجھ $P$، $Q$، $P'$ تے
    $Q'$ لئی، جتھے $P \bredpar P'$ تے $Q \bredpar Q'$۔
    استقرائی مفروضے مطابق $P \bred P'$ تے $Q \bred
    Q'$۔ سو $PQ \bred P'Q'$ لئی پہلاں تقلیل $P \bred P'$
    لاؤ، تے فیر تقلیل $Q \bred Q'$۔
  \item آخری قاعدہ \olref[pb]{defn:bredpar4} اے۔ $M$ دی شکل
    $(\lambd[x][N])Q$ تے $M'$ دی شکل $\Subst{N'}{Q'}{x}$ اے،
    کجھ $x$، $N$، $M'$، $Q$، $Q'$ لئی، جتھے $N \bredpar N'$
    تے $Q \bredpar Q'$۔ استقرائی مفروضے توں $Q \bred Q'$ تے
    $N \bred N'$ ملدے نیں۔ سو $(\lambd[x][N])Q \bred \Subst{N'}{Q'}{x}$
    لئی پہلاں $N \bred N'$، فیر $Q \bred Q'$، تے آخر
    $(\lambd[x][N'])Q'$ دا سکڑاؤ کر کے $\Subst{N'}{Q'}{x}$ حاصل کرو۔
  \end{enumerate}
\end{proof}

\begin{lem}\ollabel{lem:str}
  $\bred$ سبھ توں چھوٹا تعدّی والا تعلق اے جیہدے وچ $\bredpar$ شامل اے۔
\end{lem}

\begin{proof}
  منّ لو $\xred$ سبھ توں چھوٹا تعدّی والا تعلق اے جیہدے
  وچ $\bredpar$ شامل اے۔

  $\bred \subseteq \xred$۔ فرض کرو $M \bred M'$، یعنی $M \ident M_1
  \bredone \dots \bredone M_k \ident M'$۔ \olref{lem:one-par}
  مطابق $M
  \ident M_1 \bredpar \dots \bredpar M_k \ident M'$۔ کیوں جے $\xred$
  وچ $\bredpar$ شامل اے تے ایہ تعدّی والا اے، ایس لئی $M \xred M'$۔

  $\xred \subseteq \bred$۔ فرض کرو $M \xred M'$، یعنی $M
  \ident M_1 \bredpar \dots \bredpar M_k \ident M'$۔
  \olref{lem:par-red} مطابق $M \ident M_1
  \bred \dots \bred M_k \ident M'$۔ کیوں جے $\bred$ تعدّی والا اے،
  ایس لئی $M
  \bred M'$۔
\end{proof}

\begin{thm}\ollabel{thm:cr}
  $\bred$ وچ چرچ–روسر خاصیت اے۔
\end{thm}

\begin{proof}
  \olref[dap]{thm:str}، \olref[pb]{thm:cr} تے \olref{lem:str} توں سیدھا ملدا اے۔
\end{proof}
\endgroup
% /OLP-0370
% OLP-0371 content/lambda-calculus/church-rosser/parallel-beta-eta-reduction.tex
\begingroup
\def\olfilename{parallel-beta-eta-reduction}
\expandafter\def\csname ol@sectioncs\endcsname{section}

\olfileid{lam}{cr}{pbe}

\olsection{\texorpdfstring{متوازی $\beta\eta$-تقلیل}{متوازی βη-تقلیل}}

ایس حصے وچ اسیں متوازی $\beta\eta$-تقلیل لئی چرچ–روسر خاصیت
ثابت کردے آں۔ ایہ $\beta\eta$-تقلیل دے مطابق متوازی تقلیل اے۔

\begin{defn}[متوازی $\beta\eta$-تقلیل، $\beredpar$] \ollabel{defn:beredpar}
  ٹرماں اُتے \emph{متوازی $\beta\eta$-تقلیل} ($\beredpar$)
  استقرا نال ایویں متعین ہوندی اے۔
  \begin{enumerate}
    \item \ollabel{defn:beredpar1} $x \beredpar x$۔
    \item \ollabel{defn:beredpar2} جے $N \xrightarrow{\beta} N'$، تاں $\lambd[x][N] \beredpar
      \lambd[x][N']$۔
    \item \ollabel{defn:beredpar3} جے $P \beredpar P'$ تے $Q \beredpar Q'$، تاں $PQ \beredpar
      P'Q'$۔
    \item \ollabel{defn:beredpar4} جے $N \beredpar N'$ تے $Q \beredpar Q'$، تاں
      $(\lambd[x][N])Q \beredpar \Subst{N'}{Q'}{x}$۔
    \item \ollabel{defn:beredpar5} جے $N \beredpar N'$، تاں $\lambd[x][Nx]
      \beredpar N'$، بشرطیکہ $x \notin FV(N)$۔
  \end{enumerate}
\end{defn}

\begin{thm}\ollabel{thm:refl}
  $M \beredpar M$۔
\end{thm}

\begin{proof}
  مشق۔
\end{proof}

\begin{prob}
  \olref[lam][cr][pbe]{thm:refl} ثابت کرو۔
\end{prob}

\begin{defn}[$\beta\eta$-مکمل تقلیل]\ollabel{defn:becd}
  \emph{$\beta\eta$-مکمل تقلیل}~$\becd{M}$، جیہڑی~$M$ توں
  ملدی اے، ایویں متعین ہوندی اے۔
  \begin{align}
    \becd{x} &= x \ollabel{defn:becd1} \\
    \becd{(\lambd[x][N])} &= \lambd[x][\becd{N}] \ollabel{defn:becd2}\\
    \becd{(PQ)} &= \becd{P}\becd{Q} && \text{جے $P$، $\lambd$ تجرید نہیں}
    \ollabel{defn:becd3} \\
    \becd{((\lambd[x][N])Q)} &= \Subst{\becd{N}}{\becd{Q}}{x}
                             \ollabel{defn:becd4} \\
    \becd{(\lambd[x][Nx])} &= \becd{N} \ollabel{defn:becd5} & \text{جے $x
                                                      \notin FV(N)$}
  \end{align}
\end{defn}

\begin{lem}\ollabel{lem:comp}
  جے $M \beredpar M'$ تے $R \beredpar R'$، تاں $\Subst{M}{R}{y}
  \beredpar \Subst{M'}{R'}{y}$۔
\end{lem}

\begin{proof}
  $M \beredpar M'$ دے !!{derivation} اُتے استقرا نال۔

  پہلِیاں چار صورتاں عین \olref[pb]{lem:comp} والیاں نیں۔
  جے آخری قاعدہ \olref{defn:beredpar5} اے، تاں $M$ دی شکل
  $\lambd[x][Nx]$ تے $M'$ دی شکل $N'$ اے، کجھ $x$ تے~$N'$ لئی
  جتھے $x \notin
  FV(N)$ تے $N \beredpar N'$۔ اسیں دکھاؤنا اے کہ
  $\Subst{(\lambd[x][Nx])}{R}{y} \beredpar \Subst{N'}{R'}{y}$، یعنی
  $\lambd[x][\Subst{N}{R}{y} x] \beredpar \Subst{N'}{R'}{y}$۔ ایہ
  \olref{defn:beredpar}\olref{defn:beredpar5} تے استقرائی مفروضے توں ملدا اے۔
\end{proof}

\begin{lem}\ollabel{lem:cont}
  جے $M \beredpar M'$، تاں $M' \beredpar \becd{M}$۔
\end{lem}

\begin{proof}
  $M \beredpar M'$ دے !!{derivation} اُتے استقرا نال۔

  پہلِیاں چار صورتاں \olref[pb]{lem:cont} والیاں وانگ نیں۔
  جے آخری قاعدہ \olref{defn:beredpar5} اے، تاں $M$ دی شکل
  $\lambd[x][Nx]$ تے $M'$ دی شکل $N'$ اے، کجھ $x$، $N$، $N'$ لئی
  جتھے $x \notin FV(N)$ تے $N
  \beredpar N'$۔ اسیں دکھاؤنا اے کہ $N' \beredpar
  \becd{(\lambd[x][Nx])}$، یعنی $N' \beredpar \becd{N}$۔
  ایہ استقرائی مفروضے توں سیدھا ملدا اے۔
\end{proof}

\begin{thm}\ollabel{thm:cr}
  $\beredpar$ وچ چرچ–روسر خاصیت اے۔
\end{thm}

\begin{proof}
  \olref{lem:cont} توں سیدھا ملدا اے۔
\end{proof}
\endgroup
% /OLP-0371
% OLP-0372 content/lambda-calculus/church-rosser/beta-eta-reduction.tex
\begingroup
\def\olfilename{beta-eta-reduction}
\expandafter\def\csname ol@sectioncs\endcsname{section}

\olfileid{lam}{cr}{be}

\olsection{\texorpdfstring{$\beta\eta$-تقلیل}{βη-تقلیل}}

چرچ–روسر خاصیت $\beta\eta$-تقلیل ($\bered$) لئی قائم اے۔

\begin{lem}\ollabel{lem:one-par}
  جے $M \beredone M'$، تاں $M \beredpar M'$۔
\end{lem}

\begin{proof}
  $M \beredone M'$ دے !!{derivation} اُتے استقرا نال۔ جے $M \bredone
  M'$، $\eta$-روپ بدلن نال اے، یعنی \olref[syn][eta]{defn:beredone}
  مطابق، تاں اسیں \olref[pbe]{thm:refl} ورتدے آں۔ باقی صورتاں
  \olref[b]{lem:one-par} والیاں وانگ نیں۔
\end{proof}

\begin{lem}\ollabel{lem:par-red}
  جے $M \beredpar M'$، تاں $M \bered M'$۔
\end{lem}

\begin{proof} $M \beredpar M'$ دے !!{derivation} اُتے استقرا نال۔

  جے آخری قاعدہ \olref[pbe]{defn:beredpar5} اے، تاں $M$
  دی شکل $\lambd[x][Nx]$ تے $M'$ دی شکل $N'$ اے، کجھ $x$،
  $N$، $N'$ لئی جتھے $x
  \notin FV(N)$ تے $N \beredpar N'$۔ سو اسیں پہلاں
  $\lambd[x][Nx]$ دی تقلیل $N$ تک $\eta$-روپ بدلن نال
  کر سکدے آں، تے فیر $\beredone$ قدماں دا اوہ سلسلہ لا سکدے
  آں جیہڑا $N \bered N'$ دکھاؤندا اے۔ ایہ سلسلہ استقرائی مفروضے
  توں ملدا اے۔
\end{proof}

\begin{lem}\ollabel{lem:str}
  $\bered$ سبھ توں چھوٹا تعدّی والا تعلق اے جیہدے وچ $\beredpar$ شامل اے۔
\end{lem}

\begin{proof}
  \olref[b]{lem:str} وانگ۔
\end{proof}

\begin{thm}\ollabel{thm:cr}
  $\bered$ وچ چرچ–روسر خاصیت اے۔
\end{thm}

\begin{proof}
  \olref[dap]{thm:str}، \olref[pbe]{thm:cr} تے \olref{lem:str} توں۔
\end{proof}
\endgroup
% /OLP-0372

\OLEndChapterHook
\endgroup
% /OLP-0367

% OLP-0373 content/lambda-calculus/lambda-definability/lambda-definability.tex
\begingroup
\def\olfilename{lambda-definability}
\expandafter\def\csname ol@sectioncs\endcsname{section}

\olchapter{lam}{rep}{لامبڈا نال تعریف ہون دی اہلیت}

\begin{editorial}
  ایہ باب آزمائشی اے۔ اینوں ہور وضاحت دی لوڑ اے۔ مواد نوں
  تعریفاں، ثبوتاں والے قضیاں تے ہور مثالاں وچ بہتر ترتیب دین دی
  لوڑ اے۔
\end{editorial}

% OLP-0374 content/lambda-calculus/lambda-definability/introduction.tex
\begingroup
\def\olfilename{introduction}
\expandafter\def\csname ol@sectioncs\endcsname{section}

\olfileid{lam}{ldf}{int}
\olsection{جان پہچان}

پہلی نظر وچ لامبڈا کیلکولس بس عبارتاں دا بڑا تجریدی حسابی
نظام لگدا اے۔ ایہ عبارتاں فنکشوناں تے اوہناں نوں ہور فنکشوناں
اُتے لاؤن دا اظہار کردیاں نیں۔ لامبڈا کیلکولس دی نحو وچ کوئی
ایہو جیہی گل نہیں جیہڑی دسے کہ ایہ کسے خاص قسم دیاں شَواں
اُتے فنکشن نیں۔ خاص کر کے، نحو وچ قدرتی عدداں دا کوئی ذکر
نہیں۔ ایہدیاں بُنیادی کارروائیاں، اطلاق تے لامبڈا تجرید، ہر
فنکشن لئی ورتیاں جا سکدیاں نیں، صرف قدرتی عدداں اُتے فنکشوناں
لئی نہیں۔

فیر وی تھوڑی جیہی سُوجھ نال لامبڈا کیلکولس وچ حسابی فنکشن،
یعنی قدرتی عدداں اُتے فنکشن، متعین کیتے جا سکدے نیں۔ ایس
لئی اسیں ہر قدرتی عدد~$n \in
\Nat$ لئی اک خاص $\lambd$-ٹرم~$\num n$ متعین کردے آں، جیہڑا
~$n$ دا \emph{چرچ عدد} اے۔ (چرچ عدداں دا ناں الونزو چرچ توں آیا اے۔)

\begin{defn}
  جے $n \in \Nat$، تاں اوہدے نال جُڑیا \emph{چرچ عدد} $\num{n}$
  عدد~$n$ نوں ظاہر کردا اے۔
  \[
    \num{n} \ident \lambd[fx][f^n(x)]
  \]
  ایتھے $f^n(x)$ اوہ نتیجہ اے جیہڑا $f$ نوں~$x$ اُتے $n$
  واری لاؤن نال ملدا اے۔ مثلاً $\num{0}$، $\lambd[fx][x]$ اے،
  تے $\num{3}$، $\lambd[fx][f(f(f\,x))]$ اے۔
\end{defn}

چرچ عدد $\num n$ نوں لامبڈا ٹرم دے روپ وچ لکھیا جاندا اے۔
اوہ ٹرم اک فنکشن ظاہر کردا اے جیہڑا دو دلیلاں $f$ تے $x$
قبول کردا اے تے $f^n(x)$ واپس دیندا اے۔ صاف اے کہ چرچ
عدد نارمل روپ وچ ہوندے نیں۔

قدرتی عدداں نوں لامبڈا کیلکولس وچ ظاہر کرنا اُدوں ای کم آؤندا
اے جدوں اسیں اوہناں نال حساب کر سکیے۔ لامبڈا کیلکولس وچ
چرچ عدداں نال حساب کرن دا مطلب اے کہ اک $\lambd$-ٹرم~$F$
ایہو جیہے چرچ عدد اُتے لائیے تے سانجھے ٹرم~$F\, \num n$ دی
تقلیل نارمل روپ تک کریے۔ جے ایہ ہمیشہ نارمل روپ تک پہنچدا
اے تے اوہ نارمل روپ ہمیشہ کوئی چرچ عدد~$\num m$ ہوندا اے،
تاں اسیں عدد~$m$ نوں حساب دی حاصل قدر سمجھ سکدے آں۔
فیر اسیں سمجھ سکدے آں کہ~$F$ اوہ فنکشن $f\colon \Nat \to \Nat$
متعین کردا اے جیہدے لئی $f(n) = m$ اُدوں تے صرف اُدوں
اے جدوں $F\, \num n \red \num m$۔ چرچ–روسر خاصیت کر کے
جے نارمل روپ موجود ہووے، تاں اکّو ہوندا اے۔ سو جے
$F\, \num n \red \num m$، تاں کوئی ہور نارمل ٹرم، خاص کر کے کوئی
ہور چرچ عدد، نہیں ہو سکدا جیہدے تک $F \, \num n$ دی تقلیل ہووے۔

اُلٹے پاسے، فنکشن $f\colon \Nat \to \Nat$ دِتّا ہووے، تاں
اسیں پُچھ سکدے آں کہ کوئی ٹرم $F$ موجود اے جیہڑا~$f$ نوں
ایس طرح متعین کرے۔ اوہ صورت وچ اسیں آکھدے آں کہ $F$،
~$f$ نوں \emph{!!{lambda define}s} تے $f$، !!{lambda definable}
اے۔ ایہ گل کئی دلیلاں والے تے جزوی فنکشوناں لئی وی عام کر سکدے آں۔

\begin{defn}
فرض کرو $f\colon \Nat^k \to \Nat$۔ اسیں آکھدے آں کہ
لامبڈا ٹرم~$F$، $f$ نوں \emph{!!{lambda define}s} جے ہر
$n_0$، \dots،~$n_{k-1}$ لئی
\[
F \, \num{n_0} \, \num{n_1} \dots \num{n_{k-1}} \red \num{f(n_0, n_1, \dots,
  n_{k-1})}
\]
ہووے جدوں $f(n_0, \dots, n_{k-1})$ معین ہووے، تے ہور صورت
وچ $F \, \num{n_0} \,
\num{n_1} \dots \num{n_{k-1}}$ دا کوئی نارمل روپ نہ ہووے۔
\end{defn}

ثابت قدر والے فنکشن بڑی سادی مثال نیں۔ ٹرم $C_k \ident
\lambd[x][\num{k}]$ اوہ فنکشن $c_k\colon \Nat \to
\Nat$ نوں !!{lambda define}s جیہدے لئی $c(n) = k$۔ کیوں جے $C_k \, \num n \ident
(\lambd[x][\num{k}])\num n \redone \num{k}$، بھاویں~$n$ کوئی
وی ہووے۔ عین اوہی قدر واپس دین والا فنکشن !!{lambda defined}
اے $\lambd[x][x]$ نال۔ ہور گُنجھلدار فنکشوناں نوں متعین کرنا
اوکھا ہوندا اے تے ایس لئی اکثر بڑی سُوجھ دی لوڑ پیندی اے۔
سو ایہ گل شاید حیرانی والی اے کہ ہر قابلِ حساب فنکشن
!!{lambda definable} اے۔ اُلٹی گل وی سچ اے۔ جے کوئی فنکشن
!!{lambda definable} اے، تاں اوہ قابلِ حساب اے۔
\endgroup
% /OLP-0374
% OLP-0375 content/lambda-calculus/lambda-definability/arithmetical-functions.tex
\begingroup
\def\olfilename{arithmetical-functions}
\expandafter\def\csname ol@sectioncs\endcsname{section}

\olfileid{lam}{rep}{arf}
\olsection{\usetoken{S}{lambda definable} حسابی فنکشن}

\begin{prop}
  \ollabel{prop:succ-ld}
  جانشین فنکشن~$\Succ$، !!{lambda definable} اے۔
\end{prop}

\begin{proof}
جانشین فنکشن نوں جیہڑا ٹرم !!{lambda define}s، اوہ ایہ اے۔
\begin{align*}
  \fn{Succ} & \ident \lambd[a][\lambd[fx][f ({a} f x)]].
  \intertext{ساڈے لکھت دے رواج مطابق ایہ ہِٹھلی لکھت دا مختصر روپ اے۔}
  \fn{Succ} & \ident \lambd[a][\lambd[f][\lambd[x][(f (({a} f) x))]]].
  \intertext{$\fn{Succ}$ اوہ فنکشن اے جیہڑا دلیل وجوں عدد~$a$
    لیندا اے تے ہور فنکشن $\lambd[fx][f ({a}
      f x)]$ دیندا اے۔ ایہ فنکشن آپ چرچ عدد نہیں۔ پر جے دلیل $a$
    چرچ عدد ہووے، تاں ایہدی تقلیل نال چرچ عدد ملدا اے۔ ویکھو۔}
   (\lambd[a][\lambd[fx][f ({a} f x)]])\,\num{n} & \redone
   \lambd[fx][f ({\num{n}} f x)].
   \intertext{اندرلا ٹرم $\num{n}fx$ اک ریڈیکس اے، کیوں جے
     $\num{n}$، $\lambd[fx][f^nx]$ اے۔ سو $\num{n}fx \redone f^nx$،
     تے سارے ٹرم لئی سانوں ایہ ملدا اے۔}
   \fn{Succ}\, \num n & \red \lambd[fx][f(f^n(x))],
\end{align*}
یعنی $\num{n+1}$۔
\end{proof}

\begin{ex}
  ویکھیے کہ $\fn{Succ}$ نوں~$\num{0}$، یعنی $\lambd[fx][x]$،
  اُتے لائیے تاں کیہ ہوندا اے۔ اسیں ٹرم پورے لکھدے آں۔
  \begin{align*}
    \fn{Succ}\, \num{0} & \ident (\lambd[a][\lambd[f][\lambd[x][(f (({a} f) x))]]])(\lambd[f][\lambd[x][x]])\\
    & \redone \lambd[f][\lambd[x][(f (({(\lambd[f][\lambd[x][x]])} f) x))]] \\
    & \redone \lambd[f][\lambd[x][(f ({(\lambd[x][x])} x))]] \\
    & \redone \lambd[f][\lambd[x][(f x)]] \ident \num{1}\\
  \end{align*}
\end{ex}

\begin{prob}
  ہِٹھلا ٹرم
  \begin{align*}
    \fn{Succ}' & \ident \lambd[{n}][\lambd[fx][{n} f (f x)]]
  \end{align*}
  جانشین فنکشن نوں !!{lambda define}s۔ دسو کیوں۔
\end{prob}

\begin{prop}
  \ollabel{prop:add-ld}
  جمع دا فنکشن~$\Add$، !!{lambda definable} اے۔
\end{prop}

\begin{proof}
جمع ہِٹھلے ٹرماں نال !!{lambda defined} اے۔
\begin{align*}
  \fn{Add} & \ident \lambd[{a}{b}][\lambd[fx][{a} f ({b} f x)]]
\intertext{یا دوجے طریقے نال}
  \fn{Add}' & \ident \lambd[{a}{b}][{a}\, \fn{Succ} \, {b}].
\intertext{جمع دا پہلا طریقہ ایویں کم کردا اے۔ $\fn{Add}$ پہلاں
  دو عدد ${a}$ تے ${b}$ قبول کردا اے۔ نتیجہ اوہ فنکشن اے جیہڑا
  $f$ تے $x$ لیندا اے تے $af(bfx)$ واپس دیندا اے۔ جے $a$ تے $b$ چرچ
  عدد $\num{n}$ تے $\num{m}$ ہون، تاں ایس دی تقلیل $f^{n+m}(x)$
  تک ہوندی اے، جیہڑا $f^{n}(f^{m}(x))$ دے عین برابر اے۔ ہولی ہولی
  قدم ویکھیے۔}
  (\lambd[{a}{b}][\lambd[fx][{a} f ({b} f x)]])\num n\,\num m & \redone
  \lambd[fx][\num{n}\, f (\num {m}\, f x)] \\
  & \redone \lambd[fx][\num{n}\, f (f^m x)] \\
  & \redone \lambd[fx][f^n (f^m x)] \ident \num {n+m}.
\intertext{جمع دا دوجا روپ $\fn{Add'}$ وکھرے طریقے نال کم کردا اے۔
  اینوں دو چرچ عدداں $\num{n}$ تے $\num{m}$ اُتے لائیے، تاں}
\fn{Add}' \num n \,\num m
& \redone \num{n}\, \fn{Succ}\, \num{m}.
\intertext{پر $\num{n} f x$ دی تقلیل ہمیشہ $f^n(x)$ تک ہوندی اے۔
  ایس لئی}
  \num{n}\,  \fn{Succ}\, \num{m} & \red \fn{Succ}^n(\num{m}).
\end{align*}
تے کیوں جے $\fn{Succ}$ جانشین فنکشن نوں !!{lambda define}s تے
جانشین فنکشن $n$ واری~$m$ اُتے لاؤن نال $n+m$ ملدا اے، ایس
لئی ایس دی اگلی تقلیل~$\num{n+m}$ تک ہوندی اے۔
\end{proof}

\begin{prop}
  \ollabel{prop:mult-ld}
  ضرب ہِٹھلے ٹرم نال !!{lambda definable} اے۔
  \[
  \fn{Mult} \ident \lambd[ab][\lambd[fx][a (b f) x]]
  \]
\end{prop}

\begin{proof}
  ایہ کم کیویں ہوندا اے، ویکھن لئی فرض کرو $\fn{Mult}$ نوں چرچ
  عدداں $\num{n}$ تے $\num{m}$ اُتے لاؤندے آں۔ $\fn{Mult} \, \num{n} \, \num{m}$
  دی تقلیل $\lambd[fx][\num{n}(\num{m}\, f)x]$ تک ہوندی اے۔
  ٹرم $\num{m} f$ اوہ فنکشن متعین کردا اے جیہڑا $f$ نوں اپنی
  دلیل اُتے $m$ واری لاؤندا اے۔ سو $\num{n} (\num{m} f) x$ فنکشن
  «$f$ نوں $m$ واری لاؤ» آپ $n$ واری~$x$ اُتے لاؤندا اے۔
  دوجے لفظاں وچ اسیں $f$ نوں $x$ اُتے $n\cdot
  m$ واری لاؤندے آں۔ پر نتیجے دا نارمل ٹرم بس چرچ عدد
  $\num{nm}$ اے۔
\end{proof}

\begin{editorial}
اسیں $\eta$-تقلیل نال ایس ٹرم نوں ہور سوکھا کر سکدے آں۔
\[
  \fn{Mult} \ident \lambd[ab][\lambd[f][a (b f)]].
\]

پر فیر پہلاں $\eta$-تقلیل سمجھاؤنی پئے گی۔
\end{editorial}

\begin{prob}
ضرب ہِٹھلے ٹرم نال !!{lambda defined} ہو سکدی اے۔
\[
  \fn{Mult}' \ident \lambd[ab][a (\fn{Add}\, a) \num{0}].
\]
دسو ایہ کیوں کم کردا اے۔
\end{prob}

گھات چڑھاؤن دی تعریف $\lambd$-ٹرم دے طور اُتے حیرانی جوگی
سادی اے۔
\[
  \fn{Exp} \ident \lambd[be][e b].
\]
پہلی دلیل $b$ بُنیاد اے تے دوجی~$e$ گھات۔ ساڈی عددی لکھت
مطابق سمجھ ایہ اے کہ $e f$، $f^e$ اے۔ جے ایہ سمجھنا اوکھا
لگے، تاں گھات چڑھاؤن دی تعریف بار بار ضرب نال وی کر سکدے آں۔
\[
  \fn{Exp}' \ident \lambd[be][e (\fn{Mult}\, b) \num{1}].
\]

چرچ عدداں لئی پچھلا عدد تے گھٹاؤ اِنّے سوکھے نہیں جِنّے
لگدے نیں۔ ایہناں لئی جوڑیاں نوں اک خاص روپ وچ لکھنا پیندا اے۔
\endgroup
% /OLP-0375
% OLP-0376 content/lambda-calculus/lambda-definability/pairs.tex
\begingroup
\def\olfilename{pairs}
\expandafter\def\csname ol@sectioncs\endcsname{section}

\olfileid{lam}{ldf}{pai}
\olsection{جوڑے تے پچھلا عدد}

\begin{defn}
$M$ تے $N$ دا جوڑا، جیہڑا $\tuple{M,N}$ لکھدے آں، ایویں
متعین ہوندا اے۔
\[
\tuple{M,N} \ident \lambd[f][fMN].
\]
\end{defn}

بدیہی خیال ایہ اے کہ ایہ اوہ فنکشن اے جیہڑا اک فنکشن قبول
کردا اے تے اوہنوں جوڑے دے دوہاں رکناں اُتے لاؤندا اے۔ ایس
خیال توں ایہ جوڑا بناؤن والا فنکشن ملدا اے، جیہڑا دو ٹرم
لیندا اے تے اوہناں نوں رکھن والا جوڑا واپس دیندا اے۔
\[
  \fn{Pair} \ident \lambd[mn][\lambd[f][fmn]]
\]
جے جوڑا دِتّا ہووے، اسیں اوہدے رکن وی واپس حاصل کرنا چاہندے آں۔
ایس لئی دو رکن کڈھن والے فنکشن چاہیدے نیں، جیہڑے دلیل وجوں جوڑا
لیندے نیں تے اوہدے وچوں پہلا یا دوجا رکن واپس دیندے نیں۔
\begin{align*}
  \fn{Fst} & \ident \lambd[p][p(\lambd[mn][m])]\\
  \fn{Snd} & \ident \lambd[p][p(\lambd[mn][n])]
\end{align*}

\begin{prob}
  دسو کہ رکن کڈھن والے فنکشن $\fn{Fst}$ تے $\fn{Snd}$ کیوں
  کم کردے نیں۔
\end{prob}

ہُن جوڑیاں دی مدد نال پچھلا عدد دین والے فنکشن نوں
!!{lambda define} کر سکدے آں۔
\[
  \fn{Pred} \ident \lambd[n][\fn{Fst}(n (\lambd[p][\tuple{\fn{Snd}\, {p}, \fn{Succ}(\fn{Snd}\, {p})}]) \tuple{\num 0, \num 0})]
\]
یاد رکھو کہ $\num n\, f x$ دی تقلیل $f^{n}(x)$ تک ہوندی
اے۔ ایس صورت وچ $f$ اک فنکشن اے جیہڑا جوڑا $p$ لیندا اے
تے نواں جوڑا واپس دیندا اے۔ نویں جوڑے وچ $p$ دا دوجا رکن
تے اوس دوجے رکن دا جانشین ہوندا اے۔ $x$ جوڑا $\tuple{0,0}$
اے۔ سو نتیجہ $\tuple{0,0}$ اے جے $n=0$، تے ہور صورت وچ
$\tuple{\num{n-1}, \num n}$ اے۔ فیر $\fn{Pred}$ نتیجے دا پہلا
رکن واپس دیندا اے۔

گھٹاؤ دی تعریف $\fn{Pred}$ نوں $a$ اُتے $b$ واری لاؤن نال
ہو سکدی اے۔
\[
\fn{Sub} \ident \lambd[ab][b \fn{Pred}\, a].
\]
\endgroup
% /OLP-0376
% OLP-0377 content/lambda-calculus/lambda-definability/truth-values.tex
\begingroup
\def\olfilename{truth-values}
\expandafter\def\csname ol@sectioncs\endcsname{section}

\olfileid{lam}{ldf}{tvr}
\olsection{صداقتی قدراں تے تعلق}

خالص لامبڈا کیلکولس وچ صداقتی قدراں نوں ایویں لکھ سکدے آں۔
\begin{align*}
  \fn{true} & \ident \lambd[x][\lambd[y][x]]\\
  \fn{false} & \ident \lambd[x][\lambd[y][y]]
\end{align*}

صداقتی قدراں نوں \emph{چُنن والے فنکشوناں} دے روپ وچ لکھدے
آں، یعنی اوہ فنکشن جیہڑے دو دلیلاں لیندے نیں تے اوہناں وچوں
اک واپس دیندے نیں۔ صداقتی قدر $\fn{true}$ اپنی پہلی دلیل
چُندی اے تے $\fn{false}$ دوجی۔ مثلاً $\fn{true}\, M N$ دی
تقلیل ہمیشہ $M$ تک ہوندی اے، تے $\fn{false}\, M N$ دی
ہمیشہ~$N$ تک۔

\begin{defn}
اسیں تعلق $R \subseteq \Nat^n$ نوں !!{lambda definable}
آکھدے آں جے کوئی ٹرم~$R$ ایہو جیہا ہووے کہ
\begin{align*}
  R\, \num{n_1} \dots \num{n_k} & \bred \fn{true}
  \intertext{جدوں وی $R(n_1, \dots, n_k)$ قائم ہووے، تے}
  R\, \num{n_1} \dots \num{n_k} & \bred \fn{false}
\end{align*}
ہور صورت وچ۔
\end{defn}

مثلاً تعلق $\fn{IsZero} = \{0\}$ جیہڑا $0$ لئی تے صرف $0$
لئی قائم اے، ہِٹھلے ٹرم نال !!{lambda definable} اے۔
\[
  \fn{IsZero} \ident \lambd[n][n (\lambd[x][\fn{false}])\, \fn{true}].
\]
ایہ کیویں کم کردا اے؟ چرچ عدد دہراؤ والے فنکشن نیں جیہڑے
پہلی دلیل نوں $n$ واری دوجی اُتے لاؤندے نیں۔ سو پہلاں
قدر $\fn{true}$ رکھدے آں۔ تکرار دے ہر قدم اُتے پچھلے
قدم دا نتیجہ جو وی ہووے، $\fn{false}$ واپس دیندے آں۔
ایہ قدم مُڈھلی قدر اُتے $n$ واری لگے گا۔ نتیجہ $\fn{true}$
اُدوں تے صرف اُدوں ہووے گا جدوں ایہ قدم اک واری وی نہ لگے،
یعنی جدوں $n = 0$۔

صداقتی قدراں دی ایس نمائندگی توں اسیں ہور صداقتی فنکشن
وی متعین کر سکدے آں۔ نفی تے عطف لئی دو فنکشن ایہ نیں۔
\begin{align*}
  \fn{Not} & \ident \lambd[x][x\, \fn{false}\, \fn{true}]\\
  \fn{And} & \ident \lambd[x][\lambd[y][xy \,\fn{false}]]
\end{align*}
فنکشن «$\fn{Not}$» اک دلیل لیندا اے۔ ایہ $\fn{true}$ واپس
دیندا اے جے دلیل $\fn{false}$ ہووے، تے $\fn{false}$ واپس
دیندا اے جے دلیل~$\fn{true}$ ہووے۔ فنکشن «$\fn{And}$» دو
صداقتی قدراں دلیل وجوں لیندا اے، تے $\fn{true}$ اُدوں تے
صرف اُدوں واپس دینا چاہیدا اے جدوں دوویں دلیلاں~$\fn{true}$
ہون۔ جیویں اُتے دسیا اے، صداقتی قدراں چُنن والے فنکشن نیں۔
سو جدوں $x$ صداقتی قدر ہووے تے اینوں دو دلیلاں اُتے لائیے،
تاں پہلی دلیل نتیجہ ہووے گی جے $x$، $\fn{true}$ اے، تے ہور
صورت وچ دوجی دلیل۔ ہُن $\fn{And}$ اپنی دو دلیلاں $x$
تے $y$ لیندا اے تے $y$ تے $\fn{false}$ نوں اپنی پہلی
دلیل~$x$ اُتے لاؤندا اے۔ جے $x$ صداقتی قدر اے، تاں
نتیجہ~$y$ اے جے $x$، $\fn{true}$ اے، تے $\fn{false}$ اے
جے $x$، $\fn{false}$ اے۔ سانوں ایہو نتیجہ چاہیدا سی۔

دھیان دیو کہ ایتھے اسیں منّ کے چلدے آں کہ $\fn{And}$ دیاں
دلیلاں صرف صداقتی قدراں نیں۔ جے اینوں ہور ٹرم دِتّے جان،
تاں ہو سکدا اے کہ نتیجہ، یعنی نارمل روپ جے موجود ہووے، صداقتی
قدر نہ ہووے۔

\begin{prob}
  فنکشن $\fn{Or}$ تے $\fn{Xor}$ متعین کرو، جیہڑے شمولی تے
  اخراجی فصل دے صداقتی فنکشن ظاہر کردے نیں۔ ایس لئی صداقتی
  قدراں نوں $\lambd$-ٹرماں دے روپ وچ لکھن دا طریقہ ورتو۔
\end{prob}
\endgroup
% /OLP-0377
%\olimport{lists}
% OLP-0378 content/lambda-calculus/lambda-definability/primitive-recursive-functions.tex
\begingroup
\def\olfilename{primitive-recursive-functions}
\expandafter\def\csname ol@sectioncs\endcsname{section}

\olfileid{lam}{ldf}{prf}
\olsection{بنیادی بازگشتی فنکشن \usetoken{S}{lambda definable} نیں}

یاد کرو کہ بنیادی بازگشتی فنکشن اوہ نیں جیہڑے مُڈھلے فنکشوناں
$\Zero$، $\Succ$ تے $\Proj{n}{i}$ توں ترکیب تے بنیادی بازگشت
نال متعین کیتے جا سکدے نیں۔

\begin{lem}
  \ollabel{lem:basic}
  مُڈھلے بنیادی بازگشتی فنکشن $\Zero$، $\Succ$ تے دلیلاں
  وچوں اک چُنن والے فنکشن~$\Proj{n}{i}$، !!{lambda definable} نیں۔
\end{lem}

\begin{proof}
ایہ ہِٹھلے ٹرماں نال !!{lambda defined} نیں۔
\begin{align*}
  \fn{Zero} & \ident \lambd[a][\lambd[fx][x]]\\
  \fn{Succ} & \ident \lambd[a][\lambd[fx][f (a f x)]] \\
  \fn{Proj}^n_i & \ident \lambd[x_0\dots x_{n-1}][x_i]
\end{align*}
\end{proof}

\begin{lem}
  \ollabel{lem:comp} فرض کرو $k$ دلیلاں والا فنکشن $f$ تے
  $n$ دلیلاں والے فنکشن $g_0, \dots, g_{k-1}$، ٹرماں
  $F$، $G_0$، \dots، $G_k$ نال !!{lambda definable} نیں تے
  $h$ اوہناں دی ترکیب نال متعین اے۔ فیر $H$، !!{lambda definable} اے۔
\end{lem}

\begin{proof}
  $h$ ہِٹھلے ٹرم نال !!{lambda defined} ہو سکدا اے۔
  \[
  H \ident \lambd[x_0 \dots x_{n-1}][F\, (G_0 x_0 \dots
    x_{n-1}) \dots (G_{k-1} x_0 \dots x_{n-1})]
  \]
  ایس گل دی پڑتال مشق لئی چھڈدے آں۔
\end{proof}

\begin{prob}
  \olref[lam][ldf][prf]{lem:comp} دا ثبوت پورا کرن لئی دکھاؤ
  کہ $H\num{n_0}\dots\num{n_{n-1}} \red \num{h(n_0, \dots,
    n_{n-1})}$۔
\end{prob}

دھیان دیو کہ \olref{lem:comp} وچ ایہ شرط نہیں سی کہ $f$
تے $g_0$، \dots،~$g_{k-1}$ بنیادی بازگشتی ہون۔ شرط بس ایہ
سی کہ اوہ ہر انپٹ اُتے معین ہون تے !!{lambda definable} ہون۔

\begin{lem}\ollabel{lem:prim}
  فرض کرو $f$ اک $n$ دلیلاں والا فنکشن اے تے~$g$ اک $n+2$
  دلیلاں والا فنکشن اے۔ ایہ ٹرم $F$ تے~$G$ نال !!{lambda definable}
  نیں، تے فنکشن~$h$ نوں $f$ تے $g$ توں بنیادی بازگشت نال
  متعین کیتا اے۔ فیر $h$ وی !!{lambda definable} اے۔
\end{lem}

\begin{proof}
  یاد کرو کہ $h$ دی تعریف ایویں اے۔
  \begin{align*}
    h(x_1, \dots, x_n, 0) &= f(x_1, \dots, x_n)\\
    h(x_1, \dots, x_n, y+1) & = h(x_1, \dots, x_n, y, h(x_1, \dots, x_n, y)).
  \end{align*}
  سادے لفظاں وچ بنیادی بازگشتی تعریف فنکشن $h$ دا اطلاق
  $y$ واری دہراؤندی اے تے اینوں $f(x_1,
  \dots, x_n)$ اُتے لاؤندی اے۔ ایہ چرچ عدداں دی تعریف ورگا
  اے، جیہڑے آپ وی دہراؤ کرن والے ٹرماں وجوں متعین نیں۔

  سوکھ لئی اسیں اک ودھیک دلیل~$x$ واسطے تعریف تے ثبوت دیندے آں۔
  فنکشن $h$ ایویں !!{lambda defined} اے۔
  \begin{align*}
    H \ident & \lambd[x][\lambd[y][\fn{Snd} (y
        D \tuple{\num{0}, F x})]]
    \intertext{جتھے}
    D \ident & \lambd[p][\tuple{\fn{Succ} (\fn{Fst}\, p), (G x
          (\fn{Fst}\, p) (\fn{Snd}\, p))}]
  \end{align*}
  ہر قدم دی حالت اک جوڑا اے، جیہدا پہلا رکن اوہ ویلے دا $y$
  تے دوجا اوہدے نال جُڑی~$h$ دی قدر اے۔ تکرار دے ہر قدم
  اُتے اسیں $y$ تے $h$ دیاں نویاں قدراں دا جوڑا بناؤندے آں۔
  تکرار مکّن اُتے جوڑے دا دوجا رکن واپس دیندے آں۔ اوہی
  $h$ دی آخری قدر اے۔ ہُن ثابت کردے آں کہ ایہ واقعی بنیادی
  بازگشت دی نمائندگی اے۔

  اسیں ثابت کرنا اے کہ ہر $n$ تے $m$ لئی $H\,\num{n}\,\num{m} \red
  \num{h(n,m)}$۔ ایس لئی پہلاں دکھاؤندے آں کہ جے $D_n \ident
  \Subst{D}{\num{n}}{x}$، تاں $D_n^m \tuple{\num{0}, F\, \num{n}} \red
  \tuple{\num{m}, \num{h(n, m)}}$۔ اسیں~$m$ اُتے استقرا نال اگّے ودھدے آں۔

  جے $m=0$، تاں سانوں $D_n^0 \tuple{\num{0}, F\, \num{n}} \red
  \tuple{\num{0}, \num{h(n, 0)}}$ چاہیدا اے۔ پر $D_n^0 \tuple{\num{0}, F\,
    \num{n}}$ بس $\tuple{\num{0}, F\, \num{n}}$ اے۔ کیوں جے $F$،
  ~$f$ نوں !!{lambda define}s، ایس لئی ایہدی تقلیل $\tuple{\num{0},
    \num{f(n)}}$ تک ہوندی اے۔ تے کیوں جے $f(n) = h(n, 0)$، ایہ
  $\tuple{\num{0},
    \num{h(n,0)}}$ اے۔

  ہُن فرض کرو $D_n^m \tuple{\num{0}, F\, \num{n}} \red
  \tuple{\num{m}, \num{h(n, m)}}$۔ اسیں دکھاؤنا اے کہ
  $D_n^{m+1} \tuple{\num{0}, F\, \num{n}} \red
  \tuple{\num{m+1}, \num{h(n, m+1)}}$۔
  \begin{align*}
    D_n^{m+1} \tuple{\num{0}, F\, \num{n}}
    & \ident D_n(D_n^{m} \tuple{\num{0}, F\, \num{n}})\\
    & \red D_n\,\tuple{\num{m}, \num{h(n, m)}} \text{\qquad (استقرائی مفروضے توں)}\\
    & \ident (\lambd[p][
      \tuple{
        \fn{Succ} (\fn{Fst}\, p), (G \, \num{n} (\fn{Fst}\, p) (\fn{Snd}\, p))
    }]) \tuple{\num{m}, \num{h(n,m)}}\\
    & \redone
    \tuple{\fn{Succ} (\fn{Fst}\, \tuple{\num{m}, \num{h(n,m)}}),\\
      & \qquad (G \, \num{n} (\fn{Fst}\, \tuple{\num{m}, \num{h(n,m)}}) (\fn{Snd}\, \tuple{\num{m}, \num{h(n,m)}}))}\\
    & \red
    \tuple{\fn{Succ}\, \num{m}, (G \, \num{n} \,\num{m} \, \num{h(n,m)})}\\
    & \red \tuple{\num{m+1}, \num{g(n, m, h(n, m))}}
  \end{align*}
  کیوں جے $g(n, m, h(n, m)) = h(n, m+1)$، ایس لئی ایہ قدم پورا ہویا۔

  آخر وچ ایہ تقلیل ویکھو۔
  \begin{align*}
    H\,\num n\,\num m &\ident \lambd[x][\lambd[y][\fn{Snd} (y
        (\lambd[p].\tuple{\fn{Succ} (\fn{Fst}\, p), (G\, x\,
          (\fn{Fst}\, p)\, (\fn{Snd}\, p))}) \tuple{\num{0}, F x})]]\\
    & \qquad\,\num n\, \num m\\
    & \red \fn{Snd} (\num m \,
    \underbrace{(\lambd[p].\tuple{\fn{Succ} (\fn{Fst}\, p), (G \,\num n\,
      (\fn{Fst}\, p) (\fn{Snd}\, p))})}_{D_n} \tuple{\num{0}, F \num n})\\
    & \ident \fn{Snd} (\num m \, D_n\, \tuple{\num{0}, F \num n})\\
    & \red \fn{Snd} \,(D_n^m  \tuple{\num{0}, F \num n})\\
    & \red \fn{Snd} \,\tuple{\num{m}, \num{h(n, m)}} \\
    & \red \num{h(n,m)}.
  \end{align*}
\end{proof}

\begin{prop}
  ہر بنیادی بازگشتی فنکشن !!{lambda definable} اے۔
\end{prop}

\begin{proof}
  \olref{lem:basic} مطابق سارے مُڈھلے فنکشن !!{lambda definable}
  نیں۔ \olref{lem:comp} تے \olref{lem:prim} مطابق
  !!{lambda definable} فنکشن ترکیب تے بنیادی بازگشت ہِٹھ بند نیں۔
\end{proof}
\endgroup
% /OLP-0378
% OLP-0379 content/lambda-calculus/lambda-definability/fixpoints.tex
\begingroup
\def\olfilename{fixpoints}
\expandafter\def\csname ol@sectioncs\endcsname{section}

\olfileid{lam}{ldf}{fp}

\olsection{مقرر نقطے}

فرض کرو اسیں عاملی فنکشن بازگشت نال ٹرم $\fn{Fac}$ دے طور
اُتے متعین کرنا چاہندے آں، جیہدی خاصیت ایہ ہووے۔
\[
\fn{Fac} \ident \lambd[n][\fn{IsZero}\, n\, \num 1 (\fn{Mult}\, n (\fn{Fac}(\fn{Pred} \, n)))]
\]
یعنی $n$ دا عامل $1$ اے جے $n = 0$، تے ہور صورت وچ
$n$ گُنا $n-1$ دا عامل۔ صاف اے کہ ٹرم $\fn{Fac}$ نوں
ایویں متعین نہیں کر سکدے، کیوں جے سجے پاسے $\fn{Fac}$ آپ
وی آؤندا اے۔ اپنے آپ دا حوالہ دین والیاں ایہو جیہیاں بازگشتی
تعریفاں لامبڈا کیلکولس دا حصہ نہیں۔ مثلاً ٹرم دی ایس تعریف وچ
\[
\fn{Mult} \ident \lambd[ab][a (\fn{Add}\, a) 0]
\]
سجے پاسے صرف پہلاں توں متعین ٹرم آؤندے نیں، جیویں $\fn{Add}$۔
$\fn{Add}$ نوں اوہدی تعریف والا ٹرم رکھ کے ہمیشہ ہٹا سکدے آں۔
ایس نال $\fn{Mult}$ خالص لامبڈا ٹرم بن جاندا اے۔ جے $\fn{Add}$
وچ وی متعین ٹرم ہون، جیویں $\fn{Add}'$ وچ نیں، تاں ایس عمل نوں
جاری رکھ کے آخر اک خالص لامبڈا ٹرم تک پہنچ سکدے آں۔

پر اُپر دِتّی $\fn{Fac}$ ورگی بازگشتی تعریف لئی ایہ گل
سچ نہیں۔ جے سجے پاسے $\fn{Fac}$ دی آمد دی تھاں $\fn{Fac}$
دی اپنی تعریف رکھیے، تاں سانوں ایہ ملدا اے۔
\begin{align*}
  \fn{Fac} & \ident \lambd[n][\fn{IsZero}\, n\, \num{1}] \\
    & \qquad (\fn{Mult}\, n
  ((\lambd[n][\fn{IsZero} \, n \, \num 1\, (\fn{Mult}\, n\, (\fn{Fac}
    (\fn{Pred}\, n)))]) (\fn{Pred}\, n)))
\end{align*}
تے سجے پاسے $\fn{Fac}$ ہلے وی موجود اے۔ صاف اے کہ ایہ عمل
دہراؤن نال تعریف لمّی ہوندی جاندی اے تے کدے خالص لامبڈا ٹرم
نہیں ملدا۔ سو عاملی فنکشن، یا عام طور اُتے بازگشتی فنکشن،
ایس طریقے نال متعین نہیں کیتے جا سکدے۔

فیر وی بازگشتی تعریف توں اک گل پتا لگدی اے۔ جے $f$ عاملی
فنکشن دا ٹرم ہووے، تاں ہِٹھلا ٹرم
\[
\fn{Fac}' \ident \lambd[g][\lambd[n][\fn{IsZero} \, n \, \num 1\, (\fn{Mult} \, n\, (g (\fn{Pred} n)))]]
\]
ٹرم $f$ اُتے لاؤن نال، یعنی $\fn{Fac}'\,f$، عاملی فنکشن ای
ظاہر کردا اے۔ یعنی جے $\fn{Fac}'$ نوں ایہو جیہا فنکشن
سمجھیے جیہڑا فنکشن لیندا تے فنکشن واپس دیندا اے، تاں
$\fn{Fac'}\, f$ دی قدر بس~$f$ اے، بشرطیکہ $f$ عاملی
فنکشن ہووے۔ فنکشن~$f$ نوں $\fn{Fac}' \, f \equal[\beta] f$
ہوون اُتے $\fn{Fac}'$ دا \emph{مقرر نقطہ} آکھدے آں۔
سو عاملی فنکشن $\fn{Fac}'$ دا مقرر نقطہ اے۔

لامبڈا کیلکولس وچ ایہو جیہے ٹرم نیں جیہڑے دِتّے ٹرم دے مقرر
نقطے کڈھدے نیں۔ ایہناں نال $\fn{Fac}'$ ورگے ٹرم توں عاملی
فنکشن دی تعریف بنائی جا سکدی اے۔

\begin{defn}
  \ollabel{defn:Turing-Y}
  \emph{وائی کمبینیٹر} ایہ ٹرم اے۔
  \[
  Y \ident (\lambd[ux][x(uux)])(\lambd[ux][x(uux)]).
  \]
\end{defn}

\begin{thm}
  $Y$ دی خاصیت اے کہ $Yg \red g(Yg)$ ہر ٹرم $g$ لئی اے۔
  سو $Yg$ ہمیشہ~$g$ دا مقرر نقطہ اے۔
\end{thm}

\begin{proof}
  $(\lambd[ux][x(uux)])$ نوں مختصر~$U$ لکھ لئیے، تاں جو
  $Y \ident UU$۔ فیر
  \begin{align*}
    Y g &\ident (\lambd[ux][x(uux)])U\,g \\
    &\red (\lambd[x][x(UUx)])g\\
    & \red g(UUg) \ident g(Yg).
  \end{align*}
  کیوں جے $g(Yg)$ تے $Yg$ دوہاں دی تقلیل $g(Yg)$ تک ہوندی اے،
  ایس لئی $g(Yg) \equal[\beta]
  Yg$، تے $Yg$،~$g$ دا مقرر نقطہ اے۔
\end{proof}

کیوں جے $Yg$ اک ریڈیکس اے، ایس لئی تقلیل بغیر مُکے جاری رہ
سکدی اے۔
\begin{align*}
  Y g &\red g (Y g) \\
    &\red g (g (Y g)) \\
    &\red g(g (g (Y g)))\\
    &\ldots
\end{align*}
سو اسیں $Yg$ نوں $g$ دا اپنے آپ اُتے لامتناہی وار اطلاق
سمجھ سکدے آں۔ جے~$g$ نوں ایس اُتے اک وار ہور لائیے، تاں
ایویں آکھ لو کہ کوئی وادھا نہیں ہوندا۔ $g$ نوں~$g$ دے
~$Yg$ اُتے لامتناہی وار اطلاق دے نتیجے اُتے لاؤنا ہلے وی
$g$ دا~$Yg$ اُتے لامتناہی وار اطلاق ای اے۔

دھیان دیو کہ $\beta$-تقلیل دے قدماں دا اُپرلا سلسلہ
~$Yg$ توں شروع ہوندا اے تے لامتناہی اے۔ سو جے $Yg$ نوں
کسے ٹرم اُتے لائیے، یعنی $(Yg)N$ ویکھیے، تاں اوہ ٹرم وی
لامتناہی وکھرے ٹرماں تک گھٹدا اے، یعنی $(g(Yg))N$،
$(g(g(Yg)))N$، \dots۔ فیر وی ہو سکدا اے کہ تقلیل دے
قدماں دا کوئی \emph{ہور} سلسلہ نارمل روپ اُتے مُک جاوے۔

مثال لئی عاملی فنکشن ویکھو۔ $\fn{Fac}$ دی تعریف
$Y\, \fn{Fac}'$ دے طور اُتے کرو، یعنی $\fn{Fac'}$ دے مقرر
نقطے وجوں۔ فیر
\begin{align*}
  \fn{Fac}\,\num 3
  &\red Y\, \fn{Fac}' \, \num 3 \\
  &\red \fn{Fac}' (Y \,\fn{Fac}') \, \num 3 \\
  &\ident (\lambd[x][\lambd[n][\fn{IsZero} \, n\, \num 1\,
      (\fn{Mult}\, n\, (x (\fn{Pred}\, n)))]]) \, \fn{Fac} \, \num 3\\
  & \red \fn{IsZero} \, \num 3\, \num 1\,
      (\fn{Mult}\, \num 3\, (\fn{Fac} (\fn{Pred}\, \num 3))) \\
  &\red  \fn{Mult}\, \num 3 \, (\fn{Fac} \, \num 2).
  \intertext{ایسے طرح}
  \fn{Fac}\,\num 2
  &\red  \fn{Mult}\, \num 2 \, (\fn{Fac} \, \num 1) \\
  \fn{Fac}\,\num 1
  &\red  \fn{Mult}\, \num 1 \, (\fn{Fac} \, \num 0)
  \intertext{پر}
  \fn{Fac}\,\num 0
  &\red \fn{Fac}' (Y \,\fn{Fac}') \, \num 0 \\
  &\ident (\lambd[x][\lambd[n][\fn{IsZero} \, n\, \num 1\,
      (\fn{Mult}\, n\, (x (\fn{Pred}\, n)))]]) \, \fn{Fac} \, \num 0\\
  & \red \fn{IsZero} \, \num 0\, \num 1\,
      (\fn{Mult}\, \num 0\, (\fn{Fac} (\fn{Pred}\, \num 0))).\\
  &\red \num 1.
  \intertext{سو سارے قدم جوڑ کے}
  \fn{Fac}\,\num 3
  &\red  \fn{Mult}\, \num 3 \,
  (\fn{Mult} \, \num 2\, (\fn{Mult} \, \num 1\, \num 1)).
\end{align*}

جیہڑی گل $\fn{Fac'}$ لئی اے، اوہی کسے وی بازگشتی تعریف لئی
اے۔ فرض کرو ساڈے کول اک بازگشتی مساوات اے۔
\begin{align*}
g\,x_1\dots x_n & \equal[\beta] N
\intertext{ایتھے $N$ وچ~$g$ تے $x_1$، \dots،~$x_n$ آ سکدے نیں۔
  فیر ہمیشہ اک ایہو جیہا ٹرم~$G \ident (Y \lambd[g][\lambd[x_1\dots x_n][N]])$ اے جیہدے لئی}
G\,x_1 \dots x_n & \equal[\beta] \Subst{N}{G}{g}.
\intertext{کیوں جے مقرر نقطے دے قضیے مطابق}
G \ident (Y \lambd[g][\lambd[x_1\dots x_n][N]]) & \red \lambd[g][\lambd[x_1\dots x_n][N]](Y \lambd[g][\lambd[x_1\dots x_n][N]])\\
& \ident (\lambd[g][\lambd[x_1\dots x_n][N]])\,G
\intertext{تے ایس توں}
G\,x_1 \dots x_n
& \red (\lambd[g][\lambd[x_1\dots x_n][N]])\,G\,x_1\dots x_n\\
& \red (\lambd[x_1\dots x_n][\Subst{N}{G}{g}])\,x_1\dots x_n\\
  & \red \Subst{N}{G}{g}.
\end{align*}

$Y$ کمبینیٹر، جیہڑا \olref{defn:Turing-Y} وچ دِتّا اے، ایلن ٹیورنگ دا اے۔ الونزو
چرچ نے اک ہور روپ تجویز کیتا سی جیہنوں اسیں~$Y_C$ آکھاں گے۔
\[
Y_C \ident \lambd[g][(\lambd[x][g(xx)])(\lambd[x][g(xx)])].
\]
چرچ دا کمبینیٹر ٹیورنگ دے کمبینیٹر نالوں ذرا کمزور اے۔
ایہدے لئی $Yg
\equal[\beta] g(Yg)$ اے، پر $Yg \bred g(Yg)$ نہیں۔ منّ لو $V$
ٹرم $\lambd[x][g(xx)]$ اے، تاں جو $Y_C \ident \lambd[g][VV]$۔
فیر
\begin{align*}
  VV & \ident (\lambd[x][g(xx)])V \red g(VV)
  \text{ تے ایس لئی}\\
  Y_C g & \ident (\lambd[g][VV])g \red VV \red g(VV), \text{ پر ایہ وی}\\
  g(Y_C g) & \ident g((\lambd[g][VV])g) \red g(VV).
\end{align*}
دوجے لفظاں وچ، $Y_Cg$ تے $g(Y_Cg)$ دوہاں دی تقلیل اکّو ٹرم
$g(VV)$ تک ہوندی اے۔ سو $Y_Cg
\equal[\beta] g(Y_Cg)$۔ کئی ورتوں لئی ایہ کافی ہوندا اے۔
\endgroup
% /OLP-0379
% OLP-0380 content/lambda-calculus/lambda-definability/minimization.tex
\begingroup
\def\olfilename{minimization}
\expandafter\def\csname ol@sectioncs\endcsname{section}

\olfileid{lam}{ldf}{min}
\olsection{گھٹ توں گھٹ حل لبھنا}

عام بازگشتی فنکشن اوہ نیں جیہڑے مُڈھلے فنکشوناں $\Zero$،
$\Succ$، $\Proj{n}{i}$ توں ترکیب، بنیادی بازگشت تے باقاعدہ
گھٹ توں گھٹ حل لبھن نال حاصل کیتے جا سکدے نیں۔ سارے عام
بازگشتی فنکشوناں نوں !!{lambda definable} دکھاؤن لئی سانوں
دکھاؤنا پئے گا کہ کسے !!a{lambda definable} فنکشن توں باقاعدہ گھٹ توں گھٹ حل لبھ کے بنایا
ہویا فنکشن آپ وی !!{lambda definable} اے۔

\begin{lem}
  \ollabel{lem:min} جے $f(x_1, \dots, x_k, y)$ باقاعدہ اے تے
  !!{lambda definable} اے، تاں ہِٹھلی تعریف والا~$g$
  \[
  g(x_1, \dots, x_k) = \umin{y}{f(x_1,\dots,x_k, y) = 0}
  \]
  وی !!{lambda definable} اے۔
\end{lem}

\begin{proof}
  فرض کرو لامبڈا ٹرم~$F$، $\lambda$ نال باقاعدہ فنکشن
  $f(\vec x, y)$ دی تعریف دیندا اے۔ ~$h$ نوں
  !!{lambda define} کرن لئی اسیں تلاش دا فنکشن تے مقرر نقطے دا
  کمبینیٹر ورتدے آں۔
  \begin{align*}
    \fn{Search} & \ident \lambd[g][\lambd[f\,\vec{x}\,y][
        \fn{IsZero} (f\, \vec{x}\, y)\, y\, (g\, \vec{x} (\fn{Succ}\, y)]]\\
    H & \ident \lambd[\vec x][(Y \, \fn{Search}) F\, \vec{x}\, \num{0}],
  \end{align*}
  جتھے $Y$ کوئی وی مقرر نقطے دا کمبینیٹر اے۔ سادے لفظاں وچ،
  $\fn{Search}$ اپنے آپ نوں سدّن والا فنکشن اے۔~$y$ توں
  شروع کر کے پرکھو کہ $f\, \vec x\, y$ صفر اے کہ نہیں۔
  جے صفر اے، تاں~$y$ واپس دیو۔ ہور صورت وچ $\fn{Succ}\, y$
  نال اوہو فنکشن فیر سدو۔ سو $(Y \, \fn{Search}) F
  \num{n_1}\dots\num{n_k}\,\num{0}$ اوہ سبھ توں گھٹ~$m$
  واپس دیندا اے جیہدے لئی $f(n_1,
  \dots, n_k, m) = 0$۔

  خاص کر کے، دھیان دیو۔
  \begin{align*}
    (Y \, \fn{Search}) F \num{n_1}
    \dots\num{n_k}\, \num{m} & \red \num{m}
    \intertext{جے $f(n_1, \dots,
      n_k, m) = 0$، یا}
    & \red (Y \, \fn{Search}) F\, \num{n_1} \dots\num{n_k}\, \num{m+1}
    \intertext{ہور صورت وچ۔ کیوں جے $f$ باقاعدہ اے، $f(n_1, \dots, n_k, y)
      = 0$ کسے $y$ لئی، تے ایس لئی}
    (Y \, \fn{Search}) F \num{n_1} \dots\num{n_k}\,\num{0}
    & \red \num{h(n_1, \dots, n_k)}.
    \end{align*}
\end{proof}

\begin{prop}
  ہر عام بازگشتی فنکشن !!{lambda definable} اے۔
\end{prop}

\begin{proof}
 \olref[prf]{lem:basic} مطابق سارے مُڈھلے فنکشن
 !!{lambda definable} نیں۔ \olref[prf]{lem:comp}،
 \olref[prf]{lem:prim} تے \olref{lem:min} مطابق
 !!{lambda definable} فنکشن ترکیب، بنیادی بازگشت تے باقاعدہ
 گھٹ توں گھٹ حل لبھن ہِٹھ بند نیں۔
\end{proof}
\endgroup
% /OLP-0380
% OLP-0381 content/lambda-calculus/lambda-definability/partial-recursive-functions.tex
\begingroup
\def\olfilename{partial-recursive-functions}
\expandafter\def\csname ol@sectioncs\endcsname{section}

\olfileid{lam}{ldf}{par}
\olsection{جزوی بازگشتی فنکشن \usetoken{S}{lambda definable} نیں}

جزوی بازگشتی فنکشن اوہ نیں جیہڑے مُڈھلے فنکشوناں توں
ترکیب، بنیادی بازگشت تے بغیر حد گھٹ توں گھٹ حل لبھن نال
حاصل ہوندے نیں۔ ایہ عام بازگشتی فنکشوناں توں ایس گل وچ
وکھرے نیں کہ بغیر حد دی تلاش وچ ورتے فنکشوناں دا باقاعدہ
ہونا ضروری نہیں۔ باقاعدگی دی شرط نہ ہووے، تاں گھٹ توں
گھٹ حل لبھن نال متعین فنکشن کدے غیر معین وی ہو سکدا اے۔

پہلی نظر وچ لگ سکدا اے کہ جیہڑے طریقیاں نال (ہر انپٹ
اُتے معین) سارے عام بازگشتی فنکشن !!{lambda definable}
دکھائے سن، اوہناں نال سارے جزوی بازگشتی فنکشن وی
!!{lambda definable} ثابت ہو سکدے نیں۔ مثلاً، $f$ دی $g$
نال ترکیب $\lambd[x][F (G x)]$ نال !!{lambda defined} اے،
جے $f$ تے $g$ بالترتیب ٹرماں $F$ تے~$G$ نال
!!{lambda defined} ہون۔ پر جے فنکشن جزوی ہون، تاں ایس
وچ دِقّت اے۔ جدوں $g(x)$ غیر معین اے، اوہدا مطلب اے کہ
$G x$ دا کوئی نارمل روپ نہیں۔ زیادہ تر صورتاں وچ ایس دا
مطلب اے کہ $F (G
x)$ دا وی کوئی نارمل روپ نہیں، تے سانوں ایہو چاہیدا اے۔
پر جے $F$، $\lambd[x][\lambd[y][y]]$ ہووے، تاں $F (G x)$
دا نارمل روپ ($\lambd[y][y]$) موجود اے۔

ایہ دِقّت حل ہو سکدی اے۔ سارے جزوی بازگشتی فنکشوناں نوں
!!{lambda define} کرن دے ایہو جیہے طریقے موجود نیں جیہناں
وچ غیر معین قدراں نوں بغیر نارمل روپ والے ٹرم ظاہر کردے
نیں۔ پر ایہ طریقے اوہناں طریقیاں نالوں کجھ ودھ گُنجھلدار
نیں تے بدیہی طور اُتے گھٹ سمجھ آؤندے نیں جیہڑے اسیں عام
بازگشتی فنکشوناں لئی ورتے سن۔ ایتھے اسیں قضیہ ثبوت توں
بغیر لکھدے آں۔

\begin{thm}
  سارے جزوی بازگشتی فنکشن !!{lambda definable} نیں۔
\end{thm}
\endgroup
% /OLP-0381
% OLP-0382 content/lambda-calculus/lambda-definability/lambda-definable-recursive.tex
\begingroup
\def\olfilename{lambda-definable-recursive}
\expandafter\def\csname ol@sectioncs\endcsname{section}

\olfileid{lam}{dfl}{ldr}
\olsection{\usetoken{S}{lambda definable} فنکشن بازگشتی نیں}

سارے جزوی بازگشتی فنکشن !!{lambda definable} نیں، تے اُلٹی گل
وی سچ اے۔ یعنی سارے !!{lambda definable} فنکشن جزوی بازگشتی
نیں۔

\begin{thm}
  \ollabel{thm:lambda-computable} جے جزوی فنکشن $f$،
  !!{lambda definable} اے، تاں اوہ جزوی بازگشتی اے۔
\end{thm}

\begin{proof}
  اسیں صرف ثبوت دا خاکہ دیندے آں۔ پہلاں $\lambd$-ٹرماں نوں
  عددی روپ دیندے آں، یعنی $\lambd$-ٹرماں نوں قاعدے نال گوڈل
  نمبر دیندے آں۔ ایس لئی سلسلیاں دا عام کوڈ ورتدے آں جیہڑا
  اولی عدداں دیاں گھاتاں نال بندا اے۔ فیر اک جزوی بازگشتی
  فنکشن $\fn{normalize}(t)$ متعین کردے آں جیہڑا دلیل وجوں
  لامبڈا ٹرم دا گوڈل نمبر~$t$ لیندا اے۔ جے ٹرم دا نارمل روپ
  ہووے، تاں اوہدا گوڈل نمبر واپس دیندا اے۔ ہور صورت وچ
  غیر معین اے۔ فیر دو جزوی بازگشتی فنکشن $\fn{toChurch}$
  تے $\fn{fromChurch}$ متعین کردے آں جیہڑے قدرتی عدداں
  نوں اوہناں دے چرچ عدداں دے گوڈل نمبراں وچ تے واپس بدلدے نیں۔

  ایہناں بازگشتی فنکشوناں نال اسیں فنکشن~$f$ نوں جزوی بازگشتی
  فنکشن وجوں متعین کر سکدے آں۔ کوئی $\lambd$-ٹرم~$F$ اے
  جیہڑا~$f$ نوں !!{lambda define}s۔ $f(n_1, \dots, n_k)$ دا
  حساب کرن لئی پہلاں $\fn{toChurch}(n_i)$ نال متعلقہ چرچ
  عدداں دے گوڈل نمبر حاصل کرو۔ ایہناں نوں $\Gn{F}$ نال
  جوڑ کے ٹرم $F \num{n_1}\dots\num{n_k}$ دا گوڈل نمبر بناؤ۔
  ہُن ایس گوڈل نمبر اُتے $\fn{normalize}$ لاؤ۔ جے
  $f(n_1, \dots, n_k)$ معین اے، تاں $F \num{n_1}\dots\num{n_k}$
  دا نارمل روپ اے تے اوہ لازماً چرچ عدد اے۔ ہور صورت وچ
  اوہدا کوئی نارمل روپ نہیں، تے ایس لئی
  \[\fn{normalize}(\Gn{F\num{n_1}\dots\num{n_k}})\]
  غیر معین اے۔ آخر وچ نارمل کیتے ٹرم دے گوڈل نمبر اُتے
  $\fn{fromChurch}$ لاؤ۔
\end{proof}
\endgroup
% /OLP-0382

% OLP-0650 content/lambda-calculus/lambda-definability/lists.tex
\begingroup
\def\olfilename{lists}
\expandafter\def\csname ol@sectioncs\endcsname{section}

\olfileid{lam}{ldf}{lst}
\olsection{فہرستاں}

فہرست $[M_1, M_2, \ldots, M_n]$ دی تعریف ایویں کیتی جا سکدی اے:
\[
  [M_1, M_2, \ldots, M_n] \ident \lambd[sz][s M_1 (s M_2 (\ldots (s M_n z)))]
\]
عام بولی وچ، فہرست دی نمائندگی کرن والا ٹرم اوہدا «قدر جوڑدے جان والا» فنکشن اے۔ ایہ دو ٹرم لیندا اے۔ پہلا فنکشن اے جیہڑا ہُن تک جوڑی قدر تے اک نواں رکن لے کے نویں جوڑی قدر موڑدا اے۔ دوجا جوڑ شروع کرن والی قدر اے۔ فہرست والا فنکشن اک اک کر کے رکن جوڑدا اے تے آخر حاصل موڑدا اے۔

\begin{digress}
  جے تسی فنکشنی پروگرامنگ توں واقف او، تاں ایہ تعریف \emph{سجّے پاسے دی تہہ} ورگی لگے گی۔ فہرست نوں اوہدے رکناں دی سجّے پاسے دی تہہ والے فنکشن وجوں تعریف کیتا جاندا اے۔
\end{digress}

ایس کوڈ بندی دی بنیاد اُتے اسیں کجھ کم آؤن والے فنکشن سوکھے ای تعریف کر سکدے آں:
\begin{align*}
  \fn{Sum} & \ident
  \lambd[l][l \, (\lambd[xa][\fn{Add} \, x \, a])\,  \num{0}]\\
  \fn{Len} & \ident
  \lambd[l][l \, (\lambd[xa][\fn{Add} \, a \, \num{1}]) \num{0}]
\end{align*}

$\fn{Sum}$ چرچ دیاں عددی عبارتاں دی فہرست دا جوڑ کڈھدا اے۔ ایہ فہرست اُتے اک اک رکن جوڑ کے کم کردا اے۔ شروعاتی قدر $\num{0}$ اے تے ہر رکن لئی نویں قدر وجوں $\fn{Add} \, x\, a$ موڑدا اے (جتھے $x$ ہُن والا رکن اے تے $a$ ہُن تک جوڑی قدر اے)۔ حاصل سارے رکناں دا جوڑ اے۔

\begin{prob}
  $\fn{Len}$ کیہ کردا اے؟ سمجھاؤ۔
\end{prob}

\begin{prob}
  اک فنکشن تعریف کرو جیہڑا فہرست دی ترتیب اُلٹا دے۔
\end{prob}
\endgroup
% /OLP-0650
\OLEndChapterHook
\endgroup
% /OLP-0373

\OLEndPartHook
\endgroup
% /OLP-0341

% OLP-0383 content/many-valued-logic/many-valued-logic.tex
\begingroup
\def\olfilename{many-valued-logic}
\expandafter\def\csname ol@sectioncs\endcsname{section}

\olpart{mvl}{کئی قدراں والی منطق}

\begin{editorial}
  ایس حصے وچ کئی قدراں والیاں قضیاتی منطقاں دا مسودہ اے۔
\end{editorial}

% OLP-0384 content/many-valued-logic/syntax-and-semantics/syntax-and-semantics.tex
\begingroup
\def\olfilename{syntax-and-semantics}
\expandafter\def\csname ol@sectioncs\endcsname{section}

\olchapter{mvl}{syn}{نحو تے معنیات}

% OLP-0385 content/many-valued-logic/syntax-and-semantics/introduction.tex
\begingroup
\def\olfilename{introduction}
\expandafter\def\csname ol@sectioncs\endcsname{section}

\olfileid{mvl}{syn}{int}

\olsection{جان پہچان}

کلاسیکی منطق وچ اسیں !!{formula}s نال کم کردے آں جیہڑے
!!{propositional variable}s توں قضیاتی رابطے $\lnot$، $\land$،
$\lor$، $\lif$ تے $\liff$ ورت کے بندے نیں۔ کلاسیکی منطق دی
معنیات متعین کرن لئی دو صداقتی قدراں $\True$ تے $\False$
ورتدے آں۔ !!{propositional variable}s دی تشریح !!a{valuation}~$\pAssign{v}$
وچ کردے آں۔ ایہ صداقتی قدراں $\True$، $\False$ نوں
!!{propositional variable}s نال جوڑدی اے۔ فیر ہر
!!{valuation} صداقتی قدر $\pValue{v}(!A)$ متعین کردی اے
کسے وی !!{formula}~$!A$ لئی۔ !!^a{formula}،
!!a{valuation}~$\pAssign{v}$ وچ پورا اتردا اے، یعنی $\pSat{v}{!A}$،
اُدوں تے صرف اُدوں جدوں $\pValue{v}(!A) =
\True$۔

کئی قدراں والیاں منطقاں کلاسیکی دو قدری منطق دی وسعت نیں۔
ایہناں وچ $\True$ تے $\False$ توں علاوہ ہور صداقتی قدراں وی
ہو سکدیاں نیں۔ سو کئی قدراں والی منطق وچ !!a{valuation}~$\pAssign{v}$
اوہ فنکشن اے جیہڑا ہر !!{propositional variable}~$p$ نوں ممکن
صداقتی قدراں وچوں اک دیندا اے۔ عام طور اُتے جائز صداقتی
قدراں دے سیٹ نوں~$V$ آکھاں گے۔ کلاسیکی منطق وی کئی قدراں
والی منطق اے جیہدے لئی $V =
\{\True, \False\}$۔ اوہدے وچ صداقتی قدر~$\pValue{v}(!A)$
رابطیاں دے جانے پہچانے صداقتی جدولاں نال کڈھی جاندی اے۔

جدوں ہور صداقتی قدراں شامل کردے آں، تاں اوہناں رابطیاں لئی
~$\pValue{v}(!A)$ کڈھن دے اک توں ودھ سُبھاوک طریقے ہوندے
نیں جنہاں نوں اسیں «تے»، «یا»، «نہیں» تے «جے ... تاں»
پڑھدے آں۔ سو کئی قدراں والی منطق نوں صداقتی قدراں دا سیٹ
تے ہر رابطے لئی ساڈے چُنے ہوئے \emph{صداقتی فنکشن} رل کے
متعین کردے نیں۔ کئی قدراں والی منطق~$\Log L$ لئی ایہ
فنکشن چُنے جان، تاں صداقتی قدر $\pValue{v}(!A)[\Log L]$
قدر بندی نال اکّو طرح متعین ہوندی اے، جیویں کلاسیکی منطق
وچ۔ کلاسیکی منطق وانگ کئی قدراں والیاں منطقاں
\emph{صداقتی فنکشوناں اُتے ٹکیاں} ہوندیاں نیں۔

معنیات دیاں ایہ بنیادی گلاں ہتھ آ جان، تاں اسیں ہمیشہ صادق
ہون، منطقی نتیجہ ہون تے پورا اتر سکّن دے ملدے جُلدے مفہوم
متعین کر سکدے آں۔ کلاسیکی منطق وچ !!a{formula} ہمیشہ صادق
اے جے اوہدی صداقتی قدر $\pValue{v}(!A) = \True$ ہووے ہر~$\pAssign{v}$
لئی۔ کئی قدراں والی منطق وچ اینوں وی کجھ عام کرنا پیندا اے۔
پہلی گل، صداقتی قدراں دے سیٹ~$V$ وچ~$\True$ ہونا ضروری نہیں۔
مثلاً کجھ منطقاں صداقتی قدراں وجوں عدد ورتدیاں نیں، جیویں
$0$ تے~$1$ وچکارلے سارے ناطق عدد۔ اوہ صورت وچ $1$ عام طور
اُتے~$\True$ دا کم کردا اے۔ ہور منطقاں وچ کئی صداقتی قدراں ایہ
کم کردیاں نیں۔ سو ہر کئی قدراں والی منطق لئی اسیں
\emph{مقرر قدراں} دا سیٹ~$V^+$ ضروری منّدے آں۔ فیر آکھ سکدے آں کہ
!!a{formula}، !!a{valuation}~$\pAssign{v}$ وچ پورا اتردا اے،
یعنی $\pSat{v}{!A}[\Log L]$، اُدوں تے صرف اُدوں جدوں
$\pValue{v}(!A)[\Log L] \in V^+$۔ !!^a{formula}~$!A$ منطق
دا ہمیشہ صادق فارمولا اے، یعنی $\Entails[\Log L] !A$،
اُدوں تے صرف اُدوں جدوں $\pValue{v}(!A) \in V^+$ ہووے ہر~$\pAssign{v}$
لئی۔ آخر وچ، آکھدے آں کہ $!A$، !!{formula}s دے سیٹ دا
منطقی نتیجہ اے، یعنی $\Gamma \Entails[\Log L] !A$، جے ہر
!!{valuation} جیہدے وچ~$\Gamma$ والے سارے !!{formula}s پورے
اتردے نیں، اوہدے وچ~$!A$ وی پورا اتردا اے۔
\endgroup
% /OLP-0385

% OLP-0386 content/many-valued-logic/syntax-and-semantics/connectives.tex
\begingroup
\def\olfilename{connectives}
\expandafter\def\csname ol@sectioncs\endcsname{section}

\olfileid{mvl}{syn}{con}

\olsection{زباناں تے رابطے}

کلاسیکی قضیاتی منطق تے کئی ہور منطقاں وچ \emph{قضیاتی مستقل}
تے \emph{رابطے} اک مقرر گنتی وچ ورتے جاندے نیں۔ مثلاً اسیں
ہِٹھلے نشان مُڈھلے نشان وجوں ورتدے آں۔
\begin{enumerate}
\tagitem{prvFalse}{!!{falsity} لئی قضیاتی مستقل~$\lfalse$۔}{}
\tagitem{prvTrue}{!!{truth} لئی قضیاتی مستقل~$\ltrue$۔}{}
\item منطقی رابطے۔
  \startycommalist
  \iftag{prvNot}{\ycomma $\lnot$ (نفی)}{}%
  \iftag{prvAnd}{\ycomma $\land$ (عطف)}{}%
  \iftag{prvOr}{\ycomma $\lor$ (فصل)}{}%
  \iftag{prvIf}{\ycomma $\lif$ (!!{conditional})}{}%
  \iftag{prvIff}{\ycomma $\liff$ (!!{biconditional})}{}%
\end{enumerate}
\iftag{defNot,defOr,defAnd,defIf,defIff,defTrue,defFalse,defEx,defAll}{%
اُپرلے مُڈھلے رابطیاں دے نال اسیں اوہ نشان وی ورتدے آں
جیہڑے مختصر لکھتاں وجوں متعین ہوندے نیں۔ مثلاً
\startycommalist
  \iftag{defNot}{\ycomma $\lnot$ (نفی)}{}%
  \iftag{defAnd}{\ycomma $\land$ (عطف)}{}%
  \iftag{defOr}{\ycomma $\lor$ (فصل)}{}%
  \iftag{defIf}{\ycomma $\lif$ (!!{conditional})}{}%
  \iftag{defIff}{\ycomma $\liff$ (!!{biconditional})}{}%
  \iftag{defFalse}{\ycomma $\lfalse$ (!!{falsity})}{}%
  \iftag{defTrue}{\ycomma $\ltrue$ (!!{truth})}۔}{}

کئی قدراں والیاں منطقاں وچ وی اوہی رابطے ورتے جاندے نیں۔
پر اکثر اکّو منطق وچ، مثلاً عطف دے وکھ وکھ روپ رکھنا کم
آؤندا اے۔ روپاں نوں وکھ رکھن لئی وکھرے نشان چاہیدے نیں۔
کجھ کئی قدراں والیاں منطقاں وچ ایسے رابطے وی ہوندے نیں
جنہاں دا کلاسیکی منطق وچ کوئی ہم ارز نہیں۔ ایس لئی اسیں
عام نالوں ذرا ودھ کھلھ نال تعریف کراں گے۔

\begin{defn}
  \emph{قضیاتی زبان} \emph{رابطیاں} دے سیٹ $\Lang L$ توں
  بندی اے۔ ہر رابطے $\star$ دی اک \emph{مقاماں دی گنتی} اے۔
  جیہدے مقاماں دی گنتی~$n$ اے، اوہنوں \emph{$n$-مقامی} رابطہ
  آکھدے آں۔ مقاماں دی گنتی~$0$ ہووے تاں اوہ رابطے
  \emph{مستقل} وی اکھواؤندے نیں۔ مقاماں دی گنتی~$1$ ہووے
  تاں \emph{اک مقامی}، تے مقاماں دی گنتی~$2$ ہووے تاں
  \emph{دو مقامی} اکھواؤندے نیں۔
\end{defn}

\begin{ex}
  قضیاتی منطق دی معیاری زبان $\Lang L_0$ وچ ہِٹھلے رابطے نیں۔
  نال اوہناں دے مقاماں دی گنتی دِتّی اے۔
  $\lfalse$~($0$)،
  $\lnot$~($1$)،
  $\land$~($2$)،
  $\lor$~($2$)،
  $\lif$~($2$)۔ اسیں جِنہاں منطقاں نوں ویکھاں گے، اوہناں
  وچوں بہتیاں وچ ایہو زبان ورتی جاندی اے۔ رواج مطابق کجھ
  منطقاں وچ کجھ رابطیاں دے نشان وکھرے نیں۔ مثلاً حاصلِ ضرب
  والی منطق وچ عطف دا نشان اکثر $\odot$ ہوندا اے~$\land$
  دی تھاں۔ کدے نواں عامل شامل کرنا سوکھا پیندا اے، جیویں
  تعیّن دا عامل $\triangle$ ($1$-مقامی)۔
\end{ex}
\endgroup
% /OLP-0386

% OLP-0387 content/many-valued-logic/syntax-and-semantics/formulas.tex
\begingroup
\def\olfilename{formulas}
\expandafter\def\csname ol@sectioncs\endcsname{section}

\olfileid{mvl}{syn}{fml}

\olsection{\usetoken{P}{formula}}

\begin{defn}[فارمولا]
\ollabel{defn:formulas}
سیٹ~$\Frm[L]$ وچ قضیاتی زبان~$\Lang L$ دے \emph{!!{formula}s}
ہوندے نیں۔ ایہ سیٹ استقرا نال ایویں متعین ہوندا اے۔
\begin{enumerate}
\item ہر !!{propositional variable}~$\Obj p_i$ اک ایٹمی
  !!{formula} اے۔
\item ہر $0$-مقامی رابطہ، یعنی قضیاتی مستقل، جیہڑا~$\Lang L$
  وچ اے، اک ایٹمی !!{formula} اے۔
\item جے $\star$ اک $n$-مقامی رابطہ اے جیہڑا~$\Lang L$ وچ اے، تے $!A_1$،
\dots، $!A_n$، !!{formula}s نیں، تاں $\star(!A_1, \dots, !A_n)$
!!a{formula} اے۔
\tagitem{limitClause}{ہور کوئی شے !!a{formula} نہیں۔}{}
\end{enumerate}
جے $\star$ اک $1$-مقامی رابطہ اے، تاں $\star(!A_1)$ نوں
اکثر بس $\star !A_1$ لکھدے آں۔ جے $\star$ اک $2$-مقامی
رابطہ اے، تاں $\star(!A_1,!A_2)$ نوں اکثر $(!A_1 \star !A_2)$
لکھدے آں۔
\end{defn}

رواج مطابق اسیں اکثر سبھ توں باہر دیاں قوساں بِن دسے چھڈ
دیندے آں۔

\begin{ex}
  معیاری زبان~$\Lang{L_0}$ وچ $\Obj p_1 \lif (\Obj p_1
  \land \lnot \Obj p_2)$ اک فارمولا اے۔ حاصلِ ضرب والی منطق
  دی زبان وچ اینوں $\Obj p_1 \lif (\Obj p_1 \odot
  \lnot \Obj p_2)$ لکھیا جاندا اے۔ جے $1$-مقامی $\triangle$
  زبان وچ شامل کریے، تاں ساڈے کول $\triangle (\Obj p_1
  \land \Obj p_2) \lif (\triangle \Obj p_1 \land \triangle \Obj p_2)$
  ورگے فارمولے وی ہون گے۔
\end{ex}
\endgroup
% /OLP-0387

% OLP-0388 content/many-valued-logic/syntax-and-semantics/matrices.tex
\begingroup
\def\olfilename{matrices}
\expandafter\def\csname ol@sectioncs\endcsname{section}

\olfileid{mvl}{syn}{mat}

\olsection{منطقی میٹرکس}

کئی قدراں والی منطق اوہدی زبان، صداقتی قدراں دے سیٹ~$V$،
مقرر صداقتی قدراں دے ذیلی سیٹ تے اوہدے رابطیاں دے صداقتی
فنکشوناں نال متعین ہوندی اے۔ ایہناں ساریاں شےواں نوں رلا کے
اک \emph{میٹرکس} آکھدے آں۔

\begin{defn}[منطقی میٹرکس]
\ollabel{defn:matrix}
منطق~$\Log L$ دا \emph{میٹرکس} ایہناں شےواں توں بندا اے۔
\begin{enumerate}
\item زبان~$\Lang L$ بناؤن والے رابطیاں دا سیٹ۔
\item صداقتی قدراں دا سیٹ $V \neq \emptyset$۔
\item مقرر صداقتی قدراں دا سیٹ $V^+ \subseteq V$۔
\item ہر $n$-مقامی رابطے $\star$ لئی جیہڑا $\Lang L$ وچ اے، صداقتی
فنکشن~$\tf{\star} : V^n \to V$۔ جے $n = 0$، تاں
$\tf{\star}$ بس~$V$ دا اک رکن اے۔
\end{enumerate}
\end{defn}

\begin{ex}
کلاسیکی منطق~\LogCL{} دا میٹرکس ایہناں شےواں توں بندا اے۔
\begin{enumerate}
  \item معیاری قضیاتی زبان $\Lang L_0$ جیہدے وچ
  $\lfalse$، $\lnot$، $\land$، $\lor$، $\lif$ نیں۔
  \item صداقتی قدراں دا سیٹ $V = \{\True, \False\}$۔
  \item $\True$ اکّو مقرر قدر اے، یعنی $V^+ = \{\True\}$۔
  \item $\lfalse$ لئی ساڈے کول $\tf{\lfalse} = \False$ اے۔
  باقی صداقتی فنکشن عام صداقتی جدولاں نال ملدے نیں۔
  \olref{fig:tf-CL} ویکھو۔
\end{enumerate}
\begin{figure}
  \begin{center}
      \begin{tabular}{c|c}
        $\tf{\lnot}$ & \\
        \hline
        $\True$ & $\False$ \\
        $\False$ & $\True$
      \end{tabular}
      \quad
      \begin{tabular}{c|cc}
        $\tf{\land}$ & $\True$ & $\False$ \\
        \hline
        $\True$ & $\True$ & $\False$ \\
        $\False$ & $\False$ & $\False$
      \end{tabular}
      \quad
      \begin{tabular}{c|cc}
        $\tf{\lor}$ & $\True$ & $\False$ \\
        \hline
        $\True$ & $\True$ & $\True$ \\
        $\False$ & $\True$ & $\False$
      \end{tabular}
      \quad
      \begin{tabular}{c|cc}
        $\tf{\lif}$ & $\True$ & $\False$ \\
        \hline
        $\True$ & $\True$ & $\False$ \\
        $\False$ & $\True$ & $\True$
      \end{tabular}
    \end{center}
    \caption{کلاسیکی منطق~$\LogCL$ دے صداقتی فنکشن۔}
    \ollabel{fig:tf-CL}
  \end{figure}
  \end{ex}
\endgroup
% /OLP-0388

% OLP-0389 content/many-valued-logic/syntax-and-semantics/valuations-sat.tex
\begingroup
\def\olfilename{valuations-sat}
\expandafter\def\csname ol@sectioncs\endcsname{section}

\olfileid{mvl}{syn}{val}

\olsection{\usetoken{P}{valuation} تے پورا اترنا}

\begin{defn}[!!^{valuation}s]
منّ لو $V$ صداقتی قدراں دا سیٹ اے۔ \emph{!!{valuation}}،
جیہڑی $\Lang{L}$ توں~$V$ وچ اے، اوہ فنکشن~$\pAssign{v}$ اے
جیہڑا !!a{element}~$V$ وچوں زبان دے !!{propositional variable}s
وچوں ہر اک نوں دیندا اے، یعنی $\pAssign{v} \colon
\PVar \to V$۔
\end{defn}

\begin{defn}\ollabel{defn:pValue}
  فرض کرو !!a{valuation}~$\pAssign{v}$ صداقتی قدراں دے
  سیٹ~$V$ وچ اے، تے ایہ سیٹ کئی قدراں والی منطق~$\Log L$ دا اے۔
  فیر قدر کڈھن والا فنکشن $\pValue{v} \colon \Frm[L] \to V$
  استقرا نال ایویں متعین کرو۔
  \begin{enumerate}
    \item $\pValue{v}(\Obj p_n) = \pAssign{v}(\Obj p_n)$۔
    \item جے $\star$ اک $0$-مقامی رابطہ اے، تاں $\pValue{v}(\star)  = \tf{\star}[\Log L]$۔
    \item جے $\star$ اک $n$-مقامی رابطہ اے، تاں
    \[
      \pValue{v}(\star(!A_1, \dots, !A_n)) = \tf{\star}[\Log L]
      (\pValue{v}(!A_1), \dots, \pValue{v}(!A_n)).
    \]
  \end{enumerate}
\end{defn}

\begin{defn}[پورا اترنا]
\ollabel{defn:satisfaction} !!{formula}~$!A$،
!!a{valuation}~$\pAssign{v}$ نال \emph{پورا اتردا اے}، یعنی
$\pSat{v}{!A}[\Log L]$، اُدوں تے صرف اُدوں جدوں
$\pValue{v}(!A)[\Log L] \in V^+$، جتھے $V^+$، منطق~$\Log L$
دیاں مقرر صداقتی قدراں دا سیٹ اے۔

  اسیں $\pSat/{v}{!A}[\Log L]$ دا مطلب «$\pSat{v}{!A}[\Log
  L]$ نہیں» لیندے آں۔ جے $\Gamma$، !!{formula}s دا سیٹ اے، تاں
  $\pSat{v}{\Gamma}[\Log
  L]$ اُدوں تے صرف اُدوں جدوں $\pSat{v}{!A}[\Log L]$ ہووے
  ہر~$!A \in \Gamma$ لئی۔
\end{defn}
\endgroup
% /OLP-0389

% OLP-0390 content/many-valued-logic/syntax-and-semantics/semantic-notions.tex
\begingroup
\def\olfilename{semantic-notions}
\expandafter\def\csname ol@sectioncs\endcsname{section}

\olfileid{mvl}{syn}{sem}

\olsection{معنیاتی مفہوم}

منّ لو کئی قدراں والی منطق $\Log L$ اک میٹرکس نال دِتّی اے۔
فیر اسیں~$\Log L$ لئی عام معنیاتی مفہوم متعین کر سکدے آں۔

\begin{defn}
\begin{enumerate}
\item !!^a{formula}~$!A$ \emph{پورا اترن جوگ} اے جے
  کسے~$\pAssign{v}$ لئی $\pSat{v}{!A}$۔ ایہ
  \emph{پورا نہ اترن جوگ} اے جے کسے وی $\pAssign{v}$ لئی
  $\pSat{v}{!A}$ نہ ہووے۔
\item !!^a{formula}~$!A$ \emph{ہمیشہ صادق} اے جے $\pSat{v}{!A}$
  ساریاں !!{valuation}s~$v$ لئی ہووے۔
\item جے $\Gamma$، !!{formula}s دا سیٹ اے، تاں $\Gamma \Entails !A$،
  یعنی «$\Gamma$ دا منطقی نتیجہ $!A$ اے»، اُدوں تے صرف اُدوں
  جدوں $\pSat{v}{!A}$ ہر اوہ !!{valuation}~$\pAssign{v}$ لئی
  ہووے جیہدے لئی $\pSat{v}{\Gamma}$ اے۔
\item جے $\Gamma$، !!{formula}s دا سیٹ اے، تاں $\Gamma$
  \emph{اکٹھے پورے اترن جوگ} اے جے کوئی !!a{valuation}~$\pAssign{v}$
  ایہو جیہی ہووے کہ $\pSat{v}{\Gamma}$۔ ہور صورت وچ $\Gamma$
  \emph{اکٹھے پورے نہ اترن جوگ} اے۔
\end{enumerate}
\end{defn}

ایہناں مفہوماں بارے کجھ اوہو گلّاں سچ نیں جیہڑیاں کلاسیکی
منطق وچ سچ نیں۔

\begin{prop}
\ollabel{prop:semanticalfacts}
\begin{enumerate}
\item $!A$ اُدوں تے صرف اُدوں ہمیشہ صادق اے جدوں
  $\emptyset \Entails !A$۔
\item جے $\Gamma$ اکٹھے پورے اترن جوگ اے، تاں $\Gamma$ دا ہر
  محدود ذیلی سیٹ وی اکٹھے پورے اترن جوگ اے۔
\item\ollabel{def:monotonicity}%
یک رُخیت۔ جے $\Gamma \subseteq \Delta$ تے $\Gamma \Entails !A$،
  تاں $\Delta \Entails !A$ وی اے۔
\item\ollabel{def:Cut}%
تعدّی۔ جے $\Gamma \Entails !A$ تے
  $\Delta \cup \{ !A\} \Entails !B$، تاں $\Gamma \cup \Delta \Entails
  !B$۔
\end{enumerate}
\end{prop}

\begin{proof}
مشق۔
\end{proof}

\begin{prob}
\olref[mvl][syn][sem]{prop:semanticalfacts} ثابت کرو۔
\end{prob}

کلاسیکی منطق وچ اسیں منطقی نتیجے تے شرطیہ دا جوڑ ویکھ سکدے
آں۔ مثلاً \emph{قاعدۂ اثباتِ مقدم} معتبر اے۔ جے $\Gamma
\Entails !A$ تے $\Gamma \Entails !A \lif !B$، تاں $\Gamma \Entails
!B$۔ کلاسیکی منطق وچ $\Entails$ تے $\lif$ دا اک ہور اہم
جوڑ معنوی استنباط دا قضیہ اے۔ $\Gamma \Entails !A
\lif !B$ اُدوں تے صرف اُدوں جدوں $\Gamma \cup \{!A\} \Entails !B$۔
ایہ نتیجے کئی قدراں والیاں منطقاں وچ ہمیشہ قائم \emph{نہیں}
ہوندے۔ ایہناں دا قائم ہونا صداقتی فنکشن~$\tf{\lif}$ اُتے منحصر اے۔
\endgroup
% /OLP-0390

% OLP-0391 content/many-valued-logic/syntax-and-semantics/sublogics.tex
\begingroup
\def\olfilename{sublogics}
\expandafter\def\csname ol@sectioncs\endcsname{section}

\olfileid{mvl}{syn}{sub}

\olsection{\texorpdfstring{کئی قدراں والیاں منطقاں بطور~$\LogCL$ دیاں ذیلی منطقاں}{کئی قدراں والیاں منطقاں بطور C دیاں ذیلی منطقاں}}

عام کئی قدراں والیاں منطقاں ایہو جیہے میٹرکساں نال متعین
ہوندیاں نیں جنہاں وچ $\{\True, \False\}$ دیاں دلیلاں لئی
صداقتی فنکشن دی قدر کلاسیکی صداقتی فنکشن دی قدر نال ملدی اے۔
خاص کر کے، ایہناں منطقاں وچ جے $x \in \{\True, \False\}$،
تاں $\tf{\lnot}[\Log L](x) =
\tf{\lnot}[\LogCL](x)$۔ نالے جے $\star$، $\land$، $\lor$،
$\lif$ وچوں کوئی اے تے $x, y \in \{\True, \False\}$،
تاں $\tf{\star}[\Log L](x,y) =
\tf{\star}[\LogCL](x,y)$۔ دوجے لفظاں وچ $\lnot$، $\land$،
$\lor$، $\lif$ دے صداقتی فنکشن جدوں $\{\True,\False\}$ تک
محدود ہون، تاں اوہ عین کلاسیکی صداقتی فنکشن نیں۔

\begin{prop}\ollabel{prop:mvl-cl}
  فرض کرو کئی قدراں والی منطق~$\Log L$ دی زبان وچ رابطے
  $\lnot$، $\land$، $\lor$، $\lif$ نیں، $\True, \False \in
  V$، تے اوہدے صداقتی فنکشن ایہ شرطاں پوریاں کردے نیں۔
  \begin{enumerate}
    \item\ollabel{prop:not} $\tf{\lnot}[\Log L](x) = \tf{\lnot}[\LogCL](x)$ جے $x =
    \True$ یا $x = \False$۔
    \item\ollabel{prop:land} $\tf{\land}[\Log L](x,y) = \tf{\land}[\LogCL](x,y)$،
    \item\ollabel{prop:lor} $\tf{\lor}[\Log L](x,y) = \tf{\lor}[\LogCL](x,y)$،
    \item\ollabel{prop:lif} $\tf{\lif}[\Log L](x,y) = \tf{\lif}[\LogCL](x,y)$،
    جے $x, y \in \{\True, \False\}$۔
  \end{enumerate}
  فیر کسے وی قدر بندی $\pAssign v$ لئی جیہڑی~$V$ وچ ہووے
  تے جیہدے لئی $\pAssign
  v(p) \in \{\True,\False\}$، ساڈے کول $\pValue v[\Log L](!A) = \pValue
  v[\LogCL](!A)$ اے۔
\end{prop}

\begin{proof}
  ~$!A$ اُتے استقرا نال۔
  \begin{enumerate}
  \item جے $!A \ident p$ ایٹمی اے، تاں ساڈے کول
  $\pValue v[\Log L](!A) =
  \pAssign v(p) = \pValue
  v[\LogCL](!A)$ اے۔
  \item جے $!A \ident \lnot B$، تاں ساڈے کول
  \begin{align*}
    \pValue v[\Log L](!A) & = \tf{\lnot}[\Log L](\pValue v[\Log L](!B)) & &\text{\olref[val]{defn:pValue} مطابق}\\
    & = \tf{\lnot}[\Log L](\pValue v[\LogCL](!B)) && \text{استقرائی مفروضے مطابق}\\
    & = \tf{\lnot}[\LogCL](\pValue v[\LogCL](!B))
&&\text{مفروضے \olref{prop:not} مطابق،}\\
&&&\text{کیوں جے $\pValue v[\LogCL](!B) \in \{\True, \False\}$،}\\
& = \pValue v[\LogCL](!A) &&\text{\olref[val]{defn:pValue} مطابق}.
\end{align*}
\item جے $!A \ident (!B \land !C)$، تاں ساڈے کول
\begin{align*}
  \pValue v[\Log L](!A) & = \tf{\land}[\Log L](\pValue v[\Log L](!B), \pValue v[\Log L](!C)) & &\text{\olref[val]{defn:pValue} مطابق}\\
  & = \tf{\land}[\Log L](\pValue v[\LogCL](!B), \pValue v[\LogCL](!C)) && \text{استقرائی مفروضے مطابق}\\
  & = \tf{\land}[\LogCL](\pValue v[\LogCL](!B), \pValue v[\LogCL](!C))
&&\text{مفروضے \olref{prop:land} مطابق،}\\
&&&\text{کیوں جے $\pValue v[\LogCL](!B),\pValue v[\LogCL](!C) \in \{\True, \False\}$،}\\
& = \pValue v[\LogCL](!A) &&\text{\olref[val]{defn:pValue} مطابق}.
\end{align*}
\end{enumerate}
$!A \ident (!B \lor !C)$ تے $!A \ident (!B \lif !C)$ والیاں
صورتاں ایسے طرح نیں۔
\end{proof}

\begin{cor}
  جے کئی قدراں والی منطق \olref{prop:mvl-cl} دیاں شرطاں
  پوریاں کرے، $\True \in V^+$ تے $\False \notin V^+$، تاں
  ${\Entails[\Log L]} \subseteq {\Entails[\LogCL]}$۔ یعنی جے $\Gamma
  \Entails[\Log L] !B$، تاں $\Gamma \Entails[\LogCL] !B$۔
  خاص کر کے، $\Log L$ دا ہر ہمیشہ صادق فارمولا کلاسیکی منطق
  وچ وی ہمیشہ صادق اے۔
\end{cor}

\begin{proof}
  اسیں اُلٹ نقیض ثابت کردے آں۔ فرض کرو $\Gamma \Entails/[\LogCL] !B$۔
  فیر کوئی !!{valuation}~$\pAssign v\colon \PVar \to
  \{\True, \False\}$ ایہو جیہی اے کہ $\pValue v[\LogCL](!A) = \True$
  ہر $!A \in \Gamma$ لئی، تے $\pValue v[\LogCL](!B) = \False$۔
  کیوں جے $\True,
  \False \in V$، ایس لئی !!{valuation}~$\pAssign v$،
  ~$\Log L$ لئی وی !!a{valuation} اے۔ \olref{prop:mvl-cl}
  مطابق $\pValue v[\Log L](!A) =
  \True$ ہر $!A \in \Gamma$ لئی، تے $\pValue v[\Log L](!B) = \False$۔
  کیوں جے $\True \in V^+$ تے $\False \notin V^+$، ایس دا
  مطلب اے کہ $\pAssign v
  \Entails[\Log L] \Gamma$ تے $\pAssign v \Entails/[\Log L] !B$،
  یعنی $\Gamma \Entails/[\Log L] !B$۔
\end{proof}
\endgroup
% /OLP-0391

\OLEndChapterHook
\endgroup
% /OLP-0384

% OLP-0392 content/many-valued-logic/three-valued-logics/three-valued-logics.tex
\begingroup
\def\olfilename{three-valued-logics}
\expandafter\def\csname ol@sectioncs\endcsname{section}

\olchapter{mvl}{thr}{تِن قدراں والیاں منطقاں}

% OLP-0393 content/many-valued-logic/three-valued-logics/introduction.tex
\begingroup
\def\olfilename{introduction}
\expandafter\def\csname ol@sectioncs\endcsname{section}

\olfileid{mvl}{thr}{int}

\olsection{جان پہچان}

جے اسیں اک ہور قدر~$\Undef$، $\True$ تے $\False$ نال
شامل کریے، تاں تِن قدراں والی منطق ملدی اے۔ بھاویں صداقتی
قدر اکّو ودھدی اے، فیر وی $\lnot$، $\land$، $\lor$ تے $\lif$
دے صداقتی فنکشن متعین کرن دے کئی طریقے نیں۔ منطق وچ ایہناں
فنکشوناں دی کوئی وی ترکیب ورتی جا سکدی اے۔ نالے ایہ چُن سکدے
آں کہ صرف $\True$ مقرر ہووے، یا $\True$ تے~$\Undef$ دوویں۔

ایتھے اسیں تِن قدراں والیاں سبھ توں معروف منطقاں وچوں کجھ
منطقاں، اوہناں دے بناؤن دے سبب تے کجھ خاصیتاں پیش کردے آں۔
\endgroup
% /OLP-0393

% OLP-0394 content/many-valued-logic/three-valued-logics/lukasiewicz.tex
\begingroup
\def\olfilename{lukasiewicz}
\expandafter\def\csname ol@sectioncs\endcsname{section}

\olfileid{mvl}{thr}{luk}

\olsection{لوکاشیویچ دی منطق}

کئی قدراں والی منطق دیاں پہلِیاں چھپیاں تے تفصیل نال تیار
کیتیاں تجویزاں وچوں اک پولینڈ دے فلسفی یان لوکاشیویچ دی 1921
دی تجویز اے۔ لوکاشیویچ نوں ارسطو دی سمندری لڑائی والے مسئلے
توں خیال آیا۔ لگدا اے کہ \emph{اجّ} جملہ «بھلکے سمندر وچ
اک لڑائی ہووے گی» نہ سچ اے نہ جھوٹ، کیوں جے اوہدی صداقتی قدر
اجے طے نہیں ہوئی۔ لوکاشیویچ نے ایہو جیہے جملے، جیہڑے مستقبل
دے اوہناں واقعات بارے نیں جو ہو وی سکدے نیں تے نہیں وی،
لئی تِیجی صداقتی قدر شامل کرن دی تجویز دِتّی۔
\begin{quote}
میں کسے تضاد توں بغیر منّ سکدا آں کہ اگلے سال کسے خاص ویلے،
مثلاً 21 دسمبر دوپہر نوں، وارسا وچ میرا ہونا ہُن نہ ہاں وچ
طے اے نہ ناں وچ۔ ایس لئی ممکن اے کہ میں دِتّے ویلے وارسا
وچ ہواں، پر ضروری نہیں۔ ایس مفروضے تحت قضیہ «میں اگلے سال
21 دسمبر دوپہر نوں وارسا وچ ہواں گا» ایس ویلے نہ سچ ہو سکدا
اے نہ جھوٹ۔ کیوں جے ایہ ہُن سچ ہووے، تاں میرا اگّوں وارسا
وچ ہونا ضروری ہووے گا، جیہڑا مفروضے دے اُلٹ اے۔ دوجے پاسے،
جے ایہ ہُن جھوٹ ہووے، تاں میرا اگّوں وارسا وچ ہونا ناممکن
ہووے گا، جیہڑا وی مفروضے دے اُلٹ اے۔ سو ایہ قضیہ ایس ویلے
نہ سچ اے نہ جھوٹ تے ایہدی تِیجی قدر ہونی چاہیدی اے، جیہڑی
«0» یعنی جھوٹ تے «1» یعنی سچ توں وکھری ہووے۔ ایہ قدر
اسیں «$\frac{1}{2}$» لکھ سکدے آں۔ ایہ «ممکن» نوں ظاہر
کردی اے تے تِیجی قدر وجوں «سچ» تے «جھوٹ» دے نال جُڑدی اے۔
\end{quote}
اسیں لوکاشیویچ دی تِیجی صداقتی قدر لئی $\Undef$ ورتاں گے۔
\footnote{لوکاشیویچ ایتھے «ممکن» نوں اوہناں معنیاں وچ ورتدا
اے جیہڑے اجّ عام نہیں، یعنی ممکن پر ضروری نہیں۔}

ایس سمجھ مطابق $\lnot$، $\land$ تے $\lor$ دے صداقتی
فنکشن طے کرنے سوکھے نیں۔ مستقبل بارے غیر طے شدہ جملے دی
نفی وی مستقبل بارے غیر طے شدہ جملہ اے، ایس لئی
$\tf{\lnot}(\Undef) = \Undef$۔ جے عطف دا اک جزو غیر طے
شدہ تے دوجا سچ اے، تاں عطف وی غیر طے شدہ اے۔ مستقبل والا
جزو جیویں وی نکلے، اوہدے مطابق عطف سچ وی نکل سکدا اے تے
جھوٹ وی۔ ایس لئی
\[
    \tf{\land}(\True, \Undef) =
\tf{\land}(\Undef, \True) =
\Undef.
\]
پر جے دوجا جزو جھوٹ اے، تاں عطف سچ نہیں نکل سکدا۔ ایس لئی
\[\tf{\land}(\False, \Undef) =
\tf{\land}(\False, \Undef) = \False.\]
باقی قدراں، جدوں دلیلاں دیاں قدراں طے ہون، یعنی $\True$
یا $\False$ ہون، کلاسیکی منطق والیاں نیں۔

شرطیے لئی گل ذرا گُنجھلدار اے۔ منّ لو $q$ مستقبل بارے غیر
طے شدہ جملہ اے۔ جے $p$ جھوٹ اے، تاں $p \lif q$ سچ ہووے گا،
بھاویں $q$ جیویں وی نکلے۔ سو $\tf{\lif}(\False, \Undef) = \True$
رکھنا چاہیدا اے۔ تے جے $p$ سچ اے، تاں $q \lif p$ سچ ہووے گا،
بھاویں $q$ جیویں وی نکلے۔ ایس لئی $\tf{\lif}(\Undef, \True) = \True$۔
جے $p$ سچ اے، تاں $p
\lif q$ سچ وی نکل سکدا اے تے جھوٹ وی، سو $\tf{\lif}(\True,
\Undef) = \Undef$۔ ایسے طرح جے $p$ جھوٹ اے، تاں $q
\lif p$ سچ وی نکل سکدا اے تے جھوٹ وی، سو
$\tf{\lif}(\Undef, \False) = \Undef$۔ ہُن اوہ صورت باقی اے
جدوں $p$ تے $q$ دوویں مستقبل بارے غیر طے شدہ نیں۔ مُڈھلے
خیال مطابق ایس صورت وچ $\Undef$ دینا چاہیدا اے۔ پر ایس
نال $!A \lif !A$ ہمیشہ صادق \emph{نہیں} رہے گا۔ لوکاشیویچ
نوں $!A \lor \lnot !A$ تے $\lnot(!A
\land \lnot !A)$ چھڈن وچ کوئی دِقّت نہیں سی، پر اوہ $!A \lif !A$
چھڈن لئی تیار نہیں سی۔ ایس لئی اوہنے مقرر کیتا کہ
$\tf{\lif}(\Undef, \Undef) =
\True$۔

\begin{defn}\ollabel{def:lukasiewicz}
تِن قدراں والی لوکاشیویچ دی منطق ایس میٹرکس نال متعین اے۔
\begin{enumerate}
  \item معیاری قضیاتی زبان $\Lang L_0$ جیہدے وچ
  $\lnot$، $\land$، $\lor$، $\lif$ نیں۔
  \item صداقتی قدراں دا سیٹ $V = \{\True, \Undef, \False\}$۔
  \item $\True$ اکّو مقرر قدر اے، یعنی $V^+ = \{\True\}$۔
  \item صداقتی فنکشن ہِٹھلے جدولاں توں ملدے نیں۔
  \begin{center}
    \begin{tabular}{c|c} 
      $\tf{\lnot}$ & \\ 
      \hline  
      $\True$ & $\False$ \\ 
      $\Undef$ & $\Undef$ \\
      $\False$ & $\True$ 
    \end{tabular}
    \quad
    \begin{tabular}{c|ccc} 
      $\tf{\land}[\LogLuk[3]]$ & $\True$ & $\Undef$ & $\False$ \\ 
      \hline 
      $\True$ & $\True$ & $\Undef$ & $\False$ \\ 
      $\Undef$ & $\Undef$ & $\Undef$ & $\False$\\ 
      $\False$ & $\False$ & $\False$ & $\False$ 
    \end{tabular}
    \\[2ex]
    \begin{tabular}{c|ccc} 
      $\tf{\lor}[\LogLuk[3]]$ & $\True$ & $\Undef$ & $\False$ \\ 
      \hline 
      $\True$ & $\True$ & $\True$ & $\True$ \\ 
      $\Undef$ & $\True$ & $\Undef$ & $\Undef$ \\
      $\False$ & $\True$ & $\Undef$ & $\False$ 
    \end{tabular}
    \quad
    \begin{tabular}{c|ccc} 
      $\tf{\lif}[\LogLuk[3]]$ & $\True$ & $\Undef$ & $\False$ \\ 
      \hline 
      $\True$ & $\True$ & $\Undef$ & $\False$ \\ 
      $\Undef$ & $\True$ & $\True$ & $\Undef$  \\ 
      $\False$ & $\True$ & $\True$ & $\True$ 
    \end{tabular}
  \end{center}
\end{enumerate}
\end{defn}

سوکھے ای دِسدا اے کہ کوئی وی !!{formula}~$!A$ جیہدے وچ
صرف $\lnot$، $\land$ تے $\lor$ ہوون، صداقتی قدر~$\Undef$
لیندا اے جے اوہدے سارے !!{propositional variable}s نوں~$\Undef$
دِتّی جاوے۔ سو مثلاً کلاسیکی ہمیشہ صادق فارمولے $p \lor \lnot p$
تے $\lnot(p
\land \lnot p)$، $\LogLuk[3]$ وچ ہمیشہ صادق نہیں، کیوں جے
$\pValue{v}(!A)
= \Undef$ جدوں $\pAssign v(p) = \Undef$۔

اوہناں !!{valuation}s وچ جتھے $\pAssign v(p) = \True$ یا
$\False$، $\pValue
v(!A)$ اپنی کلاسیکی صداقتی قدر دے برابر ہووے گا۔

\begin{prop}
  جے $\pAssign v(p) \in \{\True, \False\}$ ہر اوہ $p$ لئی
  جیہڑا~$!A$ وچ اے، تاں
  $\pValue v(!A)[\LogLuk[3]] = \pValue v(!A)[\LogCL]$۔
\end{prop}

\begin{prob}\label{mvl:thr:luk:prob:luk-iff} منّ لو اسیں
  $\pValue v(!A \liff !B) = \pValue v((!A \lif !B) \land (!B \lif
  !A))$ نوں~$\LogLuk[3]$ وچ تعریف وجوں ورتدے آں۔ $\liff$ دا
  صداقتی جدول کیہ ہووے گا؟
\end{prob}

کئی کلاسیکی ہمیشہ صادق فارمولے \LogLuk[3] وچ وی \emph{ہمیشہ
صادق} نیں، مثلاً $\lnot p \lif (p \lif q)$۔ کلاسیکی منطق وانگ
اسیں ایہ گل صداقتی جدولاں نال پرکھ سکدے آں۔
\begin{center}
\begin{tabular}{cc|cccccc}
  $p$ & $q$ & $\lnot$ & $p$ & $\lif$ & $\smash{(}p$ & $\lif$ & $q\smash{)}$ \\ \hline
  \True & \True & \False & \True & \True & \True & \True & \True \\
  \True & \Undef & \False & \True & \True & \True & \Undef & \Undef \\
  \True & \False & \False & \True & \True & \True & \False & \False \\
  \Undef & \True & \Undef & \Undef & \True & \Undef & \True & \True \\
  \Undef & \Undef & \Undef & \Undef & \True & \Undef & \True & \Undef \\
  \Undef & \False & \Undef & \Undef & \True & \Undef & \Undef & \False \\
  \False & \True & \True & \False & \True & \False & \True & \True \\
  \False & \Undef & \True & \False & \True & \False & \True & \Undef \\
  \False & \False & \True & \False & \True & \False & \True & \False \\  
\end{tabular}
\end{center}

\begin{prob}
  دکھاؤ کہ ہِٹھلے فارمولے \LogLuk[3] وچ ہمیشہ صادق نیں۔
  \begin{enumerate}
    \item $p \lif (q \lif p)$
    \item\label{mvl:thr:luk:prob:luk-taut-2} $\lnot(p \land q) \liff (\lnot p \lor \lnot q)$
    \item\label{mvl:thr:luk:prob:luk-taut-3} $\lnot(p \lor q) \liff (\lnot p \land \lnot q)$
  \end{enumerate}
  (\olref[mvl][thr][luk]{prob:luk-taut-2} تے
  \olref[mvl][thr][luk]{prob:luk-taut-3} وچ $!A \liff !B$
  نوں $(!A \lif !B) \land (!B \lif !A)$ دا مختصر روپ سمجھو،
  یا \olref[mvl][thr][luk]{prob:luk-iff} دا اپنا حل ویکھو۔)
\end{prob}

\begin{prob}
  دکھاؤ کہ ہِٹھلے کلاسیکی ہمیشہ صادق فارمولے~\LogLuk[3] وچ
  ہمیشہ صادق نہیں۔
  \begin{enumerate}
    \item $(\lnot p \land p) \lif q)$
    \item $((p \lif q) \lif p) \lif p$
    \item $(p \lif (p \lif q)) \lif (p \lif q)$
  \end{enumerate}
\end{prob}

سو شاید ایہ خیال آوے کہ بھاویں سارے کلاسیکی ہمیشہ صادق
فارمولے~$\LogLuk[3]$ وچ ہمیشہ صادق نہیں، فیر وی اوہناں دی قدر~$\True$ یا~$\Undef$ ہونی چاہیدی اے
ہر !!{valuation} اُتے۔ ایہ گل سچ نہیں۔ اک ردّی مثال ایہ اے۔
\[
  \lnot(p \lif \lnot p) \lor \lnot(\lnot p \lif p)
\]
ایہدی قدر $\False$ اے جے $p$ دی قدر~$\Undef$ اے۔

\begin{prob}
  ہِٹھلے تعلقاں وچوں کیہڑے لوکاشیویچ دی منطق وچ قائم نیں؟
  ہر اک لئی صداقتی جدول دیو۔
  \begin{enumerate}
    \item $p, p \lif q \Entails q$
    \item $\lnot\lnot p \Entails p$
    \item $p \land q \Entails p$
    \item $p \Entails p \land p$
    \item $p \Entails p \lor q$
  \end{enumerate}
\end{prob}

لوکاشیویچ نوں امید سی کہ اپنے تِن قدراں والے نظام اُتے امکان
دی منطق بنا سکے گا۔ ایس لئی اوہنے اک مقامی رابطہ $\Diamond
!A$، یعنی «$!A$ ممکن اے»، تے اوہدے نال ملدا $\Box !A$،
یعنی «$!A$ ضروری اے»، شامل کیتے۔
\begin{center}
  \begin{tabular}{c|c} 
    $\tf{\Diamond}$ & \\ 
    \hline  
    $\True$ & $\True$ \\ 
    $\Undef$ & $\True$ \\
    $\False$ & $\False$ 
  \end{tabular}
  \quad
  \begin{tabular}{c|c} 
    $\tf{\Box}$ & \\ 
    \hline  
    $\True$ & $\True$ \\ 
    $\Undef$ & $\False$ \\
    $\False$ & $\False$ 
  \end{tabular}
\end{center}
دوجے لفظاں وچ، $p$ اُدوں تے صرف اُدوں ممکن اے جدوں ایہ
پہلاں ای جھوٹ طے نہ ہووے۔ تے $p$ اُدوں تے صرف اُدوں ضروری
اے جدوں ایہ پہلاں ای سچ طے ہووے۔

\begin{prob}
  دکھاؤ کہ $\Box p \liff \lnot\Diamond \lnot p$ تے $\Diamond p \liff
  \lnot \Box \lnot p$، ~\LogLuk[3] وچ ہمیشہ صادق نیں جدوں
  اینوں $\Box$ تے~$\Diamond$ دے صداقتی جدولاں نال ودھایا جاوے۔
\end{prob}

پر ایس تجویز کیتی ہوئی امکانی منطق دیاں کمزوریاں چھیتی
سامھنے آ گئیاں۔ جیویں وی گل نکلے، $p \land \lnot p$ کدے سچ
نہیں نکل سکدا۔ سو بھاویں ایہ ہُن طے نہیں ہویا، تے ایس لئی
غیر طے شدہ اے، اینوں ناممکن منّنا چاہیدا اے۔ یعنی
$\lnot \Diamond(p
\land \lnot p)$ ہمیشہ صادق ہونا چاہیدا اے۔ پر جے $\pAssign v(p) =
\Undef$، تاں $\pValue v(\lnot \Diamond(p \land \lnot p)) =
\Undef$۔ لوکاشیویچ دی ایہ گل ٹھیک سی کہ امکان تے ضرورت
ورگے فرق رکھن لئی دو صداقتی قدراں کافی نہیں۔ پر تِیجی
صداقتی قدر شامل کرنا وی کافی نہیں۔
\endgroup
% /OLP-0394

% OLP-0395 content/many-valued-logic/three-valued-logics/kleene.tex
\begingroup
\def\olfilename{kleene}
\expandafter\def\csname ol@sectioncs\endcsname{section}

\olfileid{mvl}{thr}{skl}

\olsection{کلینی دیاں منطقاں}

سٹیفن کلینی نے تِن قدراں والیاں دو منطقاں پیش کیتیاں۔ اوہناں
دا خیال اوس منطق توں آیا جیہدے وچ صداقتی قدراں نوں حسابی
طریقیاں دے نتیجے سمجھدے آں۔ کوئی طریقہ $\True$ یا $\False$
دے سکدا اے، پر ایہ وی ہو سکدا اے کہ اوہ مُکے ای نہ۔ اوہ
صورت وچ متعلقہ صداقتی قدر غیر معین ہوندی اے، جیہنوں~$\Undef$
نال لکھدے آں۔

قضیے~$!A$ دی نفی کڈھن لئی پہلاں~$!A$ دی قدر کڈھو، تے
فیر نتیجے دی اُلٹ قدر واپس دیو۔ جے $!A$ دا حساب نہیں مُکدا،
تاں سارا طریقہ وی نہیں مُکدا۔ سو $\Undef$ دی نفی $\Undef$ اے۔

عطف $!A \land !B$ دا حساب کرن دے دو طریقے نیں۔ اک ایہ اے
کہ پہلاں~$!A$ دا حساب کرو، فیر $!B$ دا۔ عطف دا نتیجہ $\True$
ہووے گا جے دوہاں دا نتیجہ~$\True$ ہووے۔ ہور صورت وچ $\False$
ہووے گا۔ جے دوہاں وچوں کوئی حساب نہ رُکے، تاں سارا طریقہ وی
نہیں رُکدا۔ سو ایس صورت وچ جے اک جزو غیر معین ہووے، تاں عطف
وی غیر معین اے۔ فصل بارے وی اوہی گل اے۔

پر جے اسیں $!A$ تے $!B$ دیاں قدراں متوازی کڈھ سکیے، تاں
ہور سوکھ اے۔ جے دو طریقیاں وچوں اک رُک کے $\False$ واپس
دے دے، تاں اسیں رک سکدے آں، کیوں جے جواب لازماً جھوٹ اے۔
سو ایس صورت وچ اک جزو جھوٹ ہووے تاں عطف جھوٹ اے، بھاویں
دوجا جزو غیر معین ہووے۔ ایسے طرح فصل دا حساب متوازی کریے،
تاں جدوں دوہاں جزواں وچوں اک دا طریقہ سچ واپس دے دے، اسیں
رک سکدے آں۔ فیر فصل لازماً سچ اے۔ سو ایس صورت وچ مرکب
قضیے دا نتیجہ پتا لگ سکدا اے، بھاویں اوہدے جزواں وچوں اک
غیر معین ہووے۔ ایس سمجھ مطابق $\Undef$ نوں «غیر معین» دی
تھاں «نامعلوم» پڑھ سکدے آں۔

ایہ دو سمجھاں توں کلینی دی مضبوط تے کمزور منطق بندی اے۔
شرطیہ نوں $\lnot !A \lor !B$ دے ہم ارز متعین کردے آں۔

\begin{defn}
\emph{کلینی دی مضبوط منطق}~$\LogKs$ ایس میٹرکس نال متعین اے۔
\begin{enumerate}
  \item معیاری قضیاتی زبان $\Lang L_0$ جیہدے وچ
  $\lnot$، $\land$، $\lor$، $\lif$ نیں۔
  \item صداقتی قدراں دا سیٹ $V = \{\True, \Undef, \False\}$۔
  \item $\True$ اکّو مقرر قدر اے، یعنی $V^+ = \{\True\}$۔
  \item صداقتی فنکشن ہِٹھلے جدولاں توں ملدے نیں۔
  \begin{center}
    \begin{tabular}{c|c} 
      $\tf{\lnot}$ & \\ 
      \hline  
      $\True$ & $\False$ \\ 
      $\Undef$ & $\Undef$ \\
      $\False$ & $\True$ 
    \end{tabular}
    \quad
    \begin{tabular}{c|ccc} 
      $\tf{\land}[\LogKs]$ & $\True$ & $\Undef$ & $\False$ \\ 
      \hline 
      $\True$ & $\True$ & $\Undef$ & $\False$ \\ 
      $\Undef$ & $\Undef$ & $\Undef$ & $\False$\\ 
      $\False$ & $\False$ & $\False$ & $\False$ 
    \end{tabular}
    \\[2ex]
    \begin{tabular}{c|ccc} 
      $\tf{\lor}[\LogKs]$ & $\True$ & $\Undef$ & $\False$ \\ 
      \hline 
      $\True$ & $\True$ & $\True$ & $\True$ \\ 
      $\Undef$ & $\True$ & $\Undef$ & $\Undef$ \\
      $\False$ & $\True$ & $\Undef$ & $\False$ 
    \end{tabular}
    \quad
    \begin{tabular}{c|ccc} 
      $\tf{\lif}[\LogKs]$ & $\True$ & $\Undef$ & $\False$ \\ 
      \hline 
      $\True$ & $\True$ & $\Undef$ & $\False$ \\ 
      $\Undef$ & $\True$ & $\Undef$ & $\Undef$  \\ 
      $\False$ & $\True$ & $\True$ & $\True$ 
    \end{tabular}
  \end{center}
\end{enumerate}
\end{defn}

\begin{defn}
\emph{کلینی دی کمزور منطق}~$\LogKw$ ایس میٹرکس نال متعین اے۔
\begin{enumerate}
  \item معیاری قضیاتی زبان $\Lang L_0$ جیہدے وچ
  $\lnot$، $\land$، $\lor$، $\lif$ نیں۔
  \item صداقتی قدراں دا سیٹ $V = \{\True, \Undef, \False\}$۔
  \item $\True$ اکّو مقرر قدر اے، یعنی $V^+ = \{\True\}$۔
  \item صداقتی فنکشن ہِٹھلے جدولاں توں ملدے نیں۔
  \begin{center}
    \begin{tabular}{c|c} 
      $\tf{\lnot}$ & \\ 
      \hline  
      $\True$ & $\False$ \\ 
      $\Undef$ & $\Undef$ \\
      $\False$ & $\True$ 
    \end{tabular}
    \quad
    \begin{tabular}{c|ccc} 
      $\tf{\land}[\LogKw]$ & $\True$ & $\Undef$ & $\False$ \\ 
      \hline 
      $\True$ & $\True$ & $\Undef$ & $\False$ \\ 
      $\Undef$ & $\Undef$ & $\Undef$ & $\Undef$\\ 
      $\False$ & $\False$ & $\Undef$ & $\False$ 
    \end{tabular}
    \\[2ex]
    \begin{tabular}{c|ccc} 
      $\tf{\lor}[\LogKw]$ & $\True$ & $\Undef$ & $\False$ \\ 
      \hline 
      $\True$ & $\True$ & $\Undef$ & $\True$ \\ 
      $\Undef$ & $\Undef$ & $\Undef$ & $\Undef$ \\
      $\False$ & $\True$ & $\Undef$ & $\False$ 
    \end{tabular}
    \quad
    \begin{tabular}{c|ccc} 
      $\tf{\lif}[\LogKw]$ & $\True$ & $\Undef$ & $\False$ \\ 
      \hline 
      $\True$ & $\True$ & $\Undef$ & $\False$ \\ 
      $\Undef$ & $\Undef$ & $\Undef$ & $\Undef$  \\ 
      $\False$ & $\True$ & $\Undef$ & $\True$ 
    \end{tabular}
  \end{center}
\end{enumerate}
\end{defn}

\begin{prop}
  $\LogKs$ تے $\LogKw$ وچ کوئی ہمیشہ صادق فارمولا نہیں۔
\end{prop}

\begin{proof}
  جے $\pAssign v(p) = \Undef$ سارے !!{propositional variable}s~$p$
  لئی، تاں ہر فارمولا~$!A$ دی صداقتی قدر~$\pValue v(!A) =
  \Undef$ ہووے گی، کیوں جے دوہاں منطقاں وچ
  \[
    \tf{\lnot}(\Undef) = \tf{\lor}(\Undef, \Undef) = \tf{\land}(\Undef,
  \Undef) = \tf{\lif}(\Undef, \Undef) = \Undef
  \]
  اے۔ کیوں جے $\Undef \notin V^+$، چاہے منطق $\LogKs$ ہووے یا
  $\LogKw$، ایس لئی ایس !!{valuation} اُتے $!A$ دی قدر مقرر نہیں۔
\end{proof}

بھاویں کلینی دیاں کمزور تے مضبوط دوہاں منطقاں وچ کوئی
ہمیشہ صادق فارمولا نہیں، اوہناں دے منطقی نتیجے دے تعلق
معمولی نہیں۔

\begin{prob}
  ہِٹھلے تعلقاں وچوں کیہڑے (الف) کلینی دی مضبوط منطق وچ تے
  (ب) اوہدی کمزور منطق وچ قائم نیں؟ ہر اک لئی صداقتی جدول دیو۔
  \begin{enumerate}
    \item $p, p \lif q \Entails q$
    \item $p \lor q, \lnot p \Entails q$
    \item $p \land q \Entails p$
    \item $p \Entails p \land p$
    \item $p \Entails p \lor q$
  \end{enumerate}
\end{prob}

دمتری بوچوار نے $\Undef$ نوں «بے معنی» سمجھیا۔ اوہنے جھوٹے
دے تناقض ورگے تناقض حل کرن دی کوشش وچ مقرر کیتا کہ تناقض
والے جملیاں دی قدر~$\Undef$ ہووے۔ اوہنے اک منطق پیش کیتی جیہڑی
بنیادی طور اُتے کلینی دی کمزور منطق اے، جیہدے وچ کجھ ہور
رابطے شامل نیں۔ اوہناں وچوں دو «باہرلی نفی» تے «غیر معین اے»
والا عامل نیں۔
\begin{center}
  \begin{tabular}{c|c} 
    $\tf{\sim}$ & \\ 
    \hline  
    $\True$ & $\False$ \\ 
    $\Undef$ & $\True$ \\
    $\False$ & $\True$ 
  \end{tabular}
\quad
  \begin{tabular}{c|c} 
    $\tf{+}$ & \\ 
    \hline  
    $\True$ & $\False$ \\ 
    $\Undef$ & $\True$ \\
    $\False$ & $\False$ 
  \end{tabular}
\end{center}

\begin{prob}
  کی تُسیں بوچوار دی منطق وچ $\sim$ نوں $\lnot$ تے~$+$ نال
  متعین کر سکدے او؟ یعنی !!a{formula} لبھو جیہدے وچ بس
  !!{propositional variable}~$p$ ہووے تے $\sim$ نہ آوے، تے جیہدی صداقتی قدر
  ہمیشہ~$\mathord{\sim}p$ دی قدر دے برابر ہووے۔ اپنی گل صحیح
  ثابت کرن لئی صداقتی جدول دیو۔
\end{prob}
\endgroup
% /OLP-0395

% OLP-0396 content/many-valued-logic/three-valued-logics/goedel.tex
\begingroup
\def\olfilename{goedel}
\expandafter\def\csname ol@sectioncs\endcsname{section}

\olfileid{mvl}{thr}{god}

\olsection{گوڈل دیاں منطقاں}

کرٹ گوڈل نے $n$ قدراں والیاں منطقاں دا اک سلسلہ پیش کیتا۔
ایہناں وچوں ہر اک وچ اوہ سارے !!{formula}s شامل نیں جیہڑے
وجدانیت پسند منطق وچ معتبر نیں، تے ہر اک آپ کلاسیکی منطق وچ
شامل اے۔ ایہ پہلی دلچسپ منطق اے۔

\begin{defn}\ollabel{defn:goedel}
\emph{$3$ قدراں والی گوڈل دی منطق}~$\LogGod$ ایس میٹرکس نال متعین اے۔
\begin{enumerate}
  \item معیاری قضیاتی زبان $\Lang L_0$ جیہدے وچ
  $\lfalse$، $\lnot$، $\land$، $\lor$، $\lif$ نیں۔
  \item صداقتی قدراں دا سیٹ $V = \{\True, \Undef, \False\}$۔
  \item $\True$ اکّو مقرر قدر اے، یعنی $V^+ = \{\True\}$۔
  \item $\lfalse$ لئی ساڈے کول $\tf{\lfalse} = \False$ اے۔
  باقی رابطیاں دے صداقتی فنکشن ہِٹھلے جدولاں توں ملدے نیں۔
  \begin{center}
    \begin{tabular}{c|c} 
      $\tf{\lnot}[\LogGod]$ & \\ 
      \hline  
      $\True$ & $\False$ \\ 
      $\Undef$ & $\False$ \\
      $\False$ & $\True$ 
    \end{tabular}
    \quad
    \begin{tabular}{c|ccc} 
      $\tf{\land}[\LogGod]$ & $\True$ & $\Undef$ & $\False$ \\ 
      \hline 
      $\True$ & $\True$ & $\Undef$ & $\False$ \\ 
      $\Undef$ & $\Undef$ & $\Undef$ & $\False$\\ 
      $\False$ & $\False$ & $\False$ & $\False$ 
    \end{tabular}
    \\[2ex]
    \begin{tabular}{c|ccc} 
      $\tf{\lor}[\LogGod]$ & $\True$ & $\Undef$ & $\False$ \\ 
      \hline 
      $\True$ & $\True$ & $\True$ & $\True$ \\ 
      $\Undef$ & $\True$ & $\Undef$ & $\Undef$ \\
      $\False$ & $\True$ & $\Undef$ & $\False$ 
    \end{tabular}
    \quad
    \begin{tabular}{c|ccc} 
      $\tf{\lif}[\LogGod]$ & $\True$ & $\Undef$ & $\False$ \\ 
      \hline 
      $\True$ & $\True$ & $\Undef$ & $\False$ \\ 
      $\Undef$ & $\True$ & $\True$ & $\False$  \\ 
      $\False$ & $\True$ & $\True$ & $\True$ 
    \end{tabular}
  \end{center}
\end{enumerate}
\end{defn}

دھیان دیو کہ $\land$ تے~$\lor$ دے صداقتی جدول لوکاشیویچ
تے کلینی دی مضبوط منطق والیاں وانگ نیں۔ پر $\lnot$ تے~$\lif$
دے جدول ہر اک وچ وکھرے نیں۔ گوڈل دی منطق وچ
$\tf{\lnot}(\Undef) = \False$۔ لوکاشیویچ تے کلینی دیاں منطقاں
توں وکھری گل ایہ اے کہ $\tf{\lif}(\Undef, \False) = \False$۔
کلینی دی منطق توں وکھری پر لوکاشیویچ دی منطق والی گل ایہ اے
کہ $\tf{\lif}(\Undef,
\Undef) = \True$۔

جیویں اُتے وجدانیت پسند منطق نال جوڑ توں دِسدا اے، $\LogGod[3]$
وجدانیت پسند منطق دے نیڑے اے۔ وجدانیت پسند منطق دیاں ساریاں سچائیاں
~$\LogGod[3]$ وچ ہمیشہ صادق نیں۔ کئی کلاسیکی ہمیشہ صادق فارمولے
جیہڑے وجدانیت پسند منطق وچ معتبر نہیں، ~$\LogGod[3]$ وچ وی
ہمیشہ صادق نہیں۔ مثلاً ہِٹھلے فارمولے ہمیشہ صادق نہیں۔
\begin{align*}
  & p \lor \lnot p && (p \lif q) \lif (\lnot p \lor q) \\
  & \lnot\lnot p \lif p && \lnot(\lnot p \land \lnot q) \lif (p \lor q) \\
  & ((p \lif q) \lif p) \lif p && \lnot(p \lif q) \lif (p \land \lnot q)
\end{align*}
پر $\LogGod[3]$ دا ہر ہمیشہ صادق فارمولا وجدانیت پسند منطق وچ
معتبر نہیں، مثلاً $\lnot\lnot p \lor \lnot p$ یا $(p
\lif q) \lor (q \lif p)$۔

\begin{prob}
  صداقتی جدول دے کے دکھاؤ کہ ہِٹھلے فارمولے~$\LogGod[3]$
  وچ ہمیشہ صادق نیں۔
  \begin{align*}
    & \lnot\lnot p \lor \lnot p\\
    & (p \lif q) \lor (q \lif p) \\
    & \lnot(p \land q) \lif (\lnot p \lor \lnot q) \\
    & (p \lif q) \lor (q \lif r) \lor (r \lif s)
  \end{align*}
\end{prob}

\begin{prob}
  صداقتی جدول دے کے دکھاؤ کہ ہِٹھلے فارمولے~$\LogGod[3]$
  وچ ہمیشہ صادق نہیں۔
  \begin{align*}
    & (p \lif q) \lif (\lnot p \lor q) \\
    & \lnot(\lnot p \land \lnot q) \lif (p \lor q) \\
    & ((p \lif q) \lif p) \lif p \\
    & \lnot(p \lif q) \lif (p \land \lnot q)
  \end{align*}
\end{prob}

\begin{prob}
  ہِٹھلے تعلقاں وچوں کیہڑے گوڈل دی منطق وچ قائم نیں؟
  ہر اک لئی صداقتی جدول دیو۔
  \begin{enumerate}
    \item $p, p \lif q \Entails q$
    \item $p \lor q, \lnot p \Entails q$
    \item $p \land q \Entails p$
    \item $p \Entails p \land p$
    \item $p \Entails p \lor q$
  \end{enumerate}
\end{prob}
\endgroup
% /OLP-0396

% OLP-0397 content/many-valued-logic/three-valued-logics/multiple-designation.tex
\begingroup
\def\olfilename{multiple-designation}
\expandafter\def\csname ol@sectioncs\endcsname{section}

\olfileid{mvl}{thr}{mul}

\olsection{\texorpdfstring{$\True$ دے نال ہور مقرر قدراں}{𝕋 دے نال ہور مقرر قدراں}}

ہُن تک جیہڑیاں منطقاں ویکھیاں نیں، اوہناں ساریاں وچ مقرر
صداقتی قدراں دا سیٹ $V^+ = \{\True\}$ سی۔ یعنی کوئی گل
اُدوں تے صرف اُدوں سچ منّی جاندی سی جدوں اوہدی صداقتی
قدر~$\True$ ہووے۔ پر کوئی گل بس~$\False$ نہ ہووے تاں وی
اوہنوں سچ منّ سکدے آں۔ فیر میٹرکس وچ، مثلاً، $V^+ = \{\True, \Undef\}$
رکھ کے اک منطق ملے گی۔

\begin{defn}
\emph{تناقض دی منطق}~$\LogLP$ ایس میٹرکس نال متعین اے۔
\begin{enumerate}
  \item معیاری قضیاتی زبان $\Lang L_0$ جیہدے وچ
  $\lnot$، $\land$، $\lor$، $\lif$ نیں۔
  \item صداقتی قدراں دا سیٹ $V = \{\True, \Undef, \False\}$۔
  \item $\True$ تے $\Undef$ مقرر نیں، یعنی $V^+ = \{\True, \Undef\}$۔
  \item صداقتی فنکشن کلینی دی مضبوط منطق والے ای نیں۔
\end{enumerate}
\end{defn}

\begin{defn}
ہالدین دی \emph{بے معنائی دی منطق}~$\LogHal$ ایس میٹرکس نال متعین اے۔
\begin{enumerate}
  \item معیاری قضیاتی زبان $\Lang L_0$ جیہدے وچ
  $\lnot$، $\land$، $\lor$، $\lif$ تے $1$-مقامی رابطہ~$+$ نیں۔
  \item صداقتی قدراں دا سیٹ $V = \{\True, \Undef, \False\}$۔
  \item $\True$ تے $\Undef$ مقرر نیں، یعنی $V^+ = \{\True, \Undef\}$۔
  \item کلینی دی کمزور منطق والے صداقتی فنکشن، نال «بے معنی اے» والا عامل۔
  \begin{center}
    \begin{tabular}{c|c} 
    $\tf{+}$ & \\ 
    \hline  
    $\True$ & $\False$ \\ 
    $\Undef$ & $\True$ \\
    $\False$ & $\False$ 
    \end{tabular}
  \end{center}
\end{enumerate}
\end{defn}

ایہناں دے صداقتی جدول کلینی دیاں منطقاں والے نیں، پر ایہناں
وچ \emph{ہمیشہ صادق فارمولے نیں}۔

\begin{prop}\ollabel{prop:LP-taut-CL}
  $\LogLP$ دے ہمیشہ صادق فارمولے اوہی نیں جیہڑے کلاسیکی
  قضیاتی منطق وچ ہمیشہ صادق نیں۔
\end{prop}

\begin{proof}
  \olref[syn][sub]{prop:mvl-cl} مطابق جے $\Entails[\LogLP] !A$،
  تاں $\Entails[\LogCL] !A$۔ اُلٹی سمت دکھاؤن لئی اسیں
  دکھاؤندے آں کہ جے کوئی !!a{valuation}~$\pAssign v\colon \PVar \to \{\False, \True,
  \Undef\}$ ایہو جیہی اے کہ $\pValue v(!A)[\LogKs] = \False$، تاں
  کوئی !!a{valuation}~$\pAssign {v'}\colon \PVar \to \{\False, \True\}$
  ایہو جیہی اے کہ $\pValue {v'}(!A)[\LogCL] = \False$۔ ایس توں
  $\LogLP$ والا نتیجہ مل جاندا اے، کیوں جے $\LogKs$ تے $\LogLP$
  دے خاص صداقتی فنکشن اکّو نیں، تے $\False$ اکّو صداقتی قدر
  اے جیہڑی $\LogLP$ وچ مقرر نہیں۔ ایہو اک فرق $\LogLP$ تے~$\LogKs$
  وچ اے۔ سو جے $\Entails/[\LogLP] !A$، تاں کسے
  !!{valuation}~$\pAssign v$ لئی $\pValue v(!A)[\LogLP] = \pValue
  v[\LogKs](!A) = \False$۔ جیہڑا دعویٰ ثابت کر رہے آں، اوہدے
  مطابق $\pValue{v'}[\LogCL](!A) = \False$، یعنی $\Entails/[\LogCL] !A$۔

  دعویٰ ثابت کرن لئی پہلاں $\pAssign {v'}$ دی تعریف ایویں کردے آں۔
  \[
    \pAssign {v'}(p) =
    \begin{cases}
    \True & \text{جے } \pAssign {v}(p) \in \{\True, \Undef\}\\
    \False & \text{ہور صورت وچ}
  \end{cases}
  \]
  ہُن $!A$ اُتے استقرا نال دکھاؤندے آں کہ (الف) جے
  $\pValue v(!A)[\LogKs]
  = \False$، تاں $\pValue {v'}(!A)[\LogCL] = \False$، تے (ب) جے
  $\pValue v(!A)[\LogKs] = \True$، تاں $\pValue {v'}(!A)[\LogCL] =
  \True$۔
  \begin{enumerate}
    \item استقرا دی مُڈھلی صورت $!A \ident p$ اے۔
    \olref[syn][val]{defn:pValue} مطابق $\pValue v(!A)[\LogKs]
    = \pAssign v(p) = \pValue {v'}(!A)[\LogCL]$، جیہدے توں
    (الف) تے (ب) دوویں ملدے نیں۔

    استقرائی قدم لئی ایہ صورتاں ویکھو۔
    \item $!A \ident \lnot !B$۔
    \begin{enumerate}
      \item منّ لو $\pValue v(\lnot!B)[\LogKs] = \False$۔
      $\tf{\lnot}[\LogKs]$ دی تعریف مطابق $\pValue v(!B)[\LogKs]  =
      \True$۔ استقرائی مفروضے دی صورت (ب) توں $\pValue
      {v'}(!B)[\LogCL]  = \True$ ملدا اے، سو $\pValue {v'}(\lnot !B)[\LogCL] = \False$۔
      \item منّ لو $\pValue v(\lnot!B)[\LogKs] = \True$۔
      $\tf{\lnot}[\LogKs]$ دی تعریف مطابق $\pValue v(!B)[\LogKs]  =
      \False$۔ استقرائی مفروضے دی صورت (الف) توں $\pValue
      {v'}(!B)[\LogCL]  = \False$ ملدا اے، سو $\pValue {v'}(\lnot !B)[\LogCL] = \True$۔
    \end{enumerate}

    \item $!A \ident (!B \land !C)$۔
    \begin{enumerate}
      \item منّ لو $\pValue v(!B \land !C)[\LogKs] = \False$۔
      $\tf{\land}[\LogKs]$ دی تعریف مطابق $\pValue v(!B)[\LogKs]  =
      \False$ یا $\pValue v(!B)[\LogKs]  = \False$۔ استقرائی
      مفروضے دی صورت (الف) توں $\pValue {v'}(!B)[\LogCL]  = \False$
      یا $\pValue {v'}(!C)[\LogCL]  = \False$ ملدا اے، سو $\pValue {v'}(!B \land
      !C)[\LogCL] = \False$۔
      \item منّ لو $\pValue v(!B \land !C)[\LogKs] = \True$۔
      $\tf{\land}[\LogKs]$ دی تعریف مطابق $\pValue v(!B)[\LogKs]  =
      \True$ تے $\pValue v(!B)[\LogKs]  = \True$۔ استقرائی
      مفروضے دی صورت (ب) توں $\pValue {v'}(!B)[\LogCL]  = \True$
      تے $\pValue {v'}(!C)[\LogCL]  = \True$ ملدا اے، سو $\pValue {v'}(!B \land
      !C)[\LogCL] = \True$۔
    \end{enumerate}
  \end{enumerate}

  باقی دو صورتاں ایسے طرح نیں تے مشق لئی چھڈیاں نیں۔ ہور
  طریقہ ایہ اے کہ اُپرلا ثبوت اوہناں سارے !!{formula}s لئی
  نتیجہ قائم کردا اے جنہاں وچ صرف $\lnot$ تے~$\land$ نیں۔
  ہُن ایہ گل ورت سکدے آں کہ $\LogKs$ تے $\LogCL$ دوہاں
  وچ ہر $\pAssign v$ لئی $\pValue v(!B \lor
  !C) = \pValue v(\lnot(\lnot!B \land \lnot !C))$ تے $\pValue v(!B \lif
  !C) = \pValue v(\lnot(!B \land \lnot !C))$۔
\end{proof}

\begin{prob}
  \olref[mvl][thr][mul]{prop:LP-taut-CL} دا ثبوت پورا کرو۔
  یعنی (الف) تے (ب) اوہناں صورتاں لئی ثابت کرو جتھے
  $!A \ident (!B \lor !C)$ تے $!A \ident (!B \lif !C)$۔
\end{prob}

\begin{prob}
ثابت کرو کہ ہر کلاسیکی ہمیشہ صادق فارمولا~$\LogHal$ وچ وی
ہمیشہ صادق اے۔
\end{prob}

بھاویں ایہناں دے ہمیشہ صادق فارمولے کلاسیکی منطق والے ای نیں،
منطقی نتیجے دے تعلق وکھرے نیں۔ مثلاً $\LogLP$
\emph{تضاد برداشت کرن والی} اے، ایس معنٰی وچ کہ
$\lnot p, p \Entails/ q$۔ ایس لئی تضاد توں کجھ وی کڈھن دا
اصول $\lnot !A, !A \Entails !B$ عام طور اُتے قائم نہیں۔
(ایہ $!A$ تے $!B$ دیاں کجھ صورتاں وچ قائم اے، مثلاً جے
$!B$ ہمیشہ صادق ہووے۔)

\begin{prob}
  ہِٹھلے تعلقاں وچوں کیہڑے (الف)~$\LogLP$ وچ تے (ب)~$\LogHal$
  وچ قائم نیں؟ ہر اک لئی صداقتی جدول دیو۔
  \begin{enumerate}
    \item $p, p \lif q \Entails q$
    \item $\lnot q, p \lif q \Entails \lnot p$
    \item $p \lor q, \lnot p \Entails q$
    \item $\lnot p, p \Entails q$
    \item $p \Entails p \lor q$
    \item $p \lif q, q\lif r \Entails p \lif r$
  \end{enumerate}
\end{prob}

جے $\Undef$ نوں $\LogLuk[3]$ وچ مقرر قدر بنا دیو، تاں کیہ ہووے گا؟

\begin{defn}
  منطق \emph{تِن قدراں والی آر-مِنگل}~$\LogRM[3]$ ایس میٹرکس نال متعین اے۔
  \begin{enumerate}
    \item معیاری قضیاتی زبان $\Lang L_0$ جیہدے وچ
    $\lfalse$، $\lnot$، $\land$، $\lor$، $\lif$ نیں۔
    \item صداقتی قدراں دا سیٹ $V = \{\True, \Undef, \False\}$۔
    \item $\True$ تے $\Undef$ مقرر نیں، یعنی $V^+ = \{\True, \Undef\}$۔
    \item صداقتی فنکشن لوکاشیویچ دی منطق~$\LogLuk[3]$ والے ای نیں۔
  \end{enumerate}
\end{defn}

\begin{prob}
  ہِٹھلے تعلقاں وچوں کیہڑے $\LogRM[3]$ وچ قائم نیں؟
  \begin{enumerate}
    \item $p, p \lif q \Entails q$
    \item $p \lor q, \lnot p \Entails q$
    \item $\lnot p, p \Entails q$
    \item $p \Entails p \lor q$
  \end{enumerate}
\end{prob}

کدے وکھ وکھ صداقتی جدولاں توں صرف مقرر قدراں بدل کے اکّو
منطق، یعنی اکّو منطقی نتیجے دا تعلق، مل سکدا اے۔ مثلاً ایہ
اُدوں ہوندا اے جدوں گوڈل دی منطق وچ $V^+ = \{\True,
\Undef\}$، ~$\{\True\}$ دی تھاں لئیے۔

\begin{prop}\ollabel{prop:gl-udes}
  جیہدے وچ $V = \{\False, \Undef,\True\}$، $V^+=\{\True,
  \Undef\}$ تے $3$ قدراں والی گوڈل دی منطق دے صداقتی فنکشن
  ہون، اوہ میٹرکس کلاسیکی منطق متعین کردا اے۔
\end{prop}

\begin{proof}
  مشق۔
\end{proof}

\begin{prob}
  \olref[mvl][thr][mul]{prop:gl-udes} ثابت کرن لئی دکھاؤ
  کہ منطق~$\Log L$، جیہڑی گوڈل دی منطق وانگ متعین اے پر
  $V^+=\{\True,\Undef\}$ اے، اوہدے لئی جے $\Gamma \Entails/[\Log L] !B$،
  تاں $\Gamma \Entails/[\LogCL] !B$۔ \olref[mvl][thr][mul]{prop:LP-taut-CL}
  دے خیال ورتو۔ پر (الف) تے (ب) ثابت کرن دی تھاں دکھاؤ کہ
  $\pValue v(!A)[\LogGod] = \False$ اُدوں تے صرف اُدوں اے
  جدوں $\pValue {v'}(!A)[\LogCL] = \False$۔ ایس توں
  $\pValue
  v(!A)[\LogGod] \in \{\True,\Undef\}$ اُدوں تے صرف اُدوں اے
  جدوں $\pValue {v'}(!A)[\LogCL] =
  \True$۔ سمجھاؤ کہ ایس توں قضیہ کیوں ملدا اے۔
\end{prob}
\endgroup
% /OLP-0397

\OLEndChapterHook
\endgroup
% /OLP-0392

% OLP-0398 content/many-valued-logic/infinite-valued-logics/infinite-valued-logics.tex
\begingroup
\def\olfilename{infinite-valued-logics}
\expandafter\def\csname ol@sectioncs\endcsname{section}

\olchapter{mvl}{inf}{لامتناہی قدراں والیاں منطقاں}

% OLP-0399 content/many-valued-logic/infinite-valued-logics/introduction.tex
\begingroup
\def\olfilename{introduction}
\expandafter\def\csname ol@sectioncs\endcsname{section}

\olfileid{mvl}{inf}{int}

\olsection{جان پہچان}

میٹرکس وچ صداقتی قدراں دی گنتی محدود ہونی ضروری نہیں۔
لامتناہی صداقتی قدراں لئی اک صاف چُون $0$ تے~$1$ دے وچکارلے
ناطق عدداں دا سیٹ $V_\infty = [0,1] \cap \Rat$ اے، یعنی
\begin{align*}
    V_\infty & = \Setabs{\frac{n}{m}}{n,m \in \Nat \text{ تے } n\le m}.
\intertext{صداقتی قدراں دے ایس لامتناہی سیٹ نوں ویکھدیاں ایہ ذیلی سیٹ وی اکثر کم آؤندے نیں۔}
V_m & = \Setabs{\frac{n}{m-1}}{n \in \Nat \text{ تے } n\le m}
\intertext{مثلاً $V_5$ اوہ سیٹ اے جیہدیاں $5$ صداقتی قدراں برابر فاصلے اُتے نیں۔}
V_5 & = \{0, \frac{1}{4}, \frac{1}{2}, \frac{3}{4}, 1\}.
\end{align*}
ایہناں صداقتی سیٹاں اُتے بنیاں منطقاں وچ عام طور اُتے صرف
$1$ مقرر ہوندا اے، یعنی $V^+ = \{1\}$۔ دوجے لفظاں وچ
$1$ نوں (پورا) سچ تے $0$ نوں پورا جھوٹ سمجھدے آں، پر
!!{formula}s~$V$ وچ کوئی وی درمیانی قدر لے سکدے نیں۔

فیر وی اسیں $V_{[0,1]} = [0,1]$، یعنی $0$ تے~$1$
دے وچکارلے سارے \emph{حقیقی} عدداں دا سیٹ، یا $[0,1]$ دے
ہور لامتناہی ذیلی سیٹ وی ویکھ سکدے آں۔ ایس صداقتی سیٹ والیاں
منطقاں نوں اکثر \emph{دھندلیاں منطقاں} آکھدے آں۔
\endgroup
% /OLP-0399

% OLP-0400 content/many-valued-logic/infinite-valued-logics/lukasiewicz.tex
\begingroup
\def\olfilename{lukasiewicz}
\expandafter\def\csname ol@sectioncs\endcsname{section}

\olfileid{mvl}{inf}{luk}

\olsection{لوکاشیویچ دی منطق}

\begin{editorial}
  ایہ لامتناہی قدراں والی لوکاشیویچ دی منطق بارے حصے دا نکا جیہا مسودہ اے۔
\end{editorial}

\begin{defn}\ollabel{def:lukasiewicz}
لامتناہی قدراں والی لوکاشیویچ دی منطق~$\LogLuk[\infty]$ ایس
میٹرکس نال متعین اے۔
\begin{enumerate}
  \item معیاری قضیاتی زبان $\Lang L_0$ جیہدے وچ
  $\lnot$، $\land$، $\lor$، $\lif$ نیں۔
  \item صداقتی قدراں دا سیٹ $V_\infty$۔
  \item $1$ اکّو مقرر قدر اے، یعنی $V^+ = \{1\}$۔
  \item صداقتی فنکشن ہِٹھ دِتّے نیں۔
  \begin{align*}
    \tf{\lnot}[\LogLuk](x) & = 1 - x\\
    \tf{\land}[\LogLuk](x,y) & = \min(x,y)\\
    \tf{\lor}[\LogLuk](x,y) & = \max(x,y)\\
    \tf{\lif}[\LogLuk](x,y) & = \min(1,1-(x-y)) = \begin{cases}
      1 & \text{جے } x \le y\\
      1-(x-y) & \text{ہور صورت وچ۔}
    \end{cases}
    \end{align*}
\end{enumerate}
$m$ قدراں والی لوکاشیویچ دی منطق ایسے طرح متعین اے، بس $V = V_m$ اے۔
\end{defn}

\begin{prop}
  \olref[thr][luk]{def:lukasiewicz} وچ متعین منطق $\LogLuk[3]$،
  \olref{def:lukasiewicz} وچ متعین منطق $\LogLuk[3]$ دے عین برابر اے۔
\end{prop}

\begin{proof}
  ایہ گل \olref[thr][luk]{def:lukasiewicz} وچ دِتّے رابطیاں دے
  صداقتی جدولاں نوں \olref{def:lukasiewicz} دیاں مساواتاں توں
  ملن والے جدولاں نال ملا کے دِسدی اے۔
  \begin{center}
    \begin{tabular}{c|c} 
      $\tf{\lnot}$ & \\ 
      \hline  
      $1$ & $0$ \\ 
      $1/2$ & $1/2$ \\
      $0$ & $1$ 
    \end{tabular}
    \quad
    \begin{tabular}{c|ccc} 
      $\tf{\land}[\LogLuk[3]]$ & $1$ & $1/2$ & $0$ \\ 
      \hline 
      $1$ & $1$ & $1/2$ & $0$ \\ 
      $1/2$ & $1/2$ & $1/2$ & $0$\\ 
      $0$ & $0$ & $0$ & $0$ 
    \end{tabular}
    \\[2ex]
    \begin{tabular}{c|ccc} 
      $\tf{\lor}[\LogLuk[3]]$ & $1$ & $1/2$ & $0$ \\ 
      \hline 
      $1$ & $1$ & $1$ & $1$ \\ 
      $1/2$ & $1$ & $1/2$ & $1/2$ \\
      $0$ & $1$ & $1/2$ & $0$ 
    \end{tabular}
    \quad
    \begin{tabular}{c|ccc} 
      $\tf{\lif}[\LogLuk[3]]$ & $1$ & $1/2$ & $0$ \\ 
      \hline 
      $1$ & $1$ & $1/2$ & $0$ \\ 
      $1/2$ & $1$ & $1$ & $1/2$  \\ 
      $0$ & $1$ & $1$ & $1$ 
    \end{tabular}
  \end{center}
\end{proof}

\begin{prop}\ollabel{prop:luk-infty-m}
  جے $\Gamma \Entails[\LogLuk[\infty]] !B$، تاں $\Gamma
  \Entails[\LogLuk[m]] !B$ ہر~$m \ge 2$ لئی۔
\end{prop}

\begin{proof}
  مشق۔
\end{proof}

\begin{prob}
  \olref[mvl][inf][luk]{prop:luk-infty-m} ثابت کرو۔
\end{prob}

اصل وچ اُلٹی گل وی قائم اے۔

لامتناہی قدراں والی لوکاشیویچ دی منطق سبھ توں مقبول دھندلی
منطق اے۔ دھندلی منطق بارے لکھتاں وچ شرطیہ اکثر $\lnot !A \lor !B$
وجوں متعین ہوندا اے۔ ایس توں لامتناہی قدراں والی کلینی دی مضبوط منطق
بنے گی۔

\begin{prob}
  دکھاؤ کہ $(p \lif q) \lor (q \lif p)$، ~$\LogLuk[\infty]$
  وچ ہمیشہ صادق اے۔
\end{prob}
\endgroup
% /OLP-0400

% OLP-0401 content/many-valued-logic/infinite-valued-logics/goedel.tex
\begingroup
\def\olfilename{goedel}
\expandafter\def\csname ol@sectioncs\endcsname{section}

\olfileid{mvl}{inf}{god}

\olsection{گوڈل دیاں منطقاں}

\begin{editorial}
  ایہ لامتناہی قدراں والی گوڈل دی منطق بارے حصے دا نکا جیہا مسودہ اے۔
\end{editorial}

\begin{defn}\ollabel{def:goedel}
لامتناہی قدراں والی گوڈل دی منطق~$\LogGod[\infty]$ ایس میٹرکس
نال متعین اے۔
\begin{enumerate}
  \item معیاری قضیاتی زبان $\Lang L_0$ جیہدے وچ
  $\lfalse$، $\lnot$، $\land$، $\lor$، $\lif$ نیں۔
  \item صداقتی قدراں دا سیٹ $V_\infty$۔
  \item $1$ اکّو مقرر قدر اے، یعنی $V^+ = \{1\}$۔
  \item صداقتی فنکشن ہِٹھ دِتّے نیں۔
  \begin{align*}
    \tf{\lfalse} & = 0\\
    \tf{\lnot}[\LogGod](x) & = \begin{cases}
      $1$ & \text{جے } x =0\\
      $0$ & \text{ہور صورت وچ}
    \end{cases}\\
    \tf{\land}[\LogGod](x,y) & = \min(x,y)\\
    \tf{\lor}[\LogGod](x,y) & = \max(x,y)\\
    \tf{\lif}[\LogGod](x,y) & = \begin{cases}
      1 & \text{جے } x \le y\\
      y & \text{ہور صورت وچ۔}
    \end{cases}
    \end{align*}
\end{enumerate}
$m$ قدراں والی گوڈل دی منطق ایسے طرح متعین اے، بس $V = V_m$ اے۔
\end{defn}

\begin{prop}
  \olref[thr][god]{defn:goedel} وچ متعین منطق $\LogGod[3]$،
  \olref{def:goedel} وچ متعین منطق $\LogGod[3]$ دے عین برابر اے۔
\end{prop}

\begin{proof}
  ایہ گل \olref[thr][god]{defn:goedel} وچ دِتّے رابطیاں دے
  صداقتی جدولاں نوں \olref{def:goedel} دیاں مساواتاں توں ملن
  والے جدولاں نال ملا کے دِسدی اے۔
  \begin{center}
    \begin{tabular}{c|c} 
      $\tf{\lnot}[\LogGod[3]]$ & \\ 
      \hline  
      $1$ & $0$ \\ 
      $1/2$ & $0$ \\
      $0$ & $1$ 
    \end{tabular}
    \quad
    \begin{tabular}{c|ccc} 
      $\tf{\land}[\LogGod]$ & $1$ & $1/2$ & $0$ \\ 
      \hline 
      $1$ & $1$ & $1/2$ & $0$ \\ 
      $1/2$ & $1/2$ & $1/2$ & $0$\\ 
      $0$ & $0$ & $0$ & $0$ 
    \end{tabular}
    \\[2ex]
    \begin{tabular}{c|ccc} 
      $\tf{\lor}[\LogGod]$ & $1$ & $1/2$ & $0$ \\ 
      \hline 
      $1$ & $1$ & $1$ & $1$ \\ 
      $1/2$ & $1$ & $1/2$ & $1/2$ \\
      $0$ & $1$ & $1/2$ & $0$ 
    \end{tabular}
    \quad
    \begin{tabular}{c|ccc} 
      $\tf{\lif}[\LogGod]$ & $1$ & $1/2$ & $0$ \\ 
      \hline 
      $1$ & $1$ & $1/2$ & $0$ \\ 
      $1/2$ & $1$ & $1$ & $0$  \\ 
      $0$ & $1$ & $1$ & $1$ 
    \end{tabular}
  \end{center}
\end{proof}

\begin{prop}\ollabel{prop:god-infty-m}
  جے $\Gamma \Entails[\LogGod[\infty]] !B$، تاں $\Gamma
  \Entails[\LogGod[m]] !B$ ہر~$m \ge 2$ لئی۔
\end{prop}

\begin{proof}
  مشق۔
\end{proof}

\begin{prob}
  \olref[mvl][inf][god]{prop:god-infty-m} ثابت کرو۔
\end{prob}

اصل وچ اُلٹی گل وی قائم اے۔

$\LogGod[3]$ وانگ، $\LogGod[\infty]$ وچ اوہ سارے !!{formula}s
ہمیشہ صادق نیں جیہڑے وجدانیت پسند منطق وچ معتبر نیں۔ اوہی
مثالاں جیہڑیاں ہمیشہ صادق نہیں سن، ~$\LogGod[\infty]$ وچ وی
ہمیشہ صادق نہیں۔
\begin{align*}
  & p \lor \lnot p && (p \lif q) \lif (\lnot p \lor q) \\
  & \lnot\lnot p \lif p && \lnot(\lnot p \land \lnot q) \lif (p \lor q) \\
  & ((p \lif q) \lif p) \lif p && \lnot(p \lif q) \lif (p \land \lnot q)
\end{align*}
اوہ !!{formula} جیہڑا وجدانیت پسند منطق وچ معتبر نہیں پر
~$\LogGod[3]$ وچ ہمیشہ صادق اے، یعنی $(p \lif q) \lor (q \lif
p)$، ~$\LogGod[\infty]$ وچ وی ہمیشہ صادق اے۔ دراصل
$\LogGod[\infty]$ نوں وجدانیت پسند منطق وچ سانچا
$(!A \lif !B) \lor (!B \lif !A)$ شامل کر کے حاصل کر سکدے آں۔
ایہ گل مائیکل ڈمٹ نے دکھائی سی۔ ایس لئی $\LogGod[\infty]$
نوں اکثر گوڈل–ڈمٹ منطق~$\Log{LC}$ آکھدے آں۔

\begin{prob}
  دکھاؤ کہ $(p \lif q) \lor (q \lif p)$، ~$\LogGod[\infty]$
  وچ ہمیشہ صادق اے۔
\end{prob}

\begin{prob}
  دکھاؤ کہ $(p \lif q) \lor (q \lif r) \lor (r \lif s)$،
  جیہڑا $\LogGod[3]$ وچ ہمیشہ صادق اے، ~$\LogGod[\infty]$ وچ
  ہمیشہ صادق نہیں۔
\end{prob}
\endgroup
% /OLP-0401

\OLEndChapterHook
\endgroup
% /OLP-0398

% OLP-0402 content/many-valued-logic/sequent-calculus/sequent-calculus.tex
\begingroup
\def\olfilename{sequent-calculus}
\expandafter\def\csname ol@sectioncs\endcsname{section}

\olchapter{mvl}{seq}{سیکوئنٹ حساب}

% OLP-0403 content/many-valued-logic/sequent-calculus/introduction.tex
\begingroup
\def\olfilename{introduction}
\expandafter\def\csname ol@sectioncs\endcsname{section}

\olfileid{mvl}{seq}{int}

\olsection{جان پہچان}

کلاسیکی منطق دا سیکوئنٹ حساب سوکھا تے کارگر !!{derivation}
دا نظام اے۔ جے کئی قدراں والی منطق ایہو جیہے میٹرکس نال متعین
ہووے جیہدیاں صداقتی قدراں محدود گنتی وچ نیں، یعنی $V$ محدود
ہووے، تاں اوہدے لئی سیکوئنٹ حساب دِتّا جا سکدا اے۔ ایہ کم کرن
دا خیال کلاسیکی صورت وچ سیکوئنٹاں دے معنیاں تے استنباط دے
قاعدیاں دے روپ نوں ویکھ کے آؤندا اے۔

ہُن یاد کرو کہ سیکوئنٹ
\begin{align*}
    \formula A_1, \dots, !A_n & \Sequent !B_1, \dots, !B_n
\intertext{نوں ہِٹھلے !!{formula} دے طور اُتے سمجھ سکدے آں۔}
    (!A_1 \land \cdots \land !A_m) & \lif (!B_1 \lor \cdots \lor
\formula B_n)
\end{align*}
دوجے لفظاں وچ، !!^a{valuation}~$\pAssign{v}$ نال سیکوئنٹ
$\Gamma \Sequent \Delta$ اُدوں تے صرف اُدوں \emph{پورا
اتردا اے} جدوں یا $\pValue{v}(!A) = \False$ ہووے کسے
$!A \in \Gamma$ لئی، یا $\pValue{v}(!A) = \True$ ہووے کسے
$!A \in
\Delta$ لئی۔ ایس مطلب مطابق ابتدائی سیکوئنٹ $!A \Sequent !A$
ہمیشہ پورے اتردے نیں، کیوں جے یا $\pValue v(!A) = \True$ اے
یا $\pValue(!A) =
\False$۔

~$\Log{LK}$ وچ شرطیے دے استنباطی قاعدے ہِٹھ دِتّے نیں۔
پاسیاں والے فارمولے $\Gamma$، $\Delta$ چھڈ دِتّے نیں۔

\begin{defish}
    \Axiom$ \fCenter !A$
    \Axiom$ !B \fCenter $
    \RightLabel{\LeftR{\lif}}
    \BinaryInf$ !A \lif !B \fCenter $
    \DisplayProof
    \hfill
    \Axiom$ !A \fCenter !B$
    \RightLabel{\RightR{\lif}}
    \UnaryInf$ \fCenter !A \lif !B $
    \DisplayProof
\end{defish}

جے سیکوئنٹ دی اُپرلی معنوی سمجھ ورتیئے، تاں $\LeftR{\lif}$
قاعدے دا مطلب ایہ اے کہ جے $\pValue v(!A) = \True$ تے
$\pValue v(!B) = \False$، تاں $\pValue v(!A \lif !B) = \False$۔
ایسے طرح $\RightR{\lif}$ قاعدہ آکھدا اے کہ جے
$\pValue v(!A)
= \False$ یا $\pValue v(!B) = \True$، تاں
$\pValue v(!A \lif !B) =
\True$۔ اصل وچ ایہ دوویں شرطیے دو شرطیے نیں۔ $\LeftR{\land}$
تے $\RightR{\lor}$ دے معیاری قاعدیاں وچ متعلقہ شرطیے دو
شرطیے نہیں ہوندے۔ پر ایہناں قاعدیاں دے ہور روپ وی نیں جنہاں
وچ اوہ دو شرطیے ہوندے نیں۔

\begin{defish}
    \Axiom$!A, !B, \Gamma \fCenter \Delta$
    \RightLabel{\LeftR{\land}}
    \UnaryInf$!A \land !B, \Gamma \fCenter \Delta$
    \DisplayProof
    \hfill
    \Axiom$ \Gamma \fCenter \Delta, !A, !B$
    \RightLabel{\RightR{\lor}}
    \UnaryInf$ \Gamma \fCenter \Delta, !A \lor !B$
    \DisplayProof
\end{defish}

ایہ بنیادی خیال $n$ قدراں والی منطق اُتے لائیے، تاں
سیکوئنٹ حساب وچ دو دیاں تھاواں~$n$ تھاواں ہوندیاں نیں،
ہر صداقتی قدر لئی اک۔ تِن قدراں والی منطق جیہدے لئی
$V = \{\False, \Undef, \True\}$، اوہدے وچ سیکوئنٹ عبارت
$\Gamma \mid \Pi \mid \Delta$ اے۔ ایہ !!a{valuation}~$\pAssign v$
وچ اُدوں تے صرف اُدوں پورا اتردا اے جدوں یا
$\pValue{v}(!A) =
\False$ کسے $!A \in \Gamma$ لئی ہووے، یا $\pValue{v}(!A) = \True$
کسے $!A \in \Delta$ لئی، یا $\pValue{v}(!A) = \Undef$ کسے
$!A \in \Pi$ لئی۔ ایس کرکے ابتدائی سیکوئنٹ $!A \mid !A \mid !A$
ہمیشہ پورے اتردے نیں۔
\endgroup
% /OLP-0403

% OLP-0404 content/many-valued-logic/sequent-calculus/rules-and-proofs.tex
\begingroup
\def\olfilename{rules-and-proofs}
\expandafter\def\csname ol@sectioncs\endcsname{section}

\olfileid{mvl}{seq}{rul}

\olsection{قاعدے تے \usetoken{P}{derivation}}

اگلی گل لئی منّ لو $\Gamma, \Delta, \Pi, \Lambda$،
!!{sentence}s دے محدود سلسلے نیں۔

\begin{defn}[سیکوئنٹ]
\emph{$n$ پاسیاں والا سیکوئنٹ} ایس روپ دی عبارت اے۔
\[
\Gamma_1 \nSequent \dots \nSequent \Gamma_n
\]
ایتھے ہر $\Gamma_1$ !!{sentence}s دا محدود سلسلہ اے، جیہڑا خالی وی ہو سکدا اے۔
ایہ جملے زبان~$\Lang L$ دے نیں۔
\end{defn}

\begin{defn}[ابتدائی سیکوئنٹ]
\emph{$n$ پاسیاں والا ابتدائی سیکوئنٹ} اوہ $n$ پاسیاں والا
سیکوئنٹ اے جیہدا روپ $!A \nSequent \dots \nSequent !A$ اے،
زبان دے کسے وی !!{sentence} $!A$ لئی۔

جے زبان وچ $0$-مقامی رابطہ~$\star$ اے، یعنی قضیاتی مستقل،
تاں اسیں سیکوئنٹ $\dots \nSequent
\star \nSequent \dots$ وی لیندے آں۔ ایس وچ $\star$ متعلقہ
صداقتی قدر~$\tf{\star} \in V$ والی تھاں اُتے آؤندا اے تے
باقی تھاواں خالیاں نیں۔
\end{defn}

$n$ قدراں والی منطق~$\Log{L}$ دے ہر رابطے لئی، ہر اوہ
صداقتی قدر دا اک منطقی قاعدہ اے جیہڑی اوہ رابطہ~$\Log{L}$
وچ لے سکدا اے۔ $n$ پاسیاں والے سیکوئنٹ حساب وچ~$\Log{L}$
دے !!^{derivation}s سیکوئنٹاں دے درخت نیں۔ سبھ توں اُتّے
ابتدائی سیکوئنٹ ہوندے نیں۔ جے کوئی سیکوئنٹ اک یا کئی ہوراں
دے ہِٹھ اے، تاں اوہ~$\Log{L}$ دے رابطیاں دے کسے استنباطی
قاعدے توں صحیح طرح نکلنا چاہیدا اے۔

\begin{defn}[قضیے]
!!^a{sentence}~$!A$، $n$ قدراں والی منطق~$\Log{L}$ دا
\emph{قضیہ} اے جے !!a{derivation} اوہ $n$ پاسیاں والے
سیکوئنٹ دا موجود ہووے جیہدے وچ $!A$ ہر اوہ تھاں اُتے آوے جیہڑی~$\Log{L}$ دی مقرر
صداقتی قدر نال جُڑی اے۔ اسیں $\Proves[\Log{L}]
!A$ لکھدے آں جے $!A$ قضیہ اے، تے جے ایہ قضیہ نہیں تاں $\Proves/[\Log{L}] !A$۔
\end{defn}

\begin{defn}[!!^{derivability}]
!!^a{sentence}~$!A$، !!{sentence}s دے سیٹ~$\Gamma$ توں
$n$ قدراں والی منطق~$\Log{L}$ وچ \emph{!!{derivable} اے}،
یعنی $\Gamma
\Proves[\Log{L}] !A$، اُدوں تے صرف اُدوں جدوں اک محدود
ذیلی سیٹ~$\Gamma_0 \subseteq
\Gamma$ تے اک سلسلہ $\Gamma_0'$ ہووے جیہڑا~$\Gamma_0$ دے
!!{sentence}s توں بنے، تے ہِٹھلے سیکوئنٹ دا
!!a{derivation} ہووے۔
\[ \Lambda_1 \nSequent \dots \nSequent \Lambda_n \]
ایتھے $\Lambda_i$، $!A$ اے جے تھاں $i$ مقرر صداقتی قدر نال
جُڑی ہووے، تے ہور صورت وچ $\Gamma_0'$ اے۔ جے $!A$،
$\Gamma$ توں !!{derivable} نہیں، تاں $\Gamma \Proves/ !A$
لکھدے آں۔
\end{defn}

مثلاً $3$ قدراں والی لوکاشیویچ دی منطق دا $3$ پاسیاں والا
سیکوئنٹ حساب اے۔ $3$ پاسیاں والے سیکوئنٹ $\Gamma \nSequent \Pi \nSequent
\Delta$ وچ $\Gamma$،~$\False$ نال، $\Delta$،~$\True$ نال،
تے $\Pi$،~$\Undef$ نال جُڑیا اے۔ مسلّمے $!A \nSequent !A \nSequent !A$
نیں۔ کیوں جے بس $\True$ مقرر اے، ایس لئی
$\Gamma \Proves[\LogLuk[3]] !A$ اُدوں تے صرف اُدوں اے جدوں
سیکوئنٹ $\Gamma \nSequent \Gamma \nSequent !A$ دا !!a{derivation}
ہووے۔ (جے $\Undef$ وی مقرر ہووے، تاں !!a{derivation} ہِٹھلے سیکوئنٹ دا چاہیدا ہووے گا۔
$\Gamma \nSequent !A \nSequent !A$۔)
\endgroup
% /OLP-0404

% OLP-0405 content/many-valued-logic/sequent-calculus/structural-rules.tex
\begingroup
\def\olfilename{structural-rules}
\expandafter\def\csname ol@sectioncs\endcsname{section}

\olfileid{mvl}{seq}{str}

\olsection{ساختی قاعدے}

$n$ پاسیاں والے سیکوئنٹ حساب دے ساختی قاعدے کلاسیکی
صورت وانگ کم کردے نیں، پر ہر تھاں~$i$ لئی وکھرے طور اُتے۔

\begin{defish}
\begin{center}
\AxiomC{$ \Gamma_1 \nSequent \dots \nSequent \phantom{!A,}\Gamma_i \nSequent \dots \nSequent \Gamma_n $}
\RightLabel{$\iR{\Weakening}{i}$}
\UnaryInfC{$\Gamma_1 \nSequent \dots \nSequent !A, \Gamma_i \nSequent \dots \nSequent \Gamma_n$}
\DisplayProof
\\[2ex]
\AxiomC{$ \Gamma_1 \nSequent \dots \nSequent !A, !A, \Gamma_i \nSequent \dots \nSequent \Gamma_n $}
\RightLabel{$\iR{\Contraction}{i}$}
\UnaryInfC{$\Gamma_1 \nSequent \dots \nSequent \phantom{!A,}!A, \Gamma_i \nSequent
\dots \nSequent \Gamma_n$}
\DisplayProof
\\[2ex]
\AxiomC{$ \Gamma_1 \nSequent \dots \nSequent \Gamma_i, !A, !B, \Gamma_i' \nSequent \dots \nSequent \Gamma_n $}
\RightLabel{$\iR{\Exchange}{i}$}
\UnaryInfC{$\Gamma_1 \nSequent \dots \nSequent \Gamma_i, !B, !A, \Gamma_i' \nSequent \dots
\nSequent \Gamma_n$}
\DisplayProof
\end{center}
\end{defish}

کمزور کرن، سکیڑن تے تھاں وٹاؤن دے کئی قدماں نوں اکثر دوہری
استنباطی لکیر نال دکھاؤندے آں۔

\Cut{} دا قاعدہ کئی روپاں وچ آؤندا اے۔ سیکوئنٹ دیاں دو
وکھریاں تھاواں $i \neq j$ دی ہر چُون لئی اک روپ اے۔
\begin{defish}
\[
\AxiomC{$ \Gamma_1 \nSequent \dots \nSequent !A, \Gamma_i \nSequent \dots \nSequent \Gamma_n $}
\AxiomC{$ \Delta_1 \nSequent \dots \nSequent !A, \Delta_j \nSequent \dots \nSequent \Delta_n $}
\RightLabel{$\iR{\Cut}{i,j}$}
\BinaryInfC{$\Gamma_1,\Delta_1 \nSequent \dots \nSequent \Gamma_n, \Delta_n$}
\DisplayProof
\]
\end{defish}
\endgroup
% /OLP-0405

% OLP-0406 content/many-valued-logic/sequent-calculus/propositional-rules.tex
\begingroup
\def\olfilename{propositional-rules}
\expandafter\def\csname ol@sectioncs\endcsname{section}

\olfileid{mvl}{seq}{prl}

\olsection{چُنیاں منطقاں دے قضیاتی قاعدے}

$n$ پاسیاں والے سیکوئنٹ حساب وچ رابطے دے استنباطی قاعدے
صرف اوس رابطے دے خاص صداقتی فنکشن اُتے ٹکے ہوندے نیں۔
سو جے وکھ وکھ منطقاں وچ کوئی رابطہ اکّو صداقتی فنکشن نال
متعین اے، تاں رابطے دے ایہ $n$ پاسیاں والے سیکوئنٹ قاعدے
اوہناں منطقاں وچ اکّو نیں۔

\subsection{$\lnot$ دے قاعدے}

$\lnot$ دے ہِٹھلے قاعدے لوکاشیویچ تے کلینی دیاں منطقاں تے
اوہناں دے ہور روپاں اُتے لاگو ہوندے نیں۔

\begin{defish}
  \begin{center}
\AxiomC{$ \Gamma \nSequent \Pi \nSequent \Delta, !A $}
\RightLabel{\iR{\lnot}{\False}}
\UnaryInfC{$ \lnot !A, \Gamma \nSequent \Pi \nSequent \Delta$}
\DisplayProof
\\[2ex]
\AxiomC{$ \Gamma \nSequent !A, \Pi \nSequent \Delta $}
\RightLabel{\iR{\lnot}{\Undef}}
\UnaryInfC{$\Gamma \nSequent \lnot !A, \Pi \nSequent \Delta$}
\DisplayProof
\\[2ex]
\AxiomC{$!A, \Gamma \nSequent \Pi \nSequent \Delta$}
\RightLabel{\iR{\lnot}{\True}}
\UnaryInfC{$\Gamma \nSequent \Pi \nSequent \Delta,  \lnot !A$}
\DisplayProof
  \end{center}
\end{defish}

$\lnot$ دے ہِٹھلے قاعدے گوڈل دی منطق اُتے لاگو ہوندے نیں۔

\begin{defish}
\AxiomC{$ \Gamma \nSequent !A, \Pi \nSequent \Delta, !A $}
\RightLabel{\iR{\lnot}{\False}[\LogGod]}
\UnaryInfC{$ \lnot !A, \Gamma \nSequent \Pi \nSequent \Delta$}
\DisplayProof
\hfill
\AxiomC{$!A, \Gamma \nSequent \Pi \nSequent \Delta$}
\RightLabel{\iR{\lnot}{\True}[\LogGod]}
\UnaryInfC{$\Gamma \nSequent \Pi \nSequent \Delta,  \lnot !A$}
\DisplayProof
\hfill
\end{defish}

(گوڈل دی منطق وچ $\lnot !A$ کدے قدر~$\Undef$ نہیں لے سکدا۔
سو وچکارلی تھاں لئی کوئی قاعدہ نہیں۔)

\subsection{$\land$ دے قاعدے}

ایہ $\land$ دے قاعدے لوکاشیویچ، کلینی دی مضبوط تے گوڈل دیاں
منطقاں لئی نیں۔

\begin{defish}
\begin{center}
  \AxiomC{$!A, !B, \Gamma \nSequent \Pi \nSequent \Delta$}
  \RightLabel{$\iR{\land}{\False}$}
  \UnaryInfC{$!A \land !B, \Gamma \nSequent \Pi \nSequent \Delta$}
  \DisplayProof
\\[2ex]
  \AxiomC{$\Gamma \nSequent !A, \Pi \nSequent !A, \Delta$}
  \AxiomC{$\Gamma \nSequent !B, \Pi \nSequent !B, \Delta$}
  \AxiomC{$\Gamma \nSequent !A, !B, \Pi \nSequent \Delta$}
  \RightLabel{$\iR\land\Undef$}
  \TrinaryInfC{$\Gamma \nSequent !A \land !B, \Pi \nSequent \Delta$}
  \DisplayProof
  \\[2ex]  
  \AxiomC{$\Gamma \nSequent \Pi \nSequent \Delta, !A$}
  \AxiomC{$\Gamma \nSequent \Pi \nSequent \Delta, !B$}
  \RightLabel{$\iR\land\True$}
  \BinaryInfC{$\Gamma \nSequent \Pi \nSequent \Delta, !A \land !B$}
  \DisplayProof
\end{center}
\end{defish}

\subsection{$\lor$ دے قاعدے}

ایہ $\lor$ دے قاعدے لوکاشیویچ، کلینی دی مضبوط تے گوڈل دیاں
منطقاں لئی نیں۔

\begin{defish}
\begin{center}
 \AxiomC{$!A, \Gamma \nSequent \Pi \nSequent \Delta$}
  \AxiomC{$!B, \Gamma \nSequent \Pi \nSequent \Delta$}
  \RightLabel{$\iR\lor\False$}
  \BinaryInfC{$!A \lor !B, \Gamma \nSequent \Pi \nSequent \Delta$}
  \DisplayProof
\\[2ex]
  \AxiomC{$!A, \Gamma \nSequent !A, \Pi \nSequent \Delta$}
  \AxiomC{$!B, \Gamma \nSequent !B, \Pi \nSequent \Delta$}
  \AxiomC{$\Gamma \nSequent !A, !B, \Pi \nSequent \Delta$}
  \RightLabel{$\iR\lor\Undef$}
  \TrinaryInfC{$\Gamma \nSequent !A \lor !B, \Pi \nSequent \Delta$}
  \DisplayProof
  \\[2ex]  
  \AxiomC{$\Gamma \nSequent \Pi \nSequent \Delta, !A, !B$}
  \RightLabel{$\iR\lor\True$}
  \UnaryInfC{$\Gamma \nSequent \Pi \nSequent \Delta, !A \lor !B$}
  \DisplayProof
\end{center}
\end{defish}

\subsection{$\lif$ دے قاعدے}

ایہ $\lif$ دے قاعدے لوکاشیویچ دی منطق لئی نیں۔

\begin{defish}
\begin{center}
  \AxiomC{$\Gamma \nSequent \Pi \nSequent \Delta, !A$}
  \AxiomC{$!B, \Gamma \nSequent \Pi \nSequent \Delta$}
  \RightLabel{$\iR\lif\False[\LogLuk[3]]$}
  \BinaryInfC{$!A \lif !B, \Gamma \nSequent \Pi \nSequent \Delta$}
  \DisplayProof
  \\[2ex]
  \AxiomC{$\Gamma \nSequent !A, !B, \Pi \nSequent \Delta$}
  \AxiomC{$!B, \Gamma \nSequent \Pi \nSequent \Delta, !A$}
  \RightLabel{$\iR\lif\Undef[\LogLuk[3]]$}
  \BinaryInfC{$\Gamma \nSequent !A \lif !B, \Pi \nSequent \Delta$}
  \DisplayProof
\\[2ex]
  \AxiomC{$!A, \Gamma \nSequent !B, \Pi \nSequent \Delta, !B$}
  \AxiomC{$!A, \Gamma \nSequent !A, \Pi \nSequent \Delta, !B$}
  \RightLabel{$\iR{\lif}{\True}[\LogLuk[3]]$}
  \BinaryInfC{$\Gamma \nSequent \Pi \nSequent \Delta, !A \lif !B$}
  \DisplayProof
\end{center}
\end{defish}

ایہ $\lif$ دے قاعدے کلینی دی مضبوط منطق لئی نیں۔

\begin{defish}
\begin{center}
  \AxiomC{$\Gamma \nSequent \Pi \nSequent \Delta, !A$}
  \AxiomC{$!B, \Gamma \nSequent \Pi \nSequent \Delta$}
  \RightLabel{$\iR\lif\False[\LogKs]$}
  \BinaryInfC{$!A \lif !B, \Gamma \nSequent \Pi \nSequent \Delta$}
  \DisplayProof
  \\[2ex]
  \AxiomC{$!B, \Gamma \nSequent !B, \Pi \nSequent \Delta$}
  \AxiomC{$\Gamma \nSequent !A, !B, \Pi \nSequent \Delta$}
  \AxiomC{$\Gamma \nSequent !A, \Pi \nSequent \Delta, !A$}
  \RightLabel{$\iR\lif\Undef[\LogKs]$}
  \TrinaryInfC{$\Gamma \nSequent !A \lif !B, \Pi \nSequent \Delta$}
  \DisplayProof
\\[2ex]
  \AxiomC{$!A, \Gamma \nSequent \Pi \nSequent \Delta, !B$}
  \RightLabel{$\iR{\lif}{\True}[\LogKs]$}
  \UnaryInfC{$\Gamma \nSequent \Pi \nSequent \Delta, !A \lif !B$}
  \DisplayProof
\end{center}
\end{defish}

ایہ $\lif$ دے قاعدے گوڈل دی منطق لئی نیں۔

\begin{defish}
\begin{center}
  \AxiomC{$\Gamma \nSequent !A, \Pi \nSequent \Delta, !A$}
  \AxiomC{$!B, \Gamma \nSequent \Pi \nSequent \Delta$}
  \RightLabel{$\iR\lif\False[\LogGod[3]]$}
  \BinaryInfC{$!A \lif !B, \Gamma \nSequent \Pi \nSequent \Delta$}
  \DisplayProof
  \\[2ex]
  \AxiomC{$\Gamma \nSequent !B, \Pi \nSequent \Delta$}
  \AxiomC{$\Gamma \nSequent \Pi \nSequent \Delta, !A$}
  \RightLabel{$\iR\lif\Undef[\LogGod[3]]$}
  \BinaryInfC{$\Gamma \nSequent !A \lif !B, \Pi \nSequent \Delta$}
  \DisplayProof
\\[2ex]
  \AxiomC{$!A, \Gamma \nSequent !B, \Pi \nSequent \Delta, !B$}
  \AxiomC{$!A, \Gamma \nSequent !A, \Pi \nSequent \Delta, !B$}
  \RightLabel{$\iR{\lif}{\True}[\LogGod[3]]$}
  \BinaryInfC{$\Gamma \nSequent \Pi \nSequent \Delta, !A \lif !B$}
  \DisplayProof
\end{center}
\end{defish}

\begin{sidewaysfigure}
  \begin{center}  
  \AxiomC{$A \nSequent A \nSequent A$}
  \RightLabel{\iR \Weakening \True}
  \UnaryInfC{$A \nSequent A \nSequent B, A$}
  \RightLabel{\iR \Weakening \Undef}
  \UnaryInfC{$A \nSequent B, A \nSequent B, A$}
  \RightLabel{\iR \Weakening \Undef}
  \UnaryInfC{$A \nSequent A, B, A \nSequent B, A$}
  \AxiomC{$A \nSequent A \nSequent A$}
  \RightLabel{\iR \Weakening \True}
  \UnaryInfC{$A \nSequent A \nSequent A, A$}
  \RightLabel{\iR \Weakening \True}
  \UnaryInfC{$A \nSequent A \nSequent B, A, A$}
  \RightLabel{\iR \Weakening \False}
  \UnaryInfC{$B, A \nSequent A \nSequent B, A, A$}
  \RightLabel{$\iR\lif\Undef$}
  \BinaryInfC{$A \nSequent A \lif B, A \nSequent B, A$}
  \AxiomC{$B \nSequent B \nSequent B$}
  \RightLabel{\iR \Weakening \Undef}
  \UnaryInfC{$B \nSequent A, B \nSequent B$}
  \RightLabel{\iR \Exchange \Undef}
  \UnaryInfC{$B \nSequent B, A \nSequent B$}
  \RightLabel{\iR \Weakening \Undef}
  \UnaryInfC{$B \nSequent A, B, A \nSequent B$}
  \RightLabel{\iR \Weakening \False}
  \UnaryInfC{$A, B \nSequent A, B, A \nSequent B$}
  \RightLabel{\iR \Exchange \False}
  \UnaryInfC{$B, A \nSequent A, B, A \nSequent B$}
  \AxiomC{$A \nSequent A \nSequent A$}
  \RightLabel{\iR \Weakening \True}
  \UnaryInfC{$A \nSequent A \nSequent B, A$}
  \RightLabel{\iR \Weakening \False}
  \UnaryInfC{$B, A \nSequent A \nSequent B, A$}
  \RightLabel{\iR \Weakening \False}
  \UnaryInfC{$B, B, A \nSequent A \nSequent B, A$}
  \RightLabel{$\iR\lif\Undef$}
  \BinaryInfC{$B, A \nSequent A \lif B, A \nSequent B$}
  \RightLabel{$\iR\lif\False$}
  \BinaryInfC{$A \lif B, A \nSequent A \lif B, A \nSequent B$}
  \DisplayProof
  \end{center}
  \caption{~$\LogLuk[3]$ وچ !!{derivation} دی مثال}
\end{sidewaysfigure}
\endgroup
% /OLP-0406

% OLP-0652 content/many-valued-logic/sequent-calculus/proof-theoretic-notions.tex
\begingroup
\def\olfilename{proof-theoretic-notions}
\expandafter\def\csname ol@sectioncs\endcsname{section}

\olfileid{mvl}{seq}{prf}

\olsection{ثبوتی نظریے دے تصور}

\begin{explain}
جیویں اسیں کئی اہم معنوی تصورات (منطقی اعتبار، منطقی نتیجہ تے پورا ہون جوگتا) دی تعریف دِتّی اے، ہُن اسیں اوہناں نال میل کھاندے \emph{ثبوتی نظریے دے تصوراں} دی تعریف دیندے آں۔ ایہ تعریف !!{sentence}s دے !!{structure}s وچ پورا اترن اُتے نہیں، کجھ سیکوئنٹاں دے !!{derivability} یا !!{nonderivability} اُتے ٹِکی اے۔ ایہ وڈی دریافت سی کہ ایہ تصور آپس وچ میل کھاندے نیں۔ اوہناں دا ایہ میل \emph{صحت مندی} تے \emph{تمامیت دے قضیے} دا مضمون اے۔
\end{explain}

\begin{defn}[قضیے]
!!^a{sentence}~$!A$ اک \emph{قضیہ} اے جے~$\Log{LK}$ وچ سیکوئنٹ $\quad \Sequent !A$ دی !!a{derivation} ہووے۔ اسیں $\Proves
!A$ لکھدے آں جے $!A$ قضیہ اے، تے $\Proves/ !A$ جے ایہ قضیہ نہیں۔
\end{defn}

\begin{defn}[!!^{derivability}]
!!^a{sentence}~$!A$، !!{sentence}s دے سیٹ~$\Gamma$ توں \emph{!!{derivable}} اے، $\Gamma \Proves !A$، اُدوں تے صرف اُدوں جدوں اک محدود ذیلی سیٹ~$\Gamma_0 \subseteq \Gamma$ تے لڑی $\Gamma_0'$ ہووے جیہڑی~$\Gamma_0$ وچلے !!{sentence}s دی لڑی ہووے، ایویں کہ $\Log{LK}$، $\Gamma_0' \Sequent !A$ نوں !!{derive}s کرے۔ جے $!A$، $\Gamma$ توں !!{derivable} نہیں تاں اسیں $\Gamma \Proves/ !A$ لکھدے آں۔
\end{defn}

سکیڑن، کمزور کرن تے تبادلے دے قاعدیاں کرکے~$\Gamma_0'$ وچ !!{sentence}s دی ترتیب تے گنتی نال فرق نہیں پیندا۔ جے سیکوئنٹ $\Gamma_0' \Sequent !A$ !!{derivable} اے، تاں $\Gamma_0'' \Sequent !A$ وی اے، ہر اوہ $\Gamma_0''$ لئی جیہدے وچ اوہی !!{sentence}s نیں جیہڑے~$\Gamma_0'$ وچ نیں۔ مثلاً جے $\Gamma_0 = \{!B, !C\}$ ہووے، تاں $\Gamma_0' = \tuple{!B, !B, !C}$ تے $\Gamma_0'' =
\tuple{!C, !C, !B}$ دوویں ایہو جیہیاں لڑیاں نیں جیہناں وچ بس اوہی !!{sentence}s نیں جیہڑے~$\Gamma_0$ وچ نیں۔ جے اِنہاں وچوں اک والا سیکوئنٹ !!{derivable} اے، تاں دوجا وی اے، مثال لئی:
\begin{prooftree}
  \AxiomC{}
  \Deduce$!B, !B, !C \fCenter !A$
  \RightLabel{\LeftR{\Contraction}}
  \UnaryInf$!B, !C \fCenter !A$
  \RightLabel{\LeftR{\Exchange}}
  \UnaryInf$!C, !B \fCenter !A$
  \RightLabel{\LeftR{\Weakening}}
  \UnaryInf$!C, !C, !B \fCenter !A$
\end{prooftree}
ہُن توں اسیں آکھاں گے کہ جے $\Gamma_0$، !!{sentence}s دا محدود سیٹ اے، تاں $\Gamma_0 \Sequent !A$ کوئی وی ایہو جیہا سیکوئنٹ اے جیہدے مقدم وچ~$\Gamma_0$ دے !!{sentence}s دی لڑی ہووے۔ جتھے لوڑ ہووے، سکیڑن، تبادلے تے کمزور کرن دے قدم بناں دسے شامل سمجھاں گے۔

\begin{defn}[بے تضادی]
جملیاں دا سیٹ~$\Gamma$ اُدوں تے صرف اُدوں \emph{متضاد} اے جدوں اوہدا کوئی محدود ذیلی سیٹ~$\Gamma_0 \subseteq \Gamma$ ہووے ایویں کہ $\Log{LK}$ !!{derive}s کرے سیکوئنٹ $\Gamma_0 \Sequent \quad$۔ جے $\Gamma$ متضاد نہیں، یعنی جے ہر محدود $\Gamma_0 \subseteq \Gamma$ لئی $\Log{LK}$ !!{derive} نہیں کردا سیکوئنٹ $\Gamma_0 \Sequent \quad$، تاں اسیں اوہنوں \emph{بے تضاد} آکھدے آں۔
\end{defn}

\begin{prop}[انعکاس]
\ollabel{prop:reflexivity}
جے $!A \in \Gamma$، تاں $\Gamma \Proves !A$۔
\end{prop}

\begin{proof}
مُڈھلا سیکوئنٹ $!A \Sequent !A$ !!{derivable} اے، تے $\{!A\}
\subseteq \Gamma$۔
\end{proof}

\begin{prop}[وڈے سیٹ ول برقرار رہنا]
\ollabel{prop:monotonicity}
جے $\Gamma \subseteq \Delta$ تے $\Gamma \Proves !A$، تاں $\Delta
\Proves !A$۔
\end{prop}

\begin{proof}
منّ لو $\Gamma \Proves !A$، یعنی کوئی محدود $\Gamma_0
\subseteq \Gamma$ اے جیہدے لئی $\Gamma_0 \Sequent !A$ !!{derivable} اے۔ چونکہ $\Gamma \subseteq \Delta$، ایس لئی $\Gamma_0$،~$\Delta$ دا وی محدود ذیلی سیٹ اے۔ سو $\Gamma_0
\Sequent !A$ دا !!{derivation} ایہ وی وکھاؤندا اے کہ $\Delta \Proves !A$۔
\end{proof}

\begin{prop}[تعدی]
\ollabel{prop:transitivity}
جے $\Gamma \Proves !A$ تے $\{!A\} \cup \Delta \Proves
!B$، تاں $\Gamma \cup \Delta \Proves !B$۔
\end{prop}

\begin{proof}
جے $\Gamma \Proves !A$، تاں اک محدود $\Gamma_0 \subseteq \Gamma$ تے اک !!a{derivation}~$\pi_0$ اے جیہڑا $\Gamma_0 \Sequent !A$ دا اے۔ جے $\{!A\}
\cup \Delta \Proves !B$، تاں کسے محدود ذیلی سیٹ $\Delta_0
\subseteq \Delta$ لئی اک !!a{derivation}~$\pi_1$ اے جیہڑا $!A, \Delta_0
\Sequent !B$ دا اے۔ ایہ !!{derivation} ویکھو:
\begin{prooftree}
  \AxiomC{}
  \RightLabel{$\pi_0$}
  \Deduce$\Gamma_0 \fCenter !A$
  \AxiomC{}
  \RightLabel{$\pi_1$}
  \Deduce$!A, \Delta_0 \fCenter !B$
  \RightLabel{\Cut}
  \BinaryInf$\Gamma_0, \Delta_0 \fCenter !B$
\end{prooftree}
چونکہ $\Gamma_0 \cup \Delta_0 \subseteq \Gamma \cup \Delta$، ایس لئی ایہ وکھاؤندا اے کہ $\Gamma \cup \Delta \Proves !B$۔
\end{proof}

دھیان دیو، ایس دا خاص مطلب ایہ اے کہ جے $\Gamma \Proves !A$ تے $!A
\Proves !B$، تاں $\Gamma \Proves !B$۔ ایس توں ایہ وی نکلدا اے کہ جے $!A_1,
\dots, !A_n \Proves !B$ تے ہر واری $\Gamma \Proves !A_i$ ہووے، ہر~$i$ لئی، تاں $\Gamma \Proves !B$۔

\begin{prop}
\ollabel{prop:incons}
$\Gamma$ اُدوں تے صرف اُدوں متضاد اے جدوں $\Gamma \Proves {!A}$ ہر جملے~$!A$ لئی۔
\end{prop}

\begin{proof}
مشق۔
\end{proof}

\begin{prob}
\olref[fol][seq][ptn]{prop:incons} ثابت کرو۔
\end{prob}

\begin{prop}[محدود گواہی]
  \ollabel{prop:proves-compact}
  \begin{enumerate}
  \item جے $\Gamma \Proves !A$، تاں اک محدود ذیلی سیٹ $\Gamma_0
    \subseteq \Gamma$ اے جیہدے لئی $\Gamma_0 \Proves !A$۔
  \item جے~$\Gamma$ دا ہر محدود ذیلی سیٹ بے تضاد اے، تاں $\Gamma$ بے تضاد اے۔
  \end{enumerate}
\end{prop}

\begin{proof}
  \begin{enumerate}
    \item جے $\Gamma \Proves !A$، تاں اک محدود ذیلی سیٹ $\Gamma_0 \subseteq \Gamma$ اے جیہدے سیکوئنٹ $\Gamma_0
      \Sequent !A$ دا !!a{derivation} اے۔ ایس لئی $\Gamma_0
      \Proves !A$۔
    \item جے $\Gamma$ متضاد اے، تاں اک محدود ذیلی سیٹ~$\Gamma_0 \subseteq \Gamma$ اے جیہدے لئی $\Log{LK}$، $\Gamma_0 \Sequent \quad$ نوں !!{derive}s کرے۔ پر ایہدا مطلب اے کہ اک محدود ذیلی سیٹ~$\Gamma_0$، جیہڑا $\Gamma$ دا اے، متضاد اے۔
  \end{enumerate}
\end{proof}
\endgroup
% /OLP-0652
\OLEndChapterHook
\endgroup
% /OLP-0402

\OLEndPartHook
\endgroup
% /OLP-0383

% OLP-0407 content/normal-modal-logic/normal-modal-logic.tex
\begingroup
\def\olfilename{normal-modal-logic}
\expandafter\def\csname ol@sectioncs\endcsname{section}

\olpart{nml}{نارمل امکانی منطقاں}

\begin{editorial}
  ایہ حصہ نارمل امکانی منطقاں دے نظریاتی مطالعے بارے اے۔
  ایس ویلے ایس وچ بنیادی امکانی منطق دی کلاسیکی مطابقت دے
  نظریے بارے آلڈو انٹونیلی دے نوٹ نیں۔
\end{editorial}

% OLP-0408 content/normal-modal-logic/syntax-and-semantics/syntax-and-semantics.tex
\begingroup
\def\olfilename{syntax-and-semantics}
\expandafter\def\csname ol@sectioncs\endcsname{section}

\olchapter{nml}{syn}{نحو تے معنیات}

% OLP-0409 content/normal-modal-logic/syntax-and-semantics/introduction.tex
\begingroup
\def\olfilename{introduction}
\expandafter\def\csname ol@sectioncs\endcsname{section}

\olfileid{nml}{syn}{int}

\olsection{جان پہچان}

امکانی منطق \emph{امکان تے ضرورت ورگی جہت والے قضیاں} تے
اوہناں دے آپس دے منطقی نتیجیاں بارے اے۔ ایہو جیہے قضیاں دیاں
کجھ مثالاں ہِٹھ نیں۔
\begin{enumerate}
\item ضروری اے کہ $2+2=4$۔
\item ایہ لازماً ممکن اے کہ بھلکے مینہ پوے۔
\item جے~$!A$ دا ہونا لازماً ممکن اے، تاں~$!A$ دا ہونا ممکن اے۔
\end{enumerate}
امکان تے ضرورت توں علاوہ ہور جہتاں وی نیں۔ ہور اک-دلیلی
رابطے وی جہتی رابطے منّے جاندے نیں، جیویں «~$!A$ ہونا چاہیدا
اے»، «~$!A$ ہووے گا»، «ڈینا جاندا اے کہ~$!A$»، یا
«ڈینا منّدا اے کہ~$!A$»۔

امکانی منطق پہلاں ارسطو دی \emph{De
  Interpretatione} وچ سامنے آؤندی اے۔ اوہنے سبھ توں پہلاں
دھیان دِتّا کہ ضرورت توں امکان نکلدا اے، پر اُلٹی گل نہیں۔
امکان تے ضرورت نوں اک دوجے دے ذریعے متعین کیتا جا سکدا اے۔
جے $!A \land !B$ دا سچ ہونا ممکن اے، تاں $!A$ دا سچ ہونا
ممکن اے تے $!B$ دا سچ ہونا وی ممکن اے، پر اُلٹی گل نہیں۔
تے جے $!A \to !B$ ضروری اے، تاں جے $!A$ ضروری اے،
~$!B$ وی ضروری اے۔

امکانی منطق دا پہلا جدید طریقہ سی۔ آئی۔ لوئس دے کم توں
آیا، جیہڑا لوئس تے لینگفورڈ دی کتاب \emph{Symbolic Logic}
(1932) تک پُہنچیا۔ لوئس تے لینگفورڈ، لزوم نوں مادی شرطیے
نال ظاہر کرن توں مطمئن نہیں سن۔ $!A \lif !B$، «$!A$ توں
$!B$ لازم اے» دی تھاں چنگی طرح کم نہیں کردا۔ ایس لئی اوہناں
لزوم نوں «ضروری اے کہ جے $!A$، تاں $!B$» وجوں متعین
کرن دی تجویز دِتّی، تے ایہنوں $!A \strictif !B$ لکھیا۔
وکھ وکھ خاصیتاں سمجھدیاں لوئس نے پنج امکانی نظام پچھانے،
\Log{S1}، \ldots، \Log{S4}، \Log{S5}۔ آخری دو اجے وی
ورتے جاندے نیں۔

لوئس تے لینگفورڈ دا طریقہ سراسر \emph{نحوی} سی۔ اوہناں
موزوں مسلّمے تے قاعدے چُنے تے پرکھیا کہ اوہناں نال کیہ
ثابت ہو سکدا اے۔ معنیات دا طریقہ لمّے چِر تک ہتھ نہ آیا۔ پہلی
کوشش روڈولف کارناپ نے \emph{Meaning and Necessity} (1947)
وچ \emph{حالت دے بیان} دا خیال ورت کے کیتی۔ ایہ ایٹمی
جملیاں دا اکٹھ اے، یعنی اوہ جملے جیہڑے اوس حالت دے بیان
وچ «سچ» نیں۔ سچ ہون دی تعریف ہر طرح دے جملے $!A$ تک
ودھا کے، کارناپ $!A$ نوں \emph{لازماً سچ} منّدا اے جے اوہ
حالتاں دے سارے بیاناں وچ سچ ہووے۔ کارناپ دا طریقہ \emph{وار
وار لگیاں} جہتاں نوں نہیں سنبھال سکدا سی۔ «ممکن اے کہ
ضروری اے \ldots ممکن اے کہ $!A$» ورگے جملے ہمیشہ
سبھ توں اندرلی جہت تک گھٹ جاندے نیں۔

امکانی معنیات وچ وڈی پیش رفت ساؤل کرپکے دے مضمون
«A Completeness Theorem in Modal Logic» (JSL 1959) نال
ہوئی۔ کرپکے دا کم لائبنز دے ایس خیال اُتے ٹکیا سی کہ جملہ
لازماً سچ اے جے اوہ «سارے ممکن جہاناں وچ» سچ اے۔ پر ایس
خیال وچ وی کارناپ والی دِقّت اے۔ کسے جہان $w$ یا حالت دے
بیان $s$ وچ جملے دا سچ ہونا $w$ اُتے بالکل منحصر نہیں۔
سو کرپکے نے منّیا کہ جہان \emph{رسائی دے تعلق} $R$ نال
جُڑے نیں۔ «ضروری اے کہ $!A$» دی شکل والا جملہ جہان $w$
وچ اُدوں تے صرف اُدوں سچ اے جدوں $!A$ سارے اوہناں جہاناں
$w'$ وچ سچ اے جنہاں تک \emph{رسائی} $w$ توں اے۔
ایس طریقے دا کوئی روپ دین والی معنیات نوں کرپکے معنیات
آکھدے آں۔ ایہناں نال کئی امکانی منطقاں تیزی نال اگّے ودھیاں۔

کرپکے معنیات مطابق سمجھیے، تاں امکانی منطق سانوں دسدِی
اے کہ \emph{تعلقی ساختاں} «اندروں» کیویں دِسدیاں نیں۔
تعلقی ساخت بس اوہ سیٹ اے جیہدے اُتے دو رکنی تعلق اے۔
مثلاً جماعت دے طالب علماں دا سیٹ، جیہڑا اوہناں دے سماجی
تحفظ دے نمبراں مطابق ترتیب دِتّا ہووے، اک تعلقی ساخت اے۔
تعلقی ساختاں کئی طرح دے میدان وچ ہوندیاں نیں۔ دنیا دیاں
حالتاں دے اک دوجی دے لحاظ نال ممکن ہون توں علاوہ، کسے
کارندے دیاں علمی حالتاں علمی امکان نال جُڑ سکدیاں نیں۔
ایسے طرح بدل دے نظام دیاں حالتاں اوہناں دیاں حالتی
تبدیلیاں نال جُڑ سکدیاں نیں۔ امکانی منطق نال ایہناں
ساریاں دا نمونہ بنا سکدے آں۔ پہلی صورت توں صدق نال جُڑی
عام امکانی منطق ملدی اے۔ ہور صورتاں توں علمی منطق، عمل تے
تبدیلی دی منطق وغیرہ ملدیاں نیں۔

ساڈی توجہ اک خاص پاسے اُتے اے جیہنوں امکانی منطق دے
ماہر «مطابقت دا نظریہ» آکھدے نیں۔ کرپکے دیاں اہم ابتدائی
لبھتاں وچوں اک ایہ اے کہ رسائی دے تعلق~$R$ دیاں کئی
خاصیتاں، جیویں ایہدا تعدّی والا یا متقابل ہونا، مناسب
«امکانی سانچیاں» نال \emph{امکانی زبان دے اندر ای} پچھانیاں
جا سکدیاں نیں۔ مثلاً امکانی منطق دے ماہر آکھدے نیں کہ
$R$ دا انعکاسی ہونا سانچے «جے لازماً~$!A$، تاں~$!A$»
نال «مطابق» اے۔ اسیں زیادہ تر امکانی منطق دے کجھ کلاسیکی
نظاماں، مثلاً \Log{S4} تے \Log{S5}، دا مطابقتی نظریہ
ویکھاں گے۔ ایہ نظام سانچیاں \Ax{D}، \Ax{T}، \Ax{B}،
\Ax{4} تے~\Ax{5} دیاں ترکیباں توں ملدے نیں۔
\endgroup
% /OLP-0409
% OLP-0410 content/normal-modal-logic/syntax-and-semantics/language-modal-logic.tex
\begingroup
\def\olfilename{language-modal-logic}
\expandafter\def\csname ol@sectioncs\endcsname{section}

\olfileid{nml}{syn}{lan}

\olsection{بنیادی امکانی منطق دی زبان}

\begin{defn}
بنیادی امکانی منطق دی زبان وچ ایہ نشان ہوندے نیں۔
\begin{enumerate}
  \tagitem{prvFalse}{!!{falsity} لئی قضیاتی مستقل~$\lfalse$۔}{}
  \tagitem{prvTrue}{!!{truth} لئی قضیاتی مستقل~$\ltrue$۔}{}
  \item !!{propositional variable}s دا !!{denumerable}s سیٹ،
    یعنی $\Obj
    p_0$، $\Obj p_1$، $\Obj p_2$، \dots۔
  \item قضیاتی رابطے۔ \startycommalist
  \iftag{prvNot}{\ycomma $\lnot$ (نفی)}{}%
  \iftag{prvAnd}{\ycomma $\land$ (عطف)}{}%
  \iftag{prvOr}{\ycomma $\lor$ (فصل)}{}%
  \iftag{prvIf}{\ycomma $\lif$ (!!{conditional})}{}%
  \iftag{prvIff}{\ycomma $\liff$ (!!{biconditional})}{}۔
  \tagitem{prvBox}{امکانی عامل $\Box$۔}{}
  \tagitem{prvDiamond}{امکانی عامل $\Diamond$۔}{}
\end{enumerate}
\end{defn}

\begin{defn}
بنیادی امکانی زبان دے \emph{!!^{formula}s} استقرا نال ایویں متعین ہوندے نیں۔
\begin{enumerate}
\tagitem{prvFalse}{$\lfalse$ اک ایٹمی !!{formula} اے۔}{}

\tagitem{prvTrue}{$\ltrue$ اک ایٹمی !!{formula} اے۔}{}

\item ہر !!{propositional variable} $\Obj p_i$ اک (ایٹمی) !!{formula} اے۔

\tagitem{prvNot}{جے $!A$ !!a{formula} اے، تاں $\lnot !A$ وی
  !!a{formula} اے۔}{}

\tagitem{prvAnd}{جے $!A$ تے $!B$، !!{formula}s نیں، تاں $(!A \land
  !B)$ وی !!a{formula} اے۔}{}

\tagitem{prvOr}{جے $!A$ تے $!B$، !!{formula}s نیں، تاں $(!A \lor !B)$
  وی !!a{formula} اے۔}{}

\tagitem{prvIf}{جے $!A$ تے $!B$، !!{formula}s نیں، تاں $(!A \lif !B)$
  وی !!a{formula} اے۔}{}

\tagitem{prvIff}{جے $!A$ تے $!B$، !!{formula}s نیں، تاں $(!A \liff !B)$
  وی !!a{formula} اے۔}{}

\tagitem{prvBox}{جے $!A$ !!a{formula} اے، تاں $\Box !A$ وی
  !!a{formula} اے۔}{}

\tagitem{prvDiamond}{جے $!A$ !!a{formula} اے، تاں $\Diamond !A$ وی
  !!a{formula} اے۔}{}

\tagitem{limitClause}{ہور کوئی شے !!a{formula} نہیں۔}{}
\end{enumerate}
\end{defn}

\begin{tagblock}{defNot,defOr,defAnd,defIf,defIff,defTrue,defFalse,defBox,defDiamond}
\begin{defn}
متعین کیتے عاملاں نال بنے فارمولے ایویں سمجھے جاندے نیں۔

\begin{tagenumerate}{defTrue,defFalse,defNot,defOr,defAnd,defIf,defIff,defBox,defDiamond}

\tagitem{defTrue}{$\ltrue$ ایہدا مختصر روپ اے۔
  \iftag{prvFalse}{$\lnot\lfalse$}{$(!A \lor \lnot !A)$، کسے
    مقرر ایٹمی !!{formula}~$!A$ لئی}۔}{}

\tagitem{defFalse}{$\lfalse$ ایہدا مختصر روپ اے۔
  \iftag{prvTrue}{$\lnot\ltrue$}{$(!A \land \lnot !A)$، کسے
    مقرر ایٹمی !!{formula}~$!A$ لئی}۔}{}

\tagitem{defNot}{$\lnot !A$، $!A \lif \lfalse$ دا مختصر روپ اے۔}{}

\tagitem{defOr}{$!A \lor !B$ ایہدا مختصر روپ اے۔
  \iftag{prvAnd}{$\lnot(\lnot !A \land \lnot !B)$}{$\lnot !A \lif
    !B$}۔}{}

\tagitem{defAnd}{$!A \land !B$ ایہدا مختصر روپ اے۔
  \iftag{prvOr}{$\lnot(\lnot !A \lor \lnot !B)$}{$\lnot (!A \lif
    \lnot !B)$}۔}{}

\tagitem{defIf}{$!A \lif !B$ ایہدا مختصر روپ اے۔
  \iftag{prvOr}{$\lnot !A \lor !B)$}{$\lnot (!A \land \lnot !B)$}۔}{}

\tagitem{defIff}{$!A \liff !B$، $(!A \lif !B) \land (!B
  \lif !A)$ دا مختصر روپ اے۔}{}

\tagitem{defBox}{$\Box !A$، $\lnot\Diamond\lnot !A$ دا مختصر روپ اے۔}{}

\tagitem{defDiamond}{$\Diamond !A$، $\lnot\Box\lnot !A$ دا مختصر روپ اے۔}{}
\end{tagenumerate}
\end{defn}
\end{tagblock}

جے !!a{formula}~$!A$ وچ $\Box$ یا $\Diamond$ نہ آوے،
تاں اوہنوں \emph{امکانی عامل توں خالی} آکھدے آں۔
\endgroup
% /OLP-0410
% OLP-0411 content/normal-modal-logic/syntax-and-semantics/substitution.tex
\begingroup
\def\olfilename{substitution}
\expandafter\def\csname ol@sectioncs\endcsname{section}

\olfileid{nml}{syn}{sub}

\olsection{اکّو ویلے بدل}

!!a{formula}~$!A$ دی \emph{بدل توں بنی مثال} اوہ نتیجہ اے جیہڑا
کسے !!a{propositional variable} دی~$!A$ وچ ہر آمد نوں
کسے ہور !!{formula} نال بدل کے ملدا اے۔ اسیں منطقی اعتبار دی گل
کردیاں وی تے !!{derivability} دی گل کردیاں وی اکثر
!!{formula}s دیاں مثالاں دا حوالہ دیاں گے۔ ایس لئی ایس خیال
دی ٹھیک تعریف کرنا فائدے مند اے۔

\begin{defn}\ollabel{def:subst-inst}
  جے $!A$ اک امکانی !!{formula} اے جیہدے سارے !!{propositional variable}s، $p_1$، \dots، $p_n$ وچوں نیں، تے $!D_1$، \dots،
  $!D_n$ وی امکانی !!{formula}s نیں، تاں
  $\SSubst{!A}{\subst{!D_1}{p_1}, \dots, \subst{!D_n}{p_n}}$
  اوہ نتیجہ اے جیہڑا ہر $!D_i$ نوں $p_i$ دی تھاں~$!A$ وچ
  اکّو ویلے رکھن نال ملدا اے۔ باقاعدہ طور اُتے ایہ~$!A$ اُتے
  استقرا والی تعریف اے۔
  \begin{enumerate}
    \tagitem{prvFalse}{\indcase{!A}{\lfalse}{%
        $\SSubst{\indfrm}{\subst{!D_1}{p_1}, \dots,
          \subst{!D_n}{p_n}}$ دا نتیجہ $\lfalse$ اے۔}}{}
    \tagitem{prvTrue}{\indcase{!A}{\ltrue}{%
        $\SSubst{\indfrm}{\subst{!D_1}{p_1}, \dots,
          \subst{!D_n}{p_n}}$ دا نتیجہ $\ltrue$ اے۔}}{}
    \item \indcase{!A}{q}{$\SSubst{\indfrm}{\subst{!D_1}{p_1}, \dots,
        \subst{!D_n}{p_n}}$ دا نتیجہ $q$ اے، بشرطیکہ $q \not\ident p_i$
      ہووے $i
      = 1$، \dots،~$n$ لئی۔}
    \item \indcase{!A}{p_i}{$\SSubst{\indfrm}{\subst{!D_1}{p_1},
        \dots, \subst{!D_n}{p_n}}$ دا نتیجہ $!D_i$ اے۔}
    \tagitem{prvNot}{\indcase{!A}{\lnot !B}{%
        $\SSubst{\indfrm}{\subst{!D_1}{p_1}, \dots, \subst{!D_n}{p_n}}$ دا نتیجہ
        $\lnot \SSubst{!B}{\subst{!D_1}{p_1}, \dots, \subst{!D_n}{p_n}}$ اے۔}}{}
    \tagitem{prvAnd}{\indcase{!A}{(!B \land
        !C)}{$\SSubst{\indfrm}{\subst{!D_1}{p_1}, \dots,
          \subst{!D_n}{p_n}}$ دا نتیجہ ایہ اے۔
        \[(\SSubst{!B}{\subst{!D_1}{p_1}, \dots, \subst{!D_n}{p_n}} \land
          \SSubst{!C}{\subst{!D_1}{p_1}, \dots, \subst{!D_n}{p_n}}).\]}}{}
    \tagitem{prvOr}{\indcase{!A}{(!B \lor
          !C)}{$\SSubst{\indfrm}{\subst{!D_1}{p_1}, \dots,
          \subst{!D_n}{p_n}}$ دا نتیجہ ایہ اے۔
        \[(\SSubst{!B}{\subst{!D_1}{p_1}, \dots, \subst{!D_n}{p_n}} \lor
          \SSubst{!C}{\subst{!D_1}{p_1}, \dots, \subst{!D_n}{p_n}}).\]}}{}
    \tagitem{prvIf}{\indcase{!A}{(!B \lif
          !C)}{$\SSubst{\indfrm}{\subst{!D_1}{p_1}, \dots,
          \subst{!D_n}{p_n}}$ دا نتیجہ ایہ اے۔
        \[(\SSubst{!B}{\subst{!D_1}{p_1}, \dots, \subst{!D_n}{p_n}} \lif
          \SSubst{!C}{\subst{!D_1}{p_1}, \dots, \subst{!D_n}{p_n}}).\]}}{}
    \tagitem{prvIf}{\indcase{!A}{(!B \liff
          !C)}{$\SSubst{\indfrm}{\subst{!D_1}{p_1}, \dots,
          \subst{!D_n}{p_n}}$ دا نتیجہ ایہ اے۔
        \[(\SSubst{!B}{\subst{!D_1}{p_1}, \dots, \subst{!D_n}{p_n}} \liff
          \SSubst{!C}{\subst{!D_1}{p_1}, \dots, \subst{!D_n}{p_n}}).\]}}{}
    \item \indcase{!A}{\Box !B}{%
        $\SSubst{\indfrm}{\subst{!D_1}{p_1}, \dots, \subst{!D_n}{p_n}}$ دا نتیجہ
        $\Box
        \SSubst{!B}{\subst{!D_1}{p_1}, \dots, \subst{!D_n}{p_n}}.$}
    \tagitem{prvDiamond}{\indcase{!A}{\Diamond !B}{%
        $\SSubst{\indfrm}{\subst{!D_1}{p_1}, \dots, \subst{!D_n}{p_n}}$ دا نتیجہ
        $\Diamond
        \SSubst{!B}{\subst{!D_1}{p_1}, \dots, \subst{!D_n}{p_n}}.$}}{}
  \end{enumerate}
  !!{formula} $\SSubst{!A}{\subst{!D_1}{p_1}, \dots,
    \subst{!D_n}{p_n}}$ نوں~$!A$ دی \emph{بدل توں بنی مثال}
  آکھدے آں۔
\end{defn}

\begin{ex}
  منّ لو $!A$، $p_1 \lif \Box(p_1 \land p_2)$ اے، $!D_1$،
  $\Diamond(p_2 \lif p_3)$ اے، تے $!D_2$، $\lnot\Box p_1$ اے۔
  فیر $\SSubst{!A}{\subst{!D_1}{p_1}, \subst{!D_2}{p_2}}$ دا
  نتیجہ ایہ اے۔
  \begin{align*}
    \Diamond(p_2 \lif p_3) &
    \lif \Box(\Diamond(p_2 \lif p_3) \land \lnot\Box p_1)
    \intertext{جدکہ $\SSubst{!A}{\subst{!D_2}{p_1}, \subst{!D_1}{p_2}}$ دا نتیجہ ایہ اے۔}
    \lnot\Box p_1 &
    \lif \Box(\lnot\Box p_1 \land \Diamond(p_2 \lif p_3))
    \intertext{دھیان دیو کہ اکّو ویلے بدل عام طور اُتے وار وار بدل ورگا نہیں۔
      مثلاً اُپرلے $\SSubst{!A}{\subst{!D_1}{p_1}, \subst{!D_2}{p_2}}$ نوں
      $\Subst{(\Subst{!A}{!D_1}{p_1})}{!D_2}{p_2}$ نال ملاؤ، جیہدا نتیجہ ایہ اے۔}
    \Diamond(p_2 \lif p_3) & \Subst{\lif \Box(\Diamond(p_2 \lif p_3) \land p_2)}{\lnot\Box p_1}{p_2}, \text{ یعنی،}\\
    \Diamond(\lnot\Box p_1 \lif p_3) &
    \lif \Box(\Diamond(\lnot\Box p_1
    \lif p_3) \land \lnot\Box p_1)\\
    \intertext{تے $\Subst{(\Subst{!A}{!D_2}{p_2})}{!D_1}{p_1}$ نال ملاؤ۔}
    p_1 & \lif \Subst{\Box(p_1 \land \lnot\Box p_1)}{\Diamond(p_2 \lif p_3)}{p_1}, \text{ یعنی،}\\
    \Diamond(p_2 \lif p_3) &
    \lif \Box(\Diamond(p_2
    \lif p_3) \land \lnot\Box\Diamond(p_2 \lif p_3)).
  \end{align*}
\end{ex}
\endgroup
% /OLP-0411
% OLP-0412 content/normal-modal-logic/syntax-and-semantics/relational-models.tex
\begingroup
\def\olfilename{relational-models}
\expandafter\def\csname ol@sectioncs\endcsname{section}

\olfileid{nml}{syn}{rel}

\olsection{تعلقی نمونے}

نارمل امکانی منطقاں دی معنیات دا بنیادی خیال \emph{تعلقی نمونہ}
اے۔ ایہ جہاناں دے اک سیٹ توں بندا اے۔ جہان دو رکنی «رسائی دے
تعلق» نال جُڑے نیں، تے اک قدر بندی دسدی اے کہ کیہڑے
!!{propositional variable}s کیہڑے جہاناں وچ «سچ» منّے جاندے نیں۔

\begin{defn}
  بنیادی امکانی زبان دا \emph{نمونہ} اک ترتیب وار تِکڑی
  $\mModel{M}
  = \tuple{W, R, V}$ اے، جتھے
  \begin{enumerate}
  \item $W$ «جہاناں» دا غیر خالی سیٹ اے،
  \item $R$، ~$W$ اُتے دو رکنی رسائی دا تعلق اے، تے
  \item $V$ اوہ فنکشن اے جیہڑا ہر !!{propositional variable}~$p$ نوں ممکن جہاناں دا سیٹ $V(p)$ دیندا اے۔
  \end{enumerate}
  جدوں $Rww'$ قائم ہووے، اسیں آکھدے آں کہ $w'$ تک
  \emph{رسائی}~$w$ توں اے۔ جدوں $w \in V(p)$، اسیں آکھدے آں
  کہ $p$،~$w$ \emph{وچ سچ} اے۔
\end{defn}

تعلقی معنیات دا وڈا فائدہ ایہ اے کہ نمونیاں نوں سادے خاکیاں
نال ظاہر کر سکدے آں، جیویں \olref{fig:simple} وچ۔ جہان گرہاں
نال ظاہر ہوندے نیں۔ جہان~$w'$ تک~$w$ توں رسائی اُدوں تے صرف
اُدوں اے جدوں~$w$ توں~$w'$ ول تیر ہووے۔ نالے اسیں گرہ،
یعنی جہان، اُتے~$\mTrue{p}$ لکھدے آں جدوں $w \in V(p)$،
تے ہور صورت وچ~\mFalse{p}۔ \olref{fig:simple} اوہ نمونہ
دکھاؤندا اے جیہدے وچ $W = \{w_1, w_2, w_3\}$،
$R = \{\tuple{w_1, w_2},\tuple{w_1,
  w_3}\}$، $V(p) = \{w_1, w_2\}$ تے $V(q) = \{w_2\}$۔

\begin{figure}
  \begin{center}
    \begin{tikzpicture}[modal]
      \node[world] (w1) [label={[align=right]right:\mTrue{p}\\ \mFalse{q}}]{$w_1$}; 
      \node[world] (w2) [label={[align=right]right:\mTrue{p}\\ \mTrue{q}},
        above right=of w1]{$w_2$}; 
      \node[world] (w3) [label={[align=right]right:\mFalse{p}\\ \mFalse{q}},
        below right=of w1] {$w_3$};
      \draw[->] (w1) to (w2);
      \draw[->] (w1) to (w3);
    \end{tikzpicture}
  \end{center}
  \caption{اک سادا نمونہ۔}
  \ollabel{fig:simple}
\end{figure}
\endgroup
% /OLP-0412
% OLP-0413 content/normal-modal-logic/syntax-and-semantics/truth-at-w.tex
\begingroup
\def\olfilename{truth-at-w}
\expandafter\def\csname ol@sectioncs\endcsname{section}

\olfileid{nml}{syn}{trw}

\olsection{کسی جہان وچ سچائی}

ہر امکانی نمونہ طے کردا اے کہ اوہدے کہڑے جہاناں وچ کہڑے امکانی !!{formula}s سچے نیں۔ «نمونہ $\mModel{M}$، !!{formula}~$!A$ نوں جہان~$w$ وچ سچا کردا اے» والا تعلق تعلقی معنیات دا بنیادی خیال اے۔ ایس تعلق دی تعریف استقرا نال ہوندی اے تے غیر امکانی عاملاں لئی ایہ سچائی دیاں سارنیاں والی عام تعریف نال میل کھاندا اے۔

\begin{defn}\ollabel{defn:mmodels}
  \emph{!!a{formula}~$!A$ دا جہان~$w$ وچ سچا ہونا}، نمونے~$\mModel M$ دے لحاظ نال، علامتاں وچ $\mSat{M}{!A}[w]$، ایس طرح استقرا نال متعین ہوندا اے:
  \begin{enumerate}
  \tagitem{prvFalse}{\indcase{!A}{\lfalse}{کدے وی $\mSat{M}{\lfalse}[w]$ نہیں ہوندا}۔}{}
  \tagitem{prvTrue}{\indcase{!A}{\ltrue}{ہمیشہ $\mSat{M}{\ltrue}[w]$ ہوندا اے}۔}{}
  \item $\mSat{M}{p}[w]$ اُدوں تے صرف اُدوں جدوں $w \in V(p)$۔
  \tagitem{prvNot}{\indcase{!A}{\lnot !B}{$\mSat{M}{\indfrm}[w]$ اُدوں تے صرف اُدوں جدوں
    $\mSat/{M}{!B}[w]$}۔}{}
  \tagitem{prvAnd}{\indcase{!A}{(!B \land !C)}{$\mSat{M}{\indfrm}[w]$ اُدوں تے صرف اُدوں جدوں
    $\mSat{M}{!B}[w]$ تے $\mSat{M}{!C}[w]$}۔}{}
  \tagitem{prvOr}{\indcase{!A}{(!B \lor !C)}{$\mSat{M}{\indfrm}[w]$ اُدوں تے صرف اُدوں جدوں
    $\mSat{M}{!B}[w]$ یا $\mSat{M}{!C}[w]$} (یا دونیں)۔}{}
  \tagitem{prvIf}{\indcase{!A}{(!B \lif !C)}{$\mSat{M}{\indfrm}[w]$ اُدوں تے صرف اُدوں جدوں
    $\mSat/{M}{!B}[w]$ یا $\mSat{M}{!C}[w]$}۔}{}
  \tagitem{prvIff}{\indcase{!A}{(!B \liff !C)}{$\mSat{M}{\indfrm}[w]$ اُدوں تے صرف اُدوں جدوں
    یا تاں $\mSat{M}{!B}[w]$ تے $\mSat{M}{!C}[w]$ دونیں ہون، یا
    نہ $\mSat{M}{!B}[w]$ ہووے تے نہ $\mSat{M}{!C}[w]$}۔}{}
  \item\ollabel{defn:sub:mmodels-box}
    \indcase{!A}{\Box !B}{$\mSat{M}{\indfrm}[w]$ اُدوں تے صرف اُدوں جدوں
    $\mSat{M}{!B}[w']$ ہر اوس $w' \in W$ لئی ہووے جیہدے لئی $Rww'$ ہووے۔}
  \item\ollabel{defn:sub:mmodels-diamond}
    \indcase{!A}{\Diamond !B}{$\mSat{M}{\indfrm}[w]$ اُدوں تے صرف اُدوں جدوں
    $\mSat{M}{!B}[w']$ گھٹ توں گھٹ اک ایسے $w' \in W$ لئی ہووے جیہدے لئی $Rww'$ ہووے۔}
  \end{enumerate}
\end{defn}

دھیان دیو کہ شق~\olref{defn:sub:mmodels-box} مطابق !!a{formula}~$\Box
!B$، جہان~$w$ وچ سچا ہوندا اے جدوں کوئی~$w'$ ایسا نہ ہووے جیہدے لئی $Rww'$ ہووے۔ ایس صورت وچ $\Box !B$، جہان~$w$ وچ \emph{خالی شرط کارن} سچا اے۔ $\Box !B$، جہان~$w$ وچ سچا ہو سکدا اے بھاویں $!B$ اُتھے سچا نہ ہووے۔ $!B$ دا جہان~$w$ وچ سچا ہونا~$\Diamond !B$ دے جہان~$w$ وچ سچا ہون دی ضمانت نہیں دیندا۔ البتہ جے $Rww$ ہووے، جیویں جدوں $R$ انعکاسی ہووے، تاں ایہ گل پکی اے۔ جے کوئی $w'$ ایسا نہ ہووے جیہدے لئی $Rww'$ ہووے، تاں $\mSat/{M}{\Diamond !A}[w]$، ہر~$!A$ لئی۔

\begin{prob}
  \olref[nml][syn][rel]{fig:simple} والے نمونے نوں ویکھو۔ ہِیٹھ لکھیاں وچوں کہڑے بیان سچے نیں؟
    \begin{enumerate}
    \item $\mSat{M}{q}[w_1]$؛
    \item $\mSat{M}{\lnot q}[w_3]$؛
    \item $\mSat{M}{p \lor q}[w_1]$؛
    \item $\mSat{M}{\Box (p \lor q)}[w_1]$؛
    \item $\mSat{M}{\Box q}[w_3]$؛
    \item $\mSat{M}{\Box \bot}[w_3]$؛
    \item $\mSat{M}{\Diamond q}[w_1]$؛
    \item $\mSat{M}{\Box q}[w_1]$؛
    \item $\mSat{M}{\lnot \Box\Box \lnot q}[w_1]$۔
    \end{enumerate}
\end{prob}

% \Box تے \Diamond دے دوہرے تعلق دا ثبوت
\begin{tagblock}{prvDiamond}
\begin{prop}\ollabel{prop:dual}
  \begin{enumerate}
  \item $\mSat{M}{\Box !A}[w]$ اُدوں تے صرف اُدوں جدوں $\mSat{M}{\lnot\Diamond\lnot !A}[w]$۔
  \item $\mSat{M}{\Diamond !A}[w]$ اُدوں تے صرف اُدوں جدوں $\mSat{M}{\lnot\Box\lnot !A}[w]$۔
  \end{enumerate}
\end{prop}

\begin{proof}
  \begin{enumerate}
    \item $\mSat{M}{\lnot\Diamond\lnot !A}[w]$ اُدوں تے صرف اُدوں جدوں $\mSat/{M}{\Diamond\lnot
      !A}[w]$، $\mSat{M}{}[w]$ دی تعریف مطابق۔ $\mSat{M}{\Diamond\lnot
      !A}[w]$ اُدوں تے صرف اُدوں جدوں کسے ایسے $w'$ لئی جیہدے لئی $Rww'$ ہووے، $\mSat{M}{\lnot
      !A}[w']$ ہووے۔ ایس لئی $\mSat/{M}{\Diamond\lnot !A}[w]$ اُدوں تے صرف اُدوں جدوں ہر ایسے
    $w'$ لئی جیہدے لئی $Rww'$ ہووے، $\mSat/{M}{\lnot !A}[w']$ ہووے۔ نالے
    $\mSat/{M}{\lnot !A}[w']$ اُدوں تے صرف اُدوں جدوں $\mSat{M}{!A}[w']$۔ اینہاں نوں جوڑ کے
    $\mSat{M}{\lnot\Diamond\lnot !A}[w]$ اُدوں تے صرف اُدوں جدوں ہر ایسے $w'$ لئی جیہدے لئی $Rww'$ ہووے،
    $\mSat{M}{!A}[w']$ ہووے۔ فیر $\mSat{M}{}[w]$ دی تعریف مطابق، ایہ اُدوں تے صرف اُدوں ہوندا اے جدوں $\mSat{M}{\Box !A}[w]$۔
    \item \iftag{probDiamond}{مشق۔}{$\mSat{M}{\lnot\Box\lnot !A}[w]$ اُدوں تے صرف اُدوں جدوں
      $\mSat/{M}{\Box\lnot !A}$۔ $\mSat{M}{\Box\lnot !A}[w]$ اُدوں تے صرف اُدوں جدوں ہر ایسے
      $w'$ لئی جیہدے لئی $Rww'$ ہووے، $\mSat{M}{\lnot !A}[w']$ ہووے۔ ایس لئی
      $\mSat/{M}{\Box\lnot !A}[w]$ اُدوں تے صرف اُدوں جدوں کسے ایسے $w'$ لئی جیہدے لئی $Rww'$ ہووے،
      $\mSat/{M}{\lnot !A}[w']$ ہووے۔ نالے $\mSat/{M}{\lnot
        !A}[w']$ اُدوں تے صرف اُدوں جدوں $\mSat{M}{!A}[w']$۔ اینہاں نوں جوڑ کے
      $\mSat{M}{\lnot\Box\lnot !A}[w]$ اُدوں تے صرف اُدوں جدوں کسے ایسے~$w'$ لئی جیہدے لئی
      $Rww'$ ہووے، $\mSat{M}{!A}[w']$ ہووے۔ فیر
      $\mSat{M}{}[w]$ دی تعریف مطابق، ایہ اُدوں تے صرف اُدوں ہوندا اے جدوں $\mSat{M}{\Diamond
        !A}[w]$۔}
  \end{enumerate}
\end{proof}

\begin{probtag}{probDiamond}
  \olref[nml][syn][trw]{prop:dual} دا ثبوت پورا کرو۔
\end{probtag}
\end{tagblock}

\begin{prob}
  منّ لو کہ $\mModel{M} = \tuple{W, R, V}$ اک نمونہ اے تے $w_1, w_2 \in
  W$ لئی ہِیٹھ لکھیاں شرطاں پوری ہوندیاں نیں:
  \begin{enumerate}
  \item $w_1 \in V(p)$ اُدوں تے صرف اُدوں جدوں $w_2 \in V(p)$ (ہر !!{propositional variable} $p$ لئی)؛ تے
  \item ہر $w \in W$ لئی: $Rw_1w$ اُدوں تے صرف اُدوں جدوں $Rw_2w$۔
  \end{enumerate}
  !!{formula}s اُتے استقرا نال وکھاؤ کہ سارے !!{formula}s $!A$ لئی
  $\mSat{M}{!A}[w_1]$ اُدوں تے صرف اُدوں جدوں $\mSat{M}{!A}[w_2]$۔
\end{prob}

\begin{probtag}{prvDiamond}
  منّ لو کہ $\mModel{M} = \tuple{W, R, V}$۔ وکھاؤ کہ
  $\mSat{M}{\lnot\Diamond !A}[w]$ اُدوں تے صرف اُدوں جدوں
  $\mSat{M}{\Box\lnot!A}[w]$۔
\end{probtag}
\endgroup
% /OLP-0413
% OLP-0414 content/normal-modal-logic/syntax-and-semantics/truth-in-model.tex
\begingroup
\def\olfilename{truth-in-model}
\expandafter\def\csname ol@sectioncs\endcsname{section}

\olfileid{nml}{syn}{tru}

\olsection{نمونے وچ سچائی}

کدی اسیں جاننا چاہندے آں کہ دتے نمونے دے ہر جہان وچ کہڑے !!{formula}s سچے نیں۔ ایس لئی اک علامتی لکھت طے کر لئیے۔

\begin{defn}
  !!^a{formula}~$!A$ اک \emph{نمونے وچ سچا اے}، جیہڑا نمونہ $M = \tuple{W, R,
    V}$ اے، تے اینوں $\mSat{M}{!A}$ لکھدے آں، اُدوں تے صرف اُدوں جدوں $\mSat{M}{!A}[w]$
  ہر $w \in W$ لئی ہووے۔
\end{defn}

\begin{prop}\ollabel{prop:truthfacts}
  \begin{enumerate}
  \item جے $\mSat{M}{!A}$ تاں $\mSat/{M}{\lnot !A}$، پر اُلٹ گل \emph{نہیں} چل دی۔
  \item جے $\mSat{M}{!A \lif !B}$ تاں $\mSat{M}{!A}$ صرف اُدوں جدوں
    $\mSat{M}{!B}$، پر اُلٹ گل \emph{نہیں} چل دی۔
  \end{enumerate}
\end{prop}

\begin{proof}
  \begin{enumerate}
  \item جے $\mSat{M}{!A}$ تاں $!A$، $W$ دے سارے جہاناں وچ سچا اے۔ کیوں جے $W \neq \emptyset$، ایس لئی $\mSat{M}{\lnot!A}$ نہیں ہو سکدا؛ نئیں تاں $!A$ نوں کسے جہان وچ سچا تے جھوٹھا دونیں ہونا پئے گا۔

    دوجے پاسے جے $\mSat/{M}{\lnot !A}$ تاں $!A$ کسے جہان $w \in W$ وچ سچا اے۔ ایس توں ایہ نہیں نکلدا کہ $\mSat{M}{!A}[w]$
    \emph{ہر}~$w \in W$ لئی ہووے۔ مثال لئی
    \olref[rel]{fig:simple} والے نمونے وچ $\mSat/{M}{\lnot p}$ تے $\mSat/{M}{p}$ وی۔
  \item منّ لو $\mSat{M}{!A \lif !B}$ تے $\mSat{M}{!A}$۔ $\mSat{M}{!B}$ وکھاؤن لئی $w \in W$ کوئی وی جہان لے لو۔ فیر
    $\mSat{M}{!A \lif !B}[w]$ تے $\mSat{M}{!A}[w]$، ایس لئی
    $\mSat{M}{!B}[w]$۔ کیوں جے $w$ کوئی وی جہان سی، ایس لئی
    $\mSat{M}{!B}$۔

    ایہ وکھاؤن لئی کہ اُلٹ گل نہیں چل دی، سانوں اک ایسا نمونہ
    $\mModel{M}$ لبھنا پئے گا جیہدے وچ $\mSat{M}{!A}$ صرف اُدوں جدوں $\mSat{M}{!B}$، پر
    $\mSat/{M}{!A \lif !B}$۔ فیر
    \olref[rel]{fig:simple} والا نمونہ ویکھو۔ $\mSat/{M}{p}$، ایس لئی خالی شرط کارن
    $\mSat{M}{p}$ صرف اُدوں جدوں $\mSat{M}{q}$۔ پر $\mSat/{M}{p \lif
      q}$، کیوں جے $p$ سچا تے $q$ جھوٹھا اے جہان~$w_1$ وچ۔
  \end{enumerate}
\end{proof}

\begin{prob}
  اوس زبان لئی ہِیٹھلا نمونہ $\mModel{M}$ ویکھو جیہدے وچ
  $p_1$، $p_2$، $p_3$ ہی !!{propositional variable}s نیں:
  \begin{center}
    \begin{tikzpicture}[modal]
      \node[world] (w1) [label={[align=right]left:\mTrue{p_1}\\\mFalse{p_2}\\\mFalse{p_3}}]
            {$w_1$} ; 
      \node[world] (w2) [label={[align=right]right:\mTrue{p_1}\\\mTrue{p_2}\\\mFalse{p_3}},
        below right=of w1]  {$w_2$}; 
      \node[world] (w3) [label={[align=right]right:\mTrue{p_1}\\\mTrue{p_2}\\\mTrue{p_3}},
        above right=of w2] {$w_3$};
      \draw[reflexive above] (w3) to (w3); 
      \draw[->] (w1) to (w2);
      \draw[->] (w2) to (w3);
      \draw[->] (w1) to (w3); 
    \end{tikzpicture}
  \end{center}
  کی ہِیٹھلے !!{formula}s تے سانچے نمونے $\mModel{M}$ وچ سچے نیں،
  یعنی $\mModel{M}$ دے ہر جہان وچ سچے نیں؟ وجہ سمجھاؤ۔
  \begin{enumerate}
  \item $p\lif \Diamond p$ ($p$ ایٹمی اے)؛
  \item $!A\lif \Diamond !A$ ($!A$ کوئی وی فارمولا اے)؛
  \item $\Box p \lif p$ ($p$ ایٹمی اے)؛
  \item $\lnot p \lif \Diamond \Box p$ ($p$ ایٹمی اے)؛
  \item $\Diamond \Box !A$ ($!A$ کوئی وی فارمولا اے)؛
  \item $\Box \Diamond p$ ($p$ ایٹمی اے)۔
  \end{enumerate}
\end{prob}
\endgroup
% /OLP-0414
% OLP-0415 content/normal-modal-logic/syntax-and-semantics/modal-validity.tex
\begingroup
\def\olfilename{modal-validity}
\expandafter\def\csname ol@sectioncs\endcsname{section}

\olfileid{nml}{syn}{val}

\olsection{منطقی اعتبار}

\begin{explain}
  !!^{formula}s جیہڑے سارے نمونیاں وچ سچے نیں، یعنی ہر نمونے دے ہر جہان وچ سچے نیں، خاص طور تے دھیان جوگ نیں۔ اوہ اوہناں امکانی قضیاں نوں ظاہر کردے نیں جیہڑے $\Box$ تے $\Diamond$ دی تعبیر جو وی ہووے سچے رہندے نیں، جے تعبیر ایس معنی وچ «نارمل» ہووے کہ اوہ ممکن جہاناں دے کسے رسائی دے تعلق توں بندی اے۔ اسیں اینج دے !!{formula}s نوں \emph{معتبر} آکھدے آں۔ مثال لئی $\Box(p \land q) \lif \Box p$ معتبر اے۔ کجھ !!{formula}s، جیہڑے $\Box$ دی سچائی نال جڑی تعبیر دے آسرے معتبر لگ سکدے نیں، جیویں $\Box p \lif p$، دراصل معتبر نہیں۔ تعلقی نمونیاں دی اک خاص گل ایہ اے کہ $\Box$ تے $\Diamond$ دیاں وکھ وکھ تعبیراں رسائی دے تعلقاں دیاں وکھ وکھ قسماں نال دتیاں جا سکدیاں نیں۔ ایس توں خیال آؤندا اے کہ منطقی اعتبار دی تعریف \emph{سارے} نمونیاں دے لحاظ نال تے \emph{کسے خاص قسم دے} سارے نمونیاں دے لحاظ نال وی کر سکدے آں۔ مثال لئی پتا لگے گا کہ $\Box p \lif p$ اوہناں سارے نمونیاں وچ سچا اے جتھے ہر جہان اپنے آپ توں رسائی وچ اے، یعنی $R$ انعکاسی اے۔ نمونیاں دیاں جماعتاں دے لحاظ نال منطقی اعتبار دی تعریف نال اسیں ایہ گل مختصر ایویں آکھ سکدے آں: $\Box p \lif p$ انعکاسی نمونیاں دی جماعت وچ معتبر اے۔
\end{explain}

\begin{defn}
  !!^a{formula}~$!A$ نمونیاں دی جماعت $\mClass{C}$ وچ \emph{معتبر} اے جے اوہ~$\mClass{C}$ دے ہر نمونے وچ سچا اے (یعنی~$\mClass{C}$ دے ہر نمونے دے ہر جہان وچ سچا اے)۔ جے $!A$،~$\mClass{C}$ وچ معتبر اے، تاں اسیں $\mClass{C} \Entails !A$ لکھدے آں، تے $\Entails !A$ لکھدے آں جے $!A$ \emph{سارے} نمونیاں دی جماعت وچ معتبر اے۔
\end{defn}

\begin{prop}\ollabel{prop:subset-class}
  جے $!A$، $\mClass{C}$ وچ معتبر اے، تاں اوہ ہر جماعت
  $\mClass{C}' \subseteq \mClass{C}$ وچ وی معتبر اے۔
\end{prop}

\begin{prop}\ollabel{prop:Nec-rule}
  جے $!A$ معتبر اے، تاں $\Box!A$ وی معتبر اے۔
\end{prop}

\begin{proof}
  منّ لو $\Entails !A$۔ $\Entails \Box!A$ وکھاؤن لئی $\mModel{M} =
  \tuple{W, R, V}$ کوئی نمونہ لے لو تے $w \in W$ کوئی جہان۔ جے $Rww'$ تاں
  $\mSat{M}{!A}[w']$، کیوں جے $!A$ معتبر اے، تے ایس لئی
  $\mSat{M}{\Box!A}[w]$ وی۔ کیوں جے $\mModel{M}$ تے $w$ آپو اپنی تھاں کوئی وی ہو سکدے سن، ایس لئی
  $\Entails \Box!A$۔
\end{proof}

\begin{prob}
  وکھاؤ کہ ہِیٹھ لکھے فارمولے معتبر نیں:
  \begin{enumerate}
  \item $\Box p \lif \Box (q \lif p)$؛
  \item $\Box \lnot \lfalse$؛
  \item $\Box p \lif (\Box q \lif \Box p)$۔
  \end{enumerate}
\end{prob}

\begin{prob}
  وکھاؤ کہ $!A \lif \Box!A$ اوہناں نمونیاں دی جماعت $\mClass{C}$ وچ معتبر اے جیہڑے
  $\mModel{M} = \tuple{W, R, V}$ نیں تے جیہناں لئی $W = \{w\}$ اے۔ ایویں ای وکھاؤ کہ
  $!B \lif \Box !A$ تے $\Diamond !A \lif !B$ اوہناں نمونیاں دی جماعت وچ معتبر نیں جیہڑے
  $\mModel{M} = \tuple{W, R, V}$ نیں تے جیہناں لئی
  $R = \emptyset$ اے۔
\end{prob}
\endgroup
% /OLP-0415
% OLP-0416 content/normal-modal-logic/syntax-and-semantics/tautological-instances.tex
\begingroup
\def\olfilename{tautological-instances}
\expandafter\def\csname ol@sectioncs\endcsname{section}

\olfileid{nml}{syn}{tau}

\olsection{ہمیشہ صادق فارمولیاں توں بدل نال بنیاں مثالاں}

\begin{explain}
امکانی عامل توں خالی فارمولا اُدوں ہمیشہ صادق اے جدوں اوہ ہر صداقتی قدر بندی تحت سچا ہووے۔ صاف اے کہ ہر ہمیشہ صادق فارمولا ہر نمونے دے ہر جہان وچ سچا اے۔ پر $\Box$ تے~$\Diamond$ والے !!{formula}s لئی ہمیشہ صادق فارمولے دی تعریف نہیں دِتّی گئی۔ مثلاً کی $\Box p \lor \lnot \Box p$ معتبر اے، جیہڑا خارج شدہ وسط دے اصول دی اک مثال اے؟ \emph{ہمیشہ صادق فارمولے توں بدل نال بنی مثال} دا خیال ایتھے کم آؤندا اے۔ ایہ !!a{formula} اے جیہڑا کسے (غیر امکانی) ہمیشہ صادق فارمولے توں بدل نال بنیا ہووے۔ ایہ حیرانی دی گل نہیں کہ ہر ایسی مثال معتبر اے، پر ایس دا ثبوت فیر وی چاہیدا اے۔
\end{explain}

\begin{defn}
  امکانی !!{formula}~$!B$، \emph{ہمیشہ صادق فارمولے توں بدل نال بنی مثال} اُدوں تے صرف اُدوں اے جدوں اک امکانی عامل توں خالی ہمیشہ صادق فارمولا~$!A$ ہووے جیہدے !!{propositional variable}s $p_1$، \dots،~$p_n$ ہون، تے !!{formula}s $!D_1$،
  \dots،~$!D_n$ ایسے ہون کہ $!B \ident \SSubst{!A}{\subst{!D_1}{p_1},
    \dots, \subst{!D_n}{p_n}}$۔
\end{defn}

\begin{lem}\ollabel{lem:valid-taut}
  منّ لو $!A$ امکانی عامل توں خالی !!{formula} اے جیہدے !!{propositional variable}s $p_1$، \dots، $p_n$ نیں، تے $!D_1$، \dots،
  $!D_n$ امکانی !!{formula}s نیں۔ فیر کسے وی قدر بندی $\pAssign{v}$،
  کسے وی نمونے $\mModel{M} = \tuple{W, R, V}$، تے کسے وی $w \in W$ لئی جے
  $\pAssign{v}(p_i) = \True$ اُدوں تے صرف اُدوں ہووے جدوں $\mSat{M}{!D_i}[w]$، تاں
  $\pSat{v}{!A}$ اُدوں تے صرف اُدوں ہوندا اے جدوں
  $\mSat{M}{\SSubst{!A}{\subst{!D_1}{p_1}, \dots,
      \subst{!D_n}{p_n}}}[w]$۔
\end{lem}

\begin{proof}
  $!A$ اُتے استقرا نال۔
  \begin{enumerate}
  \tagitem{prvFalse}{\indcase{!A}{\lfalse}{$\pSat/{v}{\lfalse}$
      تے $\mSat/{M}{\lfalse}[w]$ دونیں}.}{}
  \tagitem{prvTrue}{\indcase{!A}{\ltrue}{$\pSat{v}{\ltrue}$
      تے $\mSat{M}{\ltrue}[w]$ دونیں}.}{}
  \item \indcase{!A}{p_i}{%
    \begin{align*}
      \pSat{v}{p_i} \Leftrightarrow {} & \pAssign{v}(p_i) = \True \\
      & \qquad \text{$\pSat{v}{p_i}$ دی تعریف مطابق}\\
      \Leftrightarrow {} & \mSat{M}{!D_i}[w] \\
      &\qquad \text{مفروضے مطابق}\\
      \Leftrightarrow {} & \mSat{M}{\SSubst{p_i}{\subst{!D_1}{p_1}, \dots,
          \subst{!D_n}{p_n}}}[w]\\
      &\qquad \text{کیوں جے $\SSubst{p_i}{\subst{!D_1}{p_1}, \dots,
          \subst{!D_n}{p_n}} \ident !D_i$}.
      \end{align*}}
  \tagitem{prvFalse}{\indcase{!A}{\lnot !B}{%
      \begin{align*}
        \pSat{v}{\lnot !B} \Leftrightarrow {} & \pSat/{v}{!B}\\
        &\qquad \text{$\pSat{v}{}$ دی تعریف مطابق};\\
        \Leftrightarrow {} & \mSat/{M}{\SSubst{!B}{\subst{!D_1}{p_1}, \dots,
          \subst{!D_n}{p_n}}}[w]\\
        &\qquad \text{استقرائی مفروضے مطابق}\\
        \Leftrightarrow {} &
        \mSat{M}{\SSubst{\lnot !B}{\subst{!D_1}{p_1}, \dots,
          \subst{!D_n}{p_n}}}[w]\\
        &\qquad \text{$\pSat{v}{}$ دی تعریف مطابق}.
      \end{align*}}}{}
  \tagitem{prvAnd}{\indcase{!A}{(!B \land !C)}{%
      \begin{align*}
        \pSat{v}{!B \land !C} \Leftrightarrow {} &
        \pSat{v}{!B} \text{ تے } \pSat{v}{!C}\\
        &\qquad \text{$\pSat{v}{}$ دی تعریف مطابق}\\
        \Leftrightarrow {} &
        \mSat{M}{\SSubst{!B}{\subst{!D_1}{p_1}, \dots,
            \subst{!D_n}{p_n}}}[w] \text{ تے } \\
        & \mSat{M}{\SSubst{!C}{\subst{!D_1}{p_1}, \dots,
          \subst{!D_n}{p_n}}}[w]\\
        &\qquad \text{استقرائی مفروضے مطابق}\\
        \Leftrightarrow{} &
        \mSat{M}{\SSubst{(!B \land !C)}{\subst{!D_1}{p_1}, \dots,
          \subst{!D_n}{p_n}}}[w]\\
        &\qquad \text{$\mSat{M}{}[w]$ دی تعریف مطابق}.
      \end{align*}}}{}
  \tagitem{prvOr}{\indcase{!A}{(!B \lor !C)}{%
      \begin{align*}
        \pSat{v}{!B \lor !C} \Leftrightarrow{} &
        \pSat{v}{!B} \text{ یا } \pSat{v}{!C}\\
        &\qquad \text{$\pSat{v}{}$ دی تعریف مطابق};\\
        \Leftrightarrow{} & \mSat{M}{\SSubst{!B}{\subst{!D_1}{p_1}, \dots,
            \subst{!D_n}{p_n}}}[w] \text{ یا }\\
        & \mSat{M}{\SSubst{!C}{\subst{!D_1}{p_1}, \dots,
          \subst{!D_n}{p_n}}}[w]\\
        &\qquad \text{استقرائی مفروضے مطابق}\\
        \Leftrightarrow{} &
        \mSat{M}{\SSubst{(!B \lor !C)}{\subst{!D_1}{p_1}, \dots,
          \subst{!D_n}{p_n}}}[w]\\
        &\qquad \text{$\mSat{M}{}[w]$ دی تعریف مطابق}.
      \end{align*}}}{}
  \tagitem{prvIf}{\indcase{!A}{(!B \lif !C)}{%
      \begin{align*}
        \pSat{v}{!B \lif !C} \Leftrightarrow{} &
        \pSat/{v}{!B} \text{ یا } \pSat{v}{!C}\\
        &\qquad \text{$\pSat{v}{}$ دی تعریف مطابق}\\
        \Leftrightarrow{} &
        \mSat/{M}{\SSubst{!B}{\subst{!D_1}{p_1}, \dots,
            \subst{!D_n}{p_n}}}[w] \text{ یا }\\
        & \mSat{M}{\SSubst{!C}{\subst{!D_1}{p_1}, \dots,
          \subst{!D_n}{p_n}}}[w]\\
        &\qquad \text{استقرائی مفروضے مطابق}\\
        \Leftrightarrow{} &
        \mSat{M}{\SSubst{(!B \lif !C)}{\subst{!D_1}{p_1}, \dots,
          \subst{!D_n}{p_n}}}[w]\\
        &\qquad \text{$\mSat{M}{}[w]$ دی تعریف مطابق}.
      \end{align*}}}{}
  \tagitem{prvIff}{\indcase{!A}{(!B \liff !C)}{%
      \begin{align*}
        \pSat{v}{!B \lif !C} \Leftrightarrow {} &
        \text{یا تاں } \pSat{v}{!B} \text{ تے } \pSat{v}{!C}\\
        & \text{or } \pSat/{v}{!B} \text{ تے } \pSat/{v}{!C}\\
        &\qquad \text{$\pSat{v}{}$ دی تعریف مطابق}\\
        \Leftrightarrow {} &
        \text{یا تاں } \mSat{M}{\SSubst{!B}{\subst{!D_1}{p_1}, \dots,
            \subst{!D_n}{p_n}}}[w] \text{ تے }\\
        & \qquad \mSat{M}{\SSubst{!C}{\subst{!D_1}{p_1}, \dots,
          \subst{!D_n}{p_n}}}[w]\\ 
        &  \text{ یا } \mSat/{M}{\SSubst{!B}{\subst{!D_1}{p_1}, \dots,
            \subst{!D_n}{p_n}}}[w] \text{ تے }\\
        & \qquad \mSat/{M}{\SSubst{!C}{\subst{!D_1}{p_1}, \dots,
          \subst{!D_n}{p_n}}}[w]\\
        &\qquad\text{استقرائی مفروضے مطابق}\\
        \Leftrightarrow {} &
        \mSat{M}{\SSubst{(!B \liff !C)}{\subst{!D_1}{p_1}, \dots,
          \subst{!D_n}{p_n}}}[w]\\
        &\qquad \text{$\mSat{M}{}[w]$ دی تعریف مطابق}.
      \end{align*}}}{}
  \end{enumerate}
  \end{proof}

\begin{prop}\ollabel{prop:valid-taut}
  ہمیشہ صادق فارمولیاں توں بدل نال بنیاں ساریاں مثالاں معتبر نیں۔
\end{prop}

\begin{proof}
  اُلٹی نفی نال ثابت کرن لئی، منّ لو $!A$ ایسا اے کہ
  $\mSat/{M}{\SSubst{!A}{\subst{!D_1}{p_1}, \dots,
  \subst{!D_n}{p_n}}}[w]$، کسے نمونے $\mModel{M}$ تے جہان~$w$ لئی۔ قدر بندی
  $\pAssign{v}$ ایویں متعین کرو کہ $\pAssign{v}(p_i) =
  \True$ اُدوں تے صرف اُدوں جدوں $\mSat{M}{!D_i}[w]$ (تے $\pAssign{v}$،
  $q \notin \{p_1, \dots, p_n \}$ نوں من مرضی دیاں قدراں دیوے)۔ فیر
  \olref{lem:valid-taut} مطابق، $\pSat/{v}{!A}$، ایس لئی $!A$ ہمیشہ صادق فارمولا نہیں۔
\end{proof}
\endgroup
% /OLP-0416
% OLP-0417 content/normal-modal-logic/syntax-and-semantics/schemas.tex
\begingroup
\def\olfilename{schemas}
\expandafter\def\csname ol@sectioncs\endcsname{section}

\olfileid{nml}{syn}{sch}

\olsection{سانچے تے منطقی اعتبار}

\begin{defn}
  \emph{سانچہ} !!{formula}s دا اوہ سیٹ اے جیہدے وچ کسے امکانی !!{formula}~$!C$ توں بدل نال بنیاں ساریاں مثالاں ہون تے صرف اوہو ہون، یعنی
  \[
  \Setabs{!B}{\lexists[!D_1], \dots, \lexists[!D_n] \left(!B =
  \SSubst{!C}{\subst{!D_1}{p_1}, \dots, \subst{!D_n}{p_n}} \right) }.
  \]
  !!{formula}~$!C$ نوں سانچے دا \emph{خاص} !!{formula} آکھدے آں۔ !!{propositional variable}s دے ناں بدلّن نوں چھڈ کے، ایہ اکّو اے۔ !!^a{formula}~$!A$ سانچے دی \emph{مثال} اے جے اوہ ایس سیٹ دا رکن اے۔
\end{defn}

سانچے نوں اوہ مابعد زبانی عبارت نال لکھنا سوکھا اے جیہڑی «$!A$»، «$!B$»،~\dots، نوں ایٹمی جزاں دی تھاں رکھ کے~$!C$ توں ملدی اے۔ مثال لئی ایہ لکھتاں سانچے ظاہر کردیاں نیں۔
«$!A$»، «$!A \lif \Box!A$»، «$!A \lif (!B \lif !A)$»۔ اینہاں دے خاص !!{formula}s بالترتیب $p$، $p \lif \Box p$، $p \lif (q
\lif p)$ نیں۔ سانچہ «$!A$» \emph{سارے} !!{formula}s دا سیٹ ظاہر کردا اے۔

\begin{defn}
  کوئی سانچہ نمونے وچ \emph{سچا} اُدوں تے صرف اُدوں اے جدوں اوہدیاں ساریاں مثالاں اوس نمونے وچ سچیاں ہون۔ سانچہ \emph{معتبر} اُدوں تے صرف اُدوں اے جدوں اوہ ہر نمونے وچ سچا ہووے۔
\end{defn}

\begin{prop}\ollabel{prop:Kvalid}
  ہِیٹھلا سانچہ \Ax{K} معتبر اے۔
  \begin{equation*}
  \Box(!A \lif !B) \lif (\Box !A \lif \Box !B). \tag{\Ax{K}}
  \end{equation*}
\end{prop}

\begin{proof}
  سانوں وکھاؤنا اے کہ سانچے دیاں ساریاں مثالاں ہر نمونے دے ہر جہان وچ سچیاں نیں۔ ایس لئی $\mModel{M} = \tuple{W,R,V}$ تے $w \in
  W$ کوئی وی لے لو۔ کسے جہان وچ شرطیہ سچا وکھاؤن لئی اسیں منّ لیندے آں کہ مقدم سچا اے تے وکھاؤندے آں کہ تالی وی سچا اے۔ ایتھے منّ لو $\mSat{M}{\Box(!A \lif !B)}[w]$ تے
  $\mSat{M}{\Box !A}[w]$۔ سانوں $\mSat{M}{\Box !B}[w]$ وکھاؤنا اے۔ کوئی وی ایسا $w'$ لے لو جیہدے لئی $Rww'$ ہووے۔ پہلے مفروضے مطابق
  $\mSat{M}{!A \lif !B}[w']$ تے دوجے مفروضے مطابق
  $\mSat{M}{!A}[w']$۔ ایس توں $\mSat{M}{!B}[w']$ نکلدا اے۔ کیوں جے $w'$ کوئی وی ایسا جہان سی، ایس لئی $\mSat{M}{\Box !B}[w]$۔
\end{proof}

\begin{prop}\ollabel{prop:Dual-valid}
  ہِیٹھلا سانچہ \Dual{} معتبر اے۔
  \begin{equation*}
  \Diamond !A \liff \lnot\Box\lnot !A. \tag{\Dual}
  \end{equation*}
\end{prop}

\begin{proof}
  مشق۔
\end{proof}

\begin{prob}
  \olref[nml][syn][sch]{prop:Dual-valid} ثابت کرو۔
\end{prob}

\begin{prop}\ollabel{prop:soundMP}
  جے $!A$ تے $!A \lif !B$ کسے نمونے دے کسے جہان وچ سچے نیں، تاں $!B$ وی اُتھے سچا اے۔ ایس لئی معتبر !!{formula}s دا سیٹ قاعدۂ اثباتِ مقدم تحت بند اے۔
\end{prop}

\begin{prop}\ollabel{prop:valid-instances}
  !!^a{formula}~$!A$ اُدوں تے صرف اُدوں معتبر اے جدوں اوہدیاں بدل توں بنیاں ساریاں مثالاں معتبر ہون۔ دوجے لفظاں وچ سانچہ اُدوں تے صرف اُدوں معتبر اے جدوں اوہدا خاص !!{formula} معتبر ہووے۔
\end{prop}

\begin{proof}
  «جے» والا رخ صاف اے، کیوں جے $!A$ اپنے آپ توں بدل نال بنن والی اک مثال اے۔

  «صرف جے» والا رخ ثابت کرن لئی اسیں ہِیٹھلی گل وکھاؤندے آں۔ منّ لو
  $\mModel{M} = \tuple{W, R, V}$ اک امکانی نمونہ اے تے $!B \ident
  \SSubst{!A}{\subst{!D_1}{p_1},\dots,\subst{!D_n}{p_n}}$ بدل نال بنی اک مثال اے~$!A$ دی۔ $\mModel{M'} = \tuple{W, R,
    V'}$ دی تعریف $V'(p_i) = \Setabs{w}{\mSat{M}{!D_i}[w]}$ نال کرو۔ فیر
  $\mSat{M}{!B}[w]$ اُدوں تے صرف اُدوں جدوں $\mSat{M'}{!A}[w]$، ہر~$w \in W$ لئی۔ (ایہدا ثبوت مشق لئی چھڈدے آں۔) ہُن منّ لو $!A$ معتبر اے پر بدل توں بنی کوئی مثال $!B$،~$!A$ دی، معتبر نہیں۔ فیر کسے
  $\mModel{M} = \tuple{W, R, V}$ تے کسے $w \in W$ لئی
  $\mSat/{M}{!B}[w]$۔ پر فیر دعوے مطابق $\mSat/{M'}{!A}[w]$، تے
  $!A$ معتبر نہیں۔ ایہ تضاد اے۔
\end{proof}

\begin{prob}
  «صرف جے» والے حصے دا دعویٰ ثابت کرو، جیہڑا \olref[nml][syn][sch]{prop:valid-instances} دے ثبوت وچ اے۔ (اشارہ: استقرا کرو~$!A$ اُتے۔)
\end{prob}

دھیان دیو کہ کسے نمونے وچ سانچے دے سچا ہون تے اوہدے خاص فارمولے دے سچا ہون دی دوہاں رخاں دی شرط نہیں چلدی۔ «صرف جے» والا رخ ضرور چلدا اے۔ جے $!A$ دی ہر مثال~$\mModel{M}$ وچ سچی اے، تاں
$!A$~خود وی~$\mModel{M}$ وچ سچا اے۔ پر ہو سکدا اے کہ $!A$،
~$\mModel{M}$ وچ سچا ہووے تے~$!A$ دی کوئی مثال کسے جہان وچ جھوٹھی ہووے، اوسے~$\mModel{M}$ وچ۔ اک سادی ردّی مثال لئی~$p$ نوں ایسے نمونے وچ ویکھو جیہدے وچ صرف اک جہان~$w$ اے تے $V(p) = \{w\}$، سو $p$ سچا اے~$w$ وچ۔ پر $\lfalse$،~$p$ دی مثال اے تے~$w$ وچ سچا نہیں۔

\begin{prob}
  وکھاؤ کہ ہِیٹھ لکھیاں !!{formula}s وچوں کوئی وی معتبر نہیں۔
  \begin{enumerate}
    \item[\Ax{D}:] \quad $\Box p \lif \Diamond p$؛
    \item[\Ax{T}:] \quad $\Box p \lif p$؛
    \item[\Ax{B}:] \quad $p \lif \Box\Diamond p$؛
    \item[\Ax{4}:] \quad $\Box p \lif \Box \Box p$؛
    \item[\Ax{5}:] \quad $\Diamond p \lif \Box \Diamond
      p$۔
  \end{enumerate}
\end{prob}

\begin{table}[t]
    \centering
    \begin{tabular}{| l || l |}
      \hline
      {\emph{معتبر سانچے}} & {\emph{غیر معتبر سانچے}} \\
      \hline\hline
      $\Box(!A \lif !B) \lif (\Diamond !A \lif \Diamond !B)$
      & $\Box (!A \lor !B) \lif (\Box !A \lor \Box !B)$ \\
      $\Diamond (!A \lif !B) \lif (\Box !A \lif \Diamond
      !B)$
      & $(\Diamond !A \land \Diamond !B) \lif \Diamond (!A
      \land !B)$\\
      $\Box (!A \land !B) \liff (\Box !A \land \Box !B)$
      & $!A \lif \Box !A$ \\
      $\Box !A \lif \Box (!B \lif !A)$
      & $\Box \Diamond !A \lif !B$ \\
      $\lnot \Diamond !A \lif \Box (!A \lif !B)$
      & $\Box \Box !A \lif \Box !A$ \\
      $\Diamond (!A \lor !B) \liff (\Diamond !A \lor
      \Diamond !B)$
      & $\Box \Diamond !A \lif \Diamond \Box !A$. \\
      \hline
    \end{tabular}
    \caption{معتبر تے (یا؟) غیر معتبر سانچے۔}
    \ollabel{tab:valid-invalidSchemas}
\end{table}

\begin{prob}%\ollabel{ex:in/validSchemas}
  ثابت کرو کہ پہلے کالم دے سانچے، جیہڑے \olref[nml][syn][sch]{tab:valid-invalidSchemas} وچ نیں، معتبر نیں تے دوجے کالم دے سانچے معتبر نہیں۔
\end{prob}

\begin{prob}
  فیصلہ کرو کہ ہِیٹھلے سانچے معتبر نیں یا غیر معتبر۔
  \begin{enumerate}
  \item $(\Diamond !A \lif \Box !B) \lif (\Box !A \lif \Box
    !B)$؛
  \item $\Diamond(!A \lif !B) \lor \Box(!B \lif !A)$۔
  \end{enumerate}
\end{prob}

\begin{prob}
  ہِیٹھلے ہر سانچے لئی اک ایسا نمونہ $\mModel{M}$ لبھو جیہدے وچ !!{formula} دی ہر مثال $\mModel{M}$ وچ سچی ہووے۔
  \begin{enumerate}
  \item $p \lif \Diamond\Diamond p$؛
  \item $\Diamond p \lif \Box p$۔
  \end{enumerate}
\end{prob}
\endgroup
% /OLP-0417
% OLP-0418 content/normal-modal-logic/syntax-and-semantics/entailment.tex
\begingroup
\def\olfilename{entailment}
\expandafter\def\csname ol@sectioncs\endcsname{section}

\olfileid{nml}{syn}{ent}

\olsection{منطقی نتیجہ}

\begin{explain}
  جہان وچ سچائی دی تعریف نال اسیں !!{formula}s دے وچ منطقی نتیجے دا تعلق متعین کر سکدے آں۔ !!^a{formula}~$!B$ توں~$!A$ منطقی نتیجے وجوں نکلدا اے اُدوں تے صرف اُدوں جدوں $!B$ دے سچا ہون نال $!A$ وی سچا ہووے۔ ایتھے «جدوں» دا مطلب اے کہ نمونہ کوئی وی ہووے تے اوس نمونے وچ جہان وی کوئی وی ہووے۔
\end{explain}

\begin{defn}
  جے $\Gamma$، !!{formula}s دا سیٹ اے تے $!A$، !!a{formula} اے، تاں $\Gamma$ توں~$!A$ \emph{منطقی نتیجے وجوں نکلدا اے}، علامتاں وچ $\Gamma \Entails !A$، اُدوں تے صرف اُدوں جدوں ہر نمونے $\mModel{M} = \tuple{W, R, V}$ تے ہر جہان
  $w \in W$ لئی، جے $\mSat{M}{!B}[w]$ ہر $!B \in \Gamma$ لئی ہووے، تاں
  $\mSat{M}{!A}[w]$۔ جے $\Gamma$ وچ صرف اک !!{formula}~$!B$ اے، تاں اسیں $!B \Entails !A$ لکھدے آں۔
\end{defn}

\begin{ex}
  ایہ وکھاؤن لئی کہ اک !!a{formula} توں دوجا منطقی نتیجے وجوں نکلدا اے، سانوں سارے نمونیاں بارے $\mSat{M}{}[w]$ دی تعریف ورت کے دلیل دینی پیندی اے۔ مثال لئی
  $p \lif \Diamond p \Entails \Box\lnot p \lif \lnot p$ وکھاؤن لئی اسیں ایویں دلیل دے سکدے آں۔ اک نمونہ $\mModel{M} = \tuple{W, R,
    V}$ تے $w \in W$ لے لو تے منّ لو $\mSat{M}{p \lif \Diamond
    p}[w]$۔ سانوں $\mSat{M}{\Box\lnot p \lif \lnot
    p}[w]$ وکھاؤنا اے۔ منّ لو ایہ گل نہیں۔ فیر $\mSat{M}{\Box\lnot p}[w]$ تے
  $\mSat/{M}{\lnot p}[w]$۔ کیوں جے $\mSat/{M}{\lnot p}[w]$، ایس لئی $\mSat{M}{
    p}[w]$۔ مفروضے مطابق $\mSat{M}{p \lif \Diamond p}[w]$، ایس لئی
  $\mSat{M}{\Diamond p}[w]$۔ $\mSat{M}{\Diamond
    p}[w]$ دی تعریف مطابق کوئی ایسا $w'$ اے جیہدے لئی $Rww'$ ہووے تے
  $\mSat{M}{p}[w']$ ہووے۔ کیوں جے $\mSat{M}{\Box \lnot p}[w]$ وی اے، ایس لئی
  $\mSat{M}{\lnot p}[w']$۔ ایہ تضاد اے۔

  ایہ وکھاؤن لئی کہ کسے !!a{formula}~$!B$ توں دوجا فارمولا~$!A$ منطقی نتیجے وجوں نہیں نکلدا، سانوں اک ردّی مثال دینی پیندی اے۔ یعنی اک ایسا نمونہ $\mModel{M} = \tuple{W, R,
    V}$ دیندے آں جیہدے کسے جہان~$w \in W$ وچ $\mSat{M}{!B}[w]$ ہووے
  پر $\mSat/{M}{!A}[w]$۔ وکھائیے کہ $p \lif \Diamond p \Entails/
  \Box p \lif p$۔ \olref{fig:counterex} والا نمونہ ویکھو۔
  $\mSat{M}{\Diamond p}[w_1]$، ایس لئی $\mSat{M}{p \lif
    \Diamond p}[w_1]$۔ پر کیوں جے $\mSat{M}{\Box p}[w_1]$ تے
  $\mSat/{M}{p}[w_1]$، ایس لئی $\mSat/{M}{\Box p \lif p}[w_1]$۔
  \begin{figure}
  \begin{center}
    \begin{tikzpicture}[modal]
      \node[world] (w1) [label=right:\mFalse{p}]{$w_1$}; 
      \node[world] (w2) [label=right:\mTrue{p}, above left=of w1]{$w_2$}; 
      \node[world] (w3) [label=right:\mTrue{p}, above right=of w1] {$w_3$};
      \draw[->] (w1) to (w2);
      \draw[->] (w1) to (w3);
    \end{tikzpicture}
  \end{center}
  \caption{$p \lif \Diamond p
  \Entails \Box p \lif p$ دی ردّی مثال۔}
  \ollabel{fig:counterex}
  \end{figure}

  اکثر بڑی سادی ردّی مثال وی کافی ہوندی اے۔ نمونہ $\mModel{M'} =
  \{W', R', V'\}$، جیہدے وچ $W' = \{w\}$، $R' = \emptyset$ تے $V'(p) =
  \emptyset$ اے، وی اک ردّی مثال اے۔ کیوں جے $\mSat/{M'}{p}[w]$،
  ایس لئی $\mSat{M'}{p \lif \Diamond p}[w]$۔ کیوں جے کوئی جہان رسائی وچ نہیں~$w$ توں، ایس لئی $\mSat{M'}{\Box p}[w]$، تے فیر $\mSat/{M'}{\Box p
    \lif p}[w]$۔
\end{ex}

\begin{prob}
  وکھاؤ کہ $\Box (!A \land !B) \Entails \Box !A$۔
\end{prob}

\begin{prob}
  وکھاؤ کہ $\Box (p \lif q) \Entails/ p \lif \Box q$ تے $p \lif \Box
  q \Entails/\Box (p \lif q)$۔
\end{prob}
\endgroup
% /OLP-0418

% OLP-0654 content/normal-modal-logic/syntax-and-semantics/normal-modal-logics.tex
\begingroup
\def\olfilename{normal-modal-logics}
\expandafter\def\csname ol@sectioncs\endcsname{section}

\olfileid{nml}{syn}{nor}

\olsection{نارمل امکانی منطقاں}

امکانی منطق نوں !!{formula}s دا کوئی ایہو جیہا سیٹ سمجھیا جا سکدا اے جیہڑا ڈھکدیاں شرطاں پوریاں کردا ہووے۔ کجھ امکانی منطقاں خاص دلچسپی رکھدیاں نیں۔ اوہناں دیاں تعبیراں تے ورتوں دلچسپ نیں، پر وڈی گل ایہ اے کہ نظریاتی طور اُتے اوہناں دی چنگی سمجھ اے۔ ایہناں نوں نارمل امکانی منطقاں آکھدے نیں۔

\begin{defn}\ollabel{defn:modal-logic}
اک \emph{امکانی منطق}~$\Log L$ بنیادی امکانی زبان دے !!{formula}s دا سیٹ اے جیہڑا
\begin{enumerate}
\item سارے قضیائی ہمیشہ سچے فارمولے رکھدا اے،
\item یکساں بدلی ہیٹھ بند اے، تے
\item اثباتِ مقدم ہیٹھ بند اے، یعنی جدوں وی $!A$ تے $(!A \lif
  !B) \in \Log L$، تاں $!B \in \Log L$ وی۔
\end{enumerate}
\end{defn}

\begin{ex}
!!{formula}s دے کجھ سیٹ، جیہڑے منّیا کہ بہتے دلچسپ نہیں، ایہ تعریف پوری کردے نیں تے ایس لئی امکانی منطقاں نیں:
\begin{enumerate}
\item سارے ہمیشہ سچے فارمولیاں دا سیٹ،
\item سارے !!{formula}s دا سیٹ (متضاد امکانی منطق)۔
\end{enumerate}
\end{ex}

\begin{defn}\label{defn:normal-modal-logic}
امکانی منطق~$\Log L$ اُدوں تے صرف اُدوں \emph{نارمل} اے جدوں
\begin{enumerate}
\item ایہدے وچ دوویں نیں:
\begin{align}
\tag{K} \Box(p \lif q) & \lif (\Box p \lif \Box q) \\
\tag{Dual} \Diamond p & \liff \lnot \Box \lnot p
\end{align}
\item تے ایہ لازمیت دے قاعدے ہیٹھ بند اے، یعنی جدوں وی $!A \in
  \Log L$، تاں $\Box !A \in \Log L$ وی۔
\end{enumerate}
\end{defn}

امکانی منطقاں دی تعریف دے دو مُکھ ڈھنگ سامنے آؤن گے۔ پہلا ثبوتی نظریے والا اے۔ ایس وچ اوہ !!{formula}s دا سیٹ ہوندی اے جیہڑے کجھ !!{derivation} نظاماں وچ اخذ کیتے جا سکدے نیں۔ دوجا نمونیاں دے نظریے والا اے۔ ایس وچ اوہ !!{formula}s دے سیٹ ہوندے نیں جیہڑے نمونیاں دیاں کجھ جماعتاں وچ معتبر نیں۔ ایہ عام قضیائی منطق (تے ایس گل لئی پہلے درجے دی منطق وی) والی صورت وانگ اے۔ ایتھے وی اسیں !!{formula}s دا سیٹ دو وکھرے ڈھنگاں نال وکھ کر سکدے آں، مثلاً اوہناں ہمیشہ سچے فارمولیاں دا سیٹ جیہڑے سچائی دیاں قدراں دی ہر قدر بندی وچ سچ نیں، تے اوہناں سارے !!{derivable} !!{formula}s دا سیٹ جیہڑے کسے !!{derivation} نظام وچ نیں۔ نارمل امکانی منطقاں نوں ایہناں دوہاں ڈھنگاں نال مقرر کرنا تعلقی نمونے تے مسلّماتی (ہلبرٹ طرز دے) ثبوتی نظام ورت کے ممکن اے۔
\endgroup
% /OLP-0654
\OLEndChapterHook
\endgroup
% /OLP-0408

% OLP-0419 content/normal-modal-logic/frame-definability/frame-definability.tex
\begingroup
\def\olfilename{frame-definability}
\expandafter\def\csname ol@sectioncs\endcsname{section}

\olchapter{nml}{frd}{فریماں دیاں جماعتاں دی تعریف}

% OLP-0420 content/normal-modal-logic/frame-definability/introduction.tex
\begingroup
\def\olfilename{introduction}
\expandafter\def\csname ol@sectioncs\endcsname{section}

\olfileid{nml}{frd}{int}

\olsection{جان پہچان}

امکانی منطق دے ماہراں لئی اک دلچسپ سوال ایہ اے کہ رسائی دے تعلق تے اوس تعلق والے نمونیاں وچ کجھ !!{formula}s دی سچائی دا آپس وچ کیہ تعلق اے۔ مثال لئی منّ لو رسائی دا تعلق انعکاسی اے، یعنی ہر $w
\in W$ لئی $Rww$۔ دوجے لفظاں وچ ہر جہان اپنے آپ توں رسائی وچ اے۔ ایہدا مطلب اے کہ جدوں $\Box !A$ جہان~$w$ وچ سچا اے، تاں $w$ آپ وی رسائی والے اوہناں جہاناں وچ اے جیہناں وچ~$!A$ لازماً سچا اے۔ سو جے نمونے دے رسائی دا تعلق~$R$،~$\mModel{M}$ وچ، انعکاسی اے، تاں اسیں جیہڑا وی جہان $w$ تے فارمولا $!A$ لئیے، $\Box !A
\lif !A$ اُتھے سچا ہووے گا (یعنی سانچہ $\Box p \lif
p$ تے اوہدیاں بدل نال بنیاں ساریاں مثالاں~$\mModel{M}$ وچ سچیاں نیں)۔

پر اُلٹی گل جھوٹھ اے۔ مثال لئی ایہ گل نہیں کہ جے
$\Box p \lif p$،~$\mModel{M}$ وچ سچا اے، تاں $R$ انعکاسی اے۔ کیوں جے اک غیر انعکاسی نمونہ~$\mModel{M}$ سوکھا لبھ سکدا اے جیہدے وچ $\Box p \lif
p$ سارے جہاناں وچ سچا اے۔ اوہ نمونہ لے لو جیہدے وچ صرف اک جہان~$w$ اے، اوہ اپنے آپ توں رسائی وچ نہیں، پر $w \in V(p)$ اے۔ $p$ دی صداقتی قدر ٹھیک چُن کے اسیں $\Box !A \lif !A$ نوں غیر انعکاسی نمونے وچ سچا کر سکدے آں۔

ایہدا حل اے کہ متغیر قدر بندی~$V$ نوں ایس شرط توں کڈھ دیئے۔ جے اسیں ایہ شرط لائیے کہ $\Box p \lif p$،~$\mModel{M}$ دے سارے جہاناں وچ سچا ہووے، بھاویں~$V(p)$ وچ جیہڑے مرضی جہان ہون، تاں $R$ دا انعکاسی ہونا ضروری اے۔ کیوں جے ہر غیر انعکاسی نمونے وچ گھٹ توں گھٹ اک جہان~$w$ ہووے گا جیہدے لئی $Rww$ نہیں۔ جے اسیں $V(p) =
W \setminus \{w\}$ رکھ دیئے، تاں $p$ سارے جہاناں وچ سچا ہووے گا، سوائے~$w$ دے۔ ایس لئی اوہ سارے اوہناں جہاناں وچ وی سچا ہووے گا جیہڑے~$w$ توں رسائی وچ نیں (کیوں جے $w$ دا~$w$ توں رسائی وچ نہ ہونا پکا اے، تے $w$ اکّو جہان اے جتھے~$p$ جھوٹھا اے)۔ دوجے پاسے $p$، $w$ وچ جھوٹھا اے، ایس لئی $\Box
p \lif p$،~$w$ وچ جھوٹھا اے۔

ایس توں خیال آؤندا اے کہ قدر بندی توں بغیر نمونے دی ساخت لئی اک علامتی لکھت ورتیئے۔ ایسی ساخت نوں اسیں \emph{فریم} آکھدے آں۔ فریم $\mModel{F}$ اک جوڑا $\tuple{W, R}$ اے جیہدے وچ جہاناں دا سیٹ تے رسائی دا تعلق اے۔ فیر ہر نمونہ $\tuple{W, R, V}$، فریم $\tuple{W, R}$ \emph{اُتے ٹکیا} اے۔ اُلٹے رخ وچ اک فریم اوہدے اُتے ٹکے نمونیاں دی جماعت متعین کردا اے، تے فریماں دی اک جماعت اوہ سارے نمونے متعین کردی اے جیہڑے اوس جماعت دے کسے فریم اُتے ٹکے نیں۔ اسیں $\mModel{F} \Entails !A$، یعنی کسے فریم وچ !!a{formula} دے \emph{معتبر} ہون دی تعریف، ایویں دے سکدے آں۔ $\mSat{M}{!A}$ ہر اوس $\mModel{M}$ لئی ہووے جیہڑا~$\mModel{F}$ اُتے ٹکیا اے۔

ایس علامتی لکھت نال اسیں !!{formula}s تے فریماں دیاں جماعتاں وچ مطابقت قائم کر سکدے آں۔ مثال لئی $\mModel{F} \Entails \Box p \lif p$ اُدوں تے صرف اُدوں جدوں $\mModel{F}$ انعکاسی اے۔
\endgroup
% /OLP-0420
% OLP-0421 content/normal-modal-logic/frame-definability/properties-accessibility.tex
\begingroup
\def\olfilename{properties-accessibility}
\expandafter\def\csname ol@sectioncs\endcsname{section}

\olfileid{nml}{frd}{acc}

\olsection{رسائی دے تعلقاں دیاں خاصیتاں}

کئی امکانی !!{formula}s رسائی دے تعلق دیاں سادیاں، سگوں جانیاں پچھانیاں، خاصیتاں نوں ظاہر کردے نیں۔ اک رخ وچ ایہدا مطلب اے کہ دِتّی خاصیت والے ہر نمونے وچ اوہدے نال جڑیا !!{formula} (تے اوہدیاں بدل نال بنیاں ساریاں مثالاں) سچا ہوندا اے۔ اسیں رسائی دے تعلقاں دیاں پنج کلاسیکی قسماں تے اوہناں فارمولیاں توں شروع کردے آں جیہڑے اوہناں قسماں دے تعلق نال ضرور سچے ہوندے نیں۔

\begin{thm}\ollabel{thm:soundschemas}
  منّ لو $\mModel{M} = \tuple{W, R, V}$ اک نمونہ اے۔ جے $R$ وچ \olref{tab:five} دے کھبے پاسے والی خاصیت اے، تاں سجے پاسے والے !!{formula} دی ہر مثال~$\mModel{M}$ وچ سچی اے۔
\end{thm}

\begin{table}[t]
    \begin{tabular}{| l || l |}
      \hline
      {\emph{جے $R$ \dots اے}} & {\emph{تاں \dots،~$\mModel{M}$ وچ سچا اے۔}} \\
      \hline \hline
      \emph{ہر جہان توں رسائی والا}: $\forall u \exists v Ruv$ & \hbox to.43\textwidth
           {$\Box p \lif \Diamond p$ \hfill (\Ax{D})} \\
      \hline
      \emph{انعکاسی}: $\forall w Rww$  
      & \hbox to.43\textwidth
          {$\Box p \lif p$ \hfill (\Ax{T})} \\
      \hline
      \emph{متقابل}: &  \hbox to .43\textwidth
           {$p \lif \Box\Diamond p$ \hfill (\Ax{B})} \\
           $\forall u\forall v(Ruv \lif Rvu)$ & \\
      \hline
      \emph{متعدی}: & \hbox to.43\textwidth
           {$\Box p \lif \Box \Box p$ \hfill (\Ax{4})} \\
      $\forall u \forall v \forall w ((Ruv \land Rvw) \lif Ruw)$ & \\
      \hline 
      \emph{اقلیدسی}: & \hbox to .43\textwidth
        {$\Diamond p \lif \Box\Diamond p$ \hfill (\Ax{5})} \\
      $\forall w \forall u \forall v ((Rwu \land Rwv) \lif Ruv)$ &\\
      \hline
    \end{tabular}
    \caption{مطابقت دیاں پنج گلاں۔}
    \ollabel{tab:five}
  \end{table}

\begin{proof}
  \Ax{B} والا معاملہ ایہ اے۔ سانچہ کسے نمونے وچ سچا وکھاؤن لئی سانوں وکھاؤنا اے کہ اوہدیاں ساریاں مثالاں اوس نمونے دے سارے جہاناں وچ سچیاں نیں۔ منّ لو $!A \lif \Box\Diamond !A$، \Ax{B} دی دِتّی ہوئی مثال اے تے $w \in W$ کوئی وی جہان اے۔ منّ لو مقدم $!A$، $w$ وچ سچا اے، تاں جو وکھائیے کہ $\Box \Diamond
  !A$، $w$ وچ سچا اے۔ سانوں وکھاؤنا اے کہ $\Diamond !A$ سارے اوہناں $w'$ وچ سچا اے جیہڑے $w$ توں رسائی وچ نیں۔ ہُن جے کسے $w'$ لئی $Rww'$ اے، تاں تقابل دی شرط توں $Rw'w$ وی اے (\olref{fig:Bsymm} ویکھو)۔ کیوں جے $\mSat{M}{!A}[w]$، ایس لئی
  $\mSat{M}{\Diamond !A}[w']$۔ کیوں جے $w'$ کوئی وی ایسا جہان سی جیہدے لئی $Rww'$ اے، ایس لئی $\mSat{M}{\Box\Diamond!A}[w]$۔

  باقی معاملے مشقاں لئی چھڈدے آں۔
\end{proof}

\begin{prob}
  \olref[nml][frd][acc]{thm:soundschemas} دا ثبوت پورا کرو۔
\end{prob}

\begin{figure}
  \begin{center}
    \begin{tikzpicture}[modal]
      \node[world] (w1) [label={below:$\mSat{{}}{\formula{A}}$\\
          $\mSat{{}}{\Box\Diamond \formula{A}}$}] {$w$}; 
      \node[world] (w2) [label=below:{$\mSat{{}}{\Diamond \formula{A}}$},
        right=of w1] {$w'$}; 
      \draw[->, bend left] (w1) to (w2); 
      \draw[->, bend left] (w2) to (w1); 
    \end{tikzpicture}
  \end{center}
\caption{تقابل اُتے مبنی دلیل۔}
\ollabel{fig:Bsymm}
\end{figure}

دھیان دیو کہ \olref{thm:soundschemas} دے اُلٹے لزوم نہیں چلدے۔ ایہ گل سچی نہیں کہ جے کوئی نمونہ سانچے نوں پورا کردا اے، تاں اوس نمونے دے رسائی دے تعلق وچ سانچے نال جڑی خاصیت اے۔ \Ax{T} تے انعکاسی نمونیاں دے معاملے وچ اک نمونہ سوکھا لبھ جاندا اے جیہدے وچ \Ax{T} آپ جھوٹھا اے۔ منّ لو $W = \{w\}$ تے
$V(p) = \emptyset$۔ فیر $R$ انعکاسی نہیں، پر $\mSat{M}{\Box
  p}[w]$ تے $\mSat/{M}{p}[w]$۔ پر ایتھے \Ax{T} دی صرف اک مثال~$\mModel{M}$ وچ جھوٹھھی اے؛ ہور مثالاں، جیویں
$\Box \lnot p \lif \lnot p$، سچیاں نیں۔ ایسی مثال لبھنی اوکھی اے جیہدے وچ~\Ax{T} دی \emph{بدل نال بنی ہر مثال}~$\mModel{M}$ وچ سچی ہووے تے $\mModel{M}$ انعکاسی نہ ہووے۔ پر ایسے نمونے وی نیں۔

\begin{prop}\ollabel{prop:reflexive}
  منّ لو $\mModel{M} = \tuple{W, R, V}$ ایسا نمونہ اے جیہدے لئی $W = \{u, v
  \}$ اے، جتھے جہان $u$ تے $v$، $R$ نال جُڑے نیں، یعنی $Ruv$
  تے $Rvu$ دونیں۔ منّ لو ہر $p$ لئی $u \in V(p) \Leftrightarrow v
  \in V(p)$۔ فیر
  \begin{enumerate}
  \item ہر $!A$ لئی $\mSat{M}{!A}[u]$ اُدوں تے صرف اُدوں جدوں
    $\mSat{M}{!A}[v]$ ($!A$ اُتے استقرا ورتو)۔
  \item ~\Ax{T} دی ہر مثال $\mModel{M}$ وچ سچی اے۔
  \end{enumerate}
  کیوں جے $\mModel{M}$ انعکاسی نہیں (اصل وچ ایہ \emph{اپنے آپ نال تعلق توں خالی} اے)، ایس لئی \olref{thm:soundschemas} دا اُلٹا رخ \Ax{T} دے معاملے وچ نہیں چلدا (\olref{thm:soundschemas} وچ دسے ہور سانچیاں وچوں کجھ لئی ایہو جیہی دلیل دِتّی جا سکدی اے، پر ساریاں لئی نہیں)۔
\end{prop}

\begin{prob}
  \olref[nml][frd][acc]{prop:reflexive} دے دعوے ثابت کرو۔
\end{prob}

ساڈی توجہ پنج کلاسیکی !!{formula}s \Ax{D}،
\Ax{T}، \Ax{B}، \Ax{4}، تے \Ax{5} اُتے رہے گی، پر
\olref{tab:anotherfive} وچ رسائی دے تعلقاں دیاں کجھ ہور خاصیتاں وی دسدے آں۔ رسائی دا تعلق~$R$ جزوی فنکشنی اے جے ہر جہان توں ودھ توں ودھ اک جہان رسائی وچ اے۔ جے ہر جہان توں ٹھیک اک جہان رسائی وچ اے، تاں اینوں فنکشنی آکھدے آں۔ (ایس طرح فنکشنی تعلق اوہی نیں جیہڑے ہر جہان توں رسائی والے وی نیں تے جزوی فنکشنی وی۔) «فنکشنی» ایس لئی آکھدے آں کہ رسائی دا تعلق اک (جزوی) فنکشن وانگ کم کردا اے۔ تعلق کمزور طور تے گنجان اے جے جدوں وی $Ruv$ ہووے، کوئی~$w$، $u$ تے~$v$ دے «وچکار» ہووے۔ ایس معنی وچ کمزور طور تے گنجان تعلق متعدی تعلق دا اُلٹ اے۔ متعدی تعلق وچ جدوں وی $v$ تک~$u$ توں~$w$ راہیں پہنچ سکدے آں، تاں $v$ تک $u$ توں سدھا وی پہنچ سکدے آں؛ کمزور طور تے گنجان تعلق وچ جدوں وی $v$ تک~$u$ توں سدھا پہنچ سکدے آں، تاں کسے~$w$ راہیں وی پہنچ سکدے آں۔ تعلق کمزور طور تے ہم رُخ اے جے جدوں وی جہان $u$ تے~$v$ تک کسے جہان~$w$ توں پہنچ سکدے آں، تاں اکّو جہان~$t$ تک $u$ تے~$v$ دوناں توں پہنچ سکدے آں۔ اینوں کدی «ہیرے والی خاصیت» یا «راہواں دا سنگم» وی آکھدے نیں۔

\begin{table}[t]
    \begin{tabular}{| p{.50\textwidth} || p{.43\textwidth} |}
      \hline
      {\emph{جے $R$ \dots اے}} & {\emph{تاں \dots،~$\mModel{M}$ وچ سچا اے۔}} \\
      \hline \hline
      \emph{جزوی فنکشنی}: & 
      \multirow{2}{*}{$\Diamond p \lif \Box p$} \\
      $ \forall w \forall u \forall v ((Rwu \land Rwv) \lif u =v )$ & \\
      \hline
      \multirow{2}{*}{\emph{فنکشنی}: $\forall w \exists v \forall u(Rwu \liff \eq[u][v]) $} 
      & \multirow{2}{*}{$\Diamond p \liff \Box p$} \\
      & \\
      \hline
      \emph{کمزور طور تے گنجان}: &  \multirow{2}{*}{$\Box \Box p \lif
        \Box p$} \\
      $\forall u\forall v(Ruv \lif \exists w(Ruw \land Rwv))$ & \\
      \hline
      \emph{کمزور طور تے جُڑیا}: &
      \multirow{3}{*}{\hbox to.43\textwidth{$
        \begin{array}{@{}l@{}}
          \Box ((p \land \Box p) \lif q) \lor {} \\
          \qquad \Box ((q \land \Box q) \lif p)
        \end{array}$ \hfill (\Ax{L})}
      } \\
      $\forall w \forall u \forall v ((Rwu \land Rwv) \lif {}$ &\\
      \qquad $(Ruv \lor u=v \lor Rvu))$ & \\
      \hline 
      \emph{کمزور طور تے ہم رُخ}: & \multirow{3}{*}{\hbox to .43\textwidth
        {$\Diamond\Box p \lif \Box\Diamond p$\hfill (\Ax{G})}} \\
      $\forall w \forall u \forall v ((Rwu \land Rwv) \lif {}$ & \\
      \qquad $\exists t (Rut \land Rvt))$ & \\
      \hline
    \end{tabular}
    \caption{مطابقت دیاں پنج ہور گلاں۔}
    \ollabel{tab:anotherfive}
  \end{table}

\begin{prob}
  منّ لو $\mModel{M} = \tuple{W, R, V}$ اک نمونہ اے۔ وکھاؤ کہ جے $R$،
  \olref[nml][frd][acc]{tab:anotherfive} دے کھبے پاسے دِتّی خاصیت پوری کردا اے، تاں سجے پاسے اوہدے نال جڑے فارمولے دی ہر مثال~$\mModel{M}$ وچ سچی اے۔
\end{prob}
\endgroup
% /OLP-0421
% OLP-0422 content/normal-modal-logic/frame-definability/frames.tex
\begingroup
\def\olfilename{frames}
\expandafter\def\csname ol@sectioncs\endcsname{section}

\olfileid{nml}{frd}{fra}

\olsection{فریم}

\begin{defn}
  \emph{فریم} اک جوڑا $\mModel{F} = \tuple{W,R}$ اے جتھے $W$ جہاناں دا غیر خالی سیٹ اے تے $R$،~$W$ اُتے دو رکنی تعلق اے۔ نمونہ
  $\mModel{M}$ فریم $\mModel{F} = \tuple{W,R}$ \emph{اُتے ٹکیا} اے اُدوں تے صرف اُدوں جدوں $\mModel{M} = \tuple{W, R, V}$ کسے قدر بندی $V$ لئی ہووے۔
\end{defn}

\begin{defn}
  جے $\mModel{F}$ فریم اے، تاں اسیں آکھدے آں کہ $!A$، \emph{
    $\mModel{F}$ وچ معتبر اے}، علامتاں وچ $\mModel{F} \Entails !A$، جے $\mSat{M}{!A}$ ہر اوس نمونے~$\mModel{M}$ لئی ہووے جیہڑا~$\mModel{F}$ اُتے ٹکیا اے۔

  جے $\mClass{F}$ فریماں دی جماعت اے، تاں اسیں آکھدے آں کہ $!A$، \emph{
    $\mClass{F}$ وچ معتبر اے}، علامتاں وچ $\mClass{F} \Entails !A$، اُدوں تے صرف اُدوں جدوں $\mModel{F} \Entails
  !A$ ہر فریم~$\mModel{F} \in \mClass{F}$ لئی ہووے۔
\end{defn}

فریماں دی خاص گل ایہ اے کہ سانچیاں تے رسائی دے تعلق~$R$ دیاں خاصیتاں دی مطابقت فریماں دی سطح اُتے ہوندی اے، \emph{نمونیاں دی سطح اُتے نہیں}۔ مثال لئی \Ax{T} سارے انعکاسی نمونیاں وچ سچا اے، پر ہر اوہ نمونہ جیہدے وچ \Ax{T} سچا اے، انعکاسی نہیں۔ البتہ ایہ گل \emph{ضرور} سچی اے کہ \Ax{T} \emph{معتبر} اے سارے انعکاسی \emph{فریماں} وچ، تے ہر اوہ فریم وی جیہدے وچ \Ax{T} معتبر اے، انعکاسی اے۔

\begin{rem}
فریماں دی جماعت وچ منطقی اعتبار، نمونیاں دی جماعت وچ منطقی اعتبار دا اک خاص معاملہ اے۔ $\mClass{F} \Entails !A$ اُدوں تے صرف اُدوں جدوں
$\mClass{C} \Entails !A$، جتھے $\mClass{C}$ اوہناں سارے نمونیاں دی جماعت اے جیہڑے~$\mClass{F}$ دے کسے فریم اُتے ٹکے نیں۔

صاف اے کہ جے کوئی !!a{formula} یا سانچہ معتبر اے، یعنی \emph{سارے} نمونیاں دی جماعت دے لحاظ نال معتبر اے، تاں اوہ فریماں دی کسے وی جماعت~$\mClass{F}$ دے لحاظ نال وی معتبر اے۔
\end{rem}
\endgroup
% /OLP-0422
% OLP-0423 content/normal-modal-logic/frame-definability/definability.tex
\begingroup
\def\olfilename{definability}
\expandafter\def\csname ol@sectioncs\endcsname{section}

\olfileid{nml}{frd}{def}

\olsection{فریماں دیاں جماعتاں دی تعریف}

\olref[acc]{thm:soundschemas} دے اُلٹے لزوم بھاویں نہیں چلدے، پر جے «نمونے» دی تھاں «فریم» رکھیے تاں اوہ چلدے نیں۔ \olref[acc]{thm:soundschemas} وچ ویکھیاں خاصیتاں لئی ایہ گل \emph{ضرور} سچی اے کہ جے کوئی !!a{formula} کسے \emph{فریم} وچ معتبر اے، تاں اوس فریم دے رسائی دے تعلق وچ اوہدے نال جڑی خاصیت اے۔ ایس طرح ایہ !!{formula}s اوہناں فریماں دیاں جماعتاں \emph{متعین کردے نیں} جیہناں وچ اوہ خاصیت اے۔

\begin{defn}
  جے $\mClass{F}$ فریماں دی جماعت اے، تاں اسیں $!A$ نوں~$\mClass{F}$ دا \emph{متعین کرن والا فارمولا} آکھدے آں اُدوں تے صرف اُدوں جدوں $\mModel{F} \Entails !A$ سارے اوہناں فریماں~$\mModel{F} \in \mClass{F}$ لئی ہووے تے صرف اوہناں لئی۔
\end{defn}

ہُن اسیں فریماں دی تعریف ہو سکّن بارے پورے نتیجے ثابت کرنے آں۔

\begin{thm}\ollabel{thm:fullCorrespondence}
جے اوہ !!{formula} جیہڑا \olref[acc]{tab:five} دے سجے پاسے اے کسے فریم~$\mModel{F}$ وچ معتبر اے، تاں $\mModel{F}$ وچ کھبے پاسے والی خاصیت اے۔
\end{thm}

\begin{proof}
  \begin{enumerate}
  \item منّ لو \Ax{D}، $\mModel{F} = \tuple{W, R}$ وچ معتبر اے، یعنی
    $\mModel{F} \Entails \Box p \lif \Diamond p$۔ $\mModel{M} =
    \tuple{W, R, V}$ کوئی نمونہ لے لو جیہڑا~$\mModel{F}$ اُتے ٹکیا اے، تے $w \in
    W$ لے لو۔ سانوں وکھاؤنا اے کہ کوئی $v$ اے جیہدے لئی~$Rwv$ ہووے۔ منّ لو ایہ گل نہیں۔ فیر $\mSat{M}{\Box !A}$ تے $\mSat/{M}{\Diamond !A}[w]$ دونیں، کسے وی~$!A$ لئی،~$p$ سمیت۔ پر فیر $\mSat/{M}{\Box p \lif
      \Diamond p}[w]$، جیہڑا مفروضے $\mModel{F}
    \Entails \Box p \lif \Diamond p$ دے اُلٹ اے۔
  \item منّ لو \Ax{T}، $\mModel{F}$ وچ معتبر اے، یعنی $\mModel{F} \Entails \Box
    p \lif p$۔ کوئی وی جہان $w \in W$ لے لو۔ سانوں $Rww$ وکھاؤنا اے۔ $u \in V(p)$ اُدوں تے صرف اُدوں ہووے جدوں $Rwu$ (جدوں $q$،
    $p$ توں وکھرا اے، تاں $V(q)$ کوئی وی ہو سکدا اے، مثلاً $V(q) = \emptyset)$۔ منّ لو
    $\mModel{M} = \tuple{W, R, V}$۔ تعریف مطابق ہر اوہ $u$ جیہدے لئی
    $Rwu$ ہووے، اوہدے لئی $\mSat{M}{p}[u]$، ایس لئی $\mSat{M}{\Box p}[w]$۔ پر مفروضے مطابق $\Box p \lif p$،~$w$ وچ سچا اے، ایس لئی
    $\mSat{M}{p}[w]$، پر $V$ دی تعریف مطابق ایہ صرف اُدوں ہو سکدا اے جدوں~$Rww$ ہووے۔
  \item اُلٹی نفی نال ثابت کرنے آں۔ منّ لو $\mModel{F}$ متقابل نہیں۔ وکھاؤنے آں کہ \Ax{B}، یعنی $p \lif \Box\Diamond p$،
    $\mModel{F}= \tuple{W, R}$ وچ معتبر نہیں۔ جے $\mModel{F}$ متقابل نہیں تاں ایسے $u$، $v \in W$ نیں جیہناں لئی $Ruv$ اے پر
    $Rvu$ نہیں۔ $V$ ایویں متعین کرو کہ $w \in V(p)$ اُدوں تے صرف اُدوں جدوں $Rvw$ نہ ہووے
    (ہور تھاں $V$ کوئی وی ہو سکدا اے)۔ منّ لو $\mModel{M} =\tuple{W, R,
      V}$۔ ہُن $V$ دی تعریف مطابق $\mSat{M}{p}[w]$ ہر ایسے $w$ لئی اے جیہدے لئی $Rvw$ نہیں۔ خاص طور تے $\mSat{M}{p}[u]$، کیوں جے
    ~$Rvu$ نہیں۔ نالے کیوں جے $Rvw$ اُدوں تے صرف اُدوں جدوں $w \notin V(p)$، ایس لئی کوئی~$w$ نہیں جیہدے لئی $Rvw$ تے $\mSat{M}{p}[w]$ دونیں ہون، تے ایس لئی
    $\mSat/{M}{\Diamond p}[v]$۔ کیوں جے~$Ruv$، ایس لئی
    $\mSat/{M}{\Box\Diamond p}[u]$ وی۔ ایس توں $\mSat/{M}{p \lif
      \Box\Diamond p}[u]$ نکلدا اے، ایس لئی $\Ax{B}$،~$\mModel{F}$ وچ معتبر نہیں۔
  \item منّ لو \Ax{4}، $\mModel{F} = \tuple{W,R}$ وچ معتبر اے، یعنی
    $\mModel{F} \Entails \Box p \lif \Box\Box p$، تے $u$، $v$، $w \in W$ کوئی وی ایسے جہان نیں جیہناں لئی $Ruv$ تے $Rvw$ ہون۔ سانوں
    $Ruw$ وکھاؤنا اے۔ $V$ ایویں متعین کرو کہ $z \in V(p)$ اُدوں تے صرف اُدوں جدوں $Ruz$ ہووے (ہور تھاں
    $V$ کوئی وی ہو سکدا اے)۔ منّ لو $\mModel{M} =\tuple{W, R, V}$۔
    $V$ دی تعریف مطابق $\mSat{M}{p}[z]$ ہر اوس $z$ لئی جیہدے لئی $Ruz$، تے ایس لئی $\mSat{M}{\Box p}[u]$۔ پر مفروضے مطابق \Ax{4}، یعنی $\Box p
    \lif \Box \Box p$، $u$ وچ سچا اے، ایس لئی $\mSat{M}{\Box \Box
      p}[u]$۔ کیوں جے $Ruv$ تے $Rvw$، ایس لئی $\mSat{M}{p}[w]$، پر~$V$ دی تعریف مطابق ایہ صرف اُدوں ہو سکدا اے جدوں $Ruw$ ہووے، جیویں وکھاؤنا سی۔
  \item اُلٹی نفی ورتدے آں۔ منّ لو فریم
    $\mModel{F} = \tuple{W, R}$ اقلیدسی نہیں، تے وکھاؤنے آں کہ ایہ~\Ax{5} نوں جھوٹھا کردا اے، یعنی $\mModel{F} \Entails/ \Diamond p \lif
      \Box\Diamond p$۔ منّ لو ایسے جہان $u$، $v$، $w \in W$ نیں جیہناں لئی $Rwu$ تے $Rwv$ نیں پر~$Ruv$ نہیں۔ $V$ ایویں متعین کرو کہ سارے جہاناں $z$ لئی $z \in V(p)$ اُدوں تے صرف اُدوں جدوں~ایہ گل \emph{نہ} ہووے کہ $Ruz$۔ منّ لو $\mModel{M} =\tuple{W, R, V}$۔ فیر مفروضے مطابق $\mSat{M}{p}[v]$ تے کیوں جے $Rwv$، ایس لئی
    $\mSat{M}{\Diamond p}[w]$ وی۔ پر کوئی ایسا جہان $y$ نہیں جیہدے لئی $Ruy$ تے $\mSat{M}{p}[y]$ ہووے، ایس لئی $\mSat/{M}{\Diamond
      p}[u]$۔ کیوں جے $Rwu$، ایس توں $\mSat/{M}{\Box\Diamond
      p}[w]$، سو \Ax{5}، یعنی $\Diamond p \lif \Box\Diamond p$،~$w$ وچ جھوٹھا اے۔
  \end{enumerate}
\end{proof}

تہانوں \Ax{D} دے ثبوت تے ہور معاملیاں وچ اک فرق دِسے گا۔ قدر بندی~$V$ دا ذکر نہیں ہویا۔ اصل وچ اسیں ثابت کیتا کہ جے $\mSat{M}{\Ax{D}}$، تاں $\mModel{M}$ ہر جہان توں رسائی والا اے۔ ایس طرح \Ax{D} ہر جہان توں رسائی والے \emph{نمونیاں} دی جماعت متعین کردا اے، صرف فریماں دی جماعت نہیں۔

\begin{cor}\ollabel{prop:D-serial}
  جیہڑے وی نمونے وچ \Ax{D} سچا اے، اوہ ہر جہان توں رسائی والا اے۔
\end{cor}

\begin{cor}
ہر !!{formula} جیہڑا \olref[acc]{tab:five} دے سجے پاسے اے، اوہناں فریماں دی جماعت متعین کردا اے جیہناں وچ کھبے پاسے والی خاصیت اے۔
\end{cor}

\begin{proof}
  \olref[acc]{thm:soundschemas} وچ اسیں ثابت کیتا کہ جے کسے نمونے وچ کھبے پاسے والی خاصیت اے، تاں سجے پاسے والا !!{formula} اوس وچ سچا اے۔ سو جے فریم~$\mModel{F}$ وچ کھبے پاسے والی خاصیت اے، تاں سجے پاسے والا !!{formula}،~$\mModel{F}$ وچ معتبر اے۔
  \olref{thm:fullCorrespondence} وچ اسیں اُلٹے رخ دے لزوم ثابت کیتے۔ جے سجے پاسے والا کوئی !!a{formula}،~$\mModel{F}$ وچ معتبر اے، تاں $\mModel{F}$ وچ کھبے پاسے والی خاصیت اے۔
\end{proof}

\begin{prob}
وکھاؤ کہ جے اوہ !!{formula} جیہڑا
\olref[nml][frd][acc]{tab:anotherfive} دے سجے پاسے اے، کسے فریم~$\mModel{F}$ وچ معتبر اے، تاں
$\mModel{F}$ وچ کھبے پاسے والی خاصیت اے۔ ایس لئی اک ایسا فریم ویکھو جیہڑا کھبے پاسے والی خاصیت پوری \emph{نہ} کردا ہووے، تے کوئی مناسب~$V$ متعین کرو جیہدے نال سجے پاسے والا !!{formula} کسے جہان وچ جھوٹھا ہووے۔
\end{prob}

\olref{thm:fullCorrespondence} ایہ وی وکھاؤندا اے کہ خاصیتاں اکٹھیاں ہو سکدیاں نیں۔ مثلاً جے \Ax{B} تے \Ax{4} دونیں~$\mModel{F}$ وچ معتبر نیں، تاں فریم متقابل وی اے تے متعدی وی۔ ہور خاصیتاں لئی وی ایویں اے۔ کئی اہم امکانی منطقاں اوہناں !!{formula}s دے سیٹ وجوں پچھانیاں جاندیاں نیں جیہڑے کجھ فریمی خاصیتاں اکٹھیاں رکھن والے سارے فریماں وچ معتبر نیں۔ سو اسیں اینہاں منطقاں نوں اوہناں !!{formula}s دے سیٹ وجوں وی پچھان سکدے آں جیہڑے اوہناں سارے فریماں وچ معتبر نیں جیہناں وچ خاصیتاں نوں متعین کرن والے متعلقہ !!{formula}s معتبر نیں۔ مثال لئی کلاسیکی نظام \Log{S4} سارے اوہناں !!{formula}s دا سیٹ اے جیہڑے سارے انعکاسی تے متعدی فریماں وچ معتبر نیں، یعنی اوہناں سارے فریماں وچ جیہناں وچ \Ax{T} تے~\Ax{4} دونیں معتبر نیں۔ \Log{S5} اوہناں سارے فارمولیاں دا سیٹ اے جیہڑے سارے انعکاسی، متقابل تے اقلیدسی فریماں وچ معتبر نیں، یعنی جیہناں فریماں وچ \Ax{T}، \Ax{B} تے \Ax{5} سارے معتبر نیں۔

عام طور تے $R$ دیاں خاصیتاں دے منطقی تعلق اوہناں نوں متعین کرن والے !!{formula}s دے آپسی تعلقاں نال میل کھاندے نیں۔ مثال لئی ہر انعکاسی تعلق ہر جہان توں رسائی والا اے۔ ایس لئی جدوں وی \Ax{T} کسے فریم وچ معتبر اے،~\Ax{D} وی معتبر اے۔ (دھیان دیو کہ ایہ منطقی نتیجے دا تعلق \emph{نہیں}۔ ایہ گل نہیں کہ جدوں وی
$\mSat{M}{\Ax{T}}[w]$، تاں $\mSat{M}{\Ax{D}}[w]$۔) ایہو جیہے کجھ ہور تعلق لکھدے آں۔

\begin{prop}\ollabel{prop:relation-facts}
  منّ لو $R$ سیٹ $W$ اُتے دو رکنی تعلق اے۔ فیر
  \begin{enumerate}
  \item جے $R$ انعکاسی اے، تاں ہر جہان توں رسائی والا اے۔
  \item جے $R$ متقابل اے، تاں ایہ اُدوں تے صرف اُدوں متعدی اے جدوں اقلیدسی ہووے۔
  \item جے $R$ متقابل یا اقلیدسی اے، تاں ایہ کمزور طور تے ہم رُخ اے (ایہدے وچ «ہیرے والی خاصیت» اے)۔
  \item جے $R$ اقلیدسی اے، تاں ایہ کمزور طور تے جُڑیا اے۔
  \item جے $R$ فنکشنی اے، تاں ایہ ہر جہان توں رسائی والا اے۔
  \end{enumerate}
\end{prop}

\begin{prob}
  \olref[nml][frd][def]{prop:relation-facts} ثابت کرو۔
\end{prob}
\endgroup
% /OLP-0423
% OLP-0424 content/normal-modal-logic/frame-definability/first-order-definability.tex
\begingroup
\def\olfilename{first-order-definability}
\expandafter\def\csname ol@sectioncs\endcsname{section}

\olfileid{nml}{frd}{fol}

\olsection{پہلے درجے دی زبان وچ تعریف}

اسیں ویکھیا اے کہ فریماں دے رسائی دے تعلقاں دیاں کئی خاصیتاں امکانی !!{formula}s نال متعین ہو سکدیاں نیں۔ مثال لئی فریماں دا تقابل !!{formula}~\Ax{B}، یعنی $p \lif \Box\Diamond
p$، نال متعین ہو سکدا اے۔ ہُن تک جیہڑیاں شرطاں ویکھیاں نیں اوہ ساریاں پہلے درجے دے !!{formula}s نال ایسی !!a{language} وچ دَسیاں جا سکدیاں نیں جیہدے وچ صرف اک دو رکنی !!{predicate} ہووے۔ مثال لئی تقابل
$\lforall[x][\lforall[y][(\Atom{Q}{x,y} \lif \Atom{Q}{y, x})]]$ نال متعین ہوندا اے، ایس معنی وچ کہ پہلے درجے دی !!{structure}~$\Struct{M}$، جیہدے وچ $\Domain{M} =
W$ تے $\Assign{Q}{M} = R$ اے، اُپرلا فارمولا اُدوں تے صرف اُدوں پورا کردی اے جدوں $R$ متقابل ہووے۔ ایس توں ہِیٹھلی تعریف دا خیال آؤندا اے۔

\begin{defn}
  فریماں دی جماعت~$\mClass{F}$ \emph{پہلے درجے دی زبان وچ متعین ہو سکدی اے} جے کوئی ایسا !!a{sentence}~$!A$ ہووے، پہلے درجے دی اوس زبان وچ جیہدے وچ اک دو رکنی !!{predicate}~$Q$ ہووے، تے ایہ شرط پوری ہوندی ہووے کہ $\mModel{F} =
  \tuple{W, R} \in \mClass{F}$ اُدوں تے صرف اُدوں جدوں $\Sat{M}{!A}$ ہووے، پہلے درجے دی اوس !!{structure}~$\Struct{M}$ وچ جیہدے لئی $\Domain{M} = W$ تے $\Assign{Q}{M}
  = R$۔
\end{defn}

پتا لگدا اے کہ ہُن تک ویکھیاں خاصیتاں تے اوہناں نوں متعین کرن والے امکانی !!{formula}s خاص معاملے نیں۔ کجھ !!{formula} ایسے نیں جیہڑے فریماں دی ایسی جماعت متعین کردے نیں جیہڑی پہلے درجے دی زبان وچ متعین نہیں ہو سکدی۔ تے فریماں دیاں کجھ اوہ جماعتاں، جیہڑیاں پہلے درجے دی زبان وچ متعین ہو سکدیاں نیں، کسے امکانی فارمولے نال متعین نہیں ہو سکدیاں۔

پہلی گل دی ردّی مثال لوب دا فارمولا اے۔
\begin{equation}
\Box(\Box p \lif p) \lif \Box p. \tag{\Ax{W}}
\end{equation}
\Ax{W} اوہناں فریماں دی جماعت متعین کردا اے جیہڑے متعدی تے اُلٹے رخ توں پکّی بنیاد والے نیں۔ تعلق پکّی بنیاد والا اے جے کوئی لامتناہی سلسلہ $w_1$، $w_2$، \dots{} ایسا نہ ہووے جیہدے لئی $Rw_2w_1$، $Rw_3w_2$، \dots۔ مثلاً تعلق $<$،~$\Nat$ اُتے پکّی بنیاد والا اے، جدکہ تعلق $<$،~$\Int$ اُتے پکّی بنیاد والا نہیں۔ تعلق اُلٹے رخ توں پکّی بنیاد والا اُدوں تے صرف اُدوں اے جدوں اوہدا اُلٹا تعلق پکّی بنیاد والا ہووے۔ سو اُلٹے رخ توں پکّی بنیاد والے تعلق اوہ نیں جیہناں وچ کوئی لامتناہی سلسلہ $w_1$، $w_2$، \dots{} ایسا نہیں ہوندا جیہدے لئی $Rw_1w_2$، $Rw_2w_3$، \dots۔

پر پہلے درجے دا کوئی !!{formula} نہیں جیہڑا متعدی تے اُلٹے رخ توں پکّی بنیاد والے تعلقاں نوں متعین کرے۔ منّ لو~$\Sat{M}{!F}$ اُدوں تے صرف اُدوں جدوں $R =
\Assign{Q}{M}$ متعدی تے اُلٹے رخ توں پکّی بنیاد والا ہووے۔ منّ لو $!A_n$ ایہ !!{formula} اے۔
\[
(\Atom{Q}{a_1,a_2} \land \dots \land
    \Atom{Q}{a_{n-1},a_{n}})
\]
ہُن ایہ !!{formula}s دا سیٹ ویکھو۔
\[
\Gamma = \{!F, !A_1, !A_2, \dots\}.
\]
$\Gamma$ دا ہر محدود ذیلی سیٹ پورا ہون جوگا اے۔ سبھ توں وڈا $k$ لے لو جیہدے لئی $!A_k$ ذیلی سیٹ وچ اے، تے $\Domain{M_k} = \{1, \dots, k\}$،
$\Assign{a_i}{M_k} = i$ تے $\Assign{Q}{M_k} = <$ رکھو۔ کیوں جے $<$، $\{1,
\dots, k\}$ اُتے متعدی تے اُلٹے رخ توں پکّی بنیاد والا اے، ایس لئی
$\Sat{M_k}{!F}$۔ بناوٹ مطابق $\Sat{M_k}{!A_i}$ سارے $i \le k$ لئی۔ پہلے درجے دی منطق دے محدود گواہی دے قضیے مطابق $\Gamma$ کسے !!{structure}~$\Struct{M}$ وچ پورا ہون جوگا اے۔ مفروضے مطابق، کیوں جے
$\Sat{M}{!F}$، ایس لئی تعلق $\Assign{Q}{M}$ اُلٹے رخ توں پکّی بنیاد والا اے۔ پر صاف اے کہ $\Assign{a_1}{M}$، $\Assign{a_2}{M}$،
\dots{} اوہو جیہا لامتناہی سلسلہ بناؤندے نیں جیہڑا اُلٹے رخ توں پکّی بنیاد ہون دی شرط منع کردی اے۔

دوجے دعوے دی ردّی مثال ہمہ گیری دی خاصیت اے۔ ہر $u$ تے $v$ لئی $Ruv$۔ ہمہ گیر فریم پہلے درجے دے فارمولے
$\lforall[x][\lforall[y][\Atom{Q}{x,y}]]$ نال متعین ہوندے نیں۔ پر کوئی امکانی فارمولا ایہو جیہا نہیں جیہڑا سارے ہمہ گیر فریماں وچ، تے صرف اوہناں وچ، معتبر ہووے۔ ایہ اک نتیجے توں نکلدا اے جیہڑا آپ وی دلچسپ اے۔ ہمہ گیر فریماں وچ معتبر فارمولے عین اوہی نیں جیہڑے انعکاسی، متقابل تے متعدی فریماں وچ معتبر نیں۔ ایسے انعکاسی، متقابل تے متعدی فریم نیں جیہڑے ہمہ گیر نہیں۔ ایس لئی ہر اوہ !!{formula} جیہڑا سارے ہمہ گیر فریماں وچ معتبر اے، کجھ غیر ہمہ گیر فریماں وچ وی معتبر اے۔
\endgroup
% /OLP-0424
% OLP-0425 content/normal-modal-logic/frame-definability/equivalence-S5.tex
\begingroup
\def\olfilename{equivalence-S5}
\expandafter\def\csname ol@sectioncs\endcsname{section}

\olfileid{nml}{frd}{es5}

\olsection{ہم ارزیت دے تعلق تے \Log{S5}}

امکانی منطق \Log{S5}، اوہناں !!{formula}s دا سیٹ اے جیہڑے سارے ہمہ گیر فریماں اُتے معتبر نیں۔ یعنی ہر جہان ہر جہان توں، اپنے آپ توں وی، رسائی وچ اے۔ ایس صورت وچ $\Box$ ضرورت نال تے $\Diamond$ امکان نال میل کھاندا اے۔ $\Box !A$ سچا اے جے $!A$ \emph{ہر} جہان وچ سچا اے، تے $\Diamond !A$ سچا اے جے $!A$ \emph{کسے اک} جہان وچ سچا اے۔ پتا لگدا اے کہ \Log{S5} نوں اوہناں !!{formula}s دے سیٹ وجوں وی پچھان سکدے آں جیہڑے سارے انعکاسی، متقابل تے متعدی فریماں اُتے، یعنی سارے \emph{ہم ارزیت دے تعلقاں} اُتے معتبر نیں۔

\begin{defn}
  $R$، سیٹ $W$ اُتے \emph{ہم ارزیت دا تعلق} اُدوں تے صرف اُدوں اے جدوں ایہ دو رکنی تعلق انعکاسی، متقابل تے متعدی ہووے۔ تعلق
  $R$، $W$ اُتے \emph{ہمہ گیر} اُدوں تے صرف اُدوں اے جدوں $Ruv$ ہر $u,v \in
  W$ لئی ہووے۔
\end{defn}

کیوں جے \Ax{T}، \Ax{B} تے \Ax{4} انعکاسی، متقابل تے متعدی فریماں نوں پچھاندے نیں، ایس لئی رسائی دا تعلق اوہناں فریماں وچ ہم ارزیت دا تعلق اے جیہناں وچ ایہ تنّے !!{formula}s معتبر نیں، تے صرف اوہناں وچ۔ پتا لگدا اے کہ ہم ارزیت دے تعلقاں نوں !!{formula}s دیاں ہور ترکیباں نال وی پچھان سکدے آں، کیوں جے ہم ارزیت دے تعلقاں دی تعریف والیاں شرطاں،~$R$ دیاں ہور جانیاں شرطاں دیاں ترکیباں نال ہم ارز نیں۔

\begin{prop}\ollabel{prop:equivalences}
  ہِیٹھ لکھیاں گلاں ہم ارز نیں۔
  \begin{enumerate}
  \item $R$ ہم ارزیت دا تعلق اے؛
  \item $R$ انعکاسی تے اقلیدسی اے؛
  \item $R$ ہر جہان توں رسائی والا، متقابل تے اقلیدسی اے؛
  \item $R$ ہر جہان توں رسائی والا، متقابل تے متعدی اے۔
  \end{enumerate}
\end{prop}

\begin{proof}
  مشق۔
\end{proof}

\begin{prob}
  \olref[nml][frd][es5]{prop:equivalences} ثابت کرن لئی وکھاؤ۔
  \begin{enumerate}
  \item جے $R$ متقابل تے متعدی اے، تاں اقلیدسی اے۔
  \item جے $R$ انعکاسی اے، تاں ہر جہان توں رسائی والا اے۔
  \item جے $R$ انعکاسی تے اقلیدسی اے، تاں متقابل اے۔
  \item جے $R$ متقابل تے اقلیدسی اے، تاں متعدی اے۔
  \item جے $R$ ہر جہان توں رسائی والا، متقابل تے متعدی اے، تاں انعکاسی اے۔
  \end{enumerate}
  سمجھاؤ کہ ایہ شرطاں دی ہم ارزی ثابت کرن لئی کیوں کافی اے۔
\end{prob}

\olref{prop:equivalences}، \olref[prf][prs]{prop:S5} دا معنیاتی جوڑی دار اے، کیوں جے ایہ اوہناں فریماں دی امکانی منطق دی ہم ارز پچھان دیندا اے جیہناں اُتے $R$ ہم ارزیت دا تعلق اے (ایس منطق نوں رواج مطابق \Log{S5} آکھدے آں)۔

ہمہ گیر تعلقاں تے ہم ارزیت دے تعلقاں دا آپس وچ کیہ رشتہ اے؟ ہر ہمہ گیر تعلق ہم ارزیت دا تعلق اے، پر صاف اے کہ ہر ہم ارزیت دا تعلق ہمہ گیر نہیں۔ البتہ سارے ہمہ گیر تعلقاں اُتے معتبر !!{formula}s عین اوہی نیں جیہڑے سارے ہم ارزیت دے تعلقاں اُتے معتبر نیں۔

\begin{prop}
  منّ لو $R$ ہم ارزیت دا تعلق اے۔ ہر $w \in W$ لئی \emph{ہم ارزیت دا طبقہ}، $w$ دا، سیٹ $[w] = \{w'\in W :
  Rww'\}$ متعین کرو۔ فیر
  \begin{enumerate}
  \item $w \in [w]$؛
  \item $R$ ہر ہم ارزیت دے طبقے $[w]$ اُتے ہمہ گیر اے؛
  \item ہم ارزیت دے طبقے $W$ نوں ایسے ذیلی سیٹاں وچ ونڈدے نیں جیہڑے آپس وچ سانجھے رکن نہیں رکھدے تے رل کے سارا سیٹ بنا دیندے نیں۔
  \end{enumerate}
\end{prop}

\begin{prop}\ollabel{prop:S5=univ}
  !!^a{formula}~$!A$ سارے اوہناں فریماں $\mModel{F} =\tuple{W,R}$ وچ معتبر اے جیہناں وچ $R$ ہم ارزیت دا تعلق اے، اُدوں تے صرف اُدوں جدوں اوہ سارے اوہناں فریماں $\mModel{F} =\tuple{W,R}$ وچ معتبر اے جیہناں وچ $R$ ہمہ گیر اے۔ ایس لئی ہمہ گیر فریماں دی منطق \Log{S5} ای اے۔
\end{prop}

\begin{proof}
  ایہ پرکھنا سوکھا اے کہ ہمہ گیر تعلق~$R$،~$W$ اُتے، ہم ارزیت دا تعلق اے۔ ایس لئی جے $!A$ سارے فریماں وچ معتبر اے جیہناں وچ $R$ ہم ارزیت دا تعلق اے، تاں سارے ہمہ گیر فریماں وچ وی معتبر اے۔ دوجے رخ لئی اُلٹی نفی ورتدے آں۔ منّ لو $!B$ کوئی !!a{formula} اے جیہڑا جہان $w$ اُتے اک نمونے $\mModel{M} = \tuple{W, R, V}$ وچ جھوٹھا اے۔ نمونہ فریم $\tuple{W,R}$ اُتے ٹکیا اے جیہدے وچ $R$، $W$ اُتے ہم ارزیت دا تعلق اے۔ سو $\mSat/{M}{!B}[w]$۔ نمونہ $\mModel{M}' = \tuple{W',
    R', V'}$ ایویں متعین کرو۔
  \begin{enumerate}
  \item $W' = [w]$؛
  \item $R'$، $W'$ اُتے ہمہ گیر اے؛
  \item $V'(p) = V(p) \cap W'$۔
  \end{enumerate}
  (سو جہاناں دا سیٹ $W'$، $\mModel{M}'$ وچ، \olref{fig:partition} دے چھاں والے حصے نال وکھایا گیا اے۔) سوکھا دِسدا اے کہ $R$ تے $R'$، $W'$ اُتے اکّو جیہے نیں۔ فیر !!{formula}s اُتے استقرا نال وکھا سکدے آں کہ ہر $w' \in W'$ لئی $\mSat{M'}{!A}[w']$ اُدوں تے صرف اُدوں جدوں $\mSat{M}{!A}[w']$، ہر $!A$ لئی (ایہ گل ٹھیک بن دی اے، کیوں جے $W'
  \subseteq W$)۔ خاص طور تے $\mSat/{M'}{!B}[w]$، تے $!B$ ہمہ گیر فریم اُتے ٹکے نمونے وچ جھوٹھا اے۔
\end{proof}

\begin{figure}[t]
  \centering
  \begin{tikzpicture}[node distance=2cm, auto, thick]
    \clip [rounded corners] (0,0) -- (6,0) -- (6,4) -- (0,4) --  cycle;
    \begin{scope}
      \clip (0,0) -- (2,0)  .. controls (1.5,1.5) and (4.5,1.5)
      .. (4,4) -- (0,4) -- (0,0) -- cycle;
      \filldraw[fill=gray!40] (0,4) -- (0,2) .. controls (1.5,1.5)
      and (4.5,1.5) .. (4.5,0) -- (6,0) -- (6,4) -- (0,4) --  cycle;
    \end{scope}
    \draw [name path=line1] (2,0) .. controls (1.5,1.5) and (4.5,1.5) .. (4,4);
    \draw [name path=line2] (0,2)  .. controls (1.5,1.5) and (4.5,1.5) .. (4.5,0);
    \path [name intersections={of=line1 and line2}];
    \draw [rounded corners, use as bounding box, very thick] (0,0)
    -- (6,0) -- (6,4) -- (0,4) --  cycle; 
    \node at (2,3) {$[w]$} ;
    \node at (1,1) {$[u]$} ;
    \node at (3,0.75) {$[v]$} ;
    \node at (4.75,2) {$[z]$} ;
  \end{tikzpicture}
  \caption{$W$ دی ہم ارزیت دے طبقیاں وچ ونڈ۔}
  \ollabel{fig:partition}
\end{figure}
\endgroup
% /OLP-0425
% OLP-0426 content/normal-modal-logic/frame-definability/second-order-definability.tex
\begingroup
\def\olfilename{second-order-definability}
\expandafter\def\csname ol@sectioncs\endcsname{section}

\olfileid{nml}{frd}{st}

\olsection{دوجے درجے دی زبان وچ تعریف}

فریماں دیاں کجھ خاصیتاں امکانی فارمولیاں نال متعین ہو سکدیاں نیں پر پہلے درجے دی زبان وچ نہیں۔ پر جے اک رکنی محمولاں اُتے مسوّراں دی اجازت دیئے (یعنی اک رکنی دوجے درجے دی تسویر دی)، تاں فریم دیاں ساریاں اوہ خاصیتاں متعین کر سکدے آں جیہڑیاں امکانی فارمولیاں نال متعین ہو سکدیاں نیں۔ ایس لئی اسیں اوس باقاعدہ جوڑ نوں ورتدے آں جیہڑا امکانی !!{formula} دے کسے جہان وچ سچا ہون دیاں شرطاں تے پہلے درجے دے !!{formula}s وچ اے۔ اینوں امکانی فارمولیاں دا پہلے درجے دے فارمولیاں وچ معیاری ترجمہ آکھدے آں۔ ایہ ایسی زبان وچ اے جیہدے وچ رسائی دے تعلق لئی دو رکنی !!{predicate}~$Q$ تے اک رکنی !!{predicate}~$P_i$ وی اے، اوہناں !!{propositional variable}s~$\Obj p_i$ لئی جیہڑے~$!A$ وچ آؤندے نیں۔

\begin{defn}
  \emph{معیاری ترجمہ}~$\ST_x(!A)$ دی تعریف استقرا نال ایویں ہوندی اے۔
  \begin{enumerate}
  \tagitem{prvFalse}{\indcase{!A}{\lfalse}{$\ST_x(\indfrm) = \lfalse$.}}{}
  \tagitem{prvTrue}{\indcase{!A}{\lfalse}{$\ST_x(\indfrm) = \ltrue$.}}{}
  \item \indcase{!A}{\Obj p_i}{$\ST_x(\indfrm) = \Atom{P_i}{x}$.}
  \tagitem{prvNot}{\indcase{!A}{\lnot !B}{$\ST_x(\indfrm) = \lnot \ST_x(!B)$.}}{}
  \tagitem{prvAnd}{\indcase{!A}{(!B \land !C)}{$\ST_x(\indfrm) = (\ST_x(!B)
      \land \ST_x(!C))$.}}{}
  \tagitem{prvOr}{\indcase{!A}{(!B \lor !C)}{$\ST_x(\indfrm) = (\ST_x(!B)
      \lor \ST_x(!C))$.}}{}
  \tagitem{prvIf}{\indcase{!A}{(!B \lif !C)}{$\ST_x(\indfrm) = (\ST_x(!B)
      \lif \ST_x(!C))$.}}{}
  \tagitem{prvIff}{\indcase{!A}{(!B \liff !C)}{$\ST_x(\indfrm) = (\ST_x(!B)
      \liff \ST_x(!C))}$.}{}
  \tagitem{prvBox}{\indcase{!A}{\Box !B}{$\ST_x(\indfrm) =
      \lforall[y][(\Atom{Q}{x,y} \lif \ST_y(!B))]$.}}{}
  \tagitem{prvDiamond}{\indcase{!A}{\Diamond !B}{$\ST_x(\indfrm) =
      \lexists[y][(\Atom{Q}{x,y} \land \ST_y(!B))]$.}}{}
  \end{enumerate}
\end{defn}

مثلاً $\ST_x(\Box p \lif p)$ دا ترجمہ $\lforall[y][(\Atom{Q}{x,y}
  \lif \Atom{P}{y})] \lif \Atom{P}{x}$ اے۔ ہر !!{structure} جیہڑی~$\ST_x(!A)$ دی زبان لئی اے، اوہدے لئی اک دائرہ،~$Q$ نوں دِتّا دو رکنی تعلق، تے اک رکنی !!{predicate}s~$P_i$ نوں دِتّے دائرے دے ذیلی سیٹ چاہیدے نیں۔ دوجے لفظاں وچ ایسی !!a{structure} دے جز اوہی نیں جیہڑے~$!A$ دے نمونے دے جز نیں۔ دائرہ جہاناں دا سیٹ اے،~$Q$ نوں دِتّا دو رکنی تعلق رسائی دا تعلق اے، تے~$P_i$ نوں دِتّے ذیلی سیٹ اوہی دِتّیاں قدراں~$V(\Obj p_i)$ نیں۔ ایس لئی ایہ حیرانی دی گل نہیں کہ $!A$ دا امکانی نمونے وچ پورا اترنا تے $\ST_x(!A)$ دا اوہدے نال جڑی !!{structure} وچ پورا اترنا اک دوجے نال میل کھاندے نیں۔

\begin{prop}\ollabel{prop:st}
  منّ لو $\mModel{M} = \tuple{W, R, V}$ اے، تے $\Struct{M'}$ پہلے درجے دی !!{structure} اے جیہدے وچ $\Domain{M'} = W$، $\Assign{Q}{M'} = R$ تے
  $\Assign{P_i}{M'} = V(\Obj p_i)$ نیں، تے $s(x) = w$ اے۔ فیر
  \[
  \mSat{M}{!A}[w] \text{ اُدوں تے صرف اُدوں جدوں } \Sat{M'}{\ST_x(!A)}[s]
  \]
\end{prop}

\begin{proof}
  استقرا نال،~$!A$ اُتے۔
\end{proof}

\begin{prop}
  منّ لو $!A$ امکانی !!{formula} اے تے $\mModel{F} = \tuple{W, R}$ فریم اے۔ منّ لو $\Struct{F'}$ پہلے درجے دی !!{structure} اے جیہدے وچ $\Domain{F'} = W$ تے $\Assign{Q}{F'} = R$ اے، تے $!A'$ ایہ دوجے درجے دا فارمولا اے۔
  \[
  \lforall[X_1][\dots\lforall[X_n][\lforall[x][
        \SSubst{\ST_x(!A)}{\subst{X_1}{P_1}, \dots,
          \subst{X_n}{P_n}}]]],
  \]
  جتھے $P_1$، \dots، $P_n$ اوہ سارے اک رکنی !!{predicate}s نیں جیہڑے~$\ST_x(!A)$ وچ نیں۔ فیر
  \[
  \mModel{F} \Entails !A \text{ اُدوں تے صرف اُدوں جدوں } \Sat{F'}{!A'}
  \]
\end{prop}

\begin{proof}
  $\Sat{F'}{!A'}$ اُدوں تے صرف اُدوں جدوں ہر ایسی !!{structure}~$\mModel{M'}$ لئی، جیہدے وچ
  $\Assign{P_i}{M'} \subseteq W$ ہووے $i = 1$، \dots،~$n$ لئی، تے ہر ایسے $s$ لئی جیہدے وچ $s(x) \in W$، اوہدے لئی $\Sat{M'}{\ST_x(!A)}[s]$ ہووے۔
  \olref{prop:st} مطابق ایہ گل اُدوں تے صرف اُدوں ہوندی اے جدوں سارے نمونیاں~$\mModel{M}$ لئی جیہڑے~$\mModel{F}$ اُتے ٹکے نیں، تے ہر جہان~$w \in W$ لئی، $\mSat{M}{!A}[w]$، یعنی $\mModel{F} \Entails !A$۔
\end{proof}

\begin{defn}
  فریماں دی جماعت~$\mClass{F}$ \emph{دوجے درجے دی زبان وچ متعین ہو سکدی اے} جے کوئی ایسا !!a{sentence}~$!A$ ہووے، دوجے درجے دی اوس زبان وچ جیہدے وچ صرف اک دو رکنی !!{predicate}~$P$ ہووے تے مسوّر صرف اک رکنی سیٹ متغیراں اُتے لگن، کہ $\mModel{F} = \tuple{W, R} \in \mClass{F}$ اُدوں تے صرف اُدوں جدوں $\Sat{M}{!A}$ ہووے اوس !!{structure}~$\Struct{M}$ وچ جیہدے لئی
  $\Domain{M} = W$ تے $\Assign{P}{M} = R$۔
\end{defn}

\begin{cor}
  جے فریماں دی جماعت !!a{formula}~$!A$ نال متعین ہو سکدی اے، تاں اوہدے نال جڑی رسائی دے تعلقاں دی جماعت اک رکنی دوجے درجے دے جملے نال متعین ہو سکدی اے۔
\end{cor}

\begin{proof}
  پچھلے ثبوت دے اک رکنی دوجے درجے دے جملے~$!A'$ وچ مطلوب خاصیت اے۔
\end{proof}

مثال لئی فیر !!{formula}~$\Box p \lif p$ ویکھو۔ ایہ انعکاس نوں متعین کردا اے۔ انعکاس پہلے درجے دے !!{sentence}~$\lforall[x][\Atom{Q}{x,x}]$ نال ضرور متعین ہو سکدا اے۔ پر ایہ ہِیٹھلے اک رکنی دوجے درجے دے !!{sentence} نال وی متعین ہو سکدا اے۔
\[
\lforall[X][\lforall[x][(\lforall[y][(\Atom{Q}{x,y} \lif \Atom{X}{y})]
    \lif \Atom{X}{x})]].
\]
ایہدا مطلب اے کہ دونیں جملے ہم ارز نیں۔ ایہ گل سدھی طرح پرکھن دا اک طریقہ ایہ اے۔ پہلاں منّ لو دوجے درجے دا جملہ !!a{structure}~$\Struct{M}$ وچ سچا اے۔ کیوں جے $x$ تے $X$ دوناں اُتے کُلّی مسوّر لگے نیں، ایس لئی باقی حصہ ہر $x \in W$ تے سیٹ~$X \subseteq W$ لئی پورا اترنا چاہیدا اے، مثلاً سیٹ $\Setabs{z}{Rxz}$ لئی، جتھے $R =
\Assign{Q}{M}$۔ سو ہر ایسے~$s$ لئی جیہدے وچ $s(x) \in W$ تے $s(X) =
\Setabs{z}{Rxz}$، ساڈے کول $\Sat{M}{\lforall[y][(\Atom{Q}{x,y} \lif
    \Atom{X}{y})] \lif \Atom{X}{x}}$ اے۔ پر $s(X)$ نوں جیویں چُنیا اے، اوہدے مطابق ایس دا مطلب اے $\Sat{M}{\lforall[y][(\Atom{Q}{x,y} \lif
    \Atom{Q}{x,y})] \lif \Atom{Q}{x,x}}[s]$، جیہڑا $\Atom{Q}{x,x}$ نال ہم ارز اے کیوں جے مقدم معتبر اے۔ کیوں جے $s(x)$ کوئی وی ہو سکدا اے، ایس لئی $\Sat{M}{\lforall[x][\Atom{Q}{x,x}]}$۔

ہُن منّ لو $\Sat{M}{\lforall[x][\Atom{Q}{x,x}]}$ تے وکھاؤ کہ
$\Sat{M}{\lforall[X][\lforall[x][(\lforall[y][(\Atom{Q}{x,y} \lif
        \Atom{X}{y})] \lif \Atom{X}{x})]]}$۔ کوئی وی قدر بندی~$s$ چُنو تے منّ لو $\Sat{M}{\lforall[y][(\Atom{Q}{x,y} \lif
    \Atom{X}{y})]}[s]$۔ منّ لو $s'$ اوہ $y$-بدل قدر بندی اے جیہڑی~$s$ دی اے جیہدے لئی $s'(y)
= s(x)$۔ سو $\Sat{M}{\Atom{Q}{x,y} \lif \Atom{X}{y}}[s']$، یعنی
$\Sat{M}{\Atom{Q}{x,x} \lif \Atom{X}{x}}[s]$۔ کیوں جے
$\Sat{M}{\lforall[x][\Atom{Q}{x,x}]}$، مقدم سچا اے، تے ساڈے کول
$\Sat{M}{\Atom{X}{x}}[s]$ اے، جیہڑا اسیں وکھاؤنا سی۔

کیوں جے فریماں دیاں کجھ متعین ہو سکّن والیاں جماعتاں پہلے درجے دی زبان وچ متعین نہیں ہو سکدیاں، ایس لئی کجھ اک رکنی دوجے درجے دے !!{sentence} جیہڑے~$!A'$ دی شکل دے نیں، پہلے درجے دے !!{sentence} نال ہم ارز نہیں۔ ایہ فیصلہ کرن دا کوئی مؤثر طریقہ نہیں کہ ایہناں وچوں کہڑے ہم ارز نیں۔
\endgroup
% /OLP-0426

\OLEndChapterHook
\endgroup
% /OLP-0419

% OLP-0427 content/normal-modal-logic/axioms-systems/axioms-systems.tex
\begingroup
\def\olfilename{axioms-systems}
\expandafter\def\csname ol@sectioncs\endcsname{section}

\olchapter{nml}{prf}{مسلّماتی \usetoken{P}{derivation}}

% OLP-0428 content/normal-modal-logic/axioms-systems/introduction.tex
\begingroup
\def\olfilename{introduction}
\expandafter\def\csname ol@sectioncs\endcsname{section}

\olfileid{nml}{axs}{int}

\olsection{جان پہچان}

ساڈے کول بنیادی امکانی زبان دی معنیات امکانی نمونیاں راہیں اے۔ ساڈے کول !!a{formula} دے معتبر ہون دا خیال وی اے، یعنی سارے نمونیاں دے سارے جہاناں وچ سچا ہونا، یا نمونیاں یا فریماں دی کسے جماعت دے لحاظ نال معتبر ہونا، یعنی جماعت دے سارے نمونیاں دے سارے جہاناں وچ، یا فریم اُتے ٹکے نمونیاں دے سارے جہاناں وچ، سچا ہونا۔ منطق عام طور تے اعتبار دیاں ایہو جیہیاں معنیاتی پچھاناں نوں !!{derivability} دے ثبوتی خیال نال جوڑدی اے۔ مقصد ایہ اے کہ کسے نظام وچ !!{derivability} ایویں متعین کریے کہ !!a{formula} اُدوں تے صرف اُدوں !!{derivable} ہووے جدوں اوہ معتبر ہووے۔

سبھ توں سادے تے تاریخی طور تے سبھ توں پرانے !!{derivation} دے نظام اوہ نیں جیہناں نوں ہلبرٹ طرز دے یا مسلّماتی !!{derivation} دے نظام آکھدے آں۔ کئی امکانی منطقاں لئی ہلبرٹ طرز دے !!{derivation} دے نظام بناؤنے مقابلتاً سوکھے نیں۔ نظریاتی مطالعے وچ اینہاں نال کم کرنا سوکھا اے، جیویں اینہاں دی صحت مندی تے تمامیت ثابت کرنا۔ پر !!{formula}s ثابت کرن لئی اینہاں نوں ورتنا، مثلاً فطری استنباط دے نظاماں دی ورتوں نالوں، کِتے اوکھا اے۔

ہلبرٹ طرز دے !!{derivation} دے نظاماں وچ کسے !!a{derivation} دا، اک !!a{formula} لئی، مطلب !!{formula}s دی اک لڑی اے جیہڑی کجھ مسلّمیاں توں تھوڑے جیہے استنباطی قاعدیاں راہیں متعلقہ !!{formula} تک پہنچدی اے۔ کیوں جے اسیں چاہندے آں کہ !!{derivation} دا نظام معنیات نال میل کھاوے، ایس لئی سانوں ایہ پکا کرنا پئے گا کہ !!{derivable} فارمولیاں دے سیٹ دا ہر فارمولا ہر نمونے وچ سچا ہووے (یا ہر اوس نمونے وچ سچے ہون جیہدے وچ سارے مسلّمے سچے نیں)۔ پہلاں اسیں امکانی منطقاں دیاں اوہ خاصیتاں کڈھ کے ویکھاں گے جیہڑیاں ایس کم لئی ضروری نیں، یعنی «نارمل» امکانی منطقاں دیاں۔ نارمل امکانی منطقاں لئی صرف دو استنباطی قاعدے منّنے پیندے نیں، قاعدۂ اثباتِ مقدم تے لازمیت دا قاعدہ۔ مسلّمیاں وجوں ہمیشہ صادق فارمولیاں دیاں ساریاں (بدل توں بنیاں) مثالاں، تے متعلقہ امکانی منطق دے حساب نال کجھ امکانی مسلّمے، لیندے آں۔ جے ساڈی دلچسپی صرف سارے نمونیاں دی جماعت وچ اے، تاں وی~$\Ax{K}$\iftag{notprvDiamond}{}{ تے~$\Ax{Dual}$} دیاں بدل نال بنیاں ساریاں مثالاں مسلّمیاں وچ گنّیاں جان گیاں۔ ایہناں نال ای گھٹ توں گھٹ نارمل امکانی منطق~$\Log K$ بندی اے۔

\begin{defn}
\emph{قاعدۂ اثباتِ مقدم} استنباط دا ایہ سانچہ اے۔
\begin{prooftree}
\AxiomC{$!A$}
\AxiomC{$!A \lif !B$}
\RightLabel{\MP}
\BinaryInfC{$!B$}
\end{prooftree}
اسیں آکھدے آں کہ !!a{formula}~$!B$ \emph{نکلدا اے}~!!{formula}s $!A$، $!C$ توں قاعدۂ اثباتِ مقدم نال اُدوں تے صرف اُدوں جدوں $!C \ident !A \lif !B$۔
\end{defn}

\begin{defn}
\emph{لازمیت دا قاعدہ} استنباط دا ایہ سانچہ اے۔
\begin{prooftree}
\AxiomC{$!A$}
\RightLabel{\Nec}
\UnaryInfC{$\Box !A$}
\end{prooftree}
اسیں آکھدے آں کہ !!{formula}~$!B$، !!{formula}s $!A$ توں لازمیت دے قاعدے نال اُدوں تے صرف اُدوں نکلدا اے جدوں $!B \ident \Box !A$۔
\end{defn}

\begin{defn}
مسلّمیاں دے سیٹ~$\Sigma$ توں \emph{!!{derivation}} اوہ !!{formula}s دی لڑی اے جیہڑی $!B_1$، $!B_2$، \dots، $!B_n$ نال لکھی جاندی اے، جیہدے وچ ہر $!B_i$ لئی ہِیٹھ لکھیاں وچوں کوئی اک گل سچی اے۔
\begin{enumerate}
\item اوہ کسے ہمیشہ صادق فارمولے توں بدل نال بنی مثال اے؛ یا
\item اوہ~$\Sigma$ دے کسے !!a{formula} توں بدل نال بنی مثال اے؛ یا
\item اوہ دو !!{formula}s $!B_j$، $!B_k$ توں، جتھے $j$، $k < i$، قاعدۂ اثباتِ مقدم نال نکلدا اے؛ یا
\item اوہ کسے !!a{formula}~$!B_j$ توں، جتھے $j < i$، لازمیت دے قاعدے نال نکلدا اے۔
\end{enumerate}
جے کوئی ایسا !!a{derivation} ہووے جیہدے لئی $!B_n \ident !A$، تاں اسیں آکھدے آں کہ
$!A$، \emph{!!{derivable} اے، $\Sigma$ توں}، علامتاں وچ $\Sigma \Proves !A$۔
\end{defn}

ایس تعریف نال پتا لگے گا کہ !!{derivable} فارمولیاں دا سیٹ اک نارمل امکانی منطق بناؤندا اے، تے ہر !!{derivable} !!{formula} ہر اوس نمونے وچ سچا اے جیہدے وچ ہر مسلّمہ سچا اے۔ !!{derivation}s دی ایس خاصیت نوں \emph{صحت مندی} آکھدے آں۔ اُلٹے رخ نوں \emph{تمامیت} آکھدے آں، تے اینوں ثابت کرنا اوکھا اے۔
\endgroup
% /OLP-0428
% OLP-0429 content/normal-modal-logic/axioms-systems/normal-logics.tex
\begingroup
\def\olfilename{normal-logics}
\expandafter\def\csname ol@sectioncs\endcsname{section}

\olfileid{nml}{prf}{nor}

\olsection{نارمل امکانی منطقاں}

امکانی !!{formula}s دے کجھ سیٹ ایسے نیں جیہناں نوں کسے مسلّمیاں دے سیٹ توں اخذ ہون والے !!{formula}s وجوں پچھاننا سوکھا نہیں۔ اسیں چاہندے آں کہ امکانی منطقاں کجھ خاص قاعدے پالن۔ سبھ توں پہلاں، جو کجھ اسیں کلاسیکی قضیاتی منطق وچ !!{derive} کر سکدے آں، اوہ اجے وی !!{derivable} ہونا چاہیدا اے۔ ایہ خیال رکھنا اے کہ ہُن !!{formula}s وچ
\iftag{prvBox}{$\Box$\iftag{prvDiamond}{
    تے~}{}}{}\iftag{prvDiamond}{$\Diamond$}{} وی آ سکدے نیں۔ ایس لئی اسیں شرط لاؤندے آں کہ امکانی منطق وچ ہمیشہ صادق فارمولیاں توں بدل نال بنیاں ساریاں مثالاں ہون، تے ایہ قاعدۂ اثباتِ مقدم تحت بند ہووے۔

\begin{defn}
  \emph{امکانی منطق} امکانی !!{formula}s دا ایسا سیٹ~$\Sigma$ اے جیہڑا
  \begin{enumerate}
  \item سارے ہمیشہ صادق فارمولے رکھدا اے؛
  \item بدل تحت بند اے، یعنی جے $!A \in \Sigma$ تے
    $!D_1$، \dots، $!D_n$، !!{formula}s نیں، تاں
    \[
    \SSubst{!A}{\subst{!D_1}{p_1}, \dots, \subst{!D_n}{p_n}} \in \Sigma,
    \]
  \item \emph{قاعدۂ اثباتِ مقدم} تحت بند اے، یعنی جے $!A$ تے $!A
      \lif !B \in \Sigma$، تاں $!B \in \Sigma$۔
  \end{enumerate}
\end{defn}

امکانی منطقاں لئی تعلقی معنیات ورتن واسطے سانوں ایہ شرط وی لاؤنی پئے گی کہ سارے امکانی نمونیاں وچ معتبر سارے !!{formula}s شامل ہون۔ پتا لگدا اے کہ ایہ شرط پوری ہو جاندی اے جے \Ax{K}\iftag{prvDiamond}{ تے~\Dual{}}{} دیاں ساریاں مثالاں
!!{derivable} ہون، تے جدوں وی !!a{formula}~$!A$، !!{derivable} ہووے، تاں~$\Box !A$ وی۔ ایہ شرطاں پوریاں کرن والی امکانی منطق نوں \emph{نارمل} آکھدے آں۔ (غیر نارمل امکانی منطقاں وی نیں، پر عام تعلقی نمونے اوہناں لئی کافی نہیں۔)

\begin{defn}
  امکانی منطق $\Sigma$ \emph{نارمل} اے جے اوہدے وچ
  \begin{align*}
    \tag{\Ax{K}} & \Box(p \lif q) \lif (\Box p \lif \Box q),
    \iftag{prvDiamond}{\\
      \tag{\Dual} & \Diamond p \liff \lnot\Box\lnot p}{}
  \end{align*}
  شامل نیں، تے اوہ \emph{لازمیت دے قاعدے} تحت بند اے، یعنی جے $!A \in
  \Sigma$، تاں $\Box !A \in \Sigma$۔
\end{defn}

دھیان دیو کہ ہمیشہ صادق لزوم اینا باریک اے کہ \emph{جہان وچ سچائی} برقرار رکھدا اے، جدکہ قاعدہ \Nec{} صرف \emph{نمونے وچ سچائی} برقرار رکھدا اے (تے ایس لئی فریم یا فریماں دی جماعت وچ منطقی اعتبار وی)۔

\begin{prop}\ollabel{prop:rk}
  ہر نارمل امکانی منطق قاعدے \RK تحت بند اے۔
  \begin{prooftree}
    \AxiomC{$!A_1 \lif (!A_2 \lif \cdots (!A_{n-1} \lif !A_n)\cdots)$}
    \RightLabel{\RK}
    \UnaryInfC{$\Box!A_1 \lif (\Box!A_2 \lif \cdots (\Box!A_{n-1}
      \lif \Box!A_n)\cdots).$}
  \end{prooftree}
\end{prop}

\begin{proof}
  استقرا نال،~$n$ اُتے۔ جے $n = 1$ تاں قاعدہ \Nec ای اے، تے ہر نارمل امکانی منطق \Nec تحت بند اے۔

  ہُن منّ لو نتیجہ $n-1$ لئی چلدا اے۔ اسیں وکھاؤنے آں کہ ایہ~$n$ لئی وی چلدا اے۔

  منّ لو
  \begin{align*}
  & \formula A_1 \lif (!A_2 \lif \cdots (!A_{n-1} \lif !A_n)\cdots) \in \Sigma
  \intertext{استقرائی مفروضے مطابق ساڈے کول ایہ اے}
  & \Box!A_1 \lif (\Box!A_2 \lif \cdots \Box(!A_{n-1} \lif !A_n)\cdots)
  \in \Sigma
  \intertext{کیوں جے $\Sigma$ اک نارمل امکانی منطق اے، ایس وچ~$\Ax{K}$ دیاں ساریاں مثالاں نیں، خاص طور تے ایہ}
  & \Box(!A_{n-1} \lif !A_n) \lif (\Box!A_{n-1} \lif \Box!A_n) \in \Sigma
  \intertext{قاعدۂ اثباتِ مقدم تے ہمیشہ صادق فارمولیاں توں بدل نال بنیاں مناسب مثالاں ورت کے سانوں ایہ ملدا اے}
  & \Box!A_1 \lif (\Box!A_2 \lif \cdots (\Box!A_{n-1}
  \lif \Box!A_n)\cdots) \in \Sigma.
  \end{align*}
\end{proof}

\begin{prop}\ollabel{prop:notDiamondBot}
  ہر نارمل امکانی منطق $\Sigma$ وچ~$\lnot\Diamond\lfalse$ شامل اے۔
\end{prop}

\begin{prob}
  \olref[nml][prf][nor]{prop:notDiamondBot} ثابت کرو۔
\end{prob}

\begin{prop}
  منّ لو $!A_1$، \dots، $!A_n$، !!{formula}s نیں۔ فیر اک سبھ توں چھوٹی امکانی منطق $\Sigma$ اے جیہدے وچ
  $!A_1$، \dots، $!A_n$ دیاں ساریاں مثالاں شامل نیں۔
\end{prop}

\begin{proof}
  $!A_1$، \dots، $!A_n$ دتے ہون، تاں $\Sigma$ نوں ساریاں اوہناں نارمل امکانی منطقاں دا سانجھا حصہ متعین کرو جیہناں وچ
  $!A_1$، \dots، $!A_n$ دیاں ساریاں مثالاں شامل نیں۔ سانجھا حصہ خالی نہیں، کیوں جے
  $\Frm[L]$، یعنی سارے !!{formula}s دا سیٹ، ایسی اک امکانی منطق اے۔
\end{proof}

\begin{defn}
سبھ توں چھوٹی نارمل امکانی منطق جیہدے وچ $!A_1$، \dots، $!A_n$ شامل ہون، \emph{امکانی نظام} اکھواؤندی اے تے اینوں $\Log{K} !A_1 \dots
!A_n$ لکھدے آں۔ سبھ توں چھوٹی نارمل امکانی منطق نوں~\Log{K} لکھدے آں۔
\end{defn}
\endgroup
% /OLP-0429
% OLP-0430 content/normal-modal-logic/axioms-systems/logics-proofs.tex
\begingroup
\def\olfilename{logics-proofs}
\expandafter\def\csname ol@sectioncs\endcsname{section}

\olfileid{nml}{prf}{prf}

\olsection{\usetoken{P}{derivation} تے امکانی نظام}

پہلاں اسیں ایہ متعین کردے آں کہ نارمل امکانی منطقاں وچ !!a{derivation} کیہ ہوندا اے۔ موٹے طور تے !!a{derivation}، !!{formula}s دی لڑی اے جیہدے وچ ہر !!{element} یا تاں کجھ \emph{مسلّمیاں} وچوں کوئی اک (یا اوہدی بدل توں بنی مثال) اے، یا پہلے !!{element}s توں کجھ استنباطی قاعدیاں وچوں کسے اک نال نکلدا اے۔ نارمل امکانی منطقاں لئی ہمیشہ صادق فارمولیاں\iftag{prvDiamond}{، \Ax{K}، تے~\Dual{}}{ تے~\Ax{K}} دیاں ساریاں مثالاں مسلّمے منّیاں جاندیاں نیں۔ ایس توں امکانی نظام~$\Log{K}$، یعنی سبھ توں چھوٹی نارمل امکانی منطق، ملدا اے۔ پر ہور نظام لبھن لئی اسیں ہور مسلّمے شامل کر سکدے آں۔ قاعدے ہمیشہ قاعدۂ اثباتِ مقدم~\MP{} تے لازمیت دا قاعدہ~\Nec نیں۔

\begin{defn}
  اک امکانی نظام $\Log{K} !A_1 \dots !A_n$ تے اک !!a{formula}~$!B$ دِتّے ہون، تاں اسیں آکھدے آں کہ $!B$،
  $\Log{K} !A_1 \dots
  !A_n$ وچ \emph{!!{derivable}} اے، تے اینوں $\Log{K} !A_1 \dots !A_n \Proves !B$ لکھدے آں، اُدوں تے صرف اُدوں جدوں ایسے !!{formula}s $!C_1$، \dots، $!C_k$ ہون کہ $!C_k = !B$ تے ہر $!C_i$ یا تاں ہمیشہ صادق فارمولے توں بدل نال بنی مثال اے، یا $\Ax{K}$،\iftag{prvDiamond}{ $\Dual$،}{} $!A_1$،
  \dots، $!A_n$ وچوں کسے اک دی مثال اے، یا پہلے !!{formula}s توں قاعدیاں \MP{} یا~\Nec نال نکلدا اے۔
\end{defn}

ہِیٹھلے قضیے نال اسیں $!B \in \Sigma$ وکھا سکدے آں، جے $\Sigma$-!!{derivation} پیش کریے~$!B$ دا۔

\begin{prop}
  $\Log{K} !A_1 \dots !A_n = \Setabs{!B}{\Log{K} !A_1 \dots !A_n \Proves !B}$.
\end{prop}

\begin{proof}
  !!{derivation}s دی لمبائی اُتے استقرا ورت کے وکھاؤنے آں کہ
  $\Setabs{!B}{\Log{K} !A_1 \dots !A_n \Proves !B} \subseteq
  \Log{K} !A_1 \dots !A_n$۔

  جے~$!B$ دے !!{derivation} دی لمبائی~$1$ اے، تاں اوہدے وچ صرف اک !!{formula} اے۔ اوہ !!{formula} پچھلے فارمولیاں توں
  \MP{} یا \Nec نال نہیں نکل سکدا، ایس لئی اوہ یا تاں ہمیشہ صادق فارمولے توں بدل نال بنی مثال اے، یا
  \Ax{K}،\iftag{prvDiamond}{ $\Dual$،}{} یا $!A_1$، \dots،~$!A_n$ وچوں کسے اک دی مثال اے۔ پر $\Log{K}!A_1\dots!A_n$ وچ ایہ وی شامل نیں، ایس لئی $!B \in \Log{K}!A_1 \dots !A_n$۔

  جے $!B$ دے !!{derivation} دی لمبائی~$> 1$ اے، تاں $!B$، \MP{} یا \Nec{} نال اوہناں !!{formula}s توں وی مل سکدا اے جیہڑے !!{derivation} دی آخری سطر اُتے نہیں۔ جے $!B$،
  $!C$ تے $!C \lif !B$ توں نکلدا اے (\MP نال)، تاں $!C$ تے $!C \lif !B \in
  \Log{K}!A_1 \dots !A_n$، استقرائی مفروضے مطابق۔ پر ہر امکانی منطق قاعدۂ اثباتِ مقدم تحت بند اے، ایس لئی $!B \in \Log{K}!A_1 \dots
  !A_n$۔ جے $!B \equiv \Box !C$، $!C$ توں \Nec نال نکلدا اے، تاں $!C
  \in \Log{K}!A_1 \dots !A_n$، استقرائی مفروضے مطابق۔ پر ہر نارمل امکانی منطق $\Nec$ تحت بند اے، ایس لئی $!B \in
  \Log{K}!A_1\dots!A_n$۔

  اُلٹے رخ دا شامل ہونا وکھاؤن لئی ثابت کرو کہ
  $\Sigma = \Setabs{!B}{\Log{K} !A_1 \dots !A_n \Proves !B }$ اک نارمل امکانی منطق اے جیہدے وچ $!A_1$،
  \dots، $!A_n$ دیاں ساریاں مثالاں شامل نیں، تے دھیان دیو کہ $\Log{K} !A_1 \dots !A_n$ تعریف مطابق سبھ توں چھوٹی ایسی منطق اے۔
  \begin{enumerate}
    \item ہر ہمیشہ صادق فارمولا~$!B$ ہمیشہ صادق فارمولے توں بدل نال بنی مثال اے، ایس لئی
      $\Log{K}!A_1\dots!A_n \Proves !B$، سو $\Sigma$ وچ سارے ہمیشہ صادق فارمولے شامل نیں۔
    \item جے $\Log{K}!A_1\dots!A_n \Proves !C$ تے
      $\Log{K}!A_1\dots!A_n \Proves !C \lif !B$، تاں
      $\Log{K}!A_1\dots!A_n \Proves !B$۔ !!{derivation}، جیہڑا
      $!C$ دا اے، $!C \lif !B$ دے اخذ نال جوڑو تے سطر~$!B$ شامل کرو۔ آخری سطر دا جواز \MP{} اے۔ سو $\Sigma$ قاعدۂ اثباتِ مقدم تحت بند اے۔
    \item جے $!B$ دا کوئی !!a{derivation} اے، تاں~$!B$ توں بدل نال بنی ہر مثال دا وی اک !!a{derivation} اے۔ !!{derivation} وچ ہر !!{formula} اُتے اوہی بدل لاؤ۔ (مشق۔ !!{derivation}s دی لمبائی اُتے استقرا نال ثابت کرو کہ نتیجہ وی اک ٹھیک !!{derivation} اے۔) سو $\Sigma$ یکساں بدل تحت بند اے۔ (ہُن اسیں ثابت کر دِتّا اے کہ $\Sigma$ امکانی منطق دیاں ساریاں شرطاں پوریاں کردا اے۔)
    \item ساڈے کول $\Log{K}!A_1\dots!A_n \Proves \Ax{K}$ اے، ایس لئی $K \in \Sigma$۔
    \tagitem{prvDiamond}{ساڈے کول
      $\Log{K}!A_1\dots!A_n \Proves \Dual$ اے، ایس لئی $\Dual \in \Sigma$۔}{}
    \item جے $\Log{K}!A_1\dots!A_n \Proves !C$، تاں ودھائی سطر~$\Box !C$ دا جواز~\Nec اے۔ ایس لئی $\Sigma$،
      ~\Nec تحت بند اے۔ سو $\Sigma$ نارمل اے۔
  \end{enumerate}
\end{proof}
\endgroup
% /OLP-0430
% OLP-0431 content/normal-modal-logic/axioms-systems/proofs-in-K.tex
\begingroup
\def\olfilename{proofs-in-K}
\expandafter\def\csname ol@sectioncs\endcsname{section}

\olfileid{nml}{prf}{prk}

\olsection{\Log{K} وچ ثبوت}

سبھ توں چھوٹے امکانی نظام وچ ثبوتاں دی مشق کرن لئی اسیں وکھاؤنے آں کہ معتبر !!{formula}s، جیہڑے
\olref[syn][sch]{tab:valid-invalidSchemas} دے کھبے پاسے نیں، اوہناں ساریاں دے
\Log{K} وچ ثبوت دِتّے جا سکدے نیں۔

\begin{prop}
  $\Log{K} \Proves \Box!A \lif \Box (!B\lif !A)$
\end{prop}

\begin{proof}
  \begin{derivation}
    1. & $!A \lif (!B \lif !A)$ & \Taut \\
    2. & $\Box(!A \lif (!B \lif !A))$ & \Nec, 1 \\
    3. & $\Box(!A \lif (!B \lif !A)) \lif
    (\Box!A \lif \Box(!B \lif !A))$ & \Ax{K}\\
    4. & $\Box!A \lif \Box (!B\lif !A)$ & \MP, 2, 3
  \end{derivation}
\end{proof}

\begin{prop}
  $\Log{K} \Proves \Box(!A \land !B) \lif (\Box !A \land \Box!B)$
\end{prop}

\begin{proof}
\begin{derivation}
  1. & $(!A \land !B) \lif !A$ & \Taut \\
  2. & $\Box((!A \land !B) \lif !A)$ & \Nec \\
  3. & $\Box((!A \land !B) \lif !A) \lif (\Box(!A \land !B) \lif \Box!A)$
  & \Ax{K} \\
  4. & $\Box(!A \land !B) \lif \Box!A$ & \MP, 2, 3 \\
  5. & $(!A \land !B) \lif !B$ & \Taut \\
  6. & $\Box((!A \land !B) \lif !B)$ & \Nec \\
  7. & $\Box((!A \land !B) \lif !B) \lif (\Box(!A \land !B) \lif \Box!B)$
  & \Ax{K} \\
  8. & $\Box(!A \land !B) \lif \Box!B$ & \MP, 6, 7 \\
  9. & $(\Box(!A \land !B) \lif \Box!A) \lif{}$ \\
  & \qquad $((\Box(!A \land !B) \lif \Box!B) \lif{}$ \\
  & \qquad $(\Box(!A \land !B) \lif (\Box !A \land \Box!B)))$ & \Taut\\
  10. &  $(\Box(!A \land !B) \lif \Box!B) \lif{}$ \\
  & \qquad $(\Box(!A \land !B) \lif (\Box !A \land \Box!B))$ & \MP, 4, 9\\
  11. & $\Box(!A \land !B) \lif (\Box !A \land \Box!B)$ & \MP, 8, 10.
\end{derivation}
دھیان دیو کہ سطر~$9$ اُتے !!{formula} ہِیٹھلے ہمیشہ صادق فارمولے توں بدل نال بنی مثال اے۔
\[
(p \lif q) \lif ((p \lif r) \lif (p \lif (q \land r))).
\]
\end{proof}

\begin{prop}
  $\Log{K}\Proves (\Box!A \land \Box!B) \lif \Box (!A \land !B)$
\end{prop}

\begin{proof}
  \begin{derivation}
    1. & $!A \lif (!B \lif (!A \land !B))$ & \Taut \\
    2. & $\Box(!A \lif (!B \lif (!A \land !B)))$ & \Nec, 1 \\
    3. & $\Box(!A \lif (!B \lif (!A \land !B))) \lif (\Box!A \lif \Box(!B \lif (!A \land !B)))$ & \Ax{K}\\
    4. & $\Box!A \lif \Box(!B \lif (!A \land !B))$ & \MP, 2, 3\\
    5. & $\Box(!B \lif (!A \land !B)) \lif (\Box!B \lif \Box(!A \land !B))$ & \Ax{K} \\
    6. & $(\Box!A \lif \Box(!B \lif (!A \land !B))) \lif {}$ \\
    & \qquad $(\Box(!B \lif (!A \land !B)) \lif (\Box!B \lif \Box(!A \land !B))) \lif {}$\\
    & \qquad $(\Box!A \lif (\Box !B \lif \Box(!A \land !B))))$ & \Taut\\
    7. & $(\Box(!B \lif (!A \land !B)) \lif (\Box!B \lif \Box(!A \land !B))) \lif {}$\\
    & \qquad $(\Box!A \lif (\Box !B \lif \Box(!A \land !B)))$ & \MP, 4, 6\\
    8. & $\Box!A \lif (\Box !B \lif \Box(!A \land !B)))$ & \MP, 5, 7\\
    9. & $(\Box!A \lif (\Box !B \lif \Box(!A \land !B)))) \lif {}$\\
    & \qquad $((\Box!A \land \Box!B) \lif \Box(!A \land !B))$ & \Taut\\
    10. & $(\Box!A \land \Box!B) \lif \Box(!A \land !B)$ & \MP, 8, 9
  \end{derivation}
  سطراں $6$ تے $9$ اُتے !!{formula}s ہِیٹھلے ہمیشہ صادق فارمولیاں توں بدل نال بنیاں مثالاں نیں۔
  \begin{align*}
    (p \lif q) & \lif ((q \lif r) \lif (p \lif r)) \\
    (p \lif (q \lif r)) & \lif ((p \land q) \lif r)
  \end{align*}
\end{proof}

\begin{prop}
  \iftag{prvBox}{$\Log{K} \Proves \lnot\Box p \lif \Diamond \lnot
    p$}{$\Log{K} \Proves \Box\lnot p \lif \lnot\Diamond p$}
\end{prop}

\begin{proof}
  \iftag{prvBox}{\iftag{prvDiamond}{% لازمیت تے امکان دے نشان ابتدائی نیں۔
  \begin{derivation}
    1. & $\Diamond \lnot p \liff \lnot\Box\lnot\lnot p$ & \Dual\\
    2. & $(\Diamond \lnot p \liff \lnot\Box\lnot\lnot p) \lif {}$ \\
    & \qquad $(\lnot \Box \lnot\lnot p \lif \Diamond \lnot p)$& \Taut\\
    3. & $\lnot \Box \lnot\lnot p \lif \Diamond \lnot p$ & \MP, 1, 2\\
    4. & $\lnot\lnot p \lif p$ & \Taut\\
    5. & $\Box(\lnot\lnot p \lif p)$ & \Nec, 4\\
    6. & $\Box(\lnot\lnot p \lif p) \lif (\Box \lnot\lnot p \lif \Box p)$ & \Ax{K}\\
    7. & $(\Box \lnot\lnot p \lif \Box p)$ & \MP, 5, 6\\
    8. & $(\Box \lnot\lnot p \lif \Box p) \lif (\lnot \Box p \lif \lnot\Box\lnot\lnot p)$ & \Taut\\
    9. & $\lnot \Box p \lif \lnot\Box\lnot\lnot p$ & \MP, 7, 8\\
    10. & $(\lnot \Box p \lif \lnot\Box\lnot\lnot p) \lif {}$ \\
    & \qquad $((\lnot \Box \lnot\lnot p \lif \Diamond \lnot p) \lif (\lnot \Box p \lif \Diamond\lnot p))$ & \Taut\\
    11. & $(\lnot \Box \lnot\lnot p \lif \Diamond \lnot p) \lif (\lnot \Box p \lif \Diamond\lnot p)$ & \MP, 9, 10\\
    12. & $\lnot\Box p \lif \Diamond \lnot p$ & \MP, 3, 11\\
  \end{derivation}
  سطراں $8$ تے $10$ اُتے !!{formula}s ہِیٹھلے ہمیشہ صادق فارمولیاں توں بدل نال بنیاں مثالاں نیں۔
  \begin{align*}
    & (p \lif q) \lif (\lnot q \lif \lnot p) \\
    & (p \lif q) \lif ((q \lif r) \lif (p \lif r)).
  \end{align*}
  }{% امکان دا نشان تعریف نال ملدا اے۔
  \begin{derivation}
    1. & $\lnot\lnot p \lif p$ & \Taut\\
    2. & $\Box(\lnot\lnot p \lif p)$ & \Nec, 1\\
    3. & $\Box(\lnot\lnot p \lif p) \lif (\Box \lnot\lnot p \lif \Box p)$ & \Ax{K}\\
    4. & $(\Box \lnot\lnot p \lif \Box p)$ & \MP, 2, 3\\
    5. & $(\Box \lnot\lnot p \lif \Box p) \lif (\lnot \Box p \lif \lnot\Box\lnot\lnot p)$ & \Taut\\
    6. & $\lnot \Box p \lif \lnot\Box\lnot\lnot p$ & \MP, 4, 5\\
    & $\lnot\Box p \lif \Diamond \lnot p$
  \end{derivation}
  آخری سطر $\lnot \Box p \lif \Diamond\lnot p$ اے، کیوں جے
  $\Diamond$ دی تعریف $\lnot\Box\lnot$ اے۔ سطر~$5$ اُتے !!{formula} ہِیٹھلے ہمیشہ صادق فارمولے توں بدل نال بنی مثال اے۔
  \[
    (p \lif q) \lif (\lnot q \lif \lnot p).
  \]
  }}{% لازمیت دا نشان تعریف نال ملدا اے۔
  \begin{derivation}
    1. & $\Diamond  p \liff \lnot\Box\lnot p$ & \Dual\\
    2. & $(\Diamond p \liff \lnot\Box\lnot p) \lif
    (\Diamond p \lif \lnot \Box \lnot p)$& \Taut\\
    3. & $\Diamond p \lif \lnot \Box \lnot p$ & \MP, 1, 2\\
    4. & $(\Diamond p \lif \lnot \Box \lnot p) \lif (\Box\lnot p \lif \lnot\Diamond p)$ & \Taut\\
    5. & $\Box\lnot p \lif \lnot\Diamond p$ & \MP, 3, 4
  \end{derivation}
  سطر~$4$ اُتے !!{formula} ہِیٹھلے ہمیشہ صادق فارمولے توں بدل نال بنی مثال اے۔
  \begin{align*}
    & (p \lif \lnot q) \lif (q \lif \lnot p).
  \end{align*}
}
\end{proof}

\begin{prob}
  ~$\Log{K}$ وچ ہِیٹھ لکھیاں !!{formula}s لئی !!{derivation}s لبھو۔
  \begin{enumerate}
    \item $\Box \lnot p \lif \Box(p \lif q)$
    \item $(\Box p \lor \Box q) \lif \Box(p \lor q)$
    \item $\Diamond p \lif \Diamond(p \lor q)$
  \end{enumerate}
\end{prob}
\endgroup
% /OLP-0431
% OLP-0432 content/normal-modal-logic/axioms-systems/derived-rules.tex
\begingroup
\def\olfilename{derived-rules}
\expandafter\def\csname ol@sectioncs\endcsname{section}

\olfileid{nml}{prf}{der}

\olsection{اخذ کیتے قاعدے}

!!{derivation}s لبھن تے لکھن وچ اوکھ اے، کم لمّا اے تے گلّاں وار وار لکھنیاں پیندیاں نیں۔ مثال لئی اسیں اکثر
$!A \lif !B$ توں $\Box !A \lif \Box !B$ تک جانا چاہندے آں، یعنی قاعدہ~\RK لاؤنا چاہندے آں۔ ایس لئی پہلاں \Nec لاؤنا، فیر~\Ax{K} دی مناسب مثال لکھنی، تے فیر~\MP لاؤنا پیندا اے۔ $!A \lif
!B$ تے $!B \lif !C$ توں $!A \lif !C$ تک جان لئی ہمیشہ صادق فارمولے توں بنی ایہ (لمّی) مثال لکھنی پیندی اے۔
\[
(!A \lif !B) \lif ((!B \lif !C) \lif (!A \lif !C))
\]
تے \MP{} دو واری لاؤنا پیندا اے۔ اکثر اسیں اک ذیلی !!{formula} دی تھاں ایسا فارمولا رکھنا چاہندے آں جیہڑا اوہدے نال ہم ارز ہووے۔ مثال لئی \iftag{prvDiamond}{$\Diamond
  !A$ دی تھاں $\lnot\Box\lnot !A$}{$\Box !A$ دی تھاں $\lnot\Diamond\lnot !A$}، یا
$\lnot\lnot !A$ دی تھاں $!A$۔ ایس لئی سارا !!{derivation} لکھن دی تھاں صرف ایہ دسنا سوکھا اے کہ وچکارلے قدم کیوں !!{derivable} نیں۔ ایس کم لئی آؤ !!{derivability} بارے کجھ گلاں اکٹھیاں کریے۔

\begin{prop}
  جے $\Log{K} \Proves !A_1$، \dots، $\Log{K} \Proves !A_n$، تے $!B$،
  $!A_1$، \dots، $!A_n$ توں قضیاتی منطق دے مطابق نکلدا اے، تاں
  $\Log{K} \Proves !B$۔
\end{prop}

\begin{proof}
  جے $!B$، $!A_1$، \dots،~$!A_n$ توں قضیاتی منطق دے مطابق نکلدا اے، تاں
  \[
  !A_1 \lif (!A_2 \lif \cdots (!A_n \lif !B)\dots)
  \]
  ہمیشہ صادق فارمولے توں بدل نال بنی مثال اے۔ \MP{} نوں $n$ واری لاؤن نال~$!B$ دا !!{derivation} ملدا اے۔
\end{proof}

ایس قضیے دی ورتوں نوں اسیں~\PL نال ظاہر کراں گے۔

\begin{prop}
  جے $\Log{K} \Proves !A_1 \lif (!A_2 \lif \cdots (!A_{n-1} \lif
  !A_n)\dots)$، تاں $\Log{K} \Proves \Box !A_1 \lif (\Box !A_2 \lif
  \cdots (\Box !A_{n-1} \lif \Box !A_n)\dots)$۔
\end{prop}

\begin{proof}
  $n$ اُتے استقرا نال، جیویں \olref[nor]{prop:rk} دے ثبوت وچ۔
\end{proof}

ایس قضیے دی ورتوں نوں اسیں~\RK نال ظاہر کراں گے۔ مثال نال ویکھیے کہ !!{derivability} دے نتیجے ثابت کرن وچ ایہ گلاں کِویں سوکھ کردیاں نیں۔

\begin{prop}
  $\Log{K} \Proves (\Box!A \land \Box!B) \lif \Box (!A \land !B)$
\end{prop}

\begin{proof}
  \begin{derivation}
    1. & $\Log{K} \Proves !A \lif (!B \lif (!A \land !B))$ & \Taut \\
    2. & $\Log{K} \Proves \Box!A \lif (\Box !B \lif \Box(!A \land !B)))$ & \RK, 1\\
    3. & $\Log{K} \Proves (\Box!A \land \Box!B) \lif \Box(!A \land !B)$ & \PL, 2
  \end{derivation}
\end{proof}

\begin{prop}\ollabel{prop:rewriting}
  جے $\Log{K} \Proves !A \liff !B$ تے $\Log{K} \Proves
  \Subst{!C}{!A}{q}$، تاں $\Log{K} \Proves \Subst{!C}{B}{q}$۔
\end{prop}

\begin{proof}
  مشق۔
\end{proof}

\begin{prob}
  \olref[nml][prf][der]{prop:rewriting} ثابت کرن لئی~$!C$ دی پیچیدگی اُتے استقرا نال ایہ وکھاؤ کہ جے $\Log{K} \Proves !A \liff !B$،
  تاں $\Log{K} \Proves \Subst{!C}{!A}{q} \liff \Subst{!C}{!B}{q}$۔
\end{prob}

ایہ قضیہ خاص کر اُدوں کم آؤندا اے جدوں اسیں $\Diamond$ نوں $\Box$ وچ بدلنا ہووے (یا اُلٹے پاسے)، یا کسے !!a{formula} دے اندر دوہری نفی ہٹاؤنی ہووے۔ اگے \olref{prop:rewriting} دی ورتوں اُتے اسیں «$!A$، $!B$ دی تھاں» لکھاں گے جدوں وی اک !!a{formula}~$!C(!B)$ نوں~$!C(!A)$ دی تھاں لکھدے آں۔ دوجے لفظاں وچ «$!A$، $!B$ دی تھاں» ایہدا اختصار اے۔
  \begin{derivation}
    & $\Proves !C(!A)$ \\
    & $\Proves !A \liff !B$\\
    & $\Proves !C(!B)$ & \olref{prop:rewriting} مطابق
  \end{derivation}
مثال ویکھو۔

\begin{prop}
  $\Log{K} \Proves \lnot\Box p \lif \Diamond \lnot p$
\end{prop}

\begin{proof}
  \iftag{prvDiamond}{% امکان دا نشان ابتدائی اے۔
  \begin{derivation}
    1. & $\Log{K} \Proves \Diamond \lnot p \liff \lnot\Box\lnot\lnot p$ &
    \Dual\\
    2. & $\Log{K} \Proves \lnot \Box \lnot\lnot p \lif \Diamond \lnot p$ & \PL, 1\\
    3. & $\Log{K} \Proves \lnot\Box p \lif \Diamond \lnot p$ & $p$، $\lnot\lnot p$ دی تھاں\\
  \end{derivation}
  }{% امکان دا نشان تعریف نال ملدا اے۔
  \begin{derivation}
    1. & $\Log{K} \Proves \lnot \Box \lnot\lnot p \lif \Diamond \lnot p$ & \Taut\\
    2. & $\Log{K} \Proves \lnot\Box p \lif \Diamond \lnot p$ & $p$، $\lnot\lnot p$ دی تھاں\\
  \end{derivation}
  سطر~$1$ اُتے فارمولا $\lnot q \lif \lnot q$ دی مثال اے، جیہڑی
  $\Box\lnot\lnot p$ نوں~$q$ دی تھاں رکھ کے ملدی اے۔ $\Diamond\lnot p$ تے
  $\lnot\Box\lnot\lnot p$ عین اکّو نیں، کیوں جے $\Diamond$ دی تعریف
  ~$\lnot\Box\lnot$ اے۔}
\end{proof}

اُپرلے !!{derivation} وچ آخری قدم «$p$، $\lnot\lnot p$ دی تھاں» ایہدا اختصار اے۔
  \begin{derivation}
    & $\Log{K} \Proves \lnot \Box \lnot\lnot p \lif \Diamond \lnot
    p$\\
    & $\Log{K} \Proves \lnot\lnot p \liff p$ & \Taut\\
    & $\Log{K} \Proves \lnot\Box p \lif \Diamond \lnot p$ & \olref{prop:rewriting} مطابق
  \end{derivation}
$!C(q)$، $!A$ تے $!B$ دیاں تھاواں، \olref{prop:rewriting} مطابق، ایتھے بالترتیب $\lnot \Box q \lif \Diamond \lnot p$،
$\lnot\lnot p$ تے~$p$ نے لئیاں نیں۔

جدوں !!a{formula} وچ ذیلی !!{formula} $\lnot\Diamond !A$ ہووے، تاں اوہدے بدلے $\Box\lnot !A$ رکھ سکدے آں، \olref{prop:rewriting} ورت کے، کیوں جے
$\Log{K} \Proves \lnot\Diamond !A \liff \Box\lnot !A$۔ ایہ تے ایہو جیہے ہور بدل صرف «$\Box\lnot$،
$\lnot\Diamond$ دی تھاں» نال ظاہر کراں گے۔

ہِیٹھلا قضیہ ایس گل دا جواز دیندا اے کہ !!{derivability} دے نتیجے سانچیاں دے روپ وچ ثابت کر سکدے آں۔ مثلاً پچھلے قضیے توں صرف $\Log{K} \Proves \lnot\Box p \lif
\Diamond \lnot p$ ای نہیں، بلکہ $\Log{K} \Proves \lnot\Box !A \lif \Diamond
\lnot !A$ وی ملدا اے، کسے وی~$!A$ لئی۔

\begin{prop}
  جے $!A$، $!B$ توں بدل نال بنی مثال اے تے $\Log{K} \Proves !B$،
  تاں $\Log{K} \Proves !A$۔
\end{prop}

\begin{proof}
  ایہ پرکھنا لمّا پر عام کم اے (!!{derivation} دی لمبائی اُتے استقرا، جیہڑا~$!B$ دا اے نال) کہ سارے !!{derivation} اُتے بدل لاؤن نال وی اک ٹھیک
  !!{derivation} ملدا اے۔ خاص طور تے، ہمیشہ صادق فارمولے توں بدل نال بنیاں مثالاں دے اگے بدل وی ہمیشہ صادق فارمولیاں توں بدل نال بنیاں مثالاں نیں۔ ~\iftag{prvDiamond}{\Dual{} تے~}{}\Ax{K} دیاں مثالاں دے بدل خود~\iftag{prvDiamond}{\Dual{}
    تے~}{}\Ax{K} دیاں مثالاں نیں۔ نالے \MP{} تے~\Nec{} دی ورتوں اُدوں وی ٹھیک رہندی اے جدوں !!{formula}s نوں !!{propositional variable}s دی تھاں مقدمے (یا مقدمیاں) تے نتیجے دوناں وچ رکھیے۔
\end{proof}
\endgroup
% /OLP-0432
% OLP-0433 content/normal-modal-logic/axioms-systems/more-proofs-in-K.tex
\begingroup
\def\olfilename{more-proofs-in-K}
\expandafter\def\csname ol@sectioncs\endcsname{section}

\olfileid{nml}{prf}{mpr}

\olsection{\Log{K} وچ ہور ثبوت}

ہُن~\Log{K} وچ !!{derivability} دیاں کجھ ہور مثالاں ویکھیے، \olref[der]{sec} وچ دسے سوکھے طریقے نال۔

\begin{prop}
  $\Log{K} \Proves \Box (!A \lif !B) \lif (\Diamond !A \lif \Diamond
  !B)$
\end{prop}

\begin{proof}
\begin{derivation}
1. & $\Log{K} \Proves (!A \lif !B) \lif (\lnot !B \lif \lnot !A)$ & \PL \\
2. &  $\Log{K} \Proves \Box(!A \lif !B) \lif (\Box\lnot!B \lif \Box\lnot!A)$ & \RK, 1 \\
3. & $\Log{K} \Proves (\Box\lnot!B \lif \Box\lnot!A) \lif (\lnot
\Box\lnot!A \lif \lnot\Box\lnot!B)$ & \Taut\\
4. & $\Log{K} \Proves \Box(!A \lif !B) \lif (\lnot
\Box\lnot!A \lif \lnot\Box\lnot!B)$ & \PL, 2, 3 \\
5. &  $\Log{K} \Proves \Box(!A \lif !B) \lif (\Diamond!A \lif \Diamond!B)$ & $\Diamond$، $\lnot\Box\lnot$ دی تھاں.
\end{derivation}
\end{proof}

\begin{prop}
$\Log{K}\Proves \Box!A \lif (\Diamond(!A \lif !B) \lif
  \Diamond !B)$
\end{prop}

\begin{proof}
  \begin{derivation}
  1. & $\Log{K} \Proves !A \lif (\lnot!B \lif \lnot (!A \lif !B))$ & \Taut \\
  2. & $\Log{K} \Proves \Box!A \lif  (\Box\lnot!B \lif \Box\lnot (!A\lif !B))$ &
  \RK, 1 \\
  3. &  $\Log{K} \Proves \Box!A \lif (\lnot \Box\lnot (!A\lif !B) \lif \lnot
  \Box\lnot!B)$ & \PL, 2 \\
  4. & $\Log{K} \Proves \Box!A \lif (\Diamond(!A \lif !B) \lif
  \Diamond !B)$ & $\Diamond$، $\lnot\Box\lnot$ دی تھاں.
  \end{derivation}
\end{proof}

\begin{prop}
  $\Log{K} \Proves (\Diamond!A \lor \Diamond!B) \lif \Diamond(!A \lor !B)$
\end{prop}

\begin{proof}
  \begin{derivation}
    1. & $\Log{K} \Proves \lnot(!A \lor !B) \lif \lnot!A$ & \Taut \\
    2. & $\Log{K} \Proves \Box\lnot(!A \lor !B) \lif \Box\lnot!A$ & \RK, 1 \\
    3. & $\Log{K} \Proves \lnot\Box\lnot !A \lif \lnot\Box\lnot(!A \lor!B)$ &
    \PL, 2\\
    4. & $\Log{K} \Proves \Diamond!A \lif \Diamond(!A \lor !B)$ &
    $\Diamond$، $\lnot\Box\lnot$ دی تھاں\\
    5. & $\Log{K} \Proves \Diamond!B \lif \Diamond(!A \lor !B)$ & ایس طرح ای\\
    6. & $\Log{K} \Proves (\Diamond!A \lor\Diamond!B) \lif \Diamond(!A \lor !B)$
    & \PL, 4, 5.
  \end{derivation}
\end{proof}

\begin{prop}
  $\Log{K} \Proves \Diamond(!A \lor!B) \lif (\Diamond!A \lor \Diamond!B)$
\end{prop}

\begin{proof}
  \begin{derivation}
    1. & $\Log{K} \Proves \lnot !A \lif (\lnot !B \lif \lnot (!A \lor !B))$ & \Taut \\
    2. & $\Log{K} \Proves \Box\lnot !A \lif
    (\Box\lnot!B \lif \Box \lnot (!A \lor !B))$ & \RK\\
    3. & $\Log{K} \Proves \Box\lnot !A \lif (\lnot \Box \lnot (!A \lor!B)
    \lif \lnot\Box\lnot !B)$ & \PL, 2\\
    4. & $\Log{K} \Proves \lnot \Box \lnot(!A \lor !B) \lif (\Box \lnot !A \lif 
    \lnot\Box\lnot !B)$ & \PL, 3\\
    5. & $\Log{K} \Proves \lnot \Box \lnot(!A\lor!B) \lif (\lnot
    \lnot\Box\lnot!B \lif \lnot\Box\lnot!A)$ & \PL, 4\\
    6. & $\Log{K} \Proves \Diamond(!A \lor !B) \lif (\lnot
    \Diamond!B \lif \Diamond!A)$ & $\Diamond$، $\lnot\Box\lnot$ دی تھاں\\
    7. & $\Log{K} \Proves \Diamond(!A\lor!B) \lif (\Diamond!B \lor \Diamond!A)$ & \PL, 6. \\
  \end{derivation}
\end{proof}

\begin{prob}
  وکھاؤ کہ ہِیٹھلے !!{derivability} دے دعوے سچے نیں۔
  \begin{enumerate}
  \item $\Log{K} \Proves \Diamond \lnot \lfalse \lif (\Box !A \lif
    \Diamond !A)$؛
  \item $\Log{K} \Proves \Box(!A \lor !B) \lif (\Diamond !A \lor \Box
    !B)$؛
  \item $\Log{K} \Proves (\Diamond !A \lif \Box !B) \lif \Box(!A \lif
    !B)$۔
  \end{enumerate}
\end{prob}
\endgroup
% /OLP-0433
% OLP-0434 content/normal-modal-logic/axioms-systems/duals.tex
\begingroup
\def\olfilename{duals}
\expandafter\def\csname ol@sectioncs\endcsname{section}

\olfileid{nml}{prf}{dua}

\olsection{دوہرے \usetoken{P}{formula}}

\begin{defn}\ollabel{def:duals}
  !!{formula}s \Ax{T}، \Ax{B}، \Ax{4} تے
  \Ax{5} وچوں ہر اک دا \emph{دوہرا} اے، جیہڑا ہِیٹھلے پاسے لائے ہیرے دے نشان نال ایویں ظاہر ہوندا اے۔
  \begin{align}
    \tag{\Ax{T_\Diamond}} p & \lif \Diamond p\\
    \tag{\Ax{B_\Diamond}} \Diamond\Box p & \lif p\\
    \tag{\Ax{4_\Diamond}} \Diamond\Diamond p & \lif \Diamond p\\
    \tag{\Ax{5_\Diamond}} \Diamond\Box p & \lif \Box p
    \end{align}
\end{defn}

اُپرلے دوہرے !!{formula}s وچوں ہر اک اپنے نال جڑے !!{formula} توں ایویں ملدا اے۔ $\lnot p$ نوں $p$ دی تھاں رکھو، اُلٹی نفی بناؤ، $\lnot\Box\lnot$ دی تھاں $\Diamond$ رکھو، تے $\lnot\Diamond\lnot$ دی تھاں~$\Box$ رکھو۔ \Ax{D}، یعنی $\Box!A \lif \Diamond!A$، ایس معنی وچ اپنا آپ دوہرا اے۔

\begin{prop}\ollabel{prop:dualsys}
  ہر !!{formula}~$!A$ لئی جیہڑا \olref[nml][prf][dua]{def:duals} وچ اے،
  $\Log{K}!A = \Log{K}!A_{\Diamond}$۔
\end{prop}
\begin{proof}
  مشق۔
\end{proof}
\begin{prob}
  \olref[nml][prf][dua]{prop:dualsys} ثابت کرو۔
\end{prob}
\endgroup
% /OLP-0434
% OLP-0435 content/normal-modal-logic/axioms-systems/proofs-modal-systems.tex
\begingroup
\def\olfilename{proofs-modal-systems}
\expandafter\def\csname ol@sectioncs\endcsname{section}

\olfileid{nml}{prf}{prs}

\olsection{امکانی نظاماں وچ ثبوت}

ہُن اسیں~\Log{K} توں ہور امکانی نظاماں وچ ثبوت ویکھدے آں۔

\begin{prop}\ollabel{prop:S5facts}
ثابت ہون بارے ہِیٹھلے نتیجے ملدے نیں۔
  \begin{enumerate}
  \item $\Log{KT5} \Proves \Ax{B}$؛
  \item $\Log{KT5} \Proves \Ax{4}$؛
  \item $\Log{KDB4} \Proves \Ax{T}$؛
  \item $\Log{KB4} \Proves \Ax{5}$؛
  \item $\Log{KB5} \Proves \Ax{4}$؛
  \item \ollabel{prop:S5facts-KT-D}$\Log{KT} \Proves \Ax{D}$۔
  \end{enumerate}
\end{prop}

\begin{proof}
  اسیں ہر اک دا ثبوت دیندے آں۔
  \begin{enumerate}
    \item $\Log{KT5} \Proves \Ax{B}$:
      \begin{derivation}
        1. & $\Log{KT5} \Proves \Diamond!A \lif \Box\Diamond!A$ & \Ax{5}\\
        2. & $\Log{KT5} \Proves !A \lif \Diamond!A$ & $\Ax{T_\Diamond}$\\
        3. & $\Log{KT5} \Proves !A \lif \Box\Diamond!A$ & \PL, 2, 1.
      \end{derivation}
    \item $\Log{KT5} \Proves \Ax{4}$:
      \begin{derivation}
        1. &$\Log{KT5} \Proves \Diamond\Box!A \lif \Box\Diamond\Box!A$ & \Ax{5}
        $\Box!A$، $p$ دی تھاں\\
        2. & $\Log{KT5} \Proves \Box!A \lif \Diamond\Box!A$ & $\Ax{T_\Diamond}$
        $\Box!A$، $p$ دی تھاں\\
        3. & $\Log{KT5} \Proves \Box!A \lif \Box\Diamond\Box!A$ & \PL, 2, 1\\
        4. & $\Log{KT5} \Proves \Diamond\Box!A \lif \Box!A$ & $\Ax{5_\Diamond}$\\
        5. & $\Log{KT5} \Proves \Box\Diamond\Box!A \lif \Box\Box!A$ & \RK{}, 4 \\
        6. & $\Log{KT5} \Proves \Box!A \lif \Box\Box!A$ & \PL, 3, 5. \\
      \end{derivation}
    \item $\Log{KDB4} \Proves \Ax{T}$:
      \begin{derivation}
        1. & $\Log{KDB4} \Proves \Diamond\Box!A \lif !A $ & $\Ax{B_\Diamond}$ \\
        2. & $\Log{KDB4} \Proves \Box\Box!A \lif \Diamond\Box!A$ & $\Ax{D}$
        $\Box!A$، $p$ دی تھاں\\
        3. & $\Log{KDB4} \Proves \Box\Box!A \lif !A$ & \PL, 1, 2\\
        4. & $\Log{KDB4} \Proves \Box!A \lif \Box\Box!A$ & \Ax{4} \\
        5. & $\Log{KDB4} \Proves \Box!A \lif !A$ & \PL, 4, 3. \\
      \end{derivation}
    \item $\Log{KB4} \Proves \Ax{5}$:
      \begin{derivation}
        1. & $\Log{KB4} \Proves \Diamond!A \lif \Box \Diamond\Diamond!A$ & \Ax{B}
        $\Diamond!A$، $p$ دی تھاں\\
        2. & $\Log{KB4} \Proves \Diamond\Diamond !A \lif \Diamond !A$ &
        $\Ax{4_\Diamond}$ \\
        3. & $\Log{KB4} \Proves \Box\Diamond\Diamond !A \lif \Box \Diamond!A$ &
        \RK{}, 2 \\
        4. & $\Log{KB4} \Proves \Diamond!A \lif \Box\Diamond!A$ & \PL, 1, 3.
      \end{derivation}
    \item $\Log{KB5} \Proves \Ax{4}$:
      \begin{derivation}
        1. & $\Log{KB5} \Proves \Box!A \lif \Box\Diamond\Box !A$ & \Ax{B} $\Box!A$، $p$ دی تھاں \\
        2. & $\Log{KB5} \Proves  \Diamond\Box!A \lif \Box !A$ &
        $\Ax{5_\Diamond}$\\
        3. & $\Log{KB5} \Proves  \Box\Diamond\Box!A \lif \Box\Box !A$ & \RK{}, 2 \\
        4. & $\Log{KB5} \Proves  \Box!A \lif \Box\Box!A$ & \PL, 1, 3.
      \end{derivation}
    \item $\Log{KT} \Proves \Ax{D}$:
      \begin{derivation}
        1. & $\Log{KT} \Proves \Box !A \lif !A$ & \Ax{T} \\
        2. & $\Log{KT} \Proves !A \lif \Diamond!A$ & $\Ax{T_\Diamond}$ \\
        3. & $\Log{KT} \Proves \Box !A \lif \Diamond !A$ &  \PL, 1, 2
      \end{derivation}
  \end{enumerate}
\end{proof}

\begin{defn}
  رواج مطابق اسیں \Log{S4} نوں نظام
  \Log{KT4} تے \Log{S5} نوں نظام \Log{KTB4} متعین کردے آں۔
\end{defn}

ہِیٹھلا قضیہ وکھاؤندا اے کہ کلاسیکی نظام \Log{S5} نوں کئی ہم ارز مسلّماتی طریقیاں نال دِتّا جا سکدا اے۔ ایہ حیرانی دی گل نہیں، کیوں جے مسلّمیاں دیاں وکھ وکھ ترکیباں ساریاں ہم ارزیت دے تعلقاں نوں پچھاندیاں نیں (\olref[frd][es5]{prop:equivalences} ویکھو)۔

\begin{prop}\ollabel{prop:S5}
  $\Log{KTB4} = \Log{KT5} = \Log{KDB4} = \Log{KDB5}$.
\end{prop}

\begin{proof}
  مشق۔
\end{proof}

\begin{prob}
  \olref[nml][prf][prs]{prop:S5} ثابت کرو۔
\end{prob}
\endgroup
% /OLP-0435
% OLP-0436 content/normal-modal-logic/axioms-systems/soundness.tex
\begingroup
\def\olfilename{soundness}
\expandafter\def\csname ol@sectioncs\endcsname{section}

\olfileid{nml}{prf}{snd}

\olsection{صحت مندی}

!!^a{derivation} دا نظام صحت مند اکھواؤندا اے جے اوہدے وچ جو کجھ !!{derive} ہو سکدا اے، اوہ معتبر ہووے۔ امکانی نظاماں نوں ویکھدیاں، یعنی اوہناں !!{derivation}s نوں جیہناں وچ \Ax{K} دے نال کجھ !!{formula}s $!A_1$، \dots،~$!A_n$ دیاں مثالاں وی ورت سکدے آں، اسیں چاہندے آں کہ ہر !!{derivable} فارمولا ہر اوس نمونے وچ سچا ہووے جیہدے وچ $!A_1$، \dots، $!A_n$ سچے نیں۔

\begin{thm}[صحت مندی دا قضیہ]\ollabel{thm:soundness}
  جے $!A_1$، \dots، $!A_n$ وچوں ہر اک دی ہر مثال نمونیاں دیاں جماعتاں $\mClass{C}_1$، \dots، $\mClass{C}_n$ وچ بالترتیب معتبر اے، تاں $\Log{K}!A_1\dots !A_n
  \Proves !B$ توں ایہ نکلدا اے کہ $!B$ نمونیاں دی جماعت $\mClass{C}_1 \cap \dots \cap \mClass{C}_n$ وچ معتبر اے۔
\end{thm}

\begin{proof}
  ثبوتاں دی لمبائی اُتے استقرا نال۔ مختصر لکھت لئی منّ لو $\mClass{C} =
  \mClass{C}_1 \cap \dots \cap \mClass{C}_n$۔
  \begin{enumerate}
  \item استقرا دی بنیاد۔ جے $!B$ دا ثبوت لمبائی~$1$ دا اے، تاں اوہ یا تاں ہمیشہ صادق فارمولے توں بدل نال بنی مثال اے، یا~\Ax{K} دی مثال اے،\iftag{prvDiamond}{ یا~\Dual{} دی مثال اے،}{} یا $!A_1$، \dots،~$!A_n$ وچوں کسے اک دی مثال اے۔ پہلی صورت وچ $!B$، $\mClass{C}$ وچ معتبر اے، کیوں جے ہمیشہ صادق فارمولیاں توں بدل نال بنیاں مثالاں نمونیاں دی \emph{ہر} جماعت وچ معتبر نیں، \olref[syn][tau]{prop:valid-taut} مطابق۔ دوجی صورت وچ ایویں ای
    \olref[syn][sch]{prop:Kvalid}\iftag{prvDiamond}{ تے
      \olref[syn][sch]{prop:Dual-valid}}{} مطابق۔ آخر تِیجی صورت وچ کیوں جے $!B$، $\mClass{C}_i$ وچ معتبر اے تے $\mClass{C} \subseteq
    \mClass{C}_i$، ایس لئی $!B$، $\mClass{C}$ وچ وی معتبر اے،
    \olref[syn][val]{prop:subset-class} مطابق۔
  \item استقرائی قدم۔ منّ لو $!B$ دا ثبوت لمبائی $k>1$ دا اے۔ جے
    $!B$ ہمیشہ صادق فارمولے توں بدل نال بنی مثال اے یا $!A_1$،
    \dots، $!A_n$ وچوں کسے اک دی مثال اے، تاں پچھلے قدم وانگ اگے ودھدے آں۔ سو منّ لو $!B$،
    \MP{} نال پچھلے !!{formula}s $!C \lif !B$ تے
    $!C$ توں ملدا اے۔ فیر $!C \lif !B$ تے $!C$ دے ثبوتاں دی لمبائی $<k$ اے، تے استقرائی مفروضے مطابق اوہ~$\mClass{C}$ وچ معتبر نیں۔
    \olref[syn][sch]{prop:soundMP} مطابق $!B$، $\mClass{C}$ وچ وی معتبر اے۔ آخر منّ لو $!B$،
    \Nec{} نال $!C$ توں ملدا اے (سو
    $!B = \Box!C$)۔ استقرائی مفروضے مطابق $!C$،
    $\mClass{C}$ وچ معتبر اے، تے \olref[syn][val]{prop:Nec-rule} مطابق~$!B$ وی۔
  \end{enumerate}
\end{proof}
\endgroup
% /OLP-0436
% OLP-0437 content/normal-modal-logic/axioms-systems/systems-distinct.tex
\begingroup
\def\olfilename{systems-distinct}
\expandafter\def\csname ol@sectioncs\endcsname{section}

\olfileid{nml}{prf}{dis}

\olsection{نظاماں دا وکھرا ہونا وکھاؤنا}

\olref[prs]{sec} وچ اسیں ویکھیا سی کہ امکانی منطق دے دو نظاماں نوں اصل وچ اکّو نظام کیویں ثابت کریے۔ \olref[snd]{thm:soundness} نال اسیں ایہ وکھا سکدے آں کہ دو امکانی نظام $\Sigma$ تے $\Sigma'$ وکھرے نیں۔ ایس لئی اک ایسا !!a{formula}~$!A$ لبھدے آں جیہدے لئی $\Sigma' \Proves !A$، پر اوہ~$\Sigma$ دے کسے نمونے وچ جھوٹھا اے۔

\begin{prop}
  $\Log{KD} \subsetneq \Log{KT}$
\end{prop}

\begin{proof} ایہ اوس معنیاتی حقیقت دا نحوی جوڑی دار اے کہ ہر انعکاسی تعلق ہر جہان توں رسائی والا اے۔ $\Log{KD} \subseteq
  \Log{KT}$ وکھاؤن لئی سانوں ویکھنا اے کہ $\Log{KD} \Proves !B$ توں
  $\Log{KT} \Proves !B$ نکلدا اے۔ ایہ $\Log{KT} \Proves
  \Log{D}$ توں ملدا اے، جیویں
  \olref[prs]{prop:S5facts}\olref[prs]{prop:S5facts-KT-D} وچ وکھایا۔ ایہ وکھاؤن لئی کہ ایہ حقیقی شمولیت اے، صحت مندی مطابق
  (\olref[snd]{thm:soundness})، \Log{KD} دا ایسا نمونہ دینا کافی اے جیہدے وچ \Ax{T}، یعنی $\Box p \lif p$، جھوٹھا اے (ایہ سوکھا کم مشق لئی چھڈدے آں)۔ کیوں جے فیر صحت مندی مطابق $\Log{KD}
  \Proves/ \Box p \lif p$۔
\end{proof}

\begin{prop}
  $\Log{KB} \neq \Log{K4}$. 
\end{prop}

\begin{proof}
  اسیں اک ایسا متقابل نمونہ بناؤنے آں جیہدے وچ \Ax{4} دی کوئی مثال جھوٹھھی اے۔ ایہ مثال صاف طور تے !!{derivable} اے \Log{K4} وچ، پر \Log{KB} وچ نہیں، ایس لئی $\Log{K4} \nsubseteq \Log{KB}$ نکلے گا۔
  متقابل نمونہ $\mModel{M}$ ویکھو جیہڑا \olref{fig:Bnot4} وچ اے۔ کیوں جے نمونہ متقابل اے، \Ax{K} تے~\Ax{B}،~$\mModel{M}$ وچ سچے نیں
  (بالترتیب \olref[syn][sch]{prop:Kvalid} تے
  \olref[frd][acc]{thm:soundschemas} مطابق)۔ پر
  $\mSat/{M}{\Box p \lif \Box\Box p}[w_1]$۔
\end{proof}

\begin{figure}[htpb]
  \centering
  \begin{tikzpicture}[modal]
    \node[world] (w1) [label=above:\mFalse{p},
    label={below:$\mSat{{}}{\Box p}$\\ $\mSat/{{}}{\Box\Box p}$}] {$w_1$}; 
    \node[world] (w2) [label=above:\mTrue{p},
      label=below:$\mSat/{{}}{\Box p}$, right=of w1]{$w_2$}; 
    \draw[->, bend left] (w1) to (w2); 
    \draw[->, bend left] (w2) to (w1); 
  \end{tikzpicture}
  \caption{اک متقابل نمونہ جیہڑا \Ax{4} دی اک مثال نوں جھوٹھا کردا اے۔}
  \ollabel{fig:Bnot4}
\end{figure}

\begin{thm}\ollabel{thm:KTBnot45}
  $\Log{KTB} \Proves/ \Log{4}$ تے $\Log{KTB} \Proves/ \Log{5}$۔
\end{thm}

\begin{proof}
  \olref[frd][acc]{thm:soundschemas} توں اسیں جان دے آں کہ
  \Ax{T} تے \Ax{B} دیاں ساریاں مثالاں ہر انعکاسی متقابل نمونے وچ سچیاں نیں (بالترتیب)۔ ایس لئی صحت مندی مطابق اک ایسا انعکاسی متقابل نمونہ لبھنا کافی اے جیہدے وچ کسے جہان اُتے~\Ax{4} دی کوئی مثال جھوٹھھی ہووے، تے~\Ax{5} لئی وی ایویں۔ اسیں دوناں دعویاں لئی اکّو نمونہ ورتدے آں۔
  \olref{fig:KTBnot45} والا متقابل تے انعکاسی نمونہ ویکھو۔ فیر $\mSat/{M}{\Box p \to \Box \Box
    p}[w_1]$، ایس لئی \Ax{4}، $w_1$ اُتے جھوٹھا اے۔ ایس طرح $\mSat/{M}{\Diamond
    \lnot p \lif \Box \Diamond \lnot p}[w_2]$، سو~\Ax{5} دی $!A = \lnot p$ والی مثال~$w_2$ اُتے جھوٹھھی اے۔
\end{proof}

\begin{figure}[htpb]
  \centering
  \begin{tikzpicture}[modal]
    \node[world] (w1) [label={right:\mTrue{p}},
      label={below: $\mSat{{}}{\Box p}$\\
        $\mSat/{{}}{\Box\Box p}$\\
        $\mSat/{{}}{\Diamond\lnot p}$}] {$w_1$};
    \draw[reflexive above] (w1) to (w1);
    \node[world] (w2) [label={right:\mTrue{p}}, label={below:
        $\mSat{{}}{\Diamond\lnot p}$\\
        $\mSat/{{}}{\Box\Diamond\lnot p}$}, right=of w1] {$w_2$} ;
    \draw[reflexive above] (w2) to (w2);
    \node[world] (w3) [label={right:\mFalse{p}},right=of w2] {$w_3$};
    \draw[reflexive above] (w3) to (w3);
    \draw[->, bend left] (w1) to (w2);  
    \draw[->, bend left] (w2) to (w3);  
    \draw[->, bend left] (w3) to (w2);  
    \draw[->, bend left] (w2) to (w1);  
  \end{tikzpicture}
  \caption{\olref{thm:KTBnot45} لئی نمونہ۔}
  \ollabel{fig:KTBnot45}
\end{figure}

\begin{thm}\ollabel{thm:KD5not4}
  $\Log{KD5} \neq \Log{KT4} = \Log{S4}$.
\end{thm}

\begin{proof}
  \olref[frd][acc]{thm:soundschemas} توں اسیں جان دے آں کہ
  \Ax{D} تے~\Ax{5} دیاں ساریاں مثالاں ہر جہان توں رسائی والے سارے اقلیدسی نمونیاں وچ سچیاں نیں۔ ایس لئی ایسا ہر جہان توں رسائی والا اقلیدسی نمونہ لبھنا کافی اے جیہدے کسے جہان اُتے~\Ax{4} دی کوئی مثال جھوٹھھی ہووے۔
  \olref{fig:KD5not4} والا نمونہ ویکھو تے دھیان دیو کہ $\mSat/{M}{\Box p
  \lif \Box\Box p}[w_1]$۔
\end{proof}

\begin{figure}[t]
  \centering
  \begin{tikzpicture}[modal]
    \node[world] (w2) [label=north west:\mTrue{p}]{$w_2$} ;
    \draw[reflexive left] (w2) to (w2);
    \node[world] (w1) [label=right:\mFalse{p},
      label=below:{$\mSat{{}}{\Box p}, \mSat/{{}}{\Box\Box p}$},
      below right=of w2]{$w_1$};
    \node[world] (w3) [label=north east:\mTrue{p}, above right=of w1]{$w_3$} ;
    \draw[reflexive right] (w3) to (w3);
    \node[world] (w4) [label=right:\mFalse{p},above right=of w2]{$w_4$};
    \draw[reflexive above] (w4) to (w4);
    \draw[->] (w1) to (w2);  
    \draw[->] (w1) to (w3);  
    \draw[->, bend left=15] (w2) to (w3);
    \draw[->, bend left=15] (w3) to (w2);  
    \draw[->, bend left=15] (w2) to (w4);
    \draw[->, bend left=15] (w4) to (w2);  
    \draw[->, bend left=15] (w3) to (w4);
    \draw[->, bend left=15] (w4) to (w3);  
  \end{tikzpicture}
  \caption{\olref{thm:KD5not4} لئی نمونہ۔}
  \ollabel{fig:KD5not4}
\end{figure}

\begin{prob}
  \olref[nml][prf][dis]{thm:KD5not4} دا کوئی ہور ثبوت دیو جیہدے وچ $3$ جہاناں والا نمونہ ورتیا ہووے۔
\end{prob}

\begin{prob}
  اکّو انعکاسی متعدی نمونہ دیو جیہڑا ایہ دونیں گلاں وکھاوے۔
  $\Log{KT4} \Proves/ \Ax{B}$ تے $\Log{KT4} \Proves/ \Ax{5}$۔
\end{prob}
\endgroup
% /OLP-0437
% OLP-0438 content/normal-modal-logic/axioms-systems/provability-from-set.tex
\begingroup
\def\olfilename{provability-from-set}
\expandafter\def\csname ol@sectioncs\endcsname{section}

\olfileid{nml}{prf}{prg}

\olsection{\usetoken{S}{derivability}، \usetoken{P}{formula} دے سیٹ توں}

\olref[prf][prs]{sec} وچ اسیں !!a{formula} دے نظام~$\Sigma$ وچ ثابت ہون دی تعریف کیتی سی۔ ہُن اسیں ایس تعریف نوں ودھا کے $\Sigma$ وچ، سیٹ~$\Gamma$ دے !!{formula}s توں ثابت ہون دی تعریف دیندے آں۔

\begin{defn}\ollabel{defn:Gammaproves}
  !!^a{formula}~$!A$، !!{derivable} اے نظام~$\Sigma$ وچ، !!{formula}s دے سیٹ $\Gamma$ توں، جیہنوں $\Gamma \Proves[\Sigma] !A$ لکھدے آں، اُدوں تے صرف اُدوں جدوں ایسے $!B_1$، \dots، $!B_n \in \Gamma$ ہون کہ
  $\Sigma \Proves !B_1 \lif (!B_2 \lif \cdots (!B_n \lif !A)
  \cdots)$۔
\end{defn}
\endgroup
% /OLP-0438
% OLP-0439 content/normal-modal-logic/axioms-systems/provability-properties.tex
\begingroup
\def\olfilename{provability-properties}
\expandafter\def\csname ol@sectioncs\endcsname{section}

\olfileid{nml}{prf}{prp}

\olsection{\usetoken{S}{derivability} دیاں خاصیتاں}

\begin{prop}\ollabel{prop:derivabilityfacts}
  منّ لو $\Sigma$ امکانی نظام اے تے $\Gamma$ امکانی
  !!{formula}s دا سیٹ اے۔ ہِیٹھ لکھیاں خاصیتاں لاگو نیں۔
  \begin{enumerate}
  \item\ollabel{prop:derivabilityfacts-monotonicity} \emph{یک رُخیت}۔ جے
    $\Gamma \Proves[\Sigma] !A$ تے $\Gamma \subseteq \Delta$، تاں
    $\Delta \Proves[\Sigma] !A$؛
  \item\ollabel{prop:derivabilityfacts-reflexivity}
    \emph{انعکاس}۔ جے $!A \in \Gamma$، تاں $\Gamma
    \Proves[\Sigma] !A$؛
  \item\ollabel{prop:derivabilityfacts-cut} \emph{کٹ}۔ جے $\Gamma
    \Proves[\Sigma] !A$ تے $\Delta \cup \{!A\} \Proves[\Sigma] !B$،
    تاں $\Gamma \cup \Delta \Proves[\Sigma] !B$؛
  \item\ollabel{prop:derivabilityfacts-deduction} \emph{استنباط دا قضیہ}۔ $\Gamma \cup \{!B\}
    \Proves[\Sigma] !A$ اُدوں تے صرف اُدوں جدوں $\Gamma \Proves[\Sigma] !B \lif
      !A$؛
  \item \ollabel{prop:derivabilityfacts-ruleT}% \emph{قاعدہ ٹی}۔ جے
    $\Gamma \Proves[\Sigma] !A_1$ تے \dots تے $\Gamma
    \Proves[\Sigma] !A_n$ تے $!A_1 \to (!A_2 \lif \cdots (!A_n \lif
    !B)\cdots)$ ہمیشہ صادق فارمولے توں بدل نال بنی مثال اے، تاں $\Gamma
    \Proves[\Sigma] !B$۔
  \end{enumerate}
\end{prop}

ثبوت اک سوکھی مشق اے۔ \olref{prop:derivabilityfacts-ruleT} والا حصہ، جیہڑا \olref{prop:derivabilityfacts} وچ اے، ایہ دیندا اے کہ مثلاً جے $\Gamma \Proves[\Sigma] !A \lor !B$
تے $\Gamma \Proves[\Sigma] \lnot !A$، تاں $\Gamma \Proves[\Sigma]
!B$۔ نالے اگے اسیں $\Gamma, !A \Proves[\Sigma] !B$
لکھاں گے $\Gamma \cup \{ !A\} \Proves[\Sigma] !B$ دی تھاں۔

\begin{defn}
  سیٹ $\Gamma$ \emph{استنباط تحت بند} اے نظام
  $\Sigma$ دے لحاظ نال اُدوں تے صرف اُدوں جدوں $\Gamma \Proves[\Sigma] !A$ توں
  $!A \in \Gamma$ نکلدا اے۔
\end{defn}
\endgroup
% /OLP-0439
% OLP-0440 content/normal-modal-logic/axioms-systems/consistency.tex
\begingroup
\def\olfilename{consistency}
\expandafter\def\csname ol@sectioncs\endcsname{section}

\olfileid{nml}{prf}{con}

\olsection{بے تضادی}

بے تضادی !!{formula}s دے سیٹاں دی اک اہم خاصیت اے۔ !!{formula}s دا سیٹ تضاد والا اے جے اوہدے توں کوئی تضاد، جیویں~$\lfalse$، !!{derivable} اے۔ ہور صورت وچ ایہ بے تضاد اے۔ جے سیٹ تضاد والا اے، تاں اوہدے سارے !!{formula}s کسے نمونے دے کسے جہان وچ اکٹھے سچے نہیں ہو سکدے۔ تمامیت دے قضیے لئی اسیں اُلٹا رخ ثابت کردے آں۔ ہر بے تضاد سیٹ کسے نمونے دے کسے جہان وچ سچا ہوندا اے، یعنی «کینونیکل نمونے» وچ۔

\begin{defn}
  سیٹ $\Gamma$ نظام~$\Sigma$ دے لحاظ نال \emph{بے تضاد}، یا جیویں اسیں آکھاں گے، $\Sigma$-بے تضاد، اُدوں تے صرف اُدوں اے جدوں $\Gamma
  \Proves/[\Sigma] \lfalse$۔
\end{defn}

سو مثلاً سیٹ $\{ \Box(p \lif q), \Box p, \lnot\Box q \}$ قضیاتی منطق دے لحاظ نال بے تضاد اے، پر
\Log{K} دے لحاظ نال بے تضاد نہیں۔ ایس طرح سیٹ $\{ \Diamond p, \Box\Diamond
p \lif q, \lnot q \}$، \Log{K5} دے لحاظ نال بے تضاد نہیں۔

\begin{prop}\ollabel{prop:consistencyfacts}
  منّ لو $\Gamma$، !!{formula}s دا سیٹ اے۔ فیر
  \begin{enumerate}
  \item $\Gamma$، $\Sigma$-بے تضاد اُدوں تے صرف اُدوں اے جدوں کوئی ایسا !!{formula}~$!A$ ہووے کہ $\Gamma \Proves/[\Sigma]
    !A$۔
  \item \ollabel{prop:consistencyfacts-b}%
    $\Gamma \Proves[\Sigma] !A$ اُدوں تے صرف اُدوں جدوں $\Gamma \cup \{
    \lnot!A \}$، $\Sigma$-بے تضاد نہ ہووے۔
  \item \ollabel{prop:consistencyfacts-c}%
    جے $\Gamma$، $\Sigma$-بے تضاد اے، تاں کسے وی !!{formula}
    $!A$ لئی، یا تاں $\Gamma \cup \{ !A \}$،
    $\Sigma$-بے تضاد اے، یا $\Gamma \cup \{ \lnot!A \}$،
    $\Sigma$-بے تضاد اے۔
  \end{enumerate}
\end{prop}

\begin{proof}
  کلاسیکی قضیاتی منطق ورت کے ایہ گلاں سوکھیاں نکلدیاں نیں۔ اسیں \olref{prop:consistencyfacts-c} دی دلیل دیندے آں۔ اُلٹی نفی راہیں منّ لو $\Gamma \cup \{ !A \}$،
  $\Gamma \cup \{ \lnot!A \}$ وچوں کوئی وی $\Sigma$-بے تضاد نہیں۔ فیر
  \olref{prop:consistencyfacts-b} مطابق $\Gamma, !A \Proves[\Sigma]
  \lfalse$ تے $\Gamma, \lnot !A \Proves[\Sigma] \lfalse$ دونیں۔ استنباط دے قضیے مطابق $\Gamma \Proves[\Sigma] !A \to \lfalse$ تے
  $\Gamma \Proves[\Sigma] \lnot!A \lif \lfalse$۔ پر $(!A \lif
  \lfalse) \lif ((\lnot!A \lif \lfalse) \lif \lfalse)$ ہمیشہ صادق فارمولے توں بدل نال بنی مثال اے، سو
  \olref[prp]{prop:derivabilityfacts}\olref[prp]{prop:derivabilityfacts-ruleT} مطابق،
  $\Gamma \Proves[\Sigma] \lfalse$۔
\end{proof}
\endgroup
% /OLP-0440

\OLEndChapterHook
\endgroup
% /OLP-0427

% OLP-0441 content/normal-modal-logic/completeness/completeness.tex
\begingroup
\def\olfilename{completeness}
\expandafter\def\csname ol@sectioncs\endcsname{section}

\olchapter{nml}{com}{تمامیت تے کینونیکل نمونے}

% OLP-0442 content/normal-modal-logic/completeness/introduction.tex
\begingroup
\def\olfilename{introduction}
\expandafter\def\csname ol@sectioncs\endcsname{section}

\olfileid{nml}{com}{int}

\olsection{جان پہچان}

جے $\Sigma$ اک امکانی نظام اے، تاں صحت مندی دا قضیہ ثابت کردا اے کہ جے $\Sigma \Proves !A$، تاں $!A$ نمونیاں دی ہر اوس جماعت
$\mClass{C}$ وچ معتبر اے جیہدے وچ~$\Sigma$ دے سارے !!{formula}s دیاں ساریاں مثالاں معتبر نیں۔ خاص طور تے ایس دا مطلب اے کہ جے $\Log{K}
\Proves !A$، تاں $!A$ سارے نمونیاں وچ سچا اے؛ جے $\Log{KT} \Proves !A$،
تاں $!A$ سارے انعکاسی نمونیاں وچ سچا اے؛ جے $\Log{KD} \Proves !A$،
تاں $!A$ ہر جہان توں رسائی والے سارے نمونیاں وچ سچا اے، تے اگے وی ایویں۔

تمامیت صحت مندی دا اُلٹا رخ اے۔ مثلاً \Log{K} دے مکمل ہون دا مطلب اے کہ جے کوئی !!a{formula}~$!A$ معتبر اے، تاں $\Proves !A$۔
تمامیت ثابت کرنا صحت مندی ثابت کرن نالوں کِتے اوکھا اے۔ پہلاں اُلٹی نفی والی صورت ویکھنی کم آؤندی اے۔ \Log{K} اُدوں تے صرف اُدوں مکمل اے جے جدوں وی $\Proves/ !A$ ہووے، تاں اک ردّی نمونہ ہووے، یعنی نمونہ~$\mModel{M}$ ایسا ہووے کہ $\mSat/{M}{!A}$۔ ایہدے ہم ارز طور تے ($!A$ دی نفی کر کے)، اسیں ثابت کر سکدے آں کہ جدوں وی $\Proves/ \lnot !A$، تاں~$!A$ دا اک نمونہ اے۔ ایسا نمونہ بناؤندیاں اسیں $!A$ وچلی جانکاری ورت سکدے آں۔ خاص !!{formula}s لئی نمونے لبھدیاں اسیں اکثر ایویں ای کردے آں۔ مثلاً جے $p \lif \Box q$ دا ردّی نمونہ لبھنا اے، تاں اسیں جان دے آں کہ اوہدے وچ اک جہان ہونا چاہیدا اے جتھے $p$ سچا تے $\Box q$ جھوٹھا اے۔ تے $\Box
q$ دا کسے جہان وچ جھوٹھا ہونا ایہ منگدا اے کہ اوس جہان توں اک ایسا جہان رسائی وچ ہووے جتھے $q$ جھوٹھا اے۔ سانوں بس اینی جانکاری چاہیدی اے کہ کہڑے جہان !!{propositional variable}s نوں سچا کردے نیں، تے کہڑے جہان کنہاں جہاناں توں رسائی وچ نیں۔

پر تمامیت ثابت کردیاں ساڈے کول کوئی خاص !!{formula}~$!A$ نہیں جیہدے لئی نمونہ بنا رہے ہوئیے۔ اسیں ثابت کرنا چاہندے آں کہ ہر $!A$ لئی نمونہ اے جیہدے لئی
$\Proves/[\Sigma] \lnot !A$۔ ایہ گھٹ توں گھٹ شرط اے، کیوں جے
\emph{جے} $\Proves[\Sigma] \lnot !A$، تاں صحت مندی مطابق~$!A$ دا کوئی نمونہ نہیں (جیہدے وچ $\Sigma$ سچا اے)۔ ہُن دھیان دیو کہ $\Proves/[\Sigma]
\lnot !A$ اُدوں تے صرف اُدوں جدوں $!A$، $\Sigma$-بے تضاد اے۔ (یاد کرو کہ $\Sigma
\Proves/[\Sigma] \lnot !A$ تے $!A \Proves/[\Sigma] \lfalse$ ہم ارز نیں۔) سو ساڈا کم ہر $\Sigma$-بے تضاد !!{formula} لئی نمونہ بناؤنا اے۔

ساڈا طریقہ ایہ اے کہ اک $\Sigma$-بے تضاد سیٹ !!{formula}s دا لبھیے جیہدے وچ~$!A$ وی ہووے تے ہور فارمولے وی ہوون جیہڑے سانوں دسن کہ $!A$ نوں سچا کرن والا جہان کیہو جیہا ہونا چاہیدا اے۔ ایسے سیٹ \emph{مکمل} $\Sigma$-بے تضاد سیٹ نیں۔~$!A$ نوں سچا کرن لئی صرف اک جہان والا نمونہ بناؤنا کافی نہیں۔ نمونے وچ کئی جہان تے رسائی دا تعلق ہونا پئے گا۔مکمل $\Sigma$-بے تضاد سیٹ جیہڑا~$!A$ رکھدا اے، اوہدے وچ ہور !!{formula}s، $\Box !B$ تے~$\Diamond !C$ دی شکل دے، وی ہوون گے۔ سارے رسائی والے جہاناں وچ $!B$ سچا ہونا چاہیدا اے، تے گھٹ توں گھٹ اک وچ $!C$ سچا ہونا چاہیدا اے۔ ایہ کم کرن لئی اسیں \emph{سارے} ممکن مکمل $\Sigma$-بے تضاد سیٹ جہاناں دے سیٹ دی بنیاد بنا لواں گے۔ اک اوکھا کم ایہ لبھنا ہووے گا کہ ساڈے نمونے وچ اک مکمل $\Sigma$-بے تضاد سیٹ دوجے توں کدوں رسائی وچ منّیا جاوے۔

اسیں وکھاواں گے کہ ایویں متعین کیتے نمونے وچ $!A$ کسے جہان وچ، جیہڑا آپ وی اک مکمل $\Sigma$-بے تضاد سیٹ اے، اُدوں تے صرف اُدوں سچا اے جدوں $!A$ اوس سیٹ دا !!a{element} اے۔ جے $!A$، $\Sigma$-بے تضاد اے، تاں اوہ گھٹ توں گھٹ اک مکمل $\Sigma$-بے تضاد سیٹ دا !!a{element} ہووے گا (ایہ گل اسیں ثابت کراں گے)، تے ایس لئی اک ایسا جہان ہووے گا جتھے $!A$~سچا اے۔ سو ساڈے کول اکّو نمونہ ہووے گا جیہدے وچ ہر $\Sigma$-بے تضاد !!{formula}~$!A$ کسے جہان وچ سچا اے۔ ایہ اکّو نمونہ~$\Sigma$ دا \emph{کینونیکل} نمونہ اے۔
\endgroup
% /OLP-0442
% OLP-0443 content/normal-modal-logic/completeness/complete-consistent-sets.tex
\begingroup
\def\olfilename{complete-consistent-sets}
\expandafter\def\csname ol@sectioncs\endcsname{section}

\olfileid{nml}{com}{ccs}

\olsection{\texorpdfstring{مکمل $\Sigma$-بے تضاد سیٹ}{مکمل Σ-بے تضاد سیٹ}}

منّ لو $\Sigma$ امکانی !!{formula}s دا سیٹ اے۔ اینہاں نوں نارمل امکانی منطق دے مسلّمے یا اوہنوں متعین کرن والے اصول سمجھو۔ سیٹ $\Gamma$، $\Sigma$-بے تضاد اُدوں تے صرف اُدوں اے جدوں $\Gamma \Proves/[\Sigma] \lfalse$۔ یعنی کوئی !!{derivation} نہ ہووے~$!A_1 \lif (!A_2 \lif \cdots (!A_n \lif
\lfalse)\dots)$ دا، $\Sigma$ توں، جتھے ہر $!A_i \in \Gamma$۔ اسیں اک «کینونیکل» نمونہ بناواں گے جیہدے وچ ہر جہان اک خاص طرح دا $\Sigma$-بے تضاد سیٹ اے۔ ایہ سیٹ $\Sigma$-بے تضاد اے تے ایس معنی وچ ودھ توں ودھ بے تضاد اے کہ ہر امکانی !!{formula} دی صداقتی قدر طے کردا اے۔ یعنی ہر $!A$ لئی یا تاں $!A
\in \Gamma$ یا $\lnot !A \in \Gamma$۔

\begin{defn}
  سیٹ $\Gamma$ \emph{مکمل $\Sigma$-بے تضاد} اُدوں تے صرف اُدوں اے جدوں اوہ $\Sigma$-بے تضاد ہووے تے ہر~$!A$ لئی یا تاں
  $!A \in \Gamma$ ہووے یا $\lnot !A \in \Gamma$۔
\end{defn}

مکمل $\Sigma$-بے تضاد سیٹاں $\Gamma$ دیاں کئی کم دیاں خاصیتاں نیں۔ اک ایہ کہ اوہ استنباط تحت بند نیں، یعنی جے $\Gamma
\Proves[\Sigma] !A$، تاں $!A \in \Gamma$۔ خاص طور تے ایہدا مطلب اے کہ کسے !!a{formula}~$!A \in \Sigma$ دی ہر مثال وی $\in
\Gamma$ اے۔ نالے $\Gamma$ دی رکنیت قضیاتی رابطیاں دیاں سچائی دیاں شرطاں ورگی اے۔ جدوں اسیں «کینونیکل نمونہ» متعین کراں گے، ایہ گل کم آوے گی۔

\begin{prop}\ollabel{prop:ccs-properties}
  منّ لو $\Gamma$ مکمل $\Sigma$-بے تضاد اے۔ فیر
  \begin{enumerate}
  \item \ollabel{prop:ccs-closed}%
    $\Gamma$،~$\Sigma$ وچ استنباط تحت بند اے۔
  \item \ollabel{prop:ccs-sigma}%
    $\Sigma \subseteq \Gamma$۔
  \tagitem{prvFalse}{\ollabel{prop:ccs-lfalse}%
    $\lfalse \notin \Gamma$}{}
  \tagitem{prvTrue}{\ollabel{prop:ccs-ltrue}%
    $\ltrue \in \Gamma$}{}
  \item \ollabel{prop:ccs-lnot}%
    $\lnot!A \in \Gamma$ اُدوں تے صرف اُدوں جدوں $!A \notin
    \Gamma$۔
  \tagitem{prvAnd}{\ollabel{prop:ccs-land}%
    $!A \land !B \in \Gamma$ اُدوں تے صرف اُدوں جدوں $!A \in \Gamma$ تے $!B \in \Gamma$}{}
  \tagitem{prvOr}{\ollabel{prop:ccs-lor}%
    $!A \lor !B \in \Gamma$ اُدوں تے صرف اُدوں جدوں $!A \in \Gamma$ یا $!B \in \Gamma$}{}
  \tagitem{prvIf}{\ollabel{prop:ccs-lif}%
    $!A \lif !B \in \Gamma$ اُدوں تے صرف اُدوں جدوں $!A \notin \Gamma$ یا $!B \in \Gamma$}{}
  \tagitem{prvIff}{\ollabel{prop:ccs-liff}%
    $!A \liff !B \in \Gamma$ اُدوں تے صرف اُدوں جدوں یا تاں $!A \in \Gamma$ تے $!B \in
    \Gamma$ دونیں ہون، یا $!A \notin \Gamma$ تے $!B \notin \Gamma$ دونیں ہون}{}
  \end{enumerate}
\end{prop}

\begin{proof}
  \begin{enumerate}
  \item منّ لو $\Gamma \Proves[\Sigma] !A$ پر $!A \notin
    \Gamma$۔ فیر کیوں جے $\Gamma$ مکمل $\Sigma$-بے تضاد اے، ایس لئی
    $\lnot!A \in \Gamma$۔ ایس نال $\Gamma$ تضاد والا ہو جاندا، کیوں جے
    $!A, \lnot !A \Proves[\Sigma] \lfalse$۔

  \item جے $!A \in \Sigma$ تاں $\Gamma \Proves[\Sigma] !A$، تے $!A
    \in \Gamma$ استنباط تحت بندش کارن، یعنی
    معاملہ~\olref{prop:ccs-closed} مطابق۔

  \tagitem{prvFalse}{جے $\lfalse \in \Gamma$ تاں $\Gamma
    \Proves[\Sigma] \lfalse$، سو $\Gamma$،
    $\Sigma$ دے لحاظ نال تضاد والا ہو جاندا۔}{}

  \tagitem{prvTrue}{$\Gamma \Proves[\Sigma] \ltrue$، سو $\ltrue \in
    \Gamma$ استنباط تحت بندش کارن، یعنی
    معاملہ~\olref{prop:ccs-closed} مطابق۔}{}

\item جے $\lnot!A \in \Gamma$، تاں بے تضادی کارن $!A \notin
    \Gamma$؛ تے جے $!A \notin \Gamma$ تاں $!A \in \Gamma$، کیوں جے
    $\Gamma$ مکمل $\Sigma$-بے تضاد اے۔

  \tagitem{prvAnd}{\iftag{probAnd}{مشق۔}{منّ لو $!A \land !B \in
      \Gamma$۔ کیوں جے $(!A \land !B) \lif !A$ ہمیشہ صادق فارمولے توں بدل نال بنی مثال اے، ایس لئی $!A \in \Gamma$ استنباط تحت بندش کارن، یعنی معاملہ~\olref{prop:ccs-closed} مطابق۔ ایس طرح $!B \in \Gamma$ وی۔ دوجے پاسے منّ لو $!A \in \Gamma$ تے $!B \in
      \Gamma$ دونیں۔ فیر استنباط تحت بندش توں $(!A \land !B) \in
      \Gamma$ نکلدا اے، کیوں جے $!A \lif (!B \lif (!A \land !B))$ ہمیشہ صادق فارمولے توں بدل نال بنی مثال اے۔}}{}

  \tagitem{prvOr}{\iftag{probOr}{مشق۔}{منّ لو $!A \lor !B \in
      \Gamma$، تے $!A \notin \Gamma$ تے $!B \notin \Gamma$۔ کیوں جے
      $\Gamma$ مکمل $\Sigma$-بے تضاد اے، ایس لئی $\lnot!A \in \Gamma$
      تے $\lnot !B \in \Gamma$۔ فیر $\lnot(!A \lor !B) \in \Gamma$
      کیوں جے $\lnot !A \lif (\lnot !B \lif \lnot (!A \lor !B))$ ہمیشہ صادق فارمولے توں بدل نال بنی مثال اے۔ ایس دا مطلب ہووے گا کہ $\Gamma$، $\Sigma$ دے لحاظ نال تضاد والا اے۔ ایہ تضاد اے۔}}{}

  \tagitem{prvIf}{\iftag{probIf}{مشق۔}{منّ لو $!A \lif !B \in
      \Gamma$ تے $!A \in \Gamma$؛ فیر $\Gamma \Proves[\Sigma] !B$، جیہدے توں $!B \in \Gamma$ نکلدا اے استنباط تحت بندش کارن۔ اُلٹے پاسے جے $!A
      \lif !B \notin \Gamma$، تاں کیوں جے $\Gamma$ مکمل
      $\Sigma$-بے تضاد اے، ایس لئی $\lnot (!A \lif !B) \in \Gamma$۔ کیوں جے
      $\lnot(!A \lif !B) \lif !A$ ہمیشہ صادق فارمولے توں بدل نال بنی مثال اے، ایس لئی $!A \in
      \Gamma$ استنباط تحت بندش کارن۔ کیوں جے $\lnot(!A \lif !B) \lif
      \lnot !B$ ہمیشہ صادق فارمولے توں بدل نال بنی مثال اے، ایس لئی $\lnot !B \in
      \Gamma$۔ فیر $!B \notin \Gamma$ کیوں جے $\Gamma$،
      $\Sigma$-بے تضاد اے۔}}{}

  \tagitem{prvIff}{\iftag{probIff}{مشق۔}{منّ لو $!A \liff !B \in
      \Gamma$۔ جے $!A \in \Gamma$ تاں $!B \in \Gamma$، کیوں جے $(!A
      \liff !B) \lif (!A \lif !B)$ ہمیشہ صادق فارمولے توں بدل نال بنی مثال اے۔ ایس طرح جے $!B \in \Gamma$ تاں $!A \in \Gamma$۔
      سو یا تاں $!A \in \Gamma$ تے $!B \in \Gamma$ دونیں، یا نہ
      $!A \in \Gamma$ تے نہ $!B \in \Gamma$۔

      اُلٹے پاسے منّ لو $!A \lif !B \notin \Gamma$۔ کیوں جے $\Gamma$
      مکمل $\Sigma$-بے تضاد اے، ایس لئی $\lnot (!A \liff !B) \in
      \Gamma$۔ کیوں جے $\lnot(!A \liff !B) \lif (!A \lif \lnot !B)$ ہمیشہ صادق فارمولے توں بدل نال بنی مثال اے، ایس لئی جے $!A \in \Gamma$ تاں $\lnot !B \in
      \Gamma$، تے کیوں جے $\Gamma$، $\Sigma$-بے تضاد اے، ایس لئی $!B \notin
      \Gamma$۔ ایس طرح جے $!B \in \Gamma$، تاں $!A \notin \Gamma$۔
      سو نہ $!A \in \Gamma$ تے $!B \in \Gamma$ اکٹھے نیں، تے نہ $!A \notin
      \Gamma$ تے $!B \notin \Gamma$ اکٹھے نیں۔}}{}
  \end{enumerate}
\end{proof}

\begin{probtag}{probAnd,probOr,probIf,probIff}
  \olref[nml][com][ccs]{prop:ccs-properties} دا ثبوت پورا کرو۔
\end{probtag}
\endgroup
% /OLP-0443
% OLP-0444 content/normal-modal-logic/completeness/lindenbaums-lemma.tex
\begingroup
\def\olfilename{lindenbaums-lemma}
\expandafter\def\csname ol@sectioncs\endcsname{section}

\olfileid{nml}{com}{lin}

\olsection{لنڈن باؤم دا معاون قضیہ}

لنڈن باؤم دا معاون قضیہ ثابت کردا اے کہ !!{formula}s دا ہر $\Sigma$-بے تضاد سیٹ گھٹ توں گھٹ اک \emph{مکمل}
$\Sigma$-بے تضاد سیٹ وچ شامل اے۔ کینونیکل نمونے دی ساڈی بناوٹ وکھاوے گی کہ ہر مکمل $\Sigma$-بے تضاد سیٹ~$\Delta$ لئی کینونیکل نمونے وچ اک ایسا جہان اے جتھے ~$\Delta$ دے سارے !!{formula}s تے صرف اوہو، سچے نیں۔ سو لنڈن باؤم دا معاون قضیہ پکا کردا اے کہ ہر $\Sigma$-بے تضاد سیٹ کینونیکل نمونے دے کسے جہان وچ سچا اے۔

\begin{thm}[لنڈن باؤم دا معاون قضیہ]\ollabel{thm:lindenbaum}
  جے $\Gamma$، $\Sigma$-بے تضاد اے، تاں اک مکمل
  $\Sigma$-بے تضاد سیٹ $\Delta$ اے جیہڑا~$\Gamma$ نوں ودھاؤندا اے۔
\end{thm}

\begin{proof}
منّ لو $!A_0$، $!A_1$، \dots{} زبان دے سارے فارمولیاں دی پوری فہرست اے (دہراؤ دی اجازت اے)۔ مثلاً $\Obj p_0$ توں فہرست شروع کرو، تے ہر مرحلے $n \ge 1$ اُتے لمبائی~$n$ والے اوہ سارے محدود گنتی دے فارمولے لکھو جیہناں وچ صرف $\Obj p_0$،
\dots،~$\Obj p_n$ وچوں متغیر آؤندے نیں۔ اسیں !!{formula}s دے سیٹ $\Delta_n$،~$n$ اُتے استقرا نال متعین کردے آں، تے فیر $\Delta = \bigcup_n \Delta_n$ رکھدے آں۔ پہلاں $\Delta_0 = \Gamma$ رکھو۔ منّ لو $\Delta_n$ متعین ہو چکیا اے۔ $\Delta_{n+1}$ دی تعریف ایویں کردے آں۔
\[
\Delta_{n+1} =
\begin{cases}
  \Delta_n \cup \{!A_n\}, & \text{جے $\Delta_n \cup \{ !A_n\}$،
    $\Sigma$-بے تضاد اے؛} \\
  \Delta_n \cup \{ \lnot !A_n\}, & \text{ہور صورت وچ۔}
\end{cases}
\]
ہُن $\Delta = \bigcup_{n=0}^\infty \Delta_n$ رکھو۔

سانوں وکھاؤنا اے کہ ایس تعریف توں سچ مچ ایسا سیٹ $\Delta$ ملدا اے جیہدے وچ مطلوب خاصیتاں نیں، یعنی $\Gamma \subseteq \Delta$ تے
$\Delta$ مکمل $\Sigma$-بے تضاد اے۔

صاف اے کہ $\Gamma \subseteq \Delta$، کیوں جے $\Delta_0 \subseteq
\Delta$ بناوٹ مطابق، تے $\Delta_0 = \Gamma$۔ اصل وچ $\Delta_n
\subseteq \Delta$ ہر~$n$ لئی، کیوں جے $\Delta$ سارے~$\Delta_n$ دا اکٹھ اے۔ (بناوٹ دے ہر قدم وچ اسیں پہلاں بنے سیٹ وچ اک !!a{formula} پاؤندے آں، ایس لئی $\Delta_n \subseteq
\Delta_{n+1}$۔ سو کیوں جے $\subseteq$ متعدی اے، ایس لئی $\Delta_n \subseteq
\Delta_{m}$ جدوں وی $n \le m$۔) بناوٹ دے ہر مرحلے اسیں یا تاں $!A_n$ یا $\lnot !A_n$ پاؤندے آں، تے ہر !!{formula} سارے~$!A_n$ دی فہرست وچ (گھٹ توں گھٹ اک واری) آؤندا اے۔ سو ہر $!A$ لئی یا تاں $!A
\in \Delta$ یا $\lnot !A \in \Delta$، ایس لئی تعریف مطابق $\Delta$ مکمل اے۔

آخر وچ سانوں وکھاؤنا اے کہ $\Delta$، $\Sigma$-بے تضاد اے۔ ایس لئی اسیں دو گلاں وکھاؤنے آں۔ (الف) جے $\Delta$، $\Sigma$ دے لحاظ نال تضاد والا ہوندا، تاں کوئی $\Delta_n$ وی $\Sigma$ دے لحاظ نال تضاد والا ہوندا، تے (ب) سارے $\Delta_n$،
$\Sigma$-بے تضاد نیں۔

سو منّ لو $\Delta$، $\Sigma$ دے لحاظ نال تضاد والا ہوندا۔ فیر $\Delta
\Proves[\Sigma] \lfalse$، یعنی ایسے $!A_1$، \dots،~$!A_k \in
\Delta$ نیں جیہناں لئی $\Sigma \Proves !A_1 \lif (!A_2 \lif \cdots (!A_k
\lif \lfalse)\dots)$۔ کیوں جے $\Delta = \bigcup_{n=0}^\infty \Delta_n$، ایس لئی ہر
$!A_i \in \Delta_{n_i}$ کسے~$n_i$ لئی۔ منّ لو $n$ ایہناں وچوں سبھ توں وڈا اے۔ کیوں جے $n_i \le n$، ایس لئی $\Delta_{n_i} \subseteq \Delta_n$۔ سو سارے
$!A_i$ کسے~$\Delta_n$ وچ نیں۔ ایس دا مطلب ہووے گا $\Delta_n
\Proves[\Sigma] \lfalse$، یعنی $\Delta_n$، $\Sigma$ دے لحاظ نال تضاد والا اے۔

ہر $\Delta_n$ نوں $\Sigma$-بے تضاد وکھاؤن لئی اسیں~$n$ اُتے سادا استقرا ورتدے آں۔ $\Delta_0 = \Gamma$، تے اسیں منّیا سی کہ $\Gamma$،
$\Sigma$-بے تضاد اے۔ سو دعویٰ $n = 0$ لئی سچا اے۔ ہُن منّ لو دعویٰ
$n$ لئی سچا اے، یعنی $\Delta_n$، $\Sigma$-بے تضاد اے۔ $\Delta_{n+1}$ یا تاں
$\Delta_n \cup \{!A_n\}$ اے جے اوہ $\Sigma$-بے تضاد اے، ہور صورت وچ $\Delta_n \cup \{\lnot!A_n\}$ اے۔ پہلی صورت وچ
$\Delta_{n+1}$ صاف طور تے $\Sigma$-بے تضاد اے۔ پر
\olref[prf][con]{prop:consistencyfacts}\olref[prf][con]{prop:consistencyfacts-c} مطابق،
$\Delta_n \cup \{!A_n\}$ یا $\Delta_n \cup \{\lnot!A_n\}$ وچوں کوئی اک بے تضاد اے، سو $\Delta_{n+1}$ دوجی صورت وچ وی بے تضاد اے۔
\end{proof}

\begin{cor}\ollabel{cor:provability-characterization}
  $\Gamma \Proves[\Sigma] !A$ اُدوں تے صرف اُدوں جدوں $!A \in \Delta$ ہر اوس مکمل $\Sigma$-بے تضاد سیٹ $\Delta$ لئی ہووے جیہڑا $\Gamma$ نوں ودھاؤندا اے (اوہدوں وی جدوں $\Gamma = \emptyset$۔ اوس صورت وچ امکانی نظام $\Sigma$ دی اک ہور پچھان ملدی اے)۔
\end{cor}

\begin{proof}
  منّ لو $\Gamma \Proves[\Sigma] !A$، تے $\Delta$ کوئی وی مکمل $\Sigma$-بے تضاد سیٹ ہووے جیہڑا $\Gamma$ نوں ودھاؤندا اے۔ جے $!A
  \notin \Delta$ تاں ودھ توں ودھ بے تضادی کارن $\lnot!A \in \Delta$، تے ایس لئی
  $\Delta \Proves[\Sigma] !A$ (یک رُخیت کارن) تے $\Delta
  \Proves[\Sigma] \lnot!A$ (انعکاس کارن)، سو $\Delta$ تضاد والا اے۔ اُلٹے پاسے جے $\Gamma \Proves/[\Sigma] !A$، تاں
  $\Gamma \cup \{ \lnot!A \}$، $\Sigma$-بے تضاد اے، تے لنڈن باؤم دے معاون قضیے مطابق اک مکمل بے تضاد سیٹ $\Delta$ اے جیہڑا $\Gamma \cup \{ \lnot!A \}$ نوں ودھاؤندا اے۔ بے تضادی کارن $!A
  \notin \Delta$۔
\end{proof}
\endgroup
% /OLP-0444
% OLP-0445 content/normal-modal-logic/completeness/modalities-ccs.tex
\begingroup
\def\olfilename{modalities-ccs}
\expandafter\def\csname ol@sectioncs\endcsname{section}

\olfileid{nml}{com}{mod}

\olsection{امکانی عامل تے مکمل بے تضاد سیٹ}

\begin{explain}
جدوں اسیں نمونہ~$\mModel{M^\Sigma}$ بناؤندے آں، جیہدے جہاناں دا سیٹ مکمل $\Sigma$-بے تضاد سیٹاں~$\Delta$ نال دِتّا جاندا اے، کسے نارمل امکانی منطق~$\Sigma$ وچ، تاں اینہاں «جہاناں» وچ رسائی دا تعلق~$R^\Sigma$ وی متعین کرنا پئے گا۔ اسیں چاہندے آں کہ رسائی دا تعلق (تے قدر بندی $V^\Sigma)$ ایویں متعین ہون کہ $\mSat{M^\Sigma}{!A}[\Delta]$ اُدوں تے صرف اُدوں جدوں $!A \in
\Delta$۔ ایہ کیویں کریے؟

رسائی دا تعلق متعین ہو جان تے جہان وچ سچائی دی تعریف ایہ پکا کردی اے کہ \iftag{prvBox}{$\mSat{M^\Sigma}{\Box!A}[\Delta]$
  اُدوں تے صرف اُدوں جدوں $\mSat{M^\Sigma}{!A}[\Delta']$ ہر ایسے $\Delta'$ لئی ہووے جیہدے لئی
  $R^\Sigma\Delta\Delta'$ ہووے}{$\mSat{M^\Sigma}{\Diamond!A}[\Delta]$ اُدوں تے صرف اُدوں جدوں
  $\mSat{M^\Sigma}{!A}[\Delta']$ گھٹ توں گھٹ اک ایسے $\Delta'$ لئی ہووے جیہدے لئی
  $R^\Sigma\Delta\Delta'$ ہووے}۔ ایہ ثابت کرن لئی کہ
$\mSat{M^\Sigma}{!A}[\Delta]$ اُدوں تے صرف اُدوں جدوں $!A \in \Delta$، ضروری اے کہ ایہ گل خاص طور تے امکانی عامل نال شروع ہون والے !!{formula}s لئی وی سچی ہووے، یعنی \iftag{prvBox}{$\mSat{M^\Sigma}{\Box!A}[\Delta]$ اُدوں تے صرف اُدوں جدوں $\Box !A
  \in \Delta$}{$\mSat{M^\Sigma}{\Diamond!A}[\Delta]$ اُدوں تے صرف اُدوں جدوں $\Diamond !A
  \in \Delta$}۔ ایس شرط نوں جہان وچ سچائی دی تعریف نال، \iftag{prvBox}{$\Box !A$}{$\Diamond !A$} لئی، جوڑ کے ایہ ملدا اے۔
\[
\iftag{prvBox}{%
  \Box!A \in \Delta \text{ اُدوں تے صرف اُدوں جدوں } !A \in \Delta' \text{
    ہر ایسے $\Delta'$ لئی جیہدے لئی } R^\Sigma\Delta\Delta'}{%
  \Diamond!A \in \Delta \text{ اُدوں تے صرف اُدوں جدوں } !A \in \Delta' \text{ گھٹ توں گھٹ اک ایسے }
    \Delta' \text{ لئی جیہدے لئی } R^\Sigma\Delta\Delta'}
\]
\iftag{prvBox}{کھبے توں سجے والا رخ ویکھو۔ ایہ آکھدا اے کہ جے
  $\Box !A \in \Delta$، تاں $!A \in \Delta'$، کسے وی $!A$ تے کسے وی ایسے
  $\Delta'$ لئی جیہدے لئی $R^\Sigma\Delta\Delta'$ ہووے۔ جے اسیں ایہ شرط رکھیے کہ
  $R^\Sigma\Delta\Delta'$ اُدوں تے صرف اُدوں جدوں $!A \in \Delta'$ ہر $\Box!A \in
  \Delta$ لئی ہووے، تاں ایہ گل پوری ہو جاندی اے۔ «اُدوں تے صرف اُدوں» دے سجے پاسے والی شرط اسیں مختصر ایویں لکھ سکدے آں۔ $\Setabs{!A}{\Box!A \in \Delta}
  \subseteq \Delta'$۔}{سجے توں کھبے والا رخ ویکھو۔ ایہ آکھدا اے
  کہ جے $!A \in \Delta'$، تاں $\Diamond!A \in \Delta$، کسے وی $!A$ تے
  ایسے $\Delta'$ لئی جیہدے لئی $R^\Sigma\Delta\Delta'$ ہووے۔ جے اسیں ایہ شرط رکھیے کہ
  $R^\Sigma\Delta\Delta'$ جدوں وی $\Diamond!A \in \Delta$ ہر
  $!A \in \Delta'$ لئی ہووے، تاں ایہ گل پوری ہو جاندی اے۔ «اُدوں تے صرف اُدوں» دے سجے پاسے والی شرط اسیں مختصر ایویں لکھ سکدے آں۔
  $\Setabs{\Diamond!A}{!A \in \Delta'} \subseteq \Delta$۔}

سو سوال ایہ اے کہ $R^\Sigma$ دی ایہ تعریف سچ مچ پکا کردی اے کہ \iftag{prvBox}{$\Box !A \in \Delta$ اُدوں تے صرف اُدوں جدوں
  $\mSat{M^\Sigma}{\Box !A}[\Delta]$؟ \iftag{prvDiamond}{کی ایہ گل وی پکی اے کہ }{}}{}\iftag{prvDiamond}{$\Diamond !A \in \Delta$
  اُدوں تے صرف اُدوں جدوں $\mSat{M^\Sigma}{\Diamond !A}[\Delta]$؟}{} اگلے کجھ نتیجے ایہ گل قائم کرن گے۔
\end{explain}

\begin{defn}
  جے $\Gamma$، !!{formula}s دا سیٹ اے، تاں ایہ لکھو۔
  \begin{align*}
    \Box\Gamma & = \Setabs{\Box !B}{!B \in \Gamma}\\
    \Diamond\Gamma & = \Setabs{\Diamond !B}{!B \in \Gamma}\\
    \intertext{تے}
    \Box^{-1}\Gamma & = \Setabs{!B}{\Box !B \in \Gamma}\\
    \Diamond^{-1}\Gamma & = \Setabs{!B}{\Diamond !B \in \Gamma}\\
  \end{align*}
\end{defn}

دوجے لفظاں وچ $\Box\Gamma$، $\Gamma$ دا اوہ روپ اے جتھے $\Box$ ہر !!{formula} دے اگے اے جیہڑا~$\Gamma$ وچ اے۔ $\Box^{-1}\Gamma$ سارے اوہناں $\Box$ والے !!{formula}s توں بندا اے جیہڑے~$\Gamma$ وچ نیں، جدوں اوہناں دے شروع والا $\Box$ ہٹا دیندے آں۔ ایس تعریف دی اپنی کوئی بہتی اہمیت نہیں، پر ایس نال علامتی لکھت کِتے سوکھی ہو جاوے گی۔

دھیان دیو کہ $\Box\Box^{-1}\Gamma \subseteq \Gamma$۔
\[
\Box\Box^{-1}\Gamma = \Setabs{\Box!B}{\Box!B \in \Gamma}
\]
یعنی ایہ صرف~$\Gamma$ دے سارے اوہ !!{formula}s دا سیٹ اے جیہڑے~$\Box$ نال شروع ہوندے نیں۔

\begin{lem}\ollabel{lem:box1}
  جے $\Gamma \Proves[\Sigma] !A$، تاں $\Box\Gamma \Proves[\Sigma] \Box
  !A$۔
\end{lem}

\begin{proof}
  جے $\Gamma \Proves[\Sigma] !A$، تاں ایسے $!B_1$، \dots، $!B_k
  \in \Gamma$ نیں کہ $\Sigma \Proves !B_1 \lif (!B_2 \lif \cdots
  (!B_n \lif !A)\cdots)$۔ کیوں جے $\Sigma$ نارمل اے، قاعدہ \RK{} مطابق
  $\Sigma \Proves \Box!B_1 \lif (\Box!B_2 \lif \cdots (\Box!B_n \lif
  \Box!A)\cdots)$، جتھے صاف اے کہ $\Box!B_1$، \dots، $\Box!B_k \in
  \Box\Gamma$۔ سو تعریف مطابق $\Box\Gamma \Proves[\Sigma] \Box
  !A$۔
\end{proof}

\begin{lem}\ollabel{lem:box2}
  جے $\Box^{-1} \Gamma \Proves[\Sigma] !A$، تاں $\Gamma
  \Proves[\Sigma] \Box!A$۔
\end{lem}

\begin{proof}
  منّ لو $\Box^{-1}\Gamma \Proves[\Sigma] !A$۔ فیر
  \olref{lem:box1} مطابق $\Box\Box^{-1}\Gamma \Proves \Box!A$۔ پر کیوں جے
  $\Box\Box^{-1}\Gamma \subseteq \Gamma$، ایس لئی یک رُخیت کارن $\Gamma \Proves[\Sigma]
  \Box!A$ وی۔
\end{proof}

\iftag{prvBox}{%
% لازمیت دا نشان ابتدائی اے۔

\begin{prop}\ollabel{prop:box}
  جے $\Gamma$ مکمل $\Sigma$-بے تضاد اے، تاں $\Box !A \in
  \Gamma$ اُدوں تے صرف اُدوں جدوں ہر اوس مکمل $\Sigma$-بے تضاد
  $\Delta$ لئی، جیہدے لئی $\Box^{-1} \Gamma \subseteq \Delta$ ہووے،
  $!A \in \Delta$ ہووے۔
\end{prop}

\begin{proof}
  منّ لو $\Gamma$ مکمل $\Sigma$-بے تضاد اے۔ «صرف جے» والا رخ سوکھا اے۔ منّ لو $\Box !A \in \Gamma$ تے
  $\Box^{-1}\Gamma \subseteq \Delta$۔ کیوں جے $\Box!A \in \Gamma$، ایس لئی $!A
  \in \Box^{-1}\Gamma \subseteq \Delta$، سو $!A \in \Delta$۔

  «جے» والے رخ لئی اسیں اُلٹی نفی ثابت کردے آں۔ منّ لو
  $\Box!A \notin \Gamma$۔ کیوں جے $\Gamma$ مکمل
  $\Sigma$-بے تضاد اے، ایس لئی ایہ استنباط تحت بند اے، سو $\Gamma
  \Proves/[\Sigma] \Box !A$۔ \olref{lem:box2} مطابق
  $\Box^{-1}\Gamma \Proves/[\Sigma] !A$۔
  \olref[prf][con]{prop:consistencyfacts}\olref[prf][con]{prop:consistencyfacts-b} مطابق
  $\Box^{-1}\Gamma \cup \{ \lnot!A \}$، $\Sigma$-بے تضاد اے۔ لنڈن باؤم دے معاون قضیے مطابق اک مکمل $\Sigma$-بے تضاد
  سیٹ~$\Delta$ اے جیہدے لئی $\Box^{-1}\Gamma \cup \{ \lnot!A \} \subseteq
  \Delta$۔ بے تضادی کارن $!A \notin \Delta$۔
\end{proof}
}{%
% لازمیت دا نشان ابتدائی نہیں، صرف امکان دا نشان ابتدائی اے۔

\begin{prop}\ollabel{prop:diamond}
  جے $\Gamma$ مکمل $\Sigma$-بے تضاد اے، تاں $\Diamond !A \in
  \Gamma$ اُدوں تے صرف اُدوں جدوں کسے ایسے مکمل $\Sigma$-بے تضاد
  $\Delta$ لئی جیہدے لئی $\Diamond\Delta \subseteq
  \Gamma$ ہووے، $!A \in \Delta$ ہووے۔
\end{prop}

\begin{proof}
  منّ لو $\Gamma$ مکمل $\Sigma$-بے تضاد اے۔ سجے توں کھبے والا حصہ سوکھا اے۔ منّ لو $\Delta$ مکمل $\Sigma$-بے تضاد اے جیہدے لئی
  $\Diamond\Delta \subseteq \Gamma$ تے $!A \in \Delta$۔ فیر
  $\Diamond!A \in \Diamond\Delta \subseteq \Gamma$، یعنی $\Diamond!A
  \in \Gamma$۔

  کھبے توں سجے والے رخ لئی منّ لو $\Diamond !A \in \Gamma$۔ سانوں وکھاؤنا اے کہ اک مکمل $\Sigma$-بے تضاد $\Delta$
  اے جیہدے لئی $!A \in \Delta$ تے $\Diamond\Delta \subseteq \Gamma$۔ کیوں جے
  $\Diamond !A \in \Gamma$ تے $\Gamma$، $\Sigma$-بے تضاد اے، ایس لئی
  $\Box\lnot !A \notin \Gamma$۔ کیوں جے $\Gamma$ مکمل
  $\Sigma$-بے تضاد اے، ایس لئی ایہ استنباط تحت بند اے، سو $\Gamma
  \Proves/[\Sigma] \Box\lnot !A$۔ \olref{lem:box2} مطابق
  $\Box^{-1}\Gamma \Proves/[\Sigma] \lnot !A$۔
  \olref[prf][con]{prop:consistencyfacts}\olref[prf][con]{prop:consistencyfacts-b} مطابق
  $\Box^{-1}\Gamma \cup \{ !A \}$، $\Sigma$-بے تضاد اے۔ لنڈن باؤم دے معاون قضیے مطابق اک مکمل $\Sigma$-بے تضاد
  سیٹ~$\Delta$ اے جیہدے لئی $\Box^{-1}\Gamma \cup \{ !A \} \subseteq
  \Delta$۔ صاف اے کہ $!A \in \Delta$۔ ایہ ویکھن لئی کہ $\Diamond\Delta
  \subseteq \Gamma$، منّ لو ایہ گل نہیں۔ فیر کسے $!B \in \Delta$ لئی
  $\Diamond !B \notin \Gamma$۔ فیر کیوں جے $\Gamma$ مکمل
  $\Sigma$-بے تضاد اے، ایس لئی $\Box\lnot !B \in \Gamma$۔ پر فیر $\lnot
  !B \in \Box^{-1}\Gamma \subseteq \Delta$ وی، جیہدے نال $\Delta$
  تضاد والا ہو جاندا۔
\end{proof}
}

\begin{lem}\ollabel{lem:box-iff-diamond}
  منّ لو $\Gamma$ تے $\Delta$ مکمل
  $\Sigma$-بے تضاد نیں۔ فیر $\Box^{-1}\Gamma \subseteq \Delta$ اُدوں تے صرف اُدوں جدوں $\Diamond\Delta \subseteq \Gamma$۔
\end{lem}

\begin{proof}
  «صرف جے» والا رخ۔ منّ لو $\Box^{-1}\Gamma \subseteq \Delta$ تے
  $\Diamond!A \in \Diamond\Delta$ (یعنی $!A \in \Delta$)۔
  $\Diamond!A \in \Gamma$ وکھاؤن لئی
  $\Box\lnot!A \notin \Gamma$ وکھاؤنا کافی اے، کیوں جے فیر ودھ توں ودھ بے تضادی کارن $\lnot\Box\lnot
  !A \in \Gamma$۔ ہُن جے $\Box\lnot!A \in \Gamma$، تاں مفروضے مطابق
  $\lnot!A \in \Delta$، جیہڑی $\Delta$ دی بے تضادی دے اُلٹ اے (کیوں جے $!A
  \in \Delta$)۔ سو $\Box\lnot!A \notin \Gamma$، جیویں وکھاؤنا سی۔

  «جے» والا رخ۔ منّ لو $\Diamond\Delta
  \subseteq \Gamma$۔ اسیں اُلٹی نفی ورتدے آں۔ منّ لو $!A
  \notin \Delta$، تاں جو $\Box!A \notin \Gamma$ وکھائیے۔ جے
  $!A \notin \Delta$ تاں ودھ توں ودھ بے تضادی کارن $\lnot!A \in \Delta$
  تے ایس لئی مفروضے مطابق $\Diamond\lnot!A \in \Gamma$۔ پر کسے نارمل امکانی منطق وچ $\Diamond\lnot!A$،
  $\lnot\Box !A$ دے ہم ارز اے، تے جے دوجا $\Gamma$ وچ اے، تاں بے تضادی کارن
  $\Box!A \notin\Gamma$، جیویں وکھاؤنا سی۔
\end{proof}

\iftag{prvBox}{%
  \iftag{prvDiamond}{%
% لازمیت تے امکان دے نشان دونیں ابتدائی نیں۔
% prop:diamond نوں lem:box-iff-diamond ورت کے ثابت کرو۔

\begin{prop}\ollabel{prop:diamond}
  جے $\Gamma$ مکمل $\Sigma$-بے تضاد اے، تاں $\Diamond !A \in
  \Gamma$ اُدوں تے صرف اُدوں جدوں کسے ایسے مکمل $\Sigma$-بے تضاد
  $\Delta$ لئی جیہدے لئی $\Diamond\Delta \subseteq
  \Gamma$ ہووے، $!A \in \Delta$ ہووے۔
\end{prop}

\begin{proof}
منّ لو $\Gamma$ مکمل $\Sigma$-بے تضاد اے۔ $\Diamond!A \in
\Gamma$ اُدوں تے صرف اُدوں جدوں $\lnot\Box\lnot !A \in \Gamma$، \Dual{} تے استنباط تحت بندش کارن۔
$\lnot\Box\lnot !A \in \Gamma$ اُدوں تے صرف اُدوں جدوں $\Box\lnot !A \notin \Gamma$،
\olref[ccs]{prop:ccs-properties}\olref[ccs]{prop:ccs-lnot} مطابق، کیوں جے
$\Gamma$ مکمل $\Sigma$-بے تضاد اے۔ \olref{prop:box} مطابق
$\Box\lnot !A \notin \Gamma$ اُدوں تے صرف اُدوں جدوں کوئی مکمل
$\Sigma$-بے تضاد $\Delta$ ایسا ہووے کہ $\Box^{-1}\Gamma \subseteq \Delta$ تے
$\lnot!A \notin \Delta$۔ ہُن کوئی وی ایسا~$\Delta$ ویکھو۔
\olref{lem:box-iff-diamond} مطابق $\Box^{-1}\Gamma \subseteq \Delta$ اُدوں تے صرف اُدوں جدوں
$\Diamond\Delta \subseteq \Gamma$۔ نالے $\lnot !A \notin \Delta$ اُدوں تے صرف اُدوں جدوں
$!A \in \Delta$،
\olref[ccs]{prop:ccs-properties}\olref[ccs]{prop:ccs-lnot} مطابق۔ سو
$\Diamond!A \in \Gamma$ اُدوں تے صرف اُدوں جدوں کوئی مکمل $\Sigma$-بے تضاد
$\Delta$ ایسا ہووے کہ $\Diamond\Delta \subseteq \Gamma$ تے $!A \in \Delta$۔
\end{proof}

}{}}{}

\begin{prob}
  وکھاؤ کہ جے $\Gamma$ مکمل $\Sigma$-بے تضاد اے، تاں
  \iftag{prvBox}{$\Diamond !A \in \Gamma$ اُدوں تے صرف اُدوں جدوں کوئی مکمل
    $\Sigma$-بے تضاد $\Delta$ ایسا ہووے کہ $\Box^{-1}\Gamma
    \subseteq \Delta$ تے $!A \in \Delta$۔}{$\Box !A \in
    \Gamma$ اُدوں تے صرف اُدوں جدوں ہر اوس مکمل $\Sigma$-بے تضاد
    $\Delta$ لئی جیہدے لئی $\Box^{-1} \Gamma \subseteq \Delta$ ہووے،
    $!A \in \Delta$ ہووے۔} \emph{ایہ
    \olref[nml][com][mod]{lem:box-iff-diamond} ورتے بغیر کرو}۔
\end{prob}
\endgroup
% /OLP-0445
% OLP-0446 content/normal-modal-logic/completeness/canonical-models.tex
\begingroup
\def\olfilename{canonical-models}
\expandafter\def\csname ol@sectioncs\endcsname{section}

\olfileid{nml}{com}{cmd}

\olsection{کینونیکل نمونے}

امکانی نظام~$\Sigma$ دا \emph{کینونیکل نمونہ} اک خاص نمونہ~$\mModel{M^\Sigma}$ اے جیہدے جہان سارے مکمل $\Sigma$-بے تضاد سیٹ نیں۔ ایہدا رسائی دا تعلق~$R^\Sigma$ تے قدر بندی $V^\Sigma$ ایویں متعین کیتے جاندے نیں کہ جہان~$\Delta$ وچ سچے !!{formula}s عین اوہی !!{formula}s ہون جیہڑے~$\Delta$ وچ نیں۔

\begin{defn}
  منّ لو $\Sigma$ اک نارمل امکانی منطق اے۔ \emph{کینونیکل نمونہ}، نظام
  $\Sigma$ دا، $\mModel{M}^\Sigma = \tuple{W^\Sigma, R^\Sigma,
    V^\Sigma}$ اے، جتھے
  \begin{enumerate}
  \item $W^\Sigma = \Setabs{ \Delta }{\Delta \text{
      مکمل $\Sigma$-بے تضاد اے} }$۔
  \item $R^\Sigma \Delta\Delta'$ اُدوں تے صرف اُدوں ہوندا اے جدوں
    \iftag{prvBox}{$\Box^{-1}\Delta \subseteq
      \Delta'$}{$\Diamond\Delta' \subseteq \Delta$}۔
  \item $V^\Sigma(p) = \Setabs{\Delta}{p \in \Delta}$۔
  \end{enumerate}
\end{defn}
\endgroup
% /OLP-0446
% OLP-0447 content/normal-modal-logic/completeness/truth-lemma.tex
\begingroup
\def\olfilename{truth-lemma}
\expandafter\def\csname ol@sectioncs\endcsname{section}

\olfileid{nml}{com}{tru}

\olsection{سچائی دا معاون قضیہ}

کینونیکل نمونہ~$\mModel{M^\Sigma}$ ایویں متعین اے کہ
$\mSat{M^\Sigma}{!A}[\Delta]$ اُدوں تے صرف اُدوں جدوں $!A \in \Delta$۔ !!{propositional variable}s لئی $V^\Sigma$ دی تعریف ایہ گل سدھی دیندی اے۔ پر سانوں پرکھنا اے کہ ایہ ہم ارزی سارے !!{formula}s لئی چلدی اے۔ اسیں ایہ استقرا نال کردے آں۔ استقرائی قدم وچ قضیاتی عاملاں والے !!{formula}s لئی ہم ارزی ثابت کرنی اے، جتھے سانوں \olref[ccs]{prop:ccs-properties} ورتنا پئے گا، تے امکانی عاملاں لئی وی، جتھے اسیں \olref[mod]{sec} دے نتیجے ورتاں گے۔

\begin{prop}[سچائی دا معاون قضیہ]
  \ollabel{prop:truthlemma}
  ہر !!{formula}~$!A$ لئی $\mSat{M^\Sigma}{!A}[\Delta]$ اُدوں تے صرف اُدوں جدوں
  $!A \in \Delta$۔
\end{prop}

\begin{proof}
  استقرا نال،~$!A$ اُتے۔
  \begin{enumerate}
    \tagitem{prvFalse}{\indcase{!A}{\lfalse}{$\mSat/{M^\Sigma}{\lfalse}[\Delta]$،
        \olref[syn][trw]{defn:mmodels} مطابق، تے $\lfalse \notin \Delta$،
        \olref[ccs]{prop:ccs-properties}\olref[ccs]{prop:ccs-lfalse} مطابق۔}}{}

    \tagitem{prvTrue}{\indcase{!A}{\ltrue}{$\mSat{M^\Sigma}{\ltrue}[\Delta]$،
        \olref[syn][trw]{defn:mmodels} مطابق، تے $\ltrue \in \Delta$،
        \olref[ccs]{prop:ccs-properties}\olref[ccs]{prop:ccs-ltrue} مطابق۔}}{}

    \item \indcase{!A}{p}{$\mSat{M^\Sigma}{p}[\Delta]$ اُدوں تے صرف اُدوں جدوں $\Delta \in
      V^\Sigma(p)$، \olref[syn][trw]{defn:mmodels} مطابق۔ نالے $\Delta
      \in V^\Sigma(p)$ اُدوں تے صرف اُدوں جدوں $p \in \Delta$،
      تعریف مطابق~$V^\Sigma$ دی۔}

    \tagitem{prvNot}{\indcase{!A}{\lnot !B}{%
        \iftag{probNot}{مشق۔}{$\mSat{M^\Sigma}{\lnot !B}[\Delta]$
          اُدوں تے صرف اُدوں جدوں $\mSat/{M^\Sigma}{!B}[\Delta]$
          (\olref[syn][trw]{defn:mmodels} مطابق)، ایہ اُدوں تے صرف اُدوں جدوں $!B \notin \Delta$ (استقرائی مفروضے مطابق)، ایہ اُدوں تے صرف اُدوں جدوں $\lnot !B \in \Delta$ (
          \olref[ccs]{prop:ccs-properties}\olref[ccs]{prop:ccs-lnot} مطابق)۔}}}{}

    \tagitem{prvAnd}{\indcase{!A}{!B \land !C}{%
        \iftag{probAnd}{مشق۔}{$\mSat{M^\Sigma}{!B \land
            !C}[\Delta]$ اُدوں تے صرف اُدوں جدوں $\mSat{M^\Sigma}{!B}[\Delta]$ تے
          $\mSat{M^\Sigma}{!C}[\Delta]$ (\olref[syn][trw]{defn:mmodels} مطابق)، ایہ اُدوں تے صرف اُدوں جدوں $!B \in \Delta$ تے $!C
          \in \Delta$ (استقرائی مفروضے مطابق)، ایہ اُدوں تے صرف اُدوں جدوں $!B \land !C \in
          \Delta$ (\olref[ccs]{prop:ccs-properties}\olref[ccs]{prop:ccs-land} مطابق)۔}}}{}

    \tagitem{prvOr}{\indcase{!A}{!B \lor !C}{%
        \iftag{probOr}{مشق۔}{$\mSat{M^\Sigma}{!B \lor
            !C}[\Delta]$ اُدوں تے صرف اُدوں جدوں $\mSat{M^\Sigma}{!B}[\Delta]$ یا
          $\mSat{M^\Sigma}{!C}[\Delta]$ (\olref[syn][trw]{defn:mmodels} مطابق)، ایہ اُدوں تے صرف اُدوں جدوں $!B \in \Delta$ یا $!C
          \in \Delta$ (استقرائی مفروضے مطابق)، ایہ اُدوں تے صرف اُدوں جدوں $!B \lor !C \in
          \Delta$ (\olref[ccs]{prop:ccs-properties}\olref[ccs]{prop:ccs-lor} مطابق)۔}}}{}

    \tagitem{prvIf}{\indcase{!A}{!B \lif !C}{%
        \iftag{probIf}{مشق۔}{$\mSat{M^\Sigma}{!B \lif
            !C}[\Delta]$ اُدوں تے صرف اُدوں جدوں $\mSat/{M^\Sigma}{!B}[\Delta]$ یا
          $\mSat{M^\Sigma}{!C}[\Delta]$ (\olref[syn][trw]{defn:mmodels} مطابق)، ایہ اُدوں تے صرف اُدوں جدوں $!B \notin \Delta$ یا
          $!C \in \Delta$ (استقرائی مفروضے مطابق)، ایہ اُدوں تے صرف اُدوں جدوں $!B \lif !C
          \in \Delta$ (\olref[ccs]{prop:ccs-properties}\olref[ccs]{prop:ccs-lif} مطابق)۔}}}{}

    \tagitem{prvIff}{\indcase{!A}{!B \liff !C}{%
        \iftag{probIff}{مشق۔}{$\mSat{M^\Sigma}{!B \liff
            !C}[\Delta]$ اُدوں تے صرف اُدوں جدوں یا تاں $\mSat{M^\Sigma}{!B}[\Delta]$ تے
          $\mSat{M^\Sigma}{!C}[\Delta]$ دونیں، یا
          $\mSat/{M^\Sigma}{!B}[\Delta]$ تے
          $\mSat/{M^\Sigma}{!C}[\Delta]$ دونیں (\olref[syn][trw]{defn:mmodels} مطابق)۔ ایہ اُدوں تے صرف اُدوں جدوں یا تاں $!B \in \Delta$
          تے $!C \in \Delta$ دونیں، یا $!B \notin \Delta$ تے $!C \notin
          \Delta$ دونیں (استقرائی مفروضے مطابق)۔ ایہ اُدوں تے صرف اُدوں جدوں $!B \liff !C \in
          \Delta$ (\olref[ccs]{prop:ccs-properties}\olref[ccs]{prop:ccs-liff} مطابق)۔}}}{}

    \tagitem{prvBox}{\indcase{!A}{\Box !B}{%
        \iftag{probBox}{مشق۔}{پہلاں منّ لو
          $\mSat{M^\Sigma}{\Box !B}[\Delta]$۔
          \olref[syn][trw]{defn:mmodels} مطابق ہر ایسے $\Delta'$ لئی
          جیہدے لئی $R^\Sigma \Delta\Delta'$،
          $\mSat{M^\Sigma}{!B}[\Delta']$ اے۔ استقرائی مفروضے مطابق ہر ایسے
          $\Delta'$ لئی جیہدے لئی $R^\Sigma \Delta\Delta'$، $!B \in
          \Delta'$۔ $R^\Sigma$ دی تعریف مطابق ہر ایسے $\Delta'$ لئی
          جیہدے لئی $\Box^{-1} \Delta \subseteq \Delta'$، $!B \in
          \Delta'$۔ \olref[mod]{prop:box} مطابق $\Box !B \in
          \Delta$۔

          ہُن منّ لو $\Box !B \in \Delta$۔ کوئی ایسا $\Delta' \in W^\Sigma$
          لے لو جیہدے لئی $R^\Sigma \Delta\Delta'$، یعنی $\Box^{-1}\Delta
          \subseteq \Delta'$۔ کیوں جے $\Box !B \in \Delta$، ایس لئی $!B \in
          \Box^{-1} \Delta$۔ سو $!B \in \Delta'$۔ استقرائی مفروضے مطابق
          $\mSat{M^\Sigma}{!B}[\Delta']$۔ کیوں جے
          $\Delta'$ کوئی وی ہو سکدا سی جیہدے لئی $R^\Sigma \Delta\Delta'$، ایس لئی
          ہر ایسے $\Delta' \in W^\Sigma$ لئی جیہدے لئی $R^\Sigma
          \Delta\Delta'$، $\mSat{M^\Sigma}{!B}[\Delta']$۔ \olref[syn][trw]{defn:mmodels} مطابق $\mSat{M^\Sigma}{\Box
            !B}[\Delta]$۔}}}{}

    \tagitem{prvDiamond}{\indcase{!A}{\Diamond !B}{%
        \iftag{probDiamond}{مشق۔}{پہلاں منّ لو
          $\mSat{M^\Sigma}{\Diamond !B}[\Delta]$۔
          \olref[syn][trw]{defn:mmodels} مطابق کسے ایسے $\Delta'$ لئی جیہدے لئی
          $R^\Sigma \Delta\Delta'$، $\mSat{M^\Sigma}{!B}[\Delta']$۔ استقرائی مفروضے مطابق کسے ایسے $\Delta'$ لئی جیہدے لئی
          $R^\Sigma \Delta\Delta'$، $!B \in \Delta'$۔
          $R^\Sigma$ دی تعریف مطابق،\iftag{prvBox}{ کسے ایسے $\Delta'$ لئی جیہدے لئی
            $\Box^{-1} \Delta \subseteq \Delta'$، $!B \in \Delta'$۔
            \olref[mod]{prop:diamond} مطابق،}{} کسے ایسے $\Delta'$ لئی جیہدے لئی
          $\Diamond \Delta' \subseteq \Delta$، $!B \in
          \Delta'$۔ کیوں جے $!B \in \Delta'$، ایس لئی $\Diamond !B \in \Diamond
          \Delta'$، سو $\Diamond !B \in \Delta$۔

          ہُن منّ لو $\Diamond !B \in \Delta$۔
          \olref[mod]{prop:diamond} مطابق اک مکمل
          $\Sigma$-بے تضاد $\Delta' \in W^\Sigma$ اے جیہدے لئی
          $\Diamond \Delta' \subseteq \Delta$ تے $!B \in
          \Delta'$۔\iftag{prvBox}{ \olref[mod]{lem:box-iff-diamond} مطابق،
            ایسا $\Delta' \in W^\Sigma$ اے جیہدے لئی $\Box^{-1}
            \Delta \subseteq \Delta'$ تے $!B \in \Delta'$۔}{} 
          $R^\Sigma$ دی تعریف مطابق $R^\Sigma \Delta\Delta'$، سو
          ایسا $\Delta' \in W^\Sigma$ اے جیہدے لئی $R^\Sigma
          \Delta\Delta'$ تے $!B \in \Delta'$۔ \olref[syn][trw]{defn:mmodels} مطابق $\mSat{M^\Sigma}{\Diamond
            !B}[\Delta]$۔}}}{}
  \end{enumerate}
\end{proof}

\begin{probtag}{probFalse,probTrue,probNot,proband,probOr,probIf,probIff,probBox,probDiamond}
  \olref[nml][com][tru]{prop:truthlemma} دا ثبوت پورا کرو۔
\end{probtag}
\endgroup
% /OLP-0447
% OLP-0448 content/normal-modal-logic/completeness/completeness-K.tex
\begingroup
\def\olfilename{completeness-K}
\expandafter\def\csname ol@sectioncs\endcsname{section}

\olfileid{nml}{com}{cmk}

\olsection{\Log{K} دا تعیّن تے تمامیت}

ہُن اسیں کینونیکل نمونہ ورت کے تمامیت ثابت کر سکدے آں۔ تمامیت ایس حقیقت توں ملدی اے کہ~$\Sigma$ دے کینونیکل نمونے وچ سچے !!{formula}s عین اوہی نیں جیہڑے $\Sigma$ توں !!{derivable} نیں۔ ایس خاصیت والے نمونے~$\Sigma$ نوں \emph{متعین کردے نیں}۔

\begin{defn}
  نمونہ $\mModel{M}$ نارمل امکانی منطق~$\Sigma$ نوں \emph{متعین کردا اے} ٹھیک اُدوں جدوں $\mSat{M}{!A}$ اُدوں تے صرف اُدوں ہووے جدوں $\Sigma
  \Proves !A$، سارے !!{formula}s~$!A$ لئی۔
\end{defn}

\begin{thm}[تعیّن]\ollabel{thm:determination}
   $\mSat{M^\Sigma}{!A}$ اُدوں تے صرف اُدوں جدوں $\Sigma \Proves !A$۔
\end{thm}

\begin{proof}
  جے $\mSat{M^\Sigma}{!A}$، تاں ہر مکمل
  $\Sigma$-بے تضاد $\Delta$ لئی
  $\mSat{M^\Sigma}{!A}[\Delta]$۔ ایس لئی سچائی دے معاون قضیے مطابق $!A \in
  \Delta$ ہر مکمل $\Sigma$-بے تضاد $\Delta$ لئی۔ ایس توں
  \olref[lin]{cor:provability-characterization} مطابق ($\Gamma =
  \emptyset$ رکھ کے)، $\Sigma \Proves !A$۔

  اُلٹے پاسے جے $\Sigma \Proves !A$، تاں
  \olref[ccs]{prop:ccs-properties}\olref[ccs]{prop:ccs-closed} مطابق ہر مکمل $\Sigma$-بے تضاد $\Delta$ وچ~$!A$ شامل اے۔ ایس لئی سچائی دے معاون قضیے مطابق
  $\mSat{M^\Sigma}{!A}[\Delta]$ ہر~$\Delta \in
  W^\Sigma$ لئی، یعنی $\mSat{M^\Sigma}{!A}$۔
\end{proof}

کیوں جے \Log{K} دا کینونیکل نمونہ~\Log{K} نوں متعین کردا اے، ایس لئی~\Log{K} دی تمامیت سدھی اک نتیجے وجوں ملدی اے۔

\begin{cor}\ollabel{cor:Kcomplete}
  بنیادی امکانی منطق \Log{K} سارے نمونیاں دی جماعت دے لحاظ نال مکمل اے، یعنی جے $\Entails !A$، تاں $\Log{K}
  \Proves !A$۔
\end{cor}

\begin{proof}
  اُلٹی نفی راہیں، جے $\Log{K} \Proves/ !A$، تاں تعیّن دے قضیے مطابق
  $\mSat/{M^{\Log{K}}}{!A}$، سو $!A$ معتبر نہیں۔
\end{proof}

عام صورت وچ کسے نظام $\Sigma$ دی نمونیاں دی جماعت دے لحاظ نال تمامیت لئی، جیویں \Log{KTB4} دی سارے انعکاسی، متقابل تے متعدی نمونیاں دی جماعت دے لحاظ نال، صرف تعیّن کافی نہیں۔ سانوں ایہ وی وکھاؤنا اے کہ نظام
$\Sigma$ دا کینونیکل نمونہ اوس جماعت دا رکن اے۔ ایہ گل کینونیکل نمونے دی بناوٹ توں سدھی نہیں نکلدی، تے ایہ ہمیشہ سچی وی نہیں!
\endgroup
% /OLP-0448
% OLP-0449 content/normal-modal-logic/completeness/frame-completeness.tex
\begingroup
\def\olfilename{frame-completeness}
\expandafter\def\csname ol@sectioncs\endcsname{section}

\olfileid{nml}{com}{fra}

\olsection{فریماں دے لحاظ نال تمامیت}

\Log{K} دی تمامیت دا قضیہ ہور امکانی نظاماں تک ودھا سکدے آں جدوں اسیں وکھا دیئے کہ دِتّی منطق دے کینونیکل نمونے وچ متعلقہ فریم خاصیت اے۔

\begin{thm}\ollabel{thm:completeframeprops}
  جے نارمل امکانی منطق $\Sigma$ وچ \olref{tab:correspondencetable} دے کھبے پاسے والے !!{formula}s وچوں کوئی اک شامل اے،
  تاں~$\Sigma$ دے کینونیکل نمونے وچ سجے پاسے والی متعلقہ خاصیت اے۔
\end{thm}

\begin{table}[htp]
  \centering
    \begin{tabular}{| l || l |}
      \hline
      \emph{جے $\Sigma$ وچ \dots شامل اے} & \emph{\dots، $\Sigma$ دا کینونیکل نمونہ ایہ اے۔} \\
      \hline \hline
      \Ax{D}:\;\; $\Box!A \lif \Diamond !A$ & \quad ہر جہان توں رسائی والا؛ \\
      \hline
      \Ax{T}:\;\; $\Box!A \lif !A$ &\quad  انعکاسی؛\\
      \hline
      \Ax{B}: \;\;$!A \lif \Box\Diamond!A$ &\quad  متقابل؛ \\
      \hline
      \Ax{4}: \;\; $\Box!A \lif \Box\Box!A$ & \quad متعدی؛ \\
      \hline 
      \Ax{5}: \;\; $\Diamond !A \lif \Box\Diamond!A$& \quad اقلیدسی۔\\
      \hline
    \end{tabular}
    \caption{مطابقت دیاں بنیادی گلاں۔}
    \ollabel{tab:correspondencetable}
\end{table}

\begin{proof}
  اسیں اینہاں وچوں ہر اک نوں واری واری لیندے آں۔

  منّ لو $\Sigma$ وچ \Ax{D} شامل اے تے $\Delta \in W^\Sigma$ لے لو۔ سانوں وکھاؤنا اے کہ کوئی $\Delta'$ اے جیہدے لئی $R^\Sigma
  \Delta\Delta'$۔ ایہ وکھاؤنا کافی اے کہ $\Box^{-1}\Delta$،
  $\Sigma$-بے تضاد اے، کیوں جے فیر لنڈن باؤم دے معاون قضیے مطابق اک
  مکمل $\Sigma$-بے تضاد سیٹ $\Delta' \supseteq
  \Box^{-1}\Delta$ اے، تے $R^\Sigma$ دی تعریف
  \iftag{prvBox}{}{تے \olref[mod]{lem:box-iff-diamond} }مطابق ساڈے کول
  $R^\Sigma \Delta\Delta'$ اے۔ سو تضاد کڈھن لئی منّ لو
  $\Box^{-1}\Delta$، $\Sigma$-بے تضاد \emph{نہیں}، یعنی
  $\Box^{-1}\Delta \Proves[\Sigma] \lfalse$۔ \olref[mod]{lem:box2} مطابق
  $\Delta \Proves[\Sigma] \Box\lfalse$، تے کیوں جے $\Sigma$ وچ
  \Ax{D} شامل اے، ایس لئی $\Delta \Proves[\Sigma] \Diamond\lfalse$ وی۔ پر $\Sigma$
  نارمل اے، ایس لئی $\Sigma \Proves \lnot\Diamond \lfalse$
  (\olref[prf][nor]{prop:notDiamondBot})، جیہدے توں $\Delta
  \Proves[\Sigma] \lnot\Diamond \lfalse$ وی، جو~$\Delta$ دی بے تضادی دے اُلٹ اے۔

  ہُن منّ لو $\Sigma$ وچ \Ax{T} شامل اے، تے $\Delta \in
  W^\Sigma$ لے لو۔ اسیں $R^\Sigma \Delta\Delta$ وکھاؤنا چاہندے آں،
  یعنی\iftag{prvBox}{}{ \olref[mod]{lem:box-iff-diamond} مطابق،}
  $\Box^{-1}\Delta \subseteq \Delta$۔ پر جے $\Box!A \in \Delta$ تاں
  \Ax{T} مطابق $!A \in \Delta$ وی، جیویں وکھاؤنا سی۔

  ہُن منّ لو $\Sigma$ وچ \Ax{B} شامل اے تے $R^\Sigma
  \Delta\Delta'$ اے $\Delta$، $\Delta' \in W^\Sigma$ لئی۔ سانوں وکھاؤنا اے کہ $R^\Sigma \Delta'\Delta$، یعنی\iftag{prvBox}{
    $\Box^{-1}\Delta' \subseteq \Delta$۔
    \olref[mod]{lem:box-iff-diamond} مطابق ایہ ایس دے ہم ارز اے}{} کہ
  $\Diamond\Delta \subseteq \Delta'$۔ سو منّ لو $!A \in \Delta$۔
  \Ax{B} مطابق $\Box\Diamond!A \in \Delta$ وی۔ ایس مفروضے توں کہ
  $R^\Sigma \Delta\Delta'$\iftag{prvBox}{}{، تے
    \olref[mod]{lem:box-iff-diamond} توں}، ملدا اے کہ $\Box^{-1}\Delta
  \subseteq \Delta'$، سو $\Diamond!A \in \Delta'$، جیویں وکھاؤنا سی۔

  ہُن منّ لو $\Sigma$ وچ \Ax{4} شامل اے تے $R^\Sigma
  \Delta_1\Delta_2$ تے $R^\Sigma \Delta_2\Delta_3$ نیں۔ سانوں
  $R^\Sigma \Delta_1\Delta_3$ وکھاؤنا اے۔ مفروضے توں\iftag{prvBox}{}{ تے
    \olref[mod]{lem:box-iff-diamond} توں} سانوں $\Box^{-1}\Delta_1
  \subseteq \Delta_2$ تے $\Box^{-1}\Delta_2 \subseteq \Delta_3$ دونیں ملدے نیں۔
  $R^\Sigma \Delta_1\Delta_3$ وکھاؤن لئی
  $\Box^{-1}\Delta_1 \subseteq \Delta_3$ وکھاؤنا کافی اے۔ سو منّ لو $!B \in
  \Box^{-1}\Delta_1$، یعنی $\Box!B \in \Delta_1$۔ \Ax{4} مطابق
  $\Box\Box!B \in \Delta_1$ وی، تے مفروضے توں پہلاں
  $\Box!B \in \Delta_2$ تے فیر $!B \in \Delta_3$ ملدے نیں، جیویں وکھاؤنا سی۔

  ہُن منّ لو $\Sigma$ وچ \Ax{5} شامل اے تے $R^\Sigma
  \Delta_1\Delta_2$ تے $R^\Sigma \Delta_1\Delta_3$ نیں۔ سانوں
  $R^\Sigma \Delta_2\Delta_3$ وکھاؤنا اے۔ پہلے مفروضے توں
  $\Box^{-1}\Delta_1 \subseteq \Delta_2$ ملدا اے\iftag{prvBox}{}{،
    \olref[mod]{lem:box-iff-diamond} مطابق}، تے دوجا مفروضہ
  $\Diamond\Delta_3 \subseteq \Delta_1$\iftag{prvBox}{ دے ہم ارز اے،
    \olref[mod]{lem:box-iff-diamond} مطابق}{ دے ہم ارز اے}۔
  $R^\Sigma
  \Delta_2\Delta_3$ وکھاؤن لئی\iftag{prvBox}{،
    \olref[mod]{lem:box-iff-diamond} مطابق، ایہ کافی اے کہ}{ سانوں ایہ وکھاؤنا اے کہ}
  $\Diamond\Delta_3 \subseteq \Delta_2$۔ سو منّ لو $\Diamond!A \in
  \Diamond\Delta_3$، یعنی $!A \in \Delta_3$۔ دوجے مفروضے مطابق
  $\Diamond!A \in \Delta_1$، تے \Ax{5} مطابق $\Box\Diamond!A \in
  \Delta_1$ وی۔ ہُن پہلے مفروضے توں $\Diamond!A \in
  \Delta_2$ ملدا اے، جیویں وکھاؤنا سی۔
  \end{proof}

نتیجے وجوں سانوں کئی نظاماں دی تمامیت دے نتیجے ملدے نیں۔ مثلاً اسیں جان دے آں کہ $\Log{S5} = \Log{KT5} =
\Log{KTB4}$ سارے انعکاسی اقلیدسی نمونیاں دی جماعت دے لحاظ نال مکمل اے، تے ایہ جماعت اوہی اے جیہڑی سارے انعکاسی، متقابل تے متعدی نمونیاں دی جماعت اے۔

\begin{thm}\ollabel{thm:generaldet}
  منّ لو $\mClass{C}_\Ax{D}$، $\mClass{C}_\Ax{T}$،
  $\mClass{C}_\Ax{B}$، $\mClass{C}_\Ax{4}$، تے
  $\mClass{C}_\Ax{5}$ بالترتیب سارے ہر جہان توں رسائی والے، انعکاسی،
  متقابل، متعدی، تے اقلیدسی نمونیاں دیاں جماعتاں نیں۔ فیر کسے وی سانچیاں $!A_1$، \dots، $!A_n$ لئی، جیہڑے \Ax{D}،
  \Ax{T}، \Ax{B}، \Ax{4}، تے \Ax{5} وچوں نیں، نظام
  $\Log{K}!A_1 \dots !A_n$ نوں نمونیاں دی جماعت
  $\mClass{C} = \mClass{C}_{!A_1} \cap \dots
  \cap \mClass{C}_{!A_n}$ متعین کردی اے۔
\end{thm}

\begin{prop}
  منّ لو $\Sigma$ نارمل امکانی منطق اے۔ فیر
  \begin{enumerate}
  \item\ollabel{prop:anotherfive-a}%
    جے $\Sigma$ وچ سانچہ $\Diamond!A \lif \Box
    !A$ شامل اے، تاں $\Sigma$ دا کینونیکل نمونہ جزوی فنکشنی اے۔
  \item جے $\Sigma$ وچ سانچہ $\Diamond!A \liff \Box
    !A$ شامل اے، تاں $\Sigma$ دا کینونیکل نمونہ فنکشنی اے۔
  \item جے $\Sigma$ وچ سانچہ $\Box\Box!A \lif \Box
    !A$ شامل اے، تاں $\Sigma$ دا کینونیکل نمونہ کمزور طور تے گنجان اے۔
  \end{enumerate}
(ایہناں فریم خاصیتاں دیاں تعریفاں لئی \olref[frd][acc]{tab:anotherfive} ویکھو۔)
\end{prop}

\begin{proof}
  \begin{enumerate}
  \item منّ لو $\Sigma$ وچ سانچہ $\Diamond !A \lif
    \Box !A$ شامل اے۔ $R^\Sigma$ نوں جزوی فنکشنی وکھاؤن لئی ثابت کرنا اے کہ کسے وی $\Delta_1$، $\Delta_2$، $\Delta_3 \in
    W^\Sigma$ لئی جے $R^\Sigma \Delta_1\Delta_2$ تے $R^\Sigma
    \Delta_1\Delta_3$، تاں $\Delta_2=\Delta_3$۔ کیوں جے $R^\Sigma
    \Delta_1\Delta_2$، ایس لئی $\Box^{-1}\Delta_1 \subseteq \Delta_2$،
    تے کیوں جے $R^\Sigma \Delta_1\Delta_3$، ایس لئی $\Box^{-1}\Delta_1
    \subseteq \Delta_3$\iftag{prvBox}{}{ وی، دونیں
      \olref[mod]{lem:box-iff-diamond} مطابق}۔ برابری
    $\Delta_2=\Delta_3$ اوہدوں نکلے گی جے اسیں دونیں شمولیتاں
    $\Delta_2 \subseteq \Delta_3$ تے $\Delta_3 \subseteq
    \Delta_2$ ثابت کر سکیے۔ پہلی لئی منّ لو $!A \in \Delta_2$۔ فیر
    $\Diamond!A \in \Delta_1$، تے سانچے تے استنباط تحت بندش مطابق، جیہڑی $\Delta_1$ دی اے، $\Box!A \in \Delta_1$ وی۔ ایس توں مفروضے مطابق کہ
    $R^\Sigma \Delta_1\Delta_3$، $!A \in \Delta_3$، جیویں وکھاؤنا سی۔ دوجی شمولیت ایس طرح اے۔

  \item ایہ سدھا حصے~\olref{prop:anotherfive-a} تے
    \olref{thm:completeframeprops} وچ ہر جہان توں رسائی والے ثبوت توں نکلدا اے۔

  \item منّ لو $\Sigma$ وچ سانچہ $\Box\Box!A \lif \Box!A$ شامل اے۔ $R^\Sigma$ نوں کمزور طور تے گنجان وکھاؤن لئی منّ لو $R^\Sigma
    \Delta_1\Delta_2$۔ سانوں وکھاؤنا اے کہ اک مکمل $\Sigma$-بے تضاد سیٹ $\Delta_3$ اے جیہدے لئی $R^\Sigma
    \Delta_1\Delta_3$ تے $R^\Sigma \Delta_3\Delta_2$۔ ایہ رکھو۔
    \[
    \Gamma = \Box^{-1}\Delta_1 \cup \Diamond\Delta_2.
    \]
    ایہ وکھاؤنا کافی اے کہ $\Gamma$، $\Sigma$-بے تضاد اے۔ فیر لنڈن باؤم دے معاون قضیے مطابق ایہنوں اک مکمل
    $\Sigma$-بے تضاد سیٹ~$\Delta_3$ تک ودھا سکدے آں جیہدے لئی $\Box^{-1}\Delta_1
    \subseteq \Delta_3$ تے $\Diamond\Delta_2 \subseteq \Delta_3$،
    یعنی $R^\Sigma \Delta_1\Delta_3$\iftag{prvBox}{}{ (
      \olref[mod]{lem:box-iff-diamond} مطابق)} تے $R^\Sigma
    \Delta_3\Delta_2$\iftag{prvBox}{ (
      \olref[mod]{lem:box-iff-diamond} مطابق)}{}۔

    تضاد کڈھن لئی منّ لو $\Gamma$ بے تضاد نہیں۔ فیر ایسے !!{formula}s $\Box!A_1$، \dots، $\Box!A_n \in \Delta_1$
    تے $!B_1$، \dots،~$!B_m \in \Delta_2$ نیں کہ \[!A_1, \dots,
    !A_n, \Diamond!B_1, \dots, \Diamond!B_m \Proves[\Sigma]
    \lfalse.\] کیوں جے $\Diamond (!B_1 \land \dots \land !B_m) \to
    (\Diamond!B_1 \land \dots \land \Diamond!B_m)$ ہر نارمل امکانی منطق وچ !!{derivable} اے، ایس لئی اسیں ایویں دلیل دیندے آں، جیہڑی~$\Delta_2$ دی بے تضادی دے اُلٹ اے۔
    \begin{align*}
      \formula A_1, \dots, !A_n, & \Diamond!B_1, \dots,
      \Diamond!B_m \Proves[\Sigma] \lfalse \\
      \formula A_1, \dots,!A_n & \Proves[\Sigma] (\Diamond!B_1 \land
      \dots \land \Diamond!B_m) \lif \lfalse\\
      & \qquad\text{استنباط دے قضیے مطابق}\\
      & \qquad\text{\olref[prf][prp]{prop:derivabilityfacts}\olref[prf][prp]{prop:derivabilityfacts-deduction}، تے \Taut} \\
      \formula A_1, \dots,!A_n & \Proves[\Sigma]
      \Diamond(!B_1 \land
      \dots \land !B_m) \lif \lfalse\\
      & \qquad \text{کیوں جے $\Sigma$ نارمل اے} \\
      \formula A_1, \dots,!A_n & \Proves[\Sigma]
      \lnot\Diamond (!B_1 \land \dots \land !B_m)\\
      & \qquad\text{\PL مطابق} \\
      \formula A_1, \dots,!A_n & \Proves[\Sigma]
      \Box\lnot (!B_1 \land \dots \land !B_m)\\
      & \qquad\text{$\Box\lnot$، $\lnot\Diamond$ دی تھاں} \\
      \Box!A_1, \dots,\Box!A_n
      & \Proves[\Sigma] \Box\Box \lnot (!B_1 \land \dots \land !B_m)\\
      & \qquad\text{\olref[mod]{lem:box1} مطابق} \\
      \Box!A_1, \dots,\Box!A_n
      & \Proves[\Sigma]
      \Box\lnot (!B_1 \land \dots \land !B_m)\\
      &\qquad\text{سانچہ $\Box\Box!A \lif \Box!A$ مطابق} \\
      \Delta_1 & \Proves[\Sigma]
      \Box\lnot (!B_1 \land \dots \land !B_m)\\
      &\qquad\text{یک رُخیت مطابق، \olref[prf][prp]{prop:derivabilityfacts}\olref[prf][prp]{prop:derivabilityfacts-monotonicity}} \\
      & \Box\lnot (!B_1 \land \dots \land !B_m) \in \Delta_1\\
      &\qquad\text{استنباط تحت بندش کارن}; \\
      & \lnot (!B_1 \land \dots \land !B_m) \in \Delta_2\\
      &\qquad \text{کیوں جے }
      R^\Sigma \Delta_1\Delta_2. 
    \end{align*}
  \end{enumerate}
\end{proof}

ایہناں مثالاں توں خیال آ سکدا اے کہ امکانی منطق دا ہر نظام~$\Sigma$ \emph{مکمل} اے، ایس معنی وچ کہ اوہ ہر اوہ فارمولا ثابت کردا اے جیہڑا ہر اوس فریم وچ معتبر اے جیہدے وچ $\Sigma$ دا ہر قضیہ معتبر اے۔ پر کئی نظام ایس معنی وچ مکمل نہیں۔
\endgroup
% /OLP-0449

\OLEndChapterHook
\endgroup
% /OLP-0441

% OLP-0450 content/normal-modal-logic/filtrations/filtrations.tex
\begingroup
\def\olfilename{filtrations}
\expandafter\def\csname ol@sectioncs\endcsname{section}

\olchapter{nml}{fil}{فلٹریشن تے فیصلہ کرن دی اہلیت}

% OLP-0451 content/normal-modal-logic/filtrations/introduction.tex
\begingroup
\def\olfilename{introduction}
\expandafter\def\csname ol@sectioncs\endcsname{section}

\olfileid{nml}{fil}{int}

\olsection{جان پہچان}

کسے منطق بارے اک اہم سوال ہمیشہ ایہ ہوندا اے کہ کی اوہ فیصلہ کرن جوگ اے۔ یعنی کی کوئی مؤثر طریقہ اے جیہڑا ایس سوال دا جواب دیوے، «کی ایہ !!{formula} معتبر اے؟» قضیاتی منطق فیصلہ کرن جوگ اے۔ اسیں صداقتی جدول بنا کے مؤثر طریقے نال پرکھ سکدے آں کہ کوئی !!a{formula} ہمیشہ صادق اے یا نہیں، تے دِتّے !!{formula} لئی صداقتی جدول محدود اے۔ پر ایہ صاف نہیں کہ اسیں کسے امکانی فارمولے نوں سارے نمونیاں وچ سچا ہون لئی کیویں پرکھیے، کیوں جے نمونے لامتناہی گنتی وچ نیں۔ کسے دِتّے فارمولے نال متعلق سارے محدود نمونیاں دی فہرست بنا سکدے آں، کیوں جے صرف جہاناں دے ذیلی سیٹ اوہناں !!{propositional variable}s نوں دین دی گل متعلق اے جیہڑے ایس !!{formula} وچ سچ مچ آؤندے نیں۔ جے رسائی دا تعلق مقرر اے، تاں وکھ وکھ ممکن قدر بندیاں $V(p)$ عین~$W$ دے سارے ذیلی سیٹ نیں، تے جے $\card{W} = n$ تاں ایہناں دی گنتی $2^n$ اے۔ جے ساڈے !!{formula}~$!A$ وچ $m$ !!{propositional variable}s نیں، تاں $2^{nm}$ وکھ وکھ نمونے نیں جیہناں وچ~$n$ جہان نیں۔ ہر نمونے لئی اسیں پرکھ سکدے آں کہ~$!A$ سارے جہاناں وچ سچا اے، ہر جہان وچ~$!A$ دی صداقتی قدر کڈھ کے۔ سانوں رسائی دے سارے ممکن تعلق وی جانچنے پین گے، پر~$n$ جہاناں اُتے تعلق وی محدود گنتی وچ نیں (خاص طور تے $W \times W$ دے ذیلی سیٹاں دی گنتی، یعنی $2^{n^2}$۔

جے ساڈی دلچسپی~\Log{K} دی بجائے نمونیاں دی کسے جماعت نال متعین منطق وچ اے (مثلاً انعکاسی متعدی نمونے)، تاں سانوں ایہ وی پرکھ سکنا پئے گا کہ رسائی دا تعلق مطلوب قسم دا اے۔ ایہ اسیں اوہدوں کر سکدے آں جدوں متعلقہ فریم امکانی !!{formula}s نال متعین ہو سکدے ہون (مثلاً پرکھ کے کہ \Ax{T} تے \Ax{4} فریم وچ معتبر نیں)۔ سو طریقہ ایہ ہووے گا کہ سارے محدود فریم واری واری ویکھیے، ہر اک نوں پرکھیے کہ اوہ متعلقہ جماعت دا رکن اے، فیر اوس فریم اُتے ٹکے سارے ممکن نمونے لکھ کے پرکھیے کہ $!A$ ہر اک وچ سچا اے۔ جے نہیں، تاں رک جاؤ، $!A$ متعلقہ نمونیاں دی جماعت وچ معتبر نہیں۔

ایس خیال وچ اک اوکھ اے۔ سانوں پتا نہیں کہ تلاش کدوں مکائیے، یا کدی مکا وی سکدے آں۔ جے !!{formula} دا کوئی محدود ردّی نمونہ اے، تاں ساڈا طریقہ اوہنوں لبھ لے گا۔ پر جے اوہدا کوئی محدود ردّی نمونہ نہیں، تاں سانوں جواب نہیں ملے گا۔ !!{formula} معتبر ہو سکدا اے (کوئی ردّی نمونہ ای نہیں)، یا اوہدا صرف لامتناہی ردّی نمونہ ہو سکدا اے، جیہڑا اسیں کدی نہیں ویکھاں گے۔ ایہ اوکھ دور ہو سکدی اے جے اسیں وکھا دیئے کہ ہر !!{formula} جیہدا کوئی ردّی نمونہ اے، اوہدا اک محدود ردّی نمونہ وی اے۔ جے ایہ گل سچی اے، تاں اسیں آکھدے آں کہ منطق وچ \emph{محدود نمونے دی خاصیت} اے۔

پر اسیں کیویں وکھائیے کہ کسے منطق وچ محدود نمونے دی خاصیت اے؟ اک طریقہ ایہ ہو سکدا اے کہ~$!A$ دے لامتناہی (ردّی) نمونے نوں محدود نمونے وچ بدلّن دا طریقہ لبھیے۔ جے ایہ ہو سکدا اے، تاں جدوں وی کوئی نمونہ اے جیہدے وچ $!A$ سچا نہیں، تاں ایس توں بنیا محدود نمونہ وی $!A$ نوں سچا نہیں کردا۔ ایہ محدود نمونہ سارے محدود نمونیاں دی ساڈی فہرست وچ آوے گا، تے ہر غیر معتبر فارمولے لئی سانوں آخر پتا لگے گا کہ اوہ معتبر نہیں۔ جے فارمولا معتبر اے، تاں ساڈا طریقہ مُکے گا نہیں۔ جے اسیں نال ایہ وی وکھا دیئے کہ ساڈے طریقے توں ملن والے محدود نمونے دا کوئی ودھ توں ودھ سائز اے، تے ایہ سائز صرف !!{formula}~$!A$ اُتے نربھر اے، تاں ردّی نمونے دی تلاش وچ ساڈے کول محدود نمونیاں دی جانچ دی اک حد ہووے گی۔ جے اوس حد تک کوئی ردّی نمونہ نہیں لبھیا، تاں کوئی ہے ای نہیں۔ فیر ساڈا طریقہ «کی $!A$ معتبر اے؟» دا فیصلہ ہر !!{formula}~$!A$ لئی کرے گا۔

لامتناہی ساختاں نوں محدود ساختاں وچ بدلّن لئی اک طریقہ اکثر کم کردا اے۔ ساخت دے اوہناں !!{element}s نوں «اکّو» منّ لئیے جیہڑے متعلقہ پاسیاں توں اکّو طرح کم کردے نیں۔ جے~$\mModel{M}$ وچ لامتناہی گنتی دے جہان متعلقہ پاسیاں توں اکّو طرح کم کردے نیں، تاں ہو سکدا اے اسیں \emph{سارے} جہان ایہو جیہے جہاناں دے محدود گنتی والے (ممکن اے لامتناہی) «طبقیاں» وچ اکٹھے کر لئیے۔ دوجے لفظاں وچ جہاناں دے سیٹ نوں ٹھیک طرح ونڈیے، یعنی ہر حصے وچ لامتناہی گنتی دے جہان ہون پر حصے صرف محدود گنتی دے ہون۔ فیر نواں نمونہ~$\mModel{M^*}$ متعین کردے آں جیہدے جہان اوہ حصے نیں۔ پرانے نمونے دے محدود گنتی والے حصیاں توں نویں نمونے دے محدود گنتی والے جہان ملدے نیں، یعنی اک محدود نمونہ۔ جہان~$w$ جیہڑے حصے وچ اے، اوہنوں $[w]$ آکھ لئیے۔ اسیں چاہواں گے کہ $\mSat{M}{!A}[w]$ اُدوں تے صرف اُدوں جدوں
$\mSat{M^*}{!A}[{[w]}]$، کیوں جے اسیں چاہندے آں کہ جے پرانا نمونہ~$!A$ دا ردّی نمونہ سی، تاں نواں وی ہووے۔ ایس لئی ونڈ تے~$\mModel{M^*}$ دے رسائی دے تعلق دوناں نوں ٹھیک طرح متعین کرنا پئے گا۔

ایہ طریقہ ویکھن لئی پہلاں خیال کرو کہ ساڈے کول رسائی دا کوئی تعلق نہیں۔ $\mSat{M}{\Box !B}[w]$ اُدوں تے صرف اُدوں جدوں کسے $v \in W$ لئی
$\mSat{M}{\Box !B}[v]$، تے $\mModel{M^*}$ لئی وی ایویں ای، بس
$[w]$ تے $[v]$ نال۔ پہلے خیال وجوں آؤ آکھیے کہ دو جہان $u$ تے
$v$ ہم ارز نیں (اکّو حصے وچ نیں) جے اوہ سارے !!{propositional variable}s اُتے~$\mModel{M}$ وچ اکّو فیصلے دیندے نیں، یعنی $\mSat{M}{p}[u]$
اُدوں تے صرف اُدوں جدوں $\mSat{M}{p}[v]$۔ منّ لو $V^*(p) = \Setabs{[w]}{\mSat{M}{p}[w]}$۔ ساڈا مقصد وکھاؤنا اے کہ $\mSat{M}{!A}[w]$ اُدوں تے صرف اُدوں جدوں
$\mSat{M^*}{!A}[{[w]}]$۔ صاف اے کہ اسیں استقرا نال ثابت کراں گے۔ ابتدائی معاملہ $!A \equiv p$ ہووے گا۔ پہلاں منّ لو $\mSat{M}{p}[w]$۔ فیر تعریف مطابق
$[w] \in V^*$، سو $\mSat{M^*}{p}[{[w]}]$۔ ہُن منّ لو
$\mSat{M^*}{p}[{[w]}]$۔ ایہدا مطلب اے $[w] \in V^*(p)$، یعنی کسے $v$ لئی جیہڑا~$w$ دے ہم ارز اے، $\mSat{M}{p}[v]$۔ پر «$w$ دا
$v$ دے ہم ارز ہونا» ایہدا مطلب اے کہ «$w$ تے $v$ عین اوہی !!{propositional variable}s سچے کردے نیں»، سو $\mSat{M}{p}[w]$۔ ہُن استقرائی قدم لئی مثلاً $!A
\ident \lnot !B$۔ فیر $\mSat{M}{\lnot !B}[w]$ اُدوں تے صرف اُدوں جدوں $\mSat/{M}{!B}[w]$، ایہ اُدوں تے صرف اُدوں جدوں $\mSat/{M^*}{!B}[{[w]}]$ (استقرائی مفروضے مطابق)، ایہ اُدوں تے صرف اُدوں جدوں
$\mSat{M^*}{\lnot !B}[{[w]}]$۔ باقی غیر امکانی عاملاں لئی وی ایویں اے۔ $\Box$ لئی وی ایہ چلدا اے۔ منّ لو $\mSat{M^*}{\Box
  !B}[{[w]}]$۔ ایہدا مطلب اے کہ ہر $[u]$ لئی
$\mSat{M^*}{!B}[{[u]}]$۔ استقرائی مفروضے مطابق ہر $u$ لئی
$\mSat{M}{!B}[u]$۔ ایس لئی $\mSat{M}{\Box !B}[w]$۔

عام صورت وچ، جتھے $\mModel{M^*}$ دا رسائی دا تعلق وی متعین کرنا پئے گا، گل ہور اوکھی اے۔ نمونے~$\mModel{M^*}$ نوں اسیں \emph{فلٹریشن} آکھاں گے جے اوہدا رسائی دا تعلق~$R^*$ اوہ شرطاں پوریاں کرے جیہڑیاں اُپرلا استقرائی ثبوت چلاؤن لئی چاہیدیاں نیں۔ فیر کوئی وی فلٹریشن~$\mModel{M^*}$، $!A$ نوں $[w]$ وچ اُدوں تے صرف اُدوں سچا کرے گی جدوں $\mModel{M}$، $!A$ نوں~$w$ وچ سچا کردا اے۔ پر ہُن سانوں ایہ وی وکھاؤنا اے کہ فلٹریشن \emph{ہے}، یعنی اسیں~$R^*$ ایویں متعین کر سکدے آں کہ اوہ مطلوب شرطاں پوریاں کرے۔ ایس کم لئی سانوں شرط لاؤنی پئے گی کہ جہان $u$ تے
$v$ اُدوں ہم ارز منّے جان جدوں اوہ سارے !!{propositional variable}s تے نال ای سارے ذیلی !!{formula}s جیہڑے~$!A$ دے نیں، اُتے اکّو فیصلے دیندے ہون۔ کیوں جے $!A$ دے صرف محدود گنتی دے ذیلی !!{formula}s نیں، ایس لئی ایہ ہن وی پکا کرے گا کہ فلٹریشن محدود اے۔ فلٹریشن متعین کرن دا اکّو طریقہ نہیں، تے فلٹریشن دے رسائی دے تعلق وچ مطلوب خاصیتاں (مثلاً انعکاس، تعدّی وغیرہ) پکا کرن لئی~$R^*$ دی تعریف سوچ سمجھ کے چُننی پئے گی۔
\endgroup
% /OLP-0451
% OLP-0452 content/normal-modal-logic/filtrations/preliminaries.tex
\begingroup
\def\olfilename{preliminaries}
\expandafter\def\csname ol@sectioncs\endcsname{section}

\olfileid{nml}{fil}{pre}

\olsection{مُڈھلیاں گلاں}

فلٹریشناں نال اسیں اپنے امکانی منطق دے نظاماں دی فیصلہ کرن دی اہلیت قائم کر سکدے آں۔ ایس لئی وکھاؤنے آں کہ اوہناں وچ \emph{محدود نمونے دی خاصیت} اے، یعنی کوئی وی !!{formula} جیہڑا اک نمونے وچ سچا (جھوٹھا) اے، اک \emph{محدود} نمونے وچ وی سچا (جھوٹھا) اے۔ فلٹریشناں اوہناں !!{formula}s دے سیٹاں دے لحاظ نال متعین ہوندیاں نیں جیہڑے ذیلی فارمولیاں تحت بند ہون۔

\begin{defn}\ollabel{defn:modallyclosed}
  !!{formula}s دا سیٹ $\Gamma$ \emph{ذیلی فارمولیاں تحت بند} اے جے ہر اوہ !!a{formula} جیہڑا $\Gamma$ وچ موجود اے، اوہدا ہر ذیلی فارمولا اوہدے وچ شامل ہووے۔ نالے
  $\Gamma$ \emph{امکانی عاملاں تحت بند} اے جے ایہ ذیلی فارمولیاں تحت بند اے تے $!A \in \Gamma$ توں $\Box!A,
  \Diamond!A \in \Gamma$ نکلدا اے۔
\end{defn}

مثلاً دِتّے !!a{formula}~$!A$ دے سارے ذیلی !!{formula}s دا سیٹ ذیلی !!{formula}s تحت بند اے۔ جدوں اسیں کسے نمونے دی فلٹریشن ذیلی !!{formula}s دے سیٹ، جیہڑا~$!A$ دا اے، راہیں متعین کردے آں، تاں اوہدے وچ اوہ خاصیت ہووے گی جیہڑی سانوں چاہیدی اے۔ ایہ $!A$ نوں اُدوں تے صرف اُدوں سچا (جھوٹھا) کردی اے جدوں اصل نمونہ ایویں کردا اے۔

~$\mModel{M}$ دی~$\Gamma$ راہیں فلٹریشن دے جہاناں دا سیٹ ہِیٹھلے ہم ارزیت دے تعلق دے سارے ہم ارزیت دے طبقیاں دا سیٹ اے۔

\begin{defn}
منّ لو $\mModel{M} =\tuple{W, R, V}$ تے $\Gamma$ ذیلی !!{formula}s تحت بند اے۔ تعلق $\equiv$، $W$ اُتے، ایویں متعین کرو کہ ایہ اوہناں کسے وی دو جہاناں وچ قائم اے جیہڑے اکّو !!{formula}s، $\Gamma$ وچوں، سچے کردے نیں، یعنی
\[
u \equiv v \quad \text{اُدوں تے صرف اُدوں جدوں } \quad
\forall !A \in \Gamma : \mSat{M}{!A}[u] \Leftrightarrow \mSat{M}{!A}[v].
\]
ہم ارزیت دا طبقہ~$[w]_\equiv$، جہان~$w$ دا، یا مختصر $[w]$، سارے اوہناں جہاناں دا سیٹ اے جیہڑے $\equiv$ دے لحاظ نال~$w$ دے ہم ارز نیں۔
\[
[w] = \Setabs{v}{v \equiv w}.
\]
\end{defn}

\begin{prop}
  دِتّے $\mModel{M}$ تے $\Gamma$ لئی، اُتے متعین کیتا $\equiv$ ہم ارزیت دا تعلق اے، یعنی ایہ انعکاسی، متقابل، تے متعدی اے۔
\end{prop}

\begin{proof}
  تعلق $\equiv$ انعکاسی اے، کیوں جے $w$ عین اوہی !!{formula}s،~$\Gamma$ وچوں، سچے کردا اے جیہڑے اوہ آپ سچے کردا اے۔ ایہ متقابل اے، کیوں جے جے $u$ اوہی !!{formula}s،~$\Gamma$ وچوں، سچے کردا اے جیہڑے $v$، تاں ایہو گل $v$ تے~$u$ لئی وی اے۔ ایہ متعدی وی اے، کیوں جے جے $u$ اوہی !!{formula}s،~$\Gamma$ وچوں، سچے کردا اے جیہڑے~$v$، تے $v$ اوہی جیہڑے $w$، تاں
  $u$ اوہی !!{formula}s،~$\Gamma$ وچوں، سچے کردا اے جیہڑے~$w$۔
\end{proof}

تعلق $\equiv$، ہر ہم ارزیت دے تعلق وانگ،~$W$ نوں \emph{حصیاں} وچ ونڈدا اے۔ یعنی~$W$ دے ایسے ذیلی سیٹاں وچ جیہڑے جوڑی جوڑی بے سانجھے نیں تے رل کے سارا~$W$ پورا کردے نیں۔ ہر $w \in W$ کسے اک حصے دا !!a{element} اے، یعنی~$[w]$ دا، کیوں جے $w \equiv w$۔ سو حصے $[w]$ سارا~$W$ پورا کردے نیں۔ ایہ جوڑی جوڑی بے سانجھے نیں، کیوں جے جے
$u \in [w]$ تے $u \in [v]$، تاں $u \equiv w$ تے $u \equiv v$، تے تقابل تے تعدّی کارن $w \equiv v$، ایس لئی $[w] = [v]$۔
\endgroup
% /OLP-0452
% OLP-0453 content/normal-modal-logic/filtrations/filtrations-def.tex
\begingroup
\def\olfilename{filtrations-def}
\expandafter\def\csname ol@sectioncs\endcsname{section}

\olfileid{nml}{fil}{fil}

\olsection{فلٹریشناں}

$\mModel{M}$ دی $\Gamma$ راہیں فلٹریشن نوں اکّو مقرر نمونہ منّن دی بجائے، اسیں ایہ متعین کردے آں کہ نمونہ~$\mModel{M^*}$ کدوں~$\mModel{M}$ دی فلٹریشن گنیا جاندا اے۔ ساریاں فلٹریشناں دا جہاناں دا سیٹ~$W^*$ تے قدر بندی~$V^*$ اکّو ہوندے نیں۔ پر وکھ وکھ فلٹریشناں دے رسائی دے تعلق~$R^*$ وکھرے ہو سکدے نیں۔ فلٹریشن گنے جان لئی $R^*$ نوں کئی شرطاں پوریاں کرنیاں پیندیاں نیں۔ ایہ عین اوہی شرطاں نیں جیہڑیاں مُکھ نتیجہ ثابت کرن لئی چاہیدیاں نیں، یعنی $\mSat{M}{!A}[w]$ اُدوں تے صرف اُدوں جدوں $\mSat{M^*}{!A}[{[w]}]$، بشرطیکہ $!A \in
\Gamma$۔

\begin{defn}\ollabel{defn:filtration}
  منّ لو $\Gamma$ ذیلی فارمولیاں تحت بند اے تے $\mModel{M} = \tuple{W,
    R, V}$۔ \emph{$\mModel{M}$ دی $\Gamma$ راہیں فلٹریشن} کوئی وی نمونہ $\mModel{M^*} = \tuple{W^*,R^*,V^*}$ اے جیہدے لئی
  \begin{enumerate}
  \item $W^* = \Setabs{[w]}{w \in W}$؛
  \item \ollabel{defn:filtration-R}%
    ہر $u,v \in W$ لئی
    \begin{enumerate}
    \item \ollabel{defn:filtration-R1}%
      جے $Ruv$، تاں $R^*[u][v]$؛
    \tagitem{prvBox}{\ollabel{defn:filtration-R2}%
      جے $R^*[u][v]$، تاں ہر $\Box !A \in \Gamma$ لئی، جے
      $\mSat{M}{\Box!A}[u]$ تاں $\mSat{M}{!A}[v]$}{}؛
    \tagitem{prvDiamond}{\ollabel{defn:filtration-R3}%
      جے $R^*[u][v]$، تاں ہر $\Diamond !A \in \Gamma$ لئی، جے
      $\mSat{M}{!A}[v]$ تاں $\mSat{M}{\Diamond!A}[u]$۔}{}
    \end{enumerate}
  \item $V^*(p) = \Setabs{[u]}{u \in V(p)}$۔
  \end{enumerate}
\end{defn}

ایہ سوچنا فائدے مند اے کہ $V^*(p)$ کیہ اے۔ ایہ ہم ارزیت دے طبقیاں~$[w]$ دا سیٹ اے، سارے اوہناں جہاناں~$w$ لئی جتھے $p$،~$\mModel{M}$ وچ سچا اے۔ اک پاسے جے $w \in V(p)$، تاں $[w] \in
V^*(p)$ تعریف مطابق۔ پر ایہ ضروری نہیں کہ جے $[w] \in V^*(p)$، تاں $w \in V(p)$۔ جے $[w] \in V^*(p)$، تاں صرف ایہ گل پکی اے کہ $[w] = [u]$ \emph{کسے}~$u \in V(p)$ لئی۔ صاف اے کہ $[w] = [u]$ دا مطلب $w \equiv u$ اے۔ سو جدوں $[w] \in
V^*(p)$، اسیں (صرف) ایہ نتیجہ کڈھ سکدے آں کہ $w \equiv u$ کسے $u \in
V(p)$ لئی۔

\begin{thm}\ollabel{thm:filtrations}
  جے $\mModel{M^*}$، $\mModel{M}$ دی $\Gamma$ راہیں فلٹریشن اے، تاں ہر $!A \in \Gamma$ تے $w \in W$ لئی
  $\mSat{M}{!A}[w]$ اُدوں تے صرف اُدوں جدوں $\mSat{M^*}{!A}[{[w]}]$۔
\end{thm}

\begin{proof}
  $!A$ اُتے استقرا نال، ایس حقیقت نوں ورت کے کہ $\Gamma$ ذیلی فارمولیاں تحت بند اے۔ کیوں جے $!A \in \Gamma$ تے $\Gamma$ ذیلی !!{formula}s تحت بند اے، ایس لئی سارے ذیلی !!{formula}s جیہڑے~$!A$ دے نیں، وی $\in
  \Gamma$ نیں۔ سو ہر استقرائی قدم وچ استقرائی مفروضہ ذیلی !!{formula}s اُتے، جیہڑے~$!A$ دے نیں، لاگو ہوندا اے۔
  \begin{enumerate}
    \tagitem{prvFalse}{\indcase{!A}{\lfalse}{نہ
        $\mSat{M}{\indfrm}[w]$، تے نہ $\mSat{M^*}{\indfrm}[{[w]}]$۔}}{}
    \tagitem{prvTrue}{\indcase{!A}{\ltrue}{
        $\mSat{M}{\indfrm}[w]$ تے $\mSat{M^*}{\indfrm}[{[w]}]$ دونیں۔}}{}
    \item \indcase{!A}{p}{کھبے توں سجے والا رخ سدھا اے، کیوں جے
      $\mSat{M}{\indfrm}[w]$ صرف اُدوں جدوں $w \in V(p)$، جیہدے توں $[w]
      \in V^*(p)$ نکلدا اے، یعنی $\mSat{M^*}{\indfrm}[{[w]}]$۔ اُلٹے پاسے منّ لو
      $\mSat{M^*}{\indfrm}[{[w]}]$، یعنی $[w] \in
      V^*(p)$۔ فیر کسے $v \in V(p)$ لئی $w \equiv v$۔ صاف اے کہ فیر
      $\mSat{M}{p}[v]$ وی۔ کیوں جے $w \equiv v$، ایس لئی $w$ تے $v$،
      اکّو !!{formula}s،~$\Gamma$ وچوں، سچے کردے نیں۔ مفروضے مطابق $p
      \in \Gamma$ تے $\mSat{M}{p}[v]$، ایس لئی $\mSat{M}{\indfrm}[w]$۔}

    \tagitem{prvNot}{\iftag{probNot}{مشق۔}{\indcase{!A}{\lnot
          !B}{$\mSat{M}{\indfrm}[w]$ اُدوں تے صرف اُدوں جدوں $\mSat/{M}{!B}[w]$۔
          استقرائی مفروضے مطابق $\mSat/{M}{!B}[w]$ اُدوں تے صرف اُدوں جدوں
          $\mSat/{M^*}{!B}[{[w]}]$۔ آخر $\mSat/{M^*}{!B}[{[w]}]$
          اُدوں تے صرف اُدوں جدوں $\mSat{M^*}{\indfrm}[{[w]}]$۔}}}{}

    \tagitem{prvAnd}{\iftag{probAnd}{مشق۔}{\indcase{!A}{(!B \land
          !C)}{$\mSat{M}{\indfrm}[w]$ اُدوں تے صرف اُدوں جدوں $\mSat{M}{!B}[w]$ تے
          $\mSat{M}{!C}[w]$۔ استقرائی مفروضے مطابق $\mSat{M}{!B}[w]$
          اُدوں تے صرف اُدوں جدوں $\mSat{M^*}{!B}[{[w]}]$، تے $\mSat{M}{!C}[w]$ اُدوں تے صرف اُدوں جدوں
          $\mSat{M^*}{!C}[{[w]}]$۔ تے $\mSat{M^*}{\indfrm}[{[w]}]$
          اُدوں تے صرف اُدوں جدوں $\mSat{M^*}{!B}[{[w]}]$ تے
          $\mSat{M^*}{!C}[{[w]}]$۔}}}{}

    \tagitem{prvOr}{\iftag{probOr}{مشق۔}{\indcase{!A}{(!B \lor
          !C)}{$\mSat{M}{\indfrm}[w]$ اُدوں تے صرف اُدوں جدوں $\mSat{M}{!B}[w]$ یا
          $\mSat{M}{!C}[w]$۔ استقرائی مفروضے مطابق $\mSat{M}{!B}[w]$
          اُدوں تے صرف اُدوں جدوں $\mSat{M^*}{!B}[{[w]}]$، تے $\mSat{M}{!C}[w]$ اُدوں تے صرف اُدوں جدوں
          $\mSat{M^*}{!C}[{[w]}]$۔ تے $\mSat{M^*}{\indfrm}[{[w]}]$
          اُدوں تے صرف اُدوں جدوں $\mSat{M^*}{!B}[{[w]}]$ یا
          $\mSat{M^*}{!C}[{[w]}]$۔}}}{}

    \tagitem{prvIf}{\iftag{probIf}{مشق۔}{\indcase{!A}{(!B \lif
          !C)}{$\mSat{M}{\indfrm}[w]$ اُدوں تے صرف اُدوں جدوں $\mSat/{M}{!B}[w]$ یا
          $\mSat{M}{!C}[w]$۔ استقرائی مفروضے مطابق $\mSat{M}{!B}[w]$
          اُدوں تے صرف اُدوں جدوں $\mSat{M^*}{!B}[{[w]}]$، تے $\mSat{M}{!C}[w]$ اُدوں تے صرف اُدوں جدوں
          $\mSat{M^*}{!C}[{[w]}]$۔ تے $\mSat{M^*}{\indfrm}[{[w]}]$
          اُدوں تے صرف اُدوں جدوں $\mSat/{M^*}{!B}[{[w]}]$ یا
          $\mSat{M^*}{!C}[{[w]}]$۔}}}{}

    \tagitem{prvIff}{\iftag{probIff}{مشق۔}{\indcase{!A}{(!B \liff
          !C)}{$\mSat{M}{\indfrm}[w]$ اُدوں تے صرف اُدوں جدوں یا تاں $\mSat{M}{!B}[w]$ تے
          $\mSat{M}{!C}[w]$ دونیں، یا $\mSat/{M}{!B}[w]$ تے
          $\mSat/{M}{!C}[w]$ دونیں۔ استقرائی مفروضے مطابق $\mSat{M}{!B}[w]$
          اُدوں تے صرف اُدوں جدوں $\mSat{M^*}{!B}[{[w]}]$، تے $\mSat{M}{!C}[w]$ اُدوں تے صرف اُدوں جدوں
          $\mSat{M^*}{!C}[{[w]}]$۔ تے $\mSat{M^*}{\indfrm}[{[w]}]$
          اُدوں تے صرف اُدوں جدوں یا تاں $\mSat{M^*}{!B}[{[w]}]$ تے
          $\mSat{M^*}{!C}[{[w]}]$ دونیں، یا $\mSat/{M^*}{!B}[{[w]}]$ تے
          $\mSat/{M^*}{!C}[{[w]}]$ دونیں۔}}}{}

    \tagitem{prvBox}{\iftag{probBox}{مشق۔}{\indcase{!A}{\Box
          !B}{منّ لو $\mSat{M}{\indfrm}[w]$۔
          $\mSat{M^*}{\indfrm}[{[w]}]$ وکھاؤن لئی ایسا $v$ لے لو جیہدے لئی
          $R^*[w][v]$۔
          \olref{defn:filtration}\olref{defn:filtration-R2} توں
          $\mSat{M}{!B}[v]$ ملدا اے، تے استقرائی مفروضے توں
          $\mSat{M^*}{!B}[{[v]}]$۔ کیوں جے $v$ کوئی وی ایسا جہان ہو سکدا سی،
          ایس لئی $\mSat{M^*}{\indfrm}[{[w]}]$ نکلدا اے۔

          اُلٹے پاسے منّ لو $\mSat{M^*}{\indfrm}[{[w]}]$ تے کوئی وی ایسا $v$
          لے لو جیہدے لئی $Rwv$۔
          \olref{defn:filtration}\olref{defn:filtration-R1} توں
          $R^*[w][v]$ ملدا اے، سو $\mSat{M^*}{!B}[{[v]}]$، تے استقرائی مفروضے توں $\mSat{M}{!B}[v]$۔ کیوں جے $v$ کوئی وی ایسا جہان ہو سکدا سی،
          ایس لئی $\mSat{M}{\indfrm}[w]$۔}}}{}

    \tagitem{prvDiamond}{\iftag{probDiamond}{مشق۔}{\indcase{!A}{\Diamond
          !B}{منّ لو $\mSat{M}{\indfrm}[w]$۔ فیر کسے $v \in W$ لئی
          $Rwv$ تے $\mSat{M}{!B}[v]$۔ استقرائی مفروضے مطابق
          $\mSat{M^*}{!B}[{[v]}]$، تے
          \olref{defn:filtration}\olref{defn:filtration-R1} مطابق
          $R^*[w][v]$۔ سو $\mSat{M^*}{\indfrm}[{[w]}]$۔

          ہُن منّ لو $\mSat{M^*}{\indfrm}[{[w]}]$۔ فیر کسے $[v]
          \in W^*$ لئی جیہدے نال $R^*[w][v]$، $\mSat{M^*}{!B}[{[v]}]$۔
          استقرائی مفروضے مطابق $\mSat{M}{!B}[v]$۔
          \olref{defn:filtration}\olref{defn:filtration-R3} مطابق
          $\mSat{M}{\indfrm}[w]$۔}}}{}
  \end{enumerate}
\end{proof}

\begin{prob}
  \olref[nml][fil][fil]{thm:filtrations} دا ثبوت پورا کرو۔
\end{prob}

جیہڑی گل نمونے دے جہاناں وچ سچائی لئی قائم اے، اوہی گل نمونے وچ سچائی تے نمونیاں دی جماعت وچ اعتبار لئی وی قائم اے۔

\begin{cor}
  منّ لو $\Gamma$ ذیلی فارمولیاں تحت بند اے۔ فیر
  \begin{enumerate}
  \item جے $\mModel{M^*}$، $\mModel{M}$ دی
    $\Gamma$ راہیں فلٹریشن اے، تاں ہر $!A \in \Gamma$ لئی $\mSat{M}{!A}$ اُدوں تے صرف اُدوں جدوں $\mSat{M^*}{!A}$۔
  \item جے $\mClass{C}$ نمونیاں دی جماعت اے تے $\Gamma(\mClass{C})$،
    $\Gamma$-فلٹریشناں دی اوہ جماعت اے جیہڑیاں $\mClass{C}$ دے نمونیاں توں بنیاں نیں، تاں کوئی وی !!{formula}~$!A \in \Gamma$، $\mClass{C}$ وچ اُدوں تے صرف اُدوں معتبر اے جدوں اوہ $\Gamma(\mClass{C})$ وچ معتبر ہووے۔
  \end{enumerate}
\end{cor}
\endgroup
% /OLP-0453
% OLP-0454 content/normal-modal-logic/filtrations/examples-of-filtrations.tex
\begingroup
\def\olfilename{examples-of-filtrations}
\expandafter\def\csname ol@sectioncs\endcsname{section}

\olfileid{nml}{fil}{exf}

\olsection{فلٹریشناں دیاں مثالاں}

اسیں ہن تک ایہ نہیں وکھایا کہ کوئی فلٹریشن ہے۔ پر ہر نمونے $\mModel{M}$ لئی $\mModel{M}$ دیاں $\Gamma$ راہیں کئی فلٹریشناں نیں۔ خاص طور تے دو پچھاندے آں، سبھ توں باریک تے سبھ توں موٹی فلٹریشن۔ اکّو نمونے دیاں فلٹریشناں دے رسائی دے تعلق وکھرے ہون گے (کیوں جے \olref[fil]{defn:filtration} سدھا مقرر کردی اے کہ $W^*$ تے~$V^*$ کیہ ہون)۔ سبھ توں باریک فلٹریشن وچ تعلق والے جہاناں دیاں جوڑیاں جِنّیاں ہو سکّن گھٹ ہون گیاں، تے سبھ توں موٹی وچ جِنّیاں ہو سکّن ودھ۔

\begin{defn}
  جتھے $\Gamma$ ذیلی فارمولیاں تحت بند اے، نمونے دی \emph{سبھ توں باریک} فلٹریشن $\mModel{M^*}$، نمونے $\mModel{M}$ دی، ایویں متعین اے۔
  \[
  R^*[u][v] \quad \text{اُدوں تے صرف اُدوں جدوں} \quad \exists u'\in [u] \;
  \exists v' \in [v] : Ru'v'.
  \]
\end{defn}

\begin{prop}\ollabel{prop:finest}
  سبھ توں باریک فلٹریشن $\mModel{M^*}$ سچ مچ فلٹریشن اے۔
\end{prop}

\begin{proof}
  سانوں پرکھنا اے کہ ایویں متعین کیتا $R^*$،
  \olref[fil]{defn:filtration}\olref[fil]{defn:filtration-R} دیاں شرطاں پوریاں کردا اے۔ اسیں تینوں شرطاں واری واری پرکھدے آں۔

  جے $Ruv$، تاں کیوں جے $u \in [u]$ تے $v \in [v]$، ایس لئی $R^*[u][v]$ وی،
  سو \olref[fil]{defn:filtration-R1} پوری اے۔

  \iftag{prvBox}{\iftag{probBox}{
      \olref[fil]{defn:filtration-R2} دی پرکھ مشق لئی چھڈدے آں۔}{
        \olref[fil]{defn:filtration-R2} لئی منّ لو $\Box!A \in \Gamma$،
        $R^*[u][v]$ تے $\mSat{M}{\Box!A}[u]$۔
        $R^*$ دی تعریف مطابق ایسے $u' \equiv u$ تے $v' \equiv v$ نیں جیہناں لئی
        $Ru'v'$۔ کیوں جے $u$ تے $u'$، $\Gamma$ اُتے اکّو فیصلے دیندے نیں، ایس لئی
        $\mSat{M}{\Box!A}[u']$ وی، سو $\mSat{M}{!A}[v']$۔ کیوں جے
        $\Gamma$ ذیلی !!{formula}s تحت بند اے، $v$ تے $v'$،
        $!A$ اُتے اکّو فیصلے دیندے نیں، سو $\mSat{M}{!A}[v]$، جیویں وکھاؤنا سی۔}}{}

  \iftag{prvDiamond}{\iftag{probDiamond}{
      \olref[fil]{defn:filtration-R3} دی پرکھ مشق لئی چھڈدے آں۔}{
        \olref[fil]{defn:filtration-R3} پرکھن لئی منّ لو $\Diamond!A \in
        \Gamma$، $R^*[u][v]$ تے $\mSat{M}{!A}[v]$۔
        $R^*$ دی تعریف مطابق ایسے $u' \equiv u$ تے $v' \equiv v$ نیں جیہناں لئی
        $Ru'v'$۔ کیوں جے $v$ تے $v'$،~$\Gamma$ اُتے اکّو فیصلے دیندے نیں، تے $\Gamma$ ذیلی !!{formula}s تحت بند اے، ایس لئی $\mSat{M}{!A}[v']$ وی، سو
        $\mSat{M}{\Diamond !A}[u']$۔ کیوں جے $u$ تے $u'$ وی
        $\Gamma$ اُتے اکّو فیصلے دیندے نیں، ایس لئی $\mSat{M}{\Diamond !A}[u]$۔}}{}
\end{proof}

\begin{probtag}{probBox,probDiamond}
  \olref[nml][fil][exf]{prop:finest} دا ثبوت پورا کرو۔
\end{probtag}

\begin{defn}
  جتھے $\Gamma$ ذیلی فارمولیاں تحت بند اے، نمونے دی \emph{سبھ توں موٹی}
  فلٹریشن~$\mModel{M^*}$، نمونے~$\mModel{M}$ دی، ایویں متعین اے کہ $R^*[u][v]$ اُدوں تے صرف اُدوں ہووے جدوں
  \iftag{notprvBox,notprvDiamond}{ہِیٹھلی شرط پوری ہووے۔}{ہِیٹھلیاں \emph{دونیں} شرطاں پوریاں ہون۔}
  \begin{tagenumerate}{prvBox,prvDiamond}
  \tagitem{prvBox}{\ollabel{defn:coarsest-Box}جے $\Box!A \in \Gamma$
    تے $\mSat{M}{\Box!A}[u]$، تاں
    $\mSat{M}{!A}[v]$\iftag{prvDiamond}{؛}{۔}}{}
  \tagitem{prvDiamond}{\ollabel{defn:coarsest-Diamond}جے $\Diamond!A
    \in \Gamma$ تے $\mSat{M}{!A}[v]$، تاں
    $\mSat{M}{\Diamond!A}[u]$۔}{}
  \end{tagenumerate}
\end{defn}

\begin{prop}
  سبھ توں موٹی فلٹریشن $\mModel{M^*}$ سچ مچ فلٹریشن اے۔
\end{prop}

\begin{proof}
  $R^*$ دی تعریف دِتّی ہون تے پرکھن لئی صرف ایہ شرط رہ جاندی اے کہ $Ruv$ توں $R^*[u][v]$ نکلدا اے۔ سو منّ لو
  $Ruv$۔ \iftag{prvBox}{منّ لو $\Box!A \in \Gamma$ تے
    $\mSat{M}{\Box!A}[u]$۔ فیر صاف اے کہ
    $\mSat{M}{!A}[v]$\iftag{prvDiamond}{، تے
      \olref{defn:coarsest-Box} پوری اے۔ }{۔}}{}\iftag{prvDiamond}{منّ لو $\Diamond!A \in
    \Gamma$ تے $\mSat{M}{!A}[v]$۔ فیر $\mSat{M}{\Diamond!A}[u]$
    کیوں جے~$Ruv$\iftag{prvBox}{، تے \olref{defn:coarsest-Diamond} پوری اے}{}۔}{}
\end{proof}

\begin{ex}
  منّ لو $W = \PosInt$، $Rnm$ اُدوں تے صرف اُدوں جدوں $m = n + 1$، تے $V(p) = \Setabs{2n}{n \in
    \Nat}$۔ نمونہ $\mModel{M} = \tuple{W, R, V}$،
  \olref{fig:ex-filtration} وچ وکھایا گیا اے۔ جہان $1$، $2$ وغیرہ نیں۔ ہر جہان توں ٹھیک اک ہور جہان، یعنی اوہدا اگلا، رسائی وچ اے، تے $p$ سارے جفت عدداں اُتے سچا اے تے صرف اوہناں اُتے۔

  \begin{figure}
    \centering
    \begin{tikzpicture}[modal]
      \node[world] (1) [label=below:\mFalse{p}] {$1$};
      \node[world] (2) [label=below:\mTrue{p}, right=of 1] {$2$};
      \node[world] (3) [label=below:\mFalse{p}, right=of 2] {$3$};
      \node[world] (4) [label=below:\mTrue{p}, right=of 3] {$4$};
      \node[phantom] (5) [right=of 4] {};
      \draw[->] (1) to (2);
      \draw[->] (2) to (3);
      \draw[->] (3) to (4);
      \draw[dotted] (4) to (5);
    \end{tikzpicture}

    \begin{tikzpicture}[modal]
      \node[world] (1) [label=below:\mFalse{p}] {$[1]$};
      \node[world] (2) [label=below:\mTrue{p}, right=of 1] {$[2]$};
      \draw[->,bend left] (1) to (2);
      \draw[->,bend left] (2) to (1);
      
      \node[world] (3) [label=below:\mFalse{p},right=of 2] {$[1]$};
      \node[world] (4) [label=below:\mTrue{p}, right=of 3] {$[2]$};
      \draw[->,bend left] (3) to (4);
      \draw[->,bend left] (4) to (3);
      \draw[->,reflexive above] (4) to (4);
    \end{tikzpicture}
    \caption{اک لامتناہی نمونہ تے اوہدیاں فلٹریشناں۔}
    \ollabel{fig:ex-filtration}
  \end{figure}

  ہُن منّ لو $\Gamma$ ذیلی !!{formula}s دا سیٹ، جیہڑا~$\Box p \lif p$ دا اے،
  یعنی $\{p, \Box p, \Box p \lif p\}$۔ $p$ سارے جفت عدداں اُتے تے صرف اوہناں اُتے سچا اے، $\Box p$ سارے طاق عدداں اُتے تے صرف اوہناں اُتے سچا اے، سو $\Box p \lif p$ سارے جفت عدداں اُتے تے صرف اوہناں اُتے سچا اے۔ دوجے لفظاں وچ ہر طاق عدد $\Box p$ نوں سچا پر $p$ تے $\Box
  p \lif p$ نوں جھوٹھا کردا اے۔ ہر جفت عدد $p$ تے $\Box p \lif p$ نوں سچا پر $\Box p$ نوں جھوٹھا کردا اے۔ سو $W^* = \{ [1], [2] \}$، جتھے $[1] =
  \{1, 3, 5, \dots\}$ تے $[2] = \{2, 4, 6, \dots\}$۔ کیوں جے $2 \in
  V(p)$، ایس لئی $[2] \in V^*(p)$۔ کیوں جے $1 \notin V(p)$، ایس لئی $[1] \notin
  V^*(p)$۔ سو $V^*(p) = \{[2]\}$۔

  $W^*$ اُتے ٹکی ہر فلٹریشن دے رسائی دے تعلق وچ
  $\tuple{[1], [2]}, \tuple{[2],[1]}$ شامل ہونے ضروری نیں۔ کیوں جے $R12$، ایس لئی
  $R^*[1][2]$ ہونا چاہیدا اے،
  \olref[fil]{defn:filtration}\olref[fil]{defn:filtration-R1} مطابق۔ کیوں جے
  $R23$، ایس لئی $R^*[2][3]$ ہونا چاہیدا اے، تے $[3]=[1]$۔ اوہدے وچ
  $\tuple{[1],[1]}$ شامل نہیں ہو سکدا۔ جے ہووے تاں $R^*[1][1]$،
  $\mSat{M}{\Box p}[1]$ پر $\mSat/{M}{p}[1]$، جیہڑا
  \olref[fil]{defn:filtration-R2} دے اُلٹ اے۔ کوئی شرط ایہ نہیں منگدی تے نہ منع کردی اے کہ
  $R^*[2][2]$۔ سو~$\mModel{M}$ دیاں دو ممکن فلٹریشناں نیں، جیہڑیاں رسائی دے اینہاں دو تعلقاں نال جُڑیاں نیں۔
  \[
  \{\tuple{[1],[2]}, \tuple{[2],[1]}\} \text{ تے }
  \{\tuple{[1],[2]}, \tuple{[2],[1]}, \tuple{[2],[2]}\}.
  \]
  دوناں صورتاں وچ $p$ تے $\Box p \lif p$ جھوٹھے نیں تے $\Box p$ سچا اے~$[1]$ اُتے۔ $p$ تے $\Box p \lif p$ سچے نیں تے $\Box p$ جھوٹھا اے~$[2]$ اُتے۔
\end{ex}

\begin{prob}
  ہِیٹھلا نمونہ $\mModel{M} = \tuple{W, R, V}$ ویکھو جیہدے وچ $W
  = \Setabs{0\sigma}{\sigma \in \Bin^*}$، یعنی $0$ تے~$1$ دیاں اوہناں لڑیاں دا سیٹ جیہڑیاں~$0$ نال شروع ہوندیاں نیں، تے $R\sigma\sigma'$ اُدوں تے صرف اُدوں جدوں $\sigma' =
  \sigma 0$ یا $\sigma' = \sigma 1$، تے $V(p) = \Setabs{\sigma
    0}{\sigma \in \Bin^*}$ تے $V(q) = \Setabs{\sigma 1}{\sigma \in
    \Bin^* \setminus \{1\}}$۔ ایہدا خاکہ ایہ اے۔
  \begin{center}
    \begin{tikzpicture}[modal,
        node distance=2cm
      ]

      \node[world] (0) [label={below:\mTrue{p}\\\mFalse{q}},font=\tiny] {$0$};
      \node[world] (00) [label={below:\mTrue{p}\\\mFalse{q}},
        above right=of 0,font=\tiny] {$00$};
      \node[world] (000) [label={below:\mTrue{p}\\\mFalse{q}},
        right=of 00, yshift=.8cm, font=\tiny] {$000$};
      \node[phantom] (0000) [right=of 000, yshift=.5cm] {};
      \node[phantom] (0001) [right=of 000, yshift=-.5cm] {};
      \node[world] (001) [label={below:\mFalse{p}\\\mTrue{q}},
        right=of 00, yshift=-.8cm, font=\tiny] {$001$};
      \node[phantom](0010) [right=of 001, yshift=.5cm] {};
      \node[phantom](0011) [right=of 001, yshift=-.5cm] {};
      \node[world] (01) [label={below:\mFalse{p}\\\mTrue{q}},
        below right=of 0,font=\tiny] {$01$};
      \node[world] (010) [label={below:\mTrue{p}\\\mFalse{q}},
        right=of 01, yshift=.8cm,font=\tiny] {$010$};
      \node[phantom] (0100) [right=of 010, yshift=.5cm] {};
      \node[phantom] (0101) [right=of 010, yshift=-.5cm] {};
      \node[world] (011) [label={below:\mFalse{p}\\\mTrue{q}},
        right=of 01, yshift=-.8cm,font=\tiny] {$011$};
      \node[phantom] (0110) [right=of 011, yshift=.5cm] {};
      \node[phantom] (0111) [right=of 011, yshift=-.5cm] {};

      \draw[->] (0) to (00);
      \draw[->] (0) to (01);
      \draw[->] (00) to (000);
      \draw[dotted] (000) to (0000);
      \draw[dotted] (000) to (0001);
      \draw[->] (00) to (001);
      \draw[dotted] (001) to (0010);
      \draw[dotted] (001) to (0011);
      \draw[->] (01) to (010);
      \draw[dotted] (010) to (0100);
      \draw[dotted] (010) to (0101);
      \draw[->] (01) to (011);
      \draw[dotted] (011) to (0110);
      \draw[dotted] (011) to (0111);
    \end{tikzpicture}
  \end{center}
  $\mSat/{M}{\Box(p \lor q) \lif (\Box p \lor \Box q)}[w]$ ہر~$w$ لئی اے۔

  منّ لو $\Gamma$ ذیلی !!{formula}s دا سیٹ اے جیہڑا~$\Box(p \lor q) \lif
  (\Box p \lor \Box q)$ دا اے۔ $W^*$ تے~$V^*$ کیہ نیں؟~$\mModel{M}$ دی سبھ توں باریک فلٹریشن دا رسائی تعلق کیہ اے؟ سبھ توں موٹی دا؟
\end{prob}
\endgroup
% /OLP-0454
% OLP-0455 content/normal-modal-logic/filtrations/finite.tex
\begingroup
\def\olfilename{finite}
\expandafter\def\csname ol@sectioncs\endcsname{section}

\olfileid{nml}{fil}{fin}

\olsection{فلٹریشناں محدود ہوندیاں نیں}

اسیں ہر اوہ سیٹ~$\Gamma$ لئی فلٹریشناں متعین کیتیاں نیں جیہڑا ذیلی !!{formula}s تحت بند اے۔ تعریف آپ ایہ ضمانت نہیں دیندی کہ فلٹریشناں محدود نیں۔ اصل وچ جدوں $\Gamma$ لامتناہی اے (مثلاً سارے !!{formula}s دا سیٹ اے)، تاں فلٹریشن وی لامتناہی ہو سکدی اے۔ پر جے $\Gamma$ محدود اے (مثلاً جدوں ذیلی !!{formula}s دا سیٹ اے، دِتّے !!{formula}~$!A$ لئی)، تاں~$\Gamma$ راہیں ہر فلٹریشن وی محدود اے۔

\begin{prop}\ollabel{prop:filt-are-finite}
  جے $\Gamma$ محدود اے، تاں ہر فلٹریشن $\mModel{M^*}$، نمونے
  $\mModel{M}$ دی $\Gamma$ راہیں، وی محدود اے۔
\end{prop}

\begin{proof}
  $W^*$ دی گنتی وکھرے طبقیاں~$[w]$ دی گنتی اے، ہم ارزیت دے تعلق~$\equiv$ تحت۔ ایسے طبقے دے کوئی وی دو جہان $u$، $v$، یعنی کوئی وی $u$ تے $v$ جیہناں لئی~$u \equiv v$، سارے !!{formula}s~$!A$،~$\Gamma$ وچوں، اُتے اکّو فیصلے دیندے نیں۔ $!A \in \Gamma$ لئی یا تاں $!A$،
  $u$ تے $v$ دوناں وچ سچا اے، یا دوناں وچوں کسے وچ وی نہیں۔ سو ہر طبقے~$[w]$ نال~$\Gamma$ دا اک ذیلی سیٹ جُڑدا اے، یعنی سارے $!A \in
  \Gamma$ جیہناں لئی $!A$،~$[w]$ دے جہاناں وچ سچا اے۔ کوئی دو وکھرے طبقے $[u]$ تے $[v]$،~$\Gamma$ دے اکّو ذیلی سیٹ نال نہیں جُڑدے۔ کیوں جے، جے سچے !!{formula}s دا سیٹ، $u$ وچ، تے سچے !!{formula}s دا سیٹ، $v$ وچ، اکّو اے، تاں
  $u$ تے $v$،~$\Gamma$ دے سارے فارمولیاں اُتے اکّو فیصلے دیندے نیں، یعنی $u \equiv v$۔ پر فیر $[u] = [v]$۔ سو
  $W^*$ توں $\Pow{\Gamma}$ ول اک !!a{injective} فنکشن اے، ایس لئی $\card{W^*} \le \card{\Pow{\Gamma}}$۔ سو جے $\Gamma$ وچ $n$ جملے نیں، تاں $W^*$ دی گنتی~$2^n$ توں ودھ نہیں۔
\end{proof}
\endgroup
% /OLP-0455
% OLP-0456 content/normal-modal-logic/filtrations/S5-fmp.tex
\begingroup
\def\olfilename{S5-fmp}
\expandafter\def\csname ol@sectioncs\endcsname{section}

\olfileid{nml}{fil}{fmp}

\olsection{\Log{K} تے \Log{S5} وچ محدود نمونے دی خاصیت}

\begin{defn}
  امکانی منطق دے نظام $\Sigma$ وچ \emph{محدود نمونے دی خاصیت} ہوندی اے جے جدوں وی کوئی !!a{formula}~$!A$، $\Sigma$ دے کسے نمونے دے کسے جہان وچ سچا اے، تاں $!A$ کسے \emph{محدود} نمونے،~$\Sigma$ دے، دے کسے جہان وچ وی سچا ہووے۔
\end{defn}

\begin{prop}\ollabel{prop:K-fmp}
  \Log{K} وچ محدود نمونے دی خاصیت اے۔
\end{prop}

\begin{proof}
  \Log{K} معتبر !!{formula}s دا سیٹ اے، یعنی ہر نمونہ~\Log{K} دا نمونہ اے۔
  \olref[fil]{thm:filtrations} مطابق جے
  $\mSat{M}{!A}[w]$، تاں $\mSat{M^*}{!A}[w]$ ہر فلٹریشن لئی جیہڑی~$\mModel{M}$ دی، سیٹ $\Gamma$ راہیں، جیہڑا~$!A$ دے ذیلی !!{formula}s دا اے، بنی اے۔ ہر !!{formula} دے صرف محدود گنتی دے ذیلی !!{formula}s نیں، سو
  $\Gamma$ محدود اے۔ \olref[fin]{prop:filt-are-finite} مطابق
  $\card{W^*} \le 2^n$، جتھے $n$ اوہناں !!{formula}s دی گنتی اے جیہڑے~$\Gamma$ وچ نیں۔ تے کیوں جے \Log{K} نمونیاں اُتے کوئی پابندی نہیں لاؤندی،
  $\mModel{M^*}$ اک \Log{K} دا نمونہ اے۔
\end{proof}

فلٹریشناں نال ایہ وکھاؤن لئی کہ منطق~\Log{L} وچ محدود نمونے دی خاصیت اے، ضروری اے کہ \Log{L} دے نمونے دی فلٹریشن آپ وی \Log{L} دا نمونہ ہووے۔ ایس کم لئی اکثر خاصی محنت چاہیدی اے، تے ہر فلٹریشن \Log{L} دا نمونہ نہیں دیندی۔ پر ہمہ گیر نمونیاں لئی ایہ گل ہن وی چلدی اے۔

\begin{prop}\ollabel{prop:univ-fin}
  منّ لو $\mClass{U}$ ہمہ گیر نمونیاں دی جماعت اے (
  \olref[frd][es5]{prop:S5=univ} ویکھو)، تے $\mClass{U}_\mathrm{Fin}$ سارے محدود ہمہ گیر نمونیاں دی جماعت اے۔ فیر کوئی وی !!{formula}~$!A$، $\mClass{U}$ وچ اُدوں تے صرف اُدوں معتبر اے جدوں اوہ
  $\mClass{U}_\mathrm{Fin}$ وچ معتبر ہووے۔
\end{prop}

\begin{proof}
  محدود ہمہ گیر نمونے ہمہ گیر نمونے نیں، سو کھبے توں سجے والا رخ صاف اے۔ سجے توں کھبے والے رخ لئی منّ لو~$!A$ ہمہ گیر نمونے
  وچ کسے جہان $w$ اُتے جھوٹھا اے، نمونہ $\mModel{M}$ اے۔ منّ لو $\Gamma$ وچ $!A$ تے اوہدے سارے ذیلی فارمولے شامل نیں۔ صاف اے کہ $\Gamma$ محدود اے۔ فلٹریشن
  $\mModel{M^*}$، $\mModel{M}$ دی، لے لو۔ فیر $\mModel{M^*}$ محدود اے،
  \olref[fin]{prop:filt-are-finite} مطابق، تے \olref[fil]{thm:filtrations} مطابق $!A$، $[w]$ اُتے
  $\mModel{M^*}$ وچ جھوٹھا اے۔ ہُن بس ایہ ویکھنا باقی اے کہ $\mModel{M^*}$ ہمہ گیر وی اے۔ دِتّے $u$ تے $v$ لئی مفروضے مطابق $Ruv$، تے
  \olref[fil]{defn:filtration}\olref[fil]{defn:filtration-R} مطابق
  $R^*[u][v]$ وی۔
\end{proof}

\begin{cor}\ollabel{cor:S5fmp}
  \Log{S5} وچ محدود نمونے دی خاصیت اے۔
\end{cor}

\begin{proof}
  \olref[frd][es5]{prop:S5=univ} مطابق جے $!A$ کسے انعکاسی تے اقلیدسی نمونے دے کسے جہان وچ سچا اے، تاں ایہ کسے ہمہ گیر نمونے دے کسے جہان وچ سچا اے۔
  \olref{prop:univ-fin} مطابق ایہ کسے محدود ہمہ گیر نمونے دے کسے جہان وچ سچا اے (یعنی نمونے دی فلٹریشن، ذیلی !!{formula}s دے سیٹ~$!A$ دے راہیں)۔ ہر ہمہ گیر نمونہ انعکاسی تے اقلیدسی وی اے، سو $!A$ کسے محدود انعکاسی اقلیدسی نمونے دے کسے جہان وچ سچا اے۔
\end{proof}

\begin{prob}
  وکھاؤ کہ ہر جہان توں رسائی والے یا انعکاسی نمونے دی کوئی وی فلٹریشن بالترتیب ہر جہان توں رسائی والی یا انعکاسی ہوندی اے۔
\end{prob}

\begin{prob}
  کسے متقابل (متعدی، اقلیدسی) نمونے دی غیر متقابل (غیر متعدی، غیر اقلیدسی) فلٹریشن لبھو۔
\end{prob}
\endgroup
% /OLP-0456
% OLP-0457 content/normal-modal-logic/filtrations/S5-decidable.tex
\begingroup
\def\olfilename{S5-decidable}
\expandafter\def\csname ol@sectioncs\endcsname{section}

\olfileid{nml}{fil}{dec}

\olsection{\Log{S5} فیصلہ کرن جوگ اے}
محدود نمونے دی خاصیت سانوں اک سوکھا طریقہ دیندی اے ایہ وکھاؤن دا کہ سانچیاں نال دِتّے امکانی منطق دے نظام \emph{فیصلہ کرن جوگ} نیں، یعنی ایہ طے کرن لئی اک قابلِ حساب طریقہ اے کہ کوئی !!a{formula} نظام وچ !!{derivable} اے یا نہیں۔

\begin{thm}
  \Log{S5} فیصلہ کرن جوگ اے۔
\end{thm}

\begin{proof}
  منّ لو $!A$ دِتّا اے تے آؤن والے !!{propositional variable}s، $!A$ وچ، $p_1$، \dots، $p_k$ وچوں نیں۔ کیوں جے ہر
  $n$ لئی $n$ جہاناں والے صرف محدود گنتی دے نمونے نیں جیہڑے
  $p_1$، \dots، $p_k$ نوں قدر دیندے نیں، اسیں \emph{نال نال} دونیں فہرستاں بنا سکدے آں۔ \Log{S5} دے سارے قضیے، کسے باقاعدہ طریقے نال ثبوت بنا کے، تے $1$، $2$، \dots جہاناں والے سارے نمونے، ایہ پرکھ کے کہ کی $!A$ ایسے کسے نمونے دے کسے جہان وچ جھوٹھا اے۔ آخر ایہناں نال نال چلدے کمّاں وچوں اک جواب دیوے گا، کیوں جے
  \olref[com][fra]{thm:generaldet} تے \olref[fmp]{cor:S5fmp} مطابق یا تاں
  $!A$، !!{derivable} اے یا کسے محدود ہمہ گیر نمونے وچ جھوٹھا اے۔
\end{proof}

اُپرلا ثبوت \Log{S5} لئی چلدا اے کیوں جے ہمہ گیر نمونیاں دیاں فلٹریشناں اپنے آپ ہمہ گیر ہوندیاں نیں۔ انعکاس تے ہر جہان توں رسائی لئی وی ایہو گل اے، پر ہور خاصیتاں لئی ہور کم چاہیدا اے۔
\endgroup
% /OLP-0457
% OLP-0458 content/normal-modal-logic/filtrations/more-filtrations.tex
\begingroup
\def\olfilename{more-filtrations}
\expandafter\def\csname ol@sectioncs\endcsname{section}

\olfileid{nml}{fil}{acc}

\olsection{فلٹریشناں تے رسائی دیاں خاصیتاں}

جیویں دَسیا، ہمہ گیر، ہر جہان توں رسائی والے، تے انعکاسی نمونیاں دیاں فلٹریشناں بالترتیب ہمیشہ ہمہ گیر، ہر جہان توں رسائی والیاں، یا انعکاسی وی ہوندیاں نیں۔ پر متقابل یا متعدی نمونے دی ہر فلٹریشن بالترتیب متقابل یا متعدی نہیں ہوندی۔ کجھ صورتاں وچ فلٹریشناں ایویں متعین ہو سکدیاں نیں کہ ایہ گل قائم رہے۔ ایس لئی اسیں سبھ توں موٹی فلٹریشن دی تعریف وانگ اگے ودھدے آں، پر~$R^*$ دی تعریف وچ ہور شرطاں شامل کردے آں۔ منّ لو $\Gamma$ ذیلی !!{formula}s تحت بند اے۔ تعلقاں~$C_i(u,v)$ نوں ویکھو، جیہڑے
\olref{tab:Cn-filtrations} وچ، جہاناں $u$، $v$ دے وچ، نمونے
$\mModel{M} =\tuple{W, R, V}$ وچ نیں۔ اسیں ایہناں شرطاں دے جوڑاں نال $R^*[u][v]$ متعین کر سکدے آں۔ مثلاً جے اسیں طے کریے کہ $R^*[u][v]$ اُدوں تے صرف اُدوں اے جدوں شرط $C_1(u,v)$ قائم اے، تاں عین سبھ توں موٹی فلٹریشن ملدی اے۔ جے طے کریے کہ $R^*[u][v]$ اُدوں تے صرف اُدوں اے جدوں $C_1(u,v)$ تے $C_2(u, v)$ دونیں قائم نیں، تاں وکھری فلٹریشن ملدی اے۔ ایہ سبھ توں موٹی فلٹریشن نالوں «باریک» اے، کیوں جے $C_1(u,v)$ تے $C_2(u,v)$ دوناں نوں پورا کرن والے جہاناں دے جوڑے، صرف $C_1(u,v)$ نوں پورا کرن والیاں توں گھٹ نیں۔

\begin{table}[ht]
  \centering
  \begin{tabular}{|ll|}
    \hline
    \multirow{2}{*}{$C_1(u,v)$:}
    \iftag{prvBox}{&
      جے $\Box!A \in \Gamma$ تے $\mSat{M}{\Box!A}[u]$
      تاں $\mSat{M}{!A}[v]$;\iftag{prvDiamond}{ تے}{}\\}{}
    \iftag{prvDiamond}{&
      جے $\Diamond!A \in \Gamma$  تے $\mSat{M}{!A}[v]$
      تاں $\mSat{M}{\Diamond!A}[u]$; \\}{}
    \hline
    \multirow{2}{*}{$C_2(u,v)$:} 
    \iftag{prvBox}{&
      جے $\Box!A \in \Gamma$ تے $\mSat{M}{\Box!A}[v]$
      تاں $\mSat{M}{!A}[u]$;\iftag{prvDiamond}{ تے}{}\\}{}
    \iftag{prvDiamond}{&
      جے $\Diamond!A \in \Gamma$ تے $\mSat{M}{!A}[u]$
      تاں $\mSat{M}{\Diamond!A}[v]$; \\}{}
    \hline
    \multirow{2}{*}{$C_3(u,v)$:}
    \iftag{prvBox}{&
      جے $\Box!A \in \Gamma$ تے $\mSat{M}{\Box!A}[u]$
      تاں $\mSat{M}{\Box!A}[v]$;\iftag{prvDiamond}{ تے}{}\\}{}
    \iftag{prvDiamond}{&
      جے $\Diamond!A \in \Gamma$ تے $\mSat{M}{\Diamond!A}[v]$
      تاں $\mSat{M}{\Diamond!A}[u]$; \\}{}
    \hline
    \multirow{2}{*}{$C_4(u,v)$:}
    \iftag{prvBox}{&
      جے $\Box!A \in \Gamma$ تے $\mSat{M}{\Box!A}[v]$
      تاں $\mSat{M}{\Box!A}[u]$;\iftag{prvDiamond}{ تے}{}\\}{}
    \iftag{prvDiamond}{&
      جے $\Diamond!A \in \Gamma$ تے $\mSat{M}{\Diamond!A}[u]$
      تاں $\mSat{M}{\Diamond!A}[v]$; \\}{}
    \hline
    \end{tabular}
  \caption{فلٹریشناں متعین کرن لئی ممکن جہاناں اُتے شرطاں۔}
  \ollabel{tab:Cn-filtrations}
\end{table}

\begin{thm}\ollabel{thm:more-filtrations}
  منّ لو $\mModel{M} =\tuple{W,R,V}$ اک نمونہ اے تے $\Gamma$ ذیلی !!{formula}s تحت بند اے۔ منّ لو $W^*$ تے $V^*$ اوویں متعین نیں جیویں
  \olref[fil]{defn:filtration} وچ۔ فیر
  \begin{enumerate}
  \item منّ لو $R^*[u][v]$ اُدوں تے صرف اُدوں جدوں $C_1(u, v) \land
    C_2(u,v)$۔ فیر $R^*$ متقابل اے، تے $\mModel{M^*} =
    \tuple{W^*,R^*,V^*}$ فلٹریشن اے جے $\mModel{M}$ متقابل اے۔
  \item منّ لو $R^*[u][v]$ اُدوں تے صرف اُدوں جدوں $C_1(u, v)
    \land C_3(u,v)$۔ فیر $R^*$ متعدی اے، تے
    $\mModel{M^*}=\tuple{W^*,R^*,V^*}$ فلٹریشن اے جے $\mModel{M}$
    متعدی اے۔
  \item منّ لو $R^*[u][v]$ اُدوں تے صرف اُدوں جدوں $C_1(u, v) \land C_2(u,v)
    \land C_3(u,v) \land C_4(u,v)$۔ فیر $R^*$ متقابل تے متعدی اے، تے
    $\mModel{M^*}=\tuple{W^*,R^*,V^*}$ فلٹریشن اے جے $\mModel{M}$ متقابل تے متعدی اے۔
  \item منّ لو $R^*$ ایویں متعین اے کہ $R^*[u][v]$ اُدوں تے صرف اُدوں جدوں $C_1(u,
    v) \land C_3(u,v) \land C_4(u,v)$۔ فیر $R^*$ متعدی تے اقلیدسی اے، تے
    $\mModel{M^*}=\tuple{W^*,R^*,V^*}$ فلٹریشن اے جے $\mModel{M}$ متعدی تے اقلیدسی اے۔
  \end{enumerate}
\end{thm}

\begin{proof}
  \begin{enumerate}
    \item صاف اے کہ $R^*$ متقابل اے، کیوں جے $C_1(u,v)
      \Leftrightarrow C_2(v,u)$ تے $C_2(u,v) \Leftrightarrow
      C_1(v,u)$۔ سو ایہ وکھاؤنا باقی اے کہ جے $\mModel{M}$ متقابل اے تاں $\mModel{M^*}$،~$\Gamma$ راہیں فلٹریشن اے۔ شرط $C_1(u,v)$ پکا کردی اے کہ
      \iftag{prvBox}{%
        \iftag{prvDiamond}{%
          \olref[fil]{defn:filtration-R2} تے
          \olref[fil]{defn:filtration-R3}، جیہڑیاں
          \olref[fil]{defn:filtration} دیاں نیں،}{%
          \olref[fil]{defn:filtration}\olref[fil]{defn:filtration-R2} دی شرط}}{%
        \olref[fil]{defn:filtration}\olref[fil]{defn:filtration-R3} دی شرط}
      پوری ہوندی اے۔ سو سانوں بس
      \olref[fil]{defn:filtration}\olref[fil]{defn:filtration-R1} پرکھنی اے،
      یعنی $Ruv$ توں $R^*[u][v]$ نکلدا اے۔

      سو منّ لو $Ruv$۔ $R^*[u][v]$ وکھاؤن لئی سانوں
      $C_1(u,v)$ تے $C_2(u,v)$ ثابت کرنے پین گے۔ $C_1$ لئی \iftag{prvBox}{جے $\Box!A
        \in\Gamma$ تے $\mSat{M}{\Box!A}[u]$، تاں
        $\mSat{M}{!A}[v]$ وی (کیوں جے $Ruv$)۔\iftag{prvDiamond}{ ایس طرح،
        }{}}{}\iftag{prvDiamond}{جے $\Diamond!A \in \Gamma$ تے
        $\mSat{M}{!A}[v]$، تاں $\mSat{M}{\Diamond!A}[u]$ کیوں جے
        $Ruv$۔}{} $C_2$ لئی \iftag{prvBox}{جے $\Box!A \in \Gamma$ تے
        $\mSat{M}{\Box!A}[v]$، تاں $Ruv$ توں تقابل کارن $Rvu$ نکلدا اے، سو
        $\mSat{M}{!A}[u]$۔\iftag{prvDiamond}{ ایس طرح،
        }{}}{}\iftag{prvDiamond}{جے $\Diamond!A \in\Gamma$ تے
        $\mSat{M}{!A}[u]$، تاں $\mSat{M}{\Diamond!A}[v]$ (کیوں جے تقابل کارن $Rvu$)۔}{}
    \item مشق۔
    \item مشق۔
    \item مشق۔
  \end{enumerate}
\end{proof}

\begin{prob}
  \olref[nml][fil][acc]{thm:more-filtrations} دا ثبوت پورا کرو۔
\end{prob}
\endgroup
% /OLP-0458
% OLP-0459 content/normal-modal-logic/filtrations/euclidean-filtrations.tex
\begingroup
\def\olfilename{euclidean-filtrations}
\expandafter\def\csname ol@sectioncs\endcsname{section}

\olfileid{nml}{fil}{euc}

\olsection{اقلیدسی نمونیاں دیاں فلٹریشناں}

\olref[acc]{sec} دا طریقہ اوہناں نمونیاں لئی نہیں چلدا جیہڑے اقلیدسی، یا ہر جہان توں رسائی والے تے اقلیدسی، نیں۔ \olref{fig:ser-eucl} دے اُتلے نمونے نوں ویکھو، جیہڑا اقلیدسی تے ہر جہان توں رسائی والا دونیں اے۔ منّ لو
$\Gamma = \{p, \Box p \}$۔ $\Gamma$ راہیں فلٹریشن لیندیاں
$[w_1] = [w_3]$، کیوں جے $w_1$ تے $w_3$ ہی اوہ جہان نیں جیہڑے~$\Gamma$ اُتے اکّو فیصلے دیندے نیں۔ ہر فلٹریشن وچ $\mModel{M}$ توں ملے تیر وی ہون گے، جیویں \olref{fig:ser-eucl2} وچ وکھایا اے۔ اوہ نمونہ اقلیدسی نہیں۔ نالے اسیں اوس نمونے وچ تیر ودھا کے اینوں اقلیدسی نہیں بنا سکدے۔ سانوں
$[w_2]$ تے $[w_4]$ دے وچ دو پاسیاں والے تیر پاؤنے پین گے، تے فیر $w_2$ تے~$w_5$ دے وچ وی۔ پر $\Box
p$،~$w_2$ اُتے سچا ہونا چاہیدا اے، جدکہ $p$،~$w_5$ اُتے جھوٹھا اے۔

\begin{figure}[htpb]
  \centering
  \begin{tikzpicture}[modal]
    \node[world] (w1)
          [label={left: \mFalse{p}},
            label={below: $\mSat{{}}{\Box p}$}] {$w_1$}; 
    \node[world] (w2)
          [label={right: \mTrue{p}},
            label = {below: $\mSat{{}}{\Box p}$},
            right=of w1] {$w_2$}; 
    \draw[->] (w1) to (w2);
    \node[world] (w3)
          [label={left: \mFalse{p}},
            label={below: $\mSat{{}}{\Box p}$},
            below=of w1] {$w_3$}; 
    \node[world] (w4)
          [label={right:\mTrue{p}},
            label={below: $\mSat/{{}}{\Box p}$},
            right=of w3] {$w_4$};
    \draw[->] (w3) to (w4);
    \draw[reflexive above] (w4) to (w4);     
    \node[world] (w5)
          [label={right: \mFalse{p}},
            label=below:$\mSat/{{}}{\Box p}$,
            right=of w4] {$w_5$}; 
    \draw[->, bend left] (w4) to (w5);
    \draw[->, bend left] (w5) to (w4);
    \draw[reflexive above] (w5) to (w5);     
  \end{tikzpicture}
  \caption{ہر جہان توں رسائی والا تے اقلیدسی نمونہ۔}
  \ollabel{fig:ser-eucl}
\end{figure}

\begin{figure}[ht]
  \centering
  \begin{tikzpicture}[modal,every text node part/.style={align=center}]
    \node[world] (w1)
          [label={left: \mFalse{p}},
            label={right:$[w_1]=[w_3]$},
            label={below: $\mSat{{}}{\Box p}$}] {$[w_1]$}; 
    \node[world] (w2)
         [label={right: \mTrue{p}},
           label = {below: $\mSat{{}}{\Box p}$},
           right=of w1, yshift=1.5cm] {$[w_2]$}; 
    \draw[->] (w1) to (w2);
    \node[world] (w4)
         [label={right:\mTrue{p}},
           label={below: $\mSat/{{}}{\Box p}$},
           right=of w1,yshift=-1.5cm] {$[w_4]$};
    \draw[->] (w1) to (w4);
    \draw[reflexive above] (w4) to (w4);     
    \node[world] (w5)
         [label={right: \mFalse{p}},
             label=below:$\mSat/{{}}{\Box p}$,
             right=of w4] {$[w_5]$}; 
    \draw[->, bend left] (w4) to (w5);
    \draw[->, bend left] (w5) to (w4);
    \draw[reflexive above] (w5) to (w5);     
  \end{tikzpicture}
  \caption{\olref{fig:ser-eucl} وچلے نمونے دی فلٹریشن۔}
  \ollabel{fig:ser-eucl2}
\end{figure}

خاص طور تے اقلیدسی فلٹریشن لبھن لئی اوہناں من مرضی دے $\Gamma$ راہیں فلٹریشناں ویکھنیاں کافی نہیں جیہڑے ذیلی !!{formula}s تحت بند نیں۔ سانوں ایسے سیٹ $\Gamma$ ویکھنے پین گے جیہڑے \emph{امکانی عاملاں تحت بند} نیں (\olref[pre]{defn:modallyclosed} ویکھو)۔ جملیاں دے ایسے سیٹ لامتناہی ہوندے نیں، سو ایہ متعلقہ نظام لئی محدود نمونے دی خاصیت یا فیصلہ کرن دی اہلیت سدھی نہیں دیندے۔

\begin{thm}\ollabel{thm:modal-closed-filt}
  منّ لو $\Gamma$ امکانی عاملاں تحت بند اے، $\mModel{M}=\tuple{W,R,V}$ اے، تے
  $\mModel{M^*} = \tuple{W^*,R^*,V^*}$، $\mModel{M}$ دی سبھ توں موٹی فلٹریشن اے۔
  \begin{enumerate}
  \item جے $\mModel{M}$ متقابل اے، تاں $\mModel{M^*}$ وی اے۔
  \item جے $\mModel{M}$ متعدی اے، تاں $\mModel{M^*}$ وی اے۔
  \item جے $\mModel{M}$ اقلیدسی اے، تاں $\mModel{M^*}$ وی اے۔
  \end{enumerate}
\end{thm}

\begin{proof}
\begin{enumerate}
  \item جے $\mModel{M^*}$ سبھ توں موٹی فلٹریشن اے، تاں تعریف مطابق
    $R^*[u][v]$ اُدوں تے صرف اُدوں قائم اے جدوں $C_1(u,v)$۔ تعدّی لئی منّ لو $C_1(u,v)$ تے $C_1(v,w)$۔ سانوں
    $C_1(u,w)$ وکھاؤنا اے۔ \iftag{prvBox}{منّ لو $\mSat{M}{\Box !A}[u]$۔ فیر
      $\mSat{M}{\Box\Box!A}[u]$، کیوں جے \Ax{4} سارے متعدی نمونیاں وچ معتبر اے۔ کیوں جے بندش کارن $\Box\Box!A \in \Gamma$، تے
      $C_1(u,v)$ مطابق، $\mSat{M}{\Box!A}[v]$، تے $C_1(v,w)$ مطابق
      $\mSat{M}{!A}[w]$ وی۔ }{}\iftag{prvDiamond}{منّ لو
      $\mSat{M}{!A}[w]$۔ فیر $\mSat{M}{\Diamond !A}[v]$،
      $C_1(v,w)$ مطابق، کیوں جے امکانی بندش کارن $\Diamond !A \in \Gamma$۔
      $C_1(u,v)$ توں سانوں $\mSat{M}{\Diamond\Diamond !A}[u]$ ملدا اے، کیوں جے امکانی بندش کارن
      $\Diamond\Diamond!A \in \Gamma$۔ کیوں جے
      $\Ax{4}_\Diamond$ سارے متعدی نمونیاں وچ معتبر اے،
      $\mSat{M}{\Diamond!A}[u]$۔}{}
  \item مشق۔ ایہ گل ورتو کہ \Ax{5} تے $\Ax{5_\Diamond}$ دونیں سارے اقلیدسی نمونیاں وچ معتبر نیں۔
  \item مشق۔ ایہ گل ورتو کہ \Ax{B} تے $\Ax{B_\Diamond}$ سارے متقابل نمونیاں وچ معتبر نیں۔
\end{enumerate}
\end{proof}

\begin{prob}
  \olref[nml][fil][euc]{thm:modal-closed-filt} دا ثبوت پورا کرو۔
\end{prob}
\endgroup
% /OLP-0459

\OLEndChapterHook
\endgroup
% /OLP-0450

% OLP-0460 content/normal-modal-logic/tableaux/tableaux.tex
\begingroup
\def\olfilename{tableaux}
\expandafter\def\csname ol@sectioncs\endcsname{section}

\olchapter{nml}{tab}{امکانی \usetoken{P}{tableau}}

\begin{editorial}
  امکانی منطق دے اگیتے والے تابلوآں بارے کچا باب۔ ہور مثالاں، تمامیت دے ثبوت، تے بند تابلو دی ناکام تلاش توں ردّی نمونے لبھن دا طریقہ سمجھاؤن دی لوڑ اے۔
\end{editorial}

% OLP-0461 content/normal-modal-logic/tableaux/introduction.tex
\begingroup
\def\olfilename{introduction}
\expandafter\def\csname ol@sectioncs\endcsname{section}

\olfileid{nml}{tab}{int}

\olsection{جان پہچان}

!!^{tableau}s، !!{signed formula}s دے خاص (ہِٹھاں ول شاخاں کڈھن والے) درخت نیں، یعنی اوہناں جوڑیاں دے درخت جیہناں وچ صداقتی قدر دا نشان ($\True$ یا
$\False$) تے اک !!a{sentence} ہوندے نیں۔
\[
\sFmla{\True}{!A} \text{ یا } \sFmla{\False}{!A}.
\]
!!^a{tableau} کجھ \emph{مفروضیاں} نال شروع ہوندا اے۔ ہر اگلا
!!{signed formula} کسے استنباطی قاعدے نوں لا کے پیدا ہوندا اے۔ کجھ قاعدے درخت دے اک سِرے اُتے اک یا ودھ !!{signed formula}s لاؤندے نیں۔ ہور دو نویں سِرے لاؤندے نیں، تے ایس نال دو شاخاں بندیاں نیں۔
قاعدیاں توں ایسے !!{signed formula}s ملدے نیں جیہناں دا !!{formula} اوس !!{signed formula} دے فارمولے توں گھٹ پیچیدہ ہوندا اے جیہدے اُتے قاعدہ لگیا سی۔ جدوں کسے شاخ وچ $\sFmla{\True}{!A}$ تے
$\sFmla{\False}{!A}$ دونیں ہون، تاں اوہ شاخ \emph{بند} اکھواؤندی اے۔ جے !!a{tableau} دی ہر شاخ بند اے، تاں سارا !!{tableau} بند اے۔ بند !!{tableau} اک !!a{derivation} اے جیہڑا وکھاؤندا اے کہ !!{tableau} شروع کرن لئی ورتے گئے !!{signed formula}s دا سیٹ اکٹھا پورا نہیں ہو سکدا۔ ایس نال $\Proves$ دا تعلق متعین کر سکدے آں۔
$\Gamma \Proves !A$ اُدوں تے صرف اُدوں جدوں کوئی ایسا محدود سیٹ~$\Gamma_0 = \{!B_1,
\dots, !B_n\} \subseteq \Gamma$ ہووے جیہدے لئی ہِیٹھلے مفروضیاں دا اک بند
!!{tableau} ہووے۔
\[
\{\sFmla{\False}{!A}, \sFmla{\True}{!B_1}, \dots, \sFmla{\True}{!B_n}\}.
\]

امکانی منطقاں لئی سانوں !!{signed formula} دا خیال ودھاؤنا وی پئے گا تے اوہ قاعدے شامل کرنے پین گے جیہڑے~\iftag{prvBox}{$\Box$\iftag{prvDiamond}{ تے
    $\Diamond$}{}}{$\Diamond$} لئی نیں۔ نشان ($\True$ یا
$\False$) توں علاوہ، !!{formula}s نال امکانی !!{tableau}s وچ \emph{اگیتے}~$\sigma$ وی ہوندے نیں۔ اگیتے مثبت صحیح عدداں دیاں غیر خالی لڑیاں نیں، یعنی $\sigma \in (\PosInt)^* \setminus
\{\emptyseq\}$۔ ایسے اگیتے اسیں ارد گرد دے
$\tuple{\ }$ توں بغیر لکھدے آں، تے وکھ وکھ !!{element}s نوں~$.$ نال وکھ کردے آں، $,$ دی تھاں۔ جے $\sigma$ اگیتا اے، تاں $\sigma.n$ دا مطلب $\sigma
\concat \tuple{n}$ اے۔ مثلاً جے $\sigma = 1.2.1$، تاں $\sigma.3$، $1.2.1.3$ اے۔ سو مثلاً
\[
\sFmla{\True}{\Box !A \lif !A}[1.2]
\]
اک \emph{اگیتے والا !!{signed formula}} اے (مختصر طور تے \emph{اگیتے والا
  !!{formula}})۔

سیدھی سمجھ ایہ اے کہ اگیتا نمونے دے ایسے جہان دا ناں اے جیہڑا اوہ !!{formula}s پورے کر سکدا ہووے جیہڑے !!a{tableau} دی اک شاخ اُتے آئے نیں۔ تے جے $\sigma$ کسے جہان دا ناں اے، تاں $\sigma.n$ ایسے جہان دا ناں اے جیہڑا~$\sigma$ دے ناں والے جہان توں رسائی وچ اے۔
\endgroup
% /OLP-0461
% OLP-0462 content/normal-modal-logic/tableaux/rules-for-K.tex
\begingroup
\def\olfilename{rules-for-K}
\expandafter\def\csname ol@sectioncs\endcsname{section}

\olfileid{nml}{tab}{rul}

\olsection{\texorpdfstring{\Ax{K} دے قاعدے}{K دے قاعدے}}

عام قضیاتی رابطیاں لئی قاعدے اوہی نیں جیہڑے عام قضیاتی نشان والے
!!{tableau}s لئی نیں، بس اگیتے شامل ہو گئے نیں۔ ہر صورت وچ، نشان والے !!{formula}
$\sFmla{S}{!A}[\sigma]$ اُتے قاعدہ لاؤن نال نویں !!{formula}s ملدے نیں
جیہناں دا اگیتا وی~$\sigma$ اے۔ ایہ گل سِدھی اے۔ مثلاً جے $!A
\land !B$،~$\sigma$ دے ناں والے جہان وچ سچ اے، تاں $!A$ تے $!B$
اوسے~$\sigma$ وچ سچ نیں (کسی ہور جہان وچ نہیں)۔ قضیاتی قاعدے اسیں
\olref{tab:prop-rules} وچ اکٹھے دتے نیں۔

\begin{table}
  \[\def\arraystretch{3}\begin{array}{|c|c|}
    \hline
    \AxiomC{\sFmla{\True}{\lnot !A}[\sigma]}
    \RightLabel{\TRule{\True}{\lnot}}
    \UnaryInfC{\sFmla{\False}{!A}[\sigma]}
    \DisplayProof
    &
    \AxiomC{\sFmla{\False}{\lnot !A}[\sigma]}
    \RightLabel{\TRule{\False}{\lnot}}
    \UnaryInfC{\sFmla{\True}{!A}[\sigma]}
    \DisplayProof
    \\[1ex]
    \hline
    \AxiomC{\sFmla{\True}{!A \land !B}[\sigma]}
    \RightLabel{\TRule{\True}{\land}}
    \UnaryInfC{\sFmla{\True}{!A}[\sigma]}
    \noLine
    \UnaryInfC{\sFmla{\True}{!B}[\sigma]}
    \DisplayProof
    &
    \AxiomC{\sFmla{\False}{!A \land !B}[\sigma]}
    \RightLabel{\TRule{\False}{\land}}
    \UnaryInfC{$\sFmla{\False}{!A}[\sigma] \quad \mid \quad
      \sFmla{\False}{!B}[\sigma]$}
    \DisplayProof
    \\[2ex]
    \hline
    \AxiomC{\sFmla{\True}{!A \lor !B}[\sigma]}
    \RightLabel{\TRule{\True}{\lor}}
    \UnaryInfC{$\sFmla{\True}{!A}[\sigma] \quad \mid \quad
      \sFmla{\True}{!B}[\sigma]$}
    \DisplayProof
    &
    \AxiomC{\sFmla{\False}{!A \lor !B}[\sigma]}
    \RightLabel{\TRule{\False}{\lor}}
    \UnaryInfC{\sFmla{\False}{!A}[\sigma]}
    \noLine
    \UnaryInfC{\sFmla{\False}{!B}[\sigma]}
    \DisplayProof
    \\[2ex]
    \hline
    \AxiomC{\sFmla{\True}{!A \lif !B}[\sigma]}
    \RightLabel{\TRule{\True}{\lif}}
    \UnaryInfC{$\sFmla{\False}{!A}[\sigma] \quad \mid
      \quad \sFmla{\True}{!B}[\sigma]$}
    \DisplayProof
    &
    \AxiomC{\sFmla{\False}{!A \lif !B}[\sigma]}
    \RightLabel{\TRule{\False}{\lif}}
    \UnaryInfC{\sFmla{\True}{!A}[\sigma]}
    \noLine
    \UnaryInfC{\sFmla{\False}{!B}[\sigma]}
    \DisplayProof
    \\[2ex]
    \hline
  \end{array}\]
  \caption{قضیاتی رابطیاں لئی اگیتے والے !!{tableau} دے قاعدے}
  \ollabel{tab:prop-rules}
\end{table}

شاخ بند ہون دی شرط اوہی اے جیہڑی عام !!{tableau}s لئی اے،
پر !!{formula}s دے نال نال اگیتے وی اکو جیہے ہونے چاہیدے نیں۔
سو شاخ بند اے جے اوہدے وچ ایہ دونیں ہون:
\[
\sFmla{\True}{!A}[\sigma] \quad\text{تے}\quad \sFmla{\False}{!A}[\sigma]
\]
ایتھے $\sigma$ کوئی اگیتا اے تے~$!A$ کوئی !!{formula} اے۔

مفروضے رکھن دے قاعدے وی عام !!{tableau}s والے ای نیں، بس
مفروضیاں لئی اسیں ہمیشہ اگیتا~$1$ ورتدے آں۔ (اگیتا کوئی وی ہو سکدا اے،
جے سارے مفروضیاں دا اگیتا اکّو ہووے۔) سو مثلاً اسیں آکھدے آں کہ
\[
!B_1, \dots, !B_n \Proves !A
\]
اُدوں تے صرف اُدوں جدوں ہِیٹھلے مفروضیاں دا اک بند تابلو ہووے:
\[
\sFmla{\True}{!B_1}[1], \dots, \sFmla{\True}{!B_n}[1],
\sFmla{\False}{!A}[1].
\]

امکانی عمل کار\iftag{prvBox}{\iftag{prvDiamond}{~$\Box$
    تے}{~$\Box$}}{}\iftag{prvDiamond}{~$\Diamond$}{} لئی،
اگیتے~$\sigma$ والے !!a{formula} اُتے قاعدہ لاؤن نال جیہڑا نتیجہ
ملدا اے، اوہدا اگیتا $\sigma.n$ اے۔ پر کہڑا $n$ ورت سکدے آں،
ایہ ایس گل اُتے اے کہ نشان~$\True$ اے یا~$\False$۔

\iftag{prvBox}{قاعدہ $\TRule{\True}{\Box}$ اوہ شاخ ودھاؤندا اے
  جیہدے وچ $\sFmla{\True}{\Box !A}[\sigma]$ ہووے۔ ایہ اوہدے اُتے
  $\sFmla{\True}{!A}[\sigma.n]$ لاؤندا اے۔\iftag{prvDiamond}{ ایس طرح، }{}}{}\iftag{prvDiamond}{قاعدہ $\TRule{\False}{\Diamond}$ اوہ شاخ
  ودھاؤندا اے جیہدے وچ $\sFmla{\False}{\Diamond !A}[\sigma]$ ہووے۔
  ایہ اوہدے اُتے $\sFmla{\False}{!A}[\sigma.n]$ لاؤندا اے۔}{}
\iftag{notprvBox,notprvDiamond}{ایہ قاعدہ}{ایہ قاعدے} اوہدوں ای
لائے جا سکدے نیں جدوں اگیتا~$\sigma.n$ اوس شاخ اُتے \emph{پہلاں ای}
موجود ہووے جیہدے اُتے قاعدہ لگدا اے۔ شاخ لئی ایسے اگیتے نوں
اسیں «ورتیا ہویا» آکھدے آں۔

\iftag{prvBox}{قاعدہ $\TRule{\False}{\Box}$ اوہ شاخ ودھاؤندا اے
  جیہدے وچ $\sFmla{\False}{\Box !A}[\sigma]$ ہووے۔ ایہ اوہدے اُتے
  $\sFmla{\False}{!A}[\sigma.n]$ لاؤندا اے۔\iftag{prvDiamond}{ ایس طرح، }{}}{}\iftag{prvDiamond}{قاعدہ $\TRule{\True}{\Diamond}$ اوہ شاخ
  ودھاؤندا اے جیہدے وچ $\sFmla{\True}{\Diamond !A}[\sigma]$ ہووے۔
  ایہ اوہدے اُتے $\sFmla{\True}{!A}[\sigma.n]$ لاؤندا اے۔}{}
پر \iftag{notprvBox,notprvDiamond}{ایہ قاعدہ}{ایہ قاعدے} اوہدوں ای
لائے جا سکدے نیں جدوں اگیتا~$\sigma.n$ اوس شاخ اُتے پہلاں توں
\emph{موجود نہ} ہووے جیہدے اُتے قاعدہ لگدا اے۔ شاخ لئی ایسے اگیتے
نوں اسیں «نواں» آکھدے آں۔

قاعدے \olref{tab:rules-K} وچ دتے نیں۔

\begin{table}
  \begin{center}
    \def\arraystretch{3}\def\fCenter{}
    \begin{tabular}{|c|c|}
    \hline
    \iftag{prvBox}{
      \AxiomC{\sFmla{\True}{\Box !A}[\sigma]}
      \RightLabel{\TRule{\True}{\Box}}
      \UnaryInfC{\sFmla{\True}{!A}[\sigma.n]}
      \DisplayProof
      &
      \AxiomC{\sFmla{\False}{\Box !A}[\sigma]}
      \RightLabel{\TRule{\False}{\Box}}
      \UnaryInfC{\sFmla{\False}{!A}[\sigma.n]}
      \DisplayProof\\
      $\sigma.n$ ورتیا ہویا اے & $\sigma.n$ نواں اے
      \\[1ex]
      \hline}{}
    \iftag{prvDiamond}{
      \AxiomC{\sFmla{\True}{\Diamond !A}[\sigma]}
      \RightLabel{\TRule{\True}{\Diamond}}
      \UnaryInfC{\sFmla{\True}{!A}[\sigma.n]}
      \DisplayProof
      &
      \AxiomC{\sFmla{\False}{\Diamond !A}[\sigma]}
      \RightLabel{\TRule{\False}{\Diamond}}
      \UnaryInfC{\sFmla{\False}{!A}[\sigma.n]}
      \DisplayProof\\
      $\sigma.n$ نواں اے & $\sigma.n$ ورتیا ہویا اے
      \\[1ex]
      \hline}{}
    \end{tabular}
  \end{center}
  \caption{\Ax{K} دے امکانی قاعدے۔}
  \ollabel{tab:rules-K}
\end{table}

قاعدے \iftag{prvBox}{\TRule{\True}{\Box}}{\TRule{\False}{\Diamond}}
لئی شرط اے کہ اگیتا ورتیا ہویا ہووے۔ ایہ شرط ضروری اے، کیوں جے ایہدے
بغیر ہِیٹھلے !!{tableau} نوں بند گنیا جاندا:
\iftag{notprvBox,notprvDiamond}{%
  \iftag{prvBox}{%
  \begin{oltableau}
    [\pFmla{\True}{\Box \formula{A}}{1}, just = \TAss
      [\pFmla{\False}{\lnot\Box\lnot \formula{A}}{1}, just = \TAss
        [\pFmla{\True}{\formula{A}}{1.1}, just = {\TRule{\True}{\Box}[1]}
          [\pFmla{\True}{\Box\lnot\formula{A}}{1},
            just ={\TRule{\False}{\lnot}[2]}
            [\pFmla{\True}{\lnot\formula{A}}{1.1},
              just = {\TRule{\True}{\Box}[4]}
              [\pFmla{\False}{\formula{A}}{1.1},
                  just ={\TRule{\True}{\lnot}[5]}, close]
            ]
          ]
        ]
      ]
    ]
  \end{oltableau}
  }{
  \begin{oltableau}
    [\pFmla{\True}{\lnot\Diamond\lnot \formula{A}}{1}, just = \TAss
      [\pFmla{\False}{\Diamond\formula{A}}{1}, just = \TAss
        [\pFmla{\False}{\formula{A}}{1.1}, just = {\TRule{\False}{\Diamond}[2]}
          [\pFmla{\False}{\Diamond\lnot\formula{A}}{1},
            just ={\TRule{\True}{\lnot}[1]}
            [\pFmla{\False}{\lnot\formula{A}}{1.1},
              just = {\TRule{\False}{\Diamond}[4]}
              [\pFmla{\True}{\formula{A}}{1.1},
                  just ={\TRule{\True}{\lnot}[5]}, close]
            ]
          ]
        ]
      ]
    ]
  \end{oltableau}
}}{
  \begin{oltableau}
    [\pFmla{\True}{\Box \formula{A}}{1}, just = \TAss
      [\pFmla{\False}{\Diamond \formula{A}}{1}, just = \TAss
        [\pFmla{\True}{\formula{A}}{1.1}, just = {\TRule{\True}{\Box}[1]}
          [\pFmla{\False}{\formula{A}}{1.1},
            just = {\TRule{\False}{\Diamond}[2]}, close]
        ]
      ]
    ]
\end{oltableau}}

پر $\Box \formula{A} \Entails/ \Diamond \formula{A}$، سو ساڈا
اخذ دا نظام صحیح نہ رہندا۔ ایس طرح $\Diamond \formula{A} \Entails/
\Box \formula{A}$، پر جے قاعدے
\iftag{prvBox}{\TRule{\False}{\Box}}{\TRule{\True}{\Diamond}}
لئی اگیتا نواں ہون دی شرط نہ ہووے، تاں ایہ بند تابلو ہووے گا: \iftag{notprvBox,notprvDiamond}{%
  \iftag{prvBox}{%
  \begin{oltableau}
    [\pFmla{\True}{\lnot\Box\lnot \formula{A}}{1}, just = \TAss
      [\pFmla{\False}{\Box \formula{A}}{1}, just = \TAss
        [\pFmla{\False}{\formula{A}}{1.1}, just = {\TRule{\True}{\Box}[2]}
          [\pFmla{\False}{\Box\lnot\formula{A}}{1},
            just ={\TRule{\True}{\lnot}[1]}
            [\pFmla{\False}{\lnot\formula{A}}{1.1},
              just = {\TRule{\False}{\Box}[4]}
              [\pFmla{\True}{\formula{A}}{1.1},
                  just ={\TRule{\False}{\lnot}[5]}, close]
            ]
          ]
        ]
      ]
    ]
  \end{oltableau}
  }{
  \begin{oltableau}
    [\pFmla{\True}{\Diamond \formula{A}}{1}, just = \TAss
      [\pFmla{\False}{\lnot\Diamond\lnot\formula{A}}{1}, just = \TAss
        [\pFmla{\True}{\formula{A}}{1.1}, just = {\TRule{\True}{\Diamond}[1]}
          [\pFmla{\True}{\Diamond\lnot\formula{A}}{1},
            just ={\TRule{\False}{\lnot}[2]}
            [\pFmla{\True}{\lnot\formula{A}}{1.1},
              just = {\TRule{\True}{\Diamond}[4]}
              [\pFmla{\False}{\formula{A}}{1.1},
                  just ={\TRule{\True}{\lnot}[5]}, close]
            ]
          ]
        ]
      ]
    ]
  \end{oltableau}
}}{
  \begin{oltableau}
    [\pFmla{\True}{\Diamond \formula{A}}{1}, just = \TAss,name=one
      [\pFmla{\False}{\Box \formula{A}}{1}, just = \TAss, name=two
        [\pFmla{\True}{\formula{A}}{1.1}, just = {\TRule{\True}{\Diamond}[1]}
          [\pFmla{\False}{\formula{A}}{1.1},
            just = {\TRule{\False}{\Box}[2]}, close]
        ]
      ]
    ]
  \end{oltableau}}
\endgroup
% /OLP-0462
% OLP-0463 content/normal-modal-logic/tableaux/proofs-in-K.tex
\begingroup
\def\olfilename{proofs-in-K}
\expandafter\def\csname ol@sectioncs\endcsname{section}

\olfileid{nml}{tab}{prk}

\olsection{\Log{K} لئی \usetoken{P}{tableau}}

\iftag{prvBox}{

\begin{ex}
  اسیں اک بند تابلو دیندے آں جیہڑا وکھاؤندا اے کہ $\Proves (\Box!A \land \Box!B)
  \lif \Box (!A \land !B)$.
  \begin{oltableau}
    [\pFmla{\False}{(\Box\formula{A} \land \Box\formula{B}) \lif
        \Box (\formula{A} \land \formula{B})}{1},
      just =\TAss
      [\pFmla{\True}{\Box\formula{A} \land \Box\formula{B}}{1},
        just = {\TRule{\False}{\lif}[1]}
        [\pFmla{\False}{\Box(\formula{A} \land \formula{B})}{1},
          just = {\TRule{\False}{\lif}[1]}
          [\pFmla{\True}{\Box\formula{A}}{1},
            just = {\TRule{\True}{\land}[2]}
            [\pFmla{\True}{\Box\formula{B}}{1},
              just = {\TRule{\True}{\land}[2]}
              [\pFmla{\False}{\formula{A} \land \formula{B}}{1.1},
                just = {\TRule{\False}{\Box}[3]}
                [\pFmla{\False}{\formula{A}}{1.1},
                  just = {\TRule{\False}{\land}[6]}
                  [\pFmla{\True}{\formula{A}}{1.1},
                    just= {\TRule{\True}{\Box}[4]}, close]]
                [\pFmla{\False}{\formula{B}}{1.1},
                  just = {\TRule{\False}{\land}[6]}
                  [\pFmla{\True}{\formula{B}}{1.1},
                    just= {\TRule{\True}{\Box}[5]}, close]]
              ]
            ]
          ]
        ]
      ]
    ]
  \end{oltableau}
\end{ex}
}{}

\iftag{prvDiamond}{
\begin{ex}
  اسیں اک بند تابلو دیندے آں جیہڑا وکھاؤندا اے کہ $\Proves \Diamond(!A \lor !B)
  \lif (\Diamond !A \lor \Diamond !B)$:
  \begin{oltableau}
    [\pFmla{\False}{\Diamond(\formula{A} \lor \formula{B}) \lif
        (\Diamond \formula{A} \lor \Diamond \formula{B})}{1},
      just =\TAss
      [\pFmla{\True}{\Diamond(\formula{A} \lor \formula{B})}{1},
        just = {\TRule{\False}{\lif}[1]}
        [\pFmla{\False}{\Diamond\formula{A} \lor \Diamond\formula{B}}{1},
          just = {\TRule{\False}{\lif}[1]}
          [\pFmla{\False}{\Diamond\formula{A}}{1},
            just = {\TRule{\False}{\lor}[3]}
            [\pFmla{\False}{\Diamond\formula{B}}{1},
              just = {\TRule{\False}{\lor}[3]}
              [\pFmla{\True}{\formula{A} \lor \formula{B}}{1.1},
                just = {\TRule{\True}{\Diamond}[2]},
                [\pFmla{\True}{\formula{A}}{1.1},
                  just = {\TRule{\True}{\lor}[6]}
                  [\pFmla{\False}{\formula{A}}{1.1},
                    just= {\TRule{\False}{\Diamond}[4]}, close]]
                [\pFmla{\True}{\formula{B}}{1.1},
                  just = {\TRule{\True}{\lor}[6]}
                  [\pFmla{\False}{\formula{B}}{1.1},
                    just= {\TRule{\False}{\Diamond}[5]}, close]]
              ]
            ]
          ]
        ]
      ]
    ]
  \end{oltableau}
\end{ex}
}{}

\begin{prob}
  بند !!{tableau}s لبھو،~$\Log{K}$ وچ، ہِیٹھلے !!{formula}s لئی:
  \begin{enumerate}
    \item $\Box \lnot p \lif \Box(p \lif q)$
    \item $(\Box p \lor \Box q) \lif \Box(p \lor q)$
    \item $\Diamond p \lif \Diamond(p \lor q)$
    \item $\Box(p \land q) \lif \Box p$
  \end{enumerate}
\end{prob}
\endgroup
% /OLP-0463
% OLP-0464 content/normal-modal-logic/tableaux/soundness.tex
\begingroup
\def\olfilename{soundness}
\expandafter\def\csname ol@sectioncs\endcsname{section}

\olfileid{nml}{tab}{sou}

\olsection{\Log{K} لئی صحت مندی}

\begin{editorial}
  صحت مندی دا ایہ ثبوت کلاسیکی قضیاتی منطق والا ثبوت دوبارہ ورتدا اے،
  یعنی ہر گل مڈھ توں ثابت کردا اے۔ جے صحت مندی دا اپنے آپ وچ پورا
  ثبوت چاہیدا اے تاں ایہ ٹھیک اے۔ جے تسی عام تابلو دی صحت مندی پہلاں
  ویکھ چکے او، تاں ایہ کجھ دہرایا ہویا لگے گا۔ اگے چل کے ایس اپنے آپ
  وچ پورے روپ تے غیر امکانی صورت اُتے ٹکے روپ وچوں چُن سکنا ممکن
  بناؤن دی وِونت اے۔
\end{editorial}

\begin{explain}
  اگیتے والے !!{tableau}s دی صحت مندی وکھاؤن لئی اسیں ثابت کرنا اے کہ جے
  \[
  \sFmla{\True}{!B_1}[1], \dots, \sFmla{\True}{!B_n}[1], \sFmla{\False}{!A}[1]
  \]
  دا اک بند !!{tableau} اے، تاں $!B_1, \dots, !B_n \Entails !A$۔
  اُلٹے نکار والا بیان ثابت کرنا سوکھا اے۔ جے کسے $\mModel{M}$ تے
  جہان~$w$ لئی، $\mSat{M}{!B_i}[w]$ سارے $i=1$, \dots,~$n$ لئی ہووے
  پر $\mSat{M}{!A}[w]$ ہووے، تاں کوئی !!{tableau} بند نہیں ہو سکدا۔
  اجیہا ردّی نمونہ وکھاؤندا اے کہ !!{tableau} دے شروع والے مفروضے پورے
  ہون جوگ نیں۔ ثبوت دا طریقہ ایہ اے کہ جدوں !!a{tableau} دی کسے شاخ
  اُتے سارے اگیتے والے !!{formula}s پورے ہون جوگ ہون، تاں قاعدہ لاؤن نال
  گھٹ توں گھٹ اک ودھی ہوئی شاخ وی پورا ہون جوگ ہوندی اے۔ بند شاخاں
  پوریاں ہون جوگ نہیں، ایس لئی اگیتے والے !!{formula}s دے کسے پورا ہون
  جوگ سیٹ دے ہر !!{tableau} وچ گھٹ توں گھٹ اک کھلی شاخ ہووے گی۔

امکانی صورت وچ ایہ طریقہ ورتن لئی سانوں «پورا ہون جوگ» دی تعریف
  امکانی نمونیاں تے اگیتیاں تک ودھاؤنی پئے گی۔ اوہ تعریف ملن پچھوں
  ثبوت سِدھا اے۔
\end{explain}

\begin{defn}
  منّ لو $P$ اگیتیاں دا کوئی سیٹ اے، یعنی $P \subseteq (\PosInt)^*
  \setminus \{\emptyseq\}$، تے $\mModel{M}$ اک نمونہ اے۔
  فنکشن~$f\colon P \to W$،~$P$ دی $\mModel{M}$ وچ \emph{تعبیر} اے جے
  جدوں وی $\sigma$ تے $\sigma.n$ دونیں~$P$ وچ ہون، تاں
  $Rf(\sigma)f(\sigma.n)$ ہووے۔

اگیتیاں دے سیٹ~$P$ دی تعبیر دے لحاظ نال اسیں ایہ تعریف کر سکدے آں:
  \begin{enumerate}
  \item $\mModel{M}$، $\sFmla{\True}{!A}[\sigma]$ نوں اُدوں تے صرف اُدوں
    پورا کردا اے جدوں $\mSat{M}{!A}[f(\sigma)]$۔
  \item $\mModel{M}$، $\sFmla{\False}{!A}[\sigma]$ نوں اُدوں تے صرف اُدوں
    پورا کردا اے جدوں $\mSat/{M}{!A}[f(\sigma)]$۔
  \end{enumerate}
\end{defn}

\begin{defn}
  منّ لو $\Gamma$ اگیتے والے !!{formula}s دا اک سیٹ اے، تے $P(\Gamma)$
  اوہدے وچ آئے اگیتیاں دا سیٹ اے۔ جے $f$،~$P(\Gamma)$ دی $\mModel{M}$
  وچ تعبیر اے، تاں اسیں آکھدے آں کہ $\mModel{M}$،$\Gamma$ نوں~$f$ دے لحاظ نال پورا کردا اے، $\mSat{M}{\Gamma}[f]$، جے $\mModel{M}$،
  ~$\Gamma$ دے ہر اگیتے والے !!{formula} نوں~$f$ دے لحاظ نال پورا کردا اے۔
  $\Gamma$ \emph{پورا ہون جوگ} اے اُدوں تے صرف اُدوں جدوں کوئی
  نمونہ~$\mModel{M}$ تے کوئی تعبیر~$f$، $P(\Gamma)$ لئی، ہووے جیہدے لئی
  $\mSat{M}{\Gamma}[f]$۔
\end{defn}

\begin{prop}
  جے $\Gamma$ وچ $\sFmla{\True}{!A}[\sigma]$ تے
  $\sFmla{\False}{!A}[\sigma]$ دونیں ہون، کسے !!{formula}~$!A$ تے
  اگیتے~$\sigma$ لئی، تاں $\Gamma$ پورا نہیں ہو سکدا۔
\end{prop}

\begin{proof}
  ایسا نمونہ~$\mModel{M}$ تے تعبیر~$f$، $P(\Gamma)$ لئی، ہو ای نہیں سکدے
  جیہناں لئی $\mSat{M}{!A}[f(\sigma)]$ تے
  $\mSat/{M}{!A}[f(\sigma)]$ دونیں ہون۔
\end{proof}

\begin{thm}[صحت مندی]
  \ollabel{thm:tableau-soundness}
  جے $\Gamma$ دا کوئی بند !!{tableau} اے، تاں $\Gamma$ پورا نہیں ہو سکدا۔
\end{thm}

\begin{proof}
!!a{tableau} دی شاخ پورا ہون جوگ اُدوں تے صرف اُدوں اے جدوں اوہدے اُتے
آئے !!{signed formula}s دا سیٹ پورا ہون جوگ اے۔ !!a{tableau} نوں پورا
ہون جوگ آکھاں گے جے اوہدے وچ گھٹ توں گھٹ اک پورا ہون جوگ شاخ ہووے۔

اسیں ایہ وکھاؤنے آں کہ کسے پورا ہون جوگ !!{tableau} نوں اخذ دے
کسے قاعدے نال ودھائیے تاں ہر واری پورا ہون جوگ !!{tableau} ای ملدا اے۔
ایس نال قضیہ ثابت ہو جائے گا۔ ہر بند !!{tableau} اوہناں قاعدیاں نوں
اوس !!{tableau} اُتے لا کے ملدا اے جیہدے وچ صرف~$\Gamma$ دے مفروضے نیں۔
سو جے $\Gamma$ پورا ہون جوگ ہووے، تاں اوہدا ہر !!{tableau} پورا ہون جوگ
ہووے گا۔ پر بند !!{tableau} پورا نہیں ہو سکدا، کیوں جے اوہدیاں ساریاں
شاخاں بند ہوندیاں نیں، تے بند شاخاں پوریاں نہیں ہو سکدیاں۔

منّ لو ساڈے کول پورا ہون جوگ !!{tableau} اے، یعنی اجیہا !!a{tableau}
جیہدے وچ گھٹ توں گھٹ اک پورا ہون جوگ شاخ اے۔ اخذ دا قاعدہ لاؤن نال
یا تاں کسے شاخ اُتے !!{signed formula}s ودھدے نیں، یا شاخ دو شاخاں وچ
ونڈی جاندی اے۔ جے !!{tableau} وچ کوئی پورا ہون جوگ شاخ اے جیہڑی ایس
قاعدے نال نہیں ودھی، تاں ودھے ہوئے !!{tableau} وچ وی اوہ پورا ہون جوگ
شاخ رہندی اے۔ ایس لئی ودھیا تابلو پورا ہون جوگ اے۔ سو اسیں صرف اوہ صورت
ویکھنی اے جیہدے وچ قاعدہ کسے پورا ہون جوگ شاخ اُتے لگدا اے۔

منّ لو $\Gamma$ اوس شاخ دے !!{signed formula}s دا سیٹ اے، تے
$\sFmla{S}{!A}[\sigma] \in \Gamma$ اوہ !!{signed formula} اے جیہدے اُتے
قاعدہ لگدا اے۔ جے قاعدے نال شاخ نہیں ونڈی جاندی، تاں اسیں وکھانا اے
کہ ودھی ہوئی شاخ، یعنی قاعدے دے نتیجیاں سنے $\Gamma$، ہُن وی پورا ہون
جوگ اے۔ جے قاعدے نال شاخ دو حصیاں وچ ونڈی جاندی اے، تاں اسیں وکھانا
اے کہ ملن والیاں دونہاں شاخاں وچوں گھٹ توں گھٹ اک پورا ہون جوگ اے۔
\tagfalse{prvDiamond}
پہلاں اوہ اخذ ویکھدے آں جیہناں دا صرف اک مقدمہ اے۔
\begin{enumerate}
\item شاخ، قاعدہ $\TRule{\True}{\lnot}$ نوں
  $\sFmla{\True}{\lnot !B}[\sigma] \in \Gamma$ اُتے لا کے ودھدی اے۔
  تاں ودھی ہوئی شاخ اُتے !!{signed formula}s
  $\Gamma \cup \{\sFmla{\False}{!B}[\sigma]\}$ نیں۔ منّ لو
  $\mSat{M}{\Gamma}[f]$۔ خاص طور اُتے $\mSat{M}{\lnot !B}[f(\sigma)]$۔
  ایس لئی $\mSat/{M}{!B}[f(\sigma)]$، یعنی $\mModel{M}$،
  $\sFmla{\False}{!B}[\sigma]$ نوں~$f$ دے لحاظ نال پورا کردا اے۔
\item شاخ، قاعدہ $\TRule{\False}{\lnot}$ نوں
  $\sFmla{\False}{\lnot !B}[\sigma] \in \Gamma$ اُتے لا کے ودھدی اے۔ مشق۔
\item شاخ، قاعدہ $\TRule{\True}{\land}$ نوں
  $\sFmla{\True}{!B \land !C}[\sigma] \in \Gamma$ اُتے لا کے ودھدی اے۔
  ایس نال شاخ اُتے دو نویں !!{signed formula}s
  $\sFmla{\True}{!B}[\sigma]$ تے $\sFmla{\True}{!C}[\sigma]$ ملدے نیں۔
  منّ لو $\mSat{M}{\Gamma}[f]$، خاص کر کے $\mSat{M}{!B \land
    !C}[f(\sigma)]$۔ تاں $\mSat{M}{!B}[f(\sigma)]$ تے
  $\mSat{M}{!C}[f(\sigma)]$۔ ایس دا مطلب اے کہ $\mModel{M}$،
  $\sFmla{\True}{!B}[\sigma]$ تے $\sFmla{\True}{!C}[\sigma]$
  دونہاں نوں~$f$ دے لحاظ نال پورا کردا اے۔
\item شاخ، قاعدہ $\TRule{\False}{\lor}$ نوں
  $\sFmla{\False}{!B \lor !C} \in \Gamma$ اُتے لا کے ودھدی اے۔ مشق۔
\item شاخ، قاعدہ $\TRule{\False}{\lif}$ نوں
  $\sFmla{\False}{!B \lif !C}[\sigma] \in \Gamma$ اُتے لا کے ودھدی اے۔
  شاخ اُتے دو نویں !!{signed formula}s ملدے نیں،
  $\sFmla{\True}{!B}[\sigma]$ تے $\sFmla{\False}{!C}[\sigma]$۔
  منّ لو $\mSat{M}{\Gamma}[f]$، خاص کر کے
  $\mSat/{M}{!B \lif !C}[f(\sigma)]$۔ تاں
  $\mSat{M}{!B}[f(\sigma)]$ تے $\mSat/{M}{!C}[f(\sigma)]$۔
  ایس دا مطلب اے کہ $\mModel{M}, f$،
  $\sFmla{\True}{!B}[\sigma]$ تے $\sFmla{\False}{!C}[\sigma]$ دونہاں نوں
  پورا کردے نیں۔

\iftag{prvBox}{%
\item شاخ، قاعدہ $\TRule{\True}{\Box}$ نوں
  $\sFmla{\True}{\Box !B}[\sigma] \in \Gamma$ اُتے لا کے ودھدی اے۔
  \iftag{probBox}{مشق۔}{ایس نال شاخ اُتے نواں !!{signed formula}~$\sFmla{\True}{!B}[\sigma.n]$ ملدا اے، کسے
    $\sigma.n \in P(\Gamma)$ لئی (کیوں جے $\sigma.n$ پہلاں توں ورتیا
    ہویا ہونا چاہیدا اے)۔ منّ لو $\mSat{M}{\Gamma}[f]$، خاص کر کے
    $\mSat{M}{\Box !B}[f(\sigma)]$۔ جدوں $f$ اگیتیاں دی تعبیر اے تے
    $\sigma$, $\sigma.n \in P(\Gamma)$ دونیں ہون، تاں اسیں جاندے آں کہ
    $Rf(\sigma)f(\sigma.n)$۔ ایس لئی $\mSat{M}{!B}[f(\sigma.n)]$،
    یعنی $\mModel{M}, f$، $\sFmla{\True}{!B}[\sigma.n]$ نوں پورا کردے نیں۔}

\item شاخ، قاعدہ $\TRule{\False}{\Box}$ نوں
  $\sFmla{\False}{\Box !B}[\sigma] \in \Gamma$ اُتے لا کے ودھدی اے۔
  \iftag{probBox}{مشق۔}{ایس نال نواں !!{signed formula}~$\sFmla{\False}{!A}[\sigma.n]$ ملدا اے، جتھے $\sigma.n$
    شاخ لئی نواں اگیتا اے، یعنی $\sigma.n \notin P(\Gamma)$۔
    $\Gamma$ پورا ہون جوگ اے، ایس لئی کوئی $\Struct{M}$ تے تعبیر~$f$، $P(\Gamma)$ لئی، موجود اے جیہدے لئی $\Sat{M}{\Gamma}[f]$، خاص کر کے
    $\mSat/{M}{\Box !B}[f(\sigma)]$۔ اسیں وکھانا اے کہ
    $\Gamma \cup \{\sFmla{\False}{!B}[\sigma.n]\}$ پورا ہون جوگ اے۔
    ایس لئی اسیں $P(\Gamma) \cup \{\sigma.n\}$ دی تعبیر ایویں
    متعین کردے آں:

$\mSat/{M}{\Box !B}[f(\sigma)]$ اے، ایس لئی کوئی $w \in W$ اے
  جیہدے لئی $Rf(\sigma)w$ تے $\mSat/{M}{!B}[w]$۔ منّ لو $f'$، $f$ ورگا
  ای اے، بس $f'(\sigma.n) = w$۔ $f'(\sigma) = f(\sigma)$ تے
  $Rf(\sigma)w$ ہون کر کے، ایس لئی $Rf'(\sigma)f'(\sigma.n)$، سو $f'$،
  $P(\Gamma) \cup \{\sigma.n\}$ دی تعبیر اے۔ صاف اے کہ
  $\mSat/{M}{!B}[f'(\sigma.n)]$۔ $f(\sigma') = f'(\sigma')$
  ہر اگیتے $\sigma' \in P(\Gamma)$ لئی اے، ایس لئی
  $\mSat{M}{\Gamma}[f']$۔ سو $\mModel{M}, f'$،
  $\Gamma \cup \{\sFmla{\False}{!B}[\sigma.n]\}$ نوں پورا کردے نیں۔}
  }{}

\iftag{prvDiamond}{%
\item شاخ، قاعدہ $\TRule{\True}{\Diamond}$ نوں
  $\sFmla{\True}{\Diamond !B}[\sigma] \in \Gamma$ اُتے لا کے ودھدی اے۔
  \iftag{probDiamond}{مشق۔}{ایس نال نواں !!{signed formula}~$\sFmla{\True}{!A}[\sigma.n]$ ملدا اے، جتھے $\sigma.n$
    شاخ لئی نواں اگیتا اے، یعنی $\sigma.n \notin P(\Gamma)$۔
    $\Gamma$ پورا ہون جوگ اے، ایس لئی کوئی $\Struct{M}$ تے تعبیر~$f$، $P(\Gamma)$ لئی، موجود اے جیہدے لئی $\Sat{M}{\Gamma}[f]$، خاص کر کے
    $\mSat{M}{\Diamond !B}[f(\sigma)]$۔ اسیں وکھانا اے کہ
    $\Gamma \cup \{\sFmla{\True}{!B}[\sigma.n]\}$ پورا ہون جوگ اے۔
    ایس لئی اسیں $P(\Gamma) \cup \{\sigma.n\}$ دی تعبیر ایویں
    متعین کردے آں:

$\mSat{M}{\Diamond !B}[f(\sigma)]$ اے، ایس لئی کوئی $w \in W$ اے
  جیہدے لئی $Rf(\sigma)w$ تے $\mSat{M}{!B}[w]$۔ منّ لو $f'$، $f$ ورگا
  ای اے، بس $f'(\sigma.n) = w$۔ $f'(\sigma) = f(\sigma)$ تے
  $Rf(\sigma)w$ ہون کر کے، ایس لئی $Rf'(\sigma)f'(\sigma.n)$، سو $f'$،
  $P(\Gamma) \cup \{\sigma.n\}$ دی تعبیر اے۔ صاف اے کہ
  $\mSat{M}{!B}[f'(\sigma.n)]$۔ $f(\sigma') = f'(\sigma')$
  ہر اگیتے $\sigma' \in P(\Gamma)$ لئی اے، ایس لئی
  $\mSat{M}{\Gamma}[f']$۔ سو $\mModel{M}, f'$،
  $\Gamma \cup \{\sFmla{\True}{!B}[\sigma.n]\}$ نوں پورا کردے نیں۔}
\item شاخ، قاعدہ $\TRule{\False}{\Diamond}$ نوں
  $\sFmla{\False}{\Diamond !B}[\sigma] \in \Gamma$ اُتے لا کے ودھدی اے۔
  \iftag{probDiamond}{مشق۔}{ایس نال شاخ اُتے نواں !!{signed formula}~$\sFmla{\False}{!B}[\sigma.n]$ ملدا اے، کسے
    $\sigma.n \in P(\Gamma)$ لئی (کیوں جے $\sigma.n$ پہلاں توں ورتیا
    ہویا ہونا چاہیدا اے)۔ منّ لو $\mSat{M}{\Gamma}[f]$، خاص کر کے
    $\mSat/{M}{\Diamond !B}[f(\sigma)]$۔ جدوں $f$ اگیتیاں دی تعبیر اے تے
    $\sigma$, $\sigma.n \in P(\Gamma)$ دونیں ہون، تاں اسیں جاندے آں کہ
    $Rf(\sigma)f(\sigma.n)$۔ ایس لئی $\mSat/{M}{!B}[f(\sigma.n)]$،
    یعنی $\mModel{M}, f$، $\sFmla{\False}{!B}[\sigma.n]$ نوں پورا کردے نیں۔}
  }{}
\end{enumerate}
ہُن اوہ اخذ ویکھدے آں جیہناں دے دو مقدمے نیں۔
\begin{enumerate}
\item شاخ، قاعدہ $\TRule{\False}{\land}$ نوں
  $\sFmla{\False}{!B \land !C}[\sigma] \in \Gamma$ اُتے لا کے ودھدی اے۔
  ایس نال دو شاخاں بندیاں نیں۔ کھبی شاخ $\sFmla{\False}{!B}[\sigma]$
  راہیں تے سجی شاخ $\sFmla{\False}{!C}[\sigma]$ راہیں اگے ودھدی اے۔
  منّ لو $\mSat{M}{\Gamma}[f]$، خاص کر کے
  $\mSat/{M}{!B \land !C}[f(\sigma)]$۔ تاں
  $\mSat/{M}{!B}[f(\sigma)]$ یا $\mSat/{M}{!C}[f(\sigma)]$۔
  پہلی صورت وچ $\mModel{M}, f$، $\sFmla{\False}{!B}[\sigma]$ نوں
  پورا کردے نیں، یعنی کھبی شاخ پورا ہون جوگ اے۔ دوجی صورت وچ
  $\mModel{M}, f$، $\sFmla{\False}{!C}[\sigma]$ نوں پورا کردے نیں،
  یعنی سجی شاخ پورا ہون جوگ اے۔
\item شاخ، قاعدہ $\TRule{\True}{\lor}$ نوں
    $\sFmla{\True}{!B \lor !C}[\sigma] \in \Gamma$ اُتے لا کے ودھدی اے۔ مشق۔
\item شاخ، قاعدہ $\TRule{\True}{\lif}$ نوں
    $\sFmla{\True}{!B \lif !C}[\sigma] \in \Gamma$ اُتے لا کے ودھدی اے۔ مشق۔
\end{enumerate}
\end{proof}

\begin{prob}
\olref[nml][tab][sou]{thm:tableau-soundness} دا ثبوت پورا کرو۔
\end{prob}

\begin{cor}
\ollabel{cor:entailment-soundness}
جے $\Gamma \Proves !A$، تاں $\Gamma \Entails !A$۔
\end{cor}

\begin{proof}
  جے $\Gamma \Proves !A$، تاں کجھ $!B_1$, \dots, $!B_n \in
  \Gamma$ لئی، $\Delta = \{\sFmla{\False}{!A}[1], \sFmla{\True}{!B_1}[1],
  \dots, \sFmla{\True}{!B_n}[1]\}$ دا بند !!{tableau} اے۔
  اسیں وکھانا اے کہ $\Gamma \Entails !A$۔ منّ لو ایہ نہیں۔ تاں کسے
  $\mModel{M}$ تے $w$ لئی $\mSat{M}{!B_i}[w]$، $i=1$, \dots,~$n$
  لئی اے، پر $\mSat/{M}{!A}[w]$۔ منّ لو $f(1) = w$۔ تاں $f$،
  ~$P(\Delta)$ دی~$\mModel{M}$ وچ تعبیر اے، تے $\mModel{M}$،
  ~$\Delta$ نوں~$f$ دے لحاظ نال پورا کردا اے۔ پر
  \olref{thm:tableau-soundness} مطابق، $\Delta$ دا بند !!{tableau} ہون
  کر کے اوہ پورا نہیں ہو سکدا۔ ایہ تضاد اے۔ سو آخرکار
  $\Gamma \Proves !A$ ہونا لازمی اے۔
\end{proof}

\begin{cor}
\ollabel{cor:weak-soundness}
جے $\Proves !A$، تاں $!A$ سارے نمونیاں وچ سچ اے۔
\end{cor}
\endgroup
% /OLP-0464
% OLP-0465 content/normal-modal-logic/tableaux/more-rules.tex
\begingroup
\def\olfilename{more-rules}
\expandafter\def\csname ol@sectioncs\endcsname{section}

\olfileid{nml}{tab}{mru}

\olsection{ہور رسائی تعلقاں لئی قاعدے}

خاص رسائی تعلقاں توں متعین ہوندیاں منطقاں لئی، اسیں
\olref{tab:more-rules} وچ دتے وادھو قاعدے ویکھدے آں۔

\begin{table}
  \begin{center}
    \def\arraystretch{3}\def\fCenter{}
    \iftag{notprvBox,notprvDiamond}
          {\begin{tabular}{|c|}}
          {\begin{tabular}{|c|c|}}
    \hline
    \iftag{prvBox}{
      \AxiomC{\sFmla{\True}{\Box !A}[\sigma]}
      \RightLabel{T$\Box$}
      \UnaryInfC{\sFmla{\True}{!A}[\sigma]}
      \DisplayProof}{}
    \iftag{notprvBox,notprvDiamond}{}{&}
    \iftag{prvDiamond}{
      \AxiomC{\sFmla{\False}{\Diamond !A}[\sigma]}
      \RightLabel{T$\Diamond$}
      \UnaryInfC{\sFmla{\False}{!A}[\sigma]}
      \DisplayProof}{}
    \\[1ex]
    \hline
    \iftag{prvBox}{
      \AxiomC{\sFmla{\True}{\Box !A}[\sigma]}
      \RightLabel{D$\Box$}
      \UnaryInfC{\sFmla{\True}{\Diamond!A}[\sigma]}
      \DisplayProof}{}
    \iftag{notprvBox,notprvDiamond}{}{&}
    \iftag{prvDiamond}{
      \AxiomC{\sFmla{\False}{\Diamond !A}[\sigma]}
      \RightLabel{D$\Diamond$}
      \UnaryInfC{\sFmla{\False}{\Box!A}[\sigma]}
      \DisplayProof}{}
    \\[1ex]
    \hline
    \iftag{prvBox}{
      \AxiomC{\sFmla{\True}{\Box !A}[\sigma.n]}
      \RightLabel{B$\Box$}
      \UnaryInfC{\sFmla{\True}{!A}[\sigma]}
      \DisplayProof}{}
    \iftag{notprvBox,notprvDiamond}{}{&}
    \iftag{prvDiamond}{
      \AxiomC{\sFmla{\False}{\Diamond !A}[\sigma.n]}
      \RightLabel{B$\Diamond$}
      \UnaryInfC{\sFmla{\False}{!A}[\sigma]}
      \DisplayProof}{}
    \\[1ex]
    \hline
    \iftag{prvBox}{
      \AxiomC{\sFmla{\True}{\Box !A}[\sigma]}
      \RightLabel{4$\Box$}
      \UnaryInfC{\sFmla{\True}{\Box!A}[\sigma.n]}
      \DisplayProof}{}
    \iftag{notprvBox,notprvDiamond}{}{&}
    \iftag{prvDiamond}{
      \AxiomC{\sFmla{\False}{\Diamond !A}[\sigma]}
      \RightLabel{4$\Diamond$}
      \UnaryInfC{\sFmla{\False}{\Diamond!A}[\sigma.n]}
      \DisplayProof}{}
    \\
    \iftag{notprvBox,notprvDiamond}
         {$\sigma.n$ ورتیا ہویا اے}
         {$\sigma.n$ ورتیا ہویا اے & $\sigma.n$ ورتیا ہویا اے}
    \\[1ex]
    \hline
    \iftag{prvBox}{
      \AxiomC{\sFmla{\True}{\Box !A}[\sigma.n]}
      \RightLabel{4r$\Box$}
      \UnaryInfC{\sFmla{\True}{\Box!A}[\sigma]}
      \DisplayProof}{}
    \iftag{notprvBox,notprvDiamond}{}{&}
    \iftag{prvDiamond}{
      \AxiomC{\sFmla{\False}{\Diamond !A}[\sigma.n]}
      \RightLabel{4r$\Diamond$}
      \UnaryInfC{\sFmla{\False}{\Diamond!A}[\sigma]}
      \DisplayProof}{}
    \\[1ex]
    \hline
    \end{tabular}
  \end{center}
  \caption{ہور امکانی قاعدے۔}
  \ollabel{tab:more-rules}
\end{table}

ایہ قاعدے شامل کرن نال ملے نظام، \olref{tab:logics-rules}
وچ دتیاں منطقاں لئی صحت مند تے مکمل نیں۔

\begin{table}
  \begin{center}
    \begin{tabular}{lll}
      \hline
      منطق & $R$ دی خاصیت \dots & قاعدے\\
      \hline
      $\Log{T} = \Log{KT}$ & انعکاسی &
      \iftag{prvBox}{T$\Box$}{}%
      \iftag{notprvBox,notprvDiamond}{}{, }%
      \iftag{prvDiamond}{T$\Diamond$}{}
      \\ \hline
      $\Log{D} = \Log{KD}$ & ہر جہان توں رسائی والا &
      \iftag{prvBox}{D$\Box$}{}%
      \iftag{notprvBox,notprvDiamond}{}{, }%
      \iftag{prvDiamond}{D$\Diamond$}{}
      \\ \hline
      $\Log{K4}$ & متعدی &
      \iftag{prvBox}{4$\Box$}{}%
      \iftag{notprvBox,notprvDiamond}{}{, }%
      \iftag{prvDiamond}{4$\Diamond$}{}
      \\ \hline
      $\Log{B} = \Log{KTB}$ & انعکاسی, &
      \iftag{prvBox}{T$\Box$}{}%
      \iftag{notprvBox,notprvDiamond}{}{, }%
      \iftag{prvDiamond}{T$\Diamond$}{}\\
      & متقابل &
      \iftag{prvBox}{B$\Box$}{}%
      \iftag{notprvBox,notprvDiamond}{}{, }%
      \iftag{prvDiamond}{B$\Diamond$}{}
      \\ \hline
      $\Log{S4} = \Log{KT4}$ & انعکاسی, &
      \iftag{prvBox}{T$\Box$}{}%
      \iftag{notprvBox,notprvDiamond}{}{, }%
      \iftag{prvDiamond}{T$\Diamond$}{},\\
      & متعدی &
      \iftag{prvBox}{4$\Box$}{}%
      \iftag{notprvBox,notprvDiamond}{}{, }%
      \iftag{prvDiamond}{4$\Diamond$}{}
      \\ \hline
      $\Log{S5} = \Log{KT4B}$ & انعکاسی, &
      \iftag{prvBox}{T$\Box$}{}%
      \iftag{notprvBox,notprvDiamond}{}{, }%
      \iftag{prvDiamond}{T$\Diamond$}{},\\
      & متعدی, &
      \iftag{prvBox}{4$\Box$}{}%
      \iftag{notprvBox,notprvDiamond}{}{, }%
      \iftag{prvDiamond}{4$\Diamond$}{},\\
      &  اقلیدسی &
      \iftag{prvBox}{4r$\Box$}{}%
      \iftag{notprvBox,notprvDiamond}{}{, }%
      \iftag{prvDiamond}{4r$\Diamond$}{}
      \\ \hline
    \end{tabular}
  \end{center}
  \caption{وکھ وکھ امکانی منطقاں لئی !!^{tableau} دے قاعدے۔}
  \ollabel{tab:logics-rules}
\end{table}

\begin{ex}
  اسیں اک بند تابلو دیندے آں جیہڑا وکھاؤندا اے کہ $\Log{S5} \Proves \Ax{5}$، یعنی
  $\Box!A \lif \Box\Diamond!A$.
  \begin{oltableau}
    [\pFmla{\False}{\Box\formula{A} \lif \Box\Diamond \formula{A}}{1},
      just = \TAss
      [\pFmla{\True}{\Box \formula{A}}{1}, just = {\TRule{\False}{\lif}[1]}
        [\pFmla{\False}{\Box\Diamond \formula{A}}{1},
          just = {\TRule{\False}{\lif}[1]}
          [\pFmla{\False}{\Diamond \formula{A}}{1.1},
            just = {\TRule{\False}{\Box}[3]}
            [\pFmla{\False}{\Diamond \formula{A}}{1},
              just = {4r$\Diamond$ 4}
              [\pFmla{\False}{\formula{A}}{1.1},
                just = {\TRule{\False}{\Diamond}[5]}
                [\pFmla{\True}{\formula{A}}{1.1},
                  just = {\TRule{\True}{\Box}[2]}, close]
              ]
            ]
          ]
        ]
      ]
    ]
  \end{oltableau}
\end{ex}

\begin{prob}
ایہ گلاں وکھاؤن والے بند !!{tableau}s دیو:
  \begin{enumerate}
  \item $\Log{KT5} \Proves \Ax{B}$;
  \item $\Log{KT5} \Proves \Ax{4}$;
  \item $\Log{KDB4} \Proves \Ax{T}$;
  \item $\Log{KB4} \Proves \Ax{5}$;
  \item $\Log{KB5} \Proves \Ax{4}$;
  \item $\Log{KT} \Proves \Ax{D}$.
  \end{enumerate}
\end{prob}
\endgroup
% /OLP-0465
% OLP-0466 content/normal-modal-logic/tableaux/more-soundness.tex
\begingroup
\def\olfilename{more-soundness}
\expandafter\def\csname ol@sectioncs\endcsname{section}

\olfileid{nml}{tab}{msn}

\olsection{وادھو قاعدیاں دی صحت مندی}

قاعدہ نمونیاں دی کسے جماعت لئی صحت مند اے جے جدوں وی !!a{tableau} دی کوئی شاخ اوس جماعت دے کسے نمونے وچ پورا ہون جوگ ہووے، قاعدہ لاؤن نال ملدی شاخ وی اوس جماعت دے کسے نمونے وچ پورا ہون جوگ ہووے۔

\begin{prop}
  \ollabel{prop:soundness-T} 
  \iftag{prvBox}{\Ax{T}$\Box$}{}%
  \iftag{notprvBox,notprvDiamond}{}{ تے }%
  \iftag{prvDiamond}{\Ax{T}$\Diamond$}{}
  انعکاسی نمونیاں لئی صحت مند \iftag{notprvBox,notprvDiamond}{اے}{نیں}۔
\end{prop}

\begin{proof}
\begin{tagenumerate}{prvBox,prvDiamond}
\tagitem{prvBox}{شاخ، قاعدہ T$\Box$ نوں $\sFmla{\True}{\Box !B}[\sigma] \in \Gamma$ اُتے لا کے ودھدی اے۔
  \iftag{probBox}{مشق۔}{ایس نال شاخ اُتے نواں !!{signed formula}~$\sFmla{\True}{!B}[\sigma]$ ملدا اے۔ منّ لو $\mSat{M}{\Gamma}[f]$، خاص کر کے $\mSat{M}{\Box
      !B}[f(\sigma)]$۔ $R$ انعکاسی اے، ایس لئی اسیں جاندے آں کہ $Rf(\sigma)f(\sigma)$۔ سو $\mSat{M}{!B}[f(\sigma)]$، یعنی $\mModel{M}, f$، $\sFmla{\True}{!B}[\sigma]$ نوں پورا کردے نیں۔}}{}
\tagitem{prvDiamond}{شاخ، قاعدہ T$\Diamond$ نوں $\sFmla{\False}{\Diamond !B}[\sigma] \in \Gamma$ اُتے لا کے ودھدی اے۔
  \iftag{probDiamond}{مشق۔}{ایس نال شاخ اُتے نواں !!{signed formula}~$\sFmla{\False}{!B}[\sigma]$ ملدا اے۔ منّ لو $\mSat{M}{\Gamma}[f]$، خاص کر کے $\mSat/{M}{\Diamond
      !B}[f(\sigma)]$۔ $R$ انعکاسی اے، ایس لئی اسیں جاندے آں کہ $Rf(\sigma)f(\sigma)$۔ سو $\mSat/{M}{!B}[f(\sigma)]$، یعنی $\mModel{M}, f$، $\sFmla{\False}{!B}[\sigma]$ نوں پورا کردے نیں۔}}{}
\end{tagenumerate}
\end{proof}

\tagprob{probBox,probDiamond}
\begin{prob}
  \olref[nml][tab][msn]{prop:soundness-T} دا ثبوت پورا کرو۔
\end{prob}
\tagendprob

\begin{prop}
  \ollabel{prop:soundness-D} 
  \iftag{prvBox}{\Ax{D}$\Box$}{}%
  \iftag{notprvBox,notprvDiamond}{}{ تے }%
  \iftag{prvDiamond}{\Ax{D}$\Diamond$}{}
  ہر جہان توں رسائی والے نمونیاں لئی صحت مند \iftag{notprvBox,notprvDiamond}{اے}{نیں}۔
\end{prop}

\begin{proof}
\begin{tagenumerate}{prvBox,prvDiamond}
\tagitem{prvBox}{شاخ، قاعدہ D$\Box$ نوں $\sFmla{\True}{\Box !B}[\sigma] \in \Gamma$ اُتے لا کے ودھدی اے۔
  \iftag{probBox}{مشق۔}{ایس نال شاخ اُتے نواں !!{signed formula}~$\sFmla{\True}{\Diamond !B}[\sigma]$ ملدا اے۔ منّ لو $\mSat{M}{\Gamma}[f]$، خاص کر کے $\mSat{M}{\Box !B}[f(\sigma)]$۔ $R$ ہر جہان توں رسائی والا اے، ایس لئی کوئی $w
    \in W$ اے جیہدے لئی $Rf(\sigma)w$۔ تاں $\mSat{M}{!B}[w]$، تے ایس لئی $\mSat{M}{\Diamond !B}[f(\sigma)]$۔ سو $\mModel{M}, f$، $\sFmla{\True}{\Diamond!B}[\sigma]$ نوں پورا کردے نیں۔}}{}
\tagitem{prvDiamond}{شاخ، قاعدہ D$\Diamond$ نوں $\sFmla{\False}{\Diamond !B}[\sigma] \in \Gamma$ اُتے لا کے ودھدی اے۔
  \iftag{probDiamond}{مشق۔}{ایس نال شاخ اُتے نواں !!{signed formula}~$\sFmla{\False}{\Box !B}[\sigma]$ ملدا اے۔ منّ لو $\mSat{M}{\Gamma}[f]$، خاص کر کے $\mSat/{M}{\Diamond !B}[f(\sigma)]$۔ $R$ ہر جہان توں رسائی والا اے، ایس لئی کوئی $w
    \in W$ اے جیہدے لئی $Rf(\sigma)w$۔ تاں $\mSat/{M}{!B}[w]$، تے ایس لئی $\mSat/{M}{\Box !B}[f(\sigma)]$۔ سو $\mModel{M}, f$، $\sFmla{\False}{\Box!B}[\sigma]$ نوں پورا کردے نیں۔}}{}
\end{tagenumerate}
\end{proof}

\tagprob{probBox,probDiamond}
\begin{prob}
  \olref[nml][tab][msn]{prop:soundness-D} دا ثبوت پورا کرو۔
\end{prob}
\tagendprob

\begin{prop}
  \ollabel{prop:soundness-B} 
  \iftag{prvBox}{\Ax{B}$\Box$}{}%
  \iftag{notprvBox,notprvDiamond}{}{ تے }%
  \iftag{prvDiamond}{\Ax{B}$\Diamond$}{}
  متقابل نمونیاں لئی صحت مند \iftag{notprvBox,notprvDiamond}{اے}{نیں}۔
\end{prop}

\begin{proof}
\begin{tagenumerate}{prvBox,prvDiamond}
\tagitem{prvBox}{شاخ، قاعدہ B$\Box$ نوں $\sFmla{\True}{\Box !B}[\sigma.n] \in \Gamma$ اُتے لا کے ودھدی اے۔
  \iftag{probBox}{مشق۔}{ایس نال شاخ اُتے نواں !!{signed formula}~$\sFmla{\True}{!B}[\sigma]$ ملدا اے۔ منّ لو $\mSat{M}{\Gamma}[f]$، خاص کر کے $\mSat{M}{\Box
      !B}[f(\sigma.n)]$۔ $f$ شاخ دے اگیتیاں دی~$\mModel{M}$ وچ تعبیر اے، ایس لئی $Rf(\sigma)f(\sigma.n)$۔ $R$ متقابل اے، ایس لئی $Rf(\sigma.n)f(\sigma)$۔ $\mSat{M}{\Box !B}[f(\sigma.n)]$ ہون کر کے $\mSat{M}{!B}[f(\sigma)]$۔ سو $\mModel{M}, f$، $\sFmla{\True}{!B}[\sigma]$ نوں پورا کردے نیں۔}}{}
\tagitem{prvDiamond}{شاخ، قاعدہ B$\Diamond$ نوں $\sFmla{\False}{\Diamond !B}[\sigma.n] \in \Gamma$ اُتے لا کے ودھدی اے۔
  \iftag{probDiamond}{مشق۔}{ایس نال شاخ اُتے نواں !!{signed formula}~$\sFmla{\False}{!B}[\sigma]$ ملدا اے۔ منّ لو $\mSat{M}{\Gamma}[f]$، خاص کر کے $\mSat/{M}{\Diamond
      !B}[f(\sigma.n)]$۔ $f$ شاخ دے اگیتیاں دی~$\mModel{M}$ وچ تعبیر اے، ایس لئی $Rf(\sigma)f(\sigma.n)$۔ $R$ متقابل اے، ایس لئی $Rf(\sigma.n)f(\sigma)$۔ $\mSat/{M}{\Diamond !B}[f(\sigma.n)]$ ہون کر کے $\mSat/{M}{!B}[f(\sigma)]$۔ سو $\mModel{M}, f$، $\sFmla{\False}{!B}[\sigma]$ نوں پورا کردے نیں۔}}{}
\end{tagenumerate}
\end{proof}

\tagprob{probBox,probDiamond}
\begin{prob}
  \olref[nml][tab][msn]{prop:soundness-B} دا ثبوت پورا کرو۔
\end{prob}
\tagendprob

\begin{prop}
  \ollabel{prop:soundness-4} 
  \iftag{prvBox}{\Ax{4}$\Box$}{}%
  \iftag{notprvBox,notprvDiamond}{}{ تے }%
  \iftag{prvDiamond}{\Ax{4}$\Diamond$}{}
  متعدی نمونیاں لئی صحت مند \iftag{notprvBox,notprvDiamond}{اے}{نیں}۔
\end{prop}

\begin{proof}
\begin{tagenumerate}{prvBox,prvDiamond}
\tagitem{prvBox}{شاخ، قاعدہ 4$\Box$ نوں $\sFmla{\True}{\Box !B}[\sigma] \in \Gamma$ اُتے لا کے ودھدی اے۔
  \iftag{probBox}{مشق۔}{ایس نال شاخ اُتے نواں !!{signed formula}~$\sFmla{\True}{\Box!B}[\sigma.n]$ ملدا اے۔ منّ لو $\mSat{M}{\Gamma}[f]$، خاص کر کے $\mSat{M}{\Box !B}[f(\sigma)]$۔ $f$ شاخ دے اگیتیاں دی~$\mModel{M}$ وچ تعبیر اے تے $\sigma.n$ پہلاں توں ورتیا ہویا ہونا چاہیدا اے۔ ایس لئی $Rf(\sigma)f(\sigma.n)$۔ ہُن $w$ کوئی وی اجیہا جہان منّ لو جیہدے لئی $Rf(\sigma.n)w$۔ $R$ متعدی اے، ایس لئی $Rf(\sigma)w$۔ $\mSat{M}{\Box !B}[f(\sigma)]$ ہون کر کے $\mSat{M}{!B}[w]$۔ سو $\mSat{M}{\Box !B}[f(\sigma.n)]$، تے $\mModel{M}, f$، $\sFmla{\True}{\Box !B}[\sigma.n]$ نوں پورا کردے نیں۔}}{}
\tagitem{prvDiamond}{شاخ، قاعدہ 4$\Diamond$ نوں $\sFmla{\False}{\Diamond !B}[\sigma] \in \Gamma$ اُتے لا کے ودھدی اے۔
  \iftag{probDiamond}{مشق۔}{ایس نال شاخ اُتے نواں !!{signed formula}~$\sFmla{\False}{\Diamond!B}[\sigma.n]$ ملدا اے۔ منّ لو $\mSat{M}{\Gamma}[f]$، خاص کر کے $\mSat/{M}{\Diamond !B}[f(\sigma)]$۔ $f$ شاخ دے اگیتیاں دی~$\mModel{M}$ وچ تعبیر اے تے $\sigma.n$ پہلاں توں ورتیا ہویا ہونا چاہیدا اے۔ ایس لئی $Rf(\sigma)f(\sigma.n)$۔ ہُن $w$ کوئی وی اجیہا جہان منّ لو جیہدے لئی $Rf(\sigma.n)w$۔ $R$ متعدی اے، ایس لئی $Rf(\sigma)w$۔ $\mSat/{M}{\Diamond !B}[f(\sigma)]$ ہون کر کے $\mSat/{M}{!B}[w]$۔ سو $\mSat/{M}{\Diamond !B}[f(\sigma.n)]$، تے $\mModel{M}, f$، $\sFmla{\False}{\Diamond !B}[\sigma.n]$ نوں پورا کردے نیں۔}}{}
\end{tagenumerate}
\end{proof}

\tagprob{probBox,probDiamond}
\begin{prob}
  \olref[nml][tab][msn]{prop:soundness-4} دا ثبوت پورا کرو۔
\end{prob}
\tagendprob

\begin{prop}
  \ollabel{prop:soundness-4r} 
  \iftag{prvBox}{\Ax{4r}$\Box$}{}%
  \iftag{notprvBox,notprvDiamond}{}{ تے }%
  \iftag{prvDiamond}{\Ax{4r}$\Diamond$}{}
  اقلیدسی نمونیاں لئی صحت مند \iftag{notprvBox,notprvDiamond}{اے}{نیں}۔
\end{prop}

\begin{proof}
\begin{tagenumerate}{prvBox,prvDiamond}
\tagitem{prvBox}{شاخ، قاعدہ 4r$\Box$ نوں $\sFmla{\True}{\Box !B}[\sigma.n] \in \Gamma$ اُتے لا کے ودھدی اے۔
  \iftag{probBox}{مشق۔}{ایس نال شاخ اُتے نواں !!{signed formula}~$\sFmla{\True}{\Box!B}[\sigma]$ ملدا اے۔ منّ لو $\mSat{M}{\Gamma}[f]$، خاص کر کے $\mSat{M}{\Box
      !B}[f(\sigma.n)]$۔ $f$ شاخ دے اگیتیاں دی~$\mModel{M}$ وچ تعبیر اے، ایس لئی $Rf(\sigma)f(\sigma.n)$۔ ہُن $w$ کوئی وی اجیہا جہان منّ لو جیہدے لئی $Rf(\sigma)w$۔ $R$ اقلیدسی اے، ایس لئی $Rf(\sigma.n)w$۔ $\mSat{M}{\Box !B}[f(\sigma).n]$ ہون کر کے $\mSat{M}{!B}[w]$۔ سو $\mSat{M}{\Box !B}[f(\sigma)]$، تے $\mModel{M}, f$، $\sFmla{\True}{\Box !B}[\sigma]$ نوں پورا کردے نیں۔}}{}
\tagitem{prvDiamond}{شاخ، قاعدہ 4r$\Diamond$ نوں $\sFmla{\False}{\Diamond !B}[\sigma.n] \in \Gamma$ اُتے لا کے ودھدی اے۔
  \iftag{probDiamond}{مشق۔}{ایس نال شاخ اُتے نواں !!{signed formula}~$\sFmla{\True}{\Box!B}[\sigma]$ ملدا اے۔ منّ لو $\mSat{M}{\Gamma}[f]$، خاص کر کے $\mSat/{M}{\Diamond
      !B}[f(\sigma.n)]$۔ $f$ شاخ دے اگیتیاں دی~$\mModel{M}$ وچ تعبیر اے، ایس لئی $Rf(\sigma)f(\sigma.n)$۔ ہُن $w$ کوئی وی اجیہا جہان منّ لو جیہدے لئی $Rf(\sigma)w$۔ $R$ اقلیدسی اے، ایس لئی $Rf(\sigma.n)w$۔ $\mSat/{M}{\Diamond !B}[f(\sigma).n]$ ہون کر کے $\mSat/{M}{!B}[w]$۔ سو $\mSat/{M}{\Diamond !B}[f(\sigma)]$، تے $\mModel{M}, f$، $\sFmla{\False}{\Diamond !B}[\sigma]$ نوں پورا کردے نیں۔}}{}
\end{tagenumerate}
\end{proof}

\tagprob{probBox,probDiamond}
\begin{prob}
  \olref[nml][tab][msn]{prop:soundness-4r} دا ثبوت پورا کرو۔
\end{prob}
\tagendprob

\begin{cor}
\ollabel{cor:soundness-logics} !!{tableau} دے اوہ نظام جیہڑے
\olref[mru]{tab:logics-rules} وچ دتے نیں، آپو اپنی نمونیاں دی جماعت لئی
صحت مند نیں۔
\end{cor}
\endgroup
% /OLP-0466
% OLP-0467 content/normal-modal-logic/tableaux/simple-S5.tex
\begingroup
\def\olfilename{simple-S5}
\expandafter\def\csname ol@sectioncs\endcsname{section}

\olfileid{nml}{tab}{s5}

\olsection{\Log{S5} لئی سادے \usetoken{P}{tableau}}

\Log{S5} ہر پاسے رسائی والے نمونیاں دی جماعت لئی صحت مند تے مکمل اے،
یعنی اوہ نمونے جیہناں وچ ہر جہان ہر جہان توں رسائی وچ اے۔ ایسے
نمونیاں وچ رسائی دا تعلق فرق نہیں پاؤندا۔ «کوئی جہان~$w$ اے جتھے
$\mSat{M}{!A}[w]$» اُدوں تے صرف اُدوں سچ اے جدوں اجیہا کوئی $w$،
~$u$ توں رسائی وچ اے۔ سو \Log{S5} وچ نمونیاں دی تعریف لئی جہاناں
دا اک سیٹ تے قدر بندی~$V$ ای کافی نیں۔ ایس توں لگدا اے کہ اسیں
!!{tableau} دے قاعدے وی سوکھے کر سکدے آں۔ عام صورت وچ اگیتے مثبت
صحیح عدداں دیاں لڑیاں ہوندے نیں، تاں جو پتا رہے کہ کہڑے اگیتے اوہناں
جہاناں دے ناں نیں جیہڑے دوجے جہاناں توں رسائی وچ نیں۔ $\sigma.n$
ایسے جہان دا ناں اے جیہڑا~$\sigma$ توں رسائی وچ اے۔ پر \Log{S5} وچ
ہر جہان ہر جہان توں رسائی وچ اے، ایس لئی اوہ حساب رکھن دی لوڑ نہیں۔
اگیتیاں لئی مثبت صحیح عدد ورت سکدے آں۔ سوکھے قاعدے
\olref{tab:rules-S5} وچ دتے نیں۔

\begin{table}
  \begin{center}
    \def\arraystretch{3}\def\fCenter{}
    \begin{tabular}{|c|c|}
    \hline
    \iftag{prvBox}{
      \AxiomC{\sFmla{\True}{\Box !A}[n]}
      \RightLabel{\TRule{\True}{\Box}}
      \UnaryInfC{\sFmla{\True}{!A}[m]}
      \DisplayProof
      &
      \AxiomC{\sFmla{\False}{\Box !A}[n]}
      \RightLabel{\TRule{\False}{\Box}}
      \UnaryInfC{\sFmla{\False}{!A}[m]}
      \DisplayProof\\
      $m$ ورتیا ہویا اے & $m$ نواں اے
      \\[1ex]
      \hline}{}
    \iftag{prvDiamond}{
      \AxiomC{\sFmla{\True}{\Diamond !A}[n]}
      \RightLabel{\TRule{\True}{\Diamond}}
      \UnaryInfC{\sFmla{\True}{!A}[m]}
      \DisplayProof
      &
      \AxiomC{\sFmla{\False}{\Diamond !A}[n]}
      \RightLabel{\TRule{\False}{\Diamond}}
      \UnaryInfC{\sFmla{\False}{!A}[m]}
      \DisplayProof\\
      $m$ نواں اے & $m$ ورتیا ہویا اے
      \\[1ex]
      \hline}{}
    \end{tabular}
  \end{center}
  \caption{\Log{S5} لئی سوکھے قاعدے۔}
  \ollabel{tab:rules-S5}
\end{table}

\begin{ex}
  اسیں اک سادہ بند تابلو دیندے آں جیہڑا وکھاؤندا اے کہ $\Log{S5} \Proves
  \Ax{5}$، یعنی $\Diamond!A \lif \Box\Diamond!A$.
  \begin{oltableau}
    [\pFmla{\False}{\Diamond\formula{A} \lif \Box\Diamond \formula{A}}{1},
      just = \TAss
      [\pFmla{\True}{\Diamond \formula{A}}{1}, just = {\TRule{\False}{\lif}[1]}
        [\pFmla{\False}{\Box\Diamond \formula{A}}{1},
          just = {\TRule{\False}{\lif}[1]}
          [\pFmla{\False}{\Diamond \formula{A}}{2},
            just = {\TRule{\False}{\Box}[3]}
            [\pFmla{\True}{\formula{A}}{3},
              just = {\TRule{\True}{\Diamond}[2]}
                [\pFmla{\False}{\formula{A}}{3},
                  just = {\TRule{\False}{\Diamond}[4]}, close]
            ]
          ]
        ]
      ]
    ]
  \end{oltableau}
\end{ex}
\endgroup
% /OLP-0467
% OLP-0468 content/normal-modal-logic/tableaux/completeness.tex
\begingroup
\def\olfilename{completeness}
\expandafter\def\csname ol@sectioncs\endcsname{section}

\olfileid{nml}{tab}{cpl}

\olsection{\Log{K} لئی تمامیت}

\begin{explain}
  ایہ وکھاؤن لئی کہ !!{tableau}s دا طریقہ مکمل اے، اسیں ثابت کرنا اے
  کہ جدوں $\Gamma \Proves !A$ وکھاؤن والا بند !!{tableau} کوئی نہ ہووے،
  تاں $\Gamma \Entails/ !A$، یعنی ردّی نمونہ موجود اے۔ پر «کوئی بند
  !!{tableau} نہیں» دا مطلب اے کہ تابلو بناؤن دا ہر طریقہ بندش تک پہنچن
  وچ ناکام ہووے گا۔ گل ایہ اے کہ جے ایہ سارے طریقے بندش تک نہیں پہنچدے،
  تاں اک خاص \emph{ترتیب وار تے ساریاں صورتاں سمیٹن والا} طریقہ وی
  بندش تک نہیں پہنچدا۔ تے جے کوئی بند !!{tableau} موجود اے، تاں ایہ
  ترتیب وار تے ساریاں صورتاں نوں سمیٹن والا طریقہ بندش تک پہنچے گا۔
  اوس اکلے تابلو دیاں کھلیاں شاخاں وچ ردّی نمونہ متعین کرن لئی لوڑ والی
  ساری جانکاری ہووے گی۔ تابلو دی اک کھلی شاخ توں ملدے ردّی نمونے دے
  جہان اوہ سارے اگیتے ہون گے جیہڑے اوس شاخ اُتے ورتے گئے نیں۔ کوئی
  !!a{propositional variable}~$p$، $\sigma$ اُتے اُدوں تے صرف اُدوں سچ اے
  جدوں $\sFmla{\True}{p}[\sigma]$ شاخ اُتے آوے۔
\end{explain}

\begin{defn}
  !!a{tableau} دی شاخ مکمل اکھواؤندی اے جے جدوں وی اوہدے وچ اگیتے والا
  !!{formula} $\sFmla{S}{!A}[\sigma]$ ہووے جیہدے اُتے کوئی قاعدہ
  لگ سکدا اے، تاں اوہدے وچ ایہ وی ہون:
  \begin{enumerate}
    \item قضیاتی قاعدہ جے اکّو شاخ اُتے نتیجے لاؤندا اے، تاں اوہدے
      سارے متعلقہ نتیجے، یعنی اوہ اگیتے والے !!{formula}s۔
    \item قضیاتی قاعدہ جے شاخ ونڈدا اے، تاں اوہدے متعلقہ نتیجے والے
      !!{formula}s وچوں کوئی اک۔
    \item امکانی قاعدے نوں جے نواں اگیتا چاہیدا اے، تاں گھٹ توں گھٹ
      اک ممکن نتیجہ۔
    \item امکانی قاعدے نوں جے پہلاں ورتیا ہویا اگیتا چاہیدا اے، تاں
      شاخ اُتے آئے ہر اگیتے لئی متعلقہ نتیجہ۔
  \end{enumerate}
\end{defn}

\begin{explain}
مثلاً مکمل شاخ وچ $\sFmla{\True}{!B}[\sigma]$ تے
$\sFmla{\True}{!C}[\sigma]$ دونیں ہوندے نیں جدوں اوہدے وچ
$\sFmla{\True}{!B \land !C}$ ہووے۔ جے اوہدے وچ
$\sFmla{\True}{!B \lor !C}[\sigma]$ اے، تاں
$\sFmla{\False}{!B}[\sigma]$ تے $\sFmla{\True}{!C}[\sigma]$ وچوں
گھٹ توں گھٹ اک ہوندا اے۔ جے اوہدے وچ \iftag{prvBox}
{$\sFmla{\False}{\Box}[\sigma]$} {$\sFmla{\True}{\Diamond}[\sigma]$}
ہووے، تاں اوہدے وچ \iftag{prvBox} {$\sFmla{\False}{\Box}[\sigma.n]$}
{$\sFmla{\True}{\Diamond}[\sigma.n]$} گھٹ توں گھٹ اک~$n$ لئی وی ہوندا اے۔ تے جدوں وی اوہدے وچ \iftag{prvBox} {$\sFmla{\True}{\Box}[\sigma]$}
{$\sFmla{\False}{\Diamond}[\sigma]$} ہووے، اوہدے وچ \iftag{prvBox}
{$\sFmla{\True}{\Box}[\sigma.n]$}
{$\sFmla{\False}{\Diamond}[\sigma.n]$} ہر ایسے~$n$ لئی ہوندا اے
جیہدے لئی $\sigma.n$ شاخ اُتے ورتیا ہویا اے۔
\end{explain}

\begin{prop}\ollabel{prop:complete-tableau}
  ہر محدود $\Gamma$ دا !!a{tableau} ہوندا اے جیہدے وچ ہر شاخ مکمل اے۔
\end{prop}

\begin{proof}
  !!a{tableau} دی کوئی کھلی شاخ، جیہڑا~$\Gamma$ لئی بنیا اے، ویکھو۔ شاخ اُتے محدود
  گنتی وچ ایسے اگیتے والے !!{formula}s نیں جیہناں اُتے قاعدہ لگ سکدا
  اے۔ کوئی ترتیب مقرر کر لو (مثلاً اُتّوں ہِیٹھاں ول)۔ ایس ترتیب وچ
  اجیہے اگیتے والے !!{formula}s وچوں ہر اوس لئی جیہدے لئی شرطاں (1)--(4) پہلاں
  توں پوریاں نہیں، اوہدے اُتے لگ سکن والے قاعدے لا کے شاخ ودھاؤ۔
  کجھ صورتاں وچ شاخ ونڈی جائے گی۔ ملن والی ہر شاخ دے سِرے اُتے
  باقی سارے اگیتے والے !!{formula}s لئی قاعدہ لاؤ۔ اگیتے والے
  !!{formula}s دی گنتی محدود اے، تے شاخ اُتے ورتے ہوئے اگیتیاں دی گنتی
  وی محدود اے۔ ایس لئی ایہ طریقہ آخرکار اصل شاخ نوں ودھاؤن والیاں
  (ہو سکدا اے کئی) شاخاں دیندا اے۔ ایہ طریقہ ہر اک اُتے لاؤ تے
  دہراؤ۔ پر بناوٹ مطابق، ہر شاخ بند اے۔
\end{proof}

\begin{thm}[تمامیت]
  \ollabel{thm:tableau-completeness}
  جے $\Gamma$ دا کوئی بند !!{tableau} نہیں، تاں $\Gamma$ پورا ہون جوگ اے۔
\end{thm}

\begin{proof}
قضیے مطابق، $\Gamma$ دا اک !!a{tableau} اے جیہدے وچ ہر شاخ مکمل اے۔
اوہدا کوئی بند !!{tableau} نہیں، ایس لئی اوہدا اک !!a{tableau} اے
جیہدے وچ گھٹ توں گھٹ اک شاخ کھلی تے مکمل اے۔ منّ لو $\Delta$ اوس
شاخ اُتے اگیتے والے !!{formula}s دا سیٹ اے، تے $P(\Delta)$ اوہدے وچ
آئے اگیتیاں دا سیٹ اے۔

اسیں اک نمونہ~$\mModel{M(\Delta)} = \tuple{P(\Delta), R, V}$ متعین
کردے آں۔ اوہدے جہان~$\Delta$ وچ آئے اگیتے نیں، تے رسائی دا تعلق ایویں اے:
\[
R\sigma\sigma' \quad \text{اُدوں تے صرف اُدوں جدوں} \quad
\sigma'=\sigma.n \quad \text{کسے~$n$ لئی}
\]
تے
\[
V(p) = \Setabs{\sigma}{\sFmla{\True}{p}[\sigma] \in \Delta}.
\]
اسیں~$!A$ اُتے استقرا نال وکھاؤنے آں کہ جے $\sFmla{\True}{!A}[\sigma] \in
\Delta$، تاں $\mSat{M(\Delta)}{!A}[\sigma]$، تے جے
$\sFmla{\False}{!A}[\sigma] \in \Delta$، تاں
$\mSat/{M(\Delta)}{!A}[\sigma]$۔

\begin{enumerate}
  \item \indcase{!A}{p}{جے $\sFmla{\True}{\indfrm}[\sigma] \in \Delta$،
    تاں $\sigma \in V(p)$ (~$V$ دی تعریف مطابق)، ایس لئی
    $\mSat{M(\Delta)}{\indfrm}[\sigma]$۔

جے $\sFmla{\False}{\indfrm}[\sigma] \in \Delta$، تاں
    $\sFmla{\True}{\indfrm}[\sigma] \notin \Delta$، کیوں جے جے ایہ
    اوہدے وچ ہوندا تاں شاخ بند ہوندی۔ سو $\sigma \notin V(p)$، ایس لئی
    $\mSat/{M(\Delta)}{\indfrm}[\sigma]$۔}
  \item \indcase{!A}{\lnot !B}{\iftag{probNot}{مشق۔}{جے
      $\sFmla{\True}{\indfrm}[\sigma] \in \Delta$، تاں
      $\sFmla{\False}{!B}[\sigma] \in \Delta$، کیوں جے شاخ مکمل اے۔
      استقرائی مفروضے مطابق، $\mSat/{M(\Delta)}{!B}[\sigma]$،
      سو $\mSat{M(\Delta)}{\indfrm}[\sigma]$۔

جے $\sFmla{\False}{\indfrm}[\sigma] \in \Delta$، تاں
      $\sFmla{\True}{!B}[\sigma] \in \Delta$، کیوں جے شاخ مکمل اے۔
      استقرائی مفروضے مطابق، $\mSat{M(\Delta)}{!B}[\sigma]$،
      سو $\mSat/{M(\Delta)}{\indfrm}[\sigma]$۔}}
  \item \indcase{!A}{!B \land !C}{\iftag{probAnd}{مشق۔}{جے
      $\sFmla{\True}{\indfrm}[\sigma] \in \Delta$، تاں
      $\sFmla{\True}{!B}[\sigma] \in \Delta$ تے
      $\sFmla{\True}{!C}[\sigma] \in \Delta$ دونیں، کیوں جے شاخ مکمل اے۔
      استقرائی مفروضے مطابق، $\mSat{M(\Delta)}{!B}[\sigma]$ تے
      $\mSat{M(\Delta)}{!C}[\sigma]$۔ سو
      $\mSat{M(\Delta)}{\indfrm}[\sigma]$۔

جے $\sFmla{\False}{\indfrm}[\sigma] \in \Delta$، تاں
      $\sFmla{\False}{!B}[\sigma] \in \Delta$ یا
      $\sFmla{\False}{!C}[\sigma] \in \Delta$، کیوں جے شاخ مکمل اے۔
      استقرائی مفروضے مطابق، $\mSat/{M(\Delta)}{!B}[\sigma]$ یا
      $\mSat/{M(\Delta)}{!B}[\sigma]$۔ سو
      $\mSat/{M(\Delta)}{\indfrm}[\sigma]$۔}}
  \item \indcase{!A}{!B \lor !C}{\iftag{probOr}{مشق۔}{جے
      $\sFmla{\True}{\indfrm}[\sigma] \in \Delta$، تاں
      $\sFmla{\True}{!B}[\sigma] \in \Delta$ یا
      $\sFmla{\True}{!C}[\sigma] \in \Delta$، کیوں جے شاخ مکمل اے۔
      استقرائی مفروضے مطابق، $\mSat{M(\Delta)}{!B}[\sigma]$ یا
      $\mSat{M(\Delta)}{!C}[\sigma]$۔ سو
      $\mSat{M(\Delta)}{\indfrm}[\sigma]$۔

جے $\sFmla{\False}{\indfrm}[\sigma] \in \Delta$، تاں
      $\sFmla{\False}{!B}[\sigma] \in \Delta$ تے
      $\sFmla{\False}{!C}[\sigma] \in \Delta$ دونیں، کیوں جے شاخ مکمل اے۔
      استقرائی مفروضے مطابق، $\mSat/{M(\Delta)}{!B}[\sigma]$ تے
      $\mSat/{M(\Delta)}{!B}[\sigma]$ دونیں۔ سو
      $\mSat/{M(\Delta)}{\indfrm}[\sigma]$۔}}
  \item \indcase{!A}{!B \lif !C}{\iftag{probIf}{مشق۔}{جے
      $\sFmla{\True}{\indfrm}[\sigma] \in \Delta$، تاں
      $\sFmla{\False}{!B}[\sigma] \in \Delta$ یا
      $\sFmla{\True}{!C}[\sigma] \in \Delta$، کیوں جے شاخ مکمل اے۔
      استقرائی مفروضے مطابق، $\mSat/{M(\Delta)}{!B}[\sigma]$ یا
      $\mSat{M(\Delta)}{!C}[\sigma]$۔ سو
      $\mSat{M(\Delta)}{\indfrm}[\sigma]$۔

جے $\sFmla{\False}{\indfrm}[\sigma] \in \Delta$، تاں
      $\sFmla{\True}{!B}[\sigma] \in \Delta$ تے
      $\sFmla{\False}{!C}[\sigma] \in \Delta$ دونیں، کیوں جے شاخ مکمل اے۔
      استقرائی مفروضے مطابق، $\mSat{M(\Delta)}{!B}[\sigma]$ تے
      $\mSat/{M(\Delta)}{!B}[\sigma]$ دونیں۔ سو
      $\mSat/{M(\Delta)}{\indfrm}[\sigma]$۔}}
  \item \indcase{!A}{\Box !B}{\iftag{probBox}{مشق۔}{جے
      $\sFmla{\True}{\indfrm}[\sigma] \in \Delta$، تاں شاخ مکمل
      ہون کر کے $\sFmla{\True}{!B}[\sigma.n] \in \Delta$ شاخ اُتے
      ورتے ہوئے ہر $\sigma.n$ لئی اے، یعنی ہر اوہ $\sigma'
      \in P(\Delta)$ لئی جیہدے لئی $R\sigma\sigma'$۔ استقرائی مفروضے
      مطابق، $\mSat{M(\Delta)}{!B}[\sigma']$ ہر اوہ $\sigma'$ لئی
      اے جیہدے لئی $R\sigma\sigma'$۔ سو
      $\mSat{M(\Delta)}{\indfrm}[\sigma]$۔

جے $\sFmla{\False}{\indfrm}[\sigma] \in \Delta$، تاں کسے
      $\sigma.n$ لئی $\sFmla{\False}{!B}[\sigma.n] \in \Delta$، کیوں
      جے شاخ مکمل اے۔ استقرائی مفروضے مطابق،
      $\mSat/{M(\Delta)}{!B}[\sigma.n]$۔ $R\sigma(\sigma.n)$ ہون
      کر کے کوئی~$\sigma'$ اے جیہدے لئی
      $\mSat/{M(\Delta)}{!B}[\sigma']$۔ سو
      $\mSat/{M(\Delta)}{\indfrm}[\sigma]$۔}}
  \item \indcase{!A}{\Diamond !B}{\iftag{probDiamond}{مشق۔}{جے
      $\sFmla{\True}{\indfrm}[\sigma] \in \Delta$، تاں کسے
      $\sigma.n$ لئی $\sFmla{\True}{!B}[\sigma.n] \in \Delta$، کیوں
      جے شاخ مکمل اے۔ استقرائی مفروضے مطابق،
      $\mSat{M(\Delta)}{!B}[\sigma.n]$۔ $R\sigma(\sigma.n)$ ہون
      کر کے کوئی~$\sigma'$ اے جیہدے لئی
      $\mSat{M(\Delta)}{!B}[\sigma']$۔ سو
      $\mSat{M(\Delta)}{\indfrm}[\sigma]$۔

جے $\sFmla{\False}{\indfrm}[\sigma] \in \Delta$، تاں شاخ
      مکمل ہون کر کے $\sFmla{\False}{!B}[\sigma.n] \in \Delta$
      شاخ اُتے ورتے ہوئے ہر $\sigma.n$ لئی اے، یعنی ہر اوہ
      $\sigma' \in P(\Delta)$ لئی جیہدے لئی $R\sigma\sigma'$۔
      استقرائی مفروضے مطابق، $\mSat/{M(\Delta)}{!B}[\sigma']$
      ہر اوہ $\sigma'$ لئی اے جیہدے لئی $R\sigma\sigma'$۔ سو
      $\mSat/{M(\Delta)}{\indfrm}[\sigma]$۔}}
\end{enumerate}
$\Gamma \subseteq \Delta$ ہون کر کے $\mSat{M(\Delta)}{\Gamma}$۔
\end{proof}

\begin{prob}
\olref[nml][tab][cpl]{thm:tableau-completeness} دا ثبوت پورا کرو۔
\end{prob}

\begin{cor}
\ollabel{cor:entailment-completeness}
جے $\Gamma \Entails !A$، تاں $\Gamma \Proves !A$۔
\end{cor}

\begin{cor}
\ollabel{cor:weak-completeness}
جے $!A$ سارے نمونیاں وچ سچ اے، تاں $\Proves !A$۔
\end{cor}
\endgroup
% /OLP-0468
% OLP-0469 content/normal-modal-logic/tableaux/countermodels.tex
\begingroup
\def\olfilename{countermodels}
\expandafter\def\csname ol@sectioncs\endcsname{section}

\olfileid{nml}{tab}{cou}

\olsection{\usetoken{P}{tableau} توں ردّی نمونے}

\begin{explain}
  تمامیت قضیے دا ثبوت وکھاؤندا اے کہ جے $\Entails !A$، تاں
  $\Proves !A$۔ ایہ~$!A$ لئی ردّی نمونہ بناؤن دا طریقہ وی دیندا اے،
  جے $\Entails/ A$۔~$\Log{K}$ دی صورت وچ ایہ طریقہ \emph{فیصلہ کرن
    دا طریقہ} اے۔ منّ لو $\Entails/ !A$۔ تاں
  \olref[cpl]{prop:complete-tableau} دے ثبوت توں مکمل !!{tableau}
  بناؤن دا اک طریقہ ملدا اے۔ ایہ طریقہ ہمیشہ مک جاندا اے۔
  ~$\Log{K}$ دے قضیاتی قاعدے صرف گھٹ پیچیدگی والے اگیتے والے
  !!{formula}s جوڑدے نیں، یعنی ہر نشان والے فارمولے $\sFmla{S}{!A}[\sigma]$
  لئی اک شاخ اُتے ہر قضیاتی قاعدہ اک واری لانا کافی اے۔ نویں اگیتے صرف
  \iftag{prvBox}{$\TRule{\False}{\Box}$}{}\iftag{notprvBox,notprvDiamond}{}{
    تے }\iftag{prvDiamond}{$\TRule{\True}{\Diamond}$}{}
  \iftag{notprvBox,notprvDiamond}{قاعدے}{قاعدیاں} توں ملدے نیں۔
  ایہ وی اک واری لانے کافی نیں (تے اکّو نواں اگیتا دیندے نیں)۔
  \iftag{prvBox}{$\TRule{\True}{\Box}$}{}\iftag{notprvBox,notprvDiamond}{}{
    تے }\iftag{prvDiamond}{$\TRule{\False}{\Diamond}$}{}
  \iftag{notprvBox,notprvDiamond}{قاعدہ}{قاعدے} اک توں ودھ وار لانے پے
  سکدے نیں، پر ہر اگیتے لئی اکّو واری۔ نویں اگیتے محدود گنتی وچ ای
  بندے نیں۔ سو محدود پڑاواں پچھوں ایس بناوٹ توں یا بند شاخ ملدی اے
  یا مکمل شاخ۔

جدوں کھلی مکمل شاخ والا تابلو بن جائے، تاں
  \olref[cpl]{thm:tableau-completeness} دا ثبوت سانوں اوہ عین نمونہ
  دیندا اے جیہڑا اگیتے والے !!{formula}s دے اصل سیٹ نوں پورا کردا اے۔
  سو جے $\Gamma \Entails !A$، تاں بند !!{tableau} موجود اے تے
  $\Gamma \Proves !A$۔ تے جے بند !!{tableau} نوں ٹھیک طریقے نال لبھدیاں
  اک «مکمل» !!{tableau} ملے، تاں سانوں $\Gamma \Entails/ !A$ دا پتا
  لگے گا تے اسیں ردّی نمونہ وی بنا سکاں گے۔
\end{explain}

\iftag{prvBox}{
\begin{ex}
  سانوں پتا اے کہ $\Proves/ \Box(p \lor q) \lif (\Box p \lor \Box
  q)$۔ تابلو بناؤن دا مڈھ ایویں اے:
  \begin{oltableau}
    [\pFmla{\False}{\Box(p \lor q) \lif (\Box p \lor \Box q)}{1},
      just = \TAss, checked
      [\pFmla{\True}{\Box(p \lor q)}{1},
        just = {\TRule{\False}{\lif}[1]}, 
        [\pFmla{\False}{\Box p \lor \Box q}{1},
          just = {\TRule{\False}{\lif}[1]}, checked
          [\pFmla{\False}{\Box p}{1},
            just = {\TRule{\False}{\lor}[3]}, checked
            [\pFmla{\False}{\Box q}{1},
              just = {\TRule{\False}{\lor}[3]}, checked
              [\pFmla{\False}{p}{1.1}, 
                just = {\TRule{\False}{\Box}[4]}, checked
                [\pFmla{\False}{q}{1.2}, 
                  just = {\TRule{\False}{\Box}[5]}, checked
                ]
              ]
            ]
          ]
        ]
      ]
    ]
  \end{oltableau}
  !!{tableau} ہجے پورا نہیں ہویا۔ اگلے پڑا وچ اسیں اوہ اکّو سطر ویکھدے
  آں جیہدے اُتے صحیح دا نشان نہیں، یعنی اگیتے والا !!{formula}
  $\sFmla{\True}{\Box(p \lor q)}[1]$ سطر~$2$ اُتے۔ بند تابلو بناؤن لئی
  $\TRule{\True}{\Box}$ دا قاعدہ شاخ اُتے ورتے ہوئے ہر اگیتے لئی لانا
  اے، یعنی $1.1$ تے $1.2$ دونہاں لئی:
  \begin{oltableau}
    [\pFmla{\False}{\Box(p \lor q) \lif (\Box p \lor \Box q)}{1},
      just = \TAss, checked
      [\pFmla{\True}{\Box(p \lor q)}{1},
        just = {\TRule{\False}{\lif}[1]}, 
        [\pFmla{\False}{\Box p \lor \Box q}{1},
          just = {\TRule{\False}{\lif}[1]}, checked
          [\pFmla{\False}{\Box p}{1},
            just = {\TRule{\False}{\lor}[3]}, checked
            [\pFmla{\False}{\Box q}{1},
              just = {\TRule{\False}{\lor}[3]}, checked
              [\pFmla{\False}{p}{1.1}, 
                just = {\TRule{\False}{\Box}[4]}, checked
                [\pFmla{\False}{q}{1.2}, 
                  just = {\TRule{\False}{\Box}[5]}, checked
                  [\pFmla{\True}{p \lor q}{1.1}, 
                    just = {\TRule{\True}{\Box}[2]}
                    [\pFmla{\True}{p \lor q}{1.2}, 
                      just = {\TRule{\True}{\Box}[2]}
                    ]
                  ]
                ]
              ]
            ]
          ]
        ]
      ]
    ]
  \end{oltableau}
  ہُن سطر 2، 8 تے 9 اُتے صحیح دے نشان نہیں۔ پر کوئی نواں اگیتا نہیں
  جڑیا، ایس لئی اسیں $\TRule{\True}{\lor}$ سطر 8 تے 9 اُتے، ملن والیاں
  ساریاں شاخاں وچ (جنا چر اوہ بند نہیں ہوندیاں)، لاؤندے آں:
  \begin{oltableau}
    [\pFmla{\False}{\Box(p \lor q) \lif (\Box p \lor \Box q)}{1},
      just = \TAss, checked
      [\pFmla{\True}{\Box(p \lor q)}{1},
        just = {\TRule{\False}{\lif}[1]}, checked
        [\pFmla{\False}{\Box p \lor \Box q}{1},
          just = {\TRule{\False}{\lif}[1]}, checked
          [\pFmla{\False}{\Box p}{1},
            just = {\TRule{\False}{\lor}[3]}, checked
            [\pFmla{\False}{\Box q}{1},
              just = {\TRule{\False}{\lor}[3]}, checked
              [\pFmla{\False}{p}{1.1}, 
                just = {\TRule{\False}{\Box}[4]}, checked
                [\pFmla{\False}{q}{1.2}, 
                  just = {\TRule{\False}{\Box}[5]}, checked
                  [\pFmla{\True}{p \lor q}{1.1}, 
                    just = {\TRule{\True}{\Box}[2]}, checked
                    [\pFmla{\True}{p \lor q}{1.2}, 
                      just = {\TRule{\True}{\Box}[2]}, checked
                      [\pFmla{\True}{p}{1.1},
                        just = {\TRule{\True}{\lor}[8]}, checked, close
                      ]
                      [\pFmla{\True}{q}{1.1},
                        just = {\TRule{\True}{\lor}[8]}, checked
                        [\pFmla{\True}{p}{1.2},
                          just = {\TRule{\True}{\lor}[9]}, checked]
                        [\pFmla{\True}{q}{1.2},
                          just = {\TRule{\True}{\lor}[9]}, checked, close]
                      ]
                    ]
                  ]
                ]
              ]
            ]
          ]
        ]
      ]
    ]
  \end{oltableau}
  اک کھلی شاخ بچدی اے، تے اوہ مکمل اے۔ اوس توں اسیں نمونہ متعین کردے آں۔
  جہان $W = \{1, 1.1, 1.2\}$ نیں (کھلی شاخ اُتے آؤن والے ای اگیتے)،
  رسائی دا تعلق $R = \{\tuple{1, 1.1}, \tuple{1, 1.2}\}$ اے، تے قدر بندی
  $V(p) = \{1.2\}$ اے (کیوں جے سطر 11 اُتے $\sFmla{\True}{p}[1.2]$ اے)
  تے $V(q) = \{1.1\}$ اے (کیوں جے سطر 10 اُتے $\sFmla{\True}{q}[1.1]$ اے)۔
  نمونے دی تصویر \olref{fig:counter-Box} وچ اے، تے تسی پرکھ سکدے او کہ ایہ
  $\Box(p \lor q) \lif (\Box p \lor \Box q)$ دا ردّی نمونہ اے۔
  \begin{figure}
  \begin{center}
    \begin{tikzpicture}[modal]
      \node[world] (w1) [label={[align=right]right:\mFalse{p}\\ \mFalse{q}}]{$1$}; 
      \node[world] (w2) [label={[align=right]right:\mFalse{p}\\ \mTrue{q}},
        above left=of w1]{$1.1$}; 
      \node[world] (w3) [label={[align=right]right:\mTrue{p}\\ \mFalse{q}},
        above right=of w1] {$1.2$};
      \draw[->] (w1) to (w2);
      \draw[->] (w1) to (w3);
    \end{tikzpicture}
  \end{center}
  \caption{$\Box(p \lor q) \lif (\Box p \lor \Box
    q)$ دا اک ردّی نمونہ۔}
  \ollabel{fig:counter-Box}
\end{figure}
\end{ex}
}
{
\begin{ex}
  سانوں پتا اے کہ $\Proves/ (\Diamond p \land \Diamond q) \lif \Diamond(p
  \land q)$۔ تابلو بناؤن دا مڈھ ایویں اے:
  \begin{oltableau}
    [\pFmla{\False}{(\Diamond p \land \Diamond q) \lif \Diamond(p \land q)}{1},
      just = \TAss, checked
      [\pFmla{\True}{\Diamond p \land \Diamond q}{1},
        just = {\TRule{\False}{\lif}[1]}, checked
        [\pFmla{\False}{\Diamond(p \land q)}{1},
          just = {\TRule{\False}{\lif}[1]}
          [\pFmla{\True}{\Diamond p}{1},
            just = {\TRule{\True}{\land}[2]}, checked
            [\pFmla{\True}{\Diamond q}{1},
              just = {\TRule{\True}{\land}[2]}, checked
              [\pFmla{\True}{p}{1.1}, 
                just = {\TRule{\True}{\Diamond}[4]}, checked
                [\pFmla{\True}{q}{1.2}, 
                  just = {\TRule{\True}{\Diamond}[5]}, checked
                ]
              ]
            ]
          ]
        ]
      ]
    ]
  \end{oltableau}
  !!{tableau} ہجے پورا نہیں ہویا۔ اگلے پڑا وچ اسیں اوہ اکّو سطر ویکھدے
  آں جیہدے اُتے صحیح دا نشان نہیں، یعنی اگیتے والا !!{formula}
  $\sFmla{\True}{\Diamond(p \land q)}[1]$ سطر~$3$ اُتے۔ بند تابلو بناؤن لئی
  $\TRule{\True}{\Diamond}$ دا قاعدہ شاخ اُتے ورتے ہوئے ہر اگیتے لئی لانا
  اے، یعنی $1.1$ تے $1.2$ دونہاں لئی:
  \begin{oltableau}
    [\pFmla{\False}{\Diamond(p \land q) \lif (\Diamond p \land \Diamond q)}{1},
      just = \TAss, checked
      [\pFmla{\True}{\Diamond p \land \Diamond q}{1},
        just = {\TRule{\False}{\lif}[1]}, checked
        [\pFmla{\False}{\Diamond (p \land q)}{1},
          just = {\TRule{\False}{\lif}[1]}
          [\pFmla{\True}{\Diamond p}{1},
            just = {\TRule{\True}{\land}[2]}, checked
            [\pFmla{\True}{\Diamond q}{1},
              just = {\TRule{\True}{\land}[2]}, checked
              [\pFmla{\True}{p}{1.1}, 
                just = {\TRule{\True}{\Diamond}[4]}, checked
                [\pFmla{\True}{q}{1.2}, 
                  just = {\TRule{\True}{\Diamond}[5]}, checked
                  [\pFmla{\False}{p \land q}{1.1}, 
                    just = {\TRule{\False}{\Diamond}[3]}
                    [\pFmla{\False}{p \land q}{1.2}, 
                      just = {\TRule{\False}{\Diamond}[3]}
                    ]
                  ]
                ]
              ]
            ]
          ]
        ]
      ]
    ]
  \end{oltableau}
  ہُن سطر 3، 8 تے 9 اُتے صحیح دے نشان نہیں۔ پر کوئی نواں اگیتا نہیں
  جڑیا، ایس لئی اسیں $\TRule{\False}{\land}$ سطر 8 تے 9 اُتے، ملن والیاں
  ساریاں شاخاں وچ (جنا چر اوہ بند نہیں ہوندیاں)، لاؤندے آں:
  \begin{oltableau}
    [\pFmla{\False}{(\Diamond p \land \Diamond q) \lif \Diamond(p \land q)}{1},
      just = \TAss, checked
      [\pFmla{\True}{\Diamond p \land \Diamond q}{1},
        just = {\TRule{\False}{\lif}[1]}, checked
        [\pFmla{\False}{\Diamond(p \land q)}{1},
          just = {\TRule{\False}{\lif}[1]}
          [\pFmla{\True}{\Diamond p}{1},
            just = {\TRule{\True}{\land}[2]}, checked
            [\pFmla{\True}{\Diamond q}{1},
              just = {\TRule{\True}{\land}[2]}, checked
              [\pFmla{\True}{p}{1.1}, 
                just = {\TRule{\True}{\Diamond}[4]}, checked
                [\pFmla{\True}{q}{1.2}, 
                  just = {\TRule{\True}{\Diamond}[5]}, checked
                  [\pFmla{\False}{p \land q}{1.1}, 
                    just = {\TRule{\False}{\Diamond}[3]}, checked
                    [\pFmla{\False}{p \land q}{1.2}, 
                      just = {\TRule{\False}{\Diamond}[3]}, checked
                      [\pFmla{\False}{p}{1.1},
                        just = {\TRule{\False}{\land}[8]}, checked, close
                      ]
                      [\pFmla{\False}{q}{1.1},
                        just = {\TRule{\False}{\land}[8]}, checked
                        [\pFmla{\False}{p}{1.2},
                          just = {\TRule{\False}{\land}[9]}, checked]
                        [\pFmla{\False}{q}{1.2},
                          just = {\TRule{\False}{\land}[9]}, checked, close]
                      ]
                    ]
                  ]
                ]
              ]
            ]
          ]
        ]
      ]
    ]
  \end{oltableau}
  اک کھلی شاخ بچدی اے، تے اوہ مکمل اے۔ اوس توں اسیں نمونہ متعین کردے آں۔
  جہان $W = \{1, 1.1, 1.2\}$ نیں (کھلی شاخ اُتے آؤن والے ای اگیتے)،
  رسائی دا تعلق $R = \{\tuple{1, 1.1}, \tuple{1, 1.2}\}$ اے، تے قدر بندی
  $V(p) = \{1.1\}$ اے (کیوں جے سطر 6 اُتے $\sFmla{\True}{p}[1.1]$ اے)
  تے $V(q) = \{1.2\}$ اے (کیوں جے سطر 7 اُتے $\sFmla{\True}{q}[1.1]$ اے)۔
  نمونے دی تصویر \olref{fig:counter-Diamond} وچ اے، تے تسی پرکھ سکدے او کہ ایہ
  $(\Diamond p \land \Diamond q) \lif \Diamond (p \land q)$ دا ردّی نمونہ اے۔
  \begin{figure}
  \begin{center}
    \begin{tikzpicture}[modal]
      \node[world] (w1) [label={[align=right]right:\mFalse{p}\\ \mFalse{q}}]{$1$}; 
      \node[world] (w2) [label={[align=right]right:\mTrue{p}\\ \mFalse{q}},
        above left=of w1]{$1.1$}; 
      \node[world] (w3) [label={[align=right]right:\mFalse{p}\\ \mTrue{q}},
        above right=of w1] {$1.2$};
      \draw[->] (w1) to (w2);
      \draw[->] (w1) to (w3);
    \end{tikzpicture}
  \end{center}
  \caption{$(\Diamond p \land \Diamond q) \lif
    \Diamond (p \land q)$ دا اک ردّی نمونہ۔}
  \ollabel{fig:counter-Diamond}
\end{figure}
\end{ex}
}
\endgroup
% /OLP-0469

\OLEndChapterHook
\endgroup
% /OLP-0460

% OLP-0470 content/normal-modal-logic/sequent-calculus/sequent-calculus.tex
\begingroup
\def\olfilename{sequent-calculus}
\expandafter\def\csname ol@sectioncs\endcsname{section}

\olchapter{nml}{seq}{امکانی سیکوئنٹ حساب}
%\tagfalse{prvDiamond}
\begin{editorial}
  امکانی منطق دے سیکوئنٹ حساباں بارے کچا باب۔ ہور مثالاں تے صحت مندی
  تے تمامیت دے ثبوتاں دی لوڑ اے۔
\end{editorial}

% OLP-0471 content/normal-modal-logic/sequent-calculus/introduction.tex
\begingroup
\def\olfilename{introduction}
\expandafter\def\csname ol@sectioncs\endcsname{section}

\olfileid{nml}{seq}{int}

\olsection{جان پہچان}

قضیاتی منطق دا سیکوئنٹ حساب وادھو قاعدیاں نال ودھایا جا سکدا اے،
جیہڑے~\iftag{prvBox}{$\Box$\iftag{prvDiamond}{ تے~}{}}{}%
\iftag{prvDiamond}{$\Diamond$}{} لئی نیں۔ مثلاً~$\Log{K}$ لئی
\Log{LK} دے نال ایہ وی نیں:
\[
  \iftag{prvBox}{\iftag{prvDiamond}{% <> and [] primitive
    \Axiom$\Gamma \fCenter \Delta, !A$
    \RightLabel{$\Box$}
    \UnaryInf$\Box\Gamma \fCenter \Diamond\Delta, \Box!A$
    \DisplayProof  
    \qquad
    \Axiom$!A, \Gamma \fCenter \Delta$
    \RightLabel{$\Diamond$}
    \UnaryInf$\Diamond!A, \Box\Gamma \fCenter \Diamond\Delta$
    \DisplayProof
}{% only Box primitive
\Axiom$\Gamma \fCenter !A$
\RightLabel{$\Box$}
\UnaryInf$\Box\Gamma \fCenter \Box!A$
\DisplayProof
}}{% only <> primitive
  \Axiom$!A \fCenter \Delta$
  \RightLabel{$\Diamond$}
  \UnaryInf$\Diamond!A \fCenter \Diamond\Delta$
  \DisplayProof
}
\]
~$\Log{K}$ نوں ودھاؤن والیاں منطقاں لئی ہور قاعدے وی لانے پین گے۔

ہر امکانی منطق دا ایسا سیکوئنٹ حساب نہیں۔ ایہ گل~$\Log{S5}$ اُتے
وی لاگو اے، جیہدی معنیات سادی اے (رسائی دے تعلق ورتے بغیر وی اوہدی
تعریف ہو سکدی اے)۔ اوہدے لئی~$\Log{LK}$ توں بنیا کوئی ایسا سیکوئنٹ
حساب معلوم نہیں جیہڑا قاعدہ~\Cut توں بغیر مکمل ہووے۔ پر اوہدا اک مکمل
\emph{ہائپر سیکوئنٹ} حساب اے جیہدے وچ کٹ دا قاعدہ نہیں۔
\endgroup
% /OLP-0471
% OLP-0472 content/normal-modal-logic/sequent-calculus/rules-for-K.tex
\begingroup
\def\olfilename{rules-for-K}
\expandafter\def\csname ol@sectioncs\endcsname{section}

\olfileid{nml}{seq}{rul}

\olsection{\Log{K} دے قاعدے}

عام قضیاتی رابطیاں دے قاعدے اوہی نیں جیہڑے عام سیکوئنٹ حساب~$\Log{LK}$
دے نیں۔ مسلّمے وی اوہی نیں۔ شکل $!A \Sequent !A$ والا ہر سیکوئنٹ
اک مسلّمہ اے۔

امکانی عامل\iftag{prvBox}{\iftag{prvDiamond}{~$\Box$
تے}{~$\Box$}}{}\iftag{prvDiamond}{~$\Diamond$}{} لئی ایہ وادھو
\iftag{notprvBox,notprvDiamond}{قاعدہ اے}{قاعدے نیں}:
\[
  \iftag{prvBox}{\iftag{prvDiamond}{% <> and [] primitive
    \Axiom$\Gamma \fCenter \Delta, !A$
    \RightLabel{$\Box$}
    \UnaryInf$\Box\Gamma \fCenter \Diamond\Delta, \Box!A$
    \DisplayProof  
    \qquad
    \Axiom$!A, \Gamma \fCenter \Delta$
    \RightLabel{$\Diamond$}
    \UnaryInf$\Diamond!A, \Box\Gamma \fCenter \Diamond\Delta$
    \DisplayProof
}{% only Box primitive
\Axiom$\Gamma \fCenter !A$
\RightLabel{$\Box$}
\UnaryInf$\Box\Gamma \fCenter \Box!A$
\DisplayProof
}}{% only <> primitive
  \Axiom$!A \fCenter \Delta$
  \RightLabel{$\Diamond$}
  \UnaryInf$\Diamond!A \fCenter \Diamond\Delta$
  \DisplayProof
}
\]
ایتھے \iftag{prvBox}{$\Box\Gamma$ اوہناں !!{formula}s دی لڑی اے جیہڑے
~$\Gamma$ توں ملدے نیں جدوں $\Box$، ہر !!{formula} دے اگے لایا جائے
جیہڑا~$\Gamma$ وچ اے\iftag{prvDiamond}{ تے
}{}}{}\iftag{prvDiamond}{$\Diamond\Delta$ اوہناں !!{formula}s دی لڑی اے
جیہڑے~$\Delta$ توں ملدے نیں جدوں $\Diamond$، ہر !!{formula} دے اگے
لایا جائے جیہڑا~$\Delta$ وچ اے}{}۔ \iftag{notprvDiamond}{قاعدے $\Box$
دے مقدمے دے سجے پاسے ودھ توں ودھ اک !!{formula}~$!A$ ہونا چاہیدا اے۔ }{}%
\iftag{notprvBox}{قاعدے $\Diamond$ دے مقدمے دے کھبے پاسے ودھ توں ودھ
اک !!{formula}~$!A$ ہونا چاہیدا اے۔ }{}%
\iftag{prvBox}{$\Gamma$\iftag{prvDiamond}{~تے~}{}}{}\iftag{prvDiamond}{$\Delta$}{}
دا خالی ہونا ممکن اے۔ اوس صورت وچ نتیجے والے سیکوئنٹ دا متعلقہ حصہ
\iftag{prvBox}{$\Box\Gamma$\iftag{prvDiamond}{~تے~}{}}{}\iftag{prvDiamond}{$\Diamond\Delta$}{}
وی خالی اے۔

ایہ شرط ضروری اے کہ
\iftag{prvBox}{$\Box$ سجے پاسے\iftag{prvDiamond}{ تے
}{}}{}\iftag{prvDiamond}{$\Diamond$ کھبے پاسے}{} صرف اکّو
!!{formula}~$!A$ اُتے لایا جائے۔ جے اسیں
\iftag{prvBox}{$\Box$ سجے پاسے دے کِنّے وی !!{formula}s اُتے لاؤن دی
اجازت دیئیے\iftag{prvDiamond}{ یا }{}}{}%
\iftag{prvDiamond}{$\Diamond$ کھبے پاسے دے کِنّے وی !!{formula}s اُتے
لاؤن دی اجازت دیئیے}{}، تاں اسیں ایہ !!{derive} کر سکدے آں:
\[
    \iftag{prvBox}{
      \Axiom$!A \fCenter !A$
      \RightLabel{\RightR{\lnot}}
      \UnaryInf$\fCenter !A, \lnot !A$
      \RightLabel{$\Box*$}
      \UnaryInf$\fCenter \Box!A, \Box\lnot !A$
      \doubleLine
      \RightLabel{\RightR{\lor}}
      \UnaryInf$\fCenter \Box!A \lor \Box \lnot !A$
      \DisplayProof}{}
    \iftag{notprvBox,notprvDiamond}{}{\qquad}
    \iftag{prvDiamond}{
      \Axiom$!A \fCenter !A$
      \RightLabel{\LeftR{\lnot}}
      \UnaryInf$\lnot !A, !A \fCenter$
      \RightLabel{$\Diamond*$}
      \UnaryInf$\Diamond\lnot !A,\Diamond!A \fCenter $
      \RightLabel{\RightR{\lnot}}
      \UnaryInf$\Diamond!A \fCenter \lnot\Diamond\lnot !A$
      \RightLabel{\RightR{\lif}}
      \UnaryInf$ \fCenter \Diamond!A \lif \lnot\Diamond\lnot !A$
      \DisplayProof}{}
  \]
پر \iftag{prvBox}{$\Box!A \lor \Box \lnot !A$\iftag{prvDiamond}{ تے
$\Diamond!A \lif \lnot\Diamond\lnot !A$}{}}{$\Diamond!A \lif
\lnot\Diamond\lnot !A$} دا منطقی اعتبار~$\Log{K}$ وچ نہیں۔

جے اسیں مقدمے وچ~$!A$ توں علاوہ پاسے والے فارمولیاں دی اجازت
دیئیے، تے \iftag{prvBox}{قاعدے $\Box$ نوں $\Box$ سجے پاسے صرف~$!A$
اُتے لاؤن دی اجازت دیئیے\iftag{prvDiamond}{، یا }{}}{}%
\iftag{prvDiamond}{قاعدے $\Diamond$ نوں $\Diamond$ کھبے پاسے صرف~$!A$
اُتے لاؤن دی اجازت دیئیے}{} (پر پاسے والے فارمولیاں نوں جیویں دا تیویں
رکھیے)، تاں اسیں ایہ !!{derive} کر سکدے آں:
\[
  \iftag{prvBox}{
    \Axiom$!A \fCenter !A$
    \RightLabel{\RightR{\lnot}}
    \UnaryInf$\fCenter !A, \lnot !A$
    \RightLabel{\RightR{\Exchange}}
    \UnaryInf$\fCenter \lnot !A, !A$
    \RightLabel{$\Box*$}
    \UnaryInf$\fCenter \lnot!A, \Box !A$
    \doubleLine
    \RightLabel{\RightR{\lor}}
    \UnaryInf$\fCenter \lnot!A \lor \Box !A$
    \DisplayProof}{}
  \iftag{notprvBox,notprvDiamond}{}{\qquad}
  \iftag{prvDiamond}{
    \Axiom$!A \fCenter !A$
    \RightLabel{\LeftR{\lnot}}
    \UnaryInf$\lnot !A, !A \fCenter$
    \RightLabel{$\Diamond*$}
    \UnaryInf$\Diamond\lnot !A, !A \fCenter $
    \RightLabel{\RightR{\lnot}}
    \UnaryInf$!A \fCenter \lnot\Diamond\lnot !A$
    \RightLabel{\RightR{\lif}}
    \UnaryInf$ \fCenter !A \lif \lnot\Diamond\lnot !A$
    \DisplayProof}{}
\]
پر \iftag{prvBox}{$\lnot!A \lor \Box !A$ (جیہڑا $!A
\lif \Box!A$ دے ہم ارز اے)\iftag{prvDiamond}{ تے $!A \lif \lnot\Diamond\lnot !A$}{}}{$!A \lif \lnot\Diamond\lnot !A$}
دا منطقی اعتبار~$\Log{K}$ وچ نہیں۔
\endgroup
% /OLP-0472
% OLP-0473 content/normal-modal-logic/sequent-calculus/proofs-in-K.tex
\begingroup
\def\olfilename{proofs-in-K}
\expandafter\def\csname ol@sectioncs\endcsname{section}

\olfileid{nml}{seq}{prk}

\olsection{\Log{K} لئی سیکوئنٹ \usetoken{P}{derivation}}

\iftag{prvBox}{

\begin{ex}
  اسیں سیکوئنٹ حساب دی اک !!{derivation} دیندے آں جیہڑی وکھاؤندی اے کہ $\Proves (\Box!A \land \Box!B)
  \lif \Box (!A \land !B)$.
  \begin{prooftree}
    \Axiom$!A \fCenter !A$
    \doubleLine
    \UnaryInf$!B, !A \fCenter !A$   
    \Axiom$!B \fCenter !B$
    \doubleLine
    \UnaryInf$!B, !A \fCenter !B$   
    \RightLabel{\RightR{\land}}
    \BinaryInf$!B, !A \fCenter !A \land !B$    
    \RightLabel{$\Box$}
    \UnaryInf$\Box!B, \Box!A \fCenter \Box
    (!A \land !B)$
    \RightLabel{\LeftR{\land}}
    \UnaryInf$\Box!A \land \Box!B, \Box!A \fCenter \Box
    (!A \land !B)$
    \RightLabel{\LeftR{\Exchange}}
    \UnaryInf$\Box!A, \Box!A \land \Box!B \fCenter \Box
    (!A \land !B)$
    \RightLabel{\LeftR{\land}}
    \UnaryInf$\Box!A \land \Box!B, \Box!A \land \Box!B \fCenter \Box (!A \land !B)$
    \RightLabel{\LeftR{\Contraction}}
    \UnaryInf$\Box!A \land \Box!B \fCenter \Box (!A \land !B)$
    \RightLabel{\RightR{\lif}}
    \UnaryInf$\fCenter (\Box!A \land \Box!B)
    \lif \Box (!A \land !B)$
  \end{prooftree}
\end{ex}
}{}

\iftag{prvDiamond}{
  \begin{ex}
    اسیں سیکوئنٹ حساب دی اک !!{derivation} دیندے آں جیہڑی وکھاؤندی اے کہ $\Proves \Diamond(!A
    \lor !B) \lif (\Diamond !A \lor \Diamond!B)$.
    \begin{prooftree}
      \Axiom$!A \fCenter !A$
      \doubleLine
      \UnaryInf$!A \fCenter !A, !B$   
      \Axiom$!B \fCenter !B$
      \doubleLine
      \UnaryInf$!B \fCenter !A, !B$   
      \RightLabel{\LeftR{\lor}}
      \BinaryInf$!A \lor !B \fCenter !A, !B$    
      \RightLabel{$\Diamond$}
      \UnaryInf$\Diamond(!A \lor !B) \fCenter 
      \Diamond !A, \Diamond!B$
      \RightLabel{\RightR{\lor}}
      \UnaryInf$\Diamond(!A \lor !B) \fCenter 
      \Diamond !A, \Diamond !A \lor \Diamond!B$
      \RightLabel{\RightR{\Exchange}}
      \UnaryInf$\Diamond(!A \lor !B) \fCenter \Diamond !A \lor \Diamond!B, \Diamond !A$
      \RightLabel{\RightR{\lor}}
      \UnaryInf$\Diamond(!A \lor !B) \fCenter 
      \Diamond !A \lor \Diamond!B, \Diamond !A \lor \Diamond!B$
      \RightLabel{\RightR{\Contraction}}
      \UnaryInf$\Diamond(!A
      \lor !B) \fCenter \Diamond !A \lor \Diamond!B$
      \RightLabel{\RightR{\lif}}
      \UnaryInf$\fCenter \Diamond(!A
      \lor !B) \lif (\Diamond !A \lor \Diamond!B)$
    \end{prooftree}
  \end{ex}
}{}

\iftag{notprvBox,notprvDiamond}{}{
  ~\Dual دا !!a{derivation} ہِیٹھ اے۔
  \begin{prooftree}
    \Axiom$!A \fCenter !A$
    \RightLabel{\RightR{\lnot}}
    \UnaryInf$\lnot !A, !A \fCenter $
    \RightLabel{$\Diamond$}
    \UnaryInf$\Diamond\lnot !A, \Box !A \fCenter $
    \RightLabel{\RightR{\lnot}}
    \UnaryInf$\Box !A \fCenter \lnot\Diamond\lnot !A$
    \RightLabel{\RightR{\lif}}
    \UnaryInf$\fCenter \Box !A \lif \lnot\Diamond\lnot !A$
    \Axiom$!A \fCenter !A$
    \RightLabel{\RightR{\lnot}}
    \UnaryInf$\fCenter !A, \lnot !A $
    \RightLabel{\RightR{\Exchange}}
    \UnaryInf$\fCenter \lnot !A, !A $
    \RightLabel{$\Box$}
    \UnaryInf$\fCenter \Diamond\lnot !A, \Box !A$
    \RightLabel{\RightR{\Exchange}}
    \UnaryInf$\fCenter \Box !A, \Diamond\lnot !A$
    \RightLabel{\RightR{\lnot}}
    \UnaryInf$\lnot\Diamond\lnot !A \fCenter \Box !A$
    \RightLabel{\RightR{\lif}}
    \UnaryInf$\fCenter \lnot \Diamond \lnot !A \lif \Box !A$
    \RightLabel{\RightR{\land}}
    \BinaryInf$\fCenter \Box !A \liff \lnot\Diamond\lnot !A$ 
  \end{prooftree} 

% مشقاں اگے پائیاں جان تاں ایہ کم نہیں کردا۔
% \begin{prob}
%     $\Diamond!A \liff \lnot\Box\lnot !A$ دا !!a{derivation}
%     ~$\Log{K}$ وچ دیو۔
% \end{prob}
}

\begin{prob}
  ~$\Log{K}$ وچ ہِیٹھلے !!{formula}s لئی سیکوئنٹ حساب دے ثبوت لبھو:
  \begin{enumerate}
    \item $\Box \lnot p \lif \Box(p \lif q)$
    \item $(\Box p \lor \Box q) \lif \Box(p \lor q)$
    \item $\Diamond p \lif \Diamond(p \lor q)$
    \item $\Box(p \land q) \lif \Box p$
  \end{enumerate}
\end{prob}
\endgroup
% /OLP-0473
%\olimport{soundness}
% OLP-0474 content/normal-modal-logic/sequent-calculus/more-rules.tex
\begingroup
\def\olfilename{more-rules}
\expandafter\def\csname ol@sectioncs\endcsname{section}

\olfileid{nml}{seq}{mru}

\olsection{ہور رسائی تعلقاں لئی قاعدے}

خاص رسائی تعلقاں توں متعین ہوندیاں منطقاں لئی اسیں
\olref{tab:more-rules} وچ دتے وادھو قاعدے ویکھدے آں۔

\iftag{prvBox}{\iftag{prvDiamond}{% [] and <> primitive
\begin{table}
  \begin{center}
    \def\arraystretch{3}
    \begin{tabular}{cc}
    \hline
      \Axiom$!A, \Gamma \fCenter \Delta$
      \RightLabel{T$\Box$}
      \UnaryInf$\Box!A, \Gamma \fCenter \Delta$
      \DisplayProof
    &
      \Axiom$\Gamma \fCenter \Delta, !A$
      \RightLabel{T$\Diamond$}
      \UnaryInf$\Gamma \fCenter \Delta, \Diamond!A$
      \DisplayProof
    \\[1ex]
%    \hline
    \multicolumn{2}{c}{
      \Axiom$\Gamma \fCenter \Delta$
      \RightLabel{D}
      \UnaryInf$\Box\Gamma \fCenter \Diamond\Delta$
      \DisplayProof}
    \\[1ex]
%    \hline
      \Axiom$\Gamma, \Diamond\Pi \fCenter \Box\Delta, \Lambda, !A$
      \RightLabel{B$\Box$}
      \UnaryInf$\Box\Gamma, \Pi \fCenter \Delta, \Diamond\Lambda, \Box!A$
      \DisplayProof
      &
      \Axiom$!A, \Diamond\Gamma, \Pi \fCenter \Box\Lambda, \Delta$
      \RightLabel{B$\Diamond$}
      \UnaryInf$\Diamond!A, \Gamma, \Box\Pi \fCenter \Lambda, \Diamond\Delta$
      \DisplayProof
    \\[1ex]
%    \hline
      \Axiom$\Box\Gamma \fCenter \Diamond\Delta, !A$
      \RightLabel{4$\Box$}
      \UnaryInf$\Box\Gamma \fCenter \Diamond\Delta, \Box!A$
      \DisplayProof
      &
      \Axiom$!A, \Box\Gamma \fCenter \Diamond\Delta$
      \RightLabel{4$\Diamond$}
      \UnaryInf$\Diamond!A, \Box\Gamma \fCenter \Diamond\Delta$
      \DisplayProof
    \\[1ex]
%    \hline
      \Axiom$\Box\Gamma, \Diamond\Pi \fCenter \Box\Delta,
      \Diamond\Lambda, !A$
      \RightLabel{5$\Box$}
      \UnaryInf$\Box\Gamma, \Diamond\Pi \fCenter \Box\Delta,
      \Diamond\Lambda, \Box!A$
      \DisplayProof
&
      \Axiom$!A, \Diamond\Gamma, \Box\Pi \fCenter \Diamond\Delta, \Box\Lambda$
      \RightLabel{5$\Diamond$}
      \UnaryInf$\Diamond!A,\Diamond\Gamma,\Box\Pi \fCenter \Diamond\Delta,\Box\Lambda$
      \DisplayProof
    \\[1ex]
    \hline
    \end{tabular}
  \end{center}
  \caption{ہور امکانی قاعدے۔}
  \ollabel{tab:more-rules}
\end{table}
}{ % only [] primitive
\begin{table}
  \begin{center}
    \def\arraystretch{3}
    \begin{tabular}{cc}
    \hline
      \Axiom$!A,\Gamma \fCenter \Delta$
      \RightLabel{T$\Box$}
      \UnaryInf$\Box!A, \Gamma \fCenter \Delta$
      \DisplayProof
    &
    %\hline
      \Axiom$\Gamma \fCenter$
      \RightLabel{D$\Box$}
      \UnaryInf$\Box\Gamma \fCenter$
      \DisplayProof
    \\[1ex]
    %\hline
      \Axiom$\Gamma \fCenter \Box\Delta, !A$
      \RightLabel{B$\Box$}
      \UnaryInf$\Box\Gamma \fCenter \Delta, \Box!A$
      \DisplayProof
    &
    %\hline
      \Axiom$\Box\Gamma \fCenter !A$
      \RightLabel{4$\Box$}
      \UnaryInf$\Box\Gamma \fCenter \Box!A$
      \DisplayProof
    \\[1ex]
    %\hline
    \multicolumn{2}{c}{\Axiom$\Box\Gamma \fCenter \Box\Delta, !A$
      \RightLabel{5$\Box$}
      \UnaryInf$\Box\Gamma \fCenter \Box\Delta, \Box!A$
      \DisplayProof}
    \\[1ex]
    \hline
    \end{tabular}
  \end{center}
  \caption{ہور امکانی قاعدے۔}
  \ollabel{tab:more-rules}
\end{table}
}}{% only <> primitive
\begin{table}
  \begin{center}
    \def\arraystretch{3}
    \begin{tabular}{cc}
    \hline
      \Axiom$\Gamma \fCenter \Delta, !A$
      \RightLabel{T$\Diamond$}
      \UnaryInf$\Gamma \fCenter \Delta, \Diamond!A$
      \DisplayProof
    &
    %\hline
      \Axiom$\fCenter \Delta$
      \RightLabel{D$\Diamond$}
      \UnaryInf$\fCenter \Diamond\Delta$
      \DisplayProof
    \\[1ex]
    %\hline
      \Axiom$!A, \Diamond\Gamma \fCenter \Delta$
      \RightLabel{B$\Diamond$}
      \UnaryInf$\Diamond!A, \Gamma \fCenter \Diamond\Delta$
      \DisplayProof
    &
    %\hline
      \Axiom$!A \fCenter \Diamond\Delta$
      \RightLabel{4$\Diamond$}
      \UnaryInf$\Diamond!A \fCenter \Diamond\Delta$
      \DisplayProof
    \\[1ex]
    %\hline
    \multicolumn{2}{c}{
      \Axiom$!A, \Diamond\Gamma \fCenter \Diamond\Delta$
      \RightLabel{5$\Diamond$}
      \UnaryInf$\Diamond!A,\Diamond\Gamma \fCenter \Diamond\Delta$
      \DisplayProof}
    \\[1ex]
    \hline
    \end{tabular}
  \end{center}
  \caption{ہور امکانی قاعدے۔}
  \ollabel{tab:more-rules}
\end{table}
}

ایہ قاعدے شامل کرن نال ملدے نظام \olref{tab:logics-rules}
وچ دتیاں منطقاں لئی صحت مند تے مکمل نیں۔

\begin{table}
  \begin{center}
    \begin{tabular}{lll}
      \hline
      منطق & $R$ دی خاصیت \dots & قاعدے\\
      \hline
      $\Log{T} = \Log{KT}$ & انعکاسی & $\Box$,
      \iftag{prvBox}{T$\Box$}{}%
      \iftag{notprvBox,notprvDiamond}{}{, }%
      \iftag{prvDiamond}{T$\Diamond$}{}
      \\ \hline
      $\Log{D} = \Log{KD}$ & ہر جہان توں رسائی والا & $\Box$,
      \iftag{prvBox}{%
        \iftag{prvDiamond}{D}{D$\Box$}}{D$\Diamond$}
      \\ \hline
      $\Log{K4}$ & متعدی & $\Box$,
      \iftag{prvBox}{4$\Box$}{}% 
      \iftag{notprvBox,notprvDiamond}{}{, }%
      \iftag{prvDiamond}{4$\Diamond$}{}
      \\ \hline
      $\Log{B} = \Log{KTB}$ & انعکاسی, & $\Box$,
      \iftag{prvBox}{T$\Box$}{}%
      \iftag{notprvBox,notprvDiamond}{}{, }%
      \iftag{prvDiamond}{T$\Diamond$}{}\\
      & متقابل &
      \iftag{prvBox}{B$\Box$}{}%
      \iftag{notprvBox,notprvDiamond}{}{, }%
      \iftag{prvDiamond}{B$\Diamond$}{}
      \\ \hline
      $\Log{S4} = \Log{KT4}$ & انعکاسی, & $\Box$,
      \iftag{prvBox}{T$\Box$}{}%
      \iftag{notprvBox,notprvDiamond}{}{, }%
      \iftag{prvDiamond}{T$\Diamond$}{}\\
      & متعدی &
      \iftag{prvBox}{4$\Box$}{}%
      \iftag{notprvBox,notprvDiamond}{}{, }%
      \iftag{prvDiamond}{4$\Diamond$}{}
      \\ \hline
      $\Log{S5} = \Log{KT5}$ & انعکاسی, & $\Box$,
      \iftag{prvBox}{T$\Box$}{}%
      \iftag{notprvBox,notprvDiamond}{}{, }%
      \iftag{prvDiamond}{T$\Diamond$}{}\\
      & متعدی, &
      \iftag{prvBox}{5$\Box$}{}%
      \iftag{notprvBox,notprvDiamond}{}{, }%
      \iftag{prvDiamond}{5$\Diamond$}{}\\
      &  اقلیدسی &
      \\ \hline
    \end{tabular}
  \end{center}
  \caption{وکھ وکھ امکانی منطقاں لئی سیکوئنٹ قاعدے۔}
  \ollabel{tab:logics-rules}
\end{table}

\iftag{prvBox}{
\begin{ex}
  اسیں اک سیکوئنٹ !!{derivation} دیندے آں جیہڑی وکھاؤندی اے کہ $\Log{K4} \Proves \Ax{4}$، یعنی
  $\Box!A \lif \Box\Box!A$.
  \begin{prooftree}
    \Axiom$\Box!A \fCenter \Box!A$
    \RightLabel{4$\Box$}
    \UnaryInf$\Box!A \fCenter \Box\Box!A$
    \RightLabel{\RightR{\lif}}
    \UnaryInf$\fCenter \Box!A \lif \Box\Box!A$
  \end{prooftree}
\end{ex}
}{\iftag{prvDiamond}{
\begin{ex}
  اسیں اک سیکوئنٹ !!{derivation} دیندے آں جیہڑی وکھاؤندی اے کہ $\Log{K4} \Proves \Ax{4}$، یعنی
  $\Diamond\Diamond!A \lif \Diamond!A$.
  \begin{prooftree}
    \Axiom$\Diamond!A \fCenter \Diamond!A$
    \RightLabel{4$\Diamond$}
    \UnaryInf$\Diamond\Diamond!A \fCenter \Diamond!A$
    \RightLabel{\RightR{\lif}}
    \UnaryInf$\fCenter \Diamond\Diamond!A \lif \Diamond!A$
  \end{prooftree}
\end{ex}
}{}}

\iftag{prvBox}{\iftag{prvDiamond}{% <> and [] primitive
\begin{ex}
  اسیں اک سیکوئنٹ !!{derivation} دیندے آں جیہڑی وکھاؤندی اے کہ $\Log{S5} \Proves \Ax{5}$، یعنی
  $\Diamond!A \lif \Box\Diamond!A$.
  \begin{prooftree}
    \Axiom$\Diamond!A \fCenter \Diamond!A$
    \RightLabel{5$\Box$}
    \UnaryInf$\Diamond!A \fCenter \Box\Diamond!A$
    \RightLabel{\RightR{\lif}}
    \UnaryInf$\fCenter \Diamond!A \lif \Box\Diamond!A$
  \end{prooftree}
\end{ex}
}{% only Box primitive
\begin{ex}
  اسیں اک سیکوئنٹ !!{derivation} دیندے آں جیہڑی وکھاؤندی اے کہ $\Log{S5} \Proves \Ax{5}$، یعنی
  $\Diamond!A \lif \Box\Diamond!A$.
  \begin{prooftree}
    \Axiom$\Box\lnot!A \fCenter \Box\lnot!A$
    \RightLabel{\RightR{\lnot}}
    \UnaryInf$\fCenter \Box\lnot!A, \lnot\Box\lnot!A$
    \RightLabel{5$\Box$}
    \UnaryInf$ \fCenter \Box\lnot!A, \Box\lnot\Box\lnot!A$
    \RightLabel{\LeftR{\lnot}}
    \UnaryInf$\lnot\Box\lnot!A \fCenter \Box\lnot\Box\lnot!A$
    \RightLabel{\RightR{\lif}}
    \UnaryInf$\fCenter \lnot\Box\lnot!A \lif \Box\lnot\Box\lnot!A$
  \end{prooftree}
\end{ex}
}}{% only <> prmitive
\begin{ex}
  اسیں اک سیکوئنٹ !!{derivation} دیندے آں جیہڑی وکھاؤندی اے کہ $\Log{S5} \Proves \Ax{5}$، یعنی
  $\Diamond!A \lif \Box\Diamond!A$.
  \begin{prooftree}
    \Axiom$\Diamond!A \fCenter\Diamond!A$
    \RightLabel{\LeftR{\lnot}}
    \UnaryInf$\lnot\Diamond !A, \Diamond!A \fCenter $
    \RightLabel{5$\Diamond$}
    \UnaryInf$\Diamond\lnot\Diamond!A, \Diamond!A \fCenter $
    \RightLabel{\RightR{\lnot}}
    \UnaryInf$\Diamond!A \fCenter \lnot\Diamond\lnot\Diamond!A$
    \RightLabel{\RightR{\lif}}
    \UnaryInf$\fCenter \Diamond!A \lif \lnot\Diamond\lnot\Diamond!A$
  \end{prooftree}
\end{ex}
}
\begin{ex}
\Log{S5} دا سیکوئنٹ حساب \Cut{} قاعدے توں بغیر مکمل نہیں۔ مثلاً
\iftag{prvBox}{$\Diamond\Box !A \lif !A$}{$!A \lif
\Box\Diamond !A$}، جیہڑا~$\Log{S5}$ وچ معتبر اے، اوہدا~\Cut توں
بغیر کوئی ثبوت نہیں۔ ہِیٹھ اک !!a{derivation} اے جیہڑی~\Cut ورتدی اے:
\begin{prooftree}
\iftag{prvBox}{\iftag{prvDiamond}{% both [] and <> primitive
\Axiom$\Box !A \fCenter \Box!A$
\RightLabel{5$\Diamond$}
\UnaryInf$\Diamond\Box!A \fCenter \Box!A$

\Axiom$!A \fCenter !A$
\RightLabel{T$\Box$}
\UnaryInf$\Box!A \fCenter !A$
\RightLabel{\Cut}
\BinaryInf$\Diamond\Box!A \fCenter !A$
\RightLabel{\RightR{\lif}}
\UnaryInf$\fCenter \Diamond\Box!A \lif !A$
}{% only [] primitive
\Axiom$\Box !A \fCenter \Box!A$
\RightLabel{\RightR{\lnot}}
\UnaryInf$\fCenter \Box!A, \lnot\Box!A$
\RightLabel{5$\Box$}
\UnaryInf$\fCenter \Box!A, \Box\lnot\Box!A$
\RightLabel{\RightR{\Exchange}}
\UnaryInf$\fCenter \Box\lnot\Box!A, \Box!A$

\Axiom$!A \fCenter !A$
\RightLabel{T$\Box$}
\UnaryInf$\Box!A \fCenter !A$
\RightLabel{\Cut}
\BinaryInf$ \fCenter \Box\lnot\Box!A, !A$
\RightLabel{\LeftR{\lnot}}
\UnaryInf$\lnot\Box\lnot\Box!A \fCenter !A$
\RightLabel{\RightR{\lif}}
\UnaryInf$\fCenter \lnot\Box\lnot\Box!A \lif !A$
}}{% only <> primitive
\Axiom$!A \fCenter !A$
\RightLabel{T$\Diamond$}
\UnaryInf$!A \fCenter \Diamond!A$

\Axiom$\Diamond !A \fCenter \Diamond!A$
\RightLabel{\LeftR{\lnot}}
\UnaryInf$\lnot\Diamond!A , \Diamond!A \fCenter$
\RightLabel{5$\Diamond$}
\UnaryInf$\Diamond\lnot\Diamond!A , \Diamond!A\fCenter$
\RightLabel{\RightR{\Exchange}}
\UnaryInf$\Diamond!A, \Diamond\lnot\Diamond!A\fCenter $

\RightLabel{\Cut}
\BinaryInf$\Diamond\lnot\Diamond!A, !A\fCenter $
\RightLabel{\RightR{\lnot}}
\UnaryInf$!A \fCenter \lnot\Diamond\lnot\Diamond!A$
\RightLabel{\RightR{\lif}}
\UnaryInf$\fCenter !A \lif \lnot\Diamond\lnot\Diamond!A$
}
\end{prooftree}
\end{ex}

\begin{prob}
ہِیٹھلے بیان وکھاؤن والیاں سیکوئنٹ !!{derivation}s دیو:
  \begin{enumerate}
  \item $\Log{KT5} \Proves \Ax{B}$;
  \item $\Log{KT5} \Proves \Ax{4}$;
  \item $\Log{KDB4} \Proves \Ax{T}$;
  \item $\Log{KB4} \Proves \Ax{5}$;
  \item $\Log{KB5} \Proves \Ax{4}$;
  \item $\Log{KT} \Proves \Ax{D}$.
  \end{enumerate}
\end{prob}
\endgroup
% /OLP-0474
%\olimport{more-soundness}
%\olimport{hypersequents-S5}

\OLEndChapterHook
\endgroup
% /OLP-0470

\OLEndPartHook
\endgroup
% /OLP-0407

% OLP-0475 content/applied-modal-logic/applied-modal-logic.tex
\begingroup
\def\olfilename{applied-modal-logic}
\expandafter\def\csname ol@sectioncs\endcsname{section}

\olpart{aml}{امکانی منطق دیاں ورتوں}

\begin{editorial}
ایس حصے وچ امکانی منطق دیاں کجھ ورتوں بارے ازمائشی مسودے نیں، مثلاً
سمے دیاں تے علم دیاں منطقاں۔
\end{editorial}

% OLP-0476 content/applied-modal-logic/temporal-logic/temporal-logic.tex
\begingroup
\def\olfilename{temporal-logic}
\expandafter\def\csname ol@sectioncs\endcsname{section}

\olchapter{aml}{tl}{سمے دیاں منطقاں}

\begin{editorial}
ایہ باب سمے دیاں منطقاں بارے اے۔
\end{editorial}

% OLP-0477 content/applied-modal-logic/temporal-logic/introduction.tex
\begingroup
\def\olfilename{introduction}
\expandafter\def\csname ol@sectioncs\endcsname{section}

\olfileid{aml}{tl}{int}

\olsection{جان پہچان}

سمے دیاں منطقاں اوہناں گلاں دے بیان ویکھدیاں نیں جیہڑیاں اگے ہون
گیاں یا پہلاں ہو چکیاں نیں۔ سمے دی منطق دا مڈھ بنھن دا سیہرا آرتھر
پرائر نوں دتا جاندا اے۔ اوہ اینوں \emph{فعل دے زمانے دی منطق} آکھدا سی۔
ایتھے سمے دی منطق نوں سمجھاؤن لئی اسیں زیادہ تر پرائر دے اصل امکانی
طریقے اُتے چلاں گے۔ ایس طریقے وچ سمے دے عامل قضیاتی منطق دے بنیادی
ڈھانچے وچ شامل کیتے جاندے نیں، جیہڑا بیان عام طور اُتے فعل دے زمانے
توں بغیر لیندا اے۔

مثلاً قضیاتی منطق وچ میں بیزی ناں دے اک کُتے دی گل کر سکدا آں، جیہڑا
کدی بہندا اے تے کدی نہیں بہندا، جیویں کُتے عام کردے نیں۔ کلاسیکی منطق
وچ ایہ آکھنا تضاد ہووے گا کہ بیزی بیٹھا اے تے بیزی بیٹھا نہیں۔ پر صاف
اے کہ ایہ دونیں وکھ وکھ ویلے سچ ہو سکدے نیں۔ زبان وچ سمے دے عامل شامل
کر کے اسیں ایہ گل سوکھ نال آکھ سکدے آں۔ سمے دے عامل ایہ وی سمجھاؤندے
نیں کہ ایسے استنباط کیوں معتبر نیں جیویں «بیزی نوں انعام وچ کھان والی
کوئی شے یا اک گیند ملے گی» توں «بیزی نوں انعام وچ کھان والی کوئی شے
ملے گی یا بیزی نوں اک گیند ملے گی»۔

سمے دی منطق نال کئی فلسفیانہ سوال اُٹھدے نیں جیہڑے سانوں اک ڈھانچے
دی تھاں دوجا چُنن ول لے جا سکدے نیں۔ مثلاً آؤن والے سمے دا ممکن پر
غیر لازمی بیان اوہ ہوندا اے جیہڑا نہ لازمی اے نہ ناممکن۔ جے اسیں آکھیے
«رچرڈ بھلکے سودا لین دکان جائے گا»، تاں اسیں اجیہی گل آکھ رہے آں
جیہڑی ہجے نہیں ہوئی، تے جیہدی صداقتی قدر اُتے اختلاف ہو سکدا اے۔
ایس گل اُتے وی اختلاف ہو سکدا اے کہ اوہنوں کوئی صداقتی قدر \emph{دتی}
وی جا سکدی اے یا نہیں۔ جے اسیں پکے جبریت وادی ہوئیے، تاں شاید ایہ منّن
وچ کوئی اوکھ نہ ہووے کہ متعلقہ واقعہ ہوون توں پہلاں وی ایہ جملہ سچ اے
یا جھوٹ۔ ہو سکدا اے سانوں اوہدی صداقتی قدر دا ہجے پتا نہ ہووے۔
دوجے پاسے اسیں منّ سکدے آں کہ آؤن والا سماں واقعی کھلا اے، تے اوہدے
ممکن پر غیر لازمی بیاناں دیاں صداقتی قدراں طے نہیں ہوئیاں۔

سمے دی بناوٹ تے اوہدے سُبھاؤ بارے ایہناں وچوں کئی گلّاں ساڈے چنے
نمونیاں تے سمے دی منطق دے ڈھانچیاں وچ پہلے ای شامل ہوندیاں نیں۔
مثلاً اسیں اپنے آپ نوں پچھ سکدے آں کہ ساڈے نمونیاں وچ سماں اک سیدھی
لڑی ہووے، شاخاں کڈھے، یا چکر وچ چلے۔ سانوں ایہ وی چُننا پے سکدا اے
کہ ساڈے سمے دے نمونیاں دا مڈھ تے آخر ہووے یا نہیں، تے سماں وکھ وکھ
پلاں نال وکھایا جائے یا اک لگاتار پھیلاؤ وانگ۔
\endgroup
% /OLP-0477
% OLP-0478 content/applied-modal-logic/temporal-logic/temporal-logic-semantics.tex
\begingroup
\def\olfilename{temporal-logic-semantics}
\expandafter\def\csname ol@sectioncs\endcsname{section}

\olfileid{aml}{tl}{sem}

\olsection{سمے دی منطق دیاں معنیات}

\begin{defn}
سمے دی منطق دی بنیادی زبان وچ ایہ شامل نیں:
\begin{enumerate}
  \tagitem{prvFalse}{!!{falsity} لئی قضیاتی مستقل~$\lfalse$۔}{}
  \tagitem{prvTrue}{!!{truth} لئی قضیاتی مستقل~$\ltrue$۔}{}
  \item !!{propositional variable}s دا اک !!{denumerable}s سیٹ: $\Obj
    p_0$, $\Obj p_1$, $\Obj p_2$, \dots
  \item قضیاتی رابطے: \startycommalist
  \iftag{prvNot}{\ycomma $\lnot$ (نفی)}{}%
  \iftag{prvAnd}{\ycomma $\land$ (عطف)}{}%
  \iftag{prvOr}{\ycomma $\lor$ (فصل)}{}%
  \iftag{prvIf}{\ycomma $\lif$ (!!{conditional})}{}%
  \iftag{prvIff}{\ycomma $\liff$ (!!{biconditional})}{}۔
  \item لنگھے سمے دے عامل $\Ptemp$ تے $\Htemp$۔
  \item آؤن والے سمے دے عامل $\Ftemp$ تے $\Gtemp$۔
\end{enumerate}
\end{defn}

اگے اسیں ویکھاں گے کہ ہور قسماں دے امکانی عامل کِویں شامل ہو سکدے نیں۔

\begin{defn}
سمے دی زبان دے \emph{!!^{formula}s} استقرا نال ایویں متعین ہوندے نیں:
\begin{enumerate}
\tagitem{prvFalse}{$\lfalse$ اک ایٹمی !!{formula} اے۔}{}

\tagitem{prvTrue}{$\ltrue$ اک ایٹمی !!{formula} اے۔}{}

\item ہر قضیاتی متغیر $\Obj p_i$ اک (ایٹمی) !!{formula} اے۔

\tagitem{prvNot}{جے $!A$ اک !!a{formula} اے، تاں $\lnot !A$ وی
  اک !!a{formula} اے۔}{}

\tagitem{prvAnd}{جے $!A$ تے $!B$، !!{formula}s نیں، تاں $(!A \land
  !B)$ اک !!a{formula} اے۔}{}

\tagitem{prvOr}{جے $!A$ تے $!B$، !!{formula}s نیں، تاں $(!A \lor !B)$
  اک !!a{formula} اے۔}{}

\tagitem{prvIf}{جے $!A$ تے $!B$، !!{formula}s نیں، تاں $(!A \lif !B)$
  اک !!a{formula} اے۔}{}

\tagitem{prvIff}{جے $!A$ تے $!B$، !!{formula}s نیں، تاں $(!A \liff !B)$
  اک !!a{formula} اے۔}{}

\item جے $!A$ اک !!a{formula} اے، تاں $\Ptemp  !A$، $\Htemp  !A$،
  $F !A$ تے $\Gtemp  !A$ سارے !!{formula}s نیں۔

\tagitem{limitClause}{ہور کوئی شے !!a{formula} نہیں۔}{}
\end{enumerate}
\end{defn}

ہور قسماں دیاں مفہومی منطقاں وانگ، سمے دیاں منطقاں دیاں معنیات تعلقی نمونیاں راہیں دتیاں جاندیاں نیں۔

\begin{defn}
  سمے دی زبان دا اک \emph{نمونہ} ترتیبی تِکڑی
  $\mModel{M} = \tuple{T, \prec, V}$ اے، جتھے
  \begin{enumerate}
  \item $T$ اک غیر خالی سیٹ اے، جیہدے اجزا سمے دے نقطے نیں۔
  \item $\prec$، $T$ اُتے اک دو رکنی تعلق اے۔
  \item $V$ اوہ فنکشن اے جیہڑا ہر !!{propositional variable}~$p$ نوں سمے دے نقطیاں دا سیٹ $V(p)$ دیندا اے۔
  \end{enumerate}
  جدوں $t \prec t'$ ہووے، اسیں آکھدے آں کہ $t$،~$t'$ توں
  \emph{پہلاں اے}۔
  جدوں $t \in V(p)$ ہووے، اسیں آکھدے آں کہ $p$،~$t$ اُتے \emph{سچ اے}۔
\end{defn}

ہن لئی ویکھو کہ پہلاں ہون والے تعلق $\prec$ اُتے اسیں کوئی شرط
نہیں لائی۔ سو ہن ساڈے سمے دے نمونیاں دی بناوٹ اُتے کوئی پابندی نہیں۔
سماں اک سیدھی لڑی ہو سکدا اے، شاخاں کڈھ سکدا اے، چکر وچ چل سکدا اے
یا اوہدی کوئی وی ہور بناوٹ ہو سکدی اے۔

معیاری امکانی منطق وانگ، سمے دا ہر نمونہ طے کردا اے کہ اوہدے کہڑے
نقطیاں اُتے کہڑے !!{formula}s سچ نیں۔ تعلقی معنیات دی بنیادی گل لئی
اسیں اوہی لکھت ورتدے آں، «نمونہ $\mModel{M}$، !!{formula}~$!A$
نوں نقطے~$t$ اُتے سچ کردا اے»۔ ایہ تعلق استقرا نال متعین ہوندا اے،
تے سارے غیر امکانی عاملاں لئی معیاری امکانی منطق والا ای اے۔

\begin{defn}\ollabel{defn:tmodels}
  \emph{!!a{formula}~$!A$ دا~$t$ اُتے سچ ہونا}، نمونے~$\mModel M$
  وچ، نشانیاں وچ $\mSat{M}{!A}[t]$، استقرا نال ایویں متعین ہوندا اے:
  \begin{enumerate}
  \tagitem{prvFalse}{\indcase{!A}{\lfalse}{کدی نہیں کہ $\mSat{M}{\lfalse}[t]$}۔}{}
  \tagitem{prvTrue}{\indcase{!A}{\ltrue}{ہمیشہ $\mSat{M}{\ltrue}[t]$}۔}{}
  \item $\mSat{M}{p}[t]$ اُدوں تے صرف اُدوں جدوں $t \in V(p)$
  \tagitem{prvNot}{\indcase{!A}{\lnot !B}{$\mSat{M}{\indfrm}[t]$ اُدوں تے
    صرف اُدوں جدوں $\mSat/{M}{!B}[t]$}۔}{}
  \tagitem{prvAnd}{\indcase{!A}{(!B \land !C)}{$\mSat{M}{\indfrm}[t]$ اُدوں
    تے صرف اُدوں جدوں $\mSat{M}{!B}[t]$ تے $\mSat{M}{!C}[t]$}۔}{}
  \tagitem{prvOr}{\indcase{!A}{(!B \lor !C)}{$\mSat{M}{\indfrm}[t]$ اُدوں
    تے صرف اُدوں جدوں $\mSat{M}{!B}[t]$ یا $\mSat{M}{!C}[t]$}
    (یا دونیں)۔}{}
  \tagitem{prvIf}{\indcase{!A}{(!B \lif !C)}{$\mSat{M}{\indfrm}[t]$ اُدوں
    تے صرف اُدوں جدوں $\mSat/{M}{!B}[t]$ یا $\mSat{M}{!C}[t]$}۔}{}
  \tagitem{prvIff}{\indcase{!A}{(!B \liff !C)}{$\mSat{M}{\indfrm}[t]$ اُدوں
    تے صرف اُدوں جدوں یا تاں $\mSat{M}{!B}[t]$ تے $\mSat{M}{!C}[t]$
    دونیں ہون، یا نہ $\mSat{M}{!B}[t]$ ہووے تے نہ $\mSat{M}{!C}[t]$}۔}{}
  \item\ollabel{defn:sub:mmodels-p}
    \indcase{!A}{\Ptemp  !B}{$\mSat{M}{\indfrm}[t]$ اُدوں تے صرف اُدوں
    جدوں $\mSat{M}{!B}[t']$ کسے ایسے $t' \in T$ لئی ہووے جیہدے لئی $t' \prec t$}
  \item\ollabel{defn:sub:mmodels-h}
    \indcase{!A}{\Htemp  !B}{$\mSat{M}{\indfrm}[t]$ اُدوں تے صرف اُدوں
    جدوں $\mSat{M}{!B}[t']$ ہر ایسے $t' \in T$ لئی ہووے جیہدے لئی $t' \prec t$}
  \item\ollabel{defn:sub:mmodels-f}
    \indcase{!A}{\Ftemp  !B}{$\mSat{M}{\indfrm}[t]$ اُدوں تے صرف اُدوں
    جدوں $\mSat{M}{!B}[t']$ کسے ایسے $t' \in T$ لئی ہووے جیہدے لئی $t \prec t'$}
  \item\ollabel{defn:sub:mmodels-g}
    \indcase{!A}{\Gtemp  !B}{$\mSat{M}{\indfrm}[t]$ اُدوں تے صرف اُدوں
    جدوں $\mSat{M}{!B}[t']$ ہر ایسے $t' \in T$ لئی ہووے جیہدے لئی $t \prec t'$}
  \end{enumerate}
\end{defn}

معنیات توں تسی ویکھ سکدے او کہ عامل $\Ptemp$ تے~$\Htemp$ اک دوجے
دے دوہرے نیں، تے ایویں ای عامل $\Ftemp$ تے $\Gtemp$ وی۔ ایس لئی اسیں
$\Htemp  !A$ دی تعریف $\lnot \Ptemp \lnot !A$ دے طور اُتے کر سکدے آں،
تے $\Gtemp$ تے~$\Ftemp$ لئی وی ایویں۔
\endgroup
% /OLP-0478
% OLP-0479 content/applied-modal-logic/temporal-logic/properties-accessibility.tex
\begingroup
\def\olfilename{properties-accessibility}
\expandafter\def\csname ol@sectioncs\endcsname{section}

\olfileid{aml}{tl}{acc}

\olsection{سمے دے فریماں دیاں خاصیتاں}

ساڈے سمے دے نمونے تعلق~$\prec$ اُتے کوئی شرط نہیں لاؤندے، ایس لئی
ساڈے جانے پہچانے مسلّمیاں وچوں صرف~$K$، یا اوہدے برابر دے
$K_{\Gtemp}$ تے~$K_{\Htemp}$، سارے نمونیاں وچ لاگو ہوندے نیں:
\begin{align*}
\tag{$K_{\Gtemp}$}  \Gtemp (p \to q) & \to (\Gtemp p \to \Gtemp q)\\
\tag{$K_{\Htemp}$}  \Htemp (p \to q) & \to (\Htemp p \to \Htemp q)
\end{align*}

پر جے اسیں چاہندے آں کہ ساڈے نمونے سماں وکھاؤن اُتے ہور سخت شرطاں لاؤن،
مثلاً ایہ یقینی بنائیے کہ $\prec$ اک سیدھی ترتیب اے، تاں ساڈے
نمونیاں وچ ہور بیان وی معتبر ہون گے۔

\begin{table}[t]
  \begin{tabular}{| p{.48\textwidth} || p{.48\textwidth} |}
    \hline
    {\emph{جے $\prec$ \dots اے}} & {\emph{تاں \dots،~$\mModel{M}$ وچ سچ اے}} \\
    \hline \hline
    \emph{متعدی}: \newline
    $\forall u \forall v \forall w ((u \prec v \land v \prec w) \lif u \prec w)$ & 
    $\Ftemp \Ftemp p \lif \Ftemp p$  \\
    \hline 
    \emph{سیدھی ترتیب}: \newline
    $\forall w \forall v (w \prec v \lor w = v \lor v \prec w)$ &  
    $(\Ftemp \Ptemp  p \lor \Ptemp \Ftemp p) \lif (\Ptemp p \lor p \lor \Ftemp p)$\\
    \hline
    \emph{گنجان}: \newline
    $\forall w \forall v (w \prec v \to \exists u(w \prec u \land u \prec v))$ &  
    $\Ftemp p \lif \Ftemp \Ftemp p$ \\
    \hline
    \emph{لنگھے پاسے بے حد}: \newline
    $\forall w \exists v( v \prec w)$ &  
    $\Htemp p \to \Ptemp p$ \\
    \hline
    \emph{آؤن والے پاسے بے حد}: \newline
    $\forall w \exists v( w \prec v)$ &  
    $\Gtemp p \to \Ftemp p$ \\
    \hline
  \end{tabular}
  \caption{سمے دے فریماں دیاں کجھ خاصیتاں تے اوہناں دے برابر فارمولے۔}
  \ollabel{tab:correspondence}
\end{table} 

\olref{tab:correspondence} دیاں کئی خاصیتاں سمے نوں وکھاؤن والے
نمونے لئی چنگیاں لگ سکدیاں نیں۔ پر ایہ گل دھیان وچ رکھو کہ بھاویں
اسیں تعلق $\prec$ اُتے من چاہیاں شرطاں لا سکدے آں، ساریاں شرطاں دے
برابر ایسے !!{formula}s نہیں ہوندے جیہڑے سمے دی منطق دی زبان وچ لکھے
جا سکن۔ مثلاً غیر انعکاسی ہون دی شرط، یعنی
$\forall w \lnot (w \prec w)$، دے برابر سمے دی منطق دا کوئی فارمولا نہیں۔
\endgroup
% /OLP-0479
% OLP-0480 content/applied-modal-logic/temporal-logic/extra-temporal-operators.tex
\begingroup
\def\olfilename{extra-temporal-operators}
\expandafter\def\csname ol@sectioncs\endcsname{section}

\olfileid{aml}{tl}{ext}

\olsection{سمے دی منطق لئی ہور عامل}

لنگھے تے آؤن والے سمے دے اک رکنی عاملاں توں علاوہ، سمے دیاں منطقاں
وچ کدی کدی دو رکنی عامل $\Since $ تے~$\Until$ وی شامل ہوندے نیں،
جیہڑے «جدوں توں» تے «جدوں تائیں» دے نشان نیں۔ ایس لئی سمے دی منطق
دی زبان وچ $\Since $ تے~$\Until$ شامل کرنے پین گے، تے سمے دے
!!{formula} دی تعریف وچ ایہ شق وی لاؤنی پئے گی:

\begin{itemize}
\item[] جے $!A$ تے $!B$، !!{formula}s نیں، تاں $(\Since  !A !B)$ تے
  $(\Until !A !B)$ دونیں !!{formula}s نیں۔
\end{itemize}

فیر ایہناں عاملاں دیاں معنیات ایویں دتیاں جاندیاں نیں:

\begin{defn}\ollabel{defn:since-until}
  نمونے وچ \emph{!!a{formula}~$!A$ دا~$t$ اُتے سچ ہونا}،~$\mModel M$
  لئی:
  \begin{enumerate}
  \item\ollabel{defn:sub:mmodels-since} \indcase{!A}{\Since  !B
    !C}{$\mSat{M}{\indfrm}[t]$ اُدوں تے صرف اُدوں جدوں
    $\mSat{M}{!B}[t']$ کسے ایسے $t' \in
    T$ لئی ہووے جیہدے لئی $t' \prec t$، تے ہر ایسے $s$ لئی جیہدے لئی
    $t' \prec s \prec t$، $\mSat{M}{!C}[s]$ ہووے}
  \item\ollabel{defn:sub:mmodels-until} \indcase{!A}{\Until !B
    !C}{$\mSat{M}{\indfrm}[t]$ اُدوں تے صرف اُدوں جدوں
    $\mSat{M}{!B}[t']$ کسے ایسے $t' \in
    T$ لئی ہووے جیہدے لئی $t \prec t'$، تے ہر ایسے $s$ لئی جیہدے لئی
    $t \prec s \prec t'$، $\mSat{M}{!C}[s]$ ہووے}
  \end{enumerate}
\end{defn}

$\Since  !B !C$ نوں سِدھی طرح ایویں پڑھ سکدے آں، «جدوں توں
$!B$ سچ سی، اوہدوں توں~$!C$ سچ رہیا اے۔» تے $\Until !B !C$ نوں
ایویں پڑھ سکدے آں، «جدوں تائیں $!B$ سچ نہیں ہو جاندا، اوہدوں تائیں
~$!C$ سچ رہے گا۔»
\endgroup
% /OLP-0480
% OLP-0481 content/applied-modal-logic/temporal-logic/possible-histories.tex
\begingroup
\def\olfilename{possible-histories}
\expandafter\def\csname ol@sectioncs\endcsname{section}

\olfileid{aml}{tl}{poss}

\olsection{ممکن تاریخاں}

سمے دی منطق دے تعلقی نمونے جیہڑے اسیں ورتدے آئے آں، بہت لچک
رکھدے نیں، کیوں جے رسائی دے تعلق اُتے کوئی پابندی لاؤن دی لوڑ نہیں۔
ایس لئی سمے دے نمونے لنگھے تے آؤن والے سمے وچ شاخاں کڈھ سکدے نیں۔
پر اسیں شاخاں کڈھن دا ہور «امکانی» خیال وی ویکھ سکدے آں، جیہدے وچ
واقعات دیاں لڑیاں نوں ممکن تاریخاں سمجھدے آں۔ ایس لئی ساڈی زبان
بدلنی ضروری نہیں، بھاویں اسیں اپنے «عام» امکانی عامل $\Box$ تے~$\Diamond$
وی پا سکدے آں۔ اسیں علم دے رسائی تعلق شامل کرن بارے وی سوچ سکدے
آں، تاں جو سمے دے نال جانن والیاں دے علم وچ آؤن والیاں تبدیلیاں
وکھائیاں جا سکن۔

\begin{defn}
  سمے دی زبان لئی اک \emph{ممکن تاریخاں دا نمونہ} ترتیبی تِکڑی
  $\mModel{M} = \tuple{T, C, V}$ اے، جتھے
  \begin{enumerate}
    \item $T$ اک غیر خالی سیٹ اے، جیہدے اجزا نوں سمے دیاں حالتاں
      سمجھیا جاندا اے۔
    \item $C$ نظام دیاں حسابی راہواں، یا \emph{ممکن تاریخاں} دا سیٹ اے۔
      دوجے لفظاں وچ،~$C$ لڑیاں~$\sigma$ دا سیٹ اے جیہناں دے
      اجزا حالتاں $s_1$، $s_2$،~$s_3$، \dots نیں، جتھے ہر $s_i \in T$۔
    \item $V$ اوہ فنکشن اے جیہڑا ہر !!{propositional variable}~$p$ نوں سمے دے نقطیاں دا سیٹ~$V(p)$ دیندا اے۔
  \end{enumerate}
  سوکھ لئی اسیں عام طور اُتے ایہ وی منّاں گے کہ جے کوئی تاریخ~$C$
  وچ اے، تاں اوہدے مڈھوں حالتاں ہٹا کے بچدیاں ساریاں پچھلیاں لڑیاں
  وی اوہدے وچ نیں۔ مثلاً جے $s_1$، $s_2$،~$s_3$،~$C$ وچ اک لڑی اے،
  تاں $s_2$، $s_3$ تے~$s_3$ وی اوہدے وچ نیں۔ جدوں دو حالتاں $s_i$
  تے~$s_j$ کسے لڑی~$\sigma$ وچ آؤندیاں نیں، تاں اسیں آکھدے آں
  کہ $s_i \prec_\sigma s_j$ جدوں $i < j$۔
  جدوں $t \in V(p)$ ہووے، اسیں آکھدے آں کہ $p$،~$t$ اُتے \emph{سچ اے}۔
\end{defn}

ایتھے کم دی تبدیلی ایہ اے کہ جدوں اسیں نمونے وچ !!a{formula} دا
سچ ہونا سمے دے نقطے~$t$ اُتے~$\mModel M$ لئی پرکھدے آں، تاں
ایہ پرکھ اوس تاریخ~$\sigma$ دے لحاظ نال کردے آں جیہدے وچ $t$
اک حالت وانگ آؤندا اے۔ !!{propositional variable}s یا صداقتی فنکشن والے رابطیاں دیاں معنیات بدلنی ضروری
نہیں۔ ایہ ساریاں اوہو جیہیاں نیں جیویں \olref[sem]{defn:tmodels}
وچ سن، کیوں جے ایہناں وچوں کوئی وی~$\sigma$ دا حوالہ نہیں دیندی۔
پر ہُن اسیں آؤن والے سمے دا عامل~$\Ftemp$ مُڑ متعین کردے آں تے
عامل~$\Diamond$ شامل کردے آں، ایہناں تاریخاں دے لحاظ نال۔

\begin{defn}\ollabel{defn:phmodels}
  \emph{!!a{formula}~$!A$ دا~$t, \sigma$ اُتے سچ ہونا}، نمونے~$\mModel M$
  وچ، نشانیاں وچ $\mSat{M}{!A}[t, \sigma]$، ایویں اے:
  \begin{enumerate}
  \item\ollabel{defn:sub:phmodels-f} \indcase{!A}{\Ftemp
    !B}{$\mSat{M}{\indfrm}[t, \sigma]$ اُدوں تے صرف اُدوں جدوں
    $\mSat{M}{!B}[t', \sigma]$ کسے ایسے $t' \in T$ لئی ہووے
    جیہدے لئی $t \prec_\sigma t'$۔}
  \item\ollabel{defn:sub:phmodels-diamond} \indcase{!A}{\Diamond
    !B}{$\mSat{M}{\indfrm}[t, \sigma]$ اُدوں تے صرف اُدوں جدوں
    $\mSat{M}{!B}[t, \sigma']$ کسے ایسے $\sigma' \in C$ لئی
    ہووے جیہدے وچ $t$ آؤندا اے۔}
  \end{enumerate}
\end{defn}

ہور سمے دے تے امکانی عامل وی ایس طرح متعین ہو سکدے نیں۔ ہُن اسیں
اوہ بیان وکھا سکدے آں جیہناں وچ فعل دا زمانہ تے امکان رلے ہون۔
مثلاً «$p$ نہیں ہووے گا، پر ہو سکدا سی» نوں اسیں فارمولا
$\lnot \Ftemp p \land \Diamond \Ftemp p$ نال وکھا سکدے آں۔
ایہ ایسے نقطے تے تاریخ اُتے سچ ہووے گا جیہدے وچ $p$ اگلی کسے
حالت اُتے سچ نہیں ہوندا، پر کوئی ہور تاریخ اے جیہدے وچ $p$ سچ ہو جائے گا۔
\endgroup
% /OLP-0481

\OLEndPartHook
\endgroup
% /OLP-0476

% OLP-0482 content/applied-modal-logic/epistemic-logic/epistemic-logic.tex
\begingroup
\def\olfilename{epistemic-logic}
\expandafter\def\csname ol@sectioncs\endcsname{section}

\olchapter{aml}{el}{علم دیاں منطقاں}

\begin{editorial}
  ایہ باب علم دیاں منطقاں دے مابعد نظریے بارے اے۔ ایہدی بناوٹ آلڈو
  انتونیلی دے کلاسیکی بنیادی امکانی منطق والے نوٹاں ورگی اے۔ آڈری یاپ
  نے اینوں مُڑ لکھیا تاں جو بائی سیمولیشن تے علم دی تبدیلی دیاں منطقاں
  بارے مواد شامل ہو سکے۔
\end{editorial}

% OLP-0483 content/applied-modal-logic/epistemic-logic/introduction.tex
\begingroup
\def\olfilename{introduction}
\expandafter\def\csname ol@sectioncs\endcsname{section}

\olfileid{aml}{el}{int}

\olsection{جان پہچان}

جیویں امکانی منطق \emph{امکانی قضیاں} تے اوہناں دے آپسی استلزام
دے تعلقاں نوں ویکھدی اے، اویویں علم دی منطق \emph{علم بارے قضیاں}
تے اوہناں دے آپسی استلزام دے تعلقاں نوں ویکھدی اے۔ ایتھے اک رکنی
رابطیاں دی تعبیر علم یا منّے ہوئے خیال دے لحاظ نال ہوندی اے، تاں جو
جانن تے منّن دے نمونے بن سکن۔ عام امکانی تعبیر وچ اوہی عامل امکان
تے لزوم دسدے نیں۔ مثلاً اسیں ایہو جیہے بیان آکھنا چاہ سکدے آں:

\begin{enumerate}
\item رچرڈ نوں پتا اے کہ کیلگری البرٹا وچ اے۔
\item آڈری سمجھدی اے کہ ہو سکدا اے کوئی کُتا صوفے اُتے ہووے۔
\item رچرڈ نوں پتا اے کہ آڈری نوں پتا اے کہ اوہدی جماعت منگل والے دن ہوندی اے۔
\item ہر کسے نوں پتا اے کہ اک سال وچ 12 مہینے ہوندے نیں۔
\end{enumerate}
اجوکی علم دی منطق دا سرا اکثر یاکو ہنٹکّا دی 1962 دی کتاب
\emph{Knowledge and Belief} نال جوڑیا جاندا اے۔ ایہ اوس ویلے لکھی گئی
جدوں منطق وچ ممکن جہاناں دیاں معنیات دی ورتوں ودھ رہی سی۔ علم دیاں
منطقاں ہور امکانی منطقاں والے زیادہ تر معنیاتی وسیلے ای ورتدیاں نیں،
پر اوہناں دی تعبیر وکھری اے۔ وڈی تبدیلی ایہ اے کہ اسیں \emph{رسائی دے تعلق}
نوں کیہدا نشان سمجھدے آں۔ علم دیاں منطقاں وچ ایہ تعلق \emph{علم دے لحاظ
نال امکان} دی کوئی صورت وکھاؤندے نیں۔ اسیں ویکھاں گے کہ علم دا جیہڑا
خیال اسیں نمونے وچ لیندے آں، اوہ رسائی دے تعلق اُتے لاؤن والیاں قیداں
نوں کِویں بدلدا اے۔ اسیں ایہ وی ویکھاں گے کہ علم دے لحاظ نال تعبیر
کرن تے فریم خاصیتاں تے فارمولیاں دے میل دا نظریہ کیویں ہوندا اے۔
اُتلیاں مثالاں وچ دو فاعل، رچرڈ تے آڈری، تے اوہناں دونہاں دی جانکاری
دا آپسی تعلق اے۔ اسیں کئی فاعلاں والیاں علم دیاں منطقاں ویکھاں گے،
جیہناں وچ ایہو جیہیاں گلاں آکھیاں جا سکدیاں نیں۔ اک فاعل والی علم دی
منطق صرف اک فرد دے جانن یا منّن بارے گل کردی اے۔
\endgroup
% /OLP-0483
% OLP-0484 content/applied-modal-logic/epistemic-logic/language-epistemic-logic.tex
\begingroup
\def\olfilename{language-epistemic-logic}
\expandafter\def\csname ol@sectioncs\endcsname{section}

\olfileid{aml}{el}{lan}

\olsection{علم دی منطق دی زبان}

\begin{defn}
منّ لو $G$ فاعلاں دے نشاناں دا سیٹ اے۔ کئی فاعلاں والی علم دی منطق
دی بنیادی زبان وچ ایہ شامل نیں:
\begin{enumerate}
  \tagitem{prvFalse}{!!{falsity} لئی قضیاتی مستقل~$\lfalse$۔}{}
  \tagitem{prvTrue}{!!{truth} لئی قضیاتی مستقل~$\ltrue$۔}{}
  \item !!{propositional variable}s دا اک !!{denumerable}s سیٹ: $\Obj
    p_0$, $\Obj p_1$, $\Obj p_2$, \dots
  \item قضیاتی رابطے: \startycommalist
  \iftag{prvNot}{\ycomma $\lnot$ (نفی)}{}%
  \iftag{prvAnd}{\ycomma $\land$ (عطف)}{}%
  \iftag{prvOr}{\ycomma $\lor$ (فصل)}{}%
  \iftag{prvIf}{\ycomma $\lif$ (!!{conditional})}{}%
  \iftag{prvIff}{\ycomma $\liff$ (!!{biconditional})}{}۔
  \item علم دا عامل $\Knows_a$ جتھے $a \in G$۔
\end{enumerate}
\end{defn}

جے ساڈے نظام وچ اسیں صرف اک فاعل دا علم ویکھدے آں، تاں سیٹ~$G$
تے وکھ وکھ فاعلاں دا حوالہ چھڈ سکدے آں۔ فیر ساڈے کول صرف بنیادی
عامل~$\Knows$ اے۔

\begin{defn}
علم دی زبان دے \emph{!!^{formula}s} استقرا نال ایویں متعین ہوندے نیں:
\begin{enumerate}
\tagitem{prvFalse}{$\lfalse$ اک ایٹمی !!{formula} اے۔}{}

\tagitem{prvTrue}{$\ltrue$ اک ایٹمی !!{formula} اے۔}{}

\item ہر قضیاتی متغیر $\Obj p_i$ اک (ایٹمی) !!{formula} اے۔

\tagitem{prvNot}{جے $!A$ اک !!a{formula} اے، تاں $\lnot !A$ وی
  اک !!a{formula} اے۔}{}

\tagitem{prvAnd}{جے $!A$ تے $!B$، !!{formula}s نیں، تاں $(!A \land
  !B)$ اک !!a{formula} اے۔}{}

\tagitem{prvOr}{جے $!A$ تے $!B$، !!{formula}s نیں، تاں $(!A \lor !B)$
  اک !!a{formula} اے۔}{}

\tagitem{prvIf}{جے $!A$ تے $!B$، !!{formula}s نیں، تاں $(!A \lif !B)$
  اک !!a{formula} اے۔}{}

\tagitem{prvIff}{جے $!A$ تے $!B$، !!{formula}s نیں، تاں $(!A \liff !B)$
  اک !!a{formula} اے۔}{}

\item جے $!A$ اک !!a{formula} اے تے $a \in G$، تاں $\Knows_a !A$
  اک !!a{formula} اے۔

\tagitem{limitClause}{ہور کوئی شے !!a{formula} نہیں۔}{}
\end{enumerate}
\end{defn}

جے !!a{formula}~$!A$ وچ $\Knows_a$ نہیں آؤندا، تاں اسیں اوہنوں
\emph{امکانی عامل توں خالی} آکھدے آں۔

\begin{defn}
  عامل $\Knows$ اک فرد دا علم وکھاؤندا اے۔ $\EKnows$، جیہنوں اکثر
  «ساریاں نوں پتا اے» پڑھدے آں، جماعت دا علم وکھاؤندا اے۔ جتھے
  $G' \subseteq G$، اوتھے اسیں $\EKnows_{G'} !A$ نوں
  ایہدی مختصر لکھت وجوں متعین کردے آں:
  \[\bigwedge_{b \in G'} \Knows_b !A.\]
\end{defn}

اسیں علم دی اک ہور تگڑی صورت وی متعین کر سکدے آں، یعنی فاعلاں
دی جماعت~$G$ وچ \emph{سانجھا علم}۔ جدوں کوئی گل کسے جماعت دے
فاعلاں دا سانجھا علم اے، تاں جماعت دے فاعلاں دے ہر میل لئی ساریاں
نوں پتا اے کہ اک دوجے نوں پتا اے کہ اک دوجے نوں پتا اے \dots
تے ایہ سلسلہ چلدا رہندا اے۔ ایہ جماعت دے علم توں کِتے تگڑی شرط اے۔
ایسے تعلقی نمونے سوکھ نال بن سکدے نیں جیہناں وچ !!a{formula} جماعت
دا علم اے پر سانجھا علم نہیں۔ اسیں $\CKnows_G !A$ نال ایہ وکھاواں گے،
«جماعت~$G$ وچ ایہ سانجھا علم اے کہ~$!A$۔»
\endgroup
% /OLP-0484
% OLP-0485 content/applied-modal-logic/epistemic-logic/relational-models.tex
\begingroup
\def\olfilename{relational-models}
\expandafter\def\csname ol@sectioncs\endcsname{section}

\olfileid{aml}{el}{rel}

\olsection{تعلقی نمونے}

علم دیاں منطقاں دا بنیادی معنیاتی خیال عام امکانی منطقاں والا اے۔
تعلقی نمونیاں وچ ہُن وی جہاناں دا سیٹ تے اک قدر بندی اے جیہڑی طے کردی
اے کہ کہڑے !!{propositional variable}s کہڑے جہاناں اُتے «سچ» نیں۔ جے اسیں صرف اک فاعل دی گل
کر رہے آں، تاں معمول وانگ رسائی دا اکّو تعلق اے۔ کئی فاعلاں والی علم
دی منطق وچ رسائی دے تعلقاں دا سیٹ ہوندا اے، فاعلاں دے نشاناں دے
سیٹ وچ ہر~$a$ لئی اک۔ اوس سیٹ دا ناں~$G$ اے۔

اک \emph{تعلقی نمونے} وچ جہاناں دا سیٹ اے، جیہڑے دو رکنی رسائی
تعلقاں نال جڑے نیں، ہر فاعل لئی اک تعلق۔ نال اک قدر بندی اے جیہڑی
طے کردی اے کہ کہڑے !!{propositional variable}s کہڑے جہاناں اُتے سچ نیں۔

\begin{defn}
  کئی فاعلاں والی علم دی زبان دا اک \emph{نمونہ} ترتیبی تِکڑی
  $\mModel{M} = \tuple{W, R, V}$ اے، جتھے
  \begin{enumerate}
  \item $W$ «جہاناں» دا اک غیر خالی سیٹ اے۔
  \item ہر $a \in G$ لئی، $ {R}_a$،~$W$ اُتے رسائی دا دو رکنی تعلق اے۔
  \item $V$ اوہ فنکشن اے جیہڑا ہر !!{propositional variable}~$p$ نوں ممکن جہاناں دا سیٹ $V(p)$ دیندا اے۔
  \end{enumerate}
  جدوں $R_a ww'$ ہووے، اسیں آکھدے آں کہ $w'$، \emph{فاعل a لئی
  رسائی وچ اے}~$w$ توں۔ جدوں $w \in V(p)$ ہووے، اسیں آکھدے آں کہ
  $p$،~$w$ اُتے \emph{سچ اے}۔
\end{defn}

کم کرن دا طریقہ معیاری امکانی منطق والا اے، بس رسائی دے ہور تعلق
شامل نیں۔ کسے فاعل لئی اسیں عام طور اُتے اوہدے رسائی دے تعلق نوں
اوہدی جانکاری دی حالت بارے کجھ دسن والا سمجھاں گے۔ مثلاً اسیں اکثر
$R_a ww'$ دا مطلب ایہ لیندے آں کہ $w'$،~$a$ دی جانکاری نال میل
کھاندا اے جیہڑی اوہنوں~$w$ اُتے اے۔ دوجے لفظاں وچ،~$w$ اُتے اوہ
جہان $w$ تے جہان~$w'$ وچ فرق نہیں دس سکدا۔
\endgroup
% /OLP-0485
% OLP-0486 content/applied-modal-logic/epistemic-logic/truth-at-w.tex
\begingroup
\def\olfilename{truth-at-w}
\expandafter\def\csname ol@sectioncs\endcsname{section}

\olfileid{aml}{el}{trw}

\olsection{کسے جہان اُتے سچائی}

معیاری امکانی منطق وانگ، علم دا ہر نمونہ طے کردا اے کہ اوہدے
کہڑے جہاناں اُتے کہڑے !!{formula}s سچ نیں۔ تعلقی معنیات دی بنیادی
گل لئی اسیں اوہی لکھت ورتدے آں، «نمونہ $\mModel{M}$، !!{formula}~$!A$
نوں جہان~$w$ اُتے سچ کردا اے»۔ ایہ تعلق استقرا نال متعین ہوندا اے تے
سارے غیر امکانی عاملاں لئی معیاری امکانی منطق والا ای اے۔

\begin{defn}\ollabel{defn:mmodels}
  \emph{!!a{formula}~$!A$ دا~$w$ اُتے سچ ہونا}، نمونے~$\mModel M$
  وچ، نشانیاں وچ $\mSat{M}{!A}[w]$، استقرا نال ایویں متعین ہوندا اے:
  \begin{enumerate}
  \tagitem{prvFalse}{\indcase{!A}{\lfalse}{کدی نہیں کہ $\mSat{M}{\lfalse}[w]$}۔}{}
  \tagitem{prvTrue}{\indcase{!A}{\ltrue}{ہمیشہ $\mSat{M}{\ltrue}[w]$}۔}{}
  \item $\mSat{M}{p}[w]$ اُدوں تے صرف اُدوں جدوں $w \in V(p)$
  \tagitem{prvNot}{\indcase{!A}{\lnot !B}{$\mSat{M}{\indfrm}[w]$ اُدوں تے
    صرف اُدوں جدوں $\mSat/{M}{!B}[w]$}۔}{}
  \tagitem{prvAnd}{\indcase{!A}{(!B \land !C)}{$\mSat{M}{\indfrm}[w]$ اُدوں
    تے صرف اُدوں جدوں $\mSat{M}{!B}[w]$ تے $\mSat{M}{!C}[w]$}۔}{}
  \tagitem{prvOr}{\indcase{!A}{(!B \lor !C)}{$\mSat{M}{\indfrm}[w]$ اُدوں
    تے صرف اُدوں جدوں $\mSat{M}{!B}[w]$ یا $\mSat{M}{!C}[w]$}
    (یا دونیں)۔}{}
  \tagitem{prvIf}{\indcase{!A}{(!B \lif !C)}{$\mSat{M}{\indfrm}[w]$ اُدوں
    تے صرف اُدوں جدوں $\mSat/{M}{!B}[w]$ یا $\mSat{M}{!C}[w]$}۔}{}
  \tagitem{prvIff}{\indcase{!A}{(!B \liff !C)}{$\mSat{M}{\indfrm}[w]$ اُدوں
    تے صرف اُدوں جدوں یا تاں $\mSat{M}{!B}[w]$ تے $\mSat{M}{!C}[w]$
    دونیں ہون، یا نہ $\mSat{M}{!B}[w]$ ہووے تے نہ $\mSat{M}{!C}[w]$}۔}{}
  \item\ollabel{defn:sub:mmodels-box}
    \indcase{!A}{\Knows_a !B}{$\mSat{M}{\indfrm}[w]$ اُدوں تے صرف اُدوں
    جدوں $\mSat{M}{!B}[w']$ ہر اوہ $w' \in W$ لئی ہووے جیہدے لئی $R_a ww'$}
  \end{enumerate}
\end{defn}

پر ایتھے سانوں اپنے رسائی تعلقاں دیاں قیداں بارے سوچنا پئے گا۔
آخر شق~\olref{defn:sub:mmodels-box} مطابق، !!a{formula}~$\Knows_a !B$
~$w$ اُتے اوہدوں سچ اے جدوں کوئی~$w'$ ایسا نہیں جیہدے لئی $R_a ww'$۔
ایہ شق معیاری امکانی منطق والی ای اے۔ جے کسے جہان توں اگے کوئی جہان
نہیں، تاں سارے $\Box$ والے فارمولے اوتھے خالی شرط کر کے سچ نیں۔
جے $\Knows$ نوں \emph{علم} دا نشان سمجھدے آں، تاں ایہ گل عام سمجھ
دے بہت اُلٹ لگدی اے۔ آخر اسیں سمجھدے آں کہ ایہو جیہی \emph{کوئی}
صورت نہیں جیہدے وچ فاعل نوں $!A$ تے $\lnot !A$ دونہاں دا پتا
اکّو ویلے ہو سکدا اے۔

اک حل ایہ اے کہ علم دی منطق وچ رسائی دا تعلق ہمیشہ
\emph{انعکاسی} ہووے۔ ایہ موٹے طور اُتے ایس خیال نال میل کھاندا اے
کہ اصلی جہان فاعل دی جانکاری نال میل کھاندا اے۔ عام طور اُتے علم
دیاں منطقاں ایس 5 ورتدیاں نیں، پر دوجے کمزور نظام وی ورت سکدے نیں،
ایس گل مطابق کہ $\Knows_a$ دا تعلق عین کیہڑی گل دا نشان بننا اے۔

\begin{figure}
  \begin{center}
    \begin{tikzpicture}[modal]
      \node[world] (w1) [label={[align=right]right:\mTrue{p}\\ \mFalse{q}}]{$w_1$}; 
      \node[world] (w2) [label={[align=right]right:\mTrue{p}\\ \mTrue{q}},
        above right=of w1]{$w_2$}; 
      \node[world] (w3) [label={[align=right]right:\mFalse{p}\\ \mFalse{q}},
        below right=of w1] {$w_3$};
      \draw[<->] (w1) to node {$a$} (w2);
      \draw[->,loop left] (w1) to node {$a,b$} (w1);
      \draw[<->] (w1) to [swap] node {$b$} (w3);
      \draw[->,loop above] (w2) to node {$a,b$} (w2) ;
      \draw[->,loop below] (w3) to node {$a,b$} (w3) ;
    \end{tikzpicture}
  \end{center}
  \caption{علم دا اک سادہ نمونہ۔}
  \ollabel{fig:simple}
\end{figure}

\begin{prob}
  ویکھو ہِیٹھلیاں وچوں کہڑیاں گلاں
  \olref[aml][el][trw]{fig:simple} وچ سچ نیں:
  \begin{enumerate}
    \item $\mSat{M}{\lnot q}[w_1]$;
    \item $\mSat{M}{\Knows_a \lnot q}[w_1]$;
    \item $\mSat{M}{\Knows_b \lnot q}[w_1]$;
    \item $\mSat{M}{\Knows_b q \lor \Knows_b \lnot q}[w_2]$;
    \item $\mSat{M}{\Knows_a( \Knows_b q \lor \Knows_b \lnot q)}[w_2]$;
    \item $\mSat{M}{\EKnows_{\{a,b\}} \lnot q}[w_3]$;
    \end{enumerate}
\end{prob}

کسے جہان اُتے سچائی دی بنیادی تعریف ہُن مل گئی اے۔ معیاری
امکانی منطق دے ہور معنیاتی خیال، جیویں امکانی اعتبار تے استلزام،
ایسے ای ایتھے آ جاندے نیں۔ اوہناں نوں عاملاں دی ایس نویں تعبیر
اُتے لائی دا اے۔

ہُن اسیں سانجھے علم دے عامل~$\CKnows_G$ لئی وی سچائی دیاں شرطاں
دے سکدے آں۔ \olref[sfr][rel][ops]{sec} توں یاد کرو کہ
\emph{تعدی تحت بندش}~$R^+$، تعلق~$R$ دی، ایویں متعین ہوندی اے:
\begin{align*}
  R^+ &= \bigcup_{ n \in \mathbb{N}} R^n,
\intertext{جتھے}
  R^0 & = R \text{ تے}\\
  R^{n+1} & = \Setabs{\tuple{x, z}}{\exists y (R^n xy \land Ryz)}.
\end{align*}
فیر جتھے $G$ فاعلاں دی جماعت اے، اسیں $R_G = ( \bigcup_{b
\in G} R_b )^+$ نوں سارے فاعلاں دے رسائی تعلقاں دے میل دی تعدی
تحت بندش وجوں متعین کردے آں۔

\begin{defn}
جے $G' \subseteq G$، تاں $\mSat{M}{\CKnows_{G'} !A}[ w]$ اُدوں تے
صرف اُدوں جدوں ہر اوہ $w'$ لئی جیہدے لئی $R_{G'} w w'$، ایہ ہووے
کہ $\mSat{M}{!A}[w']$۔
\end{defn}
\endgroup
% /OLP-0486
% OLP-0487 content/applied-modal-logic/epistemic-logic/properties-accessibility.tex
\begingroup
\def\olfilename{properties-accessibility}
\expandafter\def\csname ol@sectioncs\endcsname{section}

\olfileid{aml}{el}{acc}

\olsection{رسائی دے تعلق تے علم دے اصول}

معیاری امکانی منطقاں وچ فریم دی خاصیت تے فارمولے دے میل بارے
اسیں پہلاں ای کجھ جاندے آں۔ ہُن ویکھ سکدے آں کہ علم دے لحاظ نال
تعبیر کرن تے خاصیت دی پچھان دین والے !!{formula}s کِہو جیہے ہوندے
نیں۔ اسیں آکھ چکے آں کہ علم دیاں منطقاں دی تعبیر عام طور اُتے ایس 5
وچ ہوندی اے۔ سو ویکھدے آں کہ علم دے وکھ وکھ اصول کِویں وکھائے
جاندے نیں تے فریم دیاں کیہڑیاں شرطاں نال ملدے نیں۔

معیاری امکانی منطق توں یاد کرو کہ وکھ وکھ امکانی !!{formula}s
رسائی تعلقاں دیاں وکھ وکھ خاصیتاں دی پچھان دیندے سن۔ ایس جدول وچ
خاص علم اصولاں نال ملدیاں کجھ خاصیتاں چنیاں نیں۔

\begin{table}[t]
    \begin{tabular}{| p{.48\textwidth} || p{.48\textwidth} |}
      \hline
      {\emph{جے $R$ \dots اے}} & {\emph{تاں \dots،~$\mModel{M}$ وچ سچ اے}} \\
      \hline \hline

      & $\Knows (p \lif q) \lif (\Knows p \lif \Knows q)$ 
      \hfill \newline{(استلزام تحت بندش)} \\
      \hline
      \emph{انعکاسی}: $\forall w Rww$  
      & $\Knows p \lif p$ \hfill \newline{(علم دی صداقت)} \\
      \hline
      \emph{متعدی}: \newline
      $\forall u \forall v \forall w ((Ruv \land Rvw) \lif Ruw)$ & 
      $\Knows p \lif \Knows \Knows p$ \hfill 
           \newline{(مثبت خود آگاہی)} \\
      \hline 
      \emph{اقلیدسی}: \newline
      $\forall w \forall u \forall v ((Rwu \land Rwv) \lif Ruv)$ &  
      $\lnot \Knows p \lif \Knows \neg \Knows p$ \hfill 
        \newline{(منفی خود آگاہی)} \\
      \hline
    \end{tabular}
    \caption{علم دے چار اصول۔}
    \ollabel{tab:four}
  \end{table} 

علم دی صداقت، جیہڑی مسلّمے $T$ نال ملدی اے، اکثر ایہناں اصولاں
وچوں سبھ توں گھٹ اختلاف والی سمجھی جاندی اے۔ ایہ آکھدی اے کہ جے
!!a{formula} دا پتا اے، تاں اوہ لازمی سچ اے۔ استلزام تحت بندش تے
مثبت تے منفی خود آگاہی اُتے کِتے ودھ اختلاف اے۔

استلزام تحت بندش، جیہڑی مسلّمے $K$ نال ملدی اے، ایہ خیال وکھاؤندی
اے کہ فاعل دا علم استلزام تحت بند اے۔ کجھ صورتاں وچ ایہ گل منّن
جوگی لگ سکدی اے۔ مثلاً مینوں پتا ہو سکدا اے کہ جے میں وکٹوریا وچ
آں، تاں وینکوور ٹاپو اُتے آں۔ جے شک دیاں انوکھی صورتاں پاسے رکھیے،
تاں مینوں پتا اے کہ میں وکٹوریا وچ آں۔ ایس توں ایہ وی سمجھ آؤندی اے
کہ مینوں پتا اے کہ میں وینکوور ٹاپو اُتے آں۔ سو ایس صورت وچ میرے
علم دی منطقی بندش چنگی طرح سمجھ آ سکدی اے۔ دوجے پاسے اسیں اپنے علم
دے نتیجیاں اُتے ہمیشہ پورا غور نہیں کردے۔ ایس کرکے ہور صورتاں وچ
ایسے نتیجے نکل سکدے نیں جیہڑے گھٹ سِدھے لگن۔

مثبت خود آگاہی نوں کدی کدی کے کے دا اصول آکھیا جاندا اے۔
اینوں ایویں بیان کیتا جاندا اے کہ جے مینوں کسے گل دا پتا اے، تاں
مینوں ایہ وی پتا اے کہ مینوں اوہدا پتا اے۔ ایہ مسلّمے 4 دا علم
والا روپ اے۔ ایس طرح منفی خود آگاہی آکھدی اے کہ جے مینوں کسے
گل دا پتا \emph{نہیں}، تاں مینوں پتا اے کہ مینوں اوہدا پتا نہیں۔
ایہ مسلّمے 5 دا علم والا روپ اے۔ دونہاں دیاں عام جہیاں خلاف مثالاں
لگدیاں نیں، جتھے مینوں یقین نہیں کہ مینوں کسے گل دا پتا اے یا
نہیں، بھاویں اصل وچ مینوں اوہدا پتا اے۔
\endgroup
% /OLP-0487
% OLP-0488 content/applied-modal-logic/epistemic-logic/bisimulations.tex
\begingroup
\def\olfilename{bisimulations}
\expandafter\def\csname ol@sectioncs\endcsname{section}

\olfileid{aml}{el}{bsd}

\olsection{بائی سیمولیشناں}

ساڈی علم دی زبان دی اظہاری طاقت بارے اک ہور سوال نمونیاں تے
اوہناں وچ لاگو ہون والے فارمولیاں دے تعلق بارے اے۔ فریم خاصیت تے
فارمولے دے میل والے نتیجیاں توں اسیں ویکھیا اے کہ جدوں کجھ فارمولے
کسے فریم وچ معتبر ہون، تاں اوہ ایہ وی یقینی بناؤندے نیں کہ فریم
کجھ خاصیتاں پوریاں کردا اے۔ پر مثلاً کی ساڈی امکانی زبان سانوں
اک ایسے جہان، جیہدے توں اوہدے اپنے ول تیر اے، تے جہاناں دی اک
لامتناہی لڑی، جیہدے وچ ہر جہان توں اگلے ول تیر اے، وچ فرق دسن
دیندی اے؟ یعنی کی کوئی ایسا فارمولا $A$ اے جیہڑا ایہناں دو جہاناں
وچوں صرف اک اُتے لاگو ہو سکدا اے؟

بائی سیمولیشن اک تعلق اے جیہڑا تعلقی نمونیاں وچ ایہ دسن لئی
متعین کر سکدے آں کہ اوہناں دی بناوٹ اثر دے لحاظ نال اکو جیہی اے۔
اسیں ویکھاں گے کہ ایہ نمونیاں دی اوہ ہم ارزی پھڑدا اے جیہڑی ساڈی
علم دی زبان وچ ظاہر ہو سکدی اے۔

\begin{defn}[بائی سیمولیشن]
منّ لو $M_1 = \tuple{W_1, R_1, V_1}$ تے $M_2 = \tuple{W_2, R_2, V_2}$
دو تعلقی نمونے نیں۔ منّ لو $\mathcal{R} \subseteq W_1 \times W_2$
اک دو رکنی تعلق اے۔ $\mathcal{R}$ نوں \emph{بائی سیمولیشن} آکھدے
آں جدوں ہر $\tuple{w_1, w_2} \in \mathcal{R}$ لئی ایہ ہووے:

\begin{enumerate}
  \item $w_1 \in V_1(p)$ اُدوں تے صرف اُدوں جدوں $w_2 \in V_2(p)$،
    سارے !!{propositional variable}s~$p$ لئی۔
  \item سارے فاعلاں $a \in A$ تے جہاناں $v_1 \in W_1$ لئی، جے
    $R_{1_a} w_1 v_1$، تاں کوئی $v_2 \in W_2$ اے جیہدے لئی
    $R_{2_a} w_2 v_2$ تے $\tuple{v_1, v_2} \in \mathcal{R}$۔
  \item سارے فاعلاں $a \in A$ تے جہاناں $v_2 \in W_2$ لئی، جے
    $R_{2_a} w_2 v_2$، تاں کوئی $v_1 \in W_1$ اے جیہدے لئی
    $R_{1_a} w_1 v_1$ تے $\tuple{v_1, v_2} \in \mathcal{R}$۔
\end{enumerate}

جدوں $M_1$ تے $M_2$ وچ کوئی بائی سیمولیشن اے جیہڑی جہاناں
$w_1$ تے $w_2$ نوں جوڑدی اے، تاں اسیں
$\tuple{M_1, w_1} \leftrightarroweq \tuple{M_2, w_2}$ وی لکھ سکدے
آں، تے $\tuple{M_1, w_1}$ تے $\tuple{M_2, w_2}$ نوں
\emph{بائی سیمولیشن نال جڑے} آکھدے آں۔
\end{defn}

بائی سیمولیشن تعلق دیاں وکھ وکھ شقّاں وکھ وکھ گلاں یقینی بناؤندیاں
نیں۔ شق 1 یقینی بناؤندی اے کہ بائی سیمولیشن نال جڑے جہان امکانی
عامل توں خالی اوہی فارمولے پورے کرن گے، کیوں جے سارے
!!{propositional variable}s اُتے میل اے۔ دوجیاں دو شقّاں نوں کدی
کدی واری واری «اگے» تے «پچھے» آکھدے نیں۔ اوہ یقینی بناؤندیاں نیں
کہ رسائی دے تعلقاں دی بناوٹ اکو جیہی ہووے گی۔

\begin{thm}
جے $\tuple{M_1, w_1} \leftrightarroweq \tuple{M_2, w_2}$، تاں ہر
!!{formula}~$!A$ لئی $\mSat{M_1}{!A}[w_1]$ اُدوں تے صرف اُدوں
جدوں $\mSat{M_2}{!A}[w_2]$۔
\end{thm}

\begin{figure}
  \begin{center}
    \begin{tikzpicture}[modal]
      \node[world] (w1) {$w_1$}; 
      \node[world] (w2) [above left=of w1]{$w_2$}; 
      \node[world] (w3) [above right=of w1] {$w_3$};
      \draw[<->] (w1) to node {$a$} (w2);
      \draw[->,loop below] (w1) to node {$a$} (w1);
      \draw[<->] (w1) to [swap] node {$a$} (w3);
      \draw[->,loop above] (w2) to node {$a$} (w2) ;
      \draw[->,loop above] (w3) to node {$a$} (w3) ;
      
      \node[world](v1)[right of=w1, xshift=1.5in]{$v_1$};
      \node[world](v2)[above of=v1]{$v_2$};
      \draw[<->] (v1) to node {$a$} (v2);
      \draw[->,loop below] (v1) to node {$a$} (v1);
      \draw[->,loop above] (v2) to node {$a$} (v2) ;

      \draw[-,dotted] (w1) to (v1) ;
      \draw[-,dotted,bend left=45] (w2) to (v2) ;
      \draw[-,dotted,bend left=30] (w3) to (v2) ;
    \end{tikzpicture}
  \end{center}
  \caption{بائی سیمولیشن نال جڑے دو نمونے۔}
  \ollabel{fig:bisimilar}
\end{figure}

\olref{fig:bisimilar} وچ وکھائے دونیں نمونے کجھ وکھرے نیں، پر
اک بائی سیمولیشن جہاناں $w_1$ تے~$v_1$ نوں جوڑدی اے۔ ایہ بائی
سیمولیشن $w_2$ تے $w_3$ دونہاں نوں~$v_2$ نال وی جوڑے گی۔
گل ایہ اے کہ ساڈی امکانی زبان وچ لکھن جوگ کوئی شے ایہناں وچ فرق
نہیں دس سکدی۔ پر جے $w_2$ تے~$w_3$ وکھ وکھ
!!{propositional variable}s پورے کردے ہون، تاں صورت وکھری ہووے گی۔
\endgroup
% /OLP-0488
% OLP-0489 content/applied-modal-logic/epistemic-logic/public-announcement-logic-lang.tex
\begingroup
\def\olfilename{public-announcement-logic-lang}
\expandafter\def\csname ol@sectioncs\endcsname{section}

\olfileid{aml}{el}{pal}

\olsection{عام اعلان دی منطق}

علم دی تبدیلی دیاں منطقاں سانوں وکھاؤن دیندیاں نیں کہ فاعلاں دا
علم سمے نال، یا نویں جانکاری ملن نال، کِویں بدلدا اے۔ ایہناں وچوں
کئی، علم دی تبدیلی نوں جانکاری دے \emph{واقعیاں} یا \emph{تبدیلیاں}
نال وکھاؤندیاں نیں۔ سبھ توں سادی تبدیلی عام اعلان اے۔ اوہدے وچ اک
فارمولے دا سچا اعلان ہوندا اے تے سارے فاعل ایہ ہوندا اکٹھے ویکھدے
نیں۔ ایس لئی اسیں زبان نوں ایویں ودھاؤندے آں۔

\begin{defn}
منّ لو $G$ فاعلاں دے نشاناں دا سیٹ اے۔ کئی فاعلاں والی علم دی منطق
دی عام اعلان والی بنیادی زبان وچ ایہ شامل نیں:
\begin{enumerate}
  \tagitem{prvFalse}{!!{falsity} لئی قضیاتی مستقل~$\lfalse$۔}{}
  \tagitem{prvTrue}{!!{truth} لئی قضیاتی مستقل~$\ltrue$۔}{}
  \item !!{propositional variable}s دا اک !!{denumerable}s سیٹ: $\Obj
    p_0$, $\Obj p_1$, $\Obj p_2$, \dots
  \item قضیاتی رابطے: \startycommalist
  \iftag{prvNot}{\ycomma $\lnot$ (نفی)}{}%
  \iftag{prvAnd}{\ycomma $\land$ (عطف)}{}%
  \iftag{prvOr}{\ycomma $\lor$ (فصل)}{}%
  \iftag{prvIf}{\ycomma $\lif$ (!!{conditional})}{}۔
  \item علم دا عامل $\Knows_a$ جتھے $a \in G$۔
  \item عام اعلان دا عامل $[!B]$، جتھے $!B$ اک !!a{formula} اے۔
\end{enumerate}
\end{defn}

عام اعلان دا عامل باکس والے عامل وانگ کم کردا اے۔ ایس حساب نال
ساڈی زبان دی استقرائی تعریف ایویں اے:

\begin{defn}
علم دی زبان دے \emph{!!^{formula}s} استقرا نال ایویں متعین ہوندے نیں:
\begin{enumerate}
\tagitem{prvFalse}{$\lfalse$ اک ایٹمی !!{formula} اے۔}{}

\tagitem{prvTrue}{$\ltrue$ اک ایٹمی !!{formula} اے۔}{}

\item ہر قضیاتی متغیر $\Obj p_i$ اک (ایٹمی) !!{formula} اے۔

\tagitem{prvNot}{جے $!A$ اک !!a{formula} اے، تاں $\lnot !A$ وی
  اک !!a{formula} اے۔}{}

\tagitem{prvAnd}{جے $!A$ تے $!B$، !!{formula}s نیں، تاں $(!A \land
  !B)$ اک !!a{formula} اے۔}{}

\tagitem{prvOr}{جے $!A$ تے $!B$، !!{formula}s نیں، تاں $(!A \lor !B)$
  اک !!a{formula} اے۔}{}

\tagitem{prvIf}{جے $!A$ تے $!B$، !!{formula}s نیں، تاں $(!A \lif !B)$
  اک !!a{formula} اے۔}{}

\tagitem{prvIff}{جے $!A$ تے $!B$، !!{formula}s نیں، تاں $(!A \liff !B)$
  اک !!a{formula} اے۔}{}

\item جے $!A$ اک !!a{formula} اے تے $a \in G$، تاں $\Knows_a !A$
  اک !!a{formula} اے۔

\item جے $!A$ تے $!B$، !!{formula}s نیں، تاں $[!A] !B$ اک !!a{formula} اے۔

\tagitem{limitClause}{ہور کوئی شے !!a{formula} نہیں۔}{}
\end{enumerate}
\end{defn}

!!{formula} $[!A] !B$ دی مراد ایہ اے، «$!A$ دا سچا اعلان ہوون
پچھوں~$!B$ لاگو اے۔» کدی کدی عام اعلان دے سیاق وچ سانجھے علم دی گل
کرنی وی فائدے والی اے، ایس لئی زبان وچ سانجھے علم دا عامل~$\CKnows_G
!A$ وی شامل ہو سکدا اے۔
\endgroup
% /OLP-0489
% OLP-0490 content/applied-modal-logic/epistemic-logic/public-announcement-logic-semantics.tex
\begingroup
\def\olfilename{public-announcement-logic-semantics}
\expandafter\def\csname ol@sectioncs\endcsname{section}

\olfileid{aml}{el}{psm}

\olsection{عام اعلان دی منطق دیاں معنیات}

عام اعلان دیاں منطقاں لئی تعلقی نمونے اوہی نیں جیہڑے علم دیاں
منطقاں لئی سن۔ عام اعلان دے عامل دیاں معنیات نویاں نیں۔

\begin{defn}\ollabel{defn:mmodels}
  \emph{!!a{formula}~$!A$ دا~$w$ اُتے سچ ہونا}، نمونے~$\mModel M = \tuple{W, R, V}$
  وچ، نشانیاں وچ $\mSat{M}{!A}[w]$، استقرا نال ایویں متعین ہوندا اے:
  \begin{enumerate}
  \tagitem{prvFalse}{\indcase{!A}{\lfalse}{کدی نہیں کہ $\mSat{M}{\lfalse}[w]$}۔}{}
  \tagitem{prvTrue}{\indcase{!A}{\ltrue}{ہمیشہ $\mSat{M}{\ltrue}[w]$}۔}{}
  \item $\mSat{M}{p}[w]$ اُدوں تے صرف اُدوں جدوں $w \in V(p)$
  \tagitem{prvNot}{\indcase{!A}{\lnot !B}{$\mSat{M}{\indfrm}[w]$ اُدوں تے
    صرف اُدوں جدوں $\mSat/{M}{!B}[w]$}۔}{}
  \tagitem{prvAnd}{\indcase{!A}{(!B \land !C)}{$\mSat{M}{\indfrm}[w]$ اُدوں
    تے صرف اُدوں جدوں $\mSat{M}{!B}[w]$ تے $\mSat{M}{!C}[w]$}۔}{}
  \tagitem{prvOr}{\indcase{!A}{(!B \lor !C)}{$\mSat{M}{\indfrm}[w]$ اُدوں
    تے صرف اُدوں جدوں $\mSat{M}{!B}[w]$ یا $\mSat{M}{!C}[w]$}
    (یا دونیں)۔}{}
  \tagitem{prvIf}{\indcase{!A}{(!B \lif !C)}{$\mSat{M}{\indfrm}[w]$ اُدوں
    تے صرف اُدوں جدوں $\mSat/{M}{!B}[w]$ یا $\mSat{M}{!C}[w]$}۔}{}
  \tagitem{prvIff}{\indcase{!A}{(!B \liff !C)}{$\mSat{M}{\indfrm}[w]$ اُدوں
    تے صرف اُدوں جدوں یا تاں $\mSat{M}{!B}[w]$ تے $\mSat{M}{!C}[w]$
    دونیں ہون، یا نہ $\mSat{M}{!B}[w]$ ہووے تے نہ $\mSat{M}{!C}[w]$}۔}{}
  \item\ollabel{defn:sub:mmodels-box}
    \indcase{!A}{\Knows_a !B}{$\mSat{M}{\indfrm}[w]$ اُدوں تے صرف اُدوں
    جدوں $\mSat{M}{!B}[w']$ ہر اوہ $w' \in W$ لئی ہووے جیہدے لئی $R_a ww'$}
  \item\ollabel{defn:sub:mmodels-pal}
    \indcase{!A}{[!B] !C}{$\mSat{M}{\indfrm}[w]$ اُدوں تے صرف اُدوں
    جدوں $\mSat{M}{!B}[w]$ توں $\mSat{M \mid !B}{!C}[w]$ نکلے}

ایتھے $\mModel M \mid !B = \tuple{W', R', V'}$ دی تعریف ایویں اے:
  \begin{enumerate}
    \item $W' = \Setabs{ u \in W}{\mSat{M}{!B}[u]}$۔ سو
      $\mModel M \mid !B$ دے جہان اوہو نیں جیہڑے~$\mModel M$ وچ نیں
      تے جیہناں اُتے $!B$ لاگو اے۔
    \item $R'_a = R_a \cap (W' \times W')$۔ ہر فاعل دا رسائی دا تعلق
      اوہناں جہاناں تک محدود اے جیہڑے~$W'$ وچ رہ گئے نیں۔
    \item $V'(p) = \Setabs{ u \in W'}{u \in V(p) }$۔ ایس طرح جہاناں
      اُتے قضیاتی قدراں اوہو رہندیاں نیں۔ ایہ اوس خیال نوں وکھاؤندا
      اے کہ جانکاری والے واقعے قضیاتی متغیراں دی صداقتی قدر نہیں بدلدے۔
  \end{enumerate}

  \end{enumerate} 
\end{defn}

سو عام اعلان دیاں منطقاں دی خاص گل ایہ اے کہ کدی کدی !!a{formula}
دی سچائی~$\mModel M$ وچ پرکھن لئی،~$\mModel M$ توں وکھرے نمونے دا
حوالہ لینا پیندا اے۔

ایہ وی ویکھو کہ ساڈیاں معنیات وچ اعلان دا عامل $\Box$~والا عامل
اے۔ سو جے کسے جہان اُتے !!a{formula}~$!A$ دا سچا اعلان نہیں ہو سکدا،
تاں $[!A]B$ اوتھے خالی شرط کر کے لاگو اے، جیویں سارے $\Box$
!!{formula}s آخری نقطیاں اُتے لاگو ہوندے نیں۔

\begin{figure}
  \begin{center}
    \begin{tikzpicture}[modal]
      \node[world] (w1) [label={below left:$p, \lnot q$}]{$w_1$}; 
      \node[world] (w2) [left=of w1, label={left:$\lnot p, \lnot q$}]{$w_2$}; 
      \node[world] (w3) [above=of w1, label={right:$p, q$}] {$w_3$};
      \draw[<->] (w1) to [swap] node {$b$} (w2);
      \draw[->,loop below] (w1) to node {$a, b$} (w1);
      \draw[<->] (w1) to [swap] node {$a$} (w3);
      \draw[->,loop above] (w2) to node {$a, b$} (w2) ;
      \draw[->,loop above] (w3) to node {$a, b$} (w3) ;
      
      \node[world](v1)[right of=w1, xshift=1.25in, label={right:$p, \lnot q$}]{$w'_1$};
      \node[world](v3)[above of=v1,label={right:$p,q$}]{$w'_3$};
      \draw[<->] (v1) to node {$a$} (v3);
      \draw[->,loop below] (v1) to node {$a,b$} (v1);
      \draw[->,loop above] (v3) to node {$a,b$} (v3) ;
      
      \node(m)[below=of w1]{$\mModel M$};
      \node(m')[below=of v1]{$\mModel M \mid p$}; 

      \draw[-,dotted] (w1) to node {\footnotesize $p$ دا اعلان}(v1) ;
     
    \end{tikzpicture}
  \end{center}
  \caption{$p$ دے عام اعلان توں پہلاں تے پچھوں۔}
  \ollabel{fig:announcement-example}
\end{figure}

!!a{formula} دے عام اعلان نوں اسیں نمونہ چھوٹا کرن وانگ ویکھ سکدے
آں۔ نمونہ اوہناں جہاناں تک محدود ہو جاندا اے جیہناں اُتے !!{formula}
سچ سی۔ \olref{fig:announcement-example} وچ~$p$ دے عام اعلان دا اثر
وکھایا اے۔ اوس نمونے دی خاص گل ایہ اے کہ فاعل~$b$ اعلان نال سکھدا
اے کہ~$p$ سچ اے، پر فاعل~$a$ لئی ایہ نویں گل نہیں (کیوں جے $a$
نوں پہلاں ای پتا سی کہ~$p$ سچ اے)۔

رسمی لکھت وچ $\mSat{M}{\lnot \Knows_b p}[w_1]$ اے، پر $\mSat{M
\mid p}{\Knows_b p}[w'_1]$ وی اے۔ ایس توں $\mSat{M}{[p] \Knows_b
p}[w_1]$ نکلدا اے۔ اعلان دے نتیجے بارے اسیں ہور تگڑیاں گلاں وی آکھ
سکدے آں۔ اصل وچ $\mSat{M}{[p]\CKnows_{\{a,b\}} p}[w_1]$ اے۔
دوجے لفظاں وچ $p$ دا اعلان ہوون پچھوں ایہ \emph{سانجھا علم} بن جاندا اے۔

پر سوال ایہ اے کہ عام حالت وچ وی ایہ گل لاگو اے؟ کی~$!A$ دا سچا
اعلان \emph{ہر واری} $!A$ نوں سانجھا علم بنا دیندا اے؟ شاید ایہ سن
کے اچرج ہووے کہ جواب ناں اے۔ اصل وچ ایسے !!{formula}s دا سچا اعلان
ہو سکدا اے جیہڑے اعلان ہو جان پچھوں سچ نہیں رہندے۔ مثلاً
$p \land \lnot \Knows_b p$ دا اعلان~$w_1$ اُتے کرن دا اثر ویکھو،
\olref{fig:announcement-example} وچ۔ اصل وچ $\mModel M \mid p$ تے
$\mModel M \mid (p \land \lnot \Knows_b p)$ اکّو نمونہ نیں۔
پر اسیں ویکھ چکے آں کہ $\mSat{M \mid p}{\Knows_b
p}[w'_1]$۔ سو $\mSat{M \mid (p \land \lnot \Knows_b p)}{\lnot
(p \land \lnot \Knows_b p)}[w'_1]$۔ ایس لئی ایہ اجیہا !!a{formula}
اے جیہڑا اعلان ہوون پچھوں جھوٹ ہو جاندا اے۔
\endgroup
% /OLP-0490

\OLEndPartHook
\endgroup
% /OLP-0482

\OLEndPartHook
\endgroup
% /OLP-0475

% OLP-0491 content/intuitionistic-logic/intuitionistic-logic.tex
\begingroup
\def\olfilename{intuitionistic-logic}
\expandafter\def\csname ol@sectioncs\endcsname{section}

\olpart{int}{وجدانیت پسند منطق}

% OLP-0492 content/intuitionistic-logic/introduction/introduction.tex
\begingroup
\def\olfilename{introduction}
\expandafter\def\csname ol@sectioncs\endcsname{section}

\olchapter{int}{int}{جان پہچان}

% OLP-0493 content/intuitionistic-logic/introduction/constructive-reasoning.tex
\begingroup
\def\olfilename{constructive-reasoning}
\expandafter\def\csname ol@sectioncs\endcsname{section}

\olfileid{int}{int}{cr}

\section{تعمیری دلیل}

امکانی عاملاں یا دوجے درجے دے مسوّراں نال کلاسیکی منطق ودھ جاندی
اے۔ وجدانیت پسند منطق کلاسیکی منطق اُتے قید لاؤندی اے، تے ایس لحاظ
نال «غیر کلاسیکی» اے۔ کلاسیکی منطق کئی طرحاں نال \emph{غیر تعمیری} اے۔
وجدانیت پسند منطق دا مقصد اوس تعمیری دلیل نوں رسمی شکل دینا اے جیہڑی
اک قسم دی تعمیری ریاضی وچ ورتی جاندی اے۔ ہِیٹھلیاں مثالاں ایس پچھے
کم کردیاں کجھ وجھاں سمجھا سکدیاں نیں۔

منّ لو کوئی آکھے کہ اوہنے اک قدرتی عدد~$n$ لبھیا اے جیہدی خاصیت
ایہ اے کہ جے $n$ جفت اے، تاں ریمن دا مفروضہ سچ اے، تے جے $n$ طاق
اے، تاں ریمن دا مفروضہ جھوٹ اے۔ بڑی چنگی خبر!{} ریمن دا مفروضہ سچ
اے یا نہیں، ایہ ریاضی دے وڈے اَن حل سوالاں وچوں اک اے۔ لگدا اے
اوہنے ایس مسئلے نوں اک حساب تک لے آندا اے، یعنی صرف ایہ ویکھنا
رہ گیا اے کہ خاص عدد جفت اے یا نہیں۔

~$n$ دی ایہ کراماتی قدر کیہڑی اے؟ اوہ ایویں دسدا اے کہ $n$ اوہ
قدرتی عدد اے جیہڑا $2$ دے برابر اے جے ریمن دا مفروضہ سچ اے، تے
ہور صورت وچ $3$ دے برابر اے۔

تسی غصے وچ اپنے پیسے واپس منگدے او۔ کلاسیکی نظر توں، اُتلا بیان
سچی مچی $n$ دی اکّو قدر متعین کردا اے۔ پر تہانوں $n$ دی اوہ قدر
چاہیدی اے جیہڑی \emph{صاف صاف} دتی ہووے۔

اک ہور مثال ویکھو جیہڑی شاید گھٹ بناوٹی اے۔ سانوں پتا اے کہ
غیر ناطق عدد نوں ناطق قوّت تائیں چُک کے ناطق نتیجہ آ سکدا اے۔
مثلاً $\sqrt{2}^2 = 2$۔ پر ایہ گل اینی صاف نہیں کہ غیر ناطق عدد
نوں \emph{غیر ناطق} قوّت تائیں چُک کے ناطق نتیجہ آ سکدا اے یا نہیں۔
ہِیٹھلا قضیہ دسدا اے کہ ایہ ہو سکدا اے:

\begin{thm}
ایسے غیر ناطق عدد $a$ تے $b$ موجود نیں جیہناں لئی $a^b$ ناطق اے۔
\end{thm}

\begin{proof}
$\sqrt{2}^{\sqrt{2}}$ نوں ویکھو۔ جے ایہ ناطق اے، تاں کم ہو گیا۔
اسیں $a = b = \sqrt{2}$ رکھ سکدے آں۔ ہور صورت وچ ایہ غیر ناطق
اے۔ فیر ساڈے کول
\[
(\sqrt{2}^{\sqrt{2}})^{\sqrt{2}} = \sqrt{2}^{\sqrt{2} \cdot
  \sqrt{2}} = \sqrt{2}^2 = 2,
\]
اے، جیہڑا ناطق اے۔ سو ایس صورت وچ $a$ نوں $\sqrt{2}^{\sqrt{2}}$
تے $b$ نوں~$\sqrt 2$ رکھو۔
\end{proof}

کی ایہ معتبر ثبوت اے؟ زیادہ تر ریاضی دان اینوں معتبر سمجھدے
نیں۔ پر ایتھے وی کجھ کھٹکدا اے۔ اسیں ثابت کیتا اے کہ خاص خاصیت
والی حقیقی عدداں دی جوڑی موجود اے، پر ایہ نہیں دس سکدے کہ
اوہ \emph{کیہڑی} جوڑی اے۔ ایہو نتیجہ ایویں وی ثابت ہو سکدا اے
کہ جوڑی $a$، $b$ ثبوت وچ \emph{دتی ہووے}۔ $a = \sqrt{3}$ تے
$b = \log_3 4$ رکھو۔ فیر
\[
a^b = \sqrt{3}^{\log_3 4} = 3^{1/2 \cdot \log_3 4} = (3^{\log_3
  4})^{1/2} = 4^{1/2}= 2,
\]
کیوں جے $3^{\log_3 x} = x$۔

وجدانیت پسند منطق اوس طرح دی دلیل نوں رسمی شکل دیندی اے جیہدے
وچ پہلے ثبوت ورگے قدم دی اجازت نہیں۔ ایہ ثابت کرن لئی کہ کوئی $x$
~$!A(x)$ نوں پورا کردا اے، اک خاص~$x$ دینا پیندا اے، تے ایہدا
ثبوت وی کہ اوہ $!A$ نوں پورا کردا اے، جیویں دوجے ثبوت وچ اے۔
$!A$ یا $!B$ دے لاگو ہون دا ثبوت دین لئی دونہاں وچوں اک دا ثبوت
دے سکنا چاہیدا اے۔

رسمی لحاظ نال، کلاسیکی منطق دے کسے !!a{derivation} نظام اُتے
اک خاص طرح دی قید لاؤن نال وجدانیت پسند منطق ملدی اے۔ ریاضی دی
نظر توں ایہ رسمی اخذ نظام نیں۔ پر جیویں پہلاں آکھیا اے، ایہناں
دا مقصد ریاضیاتی دلیل دی اک قسم نوں رسمی شکل دینا اے۔ کوئی
اینوں ریاضی دے کسے فلسفیانہ خیال تحت جائز دلیل سمجھ سکدا اے،
مثلاً براوَر دی وجدانیت۔ کوئی اینوں ودھ «ٹھوس» تے تسلی بخش
ریاضیاتی دلیل سمجھ سکدا اے، جیویں بِشپ دی تعمیریت وچ۔ ایس گل
اُتے بحث ہو سکدی اے کہ رسمی بیان پچھے کم کردی غیر رسمی وجھ نوں
پھڑدا اے یا نہیں۔ ساڈی فلسفیانہ سوچ بھاویں جیہڑی وی ہووے، اسیں
وجدانیت پسند منطق نوں رسمی منطق دے طور اُتے پڑھ سکدے آں۔ کئی
ریاضیاتی منطق دان ایہ پڑھائی وکھ وکھ وجھاں کرکے دلچسپ سمجھدے نیں۔
\endgroup
% /OLP-0493

% OLP-0494 content/intuitionistic-logic/introduction/syntax.tex
\begingroup
\def\olfilename{syntax}
\expandafter\def\csname ol@sectioncs\endcsname{section}

\olfileid{int}{int}{syn}

\olsection{وجدانیت پسند منطق دی نحو}

وجدانیت پسند منطق دی نحو قضیاتی منطق والی ای اے۔ کلاسیکی
قضیاتی منطق وچ کجھ رابطے ہور رابطیاں نال متعین ہو سکدے نیں۔
مثلاً $!A \lif !B$ نوں $\lnot !A \lor !B$ نال، یا $!A \lor !B$
نوں $\lnot(\lnot !A \land \lnot !B)$ نال متعین کر سکدے آں۔
ایس لئی کلاسیکی منطق نوں پیش کرن ویلے کجھ رابطے اکثر ایہناں
تعریفاں دی مختصر لکھت وجوں لیاندے جاندے نیں۔ وجدانیت پسند منطق
وچ ایہ گل دو صورتاں وچ چلدی اے۔ $\lnot !A$ نوں $!A \lif \lfalse$
دی مختصر لکھت وجوں متعین کر سکدے آں تے اکثر کردے آں۔ فیر
$\lfalse$ دی اپنی تعریف نہیں ہونی چاہیدی!{} نالے $!A \liff !B$
نوں کلاسیکی منطق وانگ $(!A \lif !B) \land (!B \lif !A)$
دے طور اُتے متعین کر سکدے آں۔

قضیاتی وجدانیت پسند منطق دے !!^{formula}s،
\emph{!!{propositional variable}s} تے قضیاتی مستقل~$\lfalse$ توں
\emph{منطقی رابطیاں} نال بندے نیں۔ ساڈے کول ایہ نیں:

\begin{enumerate}
\item !!^a{denumerable} سیٹ~$\PVar$، جیہدے !!{propositional variable}s
  $\Obj p_0$، $\Obj p_1$، \dots نیں۔
\item !!{falsity} لئی قضیاتی مستقل~$\lfalse$۔
\item منطقی رابطے $\land$ (عطف)، $\lor$ (فصل)، $\lif$ (شرطیہ)۔
\item رموزِ اوقاف، کھلدا قوس «(»، بند ہون والا قوس «)» تے کوما۔
\end{enumerate}

\begin{defn}[فارمولا]
\ollabel{defn:formulas} قضیاتی وجدانیت پسند منطق دے
\emph{!!{formula}s} دا سیٹ~$\Frm[L_0]$ استقرا نال ایویں متعین ہوندا اے:
\begin{enumerate}
\item $\lfalse$ اک ایٹمی !!{formula} اے۔

\item ہر !!{propositional variable}~$\Obj p_i$ اک ایٹمی
  !!{formula} اے۔

\item جے $!A$ تے $!B$، !!{formula}s نیں، تاں $(!A \land
  !B)$ اک !!a{formula} اے۔

\item جے $!A$ تے $!B$، !!{formula}s نیں، تاں $(!A \lor !B)$
  اک !!a{formula} اے۔

\item جے $!A$ تے $!B$، !!{formula}s نیں، تاں $(!A \lif !B)$
  اک !!a{formula} اے۔

\tagitem{limitClause}{ہور کوئی شے !!a{formula} نہیں۔}{}
\end{enumerate}
\end{defn}

اُتے دتے ابتدائی رابطیاں توں علاوہ، اسیں ایہ \emph{متعین} نشان
وی ورتدے آں، $\lnot$ (نفی) تے $\liff$ (!!{biconditional})۔
ایہناں متعین عاملاں نال بنے فارمولے ایویں سمجھے جان:
\begin{enumerate}
\item $\lnot !A$، $!A \lif \lfalse$ دی مختصر لکھت اے۔

\item $!A \liff !B$، $(!A \lif !B) \land (!B
  \lif !A)$ دی مختصر لکھت اے۔
\end{enumerate}

رسمی طور اُتے $\lnot$ اک مختصر لکھت اے، پر کدی کدی اسیں
~$\lnot$ لئی تعریف وچ صاف شقّاں تے قاعدے دیاں گے، جیویں ایہ
ابتدائی نشان ہووے۔ ایہ زیادہ تر ایس لئی اے کہ اسیں مشق دے
مسئلے بیان کر سکیے۔
\endgroup
% /OLP-0494

% OLP-0495 content/intuitionistic-logic/introduction/bhk-interpretation.tex
\begingroup
\def\olfilename{bhk-interpretation}
\expandafter\def\csname ol@sectioncs\endcsname{section}

\olfileid{int}{int}{bhk}

\olsection{براوَر--ہائٹنگ--کولموگوروف دی تعبیر}

\begin{editorial}
  بی ایچ کے تعبیر نال وجدانیت پسند قضیاں دے اعتبار دے ثبوت سمجھنے
  اوکھے نیں۔ اوہناں نوں ہور صاف سمجھاؤن دی لوڑ اے۔
\end{editorial}

وجدانیت پسند رابطیاں دی اک غیر رسمی تعمیری تعبیر اے، جیہنوں عام
طور اُتے براوَر، ہائٹنگ تے کولموگوروف دی تعبیر آکھدے نیں۔ ایہ
«تعمیر» دا خیال ورتدی اے، جیہنوں تسی تعمیری ثبوت سمجھ سکدے او۔
(بی ایچ کے تعبیر وچ اسیں «ثبوت» نہیں آکھدے، تاں جو ایس نوں رسمی
!!{derivation} نظام دے !!a{derivation} نال رلا نہ دیئیے۔) بی ایچ
کے تعبیر ایس سِدھے خیال اُتے وجدانیت پسند رابطیاں دے معنی سمجھاؤندی اے۔

\begin{enumerate}
\item اسیں منّ لیندے آں کہ سانوں پتا اے کہ ایٹمی بیان دی تعمیر
  کیہڑی شے اے۔
\item $!A_1 \land !A_2$ دی تعمیر اک جوڑی $\tuple{M_1, M_2}$ اے،
  جتھے $M_1$،~$!A_1$ دی تعمیر اے تے $M_2$،~$!A_2$ دی تعمیر اے۔
\item $!A_1 \lor !A_2$ دی تعمیر اک جوڑی $\tuple{s, M}$ اے،
  جتھے یا $s$،~$1$ اے تے $M$،~$!A_1$ دی تعمیر اے، یا $s$،~$2$
  اے تے $M$،~$!A_2$ دی تعمیر اے۔
\item $!A \lif !B$ دی تعمیر اوہ فنکشن اے جیہڑا~$!A$ دی تعمیر
  نوں~$!B$ دی تعمیر وچ بدل دیندا اے۔
\item ~$\lfalse$ (محال) دی کوئی تعمیر نہیں۔
\item $\lnot !A$ دی تعریف $!A \lif \lfalse$ دے ہم معنی وجوں ہوندی اے۔
  یعنی $\lnot !A$ دی تعمیر اوہ فنکشن اے جیہڑا~$!A$ دی تعمیر
  نوں~$\lfalse$ دی تعمیر وچ بدل دیندا اے۔
\end{enumerate}

\begin{ex}
مثال وجوں $\lnot \lfalse$ لو۔ ایہدی تعمیر ایسا فنکشن اے جیہڑا
$\lfalse$ دی کوئی وی تعمیر لے کے جواب وچ~$\lfalse$ دی تعمیر دیندا
اے۔ صاف اے کہ شناختی فنکشن~$\Id{}$ ایہو جیہی تعمیر اے۔ جے~$M$،
~$\lfalse$ دی تعمیر اے، تاں $\Id{}(M) = M$ توں~$\lfalse$ دی تعمیر ملدی اے۔
\end{ex}

عام طور اُتے $\lnot !A$ دا مطلب اے کہ $!A$ دی تعمیر ناممکن اے۔

\begin{ex}
$!A \lif \lnot \lnot !A$ ثابت کریے، کسے وی قضیے~$!A$ لئی، یعنی
$!A \lif ((!A \lif \lfalse) \lif \lfalse)$۔ تعمیر ایسا فنکشن~$f$
ہونا چاہیدا اے جیہڑا تعمیر~$M$، $!A$ دی، لے کے تعمیر~$f(M)$
دیندا اے، $(!A \lif \lfalse) \lif \lfalse$ دی۔~$f$،
$(!A \lif \lfalse) \lif \lfalse$ دی تعمیر ایویں بناؤندا اے۔
سانوں ایسا فنکشن $g$ متعین کرنا پئے گا جیہڑا تعمیر~$h$،
$!A \lif \lfalse$ دی، لے کے~$\lfalse$ دی تعمیر دیوے۔ $g$ نوں
ایویں متعین کر سکدے آں۔ ملی ہوئی~$h$ نوں تعمیر~$M$ اُتے لاؤ، جیہڑی
$!A$ دی اے (تے سانوں پہلاں ملی سی)۔ $h(M)$، $h$ توں ملدا جواب،
$\lfalse$ دی تعمیر اے۔ ایس لئی $f(M)(h) = h(M)$،~$\lfalse$ دی تعمیر اے
جے $M$،~$!A$ دی تعمیر اے۔
\end{ex}

\begin{ex}
$\lnot(!A \land \lnot !A)$ دی تعمیر دیئیے، یعنی $(!A
\land (!A \lif \lfalse)) \lif \lfalse$ دی۔ ایہ ایسا فنکشن~$f$
اے جیہڑا تعمیر~$M$، $!A \land (!A \lif \lfalse)$ دی، لے کے
~$\lfalse$ دی تعمیر دیندا اے۔ عطف $!B_1 \land !B_2$ دی تعمیر اک
جوڑی $\tuple{N_1, N_2}$ اے، جتھے $N_1$،~$!B_1$ دی تعمیر اے
تے $N_2$،~$!B_2$ دی تعمیر اے۔ اسیں فنکشن $p_1$ تے $p_2$ متعین
کر سکدے آں جیہڑے $!B_1 \land !B_2$ دی تعمیر توں واری واری
$!B_1$ تے $!B_2$ دیاں تعمیراں کڈھدے نیں:
\begin{align*}
  p_1(\tuple{N_1, N_2}) &= N_1\\
  p_2(\tuple{N_1, N_2}) &= N_2
\end{align*}
$f$ ایویں کم کردا اے۔ پہلاں $p_1$ نوں ملی ہوئی~$M$ اُتے لاؤندا اے۔
ایس توں~$!A$ دی تعمیر ملدی اے۔ فیر $p_2$ نوں~$M$ اُتے لاؤندا اے۔
ایس توں~$!A \lif \lfalse$ دی تعمیر ملدی اے۔ ایہ تعمیر آپ اک فنکشن
~$p_2(M)$ اے جیہڑا~$!A$ دی تعمیر لے کے~$\lfalse$ دی تعمیر دیندا اے۔
دوجے لفظاں وچ، جے اسیں $p_2(M)$ نوں $p_1(M)$ اُتے لائیے، تاں
~$\lfalse$ دی تعمیر ملدی اے۔ سو اسیں $f(M) = p_2(M)(p_1(M))$
متعین کر سکدے آں۔
\end{ex}

\begin{ex}
$((!A \land !B) \lif !C) \lif (!A \lif
(!B \lif !C))$ دی تعمیر دیئیے، یعنی ایسا فنکشن $f$ جیہڑا تعمیر~$g$،
$(!A \land !B) \lif !C$ دی، نوں $(!A \lif (!B \lif
!C))$ دی تعمیر وچ بدل دے۔ تعمیر $g$ آپ اک فنکشن اے
(جیہڑا $!A \land !B$ دیاں تعمیراں توں~$C$ دیاں تعمیراں دیندا اے)۔
تے جواب $f(g)$ اک فنکشن~$h_g$ اے، جیہڑا $!A$ دیاں تعمیراں لے کے
ایسے فنکشن دیندا اے جیہڑے~$!B$ دیاں تعمیراں توں~$!C$ دیاں تعمیراں دیندے نیں۔

ایہ گل سمجھنی اوکھی اے۔ سانوں اک فنکشن $h_g$ بناؤنا پئے گا جیہڑا
$f$ نوں~$g$ دین تے جواب وچ ملے۔~$h_g$ نوں دتی جاندی شے تعمیر~$M$
اے،~$!A$ دی۔ $h_g(M)$ توں ملن والا جواب ایسا فنکشن~$k_M$
ہونا چاہیدا اے جیہڑا تعمیراں~$N$،~$!B$ دیاں، توں~$!C$ دیاں
تعمیراں دیوے۔ منّ لو $k_{g,M}(N) = g(\tuple{M, N})$۔ یاد رکھو
کہ $\tuple{M,
  N}$،~$!A \land !B$ دی تعمیر اے۔ سو $k_{g,M}$،
$!B \lif !C$ دی تعمیر اے۔ ایہ تعمیراں~$N$،~$!B$ دیاں، نوں
~$!C$ دیاں تعمیراں وچ بدلدا اے۔ ہُن $h_g(M) = k_{g,M}$ رکھو۔
ایہ فنکشن تعمیراں~$M$،~$!A$ دیاں، نوں تعمیراں $k_{g,M}$،
$!B \lif !C$ دیاں، وچ بدلدا اے۔ ہُن $f(g) = h_g$ رکھو۔
ایہ فنکشن تعمیراں $g$، $(!A \land !B) \lif !C$ دیاں، نوں
$!A \lif (!B \lif !C)$ دیاں تعمیراں وچ بدلدا اے۔ لو، ہو گیا!{}
\end{ex}

بیان $!A \lor \lnot !A$ نوں تیسرے امکان دے اخراج دا اصول آکھدے آں۔
ایہ کجھ خاص $!A$ لئی ثابت کر سکدے آں (مثلاً $\lfalse \lor
\lnot \lfalse$)، پر عام صورت وچ نہیں۔ وجہ ایہ اے کہ وجدانیت پسند
فصل لئی دونہاں وچوں اک جز دی تعمیر چاہیدی اے۔ پر کجھ بیان ہن نہ
ثابت ہو سکدے نیں نہ ردّ (مثلاً گولڈباخ دا قیاس)۔ تیسرے امکان دے
اخراج دے اصول نوں ردّ وی نہیں کر سکدے، یعنی $\lnot \lnot (!A
\lor \lnot !A)$ لاگو اے۔

\begin{ex}
$\lnot \lnot (!A \lor \lnot !A)$ ثابت کرن لئی سانوں ایسا فنکشن~$f$
چاہیدا اے جیہڑا $\lnot(!A \lor \lnot !A)$ دی تعمیر، یعنی $(!A
\lor (!A \lif \lfalse)) \lif \lfalse$ دی تعمیر، نوں~$\lfalse$
دی تعمیر وچ بدلے۔ دوجے لفظاں وچ سانوں ایسا فنکشن $f$ چاہیدا اے
جیہدے لئی $f(g)$،~$\lfalse$ دی تعمیر اے جے $g$،
$\lnot(!A \lor \lnot !A)$ دی تعمیر اے۔

منّ لو $g$، $\lnot(!A \lor \lnot !A)$ دی تعمیر اے، یعنی ایسا فنکشن
جیہڑا $!A \lor \lnot !A$ دی تعمیر نوں~$\lfalse$ دی تعمیر وچ
بدلدا اے۔ $!A \lor \lnot !A$ دی تعمیر اک جوڑی $\tuple{s, M}$ اے
جتھے یا $s = 1$ اے تے $M$، $!A$ دی تعمیر اے، یا $s = 2$ اے تے
$M$، $\lnot !A$ دی تعمیر اے۔ منّ لو $h_1$ اوہ فنکشن اے جیہڑا تعمیر~$M_1$،
$!A$ دی، نوں $!A \lor \lnot !A$ دی تعمیر وچ بدلدا اے۔ ایہ $M_1$ نوں
$\tuple{1, M_2}$ دیندا اے۔ تے منّ لو $h_2$ اوہ فنکشن اے جیہڑا
تعمیر~$M_2$، $\lnot !A$ دی، نوں $!A \lor \lnot !A$ دی تعمیر وچ
بدلدا اے۔ ایہ $M_2$ نوں $\tuple{2, M_2}$ دیندا اے۔

$k$ نوں $\comp{h_1}{g}$ رکھو۔ ایہ ایسا فنکشن اے جیہڑا~$!A$ دی
تعمیر لے کے~$\lfalse$ دی تعمیر دیندا اے، یعنی ایہ~$!A \lif \lfalse$
یا $\lnot !A$ دی تعمیر اے۔ ہُن $l$ نوں $\comp{h_2}{g}$ رکھو۔
ایہ ایسا فنکشن اے جیہڑا $\lnot !A$ دی تعمیر لے کے~$\lfalse$ دی
تعمیر دیندا اے۔ $k$، $\lnot !A$ دی تعمیر اے، ایس لئی $l(k)$،
~$\lfalse$ دی تعمیر اے۔

سو اسیں ایہ دس دتا اے کہ تعمیر~$g$، $\lnot(!A \lor \lnot !A)$
دی، نوں~$\lfalse$ دی تعمیر وچ کِویں بدلیے۔ یعنی فنکشن $f$، جیہڑا
تعمیر~$g$، $\lnot(!A \lor \lnot !A)$ دی، نوں تعمیر $l(k)$،
~$\lfalse$ دی، وچ بدلدا اے، اوہ $\lnot\lnot(!A \lor \lnot !A)$
دی تعمیر اے۔
\end{ex}

تسی ویکھ سکدے او کہ بی ایچ کے تعبیر نال !!{formula}s دا وجدانیت
پسند اعتبار وکھانا چھیتی ای لمبا تے الجھواں ہو جاندا اے۔ چنگی گل اے
کہ وجدانیت پسند منطق لئی ہور چنگے !!{derivation} نظام تے ہور عین
معنیاتی تعبیراں موجود نیں۔
\endgroup
% /OLP-0495

% OLP-0496 content/intuitionistic-logic/introduction/natural-deduction.tex
\begingroup
\def\olfilename{natural-deduction}
\expandafter\def\csname ol@sectioncs\endcsname{section}

\olfileid{int}{int}{ntd}

\olsection{فطری استنباط}

$\FalseCl$ قاعدیاں توں بغیر فطری استنباط وجدانیت پسند منطق دا اک
معیاری !!{derivation} نظام اے۔ ایتھے اسیں قاعدے فیر دیندے آں تے بی
ایچ کے تعبیر نال اوہناں دی وجھ دسّنے آں۔ ہر صورت وچ اسیں قاعدے نوں
ایویں سمجھ سکدے آں کہ جے مقدمیاں دیاں تعمیراں نیں، تاں نتیجے دی
تعمیر وی اے۔

فطری استنباط دیاں !!{derivation}s وچ ہجے قائم مفروضے ہوندے نیں۔
سو ایسے !!a{derivation} نوں، مثلاً $!A$ دی !!{undischarged}
مفروضیاں~$\Gamma$ توں، اک فنکشن سمجھنا چاہیدا اے جیہڑا سارے
$!B \in \Gamma$ دیاں تعمیراں نوں~$!A$ دی تعمیر وچ بدلدا اے۔
جے !!a{derivation}~$!A$ دی کسے وی !!{undischarged} مفروضے توں بغیر
اے، تاں بی ایچ کے تعبیر دے معنی وچ~$!A$ دی تعمیر اے۔ گل سمجھاؤن
لئی جدوں لوڑ نہیں ہووے گی، اسیں~$\Gamma$ چھڈ دیاں گے۔

مفروضہ~$!A$ اپنے آپ وچ !!a{derivation} اے،~$!A$ دی
!!{undischarged} مفروضے~$!A$ توں۔ ایہ بی ایچ کے تعبیر نال میل
کھاندا اے۔ تعمیراں اُتے شناختی فنکشن~$!A$ دی ہر تعمیر نوں~$!A$
دی تعمیر وچ بدلدا اے۔

\subsection{عطف}

\begin{defish}
\AxiomC{$!A$}
\AxiomC{$!B$}
\RightLabel{\Intro{\land}}
\BinaryInfC{$!A \land !B$}
\DisplayProof
\hfill
\begin{tabular}{l}
\AxiomC{$!A \land !B$}
\RightLabel{\Elim{\land}}
\UnaryInfC{$!A$}
\DisplayProof
\\[3ex]
\AxiomC{$!A \land !B$}
\RightLabel{\Elim{\land}}
\UnaryInfC{$!B$}
\DisplayProof
\end{tabular}
\end{defish}

منّ لو ساڈے کول تعمیراں $N_1$، $N_2$ نیں، واری واری $!A_1$ تے
$!A_2$ دیاں۔ تاں ساڈے کول $!A_1 \land !A_2$ دی تعمیر وی اے،
یعنی جوڑی $\tuple{N_1, N_2}$۔

بی ایچ کے تعبیر مطابق $!A_1 \land !A_1$ دی تعمیر اک جوڑی
$\tuple{N_1, N_2}$ اے۔ سو منّ لو ساڈے کول ایہو جیہی جوڑی اے۔
تاں ساڈے کول عطف دے ہر جز دی تعمیر وی اے۔ $N_1$،~$!A_1$
دی تعمیر اے تے $N_2$، $!A_2$ دی تعمیر اے۔

\subsection{شرطیہ}

\begin{defish}
  \AxiomC{$\Discharge{!A}{u}$}
  \DeduceC{$!B$}
  \DischargeRule{$\Intro{\lif}$}{u}
  \UnaryInfC{$!A \lif !B$}
  \DisplayProof
\hfill
  \AxiomC{$!A \lif !B$}
  \AxiomC{$!A$}
  \RightLabel{$\Elim{\lif}$}
  \BinaryInfC{$!B$}
  \DisplayProof
\end{defish}

جے ساڈے کول !!a{derivation}~$!B$ دی !!{undischarged}
مفروضے~$!A$ توں اے، تاں ایسا فنکشن~$f$ موجود اے جیہڑا
~$!A$ دیاں تعمیراں نوں~$!B$ دیاں تعمیراں وچ بدلدا اے۔ اوہی فنکشن
$!A \lif !B$ دی تعمیر اے۔ سو جے $\Intro\lif$ دے مقدمے دی
تعمیر~$!A$ دی تعمیر اُتے منحصر اے، تاں نتیجے $!A
\lif !B$ دی تعمیر اے۔

دوجے پاسے منّ لو تعمیر~$N$، $!A$ دی، تے تعمیر $f$،
$!A \lif !B$ دی، موجود نیں۔ $!A \lif !B$ دی تعمیر ایسا فنکشن
اے جیہڑا $!A$ دیاں تعمیراں نوں~$!B$ دیاں تعمیراں وچ بدلدا اے۔
سو $f(N)$،~$!B$ دی تعمیر اے، یعنی $\Elim\lif$ دے نتیجے دی تعمیر اے۔

\subsection{فصل}

\begin{defish}
\begin{tabular}{l}
\AxiomC{$!A$}
\RightLabel{\Intro{\lor}}
\UnaryInfC{$!A \lor !B$}
\DisplayProof
\\[3ex]
\AxiomC{$!B$}
\RightLabel{\Intro{\lor}}
\UnaryInfC{$!A \lor !B$}
\DisplayProof
\end{tabular}
\hfill
\AxiomC{$!A \lor !B$}
\AxiomC{$\Discharge{!A}{n}$}
\DeduceC{$!C$}
\AxiomC{$\Discharge{!B}{n}$}
\DeduceC{$!C$}
\DischargeRule{\Elim{\lor}}{n}
\TrinaryInfC{$!C$}
\DisplayProof
\end{defish}

جے ساڈے کول تعمیر~$N_i$،~$!A_i$ دی، اے، تاں اینوں
تعمیر~$\tuple{i, N_i}$،~$!A_1 \lor !A_2$ دی، وچ بدل سکدے آں۔
دوجے پاسے منّ لو ساڈے کول~$!A_1 \lor !A_2$ دی تعمیر اے، یعنی
اک جوڑی $\tuple{i, N_i}$ جتھے $N_i$،~$!A_i$ دی تعمیر اے۔
نال فنکشن $f_1$، $f_2$ وی نیں جیہڑے واری واری $!A_1$، $!A_2$
دیاں تعمیراں نوں~$!C$ دیاں تعمیراں وچ بدلدے نیں۔ تاں $f_i(N_i)$،
~$!C$ دی تعمیر اے، یعنی~$\Elim\lor$ دے نتیجے دی۔

\subsection{محال}

\begin{defish}
  \AxiomC{$\lfalse$}
  \RightLabel{\FalseInt}
  \UnaryInfC{$!A$}
  \DisplayProof
\end{defish}

جے ساڈے کول !!a{derivation}~$\lfalse$ دی !!{undischarged}
مفروضیاں~$!B_1$، \dots، $!B_n$ توں اے، تاں اک فنکشن~$f(M_1,
\dots, M_n)$ اے جیہڑا $!B_1$، \dots، $!B_n$ دیاں تعمیراں نوں
~$\lfalse$ دی تعمیر وچ بدلدا اے۔ $\lfalse$ دی کوئی تعمیر نہیں،
ایس لئی ساریاں $!B_1$، \dots، $!B_n$ دیاں تعمیراں وی اکٹھیاں
نہیں ہو سکدیاں۔ سو $f$ دی ایہ خاصیت وی اے کہ \emph{جے} $M_1$، \dots،
$M_n$ واری واری $!B_1$، \dots، $!B_n$ دیاں تعمیراں نیں،
\emph{تاں} $f(M_1, \dots, M_n)$،~$!A$ دی تعمیر اے۔

\subsection{$\lnot$ دے قاعدے}

$\lnot !A$ نوں $!A \lif \lfalse$ دے طور اُتے متعین کیتا گیا اے،
سو عین لحاظ نال~$\lnot$ لئی قاعدیاں دی لوڑ نہیں۔ جے لوڑ ہوندی،
تاں اوہ ایویں ہوندے:

\begin{defish}
\AxiomC{$\Discharge{!A}{n}$}
\noLine
\DeduceC{$\lfalse$}
\DischargeRule{\Intro{\lnot}}{n}
\UnaryInfC{$\lnot !A$}
\DisplayProof
\hfill
\AxiomC{$\lnot !A$}
\AxiomC{$!A$}
\RightLabel{\Elim{\lnot}}
\BinaryInfC{$\lfalse$}
\DisplayProof
\end{defish}

\subsection{\usetoken{P}{derivation} دیاں مثالاں}

\begin{enumerate}
\item $\Proves !A \lif (\lnot !A \lif \lfalse)$, یعنی $\Proves !A
  \lif ((!A \lif \lfalse) \lif \lfalse)$
  \begin{prooftree}
    \AxiomC{$\Discharge{!A}{2}$}
    \AxiomC{$\Discharge{!A \lif \lfalse}{1}$}
    \RightLabel{\Elim\lif}
    \BinaryInfC{$\lfalse$}
    \DischargeRule{\Intro\lif}{1}
    \UnaryInfC{$(!A \lif \lfalse)\lif \lfalse$}
    \DischargeRule{\Intro\lif}{2}
    \UnaryInfC{$!A \lif (!A \lif \lfalse) \lif \lfalse$}
  \end{prooftree}

\item $\Proves ((!A \land !B) \lif !C) \lif (!A \lif (!B \lif !C))$
  \begin{prooftree}
    \AxiomC{$\Discharge{(!A \land !B) \lif !C}{3}$}
    \AxiomC{$\Discharge{!A}{2}$}
    \AxiomC{$\Discharge{!B}{1}$}
    \RightLabel{\Intro\land}
    \BinaryInfC{$!A \land !B$}
    \RightLabel{\Elim\lif}
    \BinaryInfC{$!C$}
    \DischargeRule{\Intro\lif}{1}
    \UnaryInfC{$!B \lif !C$}
    \DischargeRule{\Intro\lif}{2}
    \UnaryInfC{$!A \lif (!B \lif !C)$}
    \DischargeRule{\Intro\lif}{3}
    \UnaryInfC{$((!A \land !B) \lif !C) \lif (!A \lif (!B \lif !C))$}
  \end{prooftree}

\item $\Proves \lnot(!A \land \lnot !A)$, یعنی $\Proves (!A \land (!A
  \lif \lfalse))\lif \lfalse$
  \begin{prooftree}
    \AxiomC{$\Discharge{!A \land (!A \lif \lfalse)}{1}$}
    \RightLabel{\Elim\land}
    \UnaryInfC{$!A \lif \lfalse$}
    \AxiomC{$\Discharge{!A \land (!A \lif \lfalse)}{1}$}
    \RightLabel{\Elim\land}
    \UnaryInfC{$!A$}
    \RightLabel{\Elim\lif}
    \BinaryInfC{$\lfalse$}
    \DischargeRule{\Intro\lif}{1}
    \UnaryInfC{$(!A \land (!A \lif \lfalse))\lif \lfalse$}
  \end{prooftree}

\item $\Proves \lnot\lnot(!A \lor \lnot !A)$, یعنی $\Proves ((!A \lor
  (!A \lif \lfalse))\lif \lfalse)\lif \lfalse$
  \begin{prooftree}\hskip-3em
    \AxiomC{$\Discharge{(!A \lor (!A \lif \lfalse))\lif \lfalse}{2}$}
    \AxiomC{$\Discharge{(!A \lor (!A \lif \lfalse))\lif \lfalse}{2}$}
    \AxiomC{$\Discharge{!A}{1}$}
    \RightLabel{\Intro\lor}
    \UnaryInfC{$!A \lor (!A \lif \lfalse)$}
    \RightLabel{\Elim\lif}
    \BinaryInfC{$\lfalse$}
    \DischargeRule{\Intro\lif}{1}
    \UnaryInfC{$!A \lif \lfalse$}
    \RightLabel{\Intro\lor}
    \UnaryInfC{$!A \lor (!A \lif \lfalse)$}
    \RightLabel{\Elim\lif}
    \insertBetweenHyps{\hskip -4em}
    \BinaryInfC{$\lfalse$}
    \DischargeRule{\Intro\lif}{2}
    \UnaryInfC{$((!A \lor (!A \lif \lfalse))\lif \lfalse)\lif \lfalse$}
  \end{prooftree}
\end{enumerate}

\begin{prop}
  جے وجدانیت پسند منطق دے فطری استنباط وچ $\Gamma \Proves !A$،
  تاں کلاسیکی منطق وچ $\Gamma \Proves !A$۔ خاص کر کے جے
  $!A$ وجدانیت پسند قضیہ اے، تاں کلاسیکی قضیہ وی اے۔
\end{prop}

\begin{proof}
  فطری استنباط دا ہر قاعدہ کلاسیکی فطری استنباط دا وی قاعدہ اے۔
  سو وجدانیت پسند منطق دی ہر !!{derivation} کلاسیکی منطق وچ وی
  !!a{derivation} اے۔
\end{proof}

\begin{prob}
  وجدانیت پسند منطق وچ ہِیٹھلے !!{formula}s دیاں !!{derivation}s دیو:
  \begin{enumerate}
    \item $(\lnot !A \lor !B) \lif (!A \lif !B)$
    \item $\lnot\lnot\lnot !A \lif \lnot !A$
    \item $\lnot\lnot (!A \land !B) \liff (\lnot\lnot !A\land \lnot \lnot !B)$
    \item $\lnot (!A \lor !B) \liff (\lnot !A \land \lnot !B)$
    \item $(\lnot !A \lor \lnot !B) \lif \lnot (!A \land !B)$
    \item $\lnot\lnot ( !A \land !B) \lif  (\lnot\lnot !A \lor
    \lnot\lnot !B)$
  \end{enumerate}
\end{prob}
\endgroup
% /OLP-0496

% OLP-0497 content/intuitionistic-logic/introduction/axiomatic-derivations.tex
\begingroup
\def\olfilename{axiomatic-derivations}
\expandafter\def\csname ol@sectioncs\endcsname{section}

\olfileid{int}{int}{axd}

\olsection{مسلّمی \usetoken{P}{derivation}}

وجدانیت پسند قضیاتی منطق لئی مسلّمی !!{derivation}s خیال دے لحاظ
نال سبھ توں سادے تے تاریخ وچ سبھ توں پہلے !!{derivation} نظام نیں۔
ایہ کلاسیکی قضیاتی منطق وانگ ای کم کردے نیں۔

\begin{defn}[!!^{derivability}]
جے $\Gamma$، $\Lang L$ دے !!{formula}s دا سیٹ اے، تاں $\Gamma$
توں اک \emph{!!{derivation}}، اک محدود لڑی $!A_1$،
\dots،~$!A_n$ اے جیہڑی !!{formula}s دی اے۔ ایس وچ ہر $i \le n$ لئی ہِیٹھلیاں وچوں اک
گل لاگو اے:
\begin{enumerate}
\item $!A_i \in \Gamma$۔
\item $!A_i$ اک مسلّمہ اے۔
\item $!A_i$ کسے $!A_j$ تے $!A_k$ توں، جتھے $j < i$ تے $k <
  i$، قاعدۂ اثباتِ مقدم نال نکلدا اے، یعنی $!A_k \ident !A_j \lif !A_i$۔
\end{enumerate}
\end{defn}

\begin{defn}[مسلّمے]
وجدانیت پسند قضیاتی منطق دے \emph{مسلّمیاں} دے سیٹ $\PAx$ وچ ہِیٹھلیاں
شکلاں والے سارے !!{formula}s شامل نیں:
\begin{align}
  & (!A \land !B) \lif !A \ollabel{ax:land1}\\
  & (!A \land !B) \lif !B \ollabel{ax:land2}\\
  & \formula A \lif (!B \lif (!A \land !B)) \ollabel{ax:land3}\\
  & \formula A \lif (!A \lor !B) \ollabel{ax:lor1}\\
  & \formula A \lif (!B \lor !A) \ollabel{ax:lor2}\\
  & (!A \lif !C) \lif ((!B \lif !C) \lif ((!A \lor !B) \lif !C)) \ollabel{ax:lor3}\\
  & \formula A \lif (!B \lif !A) \ollabel{ax:lif1}\\
  & (!A \lif (!B \lif !C)) \lif ((!A \lif !B) \lif (!A \lif !C)) \ollabel{ax:lif2}\\
  & \lfalse \lif !A \ollabel{ax:lfalse1}
\end{align}
\end{defn}

\begin{defn}[!!^{derivability}]
!!^a{formula}~$!A$، $\Gamma$ توں \emph{!!{derivable}} اے،
$\Gamma \Proves !A$ لکھدے آں، جے !!a{derivation} ہووے جیہڑی
$\Gamma$ توں $!A$ اُتے مکدی اے۔
\end{defn}

\begin{defn}[قضیے]
!!^a{formula}~$!A$ اک \emph{قضیہ} اے جے !!a{derivation}~$!A$ دی
خالی سیٹ توں ہووے۔ $\Proves !A$ لکھدے آں جے $!A$ قضیہ اے،
تے $\Proves/ !A$ لکھدے آں جے نہیں۔
\end{defn}

\begin{prop}
  جے وجدانیت پسند منطق دے مسلّمی !!{derivation} نظام وچ
  $\Gamma \Proves !A$، تاں کلاسیکی منطق وچ وی $\Gamma \Proves !A$۔
  خاص کر کے، جے $!A$ وجدانیت پسند قضیہ اے، تاں کلاسیکی قضیہ وی اے۔
\end{prop}

\begin{proof}
  ہر وجدانیت پسند مسلّمہ کلاسیکی مسلّمہ وی اے۔ سو وجدانیت پسند منطق
  دی ہر !!{derivation} کلاسیکی منطق وچ وی !!a{derivation} اے۔
\end{proof}
\endgroup
% /OLP-0497

\OLEndChapterHook
\endgroup
% /OLP-0492

% OLP-0498 content/intuitionistic-logic/semantics/semantics.tex
\begingroup
\def\olfilename{semantics}
\expandafter\def\csname ol@sectioncs\endcsname{section}

\olchapter{int}{sem}{معنیات}

\begin{editorial}
  ایس باب وچ وجدانیت پسند منطق دی معنیات دیاں تعریفاں اکٹھیاں کیتیاں گئیاں نیں۔
  ہُن تک کرپکے معنیات تے ٹوپولوجی والی معنیات ای شامل نیں۔ ہجے نہ ایس
  گل دی کوئی مثال اے کہ نمونے فارمولیاں نوں کیویں سچ کردے نیں، نہ
  فارمولیاں دے معتبر ہون دے ثبوتاں دی کوئی مثال اے۔
\end{editorial}

% OLP-0499 content/intuitionistic-logic/semantics/introduction.tex
\begingroup
\def\olfilename{introduction}
\expandafter\def\csname ol@sectioncs\endcsname{section}

\olfileid{int}{sem}{int}

\olsection{جان پہچان}

معنیات توں بناں کسے منطق دا وی تسلی بخش بیان نہیں ہو سکدا۔ ایہ گل وجدانیت پسند
منطق اُتے وی لاگو ہوندی اے۔ کلاسیکی منطق وچ !!{valuation}s اُتے ٹکی معنیات
معیاری اے، پر وجدانیت پسند منطق لئی کئی وکھ وکھ معنیات پیش کیتیاں گئیاں نیں۔
ایہناں وچوں کوئی وی ایس لحاظ نال پوری طرح تسلی بخش نہیں کہ اوہ منطقی رابطیاں
دے معنی وجدانیت پسند منطق دے اصولاں مطابق قبول ہون جوگے طریقے نال سمجھاوے۔

!!{relational model}s اُتے ٹکی معنیات، جیہڑی امکانی منطق دی معنیات ورگی اے،
شاید سبھ توں ودھ رائج اے۔ ایس معنیات وچ !!{propositional variable}s نوں جہاناں
نال جوڑیا جاندا اے تے جہان آپس وچ اک رسائی تعلق نال جڑے ہوندے نیں۔
اوہ تعلق ہمیشہ جزوی ترتیب ہوندا اے، یعنی انعکاسی، ضدِ تقابلی تے متعدی۔

سمجھن لئی تسی ایہناں جہاناں نوں علم دیاں حالتاں یا ’گواہی دیاں صورتاں‘ سمجھ سکدے او۔
حالت~$w'$ تک~$w$ توں رسائی اُدوں تے صرف اُدوں اے جدوں ساڈے موجودہ علم دے لحاظ
نال $w'$ علم دی اک ممکنہ (اگلی) حالت ہووے، یعنی~$w$ وچ جو کجھ معلوم اے
اوہدے نال میل کھاوے۔ اک واری کوئی قضیہ معلوم ہو جاوے تاں اوہ مُڑ نامعلوم نہیں
ہو سکدا۔ یعنی جدوں وی $!A$ دا~$w$ وچ پتا ہووے تے~$Rww'$ ہووے، تاں $!A$ دا
~$w'$ وچ وی پتا اے۔ سو رسائی تعلق دے لحاظ نال ’علم‘ گھٹدا نہیں۔

جے اسیں علم دی منطق وانگ ’$!A$ معلوم اے‘ دی تعریف ’سارے علم دے متبادلاں وچ سچ اے‘
دے طور اُتے کریے، تاں $!A \land !B$،~$w$ وچ معلوم اے جے سارے علم دے متبادلاں وچ
$!A$ تے~$!B$ دونیں معلوم نیں۔ پر علم گھٹدا نہیں تے $R$ انعکاسی اے، ایس لئی
$!A \land !B$، $w$ وچ اُدوں تے صرف اُدوں معلوم اے جدوں $!A$ تے $!B$،~$w$
وچ معلوم نیں۔ ایس وجہ نال $!A \lor !B$ وی $w$ وچ اُدوں تے صرف اُدوں معلوم
اے جدوں دوناں وچوں گھٹ توں گھٹ اک معلوم ہووے۔ سو $\land$ تے $\lor$ دے سچ ہون
دیاں شرطاں اوہو نیں جیہڑیاں کلاسیکی منطق وچ نیں۔

پر شرطیے دے سچ ہون دیاں شرطاں کلاسیکی منطق توں وکھریاں نیں۔ $!A \lif !B$،
~$w$ وچ اُدوں تے صرف اُدوں معلوم اے جدوں کوئی $w'$، جیہڑی $Rww'$ پوری کرے، ایہو جیہی
نہ ہووے جتھے $!A$ معلوم ہووے تے $!B$ معلوم نہ ہووے۔ ایہ ایس شرط ورگی نہیں
کہ $!A$ نامعلوم ہووے یا $!B$ معلوم ہووے، ایہ گل~$w$ وچ ہووے۔ کیوں جے، جے سانوں نہ $!A$ دا پتا اے نہ $!B$ دا، ایہ گل~$w$ وچ ہووے، تاں ہو سکدا اے اگلی علمی حالت~$w'$، جیہڑی $Rww'$ پوری کرے، ایہو جیہی ہووے کہ $w'$ وچ $!A$ معلوم ہو جاوے پر~$!B$ دا پتا نہ لگے۔

سانوں $\lnot !A$ دا پتا صرف اوہدوں اے جدوں اگوں علم دی کوئی ممکنہ حالت
ایہو جیہی نہ ہووے جتھے سانوں~$!A$ دا پتا ہووے۔ گل ایہ اے کہ جے $!A$ دا
پتا لگ سکدا ہووے، تاں علم دی کسے ممکنہ اگلی حالت وچ~$!A$ معلوم ہو جاوے گا۔
چونکہ سانوں $\lfalse$ دا پتا نہیں لگ سکدا، ایس لئی اوہ اگلی علمی حالت وچ
سانوں $!A$ دا پتا ہووے گا پر~$\lfalse$ دا نہیں۔

ایس تعبیر وچ تیسرے امکان دے اخراج دا اصول نہیں چلدا۔ کیوں جے کجھ $!A$ ایسے
ہوندے نیں جیہناں دا سانوں ہجے پتا نہیں، پر پتا لگ سکدا اے۔ ایہو جیہے
!!a{formula}~$!A$ لئی $!A$ تے~$\lnot !A$ دونیں نامعلوم نیں، ایس لئی
$!A \lor \lnot !A$ وی نامعلوم اے۔ پر مثال لئی ایہ گل سانوں معلوم اے کہ
$\lnot(!A \land \lnot !A)$۔ کیوں جے کوئی اگلی حالت ممکن نہیں جتھے سانوں
$!A$ تے $\lnot !A$ دوناں دا پتا ہووے، تے ایہ گل سانوں معلوم اے بھانویں سانوں
~$!A$ یا $\lnot !A$ دا پتا ہووے بھانویں نہ ہووے۔

!!^{relational model}s توں علاوہ وجدانیت پسند منطق دی ہور معنیات وی اے۔
اک ہور ٹوپولوجی والی معنیات اے۔ ایتھے قضیاں دی تعبیر ٹوپولوجی والی فضا دے
کھلے سیٹاں دے طور اُتے ہوندی اے تے رابطیاں دی تعبیر ایہناں سیٹاں اُتے عملاں دے
طور اُتے (مثلاً $\land$ دا میل اشتراک نال اے)۔
\endgroup
% /OLP-0499
% OLP-0500 content/intuitionistic-logic/semantics/relational-models.tex
\begingroup
\def\olfilename{relational-models}
\expandafter\def\csname ol@sectioncs\endcsname{section}

\olfileid{int}{sem}{rel}

\olsection{\usetoken{P}{relational model}}

وجدانیت پسند قضیاتی منطق دی ٹھیک ٹھیک معنیات دین لئی سانوں تعریف کرنی پئے گی
کہ نمونہ کِہنوں آکھنا اے، تاں جو اوہدے لحاظ نال !!{formula}s دے سچ ہون دی پرکھ
کر سکیے۔ ایس تعریف دی بنیاد اُتے معنیاتی دھارناواں، جیویں منطقی اعتبار تے معنوی
نتیجہ، وی تعریف کیتیاں جا سکدیاں نیں۔ اک ایہو جیہی معنیات !!{relational model}s
دے ذریعے ملدی اے۔

\begin{defn}
  وجدانیت پسند قضیاتی منطق لئی !!^a{relational model} اک مرتب تِکڑی
  $\mModel{M} = \tuple{W, R, V}$ اے، جتھے
  \begin{enumerate}
  \item $W$ اک غیر خالی سیٹ اے،
  \item $R$،~$W$ اُتے اک جزوی ترتیب (یعنی اک انعکاسی، ضدِ تقابلی تے متعدی دو رکنی تعلق) اے، تے
  \item $V$ اک فنکشن اے جیہڑا ہر !!{propositional variable}~$p$ نوں~$W$ دا اک ذیلی سیٹ دیندا اے، تے
  \item $V$، $R$ دے لحاظ نال گھٹدا نہیں، یعنی جے $w
    \in V(p)$ تے $Rww'$، تاں $w' \in V(p)$۔
  \end{enumerate}
\end{defn}

\begin{defn}\ollabel{defn:true-at-w}
  اسیں $!A$ دے \emph{$w$ اُتے $\mModel{M}$ وچ سچ ہون}، یعنی
  $\mSat{M}{!A}[w]$، دی تعریف استقرائی طور اُتے ایویں کردے آں:
  \begin{enumerate}
  \item \indcase{!A}{p}{$\mSat{M}{\indfrm}[w]$ اُدوں تے صرف اُدوں، جدوں $w \in V(p)$۔}
  \item \indcase{!A}{\lfalse}{$\mSat{M}{\indfrm}[w]$ نہیں۔}
  \item \indcase{!A}{\lnot !B}{$\mSat{M}{\indfrm}[w]$ اُدوں تے صرف اُدوں، جدوں کوئی
    $w'$ ایہو جیہی نہ ہووے کہ $Rww'$ تے $\mSat{M}{!B}[w']$۔}
  \item \indcase{!A}{!B \land !C}{$\mSat{M}{\indfrm}[w]$ اُدوں تے صرف اُدوں، جدوں
    $\mSat{M}{!B}[w]$ تے $\mSat{M}{!C}[w]$ دونیں۔}
  \item \indcase{!A}{!B \lor !C}{$\mSat{M}{\indfrm}[w]$ اُدوں تے صرف اُدوں، جدوں
    $\mSat{M}{!B}[w]$ یا $\mSat{M}{!C}[w]$ (یا دونیں)۔}
  \item \indcase{!A}{!B \lif !C}{$\mSat{M}{\indfrm}[w]$ اُدوں تے صرف اُدوں، جدوں ہر
    اوہ $w'$ جیہڑی $Rww'$ پوری کرے، اوہدے لئی $\mSat{M}{!B}[w']$ نہیں
    یا $\mSat{M}{!C}[w']$ اے} (یا دونیں)۔
  \end{enumerate}
  اسیں $\mSat/{M}{!A}[w]$ لکھدے آں جے $\mSat{M}{!A}[w]$ نہیں۔ جے $\Gamma$،
  !!{formula}s دا اک سیٹ اے، تاں $\mSat{M}{\Gamma}[w]$ دا مطلب اے کہ
  $\mSat{M}{!B}[w]$ ہر $!B \in \Gamma$ لئی اے۔
\end{defn}

\begin{prob}
  وکھاؤ کہ \olref[int][sem][rel]{defn:true-at-w} دے مطابق
  $\mSat{M}{\lnot !A}[w]$ اُدوں تے صرف اُدوں اے جدوں $\mSat{M}{!A \lif \lfalse}[w]$۔
\end{prob}

\begin{prop}\ollabel{prop:true-monotonic}
  جہاناں اُتے سچائی~$R$ دے لحاظ نال گھٹدی نہیں۔ یعنی جے
  $\mSat{M}{!A}[w]$ تے $Rww'$، تاں $\mSat{M}{!A}[w']$۔
\end{prop}

\begin{proof}
  مشق۔
\end{proof}

\begin{prob}
  \olref[int][sem][rel]{prop:true-monotonic} ثابت کرو۔
\end{prob}
\endgroup
% /OLP-0500
% OLP-0501 content/intuitionistic-logic/semantics/semantic-notions.tex
\begingroup
\def\olfilename{semantic-notions}
\expandafter\def\csname ol@sectioncs\endcsname{section}

\olfileid{int}{sem}{sem}

\olsection{معنیاتی دھارناواں}

\begin{defn}
  اسیں آکھدے آں کہ $!A$ \emph{نمونے $\mModel{M} = \tuple{W,R,V}$ وچ سچ اے}، یعنی $
  \mSat{M}{!A}$، اُدوں تے صرف اُدوں جدوں $\mSat{M}{!A}[w]$ ہر $w \in W$ لئی ہووے۔
  $!A$ \emph{معتبر اے}، یعنی $\Entails !A$، اُدوں تے صرف اُدوں جدوں ایہ سارے نمونیاں وچ سچ ہووے۔
  اسیں آکھدے آں کہ !!{formula}s دے سیٹ~$\Gamma$ دا \emph{معنوی نتیجہ}~$!A$ اے، یعنی $\Gamma
  \Entails !A$، اُدوں تے صرف اُدوں جدوں ہر نمونے~$\mModel{M}$ تے ہر ایہو جیہے~$w$
  لئی، جتھے $\mSat{M}{\Gamma}[w]$ ہووے، $\mSat{M}{!A}[w]$ وی ہووے۔
\end{defn}

\begin{prop}\ollabel{prop:sat-entails}
  \begin{enumerate}
  \item\ollabel{prop:sat-entails1} جے $\mSat{M}{\Gamma}[w]$ تے
    $\Gamma \Entails !A$، تاں $\mSat{M}{!A}[w]$۔
  \item\ollabel{prop:sat-entails2} جے $\mSat{M}{\Gamma}$ تے $\Gamma
    \Entails !A$، تاں $\mSat{M}{!A}$۔
  \end{enumerate}
\end{prop}

\begin{proof}
  \begin{enumerate}
  \item فرض کرو $\mSat{M}{\Gamma}$۔ چونکہ $\Gamma \Entails !A$، اسیں
    جاندے آں کہ جے $\mSat{M}{\Gamma}[w]$، تاں $\mSat{M}{!A}[w]$۔
    چونکہ $\mSat{M}{\Gamma}[u]$ ہر $u \in W$ لئی اے، ایس لئی
    $\mSat{M}{\Gamma}[w]$۔ سو $\mSat{M}{!A}[w]$۔
  \item ایہ \olref{prop:sat-entails1} توں سدھا نکلدا اے۔
  \end{enumerate}
\end{proof}

\begin{defn}\ollabel{defn:restrict}
  فرض کرو $\mModel{M}$ اک تعلقی نمونہ اے تے $w \in W$۔ 
  \emph{محدود نمونہ} $\mModel{M}_w=\tuple{W_w, R_w, V_w}$، جیہڑا
  ~$\mModel{M}$ نوں~$w$ تک محدود کر کے ملدا اے، ایویں دتا جاندا اے:
  \begin{align*}
    W_w & = \Setabs{u \in W}{Rwu},\\
    R_w  & = R \cap (W_w)^2, \text{ تے }\\
    V_w(p) & = V(p) \cap W_w.
  \end{align*}
\end{defn}

\begin{prop}\ollabel{prop:restrict}
  $\mSat{M}{!A}[w]$ اُدوں تے صرف اُدوں اے جدوں $\mSat{M_w}{!A}$۔
\end{prop}

\begin{prob}
  \olref[int][sem][sem]{prop:restrict} ثابت کرو۔
\end{prob}

\begin{prop}
  فرض کرو ہر ایہو جیہے نمونے~$\mModel{M}$ لئی، جتھے $\mSat{M}{\Gamma}$،
  اوتھے $\mSat{M}{!A}$ وی اے۔ فیر $\Gamma \Entails !A$۔
\end{prop}

\begin{proof}
  فرض کرو $\mSat{M}{\Gamma}[w]$۔ \olref{prop:restrict} ہر~$!B \in \Gamma$ اُتے
  لاؤن نال سانوں $\mSat{M_w}{\Gamma}$ ملدا اے۔ مفروضے نال $\mSat{M_w}{!A}$۔
  \olref{prop:restrict} دوبارہ لا کے سانوں $\mSat{M}{!A}[w]$ ملدا اے۔
\end{proof}
\endgroup
% /OLP-0501
% OLP-0502 content/intuitionistic-logic/semantics/topological-semantics.tex
\begingroup
\def\olfilename{topological-semantics}
\expandafter\def\csname ol@sectioncs\endcsname{section}

\olfileid{int}{sem}{top}

\olsection{ٹوپولوجی والی معنیات}

وجدانیت پسند منطق دی معنیات دین دا اک ہور طریقہ ٹوپولوجی دی ریاضیاتی دھارنا
ورتنا اے۔

\begin{defn}
  من لو $X$ اک سیٹ اے۔ \emph{$X$ اُتے ٹوپولوجی} اک سیٹ $\Top{O}
  \subseteq \Pow{X}$ اے جیہڑا ہیٹھلے گُن رکھدا اے۔~$\Top{O}$ دے !!{element}s نوں ٹوپولوجی دے \emph{کھلے سیٹ} آکھدے آں۔
  سیٹ $X$ تے $\Top{O}$ دوناں نوں رلا کے \emph{ٹوپولوجی والی فضا} آکھدے آں۔
  \begin{enumerate}
  \item خالی سیٹ تے پوری فضا کھلے نیں: $\emptyset$, $X \in
    \Top{O}$۔
  \item محدود گنتی دے کھلے سیٹاں دا اشتراک وی کھلا اے۔ جے $U$, $V \in
    \Top{O}$، تاں $U \cap V \in \Top{O}$۔
  \item کھلے سیٹاں دے کسے وی خاندان دا اتحاد وی کھلا اے۔ جے $U_i \in
    \Top{O}$ ہر $i \in I$ لئی، تاں $\bigcup \Setabs{U_i}{i \in
      I} \in \Top{O}$۔
  \end{enumerate}
\end{defn}

جے عبارت دے سیاق توں کھلے سیٹاں دی جماعت معلوم ہو سکے، تاں اسیں ٹوپولوجی
لئی $X$ لکھ سکدے آں۔ فیر وی دھیان رہے کہ $X$ لئی کھلے سیٹ مقرر ہون توں پچھوں ای
اوہنوں ٹوپولوجی آکھیا جا سکدا اے۔

\begin{defn}
  وجدانیت پسند قضیاتی منطق دا \emph{ٹوپولوجی والا نمونہ} اک مرتب تِکڑی
  $\mModel{X} = \tuple{X, \Top{O}, V}$ اے، جتھے $\Top{O}$،~$X$ اُتے اک
  ٹوپولوجی اے تے $V$ اک فنکشن اے جیہڑا $\Top{O}$ وچوں اک کھلا سیٹ ہر
  قضیاتی متغیر نوں دیندا اے۔

  کوئی ٹوپولوجی والا نمونہ~$\mModel{X}$ دتا ہووے، تاں اسیں $\Prop{X}{!A}$ دی
  استقرائی تعریف ایویں کر سکدے آں:
  \begin{enumerate}
  \item $\Prop{X}{\lfalse} = \emptyset$
  \item $\Prop{X}{p} = V(p)$
  \item $\Prop{X}{!A \land !B} = \Prop{X}{!A} \cap \Prop{X}{!B}$
  \item $\Prop{X}{!A \lor !B} = \Prop{X}{!A} \cup \Prop{X}{!B}$
  \item $\Prop{X}{!A \lif !B} = \Interior{(X \setminus \Prop{X}{!A}) \cup
    \Prop{X}{!B}}$
  \end{enumerate}
  ایتھے $\Interior{V}$ اوہ فنکشن اے جیہڑا سیٹ $V \subseteq X$ نوں اوہدے
  \emph{اندرلے حصے} ول بھیجدا اے، یعنی اوہدے اندر پئے سارے کھلے سیٹاں دے
  اتحاد ول۔ دوجے لفظاں وچ،
  \[
  \Interior{V} = \bigcup \Setabs{U}{U \subseteq V \text{ تے } U \in \Top{O}}.
  \]
\end{defn}

دھیان دیو کہ کسے وی سیٹ دا اندرلا حصہ ہمیشہ کھلا ہوندا اے، کیوں جے اوہ
کھلے سیٹاں دا اتحاد اے۔ ایس لئی $\Prop{X}{!A}$ ہمیشہ اک کھلا سیٹ اے۔

بھانویں ٹوپولوجی والی معنیات بڑی مجرد اے، فیر وی اوہنوں سمجھن دے کجھ
طریقے نیں جیہڑے اوہدی ورتوں دی وجہ سمجھا سکدے نیں۔ فرض کرو !!{element}s، یا ’نقطے‘، جیہڑے $X$ دے نیں، اوہ ایسے نقطے نیں
جیہناں اُتے بیاناں دے سچ ہون دی
پرکھ ہو سکدی اے۔ جیہناں سارے نقطیاں اُتے $!A$ سچ اے، اوہناں دا سیٹ
~$!A$ نال ظاہر ہون والا قضیہ اے۔ نقطیاں دا ہر سیٹ ممکنہ قضیہ نہیں، صرف
$\Top{O}$ دے !!{element}s ای ہو سکدے نیں۔ $!A \Entails !B$ اُدوں تے صرف
اُدوں جدوں $!B$ ہر اوہ نقطے اُتے سچ ہووے جتھے~$!A$ سچ اے، یعنی
$\Prop{X}{!A} \subseteq \Prop{X}{!B}$، ہر~$X$ لئی۔ محال بیان~$\lfalse$
کدی سچ نہیں ہوندا، ایس لئی $\Prop{X}{\lfalse} = \emptyset$۔

وجدانیت پسند منطق وچ معتبر قاعدے پورے ہون لئی $!B \land !C$، $!B \lor !C$
تے $!B \lif !C$ نال ظاہر ہون والے قضیاں دا $!B$ تے~$!C$ نال ظاہر ہون والے
قضیاں نال کیہ تعلق ہونا چاہیدا اے؟ یعنی ایہ کیویں ہووے کہ $!A \Proves
!B$ اُدوں تے صرف اُدوں جدوں $\Prop{X}{!A} \subset \Prop{X}{!B}$؟
سانوں $\lfalse \Proves !A$ ہر $!A$ لئی چاہیدا اے۔ ایہ شرط پوری ہوندی اے،
کیوں جے $\emptyset \subseteq U$ ہر $U$ لئی اے۔ چونکہ $!B \land !C \Proves !B$،
سانوں $\Prop{X}{!B \land !C} \subseteq \Prop{X}{!B}$ چاہیدا اے، تے
ایسے طرح $\Prop{X}{!B \land !C} \subseteq \Prop{X}{!C}$۔ $W \subseteq U$
تے $W \subseteq V$ پورے کرن والا سبھ توں وڈا سیٹ $U \cap V$ اے۔
اُلٹے پاسے $!B \Proves !B \lor !C$ تے $!C \Proves !B \lor !C$، ایس لئی
سانوں $\Prop{X}{!B} \subseteq \Prop{X}{!B \lor !C}$ تے
$\Prop{X}{!C} \subseteq \Prop{X}{!B \lor !C}$ چاہیدے نیں۔ سبھ توں چھوٹا
سیٹ~$W$ جیہڑا $U \subseteq W$ تے $V \subseteq W$ پورے کرے، $U \cup V$ اے۔

$\lif$ دی تعریف ذرا اوکھی اے۔ $!A \lif !B$ اوہ سبھ توں کمزور قضیہ ظاہر
کردا اے جیہڑا $!A$ نال رل کے $!B$ نوں معنوی نتیجے دے طور اُتے دیندا اے۔
$!A \lif !B$ دا $!A$ نال رل کے~$!B$ دینا $(!A \lif !B)
\land !A \Proves !B $ توں صاف اے۔ سو $\Prop{X}{!A \lif !B}$ سبھ توں وڈا کھلا
سیٹ ہونا چاہیدا اے جیہڑا $\Prop{X}{!A \lif !B} \cap \Prop{X}{!A}
\subset \Prop{X}{!B}$ پورا کرے۔ ایس توں ساڈی تعریف ملدی اے۔
\endgroup
% /OLP-0502

% OLP-0649 content/intuitionistic-logic/semantics/propositions.tex
\begingroup
\def\olfilename{propositions}
\expandafter\def\csname ol@sectioncs\endcsname{section}

\olfileid{int}{sem}{pro}

\olsection{قضیے}

کدے کدے !!a{relational model}~$\mModel{M}$ وچ اوہناں جہاناں دے سیٹ دا ناں رکھنا کم آؤندا اے جیہناں وچ !!a{formula}~$!A$ سچ اے۔ اسیں ایہنوں~$!A$ نال تعریف کیتا \emph{قضیہ} آکھدے آں تے $\Prop{M}{!A}$ لکھدے آں۔

\begin{defn}
  اسیں فنکشن $V$ دا وادھا کر کے فنکشن $\Prop{M}{!A} \colon \Frm
  \to \Pow{W}$ استقرا نال ایویں بناؤنے آں:
  \begin{enumerate}
  \item $\Prop{M}{\lfalse} = \emptyset$
  \item $\Prop{M}{p} = V(p)$
  \item $\Prop{M}{\lnot !B} = {}$
    \[
      \Setabs{w \in W}{\text{ہر $w'$ لئی جیہدے نال $Rww'$، $w'
          \notin \Prop{M}{!B}$}}.
    \]
  \item $\Prop{M}{!B \land !C} = \Prop{M}{!B} \cap \Prop{M}{!C}$۔
  \item $\Prop{M}{!B \lor !C} = \Prop{M}{!B} \cup \Prop{M}{!C}$۔
  \item $\Prop{M}{!B \lif !C} = {}$
    \[
      \Setabs{w \in W}{\text{ہر $w'$ لئی جیہدے نال $Rww'$، جے $w'
          \in \Prop{M}{!B}$ تاں $w' \in \Prop{M}{!C}$}}.
    \]
  \end{enumerate}
\end{defn}

\begin{prop}
  \ollabel{prop:propositions-true}
  $\mModel{M}{!A}[w]$ اُدوں تے صرف اُدوں جدوں $w \in \Prop{M}{!A}$۔
\end{prop}

\begin{proof}
  مشق۔
\end{proof}

\begin{prob}
  \olref[int][sem][pro]{prop:propositions-true} ثابت کرو۔
\end{prob}
\endgroup
% /OLP-0649
\OLEndChapterHook
\endgroup
% /OLP-0498

% OLP-0503 content/intuitionistic-logic/soundness-completeness/soundness-completeness.tex
\begingroup
\def\olfilename{soundness-completeness}
\expandafter\def\csname ol@sectioncs\endcsname{section}

\olchapter{int}{sc}{صحت مندی تے تمامیت}

\begin{editorial}
  ایس باب وچ وجدانیت پسند قضیاتی منطق دی صحت مندی تے تمامیت دے نتیجے
  اکٹھے کیتے گئے نیں۔ ایس نوں اک تعارف دی لوڑ اے۔ تمامیت دا ثبوت اخذ ہون
  جوگتا بارے کجھ گلاں ورتدا اے، جیہڑیاں کسے تھاں صاف لفظاں وچ بیان تے
  ثابت کرنیاں چاہیدیاں نیں۔
\end{editorial}

% OLP-0504 content/intuitionistic-logic/soundness-completeness/soundness-axd.tex
\begingroup
\def\olfilename{soundness-axd}
\expandafter\def\csname ol@sectioncs\endcsname{section}

\olfileid{int}{sc}{sax}

\olsection{مسلّمی \usetoken{P}{derivation} دی صحت مندی}

\begin{editorial}
  صحت مندی دا ثبوت ایس گل اُتے ٹکیا اے کہ سارے مسلّمے وجدانیت پسند
  منطق وچ معتبر نیں۔ ایہ گل ہجے ثابت کرنی پئے گی، مثلاً معنیات دے باب وچ۔
\end{editorial}

\begin{thm}[صحت مندی]
  \ollabel{thm:soundness}
  جے $\Gamma \Proves !A$، تاں $ \Gamma \Entails !A$۔
\end{thm}

\begin{proof}
  اسیں ثابت کرنے آں کہ جے $\Gamma \Proves !A$، تاں $\Gamma \Entails !A$۔
  ثبوت گنتی~$n$ اُتے استقرا نال اے۔ ایہ اوہناں !!{formula}s دی گنتی اے جیہڑے اُس !!{derivation} وچ نیں
  جیہڑا~$!A$ دا~$\Gamma$ توں اے۔
  اسیں وکھاواں گے کہ جے $!A_1$،
  \dots، $!A_n = !A$، !!a{derivation} اے~$\Gamma$ توں، تاں $\Gamma
  \Entails !A_n$۔ دھیان دیو کہ جے $!A_1$، \dots، $!A_n$ اک
  !!a{derivation} اے، تاں $!A_1$، \dots، $!A_k$ وی اے، ہر $k<n$ لئی۔

  لمبائی~$0$ دے !!{derivation}s نہیں ہوندے، ایس لئی $n=0$ اُتے دعویٰ
  خالی شرط کارن پورا اے۔ ہُن دعویٰ سارے اوہناں !!{derivation}s لئی قائم اے
  جیہناں دی لمبائی~$<n$ اے۔ اسیں~$!A_n$ دی دلیل دے حساب نال وکھ وکھ صورتاں
  ویکھدے آں۔

  \begin{enumerate}
  \item $!A_n$ مسلّمہ اے۔ سارے مسلّمے معتبر نیں، سو $\Gamma \Entails
    !A_n$ ہر~$\Gamma$ لئی۔
  \item $!A_n \in \Gamma$۔ فیر ہر $\mModel{M}$ تے $w$ لئی، جے
    $\mSat{M}{\Gamma}[w]$، تاں صاف اے کہ $\mSat{M}{\Gamma}{!A_n}[w]$،
    یعنی $\Gamma \Entails !A$۔
  \item $!A_n$، \MP{} نال $!A_i$ تے $!A_j \ident !A_i \lif
    !A_n$ توں نکلدا اے۔ $!A_1$، \dots، $!A_i$ تے $!A_1$، \dots، $!A_j$،
    !!{derivation}s نیں~$\Gamma$ توں۔ سو استقرائی مفروضے مطابق $\Gamma
    \Entails !A_i$ تے $\Gamma \Entails !A_i \lif !A_n$۔

    فرض کرو $\mSat{M}{\Gamma}[w]$۔ چونکہ $\mSat{M}{\Gamma}[w]$ تے
    $\Gamma \Entails !A_i \lif !A_n$، ایس لئی $\mSat{M}{!A_i \lif
      !A_n}[w]$۔ تعریف مطابق ایہدا مطلب اے کہ ہر اوہ $w'$ جیہڑی
    $Rww'$ پوری کرے، اوہدے لئی جے $\mSat{M}{!A_i}[w']$، تاں $\mSat{M}{!A_n}[w']$۔
    چونکہ $R$ انعکاسی اے، $w$ وی اوہناں $w'$ وچ شامل اے جیہڑیاں $Rww'$ پورا
    کردیاں نیں۔ یعنی جے $\mSat{M}{!A_i}[w]$، تاں $\mSat{M}{!A_n}[w]$۔
    چونکہ $\Gamma \Entails !A_i$، ایس لئی $\mSat{M}{!A_i}[w]$۔ سو
    $\mSat{M}{!A_n}[w]$، جیہڑا اسیں وکھانا سی۔
  \end{enumerate}
\end{proof}
\endgroup
% /OLP-0504
% OLP-0505 content/intuitionistic-logic/soundness-completeness/soundness-nd.tex
\begingroup
\def\olfilename{soundness-nd}
\expandafter\def\csname ol@sectioncs\endcsname{section}

\olfileid{int}{sc}{snd}

\olsection{فطری استنباط دی صحت مندی}

ہُن اسیں تعلقی معنیات دے لحاظ نال فطری استنباط دی صحت مندی ثابت کراں گے۔
یعنی وکھاواں گے کہ جے !!a{formula} مفروضیاں دے اک سیٹ توں !!{derivable}
اے، تاں مفروضیاں دے اوس سیٹ دا معنوی نتیجہ اوہ !!{formula} اے۔

\begin{thm}[صحت مندی]
  \ollabel{thm:soundness}
  جے $\Gamma \Proves !A$، تاں $ \Gamma \Entails !A$۔
\end{thm}

\begin{proof}
  اسیں ثابت کراں گے کہ جے $\Gamma \Proves !A$، تاں $\Gamma \Entails !A$۔
  ثبوت~$!A$ دے~$\Gamma$ توں !!{derivation} اُتے استقرا نال اے۔

  \begin{enumerate}
  \item جے !!{derivation} صرف مفروضے~$!A$ دا اے، تاں ساڈے کول $!A \Proves !A$
    اے تے اسیں وکھانا اے کہ $!A \Entails !A$۔ فرض کرو $\mSat{M}{!A}[w]$۔
    فیر ظاہر اے کہ $\mSat{M}{!A}[w]$۔

  \item !!{derivation} دا آخری قاعدہ $\Intro{\land}$ اے۔
    \iftag{probAnd}{مشق۔}{~$!B$ دا !!{undischarged} مفروضیاں~$\Gamma$ توں
      تے~$!C$ دا !!{undischarged} مفروضیاں~$\Delta$ توں مقدمے والے !!{derivation}s
      وکھاندے نیں کہ $\Gamma \Proves !B$ تے $\Delta \Proves !C$۔
      استقرائی مفروضے توں ساڈے کول $\Gamma \Entails !B$ تے $\Delta \Entails !C$ نیں۔
      سانوں وکھانا اے کہ $\Gamma \cup \Delta \Entails !A \land !B$، کیوں جے
      پورے !!{derivation} دے !!{undischarged} مفروضے $\Gamma$ تے $\Delta$ نوں
      رلا کے ملدے نیں۔ سو فرض کرو $\mSat{M}{\Gamma \cup \Delta}[w]$۔
      فیر $\mSat{M}{\Gamma}[w]$ وی اے۔ چونکہ $\Gamma \Entails !B$، ایس لئی
      $\mSat{M}{!B}[w]$۔ ایس طرح $\mSat{M}{!C}[w]$۔ سو
      $\mSat{M}{!B \land !C}[w]$۔}

  \item !!{derivation} دا آخری قاعدہ $\Elim{\land}$ اے۔
    \iftag{probAnd}{مشق۔}{مقدمے $!B \land !C$ دا !!{undischarged} مفروضیاں
      $\Gamma$ توں !!{derivation} وکھاندا اے کہ $\Gamma \Proves !B \land !C$۔
      استقرائی مفروضے توں $\Gamma \Entails !B \land !C$۔
      اسیں وکھانا اے کہ $\Gamma \Entails !B$۔ سو فرض کرو $\mSat{M}{\Gamma}[w]$۔
      چونکہ $\Gamma \Entails !B \land !C$، ایس لئی $\mSat{M}{!B \land !C}[w]$۔
      فیر $\mSat{M}{!B}[w]$ وی اے۔ ایس طرح جے $\Elim\land$ دا نتیجہ~$!C$ اے،
      تاں $\Gamma \Entails !C$۔}

  \item !!{derivation} دا آخری قاعدہ $\Intro{\lor}$ اے۔
    \iftag{probOr}{مشق۔}{فرض کرو مقدمہ $!B$ اے تے~$!B$ اُتے مُکّن والے !!{derivation} دے
      !!{undischarged} مفروضے $\Gamma$ نیں۔
      فیر $\Gamma \Proves !B$ تے استقرائی مفروضے نال $\Gamma \Entails !B$۔
      اسیں وکھانا اے کہ $\Gamma \Entails !B \lor !C$۔ فرض کرو $\mSat{M}{\Gamma}[w]$۔
      چونکہ $\Gamma \Entails !B$، ایس لئی $\mSat{M}{!B}[w]$۔
      پر فیر $\mSat{M}{!B \lor !C}[w]$ وی اے۔ ایس طرح جے مقدمہ~$!C$ اے،
      تاں $\Gamma \Entails !C$۔}

  \item !!{derivation} دا آخری قاعدہ~$\Elim{\lor}$ اے۔
    \iftag{probOr}{مشق۔}{مقدمیاں اُتے مُکّن والے !!{derivation}s ایہ نیں:
      $!B \lor !C$ دا !!{undischarged} مفروضیاں~$\Gamma$ توں،
      $!D$ دا !!{undischarged} مفروضیاں $\Delta_1 \cup \{!B\}$ توں،
      تے $!D$ دا !!{undischarged} مفروضیاں~$\Delta_2 \cup \{!C\}$ توں۔
      سو ساڈے کول $\Gamma \Proves !B \lor !C$، $\Delta_1 \cup \{!B\} \Proves !D$
      تے $\Delta_2 \cup \{!C\} \Proves !D$ نیں۔ استقرائی مفروضے توں
      $\Gamma \Entails !B \lor !C$، $\Delta_1 \cup \{!B\} \Entails !D$ تے
      $\Delta_2 \cup \{!C\} \Entails !D$۔ سانوں ثابت کرنا اے کہ
      $\Gamma \cup \Delta_1 \cup \Delta_2 \Entails !D$۔

    فرض کرو $\mSat{M}{\Gamma \cup \Delta_1 \cup \Delta_2}[w]$۔ فیر
    $\mSat{M}{\Gamma}[w]$ تے چونکہ $\Gamma \Entails !B \lor !C$، ایس لئی
    $\mSat{M}{!B \lor !C}[w]$۔ $\mSat{M}{}$ دی تعریف مطابق، یا
    $\mSat{M}{!B}[w]$ یا $\mSat{M}{!C}[w]$۔ سو دو صورتاں ویکھدے آں۔
    (ا)~$\mSat{M}{!B}$[w]۔ فیر $\mSat{M}{\Delta_1 \cup \{!B\}}[w]$۔
    چونکہ $\Delta_1 \cup !B \Entails !D$، ایس لئی $\mSat{M}{!D}[w]$۔
    (ب)~$\mSat{M}{!C}[w]$۔ فیر $\mSat{M}{\Delta_2 \cup \{!C\}}[w]$۔
    چونکہ $\Delta_2 \cup !C \Entails !D$، ایس لئی $\mSat{M}{!D}[w]$۔
    سو دوناں وچوں جیہڑی وی صورت ہووے، $\mSat{M}{!D}[w]$، جیہڑا اسیں وکھانا سی۔}

  \item !!{derivation} دا آخری قاعدہ $\Intro{\lif}$ اے تے اوہدا نتیجہ~$!B\lif
    !C$ اے۔ فیر مقدمہ~$!C$ اے تے مقدمے اُتے مُکّن والے !!{derivation} دے
    !!{undischarged} مفروضے~$\Gamma \cup \{!B\}$ نیں۔ سو
    $\Gamma \cup \{!B\} \Proves !C$ تے استقرائی مفروضے نال
    $\Gamma \cup \{!B\} \Entails !C$۔ اسیں وکھانا اے کہ $\Gamma \Entails !B \lif !C$۔

    فرض کرو $\mSat{M}{\Gamma}[w]$۔ اسیں وکھانا چاہندے آں کہ ہر اوہ $w'$ جیہڑی
    $Rww'$ پوری کرے، اوہدے لئی جے $\mSat{M}{!B}[w']$، تاں $\mSat{M}{!C}[w']$۔
    سو منّ لو $Rww'$ تے $\mSat{M}{!B}[w']$۔ \olref[sem][rel]{prop:true-monotonic}
    مطابق، $\mSat{M}{\Gamma}[w']$۔ چونکہ $\Gamma \cup \{!B\} \Entails !C$،
    ایس لئی $\mSat{M}{!C}[w']$، جیہڑا اسیں وکھانا سی۔

  \item !!{derivation} دا آخری قاعدہ $\Elim{\lif}$ اے تے نتیجہ~$!C$ اے۔
    مقدمے $!B \lif !C$ تے $!B$ نیں، جیہناں دے !!{derivation}s بالترتیب
    !!{undischarged} مفروضیاں $\Gamma$، $\Delta$ توں نیں۔ سو
    $\Gamma \Proves !B \lif !C$ تے $\Delta \Proves !B$۔ استقرائی مفروضے توں
    $\Gamma \Entails !B \lif !C$ تے $\Delta \Entails !B$۔
    اسیں وکھانا اے کہ $\Gamma \cup \Delta \Entails !C$۔

    فرض کرو $\mSat{M}{\Gamma \cup \Delta}[w]$۔ چونکہ
    $\mSat{M}{\Gamma}[w]$ تے $\Gamma \Entails !B \lif !C$، ایس لئی
    $\mSat{M}{!B \lif !C}[w]$۔ تعریف مطابق ایہدا مطلب اے کہ ہر اوہ
    $w'$ جیہڑی $Rww'$ پوری کرے، اوہدے لئی جے $\mSat{M}{!B}[w']$ تاں
    $\mSat{M}{!C}[w']$۔ چونکہ $R$ انعکاسی اے، $w$ وی اوہناں $w'$
    وچ اے جیہڑیاں $Rww'$ پورا کردیاں نیں۔ یعنی جے $\mSat{M}{!B}[w]$، تاں
    $\mSat{M}{!C}[w]$۔ چونکہ $\mSat{M}{\Delta}[w]$ تے $\Delta
    \Entails !B$، ایس لئی $\mSat{M}{!B}[w]$۔ سو $\mSat{M}{!C}[w]$،
    جیہڑا اسیں وکھانا سی۔

  \item !!{derivation} دا آخری قاعدہ $\FalseInt$ اے تے اوہدا نتیجہ~$!A$ اے۔
    مقدمہ $\lfalse$ اے تے اوہدے !!{derivation} دے !!{undischarged} مفروضے~$\Gamma$
    نیں۔ فیر $\Gamma \Proves \lfalse$۔ استقرائی مفروضے توں
    $\Gamma \Entails \lfalse$۔ اسیں وکھانا اے کہ $\Gamma \Entails !A$۔

    اسیں بالواسطہ ثبوت دیندے آں۔ جے $\Gamma \Entails/ !A$، تاں اک نمونہ~$\mModel{M}$
    تے جہان~$w$ اے جتھے $\mSat{M}{\Gamma}[w]$ تے $\mSat/{M}{!A}[w]$۔
    چونکہ $\Gamma \Entails \lfalse$، ایس لئی $\mSat{M}{\lfalse}[w]$۔
    پر ایہ ناممکن اے، کیوں جے تعریف مطابق $\mSat/{M}{\lfalse}[w]$۔
    سو $\Gamma \Entails !A$۔
  \item !!{derivation} دا آخری قاعدہ $\Intro\lnot$ اے۔ مشق۔
  \item !!{derivation} دا آخری قاعدہ $\Elim\lnot$ اے۔ مشق۔
  \end{enumerate}
\end{proof}

\begin{prob}
  \olref[int][sc][snd]{thm:soundness} دا ثبوت مکمل کرو۔ $\Intro\lnot$ تے
  $\Elim\lnot$ دیاں صورتاں لئی $\mSat{M}{\lnot !A}[w]$ دی تعریف ورتو
  جیہڑی \olref[int][sem][rel]{defn:true-at-w} وچ دتی اے۔ یعنی $\lnot !A$
  نوں $!A \lif \lfalse$ نال تعریف شدہ نہ منّو۔
\end{prob}

\begin{prob}
  وکھاؤ کہ ہیٹھلے !!{formula}s وجدانیت پسند منطق وچ !!{derivable} نہیں۔
  \begin{enumerate}
    \item $(!A \lif !B) \lor (!B \lif !A)$
    \item $(\lnot\lnot !A \lif !A) \lif (!A \lor \lnot !A)$
    \item $(!A \lif !B \lor !C) \lif \bigl((!A \lif !B)\lor(!A \lif !C)\bigr)$
  \end{enumerate}
\end{prob}
\endgroup
% /OLP-0505
% OLP-0506 content/intuitionistic-logic/soundness-completeness/lindenbaum.tex
\begingroup
\def\olfilename{lindenbaum}
\expandafter\def\csname ol@sectioncs\endcsname{section}

\olfileid{int}{sc}{lin}

\olsection{لنڈن باؤم دا معاون قضیہ}

وجدانیت پسند منطق دی تمامیت دا قضیہ $\Gamma \Proves/ !A$ منّ کے ثابت
ہوندا اے۔ اسیں اک نمونہ بناندے آں جتھے $\mSat{M}{\Gamma}$ تے~$\mSat/{M}{!A}$۔

کلاسیکی منطق وچ !!{derivability} دا تعلق بے تضادی دی دھارنا راہیں سمجھایا
جا سکدا اے۔ کیوں جے !!a{formula}~$!A$، !!{formula}s دے کسے سیٹ توں !!{derivable}
اُدوں تے صرف اُدوں اے جدوں~$!A$ دی نفی نوں اُس سیٹ نال رلا کے بنا سیٹ متضاد
ہووے۔ وجدانیت پسند منطق وچ ایہ ممکن نہیں۔ ایس منطق وچ جے $\lnot!A$ متضاد
اے، تاں صرف $\Proves \lnot\lnot !A$ ملدا اے۔ چونکہ $\lnot\lnot!A \lif !A$
عام طور اُتے وجدانیت پسند منطق وچ نہیں چلدا، ایس لئی اسیں~$\Proves !A$ دا
نتیجہ نہیں کڈھ سکدے۔

سو نمونہ~$\mModel{M}$ بناندیاں اسیں !!{formula}~$!A$ دی !!{derivability} نہ ہون دی شرط
سامھنے رکھنی اے۔ ایس لئی اسیں اک مکمل سیٹ $\Gamma^* \supseteq \Gamma$ نوں
نمونہ $\mModel{M}$ بناؤن لئی نہیں ورت سکدے،
کیوں جے ہر مکمل سیٹ $\Gamma^*$ وچ $\Gamma^* \Proves !A \lor \lnot !A$ ہوندا اے۔

مکمل سیٹ~$\Gamma^*$ دی تھاں اسیں فارمولیاں دے اولی سیٹ دی دھارنا ورتاں گے۔

\begin{defn}\ollabel{defn:prime}
  !!{formula}s دا سیٹ~$\Gamma$ \emph{اولی} اُدوں تے صرف اُدوں اے جدوں
  \begin{enumerate}
  \item\ollabel{defn:prime1} $\Gamma$ بے تضاد اے، یعنی $\Gamma \Proves/ \lfalse$۔
  \item\ollabel{defn:prime2} جے $\Gamma \Proves !A$، تاں $!A \in \Gamma$۔
  \item\ollabel{defn:prime3} جے $!A \lor !B \in \Gamma$، تاں $!A \in
    \Gamma$ یا $!B \in \Gamma$۔
  \end{enumerate}
\end{defn}

\begin{lem}[لنڈن باؤم دا معاون قضیہ]
  \ollabel{lem:lindenbaum} جے $\Gamma \Proves/ !A$، تاں اک
  $\Gamma^* \supseteq \Gamma$ موجود اے جیہدے لئی $\Gamma^*$ اولی اے تے
  $\Gamma^* \Proves/ !A$۔
\end{lem}

\begin{proof}
  منّ لو $!B_1 \lor !C_1$، $!B_2 \lor !C_2$، \dots، صورت~$!B \lor !C$
  والے سارے !!{formula}s دی گنتی وار فہرست اے۔ اسیں !!{formula}s دے سیٹاں
  دی اک ودھدی لڑی~$\Gamma_n$ دی تعریف کراں گے، جیہدے وچ ہر $\Gamma_{n+1}$،
  $\Gamma_n$ نال اک نواں !!{formula} رلا کے بنے گا۔ $\Gamma^*$ سارے~$\Gamma_n$
  دا اتحاد ہووے گا۔ نویں !!{formula}s ایویں چنے جان گے کہ $\Gamma^*$ اولی ہووے
  تے $\Gamma^* \Proves/ !A$ فیر وی قائم رہے۔ یعنی ہر قدم اُتے اسیں پہلی فصل
  ~$!B_i \lor !C_i$ لبھنی اے جیہڑی ایہ شرطاں پوریاں کرے:
  \begin{enumerate}
  \item\ollabel{gamma-1} $\Gamma_n \Proves !B_i \lor !C_i$
  \item\ollabel{gamma-2} $!B_i \notin \Gamma_n$ تے $!C_i \notin \Gamma_n$
  \end{enumerate}
  اسیں $\Gamma_n$ وچ $!B_i$ شامل کردے آں جے $\Gamma_n \cup \{!B_i\}
  \Proves/ !A$، تے ہور صورت وچ $!C_i$۔ سانوں وکھانا پئے گا کہ ایہ طریقہ چلدا اے۔
  ہُن لئی $i(n)$ نوں سبھ توں چھوٹا اوہ $i$ منّ لو جیہدے لئی
  \olref{gamma-1} تے~\olref{gamma-2} پوریاں ہون۔

  $\Gamma_0 = \Gamma$ رکھو تے
  \[
  \Gamma_{n+1} = \begin{cases}
    \Gamma_n \cup \{!B_{i(n)}\} &
    \text{جے $\Gamma_n \cup \{!B_{i(n)}\} \Proves/ !A$} \\
    \Gamma_n \cup \{!C_{i(n)}\} & \text{ہور صورت وچ}
    \end{cases}
  \]
  جے $i(n)$ تعریف شدہ نہ ہووے، یعنی جدوں وی $\Gamma_n \Proves !B \lor !C$،
  اوہدوں یا $!B \in \Gamma_n$ یا $!C \in \Gamma_n$، تاں اسیں $\Gamma_{n+1}
  = \Gamma_n$ رکھدے آں۔ ہُن منّ لو $\Gamma^* = \bigcup_{n=0}^\infty \Gamma_n$۔

  پہلاں اسیں وکھاندے آں کہ ہر $n$ لئی $\Gamma_n \Proves/ !A$۔ ثبوت~$n$
  اُتے استقرا نال اے۔ $n = 0$ لئی دعویٰ قضیے دے مفروضے توں قائم اے،
  یعنی $\Gamma \Proves/ !A$۔ جے $n>0$، تاں اسیں وکھانا اے کہ جے
  $\Gamma_n \Proves/ !A$، تاں $\Gamma_{n+1} \Proves/ !A$۔ جے $i(n)$ تعریف
  شدہ نہیں، تاں $\Gamma_{n+1} = \Gamma_n$ تے ثابت کرن نوں کجھ نہیں۔
  سو فرض کرو $i(n)$ تعریف شدہ اے۔ سوکھ لئی $i = i(n)$ رکھ لو۔

  اسیں دعوے دا عکسِ نقیض ثابت کراں گے۔ فرض کرو $\Gamma_{n+1} \Proves !A$۔
  بناوٹ مطابق، $\Gamma_{n+1} = \Gamma_n \cup \{!B_i\}$ جے
  $\Gamma_n \cup \{!B_i\} \Proves/ !A$، تے ہور صورت وچ
  $\Gamma_{n+1} = \Gamma_n \cup \{!C_i\}$۔ صاف اے کہ پہلی صورت نہیں ہو سکدی،
  کیوں جے فیر $\Gamma_{n+1} \Proves/ !A$ ہوندا۔ سو
  $\Gamma_n \cup \{!B_i\} \Proves !A$ تے $\Gamma_{n+1} = \Gamma_n \cup \{!C_i\}$۔
  $i(n)$ دی تعریف مطابق، $\Gamma _n \Proves !B_i \lor !C_i$۔
  ساڈے کول $\Gamma_n \cup \{!B_i\} \Proves !A$ اے تے
  $\Gamma_{n+1} = \Gamma_n \cup \{!C_i\} \Proves !A$ وی اے۔
  سو $\Gamma_n \Proves !A$، جیہڑا اسیں وکھانا سی۔

  جے $\Gamma^* \Proves !A$ ہوندا، تاں کوئی محدود ذیلی سیٹ $\Gamma'
  \subseteq \Gamma^*$ ہوندا جیہدے لئی $\Gamma' \Proves !A$۔
  ہر $!D \in \Gamma'$ لازماً $\Gamma_i$ وچ ہووے گا، کسے~$i$ لئی۔
  ایہناں وچوں سبھ توں وڈا $n$ منّ لو۔ چونکہ $\Gamma_i \subseteq \Gamma_n$
  جے $i \le n$، ایس لئی $\Gamma' \subseteq \Gamma_n$۔ پر فیر
  $\Gamma_n \Proves !A$، جیہڑا اُتلے ثبوت دے اُلٹ اے کہ $\Gamma_n \Proves/ !A$۔

  آخر وچ اسیں وکھاندے آں کہ $\Gamma^*$ اولی اے، یعنی
  \olref{defn:prime1}، \olref{defn:prime2} تے~\olref{defn:prime3} شرطاں
  پوری کردا اے، جیہڑیاں \olref{defn:prime} وچ نیں۔

  پہلاں، $\Gamma^* \Proves/ !A$، ایس لئی $\Gamma^*$ بے تضاد اے۔
  سو \olref{defn:prime1} پوری ہوندی اے۔

  ہُن اسیں وکھاندے آں کہ جے $\Gamma^* \Proves !B \lor !C$، تاں یا
  $!B \in \Gamma^*$ یا $!C \in \Gamma^*$۔ ایس نال \olref{defn:prime3}
  ثابت ہوندی اے، کیوں جے جے $!B \lor !C \in \Gamma^*$، تاں
  $\Gamma^* \Proves !B \lor !C$ وی۔ سو فرض کرو $\Gamma^* \Proves !B \lor !C$
  پر $!B \notin \Gamma^*$ تے $!C \notin \Gamma^*$۔ چونکہ $\Gamma^* \Proves
  !B \lor !C$، ایس لئی $\Gamma_n \Proves !B \lor !C$، کسے~$n$ لئی۔
  $!B \lor !C$ ساریاں فلاں دی گنتی وار فہرست وچ اے، منّ لو $!B_j \lor !C_j$
  دے طور اُتے۔ $!B_j \lor !C_j$، $i(n)$ دی تعریف دیاں شرطاں پوریاں کردا اے۔
  یعنی $\Gamma_n \Proves !B_j \lor !C_j$، جدکہ $!B_j \notin \Gamma_n$
  تے $!C_j \notin \Gamma_n$۔ ہر قدم اُتے ایہ شرطاں پوریاں کرن والے
  $!B_i \lor !C_i$ وچوں گھٹ توں گھٹ اک گھٹ جاندا اے (کیوں جے ہر قدم اُتے اسیں
  یا $!B_i$ یا $!C_i$ شامل کردے آں)، ایس لئی کسے قدم~$m$ اُتے $j = i(m)$ ہووے گا۔
  پر فیر یا $!B \in \Gamma_{m+1}$ یا $!C \in \Gamma_{m+1}$، جیہڑا
  ایس مفروضے دے اُلٹ اے کہ $!B \notin \Gamma^*$ تے $!C \notin \Gamma^*$۔

  ہُن فرض کرو $\Gamma^* \Proves !B$۔ فیر $\Gamma^* \Proves !B \lor !B$۔
  پر اسیں ہُنے ثابت کیتا اے کہ جے $\Gamma^* \Proves !B \lor !B$، تاں
  $!B \in \Gamma^*$۔ سو $\Gamma^*$، شرط~\olref{defn:prime2}
  پوری کردا اے، جیہڑی~\olref{defn:prime} وچ اے۔
\end{proof}

\begin{prob}
وکھاؤ کہ جے $\Gamma \Proves/ \bot$، تاں $\Gamma$ کلاسیکی منطق وچ بے تضاد
اے۔ یعنی اک ایہو جیہی قدر بندی موجود اے جیہڑی~$\Gamma$ دے سارے فارمولیاں
نوں سچ کردی اے۔
\end{prob}
\endgroup
% /OLP-0506
% OLP-0507 content/intuitionistic-logic/soundness-completeness/canonical-model.tex
\begingroup
\def\olfilename{canonical-model}
\expandafter\def\csname ol@sectioncs\endcsname{section}

\olfileid{int}{sc}{mod}

\olsection{کینونیکل نمونہ}

ساڈے نمونے دے جہان قدرتی عدداں دیاں محدود لڑیاں~$\sigma$ ہون گیاں،
یعنی $\sigma \in \Nat^*$۔ دھیان دیو کہ $\Nat^*$ دی تعریف استقرائی طور اُتے ایویں اے:
\begin{enumerate}
\item $\emptyseq \in \Nat^*$۔
\item جے $\sigma \in \Nat^*$ تے $n \in \Nat$، تاں $\sigma.n \in
  \Nat^*$ (جتھے $\sigma.n$، $\sigma \concat \tuple{n}$ اے تے $\sigma \concat \sigma'$،
  $\sigma$ تے $\sigma'$ دا جوڑ اے)۔
\item $\Nat^*$ وچ ہور کجھ نہیں۔
\end{enumerate}
سو اسیں $\Nat^*$ ورت کے استقرائی تعریفاں دے سکدے آں۔

منّ لو $\tuple{!B_1, !C_1}$، $\tuple{!B_2, !C_2}$، \dots، !!{formula}s
دیاں ساریاں جوڑیاں دی گنتی وار فہرست اے۔ !!{formula}s دا سیٹ~$\Delta$ دتا ہووے،
تاں $\Delta(\sigma)$ دی تعریف استقرا نال ایویں کرو:
\begin{enumerate}
\item $\Delta(\emptyseq) = \Delta$
\item $\Delta(\sigma.n) = {}$
  \[
  \begin{cases}
    (\Delta(\sigma) \cup \{!B_n\})^* &
    \text{جے $\Delta(\sigma) \cup \{!B_n\} \Proves/ !C_n$} \\
    \Delta(\sigma) & \text{ہور صورت وچ}
  \end{cases}
  \]
\end{enumerate}
ایتھے $(\Delta(\sigma) \cup \{!B_n\})^*$ توں مراد !!{formula}s دا اوہ اولی سیٹ اے
جیہڑا \olref[lin]{lem:lindenbaum} نوں سیٹ $\Delta(\sigma) \cup \{!B_n\}$ تے
!!{formula}~$!C_n$ اُتے لاؤن نال ملدا اے۔ دھیان دیو کہ ایس تعریف مطابق جے
$\Delta(\sigma) \cup \{!B_n\} \Proves/ !C_n$، تاں $\Delta(\sigma.n)
\Proves !B_n$ تے $\Delta(\sigma.n) \Proves/ !C_n$۔ ایہ وی دھیان دیو کہ
$\Delta(\sigma) \subseteq \Delta(\sigma.n)$ ہر~$n$ لئی اے۔ جے $\Delta$
اولی اے، تاں $\Delta(\sigma)$ ہر~$\sigma$ لئی اولی اے۔

\begin{defn}\ollabel{defn:canonical-model}
  فرض کرو $\Delta$ اولی اے۔ فیر \emph{کینونیکل نمونہ} $\mModel{M(\Delta)}$،
  $\Delta$ لئی، ایویں تعریف ہوندا اے:
  \begin{enumerate}
  \item $W = \Nat^*$، قدرتی عدداں دیاں محدود لڑیاں دا سیٹ۔
  \item $R$ اوہ جزوی ترتیب اے جیہدے مطابق $R\sigma\sigma'$ اُدوں تے صرف اُدوں
    اے جدوں $\sigma$،~$\sigma'$ دا مُڈھلا ٹکڑا ہووے (یعنی $\sigma' =
    \sigma \concat \sigma''$ کسے لڑی~$\sigma''$ لئی)۔
  \item $V(p) = \Setabs{\sigma}{p \in \Delta(\sigma)}$۔
  \end{enumerate}
\end{defn}

ایہ پرکھنا سوکھا اے کہ $R$ واقعی اک جزوی ترتیب اے۔~$V$ دے نہ گھٹّن دی
شرط وی پوری اے۔ چونکہ $\Delta(\sigma) \subseteq \Delta(\sigma.n)$، ایس لئی
سانوں $\Delta(\sigma) \subseteq \Delta(\sigma')$ ملدا اے جدوں وی
$R\sigma\sigma'$ ہووے،~$\sigma$ اُتے استقرا نال۔
\endgroup
% /OLP-0507
% OLP-0508 content/intuitionistic-logic/soundness-completeness/truth-lemma.tex
\begingroup
\def\olfilename{truth-lemma}
\expandafter\def\csname ol@sectioncs\endcsname{section}

\olfileid{int}{sc}{tru}

\olsection{سچائی دا معاون قضیہ}

ہیٹھلا معاون قضیہ کینونیکل نمونے وچ پورا اترن دی گل نوں ایس گل نال جوڑدا اے
کہ کیہڑے !!{formula}s اولی سیٹ $\Delta$ دے !!{element}s نیں۔

\begin{lem}\ollabel{lem:truth}
  جے $\Delta$ اولی اے، تاں $\mSat{M(\Delta)}{!A}[\sigma]$ اُدوں تے صرف اُدوں
  اے جدوں $\Delta(\sigma) \Proves !A$۔
\end{lem}

\begin{proof}
  ~$!A$ اُتے استقرا نال۔
  \begin{enumerate}
    \item \indcase{!A}{\lfalse}{چونکہ $\Delta(\sigma)$ اولی اے، ایس لئی ایہ
      بے تضاد اے، سو $\Delta(\sigma) \Proves/ \indfrm$۔ تعریف مطابق،
      $\mSat/{M(\Delta)}{\indfrm}[\sigma]$۔}
      \item \indcase{!A}{p}{$\mSat{{}}{}$ دی تعریف مطابق،
        $\mSat{M(\Delta)}{\indfrm}[\sigma]$ اُدوں تے صرف اُدوں اے جدوں
        $\sigma \in V(p)$، یعنی $\Delta(\sigma) \Proves \indfrm$۔}
      \item \indcase!{!A}{\lnot !B}{}
      \item \indcase{!A}{!B \land
        !C}{$\mSat{M(\Delta)}{\indfrm}[\sigma]$ اُدوں تے صرف اُدوں اے جدوں
        $\mSat{M(\Delta)}{!B}[\sigma]$ تے $\mSat{M(\Delta)}{!C}[\sigma]$۔
        استقرائی مفروضے مطابق $\mSat{M(\Delta)}{!B}[\sigma]$ اُدوں تے صرف
        اُدوں اے جدوں $\Delta(\sigma) \Proves !B$، تے~$!C$ لئی وی ایسے طرح۔
        پر $\Delta(\sigma) \Proves !B$ تے $\Delta(\sigma) \Proves !C$
        دونیں اُدوں تے صرف اُدوں نیں جدوں $\Delta(\sigma) \Proves \indfrm$۔}
      \item \indcase{!A}{!B \lor
        !C}{$\mSat{M(\Delta)}{\indfrm}[\sigma]$ اُدوں تے صرف اُدوں اے جدوں
        $\mSat{M(\Delta)}{!B}[\sigma]$ یا $\mSat{M(\Delta)}{!C}[\sigma]$۔
        استقرائی مفروضے مطابق ایہ اُدوں تے صرف اُدوں اے جدوں
        $\Delta(\sigma) \Proves !B$ یا $\Delta(\sigma) \Proves !C$۔
        اسیں وکھانا اے کہ ایہ اگوں اُدوں تے صرف اُدوں اے جدوں
        $\Delta(\sigma) \Proves \indfrm$۔ کھبے توں سجے پاسے دا قدم صاف اے۔
        سجے توں کھبے پاسے دا قدم ایس توں نکلدا اے کہ $\Delta(\sigma)$ اولی اے۔}
      \item \indcase{!A}{!B \lif !C}{پہلاں کھبے توں سجے پاسے والے قدم دا عکسِ نقیض۔
        فرض کرو $\Delta(\sigma) \Proves/ !B \lif !C$۔ فیر
        $\Delta(\sigma) \cup \{!B\} \Proves/ !C$ وی۔ چونکہ $\tuple{!B, !C}$،
        $\tuple{!B_n, !C_n}$ اے، کسے~$n$ لئی، ایس لئی $\Delta(\sigma.n) =
        (\Delta(\sigma) \cup \{!B\})^*$، تے $\Delta(\sigma.n) \Proves !B$
        پر $\Delta(\sigma.n) \Proves/ !C$۔ استقرائی مفروضے مطابق
        $\mSat{M(\Delta)}{!B}[\sigma.n]$ تے $\mSat/{M(\Delta)}{!C}[\sigma.n]$۔
        چونکہ $R\sigma(\sigma.n)$، ایس لئی $\mSat/{M(\Delta)}{\indfrm}[\sigma]$۔

        ہُن فرض کرو $\Delta(\sigma) \Proves !B \lif !C$ تے منّ لو
        $R\sigma\sigma'$۔ چونکہ $\Delta(\sigma) \subseteq \Delta(\sigma')$،
        ایس لئی جے $\Delta(\sigma') \Proves !B$، تاں $\Delta(\sigma') \Proves !C$۔
        دوجے لفظاں وچ ہر اوہ $\sigma'$ جیہڑی $R\sigma\sigma'$ پوری کرے،
        اوہدے لئی یا $\Delta(\sigma') \Proves/ !B$ یا $\Delta(\sigma') \Proves !C$۔
        استقرائی مفروضے مطابق ایہدا مطلب اے کہ جدوں وی $R\sigma\sigma'$،
        یا $\mSat/{M(\Delta)}{!B}[\sigma']$ یا $\mSat{M(\Delta)}{!C}[\sigma']$،
        یعنی $\mSat{M(\Delta)}{\indfrm}[\sigma]$۔}
  \end{enumerate}
\end{proof}
\endgroup
% /OLP-0508
% OLP-0509 content/intuitionistic-logic/soundness-completeness/completeness-thm.tex
\begingroup
\def\olfilename{completeness-thm}
\expandafter\def\csname ol@sectioncs\endcsname{section}

\olfileid{int}{sc}{cpl}

\olsection{تمامیت دا قضیہ}

\begin{thm}\ollabel{thm:completeness}
  جے $\Gamma \Entails !A$، تاں $\Gamma \Proves !A$۔
\end{thm}

\begin{proof}
  اسیں عکسِ نقیض ثابت کردے آں۔ فرض کرو $\Gamma \Proves/ !A$۔ فیر
  \olref[lin]{lem:lindenbaum} مطابق اک اولی سیٹ~$\Gamma^* \supseteq \Gamma$
  اے جیہدے لئی $\Gamma^* \Proves/ !A$۔ کینونیکل نمونہ~$\mModel{M(\Gamma^*)}$،
  $\Gamma^*$ لئی، \olref[mod]{defn:canonical-model} دی تعریف مطابق لو۔
  ہر $!B \in \Gamma$ لئی، $\Gamma^* \Proves !B$۔ دھیان دیو کہ
  $\Gamma^*(\emptyseq) = \Gamma^*$۔ سچائی دے معاون قضیے
  (\olref[tru]{lem:truth}) توں $\mSat{M(\Gamma^*)}{!B}[\emptyseq]$
  ہر $!B \in \Gamma$ لئی اے، تے $\mSat/{M(\Gamma^*)}{!A}[\emptyseq]$۔
  ایس توں پتا لگدا اے کہ $\Gamma \Entails/ !A$۔
\end{proof}

\begin{prob}
  وکھاؤ کہ جے $!A$ وچ صرف !!{propositional variable}s، $\lor$ تے~$\land$
  نیں، تاں $\Entails/ !A$۔ ایس توں نتیجہ کڈھو کہ $\lif$ دی تعریف
  وجدانیت پسند منطق وچ $\lor$ تے~$\land$ توں نہیں ہو سکدی۔
\end{prob}

\begin{prob}
  تمامیت دا قضیہ ورت کے ثابت کرو کہ جے $\Proves !A \lor !B$، تاں
  $\Proves !A$ یا $\Proves !B$۔ (اشارہ۔ $\mSat/{M_1}{!A}$ تے
  $\mSat/{M_2}{!B}$ منّ لو تے اک نواں نمونہ $\mModel{M}$ ایویں بناؤ کہ
  $\mSat/{M}{!A \lor !B}$۔)
\end{prob}

\begin{prob}
  وکھاؤ کہ جے $\mModel{M}$ اک تعلقی نمونہ اے جیہڑا خطی ترتیب ورتدا اے،
  تاں $\mSat{M}{(!A \lif !B)\lor(!B \lif !A)}$۔
\end{prob}
\endgroup
% /OLP-0509
% OLP-0510 content/intuitionistic-logic/soundness-completeness/decidability.tex
\begingroup
\def\olfilename{decidability}
\expandafter\def\csname ol@sectioncs\endcsname{section}

\olfileid{int}{sc}{dec}

\olsection{فیصلہ کرن دی اہلیت}
دھیان دیو کہ تمامیت دے قضیے دا ثبوت ہر $\Gamma \Proves/ !A$ لئی سانوں اک
لامتناہی گنتی دے جہاناں والا نمونہ دیندا اے، جیہڑا ایس گل دا گواہ اے کہ
$\Gamma \Entails/ !A$۔ ہیٹھلا بیان وکھاندا اے کہ $\Entails !A$ ثابت کرن
لئی ایہ ثابت کرنا کافی اے کہ $\mSat{M}{!A}$ سارے محدود نمونیاں لئی اے
(یعنی اوہ نمونے جیہناں دے جہاناں دا سیٹ محدود اے)۔

\begin{thm}\ollabel{thm:decidability}
  جے $\Entails/ !A$، تاں اک محدود نمونہ اے جیہدے وچ $\mSat/{M'}{!A}$۔
\end{thm}

\begin{proof}
  منّ لو $\mModel{M}=\tuple{W, R, V}$ ایہو جیہا اے کہ $\mSat/{M}{!A}$ تے $P$
  اوہناں !!{propositional variable}s دا سیٹ اے جیہڑے~$!A$ وچ آندے نیں۔
  $\mModel{M'}=\tuple{W', R', V'}$ دی تعریف لئی $W' = \Setabs{[w]}{w \in W}$
  رکھو، جتھے $[w] = \Setabs{p \in P}{w \in V(p)}$، $R'$ نوں ذیلی سیٹ ہون
  دا تعلق رکھو تے $V'(p)=\Setabs{[w]}{p \in [w]}$ رکھو۔ ایہ صاف ہونا چاہیدا
  اے کہ $W'$ اک محدود سیٹ اے تے $\mModel{M'}$ !!a{relational model} اے۔

  ~$!A$ اُتے استقرا نال وکھایا جا سکدا اے کہ
  \[
    \mSat{M}{!A}[w] \text{ اُدوں تے صرف اُدوں جدوں } \mSat{M'}{!A}[{[w]}]
  \]
  اوہناں سارے !!{formula}s $!A$ لئی جیہناں دے !!{propositional variable}s
  صرف~$P$ وچوں ہون۔ ایہ پڑھّن والے لئی مشق چھڈّی اے۔
\end{proof}

\begin{prob}
\olref[int][sc][dec]{thm:decidability} دا ثبوت مکمل کرو۔ وکھاؤ کہ
$\mSat{M,w}{!A}$ اُدوں تے صرف اُدوں اے جدوں $\mSat{M',[w]}{!A}$، اوہناں
سارے !!{formula}s~$!A$ لئی جیہناں دے قضیاتی متغیر صرف~$P$ وچوں ہون۔
\end{prob}

\olref{thm:decidability} توں ایہ نتیجہ نکلدا اے کہ فیصلہ کرن لئی اک الگورتھم
موجود اے کہ آیا~$\Entails !A$۔
\endgroup
% /OLP-0510

\OLEndChapterHook
\endgroup
% /OLP-0503

% OLP-0511 content/intuitionistic-logic/tableaux/tableaux.tex
\begingroup
\def\olfilename{tableaux}
\expandafter\def\csname ol@sectioncs\endcsname{section}

\olchapter{int}{tab}{وجدانیت پسند \usetoken{P}{tableau}}

\begin{editorial}
  وجدانیت پسند منطق لئی اگیتیاں والے ٹیبلاں دا مسودہ۔ ہور مثالاں،
  تمامیت دے ثبوت، تے ایہ وکھاؤن دی لوڑ اے کہ بند ٹیبلاں دی ناکام تلاش
  توں ردّی نمونے کیویں لبھے جا سکدے نیں۔
\end{editorial}

% OLP-0512 content/intuitionistic-logic/tableaux/introduction.tex
\begingroup
\def\olfilename{introduction}
\expandafter\def\csname ol@sectioncs\endcsname{section}

\olfileid{int}{tab}{int}

\olsection{جان پہچان}

!!^{tableau}s !!{signed formula}s دے خاص درخت نیں، جیہناں دیاں ٹہنیاں تھلّے ول ودھدیاں نیں۔
یعنی ہر نوڈ اک جوڑی اے، سچائی دی قدر دا نشان ($\True$ یا $\False$) تے
!!a{sentence}:
\[
\sFmla{\True}{!A} \text{ یا } \sFmla{\False}{!A}.
\]
!!^a{tableau} کئی \emph{مفروضیاں} نال شروع ہوندا اے۔ ہر اگلا !!{signed formula}
استنباط دے قاعدیاں وچوں اک نوں لا کے بندا اے۔ کجھ قاعدے درخت دے کسے
سِرے اُتے اک یا ودھ !!{signed formula}s پا دیندے نیں۔ کجھ ہور دو نویں سِرے
بناندے نیں تے دو ٹہنیاں نکلدیاں نیں۔ قاعدے توں ملے !!{signed formula}s وچ
!!{formula} اوس !!{signed formula} دے فارمولے توں گھٹ گنجلدار ہوندا اے جیہدے
اُتے قاعدہ لگیا سی۔ جدوں ٹہنی وچ $\sFmla{\True}{!A}$ تے
$\sFmla{\False}{!A}$ دونیں ہون، اسیں ٹہنی نوں \emph{بند} آکھدے آں۔
جے !!a{tableau} دی ہر ٹہنی بند اے، تاں سارا !!{tableau} بند اے۔ بند
!!{tableau} اک !!a{derivation} اے جیہڑا وکھاندا اے کہ !!{tableau} دے شروع
وچ ورتے !!{signed formula}s دا سیٹ اکٹھا پورا ہون جوگا نہیں۔
ایس نال $\Proves$ تعلق دی تعریف کیتی جا سکدی اے۔
$\Gamma \Proves !A$ اُدوں تے صرف اُدوں اے جدوں اک محدود سیٹ~$\Gamma_0 = \{!B_1,
\dots, !B_n\} \subseteq \Gamma$ ہووے جیہدے لئی ہِیٹھلے مفروضیاں دا بند
!!{tableau} ہووے:
\[
\{\sFmla{\False}{!A}, \sFmla{\True}{!B_1}, \dots, \sFmla{\True}{!B_n}\}.
\]

وجدانیت پسند منطق لئی سانوں !!{signed formula} دی دھارنا ودھاؤنی پئے گی
تے رابطیاں دے قاعدے وی بدلنے پین گے۔ نشان ($\True$ یا $\False$) دے نال نال،
امکانی !!{tableau}s وچ !!{formula}s دے \emph{اگیتے}~$\sigma$ وی ہوندے نیں۔
اگیتے مثبت صحیح عدداں دیاں غیر خالی لڑیاں نیں، یعنی $\sigma \in
(\PosInt)^* \setminus \{\emptyseq\}$۔ اسیں ایہو جیہے اگیتے آلے دوالے
$\tuple{\ }$ توں بناں لکھدے آں، تے وکھرے وکھرے !!{element}s نوں~$.$ نال وکھ
کردے آں، $,$ نال نہیں۔ جے $\sigma$ اک اگیتا اے، تاں $\sigma.n$، $\sigma \concat \tuple{n}$
اے۔ مثلاً جے $\sigma = 1.2.1$، تاں $\sigma.3$، $1.2.1.3$ اے۔ سو مثال لئی،
\[
\sFmla{\True}{!A \lif (!B \lif !C)}[1.2]
\]
اک \emph{اگیتے والا !!{signed formula}} اے (یا مختصر طور اُتے
\emph{اگیتے والا !!{formula}})۔

سمجھن لئی اگیتے نوں نمونے وچ اک جہان دا ناں سمجھ لو۔ ایہ نمونہ
!!a{tableau} دی اک ٹہنی والے !!{formula}s پورے کر سکدا اے۔ جے $\sigma$
کسے جہان دا ناں اے، تاں $\sigma.n$ اوہ جہان دا ناں اے جیہڑا
~$\sigma$ ناں والے جہان توں رسائی وچ اے۔

وجدانیت پسند نمونیاں وچ رسائی دا تعلق انعکاسی تے متعدی اے۔ اگیتیاں دے
لحاظ نال ایہدا مطلب اے کہ $\sigma$، اپنے آپ~$\sigma$ توں رسائی وچ اے،
تے ہر اوہ اگیتا وی جیہڑا~$\sigma$ نوں ودھاوے، یعنی
$\sigma.n_1.\cdots.n_k$ شکل والا اگیتا۔ اسیں نشان $\sigma.*$ ورتاں گے،
جیہڑا خود $\sigma$ تے اوہدی ہر توسیع دسدا اے۔ دوجے لفظاں وچ
$\sigma.*$ اوہ سارے، تے صرف اوہی، اگیتے نیں جیہڑے~$\sigma$ توں رسائی وچ نیں۔
\endgroup
% /OLP-0512
% OLP-0513 content/intuitionistic-logic/tableaux/rules.tex
\begingroup
\def\olfilename{rules}
\expandafter\def\csname ol@sectioncs\endcsname{section}

\olfileid{int}{tab}{rul}

\olsection{وجدانیت پسند منطق دے قاعدے}

$\land$ تے $\lor$ دے قاعدے اوہو نیں جیہڑے عام قضیاتی نشان والے
!!{tableau}s وچ نیں۔ بس اگیتے شامل نیں۔ ہر صورت وچ نشان والے !!{formula}
$\sFmla{S}{!A}[\sigma]$ اُتے قاعدہ لاؤن نال نویں !!{formula}s ملدے نیں
جیہناں دا اگیتا وی~$\sigma$ اے۔ ایہ گل سوکھی سمجھ آندی اے۔ مثلاً جے
$!A \land !B$،~$\sigma$ ناں والے جہان اُتے سچ اے، تاں $!A$ تے $!B$
~$\sigma$ اُتے سچ نیں (تے کسے ہور جہان اُتے نہیں)۔ $\land$ تے $\lor$
دے قاعدے \olref{tab:prop-rules} وچ اکٹھے کیتے نیں۔

\begin{table}
  \[\def\arraystretch{3}\begin{array}{|c|c|}
    \hline
    \AxiomC{\sFmla{\True}{!A \land !B}[\sigma]}
    \RightLabel{\TRule{\True}{\land}}
    \UnaryInfC{\sFmla{\True}{!A}[\sigma]}
    \noLine
    \UnaryInfC{\sFmla{\True}{!B}[\sigma]}
    \DisplayProof
    &
    \AxiomC{\sFmla{\False}{!A \land !B}[\sigma]}
    \RightLabel{\TRule{\False}{\land}}
    \UnaryInfC{$\sFmla{\False}{!A}[\sigma] \quad \mid \quad
      \sFmla{\False}{!B}[\sigma]$}
    \DisplayProof
    \\[2ex]
    \hline
    \AxiomC{\sFmla{\True}{!A \lor !B}[\sigma]}
    \RightLabel{\TRule{\True}{\lor}}
    \UnaryInfC{$\sFmla{\True}{!A}[\sigma] \quad \mid \quad
      \sFmla{\True}{!B}[\sigma]$}
    \DisplayProof
    &
    \AxiomC{\sFmla{\False}{!A \lor !B}[\sigma]}
    \RightLabel{\TRule{\False}{\lor}}
    \UnaryInfC{\sFmla{\False}{!A}[\sigma]}
    \noLine
    \UnaryInfC{\sFmla{\False}{!B}[\sigma]}
    \DisplayProof
    \\[2ex]
    \hline
  \end{array}\]
  \caption{اگیتے والے !!{tableau} دے قاعدے $\land$ تے $\lor$ لئی}
  \ollabel{tab:prop-rules}
\end{table}

ٹہنی بند ہون دی شرط عام !!{tableau}s ورگی اے، پر ایتھے !!{formula}s دے
نال اگیتیاں دا میل وی چاہیدا اے۔ اصل وچ شرط کجھ کھُلّی ہو سکدی اے۔
وجدانیت پسند نمونیاں وچ !!{formula}s اک واری سچ ہو جان، تاں سچ رہندے نیں۔
ایس لئی ایہ ممکن نہیں کہ $!A$،~$\sigma$ اُتے سچ پر کسے رسائی والے اگیتے
~$\sigma.{*}$ اُتے جھوٹ ہووے۔ سو ٹہنی بند اے جے اوہدے وچ دونیں ہون:
\[
\sFmla{\True}{!A}[\sigma] \quad\text{ تے }\quad \sFmla{\False}{!A}[\sigma.{*}]
\]
کسے اگیتے $\sigma$ تے !!{formula}~$!A$ لئی۔ دھیان دیو، جے نشان اُلٹ ہون،
یعنی ٹہنی وچ
\[
\sFmla{\False}{!A}[\sigma] \quad\text{ تے }\quad \sFmla{\True}{!A}[\sigma.{*}]
\]
ہون، تاں ٹہنی صرف اوہدوں بند اے جدوں $*$ خالی لڑی اے۔

ٹہنی اوہدوں وی بند اے جدوں اوہدے وچ~$\sFmla{\True}{\bot}[\sigma]$ ہووے۔

مفروضے قائم کرن دے قاعدے وی عام !!{tableau}s والے ای نیں، فرق ایہ اے کہ
مفروضیاں لئی اسیں ہمیشہ اگیتا~$1$ ورتدے آں۔ (کیہڑا اگیتا ورتنا اے، ایس نال
فرق نہیں پیندا، بس سارے مفروضیاں دا اگیتا اکّو ہووے۔) سو مثلاً اسیں آکھدے آں کہ
\[
!B_1, \dots, !B_n \Proves !A
\]
اُدوں تے صرف اُدوں اے جدوں ہِیٹھلے مفروضیاں دا اک بند ٹیبلو ہووے:
\[
\sFmla{\True}{!B_1}[1], \dots, \sFmla{\True}{!B_n}[1],
\sFmla{\False}{!A}[1].
\]

شرطیے~$\lif$ لئی قاعدے کلاسیکی تے امکانی صورتاں توں وکھرے نیں۔
قاعدہ $\TRule{\lif}{\True}$، $\sFmla{\True}{!A \lif !B}[\sigma]$ والی ٹہنی
نوں دو وکھریاں ٹہنیاں اُتے $\sFmla{\True}{!A}[\sigma.{*}]$ تے
$\sFmla{\False}{!B}[\sigma.{*}]$ پا کے ودھاندا اے۔ ایہ قاعدہ صرف اوہ
اگیتا~$\sigma.{*}$ ورت کے لا سکدے آں جیہڑا قاعدہ لاؤن والی ٹہنی اُتے
\emph{پہلاں ای} آ چکیا ہووے۔ ایہو جیہے اگیتے نوں (ٹہنی اُتے) ’ورتیا ہویا‘
آکھیے۔ (چونکہ $\sigma.{*}$ وچ خود $\sigma$ شامل اے، ایس لئی ایہ قاعدہ
اگیتے والے نشان لائے فارمولے $\sFmla{\True}{!A}[\sigma]$ تے
$\sFmla{\False}{!B}[\sigma]$ وکھریاں ٹہنیاں اُتے پا کے ہمیشہ لایا جا سکدا اے۔)

قاعدہ $\TRule{\lif}{\False}$، $\sFmla{\False}{!A \lif !B}[\sigma]$ والی
ٹہنی نوں $\sFmla{\True}{!A}[\sigma.n]$ تے $\sFmla{\False}{!B}[\sigma.n]$
دونیں اوسے ٹہنی اُتے پا کے ودھاندا اے، جتھے $\sigma.n$ ٹہنی لئی نواں اگیتا اے۔

$\lnot$ دے قاعدے ایس طرح تعریف ہوندے نیں (اوہ تعریف ورت کے جیہڑی
$\lnot !A$ نوں $!A \lif \lfalse$ دے طور اُتے دیندی اے)۔

قاعدے \olref{tab:rules-lif-lnot} وچ دتے نیں۔

\begin{table}
  \[\def\arraystretch{3}\begin{array}{|c|c|}
    \hline
    \AxiomC{\sFmla{\True}{\lnot !A}[\sigma]}
    \RightLabel{\TRule{\True}{\lnot}}
    \UnaryInfC{$\sFmla{\False}{!A}[\sigma.{*}]$}
    \DisplayProof
    &
    \AxiomC{\sFmla{\False}{\lnot !A}[\sigma]}
    \RightLabel{\TRule{\False}{\lnot}}
    \UnaryInfC{\sFmla{\True}{!A}[\sigma.n]}
    \DisplayProof
    \\[1ex]
    \text{$\sigma.{*}$ ورتیا ہویا اے} & \text{$\sigma.n$ نواں اے}\\
    \hline
    \AxiomC{\sFmla{\True}{!A \lif !B}[\sigma]}
    \RightLabel{\TRule{\True}{\lif}}
    \UnaryInfC{$\sFmla{\False}{!A}[\sigma.{*}] \quad \mid \quad
      \sFmla{\True}{!B}[\sigma.{*}]$}
    \DisplayProof
    &
    \AxiomC{\sFmla{\False}{!A \lif !B}[\sigma]}
    \RightLabel{\TRule{\False}{\lif}}
    \UnaryInfC{\sFmla{\True}{!A}[\sigma.n]}
    \noLine
    \UnaryInfC{\sFmla{\False}{!B}[\sigma.n]}
    \DisplayProof
    \\[1ex]
    \text{$\sigma.{*}$ ورتیا ہویا اے} & \text{$\sigma.n$ نواں اے}\\
    \hline
  \end{array}\]
  \caption{اگیتے والے !!{tableau} دے قاعدے $\lnot$ تے $\lif$ لئی}
  \ollabel{tab:rules-lif-lnot}
\end{table}
\endgroup
% /OLP-0513
% OLP-0514 content/intuitionistic-logic/tableaux/proofs.tex
\begingroup
\def\olfilename{proofs}
\expandafter\def\csname ol@sectioncs\endcsname{section}

\olfileid{int}{tab}{prf}

\olsection{وجدانیت پسند منطق لئی \usetoken{P}{tableau}}

\begin{ex}
  اسیں اک بند ٹیبلو دیندے آں جیہڑا وکھاندا اے کہ $(!A \land !B) \lif !C \Proves
  !A \lif (!B \lif !C)$۔
  \begin{oltableau}
    [\pFmla{\True}{(\formula{A} \land \formula{B}) \lif \formula{C}}{1},
      just =\TAss
      [\pFmla{\False}{\formula{A} \lif (\formula{B} \lif \formula{C})}{1},
        just =\TAss
        [\pFmla{\True}{\formula{A}}{1.1},
          just = {\TRule{\False}{\lif}[2]}
          [\pFmla{\False}{\formula{B} \lif \formula{C}}{1.1},
            just = {\TRule{\False}{\lif}[2]}
            [\pFmla{\True}{\formula{B}}{1.1.1},
              just = {\TRule{\False}{\lif}[4]}
              [\pFmla{\False}{\formula{C}}{1.1.1},
                just = {\TRule{\False}{\lif}[4]}
                [\pFmla{\False}{\formula{A} \land \formula{B}}{1.1.1},
                  just = {\TRule{\True}{\lif}[1]}
                  [\pFmla{\False}{\formula{A}}{1.1.1},
                    just= {\TRule{\False}{\land}[4]}, close]
                  [\pFmla{\False}{\formula{B}}{1.1.1},
                    just= {\TRule{\False}{\land}[4]}, close]]
                [\pFmla{\True}{\formula{C}}{1.1.1},
                  just = {\TRule{\True}{\lif}[1]}, close]]
            ]
          ]
        ]
      ]
    ]
  \end{oltableau}
\end{ex}

\begin{prob}
  ہِیٹھلیاں گلاں وکھاؤن لئی بند وجدانیت پسند !!{tableau}s لبھو:
  \begin{enumerate}
    \item $\Proves !A \lif (!B \lif !A)$
    \item $\Proves \lnot(!A \land \lnot !A)$
    \item $!A \lif (!B \lif !C) \Proves (!A \land !B) \lif !C$
    \item $\lnot !A \lor \lnot !B \Proves \lnot(!A \land !B)$
  \end{enumerate}
\end{prob}
\endgroup
% /OLP-0514
% OLP-0515 content/intuitionistic-logic/tableaux/soundness.tex
\begingroup
\def\olfilename{soundness}
\expandafter\def\csname ol@sectioncs\endcsname{section}

\olfileid{int}{tab}{sou}

\olsection{وجدانیت پسند \usetoken{P}{tableau} دی صحت مندی}

\begin{explain}
  وجدانیت پسند !!{tableau}s دی صحت مندی وکھاؤن لئی سانوں ایہ ثابت کرنا اے
  کہ جے
  \[
  \sFmla{\True}{!B_1}[1], \dots, \sFmla{\True}{!B_n}[1], \sFmla{\False}{!A}[1]
  \]
  دا بند !!{tableau} اے، تاں $!B_1, \dots, !B_n \Entails !A$۔ عکسِ نقیض
  ثابت کرنا سوکھا اے۔ جے کسے $\mModel{M}$ تے جہان~$w$ لئی
  $\mSat{M}{!B_i}[w]$ ہر $i=1$، \dots،~$n$ لئی اے، پر $\mSat{M}{!A}[w]$،
  تاں کوئی !!{tableau} بند نہیں ہو سکدا۔ ایہو جیہا ردّی نمونہ وکھاندا اے
  کہ !!{tableau} دے شروع دے مفروضے اکٹھے پورے ہو سکدے نیں۔ ثبوت دا طریقہ
  ایہ اے کہ جدوں وی !!a{tableau} دی ٹہنی اُتے سارے اگیتے والے
  !!{formula}s اکٹھے پورے ہو سکن، تاں کوئی قاعدہ لاؤن نال گھٹ توں گھٹ
  اک ایہو جیہی ودھی ہوئی ٹہنی ملدی اے جیہڑی ہجے وی پوری ہو سکدی اے۔
  چونکہ بند ٹہنیاں پوریاں نہیں ہو سکدیاں، ایس لئی اگیتے والے !!{formula}s دے
  پورے ہو سکّن والے سیٹ لئی ہر !!{tableau} وچ گھٹ توں گھٹ اک کھُلّی ٹہنی ہووے گی۔

  ایس طریقے نوں امکانی صورت وچ ورتن لئی سانوں ’پورا ہون جوگا‘ دی تعریف
  تعلقی نمونیاں تے اگیتیاں تک ودھاؤنی پئے گی۔ ایہ ہو جائے تاں ثبوت سدھا اے۔
\end{explain}

\begin{defn}
  منّ لو $P$ اگیتیاں دا کوئی سیٹ اے، یعنی $P \subseteq (\PosInt)^*
  \setminus \{\emptyseq\}$، تے $\mModel{M}$ اک نمونہ اے۔
  فنکشن~$f\colon P \to W$،~$P$ دی~$\mModel{M}$ وچ \emph{تعبیر} اے جے
  جدوں وی $\sigma$ تے $\sigma.n$ دونیں~$P$ وچ ہون، اوہدوں $Rf(\sigma)f(\sigma.n)$۔

  اگیتیاں~$P$ دی تعبیر دے لحاظ نال اسیں ایویں تعریف کر سکدے آں:
  \begin{enumerate}
  \item $\mModel{M}$ وچ $\sFmla{\True}{!A}[\sigma]$ اُدوں تے صرف اُدوں
    پورا اتردا اے جدوں $\mSat{M}{!A}[f(\sigma)]$۔
  \item $\mModel{M}$ وچ $\sFmla{\False}{!A}[\sigma]$ اُدوں تے صرف اُدوں
    پورا اتردا اے جدوں $\mSat/{M}{!A}[f(\sigma)]$۔
  \end{enumerate}
\end{defn}

دھیان دیو کہ $R$ انعکاسی تے متعدی اے، تے $\sigma.{*}$ توں مراد
$\sigma$، $\sigma.n_1$، $\sigma.n_1.n_2$، \dots نیں۔ ایس لئی ساڈے کول
$Rf(\sigma)f(\sigma.{*})$ وی اے۔

\begin{defn}
  منّ لو $\Gamma$ اگیتے والے !!{formula}s دا سیٹ اے، تے $P(\Gamma)$
  اوہدے وچ آؤن والے اگیتیاں دا سیٹ اے۔ جے $f$،~$P(\Gamma)$ دی
  $\mModel{M}$ وچ تعبیر اے، تاں اسیں آکھدے آں کہ $\mModel{M}$ وچ
  $\Gamma$،~$f$ دے لحاظ نال پورا اتردا اے، یعنی $\mSat{M}{\Gamma}[f]$،
  جے $\mModel{M}$ وچ~$\Gamma$ دا ہر اگیتے والا !!{formula}،~$f$ دے لحاظ نال
  پورا اتردا اے۔ $\Gamma$ \emph{پورا ہون جوگا} اُدوں تے صرف اُدوں اے جدوں
  اک نمونہ~$\mModel{M}$ تے اک تعبیر~$f$ ہووے، جیہڑی $P(\Gamma)$ دی اے تے جیہدے لئی
  $\mSat{M}{\Gamma}[f]$۔
\end{defn}

\begin{prop}
  جے $\Gamma$ وچ دونیں $\sFmla{\True}{!A}[\sigma]$ تے
  $\sFmla{\False}{!A}[\sigma.{*}]$ ہون، کسے !!{formula}~$!A$ تے اگیتے~$\sigma$ لئی،
  یا اوہدے وچ $\sFmla{\True}{\lfalse}[\sigma]$ ہووے، تاں $\Gamma$ پورا ہون جوگا نہیں۔
\end{prop}

\begin{proof}
  چونکہ ہمیشہ $\mSat/{M}{\lfalse}[f(\sigma)]$، ایس لئی اوہ $\Gamma$
  جیہدے وچ $\sFmla{\True}{\lfalse}$ ہووے، پورا ہون جوگا نہیں۔

  کوئی نمونہ~$\mModel{M}$ تے تعبیر~$f$ وی نہیں ہو سکدی، جیہڑی $P(\Gamma)$ دی اے تے
  جتھے دونیں شرطاں پوریاں ہون۔ جے $\mSat{M}{!A}[f(\sigma)]$، تاں
  \olref[sem][rel]{prop:true-monotonic} مطابق، چونکہ $Rf(\sigma)(\sigma.{*})$،
  ایس لئی $\mSat{M}{!A}[f(\sigma)]$۔ سو دونیں
  $\mSat{M}{!A}[f(\sigma)]$ تے $\mSat/{M}{!A}[f(\sigma.{*})]$ نہیں ہو سکدے۔
\end{proof}

\begin{thm}[صحت مندی]
  \ollabel{thm:tableau-soundness}
  جے $\Gamma$ دا کوئی بند !!{tableau} اے، تاں $\Gamma$ پورا ہون جوگا نہیں۔
\end{thm}

\begin{proof}
اسیں !!a{tableau} دی ٹہنی نوں پورا ہون جوگا اُدوں تے صرف اُدوں آکھدے آں
جدوں اوہدے اُتے !!{signed formula}s دا سیٹ پورا ہون جوگا اے۔ تے
!!a{tableau} نوں پورا ہون جوگا اوہدوں آکھیے جے اوہدے وچ گھٹ توں گھٹ
اک پورا ہون جوگی ٹہنی اے۔

اسیں ایہ وکھاندے آں کہ کسے پورا ہون جوگے !!{tableau} نوں استنباط دے
اک قاعدے نال ودھاؤن توں ہمیشہ اک پورا ہون جوگا !!{tableau} ملدا اے۔
ایس نال قضیہ ثابت ہووے گا۔ کوئی وی بند !!{tableau} اوس !!{tableau}
اُتے استنباط دے قاعدے لاؤن نال ملدا اے جیہڑا صرف~$\Gamma$ دے مفروضیاں توں
بندا اے۔ سو جے $\Gamma$ پورا ہون جوگا ہووے، تاں اوہدے لئی ہر !!{tableau}
پورا ہون جوگا ہووے گا۔ پر بند !!{tableau} صاف طور اُتے پورا ہون جوگا نہیں،
کیوں جے اوہدیاں ساریاں ٹہنیاں بند نیں تے بند ٹہنیاں پوریاں نہیں ہو سکدیاں۔

فرض کرو ساڈے کول اک پورا ہون جوگا !!{tableau} اے، یعنی !!a{tableau}
جیہدے وچ گھٹ توں گھٹ اک پورا ہون جوگی ٹہنی اے۔ استنباط دا قاعدہ لاؤن نال
یا کسے ٹہنی اُتے !!{signed formula}s شامل ہوندے نیں یا اوہ ٹہنی دو حصیاں
وچ ونڈی جاندی اے۔ جے !!{tableau} دی کسے پورا ہون جوگی ٹہنی نوں ایس قاعدے دے
اطلاق نال ودھایا نہیں جاندا، تاں اوہ ودھے ہوئے !!{tableau} وچ وی پورا ہون
جوگی ٹہنی رہندی اے۔ سو ودھیا ہویا ٹیبلو پورا ہون جوگا اے۔ ایس لئی سانوں
صرف اوہ صورت ویکھنی اے جدوں قاعدہ کسے پورا ہون جوگی ٹہنی اُتے لگدا اے۔

منّ لو $\Gamma$ اوس ٹہنی دے !!{signed formula}s دا سیٹ اے تے
$\sFmla{S}{!A}[\sigma] \in \Gamma$ اوہ !!{signed formula} اے جیہدے اُتے
قاعدہ لگدا اے۔ جے قاعدہ ٹہنی نوں دو حصیاں وچ نہیں ونڈدا، تاں اسیں وکھانا اے
کہ ودھی ہوئی ٹہنی، یعنی $\Gamma$ تے قاعدے دے نتیجے رلا کے، ہجے وی پورا
ہون جوگی اے۔ جے قاعدہ ٹہنی نوں ونڈدا اے، تاں اسیں وکھانا اے کہ بنیاں دوناں
ٹہنیاں وچوں گھٹ توں گھٹ اک پورا ہون جوگی اے۔
\begin{enumerate}
\item $\TRule{\True}{\lnot}$ نوں $\sFmla{\True}{\lnot !B}[\sigma] \in \Gamma$
  اُتے لاؤن نال ٹہنی ودھدی اے۔ ودھی ہوئی ٹہنی اُتے !!{signed formula}s
  $\Gamma \cup \{\sFmla{\False}{!B}[\sigma.{*}]\}$ نیں۔
  فرض کرو $\mSat{M}{\Gamma}[f]$۔ خاص کر کے $\mSat{M}{\lnot !B}[f(\sigma)]$۔
  سو $\mSat/{M}{!B}[w]$ ہر اوہ $w$ لئی اے جیہڑا $Rf(\sigma)w$ پورا کرے،
  تے اوہناں وچ $f(\sigma.{*})$ وی اے۔ سو $\mModel{M}$ وچ
  $\sFmla{\False}{!B}[\sigma.{*}]$،~$f$ دے لحاظ نال پورا اتردا اے۔
\item $\TRule{\False}{\lnot}$ نوں $\sFmla{\False}{\lnot !B}[\sigma] \in \Gamma$
  اُتے لاؤن نال ٹہنی ودھدی اے۔ مشق۔
\item $\TRule{\True}{\land}$ نوں $\sFmla{\True}{!B \land !C}[\sigma] \in \Gamma$
  اُتے لاؤن نال ٹہنی ودھدی اے تے ٹہنی اُتے دو نویں !!{signed formula}s ملدے نیں:
  $\sFmla{\True}{!B}[\sigma]$ تے $\sFmla{\True}{!C}[\sigma]$۔ فرض کرو
  $\mSat{M}{\Gamma}[f]$، خاص کر کے $\mSat{M}{!B \land !C}[f(\sigma)]$۔
  فیر $\mSat{M}{!B}[f(\sigma)]$ تے $\mSat{M}{!C}[f(\sigma)]$۔
  یعنی $\mModel{M}$ وچ $\sFmla{\True}{!B}[\sigma]$ تے
  $\sFmla{\True}{!C}[\sigma]$ دونیں~$f$ دے لحاظ نال پورے اتردے نیں۔
\item $\TRule{\False}{\lor}$ نوں $\sFmla{\False}{!B \lor !C} \in \Gamma$
  اُتے لاؤن نال ٹہنی ودھدی اے۔ مشق۔
\item $\TRule{\False}{\lif}$ نوں $\sFmla{\False}{!B \lif !C}[\sigma] \in \Gamma$
  اُتے لاؤن نال ٹہنی ودھدی اے۔ ایس نال ٹہنی اُتے دو نویں !!{signed formula}s ملدے نیں:
  $\sFmla{\True}{!B}[\sigma.n]$ تے $\sFmla{\False}{!C}[\sigma.n]$، جتھے
  $\sigma.n$ ٹہنی اُتے نواں اگیتا اے، یعنی $\sigma.n \notin P(\Gamma)$۔

  چونکہ $\Gamma$ پورا ہون جوگا اے، اک $\mModel{M}$ تے تعبیر~$f$ اے، جیہڑی
  $P(\Gamma)$ دی اے تے جیہدے لئی $\mSat{M}{\Gamma}[f]$، خاص کر کے
  $\mSat/{M}{!B \lif !C}[f(\sigma)]$۔ اسیں وکھانا اے کہ
  $\Gamma \cup \{\sFmla{\False}{!B \lif !C}[\sigma.n]\}$ پورا ہون جوگا اے۔
  ایس لئی اسیں $P(\Gamma) \cup \{\sigma.n\}$ دی تعبیر ایویں تعریف کردے آں:

  چونکہ $\mSat/{M}{!B \lif !C}[f(\sigma)]$، ایس لئی اک $w \in W$ اے
  جیہڑا $Rf(\sigma)w$ پورا کردا اے تے جیہدے لئی $\mSat{M}{!B}[w]$ تے
  $\mSat/{M}{!C}[w]$۔ منّ لو $f'$، $f$ ورگا اے، بس $f'(\sigma.n) = w$۔
  چونکہ $f'(\sigma) = f(\sigma)$ تے $Rf(\sigma)w$، ایس لئی
  $Rf'(\sigma)f'(\sigma.n)$۔ سو $f'$، $P(\Gamma) \cup \{\sigma.n\}$ دی
  اک تعبیر اے۔ ظاہر اے کہ $\mSat{M}{!B}[f'(\sigma.n)]$ تے
  $\mSat/{M}{!C}[f'(\sigma.n)]$۔ چونکہ $f(\sigma') = f'(\sigma')$
  سارے اگیتیاں $\sigma' \in P(\Gamma)$ لئی، ایس لئی $\mSat{M}{\Gamma}[f']$۔
  سو $\mModel{M}, f'$ وچ $\Gamma \cup \{\sFmla{\False}{!B \lif !C}[\sigma.n]\}$
  پورا اتردا اے۔
\end{enumerate}
ہُن دو مقدمیاں والے ممکن استنباط ویکھیے۔
\begin{enumerate}
\item $\TRule{\False}{\land}$ نوں $\sFmla{\False}{!B \land !C}[\sigma] \in \Gamma$
  اُتے لاؤن نال ٹہنی ودھدی اے تے دو ٹہنیاں ملدیاں نیں۔ کھبے والی
  $\sFmla{\False}{!B}[\sigma]$ توں اگے جاندی اے تے سجے والی
  $\sFmla{\False}{!C}[\sigma]$ توں۔ فرض کرو $\mSat{M}{\Gamma}[f]$،
  خاص کر کے $\mSat/{M}{!B \land !C}[f(\sigma)]$۔ فیر
  $\mSat/{M}{!B}[f(\sigma)]$ یا $\mSat/{M}{!C}[f(\sigma)]$۔
  پہلی صورت وچ $\mModel{M}, f$ وچ $\sFmla{\False}{!B}[\sigma]$ پورا
  اتردا اے، یعنی کھبی ٹہنی پورا ہون جوگی اے۔ دوجی صورت وچ
  $\mModel{M}, f$ وچ $\sFmla{\False}{!C}[\sigma]$ پورا اتردا اے،
  یعنی سجی ٹہنی پورا ہون جوگی اے۔
\item $\TRule{\True}{\lor}$ نوں $\sFmla{\True}{!B \lor !C}[\sigma] \in \Gamma$
  اُتے لاؤن نال ٹہنی ودھدی اے۔ مشق۔
\item $\TRule{\True}{\lif}$ نوں $\sFmla{\True}{!B \lif !C}[\sigma] \in \Gamma$
  اُتے لاؤن نال ٹہنی ودھدی اے۔ مشق۔
\end{enumerate}
\end{proof}

\begin{prob}
\olref[int][tab][sou]{thm:tableau-soundness} دا ثبوت مکمل کرو۔
\end{prob}

\begin{cor}
\ollabel{cor:entailment-soundness}
جے $\Gamma \Proves !A$، تاں $\Gamma \Entails !A$۔
\end{cor}

\begin{proof}
  جے $\Gamma \Proves !A$، تاں کجھ $!B_1$، \dots، $!B_n \in \Gamma$ لئی،
  $\Delta = \{\sFmla{\False}{!A}[1], \sFmla{\True}{!B_1}[1],
  \dots, \sFmla{\True}{!B_n}[1]\}$ دا بند !!{tableau} اے۔ اسیں وکھانا اے
  کہ $\Gamma \Entails !A$۔ فرض کرو ایہ نہیں۔ تاں کسے $\mModel{M}$ تے $w$ لئی،
  $\mSat{M}{!B_i}[w]$ ہر $i=1$، \dots،~$n$ لئی، پر $\mSat/{M}{!A}[w]$۔
  منّ لو $f(1) = w$۔ فیر $f$،~$P(\Delta)$ دی~$\mModel{M}$ وچ تعبیر اے،
  تے $\mModel{M}$ وچ~$\Delta$،~$f$ دے لحاظ نال پورا اتردا اے۔
  پر \olref{thm:tableau-soundness} مطابق، $\Delta$ پورا ہون جوگا نہیں،
  کیوں جے اوہدا بند !!{tableau} اے۔ ایہ تضاد اے۔ سو آخر $\Gamma \Proves !A$ ہونا اے۔
\end{proof}

\begin{cor}
\ollabel{cor:weak-soundness}
جے $\Proves !A$، تاں $!A$ سارے نمونیاں وچ سچ اے۔
\end{cor}
\endgroup
% /OLP-0515
%\olimport{completeness}
%\olimport{countermodels}

\OLEndChapterHook
\endgroup
% /OLP-0511

\OLEndPartHook
\endgroup
% /OLP-0491

% OLP-0685 content/proof-theory/proof-theory.tex
\begingroup
\def\olfilename{proof-theory}
\expandafter\def\csname ol@sectioncs\endcsname{section}

\olpart{pt}{ثبوت دا نظریہ}

\begin{editorial}
ثبوت دے نظریے دا مواد ہجے ادھورا تے تجرباتی اے، تے ہلے مُکھ PDF وچ شامل نہیں۔
\end{editorial}

% OLP-0712 content/proof-theory/sequent-calculus/sequent-calculus.tex
\begingroup
\def\olfilename{sequent-calculus}
\expandafter\def\csname ol@sectioncs\endcsname{section}

\olchapter{pt}{seq}{سیکوئنٹ حساب}

% OLP-0700 content/proof-theory/sequent-calculus/introduction.tex
\begingroup
\def\olfilename{introduction}
\expandafter\def\csname ol@sectioncs\endcsname{section}

\olfileid{pt}{seq}{int}

\olsection{جان پچھان}

سیکوئنٹ حساب !!{proof} نظام نیں جیہڑے گیرہارٹ گینٹسن نے 1930 دیاں دہائیاں وچ متعارف کرائے۔ ایہناں دا مقصد !!{proving} لئی اک عملی طریقہ دینا گھٹ سی۔ گینٹسن نے ایہ ایس لئی ورتے کہ !!{proof}s دی ممکن بناوٹ تے اوہناں دی خاصیت بارے کجھ نتیجے ثابت کر سکے۔ عین ایہو جیہے سوال ثبوت دے نظریے دا موضوع نیں۔

جیویں ناں توں پتا لگدا اے، سیکوئنٹ حساب \emph{سیکوئنٹاں} اُتے کم کردے نیں، !!{formula}s اُتے نہیں۔ سیکوئنٹ !!{formula}s دیاں لڑیاں یا کثیر سیٹاں دا جوڑا اے۔ گینٹسن دے اصل نظام (\Log{LK} تے \Log{LJ}) لڑیاں ورتدے سن۔ ہُن ثبوت دے نظریے دے ماہر کثیر سیٹ ورتنا پسند کردے نیں۔ کثیر سیٹ، سیٹ ورگا اے، فرق ایہ اے کہ کثیر سیٹ کسے !!{element} نوں اک توں ودھ واری رکھ سکدا اے۔ پر سیٹاں وانگ !!{element}s دی ترتیب نال فرق نہیں پیندا۔ سو عبارتیں $\{!A, !A, !B\}$ تے $\{!A, !B, !A\}$ اوہی کثیر سیٹ دَسّدیاں نیں، پر اوہ کثیر سیٹ $\{!A,
!B\}$ تے $\{!A, !A, !B, !B\}$ توں وکھرا اے، کیونکہ پچھلے دوہاں وچ $!A$ تے $!B$ دیاں گنتیاں پہلے نالوں وکھریاں نیں۔

رواج مطابق ثبوت دے نظریے وچ !!{formula}s دے کثیر سیٹ وڈّے یونانی حروف نال لکھے جاندے نیں۔ سو جدوں اسیں $\Gamma$، $\Delta$، $\Pi$، $\Lambda$ یا $\Theta$ لکھدے آں، تاں اوہناں توں ساڈا مطلب اوس منطق دے !!{formula}s دے کثیر سیٹ نیں جیہڑی اسیں ورت رہے آں۔ رواج مطابق ایہ وی اے کہ $\tuple{\Gamma, \Delta}$ لکھن دی تھاں ثبوت دے نظریے دے ماہر دوہاں جزواں نوں سیکوئنٹ دی نشانی نال وکھ کردے نیں۔ اسیں~$\Sequent$ ورتدے آں۔

\begin{defn}[سیکوئنٹ]
اک \emph{سیکوئنٹ} ایس شکل دی عبارت اے:
\[
\Gamma \Sequent \Delta
\]
جتھے $\Gamma$ تے $\Delta$، !!{formula}s دے محدود (شاید خالی) کثیر سیٹ نیں۔ $\Gamma$ نوں \emph{کھبّا پاسا} تے $\Delta$ نوں \emph{سجّا پاسا} آکھدے آں۔
\end{defn}

منّ لو $!G$، اوہناں !!{formula}s دا عطف اے جیہڑے~$\Gamma$ وچ نیں (خالی عطف نوں~$\ltrue$ منّ کے)، تے $!D$، اوہناں !!{formula}s دی فصل اے جیہڑے~$\Delta$ وچ نیں (خالی فصل نوں~$\lfalse$ منّ کے)۔ تاں سیکوئنٹ $\Gamma \Sequent \Delta$ کسے ساخت~$\Struct{M}$ وچ اوہدوں تے صرف اوہدوں سچا اے جدوں $!G \lif !D$ سچا اے۔ کلاسیکی منطق وچ ایہ اوہو گل اے کہ یا اک !!{formula}s وچوں جیہڑا~$\Gamma$ وچ اے جھوٹھا اے، یا اک !!{formula}s وچوں جیہڑا~$\Delta$ وچ اے سچا اے۔

جے $\Gamma$، !!{formula}s دا کثیر سیٹ اے، تاں اسیں $\Gamma, !A$ (یا $!A, \Gamma$) لکھدے آں، $!A$ نوں~$\Gamma$ وچ شامل کرن دے حاصل لئی۔ جے $ \Delta$ وی !!{formula}s دا کثیر سیٹ اے، تاں $\Gamma, \Delta$ دوہاں لڑیاں دا کثیر سیٹ اتحاد اے۔ جے $\Gamma$ وچ $!A$ دی \emph{کثرت}~n اے (یعنی اوہدے وچ $n$ نقلاں~$!A$ دیاں نیں) تے $ \Delta$ وچ کثرت~$m$ اے، تاں $\Gamma, \Delta$ وچ $!A$ دی کثرت~$n+m$ اے۔ جے سیکوئنٹ دے اک پاسے دا کثیر سیٹ خالی اے، تاں اسیں عام طور اُتے اوس پاسے خالی تھاں چھڈّ دیندے آں۔ مثال لئی $\Sequent !A$ اوہ سیکوئنٹ اے جیہدا کھبّا پاسا خالی اے تے سجّے پاسے $!A$ دی کثرت~$1$ اے۔

سیکوئنٹ حساب وچ !!^a{proof} سیکوئنٹاں دا درخت اے۔ اوہدے سبھ توں اُتلے (یعنی شروع دے) سیکوئنٹ متعلقہ نظام دے مسلّمے نیں، تے جے کوئی سیکوئنٹ اک یا دو ہور سیکوئنٹاں دے تھلّے اے، تاں اوہ نظام دے کسے استنباطی قاعدے نال ٹھیک طرح نکلنا چاہیدا اے۔ اسیں پہلاں کلاسیکی منطق لئی دو وکھرے سیکوئنٹ نظام، $\Log{G1c}$ تے~$\Log{G3c}$، ویکھاں گے۔ اوہناں دے مسلّمے تے قاعدے \cref{pt:seq:int:tab:G1c,pt:seq:int:tab:G3c} وچ دِتّے نیں۔ \oliflabeldef{fol:seq::chap}{ (گینٹسن دا اصل نظام \Log{LK}، \olref[fol][seq][]{chap} وچ بیان کیتا گیا اے۔)}{}

% OLP-0704 content/proof-theory/sequent-calculus/rules-G1c.tex
\begingroup

\begin{table}
\begin{tabular}{@{}c@{}c@{}}
\multicolumn{2}{@{}c@{}}{Axioms}\\[2ex]
$!A \Sequent !A$ & $\lfalse \Sequent \quad$
\\[2ex]
\multicolumn{2}{@{}c@{}}{Structural Rules}\\[2ex]
\Axiom$ \Gamma \fCenter \Delta $
\RightLabel{\LeftR{\Weakening}}
\UnaryInf$ !A, \Gamma \fCenter \Delta$
\DisplayProof
&
\Axiom$ \Gamma \fCenter \Delta$
\RightLabel{\RightR{\Weakening}}
\UnaryInf$ \Gamma \fCenter \Delta, !A$
\DisplayProof
\\[4ex]
\Axiom$ !A, !A, \Gamma \fCenter \Delta $
\RightLabel{\LeftR{\Contraction}}
\UnaryInf$ !A, \Gamma \fCenter \Delta$
\DisplayProof
&
\Axiom$ \Gamma \fCenter \Delta, !A, !A$
\RightLabel{\RightR{\Contraction}}
\UnaryInf$ \Gamma \fCenter \Delta, !A$
\DisplayProof
\\[4ex]
\multicolumn{2}{@{}c@{}}{Logical Rules}\\[2ex]
\Axiom$ \Gamma \fCenter \Delta, !A $
\RightLabel{\LeftR{\lnot}}
\UnaryInf$ \lnot !A, \Gamma \fCenter \Delta$
\DisplayProof
&
\Axiom$!A, \Gamma \fCenter \Delta$
\RightLabel{\RightR{\lnot}}
\UnaryInf$ \Gamma \fCenter \Delta, \lnot !A $
\DisplayProof
\\[4ex]\begin{tabular}[t]{@{}l@{}}
\Axiom$ !A, \Gamma \fCenter \Delta$
\RightLabel{\LeftR{\land}}
\UnaryInf$ !A \land !B, \Gamma \fCenter \Delta$
\DisplayProof
\\[3ex]
\Axiom$!B, \Gamma \fCenter \Delta$
\RightLabel{\LeftR{\land}}
\UnaryInf$!A \land !B, \Gamma \fCenter \Delta$
\DisplayProof\\[3ex]
\end{tabular}
&
\Axiom$\Gamma \fCenter \Delta, !A$
\Axiom$ \Gamma \fCenter \Delta, !B$
\RightLabel{\RightR{\land}}
\BinaryInf$ \Gamma \fCenter \Delta, !A \land !B $
\DisplayProof
\\[4ex]\Axiom$!A, \Gamma \fCenter \Delta$
\Axiom$!B, \Gamma \fCenter \Delta$
\RightLabel{\LeftR{\lor}}
\BinaryInf$!A \lor !B, \Gamma \fCenter \Delta$
\DisplayProof
&
\begin{tabular}[t]{@{}r@{}}
\Axiom$\Gamma \fCenter \Delta, !A$
\RightLabel{\RightR{\lor}}
\UnaryInf$ \Gamma \fCenter \Delta, !A \lor !B$
\DisplayProof
\\[3ex]
\Axiom$ \Gamma \fCenter \Delta, !B$
\RightLabel{\RightR{\lor}}
\UnaryInf$ \Gamma \fCenter \Delta, !A \lor !B$
\DisplayProof\\[3ex]
\end{tabular}
\\[4ex]
\Axiom$ \Gamma \fCenter \Delta, !A$
\Axiom$ !B, \Gamma \fCenter \Delta$
\RightLabel{\LeftR{\lif}}
\BinaryInf$ !A \lif !B, \Gamma \fCenter \Delta$
\DisplayProof
&
\Axiom$ !A, \Gamma \fCenter \Delta, !B$
\RightLabel{\RightR{\lif}}
\UnaryInf$ \Gamma \fCenter \Delta, !A \lif !B $
\DisplayProof
\\[4ex]
\Axiom$ !A(t), \Gamma \fCenter \Delta$
\RightLabel{\LeftR{\lforall}}
\UnaryInf$ \lforall[x][!A(x)],\Gamma \fCenter \Delta$
\DisplayProof
&
\Axiom$ \Gamma \fCenter \Delta, !A(a) $
\RightLabel{\RightR{\lforall}}
\UnaryInf$ \Gamma \fCenter \Delta, \lforall[x][!A(x)]$
\DisplayProof
\\[4ex]
\Axiom$ !A(a), \Gamma \fCenter \Delta $
\RightLabel{\LeftR{\lexists}}
\UnaryInf$ \lexists[x][!A(x)], \Gamma \fCenter \Delta$
\DisplayProof
&
\Axiom$ \Gamma \fCenter \Delta, !A(t) $
\RightLabel{\RightR{\lexists}}
\UnaryInf$ \Gamma \fCenter \Delta, \lexists[x][!A(x)]$
\DisplayProof
\end{tabular}
\caption{\Log{G1c} دے قاعدے۔ $\RightR{\forall}$ تے
$\LeftR{\exists}$ وچ !!{constant}~$a$ نتیجے والے سیکوئنٹ وچ نہیں آؤنا چاہیدا۔ $\Gamma$ تے $ \Delta$، !!{formula}s دے کثیر سیٹ نیں۔}
\ollabel{tab:G1c}
\end{table}
\endgroup
% /OLP-0704
% OLP-0707 content/proof-theory/sequent-calculus/rules-G3c.tex
\begingroup

\begin{table}
\begin{tabular}{@{}c@{}c@{}}
\multicolumn{2}{@{}c@{}}{Axioms}\\[2ex]
$!C, \Gamma \Sequent \Delta, !C$ & $\lfalse, \Gamma \Sequent \Delta$
\\[2ex]
\multicolumn{2}{@{}c@{}}{Logical Rules}\\[2ex]
\Axiom$ \Gamma \fCenter \Delta, !A $
\RightLabel{\LeftR{\lnot}}
\UnaryInf$ \lnot !A, \Gamma \fCenter \Delta$
\DisplayProof
&
\Axiom$!A, \Gamma \fCenter \Delta$
\RightLabel{\RightR{\lnot}}
\UnaryInf$ \Gamma \fCenter \Delta, \lnot !A $
\DisplayProof
\\[4ex]
\Axiom$ !A, !B, \Gamma \fCenter \Delta$
\RightLabel{\LeftR{\land}}
\UnaryInf$ !A \land !B, \Gamma \fCenter \Delta$
\DisplayProof
&
\Axiom$\Gamma \fCenter \Delta, !A$
\Axiom$ \Gamma \fCenter \Delta, !B$
\RightLabel{\RightR{\land}}
\BinaryInf$ \Gamma \fCenter \Delta, !A \land !B $
\DisplayProof
\\[4ex]\Axiom$!A, \Gamma \fCenter \Delta$
\Axiom$!B, \Gamma \fCenter \Delta$
\RightLabel{\LeftR{\lor}}
\BinaryInf$!A \lor !B, \Gamma \fCenter \Delta$
\DisplayProof
&
\Axiom$\Gamma \fCenter \Delta, !A, !B$
\RightLabel{\RightR{\lor}}
\UnaryInf$ \Gamma \fCenter \Delta, !A \lor !B$
\DisplayProof
\\[4ex]
\Axiom$ \Gamma \fCenter \Delta, !A$
\Axiom$ !B, \Gamma \fCenter \Delta$
\RightLabel{\LeftR{\lif}}
\BinaryInf$ !A \lif !B, \Gamma \fCenter \Delta$
\DisplayProof
&
\Axiom$ !A, \Gamma \fCenter \Delta, !B$
\RightLabel{\RightR{\lif}}
\UnaryInf$ \Gamma \fCenter \Delta, !A \lif !B $
\DisplayProof
\\[4ex]
\Axiom$ !A(t), \lforall[x][!A(x)]\Gamma \fCenter \Delta$
\RightLabel{\LeftR{\lforall}}
\UnaryInf$ \lforall[x][!A(x)],\Gamma \fCenter \Delta$
\DisplayProof
&
\Axiom$ \Gamma \fCenter \Delta, !A(a) $
\RightLabel{\RightR{\lforall}}
\UnaryInf$ \Gamma \fCenter \Delta, \lforall[x][!A(x)]$
\DisplayProof
\\[4ex]
\Axiom$ !A(a), \Gamma \fCenter \Delta $
\RightLabel{\LeftR{\lexists}}
\UnaryInf$ \lexists[x][!A(x)], \Gamma \fCenter \Delta$
\DisplayProof
&
\Axiom$ \Gamma \fCenter \Delta, \lexists[x][!A(x)], !A(t) $
\RightLabel{\RightR{\lexists}}
\UnaryInf$ \Gamma \fCenter \Delta, \lexists[x][!A(x)]$
\DisplayProof
\end{tabular}
\caption{\Log{G3c} دے قاعدے۔ $\RightR{\forall}$ تے
$\LeftR{\exists}$ وچ !!{constant}~$a$ نتیجے والے سیکوئنٹ وچ نہیں آؤنا چاہیدا۔ $\Gamma$ تے $\Delta$، !!{formula}s دے کثیر سیٹ نیں۔ مسلّمیاں وچ !!{formula}~$!C$ جوہری ہونا چاہیدا اے۔}
\ollabel{tab:G3c}
\end{table}
\endgroup
% /OLP-0707

نظاماں وچ باریک فرق نیں تے اوہناں کر کے ثبوت دے نظریے نال متعلق اوہناں دیاں خاصیتاں وکھریاں نیں۔ \Log{G1c} دے مسلّمے $!A
\Sequent !A$ نیں، جیہناں وچ ہور !!{formula}s دی کھُل نہیں۔ \Log{G3c} دے مسلّمے $!C, \Gamma \Sequent \Delta, !C$ دی شکل دے نیں، جیہناں وچ «پاسے دے» کوئی وی !!{formula}s، $\Gamma$ تے~$\Delta$ وچ، ہو سکدے نیں، پر !!{formula}~$!C$ جیہڑا دوہاں پاسیاں اُتے آؤندا اے \emph{جوہری} ہونا چاہیدا اے۔ \Log{G1c} وچ $\LeftR{\land}$ تے $\RightR{\lor}$ قاعدیاں دیاں دو دو شکلاں نیں، جدکہ \Log{G3c} وچ اک اک اے۔ \Log{G1c} وچ وادھو «ساختی قاعدے» وی نیں، جنہاں نوں «سکڑاؤ» (\LeftR{\Contraction}، \RightR{\Contraction}) تے «کمزور کرنا» (\LeftR{\Weakening}، \RightR{\Weakening}) آکھدے آں۔

\Log{G1c} تے \Log{G3c} نظام کلاسیکی منطق لئی صحت مند تے کامل نیں۔ دوجے لفظاں وچ، ہر سیکوئنٹ جیہڑا ایہناں نظاماں وچ !!{provable} اے، ہر !!{structure} وچ سچا اے، تے ہر سیکوئنٹ جیہڑا ہر !!{structure} وچ سچا اے، اوہدا !!a{proof} اے۔ سو ایہ سوال کہ \emph{کی} کسے سیکوئنٹ دا !!a{proof} اے، منطق دے ہور طریقیاں نال وی حل ہو سکدا اے (جیویں نمونیاں دے نظریے دے طریقیاں نال)۔ تے کیونکہ دونے نظام عین اوہ سیکوئنٹ !!{prove} کردے نیں جیہڑے ہر !!{structure} وچ سچے نیں، ایس لئی خاص کر اوہ اوہی سیکوئنٹ ثابت کردے نیں۔

پر ثبوت دے نظریے دا واسطہ صرف ایس سوال نال نہیں کہ \emph{کی} کوئی گل ثابت ہو سکدی اے، بلکہ ایہ وی کہ \emph{کِنج}۔ ثبوت دے نظریے دے ماہر ایہو جیہے سوال ویکھدے نیں، «کی کسے خاص قاعدے نوں ہمیشہ ٹالیا جا سکدا اے؟»، «کی ایہ قاعدہ نظام وچ شامل کیتا جا سکدا اے بغیر ایہدے کہ نویں سیکوئنٹ ثابت ہون جوگ ہو جان؟»، یا «کی اک نظام دے !!{proof}s دوجے دے !!{proof}s وچ بدلے جا سکدے نیں؟» جے ایہو جیہے سوال دا جواب «ہاں» اے، تاں ایس گل دا ثبوت اکثر !!{proof}s دی گنجل اُتے استقرا نال، یا اوہو گل دوجے لفظاں وچ، !!{proof}s دی بناوٹ اُتے استقرا نال ملدا اے۔ ایس نال گل دا ثبوت ملدا اے (مثلاً جے \Log{G1c} اک سیکوئنٹ ثابت کردا اے تاں \Log{G3c} وی اوہی سیکوئنٹ ثابت کردا اے)، تے اک طریقہ وی (مثلاً !!{proof}s،~\Log{G1c} دے، نوں !!{proof}s،~\Log{G3c} دے، وچ بدلّن دا)۔

ثبوت دے نظریے دا اک مرکزی نتیجہ \emph{قطع ہٹاؤن دا قضیہ} اے۔ ایہ وکھاؤندا اے کہ اسیں قطع دا قاعدہ
\[
\Axiom$\Gamma \fCenter \Delta, !A$
\Axiom$!A, \Pi \fCenter \Lambda$
\RightLabel{\Cut}
\BinaryInf$\Gamma, \Pi \fCenter \Delta, \Lambda$
\DisplayProof
\]
سیکوئنٹ نظاماں وچ شامل کر سکدے آں، بغیر ایہدے کہ !!{provable} سیکوئنٹ وکھرے ہو جان۔ نالے ثبوت اک طریقہ دیندا اے جیہڑا کوئی وی !!{proof}، جیہڑا \Log{G1c} (یا \Log{G3c}) دے قاعدے تے نال \Cut{} ورتدا اے، ایہو جیہے ثبوت وچ بدل دیندا اے جیہڑا صرف \Log{G1c} (یا \Log{G3c}) دے قاعدے ورتدا اے۔ دوجے لفظاں وچ، ایہ طریقہ \Cut{} قاعدے نوں اوہناں !!{proof}s وچوں «ہٹا دیندا» اے جیہڑے اوہنوں ورتدے نیں۔
\endgroup
% /OLP-0700
% OLP-0711 content/proof-theory/sequent-calculus/rules-proofs.tex
\begingroup
\def\olfilename{rules-proofs}
\expandafter\def\csname ol@sectioncs\endcsname{section}

\olfileid{pt}{seq}{rp}

\olsection{قاعدے تے \usetoken{P}{proof}}

اسیں کجھ تعریفاں تے اصطلاحاں طے کرن توں شروع کردے آں۔

مثال لئی \Log{G3c} دے اوہ دو قاعدے ویکھو جیہڑے !!{proving} دی کھُل دیندے نیں اوہناں سیکوئنٹاں لئی جیہناں وچ $!A \land !B$ اے، اوہناں سیکوئنٹاں توں جیہناں وچ $!A$ تے~$!B$ نیں:
\[
\Axiom$ !A, !B, \Gamma \fCenter \Delta$
\RightLabel{\LeftR{\land}}
\UnaryInf$ !A \land !B, \Gamma \fCenter \Delta$
\DisplayProof
\qquad
\Axiom$\Gamma \fCenter \Delta, !A$
\Axiom$ \Gamma \fCenter \Delta, !B$
\RightLabel{\RightR{\land}}
\BinaryInf$ \Gamma \fCenter \Delta, !A \land !B $
\DisplayProof
\]
لکیر توں اُتلے سیکوئنٹاں نوں \emph{مقدمے} تے لکیر توں تھلّے والے سیکوئنٹ نوں \emph{نتیجہ} آکھدے آں۔ ہر قاعدے وچ !!{formula}s جیہڑے $ \Gamma$ تے $ \Delta$ وچ نیں، اوہناں نوں \emph{سیاق} یا \emph{پاسے دے !!{formula}s} آکھدے آں۔ نتیجے وچ اوہ !!{formula} جیہڑا سیاق دا حصہ نہیں، اوہ \emph{مُکھ} !!{formula} اے، تے مقدمیّاں دے اوہ !!{formula}s جیہڑے سیاق وچ نہیں، اوہ \emph{فعال} (یا \emph{معاون}) !!{formula}s نیں۔

مسوّراں دے قاعدے ذرا ہور گنجل دار نیں۔ مثال لئی \Log{G1c} وچ $\lforall$ لئی قاعدے ایہ نیں:
\[\Axiom$ !A(t), \Gamma \fCenter \Delta$
\RightLabel{\LeftR{\lforall}}
\UnaryInf$ \lforall[x][!A(x)],\Gamma \fCenter \Delta$
\DisplayProof
\qquad
\Axiom$ \Gamma \fCenter \Delta, !A(a) $
\RightLabel{\RightR{\lforall}}
\UnaryInf$ \Gamma \fCenter \Delta, \lforall[x][!A(x)]$
\DisplayProof
\]
\LeftR{\lforall} قاعدے وچ فعال !!{formula}~$!A(t)$ اے، جتھے $t$ بند ٹرم اے، یعنی متغیراں توں بناں ٹرم۔ ایہو جیہا استنباط ٹھیک اے، یعنی خاکہ وار قاعدے \LeftR{\lforall} نال میل کھاندا اے، جے کوئی !!{formula}~$!A$، کوئی متغیر~$x$، تے کوئی بند ٹرم~$t$ ایہو جیہے ہون کہ نتیجے دا مُکھ !!{formula}، $\lforall[x][!A]$ اے تے مقدمے دا فعال !!{formula}~$\Subst{!A}{t}{x}$ اے۔ \RightR{\lforall} قاعدہ وی ایویں کم کردا اے، فرق ایہ اے کہ کسے وی ٹرم~$t$ دی تھاں مقدمے دا فعال !!{formula}~$\Subst{!A}{a}{x}$ ہونا چاہیدا اے، جتھے $a$ اک !!a{constant} اے۔ ایس مستقل نوں \RightR{\lforall} استنباط دا \emph{خاص متغیر} آکھدے آں۔ استنباط صرف اوہدوں ٹھیک اے جدوں تھلّے والے سیکوئنٹ وچ~$a$ نہیں، یعنی جے $a$، $ \Gamma$، $ \Delta$ یا~$!A$ وچ نہیں آؤندا۔ ایس نوں \emph{خاص متغیر دی شرط} آکھدے آں۔

\Log{G1c} تے \Log{G3c} ورگے سیکوئنٹ نظام ایہو جیہے خاکہ وار دَسّے قاعدیاں دے مجموعے توں بندے نیں، نال کجھ سیکوئنٹ (اوہ وی خاکہ وار دَسّے ہوئے) جیہڑے مسلّمے منّے جاندے نیں۔ ایہو جیہے نظام وچ !!^{proof}s سیکوئنٹاں دے محدود درخت نیں جیہڑے مسلّمیّاں توں استنباطی قاعدے ورت کے استقرائی طور اُتے بندے نیں۔ ہر پَتّا مسلّمہ اے، تے ہر ہور سیکوئنٹ اوہدے جھٹ اُتلے سیکوئنٹاں توں نظام دے کسے اک استنباطی قاعدے نال نکلدا اے۔

\begin{defn}\ollabel{defn:proof}
سیکوئنٹ نظام وچ \emph{!!^{proof}s} سیکوئنٹاں دے محدود درخت نیں جیہڑے استقرائی طور اُتے ایویں بندے نیں:
\begin{enumerate}
  \item\ollabel{defn:prf:axiom} ہر مسلّمہ !!a{proof} اے۔
  \item\ollabel{defn:prf:one-prem} جے $\pi_1$ اک !!a{proof} اے جیہڑا سیکوئنٹ~$S_1$ اُتے مُکدا اے، تے $S$ اوہ سیکوئنٹ اے جیہڑا~$S_1$ توں اک مقدمے والے قاعدے~$R$ نال نکلدا اے، تاں
  \[
  \AxiomC{}
  \RightLabel{$\pi_1$}
  \DeduceC{$S_1$}
  \RightLabel{$R$}
  \UnaryInfC{$S$}
  \DisplayProof
  \]
  اک !!a{proof} اے۔
  \item\ollabel{defn:prf:two-prem} جے $\pi_1$ تے $\pi_2$، !!{proof}s نیں جیہڑے سیکوئنٹاں~$S_1$ تے~$S_2$ اُتے مُکدے نیں، تے $S$ اوہ سیکوئنٹ اے جیہڑا~$S_1$ تے $S_2$ توں دو مقدمیّاں والے قاعدے~$R$ نال نکلدا اے، تاں
  \[
  \AxiomC{}
  \RightLabel{$\pi_1$}
  \DeduceC{$S_1$}
  \AxiomC{}
  \RightLabel{$\pi_2$}
  \DeduceC{$S_2$}
  \RightLabel{$R$}
  \BinaryInfC{$S$}
  \DisplayProof
  \]
  اک !!a{proof} اے۔
\end{enumerate}
\end{defn}

\Olref{defn:proof}، !!{proof}s نوں سیکوئنٹاں دے درختاں دے طور اُتے متعین کردی اے۔ اسیں ایہناں درختاں نوں «اُتّے ول» ودھدا سمجھدے آں۔ ایس درخت دیاں سبھ توں اُتلیاں (پَتّیاں والیاں) گنڈھاں مسلّمے نیں تے اوہناں نوں \emph{شروع دے سیکوئنٹ} آکھدے آں۔ سبھ توں تھلّے والے سیکوئنٹ نوں \emph{آخری سیکوئنٹ} آکھدے آں۔ جے $\pi$ اک !!a{proof} اے جیہدا آخری سیکوئنٹ~$S$ اے، تاں اسیں $\pi \Proves S$ وی لکھدے آں۔ جے سیکوئنٹ~$S$ دا !!a{proof} اے، تاں اسیں $\Proves S$ لکھدے آں، تے ہو سکدا اے جیہڑے نظام وچ ثبوت اے اوہ~$\Proves$ نشانی دے کھبّے پاسے دَسّ دیئیے، جیویں $\Log{G1c} \Proves !A, !B \Sequent !A \land
!B$۔ !!a{proof} وچ ہر سیکوئنٹ، شروع دے سیکوئنٹاں توں چھُٹ، کسے قاعدے~$R$ نال ہوئے استنباط دا \emph{نتیجہ} اے۔ !!a{proof} وچ کسے سیکوئنٹ توں جھٹ اُتلے اک یا دو سیکوئنٹ (یعنی اوس گنڈھ دے ماپے) اوس استنباط دے \emph{مقدمے} نیں۔

اوہ !!{proof}s جیہڑے !!a{proof} دی استقرائی بناوٹ وچ ورتے جاندے نیں، اوہدے \emph{ذیلی !!{proof}s} نیں۔ اوہ ذیلی !!{proof}s جیہڑے استنباط دے مقدمیّاں اُتے مُکدے نیں، اوس استنباط دے \emph{سدھے ذیلی !!{proof}s} نیں۔

اکثر !!{proof}s دی گنجل نوں قدرتی عدد نال ماپنا پیندا اے۔ اسیں !!a{proof} وچ سیکوئنٹاں دی گنتی ای کر سکدے آں، پر اکثر ہور کم آؤن والی ماپ !!{height} اے، !!a{proof} دی۔ ایہ آخری سیکوئنٹ تے کسے شروع دے سیکوئنٹ دے وچکار سیکوئنٹاں دی سبھ توں وڈی گنتی اے۔ اسیں استقرائی تعریف وی دے سکدے آں۔

\begin{defn}
\emph{!!{height}}~$\pheight{\pi}$، !!a{proof}~$\pi$ دی، استقرائی طور اُتے ایویں متعین اے:
\begin{enumerate}
  \item جے $\pi$ صرف اک مسلّمہ اے، تاں $\pheight{\pi} = 0$۔
  \item جے $\pi$ اک مقدمے والے استنباط اُتے مُکدا اے، جیہدا سدھا ذیلی !!{proof}~$\pi_1$ اے، تاں $\pheight{\pi} = \pheight{\pi_1} + 1$۔
  \item جے $\pi$ دو مقدمیّاں والے استنباط اُتے مُکدا اے، جیہدے سدھے ذیلی !!{proof}s~$\pi_1$ تے~$\pi_2$ نیں، تاں $\pheight{\pi} =
  \max(\pheight{\pi_1}, \pheight{\pi_2}) + 1$۔
\end{enumerate}
جے سیکوئنٹ~$S$ دا !!a{proof} اے جیہدی اُچائی $\le n$ اے، تاں اسیں $\Proves[n] S$ لکھدے آں۔
\end{defn}

\begin{ex}
\Log{G3c} وچ $!C, \Gamma \Sequent \Delta, !C$ دی شکل دا ہر سیکوئنٹ مسلّمہ اے، جتھے $!C$ جوہری ہونا چاہیدا اے۔ منّ لو $!D$ تے $!E$ جوہری جملے نیں۔ تاں $!D, !E \Sequent !D$ تے $!D, !E \Sequent !E$، مسلّمے $!C, \Gamma \Sequent \Delta, !C$ دیاں مثالاں نیں۔ مثال لئی جے اسیں $!C \ident !E$، $\Gamma = \{!D\}$ تے $\Delta = \emptyset$ لئیے $!C, \Gamma \Sequent \Delta, !C$ وچ، تاں عین سیکوئنٹ $!D, !E
\Sequent !E$ ملدا اے۔ (چیتے رکھو کہ اسیں کثیر \emph{سیٹاں} نال کم کردے آں، سو $!D, !E \Sequent !E$ تے $!E, !D \Sequent !D$ اوہی سیکوئنٹ نیں۔)

چونکہ $!D, !E \Sequent !D$،~\Log{G3c} دے مسلّمے دی اک مثال اے، ایس لئی ایہ (اپنے آپ) !!a{proof} منّیا جاندا اے~\Log{G3c} وچ، تے ایویں ای $!D, !E
\Sequent !E$ وی، \olref{defn:proof}\olref{defn:prf:axiom} نال۔ آؤ ایہناں نوں $\pi_1$ تے~$\pi_2$ آکھیے۔ \olref{defn:prf:two-prem} نال اسیں ایہ دونے !!{proof}s لے کے اک نواں ($\pi_3$)، دو مقدمیّاں والے قاعدے~$\RightR{\land}$ نال، بنا سکدے آں:
\[
\Axiom$!D, !E \fCenter !D$
\Axiom$!D, !E \fCenter !E$
\RightLabel{\RightR{\land}}
\BinaryInf$!D, !E  \fCenter !D \land !E $
\DisplayProof
\]
$\pi_3$ توں اسیں اگّوں !!{proof}~$\pi_4$ لیندے آں
\[
\Axiom$!D, !E \fCenter !D$
\Axiom$!D, !E \fCenter !E$
\RightLabel{\RightR{\land}}
\insertBetweenHyps{\hspace{5ex}}
\BinaryInf$!D, !E  \fCenter !D \land !E $
\LeftSubproofLabel{$\pi_3$}
\RightLabel{\LeftR{\land}}
\UnaryInf$!D \land !E  \fCenter !D \land !E$
\RightSubproofLabel{$\pi_4$}
\DisplayProof
\]
اک مقدمے والا قاعدہ~$\LeftR{\land}$ ورت کے (\olref{defn:proof}\olref{defn:prf:one-prem} نال)۔ آخری استنباط دا سدھا ذیلی !!{proof}،~$\pi_4$ وچ،~$\pi_3$ اے۔ سدھے ذیلی !!{proof}s، $\RightR{\land}$ استنباط دے،، $\pi_1$ تے~$\pi_2$ نیں۔ ذیلی !!{proof}s، $ \pi_4$ دے، ایہ سارے نیں، $ \pi_1$، $ \pi_2$، $ \pi_3$، تے $ \pi_4$ آپ۔ ساڈے کول $\pheight{\pi_1} =
\pheight{\pi_2} = 0$، $\pheight{\pi_3} = 1$ تے $\pheight{\pi_4} = 2$ اے۔ اسیں وکھایا اے کہ $ \Log{G3c} \Proves[4] !D \land !E  \fCenter !D
\land !E$ کسے وی جوہری $!D$ تے~$!E$ لئی۔
\end{ex}
\endgroup
% /OLP-0711
% OLP-0702 content/proof-theory/sequent-calculus/proof-examples.tex
\begingroup
\def\olfilename{proof-examples}
\expandafter\def\csname ol@sectioncs\endcsname{section}

\olfileid{pt}{seq}{ex}

\olsection{ثبوتاں دیاں مثالاں}

سیکوئنٹاں دے !!{proof}s دین لئی عام طور اُتے سوکھا اے کہ آخری سیکوئنٹ توں اُتّے ول چلیے۔ اوہدے وچ !!a{formula} نوں فعال !!{formula} چُنو تے میل کھاندے قاعدے دے مقدمے سیکوئنٹ توں اُتّے لکھو۔ فیر ایویں کردے جاؤ جدوں تک مسلّمیّاں تک نہ پہنچ جاؤ۔ مثال لئی منّ لو اسیں ایہ سیکوئنٹ ثابت کرنا چاہنے آں:
\[(!C \land !D) \lif !E \Sequent !C \lif (!D \lif !E)\]
$!C \lif (!D \lif !E)$ ویکھو۔ ایہ $!A \lif !B$ دی شکل دا اے تے سیکوئنٹ دے سجّے پاسے آؤندا اے۔ پورے سیکوئنٹ نوں $\RightR{\lif}$ قاعدے، جیہڑا \Log{G1c} دا اے، نال ملاؤن توں
\[\Axiom$ !A, \Gamma \fCenter \Delta, !B$
\RightLabel{\RightR{\lif}}
\UnaryInf$ \Gamma \fCenter \Delta, !A \lif !B $
\DisplayProof\]
پتا لگدا اے کہ سانوں $\Gamma = \{!C \lif (!D \lif !E)\}$، $\Delta =
\emptyset$، $!A \ident !C$ تے $!B \ident !D \lif !E$ لینا چاہیدا اے۔ سو ثبوت دا آخری استنباط ایہ ہو سکدا اے:
\[\Axiom$ !C, (!C \land !D) \lif !E \fCenter !D \lif !E$
\RightLabel{\RightR{\lif}}
\UnaryInf$(!C \land !D) \lif !E \fCenter !C \lif (!D \lif !E)$
\DisplayProof\]
جے اسیں ایس خیال نوں مُڑ لائیے، $!D \lif !E$ نوں فعال فارمولا رکھ کے، تاں ایہ ملدا اے:
\[\Axiom$ !D, !C, (!C \land !D) \lif !E \fCenter !E$
\RightLabel{\RightR{\lif}}
\UnaryInf$!C, (!C \land !D) \lif !E \fCenter !D \lif !E$
\RightLabel{\RightR{\lif}}
\UnaryInf$(!C \land !D) \lif !E \fCenter !C \lif (!D \lif !E)$
\DisplayProof\]
ہُن سانوں $(!C \land !D) \lif !E$ ول دھیان دینا تے $\LeftR{\lif}$ قاعدہ، جیہڑا \Log{G1c} دا اے، ورتنا پینا اے:
\[\Axiom$ \Gamma \fCenter \Delta, !A$
\Axiom$ !B, \Gamma \fCenter \Delta$
\RightLabel{\LeftR{\lif}}
\BinaryInf$ !A \lif !B, \Gamma \fCenter \Delta$
\DisplayProof\]
ایہدے توں ایہ ملدا اے:
\[\Axiom$!D, !C \fCenter !E, !C \land !D$
\Axiom$!D, !C, !E \fCenter !E$
\RightLabel{\LeftR{\lif}}
\BinaryInf$ !D, !C, (!C \land !D) \lif !E \fCenter !E$
\RightLabel{\RightR{\lif}}
\UnaryInf$!C, (!C \land !D) \lif !E \fCenter !D \lif !E$
\RightLabel{\RightR{\lif}}
\UnaryInf$(!C \land !D) \lif !E \fCenter !C \lif (!D \lif !E)$
\DisplayProof\]
کھبّے پاسے دے سبھ توں اُتلے سیکوئنٹ وچ !!{formula} $!C \land !D$ نوں مُکھ رکھ کے تے $\RightR{\land}$ قاعدہ ورت کے سانوں ایہ ملدا اے:
\[
\Axiom$!D, !C \fCenter !E, !C$
\Axiom$!D, !C \fCenter !E, !D$
\RightLabel{\RightR{\land}}
\BinaryInf$!D, !C \fCenter !E, !C \land !D$
\Axiom$!D, !C, !E \fCenter !E$
\RightLabel{\RightR{\lif}}
\BinaryInf$ !D, !C, (!C \land !D) \lif !E \fCenter !E$
\RightLabel{\RightR{\lif}}
\UnaryInf$!C, (!C \land !D) \lif !E \fCenter !D \lif !E$
\RightLabel{\RightR{\lif}}
\UnaryInf$(!C \land !D) \lif !E \fCenter !C \lif (!D \lif !E)$
\DisplayProof
\]
ہُن تک کم ٹھیک چلدا اے۔ ویکھو کہ جیہڑا کجھ اسیں کیتا اے، اوہ \Log{G3c} وچ وی چلدا اے، کیونکہ ہُن تک ورتے قاعدے \Log{G1c} تے~\Log{G3c} وچ اوہی نیں۔ اسیں !!a{proof} دے ایس ٹوٹے تک پہنچ گئے آں جتھے سارے اُتلے سیکوئنٹاں دے کھبّے تے سجّے پاسے اوہی !!{formula} اے۔ پر اوہ ہجے پورے مسلّمے نہیں۔

\Log{G1c} وچ سانوں !!{proof}s، سبھ توں اُتلے سیکوئنٹاں دے،، $!A \Sequent !A$ دی شکل والے مسلّمیّاں توں دینے پین گے۔ ایہ کمزور کرن دے قاعدیاں نال سوکھا اے۔ مثال لئی !!{prove} ہو سکدا اے ایویں، سبھ توں کھبّا اُتلا سیکوئنٹ $!D, !C \Sequent
!E, !C$:
\[
\Axiom$!C \fCenter !C$
\RightLabel{\LeftR{\Weakening}}
\UnaryInf$!D, !C \fCenter !C$
\RightLabel{\RightR{\Weakening}}
\UnaryInf$!D, !C \fCenter !E, !C$
\]
سارا رلا کے ساڈے کول !!a{proof} اے \Log{G1c} وچ، جیہڑا ایویں دِسدا اے:
\[
\Axiom$!C \fCenter !C$
\RightLabel{\LeftR{\Weakening}}
\UnaryInf$!D, !C \fCenter !C$
\RightLabel{\RightR{\Weakening}}
\UnaryInf$!D, !C \fCenter !E, !C$
\Axiom$!D \fCenter !D$
\RightLabel{\LeftR{\Weakening}}
\UnaryInf$!D, !C \fCenter !D$
\RightLabel{\RightR{\Weakening}}
\UnaryInf$!D, !C \fCenter !E, !D$
\RightLabel{\RightR{\land}}
\BinaryInf$!D, !C \fCenter !E, !C \land !D$
\Axiom$!E \fCenter !E$
\RightLabel{\LeftR{\Weakening}}
\UnaryInf$!C, !E \fCenter !C$
\RightLabel{\LeftR{\Weakening}}
\UnaryInf$!D, !C, !E \fCenter !E$
\RightLabel{\LeftR{\lif}}
\BinaryInf$ !D, !C, (!C \land !D) \lif !E \fCenter !E$
\RightLabel{\RightR{\lif}}
\UnaryInf$!C, (!C \land !D) \lif !E \fCenter !D \lif !E$
\RightLabel{\RightR{\lif}}
\UnaryInf$(!C \land !D) \lif !E \fCenter !C \lif (!D \lif !E)$
\DisplayProof
\]
ایس ثبوت دی بناوٹ (خاص کر اُچائی) ایس اُتے منحصر نہیں کہ !!{formula}s $!C$، $!D$ تے~$!E$ کیہ نیں۔ اسیں اُتلے !!{proof} وچ ایہ نشان کسے وی !!{formula}s نال بدل سکدے آں تے حاصل ہجے وی !!{proof} ہووے گا~\Log{G1c} وچ۔ \Log{G1c} دا ثبوت \emph{خاکہ وار} اے۔

\Log{G3c} وچ مسلّمے $!A, \Gamma \Sequent
\Delta, !A$ دی شکل والے سیکوئنٹ نیں، جتھے $!A$ جوہری ہونا چاہیدا اے۔ سو بھاویں $!D, !C \Sequent
!E, !C$، \Log{G3c} دا مسلّمہ لگدا اے ($\Gamma = \{!D\}$، $\Delta =
\{!E\}$)، اوہ \emph{سچ مُچ} مسلّمہ اوہدوں ای اے جے $!C$ جوہری اے۔ ساڈا ثبوت ٹوٹا \Log{G3c} وچ پورا ثبوت اوہدوں ای اے جے $!C$، $!D$ تے $!E$ جوہری !!{formula}s نیں۔

بھاویں عین گل ایہ اے کہ اسیں ہجے ایہ نہیں وکھایا کہ
\[\Log{G1c} \Proves (!C \land !D) \lif !E \fCenter !C \lif (!D \lif
!E),\]
پر عملی طور اُتے سانوں اینا ای کرن دی لوڑ اے (تے اینا ای کر سکدے آں)۔ وجہ ایہ اے کہ $!A, \Gamma \Sequent \Delta, !A$ دی شکل دا ہر سیکوئنٹ \Log{G3c} وچ ثابت ہو سکدا اے، بھاویں $!A$ آپ جوہری نہ ہووے۔ اسیں ایہ گہرائی~$\depth{!A}$ اُتے استقرا نال وکھا سکدے آں، جیہڑی~$!A$ دی اے۔

\begin{defn}\ollabel{defn:depth}
\emph{گہرائی}~$\depth{!A}$، !!a{formula} دی، ایویں متعین اے:
\begin{enumerate}
  \item جے $!A$ جوہری اے، تاں $\depth{!A} = 0$۔
  \item جے $!A \ident \lnot !B$، $!A \ident \lforall[x][A]$، یا $!A \ident \lexists[x][A]$، تاں $\depth{!A} = \depth{!B} + 1$۔
  \item جے $!A \ident (!B \land !C)$، $(!B \lor !C)$، یا $(!B \lif
  !C)$، تاں $\depth{!A} = \max(\depth{!B},\depth{!C}) + 1$۔
\end{enumerate}
\end{defn}

ایہ ثابت کرنا اوکھا نہیں کہ $\depth{A(x)} = \depth{A(t)}$ ہر ٹرم~$t$ لئی۔ ساڈے لئی اہم گل ایہ اے کہ ہر قاعدے وچ مُکھ !!{formula} دی گہرائی ہمیشہ فعال !!{formula}s دی گہرائی توں ودھ اے۔ مثال لئی ساڈے کول ایہ اے:
\begin{align*}
\depth{P(x)} = \depth{P(a)} & = 0,\\
\depth{\lforall[x][P(x)]} & = 1,\\
\depth{(\lforall[x][P(x)] \lif P(a))} & = \max(1,0) + 1 = 2.
\end{align*}
اسیں ایس گل نوں ورت کے~$\depth{!A}$ اُتے استقرا نال کجھ نتیجے ثابت کر سکدے آں، جیویں ہِٹھلا نتیجہ۔

\begin{prop}\ollabel{prop:G1c-axioms}
ہر~$!A$ لئی !!a{proof} اے $\Log{G1c}$ وچ $!A \Sequent !A$ دا، جیہدے وچ سارے مسلّمے جوہری نیں۔
\end{prop}

\begin{proof}
  ~$ \depth{!A}$ اُتے استقرا نال۔ جے $\depth{!A} = 0$، تاں $!A$ جوہری اے تے $!A \Sequent !A$ پہلاں ای مطلوبہ شکل دا !!{proof} اے~\Log{G1c} وچ۔ جے $\depth{!A} > 0$، تاں $!A$ گھٹ گہرائی والے !!{formula}s توں بنیا اے، جیویں $!A \ident \lnot !B$، $!A \ident
  (!B \land !C)$، یا $!A \ident \lforall[x][A(x)]$۔ استقرا دا مفروضہ ایہ اے کہ ہر $!D$ لئی جیہدے لئی $\depth{!D} < \depth{!A}$، جوہری مسلّمیّاں توں $\Log{G1c} \Proves !D \Sequent !D$۔

  جے $!A \ident \lnot !B$، تاں $\depth{!B} < \depth{!A}$ تے استقرا دے مفروضے نال !!a{proof}~$\pi_1$ اے \Log{G1c} وچ، جوہری مسلّمیّاں توں $!B \Sequent !B$ دا۔ اسیں ایہنوں ایویں ودھا کے !!a{proof} بنا سکدے آں $\lnot !B \Sequent \lnot !B$ دا:
  \[
  \AxiomC{}
  \RightLabel{$\pi_1$}
  \Deduce$!B \fCenter !B$
  \RightLabel{\LeftR{\lnot}}
  \UnaryInf$\lnot !B, !B \fCenter $
  \RightLabel{\RightR{\lnot}}
  \UnaryInf$\lnot !B \fCenter \lnot !B$
  \DisplayProof
\]

  جے $!A \ident (!B \land !C)$، تاں $\depth{!B} < \depth{!A}$ تے $\depth{!C} < \depth{!A}$۔ استقرا دے مفروضے نال !!{proof}s~$\pi_1$ تے~$\pi_2$ نیں \Log{G1c} وچ، جوہری مسلّمیّاں توں $!B
  \Sequent !B$ تے $!C \Sequent !C$ دے۔ اسیں ایہناں نوں ایویں ودھا کے !!a{proof} بنا سکدے آں $!B \land !C \Sequent !B \land !C$ دا:
  \[
  \AxiomC{}
  \RightLabel{$\pi_1$}
  \Deduce$!B, \fCenter !B$
  \RightLabel{\LeftR{\Weakening}}
  \UnaryInf$!B, !C, \fCenter !B$
  \RightLabel{\LeftR{\land}}
  \UnaryInf$!B \land !C, !C\fCenter !B$
  \RightLabel{\LeftR{\land}}
  \UnaryInf$!B \land !C, !B \land !C \fCenter !B$
  \AxiomC{}
  \RightLabel{$\pi_2$}
  \Deduce$!C, \fCenter !C$
  \RightLabel{\LeftR{\Weakening}}
  \UnaryInf$!B, !C, \fCenter !C$
  \RightLabel{\LeftR{\land}}
  \UnaryInf$!B \land !C, !C \fCenter !C$
  \RightLabel{\LeftR{\land}}
  \UnaryInf$!B \land !C, !B \land !C \fCenter !C$
  \RightLabel{\RightR{\land}}
  \BinaryInf$!B \land !C \fCenter !B \land !C$
  \RightLabel{\LeftR{\Contraction}}
  \UnaryInf$!B \land !C \fCenter !B$
  \DisplayProof
  \]

  ہُن منّ لو $!A \ident \lforall[x][B(x)]$۔ اک !!a{constant}~$c$ لو جیہڑا $\lforall[x][B(x)]$ وچ نہیں آؤندا۔ فیر $\depth{B(c)} =
  \depth{!B(x)} < \depth{!A}$ تے استقرا دے مفروضے نال !!a{proof}~$\pi_1$ اے \Log{G1c} وچ، جوہری مسلّمیّاں توں $!B(c)
  \Sequent !B(c)$ دا۔ اسیں ایہنوں ایویں ودھا کے !!a{proof} بنا سکدے آں $\lforall[x][B(x)] \Sequent \lforall[x][B(x)]$ دا:
  \[
  \AxiomC{}
  \RightLabel{$\pi_1$}
  \Deduce$!B(c) \fCenter !B(c)$
  \RightLabel{\LeftR{\lforall}}
  \UnaryInf$\lforall[x][!B(x)] \fCenter !B(c)$
  \RightLabel{\RightR{\lforall}}
  \UnaryInf$\lforall[x][!B(x)] \fCenter \lforall[x][!B(x)]$
  \DisplayProof
\]

  باقی صورتاں مشقاں لئی چھڈّیاں گئیاں نیں۔
\end{proof}

\begin{prob}
  \olref{prop:G1c-axioms} دا ثبوت اوہ صورتاں کر کے پورا کرو جیہناں وچ $!A \ident (!B \lor !C)$، $!A \ident (!B \lif !C)$ تے $!A \ident \lexists[x][B(x)]$۔
\end{prob}

ساڈے مقصد لئی اسیں ایہ !!{proof} وی ورت سکدے ساں
\[
\AxiomC{}
\RightLabel{$\pi_1$}
\Deduce$!B \fCenter !B$
\RightLabel{\RightR{\lnot}}
\UnaryInf$\fCenter !B, \lnot !B$
\RightLabel{\LeftR{\lnot}}
\UnaryInf$\lnot !B \fCenter \lnot !B$
\DisplayProof
\]
اوس صورت وچ جتھے $!A \ident \lnot !B$۔ پر کجھ سیکوئنٹ نظام نیں جیہناں وچ ایہ کم نہ کردا۔ وجدانی سیکوئنٹ حساب \Log{G1i}، \Log{G1c} ورگا اے، پر اوہ ہر سیکوئنٹ دے سجّے پاسے ودھ توں ودھ اک !!{formula} دی کھُل دیندا اے۔ جے $\Delta = \emptyset$، تاں پہلا !!{proof}، !!a{proof} وی ہووے گا~\Log{G1i} وچ، پر دوجا نہیں، کیونکہ اک سیکوئنٹ دے سجّے پاسے $!B$ تے~$\lnot !B$ دونے نیں۔

ویکھو کہ \Log{G1c} وچ وی اسیں قاعدیاں دی ہور ترتیب نہیں ورت سکدے ساں اوس صورت وچ جتھے $!A \ident \lforall[x][!B(x)]$۔ ہِٹھلا \emph{نہیں} !!a{proof} اے
  \[
  \AxiomC{}
  \RightLabel{$\pi_1$}
  \Deduce$!B(c) \fCenter !B(c)$
  \RightLabel{\RightR{\lforall}}
  \UnaryInf$!B(c) \fCenter \lforall[x][!B(x)]$
  \RightLabel{\LeftR{\lforall}}
  \UnaryInf$\lforall[x][!B(x)] \fCenter \lforall[x][!B(x)]$
  \DisplayProof
\]
کیونکہ~$\RightR{\lforall}$ دی خاص متغیر والی شرط ٹُٹدی اے۔

\begin{prop}\ollabel{prop:G3c-axioms}
$!A, \Gamma
\Sequent \Delta, !A$ دی شکل دا ہر سیکوئنٹ (جتھے $!A$ دا جوہری ہونا \emph{ضروری نہیں})~\Log{G3c} وچ ثابت ہو سکدا اے۔
\end{prop}

\begin{proof}
ثبوت \olref{prop:G1c-axioms} دے ثبوت ورگا ای اے، پر استقرا دے مفروضے بارے دھیان رکھنا پیندا اے۔ جتھے $n =
\depth{!A} = 0$، اوہ صورت فیر سوکھی اے۔ جے $n=0$، تاں $!A$ جوہری اے، تے $!A,
\Gamma \Sequent \Delta, !A$،~\Log{G3c} دا مسلّمہ اے۔

ہُن منّ لو $n>0$۔ فیر اسیں صورتاں وکھ کردے آں۔ استقرا دا مفروضہ ایہ اے کہ ہر $D!$ لئی جیہدے لئی $\depth{!D} < n$، تے ہر $\Pi$ تے $\Lambda$ لئی، $\Log{G3c} \Proves !D, \Pi \Sequent \Lambda, !D$۔ ایہ صورت $!A \ident (!B \land !C)$ لئی اے۔ اسیں استقرا دا مفروضہ دو واری ورتدے آں۔ اک واری $!D = !B$، $\Pi = !C, \Gamma$ تے $\Lambda
= \Delta$ لئی، !!a{proof}~$\pi_1$ لین واسطے $!B, !C, \Gamma \Sequent
\Delta, !B$ دا، تے اک واری $!D = !C$، $\Pi = !B, \Lambda$ تے $\Lambda
= \Delta$ لئی، !!a{proof}~$\pi_2$ لین واسطے $!B, !C, \Gamma \Sequent
\Delta, !C$ دا۔ ایہ !!{proof} اے~$!B \land !C, \Gamma \Sequent
\Delta, !B \land !C$ دا:
\[
\AxiomC{}
\RightLabel{$\pi_1$}
\Deduce$!B, !C, \Gamma \fCenter \Delta, !B$
\RightLabel{\LeftR{\land}}
\UnaryInf$!B \land !C, \Gamma \fCenter \Delta, !B$
\AxiomC{}
\RightLabel{$\pi_2$}
\Deduce$!B, !C, \Gamma \fCenter \Delta, !C$
\RightLabel{\LeftR{\land}}
\UnaryInf$!B \land !C, \Gamma \fCenter \Delta, !C$
\RightLabel{\RightR{\land}}
\BinaryInf$!B \land !C, \Gamma \fCenter \Delta, !B \land !C$
\DisplayProof
\]

$!A \ident \lforall[x][!B(x)]$ لئی !!a{constant} چُنو جیہڑا~$\lforall[x][!B(x)]$، $\Gamma$ یا~$\Delta$ وچ نہیں۔ استقرا دا مفروضہ $!D = !B(c)$، $\Pi = \lforall[x][!B(x)], \Gamma$ تے $\Lambda = \Delta$ لئی لاؤ۔ ایہدے نال !!a{proof}~$\pi_1$ ملدا اے $!B(c),
\lforall[x][!B(x)], \Gamma \Sequent \Delta, !B(c)$ دا، جیہدے توں سانوں ایہ ملدا اے:
\[
\AxiomC{}
\RightLabel{$\pi_1$}
\Deduce$!B(c), \lforall[x][!B(x)], \Gamma \fCenter \Delta, !B(c)$
\RightLabel{\LeftR{\lforall}}
\UnaryInf$\lforall[x][!B(x)], \Gamma \fCenter \Delta, !B(c)$
\RightLabel{\RightR{\lforall}}
\UnaryInf$\lforall[x][!B(x)], \Gamma \fCenter \Delta, \lforall[x][!B(x)]$
\DisplayProof
\]
$\RightR{\lforall}$ دی خاص متغیر والی شرط ساڈے~$c$ چُنن دے طریقے نال پوری ہوندی اے۔
\end{proof}

\begin{prob}
  \olref{prop:G3c-axioms} دا ثبوت اوہ صورتاں کر کے اگے ودھاؤ جیہناں وچ $!A \ident (!B \lif !C)$ تے $!A \ident
  \lexists[x][B(x)]$۔
\end{prob}
\endgroup
% /OLP-0702
% OLP-0699 content/proof-theory/sequent-calculus/interpretation-rules.tex
\begingroup
\def\olfilename{interpretation-rules}
\expandafter\def\csname ol@sectioncs\endcsname{section}

\olfileid{pt}{seq}{irl}

\olsection{قاعدیاں دا مطلب}

سیکوئنٹ دا مطلب چیتے کرو۔ سیکوئنٹ $\Gamma \Sequent
\Delta$، !!a{structure}~$\Struct{M}$ وچ سچا اے جے یا گھٹّو گھٹ اک $!A \in \Gamma$ جھوٹھا اے، یا گھٹّو گھٹ اک~$!A \in \Delta$ سچا اے۔ اسیں~\Log{G3c} دے منطقی قاعدیاں نوں اوہناں دے مُکھ !!{formula}s دیاں سچ ہون دیاں شرطاں دا رسمی بیان سمجھ سکدے آں۔ \RightR{\lif} نوں ویکھو۔ اوہدا مُکھ !!{formula} $!A \lif !B$ اے۔ ایہ سچا اے جے $!A$ جھوٹھا اے یا~$!B$ سچا اے۔ سو سیکوئنٹ $\ \Sequent !A \lif !B$ اوہدوں تے صرف اوہدوں سچا اے جدوں $!A \Sequent !B$ سچا اے۔ جے اسیں سیاق دے !!{formula}s $\Gamma$، $\Delta$ ایہناں سیکوئنٹاں دے کھبّے تے سجّے پاسیاں وچ شامل کریے، تاں عین \RightR{\lif} قاعدے دا نتیجہ تے مقدمہ ملدے نیں۔~\LeftR{\lif} دی صورت رلدی ملدی اے۔ $!A
\lif !B$ جھوٹھا اے جے $!A$ سچا تے $!B$ جھوٹھا اے۔ سو $!A \lif !B
\Sequent \ $ اوہدوں تے صرف اوہدوں سچا اے جدوں $\ \Sequent !A$ تے $!B \Sequent \ $ دونے سچے نیں۔ \LeftR{\lif} دا نتیجہ اوہدے مقدمیّاں دے عطف دے برابر اے۔ ایہ بناوٹ پکی کردی اے کہ~\Log{G3c} دے قاعدے صحت مند نیں (جے سارے مقدمے سچے نیں تاں نتیجہ سچا اے)۔ ایہ تمامیت وی پکی کردی اے۔

اصل وچ کسے وی سچائی-فنکشن والے جوڑ دیاں سچ ہون دیاں شرطاں ایس طرح سیکوئنٹاں دے سیٹ دے روپ وچ بیان ہو سکدیاں نیں۔ ایہ ایس گل توں نکلدا اے کہ کلاسیکی منطق وچ ہر سچائی-فنکشن عطفی نارمل روپ والے !!a{formula} نال بیان ہو سکدا اے، یعنی جوہری تے نفی کیتے جوہری !!{formula}s دیاں فصلاں دے عطف نال۔ مثال لئی منّ لو $!A
\oplus !B$ اوہدوں تے صرف اوہدوں سچا اے جدوں $!A$ تے $!B$ وچوں عین اک سچا اے، یعنی «صرف اک والا یا»۔ فیر $!A \oplus !B$ اوہدوں تے صرف اوہدوں سچا اے جدوں $(!A \lor !B) \land
(\lnot !A \lor \lnot !B)$ سچا اے، تے ایہ جھوٹھا اے جے $(\lnot !A \lor !B)
\land (!A \lor \lnot !B)$ سچا اے۔ سو اسیں $\oplus$ لئی سیکوئنٹ حساب دے ایہ قاعدے متعارف کرا سکدے آں:
\begin{align*}
&
\Axiom$\Gamma \fCenter \Delta, !A, !B$
\Axiom$!A, !B, \Gamma \fCenter \Delta$
\RightLabel{\RightR{\oplus}}
\BinaryInf$\Gamma \fCenter \Delta, !A \oplus !B$
\DisplayProof
\\
& \Axiom$!A, \Gamma \fCenter \Delta, !B$
\Axiom$!B, \Gamma \fCenter \Delta, !A$
\RightLabel{\RightR{\oplus}}
\BinaryInf$!A \oplus !B, \Gamma \fCenter \Delta$
\DisplayProof
\end{align*}
تے ایہ قاعدے وی صحت مند تے کامل ہون گے۔ اوہناں دے نتیجے اوہدوں تے صرف اوہدوں سچے نیں جدوں سارے مقدمے سچے نیں۔

\begin{prob}
منّ لو $!A \downarrow !B$ اوہدوں تے صرف اوہدوں سچا اے جدوں نہ $!A$ سچا اے نہ~$!B$ («NOR»)۔ \Log{G3c} وچ ایہدے قاعدے کیہڑے ہون گے؟ (دو سیکوئنٹ لبھو جیہڑے اکّو ویلے اوہدوں تے صرف اوہدوں سچے ہون جدوں $!A \downarrow !B$ سچا اے۔ ایس نال \RightR{\downarrow} قاعدہ ملدا اے۔ فیر اک سیکوئنٹ لبھو جیہڑا اوہدوں تے صرف اوہدوں سچا ہووے جدوں $!A \downarrow !B$ جھوٹھا اے۔ ایس نال \LeftR{\downarrow} قاعدہ ملدا اے۔)
\end{prob}

\Log{G3c} وچ ہر جوڑ تے مسوّر دے عین دو قاعدے نیں، اک کھبّا قاعدہ تے اک سجّا۔ پر \Log{G1c} وچ دو \LeftR{\land} تے دو \RightR{\lor} قاعدے نیں۔ وجہ ایہ اے کہ گینٹسن دے اصل نظام~\Log{LK}، جیہدے اُتے \Log{G1c} بنیا اے، دی اک شکل~\Log{LJ} اہم غیر کلاسیکی منطق، یعنی وجدانی منطق، لئی سی۔ وجدانی منطق لئی صحت مند تے کامل نظام لین واسطے~\Log{LJ} وچ ہر سیکوئنٹ دا سجّا پاسا ودھ توں ودھ اک !!{formula} رکھ سکدا اے۔ ایس کرکے \RightR{\lor} قاعدہ~\Log{G3c} دا ورتیا نہیں جا سکدا، کیونکہ اوہدے مقدمے دے سجّے پاسے $!A$ تے~$!B$ دونے نیں۔ ایس لئی گینٹسن نے \RightR{\lor} قاعدیاں دا اک جوڑا \Log{LJ} لئی ورتیا تے اوہی~\Log{LK} لئی وی ورتیا۔ ایویں ای وجدانی شکل~\Log{G1i}، بس~\Log{G1c} اے جیہدے سارے سجّے پاسے ودھ توں ودھ اک !!{formula} رکھ سکدے نیں۔ (\Log{G3c} دی وی اک وجدانی شکل اے، پر اوہدے کجھ قاعدے~\Log{G3c} دے میل کھاندے قاعدیاں توں وکھرے نیں۔ مثال لئی اوہ وی \RightR{\lor} قاعدیاں دا جوڑا ورتدی اے۔)

ایہ ویکھنا سوکھا اے کہ~\Log{G1c} دے ساختی قاعدے صحت مند نیں۔ مثال لئی جے $\Gamma \Sequent \Delta$ دے سجّے پاسے کوئی جھوٹھا !!{formula} یا کھبّے پاسے کوئی سچا !!{formula} اے، تاں $\Gamma \Sequent
\Delta, !A$ وچ وی اوہی اے۔ $!A$ دے سچا یا جھوٹھا ہون نال کوئی فرق نہیں پیندا۔ سو جے~\RightR{\Weakening} دا مقدمہ سچا اے، تاں اوہدا نتیجہ وی سچا اے۔ ویکھو کہ ایتھے الٹی گل سچی نہیں۔ ہو سکدا اے نتیجہ ایس لئی سچا ہووے کہ $!A$ سچا اے، تے فیر مقدمے دے سچے ہون دی ضمانت نہیں۔

ساڈے کول ساختی قاعدے کیوں نیں؟ مثال لئی سکڑاؤ (\RightR{\Contraction}، \LeftR{\Contraction}) دی لوڑ اے !!{prove} کرن لئی $\ \Sequent !A
\lor \lnot !A$۔ جے تسی $!A \Sequent !A$ توں شروع کرو تاں پہلاں $\Sequent !A, \lnot !A$ ملدا اے~\RightR{\lnot} نال۔ فیر تسی \RightR{\lor} دیاں دونے شکلاں واری واری ورت سکدے او، اک واری $!A \lor \lnot
!A$ لین لئی $\lnot !A$ توں، تے اک واری $!A \lor \lnot !A$ لین لئی~$!A$ توں۔ تہانوں $\ \Sequent !A \lor \lnot !A, !A \lor \lnot !A$ ملدا اے۔ سو $\Sequent !A \lor \lnot !A$ لین لئی~\Log{G1c} وچ تہانوں اکّو !!{formula} دیاں دو موجودگیاں !!a{proof} وچ «سکیڑن» جوگے ہونا پینا اے:
\[
\Axiom$!A \fCenter !A$
\RightLabel{\RightR{\lnot}}
\UnaryInf$\fCenter !A, \lnot !A$
\RightLabel{\RightR{\lor}}
\UnaryInf$\fCenter !A, !A \lor \lnot !A$
\RightLabel{\RightR{\lor}}
\UnaryInf$\fCenter !A \lor \lnot !A, !A \lor \lnot !A$
\RightLabel{\RightR{\Contraction}}
\UnaryInf$\fCenter !A \lor \lnot !A$
\DisplayProof
\]
اصل وچ معتبر سیکوئنٹ وی نیں جیہڑے تسی !!{prove} نہیں کر سکدے~\Log{G1c} وچ، کمزور کرن یا سکڑاؤ توں بناں۔ مثال لئی !!a{proof} جیہڑا سیکوئنٹ $!A \Sequent !B \lif !A$ دا اے~\Log{G1c} وچ، اوہ \RightR{\lif} اُتے ضرور مُکدا اے (کیونکہ صرف $!B \lif !A$ مُکھ ہو سکدا اے)، تے ایس لئی $!B, !A \Sequent !A$ مقدمہ چاہیدا اے۔ ایہ لین لئی \LeftR{\Weakening} چاہیدا اے، مسلّمے $!A \Sequent !A$ توں:
\[
\Axiom$!A \fCenter !A$
\RightLabel{\LeftR{\Weakening}}
\UnaryInf$!B, !A \fCenter !A$
\RightLabel{\RightR{\lif}}
\UnaryInf$!A \fCenter !B \lif !A$
\DisplayProof
\]

\Log{G3c} وچ سکڑاؤ تے کمزور کرن دے قاعدیاں دی بالکل لوڑ نہیں۔ کمزور کرنا مسلّمیاں وچ ای شامل اے۔ دو مقدمیّاں والے قاعدیاں وچ سیاق $\Gamma$، $\Delta$ دیاں دو نقلاں نوں اک بنا کے سکڑاؤ زیادہ تر ٹالیا جا سکدا اے۔ \Log{G1c} تے \Log{G3c} دوناں وچ ایہو جیہے «سانجھے سیاق» والے قاعدے نیں۔ اک مقدمے والے قاعدیاں وچ سکڑاؤ صرف \RightR{\lor}، \LeftR{\land}، \RightR{\lexists} تے~\LeftR{\lforall} لئی چاہیدا اے۔~\Log{G3c} وچ ایہناں قاعدیاں دیاں شکلاں سکڑاؤ نوں بالکل بے لوڑ کر دیندیاں نیں۔ اک نظام \Log{G2c} اے (ویکھو \olref{tab:G2c})، جیہدے وچ دو مقدمیّاں والے قاعدیاں دے سیاق اک دوجے توں آزاد نیں۔ \Log{G2c} وچ سکڑاؤ ہور وی اہم اے،~\Log{G1c} نالوں۔

% OLP-0706 content/proof-theory/sequent-calculus/rules-G2c.tex
\begingroup

\begin{table}
\begin{tabular}{@{}c@{\hskip-5ex}c@{}}
\multicolumn{2}{@{}c@{}}{Axioms}\\[2ex]
$!A \Sequent !A$ & $\lfalse \Sequent \quad$
\\[2ex]
\multicolumn{2}{@{}c@{}}{Structural Rules}\\[2ex]
\Axiom$ \Gamma \fCenter \Delta $
\RightLabel{\LeftR{\Weakening}}
\UnaryInf$ !A, \Gamma \fCenter \Delta$
\DisplayProof
&
\Axiom$ \Gamma \fCenter \Delta$
\RightLabel{\RightR{\Weakening}}
\UnaryInf$ \Gamma \fCenter \Delta, !A$
\DisplayProof
\\[4ex]
\Axiom$ !A, !A, \Gamma \fCenter \Delta $
\RightLabel{\LeftR{\Contraction}}
\UnaryInf$ !A, \Gamma \fCenter \Delta$
\DisplayProof
&
\Axiom$ \Gamma \fCenter \Delta, !A, !A$
\RightLabel{\RightR{\Contraction}}
\UnaryInf$ \Gamma \fCenter \Delta, !A$
\DisplayProof
\\[4ex]
\multicolumn{2}{@{}c@{}}{Logical Rules}\\[2ex]
\Axiom$ \Gamma \fCenter \Delta, !A $
\RightLabel{\LeftR{\lnot}}
\UnaryInf$ \lnot !A, \Gamma \fCenter \Delta$
\DisplayProof
&
\Axiom$!A, \Gamma \fCenter \Delta$
\RightLabel{\RightR{\lnot}}
\UnaryInf$ \Gamma \fCenter \Delta, \lnot !A $
\DisplayProof
\\[4ex]\begin{tabular}[t]{@{}l@{}}
\Axiom$ !A, \Gamma \fCenter \Delta$
\RightLabel{\LeftR{\land}}
\UnaryInf$ !A \land !B, \Gamma \fCenter \Delta$
\DisplayProof
\\[3ex]
\Axiom$!B, \Gamma \fCenter \Delta$
\RightLabel{\LeftR{\land}}
\UnaryInf$!A \land !B, \Gamma \fCenter \Delta$
\DisplayProof\\[3ex]
\end{tabular}
&
\Axiom$\Gamma_1 \fCenter \Delta_1, !A$
\Axiom$ \Gamma_2 \fCenter \Delta_2, !B$
\RightLabel{\RightR{\land}}
\BinaryInf$ \Gamma_1, \Gamma_2 \fCenter \Delta_1, \Delta_2, !A \land !B $
\DisplayProof
\\[4ex]\Axiom$!A, \Gamma_1 \fCenter \Delta_1$
\Axiom$!B, \Gamma_2 \fCenter \Delta_2$
\RightLabel{\LeftR{\lor}}
\BinaryInf$!A \lor !B, \Gamma_1, \Gamma_2 \fCenter \Delta_1, \Delta_2$
\DisplayProof
&
\begin{tabular}[t]{@{}r@{}}
\Axiom$\Gamma \fCenter \Delta, !A$
\RightLabel{\RightR{\lor}}
\UnaryInf$ \Gamma \fCenter \Delta, !A \lor !B$
\DisplayProof
\\[3ex]
\Axiom$ \Gamma \fCenter \Delta, !B$
\RightLabel{\RightR{\lor}}
\UnaryInf$ \Gamma \fCenter \Delta, !A \lor !B$
\DisplayProof\\[3ex]
\end{tabular}
\\[4ex]
\Axiom$ \Gamma_1 \fCenter \Delta_1, !A$
\Axiom$ !B, \Gamma_2 \fCenter \Delta_2$
\RightLabel{\LeftR{\lif}}
\BinaryInf$ !A \lif !B, \Gamma_1, \Gamma_2 \fCenter \Delta_1, \Delta_2$
\DisplayProof
&
\Axiom$ !A, \Gamma \fCenter \Delta, !B$
\RightLabel{\RightR{\lif}}
\UnaryInf$ \Gamma \fCenter \Delta, !A \lif !B $
\DisplayProof
\\[4ex]
\Axiom$ !A(t), \Gamma \fCenter \Delta$
\RightLabel{\LeftR{\lforall}}
\UnaryInf$ \lforall[x][!A(x)],\Gamma \fCenter \Delta$
\DisplayProof
&
\Axiom$ \Gamma \fCenter \Delta, !A(a) $
\RightLabel{\RightR{\lforall}}
\UnaryInf$ \Gamma \fCenter \Delta, \lforall[x][!A(x)]$
\DisplayProof
\\[4ex]
\Axiom$ !A(a), \Gamma \fCenter \Delta $
\RightLabel{\LeftR{\lexists}}
\UnaryInf$ \lexists[x][!A(x)], \Gamma \fCenter \Delta$
\DisplayProof
&
\Axiom$ \Gamma \fCenter \Delta, !A(t) $
\RightLabel{\RightR{\lexists}}
\UnaryInf$ \Gamma \fCenter \Delta, \lexists[x][!A(x)]$
\DisplayProof
\end{tabular}
\caption{\Log{G2c} دے قاعدے۔ $\RightR{\forall}$ تے
$\LeftR{\exists}$ وچ !!{constant}~$a$ نتیجے والے سیکوئنٹ وچ نہیں آؤنا چاہیدا۔ $\Gamma$ تے $\Delta$، !!{formula}s دے کثیر سیٹ نیں۔}
\ollabel{tab:G2c}
\end{table}
\endgroup
% /OLP-0706
\endgroup
% /OLP-0699
% OLP-0703 content/proof-theory/sequent-calculus/quantifiers.tex
\begingroup
\def\olfilename{quantifiers}
\expandafter\def\csname ol@sectioncs\endcsname{section}

\olfileid{pt}{seq}{qua}

\olsection{باقاعدہ ثبوت تے تھانواں بدلنا}

مسوّراں دے قاعدیاں وچ کجھ اوکڑاں «خاص متغیر دیاں شرطاں» کر کے آؤندیاں نیں جیہڑیاں \LeftR{\lexists} تے~\RightR{\lforall} قاعدیاں اُتے لگدیاں نیں۔ ایہناں وچوں بہتی اوکڑ دور ہو سکدی اے جے پکا کریے کہ ایہناں استنباطاں وچ !!{constant}~$a$ کدے مُڑ نہ ورتیا جاوے۔ آؤ ہِٹھلی تعریف دیئیے۔

\begin{defn}
!!^a{proof}~$\pi$، \Log{G1c} یا \Log{G3c} وچ، \emph{باقاعدہ} اے جے
\begin{enumerate}
  \item کوئی خاص متغیر~$a$ اک توں ودھ \LeftR{\lexists} یا~\RightR{\lforall} استنباط دے خاص متغیر دے طور اُتے نہیں ورتیا جاندا، تے
  \item جے $a$، \LeftR{\lexists} یا~\RightR{\lforall} استنباط دا خاص متغیر اے، تاں اوہ~$\pi$ وچ صرف اوس استنباط توں اُتّے آؤندا اے۔
\end{enumerate}
\end{defn}

\begin{lem}\ollabel{lem:subst-regular}
جے $\pi$ باقاعدہ اے، $ \Gamma \Sequent \Delta$ اُتے مُکدا اے، $c$ اک !!a{constant} اے جیہڑا~$\pi$ وچ خاص متغیر دے طور اُتے نہیں ورتیا جاندا، تے $t$ ایہو جیہا ٹرم اے جیہدے وچ~$\pi$ دا کوئی خاص متغیر نہیں، تاں $ \Subst{\pi}{t}{c}$، جیہڑا~$\pi$ وچ~$c$ دی ہر موجودگی نوں~$t$ نال بدل کے ملدا اے، باقاعدہ !!{proof} اے۔ $ \Subst{\pi}{t}{c}$ دا آخری سیکوئنٹ $ \Subst{\Gamma}{t}{c}
\Sequent \Subst{\Delta}{t}{c}$ اے، جتھے $ \Subst{\Gamma}{t}{c} =
\Setabs{!A(t)}{!A(c) \in \Gamma}$۔
\end{lem}

\begin{proof}
  $ \pheight{\pi}$ اُتے استقرا نال۔ جے $ \pheight{\pi} = 0$، تاں ایہ بس اک مسلّمہ اے۔ فیر $ \Subst{\pi}{t}{c}$ وی اک مسلّمہ اے، کیونکہ اوہدے وچ کوئی وی !!{formula}~$!A(c)$، $!A(t)$ بن جاندا اے۔

  جے $ \pheight{\pi} > 0$، تاں~$\pi$ دے آخری استنباط دے حساب نال صورتاں وکھ کردے آں۔ ساختی قاعدیاں (\Log{G1c} وچ) تے قضیاتی جوڑاں دے منطقی قاعدیاں دیاں صورتاں سوکھیاں نیں تے مشقاں لئی چھڈّیاں گئیاں نیں۔ اسیں اوہ صورتاں ویکھدے آں جیہناں وچ $\pi$ مسوّر والے استنباط اُتے مُکدا اے۔
  \begin{enumerate}
    \item جے آخری استنباط $\RightR{\lforall}$ اے، تاں $\pi$ دی شکل ایہ اے:
    \[
    \AxiomC{}
    \RightLabel{$\pi'$}
    \Deduce$\Gamma(c) \fCenter \Delta'(c), !A(a,c)$
    \RightLabel{\RightR{\lforall}}
    \UnaryInf$\Gamma(c) \fCenter \Delta'(c), \lforall[x][!A(x,c)]$
    \DisplayProof
    \]
    استقرا دے مفروضے نال، $ \Subst{\pi'}{t}{c}$ باقاعدہ !!{proof} اے $ \Gamma(t) \fCenter \Delta'(t), !A(a,t)$ دا۔ کیونکہ $a$،~$\pi$ دا خاص متغیر اے، $a$،~$t$ وچ نہیں آؤندا۔ سو $ \Subst{\pi}{t}{c}$ دا آخری استنباط،
    \[
    \AxiomC{}
    \RightLabel{$\Subst{\pi'}{t}{c}$}
    \Deduce$\Gamma(t) \fCenter \Delta'(t), !A(a,t)$
    \RightLabel{\RightR{\lforall}}
    \UnaryInf$\Gamma(t) \fCenter \Delta'(t), \lforall[x][!A(x,t)]$
    \DisplayProof
    \]
    ٹھیک اے، کیونکہ نتیجے وچ~$a$ نہیں۔
    \item جے آخری استنباط $\LeftR{\lforall}$ اے \Log{G3c} وچ، تاں $\pi$ دی شکل ایہ اے:
    \[
    \AxiomC{}
    \RightLabel{$\pi'$}
    \Deduce$!A(s(c),c), \lforall[x][!A(x,c)], \Gamma(c) \fCenter \Delta'(c)$
    \RightLabel{\RightR{\lforall}}
    \UnaryInf$\lforall[x][!A(x,c)], \Gamma(c) \fCenter \Delta'(c)$
    \DisplayProof
    \]
    استقرا دے مفروضے نال، $ \Subst{\pi'}{t}{c}$ باقاعدہ !!{proof} اے $!A(s(t),t), \lforall[x][!A(x,t)], \Gamma(t) \fCenter
    \Delta'(t)$ دا۔ $ \Subst{\pi}{t}{c}$ دا آخری استنباط،
    \[
    \AxiomC{}
    \RightLabel{$\Subst{\pi'}{t}{c}$}
    \Deduce$!A(s(t),t), \lforall[x][!A(x,t)], \Gamma(t) \fCenter \Delta'(t)$
    \RightLabel{\RightR{\lforall}}
    \UnaryInf$\lforall[x][!A(x,t)], \Gamma(t) \fCenter \Delta'(t)$
    \DisplayProof
    \]
    ٹھیک اے۔\LeftR{\lforall} والی صورت~\Log{G1c} وچ ایہو جیہی اے، پر ذرا گھٹ گنجل دار۔
    \item آخری استنباط \RightR{\lexists} یا \LeftR{\lexists} اے۔ ایہ مشق اے۔
  \end{enumerate}
  ہر صورت وچ اسیں باقاعدہ !!{proof} دے آخر وچ اک سیکوئنٹ شامل کیتا (یا دو باقاعدہ !!{proof}s دے آخر وچ، دو مقدمیّاں والے قاعدیاں دی صورت وچ)۔ نواں سیکوئنٹ $\pi$ دے آخری سیکوئنٹ توں صرف ایس گل وچ وکھرا اے کہ $c$ نوں ٹرم~$t$ نال بدلیا گیا اے۔ نالے $\pi$ دے خاص متغیر، تے کیونکہ اسیں منّیا سی کہ $t$ وچ~$\pi$ دا کوئی خاص متغیر نہیں، تے $ \Subst{\pi}{t}{c}$ دے خاص متغیر اوہی نیں، ایس لئی نویں آخری سیکوئنٹ وچ $ \Subst{\pi}{t}{c}$ دا کوئی خاص متغیر نہیں۔ سو $ \Subst{\pi}{t}{c}$ باقاعدہ اے۔
\end{proof}

\begin{prop}
جے $\pi$، !!a{proof} اے \Log{G1c} یا \Log{G3c} وچ، تاں اک باقاعدہ !!{proof}~$\pi'$ اے جیہدی بناوٹ اوہی اے (خاص کر !!{height} اوہی اے) تے جیہڑا $\pi$ توں بس خاص متغیر بدلّن نال وکھرا اے۔
\end{prop}

\begin{proof}
صرف ایس ثبوت لئی~$\pi$ دے خاص متغیر والے استنباط نوں \emph{صاف} آکھو جے اوہدا خاص متغیر کسے ہور استنباط دا خاص متغیر نہیں تے~$\pi$ وچ اوس استنباط دے سدھے ذیلی !!{proof} توں چھُٹ کِتے ہور نہیں آؤندا۔ نئیں تاں اوہنوں \emph{گندا} آکھو۔ منّ لو $n$،~$\pi$ وچ گندے خاص متغیر والے استنباطاں دی گنتی اے۔ اسیں~$n$ اُتے استقرا نال وکھاؤنے آں کہ $\pi$ نوں مطلوبہ باقاعدہ !!{proof}~$\pi'$ وچ بدلیا جا سکدا اے۔ جے $n=0$، تاں~$\pi$ دے سارے خاص متغیر والے استنباط صاف نیں، سو $\pi$ باقاعدہ اے تے ثابت کرن نوں کجھ نہیں۔ نئیں تاں خاص متغیر والا سبھ توں اُتلا استنباط چُنو، منّ لو \RightR{\lforall}۔ذیلی !!{proof}~$\pi_1$،~$\pi$ دا، جیہڑا اوس اُتے مُکدا اے، ایہ شکل رکھدا اے:
    \[
    \AxiomC{}
    \RightLabel{$\pi_2$}
    \Deduce$\Gamma \fCenter \Delta, !A(c)$
    \RightLabel{\RightR{\lforall}}
    \UnaryInf$\Gamma \fCenter \Delta, \lforall[x][!A(x)]$
    \DisplayProof
    \]
ویکھو کہ چونکہ $c$ خاص متغیر اے، اوہ~$\Gamma$، $ \Delta$ یا $ \lforall[x][!A(x)]$ وچ نہیں آؤندا۔ ثبوت $ \pi_2$ باقاعدہ اے کیونکہ اوہدے وچ خاص متغیر والا کوئی استنباط ہے ای نہیں۔ منّ لو $a$ اک !!a{constant} اے جیہڑا~$\pi$ وچ نہیں آؤندا۔ \olref{lem:subst-regular} نال، $ \Subst{\pi_2}{a}{c}$، $ \Gamma \fCenter \Delta,
!A(a)$ دا باقاعدہ ثبوت اے جیہدے اُتے اسیں~$\RightR{\lforall}$ لا سکدے آں۔ جے اسیں $ \pi_1$ نوں~$\pi$ وچ ثبوت $ \pi_1'$ نال بدلیے،
    \[
    \AxiomC{}
    \RightLabel{$\Subst{\pi_2}{a}{c}$}
    \Deduce$\Gamma \fCenter \Delta, !A(a)$
    \RightLabel{\RightR{\lforall}}
    \UnaryInf$\Gamma \fCenter \Delta, \lforall[x][!A(x)]$
    \DisplayProof
    \]
تاں اسیں خاص متغیر والا اک گندا استنباط صاف استنباط نال بدل دِتّا اے، تے اکّو فرق خاص متغیر $c$ نوں~$a$ نال بدلنا اے۔ حاصل ہوئے ثبوت وچ $n-1$ گندے خاص متغیر والے استنباط نیں، تے نتیجہ استقرا دے مفروضے توں نکلدا اے۔
\end{proof}

اسیں ہُن ثابت کیتے قضیے دا سدھا حوالہ نہیں دیاں گے، پر منّاں گے کہ سارے !!{proof}s جیہڑے اسیں ویکھدے آں باقاعدہ نیں۔ جے اوہ باقاعدہ نہیں، تاں اسیں ہمیشہ خاص متغیر بدل کے اوہناں نوں باقاعدہ بنا سکدے آں۔
\endgroup
% /OLP-0703
% OLP-0698 content/proof-theory/sequent-calculus/admissible-derivable.tex
\begingroup
\def\olfilename{admissible-derivable}
\expandafter\def\csname ol@sectioncs\endcsname{section}

\olfileid{pt}{seq}{adm}

\olsection{قبول کرن جوگ تے اخذ ہو سکّن والے قاعدے}

ثبوت دے نظریے دے کئی نتیجے ایس گل بارے نیں، یا ایس گل بارے نتیجے ورتدے نیں، کہ جو کجھ کسے خاص قاعدے نال !!{prove} ہو سکدا اے، کی اوہ سارا اوس قاعدے توں بناں وی !!{prove} ہو سکدا اے۔ قطع ہٹاؤن دا قضیہ سبھ توں مشہور مثال اے۔ ایہ وکھاؤندا اے کہ ہر اوہ سیکوئنٹ جیہڑا~\Cut{} سمیت کسے سیکوئنٹ نظام دے قاعدیاں نال !!{prove} ہو سکدا اے، اوہ~\Cut توں بناں وی !!{prove} ہو سکدا اے۔ ایس طرح دی خاصیت اینی اہم اے کہ ایہدا اک ناں اے۔

\begin{defn}
قاعدہ~$R$ کسے سیکوئنٹ نظام وچ \emph{قبول کرن جوگ} اے جے ہر اوہ سیکوئنٹ جیہڑا اوس نظام تے قاعدے~$R$ نال ثابت ہو سکدا اے، اوہ~$R$ توں بناں وی ثابت ہو سکدا اے۔
\end{defn}

\begin{ex}\ollabel{ex:land-deriv}
قاعدے $\LeftR{\land}$ دیاں شکلاں \Log{G1c} تے \Log{G3c} وچ وکھریاں نیں۔ \Log{G1c} وچ قاعدے دیاں دو شکلاں نیں۔ اک دے مقدمے وچ کھبّا عطف جزو $!A$ اے، مُکھ !!{formula}~$!A \land !B$ دا، تے دوجے دے مقدمے وچ سجّا عطف جزو~$!B$:
\[
\Axiom$ !A, \Gamma \fCenter \Delta$
\RightLabel{$\LeftR{\land}_1$}
\UnaryInf$ !A \land !B, \Gamma \fCenter \Delta$
\DisplayProof
\qquad
\Axiom$!B, \Gamma \fCenter \Delta$
\RightLabel{$\LeftR{\land}_2$}
\UnaryInf$!A \land !B, \Gamma \fCenter \Delta$
\DisplayProof
\]
دوجے پاسے \Log{G3c} وچ اکّو قاعدہ اے:
\[
\Axiom$ !A, !B, \Gamma \fCenter \Delta$
\RightLabel{$\LeftR{\land}_3$}
\UnaryInf$ !A \land !B, \Gamma \fCenter \Delta$
\DisplayProof
\]
ہُن اسیں پُچھ سکدے آں، «کی \Log{G3c} دا قاعدہ $\LeftR{\land}_3$، \Log{G1c} وچ قبول کرن جوگ اے؟» ہاں۔ اسیں $\LeftR{\land}_3$ نال کیتا کوئی وی استنباط صرف~\Log{G1c} دے قاعدیاں نال سیدھا مُڑ بنا سکدے آں۔ $\LeftR{\land}_1$ تے $\LeftR{\land}_2$ دے نال سانوں ساختی قاعدہ~$\LeftR{\Contraction}$ وی چاہیدا اے:
\[
\Axiom$!A, !B, \Gamma \fCenter \Delta$
\RightLabel{$\LeftR{\land}_1$}
\UnaryInf$!A \land !B, !B, \Gamma \fCenter \Delta$
\RightLabel{$\LeftR{\land}_2$}
\UnaryInf$!A \land !B, !A \land !B, \Gamma \fCenter \Delta$
\RightLabel{\LeftR{\Contraction}}
\UnaryInf$!A \land !B \Gamma \fCenter \Delta$
\DisplayProof
\]
\end{ex}

اک قاعدے نال کیتے استنباط نوں ہور قاعدیاں نال مُڑ بنانا، ایہ وکھاؤن دا اک طریقہ اے کہ قاعدہ قبول کرن جوگ اے۔ ایہو اکّو طریقہ نہیں۔ کجھ قاعدے قبول کرن جوگ نیں پر ایس طرح مُڑ بنائے نہیں جا سکدے۔ جیہڑے مُڑ بنائے جا سکدے نیں، اوہناں نوں \emph{اخذ ہو سکّن والے} آکھدے آں۔

\begin{defn}
قاعدہ~$R$،
\[
\AxiomC{$S_1$}
\AxiomC{$\dots$}
\AxiomC{$S_n$}
\RightLabel{$R$}
\TrinaryInfC{$S$}
\DisplayProof
\]
کسے نظام وچ \emph{اخذ ہو سکّن والا} اے جے اک خاکہ وار !!{proof} اے، سیکوئنٹ~$S$ دا، جیہڑا اوس نظام دے قاعدیاں تے مسلّمیاں نوں، نالے~$S_1$، \dots،~$S_n$ دیاں شکلاں والے وادھو شروع دے سیکوئنٹاں نوں ورتدا اے۔
\end{defn}

!!{proof} جیہڑا \olref{ex:land-deriv} وچ دِتّا اے، وکھاؤندا اے کہ $\LeftR{\land}_3$،~\Log{G1c} وچ اخذ ہو سکّن والا قاعدہ اے۔ وجہ ایہ اے کہ اوہ !!{proof} اک خاکہ وار !!{proof} اے، $\LeftR{\land}_3$ دے نتیجے دا، $\LeftR{\land}_3$ دے مقدمیّاں نوں شروع دے سیکوئنٹ منّ کے تے صرف \Log{G1c} دے قاعدے ورت کے۔ (اوہنے \Log{G1c} دے مسلّمیاں وچوں کوئی نہیں ورتیا۔)

\begin{prop}\ollabel{prop:lor-G3-adm}
قاعدے $\LeftR{\land}$ تے $\RightR{\lor}$،~\Log{G3c} دے،~\Log{G1c} وچ اخذ ہو سکدے نیں۔
\end{prop}

\begin{proof}
  $\LeftR{\land}$ دا ثبوت \olref{ex:land-deriv} وچ دِتّا اے۔ $\RightR{\lor}$ دا ثبوت مشق لئی چھڈّیا اے۔
\end{proof}

\begin{prob}
وکھاؤ کہ قاعدہ $\RightR{\lor}$، \Log{G3c} دا،~\Log{G1c} وچ اخذ ہو سکدا اے۔
\end{prob}

\begin{prob}
منّ لو $\liff$ اک متعین !!{operator} اے۔ $!A \liff !B$، $(!A \lif !B) \land (!B \lif !A)$ دا چھوٹا روپ اے۔ وکھاؤ کہ قاعدے
\begin{enumerate}
\item \Axiom$\Gamma, !A, !B \fCenter \Delta$
\Axiom$\Gamma \fCenter \Delta, !A, !B$
\RightLabel{\LeftR{\liff}}
\BinaryInf$\Gamma \fCenter \Delta, !A \liff !B$
\DisplayProof
\item
\Axiom$!A, \Gamma \fCenter \Delta, !B$
\Axiom$!B, \Gamma \fCenter \Delta, !A$
\RightLabel{\RightR{\liff}}
\BinaryInf$\Gamma \fCenter \Delta, !A \liff !B$
\DisplayProof
\end{enumerate}
(a)~\Log{G1c} تے (b)~\Log{G3c} وچ اخذ ہو سکدے نیں۔
\end{prob}

\begin{ex}\ollabel{ex:weak-adm}
قاعدہ $\RightR{\Weakening}$، \Log{G1c} دا،
\[
\Axiom$ \Gamma \fCenter \Delta$
\RightLabel{\RightR{\Weakening}}
\UnaryInf$ \Gamma \fCenter \Delta, !A$
\DisplayProof
\]
~\Log{G3c} وچ قبول کرن جوگ اے۔ ثبوت دا خیال سدھا اے۔ منّ لو !!a{proof}~$\pi$ اے، مقدمے $\Gamma \Sequent \Delta$ دا،~\Log{G3c} وچ۔ !!{formula}~$!A$ نوں ہر سیکوئنٹ دے سجّے پاسے شامل کرو،~$\pi$ وچ۔ مسلّمے مسلّمے رہندے نیں تے درست استنباط ہجے وی درست نیں۔ آخری !!{formula}، نویں !!{proof} دا، $ \Gamma \Sequent
\Delta, !A$ اے۔

پر قاعدہ $\RightR{\Weakening}$،~\Log{G3c} وچ اخذ نہیں ہو سکدا۔ منّ لو ایہ ہو سکدا اے، یعنی کوئی !!a{proof} لبھیا جا سکدا اے $ \Gamma \Sequent
\Delta, !A$ دا، صرف~\Log{G3c} دے مسلّمے تے قاعدے ورت کے تے شاید وادھو شروع دا سیکوئنٹ~$\Gamma \Sequent \Delta$ وی ورت کے۔ جے ایہ !!{proof} پورا خاکہ وار اے، جیویں $\RightR{\Weakening}$ دے اخذ ہو سکّن لئی ضروری ہووے گا، تاں اسیں اوہنوں !!a{proof} وچ بدل سکدے آں~$\Sequent !A$ دا، $ \Gamma$ تے~$\Delta$ ہٹا کے۔ ایس طرح وادھو شروع دے سیکوئنٹ~$\Gamma \Sequent \Delta$ خالی سیکوئنٹ~$\quad \Sequent \quad$ بن جان گے۔ خالی سیکوئنٹ وچ کوئی !!{formula}s نہیں، پر~\Log{G3c} دے سارے قاعدیاں دے مقدمیّاں وچ گھٹّو گھٹ اک فعال !!{formula} چاہیدا اے۔ سو ایس خاکہ وار !!{proof} وچ خالی سیکوئنٹ کسے استنباط دا مقدمہ نہیں ہو سکدا۔ ایس لئی !!{proof} صرف اوہدوں خاکہ وار ہو سکدا اے جے ایہ وادھو سیکوئنٹ $ \Gamma \Sequent \Delta$ نوں شروع دے سیکوئنٹ دے طور اُتے اصل وچ نہ ورتے۔ دوجے لفظاں وچ، ایہ $ \Gamma \Sequent
\Delta, !A$ دا خاکہ وار اخذ ہووے گا جیہڑا صرف~\Log{G3c} دے مسلّمے تے قاعدے ورتدا اے۔ اسیں وکھا چکے ہوواں گے کہ ایس شکل دے ہر سیکوئنٹ دا !!a{proof} اے~\Log{G3c} وچ۔ پر ایہ ممکن نہیں، کیونکہ \Log{G3c} صحت مند اے تے ایس لئی صرف اوہ سیکوئنٹ !!{prove} کر سکدا اے جیہڑے ہر !!{structure} وچ سچے نیں۔ $ \Gamma = \Delta
= \emptyset$ تے $!A \ident \lfalse$ لے کے، \Log{G3c} نوں مثلاً سیکوئنٹ $ \quad \Sequent \lfalse$ !!{prove} کرنا پئے گا، جیہڑا ایہ نہیں کردا۔
\end{ex}

آؤ کمزور کرن والا قاعدہ~\Log{G3c} وچ قبول کرن جوگ ہون دا اک سوہنا استقرائی ثبوت دیئیے۔ اصل وچ ذرا مضبوط نتیجہ وی ٹھیک اے، جیویں سہج دلیل وکھاؤندی اے۔ $!A$ نوں~$\pi$ دے ہر سیکوئنٹ دے سجّے پاسے شامل کرن نال اسیں !!{proof} دی بناوٹ نہیں بدل دے، خاص کر اوہدی اُچائی نہیں ودھاؤندے۔ کیونکہ ایہ گل اہم اے، ایہدی اپنی تعریف ہونی چاہیدی اے۔

\begin{defn}
قاعدہ~$R$،
\[
\AxiomC{$S_1$}
\AxiomC{$\dots$}
\AxiomC{$S_n$}
\RightLabel{$R$}
\TrinaryInfC{$S$}
\DisplayProof
\]
کسے نظام وچ \emph{!!{height} قائم رکھ کے قبول کرن جوگ} (یا \emph{!!{hp}-قبول کرن جوگ}) اے جے جدوں وی $\Proves[h_i] S_i$، $i = 1$، \dots،~$n$ لئی، تاں $\Proves_m S$، جتھے $m = \max(h_1, \dots,
h_n)$۔
\end{defn}

ویکھو کہ قاعدہ !!{height} قائم رکھ کے قبول کرن جوگ ہون لئی اینا کافی اے کہ جدوں وی مقدمیّاں~$S_i$ دے !!{proof}s~$\pi_i$ ہون، !!{height}~$h_i = \pheight{\pi_i}$ والے، تاں نتیجے دا !!a{proof}~$\pi$ ہووے جیہدی !!{height}~$\pheight{\pi}$،~$h_i$ دے سبھ توں وڈّے عدد توں ودھ نہ ہووے۔ خاص کر جے اکّو مقدمہ~$S_1$ اے، تاں اینا کافی اے کہ $\pheight{\pi} \le \pheight{\pi_1}$۔ $\RightR{\Weakening}$ تے $\LeftR{\Weakening}$ دی صورت وچ اصل وچ $\pheight{\pi} = \pheight{\pi_1}$ اے، پر ایہ ضروری نہیں۔

\begin{prop}\ollabel{prop:weak-G3c-adm}
قاعدے $\RightR{\Weakening}$ تے $\LeftR{\Weakening}$، \Log{G1c} دے،~\Log{G3c} وچ !!{height} قائم رکھ کے قبول کرن جوگ نیں۔
\end{prop}

\begin{proof}
  سانوں وکھاؤنا پینا اے کہ جے $\Log{G3c} \Proves \Gamma \Sequent \Delta$، !!{proof}~$\pi$ نال جیہدی !!{height}~$\pheight{\pi} = n$ اے، تاں اک \Log{G3c}-!!{proof} $ \pi'$ اے $ \Gamma \Sequent \Delta, !A$ دا جیہدے لئی $\pheight{\pi'} = n$۔ اسیں~$n$ اُتے استقرا نال ایہ ثابت کردے آں۔ جے $n=0$، تاں $\pi$ صرف مسلّمہ $ \Gamma \Sequent \Delta$ اے (یعنی کوئی جوہری $!C$، $ \Gamma$ تے~$\Delta$ دوناں وچ آؤندا اے)، تے $ \Gamma
  \Sequent \Delta, !A$ وی مسلّمہ اے۔ جے $n = \pheight{\pi} >0$، تاں !!{proof}~$\pi$ دا آخری استنباط اے۔ ورتے گئے قاعدے دے حساب نال وکھ وکھ صورتاں ویکھدے آں۔ اسیں صرف اک مثال دیندے آں۔ منّ لو آخری استنباط $ \RightR{\land}$ سی۔ تاں $ \Delta = \Delta', !B \land !C$ تے آخری استنباط دے سدھے ذیلی !!{proof}s $ \pi_1$ تے~$\pi_2$ دے آخری سیکوئنٹ بالترتیب $ \Gamma \Sequent \Delta', !B$ تے $ \Gamma \Sequent
  \Delta', !C$ نیں۔ دوجے لفظاں وچ، $\pi$ ایویں دِسدا اے:
  \[
  \AxiomC{}
  \RightLabel{$\pi_1$}
  \Deduce$\Gamma \fCenter \Delta', !B$
  \AxiomC{}
  \RightLabel{$\pi_2$}
  \Deduce$\Gamma \fCenter \Delta', !C$
  \RightLabel{\RightR{\land}}
  \BinaryInf$\Gamma \fCenter \Delta', !B \land !C$
  \DisplayProof
  \]
  ایہ وی اے کہ $n =
  \max(\pheight{\pi_1}, \pheight{\pi_2}) + 1$۔ سو $\pheight{\pi_1} <
  n$ تے $\pheight{\pi_2} < n$، تے استقرا دا مفروضہ لگدا اے۔ !!{proof}s $ \pi_1'$ تے~$\pi_2'$، $ \Gamma \Sequent
  \Delta', !B, !A$ تے $ \Gamma \Sequent \Delta', !C, !A$ دے، موجود نیں۔ $\RightR{\land}$ قاعدہ ورت کے اسیں !!a{proof}~$\pi'$ لیندے آں $ \Gamma
  \Sequent \Delta', !B \land !C, !A$ دا:
  \[
  \AxiomC{}
  \RightLabel{$\pi_1'$}
  \Deduce$\Gamma \fCenter \Delta', !B, !A$
  \AxiomC{}
  \RightLabel{$\pi_2'$}
  \Deduce$\Gamma \fCenter \Delta', !C, !A$
  \RightLabel{\RightR{\land}}
  \BinaryInf$\Gamma \fCenter \Delta', !B \land !C, !A$
  \DisplayProof
  \]
  ساڈے کول $ \pheight{\pi'} =
  \max(\pheight{\pi_1'}, \pheight{\pi_2'}) + 1 = \max(\pheight{\pi_1},
  \pheight{\pi_2}) + 1 = \pheight{\pi}$ اے۔
\end{proof}

\olref{prop:weak-G3c-adm} بہت کم آؤندا اے۔ اسیں ایہنوں بار بار لاؤن لئی اک لکھت ورتاں گے۔ جے $\pi$، !!a{proof} اے~$\Gamma \Sequent \Delta$ دا، تاں $\pi[\Pi \Sequent \Lambda]$ اوہ نتیجہ اے جیہڑا کھبّے پاسے کمزور کرن دی قبولیت، سارے !!{formula}s~$\Pi$ وچلے نال، تے سجّے پاسے کمزور کرن دی قبولیت، سارے !!{formula}s~$\Lambda$ وچلے نال، لاؤن نال ملدا اے۔ ایہ !!a{proof} اے $ \Gamma, \Pi
\Sequent \Delta, \Lambda$ دا۔
\endgroup
% /OLP-0698
% OLP-0701 content/proof-theory/sequent-calculus/invertibility.tex
\begingroup
\def\olfilename{invertibility}
\expandafter\def\csname ol@sectioncs\endcsname{section}

\olfileid{pt}{seq}{inv}

\olsection{قاعدیاں نوں الٹا ورت سکنا}

\begin{defn}
قاعدہ~$R$،
\[
\AxiomC{$S_1$}
\AxiomC{$\dots$}
\AxiomC{$S_n$}
\RightLabel{$R$}
\TrinaryInfC{$S$}
\DisplayProof
\]
کسے نظام وچ \emph{الٹا ورتن جوگا} اے جے جدوں وی $\Proves S$، تاں $\Proves[h_i] S_i$، $i = 1$، \dots،~$n$ لئی۔ ایہ اوس نظام وچ \emph{!!{height} قائم رکھ کے الٹا ورتن جوگا} (یا \emph{!!{hp}-الٹ جوگا}) اے جے جدوں وی $\Proves_n S$، تاں $\Proves[n] S_i$، $i =
1$، \dots،~$n$ لئی۔
\end{defn}

\begin{lem}\ollabel{lem:G3c-invert}
~\Log{G3c} وچ $\lnot$، $\land$، $\lor$ تے $\lif$ دے قاعدے !!{height} قائم رکھ کے الٹے ورتے جا سکدے نیں، یعنی:
\begin{enumerate}
  \item \RightR{\lnot}: جے $\Log{G3c} \Proves[n] \Gamma \Sequent
  \Delta, \lnot !A$، تاں $\Log{G3c} \Proves[n] !A, \Gamma \Sequent
  \Delta$۔
  \item \LeftR{\lnot}: جے $\Log{G3c} \Proves[n] \Gamma, \lnot !A
  \Sequent \Delta$، تاں $\Log{G3c} \Proves[n] \Gamma \Sequent \Delta,
  !A$۔
  \item \RightR{\land}: جے $\Log{G3c} \Proves[n] \Gamma \Sequent
  \Delta, !A \land !B$، تاں $\Log{G3c} \Proves[n] \Gamma \Sequent
  \Delta, !A$ تے $\Log{G3c} \Proves[n] \Gamma \Sequent
  \Delta, !B$۔
  \item \LeftR{\land}: جے $\Log{G3c} \Proves[n] \Gamma, !A \land !B
  \Sequent \Delta$، تاں $\Log{G3c} \Proves[n] !A, !B, \Gamma \Sequent
  \Delta$۔
  \item \RightR{\lor}: جے $\Log{G3c} \Proves[n] \Gamma \Sequent
  \Delta, !A \lor !B$، تاں $\Log{G3c} \Proves[n] \Gamma \Sequent
  \Delta, !A, !B$۔
  \item \LeftR{\lor}: جے $\Log{G3c} \Proves[n] !A \lor !B, \Gamma \Sequent
  \Delta$، تاں $\Log{G3c} \Proves[n] !A, \Gamma \Sequent
  \Delta$ تے $\Log{G3c} \Proves[n] !B, \Gamma \Sequent
  \Delta$۔
  \item \RightR{\lif}: جے $\Log{G3c} \Proves[n] \Gamma \Sequent
  \Delta, !A \lif !B$، تاں $\Log{G3c} \Proves[n] !A, \Gamma \Sequent
  \Delta, !B$۔
  \item \LeftR{\lif}: جے $\Log{G3c} \Proves[n] !A \lif !B, \Gamma
  \Sequent \Delta$، تاں $\Log{G3c} \Proves[n] \Gamma \Sequent \Delta,
  !A$ تے $\Log{G3c} \Proves[n] !B, \Gamma \Sequent \Delta$۔
\end{enumerate}
\end{lem}

\begin{proof}
سانوں وکھاؤنا پینا اے کہ کسے وی قاعدے~$R$ لئی، جیہڑا~\Log{G3c} دا اے، جے !!a{proof}~$\pi$ اے نتیجے~$S$ دا، جیہڑا~$R$ دا اے، تاں !!{proof}s~$\pi_i$ وی نیں مقدمیّاں $S_i$ دے، جیہڑے~$R$ دے نیں، تے $\pheight{\pi_i}
\le \pheight{\pi}$۔ قاعدہ \LeftR{\land} ویکھو۔ منّ لو !!a{proof}~$\pi$ اے جیہدا آخری سیکوئنٹ $!A \land !B, \Gamma
\Sequent \Delta$ دی شکل دا اے۔ سانوں وکھاؤنا پینا اے کہ !!a{proof} $\pi'$ اے $!A, !B, \Gamma \Sequent \Delta$ دا، تے $\pheight{\pi'} \le
\pheight{\pi}$۔ اسیں ایہ $n = \pheight{\pi}$ اُتے استقرا نال کردے آں۔

جے $n = 0$، تاں $\pi$ بس اک مسلّمہ اے، $!A \land !B, \Gamma
\Sequent \Delta$۔ \Log{G3c} وچ مسلّمے $!C, \Gamma
\Sequent \Delta, !C$ دی شکل دے نیں، جتھے $!C$ جوہری اے۔ دوجے لفظاں وچ $!C$، $!A \land !B$ نہیں ہو سکدا۔ جے $!A \land !B, \Gamma \Sequent \Delta$ مسلّمہ اے، تاں کوئی جوہری $!C$، !!a{element} ضرور اے $ \Gamma$ تے~$\Delta$ دوناں دا۔ پر فیر $!A, !B, \Gamma \Sequent \Delta$ وی مسلّمہ اے تے سانوں ساڈا !!{proof}~$\pi'$ مل گیا اے (جیہدی !!{height} وی~$0$ اے)۔

جے $n>0$، تاں ایہ کسے استنباط اُتے مُکدا اے۔ سانوں دعویٰ~$n$ لئی ثابت کرنا پینا اے، پر اسیں منّ سکدے آں کہ ایہ ہر !!{proof} لئی ٹھیک اے جیہدی اُچائی~$<n$ اے، بھاویں سیاق وکھرا ہووے۔ دوجے لفظاں وچ، ہر $k<n$ تے ہر $ \Pi$ تے $ \Lambda$ لئی، جے $!A \land !B, \Pi \Sequent \Lambda$ دا !!a{proof} اے !!{height}~$k$ والا، تاں !!a{proof} اے !!{height}~$\le k$ والا، $!A, !B, \Pi \Sequent \Lambda$ دا۔

دو صورتاں نیں۔ کھبّے پاسے $!A \land
!B$ یا مُکھ اے یا نہیں۔ جے ایہ مُکھ اے، تاں آخری استنباط \LeftR{\land} اے تے سدھا ذیلی !!{proof}~$\pi_1$، $!A, !B, \Gamma \Sequent \Delta$ اُتے مُکدا اے۔ ایس صورت وچ نتیجہ ٹھیک اے، کیونکہ $\pheight{\pi_1} < \pheight{\pi}$۔

جے $!A \land !B$ مُکھ نہیں، تاں کوئی ہور !!{formula} مُکھ اے تے $!A
\land !B$ آخری استنباط دے سیاق دا حصہ اے۔ فیر استقرا دا مفروضہ آخری قاعدے~$R$ دے سدھے ذیلی !!{proof}s اُتے لگدا اے۔ ایہدے نال !!a{proof}~$\pi_1'$ یا دو !!{proof}s~$\pi_1'$ تے~$\pi_2'$ ملدے نیں، جیہناں دی !!{height}، سدھے ذیلی !!{proof}s~$\pi$ دے، دیاں اُچائیاں توں ودھ نہیں۔ $ \pi_1'$ تے $ \pi_2'$ دے آخری سیکوئنٹاں دے کھبّے سیاق وچ $!A, !B$ اے، $!A \land !B$ دی تھاں۔ قاعدہ~$R$ لا کے ثبوت~$\pi'$ لو۔ مثال لئی منّ لو $\pi$،~\LeftR{\lif} اُتے مُکدا اے:
\[
\AxiomC{}
\RightLabel{$\pi_1$}
\Deduce$!A \land !B, \Gamma' \fCenter \Delta, !C$
\AxiomC{}
\RightLabel{$\pi_2$}
\Deduce$!A \land !B, \Gamma', !D \fCenter \Delta$
\RightLabel{\LeftR{\lif}}
\BinaryInf$!A \land !B, \Gamma', !C \lif !D \fCenter \Delta$
\DisplayProof
\]
ایتھے کھبّے پاسے $!C \lif !D$ مُکھ اے، تے کھبّا سیاق $!A \land !B, \Gamma'$ اے (یعنی $\Gamma = \Gamma', !C \lif !D$)۔ استقرا دے مفروضے نال !!{proof}s $ \pi_1'$ تے $ \pi_2'$، بالترتیب $!A \land !B, \Gamma' \Sequent \Delta, !C$ تے $!A \land !B, \Gamma',
!D \Sequent \Delta$ دے، ملدے نیں۔ ایہ دوویں سیکوئنٹ فیر \LeftR{\lif} قاعدے دے مقدمے بن سکدے نیں،~$\pi'$ لین لئی:
\[
\AxiomC{}
\RightLabel{$\pi_1'$}
\Deduce$!A, !B, \Gamma' \fCenter \Delta, !C$
\AxiomC{}
\RightLabel{$\pi_2'$}
\Deduce$!A, !B, \Gamma', !D \fCenter \Delta$
\RightLabel{\LeftR{\lif}}
\BinaryInf$!A, !B, \Gamma', !C \lif !D \fCenter \Delta$
\DisplayProof
\]
سانوں ایہ ملدا اے:
\begin{align*}
\pheight{\pi} & = \max(\pheight{\pi_1}, \pheight{\pi_2}) + 1\\
\pheight{\pi'} & = \max(\pheight{\pi_1'}, \pheight{\pi_2'}) + 1
\intertext{by definition, and}
\pheight{\pi_1'} & \le \pheight{\pi_1}\\
\pheight{\pi_2'} & \le \pheight{\pi_2}
\intertext{by inductive hypothesis. So,}
\pheight{\pi'} & \le \pheight{\pi}.
\end{align*}
\end{proof}

\begin{prob}
\olref{lem:G3c-invert} دا ثبوت \RightR{\lif} لئی دیو۔
\end{prob}

\begin{lem}\ollabel{lem:invert-quant}
قاعدے \RightR{\lforall} تے \LeftR{\lexists} ہِٹھلے مطلب وچ !!{height} قائم رکھ کے الٹے ورتے جا سکدے نیں:
\begin{enumerate}
  \item جے $\Log{G3c} \Proves[n] \Gamma \Sequent \Delta,
  \lforall[x][!A(x)]$، تاں $\Log{G3c} \Proves[n] \Gamma \Sequent
  \Delta, !A(c)$، تے
  \item جے $\Log{G3c} \Proves[n] \lexists[x][!A(x)], \Gamma \Sequent
  \Delta$، تاں $\Log{G3c} \Proves[n] !A(c), \Gamma \Sequent \Delta$،
\end{enumerate}
ہر $c$ لئی جیہڑا $ \Gamma$، $ \Delta$ تے $ \lforall[x][!A(x)]$ یا $ \lexists[x][!A(x)]$ وچ، بالترتیب، نہیں آؤندا۔
\end{lem}

\begin{proof}
  اسیں (1) وکھاؤنے آں تے (2) مشق لئی چھڈّ دیندے آں۔ دلیل \olref{lem:G3c-invert} دے ثبوت ورگی اے۔ استقرائی قدم وچ فیر دو صورتاں نیں۔ جیہڑی صورت وچ $\lforall[x][!A(x)]$ مُکھ اے، اوہ فیر چھیتی ہو جاندی اے۔ آخری $ \RightR{\lforall}$ استنباط دا مقدمہ مطلوبہ شکل $ \Gamma
  \Sequent \Delta, !A(a)$ دا اے، تے ایس اُتے مُکّن والے سدھے ذیلی !!{proof}~$\pi_1$ لئی $\pheight{\pi_1} < \pheight{\pi}$ اے۔ نالے خاص متغیر دی شرط توں سانوں پتا لگدا اے کہ $a$، $ \Gamma$، $ \Delta$ یا $ \lforall[x][!A(x)]$ وچ نہیں آؤندا۔ کیونکہ اسیں منّ سکدے آں کہ $\pi_1$ باقاعدہ اے، ایس لئی $ \Subst{\pi_1}{c}{a}$ اوہو !!a{proof} اے، اوہو !!{height} والا، $ \Gamma \Sequent \Delta, !A(c)$ دا، \olref[qua]{lem:subst-regular} نال۔

  دوجی صورت وچ $\lforall[x][!A(x)]$ آخری استنباط وچ مُکھ نہیں، کوئی ہور فارمولا مُکھ اے۔ اسیں فیر اک صورت مثال دے طور اُتے ویکھدے آں۔ منّ لو آخری استنباط~\RightR{\land} اے۔ کوئی !!{formula} $\Delta$ وچ~$!B \land !C$ دی شکل دا اے تے مُکھ اے۔ !!{proof}~$\pi$ ایویں مُکدا اے:
  \[
\AxiomC{}
\RightLabel{$\pi_1$}
\Deduce$\Gamma \fCenter \Delta', \lforall[x][!A(x)], !B$
\AxiomC{}
\RightLabel{$\pi_2$}
\Deduce$\Gamma \fCenter \Delta', \lforall[x][!A(x)], !C$
\RightLabel{\RightR{\land}}
\BinaryInf$\Gamma \fCenter \Delta', \lforall[x][!A(x)], !B \land !C$
\DisplayProof
\]
جتھے $\Delta = \Delta', !B \land !C$۔ منّ لو $c$ اک !!a{constant} اے جیہڑا $ \Delta$ وچ نہیں آؤندا۔ فیر صاف اے کہ ایہ $ \Delta', !B$ وچ وی نہیں آؤندا تے $ \Delta', !C$ وچ وی نہیں۔ سو استقرا دا مفروضہ $ \pi_1$ تے~$\pi_2$ اُتے لگدا اے۔ ساڈے کول !!{proof}s~$\pi_1'$ تے $ \pi_2'$ نیں، تے $\pheight{\pi_1'} \le \pheight{\pi_1}$ تے $\pheight{\pi_2'} \le
\pheight{\pi_2}$۔ مطلوبہ !!{proof}~$\pi'$، $ \Gamma \Sequent
\Delta, !A(c)$ دا، ایہ اے:
\[
\AxiomC{}
\RightLabel{$\pi_1'$}
\Deduce$\Gamma \fCenter \Delta', !A(c), !B$
\AxiomC{}
\RightLabel{$\pi_2'$}
\Deduce$\Gamma \fCenter \Delta', !A(c), !C$
\RightLabel{\RightR{\land}}
\BinaryInf$\Gamma \fCenter \Delta', !A(c), !B \land !C$
\DisplayProof
\]
جیہدے لئی $\pheight{\pi'} = \max(\pheight{\pi_1'}, \pheight{\pi_2'})+1 \le \pheight{\pi}$۔
\end{proof}

\Olref{lem:G3c-invert} قطع ہٹاؤن دے قضیے دے ثبوت وچ اہم کم کردا اے۔ ایہ ہِٹھلی گل وکھاؤن لئی وی ورتیا جا سکدا اے۔

\begin{prop}\ollabel{prop:G3c-cont-adm}
سکڑاؤ دے قاعدے \LeftR{\Contraction} تے \RightR{\Contraction}، جیہڑے~\Log{G1c} دے نیں،~\Log{G3c} وچ !!{height} قائم رکھ کے قبول کرن جوگ نیں۔
\end{prop}

\begin{proof}
اسیں \RightR{\Contraction} تے \LeftR{\Contraction} لئی دلیل اکٹھی دیندے آں۔ منّ لو $\pi$، !!a{proof} اے $\Gamma \Sequent \Delta, !A, !A$ دا یا $!A, !A, \Gamma \Sequent
\Delta$ دا،~\Log{G3c} وچ۔ سانوں وکھاؤنا پینا اے کہ اک \Log{G3c}-!!{proof}~$\pi'$ وی اے، بالترتیب $ \Gamma \Sequent \Delta, !A$ دا یا $!A,
\Gamma \Sequent \Delta$ دا، تے $\pheight{\pi'} \le
\pheight{\pi}$۔ اسیں ایہ $n = \pheight{\pi}$ اُتے استقرا نال کردے آں۔

جے $n=0$، تاں $\pi$ بس مسلّمہ اے۔ پر جے $\Gamma \Sequent \Delta, !A,
!A$،~\Log{G3c} دا مسلّمہ اے، تاں $\Gamma \Sequent \Delta, !A$ وی اے۔ یا کوئی جوہری~$!C$، !!a{element} اے $ \Gamma$ تے~$\Delta$ دوناں دا، یا $!A$ جوہری اے تے !!a{element} اے~$\Gamma$ دا۔ اوہی دلیل لگدی اے جے آخری سیکوئنٹ $!A, !A, \Gamma \Sequent
\Delta$ اے۔

جے $n>0$، تاں $\pi$ کسے قاعدے~$R$ اُتے مُکدا اے۔ جے $!A$ ایس آخری استنباط وچ مُکھ نہیں، تاں اسیں !!a{proof}~$\pi'$ اوہو جیہا لے سکدے آں جیویں \olref{lem:G3c-invert} دے ثبوت وچ۔ استقرا دا مفروضہ~$\pi$ دے سدھے ذیلی !!{proof}s اُتے لاؤ تے اوہی قاعدہ~$R$ حاصل ہوئے !!{proof} یا ثبوتاں اُتے لاؤ۔ استقرا دا مفروضہ ایہ اے کہ جے !!a{proof} اے !!{height}~$k<n$ والا، $ \Pi \Sequent \Lambda, !D,
!D$ دا یا $!D, !D, \Pi \Sequent \Lambda$ دا، تاں !!a{proof} اے !!{height}~$\le k$ والا، بالترتیب $ \Pi \Sequent \Lambda, !D$ دا یا $!D, \Pi \Sequent
\Lambda$ دا۔ (ویکھو کہ استقرا دا مفروضہ اوہدوں وی لگدا اے جے سیاق یا سکیڑیا جان والا !!{formula}، $ \Gamma$، $ \Delta$ یا~$!A$ توں وکھرے نیں!)

مثال لئی منّ لو $R$، $\RightR{\land}$ اے تے $\pi$ ایویں مُکدا اے:
\[
\AxiomC{}
\RightLabel{$\pi_1$}
\Deduce$\Gamma \fCenter \Delta', !B, !A, !A$
\AxiomC{}
\RightLabel{$\pi_2$}
\Deduce$\Gamma \fCenter \Delta', !C, !A, !A$
\RightLabel{\RightR{\land}}
\BinaryInf$\Gamma \fCenter \Delta', !B \land !C, !A, !A$
\DisplayProof
\]
جتھے $\Delta = \Delta', !B \land !C$۔ استقرا دے مفروضے نال !!{proof}s $ \pi_1'$ تے $ \pi_2'$، بالترتیب $ \Gamma \fCenter \Delta', !B, !A$ تے $ \Gamma \fCenter \Delta', !C, !A$ دے، ملدے نیں، تے $\pheight{\pi_1'} \le  \pheight{\pi_1}$ تے $\pheight{\pi_2'} \le
\pheight{\pi_2}$۔ !!{proof}~$\pi'$ ایہ اے:
\[
\AxiomC{}
\RightLabel{$\pi_1'$}
\Deduce$\Gamma \fCenter \Delta', !B, !A$
\AxiomC{}
\RightLabel{$\pi_2'$}
\Deduce$\Gamma \fCenter \Delta', !C, !A$
\RightLabel{\RightR{\land}}
\BinaryInf$\Gamma \fCenter \Delta', !B \land !C, !A$
\DisplayProof
\]

جے $!A$ آخری استنباط وچ مُکھ اے، تاں اسیں اُلٹ لاؤن دا لیما (\olref{lem:G3c-invert})، سدھے ذیلی !!{proof}s~$\pi$ دے اُتے لاؤنے آں، فیر استقرا دا مفروضہ، تے فیر قاعدہ~$R$۔ مثال لئی منّ لو $\pi$، $ \Gamma \Sequent \Delta, !A, !A$ ثابت کردا اے، $!A
\ident (!B \land !C)$ تے آخری استنباط وچ مُکھ اے۔ فیر $\pi$ ایویں دِسدا اے:
\[
\AxiomC{}
\RightLabel{$\pi_1$}
\Deduce$\Gamma \fCenter \Delta, !B \land !C, !B$
\AxiomC{}
\RightLabel{$\pi_2$}
\Deduce$\Gamma \fCenter \Delta, !B \land !C, !C$
\RightLabel{\RightR{\land}}
\BinaryInf$\Gamma \fCenter \Delta, !B \land !C, !B \land !C$
\DisplayProof
\]
\olref{lem:G3c-invert} نوں $ \pi_1$ اُتے لاؤ تاں !!a{proof}~$\pi_1'$ ملدا اے $ \Gamma \Sequent \Delta, !B, !B$ دا۔ (ایہ !!a{proof} وی دیندا اے $ \Gamma \Sequent \Delta, !C, !B$ دا، پر سانوں اوہدی لوڑ نہیں۔) ایویں ای \olref{lem:G3c-invert} نوں~$\pi_2$ اُتے لا کے اسیں !!a{proof}~$\pi_2'$ لیندے آں $ \Gamma \Sequent \Delta, !C, !C$ دا۔ کیونکہ \RightR{\land} نوں الٹا لاؤنا !!{height} قائم رکھدا اے، $\pheight{\pi_1'} \le  \pheight{\pi_1} < \pheight{\pi}$ تے $\pheight{\pi_2'} \le  \pheight{\pi_2} < \pheight{\pi}$۔ سو استقرا دا مفروضہ $ \pi_1'$ تے $ \pi_2'$ اُتے لگدا اے۔ !!{proof}s~$\pi_1''$ تے $ \pi_2''$، بالترتیب $ \Gamma \Sequent \Delta, !B$ تے $ \Gamma \Sequent \Delta, !C$ دے، نیں، تے $\pheight{\pi_1''}
\le \pheight{\pi_1'} \le \pheight{\pi_1}$ تے $\pheight{\pi_2''} \le
\pheight{\pi_2'} \le \pheight{\pi_2}$۔ فیر !!{proof}~$\pi'$ ایہ اے:
\[
\AxiomC{}
\RightLabel{$\pi_1''$}
\Deduce$\Gamma \fCenter \Delta, !B$
\AxiomC{}
\RightLabel{$\pi_2''$}
\Deduce$\Gamma \fCenter \Delta, !C$
\RightLabel{\RightR{\land}}
\BinaryInf$\Gamma \fCenter \Delta, !B \land !C$
\DisplayProof
\]
جیہدے لئی $\pheight{\pi'} \le \pheight{\pi}$۔

جیہڑیاں صورتاں وچ !!{formula} مُکھ اے، اوہناں وچ مسوّراں دے قاعدیاں لئی خاص دھیان رکھنا پیندا اے۔

منّ لو $!A \ident \lforall[x][!B(x)]$ سجّے پاسے مُکھ اے۔ فیر $\pi$ ایویں مُکدا اے:
\[
\AxiomC{}
\RightLabel{$\pi_1$}
\Deduce$\Gamma \fCenter \Delta, \lforall[x][!B(x)], !B(a)$
\RightLabel{\RightR{\lexists}}
\UnaryInf$\Gamma \fCenter \Delta, \lforall[x][!B(x)], \lforall[x][!B(x)]$
\DisplayProof
\]
\olref{lem:invert-quant} نال !!a{proof}~$\pi_1'$ اے، !!{height} $ \le \pheight{\pi_1}$ والا، $ \Gamma \fCenter \Delta, !B(c),
!B(a)$ دا، ہر~$c$ لئی جیہڑا $ \Gamma \fCenter \Delta,
\lforall[x][!B(x)], !B(a)$ وچ نہیں آؤندا۔ اسیں منّ سکدے آں کہ $\pi_1'$ باقاعدہ اے، سو \olref[qua]{lem:subst-regular} نال !!{proof} $ \Subst{\pi_1'}{a}{c}$، !!a{proof} اے $ \Gamma \fCenter \Delta, !B(a),
!B(a)$ دا، اوہ وی !!{height} $ \le \pheight{\pi_1}$ والا۔ ہُن استقرا دا مفروضہ لگدا اے تے سانوں !!a{proof}~$\pi_1''$ دیندا اے $ \Gamma
\fCenter \Delta, !B(a)$ دا۔\RightR{\lforall} قاعدہ لا کے~$\pi'$ لو:
\[
\AxiomC{}
\RightLabel{$\pi_1''$}
\Deduce$\Gamma \fCenter \Delta, !B(a)$
\RightLabel{\RightR{\lforall}}
\UnaryInf$\Gamma \fCenter \Delta, \lforall[x][!B(x)]$
\DisplayProof
\]
جیہدے لئی $\pheight{\pi'} \le \pheight{\pi}$۔

ہُن منّ لو $!A \ident \lexists[x][!B(x)]$ تے سجّے پاسے مُکھ اے۔ فیر $\pi$ ایویں دِسدا اے:
\[
\AxiomC{}
\RightLabel{$\pi_1$}
\Deduce$\Gamma \fCenter \Delta, \lexists[x][!B(x)], \lexists[x][!B(x)], !B(t)$
\RightLabel{\RightR{\lexists}}
\UnaryInf$\Gamma \fCenter \Delta, \lexists[x][!B(x)], \lexists[x][!B(x)]$
\DisplayProof
\]
استقرا دا مفروضہ~$\pi_1$ اُتے لگدا اے تے !!a{proof}s~$\pi_1'$ اے $ \Gamma \Sequent \Delta, \lexists[x][!B(x)],
!B$ دا۔ (ویکھو کہ ایتھے اُلٹ لاؤن دے لیمے دی لوڑ نہیں!) فیر !!{proof}~$\pi'$ ایہ اے:
\[
\AxiomC{}
\RightLabel{$\pi_1'$}
\Deduce$\Gamma \fCenter \Delta,  \lexists[x][!B(x)], !B$
\RightLabel{\RightR{\lexists}}
\UnaryInf$\Gamma \fCenter \Delta, \lexists[x][!B(x)]$
\DisplayProof
\]
جیہدے لئی $\pheight{\pi'} \le \pheight{\pi}$۔
\end{proof}

\begin{prob}
\olref{prop:G3c-cont-adm} دے ثبوت دا خاکہ \LeftR{\Contraction} لئی دیو۔ $!A \ident (!B \land !C)$ تے $!A \ident (!B
\lif !C)$ والیاں صورتاں بیان کرو۔
\end{prob}

\begin{prob}
\olref{prop:G3c-cont-adm} دے ثبوت دا خاکہ \LeftR{\Contraction} لئی باقی مسوّر والیاں صورتاں وچ دیو، یعنی جدوں $!A \ident
\lforall[x][!B(x)]$ تے $!A \ident \lexists[x][!B(x)]$۔
\end{prob}
\endgroup
% /OLP-0701
% OLP-0713 content/proof-theory/sequent-calculus/translations.tex
\begingroup
\def\olfilename{translations}
\expandafter\def\csname ol@sectioncs\endcsname{section}

\olfileid{pt}{seq}{tr}

\olsection{\Log{G1c} تے \Log{G3c} وچکار ترجمہ}

اسیں دَسّیا سی کہ \Log{G1c} تے \Log{G3c} دونے کلاسیکی منطق لئی صحت مند تے کامل نیں، یعنی ہر اک عین اوہ سارے سیکوئنٹ ثابت کردا اے جیہڑے ہر !!{structure} وچ سچے نیں۔ خاص کر، اوہ اوہی سیکوئنٹ ثابت کردے نیں۔ ہُن اسیں ایہ گل صحت مندی تے تمامیت دے قضیاں دا چکر پائے بناں، صرف ثبوت دے نظریے نال ثابت کر سکدے آں۔ ایہو جیہا ثبوت \emph{بدلّن دے طریقے} دیندا اے، !!{proof}s، جیہڑے \Log{G1c} دے نیں، نوں !!{proof}s، جیہڑے \Log{G3c} دے نیں، وچ تے الٹ پاسے بدلّن لئی۔

\begin{prop}
جے $\Log{G1c} \Proves S$، تاں $\Log{G3c} \Proves S$۔
\end{prop}

\begin{proof}
  اسیں وکھاؤنے آں کہ جے $\pi$، !!a{proof} اے \Log{G1c} وچ، تاں !!a{proof}~$\pi'$ اے~$S$ دا \Log{G3c} وچ۔ اسیں ایہ $n
  = \pheight{\pi}$ اُتے استقرا نال کردے آں۔

  جے $n = 0$، تاں $\pi$ مسلّمہ اے۔ مسلّمہ $\lfalse \Sequent
  \quad$،~\Log{G3c} وچ وی مسلّمہ اے (سیاق $ \Gamma,
  \Delta$ خالی لو)۔ \Log{G3c} وچ مسلّمے $!A \Sequent !A$ ورتے جا سکدے نیں جیہناں وچ~$!A$ کوئی وی ہو سکدا اے، صرف جوہری نہیں۔ پر اسیں \olref[ex]{prop:G3c-axioms} وچ ثابت کیتا سی کہ ایہو جیہے سارے سیکوئنٹ !!{provable} نیں~\Log{G3c} وچ۔

  جے $n>0$، تاں $\pi$ قاعدے~$R$ اُتے مُکدا اے۔ استقرا دا مفروضہ پکا کردا اے کہ اوس قاعدے دے مقدمیّاں دے !!{proof}s~\Log{G3c} وچ نیں، یعنی سدھے ذیلی !!{proof}s دے ترجمے~\Log{G3c} وچ نیں۔ جے قاعدہ~$R$، \Log{G3c} وچ وی قاعدہ اے، تاں اسیں اوہنوں مقدمیّاں دے ترجمہ ہوئے !!{proof}s اُتے لاؤنے آں۔ جیہڑے قاعدے~\Log{G3c} دے قاعدے وی نہیں، اوہناں لئی اسیں پہلاں ای وکھایا اے کہ گھٹّو گھٹ اوہ قبول کرن جوگ نیں۔ سو جے قاعدہ~$R$، $\LeftR{\land}$ یا $\RightR{\lor}$ سی، تاں \olref[adm]{prop:lor-G3-adm} لاؤ۔ جے کمزور کرن والا قاعدہ سی، تاں \olref[adm]{prop:weak-G3c-adm} لاؤ۔ جے سکڑاؤ دا قاعدہ سی، تاں \olref[inv]{prop:G3c-cont-adm} لاؤ۔
\end{proof}

\begin{prop}
جے $\Log{G3c} \Proves S$، تاں $\Log{G1c} \Proves S$۔
\end{prop}

\begin{proof}
  اسیں فیر !!{height} اُتے استقرا کردے آں، !!a{proof}~$\pi$ دی جیہڑا $S$ دا اے۔ \Log{G3c} دا ہر مسلّمہ !!{prove} ہو سکدا اے~\Log{G1c} وچ، اک مسلّمہ تے \LeftR{\Weakening} تے \RightR{\Weakening} قاعدیاں دیاں کئی واری ورتوں نال:
  \[
  \Axiom$!C \fCenter !C$
  \RightLabel{\LeftR{\Weakening}}
  \Deduce$!C, \Gamma \fCenter !C$
  \RightLabel{\RightR{\Weakening}}
  \Deduce$!C, \Gamma \fCenter \Delta, !C$
  \DisplayProof
\qquad
  \Axiom$\lfalse \fCenter$
  \RightLabel{\LeftR{\Weakening}}
  \Deduce$\lfalse, \Gamma \fCenter$
  \RightLabel{\RightR{\Weakening}}
  \Deduce$\lfalse, \Gamma \fCenter \Delta$
  \DisplayProof
  \]
  جے $\pi$ قاعدے~$R$ اُتے مُکدا اے، تاں استقرا دا مفروضہ مقدمیّاں دے !!{proof}s \Log{G1c} وچ دیندا اے تے اسیں اوہی قاعدہ~$R$ اوہناں !!{proof}s اُتے لا سکدے آں۔ چھوٹ \LeftR{\land} تے \RightR{\lor} قاعدیاں دی اے، جیہڑے \Log{G1c} تے~\Log{G3c} وچ وکھرے نیں۔ \Log{G3c} دیاں شکلاں \Log{G1c} وچ ایویں اخذ ہو سکدیاں نیں:
  \[
  \Axiom$!A, !B, \Gamma \fCenter \Delta$
  \RightLabel{\LeftR{\land}}
  \UnaryInf$!A \land !B, !B, \Gamma \fCenter \Delta$
  \RightLabel{\LeftR{\land}}
  \UnaryInf$!A \land !B, !A \land !B, \Gamma \fCenter \Delta$
  \RightLabel{\LeftR{\Contraction}}
  \UnaryInf$!A \land !B, \Gamma \fCenter \Delta$
    \DisplayProof
\qquad
  \Axiom$\Gamma \fCenter \Delta, !A, !B$
  \RightLabel{\RightR{\lor}}
  \UnaryInf$\Gamma \fCenter \Delta, !A \lor !B, !B$
  \RightLabel{\RightR{\lor}}
  \UnaryInf$\Gamma \fCenter \Delta, , !A \lor !B, !A \lor !B$
  \RightLabel{\RightR{\Contraction}}
  \UnaryInf$\Gamma \fCenter \Delta, !A \lor !B$
    \DisplayProof
  \]
\end{proof}
\endgroup
% /OLP-0713

% OLP-0705 content/proof-theory/sequent-calculus/rules-G1i.tex
\begingroup
\def\olfilename{rules-G1i}
\expandafter\def\csname ol@sectioncs\endcsname{section}

\begin{table}
\begin{tabular}{@{}c@{}c@{}}
\multicolumn{2}{@{}c@{}}{Axioms}\\[2ex]
$!A \Sequent !A$ & $\lfalse \Sequent \quad$
\\[2ex]
\multicolumn{2}{@{}c@{}}{Structural Rules}\\[2ex]
\Axiom$ \Gamma \fCenter \Delta $
\RightLabel{\LeftR{\Weakening}}
\UnaryInf$ !A, \Gamma \fCenter \Delta$
\DisplayProof
&
\Axiom$ \Gamma \fCenter $
\RightLabel{\RightR{\Weakening}}
\UnaryInf$ \Gamma \fCenter !A$
\DisplayProof
\\[4ex]
\Axiom$ !A, !A, \Gamma \fCenter \Delta $
\RightLabel{\LeftR{\Contraction}}
\UnaryInf$ !A, \Gamma \fCenter \Delta$
\DisplayProof
&
\\[4ex]
\multicolumn{2}{@{}c@{}}{Logical Rules}\\[2ex]
\Axiom$ \Gamma \fCenter !A $
\RightLabel{\LeftR{\lnot}}
\UnaryInf$ \lnot !A, \Gamma \fCenter $
\DisplayProof
&
\Axiom$!A, \Gamma \fCenter $
\RightLabel{\RightR{\lnot}}
\UnaryInf$ \Gamma \fCenter \lnot !A $
\DisplayProof
\\[4ex]\begin{tabular}[t]{@{}l@{}}
\Axiom$ !A, \Gamma \fCenter \Delta$
\RightLabel{\LeftR{\land}}
\UnaryInf$ !A \land !B, \Gamma \fCenter \Delta$
\DisplayProof
\\[3ex]
\Axiom$!B, \Gamma \fCenter \Delta$
\RightLabel{\LeftR{\land}}
\UnaryInf$!A \land !B, \Gamma \fCenter \Delta$
\DisplayProof\\[3ex]
\end{tabular}
&
\Axiom$\Gamma \fCenter !A$
\Axiom$ \Gamma \fCenter !B$
\RightLabel{\RightR{\land}}
\BinaryInf$ \Gamma \fCenter !A \land !B $
\DisplayProof
\\[4ex]\Axiom$!A, \Gamma \fCenter \Delta$
\Axiom$!B, \Gamma \fCenter \Delta$
\RightLabel{\LeftR{\lor}}
\BinaryInf$!A \lor !B, \Gamma \fCenter \Delta$
\DisplayProof
&
\begin{tabular}[t]{@{}r@{}}
\Axiom$\Gamma \fCenter !A$
\RightLabel{\RightR{\lor}}
\UnaryInf$ \Gamma \fCenter !A \lor !B$
\DisplayProof
\\[3ex]
\Axiom$ \Gamma \fCenter !B$
\RightLabel{\RightR{\lor}}
\UnaryInf$ \Gamma \fCenter !A \lor !B$
\DisplayProof\\[3ex]
\end{tabular}
\\[4ex]
\Axiom$ \Gamma \fCenter !A$
\Axiom$ !B, \Gamma \fCenter \Delta$
\RightLabel{\LeftR{\lif}}
\BinaryInf$ !A \lif !B, \Gamma \fCenter \Delta$
\DisplayProof
&
\Axiom$ !A, \Gamma \fCenter !B$
\RightLabel{\RightR{\lif}}
\UnaryInf$ \Gamma \fCenter !A \lif !B $
\DisplayProof
\\[4ex]
\Axiom$ !A(t), \Gamma \fCenter \Delta$
\RightLabel{\LeftR{\lforall}}
\UnaryInf$ \lforall[x][!A(x)],\Gamma \fCenter \Delta$
\DisplayProof
&
\Axiom$ \Gamma \fCenter !A(a) $
\RightLabel{\RightR{\lforall}}
\UnaryInf$ \Gamma \fCenter \lforall[x][!A(x)]$
\DisplayProof
\\[4ex]
\Axiom$ !A(a), \Gamma \fCenter \Delta $
\RightLabel{\LeftR{\lexists}}
\UnaryInf$ \lexists[x][!A(x)], \Gamma \fCenter \Delta$
\DisplayProof
&
\Axiom$ \Gamma \fCenter !A(t) $
\RightLabel{\RightR{\lexists}}
\UnaryInf$ \Gamma \fCenter \lexists[x][!A(x)]$
\DisplayProof
\end{tabular}
\caption{وجدانی منطق لئی \Log{G1i} دے قاعدے۔
$\RightR{\forall}$ تے $\LeftR{\exists}$ وچ !!{constant}~$a$ نتیجے والے سیکوئنٹ وچ نہیں آؤنا چاہیدا۔ $\Gamma$، !!{formula}s دا کثیر سیٹ اے، تے $\Delta$ خالی ہو سکدا اے یا اکّو !!{formula} رکھ سکدا اے۔ \RightR{\Contraction} قاعدہ کوئی نہیں۔ \Log{G1m}، \Log{G1i} اے~\RightR{\Weakening} توں بناں تے مسلّمیاں $\lfalse,
\Gamma \Sequent \Delta$ توں بناں۔}
\label{pt:seq:alt:tab:G1i}
\end{table}
\endgroup
% /OLP-0705
% OLP-0708 content/proof-theory/sequent-calculus/rules-G3i.tex
\begingroup
\def\olfilename{rules-G3i}
\expandafter\def\csname ol@sectioncs\endcsname{section}

\begin{table}
\begin{tabular}{@{}c@{}c@{}}
\multicolumn{2}{@{}c@{}}{Axioms}\\[2ex]
$!C, \Gamma \Sequent !C$ & $\lfalse, \Gamma \Sequent \Delta$
\\[2ex]
\multicolumn{2}{@{}c@{}}{Logical Rules}\\[2ex]
\Axiom$\lnot !A, \Gamma \fCenter !A $
\RightLabel{\LeftR{\lnot}}
\UnaryInf$\lnot !A, \Gamma \fCenter \Delta$
\DisplayProof
&
\Axiom$!A, \Gamma \fCenter$
\RightLabel{\RightR{\lnot}}
\UnaryInf$ \Gamma \fCenter \lnot !A $
\DisplayProof
\\[4ex]
\Axiom$ !A, !B, \Gamma \fCenter \Delta$
\RightLabel{\LeftR{\land}}
\UnaryInf$ !A \land !B, \Gamma \fCenter \Delta$
\DisplayProof
&
\Axiom$\Gamma \fCenter !A$
\Axiom$ \Gamma \fCenter !B$
\RightLabel{\RightR{\land}}
\BinaryInf$ \Gamma \fCenter !A \land !B $
\DisplayProof
\\[4ex]\Axiom$!A, \Gamma \fCenter \Delta$
\Axiom$!B, \Gamma \fCenter \Delta$
\RightLabel{\LeftR{\lor}}
\BinaryInf$!A \lor !B, \Gamma \fCenter \Delta$
\DisplayProof
&
\Axiom$\Gamma \fCenter !A, !B$
\RightLabel{\RightR{\lor}}
\UnaryInf$ \Gamma \fCenter !A \lor !B$
\DisplayProof
\\[4ex]
\Axiom$!A \lif !B, \Gamma \fCenter !A$
\Axiom$!B, \Gamma \fCenter \Delta$
\RightLabel{\LeftR{\lif}}
\BinaryInf$!A \lif !B, \Gamma \fCenter \Delta$
\DisplayProof
&
\Axiom$!A, \Gamma \fCenter !B$
\RightLabel{\RightR{\lif}}
\UnaryInf$ \Gamma \fCenter !A \lif !B $
\DisplayProof
\\[4ex]
\Axiom$ !A(t), \lforall[x][!A(x)]\Gamma \fCenter \Delta$
\RightLabel{\LeftR{\lforall}}
\UnaryInf$ \lforall[x][!A(x)],\Gamma \fCenter \Delta$
\DisplayProof
&
\Axiom$ \Gamma \fCenter !A(a) $
\RightLabel{\RightR{\lforall}}
\UnaryInf$ \Gamma \fCenter \lforall[x][!A(x)]$
\DisplayProof
\\[4ex]
\Axiom$ !A(a), \Gamma \fCenter \Delta $
\RightLabel{\LeftR{\lexists}}
\UnaryInf$ \lexists[x][!A(x)], \Gamma \fCenter \Delta$
\DisplayProof
&
\Axiom$ \Gamma \fCenter !A(t) $
\RightLabel{\RightR{\lexists}}
\UnaryInf$ \Gamma \fCenter \lexists[x][!A(x)]$
\DisplayProof
\end{tabular}
\caption{وجدانی منطق لئی \Log{G3i} دے قاعدے۔
$\RightR{\forall}$ تے $\LeftR{\exists}$ وچ !!{constant}~$a$ نتیجے والے سیکوئنٹ وچ نہیں آؤنا چاہیدا۔ $\Gamma$، !!{formula}s دا کثیر سیٹ اے، تے $\Delta$ خالی ہو سکدا اے یا اکّو !!{formula} رکھ سکدا اے۔ مسلّمیاں وچ !!{formula}~$!C$ جوہری ہونا چاہیدا اے۔ \Log{G1m}، \Log{G1i} اے مسلّمیاں $\lfalse, \Gamma \Sequent
\Delta$ توں بناں۔}
\label{pt:seq:alt:tab:G3i}
\end{table}
\endgroup
% /OLP-0708
% OLP-0709 content/proof-theory/sequent-calculus/rules-LK.tex
\begingroup
\def\olfilename{rules-LK}
\expandafter\def\csname ol@sectioncs\endcsname{section}

\begin{table}
\begin{tabular}{@{}c@{}c@{}}
\multicolumn{2}{@{}c@{}}{Axioms}\\[2ex]
$!A \Sequent !A$ & $\lfalse \Sequent \quad$
\\[2ex]
\multicolumn{2}{@{}c@{}}{Structural Rules}\\[2ex]
\Axiom$ \Gamma \fCenter \Delta $
\RightLabel{\LeftR{\Weakening}}
\UnaryInf$ !A, \Gamma \fCenter \Delta$
\DisplayProof
&
\Axiom$ \Gamma \fCenter \Delta$
\RightLabel{\RightR{\Weakening}}
\UnaryInf$ \Gamma \fCenter \Delta, !A$
\DisplayProof
\\[4ex]
\Axiom$ !A, !A, \Gamma \fCenter \Delta $
\RightLabel{\LeftR{\Contraction}}
\UnaryInf$ !A, \Gamma \fCenter \Delta$
\DisplayProof
&
\Axiom$ \Gamma \fCenter \Delta, !A, !A$
\RightLabel{\RightR{\Contraction}}
\UnaryInf$ \Gamma \fCenter \Delta, !A$
\DisplayProof
\\[4ex]
\Axiom$ \Gamma, !A, !B, \Pi \fCenter \Delta $
\RightLabel{\LeftR{\Exchange}}
\UnaryInf$\Gamma, !B, !A, \Pi \fCenter \Delta$
\DisplayProof
&
\Axiom$ \Gamma \fCenter \Delta, !A, !B, \Lambda$
\RightLabel{\RightR{\Exchange}}
\UnaryInf$ \Gamma \fCenter \Delta, !B, !A, \Lambda$
\DisplayProof
\\[4ex]
\multicolumn{2}{@{}c@{}}{Logical Rules}\\[2ex]
\Axiom$ \Gamma \fCenter \Delta, !A $
\RightLabel{\LeftR{\lnot}}
\UnaryInf$ \lnot !A, \Gamma \fCenter \Delta$
\DisplayProof
&
\Axiom$!A, \Gamma \fCenter \Delta$
\RightLabel{\RightR{\lnot}}
\UnaryInf$ \Gamma \fCenter \Delta, \lnot !A $
\DisplayProof
\\[4ex]\begin{tabular}[t]{@{}l@{}}
\Axiom$ !A, \Gamma \fCenter \Delta$
\RightLabel{\LeftR{\land}}
\UnaryInf$ !A \land !B, \Gamma \fCenter \Delta$
\DisplayProof
\\[3ex]
\Axiom$!B, \Gamma \fCenter \Delta$
\RightLabel{\LeftR{\land}}
\UnaryInf$!A \land !B, \Gamma \fCenter \Delta$
\DisplayProof\\[3ex]
\end{tabular}
&
\Axiom$\Gamma \fCenter \Delta, !A$
\Axiom$ \Gamma \fCenter \Delta, !B$
\RightLabel{\RightR{\land}}
\BinaryInf$ \Gamma \fCenter \Delta, !A \land !B $
\DisplayProof
\\[4ex]\Axiom$!A, \Gamma \fCenter \Delta$
\Axiom$!B, \Gamma \fCenter \Delta$
\RightLabel{\LeftR{\lor}}
\BinaryInf$!A \lor !B, \Gamma \fCenter \Delta$
\DisplayProof
&
\begin{tabular}[t]{@{}r@{}}
\Axiom$\Gamma \fCenter \Delta, !A$
\RightLabel{\RightR{\lor}}
\UnaryInf$ \Gamma \fCenter \Delta, !A \lor !B$
\DisplayProof
\\[3ex]
\Axiom$ \Gamma \fCenter \Delta, !B$
\RightLabel{\RightR{\lor}}
\UnaryInf$ \Gamma \fCenter \Delta, !A \lor !B$
\DisplayProof\\[3ex]
\end{tabular}
\\[4ex]
\Axiom$ \Gamma \fCenter \Delta, !A$
\Axiom$ !B, \Gamma \fCenter \Delta$
\RightLabel{\LeftR{\lif}}
\BinaryInf$ !A \lif !B, \Gamma \fCenter \Delta$
\DisplayProof
&
\Axiom$ !A, \Gamma \fCenter \Delta, !B$
\RightLabel{\RightR{\lif}}
\UnaryInf$ \Gamma \fCenter \Delta, !A \lif !B $
\DisplayProof
\\[4ex]
\Axiom$ !A(t), \Gamma \fCenter \Delta$
\RightLabel{\LeftR{\lforall}}
\UnaryInf$ \lforall[x][!A(x)],\Gamma \fCenter \Delta$
\DisplayProof
&
\Axiom$ \Gamma \fCenter \Delta, !A(a) $
\RightLabel{\RightR{\lforall}}
\UnaryInf$ \Gamma \fCenter \Delta, \lforall[x][!A(x)]$
\DisplayProof
\\[4ex]
\Axiom$ !A(a), \Gamma \fCenter \Delta $
\RightLabel{\LeftR{\lexists}}
\UnaryInf$ \lexists[x][!A(x)], \Gamma \fCenter \Delta$
\DisplayProof
&
\Axiom$ \Gamma \fCenter \Delta, !A(t) $
\RightLabel{\RightR{\lexists}}
\UnaryInf$ \Gamma \fCenter \Delta, \lexists[x][!A(x)]$
\DisplayProof
\end{tabular}
\caption{\Log{LK} دے قاعدے۔ $\RightR{\forall}$ تے
$\LeftR{\exists}$ وچ !!{constant}~$a$ نتیجے والے سیکوئنٹ وچ نہیں آؤنا چاہیدا۔ $\Gamma$، $\Delta$ تے~$\Pi$، !!{formula}s دیاں لڑیاں نیں۔}
\label{pt:seq:alt:tab:LK}
\end{table}
\endgroup
% /OLP-0709
% OLP-0710 content/proof-theory/sequent-calculus/rules-mG3i.tex
\begingroup
\def\olfilename{rules-mG3i}
\expandafter\def\csname ol@sectioncs\endcsname{section}

\begin{table}
\begin{tabular}{@{}c@{}c@{}}
\multicolumn{2}{@{}c@{}}{Axioms}\\[2ex]
$!C, \Gamma \Sequent \Delta, !C$ & $\lfalse, \Gamma \Sequent \Delta$
\\[2ex]
\multicolumn{2}{@{}c@{}}{Logical Rules}\\[2ex]
\Axiom$\lnot !A, \Gamma \fCenter \Delta, !A $
\RightLabel{\LeftR{\lnot}}
\UnaryInf$\lnot !A, \Gamma \fCenter \Delta$
\DisplayProof
&
\Axiom$!A, \Gamma \fCenter$
\RightLabel{\RightR{\lnot}}
\UnaryInf$ \Gamma \fCenter \Delta, \lnot !A $
\DisplayProof
\\[4ex]
\Axiom$ !A, !B, \Gamma \fCenter \Delta$
\RightLabel{\LeftR{\land}}
\UnaryInf$ !A \land !B, \Gamma \fCenter \Delta$
\DisplayProof
&
\Axiom$\Gamma \fCenter \Delta, !A$
\Axiom$ \Gamma \fCenter \Delta, !B$
\RightLabel{\RightR{\land}}
\BinaryInf$ \Gamma \fCenter \Delta, !A \land !B $
\DisplayProof
\\[4ex]\Axiom$!A, \Gamma \fCenter \Delta$
\Axiom$!B, \Gamma \fCenter \Delta$
\RightLabel{\LeftR{\lor}}
\BinaryInf$!A \lor !B, \Gamma \fCenter \Delta$
\DisplayProof
&
\Axiom$\Gamma \fCenter \Delta, !A, !B$
\RightLabel{\RightR{\lor}}
\UnaryInf$ \Gamma \fCenter \Delta, !A \lor !B$
\DisplayProof
\\[4ex]
\Axiom$!A \lif !B, \Gamma \fCenter \Delta, !A$
\Axiom$ !B, \Gamma \fCenter \Delta$
\RightLabel{\LeftR{\lif}}
\BinaryInf$!A \lif !B, \Gamma \fCenter \Delta$
\DisplayProof
&
\Axiom$ !A, \Gamma \fCenter !B$
\RightLabel{\RightR{\lif}}
\UnaryInf$ \Gamma \fCenter \Delta, !A \lif !B $
\DisplayProof
\\[4ex]
\Axiom$ !A(t), \lforall[x][!A(x)]\Gamma \fCenter \Delta$
\RightLabel{\LeftR{\lforall}}
\UnaryInf$ \lforall[x][!A(x)],\Gamma \fCenter \Delta$
\DisplayProof
&
\Axiom$ \Gamma \fCenter !A(a) $
\RightLabel{\RightR{\lforall}}
\UnaryInf$ \Gamma \fCenter \lforall[x][!A(x)]$
\DisplayProof
\\[4ex]
\Axiom$ !A(a), \Gamma \fCenter \Delta $
\RightLabel{\LeftR{\lexists}}
\UnaryInf$ \lexists[x][!A(x)], \Gamma \fCenter \Delta$
\DisplayProof
&
\Axiom$ \Gamma \fCenter \Delta, \lexists[x][!A(x)], !A(t) $
\RightLabel{\RightR{\lexists}}
\UnaryInf$ \Gamma \fCenter \Delta, \lexists[x][!A(x)]$
\DisplayProof
\end{tabular}
\caption{وجدانی منطق لئی کئی نتیجیاں والے \Log{mG3i} دے قاعدے۔
$\RightR{\forall}$ تے $\LeftR{\exists}$ وچ !!{constant}~$a$ نتیجے والے سیکوئنٹ وچ نہیں آؤنا چاہیدا۔ $\Gamma$ تے $\Delta$، !!{formula}s دے کثیر سیٹ نیں۔ مسلّمیاں وچ !!{formula}~$!C$ جوہری ہونا چاہیدا اے۔ \Log{mG1m}، \Log{mG1i} اے مسلّمیاں $\lfalse,
\Gamma \Sequent \Delta$ توں بناں۔}
\label{pt:seq:alt:tab:mG3i}
\end{table}
\endgroup
% /OLP-0710
\OLEndChapterHook
\endgroup
% /OLP-0712
% OLP-0657 content/proof-theory/cut-elimination/cut-elimination.tex
\begingroup
\def\olfilename{cut-elimination}
\expandafter\def\csname ol@sectioncs\endcsname{section}

\olchapter{pt}{cut}{قطع ہٹاؤنا}

% OLP-0659 content/proof-theory/cut-elimination/introduction.tex
\begingroup
\def\olfilename{introduction}
\expandafter\def\csname ol@sectioncs\endcsname{section}

\olfileid{pt}{cut}{int}

\olsection{جان پچھان}

\Cut{} دا قاعدہ ایہ اے:
\[
\Axiom$\Gamma \fCenter \Delta, !A$
\Axiom$!A, \Pi \fCenter \Lambda$
\RightLabel{\Cut}
\BinaryInf$\Gamma, \Pi \fCenter \Delta, \Lambda$
\DisplayProof
\]
(!!{formula}~$!A$ نوں \emph{قطع فارمولا} آکھدے نیں۔)

\Cut{} گینٹسن دے اصلی سیکوئنٹ حساب~\Log{LK} وچ شامل سی۔ گینٹسن دا \emph{Hauptsatz} («مُکھ قضیہ») کہندا اے کہ ہر سیکوئنٹ جیہڑا~\Log{LK} وچ !!{provable} اے، \Cut{} دا قاعدہ ورتے بغیر وی !!{provable} اے۔ یعنی \Cut{} نوں \Log{LK} دے ثبوتاں وچوں «ہٹایا» جا سکدا اے۔ ساڈے نظام \Log{G1c} تے \Log{G3c} وچ \Cut نہیں، پر اوہناں لئی وی ایہ سوال پچھیا جا سکدا اے۔ کی \Cut{} اوہناں !!{proof}s وچوں ہٹایا جا سکدا اے جیہڑے \Cut{} نال \Log{G1c} دے باقی قاعدے (یا \Log{G3c} دے قاعدے) ورتدے نیں؟ ساڈیاں قائم اصطلاحاں وچ ایہو سوال ایویں اے۔ کی \Cut{}، \Log{G1c} (یا \Log{G3c}) دا قبول جوگ قاعدہ اے؟

قطع ہٹاؤن دا قضیہ شاید ثبوتی نظریے دا سبھ توں اہم تے مشہور نتیجہ اے۔ پہلاں ایہ قائم کردا اے کہ \Log{G1c} تے \Log{G3c} نوں اک ایہو جیہے قاعدے دی لوڑ نہیں جیہڑا پہلی نظرے صاف ضروری لگدا اے۔ جدوں تسی اصل ثبوتاں نوں رسمی شکل دین دی کوشش کرو، تاں دِسدا اے کہ معاون قضیے ورتن دا ڈھنگ چاہیدا اے۔ ریاضی دان نتیجے ثابت کردے نیں تے فیر ہور نتیجیاں دے ثبوتاں وچ اوہناں نوں ورتدے نیں۔ سو پہلاں نتیجے $L$ (معاون قضیے) دا ثبوت، سیکوئنٹ $\Sequent L$ دے !!{proof} وجوں، رسمی شکل وچ لکھ سکدے او۔ فیر دوجے نتیجے~$R$ دا ثبوت، جیہڑا~$L$ ورتدا اے، سیکوئنٹ $L \Sequent
R$ دے !!{proof} وجوں لکھدے او۔ تہاڈے کول~$R$ دا ثبوت اے، یعنی بس !!a{proof} اے $\Sequent R$ دا، ایہ کِویں پتا لگے؟ دوہاں !!{proof}s نوں جوڑن دا ڈھنگ چاہیدا اے، تے \Cut{} دا قاعدہ ایہ کرن دی کھُل دیندا اے۔

اسیں دسیا سی کہ \Log{G1c} تے \Log{G3c} پورے نیں، سو اوہ ہر اوہ سیکوئنٹ ثابت کردے نیں جیہڑا ثابت ہو سکنا چاہیدا اے، یعنی سارے معتبر سیکوئنٹ۔ ایہ سدھا ثابت کیتا جا سکدا اے۔ پر جدوں گینٹسن لکھدا سی، اوہدوں بس مسلّماتی اخذ دے نظاماں بارے پتا سی کہ اوہ پورے نیں۔ \Log{LK} دی تمامیت وکھاؤن لئی گینٹسن نے مسلّماتی نظام دے !!{derivation}s نوں \Log{LK} دے !!{proof}s وچ بدلّن دا طریقہ دِتّا۔ ایس تبدیلی وچ \Cut{} ضروری طور اُتے ورتیا گیا سی۔ قطع ہٹاؤنا کلاسیکی منطق لئی ای اہم نتیجہ نہیں، کئی ہور منطقاں دے وی سیکوئنٹ حساب نیں۔ کجھ منطقاں وچ اوہناں دے سیکوئنٹ حساباں دی تمامیت دا سدھا ثبوت اوکھا یا ناممکن اے۔ سو دوجے نظاماں توں تبدیلی (جیہڑی \Cut ورتدی اے) ای تمامیت وکھاؤن دا اکّو ڈھنگ اے۔ فیر ثبوتی نظریے نال قطع ہٹاؤنا ثابت کرنا ضروری اے تاں جو قائم ہووے کہ \Cut{} وادھو اے تے \Cut{} توں بناں نظام پورا اے۔

قطع ہٹاؤنا ایس لئی وی اہم اے کہ ایہنوں ہور، اپنے آپ وچ دلچسپ، نتیجے کڈھن لئی ورتیا جا سکدا اے۔ اسیں دو صورتاں ویکھاں گے، کریگ دا درمیانی فارمولا والا معاون قضیہ تے ہیربرانڈ دا قضیہ۔

قطع ہٹاؤنا ثابت کرن دے خیالاں وچ عام طور اُتے ایہ وکھایا جاندا اے کہ \Cut{} ورتن والا !!a{proof} اوس توں بناں ثبوت وچ کِویں بدلیا جا سکدا اے۔ ایہ بدلاواں عموماً «مقامی» ہوندے نیں۔ یعنی اسیں !!
a{proof} دا اک حصہ لیندے آں (عام طور اُتے \Cut{} استنباط اُتے مُکدا ذیلی !!{proof}) تے اوہدی تھاں اوسے سیکوئنٹ دا ہور ذیلی !!{proof} رکھدے آں، جیہدے وچ \Cut{} استنباط ہٹا دِتّا گیا اے یا اوہدی تھاں ہور استنباط رکھے گئے نیں۔ ایہو جیہیاں دلیلاں قدم وار الگورتھم دیندیاں نیں جیہڑے \Cut{} والے !!{proof}s نوں اوس توں بناں ثبوتاں وچ بدل دیندے نیں۔ قطع ہٹاؤن دے طریقے ورتن ویلے ایس توں ہور جانکاری ملدی اے۔ مثلاً درمیانی فارمولے دے معاون قضیے دے نمونیاں دے نظریے والے ایہو جیہے ثبوت نیں کہ جے «$!A \lif !B$» معتبر اے، تاں اک درمیانی فارمولا موجود اے (ہُن لئی ایہ چھڈو کہ درمیانی فارمولا کیہ اے)۔ قطع ہٹاؤن والی ثبوتی دلیل اک الگورتھم دیندی اے جیہڑا !!a{proof} لے کے $!A
\lif !B$ دا درمیانی فارمولا موڑدا اے۔ نمونیاں دے نظریے والی دلیل دے مقابلے وچ، ثبوتی دلیل \emph{تعمیری} اے۔

قطع ہٹاؤن دا قضیہ ثابت کرن دے دو مشہور ڈھنگ نیں۔ اک، جیہڑا گینٹسن توں آؤندا اے، وکھاؤندا اے کہ ثبوت وچوں سبھ توں اُتلے قطع ہٹا سکدے آں، تے فیر اوہ اُتلے قطع ہٹائی جاندا اے جدوں تک کوئی نہیں رہندا۔ دوجا، جیہڑا مُڈھوں شوٹے تے ٹیٹ دا اے، !!a{proof} دے سبھ توں گنجل دار قطعاں اُتے دھیان دیندا اے تے وار واری اوہ قطع ہٹاؤندا اے جیہڑے سبھ توں اُچّے درجے دے نیں، ایس معنی وچ کہ ہور کسے قطع دا قطع !!{formula} اوہناں توں ڈونگھا نہیں۔ دوہاں دے اپنے فائدے تے اوکڑاں نیں۔ کجھ نظاماں وچ اک ڈھنگ کم کردا اے تے دوجا لاؤنا اوکھا یا ناممکن اے۔ سو اسیں دوویں ڈھنگ دساں گے۔ اسیں~\Log{G3c} لئی قطع ہٹاؤنا ثابت کراں گے۔ اصل وچ اسیں \Cut{} ہٹاؤن دی تھاں \emph{سانجھے سیاق والا قطع}~\CutCS ہٹاؤنا وکھاواں گے:
\[
\Axiom$\Gamma \fCenter \Delta, !A$
\Axiom$!A, \Gamma \fCenter \Delta$
\RightLabel{\Cut}
\BinaryInf$\Gamma \fCenter \Delta$
\DisplayProof
\]
جے سیکوئنٹ حساب کمزور کرن تے سکیڑن دی کھُل دیندا اے (بھاویں \Log{G1c} وانگ اصلی قاعدیاں وجوں یا \Log{G3c} وانگ قبول جوگ قاعدیاں وجوں)، تاں \Cut{}، \CutCS دا کم کر سکدا اے، تے اُلٹا وی۔ اسیں \CutCS{} چُندے آں \Cut دی تھاں، کیوں جے پہلا طریقہ~\CutCS لئی سمجھاؤنا سوکھا اے تے دوجا طریقہ بس~\CutCS لئی کم کردا اے۔
\endgroup
% /OLP-0659
% OLP-0656 content/proof-theory/cut-elimination/ce-topmost.tex
\begingroup
\def\olfilename{ce-topmost}
\expandafter\def\csname ol@sectioncs\endcsname{section}

\olfileid{pt}{cut}{top}

\olsection{سبھ توں اُتلے قطع ہٹاؤنے}

\begin{defn}
!!a{proof}~$\pi$ ویکھو جیہڑا \CutCS{} قاعدے اُتے مُکدا اے،
\[
\AxiomC{}
\RightLabel{$\pi_1$}
\Deduce$\Gamma \fCenter \Delta, !A$
\AxiomC{}
\RightLabel{$\pi_2$}
\Deduce$!A, \Gamma \fCenter \Delta$
\RightLabel{\CutCS}
\BinaryInf$\Gamma \fCenter \Delta$
\DisplayProof
\]
تے جیہدے وچ ہور کوئی \CutCS{} استنباط نہیں۔~$\pi$ دی \emph{قطع اُچائی} $\cheight{\pi} = \pheight{\pi_1} + \pheight{\pi_2}$ اے۔~$\pi$ دا \emph{قطع درجہ} $\cutrank{\pi} = \depth{!A}$ اے۔
\end{defn}

\begin{lem}\ollabel{lem:cut-adm-G3c}
\CutCS{} قاعدہ~\Log{G3c} وچ قبول جوگ اے۔ دوجے لفظاں وچ، جے $\pi$ اک !!a{proof} اے جیہڑا اکّو~\CutCS{} اُتے مُکدا اے تے ہور بس~$\Log{G3c}$ دے قاعدے ورتدا اے، تاں اوسے آخری سیکوئنٹ دا !!a{proof}~$\pi'$،~\Log{G3c} وچ، موجود اے۔
\end{lem}

\begin{proof}
ثبوت قطع !!{height} تے درجے اُتے دوہری استقرا نال اے۔ !!{proof}~$\pi$ ایویں مُکدا اے:
\[
\AxiomC{}
\RightLabel{$\pi_1$}
\Deduce$\Gamma \fCenter \Delta, !A$
\AxiomC{}
\RightLabel{$\pi_2$}
\Deduce$!A, \Gamma \fCenter \Delta$
\RightLabel{\CutCS}
\BinaryInf$\Gamma \fCenter \Delta$
\DisplayProof
\]

استقرا دی مُڈھلی صورت۔ $\cheight{\pi} = 0$ تے $\cutrank{\pi} = 0$۔ چونکہ $\cheight{\pi} = \pheight{\pi_1} + \pheight{\pi_2} = 0$، ایس لئی $\Gamma
\Sequent \Delta, !A$ تے $!A, \Gamma \Sequent \Delta$ دوویں مسلّمے نیں۔ دو صورتاں نیں۔ (a) اوہناں وچوں اک وچ $!A$ مُکھ نہیں، یا (b) دوہاں وچ $!A$ مُکھ اے۔ تھلّے دوہاں صورتاں وچ وکھاواں گے کہ $\Gamma \Sequent \Delta$ لازماً مسلّمہ اے (تے ایس لئی استقرا دے مفروضے دی لوڑ نہیں)۔

ہُن کئی ممکن صورتاں وکھ کراں گے۔ فرق ایہ اے کہ $\pi_1$ تے $\pi_2$ مسلّمے نیں یا کسے استنباط اُتے مُکدے نیں، تے $!A$ اوس استنباط وچ مُکھ اے یا نہیں۔ صورتاں A تے B استقرا دی مُڈھلی صورت نوں سمیٹدیاں نیں۔ استقرا دے قدم وچ منّنے آں کہ $\cheight{\pi} > 0$ یا $\cutrank{\pi} > 0$ (یا دوویں)۔ اسیں استقرا دا مفروضہ ورت سکدے آں۔ یعنی آخری~\CutCS{} ہر اوہ !!{proof} وچوں ہٹایا جا سکدا اے جیہڑا اکّو~\CutCS{} اُتے مُکدا اے جیہدا قطع درجہ $= \cutrank{\pi}$ اے تے قطع !!{height} $< \cheight{\pi}$، یا جیہدا قطع درجہ~$<
\cutrank{\pi}$ اے۔ جتھے $\cheight{\pi} = 0$ (دوویں مقدمے مسلّمے نیں)، اوہ صورت A تے B وچ آؤندی اے۔ جتھے $\cheight{\pi} >
0$، اوہ صورتاں C--D وچ آؤندیاں نیں۔

\paragraph{صورت A۔}
کوئی $!B$، $\Gamma$ تے~$\Delta$ دوہاں وچ آؤندا اے۔ تاں $\Gamma \Sequent \Delta$ یا مسلّمہ اے (جے $!B$ جوہری اے) یا اوہدا \Log{G3c} وچ !!a{proof}~$\pi'$ اے، \CutCS{} توں بناں، \olref[seq][ex]{prop:G3c-axioms} مطابق۔ ایہ استقرا دی مُڈھلی صورت دا (a) اے، کیوں جے جے $!A$، $\Gamma
\Sequent \Delta, !A$ وچ مُکھ نہیں یا $!A, \Gamma \Sequent \Delta$ وچ مُکھ نہیں، تے ایہ مسلّمے نیں، تاں کوئی ہور جوہری~$!B$، $\Gamma$ تے~$\Delta$ دوہاں وچ ضرور آؤندا اے۔ ایہ استقرا دے قدم اُتے وی لگدا اے جے اک مقدمہ مسلّمہ اے جیہدے وچ $!A$ مُکھ نہیں (بھاویں دوجا مقدمہ کسے قاعدے دا نتیجہ ہووے)۔

\paragraph{صورت B۔} دوویں مقدمے مسلّمے نیں، تے $!A$ مُکھ اے، سو جوہری اے۔ کھبّا مقدمہ، $\Gamma \Sequent
\Delta, !A$، ویکھو۔ چونکہ $!A$ مُکھ اے، ایہ $\Gamma$ وچ ضرور آؤندا اے (نہیں تاں $\Gamma \Sequent \Delta, !A$ مسلّمہ نہ ہوندا)۔ ایویں ای $!A$، $\Delta$ وچ ضرور آؤندا اے، نہیں تاں $!A, \Gamma \Sequent \Delta$ مسلّمہ نہ ہوندا۔ سو $!A$ جوہری اے تے $\Gamma$ تے $\Delta$ دوہاں وچ آؤندا اے، یعنی $\Gamma \Sequent \Delta$ مسلّمہ اے۔ ایہ استقرا دی مُڈھلی صورت دا (b) اے۔

\begin{center}
صورتاں A تے B دا ہور بیان
\end{center}

\paragraph{صورت A۔} قطع دا اک مقدمہ ایس شکل دا مسلّمہ اے: $!C, \Gamma' \Sequent \Delta', !C$، جتھے $!C$ جوہری اے۔ جے $!C$ قطع !!{formula}~$!A$ نہیں، تاں $!C$، $\Gamma$ تے~$\Delta$ دوہاں وچ ضرور آؤندا اے۔ ایس صورت وچ $\Gamma \Sequent \Delta$ مسلّمہ اے تے اسیں ایہنوں~$\pi'$ لے سکدے آں۔ جے $!C \ident !A$ قطع !!{formula} اے، تاں یا (جے کھبّا مقدمہ مسلّمہ اے) $!A \in \Gamma$ تے $\Gamma = \Gamma', !A$، یا (جے سجّا مقدمہ مسلّمہ اے) $!A \in
\Delta$ تے $\Delta = \Delta', !A$۔ پہلی صورت وچ $\pi_2$ اک !!a{proof} اے $!A, !A, \Gamma' \Sequent \Delta$ دا، تے \LeftR{\Contraction} دی قبولیت نال (\olref[seq][inv]{prop:G3c-cont-adm}) اک !!a{proof} ملدا اے $\Gamma \Sequent \Delta$ دا۔ دوجی صورت وچ $\pi_2$ اک !!a{proof} اے $\Gamma \Sequent \Delta', !A, !A$ دا، تے \RightR{\Contraction} دی قبولیت نال اک !!a{proof} ملدا اے $\Gamma \Sequent
\Delta$ دا۔

\paragraph{صورت B۔} مقدمیّاں وچوں اک ایس شکل دا مسلّمہ اے:~$\lfalse, \Gamma' \Sequent \Delta'$۔ جے کھبّا مقدمہ ایہو جیہا مسلّمہ اے، تاں $\lfalse \in \Gamma$ تے $\Gamma \Sequent \Delta$ وی مسلّمہ اے۔ جے سجّا مقدمہ ایہو جیہا مسلّمہ اے، تاں یا $\lfalse \in \Gamma$ (فیر $\Gamma \Sequent \Delta$ مسلّمہ اے)، یا $!A \ident \lfalse$۔ پچھلی صورت وچ $\pi_1$ اک !!a{proof} اے $\Gamma \Sequent \Delta, \lfalse$ دا۔ اسیں ایہنوں !!a{proof} بنا سکدے آں~$\Gamma \Sequent \Delta$ دا،~$\lfalse$ دی اک نقل ہر سیکوئنٹ دے تالی وچوں مٹا کے،~$\pi_1$ وچ۔ (مشق۔ ایہ~$\pheight{\pi_1}$ اُتے استقرا نال پرکھو۔)

ہُن منّ لو صورت A تے صورت B دوہاں وچوں کوئی نہیں لگدی۔ باقی صورتاں وچ $\cheight{\pi} > 0$ اے، یعنی گھٹ توں گھٹ اک مقدمہ کسے استنباط دا نتیجہ اے۔

\paragraph{صورتاں C تے D۔ قطعاں دی تھاں بدلنا۔}
صورتاں C تے D وچ ویکھیاں جان والیاں ساریاں امکانی شکلاں دا اکّو ڈھنگ اے۔ پہلاں اک !!a{proof} لو جیہڑا~\CutCS{} اُتے مُکدا اے، جتھے اک مقدمہ قاعدے~$R$ نال ملدا اے جیہدے وچ قطع !!{formula}~$!A$ مُکھ نہیں (کوئی ہور !!{formula}~$!D$ مُکھ اے)۔ ایس !!{proof} وچ \CutCS{} قاعدے~$R$ پچھوں آؤندا اے، تے اوہدی خاکے والی شکل ایہ اے:
\[
\AxiomC{}
\RightLabel{$\pi_1$}
\Deduce$\Gamma \fCenter \Delta', !D, !A$
\AxiomC{}
\RightLabel{$\pi_3$}
\Deduce$!A, \Gamma, \Pi \fCenter \Delta', \Lambda$
\RightLabel{$R$}
\UnaryInf$!A, \Gamma \fCenter \Delta', !D$
\RightSubproofLabel{$\pi_2$}
\RightLabel{\CutCS}
\BinaryInf$\Gamma \fCenter \Delta', !D$
\DisplayProof
\]
ایہ بس اوہ صورت سمیٹدا اے جتھے $R$ اک مقدمے والا سجّا قاعدہ اے، $A$،~$R$ دا مُکھ !!{formula} نہیں، تے اوہ !!{formula}~$!D$ جیہڑا مُکھ \emph{اے}،~\CutCS دے سجّے مقدمے دے تالی وچ آؤندا اے۔ جے $!D$ مقدم وچ آوے (یعنی $R$ کھبّا قاعدہ ہووے)، یا~\CutCS دے کھبّے مقدمے وچ آوے، تاں اوہدی شکل ظاہر اے وکھری ہووے گی۔ $\Pi$ تے~$\Lambda$ وچلے !!{formula}s قاعدے~$R$ دے سرگردہ !!{formula}s نوں ظاہر کردے نیں۔

اسیں ایہ ترتیب اُلٹا کے اک !!a{proof} ویکھدے آں جتھے $R$،~\CutCS{} پچھوں آؤندا اے:
\[
\AxiomC{}
\RightLabel{$\pi_1'$}
\Deduce$\Gamma, \Pi \fCenter \Delta', !D, \Lambda, !A$
\AxiomC{}
\RightLabel{$\pi_3'$}
\Deduce$!A, \Gamma, \Pi \fCenter \Delta', !D, \Lambda$
\RightLabel{\CutCS}
\BinaryInf$\Gamma, \Pi \fCenter \Delta', !D, \Lambda$
\RightLabel{$R$}
\UnaryInf$\Gamma, \fCenter \Delta', !D, !D$
\DisplayProof
\]
اسیں !!{height} قائم رکھن والا کمزور کرن دا قاعدہ ورت کے $\pi_1'$ تے $\pi_3'$، بالترتیب $\pi_1$ تے $\pi_3$ توں لیندے آں۔

چونکہ اسیں \CutCS{} تے~$R$ دی ترتیب بدلنے آں، ایس نوں «قطع نوں~$R$ توں اُتّے چُکنا» آکھدے نیں۔ ایس نال~\CutCS دے سجّے مقدمے اُتے مُکّن والے !!{proof} دی اُچائی گھٹدی اے، سو قطع !!{height} گھٹدی اے۔ یعنی نویں \CutCS{} اُتے مُکّن والے ذیلی !!{proof}s دی !!{height}، $\pi_1$ تے $\pi_3$ توں ودھ نہیں منّ سکدے۔ سو قطع !!{height} گھٹ گئی اے (اسیں سجّے پاسے اک استنباط ہٹایا اے)۔ ہُن استقرا دا مفروضہ لگدا اے تے دسدا اے کہ \CutCS{} استنباط ہٹایا جا سکدا اے۔ آخر~$!D$ دی ودھ نقل~$\RightR{\Contraction}$ دی قبولیت نال ہٹاؤنے آں۔

\paragraph{صورت C۔}
$\pi_1$ تے $\pi_2$ وچوں گھٹ توں گھٹ اک، اک مقدمے والے قاعدے~$R$ اُتے مُکدا اے جیہدے وچ $!A$ مُکھ نہیں۔ کُل 18 صورتاں نیں۔ ایہ ایس گل اُتے نربھر اے کہ $!A$،~\CutCS دے کھبّے مقدمے وچ مُکھ نہیں یا سجّے وچ، تے ایہ وی کہ قاعدہ~$R$، اک مقدمے والے نو قاعدیاں وچوں کیہڑا اے (\LeftR{\lnot}، \RightR{\lnot}، \RightR{\lor}، \LeftR{\land}، \RightR{\lif} تے مسوّراں دے چار قاعدے)۔ اسیں بس دو مثالاں ویکھدے آں جتھے $!A$،~$\pi_2$ دے آخری استنباط وچ مُکھ نہیں، تے اوہ~$\RightR{\lif}$ یا~\LeftR{\lexists} اُتے مُکدا اے۔ باقی صورتاں ایہو جیہیاں نیں تے مشقاں وجوں چھڈ دِتّیاں نیں۔

منّ لو $\pi$ ایویں مُکدا اے:
\[
\AxiomC{}
\RightLabel{$\pi_1$}
\Deduce$\Gamma \fCenter \Delta', !B \lif !C, !A$
\AxiomC{}
\RightLabel{$\pi_3$}
\Deduce$!A, !B, \Gamma \fCenter \Delta', !C$
\RightLabel{$\RightR{\lif}$}
\UnaryInf$!A, \Gamma \fCenter \Delta', !B \lif !C$
\RightSubproofLabel{$\pi_2$}
\RightLabel{\CutCS}
\BinaryInf$\Gamma \fCenter \Delta', !B \lif !C$
\DisplayProof
\]
جتھے $\Delta = \Delta', !B \lif !C$۔ !!{height} قائم رکھن والا کمزور کرن دا قاعدہ (\olref[seq][adm]{prop:weak-G3c-adm})، $\pi_1$ اُتے لا کے !!a{proof} $\pi_1'$ ملدا اے $!B, \Gamma \fCenter \Delta', !B \lif !C, !C,
!A$ دا (کھبّے $!B$ نال تے سجّے $!C$ نال کمزور کردے آں)۔ ایویں ای \olref[seq][adm]{prop:weak-G3c-adm} نوں~$\pi_3$ اُتے لا کے !!a{proof}~$\pi_3'$ ملدا اے $!A, !B, \Gamma \fCenter
\Delta', !B \lif !C, !C$ دا (سجّے $!B \lif !C$ نال کمزور کر کے)۔ استقرا دا مفروضہ ایس اُتے لگدا اے:
\[
\AxiomC{}
\RightLabel{$\pi_1'$}
\Deduce$!B, \Gamma \fCenter \Delta', !B \lif !C, !C, !A$
\AxiomC{}
\RightLabel{$\pi_3'$}
\Deduce$!A, !B, \Gamma \fCenter \Delta', !B \lif !C, !C$
\RightLabel{\CutCS}
\BinaryInf$!B, \Gamma \fCenter \Delta', !B \lif !C, !C$
\DisplayProof
\]
قطع !!{formula}~$!A$ اے۔ ایس !!{proof} دا قطع درجہ اوہی اے تے قطع !!{height}،~$\pi$ توں گھٹ اے۔ ایس نال !!a{proof}~$\pi_4$، \Log{G3c} وچ \CutCS{} توں بناں،~$!B, \Gamma \fCenter \Delta', !B \lif
!C, !C$ دا ملدا اے۔ $\RightR{\lif}$ لاؤ تے !!a{proof} ملے گا $\Gamma \fCenter
\Delta', !B \lif !C, !B \lif !C$ دا:
\[
\AxiomC{}
\RightLabel{$\pi_4$}
\Deduce$!B, \Gamma \fCenter \Delta', !B \lif !C, !C$
\RightLabel{$\RightR{\lif}$}
\UnaryInf$\Gamma \fCenter \Delta', !B \lif !C, !B \lif !C$
\DisplayProof
\]
!!a{proof}~$\pi'$، $\Gamma \Sequent \Delta$ دا، یعنی $\Gamma
\fCenter \Delta', !B \lif !C$ دا، لین لئی \olref[seq][inv]{prop:G3c-cont-adm} (سکیڑن دی قبولیت) لاؤ۔

ہُن منّ لو $\pi$ ایویں مُکدا اے:
\[
\AxiomC{}
\RightLabel{$\pi_1$}
\Deduce$\lexists[x][!B(x)], \Gamma' \fCenter \Delta, !A$
\AxiomC{}
\RightLabel{$\pi_3$}
\Deduce$!A, !B(c), \Gamma' \fCenter \Delta$
\RightLabel{$\LeftR{\lexists}$}
\UnaryInf$!A, \lexists[x][!B(x)], \Gamma' \fCenter \Delta$
\RightSubproofLabel{$\pi_2$}
\RightLabel{\CutCS}
\BinaryInf$\Gamma' \fCenter \Delta$
\DisplayProof
\]
استقرا دا مفروضہ فیر ایس اُتے لاؤ:
\[
\AxiomC{}
\RightLabel{$\pi_1[!B(c) \Sequent {}]$}
\Deduce$\lexists[x][!B(x)], !B(c), \Gamma' \fCenter \Delta, !A$
\AxiomC{}
\LeftLabel{$\pi_3[\lexists[x][!B(x)] \Sequent {}]$}
\Deduce$!A, \lexists[x][!B(x)], !B(c), \Gamma' \fCenter \Delta$
\RightLabel{\CutCS}
\BinaryInf$\lexists[x][!B(x)], !B(c), \Gamma' \fCenter \Delta$
\RightLabel{\LeftR{\lexists}}
\UnaryInf$\lexists[x][!B(x)], \lexists[x][!B(x)], \Gamma' \fCenter \Delta$
\DisplayProof
\]
دھیان دیو، خاص متغیر دی شرط پوری اے، کیوں جے~\LeftR{\lexists} دی خاص متغیر والی شرط~$\pi_2$ وچ ایہ پکا کردی اے کہ $c$، $\lexists[x][!B(x)]$، $\Gamma'$ یا~$\Delta$ وچ نہیں آؤندا۔ آخر \olref[seq][inv]{prop:G3c-cont-adm} لاؤ۔

\paragraph{صورت D۔}
$\pi_1$ تے $\pi_2$ وچوں گھٹ توں گھٹ اک، دو مقدمیّاں والے قاعدے~$R$ اُتے مُکدا اے جیہدے وچ $!A$ مُکھ نہیں۔ 6 صورتاں نیں۔ اسیں اوہ صورت ویکھدے آں جتھے $!A$،~$\pi_2$ دے آخری استنباط وچ مُکھ نہیں، تے اوہ~$\RightR{\land}$ اُتے مُکدا اے۔ باقی صورتاں ایہو جیہیاں نیں۔

منّ لو $\pi$ ایویں مُکدا اے:
\[
\AxiomC{}
\RightLabel{$\pi_1$}
\Deduce$\Gamma \fCenter \Delta', !B \land !C, !A$
\AxiomC{}
\RightLabel{$\pi_3$}
\Deduce$!A, \Gamma \fCenter \Delta', !B$
\AxiomC{}
\RightLabel{$\pi_4$}
\Deduce$!A, \Gamma \fCenter \Delta', !C$
\RightLabel{$\RightR{\land}$}
\BinaryInf$!A, \Gamma \fCenter \Delta', !B \land !C$
\RightSubproofLabel{$\pi_2$}
\RightLabel{\CutCS}
\BinaryInf$\Gamma \fCenter \Delta$
\DisplayProof
\]
استقرا دا مفروضہ ایہناں !!{proof}s اُتے لگدا اے:
\[
\AxiomC{}
\RightLabel{$\pi_1'$}
\Deduce$\Gamma \fCenter \Delta', !B \land !C, !B, !A$
\AxiomC{}
\RightLabel{$\pi_3'$}
\Deduce$!A, \Gamma \fCenter \Delta', !B \land !C, !B$
\RightLabel{\CutCS}
\BinaryInf$\Gamma \fCenter \Delta', !B \land !C, !B$
\DisplayProof
\]
تے
\[
\AxiomC{}
\RightLabel{$\pi_1''$}
\Deduce$\Gamma \fCenter \Delta, !B \land !C, !C, !A$
\AxiomC{}
\RightLabel{$\pi_4'$}
\Deduce$!A, \Gamma \fCenter \Delta', !B \land !C, !C$
\RightLabel{\CutCS}
\BinaryInf$\Gamma \fCenter \Delta', !B \land !C, !C$
\DisplayProof
\]
ایتھے \Log{G3c} دے !!{proof}s $\pi_1'$ تے $\pi_1''$، !!{height} قائم رکھن والے کمزور کرن دے قاعدے (\olref[seq][adm]{prop:weak-G3c-adm}) نال~$\pi_1$ توں ملدے نیں (اک واری سجّے $!B$ نال تے اک واری $!C$ نال کمزور کر کے)۔ !!{proof}s $\pi_3'$ تے~$\pi_4'$، $\pi_3$ تے~$\pi_4$ توں وی !!{height} قائم رکھن والے کمزور کرن دے قاعدے نال، بالترتیب $\pi_3$ تے~$\pi_4$ توں، ملدے نیں (سجّے $!B \land !C$ نال کمزور کر کے)۔

استقرا دا مفروضہ !!{proof}s~$\pi_5$ تے~$\pi_6$، \Log{G3c} وچ، دیندا اے، $\Gamma \fCenter \Delta', !B \land !C, !B$ تے $\Gamma
\fCenter \Delta', !B \land !C, !C$ دے۔ ایہناں نوں جوڑ کے، !!a{proof} ملدا اے $\Gamma \Sequent \Delta', !B \land !C, !B \land !C$ دا، قاعدہ~$\RightR{\land}$ ورتدیاں:
\[
\AxiomC{}
\RightLabel{$\pi_5$}
\Deduce$\Gamma \fCenter \Delta', !B \land !C, !B$
\AxiomC{}
\RightLabel{$\pi_6$}
\Deduce$\Gamma \fCenter \Delta', !B \land !C, !C$
\RightLabel{\RightR{\land}}
\BinaryInf$\Gamma \fCenter \Delta, !B \land !C, !B \land !C$
\DisplayProof
\]
!!a{proof} لین لئی $\Gamma \Sequent \Delta$ دا، یعنی $\Gamma \fCenter
\Delta', !B \land !C$ دا، \olref[seq][inv]{prop:G3c-cont-adm} (سکیڑن دی قبولیت) لاؤ۔

\paragraph{صورت E۔ قطع گھٹاؤنے۔}
آخری صورت اوہ اے جتھے $\pi_1$ تے $\pi_2$ دوویں ایہو جیہے استنباطاں اُتے مُکدے نیں جیہناں وچ $!A$ مُکھ اے۔ چھے صورتاں نیں،~$!A$ دے ہر ممکن !!{main operator} لئی اک۔

ہر صورت وچ اسیں قطع !!{formula}~$!A$ والے اک \CutCS{} نوں~$!A$ دے ذیلی !!{formula}s اُتے اک یا ودھ قطعاں وچ بدلنے آں، یعنی اوہناں استنباطاں دے سرگردہ !!{formula}s اُتے جیہناں دا مُکھ !!{formula}~$!A$ اے۔ ایس لئی ایہناں صورتاں نوں اکثر قطع استنباط دی «تقلیل» یا «روپ بدلنا» آکھدے نیں۔

\begin{enumerate}
\item جے $!A \ident \lnot !B$، تاں $\pi$ ایویں مُکدا اے:
\[
\AxiomC{}
\RightLabel{$\pi_1'$}
\Deduce$!B, \Gamma \fCenter \Delta$
\RightLabel{\RightR{\lnot}}
\UnaryInf$\Gamma \fCenter \Delta, \lnot !B$
\LeftSubproofLabel{$\pi_1$}
\AxiomC{}
\RightLabel{$\pi_2'$}
\Deduce$\Gamma \fCenter \Delta, !B$
\RightLabel{\LeftR{\lnot}}
\UnaryInf$\lnot !B, \Gamma \fCenter \Delta$
\RightSubproofLabel{$\pi_2$}
\RightLabel{\CutCS}
\BinaryInf$\Gamma \fCenter \Delta$
\DisplayProof
\]
استقرا دے مفروضے نال \CutCS{} ایتھوں ہٹایا جا سکدا اے:
\[
\AxiomC{}
\RightLabel{$\pi_2'$}
\Deduce$\Gamma \fCenter \Delta, !B$
\AxiomC{}
\RightLabel{$\pi_1'$}
\Deduce$!B, \Gamma \fCenter \Delta$
\RightLabel{\CutCS}
\BinaryInf$\Gamma \fCenter \Delta$
\DisplayProof
\]

\item
جے $!A \ident (!B \lor !C)$، تاں $\pi$ ایویں مُکدا اے:
\[
\AxiomC{}
\RightLabel{$\pi_1'$}
\Deduce$\Gamma \fCenter \Delta, !B, !C$
\RightLabel{\RightR{\lor}}
\UnaryInf$\Gamma \fCenter \Delta, !B \lor !C$
\LeftSubproofLabel{$\pi_1$}
\AxiomC{}
\RightLabel{$\pi_2'$}
\Deduce$!B, \Gamma \fCenter \Delta$
\AxiomC{}
\RightLabel{$\pi_2''$}
\Deduce$!C, \Gamma \fCenter \Delta$
\RightLabel{\LeftR{\lor}}
\BinaryInf$!B \lor !C, \Gamma \fCenter \Delta$
\RightSubproofLabel{$\pi_2$}
\RightLabel{\CutCS}
\BinaryInf$\Gamma \fCenter \Delta$
\DisplayProof
\]
\CutCS اُتے مُکدا !!{proof} ویکھو:
\[
\AxiomC{}
\RightLabel{$\pi_1'$}
\Deduce$\Gamma \fCenter \Delta, !B, !C$
\AxiomC{}
\RightLabel{$\pi_2'''$}
\Deduce$!C, \Gamma \fCenter \Delta, !B$
\RightLabel{\CutCS}
\BinaryInf$\Gamma \fCenter \Delta, !B$
\DisplayProof
\]
جتھے $\pi_2'''$، $\pi_2''$ توں !!{height} قائم رکھن والے کمزور کرن دے قاعدے نال ملدا اے۔ ایہدا قطع درجہ $< \cutr{\pi}$ اے کیوں جے $\depth{!C} < \depth{!B \lor !C}$، سو استقرا دا مفروضہ \Log{G3c} وچ !!a{proof} دیندا اے، $\pi_1''$، $\Gamma \fCenter \Delta, !B$ دا۔ \CutCS اُتے مُکدا !!{proof}،
\[
\AxiomC{}
\RightLabel{$\pi_1''$}
\Deduce$\Gamma \fCenter \Delta, !B$
\AxiomC{}
\RightLabel{$\pi_2'$}
\Deduce$!B, \Gamma \fCenter \Delta$
\RightLabel{\CutCS}
\BinaryInf$\Gamma \fCenter \Delta$
\DisplayProof
\]
قطع درجہ $= \depth{!B} < \depth{!B \lor !C}$ رکھدا اے، سو استقرا دا مفروضہ فیر لگدا اے تے !!a{proof} $\pi'$ دیندا اے $\Gamma
\fCenter \Delta$ دا۔

\item $!A \ident (!B \land !C)$۔ مشق۔

\item جے $!A \ident (!B \lif !C)$، تاں $\pi$ ایویں مُکدا اے:
\[
\AxiomC{}
\RightLabel{$\pi_1'$}
\Deduce$!B, \Gamma \fCenter \Delta, !C$
\RightLabel{\RightR{\lif}}
\UnaryInf$\Gamma \fCenter \Delta, !B \lif !C$
\LeftSubproofLabel{$\pi_1$}
\AxiomC{}
\RightLabel{$\pi_2'$}
\Deduce$\Gamma \fCenter \Delta, !B$
\AxiomC{}
\RightLabel{$\pi_2''$}
\Deduce$!C, \Gamma \fCenter \Delta$
\RightLabel{\LeftR{\lif}}
\BinaryInf$!B \lif !C, \Gamma \fCenter \Delta$
\RightSubproofLabel{$\pi_2$}
\RightLabel{\CutCS}
\BinaryInf$\Gamma \fCenter \Delta$
\DisplayProof
\]
\CutCS اُتے مُکدا !!{proof}،
\[
\AxiomC{}
\RightLabel{$\pi_1'$}
\Deduce$!B, \Gamma \fCenter \Delta, !C$
\AxiomC{}
\RightLabel{$\pi_2''$}
\Deduce$!C, \Gamma \fCenter \Delta$
\RightLabel{\CutCS}
\BinaryInf$!B, \Gamma \fCenter \Delta$
\DisplayProof
\]
~$\pi$ توں گھٹ قطع درجہ رکھدا اے۔ استقرا دا مفروضہ !!a{proof}~$\pi_1''$ دیندا اے $!B, \Gamma \fCenter \Delta$ دا۔
ہُن \CutCS اُتے مُکدا !!{proof} ویکھو:
\[
\AxiomC{}
\RightLabel{$\pi_2'$}
\Deduce$\Gamma \fCenter \Delta, !B$
\AxiomC{}
\RightLabel{$\pi_1''$}
\Deduce$!B, \Gamma \fCenter \Delta$
\RightLabel{\CutCS}
\BinaryInf$\Gamma \fCenter \Delta$
\DisplayProof
\]
ایہدا قطع درجہ وی~$\pi$ دے قطع درجے توں گھٹ اے، سو استقرا دا مفروضہ \Log{G3c} دا !!{proof}~$\pi'$ دیندا اے $\Gamma \fCenter
\Delta$ دا۔
\item جے $!A \ident \lforall[x][!B(x)]$، تاں $\pi$ ایویں مُکدا اے:
\[
\AxiomC{}
\RightLabel{$\pi_1'$}
\Deduce$\Gamma \fCenter \Delta, !B(c)$
\RightLabel{\RightR{\lforall}}
\UnaryInf$\Gamma \fCenter \Delta, \lforall[x][!B(x)]$
\LeftSubproofLabel{$\pi_1$}
\AxiomC{}
\RightLabel{$\pi_2'$}
\Deduce$!B(t), \lforall[x][!B(x)], \Gamma \fCenter \Delta$
\RightLabel{\LeftR{\lforall}}
\UnaryInf$\lforall[x][!B(x)], \Gamma \fCenter \Delta$
\RightSubproofLabel{$\pi_2$}
\RightLabel{\CutCS}
\BinaryInf$\Gamma \fCenter \Delta$
\DisplayProof
\]
استقرا دے مفروضے نال \CutCS{} ایتھوں ہٹایا جا سکدا اے:
\[
\AxiomC{}
\RightLabel{$\pi_1''$}
\Deduce$!B(t), \Gamma \fCenter \Delta, \lforall[x][!B(x)]$
\AxiomC{}
\RightLabel{$\pi_2'$}
\Deduce$!B(t), \lforall[x][!B(x)], \Gamma \fCenter \Delta$
\RightLabel{\CutCS}
\BinaryInf$!B(t), \Gamma \fCenter \Delta$
\DisplayProof
\]
جتھے $\pi_1''$، $\pi_1$ توں ملدا اے (~$\pi_1'$ توں نہیں!) !!{height} قائم رکھن والے کمزور کرن دے قاعدے نال۔ (!!{proof} دا قطع درجہ اوہی اے پر قطع !!{height} گھٹ اے۔) ایس نال !!a{proof}~$\pi_2''$ ملدا اے $!B(t), \Gamma \fCenter \Delta$ دا۔

!!{proof}~$\pi_1'$ دا آخری سیکوئنٹ $\Gamma \Sequent \Delta,
!B(c)$ اے، جتھے $c$ اک~$\RightR{\lforall}$ استنباط دا خاص متغیر اے۔ سو $c$،~$!B(x)$، $\Gamma$ یا~$\Delta$ وچ نہیں آؤندا۔ اسیں منّنے آں کہ !!{proof}~$\pi$ باقاعدہ اے، سو $c$ بس اگلے~$\RightR{\lforall}$ استنباط دا خاص متغیر اے تے بس اوہدے اُتّے آؤندا اے، یعنی بس~$\pi_1'$ وچ آؤندا اے، تے~$\pi_1'$ وچ کسے استنباط دا خاص متغیر نہیں۔ چونکہ $c$،~$\pi_2$ وچ نہیں آؤندا، ایہ~$t$ وچ وی نہیں آ سکدا۔ \olref[seq][qua]{lem:subst-regular} مطابق !!{proof}~$\Subst{\pi_1'}{t}{c}$ اک !!a{proof} اے $\Gamma \Sequent \Delta,
!B(t)$ دا۔ ہُن استقرا دا مفروضہ ایس !!{proof} اُتے لاؤ:
\[
\AxiomC{}
\RightLabel{$\Subst{\pi_1'}{t}{c}$}
\Deduce$\Gamma \fCenter \Delta, !B(t)$
\AxiomC{}
\RightLabel{$\pi_2''$}
\Deduce$!B(t), \Gamma \fCenter \Delta$
\RightLabel{\CutCS}
\BinaryInf$\Gamma \fCenter \Delta$
\DisplayProof
\]
(استقرا دا مفروضہ ایس لئی لگدا اے کہ قطع !!{formula}~$!B(t)$ اے، سو !!{proof} دا قطع درجہ~$\pi$ توں گھٹ اے۔)
\item $!A \ident \lexists[x][!B(x)]$ والی صورت مشق وجوں چھڈدے آں۔
\end{enumerate}
\end{proof}

\begin{prob}
\olref{lem:cut-adm-G3c} دے ثبوت وچ صورت~D دی اوہ ذیلی صورت پرکھو جتھے $!A$،~$\pi_1$ دے آخری استنباط وچ مُکھ نہیں، تے اوہ~$\LeftR{\lif}$ اُتے مُکدا اے۔ (دھیان دیو، مشق دا اک حصہ ایہ صاف دَسنا اے کہ تسی کیہ ثابت کرنا اے، یعنی ایس صورت وچ $\pi$ دی کیہڑی شکل اے۔)
\end{prob}

\begin{prob}
\olref{lem:cut-adm-G3c} دے ثبوت وچ صورت~D دی اوہ ذیلی صورت پرکھو جتھے $!A$،~$\pi_1$ دے آخری استنباط وچ مُکھ نہیں، تے اوہ~$\LeftR{\lforall}$ اُتے مُکدا اے۔
\end{prob}

\begin{prob}
\olref{lem:cut-adm-G3c} دے ثبوت وچ صورت~E دی اوہ ذیلی صورت پرکھو جتھے $!A$، $\pi_1$ تے $\pi_2$ دوہاں دے آخری استنباطاں وچ مُکھ اے تے $!A \ident (!B \land !C)$۔
\end{prob}

\begin{prob}
\olref{lem:cut-adm-G3c} دے ثبوت وچ صورت~E دی اوہ ذیلی صورت پرکھو جتھے $!A$، $\pi_1$ تے $\pi_2$ دوہاں دے آخری استنباطاں وچ مُکھ اے تے $!A \ident \lexists[x][!B(x)]$۔
\end{prob}

دھیان دیو، صورت~E وچ \CutCS{} اُتے مُکّن والے اوس !!{proof} دی قطع !!{height}، جیہدے اُتے اسیں استقرا دا مفروضہ لاؤنے آں، اصل ثبوت دی قطع !!{height} توں ودھ ہو سکدی اے۔ مثلاً $!A
\ident (!B \lif !C)$ والی صورت وچ اسیں $\pi$ نوں !!a{proof} وچ بدلیا جیہدے وچ دو \CutCS{} استنباط نیں:
\[
\AxiomC{}
\RightLabel{$\pi_2'$}
\Deduce$\Gamma \fCenter \Delta, !B$
\AxiomC{}
\RightLabel{$\pi_1'$}
\Deduce$!B, \Gamma \fCenter \Delta, !C$
\AxiomC{}
\RightLabel{$\pi_2''$}
\Deduce$!C, \Gamma \fCenter \Delta$
\RightLabel{\CutCS}
\BinaryInf$!B, \Gamma \fCenter \Delta$
\RightSubproofLabel{$\pi_1''$}
\RightLabel{\CutCS}
\BinaryInf$!B, \Gamma \fCenter \Delta$
\DisplayProof
\]
اُتلے \CutCS{} دا، $\pi_1'$ تے $\pi_2''$ دے وچکار، قطع درجہ تے قطع !!{height} دوویں~$\pi$ توں گھٹ نیں۔ پر !!{proof}~$\pi_1''$، جیہڑا استقرا دا مفروضہ لاؤن نال ملیا، دی قطع !!{height} ہُن~$\pi$ توں گھٹ نہیں۔ کیوں جے تھلّے والے \CutCS{} دا قطع !!{formula} $!B$ اے، اوہدا قطع \emph{درجہ}، فیر وی،~$\pi$ دے قطع درجے توں گھٹ اے۔

\olref{lem:cut-adm-G3c} قائم کردا اے کہ اسیں اکّو سبھ توں اُتلا \CutCS{} استنباط !!a{proof} وچوں ہٹا سکدے آں۔ اسیں ایہ وی وکھا سکدے آں کہ کِنّے وی \CutCS{} استنباط !!a{proof} وچوں ہٹا سکدے آں، ایس معاون قضیے نوں وار واری سبھ توں اُتلے ذیلی !!{proof}s اُتے لا کے جیہڑے~\CutCS اُتے مُکدے نیں۔

\begin{thm}\ollabel{thm:cut-elim-G3c-topmost}
کوئی وی !!{proof} $\pi$، $\Log{G3c} + \CutCS$ وچ،~\Log{G3c} دے !!a{proof} وچ بدلیا جا سکدا اے۔
\end{thm}

\begin{proof}
~$n$ اُتے استقرا نال، جیہڑی~$\pi$ وچ \CutCS{} استنباطاں دی گنتی اے۔ جے $n=0$، کجھ کرنا نہیں، $\pi$ پہلاں ای~\Log{G3c} دا !!a{proof} اے۔ نہیں تاں سبھ توں اُتلے \CutCS{} استنباط اُتے مُکدا ذیلی !!{proof}~$\pi'$ ویکھو۔ اوہدے وچ آخری استنباط وجوں اکّو \CutCS{} قاعدہ اے۔ \olref{lem:cut-adm-G3c} لا کے اوہنوں \CutCS{} توں بناں ثبوت~$\pi''$ وچ بدلو تے ذیلی !!{proof}~$\pi'$ دی تھاں~$\pi''$ رکھو۔ حاصل !!a{proof} اے جیہدے وچ $n-1$ \CutCS{} استنباط نیں، تے استقرا دا مفروضہ لگدا اے۔
\end{proof}
\endgroup
% /OLP-0656
% OLP-0655 content/proof-theory/cut-elimination/ce-largest.tex
\begingroup
\def\olfilename{ce-largest}
\expandafter\def\csname ol@sectioncs\endcsname{section}

\olfileid{pt}{cut}{inv}

\olsection{سبھ توں اُچّے درجے دے قطع ہٹاؤنے}

\olref[seq][inv]{lem:G3c-invert} دے ثبوت نے وکھایا کہ جے \Log{G3c} کسے منطقی قاعدے دا نتیجہ !!{prove}s کرے (\RightR{\lexists} تے \LeftR{\lforall} توں چھُٹ)، تاں ایہ قاعدے دے ہر مقدمے نوں وی !!{prove}s کرے گا، ثبوت دی !!{height} ودھائے بغیر۔ اسیں ایہ وی وکھا سکدے آں کہ ایہ گل $\Log{G3c} +
\CutCS$ لئی وی سچ اے۔

پہلاں وانگ، \CutCS{} استنباط دا \emph{قطع درجہ} اوہدے قطع !!{formula} دی ڈونگھائی نوں آکھدے آں۔

\begin{lem}\ollabel{lem:inv-G3c-cut}
جے $\Log{G3c} + \CutCS$ !!{prove}s کرے کسے منطقی قاعدے~$R$ دا نتیجہ، جیہڑا \RightR{\lexists} یا \LeftR{\lforall} نہیں، اک !!a{proof}~$\pi$ نال، تاں ایہ ہر مقدمے نوں وی !!{prove}s کرے گا، !!{proof}~$\pi'$ نال (یا !!{proof}s~$\pi'$، $\pi''$ نال، ایس گل مطابق کہ قاعدے دا اک مقدمہ اے یا دو)۔ ہور ایہ کہ
(a)~$\pheight{\pi'}\le\pheight{\pi}$ (تے
$\pheight{\pi''}\le\pheight{\pi}$) تے (b)~\CutCS{} استنباطاں دی گنتی تے اوہناں دے درجے~$\pi'$ وچ (تے $\pi''$ وچ) اوہی نیں جیہڑے~$\pi$ وچ نیں۔
\end{lem}

\begin{proof}
\cref{pt:seq:inv:lem:G3c-invert,pt:seq:inv:lem:invert-quant} دے ثبوت ویکھو۔ اوہ ثبوت~$\pheight{\pi}$ اُتے استقرا نال سی۔ مُڈھلی صورت وچ اسیں اک مسلّمہ دوجے مسلّمے نال بدلیا سی۔ سو شرطاں (a) تے~(b) سوکھیاں ای پوریاں نیں۔ جے~$\pi$ دا آخری قاعدہ~$R$ سی، تاں اوہدے سدھے ذیلی !!{proof}s ای اوہ !!{proof}s~$\pi_i$ نیں جیہڑے سانوں چاہیدے نیں، تے (a) تے~(b) دوویں پوریاں نیں۔ ہور ساریاں صورتاں وچ اسیں استقرا دا مفروضہ~$\pi$ دے سدھے ذیلی !!{proof}s اُتے لایا، یعنی اوہناں اُتے جیہڑے آخری استنباط~$R$ دے مقدمے یا مقدمیّاں تک پہنچدے نیں، تے فیر نتیجیاں اُتے اوہی قاعدہ $R$ لایا۔ (a) تے~(b) نویں ثبوت دے سدھے ذیلی !!{proof}s لئی سچ نیں، سو (a) تے~(b) پورے ثبوت لئی سچ نیں۔ ہُن اوہ صورت پرکھنی باقی اے جتھے آخری استنباط~\CutCS اے۔ فیر بس اک مثال لینے آں۔ منّ لو $\pi$، نتیجے $!B \land !C, \Gamma \Sequent
\Delta$ دا !!a{proof} اے، جیہڑا \LeftR{\land} دا نتیجہ اے، تے~\CutCS اُتے مُکدا اے:
\[
\AxiomC{}
\RightLabel{$\pi_1$}
\Deduce$!B \land !C, \Gamma \fCenter \Delta, !A$
\AxiomC{}
\RightLabel{$\pi_2$}
\Deduce$!A, !B \land !C, \Gamma \fCenter \Delta$
\RightLabel{\CutCS}
\BinaryInf$!B \land !C, \Gamma \fCenter \Delta$
\DisplayProof
\]
استقرا دا مفروضہ $\pi_1$ تے~$\pi_2$ اُتے لگدا اے (ایہ وی $\LeftR{\land}$ دے نتیجے دیاں مثالاں دے !!{proof}s نیں) تے بالترتیب !!{proof}s $\pi_1'$ تے~$\pi_2'$ دیندا اے، جیہڑے $!B, !C, \Gamma \fCenter
\Delta, !A$ تے $!A, !B, !C, \Gamma \fCenter \Delta$ دے نیں۔ ایہناں نوں \CutCS{} نال جوڑ کے~$\pi'$ ملدا اے:
\[
\AxiomC{}
\RightLabel{$\pi_1'$}
\Deduce$!B, !C, \Gamma \fCenter \Delta, !A$
\AxiomC{}
\RightLabel{$\pi_2'$}
\Deduce$!A, !B, !C, \Gamma \fCenter \Delta$
\RightLabel{\CutCS}
\BinaryInf$!B, !C, \Gamma \fCenter \Delta$
\DisplayProof
\]
صاف اے کہ (a) تے~(b) پوریاں نیں۔
\end{proof}

دھیان دیو، جے \Cut{} ورتدے \CutCS دی تھاں، تاں ایہ کم نہیں سی کرنا۔ آخری !!{proof} دے نتیجے وچ $!B, !C, \Gamma$ دیاں دو نکلاں مقدم وچ تے $\Delta$ دیاں دو نکلاں تالی وچ ہوندیاں۔ سانوں سکیڑن دی قبولیت ورتنی پَینی سی، پر \olref[seq][inv]{prop:G3c-cont-adm} دا ثبوت اُلٹ سکّن والے معاون قضیے نوں ورتدا اے۔ ایس طرح دلیل چکر وچ پے جاندی۔

\cref{pt:cut:inv:lem:inv-G3c-cut} نال اسیں قطع ہٹاؤن دا اک ہور ثبوت دے سکدے آں۔ ایس ثبوت وچ اسیں \CutCS{} استنباطاں دیاں سبھ توں اُتلیاں موجودگیاں وار واری نہیں ہٹاؤندے، بلکہ \CutCS{} استنباطاں دیاں اوہ موجودگیاں ہٹاؤندے آں جیہناں دا درجہ سبھ توں اُچّا اے۔ یعنی اسیں !!a{proof}~$\pi$، $\Log{G3c} + \CutCS$ وچ، توں شروع کر کے اوہدے وچ کسے وی تھاں توں اک قطع ہٹاؤنے آں (ممکن اے اوہدی تھاں گھٹ درجے دے قطع رکھ کے)۔ ایہ کم کردے رہندے آں جد تک کوئی قطع نہیں رہندا۔ بس ایہ شرط اے کہ جیہڑے قطع نوں اسیں ہٹاؤندے آں، اوہدے اُتّے اوسے یا ہور اُچّے درجے دا کوئی قطع نہ ہووے۔

\begin{lem}\ollabel{lem:max-cut-red-G3c}
جے $\pi$ اک !!a{proof} اے، $\Log{G3c}+\CutCS$ وچ، جیہڑا~\CutCS{} دے درجے~$n$ والے استنباط اُتے مُکدا اے تے ہور بس~$\Log{G3c}$ دے قاعدے تے \CutCS{} دے درجے $<n$ والے استنباط ورتدا اے، تاں اوسے آخری سیکوئنٹ دا اک ثبوت~$\pi'$، $\Log{G3c} + \CutCS$ وچ، اے جیہدے وچ ہر \CutCS{} استنباط دا درجہ~$<n$ اے۔
\end{lem}

\begin{proof}
!!{proof}~$\pi$ ایویں مُکدا اے:
\[
\AxiomC{}
\RightLabel{$\pi_1$}
\Deduce$\Gamma \fCenter \Delta, !A$
\AxiomC{}
\RightLabel{$\pi_2$}
\Deduce$!A, \Gamma \fCenter \Delta$
\RightLabel{\CutCS}
\BinaryInf$\Gamma \fCenter \Delta$
\DisplayProof
\]
اسیں~$!A$ دی شکل مطابق صورتاں وکھ کردے آں۔

منّ لو $!A$ جوہری اے۔ دھیان دیو کہ~$\pi_1$ تے $\pi_2$ وچ سارے \CutCS{} استنباطاں دا درجہ $<
\depth{!A} = 0$ ہونا چاہیدا اے۔ ایس لئی~$\pi_1$ یا~$\pi_2$ وچ کوئی \CutCS{} استنباط نہیں ہو سکدا۔ اسیں وکھاؤنے آں کہ $\Gamma \Sequent \Delta$ دا قطع توں بناں ثبوت اے، $\pheight{\pi_2}$ اُتے استقرا نال۔ جے $\pheight{\pi_2} = 0$، تاں $!A,
\Gamma \Sequent \Delta$ اک مسلّمہ اے۔ جے $!A$ مُکھ فارمولا نہیں، تاں کوئی جوہری $!C$،~$\Gamma$ تے $\Delta$ دوہاں دا !!a{element} اے، تے $\Gamma
\Sequent \Delta$ آپ اک مسلّمہ اے۔ جے $!A$ مُکھ اے، تاں
$\Delta = \Delta, !A$۔ ایس صورت وچ $\pi_1$، $\Gamma \Sequent
\Delta', !A, !A$ اُتے مُکدا اے۔ \olref[seq][inv]{prop:G3c-cont-adm} (سکیڑن دی قبولیت) لا کے $\Gamma \Sequent
\Delta', !A$ دا قطع توں بناں ثبوت ملدا اے۔ جے $\pheight{\pi_2} > 0$، تاں ایہ قاعدے~$R$ اُتے مُکدا اے۔ اوس قاعدے دا اک مقدمہ لو، جیہدی شکل $!A, \Gamma',
\Pi \Sequent \Delta', \Lambda$ اے، جتھے $\Pi$ تے $\Lambda$،~$R$ دے سرگردہ !!{formula}s نیں۔ \olref[seq][adm]{prop:weak-G3c-adm} نوں~$\pi_1$ تے $\pi_2$ اُتے لاؤن نال !!{proof}s ملدے نیں $\Gamma, \Gamma',
\Pi \Sequent \Delta, \Delta', \Lambda, !A$ تے $!A, \Gamma, \Gamma',
\Pi \Sequent \Delta, \Delta', \Lambda$ دے، جیہناں اُتے استقرا دا مفروضہ لگدا اے۔ ایس نال !!a{proof} ملدا اے $\Gamma, \Gamma', \Pi
\Sequent \Delta, \Delta', \Lambda$ دا،~$R$ دے ہر مقدمے لئی۔ $R$ لاؤن نال $\Gamma, \Gamma \Sequent \Delta, \Delta$ ملدا اے تے \olref[seq][inv]{prop:G3c-cont-adm} توں قطع توں بناں ثبوت ملدا اے $\Gamma \Sequent \Delta$ دا۔

جے $!A$ جوہری نہیں، تاں~$!A$ دے !!{main operator} مطابق صورتاں وکھ کردے آں۔ ہر صورت وچ \olref{lem:inv-G3c-cut} لا کے کھبّے تے سجّے قاعدیاں دے اوہناں مقدمیّاں دے !!{proof}s لیندے آں جیہناں دا مُکھ فارمولا $!A$ اے۔ فیر \olref[top]{lem:cut-adm-G3c} دے ثبوت دی صورت~E وانگ اگے ودھدے آں۔
\begin{enumerate}
\item $!A \ident \lnot !B$۔ !!{proof}s $\pi_1$ تے $\pi_2$ بالترتیب $\Gamma \Sequent \Delta, \lnot !B$ تے $\lnot !B, \Gamma \Sequent
\Delta$ اُتے مُکدے نیں۔ \olref{lem:inv-G3c-cut} نال !!{proof}s $\pi_1'$ تے $\pi_2'$، $!B, \Gamma \Sequent
\Delta$ تے $\Gamma \Sequent \Delta, !B$ دے نیں۔ !!{proof}~$\pi'$ ایہ اے:
\[
  \AxiomC{}
  \RightLabel{$\pi_2'$}
  \Deduce$\Gamma \fCenter \Delta, !B$
  \AxiomC{}
  \RightLabel{$\pi_1'$}
  \Deduce$!B, \Gamma \fCenter \Delta$
  \RightLabel{\CutCS}
  \BinaryInf$\Gamma \fCenter \Delta$
  \DisplayProof
  \]
\item $!A \ident (!B \lor !C)$۔ !!{proof}s $\pi_1$ تے $\pi_2$،
$\Gamma \Sequent \Delta, !B \lor !C$ تے $!B \lor !C, \Gamma
\Sequent \Delta$ اُتے مُکدے نیں۔ \olref{lem:inv-G3c-cut} نال (\RightR{\lor} دی اُلٹ، جیہڑی~$\pi_1$ اُتے لائی اے) !!a{proof} $\pi_1'$ اے $\Gamma \Sequent \Delta, !B, !C$ دا۔ \olref{lem:inv-G3c-cut} نال (\LeftR{\lor} دی اُلٹ، جیہڑی~$\pi_2$ اُتے لائی اے) !!a{proof} $\pi_2'$ اے $!B, \Gamma \Sequent \Delta$ دا، تے !!a{proof} $\pi_2''$ اے $!C, \Gamma \Sequent \Delta$ دا۔
\olref[seq][adm]{prop:weak-G3c-adm} نوں $\pi_2''$ اُتے لاؤن نال !!a{proof}~$\pi_2'''$ وی ملدا اے $!C, \Gamma \Sequent \Delta, !B$ دا۔ !!{proof}~$\pi'$ ایہ اے:
\[
  \AxiomC{}
  \RightLabel{$\pi_1'$}
  \Deduce$\Gamma \fCenter \Delta, !B, !C$
  \AxiomC{}
  \RightLabel{$\pi_2'''$}
  \Deduce$!C, \Gamma \fCenter \Delta, !B$
  \RightLabel{\CutCS}
  \BinaryInf$\Gamma \fCenter \Delta, !B$
  \AxiomC{}
  \RightLabel{$\pi_2'$}
  \Deduce$!B, \Gamma \fCenter \Delta$
  \RightLabel{\CutCS}
  \BinaryInf$\Gamma \fCenter \Delta$
  \DisplayProof
  \]
\item $!A \ident (!B \land !C)$۔ مشق۔
\item $!A \ident (!B \lif !C)$۔ !!{proof}s $\pi_1$ تے $\pi_2$،
$\Gamma \Sequent \Delta, !B \lif !C$ تے $!B \lif !C, \Gamma
\Sequent \Delta$ اُتے مُکدے نیں۔ \olref{lem:inv-G3c-cut} نال (\RightR{\lif} دی اُلٹ، جیہڑی~$\pi_1$ اُتے لائی اے) !!a{proof} $\pi_1'$ اے $!B, \Gamma \Sequent \Delta, !C$ دا۔ \olref{lem:inv-G3c-cut} نال (\LeftR{\lif} دی اُلٹ، جیہڑی~$\pi_2$ اُتے لائی اے) !!a{proof} $\pi_2'$ اے $\Gamma \Sequent \Delta, !B$ دا تے !!a{proof} $\pi_2''$ اے $!C, \Gamma \Sequent \Delta$ دا۔
\olref[seq][adm]{prop:weak-G3c-adm} نوں $\pi_2''$ اُتے لاؤن نال !!a{proof}~$\pi_2'''$ وی ملدا اے $!C, !B, \Gamma \Sequent \Delta$ دا۔ !!{proof}~$\pi'$ ایہ اے:
\[
  \AxiomC{}
  \RightLabel{$\pi_2'$}
  \Deduce$\Gamma \fCenter \Delta, !B$
  \AxiomC{}
  \RightLabel{$\pi_1'$}
  \Deduce$!B, \Gamma \fCenter \Delta, !C$
  \AxiomC{}
  \RightLabel{$\pi_2'''$}
  \Deduce$!C, !B, \Gamma \fCenter \Delta$
  \RightLabel{\CutCS}
  \BinaryInf$!B, \Gamma \fCenter \Delta$
  \RightLabel{\CutCS}
  \BinaryInf$\Gamma \fCenter \Delta$
  \DisplayProof
  \]
\item $!A \ident \lforall[x][!B(x)]$۔ !!{proof}s $\pi_1$ تے
$\pi_2$، $\Gamma \Sequent \Delta, \lforall[x][!B(x)]$ تے
$\lforall[x][!B(x)], \Gamma \Sequent \Delta$ اُتے مُکدے نیں۔
\olref{lem:inv-G3c-cut} نال !!a{proof}~$\pi_1'$ اے $\Gamma
\Sequent \Delta, !B(c)$ دا۔

ہُن !!{proof}~$\pi_2$ ویکھو۔ ایہ $\lforall[x][!B(x)],
  \Gamma \Sequent \Delta$ اُتے مُکدا اے۔ آخری سیکوئنٹ~$\pi_2$ توں اُتّے ول جاندیاں، $\lforall[x][!B(x)]$ دی ہر اوہ موجودگی نشان لاؤ جیہڑی کسے سیکوئنٹ دے کھبّے اے تے جیہدے توں آخری سیکوئنٹ وچ $\lforall[x][!B(x)]$ دی موجودگی~$\pi_2$ وچ آؤندی اے۔ یعنی:
(a)~$\lforall[x][!B(x)]$ نوں~$\pi_2$ دے آخری سیکوئنٹ وچ نشان لاؤ۔ (b)~جے
$\lforall[x][!B(x)]$ کسے قاعدے دے نتیجے دے سیاق وچ آؤندا اے تے نشان لگیا اے، تاں قاعدے دے مقدمے یا مقدمیّاں وچ اوہدی سانجھی موجودگی نوں نشان لاؤ۔ (c)~جے $\lforall[x][!B(x)]$، \LeftR{\lforall} قاعدے دے نتیجے وچ مُکھ اے تے نشان لگیا اے، تاں مقدمے وچ $\lforall[x][!B(x)]$ تے $!B(t)$ دیاں میل کھاندیاں سرگردہ موجودگیاں نوں نشان لاؤ۔

ہُن $\lforall[x][!B(x)]$ دیاں ایہ ساریاں نشان والیاں موجودگیاں ویکھو۔ ایہناں وچوں کوئی \LeftR{\lforall} توں ہور کسے قاعدے وچ سرگردہ نہیں (بس سرگردہ !!{formula}s اُتے نشان لگدا اے، \LeftR{\lforall} دے)۔ $\lforall[x][!B(x)]$ دیاں ساریاں نشان والیاں موجودگیاں $\pi_2$ وچوں مٹا دیو۔ جے ایہ درست !!{proof} اے، تاں $\lforall[x][!B(x)]$ دیاں نشان والیاں موجودگیاں کدے سرگردہ نہیں سن۔ اوہ ساریاں سدھیاں مسلّمیاں توں آئیاں سن۔ !!{proof}~$\Gamma \Sequent \Delta$ اُتے مُکدا اے تے اسیں $\pi$ دی تھاں ایہ رکھ سکدے آں۔ اسیں سبھ توں اُچّے درجے دا اک قطع ہٹا دِتّا اے۔

نہیں تاں ساڈے کول درست ثبوت نہیں، کیوں جے اسیں اوہ !!{formula}s $\lforall[x][!B(x)]$ مٹائے نیں جیہڑے کجھ \LeftR{\lforall} استنباطاں وچ سرگردہ سن۔ اسیں ایہناں نوں ایس شکل دے «استنباطاں» وچ بدل دِتّا اے:
\[
  \AxiomC{}
  \RightLabel{$\pi_t$}
  \Deduce$!B(t), \Pi \fCenter \Lambda$
  \RightLabel{\LeftR{\lforall}}
  \UnaryInf$\Pi \fCenter \Lambda$
  \DisplayProof
  \]
ہر سبھ توں اُتلے ایہو جیہے «استنباط» نوں ایہدے نال بدلو:
\[
  \AxiomC{}
  \RightLabel{$\Subst{\pi_1'}{t}{c}[\Pi \Sequent \Lambda]$}
  \Deduce$\Gamma, \Pi \fCenter \Delta, \Lambda, !B(t)$
  \AxiomC{}
  \RightLabel{$\pi_t[\Gamma \Sequent \Delta]$}
  \Deduce$!B(t), \Gamma, \Pi \fCenter \Delta, \Lambda$
  \RightLabel{\CutCS}
  \BinaryInf$\Gamma, \Pi \fCenter \Delta, \Lambda$
  \DisplayProof
  \]
تے ایہدے تھلّے ہر سیکوئنٹ دے کھبّے $\Gamma$ تے سجّے $\Delta$ شامل کرو۔ ایہ کم دُہراؤ جدوں تک سارے \LeftR{\lforall} «استنباط»، جیہناں دے مقدمے وچ $!B(t)$ اُتے نشان اے، کسے~$!B(t)$ اُتے \CutCS{} استنباط نال بدل نہیں جاندے۔ حاصل !!{proof} دا آخری سیکوئنٹ (تے ہُن ایہ درست !!{proof} اے) ایس شکل دا اے:
$\Gamma, \dots, \Gamma \Sequent \Delta, \dots, \Delta$۔ سکیڑن دی قبولیت لا کے ایہنوں~$\Gamma
  \Sequent \Delta$ دا !!a{proof} بناؤ تے $\pi$ دی تھاں ایہ رکھو۔ اسیں $\lforall[x][!B(x)]$ اُتے اک قطع دی تھاں $!B(t)$ اُتے (ممکن اے کئی) قطع رکھے نیں، پر ایہ سارے گھٹ درجے دے نیں۔ جیہڑے قطع پہلاں $\pi_1$ یا $\pi_2$ وچ سن، اوہ رہندے نیں، پر اوہ وی سارے گھٹ درجے دے نیں۔
\end{enumerate}
\end{proof}

\begin{prob}
$!A \ident (!B \land !C)$ والی صورت \olref{lem:inv-G3c-cut} دے ثبوت وچ پرکھو۔ یعنی وکھاؤ کہ جے اک !!a{proof}~$\pi$ اے جیہڑا ایتھے مُکدا اے:
\[
\AxiomC{}
\RightLabel{$\pi_1$}
\Deduce$\Gamma \fCenter \Delta, !B \land !C$
\AxiomC{}
\RightLabel{$\pi_2$}
\Deduce$!B \land !C, \Gamma \fCenter \Delta$
\RightLabel{\CutCS}
\BinaryInf$\Gamma \fCenter \Delta$
\DisplayProof
\]
تے $\pi_1$ تے $\pi_2$ وچ بس درجے $<\depth{!B
\land !C}$ والے قطع نیں، تاں اک !!a{proof}~$\pi'$ اے $\Gamma \Sequent \Delta$ دا، جیہدے وچ بس درجے $<\depth{!B \land !C}$ والے قطع نیں۔
\end{prob}

% \begin{prob}
% Verify the case where $!A \ident \lexists[x][!B(x)]$ in the proof of
% \olref{lem:inv-G3c-cut}.
% \end{prob}

\olref[top]{lem:cut-adm-G3c} قائم کردا اے کہ اسیں \CutCS{} دا سبھ توں اُچّے درجے والا استنباط، !!a{proof} وچ، گھٹ درجے دے \CutCS{} استنباطاں نال بدل سکدے آں۔ ایس معاون قضیے نوں فیر ورت کے وکھا سکدے آں کہ کِنّے وی \CutCS{} استنباط !!a{proof} وچوں ہٹا سکدے آں۔ ایہنوں وار واری سبھ توں اُتلے ذیلی !!{proof}s اُتے لاؤ جیہڑے سبھ توں اُچّے درجے دے~\CutCS اُتے مُکدے نیں۔

\begin{thm}\ollabel{thm:cut-elim-G3c-largest}
کوئی وی !!{proof} $\pi$، $\Log{G3c} + \CutCS$ وچ،~\Log{G3c} دے ثبوت وچ بدلیا جا سکدا اے۔
\end{thm}

\begin{proof}
پہلاں ثابت کردے آں کہ~$\pi$ وچ سبھ توں اُچّے درجے دے سارے قطع ہٹائے جا سکدے نیں۔ منّ لو $d$ سبھ توں اُچّا درجہ اے اوہناں قطعاں دا جیہڑے~$\pi$ وچ نیں، تے $n$ درجے~$d$ والے قطعاں دی گنتی اے۔ ثبوت~$n$ اُتے استقرا نال اے۔ جے $n=0$، تاں ثابت کرن نوں کجھ نہیں۔ جے $n>0$، تاں درجے~$d$ دا سبھ توں اُتلا قطع چُنو تے $\pi_1$ اوہ ذیلی !!{proof} ہووے جیہڑا ایس اُتے مُکدا اے۔ \olref{lem:max-cut-red-G3c} نال اک !!a{proof}~$\pi_1'$ اے اوسے آخری سیکوئنٹ دا، جیہدے سارے قطع درجے $<d$ دے نیں۔ $\pi_1$ نوں $\pi_1'$ نال~$\pi$ وچ بدلو۔ ایس نال اک !!a{proof} ملدا اے جیہدے وچ $n-1$ قطع درجے~$d$ دے نیں، تے استقرا دا مفروضہ لگدا اے۔

ہُن نتیجہ~$d$ اُتے استقرا نال نکلدا اے۔ جے $d=0$، تاں سارے قطع جوہری نیں تے پچھلے نتیجے نال سارے ہٹائے جا سکدے نیں۔ جے $d>0$، تاں پچھلے نتیجے توں ثبوت ملدا اے جیہدے وچ درجے $d$ دے سارے قطع درجے~$<d$ دے قطعاں نال بدلے گئے نیں۔ چونکہ $d$،~$\pi$ وچ قطعاں دا سبھ توں اُچّا درجہ سی، ساڈے کول ہُن ایہو جیہا ثبوت اے جیہدے وچ قطعاں دا سبھ توں اُچّا درجہ $<d$ اے، تے استقرا دا مفروضہ لگدا اے۔
\end{proof}
\endgroup
% /OLP-0655
% OLP-0661 content/proof-theory/cut-elimination/midsequent.tex
\begingroup
\def\olfilename{midsequent}
\expandafter\def\csname ol@sectioncs\endcsname{section}

\olfileid{pt}{cut}{mid}

\olsection{درمیانے سیکوئنٹ دا قضیہ تے ہیربرانڈ دا قضیہ}

قطع ہٹاؤن دے قضیے دا اک اہم نتیجہ درمیانے سیکوئنٹ دا قضیہ اے۔ ایہ کہندا اے کہ جے !!a{proof}~$\pi$،~$\Gamma \Sequent \Delta$ دا، اے جیہدے وچ ہر !!{formula} دے مسوّر شروع وچ نیں، تاں اک !!a{proof}~$\pi'$ اے جیہدے وچ سیکوئنٹ $S \ident \Pi \Sequent
\Lambda$ اے (جیہنوں \emph{درمیانا سیکوئنٹ} آکھدے آں)، ایویں کہ $S$ توں اُتّے~$\pi'$ وچ سارے استنباط قضیائی نیں تے~$S$ توں تھلّے سارے استنباط مسوّراں والے نیں۔ (!!^a{formula}~$!A$، \emph{مسوّر-مُڈھلا} (یا \emph{مسوّر-مُڈھلے روپ وچ}) اے جے اوہدی شکل
\[\mathsf{Q}_1x_1\dots\mathsf{Q}_nx_n\,!B(x_1,\dots,x_n)\]
ہووے، جتھے ہر $\mathsf{Q}_i$ یا $\lforall$ اے یا~$\lexists$۔) درمیانے سیکوئنٹ دے قضیے دا اک اہم نتیجہ ہیربرانڈ دا قضیہ اے۔ اوہدی عام صورت دَسنی ذرا اوکھی اے، پر بہت عام بولی وچ ایویں اے۔ پہلے درجے دی منطق دے ہر !!{provable} !!{sentence}~$!A$ لئی، اک قضیائی ہمیشہ سچا فارمولا~$!A'$ اے جیہدے توں $!A$ نوں !!{prove} کیتا جا سکدا اے۔اوہناں !!{sentence}s لئی جیہڑے~$\lexists[x][!B(x)]$ دی شکل دے نیں، جتھے $!B$ وچ کوئی مسوّر نہیں، ایہ صورت اے:

\begin{thm}[ہیربرانڈ دا قضیہ]
منّ لو $!B(x)$ وچ کوئی مسوّر نہیں تے $\Proves 
\lexists[x][!B(x)]$۔ تاں ٹرم~$t_1$، \dots، $t_n$ ایہو جیہے نیں کہ $\Proves !B(t_1) \lor \dots \lor !B(t_n)$۔
\end{thm}

!!{formula}~$!B(t_1) \lor \dots \lor !B(t_n)$ نوں~$\lexists[x][!B(x)]$ دی \emph{ہیربرانڈ والی فصلی عبارت} آکھدے آں۔ چونکہ $!B$ وچ مسوّر نہیں، ایس فصلی عبارت وچ وی نہیں۔ سو جے ایہ !!{provable} اے، تاں قطع توں بناں ثبوت نال !!{provable} اے، تے ایس لئی ایہو جیہے !!a{proof} نال جیہڑا بس قضیائی استنباط ورتدا اے۔ سو ایہ ہمیشہ سچا فارمولا اے۔

ہیربرانڈ دے قضیے دا اوکھا حصہ ایہ وکھاؤنا اے کہ ہر !!{provable}
!!{formula}~$\lexists[x][!B(x)]$ لئی ہیربرانڈ والی فصلی عبارت موجود اے۔ اُلٹی گل، یعنی جے $\lexists[x][!B(x)]$ دی ہیربرانڈ والی فصلی عبارت اے تاں ایہ ثابت ہو سکدا اے، سوکھی اے۔ اصل وچ $\lexists[x][!B(x)]$ ہیربرانڈ والی فصلی عبارت دا منطقی نتیجہ اے۔ مثلاً ایہ !!a{proof} اے~$\Log{G3c}$ وچ اوس صورت دا جتھے $n=2$:
\[
\Axiom$!B(t_1) \fCenter \lexists[x][!B(x)], !B(t_1), !B(t_2)$
\Axiom$!B(t_2) \fCenter \lexists[x][!B(x)], !B(t_1), !B(t_2)$
\RightLabel{\LeftR{\lor}}
\BinaryInf$!B(t_1) \lor !B(t_2) \fCenter \lexists[x][!B(x)], !B(t_1), !B(t_2)$
\RightLabel{\RightR{\lexists}}
\UnaryInf$!B(t_1) \lor !B(t_2) \fCenter \lexists[x][!B(x)], !B(t_2)$
\RightLabel{\RightR{\lexists}}
\UnaryInf$!B(t_1) \lor !B(t_2) \fCenter \lexists[x][!B(x)]$
\DisplayProof
\]

ہیربرانڈ والی فصلی عبارت کڈھن لئی، !!a{proof}~$\pi$ توں جیہڑا~$\Sequent
\lexists[x][!B(x)]$ دا اے، اوہ صورت ویکھو جتھے سارے $\RightR{\lexists}$ استنباط، قطع توں بناں !!{proof}~$\pi$ وچ، آخر اُتے نیں:
\[
\AxiomC{}
\RightLabel{$\pi_1$}
\Deduce$ \fCenter \lexists[x][!B(x)], !B(t_1), \dots, !B(t_n)$
\RightLabel{\RightR{\lexists}}
\UnaryInf$ \fCenter \lexists[x][!B(x)], !B(t_2), \dots, !B(t_n)$
\RightLabel{$(n-1) \times \RightR{\lexists}$}
\Deduce$ \fCenter \lexists[x][!B(x)]$
\DisplayProof
\]
جے ایہ سارے $\RightR{\lexists}$ استنباط نیں~$\pi$ وچ جیہناں وچ $\lexists[x][!B(x)]$ سرگردہ اے، تاں~$\pi_1'$ وچ ہور کوئی مسوّر والا استنباط نہیں۔ کیوں جے $!B(x)$ وچ ہور کوئی مسوّر نہیں تے آخری سیکوئنٹ دے کھبّے کوئی !!{formula}s نہیں۔ دوجے لفظاں وچ، $\fCenter \lexists[x][!B(x)], !B(t_1),
\dots, !B(t_n)$،~$\pi$ دا درمیانا سیکوئنٹ اے۔ نالے چونکہ درمیانے سیکوئنٹ توں اُتّے کوئی مسوّر والا استنباط نہیں، ایس لئی اوہدے اُتلے سارے سیکوئنٹاں وچ $\lexists[x][!B(x)]$ اے تے ایہ موجودگیاں نہ سرگردہ نیں نہ مُکھ۔ سو ایہناں نوں !!{proof}~$\pi_1$ وچوں مٹا سکدے آں، تے!!a{proof} ملدا اے~$\fCenter !B(t_1), \dots, !B(t_n)$ دا۔ $n-1$ واری $\RightR{\lor}$ لا کے !!a{proof} ملدا اے ہیربرانڈ والی فصلی عبارت~$\Sequent !B(t_1) \lor \dots \lor !B(t_n)$ دا۔

اوکڑ ایہ اے کہ قطع توں بناں !!{proof} وچ وی، جیہڑا $\Sequent \lexists[x][!B(x)]$ دا اے، اسیں ایہ نہیں منّ سکدے کہ سارے \RightR{\lexists} استنباط آخر اُتے اکّو جُٹ وچ نیں۔ ایہناں \RightR{\lexists} استنباطاں دے وچکار ایہو جیہے استنباط ہو سکدے نیں جیہناں وچ کوئی $!B(t_i)$ مُکھ اے۔ ایس لئی درمیانے سیکوئنٹ دا قضیہ چاہیدا اے۔

\begin{thm}[درمیانے سیکوئنٹ دا قضیہ]\ollabel{thm:midsequent}
جے اک \Log{G3c}-!!{proof}~$\pi$ اے $\Gamma \Sequent \Delta$ دا، جتھے سارے !!{formula}s، $\Gamma$ تے $\Delta$ وچ، مسوّر-مُڈھلے نیں، تاں اک !!a{proof}~$\pi'$ اے جیہدے وچ ایہو جیہا سیکوئنٹ $\Pi \Sequent \Lambda$ اے کہ اوہدے تھلّے سارے استنباط مسوّراں والے تے اُتّے سارے استنباط قضیائی نیں۔
\end{thm}

\begin{proof}
~$\pi$ وچ مسوّر والے استنباط دا \emph{ترتیبی عدد} اوہدے تھلّے والے قضیائی استنباطاں دی گنتی نوں آکھو، تے~$o(\pi)$، $\pi$ دا ترتیبی عدد، اوہدے سارے مسوّر والے استنباطاں دے ترتیبی عدداں دا جوڑ اے۔ جے $o(\pi) = 0$، تاں کسے مسوّر والے استنباط دے تھلّے کوئی قضیائی استنباط نہیں، تے کم ہو گیا۔ کیوں جے سارے مسوّر والے استنباطاں دا اک ای مقدمہ اے، ایس لئی بس اک سبھ توں اُتلا مسوّر والا استنباط اے۔ اوہدا مقدمہ درمیانا سیکوئنٹ اے۔

جے $o(\pi) > 0$، تاں سبھ توں تھلّا مسوّر والا استنباط چُنو جیہدا ترتیبی عدد $>0$ اے۔ چونکہ اوہدا ترتیبی عدد $>0$ اے، اوہدے تھلّے کوئی قضیائی استنباط ضرور اے۔ چونکہ اوہ سبھ توں تھلّا اے، ایس دا نتیجہ کسے مسوّر والے استنباط دا مقدمہ نہیں ہو سکدا۔ نہیں تاں اوہدے تھلّے والے دوجے مسوّر والے استنباط دے تھلّے وی کوئی قضیائی استنباط ہوندا۔ اسیں (کئی!) صورتاں وکھ کردے آں، ایس حساب نال کہ مسوّر والا استنباط کیہڑا اے تے اوہدے تھلّے قضیائی استنباط کیہڑا اے۔ چونکہ آخری سیکوئنٹ وچ بس مسوّر-مُڈھلے فارمولے نیں، مسوّر والے استنباط دا مُکھ فارمولا اگلے قضیائی استنباط وچ سرگردہ نہیں ہو سکدا۔ نہیں تاں اوس استنباط دے مُکھ فارمولے وچ اک مسوّر قضیائی جوڑ دے دائرے اندر ہوندا۔ ذیلی !!{formula} والی خاصیت نال ایہ ذیلی !!{formula}، آخری سیکوئنٹ دے کسے !!a{formula} دا، ہونا پینا سی۔ سو مسوّر والے استنباط دا مُکھ !!{formula} تھلّے والے قضیائی استنباط دے سیاق دا !!a{element} اے۔

ایہ دو مثالاں نیں۔

پہلاں منّ لو اک $\RightR{\lforall}$ اے جیہڑا $\LeftR{\lif}$ دے سجّے مقدمے اُتے مُکدا اے۔ فیر~$\pi$ دا متعلقہ ذیلی !!{proof}، ذیلی !!{proof}~$\pi_1$ وچ مُکدا اے۔
\[
\AxiomC{}
\RightLabel{$\pi_2$}
\Deduce$\Gamma' \fCenter \Delta', \lforall[x][!A(x)], !B$
\AxiomC{}
\RightLabel{$\pi_3$}
\Deduce$!C, \Gamma' \fCenter \Delta', !A(c)$
\RightLabel{\RightR{\lforall}}
\UnaryInf$!C, \Gamma' \fCenter \Delta', \lforall[x][!A(x)]$
\RightLabel{\LeftR{\lif}}
\BinaryInf$!B \lif !C, \Gamma' \fCenter \Delta', \lforall[x][!A(x)]$
\DisplayProof
\]
اسیں منّ سکدے آں کہ $\pi$ باقاعدہ اے، سو خاص متغیر~$c$،~$\pi_3$ توں باہر نہیں آؤندا۔ خاص کر ایہ $!B
\lif !C$ یا~$\pi_2$ وچ نہیں آؤندا۔ ہُن ثبوت~$\pi_1'$ ویکھو:
\[
\AxiomC{}
\RightLabel{$\pi_2'$}
\Deduce$\Gamma' \fCenter \Delta', \lforall[x][!A(x)], !A(c), !B$
\AxiomC{}
\RightLabel{$\pi_3$}
\Deduce$!C, \Gamma' \fCenter \Delta', \lforall[x][!A(x)], !A(c)$
\RightLabel{\LeftR{\lif}}
\BinaryInf$!B \lif !C, \Gamma' \fCenter \Delta', \lforall[x][!A(x)], !A(c)$
\RightLabel{\RightR{\lforall}}
\UnaryInf$!B \lif !C, \Gamma' \fCenter \Delta', \lforall[x][!A(x)], \lforall[x][!A(x)]$
\DisplayProof
\]
جتھے $\pi_2'$،~$\pi_2$ توں سجّے پاسے~$!A(c)$ نال کمزور کر کے ملیا اے۔ (ایس نال کوئی استنباط نہیں ودھدا، خاص کر کوئی مسوّر والا استنباط نہیں جیہڑا ترتیبی عدد بدل دے۔ نالے~$\pi_2$ دی کوئی خاص متغیر والی شرط نہیں ٹُٹدی، کیوں جے~$!A(c)$ وچلا کوئی ثابت~$\pi_2$ وچ خاص متغیر وجوں نہیں ورتیا جاندا۔) \RightR{\lforall} دی خاص متغیر والی شرط ہن وی پوری اے، کیوں جے $c$،~$!B \lif !C$ وچ نہیں آؤندا۔

سکیڑن دی قبولیت ورت کے !!a{proof}~$\pi_1''$ لیندے آں $!B \lif !C, \Gamma' \fCenter \Delta',
\lforall[x][!A(x)]$ دا۔ $\RightR{\Contraction}$ دی قبولیت دا ثبوت، $\lforall[x][!A(x)]$ لئی، تے ایس ثبوت وچ ورتیا $\RightR{\lforall}$ دی اُلٹ دا ثبوت، کسے استنباط دا ترتیبی عدد نہیں بدل دے۔ ودھ توں ودھ استنباط مٹاندے نیں۔ سو~$\pi_1''$ وچ~$\pi_1$ توں ودھ مسوّر والے استنباط نہیں، تے~$\pi_1''$ دے ہر مسوّر والے استنباط دے تھلّے اوہنے ای قضیائی استنباط نیں جِنّے~$\pi_1$ وچ اوہدے میل کھاندے استنباط دے تھلّے سن۔ آخرلا \RightR{\lforall} ایس توں چھُٹ اے۔ اوہدے تھلّے اک \LeftR{\lif}،~$\pi$ وچ، سی جیہڑا اسیں اوہدے اُتّے چُک دِتّا اے،~$\pi_1''$ وچ۔ $\pi_1$ نوں $\pi$ وچ $\pi_1''$ نال بدلو۔ حاصل !!{proof} دا ترتیبی عدد گھٹ اے، کیوں جے گھٹ توں گھٹ $\RightR{\lforall}$ استنباط دا ترتیبی عدد (گھٹ توں گھٹ)~$1$ گھٹیا اے۔ ہُن استقرا دا مفروضہ لاؤ۔

جتھے مسوّر والا استنباط \RightR{\lexists} اے، صورتاں ہور وی سوکھیاں نیں, کیوں جے سکیڑن دی لوڑ نہیں تے کسے خاص متغیر والی شرط دی فکر نہیں۔ مثلاً منّ لو $\RightR{\lexists}$ پچھوں $\RightR{\land}$ آؤندا اے (ایس واری کھبّے مقدمے وجوں)۔ متعلقہ ذیلی !!{proof}، $\pi_1$، اے:
\[
\AxiomC{}
\RightLabel{$\pi_2$}
\Deduce$\Gamma' \fCenter \Delta', \lexists[x][!A(x)], !A(t), !B$
\RightLabel{\RightR{\lexists}}
\UnaryInf$!\Gamma' \fCenter \Delta', \lexists[x][!A(x)], !B$
\AxiomC{}
\RightLabel{$\pi_3$}
\Deduce$\Gamma' \fCenter \Delta', \lexists[x][!A(x)], !C$
\RightLabel{\RightR{\land}}
\BinaryInf$\Gamma' \fCenter \Delta', \lexists[x][!A(x)], !B \land !C$
\DisplayProof
\]
$\pi_1$ نوں~$\pi$ وچ $\pi_1'$ نال بدلو:
\[
\AxiomC{}
\RightLabel{$\pi_2$}
\Deduce$\Gamma' \fCenter \Delta', \lexists[x][!A(x)], !A(t), !B$
\AxiomC{}
\RightLabel{$\pi_3'$}
\Deduce$\Gamma' \fCenter \Delta', \lexists[x][!A(x)], !A(t), !C$
\RightLabel{\RightR{\land}}
\BinaryInf$\Gamma' \fCenter \Delta', \lexists[x][!A(x)], !A(t), !B \land !C$
\RightLabel{\RightR{\lexists}}
\UnaryInf$\Gamma' \fCenter \Delta', \lexists[x][!A(x)], !B \land !C$
\DisplayProof
\]
جتھے $\pi_3'$، $\pi_3$ توں سجّے پاسے~$!A(t)$ نال کمزور کر کے ملیا اے۔ ایس کمزور کرن وچ~$!A(t)$ نوں بس ہر سیکوئنٹ دے سجّے شامل کیتا جاندا اے،~$\pi_3$ وچ۔ سو کسے مسوّر والے استنباط دا ترتیبی عدد نہیں بدلدا۔~$\pi'$ وچ مسوّر والے استنباطاں دے ترتیبی عدد اوہی نیں جیہڑے~$\pi$ وچ نیں، بس \RightR{\lexists} استنباط دا ترتیبی عدد، جیہدی تھاں~\RightR{\land} نال بدلی سی،~$1$ گھٹیا اے۔ استقرا دا مفروضہ لگدا اے۔
\end{proof}

\begin{prob}
\olref{thm:midsequent} دے ثبوت دیاں اوہ صورتاں دسو جیہناں وچ \LeftR{\lexists}،~\RightR{\lif} دے مقدمے اُتے مُکدا اے، تے جیہناں وچ \LeftR{\lforall}،~\LeftR{\lor} دے کھبّے مقدمے اُتے مُکدا اے۔
\end{prob}

\begin{proof}[ہیربرانڈ دے قضیے دا ثبوت]
منّ لو $\pi$، $\Sequent \lexists[x][!A(x)]$ اُتے مُکدا اے۔ \olref{thm:midsequent} نال $\pi$ نوں !!a{proof}~$\pi'$ وچ بدل سکدے آں جیہدے وچ درمیانا سیکوئنٹ~$S$ لازماً ایس شکل دا اے:
\[\Sequent
\lexists[x][!A(x)], !A(t_1), \dots, !A(t_n).\]
چونکہ $S$ توں اُتّے سارے استنباط قضیائی نیں، $\lexists[x][!A(x)]$،~$S$ توں اُتلے کسے استنباط وچ مُکھ نہیں ہو سکدا۔ چونکہ ایہ جوہری نہیں، ایہ مسلّمے وچ مُکھ نہیں ہو سکدا۔ چونکہ $\lexists[x][!A(x)]$، $!A(t_1)$ دا حقیقی ذیلی فارمولا نہیں ہو سکدا (چیتے کرو $!A(x)$ وچ کوئی مسوّر نہیں)، ایس لئی ایہ~$S$ توں اُتّے سرگردہ وی نہیں ہو سکدا۔ سو $\lexists[x][!A(x)]$ نوں~$S$ اُتے مُکّن والے ذیلی !!{proof} وچوں مٹاؤن نال درست !!{proof} ملدا اے $\Sequent !A(t_1), \dots, !A(t_n)$ دا۔ ایہدے توں !!a{proof} مل سکدا اے $\Sequent !A(t_1) \lor \dots \lor !A(t_n)$ دا، $n-1$ واری $\RightR{\lor}$ لا کے۔
\end{proof}

% \begin{prob}
% Find a Herbrand disjunction of $\lexists[x][\lforall[y][(!A(x) \lif
% !A(y))]]$ by finding a cut-free !!{proof} of $\Sequent
% \lexists[x][\lforall[y][(!A(x) \lif !A(y))]]$ and applying the
% transformation in \olref{thm:midsequent} to find a midsequent (if
% necessary).
% \end{prob}
\endgroup
% /OLP-0661
% OLP-0658 content/proof-theory/cut-elimination/interpolation.tex
\begingroup
\def\olfilename{interpolation}
\expandafter\def\csname ol@sectioncs\endcsname{section}

\olfileid{pt}{cut}{itp}

\olsection{مائےہارا دا معاون قضیہ تے کریگ دا درمیانی فارمولے والا قضیہ}

ولیم کریگ دا درمیانی فارمولے والا قضیہ کہندا اے کہ جے $!A
\Entails !B$، تاں اک !!a{sentence}~$!C$ (\emph{درمیانی جملہ}) ایہو جیہا اے کہ $!A \Entails !C$ تے $!C \Entails !B$۔ اہم گل ایہ اے کہ درمیانی جملہ~$!C$ ایویں چُن سکدے آں کہ اوہدے سارے !!{constant}s تے !!{predicate}s جیہڑے~$!C$ وچ آؤندے نیں،~$!A$ تے~$!B$ دوہاں وچ آؤن۔ (اسیں منّنے آں کہ کوئی !!{function}s شامل نہیں۔) ایس پابندی توں بناں $!A$ تے $!B$ سوکھے ای درمیانی جملے بن سکدے نیں۔

درمیانی فارمولے والا قضیہ کلاسیکی پہلے درجے دی منطق لئی سچ اے، تے ایہنوں لبھّن والیاں خاصیتاں وچوں «پہلے درجے دی منطق دی آخری اہم خاصیت» آکھیا گیا اے۔ نمونیاں دے نظریے تے خودکار استدلال وچ ایہدی اہم ورتوں اے۔ ایہدا مُڈھلا محرک تے ورتوں سائنس دے فلسفے وچ سی، کیوں جے ایہدے نال پرکھ سکدے آں کہ کِہڑیاں بنیادی اصطلاحاں نال کسے نظریے دے مسلّمے مقرر ہو سکدے نیں یا نہیں۔ ایس توں بیتھ دا تعریف پذیری والا قضیہ وی نکلدا اے۔ اوہ کہندا اے کہ جے کوئی نظریہ کسے تصور دی تعریف لُکویں طور اُتے کر سکدا اے، تاں اوہدی کھُلّھی تعریف وی کر سکدا اے۔

ریاضیاتی منطق دے کئی نتیجیاں وانگ، درمیانی فارمولے والا قضیہ نمونیاں دے نظریے تے ثبوتی نظریے، دوہاں دے طریقیاں نال ثابت ہو سکدا اے۔ ثبوتی طریقیاں دا فائدہ ایہ اے کہ اوہ درمیانی فارمولے دی موجودگی وکھاؤن دے نال اک الگورتھم وی دیندے نیں جیہڑا !!a{proof} توں درمیانی فارمولا کڈھدا اے، ایس گل دے ثبوت توں کہ $!A \Entails !B$۔ مثلاً !!a{proof}،~\Log{G3c} وچ، $!A \Sequent !B$ دا۔

منّ لو $\Lang L(!A)$ اوہناں !!{constant}s تے !!{predicate}s دا سیٹ اے جیہڑے~$!A$ وچ نیں، تے $\Lang L(\Gamma) = \bigcup \Setabs{\Lang L(!A)}{!A \in
\Gamma}$۔ سارے حصے وچ منّنے آں کہ !!{formula}s وچ !!{function}s نہیں۔

\begin{thm}[کریگ دا درمیانی فارمولے والا قضیہ]\ollabel{thm:interpolation}
!!{language}~$!A$ دی، $\Lang L(!A)$، تے $!B$ دی زبان، $\Lang L(!B)$ منّ لو۔ جے $\Log{G3c} \Proves !A \Sequent !B$، تاں اک !!a{formula}~$!C$، !!{language}~$\Lang L(!C) \subseteq \Lang
L(!A) \cap \Lang L(!B)$ وچ اے، ایویں کہ $\Log{G3c} \Proves !A \Sequent !C$ تے $\Log{G3c} \Proves !C \Sequent !B$۔
\end{thm}

اسیں درمیانی فارمولے والے قضیے دا عام روپ ثابت کردے آں تاں جو !!{proof}s دی بناوٹ ورت سکیے،~\Log{G3c} وچ۔ ایس لئی سوچو کہ سیکوئنٹاں دے مقدم تے تالیاں دو دو حصیاں وچ وکھ کیتیاں ہوئیاں نیں۔ اسیں $\maeh{\Gamma_1}{\Gamma_2} \Sequent
\maeh{\Delta_1}{\Delta_2}$ ایہ وکھاؤن لئی لکھدے آں کہ $\Gamma_1$ تے $\Delta_1$ پہلے حصے دے نیں تے $\Gamma_2$ تے $\Delta_2$ دوجے حصے دے۔ فیر درمیانی فارمولے والا قضیہ ہِٹھلے معاون قضیے توں جھٹ نکلدا اے۔

\begin{lem}[مائےہارا دا معاون قضیہ]\ollabel{lem:maehara}
جے $\Log{G3c}
\Proves \maeh{\Gamma_1}{\Gamma_2} \Sequent \maeh{\Delta_1}{\Delta_2}$، تاں اک !!a{formula}~$!C$، !!{language}~$\Lang L(!C)
\subseteq \Lang L(\Gamma_1,
\Delta_1) \cap \Lang L(\Gamma_2, \Delta_2)$ وچ اے، ایویں کہ $\Log{G3c}
\Proves \Gamma_1 \Sequent \Delta_1, !C$ تے $\Log{G3c} \Proves !C,
\Gamma_2 \Sequent \Delta_2$۔
\end{lem}

\begin{proof}
  منّ لو $\pi \Proves \maeh{\Gamma_1}{\Gamma_2} \Sequent
  \maeh{\Delta_1}{\Delta_2}$۔ دعویٰ $n=\pheight{\pi}$ اُتے استقرا نال ثابت کردے آں۔

جے $n=0$، تاں $\maeh{\Gamma_1}{\Gamma_2} \Sequent
  \maeh{\Delta_1}{\Delta_2}$ مسلّمہ اے تے کوئی جوہری !!{formula}~$!A$، $\Gamma_1, \Gamma_2$ تے $\Delta_1,
  \Delta_2$ دوہاں وچ آؤندا اے۔ چار صورتاں نیں۔ ایہ ایس گل اُتے نربھر اے کہ $!A$ کھبّے پاسے پہلے حصے وچ اے یا دوجے وچ، تے سجّے پاسے پہلے حصے وچ اے یا دوجے وچ۔
\begin{enumerate}
\item $!A \in \Gamma_1$ تے $!A \in \Delta_1$۔ سانوں ایہو جیہا~$!C$ چاہیدا اے کہ $\Gamma_1 \Sequent \Delta_1, !C$ تے $!C, \Gamma_2
    \Sequent \Delta_2$ دوویں !!{provable} ہون۔ پہلا پہلاں ای مسلّمہ اے،~$!C$ کوئی وی چُنیے، کیوں جے $!A$ کھبّے تے سجّے دوہاں پاسیاں اُتے اے۔ $!C, \Gamma_2
    \Sequent \Delta_2$ دا !!{provable} ہونا پکا کرن لئی ساڈے کول اکّو چُون اے، $!C \ident \lfalse$ منّ لئیے۔
\item $!A \in \Gamma_1$ تے $!A \in \Delta_2$۔ تاں $\Gamma_1
    \Sequent \Delta_1, !A$ تے $!A, \Gamma_2 \Sequent \Delta_2$ دوویں مسلّمے نیں، سو $!C \ident !A$ چُن سکدے آں۔
\item $!A \in \Gamma_2$ تے $!A \in \Delta_1$۔ تاں $!A, \Gamma_1
    \Sequent \Delta_1$ تے $\Gamma_2 \Sequent \Delta_2, !A$ مسلّمے ہون گے، پر اوہناں وچ~$!A$ غلط پاسیاں اُتے اے۔ سو $!A$ درمیانی فارمولا نہیں۔ پر $\RightR{\lnot}$ تے $\LeftR{\lnot}$ ورت کے اسیں !!{prove} کر سکدے آں $\Gamma_1
    \Sequent \Delta_1, \lnot!A$ تے $\lnot !A, \Gamma_2 \Sequent
    \Delta_2$۔
\item $!A \in \Gamma_2$ تے $!A \in \Delta_2$۔ مشق۔
\end{enumerate}
اوہ صورتاں مشقاں وجوں چھڈیاں نیں جیہناں وچ $\maeh{\Gamma_1}{\Gamma_2} \Sequent
  \maeh{\Delta_1}{\Delta_2}$ ایس لئی مسلّمہ اے کہ $\lfalse \in
  \Gamma_1$ یا $\lfalse \in \Gamma_2$۔

جے $n>0$، تاں $\pi$ منطقی استنباط اُتے مُکدا اے۔ اسیں ورتے گئے قاعدے تے مُکھ فارمولے دے حصے مطابق صورتاں وکھ کرنیاں نیں۔ ہر صورت وچ استقرا دا ایہ مفروضہ ورت سکدے آں کہ معاون قضیہ سارے !!{proof}s~$\pi_1$ لئی سچ اے جیہناں لئی $\pheight{\pi_1} < n$، بھاویں~$\pi_1$ دا آخری سیکوئنٹ دو حصیاں وچ کسے وی ڈھنگ نال ونڈیا ہووے۔ آؤ کجھ خاص صورتاں ویکھیے۔
\begin{enumerate}
\item $\pi$، $\RightR{\lnot}$ اُتے مُکدا اے، مُکھ فارمولا~$\lnot
    !A$ اے۔ دو صورتاں نیں۔ $\lnot !A$ پہلے یا دوجے حصے وچ ہووے، ایس حساب نال $\pi$ ایہناں وچوں اک اُتے مُکدا اے:
\[
    \AxiomC{}
    \RightLabel{$\pi_1$}
    \Deduce$\maeh{!A, \Gamma_1}{\Gamma_2} \fCenter
      \maeh{\Delta_1}{\Delta_2}$
    \RightLabel{\RightR{\lnot}}
    \UnaryInf$\maeh{\Gamma_1}{\Gamma_2} \fCenter
      \maeh{\Delta_1, \lnot !A}{\Delta_2}$
    \DisplayProof
    \]
یا
\[
    \AxiomC{}
    \RightLabel{$\pi_1$}
    \Deduce$\maeh{\Gamma_1}{!A, \Gamma_2} \fCenter
      \maeh{\Delta_1}{\Delta_2}$
    \RightLabel{\RightR{\lnot}}
    \UnaryInf$\maeh{\Gamma_1}{\Gamma_2} \fCenter
      \maeh{\Delta_1}{\Delta_2, \lnot !A}$
    \DisplayProof
    \]
دھیان دیو، جے مُکھ !!{formula}~$\lnot !A$ پہلے حصے وچ اے، تاں اسیں سرگردہ !!{formula}~$!A$ نوں وی مقدمے دے پہلے حصے وچ رکھیا اے۔ ایہ چُون اسیں کر سکدے آں۔ استقرا دا مفروضہ اوہدوں وی لگ جاندا جے اسیں سرگردہ فارمولا دوجے حصے وچ رکھدے، پر فیر درمیانی فارمولا نہیں لبھ سکدے (مشق۔ ایہ کر کے ویکھو)۔

پہلاں پہلی صورت ویکھو، جتھے $\lnot !A$ پہلے حصے وچ اے۔ استقرا دے مفروضے نال~$\pi_1$ دے آخری سیکوئنٹ دا اک درمیانی فارمولا، یعنی !!a{formula}~$!C_1$، اے، جیہدے لئی دوویں
\begin{align*}
\formula A, \Gamma_1 & \Sequent \Delta_1, !C_1 &&\text{تے} &
\formula C_1, \Gamma_2 & \Sequent \Delta_2
\intertext{ثابت ہو سکدے نیں۔ سانوں ایہو جیہا !!a{formula}~$!C$ چاہیدا اے کہ ہِٹھلے دوویں !!{provable} ہون:}
\Gamma_1 & \Sequent \Delta_1, \lnot !A, !C &&\text{تے} &
\formula C, \Gamma_2 & \Sequent \Delta_2
\end{align*}
صاف اے کہ $!C \ident !C_1$ لے سکدے آں۔ استقرا دا مفروضہ پہلاں ای دسدا اے کہ سجّا سیکوئنٹ !!{provable} اے، تے کھبّا سیکوئنٹ اوس سیکوئنٹ توں~$\RightR{\lnot}$ نال نکلدا اے جیہڑا استقرا دے مفروضے توں ملیا سی۔ استقرا دا مفروضہ ایہ وی دسدا اے کہ $\Lang L(!C_1) \subseteq \Lang L(!A, \Gamma_1,
    \Delta_1) \cap \Lang L(\Gamma_2, \Delta_2)$۔ چونکہ $\Lang L(\lnot
    !A) = \Lang L(!A)$ تے $\Lang L(!C) = \Lang L(!C_1)$، ایس لئی
$\Lang L(!C_1) \subseteq \Lang L(\Gamma_1, \Delta_1, \lnot !A) \cap \Lang
    L(\Gamma_2, \Delta_2)$۔

دوجی صورت ذرا وکھری اے، کیوں جے $\lnot !A$ دوجے حصے وچ اے۔ مقدمے وچ سرگردہ !!{formula} نوں وی دوجے حصے وچ رکھدے آں تے استقرا دا مفروضہ ایہناں اُتے لاؤنے آں:
\begin{align*}
\Gamma_1 & \Sequent \Delta_1, !C_1 &&\text{تے} &
\formula C_1, !A, \Gamma_2 & \Sequent \Delta_2
\intertext{ثابت ہو سکدے نیں۔ جے دوجے اُتے \RightR{\lnot} قاعدہ لائیے، تاں ہِٹھلیاں دے !!{proof}s ملدے نیں:}
\Gamma_1 & \Sequent \Delta_1, !C_1 &&\text{تے} &
\formula C_1, \Gamma_2 & \Sequent \Delta_2, \lnot !A
\end{align*}
جے $!C_1$ درمیانی فارمولا اے، تاں عین ایہو سیکوئنٹ سانوں !!{provable} چاہیدے نیں، سو فیر $!C \ident !C_1$ لئیے۔ پہلاں وانگ $\Lang
    L(!C_1) \subseteq \Lang L(\Gamma_1, \Delta_1) \cap \Lang L(\Gamma_2,
    \Delta_2, \lnot !A)$۔
\item $\pi$، $\RightR{\lif}$ اُتے مُکدا اے، مُکھ !!{formula}~$!A
    \lif !B$ اے۔ فیر دو صورتاں نیں، ایس گل مطابق کہ $!A \lif
    !B$ پہلے یا دوجے حصے وچ اے۔ !!{proof}~$\pi$ ایہناں وچوں اک اُتے مُکدا اے:
\[
    \AxiomC{}
    \RightLabel{$\pi_1$}
    \Deduce$\maeh{!A, \Gamma_1}{\Gamma_2} \fCenter
      \maeh{\Delta_1, !B}{\Delta_2}$
    \RightLabel{\RightR{\lif}}
    \UnaryInf$\maeh{\Gamma_1}{\Gamma_2} \fCenter
      \maeh{\Delta_1, !A \lif !B}{\Delta_2}$
    \DisplayProof
    \]
یا
\[
    \AxiomC{}
    \RightLabel{$\pi_1$}
    \Deduce$\maeh{\Gamma_1}{!A, \Gamma_2} \fCenter
      \maeh{\Delta_1}{\Delta_2, !B}$
    \RightLabel{\RightR{\lif}}
    \UnaryInf$\maeh{\Gamma_1}{\Gamma_2} \fCenter
      \maeh{\Delta_1}{\Delta_2, !A \lif !B}$
    \DisplayProof
    \]
فیر اسیں سرگردہ !!{formula}s نوں مقدمے دے اوسے حصے وچ رکھیا اے جیہدے وچ نتیجے دا مُکھ !!{formula} اے۔ پہلی صورت وچ استقرا دا مفروضہ لاؤنے آں تے ایہناں دے !!{proof}s لیندے آں:
\begin{align*}
\formula A, \Gamma_1 & \Sequent \Delta_1, !B, !C_1 &&\text{تے} & \formula C_1, \Gamma_2 & \Sequent \Delta_2
\intertext{کھبّا سیکوئنٹ \RightR{\lif} لئی ڈھکدا مقدمہ اے۔ سو ایہ ہِٹھلے !!{provable} سیکوئنٹ نیں:}
\Gamma_1 & \Sequent \Delta_1, !A \lif !B, !C_1 &&\text{تے} & \formula C_1, \Gamma_2 & \Sequent \Delta_2
\end{align*}
ایس دا مطلب اے کہ $!C_1$،~$\pi$ دے آخری سیکوئنٹ دا وی درمیانی فارمولا اے، تے اسیں اک واری فیر $!C \ident !C_1$ لے سکدے آں۔

دوجی صورت ایس طرح ای کردی اے۔
\item $\pi$، $\LeftR{\lif}$ اُتے مُکدا اے، مُکھ !!{formula}~$!A
    \lif !B$ اے۔ جے $!A \lif !B$ پہلے حصے وچ اے، تاں !!{proof}~$\pi$ ایتھے مُکدا اے:
\[
    \AxiomC{}
    \RightLabel{$\pi_1$}
    \Deduce$\maeh{\Gamma_1}{\Gamma_2} \fCenter
      \maeh{\Delta_1, !A}{\Delta_2}$
    \AxiomC{}
    \RightLabel{$\pi_2$}
    \Deduce$\maeh{!B, \Gamma_1}{\Gamma_2} \fCenter
      \maeh{\Delta_1}{\Delta_2}$
    \BinaryInf$\maeh{!A \lif !B, \Gamma_1}{\Gamma_2} \fCenter
      \maeh{\Delta_1}{\Delta_2}$
    \DisplayProof
    \]
استقرا دا مفروضہ ہُن چار !!{provable} سیکوئنٹ دیندا اے، دو کھبّے مقدمے دے درمیانی فارمولے $!C_1$ لئی تے دو دوجے مقدمے دے درمیانی فارمولے $!C_2$ لئی:
\begin{align*}
\Gamma_1 & \Sequent \Delta_1, !A, !C_1 &&\text{(1)} &
\formula C_1, \Gamma_2 & \Sequent \Delta_2 && \text{(2)}\\
\formula B, \Gamma_1 & \Sequent \Delta_1, !C_2 &&\text{(3)} &
\formula C_2, \Gamma_2 & \Sequent \Delta_2 && \text{(4)}
\end{align*}
سیکوئنٹ (1) تے (3)،~\LeftR{\lif} دے مقدمے بن سکدے نیں، جے کمزور کرن والا قاعدہ لا کے $!C_2$ نوں (1) دے سجّے تے $!C_1$ نوں (3) دے سجّے شامل کریئے:
\[
    \AxiomC{}
    \Deduce$\Gamma_1 \fCenter \Delta_1, !A, !C_1, !C_2$
    \AxiomC{}
    \Deduce$!B, \Gamma_1 \fCenter \Delta_1, !C_1, !C_2$
    \RightLabel{\LeftR{\lif}}
    \BinaryInf$!A \lif !B, \Gamma_1 \fCenter \Delta_1, !C_1, !C_2$
    \DisplayProof
    \]
$\RightR{\lor}$ لا کے ایہ سیکوئنٹ ملدا اے:
\[
!A \lif !B, \Gamma_1 \fCenter \Delta_1, !C_1 \lor !C_2
\]
ایس توں لگدا اے کہ $!C_1 \lor !C_2$ ایس صورت وچ ساڈا درمیانی فارمولا ہو سکدا اے۔ تے اسیں باقی سیکوئنٹ (2) تے~(4) توں دوجا لوڑ والا سیکوئنٹ !!{prove} کر وی سکدے آں:
\[
    \AxiomC{}
    \Deduce$!C_1, \Gamma_2 \fCenter \Delta_2$
    \AxiomC{}
    \Deduce$!C_2, \Gamma_2 \fCenter \Delta_2$
    \RightLabel{\LeftR{\lor}}
    \BinaryInf$!C_1 \lor !C_2, \Gamma_2 \fCenter \Delta_2$
    \DisplayProof
    \]
ایہ پرکھنا باقی اے کہ $!C_1 \lor !C_2$ ٹھیک زبان وچ اے۔ استقرا دے مفروضے توں ملدا اے کہ
\begin{align*}
\Lang L(!C_1) & \subseteq
\Lang L(\Gamma_1, \Delta_1, !A) \cap \Lang L(\Gamma_2, \Delta_2) &\text{تے}\\
\Lang L(!C_2) & \subseteq
\Lang L(!B, \Gamma_1, \Delta_1) \cap \Lang L(\Gamma_2, \Delta_2)
\intertext{چونکہ}
\Lang L(!A) & \subseteq \Lang L(!A \lif !B) &\text{تے}\\
\Lang L(!B) & \subseteq \Lang L(!A \lif !B)
\intertext{ایس لئی دوویں گلاں ملدیاں نیں}
\Lang L(!C_1) & \subseteq
\Lang L(\Gamma_1, \Delta_1, !A \lif !B) \cap \Lang L(\Gamma_2, \Delta_2) &\text{تے}\\
\Lang L(!C_2) & \subseteq
\Lang L(\Gamma_1, \Delta_1, !A \lif !B) \cap \Lang L(\Gamma_2, \Delta_2)
\intertext{تے ایس توں، $\Lang L (!C_1 \lor !C_2) = {}$}
\Lang L(!C_1) \cup \Lang L(!C_2) & \subseteq
\Lang L(\Gamma_1, \Delta_1, !A \lif !B) \cap \Lang L(\Gamma_2, \Delta_2),
\end{align*}
جیویں چاہیدا سی۔

جس صورت وچ $!A \lif !B$ دوجے حصے وچ اے، گل وکھری اے پر ڈھنگ ایہو جیہا اے۔ استقرا دا مفروضہ فیر چار !!{provable} سیکوئنٹ دیندا اے:
\begin{align*}
\Gamma_1 & \Sequent \Delta_1, !C_1 &&\text{(1)} &
\formula C_1, \Gamma_2 & \Sequent \Delta_2, !A && \text{(2)}\\
\Gamma_1 & \Sequent \Delta_1, !C_2 &&\text{(3)} &
\formula C_2, !B, \Gamma_2 & \Sequent \Delta_2 && \text{(4)}
\end{align*}
ہُن سیکوئنٹ (2) تے (4)، نال کمزور کرن نال شامل کجھ !!{formula}s،~\LeftR{\lif} دے مقدمے بن جاندے نیں:
\[
    \AxiomC{}
    \Deduce$!C_1, !C_2, \Gamma_2 \fCenter \Delta_2, !A$
    \AxiomC{}
    \Deduce$!B, !C_1, !C_2, \Gamma_2 \fCenter \Delta_2$
    \RightLabel{\LeftR{\lif}}
    \BinaryInf$!A \lif !B, !C_1, !C_2, \Gamma_2 \fCenter \Delta_2$
    \DisplayProof
    \]
اسیں فیر ایہنوں ایس جیہے سیکوئنٹ وچ بدلنا چاہندے آں جیہدے وچ اکّو !!{formula}، $!C_1$ تے $!C_2$ توں بنیا ہووے۔ پر ایس واری اوہ مقدم وچ چاہیدا اے (کیوں جے ہُن اسیں دوجا حصہ ویکھ رہے آں)۔ $\LeftR{\land}$ لا کے ایہ ہو جاندا اے۔ سانوں $!C \ident
    (!C_1 \land !C_2)$ درمیانی فارمولا لینا چاہیدا اے۔ چنگی گل ایہ اے کہ باقی سیکوئنٹ (1) تے (3) ورت کے دوجے حصے دا میل کھاندا سیکوئنٹ !!{prove} کر سکدے آں:
\[
    \AxiomC{}
    \Deduce$\Gamma_1 \fCenter \Delta_1, !C_1$
    \AxiomC{}
    \Deduce$\Gamma_1 \fCenter \Delta_1, !C_2$
    \RightLabel{\RightR{\land}}
    \BinaryInf$\Gamma_1 \fCenter \Delta_1, !C_1 \land !C_2$
    \DisplayProof
    \]
پہلاں وانگ، ساڈے کول $\Lang L(!C_1 \land !C_2) \subseteq \Lang
    L(\Gamma_1, \Delta_1) \cap \Lang L(\Gamma_2, \Delta_2, !A \lif
    !B)$ اے۔

\item $\pi$، $\RightR{\lforall}$ اُتے مُکدا اے، مُکھ !!{formula}~$\lforall[x][!A(x)]$ اے:
\[
    \AxiomC{}
    \RightLabel{$\pi_1$}
    \Deduce$\maeh{\Gamma_1}{\Gamma_2} \fCenter
    \maeh{\Delta_1, !A(c)}{\Delta_2}$
    \RightLabel{\RightR{\lforall}}
    \UnaryInf$\maeh{\Gamma_1}{\Gamma_2} \fCenter
      \maeh{\Delta_1, \lforall[x][!A(x)]}{\Delta_2}$
    \DisplayProof
    \]
یا
\[
    \AxiomC{}
    \RightLabel{$\pi_1$}
    \Deduce$\maeh{\Gamma_1}{\Gamma_2} \fCenter
      \maeh{\Delta_1}{\Delta_2, !A(c)}$
    \RightLabel{\RightR{\lforall}}
    \UnaryInf$\maeh{\Gamma_1}{\Gamma_2} \fCenter
      \maeh{\Delta_1}{\Delta_2, \lforall[x][!A(x)]}$
    \DisplayProof
    \]
پہلی صورت وچ استقرا دا مفروضہ !!{proof}s دیندا اے، $\Gamma_1 \Sequent \Delta_1, !A(c), !C_1$ تے $!C_1, \Gamma_2
    \Sequent \Delta_2$ دے۔ پہلے اُتے $\RightR{\lforall}$ لا کے $\Gamma_1 \Sequent \Delta_1, \lforall[x][!A(x)], !C_1$ لیندے آں۔ (خاص متغیر دی شرط پوری اے، کیوں جے اوہ پہلاں ای \RightR{\lforall} قاعدے وچ،~$\pi$ دے آخر اُتے، پوری سی۔) سو $!C_1$ درمیانی فارمولا ہووے گا جے ایہ شرط پوری اے کہ $\Lang
    L(!C_1) \subseteq \Lang L(\Gamma_1, \Delta_1, \lforall[x][!A(x)])
    \cap \Lang L(\Gamma_2, \Delta_2)$۔ استقرا دے مفروضے نال ساڈے کول $\Lang L(!C_1) \subseteq \Lang L(\Gamma_1,
    \Delta_1, !A(c)) \cap \Lang L(\Gamma_2, \Delta_2)$ اے۔ سو سانوں~$c$ بارے سوچنا پئے گا۔ کی $c$،~$!C_1$ وچ آ سکدا اے؟ نہیں۔ جے $c$، $!C_1$ وچ آوے، تاں دوویں $c \in \Lang L(\Gamma_1, \Delta_1,
    !A(c))$ (جیہڑا سچ اے) تے $c \in \Lang L(\Gamma_2,
    \Delta_2)$ ہون گے۔ پر فیر $c$، $\RightR{\lforall}$ استنباط دے نتیجے وچ،~$\pi$ اندر، آوے گا تے خاص متغیر دی شرط ٹُٹ جائے گی۔

دوجی صورت ایہو جیہی اے۔

\item $\pi$، $\RightR{\lexists}$ اُتے مُکدا اے، مُکھ !!{formula}~$\lexists[x][!A(x)]$ اے۔ فیر دو صورتاں نیں۔ اسیں اوہ صورت ویکھدے آں جتھے $\lexists[x][!A(x)]$ پہلے حصے وچ اے تے دوجی مشق وجوں چھڈدے آں۔

پہلی صورت وچ $\pi$ ایتھے مُکدا اے:
\[
    \AxiomC{}
    \RightLabel{$\pi_1$}
    \Deduce$\maeh{\Gamma_1}{\Gamma_2} \fCenter
      \maeh{\Delta_1, \lexists[x][!A(x)], !A(t)}{\Delta_2}$
    \RightLabel{\RightR{\lexists}}
    \UnaryInf$\maeh{\Gamma_1}{\Gamma_2} \fCenter
      \maeh{\Delta_1, \lexists[x][!A(x)]}{\Delta_2}$
    \DisplayProof
    \]
استقرا دا مفروضہ !!{proof}s دیندا اے، $\Gamma_1 \Sequent
    \Delta_1, \lexists[x][!A(x)], !A(t), !C_1$ تے $!C_1, \Gamma_2
    \Sequent \Delta_2$ دے۔ پہلے اُتے $\RightR{\lexists}$ لا کے $\Gamma_1 \Sequent \Delta_1, \lexists[x][!A(x)], !C_1$ ملدا اے۔ سو $!C_1$ فیر درمیانی فارمولا ہووے گا جے ایہ شرط پوری اے کہ
$\Lang L(!C_1) \subseteq \Lang L(\Gamma_1, \Delta_1,
    \lexists[x][!A(x)]) \cap \Lang L(\Gamma_2, \Delta_2)$۔ ایتھے استقرا دا مفروضہ ایہ دیندا اے کہ $\Lang L(!C_1)
    \subseteq \Lang L(\Gamma_1, \Delta_1, \lexists[x][!A(x)], !A(t))
    \cap \Lang L(\Gamma_2, \Delta_2)$۔ چیتے کرو کہ اسیں منّیا سی کہ کوئی !!{function}s نہیں۔ ایہدا مطلب اے کہ ٹرم~$t$ بس اک !!a{constant} ہو سکدا اے، یعنی $t \ident b$۔ جے $b \notin \Lang L(!C_1)$ یا $b \in \Lang L(\Gamma_1, \Delta_1, \lexists[x][!A(x)]) \cap
    \Lang L(\Gamma_2, \Delta_2)$، تاں فکر والی کوئی گل نہیں۔ اسیں $!C_1$ نوں درمیانی فارمولا لے سکدے آں۔ پر جے $b \in \Lang
    L(!C_1)$ تے $b \notin \Lang L(\Gamma_1, \Delta_1,
    \lexists[x][!A(x)]) \cap \Lang L(\Gamma_2, \Delta_2)$، تاں کیہ ہووے گا؟ پہلاں~$b$ دیاں ساریاں موجودگیاں،~$!C_1$ وچ، اک !!a{variable}~$y$ نال بدل کے، اسیں $!C_1 \ident \Subst{!C_1'}{b}{y}$ لکھ سکدے آں، جتھے $!C_1'$ وچ~$b$ نہیں۔ ہور، یا $b \notin
    \Lang L(\Gamma_1, \Delta_1, \lexists[x][!A(x)])$ یا $b \notin
    \Lang L(\Gamma_2, \Delta_2)$۔ بالترتیب $!C \ident
    \lforall[y][!C_1'(y)]$ یا $!C \ident \lexists[y][!C_1'(y)]$ لے سکدے آں۔ پہلی صورت وچ ساڈے کول ایہ اے:
\[
    \AxiomC{}
    \Deduce$\Gamma_1 \fCenter
      \Delta_1, \lexists[x][!A(x)], !A(b), !C_1'(b)$
    \RightLabel{\RightR{\lexists}}
    \UnaryInf$\Gamma_1 \fCenter \Delta_1, \lexists[x][!A(x)], !C_1'(b)$
    \RightLabel{\RightR{\lforall}}
    \UnaryInf$\Gamma_1 \fCenter
      \Delta_1, \lexists[x][!A(x)], \lforall[y][!C_1'(y)]$
    \DisplayProof
    \]
تے
\[
    \AxiomC{}
    \Deduce$!C_1'(b), \lforall[y][!C_1'(y)], \Gamma_2 \fCenter \Delta_2$
    \RightLabel{\LeftR{\lforall}}
    \UnaryInf$\lforall[y][!C_1'(y)], \Gamma_2 \fCenter \Delta_2$
    \DisplayProof
    \]
اُتلے سیکوئنٹاں دے !!{proof}s استقرا دے مفروضے توں ملدے نیں (دوجے وچ مقدم اندر $\lforall[y][!C_1'(y)]$ شامل کرن لئی کمزور کرن دی قبولیت لاؤنے آں)۔ \RightR{\lforall} دی خاص متغیر والی شرط پہلے !!{proof} وچ پوری اے کیوں جے $b \notin \Lang L(\Gamma_1, \Delta_1,
    \lexists[x][!A(x)])$، تے $!C_1'$ دی تعریف ایویں کیتی گئی سی کہ اوہدے وچ~$b$ نہ ہووے۔

دوجی صورت وچ ساڈے کول ایہ اے (استقرا دے مفروضے تے~$\lexists[y][!C_1'(y)]$ نال کمزور کرن دی قبولیت توں):
\[
    \AxiomC{}
    \Deduce$\Gamma_1 \fCenter
      \Delta_1, \lexists[x][!A(x)], !A(b), \lexists[y][!C_1'(y)], !C_1'(b)$
    \RightLabel{\RightR{\lexists}}
    \UnaryInf$\Gamma_1 \fCenter
      \Delta_1, \lexists[x][!A(x)], \lexists[y][!C_1'(y)], !C_1'(b)$
    \RightLabel{\RightR{\lexists}}
    \UnaryInf$\Gamma_1 \fCenter
      \Delta_1, \lexists[x][!A(x)], \lexists[y][!C_1'(y)]$
    \DisplayProof
    \]
تے
\[
    \AxiomC{}
    \Deduce$!C_1'(b), \Gamma_2 \fCenter \Delta_2$
    \RightLabel{\LeftR{\lexists}}
    \UnaryInf$\lexists[y][!C_1'(y)], \Gamma_2 \fCenter \Delta_2$
    \DisplayProof
    \]
\LeftR{\lexists} دی خاص متغیر والی شرط دوجے !!{proof} وچ پوری اے، کیوں جے $b \notin \Lang L(\Gamma_2,
    \Delta_2)$، تے $!C_1'$ دی تعریف ایویں کیتی گئی سی کہ اوہدے وچ~$b$ نہ ہووے۔
\end{enumerate}
باقی صورتاں ایہو جیہیاں نیں، اوہناں نوں مشقاں وجوں چھڈدے آں۔
\end{proof}

\begin{prob}
منّ لو ساڈے کول منطقی جوڑ $\lnand$ اے، ایہناں قاعدیاں نال:
\[
\Axiom$\Gamma \fCenter \Delta, !A$
\Axiom$\Gamma \fCenter \Delta, !B$
\RightLabel{\LeftR{\lnand}}
\BinaryInf$!A \lnand !B, \Gamma \fCenter \Delta$
\DisplayProof
\quad
\Axiom$!A, !B, \Gamma \fCenter \Delta$
\RightLabel{\RightR{\lnand}}
\UnaryInf$\Gamma \fCenter \Delta, !A \lnand !B$
\DisplayProof
\]
وکھاؤ کہ مائےہارا دا معاون قضیہ ہن وی سچ اے۔ \olref{lem:maehara} دے ثبوت دیاں اوہ صورتاں کرو جیہناں وچ $\pi$ ایہناں قاعدیاں وچوں کسے اُتے مُکدا اے۔
\end{prob}

مائےہارا دا معاون قضیہ قائم کر کے، اسیں ایہنوں کریگ دا درمیانی فارمولے والا قضیہ ثابت کرن لئی ورت سکدے آں۔

\begin{proof}[\olref{thm:interpolation} دا ثبوت]
منّ لو $!A \Sequent !B$ !!{provable} اے۔ \olref{lem:maehara} نوں $\maeh{!A}{\ } \Sequent \maeh{\ }{!B}$ اُتے لا کے درمیانی فارمولا~$!C$ لیندے آں۔ ساڈے کول $\Proves !A \Sequent !C$ تے $\Proves !C
  \Sequent !B$ نیں۔
\end{proof}

منّ لو $\Gamma$ نظریہ (یعنی !!{sentence}s دا سیٹ) اے، اک !!a{language}~$\Lang L$ وچ، جیہدے وچ $n$-تھاویں والا {محمول}~$R$ اے۔ اسیں آکھدے آں کہ $\Gamma$،~$R$ دی \emph{لُکویں تعریف کردا اے}، اُدوں تے صرف اُدوں جدوں
\begin{align*}
\Gamma, \Gamma'
& \Entails \lforall[x_1][\dots\lforall[x_n][(R(x_1,\dots,x_n) \liff
R'(x_1, \dots, x_n))]],
\intertext{جتھے $\Gamma'$ اوہ $\Gamma$ اے جیہدے وچ~$R$ دیاں ساریاں موجودگیاں~$R' \notin \Lang L$ نال بدلیاں گئیاں نیں۔ $\Gamma$،~$R$ دی \emph{کھُلّھی تعریف کردا اے} اُدوں تے صرف اُدوں جدوں اک ایہو جیہا~$!A(x_1, \dots, x_n)$ اے جیہدے وچ~$R$ نہیں، ایویں کہ}
\Gamma &\Entails
\lforall[x_1][\dots\lforall[x_n][(R(x_1,\dots,x_n) \liff !A(x_1,
\dots, x_n))]].
\end{align*}
صاف اے کہ جے $\Gamma$، $R$ دی کھُلّھی تعریف کردا اے، تاں~$R$ دی لُکویں تعریف کردا اے۔ اُلٹی گل صاف نہیں، پر فیر وی سچ اے:

\begin{lem}[بیتھ دا تعریف پذیری والا قضیہ]
جے $\Gamma$،~$R$ دی لُکویں تعریف کردا اے، تاں اوہدی کھُلّھی تعریف وی کردا اے۔
\end{lem}

\begin{proof}
منّ لو $\Gamma$،~$R$ دی لُکویں تعریف کردا اے۔ $R'$، $\Gamma'$ اُتلے وانگ لو تے $\Lang L' = \Lang L(\Gamma') = \Lang L \setminus \{R\}
  \cup \{R'\}$۔ مفروضے نال ساڈے کول
\begin{align*}
    \Gamma, \Gamma' &
    \Sequent \lforall[x_1][\dots\lforall[x_n][(R(x_1,\dots,x_n)
    \liff R'(x_1, \dots, x_n))]]
\intertext{منّ لو $c_1$، \dots،~$c_n$ اوہ !!{constant}s نیں جیہڑے~$\Gamma$ وچ نہیں آؤندے۔ اُلٹ سکّن والے معاون قضیے نال ساڈے کول}
    \Gamma, \Gamma', R(c_1,\dots,c_n) & \Sequent  R'(c_1, \dots, c_n)
\intertext{مائےہارا دا معاون قضیہ ایس اُتے لاؤ}
    \maeh{\Gamma, R(c_1,\dots,c_n)}{\Gamma'} & \Sequent
      \maeh{\ }{R'(c_1, \dots, c_n)}
\intertext{اک ایہو جیہا~$!A \ident !A(x_1, \dots, x_n)$ لبھّن لئی کہ}
    \Gamma, R(c_1,\dots,c_n) & \Sequent !A(c_1, \dots, c_n) \\
    \formula A(c_1, \dots, c_n), \Gamma' & \Sequent R'(c_1, \dots, c_n)
\intertext{دے !!{proof}s موجود ہون۔ چونکہ $\Lang L(!A) = \Lang L_1 \cap
  \Lang L_2 = \Lang L(\Gamma) \setminus \{R\}$، ایس لئی $!A$ وچ~$R$ نہیں۔ $R'$ نوں ہر تھاں~$R$ نال بدلو، دوجے سیکوئنٹ دے !!{proof} وچ، تاں !!a{proof} ملدا اے ایس دا}
    \formula A(c_1, \dots, c_n), \Gamma & \Sequent R(c_1, \dots, c_n)
\end{align*}
\RightR{\lif} تے \RightR{\lforall} لاؤ تاں !!a{proof} ملدا اے ایس دا:
\[
    \Gamma \Sequent \lforall[x_1][\dots\lforall[x_n][(R(x_1,\dots,x_n) \liff !A(x_1, \dots, x_n))]].
\]
ایہ وکھاؤندا اے کہ $!A(x_1, \dots, x_n)$، $!R$ دی کھُلّھی تعریف کردا اے،~$\Gamma$ وچ۔
\end{proof}

تیجا نتیجہ جیہڑا درمیانی فارمولے والے قضیے دے برابر اے (تے اوہدے وانگ مائےہارا دے معاون قضیے نال ثابت ہو سکدا اے)، ایہ قضیہ اے۔ جے دو بے تضاد نظرییاں~$\Gamma_1$، $\Gamma_2$ وچ اک پورا نظریہ $\Gamma$ شامل اے، زبان~$\Lang L(\Gamma) \subseteq \Lang L(\Gamma_1)
\cap \Lang L(\Gamma_2)$ وچ، تاں $\Gamma_1 \cup \Gamma_2$ بے تضاد اے۔

چیتے کرو کہ نظریہ اُدوں پورا اے جدوں اوہدی زبان دے ہر !!{sentence}~$!A$ لئی یا $\Gamma \Proves !A$ یا $\Gamma \Proves \lnot !A$۔ چونکہ $\Gamma_1$ تے~$\Gamma_2$،~$\Gamma$ دے بے تضاد وادھے نیں، ایس لئی $\Gamma$ لازماً بے تضاد اے۔ ایس توں نکلدا اے کہ جے $!A$ اک !!a{sentence} اے، زبان $\Lang L(\Gamma_1) \cap \Lang
L(\Gamma_2)$ وچ، تاں ایہ نہیں ہو سکدا کہ $\Gamma_1 \Entails !A$ تے $\Gamma_2
\Entails \lnot !A$۔ ایہ ویکھن لئی منّ لو ایہو جیہا~$!A$ اے۔ جے $\Gamma_1 \Entails !A$، تاں $\Gamma \Entails !A$۔ نہیں تاں،~$\Gamma$ دی تمامیت نال، $\Gamma \Entails \lnot !A$، تے چونکہ $\Gamma \subseteq \Gamma_1$، ایس لئی $\Gamma_1 \Entails \lnot !A$، سو $\Gamma_1$ متضاد ہو جاندا اے۔ ایس لئی $\Gamma \Entails !A$ تے چونکہ $\Gamma
\subseteq \Gamma_2$، تاں $\Gamma_2 \Entails !A$۔ سو $\Gamma \Entails/
\lnot !A$، نہیں تاں $\Gamma_2$ متضاد ہو جاندا۔

ایہو جیہا !!a{sentence}~$!A$، سانجھی !!{language}~$\Lang L(\Gamma_1) \cap \Lang L(\Gamma_2)$ وچ، نہ ہونا جیہدے لئی $\Gamma_1
\Entails !A$ تے $\Gamma_2 \Entails \lnot !A$، ثبوت لئی بس اینا ای چاہیدا اے۔ ہُن اسیں مائےہارا دے معاون قضیے نال ایہدا ثبوت ثبوتی نظریے دے طریقے نال دیندے آں۔

\begin{thm}[رابنسن دا سانجھی بے تضادی والا قضیہ]
منّ لو $\Gamma_1$، $\Gamma_2$ بے تضاد نظریے نیں تے~$!A$ ایہو جیہا کوئی نہیں، !!{language} $\Lang L(\Gamma_1) \cap \Lang L(\Gamma_2)$ وچ، جیہدے لئی $\Gamma_1 \Entails !A$ تے $\Gamma_2 \Entails \lnot !A$۔ تاں $\Gamma_1 \cup \Gamma_2$ بے تضاد اے۔
\end{thm}

\begin{proof}
اُلٹا فرض کرو کہ $\Gamma_1 \cup \Gamma_2$ متضاد اے۔ تاں سیکوئنٹ $\Gamma_1, \Gamma_2 \Sequent \ $ دا !!a{proof} اے۔ مائےہارا دا معاون قضیہ $\maeh{\Gamma_1}{\Gamma_2}
  \Sequent \maeh{\ }{\ }$ اُتے لاؤ، تاں !!a{sentence}~$!A$ ملدا اے، !!{language}~$\Lang L(\Gamma_1) \cap \Lang L(\Gamma_2)$ وچ، ایویں کہ $\Proves \Gamma_1 \Sequent !A$ تے $\Proves !A, \Gamma_2 \Sequent
  \quad$۔ پچھلے توں $\Proves \Gamma_2 \Sequent \lnot !A$ ملدا اے۔ پر اسیں منّیا سی کہ ایہو جیہا کوئی !!{sentence} نہیں۔
\end{proof}

اُتّے اسیں بناں دَسے منّیا سی کہ نظریے $\Gamma$، $\Gamma_1$، $\Gamma_2$ محدود نیں۔ محدود گواہی دے قضیے نال ایہ نتیجے لامتناہی نظرییاں لئی وی سچ نیں۔
\endgroup
% /OLP-0658

% OLP-0660 content/proof-theory/cut-elimination/intuitionistic.tex
\begingroup
\def\olfilename{intuitionistic}
\expandafter\def\csname ol@sectioncs\endcsname{section}
\olfileid{pt}{cut}{intuitionistic}
جے $!A \ident (!B \lif !C)$، تاں $\pi$ ایویں مُکدا اے:
\[
\AxiomC{}
\RightLabel{$\pi_1'$}
\Deduce$!B, \Gamma \fCenter \Delta, !C$
\RightLabel{\RightR{\lif}}
\UnaryInf$\Gamma \fCenter \Delta, !B \lif !C$
\LeftSubproofLabel{$\pi_1$}
\AxiomC{}
\RightLabel{$\pi_2'$}
\Deduce$!B \lif !C, \Pi \fCenter \Lambda, !B$
\AxiomC{}
\RightLabel{$\pi_2''$}
\Deduce$!C, \Pi \fCenter \Lambda$
\RightLabel{\LeftR{\lif}}
\BinaryInf$!B \lif !C, \Pi \fCenter \Lambda$
\RightSubproofLabel{$\pi_2$}
\RightLabel{\Cut}
\BinaryInf$\Gamma, \Pi \fCenter \Delta, \Lambda$
\DisplayProof
\]
\Cut اُتے مُکدا !!{proof}،
\[
\AxiomC{}
\RightLabel{$\pi_1'$}
\Deduce$!B, \Gamma \fCenter \Delta, !C$
\RightLabel{\RightR{\lif}}
\UnaryInf$\Gamma \fCenter \Delta, !B \lif !C$
\LeftSubproofLabel{$\pi_1$}
\AxiomC{}
\RightLabel{$\pi_2'$}
\Deduce$!B \lif !C, \Pi \fCenter \Lambda, !B$
\RightLabel{\Cut}
\BinaryInf$\Gamma, \Pi \fCenter \Delta, \Lambda, !B$
\DisplayProof
\]
قطع !!{height} وچ~$\pi$ توں گھٹ اے۔ استقرا دا مفروضہ !!a{proof}~$\pi_1''$ دیندا اے $\Gamma, \Pi \fCenter \Delta, \Lambda, !B$ دا۔
\Cut اُتے مُکدا !!{proof}،
\[
\AxiomC{}
\RightLabel{$\pi_1'$}
\Deduce$B!, \Gamma \fCenter \Delta, !C$
\AxiomC{}
\RightLabel{$\pi_2''$}
\Deduce$!C, \Pi \fCenter \Lambda$
\RightLabel{\Cut}
\BinaryInf$!B, \Gamma, \Pi \fCenter \Delta, \Lambda$
\DisplayProof
\]
قطع درجے تے قطع !!{height} دوہاں وچ اوہدے قطع درجے تے !!{height} توں گھٹ اے جیہڑا~$\pi$ دا اے، سو استقرا دا مفروضہ \Log{G3c}-!!{proof}~$\pi_2'''$ دیندا اے $!B, \Gamma, \Pi, \Pi \fCenter \Delta, \Lambda,
\Lambda$ دا۔ ہُن استقرا دا مفروضہ اوس !!{proof} اُتے لاؤ جیہڑا \Cut اُتے مُکدا اے:
\[
\AxiomC{}
\RightLabel{$\pi_1''$}
\Deduce$\Gamma, \Pi \fCenter \Delta, \Lambda, !B$
\AxiomC{}
\RightLabel{$\pi_2'''$}
\Deduce$!B, \Gamma, \Pi, \Pi \fCenter \Delta, \Lambda,
\Lambda$
\RightLabel{\Cut}
\BinaryInf$!B, \Gamma, \Pi \fCenter \Delta, \Lambda$
\DisplayProof
\]
(ایہدا قطع درجہ $< \cutr{\pi}$ اے) تاں !!a{proof} ملدا اے $\Gamma,
\Gamma, \Pi, \Pi, \Pi \fCenter \Delta, \Lambda, \Lambda, \Lambda$ دا۔ !!{proof} $\pi'$، $\Gamma, \Pi \fCenter \Delta, \Lambda$ دا، \olref[seq][inv]{prop:G3c-cont-adm} لاؤن نال ملدا اے۔
\endgroup
% /OLP-0660
\OLEndChapterHook
\endgroup
% /OLP-0657
% OLP-0681 content/proof-theory/proof-search/proof-search.tex
\begingroup
\def\olfilename{proof-search}
\expandafter\def\csname ol@sectioncs\endcsname{section}

\olchapter{pt}{ps}{ثبوت دی بھال}

% OLP-0680 content/proof-theory/proof-search/introduction.tex
\begingroup
\def\olfilename{introduction}
\expandafter\def\csname ol@sectioncs\endcsname{section}

\olfileid{pt}{ps}{int}

\olsection{جان پچھان}

!!{proof} نظاماں، جیویں \Log{G3c} یا \Log{N1i}، دے مسلّمے تے استنباط دے قاعدے طے کردے نیں کہ سیکوئنٹاں یا !!{formula}s دے کیہڑے درخت درست !!{proof}s منّے جاندے نیں۔ جے سیکوئنٹاں یا !!{formula}s دا ایہو جیہا درخت دِتّا ہووے (شاید نال ہور جانکاری وی ہووے، جیویں ہر قدم اُتے ورتے قاعدے، یا مفروضیاں تے اوہناں نوں خارج کرن والے استنباطی قاعدیاں دے لیبل)، تاں مسلّمیاں تے قاعدیاں دیاں تعریفاں دی مدد نال اسیں ویکھ سکدے آں کہ دِتّا درخت درست !!{proof} اے کہ نہیں۔ !!a{proof} \emph{لبھنا}، کسے دِتّے سیکوئنٹ دا یا دِتّے مفروضیاں توں دِتّے !!{formula} دا، وکھری گل اے تے عام طور اُتے اوکھی وی۔ ایہو \emph{ثبوت لبھّن} دا مسئلہ اے۔

ثبوت لبھّن دا اک بہت سدھا طریقہ اے۔ اسیں !!{formula}s نوں اک محدود حرف مالا دیاں نشانیاں دیاں لڑیاں سمجھ سکدے آں۔ سیکوئنٹ تے سیکوئنٹاں یا !!{formula}s دے درخت وی !!{formula}s دیاں لڑیاں سمجھے جا سکدے نیں (شاید درخت دیاں شاخاں وکھاؤن لئی خاص قوساں ورت کے)۔ مسلّمیاں تے قاعدیاں دی مدد نال اسیں مشینی طریقے نال جانچ سکدے آں کہ نشانیاں دی لڑی کسے درست !!{proof} نوں ظاہر کردی اے کہ نہیں۔ سو اسیں \emph{سارے ممکن} درست !!{proof}s مشینی طریقے نال بنا سکدے آں (اوہ حرف مالا جیہڑی !!{proof}s ظاہر کرن لئی ورتدے آں، اوہدیاں نشانیاں دیاں ساریاں لڑیاں بنا کے، تے ہر لڑی لئی جانچ کے کہ اوہ کسے درست !!{proof} نوں ظاہر کردی اے کہ نہیں)۔ سدھا طریقہ ایہ اے کہ سارے درست !!{proof}s بناؤندے جاؤ، جدوں تک اک ایہو جیہا !!{proof} نہ مل جاوے جیہڑا دِتّے سیکوئنٹ دا اے (یا دِتّے !!{formula} دا، جیہدے کھُلّے مفروضے دِتّے ہویاں وچوں ہون)۔

چنگا طریقہ اوہ اے جیہڑا ہتھ نال !!a{proof} لبھّن لئی ورتی دا اے۔ آخری سیکوئنٹ توں شروع کرو، !!a{formula} چُنو تے استنباط دا قاعدہ «پچھاں ول» لاؤ، تاں جو اک یا ودھ مقدمے بنن۔ جے اوہناں دے !!{proof}s ہون، تاں اوہ آخری استنباط نال رل کے دِتّے سیکوئنٹ دا !!a{proof} بنا دیندے۔ قدرتی اخذ دے نظام وچ !!a{proof} لبھّن لئی وی رلدا ملدا، پر ذرا اوکھا، طریقہ ورتی دا اے۔

پر ایس طریقے نال کامیابی دی ضمانت نہیں۔ جے غلط قاعدہ چُنو یا !!{formula}s نوں غلط ترتیب نال چُنو، تاں اگلا راہ بند ہو سکدا اے۔ طریقہ پورا متعین وی نہیں۔ ہر قدم اُتے چُنن لئی اک توں ودھ !!{formula} یا پچھاں ول لاؤن لئی اک توں ودھ قاعدے ہو سکدے نیں۔ \emph{ثبوت لبھّن دا الگورتھم} اوہ پورا متعین طریقہ اے جیہڑا !!a{proof} ضرور لبھ لیندا اے جے اوہ موجود ہووے، یعنی اوہدی کامیابی دی ضمانت اے۔

کجھ !!{proof} نظام ہورناں نالوں سوکھے ثبوت لبھّن والے الگورتھم دی گنجائش دیندے نیں۔ سیکوئنٹ نظاماں وچ !!{main operator} جیہڑا !!a{formula} دا اے، تے ایہ گل کہ اوہ کھبّے پاسے اے یا سجّے، طے کردے نیں کہ کِہڑا قاعدہ پچھاں ول لانا اے۔ مثال لئی جے تسی !!a{proof} لبھنا چاہنے او $\Gamma \Sequent \Delta, !A \land !B$ دا، جیہدے وچ $!A \land
!B$ سجّے پاسے اے، تاں \RightR{\land} نوں پچھاں ول لاؤ تے دو میل کھاندے مقدمے $ \Gamma \Sequent \Delta, !A$ تے $ \Gamma \Sequent \Delta, !B$ بناؤ۔ پر جے تہاڈا نظام \Log{G1c} اے، تاں سکڑاؤ ہون نال کم اوکھا ہو جاندا اے کیونکہ چُون لئی ہور راہ ملدے نیں۔ تسی \RightR{\Contraction} نوں پچھاں ول لا کے اوس دی تھاں مقدمہ $ \Gamma \Sequent \Delta, !A \land
!B, !A \land !B$ وی بنا سکدے او۔ فیر مقدمہ $ \Gamma \Sequent
\Delta, !A \land !B, !A \land !B, !A \land !B$، تے ایویں اگے۔ \Log{G2c} وچ ایہ ہور وی اوکھا اے۔ قاعدہ \RightR{\land} پچھاں ول لاؤن دا ایہ وی مطلب اے کہ $ \Gamma_1$، $ \Gamma_2$ دا اندازہ لانا پینا اے، جیہڑے $ \Gamma
= \Gamma_1 \cup \Gamma_2$ پورا کرن، تے $ \Delta_1$، $ \Delta_2$ دا، جیہڑے $ \Delta
= \Delta_1 \cup \Delta_2$ پورا کرن۔ فیر مقدمیّاں $ \Gamma_1 \Sequent \Delta_1, !A$ تے $ \Gamma_2 \Sequent \Delta_2, !B$ نال کوشش کرو۔ غلط اندازے نال اگلا راہ بند ہو سکدا اے تے تہانوں شروع ول مُڑ کے ہور $ \Gamma_1$، $ \Gamma_2$، $ \Delta_1$، $ \Delta_2$ چُننے پین گے۔ جے !!{formula} $!A \lor !B$ اے، تاں \RightR{\lor} دے دو ممکناں وچوں اک چُننا پینا اے، جیہڑا دو ممکن مقدمیّاں وچوں اک $ \Gamma \Sequent \Delta,
!A$ یا $ \Gamma \Sequent \Delta, !B$ بناندا اے۔ غلط چُون نال فیر پچھاں مُڑنا پینا اے۔

\Log{G3c} نظام وچ ایہ اوکڑاں نہیں۔ اوہدے وچ ساختی قاعدے نہیں، تے استنباطی قاعدے دی چُون !!{main operator} تے اوس پاسے نال طے ہوندی اے جیہدے اُتے !!{formula} آؤندا اے۔ ایس لئی \Log{G3c}، \Log{G1c} یا \Log{G2c} نالوں ثبوت لبھّن لئی ودھ راس آؤندا اے۔ کامیاب بھال نال !!a{proof}، \Log{G3c} وچ، ملے گا، تے اوس !!{proof} دا ترجمہ فیر \Log{G1c} یا \Log{G2c} (یا \Log{N1c} وچ وی) ہو سکدا اے۔ سو \Log{G3c} لئی ثبوت لبھّن دا الگورتھم تے \Log{G3c} !!{proof}s دا ہور نظاماں دے !!{proof}s وچ ترجمہ کرن والے الگورتھم ہون، تاں اوہناں نظاماں وچ ثبوت لبھّن دا مسئلہ وی حل ہو جاندا اے۔

سانوں کیویں پتا لگے کہ ثبوت لبھّن والے الگورتھم دی کامیابی دی ضمانت اے؟ ایہدے لئی وکھاؤنا پینا اے کہ جے الگورتھم !!a{proof} نہ لبھ سکے، تاں کوئی !!{proof} موجود نہیں ہو سکدا۔ آخر جے اسیں ماڑا الگورتھم چُنیے، تاں ایہو جیہیاں صورتاں ہو سکدیاں نیں جتھے اوہ ثبوت نہ لبھے بھاویں ثبوت موجود ہووے۔ سدھے الگورتھم نال ایہ اوکڑ نہیں، کیونکہ اوہ سارے ممکن !!{proof}s وچ بھال کردا اے۔ پر ثبوت لبھّن دا الگورتھم کارگر ہونا چاہیدا اے تے گھٹّو گھٹ کم کرنا چاہیدا اے۔

ثبوت لبھّن والے الگورتھم دی کامیابی دی ضمانت وکھاؤن دا مرکزی خیال ایہ اے کہ جے الگورتھم ناکام ہووے، تاں اوہدی ناکامی دے طریقے توں وکھایا جا سکے کہ سیکوئنٹ دا !!a{proof} ہو ای نہیں سکدا۔ ایہ وکھاؤن دا طریقہ اے کہ سیکوئنٹ معتبر نہیں۔ کلاسیکی پہلے درجے دی منطق وچ معتبر سیکوئنٹ $ \Gamma \Sequent \Delta$ اوہ نیں جیہناں لئی ہر !!{structure} وچ گھٹّو گھٹ اک !!{formula}، $\Gamma$ وچلا، جھوٹھا اے یا گھٹّو گھٹ اک !!{formula}،~$\Delta$ وچلا، سچا اے۔ \Log{G3c} \emph{صحت مند} اے، یعنی ایہ صرف معتبر سیکوئنٹ !!{prove} کردا اے۔ ایہ !!{proof}s اُتے استقرا نال ثابت ہوندا اے۔ مسلّمے صاف طور اُتے معتبر نیں، تے جے کسے قاعدے دے مقدمے معتبر ہون، تاں اوہدا نتیجہ وی معتبر اے۔ ثبوت لبھّن والے الگورتھم دی کامیابی دی ضمانت وکھاؤن لئی اسیں وکھاؤنے آں کہ جے !!a{proof} لبھّن دی بھال، $\Gamma \Sequent \Delta$ لئی، ناکام ہو جاوے، تاں اسیں !!a{structure} دی تعریف دے سکدے آں جیہدے وچ سارے !!{formula}s، $\Gamma$ والے، سچے نیں تے سارے !!{formula}s،~$\Delta$ والے، جھوٹھے نیں۔ ایس نال ایہ وی ثابت ہو جاندا اے کہ \Log{G3c} \emph{کامل} اے، یعنی جے $ \Gamma \Sequent \Delta$ معتبر اے، تاں اوہدا اک \Log{G3c}-!!{proof} اے۔ کیونکہ ثبوت لبھّن دی ہر ناکام بھال وکھاؤندی اے کہ $ \Gamma \Sequent \Delta$ غیر معتبر اے، معتبر سیکوئنٹ لئی بھال ضرور کامیاب ہونی اے۔
\endgroup
% /OLP-0680
% OLP-0683 content/proof-theory/proof-search/search-algorithm.tex
\begingroup
\def\olfilename{search-algorithm}
\expandafter\def\csname ol@sectioncs\endcsname{section}

\olfileid{pt}{ps}{sal}

\olsection{\Log{G3c} لئی بھال دا الگورتھم}

اسیں \Log{G3c} لئی بھال دا اک الگورتھم دَسّنے آں۔ ایہو اکّو ممکن الگورتھم نہیں، تے سبھ توں کارگر وی نہیں۔ پر ایہ درست اے۔ جے سیکوئنٹ $\Gamma \Sequent \Delta$ دا !!a{proof} اے، تاں ایہ آخر رُک جاندا اے تے !!a{proof} لبھ لیندا اے۔ ایہ خطا وی نہیں کھاندا۔ جے !!a{proof} لبھے بغیر رُک جاوے، یا کدے رُکے ای نہ، تاں $\Gamma
\Sequent \Delta$ شروع توں ای معتبر نہیں۔

الگورتھم دی اک ضروری خاصیت ایہ اے کہ ایہ پکا کردا اے کہ ہر !!{formula} جیہڑا $\Gamma \Sequent \Delta$ وچ، یا بھال دوران بنے کسے سیکوئنٹ وچ اے، اوہدی «تقلیل» ہووے۔ یعنی اوہنوں پچھاں ول لائے گئے استنباطی قاعدے دا مُکھ فارمولا بنایا جاوے، تے ایہ جِنّی واری لوڑ ہووے اُنی واری ہووے۔ ایس خاصیت نوں \emph{نرپکھتا} آکھدے آں۔

الگورتھم پڑاؤواں وچ کم کردا اے۔ ہر پڑاؤ دا حاصل سیکوئنٹاں دا اک درخت اے، تے اگلے پڑاؤ دا حاصل، اوہناں سارے سبھ توں اُتلے سیکوئنٹاں وچ !!a{formula} دی تقلیل کر کے بناندا اے، جیہڑے پہلاں ای مسلّمے نہیں۔

پڑاؤ $0$ اُتے $\Gamma \Sequent \Delta$ توں شروع کردے آں۔ نرپکھتا پکی کرن لئی اسیں !!{formula}s، $ \Gamma$ تے $ \Delta$ وچلے، نوں انڈیکس دے طور اُتے عدد دیندے آں، ایویں کہ وکھ وکھ !!{formula}s نوں وکھ وکھ عدد ملن۔ مثال لئی !!{formula}s نوں کھبّے توں سجّے ول گنّن بارے سوچو۔ عدد اُتلے نشان دے طور اُتے لکھدے آں۔ مثال لئی جے اسیں !!a{proof} لبھ رہے آں $!A, !B \Sequent !C$ دا، تاں پڑاؤ~$0$ دا حاصل $!A^1, !B^2 \Sequent !C^3$ اے۔ ایہو جیہے انڈیکس والے سیکوئنٹ وچ اُتلے نشان دے طور اُتے آؤن والا سبھ توں وڈّا~$n$ اوہدا \emph{سبھ توں وڈّا انڈیکس} اے (مثال وچ~$3$)۔

جے سیکوئنٹ وچ مسوّر نیں، تاں ایہ وی پتا ہونا چاہیدا اے کہ سجّے پاسے $ \lexists[x][!A(x)]$ یا کھبّے پاسے $ \lforall[x][!A(x)]$ دی تقلیل ویلے کِہڑے ٹرم ورتنے نیں۔ سانوں صرف !!{constant}s $ \Obj c_0$، $ \Obj c_1$، $ \Obj c_2$، \dots، توں بنے ٹرم چاہیدے نیں، جیہڑے~$\Gamma$ تے $ \Delta$ وچ آؤن والیاں کسے وی فنکشن نشانیاں نال بنن۔ اسیں منّنے آں کہ ایہناں سارے ٹرماں دا سیٹ کسے مقرر گنتی وار لڑی وچ دِتّا اے، تاں جو اسیں مثلاً «پہلے !!{constant} جیہڑا $ \Gamma \Sequent \Delta$ وچ نہیں آؤندا» دی گل کر سکیے۔ الگورتھم دے بنائے درخت وچ ہر انڈیکس والے سیکوئنٹ نال !!{constant}s دا اک سیٹ جُڑیا اے۔ پڑاؤ~$0$ اُتے سیکوئنٹ نوں دِتّا !!{constant}s دا سیٹ اوہ سارے !!{constant}s رکھدا اے جیہڑے اوہدے وچ آؤندے نیں (جے کوئی ہون)، تے $\{\Obj c_0\}$ اے جے کوئی !!{constant}s نہیں۔

منّ لو کہ الگورتھم نے پڑاؤ~$n$ اُتے انڈیکس والے سیکوئنٹاں دا درخت بنایا اے، نال !!{constant}s دے جُڑے سیٹ۔ پڑاؤ~$n+1$ اُتے اسیں ہر سبھ توں اُتلے سیکوئنٹ لئی ہِٹھ لکھیا کردے آں۔

جے !!a{formula}~$!A$ سیکوئنٹ دے کھبّے تے سجّے دونے پاسے آؤندا اے، یا کھبّے پاسے~$\lfalse$ اے، تاں اسیں کجھ نہیں کردے۔ ایہو جیہے سیکوئنٹاں دے !!{proof}s~\Log{G3c} وچ نیں، تے ایس خاص سبھ توں اُتلے سیکوئنٹ لئی ہور کم نہیں رہندا۔

جے سیکوئنٹ وچ کوئی غیر جوہری !!{formula} نہیں، تاں الگورتھم وچّے ای روک دیو۔ $\Gamma \Sequent \Delta$ دا کوئی ثبوت نہیں۔

نئیں تاں غیر جوہری !!{formula}~$!A$ چُنو جیہدا انڈیکس~$i$ سبھ توں چھوٹا اے۔ تاں ایہ سبھ توں اُتلا سیکوئنٹ یا $!A^i, \Pi
\Sequent \Lambda$ اے یا $ \Pi \Sequent \Lambda, !A^i$۔ $k$ نوں سیکوئنٹ دا سبھ توں وڈّا انڈیکس تے $C$ نوں اوہنوں دِتّا !!{constant}s دا سیٹ منّ لو۔

اسیں $!A^i$ دی «تقلیل» کردے آں، یعنی استنباطی قاعدے دے حساب نال نویں سبھ توں اُتلے سیکوئنٹ بناؤنے آں۔ ایہ قاعدہ !!{main operator}، جیہڑا~$!A$ دا اے، تے ایہ گل کہ $!A^i$ کھبّے پاسے اے یا سجّے، طے کردے نیں۔

\begin{enumerate}
  \item $!A \ident !B \land !C$ تے کھبّے پاسے آؤندا اے۔ $(!B \land !C)^i, \Pi \Sequent \Lambda$ توں اُتّے اک نواں سبھ توں اُتلا سیکوئنٹ شامل کرو:
  \[
  \Axiom$!B^k, !C^{k+1}, \Pi \fCenter \Lambda$
  \RightLabel{\LeftR{\land}}
  \UnaryInf$(!B \land !C)^i, \Pi \fCenter \Lambda$
  \DisplayProof
  \]
  نویں سیکوئنٹ نوں دِتّا !!{constant}s دا سیٹ~$C$ اے۔
  \item $!A \ident !B \land !C$ تے سجّے پاسے آؤندا اے۔ $\Pi \Sequent \Lambda, (!B \land !C)^i$ توں اُتّے دو نویں سبھ توں اُتلے سیکوئنٹ شامل کرو:
  \[
  \Axiom$\Pi \fCenter \Lambda, !B^k$
  \Axiom$\Pi \fCenter \Lambda, !C^k$
  \RightLabel{\RightR{\land}}
  \BinaryInf$\Pi \fCenter \Lambda, (!B \land !C)^i$
  \DisplayProof
  \]
  نویں سیکوئنٹ نوں دِتّا !!{constant}s دا سیٹ~$C$ اے۔

  \item $!A \ident \lnot !B$، $!A \ident !B \lor !C$ تے $!A \ident !B \lif !C$ والیاں صورتاں ایہناں ورگیاں نیں تے مشقاں لئی چھڈّیاں گئیاں نیں۔

  \item $!A \ident \lexists[x][!B(x)]$ تے کھبّے پاسے آؤندا اے۔ $\lexists[x][!B(x)]^i, \Pi \Sequent \Lambda$ توں اُتّے اک نواں سبھ توں اُتلا سیکوئنٹ شامل کرو:
  \[
  \Axiom$!B(c)^k, \Pi \fCenter \Lambda$
  \RightLabel{\LeftR{\lexists}}
  \UnaryInf$\lexists[x][!B(x)]^i, \Pi \fCenter \Lambda$
  \DisplayProof
  \]
  جتھے $c$ پہلا !!{constant} اے جیہڑا~$C$ وچ نہیں۔ نویں سیکوئنٹ نوں دِتّا !!{constant}s دا سیٹ $C \cup \{c\}$ اے۔
  \item $!A \ident \lexists[x][!B(x)]$ تے سجّے پاسے آؤندا اے۔ $\Pi \Sequent \Lambda, \lexists[x][!B(x)]$ توں اُتّے اک نواں سبھ توں اُتلا سیکوئنٹ شامل کرو:
  \[
  \Axiom$\Pi \fCenter \Lambda, \lexists[x][!B(x)]^k, !B(t)^{k+1}$
  \RightLabel{\LeftR{\lexists}}
  \UnaryInf$\Pi \fCenter \Lambda, \lexists[x][!B(x)]^i$
  \DisplayProof
  \]
  جتھے $t$ پہلا اوہ ٹرم اے جیہدے وچ صرف اوہ !!{constant}s ہون جیہڑے~$C$ وچ نیں، تے جیہڑا ایس سیکوئنٹ توں تھلّے سجّے پاسے $ \lexists[x][!B(x)]$ دی تقلیل لئی پہلاں نہ ورتیا ہووے (یا $\Obj c_0$، جے ایہو جیہے سارے ٹرم ورتے جا چکے نیں)۔ نویں سیکوئنٹ نوں دِتّا !!{constant}s دا سیٹ $C$ اے۔
  \item $!A \ident \lforall[x][!B(x)]$ والیاں صورتاں ایہناں ورگیاں نیں تے مشقاں لئی چھڈّیاں گئیاں نیں۔
\end{enumerate}

جے پڑاؤ $n+1$ اُتے اسیں کسے وی سبھ توں اُتلے سیکوئنٹ نال کوئی نواں سیکوئنٹ نہیں جوڑیا، تاں الگورتھم کامیابی نال ختم ہو جاندا اے۔ ایس صورت وچ اسیں !!a{proof} لبھ لیا اے~$\Gamma \Sequent \Delta$ دا، جیہڑا پڑاؤ $n+1$ دے درخت توں سارے انڈیکس ہٹا کے ملدا اے۔ سبھ توں اُتلے سارے سیکوئنٹ یا $!A, \Pi \Sequent \Lambda, !A$ دی شکل دے نیں یا $\lfalse, \Pi
\Sequent \Lambda$ دی۔ سارے استنباط~\Log{G3c} دے درست استنباط نیں۔ قضیاتی استنباطاں لئی ایہ صاف اے۔ ویکھو کہ ہر قدم اُتے سیٹ~$C$، یعنی سیکوئنٹ نوں دِتّا !!{constant}s دا سیٹ، اوہ سارے !!{constant}s رکھدا اے جیہڑے سیکوئنٹ وچ آؤندے نیں۔ ایس لئی \LeftR{\lexists} تے \RightR{\lforall} نال تقلیل وچ خاص متغیر دی شرط پوری ہوندی اے۔ خاص متغیر~$c$ اک !!a{constant} سی جیہڑا نتیجے نوں دِتّے سیٹ~$C$ وچ نہیں سی، سو $c$ نتیجے وچ وی نہیں آؤندا سی۔ سانوں ایہ ملدا اے:

\begin{prop}
\Log{G3c} لئی ثبوت لبھّن دا الگورتھم درست اے۔ جے ایہ کامیابی نال ختم ہووے، تاں شروع دے سیکوئنٹ~$\Gamma \Sequent \Delta$ دا !!a{proof} اے۔
\end{prop}
\endgroup
% /OLP-0683
% OLP-0679 content/proof-theory/proof-search/completeness.tex
\begingroup
\def\olfilename{completeness}
\expandafter\def\csname ol@sectioncs\endcsname{section}

\olfileid{pt}{ps}{cpl}

\olsection{تمامیت}

ہُن اسیں وکھاؤنے آں کہ ثبوت لبھّن والا الگورتھم خطا نہیں کھاندا۔ یعنی جے ایہ وچّے رُک جاوے یا کدے ختم نہ ہووے، تاں شروع دے سیکوئنٹ~$\Gamma \Sequent \Delta$ دا کوئی !!{proof} نہیں۔ ایہ وکھاؤن لئی اسیں اک !!a{structure}~$\Struct M$ ایویں متعین کردے آں کہ ہر $!A \in \Gamma$ لئی $\Sat{M}{!A}$، تے ہر~$! \in \Delta$ لئی $\Sat/{M}{!A}$۔ دوجے لفظاں وچ، سارے !!{formula}s، جیہڑے~$\Gamma$ وچ نیں، سچے نیں تے سارے !!{formula}s، جیہڑے~$\Delta$ وچ نیں، جھوٹھے نیں~$\Struct M$ وچ۔ یعنی $\Gamma \Sequent \Delta$ معتبر نہیں۔ چونکہ \Log{G3c} صحت مند اے، ایس لئی $\Gamma \Sequent \Delta$ دا کوئی !!{proof} نہیں۔

اسیں $\Struct{M}$ نوں جوڑے $ \Theta, \Xi$ توں، جیہڑا !!{formula}s دا اے، تے سیٹ~$C$ توں، جیہڑا !!{constant}s دا اے، متعین کردے آں۔ سانوں $\Theta, \Xi$ تے~$C$ ثبوت دی ناکام بھال دی اک \emph{ناکامی والی شاخ} توں ملدے نیں، بھاویں بھال وچّے رُک جاوے یا ختم ای نہ ہووے۔ ناکامی والی شاخ سیکوئنٹاں دی لڑی $ \Pi_n \Sequent
\Lambda_n$ تے !!{constant}s دے سیٹاں~$C_n$ توں بندی اے، جتھے $\Pi_0 \Sequent \Lambda_0$ آخری سیکوئنٹ $ \Gamma \Sequent
\Delta$ اے، تے $\Pi_{n+1} \Sequent \Lambda_{n+1}$، $ \Pi_n
\Sequent \Lambda_n$ دا اوہ مقدمہ اے جیہڑا پڑاؤ~$n+1$ اُتے بنیا تے جیہدا الگورتھم !!a{proof} نہیں لبھدا۔ صاف اے کہ ہر~$n$ لئی ایہو جیہا مقدمہ ضرور موجود اے، نئیں تاں الگورتھم !!a{proof} لبھ چکیا ہوندا۔ $C_n$ نوں اوہ !!{constant}s دا سیٹ منّ لو جیہڑا~$\Pi_n
\Sequent \Delta_n$ نال جُڑیا اے۔

جے الگورتھم کسے سبھ توں اُتلے سیکوئنٹ اُتے وچّے رُک جاوے، جیہدے وچ صرف جوہری !!{formula}s نیں، تاں اوس سیکوئنٹ اُتے مُکّن والی شاخ ناکامی والی شاخ اے۔ جے الگورتھم کدے ختم نہیں ہوندا، تاں ایہ ہور وڈّے توں وڈّے درخت بناؤندا رہندا اے۔ تسی ایہناں نوں اک لامحدود درخت ودھاؤن وانگ سمجھ سکدے او (اوہ درخت جیہڑا الگورتھم دے لامحدود قدم چلن توں پچھوں ہوندا)۔ ایہ درخت لامحدود ہووے گا پر ہر تھاں اوہدیاں شاخاں محدود ہون گئیاں (ہر گنڈھ دے صرف اک یا دو بچے نیں، یعنی جھٹ اوہدے اُتّے آؤن والے سیکوئنٹ)۔ K\H{o}nig دے لیمے نال ہر لامحدود درخت جیہدی ہر تھاں شاخ بندی محدود اے، اوہدی اک لامحدود شاخ ضرور ہوندی اے۔ جے بھال دا الگورتھم ہمیشہ چلدا رہے، تاں ایہو جیہی لامحدود شاخ ناکامی والی شاخ ہووے گی۔

جے $\Pi_n \Sequent \Lambda_n$ تے $C_n$ دی لڑی ناکامی والی شاخ اے، تاں $ \Theta = \bigcup_n \Pi_n$ تے $ \Xi = \bigcup_n \Lambda_n$ رکھو (!!{formula}s دے انڈیکس تے مُڑ آؤن والیاں نقلاں ہٹا کے)، تے $C =
\bigcup_n C_n$ رکھو۔

ہُن اسیں~$\Struct{M}$ دی تعریف دے سکدے آں۔
\begin{enumerate}
\item دائرہ $\Domain{M}$ اوہناں ٹرماں دا سیٹ اے جیہڑے~$C$ تے !!{function}s، جیہڑے~$\Gamma \Sequent \Delta$ وچ ہون جے کوئی ہون، نال بن سکدے نیں، پر !!{variable}s توں بناں۔
\item جے $c \in \Domain{M}$، تاں $\Assign{c}{M} = c$۔
\item جے $f$، $m$ تھاواں والا !!{function} اے $ \Gamma \Sequent \Delta$ وچ، تے $t_1$، \dots، $t_m \in \Domain{M}$، تاں $\Assign{f}{M}(t_1, \dots, t_m) =
f(t_1, \dots, t_n)$۔
\item جے $R$، $m$ تھاواں والا !!{predicate} اے $ \Gamma \Sequent
\Delta$ وچ، تاں $\Assign{R}{M} = \Setabs{\tuple{t_1, \dots, t_m} \in
\Domain{M}^n}{R(t_1, \dots, t_n) \in \Theta}$۔
\end{enumerate}
$\Struct{M}$ اک \emph{ٹرماں توں بنیا نمونہ} اے، کیونکہ اوہدا دائرہ ٹرماں دا سیٹ اے، تے !!{constant}s تے !!{function}s دے مطلب ایویں دِتّے نیں کہ $\Value{t}{M} = t$ ضرور ہووے۔ جوہری !!{formula}s، جیہڑے~$\Theta$ وچ نیں، $\Assign{R}{M}$ متعین کرن لئی ایویں ورتے نیں کہ $\Sat{M}{R(t_1, \dots, t_n)}$ اوہدوں تے صرف اوہدوں جے $R(t_1, \dots, t_n)
\in \Theta$۔

\begin{lem}\ollabel{lem:truth}
ہر $!A \in \Theta \cup \Xi$ لئی:
\begin{enumerate}
  \item جے $!A \in \Theta$، تاں $\Sat{M}{!A}$۔
  \item جے $!A \in \Xi$، تاں $\Sat/{M}{!A}$۔
\end{enumerate}
\end{lem}

\begin{proof}
$\depth{!A}$ اُتے استقرا نال۔

پہلاں منّ لو کہ $!A$ جوہری اے، یعنی یا $!A \ident \lfalse$ یا $!A
\ident R(t_1, \dots, t_m)$۔

کیونکہ $\Pi_n \Sequent \Lambda_n$ ناکامی والی شاخ اے، کسے $ \Pi_n$ وچ $ \lfalse$ نہیں آ سکدا (نئیں تاں $ \Pi_n \Sequent \Lambda_n$ اک مسلّمہ ہوندا)۔ $\Sat{M}{R(t_1, \dots, t_m)}$ اوہدوں تے صرف اوہدوں جے $R(t_1, \dots, t_m) \in
\Theta$،~$\Struct{M}$ دی تعریف نال۔ دونے گلاں رلا کے ملدا اے کہ جے $!A \in
\Theta$، تاں $\Sat{M}{!A}$۔

جے $!A$ جوہری اے تے $!A \in \Pi_n$، تاں $!A \in \Pi_{n+1}$، تے جے $!A \in \Lambda_n$، تاں $!A \in \Lambda_{n+1}$۔ (ایہ اوس طریقے توں نکلدا اے جیہدے نال ثبوت لبھّن والا الگورتھم اگلے سیکوئنٹ بناؤندا اے۔) سو جے $!A \in \Theta \cap \Xi$، تاں کسے~$n$ لئی $!A \in \Pi_n \cap
\Lambda_n$۔ ایہدا مطلب اے کہ $\Pi_n \Sequent \Lambda_n$ اک مسلّمہ اے، تے ناکامی والی شاخ وچ کوئی مسلّمہ نہیں ہو سکدا۔ سو بناوٹ دے حساب نال $!A \notin \Theta
\cap \Xi$ جوہری~$!A$ لئی، ایس لئی جے $!A \in
\Xi$، تاں $!A \notin \Theta$۔ ایہدے توں نکلدا اے کہ جے $!A \in \Xi$، تاں $\Sat/{M}{!A}$،~$\Struct{M}$ دی تعریف نال۔

استقرائی قدم لئی منّ لو کہ (1) تے (2) اوہناں سارے !!{formula}s لئی ٹھیک نیں جیہناں دی !!{depth}~$!A$ توں گھٹ اے۔ اسیں (1) تے (2) نوں $!A$ لئی وکھ وکھ صورتاں وچ ثابت کردے آں۔

پہلاں ویکھو کہ جدوں $!A^i$، $ \Pi_n \Sequent \Lambda_n$ وچ آؤندا اے، تاں $!A^i$، $ \Pi_{n+1} \Sequent \Lambda_{n+1}$ وچ وی آؤندا اے، سوائے جدوں $i$، $ \Pi_n \Sequent \Lambda_n$ وچ سبھ توں چھوٹا انڈیکس اے۔ چونکہ ہر پڑاؤ اُتے شامل ہوئے سبھ توں اُتلے سیکوئنٹاں وچ سبھ توں چھوٹا انڈیکس ودھدا اے، آخر اسیں $ \Pi_n \Sequent \Lambda_n$ تک پہنچدے آں جیہدے وچ $!A^i$ آؤندا اے تے $i$ سبھ توں چھوٹا انڈیکس اے۔ پڑاؤ $n+1$ اُتے الگورتھم~$!A^i$ نوں گھٹاؤندا اے۔ دوجے لفظاں وچ، جے $!A \in \Theta$، تاں کسے $n$ لئی $ \Pi_n = \Pi_n',
!A^i$ تے $i$ سبھ توں چھوٹا انڈیکس اے۔ تے جے $!A \in \Xi$، تاں کسے $n$ لئی $ \Lambda_n = \Lambda_n', !A^i$ تے $i$ سبھ توں چھوٹا انڈیکس اے۔

\begin{enumerate}
  \item $!A \in \Theta$ تے $!A \ident !B \land !C$۔ تاں کسے $n$ لئی $ \Pi_n = \Pi_n', (!B \land !C)^i$۔ $\Pi_{n+1} \Sequent
  \Lambda_{n+1}$،~\LeftR{\land} دا میل کھاندا مقدمہ اے، یعنی $ \Pi_{n+1} = \Pi_n', !B^k, !C^{k+1}$۔ ایس لئی $!B, !C \in
  \Theta$۔ استقرا دے مفروضے نال $ \Sat{M}{!B}$ تے $ \Sat{M}{!C}$۔ $\Sat{M}{}$ دی تعریف نال سانوں $ \Sat{M}{!B \land !C}$ ملدا اے۔

  \item $!A \in \Xi$ تے $!A \ident !B \land !C$۔ تاں کسے $n$ لئی $ \Lambda_n = \Lambda_n', !B \land !C$۔ $\Pi_{n+1} \Sequent
  \Lambda_{n+1}$،~\RightR{\land} دے دو مقدمیّاں وچوں اک اے جیہدا نتیجہ $ \Pi_n \Sequent \Lambda'_n, !B \land !C$ اے، یعنی $ \Lambda_{n+1} = \Lambda_n', !B^k$ یا $ \Lambda_{n+1} = \Lambda_n',
  !C^k$۔ ایس لئی یا $!B \in \Xi$ یا $!C \in \Xi$۔ استقرا دے مفروضے نال یا $\Sat/{M}{!B}$ یا $\Sat/{M}{!C}$۔ $\Sat{M}{}$ دی تعریف نال سانوں $ \Sat/{M}{!B \land !C}$ ملدا اے۔

  \item $!A \ident \lnot !B$، $!A \ident !B \lor !C$ تے $!A \ident !B \lif !C$ والیاں صورتاں ایہناں ورگیاں نیں تے مشقاں لئی چھڈّیاں گئیاں نیں۔

  \item $!A \in \Theta$ تے $!A \ident \lexists[x][!B(x)]$۔ تاں کسے $n$ لئی $ \Pi_n = \Pi_n', \lexists[x][!B(x)]^i$۔ $\Pi_{n+1}
  \Sequent \Lambda_{n+1}$،~\LeftR{\lexists} دا میل کھاندا مقدمہ اے، یعنی $ \Pi_{n+1} = \Pi_n', !B(c)^k$ کسے~$c$ لئی۔ ایس لئی $!B(c) \in \Theta$ تے~$c \in C$۔ استقرا دے مفروضے نال $\Sat{M}{!B(c)}$۔ $\Sat{M}{}$ دی تعریف نال سانوں $ \Sat{M}{\lexists[x][!B(x)]}$ ملدا اے۔

  \item $!A \in \Xi$ تے $!A \ident \lexists[x][!B(x)]$۔ تاں کسے $n$ لئی $ \Lambda_n = \Lambda_n', \lexists[x][!B(x)]^i$۔ کیونکہ جدوں وی $ \lexists[x][!B(x)]$ نوں گھٹایا جاندا اے، $\lexists[x][!B(x)]$ اک نویں انڈیکس نال مقدمے والے سیکوئنٹ دے سجّے پاسے رہندا اے، ایس لئی $ \lexists[x][!B(x)]$ نوں ناکامی والی شاخ وچ سجّے پاسے لامحدود واری گھٹایا جاندا اے، جے اوہ اُتھے آؤندا اے۔ ہر ٹرم $t \in
  \Domain{M}$ آخر اک ناکامی والی شاخ اُتے «پہلا اوہ ٹرم جیہدے وچ صرف~$C_n$ والے مستقل ہون تے جیہڑا سجّے پاسے $ \lexists[x][!B(x)]$ دی تقلیل وچ پہلاں نہ ورتیا گیا ہووے» بن جاندا اے۔ سو ہر $t \in \Domain{M}$ لئی $!B(t) \in \Lambda_n$ کسے~$n$ لئی، تے ایس لئی $!B(t) \in \Xi$۔ استقرا دے مفروضے نال $ \Sat/{M}{!B(t)}$ ہر $t \in
  \Domain{M}$ لئی۔ $\Sat{M}{}$ دی تعریف نال سانوں $ \Sat/{M}{\lexists[x][!B(x)]}$ ملدا اے۔

  \item $!A \ident \lforall[x][!B(x)]$ والیاں صورتاں ایہناں ورگیاں نیں تے مشقاں لئی چھڈّیاں گئیاں نیں۔
\end{enumerate}
\end{proof}

\begin{cor}\ollabel{cor:complete}
\Log{G3c} لئی ثبوت لبھّن دا الگورتھم کامل اے۔ جے ایہ وچّے رُک جاوے یا کدے ختم نہ ہووے، تاں شروع دا سیکوئنٹ~$\Gamma \Sequent \Delta$ معتبر نہیں تے ایس لئی اوہدا کوئی !!{proof} نہیں۔
\end{cor}

\begin{proof}
جے الگورتھم وچّے رُک جاوے یا کدے ختم نہ ہووے، تاں اک ناکامی والی شاخ اے جیہدے توں $\Struct{M}$ نوں متعین کیتا جا سکدا اے۔ \olref{lem:truth} نال ہر $!A \in \Gamma$ لئی $\Sat{M}{!A}$، تے ہر $!A \in \Delta$ لئی $\Sat/{M}{!A}$۔ (چیتے رکھو کہ $\Pi_0 = \Gamma$ تے $\Lambda_0 =
\Delta$، سو $\Gamma \subseteq \Theta$ تے $\Delta \subseteq \Xi$۔) سو $ \Gamma \Sequent \Delta$ معتبر نہیں۔~\Log{G3c} دی صحت مندی کرکے $ \Gamma \Sequent \Delta$ دا کوئی !!{proof} نہیں۔
\end{proof}

\begin{cor}
\Log{G3c} کامل اے۔ جے $ \Gamma \Sequent \Delta$ معتبر اے، تاں $\Log{G3c} \Proves \Gamma \Sequent \Delta$۔
\end{cor}

\begin{proof}
اسیں عکسِ نقیض ثابت کردے آں۔ منّ لو $ \Log{G3c} \Proves/ \Gamma
\Sequent \Delta$۔ تاں ثبوت لبھّن دا الگورتھم یا وچّے رُک جاندا اے یا کدے ختم نہیں ہوندا، کیونکہ لبھّن لئی کوئی !!{proof} نہیں۔ \olref{cor:complete} نال $ \Gamma \Sequent \Delta$ معتبر نہیں۔
\end{proof}
\endgroup
% /OLP-0679
% OLP-0684 content/proof-theory/proof-search/tableaux.tex
\begingroup
\def\olfilename{tableaux}
\expandafter\def\csname ol@sectioncs\endcsname{section}

\olfileid{pt}{ps}{tab}

\olsection{تابلو}

تابلو نظام !!{proof} نظام نیں جیہڑے سیکوئنٹ حساباں نال نیڑے جُڑے نیں تے ثبوت لبھّن لئی خاص طور اُتے راس آؤندے نیں، بھاویں کم ہتھ نال ہووے یا سافٹ ویئر وچ۔ تابلو وچ !!a{proof} سیکوئنٹاں دا درخت بناؤن دی تھاں !!{formula}s دا درخت اے۔ پر قدرتی اخذ توں وکھرے، ایہناں دے قاعدے سیکوئنٹ حساب دے قاعدیاں ورگے نیں۔ تابلو !!{proof} اصل وچ کسے نمونے دی کھُلّی بھال اے۔ ایس لئی ایہناں نوں کدے \emph{معنوی} تابلو یا \emph{سچائی دے درخت} وی آکھدے نیں۔

اک \emph{نشان والا} تابلو نشان والے !!{formula}s دا درخت اے، جیہڑے $ \sFmla{\True}{!A}$ یا $ \sFmla{\False}{!A}$ دی شکل دے ہو سکدے نیں۔ درخت تھلّے ول ودھدے نیں۔ اوہ اوہناں !!{formula}s توں شروع ہوندے نیں جیہڑے دَسّدے نیں کہ اسیں کیہ !!{prove} کرنا چاہنے آں۔ مثال لئی !!{prove} کرن لئی سیکوئنٹ $!A, !B \Sequent !C, !D$، ایہ $ \sFmla{\True}{!A}, \sFmla{\True}{!B},
\sFmla{\False}{!C}, \sFmla{\True}{!D}$ توں شروع ہووے گا۔ !!^a{structure} جیہڑی $!A$ تے $!B$ نوں سچا تے $!C$ تے $!D$ نوں جھوٹھا کردی اے، اوہ ایہو جیہی !!a{structure} اے جیہڑی وکھاوے گی کہ $!A, !B \Sequent !C, !D$ معتبر نہیں۔

تابلو دی آخری (سبھ توں تھلّی) گنڈھ نوں شاخ ودھاؤن دے قاعدیاں نال ودھایا جا سکدا اے۔ ایہ قاعدے اوسے شاخ اُتے ہور گنڈھاں جوڑدے نیں یا تھلّے دو سانجھیاں گنڈھاں جوڑ کے شاخ نوں دو نویاں شاخاں وچ ونڈ دیندے نیں۔ قاعدے اوسے شاخ اُتے گنڈھ توں اُتّے کسے وی نشان والے !!{formula} اُتے لائے جا سکدے نیں۔ کلاسیکی تابلو نظام \Log{Tc} لئی شاخ ودھاؤن دے قاعدے \olref{tab:Tc} وچ دِتّے نیں۔

% OLP-0682 content/proof-theory/proof-search/rules-Tc.tex
\begingroup

\begin{table}
\begin{center}
\begin{tabular}{@{}cc@{}}
\AxiomC{\sFmla{\True}{\lnot !A}}
\RightLabel{\TRule{\True}{\lnot}}
\UnaryInfC{\sFmla{\False}{!A}}
\DisplayProof
&
\AxiomC{\sFmla{\False}{\lnot !A}}
\RightLabel{\TRule{\False}{\lnot}}
\UnaryInfC{\sFmla{\True}{!A}}
\DisplayProof
\\[6ex]
\AxiomC{\sFmla{\True}{!A \land !B}}
\RightLabel{\TRule{\True}{\land}}
\UnaryInfC{\sFmla{\True}{!A}}
\noLine
\UnaryInfC{\sFmla{\True}{!B}}
\DisplayProof
&
\AxiomC{\sFmla{\False}{!A \land !B}}
\RightLabel{\TRule{\False}{\land}}
\UnaryInfC{$\sFmla{\False}{!A} \quad \mid \quad \sFmla{\False}{!B}$}
\DisplayProof
\\[6ex]
\AxiomC{\sFmla{\True}{!A \lor !B}}
\RightLabel{\TRule{\True}{\lor}}
\UnaryInfC{$\sFmla{\True}{!A} \quad \mid \quad \sFmla{\True}{!B}$}
\DisplayProof
&
\AxiomC{\sFmla{\False}{!A \lor !B}}
\RightLabel{\TRule{\False}{\lor}}
\UnaryInfC{\sFmla{\False}{!A}}
\noLine
\UnaryInfC{\sFmla{\False}{!B}}
\DisplayProof
\\[6ex]
\AxiomC{\sFmla{\True}{!A \lif !B}}
\RightLabel{\TRule{\True}{\lif}}
\UnaryInfC{$\sFmla{\False}{!A} \quad \mid \quad \sFmla{\True}{!B}$}
\DisplayProof
&
\AxiomC{\sFmla{\False}{!A \lif !B}}
\RightLabel{\TRule{\False}{\lif}}
\UnaryInfC{\sFmla{\True}{!A}}
\noLine
\UnaryInfC{\sFmla{\False}{!B}}
\DisplayProof
\\[6ex]
\AxiomC{\sFmla{\True}{\lforall[x][!A(x)]}}
\RightLabel{\TRule{\True}{\forall}}
\UnaryInfC{\sFmla{\True}{!A(t)}}
\DisplayProof
&
\AxiomC{\sFmla{\False}{\lforall[x][!A(x)]}}
\RightLabel{\TRule{\False}{\lforall}}
\UnaryInfC{\sFmla{\False}{!A(a)}}
\DisplayProof
\\[6ex]
\AxiomC{\sFmla{\True}{\lexists[x][!A(x)]}}
\RightLabel{\TRule{\True}{\lexists}}
\UnaryInfC{\sFmla{\True}{!A(a)}}
\DisplayProof
&
\AxiomC{\sFmla{\False}{\lexists[x][!A(x)]}}
\RightLabel{\TRule{\False}{\lexists}}
\UnaryInfC{\sFmla{\False}{!A(t)}}
\DisplayProof
\end{tabular}
\end{center}
\caption{\Log{Tc} دے قاعدے۔ $\TRule{\False}{\lforall}$ تے
$\TRule{\True}{\lexists}$ وچ !!{constant}~$a$ اُتلی شاخ وچ نہیں آؤنا چاہیدا، یعنی شاخ لئی نواں ہونا چاہیدا اے۔}
\ollabel{tab:Tc}
\end{table}
\endgroup
% /OLP-0682

جیویں تسی ویکھ سکدے او، ایہ \Log{G3c} دے اوہی قاعدے نیں، الٹے پاسے لائے ہوئے۔ نشان $ \True$ تے $ \False$ بالترتیب دَسّدے نیں کہ !!a{formula} سیکوئنٹ دے کھبّے یا سجّے پاسے دا مُکھ فارمولا اے۔

تابلو دی شاخ \emph{بند} اے جے اوہدے وچ $ \sFmla{\True}{!A}$ تے $ \sFmla{\False}{!A}$ دونے ہون۔ جے ہر شاخ بند اے، تاں اسیں آکھدے آں کہ سارا تابلو بند اے۔ مثال لئی ایہ $!A \lor (!B \land !C) \Sequent (!A \lor !B) \land
(!A \lor !C)$ دا کلاز تابلو اے، یعنی $ \sFmla{\True}{!A \lor (!B \land
!C)}$ تے $ \sFmla{\False}{(!A \lor !B) \land (!A \lor !C)}$ توں شروع ہون والا:
\begin{oltableau}
  [\sFmla{\True}{\formula{A} \lor (\formula{B} \land \formula{C})},
    just=\TAss, checked
    [\sFmla{\False}{(\formula{A} \lor \formula{B}) \land
        (\formula{A} \lor \formula{C})}, just=\TAss, checked
      [\sFmla{\True}{\formula{A}}, just={\TRule{\True}{\lor}[1]}
        [\sFmla{\False}{\formula{A} \lor \formula{B}},
          just={\TRule{\False}{\land}[2]}, checked
          [\sFmla{\False}{\formula{A}},
            just={\TRule{\False}{\lor}[4]}
            [\sFmla{\False}{\formula{B}},
              just={\TRule{\False}{\lor}[4]}, close]
          ]
        ]
        [\sFmla{\False}{\formula{A} \lor \formula{C}},
          just={\TRule{\False}{\land}[2]}, checked
          [\sFmla{\False}{\formula{A}},
            just={\TRule{\False}{\lor}[4]}
            [\sFmla{\False}{\formula{C}},
              just={\TRule{\False}{\lor}[4]}, close]
          ]
        ]
      ]
      [\sFmla{\True}{\formula{B} \land \formula{C}},
        just={\TRule{\True}{\lor}[1]}, checked
        [\sFmla{\False}{\formula{A} \lor \formula{B}},
          just={\TRule{\False}{\land}[2]}, checked
          [\sFmla{\True}{\formula{B}},
            just={\TRule{\True}{\land}[3]}, move by=2
            [\sFmla{\True}{\formula{C}},
              just={\TRule{\True}{\land}[3]}
              [\sFmla{\False}{\formula{A}},
                just={\TRule{\False}{\lor}[4]}
                [\sFmla{\False}{\formula{B}},
                  just={\TRule{\False}{\lor}[4]},close
                ]
              ]
            ]
          ]
        ]
        [\sFmla{\False}{\formula{A} \lor \formula{C}},
          just={\TRule{\False}{\land}[2]}, checked
          [\sFmla{\True}{\formula{B}},
            just={\TRule{\True}{\land}[3]}, move by=2
              [\sFmla{\True}{\formula{C}},
                just={\TRule{\True}{\land}[3]}
                [\sFmla{\False}{\formula{A}},
                  just={\TRule{\False}{\lor}[4]}
                  [\sFmla{\False}{\formula{C}},
                  just={\TRule{\False}{\lor}[4]},close
                ]
              ]
            ]
          ]
        ]
      ]
    ]
  ]
\end{oltableau}

ٹِک دے نشان دَسّدے نیں کہ میل کھاندا قاعدہ !!a{formula} اُتے اوہدے تھلّے دیاں ساریاں شاخاں وچ لایا گیا اے۔ $ \otimes$ دا نشان دَسّدا اے کہ شاخ بند اے۔ قاعدیاں دے پچھوں سطراں دے نمبر دَسّدے نیں کہ قاعدہ کِہڑے نشان والے !!{formula}s اُتے لگدا اے (یعنی اوسے شاخ اُتے دَسّی گئی سطر دا !!{formula})۔

تابلو وچ ثبوت لبھنا سوکھا اے۔ اسیں پڑاؤواں وچ تابلو بناؤنے آں۔ پڑاؤ $0$ اُتے بس شروع دے !!{formula}s اُتّے لکھو۔ پڑاؤ $n+1$ اُتے ہر آخری گنڈھ لئی ایہ کرو۔ پہلا (سبھ توں اُتلا) نشان والا فارمولا لبھو جیہدے اُتے ہجے ٹِک نہیں لگی۔ اوس فارمولے اُتے شاخ ودھاؤن دا قاعدہ لاؤ۔ جدوں متعلقہ نشان والے فارمولے والی ہر شاخ اُتے قاعدہ لگ جاوے، تاں فارمولے اُتے ٹِک لا دیو۔ (اُتلی مثال ایس الگورتھم نال بنی اے۔)

نشان والے !!{formula}s جیہڑے $\sFmla{\True}{\lforall}$ تے $ \sFmla{\False}{\lexists}$ دی شکل دے نیں، اوہناں اُتے کدے ٹِک نہیں لگدی، ایس لئی اوہناں نوں خاص طریقے نال سنبھالنا پیندا اے۔ جے اوہ موجود ہون، تاں سانوں پکا کرنا پیندا اے کہ $ \TRule{\True}{\forall}$ تے $ \TRule{\False}{\lexists}$ قاعدے آخر سارے لاگو ٹرماں لئی لگن۔ مثال لئی ایہ ایویں پکا کیتا جا سکدا اے کہ ہر پڑاؤ اُتے ہر شاخ دے ایہو جیہے سارے فارمولیاں اُتے اوہ قاعدے لائیے (اوس سبھ توں اُتلے نشان والے !!{formula} دا کم کرن توں پچھوں جیہڑا ایس شکل دا نہیں)۔ ورتن والا ٹرم $t$ پہلا ایہو جیہا ٹرم~$t$ اے جیہڑا !!{constant}s، جیہڑے پہلاں ای شاخ اُتے آؤندے نیں (یا $\Obj
c_0$ جے کوئی !!{constant} نہیں آؤندا)، تے !!{function}s، جیہڑے شروع دے !!{formula}s وچ آؤندے نیں، توں بن سکدا اے، تے جیہدے لئی میل کھاندا نشان والا !!{formula} $\sFmla{\True}{!B(t)}$ یا~$\sFmla{\False}{!B(t)}$ پہلاں ای شاخ اُتے نہیں آؤندا۔

جے ایہ بھال بند تابلو نہیں بناؤندی، تاں یا اوہدے وچ اک محدود کھُلّی شاخ اے جیہدے اُتے سارے غیر جوہری !!{formula}s اُتے ٹِک لگی اے، یا بھال اک لامحدود درخت بناؤندی اے جیہدی ہر تھاں شاخ بندی محدود اے تے جیہدی اک لامحدود شاخ ضرور اے۔ عین \olref[cpl]{sec} وانگ اسیں !!a{structure}~$\Struct{M}$ ایویں متعین کر سکدے آں کہ جے $ \sFmla{\True}{!A}$ شاخ اُتے آؤندا اے تاں $\Sat{M}{!A}$، تے جے $ \sFmla{\False}{!A}$ آؤندا اے تاں $\Sat/{M}{!A}$۔ ($\Theta =
\Setabs{!A}{\sFmla{\True}{!A} \text{ شاخ اُتے}}$ تے $\Xi =
\Setabs{!A}{\sFmla{\False}{!A} \text{ شاخ اُتے}}$ لو۔)

ایس طریقے نال ثبوت لبھدیاں بنے تابلو سوکھے ای \Log{G3c} دے ثبوتاں وچ ترجمہ کیتے جا سکدے نیں، ہر گنڈھ نوں اک سیکوئنٹ دے کے۔ جے شروع دے فارمولے سیکوئنٹ $ \Gamma \Sequent
\Delta$ نال میل کھاندے نیں، تاں ہر شروع دے فارمولے اُتے $ \Gamma \Sequent
\Delta$ دا لیبل لاؤ۔ جے شاخ ودھاؤن دا قاعدہ نشان والے !!{formula} $\sFmla{\True}{!A}$ اُتے لا کے شاخ ودھائی اے، تاں آخری گنڈھ دا لیبل $!A, \Pi \Sequent \Lambda$ اے۔ جے اوہ قاعدہ نشان والے !!{formula} $\sFmla{\False}{!A}$ اُتے لایا اے، تاں آخری گنڈھ دا لیبل $ \Pi \Sequent
\Lambda, !A$ اے۔ جتھے تابلو \TRule{\True}{\star} قاعدہ لاؤندا اے، اوتھے سیکوئنٹ اُتے \LeftR{\star} قاعدہ لاؤ (مُکھ فارمولا $!A$ رکھ کے)۔ جتھے \TRule{\False}{\star} قاعدہ لاؤندا اے، اوتھے \RightR{\star} قاعدہ لاؤ۔ قاعدے دے مقدمے اوہ سیکوئنٹ نیں جیہڑے تابلو وچ جوڑیاں گئیاں میل کھاندیاں گنڈھاں دے لیبل نیں۔ جدوں تابلو $ \sFmla{\True}{!A}$ تے $ \sFmla{\False}{!A}$ نال بند ہوندا اے، تاں آخری گنڈھ دا لیبل $!A, \Pi \Sequent \Lambda,
!A$ دی شکل دا سیکوئنٹ ہووے گا۔ تابلو وچ پے در پے کجھ گنڈھاں نوں اوہی سیکوئنٹ دِتّا جاوے گا۔ ودھ نقلاں ہٹا دیو۔ مثال لئی اُتلے تابلو دا ترجمہ ایہ !!{proof} اے:
\[\small
\Axiom$!A \fCenter !A, !B$
\RightLabel{\RightR{\lor}}
\UnaryInf$!A \fCenter !A \lor !B$
\Axiom$!A \fCenter !A, !C$
\RightLabel{\RightR{\lor}}
\UnaryInf$!A \fCenter !A \lor !C$
\RightLabel{\RightR{\land}}
\BinaryInf$!A \fCenter (!A \lor !B) \land (!A \lor !C)$
\Axiom$!B, !C \fCenter !A, !B$
\RightLabel{\RightR{\lor}}
\UnaryInf$!B, !C \fCenter !A \lor !B$
\RightLabel{\LeftR{\land}}
\UnaryInf$!B \land !C \fCenter !A \lor !B$
\Axiom$!B, !C \fCenter !A, !C$
\RightLabel{\RightR{\lor}}
\UnaryInf$!B, !C \fCenter !A \lor !C$
\RightLabel{\LeftR{\land}}
\UnaryInf$!B \land !C \fCenter !A \lor !C$
\RightLabel{\RightR{\land}}
\BinaryInf$!B \land !C \fCenter (!A \lor !B) \land
(!A \lor !C)$
\RightLabel{\LeftR{\lor}}
\BinaryInf$!A \lor (!B \land !C) \fCenter (!A \lor !B) \land
(!A \lor !C)$
\DisplayProof
\]
\endgroup
% /OLP-0684

\OLEndChapterHook
\endgroup
% /OLP-0681
% OLP-0664 content/proof-theory/natural-deduction/natural-deduction.tex
\begingroup
\def\olfilename{natural-deduction}
\expandafter\def\csname ol@sectioncs\endcsname{section}

\olchapter{pt}{nat}{قدرتی اخذ}

% OLP-0663 content/proof-theory/natural-deduction/introduction.tex
\begingroup
\def\olfilename{introduction}
\expandafter\def\csname ol@sectioncs\endcsname{section}

\olfileid{pt}{nat}{int}

\olsection{جان پچھان}

قدرتی اخذ !!{proof} نظاماں دا اک ٹبّر اے جیہدے وچ ہر منطقی جوڑ یا مسوّر لئی استنباطاں دی اک جوڑی اے۔ اک استنباط عامل نوں «شامل» کردا اے (!!a{introduction} قاعدہ) تے دوجا عامل نوں «ہٹاؤندا» اے (!!a{elimination} قاعدہ)۔ مثلاً \Intro{\lif} دا نتیجہ !!{formula}s ہوندے نیں~$!A \lif !B$ دی شکل دے، تے \Elim{\lif} دے مقدمیّاں وچ !!{formula}s~$!A \lif !B$ دی شکل دے ہوندے نیں۔

قدرتی اخذ دے پہلے دو روپ 1930 دی دہائی وچ گیرہارٹ گینٹسن تے ستانیسواف یاشکوفسکی (Stanisław Jaśkowski) نے پیش کیتے۔ ثبوت دے نظریہ دان زیادہ تر گینٹسن دے نظام دی گل کردے نیں، تے یاشکوفسکی دے نظام توں اوہ نظام نکلے جیہڑے پڑھائی وچ سبھ توں ودھ ورتے جاندے نیں۔ دوہاں نظاماں وچ \emph{مفروضیاں} توں شروع کر سکدے آں۔ ایہ اوہ !!{formula}s نیں جیہناں نوں !!{prove} کرن دی لوڑ نہیں، پر اوہ !!{proving} وچ، ہور !!{formula}s دے، ورتے جا سکدے نیں۔ پورا !!{proof} ایہو جیہے مفروضیاں اُتے ٹِک سکدا اے۔ چونکہ اوہ ثابت نہیں کیتے، !!{proof} اپنا نتیجہ (اپنا \emph{آخری !!{formula}}) قائم نہیں کردا، بس ایہ قائم کردا اے کہ نتیجہ مفروضیاں توں نکلدا اے۔

کجھ استنباطی قاعدیاں دے نتیجے کجھ مفروضیاں اُتے ہور نہیں ٹِکدے۔ ایہ قاعدے مفروضے «خارج» کردے نیں۔ مثلاً \Intro{\lif} قاعدہ !!a{proof} لیندا اے نتیجے~$!B$ دا، مفروضے~$!A$ توں، تے ایہدا نتیجہ $!A \lif !B$ اے۔ !!{proof} جیہڑا $!A \lif !B$ اُتے مُکدا اے،~$!A$ اُتے ہور نہیں ٹِکدا۔ مفروضہ $!A$ خارج ہو جاندا اے۔ گینٹسن دے نظاماں وچ !!{proof}s، !!{formula}s دے درخت نیں، مفروضے سبھ توں اُتلے !!{formula}s نیں، تے~\Intro{\lif} ورگے قاعدیاں اُتے نشان نیں جیہڑے دسدے نیں کہ کیہڑے مفروضے خارج ہوندے نیں۔ یاشکوفسکی دے نظاماں وچ !!{proof} دے اوہ حصے جیہڑے کسے مفروضے اُتے ٹِکدے نیں، کسے ڈھنگ نال وکھ رکھے جاندے نیں، مثلاً کھڑی لکیر پچھے ہٹواں لکھ کے یا خانے وچ پا کے۔ خارج کیتا مفروضہ~$!A$ تے شرطیے دی تالی $!B$ خانے دیاں پہلی تے آخری سطراں نیں، تے نتیجہ~$!A \lif !B$، جیہڑا \Intro{\lif} دا اے، خانے توں باہر اے۔

نظام «قدرتی» نیں کیوں جے استنباطی قاعدے ساڈی (یا گھٹ توں گھٹ ریاضی داناں دی) قدرتی دلیل ورگے نیں۔ فرضی نتیجہ (شرطیہ) قائم کرن لئی مقدم منّ کے اوس نال تالی ثابت کردے آں۔ ایہ~\Intro{\lif} اے۔ جدوں شرطیہ $!A \lif !B$ تے اوہدا مقدم~$!A$ معلوم نیں، تاں نتیجہ~$!B$ کڈھدے آں۔ ایہ~\Elim{\lif} اے۔ فصلی عبارت~$!A \lor !B$ منّ کے کوئی نتیجہ سچ وکھاؤن لئی صورتاں وکھ کر کے ثبوت دیندے آں۔ ایہ~\Elim{\lor} اے۔ عام دعویٰ~$\lforall[x][!A(x)]$ ثابت کرن لئی وکھاؤنے آں کہ $!A(c)$ سچ اے، کسے من مرضی دے~$c$ لئی۔ ایہ~\Intro{\lforall} اے۔ ایس طرح اگے۔

مثلاً گینٹسن طرز دے قدرتی اخذ وچ سوکھا !!{proof}~$!A \lif (!A \lor !B)$ دا ایہ اے:
\[
\AxiomC{$\Discharge{!A}{x}$}
\RightLabel{\Intro{\lor}}
\UnaryInfC{$!A \lor !B$}
\DischargeRule{\Intro{\lif}}{x}
\UnaryInfC{$!A \lif (!A \lor !B)$}
\DisplayProof
\]
ایہ مفروضے~$!A$ توں شروع ہوندا اے، \Intro{\lor} لا کے~$!A
\lor !B$ کڈھدا اے، تے فیر \Intro{\lif} لا کے~$!A \lif (!A \lor !B)$ کڈھدا اے۔ \Intro{\lif} استنباط مفروضہ~$!A$ خارج کردا اے۔ ایہ نشان~$x$ دَسدا اے۔

ثبوتی نظریے دے ماہر گینٹسن دے اصلی نظام پسند کردے نیں جیہڑے !!{formula}s دے درختاں نال کم کردے نیں، بھاویں اک وکھرا روپ وی ثبوتی نظریے وچ اہم اے۔ !!^a{proof}~$!A$ دا، مفروضیاں~$\Gamma$ توں، وکھاؤندا اے کہ $!A$،~$\Gamma$ توں نکلدا اے، یعنی $\Gamma \Entails !A$۔ قدرتی اخذ دے استنباطی قاعدے قدرتی طور اُتے ایہو جیہے منطقی نتیجیاں بارے گل کردے پڑھے جا سکدے نیں۔ مثلاً \Intro{\lif} کہندا اے کہ جے $\Gamma + !A
\Entails !B$، تاں $\Gamma \Entails !B$۔ قاعدیاں نوں سدھا ایویں لکھنا قدرتی اے کہ اوہ استنباطی قاعدے دے مقدمیّاں تے نتیجے وچ، مفروضیاں تے اوہناں توں نکلن والے نتیجے دوہاں اُتے کم کرن۔ یعنی اوہ سیکوئنٹاں $\Gamma \Sequent !A$ اُتے کم کرن۔ اُتلا سوکھا ثبوت ایہو جیہے نظام وچ ایویں دِسے گا:
\[
\Axiom$x:!A \fCenter !A$
\RightLabel{\Intro{\lor}}
\UnaryInf$x:!A \fCenter !A \lor !B$
\DischargeRule{\Intro{\lif}}{x}
\UnaryInf$\fCenter !A \lif (!A \lor !B)$
\DisplayProof
\]
جیویں تسی ویکھدے او، قاعدے اوہی نیں۔ اوہ !!{formula} اُتے کم کردے نیں~$\Sequent$ دے سجّے۔ کھبّے !!{formula}s بس اوہناں مفروضیاں دی فہرست نیں جیہناں اُتے ہر قدم ٹِکدا اے۔

!!{introduction}/!!{elimination} قاعدیاں دیاں ساریاں جوڑیاں دا سیٹ اپنے آپ وچ کلاسیکی منطق لئی پورا نہیں۔ سانوں دو ہور قاعدے شامل کرنے پین گے جیہڑے~$\lfalse$ نال جڑے نیں۔ پہلا «وجدانی محال دا قاعدہ»~\FalseInt، یا «دھماکا»، کہندا اے کہ $\lfalse$ توں کجھ وی کڈھ سکدے آں۔ دوجا بالواسطہ ثبوت دا کلاسیکی اصول~\FalseCl{} (یا «کلاسیکی محال دا قاعدہ») کہندا اے کہ جے اسیں !!{prove} کریئے~$\lfalse$ نوں~$\lnot !A$ توں، تاں اسیں !!{prove} کر لَیا اے~$!A$۔ جے~\FalseCl چھڈ دیئے، تاں وجدانی منطق لئی ڈھکدا نظام ملدا اے۔ جے دوویں چھڈ دیئے، تاں «کمترین منطق» ملدی اے۔

اسیں کلاسیکی تے وجدانی منطق لئی گینٹسن طرز دے قدرتی اخذ دی گل کراں گے۔ ایہ نظام \Log{N1c} تے \Log{N1i} نیں۔ نظام \Log{N2c} تے \Log{N2i} اوہناں نال میل کھاندے نیں، پر سیکوئنٹاں $\Gamma \Sequent !A$ اُتے کم کردے نیں، اکّلے !!{formula}s~$!A$ دی تھاں۔
\endgroup
% /OLP-0663
% OLP-0668 content/proof-theory/natural-deduction/rules-proofs.tex
\begingroup
\def\olfilename{rules-proofs}
\expandafter\def\csname ol@sectioncs\endcsname{section}

\olfileid{pt}{nat}{rp}

\olsection{قاعدے تے \usetoken{P}{proof}}

اسیں پہلاں کجھ تعریفاں دیندے آں تے کجھ اصطلاحاں مقرر کردے آں۔

\Log{N1c} تے~\Log{N1i} دے اوہ دو قاعدے ویکھو جیہناں وچ $!A
\lif !B$ آؤندا اے۔
\[
\bottomAlignProof
\AxiomC{$\Discharge{!A}{x}$}
\shortDeduce
\DeduceC{$!B$}
\DischargeRule{\Intro{\lif}}{x}
\UnaryInfC{$!A \lif !B$}
\DisplayProof
\qquad
\bottomAlignProof
\AxiomC{$!A \lif !B$}
\AxiomC{$!A$}
\RightLabel{\Elim{\lif}}
\BinaryInfC{$!B$}
\DisplayProof
\]
\Elim{\lif} قاعدہ کافی سوکھا اے۔ لکیر توں اُتلے !!{formula}s \emph{مقدمے} نیں۔ کھبّے والا !!{formula}، \emph{مُکھ} !!{formula}~$!A \lif !B$ اے۔ سارے !!{elimination} قاعدیاں دا اک ایہو جیہا مقدمہ اے، ہمیشہ سبھ توں کھبّے۔ ایہنوں استنباط دا \emph{وڈا} مقدمہ آکھدے آں۔ جے ہور مقدمے ہون، جیویں ایس صورت وچ، تاں اوہناں نوں \emph{چھوٹے} مقدمے آکھدے آں۔ !!{introduction} قاعدے دے مقدمے یا مقدمیّاں نوں وی وڈا مقدمہ یا وڈے مقدمے آکھدے آں۔ لکیر توں تھلّا !!{formula} \emph{نتیجہ} اے۔

\Intro{\lif} قاعدے دا اکّو مقدمہ اے، تے ایتھے نتیجہ مُکھ فارمولا~$!A \lif !B$ اے۔ سارے !!{introduction} قاعدے ایہو جیہے نیں۔ نتیجہ مُکھ فارمولا اے، تے اوہدا مُکھ عامل اوہی اے جیہڑا «شامل» کیتا جاندا اے۔

لکھت $\Discharge{!A}{x}$ دا مطلب ایہ اے۔ !!a{proof} وچ، \Intro{\lif} استنباط !!a{formula}~$!B$ نوں مقدمے وجوں لے کے لایا جاندا اے۔~$!B$ اُتے مُکدا ذیلی !!{proof} کجھ مُڈھلے !!{formula}s رکھدا اے۔ ایہناں نوں \emph{مفروضے} آکھدے آں۔ !!^a{proof}~$!B$ دا،~$!A$ دی شکل دے اک یا ودھ مفروضیاں توں، ایس گل دا !!a{proof} اے کہ $!A$ توں~$!B$ منطقی طور اُتے نکلدا اے۔ جدوں \Intro{\lif} قاعدہ لاؤنے آں تاں~$!A
\lif !B$ نتیجہ ملدا اے، تے نتیجہ مفروضے~$!A$ اُتے ہور نہیں ٹِکدا۔ ایہ دَسن لئی~$!A$ دی شکل والے مفروضے اوس !!{proof} وچ، جیہدے وچ~\Intro{\lif} استنباط شامل اے، \emph{خارج} یا \emph{بند} کر دِتّے جاندے نیں۔ قاعدہ دسدا اے کہ کیہڑے مفروضے خارج ہو سکدے نیں۔ \Intro{\lif} دی صورت وچ، جیہدا نتیجہ $!A \lif !B$ اے، ایہ اوہ مفروضے نیں جیہڑے~$!A$ دی شکل دے نیں، یعنی شرطیے دے مقدم نال میل کھاندے نیں۔ کدے کدے ایہ دَسنا اہم اے کہ کیہڑا مفروضہ کیہڑے قاعدے اُتے خارج ہویا۔ ایس لئی اخراج دا نشان~$x$ ورتدے آں۔ اہم گل ایہ اے کہ مفروضے خارج کرن والے قاعدیاں لئی ایہ \emph{ضروری نہیں} کہ ڈھکدی شکل والے \emph{سارے} مفروضے خارج کرن، یا کوئی وی خارج کرن۔ مثلاً ہِٹھلا ثبوت ٹھیک اے، بھاویں پہلا \Intro{\lif} کوئی مفروضہ خارج نہیں کردا:
\[
\AxiomC{$\Discharge{!A}{y}$}
\DischargeRule{\Intro{\lif}}{x}
\UnaryInfC{$!B \lif !A$}
\DischargeRule{\Intro{\lif}}{y}
\UnaryInfC{$!A \lif (!B \lif !A)$}
\DisplayProof
\]
جدوں قاعدہ کوئی مفروضہ خارج نہیں کردا، اسیں اخراج دا نشان نہیں لاؤندے (جیویں ایتھے کیتا)۔

ایس !!{proof} وچ
\[
\AxiomC{$\Discharge{!A}{x}$}
\AxiomC{$\Discharge{!A}{y}$}
\RightLabel{\Intro{\land}}
\BinaryInfC{$!A \land !A$}
\RightLabel{\Elim{\land}}
\UnaryInfC{$!A$}
\DischargeRule{\Intro{\lif}}{x}
\UnaryInfC{$!A \lif !A$}
\DischargeRule{\Intro{\lif}}{y}
\UnaryInfC{$!A \lif (!A \lif !A)$}
\DisplayProof
\]
پہلا \Intro{\lif} استنباط کھبّا مفروضہ $!A^x$ خارج کردا اے، تے دوجا سجّا مفروضہ~$!A^y$۔ یعنی نشان~$x$ والے سارے مفروضے \Intro{\lif} نال خارج ہوندے نیں جیہدا نشان~$x$ اے، تے نشان~$y$ والے سارے مفروضے اوس نال جیہدا نشان~$y$ اے۔ ایہ کم کرن لئی ضروری اے کہ اکّو نشان والے سارے مفروضے اوہی !!{formula} ہون۔

مسوّراں دے قاعدے ذرا گنجل دار نیں۔ مثلاً \Log{N1i} وچ $\lexists$ دے قاعدے ایہ نیں:
\[
\bottomAlignProof
\AxiomC{$!A(t)$}
\RightLabel{\Intro{\lexists}}
\UnaryInfC{$\lexists[x][!A(x)]$}
\DisplayProof
\qquad
\bottomAlignProof
\AxiomC{$\lexists[x][!A(x)]$} 
\AxiomC{$\Discharge{!A(c)}{x}$}
\shortDeduce
\DeduceC{$!B$}
\DischargeRule{\Elim{\lexists}}{x}
\BinaryInfC{$!B$} \DisplayProof 
\bottomAlignProof
\]
\Intro{\lexists} قاعدے وچ مقدمہ~$!A(t)$ اے، جتھے $t$ بند ٹرم اے، یعنی اوہدے وچ کوئی متغیر نہیں۔ ایہو جیہا استنباط ٹھیک اے، یعنی قاعدے دے خاکے \Intro{\lexists} نال میل کھاندا اے، جے کوئی !!{formula}~$!A$، کوئی متغیر~$x$، تے کوئی بند ٹرم~$t$ ایہو جیہے ہون کہ نتیجہ $\lexists[x][!A]$ ہووے تے مقدمہ~$\Subst{!A}{t}{x}$۔ \Elim{\lexists} سانوں~$!A(c)$ دی شکل والے مفروضے خارج کرن دی کھُل دیندا اے۔ یعنی ہر استنباط لئی کوئی !!{constant}~$c$ اے تے $\Subst{!A}{c}{x}$ دی شکل والے مفروضے خارج ہو سکدے نیں۔ اک استنباط دے سارے خارج کیتے مفروضیاں وچ اوہی~$c$ ورتیا جانا چاہیدا اے۔ ایس ثابت نوں \Elim{\lexists} استنباط دا \emph{خاص متغیر} آکھدے آں۔ استنباط اوہدوں ای ٹھیک اے جدوں $c$ نہ خارج کیتے~$!A(c)$ توں وکھرے کسے کھُلّے مفروضے وچ آوے، نہ~$!B$ وچ، تے نہ~$!A$ وچ۔ ایہنوں \emph{خاص متغیر دی شرط} آکھدے آں۔

!!^{proof}s، \Log{N1c} تے~\Log{N1i} وچ، سیکوئنٹاں دے محدود درخت نیں جیہڑے مسلّمیاں توں استنباطی قاعدیاں نال استقرائی طور اُتے بندے نیں۔ ہر پتّا اک مفروضہ اے، جیہدے اُتے اخراج دا نشان ہو سکدا اے۔ باقی ہر !!{formula} اپنے عین اُتلے !!{formula}s توں نظام دے کسے استنباطی قاعدے نال نکلدا اے۔ جے کسے استنباط توں اُتلے درخت وچ مفروضے وجوں آئی !!{formula} دی موجودگی نوں اوہدے اُتلے کسے استنباط نے خارج نہیں کیتا، تاں اوہنوں (اوس استنباط اُتے) \emph{کھُلّا} آکھدے آں۔

\begin{defn}\ollabel{defn:proof}
\emph{!!^{proof}s}، \Log{N1c} تے~\Log{N1i} وچ، !!{formula}s دے محدود درخت نیں جیہڑے استقرا نال ایویں بندے نیں:
\begin{enumerate}
\item\ollabel{defn:prf:assumption} ہر !!{formula} جیہدے اُتے $x$ یا $y$ ورگا اخراج دا نشان اے، اپنے آپ وچ !!a{proof} اے۔ !!a{proof} وچ جیہڑا اکّو !!{formula} دا اے، !!{formula} کھُلّا مفروضہ منّیا جاندا اے۔
\item\ollabel{defn:prf:one-prem} جے $R$ اک، دو یا تِن مقدمیّاں $!A_1$، $!A_2$، $!A_3$ والا قاعدہ اے جیہدا نتیجہ~$!A$ اے، تے $\delta_1$، $\delta_2$، $\delta_3$ اوہ !!{proof}s نیں جیہڑے بالترتیب $!A_1$، $!A_2$، $!A_3$ اُتے مُکدے نیں، تاں
\[
  \AxiomC{}
  \RightLabel{$\delta_1$}
  \DeduceC{$!A_1$}
  \RightLabel{$R$}
  \UnaryInfC{$!A$}
  \DisplayProof
\qquad
  \AxiomC{}
  \RightLabel{$\delta_1$}
  \DeduceC{$!A_1$}
  \AxiomC{}
  \RightLabel{$\delta_2$}
  \DeduceC{$!A_2$}
  \RightLabel{$R$}
  \BinaryInfC{$!A$}
  \DisplayProof
\qquad
  \AxiomC{}
  \RightLabel{$\delta_1$}
  \DeduceC{$!A_1$}
  \AxiomC{}
  \RightLabel{$\delta_2$}
  \DeduceC{$!A_2$}
  \AxiomC{}
  \RightLabel{$\delta_3$}
  \DeduceC{$!A_2$}
  \RightLabel{$R$}
  \TrinaryInfC{$!A$}
  \DisplayProof
  \]
بالترتیب !!a{proof} اے، جے !!{proof}s دے مفروضیاں دے نشان آپس وچ میل کھاندے ہون (دو وکھرے !!{formula}s دا اکّو اخراج نشان نہ ہووے) تے استنباط~$R$ دی درست ورتوں ہون۔
\item نویں !!{proof} وچ سبھ توں اُتلیاں !!{formula} دیاں موجودگیاں !!{proof} دے \emph{کھُلّے مفروضے} منّیاں جاندیاں نیں جے اوہ $\delta_1$، $\delta_2$ یا~$\delta_3$ دے کھُلّے مفروضے ہون، قاعدے نال خارج کیتیاں موجودگیاں توں چھُٹ۔ جے قاعدے دا نشان~$x$ (یا $x\, y$) اے، تاں \Intro{\lif} تے~\Intro{\lnot} لئی ایہ~$\delta_1$ وچ~$x$ نشان والے سارے کھُلّے مفروضے نیں۔ \Elim{\lexists} لئی ایہ~$x$ نشان والے سارے کھُلّے مفروضے نیں~$\delta_2$ وچ۔ \Elim{\lor} لئی ایہ~$x$ نشان والے سارے کھُلّے مفروضے نیں~$\delta_2$ وچ تے~$y$ نشان والے سارے کھُلّے مفروضے نیں~$\delta_3$ وچ۔
\end{enumerate}
\end{defn}

!!{proof}s جیہڑے !!a{proof} دی استقرائی اُساری وچ ورتے جاندے نیں، اوہدے \emph{ذیلی !!{proof}s} آکھے جاندے نیں۔ کسے استنباط دے مقدمیّاں اُتے مُکّن والے ذیلی !!{proof}s نوں اوس استنباط دے \emph{سدھے ذیلی !!{proof}s} آکھدے آں۔ \Log{N2c} تے \Log{N2i} دے قاعدے \olref{tab:N1} وچ دِتّے نیں۔

% OLP-0666 content/proof-theory/natural-deduction/rules-N1.tex
\begingroup

\begin{table}
\begin{center}
\begin{tabular}[t]{@{}cc@{}}
\bottomAlignProof
\AxiomC{$\Discharge{!A}{x}$}
\shortDeduce
\DeduceC{$!B$}
\DischargeRule{\Intro{\lif}}{x}
\UnaryInfC{$!A \lif !B$}
\DisplayProof
&
\bottomAlignProof
\AxiomC{$!A \lif !B$}
\AxiomC{$!A$}
\RightLabel{\Elim{\lif}}
\BinaryInfC{$!B$}
\DisplayProof
\\[4ex]
\bottomAlignProof
\AxiomC{$!A$}
\AxiomC{$!B$}
\RightLabel{\Intro{\land}}
\BinaryInfC{$!A \land !B$}
\DisplayProof
&
\begin{tabular}[b]{@{}l@{}}
\AxiomC{$!A \land !B$}
\RightLabel{\Elim{\land}[1]}
\UnaryInfC{$!A$}
\DisplayProof
\\[3ex]
\bottomAlignProof
\AxiomC{$!A \land !B$}
\RightLabel{\Elim{\land}[2]}
\UnaryInfC{$!B$}
\DisplayProof
\end{tabular}
\\[4ex]
\begin{tabular}{@{}l@{}}
\AxiomC{$!A$}
\RightLabel{\Intro{\lor}[1]}
\UnaryInfC{$!A \lor !B$}
\DisplayProof
\\[3ex]
\AxiomC{$!B$}
\RightLabel{\Intro{\lor}[2]}
\UnaryInfC{$!A \lor !B$}
\DisplayProof
\end{tabular}
&
\AxiomC{$!A \lor !B$}
\AxiomC{$\Discharge{!A}{x}$}
\shortDeduce
\DeduceC{$!C$}
\AxiomC{$\Discharge{!B}{y}$}
\shortDeduce
\DeduceC{$!C$}
\DischargeRule{\Elim{\lor}}{x\,y}
\TrinaryInfC{$!C$}
\DisplayProof
\bottomAlignProof
\\[6ex]
\bottomAlignProof
\AxiomC{$\Discharge{!A}{x}$}
\shortDeduce
\DeduceC{$\lfalse$}
\DischargeRule{\Intro{\lnot}}{x}
\UnaryInfC{$\lnot !A$}
\DisplayProof
&
\bottomAlignProof
\AxiomC{$\lnot !A$}
\AxiomC{$!A$}
\RightLabel{\Elim{\lnot}}
\BinaryInfC{$\lfalse$}
\DisplayProof
\\[4ex]
\bottomAlignProof
\AxiomC{$!A(a)$}
\RightLabel{\Intro{\lforall}}
\UnaryInfC{$\lforall[x][!A(x)]$}
\DisplayProof
&
\bottomAlignProof
\AxiomC{$\lforall[x][!A(x)]$}
\RightLabel{\Elim{\lforall}}
\UnaryInfC{$!A(t)$}
\DisplayProof
\\[4ex]
\bottomAlignProof
\AxiomC{$!A(t)$}
\RightLabel{\Intro{\lexists}}
\UnaryInfC{$\lexists[x][!A(x)]$}
\DisplayProof
&
\bottomAlignProof
\AxiomC{$\lexists[x][!A(x)]$}
\AxiomC{$\Discharge{!A(a)}{x}$}
\shortDeduce
\DeduceC{$!C$}
\DischargeRule{\Elim{\lexists}}{x}
\BinaryInfC{$!C$}
\DisplayProof
\\[4ex]
\bottomAlignProof
\AxiomC{$\lfalse$}
\RightLabel{\FalseInt}
\UnaryInfC{$!A$}
\DisplayProof
&
\bottomAlignProof
\AxiomC{$\Discharge{\lnot!A}{x}$}
\shortDeduce
\DeduceC{$\lfalse$}
\DischargeRule{\FalseCl}{x}
\UnaryInfC{$!A$}
\DisplayProof
\end{tabular}
\end{center}
\caption{\Log{N1i} تے \Log{N1c} دے قاعدے۔ $\Intro{\lforall}$ وچ
!!{constant}~$a$ نہ نتیجے وچ آوے نہ کسے کھُلّے مفروضے وچ۔ $\Elim{\lexists}$ وچ
$a$ نہ مقدمیّاں وچ آوے نہ کسے کھُلّے مفروضے وچ۔ \Log{N1i} وچ \FalseCl شامل نہیں۔}
\ollabel{tab:N1}
\end{table}
\endgroup
% /OLP-0666

\begin{defn}
\emph{!!{height}}~$\pheight{\delta}$، !!a{proof}~$\delta$ دی، استقرا نال ایویں تعریف ہوندی اے۔
\begin{enumerate}
\item جے $\delta$ بس اک مفروضہ اے، تاں $\pheight{\delta} = 0$۔
\item جے $\delta$ اک مقدمے والے استنباط اُتے مُکدا اے جیہدا سدھا ذیلی !!{proof}~$\delta_1$ اے، تاں $\pheight{\delta} = \pheight{\delta_1} + 1$۔
\item جے $\delta$ دو مقدمیّاں والے استنباط اُتے مُکدا اے جیہدے سدھے ذیلی !!{proof}s~$\delta_1$ تے~$\delta_2$ نیں، تاں $\pheight{\delta} =
  \max(\pheight{\delta_1}, \pheight{\delta_2}) + 1$۔
\item جے $\delta$، \Elim{\lor} اُتے مُکدا اے جیہدے سدھے ذیلی !!{proof}s~$\delta_1$، $\delta_2$ تے~$\delta_3$ نیں، تاں \[\pheight{\delta} = \max(\pheight{\delta_1}, \pheight{\delta_2}, \pheight{\delta_3}) + 1\]۔
\end{enumerate}
جے سیکوئنٹ~$S$ دا !!a{proof} اے جیہدی اُچائی $\le n$ اے، تاں اسیں $\Proves[n] S$ لکھدے آں۔
\end{defn}

\begin{ex}
\Log{N1i} وچ ہر !!{formula} اپنے آپ توں، کھُلّے مفروضے وجوں، !!a{proof} منّیا جاندا اے۔ سو
\[
!A \land !B^x
\]
!!a{proof}~$\delta_1$ اے۔ اوہدی اُچائی~$0$ اے۔ $\delta_1$ توں دو !!{proof}s~$\delta_2$ تے~$\delta_3$ لے سکدے آں،~\Elim{\land} دے دو روپاں وچوں اک لا کے:
\[
\AxiomC{$!A \land !B^x$}
\RightLabel{\Elim{\land}[2]}
\UnaryInfC{$!B$}
\DisplayProof
\qquad
\AxiomC{$!A \land !B^x$}
\RightLabel{\Elim{\land}[1]}
\UnaryInfC{$!A$}
\DisplayProof
\]
آخری استنباط دے سدھے ذیلی !!{proof}s،~$\delta_2$ تے~$\delta_3$ وچ، اوہی نیں: $!A \land !B$۔ دوہاں !!{proof}s وچ $!A\land !B$ دی موجودگی کھُلّے مفروضے وجوں اے۔ ساڈے کول $\pheight{\delta_1} = \pheight{\delta_2} = 1$ اے۔

قاعدے \Intro{\land} نوں $!B$ تے $!A$ دی شکل دے دو مقدمے چاہیدے نیں تاں جو~$!B \land !A$ دی شکل دا نتیجہ ملے۔ سو ہِٹھلا !!a{proof}~$\delta_4$ اے:
\[
\AxiomC{$!A \land !B^x$}
\RightLabel{\Elim{\land}[2]}
\UnaryInfC{$!B$}
\LeftSubproofLabel{$\delta_2$}
\AxiomC{$!A \land !B^x$}
\RightLabel{\Elim{\land}[1]}
\UnaryInfC{$!A$}
\RightSubproofLabel{$\delta_3$}
\RightLabel{\Intro{\land}}
\BinaryInfC{$!B \land !A$}
\DisplayProof
\]
اوہدی اُچائی $\pheight{\delta_4} = \max(\pheight{\delta_2},
\pheight{\delta_3})+1 = \max(1,1) + 1 = 2$ اے۔ اوہدے وچ $!A \land !B$ دیاں دو موجودگیاں مفروضے وجوں نیں، تے دوویں کھُلّیاں نیں۔

آخر \Intro{\lif} لا کے $\delta_5$ لے سکدے آں:
\[
\AxiomC{$\Discharge{!A \land !B}{x}$}
\RightLabel{\Elim{\land}[2]}
\UnaryInfC{$!B$}
\LeftSubproofLabel{$\delta_2$}
\AxiomC{$\Discharge{!A \land !B}{x}$}
\RightLabel{\Elim{\land}[1]}
\UnaryInfC{$!A$}
\RightSubproofLabel{$\delta_3$}
\RightLabel{\Intro{\land}}
\BinaryInfC{$!B \land !A$}
\DischargeRule{\Intro{\lif}}{x}
\UnaryInfC{$(!A \land !B) \lif (!B \land !A)$}
\DisplayProof
\]
اوہدی اُچائی $\pheight{\delta_4} + 1 = 3$ اے۔ \Intro{\lif} استنباط $!A \land !B$ دی شکل والے مفروضیاں دیاں ساریاں موجودگیاں خارج کردا اے جیہناں اُتے~$x$ دا نشان اے، یعنی سدھے ذیلی ثبوت دے پہلاں کھُلّے دوویں مفروضے۔ چونکہ ایہو $\delta_5$ دے اکّو اک مفروضے نیں، اوہدا کوئی کھُلّا مفروضہ نہیں۔ ایس لئی ایہ !!a{proof} اے اپنے آخری !!{formula} \emph{دا}، جیہڑا $(!A \land !B) \lif (!B \land !A)$ اے۔
\end{ex}
\endgroup
% /OLP-0668
% OLP-0669 content/proof-theory/natural-deduction/sequents.tex
\begingroup
\def\olfilename{sequents}
\expandafter\def\csname ol@sectioncs\endcsname{section}

\olfileid{pt}{nat}{seq}

\olsection{سیکوئنٹاں نال قدرتی اخذ۔ \Log{N2c} تے~\Log{N2i}}

!!^a{proof}~$\delta$، \Log{G1c} یا \Log{G1i} وچ، وکھاؤندا اے کہ اوہدا آخری !!{formula}~$!A$ اوہدے کھُلّے مفروضیاں توں نکلدا اے۔ جے $\Gamma$ اوہناں !!{formula}s~$!B$ دا سیٹ اے جیہناں لئی $!B^x$،~$\delta$ وچ کھُلّا مفروضہ اے، تاں ایہ $\delta \Proves \Gamma
\Sequent !A$ لکھ سکدے آں۔ ایس توں لگدا اے کہ سیکوئنٹاں لئی وی قدرتی اخذ دا نظام لکھ سکدے آں۔ ایہو جیہے نظام وچ ہر سیکوئنٹ ایہ دسدا اے کہ سجّے والا !!{formula} کیہڑے کھُلّے مفروضیاں اُتے ٹِکدا اے، یعنی سیکوئنٹ اُتے مُکّن والے ذیلی !!{proof} وچ کیہڑے مفروضے کھُلّے نیں۔ ایہ سیکوئنٹاں نال قدرتی اخذ اے، یا «سیکوئنٹ طرز وچ» قدرتی اخذ۔ حاصل نظاماں نوں \Log{N2c} تے~\Log{N2i} آکھدے آں۔

\Log{N1} تے \Log{N2} دا میل ہور نیڑا کرن لئی، کھبّے پاسے دے سیٹ~$\Gamma$ نشان لگے !!{formula}s~$x:!B$ دے منّ لئیے۔ چونکہ \Log{N1} دے قاعدے بس مقدمیّاں وچلے !!{formula}s اُتے، یعنی آخری !!{formula}s سدھے ذیلی !!{proof}s دے اُتے کم کردے نیں، ایس لئی \Log{N2} دے قاعدے بس سیکوئنٹ دی تالی اُتے کم کردے نیں۔ سو سیکوئنٹاں دے کھبّے نشان لگے فارمولیاں نوں بس سیاق آکھ سکدے آں۔

مفروضیاں دی تھاں \Log{N2} دے ثبوتاں وچ سبھ توں اُتلے سیکوئنٹ ایس شکل دے مسلّمے نیں:
\[x:!A \Sequent !A.\]
قاعدے عین \Log{N1} دے قاعدیاں نال میل کھاندے نیں تے \olref{tab:N2} وچ دِتّے نیں۔

% OLP-0667 content/proof-theory/natural-deduction/rules-N2.tex
\begingroup

\begin{table}
\begin{tabular}{@{}c@{}c@{}}
\Axiom$x: !A, \Gamma \fCenter !B$
\RightLabel{\Intro{\lif}}
\UnaryInf$\Gamma \fCenter !A \lif !B $
\DisplayProof
&
\Axiom$\Gamma_1 \fCenter !A \lif !B$
\Axiom$\Gamma_2 \fCenter !A$
\RightLabel{\Elim{\lif}}
\BinaryInf$\Gamma_1, \Gamma_2 \fCenter !B$
\DisplayProof
\\[4ex]
\Axiom$\Gamma_1 \fCenter !A$
\Axiom$\Gamma_2 \fCenter !B$
\RightLabel{\Intro{\land}}
\BinaryInf$\Gamma_1, \Gamma_2 \fCenter !A \land !B $
\DisplayProof
&
\begin{tabular}[b]{@{}l@{}}
\Axiom$\Gamma \fCenter !A \land !B$
\RightLabel{\Elim{\land}[1]}
\UnaryInf$\Gamma \fCenter !A$
\DisplayProof
\\[3ex]
\Axiom$\Gamma \fCenter !A \land !B$
\RightLabel{\Elim{\land}[2]}
\UnaryInf$\Gamma \fCenter !B$
\DisplayProof\\[3ex]
\end{tabular}
\\[4ex]
\Axiom$\Gamma \fCenter !A$
\RightLabel{\Intro{\lor}[1]}
\UnaryInf$\Gamma \fCenter !A \lor !B$
\DisplayProof
&
\Axiom$ \Gamma \fCenter !B$
\RightLabel{\Intro{\lor}[2]}
\UnaryInf$\Gamma \fCenter !A \lor !B$
\DisplayProof\\[3ex]
\multicolumn{2}{c}{
\begin{tabular}[b]{@{}l@{}}
\Axiom$\Gamma_1 \fCenter !A \lor !B$
\Axiom$x:!A, \Gamma_2 \fCenter !C$
\Axiom$y:!B, \Gamma_3 \fCenter !C$
\RightLabel{\Elim{\lor}}
\TrinaryInf$\Gamma_1, \Gamma_2, \Gamma_3 \fCenter !C$
\DisplayProof
\end{tabular}
}
\\[4ex]
\Axiom$x:!A, \Gamma \fCenter \lfalse$
\RightLabel{\Intro{\lnot}}
\UnaryInf$ \Gamma \fCenter \lnot !A $
\DisplayProof
&
\Axiom$\Gamma_1 \fCenter \lnot !A $
\Axiom$\Gamma_2 \fCenter !A $
\RightLabel{\Elim{\lnot}}
\BinaryInf$\Gamma_1, \Gamma_2 \fCenter \lfalse$
\DisplayProof
\\[4ex]
\Axiom$\Gamma \fCenter !A(a) $
\RightLabel{\Intro{\lforall}}
\UnaryInf$\Gamma \fCenter \lforall[x][!A(x)]$
\DisplayProof
&
\Axiom$\Gamma \fCenter \lforall[x][!A(x)]$
\RightLabel{\Elim{\lforall}}
\UnaryInf$\Gamma \fCenter !A(t)$
\DisplayProof
\\[4ex]
\Axiom$\Gamma \fCenter !A(t) $
\RightLabel{\Intro{\lexists}}
\UnaryInf$\Gamma \fCenter \lexists[x][!A(x)]$
\DisplayProof
&
\Axiom$\Gamma_1 \fCenter  \lexists[x][!A(x)]$
\Axiom$x: !A(a), \Gamma_2 \fCenter !C$
\RightLabel{\Elim{\lexists}}
\BinaryInf$\Gamma_1, \Gamma_2 \fCenter !C$
\DisplayProof
\\[4ex]
\Axiom$\Gamma \fCenter \lfalse$
\RightLabel{\FalseInt}
\UnaryInf$\Gamma \fCenter !A$
\DisplayProof
&
\Axiom$x: \lnot !A, \Gamma \fCenter \lfalse$
\RightLabel{\FalseCl}
\UnaryInf$\Gamma \fCenter !A$
\DisplayProof
\end{tabular}
\caption{\Log{N2c} تے \Log{N2i} دے قاعدے۔ $\Intro{\forall}$ تے
$\Elim{\exists}$ وچ، !!{constant}~$a$ نتیجے دے سیکوئنٹ وچ نہیں آؤنا چاہیدا۔
$\Gamma_i$ نشان لگے !!{formula}s~$x :
!A$ دے سیٹ نیں۔ \Intro{\lforall} تے \Elim{\lexists} قاعدیاں لئی
خاص متغیر دی شرط پوری ہونی چاہیدی اے: $a$ نتیجے دے سیکوئنٹ وچ نہیں آؤنا چاہیدا۔
\Log{N2i} وچ \FalseCl شامل نہیں۔}
\ollabel{tab:N2}
\end{table}
\endgroup
% /OLP-0667

جیویں \Log{N1c} تے \Log{N1i} وچ، سانوں ایہ کھُل دینی پئے گی کہ \Intro{\lif}، \Intro{\lnot}، \Elim{\lor} تے \Elim{\lexists} قاعدے مفروضے خارج کر سکدے نیں پر ضروری نہیں کہ کرن۔ ایس لئی \Log{N2c} تے~\Log{N2i} وچ وی میل کھاندے قاعدیاں دی درست ورتوں اوہدوں منّنے آں جدوں سیاق دے متعلقہ !!{formula}s $!A$، $!B$ تے $!A(a)$ متعلقہ مقدمے دے سیاق وچ موجود نہیں۔

\begin{prop}
جے $\delta$ اک !!a{proof} اے~\Log{N1c} یا \Log{N1i} وچ،~$!A$ دا، تاں اک !!a{proof}~$\delta'$ اے $\Gamma \Sequent !A$ دا،~\Log{N2c} (\Log{N2i}) وچ، جتھے $\Gamma$ سارے~$x:!B$ دا سیٹ اے جیہناں لئی $!B^x$،~$\delta$ دا کھُلّا مفروضہ اے۔
\end{prop}

\begin{proof}
$n=\pheight{\delta}$ اُتے استقرا نال۔ جے $n=0$، تاں $\delta$ اکّو کھُلّا مفروضہ~$!A^x$ اے۔ فیر $\Gamma = \{x:!A\}$ تے $\Gamma
  \Sequent !A$، \Log{N2c} یا~\Log{N2i} دا مسلّمہ اے۔

جے $n>0$، تاں~$\delta$ دے آخری استنباط مطابق صورتاں وکھ کردے آں۔ اسیں بس اک صورت دسدے آں، باقی مشقاں وجوں چھڈیاں نیں۔

منّ لو آخری استنباط $\Intro{\lif}$ اے۔ فیر $\delta$ ایتھے مُکدا اے:
\[
  \AxiomC{$\Discharge{!B}{x}$}
  \RightLabel{$\delta_1$}
  \DeduceC{$!C$}
  \DischargeRule{\Intro{\lif}}{x}
  \UnaryInfC{$!B \lif !C$}
  \DisplayProof
  \]
استقرا دے مفروضے نال !!a{proof}~$\delta_1'$ اے~\Log{N2c} (\Log{N2i}) وچ، $\Gamma' \Sequent !C$ دا، جتھے $\Gamma'$ اوہ نشان لگے مفروضے نیں جیہڑے~$\delta_1$ وچ کھُلّے نیں۔ جے $!B^x$، $\delta_1$ دا کھُلّا مفروضہ اے، تاں \Intro{\lif} استنباط اُتے ایہ خارج ہو جاندا اے تے~$\delta$ دے کھُلّے مفروضے $\Gamma =
  \Gamma' \setminus \{x:!B\}$ نیں۔ دوجے لفظاں وچ، $\delta_1'$، !!a{proof} اے~$x:!B, \Gamma \Sequent !C$ دا تے اسیں $\delta'$ ایہ لے سکدے آں:
\[
  \AxiomC{}
  \RightLabel{$\delta_1'$}
  \Deduce$x:B, \Gamma \fCenter !C$
  \RightLabel{\Intro{\lif}}
  \UnaryInf$\Gamma \fCenter !B \lif !C$
  \DisplayProof
  \]
جے $!B^x$،~$\delta_1$ دا کھُلّا مفروضہ نہیں، تاں \Intro{\lif} کوئی مفروضہ خارج نہیں کردا تے $\delta$ تے $\delta_1$ دے کھُلّے مفروضے اوہی نیں۔ چونکہ سیاق دا فارمولا $x:!B$، \Intro{\lif} دے مقدمے وچ،~\Log{N2c} تے~\Log{N2i} وچ، ہونا ضروری نہیں، ایس لئی $\delta'$ ایہ لے سکدے آں:
\[
  \AxiomC{}
  \RightLabel{$\delta_1'$}
  \Deduce$\Gamma \fCenter !C$
  \RightLabel{\Intro{\lif}}
  \UnaryInf$\Gamma \fCenter !B \lif !C$
  \DisplayProof
  \]

جتھے $\delta$ کسے ہور قاعدے اُتے مُکدا اے، اوہ صورتاں ایہو جیہیاں نیں۔
\end{proof}

\begin{prop}
جے $\delta$ اک !!a{proof} اے \Log{N2c} یا \Log{N2i} وچ، $\Gamma
\Sequent !A$ دا، تاں اک !!a{proof}~$\delta'$ اے \Log{N1c} (\Log{N1i}) وچ، جیہدا آخری !!{formula}~$!A$ اے تے کھُلّے مفروضے $!B^x$ نیں ہر $x:!B \in \Gamma$ لئی۔
\end{prop}

\begin{proof}
فیر !!{height}~$n$ اُتے استقرا نال اگے ودھدے آں،~$\delta$ دی اُچائی اُتے۔ جے $n=0$، تاں $\delta$ مسلّمہ اے، $x:!A \Sequent !A$۔ !!{proof}~$\delta'$ مفروضہ~$!A^x$ اے۔

جے $n>0$، تاں~$\delta$ دے آخری استنباط مطابق صورتاں وکھ کردے آں۔ مثلاً جے استنباط~$\Intro{\lif}$ اے، تاں $\delta$ ایتھے مُکدا اے:
\[
  \bottomAlignProof
  \AxiomC{}
  \RightLabel{$\delta_1$}
  \Deduce$x:!B, \Gamma \fCenter !C$
  \RightLabel{\Intro{\lif}}
  \UnaryInf$\Gamma \fCenter !B \lif !C$
  \DisplayProof
  \quad\text{ or }\quad
  \bottomAlignProof
  \AxiomC{}
  \RightLabel{$\delta_1$}
  \Deduce$\Gamma \fCenter !C$
  \RightLabel{\Intro{\lif}}
  \UnaryInf$\Gamma \fCenter !B \lif !C$
  \DisplayProof
  \]
استقرا دے مفروضے نال !!a{proof} اے~\Log{N1c} (\Log{N1i}) وچ،~$!C$ دا۔ پہلی صورت وچ $!B^x$،~$\delta'$ دا کھُلّا مفروضہ اے۔ دوجی صورت وچ نہیں۔ اسیں !!{proof}~$\delta'$ بالترتیب ایہ لے سکدے آں:
\[
  \bottomAlignProof
  \AxiomC{$\Discharge{!B}{x}$}
  \RightLabel{$\delta_1$}
  \DeduceC{$!C$}
  \DischargeRule{\Intro{\lif}}{x}
  \UnaryInfC{$!B \lif !C$}
  \DisplayProof
  \quad\text{ or }\quad
  \bottomAlignProof
  \AxiomC{}
  \RightLabel{$\delta_1$}
  \DeduceC{$!C$}
  \RightLabel{\Intro{\lif}}
  \UnaryInfC{$!B \lif !C$}
  \DisplayProof
  \]
\end{proof}
\endgroup
% /OLP-0669
% OLP-0665 content/proof-theory/natural-deduction/quantifiers.tex
\begingroup
\def\olfilename{quantifiers}
\expandafter\def\csname ol@sectioncs\endcsname{section}

\olfileid{pt}{nat}{qua}

\olsection{باقاعدہ \usetoken{P}{proof} تے بدلی}

مسوّراں دے قاعدیاں وچ کجھ اوکڑاں «خاص متغیر دیاں شرطاں» کرکے آؤندیاں نیں جیہڑیاں \Elim{\lexists} تے~\Intro{\lforall} قاعدیاں اُتے لگدیاں نیں۔ ایہناں وچوں بہتی اوکڑ دور کیتی جا سکدی اے جے پکا کریئے کہ ایہناں استنباطاں دا !!{constant}~$c$ مُڑ نہ ورتیا جاوے۔ آؤ ہِٹھلی تعریف دیئے۔

\begin{defn}
قدرتی اخذ دا !!{proof}~$\delta$ \emph{باقاعدہ} اے جے
\begin{enumerate}
\item کوئی خاص متغیر~$a$، اک توں ودھ \Elim{\lexists} یا~\Intro{\lforall} استنباط دا خاص متغیر نہیں ورتیا جاندا، تے
\item جے $a$، \Elim{\lexists} یا~\Intro{\lforall} استنباط دا خاص متغیر اے، تاں ایہ بس اوس سدھے ذیلی !!{proof} وچ آؤندا اے جیہڑا، بالترتیب، چھوٹے مقدمے یا مقدمے تک پہنچدا اے۔
\end{enumerate}
\end{defn}

\begin{lem}\ollabel{lem:subst-regular}
جے $\delta$ باقاعدہ اے،~$!C$ اُتے مُکدا اے، $c$ اک !!a{constant} اے جیہڑا~$\delta$ وچ خاص متغیر نہیں ورتیا جاندا، تے $t$ اک ٹرم اے جیہدے وچ~$\delta$ دا کوئی خاص متغیر نہیں، تاں $\Subst{\delta}{t}{c}$، جیہڑا~$\delta$ وچ~$c$ دی ہر موجودگی نوں~$t$ نال بدل کے ملدا اے، اک باقاعدہ !!{proof} اے۔ آخری !!{formula}، $\Subst{\delta}{t}{c}$ دا، $\Subst{!C}{t}{c}$ اے۔ جے $!A^x$،~$\delta$ دا کھُلّا مفروضہ اے، تاں $\Subst{!A}{t}{c}^x$،~$\Subst{\delta}{t}{c}$ دا کھُلّا مفروضہ اے۔
\end{lem}

\begin{proof}
$\pheight{\delta}$ اُتے استقرا نال۔ جے $\pheight{\delta} = 0$، تاں ایہ بس مفروضہ~$!A^x$ اے۔ فیر $\Subst{\delta}{t}{c}$،~$\Subst{!A}{t}{c}^x$ اے۔

جے $\pheight{\delta} > 0$، تاں~$\delta$ دے آخری استنباط مطابق صورتاں وکھ کردے آں۔ قضیائی جوڑاں دے قاعدیاں والیاں صورتاں سوکھیاں نیں تے مشقاں وجوں چھڈیاں نیں۔ اسیں اوہ صورتاں ویکھدے آں جیہناں وچ $\delta$ مسوّر والے استنباط اُتے مُکدا اے۔
\begin{enumerate}
\item جے آخری استنباط $\Intro{\lforall}$ اے، تاں $\delta$ دی شکل ایہ اے:
\[
    \AxiomC{}
    \RightLabel{$\delta'$}
    \DeduceC{$!A(a,c)$}
    \RightLabel{\Intro{\lforall}}
    \UnaryInfC{$\lforall[x][!A(x,c)]$}
    \DisplayProof
    \]
استقرا دے مفروضے نال $\Subst{\delta'}{t}{c}$، $!A(a,t)$ دا باقاعدہ !!{proof} اے۔ چونکہ $a$،~$\delta$ دا خاص متغیر اے، ایس لئی $a$،~$t$ وچ نہیں آؤندا۔ سو $\Subst{\delta}{t}{c}$ دا آخری استنباط،
\[
    \AxiomC{}
    \RightLabel{$\Subst{\delta'}{t}{c}$}
    \DeduceC{$!A(a,t)$}
    \RightLabel{\Intro{\lforall}}
    \UnaryInfC{$\lforall[x][!A(x,t)]$}
    \DisplayProof
    \]
درست اے۔ کیوں جے $t$ وچ~$a$ نہیں، نہ نتیجے وچ تے نہ کسے کھُلّے مفروضے وچ~$a$ اے۔
\item جے آخری استنباط $\Elim{\lforall}$ اے، تاں $\delta$ دی شکل ایہ اے:
\[
    \AxiomC{}
    \RightLabel{$\delta'$}
    \DeduceC{$\lforall[x][!A(x,c)]$}
    \RightLabel{\Elim{\lforall}}
    \UnaryInfC{$!A(s(c),c)$}
    \DisplayProof
    \]
استقرا دے مفروضے نال $\Subst{\delta'}{t}{c}$، $\lforall[x][!A(x,t)]$ دا باقاعدہ !!{proof} اے۔ (نتیجے دے ٹرم~$s(c)$ وچ~$c$ ہو سکدا اے، پر ہونا ضروری نہیں۔) $\Subst{\delta}{t}{c}$ دا آخری استنباط،
\[
    \AxiomC{}
    \RightLabel{$\Subst{\delta'}{t}{c}$}
    \DeduceC{$\lforall[x][!A(x,t)]$}
    \RightLabel{\Elim{\lforall}}
    \UnaryInfC{$!A(s(t),t)$}
    \DisplayProof
    \]
درست اے، تے $!A(s(t),t) \ident \Subst{!A(s(c),c)}{t}{c}$۔ \Intro{\lexists} والی صورت ایہو جیہی اے۔
\item جے آخری استنباط $\Elim{\lexists}$ اے، تاں $\delta$ دی شکل ایہ اے:
\[
    \AxiomC{}
    \RightLabel{$\delta_1'$}
    \DeduceC{$\lexists[x][!A(x,c)]$}
    \AxiomC{$\Discharge{!A(a,c)}{x}$}
    \RightLabel{$\delta_2'$}
    \DeduceC{$!C$}
    \DischargeRule{\Elim{\lexists}}{x}
    \BinaryInfC{$!C$}
    \DisplayProof
    \]
استقرا دے مفروضے نال $\Subst{\delta_1'}{t}{c}$ تے $\Subst{\delta_2'}{t}{c}$ باقاعدہ !!{proof}s نیں، پہلا $\lexists[x][!A(x,c)]$ دا تے دوجا $!C$ دا کھُلّے مفروضے~$!A(a,t)$ توں۔ چونکہ $a$،~$\delta$ دا خاص متغیر اے، ایس لئی $a$،~$t$ وچ نہیں آؤندا۔ سو $\Subst{\delta}{t}{c}$ دا آخری استنباط،
\[
    \AxiomC{}
    \RightLabel{$\Subst{\delta'}{t}{c}$}
    \DeduceC{$!A(a,t)$}
    \RightLabel{\Intro{\lforall}}
    \UnaryInfC{$\lforall[x][!A(x,t)]$}
    \DisplayProof
    \]
درست اے۔ $\Subst{!A(x,t)}{a}{x} \ident \Subst{!A(a,c)}{t}{c}$، تے نہ نتیجے~$!C$ وچ تے نہ کسے کھُلّے مفروضے وچ~$a$ اے۔
\end{enumerate}
\end{proof}

\begin{prop}\ollabel{prop:regular}
جے $\delta$، \Log{N1c} یا \Log{N1i} وچ اک !!a{proof} اے، تاں اوسے بناوٹ والا اک باقاعدہ !!{proof}~$\delta'$ اے (خاص کر اوہی !!{height}) جیہڑا $\delta$ توں بس خاص متغیر بدلّن نال وکھرا اے۔
\end{prop}

\begin{proof}
بس ایس ثبوت لئی،~$\delta$ دے کسے خاص متغیر والے استنباط نوں \emph{صاف} آکھو جے اوہدا خاص متغیر کسے ہور استنباط دا خاص متغیر نہیں تے~$\delta$ وچ اوس استنباط دے سدھے ذیلی !!{proof} توں چھُٹ کِتے وی نہیں آؤندا۔ نہیں تاں اوہنوں \emph{میلا} آکھو۔ منّ لو $n$،~$\delta$ وچ میلے خاص متغیر والے استنباطاں دی گنتی اے۔ اسیں~$n$ اُتے استقرا نال وکھاؤنے آں کہ $\delta$ نوں لوڑ دے باقاعدہ !!{proof}~$\delta'$ وچ بدلیا جا سکدا اے۔ جے $n=0$، تاں~$\delta$ وچ سارے خاص متغیر والے استنباط صاف نیں، سو $\delta$ باقاعدہ اے تے ثابت کرن نوں کجھ نہیں۔ نہیں تاں سبھ توں اُتلا خاص متغیر والا استنباط چُنو، منّ لو \Intro{\lforall}۔ذیلی !!{proof}~$\delta_1$،~$\delta$ دا جیہڑا ایس اُتے مُکدا اے، ایس شکل دا اے:
\[
    \AxiomC{}
    \RightLabel{$\delta_2$}
    \DeduceC{$!A(c)$}
    \RightLabel{\Intro{\lforall}}
    \UnaryInfC{$\lforall[x][!A(x)]$}
    \DisplayProof
    \]
دھیان دیو، چونکہ $c$ خاص متغیر اے، ایہ~$\lforall[x][!A(x)]$ یا کسے کھُلّے مفروضے وچ نہیں آؤندا۔ ثبوت $\delta_2$ باقاعدہ اے، کیوں جے اوہدے وچ کوئی خاص متغیر والا استنباط نہیں۔ منّ لو $a$ اک !!a{constant} اے جیہڑا~$\delta$ وچ نہیں آؤندا۔ \olref{lem:subst-regular} نال $\Subst{\delta_2}{a}{c}$، $!A(a)$ دا باقاعدہ ثبوت اے۔ خاص متغیر دی شرط نال~$\delta_2$ دے کسے کھُلّے مفروضے وچ~$c$ نہیں، سو~$\Subst{\delta_2}{a}{c}$ دے کسے کھُلّے مفروضے وچ~$a$ نہیں۔ ایس لئی~$\Intro{\lforall}$ لا سکدے آں۔ جے $\delta_1$ نوں~$\delta$ وچ ثبوت $\delta_1'$ نال بدلیے،
\[
    \AxiomC{}
    \RightLabel{$\Subst{\delta_2}{a}{c}$}
    \DeduceC{$!A(a)$}
    \RightLabel{\Intro{\lforall}}
    \UnaryInfC{$\lforall[x][!A(x)]$}
    \DisplayProof
    \]
تاں اک میلا خاص متغیر والا استنباط اک صاف استنباط نال بدل دِتّا اے، تے اکّو فرق ایہ اے کہ خاص متغیر $c$ دی تھاں~$a$ آ گیا اے۔ حاصل ثبوت وچ $n-1$ میلے خاص متغیر والے استنباط نیں، تے نتیجہ استقرا دے مفروضے توں نکلدا اے۔
\end{proof}

\begin{prob}
\olref{prop:regular} دے ثبوت دی اوہ صورت دسو جتھے سبھ توں اُتلا خاص متغیر والا استنباط~\Elim{\lexists} اے۔
\end{prob}

اسیں ہُنے ثابت کیتی گل نوں کھُلھ کے حوالے وجوں نہیں ورتاں گے، پر منّاں گے کہ زیرِ غور سارے !!{proof}s باقاعدہ نیں۔ جے نہیں، تاں خاص متغیر بدل کے اوہناں نوں ہمیشہ باقاعدہ بنا سکدے آں۔

\Olref{lem:subst-regular} تے \olref{prop:regular}، \Log{N2c} تے~\Log{N2i} لئی وی سچ نیں۔ ثبوت مشقاں وجوں چھڈدے آں۔
\endgroup
% /OLP-0665
% OLP-0662 content/proof-theory/natural-deduction/grafting.tex
\begingroup
\def\olfilename{grafting}
\expandafter\def\csname ol@sectioncs\endcsname{section}

\olfileid{pt}{nat}{gra}

\olsection{\usetoken{P}{proof} جوڑنے}

!!^{proof}s، \Log{N1c} تے~\Log{N1i} وچ، !!{proof}s نیں !!{formula}s~$!A$ دے، مفروضیاں~$!B$ توں۔ منّ لو $!B$ دا اپنا !!a{proof} اے۔!!{proof}~$!B$ دا، !!{proof} نال $!A$ دے~$!B$ توں، کِویں جوڑیئے کہ !!a{proof} ملے~$!A$ دا جیہڑا مفروضہ~$!B$ نہ ورتے؟

اک ڈھنگ ایہ اے کہ مفروضہ~$!B$ نوں !!{proof} وچوں،~$!A$ دے،~\Intro{\lif} لا کے ہٹائیے۔ فیر حاصل !!{proof}،~$!B \lif !A$ دا، !!{proof} نال،~$!B$ دے، جوڑ کے ایہ ملدا اے:
\[
\AxiomC{$\Discharge{!B}{x}$}
\DeduceC{$!A$}
\DischargeRule{\Intro{\lif}}{x}
\UnaryInfC{$!B \lif !C$}
\AxiomC{}
\DeduceC{$!B$}
\RightLabel{\Elim{\lif}}
\BinaryInfC{$!A$}
\DisplayProof
\]
ایہ سدھا اے پر ذرا بھدّا اے۔ $\lif$ نوں شامل ای کیوں کریئے جے اوہنوں جھٹ ہٹا دینا اے؟ $!B \lif !A$ وچوں لنگھن والے ایہو جیہے چکر توں بچ سکنا چاہیدا اے۔ (ایہو جیہے چکراں توں ہمیشہ بچ سکنا ای نارمل روپ دے قضیے دا مضمون اے۔)

\Log{N1c} تے~\Log{N1i} وچ ایس صورت وچ اسیں ایہ چکر سوکھے ای ٹال سکدے آں۔ مفروضیاں~$!B^x$ دی تھاں بس پورا !!{proof} رکھ دیو~$!B$ دا۔ دوجے ثبوتاں دے مفروضیاں دی تھاں !!{proof}s «جوڑنا» کافی کم آؤندا اے، سو ایہنوں تعریف وچ لکھدے آں۔

\begin{defn}
جے $\delta_1$ اک \Log{N1c}-!!{proof} (یا \Log{N1i}-!!{proof}) اے،~$!A$ دا کھُلّے مفروضے~$!B^x$ توں، تے $\delta_2$ اک \Log{N1c}-!!{proof} (\Log{N1i}-!!{proof}) اے~$!B$ دا، تاں $\Subst{\delta_1}{\delta_2}{!B^x}$ اوہ حاصل اے جیہڑا~$!B^x$ دی ہر اوہ موجودگی، جیہڑی خارج نہیں ہوئی،~$\delta_1$ وچ~$\delta_2$ نال بدل کے ملدا اے۔
\end{defn}

جے $\delta_1$ تے $\delta_2$ وچ وکھرے مفروضہ !!{formula}s اُتے اکّو نشان نہ لگیا ہووے، تے~$\delta_2$ دے کسے کھُلّے مفروضے وچ~$\delta_1$ دا کوئی خاص متغیر نہ ہووے، تاں حاصل $\Subst{\delta_1}{\delta_2}{!B^x}$ درست !!{proof} اے۔ اوہدے کھُلّے مفروضے $\delta_1$ تے~$\delta_2$ دے کھُلّے مفروضے رَل کے نیں۔ !!{proof}s جوڑدیاں ایہ شرطاں عام طور اُتے پوریاں ہوندیاں نیں۔ پہلی شرط جدوں وی $\delta_1$ تے $\delta_2$ وڈے !!{proof} دے ذیلی ثبوت ہون، پوری ہوندی اے۔ جے دوجی شرط پوری نہ ہووے تاں اسیں~$\delta_1$ دے خاص متغیراں دے ناں بدل کے ایہنوں پوری کر سکدے آں۔

ایہو جیہی تعریف \Log{N2c} تے~\Log{N2i} اُتے وی لگدی اے۔

\begin{defn}
جے $\delta_1$ اک \Log{N2c}-!!{proof} (یا \Log{N2i}-!!{proof}) اے،~$x:!B, \Gamma_1 \Sequent !A$ دا، تے $\delta_2$ اک \Log{N1c}-!!{proof} (\Log{N1i}-!!{proof}) اے~$\Gamma_2 \Sequent !B$ دا، تاں $\Subst{\delta_1}{\delta_2}{x:!B}$ اوہ حاصل اے جیہڑا ہر مسلّمہ $x:!B \Sequent !B$،~$\delta_1$ وچ،~$\delta_2$ نال بدل کے تے ایہو جیہے مسلّمیاں دے تھلّے ہر سیکوئنٹ دے سیاق وچ $\Gamma_2$ شامل کر کے ملدا اے۔
\end{defn}

\begin{prop}
جے $\delta_1 \Proves x:!B, \Gamma_1 \Sequent
!A$ تے $\delta_2 \Proves \Gamma_2 \Sequent !B$، تاں
$\Subst{\delta_1}{\delta_2}{x:!B} \Proves \Gamma_1, \Gamma_2 \Sequent
!A$، بشرطیکہ~$\delta_1$ دا کوئی خاص متغیر~$\Gamma_2$ وچ نہ آوے۔
\end{prop}

\begin{proof}
$\pheight{\delta}$ اُتے استقرا نال۔
\end{proof}
\endgroup
% /OLP-0662
% OLP-0670 content/proof-theory/natural-deduction/translation-G2i.tex
\begingroup
\def\olfilename{translation-G2i}
\expandafter\def\csname ol@sectioncs\endcsname{section}

\olfileid{pt}{nat}{tgi}

\olsection{\Log{G2i} توں \Log{N2i} وچ ثبوت بدلنے}

\Log{G2i} دے سیکوئنٹاں دے مقدم بہُ-سیٹ~$\Gamma$ نیں، !!{formula}s دے، جدکہ \Log{N2i} دے سیکوئنٹاں دے مقدم نشان لگے فارمولیاں دے سیٹ~$\Gamma'$ نیں جیہناں وچ ہر نشان بس اک واری آؤندا اے۔ آؤ آکھیے کہ سیٹ~$\Gamma'$، نشان لگے !!{formula}s دا، بہُ-سیٹ~$\Gamma$ نال \emph{میل کھاندا اے} جے ہر !!{formula}~$!A$ لئی،~$\Gamma'$ دے !!{element}s جیہڑے $x : !A$ دی شکل دے نیں، اوہناں دی گنتی~$!A$ دی تکرار گنتی توں گھٹ یا برابر اے (یعنی~$!A$ دیاں نکلاں دی گنتی توں،~$\Gamma$ وچ)۔

\begin{prop}\ollabel{prop:G2i-to-N2i}
جے $\Log{G2i} + \Cut \Proves \Gamma \Sequent \Delta$، تاں $\Log{N2i}
\Proves \Gamma' \Sequent !C$، جتھے $!C \ident \lfalse$ جے $\Delta$ خالی اے تے $\Delta = \{!A\}$ باقی صورت وچ، تے $\Gamma'$،~$\Gamma$ نال میل کھاندا اے۔
\end{prop}

\begin{proof}
اسیں وکھاؤنے آں کہ جے $\pi$ اک !!a{proof} اے $\Gamma \Sequent \Delta$ دا، \Log{G2ci} وچ، تاں~$\Gamma'$ موجود اے جیہڑا~$\Gamma$ نال میل کھاندا اے تے !!a{proof}~$\delta$،~$\Gamma' \Sequent !C$ دا،،~\Log{N2i} وچ، موجود اے۔ ایہ $n = \pheight{\pi}$ اُتے استقرا نال کردے آں۔

جے $n = 0$، تاں $\pi$ مسلّمہ اے۔ جے ایہ مسلّمہ $\lfalse \Sequent
  \quad$ اے، تاں $\delta$ مسلّمہ $x:\lfalse \Sequent \lfalse$ اے، یعنی $\Gamma' = \{x:\lfalse\}$ تے $!C \ident \lfalse$۔ جے $\pi$، $!A \Sequent !A$ دی شکل دا مسلّمہ اے، تاں $\delta$، $x: !A \Sequent !A$ اے۔

جے $n>0$، تاں $\pi$ قاعدے~$R$ اُتے مُکدا اے۔ استقرا دا مفروضہ پکا کردا اے کہ اوس قاعدے دے مقدمیّاں نال میل کھاندے سیکوئنٹاں دے !!{proof}s،~\Log{N2c} وچ، نیں۔ اسیں منّ سکدے آں کہ اوہناں ثبوتاں وچ سارے اخراج دے نشان وکھرے نیں، تے جے نشان $x$ کِتے خارج ہووے، تاں بس اوہنوں خارج کرن والے استنباط توں اُتّے آؤندا اے (لوڑ پئے تاں نشان بدل کے)۔ فیر وکھاؤنے آں کہ ایہ !!{proof}s کِویں جوڑ سکدے آں تاں جو ڈھکدے $\Gamma' \Sequent !C$ دے ثبوت ملن۔~$R$ مطابق صورتاں وکھ کردے آں۔

\begin{enumerate}
\item جے $R$، \RightR{\Weakening} اے، تاں $\pi$ ایتھے مُکدا اے:
\[
    \AxiomC{}
    \RightLabel{$\pi_1$}
    \Deduce$\Gamma \fCenter$
    \RightLabel{\RightR{\Weakening}}
    \UnaryInf$\Gamma \fCenter !A$
    \DisplayProof
    \]
$\pi_1$ اُتے استقرا دا مفروضہ لا کے !!a{proof} ملدا اے~$\Gamma' \Sequent \lfalse$ دا۔ $\FalseInt$ لا کے !!a{proof} لیندے آں $\Gamma' \Sequent !A$ دا۔
\item جے $R$، \LeftR{\Weakening} اے، تاں $\pi$ ایتھے مُکدا اے:
\[
    \AxiomC{}
    \RightLabel{$\pi_1$}
    \Deduce$\Gamma \fCenter \Delta$
    \RightLabel{\RightR{\Weakening}}
    \UnaryInf$!A, \Gamma \fCenter \Delta$
    \DisplayProof
    \]
$\pi_1$ اُتے استقرا دا مفروضہ لا کے !!a{proof} ملدا اے~$\Gamma' \Sequent !C$ دا، جیہنوں $\delta$ لیندے آں۔ دھیان دیو، جے $\Gamma'$، $\Gamma$ نال میل کھاندا اے، تاں $!A,
    \Gamma$ نال وی، کیوں جے $x : !A$ دی گنتی $\Gamma'$ وچ،~$!A$ دیاں نکلاں دی گنتی $\Gamma$ وچ توں ودھ نہیں، سو $!A$ دیاں نکلاں دی گنتی $\Gamma \cup \{!A\}$ وچ توں وی ودھ نہیں۔
\item جے $R$، \LeftR{\Contraction} اے، تاں $\pi$ ایتھے مُکدا اے:
\[
    \AxiomC{}
    \RightLabel{$\pi_1$}
    \Deduce$!A, !A, \Gamma \fCenter \Delta$
    \RightLabel{\LeftR{\Contraction}}
    \UnaryInf$!A, \Gamma \fCenter \Delta$
    \DisplayProof
    \]
$\pi_1$ اُتے استقرا دا مفروضہ لا کے !!a{proof} ملدا اے~$\Pi, \Gamma' \Sequent !C$ دا، جتھے $\Pi \subseteq \{x : !A,
    y:!A\}$۔ جے $\Pi = \{x : !A, y:!A\}$، تاں سارے نشان لگے !!{formula}s $y:!A$،~$\pi_1$ وچ، $x:!A$ نال بدلو۔ چونکہ $x$،~$\pi_1$ دے آخری سیکوئنٹ دے سیاق وچ آؤندا اے، اسیں منّ سکدے آں کہ~$\pi_1$ وچ کوئی استنباط!!a{formula}،~$x:!B$ دی شکل دا، خارج نہیں کردا، یعنی کوئی $x:!B$،~$\pi_1$ دے کسے سیکوئنٹ دے کھبّے پاسے، سرگردہ نہیں۔ ایس لئی حاصل ہن وی !!a{proof} اے~\Log{N2i} وچ۔
\item جے $R$، $\RightR{\lif}$ اے، تاں $\pi$ ایتھے مُکدا اے:
\[
    \AxiomC{}
    \RightLabel{$\pi_1$}
    \Deduce$!A, \Gamma \fCenter !B$
    \RightLabel{\RightR{\lif}}
    \UnaryInf$\Gamma \fCenter !A \lif !B$
    \DisplayProof
    \]
استقرا دے مفروضے نال !!a{proof}~$\delta_1$ اے، $x:
    !A, \Gamma' \Sequent !B$ یا $\Gamma' \Sequent !B$ دا، جتھے $\Gamma'$،~$\Gamma$ نال میل کھاندا اے۔ $\delta_1$ نوں $\delta$ تے \Intro{\lif} دی اک ورتوں ملا کے لے سکدے آں۔
\item جے $R$، \LeftR{\lif} اے، تاں $\pi$ ایتھے مُکدا اے:
\[
    \AxiomC{}
    \RightLabel{$\pi_1$}
    \Deduce$\Gamma_1 \fCenter !A$
    \AxiomC{}
    \RightLabel{$\pi_2$}
    \Deduce$!B, \Gamma_2 \fCenter \Delta$
    \RightLabel{\LeftR{\lif}}
    \BinaryInf$!A \lif !B, \Gamma_1, \Gamma_2 \fCenter \Delta$
    \DisplayProof
    \]
استقرا دا مفروضہ !!{proof}s دیندا اے، $\delta_1$، $\Gamma_1'
    \Sequent !A$ دا، تے $\delta_2$، $x: !B, \Gamma_2' \Sequent !C$ یا $\Gamma_2' \Sequent !C$ دا۔ جے $\delta_2$ پچھلی شکل دا اے، تاں $\delta_2$ نوں~$\delta$ لے سکدے آں۔ نہیں تاں~$\pi_1$ وچ سارے فارمولیاں تے اخراج دے نشان بدل کے پکا کر سکدے آں کہ کوئی نشان $\delta_1$ تے~$\delta_2$ دوہاں وچ نہ آوے۔ فیر $\Gamma_1' \cup
    \Gamma_2'$، $\Gamma_1 \cup \Gamma_2$ نال میل کھاندا اے۔ چونکہ $x:!B$،~$\delta_2$ دے آخری سیکوئنٹ وچ آؤندا اے، $x:!B$،~$x$ اخراج نشان والے کسے استنباط وچ،~$\delta_2$ اندر، سرگردہ نہیں ہو سکدا۔ ہر مسلّمہ $x:!B \Sequent !B$،~$\delta_2$ وچ، ایس !!{proof}~$\delta_3$ نال بدلو:
\[
    \Axiom$x: !A \lif !B \fCenter !A \lif !B$
    \AxiomC{}
    \RightLabel{$\delta_1$}
    \Deduce$\Gamma_1' \fCenter !A$
    \RightLabel{\Elim{\lif}}
    \BinaryInf$x: !A \lif !B \fCenter !B$
    \DisplayProof
    \]
تے $x:!B$ دی ہر موجودگی،~$\delta_2$ وچ، $x:!A \lif !B$ نال بدلو۔ حاصل $\Subst{\delta_2}{\delta_3}{x:!B}$، !!a{proof} اے $x:!A \lif !B, \Gamma_1', \Gamma_2' \Sequent !C$ دا۔
\item جے $R$، \RightR{\land} اے، تاں $\pi$ ایتھے مُکدا اے:
\[
    \AxiomC{}
    \RightLabel{$\pi_1$}
    \Deduce$\Gamma_1 \fCenter !A$
    \AxiomC{}
    \RightLabel{$\pi_2$}
    \Deduce$\Gamma_2 \fCenter !B$
    \RightLabel{\RightR{\land}}
    \BinaryInf$\Gamma_1, \Gamma_2 \fCenter !A \land !B$
    \DisplayProof
    \]
استقرا دا مفروضہ ثبوت $\delta_1$ دیندا اے $\Gamma_1'
    \Sequent !A$ دا تے $\delta_2$ دیندا اے $\Gamma_2' \Sequent !A$ دا۔~$\pi_2$ وچ سارے فارمولیاں تے اخراج دے نشان بدل کے پکا کر سکدے آں کہ کوئی نشان $\delta_1$ تے $\delta_2$ دوہاں وچ نہ آوے۔ فیر $\Gamma_1' \cup \Gamma_2'$، $\Gamma_1 \cup
    \Gamma_2$ نال میل کھاندا اے۔ !!a{proof}~$\delta$ لیندے آں، $\Intro{\land}$ نوں $\delta_1$ تے~$\delta_2$ دے آخری سیکوئنٹاں اُتے لا کے۔
\item جے $R$، $\LeftR{\land}$ اے، تاں مثلاً $\pi$ ایتھے مُکدا اے:
\[
    \AxiomC{}
    \RightLabel{$\pi_1$}
    \Deduce$!A, \Gamma \fCenter \Delta$
    \RightLabel{\LeftR{\land}}
    \UnaryInf$!A \land !B, \Gamma \fCenter \Delta$
    \DisplayProof
    \]
استقرا دے مفروضے نال !!a{proof}~$\delta_1$ اے، $x:
    !A, \Gamma' \Sequent !C$ یا $\Gamma' \Sequent !C$ دا، جتھے $\Gamma'$،~$\Gamma$ نال میل کھاندا اے۔ پچھلی صورت وچ $\delta_1$ نوں $\delta$ لے سکدے آں۔ پہلی صورت وچ ویکھو، چونکہ $x:!A$ آخری سیکوئنٹ وچ آؤندا اے، $x:!A$،~$x$ اخراج نشان والے کسے استنباط وچ،~$\delta_1$ اندر، سرگردہ نہیں۔ ہر مسلّمہ $x:!A
    \Sequent !A$ نوں ایس !!{proof}~$\delta_2$ نال بدلو:
\[
    \Axiom$x: !A \land !B \fCenter !A \land !B$
    \RightLabel{\Elim{\land}}
    \UnaryInf$x: !A \land !B \fCenter !A$
    \DisplayProof
    \]
تے $x:!A$ دی ہر موجودگی $x:!A \land !B$ نال بدلو، یعنی $\Subst{\delta_1}{\delta_2}{x:!A}$ ویکھو۔ ایہ !!a{proof} اے $x:!A
    \land !B, \Gamma' \Sequent !C$ دا۔
\item جے $R$، \Cut اے، تاں $\pi$ ایتھے مُکدا اے:
\[
    \AxiomC{}
    \RightLabel{$\pi_1$}
    \Deduce$\Gamma_1 \fCenter !A$
    \AxiomC{}
    \RightLabel{$\pi_2$}
    \Deduce$!A, \Gamma_2 \fCenter \Delta$
    \RightLabel{\Cut}
    \BinaryInf$\Gamma_1, \Gamma_2 \fCenter \Delta$
    \DisplayProof
    \]
استقرا دا مفروضہ ثبوت $\delta_1$ دیندا اے $\Gamma_1'
    \Sequent !A$ دا تے $\delta_2$ دیندا اے $x:!A, \Gamma_2' \Sequent !C$ یا $\Gamma_2' \Sequent !C$ دا۔ پچھلی صورت وچ $\delta$ نوں~$\delta_2$ لو۔ نہیں تاں ہر مسلّمہ $x:!A \Sequent !A$، $\delta_2$ وچ، $\delta_1$ نال بدلو تاں جو~$\delta$ ملے، یعنی !!{proof}~$\Subst{\delta_2}{\delta_1}{x:!A}$، $\Gamma_1',
    \Gamma_2' \Sequent !C$ دا۔
\end{enumerate}
\end{proof}
\endgroup
% /OLP-0670
% OLP-0671 content/proof-theory/natural-deduction/translation-N2i.tex
\begingroup
\def\olfilename{translation-N2i}
\expandafter\def\csname ol@sectioncs\endcsname{section}

\olfileid{pt}{nat}{tni}

\olsection{\Log{N2} توں \Log{G2} وچ ثبوت بدلنے}

\begin{prop}\ollabel{prop:N2-to-G2}
جے $\Log{N2c}\ (\Log{N2i}) \Proves
\Gamma \Sequent !A$، تاں $\Log{G2c}\ (\Log{N2i}) + \Cut \Proves \Gamma'
\Sequent \Delta$، جتھے $\Gamma'$ اوہ !!{formula}s دا بہُ-سیٹ اے جیہڑا $\Gamma$ توں نشان ہٹا کے ملدا اے، تے $\Delta = \{!A\}$ جے $!A$،~$\lfalse$ نہیں، تے $\Delta = \emptyset$ جے ایہ اوہی اے۔
\end{prop}

\begin{proof}
ثبوت فیر !!{proof}~$\delta$ دی اُچائی اُتے استقرا نال اے، جیہڑا $\Gamma \Sequent !A$ دا اے۔ جے $\delta$ مسلّمہ $x:!A \Sequent !A$ اے، تاں $\pi$ اوہدے نال میل کھاندا مسلّمہ $!A \Sequent !A$ اے، یا $\lfalse \Sequent \ $ جے $!A \ident \lfalse$۔

جے $\delta$ کسے استنباطی قاعدے اُتے مُکدا اے، تاں صورتاں وکھ کردے آں:
\begin{enumerate}
\item جے $R$، \Intro{\lif} اے جیہدا مُکھ فارمولا $!B \lif !C$ اے، تے $x:!B$ مقدمے دے سیاق وچ موجود اے پر نتیجے دے سیاق وچ نہیں، تاں $\delta$ ایتھے مُکدا اے:
\[
  \AxiomC{}
  \RightLabel{$\delta_1$}
  \Deduce$x:!B, \Gamma \fCenter !C$
  \RightLabel{\Intro{\lif}}
  \UnaryInf$\Gamma \fCenter !B \lif !C$
  \DisplayProof
  \]
استقرا دا مفروضہ !!a{proof}~$\pi_1$ دیندا اے، $!B, \Gamma'
  \Sequent !C$ دا (یا $!B, \Gamma'
  \Sequent \ $ دا جے $!C \ident \lfalse$)۔ ایہدے توں~\RightR{\lif} لا کے $\pi$ لیندے آں (\RightR{\Weakening} لاؤن پچھوں، جے $!C \ident \lfalse$)۔
\item جے $R$، \Intro{\lif} اے، پر $x:!B$ مقدمے دے سیاق وچ موجود نہیں، تاں $\delta$ ایتھے مُکدا اے:
\[
  \AxiomC{}
  \RightLabel{$\delta_1$}
  \Deduce$\Gamma \fCenter !C$
  \RightLabel{\Intro{\lif}}
  \UnaryInf$\Gamma \fCenter !B \lif !C$
  \DisplayProof
  \]
استقرا دا مفروضہ !!a{proof}~$\pi_1$ دیندا اے، $\Gamma'
  \Sequent !C$ دا۔ ایہدے توں $\pi$ لیندے آں، \LeftR{\Weakening} نوں~$!B$ نال لا کے، تے فیر~\RightR{\lif} لا کے۔
\item جے $R$، \Elim{\lif} اے، تاں $\delta$ ایتھے مُکدا اے:
\[
  \AxiomC{}
  \RightLabel{$\delta_1$}
  \Deduce$\Gamma_1 \fCenter !B \lif !A$
  \AxiomC{}
  \RightLabel{$\delta_2$}
  \Deduce$\Gamma_2 \fCenter !B$
  \RightLabel{\Elim{\lif}}
  \BinaryInf$\Gamma_1, \Gamma_2 \fCenter !A$
  \DisplayProof
  \]
جتھے $\Gamma = \Gamma_1 \cup \Gamma_2$۔ استقرا دا مفروضہ !!{proof}s دیندا اے، $\pi_1$، $\Gamma_1' \Sequent !B \lif !A$ دا، تے $\pi_2$، $\Gamma_2' \Sequent !B$ دا۔ !!{proof}~$\pi$ ایویں لیندے آں:
\[
  \AxiomC{}
  \RightLabel{$\pi_1$}
  \Deduce$\Gamma_1' \fCenter !B \lif !A$
  \AxiomC{}
  \RightLabel{$\pi_2$}
  \Deduce$\Gamma_2' \fCenter !B$
  \Axiom$!A \fCenter !A$
  \RightLabel{\LeftR{\lif}}
  \BinaryInf$!B \lif !A, \Gamma_2' \fCenter !A$
  \RightLabel{\Cut}
  \BinaryInf$\Gamma_1', \Gamma_2' \fCenter !A$
  \DisplayProof
  \]
بہُ-سیٹ $\Gamma_1' \cup \Gamma_2'$ شاید $\Gamma' =
  (\Gamma_1 \cup \Gamma_2)'$ دے برابر نہ ہووے۔ مثلاً جے $\Gamma_1$ تے $\Gamma_2$ دوہاں وچ $x:!C$ اے، تاں $\Gamma_1' \cup \Gamma_2'$ وچ $!C$ دیاں دو نکلاں نیں، پر $(\Gamma_1 \cup \Gamma_2)'$ وچ اکّو نقل اے۔ ڈھکدے \LeftR{\Contraction} استنباط شامل کر کے ایہ ٹھیک کر سکدے آں۔

جے $!A \ident \lfalse$ ہووے، تاں استقرا دا مفروضہ $\delta_2$ اُتے لا کے !!a{proof} ملدا اے $\Gamma_2' \Sequent \ $ دا۔ ایہدے توں !!{proof}~$\pi$ لیندے آں، ڈھکدے \LeftR{\Weakening} استنباط تے اک \RightR{\Weakening}،~$!A$ نال، شامل کر کے۔
\item جے $R$، \FalseCl اے، تاں $\delta$ ایتھے مُکدا اے:
\[
  \AxiomC{}
  \RightLabel{$\delta_1$}
  \Deduce$x: \lnot !A, \Gamma \fCenter \lfalse$
  \RightLabel{\FalseCl}
  \UnaryInf$\Gamma \fCenter !A$
  \DisplayProof
  \]
استقرا دے مفروضے نال !!a{proof} $\pi_1$ اے $\lnot !A,
  \Gamma' \Sequent \ $ دا۔ !!{proof}~$\pi$ ایہ لیندے آں:
\[
  \Axiom$!A \fCenter !A$
  \RightLabel{\RightR{\lnot}}
  \UnaryInf$\fCenter !A, \lnot !A$
  \AxiomC{}
  \RightLabel{$\pi_1$}
  \Deduce$x: \lnot !A, \Gamma' \fCenter $
  \RightLabel{\Cut}
  \BinaryInf$\Gamma' \fCenter !A$
  \DisplayProof
  \]
\end{enumerate}
دھیان دیو، بس \FalseCl{} دی تبدیلی نال !!a{proof}~$\pi$ ملدا اے جیہدے سجّے اک توں ودھ !!{formula} اے۔ یعنی~\Log{N2i} دے !!{proof} دی تبدیلی \Log{G2i}-!!{proof} اے۔
\end{proof}
\endgroup
% /OLP-0671

\OLEndChapterHook
\endgroup
% /OLP-0664
% OLP-0674 content/proof-theory/normalization/normalization.tex
\begingroup
\def\olfilename{normalization}
\expandafter\def\csname ol@sectioncs\endcsname{section}

\olchapter{pt}{nor}{نارمل روپ وچ لیاؤنا}

% OLP-0672 content/proof-theory/normalization/introduction.tex
\begingroup
\def\olfilename{introduction}
\expandafter\def\csname ol@sectioncs\endcsname{section}

\olfileid{pt}{nor}{int}

\olsection{جان پچھان}

جے تسی شرطیہ~$!A \lif !B$ !!{prove} کر چکے او، تاں ہو سکدا اے کہ !!a{proof}~$!B$ دا دین ویلے~\Elim{\lif} لا کے ایہنوں ورت سکو۔ ایس لئی!!a{proof}~$\delta_1$،~$!A$ دا، وی چاہیدا ہووے گا۔ تہاڈا !!{proof}~$\delta_1$،~$!A \lif !B$ دا، غالباً~\Intro{\lif} اُتے مُکدا اے، جیہڑا !!a{proof}~$\delta_2$،~$!B$ دا کھُلّے مفروضے~$!A$ توں،~$!A \lif !B$ دے ثبوت وچ بدلدا اے۔ جے ایویں اے، تاں تہاڈے آخری !!{proof} دی شکل ایہ اے:
\[
\AxiomC{$\Discharge{!A}{x}$}
\RightLabel{$\delta_2$}
\DeduceC{$!B$}
\DischargeRule{\Intro{\lif}}{x}
\UnaryInfC{$!A \lif !B$}
\AxiomC{}
\RightLabel{$\delta_2$}
\DeduceC{$!A$}
\RightLabel{\Elim{\lif}}
\BinaryInfC{$!B$}
\DisplayProof
\]
بھاویں تہاڈا طریقہ ایس لحاظ نال چنگا اے کہ تسی پہلاں توں موجود شے مُڑ ورتی (!!a{proof}~$!A \lif !B$ دا)، تہاڈا !!{proof} اوہنا چنگا نہیں۔ $!A \lif !B$ نوں شامل کرنا تے جھٹ ہٹا دینا اک چکر لگدا اے۔~$!B$ نوں !!{prove} کرن دا ہور سدھا طریقہ ہونا چاہیدا اے۔

اصل وچ ایہو جیہا طریقہ ہمیشہ ہوندا اے۔ اُتلے ورگے «چکر»، جتھے !!a{introduction} قاعدے دے جھٹ پچھوں !!a{elimination} قاعدہ آؤندا اے، قدرتی اخذ دے !!{proof}s وچوں ہمیشہ ہٹائے جا سکدے نیں۔ ایس نوں \emph{نارمل روپ وچ لیاؤنا} آکھدے آں تے حاصل اک \emph{نارمل} !!{proof} (یا !!a{proof} \emph{نارمل روپ وچ}) اے۔

ایس نتیجے دا ثبوت، سیکوئنٹ حساب دے قطع ہٹاؤن والے قضیے وانگ، ذرا گنجل دار اے تے کئی صورتاں ویکھنیاں پیندیاں نیں۔ کلاسیکی \Log{N1c} لئی ثابت کرنا وجدانی \Log{N1i} نالوں اوکھا اے۔ جیویں قطع ہٹاؤنا وکھاؤندا اے کہ سیکوئنٹ حساب وچ ہور کجھ !!{prove} نہیں کر سکدے، \Cut{} دا قاعدہ ورت کے، اویویں نارمل روپ وچ لیاؤنا وکھاؤندا اے کہ چکر پا کے اوس توں ودھ نہیں ثابت کر سکدے جِنّا چکر ٹال کے۔ ہور وی سانجھاں نیں۔ نارمل !!{proof} اوسے نتیجے دے چکراں والے سبھ توں چھوٹے !!{proof} توں بہت لمّا ہو سکدا اے، جیویں سیکوئنٹ حساب وچ قطعاں والے !!{proof}s قطع توں بناں سبھ توں چھوٹے !!{proof} توں بہت چھوٹے ہو سکدے نیں۔ ذیلی !!{formula} والی خاصیت، کہ تسی ہر نتیجہ ہمیشہ !!{prove} کر سکدے او صرف اوس نتیجے دے ذیلی !!{formula}s نوں اپنے !!{proof} وچ ورت کے، قدرتی اخذ لئی نارمل روپ وچ لیاؤن توں نکلدی اے جیویں سیکوئنٹ حساب لئی قطع ہٹاؤن توں نکلدی اے۔
\endgroup
% /OLP-0672
% OLP-0677 content/proof-theory/normalization/segments.tex
\begingroup
\def\olfilename{segments}
\expandafter\def\csname ol@sectioncs\endcsname{section}

\olfileid{pt}{nor}{seg}

\olsection{ٹوٹے تے قطع}

!!^a{introduction} قاعدے توں پچھوں !!a{elimination} قاعدہ آوے، تاں ایہ !!a{proof} وچ اک چکر اے۔ !!a{proof} نوں نارمل روپ وچ لیاؤن دا مطلب ایہو جیہے چکر ہٹاؤنا اے۔ پر \Elim{\lor} تے \Elim{\lexists} کر کے «اوہ چکر جیہڑا اسیں ہٹاؤنا چاہنے آں» دی تعریف وی ذرا گنجل دار ہو جاندی اے۔ اوکھ ایہ اے کہ ایہناں قاعدیاں کر کے کسے چکر وچ ضروری نہیں کہ !!{elimination} قاعدہ !!{introduction} قاعدے دے \emph{جھٹ پچھوں} آوے۔ وچکار \Elim{\lor} یا \Elim{\lexists} آ سکدا اے، جیویں:
\[
\AxiomC{}
\DeduceC{$\lexists[x][!C(x)]$}
\AxiomC{$\Discharge{!A}{x}$}
\DeduceC{$!B$}
\DischargeRule{\Intro{\lif}}{x}
\UnaryInfC{$!A \lif !B$}
\DischargeRule{\Elim{\lexists}}{y}
\BinaryInfC{$!A \lif !B$}
\AxiomC{}
\DeduceC{$!A$}
\RightLabel{\Elim{\lif}}
\BinaryInfC{$!B$}
\DisplayProof
\]
اصل وچ کِنّے وی \Elim{\lor} یا \Elim{\lexists} قاعدے \Intro{\lif} تے \Elim{\lif} دے وچکار آ سکدے نیں، جیویں:
\[
\AxiomC{}
\DeduceC{$\lexists[x][!C(x)]$}
\AxiomC{}
\DeduceC{$!D \lor !E$}
\AxiomC{$\Discharge{!A}{x}$}
\DeduceC{$!B$}
\DischargeRule{\Intro{\lif}}{x}
\UnaryInfC{$!A \lif !B$}
\AxiomC{$\Discharge{!A}{x}$}
\DeduceC{$!B$}
\DischargeRule{\Intro{\lif}}{x}
\UnaryInfC{$!A \lif !B$}
\RightLabel{\Elim{\lor}}
\TrinaryInfC{$!A \lif !B$}
\DischargeRule{\Elim{\lexists}}{y}
\BinaryInfC{$!A \lif !B$}
\AxiomC{}
\DeduceC{$!A$}
\RightLabel{\Elim{\lif}}
\BinaryInfC{$!B$}
\DisplayProof
\]
ایہ مثال وکھاؤندی اے کہ اوہی \Elim{\lif} کئی چکراں وچ وی آ سکدا اے، کیونکہ ایہ اک توں ودھ~\Intro{\lif} قاعدیاں دے پچھوں آؤندا اے۔ ایہ ساریاں صورتاں سمیٹن لئی اسیں \Log{N1i} دے ثبوت وچ !!{formula} دے آؤن دیاں کجھ خاص لڑیاں دی مدد نال چکر دی تعریف دیندے آں۔ ایہناں لڑیاں نوں \emph{ٹوٹے} آکھدے آں۔ ساڈی مثال وچ ایہ~$!A \lif !B$ دے آؤن دی اک لڑی اے، جتھے $\Elim{\lor}$ یا $\Elim{\lexists}$ دا چھوٹا مقدمہ اوہدے نتیجے پچھوں آؤندا اے۔ مثال وچ دو ایہو جیہے ٹوٹے نیں، تے ہر ٹوٹے وچ~$!A \lif !B$ تنّ واری آؤندا اے۔ رسمی تعریف ایہ اے۔

\begin{defn}
\Log{N1i} دے !!{proof}~$\delta$ وچ اک \emph{ٹوٹا} آؤناں دی لڑی $!A_1$، \dots، $!A_n$ اے، جیہڑے اوسے !!{formula}~$!A$ دے $\delta$ وچ آؤن نیں، جتھے
\begin{enumerate}
  \item $!A_1$، \Elim{\lor} یا \Elim{\lexists} دا نتیجہ نہیں؛
  \item $!A_n$، \Elim{\lor} یا \Elim{\lexists} دا چھوٹا مقدمہ نہیں؛
  \item $i = 1$، \dots،~$n-1$ لئی، $!A_i$، \Elim{\lor} یا \Elim{\lexists} دا چھوٹا مقدمہ اے تے $A_{i+1}$ اوس اخذ دا نتیجہ اے۔
\end{enumerate}
اوس ٹوٹے دی \emph{لمبائی}~$n$ اے۔
\end{defn}

ہر ٹوٹا ہٹاؤن جوگا چکر نہیں ہوندا۔ جیہڑے ٹوٹے ہٹاؤن جوگے چکر ہوندے نیں، اوہناں نوں \emph{قطع} آکھدے آں۔

\begin{defn}
اک \emph{قطع} اوہ ٹوٹا اے جیہدے وچ $!A_n$، !!a{elimination} اخذ دا وڈّا مقدمہ اے، تے یا تاں $n=1$ اے تے $!A_1 = !A_n$، !!a{introduction} اخذ دا یا~\FalseInt دا نتیجہ اے، یا فیر~$n>1$ اے۔
\end{defn}

ویکھو کہ قطع ٹوٹے دا !!a{introduction} قاعدے دے نتیجے توں شروع ہونا ضروری نہیں۔ بس ایہ ضروری اے کہ اوہ !!a{elimination} قاعدے دے وڈّے مقدمے اُتے مُکے۔ پر لمبائی~$1$ والے ٹوٹے نوں قطع منّن لئی اوہدا !!{introduction} قاعدیاں دا نتیجہ ہونا ضروری اے۔

\begin{defn}
قطع دا \emph{درجہ}~$\depth{!A}$ اے۔ \emph{قطع درجہ}~$\cutrank{\delta}$، جیہڑا !!a{proof}~$\delta$ دا اے،~$\delta$ وچ قطعاں دا سبھ توں وڈّا درجہ اے۔~$\delta$ وچ کوئی قطع \emph{وڈّے توں وڈّے درجے والا} اے جے اوہدا درجہ~$\cutrank{\delta}$ اے، یعنی اوہدا درجہ سبھ توں وڈّا اے۔~$\delta$ دی \emph{قطع لمبائی} اوہدے سارے وڈّے توں وڈّے درجے والے قطعاں دیاں لمبائیاں دا جوڑ اے۔
\end{defn}
\endgroup
% /OLP-0677
% OLP-0675 content/proof-theory/normalization/permutations.tex
\begingroup
\def\olfilename{permutations}
\expandafter\def\csname ol@sectioncs\endcsname{section}

\olfileid{pt}{nor}{per}

\olsection{تھاں بدلّن والیاں تبدیلیاں}

\Log{N1i} وچ قطع اوہ ٹوٹے نیں جیہڑے !!a{elimination} قاعدے دے وڈّے مقدمے اُتے مُکدے نیں۔ نارمل روپ وچ لیاؤن دے ثبوت وچ اک صورت ایہ اے کہ وڈّے توں وڈّے درجے والے قطعاں دی لمبائی گھٹائیے، تے اوہدے نال !!{proof} دی قطع لمبائی وی گھٹائیے۔ لمبائی $>1$ والا قطع کئی طریقیاں نال مُک سکدا اے۔ ایہ ایس اُتے منحصر اے کہ قطع وچ~$!A$ دی آخری !!{formula} موجودگی کِہڑے قاعدے (\Elim{\lor} یا \Elim{\lexists}) دا نتیجہ اے، تے اوہ کِہڑے !!{elimination} اخذ دا وڈّا مقدمہ اے۔

مثال لئی جے متعلقہ اخذ \Elim{\lor} تے \Elim{\land} نیں، تاں $!A \ident !B \lor !C$ اے تے اک ذیلی !!{proof}~$\delta_1$ اے جیہڑا ایویں مُکدا اے:
\[
\AxiomC{}
\RightLabel{$\delta_2$}
\DeduceC{$!D \lor !E$}
\AxiomC{$\Discharge{!D}{x}$}
\RightLabel{$\delta_3$}
\DeduceC{$!B \land !C$}
\AxiomC{$\Discharge{!E}{y}$}
\RightLabel{$\delta_4$}
\DeduceC{$!B \land !C$}
\DischargeRule{\Elim{\lor}}{x\,y}
\TrinaryInfC{$!B \land !C$}
\RightLabel{\Elim{\land}[1]}
\UnaryInfC{$!B$}
\DisplayProof
\]
قطع وچ آخری !!{formula} موجودگی~$!A_n$،~\Elim{\lor} دا نتیجہ اے، یعنی تھلّا $!B \land !C$۔ اوہدے توں پہلوں !!{formula} موجودگی~$!A_{n-1}$،~\Elim{\lor} دے چھوٹے مقدمیّاں وچوں اک اے۔

ایس قطع دی لمبائی گھٹاؤن لئی اسیں \Elim{\lor} تے \Elim{\land} اخذاں دی ترتیب الٹ سکدے آں، یعنی اوہناں دیاں تھاواں بدل سکدے آں۔ مطلب ایہ کہ ذیلی ثبوت~$\delta_1$ نوں~$\delta$ وچ ایہدے نال بدل دیندے آں:
\[
\AxiomC{}
\RightLabel{$\delta_2$}
\DeduceC{$!D \lor !E$}
\AxiomC{$\Discharge{!D}{x}$}
\RightLabel{$\delta_3$}
\DeduceC{$!B \land !C$}
\RightLabel{\Elim{\land}[1]}
\UnaryInfC{$!B$}
\AxiomC{$\Discharge{!E}{y}$}
\RightLabel{$\delta_4$}
\DeduceC{$!B \land !C$}
\RightLabel{\Elim{\land}[1]}
\UnaryInfC{$!B$}
\DischargeRule{\Elim{\lor}}{x\,y}
\TrinaryInfC{$!B$}
\DisplayProof
\]
جیہڑے قطع \Elim{\land} دے مقدمے اُتے، \Elim{\lor} توں تھلّے، مُکدے سن، اوہ ہُن \Elim{\land} دے اوہناں اخذاں وچوں اک دے مقدمے اُتے مُکدے نیں جیہڑے \Elim{\lor} دے چھوٹے مقدمیّاں توں اُتّے نیں۔ یعنی ہر اک دی لمبائی~$1$ گھٹ گئی اے۔

نالے کوئی وی قطع جیہڑا سارا $\delta_2$، $\delta_3$ یا $\delta_4$ دے اندر، یا~$\delta_1$ توں پورا باہر~$\delta$ وچ اے، اوہدا اوہی رہندا اے۔ اوہدے وچ اوہی !!{formula}s نیں تے اوہ اوہناں ای اخذاں تھانیں لنگھدا اے، سو خاص کر اوہدا درجہ تے لمبائی~$\delta$ وچ اوہدے نال میل کھاندے قطع جِنّے ای نیں۔ سو نویں !!{proof} دا قطع درجہ ودھیا نہیں تے نویں !!{proof} دی قطع لمبائی پرانے !!{proof} دی قطع لمبائی توں گھٹ اے۔

\begin{prob}
تھاں بدلّن والی تبدیلی توں پچھوں گھٹّو گھٹ اک قطع دی لمبائی~$1$ گھٹدی اے۔ جے اسیں \Elim{\land} دی تھاں اک~\Elim{\lor} دے پار بدل دیئیے، تاں اصل وچ دو قطع ٹوٹیاں دی لمبائی گھٹدی اے، ایس لئی ایس صورت وچ !!{proof} دی قطع لمبائی گھٹّو گھٹ~$2$ گھٹدی اے۔ کی ایہو جیہی اک تبدیلی نال !!a{proof} دی قطع لمبائی~$2$ توں ودھ گھٹ سکدی اے؟ کی قطع \emph{درجہ} گھٹ سکدا اے؟
\end{prob}

!!a{elimination} قاعدے نوں \Elim{\lor} یا \Elim{\lexists} قاعدے توں اُتّے لے جان دے عمل نوں \emph{تھاں بدلّن والی تبدیلی} آکھدے آں۔ قاعدیاں دے کجھ جوڑیاں دے نمونے \olref{tab:permutation-conversions} وچ دِتّے نیں۔ (باقی مشقاں لئی چھڈّے گئے نیں۔)

% \begin{table}
% \begin{multline*}
% \shoveleft{\AxiomC{$!D \lor !E$}
% \AxiomC{$\Discharge{!D}{x}$}
% \shortDeduce
% \DeduceC{$!B \land !C$}
% \AxiomC{$\Discharge{!E}{y}$}
% \shortDeduce
% \DeduceC{$!B \land !C$}
% \DischargeRule{\Elim{\lor}}{x\,y}
% \TrinaryInfC{$!B \land !C$}
% \RightLabel{\Elim{\land}[1]}
% \UnaryInfC{$!B$}
% \DisplayProof}\\
% \shoveright{\longrightarrow\ 
% \AxiomC{$!D \lor !E$}
% \AxiomC{$\Discharge{!D}{x}$}
% \shortDeduce
% \DeduceC{$!B \land !C$}
% \RightLabel{\Elim{\land}[1]}
% \UnaryInfC{$!B$}
% \AxiomC{$\Discharge{!E}{y}$}
% \shortDeduce
% \DeduceC{$!B \land !C$}
% \RightLabel{\Elim{\land}[1]}
% \UnaryInfC{$!B$}
% \DischargeRule{\Elim{\lor}}{x\,y}
% \TrinaryInfC{$!B$}
% \DisplayProof}
% \\
% \shoveleft{\AxiomC{$!D \lor !E$}
% \AxiomC{$\Discharge{!D}{x}$}
% \shortDeduce
% \DeduceC{$!B \lif !C$}
% \AxiomC{$\Discharge{!E}{y}$}
% \shortDeduce
% \DeduceC{$!B \lif !C$}
% \DischargeRule{\Elim{\lor}}{x\,y}
% \TrinaryInfC{$!B \lif !C$}
% \AxiomC{$!B$}
% \RightLabel{\Elim{\lif}}
% \BinaryInfC{$!C$}
% \DisplayProof}\\
% \shoveright{\longrightarrow\ 
% \AxiomC{$!D \lor !E$}
% \AxiomC{$\Discharge{!D}{x}$}
% \shortDeduce
% \DeduceC{$!B \lif !C$}
% \AxiomC{$!B$}
% \RightLabel{\Elim{\lif}}
% \BinaryInfC{$!C$}
% \AxiomC{$\Discharge{!E}{y}$}
% \shortDeduce
% \DeduceC{$!B \lif !C$}
% \AxiomC{$!B$}
% \RightLabel{\Elim{\lif}}
% \BinaryInfC{$!C$}
% \DischargeRule{\Elim{\lor}}{x\,y}
% \TrinaryInfC{$!C$}
% \DisplayProof}
% \\
% \shoveleft{\AxiomC{$!D \lor !E$}
% \AxiomC{$\Discharge{!D}{x}$}
% \shortDeduce
% \DeduceC{$!B \lif !C$}
% \AxiomC{$\Discharge{!E}{y}$}
% \shortDeduce
% \DeduceC{$!B \lif !C$}
% \DischargeRule{\Elim{\lor}}{x\,y}
% \TrinaryInfC{$!B \lif !C$}
% \AxiomC{$!B$}
% \RightLabel{\Elim{\lif}}
% \BinaryInfC{$!C$}
% \DisplayProof}\\
% \shoveright{\longrightarrow\ 
% \AxiomC{$!D \lor !E$}
% \AxiomC{$\Discharge{!D}{x}$}
% \shortDeduce
% \DeduceC{$!B \lif !C$}
% \AxiomC{$!B$}
% \RightLabel{\Elim{\lif}}
% \BinaryInfC{$!C$}
% \AxiomC{$\Discharge{!E}{y}$}
% \shortDeduce
% \DeduceC{$!B \lif !C$}
% \AxiomC{$!B$}
% \RightLabel{\Elim{\lif}}
% \BinaryInfC{$!C$}
% \DischargeRule{\Elim{\lor}}{x\,y}
% \TrinaryInfC{$!C$}
% \DisplayProof}
% \\
% \shoveleft{\AxiomC{$!D \lor !E$}
% \AxiomC{$\Discharge{!D}{x}$}
% \shortDeduce
% \DeduceC{$\lforall[x][!B(x)]$}
% \AxiomC{$\Discharge{!E}{y}$}
% \shortDeduce
% \DeduceC{$\lforall[x][!B(x)]$}
% \DischargeRule{\Elim{\lor}}{x\,y}
% \TrinaryInfC{$\lforall[x][!B(x)]$}
% \RightLabel{\Elim{\lforall}}
% \UnaryInfC{$!B(t)$}
% \DisplayProof}\\
% \shoveright{\longrightarrow\ 
% \AxiomC{$!D \lor !E$}
% \AxiomC{$\Discharge{!D}{x}$}
% \shortDeduce
% \DeduceC{$\lforall[x][!B(x)]$}
% \RightLabel{\Elim{\lforall}}
% \UnaryInfC{$!B(t)$}
% \AxiomC{$\Discharge{!E}{y}$}
% \shortDeduce
% \DeduceC{$\lforall[x][!B(x)]$}
% \RightLabel{\Elim{\lforall}}
% \UnaryInfC{$!B(t)$}
% \DischargeRule{\Elim{\lor}}{x\,y}
% \TrinaryInfC{$!B(t)$}
% \DisplayProof}
% \end{multline*}
% \caption{Some permutation conversions for \Log{N1i}.}
% \ollabel{tab:permutation-conversions}
% \end{table}

{\small
\begin{longtable}{@{}l@{}}
\AxiomC{$!D \lor !E$}
\AxiomC{$\Discharge{!D}{x}$}
\shortDeduce
\DeduceC{$!B \land !C$}
\AxiomC{$\Discharge{!E}{y}$}
\shortDeduce
\DeduceC{$!B \land !C$}
\DischargeRule{\Elim{\lor}}{x\,y}
\TrinaryInfC{$!B \land !C$}
\RightLabel{\Elim{\land}[1]}
\UnaryInfC{$!B$}
\DisplayProof
$\longrightarrow$\\[6ex]
\quad\AxiomC{$!D \lor !E$}
\AxiomC{$\Discharge{!D}{x}$}
\shortDeduce
\DeduceC{$!B \land !C$}
\RightLabel{\Elim{\land}[1]}
\UnaryInfC{$!B$}
\AxiomC{$\Discharge{!E}{y}$}
\shortDeduce
\DeduceC{$!B \land !C$}
\RightLabel{\Elim{\land}[1]}
\UnaryInfC{$!B$}
\DischargeRule{\Elim{\lor}}{x\,y}
\TrinaryInfC{$!B$}
\DisplayProof
\\[10ex]
\AxiomC{$!D \lor !E$}
\AxiomC{$\Discharge{!D}{x}$}
\shortDeduce
\DeduceC{$!B \lif !C$}
\AxiomC{$\Discharge{!E}{y}$}
\shortDeduce
\DeduceC{$!B \lif !C$}
\DischargeRule{\Elim{\lor}}{x\,y}
\TrinaryInfC{$!B \lif !C$}
\AxiomC{$!B$}
\RightLabel{\Elim{\lif}}
\BinaryInfC{$!C$}
\DisplayProof
$\longrightarrow$\\[6ex]
\AxiomC{$!D \lor !E$}
\AxiomC{$\Discharge{!D}{x}$}
\shortDeduce
\DeduceC{$!B \lif !C$}
\AxiomC{$!B$}
\RightLabel{\Elim{\lif}}
\BinaryInfC{$!C$}
\AxiomC{$\Discharge{!E}{y}$}
\shortDeduce
\DeduceC{$!B \lif !C$}
\AxiomC{$!B$}
\RightLabel{\Elim{\lif}}
\BinaryInfC{$!C$}
\DischargeRule{\Elim{\lor}}{x\,y}
\TrinaryInfC{$!C$}
\DisplayProof
\\[10ex]
\AxiomC{$!D \lor !E$}
\AxiomC{$\Discharge{!D}{x}$}
\shortDeduce
\DeduceC{$!B \lif !C$}
\AxiomC{$\Discharge{!E}{y}$}
\shortDeduce
\DeduceC{$!B \lif !C$}
\DischargeRule{\Elim{\lor}}{x\,y}
\TrinaryInfC{$!B \lif !C$}
\AxiomC{$!B$}
\RightLabel{\Elim{\lif}}
\BinaryInfC{$!C$}
\DisplayProof
$\longrightarrow$\\[6ex]
\AxiomC{$!D \lor !E$}
\AxiomC{$\Discharge{!D}{x}$}
\shortDeduce
\DeduceC{$!B \lif !C$}
\AxiomC{$!B$}
\RightLabel{\Elim{\lif}}
\BinaryInfC{$!C$}
\AxiomC{$\Discharge{!E}{y}$}
\shortDeduce
\DeduceC{$!B \lif !C$}
\AxiomC{$!B$}
\RightLabel{\Elim{\lif}}
\BinaryInfC{$!C$}
\DischargeRule{\Elim{\lor}}{x\,y}
\TrinaryInfC{$!C$}
\DisplayProof
\\[10ex]
\AxiomC{$!D \lor !E$}
\AxiomC{$\Discharge{!D}{x}$}
\shortDeduce
\DeduceC{$\lforall[x][!B(x)]$}
\AxiomC{$\Discharge{!E}{y}$}
\shortDeduce
\DeduceC{$\lforall[x][!B(x)]$}
\DischargeRule{\Elim{\lor}}{x\,y}
\TrinaryInfC{$\lforall[x][!B(x)]$}
\RightLabel{\Elim{\lforall}}
\UnaryInfC{$!B(t)$}
\DisplayProof
\quad$\longrightarrow$\\[8ex]
\quad\AxiomC{$!D \lor !E$}
\AxiomC{$\Discharge{!D}{x}$}
\shortDeduce
\DeduceC{$\lforall[x][!B(x)]$}
\RightLabel{\Elim{\lforall}}
\UnaryInfC{$!B(t)$}
\AxiomC{$\Discharge{!E}{y}$}
\shortDeduce
\DeduceC{$\lforall[x][!B(x)]$}
\RightLabel{\Elim{\lforall}}
\UnaryInfC{$!B(t)$}
\DischargeRule{\Elim{\lor}}{x\,y}
\TrinaryInfC{$!B(t)$}
\DisplayProof
\\
\caption{\Log{N1i} لئی تھاں بدلّن دیاں کجھ تبدیلیاں۔}
\ollabel{tab:permutation-conversions}
\end{longtable}
}

\begin{prob}
جتھے \Elim{\lor} دے پچھوں \Elim{\lor} یا \Elim{\lnot} آؤندا اے، اوتھے تھاں بدلّن والیاں تبدیلیاں دیو۔
\end{prob}

\begin{prob}
جتھے \Elim{\lexists} دے پچھوں \Elim{\exists} آؤندا اے، اوتھے تھاں بدلّن والیاں تبدیلیاں دیو۔ دسو کہ پچھلی صورت وچ خاص متغیر دی کوئی شرط کیوں نہیں ٹُٹدی۔
\end{prob}

سانوں یقین کر لینا چاہیدا اے کہ تھاں بدلّن والی تبدیلی لاؤن ویلے اسیں کِتے خاص متغیر دی کوئی شرط نہ توڑ دیئیے۔ ایہ ویکھو:
\[
\AxiomC{}
\RightLabel{$\delta_2$}
\DeduceC{$!D \lor !E$}
\AxiomC{$\Discharge{!D}{x}$}
\RightLabel{$\delta_3$}
\DeduceC{$\lexists[x][!B(x)]$}
\AxiomC{$\Discharge{!E}{y}$}
\RightLabel{$\delta_4$}
\DeduceC{$\lexists[x][!B(x)]$}
\DischargeRule{\Elim{\lor}}{x\,y}
\TrinaryInfC{$\lexists[x][!B(x)]$}
\AxiomC{$\Discharge{!B(c)}{z}$}
\RightLabel{$\delta_5$}
\DeduceC{$!C$}
\DischargeRule{\Elim{\lexists}}{z}
\BinaryInfC{$!C$}
\DisplayProof
\]
جے اسیں ایہنوں ایہدے نال بدلیے
\[
\AxiomC{}
\RightLabel{$\delta_2$}
\DeduceC{$!D \lor !E$}
\AxiomC{$\Discharge{!D}{x}$}
\RightLabel{$\delta_3$}
\DeduceC{$\lexists[x][!B(x)]$}
\AxiomC{$\Discharge{!B(c)}{z}$}
\RightLabel{$\delta_5$}
\DeduceC{$!C$}
\DischargeRule{\Elim{\lexists}}{z}
\BinaryInfC{$!C$}
\AxiomC{$\Discharge{!E}{y}$}
\RightLabel{$\delta_4$}
\DeduceC{$\lexists[x][!B(x)]$}
\AxiomC{$\Discharge{!B(c)}{z}$}
\RightLabel{$\delta_5$}
\DeduceC{$!C$}
\DischargeRule{\Elim{\lexists}}{z}
\BinaryInfC{$!C$}
\DischargeRule{\Elim{\lor}}{x\,y}
\TrinaryInfC{$!C$}
\DisplayProof
\]
تاں فکر ہو سکدی اے کہ خاص متغیر~$c$ شاید، منّ لو،~$!D$ وچ آؤندا ہووے۔ آخر مفروضہ $!D$ اصل~$\Elim{\lexists}$ توں اُتّے خارج ہو جاندا اے، ایس لئی اصل !!{proof} وچ ایہ~\Elim{\exists} اخذ دے وڈّے مقدمے توں اُتّے کھُلّے کسے مفروضے وچ نہیں آؤندا، تے اوتھے خاص متغیر دی شرط پوری اے۔ پر نویں !!{proof} وچ $!D$ نوں~\Elim{\lexists} اخذاں توں پچھوں ای خارج کیتا جاندا اے، سو خاص متغیر دی شرط ٹُٹ جاؤ گی۔ پر چیتے رکھو کہ اسیں منّیا اے کہ ساڈے !!{proof}s باقاعدہ نیں۔ ہر خاص متغیر اکّو اخذ دا خاص متغیر اے تے !!{proof} وچ صرف \Elim{\lexists} دے چھوٹے مقدمے توں اُتّے آؤندا اے۔ خاص کر $c$، $\delta_3$ یا~$\delta_4$ وچ نہیں آ سکدا۔

ایہ مثال اک ہور گل وی وکھاؤندی اے۔ ویکھو کہ نویں !!{proof} وچ~$\delta_5$ دیاں دو نقلاں نیں۔ ایہدا مطلب اے کہ جے $\delta_5$ وچ وڈّے توں وڈّے درجے والا قطع اے، تاں اسیں قطع لمبائی \emph{نہیں} گھٹائی۔ ایس لئی تھاں بدلّن والی تبدیلی صرف اوہدوں لائیے جدوں سانوں پتا ہووے کہ $\delta_5$ وچ وڈّے توں وڈّے درجے والا کوئی قطع نہیں۔ ایہ پکا کرن لئی اسیں ہمیشہ وڈّے توں وڈّے درجے والا \emph{سبھ توں اُتّلا} قطع چُناں گے، یعنی وڈّے توں وڈّے درجے والا ایہو جیہا قطع جیہدے ذیلی !!{proof} وچ (ایس صورت وچ،~$\delta_1$) وڈّے توں وڈّے درجے والا کوئی قطع اوس ذیلی !!{proof} دے آخری اخذ توں اُتّے نہیں مُکدا۔ ایس نال~$\delta_5$ وچ (یا $\delta_2$، $\delta_3$ یا~$\delta_4$ وچ وی) ہور وڈّے توں وڈّے درجے والے قطع ہون دی گنجائش نہیں رہندی۔
\endgroup
% /OLP-0675
% OLP-0676 content/proof-theory/normalization/reductions.tex
\begingroup
\def\olfilename{reductions}
\expandafter\def\csname ol@sectioncs\endcsname{section}

\olfileid{pt}{nor}{red}

\olsection{تقلیل والیاں تبدیلیاں}

جے \Log{N1i} دے !!a{proof} وچ کسے قطع دی لمبائی~$1$ اے، تاں اوہ ٹوٹا اکّو !!{formula} دی موجودگی توں بنیا اے، تے اوہ !!{formula}، !!a{introduction} اخذ دا نتیجہ وی اے تے !!a{elimination} اخذ دا وڈّا مقدمہ وی۔ نارمل روپ وچ لیاؤن دے ثبوت وچ ایہو جیہے قطع «گھٹائے» جاندے نیں۔ جیہڑا ذیلی !!{proof}، !!{introduction}/!!{elimination} دے ایہو جیہے جوڑے اُتے مُکدا اے، اوہدی تھاں نواں !!{proof} لیاندا جاندا اے جیہدے وچ ایہ اخذاں دا جوڑا ہٹایا گیا ہووے۔ ایس تبدیلی نال نویں قطع بن سکدے نیں۔ پر اسیں ویکھ سکدے آں کہ نویں قطعاں دا درجہ گھٹائے جا رہے قطع توں گھٹ اے۔ ایس لئی نویں !!{proof} دی یا تاں قطع لمبائی گھٹ اے (کیونکہ لمبائی~$1$ والا وڈّے توں وڈّے درجے دا اک قطع ہٹ گیا اے)، یا !!{proof} دا قطع درجہ گھٹ گیا اے (جے اوہو ای وڈّے توں وڈّے درجے والا قطع سی)۔

مثال لئی جے متعلقہ اخذ \Intro{\lif} تے \Elim{\lif} نیں، تاں $!A \ident !B \lif !C$ اے تے اک ذیلی !!{proof}~$\delta_1$ اے جیہڑا ایویں مُکدا اے:
\[
\AxiomC{$\Discharge{!B}{x}$}
\RightLabel{$\delta_2$}
\DeduceC{$!C$}
\DischargeRule{\Intro{\lif}}{x}
\UnaryInfC{$!B \lif !C$}
\AxiomC{}
\RightLabel{$\delta_3$}
\DeduceC{$!B$}
\RightLabel{\Elim{\lif}}
\BinaryInfC{$!C$}
\DisplayProof
\]
ایہ قطع ہٹاؤن لئی اسیں ذیلی ثبوت~$\delta_1$ نوں~$\delta$ وچ ایہدے نال بدل دیندے آں:
\[
\AxiomC{}
\RightLabel{$\delta_3$}
\DeduceC{$!B$}
\RightLabel{$\Subst{\delta_2}{\delta_3}{!B^x}$}
\DeduceC{$!C$}
\DisplayProof
\]
$ \Subst{\delta_2}{\delta_3}{!B^x}$ توں ساڈا مطلب اوہ !!{proof} اے جیہڑا~$\delta_2$ توں ملدا اے، جدوں ہر اوہ مفروضہ~$!B$ جیہدے اُتے~$x$ دا لیبل اے، پورے !!{proof}~$\delta_3$ نال بدل دِتّا جاوے۔ ایس !!{proof} وچ~$!B$ دی شکل دا~$x$ لیبل والا کوئی کھُلّا مفروضہ ہُن نہیں رہندا۔~$x$ لیبل والا کوئی ہور مفروضہ وی نہیں، کیونکہ ہمیشہ وانگ اسیں منّیا اے کہ کوئی دو !!{formula}s اکّو لیبل نال مفروضے دے طور اُتے نہیں آؤندے۔ سو نویں !!{proof} دے اوہی کھُلّے مفروضے نیں جیہڑے~$\delta_1$ دے نیں۔~$\delta_2$ وچ جیہڑا وی اخذ مفروضے خارج کردا اے، اوہ~$\Subst{\delta_2}{\delta_3}{!B^x}$ وچ اوہی مفروضے خارج کردا اے۔ تے~$\delta_3$ وچ مفروضے خارج کرن والے ہر اخذ دیاں کئی نقلاں بن سکدیاں نیں جیہڑیاں~$\Subst{\delta_2}{\delta_3}{!B^x}$ وچ میل کھاندے مفروضے خارج کردیاں نیں۔

کوئی وی قطع جیہڑا سارا $\delta_2$، $\delta_3$ یا $\delta_4$ دے اندر، یا~$\delta_1$ توں پورا باہر~$\delta$ وچ اے، اوہدا اوہی رہندا اے۔ اوہدے وچ اوہی !!{formula}s نیں تے اوہ اوہناں ای اخذاں تھانیں لنگھدا اے، سو خاص کر اوہدا درجہ تے لمبائی~$\delta$ وچ اوہدے نال میل کھاندے قطع جِنّے ای نیں۔ اسیں لمبائی~$1$ والا وڈّے توں وڈّے درجے دا اک قطع ہٹایا اے، ایس لئی اسیں یا قطع لمبائی گھٹائی اے، یا قطع درجہ گھٹایا اے (جے~$\delta$ دی قطع لمبائی~$1$ سی تے ایہو ای وڈّے توں وڈّے درجے والا قطع سی)۔

پر دو گلاں ہو سکدیاں نیں:
\begin{enumerate}
  \item سارے $!B^x$ مفروضے $ \delta_3$ نال بدل کے اسیں~$\delta_3$ دے سارے قطعاں دیاں نقلاں ودھا دِتّیاں نیں۔ جے اوہناں وچوں کجھ وڈّے توں وڈّے درجے والے قطع سن، تاں نویں !!{proof} دی قطع لمبائی پرانے !!{proof} دی قطع لمبائی توں ودھ ہو سکدی اے۔ سانوں پکا کرنا پینا اے کہ ایہ نہ ہووے۔ ایس لئی تقلیل والیاں تبدیلیاں صرف \emph{سبھ توں سجّے} وڈّے توں وڈّے درجے والے قطعاں اُتے لاؤنے آں، یعنی ایہو جیہے قطعاں اُتے جنہاں لئی~$\delta$ دی کوئی شاخ، گھٹائے جا رہے قطع توں سجّے، وڈّے توں وڈّے درجے دا کوئی قطع نہیں رکھدی۔ جے~$\delta_3$ وچ کوئی قطع ہووے، تاں اوہ ایتھے گھٹائے جا رہے قطع توں سجّے ہووے گا، ایس لئی اسیں سبھ توں سجّے قطع اُتے کم نہیں کر رہے ہوواں گے۔ ہمیشہ سبھ توں سجّے قطع چُن کے اسیں پکا کردے آں کہ یا !!{proof} دی قطع لمبائی گھٹائیے یا قطع درجہ وی گھٹا دیئیے۔
  \item جے $!B^x$ مفروضہ $ \delta_2$ وچ !!a{elimination} اخذ دا وڈّا مقدمہ اے، تے~$\delta_3$ دا آخری اخذ \Elim{\lor}، \Elim{\lexists} یا !!a{introduction} اخذ اے، تاں ایہو جیہے مفروضے دی تھاں~$\delta_3$ نوں $ \delta_2$ نال جوڑ کے اسیں اک نواں قطع بنایا اے۔ جیہڑا~$\delta_2$ وچ صرف مفروضہ سی، ہُن یا لمبائی~$1$ دا قطع اے (!!a{introduction} اخذ دا نتیجہ تے !!a{elimination} اخذ دا وڈّا مقدمہ) یا قطع ٹوٹے وچ !!{formula} دی آخری موجودگی اے (!!a{elimination} اخذ دا وڈّا مقدمہ تے~\Elim{\lor} یا \Elim{\lexists} دا نتیجہ)۔ جے $B^x$، \Elim{\lor} یا \Elim{\lexists} دا چھوٹا مقدمہ سی، تاں اوہ پہلاں ای کسے ٹوٹے وچ !!{formula} دی پہلی موجودگی سی، تے ہو سکدا اے کہ اوہ قطع ٹوٹا سی۔ نویں !!{proof} وچ اوہ ٹوٹا ہُن لمّا ہو سکدا اے۔ پر ویکھو کہ ایہ نواں قطع یا پہلاں والا قطع ٹوٹا بناؤن والا !!{formula}~$!B$ اے، تے $\depth{!B} < \depth{!B \lif !C}$۔ سو بھاویں اسیں قطع شامل کیتے یا قطعاں دی لمبائی ودھائی ہووے، ایہناں وچوں کوئی وی وڈّے توں وڈّے درجے والا قطع نہیں۔ ایس طرح اسیں یا قطع درجہ گھٹایا اے، یا، جے اوہو درجے دے قطع باقی نیں، قطع لمبائی گھٹائی اے۔
\end{enumerate}

لمبائی~$1$ والے قطع نوں گھٹاؤن دے عمل نوں \emph{تقلیل والی تبدیلی} (یا \emph{چکر دی تبدیلی}) آکھدے آں۔ قاعدیاں دے کجھ جوڑیاں دے نمونے \olref{tab:reduction-conversions} وچ دِتّے نیں۔ (باقی مشقاں لئی چھڈّے گئے نیں۔)

\begin{table}
\begin{multline*}
\bottomAlignProof
\AxiomC{$\Discharge{!B}{x}$}
\RightLabel{$\delta_2$}
\shortDeduce
\DeduceC{$!C$}
\DischargeRule{\Intro{\lif}}{x}
\UnaryInfC{$!B \lif !C$}
\AxiomC{}
\RightLabel{$\delta_3$}
\shortDeduce
\DeduceC{$!B$}
\RightLabel{\Elim{\lif}}
\BinaryInfC{$!C$}
\DisplayProof
\quad \longrightarrow\quad
\shoveright{\bottomAlignProof
\AxiomC{}
\RightLabel{$\delta_3$}
\shortDeduce
\DeduceC{$!B$}
\RightLabel{$\Subst{\delta_2}{\delta_3}{!B^x}$}
\shortDeduce
\DeduceC{$!C$}
\DisplayProof}
\\[4ex]
\shoveleft{\bottomAlignProof
\AxiomC{}
\RightLabel{$\delta_2$}
\shortDeduce
\DeduceC{$!B$}
\AxiomC{}
\RightLabel{$\delta_3$}
\shortDeduce
\DeduceC{$!C$}
\RightLabel{\Intro{\land}}
\BinaryInfC{$!B \land !C$}
\RightLabel{\Elim{\land}[1]}
\UnaryInfC{$!B$}
\DisplayProof
\quad \longrightarrow\quad
\bottomAlignProof
\AxiomC{}
\RightLabel{$\delta_2$}
\shortDeduce
\DeduceC{$!B$}
\DisplayProof}\qquad
\shoveright{\bottomAlignProof
\AxiomC{}
\RightLabel{$\delta_2$}
\shortDeduce
\DeduceC{$!B$}
\AxiomC{}
\RightLabel{$\delta_3$}
\shortDeduce
\DeduceC{$!C$}
\RightLabel{\Intro{\land}}
\BinaryInfC{$!B \land !C$}
\RightLabel{\Elim{\land}[2]}
\UnaryInfC{$!C$}
\DisplayProof
\quad \longrightarrow\quad
\bottomAlignProof
\AxiomC{}
\RightLabel{$\delta_3$}
\shortDeduce
\DeduceC{$!C$}
\DisplayProof}
\\[4ex]
\bottomAlignProof
\AxiomC{}
\RightLabel{$\delta_2$}
\shortDeduce
\DeduceC{$!B$}
\RightLabel{\Intro{\lor}[1]}
\UnaryInfC{$!B \lor !C$}
\AxiomC{$\Discharge{!B}{x}$}
\RightLabel{$\delta_3$}
\shortDeduce
\DeduceC{$!D$}
\AxiomC{$\Discharge{!C}{y}$}
\RightLabel{$\delta_4$}
\shortDeduce
\DeduceC{$!D$}
\DischargeRule{\Elim{\lor}}{x\,y}
\TrinaryInfC{$!D$}
\DisplayProof
\quad \longrightarrow\quad
\shoveright{\bottomAlignProof
\AxiomC{}
\RightLabel{$\delta_2$}
\shortDeduce
\DeduceC{$!B$}
\RightLabel{$\Subst{\delta_3}{\delta_2}{!B^x}$}
\shortDeduce
\DeduceC{$!D$}
\DisplayProof}
\\[2ex]
\shoveleft{\bottomAlignProof
\AxiomC{}
\RightLabel{$\delta_2$}
\shortDeduce
\DeduceC{$!B(c)$}
\RightLabel{\Intro{\lforall}}
\UnaryInfC{$\lforall[x][!B(x)]$}
\RightLabel{\Elim{\lforall}}
\UnaryInfC{$!B(t)$}
\DisplayProof
\quad \longrightarrow\quad
\bottomAlignProof
\AxiomC{}
\RightLabel{$\Subst{\delta_2}{t}{c}$}
\shortDeduce
\DeduceC{$!B(t)$}
\DisplayProof}
\\[4ex]
\shoveright{\bottomAlignProof
\AxiomC{}
\RightLabel{$\delta_2$}
\shortDeduce
\DeduceC{$!B(t)$}
\RightLabel{\Intro{\lexists}}
\UnaryInfC{$\lexists[x][!B(x)]$}
\AxiomC{$\Discharge{!B(c)}{x}$}
\RightLabel{$\delta_3$}
\shortDeduce
\DeduceC{$!D$}
\DischargeRule{\Elim{\lexists}}{x}
\BinaryInfC{$!D$}
\DisplayProof
\quad \longrightarrow\quad
\bottomAlignProof
\AxiomC{}
\RightLabel{$\delta_2$}
\shortDeduce
\DeduceC{$!B(t)$}
\RightLabel{$\Subst{\Subst{\delta_3}{t}{c}}{\delta_2}{!B(t)^x}$}
\shortDeduce
\DeduceC{$!D$}
\DisplayProof}
\end{multline*}
\caption{\Log{N1i} لئی تقلیل دیاں کجھ تبدیلیاں۔}
\ollabel{tab:reduction-conversions}
\end{table}

\begin{prob}
جتھے \Intro{\lnot} دے پچھوں \Elim{\lnot} آؤندا اے، اوتھے تقلیل والی تبدیلی دیو۔
\end{prob}

اک صورت ہجے ویکھنی باقی اے۔ قطع دی تعریف وچ اوہ صورت وی شامل اے جدوں قطع دی لمبائی~$1$ ہووے تے $!A$،~\FalseInt{} دا نتیجہ وی ہووے تے !!a{elimination} اخذ دا وڈّا مقدمہ وی۔ ایہ صورت اوہو جیہے طریقے نال حل ہوندی اے جیویں جدوں $!A$، !!a{introduction} قاعدے دا نتیجہ ہووے۔ مثال لئی ایہ تبدیلی کرو:
\[
\bottomAlignProof
\AxiomC{}
\RightLabel{$\delta_2$}
\DeduceC{$\lfalse$}
\UnaryInfC{$!B \lif !C$}
\AxiomC{}
\RightLabel{$\delta_3$}
\DeduceC{$!B$}
\RightLabel{\Elim{\lif}}
\BinaryInfC{$!C$}
\DisplayProof
\quad\text{نال}\quad
\bottomAlignProof
\AxiomC{}
\RightLabel{$\delta_2$}
\DeduceC{$\lfalse$}
\UnaryInfC{$!C$}
\DisplayProof
\]
جدوں $!A \ident (!B \lor
!C)$ یا $!A \ident \lexists[x][!B(x)]$ ہووے، سانوں ذرا دھیان کرنا پیندا اے۔ مثال لئی جے اسیں ایہ تبدیلی کریے
\[
\bottomAlignProof
\AxiomC{}
\RightLabel{$\delta_2$}
\DeduceC{$\lfalse$}
\RightLabel{\FalseInt}
\UnaryInfC{$!B \lor !C$}
\AxiomC{$\Discharge{!B}{x}$}
\RightLabel{$\delta_3$}
\DeduceC{$!D$}
\AxiomC{$\Discharge{!C}{x}$}
\RightLabel{$\delta_4$}
\DeduceC{$!D$}
\RightLabel{\Elim{\lor}}
\TrinaryInfC{$!D$}
\DisplayProof
\quad\text{نال}\quad
\bottomAlignProof
\AxiomC{}
\RightLabel{$\delta_2$}
\DeduceC{$\lfalse$}
\RightLabel{\FalseInt}
\UnaryInfC{$!D$}
\DisplayProof
\]
تاں ہو سکدا اے کہ اسیں قطع !!{formula}~$!D$ والا نواں قطع شامل کر رہیے ہوئیے، تے ایہدی کوئی ضمانت نہیں کہ $\depth{!D} < \depth{!B \lor !C}$!{} ایس لئی ایتھے سانوں اوہو طریقہ ورتنا پینا اے جیہڑا اسیں \Intro{\lor}/\Elim{\lor} دی تقلیل لئی ورتیا سی، تے ایہدے نال بدلنا پینا اے:
\[
\bottomAlignProof
\AxiomC{}
\RightLabel{$\delta_2$}
\DeduceC{$\lfalse$}
\RightLabel{\FalseInt}
\UnaryInfC{$!B$}
\RightLabel{$\delta_3$}
\DeduceC{$!D$}
\DisplayProof
\quad\text{یا}\quad
\bottomAlignProof
\AxiomC{}
\RightLabel{$\delta_2$}
\DeduceC{$\lfalse$}
\RightLabel{\FalseInt}
\UnaryInfC{$!C$}
\RightLabel{$\delta_4$}
\DeduceC{$!D$}
\DisplayProof
\]
\endgroup
% /OLP-0676
% OLP-0673 content/proof-theory/normalization/normalization-thm.tex
\begingroup
\def\olfilename{normalization-thm}
\expandafter\def\csname ol@sectioncs\endcsname{section}

\olfileid{pt}{nor}{thm}

\olsection{نارمل روپ وچ لیاؤن دا قضیہ}

!!^a{proof} \emph{نارمل} روپ وچ اے جے اوہدے وچ کوئی قطع ٹوٹا نہیں۔ صاف اے کہ ایہ اُدوں تے صرف اُدوں اے جدوں اسیں اوہدے اُتے نہ تھاں بدلّن والی تبدیلی لا سکدے آں تے نہ تقلیل والی۔ !!a{proof} نوں \emph{نارمل روپ وچ لیاؤن} دا مطلب اے اوہنوں تھاں بدلّن تے تقلیل والیاں تبدیلیاں لا کے نارمل ثبوت وچ بدلنا، جدوں تک کوئی قطع نہیں رہندا۔ ایہ ہمیشہ ہو سکدا اے، ایہو \emph{نارمل روپ وچ لیاؤن دا قضیہ} اے۔

\begin{thm}\ollabel{thm:normalization}
کوئی وی \Log{N1i}-!!{proof} $\delta$ نارمل روپ وچ آ سکدا اے، یعنی قطعاں توں بناں !!a{proof}~$\delta'$ وچ بدلیا جا سکدا اے۔
\end{thm}

\begin{proof}
قطع درجے تے قطع لمبائی اُتے استقرا نال،~$\delta$ دے۔ استقرا دا مفروضہ ایہ اے کہ ہر اوہ !!{proof} جیہدا قطع درجہ گھٹ اے، یا قطع درجہ اوہی پر قطع لمبائی گھٹ اے، نارمل روپ وچ آ سکدا اے۔

جے قطع لمبائی~$0$ اے، تاں کوئی قطع نہیں تے $\delta$ پہلاں ای نارمل اے۔ سو منّ لو~$\delta$ دا قطع درجہ تے قطع لمبائی~$>
0$ نیں۔ سبھ توں اُتلا تے سبھ توں سجّا زیادہ توں زیادہ درجے والا قطع ٹوٹا چُنو۔ جے قطع ٹوٹے دی لمبائی $>1$ اے، تاں تھاں بدلّن والی تبدیلی لاؤ۔ جے قطع ٹوٹے دی لمبائی~$1$ اے، تاں میل کھاندی تقلیل والی تبدیلی لاؤ۔ ایس نال اک نواں !!{proof} ملدا اے جیہدا یا قطع درجہ اوہی اے پر قطع لمبائی گھٹ اے، یا قطع درجہ گھٹ اے۔ استقرا دا مفروضہ لگدا اے۔
\end{proof}

اُتلے ثبوت دا طریقہ (ہمیشہ سبھ توں اُتلا تے سبھ توں سجّا زیادہ توں زیادہ درجے والا قطع لو) ایہ اشارہ دیندا اے کہ !!a{proof} اُتے تبدیلیاں کسے خاص ترتیب نال لاؤنیاں پین گئیاں تاں جو نارمل !!{proof} ملے۔ حیرانی دی گل ایہ اے کہ ترتیب نال کوئی فرق نہیں پیندا۔ تھاں بدلّن تے تقلیل والیاں تبدیلیاں کسے وی ترتیب نال لاؤ، آخر نارمل ثبوت مل جاندا اے۔

ٹوٹے تے قطع دی تعریف \Log{N2i} لئی وی سوکھے ای دِتّی جا سکدی اے، تے \Log{N1i} لئی تھاں بدلّن تے تقلیل والیاں تبدیلیاں سدھیاں \Log{N2i} وچ ورتیاں جا سکدیاں نیں۔ ایس لئی نارمل روپ وچ لیاؤن دا قضیہ \Log{N2i} لئی وی سچ اے۔

\begin{thm}\ollabel{thm:normalization-N2i}
کوئی وی \Log{N2i}-!!{proof} $\delta$ نارمل روپ وچ آ سکدا اے، یعنی قطعاں توں بناں !!a{proof}~$\delta'$ وچ بدلیا جا سکدا اے۔
\end{thm}
\endgroup
% /OLP-0673
% OLP-0678 content/proof-theory/normalization/translations.tex
\begingroup
\def\olfilename{translations}
\expandafter\def\csname ol@sectioncs\endcsname{section}

\olfileid{pt}{nor}{tr}

\olsection{نارمل \Log{N2i}-\usetoken{P}{proof} تے \Log{G2i} وچکار ترجمہ}

!!^{proof}s جیہڑے \Log{N2i} وچ نیں، \Log{G2i} وچ !!{proof}s دے طور اُتے ترجمہ کیتے جا سکدے نیں، تے الٹ پاسے وی۔ پہلاں، قطع توں بناں \Log{2i} !!{proof}s دا ترجمہ کرن نال نارمل \Log{N2i}-!!{proof}s ملدے نیں۔

نتیجہ بیان کرن لئی سانوں کجھ اصطلاحاں چاہیدیاں نیں تاں جو~\Log{N2i} تے~\Log{G2i} وچ سیکوئنٹاں دے کھبّے پاسے آپس وچ سوکھے جوڑ سکیے۔ سیٹ~$\Gamma'$، جیہڑا لیبل والے !!{formula}s دا اے، بغیر لیبل والے !!{formula}s دے کثیر سیٹ~$\Gamma$ \emph{نال میل کھاندا اے} جے ہر !!{formula}~$!A$ لئی، اوہناں !!{element}s دی گنتی جیہڑے~$\Gamma'$ وچ $x : !A$ دی شکل دے نیں،~$!A$ دی~$\Gamma$ وچ نقلاں دی گنتی توں گھٹ یا برابر ہووے۔

\begin{cor}
جے $\Log{G2i} \Proves \Gamma \Sequent \Delta$، تاں نارمل \Log{N2c}-!!{proof}~$\delta$ موجود اے، جیہڑا $ \Gamma' \Sequent !C$ دا ثبوت اے۔ ایتھے سیٹ~$\Gamma'$ لیبل والے !!{formula}s دا اے تے کثیر سیٹ~$\Gamma$ نال میل کھاندا اے، یعنی ہر !!{formula}~$!A$ لئی، اوہناں !!{element}s دی گنتی جیہڑے~$\Gamma'$ وچ $x : !A$ دی شکل دے نیں،~$!A$ دی کثرت (یعنی~$!A$ دیاں نقلاں دی گنتی) توں گھٹ یا برابر اے،~$\Gamma$ وچ۔
\end{cor}

\begin{proof}
\olref[nat][tgi]{prop:G2i-to-N2i} دا ثبوت ویکھ کے ایہ گل نکلدی اے۔ اسیں !!{introduction} قاعدے !!{proof}s دے آخر وچ شامل کردے آں تے !!{elimination} قاعدے صرف !!{proof}s دے اُتّے۔ ایس لئی کدے وی !!{introduction} قاعدہ !!a{elimination} قاعدے توں اُتّے شامل نہیں کردے۔
\end{proof}

پر \Cut{} قاعدے دا جیہڑا ترجمہ \olref[nat][tgi]{prop:G2i-to-N2i} وچ دِتّا اے، اوہ \Log{N2i}-!!{proof}~$\delta$ وچ قطع پیدا کر سکدا اے۔ جے \Log{G2i} دا !!{proof}~\Cut{} اُتے مُکدا اے، تے اوہدے سدھے ذیلی !!{proof}s $\pi_1$ تے $\pi_2$ بالترتیب $\Gamma_1
\Sequent !A$ تے $!A, \Gamma_2 \Sequent \Delta_2$ دے نیں، تاں اسیں ترجمہ $\delta_1$، جیہڑا~$\pi_1$ دا اے، اوہناں مسلّمیاں $x:!A
\Sequent !A$ نال جوڑ دیندے آں جیہڑے ترجمے~$\delta_2$، یعنی~$\pi_2$ دے ترجمے وچ آؤندے نیں۔ جے $\delta_1$، !!a{introduction} قاعدے اُتے مُکدا اے جیہدا مُکھ !!{formula}~$!A$ اے، تے مسلّمہ $x:!A \Sequent !A$ جیہڑا~$\delta_2$ وچ اے، !!a{elimination} قاعدے دا وڈّا مقدمہ اے، تاں ایہدے نال~$\delta$ وچ اک قطع بنے گا۔

ایہ وی ممکن اے کہ !!{proof}s جیہڑے \Log{N2i} وچ نیں، اوہناں دا ترجمہ !!{proof}s دے طور اُتے~\Log{G2i} وچ ہووے۔ \olref[nat][tni]{prop:N2-to-G2} دے ثبوت وچ \Elim{\land}، \Elim{\lif}، \Elim{\lnot} تے \Elim{\lforall} دے ترجمے نال~\Cut{} استنباط، حاصل ہوئے \Log{G2i} !!{proof} وچ آ جاندے نیں۔ پر جے اسیں جیہڑے \Log{N2i}-!!{proof} توں شروع کردے آں … تاں ایہ ٹالیا جا سکدا اے۔

سیکوئنٹاں دیاں موجودگیاں دی لڑی $S_1$، \dots،~$S_n$ نوں \Log{N2i}-!!{proof}~$\delta$ دی \emph{شاخ} آکھدے آں جے $S_1$ اک مسلّمہ اے تے جدوں وی $S_i$ کسے استنباط دا وڈّا مقدمہ اے، $S_{i+1}$ اوہدا نتیجہ اے۔ دوجے لفظاں وچ، شاخ~$\delta$ وچ مسلّمیاں توں تھلّے ول سیکوئنٹاں دے نال نال جاندی اے، پر !!{elimination} قاعدیاں دے چھوٹے مقدمیّاں اُتے رُک جاندی اے۔~$\delta$ دی \emph{مُکھ شاخ} اوہ شاخ اے جیہدے وچ $S_n$ آخری سیکوئنٹ اے۔ ہر !!{proof}~$\delta$ دی گھٹّو گھٹ اک مُکھ شاخ ضرور ہوندی اے، کیونکہ ہر قاعدے دا گھٹّو گھٹ اک وڈّا مقدمہ اے (صرف \Intro{\land} دے دو نیں)۔

\begin{prop}\ollabel{prop:normal-N2i-to-G2i}
جے $\delta$، نارمل $\Log{N2c}$ یا $\Log{N2i}$-!!{proof} اے~$\Gamma \Sequent !A$ دا، تاں اک (قطع توں بناں) $\Log{G2c}$-!!{proof} ($\Log{G2i}$-!!{proof})~$\pi$ موجود اے جیہڑا $\Gamma' \Sequent \Delta$ دا ثبوت اے۔ ایتھے $\Gamma'$ اوہ کثیر سیٹ اے جیہدے !!{formula}s، $ \Gamma$ دے لیبل ہٹا کے ملدے نیں، تے $\Delta = \{!A\}$ جے $!A$،~$\lfalse$ نہیں، تے $\Delta = \emptyset$ جے اے۔
\end{prop}

\begin{proof}
ایس واری استقرا، !!{proof}~$\delta$ جیہڑا $ \Gamma \Sequent !A$ دا اے، اوہدے استنباطاں دی گنتی اُتے اے۔ جے~$\delta$ وچ کوئی استنباط نہیں، تاں اوہ مسلّمہ $x:!A \Sequent !A$ اے، تے $\pi$ اوہدے نال میل کھاندا مسلّمہ $!A \Sequent !A$ اے یا $ \lfalse \Sequent \ $ اے جے $!A
\ident \lfalse$۔

جے $\delta$ کسے استنباطی قاعدے اُتے مُکدا اے، تاں وکھ وکھ صورتاں ویکھدے آں:
\begin{enumerate}
  \item جے $R$، !!a{introduction} قاعدہ یا~\FalseInt اے، تاں ترجمہ عین اوہو جیہا چلدا اے جیہا \olref[nat][tni]{prop:N2-to-G2} دے ثبوت وچ، تے~\Cut دی لوڑ نہیں پیندی۔

  جے آخری قاعدہ~$R$،~$\delta$ دا، !!a{elimination} قاعدہ اے، تاں ہِٹھ لکھیاں گلاں ویکھو:
  \begin{enumerate}
    \item~$\delta$ دی مُکھ شاخ کسے !!{introduction} قاعدے تھانیں نہیں لنگھدی، کیونکہ نئیں تاں مُکھ شاخ اُتے !!a{introduction} قاعدے یا~\FalseInt{} توں پچھوں !!a{elimination} قاعدہ آؤندا تے $\delta$ نارمل نہ ہوندا۔
    \item چونکہ کوئی !!{introduction} قاعدہ~$\delta$ دیاں مُکھ شاخاں اُتے نہیں آؤندا، ایس لئی مُکھ شاخ بس اکّو اے۔ (صرف \Intro{\land} دے دو وڈّے مقدمے نیں، ایس لئی اک توں ودھ مُکھ شاخاں نوں کسے \Intro{\land} قاعدے دے مقدمیّاں تھانیں لنگھنا پیندا۔)
    \item جے مُکھ شاخ وچ \Elim{\lor} یا \Elim{\lexists} دا وڈّا مقدمہ اے، تاں ایہو جیہا قاعدہ اکّو اے تے اوہ آخری اے، یعنی~$R$۔ (نئیں تاں مُکھ شاخ وچ \Elim{\lor} یا \Elim{\lexists} قاعدہ ہوندا جیہدا نتیجہ !!a{elimination} قاعدے دا وڈّا مقدمہ ہوندا، تے ثبوت نارمل نہ ہوندا، کیونکہ اسیں تھاں بدلّن والی تبدیلی لا سکدے ہوندے۔)
    \item \Elim{\lor} تے \Elim{\lexists} توں چھُٹ کوئی !!{elimination} قاعدہ مفروضے خارج نہیں کردا، یعنی سیکوئنٹاں دا کھبّا سیاق نہیں بدلدا۔ \Elim{\lor} تے \Elim{\lexists} صرف چھوٹے مقدمے توں اُتّے مفروضے خارج کردے نیں، یعنی صرف اپنے چھوٹے مقدمیّاں دا کھبّا سیاق بدل دے نیں۔ دوجے لفظاں وچ، جے سیکوئنٹ $S_i$ جیہڑا~$\delta$ دی مُکھ شاخ اُتے اے، اوہدا سیاق~$\Gamma_1$ اے، تاں ~$S_{i+1}$، جیہڑا اوس استنباط دا نتیجہ اے جیہدا $S_i$ وڈّا مقدمہ اے، اوہدا سیاق $\Gamma_1' \supseteq \Gamma_1$ اے۔
    \item ایس لئی جے $x: !C \Sequent !C$ مُکھ شاخ دا سبھ توں اُتّلا سیکوئنٹ~$S_1$ اے، تاں $x:!C \in \Gamma$۔
  \end{enumerate}
  ہُن اسیں~$!C$ دی شکل دے حساب نال وکھ وکھ صورتاں ویکھدے آں۔ چونکہ ہر $S_i$ مُکھ شاخ اُتے کسے !!a{elimination} قاعدے دا وڈّا مقدمہ اے، ایس لئی ایہ گل $x:!C \Sequent !C$ اُتے وی لگدی اے۔

  \item $!C \ident !D \land !E$۔ تاں $\delta$ دی شکل ایہ اے:
  \[
  \Axiom$x:!D \land !E \fCenter !D \land !E$
  \RightLabel{\Elim{\land}[1]}
  \UnaryInf$x:!D \land !E \fCenter !D$
  \Deduce$x:!D \land !E, \Gamma_1 \fCenter !A$
  \DisplayProof
  \]
  مُکھ شاخ وچ $x:!D \land !E \Sequent !D$ آؤندا اے تے اوہدے وچ \Intro{\lif} یا \Intro{\lnot} دا کوئی وڈّا مقدمہ، یا \Elim{\lor} یا \Elim{\lexists} دا کوئی چھوٹا مقدمہ نہیں آؤندا۔ ایس لئی سیاق دے !!{formula}s، جیہڑے $x:!D
  \land !E$ وچ نیں، مُکھ شاخ اُتے نہیں بدل دے۔ ایس لئی آخری سیکوئنٹ دے سیاق وچ~$x:!D \land !E$ ضرور اے۔
  نالے اسیں !!{proof} $ \delta$ نوں !!{proof}~$\delta_1$ وچ بدل سکدے آں۔ ایس لئی جیہڑا ذیلی !!{proof}، \Elim{\land} اُتے مُکدا اے، اوہدی تھاں $x:!D \Sequent !D$ رکھو، تے $x:!D \land !E$ دی تھاں مُکھ شاخ دے ہر سیکوئنٹ دے سیاق وچ $x:!D$ رکھو، یعنی ایویں بدلو:
  \[
  \bottomAlignProof
  \Axiom$x:!D \land !E \fCenter !D \land !E$
  \RightLabel{\Elim{\land}[1]}
  \UnaryInf$x:!D \land !E \fCenter !D$
  \Deduce$x:!D \land !E, \Gamma_1 \fCenter !A$
  \DisplayProof
  \quad\text{ into }\quad
  \bottomAlignProof
  \Axiom$x:!D \fCenter !D$
  \Deduce$x:!D, \Gamma_1 \fCenter !A$
  \DisplayProof
  \]
  استقرا دا مفروضہ $\delta_1$ اُتے لگدا اے، کیونکہ $ \delta$ وچ استنباطاں دی گنتی $ \delta'$ وچ استنباطاں دی گنتی توں~$1$ ودھ اے۔

  استقرا دے مفروضے نال، \Log{G2i}-!!{proof} $ \pi_1$ موجود اے $x:!D, \Gamma_1 \Sequent !A$ دا۔
  !!{proof}~$\pi$، \LeftR{\land} لا کے ملدا اے:
  \[
  \AxiomC{}
  \RightLabel{$\pi_1$}
  \Deduce$!D, \Gamma_1' \fCenter !A$
  \RightLabel{\LeftR{\land}}
  \UnaryInf$!D \land !E, \Gamma_1' \fCenter !A$
  \DisplayProof
  \]
  جے $!A \ident \lfalse$ تاں اسیں \RightR{\Weakening} شامل کردے آں، مثلاً:
  \[
  \AxiomC{}
  \RightLabel{$\pi_1$}
  \Deduce$!D, \Gamma_1' \fCenter $
  \RightLabel{\LeftR{\land}}
  \UnaryInf$!D \land !E, \Gamma_1' \fCenter $
  \RightLabel{\RightR{\Weakening}}
  \UnaryInf$!D \land !E, \Gamma_1' \fCenter !A$
  \DisplayProof
  \]
  \item جے $!C \ident !D \lor !E$، تاں ایہ \Elim{\lor} دا وڈّا مقدمہ ضرور اے، تے ایہ استنباط مُکھ شاخ دا آخری استنباط~$R$ ضرور اے۔ دوجے لفظاں وچ، $\delta$ دی شکل ایہ اے:
  \[
  \Axiom$x:!D \lor !E \fCenter !D \lor !E$
  \AxiomC{}
  \RightLabel{$\delta_1$}
  \Deduce$y:!D, \Gamma_1 \fCenter!A$
  \AxiomC{}
  \RightLabel{$\delta_2$}
  \Deduce$z:!E, \Gamma_2 \fCenter !A$
  \RightLabel{\Elim{\lor}}
  \TrinaryInf$x:!D \lor !E, \Gamma_1, \Gamma_2 \fCenter !A$
  \DisplayProof
  \]
  استقرا دا مفروضہ !!{proof}s $ \delta_1$ تے $ \delta_2$ اُتے لگدا اے۔ سانوں \Log{G2i}-!!{proof}s $ \pi_1$ تے $ \pi_2$، بالترتیب $!D, \Gamma_1' \Sequent !A$ تے $!E, \Gamma_2' \Sequent !A$ دے، ملدے نیں۔$\LeftR{\lor}$ لا کے~$\pi$ ملدا اے:
  \[
  \AxiomC{}
  \RightLabel{$\pi_1$}
  \Deduce$!D, \Gamma_1' \fCenter !A$
  \AxiomC{}
  \RightLabel{$\pi_2$}
  \Deduce$!E, \Gamma_2' \fCenter !A$
  \RightLabel{\LeftR{\lor}}
  \BinaryInf$!D \lor !E, \Gamma_1', \Gamma_2' \fCenter !A$
  \DisplayProof
  \]
  جے $ \Elim{\lor}$، $y: !D$ یا $z: !E$ نوں خارج نہ کرے (یعنی ایہ چھوٹے مقدمیّاں دے سیاق وچ نہ ہون)، یا جے~$!A$، $ \lfalse$ ہووے، تاں سانوں کجھ \LeftR{\Weakening} تے \RightR{\Weakening} شامل کرنے پین گے۔ مثال لئی:
  \[
  \AxiomC{}
  \RightLabel{$\pi_1$}
  \Deduce$\Gamma_1' \fCenter $
  \RightLabel{\LeftR{\Weakening}}
  \UnaryInf$!D, \Gamma_1' \fCenter $
  \AxiomC{}
  \RightLabel{$\pi_2$}
  \Deduce$\Gamma_2' \fCenter $
  \RightLabel{\LeftR{\Weakening}}
  \UnaryInf$!E, \Gamma_2' \fCenter $
  \RightLabel{\LeftR{\lor}}
  \BinaryInf$!D \lor !E, \Gamma_1', \Gamma_2' \fCenter $
  \RightLabel{\RightR{\Weakening}}
  \UnaryInf$!D \lor !E, \Gamma_1', \Gamma_2' \fCenter !A$
  \DisplayProof
  \]

  \item $!C \ident !D \lif !E$۔ تاں $\delta$ دی شکل ایہ اے:
  \[
  \Axiom$x:!D \lif !E \fCenter !D \lif !E$
  \AxiomC{}
  \RightLabel{$\delta_1$}
  \Deduce$\Gamma_1 \fCenter!D$
  \RightLabel{\Elim{\lif}}
  \BinaryInf$x:!D \lif !E, \Gamma_1 \fCenter !E$
  \RightLabel{$\delta_2$}
  \Deduce$x:!D \lif !E, \Gamma_1, \Gamma_2 \fCenter !A$
  \DisplayProof
  \]
  ایتھے ویکھو کہ مُکھ شاخ، جیہدے وچ $x:!D \lif !E,
  \Gamma_1 \Sequent !E$ آؤندا اے، اوہدے وچ \Intro{\lif} یا \Intro{\lnot} دا کوئی وڈّا مقدمہ، یا \Elim{\lor} یا \Elim{\lexists} دا کوئی چھوٹا مقدمہ نہیں آؤندا۔ ایس لئی سیاق دے !!{formula}s، جیہڑے $x:!D \lif !E, \Gamma_1$ وچ نیں، مُکھ شاخ اُتے نہیں بدل دے۔ ایس لئی آخری سیکوئنٹ دے سیاق وچ~$x:!D \lif !E, \Gamma_1$ ضرور اے۔ نالے اسیں !!{proof} دے ٹوٹے $ \delta_2$ نوں درست !!{proof}~$\delta_2'$ وچ بدل سکدے آں۔ ایس لئی ذیلی !!{proof} جیہڑا \Elim{\lif} اُتے مُکدا اے، اوہدی تھاں $x:!E \Sequent !E$ رکھو، تے $x:!D \lif !E,
  \Gamma_1$ دی تھاں مُکھ شاخ دے ہر سیکوئنٹ دے سیاق وچ $x:!E$ رکھو، یعنی ایویں بدلو:
  \[
  \bottomAlignProof
  \Axiom$x:!D \lif !E, \Gamma_1 \fCenter !E$
  \RightLabel{$\delta_2$}
  \Deduce$x:!D \lif !E, \Gamma_1, \Gamma_2 \fCenter !A$
  \DisplayProof
  \quad\text{ into }\quad
  \bottomAlignProof
  \Axiom$x:!E \fCenter !E$
  \RightLabel{$\delta_2'$}
  \Deduce$x:!E, \Gamma_2 \fCenter !A$
  \DisplayProof
  \]
  استقرا دا مفروضہ $\delta_1$ تے $\delta_2'$ اُتے لگدا اے، کیونکہ $ \delta$ وچ استنباطاں دی گنتی، $ \delta_1$ تے~$\delta_2$ وچ استنباطاں دیاں گنتیاں دا جوڑ اے، نال~$1$ مُکھ شاخ دے شروع والے \Elim{\lif} استنباط لئی۔ (جے اوہو استنباط آخری استنباط~$R$ وی اے، $ \delta$ وچ، تاں $ \delta_1$ وچ $0$ استنباط نیں۔)

  استقرا دے مفروضے نال \Log{G2i}-!!{proof}s $ \pi_1$، $ \Gamma_1 \Sequent !D$ دا، تے $ \pi_2$، $!E, \Gamma_1' \Sequent !A$ دا، موجود نیں۔
  !!{proof}~$\pi$، \LeftR{\lif} لا کے ملدا اے:
  \[
  \AxiomC{}
  \RightLabel{$\pi_1$}
  \Deduce$\Gamma_1' \fCenter !D$
  \AxiomC{}
  \RightLabel{$\pi_2$}
  \Deduce$!E, \Gamma_2' \fCenter !A$
  \RightLabel{\LeftR{\lif}}
  \BinaryInf$!D \lif !E, \Gamma_1', \Gamma_2' \fCenter !A$
  \DisplayProof
  \]
  جے $!D$ یا $!A \ident \lfalse$ ہووے، یا دونے، تاں اسیں اک یا دو \RightR{\Weakening} شامل کردے آں، مثلاً:
  \[
  \AxiomC{}
  \RightLabel{$\pi_1$}
  \Deduce$\Gamma_1' \fCenter $
  \RightLabel{\RightR{\Weakening}}
  \UnaryInf$\Gamma_1' \fCenter !D$
  \AxiomC{}
  \RightLabel{$\pi_2$}
  \Deduce$!E, \Gamma_2' \fCenter $
  \RightLabel{\RightR{\Weakening}}
  \UnaryInf$!E, \Gamma_2' \fCenter !A$
  \RightLabel{\LeftR{\lif}}
  \BinaryInf$!D \lif !E, \Gamma_1', \Gamma_2' \fCenter !A$
  \DisplayProof
  \]

\end{enumerate}
باقی صورتاں دا حل وی ایویں ای اے تے اوہ مشقاں لئی چھڈّیاں گئیاں نیں۔
\end{proof}

\begin{cor}
ہر !!{formula} جیہڑا نارمل \Log{N1i} یا \Log{N2i}-!!{proof} وچ آؤندا اے، اوہ ذیلی !!{formula} اے، آخری !!{formula} یا آخری سیکوئنٹ دا۔
\end{cor}

\begin{proof}
\olref{prop:normal-N2i-to-G2i} دے حساب نال ہر نارمل \Log{N2i}-!!{proof} دا ترجمہ \Log{G2i}-!!{proof} وچ ہو سکدا اے۔ ثبوت ویکھ کے پتا لگدا اے کہ ہر !!{formula} جیہڑا \Log{N2i}-!!{proof} وچ آؤندا اے، \Log{G2i}-!!{proof} وچ وی آؤندا اے۔ \Log{G2i} وچ \Cut{} دا قاعدہ نہ ہون کر کے ذیلی !!{formula} والی خاصیت اے۔
\end{proof}
\endgroup
% /OLP-0678

\OLEndChapterHook
\endgroup
% /OLP-0674
% OLP-0690 content/proof-theory/propositions-as-types/propositions-as-types.tex
\begingroup
\def\olfilename{propositions-as-types}
\expandafter\def\csname ol@sectioncs\endcsname{section}

\olchapter{int}{pty}{قضیے بطور نوعاں}

\begin{editorial}
  ایہ کری–ہاورڈ میل اُتے اک باب دا \emph{بہت تجرباتی} مسودہ اے۔ ایہنوں ہور سمجھاؤنی تے وجہ دسّن دی لوڑ اے، تے شاید ایہدے وچ غلطیاں تے رہندیاں گلاں ہون۔ نارمل روپ وچ لیاؤن دے ثبوت نوں مُڑ ویکھنا تے ودھاؤنا چاہیدا اے۔ حاصلِ ضرب دی نوع لئی کوئی مثال نہیں۔ تھاں بدلّن والیاں تے سادہ کرن والیاں تبدیلیاں شامل نہیں۔ جدوں (نوع بندی والے) لامبڈا کیلکولس بارے مواد وی ہووے گا، جیہڑا ایتھے بنیادی طور اُتے پہلاں توں منّیا گیا اے، تاں ایہ ہور سوکھا سمجھ آوے گا۔ بہت احتیاط نال ورتو۔
\end{editorial}

% OLP-0686 content/proof-theory/propositions-as-types/introduction.tex
\begingroup
\def\olfilename{introduction}
\expandafter\def\csname ol@sectioncs\endcsname{section}

\olfileid{int}{pty}{int}

\olsection{جان پچھان}

لامبڈا کیلکولس حسابی عمل دا اک سادہ نمونہ اے۔ ایہ متغیراں توں بنے ٹرماں اُتے کم کردا اے۔ جے $N$ تے $M$ ٹرم نیں، تاں $(NM)$ وی ٹرم اے («$N$ نوں~$M$ اُتے لایا»)۔ جے $N$ ٹرم اے، تاں $\lambd[x][N]$ ٹرم اے۔ جے $N$ وچ متغیر~$x$ اے، تاں~$x$ دیاں میل کھاندیاں موجودگیاں $\lambd[x][N]$ وچ $\lambd[x]$ عمل کار نال بندھ جاندیاں نیں تے ہُن آزاد نہیں رہندیاں۔ $(\lambd[x][N]M)$ دی شکل والے ٹرم بارے آکھدے آں کہ اوہ $\beta$-سکڑاؤ نال $\Subst{N}{M}{x}$ بن جاندا اے (یعنی $x$ دی ہر آزاد موجودگی~$N$ وچ نوں $M$ نال بدلّن دا حاصل)۔
اک ٹرم $N$، $\beta$-تبدیلی نال~$N'$ بن جاندا اے جے $N'$، $N$ وچ کسے ذیلی ٹرم دا $\beta$-سکڑاؤ کر کے ملدا اے، تے اسیں $N \redone N'$ لکھدے آں۔ $\lambd$-کیلکولس وچ اک حسابی عمل لڑی $N \redone N_1
\redone N_2 \redone \dots$ اے۔ جے ایہ لڑی کسے ٹرم $N_k$ اُتے مُکدی اے جیہدے وچ کسے ذیلی ٹرم دا $\beta$-سکڑاؤ نہیں ہو سکدا، تاں ایہنوں \emph{نارمل} یا~$N$ دا نارمل روپ آکھدے آں۔ $\lambd$-کیلکولس وچ حسابی عمل نوں ایہو جیہیاں لڑیاں وانگ سمجھو۔ حسابی عمل اُدوں رُک جاندا اے جدوں اوہدا حاصل نارمل روپ ہووے، تے نارمل روپ ای حسابی عمل دا نتیجہ اے۔ حسابی عملاں دا ایہ سادہ نمونہ ٹیورنگ-کامل نکلدا اے۔ ہر جزوی بازگشتی فنکشن $\lambd$-کیلکولس دے کسے ٹرم نال حسابیا جا سکدا اے۔

ہیسکل کری تے ولیم ہاورڈ نے قدرتی اخذ تے $\lambd$-کیلکولس وچ نیڑی سانجھ آپو اپنے طور اُتے لبھی۔
سہج سمجھ نال، ٹرم $\lambd[x][N]$ اک فنکشن اے۔ $x$ اوہدا آرگومنٹ اے، تے $N$ دَسّدا اے کہ~$x$ دِتّا ہووے تاں قدر کیویں حسابنی اے۔ ٹرم $(\lambd[x][N]M)$ دا مطلب فیر فنکشن $\lambd[x][N]$ نوں آرگومنٹ~$M$ اُتے لاؤنا اے، تے $\beta$-سکڑاؤ دَسّدا اے کہ اوہدی قدر کیویں کڈھنی اے۔ $x$ نوں $N$ وچ~$M$ نال بدلو (تے جیہڑا ٹرم ملے، اوہدی قدر کڈھدے رہو)۔
ہُن ایہناں نوں مقرر دائرے تے حاصل دی حد والے فنکشن سمجھو۔ جے $\lambd[x][N]$ دا دائرہ $A$ اے تے اوہدی حاصل دی حد~$B$ اے، تاں متغیر~$x$ نوں ایویں سمجھو کہ اوہ صرف~$A$ توں آرگومنٹ لیندا اے تے $N$،~$B$ وچوں قدراں دیندا اے۔ سو $(\lambd[x][N]M)$ دا مطلب صرف اوہدوں ہونا چاہیدا اے جے $\lambd[x][N]$، $A \to B$ والا فنکشن اے تے $M \in A$۔
جے اسیں ایہ جانکاری $\lambd$-کیلکولس وچ شامل کرنی چاہیئے، تاں سانوں \emph{نوع بندی والا} $\lambd$-کیلکولس ملدا اے۔ نوع بندی والے $\lambd$-کیلکولس وچ ہر عبارت $(\lambd[x][N]M)$ جائز نہیں، صرف اوہ جائز نیں جیہناں وچ $\lambd[x][N]$ تے $M$ دیاں «نوعاں» آپس وچ میل کھاندیاں نیں، یعنی $\lambd[x][N]$ دی نوع $A \lif B$ اے تے $M$ دی نوع~$A$ اے۔ $\lambd[x][N]$ دی نوع $A \to B$ صرف اوہدوں اے جے $x$ دی نوع $A$ تے $N$ دی نوع~$B$ اے۔ عام طور اُتے ہر عبارت $(M_1M_2)$ جائز نہیں۔ صرف اوہ ٹرم $(M_1M_2)$ جائز نیں جیہناں وچ $M_1$ دی نوع $A \to B$ تے $M_2$ دی نوع~$A$ اے۔ جے ایویں اے، تاں $(M_1M_2)$ دی نوع~$B$ اے۔

ہُن صاف اے کہ ایہ نوع بندی دے قاعدے قدرتی اخذ دے !!{derivation}s نال میل کھاندے نیں۔ اُتھے نوعاں نوں !!{formula}s ظاہر کردے نیں تے !!{derivation}s آپ $\lambd$-ٹرماں نال میل کھاندے نیں۔ مثال لئی !!a{derivation}، $!A \lif !B$ دا، ٹرم $\lambd[\typeof{x}{!A}][N]$ نال میل کھاندا اے جیہدی فنکشن نوع $!A \lif !B$ اے۔ قدرتی اخذ دے استنباطی قاعدے نوع بندی والے لامبڈا کیلکولس دے نوع بندی قاعدیاں نال میل کھاندے نیں۔ مثال لئی:
\begin{prooftree}
  \AxiomC{$\Discharge{!A}{x}$}
  \DeduceC{$!B$}
  \DischargeRule{\Intro\lif}{x}
  \UnaryInfC{$!A \lif !B$}
  \DisplayProof\bottomAlignProof
  \quad\text{نال میل کھاندا اے}\quad
  \Axiom$x: !A \fCenter N: !B$
  \UnaryInf$ \fCenter \lambd[\typeof{x}{!A}][N] : !A \lif !B$
\end{prooftree}
سجّے پاسے دے قاعدے دا مطلب اے کہ جے $x$ دی نوع~$!A$ اے تے $N$ دی نوع~$!B$ اے، تاں $\lambd[\typeof{x}{!A}][N]$ دی نوع~$!A \lif
!B$ اے۔

$\Elim\lif$ قاعدہ جوڑ والے ٹرماں دے نوع بندی قاعدے نال میل کھاندا اے، یعنی:
\[
  \AxiomC{$!A \lif !B$}
  \AxiomC{$!A$}
  \RightLabel{\Elim\lif}
  \BinaryInfC{$!B$}
  \DisplayProof
  \quad\text{نال میل کھاندا اے}\quad
  \Axiom$\fCenter M_1: !A \lif !B$
  \Axiom$\fCenter M_2: !A$
  \BinaryInf$ \fCenter (M_1M_2) : !B$
  \DisplayProof
\]
جے \Intro\lif{} قاعدے دے جھٹ پچھوں \Elim\lif{} قاعدہ آوے، تاں !!{derivation} نوں سوکھا کیتا جا سکدا اے:
\begin{gather*}
  \AxiomC{$\Discharge{!A}{x}$}
  \DeduceC{$!B$}
  \DischargeRule{\Intro\lif}{x}
  \UnaryInfC{$!A \lif !B$}
  \AxiomC{}
  \DeduceC{$!A$}
  \RightLabel{\Elim\lif}
  \BinaryInfC{$!B$}
  \DisplayProof
  \quad\redone\quad
  \AxiomC{}
  \DeduceC{$!A$}
  \DeduceC{$!B$}
  \DisplayProof
\end{gather*}
ایہ لامبڈا ٹرماں دے $\beta$-سکڑاؤ نال میل کھاندا اے:
\[
(\lambd[\typeof{x}{!A}][N]M) \quad\redone\quad
\Subst{N}{M}{x}.
\]

قاعدیاں تے نوعاں دا ایہو جیہا میل $ \land$ دے قاعدیاں تے «حاصلِ ضرب» نوعاں وچ، تے $ \lor$ دے قاعدیاں تے «جوڑ» نوعاں وچ وی اے۔

سادہ نوع بندی والے لامبڈا کیلکولس دے ٹرماں تے قدرتی اخذ دے !!{derivation}s دے ایس میل نوں «کری–ہاورڈ»، «ثبوت بطور پروگرام»، یا «قضیے بطور نوعاں» دا میل آکھدے آں۔ !!{formula}s (قضیاں) دے نوعاں نال تے ثبوتاں دے ٹرماں نال میل توں علاوہ، اسیں ایہ میل ایویں اکٹھے وکھا سکدے آں:
\begin{center}
  \begin{tabular}{c c c}
    منطق & پروگرام \\
    \hline
    قضیہ/!!{formula} & نوع \\
    ثبوت & ٹرم \\
    مفروضہ & متغیر \\
    خارج کیتا مفروضہ & بندھیا متغیر \\
    کھُلّا مفروضہ & آزاد متغیر \\
    $!A \lif !B$ & فنکشن نوع \\
    $!A \land !B$ & حاصلِ ضرب دی نوع \\
    $!A \lor !B$ & جوڑ دی نوع \\
    $\lfalse$ & خالی نوع \\
    \hline
  \end{tabular}
\end{center}

کری–ہاورڈ میل خودکار ثبوت معاوناں تے پروگراماں دیاں نوعاں جانچن والے نظاماں دی اک بنیادی گل اے۔ وجہ ایہ اے کہ کسے قضیے دے گواہ ثبوت دی جانچ (جیویں اسیں اُتّے کیتی)، اوسے گل دی جانچ اے کہ کسے پروگرام (ٹرم) دی نوع اوہی اے جیہڑی اوہدے لئی دَسّی گئی اے۔
\endgroup
% /OLP-0686
%\olimport{sequent-natural-deduction}
% OLP-0688 content/proof-theory/propositions-as-types/proof-terms.tex
\begingroup
\def\olfilename{proof-terms}
\expandafter\def\csname ol@sectioncs\endcsname{section}

\olfileid{int}{pty}{ter}

\olsection{ثبوت ٹرم}

\Log{N1i} وچ اسیں مفروضیاں نوں خارج کرن دے لیبل دے طور اُتے متغیر لاؤنے آں (جیویں $!A^x$)۔ \Log{N2i} وچ ہر سیکوئنٹ $ \Gamma \Sequent !A$ دے کھبّے پاسے سیاق~$\Gamma$ وچ متغیر–!!{formula} جوڑیاں دی لڑی اے۔ سو \Log{N2i} وچ !!a{proof} دا ہر سیکوئنٹ ایہ جانکاری رکھدا اے کہ !!{proof} نے ہُن تک کیہ !!{prove} کیتا اے (یعنی $!A$)، تے ہر تھاں کِہڑے مفروضے کھُلّے نیں (یعنی $ \Gamma$)۔ جے ہر سیکوئنٹ وچ ایہ جانکاری وی پا دیئیے کہ $!A$ \emph{کِنج} !!{prove} ہویا، تاں کیہ ہووے؟ ایس لئی اسیں اک \emph{ٹرماں دے لیبل والا} نظام \Log{tN2i} متعین کر سکدے آں جیہدے وچ سیکوئنٹ $ \Gamma \Sequent N : !A$ دی شکل دے ہون۔ $ \Gamma$ تے $!A$ پہلاں وانگ ای نیں۔ پہلا $x: !B$ دی شکل والے جوڑیاں دا سیٹ اے تے دوجا اوہ !!{formula} اے جیہڑا سیکوئنٹ اُتے مُکّن والے ذیلی !!{proof} نے~$\Gamma$ دے مفروضیاں توں نتیجے دے طور اُتے وکھایا اے۔ $N$ اک ایہو جیہا ٹرم ہووے گا جیہڑا سیکوئنٹ اُتے مُکّن والے !!{proof} نوں کوڈ کردا اے۔

ہر سیکوئنٹ نوں دِتّے ٹرم متغیراں $x$، $y$، \dots، توں فنکشن نشانیاں ورت کے بنائے جاندے نیں۔ ایہ نشانیاں ہُن تک ہوئے استنباطاں نوں کوڈ کردیاں نیں۔ سانوں ہر استنباطی قاعدے لئی وکھری فنکشن نشانی چاہیدی اے۔ !!{introduction} قاعدیاں نال میل کھاندیاں فنکشن نشانیاں نوں «بناؤن والیاں» تے !!{elimination} قاعدیاں نال میل کھاندیاں نوں «کھولھن والیاں» آکھدے آں۔ ہر فنکشن نشانی اُنے آرگومنٹ لیندی اے جِنّے قاعدے دے مقدمے نیں۔ مثال لئی \Intro{\land} تے \Elim{\land}[1] لئی فنکشن نشانیاں $c^\land$ (دو تھاواں والی) تے $d^\land_1$ (اک تھاں والی) ہو سکدیاں نیں۔ \Log{tN2i} وچ فیر ایہ قاعدے ایہ بن جاندے نیں:
\[
\Axiom$\Gamma_1 \fCenter N: !A$
\Axiom$\Gamma_2 \fCenter M: !B$
\RightLabel{\Intro{\land}}
\BinaryInf$\Gamma_1, \Gamma_2 \fCenter c^\land(N, m) : !A \land !B $
\DisplayProof
\quad
\Axiom$\Gamma \fCenter N: !A \land !B$
\RightLabel{\Elim{\land}[1]}
\UnaryInf$\Gamma \fCenter d^\land_1(N) : !A$
\DisplayProof
\]
جے استنباطی قاعدہ کوئی مفروضہ خارج کردا اے، یعنی اوہدا مقدمہ لیبل والا !!{formula}~$x:!A$ دَسّدا اے جیہڑا نتیجے دے سیاق وچ نہیں، تاں ایہ گل ثبوت ٹرم وچ وی وکھاؤنی پینی اے۔ ایہو جیہے قاعدیاں نال جُڑیاں فنکشن نشانیاں میل کھاندے متغیراں نوں \emph{بندھ دیندیاں} نیں۔ مثال لئی \Intro{\lif} قاعدہ مقدمے $x:!A, \Gamma \Sequent !B$ توں $ \Gamma \Sequent !A \lif !B$ تک جاندا اے۔ بناؤن والی نشانی $c^{\lif}$ اک آرگومنٹ لیندی اے، پر ایہ وی وکھانا پیندا اے کہ متغیر $x$ بندھ جاندا اے۔ ایہ ثبوت ٹرم دا ایہ دَسّن دا طریقہ اے کہ~$x$ لیبل والا مفروضہ خارج ہو گیا اے۔ مثال لئی جے مقدمے دا ثبوت ٹرم~$N$ اے، تاں نتیجے دا ثبوت ٹرم $c^{\lif}([x]N)$ تے قاعدہ ایویں لکھ سکدے آں:
\[
\Axiom$x:!A, \Gamma \fCenter N: !B$
\RightLabel{\Intro{\lif}}
\UnaryInf$\Gamma \fCenter c^{\lif}([x]N) : !A \lif !B$
\DisplayProof
\]

!!a{proof} نوں پورا مُڑ بناؤن لئی سانوں کجھ ہور جانکاری چاہیدی اے۔ جے قاعدے دے مقدمیّاں دے !!{formula}s پوری طرح طے نہیں کردے کہ نتیجے دا !!{formula} کیہڑا اے، تاں ایہ جانکاری ٹرم وچ شامل کرنی پینی اے۔ \Intro{\lif} وچ ایویں ای اے، ایس لئی اسیں $c^{\lif}([x:!A]N)$ لکھاں گے۔ \Intro{\lor} تے~\FalseInt وچ وی ایویں ای اے، ایس لئی اسیں ایہو جیہیاں فنکشن نشانیاں ورتدے آں جیہناں وچ ایہ جانکاری اے، جیویں:
\[
\Axiom$\Gamma \fCenter N:!A$
\RightLabel{\Intro{\lor}[1]}
\UnaryInf$\Gamma \fCenter c^{\lor{!B}}_1:!A \lor !B$
\DisplayProof
\quad
\Axiom$\Gamma \fCenter N:\lfalse$
\RightLabel{\FalseInt}
\UnaryInf$\Gamma \fCenter d^\lfalse_{!A}(N):!A$
\DisplayProof
\]

ایہناں خیالاں نال سیکوئنٹ طرز دا اک قدرتی اخذ نظام بنایا جا سکدا اے، جیہدے وچ ہر سیکوئنٹ دی شکل $ \Gamma
\Sequent N : !A$ اے، جتھے $N$ اک \emph{ثبوت ٹرم} اے۔

پتا لگدا اے کہ ایہو جیہے ثبوت ٹرماں تے \emph{نوع بندی والے لامبڈا کیلکولس} دے ٹرماں وچ نیڑا جوڑ اے۔ ایس جوڑ نوں قائم رکھن لئی اُتلیاں عام بناؤن تے کھولھن والیاں نشانیاں دی تھاں لامبڈا کیلکولس دیاں لکھتاں وچ ورتیاں جاندیاں نشانیاں ورتاں گے۔ مثال لئی $c^{\lif}([x:!A]N)$ دی تھاں $\lambd[\typeof{x}{!A}][N]$ لکھدے آں، $d^{\lif}(N,M)$ دی تھاں $(NM)$، $\pair{N}{M}$ لکھدے آں $c^\land(N,M)$ دی تھاں، تے $\proj{i}{N}$ لکھدے آں~$d^\land_i(N)$ دی تھاں۔

\begin{defn}
  ثبوت ٹرم ہِٹھلے قاعدیاں نال استقرائی طور اُتے بنائے جاندے نیں:
  \begin{enumerate}
  \item ہر متغیر $x$ اک ثبوت ٹرم اے (مسلّمے)۔
  \item جے $x$ متغیر اے، $!A$ اک !!a{formula} اے، تے $N$ ثبوت ٹرم اے، تاں $\lambd[\typeof{x}{!A}][N]$ وی ثبوت ٹرم اے~(\Intro\lif)۔
  \item جے $N$ تے $M$ ثبوت ٹرم نیں، تاں $(NM)$ وی ثبوت ٹرم اے~(\Elim\lif)۔
  \item جے $N$ تے $M$ ثبوت ٹرم نیں، تاں $\pair{N}{M}$ ثبوت ٹرم اے~(\Intro\land)۔
  \item جے $N$ ثبوت ٹرم اے، تاں $\proj{i}{N}$ وی ثبوت ٹرم اے، جتھے $i$، $1$ یا~$2$ اے~(\Elim\land[i])۔
  \item جے $N$ ثبوت ٹرم اے تے $!A$ فارمولا اے، تاں $\inj[!A]{i}{!M}$ ثبوت ٹرم اے، جتھے $i$، $1$ یا~$2$ اے~(\Intro\lor[i])۔
  \item جے $M$، $N_1$ تے $N_2$ ثبوت ٹرم نیں تے $x_1$، $x_2$ متغیر نیں، تاں $\dcase{M}{x_1}{N_1}{x_2}{N_2}$ ثبوت ٹرم اے~(\Elim\lor)۔
  \item جے $N$ ثبوت ٹرم اے تے $!A$ اک !!a{formula} اے، تاں $\abort{!A}{N}$ ثبوت ٹرم اے~(\FalseInt)۔
  \end{enumerate}
\end{defn}

ہر ثبوت ٹرم !!a{proof} دا ثبوت ٹرم نہیں ہوندا۔ مثال لئی ٹرم $\pair{N}{M}$ ایہو جیہا ثبوت کوڈ کردا اے جیہڑا~\Intro{\land} استنباط اُتے مُکدا اے۔ سو اوہدا آخری !!{formula} عطف~$!A \land !B$ اے۔ ایہ~\Elim{\lif} قاعدے دا وڈّا مقدمہ نہیں ہو سکدا، ایس لئی $(\pair{N}{M}P)$ کسے درست ثبوت دا ثبوت ٹرم نہیں ہو سکدا۔ صرف اوہ ثبوت ٹرم \emph{درست} ثبوت ٹرم نیں جیہڑے \Log{tN2ip} دے قاعدیاں مطابق نیں۔ ایہ قاعدے \olref{tab:tN2ip} وچ دِتّے نیں۔

% OLP-0692 content/proof-theory/propositions-as-types/rules-tN2.tex
\begingroup

\begin{table}
\begin{tabular}{@{}c@{}c@{}}
\multicolumn{2}{l}{\textbf{مسلّمے}}\\
\multicolumn{2}{c}{$x:!A \Sequent x:!A$}\\[3ex]
\multicolumn{2}{l}{\textbf{استنباط دے قاعدے}}\\[3ex]
\multicolumn{2}{c}{\Axiom$x: !A, \Gamma \fCenter N:!B$
\RightLabel{\Intro{\lif}}
\UnaryInf$\Gamma \fCenter \lambd[\typeof{x}{!A}][N] : !A \lif !B $
\DisplayProof}
\\[4ex]
\multicolumn{2}{c}{\Axiom$\Gamma_1 \fCenter N:!A \lif !B$
\Axiom$\Gamma_2 \fCenter M:!A$
\RightLabel{\Elim{\lif}}
\BinaryInf$\Gamma_1, \Gamma_2 \fCenter (NM) : !B$
\DisplayProof}
\\[4ex]
\multicolumn{2}{c}{\Axiom$\Gamma_1 \fCenter N_1 : !A$
\Axiom$\Gamma_2 \fCenter N_2 : !B$
\RightLabel{\Intro{\land}}
\BinaryInf$\Gamma_1, \Gamma_2 \fCenter \pair{N_1}{N_2} : !A \land !B $
\DisplayProof}
\\[4ex]
\Axiom$\Gamma \fCenter N: !A \land !B$
\RightLabel{\Elim{\land}[1]}
\UnaryInf$\Gamma \fCenter \proj{1}{N} : !A$
\DisplayProof
&
\Axiom$\Gamma \fCenter N:!A \land !B$
\RightLabel{\Elim{\land}[2]}
\UnaryInf$\Gamma \fCenter \proj{2}{N} : !B$
\DisplayProof\\[3ex]
\\[4ex]
\Axiom$\Gamma \fCenter N: !A$
\RightLabel{\Intro{\lor}[1]}
\UnaryInf$\Gamma \fCenter \inj[!B]{1}{N} : !A \lor !B$
\DisplayProof
&
\Axiom$ \Gamma \fCenter N : !B$
\RightLabel{\Intro{\lor}[2]}
\UnaryInf$\Gamma \fCenter \inj[!A]{2}{N} : !A \lor !B$
\DisplayProof\\[3ex]
\multicolumn{2}{c}{
\begin{tabular}[b]{@{}l@{}}
\Axiom$\Gamma_1 \fCenter M : !A \lor !B$
\Axiom$x:!A, \Gamma_2 \fCenter N_1 : !C$
\Axiom$y:!B, \Gamma_3 \fCenter N_2: !C$
\RightLabel{\Elim{\lor}}
\TrinaryInf$\Gamma_1, \Gamma_2, \Gamma_3 \fCenter
\dcase{M}{x}{N_1}{y}{N_2} : !C$
\DisplayProof
\end{tabular}
}
\\[4ex]
\multicolumn{2}{c}{\Axiom$\Gamma \fCenter N: \lfalse$
\RightLabel{\FalseInt}
\UnaryInf$\Gamma \fCenter \abort{!A}{N} : !A$
\DisplayProof}
\end{tabular}
\caption{ٹرماں دے لیبل والے قضیاتی وجدانی قدرتی اخذ نظام \Log{tN2ip} دے قاعدے۔}
\ollabel{tab:tN2ip}
\end{table}
\endgroup
% /OLP-0692
\endgroup
% /OLP-0688
% OLP-0689 content/proof-theory/propositions-as-types/proofs-to-terms.tex
\begingroup
\def\olfilename{proofs-to-terms}
\expandafter\def\csname ol@sectioncs\endcsname{section}

\olfileid{int}{pty}{pt}

\olsection{\usetoken{P}{derivation} نوں ثبوت ٹرماں وچ بدلنا}

قدرتی اخذ دے !!{proof}s نوں ثبوت ٹرماں وچ بدلّن دا طریقہ، یعنی~\Log{N2i} دے !!{proof}s دا \Log{tN2i} دے !!{proof}s وچ ترجمہ، بہت سدھا اے۔ !!{formula} جیہڑا کسے سیکوئنٹ دے سجّے پاسے اے، اوہدے کھبّے پاسے اک ثبوت ٹرم لکھو، !!{derivation} وچ ہر ایہو جیہے فارمولے لئی۔ ایہدے نال $ \Gamma \Sequent M :
!A$ دی شکل دے سیکوئنٹ ملدے نیں۔ فیر اسیں آکھاں گے کہ $M$،~$!A$ دا \emph{گواہ اے}، سیاق~$\Gamma$ وچ۔ کِہڑا ثبوت ٹرم لکھنا اے، ایہ سارا \Log{tN2i} دے قاعدے طے کردے نیں۔

\begin{prop}
ہر !!{proof}~$\delta$، $ \Gamma \Sequent !A$ دا \Log{N2i} وچ، !!a{proof}~$\delta'$ نال میل کھاندا اے، $ \Gamma \Sequent N : !A$ دا \Log{tN2i} وچ۔ ثبوت ٹرم $N$،~$\delta$ نال اکّو طرح طے ہوندا اے۔
\end{prop}

\begin{proof}
اسیں~$\delta$ دی اُچائی اُتے استقرا کردے آں۔ جے $\delta$ مسلّمہ $x:!A \Sequent !A$ اے، تاں $\delta'$ مسلّمہ~$x:!A \Sequent
x:!A$ اے۔ جے $\delta$ استنباطی قاعدے اُتے مُکدا اے، تاں استقرا دا مفروضہ ہر مقدمے لئی ثبوت ٹرم دیندا اے، تے \Log{tN2i}-!!{proof}s وی، اوہناں سیکوئنٹاں دے جیہڑے ہر مقدمے نال میل کھاندے نیں تے ایہ ثبوت ٹرم رکھدے نیں۔ اسیں \Log{tN2i} دا میل کھاندا قاعدہ لاؤنے آں۔ میل کھاندا ثبوت ٹرم \Log{tN2i} دا استنباطی قاعدہ طے کردا اے۔
\end{proof}

\begin{ex}
  بہت سدھا !!{proof} ویکھو \Log{N1i} وچ، $(!A \land !B)
  \lif (!B \land !A)$ دا:
  \[
  \AxiomC{$\Discharge{!A\land !B}{x}$}
  \RightLabel{\Elim\land[2]}
  \UnaryInfC{$!B$}
  \AxiomC{$\Discharge{!A\land !B}{x}$}
  \RightLabel{\Elim\land[1]}
  \UnaryInfC{$!A$}
  \RightLabel{\Intro\land}
  \BinaryInfC{$!B \land !A$}
  \DischargeRule{\Intro\lif}{x}
  \UnaryInfC{$(!A \land !B) \lif (!B \land !A)$}
  \DisplayProof
  \]
  \Log{N2i} وچ ایہ ایویں دِسدا اے:
  \[
  \Axiom$x : !A\land !B \fCenter !A\land !B$
  \RightLabel{\Elim\land[2]}
  \UnaryInf$x : !A\land !B \fCenter !B$
  \Axiom$x : !A\land !B \fCenter !A\land !B$
  \RightLabel{\Elim\land[1]}
  \UnaryInf$x : !A\land !B \fCenter !A$
  \RightLabel{\Intro\land}
  \BinaryInf$x : !A\land !B \fCenter !B \land !A$
  \RightLabel{\Intro\lif}
  \UnaryInf$\fCenter (!A \land !B) \lif (!B \land !A)$
  \DisplayProof
  \]
  اسیں ثبوت ٹرم اُتّوں تھلّے ول دیندے آں۔ مسلّمیاں دے ثبوت ٹرم اوہی نیں جیہڑا سیاق وچ متغیر اے، یعنی~$x$۔ فیر $ \Elim\land[i]$ نال جُڑے ثبوت ٹرم $\proj{2}{x}$ تے~$\proj{1}{x}$ طے ہوندے نیں۔ جدوں اسیں \Intro{\land} لاؤنے آں، ایہ دونے ثبوت ٹرم رلا دیندے آں، $ \pair{\proj{2}{x}}{\proj{1}{x}}$۔
  آخر \Intro{\lif} قاعدہ $\lambd$-تجرید لاؤندا اے،~$x$ نوں بندھدا اے تے ایہ درج کردا اے کہ $x$ مفروضے~$!A \land
  !B$ دا لیبل سی، $ \lambd[\typeof{x}{!A \land
  !B}][\pair{\proj{2}{x}}{\proj{1}{x}}]$۔
  فیر !!{proof}، \Log{tN2i} وچ، ایہ اے:
  \[
  \Axiom$x : !A\land !B \fCenter x: !A\land !B$
  \RightLabel{\Elim\land[2]}
  \UnaryInf$x : !A\land !B \fCenter \proj{2}{x} : !B$
  \Axiom$x : !A\land !B \fCenter x: !A\land !B$
  \RightLabel{\Elim\land[1]}
  \UnaryInf$x : !A\land !B \fCenter \proj{1}{x} : !A$
  \RightLabel{\Intro\land}
  \BinaryInf$x : !A\land !B \fCenter \pair{\proj{2}{x}}{\proj{1}{x}} : 
    !B \land !A$
  \RightLabel{\Intro\lif}
  \UnaryInf$\fCenter \lambd[\typeof{x}{!A \land
  !B}][\pair{\proj{2}{x}}{\proj{1}{x}}] : (!A \land !B) \lif (!B \land !A)$
  \DisplayProof
  \]
\end{ex}
\endgroup
% /OLP-0689
% OLP-0695 content/proof-theory/propositions-as-types/terms-to-proofs.tex
\begingroup
\def\olfilename{terms-to-proofs}
\expandafter\def\csname ol@sectioncs\endcsname{section}

\olfileid{int}{pty}{tp}

\olsection{ثبوت ٹرماں توں \usetoken{P}{derivation} مُڑ لبھنا}

ہُن دوجے پاسے ویکھدے آں، یعنی ثبوت ٹرم~$N$ نوں مُڑ !!a{proof} وچ ترجمہ کرنا، \Log{N2i} وچ۔ عام طور اُتے ایہ صرف اوس سیاق دے لحاظ نال ہو سکدا اے جیہدے وچ $x : !A$ جوڑا ہر اوس متغیر~$x$ لئی ہووے جیہڑا دِتّے ثبوت ٹرم وچ آزاد آؤندا اے۔ پر پہلاں ایہو جیہی مثال لئیے جیہدے وچ کوئی آزاد متغیر نہیں، یعنی ایہ ٹرم:
\[
\lambd[\typeof{x}{!A \land
  !B}][\pair{\proj{2}{x}}{\proj{1}{x}}]
\]
ایہ $\lambd[\typeof{x}{!A \land !B}][N_1]$ دی شکل دا اے۔ کیونکہ \Log{tN2i} دا اکّو قاعدہ جیہڑا ایس شکل دا ٹرم بناؤندا اے،~\Intro{\lif} اے، ایس لئی سانوں پتا اے کہ !!{proof} جیہڑا~$N$ کوڈ کردا اے،~\Intro{\lif} اُتے ضرور مُکدا اے، یعنی ایویں:
\[
  \Axiom$x : !A\land !B \fCenter \pair{\proj{2}{x}}{\proj{1}{x}} : 
    !C$
  \RightLabel{\Intro\lif}
  \UnaryInf$\fCenter \lambd[\typeof{x}{!A \land
  !B}][\pair{\proj{2}{x}}{\proj{1}{x}}] : (!A \land !B) \lif !C$
  \DisplayProof
  \]
سانوں ہجے پتا نہیں کہ $!C$ کیہ اے، پر سانوں پتا اے کہ اگلا جزو $!A \land !B$ اے تے مقدمے دے سیاق وچ $x:!A \land !B$ ضرور ہووے گا۔

ہُن مقدمے نال جُڑیا ثبوت ٹرم پہلاں ای طے اے، $ \pair{\proj{2}{x}}{\proj{1}{x}}$۔ !!{proof} جیہڑا \emph{ایہ ٹرم} کوڈ کردا اے، \Intro{\land} اُتے ضرور مُکدا اے۔ سانوں ہجے پتا نہیں کہ $!C$ کیہ اے، پر کیونکہ ایہ~\Intro{\land} دا نتیجہ اے، سانوں پتا اے کہ ایہ عطف ضرور اے، یعنی $!C \ident (!C_1 \land !C_2)$۔ اسیں ثبوت مُڑ بناؤندے آں:
\[
  \Axiom$x : !A\land !B \fCenter \proj{2}{x} : !C_1$
  \Axiom$x : !A\land !B \fCenter \proj{1}{x} : !C_2$
  \RightLabel{\Intro\land}
  \BinaryInf$x : !A\land !B \fCenter \pair{\proj{2}{x}}{\proj{1}{x}} : 
    !C_1 \land !C_2$
  \RightLabel{\Intro\lif}
  \UnaryInf$\fCenter \lambd[\typeof{x}{!A \land
  !B}][\pair{\proj{2}{x}}{\proj{1}{x}}] : (!A \land !B) \lif (!C_1 \land !C_2)$
  \DisplayProof
  \]
ثبوت ٹرم $\pi_2(x)$ دا مطلب اے کہ کھبّے پاسے دا اُتلا سیکوئنٹ \Elim{\land}[2] نال ملیا سی، یعنی مقدمے دی شکل $!D
\land !C_1$ اے۔ سجّے پاسے اسیں طے کر سکدے آں کہ $!C_2$ تک پہنچاؤن والا استنباط \Elim{\land}[1] ضرور اے، تے مقدمے دی شکل $!C_2 \land !E$ اے۔
\[
  \Axiom$x : !A\land !B \fCenter x: !D\land !C_1$
  \RightLabel{\Elim\land[2]}
  \UnaryInf$x : !A\land !B \fCenter \proj{2}{x} : !C_1$
  \Axiom$x : !A\land !B \fCenter x: !C_2\land !E$
  \RightLabel{\Elim\land[1]}
  \UnaryInf$x : !A\land !B \fCenter \proj{1}{x} : !C_2$
  \RightLabel{\Intro\land}
  \BinaryInf$x : !A\land !B \fCenter \pair{\proj{2}{x}}{\proj{1}{x}} : 
    !C_1 \land !C_2$
  \RightLabel{\Intro\lif}
  \UnaryInf$\fCenter \lambd[\typeof{x}{!A \land
  !B}][\pair{\proj{2}{x}}{\proj{1}{x}}] : (!A \land !B) \lif (!C_1 \land !C_2)$
  \DisplayProof
  \]
ہُن کیونکہ سبھ توں اُتلے سیکوئنٹاں دے ثبوت ٹرم متغیر نیں، سانوں پتا اے کہ اوہ مسلّمے ضرور نیں۔ ایس نال $!C_1$، $!C_2$، $!D$ تے $!E$ طے ہوندے نیں۔ $!D \ident !A$ تے $!C_1 \ident !B$، نئیں تاں کھبّا سیکوئنٹ مسلّمہ نہیں۔ تے $!C_2 \ident !A$ تے $!E \ident !B$، نئیں تاں سجّا سیکوئنٹ مسلّمہ نہیں۔ اسیں پورا ثبوت مُڑ لبھ لیا اے:
\[
  \Axiom$x : !A\land !B \fCenter x: !A\land !B$
  \RightLabel{\Elim\land[2]}
  \UnaryInf$x : !A\land !B \fCenter \proj{2}{x} : !B$
  \Axiom$x : !A\land !B \fCenter x: !A\land !B$
  \RightLabel{\Elim\land[1]}
  \UnaryInf$x : !A\land !B \fCenter \proj{1}{x} : !A$
  \RightLabel{\Intro\land}
  \BinaryInf$x : !A\land !B \fCenter \pair{\proj{2}{x}}{\proj{1}{x}} : 
    !B \land !A$
  \RightLabel{\Intro\lif}
  \UnaryInf$\fCenter \lambd[\typeof{x}{!A \land
  !B}][\pair{\proj{2}{x}}{\proj{1}{x}}] : (!A \land !B) \lif (!B \land !A)$
  \DisplayProof
  \]
\endgroup
% /OLP-0695
% OLP-0697 content/proof-theory/propositions-as-types/types.tex
\begingroup
\def\olfilename{types}
\expandafter\def\csname ol@sectioncs\endcsname{section}

\olfileid{int}{pty}{typ}

\olsection{نوعاں}

ثبوت ٹرم، جیویں $(MN)$، اپنے آپ اوہ ثبوت اکّو طرح طے نہیں کردا جیہدا اوہ ٹرم اے۔ وجہ ایہ اے کہ آزاد متغیر، یعنی مسلّمیاں دے ثبوت ٹرم، میل کھاندے مسلّمیاں وچلے !!{formula}s نہیں دَسّدے۔ پر جے اسیں ایہناں نوں سیاق~$\Gamma$ وچ دَس دیئیے، تاں ثبوت \emph{اکّو طرح} طے ہو جاندا اے۔

\begin{defn}
ثبوت ٹرم~$N$ لئی اک \emph{سیاق} اوہ سیٹ $ \Gamma$ اے ایویں کہ~$N$ دے ہر آزاد متغیر لئی $\Gamma$ وچ عین اک جوڑا $x : !A$ اے۔
\end{defn}

ویکھو کہ تعریف ایہ شرط نہیں لاؤندی کہ $\Gamma$ وچ $x:!A$ \emph{صرف}~$N$ دے آزاد متغیراں لئی ہووے۔ مثال لئی $\Gamma =
\{x:!A, y:!B, z:!C\}$، $\pair{x}{y}$ دا سیاق اے۔

\begin{defn}\ollabel{defn:type}
  منّ لو $\Gamma$،~$N$ لئی اک سیاق اے۔~$N$ دی \emph{نوع} (~$\Gamma$ دے لحاظ نال) استقرائی طور اُتے ایویں متعین اے:
  \begin{enumerate}
  \item $x$ دی نوع $!A$ اے اوہدوں تے صرف اوہدوں جے $x:!A \in \Gamma$۔
  \item جے $N$ دی نوع~$!B$ اے $x:!A, \Gamma$ دے لحاظ نال، تاں $\lambd[\typeof{x}{!A}][N]$ دی نوع $!A \lif !B$ اے۔
  \item جے $N$ دی نوع $!A \lif !B$ اے تے $M$ دی نوع~$!A$ اے، تاں $(NM)$ دی نوع~$!B$ اے۔
  \item جے $N$ دی نوع $!A$ اے تے $M$ دی نوع~$!B$ اے، تاں $\pair{N}{M}$ دی نوع $!A \land !B$ اے۔
  \item جے $N$ دی نوع~$!A_1 \land !A_2$ اے، تاں $\proj{i}{N}$ دی نوع~$!A_i$ اے۔
  \item جے $N$ دی نوع $!A$ اے، تاں $\inj[!B]{i}{!M}$ دی نوع $!A \lor !B$ اے جے $i=1$ تے نوع $!B \lor
    !A$ اے جے $i=2$۔
  \item جے $M$ دی نوع $!A \lor !B$ اے، تے $N_1$ دی نوع $!C$ اے $x:!A, \Gamma$ دے لحاظ نال، تے $N_2$ دی نوع~$!C$ اے $y:!B,
  \Gamma$ دے لحاظ نال، تاں $\dcase{M}{x_1}{N_1}{x_2}{N_2}$ دی نوع~$!C$ اے۔
  \item جے $N$ دی نوع~$\lfalse$ اے، تاں $\abort{!A}{N}$ دی نوع~$!A$ اے۔
  \end{enumerate}
\end{defn}

\begin{prop}
ہر ثبوت ٹرم~$N$ دی عین اک نوع اے، سیاق~$\Gamma$ دے لحاظ نال جیہڑا~$N$ لئی اے۔
\end{prop}

\begin{proof}
  ~$N$ اُتے استقرا نال۔
\end{proof}

اسیں $\Gamma \Sequent N : !A$ لکھدے آں جے $N$ دی نوع $\Gamma$ دے لحاظ نال~$!A$ اے۔ \olref{defn:type} دیاں شرطاں پہلاں توں جانیاں لگن گئیاں۔ ایہ عین \Log{tN2i} دے قاعدے نیں، اک نِکّے فرق نال۔ سیاق~$\Gamma$ ہر تھاں اوہی اے۔ اسیں تعریف نوں \Log{tN3i} حساب دے روپ وچ لکھ سکدے آں، \olref{tab:tN3ip} وچ دِتّے قاعدیاں نال۔

% OLP-0693 content/proof-theory/propositions-as-types/rules-tN3.tex
\begingroup

\begin{table}
\begin{tabular}{@{}c@{}c@{}}
\multicolumn{2}{l}{\textbf{مسلّمے}}\\
\multicolumn{2}{c}{$\Gamma \Sequent x:!A$ جے $x:!A \in \Gamma$}\\[3ex]
\multicolumn{2}{l}{\textbf{استنباط دے قاعدے}}\\[3ex]
\multicolumn{2}{c}{\Axiom$x: !A, \Gamma \fCenter N:!B$
\RightLabel{\Intro{\lif}}
\UnaryInf$\Gamma \fCenter \lambd[\typeof{x}{!A}][N] : !A \lif !B $
\DisplayProof}
\\[4ex]
\multicolumn{2}{c}{\Axiom$\Gamma \fCenter N:!A \lif !B$
\Axiom$\Gamma \fCenter M:!A$
\RightLabel{\Elim{\lif}}
\BinaryInf$\Gamma \fCenter (NM) : !B$
\DisplayProof}
\\[4ex]
\multicolumn{2}{c}{\Axiom$\Gamma \fCenter N_1 : !A$
\Axiom$\Gamma \fCenter N_2 : !B$
\RightLabel{\Intro{\land}}
\BinaryInf$\Gamma \fCenter \pair{N_1}{N_2} : !A \land !B $
\DisplayProof}
\\[4ex]
\Axiom$\Gamma \fCenter N: !A \land !B$
\RightLabel{\Elim{\land}[1]}
\UnaryInf$\Gamma \fCenter \proj{1}{N} : !A$
\DisplayProof
&
\Axiom$\Gamma \fCenter N:!A \land !B$
\RightLabel{\Elim{\land}[2]}
\UnaryInf$\Gamma \fCenter \proj{2}{N} : !B$
\DisplayProof\\[3ex]
\\[4ex]
\Axiom$\Gamma \fCenter N: !A$
\RightLabel{\Intro{\lor}[1]}
\UnaryInf$\Gamma \fCenter \inj[!B]{1}{N} : !A \lor !B$
\DisplayProof
&
\Axiom$ \Gamma \fCenter N : !A$
\RightLabel{\Intro{\lor}[2]}
\UnaryInf$\Gamma \fCenter \inj[!B]{2}{N} : !B \lor !A$
\DisplayProof\\[3ex]
\multicolumn{2}{c}{
\begin{tabular}[b]{@{}l@{}}
\Axiom$\Gamma \fCenter M : !A \lor !B$
\Axiom$x:!A, \Gamma \fCenter N_1 : !C$
\Axiom$y:!B, \Gamma \fCenter N_2: !C$
\RightLabel{\Elim{\lor}}
\TrinaryInf$\Gamma \fCenter
\dcase{M}{x}{N_1}{y}{N_2} : !C$
\DisplayProof
\end{tabular}
}
\\[4ex]
\multicolumn{2}{c}{\Axiom$\Gamma \fCenter N: \lfalse$
\RightLabel{\FalseInt}
\UnaryInf$\Gamma \fCenter \abort{!A}{N} : !A$
\DisplayProof}
\end{tabular}
\caption{ٹرماں دے لیبل والے قضیاتی وجدانی قدرتی اخذ نظام \Log{tN3ip} دے قاعدے جیہڑا نوع مقرر کرن والا نظام وی اے۔}
\ollabel{tab:tN3ip}
\end{table}
\endgroup
% /OLP-0693

ثبوت ٹرم $\lambd$-کیلکولس دی اک صورت دے ٹرم نیں، یعنی \emph{حاصلِ ضرب، جوڑ تے خالی نوع والا نوع بند $\lambd$-کیلکولس}۔ کلاسیکی $\lambd$-کیلکولس دے ٹرم صرف متغیراں، \emph{$\lambd$-تجریداں}~$\lambd[x][N]$ تے لاگو کرن $(MN)$ نال بندے نیں۔ اوہدے وچ~$!A$ دی اوہ نشان دہی نہیں جیہڑی نوع بند $\lambd$-تجرید $\lambd[\typeof{x}{!A}][N]$ وچ اے۔ $\lambd$-کیلکولس حسابی عمل دا ٹیورنگ-کامل نمونہ اے (یعنی ہر جزوی قابلِ حساب فنکشن اوہدے وچ کسے ٹرم نال ظاہر ہو سکدا اے)۔ ثبوت ٹرماں تے نوع بندی دے قاعدیاں نال دِتّی $\lambd$-کیلکولس دی صورت حسابی عمل دا محدود نمونہ اے، پر بنیادی خیال اوہی نیں۔ ایہ نوع مقرر کرن نال محدود اے۔ صرف \emph{ٹھیک نوع والے} ٹرم، یعنی نوع رکھنے والے ثبوت ٹرم (سیاق~$\Gamma$ دے لحاظ نال)، جائز ٹرم نیں۔ مثال لئی $\lambd[x][(xx)]$ کلاسیکی بناں نوع بندی والے $\lambd$-کیلکولس وچ جائز ٹرم اے، پر $\lambd[\typeof{x}{!A}][(xx)]$ کسے وی $!A$ لئی ٹھیک نوع والا نہیں۔ وجہ ایہ اے کہ نوع بندی دے قاعدیاں مطابق $(xx)$ دی کوئی نوع ہون لئی پہلے~$x$ دی نوع $!A \lif !B$ تے دوجے $x$ دی نوع~$!A$ ہونی پئے گی۔ ایہ ممکن نہیں کیونکہ $!A$،~$!A \lif !B$ دا اپنے توں وکھرا ذیلی !!{formula} اے۔

سہج سمجھ نال نوع نوں اوس شے دی قسم سمجھ سکدے آں جیہدا ناں اوس نوع دا ٹرم دیندا اے۔ سو نوع $!A \land !B$ دا ٹرم اک جوڑا چُندا اے جیہدا پہلا جزو~$!A$ نوع دا تے دوجا جزو~$!B$ نوع دا اے، یعنی !!a{element} حاصلِ ضرب~$!A \times !B$ دا۔ نوع $!A \lif !B$ دا ٹرم $!A$ نوع دیاں شیواں توں~$!B$ نوع دیاں شیواں تک اک فنکشن اے۔ نوع $!A \lor !B$ دا ٹرم یا $!A$ نوع دی شے اے یا~$!B$ نوع دی، نال ایہ جانکاری کہ کیہڑی اے۔ ایہ دوواں سیٹاں $!A$ تے $!B$ دے بے سانجھے اتحاد ورگا اے، جیہڑی تعریف $ \{\tuple{i,x}
: i = 1 \text{ یا } 2, x \in !A \text{ جے } i = 1, x \in !B \text{ جے
} i=2\}$ اے۔ نوع $\lfalse$ خالی سیٹ اے، یعنی جے کسے ٹرم دی اوہ نوع اے تاں اوہ کسے شے دا ناں نہیں دے سکدا۔ کیونکہ قدرتی اخذ (\Log{tN3i} سمیت) بے تضاد اے، کوئی !!{proof} نہیں~$\lfalse$ دا، جے کھُلّے مفروضے نہ ہون، ایس لئی کوئی بند ثبوت ٹرم (آزاد متغیراں توں بناں ٹرم) نہیں جیہدی نوع~$\lfalse$ ہووے۔

نوع بندی والے $\lambd$-کیلکولس دے عمل کار ایہو جیہے ٹرم بناؤندے نیں جیہڑے سہج سمجھ نال ٹھیک قسم دی شے دا ناں دیندے نیں، جے اوہ ٹرم جنہاں اُتے ایہ لائے جاندے نیں، خاص شیواں دے ناں دیندے ہون۔ مثال لئی «جے $N_1 : !A$ تے $N_2 : !B$، تاں $\pair{N_1}{N_2} : !A \land !B$» دا مطلب اے کہ جے $N_1$،~$!A$ نوع دی شے دا ناں دیندا اے، تے $N_2$،~$!B$ نوع دی شے دا، تاں $\pair{N_1}{N_2}$، $!A \land !B$ نوع دی شے دا ناں دیندا اے۔
\endgroup
% /OLP-0697
% OLP-0691 content/proof-theory/propositions-as-types/reduction.tex
\begingroup
\def\olfilename{reduction}
\expandafter\def\csname ol@sectioncs\endcsname{section}

\olfileid{int}{pty}{red}

\olsection{تقلیل}

قدرتی اخذ دے !!{derivation}s وچ، جے !!a{introduction} قاعدے توں پچھوں !!a{elimination} قاعدہ آوے، تاں ایہ اک چکر اے جیہڑا «گھٹایا» جا سکدا اے۔ مثال لئی $\Intro{\land}$ توں پچھوں $\Elim{\land}$ آوے تاں دونے «اک دوجے نوں کٹ دیندے نیں»، کیونکہ $\Elim{\land}$ دا نتیجہ ہمیشہ اوس عطف دے دو جزواں وچوں اک اے جیہڑا \Intro{\land} بناؤندا اے، تے $\Intro{\land}$ دے مقدمے اوہ دونے جزو نیں۔ سو جے سانوں !!a{derivation} دوواں وچوں کسے جزو دا چاہیدا اے، تاں اسیں اوہی !!{derivation} ورت سکدے آں، عطف نوں شامل کر کے فیر ہٹاؤن دی لوڑ نہیں۔ ایہ ملدا اے:
\begin{gather*}
  \bottomAlignProof
  \AxiomC{}
  \RightLabel{$\delta_1$}
  \DeduceC{$!A_1$}
  \AxiomC{}
  \RightLabel{$\delta_2$}
  \DeduceC{$!A_2$}
  \RightLabel{$\Intro{\land}$}
  \BinaryInfC{$!A_1 \land !A_2$}
  \RightLabel{$\Elim{\land}_i$}
  \UnaryInfC{$!A_i$}
  \DisplayProof\bottomAlignProof
  \quad
  \redone
  \quad
  \AxiomC{}
  \RightLabel{$\delta_i$}
  \DeduceC{$!A_i$}
  \DisplayProof
\end{gather*}

ایس سادے کرن نال جُڑے ثبوت ٹرماں اُتے وی تقلیل دا ایہو جیہا تعلق سوچ سکدے آں۔ جے $N_1$،~$\delta_1$ دا ثبوت ٹرم اے تے $N_2$،~$\delta_2$ دا، تاں اسیں آکھ سکدے آں کہ کھبّے پاسے دے ثبوت دا ثبوت ٹرم (اک قدم وچ) سجّے پاسے دے ثبوت دے ٹرم تک گھٹ جاندا اے:
\[
  \proj{i}{\pair{N_1, N_2}} \redone N_i
\]
نوع بندی والے لامبڈا کیلکولس وچ ایہ حاصلِ ضرب دی نوع لئی \emph{بیٹا تقلیل} دا قاعدہ اے۔

قدرتی اخذ وچ استنباطاں دے ایہو جیہے جوڑے نوں \emph{قطع} آکھدے آں، جیویں \Intro{\land} پچھوں \Elim{\land}، یعنی ایسا جوڑا جیہڑا گھٹایا جا سکدا اے۔ نوع بندی والے لامبڈا کیلکولس وچ $\redone$ دے کھبّے پاسے دے ٹرم نوں \emph{ریڈیکس} تے سجّے پاسے دے ٹرم نوں \emph{کانٹریکٹم} آکھدے آں۔ بناں نوع بندی والے لامبڈا کیلکولس وچ صرف $(\lambd[x][N])Q$ نوں ریڈیکس منّیا جاندا اے۔ نوع بندی والے لامبڈا کیلکولس وچ نحوی بناوٹ نوں $ \pair{N}{M}$، $\proj{i}{N}$، $\inj{i}{!A}{N}$، $\dcase{N}{x_1}{M_1}{x_2}{M_2}$ تے $\abort{!A}{N}$ والے ٹرماں تک ودھایا جاندا اے، تے ایہناں نال میل کھاندے ریڈیکس وی نیں۔

ایویں ای فصل لئی وی تقلیل اے:
\begin{gather*}
  \bottomAlignProof
  \AxiomC{}
  \RightLabel{$\delta$}
  \DeduceC{$!A_i$}
  \RightLabel{$\Intro{\lor}$}
  \UnaryInfC{$!A_1 \lor !A_2$}
  \AxiomC{$\Discharge{!A_1}{x_1}$}
  \RightLabel{$\delta_1$}
  \DeduceC{$!C$}
  \AxiomC{$\Discharge{!A_2}{x_2}$}
  \RightLabel{$\delta_2$}
  \DeduceC{$!C$}
  \DischargeRule{$\Elim{\lor}$}{x_1,x_2}
  \TrinaryInfC{$!C$}
  \DisplayProof\bottomAlignProof
  \quad
  \redone
  \quad
  \AxiomC{}
  \RightLabel{$\delta$}
  \DeduceC{$!A_i$}
  \RightLabel{$\delta_i$}
  \DeduceC{$!C$}
  \DisplayProof
\end{gather*}
ایہ ثبوت ٹرماں اُتے ایس تقلیل نال میل کھاندا اے:
\begin{gather*}
  \dcase{\inj[!A_i]{i}{M}}{x_1}{N_1}{x_2}{N_2} \redone \Subst{N_i}{M}{x_i}
\end{gather*}
جتھے $M$،~$\delta$ دا ثبوت ٹرم اے تے $N_i$،~$\delta_i$ دا۔ ایہ جوڑ دیاں نوعاں لئی بیٹا تقلیل دا قاعدہ اے۔ ایتھے $\Subst{M}{N_i}{x}$ دا مطلب اے متغیر~$x$ دیاں ساریاں آزاد موجودگیاں نوں $M$ وچ~$N_i$ نال بدل دیو۔

فنکشن دی نوع دی تقلیل، اوہ چکر ہٹاؤن نال میل کھاندی اے جتھے $\Intro\lif$ توں پچھوں $\Elim\lif$ آؤندا اے۔
\begin{gather*}
  \bottomAlignProof
  \AxiomC{$\Discharge{!A}{x}$}
  \RightLabel{$\delta$}
  \DeduceC{$!B$}
  \DischargeRule{$\Intro{\lif}$}{x}
  \UnaryInfC{$!A \lif !B$}
  \AxiomC{}
  \RightLabel{$\delta'$}
  \DeduceC{$!A$}
  \RightLabel{$\Elim{\lif}$}
  \BinaryInfC{$!B$}
  \DisplayProof\bottomAlignProof
  \quad
  \redone
  \quad
  \AxiomC{}
  \RightLabel{$\delta'$}
  \DeduceC{$!A$}
  \RightLabel{$\delta$}
  \DeduceC{$!B$}
  \DisplayProof
\end{gather*}
ثبوت ٹرماں لئی ایہ عام بیٹا تقلیل اے:
\begin{gather*}
  (\lambd[\typeof{x}{!A}][N]M)
  \redone \Subst{N}{M}{x}
\end{gather*}

!!{proof}s اُتے تقلیل والیاں تبدیلیاں توں علاوہ، اسیں تھاں بدلّن والیاں تبدیلیاں وی سوچ سکدے آں۔ ایہ اوہناں (ذیلی) !!{proof}s اُتے لگدیاں نیں جیہڑے \Elim{\lor} توں پچھوں ہور !!{elimination} قاعدے اُتے مُکدے نیں۔ مثال:
\begin{multline*}
  \bottomAlignProof
  \AxiomC{}
  \RightLabel{$\delta$}
  \DeduceC{$!A_i$}
  \RightLabel{$\Intro{\lor}$}
  \UnaryInfC{$!A_1 \lor !A_2$}
  \AxiomC{$\Discharge{!A_1}{x_1}$}
  \RightLabel{$\delta_1$}
  \DeduceC{$!C_1 \land !C_2$}
  \AxiomC{$\Discharge{!A_2}{x_2}$}
  \RightLabel{$\delta_2$}
  \DeduceC{$!C_1 \land !C_2$}
  \DischargeRule{$\Elim{\lor}$}{x_1,x_2}
  \TrinaryInfC{$!C_1 \land !C_2$}
  \RightLabel{\Elim{\land}[i]}
  \UnaryInfC{$!C_i$}
  \DisplayProof\bottomAlignProof
  \quad
  \redone
  \\
  \AxiomC{}
  \RightLabel{$\delta$}
  \DeduceC{$!A_i$}
  \RightLabel{$\Intro{\lor}$}
  \UnaryInfC{$!A_1 \lor !A_2$}
  \AxiomC{$\Discharge{!A_1}{x_1}$}
  \RightLabel{$\delta_1$}
  \DeduceC{$!C_1 \land !C_2$}
  \RightLabel{\Elim{\land}[i]}
  \UnaryInfC{$!C_i$}
  \AxiomC{$\Discharge{!A_2}{x_2}$}
  \RightLabel{$\delta_2$}
  \DeduceC{$!C_1 \land !C_2$}
  \RightLabel{\Elim{\land}[i]}
  \UnaryInfC{$!C_i$}
  \DischargeRule{$\Elim{\lor}$}{x_1,x_2}
  \TrinaryInfC{$!C_i$}
  \DisplayProof
\end{multline*}
جے $ \delta$، $\delta_1$ تے $\delta_2$ دے ثبوت ٹرم بالترتیب $M$، $N_1$ تے $N_2$ نیں، تاں ثبوت ٹرماں، یعنی $\lambd$ ٹرماں لئی میل کھاندا تقلیل دا قاعدہ ایہ اے:
\[
  \proj{i}{\dcase{M}{x_1}{N_1}{x_2}{N_2}} \redone 
\dcase{M}{x_1}{\proj{i}{N_1}}{x_2}{\proj{i}{N_2}}
\]

اسیں چکر !!{proof}s دے آخر توں ہٹاؤن دے نال نال !!{proof}s دے اندر توں وی ہٹا سکدے آں۔ ایس لئی تقلیل والی تبدیلی اوس ذیلی !!{proof} اُتے لاؤ جیہڑا چکر اُتے مُکدا اے۔ ایویں ای $\beta$-تقلیل $\lambd$-ٹرماں دے ذیلی ٹرماں اُتے وی لگدی اے۔ جے $N$، $M$ دا ذیلی ٹرم اے تے $N \redone
N'$، تاں ذیلی ٹرم $N$ نوں~$M$ وچ $N'$ نال بدل کے $M'$ حاصل کر سکدے آں۔ ایس صورت وچ اسیں $M \redone M'$ وی لکھدے آں۔ جے اک لڑی $M = M_1
\redone M_2 \redone M_2 \redone \dots \redone M_k = M'$ اے (ایہ صورت وی شامل اے جتھے $k=1$، یعنی $M = M'$)، تاں اسیں $M \red M'$ لکھدے آں۔

!!{proof}، جیہدے کسے وی ذیلی !!{proof} اُتے کوئی تبدیلی نہیں لگ سکدی، نوں \emph{نارمل} آکھدے آں۔ ایویں ای $\lambd$-ٹرم جیہدے اُتے (یا جیہدے کسے ذیلی ٹرم اُتے) کوئی تبدیلی نہیں لگ سکدی، \emph{نارمل} اے۔ نارمل $\lambd$-ٹرم نوں \emph{نارمل روپ} وی آکھدے آں۔
\endgroup
% /OLP-0691
% OLP-0696 content/proof-theory/propositions-as-types/type-preservation.tex
\begingroup
\def\olfilename{type-preservation}
\expandafter\def\csname ol@sectioncs\endcsname{section}

\olfileid{int}{pty}{tpr}

\olsection{نوع دی بقا}

جیہڑیاں تقلیل تے تھاں بدلّن والیاں تبدیلیاں اسیں ویکھیاں نیں، اوہ سیدھیاں \Log{tN2i} تے \Log{tN3i} نظاماں لئی وی متعین کیتیاں جا سکدیاں نیں۔ \Log{tN3i} نظام $\lambd$-ٹرم~$N$ دی نوع، سیاق~$\Gamma$ دے لحاظ نال، اوہ !!{formula}~$!A$ قرار دیندا اے جیہدے لئی $ \Gamma \Sequent N
: !A$ دا !!a{proof}، \Log{tN3i} وچ، موجود اے۔ $\beta$-تقلیل قاعدیاں، جیہڑے $\lambd$-ٹرماں دے نیں، تے !!{proof}s دیاں تقلیل تے تھاں بدلّن والیاں تبدیلیاں دے عین میل دا اک اہم نتیجہ اے۔ $\lambd$-ٹرم دی نوع، سیاق~$\Gamma$ دے لحاظ نال،، $\beta$-تبدیلی توں پچھوں وی اوہی رہندی اے۔ ایس خاصیت نوں \emph{نوع دی بقا} آکھدے آں۔

\begin{prop}
جے $N$ دی نوع~$!A$ اے $ \Gamma$ دے لحاظ نال، تے $N \redone N'$، تاں $N'$ دی نوع~$!A$ اے~$\Gamma$ دے لحاظ نال۔
\end{prop}

\begin{proof}
  ~$N$ تے !!{proof}s~$\delta$ دی اُچائی اُتے استقرا نال، جیہڑے $ \Gamma
  \Sequent N : !A$ دے نیں، اسیں ہِٹھلی گل سوکھے ای ثابت کر سکدے آں۔ جے $M$،~$N$ دا ذیلی ٹرم اے، تاں $\delta$ وچ اک ذیلی !!{proof}~$\delta_1$ اے جیہڑا $ \Gamma' \Sequent M : !B$ اُتے مُکدا اے۔ جے $M$ ریڈیکس اے، تاں~$\delta_1$ اُتے اک تبدیلی لگدی اے۔ تقلیل $\delta_1$ اُتے لا کے اسیں !!a{proof}~$\delta_1'$ لیندے آں $ \Gamma' \Sequent M'
  : !B$ دا۔ \Log{tN3i} دیاں تبدیلیاں ویکھ کے پتا لگدا اے کہ $M'$،~$M$ دی تقلیل نال ملن والا ٹرم اے۔ $ \delta_1$ نوں~$\delta$ وچ $ \delta_1'$ نال بدلّن نال !!a{proof}~$\delta'$ ملدا اے $ \Gamma \Sequent N' : !A$ دا، جتھے $N'$،~$N$ توں ذیلی ٹرم~$M$ نوں~$M'$ نال بدل کے ملدا اے۔

  مثال لئی~\Intro{\land} توں پچھوں \Elim{\land} لاؤن والی تقلیل دی تبدیلی ویکھو:
  \[
  \bottomAlignProof
  \AxiomC{}
  \Deduce$\Gamma \fCenter N_1 : !B_1$
  \AxiomC{}
  \Deduce$\Gamma \fCenter N_2 : !B_2$
  \RightLabel{\Intro{\land}}
  \BinaryInf$\Gamma \fCenter \pair{N_1}{N_2} : !B_1 \land !B_2$
  \RightLabel{\Elim{\land}[i]}
  \UnaryInf$\Gamma \fCenter \proj{i}{\pair{N_1}{N_2}} : !B_i$
  \DisplayProof
  \quad\redone\quad
  \bottomAlignProof
  \AxiomC{}
  \Deduce$\Gamma \fCenter N_i : !B_i$
  \DisplayProof
  \]
  اسیں ویکھدے آں کہ سجّے پاسے دے !!{proof} دا ثبوت ٹرم، یعنی $ \delta_1'$ دا، $N_i$ اے۔ ایہ عین اوہ نتیجہ اے جیہڑا $\beta$-تبدیلی، ثبوت ٹرم $ \proj{i}{\pair{N_1}{N_2}}$ اُتے لاؤن نال ملدا اے۔ ایہ کھبّے پاسے دے !!{proof}~($\delta_1$) دا ٹرم اے۔

  یا \Intro{\lif} توں پچھوں \Elim{\lif} استنباط ویکھو، تے اوہدے اُتے لگی تقلیل والی تبدیلی:
  \begin{multline*}
  \bottomAlignProof
  \AxiomC{$\delta_2$}
  \Deduce$x: !B_1, \Gamma \fCenter N : !B_2$
  \RightLabel{\Intro{\lif}}
  \UnaryInf$\Gamma \fCenter \lambd[\typeof{x}{!B_1}][N] : !B_1 \lif !B_2$
  \AxiomC{}
  \RightLabel{$\delta_3$}
  \Deduce$\Gamma \fCenter M : !B_1$
  \RightLabel{\Elim{\lif}}
  \BinaryInf$\Gamma \fCenter (\lambd[\typeof{x}{!B_1}][N] M) : !B_2$
  \DisplayProof
  \quad\redone\\
  \bottomAlignProof
  \AxiomC{}
  \RightLabel{$\delta_3$}
  \Deduce$\Gamma \fCenter M : !B_1$
  \RightLabel{$\Subst{\delta_2}{\delta_3}{x:!B_1}$}
  \Deduce$\Gamma \fCenter \Subst{N}{M}{x} : !B_2$
  \DisplayProof
  \end{multline*}
  فیر وی سجّے پاسے دے !!{proof} دا ثبوت ٹرم عین اوہی اے جیہڑا $\beta$-تقلیل نال کھبّے پاسے دے !!{proof} دے ثبوت ٹرم توں ملدا اے۔ بے شک ایہ گل غور نال جانچنی پئے گی (~$\delta_2$ دی اُچائی اُتے استقرا نال) کہ جے اسیں $ \Gamma' \Sequent x:!B_1$ دی شکل والے سارے مسلّمے~$\delta_2$ وچ !!{proof}~$\delta_3$ نال بدلیے، جیہڑا $ \Gamma \Sequent M:!B_1$ دا اے، تاں ایہدے نال واقعی !!a{proof} $ \Subst{\delta_2}{\delta_3}{x:!B_1}$ ملدا اے $ \Gamma \Sequent
  \Subst{N}{M}{x} : !B_2$ دا۔
\end{proof}

!!{proof}s تے $\lambd$-ٹرماں دے نیڑے میل دا اک ہور نتیجہ اے۔ جے $N$ اک $\lambd$-ٹرم اے، $!A$ نوع دا سیاق~$\Gamma$ دے لحاظ نال، تے اوہدے وچ ریڈیکس~$M$ اے (یعنی اوہ نارمل روپ وچ نہیں)، تاں \Log{tN3i} دا میل کھاندا ثبوت~$\delta$، جیہڑا $ \Gamma \Sequent N : !A$ دا اے، اوہدے وچ ذیلی !!{proof} اے (جیہڑا $ \Gamma'
\Sequent M: !B$ اُتے مُکدا اے) جیہدے اُتے اک تبدیلی لگ سکدی اے۔ جے $N \redone
N'$، $M$ نوں اوہدی تقلیل نال ملے ٹرم نال بدل کے، تاں تقلیل~$\delta$ اُتے لاؤن دا میل کھاندا حاصل !!a{proof} اے $ \Gamma \Sequent
N' : !A$ دا۔ سو ہر $\lambd$-ٹرم جیہڑا نارمل روپ وچ نہیں، $\beta$-تقلیل نال ہور ٹرم بن جاندا اے۔ ایس خاصیت نوں \emph{اگلا قدم} آکھدے آں۔

اگلا قدم تے نوع دی بقا رل کے پکا کردے نیں کہ $\beta$-تقلیل، جیہڑی $\lambd$-ٹرماں دی اے، کدے «اٹکدی» نہیں۔ جے ٹرم دی نوع ٹھیک طے اے تے اوہ نارمل روپ وچ نہیں، تاں اوہ گھٹ کے ہور ٹھیک نوع والا ٹرم بن جاندا اے (تے اوس ٹرم دی نوع اوہی اے)۔
\endgroup
% /OLP-0696
% OLP-0687 content/proof-theory/propositions-as-types/normalization.tex
\begingroup
\def\olfilename{normalization}
\expandafter\def\csname ol@sectioncs\endcsname{section}

\olfileid{int}{pty}{nor}

\olsection{نارمل روپ وچ لیاؤنا}

ایس حصے وچ اسیں ثابت کردے آں کہ تقلیل دی کسے ترتیب نال ہر اخذ نوں نارمل اخذ وچ گھٹایا جا سکدا اے۔ ایس نوں \emph{نارمل روپ وچ لیاؤن دی خاصیت} آکھدے آں۔ اسیں قضیاں تے نوعاں دا میل ورتاں گے۔ اسیں وکھاؤنے آں کہ ہر ثبوت ٹرم نوں نارمل روپ وچ گھٹایا جا سکدا اے۔ ایہدے توں قدرتی اخذ دے !!{derivation}s لئی نارمل روپ وچ لیاؤن دا نتیجہ نکلدا اے۔

پہلاں اسیں کجھ فنکشناں دی تعریف دیندے آں جیہڑے ٹرماں دی گنجل ماپدے نیں۔ \emph{لمبائی}~$\len{!A}$، !!a{formula}s دی، ایویں متعین اے:
\begin{align*}
  \len{p} &= 0\\
  \len{!A \land !B} &= \len{!A} + \len{!B} + 1\\
  \len{!A \lor !B} &= \len{!A} + \len{!B} + 1 \\
  \len{!A \lif !B} &= \len{!A} + \len{!B} + 1.
\end{align*}
ریڈیکس~$M$ دی گنجل اوہدے \emph{قطع درجے}~$\cutrank{M}$ نال ماپی جاندی اے:
\begin{align*}
  \cutrank{(\lambd[\typeof{x}{!A}][\typeof{N}{!B}])Q} &=
  \len{!A} + \len{!B} + 1 \\
  \cutrank{\ande{i}{\andi{\typeof{M}{!A}}{\typeof{N}{!B}}}} &=
  \len{!A} + \len{!B} + 1 \\
  \cutrank{\ore{\ori{i}{!A_{i}}{\typeof{M}{!A_i}}}
    {\typeof{x_1}{!A_1}}{\typeof{N_1}{!C}}
    {\typeof{x_2}{!A_2}}{\typeof{N_2}{!C}}} &=
  \len{!A} + \len{!B} + 1
\end{align*}
ثبوت ٹرم دی گنجل اوہدے وچلے سبھ توں گنجل دار ریڈیکس نال ماپی جاندی اے، تے جے ٹرم نارمل اے، تاں $0$ اے:
\[
  \maxrank{M} = \max \{\cutrank{N} | N \text{ ذیلی ٹرم اے } M \text{ دا تے ریڈیکس اے}\}
\]

\begin{lem}\ollabel{lem:subst}
  جے $\Subst{M}{\typeof{N}{!A}}{\typeof{x}{!A}}$ ریڈیکس اے تے $M
  \not\ident x$، تاں ہِٹھلیاں صورتاں وچوں اک اے:
  \begin{enumerate}
  \item $M$ آپ ریڈیکس اے، یا
  \item $M$ دی شکل $\ande{i}{x}$ اے، تے $N$ دی شکل $\andi{P_1}{P_2}$ اے۔
  \item $M$ دی شکل $\ore{i}{x_1}{P_1}{x_2}{P_2}$ اے، تے $N$ دی شکل $\ori{i}{}{Q}$ اے۔
  \item $M$ دی شکل $x Q$ اے، تے $N$ دی شکل $\lambd[x][P]$ اے۔
  \end{enumerate}
  پہلی صورت وچ $\cutrank{\Subst{M}{N}{x}} = \cutrank{M}$، تے ہورناں صورتاں وچ $\cutrank{\Subst{M}{N}{x}} = \len{!A})$۔
\end{lem}
\begin{proof}
  ثبوت~$M$ اُتے استقرا نال اے۔
  \begin{enumerate}
  \item جے $M$ اکّو متغیر $y$ اے تے $y \not\ident x$، تاں $\Subst{y}{N}{x}$، $y$ اے، ایس لئی ریڈیکس نہیں۔
  \item جے $M$ دی شکل $\andi{N_1}{N_2}$، یا $\lambd[x][N]$، یا $\ori{i}{!A}{N}$ اے، تاں $\Subst{M}{\typeof{N}{!A}}{\typeof{x}{!A}}$ وی اوسے شکل دا اے، ایس لئی ریڈیکس نہیں۔
  \item جے $M$ دی شکل $\ande{i}{P}$ اے، تاں دو صورتاں ویکھدے آں۔
    \begin{enumerate}
      \item جے $P$ دی شکل $\andi{P_1}{P_2}$ اے، تاں $M \ident
        \ande{i}{\andi{P_1}{P_2}}$ ریڈیکس اے، تے صاف اے کہ
        \[
        \Subst{M}{N}{x} \ident
        \ande{i}{\andi{\Subst{P_1}{N}{x}}{\Subst{P_2}{N}{x}}}
        \]
        وی ریڈیکس اے۔ قطع درجے برابر نیں۔
      \item جے $P$ اکّو متغیر اے، تاں تھانواں بدلّن نال ریڈیکس بنن لئی اوہ $x$ ہونا چاہیدا اے، تے $N$ دی شکل $\andi{P_1}{P_2}$ ہونی چاہیدی اے۔ ہُن ایہ ویکھو:
        \[
        \Subst{M}{N}{x} \ident \Subst{\ande{i}{x}}{\andi{P_1}{P_2}}{x},
        \]
        جیہڑا $\ande{i}{\andi{P_1}{P_2}}$ اے۔ اوہدا قطع درجہ $\cutrank{x}$ دے برابر اے، جیہڑا $\len{!A}$ اے۔
    \end{enumerate}
  \end{enumerate}
  $\ore{N}{x_1}{N_1}{x_2}{N_2}$ تے $P Q$ دیاں صورتاں ایہو جیہیاں نیں۔
\end{proof}

\begin{lem}
  جے $M$ سکڑ کے $M'$ بن جاندا اے، تے $\cutrank{M} > \cutrank{N}$، ہر ریڈیکس ذیلی ٹرم $N$ لئی جیہڑا $M$ توں وکھرا اے، تاں $\cutrank{M} >
  \maxrank{M'}$۔
\end{lem}

\begin{proof}
  وکھ وکھ صورتاں نال ثبوت۔
  \begin{enumerate}
  \item جے $M$ دی شکل $\ande{i}{\andi{M_1}{M_2}}$ اے، تاں $M'$، $M_i$ اے۔ کیونکہ $M_i$ دا ہر ذیلی ٹرم $M$ دا وی اوہدے اپنے توں وکھرا ذیلی ٹرم اے، ایس لئی دعویٰ ٹھیک اے۔
  \item جے $M$ دی شکل $(\lambd[\typeof{x}{!A}][N])\typeof{Q}{!A}$ اے، تاں $M'$، $\Subst{N}{\typeof{Q}{!A}}{\typeof{x}{!A}}$ اے۔ $M'$ وچ کسے ریڈیکس نوں ویکھو۔ یا $N$ وچ اوہدے نال میل کھاندا اوہو قطع درجے والا ریڈیکس اے، جیہڑا مفروضے دے حساب نال $\cutrank{M}$ توں گھٹ اے، یا قطع درجہ $\len{!A}$ دے برابر اے، جیہڑا تعریف دے حساب نال $\cutrank{(\lambd[x^{!A}][N])Q}$ توں گھٹ اے۔
  \item جے $M$ دی شکل
    \[
    \ore{\ori{i}{}{\typeof{N}{!A_i}}}
        {\typeof{x_1}{!A_1}}{\typeof{N_1}{!C}}
        {\typeof{x_2}{!A_2}}{\typeof{N_2}{!C}},
    \]
    اے، تاں $M' \ident \Subst{N_i}{N}{\typeof{x_i}{!A_i}}$۔~$M'$ وچ کسے ریڈیکس نوں ویکھو۔ یا $N_i$ وچ اوہدے نال میل کھاندا اوہو قطع درجے والا ریڈیکس اے، جیہڑا مفروضے دے حساب نال $\cutrank{M}$ توں گھٹ اے، یا قطع درجہ $\len{!A_i}$ دے برابر اے، جیہڑا تعریف دے حساب نال $\cutrank{\ore{\ori{i}{}{\typeof{N}{!A_i}}}
      {\typeof{x_1}{!A_1}}{\typeof{N_1}{!C}}
      {\typeof{x_2}{!A_2}}{\typeof{N_2}{!C}}}$ توں گھٹ اے۔
  \end{enumerate}
\end{proof}

\begin{thm}
  سارے ثبوت ٹرم تقلیل نال نارمل روپ وچ آ جاندے نیں۔ سارے !!{derivation}s تقلیل نال نارمل !!{derivation}s بن جاندے نیں۔
\end{thm}

\begin{proof}
  دوجی گل پہلی توں نکلدی اے۔ اسیں پہلی گل $m=\maxrank{M}$ اُتے پورے استقرا نال ثابت کردے آں، جتھے $M$ ثبوت ٹرم اے۔
  \begin{enumerate}

  \item جے $m=0$، تاں $M$ پہلاں ای نارمل اے۔
  \item نئیں تاں اسیں~$n$ اُتے استقرا کردے آں۔ ایہ $M$ وچ اوہناں ریڈیکساں دی گنتی اے جیہناں دا قطع درجہ~$m$ دے برابر اے۔
    \begin{enumerate}
    \item جے $n=1$، تاں کوئی ریڈیکس $N$ چُنو ایویں کہ $m = \cutrank{N} >
      \cutrank{P}$، ہر اوہدے اپنے توں وکھرے ذیلی ٹرم $P$ لئی جیہڑا آپ وی ریڈیکس اے۔ ایہو جیہا ریڈیکس ضرور موجود اے، کیونکہ ہر ٹرم دے ذیلی ٹرم محدود گنتی وچ نیں۔

      منّ لو کہ $N'$، $N$ دی تقلیل نال ملن والا ٹرم اے۔ ہُن لیمے نال $\maxrank{N'} < \maxrank{N}$، ایس لئی اسیں ویکھدے آں کہ $n$، یعنی $\cutrank=m$ والے ریڈیکساں دی گنتی، گھٹ گئی اے۔ سو $m$ وی گھٹ گیا اے ($1$ یا ودھ نال)، تے اسیں~$m$ لئی استقرا دا مفروضہ لا سکدے آں۔
    \item استقرائی قدم لئی منّ لو $n>1$۔ طریقہ رلدا ملدا اے، فرق ایہ اے کہ $n$ گھٹ کے وی دھناں عدد رہندا اے، سو $m$ نہیں بدلدا۔ اسیں~$n$ لئی استقرا دا مفروضہ لا دیندے آں۔
    \end{enumerate}
\end{enumerate}
\end{proof}

نوع بندی والے $\lambd$-ٹرم نارمل روپ وچ آ سکدے نیں، یعنی ہر ایہو جیہے ٹرم دا \emph{اک نارمل روپ اے}۔ جے $\beta$-تقلیل نال $\lambd$-ٹرماں نوں نارمل روپ وچ لیاؤن نوں حسابی عمل چلاؤنا تے نارمل روپ نوں اوس عمل دی قدر سمجھیے، تاں ایہدا مطلب اے کہ نوع بندی والے $\lambd$-کیلکولس وچ ہر حسابی عمل آخر رُک جاندا اے تے اک قدر دیندا اے۔ نالے نوع دی بقا پکی کردی اے کہ قدر دی نوع ٹھیک اے۔ مثال لئی جے $N$ دی نوع $!A \lif !B$ اے (یعنی اوہ ایہو جیہا فنکشن اے جیہڑا~$!A$ نوع دے آرگومنٹاں اُتے کم کرے تے~$!B$ نوع دیاں قدراں دیوے)، تے اسیں اوہنوں~$M$ ٹرم اُتے لائیے جیہدی نوع~$!A$ اے (ان پُٹ)، تاں $(NM)$ گھٹ کے ٹرم $N'$ بن جاندا اے (آؤٹ پُٹ)، تے پکا اے کہ اوس آؤٹ پُٹ دی نوع~$!B$ اے۔

اصل وچ نوع بندی والے $\lambd$-کیلکولس دے ٹرم صرف اک طریقے نال ذیلی ٹرماں اُتے $\beta$-تقلیل لا کے نارمل روپ وچ نہیں آؤندے۔ ذیلی ٹرماں اُتے $\beta$-تقلیل لاؤن دا \emph{ہر} طریقہ آخر نارمل روپ اُتے مُکدا اے۔ ایس خاصیت نوں \emph{مضبوط نارملائزیشن} آکھدے آں۔ ایہدا مطلب اے کہ حسابی عمل نوں آخر روک کے قدر حاصل کرن دی ضمانت والا طریقہ موجود اے، تے حسابی عمل چلاؤن دا \emph{ہر} طریقہ آخر قدر دیندا اے۔

سانوں کیویں پتا اے کہ حسابی عمل چلاؤن دے وکھ وکھ طریقے \emph{اوہی} قدر دیندے نیں؟ ہُن تک سانوں بس ایہ پتا اے (نوع دی بقا توں) کہ حاصل ہوئیاں ساریاں قدراں اکّو نوع دیاں نیں۔ نوع بندی والے $\lambd$-کیلکولس دی ہور اہم خاصیت اے۔ ایہ \emph{سنگم والی} یا \emph{چرچ–راسر} خاصیت رکھدا اے۔ جے $N \red N_1$ تے $N \red
N_2$، تاں کوئی ٹرم~$N'$ اے ایویں کہ $N_1 \red N'$ تے $N_2 \red
N'$۔ چیتے رکھو کہ $N \red N'$ دا مطلب ایہ اے کہ $N'$ صفر یا ودھ $\redone$ تقلیل قدماں نال ملدا اے۔ جے $N_1$ نارمل روپ وچ اے، تاں اکّو ٹرم $N'$ جیہدے لئی $N_1 \red N'$ اے، $N_1$ آپ اے (صفر $\redone$ تقلیل قدماں نال)۔ ایس توں نکلدا اے کہ جے $N_1$ تے $N_2$ دونے نارمل روپ وچ نیں، تاں $N_1 = N' = N_2$۔ دوجے لفظاں وچ، کیونکہ $\red$ مضبوط نارملائزیشن والی اے، ٹرم دی تقلیل دا ہر طریقہ آخر نارمل روپ دیندا اے۔ کیونکہ $\red$ چرچ–راسر اے، کوئی وی دو نارمل روپ برابر نیں۔ ایس لئی نوع بندی والے $\lambd$-کیلکولس وچ حسابی عمل ہمیشہ مُک جاندا اے، تے اوہدی قدر (نارمل روپ) اوہی رہندی اے بھاویں حسابی عمل کِسے وی طرح چلے۔

ایہ ثابت کرنا عام طور اُتے اوکھا اے کہ کسے تعلق دی سنگم والی خاصیت اے۔ پر اسیں ہِٹھلی کمزور شرط سوکھے ای ثابت کر سکدے آں۔

\begin{prop}[چرچ–راسر دی کمزور خاصیت]
جے $N \redone N_1$ تے $N \redone N_2$، تاں اک ٹرم $N'$ اے ایویں کہ $N_1 \red N'$ تے $N_2 \red N'$۔
\end{prop}

\begin{proof}
  دونے ٹرم $N_1$ تے $N_2$، $N$ توں ملدے نیں، بالترتیب ریڈیکس $M_1$ یا $M_2$ دی تھاں اوہدی تقلیل نال ملن والا ٹرم رکھ کے۔ ریڈیکس $M_1$ تے $M_2$، $N$ دے ذیلی ٹرم نیں، سو اوہ اک دوجے اُتے نہیں چڑھ سکدے۔ ایہدا مطلب اے کہ یا اک دوجے دا ذیلی ٹرم اے، یا اوہ $N$ وچ دو وکھریاں، سانجھ توں بناں تھاواں اُتے نیں۔ دوجی صورت ویکھو۔ $N$ دی شکل $N(M_1,M_2)$ اے۔ $M_1$ دی تھاں اوہدی تقلیل نال ملن والا $M_1'$ رکھن نال $N_1$ ملدا اے، جیہڑا $N(M_1',M_2)$ اے۔ تے $M_2$ دی تھاں اوہدی تقلیل نال ملن والا $M_2'$ رکھن نال $N_2$ ملدا اے، جیہڑا $N(M_1,M_2')$ اے۔ دونے گھٹ کے $N(M_1',M_2')$ بن جاندے نیں۔ ہور صورت لئی منّ لو $M_2$،~$M_1$ دا ذیلی ٹرم اے (دوجی صورت ایہدے نال سم روپی اے)، تے $M_1$ کِس طرح دا ریڈیکس اے، اوس دے حساب نال صورتاں وکھ کرو۔ مثال لئی جے $M_1 \ident
  \proj{1}{\pair{O_1}{O_2}}$، تاں اوہ گھٹ کے $O_1$ بن جاندا اے۔ جے $M_2$، $O_1$ دا ذیلی ٹرم اے، تاں $M_2$ دی تبدیلی نال $O_1'$ ملدا اے۔ سو پہلاں $M_1$ تے فیر $M_2$ بدلّن نال $O_1'$ ملدا اے۔ پر پہلاں $M_2$ بدلّن نال $\proj{1}{\pair{O_1'}{O_2}}$ ملدا اے تے ایہ وی گھٹ کے~$O_2$ بن جاندا اے۔
\end{proof}

\begin{lem}[نیومین دا لیما]
جے $\red$ مضبوط نارملائزیشن والی اے تے چرچ–راسر دی کمزور خاصیت رکھدی اے، تاں اوہ سنگم والی اے۔
\end{lem}

\begin{proof}
  کیونکہ $\red$ مضبوط نارملائزیشن والی اے، ہر تقلیل لڑی اک نارمل روپ دیندی اے۔ نارمل روپ اوہدوں تے صرف اوہدوں اکّو ہوندا اے جدوں $\red$ سنگم والی اے۔ سو منّ لو کہ کسے ٹرم $N$ دے دو وکھرے نارمل روپ $N'$ تے $N''$ نیں، جیہڑے تقلیل لڑیاں توں ملے نیں:
  \begin{align*}
  N & = N_1' \redone N_2' \redone \dots \redone N_k' = N' \text{ تے}\\
  N & = N_1'' \redone N_2'' \redone \dots \redone N_l'' = N''
  \end{align*}
  (تقلیل لڑی دی لمبائی $>0$ ہونی چاہیدی اے، نئیں تاں $N$ پہلاں ای نارمل روپ وچ ہووے گا۔)

  جے $N_1' = N_1''$، تاں ایہدے دو وکھرے نارمل روپ نیں ($N'$ تے $N''$)۔ جے $N_1' \neq N_1''$، تاں چرچ–راسر دی کمزور خاصیت نال کوئی $N'''$ اے ایویں کہ $N_1' \red N'''$ تے $N_1'' \red N'''$۔ کیونکہ $N' \neq N''$، تاں یا $N''' \neq N'$ یا $N''' \neq N''$۔
  ایس لئی $N_1'$ یا $N_1''$ وچوں کسے اک دے دو وکھرے نارمل روپ نیں۔ اسیں وکھایا اے کہ جے $N$ دے دو وکھرے نارمل روپ نیں، تاں اک $N_1$ اے ایویں کہ $N \redone N_1$ تے $N_1$ دے دو وکھرے نارمل روپ نیں۔ صاف اے کہ اسیں ایویں اگے ودھ سکدے آں تے $N \redone N_1
  \redone N_2 \redone \dots$ بنا سکدے آں، پر ایہ ممکن نہیں، کیونکہ $\redone$ مضبوط نارملائزیشن والی اے۔
\end{proof}
\endgroup
% /OLP-0687

% OLP-0694 content/proof-theory/propositions-as-types/sequent-natural-deduction.tex
\begingroup
\def\olfilename{sequent-natural-deduction}
\expandafter\def\csname ol@sectioncs\endcsname{section}

\olfileid{int}{pty}{snd}

\olsection{سیکوئنٹ طرز دا قدرتی اخذ}

$ \Gamma \Sequent !A$ لکھو جے قدرتی اخذ دا !!{derivation} اے جیہدے وچ $ \Gamma$، !!{undischarged} مفروضے نیں تے $!A$ نتیجہ اے۔ یا $ \Sequent !A$ لکھو جے $ \Gamma$ خالی اے۔

$ \Gamma, !A_1, \dots, !A_n$ اسیں $ \Gamma \cup \{!A_1, \dots,
!A_n\}$ لئی لکھدے آں، تے $ \Gamma, \Delta$ نوں $ \Gamma \cup \Delta$ لئی۔

ویکھو کہ جدوں ساڈے کول $ \Gamma \Sequent !A \land !A$ اے، یعنی ساڈے کول !!a{derivation} اے جیہدے وچ $ \Gamma$، !!{undischarged} مفروضے نیں تے $!A \land !A$ آخری !!{formula} اے، تاں تھلّے $ \Elim{\land}$ لا کے اسیں !!a{derivation} لے سکدے آں اوہناں ای !!{undischarged} مفروضیاں تے نتیجے $!A$ والا۔ دوجے لفظاں وچ، جے $ \Gamma
\Sequent !A \land !B$، تاں $ \Gamma \Sequent !A$۔
\begin{prooftree}
  \Axiom$\Gamma \fCenter !A \land !B$
  \RightLabel{$\Elim{\land}$}
  \UnaryInf$\Gamma \fCenter !A$
  \DisplayProof\qquad\bottomAlignProof
  \Axiom$\Gamma \fCenter !A \land !B$
  \RightLabel{$\Elim{\land}$}
  \UnaryInf$\Gamma \fCenter !B$
\end{prooftree}
لیبل $ \Elim{\land}$، قدرتی اخذ وچ اوسے ناں والے قاعدے نال جوڑ ول اشارہ کردا اے۔

ایویں ای منّ لو کہ ساڈے کول $ \Gamma, !A \Sequent !B$ اے، یعنی ساڈے کول !!a{derivation} اے جیہدے !!{undischarged} مفروضے $ \Gamma, !A$ تے آخری !!{formula}~$!B$ اے۔ جے اسیں $ \Intro{\lif}$ قاعدہ لائیے، تاں ساڈے کول !!a{derivation} اے جیہدے وچ $ \Gamma$، !!{undischarged} مفروضے نیں تے $!A
\lif !B$ آخری !!{formula} اے، یعنی $ \Gamma \Sequent !A \lif
!B$۔ ویکھو کہ ایس نال مفروضیاں دا !!{discharge} ہور صاف ہو گیا اے۔
\begin{prooftree}
  \Axiom $\Gamma, !A \fCenter !B$
  \RightLabel{$\Intro{\lif}$}
  \UnaryInf $\Gamma \fCenter !A \lif !B$
\end{prooftree}

ہور قاعدیاں توں وی اسیں ایس طرح نتیجے کڈھ سکدے آں۔ ایہ قاعدے ہِٹھ کھولھ کے دِتّے نیں:
\begin{gather*}
  \Axiom $\Gamma \fCenter !A$
  \Axiom $\Delta \fCenter !B$
  \RightLabel{$\Intro{\land}$}
  \BinaryInf $\Gamma,\Delta \fCenter !A \land !B$
  \DisplayProof\\
  \Axiom $\Gamma \fCenter !A \land !B$
  \RightLabel{$\Elim{\land}_1$}
  \UnaryInf $\Gamma \fCenter !A$
  \DisplayProof
  \qquad
  \Axiom $\Gamma \fCenter !A \land !B$
  \RightLabel{$\Elim{\land}_2$}
  \UnaryInf $\Gamma \fCenter !B$
  \DisplayProof
  \\
  \Axiom $\Gamma \fCenter !A$
  \RightLabel{$\Intro{\lor}_1$}
  \UnaryInf $\Gamma \fCenter !A \lor !B$
  \DisplayProof\qquad
  \Axiom $\Gamma \fCenter !B$
  \RightLabel{$\Intro{\lor}_2$}
  \UnaryInf $\Gamma \fCenter !A \lor !B$
  \DisplayProof\\
  \Axiom $\Gamma \fCenter !A \lor !B$
  \Axiom $\Delta, !A \fCenter !C$
  \Axiom $\Delta', !B \fCenter !C$
  \RightLabel{$\Elim{\lor}$}
  \TrinaryInf $\Gamma, \Delta, \Delta' \fCenter !C$
  \DisplayProof
  \\
  \Axiom $\Gamma, !A \fCenter !B$
  \RightLabel{$\Intro{\lif}$}
  \UnaryInf $\Gamma \fCenter !A \lif !B$
  \DisplayProof
  \qquad
  \Axiom $\Delta \fCenter !A \lif !B$
  \Axiom $\Gamma \fCenter !A$
  \RightLabel{$\Elim{\lif}$}
  \BinaryInf $\Gamma, \Delta \fCenter !B$
  \DisplayProof
  \\
  \Axiom $\Gamma \fCenter \lfalse$
  \RightLabel{$\FalseInt$}
  \UnaryInf $\Gamma \fCenter !A$
  \DisplayProof
\end{gather*}

ہر مفروضہ اپنے آپ وچ !!a{derivation} اے $!A$ دا~$!A$ توں، یعنی ساڈے کول ہمیشہ $!A \Sequent !A$ اے۔

\begin{prooftree}
  \AxiomC{$ $}
  \UnaryInfC{$!A \fCenter !A$}
\end{prooftree}

ایہ سارے قاعدے رلا کے قدرتی اخذ دے !!{derivation}s وچوں کِہڑے موجود نیں، ایہدے بارے اک حساب منّے جا سکدے نیں۔ ایہ قدرتی اخذ دی نشانیاں والی وکھری صورت وی منّے جا سکدے نیں، جیہدے وچ ہر قدم اُتے !!{formula} جیہڑا !!{derive} ہویا، اوہ وی درج ہوندا اے تے !!{undischarged} مفروضے جیہناں توں اوہ !!{derive} ہویا، اوہ وی۔

\begin{prooftree}
  \Axiom$!A \fCenter !A$
  \UnaryInf $!A \fCenter !A \lor (!A \lif \lfalse)$
  \Axiom $!B \fCenter !B$
  \BinaryInf$!A, !B\lif \fCenter \lfalse$
  \UnaryInf$(!B \fCenter !A \lif \lfalse$
  \UnaryInf$( !B \fCenter !A \lor (!A \lif \lfalse)$
  \Axiom$(!B \fCenter !B$
  \BinaryInf$(!B \fCenter \lfalse$
  \UnaryInf$\fCenter !B \lif \lfalse$
\end{prooftree}
جتھے $!B$، $(!A \lor (!A \lif \lfalse))\lif \lfalse$ دا چھوٹا روپ اے۔
\endgroup
% /OLP-0694
\OLEndChapterHook
\endgroup
% /OLP-0690

\OLEndPartHook
\endgroup
% /OLP-0685

% OLP-0516 content/counterfactuals/counterfactuals.tex
\begingroup
\def\olfilename{counterfactuals}
\expandafter\def\csname ol@sectioncs\endcsname{section}

\olpart{cnt}{خلافِ\-واقعہ شرطیے}

% OLP-0517 content/counterfactuals/introduction/introduction.tex
\begingroup
\def\olfilename{introduction}
\expandafter\def\csname ol@sectioncs\endcsname{section}

\olchapter{cnt}{int}{جان پہچان}

% OLP-0518 content/counterfactuals/introduction/material-conditional.tex
\begingroup
\def\olfilename{material-conditional}
\expandafter\def\csname ol@sectioncs\endcsname{section}

\olfileid{cnt}{int}{mat}

\olsection{مادی شرطیہ}

انگریزی وچ اپنی سبھ توں سادی شکل وچ شرطیہ ایہو جیہا جملہ اے:
’جے \dots تاں \dots‘، جتھے \dots{} آپ جملے ہوندے نیں۔ مثلاً
’جے ایہ کم گھر دے وڈے نوکر نے کیتا اے، تاں مالی بے قصور اے۔‘ منطق
دے ابتدائی کورساں وچ اسیں $ \lif$ رابطہ ورت کے شرطیے علامتاں وچ لکھنے
سکھدے آں۔ \dots نال دسّے حصیاں نوں، مثلاً !!{formula}s $!A$ تے~$!B$ نال
لکھو، تے پورے شرطیے نوں $!A \lif !B$ نال لکھدے آں۔

رابطہ $\lif$ \emph{صداقتی-فنکشنی} اے، یعنی سچائی دی قدر، $\True$ یا $\False$،
$!A \lif !B$ لئی $!A$ تے~$!B$ دیاں سچائی قدراں نال طے ہوندی اے۔
$!A \lif !B$ اُدوں تے صرف اُدوں سچ اے جدوں $!A$ جھوٹ یا $!B$ سچ اے،
تے ہور صورت وچ جھوٹ اے۔ سچائی دی قدر بندی~$\pAssign v$ دے لحاظ نال
اسیں ایویں تعریف کردے آں کہ $\pSat{v}{!A \lif !B}$ اُدوں تے صرف اُدوں اے جدوں
$\pSat/{v}{!A}$ یا $\pSat{v}{!B}$۔ ایس معنیات والے رابطے $\lif$ نوں
\emph{مادی شرطیہ} آکھدے آں۔

ایس تعریف توں منطق دیاں کئی بنیادی گلاں نکلدیاں نیں۔ سبھ توں پہلاں، استنباط دا
قضیہ مادی شرطیے لئی قائم اے:
\begin{align}
  \text{ جے } \Gamma, !A \Entails !B \text{ تاں } \Gamma \Entails !A \lif !B
\end{align}
ایہ صداقتی-فنکشنی اے۔ $!A \lif !B$ تے $\lnot !A \lor !B$ ہم ارز نیں:
\begin{align}
  \formula A \lif !B & \Entails \lnot !A \lor !B\\
  \lnot !A \lor !B & \Entails !A \lif !B
  \intertext{مادی شرطیہ اپنی تالی دا وی منطقی نتیجہ اے تے اپنے مقدم دی نفی دا وی:}
  \formula B & \Entails !A \lif !B\\
  \lnot !A & \Entails !A \lif !B
  \intertext{جھوٹا مادی شرطیہ اپنے مقدم تے اپنی تالی دی نفی دے عطف دے ہم ارز اے۔ جے $!A \lif !B$ جھوٹ اے، تاں $!A \land \lnot !B$ سچ اے، تے ایہدا اُلٹ وی:}
  \lnot(!A \lif !B) & \Entails !A \land \lnot !B\\
  \formula A \land \lnot !B & \Entails \lnot(!A \lif !B)
  \intertext{مادی شرطیے لئی اثباتِ مقدم دا قاعدہ چلدا اے:}
  \formula A, !A \lif !B & \Entails !B
  \intertext{مادی شرطیاں دیاں تالیاں اکٹھیاں ہو سکدیاں نیں:}
  \formula A \lif !B,  !A \lif !C & \Entails !A \lif (!B \land !C)
  \intertext{اسیں مقدم نوں ہمیشہ ہور مضبوط کر سکدے آں، یعنی شرطیہ \emph{یک رُخا} اے:}
  \formula A \lif !B & \Entails (!A \land !C) \lif !B
  \intertext{مادی شرطیہ متعدی اے، یعنی لڑی والا قاعدہ معتبر اے:}
  \formula A \lif !B, !B \lif !C & \Entails !A \lif !C
  \intertext{مادی شرطیہ اپنے عکسِ نقیض دے ہم ارز اے:}
  \formula A \lif !B & \Entails \lnot !B \lif \lnot !A\\
  \lnot !B \lif \lnot !A & \Entails !A \lif !B
\end{align}

ایہ سارے استنباط ریاضیاتی دلیل وچ کم آؤندے نیں تے کوئی اوکڑ نہیں پاندے۔
پر فلسفے تے لسانیات دیاں لکھتاں وچ غیر ریاضیاتی موقعیاں اُتے اوہناں دے برابر
استنباطاں دیاں کئی مبینہ ردّی مثالاں دتیاں گئیاں نیں۔ اوہ ایہ اشارہ دیندیاں نیں
کہ انگریزی دے ’جے \dots تاں \dots‘ والے بیاناں نوں علامتاں وچ لکھن لئی
مادی شرطیہ~$\lif$ مناسب رابطہ نہیں، یا گھٹ توں گھٹ ہمیشہ مناسب نہیں۔
\endgroup
% /OLP-0518

% OLP-0519 content/counterfactuals/introduction/paradoxes-material.tex
\begingroup
\def\olfilename{paradoxes-material}
\expandafter\def\csname ol@sectioncs\endcsname{section}

\olfileid{cnt}{int}{par}

\olsection{مادی شرطیے دے منطقی معمّے}

انگریزی دے ’جے \dots تاں \dots‘ والے بیاناں نوں $!A \lif !B$ نال لکھن دے
طریقے اُتے سبھ توں پہلاں اعتراض کرن والیاں وچ سی۔ آئی۔ لوئس سی۔
لوئس دا اعتراض وائٹ ہیڈ تے رسل دی \emph{Principia Mathematica} وچ مادی
شرطیے دی ورتوں اُتے سی۔ اوہ $\lif$ نوں ’لازم کردا اے‘ پڑھدے سن۔
لوئس نے ٹھیک اعتراض کیتا کہ جے $\lif$ دا مطلب ’لازم کردا اے‘ ہووے،
تاں ہر جھوٹا قضیہ~$p$ ایہ لازم کرے گا کہ $p$ توں~$q$ لازم آندا اے،
کیوں جے $p \lif (p \lif q)$ اُدوں سچ اے جدوں~$p$ جھوٹ اے۔ تے ہر سچا
قضیہ~$q$ ایہ لازم کرے گا کہ $p$ توں~$q$ لازم آندا اے، کیوں جے
$q \lif (p \lif q)$ اُدوں سچ اے جدوں $q$ سچ اے۔

منطق دے ماہر جان دے نیں کہ \emph{استلزام}، یعنی منطقی نتیجے دا تعلق،
!!{formula}s یا بیاناں وچکار اک تعلق اے، منطقی رابطہ نہیں۔ سو رَلا نہ
پئے، ایس لئی $\lif$ نوں ’لازم کردا اے‘ نہ پڑھیے۔\footnote{ریاضی دان تے
  کمپیوٹر سائنس دان ہجے وی ’$\lif$‘ نوں ’لازم کردا اے‘ عام پڑھدے نیں۔
  فلسفی لوئس دے دسے رَلے توں بچن لئی اینوں ’صرف جے‘ پڑھدے نیں۔}
جِنّی دیر اسیں ایویں نہیں پڑھدے، لوئس دی خاص اوکڑ پیدا ہی نہیں ہوندی۔
$p$، $q$ نوں ’لازم‘ نہیں کردا، بھانویں $p$ اک جھوٹے انگریزی جملے دا نشان
ہووے۔ $p \Entails q$ اے یا نہیں، ایہ طے کرن لئی \emph{ساریاں}
!!{valuation}s ویکھنیاں پیندیاں نیں۔ تے $p \Entails/ q$ اوہدوں وی اے
جدوں اسیں $p$ نوں ایسے جملے دا نشان بنائیے جیہڑا اتفاق نال جھوٹا اے۔

پر لوئس دی ہِیٹھلی مثال ورگے ’جے \dots تاں\dots‘ بیاناں وچ ہجے وی
اک عجیب گل اے:
\begin{quote}
جے چن کچّے پنیر دا بنیا اے، تاں $2+2=4$۔
\end{quote}
ایسی گل ہِیٹھلے استنباطاں وچ وی اے:
\begin{quote}
  چن کچّے پنیر دا نہیں بنیا۔ سو جے چن کچّے پنیر دا بنیا اے، تاں $2+2=4$۔

  $2+2 = 4$۔ سو جے چن کچّے پنیر دا بنیا اے، تاں $2+2=4$۔
\end{quote}
پر جے ’جے \dots تاں \dots‘ دا مطلب صرف $\lif$ ہووے، تاں ایہ جملہ
بغیر کسے اوکڑ دے سچ، تے ایہ استنباط بغیر کسے اوکڑ دے معتبر ہون گے۔

اک ہور مثال ہمیشہ سچے فارمولے $(!A \lif !B) \lor (!B \lif
!A)$ نال جڑی اے۔ ایہدے حساب نال جے اخبار وچوں دو خبری جملے~$S$ تے $T$
اٹکل نال چن لئیے، تاں ’جے $S$، تاں $T$، یا جے $T$، تاں~$S$‘ سچ
ہونا چاہیدا اے۔
\endgroup
% /OLP-0519

% OLP-0520 content/counterfactuals/introduction/strict-conditional.tex
\begingroup
\def\olfilename{strict-conditional}
\expandafter\def\csname ol@sectioncs\endcsname{section}

\olfileid{cnt}{int}{str}

\olsection{سخت شرطیہ}

لوئس نے \emph{سخت شرطیہ}~$\strictif$ پیش کیتا تے دلیل دتی کہ استلزام
ایہدے نال میل کھاندا اے، مادی شرطیے نال نہیں۔ سچائی دی امکانی منطق وچ
$!A \strictif !B$ نوں $\Box(!A \lif !B)$ دے طور اُتے تعریف کر سکدے آں۔
سو سخت شرطیہ (کسے جہان اُتے) اُدوں تے صرف اُدوں سچ اے جدوں اوہدے نال
ملدا مادی شرطیہ لازمی اے۔

مادی شرطیے دے معمّیاں دے لحاظ نال سخت شرطیہ کیویں چلدا اے؟ جھوٹے مقدم
والا سخت شرطیہ تے سچی تالی والا سخت شرطیہ، دونیں سچ وی ہو سکدے نیں
تے جھوٹ وی۔ نالے $(!A \strictif !B) \lor (!B \strictif !A)$ معتبر نہیں۔
سخت شرطیہ $!A \strictif !B$، $\lnot !A \lor !B$ دے برابر وی نہیں،
سو ایہ صداقتی-فنکشنی نہیں۔

ساڈے کول ایہ نیں:
\begin{align}
  \formula A \strictif !B & \Entails \lnot !A \lor !B
  \text{ پر }\\
  \lnot !A \lor !B & \Entails/ !A \strictif !B\\
  \formula B & \Entails/ !A \strictif !B\\
  \lnot !A & \Entails/ !A \strictif !B\\
  \lnot(!A \lif !B) & \Entails/ !A \land \lnot !B
  \text{ پر }\\
  \formula A \land \lnot !B & \Entails \lnot(!A \strictif !B)
  \intertext{فیر وی سخت شرطیے لئی اثباتِ مقدم دا قاعدہ چلدا اے:}
  \formula A, !A \strictif !B & \Entails !B
  \intertext{سخت شرطیاں دیاں تالیاں اکٹھیاں ہو سکدیاں نیں:}
  \formula A \strictif !B,  !A \strictif !C & \Entails !A \strictif (!B \land !C)
  \intertext{سخت شرطیے لئی مقدم نوں مضبوط کرن دا قاعدہ چلدا اے:}
  \formula A \strictif !B & \Entails (!A \land !C) \strictif !B
  \intertext{سخت شرطیہ متعدی وی اے:}
  \formula A \strictif !B, !B \strictif !C & \Entails !A \strictif !C
  \intertext{آخر وچ، سخت شرطیہ اپنے عکسِ نقیض دے ہم ارز اے:}
  \formula A \strictif !B & \Entails \lnot !B \strictif \lnot !A\\
  \lnot !B \strictif \lnot !A & \Entails !A \strictif !B
\end{align}

\begin{prob}
  استلزام دے اوہناں تعلقاں لئی \Log{S5} وچ ردّی مثالاں دیو جیہڑے سخت
  شرطیے لئی قائم نہیں ہوندے، یعنی ہِیٹھلیاں لئی:
  \begin{enumerate}
  \item $\lnot p \Entails/ \Box(p \lif q)$
  \item $q \Entails/ \Box(p \lif q)$
  \item $\lnot\Box(p \lif q) \Entails/ p \land \lnot q$
  \item $\Entails/ \Box(p \lif q) \lor \Box(q \lif p)$
  \end{enumerate}
\end{prob}

\begin{prob}
  ہِیٹھلیاں گلاں دے \Log{S5} وچ ثبوت دے کے وکھاؤ کہ معتبر استلزام دے
  تعلق سخت شرطیے لئی قائم نیں:
  \begin{enumerate}
  \item  $\Box(!A \lif !B) \Entails \lnot !A \lor !B$
  \item $!A \land \lnot !B \Entails \lnot\Box(!A \lif !B)$
  \item $!A, \Box(!A \lif !B) \Entails !B$
  \item $\Box(!A \lif !B), \Box(!A \lif !C) \Entails \Box(!A \lif (!B
    \land !C))$
  \item $\Box(!A \lif !B) \Entails \Box((!A \land !C) \lif !B)$
  \item $\Box(!A \lif !B), \Box(!B \lif !C) \Entails \Box(!A
    \lif !C)$
  \item $\Box(!A \lif !B) \Entails \Box(\lnot !B \lif \lnot !A)$
  \item $\Box(\lnot !B \lif \lnot !A) \Entails \Box(!A \lif !B)$
  \end{enumerate}
\end{prob}

پر سخت شرطیے دے اپنے وی ’معمّے‘ نیں۔ جیویں جھوٹے مقدم یا سچی تالی
والا مادی شرطیہ سچ ہوندا اے، اویویں \emph{لازمی طور اُتے} جھوٹے مقدم یا
لازمی طور اُتے سچی تالی والا سخت شرطیہ سچ ہوندا اے۔ نالے ہر سچا سخت شرطیہ
لازمی طور اُتے سچ اے، تے ہر جھوٹا سخت شرطیہ لازمی طور اُتے جھوٹ۔
دوجے لفظاں وچ،
\begin{align}
  \Box \lnot !A & \Entails !A \strictif !B\\
  \Box !B & \Entails !A \strictif !B\\
  \formula A \strictif !B& \Entails \Box(!A \strictif !B)\\
  \lnot(!A \strictif !B) & \Entails \Box\lnot(!A \strictif !B)
\end{align}
جے $\strictif$ دا مطلب ’لازم کردا اے‘ لئیے، تاں ایہ اوکڑاں نہیں۔
آخر منطقی استلزام دے تعلق ریاضیاتی حقیقتاں نیں، ایس لئی اوہ اتفاقی نہیں ہو سکدے۔
پر جے اسیں $\strictif$ نوں ’جے \dots تاں \dots‘ دا مطلب دین والا منطقی
رابطہ بنائیے، تاں ایہناں نال اوکڑ پَندی اے، خاص کر کے آخری دوناں نال۔
کیوں جے ’جے \dots تاں \dots‘ والے کجھ بیان اتفاقی طور اُتے سچ یا اتفاقی
طور اُتے جھوٹ ہوندے نیں۔ اصل وچ عام طور اُتے اوہ نہ لازمی ہوندے نیں،
نہ ناممکن۔

\begin{prob}
  ہِیٹھلیاں گلاں دے \Log{S5} وچ ثبوت دیو:
  \begin{enumerate}
    \item  $\Box \lnot !A \Entails !A \strictif !B$
    \item $!A \strictif !B \Entails \Box(!A \strictif !B)$
    \item $\lnot(!A \strictif !B)  \Entails \Box\lnot(!A \strictif !B)$
  \end{enumerate}
  ایس لئی $\strictif$ دی تعریف ورتو۔
\end{prob}
\endgroup
% /OLP-0520

% OLP-0521 content/counterfactuals/introduction/counterfactuals.tex
\begingroup
\def\olfilename{counterfactuals}
\expandafter\def\csname ol@sectioncs\endcsname{section}

\olfileid{cnt}{int}{cnt}

\olsection{خلافِ واقعہ شرطیے}

انگریزی وچ ’جے \dots تاں \dots‘ دی اک بڑی عام تے ضروری شکل
\emph{to be} دے ماضی والے فرضیہ صیغے \pnbLatin{(past subjunctive)} نال بندی اے۔
ایہدی شکل اے ’جے ایہ گل ہوندی کہ \dots تاں ایہ گل ہوندی کہ \dots‘۔
ایہو جیہے شرطیے دا مقدم عام طور اُتے جھوٹا ہوندا اے، یعنی حقیقت دے اُلٹ۔
ایس لئی ایہناں نوں \emph{خلافِ واقعہ شرطیے} آکھدے نیں۔ \emph{to be}
دا فرضیہ صیغہ ورتن کارن ایہناں نوں \emph{فرضیہ شرطیے} وی آکھدے نیں۔
ایہ \emph{خبری} شرطیاں توں وکھرے نیں، جیہناں دی شکل اے ’جے ایہ گل اے
کہ \dots تاں ایہ گل اے کہ \dots‘۔ خلافِ واقعہ تے خبری شرطیاں دے سچ
ہون دیاں شرطاں وکھریاں نیں۔ ایڈمز دی مشہور مثال ویکھو:
\begin{quote}
  جے اوسوالڈ نے کینیڈی نوں نہیں ماریا، تاں کسے ہور نے ماریا اے۔

  جے اوسوالڈ نے کینیڈی نوں نہ ماریا ہوندا، تاں کسے ہور نے مار دینا سی۔
\end{quote}
پہلا خبری اے، دوجا خلافِ واقعہ۔ پہلا صاف طور اُتے سچ اے۔ سانوں پتا اے
کہ صدر جان ایف کینیڈی نوں \emph{کسے} نے ماریا سی، تے جے (وارن رپورٹ دے اُلٹ)
اوہ بندہ لی ہاروی اوسوالڈ نہیں سی، تاں کینیڈی نوں کسے ہور نے ماریا سی۔
دوجا کجھ ہور آکھدا اے۔ ایہ دعویٰ کردا اے کہ جے اوسوالڈ نے کینیڈی نوں نہ
ماریا ہوندا، یعنی جے ڈیلس دی گولی باری ٹل گئی ہوندی یا ناکام رہی ہوندی،
تاں اگوں تاریخ ایویں ودھدی کہ قتل دی اک ہور کوشش کامیاب ہو جاندی۔
ایہدے سچ ہون لئی ضروری اے کہ طاقتور دھڑیاں نے کینیڈی دی موت یقینی بناؤن
لئی سازش کیتی ہوندی (جیویں کینیڈی دے قتل دی سازش منّن والے کئی لوک
سمجھدے نیں)۔

ایہ بحث ہجے جاری اے کہ مادی شرطیہ \emph{خبری} شرطیے دا مطلب ٹھیک طرح
دیندا اے یا نہیں۔ خاص کر کے بحث ایہ اے کہ مادی شرطیے دے معمّیاں دی
’سمجھاؤنی‘ ایویں ہو سکدی اے یا نہیں کہ اوہ انگریزی دے خبری شرطیاں دے
سچ ہون دیاں شرطاں وی ٹھیک دیندا ہووے۔ البتہ ایس گل اُتے اختلاف نہیں کہ
مادی شرطیاں نال خلافِ واقعہ شرطیے ٹھیک علامتاں وچ نہیں لکھے جا سکدے۔
ایہ گل ایس توں صاف اے کہ بھانویں خلافِ واقعہ شرطیاں دے مقدم عام طور اُتے
جھوٹے نیں، جھوٹے مقدم والے سارے خلافِ واقعہ شرطیے سچ نہیں۔ مثلاً جے تسی
وارن رپورٹ منّ دے او تے کینیڈی دے قتل دی کوئی سازش نہیں سی، تاں ایڈمز دا
خلافِ واقعہ شرطیہ ایہدی مثال اے۔

سبب بارے دلیل وچ خلافِ واقعہ شرطیاں دا وڈا کم اے۔ ایہناں دی اک
اہم ورتوں سبب تے اثر دا تعلق بیان کرنا اے۔ مثلاً تیلی رگڑن نال اوہ بل پَیندی اے،
تے ایہ گل ایویں آکھ سکدے آں، ’جے ایہ تیلی رگڑی جاندی، تاں بل پَیندی۔‘
مادی، تے عام طور اُتے خبری، شرطیاں نال ایہ گل نہیں آکھی جا سکدی۔
’تیلی رگڑی جاندی اے $\lif$ تیلی بلدی اے‘ سچ اے جے تیلی کدی رگڑی
ای نہ جاوے، بھانویں رگڑن اُتے جو مرضی ہوندا۔ ہور وی ماڑی گل، ’تیلی رگڑی
جاندی اے $\lif$ تیلی پھُلّاں دے گلدستے وچ بدل جاندی اے‘ وی سچ اے جے
تیلی کدی رگڑی ای نہ جاوے۔ پر جے اوہ رگڑی جاندی، تاں اوہ یقیناً پھُلّاں
دا گلدستہ نہ بندی۔

خلافِ واقعہ شرطیاں دی صحیح منطق کیہڑی اے، ایس اُتے ہجے بحث اے۔
سٹالنیکر تے لوئس نے ایہناں شرطیاں دی اک بااثر تشریح پیش کیتی۔
اوہناں مطابق خلافِ واقعہ شرطیہ ’جے ایہ گل ہوندی کہ~$S$، تاں ایہ گل
ہوندی کہ~$T$‘ اُدوں تے صرف اُدوں سچ اے جدوں $T$ اوس خلافِ واقعہ حالت
(’ممکن جہان‘) وچ سچ ہووے جیہڑی اصل جہان دی حالت دے سبھ توں نیڑے اے
تے جتھے~$S$ سچ اے۔ ایہنوں ’وجودی‘ پڑتال آکھدے نیں، کیوں جے ایہ ممکن
جہاناں دی وجودیات دا حوالہ دیندی اے۔ ہور پڑتالاں مشروط احتمالاں یا
یقین بدلّن دے نظریے ورتدیاں نیں۔ خلافِ واقعہ شرطیاں لئی کئی وکھ وکھ
منطقاں پیش کیتیاں گئیاں نیں۔ ایہناں دی کوئی اکّو لوئس--سٹالنیکر منطق وی
نہیں۔ بھانویں سٹالنیکر تے لوئس دوناں دیاں تشریحاں سبھ توں نیڑے ممکن جہاناں
دے خیال اُتے ٹکیاں نیں، اوہناں دے مفروضیاں توں وکھ وکھ معتبر استنباط
نکلدے نیں۔
\endgroup
% /OLP-0521

\OLEndChapterHook
\endgroup
% /OLP-0517

% OLP-0522 content/counterfactuals/minimal-change-semantics/minimal-change-semantics.tex
\begingroup
\def\olfilename{minimal-change-semantics}
\expandafter\def\csname ol@sectioncs\endcsname{section}

\olchapter{cnt}{min}{گھٹ توں گھٹ تبدیلی والی معنیات}

% OLP-0523 content/counterfactuals/minimal-change-semantics/introduction.tex
\begingroup
\def\olfilename{introduction}
\expandafter\def\csname ol@sectioncs\endcsname{section}

\olfileid{cnt}{min}{int}

\olsection{جان پہچان}

سٹالنیکر تے لوئس نے ’جے تیلی رگڑی جاندی، تاں بل پَیندی‘ ورگے خلافِ واقعہ
شرطیاں دیاں تشریحاں پیش کیتیاں۔ ایہ تجویزاں ایہہ سمجھاؤن لئی سن کہ ایسے جملیاں
دے سچ ہون دیاں شرطاں کیویں سمجھیاں جان۔ دوناں دا بنیادی خیال ایہ اے۔ خلافِ
واقعہ شرطیے دا سچ پرکھن لئی اوہ ممکن جہان ویکھنے پین گے جیہڑے اصل جہان توں
اینے کو وکھرے ہون کہ مقدم سچ ہو جاوے، تے ایہ فرق گھٹ توں گھٹ ہووے۔
جے تالی اوہناں ممکن جہاناں وچ سچ اے، تاں خلافِ واقعہ شرطیہ سچ اے۔
مثلاً منّ لو میرے ہتھ وچ اک تیلی تے تیلیاں دا اک پیکٹ اے۔ اصل جہان وچ
میں بس اوہناں نوں ویکھدا آں تے سوچدا آں کہ جے میں تیلی رگڑاں تاں کی ہووے۔
اصل جہان توں گھٹ توں گھٹ تبدیلی والا جہان جتھے میں تیلی رگڑدا آں، اوہ اے
جتھے میں کرن دا فیصلہ کردا آں تے تیلی رگڑ دیندا آں۔ ایہ تبدیلی گھٹ توں گھٹ
ایس لئی اے کہ ہور کجھ نہیں بدلدا۔ میں نال ہوا وچ چھلانگ نہیں لاندا، تیلی
رگڑن نال میرے والاں نوں اگ نہیں لگدی، میریاں اُنگلاں دا زور اچانک نہیں
مکدا، تے نال ای کسے لُک کے سپر سوکر نال میرے اُتے پانی نہیں سُٹیا، وغیرہ۔
اوس ہور ممکن صورت وچ تیلی بل پَیندی اے۔ ایس لئی ایہ سچ اے کہ جے میں
تیلی رگڑدا، تاں اوہ بل پَیندی۔

ایس سوکھی سمجھ آؤن والی تشریح نوں خلافِ واقعہ شرطیاں دیاں منطقاں دی رسمی
معنیات نال جوڑ سکدے آں۔ لوئس نے خلافِ واقعہ شرطیے لئی نشان ’$\boxright$‘
پیش کیتا، تے سٹالنیکر نے نشان~’$>$‘ ورتیا۔ اسیں~$\cif$ ورتاں گے تے اینوں
قضیاتی منطق وچ دو-مقامی رابطے دے طور اُتے شامل کراں گے۔ سو !!{formula}s جیہڑے $!A \lif
!B$ دی شکل دے نیں، اوہناں توں علاوہ !!{formula}s جیہڑے~$!A \cif !B$ دی
شکل دے نیں، وی نیں۔ رسمی معنیات، امکانی منطق دی تعلقی معنیات وانگ، اوہناں نمونیاں
اُتے ٹکی اے جیہناں وچ !!{formula}s جہاناں اُتے پرکھے جاندے نیں۔ پورا اترن دی
شرط جیہڑی $\mSat{M}{!A \cif !B}[w]$ دی تعریف کردی اے، $\mSat{M}{!A}[w']$ تے $\mSat{M}{!B}[w']$ دے لحاظ نال،
کجھ (ہور) جہاناں~$w'$ لئی دتی جاندی اے۔ کیہڑے $w'$؟ سمجھ مطابق،~$w$ دے سبھ توں نیڑے اوہ جہان
جیہناں لئی~$\mSat{M}{!A}[w']$ اے۔ ایس لئی نمونے وچ ’نیڑتا‘ دا تعلق وی
شامل ہونا چاہیدا اے۔

لوئس نے خلافِ واقعہ حالتاں نوں نقشے نال وکھاؤن دا اک سمجھ آؤن والا
طریقہ پیش کیتا۔ ہر ممکن جہان اک دوجے اندر سمائے ہوئے گولیاں دے سیٹ دا مرکز
اے۔ گولیاں اندر ہور جہان نیں، تے اسیں گولیاں نوں سانجھے مرکز والے دائریاں
وانگ بنا کے وکھاندے آں۔ دو گولیاں دے وچلے جہان مرکز والے جہان توں سارے
اکّو جِنّے نیڑے نیں۔ اندرلے گولے وچلے جہان ہور نیڑے، تے آلے دوالے دے گولے
وچلے جہان ہور دُور نیں۔
\begin{center}
\begin{tikzpicture}[scale=.7]\small
  \spheresystem{5}
  \draw (0,0) node {$w$};
  \propositionintersect{0}{4}{40}{3.5}
  {\tikzset{proposition/.append={smooth,tension=1.4}}
  \proposition[shift={(1.6,-1)}]{90}{1}{30}{4}}
  \path (1.5,3) node[below] {$\formula{B}$};
  \path (3,0) node {$\formula{A}$};
\end{tikzpicture}
\end{center}
سبھ توں نیڑلے $!A$ والے جہان اوہ جہان $w'$ نیں جتھے~$!A$ پورا اتردا اے
تے جیہڑے مرکز والے جہان~$w$ دے آلے دوالے سبھ توں چھوٹے گولے وچ نیں
(سُرمئی حصہ)۔ سمجھ مطابق، $!A \cif !B$،~$w$ اُتے پورا اتردا اے جے $!B$
سارے سبھ توں نیڑلے $!A$ والے جہاناں اُتے سچ ہووے۔
\endgroup
% /OLP-0523

% OLP-0524 content/counterfactuals/minimal-change-semantics/sphere-models.tex
\begingroup
\def\olfilename{sphere-models}
\expandafter\def\csname ol@sectioncs\endcsname{section}

\olfileid{con}{min}{sph}

\olsection{گولیاں والے نمونے}

خلافِ واقعہ شرطیاں دی رسمی معنیات دین دا اک طریقہ لوئس دی غیر رسمی
تشریح نوں ریاضیاتی ساخت بنانا اے۔ فیر جہان~$w$ دے آلے دوالے دے گولے
جہاناں دے سیٹ نیں۔ چونکہ گولے اک دوجے اندر سمائے نیں، ایس لئی~$w$
دے آلے دوالے دے جہاناں دے سیٹ ذیلی سیٹ ہون دے تعلق نال خطی ترتیب وچ
ہونے چاہیدے نیں۔

\begin{defn}
  \emph{گولیاں والا نمونہ} اک مرتب تِکڑی~$\mModel{M} = \tuple{W, O, V}$ اے۔
  $W$ جہاناں دا اک غیر خالی سیٹ اے، $V\colon \PVar \to \Pow{W}$ قدر بندی
  اے، تے $O\colon W \to \Pow{\Pow{W}}$ ہر جہان~$w$ نوں اک
  \emph{گولیاں دا نظام}~$O_w$ دیندا اے۔ ہر $w$ لئی $O_w$ جہاناں دے سیٹاں دا
  اک سیٹ اے، جیہڑا ایہ شرطاں پوریاں کردا اے:
  \begin{enumerate}
  \item $O_w$ دا \emph{مرکز}~$w$ اے، یعنی $\{w\} \in O_w$۔
  \item $O_w$ دے گولے \emph{اک دوجے اندر سمائے} نیں۔ جدوں وی
    $S_1$, $S_2 \in O_w$، تاں $S_1 \subseteq S_2$ یا $S_2 \subseteq S_1$،
    یعنی $O_w$،~$\subseteq$ نال خطی ترتیب وچ اے۔
  \item $O_w$ گولیاں دے غیر خالی خانداناں دے اتحاد تحت بند اے۔
  \item $O_w$ گولیاں دے غیر خالی خانداناں دے اشتراک تحت بند اے۔
  \end{enumerate}
\end{defn}

$O_w$ پچھے خیال ایہ اے کہ $w$ دے ’آلے دوالے‘ دے جہان اوہناں دی~$w$
توں دُوری دے حساب نال تہاں وچ رکھے نیں۔ سبھ توں اندرلے گولے وچ صرف $w$
اے، یعنی سیٹ~$\{w\}$۔ $w$ خود~$w$ دے، کسے وی ہور گولے دے جہاناں نالوں،
نیڑے اے۔ جے $S \subsetneq S'$، تاں $S' \setminus S$ وچلے جہان $w$ توں
~$S$ وچلے جہاناں نالوں ودھ دُور نیں۔ $S' \setminus S$ اوہ ’تہہ‘ اے جیہڑی
$S$ تے~$S'$ توں باہرلے جہاناں دے وچکار اے۔ خاص کر کے گولے وچ اوہدی باہرلی
سطح دے اندرلے سارے جہان شامل سمجھنے پین گے۔ گولے صرف وکھ وکھ تہاں نہیں۔

\begin{figure}
\begin{center}
\begin{tikzpicture}[layerwidth=1.5,scale=.6]\tiny
  \spheresystem{3}
  \propositionintersect{0}{3}{60}{6}
  \path (0,0) node[world] {$w$};
  \spherepos{-30}{2}{node[world] {$w_2$}}
  \spherepos{220}{2}{node[world] {$w_3$}}
  \spherepos{90}{2}{node[world] {$w_1$}}
  \spherepos[shift={(.1,0)}]{10}{3}{node[world] {$w_5$}}
  \spherepos[shift={(.1,0)}]{-10}{3}{node[world] {$w_6$}}
  \spherepos{225}{3}{node[world] {$w_4$}}
  \spherepos{20}{4}{node[world] {$w_7$}}
  \path (5,0) node[left] {$p$};
\end{tikzpicture}
\caption{گولیاں والے نمونے دا نقشہ}
\ollabel{fig:sphere-model}
\end{center}
\end{figure}

\olref{fig:sphere-model} دا نقشہ اوس گولیاں والے نمونے دا اے جیہدے وچ
$W = \{w, w_1, \dots, w_7\}$، $V(p) = \{w_5, w_6, w_7\}$۔ سبھ توں
اندرلا گولا $S_1 = \{w\}$ اے۔ $w$ دے سبھ توں نیڑلے جہان $w_1, w_2, w_3$
نیں، ایس لئی اگلا وڈا گولا $S_2 = \{w, w_1, w_2, w_3\}$ اے۔ ہور
دُور دے جہان $w_4$، $w_5$، $w_6$ نیں، سو سبھ توں باہرلا گولا
$S_3 = \{w, w_1, \dots, w_6\}$ اے۔ $w$ دے آلے دوالے گولیاں دا نظام
$O_w = \{S_1, S_2, S_3\}$ اے۔ جہان~$w_7$،~$w$ دے آلے دوالے دے کسے
وی گولے وچ نہیں۔ سبھ توں نیڑلے جہان جیہناں وچ $p$ سچ اے، $w_5$ تے~$w_6$
نیں۔ سو سبھ توں چھوٹا $p$ نوں تھاں دین والا گولا~$S_3$ اے۔

فارمولا~$!A$ دا جہان~$w$ اُتے گولیاں والے نمونے~$\mModel M$ وچ پورا
اترنا، $\mSat{M}{!A}[w]$، تعریف کرن لئی اسیں امکانی !!{formula}s دی تعریف
وچ $!B \cif !C$ لئی اک شق شامل کردے آں:

\begin{defn}
  $\mSat{M}{!B \cif !C}[w]$ اُدوں تے صرف اُدوں اے جدوں ہِیٹھلیاں وچوں اک ہووے:
  \begin{enumerate}
  \item\ollabel{sphere-vac} ہر $u \in \bigcup O_w$ لئی $\mSat/{M}{!B}[u]$، یا
  \item\ollabel{sphere-nonvac} کسے $S \in O_w$ لئی،
    \begin{enumerate}
      \item $\mSat{M}{!B}[u]$ کسے $u \in S$ لئی، تے
      \item ہر $v \in S$ لئی، یا $\mSat/{M}{!B}[v]$ یا $\mSat{M}{!C}[v]$۔
    \end{enumerate}
  \end{enumerate}
\end{defn}

ایس تعریف مطابق $\mSat{M}{!B \cif !C}[w]$ اُدوں تے صرف اُدوں اے جدوں
یا تاں مقدم~$!B$،~$w$ دے آلے دوالے گولیاں وچ ہر تھاں جھوٹا ہووے، یا اک
گولا~$S$ ہووے جتھے $!B$ سچ ہووے، تے مادی شرطیہ $!B \lif !C$ ایس
’’$!B$ نوں تھاں دین والا‘‘ گولے دے سارے جہاناں اُتے سچ ہووے۔ دھیان دیو کہ
اسیں تعریف وچ ایہ شرط نہیں لائی کہ $S$، $!B$ نوں تھاں دین والا
\emph{سبھ توں اندرلا} گولا اے، بھانویں سوکھی تشریح توں ایہو اُمید لگے۔
پر جے \olref{sphere-nonvac} دی شرط کسے گولے~$S$ لئی پوری اے، تاں اوہ
~$S$ دے اندر جِنّے گولے نیں، اوہناں ساریاں لئی وی پوری اے، تے خاص کر کے سبھ توں
اندرلے گولے لئی۔

ایہ وی دھیان دیو کہ گولیاں والے نمونیاں دی تعریف ایہ شرط نہیں لاندی کہ
$!B$ نوں تھاں دین والا سبھ توں اندرلا گولا \emph{موجود اے}۔ ساڈے کول
$S_1 \supsetneq S_2 \supsetneq \dots \supsetneq \{w\}$ دی لامتناہی لڑی
ہو سکدی اے جیہدی ہر گولا $!B$ نوں تھاں دیندا ہووے۔ ایس لئی $!B$ نوں
تھاں دین والا کوئی سبھ توں اندرلا گولا نہیں ہووے گی۔ ایس صورت وچ
$\mSat{M}{!B \cif !C}[w]$ اُدوں تے صرف اُدوں اے جدوں
$!B \lif !C$، $S_i$، $S_{i+1}$، \dots گولیاں وچ ہر تھاں قائم ہووے،
کسے~$i$ لئی۔
\endgroup
% /OLP-0524

% OLP-0525 content/counterfactuals/minimal-change-semantics/true-false.tex
\begingroup
\def\olfilename{true-false}
\expandafter\def\csname ol@sectioncs\endcsname{section}

\olfileid{con}{min}{tf}

\olsection{خلافِ واقعہ شرطیاں دی سچائی تے جھوٹ}

خلافِ واقعہ شرطیہ $!A \cif !B$ سچ اے (خالی شرط کارن نہیں) جے سبھ توں
نیڑلے $!A$ والے جہان سارے $!B$ والے جہان ہون، جیویں \olref{fig:true} وچ
وکھایا اے۔
\begin{figure}
\begin{center}
\begin{tikzpicture}[scale=.7]\small
  \spheresystem{5}
  \draw (0,0) node {$w$};
  \propositionintersect{0}{4}{40}{3.5}
  {\tikzset{proposition/.append={smooth,tension=1.4}}
  \proposition[shift={(1.6,-1)}]{90}{1}{30}{4}}
  \path (1.5,3) node[below] {$\formula{B}$};
  \path (3,0) node {$\formula{A}$};
\end{tikzpicture}
\caption{غیر خالی شرط نال سچا خلافِ واقعہ شرطیہ}
\ollabel{fig:true}
\end{center}
\end{figure}
خلافِ واقعہ شرطیہ $w$ اُتے اوہدوں وی سچ اے جے~$w$ دے آلے دوالے گولیاں
دے نظام وچ $!A$ نوں تھاں دین والا کوئی گولا ای نہیں۔ ایس صورت وچ اوہ
\emph{خالی شرط کارن} سچ اے (\olref{fig:vacuous} ویکھو)۔
\begin{figure}
\begin{center}
\begin{tikzpicture}[scale=.7]\small
  \spheresystem{5}
  \draw (0,0) node {$w$};
  \proposition{0}{6.5}{40}{3.5}
  {\tikzset{proposition/.append={smooth,tension=1.4}}
  \proposition[shift={(1.2,-1)}]{90}{1}{30}{4}}
  \path (1.5,3) node[below] {$\formula{B}$};
  \spherepos{0}{7}{node {$\formula{A}$}}
\end{tikzpicture}
\caption{خالی شرط کارن سچا خلافِ واقعہ شرطیہ}
\ollabel{fig:vacuous}
\end{center}
\end{figure}

ایہ دو طریقیاں نال جھوٹا ہو سکدا اے۔ اک صورت ایہ اے کہ سبھ توں نیڑلے
$!A$ والے جہاناں وچوں سارے $!B$ والے جہان نہ ہون، پر کجھ ہون۔ ایس صورت
وچ $!A \cif \lnot !B$ وی جھوٹا اے (\olref{fig:false} ویکھو)۔
\begin{figure}
\begin{center}
\begin{tikzpicture}[scale=.7]\small
  \spheresystem{5}
  \draw (0,0) node {$w$};
  \propositionintersect{0}{4}{40}{3.5}
  {\tikzset{proposition/.append={smooth,tension=1.4}}
  \proposition[shift={(1.6,0)}]{90}{1}{30}{3}}
  \path (1.5,3) node[below] {$\formula{B}$};
  \path (3,0) node {$\formula{A}$};
\end{tikzpicture}
\caption{جھوٹا خلافِ واقعہ شرطیہ، اُلٹی تالی والا وی جھوٹا}
\ollabel{fig:false}
\end{center}
\end{figure}
جے سبھ توں نیڑلے $!A$ والے جہاناں دا $!B$ والے جہاناں نال کوئی سانجھا
جہان نہیں، تاں $!A \cif !B$ جھوٹا اے۔ پر ایس صورت وچ سبھ توں نیڑلے
سارے $!A$ والے جہان $\lnot !B$ والے جہان نیں، سو $!A \cif \lnot !B$ سچ
اے (\olref{fig:false-opposite} ویکھو)۔
\begin{figure}
\begin{center}
\begin{tikzpicture}[scale=.7]\small
  \spheresystem{5}
  \draw (0,0) node {$w$};
  \propositionintersect{0}{4}{40}{3.5}
  {\tikzset{proposition/.append={smooth,tension=1.4}}
  \proposition[shift={(-1,0)}]{90}{1}{30}{3}}
  \path (-1,3) node[below] {$\formula{B}$};
  \path (3,0) node {$\formula{A}$};
\end{tikzpicture}
\caption{جھوٹا خلافِ واقعہ شرطیہ، اُلٹی تالی والا سچا}
\ollabel{fig:false-opposite}
\end{center}
\end{figure}

سخت شرطیے دے اُلٹ، خلافِ واقعہ شرطیے اتفاقی طور اُتے سچ یا جھوٹ ہو
سکدے نیں۔ \olref{fig:contingent} والا گولیاں دا نمونہ ویکھو۔
$!A$ والے جہان جیہڑے~$u$ دے سبھ توں نیڑے نیں، اوہ سارے $!B$ والے جہان نیں، سو
$\mSat{M}{!A \cif !B}[u]$۔ پر $!A$ والے جہان جیہڑے~$v$ دے سبھ توں نیڑے نیں، اوہناں وچوں کجھ جہان
$!B$ والے جہان نہیں، سو $\mSat/{M}{!A \cif !B}[v]$۔
\begin{figure}
\begin{center}
\begin{tikzpicture}[scale=.6]\small
  \spheresystem{7}
  \spheresystem[shift={(0:3.1)},dashed]{4}
  \propositionintersect{45}{4}{40}{4.5}
  \begin{scope}
    \clip \propositionplot{45}{4}{40}{4.5} ;
    \spherelayer[shift={(0:3.1)}]{3}
  \end{scope}
  \proposition[shift={(1.4,-.5)}]{90}{1}{35}{5}
  \path (0,0) node {$u$};
  \path (0:3.1) node {$v$};
  \spherepos{35}{9}{node {$\formula{A}$}}
  \spherepos{80}{9}{node {$\formula{B}$}}
\end{tikzpicture}
\end{center}
\caption{اتفاقی خلافِ واقعہ شرطیہ}
\ollabel{fig:contingent}
\end{figure}
\endgroup
% /OLP-0525

% OLP-0526 content/counterfactuals/minimal-change-semantics/antecedent-strengthening.tex
\begingroup
\def\olfilename{antecedent-strengthening}
\expandafter\def\csname ol@sectioncs\endcsname{section}

\olfileid{cnt}{min}{agg}

\olsection{مقدم مضبوط کرنا}

’مقدم مضبوط کرنا‘ دا مطلب استنباط $!A \lif !C \Entails (!A \land !B) \lif !C$
اے۔ ایہ مادی شرطیے لئی معتبر اے، پر خلافِ واقعہ شرطیاں لئی غیر معتبر۔
منّ لو ایہ سچ اے کہ جے میں ایہ تیلی رگڑدا، تاں اوہ بل پَیندی۔ (یعنی تیلی
یا پیکٹ دی رگڑ والی تھاں وچ کوئی خرابی نہیں، میں تیلی نہیں توڑاں گا، وغیرہ۔)
پر ایہ سچ نہیں کہ جے میں ایہ تیلی بیرلی خلا وچ بالدا، تاں اوہ بل پَیندی۔
سو ہِیٹھلا استنباط غیر معتبر اے:
\begin{quote}
  جے تیلی رگڑی جاندی، تاں بل پَیندی۔

  سو جے تیلی بیرلی خلا وچ رگڑی جاندی، تاں بل پَیندی۔
\end{quote}

لوئس--سٹالنیکر دی شرطیاں والی تشریح ایہ گل سمجھاندی اے۔ سبھ توں نیڑلا
جہان جتھے میں تیلی بالدا آں تے ایہ کم بیرلی خلا وچ کردا آں، اصل جہان توں
اوس سبھ توں نیڑلے جہان نالوں کِتے دُور اے جتھے میں تیلی بالدا آں۔ سو
بھانویں پچھلے جہان وچ تیلی دا بل پینا سچ اے، پہلے وچ ایہ سچ نہیں۔
تے ایویں ای ہونا چاہیدا اے۔

\begin{ex}
  گولیاں والی معنیات وچ ایہ استنباط غیر معتبر اے، یعنی ساڈے کول $p
  \cif r \Entails/ (p \land q) \cif r$ اے۔ نمونہ $\mModel{M}
  = \tuple{W, O, V}$ ویکھو، جیہدے وچ $W = \{w, w_1, w_2\}$،
  $O_w = \{\{w\}, \{w, w_1\}, \{w, w_1, w_2\}\}$، $V(p) = \{w_1, w_2\}$،
  $V(q) = \{w_2\}$، تے $V(r) = \{w_1\}$۔ $p$ نوں تھاں دین والا گولا
  $S = \{w, w_1\}$ موجود اے تے $p \lif r$ اوہدے سارے جہاناں اُتے سچ اے،
  سو $\mSat{M}{p \cif r}[w]$۔ $(p \land q)$ نوں تھاں دین والا گولا
  $S' = \{w, w_1, w_2\}$ وی موجود اے، پر $\mSat/{M}{(p \land q) \lif r}[w_2]$،
  سو $\mSat/{M}{(p \land q) \cif r}[w]$ (\olref{fig:antecedent-strengthening} ویکھو)۔

\begin{figure}
\begin{center}
\begin{tikzpicture}[layerwidth=1.5,scale=.6]\tiny
  \spheresystem{3}
  \propositionintersect{-10}{3}{30}{5}
  \begin{scope}
    \clip plot[smooth,tension=1] coordinates
          {(0,4.5) (30:.95) (4.5,-3.5)} -- (4.5,4.5) ;
    \spherefill{2}
  \end{scope}
  \draw[proposition] plot[dashed,smooth,tension=1] coordinates
          {(0,4.5) (30:.95) (4.5,-3.5)} ;
  \proposition{30}{2}{30}{5}
  \path (0,0) node[world] {$w$};
  \spherepos{30}{2}{node[world] {$w_1$}}
  \spherepos{-10}{3}{node[world] {$w_2$}}
  \spherepos{-10}{4}{node {$q$}}
  \spherepos{30}{4}{node {$r$}}
  \draw (1,4.5) node[below] {$p$};
\end{tikzpicture}
\caption{مقدم مضبوط کرن دی ردّی مثال}
\ollabel{fig:antecedent-strengthening}
\end{center}
\end{figure}
\end{ex}
\endgroup
% /OLP-0526

% OLP-0527 content/counterfactuals/minimal-change-semantics/transitivity.tex
\begingroup
\def\olfilename{transitivity}
\expandafter\def\csname ol@sectioncs\endcsname{section}

\olfileid{cnt}{min}{tra}

\olsection{تعدی}

مادی شرطیے لئی لڑی دا قاعدہ قائم اے۔ $!A \lif !B, !B \lif !C \Entails !A \lif !C$۔
دوجے لفظاں وچ مادی شرطیہ متعدی اے۔ کی خلافِ واقعہ شرطیاں لئی وی ایہو گل سچ اے؟
سٹالنیکر دی ہِیٹھلی مثال ویکھو۔
\begin{quote}
  جے جے۔ ایڈگر ہوور روسی جمیا ہوندا، تاں اوہ کمیونسٹ ہو جاندا۔

  جے جے۔ ایڈگر ہوور کمیونسٹ ہوندا، تاں اوہ غدّار ہو جاندا۔

  سو جے جے۔ ایڈگر ہوور روسی جمیا ہوندا، تاں اوہ غدّار ہو جاندا۔
\end{quote}
جے ہوور اوسے ویلے، جدوں اوہ اصل وچ جمیا سی، ریاست ہائے متحدہ دی تھاں روس
وچ جمیا ہوندا، تاں اوہ سوویت یونین وچ وڈا ہوندا تے کمیونسٹ بن جاندا
(ایہ منّ لئیے)۔ سو پہلا مقدمہ سچ اے۔ ایس طرح دوجا مقدمہ وی، جے اوہنوں
کلّا ویکھیے، سچ اے۔ پر نتیجہ جھوٹ اے۔ جے ہوور سوویت یونین وچ جمیا ہوندا،
تاں وڈا امکان ایہ سی کہ اوہ پکا کمیونسٹ ہوندا تے (اپنے ملک دا) غدّار نہ ہوندا۔
سچائی دیاں ایہ قدراں جیہڑیاں سمجھ مطابق دتیاں نیں، سٹالنیکر--لوئس دی تشریح نال
میل کھاندیاں نیں۔ ساڈے جہان دے سبھ توں نیڑلا ممکن جہان، جیہدے وچ اکّو تبدیلی
ہوور دی جنم تھاں اے، اوہ اے جتھے ہوور وڈا ہو کے سوویت یونین دا چنگا شہری
بن جاندا اے۔ ایہ سبھ توں نیڑلا ممکن جہان اے جتھے پہلے مقدمے تے نتیجے دا
مقدم سچ اے۔ اوس جہان وچ ہوور کمیونسٹ پارٹی دا وفادار رکن اے، سو غدّار نہیں۔
پر دوجا مقدمہ پرکھن لئی سانوں اک ہور جہان ویکھنا پینا اے۔ اوہ سبھ توں نیڑلا
جہان جتھے ہوور کمیونسٹ اے، اوہ اے جتھے اوہ ریاست ہائے متحدہ وچ جمیا، پاسہ
بدل لیا تے ایس طرح غدّار بن گیا۔\footnote{مثال دا زور سمجھن لئی سانوں
  وجودیات تے سیاست بارے کجھ مفروضے منّنے پین گے۔ مثلاً ایہ ممکن اے کہ
  ہوور روسی ماپیاں دے گھر جمدا، یا ایہ کہ ویہویں صدی دے پنجاہویں ورھیاں وچ
  ریاست ہائے متحدہ دے کمیونسٹ اپنے ملک دے غدّار سن۔}

\begin{prob}
  خلافِ واقعہ شرطیاں دی تعدی ناکام ہون دی اک منن جوگی، سوکھی سمجھ آؤن والی مثال لبھو۔
\end{prob}

\begin{ex}\ollabel{ex:trans-counterex}
  گولیاں والی معنیات وچ ایہ استنباط غیر معتبر اے، یعنی ساڈے کول $p
  \cif q, q \cif r \Entails/ p \cif r$ اے۔ نمونہ $\mModel{M}
  = \tuple{W, O, V}$ ویکھو، جیہدے وچ $W = \{w, w_1, w_2\}$،
  $O_w = \{\{w\}, \{w, w_1\}, \{w, w_1, w_2\}\}$، $V(p) = \{w_2\}$،
  $V(q) = \{w_1, w_2\}$، تے $V(r) = \{w_1\}$۔ $p$ نوں تھاں دین والا
  گولا $S = \{w, w_1, w_2\}$ موجود اے تے $p \lif q$ اوہدے سارے جہاناں
  اُتے سچ اے، سو $\mSat{M}{p \cif q}[w]$۔ $q$ نوں تھاں دین والا گولا
  $S' = \{w, w_1\}$ وی موجود اے تے $\mSat/{M}{q \lif r}$ اوہدے سارے جہاناں
  اُتے سچ اے، سو $\mSat{M}{q \cif r}[w]$۔ پر $p$ نوں تھاں دین والا گولا
  $\{w, w_1, w_2\}$ وچ اک ایہو جیہا جہان اے، یعنی~$w_2$، جتھے
  $\mSat/{M}{p \lif r}[w_2]$۔
\end{ex}

\begin{prob}
\olref[cnt][min][tra]{ex:trans-counterex} دی ردّی مثال نال ملدا گولیاں دا نقشہ بناؤ۔
\end{prob}

\begin{prob}
  \olref[cnt][min][tra]{ex:trans-counterex} وچ جہان $w_2$ اوہ اے جتھے
  ہوور روس وچ جمیا، کمیونسٹ اے تے غدّار نہیں، تے $w_1$ اوہ جہان اے جتھے
  ہوور ریاست ہائے متحدہ وچ جمیا، کمیونسٹ اے تے غدّار اے۔ ایس نمونے وچ
  $w_1$،~$w$ دے $w_2$ نالوں نیڑے اے۔ کی ایہ ضروری اے؟ کی تسی اک ردّی
  مثال دے سکدے او جیہڑی ایہ نہ منّے کہ ہوور دا روس وچ جمنا اوہدے
  کمیونسٹ ہون نالوں زیادہ دُور دا امکان اے؟
\end{prob}
\endgroup
% /OLP-0527

% OLP-0528 content/counterfactuals/minimal-change-semantics/contraposition.tex
\begingroup
\def\olfilename{contraposition}
\expandafter\def\csname ol@sectioncs\endcsname{section}

\olfileid{cnt}{min}{cpo}

\olsection{عکسِ نقیض}

مادی تے سخت شرطیے اپنے عکسِ نقیض دے ہم ارز نیں۔ خلافِ واقعہ شرطیاں لئی
ایہ گل سچ نہیں۔ کریٹزر دی اک مثال ویکھو:
\begin{quote}
  جے گوئٹے 1832 وچ نہ مویا ہوندا، تاں ہُن (فیر وی) مریا ہویا ہوندا۔

  جے گوئٹے ہُن مریا ہویا نہ ہوندا، تاں اوہ 1832 وچ مویا ہوندا۔
\end{quote}
پہلا جملہ سچ اے، کیوں جے انسان سیکڑیاں ورھے نہیں جیندے۔ دوجا صاف جھوٹ اے۔
جے گوئٹے ہُن مریا ہویا نہ ہوندا، تاں اوہ ہجے جیوندا ہوندا، تے ایس لئی
1832 وچ نہیں مر سکیا ہوندا۔

\begin{ex}\ollabel{ex:contraposition-counterex}
  گولیاں والی معنیات وچ عکسِ نقیض غیر معتبر اے، یعنی ساڈے کول $p
  \cif q \Entails/ \lnot q \cif \lnot p$ اے۔ $p$ دا مطلب ’گوئٹے
  1832 وچ نہیں مویا‘ تے $q$ دا مطلب ’گوئٹے ہُن مریا ہویا اے‘ لو۔
  ایہ گل نمونے $\mModel{M_1} = \tuple{W, O, V}$ وچ دکھا سکدے آں، جیہدے وچ
  $W = \{w, w_1, w_2\}$، $O = \{\{w\}, \{w, w_1\}, \{w, w_1, w_2\}\}$،
  $V(p) = \{w_1, w_2\}$ تے $V(q) = \{w, w_1\}$۔ سو $w$ اصل جہان اے
  جتھے گوئٹے 1832 وچ مویا تے ہجے وی مریا ہویا اے۔ $w_1$ اوہ (نیڑلا)
  جہان اے جتھے گوئٹے، منّ لو، 1833 وچ مویا تے ہجے وی مریا ہویا اے۔
  تے $w_2$ اوہ (دُور دا) جہان اے جتھے گوئٹے ہجے وی جیوندا اے۔
\begin{figure}
\begin{center}
\begin{tikzpicture}[scale=.6,layerwidth=1.5]\tiny
  \spheresystem{3}
  \begin{scope}
    \clip \propositionplot{0}{.8}{50}{4.5} -- (4.5,-4.5) -- (-4.5,-4.5) -- (-4.5,4.5) -- (4.5,4.5);
    \spherelayer{3}
  \end{scope}
  \propositionintersect{0}{2}{110}{5.5}
  \proposition{0}{.8}{50}{4.5}
  \path (0,0) node[world] {$w$};
  \spherepos{0}{2}{node[world] {$w_1$}}
  \spherepos{-50}{3}{node[world] {$w_2$}}
  \spherepos{28}{3}{node {$q$}}
  \spherepos{42}{3}{node {$\lnot q$}}
  \spherepos{-55}{4}{node {$p$}}
  \spherepos{-67}{4}{node {$\lnot p$}}
\end{tikzpicture}
\caption{عکسِ نقیض دی ردّی مثال}
\ollabel{fig:contraposition}
\end{center}
\end{figure}
$p$ نوں تھاں دین والا گولا $S = \{w, w_1\}$ موجود اے تے $p \lif q$
اوہدے سارے جہاناں اُتے سچ اے، سو $\mSat{M}{p \cif q}[w]$۔ پر $\lnot q$
نوں تھاں دین والا گولا $\{w, w_1, w_2\}$ وچ اک ایہو جیہا جہان اے،
یعنی~$w_2$، جتھے $q$ جھوٹ تے $p$ سچ اے، سو $\mSat/{M}{\lnot q \lif \lnot p}[w_2]$۔
\end{ex}
\endgroup
% /OLP-0528

\OLEndChapterHook
\endgroup
% /OLP-0522

\OLEndPartHook
\endgroup
% /OLP-0516

% OLP-0529 content/set-theory/set-theory.tex
\begingroup
\def\olfilename{set-theory}
\expandafter\def\csname ol@sectioncs\endcsname{section}

\olpart{sth}{سیٹ نظریہ}

% OLP-0530 content/set-theory/story/story.tex
\begingroup
\def\olfilename{story}
\expandafter\def\csname ol@sectioncs\endcsname{section}

\olchapter{sth}{story}{مرحلہ وار تصور}

% OLP-0531 content/set-theory/story/extensionality.tex
\begingroup
\def\olfilename{extensionality}
\expandafter\def\csname ol@sectioncs\endcsname{section}

\olfileid{sth}{story}{extensionality}
\olsection{عنصراں نال تعیّن دا اصول}

سبھ توں پہلی گل ایہ اے کہ سیٹ اپنے !!{element}s نال پچھانے جاندے نیں۔
ہور ٹھیک لفظاں وچ:

\begin{axiom}[عنصراں نال تعیّن دا اصول]
جے سیٹ $A$ تے $B$ دے !!{element}s اکّو نیں، تاں $A$ تے $B$ اکّو سیٹ نیں۔
\[
  \lforall[A][\lforall[B][(\lforall[x][(x \in A \liff x \in B)] \lif
  \eq[A][B])]]
\]
\end{axiom}

اسیں ایہ گل \olref[sfr][][]{part} وچ ہر تھاں منّ کے چلے ساں۔ پر ایہ گل
دوبارہ آکھن جوگی اے۔ عنصراں نال تعیّن دا مسلّمہ ایہ بنیادی خیال بیان کردا اے
کہ سیٹ دا تعیّن اوہدے !!{element}s نال ہوندا اے۔ (ایس لحاظ نال سیٹاں نوں
\emph{تصوراں} توں وکھ سمجھ سکدے آں، کیونکہ عین اوہو چیزاں کئی وکھ وکھ
تصوراں تھلّے آ سکدیاں نیں۔)

ایہ اصول کیوں منّیے؟ ایہ کہنا واجب لگدا اے کہ جیہڑا وی عنصراں نال
تعیّن توں انکار کرے، اوہ اوہو جیہی ہر شے چھڈّن دا فیصلہ کردا اے جیہنوں
\emph{سیٹ نظریہ} آکھیا وی جا سکے۔ سیٹ نظریہ عین اوہناں اکٹھاں دا نظریہ اے
جیہناں دا تعیّن اوہناں دے عنصراں نال ہوندا اے۔

پر \olref[sth][][]{part} وچ اصلی اوکڑ ایسے اصول دین دی اے جیہڑے دسّن کہ
\emph{کیہڑے سیٹ موجود نیں}۔ تے پتا لگدا اے کہ ایس سوال دا اکّو سچ مچ
’ظاہر‘ جواب ثابت کیتا جا سکدا اے کہ غلط اے۔
\endgroup
% /OLP-0531
% OLP-0532 content/set-theory/story/russells-paradox-again.tex
\begingroup
\def\olfilename{russells-paradox-again}
\expandafter\def\csname ol@sectioncs\endcsname{section}

\olfileid{sth}{story}{rus}
\olsection{رسل دا پیراڈوکس (دوبارہ)}

\olref[sfr][][]{part} وچ اسیں اک سادہ سیٹ نظریے نال کم کیتا سی۔ پر اک
\emph{بہت ای} سادے خیال مطابق سیٹ صرف محمولاں دے دائرے ہوندے نیں۔ ایہ سادہ خیال ہِیٹھلا اصول منّن
لئی آکھدا اے:

\begin{defish}
  \emph{سادہ خاصیت توں سیٹ بناؤن دا اصول۔} $\Setabs{x}{\phi(x)}$ ہر فارمولے $\phi$ لئی موجود اے۔
\end{defish}

بھانویں ایہ اصول منّن نوں دل کردا اے، پر ثابت کیتا جا سکدا اے کہ ایس
وچ تضاد اے۔ ایہ اسیں \olref[sfr][set][rus]{sec} وچ ویکھیا سی۔ پر نتیجہ اینا
اہم تے اینا سدھا اے کہ اینوں دوبارہ کہنا بندا اے۔ عین اوہناں لفظاں وچ۔

\begin{thm}[رسل دا پیراڈوکس]
کوئی سیٹ $R = \Setabs{x}{x \notin x}$ موجود نہیں۔
\end{thm}

\begin{proof}
جے $R = \Setabs{x}{x \notin x}$ موجود ہووے، تاں
$R \in R$ اُدوں تے صرف اُدوں اے جدوں $R \notin R$، تے ایہ تضاد اے۔
\end{proof}

رسل نے ایہ نتیجہ جون 1901 وچ لبھیا۔ (پر اوہنے پیراڈوکس عین اوس صورت وچ
پیش نہیں کیتا جیویں اسیں ہُنے کیتا اے۔ اوہ فریگے دے سیٹ نظریے بارے سوچ ریا سی،
جیہڑا \emph{Grundgesetze} وچ بیان ہویا اے۔ اسیں \olref[blv]{sec} وچ ایس گل
ول مُڑاں گے۔) رسل نے 16 جون 1902 نوں فریگے نوں خط لکھ کے اوہدے نظام دا
تضاد سمجھایا۔ خط کتابت تے کجھ پچھوکڑ لئی \citet[pp.~124--8]{Heijenoort1967} ویکھو۔

ایہ گل زور دے کے آکھنی چاہیدی اے کہ دو سطراں والا ایہ ثبوت \emph{خالص منطق}
دا نتیجہ اے۔ منّیا کہ اسیں اپنی علامت $\Setabs{x}{x \notin x}$ وچ اک
(غیر منطقی؟) مسلّمہ، عنصراں نال تعیّن، لُکواں ورتیا اے۔ کیوں جے
$\Setabs{x}{\phi(x)}$ دا مطلب اے کہ جے کوئی سیٹ موجود اے تاں اوہ $\phi$
والیاں چیزاں دا \emph{اکّو اک} سیٹ اے (عنصراں نال تعیّن دے اصول مطابق)۔
پر عنصراں نال تعیّن دا اشارہ وی ہٹ سکدا اے جے نتیجہ ایویں آکھیے:
\emph{کوئی ایہو جیہا سیٹ نہیں جیہدے رکن عین اوہ سیٹ ہون جیہڑے اپنے آپ دے رکن نہیں}۔
تے ایہ گل خاص طور اُتے سیٹاں نال جڑی ہوئی نہیں۔ جیویں رسل نے آپ دھیان دوایا،
عین ایہو جیہی دلیل نال ایہ نتیجہ نکلدا اے کہ \emph{کوئی مرد عین اوہناں
مرداں دی داڑھی نہیں مونڈدا جیہڑے اپنی داڑھی آپ نہیں مونڈدے}۔ یا
\emph{کوئی پگ کُتّا عین اوہناں پگ کُتیاں نوں نہیں سُنگھدا جیہڑے اپنے آپ نوں
نہیں سُنگھدے}۔ تے ایویں ہور مثالاں۔ نتیجے دی عام صورت بس ایہ اے:
\[
\lnot \exists x \forall z(Rzx \liff \lnot R zz).
\]
تے ایہ پہلے درجے دی منطق دا اک قضیہ (سانچہ) اے۔ ایس لئی سیٹ نظریے وچ
تھوڑی جیہی ردّوبدل کر کے رسل دے پیراڈوکس توں نہیں بچ سکدے۔ ایہ سیٹ نظریے
تک \emph{پہنچن} توں پہلاں ای پیدا ہوندا اے۔ جے اسیں (کلاسیکی) پہلے درجے دی
منطق ورتنی اے، تاں سانوں ایہ \emph{منّنا} پئے گا کہ کوئی سیٹ $R = \Setabs{x}{x\notin x}$ موجود نہیں۔

نتیجہ ایہ اے۔ جے تسی سادہ خاصیت توں سیٹ بناؤن دا اصول منّنا چاہندے او
تے تضاد توں \emph{بچنا} وی، تاں تسی صرف \emph{سیٹ نظریے} وچ ردّوبدل نال
کم نہیں چلا سکدے۔ تہانوں اپنی \emph{منطق} دی بناوٹ بدلنی پئے گی۔

بے شک غیر کلاسیکی منطقاں والے سیٹ نظریے پیش کیتے گئے نیں۔ پر ایہ
گھٹ توں گھٹ اینا تاں آکھ سکدے آں کہ اوہ غیر معیاری نیں۔ رسل دے پیراڈوکس
نال نِبڑن دا معیاری طریقہ اینوں موجود نہ ہون دا سدھا ثبوت منّنا تے فیر
ایہدے نال کم کرنا سِکھنا اے۔ اسیں ایہو طریقہ ورتاں گے۔
\endgroup
% /OLP-0532
% OLP-0533 content/set-theory/story/predicativity.tex
\begingroup
\def\olfilename{predicativity}
\expandafter\def\csname ol@sectioncs\endcsname{section}

\olfileid{sth}{story}{predicative}	

\olsection{اسنادی تے غیر اسنادی تعریفاں}

رسل دے سیٹ، $R$، دی تعریف $\Setabs{x}{x \notin x}$ نال دتی گئی سی۔
ہور کھول کے دسیے، تاں $R$ اوہ سیٹ ہووے گا جیہدے رکن سارے، تے صرف اوہی،
سیٹ ہوون گے جیہناں بارے ایہ گل سچ نہیں کہ اوہ اپنے آپ دے رکن نہیں۔
سو $R$ دی تعریف وچ اسیں اوس دائرے اُتے مقداری شرط لاندے آں جیہدے وچ
$R$ (جے موجود ہوندا) شامل ہوندا۔

ایہ اک \emph{غیر اسنادی} تعریف اے۔ عام طور اُتے آکھ سکدے آں کہ تعریف
اُدوں تے صرف اُدوں غیر اسنادی اے جدوں اوہ ایسے دائرے اُتے مقداری شرط
لاوے جیہدے وچ تعریف کیتی جاندی شے آپ شامل ہووے۔

پیراڈوکساں دے پچھوں وائٹ ہیڈ، رسل، پوانکارے تے وائل نے ایہو جیہیاں
غیر اسنادی تعریفاں نوں ’’خراب چکر وچ پَون والیاں‘‘ آکھ کے ردّ کیتا:
\begin{quote}
	جیہڑے پیراڈوکس توں بچنا اے، اوہناں دی پڑتال توں پتا لگدا اے کہ سارے
	اک طرح دے خراب چکر توں نکلدے نیں۔ ایہ خراب چکر ایس مفروضے توں پیدا
	ہوندے نیں کہ چیزاں دے کسے اکٹھ وچ کجھ ایہو جیہے رکن ہو سکدے نیں
	جیہناں دی تعریف صرف پورے اکٹھ دے ذریعے ہو سکدی اے[\ldots. \textparagraph]

	جس اصول نال اسیں ناجائز کُلّاں توں بچ سکدے آں، اوہ ایویں بیان ہو
	سکدا اے۔ ’کسے اکٹھ دے \emph{سارے} رکن ورتن والی کوئی شے اوس اکٹھ دی
	رکن نہیں ہو سکدی‘۔ یا اُلٹے پاسے، ’جے کوئی اکٹھ کُلّ ہون دی صورت وچ
	ایہو جیہے رکن رکھدا ہووے جیہڑے صرف اوس کُلّ نال تعریف ہو سکن، تاں
	اوس اکٹھ دا کوئی کُلّ نہیں۔‘ اسیں ایہنوں ’خراب چکر دا اصول‘ آکھاں گے،
	کیوں جے ایہ سانوں ناجائز کُلّ منّن نال پیدا ہون والے خراب چکراں توں
	بچاندا اے۔
	\citep[p.~37]{WhiteheadRussell1910}
\end{quote}
جے اسیں \emph{خراب چکر دا اصول} ردّ کرن وچ اوہناں دے پچھے چلیے، تاں اسیں
(\olref[sth][story][rus]{sec} والے) تباہ کن سادے خاصیت توں سیٹ بناؤن دے
خاکے دی تھاں ایہو جیہا اصول رکھن دی کوشش کر سکدے آں:

\begin{defish}
\emph{اسنادی طور اُتے خاصیت توں سیٹ بنانا۔} ہر اوہ فارمولا $\phi$ جیہڑا صرف سیٹاں اُتے مقداری شرط لاندا اے، اوہدے لئی سیٹ$^\prime$ $\Setabs{x}{\phi(x)}$ موجود اے۔
\end{defish}

جِنّی دیر سیٹ$^{\prime}$ سیٹ نہیں، کوئی تضاد نہیں نکلے گا۔

بدقسمتی نال اسنادی طور اُتے خاصیت توں سیٹ بناؤن دا اصول بہتا \emph{وسیع}
نہیں۔ ایہ سانوں نویں قسم دیاں شےواں، سیٹ$^\prime$، نال ملاندا اے۔
سو سانوں اوہ فارمولے وی ویکھنے پین گے جیہڑے سیٹ$^\prime$ اُتے مقداری شرط
لاندے نیں۔ جے اوہ ہمیشہ سیٹ$^\prime$ دین، تاں رسل دا پیراڈوکس دوبارہ
پیدا ہو جائے گا۔ صرف اوہ سیٹ$^\prime$ ویکھ لو جیہدے وچ سارے اوہ سیٹ$^\prime$ ہون
جیہڑے اپنے آپ دے رکن نہیں۔ سو ایس خیال نوں اگے لے جائیے تاں آکھنا پئے گا
کہ سیٹ$^\prime$ اُتے مقداری شرط لاؤن والا فارمولا اوہدے مطابق اک سیٹ$^{\prime\prime}$
دیندا اے۔ فیر سانوں سیٹ$^{\prime\prime\prime}$، سیٹ$^{\prime\prime\prime\prime}$،
وغیرہ چاہیدے ہون گے۔ اینے سارے پرائم نشان توں بچن لئی ایہناں نوں سیٹ$_0$،
سیٹ$_1$، سیٹ$_2$، سیٹ$_3$، سیٹ$_4$، \ldots سمجھنا سوکھا اے۔
ایس طرح (سادے) قسماں دے نظریے تک اک راہ مل جاندی اے۔

ایہو جیہے نظریے اُتے کجھ ظاہر اعتراض نیں (پر ایہ ظاہر نہیں کہ ایہ
\emph{فیصلہ کن} اعتراض نیں)۔ مختصر گل ایہ اے کہ نتیجے والا نظریہ ورتنا
اوکھا اے، ایہ کئی وکھ وکھ قسماں دیاں شےواں فراخ دلی نال منّ لیندا اے،
تے آخر ایہ وی صاف نہیں کہ غیر اسنادی تعریفاں \emph{ایڈیاں ماڑیاں} نیں۔

آخری گل صاف کرن لئی \citeauthor{Ramsey1925} دا ایہ قول ویکھو:
\begin{quote}
	اسیں کسے بندے نوں اک ٹولی دا سبھ توں لمّا بندہ آکھ سکدے آں۔ ایویں
	اوہ اک ایسے کُلّ دے ذریعے پچھانیا جاندا اے جیہدا اوہ آپ رکن اے،
	تے کوئی خراب چکر وی نہیں پیندا۔ \citep{Ramsey1925}
\end{quote}
ریمزی دی گل ایہ اے کہ ’’ٹولی دا سبھ توں لمّا بندہ‘‘ \emph{واقعی} اک غیر اسنادی
تعریف، پر ایہ صاف طور اُتے بالکل جائز اے۔

کوئی ایہ جواب دے سکدا اے کہ ایس صورت وچ اسیں سبھ توں لمّے بندے دی پچھان
\emph{اسنادی} طریقے نال وی کرا سکدے آں۔ مثلاً اسیں اوس بندے ول صرف اشارہ
کر سکدے آں۔ فیر غیر اسنادی تعریفاں اُتے اعتراض صرف اوہناں شےواں تک
رکھنا پئے گا جیہناں دی پچھان \emph{صرف} غیر اسنادی طریقے نال ہو سکدی اے۔
پر فیر وی ایہ ہور سمجھنا پئے گا کہ ایہو جیہی ’’لازمی غیر اسنادیت‘‘
کاہتوں اینی ماڑی اے۔\footnote{ہور تفصیل لئی \citet{Linnebo2010} ویکھو۔}

ایہ منّن والی گل اے کہ جے اسیں تعریف توں ایہ چاہندے آں کہ اوہ سانوں
کوئی شے \emph{بنان} دا طریقہ دسے، تاں غیر اسنادی تعریفاں بڑی اوکڑ پاندیاں
نیں۔ ایہو جیہی تعریف دتی ہووے تاں بندہ سچ مچ خراب چکر وچ پھس جاندا اے۔
غیر اسنادی طور اُتے دسی شے بناؤن لئی \emph{پہلاں} ساریاں شےواں بناؤنیاں
پین گیاں (اوہ غیر اسنادی طور اُتے دسی شے وی)، کیوں جے اوہ شے ساریاں شےواں
دے لحاظ نال دسی گئی اے۔ سو اوہ شے بناؤن توں پہلاں اوہنوں آپ ای بنانا پئے گا۔
پر ’’لازمی طور اُتے غیر اسنادی‘‘ طریقے نال دسے سیٹاں خلاف ایہ سنجیدہ اعتراض
صرف اوہدوں اے جے اسیں سیٹاں نوں ایہو جیہیاں شےواں سمجھدے آں جیہڑیاں
اسیں \emph{بناندے} آں۔ تے اسیں (شاید) ایویں نہیں سمجھدے۔

ایس طرح، چنگا ہویا یا ماڑا، عام ہویا طریقہ (غیر)اسنادیت بارے سخت پینتڑا
نہیں لیندا۔ ایہدے وچ اوہ طریقہ ورتیا جاندا اے جیہڑا ہُن اکٹھا کرن والا
مرحلہ وار طریقہ سمجھیا جاندا اے۔ آخر ایہ سانوں اپنے سیٹ ’’مرحلیاں‘‘ وچ
ونڈن دیوے گا۔ ایہ \emph{کجھ کجھ} اویویں اے جیویں اسنادی طریقہ شےواں نوں
سیٹ$_0$، سیٹ$_1$، سیٹ$_2$، \ldots وچ ونڈدا اے۔ پر اسیں اوہناں دی قسم وچ
کوئی فرق فرض نہیں کراں گے۔
\endgroup
% /OLP-0533
% OLP-0534 content/set-theory/story/cumulative-approach.tex
\begingroup
\def\olfilename{cumulative-approach}
\expandafter\def\csname ol@sectioncs\endcsname{section}

\olfileid{sth}{story}{approach}
\olsection{اکٹھا کرن والا مرحلہ وار طریقہ}

ہِیٹھاں ذرا کھول کے دسیا اے کہ اسیں سیٹاں نوں مرحلیاں وچ کیویں ونڈاں گے:
\begin{quote}
	سیٹ \emph{مرحلیاں} وچ بندے نیں۔ ہر مرحلے $S$ لئی کجھ مرحلے $S$ توں
	\emph{پہلاں} آؤندے نیں۔ مرحلے $S$ اُتے، $S$ توں پہلاں دے مرحلیاں وچ بنے
	سیٹاں دے ہر اکٹھ نوں اک سیٹ بنایا جاندا اے۔ مرحلیاں وچ بنے سیٹاں توں
	بناں ہور کوئی سیٹ نہیں۔ \citep[p.~323]{Shoenfield:AST}
\end{quote}
ایہ \emph{سیٹ دے اکٹھا کرن والے مرحلہ وار تصور} دی اک مُڈھلی جھلک اے۔
\olref[sth][][]{part} وچ جیہڑا رسمی سیٹ نظریہ اسیں پیش کراں گے، اوہ
ایس اُتے ٹکیا ہووے گا۔

ایہ گل ہور کھول کے ویکھیے۔ شون فیلڈ دے بیان مطابق ہر مرحلے اُتے
اسیں نویں سیٹ اوہناں {سیٹاں} توں بناندے آں جیہڑے پچھلے مرحلیاں توں
ساڈے کول سن۔ سو شون فیلڈ دے خیال مطابق، پہلے مرحلے، یعنی مرحلے $0$ اُتے،
کوئی \emph{پچھلا} مرحلہ نہیں۔ ایس لئی \emph{ہور وی صاف} اے کہ ساڈے کول
پچھلے مرحلیاں توں کوئی سیٹ نہیں۔\footnote{اسیں کیوں منّیئے کہ پہلا مرحلہ
\emph{موجود اے}؟ \stagesord{} دا حاشیہ \olref[z][story]{sec} وچ ویکھو۔}
سو اسیں صرف اک سیٹ بناندے آں: اوہ سیٹ جیہدے وچ کوئی !!{element}s نہیں،
$\emptyset$۔ مرحلے $1$ اُتے پچھلے مرحلیاں توں ساڈے کول عین اک سیٹ اے،
سو اکّو نواں سیٹ $\{\emptyset\}$ اے۔ مرحلے $2$ اُتے پچھلے مرحلیاں توں
ساڈے کول دو سیٹ نیں تے اسیں دو نویں سیٹ $\{\{\emptyset\}\}$ تے
$\{\emptyset, \{\emptyset\}\}$ بناندے آں۔ مرحلے $3$ اُتے پچھلے مرحلیاں توں
ساڈے کول چار سیٹ نیں، سو اسیں باراں نویں سیٹ بناندے آں\ldots۔
ایس طرح سیٹاں دی اکٹھا کرن والی مرحلہ وار صورت کجھ ایہو جیہی ہووے گی
(عدد مرحلے دسدے نیں):
\begin{center}
	\begin{tikzpicture}[scale=0.6]
	\tikzset{cut_here/.style={densely dotted}}
	\draw (-4,6) -- (-1, 0)-- (2,6); 
	\draw[cut_here] (-5,8)--(-4,6);
	\draw[cut_here] (2,6)--(3,8);
	\draw(-1.5, 1)--(-0.5, 1);
	\draw(-2, 2)--(0, 2);
	\draw(-2.5, 3)--(0.5,3);
	\draw(-3, 4)--(1, 4);
	\draw(-3.5, 5)--(1.5, 5);
	\draw(-4, 6)--(2, 6);
	\node[label] at (0, 0) {\small 0};
	\node[label] at (0.5, 1) {\small 1};
	\node[label] at (1, 2) {\small 2};
	\node[label] at (1.5, 3) {\small 3};
	\node[label] at (2, 4) {\small 4};
	\node[label] at (2.5, 5) {\small 5};
	\node[label] at (3, 6) {\small 6};
	\end{tikzpicture}
\end{center}
سو اسیں ایہ بیان کیوں منّیئے؟

اک وجہ ایہ اے کہ ایہ بیان سوہنا اے تے ورتنا سوکھا اے۔ سبھ توں ظاہر
بیان، یعنی سادے خاصیت توں سیٹ بناؤن دے اصول، دے ڈِگ پَین توں پچھوں
سانوں کوئی چنگا بیان چاہیدا اے۔

پر ایہ بیان \emph{صرف} سوہنا ای نہیں۔ ساڈے کول چنگی وجہ اے ایہ منّن
دی کہ ایس بیان اُتے ٹکیا ہر سیٹ نظریہ \emph{بے تضاد} ہووے گا۔ وجہ ایہ اے۔

سیٹ دے اکٹھا کرن والے مرحلہ وار تصور مطابق اسیں سیٹ مرحلیاں وچ بناندے آں،
تے اوہناں دے !!{element}s ایسی شےواں ہونیاں چاہیدیاں نیں جیہڑیاں
\emph{پہلاں ای} ساڈے کول سن۔ سو کسے وی مرحلے~$S$ لئی اسیں ایہ سیٹ بنا سکدے آں:
\[
	R_S = \Setabs{x}{x \notin x \text{ تے $x$ مرحلے $S$ توں پہلاں موجود سی}}
\]
رسل دا پیراڈوکس ثابت کرن والی دلیل ہُن ایہ ثابت کرے گی کہ $R_S$ آپ مرحلے
$S$ توں پہلاں ساڈے کول نہیں سی۔ تے ایہ تضاد نہیں۔ نالے جے اسیں سیٹ دا
اکٹھا کرن والا مرحلہ وار تصور منّ لئیے، تاں سانوں ایہ \emph{آس} وی نہیں
سی رکھنی چاہیدی کہ اسیں رسل دا سیٹ آپ بنا سکاں گے۔ اوہ اوہناں سارے
سیٹاں دا سیٹ ہووے گا جیہڑے اپنے آپ دے رکن نہیں تے جیہڑے ’’کدی وی ساڈے
کول ہون گے‘‘۔ مختصر گل ایہ اے کہ (جیویں ثابت کیتا جا سکدا اے) رسل دا سیٹ
نہ بنا سکنا ایس بیان مطابق \emph{حیرانی} دی گل نہیں۔ اسیں تے ایہو \emph{اندازہ}
لاندے ساں۔

%In one sense, then, the cumulative-iterative conception of set yields a response to Russell's Paradox which is rather like the \emph{predicativist}'s response. After all, the predicativist said that the collection of all non-self-membered sets$_n$ is a set$_{n+1}$, and this set$_n$ will not be self-membered. (Indeed, most predicativists will treat it as \emph{ungrammatical} to try to ask whether a set is self-membered.)
%
%But there is an important difference between the cumulative-iterative approach and the predicativist approach: the cumulative-iterative approach treats all of our entities as being \emph{of the same kind}. The predicativist will presumably have an empty set$_0$ (i.e., a set$_0$ with no !!{element}s), and an empty set$_1$ (i.e., a set$_1$ with no !!{element}s), and these will be \emph{different entities}. But on the cumulative-iterative approach, there is just one empty set, $\emptyset$, and it will be available at every stage of the hierarchy. 

%Indeed, the Axiom of Extensionality, stated at the start of this chapter, will hold true of the elements in this hierarchy (so that there can be \emph{only one} ``empty'' set). But I do not intend to use the cumulative-iterative conception to justify Extensionality (see \cite{Boolos1971}). Again: I suggest that we should accept Extensionality, just because we are interested in extensional collections. 
\endgroup
% /OLP-0534
% OLP-0535 content/set-theory/story/urelements.tex
\begingroup
\def\olfilename{urelements}
\expandafter\def\csname ol@sectioncs\endcsname{section}

\olfileid{sth}{story}{urelements}
\olsection{غیر سیٹ بنیادی شےواں منّیئے؟}

اگلے کجھ باباں وچ اسیں سیٹ دے اکٹھا کرن والے مرحلہ وار تصور توں مسلّمے
کڈھن دی کوشش کراں گے۔ پر اگے ودھن توں پہلاں \emph{غیر سیٹ بنیادی شےواں}
بارے کجھ ہور کہنا پئے گا۔

\olref[approach]{sec} دی صورت سانوں نویں سیٹ صرف پرانے \emph{سیٹاں}
توں بناؤن دی اجازت دیندی سی۔ پر ہو سکدا اے اسیں منّنا چاہیئے کہ کجھ
\emph{غیر سیٹ}، گایاں، سُور، ریت دے ذرے یا ہور کجھ، سیٹاں دے !!{element}s
ہو سکن۔ اوس صورت وچ اسیں کجھ بنیادی شےواں، \emph{غیر سیٹ بنیادی شےواں}،
نال شروع کراں گے۔ فیر آکھاں گے کہ ہر مرحلے $S$ اُتے غیر سیٹ بنیادی شےواں
یا $S$ توں پہلاں بنے سیٹاں نوں (کسے وی میل نال) اکٹھا کر کے ’’سارے ممکن‘‘
سیٹ بناواں گے۔ نتیجے دی صورت کجھ ایہو جیہی ہووے گی:
\begin{center}
	\begin{tikzpicture}[scale=0.6]
	\tikzset{cut_here/.style={densely dotted}}
	\draw (-4,6) -- (-1, 0) -- (1,0) -- (4,6); 
	\draw[cut_here] (-5,8)--(-4,6);
	\draw[cut_here] (5,8)--(4,6);
	\draw(-1.5, 1)--(1.5, 1);
	\draw(-2, 2)--(2, 2);
	\draw(-2.5, 3)--(2.5,3);
	\draw(-3, 4)--(3, 4);
	\draw(-3.5, 5)--(3.5, 5);
	\draw(-4, 6)--(4, 6);
	\node[label] at (2, 0) {\small 0};
	\node[label] at (2.5, 1) {\small 1};
	\node[label] at (3, 2) {\small 2};
	\node[label] at (3.5, 3) {\small 3};
	\node[label] at (4, 4) {\small 4};
	\node[label] at (4.5, 5) {\small 5};
	\node[label] at (5, 6) {\small 6};
	\end{tikzpicture}
\end{center}
سو ہُن سانوں فیصلہ کرنا اے۔ \emph{کی اسیں غیر سیٹ بنیادی شےواں منّیئے؟}

فلسفیانہ لحاظ نال اپنے نظریے وچ غیر سیٹ بنیادی شےواں شامل کرن دی تُک بندی
اے۔ وڈی وجہ ایہ اے کہ سیٹ نظریہ \emph{لاگو کرن جوگا} ہووے۔ مثال لئی
\olref[sfr][siz][]{chap} توں یاد کرو کہ دو سیٹ $A$ تے~$B$ اکّو جِڈّے،
یعنی $\cardeq{A}{B}$، اُدوں تے صرف اُدوں نیں جدوں اوہناں وچ اک یک بہ یک
تے پورا نقشہ ہووے۔ جے کھیت دیاں گایاں تے واڑے دے سُور، دوناں دے سیٹ ہون،
تاں اسیں ’’گایاں تے سُور اکّو گنتی وچ نیں‘‘ والے دعوے نوں سیٹ نظریے نال
بیان کر سکدے آں۔ پر جے اسیں غیر سیٹ بنیادی شےواں اُتے روک لائیے، تاں گایاں
تے سُوراں دے سیٹ \emph{نہیں} بنن گے تے اوہ بیان ممکن نہیں ہووے گا۔
اصل وچ سیٹ نظریہ سیٹاں توں چھٹ ہور کسے شے اُتے سیدھے طریقے نال لاگو ای
نہیں کر سکاں گے۔ (غیر سیٹ بنیادی شےواں شامل کرن دیاں ہور وجوہاں لئی
\citealt[pp.~vi, 24, 50--1]{Potter2004} ویکھو۔)

پر ریاضی وچ غیر سیٹ بنیادی شےواں منّن دا رواج بہت گھٹ اے۔ اک وجہ ایہ
اے کہ اوہناں توں بناں سیٹ نظریہ لکھنا \emph{ماڑا جیہا} سوکھا اے۔
پر کدی کدی اوہناں نوں سیٹ نظریے توں باہر رکھن لئی ہور دلچسپ دلیلاں ملدیاں نیں:
\begin{quote}
	جے اسیں منّ دے آں کہ سیٹ نظریہ ریاضی دی بنیاد اے، تاں سانوں صرف سیٹاں
	دی گل کر کے ساری ریاضی بیان کرن جوگا ہونا چاہیدا اے۔ سو ساڈے متغیر دی
	حد وچ گایاں تے سُوراں ورگیاں شےواں نہیں ہونیاں چاہیدیاں۔
	%But if $C$ is a cow, $\{C\}$ is a set, but not a legitimate mathematical object.
	\citep[p.~8]{Kunen1980}
\end{quote}
سو جے لاگو ہون اُتے زور دیئیے، تاں غیر سیٹ بنیادی شےواں \emph{شامل کرن}
دا خیال آؤندا اے۔ جے گھٹا کے بنیاد دین والا مقصد (ریاضی نوں خالص سیٹ
نظریے تک گھٹانا) سامھنے ہووے، تاں اوہناں نوں \emph{باہر رکھن} دا خیال آ
سکدا اے۔ ماڑی جیہی سُستی وی اوہناں نوں باہر رکھن ول اشارہ کردی اے۔

اسیں سبھ توں سُستا راہ پھڑاں گے۔ پر کچھ وجہ پڑھاؤن نال وی جڑی اے۔
ساڈا مقصد تہانوں سیٹ نظریے دیاں اوہ بنیادی گلاں سمجھانا اے جیہڑیاں سیٹ
نظریے دے فلسفے دی شروعات لئی چاہیدیاں نیں۔ تے اوس فلسفیانہ لکھت وچ زیادہ تر
ایسے سیٹ نظریاں دی گل اے جیہڑے غیر سیٹ بنیادی شےواں توں \emph{بناں} لکھے
گئے نیں۔ سو ایہ کتاب شاید ہور کم آوے جے اوہ اوس لکھت دے کافی نیڑے رہے۔
\endgroup
% /OLP-0535
% OLP-0536 content/set-theory/story/grundgesetze.tex
\begingroup
\def\olfilename{grundgesetze}
\expandafter\def\csname ol@sectioncs\endcsname{section}

\olfileid{sth}{story}{blv}

\olsection{ضمیمہ۔ فریگے دا پنجواں بنیادی قانون}

\olref[rus]{sec} وچ اسیں سمجھایا سی کہ رسل نے اپنا پیراڈوکس اوس نظام لئی
اک مسئلے دے طور اُتے پیش کیتا جیہڑا فریگے نے اپنی \emph{Grundgesetze} وچ
بیان کیتا۔ فریگے دے نظام وچ سادہ خاصیت توں سیٹ بناؤن دا اصول سیدھی صورت
وچ شامل نہیں سی۔ سو ایس ضمیمے وچ اسیں بہت مختصر سمجھاواں گے کہ فریگے دے
نظام وچ \emph{کیہ شامل سی}، اوہدا سادے خاصیت توں سیٹ بناؤن دے اصول نال
کیہ تعلق اے تے رسل دے پیراڈوکس نال کیہ تعلق اے۔

فریگے دا نظام \emph{دوجے درجے دا} اے، تے ایس نوں \emph{کسے تصور دے
دائرے} دی دھارنا ٹھیک طرح لکھن لئی بنایا گیا سی۔\footnote{ٹھیک گل ایہ اے
کہ فریگے اک ہور عام دھارنا رسمی طور اُتے لکھن دی کوشش کردا اے، فنکشن
دی ’’قدراں دا دائرہ‘‘۔ تصوراں دے دائرے اوس عام دھارنا دی اک خاص صورت
نیں۔ تفصیل لئی \citet[pp.\ 8--9]{Heck2012} ویکھو۔} فریگے توں ملی علامت
ورت کے اسیں $\fregeext{x}{F(x)}$ نوں \emph{تصور~$F$ دا دائرہ} لکھاں گے۔
ایہ اک ترکیب اے جیہڑی \emph{محمول} ’’$F$‘‘ لے کے اینوں (پہلے درجے دی)
\emph{ٹرم} ’’$\fregeext{x}{F(x)}$‘‘ بنا دیندی اے۔ ایس ترکیب نال فریگے
نے رکنیت دی ہِیٹھلی \emph{تعریف} دتی:
\[
	a \in b =_\text{df} \exists G(b = \fregeext{x}{G(x)} \land Ga)
\]
موٹے طور اُتے، $a \in b$ اُدوں تے صرف اُدوں جدوں $a$ ایسے تصور تھلّے آؤندا
ہووے جیہدا دائرہ $b$ اے۔ (دھیان دیو کہ مقداری نشان ’’$\exists G$‘‘ دوجے
درجے دا اے۔) فریگے نے ہِیٹھلا اصول وی منّیا، جیہنوں \emph{پنجواں بنیادی قانون}
آکھدے نیں:
$$\fregeext{x}{F(x)} = \fregeext{x}{G(x)} \liff \forall x (Fx \liff Gx)$$
موٹے طور اُتے، تصوراں دے دائرے اکّو اُدوں تے صرف اُدوں نیں جدوں تصور
عین اوہناں ای شےواں اُتے لاگو ہوندے نیں۔ (فیر، ’’$F$‘‘ تے ’’$G$‘‘
دونیں محمول دی تھاں نیں۔) ہُن اک سدھا اصول رکنیت نوں خاصیت پوری کرن نال
جوڑدا اے:

\begin{lem}[\emph{Grundgesetze} وچ]\ollabel{lem:Fregeextensions}
$\forall F \forall a(a \in \fregeext{x}{F(x)} \liff Fa)$
\end{lem}

\begin{proof}
$F$ تے $a$ مقرر کرو۔ ہُن $a \in \fregeext{x}{F(x)}$ اُدوں تے صرف اُدوں
جدوں $\exists G(\fregeext{x}{F(x)} = \fregeext{x}{G(x)} \land Ga)$
(رکنیت دی تعریف نال)، اُدوں تے صرف اُدوں جدوں
$\exists G(\forall x(Fx \liff Gx) \land Ga)$ (پنجویں بنیادی قانون نال)،
اُدوں تے صرف اُدوں جدوں $Fa$ (دوجے درجے دی بنیادی منطق نال)۔
\end{proof}

تے ایس توں سادہ خاصیت توں سیٹ بناؤن دا اصول لگ بھگ سدھا نکلدا اے:

\begin{lem}[\emph{Grundgesetze} وچ۔]
$\forall F \exists s \forall a (a \in s \liff Fa)$
\end{lem}

\begin{proof}
$F$ مقرر کرو۔ ہُن \olref{lem:Fregeextensions} توں
$\forall a (a \in \fregeext{x}{F(x)} \liff Fa)$ ملدا اے۔ سو وجودی
تعمیم نال $\exists s\forall a(a \in s \liff Fa)$۔ نتیجہ ایس لئی نکلدا اے
کہ $F$ کوئی وی ہو سکدا سی۔
\end{proof}

$F$ نوں $\forall x(Fx \liff x \notin x)$ مطابق لین نال رسل دا پیراڈوکس نکلدا اے۔
\endgroup
% /OLP-0536

\OLEndChapterHook
\endgroup
% /OLP-0530
% OLP-0537 content/set-theory/z/z.tex
\begingroup
\def\olfilename{z}
\expandafter\def\csname ol@sectioncs\endcsname{section}

\olchapter{sth}{z}{\texorpdfstring{$\Z$ ول قدم}{Z ول قدم}}

% OLP-0538 content/set-theory/z/story.tex
\begingroup
\def\olfilename{story}
\expandafter\def\csname ol@sectioncs\endcsname{section}

\olfileid{sth}{z}{story}
\olsection{سیٹ بنن دا بیان ہور تفصیل نال}

\olref[story][approach]{sec} وچ اسیں شون فیلڈ دا سیٹ بنن دی کارروائی
بارے قول دتا سی۔ ہُن اسیں کجھ ہور اصول لکھنے چاہندے آں تاں جو ایہ بیان
ذرا ہور ٹھیک ہو جاوے۔ اصول ایہ نیں:
\begin{enumerate}
\item[] \stageshier۔ ہر سیٹ کسے مرحلے اُتے بندا اے۔
\item[] \stagesord۔ مرحلے ترتیب وچ نیں، کجھ ہورناں توں \emph{پہلاں}
آؤندے نیں۔\footnote{اسیں اصل وچ بناں آکھے ای منّاں گے کہ مرحلے
\emph{حُسنِ ترتیب} وچ نیں۔ ایہدا مطلب \olref[ordinals][]{chap} وچ
سمجھایا اے۔ ایہ اک وڈا مفروضہ اے۔ اصل وچ \citet{Scott1974} دی اک
بڑی سُجھ والی ترکیب نال ایہ مفروضہ پہلاں \emph{چھڈیا} تے فیر
\emph{اخذ کیتا} جا سکدا اے۔ (ایس نال ایہ وی سمجھ آوے گی کہ سانوں
پہلا مرحلہ موجود کیوں سمجھنا چاہیدا اے۔) ایہ گل ایتھے نہیں کھول سکدے۔
ہور لئی \citet{ButtonLT1} ویکھو۔}
\item[] \stagesacc۔ کسے وی مرحلے $S$ لئی، تے مرحلے~$S$ توں \emph{پہلاں} بنے سیٹاں دی کسے وی چُناؤ لئی، مرحلے~$S$ اُتے اک
سیٹ بندا اے جیہدے رکن عین اوہی سیٹ نیں۔ مرحلے~$S$ اُتے ہور کجھ نہیں بندا۔
\end{enumerate}
ایہ غیر رسمی اصول نیں، پر اسیں ایہناں نال زرمَیلو دے سیٹ نظریے دے کئی
مسلّمیاں دا جواز دے سکاں گے۔

(اک گل دا دھیان دئیے۔ اسیں کجھ بالکل معیاری مسلّمے، بالکل معیاری ناںواں
نال، پیش کراں گے۔ پر ہُنے دتے ٹیرھے حرفاں والے اصولاں دے علمی لکھت وچ
کوئی خاص ناں نہیں۔ اسیں بس اوہ ناں ورتدے آں جیہڑے ساڈے خیال وچ کم آؤن گے۔)
\endgroup
% /OLP-0538
% OLP-0539 content/set-theory/z/separation.tex
\begingroup
\def\olfilename{separation}
\expandafter\def\csname ol@sectioncs\endcsname{section}

\olfileid{sth}{z}{sep}
\olsection{وکھ کرن دا اصول}

سادہ خاصیت توں سیٹ بناؤن دے اصول دی تھاں اسیں ایس اصول توں شروع کردے آں:

\begin{axiom}[وکھ کرن دا خاکہ] ہر فارمولے $\phi(x)$ لئی ایہ اک مسلّمہ
اے کہ کسے وی $A$ لئی سیٹ $\Setabs{x \in A}{\phi(x)}$ موجود اے۔
\end{axiom}

دھیان دیو کہ ایہ اکّو مسلّمہ نہیں۔ ایہ مسلّمیاں دا اک \emph{خاکہ} اے۔
وکھ کرن دے \emph{لامتناہی} مسلّمے نیں، ہر فارمولے $\phi(x)$ لئی اک۔
خاکہ ایویں وی لکھ سکدے آں (تے عام طور اُتے ایویں ای لکھدے آں):

\begin{defish}
ہر اوہ فارمولا $\phi(x)$ جیہدے وچ ’’$S$‘‘ نہیں آؤندا، اوہدے لئی ایہ اک
مسلّمہ اے:
\[
	\forall A \exists S \forall x(x \in S \liff (\phi(x) \land x \in A)).
\]
\end{defish}

\olref[sth][][]{part} دے شروع وچ دسے رواج مطابق، وکھ کرن دے مسلّمیاں
وچ فارمولیاں~$\phi$ دے پیرامیٹر ہو سکدے نیں۔\footnote{ایہدا مطلب
سمجھن لئی \olref[sfr][infinite][induction]{natinductionschema} توں
جھٹ پچھوں والی گل بات ویکھو۔}

سیٹاں دا اکٹھا کرن والا مرحلہ وار تصور، جیہڑا اسیں بیان کردے آئے آں،
وکھ کرن دا جواز سدھا دیندا اے۔ ویکھن لئی منّ لو $A$ اک سیٹ اے۔
سو $A$ کسے مرحلے~$S$ اُتے بندا اے (\stageshier مطابق)۔ چونکہ $A$
مرحلے~$S$ اُتے بنیا سی، ایس لئی $A$ دے سارے رکن مرحلے $S$ توں
پہلاں بنے سن (\stagesacc مطابق)۔ ہُن خاص کر کے اوہ سارے سیٹ ویکھو
جیہڑے $A$ دے رکن وی نیں تے $\phi$ نوں پورا وی کردے نیں۔ صاف اے
کہ ایہ سارے سیٹ وی مرحلے~$S$ توں پہلاں بنے سن۔ ایس لئی ایہ وی اک سیٹ
$\Setabs{x \in A}{\phi(x)}$ بناندے نیں، مرحلے~$S$ اُتے (\stagesacc مطابق)۔

سادہ خاصیت توں سیٹ بناؤن دے اصول دے اُلٹ، ایہ رسل دے پیراڈوکس توں
بچدا اے۔ اسیں سیٹ $\Setabs{x}{x \notin x}$ دی موجودگی سدھی نہیں آکھ
سکدے۔ پر جے سانوں کوئی سیٹ~$A$ \emph{دتا ہووے}، تاں اسیں سیٹ
$R_A = \Setabs{x \in A}{x \notin x}$ دی موجودگی آکھ سکدے آں۔ پر ایس توں
اینا ای ثابت ہوندا اے کہ $R_A \notin R_A$ تے $R_A \notin A$۔ ایہناں وچ
کوئی بہتی اوکڑ والی گل نہیں۔

پر وکھ کرن دا اک سدھا تے چونکا دین والا نتیجہ اے:

\begin{thm}\ollabel{thm:NoUniversalSet}
کوئی \emph{ہمہ گیر} سیٹ نہیں، یعنی $\Setabs{x}{x = x}$ موجود نہیں۔
\end{thm}

\begin{proof}
تضاد کڈھن لئی منّ لو $V$ اک ہمہ گیر سیٹ اے۔ تاں وکھ کرن دے اصول نال
$R = \Setabs{x \in V}{x \notin x} = \Setabs{x}{x \notin x}$ موجود اے۔
ایہ رسل دے پیراڈوکس دے اُلٹ اے۔
\end{proof}

ہمہ گیر سیٹ دا نہ ہونا، بلکہ سیٹاں دی درجہ بندی دا کوئی آخری سرا نہ ہونا،
اکٹھا کرن والے مرحلہ وار تصور دیاں سبھ توں بنیادی گلاں وچوں اک اے۔
ایس لئی ایہ ویکھنا چنگا اے کہ سمجھ مطابق اسیں ہور راہ توں وی ایس گل
تک پہنچ سکدے آں۔ ہمہ گیر سیٹ لازماً اپنے آپ دا !!a{element} ہووے گا۔
پر ساڈے اکٹھا کرن والے مرحلہ وار تصور وچ ہر سیٹ درجہ بندی وچ پہلی واری
اوس پہلے مرحلے اُتے آؤندا اے جیہڑا اوہدے سارے !!{element}s توں جھٹ
پچھوں اے۔ ایس توں نکلدا اے کہ \emph{کوئی وی} سیٹ اپنے آپ دا رکن نہیں۔
اپنے آپ دا رکن سیٹ پہلی واری اوس مرحلے توں جھٹ پچھوں آؤنا سی جیہدے
اُتے اوہ پہلی واری آیا، تے ایہ بے تُکی گل اے۔ (\olref[sth][spine][rank]{defnsetrank}
وچ ویکھاں گے کہ سیٹ دے ’’درجے‘‘ دا خیال ورت کے ایس گل نوں ہور پکا
کیویں کر سکدے آں۔ پر ایہ کرن لئی سانوں پہلاں کجھ ہور مسلّمے چاہیدے ہون گے۔)

وکھ کرن تے عنصراں نال تعیّن دے کجھ ہور نتیجے ہِیٹھاں نیں۔

\begin{prop}\ollabel{prop:emptyexists}
جے کوئی سیٹ موجود اے، تاں $\emptyset$ موجود اے۔
\end{prop}

\begin{proof}
جے $A$ اک سیٹ اے، تاں $\emptyset = \Setabs{x \in A}{x \neq x}$
وکھ کرن دے اصول نال موجود اے۔
\end{proof}

\begin{prop}
$A \setminus B$ کسے وی دو سیٹاں $A$ تے $B$ لئی موجود اے۔
\end{prop}

\begin{proof}
$A \setminus B = \Setabs{x \in A}{x \notin B}$ وکھ کرن دے اصول نال موجود اے۔
\end{proof}

ایہ وی پتا لگدا اے کہ (تقریباً) کسے وی خاندان دے اشتراک موجود نیں:

\begin{prop}\ollabel{prop:intersectionsexist}
جے $A \neq \emptyset$، تاں $\bigcap A = \Setabs{x}{(\forall y \in A)x \in y}$
موجود اے۔
\end{prop}

\begin{proof}
منّ لو $A \neq \emptyset$، سو کوئی $c \in A$ اے۔ فیر $\bigcap A =
\Setabs{x}{(\forall y \in A)x \in y} = \Setabs{x \in c}{(\forall y \in
A)x \in y}$، جیہڑا وکھ کرن دے اصول نال موجود اے۔
\end{proof}

پر شرط $A \neq \emptyset$ دا دھیان دیو۔ کیوں جے $\bigcap \emptyset$
خالی شرط کارن ہمہ گیر سیٹ ہووے گا، جیہڑا \olref{thm:NoUniversalSet}
دے اُلٹ اے۔
\endgroup
% /OLP-0539
% OLP-0540 content/set-theory/z/union.tex
\begingroup
\def\olfilename{union}
\expandafter\def\csname ol@sectioncs\endcsname{section}

\olfileid{sth}{z}{union}
\olsection{اتحاد}

\olref[sth][z][sep]{prop:intersectionsexist} توں سانوں اشتراک ملے سن۔
پر جے اسیں چاہندے آں کہ کسے وی خاندان دے اتحاد موجود ہون، تاں اک ہور
مسلّمہ دینا پئے گا:

\begin{axiom}[اتحاد] ہر سیٹ $A$ لئی سیٹ $\bigcup A =
\Setabs{x}{(\exists b \in A) x \in b}$ موجود اے۔
\[
	\forall A \exists U \forall x(x \in U \liff (\exists b \in A)x \in b)
\]
\end{axiom}

اکٹھا کرن والا مرحلہ وار تصور ایس مسلّمے دا وی جواز دیندا اے۔ منّ لو
$A$ اک سیٹ اے، سو $A$ کسے مرحلے $S$ اُتے بندا اے (\stageshier مطابق)۔
$A$ دا ہر رکن $S$ توں \emph{پہلاں} بنیا سی (\stagesacc مطابق)۔ ایسے
طرح دلیل دیئیے تاں $A$ دے ہر رکن دا ہر رکن $S$ توں پہلاں بنیا سی۔
ایس لئی \emph{اوہ} سارے سیٹ $S$ توں پہلاں موجود سن، تاں جو $S$
اُتے اوہناں توں اک سیٹ بن سکے۔ تے اوہ سیٹ ای $\bigcup A$ اے۔
\endgroup
% /OLP-0540
% OLP-0541 content/set-theory/z/pairs.tex
\begingroup
\def\olfilename{pairs}
\expandafter\def\csname ol@sectioncs\endcsname{section}

\olfileid{sth}{z}{pairs}
\olsection{جوڑیاں}

اگلا مسلّمہ جیہڑا ویکھنا اے، ایہ اے:

\begin{axiom}[جوڑیاں]
کسے وی دو سیٹاں $a, b$ لئی سیٹ $\{a, b\}$ موجود اے۔
\[
	\forall a \forall b \exists P \forall x (x \in P \liff (x = a \lor x = b))
\]
\end{axiom}

مرحلہ وار تصور نال ایس مسلّمے دا جواز ایویں دیندے آں۔ منّ لو $a$
مرحلے $S$ اُتے تے $b$ مرحلے $T$ اُتے موجود اے۔ منّ لو $M$، مرحلیاں
$S$ تے $T$ وچوں جیہڑا پچھوں آؤندا اے، اوہ اے۔ چونکہ $a$ تے $b$ دونیں
مرحلے $M$ اُتے موجود نیں، ایس لئی سیٹ $\{a,b\}$ اک ممکن اکٹھ اے جیہڑا
$M$ توں پچھوں آؤن والے کسے وی مرحلے اُتے دستیاب اے (جیہڑا وی وڈا ہووے)۔

پر ذرا ٹھہرو!{} اسیں کیوں منّیئے کہ $M$ توں پچھوں کوئی مرحلے
\emph{موجود نیں}؟ جے کوئی نہیں تاں ساڈا جواز نہیں چلے گا۔ سو جوڑیاں دا
جواز دین لئی سانوں \olref[sth][z][story]{sec} دے بیان وچ اک ہور
اصول شامل کرنا پئے گا:
\begin{enumerate}
	\item[] \stagessucc۔ کوئی آخری مرحلہ نہیں۔
\end{enumerate}
کی ایس اصول دا جواز اے؟ شون فیلڈ دے بیان وچ ایہ گل \emph{صاف لفظاں وچ}
نہیں کہی گئی سی کہ کوئی آخری مرحلہ نہیں۔ فیر وی، جے ایہ (ٹھیک ٹھیک
آکھیے تاں) ساڈے بیان وچ اک وادھا اے، ایہ ایس بنیادی خیال نال میل کھاندا
اے کہ سیٹ مرحلیاں وچ بندے نیں۔ اگے اسیں اینوں منّ لواں گے۔ سو جوڑیاں دا
مسلّمہ وی منّ لواں گے۔

ایس نویں مسلّمے نال اسیں کئی ہور سیٹاں دی موجودگی ثابت کر سکدے آں۔ مثلاً:

\begin{prop}\ollabel{prop:pairsconsequences}
کسے وی دو سیٹاں $a$ تے $b$ لئی ہِیٹھلے سیٹ موجود نیں:
	\begin{enumerate}
		\item\ollabel{singleton} $\{a\}$
		\item\ollabel{binunion} $a \cup b$
		\item\ollabel{tuples} $\tuple{a, b}$
	\end{enumerate}
\end{prop}

\begin{proof}
\olref{singleton}۔ جوڑیاں دے اصول نال $\{a, a\}$ موجود اے، تے عنصراں
نال تعیّن دے اصول نال ایہ $\{a\}$ اے۔

\olref{binunion}۔ جوڑیاں دے اصول نال $\{a, b\}$ موجود اے۔
ہُن $a \cup b = \bigcup \{a, b\}$ اتحاد دے اصول نال موجود اے۔

\olref{tuples}۔ \olref{singleton} نال $\{a\}$ موجود اے۔ جوڑیاں دے
اصول نال $\{a, b\}$ موجود اے۔ ہُن $\{\{a\}, \{a, b\}\} = \tuple{a, b}$
وی جوڑیاں دے اصول نال موجود اے۔
\end{proof}

\begin{prob}
وکھاؤ کہ کسے وی سیٹاں $a, b, c$ لئی سیٹ $\{a, b, c\}$ موجود اے۔
\end{prob}

\begin{prob}
وکھاؤ کہ کسے وی سیٹاں $a_1, \ldots, a_n$ لئی سیٹ $\{a_1, \ldots, a_n\}$
موجود اے۔
\end{prob}

%\begin{proof} By two applications of Pairs, $\{\{a_1, a_2\}, \{a_1,
%   a_3\}\}$ exists. By Union and Extensionality, $\bigcup \{\{a_1,
%   a_2\}, \{a_1, a_3\}\} = \{a_1, a_2, a_3\}$ exists. Repeat this
%   trick as often as necessary. \end{proof}
\endgroup
% /OLP-0541
% OLP-0542 content/set-theory/z/powerset.tex
\begingroup
\def\olfilename{powerset}
\expandafter\def\csname ol@sectioncs\endcsname{section}

\olfileid{sth}{z}{power}
\olsection{ذیلی سیٹاں دے سیٹ}

اسیں اک ہور مسلّمے ول ودھدے آں:

\begin{axiom}[ذیلی سیٹاں دے سیٹ]
ہر سیٹ $A$ لئی سیٹ $\Pow{A} = \Setabs{x}{x \subseteq A}$ موجود اے۔
\[
	\forall A \exists P \forall x(x \in P \liff (\forall z \in x)z \in A)
\]
\end{axiom}

ایہدا جواز دینا بڑا سدھا اے۔ منّ لو $A$ مرحلے $S$ اُتے بنیا۔
تاں $A$ دے سارے رکن $S$ توں پہلاں موجود سن (\stagesacc مطابق)۔
سو وکھ کرن دا جواز دیندیاں جیہڑی دلیل دتی سی، اوہدے وانگ، $A$ دا ہر
ذیلی سیٹ مرحلے $S$ تک بن جاندا اے۔ سو $S$ توں پچھوں کسے وی مرحلے اُتے
اوہ سارے اکّو سیٹ بناؤن لئی موجود نیں۔ تے سانوں پتا اے کہ کوئی ایسا
مرحلہ اے، کیوں جے $S$ آخری مرحلہ نہیں (\stagessucc مطابق)۔ سو $\Pow{A}$
موجود اے۔

ذیلی سیٹاں دے سیٹ دے اصول دا اک چنگا نتیجہ ایہ اے:

\begin{prop}\label{thm:Products}
کسے وی دو سیٹاں $A, B$ لئی اوہناں دا کارتیسی حاصل ضرب $A \times B$ موجود اے۔
\end{prop}

\begin{proof}
سیٹ $\Pow{\Pow{A \cup B}}$، ذیلی سیٹاں دے سیٹ دے اصول تے
\olref[sth][z][pairs]{prop:pairsconsequences} نال موجود اے۔ سو وکھ کرن
دے اصول نال ایہ سیٹ موجود اے:
\[
	C = \Setabs{z \in \Pow{\Pow{A \cup B}}}{(\exists x \in A)(\exists y \in B) z = \tuple{x, y}}.
\]
ہُن کسے وی $x \in A$ تے $y \in B$ لئی سیٹ $\tuple{x, y}$،
\olref[sth][z][pairs]{prop:pairsconsequences} نال موجود اے۔
نالے چونکہ $x, y \in A \cup B$، ایس لئی $\{x\}, \{x, y\} \in \Pow{A \cup B}$
تے $\tuple{x,y} \in \Pow{\Pow{A \cup B}}$۔ سو $A \times B = C$۔
\end{proof}

ایس ثبوت وچ طاقت والے سیٹاں دا اصول وکھ کرن دے اصول نال رل کے چلدا
اے۔ ایہ حیرانی دی گل نہیں۔ وکھ کرن توں بناں طاقت والے سیٹاں دا اصول
بہتا \emph{طاقتور} نہ ہوندا۔ آخر وکھ کرن دا اصول ای دسدا اے کہ کسے
سیٹ دے کیہڑے ذیلی سیٹ موجود نیں، تے ایس طرح ایہ طے کردا اے کہ ہر
ذیلی سیٹاں دا سیٹ کِنّا ’’بھریا‘‘ ہویا اے۔

\begin{prob}
وکھاؤ کہ کسے وی دو سیٹاں $A, B$ لئی (i) دائرہ $A$ تے حاصل قدراں دا
سیٹ $B$ رکھن والے سارے تعلقاں دا سیٹ موجود اے، تے (ii) $A$ توں $B$ ول
سارے فنکشوناں دا سیٹ موجود اے۔
\end{prob}

\begin{prob}
منّ لو $A$ اک سیٹ اے تے $\sim$، $A$ اُتے اک ہم ارزیت دا تعلق اے۔
ثابت کرو کہ $\sim$ دے لحاظ نال $A$ دیاں ہم ارز جماعتاں دا سیٹ،
یعنی $\equivclass{A}{\sim}$، موجود اے۔
\end{prob}
\endgroup
% /OLP-0542
% OLP-0543 content/set-theory/z/infinity-again.tex
\begingroup
\def\olfilename{infinity-again}
\expandafter\def\csname ol@sectioncs\endcsname{section}

\olfileid{sth}{z}{infinity-again}
\olsection{لامتناہیت}

ساڈے کول پہلاں ای اینے مسلّمے نیں کہ جے کوئی سیٹ موجود اے، تاں لامتناہی
گنتی وچ سیٹ ہون گے۔ کیوں جے منّ لو کوئی سیٹ موجود اے، تاں
$\emptyset$ موجود اے (\olref[sth][z][sep]{prop:emptyexists} نال)۔
ہُن کسے وی سیٹ~$x$ لئی سیٹ $x \cup \{x\}$،
\olref[sth][z][pairs]{prop:pairsconsequences} نال موجود اے۔
سو ایہ گل کجھ واری لا کے سانوں ایہ سیٹ ملن گے:
\begin{enumerate}
	\item[0.] $\emptyset$ %& $0$ & stage $0$\\
	\item[1.] $ \{\emptyset\}$ %& $1$ & stage $1$\\
	\item[2.] $\{\emptyset, \{\emptyset\}\}$ %& $2$ & stage $2$\\
	\item[3.] $\{\emptyset, \{\emptyset\}, \{\emptyset, \{\emptyset\}\}\}$ %&$3$ & stage $3$\\
	\item[4.] $\{\emptyset, \{\emptyset\}, \{\emptyset, \{\emptyset\}\}, \{\emptyset, \{\emptyset\}, \{\emptyset, \{\emptyset\}\}\}\}$% & $4$ & stage $4$ 
\end{enumerate}
تے اسیں پرکھ سکدے آں کہ ایہناں وچوں ہر سیٹ وکھرا اے۔

اسیں گنتی $0$ توں کئی وجوہاں کر کے شروع کیتی اے۔ اوہناں وچوں اک ایہ اے۔
ایہ پرکھنا بہتا اوکھا نہیں کہ جیہڑے سیٹ نوں اسیں ’’$n$‘‘ دا نشان دتا اے،
اوہدے عین $n$ رکن نیں، تے سمجھ مطابق اوہ $n$ویں مرحلے اُتے بندا اے۔

پر ایس نال سانوں \emph{لامتناہی گنتی وچ} سیٹ ملدے نیں۔ ایہ ضمانت نہیں
ملدی کہ اک \emph{لامتناہی سیٹ} وی اے، یعنی اک سیٹ جیہدے رکن لامتناہی
گنتی وچ ہون۔ تے ایس نال واقعی فرق پیندا اے۔ جِنّی دیر سانوں اک (ڈیڈیکنڈ)
لامتناہی سیٹ نہیں لبھدا، اسیں ڈیڈیکنڈ الجبرا نہیں بنا سکدے۔ پر سانوں
ڈیڈیکنڈ الجبرا چاہیدا اے تاں جو اسیں اوہنوں قدرتی عدداں دا سیٹ منّ سکیے۔
(\olref[sfr][infinite][dedekindsproof]{sec} نال ملا کے ویکھو۔)

خاص گل ایہ اے کہ ہُن تک دتے مسلّمے کسے وی لامتناہی سیٹ دی موجودگی دی
ضمانت \emph{نہیں} دیندے۔ سو سانوں اک نواں مسلّمہ دینا پئے گا:

\begin{axiom}[لامتناہیت]
اک سیٹ $I$ موجود اے جیہدے لئی $\emptyset \in I$ تے $x \cup \{x\} \in I$ ہر اوس صورت وچ جتھے $x \in I$۔
\begin{align*}
	\exists I( & (\exists o \in I)\forall x\ x \notin o \land {}\\
	& (\forall x \in I)(\exists s \in I)\forall z(z \in s \liff (z \in x \lor z = x)))
\end{align*}
\end{axiom}

ایہ ویکھنا سوکھا اے کہ لامتناہیت دے مسلّمے توں ملیا سیٹ $I$ ڈیڈیکنڈ
لامتناہی اے۔ اوہدا خاص رکن $\emptyset$ اے تے $I$ اُتے یک بہ یک نقشہ
$s(x) = x\cup \{x\}$ اے۔ ہُن \olref[sfr][infinite][dedekind]{thm:DedekindInfiniteAlgebra}
نے وکھایا کہ ڈیڈیکنڈ لامتناہی سیٹ توں ڈیڈیکنڈ الجبرا کیویں ملدا اے۔ اسیں
ایہنوں قدرتی عدداں دا اپنا سیٹ منّاں گے۔ ہور ٹھیک لفظاں وچ:
\begin{defn}\ollabel{defnomega}
منّ لو $I$ لامتناہیت دے مسلّمے توں ملیا کوئی وی سیٹ اے۔ منّ لو $s$
فنکشن $s(x) = x \cup \{x\}$ اے۔ منّ لو $\omega = \closureofunder{s}{\emptyset}$۔
اسیں $\omega$ دے رکناں نوں \emph{قدرتی عدد} آکھدے آں، تے آکھدے آں کہ $n$
اوہ نتیجہ اے جیہڑا $n$ واری $s$ نوں $\emptyset$ اُتے لاؤن نال ملدا اے۔
\end{defn}

ہُن تسی پچھاں ویکھ کے پرکھ سکدے او کہ چند پیرا پہلاں ’’$n$‘‘ دا نشان
دتا سیٹ، عدد $n$ \emph{دے طور اُتے} ورتیا جائے گا۔

اسیں ایس مقرر کیتی گل دی اہمیت \olref[sth][z][nat]{sec} وچ ویکھاں گے۔
ہُن ایس نال اسیں سمجھ آؤن والا اک نتیجہ ثابت کر سکدے آں:

\begin{prop}\ollabel{naturalnumbersarentinfinite}
کوئی قدرتی عدد ڈیڈیکنڈ لامتناہی نہیں۔
\end{prop}

\begin{proof}
ثبوت استقرا نال اے، یعنی \olref[sfr][infinite][induction]{thm:dedinfiniteinduction} نال۔
صاف اے کہ $0 = \emptyset$ ڈیڈیکنڈ لامتناہی نہیں۔ استقرائی قدم لئی اسیں
عکسِ نقیض ثابت کراں گے۔ جے (بے تُکے مفروضے مطابق) $s(n)$ ڈیڈیکنڈ لامتناہی
اے، تاں $n$ ڈیڈیکنڈ لامتناہی اے۔

سو منّ لو $s(n)$ ڈیڈیکنڈ لامتناہی اے۔ یعنی اک !!{injection} $f$ اے
جیہدے لئی $\ran{f}\subsetneq \dom{f} = s(n) = n \cup \{n\}$۔
دونیں صورتاں ویکھدے آں۔

\emph{پہلی صورت۔ $n \notin \ran{f}$۔} سو $\ran{f} \subseteq n$ تے
$f(n) \in n$۔ منّ لو $g = \funrestrictionto{f}{n}$۔ ہُن
$\ran{g} = \ran{f} \setminus \{f(n)\} \subsetneq n = \dom{g}$۔
ایس لئی $n$ ڈیڈیکنڈ لامتناہی اے۔

\emph{دوجی صورت۔ $n \in \ran{f}$۔} $m \in \dom{f} \setminus \ran{f}$ مقرر
کرو تے فنکشن $h$ دی تعریف دائرہ $s(n) = n \cup \{n\}$ اُتے کرو:
\[
h(x) =
	\begin{cases}
		f(x) & \text{جے }f(x) \neq n\\
		m & \text{جے }f(x)=n
	\end{cases}
\]
سو $h$ تے $f$ ہر تھاں ملدے نیں، بس $h(f^{-1}(n)) = m \neq n = f(f^{-1}(n))$۔
چونکہ $f$ اک !!a{injection} اے، ایس لئی $n \notin \ran{h}$ تے
$\ran{h} \subsetneq \dom{h} = s(n)$۔ ہُن پہلی صورت والی دلیل نال
$n$ ڈیڈیکنڈ لامتناہی اے۔
\end{proof}

پر سوال رہ جاندا اے کہ اسیں لامتناہیت دے مسلّمے دا \emph{جواز} کیویں دیئیے۔
مختصر جواب ایہ اے کہ ساڈے ہُن تک دے بیان وچ اک ہور اصول شامل کرنا پئے گا۔
اصول ایہ اے:
\begin{enumerate}
	\item[] \stagesinf۔ اک لامتناہی مرحلہ اے۔ یعنی اک ایہو جیہا مرحلہ اے جیہڑا (ا) پہلا مرحلہ نہیں، (ب) جیہدے توں پہلاں کجھ مرحلے نیں، پر (ج) جیہدے توں عین پہلاں کوئی مرحلہ نہیں۔
\end{enumerate}
لامتناہیت دا مسلّمہ ایس اصول توں سدھا نکلدا اے۔ سانوں پتا اے کہ قدرتی عدد
$n$ مرحلے $n$ اُتے بندا اے۔ سو سیٹ $\omega$ پہلے لامتناہی مرحلے اُتے
بندا اے۔ تے $\omega$ آپ لامتناہیت دے مسلّمے دا گواہ اے۔

پر ایس نال اسیں واپس ایس سوال ول آ جاندے آں کہ \stagesinf دا جواز
کیویں دیئیے۔ \stagessucc وانگ ایہ ساڈے اکٹھا کرن والے مرحلہ وار درجیاں
دے بیان دا صاف حصہ نہیں سی۔ ایس توں وی ودھ گل ایہ اے کہ مرحلہ وار درجیاں
دے خیال وچ، جیہدے وچ سیٹ واری واری مرحلیاں اُتے بندے نیں، کوئی گل سانوں
ایہ منّن لئی مجبور نہیں کردی کہ کوئی \emph{لامتناہی} مرحلہ اے۔ ایہ سوچنا
بالکل ٹھیک بیٹھدا اے کہ مرحلے قدرتی عدداں وانگ ترتیب وچ نیں۔

پر ایس نال اک صاف مسئلہ پیدا ہوندا اے۔ \olref[sfr][infinite][dedekindsproof]{sec}
وچ اسیں ڈیڈیکنڈ دا ’’ثبوت‘‘ ویکھیا سی کہ (خیالاں دا) اک ڈیڈیکنڈ لامتناہی
سیٹ اے۔ ہو سکدا اے ایہ تہانوں تسلی بخش نہ لگیا ہووے۔ پر جے سیٹ دا مرحلہ
وار تصور (یا ’’سوچ دے قانون‘‘) سانوں \stagesinf{} منّن لئی ’’مجبور نہیں
کردے‘‘، تاں ڈیڈیکنڈ لامتناہی سیٹ دی موجودگی دے دعوے دا اپنا اندرونی جواز
ساڈے کول ہجے وی نہیں۔

بے شک ایتھے کہن نوں ہور بہت کجھ اے۔ پر آس اے کہ ہُن تسی سوچنا شروع کر
سکدے او کہ کسے مسلّمے (یا اصول) دا جواز دین لئی کیہدی \emph{لوڑ} پَندی اے۔
اگے اسیں \stagesinf{} نوں بس منّ لواں گے۔
\endgroup
% /OLP-0543
% OLP-0544 content/set-theory/z/milestone.tex
\begingroup
\def\olfilename{milestone}
\expandafter\def\csname ol@sectioncs\endcsname{section}

\olfileid{sth}{z}{milestone}
\olsection{\texorpdfstring{$\Z^-$۔ اک اہم پڑاؤ}{Z⁻۔ اک اہم پڑاؤ}}

اسیں اگلے حصے وچ \stagesinf{} ول مُڑاں گے۔ پر لامتناہیت دے مسلّمے
نال اسیں اک اہم پڑاؤ اُتے پہنچ گئے آں۔ ہُن ساڈے کول نظریے $ \Zminus$
لئی لوڑیندے سارے مسلّمے نیں۔ تفصیل ایہ اے:

\begin{defn}
نظریے $ \Zminus$ دے مسلّمے ایہ نیں۔ عنصراں نال تعیّن، اتحاد، جوڑیاں،
ذیلی سیٹاں دے سیٹ، لامتناہیت، تے وکھ کرن دے خاکے دیاں ساریاں صورتاں۔
\end{defn}

ناں توں مراد \emph{زرمَیلو} دا سیٹ نظریہ اے (\emph{گھٹاؤ} اوہ شے
جیہڑی اسیں اگے ویکھاں گے)۔ زرمَیلو ایس عزت دا حق دار اے، کیوں جے اوہنے
اپنے \citeyear{Zermelo1908Untersuchungen} دے کم وچ بنیادی طور اُتے ایہو
نظریہ دتا سی۔\footnote{تاریخ تے فنی گلاں اُتے دلچسپ ٹپّنیاں لئی
\citet[Appendix
A]{Potter2004} ویکھو۔}

ایہ نظریہ اینا طاقتور اے کہ اسیں ایس نال بہت وڈی گنتی وچ ریاضی کر سکدے آں۔
خاص طور اُتے تہانوں \emph{چاہیدا اے} کہ \olref[sfr][][]{part} ول مُڑ کے
ویکھو تے اپنے آپ نوں یقین دواؤ کہ جیہڑا کم اسیں سادے طریقے نال کیتا سی،
اوہ~$\Zminus$ اندر ہور رسمی طور اُتے کیتا جا سکدا سی۔ (کجھ چِر ایہ
کرن توں پچھوں ہو سکدا اے تسی اگے جا کے \olref[sth][z][arbintersections]{sec}
پڑھنا چاہو۔) سو ہُن توں اسیں ہور کجھ کہے بناں ایہ منّاں گے کہ اسیں
$ \Zminus$ وچ (گھٹ توں گھٹ) کم کر رہے آں۔
\endgroup
% /OLP-0544
% OLP-0545 content/set-theory/z/nat.tex
\begingroup
\def\olfilename{nat}
\expandafter\def\csname ol@sectioncs\endcsname{section}

\olfileid{sth}{z}{nat}
\olsection{اپنے قدرتی عدد چننے}

%Let us take it for granted that the informal principle \stagesinf{}
%can be justified. It is also, though, worth commenting on the
%particular axiom we extracted from that informal principle. 

\olref[infinity-again]{defnomega} وچ اسیں ’’قدرتی عدد‘‘ دی عبارت دی صاف
تعریف دتی۔ تسی ایس مقرر کیتی گل نوں کیویں سمجھو؟ ایہ کوئی وجودیات بارے
دعویٰ نہیں، بس کجھ سیٹاں نوں قدرتی عدد \emph{منّن} دا فیصلہ اے۔ ایس خیال
دیاں وجوہاں نوں اسیں \olref[sfr][rel][ref]{sec}، \olref[sfr][arith][ref]{sec}
تے \olref[sfr][infinite][dedekindsproof]{sec} وچ چھوہیا سی۔ پر اسیں
ایہ وجوہاں ہور وی صاف کر سکدے آں۔

ساڈا لامتناہیت دا مسلّمہ \citet{VonNeumann1925} دے کم مطابق اے۔
پر ہِیٹھلا ہور مسلّمہ اے، جیہڑا اسیں اوہدی تھاں اپنا سکدے ساں:

\begin{defish}
\emph{زرمَیلو دا \citeyear{Zermelo1908Untersuchungen} لامتناہیت دا مسلّمہ۔}
اک سیٹ $A$ اے جیہدے لئی $\emptyset \in A$ تے
$(\forall x \in A)\{x\} \in A$۔
\end{defish}

جے اسیں اپنے (فون نوئمان توں لئے) لامتناہیت دے مسلّمے دی تھاں زرمَیلو دا
مسلّمہ ورتیا ہوندا، تاں سانوں اویویں اک ڈیڈیکنڈ لامتناہی سیٹ تے ایس توں
اک ڈیڈیکنڈ الجبرا مل جاندا۔ زرمَیلو دے طریقے وچ وی ساڈے الجبرے دا خاص
رکن $\emptyset$ ہوندا (ساڈے $0$ دی تھاں)۔ پر یک بہ یک نقشہ
$x \mapsto \{x\}$ ہوندا، $x \mapsto x \cup \{x\}$ نہیں۔
ایہدا سبھ توں سدھا نتیجہ ایہ اے کہ زرمَیلو $2$ نوں $\{\{\emptyset\}\}$
منّدا اے، جدکہ اسیں (فون نوئمان دے نال) $2$ نوں $\{\emptyset, \{\emptyset\}\}$
منّ دے آں۔

لامتناہیت دے اک مسلّمے نوں دوجے دی تھاں کیوں چنیے؟ وڈی عملی وجہ ایہ اے
کہ فون نوئمان دا طریقہ لامتناہی درجیاں دے عدداں نوں سنبھالن لئی سوکھ نال
ودھایا جا سکدا اے۔ اسیں \olref[ordinals][]{chap} توں اگے ایہ گل ویکھاں گے۔
پر صرف \emph{حساب کرن} دے لحاظ نال دونیں طریقے اکّو جِنّا کم کرن گے۔
سو جے کوئی تہانوں آکھے کہ قدرتی عدد سیٹ \emph{ای نیں}، تاں ظاہر سوال اے:
\emph{اوہ کیہڑے سیٹ نیں؟}

ایس عین سوال نوں \citet{Benacerraf1965} نے مشہور کیتا۔ پر ایہ گل زور دے
کے کہنی چاہیدی اے کہ ایہ اوس گل دی بس سبھ توں مشہور مثال اے جیہڑی اسیں
پہلاں کئی واری ویکھ چکے آں۔ بنیادی گل ایہ اے کہ سیٹ نظریہ سانوں کئی
’’سمجھ آؤن والیاں‘‘ قسماں دیاں شےواں دی \emph{نقل بناؤن} دا طریقہ دیندا اے۔
حقیقی، ناطق، صحیح تے قدرتی عدد تاں نیں ای، پر مرتب جوڑیاں، فنکشن تے تعلق وی۔
پر سیٹ نظریہ سانوں نقل دا \emph{اکّو اک} طریقہ کدی نہیں دیندا۔
\emph{ہمیشہ} ہور طریقے ہوندے نیں جیہڑے سیدھی طرح ساڈے اوہناں ای کم آ سکدے سن۔
\endgroup
% /OLP-0545
% OLP-0546 content/set-theory/z/arbintersections.tex
\begingroup
\def\olfilename{arbintersections}
\expandafter\def\csname ol@sectioncs\endcsname{section}

\olfileid{sth}{z}{arbintersections}
\olsection{ضمیمہ۔ بندش، خاصیت توں سیٹ بنانا تے اشتراک}

\olref[sth][z][milestone]{sec} وچ اسیں آکھیا سی کہ تسی
\olref[sfr][][]{part} دے سادے طریقے نال کیتے کم ول مُڑ کے ویکھو تے
پرکھو کہ اوہ~$\Zminus$ وچ کیتا جا سکدا اے۔ جے تسی ایویں کیتا، تاں
\emph{اک} گل وچ اٹک سکدے سی۔ ڈیڈیکنڈ دے \emph{بندشاں} والے بیان وچ
\emph{اشتراک} دی ورتوں۔

\olref[sfr][infinite][dedekind]{Closure} توں یاد کرو کہ
\begin{align*}
  \closureofunder{f}{o} & = \bigcap\Setabs{X}{o \in X \text{ تے $X$، $f$ تحت
بند اے}}.
\intertext{ایہدی عام صورت ایہو جیہی تعریف اے:}
  C & = \bigcap\Setabs{X}{\phi(X)}.
\end{align*}
پر ایس گل نال سانوں خبردار ہو جانا چاہیدا اے۔ سادہ خاصیت توں سیٹ بناؤن
دا اصول نہیں چلدا، سو کوئی ضمانت نہیں کہ $\Setabs{X}{\phi(X)}$ موجود اے۔
فیر ایویں لگدا اے کہ ایہو جیہیاں تعریفاں وچ شاید اسیں \emph{ہیرا پھیری}
کر رہے آں۔

چنگی گل ایہ اے کہ ایہ ہیرا پھیری نہیں۔ یا، جے ایہ ایس صورت وچ
\emph{واقعی} ہیرا پھیری اے، تاں اسیں تھوڑی جیہی سچّی محنت نال ایہناں نوں ٹھیک کر
سکدے آں۔ اوس محنت دی جھلک \olref[sth][z][sep]{prop:intersectionsexist}
وچ سی، جدوں اسیں سمجھایا سی کہ $\bigcap A$ کسے وی $A \neq \emptyset$
لئی کیوں موجود اے۔ پر ہُن اسیں ایہ گل کھول کے دساں گے۔

عنصراں نال تعیّن دے اصول نال، جے اسیں $C$ نوں
$\bigcap\Setabs{X}{\phi(X)}$ دے طور اُتے تعریف کرن دی کوشش کریے،
تاں اصل وچ سانوں بس اک ایسی شے $C$ چاہیدی اے جیہڑی ایہ شرط پوری کرے:
\begin{equation}\ollabel{bicondelimarbintersection}
	\forall x(x \in C \liff \forall X(\phi(X) \lif x \in X))\tag{*}
\end{equation}
ہُن منّ لو \emph{کوئی} سیٹ $S$ اے جیہدے لئی $\phi(S)$ اے۔ فیر
\olref{bicondelimarbintersection} پوری کرن لئی اسیں $C$ نوں
\emph{وکھ کرن دے اصول} نال ایویں تعریف کر سکدے آں:
\[
	C = \Setabs{x \in S}{\forall X(\phi(X) \lif x \in X)}.
\]
ایہ پرکھنا مشق چھڈدے آں کہ ایس تعریف نال \olref{bicondelimarbintersection}
ملدی اے، جیویں سانوں چاہیدی سی۔
%  fix $x$. First suppose that $x \in C$, as (re)defined; then both $x
%  \in S$ and $\forall X(\phi(X) \lif x \in X)$; but since
%  $\phi(S)$, this reduces to the condition that $\forall X(\phi(X)
%  \lif x \in X)$. Conversely, suppose $\forall X(\phi(X) \lif
%  x \in X)$; then, since $\phi(S)$, we also have that $x \in S$; so
%  $x \in c$, as (re)defined. 
ایہ عام طریقہ سانوں اشتراکاں دی تعریف وچ سادہ خاصیت توں سیٹ بناؤن دے
اصول دی ہر ظاہری ورتوں توں بچا دیوے گا۔ جیہڑی خاص صورت نال ایہ گل
شروع ہوئی، یعنی $\closureofunder{f}{o}$، اوہدے لئی ایہ طریقہ ایویں چلے گا۔
اسیں \olref[sfr][infinite][dedekind]{closureproperties} دا ثبوت ایہ گل
ویکھ کے شروع کیتا سی کہ $o \in \ran{f}\cup\{o\}$ تے $\ran{f} \cup \{o\}$،
$f$ تحت بند اے۔ سو جیہڑی شے سانوں چاہیدی اے، اوہدی تعریف ایویں کر سکدے آں:
\[
	\closureofunder{f}{o} = \Setabs{x \in \ran{f} \cup \{o\}}{(\forall X \ni o)(X \text{ $f$ تحت بند اے} \lif x \in X)}.
\]
\endgroup
% /OLP-0546

\OLEndChapterHook
\endgroup
% /OLP-0537
% OLP-0547 content/set-theory/ordinals/ordinals.tex
\begingroup
\def\olfilename{ordinals}
\expandafter\def\csname ol@sectioncs\endcsname{section}

\olchapter{sth}{ordinals}{ترتیبی عدد}

% OLP-0548 content/set-theory/ordinals/introduction.tex
\begingroup
\def\olfilename{introduction}
\expandafter\def\csname ol@sectioncs\endcsname{section}

\olfileid{sth}{ordinals}{intro}

\olsection{جان پہچان}

\olref[z][]{chap} وچ اسیں \stagesinf{} دے روپ وچ فرض کیتا سی کہ
درجہ بندی دا اک لامتناہی مرحلہ اے (ساڈا لامتناہیت دا مسلّمہ وی ویکھو)۔
پر \stagessucc{} مطابق اسیں لامتناہی مرحلے اُتے رُک نہیں سکدے۔
سانوں اگے ودھنا پئے گا۔ سو پہلے لامتناہی مرحلے توں اگلے مرحلے اُتے
اسیں اوہناں سیٹاں دے سارے ممکن اکٹھ بناندے آں جیہڑے پہلے لامتناہی
مرحلے اُتے موجود سن۔ تے فیر ایہو کم کردے آں، تے فیر، تے فیر، \dots

ایس وچ ان کہی گل ایہ اے کہ اسیں عدد دا اک ’’سوجھ والا‘‘ تصور ورتنا
شروع کر دتا اے، جیہدے مطابق سارے قدرتی عدداں توں \emph{پچھوں} وی عدد
ہو سکدے نیں۔ خاص کر کے ایہ \emph{لامتناہی درجے دے ترتیبی عدد} دا خیال
اے۔ ایس باب دا مقصد ایس خیال نوں ہور ٹھیک بنانا اے۔ اسیں ترتیبی عدد
دا عام تصور ویکھاں گے تے فیر کجھ خاص سیٹاں نوں اپنے ترتیبی عدد دے طور
اُتے صاف تعریف کراں گے۔
\endgroup
% /OLP-0548
% OLP-0549 content/set-theory/ordinals/idea.tex
\begingroup
\def\olfilename{idea}
\expandafter\def\csname ol@sectioncs\endcsname{section}

\olfileid{sth}{ordinals}{idea}
\olsection{ترتیبی عدد دا عام خیال}

قدرتی عدداں نوں اوہناں دی عام ترتیب وچ ویکھو:
\begin{pnbDiagramLine}
	\begin{tikzpicture}
	\foreach \x/\xtext in {0, 1, 2, 3, 4, 5}
	{
		\node (\x a) at (\x, 1) {\small{$\x$}};
		\node (\x a) at (\x.5, 1) {\small{$<$}};
	}
	\node (ldots) at (6, 1) {\small{$\ldots$}};
	\end{tikzpicture}
\end{pnbDiagramLine}
فنی بولی وچ اینوں $\omega$ والی لڑی آکھدے آں۔ واقعی ایہ عام ترتیب سیٹاں
دی درجہ بندی دے مرحلیاں دی مُڈھلی بناوٹ وچ وی دِسدی اے۔ پر ہُن منّ لو اسیں
$0$ نوں لڑی دے آخری سرے اُتے لے جائیے، تاں جو اوہ سارے ہور عدداں توں
پچھوں آوے:
\begin{pnbDiagramLine}
	\begin{tikzpicture}
	\foreach \x/\xtext in {1, 2, 3, 4, 5}
	{
		\node (\x a) at (\x, 1) {\small{$\x$}};
		\node (\x a) at (\x.5, 1) {\small{$<$}};
	}
	\node (ldots) at (6, 1) {\small{$\ldots$}};
	\node (bea) at (6.5, 1) {\small{$<$}};
	\node (noa) at (7, 1) {\small{$0$}};
	\end{tikzpicture}
\end{pnbDiagramLine}
ایتھے شےواں اوہو نیں، پر اوہناں دی ترتیب بنیادی طور اُتے وکھری اے۔
پہلی ترتیب دا کوئی آخری عنصر نہیں سی، نویں ترتیب دا اے۔ واقعی نویں
ترتیب وچ شےواں دی اک $\omega$ والی لڑی ($1, 2, 3, 4, 5, \ldots$) اے
تے اوہدے پچھوں اک ہور شے اے۔ ایہ $\omega+1$ والی لڑی ہووے گی۔

صرف ایہو شےواں ورت کے اسیں ہور وی کئی قسم دیاں ترتیباں بنا سکدے آں۔
مثلاً سارے جفت عدد (اوہناں دی قدرتی ترتیب وچ)، تے اوہناں توں پچھوں
سارے طاق عدد (اوہناں دی قدرتی ترتیب وچ) ویکھو:
\begin{pnbDiagramLine}
	\begin{tikzpicture}
	\node(a1) at (1,1) {\small{$0$}};
	\node(a2) at (1.5,1) {\small{$<$}};
	\node(a3) at (2,1) {\small{$2$}};
	\node(a4) at (2.5,1) {\small{$<$}};
	\node(a5) at (3,1) {\small{$4$}};
	\node(a6) at (3.5,1) {\small{$<$}};
	\node(dots) at (4,1) {\small{$\ldots$}};
	\node(aaoe) at (4.5,1) {\small{$<$}};
	\node(b1) at (5,1) {\small{$1$}};
	\node(b2) at (5.5,1) {\small{$<$}};
	\node(b3) at (6,1) {\small{$3$}};
	\node(b4) at (6.5,1) {\small{$<$}};
	\node(b5) at (7,1) {\small{$\ldots$}};
	\end{tikzpicture}
\end{pnbDiagramLine}
ایہ اک $\omega$ والی لڑی اے جیہدے پچھوں دوجی $\omega$ والی لڑی اے،
یعنی $\omega+\omega$ والی لڑی۔

اسیں اگے ودھ سکدے آں۔ پر سانوں \emph{ترتیباں} دی ایس گل بات نوں
سمجھن دا اک عام طریقہ چاہیدا اے۔
\endgroup
% /OLP-0549
% OLP-0550 content/set-theory/ordinals/wo.tex
\begingroup
\def\olfilename{wo}
\expandafter\def\csname ol@sectioncs\endcsname{section}

\olfileid{sth}{ordinals}{wo}
\olsection{حُسنِ ترتیب}

بنیادی تصور ایہ اے:

\begin{defn}
تعلق $<$، $A$ نوں \emph{حُسنِ ترتیب وچ رکھدا اے} اُدوں تے صرف اُدوں، جدوں ایہ دو شرطاں پوری کرے:
\begin{enumerate}
	\item $<$ جُڑیا ہویا اے، یعنی ہر $a, b \in A$ لئی یا $a < b$، یا
	$a = b$، یا $b < a$۔
	\item $A$ دے ہر غیر خالی ذیلی سیٹ وچ $<$ دے لحاظ نال اک ایسا !!{element} جیہدے توں کوئی ہور رکن چھوٹا نہیں اے۔ یعنی جے $\emptyset \neq X \subseteq A$، تاں
	$(\exists m \in X)(\forall z \in X)z \nless m$۔
\end{enumerate}
\end{defn}

ایہ ویکھنا سوکھا اے کہ ہُنے ویکھیاں تنّے مثالاں واقعی حُسنِ ترتیب والے تعلق سن۔

\begin{prob}
\Olref[sth][ordinals][idea]{sec} وچ قدرتی عدداں دیاں تن ترتیباں مثال دے طور اُتے
دتیاں سن۔ پرکھو کہ ہر اک حُسنِ ترتیب اے۔
\end{prob}

حُسنِ ترتیب بارے کجھ بنیادی پر بہت ضروری گلاں ہِیٹھاں نیں۔

\begin{prop}\ollabel{wo:strictorder}
جے $<$، $A$ نوں حُسنِ ترتیب وچ رکھدا اے، تاں $A$ دے ہر غیر خالی ذیلی
سیٹ وچ $<$ دے لحاظ نال اکّو اک سبھ توں چھوٹا رکن اے، تے $<$ غیر انعکاسی،
غیر متقابل تے متعدی اے۔
\end{prop}

\begin{proof}
جے $X$، $A$ دا اک غیر خالی ذیلی سیٹ اے، تاں اوہدے وچ $<$ دے لحاظ نال
اک ایسا !!{element} $m$ اے جیہدے توں کوئی ہور رکن چھوٹا نہیں، یعنی $(\forall z \in X)z \nless m$۔
چونکہ $<$ جُڑیا ہویا اے، $(\forall z \in X)m \leq z$۔
سو $m$، $<$ دے لحاظ نال $X$ دا سبھ توں چھوٹا !!{element} اے۔

غیر انعکاسیت لئی $a \in A$ مقرر کرو۔ $<$ دے لحاظ نال $\{a\}$ دا سبھ توں
چھوٹا !!{element} $a$ اے، سو $a \nless a$۔ تعدّی لئی، جے $a < b < c$،
تاں چونکہ $\{a, b, c\}$ وچ $<$ دے لحاظ نال سبھ توں چھوٹا !!{element} اے،
ایس لئی $a < c$۔ غیر تقابل، غیر انعکاسیت تے تعدّی توں نکلدا اے۔
\end{proof}

\begin{prop}\ollabel{propwoinduction}
جے $<$، $A$ نوں حُسنِ ترتیب وچ رکھدا اے، تاں ہر فارمولے $\phi(x)$ لئی:
\[
	\text{ جے }(\forall a \in A)((\forall b < a)\phi(b) \lif
		\phi(a))\text{، تاں }(\forall a \in A)\phi(a).
\]
\end{prop}

\begin{proof}
اسیں عکسِ نقیض ثابت کراں گے۔ منّ لو $\lnot(\forall a \in A)\phi(a)$، یعنی
$X = \Setabs{x \in A}{\lnot\phi(x)} \neq \emptyset$۔ تاں $X$ دا
$<$ دے لحاظ نال اک ایسا !!{element}، $a$، اے جیہدے توں کوئی ہور رکن چھوٹا نہیں۔ سو
$(\forall b < a)\phi(b)$، پر $\lnot \phi(a)$۔ \end{proof}\noindent
ایس آخری خاصیت توں قدرتی عدداں اُتے مضبوط استقرا دا اصول یاد آؤنا چاہیدا اے۔
یعنی جے $(\forall n \in \omega)((\forall m < n)\phi(m) \lif \phi(n))$،
تاں $(\forall n \in \omega)\phi(n)$۔ ایہ خاصیت حُسنِ ترتیب نوں
اک بڑا \emph{مضبوط} تصور بناندی اے۔\footnote{یاد رہے کہ سارے فارمولیاں دے
پیرامیٹر ہو سکدے نیں (جدوں تک صاف طور اُتے کجھ ہور نہ آکھیا جاوے)۔}
\endgroup
% /OLP-0550
% OLP-0551 content/set-theory/ordinals/iso.tex
\begingroup
\def\olfilename{iso}
\expandafter\def\csname ol@sectioncs\endcsname{section}

\olfileid{sth}{ordinals}{iso}
\olsection{ترتیب دے ہم ساخت نقشے}

حُسنِ ترتیب \emph{کِس طرح} مضبوط اے، ایہ سمجھاؤن لئی اسیں پہلاں حُسنِ
ترتیب والیاں ساختاں دی آپس وچ پرکھ دا اک طریقہ دساں گے۔

\begin{defn}
\emph{حُسنِ ترتیب والی ساخت} اک جوڑی $\tuple{A, <}$ اے، جیہدے وچ $<$،
$A$ نوں حُسنِ ترتیب وچ رکھدا اے۔ حُسنِ ترتیب والیاں ساختاں $\tuple{A, <}$
تے $\tuple{B, \lessdot}$ \emph{ترتیب دے لحاظ نال ہم ساخت} اُدوں تے صرف اُدوں
نیں جدوں اک !!a{bijection} $f \colon A \to B$ ہووے جیہدے لئی
$x < y$ اُدوں تے صرف اُدوں اے جدوں $f(x) \lessdot  f(y)$۔
ایس صورت وچ اسیں $\ordeq{\tuple{A, <}}{\tuple{B, \lessdot}}$ لکھدے آں
تے $f$ نوں \emph{ترتیب دا ہم ساخت نقشہ} آکھدے آں۔
\end{defn}
\noindent
اگے مختصر گل کرن لئی اسیں ’’ترتیب دے ہم ساخت نقشیاں‘‘ دی تھاں ’’ہم ساخت
نقشیاں‘‘ آکھاں گے۔ سمجھ مطابق، ہم ساخت نقشے اوہ !!{bijection}s نیں
جیہڑے ساخت نوں قائم رکھدے نیں۔ ایہناں بارے کجھ سادیاں گلاں ہِیٹھاں نیں۔

\begin{lem}\ollabel{isoscompose}
ہم ساخت نقشیاں دے مرکب وی ہم ساخت نقشے نیں۔ یعنی جے $f \colon A \to B$
تے $g \colon B \to C$ ہم ساخت نقشے نیں، تاں $(g \circ f) \colon A \to C$
اک ہم ساخت نقشہ اے۔
\end{lem}
\begin{prob}
	\olref[sth][ordinals][iso]{isoscompose} ثابت کرو۔
\end{prob}
\begin{proof}
ایہ مشق چھڈّی اے۔
\end{proof}
\begin{cor}\ollabel{ordisoisequiv}
	$\ordeq{X}{Y}$ اک ہم ارزی دا تعلق اے۔
\end{cor}
\begin{prop}\ollabel{ordisounique}
جے $\tuple{A, <}$ تے $\tuple{B, \lessdot}$ حُسنِ ترتیب والیاں ہم ساخت
ساختاں نیں، تاں اوہناں وچکار اکّو اک ہم ساخت نقشہ اے۔
\end{prop}

\begin{proof}
منّ لو $f$ تے $g$، $A \to B$ دے ہم ساخت نقشے نیں۔ اسیں نتیجہ استقرا
نال ثابت کراں گے، یعنی \olref[wo]{propwoinduction} نال۔
$a\in A$ مقرر کرو تے استقرا لئی منّ لو کہ $(\forall b < a)f(b) = g(b)$۔
$x \in B$ مقرر کرو۔

جے $x \lessdot f(a)$، تاں $f^{-1}(x) < a$، سو
$g(f^{-1}(x)) \lessdot g(a)$، کیوں جے $f$ تے $g$ ہم ساخت نقشے نیں۔
پر چونکہ $f^{-1}(x) < a$، ساڈے مفروضے مطابق
$x =f(f^{-1}(x)) = g(f^{-1}(x))$۔ سو $x \lessdot g(a)$۔
ایسے طرح، جے $x \lessdot g(a)$، تاں $x \lessdot f(a)$۔

عام صورت وچ $(\forall x \in B)(x \lessdot f(a) \liff x \lessdot g(a))$۔
ایس توں $f(a) = g(a)$ نکلدا اے، \olref[sfr][rel][ord]{prop:extensionality-strictlinearorders}
نال۔ سو $(\forall a \in A)f(a) = g(a)$، \olref[wo]{propwoinduction} نال۔
\end{proof}\noindent
ایس توں کجھ اندازہ ہوندا اے کہ حُسنِ ترتیب والیاں ساختاں مضبوط نیں۔
پر ایس گل نوں اگے سمجھاؤن لئی کجھ ہور نشان دسنے کم آؤن گے۔

\begin{defn}
جدوں $\tuple{A, <}$ حُسنِ ترتیب والی ساخت اے تے $a \in A$، تاں منّ لو
$A_a = \Setabs{x \in A}{x < a}$۔ اسیں $A_a$ نوں $A$ دا حقیقی
\emph{مُڈھلا ٹکڑا} آکھدے آں (تے $A$ نوں وی $A$ دا اک غیر حقیقی مُڈھلا
ٹکڑا منّ دے آں)۔ منّ لو $<_a$، $<$ دا اوس مُڈھلے ٹکڑے تک محدود روپ اے،
یعنی $\funrestrictionto{\mathord{<}}{A_a^2}$۔
\end{defn}
\noindent
ایہ نشان ورت کے اسیں آکھ تے ثابت کر سکدے آں کہ کوئی حُسنِ ترتیب والی
ساخت اپنے کسے حقیقی مُڈھلے ٹکڑے نال ہم ساخت نہیں۔

\begin{lem}\ollabel{wellordnotinitial}
جے $\tuple{A, <}$ حُسنِ ترتیب والی ساخت اے تے $a \in A$، تاں
$\ordneq{\tuple{A, <}}{\tuple{A_a, <_a}}$۔
\end{lem}

\begin{proof}
تضاد کڈھن لئی منّ لو $f \colon A \to A_a$ اک ہم ساخت نقشہ اے۔
چونکہ $f$ اک یک بہ یک تے پورا نقشہ اے تے $A_a \subsetneq A$،
\olref[wo]{wo:strictorder} ورت کے $b \in A$ نوں $<$ دے لحاظ نال $A$ دا
سبھ توں چھوٹا !!{element} منّ لو جیہدے لئی $b \neq f(b)$۔ اسیں وکھاواں گے
کہ $(\forall x \in A)(x<b \liff x < f(b))$، جیہدے توں
\olref[sfr][rel][ord]{prop:extensionality-strictlinearorders} نال
$b = f(b)$ نکلے گا تے تضاد مل جائے گا۔

منّ لو $x < b$۔ سو $x = f(x)$، $b$ دی چُون کر کے۔ تے $f(x) < f(b)$،
کیوں جے $f$ اک ہم ساخت نقشہ اے۔ سو $x < f(b)$۔

منّ لو $x < f(b)$۔ سو $f^{-1}(x) < b$، کیوں جے $f$ اک ہم ساخت نقشہ اے،
تے ایس لئی $f^{-1}(x) = x$، $b$ دی چُون کر کے۔ سو $x < b$۔
\end{proof}

اگلا نتیجہ، موٹی گل آکھیے تاں، ایہ وکھاندا اے کہ ہم ساخت نقشے دا
’’مُڈھلا ٹکڑا‘‘ وی ہم ساخت نقشہ اے:

\begin{lem}\ollabel{wellordinitialsegment}
منّ لو $\tuple{A, <}$ تے $\tuple{B, \lessdot}$ حُسنِ ترتیب والیاں ساختاں
نیں۔ جے $f \colon A \to B$ اک ہم ساخت نقشہ اے تے $a \in A$، تاں
$\funrestrictionto{f}{A_{a}} : A_a \to B_{f(a)}$ اک ہم ساخت نقشہ اے۔
\end{lem}

\begin{proof}
چونکہ $f$ اک ہم ساخت نقشہ اے:
	%Since $f$ is an isomorphism, $b < a$ iff $f(b) \lessdot f(a)$, so that
\begin{align*}
	\funimage{f}{A_a} &= \funimage{f}{\Setabs{x \in A}{x < a}}\\
	&= \funimage{f}{\Setabs{f^{-1}(y) \in A}{f^{-1}(y) < a}} \\
	&= \Setabs{y \in B}{y \lessdot f(a)} \\
	&=B_{f(a)}
\end{align*}
تے $\funrestrictionto{f}{A_a}$ ترتیب قائم رکھدا اے کیوں جے $f$ اینوں قائم رکھدا اے۔
\end{proof}

اگلے دو نتیجے ثابت کردے نیں کہ حُسنِ ترتیب والیاں ساختاں دی ہمیشہ آپس وچ پرکھ ہو سکدی اے:

\begin{lem}\ollabel{lemordsegments}
منّ لو $\tuple{A, <}$ تے $\tuple{B, \lessdot}$ حُسنِ ترتیب والیاں ساختاں نیں۔
جے $\ordeq{\tuple{A_{a_1}, <_{a_1}}}{\tuple{B_{b_1}, \lessdot_{b_1}}}$
تے $\ordeq{\tuple{A_{{a_2}}, <_{a_2}}}{\tuple{B_{{b_2}}, \lessdot_{b_2}}}$،
تاں ${a_1}  < {a_2} \text{ اُدوں تے صرف اُدوں جدوں }{b_1} \lessdot {b_2}$۔
\end{lem}

\begin{proof}
اسیں \emph{کھبے توں سجے پاسے} ثابت کراں گے۔ دوجے پاسے وی ایہو جیہا ثبوت اے۔
منّ لو $\ordeq{\tuple{A_{a_1}, <_{a_1}}}{\tuple{B_{b_1}, \lessdot_{b_1}}}$
تے $\ordeq{\tuple{A_{{a_2}}, <_{a_2}}}{\tuple{B_{{b_2}}, \lessdot_{b_2}}}$
دونیں نیں، تے $f \colon A_{{a_2}} \to B_{{b_2}}$ ساڈا ہم ساخت نقشہ اے۔
منّ لو ${a_1} < {a_2}$۔ فیر
$\ordeq{\tuple{A_{a_1}, <_{a_1}}}{\tuple{B_{f({a_1})}, \lessdot_{f({a_1})}}}$
اے، \olref{wellordinitialsegment} نال۔ سو
$\ordeq{\tuple{B_{b_1}, \lessdot_{b_1}}}{\tuple{B_{f({a_1})}, \lessdot_{f({a_1})}}}$،
تے ایس لئی ${b_1} = f({a_1})$، \olref{wellordnotinitial} نال۔
ہُن ${b_1} \lessdot {b_2}$، کیوں جے $f$ دا دائرہ $B_{{b_2}}$ اے۔
\end{proof}

\begin{thm}\ollabel{thm:woalwayscomparable}
کسے وی دو حُسنِ ترتیب والیاں ساختاں وچوں اک دوجی دے کسے مُڈھلے ٹکڑے
نال ہم ساخت اے۔ ایہ ٹکڑا حقیقی ہونا ضروری نہیں۔
\end{thm}

\begin{proof}
منّ لو $\tuple{A, <}$ تے $\tuple{B, \lessdot}$ حُسنِ ترتیب والیاں ساختاں نیں۔
وکھ کرن دے اصول نال منّ لو
\[
	f = \Setabs{\tuple{a, b} \in A \times B}{
		\ordeq{\tuple{A_a, <_a}}{\tuple{B_b, \lessdot_b}}}.
\]
\olref{lemordsegments} نال، $a_1 < a_2$ اُدوں تے صرف اُدوں اے جدوں
$b_1 \lessdot b_2$، ساریاں $\tuple{a_1, b_1}, \tuple{a_2, b_2} \in f$
لئی۔ سو $f \colon \dom{f} \to \ran{f}$ اک ہم ساخت نقشہ اے۔

جے $a_2 \in \dom{f}$ تے $a_1 < a_2$، تاں $a_1 \in \dom{f}$،
\olref{wellordinitialsegment} نال۔ سو $\dom{f}$، $A$ دا مُڈھلا ٹکڑا اے۔
ایسے طرح $\ran{f}$، $B$ دا مُڈھلا ٹکڑا اے۔ تضاد کڈھن لئی منّ لو دونیں
\emph{حقیقی} مُڈھلے ٹکڑے نیں۔ فیر $a$ نوں $<$ دے لحاظ نال
$A \setminus \dom{f}$ دا سبھ توں چھوٹا !!{element} منّ لو، سو
$\dom{f} = A_a$۔ تے $b$ نوں $\lessdot$ دے لحاظ نال
$B \setminus \ran{f}$ دا سبھ توں چھوٹا !!{element} منّ لو، سو
$\ran{f} = B_b$۔ سو $f \colon A_a \to B_b$ اک ہم ساخت نقشہ اے،
تے ایس لئی $\tuple{a, b} \in f$، جیہڑا تضاد اے۔
\end{proof}
\endgroup
% /OLP-0551
% OLP-0552 content/set-theory/ordinals/vn.tex
\begingroup
\def\olfilename{vn}
\expandafter\def\csname ol@sectioncs\endcsname{section}

\olfileid{sth}{ordinals}{vn}
\olsection[فون نوئمان دی بناوٹ]{ترتیبی عدداں دی فون نوئمان والی بناوٹ}

\olref[sth][ordinals][iso]{thm:woalwayscomparable} توں اک خیال آؤندا اے۔
اسیں کجھ شےواں متعارف کرا سکدے آں، جیہناں نوں \emph{ترتیب دیاں قسماں}
آکھاں گے، تے جیہڑیاں حُسنِ ترتیب والیاں ساختاں دی تھاں ورتیاں جان۔
جے حُسنِ ترتیب والی ساخت دی ترتیب دی قسم لئی $\ordtype{A, <}$ لکھدے آں،
جیہڑی ساخت $\tuple{A, <}$ اے، تاں سانوں ہِیٹھلے دو اصول ملن دی آس ہووے گی:
\begin{align*}
	\ordtype{A, <} = \ordtype{B, \lessdot} &
	\text{ اُدوں تے صرف اُدوں جدوں } \ordeq{\tuple{A, <}}{\tuple{B, \lessdot}}\\
	\ordtype{A, <} < \ordtype{B, \lessdot}&
	\text{ اُدوں تے صرف اُدوں جدوں }\ordeq{\tuple{A, <}}{\tuple{B_b, \lessdot_b}}\text{ کسے ایہو جیہے }b \in B
\end{align*}
نالے اسیں ترتیب دیاں قسماں نوں \emph{کجھ خاص سیٹاں دے طور اُتے}
متعارف کراؤن دی آس رکھ سکدے آں، جیویں قدرتی عدداں نوں کجھ خاص سیٹاں
دے طور اُتے متعارف کرا سکدے آں۔

ایہ کرن دا سبھ توں عام طریقہ، تے اوہی طریقہ جیہڑا اسیں ورتاں گے،
ترتیب دیاں قسماں نوں کجھ \emph{معیاری} حُسنِ ترتیب والے سیٹاں نال
تعریف کرنا اے۔ ایہ معیاری سیٹ سبھ توں پہلاں فون نوئمان نے متعارف کرائے:

\begin{defn}
سیٹ $A$ \emph{متعدی} اُدوں تے صرف اُدوں اے جدوں
$(\forall x \in A)x \subseteq A$۔ فیر $A$ \emph{ترتیبی عدد} اُدوں تے
صرف اُدوں اے جدوں $A$ متعدی اے تے $\in$ اینوں حُسنِ ترتیب وچ رکھدا اے۔
\end{defn}
\noindent
اگے اسیں ترتیبی عدداں لئی یونانی حرف ورتاں گے۔ تعریف توں جھٹ نکلدا اے
کہ جے $\alpha$ ترتیبی عدد اے، تاں $\tuple{\alpha, \in_\alpha}$
حُسنِ ترتیب والی ساخت اے، جتھے $\in_\alpha =
\Setabs{\tuple{x, y} \in \alpha^2}{x \in y}$۔ سو نشان نوں ذرا کھُلّ
نال ورت کے آکھ سکدے آں کہ $\alpha$ \emph{آپ} حُسنِ ترتیب والی ساخت اے۔

ساڈے پہلے کجھ ترتیبی عدد ایہ نیں:
\[
	\emptyset, \{\emptyset\},
	\{\emptyset, \{\emptyset\}\},
	\{\emptyset, \{\emptyset\}, \{\emptyset, \{\emptyset\}\}\}, \ldots
\]
تسی ویکھو گے کہ ایہ اوہی پہلے کجھ ترتیبی عدد نیں جیہڑے ساڈے لامتناہیت دے
مسلّمے وچ آئے سن، یعنی فون نوئمان دی $\omega$ دی تعریف وچ
(\olref[sth][z][infinity-again]{sec} ویکھو)۔ ایہ اتفاق نہیں۔ فون نوئمان دی
ترتیبی عدداں دی تعریف قدرتی عدداں نوں ترتیبی عدد منّدی اے، تے لامتناہی
درجے دے ترتیبی عدداں دی وی تھاں رکھدی اے۔

ہمیشہ وانگ ہُن پُچھ سکدے آں کہ کی ایہ \emph{واقعی} ترتیبی عدد نیں؟
یا فون نوئمان نے بس کجھ سیٹ دتے نیں جیہناں نوں اسیں ترتیبی عدد
\emph{منّ} سکدے آں؟ ایس سوال اُتے جیہڑی گل بات ہو سکدی اے، اوہ
\olref[sfr][rel][ref]{sec}، \olref[sfr][arith][ref]{sec}،
\olref[sfr][infinite][dedekindsproof]{sec} تے \olref[sth][z][nat]{sec}
والی گل بات ورگی اے، سو اسیں ایس نوں نہیں کھچاں گے۔ اگے اسیں بس
’’ترتیبی عدد‘‘ نال ’’فون نوئمان دے ترتیبی عدد‘‘ آکھاں گے۔
\endgroup
% /OLP-0552
% OLP-0553 content/set-theory/ordinals/basic.tex
\begingroup
\def\olfilename{basic}
\expandafter\def\csname ol@sectioncs\endcsname{section}

\olfileid{sth}{ordinals}{basic}
\olsection{ترتیبی عدداں دیاں بنیادی خاصیتاں}

اسیں ویکھیا سی کہ پہلے کجھ ترتیبی عدد قدرتی عدد نیں۔ ترتیبی عدداں دا
نظریہ بناؤن دی وڈی وجہ اوہ استقرائی اصول ودھاؤنا اے جیہڑا قدرتی عدداں
لئی قائم اے۔ اسیں کجھ بنیادی نتیجیاں دی لڑی راہیں اوس اصول تک پہنچاں گے۔

\begin{lem}\ollabel{ordmemberord}
ترتیبی عدد دا ہر !!{element} آپ اک ترتیبی عدد اے۔
\end{lem}

\begin{proof}
منّ لو $\alpha$ ترتیبی عدد اے تے $b \in \alpha$۔ چونکہ $\alpha$
متعدی اے، $b \subseteq \alpha$۔ سو $\in$، $b$ نوں حُسنِ ترتیب وچ
رکھدا اے، کیوں جے $\in$، $\alpha$ نوں حُسنِ ترتیب وچ رکھدا اے۔

ایہ ویکھن لئی کہ $b$ متعدی اے، منّ لو $x \in c \in b$۔ سو $c \in \alpha$،
کیوں جے $b \subseteq \alpha$۔ فیر چونکہ $\alpha$ متعدی اے،
$c \subseteq \alpha$، تے ایس لئی $x \in \alpha$۔ سو $x, c, b \in \alpha$۔
پر $\in$، $\alpha$ نوں حُسنِ ترتیب وچ رکھدا اے، ایس لئی $\in$، $\alpha$
اُتے متعدی تعلق اے، \olref[sth][ordinals][wo]{wo:strictorder} نال۔
سو چونکہ $x \in c \in b$، ایس لئی $x \in b$۔ عام صورت وچ $c \subseteq b$۔
\end{proof}

\begin{cor}\ollabel{ordissetofsmallerord}
$\alpha = \Setabs{\beta \in \alpha}{\beta \text{ ترتیبی عدد اے}}$،
ہر ترتیبی عدد~$\alpha$ لئی۔
\end{cor}

\begin{proof}
\olref{ordmemberord} توں سدھا نکلدا اے۔
\end{proof}

اگلے دو اہم نتیجیاں، \olref{ordinductionschema} تے \olref{ordtrichotomy}،
دی موٹی گل ایہ اے کہ رکنیت آپ ترتیبی عدداں نوں حُسنِ ترتیب وچ رکھدی اے:

\begin{thm}[لامتناہی درجیاں اُتے استقرا]\ollabel{ordinductionschema}
ہر فارمولے $\phi(x)$ لئی:
\[
	\text{ جے }\exists \alpha \phi(\alpha)\text{، تاں }\exists \alpha(\phi(\alpha)
	\land  (\forall \beta \in \alpha) \lnot \phi(\beta))
\]
ایتھے دِتے مقداری نشان ان کہے طور اُتے ترتیبی عدداں تک محدود نیں۔
\end{thm}

\begin{proof}
%\olref{ordinductionschema1}. 
منّ لو $\phi(\alpha)$ کسے ترتیبی عدد $\alpha$ لئی۔ جے $ (\forall \beta
\in \alpha) \lnot \phi(\beta)$، تاں ثبوت پورا ہویا۔ نہیں تاں چونکہ
$\alpha$ ترتیبی عدد اے، اوہدے وچ $\in$ دے لحاظ نال اک سبھ توں چھوٹا
!!{element} اے جیہڑا $\phi$ پورا کردا اے، تے \olref{ordmemberord} نال
ایہ ترتیبی عدد اے۔
%	\olref{ordinductionschema2}. Suppose there is some ordinal
%	$\gamma$ such that $\lnot\phi(\gamma)$. Then by
%	\olref{ordinductionschema1}, there is an $\in$-minimal ordinal
%	$\alpha$ for which $\lnot\phi(\alpha)$. So $(\forall \beta <
%	\alpha) \phi(\beta)$, rendering the antecedent of the conditional
%	false.
\end{proof}
\noindent
دھیان دیو کہ \olref{ordinductionschema} نوں ایس خاکے دے روپ وچ وی لکھ سکدے آں:
\[
\text{ جے }\forall \alpha((\forall \beta \in \alpha)\phi(\beta) \lif 
\phi(\alpha))\text{، تاں }\forall \alpha\phi(\alpha)
\]
بس $\lnot\phi(\alpha)$ نوں \olref{ordinductionschema} وچ لینا اے تے فیر بنیادی منطقی تبدیلیاں کرنیاں نیں۔

\begin{thm}[تِن صورتاں دا اصول]\ollabel{ordtrichotomy}
$\alpha \in \beta \lor \alpha = \beta \lor \beta \in \alpha$، کسے وی
دو ترتیبی عدداں $\alpha$ تے $\beta$ لئی۔
\end{thm}

\begin{proof}
ثبوت دوہرے استقرا نال اے، یعنی \olref{ordinductionschema} نوں دو واری
ورت کے۔ آکھو کہ $x$، $y$ نال \emph{پرکھیا جا سکدا اے} اُدوں تے صرف اُدوں
جدوں $x \in y \lor x = y \lor y \in x$۔

استقرا لئی منّ لو کہ~$\alpha$ دے ہر ترتیبی عدد دی \emph{ہر} ترتیبی
عدد نال پرکھ ہو سکدی اے۔ ہور استقرا لئی منّ لو کہ $\alpha$ دی~$\beta$
دے ہر ترتیبی عدد نال پرکھ ہو سکدی اے۔ اسیں وکھاواں گے کہ $\alpha$
دی~$\beta$ نال پرکھ ہو سکدی اے۔~$\beta$ اُتے استقرا نال نکلے گا کہ $\alpha$
دی ہر ترتیبی عدد نال پرکھ ہو سکدی اے۔ فیر~$\alpha$ اُتے استقرا نال،
\emph{ہر} ترتیبی عدد دی \emph{ہر} ترتیبی عدد نال پرکھ ہو سکدی اے،
جیہڑی سانوں چاہیدی سی۔ ایہ کافی اے کہ $\alpha \notin \beta$ تے
$\beta \notin \alpha$ منّ کے وکھائیے کہ $\alpha = \beta$۔

ایہ وکھاؤن لئی کہ $\alpha \subseteq \beta$، $\gamma \in \alpha$ مقرر کرو۔
\olref{ordmemberord} نال ایہ اک ترتیبی عدد اے۔ سو پہلے استقرائی مفروضے
نال $\gamma$ دی $\beta$ نال پرکھ ہو سکدی اے۔ پر جے یا تاں $\gamma = \beta$
یا $\beta \in \gamma$، تاں $\beta \in \alpha$ (لوڑ پئے تاں ایس گل نوں
ورت کے کہ $\alpha$ متعدی اے)۔ ایہ ساڈے مفروضے دے اُلٹ اے۔ سو
$\gamma \in \beta$۔ عام صورت وچ $\alpha \subseteq \beta$۔

دوجے استقرائی مفروضے نال عین ایہو جیہی دلیل وکھاندی اے کہ
$\beta \subseteq \alpha$۔ سو $\alpha = \beta$۔
\end{proof}\noindent ایس لئی کدی اسیں $\alpha <\beta$ لکھاں گے،
$\alpha \in \beta$ دی تھاں، کیوں جے $\in$ ترتیب دے تعلق وانگ کم کردا اے۔
ایہدی کوئی ڈونگھی وجہ نہیں، بس ایہ جانیا پچھانیا اے تے
$\alpha \leq \beta$ لکھنا $\alpha \in \beta \lor \alpha = \beta$
لکھن نالوں سوکھا اے۔\footnote{اسیں $\alpha \mathrel{\underline{\in}} \beta$
لکھ سکدے ساں، پر ایہ بالکل غیر معیاری ہوندا۔}

ساڈے پچھلے نتیجیاں دے دو سدھے حاصل ایہ نیں۔ پہلا ساڈے نویں نشاناں نوں کم وچ لیاندا اے:

\begin{cor}\ollabel{ordordered}
جے $\exists \alpha\phi(\alpha)$، تاں $\exists \alpha(\phi(\alpha)
\land \forall \beta(\phi(\beta) \lif \alpha \leq \beta))$۔
نالے کسے وی ترتیبی عدداں $\alpha, \beta, \gamma$ لئی،
$\alpha \notin \alpha$ تے $\alpha \in \beta \in \gamma \lif \alpha \in \gamma$
دونیں قائم نیں۔
\end{cor}

\begin{proof}
بالکل \olref[wo]{wo:strictorder} وانگ۔
\end{proof}

\begin{prob}
\olref[sth][ordinals][basic]{ordtrichotomy} دے ثبوت وچ ’’عین ایہو جیہی دلیل‘‘
نوں پورا کرو۔
\end{prob}

\begin{cor}\ollabel{corordtransitiveord}
$A$ ترتیبی عدد اُدوں تے صرف اُدوں اے جدوں $A$ ترتیبی عدداں دا متعدی سیٹ اے۔
\end{cor}

\begin{proof}
\emph{کھبے توں سجے پاسے۔} \olref{ordmemberord} نال۔
\emph{سجے توں کھبے پاسے۔} جے $A$ ترتیبی عدداں دا متعدی سیٹ اے، تاں
$\in$، $A$ نوں حُسنِ ترتیب وچ رکھدا اے،
\olref{ordinductionschema} تے \olref{ordtrichotomy} نال۔
\end{proof}

اسیں \olref{ordinductionschema} تے \olref{ordtrichotomy} دا مطلب ایہ دسیا سی
کہ $\in$ ترتیبی عدداں نوں حُسنِ ترتیب وچ رکھدا اے۔ پر ہِیٹھلے نتیجے کر کے
ایہو جیہے دعوے بارے سانوں \emph{بڑی احتیاط} کرنی پئے گی:

\begin{thm}[بورالی فورتی دا پیراڈوکس]\ollabel{buraliforti}
سارے ترتیبی عدداں دا کوئی سیٹ نہیں۔
\end{thm}

\begin{proof}
تضاد کڈھن لئی منّ لو $O$ سارے ترتیبی عدداں دا سیٹ اے۔ جے
$\alpha \in \beta \in O$، تاں $\alpha$ ترتیبی عدد اے، \olref{ordmemberord}
نال، سو $\alpha \in O$۔ ایس لئی $O$ متعدی اے، تے \olref{corordtransitiveord}
نال $O$ ترتیبی عدد اے۔ ایس لئی $O \in O$، جیہڑا
\olref{ordordered} دے اُلٹ اے۔
\end{proof}
\noindent
ایہ نتیجہ \citeauthor{Burali-Forti1897} دے ناں اُتے اے۔ پر 1899 وچ ڈیڈیکنڈ
نوں اک خط وچ کانتور نے سبھ توں پہلاں صاف ویکھیا کہ سارے ترتیبی عدداں دا
سیٹ منّن وچ \emph{تضاد} اے۔ جیویں فان ہائینورٹ دسدا اے:
\begin{quote}
  بورالی فورتی آپ ایس تضاد نوں \emph{برہانِ خلف} نال ایس نتیجے دا ثبوت
  سمجھدا سی کہ ترتیبی عدداں دی قدرتی ترتیب بس اک جزوی ترتیب اے۔
  \citep[p.~105]{Heijenoort1967}
\end{quote}
بورالی فورتی دی غلطی اک پاسے رکھ کے، پچھلی گل ایویں سمیٹ سکدے آں۔
ترتیبی عدد سیٹ نیں۔ رکنیت ہر اک نوں حُسنِ ترتیب وچ رکھدی اے تے
ساریاں نوں رلا کے وی حُسنِ ترتیب وچ رکھدی اے (بھانویں سارے رل کے سیٹ
نہیں بناندے)۔

گل پوری کرن لئی ترتیبی عدداں دیاں کجھ ہور بنیادی خاصیتاں ہِیٹھاں نیں،
جیہڑیاں \olref{ordinductionschema} تے \olref{ordtrichotomy} توں نکلدیاں نیں۔

\begin{prop}
ترتیبی عدداں دی ہر سختی نال گھٹدی لڑی محدود اے۔
\end{prop}

\begin{proof}
ترتیبی عدداں دی کسے وی لامتناہی سختی نال گھٹدی لڑی
$\alpha_0 > \alpha_1 > \alpha_2 > \ldots$ وچ $<$ دے لحاظ نال
کوئی ایسا رکن نہیں جیہدے توں ہور کوئی رکن چھوٹا نہ ہووے۔
ایہ \olref{ordinductionschema} دے اُلٹ اے۔
\end{proof}

\begin{prop}\ollabel{ordinalsaresubsets}
$\alpha \subseteq \beta \lor \beta \subseteq \alpha$، کسے وی
ترتیبی عدداں $\alpha, \beta$ لئی۔
\end{prop}

\begin{proof}
جے $\alpha \in \beta$، تاں $\alpha \subseteq \beta$، کیوں جے $\beta$
متعدی اے۔ ایسے طرح جے $\beta \in \alpha$، تاں $\beta \subseteq \alpha$۔
تے جے $\alpha = \beta$، تاں $\alpha \subseteq \beta$ تے $\beta \subseteq \alpha$۔
سو \olref{ordtrichotomy} نال ثبوت پورا ہویا۔
\end{proof}

\begin{prop}\ollabel{ordisoidentity}
$\alpha = \beta$ اُدوں تے صرف اُدوں اے جدوں $\ordeq{\alpha}{\beta}$،
کسے وی ترتیبی عدداں $\alpha, \beta$ لئی۔
\end{prop}

\begin{proof}
ترتیبی عدد حُسنِ ترتیب والیاں ساختاں نیں۔ سو ایہ تِن صورتاں دے اصول
(\olref[basic]{ordtrichotomy}) تے \olref[iso]{wellordnotinitial}
توں سدھا نکلدا اے۔
\end{proof}

\begin{prob}\ollabel{probunionordinalsordinal}
ثابت کرو کہ جے $X$ دا ہر رکن ترتیبی عدد اے، تاں $\bigcup X$ وی ترتیبی عدد اے۔
\end{prob}
\endgroup
% /OLP-0553
% OLP-0554 content/set-theory/ordinals/replacement.tex
\begingroup
\def\olfilename{replacement}
\expandafter\def\csname ol@sectioncs\endcsname{section}

\olfileid{sth}{ordinals}{replacement}
\olsection{بدل رکھن دا اصول}

\olref[sth][ordinals][vn]{sec} وچ اسیں ترتیبی عدداں نوں متعارف کراؤن دا
جواز ایہ دتا سی کہ اسیں ایہناں نوں ترتیب دیاں قسماں، یعنی حُسنِ ترتیب
والیاں ساختاں دے معیاری نمائندے، منّ سکدے آں۔ ایس لئی سانوں ثابت کرنا
پئے گا کہ \emph{ہر حُسنِ ترتیب والی ساخت کسے ترتیبی عدد نال ہم ساخت اے}۔
ایس نال اسیں $\ordtype{A, <}$ نوں اوہ ترتیبی عدد $\alpha$ تعریف کر سکاں گے
جیہدے لئی $\tuple{A, <} \isomorphic \alpha$۔

بدقسمتی نال ہُن تک دتے مسلّمیاں نال ای اسیں مطلوبہ نتیجہ ثابت
\emph{نہیں کر سکدے}۔ (\olref[replacement][strength]{sec} وچ وجہ ویکھاں گے۔
ہُن لئی گل اینی اے کہ نہیں کر سکدے۔) سانوں اک نویں خیال دی لوڑ اے۔ اوہ ایہ اے:

\begin{axiom}[بدل رکھن دا خاکہ]
ہر فارمولے $\phi(x, y)$ لئی ہِیٹھلی گل اک مسلّمہ اے:
\begin{quote}
	کسے وی $A$ لئی، جے $(\forall x \in A)\lexists![y][\phi(x,y)]$، تاں
	$\Setabs{y}{(\exists x \in A)\phi(x,y)}$ موجود اے۔
\end{quote}
\end{axiom}
\noindent
وکھ کرن دے اصول وانگ ایہ اک خاکہ اے۔ ایہ لامتناہی گنتی وچ مسلّمے دیندا اے،
لامتناہی گنتی وچ وکھ وکھ $\phi$ وچوں ہر اک لئی اک۔ اینوں ایویں وی لکھ
سکدے آں (تے عام طور اُتے ایویں ای لکھدے آں):

\begin{defish}
ہر اوہ فارمولا $\phi(x,y)$ جیہدے وچ ’’$B$‘‘ نہیں آؤندا، اوہدے لئی
ہِیٹھلی گل اک مسلّمہ اے:
\[
\forall A[(\forall x \in A)\lexists![y][\phi(x,y)] \lif \exists B\forall y (y \in B \liff (\exists x \in A)\phi(x,y))]
\]
\end{defish}

پر پہلی واری ویکھن اُتے ایہ خاصا اُلجھیا فارمولا اے۔ بدل رکھن دے اصول
دا ہِیٹھلا سدھا نتیجہ شاید ساڈے بنیادی خیال نوں \emph{ہور صاف} طرح
بیان کردا اے:\footnote{ٹرم اوہ عبارت اے جیہڑی اکّو اک شے چُندی اے
(اوہدے آزاد متغیراں دیاں کوئی وی قدراں مقرر ہون اُتے)۔ مثلاً ’’$\emptyset$‘‘
اک ٹرم اے جیہڑی خالی سیٹ چُندی اے۔ ’’$\{x\}$‘‘ اک ٹرم اے جیہڑی
$x$ دا اک رکن والا سیٹ چُندی اے (بھانویں $x$ کوئی وی ہووے)۔
’’$x \cup y$‘‘ اک ٹرم اے جیہڑی $x$ تے $y$ دا اتحاد چُندی اے
(بھانویں اوہ کوئی وی ہون)۔}

\begin{cor}
کسے وی ٹرم $\tau(x)$ تے کسے وی سیٹ $A$ لئی ایہ سیٹ موجود اے:
\[
	\Setabs{\tau(x)}{x \in A} = \Setabs{y}{(\exists x \in A)y = \tau(x)}.
\]
\end{cor}

\begin{proof}
چونکہ $\tau$ اک \emph{ٹرم} اے، $\forall x \lexists![y][\tau(x) = y]$۔
ایس توں ہور وی صاف اے کہ $(\forall x \in A)\lexists![y][\tau(x) = y]$۔
سو $\Setabs{y}{(\exists x \in A)\tau(x) = y}$ بدل رکھن دے اصول نال موجود اے۔
\end{proof}
\noindent
ایس توں پتا لگدا اے کہ ’’بدل رکھنا‘‘ مسلّمے دا چنگا ناں اے۔ سیٹ $A$
دتا ہووے، تاں تسی نواں سیٹ $\Setabs{\tau(x)}{x \in A}$ بنا سکدے او،
$A$ دے ہر رکن دی تھاں اوہدا~$\tau$ تحت عکس رکھ کے۔ واقعی فنکشن تحت
سیٹ دے عکس والا نشان ورت کے اسیں $\funimage{\tau}{A}$ نوں
$\Setabs{\tau(x)}{x \in A}$ لئی لکھ سکدے آں۔

پر ضروری گل ایہ اے کہ $\tau$ اک \emph{ٹرم} اے۔ ایہ کوئی \emph{فنکشن}
(دا ناں) ہونا ضروری نہیں، جیویں \olref[sfr][fun][rel]{sec} وچ فنکشن
دا مطلب اے، یعنی مرتب جوڑیاں دا اک خاص سیٹ۔ آخر جے $f$ اوس معنی وچ
فنکشن اے، تاں سیٹ $\funimage{f}{A} = \Setabs{f(x)}{x \in A}$
بس $\ran{f}$ دا اک خاص ذیلی سیٹ اے، تے اوہدی موجودگی دی ضمانت
صرف~$\Zminus$ دے مسلّمے ای دیندے نیں۔\footnote{بس
$\Setabs{y \in \bigcup \bigcup f}{(\exists x \in A)y = f(x)}$ ویکھو۔}
اِس دے اُلٹ، بدل رکھن دا اصول ساڈے مسلّمیاں وچ اک \emph{طاقتور} وادھا اے،
جیویں اسیں \olref[sth][replacement][]{chap} وچ ویکھاں گے۔
\endgroup
% /OLP-0554
% OLP-0555 content/set-theory/ordinals/milestone.tex
\begingroup
\def\olfilename{milestone}
\expandafter\def\csname ol@sectioncs\endcsname{section}
\olfileid{sth}{ordinals}{zfm}
\olsection{\texorpdfstring{$\ZFminus$۔ اک اہم پڑاؤ}{ZF⁻۔ اک اہم پڑاؤ}}

بدل رکھن دے اصول دا جواز کیویں دینا اے (جے دے وی سکدے آں)، ایہ سدھا
سوال نہیں۔ سو ایہ گل اسیں \olref[replacement][]{chap} لئی رکھاں گے۔
پر بدل رکھن دا اصول شامل کر کے اسیں اک ہور اہم پڑاؤ اُتے پہنچ گئے آں۔
ہُن ساڈے کول نظریے $\ZFminus$ لئی لوڑیندے سارے مسلّمے نیں۔ تفصیل ایہ اے:

\begin{defn}
نظریے $\ZFminus$ دے مسلّمے ایہ نیں۔ عنصراں نال تعیّن، اتحاد، جوڑیاں،
ذیلی سیٹاں دے سیٹ، لامتناہیت، تے وکھ کرن تے بدل رکھن دے خاکیاں دیاں
ساریاں صورتاں۔ دوجے لفظاں وچ، $\ZFminus$، $\Zminus$ وچ بدل رکھن
دا اصول شامل کردا اے۔
\end{defn}
\noindent
ایہ ناں \emph{زرمَیلو--فرینکل} دے سیٹ نظریے لئی اے (\emph{گھٹاؤ}
اوہ شے جیہڑی اسیں اگے ویکھاں گے)۔ فرینکل ایس عزت دا حق دار اے، کیوں
جے \citeyear{Fraenkel1922} وچ بدل رکھن دا اصول بیان کرن دا سہرا اوہدے
سر اے، بھانویں پہلا ٹھیک ٹھیک بیان \citet{Skolem1922} نے دتا سی۔
\endgroup
% /OLP-0555
% OLP-0556 content/set-theory/ordinals/ordtype.tex
\begingroup
\def\olfilename{ordtype}
\expandafter\def\csname ol@sectioncs\endcsname{section}

\olfileid{sth}{ordinals}{ordtype}
\olsection{ترتیب دیاں قسماں دے طور اُتے ترتیبی عدد}

بدل رکھن دے اصول نال، تے ہُن $\ZFminus$ وچ کم کردیاں، اسیں آخر
اوہ نتیجہ ثابت کر سکدے آں جیہڑا ساڈا مقصد سی:

\begin{thm}\ollabel{thmOrdinalRepresentation}
ہر حُسنِ ترتیب والی ساخت اکّو اک ترتیبی عدد نال ہم ساخت اے۔
\end{thm}

\begin{proof}
منّ لو $\tuple{A, <}$ حُسنِ ترتیب والی ساخت اے۔
\olref[basic]{ordisoidentity} نال ایہ ودھ توں ودھ اک ترتیبی عدد نال
ہم ساخت اے۔ سو تضاد کڈھن لئی منّ لو $\tuple{A, <}$ \emph{کسے وی}
ترتیبی عدد نال ہم ساخت نہیں۔ اسیں پہلاں ’’$\tuple{A, <}$ نوں جِنّا ہو
سکے چھوٹا کراں گے‘‘۔ تفصیل ایہ اے۔ جے کوئی حقیقی مُڈھلا ٹکڑا
$\tuple{A_a, <_a}$ کسے وی ترتیبی عدد نال ہم ساخت نہیں، تاں ایس خاصیت
والا سبھ توں چھوٹا $a \in A$ اے۔ فیر منّ لو $B = A_a$ تے
$\mathord{\lessdot} = \mathord{<_a}$۔ نہیں تاں منّ لو $B = A$ تے
$\mathord{\lessdot} = \mathord{<}$۔

تعریف مطابق $B$ دا ہر حقیقی مُڈھلا ٹکڑا کسے ترتیبی عدد نال ہم ساخت اے،
تے اُتلی دلیل نال ایہ عدد اکّو اک اے۔ سو بدل رکھن دے اصول نال ہِیٹھلا
سیٹ موجود اے تے اک فنکشن اے:
\[
	f = \Setabs{\tuple{\beta, b}}{b \in B\text{ تے }\ordeq{\beta}{\tuple{B_b, \lessdot_b}}}
\]
تضاد پورا کرن لئی اسیں وکھاواں گے کہ $f$، $\alpha \to B$ دا اک
ہم ساخت نقشہ اے، کسے ترتیبی عدد $\alpha$ لئی۔

صاف اے کہ $\ran{f} = B$۔ تے \olref[iso]{lemordsegments} نال
$f$ ترتیب قائم رکھدا اے، یعنی $\gamma \in \beta$ اُدوں تے صرف
اُدوں جدوں $f(\gamma) \lessdot f(\beta)$۔ ایہ وکھاؤن لئی کہ $\dom{f}$
ترتیبی عدد اے، \olref[basic]{corordtransitiveord} نال اینا وکھاؤنا
کافی اے کہ $\dom{f}$ متعدی اے۔ سو $\beta \in \dom{f}$ مقرر کرو، یعنی
$\ordeq{\beta}{\tuple{B_b, \lessdot_b}}$ کسے $b$ لئی۔ جے $\gamma \in \beta$،
تاں $\gamma \in \dom{f}$، \olref[iso]{wellordinitialsegment} نال۔
عام صورت وچ $\beta \subseteq \dom{f}$۔
\end{proof}

ایہ نتیجہ سانوں ہِیٹھلی تعریف دین دی اجازت دیندا اے، جیہڑی اسیں
\olref[vn]{sec} توں دینی چاہندے ساں:

\begin{defn}
جے $\tuple{A, <}$ حُسنِ ترتیب والی ساخت اے، تاں اوہدی ترتیب دی قسم
$\ordtype{A, <}$ اوہ اکّو اک ترتیبی عدد $\alpha$ اے جیہدے لئی
$\ordeq{\tuple{A, < }}{\alpha}$۔
\end{defn}

نالے ایس تعریف توں دو چنگے اصول ملدے نیں:

\begin{cor}\ollabel{ordtypesworklikeyouwant}
جے $\tuple{A, <}$ تے $\tuple{B, \lessdot}$ حُسنِ ترتیب والیاں ساختاں نیں:
\begin{align*}
	\ordtype{A, <} = \ordtype{B, \lessdot}&\text{ اُدوں تے صرف اُدوں جدوں }\ordeq{\tuple{A, <}}{\tuple{B, \lessdot}}\\
	\ordtype{A, <} \in \ordtype{B, \lessdot}&\text{ اُدوں تے صرف اُدوں جدوں }\ordeq{\tuple{A, <}}{\tuple{B_b, \lessdot_b}}\text{ کسے ایہو جیہے }b \in B
\end{align*}
\end{cor}

\begin{proof}
برابری \olref[basic]{ordisoidentity} نال قائم اے۔ دوجی گل ثابت کرن لئی
منّ لو $\ordtype{A, <} = \alpha$ تے $\ordtype{B, \lessdot} = \beta$،
تے منّ لو $f \colon \beta \to \tuple {B, \lessdot}$ ساڈا ہم ساخت
نقشہ اے۔ فیر:
\begin{align*}
	\alpha \in \beta&\text{ اُدوں تے صرف اُدوں جدوں }\funrestrictionto{f}{\alpha} \colon \alpha \to B_{f(\alpha)}\text{ اک ہم ساخت نقشہ اے}\\
	&\text{ اُدوں تے صرف اُدوں جدوں }\ordeq{\tuple{A, <}}{\tuple{B_{f(\alpha)}, \lessdot_{f(\alpha)}}}\\
	&\text{ اُدوں تے صرف اُدوں جدوں }\ordeq{\tuple{A, <}}{\tuple{B_b, \lessdot_b}}\text{ کسے $b \in B$ لئی}
\end{align*}
\olref[iso]{ordisounique}، \olref[iso]{wellordinitialsegment} تے
\olref[basic]{ordissetofsmallerord} نال۔
\end{proof}
\endgroup
% /OLP-0556
% OLP-0557 content/set-theory/ordinals/opps.tex
\begingroup
\def\olfilename{opps}
\expandafter\def\csname ol@sectioncs\endcsname{section}

\olfileid{sth}{ordinals}{opps}
\olsection{جانشین تے حدّی ترتیبی عدد}

اگلے کجھ باباں وچ اسیں ترتیبی عدد بہت ورتاں گے۔ ایس لئی کجھ سادے تصور
متعارف کرانے کم آؤن گے۔

\begin{defn}
کسے وی ترتیبی عدد $\alpha$ دا \emph{جانشین}
$\ordsucc{\alpha} =\alpha \cup \{\alpha\}$ اے۔
اسیں $\alpha$ نوں \emph{جانشین} ترتیبی عدد آکھدے آں جے
$\ordsucc{\beta} = \alpha$ کسے ترتیبی عدد~$\beta$ لئی۔
اسیں $\alpha$ نوں \emph{حدّی} ترتیبی عدد اُدوں تے صرف اُدوں آکھدے آں
جدوں $\alpha$ نہ خالی ہووے، نہ جانشین ترتیبی عدد۔
\end{defn}
\noindent
ہِیٹھلا نتیجہ وکھاندا اے کہ \emph{جانشین} دا ایہو ٹھیک تصور اے:

\begin{prop}
کسے وی ترتیبی عدد $\alpha$ لئی:
\begin{enumerate}
	\item $\alpha \in \ordsucc{\alpha}$۔
	\item $\ordsucc{\alpha}$ ترتیبی عدد اے۔
	\item کوئی ترتیبی عدد $\beta$ نہیں جیہدے لئی
	$\alpha \in \beta \in \ordsucc{\alpha}$۔
\end{enumerate}
\end{prop}

\begin{proof}
ظاہر اے کہ $\alpha \in \alpha \cup \{\alpha\} = \ordsucc{\alpha}$۔
ایسے طرح $\ordsucc{\alpha}$ ترتیبی عدداں دا متعدی سیٹ اے، تے ایس لئی
\olref[basic]{corordtransitiveord} نال ترتیبی عدد اے۔
تے $\alpha \in \beta \in \ordsucc{\alpha}$ ہونا ناممکن اے، کیوں جے
فیر یا $\beta \in \alpha$ یا $\beta = \alpha$، جیہڑا
\olref[basic]{ordordered} دے اُلٹ اے۔
\end{proof}
\noindent
ایس نال لامتناہی درجیاں اُتے استقرا نال ثبوت دا اک ہور طریقہ وی ملدا اے:

\begin{thm}[لامتناہی درجیاں اُتے سادہ استقرا]\ollabel{simpletransrecursion}
منّ لو $\phi(x)$ اک فارمولا اے جیہدے لئی:
\begin{enumerate}
	\item $\phi(\emptyset)$، تے
	\item کسے وی ترتیبی عدد $\alpha$ لئی، جے $\phi(\alpha)$، تاں
	$\phi(\ordsucc{\alpha})$، تے
	\item جے $\alpha$ حدّی ترتیبی عدد اے تے
	$(\forall \beta \in \alpha)\phi(\beta)$، تاں $\phi(\alpha)$۔
\end{enumerate}
فیر $\forall \alpha \phi(\alpha)$۔
\end{thm}

\begin{proof}
اسیں عکسِ نقیض ثابت کردے آں۔ سو منّ لو کوئی ترتیبی عدد اے جیہدے لئی
$\lnot\phi$۔ منّ لو $\gamma$ اوہ سبھ توں چھوٹا ترتیبی عدد اے۔
فیر یا $\gamma = \emptyset$، یا $\gamma = \ordsucc{\alpha}$ کسے ایہو
جیہے $\alpha$ لئی جیہدے لئی $\phi(\alpha)$، یا $\gamma$ حدّی ترتیبی عدد
اے تے $(\forall \beta \in \gamma)\phi(\beta)$۔
\end{proof}
\noindent
اک آخری نشان اگے کم آوے گا:

\begin{defn}\ollabel{defsupstrict}
جے $X$ ترتیبی عدداں دا سیٹ اے، تاں
$\supstrict(X) = \bigcup_{\alpha \in X} \ordsucc{\alpha}$۔
\end{defn}
\noindent
ایتھے ’’lsub‘‘ دا مطلب ’’سبھ توں چھوٹی سخت اُپَرلی حد‘‘ اے۔\footnote{کجھ
کتاباں ایس لئی ’’$\text{sup}(X)$‘‘ ورتدیاں نیں۔ پر ہور کتاباں
’’$\text{sup}(X)$‘‘ سبھ توں چھوٹی \emph{غیر سخت} اُپَرلی حد لئی
ورتدیاں نیں، یعنی بس $\bigcup X$۔ جے $X$ دا سبھ توں وڈا عنصر $\alpha$
ہووے، تاں ایہ تصور وکھرے ہو جاندے نیں۔ سبھ توں چھوٹی \emph{سخت} اُپَرلی
حد $\ordsucc{\alpha}$ اے، جدکہ سبھ توں چھوٹی \emph{غیر سخت} اُپَرلی حد
بس $\alpha$ اے۔} ہِیٹھلا نتیجہ ایس گل نوں سمجھاندا اے:

\begin{prop}
جے $X$ ترتیبی عدداں دا سیٹ اے، تاں $\supstrict(X)$ اوہ سبھ توں چھوٹا
ترتیبی عدد اے جیہڑا $X$ دے ہر ترتیبی عدد توں وڈا اے۔
\end{prop}

\begin{proof}
منّ لو $Y = \Setabs{\ordsucc{\alpha}}{\alpha \in X}$، سو
$\supstrict(X) = \bigcup Y$۔ چونکہ ترتیبی عدد متعدی نیں تے ترتیبی عدد
دا ہر رکن ترتیبی عدد اے، $\supstrict(X)$ ترتیبی عدداں دا متعدی سیٹ اے۔
سو \olref[basic]{corordtransitiveord} نال ایہ ترتیبی عدد اے۔

جے $\alpha \in X$، تاں $\ordsucc{\alpha} \in Y$، سو
$\ordsucc{\alpha} \subseteq \bigcup Y = \supstrict(X)$، تے ایس لئی
$\alpha \in \supstrict(X)$۔ سو $\supstrict(X)$، $X$ دے ہر ترتیبی عدد
توں سختی نال وڈا اے۔

اُلٹے پاسے، جے $\alpha \in \supstrict(X)$، تاں
$\alpha \in \ordsucc{\beta} \in Y$ کسے $\beta \in X$ لئی، سو
$\alpha \leq \beta \in X$۔ سو $\supstrict(X)$، $X$ دی \emph{سبھ توں
چھوٹی} سخت اُپَرلی حد اے۔
\end{proof}
\endgroup
% /OLP-0557

\OLEndChapterHook
\endgroup
% /OLP-0547
% OLP-0558 content/set-theory/spine/spine.tex
\begingroup
\def\olfilename{spine}
\expandafter\def\csname ol@sectioncs\endcsname{section}

\olchapter{sth}{spine}{مرحلے تے درجے}

% OLP-0559 content/set-theory/spine/idea.tex
\begingroup
\def\olfilename{idea}
\expandafter\def\csname ol@sectioncs\endcsname{section}

\olfileid{sth}{spine}{valpha}

\olsection{\texorpdfstring{مرحلیاں دی $V_\alpha$ دے روپ وچ تعریف}{مرحلیاں دی V\textunderscore{}α دے روپ وچ تعریف}}

\olref[sth][ordinals][]{chap} وچ اسیں حُسنِ ترتیب والیاں ساختاں تے
(فون نوئمان دے) ترتیبی عدد تعریف کیتے سن۔ ایس باب وچ اسیں ایہناں نوں
سیٹاں دی درجہ بندی \emph{آپ} دی خاصیت دسن لئی ورتاں گے۔ ایس لئی یاد کرو کہ
\olref[sth][ordinals][opps]{sec} وچ اسیں جانشین تے حدّی ترتیبی عدداں دا تصور
تعریف کیتا سی۔ ہِیٹھلی تعریف وچ اسیں ایہ تصور ورتدے آں:

\begin{defn}\ollabel{defValphas}
	\begin{align*}
	V_\emptyset &\defis \emptyset\\
	V_{\ordsucc{\alpha}} &\defis \Pow{V_\alpha} & &
	\text{کسے وی ترتیبی عدد لئی }\alpha\\
	V_{\alpha} &\defis \bigcup_{\gamma < \alpha} V_\gamma & &
	\text{جدوں }\alpha\text{ حدّی ترتیبی عدد ہووے}
\end{align*}
\end{defn}
\noindent
ایہ ترتیبی عدداں اُتے \emph{لامتناہی درجیاں والی بازگشت} نال تعریف ہووے گی۔
ایس لحاظ نال سانوں ایہدا قدرتی عدداں اُتے فنکشوناں دیاں بازگشتی تعریفاں
نال ٹاکرا کرنا چاہیدا اے۔\footnote{\olref[sfr][infinite][dedekind]{sec}
وچ جمع، ضرب تے قوت دیاں تعریفاں نال ملا کے ویکھو۔} قدرتی عدداں وانگ
اسیں اک مُڈھلی صورت تے جانشین صورتاں تعریف کردے آں۔ پر ترتیبی عدداں
نال کم کردیاں سانوں \emph{حدّی} صورتاں وی دسّنیاں پیندیاں نیں۔

$V_\alpha$ والیاں ایہ تعریفاں اک اہم پڑاؤ ہون گیاں۔ اسیں غیر رسمی طور
اُتے سیٹاں دی درجہ بندی دی وجہ ایویں دسی اے کہ سیٹ \emph{مرحلیاں} وچ بندے
نیں۔ $V_\alpha$ والے سیٹ حقیقت وچ اوہی مرحلے نیں۔ پر ضروری گل ایہ اے
کہ ایہ مرحلیاں دی \emph{اندرلی} خاصیت اے۔ اسیں نظریے دا اک ممکن
\emph{نمونہ} پیش کرن دی تھاں مرحلے اپنے سیٹ نظریے \emph{اندر} تعریف
کر دتے ہون گے۔
\endgroup
% /OLP-0559
% OLP-0560 content/set-theory/spine/recursion.tex
\begingroup
\def\olfilename{recursion}
\expandafter\def\csname ol@sectioncs\endcsname{section}

\olfileid{sth}{spine}{recursion}
\olsection{لامتناہی درجیاں والی بازگشت دے قضیے}

پر سانوں سبھ توں پہلاں ایہ پرکھنا اے کہ \olref[valpha]{defValphas}
اک کامیاب تعریف اے۔ ہور عام طور اُتے سانوں ثابت کرنا اے کہ لامتناہی
درجیاں والی بازگشت نال تعریف دین دی ہر کوشش کامیاب ہووے گی۔ ایس حصے دا ایہو مقصد اے۔

\emph{دھیان رہے۔ ایہ کم اوکھا اے}۔ پر ساری گل دا حاصل بڑا سدھا اے۔
لامتناہی درجیاں اُتے استقرا تے بدل رکھن دا اصول، لامتناہی درجیاں والی
بازگشت دے (کئی روپاں دے) جائز ہون دی ضمانت دیندے نیں۔\footnote{یاد رہے
کہ سارے فارمولیاں تے ٹرماں دے پیرامیٹر ہو سکدے نیں (جدوں تک صاف طور اُتے
کجھ ہور نہ آکھیا جاوے)۔}

\begin{defn}
منّ لو $\tau(x)$ اک ٹرم اے، $f$ اک فنکشن اے، تے $\alpha$ اک ترتیبی عدد اے۔
اسیں $f$ نوں \emph{$\alpha$ تک دا لگ بھگ روپ} آکھدے آں، $\tau$ لئی
اُدوں تے صرف اُدوں جدوں $\dom{f} = \alpha$ تے
$(\forall \beta \in \alpha)f(\beta) = \tau(\funrestrictionto{f}{\beta})$
دونیں شرطاں پوریاں ہون۔
\end{defn}

\begin{lem}[حدّ بندی والی بازگشت]\ollabel{transrecursionfun}
کسے وی ٹرم $\tau(x)$ تے کسے وی ترتیبی عدد $\alpha$ لئی، $\alpha$ تک دا
اکّو اک لگ بھگ روپ اے، $\tau$ لئی۔
\end{lem}
\begin{proof}
اسیں وکھاواں گے کہ کسے وی $\gamma \leq \alpha$ لئی، $\gamma$ تک دا
اکّو اک لگ بھگ روپ اے۔

پہلاں اسیں یکتائی ثابت کردے آں۔ منّ لو $g$ تے $h$ بالترتیب $\gamma$ تے
$\delta$ تک دے لگ بھگ روپ نیں۔ اوہناں دے مدخلاں اُتے لامتناہی درجیاں
والا استقرا وکھاندا اے کہ $g(\beta) = h(\beta)$ کسے وی
$\beta \in \dom{g} \cap \dom{h} = \gamma \cap \delta = \min(\gamma, \delta)$
لئی۔ سو ساڈے لگ بھگ روپ (جے موجود ہون) یکتا نیں تے ساریاں قدراں اُتے
اکّو نتیجہ دیندے نیں۔

موجودگی ثابت کرن لئی ہُن اسیں سادہ لامتناہی استقرا
(\olref[ordinals][opps]{simpletransrecursion}) ترتیبی عدداں
$\delta \leq \alpha$ اُتے ورتدے آں۔

صاف اے کہ خالی فنکشن $\emptyset$ تک دا لگ بھگ روپ اے۔

جے $g$، $\gamma$ تک دا لگ بھگ روپ اے، تاں $g \cup
\{\tuple{\gamma, \tau(g)}\}$، $\ordsucc{\gamma}$ تک دا لگ بھگ روپ اے۔

جے $\gamma$ حدّی ترتیبی عدد اے تے $g_\delta$، $\delta$ تک دا لگ بھگ
روپ اے سارے $\delta < \gamma$ لئی، تاں منّ لو
$g = \bigcup_{\delta \in \gamma} g_\delta$۔ ایہ فنکشن اے، کیوں جے ساڈے
وکھ وکھ $g_\delta$ ساریاں قدراں اُتے اکّو نتیجہ دیندے نیں۔ تے جے
$\delta \in \gamma$، تاں $g(\delta) = g_{\ordsucc{\delta}}(\delta) =
\tau(\funrestrictionto{g_{\ordsucc{\delta}}}{\delta}) =
\tau(\funrestrictionto{g}{\delta})$۔

ایس نال لامتناہی درجیاں اُتے استقرا والا ثبوت پورا ہویا۔
\end{proof}

جے اسیں فنکشن دی تھاں اک \emph{ٹرم} تعریف کرن دی اجازت دیئیے، تاں
پچھلے نتیجے توں حدّ $\alpha$ ہٹا سکدے آں۔ ہِیٹھلے نتیجے دے بیان تے ثبوت
وچ جدوں $\sigma$ اک ٹرم ہووے، اسیں
$\funrestrictionto{\sigma}{\alpha} = \Setabs{\tuple{\beta, \sigma(\beta)}}{\beta \in \alpha}$
رکھدے آں۔

\begin{thm}[عام بازگشت]\ollabel{transrecursionschema}
کسے وی ٹرم $\tau(x)$ لئی اسیں اک ٹرم $\sigma(x)$ صاف طور اُتے تعریف
کر سکدے آں، جیہدے لئی $\sigma(\alpha) = \tau(\funrestrictionto{\sigma}{\alpha})$
کسے وی ترتیبی عدد~$\alpha$ لئی۔
\end{thm}

\begin{proof}
ہر $\alpha$ لئی، \olref{transrecursionfun} نال
$\alpha$ تک دا اکّو اک لگ بھگ روپ، $f_\alpha$، اے،~$\tau$ لئی۔ $\sigma(\alpha)$ نوں
$f_{\ordsucc{\alpha}}(\alpha)$ تعریف کرو۔ ہُن:
	\begin{align*}
		\sigma(\alpha) &=
		f_{\ordsucc{\alpha}}(\alpha) \\&=
		\tau(\funrestrictionto{f_{\ordsucc{\alpha}}}{\alpha}) \\&=
		\tau(\Setabs{\tuple{\beta, f_{\ordsucc{\alpha}}(\beta)}}{\beta \in \alpha}) \\&=
		\tau(\Setabs{\tuple{\beta, f_{\ordsucc{\beta}}(\beta)}}{\beta \in \alpha})\\&=
		\tau(\funrestrictionto{\sigma}{\alpha})
	\end{align*}
دھیان دیو کہ $f_{\ordsucc{\beta}}(\beta) = f_{\ordsucc{\alpha}}(\beta)$
سارے $\beta < \alpha$ لئی، جیویں \olref{transrecursionfun} وچ۔
\end{proof}
\noindent
دھیان دیو کہ \olref{transrecursionschema} اک \emph{خاکہ} اے۔ ضروری گل
ایہ اے کہ اسیں $\sigma$ توں اک فنکشن، یعنی اک خاص قسم دا \emph{سیٹ}،
تعریف کرن دی اُمید نہیں رکھ سکدے، کیوں جے فیر $\dom{\sigma}$ سارے
ترتیبی عدداں دا سیٹ ہووے گا، جیہڑا بورالی فورتی دے پیراڈوکس
(\olref[ordinals][basic]{buraliforti}) دے اُلٹ اے۔

پر ہجے وکھاؤنا اے کہ \olref{transrecursionschema}، ساڈی $V_\alpha$ والے
سیٹاں دی تعریف دا جواز دیندا اے۔ ایہ جھٹ صاف نہ ہووے، پر لامتناہی
درجیاں والی بازگشت دے اک آخری، سادے روپ نال صاف ہو جاوے گا۔

\begin{thm}[سادہ بازگشت]\ollabel{simplerecursionschema}
کسے وی ٹرماں $\tau(x)$ تے $\theta(x)$ تے کسے وی سیٹ $A$ لئی، اسیں
اک ٹرم $\sigma(x)$ صاف طور اُتے تعریف کر سکدے آں جیہدے لئی:
\begin{align*}
	\sigma(\emptyset) &= A\\
	\sigma(\ordsucc{\alpha}) &= \tau(\sigma(\alpha)) &&
		\text{کسے وی ترتیبی عدد لئی }\alpha\\
	\sigma(\alpha) &= \theta(\ran{\funrestrictionto{\sigma}{\alpha}})&&
	\text{جدوں }\alpha\text{ حدّی ترتیبی عدد ہووے}
\end{align*}
\end{thm}

\begin{proof}
اسیں ٹرم $\xi(x)$ دی ایویں تعریف توں شروع کردے آں:
%t $\xi(x) = A$ if $x$ is not a function whose domain is an ordinal;
%otherwise let:
\[
	\xi(x) =
	\begin{cases}
			A &\text{جے $x$ ایہو جیہا فنکشن نہیں جیہدا}\\
			&\hspace{1em}\text{دائرہ ترتیبی عدد ہووے۔ نہیں تاں:}\\
			\tau(x(\alpha)) & \text{جے $\dom{x} = \ordsucc{\alpha}$}\\
			\theta(\ran{x}) & \text{جے $\dom{x}$ حدّی ترتیبی عدد اے}
	\end{cases}
\]
\olref{transrecursionschema} نال اک ٹرم $\sigma(x)$ اے جیہدے لئی
$\sigma(\alpha) = \xi(\funrestrictionto{\sigma}{\alpha})$ ہر ترتیبی
عدد $\alpha$ لئی۔ نالے $\funrestrictionto{\sigma}{\alpha}$ اک فنکشن
اے جیہدا دائرہ $\alpha$ اے۔ اسیں وکھاندے آں کہ $\sigma$ دیاں لوڑیندیاں خاصیتاں نیں،
سادہ لامتناہی استقرا (\olref[ordinals][opps]{simpletransrecursion}) نال۔

پہلاں، $\sigma(\emptyset) = \xi(\emptyset) = A$۔

فیر، $\sigma(\ordsucc{\alpha}) = \xi(\funrestrictionto{\sigma}{\ordsucc{\alpha}}) = \tau(\funrestrictionto{\sigma}{\ordsucc{\alpha}}(\alpha)) = \tau(\sigma(\alpha))$۔

آخر وچ، $\sigma(\alpha) = \xi(\funrestrictionto{\sigma}{\alpha}) = \theta(\ran{\funrestrictionto{\sigma}{\alpha}})$، جدوں $\alpha$ حدّی ترتیبی عدد اے۔
\end{proof}
\noindent
ہُن \olref[valpha]{defValphas} دا جواز دین لئی بس $A = \emptyset$،
$\tau(x) = \Pow{x}$ تے $\theta(x) = \bigcup x$ رکھو۔ آخر ایس نال
$V_\alpha$ والے سیٹاں دی تعریف دا جواز مل گیا!{}
\endgroup
% /OLP-0560
% OLP-0561 content/set-theory/spine/stagesbasics.tex
\begingroup
\def\olfilename{stagesbasics}
\expandafter\def\csname ol@sectioncs\endcsname{section}

\olfileid{sth}{spine}{Valphabasic}
\olsection{مرحلیاں دیاں بنیادی خاصیتاں}

$V_\alpha$ والے سیٹاں دی تعریف دی بنیادی اہمیت سمجھاؤن لئی اسیں اوہناں
بارے کجھ بنیادی نتیجے دیواں گے۔ پہلاں اک تعریف دیندے آں:\footnote{’’ذیلی
سیٹاں تحت بند‘‘ لئی کوئی معیاری اصطلاح نہیں۔ اصل وچ ایہدے لئی انگریزی ناں «potent» اے، جیہڑا
\citet{ButtonLT1} ورتدا اے۔}
\begin{defn}
سیٹ $A$ \emph{اپنے رکناں دے ذیلی سیٹاں تحت بند} اُدوں تے صرف اُدوں اے
جدوں $\forall x((\exists y \in A)x \subseteq y \lif x \in A)$۔
\end{defn}
\begin{lem}\ollabel{Valphabasicprops}
ہر ترتیبی عدد $\alpha$ لئی:
\begin{enumerate}
	\item\ollabel{Valphatrans} ہر $V_\alpha$ متعدی اے۔
	\item\ollabel{Valphapotent} ہر $V_\alpha$ اپنے رکناں دے ذیلی سیٹاں تحت بند اے۔
	\item\ollabel{Valphacum} جے $\gamma \in \alpha$، تاں $V_\gamma
	\in V_\alpha$ (تے ایس لئی $V_\gamma \subseteq V_\alpha$ وی،
	\olref{Valphatrans} نال)۔
\end{enumerate}
\end{lem}

\begin{proof}
اسیں ایہ لامتناہی درجیاں اُتے (اکو ویلے) استقرا نال ثابت کردے آں۔
استقرا لئی منّ لو \olref{Valphatrans}--\olref{Valphacum} ہر ترتیبی
عدد $\beta < \alpha$ لئی قائم نیں۔

$\alpha = \emptyset$ والی صورت ظاہر اے۔

منّ لو $\alpha = \ordsucc{\beta}$۔ \olref{Valphacum} وکھاؤن لئی، جے
$\gamma \in \alpha$، تاں مفروضے نال $V_\gamma \subseteq V_\beta$،
سو $V_\gamma \in \Pow{V_\beta} = V_\alpha$۔
\olref{Valphapotent} وکھاؤن لئی منّ لو $A \subseteq B \in V_\alpha$،
یعنی $A \subseteq B \subseteq V_\beta$۔ تاں $A \subseteq V_\beta$،
سو $A \in V_\alpha$۔ \olref{Valphatrans} وکھاؤن لئی دھیان دیو کہ
جے $x \in A \in V_\alpha$، تاں $A \subseteq V_\beta$، سو $x \in V_\beta$،
سو $x \subseteq V_\beta$ کیوں جے مفروضے نال $V_\beta$ متعدی اے۔
ایس لئی $x \in V_\alpha$۔

منّ لو $\alpha$ حدّی ترتیبی عدد اے۔ \olref{Valphacum} وکھاؤن لئی،
جے $\gamma \in \alpha$، تاں $\gamma \in \ordsucc{\gamma} \in \alpha$۔
سو مفروضے نال $V_\gamma \in V_{\ordsucc{\gamma}}$، تے ایس لئی
$V_\gamma \in \bigcup_{\beta \in \alpha} V_\beta = V_\alpha$۔
\olref{Valphatrans} تے \olref{Valphapotent} وکھاؤن لئی بس ویکھو کہ
متعدی (یا اپنے رکناں دے ذیلی سیٹاں تحت بند) سیٹاں دا اتحاد وی
متعدی (بالترتیب اپنے رکناں دے ذیلی سیٹاں تحت بند) اے۔
\end{proof}

\begin{lem}\ollabel{Valphanotref}
ہر ترتیبی عدد $\alpha$ لئی، $V_\alpha \notin V_\alpha$۔
\end{lem}

\begin{proof}
لامتناہی درجیاں اُتے استقرا نال۔ ظاہر اے کہ $V_\emptyset \notin V_\emptyset$۔

جے $V_{\ordsucc{\alpha}} \in V_{\ordsucc{\alpha}} = \Pow{V_\alpha}$،
تاں $V_{\ordsucc{\alpha}} \subseteq V_\alpha$۔ تے چونکہ
$V_\alpha \in V_{\ordsucc{\alpha}}$، \olref{Valphabasicprops} نال،
ایس لئی $V_\alpha \in V_\alpha$۔ اُلٹے پاسے، جے $V_\alpha \notin V_\alpha$،
تاں $V_{\ordsucc{\alpha}} \notin V_{\ordsucc{\alpha}}$۔

جے $\alpha$ حدّی عدد اے تے $V_\alpha \in V_\alpha = \bigcup_{\beta \in \alpha}V_\beta$،
تاں $V_\alpha \in V_\beta$ کسے $\beta \in \alpha$ لئی۔ پر فیر
$V_\beta \in V_\alpha$ وی، ایس لئی $V_\beta \in V_\beta$،
\olref{Valphabasicprops} دو واری ورت کے۔ اُلٹے پاسے، جے $V_\beta
\notin V_\beta$ ہر $\beta \in \alpha$ لئی، تاں $V_\alpha \notin V_\alpha$۔
\end{proof}

\begin{cor}
کسے وی ترتیبی عدداں $\alpha, \beta$ لئی، $\alpha \in \beta$ اُدوں
تے صرف اُدوں اے جدوں $V_\alpha \in V_\beta$۔
\end{cor}

\begin{proof}
\Olref{Valphabasicprops} اک پاسا دیندا اے۔ اُلٹے پاسے منّ لو
$V_\alpha \in V_\beta$۔ تاں $\alpha \neq \beta$، \olref{Valphanotref}
نال۔ تے $\beta \notin \alpha$، کیوں جے نہیں تاں $V_\beta \in V_\alpha$
ہوندا، تے ایس لئی $V_\beta \in V_\beta$، \olref{Valphabasicprops}
دو واری ورت کے، جیہڑا \olref{Valphanotref} دے اُلٹ اے۔
سو $\alpha \in \beta$ تِن صورتاں دے اصول نال۔
\end{proof}
\noindent
ایس سارے توں اسیں ہر $V_\alpha$ نوں درجہ بندی دا $\alpha$واں مرحلہ
منّ سکدے آں۔ وجہ ایہ اے۔

بے شک ساڈے $V_\alpha$ والے سیٹ \emph{مرحلہ وار} بندے سمجھ سکدے آں، کیوں جے
ترتیبی عدداں دی ساڈی ورتوں مرحلیاں دی دہرائی نوں ظاہر کردی اے۔ نالے جے اک
مرحلہ دوجے توں پہلاں بندا اے، یعنی $V_\beta \in V_\alpha$، یعنی $\beta \in \alpha$،
تاں سیٹ بنن دی ساڈی کارروائی \emph{اکٹھی کردی جاندی} اے، کیوں جے
$V_\beta \subseteq V_\alpha$۔ آخر اسیں اوہناں سیٹاں دے \emph{سارے} ممکن
اکٹھ بناندے آں جیہڑے کسے پچھلے مرحلے اُتے موجود سن، کیوں جے کوئی وی
جانشین مرحلہ $V_{\ordsucc{\alpha}}$ اپنے پچھلے مرحلے $V_\alpha$ دے ذیلی سیٹاں
دا سیٹ اے۔

مختصر گل ایہ اے کہ $\ZFminus$ نال اسیں سیٹاں دی اکٹھا کرن والی مرحلہ وار
درجہ بندی دا اپنا خیال بیان کرن دا کم \emph{تقریباً} پورا کر چکے آں۔
(بھانویں، بے شک، ہجے بدل رکھن دے اصول دا جواز دینا اے۔)
\endgroup
% /OLP-0561
% OLP-0562 content/set-theory/spine/foundation.tex
\begingroup
\def\olfilename{foundation}
\expandafter\def\csname ol@sectioncs\endcsname{section}

\olfileid{sth}{spine}{foundation}

\olsection{بنیاد دا مسلّمہ}

کم ہجے \emph{تقریباً} پورا اے، \emph{بالکل} پورا نہیں، کیوں جے
$\ZFminus$ وچ کوئی گل ضمانت نہیں دیندی کہ \emph{ہر} سیٹ کسے $V_\alpha$
وچ اے، یعنی ہر سیٹ کسے مرحلے اُتے بندا اے۔

ہُن اک کافی سدھے (ریاضیاتی) مطلب وچ سانوں ایس گل دی \emph{پرواہ} نہیں
کہ درجہ بندی توں باہر وی سیٹ نیں یا نہیں۔ (جے باہر کوئی سیٹ نیں، تاں
اسیں اوہناں نوں چھڈ سکدے آں۔) پر اسیں سیٹ دے اپنے \emph{تصور} دا جواز
ایس خیال نال دتا اے کہ ہر سیٹ کسے مرحلے اُتے بندا اے
(\stageshier{} نوں \olref[z][story]{sec} وچ ویکھو)۔ سو اسیں ایہ امکان
روکنا چاہواں گے کہ کجھ سیٹ درجہ بندی توں باہر ہون۔ ایس لئی اک نواں مسلّمہ
شامل کرنا پئے گا جیہڑا یقینی بناوے کہ ہر سیٹ درجہ بندی وچ کِتے نہ کِتے اے۔

چونکہ $V_\alpha$ والے سیٹ ساڈے مرحلے نیں، اسیں ہِیٹھلی گل نوں مسلّمہ
بناؤن بارے سوچ سکدے آں:

\begin{defish}
\emph{باقاعدگی۔} $\forall A \exists \alpha\, A \subseteq V_\alpha$
\end{defish}

ایہ بالکل معقول طریقہ ہووے گا۔ پر اوہناں وجوہاں کر کے جیہڑیاں اگلے حصے وچ
دساں گے، اسیں ایہدی تھاں اک ہور مسلّمہ منّاں گے:

\begin{axiom}[بنیاد]
$(\forall A \neq \emptyset)(\exists B \in A)A \cap B = \emptyset$۔
\end{axiom}

کجھ محنت نال اسیں ($\ZFminus$ وچ) وکھا سکدے آں کہ بنیاد توں باقاعدگی نکلدی اے:
\begin{defn}
ہر سیٹ $A$ لئی منّ لو:
\begin{align*}
	\text{cl}_0(A) &= A,\\
	\text{cl}_{n+1}(A) &= \bigcup \text{cl}_n(A),\\
	\text{trcl}(A) &= \bigcup_{n < \omega} \text{cl}_{n}(A).
\end{align*}
اسیں $\text{trcl}(A)$ نوں $A$ دی \emph{متعدی بندش} آکھدے آں۔
\end{defn}
\noindent
’’متعدی بندش‘‘ ناں ٹھیک اے:
\begin{prop}\ollabel{subsetoftrcl}
$A \subseteq \trcl{A}$ تے $\trcl{A}$ اک متعدی سیٹ اے۔
\end{prop}

\begin{proof}
ظاہر اے کہ $A = \text{cl}_0(A) \subseteq \text{trcl}(A)$۔ تے جے
$x \in b \in \trcl{A}$، تاں $b \in \text{cl}_n(A)$ کسے $n$ لئی،
سو $x \in \text{cl}_{n+1}(A) \subseteq \text{trcl}(A)$۔
\end{proof}

\begin{lem}\ollabel{lem:TransitiveWellFounded}
جے $A$ متعدی سیٹ اے، تاں کوئی $\alpha$ اے جیہدے لئی $A \subseteq V_\alpha$۔
\end{lem}

\begin{proof}
’’$\supstrict(X)$‘‘ دی تعریف \olref[ordinals][opps]{defsupstrict} توں
یاد کر کے دو سیٹ تعریف کرو:
\begin{align*}
	D  &= \Setabs{x \in A}{\forall \delta\ x \nsubseteq V_\delta}\\
	\alpha &= \supstrict\Setabs{\delta}{(\exists x \in A)
	(x \subseteq V_\delta \land (\forall \gamma \in \delta)x \nsubseteq V_\gamma)}
\end{align*}
منّ لو $D = \emptyset$۔ سو جے $x \in A$، تاں کوئی $\delta$ اے
جیہدے لئی $x \subseteq V_\delta$ تے، ترتیبی عدداں دی حُسنِ ترتیب نال،
$(\forall \gamma \in \delta)x \nsubseteq V_\gamma$۔ ایس لئی $\delta \in \alpha$
تے سو $x \in V_\alpha$، \olref[spine][Valphabasic]{Valphabasicprops} نال۔
سو $A \subseteq V_{\alpha}$، جیویں سانوں چاہیدا سی۔

سو $D = \emptyset$ وکھاؤنا کافی اے۔ تضاد کڈھن لئی منّ لو ایویں نہیں۔
بنیاد نال کوئی $B \in D \subseteq A$ اے جیہدے لئی $D \cap B = \emptyset$۔
جے $x \in B$، تاں $x \in A$، کیوں جے $A$ متعدی اے۔ تے چونکہ
$x \notin D$، ایس لئی $\exists \delta\ x \subseteq V_\delta$۔ سو ہُن منّ لو
\[
	\beta = \supstrict\Setabs{\delta}{(\exists x \in b)(x \subseteq V_\delta \land (\forall \gamma < \delta)x \nsubseteq V_\gamma)}.
\]
پہلاں وانگ، $B \subseteq V_\beta$، جیہڑا $B \in D$ دے دعوے دے اُلٹ اے۔
\end{proof}

\begin{thm}\ollabel{zfentailsregularity}
باقاعدگی قائم اے۔
\end{thm}

\begin{proof}
$A$ مقرر کرو۔ ہُن $A \subseteq \trcl{A}$، \olref{subsetoftrcl} نال،
تے ایہ متعدی اے۔ سو کوئی $\alpha$ اے جیہدے لئی
$A \subseteq \trcl{A} \subseteq V_\alpha$، \olref{lem:TransitiveWellFounded} نال۔
\end{proof}

ایہناں نتیجیاں توں پتا لگدا اے کہ $\ZFminus$ شرطیہ
$\emph{بنیاد}\Rightarrow\emph{باقاعدگی}$ ثابت کردا اے۔
\olref[spine][rank]{zfminusregularityfoundation} وچ اسیں وکھاواں گے کہ
$\ZFminus$، $\emph{باقاعدگی}\Rightarrow\emph{بنیاد}$ ثابت کردا اے۔
سو بنیاد تے باقاعدگی ($\ZFminus$ دے لحاظ نال) \emph{ہم ارز} نیں۔
ایہدا مطلب اے کہ $\ZFminus$ دتا ہووے، تاں اسیں بنیاد دا جواز ایہ
ویکھ کے دے سکدے آں کہ ایہ باقاعدگی دے ہم ارز اے۔ تے باقاعدگی دا جواز
\stageshier{} دی بنیاد اُتے سدھا دے سکدے آں۔
\endgroup
% /OLP-0562
% OLP-0563 content/set-theory/spine/zf.tex
\begingroup
\def\olfilename{zf}
\expandafter\def\csname ol@sectioncs\endcsname{section}

\olfileid{sth}{spine}{zf}
\olsection{\texorpdfstring{$\Z$ تے $\ZF$۔ اک اہم پڑاؤ}{Z تے ZF۔ اک اہم پڑاؤ}}

بنیاد دے مسلّمے نال اسیں اک ہور اہم پڑاؤ اُتے پہنچدے آں۔ اسیں
$\Zminus$ تے $\ZFminus$ نظریے ویکھے نیں تے آکھیا سی کہ ایہ کجھ
نظریے نیں جیہناں وچوں اک خاص شے ’’گھٹائی‘‘ اے۔ اوہ خاص شے بنیاد دا
مسلّمہ اے۔ سو:

\begin{defn}
نظریہ $\Z$، $\Zminus$ وچ بنیاد دا مسلّمہ شامل کردا اے۔ سو اوہدے
مسلّمے ایہ نیں۔ عنصراں نال تعیّن، اتحاد، جوڑیاں، ذیلی سیٹاں دے سیٹ،
لامتناہیت، بنیاد، تے وکھ کرن دے خاکے دیاں ساریاں صورتاں۔

نظریہ $\ZF$، $\ZFminus$ وچ بنیاد دا مسلّمہ شامل کردا اے۔
دوجے لفظاں وچ، $\ZF$،~$\Z$ وچ بدل رکھن دے خاکے دیاں ساریاں صورتاں
شامل کردا اے۔
\end{defn}

فیر وی اک سوال تہانوں سُجھیا ہو سکدا اے۔ جے باقاعدگی، $\ZFminus$ دے
لحاظ نال، بنیاد دے ہم ارز اے تے باقاعدگی دا جواز صاف اے، تاں
اینی لمّی راہ کیوں لئے، تے باقاعدگی دی تھاں بنیاد نوں اپنا بنیادی مسلّمہ کیوں منّیئے؟

تاریخی وجوہاں (کس نے کیہ تے کدوں بیان کیتا) اک پاسے رکھ کے، بنیادی وجہ
ایہ اے کہ بنیاد دا مسلّمہ~$V_\alpha$ والے سیٹاں دی تعریف ورتے بناں پیش
ہو سکدا اے۔ اوہ تعریف \olref[recursion]{sec} دے سارے کم اُتے ٹکی سی۔
اوہدا جواز وکھاؤن لئی سانوں لامتناہی درجیاں والی بازگشت ثابت کرنی پئی سی۔
پر اوہ بازگشت ثابت کرن لئی اسیں \emph{بدل رکھن دا اصول} ورتیا سی۔
سو بھانویں بنیاد تے باقاعدگی $\ZFminus$ دے لحاظ نال ہم ارز نیں،
$\Zminus$ دے لحاظ نال ہم ارز نہیں۔

اصل وچ گل ایس سادی ٹِپّن توں وی اگے جاندی اے۔ بھانویں ایہ کتاب دے
دائرے توں کافی باہر اے، پتا لگدا اے کہ $\Zminus$ تے $\Z$ دونیں
$V_\alpha$ والے سیٹاں دی تعریف لئی بہتے کمزور نیں۔ سو جے تسی صرف
$\Z$ وچ کم کر رہے او، تاں باقاعدگی (جیویں اسیں بیان کیتی اے) دا
\emph{مطلب} ای نہیں بندا۔ ایس لئی ساڈا مقرر کیتا مسلّمہ بنیاد اے، باقاعدگی نہیں۔

ہُن توں اسیں $\ZF$ وچ کم کراں گے (جے کجھ ہور نہ آکھیا جاوے)، تے
ایس بارے ہور کوئی ٹِپّن نہیں کراں گے۔
\endgroup
% /OLP-0563
% OLP-0564 content/set-theory/spine/rank.tex
\begingroup
\def\olfilename{rank}
\expandafter\def\csname ol@sectioncs\endcsname{section}

\olfileid{sth}{spine}{rank}
\olsection{درجہ}

ہُن جدوں اسیں مرحلیاں نوں $V_\alpha$ والے سیٹ تعریف کر چکے آں تے سانوں
پتا اے کہ ہر سیٹ کسے مرحلے دا ذیلی سیٹ اے، اسیں سیٹ دا \emph{درجہ}
تعریف کر سکدے آں۔ سمجھ مطابق $A$ دا درجہ اوہ پہلا ویلا اے جدوں $A$
بندا اے۔ ہور ٹھیک لفظاں وچ:

\begin{defn}\ollabel{defnsetrank}
ہر سیٹ $A$ لئی، $\setrank{A}$ اوہ سبھ توں چھوٹا ترتیبی عدد $\alpha$ اے
جیہدے لئی $A \subseteq V_\alpha$۔
\end{defn}
\begin{prop}\ollabel{ranksexist}
	$\setrank{A}$ کسے وی $A$ لئی موجود اے۔
\end{prop}
\begin{proof}
	ایہ مشق چھڈّی اے۔
\end{proof}
\begin{prob}
	\olref[sth][spine][rank]{ranksexist} ثابت کرو۔
\end{prob}
\noindent
درجیاں دا حُسنِ ترتیب سانوں کجھ اہم نتیجے ثابت کرن دیوے گا:
\begin{prop}\ollabel{valphalowerrank}
کسے وی ترتیبی عدد $\alpha$ لئی، $V_\alpha = \Setabs{x}{\setrank{x} \in \alpha}$۔
\end{prop}

\begin{proof}
جے $\setrank{x} \in \alpha$، تاں $x \subseteq V_{\setrank{x}} \in V_\alpha$،
سو $x \in V_\alpha$، کیوں جے $V_\alpha$ اپنے رکناں دے ذیلی سیٹاں تحت
بند اے (\olref[Valphabasic]{Valphabasicprops} کئی واری ورت کے)۔
اُلٹے پاسے، جے $x \in V_\alpha$، تاں $x \subseteq V_\alpha$، سو
$\setrank{x} \leq \alpha$۔ ہُن سادہ لامتناہی استقرا وکھاندا اے کہ
$x \notin V_\alpha$۔
\end{proof}
\begin{prob}
	\olref[sth][spine][rank]{valphalowerrank} وچ دسیا سادہ لامتناہی استقرا پورا کرو۔
\end{prob}

\begin{prop}\ollabel{rankmemberslower}
جے $B \in A$، تاں $\setrank{B} \in \setrank{A}$۔
\end{prop}

\begin{proof}
$A \subseteq V_{\setrank{A}} = \Setabs{x}{\setrank{x} \in \setrank{A}}$،
\olref{valphalowerrank} نال۔
\end{proof}
\noindent
ایس گل نوں ورت کے اسیں اک نتیجہ ثابت کر سکدے آں جیہڑا سانوں استقرا دی اک
صورت نال \emph{سارے سیٹاں} بارے گلاں ثابت کرن دیوے گا:

\begin{thm}[$\in$ اُتے استقرا دا خاکہ]
ہر فارمولے $\phi$ لئی:
\[
	\forall A((\forall x \in A)\phi(x) \lif \phi(A)) \lif \forall A \phi(A).
\]
\end{thm}

\begin{proof}
اسیں عکسِ نقیض ثابت کراں گے۔ سو منّ لو $\lnot \forall A \phi(A)$۔
لامتناہی درجیاں اُتے استقرا (\olref[ordinals][basic]{ordinductionschema})
نال کوئی سبھ توں گھٹ ممکن درجے والا غیر-$\phi$ اے، یعنی کوئی $A$
جیہدے لئی $\lnot \phi(A)$ تے $\forall x(\setrank{x} \in \setrank{A} \lif \phi(x))$۔
ہُن جے $x \in A$، تاں $\setrank{x} \in \setrank{A}$،
\olref{rankmemberslower} نال، ایس لئی $\phi(x)$۔ یعنی
$(\forall x \in A)\phi(x) \land \lnot \phi(A)$۔
\end{proof}\noindent ایس طاقتور نتیجے دی غیر رسمی سمجھ ایویں اے۔
$\phi$ نوں \emph{وراثتی} آکھو اُدوں تے صرف اُدوں جدوں کسے سیٹ دا ہر
!!{element}، $\phi$ ہووے تاں سیٹ آپ وی $\phi$ ہووے۔
فیر $\in$ اُتے استقرا تہانوں ایہ دسدا اے۔ جے $\phi$ وراثتی اے،
تاں ہر سیٹ $\phi$ اے۔

درجیاں دی گل (ہُن لئی) مکاؤن واسطے اسیں کجھ دعوے ثابت کراں گے
جیہناں ول پہلے کئی واری اشارہ کیتا اے۔

\begin{prop}\ollabel{ranksupstrict}
$\setrank{A} = \supstrict_{x \in A}\setrank{x}$۔
\end{prop}

\begin{proof}
منّ لو $\alpha = \supstrict_{x \in A}\setrank{x}$۔
\olref{rankmemberslower} نال $\alpha \leq \setrank{A}$۔ پر جے $x \in A$،
تاں $\setrank{x} \in \alpha$، تے ایس لئی $x \in V_\alpha$،
\olref{valphalowerrank} نال۔ سو $A \subseteq V_\alpha$، یعنی
$\setrank{A} \leq \alpha$۔ ایس لئی $\setrank{A} = \alpha$۔
\end{proof}

\begin{cor}\ollabel{ordsetrankalpha}
کسے وی ترتیبی عدد $\alpha$ لئی، $\setrank{\alpha} = \alpha$۔
\end{cor}

\begin{proof}
لامتناہی درجیاں اُتے استقرا لئی منّ لو کہ $\setrank{\beta} = \beta$ سارے
$\beta \in \alpha$ لئی۔ ہُن $\setrank{\alpha} = \supstrict_{\beta \in
\alpha}\setrank{\beta} = \supstrict_{\beta \in \alpha}\beta = \alpha$،
\olref{ranksupstrict} نال۔
%	First note that $\setrank{\alpha} \neq \beta$ for any $\beta \in
%	\alpha$. For otherwise we would have $\beta = \setrank{\beta} \in
%	\setrank{\alpha} = \beta$ by \olref{rankmemberslower}, i.e.,
%	$\beta \in \beta$, a contradiction. 
%
%	Now note that $\alpha \subseteq V_\alpha$. For $\alpha =
%	\Setabs{\beta}{\beta \in \alpha}$ by
%	\olref[sfr][ordinals][basic]{ordissetofsmallerord}. And if
%	$\beta \in \alpha$ then $\beta \subseteq V_{\setrank{\beta}} =
%	V_\beta \in V_\alpha$, hence $\beta \in V_\alpha$, by
%	\olref[sfr][spine][Valphabasic]{Valphabasicprops} twice. 
\end{proof}

آخر وچ، \olref[foundation]{sec} دے آخر وچ جیہڑا نتیجہ دین دی گل
کیتی سی، اوہدا اک سدھا ثبوت ایہ اے، یعنی $\ZFminus$ شرطیہ
$\emph{باقاعدگی} \Rightarrow \emph{بنیاد}$ ثابت کردا اے۔ (دھیان دیو کہ
’’درجے‘‘ دا خیال تے \olref{rankmemberslower} ایس ثبوت وچ ورت سکدے آں،
کیوں جے، جیویں ایس حصے دے شروع وچ دسیا سی، اوہ
$\ZFminus + \text{باقاعدگی}$ ورت کے پیش ہو سکدے نیں۔)

\begin{prop}[$\ZFminus + \text{باقاعدگی}$ وچ کم کردیاں]\ollabel{zfminusregularityfoundation}
بنیاد دا مسلّمہ قائم اے۔
\end{prop}

\begin{proof}
$A \neq \emptyset$ مقرر کرو، تے $B \in A$ جیہدا درجہ سبھ توں گھٹ ممکن اے۔
جے $c \in B$، تاں $\setrank{c} \in \setrank{B}$،
\olref{rankmemberslower} نال، ایس لئی $c \notin A$، $B$ دی چُون کر کے۔
\end{proof}
\endgroup
% /OLP-0564

\OLEndChapterHook
\endgroup
% /OLP-0558
% OLP-0565 content/set-theory/replacement/replacement.tex
\begingroup
\def\olfilename{replacement}
\expandafter\def\csname ol@sectioncs\endcsname{section}

\olchapter{sth}{replacement}{بدل رکھن دا اصول}

% OLP-0566 content/set-theory/replacement/introduction.tex
\begingroup
\def\olfilename{introduction}
\expandafter\def\csname ol@sectioncs\endcsname{section}

\olfileid{sth}{replacement}{intro}
\olsection{جان پہچان}

بدل رکھن دا اصول اوہ مسلّمی خاکہ اے جیہڑا $\ZF$ تے~$\Z$ وچ فرق
پاندا اے۔ اسیں \crefrange{sth:ordinals::chap}{sth:spine::chap} وچ
ایہنوں لگاتار ورتیا اے۔ ایس باب وچ اسیں آخر ایہ سوال ویکھاں گے کہ
کی بدل رکھن دے اصول دا جواز اے؟

سوال نوں ہور صاف کرن لئی ایہ دھیان دین جوگی گل اے کہ بدل رکھن دا اصول
واقعی کافی \emph{طاقتور} اے۔ ایس باب وچ سانوں اندازہ ہووے گا کہ ایہ کِنّا
طاقتور اے (تے فیر \olref[sth][card-arithmetic][fix]{sec} وچ وی)۔
ایس توں سُجھدا اے کہ جواز دی واقعی لوڑ اے۔

اسیں دو قسماں دے جواز دی گل کراں گے۔ موٹی گل آکھیے تاں \emph{باہرلا}
جواز کسے مسلّمے نوں اوہدے حاصلاں راہیں جائز وکھاؤن دی کوشش اے۔
\emph{اندرلا} جواز ایہ وکھاؤن دی کوشش اے کہ زیر بحث ریاضیاتی تصور آپ
مسلّمے دی تائید کردے نیں۔ ایس باب وچ سانوں ایس گل دی ہور سمجھ آوے گی،
پر ایہ بس اک وڈی بحث دی شروعات اے۔ ہور لئی خاص کر کے
\citeauthor{Maddy1988a} (\citeyear{Maddy1988a} تے \citeyear{Maddy1988b})
ویکھو۔
\endgroup
% /OLP-0566
% OLP-0567 content/set-theory/replacement/strength.tex
\begingroup
\def\olfilename{strength}
\expandafter\def\csname ol@sectioncs\endcsname{section}

\olfileid{sth}{replacement}{strength}
\olsection{بدل رکھن دے اصول دی طاقت}

اسیں بدل رکھن دے اصول دی طاقت بارے اک سوکھی گل توں شروع کردے آں۔ جے اسیں $\Z$ توں اگے نہ ودھیے، تاں $\omega + \omega$ دے برابر یا اوہتوں وڈے کسے وی فون نوئمان ترتیبی عدد دی موجودگی ثابت نہیں کر سکدے۔

ایہدی وجہ دا اک خاکہ ویکھو۔ $\ZF$ وچ کم کردیاں سیٹ $V_{\omega+\omega}$ نوں ویکھو۔ ایہ سیٹ~$\Z$ دے اک \emph{نمونے} دا دائرہ بن دا اے۔ ایہ ویکھن لئی اسیں فارمولے نوں کسے دائرے \emph{تک محدود کرن} دی کجھ لکھت متعارف کراندے آں۔
\begin{defn}\ollabel{formularelativization}
	کسے وی سیٹ $M$ تے کسے وی فارمولے $ \phi$ لئی، $\phi^M$ اوہ فارمولا اے جیہڑا $\phi$ دے سارے مسوّراں نوں $M$ تک محدود کر کے ملدا اے۔ یعنی «$\lexists[x]$» دی تھاں «$(\lexists[x \in M])$» لکھو، تے
	«$\lforall[x]$» دی تھاں «$(\lforall[x \in M])$» لکھو۔
\end{defn}\noindent
وکھایا جا سکدا اے کہ ہر مسلّمے $\phi$ لئی جیہڑا~$\Z$ دا اے، $ \ZF \vdash
\phi^{V_{\omega+\omega}}$۔ پر \olref[spine][rank]{ordsetrankalpha} مطابق $\omega+\omega$، $V_{\omega+\omega}$ \emph{دا رکن نہیں}۔ سو $\Z$، $\omega+\omega$ دی غیرموجودگی نال سازگار اے۔

ایسے لئی اسیں \olref[ordinals][replacement]{sec} وچ آکھیا سی کہ
\olref[ordinals][ordtype]{thmOrdinalRepresentation} بدل رکھن دے اصول توں بناں ثابت نہیں ہو سکدا۔ کیوں جے~$\Z$ وچ اک صاف لکھی حُسنِ ترتیب تعریف کرنی سوکھی اے جیہدی ترتیب دی قسم، ساڈی سمجھ مطابق، $\omega+\omega$ \emph{ہونی چاہیدی اے}۔ اسیں \olref[ordinals][idea]{sec} وچ ایس دی اک غیر رسمی مثال دتی سی جدوں اسیں قدرتی عدداں اُتے ایہ ترتیب پیش کیتی سی۔
\begin{align*}
	n \lessdot m \text{ اُدوں تے صرف اُدوں جدوں }&\text{یا تاں }n < m\text{ تے }m-n\text{ جفت اے،}\\
	& \text{یا $n$ جفت اے تے $m$ طاق اے۔}
\end{align*}
پر جے $\omega+\omega$ موجود نہیں، تاں ایہ حُسنِ ترتیب کسے وی ترتیبی عدد نال ہم ساخت نہیں۔ سو $\Z$ توں
\olref[ordinals][ordtype]{thmOrdinalRepresentation} ثابت \emph{نہیں ہُندا}۔

ہُن گل الٹ کے ویکھو۔ بدل رکھن دا اصول $\omega+\omega$ دی موجودگی ثابت کرن دیندا اے، تے ایس لئی $V_{\omega+\omega}$ دی موجودگی وی ثابت کرن دیندا اے۔ گل ایتھے ای نہیں مکدی۔ اسیں جیہڑی \emph{وی} حُسنِ ترتیب تعریف کر سکیے، \olref[ordinals][ordtype]{thmOrdinalRepresentation}
دسدا اے کہ اوہدے نال ہم ساخت کوئی $\alpha$ موجود اے، تے ایس لئی $V_\alpha$ وی موجود اے۔ سو ایویں سیدھے ای بدل رکھن دا اصول ایہ پکا کردا اے کہ سیٹاں دا درجہ وار سلسلہ \emph{بہت اُچا} اے۔

اگلے کجھ حصیاں وچ، تے فیر \olref[card-arithmetic][fix]{sec} وچ، سانوں ہور چنگی طرح پتا لگے گا کہ بدل رکھن دا اصول درجہ وار سلسلے نوں \emph{کِنّا} اُچا بناؤن تے مجبور کردا اے۔ ہن لئی سیدھی گل ایہ اے کہ بدل رکھن دے اصول دا جواز \emph{واقعی} دینا پیندا اے!
\endgroup
% /OLP-0567
% OLP-0568 content/set-theory/replacement/extrinsic.tex
\begingroup
\def\olfilename{extrinsic}
\expandafter\def\csname ol@sectioncs\endcsname{section}

\olfileid{sth}{replacement}{extrinsic}
\olsection[باہرلے پہلو]{بدل رکھن دے اصول دے باہرلے پہلو}

اسیں بدل رکھن دے اصول دا \emph{باہرلا} جواز دین دی اک کوشش توں شروع کردے آں۔ بُولوس دی پیشکش ایہ اے۔
\begin{quote}
  [\ldots] بدل رکھن دے مسلّمے منّن دی وجہ بڑی سادہ اے۔ اوہناں توں کئی من پسند نتیجے نکلدے نیں تے (جیویں ہن تک دِسدا اے) کوئی ناپسند نتیجہ نہیں نکلدا۔ اکٹھا کرن والے مرحلہ وار تصور بارے قضیاں توں علاوہ، ایہناں نتیجیاں وچ لامتناہی عدداں دا اک تسلّی بخش، بھاویں مثالی نہیں، نظریہ تے اک بڑا من پسند نتیجہ شامل اے جیہڑا بنیاد والے تعلقاں اُتے استقرائی تعریفاں دا جواز دیندا اے۔
  \citep[229]{Boolos1971}
\end{quote}
بُولوس دی گل دا نچوڑ ایہ اے کہ سانوں بدل رکھن دے اصول دا جواز اوہدے حاصلاں توں دینا چاہیدا اے۔ جیہڑے خاص حاصل اوہ گِندا اے، اوہناں دی گل اسیں پچھلے کجھ باباں وچ کیتی اے۔ بدل رکھن دے اصول نال اسیں ثابت کیتا کہ فون نوئمان ترتیبی عدد حُسنِ ترتیب دی قسم دا خیال ورتن لئی بڑے چنگے بدل نیں (ایہو لامتناہی عدداں دا ساڈا «تسلّی بخش، بھاویں مثالی نہیں، نظریہ» اے)۔ بدل رکھن دے اصول نال اسیں $V_\alpha$ والے سیٹ وی تعریف کر سکے، درجے دا تصور قائم کر سکے، تے $\in$-استقرا ثابت کر سکے (ایہو ساڈے «اکٹھا کرن والے مرحلہ وار تصور بارے قضیے» نیں)۔ آخر وچ، بدل رکھن دا اصول سانوں لامتناہی درجیاں والی بازگشت دا قضیہ ثابت کرن دیندا اے (ایہو «بنیاد والے تعلقاں اُتے استقرائی تعریفاں» نیں)۔

ایہ نتیجے واقعی من پسند نیں۔ پر کی ایہ من پسند نتیجے بدل رکھن دے اصول دا \emph{جواز} دین لئی کافی نیں؟ \emph{نہیں}۔ گھٹ توں گھٹ سیدھے ای کافی نہیں۔

اک سادہ اوکڑ ویکھو۔ اسیں بدل رکھن دے اصول دے کجھ من پسند نتیجے دسے نیں، پر ایہناں وچوں کئی ہور طریقیاں نال وی لبھ سکدے ساں۔ ایہ گل جِنّی لوکاں نوں پتا ہونی چاہیدی سی، اِنّی پتا نہیں۔ سو سانوں ایتھے رُک کے حال سمجھانا چاہیدا اے۔

سیٹاں دا اک سادہ نظریہ اے، سطحاں دا نظریہ، جیہنوں مختصر کر کے $\LT$ آکھدے آں۔\footnote{$\LT$ دے پہلے روپ \citet{Montague1965} تے \citet{Scott1974} نے پیش کیتے۔ \citet{Potter2004} نے ایہنوں سوکھا کیتا تے پوری کتاب وچ اوہدی تفصیل دتی۔ \citet{ButtonLT1} نے حال ای وچ $\LT$ نوں ہور سوکھا کیتا اے۔} $\LT$ دے مسلّمے بس عنصراں نال تعیّن، وکھ کرن دا اصول، تے ایہ دعویٰ نیں کہ ہر سیٹ کسے \emph{سطح} دا ذیلی سیٹ اے۔ «سطح» دی تعریف ایہو جیہی سوچ سمجھ کے دتی گئی اے کہ سطحاں ساڈے جانے پچھانے $V_\alpha$ سیٹاں وانگ کم کردیاں نیں۔ سو $\ZF$ توں $\LT$ ثابت ہُندا اے، پر $\LT$، $\ZF$ توں \emph{کِتے} کمزور اے۔ اصل وچ $\LT$ توں جوڑیاں، ذیلی سیٹاں دے سیٹ، لامتناہیت، یا بدل رکھن دا اصول نہیں ملدا۔ $\Zr$ اوہ نظریہ ہووے جیہڑا $\LT$ وچ لامتناہیت تے ذیلی سیٹاں دے سیٹ دے مسلّمے پا کے ملدا اے۔ ایس توں جوڑیاں وی مل جاندیاں نیں، سو $\Zr$ گھٹ توں گھٹ $\Z$ جِنّا طاقتور اے۔ پر اصل وچ $\Zr$، $\Z$ توں ٹھیک ٹھیک ودھ طاقتور اے، کیوں جے ایہ دعویٰ وی شامل کردا اے کہ ہر سیٹ دا درجہ اے (ایسے لئی میرا سجھا اے کہ اینوں $\Zr$ آکھیے)۔ سچ مُچ $\Zr$ توں ترتیبی عدداں دا پورا تسلّی بخش نظریہ، درجہ وار سلسلے نوں حُسنِ ترتیب والے مرحلیاں وچ ونڈن والے نتیجے، $\in$-استقرا دا ثبوت، تے لامتناہی درجیاں والی بازگشت دا اک \emph{روپ} ملدے نیں۔

مختصر گل ایہ اے کہ بھاویں بُولوس نوں ایہ پتا نہیں سی، اوہدے دسے سارے من پسند نتیجے بدل رکھن دے اصول \emph{توں بناں} لبھے جا سکدے سن۔ اوہنوں بس $\Zr$ ورتن دی لوڑ سی، $\Z$ دی تھاں۔

(ایہ سارا کجھ ہون اُتے وی اسیں $\ZF$ سکھاؤن دا رواجی راہ کیوں چُنیا، $\LT$ تے $\Zr$ کیوں نہیں سکھائے؟ دو وجوہاں نیں۔ پہلی، نِریاں تاریخی وجوہاں کرکے $\LT$ توں شروع کرنا عام طریقہ نہیں۔ اسیں تہانوں سیٹ نظریے دیاں ہور معیاری پیشکشاں پڑھن جوگا کرنا چاہندے ساں۔ دوجی، جدوں تسی $\LT$ تے $\Zr$ دی قدر سمجھ سکدے او، تسی سیدھے ای \citealt{Potter2004} تے \citealt{ButtonLT1} پڑھ سکدے او۔)

بے شک، کیوں جے $\Zr$، $\ZF$ توں ٹھیک ٹھیک کمزور اے، کجھ نتیجے $\ZF$ توں ثابت ہُندے نیں جیہڑے $\Zr$ وچ کھلے رہ جاندے نیں۔ سو کوئی ایہناں وچوں اک نتیجے ول اشارہ کر کے بدل رکھن دے اصول دا باہرلا جواز دین دی کوشش کر سکدا اے۔ پر جدوں تہانوں $\Zr$ ورتنا آ جاندا اے، ایہو جیہیاں بہتیاں مثالاں لبھنیاں اوکھیاں نیں جیہڑیاں (ا) بدل رکھن دے اصول نال طے ہو جاندیاں نیں تے ہور طریقے نال نہیں، تے (ب) سہج سمجھ مطابق سچ نیں۔ (ہور تفصیل لئی \citealt[\S13.2]{Potter2004} ویکھو۔)

اصل گل ایہ اے۔ بدل رکھن دے اصول دا مناؤن جوگا باہرلا جواز دین لئی کوئی ایہو جیہا نتیجہ لبھنا پئے گا جیہڑا اوہدے بناں حاصل \emph{نہیں ہو سکدا}۔ تے ایہ سوکھا کم نہیں۔

ہُن اک ہور اوکڑ ویکھو جیہڑی بدل رکھن دے اصول دا نِرا باہرلا جواز دین دی ہر کوشش وچ اُٹھدی اے۔ (ایہ اوکڑ شاید پہلی توں وی بنیادی اے۔) بُولوس صرف ایہ نہیں دسدا کہ بدل رکھن دے اصول توں کئی من پسند نتیجے نکلدے نیں۔ اوہ ایہ وی آکھدا اے کہ اوہدے «(جیویں ہن تک دِسدا اے) کوئی ناپسند» نتیجے نہیں۔ پر بریکٹاں والی ایہ احتیاط، «جیویں ہن تک دِسدا اے»، پکی گل اے کہ بڑی ضروری اے۔

چیتے کرو کہ اسیں ایتھے کیویں پہنچے۔ سادہ خاصیت توں سیٹ بناؤن دے اصول نال ٹکراؤ نکل آیا سی، تے اسیں ایس ٹکراؤ دے جواب وچ سیٹ دا اکٹھا کرن والا مرحلہ وار تصور اپنایا۔ ایس تصور نال اک کہانی جڑی اے جیہڑی، سانوں آس اے، اوہدے بے ٹکراؤ ہون دا بھروسا دیندی اے۔ پر جے اسیں اوہو کہانی دے اندر رہ کے بدل رکھن دے اصول دا جواز نہیں دے سکدے، تاں ہن تک ساڈے کول $\ZF$ نوں بے ٹکراؤ منّن دی کوئی وجہ نہیں۔ ہور ٹھیک آکھیے تاں ساڈے کول $\ZF$ نوں بے ٹکراؤ منّن دی وجہ بس ایہ (شاید اتفاقی) گل اے کہ کسے نوں \emph{ہن تک} کوئی ٹکراؤ نہیں لبھیا۔ این ٹھیک ایس معنی وچ بُولوس دی گل ایہو رہ جاندی لگدی اے کہ «(جیویں ہن تک دِسدا اے) $\ZF$ بے ٹکراؤ اے»۔ سانوں بے ٹکراؤ ہون دا ایہتوں ودھ بھروسا منگنا چاہیدا اے۔

ایہ اوکڑ بدل رکھن دے اصول دا \emph{نِرا} باہرلا جواز دین دی ہر کوشش وچ آوے گی، یعنی ہر اوہ جواز جیہڑا بس $\ZF$ دے (جانے پچھانے) نتیجیاں اُتے کھڑا اے۔ ایس لئی اسیں بدل رکھن دے اصول دا \emph{اندرلا} جواز لبھنا چاہواں گے، یعنی ایہ دکھاؤن والا جواز کہ سیٹاں بارے ساڈی دسی کہانی کسے طرح «پہلاں ای» سانوں بدل رکھن دا اصول منّن لئی پابند کردی اے۔
\endgroup
% /OLP-0568
% OLP-0569 content/set-theory/replacement/limofsize.tex
\begingroup
\def\olfilename{limofsize}
\expandafter\def\csname ol@sectioncs\endcsname{section}

\olfileid{sth}{replacement}{limofsize}
\olsection{وڈائی دی حد}

شاید بدل رکھن دے اصول دا «اندرلا» جواز دین دی سب توں عام کوشش ایس خیال توں شروع ہُندی اے۔
\begin{enumerate}
	\item[] \limofsize. کوئی وی شےواں رَل کے اک سیٹ بناندیاں نیں، بشرطیکہ اوہ حد توں ودھ نہ ہون۔
\end{enumerate}
ایہ اصول بدل رکھن دے اصول دی جھٹ تائید کرے گا۔ آخر بدل رکھن نال بنیا کوئی سیٹ، اوہ سیٹ جیہدے توں ایہ بنیا اے، اوہتوں وڈا نہیں ہو سکدا۔ ہور ٹھیک آکھیے تاں منّ لو تسی بدل رکھن دا اصول ورت کے سیٹ
$\funimage{\tau}{A} = \Setabs{\tau(x)}{x \in A}$ بناندے او۔
تاں $\cardle{\funimage{\tau}{A}}{A}$۔ سو جے $A$ دے !!{element}s دی گنتی اک سیٹ بنن لئی حد توں ودھ نہیں سی، تاں اوہناں دے حاصل وی $\funimage{\tau}{A}$ بناؤن لئی حد توں ودھ نہیں۔

بدل رکھن دے اصول دا جواز دین لئی \limofsize{} ورتن وچ صاف اوکڑ ایہ اے کہ اسیں ہن تک \limofsize ورگا کوئی اصول دتا ای \emph{نہیں}۔ نالے جدوں اسیں \crefrange{sth:story::chap}{sth:z::chap} وچ سیٹ دے اکٹھا کرن والے مرحلہ وار تصور دی کہانی دسی سی، اوہدے وچ \limofsize ول کوئی \emph{اشارہ} وی نہیں سی۔ این ایسے لئی بُولوس نے اک تھاں لکھیا سی «شاید ایہ نتیجہ کڈھیا جا سکدا اے کہ سیٹ نظریے دے ‘پچھے’ گھٹ توں گھٹ دو خیال نیں» \citeyearpar[p.~19]{Boolos1989}۔ اک پاسے سیٹ دے اکٹھا کرن والے مرحلہ وار تصور نال جڑے خیالاں دا کم~$\Z$ دا جواز دینا اے۔ دوجے پاسے \limofsize{} دا کم بدل رکھن دے اصول دا جواز دینا اے۔

گل صرف ایہ نہیں کہ اسیں ہن تک \limofsize بارے \emph{چُپ} رہے آں۔ گل ایہ اے کہ \limofsize{} (جیویں ہن بیان کیتا اے) سیٹ دے اکٹھا کرن والے مرحلہ وار تصور نال چنگی طرح رَلدا نہیں لگدا۔ آخر ایہ سیٹاں دے \emph{مرحلیاں} وچ بنن دے خیال دا ذکر ای نہیں کردا۔

ایہناں «دو خیالاں» (یعنی سیٹ دے اکٹھا کرن والے مرحلہ وار تصور تے \limofsize) دا پچھوکڑ ویکھ کے ایہ گل کوئی بہت حیرانی والی نہیں۔ آخرکار ایہ «دو خیال» سیٹ نظریے دے پیراڈوکساں نوں روکن دے دو خاصے وکھرے جتن نیں۔ سیٹ دا اکٹھا کرن والا مرحلہ وار تصور رسل دے پیراڈوکس نوں ایہ آکھ کے روکدا اے۔ موٹی گل ایہ اے کہ \emph{سانوں رسل سیٹ دی موجودگی دی آس رکھنی ای نہیں چاہیدی سی، کیوں جے ایہ کسے مرحلے اُتے «بننا» ای نہیں سی}۔ دوجے پاسے \limofsize{} رسل سیٹ نوں موٹی گل ایہ آکھ کے ردّ کردا اے کہ \emph{سانوں رسل سیٹ دی موجودگی دی آس رکھنی ای نہیں چاہیدی سی، کیوں جے اوہ حد توں ودھ وڈا ہونا سی}۔

سو گل ایس طرح رکھیے تاں صاف آکھنا پئے گا کہ پیراڈوکساں دے جواب وجوں \limofsize{} نوں \emph{بہت} ہور جواز دی لوڑ اے۔ مثلاً رسل دے پیراڈوکس دا ایہ روپ ویکھو۔ \emph{کوئی پگ کُتّا این ٹھیک اوہناں پگ کُتیاں نوں نہیں سُنگھدا جیہڑے اپنے آپ نوں نہیں سُنگھدے} (\olref[sth][story][rus]{sec} ویکھو)۔ جے تسی پُچھو «ایہو جیہا کوئی پگ کُتّا کیوں نہیں؟»، تاں ایہ چنگا جواب نہیں کہ اوہو جیہے کُتے نوں حد توں ودھ پگ کُتے سُنگھنے پینے سن۔ فیر ایہ آکھنا کہ رسل سیٹ موجود ہون لئی «حد توں ودھ وڈا» ہووے گا، اوہدی غیرموجودگی دی سہج سمجھ والی چنگی وضاحت کیویں ہو سکدی اے؟

مختصر گل ایہ اے کہ جے \limofsize دی «سہج سمجھ» والی وجہ تہانوں تھوڑی اوپری لگے تاں ایہ سمجھ آؤن والی گل اے۔
\endgroup
% /OLP-0569
% OLP-0570 content/set-theory/replacement/absinf.tex
\begingroup
\def\olfilename{absinf}
\expandafter\def\csname ol@sectioncs\endcsname{section}

\olfileid{sth}{replacement}{absinf}
\olsection{بدل رکھن دا اصول تے «مطلق لامتناہیت»}

ہُن اسیں \limofsize{} دی گل پچھے چھڈ کے بدل رکھن دے اصول دے (اندرلے) جواز دیاں کجھ ہور کوششاں ویکھاں گے، جیہڑیاں سیٹاں دے مرحلیاں وچ بنن دے خیال نوں سنجیدگی نال لیندیاں نیں۔

جدوں اسیں پہلی واری مرحلہ وار طریقے دا خاکہ دتا سی، اسیں کجھ اصول دتے سن جیہڑے سمجھاندے سن کہ ہر مرحلے اُتے کی ہُندا اے۔ ایہ \stageshier، \stagesord، تے \stagesacc سن۔ بعد وچ اسیں کجھ ہور اصول پائے جیہڑے مرحلیاں دی گنتی بارے دسدے سن۔ \stagessucc{} دسدا سی کہ سیٹ بنن دا کم کدی مکدا نہیں، تے \stagesinf{} دسدا سی کہ ایہ کم لامتناہی نمبر والے اک مرحلے توں وی لنگھدا اے۔

ایہ سوچنا واجب اے کہ ایہ پچھلے دو اصول کسے ہور کھلّھے اصول توں نکلدے نیں، جیویں کہ
\begin{enumerate}
	\item[] \stagesinex. مرحلے مطلق طور اُتے لامتناہی نیں۔ درجہ وار سلسلہ جِنّا اُچا ہو سکدا اے، اِنّا اُچا اے۔
\end{enumerate}
صاف اے کہ ایہ اک غیر رسمی اصول اے۔ پر بھاویں ایہ سیٹ دے اکٹھا کرن والے مرحلہ وار تصور توں جھٹ \emph{لازم نہیں آؤندا}، اوہدے نال \emph{میل} ضرور کھاندا لگدا اے۔ گھٹ توں گھٹ، \limofsize دے اُلٹ، ایہ خیال قائم رکھدا اے کہ سیٹ مرحلہ وار بندے نیں۔

ہُن آس ایہ اے کہ \stagesinex{} نوں ورت کے اسیں بدل رکھن دے اصول دا جواز دے سکدے آں۔ آؤ ویکھیے کہ ایہ کم کیویں ہو سکدا اے۔

\olref[ordinals][idea]{sec} وچ اسیں ویکھیا کہ اک ایہو جیہی حُسنِ ترتیب بنانی سوکھی اے جیہڑی (سہج سمجھ مطابق) $\omega+\omega$ نال ہم ساخت ہونی چاہیدی اے۔ دوجے لفظاں وچ، اسیں سوکھ نال اک مرحلہ وار طریقہ سوچ سکدے آں جیہدی ترتیب دی قسم (سہج سمجھ مطابق) $\omega+\omega$ اے۔ سو جے اسیں \stagesinex منّ لیا اے، تاں سانوں ضرور منّنا چاہیدا اے کہ درجہ وار سلسلے دا گھٹ توں گھٹ $\omega+\omega$ واں مرحلہ، یعنی $V_{\omega+\omega}$، اے۔ کیوں جے سلسلہ ایتھوں تک ضرور چل \emph{سکدا} اے۔

ایس خیال نوں عام کر کے ایویں آکھ سکدے آں۔ جیہڑی وی حُسنِ ترتیب ہووے، مرحلہ وار درجہ بندی بنن دا کم گھٹ توں گھٹ اوہدی حد تک چلنا چاہیدا اے۔ تے ایہدی ضمانت اسیں بس \olref[ordinals][ordtype]{thmOrdinalRepresentation} نوں اک \emph{مسلّمہ} منّ کے دے سکدے آں۔ ایہ سانوں دسے گا کہ ہر حُسنِ ترتیب کسے فون نوئمان ترتیبی عدد نال ہم ساخت اے۔ کیوں جے ہر فون نوئمان ترتیبی عدد اپنے درجے دے برابر ہووے گا، \olref[ordinals][ordtype]{thmOrdinalRepresentation} سانوں دسے گا کہ جدوں وی اسیں اپنے سیٹ نظریے وچ کسے حُسنِ ترتیب دا بیان دے سکیے، سیٹ بنن دا مرحلہ وار کم اوہ حُسنِ ترتیب توں اگے ودھے گا۔

ایہ خیال ضرور \stagesinex توں نکلدا اک نتیجہ لگدا اے۔ افسوس، جے ساڈا مقصد ایس خیال توں بدل رکھن دا اصول کڈھنا اے، تاں اک سادہ فنی اڑچن اے۔ بدل رکھن دا اصول \olref[ordinals][ordtype]{thmOrdinalRepresentation} توں ٹھیک ٹھیک ودھ طاقتور اے۔ (ایہ گل \citet[\S13.2]{Potter2004} نے دسی اے۔ اسیں ایہنوں \olref[sth][replacement][finiteaxiomatizability]{sec} وچ ثابت کراں گے۔)

نتیجہ ایہ اے کہ جے اسیں \stagesinex{} نوں ایویں سمجھنا اے کہ اوہتوں بدل رکھن دا اصول نکلے، تاں ایہ \emph{بس} ایہ نہیں آکھ سکدا کہ درجہ وار سلسلہ ہر حُسنِ ترتیب توں اگے ودھدا اے۔ ایہنوں ایہتوں وڈا دعویٰ کرنا پئے گا۔ ایس کم لئی \cite{Shoenfield:AST} نے ایس خیال دا اک بڑا سہج، ہور طاقتور روپ پیش کیتا۔ درجہ وار سلسلہ کسے وی سیٹ نال \emph{ہم حد} نہیں۔\footnote{گوڈل نے وی ایہو جیہا خیال دتا لگدا اے۔ \citet[p.~223]{Potter2004} ویکھو۔ گوڈل تے \citeauthor{Shoenfield:AST} بارے بحث لئی \citet[90--5]{Incurvati2020} ویکھو۔} ذرا ہور تفصیل نال، جے $\tau$ سیٹاں نوں درجہ وار سلسلے دے مرحلیاں ول بھیجن والا نقشہ اے، تاں کسے وی سیٹ $A$ دا $\tau$ تھلّے عکس سارے سلسلے نوں نہیں گھیر سکدا۔ خاکے دی صورت وچ ایویں آکھو۔
\begin{enumerate}
	\item[] \stagescofin. جے $A$ اک سیٹ اے تے $\tau(x)$ ہر $x \in A$ لئی اک مرحلہ اے، تاں اک مرحلہ اے جیہڑا ہر $\tau(x)$ توں پچھوں آؤندا اے، سارے $x \in A$ لئی۔
\end{enumerate}
صاف اے کہ $\ZF$ توں \stagescofin دا ڈھکواں رسمی روپ ثابت ہُندا اے۔ اُلٹے پاسے اسیں غیر رسمی طور اُتے دلیل دے سکدے آں کہ \stagescofin{} بدل رکھن دے اصول دا جواز دیندا اے۔\footnote{\stagescofin دے کسے رسمی روپ توں بدل رکھن دا اصول ثابت کرنا ہور اوکھا ہووے گا، کیوں جے اکلّا $\Z$ مرحلیاں دی تعریف کرن جوگا طاقتور نہیں۔ ایس لئی صاف نہیں کہ \stagescofin نوں رسمی صورت کیویں دیئیے۔ اک طریقہ ایہ اے کہ $\LT$ دی کسے توسیع وچ کم کریے، جیویں \olref[sth][replacement][extrinsic]{sec} وچ دَسیا اے۔} منّ لو $(\forall x \in A)\lexists![y][\phi(x,y)]$۔ فیر ہر $x \in A$ لئی، $\sigma(x)$ اوہ $y$ ہووے جیہدے لئی $\phi(x,y)$، تے $\tau(x)$ اوہ مرحلہ ہووے جتھے $\sigma(x)$ پہلی واری بندا اے۔ \stagescofin مطابق اک مرحلہ $V$ اے جیہدے لئی $(\forall x \in A)\tau(x)\in V$۔ ہُن کیوں جے ہر $\tau(x) \in V$ تے $\sigma(x) \subseteq \tau(x)$، اسیں وکھ کرن دے اصول نال $ \Setabs{y
\in V}{(\exists x \in A)\sigma(x) = y} = \Setabs{y}{(\exists x \in
A)\phi(x,y)}$ حاصل کر سکدے آں۔

\begin{prob}
	\stagescofin{} نوں $\ZF$ وچ رسمی صورت دیو۔
\end{prob}

سو \stagescofin{} بدل رکھن دے اصول دا جواز دیندا اے۔ تے ایہ گھٹ توں گھٹ منّن جوگی گل اے کہ \stagesinex{} توں \stagescofin دا جواز ملدا اے۔ کیوں جے منّ لو \stagescofin{} نئیں چلدا۔ فیر سلسلہ کسے سیٹ~$A$ نال ہم حد اے، یعنی اک نقشہ $\tau$ اے جیہدے لئی ہر مرحلے~$S$ واسطے کوئی $x \in A$ اے جس لئی $S \in \tau(x)$۔ ایس صورت وچ سلسلے دی منّی ہوئی «مطلق لامتناہیت» نوں فڑن دا طریقہ ساڈے کول اے۔ $\tau$ نوں $A$ اُتے لاؤن نال جیہڑی اقدار دا سیٹ ملدی اے، اوہ سارے سلسلے نوں \emph{گھیر لیندا اے}۔ تے ایہ گل سلسلے دے «مطلق طور اُتے لامتناہی» ہون دے خیال نوں کمزور کردی اے۔ عکسِ نقیض لئیے تاں \stagesinex{} توں \stagescofin لازم آؤندا اے، جیہڑا اگوں بدل رکھن دے اصول دا جواز دیندا اے۔

ایہ بدل رکھن دے اصول دا اندرلا جواز دین دی واقعی آس والی کوشش اے۔ پر آخر ایہ کوشش کامیاب ہُندی اے یا نہیں، ایہ فیصلہ اسیں تہاڈے اُتے چھڈاں گے۔
\endgroup
% /OLP-0570
% OLP-0571 content/set-theory/replacement/ref.tex
\begingroup
\def\olfilename{ref}
\expandafter\def\csname ol@sectioncs\endcsname{section}

\olfileid{sth}{replacement}{ref}
\olsection{بدل رکھن دا اصول تے انعکاس}

بدل رکھن دے اصول دا \stagesinex راہیں جواز دین دی ساڈی آخری کوشش اک ڈونگھے تے سوہنے نتیجے توں شروع ہُندی اے۔\footnote{چیتے رہے کہ سارے فارمولیاں وچ پیرامیٹر ہو سکدے نیں (جے کھولھ کے ہور گل نہ آکھی ہووے)۔}
%In this, I will use overlining, such as $x_1, \ldots, x_n$, to
%abbreviate ``$x_1, \ldots, x_n$'':
\begin{thm}[انعکاس دا خاکہ]\ollabel{reflectionschema}
کسے وی فارمولے $\phi$ لئی:
\[
\forall \alpha \exists \beta > \alpha (\forall x_1 \ldots, x_n \in
V_\beta)(\phi(x_1, \ldots, x_n) \liff \phi^{V_\beta}(x_1, \ldots, x_n))
\]
\end{thm}
\noindent
جیویں \olref[sth][replacement][strength]{formularelativization} وچ، $\phi^{V_\beta}$ اوہ فارمولا اے جیہڑا $\phi$ دے ہر مسوّر نوں سیٹ~$V_\beta$ تک محدود کر کے ملدا اے۔ سو سہج سمجھ مطابق انعکاس ایہ آکھدا اے۔ جے $\phi$ سارے درجہ وار سلسلے وچ سچ اے، تاں $\phi$ سلسلے دے من چاہی اونچائی اُتے ملدے \emph{مُڈھلے ٹکڑیاں} وچ وی سچ اے۔

\citet{Montague1961} تے \citet{Levy1960} نے وکھایا کہ (ڈھکویں روپ وچ) بدل رکھن دا اصول تے انعکاس~$\Z$ دے اندر برابر نیں۔ ایس لئی دوہاں وچوں کوئی وی پا دیئیے تاں~$\ZF$ ملدا اے۔ (اسیں ایہ نتیجے \olref[sth][replacement][refproofs]{sec} وچ ثابت کردے آں۔) ایس برابری نال کوئی انعکاس تے بدل رکھن دے اصول دا جواز \stagesinex{} راہیں ایویں دین دی آس رکھ سکدا اے۔ \stagesinex منّ کے درجہ وار سلسلہ بڑا ای اُچا ہووے گا۔ اینا اُچا کہ اوہدے بارے اسیں جو کجھ آکھ سکیے، اوہ اوہدی اُچائی دی حد بنّن لئی کافی نہیں۔ تے اسیں ایہ گل ایویں سمجھ سکدے آں کہ جے کوئی وی جملہ~$\phi$ سارے سلسلے وچ سچ اے، تاں اوہ سلسلے دے من چاہی اونچائی اُتے ملدے مُڈھلے ٹکڑیاں وچ سچ اے۔ تے ایہو انعکاس اے۔

ایہ وی بدل رکھن دے اصول دا اندرلا جواز دین دی واقعی آس والی کوشش لگدی اے۔ پر ایس بارے کہن نوں ایتھے گنجائش توں کِتے ودھ گلّاں نیں۔ ہُن تسی آپ فیصلہ کرو کہ ایہ کامیاب اے یا نہیں۔\footnote{اگے تسی \citet[95--100]{Incurvati2020} پڑھنا چاہو گے۔}
\endgroup
% /OLP-0571
% OLP-0572 content/set-theory/replacement/refproofs.tex
\begingroup
\def\olfilename{refproofs}
\expandafter\def\csname ol@sectioncs\endcsname{section}

\olfileid{sth}{replacement}{refproofs}
\olsection{ضمیمہ: بدل رکھن دے اصول نال جڑے نتیجے}

ایس حصے وچ اسیں $\ZF$ دے اندر انعکاس ثابت کراں گے۔ اسیں اوہ معنی وی ثابت کراں گے جیہدے وچ انعکاس بدل رکھن دے اصول دے برابر اے۔ تے ایس سارے کم دا اک دلچسپ نتیجہ ثابت کراں گے، جیہڑا انعکاس تے بدل رکھن دے اصول دی طاقت بارے اے۔ \emph{خبردار: ایہ ایس کتاب وچ ریاضی دا سبھ توں اوکھا حصہ اے۔}

اسیں اک معاون قضیے توں شروع کراں گے، جیہدے وچ متغیراں یا شےواں دیاں لڑیواں نوں مختصر لکھن لئی \emph{اُتلی لکیر} ورتی اے۔ سو «$\overline{a}_k$» دا مختصر مطلب «$a_{k_1}$، \dots،
$a_{k_n}$» اے، جتھے $n$ دا پتا سیاق توں لگدا اے۔

\begin{lem}\ollabel{lemreflection}
ہر $1 \leq i \leq k$ لئی $\phi_i(\overline{v}_i, x)$ اک فارمولا ہووے۔ فیر ہر $\alpha$ لئی کوئی $\beta > \alpha$ اے جیہدے لئی، کسے وی $\overline{a}_1,
\ldots, \overline{a}_k \in V_\beta$ تے ہر $1 \leq i \leq k$ لئی:
\[
	\exists x\phi_i(\overline{a}_i, x) \rightarrow (\exists x \in V_\beta) \phi_i(\overline{a}_i, x)
\]
\end{lem}

\begin{proof}
اسیں اک ٹرم $\mu$ ایویں تعریف کردے آں۔ $\mu(\overline{a}_1, \ldots,
\overline{a}_k)$ اوہ سبھ توں پہلا مرحلہ، $V$، اے جیہدے لئی $1 \leq i \leq k$ واسطے ہِیٹھلے سارے شرطیے سچ نیں۔
\[
\exists x\phi_i(\overline{a}_i, x) \rightarrow (\exists x \in V) \phi_i(\overline{a}_i, x))
\]
ایہ تصدیق کرنی سوکھی اے کہ $\mu(\overline{a}_1, \ldots, \overline{a}_k)$ سارے $\overline{a}_1, \ldots, \overline{a}_k$ لئی موجود اے۔ ہُن بدل رکھن دا اصول تے بازگشت دا ساڈا قضیہ ورت کے ایویں تعریف کرو۔
\begin{align*}
	S_0 & = V_{\alpha+1}\\
	S_{n+1} & = S_n \cup \bigcup
	\Setabs{\mu(\overline{a}_1, \ldots, \overline{a}_k)}
	{\overline{a}_1, \ldots, \overline{a}_k \in S_n} \\
	S &= \bigcup_{m < \omega} S_n.
\end{align*}
ہر $S_n$، تے ایس لئی $S$ آپ وی، $V_\alpha$ توں بعد دا اک مرحلہ اے۔ ہُن
$\overline{a}_1$، \dots،~$\overline{a}_k \in S$ مقرر کرو۔ سو کوئی $n <
\omega$ اے جیہدے لئی $\overline{a}_1$، \dots، $\overline{a}_k \in S_n$۔
کوئی $1 \leq i \leq k$ مقرر کرو تے منّ لو $\exists x
\phi_i(\overline{a}_i,x)$۔ سو بناوٹ مطابق $(\exists x \in \mu(\overline{a}_1,
\ldots, \overline{a}_k))\phi_i(\overline{a}_i, x)$، ایس لئی
$(\exists x \in S_{n+1})\phi_i(\overline{a}_i, x)$ تے اوہتوں
$(\exists x \in S)\phi_i(\overline{a}_i, x)$۔ سو $S$ ساڈا $V_\beta$ اے۔
\end{proof}
\noindent
ہُن اسیں \olref[sth][replacement][ref]{reflectionschema} بڑے سیدھے طریقے نال ثابت کر سکدے آں۔

\begin{proof}[ثبوت]
$\alpha$ مقرر کرو۔ عمومیت گنوائے بناں اسیں منّ سکدے آں کہ $\phi$ دے جوڑک بس $\exists$، $\lnot$ تے $\land$ نیں (کیوں جے ایہ اظہار لئی کافی نیں)۔ $\psi_1, \ldots, \psi_k$، $\phi$ دے سارے ذیلی فارمولاں نوں پیچیدگی دی ترتیب وچ گِنن، ایویں کہ $\psi_k = \phi$۔
\olref{lemreflection} مطابق اک $\beta > \alpha$ اے جیہدے لئی کسے وی $\overline{a}_i \in V_\beta$ تے ہر $1 \leq i \leq k$ واسطے:
\begin{align}\label{reflectionnicelybehaved}
	\exists x\psi_i(\overline{a}_i, x) \rightarrow
	(\exists x \in V_\beta) \psi_i(\overline{a}_i, x)\tag{*}
\end{align}
$\psi_i$ دی پیچیدگی اُتے استقرا نال اسیں وکھاواں گے کہ
$\psi_i(\overline{a}_i) \leftrightarrow
\psi_i^{V_\beta}(\overline{a}_i)$، ہر $\overline{a}_i \in
V_\beta$ لئی۔ جے $\psi_i$ جوہری اے، تاں ایہ گل صاف اے۔ دوپاسے دا شرطیہ ایہ وی طے کردا اے کہ جدوں $\psi_i$ ایس خاصیت والے ذیلی فارمولاں دی نفی یا عطف اے، $\psi_i$ وچ وی ایہ خاصیت اے۔ سو بس مسوّرے والی صورت ویکھن جوگی اے۔
$\overline{a}_i \in V_\beta$ مقرر کرو۔ فیر:
\begin{align*}
	(\exists x \psi_i(\overline{a}_i, x))^{V_\beta}
	&\text{ اُدوں تے صرف اُدوں جدوں }
	(\exists x \in V_\beta) \psi_i^{V_\beta}(\overline{a}_i, x)
	&&\text{تعریف مطابق}\\
	&\text{ اُدوں تے صرف اُدوں جدوں }
	(\exists x \in V_\beta) \psi_i(\overline{a}_i,  x)
	&&\text{مفروضے مطابق}\\
	&\text{ اُدوں تے صرف اُدوں جدوں }
	\exists x \psi_i(\overline{a}_i, x)
	&&\text{\eqref{reflectionnicelybehaved} مطابق}
\end{align*}
ایس نال استقرا پورا ہویا۔ نتیجہ نکل آیا، کیوں جے $\psi_k = \phi$۔
\end{proof}

اسیں $\ZF$ وچ انعکاس ثابت کر دتا اے۔ ساڈا ثبوت بنیادی طور اُتے \citet{Montague1961} دے ثبوت دے پچھے چلدا اے۔ ہُن اسیں $\Z$ وچ ثابت کرنا چاہندے آں کہ انعکاس توں بدل رکھن دا اصول لازم آؤندا اے۔ ثبوت \citet{Levy1960} دے طریقے مطابق اے، پر اک سادہ کیتی صورت وچ۔

کیوں جے اسیں $\Z$ وچ کم کردے آں، اسیں انعکاس نوں این اوہو روپ وچ پیش نہیں کر سکدے جیہڑا اُتے دتا اے۔ آخر اسیں انعکاس «$V_\alpha$» والی لکھت نال بیان کیتا سی، تے $\Z$ وچ ایہدی تعریف نہیں ہو سکدی (\olref[sth][spine][zf]{sec} ویکھو)۔ سو اسیں ایس دی تھاں بدل رکھن دے اصول دا بظاہر کمزور روپ ایویں پیش کراں گے۔

\begin{defish}
\emph{کمزور انعکاس۔} کسے وی فارمولے $\phi$ لئی اک متعدی سیٹ $S$ اے، ایویں کہ $0$، $1$ تے $\phi$ دے سارے پیرامیٹر $S$ دے !!{element}s نیں، تے $(\forall \overline{x} \in S)(\phi \liff
\phi^S)$۔
\end{defish}

ایہ ورت کے بدل رکھن دا اصول ثابت کرن لئی پہلاں اسیں \citet[first
part of Theorem 2]{Levy1960} دے طریقے تے چل کے وکھاواں گے کہ اسیں دو فارمولیاں دا اکّو ویلے «انعکاس» کر سکدے آں۔

\begin{lem}[$\Z + \text{کمزور انعکاس}$ وچ۔]\ollabel{lem:reflect}
کسے وی دو فارمولیاں $\psi, \chi$ لئی اک متعدی سیٹ $S$ اے، ایویں کہ $0$ تے $1$ (تے فارمولیاں دے سارے پیرامیٹر)
$S$ دے !!{element}s نیں، تے $(\forall \overline{x} \in S)((\psi \liff \psi^S) \land (\chi
\liff \chi^S))$۔
\end{lem}

\begin{proof}
$\phi$ فارمولا $(z = 0 \land \psi) \lor (z = 1 \land \chi)$ ہووے۔

ایتھے اسیں مختصر لکھت ورتدے آں۔ «$z = 0$» نوں کھولھ کے «$\forall t\, t \notin z$» لکھنا چاہیدا اے، تے «$z =1$» نوں «$\forall s(s \in z
\liff \forall t\, t \notin s)$»۔ پر کیوں جے $0, 1 \in S$ تے $S$ متعدی اے، ایہ فارمولا $S$ لئی \emph{مطلق} نیں۔ یعنی اوہ اوسے شے اُتے لاگو ہون گے، بھاویں اسیں اوہناں دے مسوّر $S$ تک محدود کر دیئیے۔\footnote{ہور رسمی طور اُتے، $\xi$ ایہناں دوہاں وچوں کوئی وی فارمولا ہووے، تاں $\xi(z) \liff \xi^S(z)$۔}

کمزور انعکاس نال کوئی ڈھکواں $S$ ملدا اے جیہدے لئی:
\begin{align*}
	(\forall z, \overline{x} \in S)(&\phi \liff \phi^S)\\
	\text{یعنی }(\forall z, \overline{x} \in S)(&((z = 0 \land \psi) \lor (z = 1 \land \chi)) \liff {}\\
	&\phantom{(}((z = 0 \land \psi) \lor (z = 1 \land \chi))^S)\\
	\text{یعنی }(\forall z, \overline{x} \in S)(&((z = 0 \land \psi) \lor (z = 1 \land \chi))\liff {}\\
	&\phantom{(}((z = 0 \land \psi^S) \lor (z = 1 \land \chi^S)))\\
	\text{یعنی }(\forall \overline{x} \in S)(&(\psi \liff \psi^S) \land (\chi \liff \chi^S))
\end{align*}
دوجے دعوے توں تیجا لازم آؤندا اے کیوں جے «$z = 0$» تے «$z=1$» $S$ لئی مطلق نیں۔ چوتھا دعویٰ نکلدا اے کیوں جے $0 \neq 1$۔
\end{proof}\noindent ہُن اسیں \citet[Theorem 6]{Levy1960} دے طریقے تے چل کے، اوہنوں سادہ کردیاں، بدل رکھن دا اصول حاصل کر سکدے آں۔

\begin{thm}[$\Z$ + کمزور انعکاس وچ]\label{thm:replacement}
کسے وی فارمولے $\phi(v,w)$ تے کسے وی $A$ لئی، جے $(\forall x \in A)\lexists![y][\phi(x,y)]$، تاں
$\Setabs{y}{(\exists x \in A)\phi(x,y}$ موجود اے۔
\end{thm}

\begin{proof}
$A$ ایہو جیہا مقرر کرو کہ $(\forall x \in A)\lexists![y][\phi(x,y)]$، تے فارمولے ایویں تعریف کرو۔
\begin{align*}
	\psi &\text{ ایہ اے } (\phi(x, z) \land A = A)\\
	\chi &\text{ ایہ اے } \lexists[y][\phi(x, y)]
\end{align*}
\olref{lem:reflect} ورت کے، کیوں جے $A$، $\psi$ دا اک پیرامیٹر اے، اک متعدی~$S$ اے جیہدے لئی $0, 1, A \in S$ (ہور سارے پیرامیٹراں سمیت) تے:
\[
	(\forall x,z \in S)((\psi \liff \psi^S) \land (\chi \liff \chi^S))
\]
سو خاص طور اُتے:
\begin{align*}
	(\forall  x, z \in S)(&\phi(x, z) \liff \phi^S(x, z))\\
	(\forall x \in S)(&\exists y\phi(x, y) \liff (\exists y \in S)\phi^S(x, y))
\end{align*}
ایہناں نوں جوڑ کے تے ایہ ویکھ کے کہ $A \subseteq S$ کیوں جے $A \in S$ تے $S$ متعدی اے:
\begin{align*}
	(\forall x \in A)(&\exists y\phi(x, y) \liff (\exists y \in S)\phi(x, y))
\end{align*}
ہُن $(\forall x \in A)(\lexists![y \in S])\phi(x, y)$، کیوں جے
$(\forall x \in A)\lexists![y][\phi(x,y)]$۔ ہُن وکھ کرن دے اصول توں
$\Setabs{y \in S}{(\exists x \in A) \phi(x, y)} = \Setabs{y}{(\exists x \in
A)\phi(x, y)}$ ملدا اے۔
\end{proof}
\endgroup
% /OLP-0572
% OLP-0573 content/set-theory/replacement/finiteaxiomatizability.tex
\begingroup
\def\olfilename{finiteaxiomatizability}
\expandafter\def\csname ol@sectioncs\endcsname{section}

\olfileid{sth}{replacement}{finiteaxiomatizability}
\olsection{ضمیمہ: محدود گنتی دے مسلّمیاں نال بیان}
اسیں بدل رکھن دے اصول توں کجھ نتیجے کڈھ کے ایہ باب مکاواں گے۔ پہلا نتیجہ \citet{Montague1961} دا اے۔ دھیان دیو کہ ایہ $\ZF$ دے \emph{اندر} دتا ثبوت نہیں، ایہ $\ZF$ \emph{بارے} ثبوت اے۔
\begin{thm}\ollabel{zfnotfinitely}
	$\ZF$ نوں محدود گنتی دے مسلّمیاں نال بیان نہیں کیتا جا سکدا۔ ہور عام طور اُتے، جے $\Th{T}$ محدود اے تے $\Th{T} \Proves \ZF$، تاں $\Th{T}$ وچ ٹکراؤ اے۔

	(ایتھے بناں وکھرا آکھے اسیں اپنا کم پہلے درجے دے اوہناں جملیاں تک محدود کردے آں جیہناں دا اکّو اک غیر منطقی بنیادی نشان $\in$ اے۔ تے اسیں $\Th{T} \Proves \ZF$ ایہ دسّن لئی لکھدے آں کہ $\Th{T} \Proves \phi$ ہر $\phi \in \ZF$ لئی۔)
\end{thm}
\begin{proof}
	کوئی محدود $\Th{T}$ مقرر کرو جیہدے لئی $\Th{T} \Proves \ZF$۔ سو $\Th{T}$ توں انعکاس، یعنی \olref[sth][replacement][ref]{reflectionschema}، ثابت ہُندا اے۔ کیوں جے $\Th{T}$ محدود اے، اسیں ایہنوں اکّو عطف $\theta$ وجوں فیر لکھ سکدے آں۔ ایس فارمولے دا انعکاس لئیے تاں $\Th{T} \Proves \exists \beta(\theta \liff \theta^{V_\beta})$۔ کیوں جے صاف اے کہ $\Th{T} \Proves \theta$، سانوں $\Th{T} \Proves \exists \beta\ \theta^{V_\beta}$ ملدا اے۔

	ہُن $\psi(X)$ ایہدا مختصر روپ ہووے:
	\[
		\theta^X \land X\text{ متعدی اے} \land (\forall Y \in X)(Y\text{ متعدی اے}\lif \lnot \theta^{Y})
	\]
	موٹی گل ایہ اے کہ $X$، $\theta$ دا متعدی نمونہ اے تے ایس لحاظ نال $\in$ مطابق گھٹ توں گھٹ اے۔ ہُن چیتے کرو کہ $\Th{T} \Proves \exists \beta\ \theta^{V_\beta}$۔ درجیاں بارے $\ZF$ دے اندرلے، تے ایس لئی $\Th{T}$ دے اندرلے، بنیادی حقائق توں سانوں ملدا اے:
	\begin{equation}
		\Th{T} \Proves \exists M \psi(M). \tag{*}\label{Mpsi}
	\end{equation}
	$\psi(X)$ دے پہلے عطف والے جز نوں ورت کے، جدوں وی $\Th{T} \Proves \sigma$، سانوں $\Th{T} \Proves \forall X(\psi(X) \lif \sigma^X)$ ملدا اے۔ سو \eqref{Mpsi} مطابق:
	\begin{align*}
		\Th{T} &\Proves \forall X(\psi(X) \lif (\exists N \psi(N))^X)\\
	\intertext{ایہ ورت کے، تے \eqref{Mpsi} فیر ورت کے:}
		\Th{T} &\Proves \exists M(\psi(M) \land (\exists N \psi(N))^M)
	\intertext{سو خاص طور اُتے:}
		\Th{T} &\Proves \exists M(\psi(M) \land (\exists N \in M)((N\text{ متعدی اے})^N \land (\theta^N)^M))
	\intertext{سو تعدی بارے بنیادی دلیل نال:}
		\Th{T} &\Proves \exists M(\psi(M) \land (\exists N \in M)(N\text{ متعدی اے} \land \theta^N))
	\end{align*}
	سو $\Th{T}$ وچ ٹکراؤ اے۔\footnote{ایس «بنیادی دلیل» وچ متعدی سیٹاں لئی کجھ «مطلق ہون دے حقائق» ثابت کرنے پیندے نیں۔}
\end{proof}

ایہو جیہا اک ہور نتیجہ \citet[223]{Potter2004} نے دَسیا اے۔

\begin{prop}\ollabel{finiteextensionofZ}
	$\Th{T}$، $\Z$ وچ محدود گنتی دے نویں مسلّمے پا کے ملی توسیع ہووے۔ جے $\Th{T} \Proves \ZF$، تاں $\Th{T}$ وچ ٹکراؤ اے۔ (ایتھے اوہو حدّاں چُپ چاپ منّیاں نیں جیہڑیاں \olref{zfnotfinitely} لئی منّیاں سن۔)
\end{prop}
\begin{proof}
	$\theta$ نوں $\Th{T}$ دے سارے مسلّمیاں دا عطف منّو، وکھ کرن دے خاکے دیاں (لامتناہی گنتی والیاں) صورتاں \emph{چھڈ کے}۔ \olref{zfnotfinitely} وانگ $\psi$ دی تعریف $\theta$ توں دے کے اسیں وکھا سکدے آں کہ $\Th{T} \Proves \exists M \psi(M)$۔

	\olref{zfnotfinitely} وانگ اسیں ایہ خاکہ ثابت کر سکدے آں کہ جدوں وی $\Th{T} \Proves \sigma$، سانوں $\Th{T} \Proves \forall X(\psi(X) \lif \sigma^X)$ ملدا اے۔ فیر اسیں ثبوت این اویویں پورا کردے آں جیویں \olref{zfnotfinitely} وچ کیتا سی۔

	پر خاکہ ثابت کرن لئی \olref{zfnotfinitely} نالوں ذرا ہور کم کرنا پیندا اے۔ آخر وکھ کرن دے خاکے دیاں صورتاں $\Th{T}$ وچ شامل نیں، پر $\theta$ دے عطف والے جز نہیں۔ پر اسیں ایہ اوکڑ ایہ ثابت کر کے ٹپ سکدے آں کہ $\Th{T} \Proves \forall X(X\text{ متعدی اے} \lif \sigma^X)$ وکھ کرن دی ہر صورت $\sigma$ لئی۔ ایہ کم اسیں پڑھن والے اُتے چھڈدے آں۔
\end{proof}
\begin{prob}
	وکھ کرن دے خاکے دی ہر صورت $\sigma$ لئی وکھاؤ کہ $\Z \Proves \forall X(X\text{ متعدی اے} \lif \sigma^X)$۔ (اسیں ایہ خاکہ \olref[sth][replacement][finiteaxiomatizability]{finiteextensionofZ} وچ ورتیا سی۔)
\end{prob}
\begin{prob}
	وکھاؤ کہ ہر $\phi \in \Z$ لئی $\ZF \Proves \phi^{V_{\omega+ \omega}}$۔
\end{prob}
\begin{prob}
	\olref[sth][replacement][finiteaxiomatizability]{zfnotfinitely} تے \olref[sth][replacement][finiteaxiomatizability]{finiteextensionofZ} دے ثبوتاں وچ ورتے باقی خاکے والے نتیجیاں دی تصدیق کرو۔
\end{prob}

جیویں \olref[sth][replacement][absinf]{sec} وچ آکھیا سی، ایہ وکھاندا اے کہ بدل رکھن دا اصول \olref[ordinals][ordtype]{thmOrdinalRepresentation} توں ٹھیک ٹھیک ودھ طاقتور اے۔ یا ذرا ہور ٹھیک آکھیے، جے $\Z$ + «ہر حُسنِ ترتیب اکّو اک ترتیبی عدد نال ہم ساخت اے» بے ٹکراؤ اے، تاں ایہ بدل رکھن دے خاکے دی گھٹ توں گھٹ اک صورت ثابت نہیں کر سکدا۔

	%	By assumption, $\psi^{V_\alpha}$ has the form: 
	%	\[
	%		(\forall A \in V_\alpha)(\exists S \in V_\alpha)(\forall x \in V_\alpha)(x \in S \liff (\phi^{V_\alpha}(x) \land x \in A))
	%	\]
	%	To establish this holds, fix $A \in V_\alpha$. Using Separation, obtain: 
	%	\[
	%		S = \Setabs{x \in A}{\phi^{V_\alpha}(x)}
	%	\]
	%	Now $S \in V_\alpha$, since $S \subseteq A \in V_\alpha$, and clearly $(\forall x \in V_\alpha)(x \in S \liff (\phi^{V_\alpha}(x) \land x \in A)$.
\endgroup
% /OLP-0573

\OLEndChapterHook
\endgroup
% /OLP-0565
% OLP-0574 content/set-theory/ord-arithmetic/ord-arithmetic.tex
\begingroup
\def\olfilename{ord-arithmetic}
\expandafter\def\csname ol@sectioncs\endcsname{section}

\olchapter{sth}{ord-arithmetic}{ترتیبی عدداں دا حساب}

% OLP-0575 content/set-theory/ord-arithmetic/introduction.tex
\begingroup
\def\olfilename{introduction}
\expandafter\def\csname ol@sectioncs\endcsname{section}
\olfileid{sth}{ord-arithmetic}{intro}

\olsection{جان پہچان}

\olref[ordinals][]{chap} وچ اسیں ترتیبی عدداں دا نظریہ بنایا سی۔
\olref[spine][]{chap} وچ اسیں ویکھیا کہ ترتیبی عدداں نوں ریڑھ دی ہڈی وانگ سمجھ سکدے آں، جیہدے دوالے درجہ وار سلسلے دا باقی حصہ بندا اے۔ پر ترتیبی عدداں دا کم بس ایہو نہیں۔ اوہناں اُتے حساب کرن دا کم وی اے۔

اسیں ایس ول پہلاں ای \olref[ordinals][idea]{sec} وچ اشارہ کیتا سی، جدوں اسیں $\omega$، $\omega+1$ تے $\omega+\omega$ دی گل کیتی سی۔ اوس ویلے اسیں غیر رسمی طور اُتے گل کیتی سی۔ ہُن ایس نوں ٹھیک طرح کھولھ کے دسّن دا ویلا آ گیا اے۔ پر ایہ دس دیئیے کہ ایس باب وچ فلسفہ بہت نہیں۔ ایتھے فنی گلّاں اگے ودھدیاں نیں، تے نال اک (کجھ) دلچسپ گل ویکھاں گے کہ اسیں اکّو کم دو وکھرے طریقیاں نال کر سکدے آں۔
\endgroup
% /OLP-0575
% OLP-0576 content/set-theory/ord-arithmetic/addition.tex
\begingroup
\def\olfilename{addition}
\expandafter\def\csname ol@sectioncs\endcsname{section}

\olfileid{sth}{ord-arithmetic}{add}
\olsection{ترتیبی جوڑ}

منّ لو اسیں $\alpha$ تے $\beta$ نوں جوڑنا چاہندے آں۔ اسیں $\beta$ دی اک {نقل}، $\alpha$ دی اک نقل توں عین پچھوں رکھ سکدے آں۔ (سانوں \emph{نقلاں} لینیاں پیندیاں نیں، کیوں جے \olref[ordinals][basic]{ordinalsaresubsets} توں سانوں پتا اے کہ یا تاں $\alpha
\subseteq \beta$ یا $\beta \subseteq \alpha$۔) سہج سمجھ مطابق ایہدا اثر ایہ اے کہ اسیں قدماں دی اک $\alpha$ لڑی توں لنگھدے آں \emph{تے فیر} اک $\beta$ لڑی توں لنگھدے آں۔ جیہڑی لڑی بنے گی اوہ حُسنِ ترتیب والی ہووے گی۔ سو \olref[ordinals][ordtype]{thmOrdinalRepresentation} مطابق اوہ اک (اکّو اک) ترتیبی عدد نال ہم ساخت اے۔ اوس ترتیبی عدد نوں $\alpha$ تے $\beta$ دا \emph{جوڑ} سمجھ سکدے آں۔

ایہ ترتیبی جوڑ دے پچھے سہج خیال اے۔ اینوں ٹھیک طرح تعریف کرن لئی پہلاں اسیں سیٹاں دیاں \emph{نقلاں} لین دا خیال ورتدے آں۔ طریقہ ایہ اے کہ اپنی مرضی دے نشان، $0$ تے $1$، لا کے چیتا رکھیے کہ کیہڑی شے کتھوں آئی۔

\begin{defn}\ollabel{defdissum}
$A$ تے $B$ دا \emph{وکھ وکھ حصیاں والا جوڑ} $A \disjointsum B = (A\times
\{0\}) \cup (B \times \{1\})$ اے۔
\end{defn}

اگے اسیں ترتیبی عدداں دیاں جوڑیاں اُتے اک ترتیب تعریف کردے آں۔

\begin{defn}
کسے وی ترتیبی عدداں $\alpha_1, \alpha_2, \beta_1, \beta_2$ لئی ایویں آکھو:
\begin{align*}
	\tuple{\alpha_1, \alpha_2} \rlexless \tuple{\beta_1, \beta_2}\text{ اُدوں تے صرف اُدوں جدوں }&
	\text{یا تاں $\alpha_2 \in \beta_2$}\\
	& \text{یا دوویں $\alpha_2 = \beta_2$ تے $\alpha_1 \in \beta_1$}
\end{align*}
\end{defn}

ایہ \emph{اُلٹی لغتی ترتیب} اے، کیوں جے پہلاں دوجا رکن ویکھ کے ترتیب لاندے او، تے فیر پہلا۔ چیتے کرو کہ اسیں $\alpha \ordplus \beta$ نوں $\alpha$ دی اک نقل تے اوہدے پچھوں $\beta$ دی اک نقل دی ترتیب دی قسم وجوں تعریف کرنا چاہندے ساں۔ ایہ کرن لئی اسیں آکھدے آں:

\begin{defn}\ollabel{defordplus}
کسے وی ترتیبی عدداں $\alpha$، $\beta$ لئی، اوہناں دا جوڑ $ \alpha \ordplus
\beta = \ordtype{\alpha \disjointsum \beta, \rlexless}$ اے۔
\end{defn}
\noindent
دھیان دیو کہ ایتھے اسیں لکھت نوں ذرا کھلّھ نال ورتیا اے۔ پورا ٹھیک لکھن لئی «$\Setabs{\tuple{x,y}\in \alpha\disjointsum\beta}{x \rlexless y}$» لکھنا چاہیدا اے، «$\rlexless$» دی تھاں۔ پر مختصر رکھن لئی اسیں اگے وی ایویں ای لکھاں گے۔

ہِیٹھلا نتیجہ، \olref[ordinals][ordtype]{thmOrdinalRepresentation} نال رَل کے، پکا کردا اے کہ ساڈی تعریف ٹھیک بنی اے۔

\begin{lem}\ollabel{ordsumlessiswo}
$\tuple{\alpha \disjointsum \beta, \rlexless}$ حُسنِ ترتیب والی ساخت اے، کسے وی ترتیبی عدداں $\alpha$ تے $\beta$ لئی۔
\end{lem}

\begin{proof}
صاف اے کہ $\rlexless$، $\alpha \disjointsum \beta$ اُتے جُڑیا ہویا اے۔ ایہ وکھاؤن لئی کہ اوہ بنیاد والا تعلق اے، کوئی غیر خالی $X \subseteq \alpha
\disjointsum \beta$ مقرر کرو۔ $Y$، $X$ دا اوہ ذیلی سیٹ ہووے جیہدے رکناں دا دوجا محدد جِنّا ہو سکے چھوٹا اے، یعنی $Y = \Setabs{\tuple{\gamma, i} \in X}{(\forall \tuple{\delta, j} \in X)i \leq j}$۔ ہُن $Y$ دا اوہ رکن چُنو جیہدا پہلا محدد سبھ توں چھوٹا اے۔
\end{proof}
\noindent
سو ساڈے کول ترتیبی جوڑ دی اک صاف تے کھلّھی تعریف اے۔ ہِیٹھلی گل کوئی حیرانی والی نہیں (چیتے کرو کہ $1 = \{0\}$، \olref[z][infinity-again]{defnomega} مطابق):
\begin{prop}
$\alpha \ordplus  1 = \ordsucc{\alpha}$، کسے وی ترتیبی عدد $\alpha$ لئی۔
\end{prop}

\begin{proof}
ہم ساخت نقشہ $f$ ویکھو، جیہڑا $\ordsucc{\alpha} = \alpha \cup
\{\alpha\}$ توں $\alpha\disjointsum1 = (\alpha \times \{0\})
\disjointsum (\{0\} \times \{1\})$ ول اے، تے $f(\gamma) =
\tuple{\gamma, 0}$ نال دتا اے $\gamma \in \alpha$ لئی، تے $f(\alpha) = \tuple{0,
1}$۔
\end{proof}
\noindent
نالے وکھانا سوکھا اے کہ جوڑ کجھ بازگشتی شرطاں اُتے پورا اتردا اے۔

\begin{lem}\ollabel{ordadditionrecursion}
کسے وی ترتیبی عدداں $\alpha, \beta$ لئی:
\begin{align*}
	\alpha\ordplus 0 &= \alpha\\
	\alpha \ordplus  (\beta\ordplus 1) &= (\alpha \ordplus  \beta) \ordplus  1\\
	\alpha  \ordplus  \beta &= \supstrict_{\delta < \beta}(\alpha \ordplus  \delta) && \text{جے $\beta $ حدّی ترتیبی عدد اے}
\end{align*}
\end{lem}

\begin{proof}
اسیں واری واری ہر صورت ویکھدے آں۔ پہلاں:
\begin{align*}
	\alpha \ordplus  0 
	& = \ordtype{(\alpha \times \{0\}) \cup (0 \times \{1\}), \rlexless} \\
	&= \ordtype{(\alpha \times \{0\}) \cup \{0\}, \rlexless}\\
	&= \alpha\\
	\alpha \ordplus (\beta \ordplus  1) 
	%&= \ordtype{(\alpha\times \{0\}) \cup (\ordtype{\beta \ordplus  1}\times \{1\}), \rlexless} \\
	&= \ordtype{(\alpha\times \{0\}) \cup (\ordsucc{\beta}\times \{1\}), \rlexless} \\
	&= \ordtype{(\alpha\times \{0\}) \cup (\beta \times \{1\}), \rlexless} \ordplus 1\\
	&= (\alpha \ordplus  \beta) \ordplus  1
\end{align*}
ہُن $\beta \neq \emptyset$ حدّی ہووے۔ جے $\delta < \beta$ تاں $\delta\ordplus 1 < \beta$ وی، سو $\alpha \ordplus  \delta$، $\alpha \ordplus  \beta$ دا حقیقی مُڈھلا ٹکڑا اے۔ سو $\alpha
\ordplus  \beta$، $X = \Setabs{\alpha
\ordplus  \delta}{\delta < \beta}$ دی سخت اُپَرلی حد اے۔ نالے جے $\alpha \leq \gamma <
\alpha \ordplus  \beta$ تاں صاف اے کہ $\gamma = \alpha \ordplus
\delta$ کسے $\delta < \beta$ لئی۔ سو $\alpha \ordplus  \beta =
\supstrict_{\delta< \beta}(\alpha\ordplus \delta)$۔
\end{proof}

پر ایہ اک خاص گل اے۔ ترتیبی جوڑ دی تعریف کرن لئی اسیں \emph{اِس دی تھاں} بس لامتناہی درجیاں والی بازگشت دا قضیہ ورت سکدے ساں، تے بازگشت دیاں مساواتاں این اویویں رکھ سکدے ساں جیویں \olref{ordadditionrecursion} وچ دتیاں نیں (بھاویں «$\ordsucc{\beta}$» ورتدے، «$\beta \ordplus 1$» دی تھاں)۔

سو ترتیبی عدداں اُتے عملاں دی تعریف دے دو وکھرے طریقے نیں۔ اسیں اک حُسنِ ترتیب والا سیٹ صاف صاف بنا کے تے اوہدی ترتیب دی قسم ویکھ کے، \emph{بناوٹ راہیں} عمل دی تعریف کر سکدے آں۔ یا بس بازگشت دیاں مساواتاں رکھ کے، \emph{بازگشت راہیں} تعریف کر سکدے آں۔ پر ٹھیک طرح کریے تاں حاصل اکّو اے۔ کیوں جے \olref[ordinals][ordtype]{thmOrdinalRepresentation} ضمانت دیندا اے کہ ایہ مساواتاں \emph{اکّو اک} ترتیبی عدد طے کردیاں نیں۔

کئی طرحاں ترتیبی حساب، قدرتی عدداں دے جوڑ وانگ ای چلدا اے۔ مثلاً اسیں ایہ ثابت کر سکدے آں:

\begin{lem}\ollabel{ordinaladditionisnice}
جے $\alpha, \beta, \gamma$ ترتیبی عدد نیں، تاں:
\begin{enumerate}
	\item\ollabel{ordaddition1} جے $\beta < \gamma$ تاں $\alpha
	\ordplus  \beta < \alpha \ordplus  \gamma$
	\item\ollabel{ordaddition2} جے $\alpha \ordplus  \beta =
	\alpha\ordplus \gamma$ تاں $\beta = \gamma$
	\item\ollabel{ordaddition3} $\alpha \ordplus  (\beta \ordplus
	\gamma) = (\alpha \ordplus  \beta) \ordplus  \gamma$، یعنی جوڑ شرکتیت رکھدا اے۔
	\item\ollabel{ordaddition4} جے $\alpha \leq \beta$ تاں $\alpha
	\ordplus  \gamma \leq \beta \ordplus \gamma$
\end{enumerate}
\end{lem}

\begin{proof}
اسیں \olref{ordaddition3} ثابت کردے آں تے باقی مشق لئی چھڈدے آں۔ ثبوت $\gamma$ اُتے سادہ لامتناہی درجیاں والے استقرا نال اے، \olref{ordadditionrecursion} ورت کے۔ جدوں $\gamma = 0$:
\[
(\alpha \ordplus  \beta) \ordplus  0 = \alpha \ordplus  \beta  = \alpha \ordplus  (\beta \ordplus  0)
\]
جدوں $\gamma = \delta\ordplus 1$، استقرا لئی منّ لو $(\alpha
\ordplus  \beta) \ordplus  \delta = \alpha \ordplus  (\beta \ordplus
\delta)$۔ ہُن \olref{ordadditionrecursion} تِنّ وار ورت کے:
\begin{align*}
	(\alpha \ordplus  \beta) \ordplus  (\delta \ordplus  1) & = ((\alpha \ordplus  \beta) \ordplus  \delta)\ordplus 1\\
	& = (\alpha \ordplus  (\beta \ordplus  \delta)) \ordplus  1\\
	& = \alpha \ordplus  ((\beta \ordplus  \delta)\ordplus 1)\\
	& = \alpha \ordplus  (\beta \ordplus  (\delta\ordplus 1))
\end{align*}
جدوں $\gamma$ حدّی ترتیبی عدد اے، استقرا لئی منّ لو کہ جے $\delta \in \gamma$ تاں $(\alpha \ordplus  \beta) \ordplus  \delta =
\alpha \ordplus  (\beta \ordplus  \delta)$۔ ہُن:
\begin{align*}
	(\alpha \ordplus  \beta) \ordplus  \gamma & = \supstrict_{\delta < \gamma}((\alpha \ordplus  \beta) \ordplus  \delta) \\
	&= \supstrict_{\delta < \gamma}(\alpha \ordplus  (\beta \ordplus  \delta))\\
	&= \alpha \ordplus  \supstrict_{\delta < \gamma}(\beta \ordplus  \delta)\\
	& = \alpha \ordplus  (\beta \ordplus  \gamma)
\end{align*}
\end{proof}

\begin{prob}
\olref[sth][ord-arithmetic][add]{ordinaladditionisnice} دا باقی حصہ ثابت کرو۔
\end{prob}

ایہناں طرحاں ترتیبی جوڑ تہانوں بڑا جانیا پچھانیا لگنا چاہیدا اے۔ پر اک بہت ضروری گل وچ ترتیبی جوڑ قدرتی عدداں دے جوڑ ورگا \emph{نہیں} اے۔

\begin{prop}\ollabel{ordsumnotcommute}
ترتیبی جوڑ {تبادلی نہیں} اے۔ $1 \ordplus  \omega = \omega <
\omega \ordplus  1$۔
\end{prop}

\begin{proof}
دھیان دیو کہ $1 \ordplus  \omega = \supstrict_{n < \omega} (1 \ordplus n)
= \omega \in \omega \cup \{\omega\} = \ordsucc{\omega} = \omega
\ordplus  1$۔
\end{proof}
\noindent
بھاویں پہلی واری ایہ گل حیران کرے، \emph{ایہ گل حیرانی والی نہیں ہونی چاہیدی}۔ اک پاسے، جدوں تسی $1 \ordplus  \omega$ ویکھدے او، تسی اوہ ترتیب دی قسم سوچدے او جیہڑی سارے قدرتی عدداں توں \emph{پہلاں} اک ہور رکن رکھن نال ملدی اے۔ \olref[sfr][infinite][hilbert]{sec} وچ ہلبرٹ دے ہوٹل والی دلیل وانگ سوچیے، تاں سہج سمجھ مطابق ایس وادھو پہلے رکن نال پوری ترتیب دی قسم نوں کوئی فرق نہیں پینا چاہیدا۔ دوجے پاسے، جدوں تسی $\omega \ordplus  1$ ویکھدے او، تسی اوہ ترتیب دی قسم سوچدے او جیہڑی سارے قدرتی عدداں توں \emph{پچھوں} اک ہور رکن رکھن نال ملدی اے۔ تے ایہ بالکل وکھری گل اے!
\endgroup
% /OLP-0576
% OLP-0577 content/set-theory/ord-arithmetic/using-addition.tex
\begingroup
\def\olfilename{using-addition}
\expandafter\def\csname ol@sectioncs\endcsname{section}

\olfileid{sth}{ord-arithmetic}{using-addition}
\olsection{ترتیبی جوڑ دی ورتوں}

ترتیبی عدداں دا جوڑ ورت کے اسیں وکھ وکھ سیٹاں دے درجے، \olref[spine][rank]{defnsetrank} دے معنی وچ، سِدھے ای حساب لا کے کڈھ سکدے آں۔

\begin{lem}\ollabel{rankcomputation}
جے $\setrank{A} = \alpha$ تے $\setrank{B} = \beta$، تاں:
\begin{enumerate}
	\item\ollabel{exrankpow} $\setrank{\Pow{A}} = \alpha\ordplus 1$
	\item\ollabel{exrankpair} $\setrank{\{A, B\}} = \max(\alpha,
	\beta) \ordplus 1$
	\item\ollabel{exrankcup} $\setrank{A \cup B} = \max(\alpha,
	\beta)$
	\item\ollabel{exranktuple} $\setrank{\tuple{A,B}} = \max(\alpha,
	\beta) \ordplus  2$
	\item\ollabel{exranktimes} $\setrank{A \times B} \leq \max(\alpha,
	\beta) \ordplus  2$
	\item\ollabel{exrankunion} $\setrank{\bigcup A} = \alpha$ جدوں
	$\alpha$ خالی یا حدّی اے۔ $\setrank{\bigcup A} = \gamma$ جدوں
	$\alpha = \gamma\ordplus 1$۔
\end{enumerate}
\end{lem}

\begin{proof}
سارے ثبوت وچ اسیں \olref[spine][rank]{ranksupstrict} وار وار ورتدے آں۔

\emph{\olref{exrankpow}۔} جے $x \subseteq A$ تاں $\setrank{x} \leq
\setrank{A}$۔ سو $\setrank{\Pow{A}} \leq \alpha \ordplus  1$۔ کیوں جے خاص طور اُتے $A
\in \Pow{A}$، ایس لئی $\setrank{\Pow{A}} = \alpha \ordplus  1$۔

\emph{\olref{exrankpair}۔} \olref[spine][rank]{ranksupstrict} مطابق۔

\emph{\olref{exrankcup}۔} \olref[spine][rank]{ranksupstrict} مطابق۔

\emph{\olref{exranktuple}۔} \olref{exrankpair} دو واری ورتو۔

\emph{\olref{exranktimes}۔} دھیان دیو کہ $A \times B \subseteq
\Pow{\Pow{A \cup B}}$، تے \olref{exranktuple} ورتو۔

\emph{\olref{exrankunion}۔} جے $\alpha = \gamma\ordplus 1$، تاں کوئی
$c \in A$ اے جیہدے لئی $\setrank{c} = \gamma$، تے $A$ دا کوئی !!{element} ایہتوں ودھ درجے دا نہیں۔ سو $\setrank{\bigcup A} = \gamma$۔ جے $\alpha$ حدّی ترتیبی عدد اے، تاں $A$ وچ ایہو جیہے !!{element}s نیں جیہناں دا درجہ جِنّا چاہو $\alpha$ دے نیڑے (پر سختی نال اوہتوں گھٹ) اے۔ ایس لئی $\bigcup A$ وچ وی ایہو جیہے !!{element}s نیں جیہناں دا درجہ جِنّا چاہو $\alpha$ دے نیڑے (پر سختی نال اوہتوں گھٹ) اے۔ سو $\setrank{\bigcup A} = \alpha$۔
\end{proof}
\noindent
\olref{exranktimes} وچ \emph{نا}برابری کیوں اے، ایہ وکھانا اسیں مشق لئی چھڈدے آں۔

\begin{prob}
ایہو جیہے سیٹ $A$ تے $B$ دیو کہ $\setrank{A \times B}=
\max(\setrank{A}, \setrank{B})$۔ ایہو جیہے سیٹ $A$ تے $B$ دیو کہ
$\setrank{A \times B}\max(\setrank{A}, \setrank{B}) \ordplus  2$۔
کی کوئی ہور درجہ ممکن اے؟
\end{prob}

ہُن اسیں ایہ وی وکھا سکدے آں کہ کسے ترتیبی عدد نوں «محدود» یا «لامتناہی» آکھن دے کئی واجب معنی آپس وچ ملدے نیں۔

\begin{lem}\ollabel{ordinfinitycharacter}
کسے وی ترتیبی عدد $\alpha$ لئی ہِیٹھلیاں گلاں برابر نیں:
\begin{enumerate}
	\item\ollabel{ord:notinomega} $\alpha\notin \omega$، یعنی
	$\alpha$ قدرتی عدد نہیں
	\item\ollabel{ord:omegaplus} $\omega \leq \alpha$
	\item\ollabel{ord:oneplus} $1 \ordplus  \alpha = \alpha$
	%\alpha \approx \alpha \ordplus 1$
	\item\ollabel{ord:plusone} $\alpha \approx \alpha\ordplus 1$،
	یعنی $\alpha$ تے $\alpha\ordplus 1$ دے رکناں دی گنتی برابر اے
	\item\ollabel{ord:infinite} $\alpha$ ڈیڈیکنڈ لامتناہی اے
	\end{enumerate}
\end{lem}
\noindent
سو ساڈے کول اک ترتیبی عدد دے محدود یا لامتناہی ہون نوں سمجھن دے پنج طریقے نیں، جیہناں دی برابری ثابت کیتی جا سکدی اے۔

\begin{proof}
\emph{\olref{ord:notinomega} $\Rightarrow$ \olref{ord:omegaplus}۔}
تِن صورتاں دے اصول مطابق۔

\emph{\olref{ord:omegaplus} $\Rightarrow$ \olref{ord:oneplus}۔}
$\alpha \geq \omega$ مقرر کرو۔ لامتناہی درجیاں اُتے استقرا نال اک سبھ توں چھوٹا ترتیبی عدد $\gamma$ اے (شاید $0$)، جیہدے لئی اک حدّی ترتیبی عدد $\beta$ اے ایویں کہ $\alpha = \beta \ordplus \gamma$۔ ہُن:
\[
	1 \ordplus \alpha =
	1 \ordplus (\beta \ordplus \gamma) =
	(1 \ordplus \beta) \ordplus \gamma =
	\supstrict_{\delta < \beta} (1 \ordplus  \delta) \ordplus  \gamma =
	\beta \ordplus  \gamma =
	\alpha.
\]
\emph{\olref{ord:oneplus} $\Rightarrow$ \olref{ord:plusone}۔}
صاف اے کہ اک !!a{bijection} $f \colon (\alpha \disjointsum 1) \to (1
\disjointsum \alpha)$ اے۔ جے $1 \ordplus \alpha = \alpha$، تاں اک ہم ساخت نقشہ $g \colon (1 \disjointsum \alpha) \to \alpha$ اے۔ ہُن $\comp{f}{g}$ ویکھو۔

\emph{\olref{ord:plusone} $\Rightarrow$ \olref{ord:infinite}۔}
جے $\alpha \approx \alpha \ordplus 1$، تاں اک !!a{bijection} $f \colon
(\alpha \disjointsum 1) \to \alpha$ اے۔ $g(\gamma) = f(\gamma, 0)$ تعریف کرو ہر $\gamma < \alpha$ لئی۔ ایہ !!{injection} گواہ اے کہ $\alpha$ ڈیڈیکنڈ لامتناہی اے، کیوں جے $f(0,1) \in \alpha \setminus \ran{g}$۔

\emph{\olref{ord:infinite} $\Rightarrow$ \olref{ord:notinomega}۔}
ایہ \olref[z][infinity-again]{naturalnumbersarentinfinite} اے۔
\end{proof}
\endgroup
% /OLP-0577
% OLP-0578 content/set-theory/ord-arithmetic/multiplication.tex
\begingroup
\def\olfilename{multiplication}
\expandafter\def\csname ol@sectioncs\endcsname{section}

\olfileid{sth}{ord-arithmetic}{mult}

\olsection{ترتیبی ضرب}

ہُن اسیں ترتیبی ضرب ول آؤندے آں۔ ایہنوں تقریباً ترتیبی جوڑ وانگ سمجھاں گے۔ سو منّ لو اسیں $\alpha$ نوں~$\beta$ نال ضرب دینا چاہندے آں۔ ایس لئی اک مستطیل جال دی کلپنا کرو، جیہدی چوڑائی $\alpha$ تے اُچائی $\beta$ اے۔ $\alpha$ تے $\beta$ دا حاصلِ ضرب ایویں ملدا اے کہ اک قطار دے سرے توں دوجے سرے تک لنگھو، فیر اگلی قطار وچ ایویں ای لنگھو\ldots جدوں تک سارے جال وچوں نہ لنگھ جاؤ۔ دوجے لفظاں وچ، $\alpha$ تے $\beta$ دا حاصلِ ضرب $\beta$ دے \emph{ہر} رکن دی تھاں $\alpha$ دی اک نقل رکھن نال بندا اے۔

اینوں رسمی صورت دین لئی اسیں $\alpha$ تے $\beta$ دے کارتیسی حاصلِ ضرب اُتے اُلٹی لغتی ترتیب ورتدے آں۔

\begin{defn}
کسے وی ترتیبی عدداں $\alpha, \beta$ لئی اوہناں دا حاصلِ ضرب $\alpha \ordtimes \beta = \ordtype{\alpha \times \beta, \rlexless}$ اے۔
\end{defn}
\noindent
اسیں فیر تصدیق کرنی اے کہ ایہ تعریف ٹھیک بنی اے۔

\begin{lem}\ollabel{ordtimeslessiswo}
$\tuple{\alpha \times \beta, \rlexless}$ حُسنِ ترتیب والی ساخت اے، کسے وی ترتیبی عدداں $\alpha$ تے $\beta$ لئی۔
\end{lem}

\begin{proof}
این اویویں جیویں \olref[add]{ordsumlessiswo} لئی۔
\end{proof}
\noindent
تے ایہ ثابت کرنا اوکھا نہیں کہ ضرب ایس طرح چل دی اے۔

\begin{lem}\ollabel{ordtimesrecursion}
کسے وی ترتیبی عدداں $\alpha, \beta$ لئی:
\begin{align*}
	\alpha \ordtimes 0 &= 0\\
	\alpha \ordtimes (\beta \ordplus 1) &=
		(\alpha \ordtimes \beta) \ordplus \alpha\\
	\alpha  \ordtimes \beta &=
		\supstrict_{\delta < \beta}(\alpha \ordtimes \delta) &&
		\text{جدوں $\beta$ حدّی ترتیبی عدد اے}.
\end{align*}
\end{lem}

\begin{proof}
مشق لئی چھڈیا اے۔
\end{proof}

جوڑ وانگ ای، اسیں سِدھی تعریف دین دی تھاں ترتیبی ضرب دی تعریف ایہناں بازگشتی مساواتاں نال وی کر سکدے ساں۔ تے جوڑ وانگ ای، ضرب وچ وی کجھ جانیاں پچھانیاں خاصیتاں نیں۔

\begin{lem}\ollabel{ordinalmultiplicationisnice}
جے $\alpha, \beta, \gamma$ ترتیبی عدد نیں، تاں:
\begin{enumerate}
	\item\ollabel{ordtimes1} جے $\alpha \neq 0$ تے $\beta < \gamma$،
	تاں $\alpha \ordtimes \beta < \alpha \ordtimes \gamma$۔
	\item\ollabel{ordtimes2} جے $\alpha \neq 0$ تے $\alpha \ordtimes
	\beta = \alpha\ordtimes\gamma$، تاں $\beta = \gamma$۔
	\item\ollabel{ordtimes3}  $\alpha \ordtimes (\beta \ordtimes
	\gamma) = (\alpha \ordtimes \beta) \ordtimes \gamma$۔
	\item\ollabel{ordtimes4} جے $\alpha \leq \beta$، تاں $\alpha
	\ordtimes \gamma \leq \beta \ordtimes \gamma$۔
	\item\ollabel{ordtimes5}  $\alpha \ordtimes (\beta \ordplus
	\gamma) = (\alpha \ordtimes \beta )\ordplus (\alpha\ordtimes
	\gamma)$۔
\end{enumerate}
\end{lem}

\begin{proof}
مشق لئی چھڈیا اے۔
\end{proof}

تسی جِنّے چاہو ہور نتیجے ثابت کر سکدے او (یا لبھ کے پڑھ سکدے او)۔ پر \olref[ord-arithmetic][add]{ordsumnotcommute} نوں چیتے رکھو، تاں ہِیٹھلی گل اُتے حیرانی نہیں ہونی چاہیدی۔

\begin{prop}
ترتیبی ضرب تبادلی نہیں اے۔ $2 \ordtimes \omega  =
\omega < \omega \ordtimes 2$
\end{prop}

\begin{proof}
$2 \ordtimes \omega = \supstrict_{n < \omega}(2\ordtimes  n) = \omega \in \supstrict_{n < \omega}(\omega \ordplus n) = \omega \ordplus \omega = \omega \ordtimes 2$۔
\end{proof}
\noindent
ایہدی سہج سمجھ والی وجہ فیر بڑی سِدھی اے۔ $2
\ordtimes \omega$ کڈھن لئی ہر قدرتی عدد دی تھاں دو شےواں رکھو۔ جے تسی قدرتی عدداں وچ سارے «ادھ» والے عدد پا دیو، یعنی $\Setabs{\nicefrac{n}{2}}{n \in \omega}$ اُتے قدرتی ترتیب ویکھو، تاں اوہو ترتیب دی قسم ملے گی۔ ایویں آکھیاں صاف اے کہ ترتیب دی قسم آپ $\omega$ دی ترتیب دی قسم ورگی اے۔ پر $\omega \ordtimes 2$ کڈھن لئی $\omega$ دیاں دو نقلاں اک دوجی دے پچھوں رکھو۔

\begin{prob}
\olref[sth][ord-arithmetic][mult]{ordtimeslessiswo}،
\olref[sth][ord-arithmetic][mult]{ordtimesrecursion}،
تے
\olref[sth][ord-arithmetic][mult]{ordinalmultiplicationisnice}
ثابت کرو۔
\end{prob}
\endgroup
% /OLP-0578
% OLP-0579 content/set-theory/ord-arithmetic/exponentiation.tex
\begingroup
\def\olfilename{exponentiation}
\expandafter\def\csname ol@sectioncs\endcsname{section}

\olfileid{sth}{ord-arithmetic}{expo}
\olsection{ترتیبی گھات}

ہُن اسیں ترتیبی گھات ول آؤندے آں۔ افسوس، ترتیبی گھات دی بناوٹ راہیں کوئی \emph{سوکھی} تعریف نہیں۔

بے شک بناوٹ راہیں صاف تعریفاں \emph{موجود} نیں۔ اک تعریف ایہ اے۔ $\text{finfun}(\alpha,\beta)$ اوہناں سارے فنکشناں $f
\colon \alpha \to \beta$ دا سیٹ ہووے، جیہناں لئی $\Setabs{\gamma \in
\alpha}{f(\gamma) \neq 0}$ دے رکناں دی گنتی کسے قدرتی عدد دے رکناں دی گنتی دے برابر اے۔ $\text{finfun}(\alpha,\beta)$ اُتے حُسنِ ترتیب ایویں تعریف کرو۔ $f
\sqsubset g$ اُدوں تے صرف اُدوں جدوں $f \neq g$ تے $f(\gamma_0) < g(\gamma_0)$، جتھے
$\gamma_0 = \text{max}\Setabs{\gamma \in \alpha}{f(\gamma) \neq
g(\gamma)}$۔ فیر $\ordexpo{\alpha}{\beta}$ دی تعریف
$\ordtype{\text{finfun}(\alpha, \beta), \sqsubset}$ وجوں دے سکدے آں۔ پوٹر ایہ صاف تعریف ورتدا اے تے نال ای سمجھاندا اے:
\begin{quote}
	ایہ ترتیب بس ایس خواہش کرکے چُنی اے کہ ترتیبی گھات دی ایہو جیہی تعریف ملے جیہڑی ڈھکویں بازگشتی شرط اُتے پوری اُترے\ldots، تے ایہدی صورت ذہن وچ بناؤنی ترتیب والے جوڑ یا ترتیب والی ضرب نالوں کِتے اوکھی اے۔
	\citep[p.~199]{Potter2004}
\end{quote}
ایہ گل ٹھیک اے۔ اسیں جوڑ نوں «$\alpha$ دی اک نقل تے اوہدے پچھوں $\beta$ دی اک نقل» آکھ کے سمجھایا سی، تے ضرب نوں «$\beta$ جِنّی لمّی لڑی، $\alpha$ دیاں نقلاں دی» آکھ کے۔ پر $\text{finfun}(\alpha, \gamma)$ بارے ساڈے کول ایہو جیہی کوئی نِکّی جیہی گل نہیں۔ سو اسیں ترتیبی گھات دی تعریف سِدھی لامتناہی درجیاں والی بازگشت \emph{راہیں} دیواں گے، یعنی:

\begin{defn}\ollabel{ordexporecursion}
\begin{align*}
	\ordexpo{\alpha}{0} &= 1\\
	\ordexpo{\alpha}{\beta\ordplus 1} &=\ordexpo{\alpha}{\beta} \ordtimes \alpha\\
	\ordexpo{\alpha}{\beta} &= \bigcup_{\delta < \beta}\ordexpo{\alpha}{\delta}& & \text{جدوں $\beta$ حدّی ترتیبی عدد اے}
\end{align*}
\end{defn}

جے اسیں سیٹ نظریے دے ماہر \emph{وجوں} کم کردے ہُندے تاں شاید ترتیبی گھات دیاں کجھ خاصیتاں ہور ویکھنا چاہندے۔ پر ساڈے کول ہور بہت کجھ کہن نوں نہیں، سوائے ایس جانی پچھانی گل دے کہ ترتیبی گھات تبادلی نہیں۔ سو $\ordexpo{2}{\omega} =
\bigcup_{\delta < \omega}\ordexpo{2}{\delta} = \omega$، جدکہ
$\ordexpo{\omega}{2} = \omega \ordtimes \omega$۔ پر سانوں ایہ \emph{آس} وی نہیں رکھنی چاہیدی کہ گھات تبادلی ہووے گی، کیوں جے قدرتی عدداں وچ وی ایہ تبادلی نہیں۔ $\ordexpo{2}{3} = 8 < 9 = \ordexpo{3}{2}$۔

\begin{prob}
لامتناہی درجیاں اُتے استقرا ورت کے ثابت کرو کہ جے اسیں
$\ordexpo{\alpha}{\beta} = \ordtype{\text{finfun}(\alpha, \beta),
\sqsubset}$ تعریف کریے، تاں
\olref[sth][ord-arithmetic][expo]{ordexporecursion} دیاں بازگشتی مساواتاں ملدیاں نیں۔
\end{prob}
\endgroup
% /OLP-0579

\OLEndChapterHook
\endgroup
% /OLP-0574
% OLP-0580 content/set-theory/cardinals/cardinals.tex
\begingroup
\def\olfilename{cardinals}
\expandafter\def\csname ol@sectioncs\endcsname{section}

\olchapter{sth}{cardinals}{کارڈینل عدد}

% OLP-0581 content/set-theory/cardinals/cp.tex
\begingroup
\def\olfilename{cp}
\expandafter\def\csname ol@sectioncs\endcsname{section}

\olfileid{sth}{cardinals}{cp}
\olsection{کانتور دا اصول}

\olref[ordinals][vn]{sec} چیتے کرو۔ اوتھے اسیں حُسنِ ترتیب والے سیٹاں دی گل کردے ساں تے آکھیا سی کہ چنگا ہووے جے ساڈے کول ایہو جیہیاں شےواں ہون جیہڑیاں حُسنِ ترتیب والیاں ساختاں دی نمائندگی کرن۔ ایس گل نوں سامنے رکھ کے اسیں ترتیبی عدد پیش کیتے، تے فیر \olref[ordinals][ordtype]{ordtypesworklikeyouwant} وچ وکھایا کہ اوہ اینج ای کم کردے نیں جیویں اسیں چاہندے ساں، یعنی:
\[
	\ordtype{A, <} = \ordtype{B, \lessdot}
	\text{ اُدوں تے صرف اُدوں جدوں } \tuple{A, <} \isomorphic \tuple{B, \lessdot}.
\]
ہُن ہور پچھے \olref[sfr][siz][equ]{sec} چیتے کرو۔ اوتھے اسیں سادہ طریقے نال کم کردیاں سیٹ دی «وڈائی» دا خیال پیش کیتا سی۔ خاص طور اُتے اسیں آکھیا سی کہ دو سیٹ ہم تعداد نیں، $\cardeq{A}{B}$، اُدوں تے صرف اُدوں جدوں اک !!a{bijection} $f \colon A \to
B$ اے۔ ایہ خیال اپنے آپ وچ حُسنِ ترتیب دے خیال نالوں {سوکھا} اے۔ ساڈا کم صرف !!{bijection}s نال اے۔ ایہناں !!{bijection}s نے (ترتیب دے ہم ساخت نقشیاں وانگ) «کوئی ساخت وی سنبھال کے رکھنی» اے یا نہیں، ایہ گل ایتھے ضروری نہیں۔

ایہ ساریاں گلاں توں اک سِدھا خیال نکلدا اے۔ جیویں اسیں حُسنِ ترتیب والیاں ساختاں دا ماپ لین لئی کجھ شےواں، یعنی \emph{ترتیبی عدد}، پیش کیتے سن، اویویں اسیں وڈائی دا ماپ لین لئی کجھ شےواں، یعنی \emph{کارڈینل عدد}، پیش کر سکدے آں۔ ایس باب دا مقصد ایہو اے۔

ایہ دسن توں پہلاں کہ ایہ کارڈینل عدد کیہ ہون گے، سانوں اک اصول دَسنا چاہیدا اے جیہدے اُتے اوہ پورے اُترن۔ $\card{X}$ نوں سیٹ $X$ دی کارڈینیلٹی لئی ورتدے ہوئیاں، اسیں چاہواں گے کہ:
\[
	\card{A} = \card{B} \text{ اُدوں تے صرف اُدوں جدوں } \cardeq{A}{B}.
\]
اسیں اینوں \emph{کانتور دا} اصول آکھاں گے، کیوں جے شاید کانتور پہلا سی جیہدے ذہن وچ ایہ گل بڑی صاف سی۔ (ایہدا \emph{ہیوم دے} اصول نال تعلق \olref[hp]{sec} وچ ہور سمجھاواں گے۔) سو ساڈا مقصد $\card{X}$ دی تعریف ہر $X$ لئی ایہو جیہی دینا اے جیہدے توں کانتور دا اصول ملے۔
\endgroup
% /OLP-0581
% OLP-0582 content/set-theory/cardinals/cardsasords.tex
\begingroup
\def\olfilename{cardsasords}
\expandafter\def\csname ol@sectioncs\endcsname{section}

\olfileid{sth}{cardinals}{cardsasords}
\olsection{کارڈینل عدد ترتیبی عدداں وجوں}

اصل وچ کارڈینل عدداں دا ساڈا نظریہ ترتیبی عدداں دا نظریہ (بے جھجک) ورتے گا۔ یعنی اسیں کارڈینل عدداں دی تعریف کجھ خاص ترتیبی عدداں وجوں ای دیواں گے۔ خاص طور اُتے ایہ تعریف دیواں گے:

\begin{defn}\ollabel{defcardinalasordinal}
جے $A$ نوں حُسنِ ترتیب دتی جا سکدی اے، تاں $\card{A}$ سبھ توں چھوٹا ترتیبی عدد $\gamma$ اے جیہدے لئی $\cardeq{A}{\gamma}$۔ کسے وی ترتیبی عدد $\gamma$ لئی، اسیں آکھدے آں کہ $\gamma$ اک \emph{کارڈینل عدد} اے اُدوں تے صرف اُدوں جدوں $\gamma = \card{\gamma}$۔
\end{defn}

اسیں ہُن ای عبارت «$A$ نوں حُسنِ ترتیب دتی جا سکدی اے» ورتی اے۔ جیویں ریاضی وچ لگ بھگ ہمیشہ ہُندا اے، ایتھے «سکدی اے» والی گل موجودگی دے دعوے نوں ذرا ڈھلّے طریقے نال آکھدی اے۔ «$A$ نوں حُسنِ ترتیب دتی جا سکدی اے» دا بس مطلب اے «اک تعلق موجود اے جیہڑا~$A$ نوں حُسنِ ترتیب دیندا اے»۔

پر \olref{defcardinalasordinal} وچ اک اڑچن اے۔ اسیں چاہندے آں کہ \emph{ہر} سیٹ دی وڈائی ہووے، یعنی $\card{A}$ ہر~$A$ لئی موجود ہووے۔ پر ہُن دتی تعریف اک شرط نال شروع ہُندی اے، «\emph{جے} $A$ نوں حُسنِ ترتیب دتی جا سکدی اے\ldots»۔ جے کوئی ایہو جیہا سیٹ $A$ اے جیہنوں حُسنِ ترتیب نہیں دتی جا سکدی، تاں ساڈی تعریف نال شے~$\card{A}$ دی تعریف ای نہیں ہووے گی۔

سو \olref{defcardinalasordinal} ورتن لئی سانوں ضمانت چاہیدی اے کہ ہر سیٹ نوں حُسنِ ترتیب دتی جا سکدی اے۔ افسوس، ایہ ضمانت~$\ZF$ وچ نہیں ملدی۔ سو جے اسیں \olref{defcardinalasordinal} ورتنا اے، تاں اک نواں مسلّمہ پاؤن توں بناں چارہ نہیں، مثلاً:
\begin{axiom}[حُسنِ ترتیب]
ہر سیٹ نوں حُسنِ ترتیب دتی جا سکدی اے۔
\end{axiom}
حُسنِ ترتیب دا مسلّمہ منّن جوگا اے یا نہیں، ایہدی گل اسیں \olref[choice][]{chap} وچ کراں گے۔ پر اگے اسیں ایہنوں ورتاں گے۔ ایہ ورت کے ثابت کرنا بڑا سِدھا اے کہ کارڈینل عدد (\olref{defcardinalasordinal} دی تعریف مطابق) موجود نیں تے چنگیاں خاصیتاں رکھدے نیں:

\begin{lem}\ollabel{lem:CardinalsExist}
ہر سیٹ $A$ لئی:
\begin{enumerate}
	\item\ollabel{cardaexists} $\card{A}$ موجود اے تے اکّو اک اے۔
	\item\ollabel{cardaapprox}  $\cardeq{\card{A}}{A}$۔
	\item\ollabel{cardaidem}  $\card{A}$ کارڈینل عدد اے، یعنی
	$\card{A} = \card{\card{A}}$۔
\end{enumerate}
\end{lem}

\begin{proof}
$A$ مقرر کرو۔ حُسنِ ترتیب دے مسلّمے نال اک حُسنِ ترتیب $\tuple{A, R}$ موجود اے۔ \olref[ordinals][ordtype]{thmOrdinalRepresentation} مطابق $\tuple{A,
R}$ اکّو اک ترتیبی عدد $\beta$ نال ہم ساخت اے۔ سو $\cardeq{A}{\beta}$۔ لامتناہی درجیاں اُتے استقرا نال، اکّو اک سبھ توں چھوٹا ترتیبی عدد $\gamma$ اے جیہدے لئی $\cardeq{A}{\gamma}$۔ سو $\card{A}
= \gamma$، تے \olref{cardaexists} تے \olref{cardaapprox} ثابت ہو گئے۔ \olref{cardaidem} ثابت کرن لئی دھیان دیو کہ جے $\delta \in \gamma$، تاں $\cardless{\delta}{A}$، جیویں اسیں $\gamma$ چُنیا سی۔ ایس لئی $\cardless{\delta}{\gamma}$ وی، کیوں جے ہم تعدادیت اک ہم ارزی دا تعلق اے (\olref[sfr][siz][equ]{equinumerosityisequi})۔ سو $\gamma =
\card{\gamma}$۔
%So, for reductio, suppose that there is some ordinal $\delta \in \gamma$ such that $\cardeq{\gamma}{\delta}$. Then, $\cardeq{\cardeq{A}{\gamma}}{\delta}$ so that $\cardeq{A}{\delta}$ by \olref[sfr][set][equ]{equinumerosityisequi}, which contradicts the choice of $\gamma$ as the least ordinal such that $\cardeq{A}{\gamma}$. 
\end{proof}

اگلا نتیجہ کانتور دے اصول دی ضمانت دیندا اے تے کجھ ہور وی دسدا اے۔
(دھیان دیو کہ کارڈینل عدداں نوں اپنی ترتیب ترتیبی عدداں توں ملدی اے، یعنی $\cardfont{a} < \cardfont{b}$ اُدوں تے صرف اُدوں جدوں $\cardfont{a} \in \cardfont{b}$۔ ہِیٹھ ایہ لکھدیاں، جنہاں شےواں دا سانوں پتا اے کہ اوہ کارڈینل عدد نیں، اوہناں لئی اسیں فراکتور اکھر ورتاں گے۔ ایہ کافی رواجی اے۔ اک ہور عام طریقہ یونانی اکھر ورتنا اے، کیوں جے کارڈینل عدد ترتیبی عدد نیں، پر یونانی ورن مالا دے وچکارلے اکھر چُنو، جیویں $\kappa, \lambda$۔)
\begin{lem}\ollabel{lem:CardinalsBehaveRight}
کسے وی دو سیٹاں $A$ تے $B$ لئی:
\begin{align*}
	\cardeq{A}{B} &\text{ اُدوں تے صرف اُدوں جدوں } \card{A} = \card{B}\\
	\cardle{A}{B} &\text{ اُدوں تے صرف اُدوں جدوں } \card{A} \leq \card{B}\\
	\cardless{A}{B}&\text{ اُدوں تے صرف اُدوں جدوں } \card{A} < \card{B}
\end{align*}
\end{lem}

\begin{proof}
اسیں دوجے دعوے دی کھبّے توں سجّے ول سمت ثابت کراں گے (ہور صورتاں ایس ورگیاں نیں تے مشق لئی چھڈیاں نیں)۔ سو ہِیٹھلا خاکہ ویکھو:
\begin{center}
	\begin{tikzpicture}
	\node (nodea) {$A$};
	\node[right = 6em of nodea] (nodeb) {$B$};
	\node[below = 2em of nodea] (nodecarda) {$\card{A}$};
	\node[below = 2em of nodeb] (nodecardb) {$\card{B}$};
	\draw[->] (nodea)--(nodeb);
	\draw[<->] (nodea)--(nodecarda);
	\draw[<->] (nodeb)--(nodecardb);
	\draw[->, dashed] (nodecarda)--(nodecardb);
\end{tikzpicture}
\end{center}
دو سِریاں والے تیر !!{bijection}s دسّدے نیں، جیہناں دی موجودگی دی ضمانت \olref{lem:CardinalsExist} دیندا اے۔ $\cardle{A}{B}$ منّن نال اک !!a{injection} $A\to B$ اے۔ ہُن $\card{A}$ توں $A$، فیر $B$، تے فیر $\card{B}$ ول تیراں دے پچھے چل کے سانوں اک !!a{injection} $\card{A} \to \card{B}$ ملدی اے (ٹُٹّواں تیر)۔
\end{proof}\noindent اسیں \olref{lem:CardinalsBehaveRight} ورت کے شروڈر–برنسٹین دا قضیہ وی فیر ثابت کر سکدے آں۔ ایہ دعویٰ اے کہ جے $\cardle{A}{B}$ تے $\cardle{B}{A}$، تاں $\cardeq{A}{B}$۔ اسیں ایہنوں \olref[sfr][siz][sb]{thm:schroder-bernstein} وجوں بیان کیتا سی، پر پہلا ثبوت، کجھ محنت نال، \olref[sfr][infinite][card-sb]{sec} وچ دتا سی۔ ہُن ویکھو:

\begin{proof}[شروڈر–برنسٹین دا دوبارہ ثبوت]
جے $\cardle{A}{B}$ تے $\cardle{B}{A}$، تاں $\card{A} \leq \card{B}$ تے $\card{B} \leq \card{A}$، \olref{lem:CardinalsBehaveRight} مطابق۔ سو $\card{A} = \card{B}$ تے $\cardeq{A}{B}$، تِن صورتاں دے اصول تے \olref{lem:CardinalsBehaveRight} مطابق۔
\end{proof}
\noindent
بھاویں ایہ بڑا سوکھا ثبوت اے، ایہ لُکّویں طور اُتے بدل رکھن دے اصول (\olref[ordinals][ordtype]{thmOrdinalRepresentation} لئی) تے حُسنِ ترتیب دے مسلّمے (\olref{lem:CardinalsBehaveRight} دی ضمانت لئی) دوہاں اُتے ٹکدا اے۔ دوجے پاسے \olref[sfr][infinite][card-sb]{sec} دا ثبوت اپنے زور اُتے کِتے ودھ کھلوتا سی (اصل وچ ایہ~$\Zminus$ وچ دِتا جا سکدا اے)۔
\endgroup
% /OLP-0582
% OLP-0583 content/set-theory/cardinals/milestone.tex
\begingroup
\def\olfilename{milestone}
\expandafter\def\csname ol@sectioncs\endcsname{section}

\olfileid{sth}{cardinals}{zfc}
\olsection{\texorpdfstring{$\ZFC$ تک پہنچ}{ZFC تک پہنچ}}

حُسنِ ترتیب دا مسلّمہ پا کے اسیں نظریے دے آخری پڑاؤ تک پہنچ گئے آں۔ ہُن~$\ZFC$ لئی لوڑیندے سارے مسلّمے ساڈے کول نیں۔ تفصیل ایہ اے:

\begin{defn}
نظریے $\ZFC$ دے مسلّمے ایہ نیں۔ عنصراں نال تعیّن، اتحاد، جوڑیاں، ذیلی سیٹاں دے سیٹ، لامتناہیت، بنیاد، حُسنِ ترتیب، تے وکھ کرن تے بدل رکھن دے خاکیاں دیاں ساریاں صورتاں۔ دوجے لفظاں وچ، $\ZFC$ وچ~$\ZF$ دے نال حُسنِ ترتیب دا مسلّمہ شامل اے۔
\end{defn}

$\ZFC$ توں مراد \emph{زرمَیلو–فرینکل} دا سیٹ نظریہ اے جیہدے وچ \emph{چون} دا مسلّمہ شامل اے۔ ایہ گل ذرا اوپری لگ سکدی اے، کیوں جے اسیں جیہڑا مسلّمہ پایا اوہدا ناں «حُسنِ ترتیب» سی، «چون» نہیں۔ پر جدوں اسیں اگے {چون} دا مسلّمہ بیان کراں گے، پتا لگے گا کہ حُسنِ ترتیب دا مسلّمہ (~$\ZF$ دے اندر) چون دے مسلّمے دے برابر اے (\olref[choice][woproblem]{thmwochoice} ویکھو)۔ ایس لئی ایہناں وچوں جیہڑا مرضی اپنا «بنیادی» مسلّمہ منّ لئیے۔ تے «$\ZFC$» ناں علمی لکھتاں وچ پوری طرح رواجی اے۔
\endgroup
% /OLP-0583
% OLP-0584 content/set-theory/cardinals/classing.tex
\begingroup
\def\olfilename{classing}
\expandafter\def\csname ol@sectioncs\endcsname{section}

\olfileid{sth}{cardinals}{classing}
\olsection{محدود، \usetoken{S}{enumerable}، \usetoken{S}{nonenumerable}}

کارڈینل عدداں دی جان پہچان پچھوں اوہناں دیاں وکھ وکھ قسماں دی تھوڑی گل کرنی چنگی اے، خاص کر کے محدود، !!{enumerable} تے !!{nonenumerable} کارڈینل عدداں دی۔

ساڈے پہلے دو نتیجیاں توں نکلدا اے کہ محدود کارڈینل عدد این ٹھیک محدود ترتیبی عدد نیں۔ ایہناں نوں اسیں \olref[z][infinity-again]{defnomega} وچ اپنے \emph{قدرتی عدد} وجوں تعریف کیتا سی۔

\begin{prop}\ollabel{finitecardisoequal}
$n, m \in \omega$ ہون۔ تاں $n = m$ اُدوں تے صرف اُدوں جدوں $\cardeq{n}{m}$۔
\end{prop}

\begin{proof}
\emph{کھبّے توں سجّے ول} صاف اے۔ \emph{سجّے توں کھبّے ول} ثابت کرن لئی منّ لو $\cardeq{n}{m}$، بھاویں $n \neq m$۔ تِن صورتاں دے اصول مطابق یا $n
\in m$ یا $m \in n$۔ عمومیت گنوائے بناں منّ لو $n \in m$۔ تاں $n \subsetneq m$ تے اک !!a{bijection} $f \colon m \to n$ اے۔ سو $m$ ڈیڈیکنڈ لامتناہی اے، جیہڑی گل \olref[z][infinity-again]{naturalnumbersarentinfinite} نال ٹکراؤندی اے۔
\end{proof}

\begin{cor}\ollabel{naturalsarecardinals}
جے $n \in \omega$، تاں $n$ اک کارڈینل عدد اے۔
\end{cor}

\begin{proof}
سِدھا ای نکلدا اے۔
\end{proof}
\noindent
ایہ وی نکلدا اے کہ کسے کارڈینل عدد نوں «محدود» یا «لامتناہی» آکھن دے کئی واجب معنی آپس وچ ملدے نیں۔
\begin{thm}\ollabel{generalinfinitycharacter}کسے وی سیٹ $A$ لئی ہِیٹھلیاں گلاں برابر نیں:
\begin{enumerate}
	\item\ollabel{card:notinomega} $\card{A} \notin \omega$، یعنی
	$A$ قدرتی عدد نہیں۔
	\item\ollabel{card:omegaplus} $\omega \leq \card{A}$۔
	\item\ollabel{card:infinite} $A$ ڈیڈیکنڈ لامتناہی اے۔
\end{enumerate}
\end{thm}

\begin{proof}
\olref[ord-arithmetic][using-addition]{ordinfinitycharacter}،
\olref[cardinals][cardsasords]{lem:CardinalsBehaveRight} تے
\olref{naturalsarecardinals} توں۔
\end{proof}

ایس توں سانوں اوہناں کجھ خیالاں دی ہِیٹھلی \emph{تعریف} دین دا جواز ملدا اے، جیہڑے \olref[sfr][][]{part} وچ اسیں ذرا غیر رسمی طور اُتے ورتے سن۔

\begin{defn}\ollabel{defnfinite}
اسیں آکھدے آں کہ $A$ \emph{محدود} اے اُدوں تے صرف اُدوں جدوں $\card{A}$ قدرتی عدد اے، یعنی $\card{A} \in \omega$۔ ہور صورت وچ اسیں آکھدے آں کہ $A$ \emph{لامتناہی} اے۔
\end{defn}
\noindent
پر دھیان دیو کہ ایہ تعریف $\ZFC$ دے پس منظر وچ دتی اے۔ آخر سانوں ایہ ضمانت دین لئی حُسنِ ترتیب دا مسلّمہ چاہیدا سی کہ ہر سیٹ دی کارڈینیلٹی اے۔ تے حُسنِ ترتیب دے مسلّمے توں بناں واقعی کوئی ایہو جیہا سیٹ ہو سکدا اے جیہڑا نہ محدود ہووے نہ ڈیڈیکنڈ لامتناہی۔ ایس طرح دے سوال ول اسیں \olref[choice][]{chap} وچ مُڑاں گے۔ ہن لئی اسیں حُسنِ ترتیب دا مسلّمہ ورتدے رہندے آں۔

ہُن اسیں محدود کارڈینل عدداں توں لامتناہی کارڈینل عدداں ول آؤندے آں۔ دو بنیادی گلاں ایہ نیں۔

\begin{cor}\ollabel{omegaisacardinal}
$\omega$ سبھ توں چھوٹا لامتناہی کارڈینل عدد اے۔
\end{cor}

\begin{proof}
$\omega$ کارڈینل عدد اے، کیوں جے $\omega$ ڈیڈیکنڈ لامتناہی اے تے جے $\cardeq{\omega}{n}$ کسے $n \in \omega$ لئی، تاں $n$ ڈیڈیکنڈ لامتناہی ہووے گا، جیہڑی گل \olref[z][infinity-again]{naturalnumbersarentinfinite} نال ٹکراؤندی اے۔ ہُن $\omega$ تعریف مطابق سبھ توں چھوٹا لامتناہی کارڈینل عدد اے۔
\end{proof}

\begin{cor}
ہر لامتناہی کارڈینل عدد حدّی ترتیبی عدد اے۔
\end{cor}

\begin{proof}
$\alpha$ اک لامتناہی جانشین ترتیبی عدد ہووے، سو $\alpha = \beta
\ordplus 1$ کسے $\beta$ لئی۔ \olref{finitecardisoequal} مطابق $\beta$ وی لامتناہی اے، سو $\cardeq{\beta}{\beta \ordplus 1}$، \olref[ord-arithmetic][using-addition]{ordinfinitycharacter} مطابق۔ ہُن $\card{\beta} = \card{\beta\ordplus 1} = \card{\alpha}$، \olref[cardinals][cardsasords]{lem:CardinalsBehaveRight} مطابق، سو $\alpha \neq \card{\alpha}$۔
\end{proof}

پہلاں ای \olref[sfr][siz][enm-alt]{defn:enumerable} وچ اسیں دَسیا سی کہ اسیں !!{enumerable} تے !!{nonenumerable} لامتناہی سیٹاں وچ فرق کر سکدے آں۔ ایس تعریف توں سہج ای ہِیٹھلی گل نکلدی اے۔

\begin{prop}
$A$ !!{enumerable} اے اُدوں تے صرف اُدوں جدوں $\card{A} \leq \omega$، تے $A$ !!{nonenumerable} اے اُدوں تے صرف اُدوں جدوں $\omega < \card{A}$۔
\end{prop}

\begin{proof}
تِن صورتاں دے اصول مطابق دوہاں دعویاں دا مطلب اکّو اے۔ سو ایہ ثابت کرنا کافی اے کہ $A$ !!{enumerable} اے اُدوں تے صرف اُدوں جدوں $\card{A} \leq \omega$۔ \emph{سجّے توں کھبّے ول}، جے $\card{A} \leq \omega$، تاں $\cardle{A}{\omega}$، \olref[cardinals][cardsasords]{lem:CardinalsBehaveRight} تے \olref{omegaisacardinal} مطابق۔ \emph{کھبّے توں سجّے ول}، منّ لو $A$ !!{enumerable} اے۔ تاں \olref[sfr][siz][enm-alt]{defn:enumerable} مطابق تِن صورتاں ممکن نیں:
\begin{enumerate}
	\item جے $A = \emptyset$، تاں $\card{A} = 0 \in \omega$، \olref{naturalsarecardinals} تے \olref[cardinals][cardsasords]{lem:CardinalsBehaveRight} مطابق۔
	\item جے $\cardeq{n}{A}$، تاں $\card{A} = n \in \omega$، \olref{naturalsarecardinals} تے \olref[cardinals][cardsasords]{lem:CardinalsBehaveRight} مطابق۔
	\item جے $\cardeq{\omega}{A}$، تاں $\card{A} = \omega$، \olref{omegaisacardinal} مطابق۔
\end{enumerate}
سو ساریاں صورتاں وچ $\card{A} \leq \omega$۔
\end{proof}
\noindent
واقعی $\omega$ دی اک خاص تھاں اے۔ گنتی وار فہرست بنن جوگ ترتیبی عدد کئی نیں، پر:

\begin{cor}
$\omega$ اکّو اک !!{enumerable} لامتناہی کارڈینل عدد اے۔
\end{cor}

\begin{proof}
$\cardfont{a}$ اک !!a{enumerable} لامتناہی کارڈینل عدد ہووے۔ کیوں جے $\cardfont{a}$ لامتناہی اے، $\omega \leq \cardfont{a}$۔ کیوں جے $\cardfont{a}$ اک !!a{enumerable} کارڈینل عدد اے، $\cardfont{a} =
\card{\cardfont{a}} \leq \omega$۔ سو $\cardfont{a} = \omega$، تِن صورتاں دے اصول مطابق۔ \end{proof}

بے شک کارڈینل عدد لامتناہی گنتی وچ نیں۔ سو اسیں پُچھ سکدے آں، \emph{کارڈینل عدد کِنّے نیں؟} ہِیٹھلے نتیجے وکھاندے نیں کہ شاید سانوں ایس سوال بارے فیر سوچنا چاہیدا اے۔

\begin{prop}\ollabel{unioncardinalscardinal}
جے $X$ دا ہر رکن اک کارڈینل عدد اے، تاں $\bigcup X$ اک کارڈینل عدد اے۔
\end{prop}

\begin{proof}
ایہ جانچنا سوکھا اے کہ $\bigcup X$ اک ترتیبی عدد اے۔ $\alpha \in
\bigcup X$ اک ترتیبی عدد ہووے۔ تاں $\alpha \in \cardfont{b} \in X$ کسے کارڈینل عدد $\cardfont{b}$ لئی۔ کیوں جے $\cardfont{b}$ اک کارڈینل عدد اے، $\cardless{\alpha}{\cardfont{b}}$۔ کیوں جے $\cardfont{b} \subseteq
\bigcup X$، سانوں $\cardle{\cardfont{b}}{\bigcup X}$ ملدا اے، تے ایس لئی $\cardneq{\alpha}{\bigcup X}$۔ ایس گل نوں عام کریے تاں $\bigcup X$ اک کارڈینل عدد اے۔
\end{proof}

\begin{thm}\ollabel{lem:NoLargestCardinal}
کوئی سبھ توں وڈا کارڈینل عدد نہیں۔
\end{thm}

\begin{proof}
کسے وی کارڈینل عدد $\cardfont{a}$ لئی، کانتور دے قضیے (\olref[sfr][siz][car]{thm:cantor}) تے \olref[cardinals][cardsasords]{lem:CardinalsExist} توں لازم آؤندا اے کہ $\cardfont{a} < \card{\Pow{\cardfont{a}}}$۔
\end{proof}

\begin{thm}
سارے کارڈینل عدداں دا سیٹ موجود نہیں۔
\end{thm}

\begin{proof}
ٹکراؤ کڈھن لئی منّ لو $C = \Setabs{\cardfont{a}}{\cardfont{a} \text{
اک کارڈینل عدد اے}}$۔ ہُن $\bigcup C$ اک کارڈینل عدد اے، \olref{unioncardinalscardinal} مطابق۔ سو \olref{lem:NoLargestCardinal} مطابق اک کارڈینل عدد $\cardfont{b} > \bigcup C$ اے۔ تعریف مطابق $\cardfont{b} \in C$، سو $\cardfont{b} \subseteq \bigcup{C}$، ایس لئی $\cardfont{b} \leq \bigcup C$۔ ایہ ٹکراؤ اے۔
\end{proof}

ایہدا رسل دے پیراڈوکس تے بورالی فورتی دے پیراڈوکس، دوہاں نال مقابلہ کرو۔
\endgroup
% /OLP-0584
% OLP-0585 content/set-theory/cardinals/hp.tex
\begingroup
\def\olfilename{hp}
\expandafter\def\csname ol@sectioncs\endcsname{section}

\olfileid{sth}{cardinals}{hp}
\olsection{ضمیمہ: ہیوم دا اصول}

\olref[cp]{sec} وچ اسیں کانتور دا اصول دَسیا سی۔ اوہ ایہ سی:
\begin{align*}
	\card{A} = \card{B} & \text{ اُدوں تے صرف اُدوں جدوں } A \approx B.
\intertext{ایہ اوہدی طرح دا ای اے جیہنوں ہُن \emph{ہیوم دا اصول} آکھدے نیں۔ اوہ ایہ آکھدا اے:}
	\fregenum{x} {F(x)} = \fregenum{x}{G(x)} & \text{ اُدوں تے صرف اُدوں جدوں } F \sim G
\end{align*}
جتھے «$F \sim G$» دا مختصر مطلب ایہ اے کہ $F$ شےواں این اونیاں ای نیں جِنّیاں $G$ شےواں، یعنی $F$ شےواں نوں $G$ شےواں نال اک دوطرفی نقشہ بندی وچ رکھ سکدے آں، یعنی:
\begin{align*}
	\exists R(&\forall v\forall y(Rvy \lif (Fv \land Gy)) \land {}\\
		&\forall v(Fv \lif \lexists![y][Rvy]) \land {}\\
		&\forall y(Gy \lif \lexists![v][Rvy]))
\end{align*}
پر ہیوم دے اصول تے کانتور دے اصول وچ نوع دا فرق اے۔ کانتور دے اصول دے بیان وچ متغیر «$A$» تے «$B$» پہلے درجے دیاں ٹرماں نیں جیہڑیاں \emph{سیٹاں} لئی کھڑیاں نیں۔ ہیوم دے اصول دے بیان وچ «$F$»، «$G$» تے «$R$» پہلے درجے دیاں ٹرماں \emph{نہیں}۔ اوہ \emph{محمول دی تھاں} نیں۔ (شاید اوہ \emph{خاصیتاں} لئی کھڑے نیں۔) سو ہیوم دے اصول نوں عام بولی وچ ایویں آکھ سکدے آں۔ $F$ شےواں دی گنتی $G$ شےواں دی گنتی اے اُدوں تے صرف اُدوں جدوں $F$ شےواں تے~$G$ شےواں وچ دوطرفی نقشہ بندی اے۔ ایس نوں \emph{ہیوم دا اصول} آکھدے نیں، کیوں جے ہیوم نے اک واری ایہ لکھیا:
\begin{quote}
  جدوں دو عدد ایس طرح جوڑے جان کہ اک دے ہر اک جُز دے سامنے دوجے دا اک جُز ہووے، تاں اسیں اوہناں نوں برابر آکھدے آں۔
  \citep[Pt.III Bk.1 \S1]{Hume1740}
\end{quote}
تے اج دی ریاضیاتی منطق وچ ہیوم دے اصول نوں اہم تھاں \citet[\S63]{Frege1884} نے دِوائی۔ اوہنے ہیوم دی ایہ عبارت نقل کیتی تے فیر (عملی طور اُتے) اوہدا خاکہ دتا جیہنوں اسیں ہیوم دا اصول آکھیا اے۔

ہیوم دے اصول تے پنجویں بنیادی قانون دی ساخت وچ مشابہت ول دھیان دیو۔ اسیں اوہ قانون \olref[story][blv]{sec} وچ ایویں لکھیا سی:
\[
	\fregeext{x}{F(x)} = \fregeext{x}{G(x)} \text{ اُدوں تے صرف اُدوں جدوں } \lforall[x][(F(x) \liff G(x))].
\]
ایتھے کجھ وضاحت تے مقابلہ کم آ سکدے نیں۔

ہیوم دے اصول یا پنجویں بنیادی قانون ورگے اصول نوں لین دے دو طریقے نیں، \emph{اسنادی طور اُتے} یا \emph{غیر اسنادی طور اُتے} (\olref[story][predicative]{sec} چیتے کرو)۔ پنجویں بنیادی قانون دی غیر اسنادی پڑھت وچ، ہر~$F$ لئی، شے $\fregeext{x}{F(x)}$ اوہناں مسوّراں دے دائرے وچ آؤندی اے جیہڑے اسیں پنجواں بنیادی قانون آپ بیان کرن لئی ورتے۔ ایویں ای ہیوم دے اصول دی غیر اسنادی پڑھت وچ، ہر~$F$ لئی، شے $\fregenum{x}{F(x)}$ اوہناں مسوّراں دے دائرے وچ آؤندی اے جیہڑے اسیں ہیوم دا اصول بیان کرن لئی ورتے۔ دوجے پاسے، \emph{اسنادی} سمجھ مطابق، شےواں $\fregeext{x}{F(x)}$ تے~$\fregenum{x}{F(x)}$ کسے \emph{وکھرے} دائرے دیاں شےواں ہون گیاں۔

ہُن جے اسیں پنجواں بنیادی قانون غیر اسنادی طور اُتے پڑھیے، تاں سادہ خاصیت توں سیٹ بناؤن دے اصول راہیں ٹکراؤ نکلدا اے (تفصیل لئی \olref[story][blv]{sec} ویکھو)۔ سادہ خاصیت توں سیٹ بناؤن دے اصول وانگ، اینوں \emph{اسنادی طور اُتے} پڑھ کے ٹکراؤ توں بچایا جا سکدا اے۔ پر شاید ایہ اوہ سارا کم نہیں کرے گا جیہڑا اسیں چاہندے ساں۔

پر ہیوم دا اصول غیر اسنادی طور اُتے بے ٹکراؤ پڑھیا جا \emph{سکدا} اے۔ تے ایس پڑھت وچ ایہ خاصا طاقتور اے۔

مثال لئی محمول «$x \neq x$» ویکھو، جیہدے اُتے صاف اے کہ کوئی شے پوری نہیں اُتر دی۔ ہُن ہیوم دے اصول توں شے $\# x( x\neq
x)$ ملدی اے۔ اسیں اینوں عدد~$0$ منّ سکدے آں۔ ہُن \emph{غیر اسنادی} سمجھ مطابق، تے \emph{صرف} غیر اسنادی سمجھ مطابق، ایہ شے $0$ ساڈے اصل مسوّراں دے دائرے وچ آؤندی اے۔ سو محمول «$x = 0$» نال ہیوم دا اصول لانا معنی رکھدا اے، تے ایویں سانوں شے $\#x (x =
0)$ ملدی اے۔ اسیں اینوں عدد~$1$ منّ سکدے آں۔ نالے ہیوم دے اصول توں لازم آؤندا اے کہ $0 \neq 1$، کیوں جے اپنے آپ نال اکّو شے نہ ہون والیاں شےواں توں $0$ نال اکّو شے ہون والیاں شےواں ول دوطرفی نقشہ بندی نہیں ہو سکدی (پہلی قسم وچ کوئی شے نہیں، دوجی وچ اک اے)۔ ہُن فیر غیر اسنادی طور اُتے کم کردیاں، $1$~ساڈے اصل مسوّراں دے دائرے وچ آؤندا اے۔ سو محمول «$(x = 0 \lor x = 1)$» نال ہیوم دا اصول لانا معنی رکھدا اے، تے ایویں سانوں شے $\#x(x = 0 \lor
x = 1)$ ملدی اے۔ اسیں اینوں عدد~$2$ منّ سکدے آں، تے وکھا سکدے آں کہ $0\neq 2$ تے $1 \neq 2$، تے ایس طرح اگے۔

سو مختصر گل ایہ اے کہ غیر اسنادی طور اُتے لئیے تاں ہیوم دے اصول توں لازم آؤندا اے کہ شےواں \emph{لامتناہی گنتی وچ} نیں۔ ایس گل نے \emph{نویں فریگے والے منطقیت پسنداں} نوں حساب دی بنیاد وجوں ہیوم دا اصول لین ول مائل کیتا اے۔

پر فریگے نے \emph{آپ} ہیوم دا اصول اپنے حساب دی بنیاد نہیں منّیا۔ اوہنے ہیوم دا اصول اک صاف تعریف توں ثابت کیتا۔ $\fregenum{x}{F(x)}$ دی تعریف تصور $F \sim \Phi$ دے دائرے وجوں دتی اے۔ اج دی لکھت وچ اسیں ایہنوں $\fregenum{x}{F(x)} = \Setabs{G}{F \sim G}$ وجوں لکھن دی کوشش کر سکدے آں، پر ایہ سانوں سادہ خاصیت توں سیٹ بناؤن دے اصول دیاں اوکڑاں ول موڑ لے جائے گا۔
\endgroup
% /OLP-0585

\OLEndChapterHook
\endgroup
% /OLP-0580
% OLP-0586 content/set-theory/card-arithmetic/card-arithmetic.tex
\begingroup
\def\olfilename{card-arithmetic}
\expandafter\def\csname ol@sectioncs\endcsname{section}

\olchapter{sth}{card-arithmetic}{کارڈینل عدداں دا حساب}

% OLP-0587 content/set-theory/card-arithmetic/opps.tex
\begingroup
\def\olfilename{opps}
\expandafter\def\csname ol@sectioncs\endcsname{section}

\olfileid{sth}{card-arithmetic}{opps}
\olsection{بنیادی عملاں دی تعریف}

کارڈینل عدداں دا حساب تعریف کرنا ترتیبی عدداں دے حساب نالوں سوکھا اے، کیوں جے ایتھے ترتیب دا حساب رکھن دی لوڑ نہیں۔ اسیں جوڑ، ضرب تے گھات دی تعریف اکّو واری دیواں گے۔

\begin{defn}
جدوں $\cardfont{a}$ تے $\cardfont{b}$ کارڈینل عدد نیں:
\begin{align*}
	\cardfont{a} \cardplus \cardfont{b} &\defis
	\card{\cardfont{a} \disjointsum \cardfont{b}}\\
	\cardfont{a} \cardtimes \cardfont{b} &\defis
	\card{\cardfont{a} \times \cardfont{b}}\\
	\cardexpo{\cardfont{a}}{\cardfont{b}} &\defis
	\card{\funfromto{\cardfont{b}}{\cardfont{a}}}
\end{align*}
جتھے $\funfromto{X}{Y} = \Setabs{f}{f \text{ اک فنکشن اے } X \to Y}$۔
(ایہ وکھانا سوکھا اے کہ $\funfromto{X}{Y}$ کسے وی دو سیٹاں $X$ تے $Y$ لئی موجود اے۔ اسیں ایہ مشق لئی چھڈدے آں۔)
\end{defn}

\begin{prob}
$\Zminus$ وچ ثابت کرو کہ $\funfromto{X}{Y}$ کسے وی دو سیٹاں $X$ تے~$Y$ لئی موجود اے۔ $\ZF$ وچ کم کردیاں، $\setrank{\funfromto{X}{Y}}$ دا حساب $\setrank{X}$ تے $\setrank{Y}$ توں لاؤ، جیویں \olref[sth][ord-arithmetic][using-addition]{rankcomputation} وچ کیتا اے۔
\end{prob}

ایہ تعریف ذرا سمجھا دیئیے۔ جوڑ لئی وکھ وکھ حصیاں والے جوڑ $\disjointsum$ دا خیال ورتیا اے، جیہدی تعریف \olref[ord-arithmetic][add]{defdissum} وچ دتی اے۔ محدود صورتاں وچ ایس تعریف توں ٹھیک حاصل نکلدا اے، ایہ ویکھنا سوکھا اے۔ ضرب بارے، \olref[sfr][set][pai]{cardnmprod} دسدا اے کہ جے $A$ دے $n$ رکن نیں تے $B$ دے $m$ رکن نیں، تاں $A \times B$ دے $n \cdot m$ رکن نیں۔ سو ساڈی تعریف ایس خیال نوں لامتناہی درجیاں والی ضرب تک ودھا دیندی اے۔ گھات وی ایس طرح اے۔ اسیں بس محدود صورت دا خیال لامتناہی درجیاں تک ودھاؤندے آں۔ سچ مُچ کجھ گلاں وچ لامتناہی درجیاں والے کارڈینل عدداں دا حساب «عام» حساب نال کِتے ودھ ملدا اے، لامتناہی درجیاں والے ترتیبی عدداں دے حساب نالوں۔

\begin{prop}\ollabel{cardplustimescommute}
$\cardplus$ تے $\cardtimes$ تبادلی نیں تے شرکتیت رکھدے نیں۔
\end{prop}

\begin{proof}
تبادلیت لئی \olref[cardinals][cardsasords]{lem:CardinalsBehaveRight} مطابق اینا ویکھنا کافی اے کہ $\cardeq{(\cardfont{a} \disjointsum
\cardfont{b})}{(\cardfont{b} \disjointsum \cardfont{a})}$ تے
$\cardeq{(\cardfont{a} \times \cardfont{b})}{(\cardfont{b} \times
\cardfont{a})}$۔ شرکتیت اسیں مشق لئی چھڈدے آں۔
\end{proof}

\begin{prob}
ثابت کرو کہ $\cardplus$ تے $\cardtimes$ شرکتیت رکھدے نیں۔
\end{prob}

\begin{prop}
$A$ لامتناہی اے اُدوں تے صرف اُدوں جدوں $\card{A} \cardplus 1 = 1 \cardplus \card{A} = \card{A}$۔
\end{prop}

\begin{proof}
جیویں \olref[cardinals][classing]{generalinfinitycharacter} وچ کیتا اے، \olref[ord-arithmetic][using-addition]{ordinfinitycharacter} تے \olref[cardinals][cardsasords]{lem:CardinalsBehaveRight} توں۔
\end{proof}

ایس توں سمجھ آؤندی اے کہ ترتیبی تے کارڈینل جوڑ یا ضرب لئی وکھرے نشان کیوں ورتنے پَیندے نیں۔ ایہ عمل سچ مُچ \emph{وکھرے} نیں۔ اگلے دو نتیجے وکھاندے نیں کہ ترتیبی گھات تے کارڈینل گھات وی وکھرے عمل نیں۔ (چیتے کرو کہ \olref[z][infinity-again]{defnomega} توں لازم آؤندا اے کہ $2 = \{0,
1\}$۔)

\begin{lem}\ollabel{lem:SizePowerset2Exp}
$\card{\Pow{A}} = \cardexpo{2}{\card{A}}$، کسے وی $A$ لئی۔
\end{lem}

\begin{proof}
ہر ذیلی سیٹ $B \subseteq A$ لئی $\chi_B \in \funfromto{A}{2}$ ایویں دِتا ہووے:
\begin{align*}
	\chi_{B}(x) &\defis
	\begin{cases}
		1 & \text{جے }x\in B\\
		0 & \text{ہور صورت وچ۔}
	\end{cases}
\end{align*}
ہُن $f(B) = \chi_B$ ہووے۔ ایس نال اک !!a{bijection} $f \colon \Pow{A}
\to \funfromto{A}{2}$ تعریف ہُندی اے۔ سو $\cardeq{\Pow{A}}{\funfromto{A}{2}}$۔
اِس توں $\cardeq{\Pow{A}}{\funfromto{\card{A}}{2}}$، ایس لئی
$\card{\Pow{A}} = \card{\funfromto{\card{A}}{2}} =
2^{\card{A}}$۔
\end{proof}

ایس نِکّے جیہے ثبوت وچ بنیادی طور اُتے \olref[sfr][siz][red-alt]{sec} دی گل وی آ جاندی اے۔ اوتھے اسیں وکھایا سی کہ $\Pow{\omega}$ دا گنتی وار فہرست نہ بن سکنا، لامتناہی دو ہندسی لڑیواں دے سیٹ $\Bin^\omega$ دا گنتی وار فہرست نہ بن سکنا وکھا کے کیویں ثابت ہو سکدا اے۔ اصل وچ $\Bin^{\omega}$ بس $\funfromto{\omega}{2}$ اے۔ تے پچھلے ثبوت نے وکھایا کہ \olref[sfr][siz][red-alt]{sec} والی دلیل کسے وی سیٹ~$A$ نوں~$\omega$ دی تھاں رکھ کے چل پَئے گی۔ ایس نتیجے توں کارڈینل گھات بارے اک گل جھٹ نکلدی اے۔

\begin{cor}\ollabel{cantorcor}
$\cardfont{a} < \cardexpo{2}{\cardfont{a}}$ کسے وی کارڈینل عدد~$\cardfont{a}$ لئی۔
\end{cor}

\begin{proof}
کانتور دے قضیے (\olref[sfr][siz][car]{thm:cantor}) تے \olref{lem:SizePowerset2Exp} توں۔
\end{proof}
\noindent
سو $\omega < \cardexpo{2}{\omega}$۔ پر دھیان دیو کہ ایہ \emph{کارڈینل} گھات دا نتیجہ اے۔ اینوں \emph{ترتیبی} گھات نال ملا کے ویکھو، کیوں جے اوس صورت وچ $\omega =
\ordexpo{2}{\omega}$ (\olref[ord-arithmetic][expo]{sec} ویکھو)۔

کارڈینل گھات دی گل کردیاں اسیں ذرا ہور ٹھیک طرح دس سکدے آں کہ $\Real$ کیہڑے «معنی وچ» !!{nonenumerable} اے۔

\begin{thm}\ollabel{continuumis2aleph0}
$\card{\Real} = \cardexpo{2}{\omega}$
\end{thm}

\begin{proof}[ثبوت دا خاکہ]
ایہ ثابت کرن دے کئی طریقے نیں۔ سبھ توں سِدھا طریقہ ایہ اے کہ $\cardle{\Pow{\omega}}{\Real}$ تے $\cardle{\Real}{\Pow{\omega}}$ وکھائیے، فیر شروڈر–برنسٹین توں کڈھیے کہ $\cardeq{\Real}{\Pow{\omega}}$، تے \olref[card-arithmetic][opps]{lem:SizePowerset2Exp} توں کڈھیے کہ $\card{\Real} = \cardexpo{2}{\omega}$۔ انجیکشن $f \colon
\Pow{\omega} \to \Real$ تے $g \colon \Real \to \Pow{\omega}$ دی تعریف اسیں اک (سمجھ ودھاؤن والی) مشق لئی چھڈدے آں۔
\end{proof}

\begin{prob}
\olref[sth][card-arithmetic][opps]{continuumis2aleph0} دا ثبوت پورا کرو۔ وکھاؤ کہ $\cardle{\Pow{\omega}}{\Real}$ تے
$\cardle{\Real}{\Pow{\omega}}$۔
\end{prob}
\endgroup
% /OLP-0587
% OLP-0588 content/set-theory/card-arithmetic/simp.tex
\begingroup
\def\olfilename{simp}
\expandafter\def\csname ol@sectioncs\endcsname{section}
\olfileid{sth}{card-arithmetic}{simp}

\olsection{جوڑ تے ضرب سادے کرنے}

پتا لگدا اے کہ لامتناہی درجیاں والے کارڈینل عدداں دا جوڑ تے ضرب \emph{بے حد} سوکھے نیں۔ ایہ ایس گل توں نکلدا اے کہ کارڈینل عدد (کجھ خاص) ترتیبی عدد نیں، ایس لئی اوہ حُسنِ ترتیب والے نیں، تے اوہناں نال اک خاص طرح دا کم ہو سکدا اے۔ پر ایہ وکھانا اینا سوکھا \emph{نہیں}۔ شروع کرن لئی سانوں اک چترائی والی تعریف چاہیدی اے۔

\begin{defn}
اسیں ترتیبی عدداں دیاں جوڑیاں اُتے \emph{معیاری ترتیب} $\canonord$ ایویں تعریف کردے آں۔ $\tuple{\alpha_1, \alpha_2} \canonord
\tuple{\beta_1, \beta_2}$ اُدوں تے صرف اُدوں جدوں ہِیٹھلیاں وچوں اک گل ہووے:
\begin{enumerate}
	\item $\max(\alpha_1, \alpha_2) < \max(\beta_1, \beta_2)$، یا
	\item $\max(\alpha_1, \alpha_2) = \max(\beta_1, \beta_2)$ تے
	$\alpha_1 < \beta_1$، یا
	\item $\max(\alpha_1, \alpha_2) = \max(\beta_1, \beta_2)$ تے
	$\alpha_1 = \beta_1$ تے $\alpha_2 < \beta_2$
\end{enumerate}
\end{defn}

\begin{lem}
$\tuple{\alpha \times \alpha, \canonord}$ حُسنِ ترتیب والی ساخت اے، کسے وی ترتیبی عدد $\alpha$ لئی۔
\end{lem}

\begin{proof}
صاف اے کہ $\canonord$، $\alpha \times \alpha$ اُتے جُڑیا ہویا اے۔ کیوں جے منّ لو کہ نہ $\tuple{\alpha_1, \alpha_2}$ تے نہ $\tuple{\beta_1,
\beta_2}$، $\canonord$ مطابق دوجے توں چھوٹا اے۔ تاں $\max(\alpha_1,
\alpha_2) = \max(\beta_1, \beta_2)$ تے $\alpha_1 = \beta_1$ تے
$\alpha_2 = \beta_2$۔ سو $\tuple{\alpha_1, \alpha_2} =
\tuple{\beta_1,  \beta_2}$۔

حُسنِ ترتیب وکھاؤن لئی منّ لو $X \subseteq \alpha\times\alpha$ غیر خالی اے۔ کیوں جے $\alpha$ اک ترتیبی عدد اے، کوئی $\delta$، $\Setabs{\max(\gamma_1, \gamma_2)}{\tuple{\gamma_1,
\gamma_2} \in X}$ دا سبھ توں چھوٹا رکن اے۔ ہُن
$\Setabs{\tuple{\gamma_1,\gamma_2} \in X}{\max(\gamma_1, \gamma_2) =
\delta}$ وچوں بس اوہ جوڑیاں رکھو جیہناں دا پہلا محدد سبھ توں چھوٹا اے۔ ایہناں وچوں جیہڑی جوڑی دا دوجا محدد سبھ توں چھوٹا اے، اوہ $\canonord$ مطابق $X$ دا سبھ توں چھوٹا رکن اے۔
\end{proof}
\noindent
ہُن اک نِکّی جیہی، سوکھی گل ویکھو۔

\begin{prop}\ollabel{simplecardproduct}
جے $\cardeq{\alpha}{\beta}$، تاں $\cardeq{\alpha \times \alpha}{\beta
\times \beta}$۔
\end{prop}

\begin{proof}
بس $f \colon \alpha \to \beta$ توں نقشہ $\tuple{\gamma_1,
\gamma_2} \mapsto \tuple{f(\gamma_1), f(\gamma_2)}$ بناؤ۔
\end{proof}
\noindent
ہُن اسیں ایہ سارا کجھ اک ضروری معاون قضیہ ثابت کرن لئی ورتاں گے:
\begin{lem}\ollabel{alphatimesalpha}
$\cardeq{\alpha}{\alpha \times \alpha}$، کسے وی لامتناہی ترتیبی عدد $\alpha$ لئی۔
\end{lem}

\begin{proof}
ٹکراؤ کڈھن لئی $\alpha$ نوں سبھ توں چھوٹا لامتناہی ترتیبی عدد منّ لو جیہدے لئی ایہ گل غلط اے۔ \olref[sfr][siz][zigzag]{natsquaredenumerable} وکھاندا اے کہ $\cardeq{\omega}{\omega\times\omega}$، سو $\omega \in \alpha$۔ نالے $\alpha$ اک کارڈینل عدد اے۔ ٹکراؤ کڈھن لئی منّ لو کہ ایہ نہیں۔ تاں $\card{\alpha} \in \alpha$، سو مفروضے مطابق $\cardeq{\card{\alpha}}{\card{\alpha} \times \card{\alpha}}$، تے تعریف مطابق $\cardeq{\card{\alpha}}{\alpha}$۔ سو $\cardeq{\alpha}{\alpha\times\alpha}$، \olref{simplecardproduct} مطابق۔

ہُن ہر $\tuple{\gamma_1, \gamma_2} \in \alpha \times \alpha$ لئی ایہ ٹکڑا ویکھو:
\begin{align*}
	\text{Seg}(\gamma_1, \gamma_2) &= \Setabs{\tuple{\delta_1, \delta_2} \in \alpha \times \alpha}{\tuple{\delta_1, \delta_2} \canonord \tuple{\gamma_1, \gamma_2}}
\end{align*}
$\gamma = \max(\gamma_1, \gamma_2)$ رکھ کے دھیان دیو کہ $\tuple{\gamma_1, \gamma_2} \canonord \tuple{\gamma+1, \gamma + 1}$۔ سو جدوں $\gamma$ لامتناہی اے، ویکھو:
\begin{align*}
	\text{Seg}(\gamma_1, \gamma_2) &
	\precsim ((\gamma \ordplus 1)\ordtimes (\gamma \ordplus 1))\\
	&\approx (\gamma \ordtimes \gamma)
	\text{، \olref[ord-arithmetic][using-addition]{ordinfinitycharacter} تے
	\olref{simplecardproduct} مطابق}\\
	&\approx \gamma \text{، استقرا دے مفروضے مطابق}\\
	& \prec \alpha\text{، کیوں جے $\alpha$ اک کارڈینل عدد اے}
\end{align*}
سو $\ordtype{\alpha\times \alpha, \canonord} \leq \alpha$، تے ایس لئی $\cardle{\alpha \times \alpha}{\alpha}$۔ کیوں جے $\cardle{\alpha}{\alpha \times \alpha}$ وی، شروڈر–برنسٹین توں نتیجہ نکلدا اے۔
\end{proof}

آخر اسیں اوہ نتیجہ لئیے جیہڑا حساب سادہ کردا اے۔

\begin{thm}\ollabel{cardplustimesmax}
جے $\cardfont{a}, \cardfont{b}$ لامتناہی کارڈینل عدد نیں، تاں:
\[
	\cardfont{a}
\cardtimes \cardfont{b} = \cardfont{a} \cardplus \cardfont{b} =
\text{max}(\cardfont{a}, \cardfont{b}).
\]
\end{thm}

\begin{proof}
عمومیت گنوائے بناں منّ لو $\cardfont{a} = \max(\cardfont{a},
\cardfont{b})$۔ فیر \olref{alphatimesalpha} نوں ورت کے،
$\cardfont{a}\cardtimes\cardfont{a} = \cardfont{a} \leq \cardfont{a}
\cardplus \cardfont{b} \leq \cardfont{a} \cardplus \cardfont{a} \leq
\cardfont{a} \cardtimes \cardfont{a}$۔ \end{proof}\noindent ایویں ای جے $\cardfont{a}$ لامتناہی اے، تاں $\cardfont{a}$ جِنّی گنتی دے اوہناں سیٹاں دا اتحاد، جیہناں وچوں ہر اک دی وڈائی $\leq\cardfont{a}$ اے، آپ $\leq\cardfont{a}$ وڈائی دا اے۔

\begin{prop}\ollabel{kappaunionkappasize}
$\cardfont{a}$ اک لامتناہی کارڈینل عدد ہووے۔ ہر ترتیبی عدد $\beta
\in \cardfont{a}$ لئی $X_\beta$ اک ایہو جیہا سیٹ ہووے جیہدے لئی $\card{X_\beta} \leq
\cardfont{a}$۔ تاں $\card{\bigcup_{\beta \in \cardfont{a}} X_\beta}
\leq \cardfont{a}$۔
\end{prop}

\begin{proof}
ہر $\beta \in \cardfont{a}$ لئی اک !!a{injection} $f_\beta \colon
X_\beta \to \cardfont{a}$ مقرر کرو۔\footnote{ایہ کیویں «مقرر» ہُندے نیں؟ \olref[sth][choice][countablechoice]{sec} ویکھو۔} اک !!a{injection} $g \colon
\bigcup_{\beta \in \cardfont{a}} X_\beta \to \cardfont{a} \times
\cardfont{a}$ دی تعریف $g(v) = \tuple{\beta, f_\beta(v)}$ نال دیو، جتھے $v \in
X_\beta$ تے $v \notin X_\gamma$ کسے وی $\gamma \in \beta$ لئی۔ ہُن $\bigcup_{\beta \in \cardfont{a}} X_\beta \preceq \cardfont{a} \times
\cardfont{a} \approx \cardfont{a}$، \olref{cardplustimesmax} مطابق۔
\end{proof}
\endgroup
% /OLP-0588
% OLP-0589 content/set-theory/card-arithmetic/expotough.tex
\begingroup
\def\olfilename{expotough}
\expandafter\def\csname ol@sectioncs\endcsname{section}
\olfileid{sth}{card-arithmetic}{expotough}

\olsection[کجھ سادگیاں]{کارڈینل گھات وچ کجھ سادگیاں}

بھاویں $\canonord$ دی تعریف ذرا گنجل دار سی، اوہتوں کارڈینل جوڑ تے ضرب بارے کم دا نتیجہ \olref[simp]{cardplustimesmax} نکلدا اے۔ پر لامتناہی درجیاں والی گھات اینی سِدھی طرح سادی نہیں ہو سکدی۔ وجہ سمجھاؤن لئی اک نتیجے توں شروع کردے آں جیہڑا محدود صورت دی جانی پچھانی خاصیت نوں ودھاؤندا اے (بھاویں ثبوت کافی تجریدی اے)۔

\begin{prop}\ollabel{simplecardexpo}
$\cardexpo{\cardfont{a}}{\cardfont{b} \cardplus \cardfont{c}} =
\cardexpo{\cardfont{a}}{\cardfont{b}} \cardtimes
\cardexpo{\cardfont{a}}{\cardfont{c}}$ تے
$\cardexpo{(\cardexpo{\cardfont{a}}{\cardfont{b}})}{\cardfont{c}} =
\cardexpo{\cardfont{a}}{\cardfont{b} \cardtimes \cardfont{c}}$،
کسے وی کارڈینل عدداں $\cardfont{a}, \cardfont{b}, \cardfont{c}$ لئی۔
\end{prop}

\begin{proof}
پہلے دعوے لئی اک فنکشن $f \colon
(\cardfont{b}\disjointsum\cardfont{c}) \to \cardfont{a}$ ویکھو۔ ہُن ایہنوں «دو حصیاں وچ ونڈو»۔ $f_\cardfont{b}(\beta) = f(\beta, 0)$ ہر $\beta \in \cardfont{b}$ لئی، تے $f_\cardfont{c}(\gamma) = f(\gamma, 1)$ ہر $\gamma \in \cardfont{c}$ لئی تعریف کرو۔ نقشہ $f \mapsto
(f_{\cardfont{b}} \times f_\cardfont{c})$ اک !!a{bijection}
$\funfromto{\cardfont{b} \disjointsum \cardfont{c}}{\cardfont{a}} \to
(\funfromto{\cardfont{b}}{\cardfont{a}} \times
\funfromto{\cardfont{c}}{\cardfont{a}})$ اے۔

دوجے دعوے لئی اک فنکشن $f \colon \cardfont{c} \to
(\funfromto{\cardfont{b}}{\cardfont{a}})$ ویکھو۔ سو ہر $\gamma \in
\cardfont{c}$ لئی ساڈے کول کوئی فنکشن $f(\gamma) \colon \cardfont{b} \to
\cardfont{a}$ اے۔ ہُن $f^*(\beta, \gamma) = (f(\gamma))(\beta)$ ہر
$\tuple{\beta, \gamma} \in \cardfont{b} \times \cardfont{c}$ لئی تعریف کرو۔ نقشہ $f \mapsto f^*$ اک !!a{bijection}
$\funfromto{\cardfont{c}}{(\funfromto{\cardfont{b}}{\cardfont{a}})}
\to \funfromto{\cardfont{b} \cardtimes \cardfont{c}}{\cardfont{a}}$ اے۔
\end{proof}

ہُن اسیں ایہ \emph{چاہندے} آں کہ لامتناہی کارڈینل عدداں لئی
$\cardexpo{\cardfont{a}}{\cardfont{b}}$ کڈھن دا سوکھا طریقہ ہووے۔ ایس پاسے اک چنگا قدم ایہ اے۔

\begin{prop}\ollabel{cardexpo2reduct}
جے $2 \leq \cardfont{a} \leq \cardfont{b}$ تے $\cardfont{b}$ لامتناہی اے، تاں $\cardexpo{\cardfont{a}}{\cardfont{b}} =
\cardexpo{2}{\cardfont{b}}$۔
\end{prop}

\begin{proof}
\begin{align*}
	\cardexpo{2}{\cardfont{b}} &\leq
	\cardexpo{\cardfont{a}}{\cardfont{b}}\text{، کیوں جے $2 \leq \cardfont{a}$}\\
	&\leq \cardexpo{(2^\cardfont{a})}{\cardfont{b}}
	\text{، \olref[opps]{lem:SizePowerset2Exp} مطابق}\\
	&= \cardexpo{2}{\cardfont{a} \cardtimes \cardfont{b}}
	\text{، \olref{simplecardexpo} مطابق}\\
	&= \cardexpo{2}{\cardfont{b}}
	\text{، \olref[simp]{cardplustimesmax} مطابق}
\end{align*}
\end{proof}

سانوں ایہ آس نہیں رکھنی چاہیدی کہ اسیں ایہنوں ہور سادہ کر سکدے آں، کیوں جے $\cardfont{b} < \cardexpo{2}{\cardfont{b}}$، \olref[card-arithmetic][opps]{lem:SizePowerset2Exp} مطابق۔ پر ایس توں سانوں پتا نہیں لگدا کہ $\cardexpo{\cardfont{a}}{\cardfont{b}}$ بارے کی آکھیے جدوں $\cardfont{b} <
\cardfont{a}$۔ بے شک جے $\cardfont{b}$ \emph{محدود} اے، سانوں پتا اے کہ کی کرنا اے۔

\begin{prop}
جے $\cardfont{a}$ لامتناہی اے تے $n \in \omega$، تاں
$\cardexpo{\cardfont{a}}{n} = \cardfont{a}$۔
\end{prop}

\begin{proof}
$\cardexpo{\cardfont{a}}{n} = \cardfont{a} \cardtimes \cardfont{a}
\cardtimes \ldots \cardtimes \cardfont{a} = \cardfont{a}$، \olref[simp]{cardplustimesmax} مطابق۔
\end{proof}
\noindent
نالے کجھ ہور صورتاں وچ اسیں $\cardexpo{\cardfont{a}}{\cardfont{b}}$ دی وڈائی قابو وچ رکھ سکدے آں۔

\begin{prop}
جے $2 \leq \cardfont{b} < \cardfont{a} \leq
\cardexpo{2}{\cardfont{b}}$ تے $\cardfont{b}$ لامتناہی اے، تاں
$\cardexpo{\cardfont{a}}{\cardfont{b}} = \cardexpo{2}{\cardfont{b}}$۔
\end{prop}

\begin{proof}
$\cardexpo{2}{\cardfont{b}}\leq \cardexpo{\cardfont{a}}{\cardfont{b}}
\leq \cardexpo{(\cardexpo{2}{\cardfont{b}})}{\cardfont{b}} =
\cardexpo{2}{\cardfont{b}\cardtimes\cardfont{b}} =
\cardexpo{2}{\cardfont{b}}$، جیویں \olref{cardexpo2reduct} وچ دلیل دتی سی۔
\end{proof}
\noindent
پر ایتھوں اگے گل کافی گنجل دار ہو جاندی اے۔
\endgroup
% /OLP-0589
% OLP-0590 content/set-theory/card-arithmetic/ch.tex
\begingroup
\def\olfilename{ch}
\expandafter\def\csname ol@sectioncs\endcsname{section}

\olfileid{sth}{card-arithmetic}{ch}

\olsection{استمرار دا مفروضہ}

پچھلا نتیجہ (ٹھیک ای) اشارہ کردا اے کہ کارڈینل گھات خاصی \emph{سوکھی} ہووے گی جے ضمانت ہووے کہ لامتناہی کارڈینل عدد $2$ دیاں گھاتاں نال، یعنی (\olref[opps]{lem:SizePowerset2Exp} مطابق) ذیلی سیٹاں دے سیٹ لین نال، «سِدھے طریقے نال چلن»۔ پر اسیں ایہ منّ نہیں سکدے کہ لامتناہی کارڈینل عدد ذیلی سیٹاں دے سیٹاں نال \emph{واقعی} سِدھے طریقے نال چلدے نیں۔

ایہ گل کھولھنی شروع کرن لئی اسیں کجھ سوہنی لکھت پیش کردے آں۔

\begin{defn}
$\cardsucc{\cardfont{a}}$ سبھ توں چھوٹا کارڈینل عدد ہووے جیہڑا $\cardfont{a}$ توں سختی نال وڈا اے۔ اسیں دو لامتناہی لڑیواں ایویں تعریف کردے آں:
\begin{align*}
	\aleph_{0} &\defis \omega &
	\beth_{0} &\defis \omega\\
	\aleph_{\alpha \ordplus 1} &\defis \cardsucc{(\aleph_{\alpha})} &
	\beth_{\alpha+1} &\defis \cardexpo{2}{\beth_{\alpha}}\\
	\aleph_{\alpha} &\defis \bigcup_{\beta< \alpha} \aleph_{\beta} &
	\beth_{\alpha} &\defis \bigcup_{\beta < \alpha}\beth_{\beta} & \text{جدوں $\alpha$ حدّی ترتیبی عدد اے}.
	\end{align*}
\end{defn}

$\cardsucc{\cardfont{a}}$ دی تعریف ٹھیک اے، کیوں جے \olref[cardinals][classing]{lem:NoLargestCardinal} دسدا اے کہ ہر کارڈینل عدد $\cardfont{a}$ لئی کوئی کارڈینل عدد $\cardfont{a}$ توں وڈا اے، تے لامتناہی درجیاں اُتے استقرا ضمانت دیندا اے کہ $\cardfont{a}$ توں وڈا \emph{سبھ توں چھوٹا} کارڈینل عدد موجود اے۔ $\cardfont{a}$ دی باقی تعریف لامتناہی درجیاں والی بازگشت نال ملدی اے۔

کانتور نے «$\aleph$» والی لکھت پیش کیتی۔ ایہ \emph{الف} اے، عبرانی ورن مالا دا پہلا اکھر، تے «لامتناہی» لئی عبرانی لفظ دا وی پہلا اکھر۔ پئرس نے «$\beth$» والی لکھت پیش کیتی۔ ایہ \emph{بیت} اے، عبرانی ورن مالا دا دوجا اکھر۔\footnote{پئرس نے ایہ لکھت دسمبر 1900 وچ کانتور نوں لکھے اک خط وچ ورتی سی۔ افسوس، اوہنے اوتھے اک غلط دلیل وی دتی کہ $\beth_\alpha$ موجود نہیں جدوں $\alpha \geq
\omega$۔} ہُن ایہ لکھتاں سانوں لامتناہی کارڈینل عدد دیندیاں نیں۔

\begin{prop}
$\aleph_\alpha$ تے $\beth_\alpha$، ہر ترتیبی عدد $\alpha$ لئی کارڈینل عدد نیں۔
\end{prop}

\begin{proof}
دوویں نتیجے لامتناہی درجیاں اُتے سادے استقرا نال نکلدے نیں۔ $\aleph_0 =
\beth_0 = \omega$ کارڈینل عدد اے، \olref[cardinals][classing]{omegaisacardinal} مطابق۔ جے منّ لئیے کہ $\aleph_\alpha$ تے $\beth_\alpha$ دوویں کارڈینل عدد نیں، تاں $\aleph_{\alpha+1}$ تے $\beth_{\alpha+1}$ تعریف وچ صاف کارڈینل عدد نیں۔ تے کارڈینل عدداں دے سیٹ دا اتحاد اک کارڈینل عدد اے، \olref[cardinals][classing]{unioncardinalscardinal} مطابق۔
\end{proof}
\noindent
نالے ہر لامتناہی کارڈینل عدد اک $\aleph$ اے۔

\begin{prop}
جے $\cardfont{a}$ اک لامتناہی کارڈینل عدد اے، تاں $\cardfont{a} =
\aleph_\gamma$ کسے اکّو اک $\gamma$ لئی۔
\end{prop}

\begin{proof}
کارڈینل عدداں اُتے لامتناہی درجیاں والے استقرا نال۔ استقرا لئی منّ لو کہ جے $\cardfont{b} < \cardfont{a}$ تاں $\cardfont{b} =
\aleph_{\gamma_\cardfont{b}}$۔ جے $\cardfont{a} =
\cardsucc{\cardfont{b}}$ کسے $\cardfont{b}$ لئی، تاں $\cardfont{a} =
\cardsucc{(\aleph_{\gamma_\cardfont{b}})}=
\aleph_{\gamma_\cardfont{b}+1}$۔ جے $\cardfont{a}$ کسے وی کارڈینل عدد دا جانشین نہیں، تاں کیوں جے کارڈینل عدد ترتیبی عدد نیں،
$\cardfont{a} = \bigcup_{\cardfont{b} < \cardfont{a}} \cardfont{b} =
\bigcup_{\cardfont{b} < \cardfont{a}}{\aleph_{\gamma_\cardfont{b}}}$،
سو $\cardfont{a} = \aleph_\gamma$ جتھے $\gamma =
\bigcup_{\cardfont{b} < \cardfont{a}}\gamma_\cardfont{b}$۔
\end{proof}

ہر لامتناہی کارڈینل عدد اک $\aleph$ اے، سو سوال اُٹھدا اے کہ کی ہر لامتناہی کارڈینل عدد اک~$\beth$ اے؟ جے ایہ گل \emph{ہووے} تاں لامتناہی کارڈینل عدد ذیلی سیٹاں دے سیٹ لین والے عمل نال «سِدھے طریقے نال چلن»۔ سچ مُچ، سانوں ہِیٹھلی گل مل جائے گی۔

\begin{defish}
\emph{استمرار دا عام مفروضہ} (GCH)۔ $\aleph_\alpha  = \beth_\alpha$ سارے $\alpha$ لئی۔
\end{defish}

نالے جے GCH سچ ہووے تاں اسیں کارڈینل گھات نوں کافی سادہ کر سکدے آں۔ خاص طور اُتے، اسیں وکھا سکدے آں کہ جدوں $\cardfont{b} < \cardfont{a}$، تاں $\cardexpo{\cardfont{a}}{\cardfont{b}}$ دی قدر ایہناں حداں اندر اے، $\cardfont{a}\leq \cardexpo{\cardfont{a}}{\cardfont{b}} \leq
\cardsucc{\cardfont{a}}$۔ فیر اسیں اوہ عین شرطاں وی دَس سکدے آں جیہڑیاں طے کردیاں نیں کہ دوہاں وچوں کیہڑی صورت اے (یعنی $\cardfont{a} = \cardexpo{\cardfont{a}}{\cardfont{b}}$ یا $\cardexpo{\cardfont{a}}{\cardfont{b}} =
\cardsucc{\cardfont{a}}$)۔\footnote{شرط دا فیصلہ \emph{ہم حدی} کردی اے۔}

پر GCH اک \emph{مفروضہ} اے، \emph{قضیہ} نہیں۔ \citet{Godel1938} نے ثابت کیتا کہ جے $\ZFC$ بے ٹکراؤ اے تاں $\ZFC + \text{GCH}$ وی بے ٹکراؤ اے۔ پر بعد وچ پتا لگیا کہ اسیں $\lnot$GCH نوں وی $\ZFC$ وچ پا سکدے آں۔ اصل وچ GCH دی سبھ توں سوکھی، پر غیر معمولی \emph{صورت} ویکھو، یعنی:

\begin{defish}
\emph{استمرار دا مفروضہ} (CH)۔ $\aleph_1 = \beth_1$۔
\end{defish}

\citet{Cohen1963} نے ثابت کیتا کہ جے $\ZFC$ بے ٹکراؤ اے تاں $\ZFC
+ \lnot\text{CH}$ وی بے ٹکراؤ اے۔ سو استمرار دا مفروضہ $\ZFC$ توں آزاد اے۔

ایہنوں استمرار دا مفروضہ ایس لئی آکھدے نیں کہ «استمرار» حقیقی عددی لکیر $\Real$ دا اک ہور ناں اے۔ \olref[opps]{continuumis2aleph0} دسدا اے کہ $\card{\Real} =
\beth_1$۔ سو استمرار دا مفروضہ آکھدا اے کہ قدرتی عدداں دی کارڈینیلٹی $\aleph_0 = \beth_0$ تے استمرار دی کارڈینیلٹی $\beth_1$ دے وچکار کوئی کارڈینل عدد نہیں۔

(G)CH دی $\ZFC$ توں \emph{آزادی} نوں ویکھدیاں، اوہناں دے \emph{سچ} ہون بارے کی آکھیے؟ کہن نوں \emph{بہت} کجھ اے۔ اصل وچ سیٹ نظریے دے حال دے بوہتے نتیجہ خیز کم دا رخ ایہناں مسئلیاں دی پڑتال ول رہیا اے۔ پر دو چھوٹیاں گلاں اُتے ضرور زور دینا چاہیدا اے۔

پہلی، ایہناں رسمی آزادی دے نتیجیاں توں \emph{سِدھا ای} ایہ نہیں نکلدا کہ GCH یا CH دی سچ قدر \emph{غیر معین} اے۔ آخر ہو سکدا اے سانوں بس ہور مسلّمے پاؤن دی لوڑ ہووے جیہڑے سانوں سہج لگن تے سوال دا فیصلہ اک پاسے یا دوجے پاسے کر دین۔ گوڈل نے آپ آکھیا سی کہ ایہ ٹھیک جواب اے۔

دوجی، CH دی $\ZFC$ توں آزادی ضرور \emph{دھیان کھچن والی} اے، پر (لفظی معنی وچ) ایہ \emph{ناقابلِ یقین} ضرور نہیں۔ سِدھی گل ایہ اے کہ $\ZFC$ سانوں اینا کجھ دسدا اے، پر کارڈینل عدداں توں اوہناں دے جانشیناں تک جان لئی شاید ذیلی سیٹاں دے سیٹ لین نالوں کوئی ہور باریک طریقہ چاہیدا اے۔
 %The operation of taking powersets moves us from one stage of the
 %hierarchy to its successor stage (from $V_{\alpha}$ to
 %$V_{\alpha+1}$). 

ایہ دو گلاں آکھ کے، جے تسی ہور جاننا چاہندے او تاں ہُن تہانوں اوہناں وکھ وکھ فلسفیاں تے ریاضی داناں ول جانا پئے گا جیہڑے ایس بحث وچ حصہ لیندے نیں۔\footnote{شروع کرن لئی تسی \citet[\S15.6]{Potter2004} پڑھنا چاہو گے۔}
\endgroup
% /OLP-0590
% OLP-0591 content/set-theory/card-arithmetic/fix.tex
\begingroup
\def\olfilename{fix}
\expandafter\def\csname ol@sectioncs\endcsname{section}

\olfileid{sth}{card-arithmetic}{fix}
\olsection{\texorpdfstring{$\aleph$ دے قائم نقطے}{ℵ دے قائم نقطے}}

\olref[spine][]{chap} وچ اسیں آکھیا سی کہ بدل رکھن دے اصول دا جواز دین دی لوڑ اے، کیوں جے ایہ درجہ وار سلسلے نوں خاصا اُچا کر دیندا اے۔ ہُن کجھ کارڈینل حساب کر کے اسیں سلسلے دی اُچائی دی اک نِکّی جیہی مثال دے سکدے آں۔

صاف اے کہ $0 < \aleph_0$، تے $1 < \aleph_1$، تے $2 < \aleph_2$\ldots تے سچ مُچ ہر قدم اُتے وڈائی دا فرق ہور \emph{ودھدا} جاندا اے۔ سو ایہ اندازہ لاؤن نوں جی کردا اے کہ $\kappa< \aleph_\kappa$ ہر ترتیبی عدد $\kappa$ لئی۔

پر $\ZFC$ منّ کے ایہ اندازہ \emph{غلط} اے۔ اصل وچ اسیں ثابت کر سکدے آں کہ \emph{$\aleph$ دے قائم نقطے} موجود نیں، یعنی ایہو جیہے کارڈینل عدد $\kappa$ جنہاں لئی $\kappa=\aleph_\kappa$۔

\begin{prop}\ollabel{alephfixed}
اک $\aleph$ دا قائم نقطہ موجود اے۔
\end{prop}

\begin{proof}
بازگشت ورت کے ایویں تعریف کرو:
\begin{align*}
	\kappa_0 &= 0\\
	\kappa_{n+1} &= \aleph_{\kappa_n}\\
	\kappa&= \bigcup_{n < \omega}\kappa_n
\end{align*}
ہُن $\kappa$ اک کارڈینل عدد اے، \olref[cardinals][classing]{unioncardinalscardinal} مطابق۔ پر ہُن:
\[
	\kappa= \bigcup_{n < \omega} \kappa_{n+1} =
	\bigcup_{n < \omega}\aleph_{\kappa_n} =
	\bigcup_{\alpha < \kappa}\aleph_\alpha = \aleph_\kappa
\]
%	By construction, $\kappa$ is the least cardinal greater than each
%	$\kappa_n$. So $\aleph_\kappa$ is the least cardinal greater than
%	each $\aleph_{\kappa_n}$, and hence greater than each
%	$\kappa_{n+1}$. But equally $\kappa$ is the least cardinal greater
%	than each $\kappa_{n+1} = \aleph_{\kappa_n}$. So $\kappa=
%	\aleph_\kappa$.
\end{proof}

بُولوس نے اک واری این اوس $\aleph$ دے قائم نقطے بارے مضمون لکھیا سی جیہڑا اسیں ہُن بنایا اے۔ مضمون دے مُڈھ وچ $\kappa$ دی موجودگی دسّن پچھوں اوہنے آکھیا:
\begin{quote}
	[$\kappa$] اوہناں لوکاں دی سمجھ مطابق جیہناں دا سیٹ نظریے نال پہلاں کوئی واسطہ نہیں پیا، \emph{خاصا وڈا} عدد اے۔ مینوں لگدا اے، اینا وڈا کہ ایہ کسے وی اوہ نظریے دے سچ اُتے سوال کھڑا کردا اے جیہدے دعویاں وچ ایہ دعویٰ وی اے کہ گھٹ توں گھٹ $\kappa$ شےواں موجود نیں۔
	\citep[p.~257]{Boolos2000}
\end{quote}
تے آخر اوہنے اپنا مضمون ایہ پُچھ کے مکایا:
\begin{quote}
	[کی] سانوں شک اے کہ بھاویں کہانی دے مُڈھ وچ جو وی حال سی، جدوں اسیں اینی دور آ گئے آں، پہیے خالی گھیڑی کھا رہے نیں تے اسیں ہُن کسے واقعی موجود شے دا بیان نہیں سُن رہے؟
	\citep[p.~268]{Boolos2000}
\end{quote}
جے اسیں واقعی «کسے وی حقیقت» توں اگے نکل گئے آں، تاں قصور سِدھا بدل رکھن دے اصول دے سر رکھنا پئے گا۔ کیوں جے ایس مسلّمے نال ای ساڈا ثبوت چلدا اے۔ ایس صورت وچ ایہ سمجھ آؤندی اے کہ بُولوس نوں کئی دہاکے پہلاں کیتے اپنے دعوے ول مُڑنا پئے گا کہ بدل رکھن دے اصول دے «کوئی ناپسند» نتیجے نہیں (\olref[replacement][extrinsic]{sec} ویکھو)۔

پر کی $\kappa$ دی موجودگی اینی ماڑی گل اے؟ ایتھے رسل دے \emph{ٹرسٹرم شینڈی دے پیراڈوکس} نوں ویکھنا کم آ سکدا اے۔ ٹرسٹرم شینڈی اپنی زندگی اپنی ڈائری وچ لکھدا اے، پر اوہنوں اک دن لکھن اُتے اک سال لگدا اے۔ ہر لنگھدے سال نال ٹرسٹرم ہور پچھے رہ جاندا اے۔ اک سال پچھوں اوہنے بس اک دن لکھیا اے تے 364 دن لکھے بناں جی چکیا اے۔ دو سالاں پچھوں اوہنے بس دو دن لکھے نیں تے 728 دن لکھے بناں جی چکیا اے۔ تِن سالاں پچھوں اوہنے بس تِن دن لکھے نیں تے 1092 دن لکھے بناں جی چکیا اے \dots\footnote{لیپ سال چھڈ کے۔} فیر وی جے ٹرسٹرم \emph{امر} اے، تاں اوہ ہر دن لکھ لوے گا، کیوں جے اوہ اپنی زندگی دے $n$ ویں دن نوں اپنے $n$ ویں سال وچ لکھے گا۔ سو «سمیں دے آخر اُتے» ٹرسٹرم دی ڈائری پوری ہووے گی۔

ہُن ایہ اوس خیال توں اینی وکھری کیویں اے کہ $\alpha$، $\aleph_\alpha$ توں چھوٹا اے، تے ہور تے ہور چھوٹا ہُندا جاندا اے، جدوں تک اسیں $\kappa$ تک نہیں پہنچدے؟ اوس تھاں فرق پورا ہو جاندا اے تے $\kappa
= \aleph_\kappa$۔

ایہ گل لاںبھے رکھ کے تے $\ZFC$ منّ کے، آؤ قائم نقطے بناؤن دیاں کجھ ہور دلچسپ گلاں نال باب مکائیے۔ اگلے تِن نتیجے سہج سمجھ وچ وکھاندے نیں کہ اک (غیر معمولی) نقطہ اے جتھے درجہ وار سلسلہ جِنّا اُچا اے، اِنّا ای چوڑا اے۔

\begin{prop}\ollabel{bethfixed}
اک $\beth$ دا قائم نقطہ اے، یعنی ایہو جیہا $\kappa$ کہ $\kappa=
\beth_\kappa$۔
\end{prop}

\begin{proof}
\olref{alephfixed} وانگ، «$\beth$» نوں «$\aleph$» دی تھاں ورت کے۔
\end{proof}

\begin{prop}\ollabel{stagesize}
$\card{V_{\omega+\alpha}} = \beth_{\alpha}$۔ جے $\omega \ordtimes
\omega \leq \alpha$، تاں $\card{V_\alpha} = \beth_\alpha$۔
\end{prop}

\begin{proof}
پہلا دعویٰ لامتناہی درجیاں اُتے سادے استقرا نال نکلدا اے۔ دوجا دعویٰ نکلدا اے کیوں جے جدوں $\omega \ordtimes \omega \leq \alpha$ تاں
$\omega + \alpha = \alpha$۔ ایہ ثابت کرن لئی اسیں \olref[ord-arithmetic][]{chap} توں ترتیبی حساب دے حقائق ورتدے آں۔ پہلاں دھیان دیو کہ $\omega \ordtimes \omega = \omega \ordtimes (1 \ordplus \omega) =
(\omega  \ordtimes 1) \ordplus (\omega\ordtimes\omega) = \omega
\ordplus (\omega \ordtimes \omega)$۔ ہُن جے $\omega \ordtimes \omega
\leq \alpha$، یعنی $\alpha = (\omega\ordtimes\omega) \ordplus \beta$
کسے $\beta$ لئی، تاں $\omega \ordplus \alpha = \omega \ordplus
((\omega \ordtimes \omega) \ordplus \beta) = (\omega \ordplus (\omega
\ordtimes \omega)) \ordplus \beta = (\omega \ordtimes \omega) \ordplus
\beta = \alpha$۔
\end{proof}

\begin{cor}
اک ایہو جیہا $\kappa$ اے کہ $\card{V_\kappa} = \kappa$۔
\end{cor}

\begin{proof}
$\kappa$ اک $\beth$ دا قائم نقطہ ہووے، جیویں \olref{bethfixed} وچ دِتا اے۔ صاف اے کہ $\omega \ordtimes \omega < \kappa$۔ سو $\card{V_\kappa} =
\beth_\kappa= \kappa$، \olref{stagesize} مطابق۔
\end{proof}

$V_\kappa$ توں تھلّے اینے ای مرحلے نیں جِنّے $V_\kappa$ دے !!{element}s نیں۔ سو سہج سمجھ مطابق $V_\kappa$ جِنّا اُچا اے، اِنّا ای چوڑا اے۔ ایہ ٹرسٹرم شینڈی ورگی گل اے۔ اسیں اک مرحلے توں اگلے ول \emph{ذیلی سیٹاں دے سیٹ} لے کے جاندے آں، تے ہر قدم اُتے درجہ وار سلسلہ \emph{کِتے} وڈا ہو جاندا اے۔ پر «آخر وچ»، یعنی مرحلے $\kappa$ اُتے، سلسلے دی چوڑائی اوہدی اُچائی نوں جا ملدی اے۔

کوئی پُچھ سکدا اے، \emph{سلسلے دی چوڑائی اوہدی اُچائی دے برابر کِنّی واری ہُندی اے؟} جواب اے، \emph{اِنّی واری جِنّے ترتیبی عدد نیں۔} پر ایہدی تھوڑی وضاحت چاہیدی اے۔

اسیں اک ٹرم $\tau$ ایویں تعریف کردے آں۔ کسے وی $A$ لئی:
\begin{align*}
	\tau_0(A) & \defis \card{A}\\
	\tau_{n+1}(A) & \defis \beth_{\tau_n(A)}\\
	\tau(A) & \defis \bigcup_{n < \omega}\tau_n(A)
\intertext{\olref{bethfixed} وانگ، $\tau(A)$ اک $\beth$ دا قائم نقطہ اے، کسے وی $A$ لئی، تے صاف اے کہ $\card{A} < \tau(A)$۔
سو ہُن ایہ بازگشتی تعریف ویکھو:}
	W_0 &\defis 0\\
	W_{\alpha + 1} & \defis \tau(W_\alpha)\\
	W_\alpha & \defis \bigcup_{\beta < \alpha} W_\beta
	\text{، جدوں $\alpha$ حدّی اے}
\end{align*}
ایہ بناوٹ سارے ترتیبی عدداں لئی تعریف کیتی اے۔ سو سہج سمجھ مطابق $W$ ترتیبی عدداں توں $\beth$ دے قائم نقطیاں ول «!!a{injection}» اے۔ تے این پہلاں وانگ، $V_{W_\alpha}$ جِنّا اُچا اے، اِنّا ای چوڑا اے، کسے وی $\alpha$ لئی۔
\endgroup
% /OLP-0591

\OLEndChapterHook
\endgroup
% /OLP-0586
% OLP-0592 content/set-theory/choice/choice.tex
\begingroup
\def\olfilename{choice}
\expandafter\def\csname ol@sectioncs\endcsname{section}

\olchapter{sth}{choice}{چُون دا مسلّمہ}

% OLP-0593 content/set-theory/choice/introduction.tex
\begingroup
\def\olfilename{introduction}
\expandafter\def\csname ol@sectioncs\endcsname{section}

\olfileid{sth}{choice}{intro}

\olsection{جان پہچان}

\crefrange{sth:cardinals::chap}{sth:card-arithmetic::chap} وچ اسیں کارڈینل عدداں نوں ترتیبی عدد منّ کے اوہناں دا نظریہ بنایا سی۔ ایہ طریقہ حُسنِ ترتیب دے مسلّمے اُتے ٹکدا اے۔ پتا لگدا اے کہ حُسنِ ترتیب اک ہور اصول، چُون دے مسلّمے، دے برابر اے، تے ایہ مسلّمہ منّن جوگا اے یا نہیں، ایس بارے سنجیدہ فلسفیانہ بحث ہوئی اے۔ ایس باب دے سوال ایہ نیں۔ مسلّمہ کیویں ورتیا جاندا اے تے کی ایہدا جواز دتا جا سکدا اے؟
\endgroup
% /OLP-0593
% OLP-0594 content/set-theory/choice/tarskiscott.tex
\begingroup
\def\olfilename{tarskiscott}
\expandafter\def\csname ol@sectioncs\endcsname{section}

\olfileid{sth}{choice}{tarskiscott}
\olsection{تارسکی–سکاٹ دی ترکیب}

\olref[cardinals][cardsasords]{defcardinalasordinal} وچ اسیں کارڈینل عدداں دی تعریف ترتیبی عدداں وجوں دتی سی۔ ایس کم لئی اسیں حُسنِ ترتیب دا مسلّمہ منّیا سی۔ ایس دی وجہ بس ایہ سی کہ شعبے وچ ایہو «چلدا معیار» اے۔

سو حُسنِ ترتیب نال جڑے فلسفیانہ مسئلیاں دی گل توں پہلاں، ایہ صاف ہونا ضروری اے کہ اسیں چلدا معیار چھڈ وی \emph{سکدے آں} تے حُسنِ ترتیب منّے \emph{بناں} کارڈینل عدداں دا نظریہ بنا سکدے آں۔ اسیں ہن وی $\cardeq{A}{B}$، $\cardle{A}{B}$ تے $\cardless{A}{B}$ دیاں اوہ تعریفاں ورت سکدے آں جیہڑیاں \olref[sfr][siz][]{chap} وچ دتیاں سن۔ سانوں بس \emph{کارڈینل عدد} دا نواں تصور چاہیدا اے۔

اک سادہ خیال ایہ ہو سکدا اے کہ $A$ دی کارڈینیلٹی ایویں تعریف کرن دی کوشش کریے:
\[
	\Setabs{x}{\cardeq{A}{x}}.
\]
تسی ایہدا مقابلہ فریگے دی $\# x Fx$ والی تعریف نال کر سکدے او، جیہدا خاکہ \olref[cardinals][hp]{sec} دے عین آخر وچ دتا سی۔ تے جیہڑیاں وجوہاں دا اوتھے اشارہ دتا سی، اوہناں کرکے ایہ تعریف نہیں چل دی۔ ہر اک رکنی سیٹ $\{\emptyset\}$ نال ہم تعداد اے۔ پر درجہ وار سلسلے دے ہر جانشین مرحلے اُتے نویں اک رکنی سیٹ بندے نیں (بس اوہ اک رکنی سیٹ ویکھو جیہدا اکّو رکن پچھلا مرحلہ اے)۔ سو $\Setabs{x}{\cardeq{A}{x}}$ موجود نہیں، کیوں جے ایہدا کوئی درجہ نہیں ہو سکدا۔

ایس اوکڑ توں لنگھن لئی اسیں تارسکی تے سکاٹ دی اک ترکیب ورتدے آں۔\footnote{چیتے رہے کہ سارے فارمولیاں وچ پیرامیٹر ہو سکدے نیں (جے صاف ہور گل نہ آکھی ہووے)۔}

\begin{defn}[تارسکی–سکاٹ]
کسے وی فارمولے $\phi(x)$ لئی، $[ x : \phi(x)] $ اوہناں سارے $x$ دا سیٹ ہووے جیہناں دا درجہ سبھ توں گھٹ ممکن درجہ اے تے جیہناں لئی $\phi(x)$ اے (یا $\emptyset$، جے کوئی $\phi$ پوری کرن والی شے نہیں)۔
\end{defn}

سانوں جانچنا چاہیدا اے کہ ایہ تعریف جائز اے۔ $\ZF$ وچ کم کردیاں، \olref[spine][foundation]{zfentailsregularity} ضمانت دیندا اے کہ $\setrank{x}$ ہر $x$ لئی موجود اے۔ ہُن جے کوئی شےواں $\phi$ اُتے پوریاں اُتردیاں نیں، تاں $\alpha$ اوہ سبھ توں گھٹ درجہ ہووے جیہدے لئی $(\exists x\subseteq V_\alpha)\phi(x)$، یعنی $(\forall \beta
\in \alpha)(\forall x \subseteq V_\beta)\lnot \phi(x)$۔ فیر وکھ کرن دے اصول نال $[x : \phi(x)]$ دی تعریف $\Setabs{x \in
V_{\alpha+1}}{\phi(x)}$ وجوں دے سکدے آں۔

تارسکی–سکاٹ دی ترکیب دا جواز دین پچھوں، ہُن اسیں ایہنوں کارڈینیلٹی دا اک تصور تعریف کرن لئی ورت سکدے آں۔

\begin{defn}
$A$ دی \textsc{ts}-کارڈینیلٹی $\text{tsc}(A) = [x :
\cardeq{A}{x}]$ اے۔
\end{defn}

اک \textsc{ts}-کارڈینل عدد دی تعریف وچ حُسنِ ترتیب دا مسلّمہ نہیں ورتیا جاندا۔ پر ایس مسلّمے توں بناں وی اسیں وکھا سکدے آں کہ \emph{\textsc{ts}-کارڈینل عدد}، \olref[cardinals][cardsasords]{defcardinalasordinal} وچ دتی تعریف والے \emph{کارڈینل عدداں} وانگ کم کردے نیں۔ مثلاً جے اسیں \olref[cardinals][cardsasords]{lem:CardinalsBehaveRight} تے \olref[card-arithmetic][opps]{lem:SizePowerset2Exp} نوں \textsc{ts}-کارڈینل عدداں دے حساب نال فیر لکھیے، تاں ثبوت $\ZF$ وچ ٹھیک چلدے نیں، حُسنِ ترتیب دا مسلّمہ منّے بناں۔

ایہدی گل کردیاں دھیان دیو کہ اسیں تارسکی–سکاٹ دی ترکیب ورت کے ترتیبی عدداں دا نظریہ وی بنا سکدے آں۔ جے $\tuple{A, <}$ اک حُسنِ ترتیب اے، تاں $\text{tso}(A, <) = [\tuple{X,
R} : \ordeq{\tuple{A, <}}{\tuple{X, R}}]$ ہووے۔ کارڈینل تے ترتیبی عدداں دے ایس بیان بارے ہور لئی \citet[chs.~9--12]{Potter2004} ویکھو۔
\endgroup
% /OLP-0594
% OLP-0595 content/set-theory/choice/hartogs.tex
\begingroup
\def\olfilename{hartogs}
\expandafter\def\csname ol@sectioncs\endcsname{section}

\olfileid{sth}{choice}{hartogs}
\olsection{موازنے دی گنجائش تے ہارٹوگس دا معاون قضیہ}

ایہ چنگا پاسا سی۔ ہُن ماڑا پاسا ویکھو۔ چُون دے مسلّمے توں بناں گل \emph{گنجل دار} ہو جاندی اے۔ وجہ ویکھن لئی \cite{Hartogs1915} دا اک سوہنا نتیجہ ویکھو:

\begin{lem}[\emph{$\ZF$ وچ}]\ollabel{HartogsLemma}
کسے وی سیٹ $A$ لئی اک ترتیبی عدد $\alpha$ اے جیہدے لئی $\cardnless{\alpha}{A}$
\end{lem}

\begin{proof}
جے $B \subseteq A$ تے $R \subseteq B^2$، تاں $\tuple{B, R} \subseteq
V_{\setrank{A}+4}$، \olref[ord-arithmetic][using-addition]{rankcomputation} مطابق۔ سو وکھ کرن دا اصول ورت کے ایہ ویکھو:
\[
	C = \Setabs{\tuple{B, R} \in V_{\setrank{A}+5}}{B\subseteq A
	\text{ تے $\tuple{B, R}$ اک حُسنِ ترتیب اے}}
\]
بدل رکھن دا اصول تے \olref[ordinals][ordtype]{thmOrdinalRepresentation} ورت کے ایہ سیٹ بناؤ:
\[
	\alpha = \Setabs{\ordtype{B, R}}{\tuple{B, R} \in C}.
\]
\olref[ordinals][basic]{corordtransitiveord} مطابق $\alpha$ اک ترتیبی عدد اے، کیوں جے ایہ ترتیبی عدداں دا متعدی سیٹ اے۔ آخر جے $\gamma \in \beta \in \alpha$، تاں $\beta = \ordtype{B, R}$ کسے $B \subseteq R$ لئی، جیہدے توں $\gamma = \ordtype{B_b, R_b}$ کسے $b
\in B$ لئی، \olref[ordinals][iso]{wellordinitialsegment} مطابق۔ سو $\gamma \in \alpha$۔

ٹکراؤ کڈھن لئی منّ لو کہ اک !!a{injection} $f \colon \alpha \to A$ اے۔ فیر ایہ تعریفاں لئیے:
\begin{align*}
	B &= \ran{f}\\
	R &= \Setabs{\tuple{f(\alpha), f(\beta)} \in A \times A}{\alpha \in \beta}.
\end{align*}
صاف اے کہ $\alpha = \ordtype{B, R}$ تے $\tuple{B, R} \in C$۔ سو $\alpha
\in \alpha$، جیہڑا اک ٹکراؤ اے۔
\end{proof}

ایس توں اک ڈونگھا نتیجہ نکلدا اے۔

\begin{thm}[\emph{$\ZF$ وچ}]
ہِیٹھلے دعوے برابر نیں:
\begin{enumerate}
	\item\ollabel{equivwo} حُسنِ ترتیب دا مسلّمہ
	\item\ollabel{equivcompare} کسے وی دو سیٹاں لئی یا $\cardle{A}{B}$ یا
	$\cardle{B}{A}$، جتھے $A$ تے $B$ اوہ سیٹ نیں
\end{enumerate}
\end{thm}

\begin{proof}
\emph{\olref{equivwo} $\Rightarrow$ \olref{equivcompare}۔} $A$ تے $B$ مقرر کرو۔ \olref{equivwo} مطابق $\tuple{A, R}$ تے $\tuple{B, S}$ حُسنِ ترتیب والیاں ساختاں موجود نیں۔ \olref[ordinals][ordtype]{thmOrdinalRepresentation} ورت کے $f \colon
\alpha \to \tuple{A, R}$ تے $g \colon \beta \to \tuple{B, S}$ ہم ساخت نقشے ہون۔ \olref[sth][ordinals][basic]{ordinalsaresubsets} مطابق یا $\alpha \subseteq \beta$ یا $\beta \subseteq \alpha$۔ جے $\alpha \subseteq \beta$، تاں $\comp{f^{-1}}{g} \colon A \to B$ اک !!a{injection} اے، سو $\cardle{A}{B}$۔ ایویں ای جے $\beta \subseteq \alpha$، تاں $\cardle{B}{A}$۔

\emph{\olref{equivcompare} $\Rightarrow$ \olref{equivwo}۔} $A$ مقرر کرو۔ \olref{HartogsLemma} مطابق کوئی ترتیبی عدد $\beta$ اے جیہدے لئی $\cardnless{\beta}{A}$۔ \olref{equivcompare} توں سانوں $\cardle{A}{\beta}$ ملدا اے۔ سو کوئی !!{injection} $f \colon A \to
\beta$ اے، تے ایہ انجیکشن ورت کے اسیں $A$ دے رکناں نوں حُسنِ ترتیب دے سکدے آں، ایہ ترتیب تعریف کر کے: $\Setabs{\tuple{a, b} \in A \times A}{f(a)
\in f(b)}$۔
\end{proof}
\noindent
سِدھا نتیجہ ایہ اے کہ جے حُسنِ ترتیب دا مسلّمہ نہیں چلدا، تاں کجھ سیٹاں دا اوہناں دی وڈائی دے لحاظ نال \emph{سچ مُچ موازنہ نہیں ہو سکدا}۔ سو جے حُسنِ ترتیب دا مسلّمہ نہیں چلدا، تاں لامتناہی درجیاں والے کارڈینل عدداں دا حساب گنجل دار ہووے گا۔ مثلاً سانوں ایہ خیال چھڈنا پئے گا کہ جے $A$ تے $B$ لامتناہی نیں تاں $\cardeq{\cardeq{A \disjointsum B}{A \times
B}}{M}$، جتھے $M$، $A$ تے $B$ وچوں وڈا سیٹ اے (\olref[card-arithmetic][simp]{cardplustimesmax} ویکھو)۔ اوکڑ سِدھی اے۔ جے اسیں $A$ تے $B$ دی وڈائی دا \emph{موازنہ} ای نہیں کر سکدے، تاں کیہڑا وڈا اے، ایہ پُچھن دا کوئی مطلب نہیں۔
\endgroup
% /OLP-0595
% OLP-0596 content/set-theory/choice/wellorderingproblem.tex
\begingroup
\def\olfilename{wellorderingproblem}
\expandafter\def\csname ol@sectioncs\endcsname{section}

\olfileid{sth}{choice}{woproblem}
\olsection{حُسنِ ترتیب دا مسئلہ}

صاف اے کہ ساڈے حُسنِ ترتیب نوں منّن یا نہ منّن اُتے کافی کجھ ٹکدا اے۔ پر ایس اصول دی بحث دا دھیان اک برابر دے اصول، چُون دے مسلّمے، ول رہیا اے۔ سو ہُن اسیں اوہدی گل کراں گے (تے دوہاں دی برابری ثابت کراں گے)۔

\citeyear{Cantor1883} وچ کانتور نے حُسنِ ترتیب دے مسلّمے دی تائید کیتی تے اینوں «سوچ دا اک قانون جیہڑا مینوں بنیادی، نتیجیاں وچ امیر، تے ہر تھاں لاگو ہون کرکے خاص دھیان جوگا لگدا اے» آکھیا (\citeauthor{Potter2004} \citeyear[p.~243]{Potter2004} وچ ایہ قول دِتا اے)۔ پر آخر کانتور نوں یقین ہو گیا کہ ایس «مسلّمے» دا ثبوت چاہیدا اے۔ ریاضی دان وی ایس نتیجے اُتے پہنچے۔

زرمَیلو نے \citeyear{Zermelo1904} وچ ایہ مسئلہ «حل» کیتا۔ اوہدا حل سمجھاؤن لئی سانوں کجھ تعریفاں چاہیدیاں نیں۔

\begin{defn}
اک فنکشن $f$ اک \emph{چُون والا فنکشن} اے اُدوں تے صرف اُدوں جدوں $f(x) \in x$ سارے $x \in \dom{f}$ لئی۔ اسیں آکھدے آں کہ $f$ اک \emph{$A$ لئی چُون والا فنکشن} اے اُدوں تے صرف اُدوں جدوں $f$ اک چُون والا فنکشن اے جیہدے لئی $\dom{f} = A \setminus \{\emptyset\}$۔
\end{defn}

سہج سمجھ مطابق، ہر (غیر خالی) سیٹ $x \in A$ لئی، $A$ دا چُون والا فنکشن اک خاص رکن، $f(x)$، $x$ وچوں \emph{چُندا} اے۔ سو چُون دا مسلّمہ ایہ اے:

\begin{axiom}[چُون]
	ہر سیٹ لئی اک چُون والا فنکشن اے۔
\end{axiom}

زرمَیلو نے وکھایا کہ چُون توں حُسنِ ترتیب لازم آؤندی اے، تے اُلٹے پاسے وی:

\begin{thm}[$\ZF$ وچ]\ollabel{thmwochoice}
حُسنِ ترتیب تے چُون دے مسلّمے برابر نیں۔
\end{thm}

\begin{proof}
\emph{کھبّے توں سجّے ول۔} $A$ سیٹاں دا اک سیٹ ہووے۔ تاں اتحاد دے مسلّمے نال $\bigcup A$ موجود اے، تے حُسنِ ترتیب دے مسلّمے نال کوئی $<$ اے جیہڑا $\bigcup A$ نوں حُسنِ ترتیب دیندا اے۔ ہُن ایہ تعریف لو:
\[
f(x) = \text{$<$ مطابق سبھ توں چھوٹا رکن، جیہڑا ایس سیٹ وچ اے }x
\] ہووے۔ ایہ $A$ لئی چُون والا فنکشن اے۔

\emph{سجّے توں کھبّے ول۔} $A$ مقرر کرو۔ چُون دے مسلّمے نال اک چُون والا فنکشن، $f$، اے، جیہڑا $\Pow{A} \setminus \{\emptyset\}$ لئی اے۔ لامتناہی درجیاں والی بازگشت ورت کے اک فنکشن ایویں تعریف کرو:
\begin{align*}
	g(0) &= f(A)\\
	g(\alpha) &=
		\begin{cases}
			\text{رُک جاؤ!{}} &\text{جے }A = \funimage{g}{\alpha}\\
			f(A \setminus \funimage{g}{\alpha}) & \text{ہور صورت وچ}\\
		\end{cases}
\end{align*}
«رُک جاؤ!» بس اوہ تعریف مختصر کرن لئی آکھیا اے جیہڑی ہور طرح لمّی ہو جانی سی۔ یعنی جدوں پہلی واری $A =
\funimage{g}{\alpha}$ ہووے، تاں $g(\delta) = A$ سارے $\delta \leq \alpha$ لئی ہووے۔ ہُن پہلاں اسیں صرف اینا یقین کر سکدے آں کہ ایس نال اک \emph{ٹرم} تعریف ہُندی اے (\olref[sth][spine][recursion]{transrecursionschema} توں بعد دیاں گلاں ویکھو)۔ پر اسیں وکھاواں گے کہ ساڈے کول واقعی اک فنکشن اے۔

کیوں جے $f$ اک چُون والا فنکشن اے، ہر $\alpha$ لئی (جدوں تعریف ہوئی ہووے) سانوں $g(\alpha) =
f(A \setminus \funimage{g}{\alpha}) \in A \setminus
\funimage{g}{\alpha}$ ملدا اے، یعنی $g(\alpha) \notin \funimage{g}{\alpha}$۔ سو جے $g(\alpha) = g(\beta)$، تاں $g(\beta) \notin
\funimage{g}{\alpha}$، یعنی $\beta \notin \alpha$، تے ایویں ای $\alpha \notin \beta$۔ سو تِن صورتاں دے اصول مطابق $\alpha = \beta$۔ سو $g$ !!{injective} اے۔

اگے دھیان دیو کہ اسیں واقعی رُکدے آں!{} یعنی اک (سبھ توں چھوٹا) ترتیبی عدد $\alpha$ اے جیہدے لئی $A = g[\alpha]$۔ کیوں جے منّ لو ایہ گل نہ ہووے، تاں $g$ دے !!{injective} ہون کرکے $\cardless{\alpha}{\Pow{A} \setminus \{\emptyset\}}$ ہر ترتیبی عدد $\alpha$ لئی ہووے گا، تے ایہ \olref[hartogs]{HartogsLemma} نال ٹکراؤندا اے۔ ایس توں $\ran{g} = A$ وی۔

ایہناں گلاں نوں جوڑیے تاں $g$ کسے ترتیبی عدد توں $A$ ول اک !!a{bijection} اے۔ ہُن $g$ نوں $A$ نوں حُسنِ ترتیب دین لئی ورت سکدے آں۔
\end{proof}

سو حُسنِ ترتیب تے چُون دے مسلّمے اکٹھے منّے یا ردّ کیتے جاندے نیں۔ پر سوال رہندا اے کہ اوہ منّے جان یا ردّ کیتے جان؟
\endgroup
% /OLP-0596
% OLP-0597 content/set-theory/choice/countablechoice.tex
\begingroup
\def\olfilename{countablechoice}
\expandafter\def\csname ol@sectioncs\endcsname{section}

\olfileid{sth}{choice}{countablechoice}
\olsection{گِنن جوگ سیٹاں لئی چُون}

چُون یا حُسنِ ترتیب دے مسلّمے ورتے بناں ای ہِیٹھلی گل ثابت کرنی سوکھی اے:

\begin{lem}[$\Zminus$ وچ]
ہر محدود سیٹ لئی اک چُون والا فنکشن اے۔
\end{lem}

\begin{proof}
$a = \{b_1, \ldots, b_n\}$ ہووے۔ سادگی لئی منّ لو کہ ہر $b_i
\neq \emptyset$۔ سو کجھ شےواں $c_1, \ldots, c_n$ نیں جیہناں لئی
$c_1 \in b_1, \ldots, c_n \in b_n$۔ ہُن \olref[z][pairs]{prop:pairsconsequences} مطابق سیٹ $\{\langle b_1,
c_1\rangle , \ldots, \langle b_n, c_n\rangle\}$ موجود اے، تے ایہ~$a$ لئی اک چُون والا فنکشن اے۔
\end{proof}

پر لامتناہی سیٹ ویکھدیاں ای گل گنجل دار ہو جاندی اے۔ مثلاً اُتے دتے نتیجے دی ایہ «گھٹ توں گھٹ» توسیع ویکھو:

\begin{defish}
\emph{گِنن جوگ سیٹاں لئی چُون۔} ہر \emph{گِنن جوگ} سیٹ لئی اک چُون والا فنکشن اے۔
\end{defish}

ایہ چُون دے مسلّمے دی اک خاص صورت اے۔ تے پتا لگدا اے کہ ایہ اصول کافی واری ورتیا گیا سی، پر ورتن والیاں نوں ایہدا صاف پتا نہیں سی۔ دو سوہنیاں مثالاں ویکھو۔\footnote{ایہ مثالاں \citet[\S9.4]{Potter2004} تے لوکا اِنکُرواتی دیاں نیں۔}

\begin{ex}
اک سہج خیال ایہ اے۔ کسے وی سیٹ $A$ لئی یا $\cardle{\omega}{A}$، یا $\cardeq{A}{n}$ کسے $n \in \omega$ لئی۔ ایہ ایس سہج خیال نوں آکھن دا اک طریقہ اے کہ ہر سیٹ یا محدود اے یا لامتناہی۔ کانتور تے کئی ہور ریاضی داناں نے ایہ دعویٰ ثبوت دتے بناں کیتا سی۔ احتیاط ورتدیاں اسیں ایہ گل \olref[cardinals][classing]{generalinfinitycharacter} وچ ثابت کیتی سی۔ پر اوس ثبوت وچ اسیں $\ZFC$ اندر کم کردے ساں، کیوں جے اسیں منّ رہے ساں کہ کسے وی سیٹ $A$ نوں حُسنِ ترتیب دتی جا سکدی اے، تے ایس لئی $\card{A}$ دی موجودگی دی ضمانت اے۔ یعنی اسیں چُون دا مسلّمہ صاف منّیا سی۔

اصل وچ \citet{Dedekind1888} نے ایس دعوے دا اپنا ثبوت ایویں دِتا:

\begin{thm}[$\Zminus + \text{گِنن جوگ سیٹاں لئی چُون}$ وچ]
کسے وی $A$ لئی یا $\cardle{\omega}{A}$ یا $\cardeq{A}{n}$ کسے $n \in \omega$ لئی۔
\end{thm}

\begin{proof}
منّ لو $\cardneq{A}{n}$ سارے $n \in \omega$ لئی۔ تاں خاص طور اُتے ہر $n < \omega$ لئی اک ذیلی سیٹ $A_n \subseteq A$ اے جیہدے این ٹھیک $\cardexpo{2}{n}$ رکن نیں۔ ایس لڑی $A_0, A_1, A_2,
\ldots$ نوں ورت کے اسیں ہر $n$ لئی ایویں تعریف کردے آں:
\[
	B_n = A_n \setminus \bigcup_{i < n} A_i.
\]
ہُن ایہ ویکھو:
\begin{align*}
	\card{\bigcup_{i < n}A_n}
	&\leq \card{A_0} + \card{A_1} + \ldots + \card{A_{n-1}}\\
	&=1 + 2 + \ldots + 2^{n-1}\\
	& = 2^n - 1\\
	& < 2^n = \card{A_n}
\end{align*}
ایس لئی ہر $B_n$ دا گھٹ توں گھٹ اک رکن، $c_n$، اے۔ نالے $B_n$ والے سیٹ جوڑی جوڑی کرکے وکھرے نیں، سو جے $c_n = c_m$ تاں $n = m$۔ پر ہر $c_n
\in A$۔ سو فنکشن $f(n) = c_n$ اک انجیکشن $\omega \to A$ اے۔
\end{proof}
\noindent
ڈیڈیکنڈ نے ایہ نہیں دسیا سی کہ اوہنے گِنن جوگ سیٹاں لئی چُون دا اصول ورتیا۔ پر کی \emph{تسی} اوہدی ورتوں ویکھی؟ فیر ویکھو۔ (سچ مُچ \emph{فیر ویکھو}۔)

ثبوت وچ گِنن جوگ سیٹاں لئی چُون دا اصول دو واری ورتیا گیا اے۔ اک واری اسیں اپنے سیٹاں دی لڑی $A_0$، $A_1$، $A_2$، \dots\@ لین لئی ورتیا۔ فیر رکن $c_n$ ہر~$B_n$ وچوں چُنن لئی ورتیا۔ نالے چُون دی ایس ورتوں نوں کڈھیا نہیں جا سکدا۔ \citet[p.~138]{Cohen1966} نے ثابت کیتا کہ جے ساڈے کول چُون دا کوئی وی روپ نہیں، تاں نتیجہ نہیں چلدا۔ یعنی $\ZF$ نال ایہ گل بے ٹکراؤ اے کہ کجھ سیٹاں دا~$\omega$ نال \emph{موازنہ نہیں ہو سکدا}۔
\end{ex}

\begin{ex}
\citeyear{Cantor1878} وچ کانتور نے آکھیا سی کہ گِنن جوگ سیٹاں دی گِنن جوگ گنتی دا اتحاد وی گِنن جوگ اے۔ اوہنے ثبوت نہیں دتا، شاید ایس کرکے کہ ثبوت اوہنوں صاف لگدا سی۔ اسیں احتیاط ورتدیاں ایس نتیجے دا ہور عام روپ \olref[card-arithmetic][simp]{kappaunionkappasize} وچ ثابت کیتا سی۔ پر ساڈے ثبوت نے چُون دا مسلّمہ صاف منّیا سی۔ تے گھٹ عام نتیجے دے ثبوت نوں وی گِنن جوگ سیٹاں لئی چُون دا اصول چاہیدا اے۔

\begin{thm}[$\Zminus + \text{گِنن جوگ سیٹاں لئی چُون}$ وچ]
جے $A_n$ گِنن جوگ اے ہر $n \in \omega$ لئی، تاں $\bigcup_{n <
\omega} A_n$ گِنن جوگ اے۔
\end{thm}

\begin{proof}
عمومیت گنوائے بناں منّ لو کہ ہر $A_n \neq \emptyset$۔ سو ہر $n \in \omega$ لئی اک !!a{surjection} $f_n \colon \omega
\to A_n$ اے۔ $f \colon \omega \times \omega \to \bigcup_{n <
\omega} A_n$ دی تعریف $f(m, n) = f_n(m)$ نال دیو۔ نتیجہ نکلدا اے، کیوں جے $\omega
\times \omega$ گِنن جوگ اے (\olref[sfr][siz][zigzag]{natsquaredenumerable}) تے $f$ اک !!a{surjection} اے۔
\end{proof}
\noindent
کی تسی گِنن جوگ سیٹاں لئی چُون دی ورتوں ویکھی؟ ایہ ساڈے فنکشناں دی لڑی $f_0$، $f_1$، $f_2$، \dots\footnote{چُون دی ایہو جیہی ورتوں \olref[card-arithmetic][simp]{kappaunionkappasize} وچ وی ہوئی سی، جدوں اسیں آکھیا سی «ہر $\beta \in \cardfont{a}$ لئی اک !!a{injection}
$f_\beta$ مقرر کرو»۔} چُنن لئی ورتیا اے۔ تے ایتھے وی، جے چُون دا کوئی اصول نہ ہووے تاں نتیجہ نہیں چلدا۔ خاص طور اُتے \citet{FefermanLevy1963} نے ثابت کیتا کہ $\ZF$ نال ایہ گل بے ٹکراؤ اے کہ گِنن جوگ سیٹاں دی گِنن جوگ گنتی دے اتحاد دی کارڈینیلٹی~$\beth_1$ ہووے۔ پر رسل نے ایہو گل بڑی ہاسے والی طرح آکھی اے:
\begin{quote}
  ایس دی مثال اوہ لکھپتی اے جیہڑا جدوں وی اک جوڑی بوٹ خریددا سی نال جراباں دی اک جوڑی خریددا سی، تے ہور کدی نہیں۔ دوویں خریدن دا اوہنوں اینا شوق سی کہ آخر اوہدے کول بوٹاں دیاں $\aleph_0$ جوڑیاں تے جراباں دیاں $\aleph_0$ جوڑیاں ہو گئیاں\dots\@ بوٹاں وچ اسیں سجّے تے کھبّے دا فرق کر سکدے آں۔ سو ہر جوڑی وچوں اک چُن سکدے آں، یعنی سارے سجّے بوٹ یا سارے کھبّے بوٹ۔ پر جراباں لئی چُون دا کوئی ایہو جیہا اصول آپ نہیں سُجھدا۔ تے جدوں تک اسیں ضرب والا مسلّمہ [یعنی عملی طور اُتے چُون دا مسلّمہ] نہ منّیئے، اسیں یقین نہیں کر سکدے کہ کوئی ایہو جیہی جماعت اے جیہدے وچ ہر جوڑی وچوں اک جراب اے۔
  \citep[p.~126]{Russell1919}
\end{quote}
مختصر گل ایہ اے کہ ہِیٹھلی گل ثابت کرن لئی چُون دا کوئی روپ چاہیدا اے۔ جے تہاڈے کول جراباں دیاں گِنن جوگ گنتی جوڑیاں نیں، تاں تہاڈے کول (بس) گِنن جوگ گنتی جراباں نیں۔ تے اصل وچ گِنن جوگ سیٹاں لئی چُون دے اصول (یا اوہدے برابر دے اصول) توں بناں گِنن جوگ سیٹاں دی گِنن جوگ گنتی دا اتحاد گِنن جوگ ہون توں رہ سکدا اے۔
\end{ex}

سبق ایہ اے کہ گِنن جوگ سیٹاں لئی چُون دا اصول وار وار ورتیا گیا، پر ورتن والیاں نوں ایہدا بہت پتا نہیں سی۔ فلسفیانہ سوال ایہ اے۔ اسیں گِنن جوگ سیٹاں لئی چُون دا \emph{جواز} کیویں دے سکدے آں؟

سہج جواز دین دی اک کوشش وچ اسیں اک ایہو جیہے کم دی گل کر سکدے آں جیہدے لامتناہی قدم ہون۔ منّ لو اسیں پہلی چُون $\nicefrac{1}{2}$ منٹ وچ کردے آں، دوجی چُون $\nicefrac{1}{4}$ منٹ وچ، \dots، اپنی $n$ ویں چُون $\nicefrac{1}{2^n}$ منٹ وچ، \dots\@ تاں $1$~منٹ اندر اسیں چُونواں دی اک $\omega$ لڑی بنا لئی ہووے گی، تے اک چُون والا فنکشن تعریف کر لِتا ہووے گا۔

پر ایہ خیالی تجربہ سانوں واقعی کی دس سکدا اے؟ پہلاں تاں ایہ «چُنن» دے خیال نوں کافی لفظی معنی وچ لین اُتے ٹکدا اے۔ دوجا ایہ ریاضی نوں مابعدالطبیعیاتی امکان نال بنّھدا لگدا اے۔

ہور ضروری گل ایہ اے کہ ایہ کوشش \emph{بس} گِنن جوگ سیٹاں دی چُون توں اگے، \emph{پورے} چُون دے مسلّمے دا کوئی جواز نہیں دیندی۔ جے سانوں \emph{ہر} سیٹ لئی چُون والا فنکشن چاہیدا اے، تاں سانوں «کسے وی ترتیبی عدد جِنّی لمبائی والا لامتناہی قدماں دا کم» کر سکنا پئے گا۔ صاف آکھیے تاں ایہ ہاسے والی گل اے۔
\endgroup
% /OLP-0597
% OLP-0598 content/set-theory/choice/justifications.tex
\begingroup
\def\olfilename{justifications}
\expandafter\def\csname ol@sectioncs\endcsname{section}

\olfileid{sth}{choice}{justifications}
\olsection{چُون دے مسلّمے دا اندرلا جواز}

سو وڈا سوال ایہ اے کہ کی حُسنِ ترتیب، چُون، یا سارے سیٹاں دی وڈائی دا آپس وچ موازنہ کر سکّن دا جواز دتا جا سکدا اے۔ ایہناں وچوں جیہڑا وی لئیے، گل اکّو اے۔

ایتھے \emph{اندرلا} جواز دین دی اک کوشش اے۔ \olref[z][story]{sec} وچ اسیں درجہ وار سلسلے بارے کئی اصول پیش کیتے سن۔ اوہناں وچوں اک فیر دسنا چنگا اے:
\begin{enumerate}
	\item[] \stagesacc. کسے وی مرحلے $S$ لئی، تے کسے وی اوہناں سیٹاں لئی جیہڑے مرحلے $S$ توں \emph{پہلاں} بنے سن، مرحلے $S$ اُتے اک سیٹ بندا اے جیہدے رکن این اوہو سیٹ نیں۔ مرحلے~$S$ اُتے ہور کوئی سیٹ نہیں بندا۔
\end{enumerate}
اصل وچ کئی لکھاریاں نے سجھایا اے کہ چُون دے مسلّمے دا جواز ایس اصول (یا ایہو جیہے اصول) نال دتا جا سکدا اے۔ اسیں مختصر طور اُتے ایہ طریقہ سمجھاواں گے۔

اسیں اک سوکھے جیہے نتیجے توں شروع کراں گے جیہڑا چُون دے برابر دا \emph{اک ہور} اصول دیندا اے۔

\begin{thm}[$\ZF$ وچ]\ollabel{choiceset}
چُون ہِیٹھلے اصول دے برابر اے۔ جے $A$ دے !!{element}s غیر خالی نیں تے اوہناں وچ کوئی سانجھا رکن نہیں، تاں کوئی $C$ اے جیہدے لئی $C
\cap x$ ہر $x \in A$ لئی اک رکنی سیٹ اے۔ (ایہو جیہے $C$ نوں $A$ لئی {چُون دا سیٹ} آکھدے آں۔)
\end{thm}

ایس نتیجے دا ثبوت سِدھا اے، تے اسیں ایہ پڑھن والے لئی مشق وجوں چھڈدے آں۔

\begin{prob}
\olref[sth][choice][justifications]{choiceset} ثابت کرو۔ جے اوکڑ آوے، تاں \cite[pp.~242--3]{Potter2004} وچ ثبوت لبھ سکدے او۔
\end{prob}

ضروری گل ایہ اے کہ $A$ لئی چُون دا سیٹ بس $A$ لئی چُون والے فنکشن دیاں حاصل قدراں دا سیٹ اے۔ سو چُون دا جواز دین لئی اسیں بس اوہدے برابر دے اصول دا جواز دین دی کوشش کر سکدے آں، جیہڑا چُون دے سیٹاں دی موجودگی بارے اے۔ ہُن اسیں این ایہو کرن دی کوشش کراں گے۔

$A$ دے !!{element}s غیر خالی ہون تے اوہناں وچ کوئی سانجھا رکن نہ ہووے۔ \stageshier{} مطابق (\olref[z][story]{sec} ویکھو)، $A$ کسے مرحلے~$S$ اُتے بندا اے۔ دھیان دیو کہ $\bigcup A$ دے سارے !!{element}s مرحلے $S$ توں پہلاں موجود نیں۔ ہُن \stagesacc{} مطابق، ~$S$ توں پہلاں بنے سیٹاں دے \emph{کسے وی} چُناؤ لئی اک ایہو جیہا سیٹ بندا اے جیہدے رکن این اوہو سیٹ نیں۔ دوجے لفظاں وچ، پہلاں توں موجود سیٹاں دا ہر \emph{ممکن} اکٹھ~$S$ اُتے موجود ہووے گا۔ پر ایہ ضرور \emph{ممکن} اے کہ اوہ شےواں چُنیے جیہڑیاں~$A$ لئی اک چُون دا سیٹ بنا سکن۔ ایہ بس $\bigcup A$ دا اک خاص ذیلی سیٹ اے۔ سو لوڑ مطابق کوئی ایہو جیہا چُون دا سیٹ موجود اے۔

ایہ چُون دا اندرلا جواز دین دی \emph{بڑی} مختصر کوشش اے۔ ایس خیال نوں اگے سمجھنا ہووے تاں پوٹر دا سوہنا بیان (\citeyear[\S14.8]{Potter2004}) پڑھو۔
\endgroup
% /OLP-0598
% OLP-0599 content/set-theory/choice/banach.tex
\begingroup
\def\olfilename{banach}
\expandafter\def\csname ol@sectioncs\endcsname{section}

\olfileid{sth}{choice}{banach}
\olsection{باناخ–تارسکی دا پیراڈوکس}

جیویں بُولوس نے بدل رکھن دے اصول دا \emph{باہرلا} جواز دین دی کوشش کیتی، اسیں چُون دا جواز وی اویویں باہرلے پہلوآں توں دین دی کوشش کر سکدے آں (\olref[replacement][extrinsic]{sec} ویکھو)۔ آخر چُون نوں منّن دے کئی من پسند نتیجے نیں۔ ہر کارڈینل عدد دا موازنہ، ہر سیٹ نوں حُسنِ ترتیب، کارڈینل عدداں نوں خاص قسم دے ترتیبی عدد سمجھ سکنا، تے ہور وی۔

پر کدی کدی ایہ دعویٰ ہُندا اے کہ چُون دے \emph{ناپسند} نتیجے وی نیں۔ ایہ دعویٰ بہت کرکے \cite{BanachTarski1924} دے اک نتیجے کارن ہُندا اے۔

\begin{thm}[باناخ–تارسکی دا پیراڈوکس ($\ZFC$ وچ)]
کسے وی ٹھوس گولے نوں محدود گنتی دے ٹکڑیاں وچ ونڈیا جا سکدا اے۔ اوہ ٹکڑے فیر (گھما کے تے تھاں بدل کے) جوڑے جا سکدے نیں تاں جو اوس گولے دیاں دو نقلاں بن جان۔
\end{thm}
\noindent
پہلی نظر وچ ایہ گل حیرانی والی اے۔ صاف اے کہ دو گولیاں دا حجم پہلے گولے دے حجم توں \emph{دوگنا} اے۔ پر شکل قائم رکھن والیاں حرکتاں، گھماؤ تے تھاں بدلنا، حجم نہیں بدلدیاں۔ سو ایویں لگدا اے جیویں باناخ–تارسکی سانوں جادو نال نواں مادہ پیدا کرن دیندا اے۔

گل ہور وی حیرانی والی اے۔\footnote{\citet[Theorem 3.12]{Wagon2016} ویکھو۔} ایہو جیہی دلیل وکھاندی اے کہ مٹر دے اک دانے نوں محدود گنتی دے ٹکڑیاں وچ ونڈ کے اوہناں نوں فیر (گھما کے تے تھاں بدل کے) جوڑیا جا سکدا اے، تاں جو بِگ بین دی شکل تے وڈائی والی شے بن جائے۔

پر ایہناں وچوں کوئی گل اکلّے $\ZF$ وچ نہیں چل دی۔\footnote{بھاویں باناخ–تارسکی اوہناں اصولاں نال ثابت ہو سکدا اے جیہڑے چُون توں ٹھیک ٹھیک کمزور نیں۔ \citet[303]{Wagon2016} ویکھو۔} سو ساڈے اگے اک فیصلہ اے۔ چُون ردّ کریے یا ایس «پیراڈوکس» نوں قبول کرنا سکھیے۔

اسیں آکھاں گے کہ سانوں «پیراڈوکس» قبول کرنا سکھنا چاہیدا اے۔ اصل وچ سانوں ایہ کوئی خاص پیراڈوکس لگدا ای نہیں۔ خاص طور اُتے سانوں کوئی وجہ نہیں دِسدی کہ ایہ ہِیٹھلے نتیجیاں نالوں ودھ یا گھٹ پیراڈوکس کیوں اے۔\footnote{\citet[276--7]{Potter2004}، \citet[16]{Weston2003}، تے \citet[31, 308--9]{Wagon2016} ہور مثالاں ورت کے ایہو جیہیاں گلاں کردے نیں۔}
\begin{enumerate}
	\item وقفے $(0,1)$ وچ اینے ای نقطے نیں جِنّے $\Real$ وچ نیں۔
	\\\emph{ثبوت}۔ $\tan(\pi(r-\nicefrac{1}{2})))$ ویکھو۔
	\item اک لکیر وچ اینے ای نقطے نیں جِنّے اک مربع وچ۔
	\\See \olref[his][set][pathology]{sec} تے \olref[his][set][cantorplane]{sec} ویکھو۔
	\item تھاں بھرن والیاں منحنی لکیراں موجود نیں۔
	\\See \olref[his][set][pathology]{sec} تے \olref[his][set][hilbertcurve]{sec} ویکھو۔
\end{enumerate}
ایہناں تِنّاں نتیجیاں وچوں کسے لئی چُون دی لوڑ نہیں۔ اصل وچ ہُن اسیں بس ایہناں نوں ریاضی دیاں حیرانی والیاں، سوہنیاں گلاں سمجھدے آں۔ شاید باناخ–تارسکی دے پیراڈوکس نال وی سانوں ایہو رویہ رکھنا چاہیدا اے۔

ایتھے اک فنی گل دا دھیان دینا ضروری اے، پر ایہ بس ٹھنڈے دل نال سوچن دی گل اے۔ شکل قائم رکھن والیاں حرکتاں حجم وی قائم رکھدیاں نیں۔ ایس لئی پنجے\footnote{اسیں پیراڈوکس نوں «محدود گنتی دے ٹکڑیاں» نال بیان کیتا سی۔ اصل وچ \citet{Robinson1947} نے ثابت کیتا کہ گولے نوں \emph{پنج} ٹکڑیاں وچ ونڈنا کافی اے (ایہتوں گھٹ وچ نہیں ہو سکدا)۔ ثبوت لئی \citet[pp.~66--7]{Wagon2016} ویکھو۔} ٹکڑے، جنہاں وچ گولا ونڈیا اے، سارے دے سارے \emph{ناپن جوگ} نہیں ہو سکدے۔ موٹی گل ایہ اے کہ ایہناں وکھ وکھ ٹکڑیاں دا حجم مقرر کرن دا کوئی مطلب نہیں۔ ایہناں نوں نقطیاں دے ایہو جیہے «لامتناہی کھلارے» سمجھو جنہاں دی تصویر ذہن وچ نہیں بن سکدی۔ ہُن ایہ ہو سکدا اے کہ ایہو جیہے «لامتناہی کھلرے» سیٹ سوچنا اوپرا لگے۔ پر اوہناں دی موجودگی \stagesacc{} وچ دتے حکم توں نکلدی لگدی اے کہ پہلاں موجود سیٹاں دے \emph{سارے ممکن} اکٹھ بناؤ۔

جے ایہناں وچوں کوئی گل تہانوں نہیں مناؤندی، تاں باناخ–تارسکی دا پیراڈوکس قبول کرن دے حق وچ اک آخری (باہرلی) دلیل ویکھو۔ ایہ توں ریاضی دا اج تک دا سبھ توں وڈیا لطیفہ جھٹ نکلدا اے:
\begin{enumerate}
	\item[] \emph{سوال}۔ «باناخ–تارسکی» دے اکھر اگے پچھے کر کے کیہ بن دا اے؟
	\item[] \emph{جواب}۔ «باناخ–تارسکی باناخ–تارسکی»۔
\end{enumerate}
\endgroup
% /OLP-0599
% OLP-0600 content/set-theory/choice/vitali.tex
\begingroup
\def\olfilename{vitali}
\expandafter\def\csname ol@sectioncs\endcsname{section}

\olfileid{sth}{choice}{vitali}
\olsection{ضمیمہ: ویتالی دا پیراڈوکس}

باناخ–تارسکی دی بناوٹ قبول کرنی ٹھیک اے یا نہیں، ایہ چنگی طرح سمجھن لئی سانوں اوہدا \emph{ثبوت} ویکھنا چاہیدا اے۔ افسوس، اوہدے لئی اینی الجبرا چاہیدی اے جِنّی اسیں ایتھے نہیں دس سکدے۔ پر ہِیٹھلے نتیجے اُتے دھیان دے کے اسیں کجھ مختصر گلاں دس سکدے آں جیہڑیاں باناخ–تارسکی دے ثبوت دی کجھ سمجھ دے سکدیاں نیں۔\footnote{کِتے ودھ تفصیل لئی \cite{Weston2003} یا \cite{Wagon2016} ویکھو۔}

\begin{thm}[ویتالی دا پیراڈوکس ($\ZFC$ وچ)]\ollabel{vitaliparadox}
کسے وی دائرے نوں گِنن جوگ گنتی دے ٹکڑیاں وچ ونڈیا جا سکدا اے۔ اوہ ٹکڑے فیر (گھما کے تے تھاں بدل کے) جوڑے جا سکدے نیں تاں جو اوس دائرے دیاں دو نقلاں بن جان۔
\end{thm}

ویتالی دا پیراڈوکس باناخ–تارسکی دے پیراڈوکس نالوں ثابت کرنا کِتے سوکھا اے۔ اسیں اینوں «ویتالی دا پیراڈوکس» آکھیا اے، کیوں جے ایہ ویتالی دی \citeyear{Vitali1905} وچ اک ناپن جوگ نہ ہون والے سیٹ دی بناوٹ توں نکلدا اے۔ پر ویتالی تے باناخ–تارسکی دے پیراڈوکس دے ثبوتاں وچ سیٹ نظریے والے پاسے بڑے ملدے نیں۔ نتیجیاں وچ اصلی فرق بس ایہ اے کہ باناخ–تارسکی \emph{محدود} گنتی دے ٹکڑیاں والی ونڈ ویکھدا اے، جدکہ ویتالی دا پیراڈوکس \emph{گِنن جوگ لامتناہی} گنتی دے ٹکڑیاں والی ونڈ ویکھدا اے۔ جیویں \citet{Weston2003} آکھدا اے، ویتالی دا پیراڈوکس «باناخ–تارسکی دے پیراڈوکس جِنّا حیران کن ہرگز نہیں، پر ایہ وکھاندا اے کہ ہندسی پیراڈوکس ‘سادیاں’ صورتاں وچ وی ہو سکدے نیں»۔

ویتالی دا پیراڈوکس اک دو جہتی شکل، دائرے، بارے اے۔ سو اسیں سطح $\Real^2$ اُتے کم کراں گے۔ $\rotationsgroup$ نقطیاں دے مبدأ دوالے (گھڑی دے کانٹیاں دی سمت) گھماؤواں دا سیٹ ہووے، جیہناں دے زاویے \emph{ناطق} تے $[0,2\pi)$ وچ ریڈیئن قدراں نیں۔ $\rotationsgroup$ بارے الجبرا دیاں کجھ گلاں ایہ نیں (جے نتیجے دا بیان سمجھ نہ آوے تاں ثبوت اوہدا مطلب صاف کر دے گا):

\begin{lem}\ollabel{rotationsgroupabelian}
$\rotationsgroup$ فنکشناں دی ترکیب تحت اک ابیلی {گروپ} بندا اے۔
\end{lem}

\begin{proof}
$0_{\rotationsgroup}$ نال $0$ ریڈیئن دا گھماؤ لکھدے ہوئیاں، ایہ $\rotationsgroup$ دا شناختی رکن اے، کیوں جے
$\comp{0_{\rotationsgroup}}{\rho} = \comp{\rho}{0_{\rotationsgroup}} =
\rho$ کسے وی $\rho \in \rotationsgroup$ لئی۔

ہر رکن دا اک اُلٹ اے۔ جے $\rho \in \rotationsgroup$، $r$ ریڈیئن گھماؤندا اے، تاں $\rho^{-1} \in \rotationsgroup$، $2\pi - r$ ریڈیئن گھماؤندا اے، سو $\rho \circ \rho^{-1} = 0_\rotationsgroup$۔

ترکیب وچ شرکتیت اے۔ $\comp{\rho}{(\comp{\sigma}{\tau})} =
\comp{(\comp{\rho}{\sigma})}{\tau}$ کسے وی $\rho, \sigma, \tau \in
\rotationsgroup$ لئی۔

ترکیب تبادلی اے۔ $\comp{\rho}{\sigma} =
\comp{\sigma}{\rho}$ کسے وی $\rho, \sigma \in \rotationsgroup$ لئی۔
\end{proof}

اصل وچ اسیں اپنے گروپ $\rotationsgroup$ نوں دو ادھّاں وچ ونڈ سکدے آں، تے فیر کسے وی اک ادھ نال سارا گروپ فیر حاصل کر سکدے آں۔

\begin{lem}\ollabel{disjointgroup}
$\rotationsgroup$ دی دو وکھرے سیٹاں، $\rotationsgroup_{1}$ تے $\rotationsgroup_{2}$، وچ ونڈ اے، تے دوویں $\rotationsgroup$ دی بنیاد نیں۔
\end{lem}

\begin{proof}
$\rotationsgroup_{1}$ وچ اوہ گھماؤ ہون جیہناں دیاں ناطق ریڈیئن قدراں $[0, \pi)$ وچ نیں۔ $\rotationsgroup_{2} = \rotationsgroup
\setminus \rotationsgroup_1$ ہووے۔ بنیادی الجبرا توں
$\Setabs{\comp{\rho}{\rho}}{\rho \in \rotationsgroup_1} =
\rotationsgroup$۔ $\rotationsgroup_2$ لئی وی ایہو جیہا نتیجہ مل سکدا اے۔
\end{proof}

اسیں گروپاں بارے ایس گل نوں \olref{vitaliparadox} ثابت کرن لئی ورتاں گے۔ $\onesphere$ اکائی دائرہ ہووے، یعنی اوہناں نقطیاں دا سیٹ جیہڑے سطح دے مبدأ توں این ٹھیک~$1$ اکائی دور نیں، یعنی
$\Setabs{\tuple{r,s} \in \Real^2}{\sqrt{r^2+s^2}=1}$۔ اسیں $\onesphere$ نوں حصیاں وچ ونڈاں گے، $\onesphere$ اُتے ہِیٹھلا تعلق ویکھ کے:
\[
	r \sim s \emph{ اُدوں تے صرف اُدوں جدوں }(\exists \rho \in \rotationsgroup)\rho(r) = s.
\]
یعنی $\onesphere$ دے دو نقطے ایس تعلق نال جڑے نیں اُدوں تے صرف اُدوں جدوں مبدأ دوالے ناطق قدر والے گھماؤ نال اک توں دوجے تک پہنچ سکدے او۔ ہُن ایہ گل حیرانی والی نہیں:

\begin{lem}
$\sim$ اک ہم ارزی دا تعلق اے۔
\end{lem}

\begin{proof}
\olref{rotationsgroupabelian} ورت کے ایہ گل صاف اے۔
\end{proof}

ہُن اسیں چُون ورت کے اک سیٹ، $C$، حاصل کردے آں جیہدے وچ $\onesphere$ دی $\sim$ تحت ہر ہم ارزی جماعت وچوں این ٹھیک اک رکن اے۔ یعنی اسیں ہم ارزی جماعتاں دے سیٹ اُتے اک چُون والا فنکشن $f$ لیندے آں۔\footnote{کیوں جے $\rotationsgroup$ !!{enumerable} اے، $E$ دا ہر !!{element} !!{enumerable} اے۔ کیوں جے $\onesphere$ !!{nonenumerable} اے، \olref{vitalicover} تے \olref[card-arithmetic][simp]{kappaunionkappasize} توں نکلدا اے کہ $E$ !!{nonenumerable} اے۔ سو ایتھے \emph{گِنن جوگ نہ ہون والے} سیٹ لئی چُون ورتی اے۔}
\[
	E = \Setabs{\equivrep{r}{\sim}}{r \in \onesphere},
\]
تے $C = \ran{f}$ ہووے۔ ہر گھماؤ $\rho \in \rotationsgroup$ لئی، سیٹ $\funimage{\rho}{C}$ اوہناں نقطیاں توں بندا اے جیہڑے گھماؤ $\rho$ نوں $C$ دے ہر نقطے اُتے لاؤن نال ملدے نیں۔ اگلے دو نتیجے وکھاندے نیں کہ ایہ سیٹ دائرے نوں پورا ڈھکدے نیں تے اوہناں وچ کوئی سانجھا نقطہ نہیں۔

\begin{lem}\ollabel{vitalicover}
$\onesphere = \bigcup_{\rho \in \rotationsgroup} \funimage{\rho}{C}$۔
\end{lem}

\begin{proof}
$s \in \onesphere$ مقرر کرو۔ کوئی $r \in C$ اے جیہدے لئی $r \in
\equivrep{s}{\sim}$، یعنی $r \sim s$، یعنی $\rho(r) = s$ کسے $\rho \in \rotationsgroup$ لئی۔
\end{proof}

\begin{lem}\ollabel{vitalinooverlap}
جے $\rho_1 \neq \rho_2$، تاں $\funimage{\rho_1}{C} \cap \funimage{\rho_2}{C} = \emptyset$۔
\end{lem}

\begin{proof}
منّ لو $s \in \funimage{\rho_{1}}{C} \cap \funimage{\rho_{2}}{C}$۔ سو
$s = \rho_{1}(r_{1}) = \rho_{2}(r_{2})$ کجھ $r_{1}, r_{2} \in C$ لئی۔
ایس لئی $\rho^{-1}_2(\rho_1(r_1)) = r_2$، تے
$\comp{\rho_1}{\rho^{-1}_2} \in \rotationsgroup$، سو $r_{1} \sim
r_{2}$۔ سو $r_1 = r_2$، کیوں جے $C$، $\sim$ تحت ہر ہم ارزی جماعت وچوں این ٹھیک اک رکن چُندا اے۔ سو $s = \rho_1(r_1) = \rho_2(r_1)$، تے ایس لئی $\rho_1 = \rho_2$۔
\end{proof}

ہُن اسیں الجبرا دیاں پہلاں والیاں گلاں اپنے دائرے اُتے لاؤندے آں۔

\begin{lem}\ollabel{pseudobanachtarski}
$\onesphere$ دی دو وکھرے سیٹاں، $D_{1}$ تے $D_{2}$، وچ ونڈ اے، ایویں کہ $D_{1}$ نوں گِنن جوگ گنتی دے سیٹاں وچ ونڈیا جا سکدا اے جیہناں نوں گھما کے $\onesphere$ دی اک نقل بنائی جا سکدی اے (تے $D_{2}$ لئی وی ایہو گل اے)۔
\end{lem}

\begin{proof}
\olref{disjointgroup} والے $\rotationsgroup_{1}$ تے $\rotationsgroup_{2}$ ورت کے ایویں لئیے:
\begin{align*}
	D_{1} &= \bigcup_{\rho \in \rotationsgroup_1} \funimage{\rho}{C} &
	D_{2} &= \bigcup_{\rho \in \rotationsgroup_2} \funimage{\rho}{C}
\end{align*}
\olref{vitalicover} مطابق ایہ $\onesphere$ دی ونڈ اے، تے $D_1$ تے $D_2$، \olref{vitalinooverlap} مطابق وکھرے نیں۔ بناوٹ مطابق $D_1$ نوں گِنن جوگ گنتی دے سیٹاں وچ ونڈیا جا سکدا اے، یعنی $\funimage{\rho}{C}$ ہر $\rho \in R_1$ لئی۔ تے ایہناں نوں گھما کے $\onesphere$ دی نقل بنائی جا سکدی اے، کیوں جے $\onesphere = \bigcup_{\rho \in
\rotationsgroup}\funimage{\rho}{C} = \bigcup_{\rho \in
\rotationsgroup_1}\funimage{(\comp{\rho}{\rho})}{C}$، \olref{disjointgroup} تے \olref{vitalicover} مطابق۔ اوہی دلیل $D_2$ لئی چل دی اے۔ \end{proof}\noindent ایس توں ویتالی دا پیراڈوکس جھٹ نکلدا اے۔ کیوں جے اسیں $\onesphere$ دیاں \emph{دو} نقلاں $\onesphere$ توں بنا سکدے آں۔ بس اینوں گِنن جوگ گنتی دے ٹکڑیاں (وکھ وکھ $\funimage{\rho}{C}$ والے سیٹاں) وچ ونڈو تے فیر اوہناں دی شکل قائم رکھدیاں ہلاؤ ($D_1$ دا ہر ٹکڑا بس گھماؤ، تے $D_2$ دے ہر ٹکڑے دی پہلاں تھاں بدلو تے فیر گھماؤ)۔

ثبوت دا طریقہ فیر چیتے کریے۔ اسیں سطح اُتے گھماؤواں دے گروپ بارے الجبرا دیاں کجھ گلاں توں شروع کیتا۔ ایس گروپ نال اسیں $\onesphere$ نوں ہم ارزی جماعتاں وچ ونڈیا۔ فیر چُون ورت کے ہر جماعت وچوں رکن چُنے تے ایس طرح «پیراڈوکس» تک پہنچے۔

باناخ–تارسکی ثابت کرن لئی این اوہی طریقہ ورتدے آں۔ وڈا فرق ایہ اے کہ باناخ–تارسکی دے ثبوت وچ ورتیاں الجبرا دیاں گلاں ویتالی دے پیراڈوکس دے ثبوت والیاں نالوں کِتے گنجل دار نیں۔ پر اوہناں الجبرا دیاں گلاں دا چُون نال کوئی تعلق نہیں۔ اسیں اوہناں دا مختصر خلاصہ دیواں گے۔

باناخ–تارسکی ثابت کرن لئی اسیں \olref{disjointgroup} ورگی اک گل ثابت کرن توں شروع کردے آں۔ کسے وی \emph{آزاد گروپ} نوں چار ٹکڑیاں وچ ونڈیا جا سکدا اے۔ سہج سمجھ مطابق اسیں ایہناں نوں «اگے پچھے ہلا» کے سارے گروپ دیاں دو نقلاں فیر حاصل کر سکدے آں۔\footnote{ایہ گل کہ اسیں \emph{چار} ٹکڑے ورت سکدے آں، \cite{Robinson1947} دی اے۔ اک حالیہ ثبوت لئی \citet[Theorem 5.2]{Wagon2016} ویکھو۔ گروپ دے ٹکڑیاں نوں «ہلاؤن» والی گل آکھدیاں اسیں \citet[p.~3]{Weston2003} دے پچھے چلدے آں۔} فیر اسیں وکھاندے آں کہ $\Real^3$ دے مبدأ دوالے دو خاص گھماؤ ورت کے گھماؤواں دا اک آزاد گروپ، $F$، پیدا کر سکدے آں۔\footnote{\citet[Theorem 2.1]{Wagon2016} ویکھو۔} (ہُن تک چُون نہیں ورتی۔) ہُن اسیں گولے دی سطح دے نقطیاں نوں «ملدے» سمجھدے آں اُدوں تے صرف اُدوں جدوں اک توں دوجا~$F$ دے کسے گھماؤ نال مل سکدا اے۔ فیر اسیں «ملدے» نقطیاں دی ہر ہم ارزی جماعت وچوں این ٹھیک اک نقطہ چُنن لئی \emph{چُون ورتدے آں}۔ $F$ دی اپنی ونڈ نوں گولے دی سطح اُتے لا کے، جیویں \olref{pseudobanachtarski} وچ، اسیں سطح نوں چار ٹکڑیاں وچ ونڈدے آں جیہناں نوں «اگے پچھے ہلا» کے گولے دی سطح دیاں دو نقلاں بنا سکدے آں۔ تے ایس توں ایہ ثابت ہُندا اے \citep{Hausdorff1914}:

\begin{thm}[ہاؤسڈورف دا پیراڈوکس ($\ZFC$ وچ)]
کسے وی گولے دی سطح نوں محدود گنتی دے ٹکڑیاں وچ ونڈیا جا سکدا اے۔ اوہ ٹکڑے فیر (گھما کے تے تھاں بدل کے) جوڑے جا سکدے نیں تاں جو اوس گولے دیاں دو وکھریاں نقلاں بن جان۔
\end{thm}

پورا باناخ–تارسکی قضیہ حاصل کرن لئی (جیہڑا گولے دی سطح دے نال اوہدے اندرلے حصے بارے وی اے) الجبرا دیاں دو کُ ہور ترکیباں چاہیدیاں نیں۔ پر سِدھی گل ایہ اے کہ الجبرا دا اصل کم ہو چکیا اے۔ ایس لئی ویسٹن لکھدا اے:
\begin{quote}
  [\ldots] آزاد گروپاں والا نتیجہ باناخ–تارسکی دے پیراڈوکس دے ثبوت دا \emph{ضروری قدم} اے۔ ایس نظر توں باناخ–تارسکی دا پیراڈوکس $\Real^3$ بارے گھٹ، تے گروپ دی پیچیدگی بارے ودھ اے [یعنی $\Real^3$ وچ تھاں بدلن تے گھماؤن دا گروپ]۔ \cite[p.~16]{Weston2003}
\end{quote}
یعنی اسیں \emph{محدود} گنتی دے ٹکڑیاں والی ونڈ دے سکدے آں (جیویں باناخ–تارسکی وچ)، یا \emph{گِنن جوگ لامتناہی} گنتی دے ٹکڑیاں والی ونڈ (جیویں ویتالی دے پیراڈوکس وچ)، ایہ دو یا تِن جہتاں وچ کم کرن بارے گروپ نظریے دیاں کجھ گلاں اُتے ٹکدا اے۔

منّنا پئے گا کہ ایہ آخری گل \olref[banach]{sec} دے آخر والے لطیفے نوں ذرا خراب کردی اے۔ دو جہتی ہون کرکے «باناخ–تارسکی» نوں گِنن جوگ \emph{لامتناہی} گنتی دے ٹکڑیاں وچ ونڈنا پئے گا، جے اوہ ٹکڑے فیر جوڑ کے «باناخ–تارسکی باناخ–تارسکی» بنانا ہووے۔ لطیفہ ٹھیک کرن لئی تِن جہتاں وچ لکھنا پئے گا۔ اسیں ایہ پڑھن والے لئی مشق وجوں چھڈدے آں۔

اک آخری گل۔ \olref[banach]{sec} وچ اسیں دسیا سی کہ گولے توں ملے «ٹکڑے» \emph{ناپن جوگ} نہیں ہو سکدے، پر نقطیاں دے ایہو جیہے «لامتناہی کھلارے» ہون گے جنہاں دی تصویر ذہن وچ نہیں بن سکدی۔ ایہو گل \olref{pseudobanachtarski} حاصل کرن لئی چُون دی ساڈی ورتوں اُتے وی لاگو ہُندی اے۔ تے ایس سارے دی وضاحت کرنی چنگی اے۔

فیر سانوں کجھ پچھوکڑ دسنا پئے گا (پر ایہ \emph{بس} اک خاکہ اے۔ تسی \emph{ناپ} بارے کسے درسی کتاب دا حصہ ویکھ سکدے او)۔ سیٹ $X$ لئی اک ناپ تعریف کرن دا مطلب اے کہ اک قدر $\mu(E) \in
\Real$ ہر $E$ نوں دِتی جائے، جیہڑا کسے «$\sigma$-الجبرا» وچ اے جیہڑا $X$ اُتے اے۔ ایتھے تفصیل ضروری نہیں، بس فنکشن $\mu$ نوں گِنن جوگ جوڑ دے اصول اُتے پورا اترنا چاہیدا اے۔ وکھرے سیٹاں دی گِنن جوگ گنتی دے اتحاد دا ناپ اوہناں دے وکھ وکھ ناپاں دا جوڑ اے، یعنی
$\mu(\bigcup_{n < \omega} X_n) = \sum_{n < \omega}\mu(X_n)$ جدوں $X_n$ والے سیٹ وکھرے ہون۔ کسے سیٹ نوں «ناپن جوگ نہیں» آکھن دا مطلب اے کہ اوہنوں کوئی ڈھکواں ناپ نہیں دِتا جا سکدا۔ ہُن پہلاں والا $\rotationsgroup$ ورتدے آں:

\begin{cor}[ویتالی]
$\mu$ اک ایہو جیہا ناپ ہووے کہ $\mu(\onesphere) = 1$، تے $\mu(X) = \mu(Y)$ جے $X$ تے $Y$ شکل تے وڈائی وچ برابر نیں۔ تاں
$\funimage{\rho}{C}$ ناپن جوگ نہیں، سارے $\rho \in
\rotationsgroup$ لئی۔
\end{cor}

\begin{proof}
ٹکراؤ کڈھن لئی اُلٹ منّ لو۔ سو $\mu(\funimage{\sigma}{C}) =
r$ ہووے کسے $\sigma \in \rotationsgroup$ تے کسے $r \in \Real$ لئی۔ کسے وی $\rho \in C$ لئی، $\funimage{\rho}{C}$ تے $\funimage{\sigma}{C}$ شکل تے وڈائی وچ برابر نیں، تے ایس لئی $\mu(\funimage{\rho}{C}) = r$ کسے وی $\rho \in
C$ لئی۔ \olref{vitalicover} تے \olref{vitalinooverlap} مطابق، $\onesphere =
\bigcup_{\rho \in \rotationsgroup}\funimage{\rho}{C}$ جوڑی جوڑی کرکے وکھرے سیٹاں دی گِنن جوگ گنتی دا اتحاد اے۔ سو گِنن جوگ جوڑ دا اصول دسدا اے کہ $\mu(\onesphere) = 1$ ہر $\funimage{\rho}{C}$ دے ناپاں دا جوڑ اے، یعنی:
\[
	1 = \mu(\onesphere) = \sum_{\rho \in \rotationsgroup}\mu(\funimage{\rho}{C}) = \sum_{\rho \in \rotationsgroup}r
\]
پر جے $r = 0$، تاں $\sum_{\rho \in \rotationsgroup}r = 0$، تے جے $r
> 0$، تاں $\sum_{\rho \in \rotationsgroup}r = \infty$۔
\end{proof}
\endgroup
% /OLP-0600

\OLEndChapterHook
\endgroup
% /OLP-0592
\endgroup
% /OLP-0529

% OLP-0601 content/methods/methods.tex
\begingroup
\def\olfilename{methods}
\expandafter\def\csname ol@sectioncs\endcsname{section}

\olpart{mth}{طریقے}

\begin{editorial}
  ایس حصے وچ عام گلاں تے کم کرن دے طریقے نیں، خاص کر کے ثبوت دے کئی طریقیاں دی وضاحت جیہناں نال ریاضی توں ہٹ کے ہور مضمون پڑھن والا طالب علم شاید واقف نہ ہووے۔ ہُن ایس وچ اک باب ثبوت لکھن بارے تے اک باب استقرا بارے اے۔ ایہناں لئی ہور حصے تے مشقاں، تے ریاضی دیاں اصطلاحاں بارے اک باب وی منصوبے وچ نیں۔
\end{editorial}

% OLP-0602 content/methods/proofs/proofs.tex
\begingroup
\def\olfilename{proofs}
\expandafter\def\csname ol@sectioncs\endcsname{section}

\olchapter{mth}{prf}{ثبوت}

% OLP-0603 content/methods/proofs/introduction.tex
\begingroup
\def\olfilename{introduction}
\expandafter\def\csname ol@sectioncs\endcsname{section}

\olfileid{mth}{prf}{int}

\olsection{جان پہچان}

منطق دی جان پہچان وچ اپنے تجربے کارن شاید تسی کسے !!a{derivation} نظام نوں سوکھ نال ورت لیندے او، شاید قدرتی اخذ یا فِچ ڈھنگ دا !!{derivation} نظام، یا شاید ثبوت دے رُکھاں والا نظام۔ تہانوں شاید یاد ہووے کہ ایہناں نظاماں وچ تسی ثبوت کیویں کردے سو، یا !!a{formula} ثابت کردے سو یا وکھاؤندے سو کہ دتی ہوئی دلیل جائز اے۔ ایہ کم کرن لئی تسی نظام دے قاعدے لاؤندے سو جدوں تک لوڑیندا آخری نتیجہ نہیں مل جاندا سی۔ منطق \emph{بارے} دلیل دیندیاں وی اسیں گلاں ثابت کردے آں، پر بہتے ویلے اسیں کوئی !!a{derivation} نظام نہیں ورتدے۔ اصل وچ، اسیں جیہڑے ثبوت ویکھدے آں، اوہناں وچوں بہتے انگریزی وچ دتے جاندے نیں (شاید کجھ علامتی زبان وی وچ پا دتی جائے)، سارے دے سارے پہلے درجے دی منطق دی زبان وچ نہیں۔ ایہو جیہے ثبوت بناندیاں پہلاں شاید تہانوں سمجھ نہ آوے۔ میں !!a{derivation} نظام توں بناں کوئی گل کیویں ثابت کراں؟ میں شروع کتھوں کراں؟ مینوں کیویں پتا لگے کہ میرا ثبوت ٹھیک اے؟

ثبوت کرن دی کوشش توں پہلاں ایہ جاننا ضروری اے کہ ثبوت کیہ اے تے کیویں بنائی دا اے۔ جیویں ناں توں ای پتا لگدا اے، \emph{ثبوت} دا کم وکھاؤنا اے کہ کوئی گل سچ اے۔ تسی ایس نوں اک گل بات وانگ سوچ سکدے او۔ کوئی تہانوں پُچھدا اے کہ کوئی گل سچ اے، مثلاً کی دو توں چھڈ کے ہر اوّلی عدد طاق اے۔ جواب وچ «ہاں» آکھنا کافی نہیں۔ اوہ جاننا چاہ سکدا اے کہ \emph{کیوں}۔ ایس صورت وچ تسی اوہنوں اک ثبوت دیو گے۔

روز دی گل بات وچ جواب ول اشارہ کر دینا یا ادھورا جواب دینا کافی ہو سکدا اے۔ پر منطق تے ریاضی وچ سانوں پکا ثبوت چاہیدا اے۔ اسیں وکھاؤنا چاہندے آں کہ گل \emph{کسے وی} شک توں پرے سچ اے۔ ایس دا مطلب اے کہ ساڈے ثبوت دا ہر قدم جائز ہونا چاہیدا اے، تے اوہدا جواز ٹھوس ہووے۔ یعنی جیہڑی گل تسی منّ کے ورتدے او، اوہ قضیے دے بیان وچ واقعی منّی گئی ہووے، تسی تعریفاں ٹھیک طرح لاؤ، تے جیہڑیاں دلیلاں نوں جواز بناؤندے او اوہ ٹھیک استنباط ہون، وغیرہ۔

عام طور اُتے اسیں کوئی بیان ثابت کردے آں۔ جیہڑے بیان ثابت کردے آں اوہناں دے کئی ناں نیں، قضیہ، وڈا قضیہ، معاون قضیہ، یا نتیجہ۔ قضیہ اک بنیادی، ثابت کرن جوگا بیان اے۔ ایہ درج کرن جوگا اے، پر شاید خاص ڈونگھا نہ ہووے یا بہت واری ورتیا نہ جاوے۔ وڈا قضیہ اک خاص اہم قضیہ اے۔ اوہدا ثبوت اکثر کئی قدماں وچ ونڈیا ہُندا اے، تے کدی اوہدا ناں اوہنوں پہلاں ثابت کرن والے دے ناں اُتے ہُندا اے (جیویں کانتور دا قضیہ، یا لووَن ہائم–سکولم دا قضیہ)، یا جیہڑی گل اوہ کردا اے اوہ دے ناں اُتے (جیویں تکمیل دا قضیہ)۔ معاون قضیہ اوہ قضیہ اے جیہڑا کسے ہور اہم نتیجے دے ثبوت وچ ورتیا جاندا اے۔ پر ایس گل نال بھلیکھا پے سکدا اے۔ کدی معاون قضیے اپنے آپ وچ اہم نتیجے ہُندے نیں تے پیش کرن والے دے ناں اُتے وی جاندے نیں (جیویں زورن دا معاون قضیہ)۔ نتیجہ اوہ گل اے جیہڑی کسے ہور نتیجے توں سوکھ نال نکلدی اے۔

جیہڑا بیان ثابت کرنا ہووے، اوہدے وچ اکثر منّیاں ہوئیاں شرطاں ہُندیاں نیں جیہڑیاں صاف کردیاں نیں کہ اسیں کیہڑی قسم دیاں شےواں بارے کوئی گل ثابت کردے آں۔ اوہ ایویں شروع ہو سکدا اے، «$!A$، $!B \lif !C$ دی صورت دا !!a{formula} ہووے» یا «منّ لو $\Gamma \Proves !A$»، یا ایہو جیہی کوئی گل۔ ایہ قضیے، وڈے قضیے، یا معاون قضیے دیاں \emph{مفروضہ شرطاں} نیں، تے تسی اپنے ثبوت وچ ایہناں نوں سچ منّ سکدے او۔ ایہ حد لاؤندیاں نیں کہ اسیں کی ثابت کردے آں، تے جنہاں شےواں دی گل کردے آں اوہناں دے کجھ ناں وی دیندیاں نیں۔ مثلاً جے تہاڈا قضیہ ایویں شروع ہُندا اے، «$!A$، $!B \lif !C$ دی صورت دا !!a{formula} ہووے»، تاں تسی بس اک خاص قسم دے سارے فارمولیاں (یعنی شرطیاں) بارے کوئی گل ثابت کردے او۔ تے ایہ سمجھیا جاندا اے کہ $!B \lif !C$ کوئی وی شرطیہ اے جیہدی گل تہاڈا ثبوت کرے گا۔
\endgroup
% /OLP-0603

% OLP-0604 content/methods/proofs/starting-proofs.tex
\begingroup
\def\olfilename{starting-proofs}
\expandafter\def\csname ol@sectioncs\endcsname{section}

\olfileid{mth}{prf}{str}

\olsection{ثبوت شروع کرنا}

پر شروع کتھوں کریے؟

تہانوں کوئی گل ثابت کرن لئی دتی گئی اے، سو ثبوت وچ سبھ توں آخر اوسے گل دا ذکر ہونا چاہیدا اے (بے شک تسی شروع وچ \emph{دس سکدے او} کہ تسی ایہ گل ثابت کرو گے، پر اینوں منّی ہوئی شرط وجوں ورتنا ٹھیک نہیں)۔ جو کجھ تسی ثابت کرنا چاہندے او، اوہ اک نویں ورق دے تھلّے لکھ لو۔ ایویں تہاڈا مقصد تہاڈی نظر وچ رہے گا۔

اگے، ہو سکدا اے تہاڈے کول کجھ منّن والیاں شرطاں ہون جیہڑیاں تسی ورت سکدے او (جدوں اسیں اگلے حصے وچ تہاڈے ثبوت دی \emph{قسم} دی گل کراں گے تاں ایہ ہور صاف ہو جائے گا)۔ ایہناں نوں ورق دے اُتے لکھو تے صاف دسو کہ ایہ منّیاں ہوئیاں شرطاں نیں (یعنی جے تسی $p$ منّ رہے او، تاں «$p$ منّ لو»، یا «فرض کرو کہ $p$» لکھو)۔ آخر، ہو سکدا اے سوال وچ کجھ تعریفاں ہون جیہڑیاں تہانوں جانن دی لوڑ اے۔ شاید تہانوں اک خاص تعریف ورتن لئی آکھیا جائے، یا تہاڈیاں منّیاں شرطاں یا لوڑیندے نتیجے وچ کئی تعریفاں ہون۔ \emph{ایہناں نوں لکھ لو تے پکا کر لو کہ تہانوں ایہناں دا مطلب سمجھ آ گیا اے۔}

تسی ثبوت دی بناوٹ کیویں رکھدے او، ایہ سوال دی صورت اُتے وی ٹکدا اے۔ اگلا حصہ دسدا اے کہ جملے دی قسم مطابق ثبوت دی بناوٹ کیویں رکھنی اے۔
\endgroup
% /OLP-0604

% OLP-0605 content/methods/proofs/using-definitions.tex
\begingroup
\def\olfilename{using-definitions}
\expandafter\def\csname ol@sectioncs\endcsname{section}

\olfileid{mod}{prf}{def}

\olsection{تعریفاں ورتنیاں}

اسیں دسیا سی کہ ثبوت وچ جیہڑیاں تعریفاں ورتیاں جا سکدیاں نیں، تہانوں اوہناں دا پتا ہونا چاہیدا اے، تے تسی اوہناں نوں ٹھیک طرح لا سکو۔ ایہ بڑی ضروری گل اے، تے ایہنوں ذرا ہور تفصیل نال ویکھنا چنگا اے۔ تعریفاں نال اسیں خاصیتاں تے تعلقاں نوں مختصر لکھدے آں تاں جو اوہناں دی گل گھٹ لفظاں وچ کر سکیے۔ جیہڑی مختصر عبارت پیش کیتی جاندی اے، اوہنوں \emph{تعریف ہون والی عبارت} آکھدے نیں، تے جیہڑی عبارت نوں اوہ مختصر کردی اے، اوہنوں \emph{تعریف دین والی عبارت}۔ ثبوتاں وچ سانوں اکثر مُڑ ویکھنا پیندا اے کہ تعریف ہون والی عبارت کیویں پیش کیتی گئی سی، کیوں جے ثبوت اگے ودھاؤن لئی سانوں تعریف دین والی عبارت دی منطقی بناوٹ ورتنی پیندی اے۔ ایہ اوہ پورا بیان اے جیہنوں تعریف ہون والی اصطلاح مختصر کردی اے۔ تعریفاں کھولھ کے تسی اوہ تھاں لبھ لیندے او جتھے منطقی دلیل کم کردی اے۔

اسیں اک مثال توں شروع کراں گے۔ منّ لو تسی ہِیٹھلی گل ثابت کرنا چاہندے او:

\begin{prop}
کسے وی سیٹاں $A$ تے $B$ لئی، $A \cup B = B \cup A$۔
\end{prop}

ثبوت شروع کرن لئی وی سانوں پتا ہونا چاہیدا اے کہ دو سیٹاں دے اکّو ہون دا مطلب کیہ اے۔ یعنی سانوں پتا ہونا چاہیدا اے کہ اوس مساوات وچ «$=$» دا سیٹاں لئی کیہ مطلب اے۔ تعریف مطابق سیٹ اُدوں اکّو ہُندے نیں جدوں اوہناں دے !!{element}s اکّو ہون۔ سو سانوں ہِیٹھلی تعریف کھولھنی پئے گی۔

\begin{defn}
سیٹ $A$ تے $B$ \emph{اکّو} نیں، $A = B$، اُدوں تے صرف اُدوں جدوں ہر !!{element}، جیہڑا~$A$ دا اے، !!a{element}~$B$ دا وی اے، تے اُلٹے پاسے وی۔
\end{defn}

ایس تعریف وچ $A$ تے~$B$ کسے وی سیٹاں دی تھاں رکھے نشان نیں۔ ایہ جیہڑی عبارت دی تعریف دیندی اے، یعنی \emph{تعریف ہون والی عبارت}، اوہ «$A = B$» اے۔ ایہ اوہ شرط دَس کے تعریف دیندی اے جیہدے ہون اُتے $A = B$ سچ اے۔ ایہ شرط، «ہر !!{element}، جیہڑا~$A$ دا اے، !!a{element}~$B$ دا وی اے، تے اُلٹے پاسے وی»، \emph{تعریف دین والی عبارت} اے۔\footnote{ایس خاص صورت وچ، تے ایہ گل بھلیکھا پا سکدی اے، جدوں $A = B$، تاں سیٹ $A$ تے $B$ بس اکّو سیٹ نیں، بھاویں اسیں کھبّے تے سجّے پاسے اوہدے لئی وکھرے اکھر ورتدے آں۔ پر اوہ سیٹ کیہڑے طریقیاں نال دَسیا گیا اے، اوہ طریقے وکھرے ہو سکدے نیں، تے ایس کرکے تعریف محض ظاہر گل نہیں۔} تعریف دسدی اے کہ $A = B$ سچ اے جے، تے صرف جے، ایہ شرط پوری ہووے (اسیں ایہنوں مختصر «اُدوں تے صرف اُدوں جدوں» لکھدے آں)۔

تعریف لاؤندیاں تہانوں تعریف والے $A$ تے $B$ نوں اوس صورت نال ملاؤنا پیندا اے جیہڑی تسی ویکھ رہے او۔ ساڈی صورت وچ ایہدا مطلب اے کہ $A \cup B = B \cup A$ سچ ہون لئی، ہر $z \in A \cup
B$ نوں $B \cup A$ وچ وی ہونا پئے گا، تے اُلٹے پاسے وی۔ قضیے وچ عبارت $A \cup
B$ تعریف والے~$A$ دی تھاں اے، تے $B
\cup A$ تعریف والے~$B$ دی تھاں۔ کیوں جے $A$ تے $B$ تعریف وچ وی تے اوہ قضیہ جیہڑا اسیں ثابت کردے آں اوہدے بیان وچ وی ورتے گئے نیں، پر وکھ وکھ طرح، تہانوں دھیان رکھنا پئے گا کہ دونیں ورتواں رَل نہ جان۔ مثلاً ایہ سوچنا غلط ہووے گا کہ اسیں ہر !!{element} جیہڑا~$A$ دا اے، اوہنوں !!a{element}~$B$ دا، تے اُلٹے پاسے وی، وکھا کے قضیہ ثابت کر سکدے آں۔ ایس نال $A = B$ ثابت ہووے گا، $A \cup B = B \cup A$ نہیں۔ (نالے $A$ تے $B$ کوئی وی دو سیٹ ہو سکدے نیں، سو تہاڈی دلیل بہتی اگے نہیں جائے گی، کیوں جے جے $A$ تے~$B$ بارے کجھ وی منّیا نہیں گیا تاں اوہ وکھرے سیٹ ہو سکدے نیں۔)

ثبوت وچ اسیں سیٹ نظریے دے اتحاد ورگے خیالاں نال کم کردے آں۔ سو ایہ سمجھن لئی کہ ثبوت اگے کیویں ودھے، سانوں نشان $\cup$ دا مطلب وی پتا ہونا چاہیدا اے۔ تے کدی کدی اک تعریف کھولھن نال اگوں ہور تعریفاں کھولھنیاں پیندیاں نیں۔ مثلاً $A \cup B$ دی تعریف $\Setabs{z}{z \in A
  \text{ یا } z \in B}$ اے۔ سو جے تسی $x \in A \cup
B$ ثابت کرنا چاہندے او، تاں $\cup$ دی تعریف کھولھن نال پتا لگدا اے کہ تہانوں $x \in \Setabs{z}{z \in A \text{ یا } z \in B}$ ثابت کرنا پئے گا۔ ہُن ایہ وی چیتے رکھنا پئے گا کہ $x \in \Setabs{z}{\dots z\dots}$ اُدوں تے صرف اُدوں جدوں $\dots
x\dots$۔ سو $\Setabs{z}{\dots z \dots}$ والی لکھت دی تعریف ہور کھولھیے تاں تہانوں ایہ وکھانا اے کہ $x \in
A$ یا $x \in B$۔ سو «ہر !!{element} جیہڑا $A \cup B$ دا اے، اوہ !!a{element} $B \cup A$ دا وی اے» دا اصل مطلب اے «ہر $x$ لئی، جے $x \in
A$ یا $x \in B$، تاں $x \in B$ یا $x \in A$۔» جے اسیں قضیے وچ تعریفاں پوریاں کھولھیے، تاں پتا لگدا اے کہ سانوں ایہ وکھانا اے:

\begin{prop}
کسے وی سیٹاں $A$ تے $B$ لئی، (ا) ہر $x$ لئی، جے $x \in A$ یا $x \in
B$، تاں $x \in B$ یا $x \in A$، تے (ب) ہر $x$ لئی، جے $x \in B$
یا $x \in A$، تاں $x \in A$ یا $x \in B$۔
\end{prop}

ضروری گل ایہ اے کہ تعریفاں کھولھنا ثبوت بناؤن دا لازمی حصہ اے۔ ایہ کم ٹھیک طرح کرنا کدی اوکھا ہُندا اے۔ تہانوں تعریف دے متغیراں تے جیہڑا دعویٰ تسی ثابت کردے او اوہدیاں اصطلاحاں وچ فرق رکھنا تے اوہناں نوں ٹھیک طرح ملاؤنا پیندا اے۔ کامیاب ہون لئی تہانوں پتا ہونا چاہیدا اے کہ سوال کی پُچھدا اے تے اوہدے وچ ورتی ہر اصطلاح دا مطلب کیہ اے۔ تہانوں اکثر اک توں ودھ تعریف کھولھنی پئے گی۔ ہِیٹھلے ثبوتاں ورگے سادے ثبوتاں وچ حل تعریفاں توں لگ بھگ جھٹ نکل آؤندا اے۔ بے شک ہر واری ایہ اینا سوکھا نہیں ہووے گا۔

\begin{prob}
منّ لو تہانوں $A \cap B \neq \emptyset$ ثابت کرن لئی آکھیا گیا اے۔ ایتھے آئیاں ساریاں تعریفاں کھولھو، یعنی ایہ گل ایویں فیر آکھو کہ اوہدے وچ «$\cap$»، «=»، یا «$\emptyset$» دا ذکر نہ ہووے۔
\end{prob}
\endgroup
% /OLP-0605

% OLP-0606 content/methods/proofs/inference-patterns.tex
\begingroup
\def\olfilename{inference-patterns}
\expandafter\def\csname ol@sectioncs\endcsname{section}

\olfileid{mth}{prf}{pat}

\olsection{استنباط دے ڈھنگ}

ثبوت وکھ وکھ استنباطاں توں بندے نیں۔ جدوں اسیں کوئی استنباط کردے آں، تاں عام طور اُتے «سو»، «ایس لئی»، یا «ایس کرکے» ورگا لفظ لا کے دس دیندے آں۔ استنباط اکثر اک یا دو اوہناں گلاں اُتے ٹکدا اے جیہڑیاں ثبوت وچ پہلاں ای ساڈے کول نیں۔ ایہ کوئی منّی ہوئی گل ہو سکدی اے یا کوئی نتیجہ جیہڑا اسیں پہلاں کڈھیا ہووے۔ گل صاف رکھن لئی اسیں ایہناں گلاں نوں نشان لا سکدے آں، تے استنباط کردیاں دس سکدے آں کہ اوہدے وچ کیہڑے ہور بیان ورتدے آں۔ استنباط وچ اکثر ایہ وضاحت وی ہُندی اے کہ ساڈا نواں نتیجہ پہلاں والیاں گلاں توں \emph{کیوں} نکلدا اے۔ استنباط دے کجھ عام ڈھنگ ثبوتاں وچ بہت واری ورتے جاندے نیں۔ ہِیٹھ اسیں کجھ ویکھاں گے۔ کجھ ڈھنگ، جیویں استقرا نال ثبوت، ذرا گنجل دار نیں (اوہناں دی گل اگے ہووے گی)۔

اسیں استنباط دے اک ڈھنگ دی گل پہلاں ای کیتی اے، یعنی تعریف کھولھنی یا لاؤنی۔ جدوں اسیں اک تعریف کھولھدے آں، تاں جیہڑی گل وچ تعریف ہون والی عبارت آؤندی اے، اوہنوں تعریف دین والی عبارت نال فیر آکھدے آں۔ مثلاً منّ لو اسیں ثبوت وچ پہلاں ای ثابت کر چکے آں کہ $D = E$ (ا)۔ تاں اسیں سیٹاں لئی $=$ دی تعریف لا کے ایہ نتیجہ کڈھ سکدے آں۔ «سو (ا) توں تعریف مطابق ہر !!{element} جیہڑا~$D$ دا اے، اوہ !!a{element}~$E$ دا وی اے، تے اُلٹے پاسے وی۔»

ایہ گل تھوڑا بھلیکھا پا سکدی اے کہ اسیں اکثر استنباط دا جواز اوہ کردیاں نہیں، اوہتوں پہلاں لکھدے آں۔ منّ لو اسیں ہن تک $D = E$ ثابت نہیں کیتا، پر کرنا چاہندے آں۔ جے $D = E$ ساڈا لوڑیندا نتیجہ اے، تاں ایہ مقصد وی تعریف لا کے فیر دس سکدے آں۔ $D = E$ ثابت کرن لئی سانوں ثابت کرنا پئے گا کہ ہر !!{element} جیہڑا~$D$ دا اے، اوہ !!a{element}~$E$ دا وی اے، تے اُلٹے پاسے وی۔ سو ساڈے ثبوت دی صورت ایہ ہووے گی۔ (ا) ثابت کرو کہ ہر !!{element} جیہڑا~$D$ دا اے، اوہ !!a{element}~$E$ دا وی اے۔ (ب) ہر !!{element} جیہڑا~$E$ دا اے، اوہ !!a{element}~$D$ دا وی اے۔ (ج) سو (ا) تے (ب) توں، $=$ دی تعریف مطابق، $D = E$۔ پر اسیں عام طور اُتے ایویں نہیں لکھدے۔ اسیں ایہو جیہی گل لکھ سکدے آں۔
\begin{quote}
اسیں $D = E$ وکھاؤنا چاہندے آں۔~$=$ دی تعریف مطابق ایہ وکھاؤنا اے کہ ہر !!{element} جیہڑا~$D$ دا اے، اوہ !!a{element}~$E$ دا وی اے، تے اُلٹے پاسے وی۔

(ا) \dots (ثبوت کہ ہر !!{element} جیہڑا~$D$ دا اے، اوہ !!a{element}~$E$ دا وی اے) \dots

(ب) \dots (ثبوت کہ ہر !!{element} جیہڑا~$E$ دا اے، اوہ !!a{element}~$D$ دا وی اے) \dots
\end{quote}

\subsection{عطف ورتنا}

شاید استنباط دا سبھ توں سوکھا ڈھنگ ایہ اے کہ عطف دے دو جزاں وچوں اک نوں نتیجے وجوں کڈھ لئیے۔ یعنی جے اسیں منّیا یا پہلاں ای ثابت کیتا اے کہ $p$ تے~$q$، تاں اسیں~$p$ کڈھ سکدے آں (تے~$q$ وی)۔ ایہ اینا بنیادی استنباط اے کہ اکثر ایس دا ذکر وی نہیں ہُندا۔ مثلاً $D = E$ دی تعریف کھولھ کے اسیں ثابت کر لیا کہ ہر !!{element} جیہڑا~$D$ دا اے، اوہ !!a{element}~$E$ دا وی اے، تے اُلٹے پاسے وی۔ ایس توں نتیجہ نکلدا اے کہ ہر !!{element} جیہڑا~$E$ دا اے، اوہ !!a{element}~$D$ دا وی اے (ایہو «اُلٹے پاسے وی» والا حصہ اے)۔

\subsection{عطف ثابت کرنا}

کدی جیہڑی گل تہانوں ثابت کرن لئی آکھن گے، اوہ عطف دی صورت وچ ہووے گی، «$p$ تے $q$ ثابت کرو»۔ ایس صورت وچ بس دو کم کرنے نیں۔ $p$ ثابت کرو تے فیر $q$ ثابت کرو۔ تسی ثبوت دو حصیاں وچ ونڈ سکدے او تے صفائی لئی اوہناں نوں نشان لا سکدے او۔ پہلاں نوٹ بناندیاں ورق دے اُتے «(۱) $p$ ثابت کرو» تے ورق دے وچکار «(۲) $q$ ثابت کرو» لکھ سکدے او۔ (ہو سکدا اے تہانوں عطف ثابت کرن لئی صاف نہ آکھیا ہووے پر ثبوت لئی عطف ثابت کرنا پئے۔ مثلاً جے $D = E$ ثابت کرن لئی آکھیا اے، تاں~$=$ دی تعریف کھولھ کے تہانوں ثابت کرنا پئے گا کہ ہر !!{element} جیہڑا~$D$ دا اے، اوہ !!a{element}~$E$ دا وی اے، \emph{تے} ہر !!{element} جیہڑا~$E$ دا اے، اوہ !!a{element}~$D$ دا وی اے۔)

\subsection{فصل ثابت کرنا}

جدوں جیہڑی گل تسی ثابت کردے او اوہ فصل دی صورت وچ ہووے (یعنی «$p$ یا $q$» دی صورت دا بیان ہووے)، تاں فصل دے جزاں وچوں اک نوں سچ وکھاؤنا کافی اے۔ پر ایہ لگ بھگ کدی نہیں ہُندا کہ کوئی اک جز تہاڈے قضیے دیاں منّیاں شرطاں توں سِدھا ای نکل آوے۔ اکثر قضیے دیاں منّیاں شرطاں آپ فصل دی صورت وچ ہُندیاں نیں، یا تسی وکھاؤندے او کہ اک قسم دیاں ساریاں شےواں وچ دو خاصیتاں وچوں اک اے، پر کجھ شےواں وچ پہلی تے ہورناں وچ دوجی خاصیت اے۔ ایتھے صورتاں وکھ کر کے ثبوت دین دا ڈھنگ کم آؤندا اے (ہِیٹھ ویکھو)۔

\subsection{شرطیہ ثبوت}

کئی قضیے جیہڑے تسی ویکھو گے شرطیہ صورت وچ ہون گے (یعنی وکھاؤ کہ جے $p$ پورا ہووے تاں $q$ وی سچ اے)۔ ایہناں صورتاں وچ ثبوت دی بناوٹ سوکھی اے۔ بس شرطیے دا مقدم منّ لو (ایس صورت وچ $p$) تے اوہتوں نتیجہ~$q$ ثابت کرو۔ سو جے قضیہ ایہ آکھدا اے، «جے $p$ تاں $q$»، تاں ثبوت «$p$ منّ لو» توں شروع کرو تے آخر تک~$q$ ثابت کر لو۔

شرطیہ کئی طرحاں نال دَسے جا سکدے نیں۔ سو «جے $p$ تاں $q$» دی تھاں قضیہ آکھ سکدا اے «$p$ صرف اُدوں جے $q$»، «$q$ جے $p$»، یا «$q$، بشرطیکہ $p$»۔ ایہناں ساریاں دا مطلب اکّو اے، تے ہر واری $p$ منّ کے اوس منّی ہوئی گل توں~$q$ ثابت کرنا پئے گا۔ چیتے کرو کہ دوپاسا شرطیہ («$p$ اُدوں تے صرف اُدوں جدوں $q$») اصل وچ دو شرطیاں دا جوڑ اے، جے $p$ تاں $q$، تے جے $q$ تاں~$p$۔ سو تہانوں بس دو شرطیہ ثبوت دینے پَین گے، اک پہلے شرطیے لئی تے اک دوجے لئی۔ پر کدی «اُدوں تے صرف اُدوں جدوں» والا بیان ایویں وی ثابت ہو سکدا اے کہ ایس طرح دے کئی ہور بیان اک لڑی وچ جوڑ دیو، «$p$» توں شروع کرو تے «$q$» اُتے مُکو۔ پر اوس صورت وچ پکا کر لو کہ ہر قدم واقعی «اُدوں تے صرف اُدوں جدوں» والا اے۔

\subsection{ساریاں لئی دعوے}

ساریاں لئی کوئی دعویٰ ورتنا سوکھا اے۔ جے کوئی گل ہر شے لئی سچ اے، تاں ہر خاص شے لئی وی سچ اے۔ سو مثلاً جے تہاڈے ثبوت دی مفروضہ شرط $A \subseteq B$ اے، تاں (تعریف~$\subseteq$ کھولھ کے) ایہدا مطلب اے کہ ہر $x \in A$ لئی، $x \in B$۔ سو جے تہانوں پہلاں پتا اے کہ $z \in A$، تاں نتیجہ کڈھ سکدے او کہ $z \in B$۔

ساریاں لئی کوئی دعویٰ ثابت کرنا ذرا اوکھا لگ سکدا اے۔ عام طور اُتے ایہ بیان ایہو جیہی صورت وچ ہُندے نیں۔ «جے $x$ وچ~$P$ اے، تاں اوہدے وچ~$Q$ وی اے» یا «سارے $P$، $Q$ نیں»۔ بے شک ہو سکدا اے ایہ عین ایس صورت وچ نہ ہووے، تے عین کیہڑی گل ثابت کرن لئی آکھی گئی اے ایہ سمجھنا تھوڑی مشق منگدا اے۔ پر سانوں اکثر ثابت کرنا پیندا اے کہ اک خاصیت والیاں ساریاں شےواں وچ اک ہور خاصیت وی اے۔

ساریاں لئی دعویٰ ثابت کرن دا طریقہ ایہ اے کہ اوہناں شےواں لئی ناں یا متغیر پیش کرو جیہناں وچ پہلی خاصیت اے، تے فیر وکھاؤ کہ اوہناں وچ دوجی خاصیت وی اے۔ ایہ گل ایویں آکھ سکدے آں کہ \emph{سارے}~$P$ لئی کوئی گل ثابت کرن واسطے اک \emph{من مانے}~$P$ لئی اوہ گل ثابت کرنی پئے گی۔ جیہڑا ناں پیش کیتا اے، اوہ اک من مانے~$P$ دا ناں اے۔ من مانیاں شےواں دے ناں لئی اسیں عام طور اُتے اکّو اکّو اکھر ورتدے آں، تے اکھراں دی چُون اکثر رواج مطابق ہُندی اے۔ مثلاً قدرتی عدداں لئی $n$، !!{formula}s لئی $!A$، سیٹاں لئی $A$، فنکشناں لئی $f$، وغیرہ۔

کم ایہ اے کہ سارے ثبوت وچ عمومیت قائم رکھو۔ شروع وچ منّ لو کہ اک من مانی شے («$x$») وچ خاصیت~$P$ اے، تے وکھاؤ (بس تعریفاں یا اوہناں گلاں توں جیہڑیاں تہانوں منّن دی کھُل اے) کہ $x$ وچ خاصیت~$Q$ اے۔ کیوں جے تسی $x$ بارے خاص طور اُتے کجھ نہیں دسیا، صرف ایہ کہ اوہدے وچ خاصیت $P$ اے، تسی آکھ سکدے او کہ ہر $P$ والی شے وچ خاصیت~$Q$ اے۔ مختصر گل ایہ اے کہ $x$ خاصیت~$P$ والیاں \emph{ساریاں} شےواں دی نمائندگی کردا اے۔

\begin{prop}
  سارے سیٹاں $A$ تے $B$ لئی، $A \subseteq A \cup B$۔
\end{prop}

\begin{proof}
  $A$ تے $B$ من مانے سیٹ ہون۔ اسیں وکھاؤنا چاہندے آں کہ $A
  \subseteq A \cup B$۔ $\subseteq$ دی تعریف مطابق ایہ وکھاؤنا اے کہ ہر $x$ لئی، جے $x \in A$ تاں $x \in A \cup B$۔ سو $x \in A$، سیٹ~$A$ دا اک من مانا !!{element} ہووے۔ سانوں وکھاؤنا اے کہ $x \in A
  \cup B$۔ کیوں جے $x \in A$، ایس لئی $x \in A$ یا $x \in B$۔ سو $x \in
  \Setabs{x}{x \in A \lor x \in B}$۔ پر $\cup
  $ دی تعریف مطابق ایہدا مطلب $x \in A \cup B$ اے۔
\end{proof}

\subsection{صورتاں وکھ کر کے ثبوت}

منّ لو فصل تہاڈی منّی ہوئی گل اے یا پہلاں ای ثابت کیتا نتیجہ اے، یعنی تسی منّیا یا ثابت کیتا اے کہ $p$ یا $q$ سچ اے۔ تسی $r$ ثابت کرنا چاہندے او۔ ایہ دو قدماں وچ کرو۔ پہلاں $p$ نوں سچ منّ کے~$r$ ثابت کرو۔ فیر $q$ نوں سچ منّ کے~$r$ فیر ثابت کرو۔ ایہ طریقہ ایس لئی چلدا اے کہ اسیں منّیا اے یا سانوں پتا اے کہ دوناں وچوں اک صورت پوری ہُندی اے۔ دو قدم ثابت کردے نیں کہ دوہاں وچوں کوئی وی صورت~$r$ دے سچ ہون لئی کافی اے۔ (جے دوویں سچ نیں، تاں $r$~دے سچ ہون دی اک دی بجائے دو وجوہاں نیں۔ $r$~دا سچ ہونا، $p$ تے~$q$ دوویں منّ کے وکھرا ثابت کرن دی لوڑ نہیں۔) اپنا طریقہ دسّن لئی اسیں آکھدے آں کہ «اسیں صورتاں وکھ کردے آں»۔ مثلاً منّ لو سانوں پتا اے کہ $x \in B \cup C$۔ $B \cup C$ دی تعریف $\Setabs{x}{x \in B \text{ یا } x \in C}$ اے۔ دوجے لفظاں وچ تعریف مطابق $x \in B$ یا $x \in C$۔ ایس توں $x \in A$ ثابت کرن لئی پہلاں $x \in B$ منّو، تے ایس منّی ہوئی گل توں $x \in A$ ثابت کرو۔ فیر $x \in C$ منّو، تے ایس توں فیر $x \in A$ ثابت کرو۔ منّی ہوئی گل تھلّے «اسیں صورتاں وکھ کردے آں» لکھو، اوہدے تھلّے «صورت (۱) $x \in B$»، تے ورق دے وچکار «صورت (۲) $x \in C$» لکھو۔ فیر ورق دا اُتلا تے تھلّلا ادھ بھرو۔

صورتاں وکھ کر کے ثبوت اوہدوں خاص کم آؤندا اے جدوں ثابت کرن والی گل آپ فصل دی صورت وچ ہووے۔ اک سوکھی مثال ایہ اے۔

\begin{prop}
منّ لو $B \subseteq D$ تے $C \subseteq E$۔ تاں $B \cup C \subseteq
D \cup E$۔
\end{prop}

\begin{proof}
  منّ لو (ا) $B \subseteq D$ تے (ب) $C \subseteq E$۔ تعریف مطابق کوئی وی $x \in B$، $\in D$ وی اے (ج)، تے کوئی وی $x \in C$، $\in E$ وی اے (د)۔ $B \cup C \subseteq D \cup E$ وکھاؤن لئی سانوں وکھاؤنا پئے گا کہ جے $x \in B \cup C$ تاں $x \in D \cup E$ ($\subseteq$ دی تعریف مطابق)۔ $x \in B \cup C$ اُدوں تے صرف اُدوں جدوں $x \in B$ یا $x \in
  C$ (تعریف~$\cup$ مطابق)۔ ایویں ای $x \in D \cup E$ اُدوں تے صرف اُدوں جدوں $x \in
  D$ یا $x \in E$۔ سو سانوں ایہ وکھاؤنا اے کہ کسے وی $x$ لئی، جے $x \in B$ یا $x \in C$، تاں $x \in D$ یا $x \in E$۔

  \begin{quote}
  ہن تک اسیں صرف تعریفاں کھولھیاں نیں!{} اسیں اپنا قضیہ $\subseteq$ تے $\cup$ توں بناں فیر بیان کیتا اے، تے ہن ساریاں لئی اک شرطیہ دعویٰ ثابت کرنا رہ گیا اے۔ اُتلی گل مطابق، ایہ کم کرن لئی $x$ نوں اک ایہو جیہی شے منّو جیہدے بارے اسیں «جے» والا حصہ سچ منّدے آں۔ اسیں اگے وکھاواں گے کہ «تاں» والا حصہ وی سچ اے۔ دوجے لفظاں وچ، اسیں منّاں گے کہ $x
  \in B$ یا $x \in C$، تے وکھاواں گے کہ $x \in D$ یا $x \in
  E$۔\footnote{ایہ پیرا بس سمجھاؤندا اے کہ اسیں کی کردے آں۔ ایہ ثبوت دا حصہ نہیں، تے اپنے ثبوت لکھدیاں تہانوں اینی تفصیل وچ جان دی لوڑ نہیں۔}
  \end{quote}

  منّ لو $x \in B$ یا $x \in C$۔ سانوں وکھاؤنا اے کہ $x \in D$ یا $x \in E$۔ اسیں صورتاں وکھ کردے آں۔

  صورت ۱۔ $x \in B$۔ (ج) مطابق $x \in D$۔ سو $x \in D$ یا $x \in
  E$۔ (ایتھے اسیں اوہ استنباط کیتا اے جیہدی گل پچھلے ذیلی حصے وچ کیتی سی!)

  صورت ۲۔ $x \in C$۔ (د) مطابق $x \in E$۔ سو $x \in D$ یا $x \in E$۔
 \end{proof}

\subsection{موجودگی دا دعویٰ ثابت کرنا}

جدوں موجودگی دا دعویٰ ثابت کرن لئی آکھیا جائے، تاں سوال عام طور اُتے ایس صورت وچ ہووے گا۔ «ثابت کرو کہ اک~$x$ اے جیہدے لئی $\dots x \dots$»، یعنی کوئی شے اے جیہدے وچ «$\dots x
\dots$» نال دسی خاصیت اے۔ ایس صورت وچ تہانوں کوئی ڈھکواں شے لبھنی تے وکھاؤنا پئے گا کہ اوہدے وچ لوڑیندی خاصیت اے۔ گل سِدھی لگدی اے، پر ایس قسم دا ثبوت اوکھا ہو سکدا اے۔ عام طور اُتے اک شے \emph{بنانی} یا \emph{تعریف کرنی} تے ثابت کرنا پیندا اے کہ ایویں تعریف کیتی شے وچ لوڑیندی خاصیت اے۔ ٹھیک شے لبھنی اوکھی ہو سکدی اے، اوہدے وچ لوڑیندی خاصیت ثابت کرنی اوکھی ہو سکدی اے، تے کدی ایہ وکھاؤنا وی اوکھا ہو سکدا اے کہ تسی واقعی کوئی شے تعریف کر لئی اے!{}

عام طور اُتے تسی شے دس کے ایہ کم لکھو گے، مثلاً «$x$، \dots ہووے» (جتھے \dots{} دسدا اے کہ تہاڈے ذہن وچ کیہڑی شے اے)۔ ہو سکدا اے پہلاں ثابت کرنا پئے کہ $\dots$ واقعی اک موجود شے دا بیان اے، تے فیر وکھاؤ کہ $x$ وچ خاصیت~$Q$ اے۔ اک سوکھی مثال ایہ اے۔

\begin{prop}
  منّ لو $x \in B$۔ تاں اک~$A$ اے جیہدے لئی $A \subseteq
  B$ تے $A \neq \emptyset$۔
\end{prop}

\begin{proof}
  $x \in B$ منّ لو۔ $A = \{x\}$ ہووے۔
  \begin{quote}
    ایتھے اسیں سیٹ~$A$ دے !!{element}s گِنوا کے اوہدی تعریف دتی اے۔ کیوں جے اسیں منّیا اے کہ $x$ اک شے اے، تے اسیں !!{element}s گِنوا کے ہمیشہ اک سیٹ بنا سکدے آں، سانوں ایہ وکھاؤن دی لوڑ نہیں کہ ایتھے اسیں واقعی اک سیٹ~$A$ تعریف کیتا اے۔ پر سانوں ہن وی وکھاؤنا اے کہ $A$ وچ قضیے دیاں لوڑیندیاں خاصیتاں نیں۔ ایس توں بناں ثبوت پورا نہیں!{}
  \end{quote}
  کیوں جے $x \in A$، ایس لئی $A \neq \emptyset$۔
  \begin{quote}
    ایہ $A$ دی $\{x\}$ وجوں تعریف تے ایہناں صاف گلاں اُتے ٹکدا اے کہ $x \in \{x\}$ تے $x \notin \emptyset$۔
  \end{quote}
  کیوں جے $x$ اکّو اک !!{element} اے~$\{x\}$ دا، تے $x \in B$، سو ہر !!{element} جیہڑا~$A$ دا اے، اوہ !!a{element}~$B$ دا وی اے۔~$\subseteq$ دی تعریف مطابق $A \subseteq B$۔
\end{proof}

\subsection{موجودگی دے دعوے ورتنے}

منّ لو تہانوں پتا اے کہ موجودگی دا کوئی دعویٰ سچ اے (تسی اوہنوں ثابت کر چکے او یا ایہ اک منّی ہوئی شرط اے جیہڑی تسی ورت سکدے او)، مثلاً «کسے~$x$ لئی، $x \in A$» یا «اک $x \in A$ اے»۔ جے تسی اینوں ثبوت وچ ورتنا چاہندے او، تاں منّ سکدے او کہ تہاڈے کول اوہناں شےواں وچوں اک دا ناں اے جیہڑیاں تہاڈی شرط مطابق موجود نیں۔ کیوں جے $A$ وچ گھٹ توں گھٹ اک شے اے، ایہو جیہیاں شےواں نیں جیہناں ول اوہ ناں اشارہ کر سکدا اے۔ ہو سکدا اے تسی اوہناں وچوں کوئی اک شے چُن نہ سکو یا اوہدا ہور بیان نہ دے سکو (بس اینا کہ اوہ $\in A$ اے)۔ پر ثبوت لئی منّ سکدے او کہ تسی اوہ شے چُن لئی اے تے اوہدا ناں رکھ سکدے او۔ ضروری اے کہ اوہ ناں چُنو جیہڑا تسی پہلاں نہیں ورتیا تے جیہڑا تہاڈیاں منّیاں شرطاں وچ نہیں آؤندا، نہیں تاں گل وگڑ سکدی اے۔ ثبوت وچ «کسے $x$ لئی، $x \in A$» توں «$a
\in A$ ہووے» ول جا کے ایہ دس دیو۔ ہُن تسی~$a$ بارے دلیل دے سکدے او، ہور منّیاں شرطاں ورت سکدے او، وغیرہ، جدوں تک تسی نتیجے $p$ تک نہ پہنچ جاؤ۔ جے $p$ وچ ہن~$a$ دا ذکر نہیں، تاں $p$ ایس منّی ہوئی گل توں آزاد اے کہ $a \in A$، تے تسی وکھا دتا اے کہ اوہ بس ایس منّی ہوئی گل توں نکلدا اے کہ «کسے $x$ لئی، $x \in A$»۔

\begin{prop}
جے $A \neq \emptyset$، تاں $A \cup B \neq \emptyset$۔
\end{prop}

\begin{proof}
  منّ لو $A \neq \emptyset$۔ سو کسے $x$ لئی، $x \in A$۔
  \begin{quote}
    ایتھے اسیں پہلاں بس قضیے دی مفروضہ شرط فیر دسی اے۔ ایس شرط، یعنی $A \neq \emptyset$، وچ موجودگی دا اک دعویٰ لُکیا اے، جیہڑا کجھ تعریفاں کھولھ کے ای لبھدا اے۔ $=$ دی تعریف دسدی اے کہ $A = \emptyset$ اُدوں تے صرف اُدوں جدوں ہر $x \in A$، $\in \emptyset$ وی اے، تے ہر $x \in
    \emptyset$، $\in A$ وی اے۔ دوہاں پاسیاں دی نفی لئیے تاں $A \neq
    \emptyset$ اُدوں تے صرف اُدوں جدوں یا کوئی $x \in A$، $\notin \emptyset$ اے، یا کوئی $x \in \emptyset$، $\notin A$ اے۔ کیوں جے کوئی شے $\in \emptyset$ نہیں، فصل دا دوجا جز کدی سچ نہیں ہو سکدا، تے «$x \in A$ تے $x
    \notin \emptyset$» بس $x \in A$ رہ جاندا اے۔ سو $x \neq
    \emptyset$ اُدوں تے صرف اُدوں جدوں کسے $x$ لئی، $x \in A$۔ ایہ موجودگی دا دعویٰ اے۔ ہُن اسیں~$A$ دے !!{element}s وچوں اک دا ناں پیش کر کے اوہ دعویٰ ورتدے آں:
  \end{quote}
  $a \in A$ ہووے۔
  \begin{quote}
    ہُن اسیں اوہناں شےواں وچوں اک دا ناں پیش کیتا اے جیہڑیاں~$\in A$ نیں۔ اسیں اگے~$a$ بارے دلیل دیواں گے، پر دھیان رکھاں گے کہ بس $a \in A$ منّیئے، تے ہور کجھ نہ منّیئے:
  \end{quote}
  کیوں جے $a \in A$، ایس لئی $a \in A \cup B$، تعریف~$\cup$ مطابق۔ سو کسے $x$ لئی، $x \in A \cup B$، یعنی $A \cup B \neq \emptyset$۔
  \begin{quote}
    ایس آخری قدم وچ اسیں «$a \in A \cup B$» توں «کسے $x$ لئی، $x \in A \cup B$» تک گئے آں۔ ایس وچ ہن $a$ دا ذکر نہیں، سو سانوں پتا اے کہ «کسے $x$ لئی، $x \in A \cup B$» بس «کسے $x$ لئی، $x \in A$» توں نکلدا اے۔ پر ایہدا مطلب اے کہ $A \cup B \neq
    \emptyset$۔
  \end{quote}
\end{proof}

شاید چنگا طریقہ اے کہ «$x$» ورگے بندھے متغیر، $a$ ورگے منّے ہوئے ناںواں توں وکھ رکھو، جیویں اسیں کیتا۔ پر کم کردیاں اسیں اکثر ایہ نہیں کردے تے بس $x$ ای ورتدے آں، ایویں:
\begin{quote}
منّ لو $A \neq \emptyset$، یعنی اک $x \in A$ اے۔ $\cup$ دی تعریف مطابق $x \in A \cup B$۔ سو $A \cup B \neq \emptyset$۔
\end{quote}
پر جدوں تسی ایویں کردے او، تاں اچھی طرح دھیان رکھنا پئے گا کہ وکھ وکھ موجودگی دے دعویاں لئی وکھرے $x$ تے $y$ ورتو۔ مثلاً ہِیٹھلا «جے $A \neq
\emptyset$ تے $B \neq \emptyset$ تاں $A \cap B \neq \emptyset$» دا ٹھیک ثبوت \emph{نہیں} (تے ایہ دعویٰ سچ نہیں)۔
\begin{quote}
منّ لو $A \neq \emptyset$ تے $B \neq \emptyset$۔ سو کسے $x$ لئی، $x
\in A$، تے کسے $x$ لئی، $x \in B$ وی۔ کیوں جے $x \in A$ تے $x \in
B$، ایس لئی $x \in A \cap B$، تعریف~$\cap$ مطابق۔ سو $A \cap B \neq
\emptyset$۔
\end{quote}
کی تسی غلط قدم لبھ کے سمجھا سکدے او کہ نتیجہ کیوں نہیں نکلدا؟
\endgroup
% /OLP-0606

% OLP-0607 content/methods/proofs/example-1.tex
\begingroup
\def\olfilename{example-1}
\expandafter\def\csname ol@sectioncs\endcsname{section}

\olfileid{mth}{prf}{ex1}

\olsection{اک مثال}

ساڈی پہلی مثال سیٹاں دے اتحاد تے سانجھے حصے بارے ہِیٹھلی سوکھی گل اے۔ ایہ تعریفاں کھولھنا، عطف تے ساریاں لئی دعوے ثابت کرنا، تے صورتاں وکھ کر کے ثبوت دینا وکھاوے گی۔

\begin{prop}
کسے وی سیٹاں $A$، $B$ تے $C$ لئی، $A \cup (B \cap C) = (A \cup B)
\cap (A \cup C)$
\end{prop}

آؤ ایہنوں ثابت کریے!{}

\begin{proof}
اسیں وکھاؤنا چاہندے آں کہ کسے وی سیٹاں $A$، $B$ تے $C$ لئی، $A \cup (B \cap
C) = (A \cup B) \cap (A \cup C)$
\begin{quote}
پہلاں قضیے دے بیان وچ «$=$» دی تعریف کھولھدے آں۔ چیتے کرو کہ سیٹاں دا اکّو ہونا ثابت کرن دا مطلب ایہ وکھاؤنا اے کہ سیٹاں دے !!{element}s اکّو نیں۔ یعنی $A \cup
(B \cap C)$ دے سارے !!{element}s، $(A \cup B) \cap (A \cup C)$ دے !!{element}s وی نیں، تے اُلٹے پاسے وی۔ «اُلٹے پاسے وی» دا مطلب اے کہ $(A
\cup B) \cap (A \cup C)$ دا ہر !!{element}، $A \cup (B \cap
C)$ دا !!a{element} وی ہونا چاہیدا اے۔ سو تعریف کھولھ کے سانوں پتا لگدا اے کہ اسیں اک عطف ثابت کرنا اے۔ آؤ ایہ لکھ لئیے:
\end{quote}
تعریف مطابق، $A \cup (B \cap C) = (A \cup B) \cap (A \cup C)$ اُدوں تے صرف اُدوں جدوں $A \cup (B \cap C)$ دا ہر !!{element}، $(A
\cup B) \cap (A \cup C)$ دا !!a{element} وی اے، تے $(A \cup B) \cap (A
\cup C)$ دا ہر !!{element}، $A \cup (B \cap C)$ دا !!a{element} اے۔
\begin{quote}
کیوں جے ایہ عطف اے، سانوں اوہدے دوویں جز وکھ وکھ ثابت کرنے پَین گے۔ پہلاں پہلا جز لئیے۔ آؤ ثابت کریے کہ $A \cup (B \cap C)$ دا ہر !!{element}، $(A
\cup B) \cap (A \cup C)$ دا !!a{element} وی اے۔

ایہ ساریاں لئی دعویٰ اے۔ سو اسیں $A \cup (B \cap C)$ دا اک من مانا !!{element} لے کے وکھاؤندے آں کہ اوہ $(A \cup B) \cap (A \cup C)$ دا !!a{element} وی ہونا چاہیدا اے۔ ایس من مانے !!{element} دا ناں رکھن لئی اک متغیر چُناں گے، منّ لو~$z$۔ ساڈا ثبوت اگے ایویں چلدا اے:
\end{quote}
پہلاں اسیں ثابت کردے آں کہ $A \cup (B \cap C)$ دا ہر !!{element}، $(A \cup B) \cap (A \cup C)$ دا !!a{element} وی اے۔ $z \in A \cup (B
\cap C)$ ہووے۔ سانوں وکھاؤنا اے کہ $z \in (A \cup B) \cap (A \cup C)$۔
\begin{quote}
ہُن $\cup$ تے~$\cap$ دی تعریف کھولھن دا ویلا اے۔ مثلاً $\cup$ دی تعریف ایہ اے، $A \cup B = \Setabs{z}{z \in A
  \text{ یا } z \in B}$۔ جدوں اسیں «$A \cup (B
\cap C)$» اُتے ایہ تعریف لاؤندے آں، تاں تعریف وچ «$B$» والی تھاں ہُن «$B \cap C$» اے، سو $A \cup (B \cap C) = \Setabs{z}{z \in A \text{ یا }
  z \in B \cap C}$۔ سو ساڈے منّن دا مطلب کہ $z \in A \cup (B \cap C)$، ایہ اے کہ $z \in \Setabs{z}{z \in A \text{ یا } z \in B \cap
  C}$۔ تے $z \in \Setabs{z}{\dots z\dots}$ اُدوں تے صرف اُدوں جدوں \dots $z$ \dots، یعنی ایس صورت وچ، $z \in A$ یا $z \in B \cap C$۔
\end{quote}
$\cup$ دی تعریف مطابق یا $z \in A$ یا $z \in B \cap C$۔
\begin{quote}
ایہ فصل اے، سو صورتاں وکھ کر کے ثبوت دینا کم آوے گا۔ اسیں دوویں صورتاں لیندے آں تے ہر اک وچ وکھاؤندے آں کہ ساڈا لوڑیندا نتیجہ، یعنی «$z \in (A \cup B) \cap (A \cup
C)$»، نکلدا اے۔
\end{quote}
صورت ۱۔ منّ لو $z \in A$۔
\begin{quote}
منّیاں شرطاں توں ہور بہت کجھ کرن نوں نہیں۔ سو ویکھیے کہ نتیجے وچ ساڈے کم دی کیہڑی گل اے۔ اسیں وکھاؤنا چاہندے آں کہ $z \in (A \cup B) \cap (A \cup C)$۔ $\cap$ دی تعریف مطابق جے اسیں $z \in (A \cup B) \cap (A \cup C)$ وکھاؤنا اے، تاں وکھاؤنا پئے گا کہ ایہ $(A \cup B)$ تے $(A \cup C)$ دوہاں وچ اے۔ پر $z
\in A \cup B$ اُدوں تے صرف اُدوں جدوں $z \in A$ یا $z \in B$، تے ساڈے کول پہلاں ای صورت~۱ دی شرط اے کہ $z \in A$۔ ایس دلیل وچ $C$ نوں $B$ دی تھاں رکھیے تاں $z \in A \cup C$۔ ایہ دلیل اُلٹی سمت وچ چلی اے، سو آؤ اینوں ثبوت لئی لوڑیندی سمت وچ لکھیے۔
\end{quote}
کیوں جے $z \in A$، سو $z \in A$ یا $z \in B$، تے ایس لئی تعریف~$\cup$ مطابق $z \in A \cup B$۔ ایویں ای $z \in A \cup C$۔ پر ایہدا مطلب اے کہ $z \in (A \cup B) \cap (A \cup C)$، تعریف~$\cap$ مطابق۔
\begin{quote}
ایس نال صورتاں وکھ کر کے دتے ثبوت دی پہلی صورت پوری ہوئی۔ ہُن اسیں دوجی صورت وچ نتیجہ کڈھنا چاہندے آں، جتھے $z \in B \cap C$۔
\end{quote}
صورت ۲۔ منّ لو $z \in B \cap C$۔
\begin{quote}
اسیں فیر دو سیٹاں دے سانجھے حصے نال کم کردے آں۔ آؤ~$\cap$ دی تعریف لائیے:
\end{quote}
کیوں جے $z \in B \cap C$، سو $z$، $B$ تے $C$ دوہاں دا !!a{element} ہونا چاہیدا اے، تعریف~$\cap$ مطابق۔
\begin{quote}
ہُن فیر اپنے نتیجے ول ویکھن دا ویلا اے۔ سانوں وکھاؤنا اے کہ $z$، $(A \cup B)$ تے $(A \cup C)$ دوہاں وچ اے۔ فیر حل سِدھا ای نکلدا اے۔
\end{quote}
کیوں جے $z \in B$، سو $z \in (A \cup B)$۔ کیوں جے $z \in C$، سو $z \in (A
\cup C)$ وی۔ ایس لئی $z \in (A \cup B) \cap (A \cup C)$۔
\begin{quote}
ایتھے اسیں $\cup$ تے $\cap$ دیاں تعریفاں فیر لائیاں، پر کیوں جے اسیں اوہ پہلاں چیتے کر چکے آں تے وکھا چکے آں کہ جے $z$ دو سیٹاں وچوں کسے اک وچ اے تاں اوہ اوہناں دے اتحاد وچ اے، ایتھے ہر قدم کھولھ کے دسّن دی لوڑ نہیں۔

صورتاں وکھ کر کے ثبوت دی دوجی صورت وی پوری ہو گئی، سو ہُن اسیں اپنا پہلا نتیجہ دَس سکدے آں۔
\end{quote}
سو جے $z \in A \cup (B \cap C)$ تاں $z \in (A \cup B) \cap (A \cup C)$۔
\begin{quote}
ہُن بس دوجی سمت وکھاؤنی اے کہ $(A \cup B) \cap (A \cup C)$ دا ہر !!{element}، $A \cup (B \cap
C)$ دا !!a{element} اے۔ پہلاں وانگ ای اسیں پہلے سیٹ دا اک من مانا !!{element} منّ کے ساریاں لئی ایہ دعویٰ ثابت کردے آں، تے وکھاؤندے آں کہ اوہ دوجے سیٹ وچ وی ہونا چاہیدا اے۔ آؤ ایہ گل لکھ لئیے۔
\end{quote}
ہُن منّ لو $z \in (A \cup B) \cap (A \cup C)$۔ اسیں وکھاؤنا چاہندے آں کہ $z \in A \cup (B \cap C)$۔
\begin{quote}
ہُن اسیں ایس مفروضہ شرط توں کم کردے آں کہ $z \in (A \cup B) \cap (A
\cup C)$۔ آس اے کہ پہلے حصے والا اوہو~$z$ ورتن نال بھلیکھا نہیں پئے گا۔ جدوں اوہ حصہ مکیا سی، اوتھے منّیاں ساریاں شرطاں اوتھے ای مُک گئیاں۔ سو ہُن اسیں $z$ بارے نویاں شرطاں منّ سکدے آں۔ جے ایہ نال بھلیکھا پئے، تاں اگے $z$ دی تھاں کوئی ہور متغیر رکھ لو۔

سانوں پتا اے کہ $z$، $A \cup B$ تے $A \cup C$ دوہاں وچ اے، تعریف~$\cap$ مطابق۔ تے $\cup$ دی تعریف نال ایہنوں ہور کھولھ سکدے آں۔ یا $z \in A$ یا $z \in B$، تے نال یا $z \in A$ یا $z \in
C$۔ ایہ فیر صورتاں وکھ کر کے ثبوت لگدا اے، پر «تے» نال بھلیکھا پے سکدا اے۔ تسی سوچ سکدے او کہ ایہدا مطلب تِن صورتاں نیں، $z$ یا $A$، یا $B$، یا $C$ وچ اے۔ پر ایہ غلط ہووے گا۔ سانوں دھیان نال ہر فصل واری واری ویکھنی چاہیدی اے۔
\end{quote}
$\cap$ دی تعریف مطابق $z \in A \cup B$ تے $z \in A \cup C$۔ $\cup$ دی تعریف مطابق $z \in A$ یا $z \in B$۔ اسیں صورتاں وکھ کردے آں۔
\begin{quote}
کیوں جے اسیں پہلی فصل ول دھیان دیندے آں، اسیں (تعریف $A \cup C$ کھولھ کے) دوجی فصل ہن تک نہیں کڈھی۔ اصل وچ ہن تک اوہدی لوڑ وی نہیں۔ پہلی صورت $z \in A$ اے، تے کسے سیٹ دا !!a{element} اوس سیٹ دے کسے ہور سیٹ نال اتحاد دا !!a{element} وی اے۔ سو صورت~۱ سوکھی اے:
\end{quote}
صورت ۱۔ منّ لو $z \in A$۔ ایس توں نکلدا اے کہ $z \in A \cup (B \cap
C)$۔
\begin{quote}
ہُن دوجی صورت، $z \in B$۔ ایتھے دوجے $\cup$ دی تعریف کھولھاں گے تے فیر صورتاں وکھ کر کے ثبوت دیواں گے:
\end{quote}
صورت ۲۔ منّ لو $z \in B$۔ کیوں جے $z \in A \cup C$، سو یا $z \in
A$ یا $z \in C$۔ اسیں اگوں صورتاں وکھ کردے آں:

صورت ۲ا۔ $z \in A$۔ تاں فیر $z \in A \cup (B \cap C)$۔
\begin{quote}
ایہ ذرا اوپری گل سی۔ ایس صورت لئی سانوں ایہ منّن دی لوڑ اصل وچ نہیں پئی کہ~$z \in B$، پر ایس نال کوئی اوکڑ نہیں۔
\end{quote}
صورت ۲ب۔ $z \in C$۔ تاں $z \in B$ تے $z \in C$، سو $z \in B \cap C$، تے ایس لئی $z \in A \cup (B \cap C)$۔
\begin{quote}
ایس نال صورتاں وکھ کر کے دتے دوویں ثبوت پورے ہو گئے، سو دوجا ادھ وی پورا ہو گیا۔
\end{quote}
سو جے $z \in (A \cup B) \cap (A \cup C)$ تاں $z \in A \cup (B \cap C)$۔
\end{proof}
\endgroup
% /OLP-0607

% OLP-0608 content/methods/proofs/example-2.tex
\begingroup
\def\olfilename{example-2}
\expandafter\def\csname ol@sectioncs\endcsname{section}

\olfileid{mth}{prf}{ex2}

\olsection{اک ہور مثال}

\begin{prop}
جے $A \subseteq C$، تاں $A \cup (C \setminus A) = C$۔
\end{prop}

\begin{proof}
  منّ لو $A \subseteq C$۔ اسیں وکھاؤنا چاہندے آں کہ $A \cup (C
  \setminus A) = C$۔
  \begin{quote}
    پہلاں اسیں ویکھدے آں کہ ایہ اک شرطیہ بیان اے۔ ایتھے ساریاں لئی ہون دی گل اَن کہی اے۔ قضیہ سارے سیٹاں $A$ تے $C$ لئی اے۔ سو $A$ تے $C$ من مانے سیٹاں لئی متغیر نیں۔ ایہو جیہا بیان ثابت کرن لئی اسیں مقدم منّدے آں تے تالی ثابت کردے آں۔

    اسیں مفروضہ شرط $A \subseteq C$ ورت کے اگے ودھدے آں۔ آؤ~$\subseteq$ دی تعریف کھولھیے۔ شرط دا مطلب اے کہ~$A$ دے سارے !!{element}s،~$C$ دے !!{element}s وی نیں۔ آؤ ایہ گل لکھ لئیے، ایہ اہم گل اے جیہڑی سارے ثبوت وچ ورتاں گے۔
  \end{quote}
  تعریف~$\subseteq$ مطابق، کیوں جے $A \subseteq C$، سو سارے $z$ لئی، جے $z \in A$ تاں $z \in C$۔
  \begin{quote}
    اسیں منّی ہوئی شرط وچ دتیاں ساریاں تعریفاں کھولھ لئیاں نیں۔ ہُن نتیجے ول جا سکدے آں۔ اسیں وکھاؤنا چاہندے آں کہ $A \cup (C \setminus A) = C$۔ سو پچھلی مثال وانگ ثبوت دی بناوٹ رکھدے آں۔ وکھاؤ کہ $A \cup (C \setminus A)$ دا ہر !!{element}،~$C$ دا !!a{element} وی اے، تے اُلٹے پاسے، $C$ دا ہر !!{element}، $A \cup (C
    \setminus A)$ دا !!a{element} اے۔ ایہنوں مختصر ایویں لکھ سکدے آں، $A \cup (C \setminus A)
    \subseteq C$ تے $C \subseteq A \cup (C \setminus A)$۔ (ایتھے اسیں تعریف کھولھن توں اُلٹ کم کردے آں، پر ایس نال ثبوت پڑھنا ذرا سوکھا ہُندا اے۔) کیوں جے ایہ عطف اے، سانوں دوویں حصے ثابت کرنے پَین گے۔ پہلا حصہ، یعنی $A \cup (C \setminus A)$ دا ہر !!{element}،~$C$ دا !!a{element} وی اے، وکھاؤن لئی اسیں منّدے آں کہ $z \in A \cup (C \setminus A)$، جتھے~$z$ من مانا اے، تے وکھاؤندے آں کہ $z \in C$۔ $\cup$ دی تعریف مطابق اسیں نتیجہ کڈھ سکدے آں کہ $z \in A$ یا $z \in C \setminus A$، کیوں جے $z \in A \cup (C \setminus A)$۔ ہُن تہانوں ایس طریقے دی سمجھ آؤندی ہونی چاہیدی اے۔
    \end{quote}
  $A \cup (C \setminus A) = C$ اُدوں تے صرف اُدوں جدوں $A \cup (C \setminus A) \subseteq
  C$ تے $C \subseteq (A \cup (C \setminus A)$۔ پہلاں اسیں ثابت کردے آں کہ $A \cup (C \setminus A) \subseteq C$۔ $z \in A \cup (C
  \setminus A)$ ہووے۔ سو یا $z \in A$ یا $z \in (C \setminus A)$۔
  \begin{quote}
    اسیں اک فصل تک پہنچے آں تے اوہتوں ثابت کرنا چاہندے آں کہ $z \in C$۔ اسیں ایہ کم صورتاں وکھ کر کے کردے آں۔
  \end{quote}
  صورت ۱۔ $z \in A$۔ کیوں جے سارے $z$ لئی جے $z \in A$ تاں $z \in C$، سو $z \in C$۔
  \begin{quote}
    ایتھے اسیں پہلاں لکھی گل ورتی اے جیہڑی قضیے دی شرط $A \subseteq C$ توں نکلی سی۔ پہلی صورت پوری ہوئی۔ ہُن دوجی صورت $z \in (C
    \setminus A)$ ول آؤ۔ چیتے کرو کہ $C \setminus A$ دوہاں سیٹاں دے \emph{فرق} نوں دسدا اے، یعنی~$C$ دے اوہناں سارے !!{element}s دا سیٹ جیہڑے~$A$ دے !!{element}s نہیں۔ پر $C$ دا ہر !!{element} جیہڑا~$A$ وچ نہیں، خاص کر کے~$C$ دا !!a{element} اے۔
  \end{quote}
  صورت ۲۔ $z \in (C \setminus A)$۔ ایہدا مطلب اے کہ $z \in C$ تے $z
  \notin A$۔ سو خاص کر کے $z \in C$۔
  \begin{quote}
    چنگا، پہلی سمت ثابت ہو گئی۔ ہُن دوجی سمت لئیے۔ ایتھے اسیں ثابت کردے آں کہ $C \subseteq A \cup (C \setminus
    A)$۔ سو منّدے آں کہ $z \in C$ تے ثابت کردے آں کہ $z \in A \cup (C
    \setminus A)$۔
  \end{quote}
  ہُن $z \in C$ ہووے۔ اسیں وکھاؤنا چاہندے آں کہ $z \in A$ یا $z \in C
  \setminus A$۔
  \begin{quote}
    کیوں جے $A$ دے سارے !!{element}s، $C$ دے !!{element}s وی نیں، تے $C
    \setminus A$ اوہناں ساریاں شےواں دا سیٹ اے جیہڑیاں $C$ دے !!{element}s نیں پر $A$ دے نہیں، سو $z$ یا $A$ وچ اے یا $C
    \setminus A$ وچ۔ جے تہانوں پہلاں ای پتا نہیں کہ نتیجہ کیوں سچ اے، تاں ایہ گل ذرا دھندلی لگ سکدی اے۔ ایہنوں قدم بہ قدم ثابت کرنا چنگا اے۔ اک سوکھی گل ورتنی کم آوے گی جیہڑی اسیں ثبوت توں بناں دس سکدے آں، $z \in A$ یا $z \notin A$۔ اینوں «تیجی صورت نہ ہون دا اصول» آکھدے نیں۔ کسے وی بیان~$p$ لئی، یا $p$ سچ اے یا اوہدی نفی سچ اے۔ (ایتھے $p$ ایہ بیان اے کہ $z \in A$۔) ایہ فصل اے، سو اسیں فیر صورتاں وکھ کر کے ثبوت دے سکدے آں۔
  \end{quote}
  یا $z \in A$ یا $z \notin A$۔ پہلی صورت وچ $z \in A \cup
  (C \setminus A)$۔ دوجی صورت وچ $z \in C$ تے $z \notin A$، سو $z \in C \setminus A$۔ پر فیر $z \in A \cup (C \setminus A)$۔
  \begin{quote}
    ساڈا ثبوت پورا اے۔ اسیں وکھا دتا اے کہ $A \cup (C \setminus A) = C$۔
  \end{quote}
\end{proof}
\endgroup
% /OLP-0608

% OLP-0609 content/methods/proofs/proof-by-contradiction.tex
\begingroup
\def\olfilename{proof-by-contradiction}
\expandafter\def\csname ol@sectioncs\endcsname{section}

\olfileid{mth}{prf}{con}

\olsection{ٹکراؤ توں ثبوت}

پہلاں ایہ ویکھو کہ ٹکراؤ توں ثبوت، استنباط دا اوہ ڈھنگ اے جیہڑا نفی والے دعوے ثابت کرن لئی ورتیا جاندا اے۔ منّ لو تسی وکھاؤنا چاہندے او کہ کوئی دعویٰ~$p$ \emph{غلط} اے، یعنی تسی~$\lnot p$ وکھاؤنا چاہندے او۔ سبھ توں آس والا طریقہ ایہ اے کہ (ا) $p$ نوں سچ منّو، تے (ب) وکھاؤ کہ ایس منّی ہوئی گل توں کوئی ایہو جیہی گل نکلدی اے جیہڑی تہانوں پتا اے کہ غلط اے۔ «ایہو جیہی گل جیہڑی پتا اے کہ غلط اے» کوئی ایہو جیہا نتیجہ ہو سکدا اے جیہڑا $p$ آپ نال، یا سارے دعوے دی کسے ہور مفروضہ شرط نال ٹکراؤندا اے۔ مثلاً «جے $q$ تاں $\lnot p$» دے ثبوت وچ $q$ نوں سچ منّنا تے اوہتوں~$\lnot
p$ ثابت کرنا پیندا اے۔ جے تسی $\lnot p$ ٹکراؤ توں ثابت کردے او، تاں $p$ وی منّنا پئے گا،~$q$ دے نال۔ جے تسی $\lnot q$، $p$ توں ثابت کر سکو، تاں تسی وکھا دتا کہ مفروضہ~$p$ توں ایہو جیہی گل نکلدی اے جیہڑی تہاڈے دوجے مفروضے~$q$ نال ٹکراؤندی اے، کیوں جے $q$~تے $\lnot q$ دوویں سچ نہیں ہو سکدے۔ بے شک ٹکراؤ ثابت کردیاں تہانوں استنباط دے ہور ڈھنگ ورتنے پَین گے تے تعریفاں وی کھولھنیاں پَین گئیاں۔ اک مثال ویکھیے۔

\begin{prop}
  جے $A \subseteq B$ تے $B = \emptyset$، تاں $A$ دے !!{element}s کوئی نہیں۔
\end{prop}

\begin{proof}
  منّ لو $A \subseteq B$ تے $B = \emptyset$۔ اسیں وکھاؤنا چاہندے آں کہ $A$ دے !!{element}s کوئی نہیں۔
  \begin{quote}
    ایہ شرطیہ دعویٰ اے، سو مقدم منّ کے تالی ثابت کرنا چاہندے آں۔ تالی ایہ اے کہ $A$ دے !!{element}s کوئی نہیں۔ ایہ گل ذرا کھولھ کے آکھ سکدے آں کہ کوئی ایہو جیہا~$x \in A$ نہیں۔
  \end{quote}
  $A$ دے !!{element}s کوئی نہیں اُدوں تے صرف اُدوں جدوں کوئی ایہو جیہا~$x$ نہیں جیہدے لئی $x \in A$۔
  \begin{quote}
    سو اسیں طے کیتا اے کہ ثابت کرن والی گل اصل وچ نفی والا دعویٰ~$\lnot p$ اے، یعنی کوئی $x \in A$ نہیں۔ ٹکراؤ توں ثبوت لئی سانوں اوہدے نال ملدا اثباتی دعویٰ~$p$ منّنا پئے گا، یعنی کوئی $x \in A$ اے، تے اوہتوں ٹکراؤ ثابت کرنا پئے گا۔ ایہ دسّن لئی کہ اسیں ٹکراؤ توں ثبوت کردے آں، «ٹکراؤ کڈھن لئی منّ لو» یا بس «اُلٹ منّ لو» لکھدے آں تے فیر مفروضہ~$p$ دس دیندے آں۔
  \end{quote}
  اُلٹ منّ لو، کوئی $x \in A$ اے۔
  \begin{quote}
    ہُن ایہ اوہ نواں مفروضہ اے جیہڑا اسیں ٹکراؤ کڈھن لئی ورتاں گے۔ ساڈے دو ہور مفروضے نیں، $A \subseteq B$ تے $B = \emptyset$۔ پہلے توں ملدا اے کہ $x \in B$:
  \end{quote}
  کیوں جے $A \subseteq B$، سو $x \in B$۔
  \begin{quote}
    پر کیوں جے $B = \emptyset$، سو $B$ دا ہر !!{element} (جیویں $x$) لازماً~$\emptyset$ دا وی !!a{element} اے۔
  \end{quote}
  کیوں جے $B = \emptyset$، سو $x \in \emptyset$۔ ایہ ٹکراؤ اے، کیوں جے تعریف مطابق $\emptyset$ دے !!{element}s کوئی نہیں۔
  \begin{quote}
    ایس نال ثبوت پہلاں ای پورا اے۔ اسیں اپنے مفروضیاں توں لوڑیندی گل (اک ٹکراؤ) تک پہنچ گئے آں، تے ایس دا مطلب اے کہ سارے مفروضے سچ نہیں ہو سکدے۔ پہلے دو مفروضیاں ($A \subseteq B$ تے $B = \emptyset$) اُتے کوئی اعتراض نہیں۔ سو آخر وچ پیش کیتا مفروضہ (کوئی $x \in A$ اے) غلط ہونا پئے گا۔ پر جے اسیں ثبوت نوں ہور صاف رکھنا چاہیئے، تاں ایہ گل کھولھ کے آکھ سکدے آں۔
  \end{quote}
  سو ساڈا مفروضہ کہ کوئی $x \in A$ اے لازماً غلط اے۔ ایس لئی ٹکراؤ توں ثبوت مطابق $A$ دے !!{element}s کوئی نہیں۔
\end{proof}

ہر اثباتی دعویٰ سِدھے ای اک نفی والے دعوے دے برابر اے۔ $p$ اُدوں تے صرف اُدوں جدوں $\lnot\lnot p$۔ سو ٹکراؤ توں ثبوت اثباتی دعوے «بالواسطہ» ثابت کرن لئی وی ورتیا جا سکدا اے۔~$p$ ثابت کرن لئی اینوں نفی والا دعویٰ $\lnot\lnot p$ سمجھو۔ جے اسیں $\lnot p$ توں ٹکراؤ ثابت کر سکیے، تاں ٹکراؤ توں ثبوت نال $\lnot\lnot p$ ثابت ہو گیا، تے ایس لئی~$p$ وی۔

پچھلی مثال وچ اسیں نفی والا دعویٰ ثابت کرنا چاہندے ساں، یعنی $A$ دے !!{element}s کوئی نہیں۔ سو ٹکراؤ توں ثبوت لئی جیہڑا مفروضہ منّیا سی (یعنی کوئی $x \in A$ اے) اوہ اثباتی دعویٰ سی۔ ایس نال سانوں دلیل لئی اک گل ملی، یعنی منّیا ہویا $x \in A$، جیہدے بارے دلیل دے کے اسیں $x \in \emptyset$ تک پہنچے۔

اثباتی دعوے دا بالواسطہ ثبوت کردیاں ٹکراؤ کڈھن لئی تہاڈا مفروضہ نفی والا ہووے گا۔ پر اکثر تسی اثباتی دعوے نوں سوکھ نال نفی والی صورت وچ، تے نفی والے دعوے نوں اثباتی صورت وچ فیر لکھ سکدے او۔ جے اسیں «$A = \emptyset$» ثابت کردے، نفی والی تالی «$A$ دے !!{element}s کوئی نہیں» دی تھاں، تاں ساڈا پچھلا ثبوت بنیادی طور اُتے اوہو رہنا سی۔ ($=$ دی تعریف مطابق «$A = \emptyset$» اک عام دعویٰ اے، کیوں جے کھولھ کے ایہدا مطلب اے «ہر !!{element} جیہڑا~$A$ دا اے، اوہ~$\emptyset$ دا وی !!{element} اے، تے اُلٹے پاسے وی»۔) پر ایہ سوکھ نال ویکھ سکدے آں کہ ایہ نفی والے دعوے «ایہ گل نہیں کہ کوئی $x \in A$ اے» دے برابر اے۔

سو کدی $\lnot p$ نوں مفروضہ منّ کے کم کرنا،~$p$ نوں سِدھا ثابت کرن نالوں سوکھا اے۔ بھاویں سِدھا ثبوت اوہنا ای سوکھا ہووے یا ہور وی سوکھا (جیویں اگلیاں مثالاں وچ)، کجھ لوک بالواسطہ طریقہ پسند کردے نیں۔ جے دوہری نفی نال تہانوں بھلیکھا پئے، تاں کسے دعوے دے ٹکراؤ توں ثبوت نوں اوہدے \emph{اُلٹ} دعوے توں ٹکراؤ دا ثبوت سمجھو۔ سو $\lnot
p$ دا ٹکراؤ توں ثبوت مفروضے~$p$ توں ٹکراؤ دا ثبوت اے، تے~$p$ دا ٹکراؤ توں ثبوت~$\lnot p$ توں ٹکراؤ دا ثبوت اے۔

\begin{prop}
$A \subseteq A \cup B$.
\end{prop}

\begin{proof}
  اسیں وکھاؤنا چاہندے آں کہ $A \subseteq A \cup B$۔
  \begin{quote}
    پہلاں ویکھن نوں ایہ اثباتی دعویٰ اے۔ ہر $x \in A$، $A \cup B$ وچ وی اے۔ اوہدی نفی ایہ اے کہ کوئی $x \in A$، $\notin A \cup B$ اے۔ سو اسیں ایس نفی والے دعوے نوں منّ کے تے اوہتوں ٹکراؤ وکھا کے دعویٰ بالواسطہ ثابت کر سکدے آں۔
    \end{quote}
  اُلٹ منّ لو، یعنی $A \nsubseteq A \cup B$۔
  \begin{quote}
    ساڈے کول $A \subseteq A \cup B$ دی تعریف اے، ہر $x \in A$، $\in A \cup B$ وی اے۔ $A \nsubseteq A \cup B$ دا مطلب سمجھن لئی سانوں کھولھی ہوئی تعریف اُتے بنیادی منطقی کم کرنا پئے گا۔ ایہ غلط اے کہ ہر $x \in A$، $\in
    A \cup B$ وی اے، اُدوں تے صرف اُدوں جدوں \emph{کوئی}~$x \in A$ اے جیہڑا $\notin C$ اے۔ (ایتھے بڑا دھیان رکھو۔ کئی طالب علماں دے ٹکراؤ توں ثبوت دی کوشش ایس لئی نہیں چل دی کہ اوہ «سارے $A$، $B$ نیں» ورگے دعوے دی نفی غلط سمجھدے نیں۔) دوجے لفظاں وچ، $A \nsubseteq A \cup B$ اُدوں تے صرف اُدوں جدوں کوئی $x$ اے جیہدے لئی $x \in A$ تے $x \notin A \cup B$۔ ایتھوں اگے گل سوکھی اے۔
  \end{quote}
  سو کوئی $x \in A$ اے جیہدے لئی $x \notin A \cup B$۔ $\cup$ دی تعریف مطابق $x \in A \cup B$ اُدوں تے صرف اُدوں جدوں $x \in A$ یا $x \in
  B$۔ کیوں جے $x \in A$، سو $x \in A \cup B$۔ ایہ مفروضے $x \notin A \cup B$ نال ٹکراؤندا اے۔
\end{proof}

\begin{prob}
\emph{بالواسطہ} ثابت کرو کہ $A \cap B \subseteq A$۔
\end{prob}

\begin{prop}
جے $A \subseteq B$ تے $B \subseteq C$، تاں $A \subseteq C$۔
\end{prop}

\begin{proof}
  منّ لو $A \subseteq B$ تے $B \subseteq C$۔ اسیں وکھاؤنا چاہندے آں کہ $A
  \subseteq C$۔
  \begin{quote}
    آؤ بالواسطہ طریقہ ورتیے۔ جیہڑی گل ثابت کرنا چاہندے آں، اوہدی نفی منّ لیندے آں۔
  \end{quote}
  اُلٹ منّ لو، یعنی $A \nsubseteq C$۔
  \begin{quote}
    پہلاں وانگ، اسیں دلیل دیندے آں کہ $A \nsubseteq C$ اُدوں تے صرف اُدوں جدوں ایہ گل نہیں کہ ہر $x \in A$، $\in C$ وی اے، یعنی کوئی $x \in A$، $\notin C$ اے۔ فکر نہ کرو، مشق نال ایہو جیہیاں نفیاں کھولھن لئی بوہتا سوچنا نہیں پئے گا۔
  \end{quote}
  دوجے لفظاں وچ، کوئی~$x$ اے جیہدے لئی $x \in A$ تے $x \notin C$۔
  \begin{quote}
    ہُن ایہ گل ٹکراؤ تک پہنچن لئی ورت سکدے آں۔ بے شک ایس کم لئی دوویں ہور مفروضے وی ورتنے پَین گے۔
  \end{quote}
  کیوں جے $A \subseteq B$، سو $x \in B$۔ کیوں جے $B \subseteq C$، سو $x \in
  C$۔ پر ایہ $x \notin C$ نال ٹکراؤندا اے۔
\end{proof}

\begin{prop}
جے $A \cup B = A \cap B$، تاں $A = B$۔
\end{prop}

\begin{proof}
  منّ لو $A \cup B = A \cap B$۔ اسیں وکھاؤنا چاہندے آں کہ $A = B$۔
  \begin{quote}
    ہُن شروع دی گل جانی پچھانی اے:
  \end{quote}
  ٹکراؤ کڈھن لئی منّ لو کہ $A \neq B$۔
  \begin{quote}
    ٹکراؤ توں ثبوت لئی ساڈا مفروضہ $A \neq
    B$ اے۔ کیوں جے $A = B$ اُدوں تے صرف اُدوں جدوں $A \subseteq B$ تے $B \subseteq A$، ایس توں ملدا اے کہ $A \neq B$ اُدوں تے صرف اُدوں جدوں $A \nsubseteq B$ \emph{یا} $B \nsubseteq
    A$۔ (ویکھو نفیاں نال کم کردیاں دھیان رکھنا کِنّا ضروری اے!{}) ایس فصل توں ٹکراؤ ثابت کرن لئی صورتاں وکھ کر کے ثبوت دیندے آں تے وکھاؤندے آں کہ ہر صورت وچ ٹکراؤ نکلدا اے۔
  \end{quote}
  $A \neq B$ اُدوں تے صرف اُدوں جدوں $A \nsubseteq B$ یا $B \nsubseteq A$۔ اسیں صورتاں وکھ کردے آں۔
  \begin{quote}
    پہلی صورت وچ اسیں $A \nsubseteq B$ منّدے آں، یعنی کسے $x$ لئی، $x \in A$ پر $\notin B$۔ $A \cap B$ دی تعریف اوہناں !!{element}s نال دتی اے جیہڑے $A$ تے $B$ دے سانجھے نیں۔ سو جے کوئی شے کسے اک وچ نہیں، تاں سانجھے حصے وچ نہیں۔ $A \cup B$، $A$ تے $B$ دا اکٹھ اے، سو جیہڑی شے کسے اک وچ اے، اوہ اتحاد وچ وی اے۔ ایس توں پتا لگدا اے کہ $x \in A \cup B$ پر $x \notin A \cap
    B$، تے ایس لئی $A \cap B \neq A \cup B$۔
  \end{quote}

  صورت ۱۔ $A \nsubseteq B$۔ تاں کسے $x$ لئی، $x \in A$ پر $x \notin
  B$۔ کیوں جے $x \notin B$، سو $x \notin A \cap B$۔ کیوں جے $x \in A$، سو $x \in A \cup B$۔ ایس لئی $A \cap B \neq A \cup B$، جیہڑا مفروضے $A \cap B = A \cup B$ نال ٹکراؤندا اے۔

  صورت ۲۔ $B \nsubseteq A$۔ تاں کسے $y$ لئی، $y \in B$ پر $y \notin
  A$۔ پہلاں وانگ، سانوں $y \in A \cup B$ پر $y \notin A \cap B$ ملدا اے، سو $A \cap B \neq A \cup B$، جیہڑا فیر $A \cap B = A \cup
  B$ نال ٹکراؤندا اے۔
\end{proof}
\endgroup
% /OLP-0609

% OLP-0610 content/methods/proofs/reading-proofs.tex
\begingroup
\def\olfilename{reading-proofs}
\expandafter\def\csname ol@sectioncs\endcsname{section}

\olfileid{mth}{prf}{rea}
\olsection{Reading Proofs}

درسی کتاباں تے مضموناں وچ ملدے ثبوتاں وچ ایہ ساری تفصیل وِرلی ای ہُندی اے جیہڑی اسیں ہُن تک اپنی مثالاں وچ دتی اے۔ لکھاری اکثر ایہ نہیں دسدے کہ اوہ کدوں صورتاں وکھ کردے نیں، کدوں بالواسطہ ثبوت دیندے نیں، یا کدوں کوئی تعریف ورتدے نیں۔ سو جدوں تسی درسی کتاب وچ ثبوت پڑھدے او، اوہنوں سمجھن لئی تہانوں اکثر ایہ تفصیل آپ پوری کرنی پَندی اے۔ ثبوت وچ جیہڑے وکھ وکھ قدم چُکنے پَیندے نیں، اوہناں دی سمجھ لئی ایہ چنگی مشق وی اے۔ اک مثال ویکھیے۔

\begin{prop}[جذب دا قانون]
سارے سیٹاں $A$، $B$ لئی،
\[
A \cap (A \cup B) = A
\]
\end{prop}

\begin{proof}
جے $z \in A \cap (A \cup B)$، تاں $z \in A$، سو $A \cap (A \cup B)
\subseteq A$۔ ہُن منّ لو $z \in A$۔ تاں $z \in A \cup B$ وی، تے ایس لئی $z \in A \cap (A \cup B)$ وی۔
\end{proof}

جذب دے قانون دا اُتلا ثبوت بڑا مختصر اے۔ جیہڑیاں تعریفاں ورتیاں نیں اوہناں دا ذکر نہیں، ثابت کرن توں پہلاں «سانوں ایہ ثابت کرنا اے کہ» نہیں لکھیا، وغیرہ۔ آؤ اینوں کھولھیے۔ ثابت کیتا قضیہ سارے سیٹاں $A$ تے $B$ بارے عام دعویٰ اے۔ جدوں ثبوت وچ $A$ یا $B$ آؤندا اے، ایہ من مانے سیٹاں لئی متغیر نیں۔ ثبوت جیہڑے عام دعوے ثابت کردا اے، سیٹاں دا اکّو ہونا ثابت کرن لئی اوہناں دی لوڑ اے، یعنی مساوات دے کھبّے پاسے دا ہر !!{element} سجّے پاسے دا !!a{element} وی اے، تے اُلٹے پاسے وی۔

\begin{quote}
«جے $z \in A \cap (A \cup B)$، تاں $z \in A$، سو $A \cap (A \cup B)
  \subseteq A$۔»
\end{quote}

ایہ سیٹاں دے اکّو ہون دے ثبوت دا پہلا ادھ اے۔ ایہ ثابت کردا اے کہ جے من مانا~$z$ کھبّے پاسے دا !!a{element} اے، تاں سجّے پاسے دا !!a{element} وی اے، یعنی $A \cap (A \cup B) \subseteq A$۔
منّ لو $z \in A \cap (A \cup B)$۔ کیوں جے $z$ دو سیٹاں دے سانجھے حصے دا !!{element} اُدوں تے صرف اُدوں اے جدوں دوہاں سیٹاں دا !!{element} اے، اسیں نتیجہ کڈھ سکدے آں کہ $z \in A$ تے $z \in A \cup B$ وی۔ خاص کر کے $z \in A$، ایہو اسیں وکھاؤنا چاہندے ساں۔ کیوں جے پہلے ادھ لئی بس ایہو کرنا سی، سانوں پتا اے کہ باقی ثبوت دوجے ادھ دا ثبوت اے، یعنی $A \subseteq A \cap
(A \cup B)$ دا ثبوت۔

\begin{quote}
«ہُن منّ لو $z \in A$۔ تاں $z \in A \cup B$ وی، تے ایس لئی $z \in A \cap (A \cup B)$ وی۔»
\end{quote}

اسیں $z \in A$ منّ کے شروع کردے آں، کیوں جے اسیں وکھاؤندے آں کہ کسے وی~$z$ لئی، جے $z \in A$ تاں $z \in A \cap (A \cup B)$۔ $z
\in A \cap (A \cup B)$ وکھاؤن لئی سانوں («$\cap$» دی تعریف مطابق) وکھاؤنا اے کہ (۱) $z \in A$ تے نال (۲) $z \in A \cup B$۔ ایتھے (۱) ساڈا مفروضہ ای اے، سو ہور کجھ ثابت کرنا نہیں، تے ایس لئی ثبوت فیر اوہدا ذکر نہیں کردا۔ (۲) لئی چیتے کرو کہ $z$، سیٹاں دے اتحاد دا !!a{element} اُدوں تے صرف اُدوں اے جدوں اوہناں وچوں گھٹ توں گھٹ اک سیٹ دا !!{element} اے۔ کیوں جے $z \in A$، تے $A \cup B$، $A$ تے $B$ دا اتحاد اے، ایتھے ایہ شرط پوری اے۔ سو $z \in A \cup B$۔ اسیں (۱) $z \in A$ تے (۲) $z \in A \cup B$ دوویں وکھا دتے نیں، ایس لئی «$\cap$» دی تعریف مطابق $z \in A \cap (A \cup B)$۔ ثبوت اوہناں تعریفاں دا ذکر نہیں کردا۔ ایہ منّیا گیا اے کہ پڑھن والے نوں اوہ پہلاں ای چنگی طرح سمجھ آ چکیاں نیں۔ جے نہیں، تاں تہانوں پچھے جا کے اوہ فیر چیتے کرنیاں پَین گئیاں۔ فیر ایہ وی پہچاننا پئے گا کہ $z
\in A$ توں $z \in A \cup B$ کیوں نکلدا اے، تے $z \in A$ تے $z \in A \cup B$ توں $z \in A \cap (A \cup B)$ کیوں نکلدا اے۔

اُتلے ثبوت دا اک ہور روپ ایہ اے، جیہدے وچ ہر گل کھولھ کے دسی اے:
\begin{proof}{}
[سیٹاں لئی $=$ دی تعریف مطابق، $A \cap (A \cup B) = A$ لئی سانوں وکھاؤنا اے کہ (ا) $A \cap (A \cup B) \subseteq A$ تے (ب) $A \cap (A \cup B)
  \subseteq A$۔ (ا) $\subseteq$ دی تعریف مطابق، سانوں وکھاؤنا اے کہ جے $z \in A \cap (A \cup B)$، تاں $z \in A$۔] جے $z \in A
\cap (A \cup B)$، تاں $z \in A$ [کیوں جے $\cap$ دی تعریف مطابق $z
  \in A \cap (A \cup B)$ اُدوں تے صرف اُدوں جدوں $z \in A$ تے $z \in A \cup B$]، سو $A
\cap (A \cup B) \subseteq A$۔ [(ب) $\subseteq$ دی تعریف مطابق، سانوں وکھاؤنا اے کہ جے $z \in A$، تاں $z \in A \cap (A \cup B)$۔] ہُن منّ لو [(۱)] $z \in A$۔ تاں [(۲)] $z \in A \cup B$ وی [کیوں جے (۱) مطابق $z \in A$ یا $z \in B$، جیہدا $\cup$ دی تعریف مطابق مطلب $z
  \in A \cup B$ اے]، تے ایس لئی $z \in A \cap (A \cup B)$ وی [کیوں جے $\cap$ دی تعریف لئی ضروری اے کہ $z \in A$، یعنی (۱)، تے $z
  \in A \cup B)$، یعنی (۲)]۔
\end{proof}

\begin{prob}
$A \cup (A \cap B) = A$ دے ہِیٹھلے ثبوت نوں کھولھ کے لکھو۔ سارے ورتے استنباط دے ڈھنگ دسو، ایہ دسو کہ ہر قدم مفروضیاں یا پہلاں ثابت دعویاں توں کیوں نکلدا اے، تے کتھے کیہڑی تعریف ورتنی پَندی اے۔
\begin{proof}
  جے $z \in A \cup (A \cap B)$ تاں $z \in A$ یا $z \in A \cap B$۔ جے $z \in A \cap B$، تاں $z \in A$۔ کوئی وی $z \in A$، $\in A \cup (A
  \cap B)$ وی اے۔
\end{proof}
\end{prob}
\endgroup
% /OLP-0610

% OLP-0611 content/methods/proofs/cant-do-it.tex
\begingroup
\def\olfilename{cant-do-it}
\expandafter\def\csname ol@sectioncs\endcsname{section}

\olfileid{mth}{prf}{cnt}

\olsection{میرے کولوں نہیں ہُندا!{}}

اسیں سارے کدی اوس تھاں پہنچ جاندے آں جتھے کم چھڈن نوں جی کردا اے۔ پر تسی ایہ کر \emph{سکدے} او۔ تہاڈے استاد، پڑھاؤن والے معاون تے جماعت دے ساتھی مدد کر سکدے نیں۔ اوہناں توں مدد منگو!{} ہِیٹھ کجھ گلّاں نیں جیہڑیاں اوکڑ توں بچن وچ مدد کرن گئیاں تے دسن گئیاں کہ جے کم چھڈن نوں جی کرے تاں کی کریے۔

مسئلے حل کر سکّن لئی ایہ کم کرو:
\begin{enumerate}
\item \emph{جِنّا ہو سکے پہلاں کم شروع کرو۔} سمسٹر وچ اسیں رُجھے رہندے آں تے ساڈے وچوں کئی کم اگے پاؤندے رہندے آں۔ تسی جیہڑے سبھ توں چنگے کم کر سکدے او اوہناں وچوں اک ایہ اے کہ گھر دا کم سویرے شروع کرو۔ ایویں جے پھس جاؤ، تاں حل لبھن لئی تہاڈے کول ویلا ہووے گا (رون توں علاوہ)۔
\item \emph{اپنی جماعت دے ساتھیاں نال گل کرو}۔ تسی اکلّے نہیں۔ جماعت وچ ہوراں نوں وی اوکڑ ہو سکدی اے، پر شاید وکھ وکھ گلاں وچ۔ ساتھیاں نال گل کرن نال مسئلہ کسے ہور پاسے توں دِس سکدا اے تے حل دا راہ کھل سکدا اے۔ بے شک اوہناں دا حل بس نقل نہ کرو۔ اوہناں توں اشارہ منگو، یا دسو کہ کتھے پھسے او تے اگلا قدم پُچھو۔ تے جدوں تہانوں سمجھ آ جائے، تسی وی اوہناں دی مدد کرو۔ کسے ہور نوں اگے ودھاؤن تے گل سمجھاؤن نال تہانوں آپ وی ہور سمجھ آوے گی۔
\item \emph{مدد منگو۔} مدد لئی تہاڈے کول کئی وسیلے نیں۔ تہاڈے استاد تے پڑھاؤن والے معاون تہاڈی مدد لئی نیں تے \emph{چاہندے} نیں کہ تسی کامیاب ہووو۔ اوہ تہانوں مسئلہ حل کرن وچ مدد دے سکن تے لبھ سکن کہ تہانوں کیہڑے قدم اُتے اوکڑ آؤندی اے۔
\item \emph{تھوڑا آرام کرو۔} جے تسی پھسے او، تاں \emph{ہو سکدا اے} تسی بڑی دیر توں مسئلے ول تکّ رہے او۔ تھوڑی دیر آرام کرو، چاء دا اک کپ پیو، یا کجھ دیر کسے ہور مسئلے اُتے کم کرو، تے فیر تازہ ذہن نال ایس مسئلے ول مُڑو۔ اک رات سَو کے فیر ویکھ لو۔
\end{enumerate}

دھیان دتا کہ ایہناں طریقیاں لئی ثبوت اُتے کم کافی پہلاں شروع کرنا ضروری اے؟ جے تسی جمع کراؤن توں پہلی رات دو وجے ثبوت شروع کیتا، تاں شاید ایہ اینے کم نہ آون۔

ایہ گل شاید بڑی مایوسی والی لگے، پر ثبوت لبھنا اک اوکھا کم اے جیہدا پھل آخر ملدا اے۔ کجھ لوک ایہ کم پیشے وجوں کردے نیں، سو ایہدے وچ کوئی سواد ضرور ہووے گا۔ لگ بھگ ہر کم وانگ مسئلے حل کرن تے ثبوت دین لئی مشق چاہیدی اے۔ تہاڈی جماعت وچ کجھ ساتھیاں نوں ایہ کم سوکھا لگ سکدا اے۔ اوہناں نے شاید پہلاں ای بہت مشق کیتی اے۔ چھیتی ہار نہ منّو۔

جے کسے مسئلے لئی تہاڈا ویلا (یا صبر) مُک جائے، تاں کوئی گل نہیں۔ ایس دا مطلب ایہ نہیں کہ تسی بے وقوف او یا تہانوں کدی سمجھ نہیں آوے گی۔ (اپنے استاد یا کسے ہور طالب علم توں) پتا کرو کہ کم کیویں ہُندا اے، تے لبھو کہ کتھے غلطی ہوئی یا کتھے پھسے سو، تاں جو اگلی واری ایہو جیہی اوکڑ آوے تاں اوہ گل فیر نہ ہووے۔ فیر حل ویکھے بناں ایہ کم کرن دی کوشش کرو۔ تے اگلی واری پہلاں شروع کرو تے پہلاں ای مدد منگو۔
\endgroup
% /OLP-0611

% OLP-0612 content/methods/proofs/resources.tex
\begingroup
\def\olfilename{resources}
\expandafter\def\csname ol@sectioncs\endcsname{section}

\olfileid{mth}{prf}{res}

\section{ہور وسیلے}

ریاضی وچ ثبوت دین بارے کئی کتاباں نیں جیہڑیاں تہاڈے کم آ سکدیاں نیں۔ خاص کر کے \emph{How to Read and do Proofs: An Introduction to
Mathematical Thought Processes} \citep{Solow2013} تے \emph{How to
Prove It: A Structured Approach} \citep{Velleman2019} ویکھو۔
\href{http://www.people.vcu.edu/~rhammack/BookOfProof/BookOfProof.pdf}{\emph{Book
of Proof}} \citep{Hammack2013} تے
\href{https://scholarworks.gvsu.edu/books/7/}{\emph{Mathematical
Reasoning}} \citep{Sandstrum2019} ثبوت بارے کتاباں نیں جیہڑیاں انٹرنیٹ اُتے مفت ملدیاں نیں۔ فلسفہ پڑھن والیاں نوں \emph{More
Precisely: The Math you need to do Philosophy} \citep{Steinhart2018} ریاضیاتی دلیل دی چنگی مُڈھلی کتاب لگ سکدی اے۔

انٹرنیٹ اُتے ثبوتاں دیاں کئی ہور مختصر رہنما لکھتاں وی ملدیاں نیں، مثلاً
\href{https://math.berkeley.edu/~hutching/teach/proofs.pdf}{«Introduction
  to Mathematical Arguments»} \citep{Hutchings2003} تے
\href{https://eugeniacheng.com/wp-content/uploads/2017/02/cheng-proofguide.pdf}{«How to
  write proofs»} \citep{Cheng2004}۔

\subsection{حوصلہ دین والیاں ویڈیوز}

گھر دا کم کرن نوں جی نہیں کردا؟ دل اُداس اے؟ ایہ ویڈیوز کم آ سکدیاں نیں!

\begin{itemize}
\item \url{https://www.youtube.com/watch?v=ZXsQAXx_ao0}
\item \url{https://www.youtube.com/watch?v=BQ4yd2W50No}
\item \url{https://www.youtube.com/watch?v=StTqXEQ2l-Y}
\end{itemize}
\endgroup
% /OLP-0612

\OLEndChapterHook
\endgroup
% /OLP-0602

% OLP-0613 content/methods/induction/induction.tex
\begingroup
\def\olfilename{induction}
\expandafter\def\csname ol@sectioncs\endcsname{section}

\olchapter{mth}{ind}{استقرا}

% OLP-0614 content/methods/induction/introduction.tex
\begingroup
\def\olfilename{introduction}
\expandafter\def\csname ol@sectioncs\endcsname{section}

\olfileid{mth}{ind}{int}

\olsection{جان پہچان}

استقرا ثبوت دا اک ضروری طریقہ اے جیہڑا وکھ وکھ صورتاں وچ منطق، نظری کمپیوٹر سائنس، تے ریاضی دے لگ بھگ سارے شعبیاں وچ ورتیا جاندا اے۔ منطق دے کئی نتیجے ثابت کرن لئی ایہدی لوڑ اے۔

استقرا نوں اکثر استخراج دے مقابلے وچ رکھیا جاندا اے تے خاص صورتاں توں عام دعوے تک پہنچن والا استنباط دسیا جاندا اے۔ مثلاً جے اسیں کئی ہرے زمرد ویکھیے تے کوئی ایہو جیہی شے نہ ویکھیے جیہنوں زمرد آکھیے پر اوہ ہری نہ ہووے، تاں اسیں نتیجہ کڈھ سکدے آں کہ سارے زمرد ہرے نیں۔ ایہ استقرائی استنباط اے، کیوں جے ایہ کئی خاص صورتاں (ایہ زمرد ہرا اے، اوہ زمرد ہرا اے، وغیرہ) توں اک عام دعوے (سارے زمرد ہرے نیں) تک جاندا اے۔ \emph{ریاضیاتی} استقرا وی اوہ استنباط اے جیہدا نتیجہ عام دعویٰ اے، پر ایہ ایس «سادے استقرا» توں بڑی وکھری قسم دا اے۔

موٹی گل ایہ اے کہ ریاضی وچ استقرائی ثبوت ایہ نتیجہ کڈھدا اے کہ اک خاص قسم دیاں ساریاں ریاضیاتی شےواں وچ کوئی خاصیت اے۔ سبھ توں سوکھی صورت وچ استقرائی ثبوت جیہناں شےواں بارے اے، اوہ قدرتی عدد نیں۔ ایس صورت وچ استقرائی ثبوت ایہ ثابت کرن لئی ورتیا جاندا اے کہ سارے قدرتی عدداں وچ کوئی خاصیت اے۔ ایہ کم ایویں ہُندا اے:
\begin{enumerate}
    \item $0$ وچ ایہ خاصیت اے، تے
    \item جدوں وی کسے عدد~$k$ وچ ایہ خاصیت اے، تاں~$k+1$ وچ وی اے۔
\end{enumerate}
قدرتی عدداں اُتے استقرا نوں اوہناں ریاضیاتی شےواں بارے عام دعوے ثابت کرن لئی وی اکثر ورت سکدے آں جیہناں نوں عدد دتے جا سکدے نیں۔ مثلاً ہر محدود سیٹ دے رکناں دی محدود گنتی~$n$ اے۔ جے اسیں استقرا نال وکھا سکیے کہ ہر عدد~$n$ وچ ایہ خاصیت اے کہ «وڈائی~$n$ والے سارے محدود سیٹ \dots»، تاں اسیں سارے محدود سیٹاں بارے کوئی گل وکھا دتی ہووے گی۔

استقرا نوں اوہناں ریاضیاتی شےواں تک وی ودھایا جا سکدا اے جیہڑیاں \emph{استقرائی طریقے نال تعریف} کیتیاں گئیاں نیں۔ مثلاً رسمی زبان دیاں عبارتاں، جیویں پہلے درجے دی منطق دیاں عبارتاں، استقرائی طریقے نال تعریف ہُندیاں نیں۔ \emph{ساختی استقرا} ایہو جیہیاں ساریاں عبارتاں بارے نتیجے ثابت کرن دا اک طریقہ اے۔ خاص کر کے ساختی استقرا منطق وچ بڑا کم آؤندا اے تے بہت ورتیا جاندا اے۔
\endgroup
% /OLP-0614

% OLP-0615 content/methods/induction/induction-on-N.tex
\begingroup
\def\olfilename{induction-on-N}
\expandafter\def\csname ol@sectioncs\endcsname{section}

\olfileid{mth}{ind}{inN}

\olsection{\texorpdfstring{~$\Nat$ اُتے استقرا}{ℕ اُتے استقرا}}

اپنی سبھ توں سوکھی صورت وچ استقرا سارے قدرتی عدداں لئی نتیجے ثابت کرن دا طریقہ اے۔ ایہ ایس گل نوں ورتدا اے کہ $0$ توں شروع کر کے وار وار~$1$ پائیے تاں آخر ہر قدرتی عدد تک پہنچ جاندے آں۔ سو ہر عدد لئی کوئی گل سچ ثابت کرن واسطے اسیں (۱) ثابت کر سکدے آں کہ ایہ $0$ لئی سچ اے، تے (۲) وکھا سکدے آں کہ جدوں وی ایہ کسے عدد~$n$ لئی سچ اے، تاں اگلے عدد~$n+1$ لئی وی سچ اے۔ جے اسیں «عدد~$n$ وچ خاصیت $P$ اے» نوں مختصر $P(n)$ لکھیے (تے «عدد~$k$ وچ خاصیت $P$ اے» نوں $P(k)$، وغیرہ)، تاں ایہ ثابت کرن لئی کہ $P(n)$ سارے $n \in \Nat$ لئی، استقرائی ثبوت وچ ایہ ہُندا اے:
\begin{enumerate}
\item $P(0)$ دا ثبوت، تے
\item ایہدا ثبوت کہ کسے وی~$k$ لئی، جے $P(k)$ تاں $P(k+1)$۔
\end{enumerate}
گل پوری صاف کرن لئی منّ لو ساڈے کول (۱) تے~(۲) دوویں نیں۔ (۱) دسدا اے کہ $P(0)$ سچ اے۔ جے~(۲) وی اے، تاں سانوں خاص کر کے پتا اے کہ جے $P(0)$ تاں $P(0+1)$، یعنی $P(1)$۔ ایہ عام بیان «کسے وی~$k$ لئی، جے $P(k)$ تاں $P(k+1)$» توں~$0$ نوں~$k$ دی تھاں رکھ کے نکلدا اے۔ سو شرطیہ تے اوہدے مقدم توں نتیجہ کڈھن دے قاعدے نال سانوں~$P(1)$ ملدا اے۔ (۲) توں فیر، ہُن $1$ نوں~$n$ دی تھاں رکھ کے، سانوں ملدا اے کہ جے $P(1)$ تاں~$P(2)$۔ کیوں جے اسیں ہُن ای~$P(1)$ ثابت کیتا اے، اوسے قاعدے نال سانوں~$P(2)$ ملدا اے۔ تے ایس طرح اگے۔ کسے وی عدد~$n$ لئی، ایہ کم $n$~واری کر کے آخر~$P(n)$ تک پہنچدے آں۔ سو (۱) تے~(۲) رَل کے~$P(n)$ ثابت کردے نیں کسے وی $n
\in \Nat$ لئی۔

اک مثال ویکھیے۔ منّ لو اسیں پتا لاؤنا چاہندے آں کہ $n$ پاسے سُٹ کے کِنّے وکھرے جوڑ کڈھ سکدے آں۔ بھاویں ایہ ذرا اوپرا لگے، آؤ $0$ پاسیاں توں شروع کریے۔ جے تہاڈے کول کوئی پاسا نہیں، تاں بس اک ممکن جوڑ «سُٹ» سکدے او۔ کوئی نقطہ نہیں، جیہدا جوڑ~$0$ اے۔ سو وکھرے ممکن حاصلاں دی گنتی~$1$ اے۔ جے اکّو پاسا اے، یعنی $n=1$، تاں چھ ممکن قدراں نیں، $1$ توں~$6$ تک۔ دو پاسیاں نال $2$ توں $12$ تک کوئی وی جوڑ سُٹ سکدے آں، یعنی $11$ صورتاں۔ تِن پاسیاں نال $3$ توں $18$ تک کوئی وی عدد سُٹ سکدے آں، یعنی $16$ وکھریاں صورتاں۔ $1$، $6$، $11$، $16$، اک ترتیب لگدی اے۔ شاید جواب $5n+1$ اے؟ بے شک $5n+1$ سبھ توں ودھ ممکن گنتی اے، کیوں جے بس $5n+1$ عدد نیں $n$ توں، جیہڑی $n$ پاسیاں نال سبھ توں گھٹ قدر اے (سارے $1$)، $6n$ تک، جیہڑی سبھ توں ودھ قدر اے (سارے $6$)۔

\begin{thm}
  $n$ پاسیاں نال ساریاں $5n+1$ ممکن قدراں سُٹ سکدے آں جیہڑیاں $n$ تے $6n$ دے وچکار نیں۔
\end{thm}

\begin{proof}
$P(n)$ ایہ دعویٰ ہووے، «$n$ تے~$6n$ دے وچکار کوئی وی عدد $n$~پاسیاں نال سُٹ سکدے آں»۔ استقرا ورتن لئی ایہ ثابت کردے آں:
\begin{enumerate}
\item \emph{استقرا دی بنیاد} $P(1)$، یعنی اکّو پاسے نال $1$ تے $6$ دے وچکار کوئی وی عدد سُٹ سکدے او۔
\item \emph{استقرا دا قدم}، سارے $k$ لئی، جے $P(k)$ تاں~$P(k+1)$۔
\end{enumerate}

(۱) اک $6$ رُخاں والا پاسا ویکھ کے ثابت ہُندا اے۔ اوہدے سارے 6 رُخ نیں، تے $1$ توں~$6$ تک ہر عدد کسے اک رُخ اُتے اے۔ سو اکّو پاسے نال $1$ تے~$6$ دے وچکار کوئی وی عدد سُٹ سکدے آں۔

(۲) ثابت کرن لئی اسیں شرطیے دا مقدم منّدے آں، یعنی $P(k)$۔ ایس منّی ہوئی گل نوں \emph{استقرائی مفروضہ} آکھدے نیں۔ اسیں ایہنوں~$P(k+1)$ ثابت کرن لئی ورتدے آں۔ اوکھا حصہ ایہ لبھنا اے کہ $k+1$ پاسیاں دے سُٹّن دیاں ممکن قدراں نوں، $k$~پاسیاں دے سُٹّن تے وادھو $k+1$ ویں پاسے دے سُٹّن دیاں قدراں نال کیویں جوڑیے۔ جے استقرائی مفروضہ ورتنا اے تاں ایہ کم کرنا پئے گا۔

استقرائی مفروضہ آکھدا اے کہ $k$ تے~$6k$ دے وچکار کوئی وی عدد $k$~پاسیاں نال کڈھ سکدے آں۔ جے~$1$ اپنے $(k+1)$ ویں پاسے نال سُٹیے، تاں کُل جوڑ وچ $1$ ودھ جاندا اے۔ سو $k+1$ توں $6k+1$ تک کوئی وی قدر سُٹ سکدے آں، $k$~پاسے سُٹ کے تے فیر~$1$ نوں $(k+1)$ ویں پاسے نال سُٹ کے۔ کیہ رہ گیا؟ $6k+2$ توں $6k+6$ تک دیاں قدراں۔ ایہناں لئی $k$ واری $6$ سُٹو تے فیر $2$ تے $6$ دے وچکار کوئی عدد اپنے $(k+1)$ ویں پاسے نال سُٹو۔ رَل کے ایہدا مطلب اے کہ $k+1$ پاسیاں نال $k+1$ تے $6(k+1)$ دے وچکار کوئی وی عدد سُٹ سکدے آں، یعنی اسیں~$P(k+1)$ مفروضے~$P(k)$، یعنی استقرائی مفروضے، نال ثابت کر دتا اے۔
\end{proof}

اسیں استقرا اکثر اوہدوں ورتدے آں جدوں شےواں (عدد، سیٹ، وغیرہ) دی کسے لڑی بارے کوئی گل ثابت کرنی ہووے تے لڑی آپ «استقرائی طریقے نال» تعریف ہوئی ہووے، یعنی $(n+1)$ ویں شے دی تعریف $n$ ویں شے نال دتی ہووے۔ مثلاً قدرتی عدداں دا جوڑ~$s_n$،~$n$ تک، ایویں تعریف کر سکدے آں:
\begin{align*}
  s_0 & = 0\\
  s_{n+1} & = s_n + (n+1)
\end{align*}
ایس تعریف توں ملدا اے:
\begin{align*}
  s_0 & = 0,\\
  s_1 & = s_0 + 1 && = 1,\\
  s_2 & = s_1 + 2 && = 1 + 2 = 3\\
  s_3 & = s_2 + 3 && = 1 + 2 + 3 = 6, \text{ وغیرہ۔}
\end{align*}
ہُن اسیں استقرا نال ثابت کر سکدے آں کہ $s_n = n(n+1)/2$۔

\begin{prop}
  $s_n = n(n+1)/2$.
\end{prop}

\begin{proof}
  سانوں ثابت کرنا اے (۱) کہ $s_0 = 0\cdot(0 + 1)/2$، تے (۲) جے $s_k =
  k(k+1)/2$ تاں $s_{k+1} = (k+1)(k+2)/2$۔ (۱) صاف اے۔ (۲) ثابت کرن لئی اسیں استقرائی مفروضہ منّدے آں، $s_k = k(k+1)/2$۔ ایہ ورت کے سانوں وکھاؤنا اے کہ $s_{k+1} = (k+1)(k+2)/2$۔

  $s_{k+1}$ کیہ اے؟ تعریف مطابق $s_{k+1} = s_k + (k+1)$۔ استقرائی مفروضے مطابق $s_k = k(k+1)/2$۔ ایہ پچھلی مساوات وچ پا سکدے آں، تے فیر بس کسراں دا تھوڑا حساب چاہیدا اے:
  \begin{align*}
    s_{k+1} & = \frac{k(k+1)}{2} + (k+1) = {}\\
    & = \frac{k(k+1)}{2} + \frac{2(k+1)}{2} = {}\\
    & = \frac{k(k+1) + 2(k+1)}{2} = {}\\
    & = \frac{(k+2)(k+1)}{2}.
  \end{align*}
\end{proof}

ایتھے ضروری گل ایہ اے کہ جے تسی استقرائی طریقے نال تعریف ہوئی کسے لڑی $a_n$ بارے کوئی گل ثابت کردے او، تاں استقرا ورتنا سِدھا طریقہ اے۔ جے ایہ گل سِدھی نہ دِسے (جیویں پاسے سُٹّن دیاں ممکن صورتاں وچ)، تاں وی استقرا ورت سکدے او جے~$k+1$ والی صورت نوں~$k$ والی صورت نال کسے طرح جوڑ سکو۔
\endgroup
% /OLP-0615

% OLP-0616 content/methods/induction/strong-induction.tex
\begingroup
\def\olfilename{strong-induction}
\expandafter\def\csname ol@sectioncs\endcsname{section}

\olfileid{mth}{ind}{str}

\olsection{مضبوط استقرا}

اُتے دسے استقرا دے اصول وچ اسیں $P(0)$ ثابت کردے آں تے ایہ وی کہ جے $P(k)$ تاں $P(k+1)$۔ دوجے حصے وچ منّدے آں کہ $P(k)$ سچ اے تے ایس مفروضے نال $P(k+1)$ ثابت کردے آں۔ ایہو گل ثابت کرن لئی اسیں $P(k-1)$ منّ کے اوہنوں $P(k)$ ثابت کرن لئی وی ورت سکدے آں۔ ضروری حصہ ایہ اے کہ کسے وی عدد توں اوہدے جانشین تک استنباط کر سکیے، یعنی کسے وی عدد لئی دعویٰ ثابت کر سکیے، ایس شرط نال کہ دعویٰ اوہتوں پہلاں والے عدد لئی پورا ہُندا اے۔

استقرا دے اصول دا اک ہور روپ اے جیہدے وچ اسیں بس ایہ نہیں منّدے کہ دعویٰ $k-1$ لئی، جیہڑا $k$ توں پہلاں اے، پورا ہُندا اے۔ اسیں منّدے آں کہ ایہ~$k$ توں چھوٹے سارے عدداں لئی پورا ہُندا اے، تے ایس مفروضے نال~$k$ لئی دعویٰ ثابت کردے آں۔ ایہ وی سانوں $P(n)$ سارے~$n \in
\Nat$ لئی دیندا اے۔ کیوں جے جدوں $P(0)$ ثابت کر لئیے، تاں اسیں ثابت کر لیا اے کہ $P$،~$1$ توں چھوٹے سارے عدداں لئی پورا ہُندا اے۔ تے جے پتا اے کہ جے $P(l)$ سارے $l<k$ لئی، تاں $P(k)$، سو ایہ گل خاص کر کے $k=1$ لئی پتا اے۔ ایس لئی اسیں $P(1)$ کڈھ سکدے آں۔ ایس نال اسیں $P(0)$ تے $P(1)$ ثابت کر لئے، یعنی $P(l)$ سارے $l<2$ لئی۔ تے ساڈے کول شرطیہ وی اے کہ جے $P(l)$ سارے $l<2$ لئی تاں $P(2)$، سو اسیں~$P(2)$ کڈھ سکدے آں، تے ایس طرح اگے۔

اصل وچ جے اسیں عام شرطیہ «سارے $k$ لئی، جے $P(l)$ سارے $l<k$ لئی تاں $P(k)$» ثابت کر سکیے، تاں $P(0)$ ہور ثابت کرن دی لوڑ نہیں، کیوں جے ایہ اوہتوں نکلدا اے۔ چیتے کرو کہ «سارے $l<k$ لئی، $P(l)$» ورگا عام دعویٰ سچ اے جے کوئی $l<k$ نہیں۔ ایہ خالی صورت وچ مسوّر دی ورتوں اے۔ «سارے $A$، $B$ نیں» سچ اے جے کوئی $A$ نہیں، تے $\lforall[x][(!A(x) \lif !B(x))]$ سچ اے جے کوئی $x$،~$!A(x)$ اُتے پورا نہیں اتردا۔ ایس صورت وچ رسمی روپ «$\lforall[l][(l < k \lif P(l))]$» ہووے گا، تے ایہ سچ اے جے کوئی $l < k$ نہیں۔ تے جے $k=0$، تاں ایہو صورت اے۔ کوئی $l<0$ نہیں، ایس لئی «سارے $l<0$ لئی، $P(0)$» سچ اے، $P$ بھاویں جو وی ہووے۔ سو «جے $P(l)$ سارے $l<k$ لئی تاں $P(k)$» دا ثبوت آپے~$P(0)$ ثابت کردا اے۔

ایہ روپ اوہدوں کم آؤندا اے جدوں~$k$ لئی دعویٰ ثابت کرنا بس $k-1$ والے دعوے اُتے نہ ٹک سکے، بلکہ ایہ منّن دی لوڑ پئے کہ اوہ اک یا کئی $l<k$ لئی سچ اے۔
\endgroup
% /OLP-0616

% OLP-0617 content/methods/induction/inductive-definitions.tex
\begingroup
\def\olfilename{inductive-definitions}
\expandafter\def\csname ol@sectioncs\endcsname{section}

\olfileid{mth}{ind}{idf}

\olsection{استقرائی تعریفاں}

منطق وچ اسیں اکثر شےواں دیاں قسماں \emph{استقرائی طریقے نال} تعریف کردے آں۔ یعنی اسیں قاعدے دیندے آں کہ کیہڑی شے ایس قسم دی منّی جائے، تے قاعدے سمجھاؤندے نیں کہ اوس قسم دیاں پرانیاں شےواں توں اوسے قسم دیاں نویاں شےواں کیویں بنن۔ مثلاً اسیں نشاناں دیاں خاص لڑیواں، جیویں کسے زبان دے ٹرم تے !!{formula}s، اکثر استقرا نال تعریف کردے آں۔ اک سوکھی مثال لئی اوہ لڑیواں ویکھو جیہڑیاں اکھراں $\mathrm{a}$، $\mathrm{b}$، $\mathrm{c}$، $\mathrm{d}$، نشان~$\circ$، تے بریکٹاں $[$ تے $]$ نال بندیاں نیں، جیویں «$[[\mathrm{c} \circ \mathrm{d}][$»، «$[\mathrm{a}[]\circ]$»، «$\mathrm{a}$»، یا «$[[\mathrm{a} \circ \mathrm{b}]\circ
\mathrm{d}]$»۔ تہانوں شاید لگے کہ پہلیاں دوہاں لڑیواں وچ کوئی «غلطی» اے۔ پہلی وچ بریکٹاں دا میل بالکل نہیں، تے تہانوں لگ سکدا اے کہ «$\circ$» نوں اوہ عبارتاں «جوڑنیاں» چاہیدیاں نیں جیہڑیاں آپ ٹھیک ہون۔ تیجی تے چوتھی لڑی ٹھیک لگدیاں نیں۔ ہر «$[$» لئی اک بند کرن والا «$]$» اے (جے کوئی بریکٹ ہے وی)، تے ہر $\circ$ دے دوہاں پاسیاں اُتے «ٹھیک» عبارتاں لبھدیاں نیں، جیہناں دے دوالے بریکٹاں دی اک جوڑی اے۔

اسیں این ٹھیک دسنا چاہندے آں کہ «ٹھیک ٹرم» کیہنوں آکھیے۔ پہلاں، ہر اکھر آپے ٹھیک اے۔ جیہڑی شے بس اکلّا اکھر نہیں، اوہ «$[t \circ s]$» دی صورت وچ ہووے، جتھے $s$ تے $t$ آپ ٹھیک ہون۔ اُلٹے پاسے، جے $t$ تے $s$ ٹھیک نیں، تاں اوہناں دے وچکار $\circ$ رکھ کے تے دوالے بریکٹاں دی اک جوڑی لا کے اسیں نواں ٹھیک ٹرم بنا سکدے آں۔ ایہ عمل ورت کے اسیں ٹھیک ٹرماں دا سیٹ \emph{تعریف} کر سکدے آں۔ ایہ \emph{استقرائی تعریف} اے۔

\begin{defn}[ٹھیک ٹرم]
  \emph{ٹھیک ٹرماں} دا سیٹ استقرائی طریقے نال ایویں تعریف ہُندا اے:
  \begin{enumerate}
  \item کوئی وی اکھر $\mathrm{a}$، $\mathrm{b}$، $\mathrm{c}$، $\mathrm{d}$ اک ٹھیک ٹرم اے۔
  \item جے $s_1$ تے $s_2$ ٹھیک ٹرم نیں، تاں $[s_1 \circ s_2]$ وی ٹھیک ٹرم اے۔
  \item ہور کوئی شے ٹھیک ٹرم نہیں۔
  \end{enumerate}
\end{defn}

ایہ تعریف دسدی اے کہ کوئی شے اُدوں تے صرف اُدوں ٹھیک ٹرم اے جدوں اوہ دو شرطاں (۱) تے~(۲) مطابق محدود گنتی دے قدماں وچ بن سکے۔ پہلے قدم وچ اسیں اوہ سارے ٹھیک ٹرم بناؤندے آں جیہڑے بس اکلّے اکھر نیں، یعنی:
\[
\mathrm{a}, \mathrm{b}, \mathrm{c}, \mathrm{d}
\]
دوجے قدم وچ اسیں اپنے بنائے ٹرماں اُتے (۲) لاؤندے آں۔ سانوں ایہ ملن گے:
\[
[\mathrm{a} \circ \mathrm{a}], [\mathrm{a} \circ \mathrm{b}],
[\mathrm{b} \circ \mathrm{a}], \dots, [\mathrm{d} \circ \mathrm{d}]
\]
دو اکھراں دے سارے جوڑاں لئی۔ تیجے قدم وچ اسیں (۲) فیر ہُن تک بنائے کسے وی دو ٹھیک ٹرماں اُتے لاؤندے آں۔ سانوں نویں ٹھیک ٹرم ملدے نیں، جیویں $[\mathrm{a} \circ [\mathrm{a} \circ
    \mathrm{a}]]$، جتھے $t$ پہلے قدم والا $\mathrm{a}$ اے تے $s$ دوجے قدم والا $[\mathrm{a} \circ \mathrm{a}]$ اے، تے $[[\mathrm{b} \circ
    \mathrm{c}] \circ [\mathrm{d} \circ \mathrm{b}]]$، جیہڑا دوجے قدم دے دو ٹرماں $[\mathrm{b} \circ \mathrm{c}]$ تے $[\mathrm{d}
  \circ \mathrm{b}]$ توں بنیا اے۔ تے ایس طرح اگے۔ شق (۳) ایہ روکدی اے کہ جیہڑی شے ایویں نہیں بنی، اوہ ٹھیک ٹرماں دے سیٹ وچ آ وڑے۔

دھیان دیو کہ اسیں ہن تک ایہ ثابت نہیں کیتا کہ نشاناں دی ہر «ٹھیک لگدی» لڑی ایس تعریف مطابق ٹھیک اے۔ پر ایہ صاف ہونا چاہیدا اے کہ اسیں جیہڑی وی لڑی بنا سکدے آں اوہ واقعی «ٹھیک لگدی» اے۔ بریکٹاں دا میل اے تے $\circ$ اوہ حصے جوڑدا اے جیہڑے آپ ٹھیک نیں۔

استقرائی تعریفاں دی خاص گل ایہ اے کہ جے تسی سارے ٹھیک ٹرماں بارے کوئی گل ثابت کرنا چاہندے او، تاں تعریف دسدی اے کہ کیہڑیاں صورتاں ویکھنیاں پَین گئیاں۔ مثلاً جے تہانوں دسیا جائے کہ $t$ اک ٹھیک ٹرم اے، تاں استقرائی تعریف دسدی اے کہ $t$ کیہو جیہا ہو سکدا اے۔ $t$ اک اکھر ہو سکدا اے، یا $[s_1 \circ s_2]$، کسے دو ٹھیک ٹرماں $s_1$ تے~$s_2$ لئی۔ شق (۳) کرکے بس ایہو صورتاں ممکن نیں۔

جدوں استقرائی طریقے نال تعریف ہوئے سیٹ دے سارے رکناں بارے دعویٰ ثابت کرنا ہووے، تاں مضبوط استقرا خاص ضروری ہو جاندا اے۔ مثلاً منّ لو اسیں ثابت کرنا چاہندے آں کہ لمبائی~$n$ والے ہر ٹھیک ٹرم وچ $[$ دی گنتی~$< n/2$ اے۔ ایس نوں سارے~$n$ بارے دعویٰ سمجھ سکدے آں۔ ہر $n$ لئی، $[$ دی گنتی، لمبائی~$n$ والے کسے وی ٹھیک ٹرم وچ، $< n/2$ اے۔

\begin{prop}
  کسے وی $n$ لئی، $[$ دی گنتی، لمبائی~$n$ والے کسے ٹھیک ٹرم وچ، $< n/2$ اے۔
\end{prop}

\begin{proof}
ایہ نتیجہ (مضبوط) استقرا نال ثابت کرن لئی سانوں وکھاؤنا پئے گا کہ ہِیٹھلا شرطیہ دعویٰ سچ اے:
\begin{quote}
  جے ہر $l < k$ لئی، لمبائی~$l$ والے کسے وی ٹھیک ٹرم وچ $< l/2$ وار $[$ آؤندا اے، تاں لمبائی~$k$ والے کسے وی ٹھیک ٹرم وچ $< k/2$ وار $[$ آؤندا اے۔
\end{quote}
ایہ شرطیہ وکھاؤن لئی اوہدا مقدم سچ منّو، یعنی منّ لو کہ کسے وی $l<k$ لئی، لمبائی~$l$ والے ٹھیک ٹرماں وچ $< l/2$ وار $[$ آؤندا اے۔ ایس منّی ہوئی گل نوں استقرائی مفروضہ آکھدے آں۔ اسیں وکھاؤنا چاہندے آں کہ ایہو گل لمبائی~$k$ والے ٹھیک ٹرماں لئی وی سچ اے۔

سو منّ لو $t$ لمبائی~$k$ والا اک ٹھیک ٹرم اے۔ کیوں جے ٹھیک ٹرم استقرائی طریقے نال تعریف ہوئے نیں، ساڈے کول دو صورتاں نیں۔ (۱)~$t$ بس اک اکھر اے، یا (۲)~$t$، $[s_1 \circ s_2]$ اے، کسے ٹھیک ٹرماں $s_1$ تے~$s_2$ لئی۔
\begin{enumerate}
\item $t$ اک اکھر اے۔ تاں $k = 1$ تے $[$ دی گنتی $t$ وچ~$0$ اے۔ کیوں جے $0 < 1/2$، دعویٰ پورا ہُندا اے۔
\item $t$، $[s_1 \circ s_2]$ اے، کسے ٹھیک ٹرماں $s_1$ تے~$s_2$ لئی۔ $l_1$،~$s_1$ دی لمبائی تے $l_2$،~$s_2$ دی لمبائی ہووے۔ تاں لمبائی~$k$، جیہڑی $t$ دی اے، $l_1+l_2+3$ اے ($s_1$ تے $s_2$ دیاں لمبائیاں تے تِن نشان $[$، $\circ$، $]$)۔ کیوں جے $l_1+l_2+3$ ہمیشہ $l_1$ توں وڈا اے، $l_1 < k$۔ ایویں ای $l_2
  < k$۔ ایہدا مطلب اے کہ استقرائی مفروضہ ٹرماں $s_1$ تے~$s_2$ اُتے لاگو ہُندا اے۔ گنتی~$m_1$ جیہڑی $[$ دی $s_1$ وچ اے، $< l_1/2$ اے، تے گنتی~$m_2$ جیہڑی $[$ دی~$s_2$ وچ اے، $< l_2/2$ اے۔

  $[$ دی گنتی $t$ وچ، $[$ دی~$s_1$ والی گنتی تے $[$ دی~$s_2$ والی گنتی دا جوڑ، نال~$1$ اے، یعنی $m_1 + m_2 + 1$۔ کیوں جے $m_1
  < l_1/2$ تے $m_2 < l_2/2$، سانوں ملدا اے:
  \[
  m_1 + m_2 + 1 < \frac{l_1}{2} + \frac{l_2}{2} + 1 = \frac{l_1+l_2+2}{2} < \frac{l_1+l_2+3}{2} = k/2.
  \]
\end{enumerate}
ہر صورت وچ اسیں وکھا دتا اے کہ $[$ دی گنتی $t$ وچ $< k/2$ اے (استقرائی مفروضے دی بنیاد اُتے)۔ مضبوط استقرا نال قضیہ نکلدا اے۔
\end{proof}

\begin{prob}
  خاص ٹھیک ٹرماں دا سیٹ ایویں تعریف کرو:
  \begin{enumerate}
  \item کوئی وی اکھر $\mathrm{a}$، $\mathrm{b}$، $\mathrm{c}$، $\mathrm{d}$ اک خاص ٹھیک ٹرم اے۔
  \item جے $s$ اک خاص ٹھیک ٹرم اے، تاں $[s]$ وی خاص ٹھیک ٹرم اے۔
  \item جے $s_1$ تے $s_2$ خاص ٹھیک ٹرم نیں، تاں $[s_1 \circ s_2]$ اک خاص ٹھیک ٹرم اے۔
  \item ہور کوئی شے خاص ٹھیک ٹرم نہیں۔
  \end{enumerate}
  وکھاؤ کہ $[$ دی گنتی، کسے خاص ٹھیک ٹرم~$t$ وچ جیہدی لمبائی~$n$ اے، $\le n/2 +1$ اے۔
\end{prob}
\endgroup
% /OLP-0617

% OLP-0618 content/methods/induction/structural-induction.tex
\begingroup
\def\olfilename{structural-induction}
\expandafter\def\csname ol@sectioncs\endcsname{section}

\olfileid{mth}{ind}{sti}

\olsection{ساختی استقرا}

ہن تک اسیں سارے قدرتی عدداں بارے نتیجے ثابت کرن لئی استقرا ورتیا اے۔ پر ایس ورگا اصول استقرائی طریقے نال تعریف ہوئے سیٹ دے سارے !!{element}s بارے نتیجے سِدھے ای ثابت کرن لئی ورتیا جا سکدا اے۔ اینوں اکثر \emph{ساختی} استقرا آکھدے نیں، کیوں جے ایہ استقرائی طریقے نال تعریف ہوئیاں شےواں دی ساخت اُتے ٹکدا اے۔

عام طور اُتے استقرائی تعریف ایویں دتی جاندی اے۔ (ا) سیٹ دے «مُڈھلے» !!{element}s دی فہرست تے (ب) اوہناں عملاں دی فہرست جیہڑے پرانے رکناں توں سیٹ دے نویں !!{element}s بناؤندے نیں۔ مثلاً ٹھیک ٹرماں دی صورت وچ مُڈھلیاں شےواں اکھر نیں۔ ساڈے کول بس اک عمل اے، ایہ عمل ایہ اے:
\begin{align*}
  o(s_1, s_2) = & [s_1 \circ s_2]
\end{align*}
تسی قدرتی عدداں~$\Nat$ نوں وی استقرائی تعریف نال دتے ہوئے سمجھ سکدے او۔ مُڈھلی شے~$0$ اے تے عمل جانشین فنکشن~$x + 1$ اے۔

استقرائی طریقے نال تعریف ہوئے سیٹ دے سارے رکناں بارے کوئی گل ثابت کرن لئی، یعنی سیٹ دے ہر !!{element} وچ خاصیت~$P$ اے، سانوں ایہ کرنا پئے گا:
\begin{enumerate}
\item ثابت کرو کہ مُڈھلیاں شےواں وچ~$P$ اے۔
\item ثابت کرو کہ ہر عمل~$o$ لئی، جے اوہدے آرگومنٹاں وچ~$P$ اے، تاں نتیجے وچ وی اے۔
\end{enumerate}
مثلاً سارے ٹھیک ٹرماں بارے کوئی گل ثابت کرن لئی، اسیں ثابت کراں گے کہ ایہ سارے اکھراں بارے سچ اے، تے $[s_1 \circ s_2]$ بارے سچ اے بشرطیکہ $s_1$ تے $s_2$ بارے وکھ وکھ سچ ہووے۔

\begin{prop}
  $[$ دی گنتی $]$ دی گنتی برابر اے، کسے وی ٹھیک ٹرم~$t$ وچ۔
\end{prop}

\begin{proof}
اسیں ساختی استقرا ورتدے آں۔ ٹھیک ٹرم استقرائی طریقے نال تعریف ہوئے نیں، اکھر مُڈھلیاں شےواں نیں تے عمل $o$ پرانیاں توں نویں ٹھیک ٹرم بناؤندا اے۔
\begin{enumerate}
\item دعویٰ ہر اکھر لئی سچ اے، کیوں جے اکلّے اکھر وچ $[$ دی گنتی~$0$ اے تے $]$ دی گنتی وی~$0$ اے۔
\item منّ لو $[$ دی گنتی $s_1$ وچ $]$ دی گنتی برابر اے، تے ایہو گل $s_2$ لئی وی سچ اے۔ $[$ دی گنتی $o(s_1, s_2)$ وچ، یعنی $[s_1 \circ s_2]$ وچ، $[$ دیاں $s_1$ تے $s_2$ وچ دیاں گنتیاں دا جوڑ تے نال اک اے۔ $]$ دی گنتی $o(s_1, s_2)$ وچ، $]$ دیاں $s_1$ تے $s_2$ وچ دیاں گنتیاں دا جوڑ تے نال اک اے۔ سو $[$ دی گنتی $o(s_1, s_2)$ وچ $]$ دی گنتی $o(s_1,s_2)$ وچ دے برابر اے۔
\end{enumerate}
\end{proof}

\begin{prob}
  ساختی استقرا نال ثابت کرو کہ کوئی ٹھیک ٹرم~$]$ نال شروع نہیں ہُندا۔
\end{prob}

آؤ ساختی استقرا نال اک ہور ثبوت دیئیے۔ نشاناں دی لڑی~$t$ دا حقیقی مُڈھلا ٹکڑا اوہ لڑی~$s$ اے جیہڑی $t$ نال، کھبّے توں پڑھدیاں، نشان بہ نشان ملدی اے، پر $t$ لمّی اے۔ سو مثلاً $[a \circ {}$، $[a \circ b]$ دا حقیقی مُڈھلا ٹکڑا اے، پر نہ $[b \circ {}$ اے (دوجے نشان اُتے فرق اے) تے نہ $[a
\circ b]$ اے (دوہاں دی لمبائی اکّو اے)۔

\begin{prop}\ollabel{prop:initial}
  کسے ٹھیک ٹرم~$t$ دے ہر حقیقی مُڈھلے ٹکڑے وچ $[$ دی گنتی $]$ دی گنتی توں ودھ اے۔
\end{prop}

\begin{proof}
  ~$t$ اُتے استقرا نال:
  \begin{enumerate}
  \item $t$ اکلّا اکھر اے۔ تاں $t$ دا کوئی حقیقی مُڈھلا ٹکڑا نہیں۔
  \item $t = [s_1 \circ s_2]$ کسے ٹھیک ٹرماں $s_1$ تے~$s_2$ لئی۔ جے $r$، $t$ دا حقیقی مُڈھلا ٹکڑا اے، تاں کئی صورتاں ممکن نیں:
    \begin{enumerate}
    \item $r$ بس $[$ اے۔ تاں $r$ وچ اک $[$،~$]$ نالوں ودھ اے۔
    \item $r$، $[r_1$ اے، جتھے $r_1$،~$s_1$ دا حقیقی مُڈھلا ٹکڑا اے۔ کیوں جے $s_1$ ٹھیک ٹرم اے، استقرائی مفروضے مطابق $r_1$ وچ $[$، $]$ نالوں ودھ نیں، تے ایہو گل $[r_1$ لئی وی سچ اے۔
    \item $r$، $[s_1$ یا $[s_1 \circ {}$ اے۔ پچھلے نتیجے مطابق $[$ تے $]$ دیاں گنتیاں~$s_1$ وچ برابر نیں۔ سو $[$ دی گنتی $[s_1$ یا $[s_1 \circ {}$ وچ،~$]$ دی گنتی نالوں اک ودھ اے۔
    \item $r$، $[s_1 \circ r_2$ اے، جتھے $r_2$،~$s_2$ دا حقیقی مُڈھلا ٹکڑا اے۔ استقرائی مفروضے مطابق $r_2$ وچ $[$، $]$ نالوں ودھ نیں۔ پچھلے نتیجے مطابق $[$ تے $]$ دیاں گنتیاں~$s_1$ وچ برابر نیں۔ سو $[$ دی گنتی $[s_1 \circ r_2$ وچ،~$]$ دی گنتی توں ودھ اے۔
    \item $r$، $[s_1 \circ s_2$ اے۔ پچھلے نتیجے مطابق $[$ تے $]$ دیاں گنتیاں $s_1$ وچ برابر نیں، تے~$s_2$ لئی وی ایہو گل اے۔ سو اک $[$، $[s_1 \circ s_2$ وچ،~$]$ دی گنتی نالوں ودھ اے۔
    \end{enumerate}
  \end{enumerate}
\end{proof}
\endgroup
% /OLP-0618

% OLP-0619 content/methods/induction/relations.tex
\begingroup
\def\olfilename{relations}
\expandafter\def\csname ol@sectioncs\endcsname{section}

\olfileid{mth}{ind}{rel}

\olsection{تعلق تے فنکشن}

جدوں اسیں شےواں دا کوئی سیٹ (جیویں قدرتی عدد یا ٹھیک ٹرم) استقرائی طریقے نال تعریف کر لئیے، تاں ایہناں شےواں \emph{اُتے تعلقاں} دی تعریف وی استقرا نال دے سکدے آں۔ مثلاً ایہ خیال ویکھو۔ اک ٹھیک ٹرم~$t_1$، اک ٹھیک ٹرم~$t_2$ دا ذیلی ٹرم اے جے اوہ اوہدے اک حصے وجوں آؤندا اے۔ ایس لئی اک نشان ورتیے، $t_1 \sqsubseteq t_2$۔ بے شک ہر ٹھیک ٹرم اپنا ذیلی ٹرم اے، $t \sqsubseteq t$۔ ایس تعلق دی استقرائی تعریف ایویں دے سکدے آں:

\begin{defn}
  ایہ تعلق کہ ٹھیک ٹرم~$t_1$،~$t_2$ دا ذیلی ٹرم اے، $t_1
  \sqsubseteq t_2$،~$t_2$ اُتے استقرا نال ایویں تعریف ہُندا اے:
  \begin{enumerate}
  \item جے $t_2$ اکھر اے، تاں $t_1 \sqsubseteq t_2$ اُدوں تے صرف اُدوں جدوں $t_1 = t_2$۔
  \item جے $t_2$، $[s_1 \circ s_2]$ اے، تاں $t_1 \sqsubseteq t_2$ اُدوں تے صرف اُدوں جدوں $t_1 =
    t_2$، یا $t_1 \sqsubseteq s_1$، یا $t_1 \sqsubseteq s_2$۔
  \end{enumerate}
\end{defn}

مثلاً ایہ تعریف دسدی اے کہ $\mathrm{a}
\sqsubseteq [\mathrm{b} \circ \mathrm{a}]$۔ کیوں جے (۲) دسدا اے کہ
$\mathrm{a} \sqsubseteq [\mathrm{b} \circ \mathrm{a}]$ اُدوں تے صرف اُدوں جدوں $\mathrm{a}
= [\mathrm{b} \circ \mathrm{a}]$، یا $\mathrm{a} \sqsubseteq b$، یا
$\mathrm{a} \sqsubseteq \mathrm{a}$۔ پہلیاں دو گلاں غلط نیں۔ صاف اے کہ $\mathrm{a}$، $[\mathrm{b} \circ
  \mathrm{a}]$ دے اکّو نہیں، تے~(۱) مطابق $\mathrm{a} \sqsubseteq \mathrm{b}$ اُدوں تے صرف اُدوں جدوں $\mathrm{a} = \mathrm{b}$، جیہڑی گل وی غلط اے۔ پر~(۱) مطابق ای، $\mathrm{a} \sqsubseteq \mathrm{a}$ اُدوں تے صرف اُدوں جدوں $\mathrm{a} = \mathrm{a}$، جیہڑی گل سچ اے۔

ایہ دھیان رکھنا ضروری اے کہ تعریف دی کامیابی اک ایہو جیہی گل اُتے ٹکدی اے جیہڑی اسیں ہن تک ثابت نہیں کیتی۔ ہر ٹھیک ٹرم~$t$ یا بس اک اکھر اے، یا \emph{اکّو طریقے نال طے ہون والے} ٹھیک ٹرم $s_1$ تے $s_2$ نیں جیہناں لئی $t = [s_1 \circ s_2]$۔ «اکّو طریقے نال طے ہون والے» دا مطلب اے کہ جے $t = [s_1 \circ s_2]$، تاں ایہ \emph{نال ای} $= [r_1 \circ r_2]$ نہیں ہو سکدا، جیہدے وچ $s_1 \neq r_1$ یا $s_2 \neq r_2$ ہووے۔ جے ایہ ہو سکدا، تاں شق~(۲) دا اپنے آپ نال ٹکراؤ ہو سکدا سی۔ $t_2$ نوں $[s_1 \circ s_2]$ پڑھ کے شاید سانوں $t_1
\sqsubseteq t_2$ ملے، پر $t_2$ نوں $[r_1 \circ r_2]$ پڑھ کے شاید سانوں $t_1 \sqsubseteq t_2$ دی نفی ملے۔ ایہ ثابت کرن توں پہلاں کہ ایہ نہیں ہو سکدا، اک مثال ویکھیے جتھے ایہ ہو \emph{سکدا} اے۔

\begin{defn}
  \emph{بریکٹاں توں بناں ٹرم} استقرائی طریقے نال ایویں تعریف کرو:
  \begin{enumerate}
  \item ہر اکھر بریکٹاں توں بناں اک ٹرم اے۔
  \item جے $s_1$ تے $s_2$ بریکٹاں توں بناں ٹرم نیں، تاں $s_1 \circ s_2$ بریکٹاں توں بناں اک ٹرم اے۔
  \item ہور کوئی شے بریکٹاں توں بناں ٹرم نہیں۔
  \end{enumerate}
\end{defn}

بریکٹاں توں بناں ٹرماں دیاں مثالاں $\mathrm{a}$، $\mathrm{b} \circ
\mathrm{d}$، $\mathrm{b} \circ \mathrm{a} \circ \mathrm{b}$ نیں۔ ہُن جے بریکٹاں توں بناں ٹرماں لئی «ذیلی ٹرم» اُتے والے طریقے نال تعریف کریے، تاں دوجی شق ایویں ہووے گی:
\begin{center}
  جے $t_2 = s_1 \circ s_2$، تاں $t_1 \sqsubseteq t_2$ اُدوں تے صرف اُدوں جدوں $t_1 = t_2$، یا $t_1
  \sqsubseteq s_1$، یا $t_1 \sqsubseteq s_2$۔
\end{center}

ہُن $\mathrm{b} \circ \mathrm{a} \circ \mathrm{b}$، $s_1
\circ s_2$ دی صورت وچ اے، جتھے
\begin{align*}
  s_1 & = \mathrm{b} \text{ تے} & s_2 & = \mathrm{a} \circ \mathrm{b}.
\intertext{ایہ $r_1 \circ r_2$ دی صورت وچ وی اے، جتھے}
  r_1 & = \mathrm{b} \circ \mathrm{a} \text{ تے} & r_2 &= \mathrm{b}.
\end{align*}
ہُن کی $\mathrm{a} \circ \mathrm{b}$، $\mathrm{b} \circ
\mathrm{a} \circ \mathrm{b}$ دا ذیلی ٹرم اے؟ پہلی پڑھت لئی جواب ہاں اے تے دوجی لئی نہیں۔

ایس خاصیت نوں \emph{اکّو طرح پڑھے جانا} آکھدے نیں کہ ہور ٹھیک ٹرماں توں کسے ٹھیک ٹرم دی بناوٹ دا طریقہ اکّو اے۔ کیوں جے ایہو جیہیاں استقرائی طریقے نال تعریف ہوئیاں شےواں لئی تعلقاں دیاں استقرائی تعریفاں ضروری نیں، سانوں ثابت کرنا پئے گا کہ ایہ خاصیت پوری ہُندی اے۔

\begin{prop}
  منّ لو $t$ اک ٹھیک ٹرم اے۔ تاں یا $t$ بس اکلّا اکھر اے، یا اکّو طریقے نال طے ہون والے ٹھیک ٹرم $s_1$، $s_2$ نیں جیہناں لئی~$t =
  [s_1 \circ s_2]$۔
\end{prop}

\begin{proof}
  جے $t$ اکلّا اکھر اے، تاں شرط پوری اے۔ سو منّ لو $t$ اکلّا اکھر نہیں۔ استقرائی تعریف توں دس سکدے آں کہ فیر $t$ لازماً $[s_1 \circ s_2]$ دی صورت وچ اے، کسے ٹھیک ٹرماں $s_1$ تے~$s_2$ لئی۔ ہن بس ایہ وکھاؤنا اے کہ ایہ دوویں اکّو طریقے نال طے نیں، یعنی جے $t = [r_1 \circ r_2]$، تاں $s_1 = r_1$ تے $s_2 = r_2$۔

  سو منّ لو $t = [s_1 \circ s_2]$ تے $t = [r_1 \circ r_2]$ وی، کسے ٹھیک ٹرماں $s_1$، $s_2$، $r_1$، $r_2$ لئی۔ سانوں وکھاؤنا اے کہ $s_1 = r_1$ تے $s_2 = r_2$۔ پہلاں، $s_1$ تے $r_1$ لازماً اکّو نیں، کیوں جے نہیں تاں اک دوجے دا حقیقی مُڈھلا ٹکڑا اے۔ پر \olref[sti]{prop:initial} مطابق جے $s_1$ تے~$r_1$ دوویں ٹھیک ٹرم نیں تاں ایہ ناممکن اے۔ پر جے $s_1 = r_1$، تاں صاف اے کہ $s_2 = r_2$ وی۔
\end{proof}

اسیں فنکشن وی استقرائی طریقے نال تعریف کر سکدے آں۔ مثلاً فنکشن~$f$ دی تعریف ایویں دے سکدے آں جیہڑا کسے وی ٹھیک ٹرم نوں اوہدے اندر اک دوجے وچ پئے $[\dots]$ دی سبھ توں ودھ ڈونگھائی ول بھیجدا اے:

\begin{defn}
  \ollabel{defn:depth} ٹھیک ٹرم دی \emph{ڈونگھائی}، $f(t)$، استقرائی طریقے نال ایویں تعریف ہُندی اے:
  \[
  f(t) = \begin{cases}
    0 & \text{ جے $t$ اک اکھر اے}\\
    \max(f(s_1), f(s_2)) + 1 & \text{ جے $t = [s_1 \circ s_2]$۔}
  \end{cases}
  \]
\end{defn}

مثلاً
\begin{align*}
  f([\mathrm{a} \circ \mathrm{b}]) & =
    \max(f(\mathrm{a}),f(\mathrm{b})) + 1 = \\
  &= \max(0, 0) + 1 = 1, \text{ تے}\\
  f([[\mathrm{a} \circ \mathrm{b}] \circ \mathrm{c}]) & =
    \max(f([\mathrm{a} \circ \mathrm{b}]), f(\mathrm{c})) + 1 = \\
  & = \max(1,0) + 1 = 2.
\end{align*}

بے شک ایتھے اسیں منّدے آں کہ $s_1$ تے $s_2$ ٹھیک ٹرم نیں، تے ایہ گل ورتدے آں کہ ہر ٹھیک ٹرم یا اکھر اے یا $[s_1 \circ s_2]$ دی صورت وچ اے۔ فیر ایہ گل ضروری اے کہ اوہ بس اکّو طریقے نال ایس صورت وچ ہو سکدا اے۔ وجہ ویکھن لئی بریکٹاں توں بناں اوہ ٹرم فیر ویکھو جیہڑے اسیں پہلاں تعریف کیتے سن۔ اوہناں لئی ایس ورگی «تعریف» ایہ ہووے گی:
\[
  g(t) =
  \begin{cases}
    0 & \text{ جے $t$ اک اکھر اے}\\
   \max(g(s_1), g(s_2)) + 1 & \text{ جے $t = s_1 \circ s_2$۔}
  \end{cases}
\]
ہُن بریکٹاں توں بناں ٹرم $\mathrm{a} \circ \mathrm{b} \circ
\mathrm{c} \circ \mathrm{d}$ ویکھو۔ ایہنوں کئی طرحاں نال پڑھ سکدے آں، مثلاً $s_1 \circ s_2$ وجوں، جتھے
\begin{align*}
  s_1 & = \mathrm{a} \text{ تے} &
  s_2 & = \mathrm{b} \circ \mathrm{c} \circ \mathrm{d},
\intertext{یا $r_1 \circ r_2$ وجوں، جتھے}
  r_1 & = \mathrm{a} \circ b \text{ تے} &
  r_2 &= \mathrm{c} \circ \mathrm{d}.
\end{align*}
پہلی پڑھت مطابق $g$ دا حساب لائیے تاں ایہ ملے گا:
\begin{align*}
  g(s_1 \circ s_2) & =
  \max(g(\mathrm{a}), g(\mathrm{b} \circ \mathrm{c} \circ \mathrm{d})) + 1 =\\ & =
  \max(0,2) + 1 = 3
  \intertext{جدکہ دوجی پڑھت مطابق سانوں ایہ ملدا اے}
  g(r_1 \circ r_2) & =
  \max(g(\mathrm{a} \circ \mathrm{b}), g(\mathrm{c} \circ \mathrm{d})) + 1=\\ &
  = \max(1,1) + 1 = 2
\end{align*}
پر فنکشن نوں ہر واری اکّو قدر دینی پَندی اے۔ سو~$g$ دی ساڈی «تعریف» کوئی فنکشن تعریف نہیں کردی۔

\begin{prob}
  فنکشن $l$ دی استقرائی تعریف دیو، جتھے $l(t)$، ٹھیک ٹرم~$t$ وچ نشاناں دی گنتی اے۔
\end{prob}

\begin{prob}
  ٹھیک ٹرماں $t$ اُتے ساختی استقرا نال ثابت کرو کہ $f(t) < l(t)$ (جتھے $l(t)$،~$t$ وچ نشاناں دی گنتی اے تے $f(t)$، $t$ دی ڈونگھائی اے جیہدی تعریف \olref[mth][ind][rel]{defn:depth} وچ دتی اے)۔
\end{prob}
\endgroup
% /OLP-0619

% OLP-0718 content/sets-functions-relations/inductive-defs-proofs/introduction.tex
\begingroup
\def\olfilename{introduction}
\expandafter\def\csname ol@sectioncs\endcsname{section}

\olfileid{sfr}{ind}{int}

\olsection{جان پچھان}

\begin{explain}
استقرا ریاضی دی اک ثبوتی چال اے جیہڑی منطق وچ عام ورتی جاندی اے۔ ہو سکدا اے تسی ریاضی توں ایہدے نال واقف ہوو، جتھے قدرتی عدداں اُتے استقرا کردے او۔ اصل وچ ریاضی تے منطق دیاں شیواں دے کئی وکھ وکھ طرح دے سیٹ نیں جنہاں اُتے استقرا لگ سکدا اے۔

عام طور اُتے استقرا نال تسی اک کلی دعویٰ ثابت کر سکدے او، یعنی وکھا سکدے او کہ کسے خاص قسم دی ہر شے وچ کوئی خاصیت اے۔ خاص کر استقرا اکثر اوہدوں کم آؤندا اے جدوں «کلی شامل کرن» والا ثبوتی طریقہ بہت گنجل دار ہووے۔ استقرا صرف ایہو جیہیاں ریاضیاتی شیواں اُتے لگدا اے جیہڑیاں خاص طریقیاں نال بنیاں ہون۔ جے سیٹ دی ہر شے یا بنیادی اے یا بنیادی شیواں توں بنی اے، تاں اسیں اوس اُتے استقرا ورت سکدے آں۔ ہور ٹھیک طرح ایہ گل ایویں اے:
\end{explain}

\begin{defn}
سیٹ $S$ کسے فنکشن $f$ ہیٹھ \emph{بند} اے اوہدوں تے صرف اوہدوں جدوں ہر واری جے $x \in
S$، تاں $f(x) \in S$۔
\end{defn}

\begin{explain}
مثال لئی قدرتی عدد ($\mathbb{N}$) اگلا عدد دین والے فنکشن $f(x) = x+1$ ہیٹھ بند نیں۔ جے $x$ قدرتی عدد اے، تاں $x+1$ وی اے۔ پر قدرتی عدد پچھلا عدد دین والے فنکشن $f(x) = x-1$ ہیٹھ بند نہیں، کیونکہ $0$ قدرتی عدد اے، پر $-1$ نہیں۔
\end{explain}

\begin{defn}
آکھدے آں کہ خاصیت $P$، فنکشن $f(x)$ ہیٹھ \emph{قائم رہندی} اے جدوں کسے وی شے $a$ لئی جیہڑی $f$ دے دائرے وچ اے، $P(a)$ توں $P(f(a))$ نکلدا اے (جے $a$ وچ خاصیت $P$ اے، تاں $f(a)$ وچ وی اے)۔
\end{defn}

\begin{explain}
«جفت عدد ہونا» والی خاصیت فنکشن $f(x) = x+2$ ہیٹھ قائم رہندی اے، کیونکہ جے $x$ جفت اے، تاں $x+2$ وی جفت اے۔ جفت عدد ہون دی خاصیت اگلا عدد دین والے فنکشن ہیٹھ قائم نہیں رہندی، کیونکہ جے $x$ جفت اے، تاں $x+1$ طاق اے۔

استقرا قدرتی عدداں اُتے ایس لئی کم کردا اے کہ ہر قدرتی عدد تک 0 توں، اگلا عدد دین والا فنکشن محدود واریاں لا کے، پہنچیا جا سکدا اے۔ ایہ ہر اوس سیٹ دی خاص گل اے جیہدے اُتے اسیں استقرا کر سکدے آں۔
\end{explain}

\begin{defn}[$\Nat$ اُتے استقرا]
جے 0 وچ خاصیت $P$ اے تے جے $P$، اگلا عدد دین والے فنکشن ہیٹھ قائم رہندی اے، تاں ہر قدرتی عدد وچ خاصیت $P$ اے۔
\end{defn}

\begin{ex} پہلے $n$ قدرتی عدداں دا جوڑ $\frac{n(n+1)}{2}$ اے۔ \\

\textbf{مُڈھلی صورت۔} پہلے 0 قدرتی عدداں دا جوڑ $0=
\frac{0\cdot 1}{2}$ اے۔\\

\textbf{استقرائی قدم۔} منّ لو کہ خاصیت $k$ لئی ٹھیک اے۔ اسیں وکھاواں گے کہ ایہ $k+1$ لئی وی ٹھیک اے۔ دوجے لفظاں وچ، منّ لو کہ $P(k)$ ٹھیک اے، یعنی
\[ 0 + 1 + \ldots + k = \frac{k(k+1)}{2} \]
اسیں وکھاواں گے کہ $P(k+1)$ ٹھیک اے، یعنی
\[ 0 + 1 + \ldots + k + (k+1) = \frac{(k+1)(k+2)}{2} \]
ایہ مفروضہ ورت کے کہ $P(k)$ ٹھیک اے۔

\begin{align*}
0 + 1 + \ldots + k &= \frac{k(k+1)}{2} \tag{مفروضہ} \\
0 + 1 + \ldots + k + (k+1) &= \frac{k(k+1)}{2} + (k+1)
\tag{$k+1$ دوناں پاسیاں وچ پاؤ} \\
&= \frac{k(k+1) + 2(k+1)}{2} \\
&= \frac{2k + k + 2k +2}{2} \\
&=\frac{(k+1)(k+2)}{2}
\end{align*}
سو استقرائی قدم پورا ہو گیا، تے اسیں دعویٰ ثابت کر دِتّا اے۔
\end{ex}

\begin{explain}
!!{formula}s لئی بنیادی اجزا جوہری !!{formula}s نیں۔ ہور گنجل دار !!{formula}s گھٹ گنجل دار !!{formula}s نوں عمل کاراں نال جوڑ کے ملدے نیں۔ ہر عمل کار نوں !!{formula}s دے کسے فنکشن نال میل کھاندا سمجھ سکدے آں۔ اوہ فنکشن اک نواں !!{formula} دیندا اے جیہڑا ایہناں دوہاں نوں متعلقہ عمل کار نال جوڑدا اے۔ مثال لئی دو !!{formula}s $!A$ تے $!B$ نوں عطف وچ جوڑن والا فنکشن $ \mathcal E _\land (!A,
!B) = !A \land !B$ لکھیا جا سکدا اے۔ ایہناں نوں \emph{فارمولے بناؤن والے فنکشن} آکھ سکدے آں۔ !!{formula}s دا سیٹ ایہناں فنکشناں ہیٹھ بند اے۔ جدوں وی $!A$ تے $!B$، !!{formula}s نیں، تاں $\lnot !A$، $!A \land
!B$ وغیرہ وی نیں۔
\end{explain}

\begin{defn}[!!{formula}s اُتے استقرا]
جے ہر جوہری !!{formula} وچ خاصیت $P$ اے تے $P$، !!{formula} بناؤن والے فنکشناں ہیٹھ قائم رہندی اے، تاں ہر !!{formula} وچ خاصیت $P$ اے۔
\end{defn}

\begin{explain}
ایہ تعریف !!{formula}s اُتے استقرائی ثبوت دا اک «طریقہ» دَسّدی اے۔ جے سانوں وکھاؤن نوں آکھیا جاوے کہ ہر !!{formula} وچ خاصیت $P$ اے، تاں ہِٹھلے قدم چُکّو:\\

\textbf{مُڈھلی صورت۔} منّ لو $!A$ اک جوہری !!{formula} اے۔ [\ldots] سو $!A$ وچ خاصیت $P$ اے۔

\textbf{استقرائی قدم۔} منّ لو $!A$ تے $!B$، !!{formula}s نیں تے دوناں وچ خاصیت $P$ اے۔

صورت $\lnot$۔ [\ldots] سو $\lnot !A$ وچ خاصیت $P$ اے۔

صورت $\land$۔ [\ldots] سو $!A \land !B$ وچ خاصیت $P$ اے۔

صورت $\lor$۔ [\ldots] سو $!A \lor !B$ وچ خاصیت $P$ اے۔

صورت $\lif$۔ [\ldots] سو $!A \lif !B$ وچ خاصیت $P$ اے۔

صورت $\liff$۔ [\ldots] سو $!A \liff !B$ وچ خاصیت $P$ اے۔

صورت $\lforall[]$۔ [\ldots] سو $\lforall[x] !A$ وچ خاصیت $P$ اے۔

صورت $\lexists[]$۔ [\ldots] سو $\lexists[x] !A$ وچ خاصیت $P$ اے۔ \\

سو ہر !!{formula} وچ خاصیت $P$ اے۔
\end{explain}
\endgroup
% /OLP-0718
\OLEndChapterHook
\endgroup
% /OLP-0613

\OLEndPartHook
\endgroup
% /OLP-0601

% OLP-0620 content/history/history.tex
\begingroup
\def\olfilename{history}
\expandafter\def\csname ol@sectioncs\endcsname{section}

\olpart{his}{تاریخ}

% OLP-0621 content/history/biographies/biographies.tex
\begingroup
\def\olfilename{biographies}
\expandafter\def\csname ol@sectioncs\endcsname{section}

\olchapter{his}{bio}{سوانح عمریاں}

% OLP-0622 content/history/biographies/georg-cantor.tex
\begingroup
\def\olfilename{georg-cantor}
\expandafter\def\csname ol@sectioncs\endcsname{section}

\olfileid{his}{bio}{can}

\olsection{گیورگ کانتور}

\olphoto{cantor-georg}{گیورگ کانتور}

گیورگ کانتور (\textsc{گے} اورگ \textsc{کان} تور) دی اک مُڈھلی سوانح عمری وچ دعویٰ کیتا گیا سی کہ اوہ روس دے سینٹ پیٹرزبرگ ول جاندے اک بحری جہاز وچ جمیا تے اوتھے ای لبھیا سی، تے اوہدے ماں پیو دا پتا نہیں سی۔ پر ایہ گل سچ نہیں۔ اوہ 1845 وچ سینٹ پیٹرزبرگ وچ ضرور جمیا سی۔

کانتور نے 1867 وچ برلن یونیورسٹی توں ریاضی وچ ڈاکٹریٹ لئی۔ اوہ سیٹ نظریے وچ اپنے کم لئی جانیا جاندا اے، تے سیٹ نظریے نوں تحقیق دے اک وکھرے شعبے وجوں قائم کرن دا سہرا اوہدے سر اے۔ اوہ پہلا بندہ سی جیہنے ثابت کیتا کہ وکھ وکھ وڈائی دے لامتناہی سیٹ ہُندے نیں۔ اوہدے نظریاں، خاص کر کے لامتناہیت دے نظریے، نے اوس ویلے دے ریاضی داناں وچ بڑی بحث چھیڑی، تے اوہدا کم اختلاف دا باعث رہیا۔

کانتور دے مذہبی یقین تے اوہدا ریاضیاتی کم اینے جڑے سن کہ اوہناں نوں وکھ نہیں کیتا جا سکدا۔ اوہنے ایہ دعویٰ وی کیتا کہ لامتناہی درجیاں والے عدداں دا نظریہ رب نے اوہنوں سِدھا دسیا سی۔ پچھلی عمر وچ کانتور ذہنی بیماری دا شکار رہیا۔ 1894 توں اوہ ہسپتال داخل ہون لگ پیا، تے عمر دے آخری سالاں وچ ایہ ہور واری واری ہویا۔ اوہدے کم اُتے سخت تنقید، جیہدے وچ ریاضی دان لیوپولڈ کرونیکر نال ان بن وی شامل سی، اوہدی افسردگی تے ریاضی وچ دلچسپی گھٹ جان دا سبب بنی۔ افسردگی دے دوراں وچ کانتور فلسفے تے ادب ول مُڑ جاندا سی۔ اوہنے ایہ نظریہ وی چھپوایا کہ شیکسپیئر دے ڈرامے فرانسس بیکن نے لکھے سن۔

کانتور 6 جنوری 1918 نوں ہالے دے اک علاج گھر وچ وفات پا گیا۔

\begin{reading}
کانتور دیاں پوریاں سوانح عمریاں لئی \citet{Dauben1990} تے \citet{Grattan-Guinness1971} ویکھو۔ کانتور دے رواج توں ہٹ کے خیالات دا بیان بی بی سی ریڈیو 4 دے پروگرام \emph{A Brief History of
  Mathematics} وچ وی ملدا اے \citep{Sautoy2014}۔ جے تسی کانتور دے نظریے ریپ دی صورت وچ سُننا چاہندے او، تاں \citet{Rose2012} ویکھو۔
\end{reading}
\endgroup
% /OLP-0622

% OLP-0623 content/history/biographies/alonzo-church.tex
\begingroup
\def\olfilename{alonzo-church}
\expandafter\def\csname ol@sectioncs\endcsname{section}

\olfileid{his}{bio}{chu}

\olsection{الونزو چرچ}

\olphoto{church-alonzo}{الونزو چرچ}

الونزو چرچ 14 جون 1903 نوں واشنگٹن ڈی سی وچ جمیا۔ نِکّی عمر وچ ہوا والی بندوق دے اک حادثے نال چرچ دی اک اکھ دی نظر چلی گئی۔ اوہنے 1920 وچ کنیکٹیکٹ وچ یونیورسٹی لئی تیاری والا سکول پورا کیتا تے اوسے سال پرنسٹن وچ یونیورسٹی دی پڑھائی شروع کیتی۔ اوہنے 1927 وچ اپنی ڈاکٹریٹ دی پڑھائی پوری کیتی۔ پردیس وچ دو کُ سال گزار کے چرچ پرنسٹن مُڑ آیا۔ چرچ بڑی شائستگی تے احتیاط لئی مشہور سی۔ تختے اُتے اوہدی لکھائی بے عیب ہُندی سی۔ ضروری کاغذاں نوں سنبھالن لئی اوہ احتیاط نال اوہناں اُتے ڈوکو سیمنٹ (اک شفاف گوند) دی تہہ لا دیندا سی۔ علمی کم توں ہٹ کے اوہ سائنس فکشن دے رسالے پڑھنا پسند کردا سی، تے جے لکھت وچ کوئی غلطی دِس پَندی تاں مدیران نوں لکھن توں جھجکدا نہیں سی۔

چرچ دیاں علمی کامیابیاں وڈیاں سن۔ اپنے شاگرداں سٹیفن کلینی تے بارکلے راسر نال رَل کے اوہنے مؤثر حساب دا نظریہ، لیمبڈا کیلکولس، بنایا۔ ایہ کم ایلن ٹیورنگ دی ٹیورنگ مشین دی بناوٹ توں آزاد سی۔ قابلِ حساب ہون دیاں دوویں تعریفاں برابر نیں، تے اوہناں توں اوہ گل نکلدی اے جیہنوں ہُن \emph{چرچ–ٹیورنگ دا دعویٰ} آکھدے نیں، یعنی قدرتی عدداں دا کوئی فنکشن مؤثر طور اُتے حساب جوگا اے اُدوں تے صرف اُدوں جدوں اوہ ٹیورنگ مشین (یا لیمبڈا کیلکولس) نال حساب جوگا اے۔ اوہنے اوہ گل وی ثابت کیتی جیہنوں ہُن \emph{چرچ دا قضیہ} آکھدے نیں۔ پہلے درجے دے فارمولیاں دے جائز ہون دا فیصلہ کرن والا مسئلہ حل نہیں ہو سکدا۔

چرچ وڈی عمر تک اپنا کم کردا رہیا۔ 1967 وچ اوہ پرنسٹن چھڈ کے یو سی ایل اے چلا گیا، جتھے اوہ 1990 وچ ریٹائر ہون تک پروفیسر رہیا۔ چرچ 1 اگست 1995 نوں 92 سال دی عمر وچ وفات پا گیا۔

\begin{reading}
چرچ دی مختصر سوانح عمری لئی \citet{EndertonND} ویکھو۔ لیمبڈا کیلکولس تے فیصلہ مسئلے، Entscheidungsproblem، (چرچ دے دعوے) بارے چرچ دیاں اصلی لکھتاں \citet{Church1936,Church1936a} نیں۔ \citet{Aspray1984} وچ 1930 دی دہائی دے پرنسٹن دے ریاضی داناں بارے چرچ نال اک انٹرویو درج اے۔ چرچ نے 1936 توں 1979 تک \emph{Journal of
Symbolic Logic} وچ کتاباں اُتے تبصریاں دی اک لڑی لکھی۔ ایہ سارے جان مکفارلین دی ویب سائٹ اُتے محفوظ نیں \citep{MacFarlane2015}۔
\end{reading}
\endgroup
% /OLP-0623

% OLP-0624 content/history/biographies/gerhard-gentzen.tex
\begingroup
\def\olfilename{gerhard-gentzen}
\expandafter\def\csname ol@sectioncs\endcsname{section}

\olfileid{his}{bio}{gen}

\olsection{گیرہارٹ گینٹسن}

\olphoto{gentzen-gerhard}{گیرہارٹ گینٹسن}

گیرہارٹ گینٹسن دی پہچان خاص کر کے ساختی نظریۂ ثبوت دے بانی وجوں اے۔ خاص طور اُتے اوہنے قدرتی اخذ تے سیکوئنٹ کیلکولس دے !!{derivation} نظام بنائے۔ اوہ 24 نومبر 1909 نوں جرمنی دے گرائفسوالڈ وچ جمیا۔ گیرہارٹ نے تیاری والے سکول جان توں پہلاں تِن سال گھر وچ پڑھائی کیتی۔ سکول وچ اوہ پڑھائی پکھوں اپنے بہتے جماعتیاں توں پچھے سی۔ ایس دے باوجود اوہ بڑا ہونہار طالب علم سی تے ریاضی وچ اوہدی خاص قابلیت دِسدی سی۔ اوہ کئی طرح دے کم پسند کردا سی، مثلاً اپنی ماں لئی نظماں تے سکول دے تھیٹر لئی ڈرامے وی لکھدا سی۔

گینٹسن نے یونیورسٹی دی پڑھائی گرائفسوالڈ یونیورسٹی توں شروع کیتی، پر فیر گوٹنگن، میونخ تے برلن وی گیا۔ اوہنے 1933 وچ گوٹنگن یونیورسٹی توں ہرمان وائل دی نگرانی تھلّے ڈاکٹریٹ لئی۔ (اوہدے کم دے بہتے حصے دی نگرانی پاؤل بیرنائز نے کیتی، پر نازیاں نے بیرنائز نوں یونیورسٹی توں کڈھ دتا سی۔) 1934 وچ گینٹسن نے ڈیوڈ ہلبرٹ دے معاون وجوں کم شروع کیتا۔ اوسے سال اوہنے اپنے مقالیاں \emph{Untersuchungen \"{u}ber das logische Schlie\ss en I--II
  [منطقی اخذ دی پڑتال، حصے ۱ تے ۲]} وچ سیکوئنٹ کیلکولس تے قدرتی اخذ دے !!{derivation} نظام بنائے۔ اوہنے 1936 وچ پیانو دے مسلّمیاں دا بے ٹکراؤ ہونا ثابت کیتا۔

گینٹسن دا نازیاں نال تعلق گنجل دار اے۔ جیہڑے ویلے اوہدے استاد بیرنائز نوں جرمنی چھڈن لئی مجبور کیتا گیا، اوسے ویلے گینٹسن نازی نیم فوجی تنظیم ایس اے دی یونیورسٹی شاخ وچ شامل ہویا۔ کئی ہور جرمن لوکاں وانگ اوہ نازی پارٹی دا رکن سی۔ جنگ دوران اوہ فضائی انٹیلیجنس یونٹ وچ مواصلات دا افسر سی۔ پر 1942 وچ شدید ذہنی پریشانی کارن اوہنوں ڈیوٹی توں فارغ کر دتا گیا۔ ایہ صاف نہیں کہ گینٹسن نازی پارٹی دا وفادار سی یا اپنی علمی ترقی پکی کرن لئی پارٹی وچ شامل ہویا۔

1943 وچ پراگ دی جرمن یونیورسٹی دے ریاضی ادارے وچ گینٹسن نوں اک علمی عہدے دی پیشکش ہوئی، جیہڑی اوہنے منّ لئی۔ پر 1945 وچ پراگ دے شہری جرمن قبضے دے خلاف اُٹھ کھلوتے۔ سوویت فوج شہر وچ آئی تے یونیورسٹی دے سارے پروفیسراں نوں گرفتار کر لیا۔ نازی تنظیماں دی رکنیت کارن گینٹسن نوں جبری مزدوری دے کیمپ لے گئے۔ اوہ 4 اگست 1945 نوں 35 سال دی عمر وچ اپنی کوٹھڑی وچ خوراک دی کمی نال مر گیا۔

\begin{reading}
گینٹسن دی پوری سوانح عمری لئی \citet{Menzler-Trott2007} ویکھو۔ نازی راج تھلّے ریاضی داناں بارے اک دلچسپ لکھت \citet{Segal2014} دی اے، جیہدے وچ گینٹسن دی زندگی دا مختصر ذکر وی اے۔ منطقی اخذ اُتے گینٹسن دے مقالے اصل جرمن وچ ملدے نیں \citep{Gentzen1935a,Gentzen1935b}۔ \citet{Gentzen1969} وچ گینٹسن دے مقالیاں دے انگریزی ترجمے اکّو جلد وچ اکٹھے کیتے گئے نیں تے اوہدے وچ اوہدی زندگی دا مختصر خاکہ وی اے۔
\end{reading}
\endgroup
% /OLP-0624

% OLP-0625 content/history/biographies/kurt-goedel.tex
\begingroup
\def\olfilename{kurt-goedel}
\expandafter\def\csname ol@sectioncs\endcsname{section}

\olfileid{his}{bio}{god}

\olsection{کرٹ گوڈل}

\olphoto{goedel-kurt}{کرٹ گوڈل}

کرٹ گوڈل (\textsc{گَر} ڈل) 28 اپریل 1906 نوں آسٹرو ہنگیرین سلطنت دے برُن شہر (ہُن چیک جمہوریہ دا برنو) وچ جمیا۔ نِکّے کرٹیلے دی جانن دی تانگھ تے ذہانت کرکے گھر دے لوک اکثر اوہنوں «Der kleine Herr Warum» (نِکّا مسٹر کیوں) آکھدے سن۔ پرائمری سکول توں ای اوہ پڑھائی وچ اگے سی۔ صرف ریاضی وچ اوہنوں سبھ توں اُچا گریڈ نہیں سی ملدا۔ ماڑی صحت کارن گوڈل اکثر سکول توں غیر حاضر رہندا سی تے اوہنوں جسمانی تعلیم توں چھوٹ سی۔ بچپن وچ اوہنوں گٹھیا دا بخار تشخیص ہویا۔ ڈاکٹری معائنے دے اُلٹ، اوہ ساری عمر ایہ منّدا رہیا کہ ایس نال اوہدے دل اُتے پکا اثر پیا سی۔

گوڈل نے 1924 وچ ویانا یونیورسٹی وچ پڑھائی شروع کیتی تے~1929 وچ ڈاکٹریٹ دی پڑھائی پوری کیتی۔ اوہدا پہلا ارادہ طبیعیات پڑھن دا سی، پر چھیتی اوہدی دلچسپی ریاضی ول، خاص کر کے منطق ول، مُڑ گئی۔ ایہدے وچ فلسفی روڈولف کارناپ دا وی کجھ اثر سی۔ ہانس ہان دی نگرانی تھلّے لکھے اوہدے مقالے نے شناخت والی پہلے درجے دی محمولی منطق دی تکمیل دا قضیہ ثابت کیتا \citep{Godel1929}۔ بس اک سال پچھوں اوہنے اپنے سبھ توں مشہور نتیجے، نامکملیت دے پہلے تے دوجے قضیے، حاصل کیتے (ایہ \citealt{Godel1931} وچ چھپے)۔ ویانا وچ رہندیاں گوڈل ویانا سرکل نال ڈونگھا جڑیا ہویا سی۔ ایہ سائنسی سوچ والے فلسفیاں دی جماعت سی جیہدے وچ کارناپ وی سی، تے گوڈل دے نتیجیاں دا کارناپ دے کم اُتے خاص اثر پیا۔

1938 وچ گوڈل نے اڈیل نمبورْسکی نال ویاہ کیتا۔ اوہدے ماں پیو خوش نہیں سن۔ اوہ گوڈل توں چھ سال وڈی سی، اوہنوں طلاق ہو چکی سی، تے اوہ رات دے کلب وچ رقاصہ سی۔ پر سماجی دباؤ دا گوڈل اُتے اثر نہیں پیا۔ اوہ گوڈل دی وفات تک خوشی نال ویاہے رہے۔

1938 وچ نازی جرمنی دے آسٹریا نوں اپنے نال ملاؤن پچھوں گوڈل تے اڈیل امریکہ ہجرت کر گئے۔ اوتھے گوڈل نے پرنسٹن، نیو جرسی دے انسٹیٹیوٹ فار ایڈوانسڈ سٹڈی وچ عہدہ سنبھالیا۔ گھٹ میل جول تے وکھرے مزاج دے باوجود پرنسٹن وچ اوہنے ہورناں نال رَل کے بہت نتیجہ خیز کم کیتا۔ اوہنے سیٹ نظریے، فلسفے تے طبیعیات اُتے مضمون چھپوائے۔ خاص گل ایہ کہ انسٹیٹیوٹ دے اپنے ساتھی البرٹ آئن سٹائن نال اوہدی بڑی گُوڑھی دوستی ہو گئی۔

آخری سالاں وچ گوڈل دی ذہنی صحت وگڑ گئی۔ 1977 وچ اوہدی گھر والی دے ہسپتال داخل ہون نال اوہ اوہدے لئی کھانا نہیں سی پکا سکدی۔ گوڈل ساری عمر ذہنی صحت دیاں اوکڑاں نال لڑدا رہیا سی تے آخر اوہنوں شدید وہم ہو گیا۔ زہر دتے جان دے ڈر نال اوہنے کھان توں انکار کر دتا۔ اوہ 14 جنوری 1978 نوں پرنسٹن وچ بھکھ نال مر گیا۔

\begin{reading}
گوڈل دی پوری سوانح عمری لئی \citet{Dawson1997} ویکھو۔ اوہدی زندگی بارے ہور لکھتاں تے منطق تے فلسفے وچ اوہدے کم بارے مضموناں لئی \citet{Wang1990}، \citet{Baaz2011}، \citet{Takeuti2003} تے \citet{Sigmund2007} ویکھو۔

گوڈل دا ڈاکٹریٹ مقالہ اصل جرمن وچ ملدا اے \citep{Godel1929}۔ نامکملیت دے قضیاں دی اصل لکھت \citep{Godel1931} اے۔ گوڈل دیاں ساریاں چھپیاں تے ان چھپیاں لکھتاں تے چُنے ہوئے خط انگریزی وچ اوہدیاں \emph{Collected Papers} وچ ملدے نیں \citet{Godel1986,Godel1990}۔

گوڈل دے نامکملیت دے قضیاں دے تفصیلی بیان لئی \citet{Smith2013} ویکھو۔ اوہدے قضیاں دی غیر رسمی، فلسفیانہ گل بات لئی مارک لِنسن مائر دا پوڈکاسٹ ویکھو \citep{Linsenmayer2014}۔
\end{reading}
\endgroup
% /OLP-0625

% OLP-0626 content/history/biographies/emmy-noether.tex
\begingroup
\def\olfilename{emmy-noether}
\expandafter\def\csname ol@sectioncs\endcsname{section}

\olfileid{his}{bio}{noe}

\olsection{ایمی نوئتر}

\olphoto{noether-emmy}{ایمی نوئتر}

ایمی نوئتر (\textsc{نَر} ٹر) 23 مارچ 1882 نوں جرمنی دے ایرلانگن شہر وچ پڑھے لکھے اُچلے درمیانے طبقے دے ٹبّر وچ جمی۔ اوہنوں «اج دی الجبرا دی ماں» آکھیا جاندا اے۔ عورتاں دی تعلیم اگے وڈیاں رکاوٹاں دے باوجود نوئتر نے ریاضی تے طبیعیات دوہاں وچ نئی راہ کھولھن والا کم کیتا۔ اوس ویلے جرمنی وچ کُڑیاں نوں فن پڑھائے جاندے سن تے اوہناں نوں کالج دی تیاری والے سکولاں وچ جان دی کھُل نہیں سی۔ پر گوٹنگن تے ایرلانگن دیاں یونیورسٹیاں وچ داخلے توں بناں جماعتاں سُنن پچھوں (ایرلانگن وچ اوہدا پیو ریاضی دا پروفیسر سی)، نوئتر آخر 1904 وچ ایرلانگن وچ طالب علم وجوں داخل ہو سکی، جدوں اوتھوں دے قاعدے بدل کے عورتاں نوں داخلے دی کھُل دتی گئی۔ اوہنے 1907 وچ ریاضی وچ ڈاکٹریٹ لئی۔

اپنی قابلیت دے باوجود نوئتر نوں اپنے پیشے وچ بہت مخالفت دا سامنا کرنا پیا۔ 1908 توں 1915 تک اوہنے ایرلانگن وچ تنخواہ توں بناں پڑھایا۔ ایس سمیں اوہ ڈیوڈ ہلبرٹ دی نظر چڑھی، جیہڑا اوس ویلے دنیا دے وڈے ریاضی داناں وچوں اک سی۔ ہلبرٹ نے اوہنوں گوٹنگن سَدیا۔ پر عورتاں نوں پروفیسری لین دی کھُل نہیں سی۔ نوئتر بس ہلبرٹ دے ناں تھلّے، فیر تنخواہ توں بناں، لیکچر دے سکدی سی۔ ایس سمیں اوہنے اوہ گل ثابت کیتی جیہنوں ہُن نوئتر دا قضیہ آکھدے نیں، تے جیہڑا اج وی نظری طبیعیات وچ ورتیا جاندا اے۔ آخر 1919 وچ نوئتر نوں پڑھاؤن دا حق مل گیا۔ دسیا جاندا اے کہ یونیورسٹی دے ساتھی جدوں ہن وی مخالفت کردے رہے تاں ہلبرٹ نے جواب دِتا، «حضرات، اساتذہ دی سینیٹ کوئی حمام نہیں۔»

1920 دی دہائی دے پچھلے سالاں وچ اوہنے تجریدی الجبرا اُتے دھیان دِتا، تے اوہدے کم نے ایس شعبے نوں بدل کے رکھ دِتا۔ اپنے ثبوتاں وچ اوہ اکثر چڑھدی زنجیر دی شرط ورتدی سی، جیہڑی دسدی اے کہ کجھ خاص سیٹاں دی سختی نال ودھدی کوئی لامتناہی زنجیر نہیں۔ مثلاً ہُن نوئتری حلقے آکھیاں جاندیاں کجھ الجبرائی ساختاں دی خاصیت اے کہ اوہناں وچ آئیڈیلاں دی کوئی لامتناہی لڑی $I_1
\subsetneq I_2 \subsetneq \dots$ نہیں۔ ایہ شرط کسے وی جزوی ترتیب تک ودھائی جا سکدی اے (الجبرا وچ ایہ ذیلی سیٹ دے تعلق نال ترتیب دتے آئیڈیلاں دی خاص صورت بارے اے)۔ اسیں ایس دی دوہری، لہندی زنجیر دی شرط وی ویکھ سکدے آں، جتھے جزوی ترتیب وچ سختی نال \emph{گھٹدی} ہر لڑی آخر مُک جاندی اے۔ جے کوئی جزوی ترتیب لہندی زنجیر دی شرط اُتے پوری اُترے، تاں اوس ترتیب اُتے استقرا اویویں ورت سکدے آں جیویں $<$ ترتیب اُتے~$\Nat$ وچ ورتدے آں۔ ایہو جیہیاں ترتیباں نوں \emph{بنیاد والیاں} یا \emph{نوئتری} آکھدے نیں، تے اوہناں دے ثبوت دے اصول نوں \emph{نوئتری استقرا} آکھدے نیں۔

نوئتر یہودی سی۔ جدوں 1933 وچ نازی راج وچ آئے تاں اوہنوں عہدے توں کڈھ دِتا گیا۔ چنگے بھاگیں نوئتر امریکہ ہجرت کر سکی، جتھے برِن مار، پنسلوانیا وچ اوہنوں عارضی عہدہ مل گیا۔ اوتھے رہندیاں اوہنے پرنسٹن وچ وی لیکچر دتے، بھاویں اوہنوں لگدا سی کہ یونیورسٹی عورتاں نوں جی آیاں نہیں آکھدی \citep[81]{Dick1981}۔ 1935 وچ بچہ دانی دی رسولی کڈھن لئی نوئتر دا آپریشن ہویا۔ اوہ آپریشن توں ہوئی انفیکشن نال مر گئی تے برِن مار وچ دفنائی گئی۔

\begin{reading}
نوئتر دی سوانح عمری لئی \citet{Dick1981} ویکھو۔ پیریمیٹر انسٹیٹیوٹ فار تھیوریٹیکل فزکس دے نوئتر دی زندگی تے اثر بارے لیکچر انٹرنیٹ اُتے ملدے نیں \citep{Perimeter2015}۔ جے تسی پڑھ پڑھ کے تھک گئے او، تاں \emph{Stuff You Missed in History Class} دا نوئتر دی زندگی تے اثر بارے اک پوڈکاسٹ اے \citep{Frey2015}۔ نوئتر دیاں اکٹھیاں لکھتاں اصل جرمن وچ ملدیاں نیں \citep{Noether1983}۔
\end{reading}
\endgroup
% /OLP-0626

% OLP-0627 content/history/biographies/rozsa-peter.tex
\begingroup
\def\olfilename{rozsa-peter}
\expandafter\def\csname ol@sectioncs\endcsname{section}

\olfileid{his}{bio}{pet} 

\olsection{روژا پیتر}

\olphoto{peter-rozsa}{روژا پیتر}

روژا پیتر 17 فروری 1905 نوں بڈاپسٹ، ہنگری وچ روژا پولٹسر دے ناں نال جمی۔ اوہ سبھ توں ودھ بازگشتی فنکشناں اُتے اپنے کم لئی جانی جاندی اے۔ ایہ کم بازگشت دا نظریہ اک شعبے وجوں بناؤن لئی ضروری سی۔

پیتر سخت سیاسی حالات وچ وڈی ہوئی۔ اوہدی چڑھدی عمر وچ پہلی عالمی جنگ چل رہی سی۔ پر اوہ بڈاپسٹ دے خوشحال ماریا تیریزیا گرلز سکول وچ پڑھ سکی تے 1922 وچ اوتھوں پڑھائی پوری کیتی۔ فیر اوہنے بڈاپسٹ دی پازمانی پیتر یونیورسٹی (جیہدا ناں بعد وچ لوراند ایوتووش یونیورسٹی رکھیا گیا) وچ پڑھیا۔ اپنے پیو دے زور دین اُتے اوہنے کیمیا پڑھنی شروع کیتی، پر فیر ریاضی ول آ گئی تے 1927 وچ ڈگری لئی۔ بھاویں اوہدے کول ہائی سکول وچ ریاضی پڑھاؤن دی اہلیت سی، اوس ویلے مالی حال بڑا ماڑا سی، کیوں جے وڈی کساد بازاری دا اثر سارے جہان دی معیشت اُتے پیا سی۔ ایس سمیں پیتر نے ریاضی دی نجی استاد وجوں تھوڑے موٹے کم کیتے۔ آخر اوہ ریاضی دی اعلیٰ پڑھائی کرن لئی یونیورسٹی مُڑ آئی۔ اوہدا پہلا ارادہ عدد نظریے وچ کم کرن دا سی، پر جدوں پتا لگیا کہ اوہدے نتیجے پہلاں ای ثابت ہو چکے سن، تاں اوہ ریاضی چھڈن لگی سی۔ اوہنوں گوڈل دے عدم تکمیل دے قضیاں اُتے کم کرن دی ہمت دِتی گئی، تے اوہنے انجانے وچ اوہدے کئی نتیجے وکھرے طریقیاں نال ثابت کر دتے۔ ایس نال اوہدا بھروسا مُڑ آیا، تے ڈیوڈ ہلبرٹ دے ریاضی دیاں بنیاداں والے پروگرام توں تحریک لے کے اوہنے بازگشت دے نظریے اُتے اپنے پہلے مقالے لکھے۔ اوہنے 1935 وچ پی ایچ ڈی لئی، تے 1937 وچ \emph{Journal of
  Symbolic Logic} دی مدیر بنی۔

پیتر دے مُڈھلے مقالے بڑے پیمانے اُتے بازگشتی فنکشناں دے نظریے دی بنیاد رکھن والا کم منّے جاندے نیں۔ \cite{Peter1935a} وچ اوہنے وکھ وکھ قسماں دی بازگشت دا آپسی تعلق ویکھیا۔ \cite{Peter1935b} وچ اوہنے وکھایا کہ بازگشت نال تعریف کیتا اک خاص فنکشن ابتدائی بازگشتی نہیں۔ ایس نال ولہلم آکَرمان دا پہلا نتیجہ سوکھا ہو گیا۔ پیتر دا سوکھا کیتا فنکشن ہُن اکثر آکَرمان فنکشن آکھیا جاندا اے، تے کدی ہور ٹھیک ناں، آکَرمان–پیتر فنکشن، نال وی۔ اوہنے بازگشتی فنکشناں دے نظریے اُتے پہلی کتاب لکھی \citep{Peter1951}۔

اپنے کم دی اہمیت تے اثر دے باوجود پیتر نوں 1945 تک پورے ویلے دی پڑھاؤن والی نوکری نہ ملی۔ دوجی عالمی جنگ وچ ہنگری اُتے نازی قبضے ویلے یہودی مخالف قاعدیاں کرکے پیتر نوں پڑھاؤن دی کھُل نہیں سی۔ 1944 وچ سرکار نے بڈاپسٹ وچ یہودیاں لئی اک گھِری ہوئی بستی بنائی۔ ایہ بستی باقی شہر توں کٹّی ہوئی سی تے اوتھے ہتھیاربند پہرا سی۔ پیتر نوں 1945 وچ بستی آزاد ہون تک اوتھے رہن لئی مجبور کیتا گیا۔ فیر اوہنے بڈاپسٹ ٹیچرز ٹریننگ کالج وچ پڑھایا تے 1955 توں ایوتووش لوراند یونیورسٹی وچ۔ اوہ ریاضی دی اکادمی والی ڈاکٹریٹ لین والی پہلی ہنگیرین عورت ریاضی دان تے ہنگیرین اکادمی آف سائنسز لئی چُنی جان والی پہلی عورت سی۔

پیتر ریاضی دی پرجوش استاد وجوں جانی جاندی سی۔ اوہ بس لیکچر دین دی بجائے اپنے طالب علماں نال ریاضیاتی مسئلیاں دی نوعیت تے سوہپ لبھنا پسند کردی سی۔ ایس لئی اوہدے طالب علم پیار نال اوہنوں «ماسی روزا» آکھدے سن۔ پیتر 1977 وچ~71 سال دی عمر وچ وفات پا گئی۔

\begin{reading}
اوہدی زندگی بارے ہور پڑھن لئی \citep{Oconnor2014} تے \citep{Andrasfai1986} ویکھو۔ \citet{Tamassy1994} نے پیتر نال اک مختصر انٹرویو کیتا۔ ریاضی بارے دلچسپ پڑھائی لئی پیتر دی کتاب \emph{Playing With Infinity} ویکھو \citep{Peter2010}۔
\end{reading}
\endgroup
% /OLP-0627

% OLP-0628 content/history/biographies/julia-robinson.tex
\begingroup
\def\olfilename{julia-robinson}
\expandafter\def\csname ol@sectioncs\endcsname{section}

\olfileid{his}{bio}{rob}

\olsection{جولیا رابنسن}

\olphoto{robinson-julia}{جولیا رابنسن}

جولیا بومین رابنسن امریکی ریاضی دان سی۔ اوہ خاص کر کے فیصلہ کرن والے مسئلیاں اُتے اپنے کم لئی جانی جاندی اے، تے سبھ توں ودھ ہلبرٹ دے دسویں مسئلے دے حل وچ اپنے حصے لئی۔ رابنسن 8 دسمبر 1919 نوں سینٹ لوئس، میزوری وچ جمی۔ رابنسن نوں چیتا سی کہ اوہ بچپن توں ای عدداں وچ دلچسپی لیندی سی \citep[4]{Reid1986}۔ نو سال دی عمر وچ اوہنوں لال بخار ہویا تے فیر گٹھیا دے بخار دے کئی دورے پئے۔ ایس کرکے اوہنوں بہت ویلا منجے اُتے گزارنا پیا تے اوہ پڑھائی وچ پچھے رہ گئی۔ نجی استاداں دی مدد نال اوہ پڑھائی وچ فیر نال رَل گئی، پر بیماری دے جسمانی اثر اوہدی ساری زندگی اُتے رہے۔

بچپن دیاں اوکڑاں دے باوجود رابنسن نے ریاضی تے سائنس دے کئی انعاماں نال ہائی سکول پورا کیتا۔ یونیورسٹی دی پڑھائی اوہنے سین ڈیاگو سٹیٹ کالج توں شروع کیتی تے آخری سال برکلی دی یونیورسٹی آف کیلیفورنیا چلی گئی۔ اوتھے ریاضی دان رافائل رابنسن دا اوہدے اُتے اثر پیا۔ اوہ چنگے دوست بن گئے تے 1941 وچ ویاہ کر لیا۔ استاداں وچوں اک دی گھر والی ہون کرکے رابنسن نوں برکلی دے ریاضی شعبے وچ پڑھاؤن دی کھُل نہ ملی۔ بھاویں اوہ بناں داخلے ریاضی دیاں جماعتاں سُندی رہی، اوہ یونیورسٹی چھڈ کے ٹبّر شروع کرنا چاہندی سی۔ پر ویاہ توں تھوڑی دیر پچھوں رابنسن نوں نمونیا ہو گیا۔ اوہنوں دسیا گیا کہ بچپن دے گٹھیا بخار کارن اوہدے دل اُتے زخم دے نشان والا ٹشو کافی جم چکیا سی۔ ایس ٹشو دی سنگینی کارن ڈاکٹر نے اندازہ لایا کہ اوہ چالی سال توں اگے نہیں جی سکدی، تے اوہنوں بچے نہ کرن دی صلاح دتی گئی \citep[13]{Reid1986}۔

رابنسن لمّا سماں افسردہ رہی، پر آخر ریاضی دی پڑھائی جاری رکھن دا فیصلہ کیتا۔ اوہ برکلی مُڑ آئی تے 1948 وچ الفریڈ تارسکی دی نگرانی تھلّے پی ایچ ڈی پوری کیتی۔ تارسکی وکھا چکیا سی کہ حقیقی عدداں دے پہلے درجے دے نظریے دا فیصلہ ہو سکدا اے، تے گوڈل دے کم توں نکلدا سی کہ قدرتی عدداں دے پہلے درجے دے نظریے دا فیصلہ نہیں ہو سکدا۔ ایہ اک وڈا کھُلّا سوال سی کہ ناطق عدداں دے پہلے درجے دے نظریے دا فیصلہ ہو سکدا اے یا نہیں۔ اپنے مقالے \citeyearpar{Robinson1949} وچ رابنسن نے ثابت کیتا کہ ایہدا فیصلہ نہیں ہو سکدا۔

فیصلہ کرن والے مسئلیاں وچ دلچسپی کرکے رابنسن نے اگے ہلبرٹ دا دسواں مسئلہ حل کرن دی کوشش کیتی۔ ایہ مسئلہ 23 مشہور ریاضیاتی مسئلیاں دی اوہ فہرست وچوں سی جیہڑی ڈیوڈ ہلبرٹ نے 1900 وچ پیش کیتی سی۔ دسواں مسئلہ پُچھدا اے کہ کی کوئی ایہو جیہا الگورتھم اے جیہڑا محدود سمیں وچ جواب دے سکے کہ صحیح عددی ضریباں والی کثیر رکنی مساوات، جیویں $3x^2 - 2y +3 = 0$، دا صحیح عدداں وچ حل اے یا نہیں۔ ایہو جیہے سوالاں نوں \emph{ڈیوفانتی مسئلے} آکھدے نیں۔ کجھ مُڈھلیاں کامیابیاں پچھوں رابنسن مارٹن ڈیوس تے ہلیری پٹنم نال رَل گئی، جیہڑے ایس مسئلے اُتے کم کردے پئے سن۔ اوہناں نے وکھایا کہ گھات والے ڈیوفانتی مسئلیاں (جتھے انجان عدد گھات دی تھاں وی آ سکدے نیں) دا فیصلہ نہیں ہو سکدا۔ اوہناں نے وی وکھایا کہ اک خاص قیاس (جیہنوں بعد وچ «J.R.» آکھیا گیا) توں لازم آؤندا اے کہ ہلبرٹ دے دسویں مسئلے دا فیصلہ نہیں ہو سکدا \citep{DavisPutnamRobinson1961}۔ رابنسن 1960 دی ساری دہائی ایس مسئلے اُتے کم کردی رہی۔ 1970 وچ نوجوان روسی ریاضی دان یوری ماتیاسےوِچ نے آخر J.R. دا مفروضہ ثابت کر دتا۔ رَلواں نتیجہ ہُن ماتیاسےوِچ–رابنسن–ڈیوس–پٹنم قضیہ، یا مختصر MRDP قضیہ، کہلاندا اے۔ ماتیاسےوِچ تے رابنسن دوست بن گئے تے کئی مقالیاں اُتے رَل کے کم کیتا۔ ماتیاسےوِچ نوں اک خط وچ رابنسن نے اک واری لکھیا، «مینوں سچ مُچ بڑی خوشی اے کہ رَل کے کم کردیاں (ہزاراں میل دور رہندیاں) اسیں صاف ای اوہتوں ودھ اگے ودھ رہے آں جِنّا ساڈے وچوں کوئی اک اکلّا ودھ سکدا سی» \citep[45]{Matijasevich1992}۔

رابنسن امریکن میتھمیٹیکل سوسائٹی دی پہلی عورت صدر سی تے نیشنل اکادمی آف سائنس لئی چُنی جان والی پہلی عورت سی۔ لیوکیمیا دی تشخیص پچھوں اوہ 30 جولائی 1985 نوں 65 سال دی عمر وچ وفات پا گئی۔

\begin{reading}
رابنسن دے ریاضیاتی مقالے اوہدیاں \textit{Collected
  Works} وچ ملدے نیں \citep{Robinson1996}۔ اوہدے وچ نیشنل اکادمی آف سائنسز لئی اوہدی سوانح دا مُڑ چھاپیا روپ وی اے \citep{Feferman1994}۔ رابنسن دی وڈی بھین کونسٹینس ریڈ نے انٹرویوواں دی بنیاد اُتے «Autobiography of Julia» چھپوائی \citep{Reid1986} تے پوری یادداشت وی لکھی \citep{Reid1996}۔ رابنسن تے ہلبرٹ دے دسویں مسئلے بارے اک مختصر دستاویزی فلم جارج سِسری نے ہدایت کیتی \citep{Csicsery2016}۔ رابنسن نال یوری ماتیاسےوِچ دے سانجھے کم تے اوہدے کم اُتے رابنسن دے اثر بارے مختصر یادداشت لئی \citep{Matijasevich1992} ویکھو۔
\end{reading}
\endgroup
% /OLP-0628

% OLP-0629 content/history/biographies/bertrand-russell.tex
\begingroup
\def\olfilename{bertrand-russell}
\expandafter\def\csname ol@sectioncs\endcsname{section}

\olfileid{his}{bio}{rus} 

\olsection{برٹرینڈ رسل}

\olphoto{russell-bertrand}{برٹرینڈ رسل}

برٹرینڈ رسل نوں اج دے تحلیلی فلسفے دے بانیواں وچوں اک منّیا جاندا اے۔ رسل 18 مئی 1872 نوں جمیا۔ اوہ فلسفے تے منطق وچ اپنے کم لئی جانیا جاندا سی تے اوہنے کئی وکھرے موضوعاں اُتے عام لوکاں لئی مشہور کتاباں وی لکھیاں۔ اوہ ساری عمر جوش نال سیاسی سرگرمیاں وچ وی حصہ لیندا رہیا۔

رسل ویلز دے مونمتھ شائر وچ ٹریلیک ناں دی تھاں اُتے جمیا۔ اوہدے ماں پیو برطانوی اشرافیہ دے رکن سن۔ اوہ آزاد خیال سن تے اوہناں نے اوس ویلے بوسٹن دے انقلابی خیال رکھن والے لوکاں نال وی دوستی پائی۔ پر جدوں رسل نِکّا سی تاں اوہدے ماں پیو فوت ہو گئے، تے رسل نوں اپنے دادا دادی کول رہن لئی بھیج دِتا گیا۔ اوتھے اوہدی پرورش مذہبی ڈھنگ نال ہوئی (جیہڑی گل اوہدے ماں پیو ہر حال وچ ٹالنا چاہندے سن)۔ اوہدی دادی اخلاق دیاں ساریاں گلاں وچ بڑی سخت سی۔ چڑھدی عمر وچ اوہ زیادہ تر گھر وچ نجی استاداں کولوں پڑھیا۔

تحلیلی فلسفے، تے خاص کر کے منطق، اُتے رسل دا اثر بڑا وڈا اے۔ اوہنے کیمبرج دے ٹرینیٹی کالج وچ ریاضی تے فلسفہ پڑھیا، جتھے ریاضی دان تے فلسفی الفریڈ نارتھ وائٹ ہیڈ دا اوہدے اُتے اثر پیا۔ 1910 وچ رسل تے وائٹ ہیڈ نے \emph{Principia Mathematica} دی پہلی جلد چھپوائی۔ اوہدے وچ اوہناں نے ایہ خیال اگے رکھیا کہ ریاضی نوں منطق تک گھٹایا جا سکدا اے۔ رسل نے اگے سیکڑیاں کتاباں، مضمون تے سیاسی رسالے چھپوائے۔ 1950 وچ اوہنوں ادب دا نوبیل انعام ملیا۔

رسل سیاست تے سماجی سرگرمیاں وچ ڈونگھا جڑیا ہویا سی۔ پہلی عالمی جنگ دوران امن دے حق وچ سرگرمیاں تے احتجاج کارن اوہنوں گرفتار کر کے چھ مہینیاں لئی جیل بھیج دتا گیا۔ جیل وچ اوہ لکھ پڑھ سکدا سی، تے اوہدا آکھنا سی کہ ایہ تجربہ اوہنوں «خاصا چنگا» لگیا۔ اوہ ساری عمر جنگ دا مخالف رہیا۔ 1961 وچ ایٹمی ہتھیار ختم کرن دے حق والی ریلی وچ جان کرکے اوہنوں فیر جیل ہوئی۔ اوہ 1948 وچ اک ہوائی جہاز دے حادثے وچ وی بچ گیا۔ اوس حادثے وچ بس اوہ لوک بچے جیہڑے سگریٹ پین والیاں دے حصے وچ بیٹھے سن۔ ایس لئی رسل آکھدا سی کہ اوہدی جان سگریٹ پین نال بچی۔ رسل نے چار ویاہ کیتے، پر ویاہ توں باہر اوہدے تعلقاں دی وی شہرت سی۔ اوہ 2 فروری 1970 نوں 97 سال دی عمر وچ ویلز دے پینرِن دیودرائتھ وچ فوت ہویا۔

\begin{reading}
رسل نے اپنی زندگی 1872 توں 1967 تک دسن والی سوانح عمری تِن حصیاں وچ لکھی \citep{Russell1967,Russell1968,Russell1969}۔ مک ماسٹر یونیورسٹی دے برٹرینڈ رسل ریسرچ سینٹر وچ برٹرینڈ رسل دا دستاویزی ذخیرہ اے۔ رسل دیاں اکٹھیاں لکھتاں دیاں جلداں (لبھن جوگ اشارییاں سمیت) تے دستاویزی منصوبیاں دی جانکاری لئی اوہناں دی ویب سائٹ \citet{Duncan2015} ویکھو۔ رسل دا مقالہ \emph{On
  Denoting} \citep{Russell1905} ویہویں صدی دے تحلیلی فلسفے دی اک کلاسیکی لکھت اے۔

سٹینفورڈ انسائیکلوپیڈیا آف فلاسفی دے رسل والے مضمون \citep{Irvine2015} وچ خواہش تے سیاسی نظریے اُتے رسل دی گل بات دیاں آڈیو ریکارڈنگاں نیں۔ رسل نال کئی ویڈیو انٹرویو انٹرنیٹ اُتے ملدے نیں۔ مثلاً اوہنوں سگریٹ پین تے ہوائی جہاز دے حادثے بارے گل کردیاں ویکھن لئی \citet{RussellND} ویکھو۔ رسل دیاں کجھ لکھتاں، جیہناں وچ اوہدی \emph{Introduction to Mathematical Philosophy} وی اے، \citet{LibriVoxND} اُتے مفت سُنن والیاں کتاباں وجوں ملدیاں نیں۔
\end{reading}
\endgroup
% /OLP-0629

% OLP-0630 content/history/biographies/alfred-tarski.tex
\begingroup
\def\olfilename{alfred-tarski}
\expandafter\def\csname ol@sectioncs\endcsname{section}

\olfileid{his}{bio}{tar}

\olsection{الفریڈ تارسکی}

\olphoto{tarski-alfred}{الفریڈ تارسکی}

الفریڈ تارسکی 14 جنوری 1901 نوں وارسا، پولینڈ (جیہڑا اوس ویلے روسی سلطنت دا حصہ سی) وچ جمیا۔ تارسکی نوں «نپولین ورگا» آکھیا جاندا سی۔ اوہ شور مچاؤن والا، بہت بولن والا تے تِکھے مزاج دا سی۔ اوہدی توانائی اکثر اوہدے لیکچراں وچ دِسدی سی۔ اک واری لیکچر دوران سگریٹ سُٹدیاں اوہنے کچرے دی ٹوکری نوں اگ لا دتی، تے اوہنوں فیر اوس عمارت وچ لیکچر دین توں روک دتا گیا۔

تارسکی نوں نِکّی عمر توں ای جانن دی تانگھ سی۔ بھاویں پچھلی عمر وچ اوہ طالب علماں نوں دسدا سی کہ اوہنے منطق ایس لئی پڑھی کیوں جے بس ایس مضمون وچ اوہنوں~B گریڈ ملیا سی، اوہدے ہائی سکول دے ریکارڈ توں پتا لگدا اے کہ اوہنوں ہر مضمون، منطق سمیت، وچ A گریڈ ملیا سی۔ اوہ 1918 توں 1924 تک وارسا یونیورسٹی وچ پڑھیا۔ پہلاں تارسکی حیاتیات پڑھنا چاہندا سی، پر فیر اوہنوں ریاضی، فلسفے تے منطق وچ دلچسپی ہو گئی، کیوں جے ایہ یونیورسٹی منطق تے فلسفے دے وارسا مکتب دا مرکز سی۔ تارسکی نے 1924 وچ ستانیسواف لیشنیفسکی (Stanis\l{}aw Le\'{s}niewski) دی نگرانی تھلّے ڈاکٹریٹ لئی۔

1939 وچ امریکہ ہجرت کرن توں پہلاں تارسکی نے وارسا دے سیکنڈری سکول وچ استاد ہُندیاں اپنا کجھ سبھ توں اہم کم پورا کیتا۔ منطقی نتیجے تے منطقی سچائی اُتے اوہدا کم ایس سمیں لکھیا گیا۔ 1939 وچ تارسکی لیکچر دین لئی امریکہ دے دورے اُتے سی۔ اوہدے دورے دوران جرمنی نے پولینڈ اُتے حملہ کر دتا، تے اپنے یہودی پس منظر کرکے تارسکی واپس نہ جا سکیا۔ اوہدی گھر والی تے بچے جنگ دے آخر تک پولینڈ وچ رہے، پر فیر اوہ وی امریکہ ہجرت کر سکے۔ تارسکی نے ہارورڈ، کالج آف دی سٹی آف نیو یارک، پرنسٹن دے انسٹیٹیوٹ فار ایڈوانسڈ سٹڈی، تے آخر برکلی دی یونیورسٹی آف کیلیفورنیا وچ پڑھایا۔ اوتھے اوہنے منطق تے سائنس دے طریقہ کار دا کئی شعبیاں نوں جوڑن والا پروگرام شروع کیتا۔ تارسکی 26 اکتوبر 1983 نوں 82 سال دی عمر وچ فوت ہویا۔

\begin{reading}
تارسکی دی زندگی بارے ہور لئی سوانح عمری \emph{Alfred Tarski: Life
  and Logic} ویکھو \citep{Feferman2004}۔ منطقی نتیجے تے سچائی اُتے تارسکی دیاں بنیادی لکھتاں انگریزی وچ \citep{Tarski1983} اندر ملدیاں نیں۔ تارسکی دیاں ساریاں اصل لکھتاں چار جلداں وچ اکٹھیاں کیتیاں گئیاں نیں \citep{Tarski1981}۔
\end{reading}
\endgroup
% /OLP-0630

% OLP-0631 content/history/biographies/alan-turing.tex
\begingroup
\def\olfilename{alan-turing}
\expandafter\def\csname ol@sectioncs\endcsname{section}

\olfileid{his}{bio}{tur} 

\olsection{ایلن ٹیورنگ}

ایلن ٹیورنگ 23 جون 1912 نوں مَیڈا ویل، لندن وچ جمیا۔ اوہنوں نظری کمپیوٹر سائنس دا بانی منّیا جاندا اے۔ طبیعی سائنس تے ریاضی وچ ٹیورنگ دی دلچسپی نِکّی عمر توں شروع ہوئی۔ پر اوہ جیہڑے سکولاں وچ پڑھیا، اوتھے اوہدی دلچسپی دے مضمون گھٹ پڑھائے جاندے سن۔ اوتھے ادب تے کلاسیکی لکھتاں اُتے زور سی۔ ایس لئی سکول وچ اوہدی کارکردگی ماڑی رہی تے کئی استاداں نے اوہنوں جھاڑ پائی۔

\olphoto{turing-alan}{ایلن ٹیورنگ}

ٹیورنگ نے کیمبرج دے کنگز کالج وچ پہلی ڈگری لئی ریاضی پڑھی۔ 1936 وچ ٹیورنگ نے اوہ شے بنائی جیہنوں ہُن ٹیورنگ مشین آکھدے نیں۔ اوہدی کوشش سی کہ حساب جوگ فنکشن دی عین تعریف دے سکے تے ثابت کرے کہ فیصلہ مسئلے دا فیصلہ نہیں ہو سکدا۔ الونزو چرچ ایس نتیجے تک پہلاں پہنچ گیا۔ اوہنے اپنے لیمبڈا کیلکولس نال ایہ نتیجہ ثابت کیتا سی۔ ٹیورنگ دا مقالہ فیر وی چرچ دے نتیجے دا حوالہ دے کے چھپیا۔ چرچ نے ٹیورنگ نوں پرنسٹن سَدیا۔ اوہنے 1936 توں 1938 تک اوتھے وقت گزاریا تے چرچ دی نگرانی تھلّے ڈاکٹریٹ لئی۔

منطق وچ دلچسپی دے باوجود ٹیورنگ دی طبیعی سائنس وچ پہلی دلچسپی قائم رہی۔ دوجی عالمی جنگ وچ بلیچلی پارک دے برطانوی خفیہ رمزاں دی پڑتال والے محکمے وچ کم کردیاں اوہدی عملی مہارت کم آئی۔ جرمن بحری فوج دے پیغاماں لئی ورتے جاندے اینگما رمز نوں توڑن وچ ٹیورنگ دا مرکزی کردار سی۔ شماریات تے خفیہ رمزاں وچ ٹیورنگ دی مہارت تے الیکٹرانک مشیناں دی آمد نال ٹیم «بومب» ناں دی رمز کھولھن والی مشین بنا کے رمز توڑ سکی۔ اوہدے خیالات نے دنیا دا پہلا پروگرام کیتا جا سکّن والا الیکٹرانک کمپیوٹر، کولاسس، بناؤن وچ وی مدد دتی۔ اوہ وی بلیچلی پارک وچ جرمن لورینز رمز توڑن لئی ورتیا گیا۔

ٹیورنگ ہم جنس پسند سی۔ فیر وی 1942 وچ اوہنے بلیچلی پارک دی اپنی ساتھی جون کلارک نوں ویاہ دی پیشکش کیتی۔ پر بعد وچ اوہنے منگنی توڑ دتی تے کلارک نوں دسیا کہ اوہ ہم جنس پسند اے۔ اوہدی زندگی وچ کئی محبوب سن، بھاویں اوس ویلے برطانیہ وچ ہم جنس جنسی عمل جرم سن۔ 1952 وچ ٹیورنگ دے اوس ویلے دے محبوب دے اک دوست نے ٹیورنگ دے گھر چوری کیتی۔ پولیس رپورٹ لکھاؤندیاں ٹیورنگ نے ہم جنس تعلق دا اقرار کیتا، کیوں جے اوہنوں لگدا سی کہ سرکار ہم جنس جنسی عملاں نوں قانونی کرن والی اے۔ ایہ گل سچ نہیں سی، تے اوہدے اُتے سنگین بے حیائی دا الزام لگایا گیا۔ جیل جان دی بجائے ٹیورنگ نے ہارمون والا علاج چُنیا جیہڑا جنسی خواہش گھٹاؤندا سی۔ ٹیورنگ 8 جون 1954 نوں مردہ لبھیا۔ موت سائینائیڈ دی ودھ مقدار نال ہوئی، غالباً خودکشی سی۔ 2013 وچ ملکہ الزبتھ~دوجی نے اوہنوں شاہی معافی دتی۔

\begin{reading}
ایلن ٹیورنگ دی تفصیلی سوانح عمری لئی \citet{Hodges2014} ویکھو۔ ٹیورنگ دی زندگی تے کم توں \emph{Breaking the Code} ناں دا ڈراما بنیا۔ 1996 وچ ایہ ٹی وی لئی بنیا، جیہدے وچ ڈیرک جیکوبی نے ٹیورنگ دا کردار کیتا۔ اکیڈمی انعام لئی نامزد فلم \emph{The Imitation Game} وچ بینیڈکٹ کمبربیچ تے کیرا نائٹلی نیں۔ ایہ وی ٹیورنگ دی زندگی تے بلیچلی پارک وچ اوہدے وقت اُتے ڈھیلے ڈھالے طور اُتے مبنی اے \citep{Imitation2014}۔

\citet{Radiolab2012} وچ ٹیورنگ دی زندگی تے کم بارے کئی پوڈکاسٹ نیں۔ بی بی سی ہورائزن دی دستاویزی فلم \emph{The Strange Life and Death of
  Dr.~Turing} انٹرنیٹ اُتے ویکھی جا سکدی اے \citep{Sykes1992}۔ \citep{Theelen2012} وچ لیگو دی کم کردی ٹیورنگ مشین دی مختصر ویڈیو اے۔ ایہ 2012 وچ ٹیورنگ دے جنم دے سو سال پورے ہون اُتے اوہدی یاد وچ بنائی گئی۔

ٹیورنگ مشیناں تے فیصلہ مسئلے اُتے ٹیورنگ دا اصل مقالہ \citet{Turing1937} اے۔
\end{reading}
\endgroup
% /OLP-0631

% OLP-0632 content/history/biographies/ernst-zermelo.tex
\begingroup
\def\olfilename{ernst-zermelo}
\expandafter\def\csname ol@sectioncs\endcsname{section}

\olfileid{his}{bio}{zer}

\olsection{ایرنسٹ زرمَیلو}

\olphoto{zermelo-ernst}{ایرنسٹ زرمَیلو}

ایرنسٹ زرمَیلو 27 جولائی 1871 نوں جرمنی دے برلن شہر وچ جمیا۔ اوہدیاں پنج بھیناں سن، پر ٹبّر دی ماڑی صحت کرکے بس تِن وڈی عمر تک پہنچیاں۔ اوہدے ماں پیو وی اوہدی نِکّی عمر وچ فوت ہو گئے، تے جدوں اوہ ستاراں سالاں دا سی تاں اوہ تے اوہدے بھین بھرا یتیم ہو گئے۔ زرمَیلو نوں فنّاں وچ، خاص کر کے شاعری وچ، ڈونگھی دلچسپی سی۔ اوہ تِکھّی سمجھ، حاضر جوابی تے تنقید کرن لئی مشہور سی۔ اوہدیاں سبھ توں مشہور ریاضیاتی کامیابیاں وچ چُون دا مسلّمہ پیش کرنا (1904 وچ) تے سیٹ نظریے نوں مسلّمیاں نال بیان کرنا (1908 وچ) شامل نیں۔

یونیورسٹی وچ زرمَیلو دیاں دلچسپیاں کئی پاسے سن۔ اوہنے طبیعیات، ریاضی تے فلسفے دے مضمون پڑھے۔ ہرمان شوارٹس دی نگرانی تھلّے زرمَیلو نے اپنا مقالہ \emph{Investigations in
  the Calculus of Variations} 1894 وچ برلن یونیورسٹی توں پورا کیتا۔ 1897 وچ اوہنے گوٹنگن یونیورسٹی وچ ہور پڑھائی دا فیصلہ کیتا، جتھے ڈیوڈ ہلبرٹ دے ریاضی دیاں بنیاداں اُتے کم دا اوہدے اُتے بڑا اثر پیا۔ 1899 وچ اوہ پروفیسری لئی اہل ہو گیا، پر گیاراں سال پچھوں تک اوہنوں پروفیسری نہ ملی، شاید اوہدے اوپرے رویے تے «گھبراہٹ بھری جلدبازی» کرکے۔

آخر 1910 وچ زرمَیلو نوں زیورخ یونیورسٹی وچ تنخواہ والی پروفیسری ملی، پر تپِ دق کرکے 1916 وچ اوہنوں ریٹائر ہونا پیا۔ صحت ٹھیک ہون پچھوں 1921 وچ فرائی برگ یونیورسٹی وچ اوہنوں اعزازی پروفیسری دتی گئی۔ ایس سمیں اوہنے ریاضی دیاں بنیاداں اُتے کم کیتا۔ اوہ تھورالف سکولم تے کرٹ گوڈل دیاں لکھتاں توں کھِجھ گیا تے اپنے مقالیاں وچ اوہناں دے طریقیاں اُتے کھُلّھی تنقید کیتی۔ 1935 وچ اوہنوں فرائی برگ دے عہدے توں ہٹا دتا گیا، کیوں جے اوہ مقبول نہیں سی تے جرمنی وچ ہٹلر دے اقتدار وچ آؤن دا مخالف سی۔

زرمَیلو دی زندگی دے آخری سال تنہائی وچ لنگھے۔ 1935 وچ عہدے توں ہٹائے جان پچھوں اوہنے ریاضی چھڈ دتی۔ اوہ پینڈو علاقے وچ چلا گیا تے سادے ڈھنگ نال رہیا۔ اوہنے 1944 وچ ویاہ کیتا، تے نظر گھٹدی جان کرکے اوہ پوری طرح اپنی گھر والی اُتے ٹکن لگ پیا۔ 1951 تک زرمَیلو دی نظر پوری طرح چلی گئی۔ اوہ 21 مئی 1953 نوں جرمنی دے گُنٹرشٹال وچ وفات پا گیا۔

\begin{reading}
زرمَیلو دی پوری سوانح عمری لئی \citet{Ebbinghaus2015} ویکھو۔ زرمَیلو دے 1904 تے 1908 والے بنیادی مقالے اصل جرمن وچ پڑھے جا سکدے نیں \citep{Zermelo1904,Zermelo1908}۔ زرمَیلو دیاں اکٹھیاں لکھتاں، طبیعیات اُتے اوہدی لکھت سمیت، انگریزی ترجمے وچ \citep{Ebbinghaus2010,Ebbinghaus2013} اندر ملدیاں نیں۔
\end{reading}
\endgroup
% /OLP-0632

\OLEndChapterHook
\endgroup
% /OLP-0621

% OLP-0633 content/history/set-theory/set-theory.tex
\begingroup
\def\olfilename{set-theory}
\expandafter\def\csname ol@sectioncs\endcsname{section}

\olchapter{his}{set}{سیٹ نظریے دی تاریخ تے متھاں}

\begin{editorial}
ایس باب وچ ٹِم بٹن دی کتاب «Open Set Theory» دا تاریخی مُڈھلا بیان شامل اے۔
\end{editorial}

% OLP-0634 content/history/set-theory/infinitesimals.tex
\begingroup
\def\olfilename{infinitesimals}
\expandafter\def\csname ol@sectioncs\endcsname{section}

\olfileid{his}{set}{infinitesimals}
\olsection{لامتناہی طور اُتے چھوٹیاں قدراں تے تفریق}

نیوٹن تے لائبنز نے ستارویں صدی دے آخر وچ حسابان وکھو وکھ لبھیا۔ حسابان دی اک خاص اہم ورتوں \emph{تفریق} سی۔ موٹے طور اُتے تفریق دا مقصد ایہ دَسنا اے کہ کسے نقطے اُتے فنکشن دی «بدلن دی شرح»، یا ڈھلوان، کیہ اے۔

ایہ خیال اک صاف مثال نال سمجھ آ سکدا اے۔ تھلّے کالے رنگ وچ وکھایا فنکشن
$f(x) = \nicefrac{x^2}{4} + \nicefrac{1}{2}$ لو:
\begin{center}
	\begin{tikzpicture}[scale=1]
	\draw[->, gray] (-1,0) -- (4.25,0) node[right] {$x$};
	\draw[->, gray] (0,-0.25) -- (0,5.25) node[above] {$f(x)$};
	
	\foreach \x/\xtext in {1/1, 2/2, 3/3, 4/4}
	\draw[shift={(\x,0)}, gray] (0pt,2pt) -- (0pt,-2pt) node[below] {$\xtext$};
	
	\foreach \y/\ytext in {1/1, 2/2, 3/3, 4/4, 5/5}
	\draw[shift={(0,\y)}, gray] (2pt,0pt) -- (-2pt,0pt) node[left] {$\ytext$};
	
	\draw[oldiagcolorC] (0.5, 9/16) -- (3.5, 57/16) -- (3.5, 9/16)--cycle;
	\draw[oldiagcolorD] (.5, 9/16) -- (2.5, 33/16) -- (2.5, 9/16) -- cycle;
	\draw[oldiagcolorE] (.5, 9/16) -- (1.5, 17/16) -- (1.5, 9/16) -- cycle;
	\draw[thick] (-1,.75) parabola bend (0,0.5) (4,4.5);
	\end{tikzpicture}
\end{center}
فرض کرو اسیں $c =
\nicefrac{1}{2}$ اُتے فنکشن دی ڈھلوان لبھنی اے۔ اسیں پہلاں اک تکون بناؤنے آں جیہدے وتر دی ڈھلوان اوس نقطے اُتے فنکشن دی ڈھلوان دے نیڑے ہووے۔ اُتلی لال تکون لے لو۔ جے ساڈی تکون دے پاوے دی لمبائی $\beta$ اے، تاں اوہدی اُچائی
$f(\nicefrac{1}{2}+\beta) - f(\nicefrac{1}{2})$ اے، سو وتر دی ڈھلوان اے:
\[
\frac{f(\nicefrac{1}{2}+\beta) - f(\nicefrac{1}{2})}{\beta}.
\]
سو !!{colorC} تکون دے پاوے دی لمبائی~$3$ اے تے اوہدی ڈھلوان عین~$1$ اے۔ اوس توں چھوٹی !!{colorD} تکون دا پاوا~$2$ اے۔ اوہدا وتر چنگا اندازہ دیندا اے، اوہدی ڈھلوان $\nicefrac{3}{4}$ اے۔ ہور وی چھوٹی !!{colorE} تکون دا پاوا~$1$ اے۔ ایہدے نال اندازہ ہور وی نیڑے اے، تے ڈھلوان $\nicefrac{1}{2}$ اے۔

جِنّی تکون چھوٹی، اُنا ای اندازہ اصل دے نیڑے۔ سو اسیں ایویں آکھ سکدے آں کہ جیہڑی تکون دے پاوے دی لمبائی \emph{لامتناہی طور اُتے چھوٹی} ہووے، اوہدے وتر دی ڈھلوان $c =
\nicefrac{1}{2}$ اُتے عین فنکشن دی ڈھلوان اے۔ ایس طرح فنکشن $f$ دے (پہلے) مشتق لئی، نقطے $c$ اُتے، ایہ فارمولا ملدا اے:
\[
{f'}(c) = \frac{f(c+\beta) - f(c)}{\beta} \text{ جتھے $\beta$ لامتناہی طور اُتے چھوٹا اے۔}
\]
نیوٹن تے لائبنز نے موٹے طور اُتے ایہو گل آکھی۔

پر جدوں اوہناں ایہ گل آکھی، تاں اوہناں توں پچھنا پیندا اے کہ \emph{لامتناہی طور اُتے چھوٹی قدر} کیہ اے؟ ایتھے وڈی دُبدھا پیدا ہوندی اے۔ جے $\beta = 0$ ہووے، تاں $f'$ دی تعریف ٹھیک نہیں بندی، کیوں جے $0$ نال تقسیم کرنا پیندا اے۔ پر جے $\beta >
0$ ہووے، تاں ڈھلوان دا بس اک \emph{اندازہ} ملدا اے، عین ڈھلوان نہیں۔

ایہ فِکر بعدلے زمانے دی سوچ نوں پچھلے ویلے اُتے لاگو کر کے نہیں اُٹھدی۔ ویکھو برکلے نیوٹن دے پیروکاراں دی تنقید کِویں کردا اے:
\begin{quote}
میں منّدا آں کہ علامتاں کسے شے نوں وی دس سکدیاں نیں تے کجھ نہ ہون نوں وی۔ سو اصلی لکھت $c + \beta$ وچ $\beta$ اک وادھے یا کجھ نہ ہون دا نشان ہو سکدا سی۔ پر ایہناں وچوں جیہڑا وی مطلب تسی لوو، دلیل اوسے مطلب اُتے قائم رکھو۔ دُہرے مطلب نال اگے ودھنا صاف مغالطہ ہووے گا۔ (\citealt[\S{}XIII]{Berkeley1734}، متغیر پچھلی لکھت نال میل کرن لئی بدلے گئے نیں۔)
\end{quote}
برکلے دے اعتراض توں لامتناہی طور اُتے چھوٹیاں قدراں والے حسابان دا بچاؤ کردیاں کوئی آکھ سکدا اے کہ «لامتناہی طور اُتے چھوٹیاں قدراں» بس مجازی گل اے۔ کوئی ایہ وی آکھ سکدا اے کہ جد تک تکون \emph{سچ مُچ چھوٹی} لئے، سانوں مماس دا \emph{کم چلاؤن جوگا} اندازہ مل جاندا اے۔ برکلے دا ایہدا وی جواب سی۔ اوہ آکھدا سی کہ انجینئرنگ وچ ایہ کم چلاؤن جوگا ہو سکدا اے، پر ایہ ریاضی دی \emph{حیثیت} نوں کھورا لاؤندا اے، کیوں جے:
\begin{quote}
سانوں دَسیا جاندا اے کہ \emph{in rebus mathematicis errores qu\`{a}m minimi
non sunt contemnendi}. [ریاضی وچ نِکّی توں نِکّی غلطی وی نظرانداز نہیں کرنی چاہیدی۔] \citep[\S{}IX]{Berkeley1734}
\end{quote}
ترچھے حرفاں والا ٹوٹا نیوٹن دی اپنی کتاب \emph{Quadrature of Curves} (1704) توں لگ بھگ لفظ بہ لفظ لیا گیا اے۔

برکلے دے فلسفیانہ اعتراض ڈونگھائی نال گل دی جڑ تک جاندے نیں۔ فیر وی حسابان دی ورتوں بے حد کامیاب رہی تے ریاضی دان غلطی وچ پئے بغیر اوہنوں ورتدے رہے۔
\endgroup
% /OLP-0634
% OLP-0635 content/history/set-theory/limits.tex
\begingroup
\def\olfilename{limits}
\expandafter\def\csname ol@sectioncs\endcsname{section}

\olfileid{his}{set}{limits}	

\olsection{حدّاں دی دقیق تعریف}

اج کل پچھلے مسئلے دا معیاری حل ایہ اے کہ لامتناہی طور اُتے چھوٹیاں قدراں دی لوڑ مکائی جاوے۔ ایہ کم ایویں ہوندا اے۔

اسیں ویکھیا اے کہ جیویں $\beta$ چھوٹا ہوندا اے، ڈھلوان دا اندازہ ہور نیڑے ہوندا اے۔ اصل وچ، جیویں $\beta$، $0$ دے جِنّا مرزی نیڑے جاندا اے، $f'(c)$ دی قدر اوس ڈھلوان ول «بناں روک ودھدی جاندی اے» جیہڑی سانوں چاہیدی اے۔ سو ایہ ویکھن دی تھاں کہ \emph{عین} $\beta = 0$ اُتے کیہ ہوندا اے، سانوں بس $f'(c)$ دا \emph{رُخ} ویکھنا چاہیدا اے، جدوں $\beta$، $0$ ول جاندا اے۔

ایس طرح آکھیے تاں عام مسئلہ ایس شکل دے دعویاں نوں سمجھنا اے:
\begin{center}
جیویں $x$، $c$ دے نیڑے جاندا اے، $g(x)$ بناں روک $\ell$ ول جاندا اے۔
\end{center}
ایہ گل ہور مختصر ایویں لکھ سکدے آں:
\[
\lim_{x \rightarrow c}g(x) = \ell.
\]
انیسویں صدی وچ وائرشٹراس نے کوشی دے پچھلے کم دی بنیاد اُتے ایس عبارت دی پوری طرح دقیق تعریف دتی۔ خیال ایہ اے کہ اسیں $g(x)$ نوں~$\ell$ دے جِنّا چاہو اُنا نیڑے لیا سکدے آں، جے $x$ نوں~$c$ دے مناسب نیڑے لے آئیے۔ ہور ٹھیک ٹھیک آکھیے تاں $\lim_{x \rightarrow c} g(x) = \ell$ دا مطلب ایہ اے:
\[
(\forall\epsilon > 0)(\exists \delta > 0)\forall x \left(|x - c| < \delta \lif |g(x) - \ell| < \epsilon \right).
\]
ایتھے کھڑیاں لکیراں مطلق قدر دسدیاں نیں۔ یعنی $|x| = x$
جدوں $x \geq 0$ ہووے، تے $|x| =-x$ جدوں $x < 0$ ہووے۔ \emph{اوس} فنکشن دا گراف ایویں بنا سکدے او:
\begin{center}
	\begin{tikzpicture}[scale=1]
	\draw[->, gray] (-2.5,0) -- (2.5,0) node[right] {$x$};
	\draw[->, gray] (0,-.2) -- (0,2.5) node[above] {$|x|$};
	
	\foreach \x/\xtext in {-2/-2, -1,1, 1/1, 2/2}
	\draw[shift={(\x,0)}] (0pt,2pt) -- (0pt,-2pt) node[below] {{$\xtext$}};
	
	\foreach \y/\ytext in {1/1, 2/2}
	\draw[shift={(0,\y)}] (2pt,0pt) -- (-2pt,0pt) node[left] {{$\ytext$}};
	
	\draw[thick] (-2.5, 2.5)--(0,0)--(2.5, 2.5);
	\end{tikzpicture}
\end{center}
سو تعریف دا موٹا جیہا مطلب ایہ اے کہ تسی اپنی «غلطی» $\epsilon$ توں گھٹ (یعنی $|g(x) - \ell| < \epsilon$) کر سکدے او۔ ایس لئی ایہو جیہیاں مدخلی قدراں چُنو جیہڑیاں $\delta$ توں ودھ دُور~$c$ توں نہ ہون (یعنی $|x -
c| < \delta$)۔

حد دی تعریف ہو جان پچھوں، اسیں لامتناہی طور اُتے چھوٹیاں قدراں نوں پورا ای پاسے رکھ سکدے آں، تے ایہ طے کر سکدے آں کہ $f$ دی $c$ اُتے ڈھلوان ایویں ملدی اے:
\[
	{f}'(c) = \lim_{x \rightarrow 0}\left(\frac{f(c +x) - f(c)}{x}\right) \text{ جتھے حد موجود ہووے}.
\]
ایہ سمجھنا ضروری اے کہ ساڈی تعریف وچ «جتھے حد موجود ہووے» والی شرط کیوں چاہیدی اے۔ اک سوکھی مثال لئی $f(x) = |x|$ لو، جیہدا گراف اسیں ہُنے ویکھیا اے۔ صاف اے کہ $f'(0)$ دی تعریف نہیں بندی۔ جے «سجّے پاسے توں» $0$ ول آئیے، تاں ڈھلوان ہمیشہ $1$ اے۔ جے «کھبّے پاسے توں» $0$ ول آئیے، تاں ڈھلوان ہمیشہ $-1$ اے۔ سو حد دی تعریف نہیں بندی۔ ایس حساب نال اسیں ایہ وی آکھ سکدے آں کہ فنکشن~$f$،~$x$ اُتے اُدوں تے صرف اُدوں \emph{تفریق جوگ} اے جدوں ایہو جیہی حد موجود ہووے۔

اسیں ویکھیا اے کہ \emph{حد} دے خیال نال تفریق کِویں سمجھائی جاندی اے۔ ایہو خیال \emph{مسلسل} فنکشن دی تعریف کرن لئی وی ورت سکدے آں۔ (بولزانو نوں ایہ گل 1817 تک، عملی طور اُتے، سمجھ آ چکی سی۔) کوشی--وائرشٹراس دے طریقے وچ تسلسل نوں ایویں سمجھیا جاندا اے۔ موٹے طور اُتے فنکشن~$f$ (کسے نقطے اُتے) مسلسل اے جے تسی اوہدے حاصل وچ جِنّی درستی منگو، مدخل اُتے درستی دی ڈھکدی شرط لا کے اُنی درستی پکی کیتی جا سکے۔ ہور ٹھیک ٹھیک آکھیے تاں $f$، $c$ اُتے مسلسل اے جے جیویں $x$ صفر ول جاندا اے، $f(c + x)$ تے $f(c)$ دا فرق آپ $0$ ول جاوے۔ دوجے لفظاں وچ $f$، $c$ اُتے اُدوں تے صرف اُدوں \emph{مسلسل} اے جدوں $f(c) = \lim_{x \rightarrow c} f(x)$۔

ہور اگے جائیے تاں گل حقیقی تجزیے ول ٹُر پئے گی، حالانکہ ساڈا موضوع سیٹ نظریہ اے۔ ایس لئی ہُن ایتھے رک کے ایس گل دا نچوڑ دَس دیئے۔ انیسویں صدی وچ ریاضی داناں نے دقیق تعریف والی \emph{حد} ورت کے لامتناہی طور اُتے چھوٹیاں قدراں توں بناں کم کرنا سِکھ لیا۔
\endgroup
% /OLP-0635
% OLP-0636 content/history/set-theory/pathologies.tex
\begingroup
\def\olfilename{pathologies}
\expandafter\def\csname ol@sectioncs\endcsname{section}

\olfileid{his}{set}{pathology}
\olsection{اجیب اُساریاں}

پر پتا لگیا کہ \emph{حد} دی تعریف نال کجھ خاصیاں «اجیب» اُساریاں وی ممکن ہو جاندیاں نیں۔

انیسویں صدی دی چوتھی دہائی دے نیڑے بولزانو نے اک فنکشن لبھیا جیہڑا \emph{ہر تھاں مسلسل} سی پر \emph{کسے تھاں وی تفریق جوگ نہیں} سی۔ (بدقسمتی نال بولزانو نے ایہ گل کدے چھپوائی نہیں۔ ریاضی داناں نوں ایہ خیال پہلی واری 1872 وچ پتا لگیا، جدوں وائرشٹراس نے ایہو گل وکھرے طور اُتے لبھی۔)\footnote{ایہ تاریخ وکی پیڈیا دے \href{http://en.wikipedia.org/wiki/Weierstrass_function}{وائرشٹراس فنکشن} والے مضمون دیاں بہت تفصیلی حاشیاں وچ دَسی گئی اے۔} گھٹ توں گھٹ اینا تاں آکھیا جا سکدا اے کہ ایہ بڑی حیرانی والی گل سی۔ $|x|$ ورگے فنکشن لبھنے سوکھے نیں جیہڑے ہر تھاں مسلسل نیں پر کسے خاص نقطے اُتے تفریق جوگ نہیں۔ پر جیہڑا فنکشن ہر تھاں مسلسل ہووے تے \emph{کسے تھاں وی} تفریق جوگ نہ ہووے، اوہ بالکل وکھری گل اے۔ ذرا سوچو ایہو جیہا فنکشن کِویں کھچو گے۔ تسلسل لئی لازم اے کہ کاغذ توں قلم چُکے بغیر سارا گراف کھچیا جا سکے۔ پر کسے تھاں وی تفریق جوگ نہ ہون لئی قلم دی سمت لگاتار تے اچانک بدلنی پئے گی۔

ہور «اجیب» مثالاں وی سامنے آئیاں۔ 5 جنوری 1874 نوں کانتور نے ڈیڈیکائنڈ نوں خط وچ ایہ سوال پُچھیا:
\begin{quote}
کی کسے سطح (مثلاً چُفیرے دی حد سمیت اک مربع) نوں کسے لکیر (مثلاً دوہاں سِریاں سمیت اک سدھی لکیر) نال اک بہ اک جوڑیا جا سکدا اے، ایویں کہ سطح دے ہر نقطے نال لکیر دا اک نقطہ تے، اُلٹے پاسے، لکیر دے ہر نقطے نال سطح دا اک نقطہ جُڑے؟

ایس ویلے وی مینوں لگدا اے کہ ایس سوال دا جواب بڑا اوکھا اے۔ بھاویں ایتھے وی \emph{نہ} آکھن نوں اینا چِت کردا اے کہ ثبوت لگ بھگ وادھو ای لگے۔ [حوالہ \citealt{Gouvea2011} وچ اے۔]
\end{quote}
پر 1877 وچ کانتور نے ثابت کیتا کہ اوہ غلط سی۔ لکیر تے مربع وچ نقطیاں دی گنتی عین اِکّو اے۔ اوہنے 29 جون 1877 نوں ڈیڈیکائنڈ نوں لکھیا، «\emph{je le vois, mais je ne le crois pas}»، یعنی «میں ایہ ویکھدا آں، پر مینوں یقین نہیں آؤندا۔» ریاضی دی «رائج تاریخ» وچ ایس گل توں اکثر ایہ مطلب لیا جاندا اے کہ اوس ویلے دے ریاضی داناں نوں ایہ نویں نتیجے \emph{لفظی معنیاں وچ ناقابلِ یقین} لگے سن۔ (خطاں دی گل بات \citet{Gouvea2011} وچ ملدی اے، تے اسیں \olref[his][set][mythology]{sec} وچ اوس ول مُڑاں گے۔ کانتور دے ثبوت دا خاکہ \olref[his][set][cantorplane]{sec} وچ اے۔)

کانتور دے نتیجے توں تحریک لے کے پیانو ایہ سوچن لگ پیا کہ کی لکیر نوں کسے مستوی اُتے \emph{ہمواری نال} بھیجیا جا سکدا اے۔ ایہ اک \emph{ساری تھاں بھرن والی منحنی} ہووے گی۔ \citeyear{Peano1890} وچ پیانو نے ایہو جیہی منحنی بنائی۔ ایہ سچ مُچ عام سہج سمجھ دے اُلٹ اے۔ اقلیدس نے لکیر نوں «چوڑائی توں بناں لمبائی» آکھیا سی (پہلی کتاب، دوجی تعریف)، پر پیانو نے وکھا دِتا کہ لکیر نوں ڈھکدی طرح وَل دے کے اوہدی لمبائی توں چوڑائی بنائی جا سکدی اے۔ \citeyear{Hilbert1891} وچ ہلبرٹ نے سمجھ وچ ذرا سوکھی آؤن والی ساری تھاں بھرن والی منحنی دَسی تے اوہنوں وکھاؤن لئی کجھ تصویراں وی دتیاں۔ منحنی اک ترتیب نال اُسردی اے۔ اُساری دے پہلے چھے پڑا ایہ نیں:
\begin{center}
\begin{tikzpicture}
\begin{scope}[xshift=20pt, yshift=20pt]
	\draw [oldiagcolorC, l-system={Hilbert curve, step=40pt, angle=90, axiom=L, order=1}] lindenmayer system;
\end{scope}
\begin{scope}[xshift=110pt, yshift=10pt]
	\draw [oldiagcolorC, l-system={Hilbert curve, step=20pt, angle=90, axiom=L, order=2}] lindenmayer system;
\end{scope}
\begin{scope}[xshift=205pt, yshift=5pt]
	\draw [oldiagcolorC, l-system={Hilbert curve, step=10pt, angle=90, axiom=L, order=3}] lindenmayer system;
\end{scope}
\begin{scope}[xshift=2.5pt, yshift=-97.5pt]
	\draw [oldiagcolorC, l-system={Hilbert curve, step=5pt, angle=90, axiom=L, order=4}] lindenmayer system;
\end{scope}
\begin{scope}[xshift=101.25pt, yshift=-98.75pt]
	\draw [oldiagcolorC, l-system={Hilbert curve, step=2.5pt, angle=90, axiom=L, order=5}] lindenmayer system;
\end{scope}
\begin{scope}[xshift=200.625pt, yshift=-99.375pt]
	\draw [oldiagcolorC, l-system={Hilbert curve, step=1.25pt, angle=90, axiom=L, order=6}] lindenmayer system;
\end{scope}
\draw[gray] (0pt,0pt) rectangle (80pt, 80pt);
\draw[gray] (100pt,0pt) rectangle (180pt, 80pt);
\draw[gray] (200pt,0pt) rectangle (280pt, 80pt);
\draw[gray] (0pt,-100pt) rectangle (80pt, -20pt);
\draw[gray] (100pt,-100pt) rectangle (180pt, -20pt);
\draw[gray] (200pt,-100pt) rectangle (280pt, -20pt);
\end{tikzpicture}  
\end{center}
حد وچ، جیہدی ہُن دقیق تعریف بن چکی سی، سارا مربع !!{colorC} رنگ نال پُورا بھر جاندا اے۔ نال ایہ وی دَس دیئے کہ ہلبرٹ دی منحنی ہر تھاں مسلسل اے پر کسے تھاں وی تفریق جوگ نہیں۔ عام سمجھ نال ایہ ایویں اے کہ لامتناہی حد وچ فنکشن ہر پل اچانک پاسا بدلدا اے۔ (ہلبرٹ دی اُساری دا ہور تفصیلی خاکہ اسیں \olref[his][set][hilbertcurve]{sec} وچ دیاں گے۔)

چنگا ہویا یا ماڑا، ایہناں «اجیب» ہندسی اُساریاں نوں ہندسے دی سہج سمجھ اُتے بھروسا کرن دے خلاف اک وجہ سمجھیا گیا۔ ایہ ریاضی کرن دے اک نویں طریقے دا لگ بھگ \emph{پرچار} بن گئیاں، جیہڑا اگے چل کے سیٹ نظریے تک پہنچیا۔ ایس موضوع دیاں بعد وچ گھڑیاں گئیاں متھاں وچ وار وار دُہرایا گیا کہ ایہ نتیجے پوری طرح دقیق تے بالکل حیران کر دین والے سن۔ ایس طرح ایہناں دا دوہرا کم سی۔ ہندسی سہج سمجھ اُتے بھروسا کرن توں خبردار کرنا تے وکھاؤنا کہ نویں ڈھنگ نال سوچن توں کِنّے نتیجے نکل سکدے نیں۔
\endgroup
% /OLP-0636
% OLP-0637 content/history/set-theory/mythology.tex
\begingroup
\def\olfilename{mythology}
\expandafter\def\csname ol@sectioncs\endcsname{section}

\olfileid{his}{set}{mythology}

\olsection{تاریخ توں ودھ متھ؟}

ہُن اک صدی توں وی ودھ پچھوں ایہناں واقعات ول مڑ کے ویکھدیاں، سانوں ایہ فیصلے سُن کے ای منّ نہیں لینے چاہیدے۔ نتیجے سچ مُچ نویں، دلچسپ تے حیرانی والے سن۔ پر ایہ کِنّے کُ حیران کر دین والے سن؟ تے کی ایہناں سچ مُچ وکھایا کہ سانوں ہندسی سہج سمجھ اُتے بھروسا نہیں کرنا چاہیدا؟

حیرانی والی گل بارے \citet{Gouvea2011} دسدا اے کہ ڈیڈیکائنڈ نوں کانتور دی مشہور گل، «\emph{je le vois, mais je ne le crois
pas}»، اپنے اصل سیاق توں کافی وکھ کر کے پیش کیتی جاندی اے۔ اوہ سیاق ہور کھولھ کے ایویں اے (\citeauthor{Gouvea2011} توں لیا گیا):
\begin{quote}
جے میں تہاڈی مہربانی تے صبر توں اینا کم لینا آں، تاں ایس موضوع لئی میرا شوق معاف کرنا۔ میں جیہڑیاں گلّاں ہُنے تہاڈے ول بھیجیاں، اوہ میرے اپنے لئی وی اینیاں اچانک تے اینیاں نویاں نیں کہ، میرے محترم سجن، جدوں تک تہاڈے کولوں اوہناں دے ٹھیک ہون بارے فیصلہ نہیں ملدا، مینوں چین نہیں آ سکدا۔ جدوں تک تسی میرے نال سہمت نہیں ہوندے، میں بس ایہ آکھ سکدا آں:
\emph{je le vois, mais je ne le crois pas.}
\end{quote}
کانتور جاندا سی کہ اوہدا نتیجہ «اینا اچانک، اینا نواں» سی۔ پر ایہ شک والی گل اے کہ اوہنوں اپنا نتیجہ کدے \emph{ناقابلِ یقین} لگیا ہووے۔ جیویں \citeauthor{Gouvea2011} دسدا اے، کانتور ڈیڈیکائنڈ نوں بس اپنا پیش کیتا ثبوت پرکھن لئی آکھ رِہا سی۔

ہندسی سہج سمجھ دی گل ایہ اے کہ پیانو نے اپنی ساری تھاں بھرن والی منحنی کسے خاکے توں بناں چھپوائی۔ پر ہلبرٹ نے اپنی منحنی چھپواؤن ویلے اپنا مقصد سمجھایا۔ جے قاری «اگّے دِتّی ہندسی سہج سمجھ دی مدد لین»، تاں اوہناں کول پیانو دے نتیجے نوں سمجھن دا صاف طریقہ ہووے گا۔ فیر اوہنے \olref[his][set][pathology]{sec} وچلے ورگے \emph{خاکیاں} دی اک لڑی دِتّی۔

ہور عام گل ایہ اے کہ بھاویں چھپے ہوئے ثبوتاں وچ خاکیاں دا رواج کافی گھٹ گیا اے، ایہ حقیقت نہیں ٹل سکدی کہ ریاضی دان ثبوت دیندیاں \emph{اکثر} خاکے ورتدے نیں۔ (موٹے طور اُتے آکھیے، چنگے ریاضی دان جاندے نیں کہ ہندسی سہج سمجھ اُتے کدوں بھروسا کیتا جا سکدا اے۔)

مُکدی گل ایہ اے کہ رولا سُن کے ای یقین نہ کرو، یا گھٹ توں گھٹ ایہنوں بناں پرکھے نہ منّو۔ ہور پڑھنا ہووے تاں \citet{Giaquinto2007} ویکھو۔
\endgroup
% /OLP-0637
% OLP-0638 content/history/set-theory/cantor-plane.tex
\begingroup
\def\olfilename{cantor-plane}
\expandafter\def\csname ol@sectioncs\endcsname{section}

\olfileid{his}{set}{cantorplane}
\olsection{لکیر تے مستوی بارے کانتور}

شروڈر-برنسٹین دے ثبوت نال جڑے کجھ حالات دا تعلق اوس تاریخ نال اے جیہڑی اسیں \olref[his][set][pathology]{sec} وچ دَسی سی۔ چیتے کرو کہ 1877 وچ کانتور نے ثابت کیتا کہ مربع وچ عین اوہنے ای نقطے نیں جِنّے اوہدی اک بھجّا اُتے نیں۔ ایتھے اسیں اوہدے ثبوت دی پہلی کوشش پیش کراں گے۔

منّ لو $\unitline$ اکائی لکیر اے، یعنی نقطیاں دا سیٹ $[0,1]$۔ منّ لو $\unitsquare$ اکائی مربع اے، یعنی نقطیاں دا سیٹ $\unitline
\times \unitline$۔ ایہناں اصطلاحاں وچ کانتور نے ثابت کیتا کہ $\cardeq{\unitline}{\unitsquare}$۔ اوہنے ڈیڈیکائنڈ نوں اک سنیہا لکھیا جیہدے وچ بنیادی طور اُتے ایہ دلیل سی۔

\begin{thm}\ollabel{thm:cantorplane}
$\cardeq{\unitline}{\unitsquare}$
\end{thm}

\begin{proof}[ثبوت دا پہلا حصہ]
$a, b \in \unitline$ چُن لو۔ اوہناں نوں دوئی لکھت وچ لکھو۔ ایس طرح $0$ تے $1$ دیاں دو لامتناہی لَڑیاں، $a_1$، $a_2$، \dots، تے $b_1$، $b_2$، \dots، ملدیاں نیں، جیہڑیاں ایویں نیں:
\begin{align*}
a &= 0.a_1a_2a_3a_4\dots\\
b &= 0.b_1b_2b_3b_4\dots
\intertext{ہُن ایویں دِتّا فنکشن $f \colon \unitsquare \to \unitline$ لو:}
f(a, b) & = 0.a_1b_1a_2b_2a_3b_3a_4b_4\dots
\end{align*}
ہُن $f$ اک !!a{injection} اے، کیوں جے جے $f(a, b) = f(c,d)$، تاں $a_n =
c_n$ تے $b_n = d_n$ سارے $n \in \Nat$ لئی، سو $a = c$ تے $b =
d$۔
\end{proof}

بدقسمتی نال، جیویں ڈیڈیکائنڈ نے کانتور نوں دسیا، ایس نال اصل سوال دا جواب نہیں ملدا۔ $0.\dot{1}\dot{0} =
0.1010101010\ldots$ لو۔ سانوں چاہیدا اے کہ $f(a,b) = 0.\dot{1}\dot{0}$ ہووے، جتھے:
\begin{align*}
a&= 0.\dot{1}\dot{1} = 0.111111\ldots\\
b&= 0
\end{align*}
پر $a = 0.\dot{1}\dot{1} = 1$ اے۔ سو جدوں اسیں آکھدے آں کہ «$a$ تے $b$ نوں دوئی لکھت وچ لکھو»، تاں سانوں چُننا پیندا اے کہ \emph{کیہڑی} لکھت ورتنی اے۔ تے چونکہ $f$ اک \emph{فنکشن} اے، اسیں دوہاں ممکن لکھتاں وچوں بس \emph{اک} ای ورت سکدے آں۔ پر جے، مثال لئی، اسیں سوکھی لکھت ورت کے $a$ نوں «$1.000\ldots$» لکھیے، تاں کوئی ایہو جیہی جوڑی $\tuple{a, b}$ نہیں لبھدی جیہدے لئی $f(a, b) = 0.\dot{1}\dot{0}$ ہووے۔

مُکدی گل ایہ اے کہ ڈیڈیکائنڈ نے دسیا کہ کجھ دُہراندے اعشاری پھیلاؤ ممکن ہون کرکے کانتور دا فنکشن $f$ !!a{injection} اے پر !!a{surjection} \emph{نہیں}۔ سو کانتور نے بس $\cardle{\unitsquare}{\unitline}$ وکھایا اے تے $\cardeq{\unitsquare}{\unitline}$ \emph{نہیں} وکھایا۔

کانتور نے لگ بھگ جھٹ ڈیڈیکائنڈ نوں جواب لکھیا تے بنیادی طور اُتے ایہ دَسیا کہ ثبوت ایویں پورا ہو سکدا اے۔

\begin{proof}[پورا ثبوت]
سو اسیں وکھایا اے کہ $\cardle{\unitsquare}{\unitline}$۔ پر صاف اے کہ اک !!a{injection} اے جیہڑا $\unitline$ توں $\unitsquare$ ول جاندا اے۔ بس لکیر نوں مربع دی اک بھجّا دے نال سدھا رکھ دیو۔ سو $\cardle{\unitline}{\unitsquare}$ تے $\cardle{\unitsquare}{\unitline}$۔ شروڈر--برنسٹین دے قضیے نال
(\olref[sfr][siz][sb]{thm:schroder-bernstein})،
$\cardeq{\unitline}{\unitsquare}$۔
\end{proof}

پر ظاہر اے کہ کانتور آخری سطر ایس طرح پوری نہیں سی کر سکدا، کیوں جے شروڈر-برنسٹین دا قضیہ ہجے ثابت نہیں سی ہویا۔ بھاویں کانتور نے بعد وچ ایہنوں عام قیاس دی شکل دِتّی، ایہدا تسلّی بخش ثبوت 1897 تک نہیں ملیا۔ (ایس لئی کانتور نے 1877 وچ ای بعد وچ \olref{thm:cantorplane} دا اک ہور ثبوت دِتّا، جیہڑا شروڈر--برنسٹین اُتے نہیں سی ٹکدا۔)
\endgroup
% /OLP-0638
% OLP-0639 content/history/set-theory/hilbert-curve.tex
\begingroup
\def\olfilename{hilbert-curve}
\expandafter\def\csname ol@sectioncs\endcsname{section}

\olfileid{his}{set}{hilbertcurve}

\olsection{ضمیمہ۔ ہلبرٹ دیاں ساری تھاں بھرن والیاں منحنیواں}

باب \olref[pathology]{sec} وچ اسیں دسیا سی کہ کانتور دے ایس ثبوت توں پچھوں کہ اک لکیر تے اک مربع وچ نقطیاں دی گنتی عین اِکّو اے (\olref[his][set][cantorplane]{thm:cantorplane})، پیانو نے پُچھیا کہ کی کوئی ایہو جیہی \emph{منحنی} ہو سکدی اے جیہڑی ساری تھاں بھر دے۔ 1890 وچ اوہنوں ہاں وچ جواب ملیا۔ ایس حصے وچ اسیں سمجھاؤنے آں کہ ہلبرٹ دی ساری تھاں بھرن والی منحنی کِویں بنائی جاندی اے۔ بیان موٹا جیہا اے تے اُساری وچ اک ذرا جیہی تبدیلی اے۔\footnote{ہور دقیق بیان لئی \cite{Rose2010} ویکھو۔ تبدیلی بس تھلّے دتیاں منحنیواں دے لال حصے شامل کرنا اے۔ ایس نال ایہ پرکھنا ذرا سوکھا ہو جاندا اے کہ منحنی مسلسل اے۔}

سانوں فنکشن $h$ نوں فنکشاں دی لڑی $h_1$، $h_2$، $h_3$، \dots\@ دی حد وجوں تعریف کرنا اے۔ اسیں پہلاں اُساری دساں گے۔ فیر وکھاواں گے کہ ایہ ساری تھاں بھر دی اے۔ فیر وکھاواں گے کہ ایہ اک منحنی اے۔

اسیں $h$ دا حاصل سیٹ اکائی مربع $\unitsquare$ منّاں گے۔ $h$ دا پہلا اندازہ، یعنی $h_1$، ایہ اے:
\begin{center}
	\begin{tikzpicture}
	\draw[gray] (0pt, 0pt) rectangle (80pt, 80pt);
	\draw[step = 40pt, gray, very thin] (0pt, 0pt) grid (80pt, 80pt);
	\draw[oldiagcolorC, thick] (20pt, 0)--(20pt, 20pt);
	\draw[oldiagcolorC, thick] (60pt, 0)--(60pt, 20pt);		
	\begin{scope}[xshift=20pt, yshift=20pt]
	\draw [black, thick, l-system={Hilbert curve, step=40pt, angle=90, axiom=L, order=1}]  lindenmayer system;
	\end{scope}
	\end{tikzpicture}
\end{center}
ہر حصے دا پتا رکھن لئی اسیں مربع اُتے $2 \times 2$ دی جالی پائی اے۔ اسیں سوچ سکدے آں کہ منحنی تھلّے کھبّے والے چوتھائی حصے توں شروع ہوندی اے، فیر اُتّے کھبّے، فیر اُتّے سجّے تے آخر تھلّے سجّے جاندی اے۔ اُساری دا دوجا پڑا، یعنی $h_2$، ایہ اے:
\begin{center}
	\begin{tikzpicture}
	\draw[gray] (0pt, 0pt) rectangle (80pt, 80pt);
	\draw[step = 20pt, gray, very thin] (0pt, 0pt) grid (80pt, 80pt);
	\draw[oldiagcolorC, thick] (0pt, 10pt)--(10pt, 10pt);
	\draw[oldiagcolorC, thick] (70pt, 10pt)--(80pt, 10pt);
	\draw[thick, oldiagcolorE] (10pt, 30pt)--(10pt, 50pt); 
	\draw[thick, oldiagcolorE] (30pt, 50pt)--(50pt, 50pt); 
	\draw[thick, oldiagcolorE] (70pt, 50pt)--(70pt, 30pt); \begin{scope}[xshift=10pt, yshift=30pt]
	\draw [thick, rotate=270,black, l-system={Hilbert curve, step=20pt, angle=90, axiom=L, order=1}]  lindenmayer system;
	\end{scope}
	\begin{scope}[xshift=10pt, yshift=50pt]
	\draw [thick, black, l-system={Hilbert curve, step=20pt, angle=90, axiom=L, order=1}]  lindenmayer system;
	\end{scope}
	\begin{scope}[xshift=50pt, yshift=50pt]
	\draw [thick, black, l-system={Hilbert curve, step=20pt, angle=90, axiom=L, order=1}]  lindenmayer system;
	\end{scope}
	\begin{scope}[xshift=70pt, yshift=10pt]
	\draw [thick, rotate=90,black, l-system={Hilbert curve, step=20pt, angle=90, axiom=L, order=1}]  lindenmayer system;
	\end{scope}
	\end{tikzpicture}
\end{center}
وکھ وکھ رنگاں توں ایہ سمجھن وچ مدد ملدی اے کہ $h_2$ کِویں بنیا۔ اسیں پہلاں $h_1$ دے اوہ حصے دیاں چھوٹیاں کیتیاں نکلاں، جیہڑا !!{colorC} رنگ دا نہیں، مربع دے تھلّے کھبّے، اُتّے کھبّے، اُتّے سجّے تے تھلّے سجّے رکھدے آں (ایہ کالے رنگ وچ نیں)۔ فیر ایہناں چار شکلاں نوں جوڑدے آں (!!{colorE} لکیراں نال)۔ آخر اپنی شکل نوں مربع دی باہری حد نال جوڑدے آں (!!{colorC} لکیراں نال)۔

ہُن $h_3$ ول آؤ۔ جیویں $h_2$، $h_1$ دے لال توں چھُٹ باقی حصے دیاں چار چھوٹیاں کیتیاں تے جوڑیاں نکلاں توں بنیا، اویویں $h_3$، $h_2$ دے لال توں چھُٹ باقی حصے دیاں چار چھوٹیاں کیتیاں نکلاں توں بندا اے (کالے رنگ وچ)۔ اوہناں نوں اک دوجے نال جوڑدے آں (!!{colorE} لکیراں نال) تے آخر مربع دی باہری حد نال جوڑدے آں (!!{colorC} لکیراں نال)۔
\begin{center}
	\begin{tikzpicture}
	\draw[gray] (0pt, 0pt) rectangle (80pt, 80pt);
	\draw[step = 10pt, gray, very thin] (0pt, 0pt) grid (80pt, 80pt);
	\draw[thick, oldiagcolorC](5pt, 0pt)--(5pt, 5pt); 
	\draw[thick, oldiagcolorC] (75pt, 0pt)--(75pt, 5pt); 
	\draw[thick, oldiagcolorE] (75pt, 45pt)--(75pt, 35pt);
	\draw[thick, oldiagcolorE] (5pt, 35pt)--(5pt, 45pt); 
	\draw[thick, oldiagcolorE] (35pt, 45pt)--(45pt, 45pt); 
	\draw[thick, oldiagcolorE] (75pt, 45pt)--(75pt, 35pt); \begin{scope}[xshift=5pt, yshift=35pt]
	\draw [rotate=270, thick, black, l-system={Hilbert curve, step=10pt, angle=90, axiom=L, order=2}]  lindenmayer system;
	\end{scope}
	\begin{scope}[xshift=5pt, yshift=45pt]
	\draw [thick, black, l-system={Hilbert curve, step=10pt, angle=90, axiom=L, order=2}]  lindenmayer system;
	\end{scope}
	\begin{scope}[xshift=45pt, yshift=45pt]
	\draw [thick, black, l-system={Hilbert curve, step=10pt, angle=90, axiom=L, order=2}]  lindenmayer system;
	\end{scope}
	\begin{scope}[xshift=75pt, yshift=5pt]
	\draw [thick, rotate=90,black, l-system={Hilbert curve, step=10pt, angle=90, axiom=L, order=2}]  lindenmayer system;
	\end{scope}
	\end{tikzpicture}
\end{center}
ہُن سانوں عام ڈھنگ دِسدا اے، جیہدے نال $h_{n+1}$ دی تعریف $h_n$ توں ہوندی اے۔ آخر اسیں \emph{آپ} منحنی $h$ دی تعریف کردے آں۔ ایہدے لئی وارو واری آؤن والے فنکشاں $h_1$، $h_2$، \dots\@ دی ہر نقطے اُتے حد لینے آں۔ یعنی ہر $x \in \unitsquare$ لئی:
\begin{align*}
	h(x) &= \lim_{n \rightarrow \infty} h_n(x)
\end{align*}
ہُن اسیں وکھاؤنے آں کہ ایہ منحنی ساری تھاں بھر دی اے۔ جدوں اسیں منحنی $h_n$ کھچدے آں، اسیں $2^n \times 2^n$ دی جالی $\unitsquare$ اُتے پاؤنے آں۔ فیثاغورث دے قضیے مطابق جالی دے ہر خانے دے قطر دی لمبائی ایہ اے:
\[
\sqrt{\left(\nicefrac{1}{2^{n}}\right)^2+\left(\nicefrac{1}{2^{n}}\right)^2} = 2^{(\frac{1}{2}-n)}
\]
تے صاف اے کہ $h_n$ جالی دے ہر خانے وچوں لنگھدا اے۔ سو $\unitsquare$ دا ہر نقطہ $2^{(\frac{1}{2}-n)}$ دی دُوری دے \emph{اندر} اے، $h_n$ دے کسے نہ کسے نقطے توں۔ ہُن $h$ دی تعریف فنکشاں دی لڑی $h_1$، $h_2$، $h_3$، \dots\@ دی حد وجوں کیتی گئی اے۔ سو کسے وی نقطے دی $h$ توں ودھ توں ودھ دُوری ایہ اے:
\[
\lim_{n \rightarrow \infty} 2^{(\frac{1}{2}-n)} = 0.
\]
یعنی $\unitsquare$ دا ہر نقطہ $0$ دُور اے~$h$ توں۔ دوجے لفظاں وچ، $\unitsquare$ دا ہر نقطہ منحنی \emph{اُتے} اے۔ سو $h$ ساری تھاں بھر دی اے!{}

ہُن ایہ وکھاؤنا باقی اے کہ $h$ سچ مُچ اک \emph{منحنی} اے۔ ایس لئی سانوں منحنی دی تعریف دینی پئے گی۔ اج دی تعریف اوس تعریف اُتے ٹِکی اے جیہڑی ژوردان نے 1887 وچ دِتّی سی (یعنی پہلی ساری تھاں بھرن والی منحنی توں بس کجھ سال پہلاں)۔

\begin{defn}
منحنی $\unitline$ توں $\Real^2$ ول اک مسلسل نقشہ بندی اے۔
\end{defn}

ایہ گل عام سہج سمجھ نال کافی صاف اے۔ ایس سمجھ مطابق منحنی اک «ہموار» نقشہ بندی اے جیہڑی اک معیاری لکیر نوں مستوی $\Real^2$ اُتے لیاندی اے۔ ساڈا فنکشن $h$ سچ مُچ $\unitline$ توں $\Real^2$ ول اک نقشہ بندی اے۔ سو سانوں بس وکھاؤنا اے کہ $h$ مسلسل اے۔ اسیں \olref[limits]{sec} وچ $\epsilon$/$\delta$ والی لکھت نال تسلسل دی تعریف دِتّی سی۔ عام بولی وچ اسیں ایہ ثابت کرنا چاہندے آں: \emph{جے تسی اک نقطہ $p$، $\unitsquare$ وچ، تے اپنی مرضی دی درستی $\epsilon$ دسو، تاں اسیں $\unitline$ دا اک کھُلّا حصہ لبھ سکدے آں، ایویں کہ اوس کھُلّے حصے دے کسے وی $x$ لئی $h(x)$ دی دُوری $\epsilon$ دے اندر ہووے، $p$ توں۔}

سو فرض کرو کہ تسی $p$ تے $\epsilon$ دس دِتّے نیں۔ ایہ اصل وچ $p$ مرکز تے $\epsilon$ نیم قطر والا دائرہ $\unitsquare$ اُتے کھچن دے برابر اے۔ (دائرہ $\unitsquare$ دی حد توں باہر وی جا سکدا اے، پر ایہدے نال کوئی فرق نہیں پیندا۔) ہُن چیتے کرو کہ فنکشن $h_n$ دا بیان کردیاں اسیں $2^n \times 2^n$ دی جالی $\unitsquare$ اُتے کھچّی سی۔ صاف اے کہ $\epsilon$ جِنّا وی چھوٹا ہووے، کوئی $n$ ایہو جیہا اے کہ $2^n \times 2^n$ دی جالی دا اک خانہ، جیہڑی $\unitsquare$ اُتے اے، $p$ مرکز تے $\epsilon$ نیم قطر والے دائرے دے اندر پورے دا پورا آ جاندا اے۔

سو اوہ $n$ لو تے $I$، $\unitline$ دا سبھ توں وڈا کھُلّا حصہ ہووے، جیہنوں $h_n$ پورے دا پورا جالی دے اوس خانے اندر بھیجدا اے۔ (صاف اے کہ $(a,b)$ موجود اے، کیوں جے اسیں پہلاں ای ویکھیا کہ $h_n$، $2^n\times 2^n$ دی جالی دے ہر خانے وچوں لنگھدا اے۔) ہُن اینا وکھاؤنا کافی اے کہ جدوں وی $x \in I$ ہووے، نقطہ $h(x)$ جالی دے اوسے خانے وچ اے۔ تے \emph{ایہ} وکھاؤن لئی اینا کافی اے کہ $h_m(x)$ اوسے خانے وچ اے، ہر $m > n$ لئی۔ پر ایہ صاف اے۔ جے اسیں ویکھیے کہ $h_m$ نال کیہ ہوندا اے، $m > n$ لئی، تاں دِسدا اے کہ اکائی وقفے دا عین «اوہی حصہ» اوسے خانے وچ بھیجیا جاندا اے۔ اسیں بس اوہنوں اوس خانے وچ ہور ودھ کھچ کے تے ہور ودھ ول دے کے بھیجنے آں۔
\endgroup
% /OLP-0639
\endgroup
% /OLP-0633

\OLEndPartHook
\endgroup
% /OLP-0620

% OLP-0640 content/reference/reference.tex
\begingroup
\def\olfilename{reference}
\expandafter\def\csname ol@sectioncs\endcsname{section}

\olpart{ref}{حوالہ جاتی مواد}

\begin{editorial}
ایس حصے وچ حرفاں، نشانیاں، تعریفاں تے قاعدیاں دیاں وکھ وکھ فہرستاں نیں، جیہڑیاں درسی کتاب وچ ضمیمیاں وجوں شامل کیتیاں جا سکدیاں نیں۔
\end{editorial}

% OLP-0641 content/reference/greek-alphabet/greek-alphabet.tex
\begingroup
\def\olfilename{greek-alphabet}
\expandafter\def\csname ol@sectioncs\endcsname{section}

\olchapter{ref}{grk}{یونانی ورن مالا}

\begin{center}
    \begin{tabular}{lll@{\qquad}lll}
    الفا &	$\alpha$ & $A$ &
    نیو & $\nu$ & $N$ \\
    بیٹا & $\beta$ & $B$ &
    کسائی & $\xi$ & $\Xi$ \\
    گاما &	$\gamma$ & $\Gamma$ &
    اومیکرون & $o$ & $O$ \\
    ڈیلٹا & $\delta$ & $\Delta$ &
    پائی & $\pi$ & $\Pi$ \\
    ایپسیلون & $\epsilon$ & $E$ &
    رو & $\rho$ & $P$ \\
    زیٹا & $\zeta$ & $Z$ &
    سگما &	$\sigma$ & $\Sigma$ \\
    ایٹا & $\eta$ & $H$ &
    ٹاؤ & $\tau$ & $T$ \\
    تھیٹا & $\theta$ & $\Theta$ &
    اُپسیلون & $\upsilon$ & $\Upsilon$ \\
    آیوٹا & $\iota$ & $I$ &
    فائی & $\varphi$ & $\Phi$ \\
    کاپا & $\kappa$ & $K$ &
    کائی & $\chi$ & $X$ \\
    لیمبڈا & $\lambda$ & $\Lambda$ &
    پسائی & $\psi$ & $\Psi$ \\
    میو & $\mu$ & $M$ &
    اومیگا & $\omega$ & $\Omega$
    \end{tabular}
\end{center}
\endgroup
% /OLP-0641

% OLP-0642 content/reference/fraktur-alphabet/fraktur-alphabet.tex
\begingroup
\def\olfilename{fraktur-alphabet}
\expandafter\def\csname ol@sectioncs\endcsname{section}

\olchapter{ref}{frk}{فراکتور ورن مالا}

\begin{center}
    \begin{tabular}{llll@{\qquad}llll}
    A &	$\mathfrak{A}$ & a & $\mathfrak{a}$ &
    N &	$\mathfrak{N}$ & n & $\mathfrak{n}$ \\
    B &	$\mathfrak{B}$ & b & $\mathfrak{b}$ &
    O &	$\mathfrak{O}$ & o & $\mathfrak{o}$ \\
    C &	$\mathfrak{C}$ & c & $\mathfrak{c}$ &
    P &	$\mathfrak{P}$ & p & $\mathfrak{p}$ \\
    D &	$\mathfrak{D}$ & d & $\mathfrak{d}$ &
    Q &	$\mathfrak{Q}$ & q & $\mathfrak{q}$ \\
    E &	$\mathfrak{E}$ & e & $\mathfrak{e}$ &
    R &	$\mathfrak{R}$ & r & $\mathfrak{r}$ \\
    F &	$\mathfrak{F}$ & f & $\mathfrak{f}$ &
    S &	$\mathfrak{S}$ & s & $\mathfrak{s}$ \\
    G &	$\mathfrak{G}$ & g & $\mathfrak{g}$ &
    T &	$\mathfrak{T}$ & t & $\mathfrak{t}$ \\
    H &	$\mathfrak{H}$ & h & $\mathfrak{h}$ &
    U &	$\mathfrak{U}$ & u & $\mathfrak{u}$ \\
    I &	$\mathfrak{I}$ & i & $\mathfrak{i}$ &
    V &	$\mathfrak{V}$ & v & $\mathfrak{v}$ \\
    J &	$\mathfrak{J}$ & j & $\mathfrak{j}$ &
    W &	$\mathfrak{W}$ & w & $\mathfrak{w}$ \\
    K &	$\mathfrak{K}$ & k & $\mathfrak{k}$ &
    X &	$\mathfrak{X}$ & x & $\mathfrak{x}$ \\
    L &	$\mathfrak{L}$ & l & $\mathfrak{l}$ &
    Y &	$\mathfrak{Y}$ & y & $\mathfrak{y}$ \\
    M &	$\mathfrak{M}$ & m & $\mathfrak{m}$ &
    Z &	$\mathfrak{Z}$ & z & $\mathfrak{z}$ 
    \end{tabular}
    \end{center}
\endgroup
% /OLP-0642

\OLEndPartHook
\endgroup
% /OLP-0640
% /OLP-0002
\backmatter
\bibliographystyle{\olpath/bib/natbib-oup}
\begin{otherlanguage}{english}
\renewcommand{\bibname}{\foreignlanguage{urdu}{حوالیاں دی فہرست}}
\nocite{Frege1953,Peter1967}
\bibliography{\olpath/bib/open-logic}
\end{otherlanguage}
\end{document}
